\IfFileExists{stacks-project-book.cls}{%
\documentclass{stacks-project-book}
}{%
\documentclass{amsbook}
}

\usepackage{verbatim}
\newenvironment{reference}{\comment}{\endcomment}
\newenvironment{slogan}{\comment}{\endcomment}
\newenvironment{history}{\comment}{\endcomment}

\usepackage[all]{xy}

\xyoption{2cell}
\UseAllTwocells

\usepackage{multicol}

\usepackage{lmodern}
\usepackage[T1]{fontenc}
\usepackage[french]{babel}
\AtBeginDocument{\shorthandoff{?}}


\usepackage{hyperref}


\theoremstyle{plain}
\newtheorem{theorem}[subsection]{Théorème}
\newtheorem{proposition}[subsection]{Proposition}
\newtheorem{lemma}[subsection]{Lemme}

\theoremstyle{definition}
\newtheorem{definition}[subsection]{Définition}
\newtheorem{example}[subsection]{Exemple}
\newtheorem{exercise}[subsection]{Exercice}
\newtheorem{situation}[subsection]{Situation}

\theoremstyle{remark}
\newtheorem{remark}[subsection]{Remarque}
\newtheorem{remarks}[subsection]{Remarques}

\renewcommand{\proofname}{Démonstration}

\numberwithin{equation}{subsection}

\def\lim{\mathop{\mathrm{lim}}\nolimits}
\def\colim{\mathop{\mathrm{colim}}\nolimits}
\def\Spec{\mathop{\mathrm{Spec}}}
\def\Hom{\mathop{\mathrm{Hom}}\nolimits}
\def\Ext{\mathop{\mathrm{Ext}}\nolimits}
\def\SheafHom{\mathop{\mathcal{H}\!\mathit{om}}\nolimits}
\def\SheafExt{\mathop{\mathcal{E}\!\mathit{xt}}\nolimits}
\def\Sch{\mathit{Sch}}
\def\Mor{\mathop{\mathrm{Mor}}\nolimits}
\def\Ob{\mathop{\mathrm{Ob}}\nolimits}
\def\Sh{\mathop{\mathit{Sh}}\nolimits}
\def\NL{\mathop{N\!L}\nolimits}
\def\CH{\mathop{\mathrm{CH}}\nolimits}
\def\proetale{{pro\text{-}\acute{e}tale}}
\def\etale{{\acute{e}tale}}
\def\QCoh{\mathit{QCoh}}
\def\Ker{\mathop{\mathrm{Ker}}}
\def\Im{\mathop{\mathrm{Im}}}
\def\Coker{\mathop{\mathrm{Coker}}}
\def\Coim{\mathop{\mathrm{Coim}}}

\DeclareMathSymbol{\boxtimes}{\mathbin}{AMSa}{"02}

\def\QCohstack{\mathcal{QC}\!\mathit{oh}}
\def\Cohstack{\mathcal{C}\!\mathit{oh}}
\def\Spacesstack{\mathcal{S}\!\mathit{paces}}
\def\Quotfunctor{\mathrm{Quot}}
\def\Hilbfunctor{\mathrm{Hilb}}
\def\Curvesstack{\mathcal{C}\!\mathit{urves}}
\def\Polarizedstack{\mathcal{P}\!\mathit{olarized}}
\def\Complexesstack{\mathcal{C}\!\mathit{omplexes}}
\def\Pic{\mathop{\mathrm{Pic}}\nolimits}
\def\Picardstack{\mathcal{P}\!\mathit{ic}}
\def\Picardfunctor{\mathrm{Pic}}
\def\Deformationcategory{\mathcal{D}\!\mathit{ef}}
\begin{document}
\title{Le projet Stacks --- édition française, partie 4}
\maketitle
\tableofcontents

% OK, start here.
%

\chapter{Schémas}



\phantomsection
\label{schemes-section-phantom}


\section{Introduction}
\label{schemes-section-introduction}

\noindent
Dans ce document, nous définissons les schémas.
Une référence fondamentale est \cite{EGA}.









\section{Espaces annelés en anneaux locaux}
\label{schemes-section-locally-ringed-spaces}

\noindent
Rappelons que nous avons défini les espaces annelés
dans Faisceaux, section \ref{sheaves-section-ringed-spaces}.
En bref, un espace annelé est un couple $(X, \mathcal{O}_X)$ constitué
d'un espace topologique $X$ et d'un faisceau d'anneaux $\mathcal{O}_X$.
Un morphisme d'espaces annelés $f : (X, \mathcal{O}_X) \to (Y, \mathcal{O}_Y)$
est donné par une application continue $f : X \to Y$ et un $f$-morphisme de faisceaux
d'anneaux $f^\sharp : \mathcal{O}_Y \to \mathcal{O}_X$. On peut
considérer $f^\sharp$ comme un morphisme $\mathcal{O}_Y \to f_*\mathcal{O}_X$; voir
Faisceaux, définition \ref{sheaves-definition-f-map} et le
lemme \ref{sheaves-lemma-f-map}.

\medskip\noindent
Un bon exemple géométrique à garder à l'esprit est celui des
variétés $\mathcal{C}^\infty$ et des morphismes de
variétés $\mathcal{C}^\infty$. En effet, si $M$ est une
variété $\mathcal{C}^\infty$, alors le faisceau $\mathcal{C}^\infty_M$
des fonctions lisses est un faisceau d'anneaux sur $M$. Et toute
application $f : M \to N$ entre variétés est lisse si et seulement si
pour toute section locale $h$ de $\mathcal{C}^\infty_N$
la composée $h \circ f$ est une section locale de $\mathcal{C}^\infty_M$.
Ainsi, une application lisse $f$ définit naturellement un morphisme
d'espaces annelés
$$
f : (M , \mathcal{C}^\infty_M) \longrightarrow (N, \mathcal{C}^\infty_N)
$$
voir Faisceaux, exemple \ref{sheaves-example-continuous-map-ringed}.
Il est instructif d'examiner ce qui arrive aux fibres. Soit donc
$m \in M$ d'image $f(m) = n \in N$. Rappelons que la fibre
$\mathcal{C}^\infty_{M, m}$ est l'anneau des germes de fonctions lisses
en $m$; voir
Faisceaux, exemple \ref{sheaves-example-germs-functions}.
L'anneau des germes de fonctions sur $(M, m)$ est un anneau local dont l'idéal
maximal est formé des fonctions qui s'annulent en $m$.
Il en va de même pour $\mathcal{C}^\infty_{N, n}$. L'application sur les fibres
$f^\sharp : \mathcal{C}^\infty_{N, n} \to \mathcal{C}^\infty_{M, m}$
envoie l'idéal maximal dans l'idéal maximal, tout simplement
parce que $f(m) = n$.

\medskip\noindent
En géométrie algébrique, nous étudions les schémas. Sur un schéma, le faisceau d'anneaux
n'est pas déterminé par une propriété intrinsèque de l'espace.
Le spectre d'un anneau $R$
(voir Algèbre, section \ref{algebra-section-spectrum-ring}) muni
d'un faisceau d'anneaux construit à partir de $R$ (voir ci-dessous),
sera notre brique de base. Il apparaîtra que
les fibres de $\mathcal{O}$ sur $\Spec(R)$ sont les anneaux locaux
de $R$ en ses idéaux premiers. Il y a deux raisons d'introduire
les espaces annelés en anneaux locaux dans ce contexte : (1) Il n'existe en général
aucun mécanisme associant à une application continue entre spectres un
homomorphisme entre les anneaux correspondants. C'est pourquoi nous ajoutons comme donnée
supplémentaire le morphisme $f^\sharp$. (2) Si nous considérons les morphismes
de ces spectres dans la catégorie des espaces annelés, alors les
homomorphismes sur les fibres peuvent ne pas être locaux. Puisque notre
intuition géométrique exige qu'ils le soient, nous introduisons les espaces annelés en anneaux locaux
comme suit.

\begin{definition}
\label{schemes-definition-locally-ringed-space}
Espaces annelés en anneaux locaux.
\begin{enumerate}
\item Un {\it espace annelé en anneaux locaux $(X, \mathcal{O}_X)$}
est un couple formé d'un
espace topologique $X$ et d'un faisceau d'anneaux $\mathcal{O}_X$ dont toutes les fibres
sont des anneaux locaux.
\item Étant donné un espace annelé en anneaux locaux $(X, \mathcal{O}_X)$, nous disons que
$\mathcal{O}_{X, x}$ est {\it l'anneau local de $X$ en $x$}.
Nous notons $\mathfrak{m}_{X, x}$ ou simplement $\mathfrak{m}_x$
l'idéal maximal de $\mathcal{O}_{X, x}$. En outre, le
{\it corps résiduel de $X$ en $x$} est le corps résiduel
$\kappa(x) = \mathcal{O}_{X, x}/\mathfrak{m}_x$.
\item Un {\it morphisme d'espaces annelés en anneaux locaux}
$(f, f^\sharp) : (X, \mathcal{O}_X) \to (Y, \mathcal{O}_Y)$
est un morphisme d'espaces annelés tel que, pour tout $x \in X$,
l'homomorphisme d'anneaux induit $\mathcal{O}_{Y, f(x)} \to \mathcal{O}_{X, x}$ soit un
homomorphisme local.
\end{enumerate}
\end{definition}

\noindent
Nous omettrons habituellement le faisceau d'anneaux $\mathcal{O}_X$
de la notation lorsque nous parlerons d'espaces annelés en anneaux locaux. Nous dirons
simplement « l'espace annelé en anneaux locaux $X$ ».
Par abus de notation, nous identifierons aussi $X$ à
son espace topologique sous-jacent. Enfin, nous appellerons
le faisceau d'anneaux correspondant
$\mathcal{O}_X$ le {\it faisceau structural de $X$}.
De plus, il est d'usage de noter l'idéal maximal
de l'anneau local $\mathcal{O}_{X, x}$ par
$\mathfrak{m}_{X, x}$ ou simplement $\mathfrak{m}_x$.
Nous dirons « soit $f : X \to Y$ un morphisme d'espaces annelés en anneaux
locaux », en omettant ainsi encore davantage les faisceaux structuraux.
Dans ce cas, par abus de notation, nous identifierons aussi $f : X\to Y$
à l'application continue sous-jacente entre espaces topologiques.
Le $f$-morphisme correspondant à $f$ sera habituellement
noté $f^\sharp$. La condition que $f$ soit un morphisme
d'espaces annelés en anneaux locaux s'exprime alors en disant que,
pour tout $x\in X$, l'homomorphisme sur les fibres
$$
f^\sharp_x : \mathcal{O}_{Y, f(x)} \longrightarrow \mathcal{O}_{X, x}
$$
envoie l'idéal maximal $\mathfrak m_{Y, f(x)}$ dans
$\mathfrak m_{X, x}$.

\medskip\noindent
Employons ces conventions de notation pour montrer que la
collection des espaces annelés en anneaux locaux et de leurs morphismes
forme une catégorie. Pour le voir, il nous faut
montrer que la composée de morphismes d'espaces annelés en anneaux locaux
est un morphisme d'espaces annelés en anneaux locaux. Soient donc
$f : X \to Y$ et $g : Y \to Z$ des morphismes d'espaces annelés en anneaux
locaux. La composée de $f$ et de $g$ est définie dans
Faisceaux, définition \ref{sheaves-definition-composition-maps-ringed-spaces}.
Soit $x \in X$. D'après
Faisceaux, lemme \ref{sheaves-lemma-compose-f-maps-stalks},
la composée
$$
\mathcal{O}_{Z, g(f(x))}
\xrightarrow{g^\sharp}
\mathcal{O}_{Y, f(x)}
\xrightarrow{f^\sharp}
\mathcal{O}_{X, x}
$$
est l'homomorphisme associé sur les fibres du morphisme $g \circ f$.
Le résultat découle de ce que la composée d'homomorphismes locaux
est un homomorphisme local.

\medskip\noindent
Un trait agréable de la définition est que le foncteur
$$
\textit{Espaces annelés en anneaux locaux}
\longrightarrow
\textit{Espaces annelés}
$$
reflète les isomorphismes (et même davantage).
Voici un énoncé moins abstrait.

\begin{lemma}
\label{schemes-lemma-isomorphism-locally-ringed}
\begin{slogan}
Un isomorphisme d'espaces annelés entre espaces annelés en anneaux locaux est un
isomorphisme d'espaces annelés en anneaux locaux.
\end{slogan}
Soient $X$, $Y$ des espaces annelés en anneaux locaux.
Si $f : X \to Y$ est un isomorphisme d'espaces annelés,
alors $f$ est un isomorphisme
d'espaces annelés en anneaux locaux.
\end{lemma}

\begin{proof}
Cela découle immédiatement du fait algébrique correspondant :
soient $A$, $B$ des anneaux locaux. Tout isomorphisme d'anneaux
$A \to B$ est un isomorphisme d'anneaux locaux.
\end{proof}













\section{Immersions ouvertes d'espaces annelés en anneaux locaux}
\label{schemes-section-open-immersion}

\begin{definition}
\label{schemes-definition-immersion-locally-ringed-spaces}
Soit $f : X \to Y$ un morphisme d'espaces annelés en anneaux locaux.
Nous disons que $f$ est une {\it immersion ouverte} si
$f$ est un homéomorphisme de $X$ sur un sous-ensemble ouvert
de $Y$, et que l'application $f^{-1}\mathcal{O}_Y \to \mathcal{O}_X$
est un isomorphisme.
\end{definition}

\noindent
La construction suivante est parallèle à
Faisceaux, définition \ref{sheaves-definition-restriction} (3).

\begin{example}
\label{schemes-example-open-subspace}
Soit $X$ un espace annelé en anneaux
locaux. Soit $U \subset X$ un ouvert.
Soit $\mathcal{O}_U = \mathcal{O}_X|_U$
la restriction de $\mathcal{O}_X$ à $U$. Pour $u
\in U$, la fibre $\mathcal{O}_{U, u}$ est égale
à la fibre $\mathcal{O}_{X, u}$, et est donc un anneau local.
Ainsi $(U, \mathcal{O}_U)$ est un espace annelé en anneaux
locaux et le morphisme $j : (U, \mathcal{O}_U) \to (X, \mathcal{O}_X)$ est
une immersion ouverte.
\end{example}

\begin{definition}
\label{schemes-definition-open-subspace}
Soit $X$ un espace annelé en anneaux
locaux. Soit $U \subset X$ un ouvert.
L'espace annelé en anneaux locaux $(U, \mathcal{O}_U)$
de l'exemple \ref{schemes-example-open-subspace} ci-dessus
est le sous-espace ouvert {\it de $X$ associé à $U$}.
\end{definition}

\begin{lemma}
\label{schemes-lemma-open-immersion}
Soit $f : X \to Y$ une immersion ouverte d'espaces
annelés en anneaux locaux. Soit $j : V = f(X)
\to Y$ le sous-espace ouvert de $Y$ associé à l'image de $f$. Il
existe un isomorphisme unique $f' : X \cong V$ d'espaces
annelés en anneaux locaux tel que $f = j \circ f'$.
\end{lemma}

\begin{proof}
Soit $f'$ l'homéomorphisme entre $X$ et $V$ induit par $f$. Alors $f = j \circ
f'$ en tant qu'applications d'espaces topologiques. Comme il
existe un isomorphisme de
faisceaux $f^\sharp : f^{-1}(\mathcal{O}_Y) \to \mathcal{O}_X$, il existe un isomorphisme
d'anneaux
$f^\sharp : \Gamma(U, f^{-1}(\mathcal{O}_Y)) \to \Gamma(U, \mathcal{O}_X)$
pour chaque ouvert $U \subset X$. Puisque
$\mathcal{O}_V = j^{-1}\mathcal{O}_Y$ et $f^{-1} = f'^{-1} j^{-1}$ (Faisceaux,
lemme \ref{sheaves-lemma-pullback-composition}) nous voyons que
$f^{-1}\mathcal{O}_Y = f'^{-1}\mathcal{O}_V$, donc $\Gamma(U,
f'^{-1}(\mathcal{O}_V)) \to \Gamma(U, f^{-1}(\mathcal{O}_Y))$ est un isomorphisme
pour chaque $U
\subset X$ ouvert. En composant ces derniers, nous obtenons un isomorphisme d'anneaux
$$
\Gamma(U, f'^{-1}(\mathcal{O}_V)) \to \Gamma(U, \mathcal{O}_X)
$$
pour chaque ouvert $U \subset X$, et donc un isomorphisme de faisceaux $f^{-1}(\mathcal{O}_V)
\to \mathcal{O}_X$. En d'autres termes, nous avons un isomorphisme
$f'^{\sharp} : f'^{-1}(\mathcal{O}_V) \to \mathcal{O}_X$ et donc un isomorphisme
d'espaces annelés en anneaux locaux $(f', f'^{\sharp})
: (X, \mathcal{O}_X) \to (V, \mathcal{O}_V)$ (utilisez le lemme
\ref{schemes-lemma-isomorphism-locally-ringed}). Notez que $f =
j \circ f'$ en tant que morphismes d'espaces annelés en anneaux locaux
par construction.

\medskip\noindent
Supposons que nous ayons un autre
morphisme $f'' : (X, \mathcal{O}_X) \to (V, \mathcal{O}_V)$ tel que $f = j \circ f''$. En
tout point $x \in X$, nous avons $j(f'(x)) = j(f''(x))$ d'où il découle que $f'(x) =
f''(x)$ puisque $j$ est l'application d'inclusion ; par conséquent $f'$ et $f''$ sont les
mêmes en tant que morphismes d'espaces topologiques. Sur les faisceaux structuraux,
pour chaque ouvert $U \subset X$ nous avons un diagramme commutatif
$$
\xymatrix @R=5em{
\Gamma(U, f^{-1}(\mathcal{O}_Y)) \ar[d]_\cong\ar[r]^\cong &
\Gamma(U, \mathcal{O}_X) \\
\Gamma(U, f'^{-1}(\mathcal{O}_V)) \ar@/^/[ru]^{f'^\sharp}
\ar@/_/[ru]_{f''^\sharp} &
}
$$
d'où nous voyons que $f'^\sharp$ et $f''^\sharp$ définissent
le même morphisme de faisceaux.
\end{proof}

\noindent
Dorénavant, nous ne distinguons plus entre les ouverts
et leurs sous-espaces associés.

\begin{lemma}
\label{schemes-lemma-restrict-map-to-opens}
Soit $f : X \to Y$ un morphisme d'espaces annelés en anneaux
locaux. Soient $U \subset X$ et $V \subset Y$ des ouverts.
Supposons que $f(U) \subset V$. Il existe un unique morphisme
d'espaces annelés en anneaux locaux $f|_U : U \to V$ tel que
le diagramme suivant soit un carré commutatif d'espaces
annelés en anneaux locaux
$$
\xymatrix{
U \ar[d]_{f|_U} \ar[r] & X \ar[d]^f \\ V
\ar[r] & Y
}
$$
\end{lemma}

\begin{proof}
Omis.
\end{proof}

\noindent
Dans ce qui suit, nous utiliserons sans autre mention le fait suivant
qui découle du lemme ci-dessus. Étant donné un morphisme $f :
Y \to X$ d'espaces annelés en anneaux locaux, et tout ouvert
$U \subset X$ tel que $f(Y) \subset U$, alors il existe un unique
morphisme d'espaces annelés en anneaux locaux $Y \to U$ tel
que la composée $Y \to U \to X$ soit égale à $f$. En fait, nous
écrirons même par abus de langage $f : Y \to U$ puisque ceci donne
rarement lieu à confusion.









\section{Immersions fermées d'espaces annelés en anneaux locaux}
\label{schemes-section-closed-immersion}

\noindent
Nous suivons nos conventions introduites
dans Modules, définition \ref{modules-definition-closed-immersion}.

\begin{definition}
\label{schemes-definition-closed-immersion-locally-ringed-spaces}
Soit $i : Z \to X$ un morphisme d'espaces annelés en anneaux locaux.
Nous disons que $i$ est une {\it immersion fermée} si :
\begin{enumerate}
\item L'application $i$ est un homéomorphisme de $Z$ sur un sous-ensemble fermé
de $X$. \item L'application $\mathcal{O}_X \to i_*\mathcal{O}_Z$ est
surjective ; notons $\mathcal{I}$ le noyau.
\item Le $\mathcal{O}_X$-module $\mathcal{I}$ est
localement engendré par des sections.
\end{enumerate}
\end{definition}

\begin{lemma}
\label{schemes-lemma-closed-local-target}
Soit $f : Z \to X$ un morphisme d'espaces annelés en anneaux
locaux. Pour que $f$ soit une immersion fermée, il suffit
qu'il existe un recouvrement ouvert $X = \bigcup U_i$ tel que
chaque $f : f^{-1}U_i \to U_i$ soit une immersion fermée.
\end{lemma}

\begin{proof}
Omis.
\end{proof}

\begin{example}
\label{schemes-example-closed-subspace}
Soit $X$ un espace annelé en anneaux
locaux. Soit $\mathcal{I} \subset \mathcal{O}_X$ un faisceau
d'idéaux qui est localement engendré par des sections en tant
que faisceau de $\mathcal{O}_X$-modules. Soit $Z$ le support
du faisceau d'anneaux $\mathcal{O}_X/\mathcal{I}$. Il
s'agit d'un sous-ensemble fermé de
$X$, d'après Modules, lemme \ref{modules-lemma-support-sheaf-rings-closed}.
On note $i : Z \to X$ l'application
d'inclusion. D'après Modules, lemme \ref{modules-lemma-i-star-exact},
il existe un unique faisceau d'anneaux $\mathcal{O}_Z$
sur $Z$ avec $i_*\mathcal{O}_Z = \mathcal{O}_X/\mathcal{I}$.
Pour tout $z \in Z$, la fibre $\mathcal{O}_{Z,
z}$ est égale à un quotient $\mathcal{O}_{X, i(z)}/\mathcal{I}_{i(z)}$ d'un
anneau local et non nul, donc un anneau local. Ainsi,
$i : (Z, \mathcal{O}_Z) \to (X, \mathcal{O}_X)$ est une
immersion fermée d'espaces annelés en anneaux locaux.
\end{example}

\begin{definition}
\label{schemes-definition-closed-subspace}
Soit $X$ un espace annelé en anneaux
locaux. Soit $\mathcal{I}$ un faisceau d'idéaux sur
$X$ qui est localement engendré par des
sections. L'espace annelé en anneaux locaux $(Z,
\mathcal{O}_Z)$ de l'exemple \ref{schemes-example-closed-subspace}
ci-dessus est le sous-espace fermé {\it de $X$ associé
au faisceau d'idéaux $\mathcal{I}$}.
\end{definition}

\begin{lemma}
\label{schemes-lemma-closed-immersion}
Soit $f : X \to Y$ une immersion fermée d'espaces
annelés en anneaux locaux. Soit $\mathcal{I}$
le noyau de l'application $\mathcal{O}_Y \to f_*\mathcal{O}_X$.
Soit $i : Z \to Y$ le sous-espace fermé de $Y$
associé à $\mathcal{I}$. Il
existe un isomorphisme unique $f' : X \cong Z$ d'espaces
annelés en anneaux locaux tel que $f = i \circ f'$.
\end{lemma}

\begin{proof}
Omis.
\end{proof}

\begin{lemma}
\label{schemes-lemma-characterize-closed-subspace}
Soient $X$, $Y$ des espaces annelés en anneaux
locaux. Soit $\mathcal{I} \subset \mathcal{O}_X$ un faisceau d'idéaux localement
engendré par des sections. Soit $i : Z \to X$ le sous-espace fermé associé.
Un morphisme $f : Y \to X$ se factorise par $Z$ si et seulement si
l'application $f^*\mathcal{I} \to f^*\mathcal{O}_X = \mathcal{O}_Y$
est nulle. Si tel est le cas, le morphisme $g : Y \to
Z$ tel que $f = i \circ g$ est unique.
\end{lemma}

\begin{proof}
De toute évidence, si $f$ se factorise en $Y \to Z \to X$
alors l'application $f^*\mathcal{I} \to \mathcal{O}_Y$ est nulle.
Inversement, supposons que $f^*\mathcal{I} \to \mathcal{O}_Y$
soit nulle. Choisissez un $y \in Y$ et considérez
l'application d'anneaux $f^\sharp_y : \mathcal{O}_{X, f(y)} \to
\mathcal{O}_{Y, y}$.
Puisque la composée $\mathcal{I}_{f(y)} \to \mathcal{O}_{X, f(y)} \to
\mathcal{O}_{Y, y}$ est nulle par hypothèse et puisque
$f^\sharp_y(1) = 1$, nous voyons que $1 \not \in \mathcal{I}_{f(y)}$,
c'est-à-dire $\mathcal{I}_{f(y)} \not = \mathcal{O}_{X, f(y)}$.
Nous en concluons que $f(Y) \subset Z = \text{Supp}(\mathcal{O}_X/\mathcal{I})$.
Ainsi, $f = i \circ g$ où $g : Y \to Z$ est
continue. Considérons l'application $f^\sharp : \mathcal{O}_X \to
f_*\mathcal{O}_Y$. L'hypothèse que $f^*\mathcal{I} \to \mathcal{O}_Y$ est nulle implique
que la composée $\mathcal{I} \to \mathcal{O}_X \to f_*\mathcal{O}_Y$ est
nulle par adjonction de $f_*$ et $f^*$. En
d'autres termes, nous obtenons un morphisme de faisceaux
d'anneaux $\overline{f^\sharp} : \mathcal{O}_X/\mathcal{I} \to f_*\mathcal{O}_Y$.
Notez que $f_*\mathcal{O}_Y = i_*g_*\mathcal{O}_Y$ et que
$\mathcal{O}_X/\mathcal{I} = i_*\mathcal{O}_Z$. D'après
Faisceaux, lemme \ref{sheaves-lemma-equivalence-categories-closed-structures}, nous obtenons
un morphisme unique de faisceaux d'anneaux $g^\sharp
: \mathcal{O}_Z \to g_*\mathcal{O}_Y$ dont l'image directe sous $i$
est $\overline{f^\sharp}$. Nous omettons la vérification que $(g, g^\sharp)$
définit un morphisme d'espaces annelés en anneaux locaux et
que $f = i \circ g$ en tant que morphisme d'espaces annelés en anneaux
locaux. L'unicité de $(g, g^\sharp)$ a été signalée ci-dessus.
\end{proof}

\begin{lemma}
\label{schemes-lemma-restrict-map-to-closed}
Soit $f : X \to Y$ un morphisme d'espaces annelés en anneaux
locaux. Soit $\mathcal{I} \subset \mathcal{O}_Y$ un faisceau
d'idéaux qui est localement engendré par des
sections. Soit $i : Z \to Y$ le sous-espace fermé associé au
faisceau d'idéaux $\mathcal{I}$. Soit
$\mathcal{J}$ l'image de l'application
$f^*\mathcal{I} \to f^*\mathcal{O}_Y = \mathcal{O}_X$. Alors
cet idéal est localement engendré par des sections. De
plus, soit $i' : Z' \to X$ le sous-espace fermé associé
de $X$. Il existe un unique morphisme
d'espaces annelés en anneaux locaux $f' : Z' \to Z$ tel que le
diagramme suivant soit un carré commutatif d'espaces
annelés en anneaux locaux
$$
\xymatrix{
Z' \ar[d]_{f'} \ar[r]_{i'} & X \ar[d]^f \\ Z
\ar[r]^{i} & Y
}
$$
De plus, ce diagramme est un carré cartésien dans la catégorie des espaces
annelés en anneaux locaux.
\end{lemma}

\begin{proof}
L'idéal $\mathcal{J}$ est localement engendré par des sections
d'après Modules, lemme \ref{modules-lemma-pullback-locally-generated}.
Le reste du lemme découle de la caractérisation, dans le lemme
\ref{schemes-lemma-characterize-closed-subspace} ci-dessus, de ce
que signifie qu'un morphisme se factorise à travers un sous-espace
fermé.
\end{proof}














\section{Schémas affines}
\label{schemes-section-affine-schemes}

\noindent
Soit $R$ un anneau. Considérons l'espace topologique $\Spec(R)$
associé à $R$, voir Algèbre, section \ref{algebra-section-spectrum-ring}.
Nous munirons cet espace d'un faisceau d'anneaux $\mathcal{O}_{\Spec(R)}$
et la paire résultante $(\Spec(R), \mathcal{O}_{\Spec(R)})$
sera un schéma affine.

\medskip\noindent
Rappelons que $\Spec(R)$ possède une base d'ouverts $D(f)$,
$f \in R$ que nous appelons ouverts standards, voir
Algèbre, définition \ref{algebra-definition-Zariski-topology}.
De plus, l'intersection de deux ouverts standards est un autre
ouvert standard : $D(f) \cap D(g) = D(fg)$, $f, g\in R$.

\begin{lemma}
\label{schemes-lemma-standard-open}
Soit $R$ un anneau. Soit $f \in R$.
\begin{enumerate}
\item Si $g\in R$ et $D(g) \subset D(f)$, alors
\begin{enumerate}
\item $f$ est inversible dans $R_g$,
\item $g^e = af$ pour certains $e \geq 1$ et $a \in R$, \item
il existe un morphisme d'anneaux canonique $R_f \to R_g$, et
\item il existe un morphisme de $R_f$-modules canonique
$M_f \to M_g$ pour tout $R$-module $M$.
\end{enumerate}
\item Tout recouvrement ouvert de $D(f)$ peut être affiné en un
recouvrement ouvert fini de la forme $D(f) = \bigcup_{i = 1}^n D(g_i)$.
\item Si $g_1, \ldots, g_n \in R$, alors $D(f) \subset \bigcup D(g_i)$
si et seulement si $g_1, \ldots, g_n$ engendrent l'idéal unité dans $R_f$.
\end{enumerate}
\end{lemma}

\begin{proof}
Rappelons que $D(g) = \Spec(R_g)$ (voir
Algèbre, lemme \ref{algebra-lemma-standard-open}). Ainsi
(a) est valable parce que l'image
de $f$ dans $R_g$ n'appartient à aucun
idéal premier, et est donc inversible, voir Algèbre, lemme
\ref{algebra-lemma-Zariski-topology}. Écrivons l'inverse de $f$
dans $R_g$ comme $a/g^d$. Cela signifie que $g^d -
af$ est annihilé par une puissance de $g$, d'où (b). Pour (c), l'application
$R_f \to R_g$ existe par (a) à cause de la propriété universelle de la localisation,
ou nous pouvons la définir en envoyant $b/f^n$ sur
$a^nb/g^{ne}$. L'égalité $M_f = M \otimes_R R_f$ peut être utilisée
pour obtenir l'application sur les modules, ou nous
pouvons définir $M_f \to M_g$ en envoyant
$x/f^n$ sur $a^nx/g^{ne}$.

\medskip\noindent
Rappelons que $D(f)$ est quasi-compact, voir
Algèbre, lemme \ref{algebra-lemma-qc-open}.
Ainsi, la deuxième affirmation découle directement
du fait que les ouverts standards forment
une base pour la topologie.

\medskip\noindent
La troisième affirmation découle directement
d'Algèbre, lemme \ref{algebra-lemma-Zariski-topology}.
\end{proof}

\noindent
Dans Faisceaux, section \ref{sheaves-section-bases}, nous avons
défini la notion de faisceau sur une base, et nous avons montré qu'elle
est essentiellement équivalente à la notion de faisceau sur l'espace,
voir Faisceaux, lemmes \ref{sheaves-lemma-extend-off-basis} et
\ref{sheaves-lemma-extend-off-basis-structures}. De plus, nous avons
montré dans
Faisceaux, lemme \ref{sheaves-lemma-cofinal-systems-coverings-standard-case} qu'il
suffit de vérifier la condition de faisceau sur
un système cofinal de recouvrements ouverts pour chaque ouvert
standard. Par le lemme ci-dessus, il suffit de vérifier sur les
recouvrements finis par des ouverts standards.

\begin{definition}
\label{schemes-definition-standard-covering}
Soit $R$ un anneau.
\begin{enumerate}
\item Un {\it recouvrement ouvert standard} de $\Spec(R)$
est un recouvrement $\Spec(R) = \bigcup_{i = 1}^n D(f_i)$,
où $f_1, \ldots, f_n \in R$.
\item Supposons que $D(f) \subset \Spec(R)$ soit un ouvert
standard. Un {\it recouvrement ouvert standard} de $D(f)$
est un recouvrement $D(f) = \bigcup_{i = 1}^n
D(g_i)$, où $g_1, \ldots, g_n \in R$.
\end{enumerate}
\end{definition}

\noindent
Soit $R$ un anneau. Soit $M$ un $R$-module. Nous définirons un préfaisceau
$\widetilde M$ sur la base des ouverts standards. Supposons que
$U \subset \Spec(R)$ soit un ouvert standard. Si $f, g \in R$
sont tels que $D(f) = D(g)$, alors par le lemme
\ref{schemes-lemma-standard-open} ci-dessus, il existe des applications canoniques $M_f
\to M_g$ et $M_g \to M_f$ qui sont mutuellement inverses. Par conséquent,
nous pouvons choisir n'importe quel $f$ tel que $U = D(f)$
et définir
$$
\widetilde M(U) = M_f.
$$
Notez que si $D(g) \subset D(f)$, alors par le lemme
\ref{schemes-lemma-standard-open} ci-dessus, nous avons une
application canonique
$$
\widetilde M(D(f)) = M_f \longrightarrow M_g = \widetilde M(D(g)).
$$
Clairement, cela définit un préfaisceau de groupes abéliens sur la base
des ouverts standards. Si $M = R$, alors $\widetilde R$ est un préfaisceau
d'anneaux sur la base des ouverts standards.

\medskip\noindent
Calculons la fibre de $\widetilde M$ en un point $x \in \Spec(R)$. Supposons
que $x$ corresponde à l'idéal premier $\mathfrak p \subset R$. Par
définition de la fibre, nous voyons que
$$
\widetilde M_x = \colim_{f\in R, f\not\in \mathfrak p} M_f
$$
Ici, l'ensemble $\{f \in R, f \not \in \mathfrak p\}$ est préordonné par
la règle $f \geq f' \Leftrightarrow D(f) \subset D(f')$.
Si $f_1, f_2 \in R \setminus \mathfrak p$, alors nous avons
$f_1f_2 \geq f_1$ dans cet ordre. Ainsi, par
Algèbre, lemme \ref{algebra-lemma-localization-colimit},
nous concluons que
$$
\widetilde M_x = M_{\mathfrak p}.
$$

\medskip\noindent
Ensuite, nous vérifions la condition de faisceau pour les recouvrements ouverts
standards. Si $D(f) = \bigcup_{i = 1}^n D(g_i)$, alors la condition
de faisceau pour ce recouvrement est équivalente à l'exactitude de
la suite
$$
0 \to M_f \to \bigoplus M_{g_i} \to \bigoplus M_{g_ig_j}.
$$
Notez que $D(g_i) = D(fg_i)$, et donc nous pouvons réécrire
cette suite comme la suite
$$
0 \to M_f \to \bigoplus M_{fg_i} \to \bigoplus M_{fg_ig_j}.
$$
De plus, d'après le lemme \ref{schemes-lemma-standard-open} ci-dessus,
nous voyons que $g_1, \ldots, g_n$ engendrent l'idéal unité
dans $R_f$. Nous pouvons donc
appliquer Algèbre, lemme \ref{algebra-lemma-cover-module},
au module $M_f$ sur $R_f$ et aux éléments $g_1, \ldots, g_n$. Nous en
concluons que la suite est exacte. D'après les remarques faites
ci-dessus, nous voyons que $\widetilde M$ est un faisceau
sur la base des ouverts standards.

\medskip\noindent
Ainsi, nous concluons à partir des résultats
de Faisceaux, section \ref{sheaves-section-bases}
qu'il existe un unique
faisceau d'anneaux $\mathcal{O}_{\Spec(R)}$ qui
coïncide avec $\widetilde R$ sur les ouverts standards. Notez que
d'après notre calcul des fibres ci-dessus, toutes les
fibres de ce faisceau d'anneaux sont des anneaux locaux.

\medskip\noindent
De même, pour tout $R$-module $M$, il existe un
unique faisceau de $\mathcal{O}_{\Spec(R)}$-modules
$\mathcal{F}$ qui coïncide avec $\widetilde M$ sur les
ouverts standards, voir
Faisceaux, lemme \ref{sheaves-lemma-extend-off-basis-module}.

\begin{definition}
\label{schemes-definition-structure-sheaf}
Soit $R$ un anneau.
\begin{enumerate}
\item Le {\it faisceau structural $\mathcal{O}_{\Spec(R)}$ du spectre
de $R$} est l'unique faisceau d'anneaux $\mathcal{O}_{\Spec(R)}$ dont la
restriction à la base des ouverts standards coïncide avec $\widetilde
R$. \item L'espace annelé en
anneaux locaux $(\Spec(R), \mathcal{O}_{\Spec(R)})$
est appelé le {\it spectre} de $R$ et désigné $\Spec(R)$.
\item Le faisceau de $\mathcal{O}_{\Spec(R)}$-modules
étendant $\widetilde M$ à tous les ouverts de
$\Spec(R)$ est appelé le faisceau de $\mathcal{O}_{\Spec(R)}$-modules
associé à $M$. Ce faisceau est également noté $\widetilde
M$.
\end{enumerate}
\end{definition}

\noindent
Nous résumons les résultats obtenus jusqu'à présent.

\begin{lemma}
\label{schemes-lemma-spec-sheaves}
Soit $R$ un anneau. Soit $M$ un $R$-module. Soit $\widetilde M$
le faisceau de $\mathcal{O}_{\Spec(R)}$-modules
associé à $M$.
\begin{enumerate}
\item Nous avons $\Gamma(\Spec(R), \mathcal{O}_{\Spec(R)}) = R$. \item
Nous avons $\Gamma(\Spec(R), \widetilde M) = M$ en tant que $R$-module.
\item Pour chaque $f \in R$ nous avons
$\Gamma(D(f), \mathcal{O}_{\Spec(R)}) = R_f$.
\item Pour chaque $f\in R$ nous avons $\Gamma(D(f), \widetilde M) = M_f$ en
tant que $R_f$-module.
\item Chaque fois que $D(g) \subset D(f)$, les applications de
restriction sur $\mathcal{O}_{\Spec(R)}$ et $\widetilde
M$ sont
les applications $R_f \to R_g$ et $M_f \to M_g$
du lemme \ref{schemes-lemma-standard-open}.
\item Soit $\mathfrak p$ un idéal premier de $R$, et soit $x \in
\Spec(R)$ le point correspondant. Nous
avons $\mathcal{O}_{\Spec(R), x} = R_{\mathfrak
p}$. \item Soit $\mathfrak p$ un idéal premier de $R$, et soit $x \in \Spec(R)$
le point correspondant. Nous avons $\widetilde M_x = M_{\mathfrak p}$ en
tant que $R_{\mathfrak p}$-module.
\end{enumerate}
De plus, toutes ces identifications sont fonctorielles par rapport
au $R$-module $M$. En particulier, le foncteur $M \mapsto \widetilde
M$ est un foncteur exact de la catégorie des $R$-modules
à la catégorie des $\mathcal{O}_{\Spec(R)}$-modules.
\end{lemma}

\begin{proof}
Les assertions (1) - (7) sont claires d'après la discussion
ci-dessus. L'exactitude du foncteur $M \mapsto \widetilde
M$ découle du fait que le foncteur $M \mapsto M_{\mathfrak p}$ est exact
et du fait que l'exactitude des suites exactes courtes peut
être vérifiée sur les fibres, voir
Modules, lemme \ref{modules-lemma-abelian}.
\end{proof}

\begin{definition}
\label{schemes-definition-affine-scheme}
Un {\it schéma affine} est un espace annelé en anneaux locaux isomorphe en
tant qu'espace annelé en anneaux locaux à $\Spec(R)$ pour un certain anneau
$R$. Un {\it morphisme de schémas affines} est un morphisme dans la catégorie des
espaces annelés en anneaux locaux.
\end{definition}

\noindent
Il s'avère que les schémas affines jouent un rôle particulier
parmi tous les espaces annelés en anneaux locaux, ce dont traite
la section suivante.


















\section{La catégorie des schémas affines}
\label{schemes-section-category-affine-schemes}

\noindent
Notez que si $Y$ est un schéma affine, alors ses points
sont en bijection canonique $1-1$ (une correspondance biunivoque)
avec les idéaux premiers de $\Gamma(Y, \mathcal{O}_Y)$.

\begin{lemma}
\label{schemes-lemma-morphism-into-affine-where-point-goes}
Soit $X$ un espace annelé en anneaux locaux.
Soit $Y$ un schéma affine. Soit $f
\in \Mor(X, Y)$ un morphisme d'espaces annelés
en anneaux locaux. Étant donné un point $x \in X$, considérons
les homomorphismes d'anneaux
$$
\Gamma(Y, \mathcal{O}_Y) \xrightarrow{f^\sharp}
\Gamma(X, \mathcal{O}_X) \to \mathcal{O}_{X, x}
$$
Soit $\mathfrak p \subset \Gamma(Y, \mathcal{O}_Y)$ l'image
réciproque de $\mathfrak m_x$. Soit $y \in Y$ le point
correspondant. Alors $f(x) = y$.
\end{lemma}

\begin{proof}
Considérons le diagramme commutatif
$$
\xymatrix{
\Gamma(X, \mathcal{O}_X) \ar[r] &
\mathcal{O}_{X, x} \\
\Gamma(Y, \mathcal{O}_Y) \ar[r] \ar[u] &
\mathcal{O}_{Y, f(x)} \ar[u]
}
$$
(voir la discussion des $f$-morphismes qui
suit Faisceaux, définition \ref{sheaves-definition-f-map}).
Puisque la flèche verticale de droite est locale,
nous voyons que $\mathfrak m_{f(x)}$ est
l'image réciproque de $\mathfrak m_x$. Le résultat
s'ensuit.
\end{proof}

\begin{lemma}
\label{schemes-lemma-f-open}
Soit $X$ un espace annelé en anneaux locaux.
Soit $f \in \Gamma(X, \mathcal{O}_X)$.
L'ensemble
$$
D(f) = \{x \in X \mid \text{l'image de }f \not\in \mathfrak m_x\}
$$
est ouvert. De plus, $f|_{D(f)}$ est inversible.
\end{lemma}

\begin{proof}
Ceci est un cas particulier de Modules, lemme \ref{modules-lemma-s-open}, mais
nous donnons également une
démonstration directe. Supposons que $U \subset X$ et $V
\subset X$ soient deux sous-ensembles ouverts tels que
$f|_U$ ait un inverse $g$ et que $f|_V$ ait un inverse
$h$. Alors clairement $g|_{U\cap V} = h|_{U\cap
V}$. Il suffit donc de montrer que $f$ est inversible dans un
voisinage ouvert de tout $x \in D(f)$. Cela est
clair car $f \not \in \mathfrak m_x$ implique que $f \in \mathcal{O}_{X,
x}$ a un inverse $g \in \mathcal{O}_{X, x}$, ce qui signifie qu'il
existe un certain voisinage ouvert $x \in U \subset
X$ tel que $g \in \mathcal{O}_X(U)$ et $g\cdot f|_U = 1$.
\end{proof}

\begin{lemma}
\label{schemes-lemma-f-open-affine}
Dans le lemme \ref{schemes-lemma-f-open} ci-dessus, si $X$ est un schéma
affine, alors l'ouvert $D(f)$ correspond à l'ouvert standard
$D(f)$ défini précédemment
(en Algèbre, définition \ref{algebra-definition-spectrum-ring}).
\end{lemma}

\begin{proof}
Omis.
\end{proof}

\begin{lemma}
\label{schemes-lemma-morphism-into-affine}
\begin{reference}
Une référence pour ce fait est \cite[II, Err 1, Prop. 1.8.1]{EGA}
où il est attribué à J. Tate.
\end{reference}
Soit $X$ un espace annelé en anneaux locaux.
Soit $Y$ un schéma affine.
L'application
$$
\Mor(X, Y)
\longrightarrow
\Hom(\Gamma(Y, \mathcal{O}_Y), \Gamma(X, \mathcal{O}_X))
$$
qui envoie $f$ sur $f^\sharp$ (au niveau des sections globales) est bijective.
\end{lemma}

\begin{proof}
Étant donné que $Y$ est affine,
nous avons $(Y, \mathcal{O}_Y) \cong (\Spec(R), \mathcal{O}_{\Spec(R)})$ pour
un certain anneau $R$.
Au cours de la démonstration, nous utiliserons des faits
sur $Y$ et son faisceau structural qui sont des conséquences directes de ce
que nous savons sur le spectre d'un anneau, voir par exemple\ 
lemme \ref{schemes-lemma-spec-sheaves}.

\medskip\noindent
Motivés par les lemmes ci-dessus, nous construisons l'application inverse.
Soit $\psi_Y : \Gamma(Y, \mathcal{O}_Y) \to \Gamma(X, \mathcal{O}_X)$ un morphisme
d'anneaux. Tout d'abord, nous définissons l'application correspondante des
espaces
$$
\Psi : X \longrightarrow Y
$$
selon la règle
du lemme \ref{schemes-lemma-morphism-into-affine-where-point-goes}. En
d'autres termes, étant donné $x \in X$, nous définissons
$\Psi(x)$ comme le point de $Y$ correspondant à l'idéal
premier dans $\Gamma(Y, \mathcal{O}_Y)$ qui est
l'image réciproque de $\mathfrak m_x$ par la composée
$
\Gamma(Y, \mathcal{O}_Y) \xrightarrow{\psi_Y}
\Gamma(X, \mathcal{O}_X) \to
\mathcal{O}_{X, x}
$.

\medskip\noindent
Nous affirmons que l'application $\Psi : X \to Y$ est continue.
Les ouverts standards $D(g)$, pour $g \in \Gamma(Y, \mathcal{O}_Y)$
sont une base pour la topologie de $Y$. Ainsi, il suffit de prouver
que $\Psi^{-1}(D(g))$ est ouvert. Par construction de $\Psi$,
l'image inverse $\Psi^{-1}(D(g))$ est exactement l'ensemble
$D(\psi_Y(g)) \subset X$ qui est ouvert par le lemme \ref{schemes-lemma-f-open}. Par
conséquent, $\Psi$ est continue.

\medskip\noindent
Ensuite, nous construisons un $\Psi$-morphisme de
faisceaux de $\mathcal{O}_Y$ à $\mathcal{O}_X$.
Par Faisceaux, lemme \ref{sheaves-lemma-f-map-basis-below-structures}, il
suffit de définir des homomorphismes
d'anneaux $\psi_{D(g)} : \Gamma(D(g), \mathcal{O}_Y)
\to \Gamma(\Psi^{-1}(D(g)), \mathcal{O}_X)$
compatibles avec les applications de
restriction. Nous avons un isomorphisme
canonique $\Gamma(D(g), \mathcal{O}_Y) = \Gamma(Y, \mathcal{O}_Y)_g$,
car $Y$ est un schéma affine. Comme
$\psi_Y(g)$ est inversible sur $D(\psi_Y(g))$, nous voyons
qu'il existe une application canonique
$$
\Gamma(Y, \mathcal{O}_Y)_g
\longrightarrow
\Gamma(\Psi^{-1}(D(g)), \mathcal{O}_X)
=
\Gamma(D(\psi_Y(g)), \mathcal{O}_X)
$$
en étendant l'application $\psi_Y$
par la propriété universelle de la localisation.
Notez qu'il n'y a pas d'autre choix que de prendre l'application canonique
ici ! Après l'identification
canonique $\Gamma(D(g), \mathcal{O}_Y)
= \Gamma(Y, \mathcal{O}_Y)_g$, nous prenons cette application pour
définir $\psi_{D(g)}$. Cela est compatible avec la localisation puisque les
applications de restriction sur les schémas affines sont également
définies en termes des propriétés universelles de la localisation,
voir les lemmes \ref{schemes-lemma-spec-sheaves} et \ref{schemes-lemma-standard-open}.

\medskip\noindent
Ainsi, nous avons défini un morphisme d'espaces annelés $(\Psi,
\psi) : (X, \mathcal{O}_X) \to (Y, \mathcal{O}_Y)$ induisant $\psi_Y$
sur les sections globales. Pour montrer qu'il s'agit d'un morphisme
d'espaces annelés en anneaux locaux, nous devons démontrer que les
applications induites sur les anneaux locaux
$$
\psi_x : \mathcal{O}_{Y, \Psi(x)} \longrightarrow \mathcal{O}_{X, x}
$$
sont locales. Cela découle immédiatement du diagramme commutatif de
la démonstration du lemme \ref{schemes-lemma-morphism-into-affine-where-point-goes}
et de la définition de $\Psi$.

\medskip\noindent
Enfin, nous devons montrer que les constructions
$(\Psi, \psi) \mapsto \psi_Y$ et la construction $\psi_Y
\mapsto (\Psi, \psi)$ sont inverses l'une de l'autre. Clairement,
$\psi_Y \mapsto (\Psi, \psi) \mapsto \psi_Y$. Par conséquent,
la seule chose à prouver est que, étant donné $\psi_Y$,
il existe au plus une paire $(\Psi, \psi)$ donnant
lieu à celle-ci. L'unicité de $\Psi$ a été montrée dans
le lemme \ref{schemes-lemma-morphism-into-affine-where-point-goes}
et étant donné l'unicité de $\Psi$, l'unicité de l'application
$\psi$ a été soulignée au cours de la démonstration
ci-dessus.
\end{proof}

\begin{lemma}
\label{schemes-lemma-category-affine-schemes}
La catégorie des schémas affines est équivalente à l'opposée de la
catégorie des anneaux. L'équivalence est donnée par le foncteur qui associe à un
schéma affine les sections globales de son faisceau structural.
\end{lemma}

\begin{proof}
Cela découle maintenant clairement de la définition \ref{schemes-definition-affine-scheme}
et du lemme \ref{schemes-lemma-morphism-into-affine}.
\end{proof}

\begin{lemma}
\label{schemes-lemma-standard-open-affine}
Soit $Y$ un schéma affine. Soit
$f \in \Gamma(Y, \mathcal{O}_Y)$. Le
sous-espace ouvert $D(f)$ est un schéma affine.
\end{lemma}

\begin{proof}
Nous pouvons supposer que $Y = \Spec(R)$ et $f \in R$.
Considérons le morphisme de schémas affines
$\phi : U = \Spec(R_f) \to \Spec(R) = Y$ induit par l'homomorphisme
d'anneaux $R \to R_f$. D'après Algèbre, lemme \ref{algebra-lemma-standard-open},
nous savons que c'est un homéomorphisme sur $D(f)$.
D'autre part, l'application $\phi^{-1}\mathcal{O}_Y \to \mathcal{O}_U$ est un
isomorphisme sur les fibres, donc un isomorphisme. Ainsi, nous voyons
que $\phi$ est une immersion ouverte. Nous concluons que $D(f)$
est isomorphe à $U$ d'après le lemme \ref{schemes-lemma-open-immersion}.
\end{proof}

\begin{lemma}
\label{schemes-lemma-fibre-product-affine-schemes}
La catégorie des schémas affines a des produits finis et des produits fibrés. En d'autres
termes, elle possède des limites finies. De plus, les produits et
produits fibrés dans la catégorie des schémas affines sont les mêmes
que dans la catégorie des espaces annelés en anneaux locaux.
Formellement, nous avons (dans la catégorie des espaces annelés en anneaux locaux)
$$
\Spec(R) \times \Spec(S) =
\Spec(R \otimes_{\mathbf{Z}} S)
$$
et étant donné des homomorphismes d'anneaux $R \to A$, $R \to B$ nous avons
$$
\Spec(A) \times_{\Spec(R)} \Spec(B) =
\Spec(A
\otimes_R B).
$$
\end{lemma}

\begin{proof}
Ceci est simplement une application du lemme \ref{schemes-lemma-morphism-into-affine}. Tout
d'abord, d'après ce lemme, le schéma affine
$\Spec(\mathbf{Z})$ est l'objet final dans la catégorie des espaces
annelés en anneaux locaux. Ainsi, la première formule affichée découle
de la seconde. Pour prouver la seconde, notez que pour tout
espace annelé en anneaux locaux $X$, nous avons
\begin{eqnarray*}
\Mor(X, \Spec(A \otimes_R B)) &
= &
\Hom(A \otimes_R B, \mathcal{O}_X(X)) \\ &
= &
\Hom(A, \mathcal{O}_X(X))
\times_{\Hom(R, \mathcal{O}_X(X))}
\Hom(B, \mathcal{O}_X(X)) \\ &
= &
\Mor(X, \Spec(A))
\times_{\Mor(X, \Spec(R))}
\Mor(X, \Spec(B))
\end{eqnarray*}
Cela prouve la formule. Voir
Catégories, section \ref{categories-section-fibre-products} pour les
définitions pertinentes.
\end{proof}

\begin{lemma}
\label{schemes-lemma-disjoint-union-affines}
Soit $X$ un espace annelé en anneaux
locaux. Supposons $X = U \amalg V$ avec $U$ et $V$ ouverts et
tels que $U$, $V$ soient des schémas affines. Alors $X$ est un schéma affine.
\end{lemma}

\begin{proof}
Fixez $R = \Gamma(X, \mathcal{O}_X)$.
Notez que $R = \mathcal{O}_X(U) \times \mathcal{O}_X(V)$ par la
propriété du faisceau. Par le lemme \ref{schemes-lemma-morphism-into-affine}, il
existe un morphisme canonique d'espaces annelés en anneaux
locaux $X \to \Spec(R)$. Par Algèbre, lemme \ref{algebra-lemma-spec-product},
nous voyons que comme espace topologique
$\Spec(\mathcal{O}_X(U)) \amalg \Spec(\mathcal{O}_X(V)) =
\Spec(R)$
avec les applications provenant des homomorphismes d'anneaux
$R \to \mathcal{O}_X(U)$ et $R \to \mathcal{O}_X(V)$.
Cela signifie bien sûr que $\Spec(R)$ est également le coproduit
dans la catégorie des espaces annelés en anneaux
locaux. Par hypothèse, le morphisme $X \to \Spec(R)$ induit un isomorphisme
de $\Spec(\mathcal{O}_X(U))$ avec $U$ et de même pour $V$.
Par conséquent, $X \to \Spec(R)$ est un isomorphisme.
\end{proof}















\section{Faisceaux quasi-cohérents sur les schémas affines}
\label{schemes-section-quasi-coherent-affine}

\noindent
Rappelons que nous avons défini la notion abstraite de faisceau quasi-cohérent
dans Modules, définition \ref{modules-definition-quasi-coherent}.
Dans cette section, nous montrons que tout faisceau quasi-cohérent sur
un schéma affine $\Spec(R)$ correspond au faisceau $\widetilde M$ associé à
un $R$-module $M$.

\begin{lemma}
\label{schemes-lemma-compare-constructions}
Soit $(X, \mathcal{O}_X) = (\Spec(R), \mathcal{O}_{\Spec(R)})$
un schéma affine. Soit $M$ un $R$-module. Il existe un isomorphisme
canonique entre le faisceau $\widetilde M$ associé au $R$-module $M$
(définition \ref{schemes-definition-structure-sheaf}) et le faisceau
$\mathcal{F}_M$ associé au $R$-module $M$ (Modules,
définition \ref{modules-definition-sheaf-associated}). Cet isomorphisme
est fonctoriel en $M$. En particulier, les
faisceaux $\widetilde M$ sont quasi-cohérents. De plus, ils sont
caractérisés par la propriété universelle suivante
$$
\Hom_{\mathcal{O}_X}(\widetilde M, \mathcal{F}) =
\Hom_R(M,
\Gamma(X, \mathcal{F}))
$$
pour tout faisceau de $\mathcal{O}_X$-modules $\mathcal{F}$. Ici,
une application $\alpha : \widetilde M \to \mathcal{F}$ correspond à
son effet sur les sections globales.
\end{lemma}

\begin{proof}
Par Modules, lemme \ref{modules-lemma-construct-quasi-coherent-sheaves}, nous
avons un morphisme $\mathcal{F}_M \to \widetilde M$ correspondant à
l'application $M \to \Gamma(X, \widetilde M) = M$. Soit $x \in X$
correspondant à l'idéal premier $\mathfrak p \subset
R$. L'homomorphisme induit sur les fibres
est le morphisme $\mathcal{O}_{X, x} \otimes_R M \to
M_{\mathfrak p}$ qui est un isomorphisme
car $R_{\mathfrak p} \otimes_R M = M_{\mathfrak
p}$. Par conséquent, l'application $\mathcal{F}_M \to \widetilde M$
est un isomorphisme. La propriété universelle découle de celle
des faisceaux $\mathcal{F}_M$.
\end{proof}

\begin{lemma}
\label{schemes-lemma-widetilde-constructions}
Soit $(X, \mathcal{O}_X) = (\Spec(R), \mathcal{O}_{\Spec(R)})$ un
schéma affine. Il existe des isomorphismes canoniques
\begin{enumerate}
\item
$
\widetilde{M \otimes_R N}
\cong
\widetilde M \otimes_{\mathcal{O}_X} \widetilde N $,
voir
Modules, section \ref{modules-section-tensor-product}.
\item $
\widetilde{\text{T}^n(M)}
\cong
\text{T}^n(\widetilde
M) $,
$
\widetilde{\text{Sym}^n(M)}
\cong
\text{Sym}^n(\widetilde
M)
$, et $
\widetilde{\wedge^n(M)}
\cong
\wedge^n(\widetilde
M)
$,
voir
Modules, section \ref{modules-section-symmetric-exterior}.
\item si $M$ est un $R$-module de présentation finie, alors
$
\SheafHom_{\mathcal{O}_X}(\widetilde M, \widetilde N)
\cong
\widetilde{\Hom_R(M, N)}
$,
voir
Modules, section \ref{modules-section-internal-hom}.
\end{enumerate}
\end{lemma}

\begin{proof}[Première démonstration]
En utilisant le lemme \ref{schemes-lemma-compare-constructions}
et Modules, lemme \ref{modules-lemma-construct-quasi-coherent-sheaves}, nous
voyons que le foncteur $M \mapsto \widetilde M$ peut être vu comme
$\pi^*$ pour un morphisme $\pi$ d'espaces annelés. Le foncteur
image inverse sur les modules commute avec les constructions
tensorielles d'après Modules, lemmes \ref{modules-lemma-tensor-product-pullback}
et \ref{modules-lemma-pullback-tensor-algebra}.
Le morphisme $\pi : (X, \mathcal{O}_X) \to (\{*\}, R)$ est plat
par exemple parce que les fibres de $\mathcal{O}_X$ sont des
localisations de $R$ (lemme \ref{schemes-lemma-spec-sheaves}) et donc
plates sur $R$. Ainsi, l'image inverse par $\pi$ commute avec le
faisceau Hom interne si le premier module est de présentation finie,
d'après Modules, lemme \ref{modules-lemma-pullback-internal-hom}.
\end{proof}

\begin{proof}[Deuxième démonstration]
Démonstration de (1). D'après le lemme \ref{schemes-lemma-compare-constructions}, pour
donner une application $\widetilde{M
\otimes_R N}$ vers $\widetilde M \otimes_{\mathcal{O}_X}
\widetilde N$, nous devons donner une
application sur les sections
globales $M \otimes_R N \to \Gamma(X, \widetilde M \otimes_{\mathcal{O}_X}
\widetilde N)$ qui existe par définition du produit
tensoriel des faisceaux de modules. Pour voir que cette
application est un isomorphisme, il suffit de vérifier qu'elle est
un isomorphisme sur les fibres. Et cela découle de la
description des fibres de $\widetilde{M}$ (soit dans le lemme \ref{schemes-lemma-spec-sheaves}
ou dans Modules, lemme \ref{modules-lemma-construct-quasi-coherent-sheaves}),
du fait que le produit tensoriel commute avec la
localisation (Algèbre, lemme \ref{algebra-lemma-tensor-product-localization})
et Modules, lemme \ref{modules-lemma-stalk-tensor-product}.

\medskip\noindent
Démonstration de (2). C'est similaire à la preuve de (1), en
utilisant Algèbre, lemme \ref{algebra-lemma-tensor-algebra-localization}
et Modules, lemme \ref{modules-lemma-stalk-tensor-algebra}.

\medskip\noindent
Démonstration de (3). Puisque la construction $M \mapsto \widetilde{M}$ est fonctorielle,
il existe une application
$R$-linéaire $\Hom_R(M, N) \to \Hom_{\mathcal{O}_X}(\widetilde{M}, \widetilde{N})$.
La cible de cette application est le module des sections
globales de $\SheafHom_{\mathcal{O}_X}(\widetilde M, \widetilde
N)$. Ainsi, par le lemme \ref{schemes-lemma-compare-constructions}, nous
obtenons une application de $\mathcal{O}_X$-modules $\widetilde{\Hom_R(M,
N)} \to \SheafHom_{\mathcal{O}_X}(\widetilde M, \widetilde
N)$. Nous vérifions que c'est un isomorphisme en comparant les
fibres. Si $M$ est de présentation finie en tant que
$R$-module, alors $\widetilde M$ a une présentation globale finie
en tant que $\mathcal{O}_X$-module. Ainsi, nous concluons
en utilisant Algèbre, lemme \ref{algebra-lemma-hom-from-finitely-presented}
et Modules, lemme \ref{modules-lemma-stalk-internal-hom}.
\end{proof}

\begin{proof}[Troisième démonstration de la partie (1)]
Pour tout $\mathcal{O}_X$-module $\mathcal{F}$,
nous avons les isomorphismes suivants fonctoriels en $M$, $N$, et $\mathcal{F}$
\begin{align*}
\Hom_{\mathcal{O}_X}(\widetilde{M} \otimes _{\mathcal{O} _X} \widetilde{N},
\mathcal{F}) &
=
\Hom_{\mathcal{O}_X}(\widetilde{M},
\SheafHom_{\mathcal{O} _X} (\widetilde{N}, \mathcal{F})) \\ &
=
\Hom_R(M, \Gamma(X,
\SheafHom_{\mathcal{O}_X}(\widetilde{N}, \mathcal{F}))) \\ &
=
\Hom_R(M, \Hom_{\mathcal{O}_X}(\widetilde{N}, \mathcal{F})) \\ &
=
\Hom_R(M, \Hom_R(N, \Gamma(X, \mathcal{F}))) \\
& =
\Hom_R(M \otimes_R N, \Gamma(X, \mathcal{F})) \\
& =
\Hom_{\mathcal{O}_X}(\widetilde{M \otimes_R N}, \mathcal{F})
\end{align*}
La première égalité est Modules, lemme \ref{modules-lemma-internal-hom}. La
deuxième égalité est la propriété universelle de $\widetilde{M}$, voir
lemme \ref{schemes-lemma-compare-constructions}. La troisième égalité vaut par
définition de $\SheafHom$. La quatrième égalité est la propriété universelle
de $\widetilde{N}$. La cinquième égalité est
Algèbre, lemme \ref{algebra-lemma-hom-from-tensor-product}.
L'égalité finale est la propriété universelle de
$\widetilde{M \otimes_R N}$. Par
le lemme de Yoneda (Catégories, lemme \ref{categories-lemma-yoneda}) nous
obtenons (1).
\end{proof}

\begin{lemma}
\label{schemes-lemma-widetilde-pullback}
Soient
$(X, \mathcal{O}_X) = (\Spec(S), \mathcal{O}_{\Spec(S)})$,
$(Y, \mathcal{O}_Y) = (\Spec(R), \mathcal{O}_{\Spec(R)})$ des
schémas affines.
Soit $\psi : (X, \mathcal{O}_X) \to (Y, \mathcal{O}_Y)$ un morphisme
de schémas affines, correspondant à l'homomorphisme d'anneaux
$\psi^\sharp : R \to S$ (voir lemme \ref{schemes-lemma-category-affine-schemes}).
\begin{enumerate}
\item Nous avons $\psi^* \widetilde M = \widetilde{S \otimes_R M}$
fonctoriellement par rapport au $R$-module $M$.
\item Nous avons $\psi_* \widetilde N = \widetilde{N_R}$ fonctoriellement par
rapport au $S$-module $N$.
\end{enumerate}
\end{lemma}

\begin{proof}
La première assertion découle de l'identification dans
le lemme \ref{schemes-lemma-compare-constructions}
et du résultat des Modules, lemme \ref{modules-lemma-restrict-quasi-coherent}.
La deuxième assertion découle du fait
que $\psi^{-1}(D(f)) = D(\psi^\sharp(f))$ et donc
$$
\psi_* \widetilde N(D(f)) = \widetilde N(D(\psi^\sharp(f))) =
N_{\psi^\sharp(f)} = (N_R)_f = \widetilde{N_R}(D(f))
$$
comme voulu.
\end{proof}

\noindent
Le lemme \ref{schemes-lemma-widetilde-pullback} ci-dessus dit en particulier que
si vous restreignez le
faisceau $\widetilde M$ à un sous-espace ouvert affine standard
$D(f)$, alors vous obtenez $\widetilde{M_f}$. Nous utiliserons ceci à partir
de maintenant sans autre mention.

\begin{lemma}
\label{schemes-lemma-quasi-coherent-affine}
Soit $(X, \mathcal{O}_X) = (\Spec(R), \mathcal{O}_{\Spec(R)})$
un schéma affine. Soit $\mathcal{F}$ un
$\mathcal{O}_X$-module quasi-cohérent.
Alors $\mathcal{F}$ est isomorphe au faisceau associé
au $R$-module $\Gamma(X, \mathcal{F})$.
\end{lemma}

\begin{proof}
Soit $\mathcal{F}$ un $\mathcal{O}_X$-module quasi-cohérent. Comme chaque
ouvert standard $D(f)$ est quasi-compact, nous voyons que $X$ est
localement quasi-compact, c'est-à-dire que chaque point possède un système
fondamental de voisinages quasi-compacts, voir Topologie,
définition \ref{topology-definition-locally-quasi-compact}. Ainsi, d'après
Modules, lemme \ref{modules-lemma-quasi-coherent-module}, pour chaque idéal
premier $\mathfrak p \subset R$ correspondant à $x \in X$, il existe un
voisinage ouvert $x \in U \subset X$ tel que $\mathcal{F}|_U$ soit isomorphe
au faisceau quasi-cohérent associé à un certain
$\mathcal{O}_X(U)$-module $M$. En d'autres termes, nous obtenons un
recouvrement ouvert par des $U$ possédant cette propriété. Par le lemme
\ref{schemes-lemma-standard-open}, par exemple, nous pouvons affiner ce recouvrement en un
recouvrement ouvert standard. Ainsi, nous obtenons
un recouvrement $\Spec(R) = \bigcup D(f_i)$ et des
$R_{f_i}$-modules $M_i$ et des isomorphismes $\varphi_i
: \mathcal{F}|_{D(f_i)} \to \mathcal{F}_{M_i}$ pour un certain
$R_{f_i}$-module $M_i$. Sur les intersections, nous obtenons
des isomorphismes
$$
\xymatrix{
\mathcal{F}_{M_i}|_{D(f_if_j)}
\ar[rr]^{\varphi_i^{-1}|_{D(f_if_j)}} &
&
\mathcal{F}|_{D(f_if_j)}
\ar[rr]^{\varphi_j|_{D(f_if_j)}} &
&
\mathcal{F}_{M_j}|_{D(f_if_j)}.
}
$$
Notons-les $\psi_{ij}$. Il est clair que nous
avons la condition de cocycle
$$
\psi_{jk}|_{D(f_if_jf_k)}
\circ
\psi_{ij}|_{D(f_if_jf_k)}
=
\psi_{ik}|_{D(f_if_jf_k)}
$$
sur les triples intersections.

\medskip\noindent
Rappelons que chacun des sous-espaces ouverts $D(f_i)$, $D(f_if_j)$,
$D(f_if_jf_k)$ est un schéma affine. Par conséquent, les faisceaux $\mathcal{F}_{M_i}$
sont isomorphes aux faisceaux $\widetilde M_i$ d'après le lemme
\ref{schemes-lemma-compare-constructions} ci-dessus. En particulier, nous voyons que
$\mathcal{F}_{M_i}(D(f_if_j)) = (M_i)_{f_j}$, etc. De plus,
d'après le lemme \ref{schemes-lemma-compare-constructions} ci-dessus, nous
voyons que $\psi_{ij}$ correspond à un unique isomorphisme de $R_{f_if_j}$-modules
$$
\psi_{ij} : (M_i)_{f_j} \longrightarrow (M_j)_{f_i}
$$
à savoir, l'effet de $\psi_{ij}$ sur les sections au-dessus de $D(f_if_j)$. De
plus, celles-ci satisfont alors la condition de cocycle selon laquelle
$$
\xymatrix{
(M_i)_{f_jf_k}
\ar[rd]_{\psi_{ij}}
\ar[rr]^{\psi_{ik}} &
&
(M_k)_{f_if_j} \\
&
(M_j)_{f_if_k} \ar[ru]_{\psi_{jk}}
}
$$
commute (pour tout triple $i, j, k$).

\medskip\noindent
Algèbre, lemme \ref{algebra-lemma-glue-modules} montre
qu'il existe un $R$-module $M$ et des identifications $M_i
= M_{f_i}$ compatibles avec les morphismes $\psi_{ij}$. Considérons
$\mathcal{F}_M = \widetilde M$. À ce stade, c'est une formalité
de montrer que $\widetilde M$ est isomorphe au faisceau
quasi-cohérent $\mathcal{F}$ avec lequel nous avons commencé. À
savoir, les faisceaux $\mathcal{F}$ et $\widetilde M$ donnent lieu
à des ensembles de données de recollement de faisceaux de $\mathcal{O}_X$-modules
isomorphes par rapport au recouvrement $X = \bigcup D(f_i)$,
voir Faisceaux, section \ref{sheaves-section-glueing-sheaves}
et en particulier le lemme \ref{sheaves-lemma-mapping-property-glue}.
Explicitement, dans la situation actuelle, cela se ramène à
l'argument suivant : construisons un homomorphisme de $R$-modules
$$
M \longrightarrow \Gamma(X, \mathcal{F}).
$$
À savoir, étant donné $m \in M$, nous obtenons $m_i = m/1 \in M_{f_i}
= M_i$ par la construction de $M$. Par la construction de $M_i$, cela
correspond à une section $s_i \in \mathcal{F}(U_i)$. (À savoir, $\varphi^{-1}_i(m_i)$.)
Nous affirmons que $s_i|_{D(f_if_j)} = s_j|_{D(f_if_j)}$.
Cela résulte de la construction de $M$, qui donne $\psi_{ij}(m_i) = m_j$, et
de celle des $\psi_{ij}$. Par la condition de faisceau de $\mathcal{F}$,
cette collection de sections donne lieu à une section unique $s$ de
$\mathcal{F}$ sur $X$. Nous laissons au lecteur le soin de montrer que $m
\mapsto s$ est un morphisme de $R$-modules. Par le
lemme \ref{schemes-lemma-compare-constructions}, nous obtenons un morphisme de
$\mathcal{O}_X$-modules associé
$$
\widetilde M \longrightarrow \mathcal{F}.
$$
Par construction, ce morphisme se réduit aux isomorphismes
$\varphi_i^{-1}$ sur chaque $D(f_i)$ et est donc un isomorphisme.
\end{proof}

\begin{lemma}
\label{schemes-lemma-equivalence-quasi-coherent}
Soit $(X, \mathcal{O}_X) = (\Spec(R), \mathcal{O}_{\Spec(R)})$ un
schéma affine. Les
foncteurs $M \mapsto \widetilde M$ et $\mathcal{F}
\mapsto \Gamma(X, \mathcal{F})$ définissent des équivalences
quasi-inverses de catégories
$$
\xymatrix{
\QCoh(\mathcal{O}_X)
\ar@<1ex>[r]
&
\text{Mod}_R
\ar@<1ex>[l]
}
$$
entre la catégorie des $\mathcal{O}_X$-modules quasi-cohérents
et la catégorie des $R$-modules.
\end{lemma}

\begin{proof}
Voir les lemmes \ref{schemes-lemma-compare-constructions}
et \ref{schemes-lemma-quasi-coherent-affine} ci-dessus.
\end{proof}

\noindent
Dorénavant, nous ne distinguerons pas entre les faisceaux quasi-cohérents
sur les schémas affines et les faisceaux de la forme $\widetilde M$.

\begin{lemma}
\label{schemes-lemma-kernel-cokernel-quasi-coherent}
Soit $X = \Spec(R)$ un schéma affine. Les
noyaux et conoyaux des morphismes de $\mathcal{O}_X$-modules
quasi-cohérents sont quasi-cohérents.
\end{lemma}

\begin{proof}
Cela découle de l'exactitude du foncteur $\widetilde{\ }$ puisque d'après
le lemme \ref{schemes-lemma-compare-constructions}, nous savons que tout morphisme
$\psi : \widetilde{M} \to \widetilde{N}$ provient d'un
morphisme de $R$-modules $\varphi : M \to N$. (Ainsi
nous avons $\Ker(\psi) = \widetilde{\Ker(\varphi)}$
et $\Coker(\psi) = \widetilde{\Coker(\varphi)}$.)
\end{proof}

\begin{lemma}
\label{schemes-lemma-colimit-quasi-coherent}
Soit $X = \Spec(R)$ un schéma affine. La somme
directe d'une collection quelconque de faisceaux quasi-cohérents sur $X$
est quasi-cohérente. Il en va de même pour les colimites.
\end{lemma}

\begin{proof}
Supposons que les $\mathcal{F}_i$, $i \in I$, forment une collection de faisceaux
quasi-cohérents sur $X$. D'après le lemme \ref{schemes-lemma-equivalence-quasi-coherent}
ci-dessus, nous pouvons écrire $\mathcal{F}_i = \widetilde{M_i}$ pour un certain $R$-module
$M_i$. Posons $M = \bigoplus M_i$. Considérons le faisceau $\widetilde{M}$.
Pour chaque ouvert standard $D(f)$, nous avons
$$
\widetilde{M}(D(f)) = M_f =
\left(\bigoplus M_i\right)_f =
\bigoplus M_{i, f}.
$$
Ainsi nous voyons que le $\mathcal{O}_X$-module quasi-cohérent
$\widetilde{M}$ est la somme directe des faisceaux $\mathcal{F}_i$. Un argument
similaire fonctionne pour les colimites générales.
\end{proof}

\begin{lemma}
\label{schemes-lemma-extension-quasi-coherent}
Soit $(X, \mathcal{O}_X) = (\Spec(R), \mathcal{O}_{\Spec(R)})$
un schéma affine. Supposons que
$$
0 \to
\mathcal{F}_1 \to
\mathcal{F}_2 \to
\mathcal{F}_3 \to
0
$$
soit une suite exacte courte de faisceaux de $\mathcal{O}_X$-modules. Si
deux sur trois sont quasi-cohérents, alors le troisième l'est aussi.
\end{lemma}

\begin{proof}
Cela est clair dans le cas où à la fois $\mathcal{F}_1$ et $\mathcal{F}_2$
sont quasi-cohérents car le foncteur $M \mapsto \widetilde M$ est
exact, voir le lemme \ref{schemes-lemma-spec-sheaves}. De
même dans le cas où à la fois $\mathcal{F}_2$ et $\mathcal{F}_3$ sont
quasi-cohérents. Maintenant, supposons que $\mathcal{F}_1 = \widetilde M_1$ et
$\mathcal{F}_3 = \widetilde M_3$ soient quasi-cohérents.
Posons $M_2 = \Gamma(X, \mathcal{F}_2)$. Nous affirmons qu'il suffit de montrer
que la suite
$$
0 \to M_1 \to M_2 \to M_3 \to 0
$$
est exacte. En effet, si tel est le cas, alors (en utilisant la propriété
universelle du lemme \ref{schemes-lemma-compare-constructions}) nous obtenons un diagramme
commutatif
$$
\xymatrix{
0 \ar[r] &
\widetilde M_1 \ar[r] \ar[d] &
\widetilde M_2 \ar[r] \ar[d] &
\widetilde M_3 \ar[r] \ar[d] & 0
\\ 0
\ar[r] &
\mathcal{F}_1 \ar[r] &
\mathcal{F}_2 \ar[r] &
\mathcal{F}_3 \ar[r] &
0
}
$$
et le lemme du serpent permet alors de conclure.

\medskip\noindent
L'argument ``correct'' ici serait de montrer d'abord que
$H^1(X, \mathcal{F}) = 0$ pour tout faisceau quasi-cohérent $\mathcal{F}$. Cela n'est
en fait pas si difficile, mais il est peut-être préférable de le reporter à plus
tard. À la place, nous utilisons une petite astuce.

\medskip\noindent
Choisissez $m \in M_3 = \Gamma(X, \mathcal{F}_3)$.
Considérons l'ensemble suivant
$$
I = \{ f \in R \mid \text{l'élément }fm\text{ provient de }M_2\}.
$$
Clairement, il s'agit d'un idéal. Il suffit de montrer $1 \in I$.
Par conséquent, il suffit de montrer que pour tout idéal premier
$\mathfrak p$, il existe un $f \in I$, $f \not\in \mathfrak
p$. Soit $x \in X$ le point correspondant à $\mathfrak p$. Comme
la surjectivité peut être vérifiée sur les fibres,
il existe un voisinage ouvert $U$ de $x$ tel que $m|_U$
provienne d'une section locale $s \in \mathcal{F}_2(U)$. En fait, nous
pouvons supposer que $U = D(f)$ est un ouvert standard, c'est-à-dire
$f \in R$, $f \not \in \mathfrak p$. Nous allons montrer
que pour un certain $N \gg 0$ nous avons $f^N \in I$, ce
qui terminera la preuve.

\medskip\noindent
Prenons un point quelconque $z \in V(f)$, disons correspondant
à l'idéal premier $\mathfrak q \subset R$. Nous pouvons également
trouver un $g \in R$, $g \not \in \mathfrak q$ tel que
$m|_{D(g)}$ se relève en un certain $s' \in
\mathcal{F}_2(D(g))$. Considérons la différence $s|_{D(fg)}
- s'|_{D(fg)}$. C'est un élément $m'$ de $\mathcal{F}_1(D(fg)) =
(M_1)_{fg}$. Pour un certain entier $n = n(z)$, l'élément $f^n m'$
provient d'un certain $m'_1 \in (M_1)_g$.
Nous voyons que $f^n s$ s'étend à une section $\sigma$ de $\mathcal{F}_2$ sur $D(f) \cup
D(g)$ parce qu'elle coïncide avec la restriction
de $f^n s' + m'_1$ sur $D(f) \cap D(g) = D(fg)$.
De plus, l'image de $\sigma$ est la restriction de $f^n
m$ à $D(f) \cup D(g)$.

\medskip\noindent
Étant donné que $V(f)$ est quasi-compact, il existe une liste
finie d'éléments $g_1, \ldots, g_m \in R$ telle que
$V(f) \subset \bigcup D(g_j)$, un entier $n > 0$ et des sections
$\sigma_j \in \mathcal{F}_2(D(f) \cup D(g_j))$ telles que
$\sigma_j|_{D(f)} = f^n s$ et que chaque $\sigma_j$ ait pour image la section
$f^nm|_{D(f) \cup D(g_j)}$ de $\mathcal{F}_3$.
Considérons les différences
$$
\sigma_j|_{D(f) \cup D(g_jg_k)}
-
\sigma_k|_{D(f) \cup D(g_jg_k)}.
$$
Cela correspond à des sections de $\mathcal{F}_1$ sur
$D(f) \cup D(g_jg_k)$ qui sont nulles sur
$D(f)$. En particulier, leurs images dans
$\mathcal{F}_1(D(g_jg_k)) = (M_1)_{g_jg_k}$ sont
nulles dans $(M_1)_{g_jg_kf}$. Ainsi,
une certaine puissance élevée de $f$ annule chacune d'entre elles.
En d'autres termes, les éléments $f^N \sigma_j$, pour un certain $N
\gg 0$, satisfont à la condition de recollement du faisceau et
donnent lieu à une section $\sigma $ de $\mathcal{F}_2$
sur $\bigcup (D(f) \cup D(g_j)) = X$ comme souhaité.
\end{proof}







\section{Sous-espaces fermés des schémas affines}
\label{schemes-section-closed-in-affine}

\begin{example}
\label{schemes-example-closed-immersion-affines}
Soit $R$ un anneau.
Soit $I \subset R$ un idéal. Considérons
le morphisme de schémas affines $i : Z =
\Spec(R/I) \to \Spec(R) = X$. D'après Algèbre,
lemme \ref{algebra-lemma-spec-closed}, il s'agit
d'un homéomorphisme de $Z$ sur un sous-ensemble fermé
de $X$. De plus, si $I \subset \mathfrak p \subset R$ est un idéal
premier correspondant à un point $x = i(z)$, $x \in X$, $z \in Z$,
alors sur les fibres, on obtient l'application
$$
\mathcal{O}_{X, x} = R_{\mathfrak p}
\longrightarrow
R_{\mathfrak p}/IR_{\mathfrak p} = \mathcal{O}_{Z, z}
$$
Ainsi, nous voyons que $i$ est une immersion fermée d'espaces annelés en
anneaux locaux, voir la définition \ref{schemes-definition-closed-immersion-locally-ringed-spaces}.
Clairement, cela est (isomorphe) au sous-espace fermé associé
au faisceau quasi-cohérent d'idéaux $\widetilde I$, comme dans
l'exemple \ref{schemes-example-closed-subspace}.
\end{example}

\begin{lemma}
\label{schemes-lemma-closed-immersion-affine-case}
\begin{slogan}
Pour les schémas affines, les immersions fermées correspondent aux idéaux.
\end{slogan}
Soit $(X, \mathcal{O}_X) = (\Spec(R), \mathcal{O}_{\Spec(R)})$ un
schéma affine. Soit $i : Z \to X$ une quelconque immersion fermée
d'espaces annelés en anneaux locaux. Alors il existe un unique
idéal $I \subset R$ tel que le morphisme $i : Z \to X$ puisse être identifié
avec l'immersion fermée $\Spec(R/I) \to \Spec(R)$
construite dans l'exemple \ref{schemes-example-closed-immersion-affines} ci-dessus.
\end{lemma}

\begin{proof}
La vérification est immédiate : d'après le lemme \ref{schemes-lemma-closed-immersion},
nous pouvons identifier $Z \to X$ avec le sous-espace fermé associé à
un faisceau d'idéaux $\mathcal{I} \subset \mathcal{O}_X$ comme dans
la définition \ref{schemes-definition-closed-subspace} et
l'exemple \ref{schemes-example-closed-subspace}. Selon
nos conventions, ce faisceau d'idéaux est localement engendré par
des sections en tant que faisceau de $\mathcal{O}_X$-modules.
Par conséquent, le faisceau quotient $\mathcal{O}_X /
\mathcal{I}$ est localement sur $X$
le conoyau d'une application $\bigoplus_{j \in J} \mathcal{O}_U \to
\mathcal{O}_U$. Ainsi, par définition, $\mathcal{O}_X / \mathcal{I}$ est
quasi-cohérent. D'après nos résultats de la section \ref{schemes-section-quasi-coherent-affine},
il est de la forme $\widetilde S$ pour un certain
$R$-module $S$. De plus, comme $\mathcal{O}_X = \widetilde R \to \widetilde
S$ est surjectif, nous voyons d'après le lemme \ref{schemes-lemma-extension-quasi-coherent}
que $\mathcal{I}$ est également quasi-cohérent, disons $\mathcal{I} =
\widetilde I$. Bien sûr, $I \subset R$ et $S = R/I$ et tout est clair.
\end{proof}













\section{Schémas}
\label{schemes-section-schemes}

\begin{definition}
\label{schemes-definition-scheme}
\begin{history}
Dans \cite{EGA1} ce que nous appelons un schéma était appelé un ``pré-schéma''
et le nom ``schéma'' était réservé pour ce qui est un schéma séparé
dans le projet Stacks. Dans la deuxième édition \cite{EGA1-second} la terminologie
a été remplacée par celle qui est désormais standard. Cependant, on
peut parfois rencontrer la terminologie ``préschéma'', par exemple dans
\cite{Murre-lectures}.
\end{history}
Un {\it schéma} est un espace annelé en anneaux locaux avec la propriété
que chaque point possède un voisinage ouvert qui est un schéma affine.
Un {\it morphisme de schémas} est un morphisme d'espaces
annelés en anneaux locaux. La catégorie des schémas sera notée
$\Sch$.
\end{definition}

\noindent
Soit $X$ un schéma. Nous
utiliserons le (très léger) abus de langage suivant. Nous
dirons que $U \subset X$ est une {\it partie ouverte affine}, ou un {\it ouvert
affine}, si le sous-espace ouvert $U$ est un schéma affine. Nous
écrirons souvent $U = \Spec(R)$ pour indiquer que $U$ est
isomorphe à $\Spec(R)$ et de plus que nous identifierons (temporairement)
$U$ et $\Spec(R)$.

\begin{lemma}
\label{schemes-lemma-open-subspace-scheme}
Soit $X$ un schéma. Soit $j : U \to X$ une immersion ouverte d'espaces
annelés en anneaux locaux. Alors $U$ est un schéma. En particulier, tout
sous-espace ouvert de $X$ est un schéma.
\end{lemma}

\begin{proof}
Soit $U \subset X$. Soit $u \in U$.
Choisissez un voisinage ouvert affine $u \in V \subset X$.
Comme les ouverts standards de $V$ forment une base de la
topologie sur $V$, nous voyons qu'il existe un $f\in \mathcal{O}_V(V)$
tel que $u \in D(f) \subset U$. Et $D(f)$ est un schéma affine
d'après le lemme \ref{schemes-lemma-standard-open-affine}. Cela prouve que chaque point
de $U$ possède un voisinage ouvert qui est affine.
\end{proof}

\noindent
De toute évidence, le lemme (ou sa preuve) montre que tout
schéma $X$ a une base (voir Topologie, section \ref{topology-section-bases})
pour la topologie constituée d'ouverts affines.

\begin{example}
\label{schemes-example-not-affine}
Soit $k$ un corps. Un
exemple de schéma qui n'est pas affine est donné
par le sous-espace ouvert $U
= \Spec(k[x, y]) \setminus \{ (x, y)\}$ du
schéma affine $X =\Spec(k[x, y])$. Il est
recouvert par deux ouverts affines, à savoir $D(x) = \Spec(k[x, y, 1/x])$
et $D(y) = \Spec(k[x, y, 1/y])$ dont l'intersection est
$D(xy) = \Spec(k[x, y, 1/xy])$. Par la propriété du faisceau
pour $\mathcal{O}_U$, il existe une suite exacte
$$
0 \to
\Gamma(U, \mathcal{O}_U) \to
k[x, y, 1/x] \times k[x, y, 1/y] \to
k[x, y, 1/xy]
$$
Nous concluons que l'application $k[x, y] \to \Gamma(U, \mathcal{O}_U)$
(provenant du morphisme $U \to X$) est un isomorphisme. Par
conséquent, $U$ ne peut pas être affine puisqu'alors, d'après
le lemme \ref{schemes-lemma-category-affine-schemes}, nous aurions $U \cong X$.
\end{example}











\section{Immersions de schémas}
\label{schemes-section-immersions}

\noindent
Dans le lemme \ref{schemes-lemma-open-subspace-scheme}, nous avons vu que tout sous-espace ouvert
d'un schéma est un schéma. Ci-dessous, nous prouverons que la même chose vaut
pour un sous-espace fermé d'un schéma.

\medskip\noindent
Notez que la notion de faisceau quasi-cohérent de $\mathcal{O}_X$-modules est définie
pour tout espace annelé $X$, en particulier lorsque $X$ est un schéma. D'après
les résultats de la section \ref{schemes-section-quasi-coherent-affine}, nous
savons qu'un tel faisceau sur tout ouvert affine $U \subset X$ est
de la forme $\widetilde M$ pour un certain $\mathcal{O}_X(U)$-module $M$.

\begin{lemma}
\label{schemes-lemma-closed-subspace-scheme}
Soit $X$ un schéma. Soit $i : Z \to X$ une immersion fermée d'espaces
annelés en anneaux locaux.
\begin{enumerate}
\item L'espace annelé en anneaux locaux $Z$ est un
schéma, \item Le noyau $\mathcal{I}$ de l'application
$\mathcal{O}_X \to i_*\mathcal{O}_Z$ est un faisceau
quasi-cohérent
d'idéaux, \item pour tout ouvert affine $U = \Spec(R)$ de
$X$ le morphisme $i^{-1}(U) \to U$ peut être identifié à
$\Spec(R/I) \to \Spec(R)$ pour un certain idéal $I \subset R$, et
\item nous avons $\mathcal{I}|_U = \widetilde I$.
\end{enumerate}
En particulier, tout faisceau d'idéaux localement engendré par des sections
est un faisceau quasi-cohérent d'idéaux (et réciproquement),
et tout sous-espace fermé de $X$ est un schéma.
\end{lemma}

\begin{proof}
Soit $i : Z \to X$ une immersion fermée. Soit
$z \in Z$ un point. Choisissez un voisinage ouvert
affine $i(z) \in U \subset X$ quelconque. Disons $U =
\Spec(R)$. D'après le lemme \ref{schemes-lemma-closed-immersion-affine-case},
nous savons que $i^{-1}(U) \to U$ peut être identifié au
morphisme de schémas affines $\Spec(R/I) \to
\Spec(R)$. Tout d'abord, cela implique que $z \in i^{-1}(U) \subset Z$
est un voisinage affine de $z$. Ainsi, $Z$ est un schéma.
Deuxièmement, cela implique que $\mathcal{I}|_U$ est $\widetilde
I$. En d'autres termes, pour chaque point $x \in i(Z)$, il
existe un voisinage ouvert tel que $\mathcal{I}$ soit quasi-cohérent
dans ce voisinage. Notez que $\mathcal{I}|_{X \setminus i(Z)}
\cong \mathcal{O}_{X \setminus i(Z)}$. Ainsi, la restriction du
faisceau d'idéaux est quasi-cohérente sur $X \setminus i(Z)$
également. Nous concluons que $\mathcal{I}$ est quasi-cohérent.
\end{proof}

\begin{definition}
\label{schemes-definition-immersion}
Soit $X$ un schéma.
\begin{enumerate}
\item Un morphisme de schémas est appelé une {\it immersion ouverte}
s'il s'agit d'une immersion ouverte d'espaces annelés en anneaux
locaux (voir la définition \ref{schemes-definition-immersion-locally-ringed-spaces}).
\item Un {\it sous-schéma ouvert} de $X$ est un sous-espace ouvert
de $X$ au sens de la définition \ref{schemes-definition-open-subspace} ; un sous-schéma ouvert
de $X$ est un schéma d'après le lemme \ref{schemes-lemma-open-subspace-scheme}.
\item Un morphisme de schémas est appelé une {\it immersion fermée}
s'il s'agit d'une immersion fermée d'espaces annelés en anneaux
locaux (voir la définition \ref{schemes-definition-closed-immersion-locally-ringed-spaces}).
\item Un {\it sous-schéma fermé} de $X$ est un sous-espace fermé
de $X$ au sens de la définition \ref{schemes-definition-closed-subspace} ; un sous-schéma fermé est
un schéma d'après le lemme \ref{schemes-lemma-closed-subspace-scheme}.
\item Un morphisme de schémas $f : X \to Y$ est appelé une {\it immersion}, ou une
{\it immersion localement fermée} s'il peut se factoriser sous la
forme $j \circ i$ où $i$ est une immersion fermée et $j$ est une immersion
ouverte.
\end{enumerate}
\end{definition}

\noindent
Il résulte des lemmes des sections \ref{schemes-section-open-immersion} et
\ref{schemes-section-closed-immersion} que toute immersion ouverte (resp. \ fermée) de
schémas est isomorphe à l'inclusion d'un sous-schéma ouvert (resp. \ fermé)
de la cible.

\medskip\noindent
Notre définition d'une immersion fermée se situe à mi-chemin entre Hartshorne
et l'EGA. Hartshorne définit une immersion fermée comme un morphisme $f : X \to Y$ de
schémas qui induit un homéomorphisme de $X$ sur un sous-ensemble fermé de $Y$
et tel que l'application $f^\# : \mathcal{O}_Y \to f_*\mathcal{O}_X$ soit surjective,
voir \cite[Page 85]{H}. Nous montrerons que cela est équivalent à notre
notion dans le lemme \ref{schemes-lemma-characterize-closed-immersions}. Dans
\cite{EGA}, Grothendieck et Dieudonné définissent d'abord les sous-schémas
fermés via la construction de l'exemple \ref{schemes-example-closed-subspace}
en utilisant des faisceaux quasi-cohérents d'idéaux, puis définissent une
immersion fermée comme un morphisme $f : X \to Y$ qui induit un isomorphisme
avec un sous-schéma fermé. Il découle du lemme \ref{schemes-lemma-closed-subspace-scheme}
que cela correspond à notre notion.

\medskip\noindent
D'un point de vue pédagogique, la définition ci-dessus est un désastre/cauchemar.
Lors de l'enseignement de ce matériel aux étudiants, nous
avons souvent trouvé pratique de définir une immersion fermée comme
un morphisme affine $f : X \to Y$ de schémas tel que $f^\# : \mathcal{O}_Y \to
f_*\mathcal{O}_X$ soit surjectif. En effet, il s'avère que la notion de
morphisme affine (Morphismes, section \ref{morphisms-section-affine})
est assez naturelle et facile à comprendre.

\medskip\noindent
Pour plus d'informations sur les immersions fermées, nous suggérons au lecteur
de consulter Morphismes, sections \ref{morphisms-section-closed-immersions}
et \ref{morphisms-section-closed-immersions-quasi-coherent}.

\medskip\noindent
Nous discuterons des sous-schémas localement fermés et des immersions
à la fin de cette section.

\begin{remark}
\label{schemes-remark-not-reverse-open-closed}
Si $f : X \to Y$ est une immersion de schémas, il n'est généralement pas
possible de factoriser $f$ en une immersion ouverte suivie
d'une immersion fermée. Voir Morphismes, exemple \ref{morphisms-example-thibaut}.
\end{remark}

\begin{lemma}
\label{schemes-lemma-immersion-when-closed}
Soit $f : Y \to X$ une immersion de schémas. Alors $f$ est une immersion
fermée si et seulement si $f(Y) \subset X$ est un sous-ensemble fermé.
\end{lemma}

\begin{proof}
Si $f$ est une immersion fermée, alors $f(Y)$ est fermé par définition.
Réciproquement, supposons que $f(Y)$ soit fermé. Par
définition, il existe un sous-schéma ouvert $U \subset X$ tel que $f$ soit la composée
d'une immersion fermée $i : Y \to U$ et de l'immersion ouverte
$j : U \to X$. Soit $\mathcal{I} \subset \mathcal{O}_U$ le
faisceau d'idéaux quasi-cohérent associé à l'immersion fermée $i$. Notez
$\mathcal{I}|_{U
\setminus i(Y)} = \mathcal{O}_{U
\setminus i(Y)} = \mathcal{O}_{X
\setminus i(Y)}|_{U \setminus i(Y)}$. Ainsi, nous pouvons
recoller (voir Faisceaux, section \ref{sheaves-section-glueing-sheaves}) $\mathcal{I}$
et $\mathcal{O}_{X \setminus i(Y)}$ en un faisceau d'idéaux
$\mathcal{J} \subset \mathcal{O}_X$. Puisque chaque point de
$X$ possède un voisinage où $\mathcal{J}$ est quasi-cohérent,
nous voyons que $\mathcal{J}$ est quasi-cohérent (en particulier
localement engendré par des sections). Par
construction, $\mathcal{O}_X/\mathcal{J}$ est supporté sur $U$
et sa restriction à cet ouvert est égale à $\mathcal{O}_U/\mathcal{I}$.
Ainsi, nous voyons que les sous-espaces fermés associés à $\mathcal{I}$
et $\mathcal{J}$ sont canoniquement isomorphes,
voir l'exemple \ref{schemes-example-closed-subspace}.
En particulier, le sous-espace fermé de $U$ associé à $\mathcal{I}$
est isomorphe à un sous-espace fermé de $X$.
Puisque $Y \to U$ est identifié avec le sous-espace fermé
associé à $\mathcal{I}$, voir le lemme \ref{schemes-lemma-closed-immersion}, nous
concluons que $Y \to U \to X$ est
une immersion fermée.
\end{proof}

\noindent
Soit $f : Y \to X$ une immersion. Soit
$Z = \overline{f(Y)} \setminus f(Y)$ qui est un sous-ensemble fermé de $X$. Soit $U
= X \setminus Z$. Le lemme
implique que $U$ est le plus grand sous-espace ouvert de $X$ tel que $f
: Y \to X$ se factorise par une immersion fermée dans $U$. Nous définissons
un {\it sous-schéma localement fermé de $X$} comme une paire $(Z, U)$
constituée d'un sous-schéma fermé $Z$ d'un sous-schéma ouvert $U$ de $X$ tel
qu'en outre $\overline{Z} \cup U = X$. Nous disons généralement simplement «
soit $Z$ un sous-schéma localement fermé de $X$ » puisque nous pouvons récupérer
$U$ à partir du morphisme $Z \to X$. Ce qui précède
montre alors que toute immersion $f : Y \to X$ se factorise de manière
unique en $Y \to Z \to X$ où $Z$ est un sous-schéma localement fermé de
$X$ et $Y \to Z$ est un isomorphisme.

\medskip\noindent
L'intérêt de ceci est que la collection de sous-schémas localement fermés de $X$ forme
un ensemble. Nous pouvons définir un ordre partiel sur cet ensemble, que
nous appelons inclusion pour des raisons évidentes. Pour être précis, si
$Z \to X$ et $Z' \to X$ sont deux sous-schémas localement fermés de $X$, alors nous
disons que {\it $Z$ est contenu dans $Z'$} simplement si le morphisme $Z
\to X$ se factorise à travers $Z'$. Si c'est le cas, alors bien sûr $Z$ est identifié à
un sous-schéma localement fermé unique de $Z'$, et ainsi de suite.

\medskip\noindent
Pour plus d'informations sur les immersions, nous renvoyons le lecteur
à Morphismes, section \ref{morphisms-section-immersions}.










\section{Topologie de Zariski des schémas}
\label{schemes-section-topology}

\noindent
Voir Topologie, section \ref{topology-section-introduction} pour
quelques notions de base en topologie adaptées à la topologie de
Zariski des schémas.

\begin{lemma}
\label{schemes-lemma-scheme-sober}
Soit $X$ un schéma. Tout
sous-ensemble fermé irréductible de $X$ a un point générique unique.
En d'autres termes, $X$ est un espace topologique sobre,
voir Topologie, définition \ref{topology-definition-generic-point}.
\end{lemma}

\begin{proof}
Soit $Z \subset X$ un sous-ensemble fermé irréductible.
Pour chaque ouvert affine $U \subset X$, $U = \Spec(R)$,
nous savons que $Z \cap U = V(I)$ pour un idéal
radical unique $I \subset R$. Notez que $Z \cap U$ est soit vide,
soit irréductible. Dans le second cas (qui se produit pour
au moins un $U$), nous voyons que $I = \mathfrak p$ est
un idéal premier, qui est un point générique $\xi$ de $Z \cap U$. Il
s'ensuit que $Z = \overline{\{\xi\}}$, en d'autres termes $\xi$
est un point générique de $Z$. Si $\xi'$ était un deuxième
point générique, alors $\xi' \in Z \cap U$ et il en découle
immédiatement que $\xi' = \xi$.
\end{proof}

\begin{lemma}
\label{schemes-lemma-basis-affine-opens}
Soit $X$ un schéma. La collection des ouverts affines
de $X$ forme une base pour la topologie sur $X$.
\end{lemma}

\begin{proof}
Cela découle de la discussion sur les sous-schémas ouverts
dans la section \ref{schemes-section-schemes}.
\end{proof}

\begin{remark}
\label{schemes-remark-intersection-affine-opens}
En général, l'intersection de deux ouverts affines dans $X$
n'est pas un ouvert affine. Voir l'exemple \ref{schemes-example-affine-space-zero-doubled}.
\end{remark}

\begin{lemma}
\label{schemes-lemma-locally-quasi-compact}
L'espace topologique sous-jacent de tout schéma est
localement quasi-compact, voir
Topologie, définition \ref{topology-definition-locally-quasi-compact}.
\end{lemma}

\begin{proof}
Cela découle du lemme \ref{schemes-lemma-basis-affine-opens} ci-dessus
et du fait que le spectre d'un anneau est quasi-compact, voir
Algèbre, lemme \ref{algebra-lemma-quasi-compact}.
\end{proof}

\begin{lemma}
\label{schemes-lemma-standard-open-two-affines}
Soit $X$ un schéma.
Soient $U, V$ des ouverts affines de $X$, et soit $x \in U \cap V$.
Il existe un voisinage ouvert affine $W$ de $x$ tel que
$W$ soit un ouvert standard à la fois de $U$ et de $V$.
\end{lemma}

\begin{proof}
Posons $U = \Spec(A)$ et $V = \Spec(B)$. Supposons
que $x$ corresponde aux idéaux premiers $\mathfrak p \subset
A$ et $\mathfrak q \subset B$. Nous pouvons
choisir un $f \in A$, $f \not \in \mathfrak p$ tel que $D(f) \subset
U \cap V$. Remarquons que tout ouvert standard de $D(f)$ est un ouvert
standard de $\Spec(A) = U$. Ainsi, nous pouvons supposer que $U
\subset V$. En d'autres termes, maintenant nous pouvons considérer
$U$ comme un ouvert affine de $V$. Ensuite, nous
choisissons un $g \in B$, $g \not \in \mathfrak
q$ tel que $D(g) \subset U$. Dans ce cas, nous voyons que $D(g) =
D(g_A)$ où $g_A \in A$ désigne l'image de $g$ par l'application $B \to A$.
Ainsi le lemme est prouvé.
\end{proof}

\begin{lemma}
\label{schemes-lemma-good-subcover}
Soit $X$ un schéma.
Soit $X = \bigcup_i U_i$ un recouvrement ouvert affine.
Soit $V \subset X$ un ouvert affine.
Il existe un recouvrement ouvert standard
$V = \bigcup_{j = 1, \ldots, m} V_j$ (voir
définition \ref{schemes-definition-standard-covering})
tel que chaque $V_j$ soit un ouvert standard de l'un des $U_i$.
\end{lemma}

\begin{proof}
Choisissez $v \in V$. Alors $v \in U_i$ pour un certain
$i$. D'après le lemme \ref{schemes-lemma-standard-open-two-affines} ci-dessus, il existe un
ouvert $v \in W_v \subset V \cap U_i$ tel que
$W_v$ soit un ouvert standard à la fois dans $V$ et
$U_i$. Comme $V$ est quasi-compact, le lemme s'ensuit.
\end{proof}

\begin{lemma}
\label{schemes-lemma-sheaf-on-affines}
Soit $X$ un schéma. Soit $\mathcal{B}$ l'ensemble des ouverts affines de $X$. Soit
$\mathcal{F}$ un préfaisceau d'ensembles sur $\mathcal{B}$, voir
Faisceaux, définition \ref{sheaves-definition-presheaf-basis}. Les énoncés suivants
sont équivalents
\begin{enumerate}
\item $\mathcal{F}$ est la restriction d'un faisceau sur $X$ à $\mathcal{B}$, \item
$\mathcal{F}$ est un faisceau sur $\mathcal{B}$, et
\item $\mathcal{F}(\emptyset)$ est un singleton et chaque
fois que $U = V \cup W$ avec $U, V, W \in \mathcal{B}$ et $V,
W \subset U$ ouverts standards
(Algèbre, définition \ref{algebra-definition-Zariski-topology})
l'application
$$
\mathcal{F}(U) \longrightarrow \mathcal{F}(V) \times \mathcal{F}(W)
$$
est injective avec pour image l'ensemble des paires $(s,
t)$ telles que $s|_{V \cap W} = t|_{V \cap W}$.
\end{enumerate}
\end{lemma}

\begin{proof}
L'équivalence de (1) et (2) est Faisceaux,
lemme \ref{sheaves-lemma-restrict-basis-equivalence}. Il est clair
que (2) implique (3). Par conséquent,
il suffit de prouver que (3) implique (2). D'après
Faisceaux, lemme \ref{sheaves-lemma-cofinal-systems-coverings-standard-case}, et le
lemme \ref{schemes-lemma-standard-open} ci-dessus, il suffit de prouver que la condition
de faisceau est satisfaite pour les recouvrements
ouverts standards (définition
\ref{schemes-definition-standard-covering}) des éléments
de $\mathcal{B}$. Soit $U = U_1 \cup \ldots \cup U_n$ un recouvrement
ouvert standard avec $U \subset X$
ouvert affine. Nous prouverons la
condition de faisceau pour ce recouvrement par
récurrence sur $n$. Si $n = 0$, alors $U$ est vide et nous obtenons
la condition de faisceau par hypothèse. Si $n = 1$, alors il
n'y a rien à prouver. Si $n = 2$, alors il
s'agit de l'hypothèse (3). Si $n > 2$, alors nous écrivons $U_i = D(f_i)$ pour $f_i \in
A = \mathcal{O}_X(U)$. Supposons que les $s_i \in \mathcal{F}(U_i)$
soient des sections telles que $s_i|_{U_i \cap U_j} = s_j|_{U_i \cap U_j}$ pour tout
$1 \leq i < j \leq n$. Comme $U = U_1 \cup \ldots
\cup U_n$, nous avons $1 = \sum_{i = 1, \ldots, n} a_i f_i$ dans $A$ pour certains
$a_i \in A$, voir Algèbre, lemme \ref{algebra-lemma-Zariski-topology}.
Posons $g = \sum_{i = 1, \ldots,
n - 1} a_if_i$. Alors $U = D(g)
\cup D(f_n)$. Remarquons que $D(g) = D(gf_1) \cup \ldots \cup D(gf_{n
- 1})$ est un recouvrement ouvert standard. Par récurrence, il existe une
section unique $s' \in \mathcal{F}(D(g))$ qui
concorde avec $s_i|_{D(gfi)}$ pour $i = 1, \ldots,
n - 1$. Nous affirmons que $s'$ et $s_n$ ont la même restriction
à $D(gf_n)$. Cela est vrai par récurrence, en utilisant le
recouvrement $D(gf_n) = D(gf_nf_1) \cup \ldots \cup D(gf_nf_{n
- 1})$. Ainsi, il existe une section unique $s \in \mathcal{F}(U)$
dont la restriction à $D(g)$ est $s'$ et dont la restriction
à $D(f_n)$ est $s_n$. Nous omettons la vérification que
$s$ se restreint à $s_i$ sur $D(f_i)$ pour $i = 1, \ldots, n - 1$
et nous omettons la vérification que $s$ est unique.
\end{proof}

\begin{lemma}
\label{schemes-lemma-scheme-finite-discrete-affine}
Soit $X$ un schéma dont l'espace topologique sous-jacent est
un ensemble discret fini.
Alors $X$ est affine.
\end{lemma}

\begin{proof}
Prenons $X = \{x_1, \ldots, x_n\}$. Alors $U_i = \{x_i\}$ est un voisinage ouvert de
$x_i$. Par le
lemme \ref{schemes-lemma-basis-affine-opens} il est
affine. Par conséquent, $X$ est une union disjointe finie de schémas affines, et
est donc affine par
le lemme \ref{schemes-lemma-disjoint-union-affines}.
\end{proof}

\begin{example}
\label{schemes-example-scheme-without-closed-points}
Il existe un schéma sans points fermés. À savoir,
soit $R$ un anneau local intègre dont le spectre ressemble
à $(0) =
\mathfrak p_0 \subset \mathfrak p_1 \subset \mathfrak p_2 \subset \ldots
\subset \mathfrak m$. Alors le sous-schéma ouvert $\Spec(R) \setminus
\{\mathfrak m\}$ n'a pas de point fermé. Pour voir qu'un
tel anneau $R$ existe, nous utilisons le fait que pour tout groupe
totalement ordonné $(\Gamma, \geq)$, il existe un anneau
de valuation $A$ dont le groupe de valuation est $(\Gamma, \geq)$, voir \cite{Krull}.
Voir Algèbre, section \ref{algebra-section-valuation-rings} pour la notation. Nous
prenons $\Gamma =
\mathbf{Z}x_1 \oplus \mathbf{Z}x_2 \oplus \mathbf{Z}x_3 \oplus \ldots$ et nous définissons
$\sum_i a_i x_i \geq 0$ si et seulement si le premier $a_i$ non
nul est $> 0$, ou si tous les coefficients satisfont $a_i =
0$. Donc $x_1 \geq x_2 \geq x_3 \geq \ldots \geq 0$.
Les sous-ensembles $x_i + \Gamma_{\geq 0}$ sont des idéaux
premiers de $(\Gamma, \geq)$, voir Algèbre, notation
au-dessus du lemme \ref{algebra-lemma-ideals-valuation-ring}. Ceux-ci,
ainsi que $\emptyset$ et $\Gamma_{\geq 0}$, sont les seuls idéaux premiers.
Par conséquent, $A$ est un exemple d'un anneau avec la structure donnée
de son spectre,
selon Algèbre, lemme \ref{algebra-lemma-ideals-valuation-ring}.
\end{example}



















\section{Schémas réduits}
\label{schemes-section-reduced}

\begin{definition}
\label{schemes-definition-reduced}
Soit $X$ un schéma. Nous disons que $X$ est {\it réduit} si chaque anneau local
$\mathcal{O}_{X, x}$ est réduit.
\end{definition}

\begin{lemma}
\label{schemes-lemma-reduced}
Un schéma $X$ est réduit si et seulement si $\mathcal{O}_X(U)$ est
un anneau réduit pour tout $U \subset X$ ouvert.
\end{lemma}

\begin{proof}
Supposons que $X$ soit réduit.
Soit $f \in \mathcal{O}_X(U)$ une section telle que $f^n = 0$. Alors
l'image de $f$ dans $\mathcal{O}_{U, u}$ est nulle pour tout $u
\in U$. Par conséquent, $f$ est nul, voir
Faisceaux, lemme \ref{sheaves-lemma-sheaf-subset-stalks}.
Inversement, supposons que $\mathcal{O}_X(U)$ soit réduit
pour tous les ouverts $U$. Prenez un élément non nul $f \in \mathcal{O}_{X, x}$.
Tout représentant $(U, f \in \mathcal{O}(U))$ de $f$ est non nul et donc non
nilpotent. Par conséquent, $f$ n'est pas nilpotent dans $\mathcal{O}_{X, x}$.
\end{proof}

\begin{lemma}
\label{schemes-lemma-affine-reduced}
Un schéma affine $\Spec(R)$ est réduit si
et seulement si $R$ est réduit.
\end{lemma}

\begin{proof}
L'implication directe suit immédiatement du lemme
\ref{schemes-lemma-reduced} ci-dessus. Dans
l'autre sens, cela découle du fait que toute localisation d'un
anneau réduit est réduite, et en particulier que les anneaux locaux
d'un anneau réduit sont réduits.
\end{proof}

\begin{lemma}
\label{schemes-lemma-reduced-closed-subscheme}
Soit $X$ un schéma. Soit $T \subset X$ un sous-ensemble fermé.
Il existe un sous-schéma fermé unique $Z \subset X$ possédant
les propriétés suivantes : (a) l'espace topologique
sous-jacent de $Z$ est égal à $T$, et (b) $Z$ est réduit.
\end{lemma}

\begin{proof}
Soit $\mathcal{I} \subset \mathcal{O}_X$ le sous-préfaisceau
défini par la règle
$$
\mathcal{I}(U) = \{f \in \mathcal{O}_X(U) \mid
f(t) = 0\text{ pour tout }t \in T\cap U\}
$$
Ici, nous utilisons $f(t)$ pour indiquer l'image
de $f$ dans le corps résiduel $\kappa(t)$ de $X$ en $t$.
En raison de la nature locale de la condition, il est
clair que $\mathcal{I}$ est un faisceau d'idéaux. De plus,
soit $U = \Spec(R)$ un ouvert affine. Nous
pouvons écrire $T \cap U = V(I)$ pour un idéal radical
unique $I \subset R$. Étant donné un $\mathfrak p \in V(I)$
premier correspondant à $t \in T \cap U$ et un élément $f \in R$, nous avons
$f(t) = 0 \Leftrightarrow f \in \mathfrak p$.
Ainsi $\mathcal{I}(U) = \bigcap_{\mathfrak p \in V(I)} \mathfrak p = I$ par
Algèbre, lemme \ref{algebra-lemma-Zariski-topology}. De plus, pour
tout ouvert standard $D(g) \subset \Spec(R) = U$, nous avons
$\mathcal{I}(D(g)) = I_g$ par le même raisonnement. Ainsi $\widetilde
I$ et $\mathcal{I}|_U$ coïncident (en tant qu'idéaux) sur
une base d'ouverts et sont donc égaux. Par conséquent,
$\mathcal{I}$ est un faisceau quasi-cohérent d'idéaux.

\medskip\noindent
À ce stade, nous pouvons définir $Z$ comme le sous-espace fermé
associé au faisceau d'idéaux $\mathcal{I}$. Pour chaque ouvert
affine $U = \Spec(R)$ de $X$, nous voyons que $Z \cap
U = \Spec(R/I)$ où $I$ est un idéal radical et donc $Z$ est
réduit (d'après le lemme \ref{schemes-lemma-affine-reduced} ci-dessus). Par construction,
le sous-ensemble fermé sous-jacent de $Z$ est $T$. Par conséquent,
nous avons trouvé un sous-schéma fermé avec les propriétés (a) et (b).

\medskip\noindent
Soit $Z' \subset X$ un second sous-schéma fermé avec les propriétés (a) et (b). Pour
chaque ouvert affine $U = \Spec(R)$ de $X$, nous voyons que $Z'
\cap U = \Spec(R/I')$ pour un certain idéal $I' \subset R$.
D'après le lemme \ref{schemes-lemma-affine-reduced}, l'anneau $R/I'$ est
réduit et donc $I'$ est radical. Puisque $V(I') = T \cap U = V(I)$, nous
en déduisons que $I = I'$ d'après
Algèbre, lemme \ref{algebra-lemma-Zariski-topology}. Par
conséquent, $Z'$ et $Z$ sont définis par le même faisceau d'idéaux
et donc sont égaux.
\end{proof}

\begin{definition}
\label{schemes-definition-reduced-induced-scheme}
Soit $X$ un schéma. Soit $Z \subset X$ un sous-ensemble fermé. Une
{\it structure de schéma sur $Z$} est donnée par un sous-schéma fermé $Z'$
de $X$ dont l'ensemble sous-jacent est égal à $Z$. Nous disons
souvent « soit $(Z, \mathcal{O}_Z)$ une structure de schéma sur
$Z$ » pour indiquer cela. La {\it structure de schéma réduite
induite} sur $Z$ est celle construite dans le lemme \ref{schemes-lemma-reduced-closed-subscheme}.
La {\it réduction $X_{red}$ de $X$} est la structure de schéma réduite
induite sur $X$ lui-même.
\end{definition}

\noindent
Souvent, lorsque nous disons « soit $Z \subset X$ une composante irréductible de $X$ »,
nous pensons à $Z$ comme à un sous-schéma fermé réduit de $X$ en utilisant la structure de schéma
réduite induite.

\begin{remark}
\label{schemes-remark-reduced-induced-locally-closed}
Soit $X$ un schéma. Soit $T \subset X$ un sous-ensemble localement fermé.
Dans cette situation, nous utilisons parfois aussi l'expression
« structure de schéma réduite induite sur $T$ ». Cela se
réfère à la structure de schéma réduite induite
de la définition \ref{schemes-definition-reduced-induced-scheme}
lorsque nous considérons $T$ comme un sous-ensemble fermé du
sous-schéma ouvert $X \setminus \partial T$ de
$X$. Ici, $\partial T = \overline{T} \setminus T$ est la « frontière » de
$T$ dans l'espace topologique de $X$.
\end{remark}

\begin{lemma}
\label{schemes-lemma-map-into-reduction}
Soit $X$ un schéma. Soit
$Z \subset X$ un sous-schéma fermé. Soit $Y$
un schéma réduit. Un morphisme
$f : Y \to X$ se factorise à travers $Z$ si et seulement si $f(Y)
\subset Z$ (au sens des ensembles). En particulier, tout
morphisme $Y \to X$ se factorise comme $Y \to X_{red} \to X$.
\end{lemma}

\begin{proof}
Supposons $f(Y) \subset Z$ (au sens des ensembles).
Soit $\mathcal{I} \subset \mathcal{O}_X$ le faisceau d'idéaux de $Z$.
Pour tous ouverts affines $U \subset X$ et $\Spec(B) = V \subset Y$
tels que $f(V) \subset U$, et pour tout $g \in \mathcal{I}(U)$,
l'image inverse $b = f^\sharp(g) \in \Gamma(V, \mathcal{O}_Y) =
B$ a une image nulle dans le corps résiduel de tout $y \in
V$. En d'autres termes $b \in \bigcap_{\mathfrak p \subset B} \mathfrak
p$. Cela implique $b = 0$ puisque $B$ est réduit (lemme \ref{schemes-lemma-reduced} et
Algèbre, lemme \ref{algebra-lemma-Zariski-topology}). Ainsi
$f$ se factorise à travers
$Z$ selon le lemme \ref{schemes-lemma-characterize-closed-subspace}.
\end{proof}










\section{Points des schémas}
\label{schemes-section-points}

\noindent
Étant donné un schéma $X$, nous pouvons définir un foncteur
$$
h_X : \Sch^{opp}
\longrightarrow
\textit{Ensembles}, \quad
T \longmapsto \Mor(T, X).
$$
Voir Catégories, exemple \ref{categories-example-hom-functor}. On l'appelle
le foncteur des points {\it de $X$}. Un aspect intéressant
de la théorie des schémas est de trouver des descriptions de
la géométrie interne de $X$ en termes de ce foncteur $h_X$. Dans cette
section, nous trouvons une manière simple de décrire les
points de $X$.

\medskip\noindent
Soit $X$ un schéma. Soit $R$ un anneau local avec un idéal maximal
$\mathfrak m \subset R$. Supposons que $f : \Spec(R) \to
X$ soit un morphisme de schémas. Soit $x \in X$ l'image du point fermé
$\mathfrak m \in \Spec(R)$. Alors nous obtenons un homomorphisme
local d'anneaux
$$
f^\sharp :
\mathcal{O}_{X, x}
\longrightarrow
\mathcal{O}_{\Spec(R), \mathfrak m} = R.
$$

\begin{lemma}
\label{schemes-lemma-morphism-from-spec-local-ring}
Soit $X$ un schéma. Soit $R$ un anneau local. La
construction ci-dessus donne une correspondance bijective
entre les morphismes $\Spec(R) \to X$ et les
paires $(x, \varphi)$ constituées d'un point $x \in X$
et d'un homomorphisme local d'anneaux $\varphi : \mathcal{O}_{X, x} \to R$.
\end{lemma}

\begin{proof}
Soit $A$ un anneau. Pour tout homomorphisme d'anneaux $\psi : A \to
R$, il existe un idéal premier unique $\mathfrak p \subset A$
et une factorisation $A \to A_{\mathfrak p} \to R$ où la dernière
application est un homomorphisme local d'anneaux. À savoir,
$\mathfrak p = \psi^{-1}(\mathfrak m)$. Via
le lemme \ref{schemes-lemma-morphism-into-affine},
cela prouve que le lemme est vrai si $X$ est un schéma affine.

\medskip\noindent
Soit $X$ un schéma général. Tout $x \in X$ est contenu dans
un ouvert affine $U \subset X$. Par le cas affine, nous concluons que
chaque paire $(x, \varphi)$ résulte de la construction qui
précède le lemme.

\medskip\noindent
Pour terminer la preuve, il suffit de montrer que tout morphisme
$f : \Spec(R) \to X$ a son image contenue dans tout
ouvert affine contenant l'image $x$ du point
fermé de $\Spec(R)$. En fait, soit $x \in V \subset X$
un voisinage ouvert quelconque contenant $x$.
Alors $f^{-1}(V) \subset \Spec(R)$ est un ouvert
contenant le point fermé unique et donc égal à $\Spec(R)$.
\end{proof}

\noindent
Comme cas particulier du lemme ci-dessus, nous obtenons pour tout
point $x$ d'un schéma $X$ un morphisme canonique
\begin{equation}
\label{schemes-equation-canonical-morphism}
\Spec(\mathcal{O}_{X, x}) \longrightarrow X
\end{equation}
correspondant à l'application identité sur l'anneau local de $X$ en $x$.
On peut reformuler le lemme ci-dessus en disant que pour tout
morphisme $f : \Spec(R) \to X$, il existe un point unique $x
\in X$ tel que $f$ se factorise comme
$\Spec(R) \to \Spec(\mathcal{O}_{X, x}) \to X$ où
la première application provient d'un homomorphisme local
$\mathcal{O}_{X, x} \to R$.

\medskip\noindent
Dans le cas où nous avons un morphisme de schémas $f : X \to S$,
et un point $x$ dont l'image est un point $s \in S$, nous obtenons
un diagramme commutatif
$$
\xymatrix{
\Spec(\mathcal{O}_{X, x}) \ar[r] \ar[d] & X \ar[d] \\
\Spec(\mathcal{O}_{S, s}) \ar[r] & S
}
$$
où l'application verticale gauche correspond à l'homomorphisme local
$f^\sharp_x : \mathcal{O}_{S, s} \to \mathcal{O}_{X, x}$.

\begin{lemma}
\label{schemes-lemma-specialize-points}
Soit $X$ un schéma.
Soient $x, x' \in X$ des points de $X$.
Alors $x' \in X$ est une générisation de $x$ si et seulement
si $x'$ est dans l'image du morphisme canonique
$\Spec(\mathcal{O}_{X, x}) \to X$.
\end{lemma}

\begin{proof}
Une application continue préserve la relation de spécialisation/générisation.
Puisque chaque point de $\Spec(\mathcal{O}_{X, x})$ est une
générisation du point fermé, nous voyons que chaque point dans l'image
de $\Spec(\mathcal{O}_{X, x}) \to X$ est une générisation de $x$.
Inversement, supposons que $x'$ soit une générisation de $x$.
Choisissez un voisinage ouvert affine $U = \Spec(R)$ de $x$.
Alors $x' \in U$. Disons que $\mathfrak p \subset R$ et
$\mathfrak p' \subset R$ sont les idéaux premiers
correspondant à $x$ et $x'$. Puisque $x'$ est une générisation
de $x$, nous voyons que $\mathfrak p' \subset \mathfrak p$. Cela
signifie que $\mathfrak p'$ est dans l'image du morphisme
$\Spec(\mathcal{O}_{X, x}) = \Spec(R_{\mathfrak p})
\to \Spec(R) = U \subset X$ comme voulu.
\end{proof}

\noindent
Maintenant, discutons des morphismes provenant des spectres des
corps. Soit $(R, \mathfrak m, \kappa)$ un anneau local
avec un idéal maximal $\mathfrak m$ et un corps résiduel $\kappa$.
Soit $K$ un corps. Un homomorphisme local $R \to K$, par définition, se
factorise comme $R \to \kappa \to K$, c'est-à-dire qu'il revient à un
morphisme $\kappa \to K$. Ainsi, nous voyons que les morphismes
$$
\Spec(K) \longrightarrow X
$$
correspondent à des paires $(x, \kappa(x) \to K)$. Nous pouvons
définir un préordre sur les morphismes des spectres
des corps vers $X$ en disant que $\Spec(K) \to
X$ domine $\Spec(L) \to X$ si $\Spec(K)
\to X$ se factorise par $\Spec(L) \to X$. Cela
suggère la notion suivante : disons temporairement
que deux morphismes $p : \Spec(K) \to X$ et
$q : \Spec(L) \to X$ sont {\it équivalents} s'il existe un
troisième corps $\Omega$ et un diagramme commutatif
$$
\xymatrix{
\Spec(\Omega) \ar[r] \ar[d] &
\Spec(L) \ar[d]^q \\
\Spec(K) \ar[r]^p &
X
}
$$
Bien sûr, cela implique immédiatement que les points uniques des
trois schémas $\Spec(K)$, $\Spec(L)$ et
$\Spec(\Omega)$ se projettent sur
le même $x \in X$. Ainsi, un diagramme (d'après les remarques ci-dessus)
correspond à un point $x \in X$ et à un diagramme commutatif
$$
\xymatrix{
\Omega &
L \ar[l] \\
K \ar[u] &
\kappa(x) \ar[l] \ar[u]
}
$$
de corps. Cela définit une relation d'équivalence, car étant donné un
ensemble d'extensions de corps $K_i/\kappa$, il existe
une extension de corps $\Omega/\kappa$ telle que toutes
les extensions de corps $K_i$ soient contenues dans l'extension $\Omega$.

\begin{lemma}
\label{schemes-lemma-characterize-points}
Soit $X$ un schéma. Les points de $X$ correspondent bijectivement
aux classes d'équivalence de morphismes des spectres
de corps dans $X$. De plus, chaque classe d'équivalence contient
un élément minimal (unique à isomorphisme unique près)
$\Spec(\kappa(x)) \to X$.
\end{lemma}

\begin{proof}
Cela découle de la discussion ci-dessus.
\end{proof}

\noindent
Bien sûr, les morphismes $\Spec(\kappa(x)) \to X$ se
factorisent par les morphismes canoniques
$\Spec(\mathcal{O}_{X, x}) \to X$. Et le
contenu du lemme \ref{schemes-lemma-specialize-points} est, dans ce cadre,
que le morphisme $\Spec(\kappa(x')) \to X$ se factorise comme
$\Spec(\kappa(x'))
\to \Spec(\mathcal{O}_{X, x}) \to X$ chaque fois que
$x'$ est une générisation de $x$. Dans le cas où
nous avons un morphisme de schémas $f : X \to S$, et un
point $x$ dont l'image est un point $s \in S$, nous obtenons
un diagramme commutatif
$$
\xymatrix{
\Spec(\kappa(x)) \ar[r] \ar[d] &
\Spec(\mathcal{O}_{X, x}) \ar[r] \ar[d] & X
\ar[d] \\
\Spec(\kappa(s)) \ar[r] &
\Spec(\mathcal{O}_{S, s}) \ar[r] &
S.
}
$$










\section{Recollement de schémas}
\label{schemes-section-glueing-schemes}

\noindent
Soit $I$ un ensemble.
Pour chaque $i \in I$, soit $(X_i, \mathcal{O}_i)$ un
espace annelé en anneaux locaux. (En réalité, la construction qui
suit fonctionne tout aussi bien pour les espaces annelés.)
Pour chaque paire $i, j \in I$ soit $U_{ij} \subset X_i$
un sous-espace ouvert.
Pour chaque paire $i, j \in I$, soit
$$
\varphi_{ij} : U_{ij} \to U_{ji}
$$
un isomorphisme d'espaces annelés en anneaux locaux.
Pour plus de commodité, nous supposons que $U_{ii} = X_i$.
Pour chaque triple $i, j, k \in I$, supposons que
\begin{enumerate}
\item nous avons
$\varphi_{ij}^{-1}(U_{ji} \cap U_{jk}) = U_{ij} \cap U_{ik}$, et
\item le diagramme
$$
\xymatrix{
U_{ij} \cap U_{ik} \ar[rr]_{\varphi_{ik}} \ar[rd]_{\varphi_{ij}} & &
U_{ki} \cap U_{kj} \\
& U_{ji} \cap U_{jk} \ar[ru]_{\varphi_{jk}}
}
$$
est commutatif.
\end{enumerate}
(En particulier, il en découle que $\varphi_{ii} = \text{id}_{X_i}$.)
Appelons une collection $(I,
(X_i)_{i\in I}, (U_{ij})_{i, j\in I}, (\varphi_{ij})_{i, j\in I})$ satisfaisant
les conditions ci-dessus des données de recollement.

\begin{lemma}
\label{schemes-lemma-glue}
\begin{slogan}
Si deux espaces annelés en anneaux locaux possèdent des sous-espaces
isomorphes, alors on peut les recoller pour former un grand
espace annelé en anneaux locaux.
\end{slogan}
Étant donnée une famille de données de recollement d'espaces annelés en
anneaux locaux, il existe un espace annelé en anneaux locaux $X$ et
des sous-espaces ouverts $U_i \subset X$ avec des
isomorphismes $\varphi_i : X_i \to U_i$ d'espaces annelés en anneaux locaux tels que
\begin{enumerate}
\item $X=\bigcup_{i\in I} U_i$,
\item $\varphi_i(U_{ij}) = U_i \cap U_j$, et
\item $\varphi_{ij} =
\varphi_j^{-1}|_{U_i \cap U_j} \circ \varphi_i|_{U_{ij}}$.
\end{enumerate}
L'espace annelé en anneaux locaux $X$ est caractérisé par les propriétés universelles
suivantes : étant donné un espace annelé en anneaux locaux $Y$, nous avons
\begin{eqnarray*}
\Mor(X, Y) & = & \{ (f_i)_{i\in I} \mid
f_i : X_i \to Y, \ f_j \circ \varphi_{ij} = f_i|_{U_{ij}}\} \\ f
& \mapsto & (f|_{U_i} \circ \varphi_i)_{i \in I} \\
\Mor(Y, X) & = &
\left\{
\begin{matrix}
\text{recouvrement ouvert }Y = \bigcup\nolimits_{i \in I} V_i\text{ et
} (g_i : V_i \to X_i)_{i \in
I} \text{ tels que}\\
g_i^{-1}(U_{ij}) = V_i \cap V_j
\text{ et
} g_j|_{V_i \cap V_j} = \varphi_{ij} \circ g_i|_{V_i \cap V_j}
\end{matrix}
\right\} \\
g & \mapsto &
V_i = g^{-1}(U_i), \ g_i = \varphi_i^{-1} \circ g|_{V_i}
\end{eqnarray*}
\end{lemma}

\begin{proof}
Nous construisons $X$ en plusieurs étapes.
Comme ensemble, nous prenons
$$
X = (\coprod X_i) / \sim.
$$
Ici, étant donné $x \in X_i$ et $x' \in X_j$, nous
disons $x \sim x'$ si et seulement si $x \in U_{ij}$, $x' \in
U_{ji}$ et $\varphi_{ij}(x) = x'$. Il s'agit d'une relation
d'équivalence car si $x \in X_i$, $x' \in X_j$, $x'' \in X_k$, $x \sim x'$ et $x'
\sim x''$, alors $x' \in U_{ji} \cap U_{jk}$, donc, d'après la condition (1) des
données de recollement, $x \in U_{ij} \cap U_{ik}$ et $x''
\in U_{ki} \cap U_{kj}$ ; d'après la condition (2),
nous voyons que $\varphi_{ik}(x) = x''$. (La réflexivité et la symétrie
découlent de nos hypothèses que $U_{ii} = X_i$ et
$\varphi_{ii} = \text{id}_{X_i}$.) On
note $\varphi_i : X_i \to X$ les
applications naturelles. On note $U_i = \varphi_i(X_i) \subset X$.
Notez que $\varphi_i : X_i \to U_i$ est une bijection.

\medskip\noindent
La topologie sur $X$ est définie par la règle selon
laquelle $U \subset X$ est ouvert si et seulement si $\varphi_i^{-1}(U)$
est ouvert pour tout $i$. Nous laissons au lecteur le soin
de vérifier que cela définit effectivement une
topologie. Notez qu'en particulier $U_i$ est ouvert puisque $\varphi_j^{-1}(U_i)
= U_{ji}$ qui est ouvert dans $X_j$ pour tout
$j$. De plus, pour tout ensemble ouvert $W \subset X_i$,
l'image $\varphi_i(W) \subset U_i$ est ouverte
parce que $\varphi_j^{-1}(\varphi_i(W)) = \varphi_{ji}^{-1}(W \cap U_{ij})$.
Par conséquent, $\varphi_i : X_i \to U_i$ est un homéomorphisme.

\medskip\noindent
Pour obtenir un espace annelé en anneaux locaux, nous devons construire
le faisceau d'anneaux $\mathcal{O}_X$. Nous le faisons en recollant
les faisceaux d'anneaux $\mathcal{O}_{U_i} := \varphi_{i, *} \mathcal{O}_i$. Plus
précisément, dans le diagramme commutatif
$$
\xymatrix{
U_{ij} \ar[rr]_{\varphi_{ij}} \ar[rd]_{\varphi_i|_{U_{ij}}} &
&
U_{ji} \ar[ld]^{\varphi_j|_{U_{ji}}} \\ &
U_i \cap U_j &
}
$$
la flèche du haut est un isomorphisme d'espaces annelés et, par
conséquent, nous obtenons des isomorphismes uniques de faisceaux d'anneaux
$$
\mathcal{O}_{U_i}|_{U_i \cap U_j}
\longrightarrow
\mathcal{O}_{U_j}|_{U_i \cap U_j}.
$$
Ces isomorphismes satisfont une condition de cocycle comme
dans Faisceaux, section \ref{sheaves-section-glueing-sheaves}.
D'après les résultats de cette section, nous obtenons un faisceau
d'anneaux $\mathcal{O}_X$ sur $X$ tel que $\mathcal{O}_X|_{U_i}$
soit isomorphe à $\mathcal{O}_{U_i}$ de manière compatible
avec les morphismes de recollement affichés
ci-dessus. En particulier, $(X, \mathcal{O}_X)$ est un espace
annelé en anneaux locaux puisque les fibres de $\mathcal{O}_X$
sont égales aux fibres de $\mathcal{O}_i$ aux points
correspondants.

\medskip\noindent
La démonstration des propriétés universelles est omise.
\end{proof}

\begin{lemma}
\label{schemes-lemma-glue-schemes}
\begin{slogan}
Les schémas peuvent être recollés pour donner de nouveaux schémas.
\end{slogan}
Dans le lemme \ref{schemes-lemma-glue} ci-dessus, supposons que tous les
$X_i$ soient des schémas. Alors l'espace annelé en anneaux locaux
résultant $X$ est un schéma.
\end{lemma}

\begin{proof}
Cela est clair puisque chacun des $U_i$ est un schéma
et donc chaque $x \in X$ possède un voisinage affine.
\end{proof}

\noindent
Il est d'usage de considérer $X_i$ comme un sous-espace ouvert de
$X$ via les isomorphismes $\varphi_i$. Nous ferons cela dans les
deux exemples suivants.

\begin{example}[Espace affine à origine dédoublée]
\label{schemes-example-affine-space-zero-doubled}
Soit $k$ un corps. Soit $n \geq 1$.
Soit $X_1 = \Spec(k[x_1, \ldots, x_n])$, soit
$X_2 = \Spec(k[y_1, \ldots, y_n])$. Soit
$0_1 \in X_1$ le point correspondant à l'idéal maximal $(x_1, \ldots,
x_n) \subset k[x_1, \ldots, x_n]$. Soit $0_2 \in
X_2$ le point correspondant à l'idéal maximal $(y_1, \ldots, y_n)
\subset k[y_1, \ldots, y_n]$. Soit $U_{12} = X_1 \setminus
\{0_1\}$ et soit $U_{21} = X_2 \setminus
\{0_2\}$. Soit $\varphi_{12} : U_{12} \to
U_{21}$ l'isomorphisme provenant de l'isomorphisme des
$k$-algèbres $k[y_1, \ldots, y_n] \to k[x_1,
\ldots, x_n]$ envoyant $y_i$ sur $x_i$ (ce qui
induit $X_1 \cong X_2$ envoyant $0_1$ sur $0_2$). Soit $X$ le schéma
obtenu à partir
des données de recollement $(X_1, X_2, U_{12}, U_{21},
\varphi_{12}, \varphi_{21} = \varphi_{12}^{-1})$.
Par le léger abus de notation introduit ci-dessus dans
l'exemple, nous considérons $X_1, X_2 \subset X$ comme
des sous-schémas ouverts. Il existe un
morphisme $f : X \to \Spec(k[t_1, \ldots, t_n])$ qui sur $X_1$
(resp. \ $X_2$) correspond à l'homomorphisme de $k$-algèbres
$k[t_1, \ldots, t_n] \to k[x_1, \ldots, x_n]$ (resp.
\ $k[t_1, \ldots, t_n] \to k[y_1, \ldots, y_n]$) envoyant
$t_i$ vers $x_i$ (resp. \ $t_i$ vers $y_i$). Il
est facile de voir que ce morphisme identifie $k[t_1,
\ldots, t_n]$ avec $\Gamma(X, \mathcal{O}_X)$. Puisque $f(0_1)
= f(0_2)$, nous voyons que $X$ n'est pas affine.

\medskip\noindent
Notez que $X_1$ et $X_2$ sont des ouverts affines de
$X$. Mais, si $n = 2$, alors $X_1 \cap X_2$ est le schéma
décrit dans l'exemple \ref{schemes-example-not-affine} et donc pas affine.
Ainsi, en général, l'intersection des ouverts affines d'un schéma
n'est pas affine. (Ce fait vaut plus généralement pour tout $n > 1$.)

\medskip\noindent
Une autre caractéristique curieuse de cet exemple est la suivante.
Si $n > 1$ il y a de nombreux sous-ensembles fermés irréductibles $T \subset
X$ (prenez par exemple l'adhérence de tout point non fermé dans $X_1$).
Mais à moins que $T = \{0_1\}$, ou $T = \{0_2\}$ nous avons
$0_1 \in T \Leftrightarrow 0_2 \in T$. Démonstration omise.
\end{example}

\begin{example}[Droite projective]
\label{schemes-example-projective-line}
Soit $k$ un corps.
Soit $X_1 = \Spec(k[x])$,
soit $X_2 = \Spec(k[y])$.
Soit $0 \in X_1$ le point correspondant à l'idéal maximal $(x) \subset
k[x]$. Soit $\infty
\in X_2$ le point correspondant à l'idéal maximal $(y) \subset k[y]$.
Soit $U_{12} = X_1
\setminus \{0\} = D(x) = \Spec(k[x, 1/x])$ et soit $U_{21} = X_2
\setminus \{\infty\} = D(y) = \Spec(k[y, 1/y])$. Soit $\varphi_{12} :
U_{12} \to U_{21}$ l'isomorphisme provenant de l'isomorphisme
des $k$-algèbres $k[y, 1/y] \to k[x, 1/x]$
envoyant $y$ sur $1/x$. Soit $\mathbf{P}^1_k$ le
schéma obtenu à partir des données de recollement $(X_1, X_2, U_{12},
U_{21}, \varphi_{12}, \varphi_{21} =
\varphi_{12}^{-1})$. Via le léger abus de notation introduit ci-dessus
dans l'exemple, nous considérons $X_i \subset
\mathbf{P}^1_k$ comme des sous-schémas ouverts. Dans ce cas, nous
voyons que $\Gamma(\mathbf{P}^1_k, \mathcal{O}) = k$ car les seuls
polynômes $g(x)$ en $x$ tels que $g(1/y)$ soit également un
polynôme en $y$ sont des polynômes constants. Puisque
$\mathbf{P}^1_k$ est infini, nous voyons que $\mathbf{P}^1_k$ n'est pas affine.

\medskip\noindent
Nous affirmons qu'il existe un ouvert affine $U \subset \mathbf{P}^1_k$
qui contient à la fois $0$ et $\infty$. À savoir,
soit $U = \mathbf{P}^1_k \setminus \{1\}$, où $1$ est le point
de $X_1$ correspondant à l'idéal maximal $(x - 1)$ et
également le point de $X_2$ correspondant à l'idéal
maximal $(y - 1)$. Il est alors facile de voir que
$s = 1/(x - 1) = y/(1 - y) \in \Gamma(U, \mathcal{O}_U)$. En
fait, vous pouvez montrer que $\Gamma(U, \mathcal{O}_U)$
est égal à l'anneau de polynômes $k[s]$ et que le
morphisme correspondant $U \to \Spec(k[s])$
est un isomorphisme de schémas. Détails omis.
\end{example}




\section{Un critère de représentabilité}
\label{schemes-section-representable}

\noindent
Dans cette section, nous reformulons le lemme de recollement
de la section \ref{schemes-section-glueing-schemes} en termes
de foncteurs. Nous rappelons certains résultats de
Catégories, section \ref{categories-section-opposite}. Rappelons
qu'étant donné un schéma $X$, nous pouvons
définir un foncteur
$$
h_X : \Sch^{opp}
\longrightarrow
\textit{Ensembles}, \quad
T \longmapsto \Mor(T, X).
$$
On l'appelle le foncteur des points {\it de $X$}.

\medskip\noindent
Soit $F$ un foncteur contravariant de la catégorie des schémas
vers la catégorie des ensembles. Dans une formule
$$
F : \Sch^{opp}
\longrightarrow
\textit{Ensembles}.
$$
Nous utiliserons la même terminologie que dans Sites,
section \ref{sites-section-presheaves}. À savoir, étant donné
un schéma $T$, un élément $\xi \in F(T)$ et un morphisme $f : T'
\to T$, nous noterons $f^*\xi$ l'élément $F(f)(\xi)$, et parfois
nous utiliserons même la notation $\xi|_{T'}$

\begin{definition}
\label{schemes-definition-representable-functor}
(Voir Catégories, définition \ref{categories-definition-representable-functor}.) Soit
$F$ un foncteur contravariant de la catégorie des schémas
vers la catégorie des ensembles (comme ci-dessus).
Nous disons que $F$ est {\it représentable par un
schéma} ou {\it représentable} s'il existe un schéma
$X$ tel que $h_X \cong F$.
\end{definition}

\noindent
Supposons que $F$ soit représentable par le schéma $X$ et que $s : h_X
\to F$ soit un isomorphisme. D'après le
lemme de Yoneda (Catégories, lemme \ref{categories-lemma-yoneda}),
la paire $(X, s : h_X \to F)$ est unique à isomorphisme
unique près si elle existe. De
plus, le lemme de Yoneda affirme que pour
tout foncteur contravariant $F$ comme ci-dessus
et tout schéma $Y$, nous avons une bijection
$$
\Mor_{\text{Fun}(\Sch^{opp}, \textit{Ensembles})} (h_Y, F)
\longrightarrow
F(Y), \quad
s \longmapsto s(\text{id}_Y).
$$
Voici la construction inverse. Étant donné un $\xi \in F(Y)$,
la transformation de foncteurs $s_\xi : h_Y \to
F$ associe à tout morphisme $f : T \to Y$
l'élément $f^*\xi \in F(T)$.

\medskip\noindent
En particulier, dans le cas où $F$ est représentable, il existe un schéma
$X$ et un élément $\xi \in F(X)$ tels que le morphisme correspondant
$h_X \to F$ soit un isomorphisme. Dans ce
cas, nous disons aussi {\it que la paire $(X, \xi)$ représente $F$}.
L'élément $\xi \in F(X)$ est
souvent appelé la {\it ``famille universelle''} pour des raisons qui deviendront
plus claires lorsque nous parlerons des champs algébriques (insérer référence
future ici). Pour le moment, observons simplement que, si la paire $(X, \xi)$
représente $F$, alors tout élément $\xi' \in F(T)$ pour tout $T$ est de
la forme $\xi' = f^*\xi$ pour un morphisme unique $f : T \to X$.

\begin{example}
\label{schemes-example-global-sections}
Considérons la règle qui associe à chaque schéma $T$ l'ensemble $F(T) =
\Gamma(T, \mathcal{O}_T)$. Nous pouvons en faire un foncteur contravariant en
utilisant, pour un morphisme $f : T' \to T$, l'homomorphisme d'image inverse
$f^\sharp : \Gamma(T, \mathcal{O}_T) \to \Gamma(T', \mathcal{O}_{T'})$. Étant
donné un anneau $R$ et un élément $t \in R$, il existe un unique homomorphisme
d'anneaux $\mathbf{Z}[x] \to R$ qui envoie $x$ sur $t$. Ainsi, en
utilisant le lemme \ref{schemes-lemma-morphism-into-affine}, nous constatons que
$$
\Mor(T, \Spec(\mathbf{Z}[x])) =
\Hom(\mathbf{Z}[x], \Gamma(T, \mathcal{O}_T)) =
\Gamma(T, \mathcal{O}_T).
$$
Cela donne effectivement
un isomorphisme $h_{\Spec(\mathbf{Z}[x])} \to F$.
Quelle est la « famille universelle » $\xi$ ? Pour l'obtenir, nous
devons appliquer les identifications ci-dessus à $\text{id}_{\Spec(\mathbf{Z}[x])}$.
Clairement, sous les identifications ci-dessus,
cela
donne que $\xi = x \in \Gamma(\Spec(\mathbf{Z}[x]),
\mathcal{O}_{\Spec(\mathbf{Z}[x])}) = \mathbf{Z}[x]$
comme prévu.
\end{example}

\begin{definition}
\label{schemes-definition-representable-by-open-immersions}
Soit $F$ un foncteur contravariant sur la catégorie des
schémas à valeurs dans les ensembles.
\begin{enumerate}
\item On dit que $F$ {\it satisfait la propriété de faisceau pour la
topologie de Zariski} si pour tout schéma $T$ et tout recouvrement ouvert
$T = \bigcup_{i \in I} U_i$, et pour toute collection d'éléments
$\xi_i \in F(U_i)$ telle que $\xi_i|_{U_i \cap U_j} =
\xi_j|_{U_i \cap U_j}$, il existe un élément unique $\xi \in
F(T)$ tel que $\xi_i = \xi|_{U_i}$ dans $F(U_i)$. \item Un {\it
sous-foncteur $H \subset F$} est une règle qui associe à chaque schéma
$T$ un sous-ensemble $H(T) \subset F(T)$ tel que les applications
$F(f) : F(T) \to F(T')$ envoient $H(T)$ dans $H(T')$
pour tous les morphismes de schémas $f : T' \to T$.
\item Soit $H \subset F$ un sous-foncteur. On dit que $H
\subset F$ est {\it représentable par des immersions ouvertes}
si pour tous les couples $(T, \xi)$, où $T$ est un schéma et $\xi \in
F(T)$, il existe un sous-schéma ouvert $U_\xi \subset T$ avec la propriété
suivante :
\begin{itemize}
\item[(*)] Un morphisme $f : T' \to T$ se factorise à travers $U_\xi$ si et seulement
si $f^*\xi \in H(T')$.
\end{itemize}
\item Soit $I$ un ensemble. Pour chaque $i \in I$, soit $H_i \subset
F$ un sous-foncteur. Nous disons que la collection $(H_i)_{i \in I}$
{\it recouvre $F$} si et seulement si pour chaque $\xi \in
F(T)$, il existe un recouvrement ouvert $T = \bigcup U_i$ tel que
$\xi|_{U_i} \in H_i(U_i)$.
\end{enumerate}
\end{definition}

\noindent
Dans la condition (4), si $H_i \subset F$ est représentable par des immersions ouvertes
pour tous les $i$, alors pour vérifier que $(H_i)_{i \in I}$ recouvre $F$, il suffit
de vérifier $F(T) = \bigcup H_i(T)$ chaque fois que $T$ est le spectre d'un corps.

\begin{lemma}
\label{schemes-lemma-glue-functors}
Soit $F$ un foncteur contravariant sur la catégorie des schémas à
valeurs dans la catégorie des ensembles. Supposons que
\begin{enumerate}
\item $F$ satisfasse la propriété de faisceau pour la topologie de Zariski.
\item il existe un ensemble $I$ et une collection de sous-foncteurs
$F_i \subset F$ tels que
\begin{enumerate}
\item chaque $F_i$ est représentable,
\item chaque $F_i \subset F$ est représentable par des immersions ouvertes,
et \item la collection $(F_i)_{i \in I}$ recouvre $F$.
\end{enumerate}
\end{enumerate}
Alors $F$ est représentable.
\end{lemma}

\begin{proof}
Soit $X_i$ un schéma représentant $F_i$ et soit $\xi_i
\in F_i(X_i) \subset F(X_i)$ la ``famille universelle''. Puisque $F_j
\subset F$ est représentable par des immersions ouvertes, il
existe un ouvert $U_{ij} \subset X_i$ tel que $T \to X_i$
se factorise par $U_{ij}$ si et seulement si $\xi_i|_T
\in F_j(T)$. En particulier $\xi_i|_{U_{ij}}
\in F_j(U_{ij})$ et donc nous obtenons un morphisme canonique
$\varphi_{ij} : U_{ij} \to X_j$ tel que $\varphi_{ij}^*\xi_j =
\xi_i|_{U_{ij}}$. Par définition de $U_{ji}$, cela implique que $\varphi_{ij}$
se factorise par $U_{ji}$. Comme $(\varphi_{ij} \circ
\varphi_{ji})^*\xi_j =\varphi_{ji}^*(\varphi_{ij}^*\xi_j)
= \varphi_{ji}^*\xi_i = \xi_j$, nous
concluons que $\varphi_{ij} \circ \varphi_{ji} =
\text{id}_{U_{ji}}$ parce que le couple $(X_j, \xi_j)$ représente
$F_j$. En particulier, les applications $\varphi_{ij}
: U_{ij} \to U_{ji}$ sont des isomorphismes de schémas.
Ensuite, nous devons montrer que
$\varphi_{ij}^{-1}(U_{ji}
\cap U_{jk}) = U_{ij} \cap U_{ik}$. Cela est vrai parce que (a)
$U_{ji} \cap U_{jk}$ est le plus grand ouvert de $U_{ji}$ tel que
$\xi_j$ se restreigne à un élément de $F_k$, (b) $U_{ij} \cap U_{ik}$
est le plus grand ouvert de $U_{ij}$ tel que $\xi_i$
se restreigne à un élément de $F_k$, et (c) $\varphi_{ij}^*\xi_j
= \xi_i$. De plus, la condition de cocycle dans la section
\ref{schemes-section-glueing-schemes} suit parce que $\varphi_{jk}|_{U_{ji}
\cap U_{jk}} \circ
\varphi_{ij}|_{U_{ij} \cap U_{ik}}$ et
$\varphi_{ik}|_{U_{ij} \cap U_{ik}}$ ont tous deux
pour image inverse de $\xi_k$ l'élément $\xi_i$. Ainsi,
nous pouvons appliquer le
lemme \ref{schemes-lemma-glue-schemes} pour obtenir un schéma
$X$ avec un recouvrement ouvert $X =
\bigcup U_i$ et des isomorphismes $\varphi_i : X_i \to U_i$ avec des propriétés
comme dans le lemme \ref{schemes-lemma-glue}. Soit
$\xi_i' = (\varphi_i^{-1})^* \xi_i$. Les conditions
du lemme \ref{schemes-lemma-glue} impliquent que $\xi_i'|_{U_i
\cap U_j} = \xi_j'|_{U_i \cap U_j}$. Par conséquent,
d'après la condition que $F$ satisfait la condition de
faisceau dans la topologie de Zariski, nous voyons
qu'il existe un élément $\xi' \in F(X)$ tel que $\xi_i
= \varphi_i^*\xi'|_{U_i}$ pour tout $i$. Puisque
$\varphi_i$ est un isomorphisme, nous obtenons également
que $(U_i, \xi'|_{U_i})$ représente le foncteur $F_i$.

\medskip\noindent
Nous affirmons que la paire $(X, \xi')$ représente le foncteur $F$. Pour
le montrer, soit $T$ un schéma et soit $\xi \in F(T)$. Nous
allons construire un morphisme unique $g : T \to X$ tel que $g^*\xi'
= \xi$. À savoir, d'après la condition que les sous-foncteurs $F_i$
couvrent $F$, il existe un recouvrement ouvert $T = \bigcup V_i$ tel que
pour chaque $i$, la restriction $\xi|_{V_i} \in F_i(V_i)$. De plus, comme
chacune des inclusions $F_i \subset F$ est représentable par des immersions ouvertes,
nous pouvons supposer que chaque $V_i \subset T$ est le plus grand ouvert
possédant cette propriété.
Puisque $(U_i, \xi'|_{U_i})$ représente le foncteur $F_i$, nous obtenons
un morphisme unique $g_i : V_i \to U_i$ tel que
$g_i^*\xi'|_{U_i} = \xi|_{V_i}$. Sur les intersections $V_i \cap V_j$, les
morphismes $g_i$ et $g_j$ coïncident, par exemple parce qu'ils ont tous deux
pour image inverse de $\xi'|_{U_i \cap U_j} \in F_i(U_i \cap
U_j)$ le même élément. Ainsi les morphismes $g_i$ se recollent en un unique morphisme
$T \to X$ comme souhaité.
\end{proof}

\begin{remark}
\label{schemes-remark-representable-locally-ringed}
Supposons que le foncteur $F$ soit défini sur tous les espaces annelés en anneaux
locaux, et que les conditions du lemme \ref{schemes-lemma-glue-functors} soient remplacées
par ce qui suit :
\begin{enumerate}
\item $F$ satisfait la propriété de faisceau sur la catégorie des espaces annelés en
anneaux
locaux, \item il existe un ensemble $I$ et une collection de sous-foncteurs
$F_i \subset F$ tels que
\begin{enumerate}
\item chaque $F_i$ soit représentable par un schéma, \item
chaque $F_i \subset F$ soit représentable par des immersions ouvertes sur
la catégorie des espaces annelés en anneaux locaux, et
\item la collection $(F_i)_{i \in I}$ recouvre $F$ en tant que
foncteur sur la catégorie des espaces annelés en anneaux locaux.
\end{enumerate}
\end{enumerate}
Nous laissons au lecteur le soin de développer cela davantage.
Le résultat final est alors que le foncteur $F$ est
représentable dans la catégorie des espaces annelés en anneaux
locaux et que l'objet représentant est un schéma.
\end{remark}



\section{Existence des produits fibrés de schémas}
\label{schemes-section-existence-fibre-products}

\noindent
Une question élémentaire est de savoir si les produits et les
produits fibrés existent dans la catégorie des schémas. Nous
prouvons d'abord leur existence de manière abstraite ;
dans la section suivante, nous en donnons une
description concrète.

\begin{lemma}
\label{schemes-lemma-fibre-products}
La catégorie des schémas possède un objet final, des produits et des produits
fibrés. En d'autres termes, la catégorie des schémas a des limites finies,
voir Catégories, lemme \ref{categories-lemma-finite-limits-exist}.
\end{lemma}

\begin{proof}
On peut omettre cette démonstration. Il est plus important
d'apprendre à travailler avec le produit fibré, qui est expliqué dans
la section suivante.

\medskip\noindent
D'après le lemme \ref{schemes-lemma-morphism-into-affine},
le schéma $\Spec(\mathbf{Z})$ est un objet final
dans la catégorie des espaces annelés en anneaux locaux. Ainsi,
il suffit de prouver que les produits fibrés existent.

\medskip\noindent
Soient $f : X \to S$ et $g : Y \to S$ des morphismes de schémas. Nous devons montrer
que le foncteur
\begin{eqnarray*}
F : \Sch^{opp} & \longrightarrow & \textit{Ensembles}
\\ T & \longmapsto
& \Mor(T, X) \times_{\Mor(T, S)} \Mor(T, Y)
\end{eqnarray*}
est représentable. Nous affirmons que le lemme \ref{schemes-lemma-glue-functors}
s'applique au foncteur $F$. Si nous prouvons cela, alors le lemme est prouvé.

\medskip\noindent
Tout d'abord, nous montrons que $F$ satisfait la propriété de
faisceau pour la topologie de Zariski. À savoir, supposons que $T$
soit un schéma, $T = \bigcup_{i \in I} U_i$ soit un
recouvrement ouvert, et $\xi_i \in
F(U_i)$ tel que $\xi_i|_{U_i \cap U_j} = \xi_j|_{U_i \cap
U_j}$ pour toutes les paires $i, j$. Par définition, $\xi_i$
correspond à une paire $(a_i, b_i)$ où $a_i : U_i \to X$ et $b_i : U_i \to
Y$ sont des morphismes de schémas tels que $f \circ a_i = g \circ b_i$.
La condition de recollement dit
que $a_i|_{U_i \cap U_j} = a_j|_{U_i \cap U_j}$ et
$b_i|_{U_i
\cap U_j} = b_j|_{U_i \cap U_j}$. Ainsi,
en recollant les morphismes $a_i$, nous obtenons un morphisme
d'espaces annelés en anneaux locaux (c'est-à-dire un morphisme
de schémas) $a : T \to X$ et de même $b : T \to Y$ (voir par
exemple la propriété universelle du lemme \ref{schemes-lemma-glue}). De plus,
sur les éléments d'un recouvrement ouvert, les compositions
$f \circ a$ et $g \circ b$ coïncident. Par
conséquent, $f \circ a = g \circ b$ et la paire $(a, b)$ définit
un élément de $F(T)$ qui se restreint aux paires
$(a_i, b_i)$ sur chaque $U_i$. La condition de faisceau est vérifiée.

\medskip\noindent
Ensuite, nous construisons la famille de sous-foncteurs.
Choisissez un recouvrement ouvert par des
ouverts affines $S = \bigcup\nolimits_{i \in
I} U_i$. Pour chaque $i \in I$, choisissez des recouvrements
ouverts par des ouverts affines $f^{-1}(U_i) = \bigcup\nolimits_{j
\in J_i} V_j$ et $g^{-1}(U_i) = \bigcup\nolimits_{k \in
K_i} W_k$. Notez que $X = \bigcup_{i \in I} \bigcup_{j \in J_i}
V_j$ est un recouvrement ouvert et de même pour
$Y$. Pour tout $i \in I$ et chaque paire $(j, k) \in J_i \times K_i$,
nous avons un diagramme commutatif
$$
\xymatrix{
& W_k \ar[d] \ar[rd] & \\ V_j
\ar[rd] \ar[r] & U_i \ar[rd] & Y \ar[d]
\\ & X \ar[r] & S
}
$$
où toutes les flèches obliques sont des immersions ouvertes. Pour un tel
triple, nous obtenons un foncteur
\begin{eqnarray*}
F_{i, j, k} : \Sch^{opp} & \longrightarrow & \textit{Ensembles}
\\ T & \longmapsto
& \Mor(T, V_j) \times_{\Mor(T, U_i)} \Mor(T, W_k).
\end{eqnarray*}
Il existe une transformation évidente de foncteurs $F_{i, j, k} \to F$
(provenant du grand diagramme commutatif ci-dessus) qui est
injective, donc nous pouvons considérer $F_{i, j, k}$ comme un sous-foncteur
de $F$.

\medskip\noindent
Nous vérifions la condition (2)(a) du lemme \ref{schemes-lemma-glue-functors}. Cela
découle directement du lemme \ref{schemes-lemma-fibre-product-affine-schemes}. (Notez que nous
utilisons ici le fait que les produits fibrés dans la catégorie des schémas affines
sont également des produits fibrés dans la catégorie tout entière des espaces annelés en
anneaux locaux.)

\medskip\noindent
Nous vérifions la condition (2)(b) du lemme \ref{schemes-lemma-glue-functors}.
Soit $T$ un schéma et soit $\xi \in F(T)$. En d'autres termes,
$\xi = (a, b)$ où $a : T \to X$ et $b : T \to Y$ sont des
morphismes de schémas tels que $f \circ a = g \circ b$. Posons
$V_{i, j, k} = a^{-1}(V_j) \cap b^{-1}(W_k)$. Pour tout autre
morphisme $h : T' \to T$ nous avons $h^*\xi
= (a \circ h, b \circ h)$. Par conséquent, nous voyons que
$h^*\xi \in F_{i, j, k}(T')$ si et seulement si
$a(h(T')) \subset V_j$ et $b(h(T')) \subset W_k$. En
d'autres termes, si et seulement si $h(T') \subset V_{i, j, k}$.
Cela prouve la condition (2)(b).

\medskip\noindent
Nous vérifions la condition (2)(c) du lemme \ref{schemes-lemma-glue-functors}.
Soit $T$ un schéma et soit $\xi = (a, b) \in F(T)$ comme ci-dessus. Posons
$V_{i, j, k} = a^{-1}(V_j) \cap b^{-1}(W_k)$ comme ci-dessus. La
condition (2)(c) signifie simplement que $T = \bigcup V_{i, j, k}$ ce
qui est évident. Ainsi, le lemme est prouvé et les produits fibrés
existent.
\end{proof}

\begin{remark}
\label{schemes-remark-fibre-product-schemes-locally-ringed}
En utilisant la remarque \ref{schemes-remark-representable-locally-ringed},
on peut montrer que le produit fibré de morphismes de schémas existe
dans la catégorie des espaces annelés en anneaux locaux et est un
schéma.
\end{remark}




\section{Produits fibrés de schémas}
\label{schemes-section-fibre-products}

\noindent
Voici un rappel de la définition générale, même si nous
avons déjà montré que les produits fibrés de schémas existent.

\begin{definition}
\label{schemes-definition-fibre-product}
Étant donnés des morphismes de schémas $f : X \to S$ et $g : Y
\to S$, le {\it produit fibré} est un schéma $X \times_S Y$ muni
des morphismes de projection $p : X \times_S Y \to X$
et $q : X \times_S Y \to Y$ qui s'insèrent dans le diagramme
commutatif suivant
$$
\xymatrix{
X \times_S Y \ar[r]_q \ar[d]_p & Y \ar[d]^g \\ X
\ar[r]^f & S
}
$$
et qui est universel parmi tous les diagrammes de ce
type, voir Catégories, définition \ref{categories-definition-fibre-products}.
\end{definition}

\noindent
En d'autres termes, étant donné tout diagramme commutatif de morphismes de schémas dont
les flèches pleines sont fixées
$$
\xymatrix{ T
\ar[rrrd] \ar@{-->}[rrd] \ar[rrdd] & & \\
&
&
X
\times_S
Y
\ar[d] \ar[r] & Y
\ar[d]
\\
&
&
X
\ar[r]
&
S
}
$$
il existe une flèche pointillée unique faisant commuter le diagramme.
Nous allons prouver quelques lemmes qui nous aideront à comprendre les
produits fibrés.

\begin{lemma}
\label{schemes-lemma-fibre-product-affines}
Soient $f : X \to S$ et $g : Y \to S$ des morphismes de schémas
ayant la même cible. Si $X, Y, S$ sont tous affines alors
$X \times_S Y$ est affine.
\end{lemma}

\begin{proof}
Supposons que $X = \Spec(A)$, $Y = \Spec(B)$ et $S
= \Spec(R)$. D'après le lemme \ref{schemes-lemma-fibre-product-affine-schemes}, le
schéma affine $\Spec(A \otimes_R B)$ est le produit
fibré $X \times_S Y$ dans la catégorie des espaces
annelés en anneaux locaux. Par conséquent, il est a fortiori
le produit fibré dans la catégorie des schémas.
\end{proof}

\begin{lemma}
\label{schemes-lemma-open-fibre-product}
Soient $f : X \to S$ et $g : Y \to S$ des morphismes de schémas
ayant la même cible. Soient $X \times_S Y$, $p$, $q$ le produit fibré.
Supposons que $U \subset S$, $V
\subset X$, $W \subset Y$ soient des sous-schémas
ouverts tels que $f(V) \subset U$ et $g(W) \subset
U$. Alors le morphisme canonique
$V \times_U W \to X \times_S Y$ est une immersion
ouverte qui identifie $V \times_U W$ avec $p^{-1}(V) \cap q^{-1}(W)$.
\end{lemma}

\begin{proof}
Soit $T$ un schéma.
Supposons que $a : T \to V$ et $b : T \to W$ soient des
morphismes tels que $f \circ a = g \circ b$ comme morphismes dans
$U$. Alors ils coïncident comme morphismes
dans $S$. Par la propriété universelle du produit fibré,
nous obtenons un morphisme unique $T \to X \times_S Y$. Bien sûr, ce
morphisme a son image contenue dans l'ouvert $p^{-1}(V) \cap
q^{-1}(W)$. Ainsi $p^{-1}(V) \cap q^{-1}(W)$ est un produit
fibré de $V$ et $W$ au-dessus de $U$. Le résultat découle de l'unicité
des produits fibrés, voir Catégories, section
\ref{categories-section-fibre-products}.
\end{proof}

\noindent
En particulier, cela montre que $V \times_U W = V \times_S W$ dans
la situation du lemme. De plus, si $U, V, W$ sont tous affines, alors
nous savons que $V \times_U W$ est affine. Et bien sûr, nous pouvons
recouvrir $X \times_S Y$ par de tels ouverts affines $V \times_U W$. Nous
formulons cela comme un lemme.

\begin{lemma}
\label{schemes-lemma-affine-covering-fibre-product}
\begin{slogan}
Construction explicite des produits fibrés : un recouvrement ouvert affine d'un
produit fibré de schémas peut être assemblé à partir de recouvrements
ouverts affines compatibles des facteurs.
\end{slogan}
Soient $f : X \to S$ et $g : Y \to S$ des morphismes de schémas
ayant la même cible. Soit $S = \bigcup U_i$ un recouvrement ouvert
affine quelconque de $S$. Pour chaque $i \in I$,
soit $f^{-1}(U_i) = \bigcup_{j \in J_i} V_j$ un recouvrement ouvert
affine de $f^{-1}(U_i)$ et
soit $g^{-1}(U_i) = \bigcup_{k \in K_i} W_k$ un recouvrement ouvert affine
de $g^{-1}(U_i)$. Alors
$$
X \times_S Y =
\bigcup\nolimits_{i \in I}
\bigcup\nolimits_{j \in J_i, \ k \in K_i}
V_j \times_{U_i} W_k
$$
est un recouvrement ouvert affine de $X \times_S Y$.
\end{lemma}

\begin{proof}
Voir la discussion qui précède le lemme.
\end{proof}

\noindent
En d'autres termes, nous aurions pu utiliser le lemme précédent pour construire
le produit fibré directement
en recollant les schémas affines. (Ce qui est bien sûr exactement
ce que nous avons fait dans la preuve du lemme \ref{schemes-lemma-fibre-products}.)
Voici une manière de décrire l'ensemble des points d'un produit fibré de schémas.

\begin{lemma}
\label{schemes-lemma-points-fibre-product}
Soient $f : X \to S$ et $g : Y \to S$ des morphismes de schémas
ayant la même cible. Les points $z$ de $X \times_S Y$ sont en correspondance
bijective avec les quadruples
$$
(x, y, s, \mathfrak p)
$$
où $x \in X$, $y \in Y$, $s \in S$ sont des points tels
que $f(x) = s$, $g(y) = s$ et $\mathfrak p$ est un idéal premier
de l'anneau $\kappa(x) \otimes_{\kappa(s)} \kappa(y)$.
Le corps résiduel de $z$ correspond au
corps résiduel de l'idéal premier $\mathfrak p$.
\end{lemma}

\begin{proof}
Soit $z$ un point de $X \times_S Y$ et construisons un quadruple comme
ci-dessus. Rappelons que nous pouvons considérer $z$ comme un morphisme
$\Spec(\kappa(z)) \to X \times_S Y$, voir le
lemme \ref{schemes-lemma-characterize-points}. Ce morphisme correspond aux
morphismes $a : \Spec(\kappa(z)) \to X$ et $b :
\Spec(\kappa(z)) \to Y$ tels que $f \circ a = g \circ
b$. Par le même lemme, nous obtenons à nouveau les
points $x \in X$, $y \in Y$ au-dessus du même point $s \in S$ ainsi
que les homomorphismes de corps $\kappa(x) \to \kappa(z)$, $\kappa(y)
\to \kappa(z)$ tels que les compositions $\kappa(s) \to
\kappa(x) \to \kappa(z)$ et $\kappa(s) \to
\kappa(y)
\to \kappa(z)$ soient les mêmes. En
d'autres termes, nous obtenons un homomorphisme d'anneaux
$\kappa(x) \otimes_{\kappa(s)} \kappa(y) \to \kappa(z)$. Nous
notons $\mathfrak p$ le noyau de cette application.

\medskip\noindent
Inversement, étant donné un quadruple $(x, y, s, \mathfrak p)$, nous obtenons un diagramme
commutatif aux flèches pleines
$$
\xymatrix{
X \times_S Y
\ar@/_/[dddr] \ar@/^/[rrrd] &
& & \\
&
\Spec(\kappa(x) \otimes_{\kappa(s)} \kappa(y)/\mathfrak p)
\ar[r] \ar[d] \ar@{-->}[lu] &
\Spec(\kappa(y))
\ar[d] \ar[r] & Y
\ar[dd] \\ &
\Spec(\kappa(x))
\ar[r] \ar[d] &
\Spec(\kappa(s)) \ar[rd] & \\ &
X
\ar[rr]
&
&
S
}
$$
voir la discussion dans la section \ref{schemes-section-points}. Ainsi, nous obtenons
la flèche en pointillés. Le point correspondant $z$ de $X \times_S Y$
est l'image du point générique de
$\Spec(\kappa(x) \otimes_{\kappa(s)} \kappa(y)/\mathfrak p)$. Nous
omettons la vérification que les deux constructions sont réciproques
l'une de l'autre.
\end{proof}

\begin{lemma}
\label{schemes-lemma-fibre-product-immersion}
Soient $f : X \to S$ et $g : Y \to S$ des morphismes de schémas
ayant la même cible.
\begin{enumerate}
\item Si $f : X \to S$ est une immersion fermée,
alors $X \times_S Y \to Y$ est une immersion fermée.
De plus, si $X \to S$ correspond au faisceau quasi-cohérent
d'idéaux $\mathcal{I} \subset \mathcal{O}_S$, alors $X
\times_S Y \to Y$ correspond au faisceau d'idéaux
$\Im(g^*\mathcal{I} \to \mathcal{O}_Y)$. \item Si $f
: X \to S$ est une immersion ouverte, alors
$X \times_S Y \to Y$ est une immersion ouverte.
\item Si $f : X \to S$ est une immersion,
alors $X \times_S Y \to Y$ est une immersion.
\end{enumerate}
\end{lemma}

\begin{proof}
Supposons que $X \to S$ soit une immersion fermée correspondant
au faisceau quasi-cohérent d'idéaux $\mathcal{I} \subset \mathcal{O}_S$. D'après le
lemme \ref{schemes-lemma-restrict-map-to-closed}, le sous-espace fermé $Z \subset
Y$ défini par le faisceau d'idéaux $\Im(g^*\mathcal{I}
\to \mathcal{O}_Y)$ est le produit fibré dans la catégorie
des espaces annelés en anneaux locaux. D'après
le lemme \ref{schemes-lemma-closed-subspace-scheme}, $Z$ est un schéma. Par
conséquent, $Z = X \times_S Y$ et la première affirmation en
découle. La deuxième découle par exemple du lemme \ref{schemes-lemma-open-fibre-product}.
La troisième est une combinaison des
deux premières.
\end{proof}

\begin{definition}
\label{schemes-definition-inverse-image-closed-subscheme}
Soit $f : X \to Y$ un morphisme de schémas. Soit $Z \subset Y$ un
sous-schéma fermé de $Y$. L'{\it image réciproque $f^{-1}(Z)$ du sous-schéma
fermé $Z$} est le sous-schéma fermé $Z \times_Y X$ de $X$.
Voir le lemme \ref{schemes-lemma-fibre-product-immersion} ci-dessus.
\end{definition}

\noindent
Nous pouvons également utiliser occasionnellement cette terminologie pour les sous-schémas localement
fermés et ouverts.









\section{Changement de base en géométrie algébrique}
\label{schemes-section-base-change}

\noindent
Une motivation pour l'introduction du langage des schémas est qu'il
donne une notion très précise de ce que signifie définir une variété
sur un corps particulier. Par exemple, une variété $X$ sur $\mathbf{Q}$
est, par définition (Variétés,
définition \ref{varieties-definition-variety}), un schéma muni d'un
morphisme $X \to \Spec(\mathbf{Q})$ de type fini et séparé,
et qui est lui-même irréductible et réduit\footnote{Bien sûr, les géomètres
algébriques discutent encore pour savoir s'il faut exiger que $X$
soit géométriquement irréductible sur $\mathbf{Q}$.}. Dans tous les cas, l'idée
est plus généralement de travailler avec des schémas sur un {\it schéma de
base} donné, souvent noté $S$. Nous dirons « soit $X$ un schéma sur $S$ » pour
signifier simplement que $X$ est muni d'un morphisme $X \to S$. Dans les
diagrammes, nous représenterons le {\it morphisme structural} $X \to S$
par une flèche descendante de $X$ vers $S$. Nous
nous intéressons souvent davantage aux propriétés de $X$ relativement
à $S$ qu'à la géométrie interne de $X$. Par exemple,
nous aimerions connaître des choses sur les fibres de $X \to S$, ce qui se passe
pour $X$ après changement de base, et ainsi de suite.

\medskip\noindent
Nous introduisons une partie du langage couramment utilisé. Bien
sûr, ce langage n'est qu'un cas particulier de réflexion sur
la catégorie des objets au-dessus d'un objet donné dans une catégorie,
voir Catégories, exemple \ref{categories-example-category-over-X}.

\begin{definition}
\label{schemes-definition-base-change}
Soit $S$ un schéma.
\begin{enumerate}
\item Nous disons que $X$ est un {\it schéma sur $S$} si $X$ est
muni d'un morphisme de schémas $X \to S$. Le morphisme
$X \to S$ est parfois appelé le {\it morphisme
structural}. \item Si $R$ est
un anneau, nous disons que $X$
est un schéma {\it sur $R$} pour dire que
$X$ est un schéma sur $\Spec(R)$.
\item Un {\it morphisme $f : X \to Y$ de schémas sur $S$}
est un morphisme de schémas tel que la composée $X
\to Y \to S$ de $f$ avec le morphisme structural de $Y$ soit égale
au morphisme structural de $X$. \item
Nous notons $\Mor_S(X, Y)$ l'ensemble de tous les morphismes
de $X$ vers $Y$ sur $S$.
\item Soit $X$ un schéma sur $S$. Soit $S' \to S$ un morphisme
de schémas. Le {\it changement de base} de $X$
est le schéma $X_{S'} = S' \times_S X$ sur $S'$. \item
Soit $f : X \to Y$ un morphisme de schémas sur $S$. Soit $S' \to S$ un
morphisme de schémas. Le {\it changement de base} de $f$ est le
morphisme induit $f' : X_{S'} \to Y_{S'}$ (à savoir le morphisme
$\text{id}_{S'} \times_{\text{id}_S} f$). \item
Soit $R$ un anneau. Soit $X$ un schéma sur $R$. Soit $R \to
R'$ un morphisme d'anneaux. Le {\it changement de base} $X_{R'}$
est le schéma $\Spec(R') \times_{\Spec(R)} X$
sur $R'$.
\end{enumerate}
\end{definition}

\noindent
Voici un résultat typique.

\begin{lemma}
\label{schemes-lemma-base-change-immersion}
Soit $S$ un schéma. Soit $f : X \to Y$ une immersion
(resp. \ immersion fermée, resp. immersion ouverte) de
schémas sur $S$. Alors tout changement de base de $f$ est une
immersion (resp. \ immersion fermée, resp. immersion ouverte).
\end{lemma}

\begin{proof}
On peut penser au changement de base de $f$ via le morphisme $S'
\to S$ comme la flèche verticale en haut à gauche dans le diagramme
commutatif suivant :
$$
\xymatrix{
X_{S'} \ar[r] \ar[d] & X \ar[d] \ar@/^4ex/[dd] \\
Y_{S'} \ar[r] \ar[d] & Y \ar[d] \\
S' \ar[r] & S
}
$$
Le diagramme implique $X_{S'} \cong Y_{S'} \times_Y
X$, et le lemme découle du lemme \ref{schemes-lemma-fibre-product-immersion}.
\end{proof}

\noindent
En fait, ce type de résultat est tellement typique qu'il
existe une expression pour l'indiquer. La voici.

\begin{definition}
\label{schemes-definition-preserved-by-base-change}
Propriétés et changement de base.
\begin{enumerate}
\item Soit $\mathcal{P}$ une propriété des schémas sur une base. On dit
que $\mathcal{P}$ est {\it préservée par un changement de base arbitraire}, ou simplement
que $\mathcal{P}$ est {\it préservée par changement de base} si, chaque
fois que $X/S$ possède
$\mathcal{P}$, tout changement de base $X_{S'}/S'$ possède $\mathcal{P}$.
\item Soit $\mathcal{P}$ une propriété des morphismes de schémas sur une base. On dit
que $\mathcal{P}$ est {\it préservée par un changement de base arbitraire}, ou simplement
qu'elle est {\it préservée par changement de base} si, chaque fois
que $f : X \to Y$ sur $S$ possède $\mathcal{P}$, tout changement de
base $f' : X_{S'} \to Y_{S'}$ sur $S'$ possède $\mathcal{P}$.
\end{enumerate}
\end{definition}

\noindent
À ce stade, nous pouvons dire que la propriété « être une immersion fermée » est
préservée par un changement de base arbitraire.

\begin{definition}
\label{schemes-definition-fibre}
Soit $f : X \to S$ un morphisme de schémas. Soit
$s \in S$ un point. La {\it
fibre schématique $X_s$ de $f$ au-dessus de $s$}, ou simplement
la {\it fibre de $f$ au-dessus de $s$}, est
le schéma qui s'insère dans le diagramme de produit fibré suivant
$$
\xymatrix{
X_s = \Spec(\kappa(s)) \times_S X \ar[r] \ar[d] & X
\ar[d] \\
\Spec(\kappa(s)) \ar[r] &
S
}
$$
Nous considérons toujours la fibre $X_s$ comme un schéma sur $\kappa(s)$.
\end{definition}

\begin{lemma}
\label{schemes-lemma-fibre-topological}
Soit $f : X \to S$ un morphisme de schémas.
Considérons les diagrammes
$$
\vcenter{
\xymatrix{
X_s \ar[r] \ar[d] & X \ar[d] \\
\Spec(\kappa(s)) \ar[r] & S
}
}
\quad\text{et}\quad
\vcenter{
\xymatrix{
\Spec(\mathcal{O}_{S, s}) \times_S X \ar[r] \ar[d] & X \ar[d] \\
\Spec(\mathcal{O}_{S, s}) \ar[r] & S
}
}
$$
Dans les deux cas, ces diagrammes induisent des carrés fibrés d'espaces
topologiques et en particulier, la flèche horizontale supérieure est un homéomorphisme
sur son image.
\end{lemma}

\begin{proof}
Choisissez un ouvert affine $U \subset S$ qui contient $s$. Les
morphismes horizontaux inférieurs se factorisent par $U$, voir
le lemme \ref{schemes-lemma-morphism-from-spec-local-ring} par exemple. Ainsi,
nous pouvons supposer que $S$ est affine. Si $X$ est également affine,
alors le résultat découle d'
Algèbre, remarque \ref{algebra-remark-fundamental-diagram}. Dans
le cas général, le résultat découle en recouvrant $X$ par des ouverts affines.
\end{proof}

\begin{lemma}
\label{schemes-lemma-local-ring-fibre}
Soit $f : X \to S$ un morphisme de schémas. Soit $x \in X$ avec
image $s \in S$. En considérant $x$ comme un point de $X_s$ (voir
le lemme \ref{schemes-lemma-fibre-topological}), nous avons
$$
\mathcal{O}_{X_s, x} \cong
\mathcal{O}_{X, x}/\mathfrak m_s\mathcal{O}_{X, x} \cong
\mathcal{O}_{X, x} \otimes_{\mathcal{O}_{S, s}} \kappa(s)
$$
\end{lemma}

\begin{proof}
Comme dans la preuve du lemme \ref{schemes-lemma-fibre-topological}, cela se
réduit au cas où $X$ et $S$ sont affines. Dans le cas affine, l'énoncé se
traduit par Algèbre, remarque \ref{algebra-remark-local-ring-fibre}.
\end{proof}







\section{Morphismes quasi-compacts}
\label{schemes-section-quasi-compact}

\noindent
Un schéma est {\it quasi-compact} si son espace topologique sous-jacent est
quasi-compact. Il existe une notion relative qui est définie comme suit.

\begin{definition}
\label{schemes-definition-quasi-compact}
Un morphisme de schémas est dit {\it quasi-compact} si
l'application continue sous-jacente entre espaces topologiques
est quasi-compacte,
voir Topologie, définition \ref{topology-definition-quasi-compact}.
\end{definition}

\begin{lemma}
\label{schemes-lemma-quasi-compact-affine}
Soit $f : X \to S$ un morphisme de schémas. Les
énoncés suivants sont équivalents
\begin{enumerate}
\item $f : X \to S$ est quasi-compact,
\item l'image réciproque de chaque ouvert affine est quasi-compacte,
et \item il existe un recouvrement ouvert affine $S = \bigcup_{i \in I} U_i$
tel que $f^{-1}(U_i)$ soit quasi-compact pour tout $i$.
\end{enumerate}
\end{lemma}

\begin{proof}
Supposons donné un recouvrement $S = \bigcup_{i \in I} U_i$ comme dans
(3). Tout d'abord, soit $U \subset S$ un ouvert affine quelconque. Pour
tout $u \in U$, nous pouvons trouver un indice $i(u) \in I$ tel que $u
\in U_{i(u)}$. Comme les ouverts standards forment une base pour la topologie sur
$U_{i(u)}$, nous pouvons trouver $W_u \subset U \cap U_{i(u)}$ qui est un ouvert
standard dans $U_{i(u)}$. Par quasi-compacité, nous pouvons trouver un nombre fini de
points $u_1, \ldots, u_n \in U$ tels que $U = \bigcup_{j = 1}^n W_{u_j}$.
Pour chaque $j$, écrivons $f^{-1}U_{i(u_j)} = \bigcup_{k \in K_j}
V_{jk}$ comme une union finie d'ouverts affines. Puisque $W_{u_j} \subset U_{i(u_j)}$
est un ouvert standard, nous voyons que $f^{-1}(W_{u_j}) \cap V_{jk}$
est un ouvert standard de $V_{jk}$, voir Algèbre, lemme \ref{algebra-lemma-spec-functorial}.
Ainsi $f^{-1}(W_{u_j}) \cap V_{jk}$ est affine, et
donc $f^{-1}(W_{u_j})$ est une union finie d'ouverts affines. Cela prouve que
l'image réciproque de tout ouvert affine est une union finie d'ouverts affines.

\medskip\noindent
Ensuite, supposons que l'image réciproque de tout ouvert affine soit une union
finie d'ouverts affines. Soit
$K \subset S$ un ouvert quasi-compact quelconque. Comme $S$ a une base de la
topologie constituée d'ouverts affines, nous voyons que $K$ est une union finie
d'ouverts affines. Par conséquent, l'image réciproque de $K$ est une union finie
d'ouverts affines. Par conséquent, $f$ est quasi-compact.

\medskip\noindent
Enfin, supposons que $f$ soit quasi-compact. Dans ce cas, l'argument du
paragraphe précédent montre que l'image réciproque de tout ouvert affine est
une union finie d'ouverts affines.
\end{proof}

\begin{lemma}
\label{schemes-lemma-quasi-compact-preserved-base-change}
Être quasi-compact est une propriété des morphismes de schémas
sur une base qui est préservée par tout changement de base.
\end{lemma}

\begin{proof}
Omis.
\end{proof}

\begin{lemma}
\label{schemes-lemma-composition-quasi-compact}
Le composé de morphismes quasi-compacts est quasi-compact.
\end{lemma}

\begin{proof}
Cela découle des définitions et de
Topologie, lemme \ref{topology-lemma-composition-quasi-compact}.
\end{proof}

\begin{lemma}
\label{schemes-lemma-closed-immersion-quasi-compact}
Une immersion fermée est quasi-compacte.
\end{lemma}

\begin{proof}
Cela découle des définitions et de
Topologie, lemme \ref{topology-lemma-closed-in-quasi-compact}.
\end{proof}

\begin{example}
\label{schemes-example-open-immersion-not-quasi-compact}
Une immersion ouverte n'est généralement pas
quasi-compacte. L'exemple standard de ceci est le sous-espace
ouvert $U \subset X$, où $X = \Spec(k[x_1, x_2, x_3,
\ldots])$, où $U$ est $X \setminus \{0\}$, et où $0$ est le point
de $X$ correspondant à l'idéal maximal
$(x_1, x_2, x_3, \ldots)$.
\end{example}

\begin{lemma}
\label{schemes-lemma-image-quasi-compact-closed}
Soit $f : X \to S$ un morphisme quasi-compact de schémas. Les affirmations
suivantes sont équivalentes
\begin{enumerate}
\item $f(X) \subset S$ est fermé, et
\item $f(X) \subset S$ est stable par spécialisation.
\end{enumerate}
\end{lemma}

\begin{proof}
Nous avons (1) $\Rightarrow$ (2) par
Topologie, lemme \ref{topology-lemma-open-closed-specialization}.
Supposons (2). Soit $U \subset S$ un ouvert affine. Il suffit de prouver que
$f(X) \cap U$ est fermé. Puisque $U \cap f(X)$ est stable par
spécialisations dans $U$, nous avons réduit au cas où $S$ est affine. Comme $f$ est
quasi-compact, nous en déduisons que $X = f^{-1}(S)$ est quasi-compact
puisque $S$ est affine. Ainsi nous pouvons écrire
$X = \bigcup_{i = 1}^n U_i$ avec $U_i \subset X$ ouvert affine.
Disons $S = \Spec(R)$
et $U_i = \Spec(A_i)$ pour une certaine $R$-algèbre
$A_i$. Alors $f(X) = \Im(\Spec(A_1 \times \ldots \times
A_n) \to \Spec(R))$. Le résultat découle donc
d'Algèbre, lemme \ref{algebra-lemma-image-stable-specialization-closed}.
\end{proof}

\begin{lemma}
\label{schemes-lemma-quasi-compact-closed}
Soit $f : X \to S$ un morphisme quasi-compact de schémas. Alors
$f$ est fermé si et seulement si les spécialisations se relèvent
le long de $f$,
voir Topologie, définition \ref{topology-definition-lift-specializations}.
\end{lemma}

\begin{proof}
D'après Topologie,
lemme \ref{topology-lemma-closed-open-map-specialization}, si $f$ est
fermé, alors les spécialisations se relèvent le long de $f$.
Inversement, supposons que les spécialisations se relèvent le long de
$f$. Soit $Z \subset X$ un sous-ensemble fermé. Nous pouvons
considérer $Z$ comme un schéma avec la structure de schéma réduite
induite, voir définition \ref{schemes-definition-reduced-induced-scheme}.
Puisque $Z \subset X$ est fermé, la restriction
de $f$ à $Z$ est toujours quasi-compacte. De plus, les spécialisations se
relèvent également le long de $Z
\to S$, voir Topologie, lemme \ref{topology-lemma-lift-specialization-composition}.
Par conséquent, il suffit de prouver que $f(X)$ est fermé si les spécialisations se relèvent
le long de $f$. En particulier, $f(X)$ est stable par spécialisation,
voir Topologie, lemme \ref{topology-lemma-lift-specializations-images}.
Ainsi $f(X)$ est fermé
par le lemme \ref{schemes-lemma-image-quasi-compact-closed}.
\end{proof}










\section{Critère valuatif de fermeture universelle}
\label{schemes-section-valuative-criterion-universal-closedness}

\noindent
Topologie, section \ref{topology-section-proper}, étudie les applications
propres comme des applications fermées entre espaces topologiques dont toutes
les fibres sont quasi-compactes, ou comme des applications telles que tous
les changements de base sont des applications fermées. Voici la notion correspondante
en géométrie algébrique.

\begin{definition}
\label{schemes-definition-universally-closed}
Un morphisme de schémas $f : X \to S$ est dit {\it
universellement fermé} si chaque changement de base
$f' : X_{S'} \to S'$ est fermé.
\end{definition}

\noindent
En fait, l'adjectif « universel » est souvent utilisé de cette manière. En
d'autres termes, étant donné une propriété $\mathcal{P}$ des morphismes,
on dit que « $X \to S$ est universellement $\mathcal{P}$ » si et
seulement si chaque changement de base $X_{S'} \to S'$ possède $\mathcal{P}$.

\medskip\noindent
Voir Morphismes, section \ref{morphisms-section-proper}, pour une
discussion plus détaillée des propriétés des morphismes universellement fermés.
Dans cette section, nous limitons la discussion à la relation entre les
morphismes universellement fermés et les morphismes satisfaisant
la partie d'existence du critère valuatif.

\begin{lemma}
\label{schemes-lemma-specializations-lift}
Soit $f : X \to S$ un morphisme de schémas.
\begin{enumerate}
\item Si $f$ est universellement fermé, alors les spécialisations se relèvent
le long de tout changement de base de
$f$, voir Topologie, définition \ref{topology-definition-lift-specializations}. \item
Si $f$ est quasi-compact et que les spécialisations se relèvent le
long de tout changement de base de $f$, alors $f$ est universellement fermé.
\end{enumerate}
\end{lemma}

\begin{proof}
La partie (1) est une conséquence directe
de Topologie, lemme \ref{topology-lemma-closed-open-map-specialization}.
La partie (2) découle
des lemmes \ref{schemes-lemma-quasi-compact-closed}
et \ref{schemes-lemma-quasi-compact-preserved-base-change}.
\end{proof}

\begin{definition}
\label{schemes-definition-valuative-criterion}
Soit $f : X \to S$ un morphisme de schémas. Nous disons que $f$ {\it
satisfait la partie d'existence du critère valuatif} si, étant donné tout diagramme
commutatif dont les flèches pleines sont données
$$
\xymatrix{
\Spec(K) \ar[r] \ar[d] & X \ar[d] \\
\Spec(A) \ar[r] \ar@{-->}[ru] & S
}
$$
où $A$ est un anneau de valuation avec corps des fractions $K$, la
flèche pointillée existe. Nous disons que $f$ {\it satisfait la
partie d'unicité du critère valuatif} s'il existe au plus une flèche
pointillée pour tout diagramme comme ci-dessus (sans exiger
l'existence bien sûr).
\end{definition}

\noindent
Un {\it anneau de valuation} est un anneau local intègre maximal pour la relation
de domination dans son corps des fractions,
voir Algèbre, définition \ref{algebra-definition-valuation-ring}.
Par conséquent, le spectre d'un anneau de valuation possède un point générique
unique $\eta$ et un point fermé unique $0$, et bien sûr nous avons la
spécialisation $\eta \leadsto 0$.
L'importance des anneaux de valuation réside dans le fait que toute spécialisation
de points dans tout schéma est l'image de $\eta \leadsto
0$ sous un certain morphisme du spectre d'un certain anneau de valuation.
Voici le résultat précis.

\begin{lemma}
\label{schemes-lemma-points-specialize}
\begin{slogan}
Les spécialisations se réalisent au moyen d'anneaux de valuation.
\end{slogan}
Soit $S$ un schéma. Soit $s' \leadsto s$ une spécialisation de points de $S$.
Alors
\begin{enumerate}
\item il existe un anneau de valuation $A$ et un morphisme
$f : \Spec(A) \to S$ tels que le point générique $\eta$ de
$\Spec(A)$ ait pour image $s'$ et le point fermé pour image $s$, et \item
étant donnée une extension de corps $K/\kappa(s')$,
on peut faire en sorte que l'extension
$\kappa(\eta)/\kappa(s')$ induite par
$f$ soit isomorphe à l'extension donnée.
\end{enumerate}
\end{lemma}

\begin{proof}
Soit $s' \leadsto s$ une spécialisation dans $S$, et soit
$K/\kappa(s')$ une extension de corps. D'après le lemme
\ref{schemes-lemma-specialize-points} et la
discussion suivant le lemme
\ref{schemes-lemma-characterize-points}, cela conduit
à des homomorphismes d'anneaux $\mathcal{O}_{S, s} \to \kappa(s') \to
K$. Soit $A \subset K$ un quelconque anneau de valuation dont le corps des
fractions est $K$ et qui domine l'image de $\mathcal{O}_{S, s} \to K$, voir
Algèbre, lemme \ref{algebra-lemma-dominate}.
L'homomorphisme d'anneaux $\mathcal{O}_{S, s} \to A$ induit le morphisme
$f : \Spec(A) \to S$, voir
lemme \ref{schemes-lemma-morphism-from-spec-local-ring}. Ce
morphisme possède toutes les propriétés souhaitées par construction.
\end{proof}

\begin{lemma}
\label{schemes-lemma-lift-specializations-valuative}
Soit $f : X \to S$ un morphisme de schémas. Les
énoncés suivants sont équivalents :
\begin{enumerate}
\item Les spécialisations se relèvent par tout changement de
base de $f$. \item Le morphisme $f$ satisfait la partie d'existence
du critère valuatif.
\end{enumerate}
\end{lemma}

\begin{proof}
Supposons que (1) soit vrai. Considérons un diagramme
dont les flèches pleines sont données, comme dans la définition
\ref{schemes-definition-valuative-criterion}. Pour trouver la flèche en
pointillés, nous pouvons remplacer $X \to S$ par
$X_{\Spec(A)} \to \Spec(A)$ puisque, après tout,
l'hypothèse est stable par changement de base.
Ainsi, nous pouvons supposer $S = \Spec(A)$. Soit $x' \in
X$ l'image de $\Spec(K) \to X$, de sorte que
nous avons $\kappa(x') \subset K$, voir le
lemme \ref{schemes-lemma-characterize-points}. D'après l'hypothèse, il existe
une spécialisation $x' \leadsto x$ dans $X$ telle que $x$ ait pour image
le point fermé de $S = \Spec(A)$. Nous obtenons un morphisme d'anneaux
locaux $A \to \mathcal{O}_{X, x}$ et un morphisme
d'anneaux $\mathcal{O}_{X, x} \to \kappa(x')$, voir le lemme
\ref{schemes-lemma-specialize-points} et la discussion qui suit le lemme
\ref{schemes-lemma-characterize-points}. La composée $A \to \mathcal{O}_{X, x}
\to \kappa(x') \to K$ est l'injection donnée $A \to K$. Comme $A \to
\mathcal{O}_{X, x}$ est local, l'image de $\mathcal{O}_{X,
x} \to K$ domine $A$ et est donc égale à $A$,
d'après Algèbre, définition \ref{algebra-definition-valuation-ring}.
Ainsi, nous obtenons un homomorphisme d'anneaux $\mathcal{O}_{X,
x} \to A$ et donc un morphisme $\Spec(A)
\to X$ (voir le lemme \ref{schemes-lemma-morphism-from-spec-local-ring}
et la discussion qui suit). Cela prouve (2).

\medskip\noindent
En revanche, supposons que (2) soit vérifié. Il est immédiat
que la partie d'existence du critère valuatif est vérifiée
pour tout changement de base $X_{S'} \to S'$ de $f$ en considérant
le diagramme commutatif suivant
$$
\xymatrix{
\Spec(K) \ar[r] \ar[d] & X_{S'} \ar[r] \ar[d] & X \ar[d] \\
\Spec(A) \ar[r] \ar@{-->}[ru] \ar@{-->}[rru] & S' \ar[r] & S
}
$$
En d'autres termes, la flèche pointillée la plus horizontale donnera
l'autre par définition du produit fibré. Il suffit donc de montrer
que les spécialisations se relèvent le long de $f$. Soit
$s' \leadsto s$ une spécialisation dans $S$, et soit $x' \in
X$ un point situé au-dessus de $s'$. Appliquons
le lemme \ref{schemes-lemma-points-specialize}
à $s' \leadsto s$ et à l'extension de corps $K =
\kappa(x')/\kappa(s')$. Nous
obtenons un diagramme commutatif
$$
\xymatrix{
\Spec(K) \ar[rr] \ar[d] & & X \ar[d] \\
\Spec(A) \ar[r] \ar@{-->}[rru] &
\Spec(\mathcal{O}_{S, s}) \ar[r] & S
}
$$
et, par la condition (2), nous obtenons la flèche
pointillée. L'image $x$ du point fermé de $\Spec(A)$ dans
$X$ sera une solution à notre problème, c'est-à-dire que
$x$ est une spécialisation de $x'$ et a pour image $s$.
\end{proof}

\begin{proposition}[Critère valuatif de fermeture universelle]
\label{schemes-proposition-characterize-universally-closed}
Soit $f$ un morphisme quasi-compact de schémas.
Alors $f$ est universellement fermé si et seulement
si $f$ satisfait la partie d'existence du critère valuatif.
\end{proposition}

\begin{proof}
Ceci est une conséquence formelle
des lemmes \ref{schemes-lemma-specializations-lift} et
\ref{schemes-lemma-lift-specializations-valuative} ci-dessus.
\end{proof}

\begin{example}
\label{schemes-example-projective-line-universally-closed}
Soit $k$ un corps. Considérons le morphisme structural $p
: \mathbf{P}^1_k \to \Spec(k)$ de la droite projective
sur $k$, voir l'exemple \ref{schemes-example-projective-line}. Utilisons
le critère valuatif ci-dessus pour prouver que $p$ est universellement
fermé. Par
construction, $\mathbf{P}^1_k$ est recouvert par deux ouverts
affines et donc $p$ est quasi-compact. Supposons
qu'un diagramme commutatif
$$
\xymatrix{
\Spec(K) \ar[r]_\xi \ar[d] & \mathbf{P}^1_k \ar[d] \\
\Spec(A) \ar[r]^\varphi & \Spec(k)
}
$$
soit donné, où $A$ est un anneau de valuation et $K$ est son corps
des fractions. Rappelons que $\mathbf{P}^1_k$ est obtenu en recollant
$\Spec(k[x])$ à $\Spec(k[y])$ en identifiant $D(x)$
à $D(y)$ via $x = y^{-1}$ (ou plus symétriquement $xy = 1$). Pour montrer
qu'il existe un morphisme $\Spec(A) \to \mathbf{P}^1_k$ s'insérant
en diagonale dans le diagramme ci-dessus, nous pouvons supposer que $\xi$
ait son image dans l'ouvert $\Spec(k[x])$ (par symétrie). Cela donne le diagramme
commutatif d'anneaux suivant
$$
\xymatrix{
K & k[x] \ar[l]^{\xi^\sharp} \\
A \ar[u] & k \ar[u] \ar[l]_{\varphi^\sharp}
}
$$
Par Algèbre, lemme \ref{algebra-lemma-valuation-ring-x-or-x-inverse}, nous voyons
que soit $\xi^\sharp(x)
\in A$, soit $\xi^\sharp(x)^{-1} \in A$. Dans le premier
cas, nous obtenons un morphisme d'anneaux
$$
k[x] \to A,
\ \lambda \mapsto \varphi^\sharp(\lambda),
\ x \mapsto \xi^\sharp(x)
$$
s'insérant dans le diagramme d'anneaux ci-dessus, ce qui conclut. Dans
le deuxième cas, nous voyons que nous obtenons un morphisme d'anneaux
$$
k[y] \to A,
\ \lambda \mapsto \varphi^\sharp(\lambda),
\ y \mapsto \xi^\sharp(x)^{-1}.
$$
Cela donne un morphisme
$\Spec(A) \to \Spec(k[y]) \to \mathbf{P}^1_k$ qui
s'insère en diagonale dans le diagramme commutatif initial de cet exemple
(vérification omise).
\end{example}





























\section{Axiomes de séparation}
\label{schemes-section-separation-axioms}

\noindent
Un espace topologique $X$ est séparé si et seulement si la
diagonale $\Delta \subset X \times X$ est un sous-ensemble fermé.
L'analogue en géométrie algébrique est, étant donné un schéma $X$ sur
un schéma de base $S$, de considérer le
morphisme diagonal
$$
\Delta_{X/S} : X \longrightarrow X \times_S X.
$$
C'est le morphisme unique de schémas tel que
$\text{pr}_1 \circ \Delta_{X/S} = \text{id}_X$ et
$\text{pr}_2 \circ \Delta_{X/S} = \text{id}_X$ (il existe dans toute
catégorie avec produits fibrés).

\begin{lemma}
\label{schemes-lemma-diagonal-affines-closed}
Le morphisme diagonal d'un morphisme entre schémas affines est fermé.
\end{lemma}

\begin{proof}
Le morphisme diagonal associé au morphisme $\Spec(S) \to
\Spec(R)$ est le morphisme entre spectres correspondant
à l'homomorphisme d'anneaux
$S \otimes_R S \to S$, $a \otimes b \mapsto ab$. Cette
application est clairement surjective, donc $S \cong S \otimes_R S/J$
pour un certain idéal $J \subset S \otimes_R S$. Par
conséquent, $\Delta$ est une immersion fermée
selon l'exemple \ref{schemes-example-closed-immersion-affines}.
\end{proof}

\begin{lemma}
\label{schemes-lemma-diagonal-immersion}
\begin{slogan}
Le morphisme diagonal d'un schéma relatif est une immersion.
\end{slogan}
Soit $X$ un schéma sur $S$. Le
morphisme diagonal $\Delta_{X/S}$ est une immersion.
\end{lemma}

\begin{proof}
Rappelons que si $V \subset X$ est ouvert affine et se projette
dans $U \subset S$ ouvert affine, alors $V \times_U V$ est ouvert
affine dans $X \times_S X$, voir les lemmes \ref{schemes-lemma-fibre-product-affines}
et \ref{schemes-lemma-open-fibre-product}.
Considérons le sous-schéma ouvert $W$ de $X \times_S X$
qui est l'union de ces ouverts affines $V \times_U V$.
D'après le lemme \ref{schemes-lemma-closed-local-target}, il suffit
de montrer que chaque morphisme
$\Delta_{X/S}^{-1}(V \times_U V) \to V \times_U V$ est
une immersion fermée. Comme $V = \Delta_{X/S}^{-1}(V \times_U V)$,
nous vérifions simplement que $\Delta_{V/U}$ est une immersion
fermée, ce qui est le lemme \ref{schemes-lemma-diagonal-affines-closed}.
\end{proof}

\begin{definition}
\label{schemes-definition-separated}
Soit $f : X \to S$ un morphisme de schémas.
\begin{enumerate}
\item Nous disons que $f$ est {\it séparé} si le morphisme diagonal $\Delta_{X/S}$
est une immersion
fermée. \item On dit que $f$ est {\it quasi-séparé} si le morphisme
diagonal $\Delta_{X/S}$ est un morphisme quasi-compact.
\item On dit qu'un schéma $Y$ est {\it séparé} si le
morphisme $Y \to \Spec(\mathbf{Z})$ est séparé.
\item On dit qu'un schéma $Y$ est {\it quasi-séparé} si le morphisme
$Y \to \Spec(\mathbf{Z})$ est quasi-séparé.
\end{enumerate}
\end{definition}

\noindent
D'après les lemmes \ref{schemes-lemma-diagonal-immersion} et \ref{schemes-lemma-immersion-when-closed}, nous
voyons que $\Delta_{X/S}$ est une immersion fermée si et seulement si
$\Delta_{X/S}(X) \subset X \times_S X$ est un sous-ensemble fermé. De plus, d'après le
lemme \ref{schemes-lemma-closed-immersion-quasi-compact}, nous voyons qu'un morphisme
séparé est quasi-séparé. La raison d'introduire les morphismes quasi-séparés est
que les morphismes non séparés apparaissent naturellement dans l'étude des variétés
algébriques (en particulier lorsqu'on étudie les problèmes de modules, les champs
algébriques, etc.). Mais le plus souvent, ils restent tout de même quasi-séparés.

\begin{example}
\label{schemes-example-not-quasi-separated}
Voici un exemple de morphisme non quasi-séparé. Supposons
$X = X_1 \cup X_2 \to S = \Spec(k)$ avec $X_1
= X_2 = \Spec(k[t_1, t_2, t_3, \ldots])$ recollés
le long du complément de $\{0\} = \{(t_1, t_2, t_3, \ldots)\}$
(recollés comme dans l'exemple \ref{schemes-example-affine-space-zero-doubled}).
Dans ce cas, l'image inverse du schéma affine
$X_1 \times_S X_2$ sous $\Delta_{X/S}$ est le schéma
$\Spec(k[t_1, t_2, t_3, \ldots]) \setminus \{0\}$
qui n'est pas quasi-compact.
\end{example}

\begin{lemma}
\label{schemes-lemma-where-are-they-equal}
Soient $X$, $Y$ des schémas sur $S$.
Soient $a, b : X \to Y$ des morphismes de schémas sur
$S$. Il existe un plus grand sous-schéma localement
fermé $Z \subset X$ tel que $a|_Z = b|_Z$. En fait, $Z$
est l'égalisateur de $(a, b)$. De plus, si $Y$ est séparé
sur $S$, alors $Z$ est un sous-schéma fermé.
\end{lemma}

\begin{proof}
L'égalisateur de $(a, b)$ est, pour des raisons catégoriques,
le produit fibré $Z$ dans le diagramme suivant.
$$
\xymatrix{
Z = Y \times_{(Y \times_S Y)} X \ar[r] \ar[d] &
 X \ar[d]^{(a , b)} \\
Y \ar[r]^-{\Delta_{Y/S}} & Y \times_S Y
}
$$
Ainsi, le lemme découle des lemmes
\ref{schemes-lemma-base-change-immersion}, \ref{schemes-lemma-diagonal-immersion} et de la
définition \ref{schemes-definition-separated}.
\end{proof}

\begin{lemma}
\label{schemes-lemma-characterize-quasi-separated}
Soit $f : X \to S$ un morphisme de schémas. Les affirmations
suivantes sont équivalentes :
\begin{enumerate}
\item Le morphisme $f$ est quasi-séparé. \item
Pour chaque paire d'ouverts affines $U, V \subset X$ qui se
projettent dans un ouvert affine commun de $S$, l'intersection $U
\cap V$ est une union finie d'ouverts affines de $X$. \item
Il existe un recouvrement par ouverts affines $S = \bigcup_{i \in I} U_i$
et pour chaque $i$ un recouvrement par ouverts affines $f^{-1}U_i = \bigcup_{j \in I_i}
V_j$ tel que pour chaque $i$ et chaque paire $j, j' \in I_i$
l'intersection $V_j \cap V_{j'}$ soit une union finie d'ouverts
affines de $X$.
\end{enumerate}
\end{lemma}

\begin{proof}
Prouvons que (3) implique (1). D'après
le lemme \ref{schemes-lemma-affine-covering-fibre-product},
le recouvrement $X \times_S X = \bigcup_i \bigcup_{j, j'} V_j \times_{U_i} V_{j'}$
est un recouvrement par ouverts affines de $X
\times_S X$. De plus, $\Delta_{X/S}^{-1}(V_j \times_{U_i} V_{j'}) = V_j \cap
V_{j'}$. Ainsi, l'implication découle du lemme \ref{schemes-lemma-quasi-compact-affine}.

\medskip\noindent
L'implication (1) $\Rightarrow$ (2) découle du fait que,
sous les hypothèses de (2), le produit fibré $U
\times_S V$ est un ouvert affine de $X \times_S X$.
L'implication (2) $\Rightarrow$ (3) est triviale.
\end{proof}

\begin{lemma}
\label{schemes-lemma-characterize-separated}
Soit $f : X \to S$ un morphisme de schémas.
\begin{enumerate}
\item Si $f$ est séparé alors pour chaque paire d'ouverts affines
$(U, V)$ de $X$ qui se projettent dans un
ouvert affine commun de $S$ nous avons
\begin{enumerate}
\item l'intersection $U \cap V$ est affine. \item
l'homomorphisme
d'anneaux $\mathcal{O}_X(U) \otimes_{\mathbf{Z}} \mathcal{O}_X(V)
\to \mathcal{O}_X(U \cap V)$
est surjectif.
\end{enumerate}
\item Si toute paire de points $x_1, x_2 \in X$ au-dessus d'un point commun
$s \in S$ est contenue dans des ouverts affines $x_1 \in U$, $x_2
\in V$ qui se projettent dans un ouvert affine commun de $S$ tels
que (a), (b) sont vérifiés, alors $f$ est séparé.
\end{enumerate}
\end{lemma}

\begin{proof}
Supposons $f$ séparé. Supposons que $(U, V)$ soit une paire comme
dans (1). Soit $W = \Spec(R)$ un ouvert affine de $S$ contenant
à la fois $f(U)$ et $f(V)$. Notons $U = \Spec(A)$
et $V = \Spec(B)$ pour les $R$-algèbres $A$ et
$B$. D'après le lemme \ref{schemes-lemma-open-fibre-product}, on
voit que $U \times_S V = U \times_W V = \Spec(A \otimes_R
B)$ est un ouvert affine de $X \times_S X$. Par
conséquent, d'après le lemme \ref{schemes-lemma-closed-subspace-scheme},
on voit que $\Delta^{-1}(U \times_S V) \to U
\times_S V$ peut être identifié à $\Spec((A \otimes_R B)/J) \to \Spec(A \otimes_R
B)$ pour un certain idéal $J \subset A \otimes_R
B$. Ainsi $U \cap V = \Delta^{-1}(U \times_S V)$ est affine.
L'assertion (1)(b) est vérifiée
car $A \otimes_{\mathbf{Z}} B \to (A \otimes_R B)/J$ est surjective.

\medskip\noindent
Supposons que l'hypothèse formulée en (2) soit vérifiée.
Il est clair que la collection d'ouverts affines $U \times_S
V$ pour les paires $(U, V)$ comme en (2) forme un recouvrement affine
ouvert de $X \times_S X$ (voir par exemple\ le lemme \ref{schemes-lemma-affine-covering-fibre-product}).
Il suffit donc de montrer que chaque morphisme
$U \cap V = \Delta_{X/S}^{-1}(U \times_S V) \to U \times_S V$ est une
immersion fermée, voir le lemme \ref{schemes-lemma-closed-local-target}. D'après
l'hypothèse (a), nous avons $U \cap V = \Spec(C)$ pour un certain anneau $C$. Après
avoir choisi un ouvert affine $W = \Spec(R)$ de $S$ dans
lequel les images de $U$ et de $V$ sont contenues et en écrivant $U =
\Spec(A)$, $V = \Spec(B)$, nous voyons que l'hypothèse (b)
signifie que la composée
$$
A \otimes_{\mathbf{Z}} B \to A \otimes_R B \to C
$$
est surjective. Ainsi $A \otimes_R B \to C$ est surjective et nous
concluons que $\Spec(C) \to \Spec(A \otimes_R B)$ est
une immersion fermée.
\end{proof}

\begin{example}
\label{schemes-example-projective-line-separated}
Soit $k$ un corps. Considérons le morphisme structural $p
: \mathbf{P}^1_k \to \Spec(k)$ de la droite projective
sur $k$, voir l'exemple \ref{schemes-example-projective-line}. Utilisons
le lemme ci-dessus pour prouver que $p$ est
séparé. Par construction, $\mathbf{P}^1_k$ est recouvert par deux ouverts
affines $U = \Spec(k[x])$ et $V = \Spec(k[y])$ avec
intersection $U \cap V = \Spec(k[x, y]/(xy - 1))$ (en utilisant
la notation évidente). Il suffit donc de vérifier que les
conditions (2)(a) et (2)(b) du lemme \ref{schemes-lemma-characterize-separated} sont
vérifiées pour les paires d'ouverts affines $(U, U)$, $(U, V)$, $(V, U)$
et $(V, V)$. Pour les paires $(U, U)$ et $(V, V)$, c'est trivial. Pour
la paire $(U, V)$, cela revient à prouver que $U
\cap V$ est affine, ce qui est vrai, et que l'homomorphisme d'anneaux
$$
k[x] \otimes_{\mathbf{Z}} k[y] \longrightarrow k[x, y]/(xy - 1)
$$
est surjectif. C'est évident car tout élément du côté
droit peut être écrit comme la somme d'un polynôme en
$x$ et d'un polynôme en $y$.
\end{example}

\begin{lemma}
\label{schemes-lemma-fibre-product-after-map}
Soient $f : X \to T$ et $g : Y \to T$ des morphismes de schémas
ayant la même cible. Soit $h : T \to S$ un morphisme de schémas.
Alors le morphisme induit $i : X \times_T Y \to X \times_S Y$ est
une immersion. Si $T \to S$ est séparé, alors $i$ est une immersion
fermée. Si $T \to S$ est quasi-séparé, alors $i$ est un
morphisme quasi-compact.
\end{lemma}

\begin{proof}
Par la théorie générale des catégories, le diagramme suivant
$$
\xymatrix{
X \times_T Y \ar[r] \ar[d] & X \times_S Y \ar[d] \\
T \ar[r]^{\Delta_{T/S}} \ar[r] & T \times_S T
}
$$
est un diagramme de produit fibré. Le lemme
découle des lemmes \ref{schemes-lemma-diagonal-immersion},
\ref{schemes-lemma-fibre-product-immersion} et
\ref{schemes-lemma-quasi-compact-preserved-base-change}.
\end{proof}

\begin{lemma}
\label{schemes-lemma-semi-diagonal}
Soit $g : X \to Y$ un morphisme de schémas sur $S$. Le
morphisme $i : X \to X \times_S Y$ est une immersion.
Si $Y$ est séparé sur $S$, c'est une immersion fermée.
Si $Y$ est quasi-séparé sur $S$, ce morphisme est quasi-compact.
\end{lemma}

\begin{proof}
Il s'agit d'un cas particulier du lemme \ref{schemes-lemma-fibre-product-after-map}
appliqué au morphisme $X = X \times_Y Y \to X \times_S Y$.
\end{proof}

\begin{lemma}
\label{schemes-lemma-section-immersion}
Soit $f : X \to S$ un morphisme de schémas.
Soit $s : S \to X$ une section de $f$ (dans une formule $f \circ s = \text{id}_S$).
Alors $s$ est une immersion.
Si $f$ est séparé, alors $s$ est une immersion fermée.
Si $f$ est quasi-séparé, alors $s$ est quasi-compact.
\end{lemma}

\begin{proof}
Il s'agit d'un cas particulier du lemme \ref{schemes-lemma-semi-diagonal} appliqué
à $g =s$, donc au morphisme $i = s : S \to S \times_S X$.
\end{proof}

\begin{lemma}
\label{schemes-lemma-separated-permanence}
Propriétés de permanence.
\begin{enumerate}
\item Le composé de deux morphismes séparés est séparé.
\item Le composé de deux morphismes quasi-séparés est quasi-séparé.
\item Le changement de base d'un morphisme séparé est séparé.
\item Le changement de base d'un morphisme quasi-séparé est quasi-séparé.
\item Un produit (fibré) de morphismes séparés est séparé.
\item Un produit (fibré) de morphismes quasi-séparés est quasi-séparé.
\end{enumerate}
\end{lemma}

\begin{proof}
Soient $X \to Y \to Z$ des morphismes. Supposons que $X \to Y$
et $Y \to Z$ soient séparés. Le composé
$$
X \to X \times_Y X \to X \times_Z X
$$
est fermée parce que le premier morphisme l'est par hypothèse et le
second d'après le lemme \ref{schemes-lemma-fibre-product-after-map}. Le même argument
fonctionne pour « quasi-séparé » (avec les mêmes références).

\medskip\noindent
Soit $f : X \to Y$ un morphisme de schémas sur une base $S$. Soit
$S' \to S$ un morphisme de schémas. Soit $f' : X_{S'} \to Y_{S'}$ le
changement de base de $f$. Alors le morphisme diagonal de
$f'$ est un morphisme
$$
\Delta_{f'} :
X_{S'} = S' \times_S X
\longrightarrow
X_{S'} \times_{Y_{S'}} X_{S'} = S' \times _S (X \times_Y X)
$$
qui est facilement vu comme étant le changement de base de
$\Delta_f$. Ainsi (3) et (4) découlent du fait que
les immersions fermées et les morphismes quasi-compacts sont
préservés par tout changement de base
(Lemmes \ref{schemes-lemma-fibre-product-immersion}
et \ref{schemes-lemma-quasi-compact-preserved-base-change}).

\medskip\noindent
Si $f : X \to Y$ et $g : U \to V$ sont des morphismes de schémas sur une base $S$,
alors $f \times g$ est le composé de $X \times_S U \to X \times_S V$ (un
changement de base de $g$) et $X \times_S V \to Y \times_S V$ (un changement de
base de $f$). Ainsi, (5) et (6) suivent de (1) -- (4).
\end{proof}

\begin{lemma}
\label{schemes-lemma-compose-after-separated}
\begin{slogan}
Les morphismes séparés et quasi-séparés satisfont la propriété de simplification.
\end{slogan}
Soient $f : X \to Y$ et $g : Y \to Z$ des morphismes de schémas. Si $g
\circ f$ est séparé, alors $f$ l'est aussi. Si
$g \circ f$ est quasi-séparé, alors $f$ l'est également.
\end{lemma}

\begin{proof}
Supposons que $g \circ f$ soit séparé.
Considérons la factorisation $X \to X \times_Y X \to X \times_Z X$ du morphisme
diagonal de $g \circ f$. D'après le lemme
\ref{schemes-lemma-fibre-product-after-map}, le dernier
morphisme est une immersion. D'après l'hypothèse, l'image de $X$ dans
$X \times_Z X$ est fermée. Par conséquent, elle est aussi fermée
dans $X \times_Y X$. Ainsi, nous voyons que $X \to X \times_Y X$
est une immersion fermée selon le lemme \ref{schemes-lemma-immersion-when-closed}.

\medskip\noindent
Supposons que $g \circ f$ soit quasi-séparé. Soit $V
\subset Y$ un ouvert affine dont l'image est contenue dans un ouvert affine
de $Z$. Soient $U_1, U_2 \subset X$ des ouverts affines dont les images
sont contenues dans $V$. Alors $U_1 \cap U_2$ est une union finie d'ouverts
affines car $U_1, U_2$ ont leurs images dans un ouvert affine commun
de $Z$. Puisque nous pouvons recouvrir $Y$ par des ouverts affines
comme $V$, nous déduisons le lemme du lemme \ref{schemes-lemma-characterize-quasi-separated}.
\end{proof}

\begin{lemma}
\label{schemes-lemma-quasi-compact-permanence}
Soient $f : X \to Y$ et $g : Y \to Z$ des morphismes de schémas.
Si $g \circ f$ est quasi-compact et $g$ est quasi-séparé,
alors $f$ est quasi-compact.
\end{lemma}

\begin{proof}
Ceci est vrai parce que $f$ est égal au composé
$(1, f) : X \to X \times_Z Y \to Y$. La première application
est quasi-compacte d'après le lemme \ref{schemes-lemma-section-immersion}
car elle est une section du morphisme quasi-séparé $X \times_Z Y \to X$ (un changement
de base de $g$, voir le lemme \ref{schemes-lemma-separated-permanence}). La
deuxième application est quasi-compacte car
elle est le changement de base de $g \circ f$,
voir le lemme \ref{schemes-lemma-quasi-compact-preserved-base-change}.
Et les composés de morphismes
quasi-compacts sont quasi-compactes, voir le lemme \ref{schemes-lemma-composition-quasi-compact}.
\end{proof}

\begin{lemma}
\label{schemes-lemma-affine-separated}
Un schéma affine est séparé. Un morphisme d'un schéma affine
vers un autre schéma est séparé.
\end{lemma}

\begin{proof}
Soit $U = \Spec(A)$ un schéma affine. Alors $U \to \Spec(\mathbf{Z})$ a une
diagonale fermée d'après le lemme \ref{schemes-lemma-diagonal-affines-closed}.
Ainsi $U$ est séparé selon la définition \ref{schemes-definition-separated}.
Si $U \to X$ est un morphisme de schémas, alors nous pouvons
appliquer le lemme \ref{schemes-lemma-compose-after-separated}
aux morphismes $U \to X \to \Spec(\mathbf{Z})$ pour
conclure que $U \to X$ est séparé.
\end{proof}

\noindent
Vous vous êtes peut-être demandé si la condition de ne
considérer que des paires d'ouverts affines dont l'image est contenue dans un
ouvert affine est vraiment nécessaire pour pouvoir conclure que leur
intersection est affine. Souvent, ce n'est pas le cas !

\begin{lemma}
\label{schemes-lemma-curiosity}
Soit $f : X \to S$ un morphisme.
Supposons que $f$ soit séparé et que $S$ soit un schéma séparé.
Supposons que $U \subset X$ et $V \subset X$ soient des ouverts
affines. Alors $U \cap V$ est affine (et un sous-schéma fermé de $U \times V$).
\end{lemma}

\begin{proof}
Dans ce cas, $X$ est séparé par le lemme \ref{schemes-lemma-separated-permanence}.
Ainsi, $U \cap V$ est affine en
appliquant le lemme \ref{schemes-lemma-characterize-separated} au
morphisme $X \to \Spec(\mathbf{Z})$.
\end{proof}

\noindent
D'autre part, l'exemple suivant montre que nous ne pouvons pas nous attendre
à ce que l'image d'un ouvert affine soit contenue dans un ouvert affine.

\begin{example}
\label{schemes-example-image-affine-projective}
Considérons le schéma non affine
$U = \Spec(k[x, y]) \setminus \{(x, y)\}$ de
l'exemple \ref{schemes-example-not-affine}. D'autre part, considérons le
schéma
$$
\mathbf{GL}_{2, k} = \Spec(k[a, b, c, d, 1/ad - bc]).
$$
Il existe un morphisme $\mathbf{GL}_{2, k} \to U$ correspondant
à l'homomorphisme d'anneaux $x \mapsto a$, $y \mapsto b$. Il est facile
de voir qu'il s'agit d'un morphisme surjectif, et donc l'image n'est
contenue dans aucun ouvert affine de $U$. En fait, le schéma
affine $\mathbf{GL}_{2, k}$ surjecte également sur $\mathbf{P}^1_k$,
et $\mathbf{P}^1_k$ n'admet même d'immersion dans {\it aucun} schéma affine.
\end{example}

\begin{remark}
\label{schemes-remark-quasi-compact-and-quasi-separated}
Soit $P$ l'une des $4$ propriétés suivantes : «
quasi-séparé », « séparé », «
quasi-compact et quasi-séparé » ou « quasi-compact et séparé ». Alors les
énoncés suivants sont vrais
\begin{enumerate}
\item Tout schéma affine est $P$.
\item Étant donné un morphisme $f : X \to Y$ de schémas où $Y$ est $P$, alors
$f$ est $P$ si et seulement si $X$ est
$P$. \item Si $X \to Y$ et $Z \to Y$ sont des morphismes de
schémas et $X$, $Y$, $Z$ ont $P$, alors $X \times_Y Z$ a $P$.
\end{enumerate}
L'énoncé (1) est clair.
Soit $f : X \to Y$ comme dans (2). Si $f$ a $P$, alors $X$ l'a aussi
par les lemmes \ref{schemes-lemma-separated-permanence} et
\ref{schemes-lemma-composition-quasi-compact} appliqués à
$X \to Y \to \Spec(\mathbf{Z})$. Si $X$ a $P$, alors $f$ l'a aussi par les
lemmes \ref{schemes-lemma-compose-after-separated} et
\ref{schemes-lemma-quasi-compact-permanence} appliqués à $X
\to Y \to \Spec(\mathbf{Z})$. Soient
$X \to Y$ et $Z \to Y$ comme dans (3). Alors
la projection $X \times_Y Z \to Z$ a $P$ comme changement de base de
la flèche $X \to Y$ qui a $P$ par (2), voir les lemmes
\ref{schemes-lemma-separated-permanence} et
\ref{schemes-lemma-quasi-compact-preserved-base-change}. Par
conséquent, $X \times_Y Z$ a $P$ par (2).
\end{remark}





\section{Critère valuatif de séparation}
\label{schemes-section-valuative-separatedness}

\begin{lemma}
\label{schemes-lemma-separated-implies-valuative}
Soit $f : X \to S$ un morphisme de schémas.
Si $f$ est séparé, alors $f$ satisfait à la partie
d'unicité du critère valuatif.
\end{lemma}

\begin{proof}
Soit un diagramme comme dans la définition \ref{schemes-definition-valuative-criterion}
donné. Supposons qu'il existe deux morphismes
$a, b : \Spec(A) \to X$ s'insérant dans le diagramme.
Soit $Z \subset \Spec(A)$ l'égalisateur de $a$ et $b$. D'après le
lemme \ref{schemes-lemma-where-are-they-equal}, ceci est un sous-schéma
fermé de $\Spec(A)$. Par hypothèse, il contient le point
générique de $\Spec(A)$. Comme $A$ est un anneau intègre, cela
implique $Z = \Spec(A)$. Par conséquent, $a = b$ comme voulu.
\end{proof}

\begin{lemma}[Critère valuatif de séparation]
\label{schemes-lemma-valuative-criterion-separatedness}
\begin{reference}
\cite[II Proposition 7.2.3]{EGA}
\end{reference}
Soit $f : X \to S$ un morphisme.
Supposons
\begin{enumerate}
\item le morphisme $f$ est quasi-séparé, et
\item le morphisme $f$ satisfait la partie
d'unicité du critère valuatif.
\end{enumerate}
Alors $f$ est séparé.
\end{lemma}

\begin{proof}
Par l'hypothèse (1),
la proposition \ref{schemes-proposition-characterize-universally-closed}, et
les lemmes \ref{schemes-lemma-diagonal-immersion} et \ref{schemes-lemma-immersion-when-closed},
nous voyons qu'il suffit de prouver que le
morphisme $\Delta_{X/S} : X \to X \times_S X$ satisfait la partie
d'existence du critère valuatif.
Considérons le diagramme commutatif
$$
\xymatrix{
\Spec(K) \ar[r] \ar[d] & X \ar[d] \\
\Spec(A) \ar[r] \ar@{-->}[ru] & X \times_S X
}
$$
dont les flèches pleines sont données. La flèche en bas à droite
correspond à une paire de morphismes $a, b : \Spec(A) \to X$ au-dessus
de $S$. Par (2), nous voyons que $a = b$. Ainsi, $a$ fournit la flèche
pointillée cherchée.
\end{proof}





\section{Monomorphismes}
\label{schemes-section-monomorphisms}

\begin{definition}
\label{schemes-definition-monomorphism}
Un morphisme de schémas est appelé {\it monomorphisme} s'il est
un monomorphisme dans la catégorie des schémas,
voir Catégories, définition \ref{categories-definition-mono-epi}.
\end{definition}

\begin{lemma}
\label{schemes-lemma-monomorphism}
\begin{slogan}
Un morphisme de schémas est un monomorphisme si et seulement si sa diagonale est un isomorphisme.
\end{slogan}
Soit $j : X \to Y$ un morphisme de schémas. Alors
$j$ est un monomorphisme si et seulement si le
morphisme diagonal $\Delta_{X/Y} : X \to X \times_Y X$ est
un isomorphisme.
\end{lemma}

\begin{proof}
Ceci est vrai dans toute catégorie avec des produits fibrés.
\end{proof}

\begin{lemma}
\label{schemes-lemma-monomorphism-separated}
Un monomorphisme de schémas est séparé.
\end{lemma}

\begin{proof}
Cela découle du fait qu'un isomorphisme est une immersion fermée et
du lemme \ref{schemes-lemma-monomorphism} ci-dessus.
\end{proof}

\begin{lemma}
\label{schemes-lemma-composition-monomorphism}
Le composé de monomorphismes est un monomorphisme.
\end{lemma}

\begin{proof}
Vrai dans toute catégorie.
\end{proof}

\begin{lemma}
\label{schemes-lemma-base-change-monomorphism}
Le changement de base d'un monomorphisme est un monomorphisme.
\end{lemma}

\begin{proof}
Vrai dans toute catégorie avec des produits fibrés.
\end{proof}

\begin{lemma}
\label{schemes-lemma-injective-points}
Soit $j : X \to Y$ un morphisme de schémas. Si
$j$ est injectif sur les points, alors $j$ est séparé.
\end{lemma}

\begin{proof}
Soit $z$ un point de $X \times_Y X$. Alors $x = \text{pr}_1(z)$ et
$\text{pr}_2(z)$ sont les mêmes parce que $j$ envoie ces points sur le
même point $y$ de $Y$. Alors nous pouvons choisir un voisinage
ouvert affine $V \subset Y$ de $y$ et un voisinage ouvert
affine $U \subset X$ de $x$ avec $j(U) \subset V$. Alors $z
\in U \times_V U \subset X \times_Y X$. Ainsi $X
\times_Y X$ est l'union des ouverts affines $U
\times_V U$. Puisque $\Delta_{X/Y}^{-1}(U \times_V U) = U$ et puisque $U
\to U \times_V U$ est une immersion fermée, nous concluons que
$\Delta_{X/Y}$ est une immersion fermée (voir l'argument dans la
démonstration du lemme \ref{schemes-lemma-diagonal-immersion}).
\end{proof}

\begin{lemma}
\label{schemes-lemma-injective-points-surjective-stalks}
Soit $j : X \to Y$ un morphisme de schémas.
Si
\begin{enumerate}
\item $j$ est injectif sur les points, et
\item pour tout $x \in X$ l'homomorphisme
d'anneaux $j^\sharp_x : \mathcal{O}_{Y, j(x)} \to \mathcal{O}_{X, x}$
est surjectif,
\end{enumerate}
alors $j$ est un monomorphisme.
\end{lemma}

\begin{proof}
Soient $a, b : Z \to X$ deux morphismes de schémas tels que $j \circ
a = j \circ b$. Alors (1)
implique $a = b$ en tant qu'applications
continues sous-jacentes entre espaces topologiques.
Pour tout $z \in Z$, nous
avons $a^\sharp_z \circ j^\sharp_{a(z)} = b^\sharp_z \circ j^\sharp_{b(z)}$ en
tant qu'applications $\mathcal{O}_{Y, j(a(z))} \to \mathcal{O}_{Z,
z}$. La surjectivité des applications
$j^\sharp_x$ impose $a^\sharp_z = b^\sharp_z$, $\forall z \in Z$. Cela
implique que $a^\sharp = b^\sharp$. Ainsi,
nous concluons $a = b$ en tant que morphismes de schémas
comme voulu.
\end{proof}

\begin{lemma}
\label{schemes-lemma-immersions-monomorphisms}
Une immersion de schémas est un monomorphisme.
En particulier, toute immersion est séparée.
\end{lemma}

\begin{proof}
Nous pouvons le vérifier en constatant que le critère du
lemme \ref{schemes-lemma-injective-points-surjective-stalks} s'applique. Plus
élégamment peut-être, nous pouvons utiliser que
les lemmes \ref{schemes-lemma-restrict-map-to-opens} et
\ref{schemes-lemma-characterize-closed-subspace} impliquent que les
immersions ouvertes et fermées sont des monomorphismes et donc que
toute immersion (qui est un composé de celles-ci)
est un monomorphisme.
\end{proof}

\begin{lemma}
\label{schemes-lemma-subscheme-of-separated-scheme}
Soit $f : X \to S$ un morphisme séparé. Tout
sous-schéma localement fermé $Z \subset X$ est séparé sur $S$.
\end{lemma}

\begin{proof}
Résulte du lemme \ref{schemes-lemma-immersions-monomorphisms}
et du fait que le composé de morphismes séparés est séparé
(lemme \ref{schemes-lemma-separated-permanence}).
\end{proof}

\begin{example}
\label{schemes-example-Q-over-Z}
Le morphisme $\Spec(\mathbf{Q}) \to \Spec(\mathbf{Z})$ est
un monomorphisme. Cela est vrai parce que
$\mathbf{Q} \otimes_{\mathbf{Z}} \mathbf{Q} = \mathbf{Q}$.
Plus généralement, pour tout schéma $S$ et tout point $s \in S$, le
morphisme canonique
$$
\Spec(\mathcal{O}_{S, s}) \longrightarrow S
$$
est un monomorphisme.
\end{example}

\begin{lemma}
\label{schemes-lemma-mono-towards-spec-field}
Soient $k_1, \ldots, k_n$ des corps.
Pour tout monomorphisme de schémas
$X \to \Spec(k_1 \times \ldots \times k_n)$,
il existe un sous-ensemble $I \subset \{1, \ldots, n\}$ tel
que $X \cong \Spec(\prod_{i \in I} k_i)$ en tant
que schémas sur $\Spec(k_1 \times \ldots \times k_n)$.
Plus généralement, si $X = \coprod_{i \in I} \Spec(k_i)$ est
une union disjointe de spectres de corps et que $Y \to X$ est un monomorphisme,
alors il existe un sous-ensemble $J \subset I$ tel
que $Y = \coprod_{i \in J} \Spec(k_i)$.
\end{lemma}

\begin{proof}
Réduisons d'abord au cas $n = 1$ (ou $\# I = 1$) en
prenant les images inverses des sous-schémas
ouverts et fermés $\Spec(k_i)$. Dans ce cas,
$X$ n'a qu'un seul point, donc il est affine. Le problème
algébrique correspondant est le suivant : si $k
\to R$ est un homomorphisme d'algèbres
avec $R \otimes_k R \cong R$, alors $R \cong k$ ou $R = 0$. Ceci
est vrai pour des raisons de dimension.
Voir
également Algèbre, lemme \ref{algebra-lemma-epimorphism-over-field}
\end{proof}










\section{Fonctorialité pour les modules quasi-cohérents}
\label{schemes-section-quasi-coherent}

\noindent
Soit $X$ un schéma. Nous notons $\QCoh(\mathcal{O}_X)$ la
catégorie des $\mathcal{O}_X$-modules quasi-cohérents telle qu'elle est définie dans
Modules, définition \ref{modules-definition-quasi-coherent}. Nous
avons vu dans
la section \ref{schemes-section-quasi-coherent-affine}
que la catégorie $\QCoh(\mathcal{O}_X)$ a de nombreuses bonnes propriétés
lorsque $X$ est affine. Puisque la propriété d'être quasi-cohérent est
locale sur $X$, ces propriétés sont héritées par la catégorie des
faisceaux quasi-cohérents sur tout schéma $X$. Nous les énumérons ici.
\begin{enumerate}
\item Un faisceau de $\mathcal{O}_X$-modules $\mathcal{F}$ est
quasi-cohérent si et seulement si la restriction de $\mathcal{F}$ à
chaque ouvert affine $U = \Spec(R)$ est de la forme
$\widetilde M$ pour un certain $R$-module $M$.
\item Un faisceau de $\mathcal{O}_X$-modules $\mathcal{F}$ est
quasi-cohérent si et seulement si la restriction de $\mathcal{F}$ à
chacun des membres d'un recouvrement ouvert affine est quasi-cohérente. \item
Toute somme directe de faisceaux quasi-cohérents est quasi-cohérente.
\item Toute colimite de faisceaux quasi-cohérents est quasi-cohérente.
\item
\label{schemes-item-kernel}
Le noyau et le conoyau d'un morphisme de faisceaux quasi-cohérents sont
quasi-cohérents.
\item Étant donné une suite exacte courte de $\mathcal{O}_X$-modules
$0 \to \mathcal{F}_1 \to \mathcal{F}_2 \to \mathcal{F}_3 \to 0$, si
deux des trois sont quasi-cohérents, le troisième l'est aussi.
\item Étant donné un morphisme de schémas $f : Y \to X$, l'image
inverse d'un $\mathcal{O}_X$-module quasi-cohérent est un
$\mathcal{O}_Y$-module quasi-cohérent.
Voir Modules, lemme \ref{modules-lemma-pullback-quasi-coherent}.
\item Étant donné deux $\mathcal{O}_X$-modules quasi-cohérents,
le produit tensoriel est quasi-cohérent,
voir Modules, lemme \ref{modules-lemma-tensor-product-permanence}.
\item Étant donné un $\mathcal{O}_X$-module quasi-cohérent $\mathcal{F}$,
les algèbres tensorielle, symétrique et extérieure sur $\mathcal{F}$
sont quasi-cohérentes,
voir Modules, lemme \ref{modules-lemma-whole-tensor-algebra-permanence}.
\item Étant donné deux $\mathcal{O}_X$-modules quasi-cohérents
$\mathcal{F}$ et $\mathcal{G}$ tels que $\mathcal{F}$
soit de présentation finie, le faisceau Hom
interne $\SheafHom_{\mathcal{O}_X}(\mathcal{F}, \mathcal{G})$
est quasi-cohérent,
voir Modules, lemme \ref{modules-lemma-internal-hom-locally-kernel-direct-sum}
et (\ref{schemes-item-kernel}) ci-dessus.
\end{enumerate}
En revanche, l'image directe d'un module
quasi-cohérent n'est pas quasi-cohérente en général.
Voici un cas où elle l'est.

\begin{lemma}
\label{schemes-lemma-push-forward-quasi-coherent}
Soit $f : X \to S$ un morphisme de schémas.
Si $f$ est quasi-compact et quasi-séparé, alors
$f_*$ envoie les $\mathcal{O}_X$-modules quasi-cohérents
en $\mathcal{O}_S$-modules quasi-cohérents.
\end{lemma}

\begin{proof}
La question est locale sur $S$ et donc nous pouvons supposer que
$S$ est affine. Comme $X$ est quasi-compact, nous pouvons
écrire $X = \bigcup_{i = 1}^n U_i$ avec chaque $U_i$ ouvert affine.
Comme $f$ est quasi-séparé, nous pouvons
écrire $U_i \cap U_j = \bigcup_{k = 1}^{n_{ij}} U_{ijk}$ pour des
ouverts affines $U_{ijk}$, voir le lemme \ref{schemes-lemma-characterize-quasi-separated}.
On note $f_i : U_i \to S$ et $f_{ijk} : U_{ijk} \to S$ les
restrictions de $f$. Pour tout ouvert $V$ de $S$ et tout faisceau
$\mathcal{F}$ sur $X$, nous avons
\begin{eqnarray*}
f_*\mathcal{F}(V) & = & \mathcal{F}(f^{-1}V) \\ &
= &
\Ker\left(
\bigoplus\nolimits_i \mathcal{F}(f^{-1}V \cap U_i)
\to
\bigoplus\nolimits_{i, j, k} \mathcal{F}(f^{-1}V \cap U_{ijk})\right) \\
& = &
\Ker\left(
\bigoplus\nolimits_i f_{i, *}(\mathcal{F}|_{U_i})(V)
\to
\bigoplus\nolimits_{i, j, k} f_{ijk, *}(\mathcal{F}|_{U_{ijk}})(V)\right) \\
& = &
\Ker\left(
\bigoplus\nolimits_i f_{i, *}(\mathcal{F}|_{U_i})
\to
\bigoplus\nolimits_{i, j, k} f_{ijk, *}(\mathcal{F}|_{U_{ijk}})\right)(V)
\end{eqnarray*}
En d'autres termes, il existe une suite exacte de faisceaux
$$
0 \to f_*\mathcal{F}
\to \bigoplus f_{i, *}\mathcal{F}_i
\to \bigoplus f_{ijk, *}\mathcal{F}_{ijk}
$$
où $\mathcal{F}_i, \mathcal{F}_{ijk}$ désignent les
restrictions de $\mathcal{F}$ à l'ouvert correspondant. Si
$\mathcal{F}$ est un $\mathcal{O}_X$-module quasi-cohérent alors
$\mathcal{F}_i$ est un $\mathcal{O}_{U_i}$-module quasi-cohérent et
$\mathcal{F}_{ijk}$ est un $\mathcal{O}_{U_{ijk}}$-module quasi-cohérent. Ainsi, par le
lemme \ref{schemes-lemma-widetilde-pullback}, nous voyons que le deuxième et le
troisième terme de la suite exacte sont des $\mathcal{O}_S$-modules
quasi-cohérents. Nous concluons donc que
$f_*\mathcal{F}$ est un $\mathcal{O}_S$-module quasi-cohérent.
\end{proof}

\noindent
En utilisant cela, nous pouvons caractériser les immersions (fermées) de schémas
comme suit.

\begin{lemma}
\label{schemes-lemma-characterize-closed-immersions}
Soit $f : X \to Y$ un morphisme de schémas.
Supposons que
\begin{enumerate}
\item $f$ induise un homéomorphisme de $X$ sur un
sous-ensemble fermé de $Y$,
et \item $f^\sharp : \mathcal{O}_Y \to f_*\mathcal{O}_X$
soit surjectif.
\end{enumerate}
Alors $f$ est une immersion fermée de schémas.
\end{lemma}

\begin{proof}
Supposons (1) et (2). D'après (1), le morphisme $f$ est quasi-compact
(voir Topologie, lemme \ref{topology-lemma-closed-in-quasi-compact}).
Les conditions (1) et (2) impliquent les conditions (1) et (2)
du lemme \ref{schemes-lemma-injective-points-surjective-stalks}.
Par conséquent, $f : X \to Y$ est un monomorphisme. En particulier,
$f$ est séparé, voir le lemme \ref{schemes-lemma-monomorphism-separated}.
Ainsi, le lemme \ref{schemes-lemma-push-forward-quasi-coherent}
ci-dessus s'applique et nous concluons que $f_*\mathcal{O}_X$ est un
$\mathcal{O}_Y$-module quasi-cohérent. Par conséquent, le noyau
de $\mathcal{O}_Y \to f_*\mathcal{O}_X$ est quasi-cohérent
selon le lemme \ref{schemes-lemma-extension-quasi-coherent}. Comme un faisceau
quasi-cohérent est localement engendré par des
sections (voir Modules, définition \ref{modules-definition-quasi-coherent}),
cela implique que $f$ est une immersion fermée,
voir définition \ref{schemes-definition-closed-immersion-locally-ringed-spaces}.
\end{proof}

\noindent
Nous pouvons utiliser ce lemme pour démontrer le lemme suivant.

\begin{lemma}
\label{schemes-lemma-composition-immersion}
Le composé d'immersions de schémas est une immersion, le
composé d'immersions fermées de schémas est une immersion fermée, et le
composé d'immersions ouvertes de schémas est une immersion ouverte.
\end{lemma}

\begin{proof}
Cela est évident dans le cas des immersions ouvertes puisqu'un sous-espace ouvert d'un
sous-espace ouvert est également un sous-espace ouvert.

\medskip\noindent
Supposons que $a : Z \to Y$ et $b : Y \to X$ soient des immersions fermées de schémas.
Nous allons vérifier que $c = b \circ a$ est également une immersion
fermée. L'hypothèse implique que $a$ et $b$ sont des homéomorphismes sur des sous-ensembles
fermés, et donc $c = b \circ a$ est aussi un homéomorphisme sur un sous-ensemble
fermé. De plus, l'application $\mathcal{O}_X \to c_*\mathcal{O}_Z$ est surjective
puisqu'elle se factorise comme la composée des applications
surjectives $\mathcal{O}_X \to b_*\mathcal{O}_Y$
et $b_*\mathcal{O}_Y \to b_*a_*\mathcal{O}_Z$ (surjectives car $b_*$ est exact,
voir Modules, lemme \ref{modules-lemma-i-star-exact}). Ainsi,
d'après le lemme \ref{schemes-lemma-characterize-closed-immersions}
ci-dessus, $c$ est une immersion fermée.

\medskip\noindent
Enfin, nous arrivons au cas des immersions. Supposons
que $a : Z \to Y$ et $b : Y \to X$ soient des immersions de schémas. Cela
signifie qu'il existe des sous-schémas ouverts
$V \subset Y$ et $U \subset X$ tels que $a(Z)
\subset V$, $b(Y) \subset U$ et $a : Z \to V$ et $b : Y \to U$
soient des immersions fermées. Puisque la
topologie sur $Y$ est induite à partir de la topologie sur $U$, nous
pouvons trouver un ouvert $U' \subset U$ tel que $V =
b^{-1}(U')$. Ensuite, nous voyons que $Z \to
V = b^{-1}(U') \to U'$ est un composé d'immersions
fermées et est donc une immersion fermée. Cela prouve
que $Z \to X$ est une immersion, ce qui achève la démonstration.
\end{proof}















% OK, start here.
%

\chapter{Constructions de schémas}



\phantomsection
\label{constructions-section-phantom}


\section{Introduction}
\label{constructions-section-introduction}

\noindent
Le présent chapitre expose diverses méthodes permettant de construire des
schémas à partir d'autres schémas. On pourra consulter \cite{EGA}.





\section{Recollement relatif}
\label{constructions-section-relative-glueing}

\noindent
Le lemme suivant est utile lorsque l'on cherche à construire un schéma $X$
sur $S$ et que l'on sait déjà construire la restriction de $X$ aux ouverts affines
de $S$. L'énoncé est en fait tout à fait général et vaut pour les espaces
(localement) annelés, bien que la démonstration soit rédigée dans le langage
des schémas.

\begin{lemma}
\label{constructions-lemma-relative-glueing}
Soit $S$ un schéma. Soit $\mathcal{B}$ une base de la topologie de $S$.
Supposons données les familles suivantes :
\begin{enumerate}
\item pour tout $U \in \mathcal{B}$, un schéma $f_U : X_U \to U$ sur $U$ ;
\item pour $U, V \in \mathcal{B}$ avec $V \subset U$, un morphisme
$\rho^U_V : X_V \to X_U$ sur $U$.
\end{enumerate}
Supposons que
\begin{enumerate}
\item[(a)] chaque $\rho^U_V$ induise un isomorphisme
$X_V \to f_U^{-1}(V)$ de schémas sur $V$ ;
\item[(b)] pour $W, V, U \in \mathcal{B}$ tels que
$W \subset V \subset U$, on ait $\rho^U_W = \rho^U_V \circ \rho ^V_W$.
\end{enumerate}
Il existe alors un morphisme de schémas
$f : X \to S$ et des isomorphismes
$i_U : f^{-1}(U) \to X_U$ au-dessus de $U \in \mathcal{B}$ tels que,
pour $V, U \in \mathcal{B}$ avec $V \subset U$, la composée
$$
\xymatrix{
X_V \ar[r]^{i_V^{-1}} &
f^{-1}(V) \ar[rr]^{inclusion} & &
f^{-1}(U) \ar[r]^{i_U} &
X_U
}
$$
soit $\rho^U_V$. De plus, $X$ est unique
à unique isomorphisme près sur $S$.
\end{lemma}

\begin{proof}
Nous allons appliquer Schémas, Lemme
\ref{schemes-lemma-glue-functors}.
Définissons d'abord un foncteur contravariant $F$ de la catégorie des
schémas vers celle des ensembles. Pour un schéma $T$, posons
$$
F(T) =
\left\{
\begin{matrix}
(g, \{h_U\}_{U \in \mathcal{B}}),
\ g : T \to S, \ h_U : g^{-1}(U) \to X_U, \\
f_U \circ h_U = g|_{g^{-1}(U)},
\ h_U|_{g^{-1}(V)} = \rho^U_V \circ h_V
\ \forall\ V, U \in \mathcal{B}, V \subset U
\end{matrix}
\right\}.
$$
L'application de restriction $F(T) \to F(T')$ associée à un morphisme
$T' \to T$ est donnée par composition. Pour $W \in \mathcal{B}$,
considérons le sous-foncteur $F_W \subset F$ constitué des systèmes
$(g, \{h_U\})$ tels que $g(T) \subset W$.

\medskip\noindent
Montrons d'abord que $F$ est un faisceau pour la topologie de Zariski.
Soient $T$ un schéma, $T = \bigcup V_i$ un recouvrement ouvert et
$\xi_i \in F(V_i)$ des éléments tels que
$\xi_i|_{V_i \cap V_j} = \xi_j|_{V_i \cap V_j}$.
Écrivons $\xi_i = (g_i, \{h_{i, U}\})$.
Les morphismes $g_i$ se recollent en un unique morphisme global
$g : T \to S$. De plus,
$g^{-1}(U) = \bigcup g_i^{-1}(U)$ ; les morphismes
$h_{i, U} : g_i^{-1}(U) \to X_U$ se recollent donc en un unique morphisme
$h_U : g^{-1}(U) \to X_U$. On vérifie aisément que le système
$(g, \{h_U\})$ appartient à $F(T)$. Le foncteur $F$ satisfait donc la
condition de faisceau pour la topologie de Zariski.

\medskip\noindent
Vérifions ensuite que chaque $F_W$,
$W \in \mathcal{B}$, est représentable.
Nous affirmons que la transformation de foncteurs
$$
F_W \longrightarrow \Mor(-, X_W), \ (g, \{h_U\}) \longmapsto h_W
$$
est un isomorphisme. Soient $T$ un schéma et $\alpha : T \to X_W$ un
morphisme. Posons $g = f_W \circ \alpha$.
Pour $U \in \mathcal{B}$ avec $U \subset W$, définissons
$h_U : g^{-1}(U) \to X_U$ comme la composée
$(\rho^W_U)^{-1} \circ \alpha|_{g^{-1}(U)}$.
C'est possible puisque $\alpha(g^{-1}(U))$ est contenu dans $f_W^{-1}(U)$,
grâce à la condition (a).
$f_U \circ h_U = g|_{g^{-1}(U)}$ pour un tel $U$.
Si, de plus, $V \in \mathcal{B}$ et $V \subset U \subset W$, la condition
(b) donne $\rho^U_V \circ h_V = h_U|_{g^{-1}(V)}$.
Il reste à définir $h_U$ pour un élément quelconque $U \in \mathcal{B}$.
Comme $\mathcal{B}$ est une base de la topologie de $S$, on peut choisir
un recouvrement ouvert $U \cap W = \bigcup U_i$
avec $U_i \in \mathcal{B}$.
Puisque $g$ est à valeurs dans $W$, on a
$g^{-1}(U) = g^{-1}(U \cap W) = \bigcup g^{-1}(U_i)$.
Considérons les morphismes
$h_i = \rho^U_{U_i} \circ h_{U_i} : g^{-1}(U_i) \to X_U$.
La condition (b) permet aisément de vérifier
$h_i|_{g^{-1}(U_i) \cap g^{-1}(U_j)} = h_j|_{g^{-1}(U_i) \cap g^{-1}(U_j)}$.
Ces morphismes se recollent donc en le morphisme cherché
$h_U : g^{-1}(U) \to X_U$. Nous omettons la vérification élémentaire que
$(g, \{h_U\})$ appartient à $F_W(T)$ et s'envoie sur $\alpha$ par
l'application affichée ci-dessus.

\medskip\noindent
Vérifions ensuite que chaque inclusion $F_W \subset F$ est représentable
par des immersions ouvertes. Cela résulte directement des définitions.

\medskip\noindent
Enfin, vérifions que la famille $(F_W)_{W \in \mathcal{B}}$ recouvre $F$.
Cela résulte de la construction et du
fait que $\mathcal{B}$ est une base
de la topologie de $S$.

\medskip\noindent
Soit $X$ un schéma représentant $F$, et soit
$(f, \{i_U\}) \in F(X)$ une « famille universelle ».
Chaque $F_W$ étant représentable par $X_W$ via la transformation de
foncteurs ci-dessus, $i_W : f^{-1}(W) \to X_W$ est un isomorphisme.
Le lemme est démontré.
\end{proof}

\begin{lemma}
\label{constructions-lemma-relative-glueing-sheaves}
Soit $S$ un schéma. Soit $\mathcal{B}$ une base de la topologie de $S$.
Supposons données les familles suivantes :
\begin{enumerate}
\item pour tout $U \in \mathcal{B}$,
un schéma $f_U : X_U \to U$ sur $U$ ;
\item pour tout $U \in \mathcal{B}$, un faisceau quasi-cohérent
$\mathcal{F}_U$ sur $X_U$ ;
\item pour $U, V \in \mathcal{B}$ avec $V \subset U$, un morphisme
$\rho^U_V : X_V \to X_U$ ;
\item pour $U, V \in \mathcal{B}$ avec $V \subset U$, un morphisme
$\theta^U_V : (\rho^U_V)^*\mathcal{F}_U \to \mathcal{F}_V$.
\end{enumerate}
Supposons que
\begin{enumerate}
\item[(a)] chaque $\rho^U_V$ induise un isomorphisme
$X_V \to f_U^{-1}(V)$ de schémas sur $V$ ;
\item[(b)] chaque $\theta^U_V$ soit un isomorphisme ;
\item[(c)] pour $W, V, U \in \mathcal{B}$
avec $W \subset V \subset U$,
on ait $\rho^U_W = \rho^U_V \circ \rho ^V_W$ ;
\item[(d)] pour $W, V, U \in \mathcal{B}$
avec $W \subset V \subset U$,
on ait $\theta^U_W = \theta^V_W \circ (\rho^V_W)^*\theta^U_V$.
\end{enumerate}
Il existe alors un morphisme de schémas $f : X \to S$, un faisceau
quasi-cohérent $\mathcal{F}$ sur $X$ et des isomorphismes
$i_U : f^{-1}(U) \to X_U$ et
$\theta_U : i_U^*\mathcal{F}_U \to \mathcal{F}|_{f^{-1}(U)}$
au-dessus de $U \in \mathcal{B}$ tels que, pour $V, U \in \mathcal{B}$
avec $V \subset U$, la composée
$$
\xymatrix{
X_V \ar[r]^{i_V^{-1}} &
f^{-1}(V) \ar[rr]^{inclusion} & &
f^{-1}(U) \ar[r]^{i_U} &
X_U
}
$$
soit $\rho^U_V$ et que la composée
\begin{equation}
\label{constructions-equation-glue}
(\rho^U_V)^*\mathcal{F}_U
=
(i_V^{-1})^*((i_U^*\mathcal{F}_U)|_{f^{-1}(V)})
\xrightarrow{\theta_U|_{f^{-1}(V)}}
(i_V^{-1})^*(\mathcal{F}|_{f^{-1}(V)})
\xrightarrow{\theta_V^{-1}}
\mathcal{F}_V
\end{equation}
soit $\theta^U_V$. De plus, $(X, \mathcal{F})$ est unique à unique
isomorphisme près sur $S$.
\end{lemma}

\begin{proof}
Le Lemme \ref{constructions-lemma-relative-glueing} fournit le schéma $X$ sur $S$ et les
isomorphismes $i_U$. Posons $\mathcal{F}'_U = i_U^*\mathcal{F}_U$
pour $U \in \mathcal{B}$ ; c'est un
$\mathcal{O}_{f^{-1}(U)}$-module quasi-cohérent. Les morphismes
$$
\mathcal{F}'_U|_{f^{-1}(V)} =
i_U^*\mathcal{F}_U|_{f^{-1}(V)} =
i_V^*(\rho^U_V)^*\mathcal{F}_U \xrightarrow{i_V^*\theta^U_V}
i_V^*\mathcal{F}_V = \mathcal{F}'_V
$$
définissent des isomorphismes
$(\theta')^U_V : \mathcal{F}'_U|_{f^{-1}(V)} \to \mathcal{F}'_V$
pour $V \subset U$ dans $\mathcal{B}$.
La condition (d) exprime précisément
leur compatibilité pour un triplet $W \subset V \subset U$ d'éléments de
$\mathcal{B}$. On peut ainsi définir des isomorphismes
$$
\varphi_{12} :
\mathcal{F}'_{U_1}|_{f^{-1}(U_1 \cap U_2)}
\longrightarrow
\mathcal{F}'_{U_2}|_{f^{-1}(U_1 \cap U_2)}
$$
pour $U_1, U_2 \in \mathcal{B}$, en recouvrant l'intersection
$U_1 \cap U_2 = \bigcup V_j$ par des éléments $V_j$ de $\mathcal{B}$ et en
posant
$$
\varphi_{12}|_{V_j} =
\left((\theta')^{U_2}_{V_j}\right)^{-1}
\circ
(\theta')^{U_1}_{V_j}.
$$
Nous omettons la vérification que ces morphismes se recollent bien en un
morphisme $\varphi_{12}$ et satisfont
la condition de cocycle d'une donnée
de recollement de faisceaux (au sens de Faisceaux,
section \ref{sheaves-section-glueing-sheaves}). Par Faisceaux,
Lemme \ref{sheaves-lemma-glue-sheaves}, on obtient le faisceau
$\mathcal{F}$ sur $X$. Nous omettons la vérification de
(\ref{constructions-equation-glue}).
\end{proof}

\begin{remark}
\label{constructions-remark-relative-glueing-functorial}
Les constructions des Lemmes \ref{constructions-lemma-relative-glueing} et
\ref{constructions-lemma-relative-glueing-sheaves} sont fonctorielles.
Supposons en effet données deux familles
$(f_U : X_U \to U, \rho^U_V)$ et
$(g_U : Y_U \to U, \sigma^U_V)$ comme dans le
Lemme \ref{constructions-lemma-relative-glueing}. Supposons que, pour chaque
$U \in \mathcal{B}$, on dispose d'un morphisme $h_U : X_U \to Y_U$ sur
$U$, compatible avec les restrictions $\rho^U_V$ et $\sigma^U_V$.
La fonctorialité signifie que l'on obtient un morphisme de schémas
$h : X \to Y$ sur $S$ dont les restrictions sont les $h_U$, où
$f : X \to S$ est construit à partir de
$(f_U : X_U \to U, \rho^U_V)$ et $g : Y \to S$ à partir de
$(g_U : Y_U \to U, \sigma^U_V)$.

\medskip\noindent
De même, supposons données deux familles
$(f_U : X_U \to U, \mathcal{F}_U, \rho^U_V, \theta^U_V)$ et
$(g_U : Y_U \to U, \mathcal{G}_U, \sigma^U_V, \eta^U_V)$ comme dans le
Lemme \ref{constructions-lemma-relative-glueing-sheaves}.
Supposons que, pour chaque
$U \in \mathcal{B}$, on dispose d'un morphisme
$h_U : X_U \to Y_U$ sur $U$, compatible
avec $\rho^U_V$ et $\sigma^U_V$,
et d'un morphisme $\tau_U : h_U^*\mathcal{G}_U \to \mathcal{F}_U$
compatible avec $\theta^U_V$ et $\eta^U_V$.
La fonctorialité signifie que ces données définissent un morphisme de
schémas $h : X \to Y$ sur $S$ dont les
restrictions sont les $h_U$, ainsi
qu'un morphisme $h^*\mathcal{G} \to \mathcal{F}$ dont les restrictions
sont les $h_U$. Ici, $(f : X \to S, \mathcal{F})$ est construit à
partir de $(f_U : X_U \to U, \mathcal{F}_U, \rho^U_V, \theta^U_V)$ et
$(g : Y \to S, \mathcal{G})$ à partir de
$(g_U : Y_U \to U, \mathcal{G}_U, \sigma^U_V, \eta^U_V)$.

\medskip\noindent
Nous omettons les vérifications, ainsi qu'une formulation convenable de
l'« équivalence de catégories » entre données de recollement relatif et
objets relatifs.
\end{remark}















\section{Spectre relatif par recollement}
\label{constructions-section-spec-via-glueing}

\begin{situation}
\label{constructions-situation-relative-spec}
Ici, $S$ est un schéma et $\mathcal{A}$ une $\mathcal{O}_S$-algèbre
quasi-cohérente. Cela signifie que $\mathcal{A}$
est un faisceau de $\mathcal{O}_S$-algèbres qui est
quasi-cohérent comme $\mathcal{O}_S$-module.
\end{situation}

\noindent
Dans cette section, nous indiquons comment construire un morphisme de schémas
$$
\underline{\Spec}_S(\mathcal{A}) \longrightarrow S
$$
en recollant les spectres $\Spec(\Gamma(U, \mathcal{A}))$, où $U$ parcourt
les ouverts affines de $S$. Montrons d'abord que les spectres des algèbres
de sections de $\mathcal{A}$ sur ces ouverts affines forment une famille
de schémas satisfaisant les conditions
du Lemme \ref{constructions-lemma-relative-glueing}.

\begin{lemma}
\label{constructions-lemma-spec-inclusion}
Dans la Situation \ref{constructions-situation-relative-spec}, soient
$U \subset U' \subset S$ des ouverts affines. Posons
$A = \mathcal{A}(U)$ et $A' = \mathcal{A}(U')$.
L'homomorphisme d'anneaux $A' \to A$ induit un morphisme
$\Spec(A) \to \Spec(A')$, et le diagramme
$$
\xymatrix{
\Spec(A) \ar[r] \ar[d] &
\Spec(A') \ar[d] \\
U \ar[r] &
U'
}
$$
est cartésien.
\end{lemma}

\begin{proof}
Posons $R = \mathcal{O}_S(U)$ et $R' = \mathcal{O}_S(U')$.
L'homomorphisme $R \otimes_{R'} A' \to A$ est un isomorphisme puisque
$\mathcal{A}$ est quasi-cohérente (voir, par exemple, Schémas,
Lemme \ref{schemes-lemma-widetilde-pullback}).
Le résultat découle de la
description du produit fibré de schémas affines dans Schémas,
Lemme \ref{schemes-lemma-fibre-product-affine-schemes}.
\end{proof}

\noindent
En particulier, le morphisme $\Spec(A) \to \Spec(A')$ du lemme est une
immersion ouverte.

\begin{lemma}
\label{constructions-lemma-transitive-spec}
Dans la Situation \ref{constructions-situation-relative-spec}, soient
$U \subset U' \subset U'' \subset S$ des ouverts affines. Posons
$A = \mathcal{A}(U)$, $A' = \mathcal{A}(U')$,
$A'' = \mathcal{A}(U'')$.
La composée des morphismes $\Spec(A) \to \Spec(A')$ et
$\Spec(A') \to \Spec(A'')$ du Lemme
\ref{constructions-lemma-spec-inclusion} est le
morphisme $\Spec(A) \to \Spec(A'')$
du Lemme \ref{constructions-lemma-spec-inclusion}.
\end{lemma}

\begin{proof}
Cela résulte de ce que $A'' \to A$ est la composée de $A'' \to A'$ et de
$A' \to A$, puisque $\mathcal{A}$ est un faisceau.
\end{proof}

\begin{lemma}
\label{constructions-lemma-glue-relative-spec}
Dans la Situation \ref{constructions-situation-relative-spec}, il existe un morphisme
de schémas
$$
\pi : \underline{\Spec}_S(\mathcal{A}) \longrightarrow S
$$
ayant les propriétés suivantes :
\begin{enumerate}
\item pour tout ouvert affine $U \subset S$, il existe un isomorphisme
$i_U : \pi^{-1}(U) \to \Spec(\mathcal{A}(U))$ sur $U$ ;
\item pour des ouverts affines $U \subset U' \subset S$, la composée
$$
\xymatrix{
\Spec(\mathcal{A}(U)) \ar[r]^{i_U^{-1}} &
\pi^{-1}(U) \ar[rr]^{inclusion} & &
\pi^{-1}(U') \ar[r]^{i_{U'}} &
\Spec(\mathcal{A}(U'))
}
$$
est l'immersion ouverte du Lemme \ref{constructions-lemma-spec-inclusion}.
\end{enumerate}
De plus, $\underline{\Spec}_S(\mathcal{A})$ est unique à unique
isomorphisme près sur $S$.
\end{lemma}

\begin{proof}
Cela résulte immédiatement des
Lemmes \ref{constructions-lemma-relative-glueing},
\ref{constructions-lemma-spec-inclusion} et
\ref{constructions-lemma-transitive-spec}.
L'unicité est énoncée dans la dernière phrase du
Lemme \ref{constructions-lemma-relative-glueing}.
\end{proof}












\section{Le spectre relatif comme foncteur}
\label{constructions-section-spec}

\noindent
Plaçons-nous dans la Situation \ref{constructions-situation-relative-spec}, c'est-à-dire que
$S$ est un schéma et que $\mathcal{A}$ est un faisceau quasi-cohérent de
$\mathcal{O}_S$-algèbres.

\medskip\noindent
Pour tout $f : T \to S$, l'image inverse
$f^*\mathcal{A}$ est un faisceau quasi-cohérent de $\mathcal{O}_T$-algèbres.
Nous allons considérer les paires $(f : T \to S, \varphi)$ où
$f$ est un morphisme de schémas et $\varphi : f^*\mathcal{A} \to \mathcal{O}_T$
un morphisme de $\mathcal{O}_T$-algèbres. Remarquons que cela revient à se
donner un homomorphisme de $f^{-1}\mathcal{O}_S$-algèbres
$\varphi : f^{-1}\mathcal{A} \to \mathcal{O}_T$ ; voir
Faisceaux, Lemme \ref{sheaves-lemma-adjointness-tensor-restrict}.
Cela revient également à se donner un
homomorphisme de $\mathcal{O}_S$-algèbres
$\varphi : \mathcal{A} \to f_*\mathcal{O}_T$ ; voir
Faisceaux, Lemme \ref{sheaves-lemma-adjoint-push-pull-modules}.
Nous utiliserons sans autre mention ces trois façons d'envisager $\varphi$.

\medskip\noindent
Étant donnés une telle paire
$(f : T \to S, \varphi)$ et un morphisme $a : T' \to T$, on obtient
une seconde paire $(f' = f \circ a, \varphi' = a^*\varphi)$, appelée
l'image inverse de $(f, \varphi)$. On peut notamment décrire
$\varphi' = a^*\varphi$ comme la composée
$\mathcal{A} \to f_*\mathcal{O}_T \to f'_*\mathcal{O}_{T'}$,
où la seconde application est $f_*a^\sharp$ avec
$a^\sharp : \mathcal{O}_T \to a_*\mathcal{O}_{T'}$.
On a ainsi défini un foncteur
\begin{eqnarray}
\label{constructions-equation-spec}
F : \Sch^{opp} & \longrightarrow & \textit{Ensembles} \\
T & \longmapsto & F(T) = \{\text{pairs }(f, \varphi) \text{ as above}\}
\nonumber
\end{eqnarray}

\begin{lemma}
\label{constructions-lemma-spec-base-change}
Dans la Situation \ref{constructions-situation-relative-spec},
soit $F$ le foncteur
associé ci-dessus à $(S, \mathcal{A})$.
Soit $g : S' \to S$ un morphisme de schémas.
Posons $\mathcal{A}' = g^*\mathcal{A}$ et soit $F'$ le
foncteur associé ci-dessus à $(S', \mathcal{A}')$.
Il existe alors un isomorphisme canonique
$$
F' \cong h_{S'} \times_{h_S} F
$$
de foncteurs.
\end{lemma}

\begin{proof}
Se donner une paire $(f' : T \to S', \varphi' : (f')^*\mathcal{A}' \to \mathcal{O}_T)$
revient à se donner une paire $(f, \varphi : f^*\mathcal{A} \to \mathcal{O}_T)$
ainsi qu'une factorisation de $f$ sous la forme
$f = g \circ f'$. En effet, avec ces notations,
$(f')^* \mathcal{A}' = (f')^*g^*\mathcal{A} = f^*\mathcal{A}$.
D'où le lemme.
\end{proof}

\begin{lemma}
\label{constructions-lemma-spec-affine}
Dans la Situation \ref{constructions-situation-relative-spec}, soit $F$ le foncteur
associé ci-dessus à $(S, \mathcal{A})$.
Si $S$ est affine, alors $F$ est représentable par le
schéma affine $\Spec(\Gamma(S, \mathcal{A}))$.
\end{lemma}

\begin{proof}
Écrivons $S = \Spec(R)$ et $A = \Gamma(S, \mathcal{A})$.
Alors $A$ est une $R$-algèbre et $\mathcal{A} = \widetilde A$.
L'homomorphisme d'anneaux $R \to A$ donne une application canonique
$$
f_{univ} : \Spec(A)
\longrightarrow
S = \Spec(R).
$$
On a
$f_{univ}^*\mathcal{A} =  \widetilde{A \otimes_R A}$
par Schémas, Lemme \ref{schemes-lemma-widetilde-pullback}.
Il existe donc une application canonique
$$
\varphi_{univ} :
f_{univ}^*\mathcal{A} = \widetilde{A \otimes_R A}
\longrightarrow
\widetilde A = \mathcal{O}_{\Spec(A)}
$$
provenant de l'application de $A$-modules $A \otimes_R A \to A$,
$a \otimes a' \mapsto aa'$. Nous affirmons que la paire
$(f_{univ}, \varphi_{univ})$ représente $F$ dans ce cas.
Autrement dit, pour tout schéma $T$, l'application
$$
\Mor(T, \Spec(A)) \longrightarrow \{\text{pairs } (f, \varphi)\},\quad
a \longmapsto (f_{univ} \circ a, a^*\varphi_{univ})
$$
est bijective.

\medskip\noindent
Construisons l'application inverse.
À toute paire $(f : T \to S, \varphi)$
est associé l'homomorphisme d'anneaux
$$
\xymatrix{
A = \Gamma(S, \mathcal{A}) \ar[r]^{f^*} &
\Gamma(T, f^*\mathcal{A}) \ar[r]^{\varphi} &
\Gamma(T, \mathcal{O}_T)
}
$$
Celui-ci induit un morphisme de schémas $T \to \Spec(A)$
par Schémas, Lemme \ref{schemes-lemma-morphism-into-affine}.

\medskip\noindent
Nous omettons de vérifier que cette application est inverse de celle
affichée ci-dessus.
\end{proof}

\begin{lemma}
\label{constructions-lemma-spec}
Dans la Situation \ref{constructions-situation-relative-spec},
le foncteur $F$ est représentable par un schéma.
\end{lemma}

\begin{proof}
Nous allons appliquer Schémas, Lemme \ref{schemes-lemma-glue-functors}.

\medskip\noindent
Vérifions d'abord que $F$ satisfait la condition de faisceau pour la
topologie de Zariski. Supposons que $T$ soit un schéma,
que $T = \bigcup_{i \in I} U_i$ soit un recouvrement ouvert
et que $(f_i, \varphi_i) \in F(U_i)$ vérifient
$(f_i, \varphi_i)|_{U_i \cap U_j} = (f_j, \varphi_j)|_{U_i \cap U_j}$.
Il en résulte que les morphismes $f_i : U_i \to S$
se recollent en un morphisme de schémas $f : T \to S$ tel que
$f|_{U_i} = f_i$ ; voir Schémas, section \ref{schemes-section-glueing-schemes}.
Ainsi, $f_i^*\mathcal{A} = f^*\mathcal{A}|_{U_i}$ et, par hypothèse, les
morphismes $\varphi_i$ coïncident sur $U_i \cap U_j$. Ils se recollent donc,
par Faisceaux, section \ref{sheaves-section-glueing-sheaves}, en un
morphisme de $\mathcal{O}_T$-algèbres $f^*\mathcal{A} \to \mathcal{O}_T$.
Cela prouve que $F$ satisfait la condition de faisceau pour la topologie
de Zariski.

\medskip\noindent
Soit $S = \bigcup_{i \in I} U_i$ un recouvrement par des ouverts affines.
Soit $F_i \subset F$ le sous-foncteur formé des
paires $(f : T \to S, \varphi)$ telles que
$f(T) \subset U_i$.

\medskip\noindent
Il faut montrer que chaque $F_i$ est représentable.
C'est le cas, car $F_i$ s'identifie
au foncteur associé à $U_i$ muni de
la $\mathcal{O}_{U_i}$-algèbre quasi-cohérente $\mathcal{A}|_{U_i}$,
par le Lemme \ref{constructions-lemma-spec-base-change}.
Le résultat découle donc du Lemme \ref{constructions-lemma-spec-affine}.

\medskip\noindent
Montrons ensuite que $F_i \subset F$ est représentable par des immersions ouvertes.
Soit $(f : T \to S, \varphi) \in F(T)$. Posons $V_i = f^{-1}(U_i)$.
La définition de $F_i$ montre que, pour tout $a : T' \to T$,
on a $a^*(f, \varphi) \in F_i(T')$ si et seulement si $a(T') \subset V_i$.
C'est précisément ce qu'il fallait démontrer.

\medskip\noindent
Enfin, il faut montrer que la famille $(F_i)_{i \in I}$
recouvre $F$. Soit $(f : T \to S, \varphi) \in F(T)$.
Posons $V_i = f^{-1}(U_i)$. Puisque $S = \bigcup_{i \in I} U_i$
est un recouvrement ouvert de $S$, on voit que $T = \bigcup_{i \in I} V_i$
est un recouvrement ouvert de $T$. De plus, $(f, \varphi)|_{V_i} \in F_i(V_i)$.
Le lemme est ainsi démontré.
\end{proof}

\begin{lemma}
\label{constructions-lemma-glueing-gives-functor-spec}
Dans la Situation \ref{constructions-situation-relative-spec}, le schéma
$\pi : \underline{\Spec}_S(\mathcal{A}) \to S$
construit dans le Lemme \ref{constructions-lemma-glue-relative-spec} et le schéma
représentant le foncteur $F$ sont canoniquement isomorphes comme schémas
au-dessus de $S$.
\end{lemma}

\begin{proof}
Soit $X \to S$ le schéma représentant le foncteur $F$.
Considérons le faisceau de $\mathcal{O}_S$-algèbres
$\mathcal{R} = \pi_*\mathcal{O}_{\underline{\Spec}_S(\mathcal{A})}$.
Par construction de $\underline{\Spec}_S(\mathcal{A})$,
on dispose d'isomorphismes $\mathcal{A}(U) \to \mathcal{R}(U)$
pour tout ouvert affine $U \subset S$ ; cela résulte de la partie (1) du
Lemme \ref{constructions-lemma-glue-relative-spec}.
Pour des ouverts $U \subset U' \subset S$, ces isomorphismes sont
compatibles aux applications de restriction ; cela résulte de la partie (2) du
Lemme \ref{constructions-lemma-glue-relative-spec}.
Par Faisceaux, Lemme \ref{sheaves-lemma-restrict-basis-equivalence-modules},
ils proviennent donc d'un isomorphisme de $\mathcal{O}_S$-algèbres
$\varphi : \mathcal{A} \to \mathcal{R}$. On obtient ainsi un élément
$(\pi, \varphi) \in F(\underline{\Spec}_S(\mathcal{A}))$.
Puisque $X$ représente le foncteur $F$, il lui correspond un
morphisme de schémas $can : \underline{\Spec}_S(\mathcal{A}) \to X$
au-dessus de $S$.

\medskip\noindent
Soit $U \subset S$ un ouvert affine. Soit $F_U \subset F$ le
sous-foncteur de $F$ correspondant aux paires $(f, \varphi)$ définies sur des
schémas $T$ tels que $f(T) \subset U$. Le changement de base
$X_U$ représente manifestement $F_U$. D'autre part, d'après le
Lemme \ref{constructions-lemma-spec-affine}, $F_U$ est représenté par
$\Spec(\mathcal{A}(U)) = \pi^{-1}(U)$. Autrement dit, $X_U \cong \pi^{-1}(U)$.
Nous omettons de vérifier que cette identification est induite par le
changement de base du morphisme $can$ à $U$.
\end{proof}

\begin{definition}
\label{constructions-definition-relative-spec}
Soit $S$ un schéma. Soit $\mathcal{A}$ un faisceau quasi-cohérent de
$\mathcal{O}_S$-algèbres. Le {\it spectre relatif de $\mathcal{A}$ au-dessus de
$S$}, ou simplement le {\it spectre de $\mathcal{A}$ au-dessus de $S$}, est le schéma
construit dans le Lemme \ref{constructions-lemma-glue-relative-spec} qui représente le
foncteur $F$ (\ref{constructions-equation-spec}) ; voir le
Lemme \ref{constructions-lemma-glueing-gives-functor-spec}.
On le note $\pi : \underline{\Spec}_S(\mathcal{A}) \to S$.
La « famille universelle » est un morphisme de $\mathcal{O}_S$-algèbres
$$
\mathcal{A}
\longrightarrow
\pi_*\mathcal{O}_{\underline{\Spec}_S(\mathcal{A})}
$$
\end{definition}

\noindent
Le lemme suivant affirme notamment que la formation du spectre relatif
commute aux changements de base.

\begin{lemma}
\label{constructions-lemma-spec-properties}
Soit $S$ un schéma. Soit $\mathcal{A}$ un faisceau quasi-cohérent
de $\mathcal{O}_S$-algèbres. Soit
$\pi : \underline{\Spec}_S(\mathcal{A}) \to S$
le spectre relatif de $\mathcal{A}$ au-dessus de $S$.
\begin{enumerate}
\item Pour tout ouvert affine $U \subset S$, l'image inverse
$\pi^{-1}(U)$ est affine.
\item Pour tout morphisme $g : S' \to S$, on a
$S' \times_S \underline{\Spec}_S(\mathcal{A}) =
\underline{\Spec}_{S'}(g^*\mathcal{A})$.
\item
L'application universelle
$$
\mathcal{A}
\longrightarrow
\pi_*\mathcal{O}_{\underline{\Spec}_S(\mathcal{A})}
$$
est un isomorphisme de $\mathcal{O}_S$-algèbres.
\end{enumerate}
\end{lemma}

\begin{proof}
La partie (1) résulte de la description du spectre relatif par recollement ;
voir le Lemme \ref{constructions-lemma-glue-relative-spec}.
La partie (2) découle immédiatement du Lemme \ref{constructions-lemma-spec-base-change}.
La partie (3) est locale sur $S$ et, lorsque $S$ est affine, elle résulte
par exemple du Lemme \ref{constructions-lemma-spec-affine}.
\end{proof}

\begin{lemma}
\label{constructions-lemma-canonical-morphism}
Soit $f : X \to S$ un morphisme quasi-compact et quasi-séparé de schémas.
Par Schémas, Lemme \ref{schemes-lemma-push-forward-quasi-coherent},
le faisceau $f_*\mathcal{O}_X$ est un faisceau quasi-cohérent de
$\mathcal{O}_S$-algèbres. Il existe un morphisme canonique
$$
can : X \longrightarrow \underline{\Spec}_S(f_*\mathcal{O}_X)
$$
de schémas au-dessus de $S$.
Pour tout ouvert affine $U \subset S$, la restriction $can|_{f^{-1}(U)}$
s'identifie au morphisme canonique
$$
f^{-1}(U) \longrightarrow \Spec(\Gamma(f^{-1}(U), \mathcal{O}_X))
$$
provenant de Schémas, Lemme \ref{schemes-lemma-morphism-into-affine}.
\end{lemma}

\begin{proof}
Par la définition de $\underline{\Spec}$ comme schéma représentant le
foncteur $F$, le morphisme provient de l'application canonique
$\varphi : f^*f_*\mathcal{O}_X \to \mathcal{O}_X$ (qui, par adjonction entre
image inverse et image directe, correspond à
$\text{id} : f_*\mathcal{O}_X \to f_*\mathcal{O}_X$).
L'assertion concernant la restriction à $f^{-1}(U)$
résulte de la description du spectre relatif au-dessus d'un ouvert affine ;
voir le Lemme \ref{constructions-lemma-spec-affine}.
\end{proof}











\section{Espace affine de dimension n}
\label{constructions-section-affine-n-space}

\noindent
Comme application du spectre relatif, définissons l'espace affine de
dimension $n$ au-dessus d'un schéma de
base $S$. Pour tout entier $n \geq 0$,
on peut considérer le faisceau quasi-cohérent de $\mathcal{O}_S$-algèbres
$\mathcal{O}_S[T_1, \ldots, T_n]$. Il est quasi-cohérent car, comme faisceau
de $\mathcal{O}_S$-modules, il est simplement la somme directe de copies de
$\mathcal{O}_S$ indexées par les multi-indices.

\begin{definition}
\label{constructions-definition-affine-n-space}
Soient $S$ un schéma et $n \geq 0$.
Le schéma
$$
\mathbf{A}^n_S =
\underline{\Spec}_S(\mathcal{O}_S[T_1, \ldots, T_n])
$$
au-dessus de $S$ est appelé {\it espace affine de dimension $n$ au-dessus de $S$}.
Si $S = \Spec(R)$ est affine, on l'appelle aussi
{\it espace affine de dimension $n$ au-dessus de $R$} et on le note $\mathbf{A}^n_R$.
\end{definition}

\noindent
Remarquons que $\mathbf{A}^n_R = \Spec(R[T_1, \ldots, T_n])$.
Pour tout morphisme de schémas $g : S' \to S$, on a
$g^*\mathcal{O}_S[T_1, \ldots, T_n] = \mathcal{O}_{S'}[T_1, \ldots, T_n]$ ;
ainsi, $\mathbf{A}^n_{S'} = S' \times_S \mathbf{A}^n_S$ est le changement
de base. Une définition équivalente de l'espace affine de dimension $n$
est donc donnée par la formule
$$
\mathbf{A}^n_S = S \times_{\Spec(\mathbf{Z})} \mathbf{A}^n_{\mathbf{Z}}.
$$
En outre, un morphisme d'un $S$-schéma $f : X \to S$
vers $\mathbf{A}^n_S$ est déterminé par un homomorphisme de
$\mathcal{O}_S$-algèbres
$\mathcal{O}_S[T_1, \ldots, T_n] \to f_*\mathcal{O}_X$.
Cela revient manifestement à se donner les images des $T_i$.
Autrement dit, se donner un morphisme de $X$ vers $\mathbf{A}^n_S$ au-dessus de $S$
revient à se donner $n$ éléments
$h_1, \ldots, h_n \in \Gamma(X, \mathcal{O}_X)$.

















\section{Fibrés vectoriels}
\label{constructions-section-vector-bundle}

\noindent
Soit $S$ un schéma.
Soit $\mathcal{E}$ un faisceau quasi-cohérent de $\mathcal{O}_S$-modules.
Par Modules, Lemme \ref{modules-lemma-whole-tensor-algebra-permanence},
l'algèbre symétrique $\text{Sym}(\mathcal{E})$ de
$\mathcal{E}$ sur $\mathcal{O}_S$
est un faisceau quasi-cohérent de $\mathcal{O}_S$-algèbres.
On peut donc lui appliquer la construction du spectre relatif.

\begin{definition}
\label{constructions-definition-vector-bundle}
Soit $S$ un schéma. Soit $\mathcal{E}$ un
$\mathcal{O}_S$-module quasi-cohérent\footnote{Le lecteur pourrait s'attendre
à ce que l'on suppose ici $\mathcal{E}$ localement libre de type fini. Nous ne
le faisons pas, afin de rester cohérents avec \cite[II, définition 1.7.8]{EGA}.}.
Le {\it fibré vectoriel associé à $\mathcal{E}$} est
$$
\mathbf{V}(\mathcal{E}) = \underline{\Spec}_S(\text{Sym}(\mathcal{E})).
$$
\end{definition}

\noindent
Le fibré vectoriel associé à $\mathcal{E}$ est muni d'une structure
supplémentaire. En effet, on a une graduation
$$
\pi_*\mathcal{O}_{\mathbf{V}(\mathcal{E})} =
\bigoplus\nolimits_{n \geq 0} \text{Sym}^n(\mathcal{E}).
$$
qui fait de $\pi_*\mathcal{O}_{\mathbf{V}(\mathcal{E})}$
une $\mathcal{O}_S$-algèbre graduée. Réciproquement, on retrouve
$\mathcal{E}$ comme sa partie de degré $1$.
Nous définissons donc un fibré vectoriel abstrait comme suit.

\begin{definition}
\label{constructions-definition-abstract-vector-bundle}
Soit $S$ un schéma. Un {\it fibré vectoriel $\pi : V \to S$ au-dessus de $S$}
est un morphisme affine de schémas tel que $\pi_*\mathcal{O}_V$ soit muni
d'une structure de $\mathcal{O}_S$-algèbre graduée
$\pi_*\mathcal{O}_V = \bigoplus\nolimits_{n \geq 0} \mathcal{E}_n$
telle que $\mathcal{E}_0 = \mathcal{O}_S$ et que les applications
$$
\text{Sym}^n(\mathcal{E}_1) \longrightarrow \mathcal{E}_n
$$
soient des isomorphismes pour tout $n \geq 0$. Un {\it morphisme de fibrés
vectoriels au-dessus de $S$} est un morphisme $f : V \to V'$ tel que
l'application induite
$$
f^* : \pi'_*\mathcal{O}_{V'} \longrightarrow \pi_*\mathcal{O}_V
$$ soit compatible avec les graduations données.
\end{definition}

\noindent
Un exemple de fibré vectoriel sur $S$ est l'espace affine de dimension $n$,
$\mathbf{A}^n_S$, sur $S$ ;
voir la Définition \ref{constructions-definition-affine-n-space}. En effet,
$\mathcal{O}_S[T_1, \ldots, T_n] = \text{Sym}(\mathcal{O}_S^{\oplus n})$.

\begin{lemma}
\label{constructions-lemma-category-vector-bundles}
La catégorie des fibrés vectoriels au-dessus d'un schéma $S$ est
anti-équivalente à la catégorie des $\mathcal{O}_S$-modules quasi-cohérents.
\end{lemma}

\begin{proof}
Omis. Indication : dans un sens, on utilise le foncteur
$\underline{\Spec}_S(\text{Sym}^*_{\mathcal{O}_S}(-))$,
et dans l'autre le foncteur
$(\pi : V \to S) \leadsto (\pi_*\mathcal{O}_V)_1$, où l'indice
signifie que l'on prend la partie de degré $1$.
\end{proof}




\section{Cônes}
\label{constructions-section-cone}

\noindent
En géométrie algébrique, les cônes correspondent aux algèbres graduées.
Selon nos conventions, un anneau ou une algèbre gradué $A$ est muni d'une graduation
$A = \bigoplus_{d \geq 0} A_d$ par les entiers positifs ou nuls ; voir
Algèbre, section \ref{algebra-section-graded}.

\begin{definition}
\label{constructions-definition-cone}
Soit $S$ un schéma. Soit $\mathcal{A}$ une
$\mathcal{O}_S$-algèbre graduée quasi-cohérente. Supposons que
$\mathcal{O}_S \to \mathcal{A}_0$ soit un isomorphisme\footnote{On impose
souvent que $\mathcal{A}$ soit engendrée par $\mathcal{A}_1$ sur
$\mathcal{O}_S$. Nous ne le supposons pas, afin de rester cohérents avec
\cite[II, (8.3.1)]{EGA}.}.
Le {\it cône associé à $\mathcal{A}$}, ou le
{\it cône affine associé à $\mathcal{A}$}, est
$$
C(\mathcal{A}) = \underline{\Spec}_S(\mathcal{A}).
$$
\end{definition}

\noindent
Le cône associé à un faisceau gradué de $\mathcal{O}_S$-algèbres
est muni d'une structure supplémentaire : on obtient une graduation
$$
\pi_*\mathcal{O}_{C(\mathcal{A})} =
\bigoplus\nolimits_{n \geq 0} \mathcal{A}_n
$$
On peut donc définir un cône abstrait comme suit.

\begin{definition}
\label{constructions-definition-abstract-cone}
Soit $S$ un schéma. Un {\it cône $\pi : C \to S$ au-dessus de $S$} est un
morphisme affine de schémas tel que $\pi_*\mathcal{O}_C$ soit muni d'une
structure de $\mathcal{O}_S$-algèbre graduée
$\pi_*\mathcal{O}_C = \bigoplus\nolimits_{n \geq 0} \mathcal{A}_n$
telle que $\mathcal{A}_0 = \mathcal{O}_S$. Un {\it morphisme de cônes}
de $\pi : C \to S$ vers $\pi' : C' \to S$
est un morphisme $f : C \to C'$ tel que l'application induite
$$
f^* : \pi'_*\mathcal{O}_{C'} \longrightarrow \pi_*\mathcal{O}_C
$$
soit compatible avec les graduations données.
\end{definition}

\noindent
Tout fibré vectoriel est un exemple de cône. En fait, la catégorie des
fibrés vectoriels au-dessus de $S$ est une sous-catégorie pleine de la
catégorie des cônes au-dessus de $S$.










\section{Proj d'un anneau gradué}
\label{constructions-section-proj}

\noindent
Dans cette section, nous construisons le Proj d'un anneau gradué
en suivant \cite[II, section 2]{EGA}.

\medskip\noindent
Soit $S$ un anneau gradué. Considérons l'espace topologique $\text{Proj}(S)$
associé à $S$ ; voir Algèbre, section \ref{algebra-section-proj}.
Nous allons munir cet espace d'un faisceau d'anneaux $\mathcal{O}_{\text{Proj}(S)}$
de sorte que la paire $(\text{Proj}(S), \mathcal{O}_{\text{Proj}(S)})$
ainsi obtenue soit un schéma.

\medskip\noindent
Rappelons que $\text{Proj}(S)$ possède une base d'ouverts $D_{+}(f)$,
$f \in S_d$, $d \geq 1$, que nous appelons {\it ouverts standards} ; voir
Algèbre, section \ref{algebra-section-proj}. Cette terminologie sous-entendra
toujours que $f$ est homogène de degré strictement positif, même si nous
omettons de le préciser. En outre, l'intersection de deux ouverts standards
est encore un ouvert standard :
$D_{+}(f) \cap D_{+}(g) = D_{+}(fg)$ pour $f, g \in S$ homogènes de degré strictement positif.

\begin{lemma}
\label{constructions-lemma-standard-open}
Soit $S$ un anneau gradué. Soit $f \in S$
homogène de degré strictement positif.
\begin{enumerate}
\item Si $g\in S$ est homogène de degré strictement positif
et $D_{+}(g) \subset D_{+}(f)$, alors
\begin{enumerate}
\item $f$ est inversible dans $S_g$, et
$f^{\deg(g)}/g^{\deg(f)}$ est inversible dans $S_{(g)}$,
\item $g^e = af$ pour un certain $e \geq 1$ et un certain $a \in S$ homogène,
\item il existe un homomorphisme canonique de $S$-algèbres $S_f \to S_g$,
\item il existe un homomorphisme canonique de $S_0$-algèbres
$S_{(f)} \to S_{(g)}$, compatible avec l'application $S_f \to S_g$,
\item l'application $S_{(f)} \to S_{(g)}$ induit un isomorphisme
$$
(S_{(f)})_{g^{\deg(f)}/f^{\deg(g)}} \cong S_{(g)},
$$
\item ces applications induisent un diagramme commutatif d'espaces
topologiques
$$
\xymatrix{
D_{+}(g) \ar[d] &
\{\mathbf{Z}\text{-graded primes of }S_g\} \ar[l] \ar[r] \ar[d] &
\Spec(S_{(g)}) \ar[d] \\
D_{+}(f) &
\{\mathbf{Z}\text{-graded primes of }S_f\} \ar[l] \ar[r] &
\Spec(S_{(f)})
}
$$
où les applications horizontales sont des homéomorphismes et les applications
verticales des immersions ouvertes,
\item il existe des applications canoniques compatibles de $S_f$-modules et
de $S_{(f)}$-modules $M_f \to M_g$ et $M_{(f)} \to M_{(g)}$ pour tout
$S$-module gradué $M$, et
\item l'application $M_{(f)} \to M_{(g)}$ induit un isomorphisme
$$
(M_{(f)})_{g^{\deg(f)}/f^{\deg(g)}} \cong M_{(g)}.
$$
\end{enumerate}
\item Tout recouvrement ouvert de $D_{+}(f)$ admet un raffinement qui est un
recouvrement ouvert fini de la forme $D_{+}(f) = \bigcup_{i = 1}^n D_{+}(g_i)$.
\item Soient $g_1, \ldots, g_n \in S$ homogènes de degré strictement positif.
Alors $D_{+}(f) \subset \bigcup D_{+}(g_i)$ si et seulement si
$g_1^{\deg(f)}/f^{\deg(g_1)}, \ldots, g_n^{\deg(f)}/f^{\deg(g_n)}$
engendrent l'idéal unité de $S_{(f)}$.
\end{enumerate}
\end{lemma}

\begin{proof}
Rappelons que $D_{+}(g) = \Spec(S_{(g)})$, l'identification étant donnée par
les homomorphismes d'anneaux $S \to S_g \leftarrow S_{(g)}$ ; voir
Algèbre, Lemme \ref{algebra-lemma-topology-proj}.
Ainsi, $f^{\deg(g)}/g^{\deg(f)}$ est un élément de $S_{(g)}$ qui n'appartient
à aucun idéal premier ; il est donc inversible, par
Algèbre, Lemme \ref{algebra-lemma-Zariski-topology}.
On en déduit (a).
Écrivons l'inverse de $f$ dans $S_g$ sous la forme $a/g^d$.
On peut remplacer $a$ par sa partie homogène de degré
$d\deg(g) - \deg(f)$.
Cela signifie que $g^d - af$ est annulé par une puissance de $g$, d'où
$g^e = af$ pour un certain $a \in S$ homogène de degré
$e\deg(g) - \deg(f)$. Ceci prouve (b).
Pour (c), l'application $S_f \to S_g$ existe, d'après (a), par la propriété
universelle de la localisation ; on peut aussi la définir en envoyant $b/f^n$
sur $a^nb/g^{ne}$. Elle induit manifestement une application entre les
sous-anneaux d'éléments de degré zéro $S_{(f)} \to S_{(g)}$.
De même, on peut définir $M_f \to M_g$ et $M_{(f)} \to M_{(g)}$ en envoyant
$x/f^n$ sur $a^nx/g^{ne}$. Les assertions qui présentent $S_{(g)}$, resp.\ $M_{(g)}$, comme une
localisation principale de $S_{(f)}$, resp.\ $M_{(f)}$, résultent aussitôt des formules précédentes.
Les applications du diagramme commutatif d'espaces topologiques correspondent
aux homomorphismes d'anneaux ci-dessus. Les flèches horizontales sont des
homéomorphismes par Algèbre, Lemme \ref{algebra-lemma-topology-proj}.
Les flèches verticales sont des immersions ouvertes, puisque celle de gauche
est l'inclusion d'un sous-ensemble ouvert.

\medskip\noindent
L'ouvert $D_{+}(f)$ est quasi-compact, car il est homéomorphe à
$\Spec(S_{(f)})$ ; voir Algèbre, Lemme
\ref{algebra-lemma-quasi-compact}.
La deuxième assertion résulte donc directement du fait que les ouverts
standards forment une base de la topologie.

\medskip\noindent
La troisième assertion découle directement de
Algèbre, Lemme \ref{algebra-lemma-Zariski-topology}.
\end{proof}

\noindent
Dans Faisceaux, section \ref{sheaves-section-bases}, nous avons défini
la notion de faisceau sur une base et montré qu'elle est essentiellement
équivalente à celle de faisceau sur l'espace ; voir Faisceaux, Lemmes
\ref{sheaves-lemma-extend-off-basis} et
\ref{sheaves-lemma-extend-off-basis-structures}. En outre, nous avons montré
dans Faisceaux, Lemme
\ref{sheaves-lemma-cofinal-systems-coverings-standard-case}, qu'il suffit de
vérifier la condition de faisceau sur un système cofinal de recouvrements
ouverts de chaque ouvert standard. D'après
le lemme précédent, il suffit donc
de la vérifier sur les recouvrements finis par des ouverts standards.

\begin{definition}
\label{constructions-definition-standard-covering}
Soit $S$ un anneau gradué. Supposons que
$D_{+}(f) \subset \text{Proj}(S)$ soit un ouvert standard. Un
{\it recouvrement ouvert standard} de $D_{+}(f)$ est un recouvrement
$D_{+}(f) = \bigcup_{i = 1}^n D_{+}(g_i)$, où
$g_1, \ldots, g_n \in S$ sont homogènes de degré strictement positif.
\end{definition}

\noindent
Soit $S$ un anneau gradué et soit $M$ un $S$-module gradué. Nous allons
définir un préfaisceau $\widetilde M$ sur la base des ouverts standards.
Supposons que $U \subset \text{Proj}(S)$ soit un ouvert standard.
Si $f, g \in S$ sont homogènes de
degré strictement positif et vérifient
$D_{+}(f) = D_{+}(g)$, le Lemme \ref{constructions-lemma-standard-open} fournit des
applications canoniques $M_{(f)} \to M_{(g)}$ et
$M_{(g)} \to M_{(f)}$, inverses l'une de l'autre. On peut donc choisir
n'importe quel $f$ tel que $U = D_{+}(f)$ et poser
$$
\widetilde M(U) = M_{(f)}.
$$
Remarquons que, si $D_{+}(g) \subset D_{+}(f)$, le
Lemme \ref{constructions-lemma-standard-open} donne une application canonique
$$
\widetilde M(D_{+}(f)) = M_{(f)} \longrightarrow
M_{(g)} = \widetilde M(D_{+}(g)).
$$
On obtient ainsi un préfaisceau de
groupes abéliens sur la base des ouverts
standards. Si $M = S$, alors $\widetilde S$
est un préfaisceau d'anneaux sur
cette base. Pour $M$ quelconque, $\widetilde M$ est un préfaisceau de
$\widetilde S$-modules sur la même base.

\medskip\noindent
Calculons la fibre de $\widetilde M$ en un point
$x \in \text{Proj}(S)$. Supposons que $x$ corresponde à l'idéal premier
homogène $\mathfrak p \subset S$. Par définition de la fibre, on a
$$
\widetilde M_x
=
\colim_{f\in S_d, d > 0, f\not\in \mathfrak p} M_{(f)}
$$
L'ensemble $\{f \in S_d, d > 0, f \not \in \mathfrak p\}$ est ici préordonné
par la règle
$f \geq f' \Leftrightarrow D_{+}(f) \subset D_{+}(f')$.
Si $f_1, f_2 \in S \setminus \mathfrak p$ sont homogènes de degré strictement
positif, alors $f_1f_2 \geq f_1$ pour cet ordre. Dans Algèbre,
section \ref{algebra-section-proj}, nous avons défini
$M_{(\mathfrak p)}$ comme le module dont les éléments sont les fractions
$x/f$, avec $x, f$ homogènes,
$\deg(x) = \deg(f)$ et $f \not \in \mathfrak p$.
Comme $\mathfrak p \in \text{Proj}(S)$, il existe au
moins un $f_0 \in S$ homogène de degré strictement
positif tel que $f_0 \not\in \mathfrak p$. Puisque
$x/f = f_0x/ff_0$, on peut toujours supposer que le
dénominateur d'un élément de $M_{(\mathfrak p)}$ est de degré strictement
positif. Ces remarques donnent aussitôt
$$
\widetilde M_x = M_{(\mathfrak p)}.
$$

\medskip\noindent
Vérifions maintenant la condition de
faisceau pour les recouvrements ouverts
standards. Si
$D_{+}(f) = \bigcup_{i = 1}^n D_{+}(g_i)$, cette condition équivaut à
l'exactitude de la suite
$$
0 \to M_{(f)} \to \bigoplus M_{(g_i)} \to \bigoplus M_{(g_ig_j)}.
$$
Comme $D_{+}(g_i) = D_{+}(fg_i)$, on
peut récrire cette suite sous la forme
$$
0 \to M_{(f)} \to \bigoplus M_{(fg_i)} \to \bigoplus M_{(fg_ig_j)}.
$$
D'après le Lemme \ref{constructions-lemma-standard-open}, les éléments
$g_1^{\deg(f)}/f^{\deg(g_1)}, \ldots, g_n^{\deg(f)}/f^{\deg(g_n)}$
engendrent l'idéal unité de $S_{(f)}$, et les modules
$M_{(fg_i)}$ et $M_{(fg_ig_j)}$ sont les localisations principales du
$S_{(f)}$-module $M_{(f)}$ par rapport à ces éléments et à leurs produits.
On peut donc appliquer Algèbre, Lemme \ref{algebra-lemma-cover-module},
au module $M_{(f)}$ sur $S_{(f)}$ et aux éléments
$g_1^{\deg(f)}/f^{\deg(g_1)}, \ldots, g_n^{\deg(f)}/f^{\deg(g_n)}$.
La suite est ainsi exacte. D'après les remarques précédentes,
$\widetilde M$ est donc un faisceau sur la base des ouverts standards.

\medskip\noindent
Il résulte alors de Faisceaux, section \ref{sheaves-section-bases}, qu'il
existe un unique faisceau d'anneaux
$\mathcal{O}_{\text{Proj}(S)}$ qui coïncide
avec $\widetilde S$ sur les ouverts
standards. Le calcul des fibres effectué
ci-dessus et Algèbre, Lemme
\ref{algebra-lemma-proj-prime}, montrent que toutes
les fibres de ce faisceau d'anneaux sont des anneaux locaux.

\medskip\noindent
De même, pour tout $S$-module gradué $M$, il existe un
unique faisceau de $\mathcal{O}_{\text{Proj}(S)}$-modules
$\mathcal{F}$ qui coïncide avec
$\widetilde M$ sur les ouverts standards ; voir Faisceaux,
Lemme \ref{sheaves-lemma-extend-off-basis-module}.

\begin{definition}
\label{constructions-definition-structure-sheaf}
Soit $S$ un anneau gradué.
\begin{enumerate}
\item Le {\it faisceau structural $\mathcal{O}_{\text{Proj}(S)}$ du spectre
homogène de $S$} est l'unique faisceau d'anneaux
$\mathcal{O}_{\text{Proj}(S)}$ qui coïncide avec $\widetilde S$ sur la base
des ouverts standards.
\item L'espace localement annelé
$(\text{Proj}(S), \mathcal{O}_{\text{Proj}(S)})$ est appelé le
{\it spectre homogène} de $S$ et noté $\text{Proj}(S)$.
\item Le faisceau de $\mathcal{O}_{\text{Proj}(S)}$-modules qui prolonge
$\widetilde M$ à tous les ouverts de
$\text{Proj}(S)$ est appelé le faisceau
de $\mathcal{O}_{\text{Proj}(S)}$-modules associé à $M$. Lui aussi est noté
$\widetilde M$.
\end{enumerate}
\end{definition}

\noindent
Résumons les résultats obtenus jusqu'ici.

\begin{lemma}
\label{constructions-lemma-proj-sheaves}
Soit $S$ un anneau gradué, soit $M$ un $S$-module gradué et soit
$\widetilde M$ le faisceau de $\mathcal{O}_{\text{Proj}(S)}$-modules associé
à $M$.
\begin{enumerate}
\item Pour tout $f \in S$ homogène de degré strictement positif, on a
$$
\Gamma(D_{+}(f), \mathcal{O}_{\text{Proj}(S)}) = S_{(f)}.
$$
\item Pour tout $f\in S$ homogène de degré strictement positif, on a
$\Gamma(D_{+}(f), \widetilde M) = M_{(f)}$ en tant que $S_{(f)}$-module.
\item Si $D_{+}(g) \subset D_{+}(f)$, les applications de restriction de
$\mathcal{O}_{\text{Proj}(S)}$ et de $\widetilde M$ sont les applications
$S_{(f)} \to S_{(g)}$ et $M_{(f)} \to M_{(g)}$ du
Lemme \ref{constructions-lemma-standard-open}.
\item Soit $\mathfrak p$ un idéal premier
homogène de $S$ qui ne contient pas
$S_{+}$, et soit $x \in \text{Proj}(S)$ le point correspondant. On a
$\mathcal{O}_{\text{Proj}(S), x} = S_{(\mathfrak p)}$.
\item Soit $\mathfrak p$ un idéal premier
homogène de $S$ qui ne contient pas
$S_{+}$, et soit $x \in \text{Proj}(S)$ le point correspondant. On a
$(\widetilde M)_x = M_{(\mathfrak p)}$ en tant que
$S_{(\mathfrak p)}$-module.
\item
\label{constructions-item-map}
Il existe un homomorphisme canonique d'anneaux
$
S_0 \longrightarrow \Gamma(\text{Proj}(S), \widetilde S)
$
et un morphisme canonique de $S_0$-modules
$
M_0 \longrightarrow \Gamma(\text{Proj}(S), \widetilde M)
$
compatibles avec les descriptions ci-dessus des sections sur les ouverts
standards et des fibres.
\end{enumerate}
Toutes ces identifications sont en outre fonctorielles
en le $S$-module gradué $M$. En particulier,
le foncteur $M \mapsto \widetilde M$ est un
foncteur exact de la catégorie des $S$-modules gradués vers la catégorie des
$\mathcal{O}_{\text{Proj}(S)}$-modules.
\end{lemma}

\begin{proof}
Les assertions (1) à (5) résultent de la discussion précédente.
L'assertion (6) découle de l'existence des applications canoniques
$M_0 \to M_{(f)}$, $x \mapsto x/1$, compatibles avec les applications de
restriction décrites en (3). L'exactitude du foncteur
$M \mapsto \widetilde M$ résulte de celle du foncteur
$M \mapsto M_{(\mathfrak p)}$ (voir Algèbre,
Lemme \ref{algebra-lemma-proj-prime}) et du fait que l'exactitude d'une suite
courte de faisceaux se vérifie sur les fibres ; voir Modules,
Lemme \ref{modules-lemma-abelian}.
\end{proof}

\begin{remark}
\label{constructions-remark-global-sections-not-isomorphism}
L'application de $M_0$ dans les sections globales de $\widetilde M$ est en
général loin d'être un isomorphisme. Un exemple élémentaire consiste à prendre
$S = k[x, y, z]$ avec $1 = \deg(x) = \deg(y) = \deg(z)$ (ou un nombre
quelconque de variables) et $M = S/(x^{100}, y^{100}, z^{100})$. On voit facilement que
$\widetilde M = 0$, tandis que $M_0 = k$.
\end{remark}

\begin{lemma}
\label{constructions-lemma-standard-open-proj}
Soit $S$ un anneau gradué et soit $f \in S$ homogène de degré strictement
positif. Supposons que $D(g) \subset \Spec(S_{(f)})$ soit un ouvert standard.
Il existe alors un élément $h \in S$ homogène de degré strictement positif tel
que $D(g)$ corresponde à $D_{+}(h) \subset D_{+}(f)$ par l'homéomorphisme
d'Algèbre, Lemme \ref{algebra-lemma-topology-proj}. On peut même choisir $h$
de sorte que $g = h/f^n$ pour un certain $n$.
\end{lemma}

\begin{proof}
Écrivons $g = h/f^n$, où $h$ est homogène de degré strictement positif et
$n \geq 1$. Si $D_{+}(h)$ n'est pas contenu dans $D_{+}(f)$, remplaçons $h$
par $hf$ et $n$ par $n + 1$. L'élément $h$ a alors la forme voulue et
$D_{+}(h) \subset D_{+}(f)$ correspond à
$D(g) \subset \Spec(S_{(f)})$.
\end{proof}

\begin{lemma}
\label{constructions-lemma-proj-scheme}
Soit $S$ un anneau gradué. L'espace
localement annelé $\text{Proj}(S)$ est un
schéma. Les ouverts standards $D_{+}(f)$
sont des ouverts affines. Pour tout
$S$-module gradué $M$, le faisceau $\widetilde M$ est un faisceau
quasi-cohérent de $\mathcal{O}_{\text{Proj}(S)}$-modules.
\end{lemma}

\begin{proof}
Considérons un ouvert standard $D_{+}(f) \subset \text{Proj}(S)$.
D'après les Lemmes \ref{constructions-lemma-standard-open} et \ref{constructions-lemma-proj-sheaves},
on a $\Gamma(D_{+}(f), \mathcal{O}_{\text{Proj}(S)}) = S_{(f)}$, et il existe
un homéomorphisme $\varphi : D_{+}(f) \to \Spec(S_{(f)})$.
Pour tout ouvert standard $D(g) \subset \Spec(S_{(f)})$, choisissons
$h \in S_{+}$ comme dans le Lemme \ref{constructions-lemma-standard-open-proj}.
Alors $\varphi^{-1}(D(g)) = D_{+}(h)$ et, d'après les Lemmes
\ref{constructions-lemma-proj-sheaves} et \ref{constructions-lemma-standard-open},
$$
\Gamma(D_{+}(h), \mathcal{O}_{\text{Proj}(S)})
=
S_{(h)}
=
(S_{(f)})_{h^{\deg(f)}/f^{\deg(h)}}
=
(S_{(f)})_g
=
\Gamma(D(g), \mathcal{O}_{\Spec(S_{(f)})}).
$$
La restriction de $\mathcal{O}_{\text{Proj}(S)}$
à $D_{+}(f)$ correspond donc,
par l'homéomorphisme $\varphi$, exactement au faisceau
$\mathcal{O}_{\Spec(S_{(f)})}$ défini dans Schémas,
section \ref{schemes-section-affine-schemes}. Il s'ensuit que $D_{+}(f)$ est
un schéma affine isomorphe à $\Spec(S_{(f)})$ par $\varphi$, et donc que
$\text{Proj}(S)$ est un schéma.

\medskip\noindent
On montre exactement de la même manière que $\widetilde M$ est un faisceau
quasi-cohérent de $\mathcal{O}_{\text{Proj}(S)}$-modules. En effet, l'argument
précédent donne
$$
\widetilde M|_{D_{+}(f)} \cong \varphi^*\left(\widetilde{M_{(f)}}\right)
$$
ce qui montre que $\widetilde M$ est quasi-cohérent.
\end{proof}

\begin{lemma}
\label{constructions-lemma-proj-separated}
Soit $S$ un anneau gradué. Le
schéma $\text{Proj}(S)$ est séparé.
\end{lemma}

\begin{proof}
Il faut montrer que le morphisme canonique
$\text{Proj}(S) \to \Spec(\mathbf{Z})$ est séparé. Utilisons Schémas,
Lemme \ref{schemes-lemma-characterize-separated}. Il suffit de montrer que,
pour toute paire d'ouverts standards $D_{+}(f)$ et $D_{+}(g)$,
l'intersection $D_{+}(f) \cap D_{+}(g) = D_{+}(fg)$ est affine (ce qui est
clair) et que l'homomorphisme d'anneaux
$$
S_{(f)} \otimes_{\mathbf{Z}} S_{(g)} \longrightarrow S_{(fg)}
$$
est surjectif. Tout élément $s$ de $S_{(fg)}$ est de la forme
$s = h/(f^ng^m)$, où $h \in S$ est homogène de degré
$n\deg(f) + m\deg(g)$. En multipliant $h$ par un monôme convenable
$f^ig^j$, on peut supposer que $n = n' \deg(g)$ et
$m = m' \deg(f)$. On peut alors récrire $s$ sous la forme
$s = h/f^{(n' + m')\deg(g)} \cdot f^{m'\deg(g)}/g^{m'\deg(f)}$.
Ainsi, $s$ appartient bien à l'image de la
flèche affichée.
\end{proof}

\begin{lemma}
\label{constructions-lemma-proj-quasi-compact}
Soit $S$ un anneau gradué. Le schéma
$\text{Proj}(S)$ est quasi-compact si et
seulement s'il existe un nombre fini d'éléments homogènes
$f_1, \ldots, f_n \in S_{+}$ tels que
$S_{+} \subset \sqrt{(f_1, \ldots, f_n)}$. Dans ce cas,
$\text{Proj}(S) = D_+(f_1) \cup \ldots \cup D_+(f_n)$.
\end{lemma}

\begin{proof}
Si une telle famille d'éléments existe, les ouverts affines standards
$D_{+}(f_i)$ recouvrent $\text{Proj}(S)$ d'après Algèbre,
Lemme \ref{algebra-lemma-topology-proj}. Réciproquement, si
$\text{Proj}(S)$ est quasi-compact, on
peut le recouvrir par un nombre fini
d'ouverts standards $D_{+}(f_i)$, $i = 1, \ldots, n$, et le lemme que nous
venons de citer donne
$S_{+} \subset \sqrt{(f_1, \ldots, f_n)}$.
\end{proof}

\begin{lemma}
\label{constructions-lemma-structure-morphism-proj}
Soit $S$ un anneau gradué. Le schéma $\text{Proj}(S)$ possède un morphisme
canonique vers le schéma affine $\Spec(S_0)$, qui coïncide sur les espaces
topologiques avec l'application d'Algèbre,
Définition \ref{algebra-definition-proj}.
\end{lemma}

\begin{proof}
Nous avons vu plus haut que notre construction de $\widetilde S$, resp.\ $\widetilde M$,
donne un faisceau de $S_0$-algèbres, resp.\ $S_0$-modules.
On obtient donc un morphisme par Schémas,
Lemme \ref{schemes-lemma-morphism-into-affine}. Restreint à $D_{+}(f)$, ce
morphisme provient de l'homomorphisme canonique d'anneaux
$S_0 \to S_{(f)}$. Les applications $S \to S_f$ et $S_{(f)} \to S_f$ sont
des homomorphismes de $S_0$-algèbres ; voir le
Lemme \ref{constructions-lemma-standard-open}. Par conséquent, si l'idéal premier homogène
$\mathfrak p \subset S$ correspond à l'idéal premier $\mathbf{Z}$-gradué
$\mathfrak p' \subset S_f$ et à l'idéal premier ordinaire
$\mathfrak p'' \subset S_{(f)}$, ces trois idéaux ont la même image inverse
dans $S_0$.
\end{proof}

\begin{lemma}
\label{constructions-lemma-proj-valuative-criterion}
Soit $S$ un anneau gradué. Si $S$ est une algèbre de type fini sur $S_0$,
le morphisme $\text{Proj}(S) \to \Spec(S_0)$
satisfait les parties existence
et unicité du critère valuatif ; voir Schémas,
Définition \ref{schemes-definition-valuative-criterion}.
\end{lemma}

\begin{proof}
L'unicité résulte de ce que $\text{Proj}(S)$ est séparé
(Lemme \ref{constructions-lemma-proj-separated} et Schémas,
Lemme \ref{schemes-lemma-separated-implies-valuative}).
Choisissons des éléments homogènes $x_i \in S_{+}$, $i = 1, \ldots, n$,
qui engendrent $S$ sur $S_0$. Posons $d_i = \deg(x_i)$ et
$d = \text{lcm}\{d_i\}$. Supposons donné un diagramme
$$
\xymatrix{
\Spec(K) \ar[r] \ar[d] & \text{Proj}(S) \ar[d] \\
\Spec(A) \ar[r] & \Spec(S_0)
}
$$
comme dans Schémas, Définition \ref{schemes-definition-valuative-criterion}.
Notons $v : K^* \to \Gamma$ la valuation de $A$ ; voir Algèbre,
Définition \ref{algebra-definition-value-group}. On peut choisir un élément
homogène $f \in S_{+}$ tel que $\Spec(K)$ s'envoie dans $D_{+}(f)$. On obtient
alors un diagramme commutatif d'homomorphismes d'anneaux
$$
\xymatrix{
K & S_{(f)} \ar[l]^{\varphi} \\
A \ar[u] & S_0 \ar[l] \ar[u]
}
$$
Après renumérotation, on peut supposer que
$\varphi(x_i^{\deg(f)}/f^{d_i})$ est non nul pour $i = 1, \ldots, r$ et nul
pour $i = r + 1, \ldots, n$. Comme les ouverts $D_{+}(x_i)$ recouvrent
$\text{Proj}(S)$, on a $r \geq 1$. Choisissons
$i_0 \in \{1, \ldots, r\}$ de sorte que
$\gamma_i = (d/d_i)v(\varphi(x_i^{\deg(f)}/f^{d_i}))$ soit minimal dans
$\Gamma$. Pour alléger les notations, posons $x_0 = x_{i_0}$ et
$d_0 = d_{i_0}$. L'homomorphisme $\varphi$ se factorise par un homomorphisme
$\varphi' : S_{(fx_0)} \to K$, qui fournit un homomorphisme
$S_{(x_0)} \to S_{(fx_0)} \to K$.
L'algèbre $S_{(x_0)}$ est engendrée sur $S_0$ par les éléments
$x_1^{e_1} \ldots x_n^{e_n}/x_0^{e_0}$ tels que
$\sum e_i d_i = e_0 d_0$. Si $e_i > 0$ pour un certain $i > r$, alors
$\varphi'(x_1^{e_1} \ldots x_n^{e_n}/x_0^{e_0}) = 0$.
Si $e_i = 0$ pour $i > r$, on a
\begin{align*}
d \deg(f) v(\varphi'(x_1^{e_1} \ldots x_r^{e_r}/x_0^{e_0}))
& =
d v(\varphi'(x_1^{e_1 \deg(f)} \ldots x_r^{e_r \deg(f)}/x_0^{e_0 \deg(f)})) \\
& =
d \sum e_i v(\varphi'(x_i^{\deg(f)}/f^{d_i}))
- e_0 v(\varphi'(x_0^{\deg(f)}/f^{d_0})) \\
& =
\sum e_i d_i \gamma_i - e_0 d_0 \gamma_0 \\
& \geq
\sum e_i d_i \gamma_0 - e_0 d_0 \gamma_0 = 0
\end{align*}
car $\gamma_0$ est minimal parmi les $\gamma_i$. Il en résulte que
$S_{(x_0)}$ s'envoie dans $A$ par $\varphi'$. Le morphisme de schémas correspondant
$\Spec(A) \to \Spec(S_{(x_0)}) = D_{+}(x_0)
\subset \text{Proj}(S)$ est le morphisme qui complète le premier diagramme
commutatif de cette démonstration.
\end{proof}

\noindent
Nous avons aussi vu dans la démonstration du
Lemme \ref{constructions-lemma-proj-valuative-criterion} que, sous les hypothèses de ce
lemme, le morphisme $\text{Proj}(S) \to \Spec(S_0)$ est quasi-compact. Il résulte donc de Schémas,
Proposition \ref{schemes-proposition-characterize-universally-closed}, que
$\text{Proj}(S) \to \Spec(S_0)$ est universellement fermé dans cette
situation. Les exemples suivants montrent que ces résultats ne subsistent
pas sans hypothèse sur l'anneau gradué $S$.

\begin{example}
\label{constructions-example-not-existence-valuative-big-proj}
Soit $\mathbf{C}[X_1, X_2, X_3, \ldots]$ la $\mathbf{C}$-algèbre graduée dans
laquelle chaque $X_i$ est de degré $1$. Considérons l'homomorphisme d'anneaux
$$
\mathbf{C}[X_1, X_2, X_3, \ldots]
\longrightarrow
\mathbf{C}[t^\alpha ; \alpha \in \mathbf{Q}_{\geq 0}]
$$
qui envoie $X_i$ sur $t^{1/i}$. Le membre de droite devient un anneau de
valuation $A$ après localisation en l'idéal
$\mathfrak m = (t^\alpha ; \alpha > 0)$. Soit $K$ le corps des fractions de
$A$. L'homomorphisme ci-dessus donne un morphisme
$\Spec(K) \to \text{Proj}(\mathbf{C}[X_1, X_2, X_3, \ldots])$ qui ne se
prolonge pas en un morphisme défini sur tout $\Spec(A)$. En effet, l'image de
$\Spec(A)$ serait contenue dans l'un des $D_{+}(X_i)$ ; mais
$X_{i + 1}/X_i$ devrait alors s'envoyer sur un élément de $A$, ce qui est
impossible puisqu'il s'envoie sur $t^{1/(i + 1) - 1/i}$.
\end{example}

\begin{example}
\label{constructions-example-not-existence-valuative-small-proj}
Soit $R = \mathbf{C}[t]$ et
$$
S = R[X_1, X_2, X_3, \ldots]/(X_i^2 - tX_{i + 1}).
$$
La graduation est définie par $R = S_0$ et
$\deg(X_i) = 2^{i - 1}$. Remarquons que, si
$\mathfrak p \in \text{Proj}(S)$, alors $t \not \in \mathfrak p$ ; sinon,
$\mathfrak p$ contiendrait tous les $X_i$,
ce qui est impossible pour un point
du spectre homogène. On a donc
$D_{+}(X_i) = D_{+}(X_{i + 1})$ pour tout $i$. Ainsi,
$\text{Proj}(S)$ est quasi-compact ; il est même affine, puisqu'il est égal à
$D_{+}(X_1)$. On voit facilement que l'image de
$\text{Proj}(S) \to \Spec(R)$ est $D(t)$. Le morphisme
$\text{Proj}(S) \to \Spec(R)$ n'est donc pas fermé.
Le critère valuatif ne peut par conséquent pas s'appliquer, puisqu'il
entraînerait que le morphisme est fermé (voir Schémas,
Proposition \ref{schemes-proposition-characterize-universally-closed}).
\end{example}

\begin{example}
\label{constructions-example-trivial-proj}
Soit $A$ un anneau. Munissons $S = A[T]$ de sa structure de $A$-algèbre
graduée pour laquelle $T$ est de degré
$1$. Alors le morphisme canonique
$\text{Proj}(S) \to \Spec(A)$ (voir le
Lemme \ref{constructions-lemma-structure-morphism-proj}) est un isomorphisme.
\end{example}

\begin{example}
\label{constructions-example-open-subset-proj}
Soit $X = \Spec(A)$ un schéma affine et soit $U \subset X$ un sous-schéma
ouvert. Munissons $A[T]$ de la graduation définie par $\deg T = 1$.
Soit $S$ le sous-anneau de $A[T]$ engendré par $A$ et par tous les $fT^i$,
où $i \ge 0$ et où $f \in A$ vérifie $D(f) \subset U$. Nous affirmons que
$S$ est un anneau gradué tel que $S_0 = A$ et
$\text{Proj}(S) \cong U$, et que cet isomorphisme identifie le morphisme
canonique $\text{Proj}(S) \to \Spec(A)$ du
Lemme \ref{constructions-lemma-structure-morphism-proj} à l'inclusion $U \subset X$.

\medskip\noindent
Supposons que $\mathfrak p \in \text{Proj}(S)$ soit tel que tout élément
$fT \in S_1$ appartienne à $\mathfrak p$. Alors tout générateur $fT^i$ avec $i \ge 1$ appartient à $\mathfrak p$, car
$(fT^i)^2 = (fT)(fT^{2i-1}) \in \mathfrak p$ et $\mathfrak p$ est radical.
On aurait donc $\mathfrak p \supset S_+$, ce qui est impossible.
Par conséquent, $\text{Proj}(S)$ est recouvert par les sous-schémas
affines ouverts standards $\{D_+(fT)\}_{fT \in S_1}$.

\medskip\noindent
Remarquons que, si $fT \in S_1$, l'inclusion $S \subset A[T]$ induit un
isomorphisme gradué de $S[(fT)^{-1}]$ sur $A[T, T^{-1}, f^{-1}]$. L'ouvert standard
$D_+(fT) \cong \Spec(S_{(fT)})$ est donc isomorphe à
$\Spec(A[T, T^{-1}, f^{-1}]_0) = \Spec(A[f^{-1}])$.
Il est clair que cet isomorphisme est la restriction du morphisme canonique
$\text{Proj}(S) \to \Spec(A)$. Si, de plus, $gT \in S_1$, alors
$S[(fT)^{-1}, (gT)^{-1}] \cong A[T, T^{-1}, f^{-1}, g^{-1}]$
comme anneaux gradués, de sorte que $D_+(fT) \cap D_+(gT) \cong \Spec(A[f^{-1}, g^{-1}])$.
Ainsi, $\text{Proj}(S)$ est la réunion des sous-schémas ouverts $D_+(fT)$,
qui sont isomorphes par le morphisme canonique aux sous-schémas ouverts
$D(f) \subset X$, et leurs intersections dans $\text{Proj}(S)$ correspondent
à leurs intersections dans $X$. Le morphisme canonique est donc un
isomorphisme de $\text{Proj}(S)$ sur la réunion de tous les
$D(f) \subset U$, c'est-à-dire sur $U$.
\end{example}





















\section{Faisceaux quasi-cohérents sur Proj}
\label{constructions-section-quasi-coherent-proj}

\noindent
Soit $S$ un anneau gradué et soit $M$ un $S$-module gradué. Nous avons vu au
Lemme \ref{constructions-lemma-proj-sheaves} comment construire un faisceau quasi-cohérent de modules $\widetilde{M}$
sur $\text{Proj}(S)$ ainsi qu'une application
\begin{equation}
\label{constructions-equation-map-global-sections}
M_0 \longrightarrow \Gamma(\text{Proj}(S), \widetilde{M})
\end{equation}
de la partie de degré $0$ de $M$ vers les sections globales de
$\widetilde{M}$. La partie de degré $0$ du $n$-ième décalé $M(n)$ du module
gradué $M$ (voir Algèbre, section \ref{algebra-section-graded}) est égale à
$M_n$. On obtient donc des applications
\begin{equation}
\label{constructions-equation-map-global-sections-degree-n}
M_n \longrightarrow \Gamma(\text{Proj}(S), \widetilde{M(n)}).
\end{equation}
Nous voudrions pouvoir effectuer cette opération pour tout faisceau
quasi-cohérent $\mathcal{F}$ sur $\text{Proj}(S)$. Pour ce faire, nous
prendrons son produit tensoriel avec le $n$-ième tordu du faisceau structural ;
voir la Définition \ref{constructions-definition-twist}. Le lemme suivant relie les deux
notions.

\begin{lemma}
\label{constructions-lemma-widetilde-tensor}
Soit $S$ un anneau gradué. Soit
$(X, \mathcal{O}_X) = (\text{Proj}(S), \mathcal{O}_{\text{Proj}(S)})$ le
schéma du Lemme \ref{constructions-lemma-proj-scheme}. Soit $f \in S_{+}$ homogène.
Soit $x \in X$ un point correspondant à l'idéal premier homogène
$\mathfrak p \subset S$. Soient $M$ et $N$ des $S$-modules gradués.
Il existe un morphisme canonique de $\mathcal{O}_{\text{Proj}(S)}$-modules
$$
\widetilde M \otimes_{\mathcal{O}_X} \widetilde N
\longrightarrow
\widetilde{M \otimes_S N}
$$
qui induit l'application canonique
$
M_{(f)} \otimes_{S_{(f)}} N_{(f)}
\to
(M \otimes_S N)_{(f)}
$
sur les sections au-dessus de $D_{+}(f)$ et l'application canonique
$
M_{(\mathfrak p)} \otimes_{S_{(\mathfrak p)}} N_{(\mathfrak p)}
\to
(M \otimes_S N)_{(\mathfrak p)}
$
sur les fibres en $x$
En outre, le diagramme suivant est commutatif :
$$
\xymatrix{
M_0 \otimes_{S_0} N_0 \ar[r] \ar[d] &
(M \otimes_S N)_0 \ar[d] \\
\Gamma(X, \widetilde M \otimes_{\mathcal{O}_X} \widetilde N) \ar[r] &
\Gamma(X, \widetilde{M \otimes_S N})
}
$$
où les applications verticales sont données par
(\ref{constructions-equation-map-global-sections}).
\end{lemma}

\begin{proof}
Construire le morphisme affiché revient à construire une application
$\mathcal{O}_X$-bilinéaire
$$
\widetilde M \times \widetilde N
\longrightarrow
\widetilde{M \otimes_S N}
$$
voir Modules, section \ref{modules-section-tensor-product}. Il suffit de la
définir sur les sections au-dessus des ouverts $D_{+}(f)$, de manière
compatible avec les applications de restriction. Sur $D_{+}(f)$, on utilise
l'application $S_{(f)}$-bilinéaire
$M_{(f)} \times N_{(f)} \to (M \otimes_S N)_{(f)}$,
$(x/f^n, y/f^m) \mapsto (x \otimes y)/f^{n + m}$. Nous omettons les détails.
\end{proof}

\begin{remark}
\label{constructions-remark-not-isomorphism}
En général, le morphisme construit au Lemme \ref{constructions-lemma-widetilde-tensor}
n'est pas un isomorphisme. Voici un exemple. Soit $k$ un corps. Prenons
$S = k[x, y, z]$, où $k$ est placé en degré $0$ et
$\deg(x) = 1$, $\deg(y) = 2$, $\deg(z) = 3$.
Posons $M = S(1)$ et $N = S(2)$ ; pour les notations, voir Algèbre,
section \ref{algebra-section-graded}. Alors $M \otimes_S N = S(3)$.
Remarquons que
\begin{eqnarray*}
S_z
& = &
k[x, y, z, 1/z] \\
S_{(z)}
& = &
k[x^3/z, xy/z, y^3/z^2]
\cong
k[u, v, w]/(uw - v^3) \\
M_{(z)} & = & S_{(z)} \cdot x + S_{(z)} \cdot y^2/z \subset S_z \\
N_{(z)} & = & S_{(z)} \cdot y + S_{(z)} \cdot x^2 \subset S_z \\
S(3)_{(z)} & = & S_{(z)} \cdot z \subset S_z
\end{eqnarray*}
Considérons l'idéal maximal $\mathfrak m = (u, v, w) \subset S_{(z)}$.
On voit sans difficulté que $M_{(z)}/\mathfrak mM_{(z)}$ et
$N_{(z)}/\mathfrak mN_{(z)}$ sont tous deux de dimension $2$ sur
$\kappa(\mathfrak m)$. En revanche,
$S(3)_{(z)}/\mathfrak mS(3)_{(z)}$ est de dimension $1$.
L'application $M_{(z)} \otimes N_{(z)} \to S(3)_{(z)}$
n'est donc pas un
isomorphisme.
\end{remark}









\section{Faisceaux inversibles sur Proj}
\label{constructions-section-invertible-on-proj}

\noindent
Rappelons la construction, donnée dans Algèbre,
section \ref{algebra-section-graded}, du module tordu $M(n)$ associé à un
module gradué sur un anneau gradué.

\begin{definition}
\label{constructions-definition-twist}
Soit $S$ un anneau gradué et soit $X = \text{Proj}(S)$.
\begin{enumerate}
\item Nous posons $\mathcal{O}_X(n) = \widetilde{S(n)}$. Ce faisceau est
appelé le $n$-ième {\it tordu du faisceau structural de $\text{Proj}(S)$}.
\item Pour tout faisceau de
$\mathcal{O}_X$-modules $\mathcal{F}$, nous posons
$\mathcal{F}(n) = \mathcal{F} \otimes_{\mathcal{O}_X} \mathcal{O}_X(n)$.
\end{enumerate}
\end{definition}

\noindent
Nous allons utiliser le Lemme \ref{constructions-lemma-widetilde-tensor} pour construire
plusieurs morphismes canoniques. Comme
$S(n) \otimes_S S(m) = S(n + m)$, il existe des morphismes canoniques
\begin{equation}
\label{constructions-equation-multiply}
\mathcal{O}_X(n) \otimes_{\mathcal{O}_X} \mathcal{O}_X(m)
\longrightarrow
\mathcal{O}_X(n + m).
\end{equation}
Ces morphismes ne sont pas des isomorphismes
en général ; voir l'exemple de la
Remarque \ref{constructions-remark-not-isomorphism}. Le même exemple montre que
$\mathcal{O}_X(n)$ n'est en général {\it
pas} un faisceau inversible sur $X$.
En prenant le produit tensoriel avec un $\mathcal{O}_X$-module quelconque
$\mathcal{F}$, on obtient des morphismes
\begin{equation}
\label{constructions-equation-multiply-on-sheaf}
\mathcal{O}_X(n) \otimes_{\mathcal{O}_X} \mathcal{F}(m)
\longrightarrow
\mathcal{F}(n + m).
\end{equation}
Les applications (\ref{constructions-equation-map-global-sections-degree-n}) sur les
sections globales donnent un homomorphisme d'anneaux gradués
\begin{equation}
\label{constructions-equation-global-sections}
S \longrightarrow \bigoplus\nolimits_{n \geq 0} \Gamma(X, \mathcal{O}_X(n)).
\end{equation}
Pour un $\mathcal{O}_X$-module quelconque $\mathcal{F}$, les morphismes
(\ref{constructions-equation-multiply-on-sheaf}) donnent une structure de module gradué
\begin{equation}
\label{constructions-equation-global-sections-module}
\bigoplus\nolimits_{n \geq 0} \Gamma(X, \mathcal{O}_X(n))
\times
\bigoplus\nolimits_{m \in \mathbf{Z}} \Gamma(X, \mathcal{F}(m))
\longrightarrow
\bigoplus\nolimits_{m \in \mathbf{Z}} \Gamma(X, \mathcal{F}(m))
\end{equation}
et, par (\ref{constructions-equation-global-sections}), une structure de $S$-module.
Plus généralement, étant donné un $S$-module gradué $M$, on a
$M(n) = M \otimes_S S(n)$. On obtient donc des morphismes
\begin{equation}
\label{constructions-equation-multiply-more-generally}
\widetilde M(n)
=
\widetilde M
\otimes_{\mathcal{O}_X}
\mathcal{O}_X(n)
\longrightarrow
\widetilde{M(n)}.
\end{equation}
Sur les sections globales, (\ref{constructions-equation-map-global-sections-degree-n})
définit un morphisme de $S$-modules gradués
\begin{equation}
\label{constructions-equation-global-sections-more-generally}
M \longrightarrow
\bigoplus\nolimits_{n \in \mathbf{Z}} \Gamma(X, \widetilde{M(n)}).
\end{equation}
Voici un fait important qui résulte presque immédiatement des définitions.

\begin{lemma}
\label{constructions-lemma-when-invertible}
Soit $S$ un anneau gradué et posons $X = \text{Proj}(S)$. Soit
$f \in S$ homogène de degré $d > 0$. Les faisceaux
$\mathcal{O}_X(nd)|_{D_{+}(f)}$ sont inversibles, et même triviaux, pour tout
$n \in \mathbf{Z}$ (voir Modules,
Définition \ref{modules-definition-invertible}). Les morphismes
(\ref{constructions-equation-multiply}) restreints à $D_{+}(f)$
$$
\mathcal{O}_X(nd)|_{D_{+}(f)} \otimes_{\mathcal{O}_{D_{+}(f)}}
\mathcal{O}_X(m)|_{D_{+}(f)}
\longrightarrow
\mathcal{O}_X(nd + m)|_{D_{+}(f)},
$$
les morphismes (\ref{constructions-equation-multiply-on-sheaf}) restreints à $D_+(f)$
$$
\mathcal{O}_X(nd)|_{D_{+}(f)} \otimes_{\mathcal{O}_{D_{+}(f)}}
\mathcal{F}(m)|_{D_{+}(f)}
\longrightarrow
\mathcal{F}(nd + m)|_{D_{+}(f)},
$$
et les morphismes (\ref{constructions-equation-multiply-more-generally}) restreints à
$D_{+}(f)$
$$
\widetilde M(nd)|_{D_{+}(f)}
=
\widetilde M|_{D_{+}(f)}
\otimes_{\mathcal{O}_{D_{+}(f)}}
\mathcal{O}_X(nd)|_{D_{+}(f)}
\longrightarrow
\widetilde{M(nd)}|_{D_{+}(f)}
$$
sont des isomorphismes pour tous $n, m \in \mathbf{Z}$.
\end{lemma}

\begin{proof}
Les homomorphismes de $S$-modules gradués $S \to S(nd)$ et
$M \to M(nd)$ définis par $x \mapsto f^n x$ deviennent des isomorphismes
après inversion de $f$. Le premier donne $S_{(f)} \cong S(nd)_{(f)}$, d'où un isomorphisme
$\mathcal{O}_{D_{+}(f)} \cong \mathcal{O}_X(nd)|_{D_{+}(f)}$.
Le second montre que l'application
$S(nd)_{(f)} \otimes_{S_{(f)}} M_{(f)} \to M(nd)_{(f)}$ est un isomorphisme.
Le cas du morphisme (\ref{constructions-equation-multiply-on-sheaf}) résulte du cas du
morphisme (\ref{constructions-equation-multiply}).
\end{proof}

\begin{lemma}
\label{constructions-lemma-apply-modules}
Soit $S$ un anneau gradué, soit $M$ un $S$-module gradué et posons
$X = \text{Proj}(S)$. Supposons que $X$ soit recouvert par les ouverts
standards $D_+(f)$ avec $f \in S_1$, ce qui est par exemple le cas si $S$ est
engendré par $S_1$ sur $S_0$. Alors les faisceaux $\mathcal{O}_X(n)$ sont
inversibles et les morphismes (\ref{constructions-equation-multiply}),
(\ref{constructions-equation-multiply-on-sheaf}) et
(\ref{constructions-equation-multiply-more-generally}) sont des isomorphismes.
En particulier, ces morphismes induisent des isomorphismes
$$
\mathcal{O}_X(1)^{\otimes n} \cong
\mathcal{O}_X(n)
\quad
\text{et}
\quad
\widetilde{M} \otimes_{\mathcal{O}_X} \mathcal{O}_X(n) =
\widetilde{M}(n) \cong \widetilde{M(n)}
$$
Ainsi, (\ref{constructions-equation-map-global-sections-degree-n}) devient une application
\begin{equation}
\label{constructions-equation-map-global-sections-degree-n-simplified}
M_n \longrightarrow \Gamma(X, \widetilde{M}(n))
\end{equation}
et (\ref{constructions-equation-global-sections-more-generally}) devient une application
\begin{equation}
\label{constructions-equation-global-sections-more-generally-simplified}
M \longrightarrow
\bigoplus\nolimits_{n \in \mathbf{Z}} \Gamma(X, \widetilde{M}(n)).
\end{equation}
\end{lemma}

\begin{proof}
Sous les hypothèses du lemme, $X$ est recouvert par les ouverts
$D_{+}(f)$ avec $f \in S_1$ ; le résultat découle donc du
Lemme \ref{constructions-lemma-when-invertible}.
\end{proof}

\begin{lemma}
\label{constructions-lemma-where-invertible}
Soit $S$ un anneau gradué, posons $X = \text{Proj}(S)$ et fixons un entier
$d \geq 1$. Les deux ouverts suivants de $X$ sont égaux :
\begin{enumerate}
\item le plus grand ouvert $W = W_d \subset X$ tel que chaque
$\mathcal{O}_X(dn)|_W$ soit inversible et que tous les morphismes de
multiplication
$\mathcal{O}_X(nd)|_W \otimes_{\mathcal{O}_W} \mathcal{O}_X(md)|_W
\to \mathcal{O}_X(nd + md)|_W$
(voir \ref{constructions-equation-multiply}) soient des isomorphismes ;
\item la réunion des ouverts $D_{+}(fg)$, où $f, g \in S$ sont homogènes et
$\deg(f) = \deg(g) + d$.
\end{enumerate}
En outre, tous les morphismes
$\widetilde M(nd)|_W = \widetilde M|_W \otimes_{\mathcal{O}_W}
\mathcal{O}_X(nd)|_W \to \widetilde{M(nd)}|_W$
(voir \ref{constructions-equation-multiply-more-generally}) sont des isomorphismes.
\end{lemma}

\begin{proof}
Si $x \in D_{+}(fg)$ avec $\deg(f) = \deg(g) + d$, alors, sur
$D_{+}(fg)$, les faisceaux $\mathcal{O}_X(dn)$ sont engendrés par l'élément
$(f/g)^n = f^{2n}/(fg)^n$. En raisonnant comme dans la démonstration du
Lemme \ref{constructions-lemma-when-invertible}, on en déduit que $x$ appartient à l'ouvert
$W$ défini en (1).

\medskip\noindent
Réciproquement, supposons que $\mathcal{O}_X(dn)$ soit libre de rang 1 pour
tout $n$, sur un voisinage ouvert $V$ de
$x \in X$, et que tous les morphismes
de multiplication
$\mathcal{O}_X(nd)|_V \otimes_{\mathcal{O}_V} \mathcal{O}_X(md)|_V
\to \mathcal{O}_X(nd + md)|_V$ soient des isomorphismes.
On peut choisir $h \in S_{+}$ homogène tel que
$x \in D_{+}(h) \subset V$. Par définition
des tordus du faisceau structural,
il existe un élément $s$ de $(S_h)_d$ tel que $s^n$ soit une base de
$(S_h)_{nd}$ comme $S_{(h)}$-module pour tout $n \in \mathbf{Z}$.
Écrivons $s = f/h^m$ avec $m \geq 1$ et
$f \in S_{d + m \deg(h)}$. Posons $g = h^m$, de sorte que $s = f/g$.
Par construction, $x \in D_{+}(g)$. Remarquons que
$g^d \in (S_h)_{d\deg(g)}$. Par hypothèse, cet élément est un multiple de
$s^{\deg(g)} = f^{\deg(g)}/g^{\deg(g)}$ ; écrivons
$g^d = a/g^e \cdot f^{\deg(g)}/g^{\deg(g)}$. On en déduit
$g^{d + e + \deg(g)} = a f^{\deg(g)}$, puis $x \in D_{+}(f)$.
Ainsi, $x$ appartient à l'ensemble défini en (2).

\medskip\noindent
L'existence de la section génératrice $s = f/g$ sur l'ouvert affine
$D_{+}(fg)$ dont les puissances engendrent
librement les faisceaux de modules
$\mathcal{O}_X(nd)$ montre aisément que les morphismes
$\widetilde M(nd)|_W = \widetilde M|_W \otimes_{\mathcal{O}_W}
\mathcal{O}_X(nd)|_W \to \widetilde{M(nd)}|_W$
(voir \ref{constructions-equation-multiply-more-generally}) sont des isomorphismes.
Comparer avec la démonstration du Lemme \ref{constructions-lemma-when-invertible}.
\end{proof}

\noindent
Rappelons que, d'après Modules, Lemme \ref{modules-lemma-s-open}, si
$\mathcal{L}$ est un faisceau inversible
sur un espace localement annelé $X$
et si $s$ est une section globale de $\mathcal{L}$, alors l'ensemble
$X_s = \{x \in X \mid s \not \in \mathfrak m_x\mathcal{L}_x\}$ est ouvert.

\begin{lemma}
\label{constructions-lemma-principal-open}
Soit $S$ un anneau gradué, posons $X = \text{Proj}(S)$ et fixons un entier
$d \geq 1$. Soit $W = W_d \subset X$ le sous-schéma ouvert défini dans le
Lemme \ref{constructions-lemma-where-invertible}. Soient $n \geq 1$ et $f \in S_{nd}$.
Notons $s \in \Gamma(W, \mathcal{O}_W(nd))$ l'image de $f$ par
(\ref{constructions-equation-global-sections}), restreinte à $W$. Alors
$$
W_s = D_{+}(f) \cap W.
$$
\end{lemma}

\begin{proof}
Soit $D_{+}(ab) \subset W$ un ouvert affine standard, où $a, b \in S$ sont
homogènes et $\deg(a) = \deg(b) + d$. Remarquons que
$D_{+}(ab) \cap D_{+}(f) = D_{+}(abf)$. D'autre part, la restriction de $s$
à $D_{+}(ab)$ correspond à l'élément
$f/1 = b^nf/a^n (a/b)^n \in (S_{ab})_{nd}$.
Nous avons vu dans la démonstration du
Lemme \ref{constructions-lemma-where-invertible} que
$(a/b)^n$ engendre $\mathcal{O}_W(nd)$ au-dessus de $D_{+}(ab)$.
Il s'ensuit que $W_s \cap D_{+}(ab)$ est l'ouvert principal associé à
$b^nf/a^n \in \mathcal{O}_X(D_{+}(ab))$. Le résultat est alors clair.
\end{proof}

\noindent
Le lemme suivant énonce les propriétés que nous utiliserons plus loin pour
caractériser les schémas munis d'un faisceau inversible ample.

\begin{lemma}
\label{constructions-lemma-ample-on-proj}
Soit $S$ un anneau gradué, soit $X = \text{Proj}(S)$ et soit
$Y \subset X$ un sous-schéma ouvert quasi-compact. Notons
$\mathcal{O}_Y(n)$ la restriction de $\mathcal{O}_X(n)$ à $Y$.
Il existe un entier $d \geq 1$ tel que
\begin{enumerate}
\item le sous-schéma $Y$ est contenu dans l'ouvert $W_d$ défini au
Lemme \ref{constructions-lemma-where-invertible} ;
\item le faisceau $\mathcal{O}_Y(dn)$ est inversible pour tout
$n \in \mathbf{Z}$ ;
\item tous les morphismes
$\mathcal{O}_Y(nd) \otimes_{\mathcal{O}_Y} \mathcal{O}_Y(m)
\longrightarrow
\mathcal{O}_Y(nd + m)$
de l'Équation (\ref{constructions-equation-multiply}) sont des isomorphismes ;
\item tous les morphismes
$\widetilde M(nd)|_Y = \widetilde M|_Y \otimes_{\mathcal{O}_Y}
\mathcal{O}_X(nd)|_Y \to \widetilde{M(nd)}|_Y$
(voir \ref{constructions-equation-multiply-more-generally}) sont des isomorphismes ;
\item si $f \in S_{nd}$ et si
$s \in \Gamma(Y, \mathcal{O}_Y(nd))$ désigne
l'image de $f$ par (\ref{constructions-equation-global-sections}),
restreinte à $Y$, alors
$D_{+}(f) \cap Y = Y_s$ ;
\item une base de la topologie de $Y$ est formée par les ouverts $Y_s$,
où $s \in \Gamma(Y, \mathcal{O}_Y(nd))$ et $n \geq 1$ ;
\item une base de la topologie de $Y$ est formée par les ouverts affines
$Y_s \subset Y$, pour
$s \in \Gamma(Y, \mathcal{O}_Y(nd))$ et $n \geq 1$.
\end{enumerate}
\end{lemma}

\begin{proof}
Comme $Y$ est quasi-compact, il existe un nombre fini d'éléments homogènes
$f_i \in S_{+}$, $i = 1, \ldots, n$, tels que les ouverts standards
$D_{+}(f_i)$ recouvrent $Y$. Posons $d_i = \deg(f_i)$ et
$d = d_1 \ldots d_n$. Comme
$D_{+}(f_i) = D_{+}(f_i^{d/d_i})$, la caractérisation (2) du
Lemme \ref{constructions-lemma-where-invertible}, ou bien (1) et le
Lemme \ref{constructions-lemma-when-invertible}, donne immédiatement $Y \subset W_d$.
Le Lemme \ref{constructions-lemma-where-invertible} montre alors que (1) entraîne (2), (3)
et (4). Remarquons que (3) est un cas particulier de (4).
L'assertion (5) résulte du Lemme \ref{constructions-lemma-principal-open}.
Enfin, (6) et (7) résultent du fait que les ouverts $D_{+}(f)$ forment une
base de la topologie de $X$ et sont affines.
\end{proof}

\begin{lemma}
\label{constructions-lemma-comparison-proj-quasi-coherent}
Soit $S$ un anneau gradué, posons $X = \text{Proj}(S)$ et soit
$\mathcal{F}$ un $\mathcal{O}_X$-module quasi-cohérent. Munissons
$M = \bigoplus_{n \in \mathbf{Z}} \Gamma(X, \mathcal{F}(n))$ de sa structure
de $S$-module gradué à l'aide de
(\ref{constructions-equation-global-sections-module}) et (\ref{constructions-equation-global-sections}).
Il existe alors un morphisme canonique de $\mathcal{O}_X$-modules
$$
\widetilde{M} \longrightarrow \mathcal{F}
$$
fonctoriel en $\mathcal{F}$, tel que l'application induite
$M_0 \to \Gamma(X, \mathcal{F})$ soit l'identité.
\end{lemma}

\begin{proof}
Soit $f \in S$ homogène de degré $d > 0$. Rappelons que, d'après le
Lemme \ref{constructions-lemma-proj-sheaves}, la restriction
$\widetilde{M}|_{D_{+}(f)}$ correspond au $S_{(f)}$-module $M_{(f)}$.
On peut donc définir une application canonique
$$
M_{(f)} \longrightarrow \Gamma(D_+(f), \mathcal{F}),\quad
m/f^n \longmapsto m|_{D_+(f)} \otimes f|_{D_+(f)}^{-n}
$$
qui a un sens parce que $f|_{D_+(f)}$ est une section trivialisante du
faisceau inversible $\mathcal{O}_X(d)|_{D_+(f)}$ ; voir le
Lemme \ref{constructions-lemma-when-invertible} et sa démonstration. Comme $\widetilde{M}$
est quasi-cohérent, on obtient ainsi un morphisme canonique
$$
\widetilde{M}|_{D_+(f)} \longrightarrow \mathcal{F}|_{D_+(f)}
$$
par Schémas, Lemme \ref{schemes-lemma-compare-constructions}.
Ces morphismes se recollent sur les intersections et définissent donc un
morphisme global. Nous omettons la vérification du recollement ainsi que la
démonstration de la dernière assertion.
\end{proof}










\section{Fonctorialité de Proj}
\label{constructions-section-functoriality-proj}

\noindent
Un homomorphisme d'anneaux gradués $\psi : A \to B$ ne donne pas toujours lieu à
un morphisme entre les spectres homogènes projectifs associés. La raison en est que
l'image réciproque $\psi^{-1}(\mathfrak q)$
d'un idéal premier homogène $\mathfrak q \subset B$ peut
contenir l'idéal $A_{+}$ des éléments de degré strictement positif même si $\mathfrak q$ ne
contient pas $B_{+}$.
L'énoncé correct est le suivant.

\begin{lemma}
\label{constructions-lemma-morphism-proj}
Soient $A$, $B$ deux anneaux gradués.
Posons $X = \text{Proj}(A)$ et $Y = \text{Proj}(B)$.
Soit $\psi : A \to B$ un homomorphisme d'anneaux gradués.
Posons
$$
U(\psi)
=
\bigcup\nolimits_{f \in A_{+}\ \text{homogeneous}} D_{+}(\psi(f))
\subset Y.
$$
Il existe alors un morphisme canonique de schémas
$$
r_\psi :
U(\psi)
\longrightarrow
X
$$
et un homomorphisme d'algèbres $\mathbf{Z}$-graduées sur
$\mathcal{O}_{U(\psi)}$
$$
\theta = \theta_\psi :
r_\psi^*\left(
\bigoplus\nolimits_{d \in \mathbf{Z}} \mathcal{O}_X(d)
\right)
\longrightarrow
\bigoplus\nolimits_{d \in \mathbf{Z}} \mathcal{O}_{U(\psi)}(d).
$$
Le triplet $(U(\psi), r_\psi, \theta)$ est
caractérisé par les propriétés suivantes :
\begin{enumerate}
\item Pour tout $d \geq 0$, le diagramme
$$
\xymatrix{
A_d \ar[d] \ar[rr]_{\psi} & &
B_d \ar[d] \\
\Gamma(X, \mathcal{O}_X(d)) \ar[r]^-\theta &
\Gamma(U(\psi), \mathcal{O}_Y(d)) &
\Gamma(Y, \mathcal{O}_Y(d)) \ar[l]
}
$$ est commutatif.
\item Pour tout $f \in A_{+}$ homogène,
on a $r_\psi^{-1}(D_{+}(f)) = D_{+}(\psi(f))$ et
la restriction de $r_\psi$ à $D_{+}(\psi(f))$
correspond à l'homomorphisme d'anneaux
$A_{(f)} \to B_{(\psi(f))}$ induit par $\psi$.
\end{enumerate}
\end{lemma}

\begin{proof}
Il est clair que la condition (2) détermine de manière unique le morphisme de schémas
et l'ouvert $U(\psi)$. Fixons $f \in A_d$ avec $d \geq 1$.
Notons que
$\mathcal{O}_X(n)|_{D_{+}(f)}$ correspond au
$A_{(f)}$-module $(A_f)_n$ et que
$\mathcal{O}_Y(n)|_{D_{+}(\psi(f))}$ correspond au
$B_{(\psi(f))}$-module $(B_{\psi(f)})_n$. Autrement dit, la restriction de $\theta$
à $D_{+}(\psi(f))$ correspond à un homomorphisme d'algèbres
$\mathbf{Z}$-graduées sur $B_{(\psi(f))}$
$$
A_f \otimes_{A_{(f)}} B_{(\psi(f))}
\longrightarrow
B_{\psi(f)}
$$
La condition (1) détermine les images de tous les éléments de $A$.
Puisque $f$ est un élément inversible envoyé sur $\psi(f)$,
on voit que $1/f^m$ est envoyé sur $1/\psi(f)^m$. Il en résulte aussitôt
que $\theta$ est uniquement déterminé, à savoir qu'il est
donné par la formule
$$
a/f^m \otimes b/\psi(f)^e \longmapsto \psi(a)b/\psi(f)^{m + e}.
$$
Pour démontrer l'existence, remarquons que la preuve d'unicité ci-dessus fournit
une prescription bien définie pour le morphisme $r$ et l'homomorphisme $\theta$
sur tout ouvert standard de la forme
$D_{+}(\psi(f)) \subset U(\psi)$ à valeurs dans $D_{+}(f)$.
Notons-les $r_f$ et $\theta_f$.
Il suffit donc de vérifier que, si $D_{+}(f) \subset D_{+}(g)$
pour des éléments homogènes $f, g \in A_{+}$, alors la restriction de
$r_g$ à $D_{+}(\psi(f))$ coïncide avec $r_f$. Cela résulte clairement des
formules données ci-dessus pour $r$ et $\theta$.
\end{proof}

\begin{lemma}
\label{constructions-lemma-morphism-proj-transitive}
Soient $A$, $B$ et $C$ des anneaux gradués.
Posons $X = \text{Proj}(A)$, $Y = \text{Proj}(B)$ et $Z = \text{Proj}(C)$.
Soient $\varphi : A \to B$, $\psi : B \to C$ des homomorphismes d'anneaux gradués.
Alors on a
$$
U(\psi \circ \varphi) = r_\psi^{-1}(U(\varphi))
\quad
\text{et}
\quad
r_{\psi \circ \varphi}
=
r_\varphi \circ r_\psi|_{U(\psi \circ \varphi)}.
$$
De plus, on a
$$
\theta_\psi \circ r_\psi^*\theta_\varphi
=
\theta_{\psi \circ \varphi}
$$
avec les notations évidentes.
\end{lemma}

\begin{proof}
La démonstration est omise.
\end{proof}

\begin{lemma}
\label{constructions-lemma-surjective-graded-rings-map-proj}
Avec les hypothèses et les notations du lemme \ref{constructions-lemma-morphism-proj} ci-dessus.
Supposons que $A_d \to B_d$ soit surjectif pour tout $d \gg 0$. Alors
\begin{enumerate}
\item $U(\psi) = Y$,
\item $r_\psi : Y \to X$ est une immersion fermée, et
\item les homomorphismes $\theta : r_\psi^*\mathcal{O}_X(n) \to \mathcal{O}_Y(n)$
sont surjectifs, mais ne sont pas des isomorphismes en général (même si $A \to B$ est
surjectif).
\end{enumerate}
\end{lemma}

\begin{proof}
La partie (1) résulte de la définition de $U(\psi)$ et du fait que
$D_{+}(f) = D_{+}(f^n)$ pour tout $n > 0$. Pour $f \in A_{+}$ homogène,
on voit que $A_{(f)} \to B_{(\psi(f))}$ est surjectif car
tout élément de $B_{(\psi(f))}$ peut être représenté par une fraction
$b/\psi(f)^n$ avec $n$ arbitrairement grand (ce qui force le degré de
$b \in B$ à être grand). Cela démontre (2).
Le même argument montre que l'homomorphisme
$$
A_f \to B_{\psi(f)}
$$
est surjectif, ce qui démontre la surjectivité de $\theta$.
Pour un exemple où cet homomorphisme n'est pas un isomorphisme,
considérons l'anneau gradué $A = k[x, y]$, où $k$ est un corps,
avec $\deg(x) = 1$, $\deg(y) = 2$. Posons $I = (x)$, de sorte que
$B = k[y]$. Notons que $\mathcal{O}_Y(1) = 0$ dans ce cas.
Mais on voit facilement que $r_\psi^*\mathcal{O}_X(1)$
n'est pas nul. (Il existe des exemples moins artificiels.)
\end{proof}

\begin{lemma}
\label{constructions-lemma-eventual-iso-graded-rings-map-proj}
Avec les hypothèses et les notations du lemme \ref{constructions-lemma-morphism-proj} ci-dessus.
Supposons que $A_d \to B_d$ soit un isomorphisme pour tout $d \gg 0$. Alors
\begin{enumerate}
\item $U(\psi) = Y$,
\item $r_\psi : Y \to X$ est un isomorphisme, et
\item les homomorphismes $\theta : r_\psi^*\mathcal{O}_X(n) \to \mathcal{O}_Y(n)$
sont des isomorphismes.
\end{enumerate}
\end{lemma}

\begin{proof}
On a (1) par le lemme \ref{constructions-lemma-surjective-graded-rings-map-proj}.
Soit $f \in A_{+}$ homogène. L'hypothèse sur $\psi$ implique que
$A_f \to B_f$ est un isomorphisme (détails omis). Il est donc clair que
les restrictions de $r_\psi$ et de $\theta$ au-dessus de $D_{+}(f)$ sont des isomorphismes.
Le lemme en résulte.
\end{proof}

\begin{lemma}
\label{constructions-lemma-surjective-graded-rings-generated-degree-1-map-proj}
Avec les hypothèses et les notations du lemme \ref{constructions-lemma-morphism-proj} ci-dessus.
Supposons que $A_d \to B_d$ soit surjectif pour $d \gg 0$ et que $A$ soit engendré
par $A_1$ comme $A_0$-algèbre. Alors
\begin{enumerate}
\item $U(\psi) = Y$,
\item $r_\psi : Y \to X$ est une immersion fermée, et
\item les homomorphismes $\theta : r_\psi^*\mathcal{O}_X(n) \to \mathcal{O}_Y(n)$
sont des isomorphismes.
\end{enumerate}
\end{lemma}

\begin{proof}
Par les lemmes \ref{constructions-lemma-eventual-iso-graded-rings-map-proj} et
\ref{constructions-lemma-morphism-proj-transitive},
on peut remplacer $B$ par l'image de $A \to B$
sans modifier $X$ ni les faisceaux $\mathcal{O}_X(n)$.
On peut donc supposer que $A \to B$ est surjectif. Par le
lemme \ref{constructions-lemma-surjective-graded-rings-map-proj}, on obtient (1) et (2),
ainsi que la surjectivité dans (3).
Par le lemme \ref{constructions-lemma-apply-modules}, on voit que
$\mathcal{O}_X(n)$ et $\mathcal{O}_Y(n)$
sont tous deux inversibles. Ainsi, $\theta$ est un isomorphisme.
\end{proof}

\begin{lemma}
\label{constructions-lemma-base-change-map-proj}
\begin{reference}
\cite[II, Proposition 2.8.10]{EGA}
\end{reference}
Avec les hypothèses et les notations du lemme \ref{constructions-lemma-morphism-proj} ci-dessus.
Supposons qu'il existe un homomorphisme d'anneaux $R \to A_0$ et un homomorphisme d'anneaux
$R \to R'$ tels que $B = R' \otimes_R A$. Alors
\begin{enumerate}
\item $U(\psi) = Y$,
\item le diagramme
$$
\xymatrix{
Y = \text{Proj}(B) \ar[r]_{r_\psi} \ar[d] &
\text{Proj}(A) = X \ar[d] \\
\Spec(R') \ar[r] &
\Spec(R)
}
$$
est un carré cartésien, et
\item les homomorphismes $\theta : r_\psi^*\mathcal{O}_X(n) \to \mathcal{O}_Y(n)$
sont des isomorphismes.
\end{enumerate}
\end{lemma}

\begin{proof}
Cela résulte immédiatement de l'examen de ce qui se passe sur les ouverts
standards $D_{+}(f)$ pour $f \in A_{+}$.
\end{proof}

\begin{lemma}
\label{constructions-lemma-localization-map-proj}
Avec les hypothèses et les notations du lemme \ref{constructions-lemma-morphism-proj} ci-dessus.
Supposons qu'il existe $g \in A_0$ tel que $\psi$ induise un
isomorphisme $A_g \to B$. Alors
$U(\psi) = Y$, $r_\psi : Y \to X$ est une immersion ouverte
qui identifie $Y$ à l'image réciproque
de $D(g) \subset \Spec(A_0)$. De plus, l'homomorphisme $\theta$
est un isomorphisme.
\end{lemma}

\begin{proof}
C'est un cas particulier du lemme \ref{constructions-lemma-base-change-map-proj} ci-dessus.
\end{proof}

\begin{lemma}
\label{constructions-lemma-d-uple}
Soit $S$ un anneau gradué. Soit $d \geq 1$. Posons $S' = S^{(d)}$ avec les notations
d'Algèbre, section \ref{algebra-section-graded}. Posons
$X = \text{Proj}(S)$ et $X' = \text{Proj}(S')$. Il existe un
isomorphisme canonique de schémas $i : X \to X'$ tel que
\begin{enumerate}
\item pour tout $S$-module gradué $M$, en posant $M' = M^{(d)}$,
on ait un isomorphisme canonique $\widetilde{M} \to i^*\widetilde{M'}$,
\item on ait des isomorphismes canoniques
$\mathcal{O}_{X}(nd) \to i^*\mathcal{O}_{X'}(n)$
\end{enumerate}
et ces isomorphismes sont compatibles avec les homomorphismes de multiplication
du lemme \ref{constructions-lemma-widetilde-tensor}, et donc avec les homomorphismes
(\ref{constructions-equation-multiply}),
(\ref{constructions-equation-multiply-on-sheaf}),
(\ref{constructions-equation-global-sections}),
(\ref{constructions-equation-global-sections-module}),
(\ref{constructions-equation-multiply-more-generally}) et
(\ref{constructions-equation-global-sections-more-generally}) (voir la démonstration pour les
énoncés précis).
\end{lemma}

\begin{proof}
L'homomorphisme injectif d'anneaux $S' \to S$ (qui n'est pas un homomorphisme d'anneaux gradués
en raison de nos conventions) induit un morphisme $j : \Spec(S) \to \Spec(S')$.
Étant donné un idéal premier homogène $\mathfrak p \subset S$, on voit que
$\mathfrak p' = j(\mathfrak p) = S' \cap \mathfrak p$
est un idéal premier homogène de $S'$.
De plus, si $f \in S_+$ est homogène et $f \not \in \mathfrak p$, alors
$f^d \in S'_+$ et $f^d \not \in \mathfrak p'$. Réciproquement, si
$\mathfrak p' \subset S'$ est un idéal premier homogène ne contenant pas un
élément homogène $f \in S'_+$, alors
$\mathfrak p = \{g \in S \mid g^d \in \mathfrak p'\}$ est un
idéal premier homogène de $S$ ne contenant pas $f$, dont l'image par $j$
est $\mathfrak p'$. Pour voir que $\mathfrak p$ est un idéal, notons
que, si $g, h \in \mathfrak p$, alors
$(g + h)^{2d} \in \mathfrak p'$ par la formule du binôme.
On a donc $g + h \in \mathfrak p'$ puisque $\mathfrak p'$ est premier.
On voit ainsi que $j$ induit un homéomorphisme $i : X \to X'$.
De plus, étant donné $f \in S_+$ homogène, on a
$S_{(f)} \cong S'_{(f^d)}$. Comme ces isomorphismes sont compatibles
avec les homomorphismes de restriction du
lemme \ref{constructions-lemma-standard-open}, on voit qu'il existe un
isomorphisme $i^\sharp : i^{-1}\mathcal{O}_{X'} \to \mathcal{O}_X$ de
faisceaux structuraux sur $X$ et $X'$ ; ainsi, $i$ est un isomorphisme
de schémas.

\medskip\noindent
Soit $M$ un $S$-module gradué. Étant donné $f \in S_+$ homogène, on a
$M_{(f)} \cong M'_{(f^d)}$ ; exactement de la même manière que ci-dessus,
on obtient donc l'isomorphisme de (1). Les isomorphismes de (2) sont un cas particulier
de (1) pour $M = S(nd)$, ce qui donne $M' = S'(n)$. Soient $M$ et $N$
des $S$-modules gradués. Alors on a
$$
M' \otimes_{S'} N' =
(M \otimes_S N)^{(d)} =
(M \otimes_S N)'
$$
comme on le vérifie directement à partir des définitions. Cela étant,
la compatibilité avec les homomorphismes de multiplication du
lemme \ref{constructions-lemma-widetilde-tensor} équivaut à la commutativité du diagramme
$$
\xymatrix{
\widetilde M \otimes_{\mathcal{O}_X} \widetilde N
\ar[d]_{(1) \otimes (1)} \ar[r] &
\widetilde{M \otimes_S N} \ar[d]^{(1)} \\
i^*(\widetilde{M'} \otimes_{\mathcal{O}_{X'}} \widetilde{N'}) \ar[r] &
i^*(\widetilde{M' \otimes_{S'} N'})
}
$$
Cela se voit en examinant la construction des homomorphismes
sur l'ouvert $D_+(f) = D_+(f^d)$, où la flèche horizontale supérieure
est donnée par l'application
$M_{(f)} \times N_{(f)} \to (M \otimes_S N)_{(f)}$
et la flèche horizontale inférieure par l'application
$M'_{(f^d)} \times N'_{(f^d)} \to (M' \otimes_{S'} N')_{(f^d)}$.
Puisque ces applications coïncident via les identifications
$M_{(f)} = M'_{(f^d)}$, etc., on obtient la compatibilité voulue.
Nous omettons la démonstration des autres compatibilités.
\end{proof}











\section{Morphismes vers Proj}
\label{constructions-section-morphisms-proj}

\noindent
Soit $S$ un anneau gradué. Soit $X = \text{Proj}(S)$ le spectre homogène de
$S$ et soit $d \geq 1$ un entier. Considérons le sous-schéma ouvert
\begin{equation}
\label{constructions-equation-Ud}
U_d = \bigcup\nolimits_{f  \in S_d} D_{+}(f)
\quad\subset\quad
X = \text{Proj}(S)
\end{equation}
Remarquons que $d | d' \Rightarrow U_d \subset U_{d'}$ et que
$X = \bigcup_d U_d$. Ni $X$ ni $U_d$ ne
sont nécessairement quasi-compacts ;
voir Algèbre, Lemme \ref{algebra-lemma-topology-proj}. Écrivons
$\mathcal{O}_{U_d}(n) = \mathcal{O}_X(n)|_{U_d}$.
D'après le Lemme \ref{constructions-lemma-when-invertible},
$\mathcal{O}_{U_d}(nd)$, pour $n \in \mathbf{Z}$, est un module inversible
sur $\mathcal{O}_{U_d}$, et tous
les morphismes de multiplication
$\mathcal{O}_{U_d}(nd) \otimes_{\mathcal{O}_{U_d}} \mathcal{O}_{U_d}(m)
\to \mathcal{O}_{U_d}(nd + m)$ de (\ref{constructions-equation-multiply}) sont des
isomorphismes. En particulier,
$\mathcal{O}_{U_d}(nd) \cong \mathcal{O}_{U_d}(d)^{\otimes n}$.
L'homomorphisme d'anneaux gradués (\ref{constructions-equation-global-sections}), composé
avec la restriction à $U_d$, donne un homomorphisme d'anneaux gradués
\begin{equation}
\label{constructions-equation-psi-d}
\psi^d : S^{(d)} \longrightarrow \Gamma_*(U_d, \mathcal{O}_{U_d}(d)).
\end{equation}
Pour la notation $S^{(d)}$, voir Algèbre,
section \ref{algebra-section-graded}. Pour la notation $\Gamma_*$, voir
Modules, Définition \ref{modules-definition-gamma-star}. En
outre, puisque $U_d$ est recouvert par les ouverts $D_{+}(f)$,
$f \in S_d$, le faisceau
$\mathcal{O}_{U_d}(d)$ est engendré par
les sections appartenant à l'image de
$\psi^d_1 : S^{(d)}_1 = S_d \to \Gamma(U_d, \mathcal{O}_{U_d}(d))$ ; voir
Modules, Définition \ref{modules-definition-globally-generated}.

\medskip\noindent
Soit $Y$ un schéma et soit $\varphi : Y \to X$ un morphisme de schémas.
Supposons que l'image $\varphi(Y)$ soit
contenue dans le sous-schéma ouvert
$U_d$ de $X$. La discussion qui suit Modules,
Définition \ref{modules-definition-gamma-star},
fournit un homomorphisme
d'anneaux gradués
$$
\Gamma_*(U_d, \mathcal{O}_{U_d}(d))
\longrightarrow
\Gamma_*(Y, \varphi^*\mathcal{O}_X(d)).
$$
Sa composée avec $\psi^d$ donne un homomorphisme d'anneaux gradués
\begin{equation}
\label{constructions-equation-psi-phi-d}
\psi_\varphi^d :
S^{(d)}
\longrightarrow
\Gamma_*(Y, \varphi^*\mathcal{O}_X(d))
\end{equation}
tel que le faisceau inversible
$\varphi^*\mathcal{O}_X(d)$ soit engendré par
les sections appartenant à l'image de
$(S^{(d)})_1 = S_d \to \Gamma(Y, \varphi^*\mathcal{O}_X(d))$.

\begin{lemma}
\label{constructions-lemma-converse-construction}
Soient $S$ un anneau gradué, $X = \text{Proj}(S)$, $d \geq 1$ et
$U_d \subset X$ comme ci-dessus. Soit
$Y$ un schéma, soit $\mathcal{L}$ un
faisceau inversible sur $Y$ et soit
$\psi : S^{(d)} \to \Gamma_*(Y, \mathcal{L})$ un homomorphisme d'anneaux
gradués tel que $\mathcal{L}$ soit
engendré par les sections appartenant à
l'image de $\psi|_{S_d} : S_d \to \Gamma(Y, \mathcal{L})$.
Il existe alors un morphisme $\varphi : Y \to X$ tel que
$\varphi(Y) \subset U_d$ et un isomorphisme
$\alpha : \varphi^*\mathcal{O}_{U_d}(d) \to \mathcal{L}$ tels que
$\psi_\varphi^d$ coïncide avec $\psi$ par l'intermédiaire de $\alpha$ :
$$
\xymatrix{
\Gamma_*(Y, \mathcal{L}) &
\Gamma_*(Y, \varphi^*\mathcal{O}_{U_d}(d)) \ar[l]^-\alpha &
\Gamma_*(U_d, \mathcal{O}_{U_d}(d)) \ar[l]^-{\varphi^*} \\
S^{(d)} \ar[u]^\psi & &
S^{(d)} \ar[u]^{\psi^d} \ar[ul]^{\psi^d_\varphi} \ar[ll]_{\text{id}}
}
$$
est commutatif. De plus, le couple $(\varphi, \alpha)$ est unique.
\end{lemma}

\begin{proof}
Choisissons $f \in S_d$ et notons
$s = \psi(f) \in \Gamma(Y, \mathcal{L})$. Sur l'ouvert $Y_s$ où $s$ ne
s'annule pas, la multiplication par $s$ induit un isomorphisme
$\mathcal{O}_{Y_s} \to \mathcal{L}|_{Y_s}$ ; voir Modules,
Lemme \ref{modules-lemma-s-open}. Nous noterons l'inverse de cet isomorphisme
$x \mapsto x/s$, et procéderons de même pour les puissances de
$\mathcal{L}$. On définit ainsi un homomorphisme d'anneaux
$\psi_{(f)} : S_{(f)} \to \Gamma(Y_s, \mathcal{O}_Y)$ en envoyant la fraction
$a/f^n$ sur $\psi(a)/s^n$. Par Schémas,
Lemme \ref{schemes-lemma-morphism-into-affine},
cet homomorphisme correspond à
un morphisme $\varphi_f : Y_s \to \Spec(S_{(f)}) = D_{+}(f)$.
Introduisons aussi l'isomorphisme
$\alpha_f : \varphi_f^*\mathcal{O}_{D_{+}(f)}(d) \to \mathcal{L}|_{Y_s}$ qui
envoie l'image inverse de la section trivialisante $f$ sur $D_{+}(f)$ vers la
section trivialisante $s$ sur $Y_s$. Avec ce choix, le diagramme du lemme est
commutatif avec $Y$ remplacé par $Y_s$, $\varphi$ remplacé par $\varphi_f$
et $\alpha$ remplacé par $\alpha_f$ ; nous omettons la vérification.

\medskip\noindent
Soit $f' \in S_d$ un second élément et notons
$s' = \psi(f') \in \Gamma(Y, \mathcal{L})$. Alors
$Y_s \cap Y_{s'} = Y_{ss'}$, et de même
$D_{+}(f) \cap D_{+}(f') = D_{+}(ff')$. Nous avons vu au
Lemme \ref{constructions-lemma-ample-on-proj} que $D_{+}(f') \cap D_{+}(f)$ est l'ensemble
des points de $D_{+}(f)$ où la section de $\mathcal{O}_X(d)$ définie par $f'$
ne s'annule pas. Ainsi,
$\varphi_f^{-1}(D_{+}(f') \cap D_{+}(f)) = Y_s \cap Y_{s'}
= \varphi_{f'}^{-1}(D_{+}(f') \cap D_{+}(f))$.
Sur $D_{+}(f) \cap D_{+}(f')$, la fraction $f/f'$ est une section inversible
du faisceau structural, d'inverse $f'/f$. Remarquons que
$\psi_{(f')}(f/f') = \psi(f)/s' = s/s'$ et
$\psi_{(f)}(f'/f) = \psi(f')/s = s'/s$. Nous affirmons qu'il existe un unique
homomorphisme d'anneaux
$S_{(ff')} \to \Gamma(Y_{ss'}, \mathcal{O}_Y)$
rendant commutatif le diagramme
$$
\xymatrix{
\Gamma(Y_s, \mathcal{O}_Y) \ar[r] &
\Gamma(Y_{ss'}, \mathcal{O}_Y) &
\Gamma(Y_{s, '} \mathcal{O}_Y) \ar[l]\\
S_{(f)} \ar[r] \ar[u]^{\psi_{(f)}} &
S_{(ff')} \ar[u] &
S_{(f')} \ar[l] \ar[u]^{\psi_{(f')}}
}
$$
Il existe, car on peut utiliser la règle
$x/(ff')^n \mapsto \psi(x)/(ss')^n$, qui
est bien définie d'après les formules
précédentes. L'unicité résulte de ce que
$\text{Proj}(S)$ est séparé ; voir le
Lemme \ref{constructions-lemma-proj-separated} et sa démonstration. Ainsi, les morphismes
$\varphi_f$ et $\varphi_{f'}$ coïncident sur $Y_s \cap Y_{s'}$.
Les restrictions de $\alpha_f$ et $\alpha_{f'}$ coïncident aussi sur
$Y_s \cap Y_{s'}$, puisque
les fonctions régulières $s/s'$ et $\psi_{(f')}(f/f')$ coïncident.
Les morphismes $\psi_f$ se recollent donc en un morphisme global de $Y$
vers $U_d \subset X$, et les morphismes $\alpha_f$ se recollent en un
isomorphisme satisfaisant les conditions du lemme.

\medskip\noindent
Il reste à montrer l'unicité du couple $(\varphi, \alpha)$. Supposons que
$(\varphi', \alpha')$ soit un second couple satisfaisant les mêmes conditions.
Soit $f \in S_d$. La commutativité des diagrammes du lemme montre que les
images inverses de $D_{+}(f)$ par $\varphi$ et par $\varphi'$ sont toutes deux
égales à $Y_{\psi(f)}$. Comme les ouverts $D_{+}(f)$ forment une base de la
topologie de $X$ et que l'espace topologique $X$ est sobre (voir Schémas,
Lemme \ref{schemes-lemma-scheme-sober}),
les applications continues sous-jacentes
à $\varphi$ et à $\varphi'$ sont égales. Utilisons $s = \psi(f)$ pour
trivialiser le faisceau inversible $\mathcal{L}$ sur $Y_{\psi(f)}$.
La commutativité des diagrammes donne
$\alpha^{\otimes n}(\psi^d_{\varphi}(x)) =
\psi(x) = (\alpha')^{\otimes n}(\psi^d_{\varphi'}(x))$
pour tout $x \in S_{nd}$. Par construction de $\psi^d_{\varphi}$ et de
$\psi^d_{\varphi'}$, on a
$\psi^d_{\varphi}(x) = \varphi^\sharp(x/f^n) \psi^d_{\varphi}(f^n)$
sur $Y_{\psi(f)}$, et de même pour $\psi^d_{\varphi'}$.
La commutativité des diagrammes du lemme entraîne alors
$\varphi^\sharp(x/f^n) = (\varphi')^\sharp(x/f^n)$.
Ainsi, $\varphi$ et $\varphi'$ induisent le même morphisme de
$Y_{\psi(f)}$ vers le schéma affine
$D_{+}(f) = \Spec(S_{(f)})$ . Ainsi, $\varphi$ et $\varphi'$
sont égaux comme morphismes. Enfin, il reste
à montrer que, $\varphi$ étant donné, la
commutativité du diagramme du lemme détermine un unique $\alpha$. Nous omettons
la vérification.
\end{proof}

\noindent
Reprenons la discussion qui précède le lemme. Soit $S$ un anneau gradué et
soit $Y$ un schéma. Nous considérerons les {\it triplets}
$(d, \mathcal{L}, \psi)$ où
\begin{enumerate}
\item $d \geq 1$ est un entier ;
\item $\mathcal{L}$ est un $\mathcal{O}_Y$-module inversible ;
\item $\psi : S^{(d)} \to \Gamma_*(Y, \mathcal{L})$ est un homomorphisme
d'anneaux gradués tel que $\mathcal{L}$ soit engendré par les sections
globales $\psi(f)$, avec $f \in S_d$.
\end{enumerate}
Étant donnés un morphisme $h : Y' \to Y$ et un triplet
$(d, \mathcal{L}, \psi)$ sur $Y$, on
peut en prendre l'image inverse : c'est
le triplet $(d, h^*\mathcal{L}, h^* \circ \psi)$.
Deux triplets $(d, \mathcal{L}, \psi)$ et
$(d, \mathcal{L}', \psi')$ associés au même entier $d$ sont dits
{\it strictement équivalents} s'il existe un isomorphisme
$\beta : \mathcal{L} \to \mathcal{L}'$ tel que
$\beta \circ \psi = \psi'$ comme homomorphismes d'anneaux gradués
$S^{(d)} \to \Gamma_*(Y, \mathcal{L}')$.

\medskip\noindent
Pour tout entier $d \geq 1$, nous définissons
\begin{eqnarray*}
F_d : \Sch^{opp} & \longrightarrow & \textit{Ensembles}, \\
Y & \longmapsto &
\{\text{strict equivalence classes of triples }
(d, \mathcal{L}, \psi)
\text{ as above}\}
\end{eqnarray*}
les images inverses étant définies comme ci-dessus.

\begin{lemma}
\label{constructions-lemma-proj-functor-strict}
Soit $S$ un anneau gradué et soit
$X = \text{Proj}(S)$. Le sous-schéma ouvert
$U_d \subset X$ défini en (\ref{constructions-equation-Ud}) représente le foncteur $F_d$,
et le triplet $(d, \mathcal{O}_{U_d}(d), \psi^d)$ défini ci-dessus est la
famille universelle (voir Schémas,
section \ref{schemes-section-representable}).
\end{lemma}

\begin{proof}
C'est une reformulation du Lemme \ref{constructions-lemma-converse-construction}.
\end{proof}

\begin{lemma}
\label{constructions-lemma-apply}
Soit $S$ un anneau gradué engendré, comme $S_0$-algèbre, par les éléments de
$S_1$. Alors le schéma $X = \text{Proj}(S)$
représente le foncteur qui associe
à un schéma $Y$ l'ensemble des couples $(\mathcal{L}, \psi)$ où
\begin{enumerate}
\item $\mathcal{L}$ est un $\mathcal{O}_Y$-module inversible ;
\item $\psi : S \to \Gamma_*(Y, \mathcal{L})$ est un homomorphisme d'anneaux
gradués tel que $\mathcal{L}$ soit engendré par les sections globales
$\psi(f)$, avec $f \in S_1$,
\end{enumerate}
à équivalence stricte près, au sens défini ci-dessus.
\end{lemma}

\begin{proof}
Sous les hypothèses du lemme, on a $X = U_1$ ; le résultat est donc une
reformulation du Lemme \ref{constructions-lemma-proj-functor-strict}.
\end{proof}

\noindent
Terminons cette section par la description d'un foncteur correspondant à
$\text{Proj}(S)$ pour un anneau gradué quelconque $S$. Nous conseillons au
lecteur de ne pas lire la suite de cette section.

\medskip\noindent
Fixons un anneau gradué quelconque $S$ et soit $T$ un schéma. Deux triplets
$(d, \mathcal{L}, \psi)$ et $(d', \mathcal{L}', \psi')$ sur $T$, où les
entiers $d$ et $d'$ peuvent être distincts, sont
dits {\it équivalents} s'il existe un isomorphisme
$\beta : \mathcal{L}^{\otimes d'} \to (\mathcal{L}')^{\otimes d}$
de faisceaux inversibles sur $T$ tel que
$\beta \circ \psi|_{S^{(dd')}}$ et $\psi'|_{S^{(dd')}}$ coïncident comme
homomorphismes d'anneaux gradués
$S^{(dd')} \to \Gamma_*(Y, (\mathcal{L}')^{\otimes dd'})$.

\begin{lemma}
\label{constructions-lemma-equivalent}
Soit $S$ un anneau gradué, posons $X = \text{Proj}(S)$ et soit $T$ un schéma.
Soient $(d, \mathcal{L}, \psi)$ et $(d', \mathcal{L}', \psi')$ deux triplets
sur $T$. Les assertions suivantes sont équivalentes :
\begin{enumerate}
\item Posons $n = \text{lcm}(d, d')$ et écrivons $n = ad = a'd'$.
Il existe un isomorphisme
$\beta : \mathcal{L}^{\otimes a} \to (\mathcal{L}')^{\otimes a'}$ tel que
$\beta \circ \psi|_{S^{(n)}}$ et $\psi'|_{S^{(n)}}$ coïncident comme
homomorphismes d'anneaux gradués
$S^{(n)} \to \Gamma_*(Y, (\mathcal{L}')^{\otimes n})$.
\item Les triplets $(d, \mathcal{L}, \psi)$ et
$(d', \mathcal{L}', \psi')$ sont équivalents.
\item Pour un certain entier strictement positif $n = ad = a'd'$, il existe un
isomorphisme
$\beta : \mathcal{L}^{\otimes a} \to (\mathcal{L}')^{\otimes a'}$ tel que
$\beta \circ \psi|_{S^{(n)}}$ et $\psi'|_{S^{(n)}}$ coïncident comme
homomorphismes d'anneaux gradués
$S^{(n)} \to \Gamma_*(Y, (\mathcal{L}')^{\otimes n})$.
\item Les morphismes $\varphi : T \to X$ et $\varphi' : T \to X$ associés
respectivement à $(d, \mathcal{L}, \psi)$ et
$(d', \mathcal{L}', \psi')$ sont égaux.
\end{enumerate}
\end{lemma}

\begin{proof}
Il est clair que (1) entraîne (2), et
(2) entraîne (3) en se restreignant à
des degrés plus divisibles et en prenant des puissances des faisceaux
inversibles. En outre, (3) entraîne (4) par l'assertion d'unicité du
Lemme \ref{constructions-lemma-converse-construction}. Il reste donc à montrer que (4)
entraîne (1). Supposons (4), c'est-à-dire $\varphi = \varphi'$.
On peut alors écrire
$\mathcal{L} = \varphi^*\mathcal{O}_X(d)$ et
$\mathcal{L}' = \varphi^*\mathcal{O}_X(d')$.
En outre, par ces identifications,
les homomorphismes d'anneaux gradués $\psi$ et $\psi'$ correspondent à la
restriction de l'homomorphisme canonique d'anneaux gradués
$$
S \longrightarrow \bigoplus\nolimits_{n \geq 0} \Gamma(X, \mathcal{O}_X(n))
$$
à $S^{(d)}$ et à $S^{(d')}$, composée avec l'image inverse par $\varphi$
(encore une fois par le Lemme
\ref{constructions-lemma-converse-construction}). Il suffit
donc de prendre pour $\beta$ l'isomorphisme
$$
(\varphi^*\mathcal{O}_X(d))^{\otimes a} =
\varphi^*\mathcal{O}_X(n) =
(\varphi^*\mathcal{O}_X(d'))^{\otimes a'}
$$
\end{proof}

\noindent
Soit $S$ un anneau gradué et soit $X = \text{Proj}(S)$. Sur le sous-schéma
ouvert $U_d \subset X = \text{Proj}(S)$ défini en (\ref{constructions-equation-Ud}), nous
disposons du triplet $(d, \mathcal{O}_{U_d}(d), \psi^d)$. Si $d | d'$, les
triplets $(d, \mathcal{O}_{U_d}(d), \psi^d)$ et
$(d', \mathcal{O}_{U_{d'}}(d'), \psi^{d'})$, restreints à l'ouvert $U_d$
(qui est contenu dans $U_{d'}$), sont
manifestement équivalents. Ce fait et le
Lemme \ref{constructions-lemma-converse-construction} montrent que les morphismes
$Y \to X$ correspondent approximativement aux classes d'équivalence de
triplets sur $Y$. Cette affirmation n'est pas tout à fait exacte : si $Y$
n'est pas quasi-compact, il se peut qu'aucun triplet unique ne convienne.
Il faut donc prendre quelques précautions pour définir le foncteur
correspondant.

\medskip\noindent
Voici une façon de procéder. Supposons
$d' = ad$. Considérons la transformation
de foncteurs $F_d \to F_{d'}$ qui associe au triplet
$(d, \mathcal{L}, \psi)$ sur $T$ le triplet
$(d', \mathcal{L}^{\otimes a}, \psi|_{S^{(d')}})$.
L'une des implications du Lemme \ref{constructions-lemma-equivalent} montre que cette
transformation $F_d \to F_{d'}$ est injective. Pour un schéma quasi-compact
$T$, posons
$$
F(T) = \bigcup\nolimits_{d \in \mathbf{N}} F_d(T)
$$
Les applications de transition sont celles décrites ci-dessus. On obtient
ainsi un foncteur contravariant de la
catégorie des schémas quasi-compacts vers
la catégorie des ensembles. Pour un schéma quelconque $T$, posons
$$
F(T)
=
\lim_{V \subset T\text{ quasi-compact open}} F(V).
$$
Autrement dit, un élément $\xi$ de $F(T)$ correspond au choix compatible
d'éléments $\xi_V \in F(V)$, où $V$ parcourt les ouverts quasi-compacts de
$T$. Nous omettons la définition de l'application d'image inverse
$F(T) \to F(T')$ associée à un morphisme de schémas $T' \to T$.
Nous avons ainsi défini le foncteur
\begin{eqnarray*}
F : \Sch^{opp} & \longrightarrow & \textit{Ensembles}
\end{eqnarray*}

\begin{lemma}
\label{constructions-lemma-proj-functor}
Soit $S$ un anneau gradué et soit
$X = \text{Proj}(S)$. Le foncteur $F$ défini
ci-dessus est représentable par le schéma $X$.
\end{lemma}

\begin{proof}
Nous avons vu que le foncteur $F_d$ correspond au sous-schéma ouvert
$U_d \subset X$. En outre, la transformation de foncteurs
$F_d \to F_{d'}$ (si $d | d'$) définie ci-dessus correspond au morphisme
d'inclusion $U_d \to U_{d'}$ ; voir la discussion précédente. Pour montrer
que $F$ est représenté par $X$, il suffit donc de vérifier que tout morphisme
$T \to X$, où $T$ est quasi-compact, se
factorise par un certain $U_d$, et que,
pour tout schéma $T$, on a
$$
\Mor(T, X)
=
\lim_{V \subset T\text{ quasi-compact open}} \Mor(V, X).
$$
Nous omettons ces vérifications.
\end{proof}







\section{Espace projectif}
\label{constructions-section-projective-space}

\noindent
L'espace projectif est l'un des objets fondamentaux de la géométrie
algébrique. Dans cette section, nous nous contentons de le construire comme le
$\text{Proj}$ d'un anneau de polynômes. Nous découvrirons plus loin nombre de
ses belles propriétés.

\begin{lemma}
\label{constructions-lemma-projective-space}
Soit $S = \mathbf{Z}[T_0, \ldots, T_n]$ avec $\deg(T_i) = 1$. Le schéma
$$
\mathbf{P}^n_{\mathbf{Z}} = \text{Proj}(S)
$$
représente le foncteur qui associe à un schéma $Y$ les couples
$(\mathcal{L}, (s_0, \ldots, s_n))$ tels que
\begin{enumerate}
\item $\mathcal{L}$ soit un $\mathcal{O}_Y$-module inversible ;
\item $s_0, \ldots, s_n$ soient des
sections globales de $\mathcal{L}$ qui
engendrent $\mathcal{L}$,
\end{enumerate}
à l'équivalence suivante près :
$(\mathcal{L}, (s_0, \ldots, s_n)) \sim
(\mathcal{N}, (t_0, \ldots, t_n))$ $\Leftrightarrow$ il existe un
isomorphisme $\beta : \mathcal{L} \to \mathcal{N}$ tel que
$\beta(s_i) = t_i$ pour $i = 0, \ldots, n$.
\end{lemma}

\begin{proof}
C'est un cas particulier du Lemme \ref{constructions-lemma-apply}. En effet, pour tout
anneau gradué $A$, on a
\begin{eqnarray*}
\Mor_{graded rings}(\mathbf{Z}[T_0, \ldots, T_n], A)
& = &
A_1 \times \ldots \times A_1 \\
\psi & \mapsto & (\psi(T_0), \ldots, \psi(T_n))
\end{eqnarray*}
et la partie de degré $1$ de $\Gamma_*(Y, \mathcal{L})$ n'est autre que
$\Gamma(Y, \mathcal{L})$.
\end{proof}

\begin{definition}
\label{constructions-definition-projective-space}
Le schéma
$\mathbf{P}^n_{\mathbf{Z}} = \text{Proj}(\mathbf{Z}[T_0, \ldots, T_n])$ est
appelé l'{\it espace projectif de dimension $n$ sur $\mathbf{Z}$}.
Son changement de base $\mathbf{P}^n_S$ à un schéma $S$ est appelé l'
{\it espace projectif de dimension $n$ sur $S$}. Si $R$ est un anneau, le
changement de base à $\Spec(R)$ est noté $\mathbf{P}^n_R$ et appelé l'
{\it espace projectif de dimension $n$ sur $R$}.
\end{definition}

\noindent
Soient $Y$ un schéma sur $S$ et
$(\mathcal{L}, (s_0, \ldots, s_n))$ un couple comme dans le
Lemme \ref{constructions-lemma-projective-space}. Le morphisme induit vers
$\mathbf{P}^n_S$ est noté
$$
\varphi_{(\mathcal{L}, (s_0, \ldots, s_n))} :
Y \longrightarrow \mathbf{P}^n_S
$$
Cette notation a un sens : le couple définit un morphisme vers
$\mathbf{P}^n_{\mathbf{Z}}$, tandis que le morphisme structural vers $S$ est
déjà donné ; ensemble, ils définissent un morphisme vers
$\mathbf{P}^n_S = \mathbf{P}^n_{\mathbf{Z}} \times S$.
Par le Lemme \ref{constructions-lemma-relative-proj-base-change}, on peut identifier
$\mathbf{P}^n_S$ au Proj relatif de la $\mathcal{O}_S$-algèbre graduée
$\mathcal{O}_S[T_0, \ldots, T_n]$. On peut alors caractériser
$\varphi_{(\mathcal{L}, (s_0, \ldots, s_n))}$ comme l'unique $S$-morphisme tel
que
$$
\mathcal{L} =
\varphi_{(\mathcal{L}, (s_0, \ldots, s_n))}^*\mathcal{O}_{\mathbf{P}^n_S}(1)
\quad
\text{et}
\quad
s_i = \varphi_{(\mathcal{L}, (s_0, \ldots, s_n))}^*T_i
$$
où l'on considère $T_i$ comme une section globale de
$\mathcal{O}_{\mathbf{P}^n_S}(1)$ par l'intermédiaire de l'application $\psi$
du Lemme \ref{constructions-lemma-glue-relative-proj-twists} ; nous omettons les détails.

\begin{lemma}
\label{constructions-lemma-standard-covering-projective-space}
L'espace projectif de dimension $n$ sur $\mathbf{Z}$ est recouvert par les
$n + 1$ ouverts standards
$$
\mathbf{P}^n_{\mathbf{Z}} =
\bigcup\nolimits_{i = 0, \ldots, n} D_{+}(T_i)
$$
chacun des $D_{+}(T_i)$ étant isomorphe à $\mathbf{A}^n_{\mathbf{Z}}$,
l'espace affine de dimension $n$ sur $\mathbf{Z}$.
\end{lemma}

\begin{proof}
Cela résulte de
$\mathbf{Z}[T_0, \ldots, T_n]_{+} = (T_0, \ldots, T_n)$
et de l'isomorphisme
évident
$$
\Spec
\left(
\mathbf{Z}
\left[\frac{T_0}{T_i}, \ldots, \frac{T_n}{T_i}
\right]
\right)
\cong
\mathbf{A}^n_{\mathbf{Z}}
$$
\end{proof}

\begin{lemma}
\label{constructions-lemma-projective-space-separated}
Soit $S$ un schéma. Le morphisme structural $\mathbf{P}^n_S \to S$
possède les propriétés suivantes :
\begin{enumerate}
\item il est séparé ;
\item il est quasi-compact ;
\item il satisfait les parties existence
et unicité du critère valuatif ;
\item il est universellement fermé.
\end{enumerate}
\end{lemma}

\begin{proof}
Toutes ces propriétés sont stables par changement de base (c'est clair pour
les deux dernières ; pour les deux premières, voir Schémas, Lemmes
\ref{schemes-lemma-separated-permanence} et
\ref{schemes-lemma-quasi-compact-preserved-base-change}). Il suffit donc de les
démontrer pour le morphisme
$\mathbf{P}^n_{\mathbf{Z}} \to \Spec(\mathbf{Z})$.
La séparation est le Lemme \ref{constructions-lemma-proj-separated}. La quasi-compacité
résulte du Lemme \ref{constructions-lemma-standard-covering-projective-space}.
Les parties existence et unicité du critère valuatif résultent du
Lemme \ref{constructions-lemma-proj-valuative-criterion}.
Enfin, le caractère universellement
fermé découle des assertions précédentes et de Schémas,
Proposition \ref{schemes-proposition-characterize-universally-closed}.
\end{proof}

\begin{remark}
\label{constructions-remark-missing-finite-type}
Que manque-t-il à la liste de propriétés ci-dessus ? Certainement la propriété
d'être de type fini. Si nous ne l'énonçons pas ici, c'est que nous n'avons pas
encore défini la notion de type fini à ce stade. Une autre propriété absente
est la « lissité » ; il y en a sans doute beaucoup d'autres auxquelles le
lecteur pourra penser.
\end{remark}

\begin{lemma}[Plongement de Segre]
\label{constructions-lemma-segre-embedding}
Soit $S$ un schéma. Il existe une immersion fermée
$$
\mathbf{P}^n_S \times_S \mathbf{P}^m_S
\longrightarrow
\mathbf{P}^{nm + n + m}_S
$$
appelée le {\it plongement de Segre}.
\end{lemma}

\begin{proof}
Il suffit de le démontrer lorsque
$S = \Spec(\mathbf{Z})$. Nous omettrons donc
l'indice $S$ et travaillerons dans la situation absolue. Écrivons
$\mathbf{P}^n = \text{Proj}(\mathbf{Z}[X_0, \ldots, X_n])$,
$\mathbf{P}^m = \text{Proj}(\mathbf{Z}[Y_0, \ldots, Y_m])$ et
$\mathbf{P}^{nm + n + m} =
\text{Proj}(\mathbf{Z}[Z_0, \ldots, Z_{nm + n + m}])$.
Pour définir un morphisme vers
$\mathbf{P}^{nm + n + m}$, il faut se donner
sur le membre de gauche un faisceau inversible $\mathcal{L}$ et
$(n + 1)(m + 1)$ sections $s_i$ qui l'engendrent ; voir le
Lemme \ref{constructions-lemma-projective-space}. Prenons le faisceau inversible
$$
\mathcal{L} =
\text{pr}_1^*\mathcal{O}_{\mathbf{P}^n}(1)
\otimes
\text{pr}_2^*\mathcal{O}_{\mathbf{P}^m}(1)
$$
et les sections
$$
s_0 = X_0Y_0, \ s_1 = X_1Y_0, \ldots, \ s_n = X_nY_0,
\ s_{n + 1} = X_0Y_1, \ldots, \ s_{nm + n + m} = X_nY_m.
$$
Elles engendrent $\mathcal{L}$, puisque les sections $X_i$ engendrent
$\mathcal{O}_{\mathbf{P}^n}(1)$ et les sections $Y_j$ engendrent
$\mathcal{O}_{\mathbf{P}^m}(1)$. Le morphisme induit $\varphi$ vérifie
$$
\varphi^{-1}(D_{+}(Z_{i + (n + 1)j})) = D_{+}(X_i) \times D_{+}(Y_j).
$$
C'est donc un morphisme affine. Dans le cas $(i, j) = (0, 0)$,
l'homomorphisme d'anneaux correspondant est
$$
\mathbf{Z}[Z_1/Z_0, \ldots, Z_{nm + n + m}/Z_0]
\longrightarrow
\mathbf{Z}[X_1/X_0, \ldots, X_n/X_0, Y_1/Y_0, \ldots, Y_n/Y_0]
$$
qui envoie $Z_i/Z_0$ sur $X_i/X_0$ pour $i \leq n$ et
$Z_{(n + 1)j}/Z_0$ sur $Y_j/Y_0$. Il est
donc surjectif. Le même raisonnement
s'applique aux autres ouverts affines.
Ainsi, $\varphi$ est une immersion
fermée (voir Schémas, Lemme \ref{schemes-lemma-closed-local-target} et
Exemple \ref{schemes-example-closed-immersion-affines}).
\end{proof}

\noindent
Les deux lemmes suivants sont des cas particuliers de résultats plus généraux
établis plus loin, mais il est sans doute utile de les démontrer directement
ici. Le premier est un cas particulier de Diviseurs, Lemme \ref{divisors-lemma-closed-subscheme-proj}.

\begin{lemma}
\label{constructions-lemma-closed-in-projective-space}
Soit $R$ un anneau et soit $Z \subset \mathbf{P}^n_R$ un sous-schéma fermé.
Posons
$$
I_d = \Ker\left(
R[T_0, \ldots, T_n]_d
\longrightarrow
\Gamma(Z, \mathcal{O}_{\mathbf{P}^n_R}(d)|_Z)\right)
$$
Alors $I = \bigoplus I_d \subset R[T_0, \ldots, T_n]$ est un idéal gradué et
$Z = \text{Proj}(R[T_0, \ldots, T_n]/I)$.
\end{lemma}

\begin{proof}
Il est clair que $I$ est un idéal gradué. Posons
$Z' = \text{Proj}(R[T_0, \ldots, T_n]/I)$. D'après le
Lemme \ref{constructions-lemma-surjective-graded-rings-generated-degree-1-map-proj},
$Z'$ est
un sous-schéma fermé de $\mathbf{P}^n_R$. Pour vérifier l'égalité $Z = Z'$,
il suffit de se placer sur un ouvert affine standard $D_{+}(T_i)$.
Après renumérotation des coordonnées homogènes, on peut supposer $i = 0$.
Écrivons que $Z \cap D_{+}(T_0)$, resp.\ $Z' \cap D_{+}(T_0)$,
est défini par l'idéal $J$, resp.\ $J'$, de
$R[T_1/T_0, \ldots, T_n/T_0]$. Alors $J'$ est engendré par les éléments
$F/T_0^{\deg(F)}$, où $F \in I$ est homogène.
Supposons que $F \in I$ soit de degré $d$. Comme $F$ s'annule en tant que
section de $\mathcal{O}_{\mathbf{P}^n_R}(d)$ restreinte à $Z$, l'élément
$F/T_0^d$ appartient à $J$. Ainsi, $J' \subset J$.

\medskip\noindent
Réciproquement, soit $f \in J$. Si $f$ est de degré total $d$ en
$T_1/T_0, \ldots, T_n/T_0$, on peut écrire
$f = F/T_0^d$ pour un certain $F \in R[T_0, \ldots, T_n]_d$.
Choisissons $i \in \{1, \ldots, n\}$. Le sous-schéma
$Z \cap D_{+}(T_i)$ est défini par un certain idéal
$J_i \subset R[T_0/T_i, \ldots, T_n/T_i]$. En outre,
$$
J \cdot
R\left[
\frac{T_1}{T_0}, \ldots, \frac{T_n}{T_0},
\frac{T_0}{T_i}, \ldots, \frac{T_n}{T_i}
\right]
=
J_i \cdot
R\left[
\frac{T_1}{T_0}, \ldots, \frac{T_n}{T_0},
\frac{T_0}{T_i}, \ldots, \frac{T_n}{T_i}
\right]
$$
Le membre de gauche est le localisé de
$J$ par rapport à $T_i/T_0$, et celui
de droite le localisé de $J_i$ par rapport à $T_0/T_i$. Il s'ensuit que
$T_0^{d_i}F/T_i^{d + d_i}$ appartient à $J_i$ pour un entier $d_i$
suffisamment grand. Par conséquent,
$T_0^{\max(d_i)}F$ appartient à $I$,
car sa restriction à chacun des ouverts affines standards $D_{+}(T_i)$
s'annule sur le sous-schéma fermé $Z \cap D_{+}(T_i)$.
Ainsi, $f \in J'$, d'où $J \subset J'$, comme voulu.
\end{proof}

\noindent
Le lemme suivant est un cas particulier des résultats plus généraux de
Propriétés, Lemmes \ref{properties-lemma-ample-quasi-coherent} ou
\ref{properties-lemma-proj-quasi-coherent}.

\begin{lemma}
\label{constructions-lemma-quasi-coherent-projective-space}
Soit $R$ un anneau et soit $\mathcal{F}$ un faisceau quasi-cohérent sur
$\mathbf{P}^n_R$. Pour $d \geq 0$, posons
$$
M_d
=
\Gamma(\mathbf{P}^n_R,
\mathcal{F} \otimes_{\mathcal{O}_{\mathbf{P}^n_R}}
\mathcal{O}_{\mathbf{P}^n_R}(d))
=
\Gamma(\mathbf{P}^n_R, \mathcal{F}(d))
$$
Alors $M = \bigoplus_{d \geq 0} M_d$ est un
$R[T_0, \ldots, R_n]$-module gradué et il existe un isomorphisme canonique
$\mathcal{F} = \widetilde{M}$.
\end{lemma}

\begin{proof}
Les applications de multiplication
$$
R[T_0, \ldots, R_n]_e \times M_d \longrightarrow M_{d + e}
$$
proviennent des isomorphismes naturels
$$
\mathcal{O}_{\mathbf{P}^n_R}(e)
\otimes_{\mathcal{O}_{\mathbf{P}^n_R}}
\mathcal{F}(d)
\longrightarrow
\mathcal{F}(e + d)
$$
voir l'Équation (\ref{constructions-equation-global-sections-module}). Construisons le
morphisme $c : \widetilde{M} \to \mathcal{F}$. Sur chacun des ouverts affines
standards $U_i = D_{+}(T_i)$, on a
$\Gamma(U_i, \widetilde{M}) = (M[1/T_i])_0$, où l'indice ${}_0$ désigne la
partie de degré $0$. Tout élément de ce module s'écrit $m/T_i^d$ avec
$m \in M_d$. Comme $T_i$ engendre $\mathcal{O}(1)$ au-dessus de $U_i$, on
peut toujours écrire $m|_{U_i} = m_i \otimes T_i^d$, avec une unique section
$m_i \in \Gamma(U_i, \mathcal{F})$. Il est donc naturel de poser
$c(m/T_i^d) = m_i$. Un raisonnement élémentaire, omis ici, montre que cela
définit un morphisme $c : \widetilde{M} \to \mathcal{F}$ dès que l'on vérifie
$$
(T_i/T_j)^d m_i|_{U_i \cap U_j} = m_j|_{U_i \cap U_j}
$$
dans $M[1/T_iT_j]$. Or cela est clair,
car sur l'intersection les générateurs
$T_i$ et $T_j$ de $\mathcal{O}(1)$ diffèrent par la fonction inversible
$T_i/T_j$.

\medskip\noindent
Injectivité de $c$. On peut la vérifier sur les ouverts affines $U_i$.
Soit $i \in \{0, \ldots, n\}$ et soit l'élément $s$ donné par
$s = m/T_i^d \in \Gamma(U_i, \widetilde{M})$ tel que
$c(m/T_i^d) = 0$. D'après la description de $c$, on a $m_i = 0$, donc
$m|_{U_i} = 0$. On a alors $T_i^em = 0$ dans $M$ pour un certain $e$.
Par conséquent, $s = m/T_i^d = T_i^e/T_i^{e + d} = 0$, comme voulu.

\medskip\noindent
Surjectivité de $c$. On peut la vérifier sur les ouverts affines $U_i$ et,
après renumérotation, il suffit de la vérifier sur $U_0$.
Soit $s \in \mathcal{F}(U_0)$. Écrivons
$\mathcal{F}|_{U_i} = \widetilde{N_i}$ pour un certain
$R[T_0/T_i, \ldots, T_0/T_i]$-module $N_i$ ; c'est possible parce que
$\mathcal{F}$ est quasi-cohérent. La section $s$ correspond alors à un élément
$x \in N_0$. On a
$$
(N_i)_{T_j/T_i} \cong (N_j)_{T_i/T_j}
$$
où les indices désignent la localisation principale, comme modules sur
l'anneau
$$
R\left[
\frac{T_0}{T_i}, \ldots, \frac{T_n}{T_i},
\frac{T_0}{T_j}, \ldots, \frac{T_n}{T_j}
\right].
$$
Il existe donc un entier $d$ assez grand et des éléments
$s_i \in N_i$, $i = 1, \ldots, n$, tels que
$$
s = (T_i/T_0)^d s_i
$$
sur $U_0 \cap U_i$. Considérons ensuite la différence
$$
t_{ij} = s_i - (T_j/T_i)^d s_j
$$
sur $U_i \cap U_j$, $0 < i < j$. Par le choix des $s_i$, on a
$t_{ij}|_{U_0 \cap U_i \cap U_j} = 0$. Il existe donc un entier $e$ assez
grand tel que $(T_0/T_i)^et_{ij} = 0$. Posons
$s_i' = (T_0/T_i)^es_i$ et $s_0' = s$. On a alors
$$
s_a' = (T_b/T_a)^{e + d} s_b'
$$
sur $U_a \cap U_b$ pour tous $a, b$. C'est exactement la condition pour que
les éléments $s'_a$ se recollent en une section globale
$m \in \Gamma(\mathbf{P}^n_R, \mathcal{F}(e + d))$.
Par construction, $c(m/T_0^{e + d}) = s$. Le morphisme $c$ est donc surjectif,
ce qui achève la démonstration.
\end{proof}

\begin{lemma}
\label{constructions-lemma-globally-generated-omega-twist-1}
Soit $X$ un schéma. Soit $\mathcal{L}$ un faisceau inversible, et soient
$s_0, \ldots, s_n$ des sections globales de $\mathcal{L}$ qui l'engendrent.
Soit $\mathcal{F}$ le noyau du morphisme induit
$\mathcal{O}_X^{\oplus n + 1} \to \mathcal{L}$. Alors
$\mathcal{F} \otimes \mathcal{L}$ est engendré par ses sections globales.
\end{lemma}

\begin{proof}
Le résultat est même vrai lorsque $X$ est un espace localement annelé
quelconque. Le faisceau $\mathcal{F}$ est un $\mathcal{O}_X$-module localement
libre de rang fini $n$. Les éléments
$$
s_{ij} = (0, \ldots, 0, s_j, 0, \ldots, 0, -s_i, 0, \ldots, 0)
\in \Gamma(X, \mathcal{L}^{\oplus n + 1})
$$
où $s_j$ occupe la $i$-ième place et $-s_i$ la $j$-ième s'envoient sur zéro
dans $\mathcal{L}^{\otimes 2}$. Ainsi,
$s_{ij} \in \Gamma(X, \mathcal{F} \otimes_{\mathcal{O}_X} \mathcal{L})$.
Un calcul local montre que ces sections engendrent
$\mathcal{F} \otimes \mathcal{L}$.

\medskip\noindent
Autre démonstration. Considérons le morphisme
$\varphi : X \to \mathbf{P}^n_\mathbf{Z}$ associé au couple
$(\mathcal{L}, (s_0, \ldots, s_n))$. Comme l'image inverse de
$\mathcal{O}(1)$ est $\mathcal{L}$ et
celle de $T_i$ est $s_i$, il suffit de
démontrer le lemme pour $\mathbf{P}^n_\mathbf{Z}$.
Dans ce cas, le faisceau $\mathcal{F}$ correspond au
$S = \mathbf{Z}[T_0, \ldots, T_n]$-module gradué
$M$ qui s'insère dans la suite exacte courte
$$
0 \to M \to S^{\oplus n + 1} \to S(1) \to 0
$$
où le second morphisme est donné par $T_0, \ldots, T_n$.
L'assertion précédente se traduit alors par le fait que les éléments
$$
T_{ij} = (0, \ldots, 0, T_j, 0, \ldots, 0, -T_i, 0, \ldots, 0)
\in M(1)_0
$$
engendrent le module gradué $M(1)$ sur $S$. Nous omettons les détails.
\end{proof}








\section{Faisceaux inversibles et morphismes vers Proj}
\label{constructions-section-invertible-proj}

\noindent
Soit $T$ un schéma et soit $\mathcal{L}$ un faisceau inversible sur $T$.
Pour une section $s \in \Gamma(T, \mathcal{L})$, on note $T_s$ l'ouvert
des points où $s$ ne s'annule pas. Voir Modules,
Lemme \ref{modules-lemma-s-open}. On peut considérer le lemme suivant comme
une légère généralisation du Lemme \ref{constructions-lemma-apply}. C'est également une
généralisation du Lemme \ref{constructions-lemma-morphism-proj}.

\begin{lemma}
\label{constructions-lemma-invertible-map-into-proj}
Soit $A$ un anneau gradué et posons $X = \text{Proj}(A)$.
Soient $T$ un schéma, $\mathcal{L}$ un
$\mathcal{O}_T$-module inversible et
$\psi : A \to \Gamma_*(T, \mathcal{L})$
un homomorphisme d'anneaux gradués.
Posons
$$
U(\psi) = \bigcup\nolimits_{f \in A_{+}\text{ homogeneous}} T_{\psi(f)}
$$
Le morphisme $\psi$ induit un morphisme canonique de schémas
$$
r_{\mathcal{L}, \psi} :
U(\psi) \longrightarrow X
$$
muni d'un morphisme d'algèbres $\mathbf{Z}$-graduées sur $\mathcal{O}_T$
$$
\theta :
r_{\mathcal{L}, \psi}^*\left(
\bigoplus\nolimits_{d \in \mathbf{Z}} \mathcal{O}_X(d)
\right)
\longrightarrow
\bigoplus\nolimits_{d \in \mathbf{Z}} \mathcal{L}^{\otimes d}|_{U(\psi)}.
$$
Le triplet $(U(\psi), r_{\mathcal{L}, \psi}, \theta)$ est caractérisé par
les propriétés suivantes :
\begin{enumerate}
\item Pour tout $f \in A_{+}$ homogène, on a
$r_{\mathcal{L}, \psi}^{-1}(D_{+}(f)) = T_{\psi(f)}$.
\item Pour tout $d \geq 0$, le diagramme
$$
\xymatrix{
A_d \ar[d]_{(\ref{constructions-equation-global-sections})} \ar[r]_{\psi} &
\Gamma(T, \mathcal{L}^{\otimes d}) \ar[d]^{restrict} \\
\Gamma(X, \mathcal{O}_X(d)) \ar[r]^{\theta} &
\Gamma(U(\psi), \mathcal{L}^{\otimes d})
}
$$
est commutatif.
\end{enumerate}
De plus, pour tout $d \geq 1$ et tout
sous-schéma ouvert $V \subset T$ tels
que les sections de $\psi(A_d)$ engendrent $\mathcal{L}^{\otimes d}|_V$,
le morphisme $r_{\mathcal{L}, \psi}|_V$ coïncide avec le morphisme
$\varphi : V \to \text{Proj}(A)$ et le morphisme $\theta|_V$ coïncide avec
$\alpha : \varphi^*\mathcal{O}_X(d) \to \mathcal{L}^{\otimes d}|_V$,
où $(\varphi, \alpha)$ est le couple du Lemme
\ref{constructions-lemma-converse-construction} associé à
$\psi|_{A^{(d)}} : A^{(d)} \to \Gamma_*(V, \mathcal{L}^{\otimes d})$.
\end{lemma}

\begin{proof}
Supposons que deux triplets $(U, r : U \to X, \theta)$ et
$(U', r' : U' \to X, \theta')$ satisfassent (1) et (2).
La propriété (1) entraîne $U = U' = U(\psi)$ ainsi que l'égalité
$r = r'$ des applications continues
sous-jacentes : les ouverts $D_{+}(f)$
forment en effet une base de la topologie de $X$, et $X$ est un espace
topologique sobre (voir Algèbre, section \ref{algebra-section-proj} et
Schémas, Lemme \ref{schemes-lemma-scheme-sober}).
Soit $f \in A_{+}$ homogène. On a
$\Gamma(D_{+}(f), \bigoplus_{n \in \mathbf{Z}} \mathcal{O}_X(n)) = A_f$
comme algèbres $\mathbf{Z}$-graduées.
Considérons les deux homomorphismes
d'anneaux $\mathbf{Z}$-gradués
$$
\theta, \theta' :
A_f
\longrightarrow
\Gamma(T_{\psi(f)}, \bigoplus \mathcal{L}^{\otimes n}).
$$
La multiplication par $f$ (resp.\ $\psi(f)$) est un isomorphisme dans le
membre de gauche (resp.\ de droite). On sait aussi, par (2), que
$\theta(x/1) = \theta'(x/1) = \psi(x)|_{T_{\psi(f)}}$
pour tout $x \in A$.
On en déduit aisément $\theta = \theta'$. En considérant les composantes
de degré $0$, on obtient $r^\sharp = (r')^\sharp$, c'est-à-dire
$r = r'$ comme morphismes de schémas. Cela prouve l'unicité.

\medskip\noindent
Passons à l'existence. Grâce à l'unicité que nous venons de démontrer,
il suffit de construire le couple $(r, \theta)$ localement sur $T$.
On peut donc supposer $T = \Spec(R)$ affine, $\mathcal{L} = \mathcal{O}_T$,
et qu'il existe $f \in A_{+}$ homogène tel que $\psi(f)$ engendre
$\mathcal{O}_T = \mathcal{O}_T^{\otimes \deg(f)}$.
Autrement dit, $\psi(f) = u \in R^*$ est inversible. Dans ce cas, $\psi$
est un homomorphisme d'anneaux gradués
$$
A \longrightarrow R[x] = \Gamma_*(T, \mathcal{O}_T)
$$ qui envoie $f$ sur $ux^{\deg(f)}$. Il se prolonge de façon unique en un
homomorphisme d'anneaux $\mathbf{Z}$-gradués
$\theta : A_f \to R[x, x^{-1}]$ en envoyant $1/f$ sur
$u^{-1}x^{-\deg(f)}$. En degré zéro, on obtient l'homomorphisme d'anneaux
$A_{(f)} \to R$ qui définit le morphisme
$r : T = \Spec(R) \to \Spec(A_{(f)}) = D_{+}(f) \subset X$.
Nous avons donc construit $(r, \theta)$ dans ce cas particulier.

\medskip\noindent
Montrons la dernière assertion du lemme. D'après le
Lemme \ref{constructions-lemma-converse-construction},
le morphisme qui y est construit
est l'unique morphisme rendant commutatif le diagramme de son énoncé.
La commutativité du diagramme du présent
lemme entraîne celle du diagramme
restreint à $V$ et à $A^{(d)}$. D'où le résultat.
\end{proof}

\begin{remark}
\label{constructions-remark-not-in-invertible-locus}
Sous les hypothèses du Lemme \ref{constructions-lemma-invertible-map-into-proj}, l'image
du morphisme $r_{\mathcal{L}, \psi}$ n'est pas nécessairement contenue dans
le lieu où le faisceau $\mathcal{O}_X(1)$ est inversible. Voici un exemple.
Soit $k$ un corps et soit $S = k[A, B, C]$ gradué par
$\deg(A) = 1$, $\deg(B) = 2$, $\deg(C) = 3$.
Posons $X = \text{Proj}(S)$ et soit
$T = \mathbf{P}^2_k = \text{Proj}(k[X_0, X_1, X_2])$.
Rappelons que $\mathcal{L} = \mathcal{O}_T(1)$ est inversible et que
$\mathcal{O}_T(n) = \mathcal{L}^{\otimes n}$.
Considérons la composée $\psi$ des homomorphismes
$$
S \to k[X_0, X_1, X_2] \to \Gamma_*(T, \mathcal{L}).
$$
Le premier est $A \mapsto X_0$, $B \mapsto X_1^2$,
$C \mapsto X_2^3$ ;
le second est (\ref{constructions-equation-global-sections}). D'après le lemme, cette
composée correspond à un morphisme
$r_{\mathcal{L}, \psi} : T \to X = \text{Proj}(S)$, dont on voit aisément
qu'il est surjectif. Or, dans la Remarque \ref{constructions-remark-not-isomorphism},
nous avons montré que $\mathcal{O}_X(1)$
n'est pas inversible en tout point
de $X$.
\end{remark}










\section{Proj relatif par recollement}
\label{constructions-section-relative-proj-via-glueing}

\begin{situation}
\label{constructions-situation-relative-proj}
Ici, $S$ est un schéma et $\mathcal{A}$ une $\mathcal{O}_S$-algèbre graduée
quasi-cohérente.
\end{situation}

\noindent
Dans cette section, nous indiquons comment construire un morphisme de schémas
$$
\underline{\text{Proj}}_S(\mathcal{A}) \longrightarrow S
$$
en recollant les spectres homogènes $\text{Proj}(\Gamma(U, \mathcal{A}))$,
où $U$ parcourt les ouverts affines de $S$. Montrons d'abord que les spectres
homogènes des algèbres de sections de $\mathcal{A}$ sur les ouverts affines
forment une famille de schémas satisfaisant les conditions du
Lemme \ref{constructions-lemma-relative-glueing}.

\begin{lemma}
\label{constructions-lemma-proj-inclusion}
Plaçons-nous dans la Situation \ref{constructions-situation-relative-proj}.
Soient $U \subset U' \subset S$ des ouverts affines et posons
$A = \mathcal{A}(U)$, $A' = \mathcal{A}(U')$.
L'homomorphisme d'anneaux gradués $A' \to A$ induit un morphisme
$r : \text{Proj}(A) \to \text{Proj}(A')$, et le diagramme
$$
\xymatrix{
\text{Proj}(A) \ar[r] \ar[d] &
\text{Proj}(A') \ar[d] \\
U \ar[r] &
U'
}
$$ est cartésien. De plus, il existe des isomorphismes canoniques
$\theta : r^*\mathcal{O}_{\text{Proj}(A')}(n) \to
\mathcal{O}_{\text{Proj}(A)}(n)$ compatibles avec les morphismes de
multiplication.
\end{lemma}

\begin{proof}
Posons $R = \mathcal{O}_S(U)$ et $R' = \mathcal{O}_S(U')$.
L'homomorphisme $R \otimes_{R'} A' \to A$ est un isomorphisme puisque
$\mathcal{A}$ est quasi-cohérente (voir, par exemple, Schémas,
Lemme \ref{schemes-lemma-widetilde-pullback}).
Le résultat découle donc du
Lemme \ref{constructions-lemma-base-change-map-proj}.
\end{proof}

\noindent
En particulier, le morphisme $\text{Proj}(A) \to \text{Proj}(A')$ du lemme
est une immersion ouverte.

\begin{lemma}
\label{constructions-lemma-transitive-proj}
Plaçons-nous dans la Situation \ref{constructions-situation-relative-proj}.
Soient $U \subset U' \subset U'' \subset S$
des ouverts affines et posons
$A = \mathcal{A}(U)$, $A' = \mathcal{A}(U')$, $A'' = \mathcal{A}(U'')$.
La composée des morphismes $r : \text{Proj}(A) \to \text{Proj}(A')$ et
$r' : \text{Proj}(A') \to \text{Proj}(A'')$ du
Lemme \ref{constructions-lemma-proj-inclusion} est le morphisme
$r'' : \text{Proj}(A) \to \text{Proj}(A'')$ du
Lemme \ref{constructions-lemma-proj-inclusion}. Une assertion analogue vaut pour les
isomorphismes $\theta$.
\end{lemma}

\begin{proof}
Cela résulte du Lemme
\ref{constructions-lemma-morphism-proj-transitive}, puisque
$A'' \to A$ est la composée de $A'' \to A'$ et de $A' \to A$.
\end{proof}

\begin{lemma}
\label{constructions-lemma-glue-relative-proj}
Dans la Situation \ref{constructions-situation-relative-proj}, il existe un morphisme
de schémas
$$
\pi : \underline{\text{Proj}}_S(\mathcal{A}) \longrightarrow S
$$
ayant les propriétés suivantes :
\begin{enumerate}
\item pour tout ouvert affine $U \subset S$, il existe un isomorphisme
$i_U : \pi^{-1}(U) \to \text{Proj}(A)$, où $A = \mathcal{A}(U)$ ;
\item pour des ouverts affines $U \subset U' \subset S$, la composée
$$
\xymatrix{
\text{Proj}(A) \ar[r]^{i_U^{-1}} &
\pi^{-1}(U) \ar[rr]^{inclusion} & &
\pi^{-1}(U') \ar[r]^{i_{U'}} &
\text{Proj}(A')
}
$$
où $A = \mathcal{A}(U)$, $A' = \mathcal{A}(U')$, est l'immersion ouverte du
Lemme \ref{constructions-lemma-proj-inclusion}.
\end{enumerate}
\end{lemma}

\begin{proof}
Cela résulte immédiatement des
Lemmes \ref{constructions-lemma-relative-glueing},
\ref{constructions-lemma-proj-inclusion} et
\ref{constructions-lemma-transitive-proj}.
\end{proof}

\begin{lemma}
\label{constructions-lemma-glue-relative-proj-twists}
Dans la Situation \ref{constructions-situation-relative-proj}, le morphisme
$\pi : \underline{\text{Proj}}_S(\mathcal{A}) \to S$ du
Lemme \ref{constructions-lemma-glue-relative-proj}
est muni de la structure supplémentaire
suivante. Il existe un faisceau quasi-cohérent $\mathbf{Z}$-gradué
d'algèbres sur $\mathcal{O}_{\underline{\text{Proj}}_S(\mathcal{A})}$
$\bigoplus\nolimits_{n \in \mathbf{Z}}
\mathcal{O}_{\underline{\text{Proj}}_S(\mathcal{A})}(n)$
et un morphisme de $\mathcal{O}_S$-algèbres graduées
$$
\psi :
\mathcal{A}
\longrightarrow
\bigoplus\nolimits_{n \geq 0}
\pi_*\left(\mathcal{O}_{\underline{\text{Proj}}_S(\mathcal{A})}(n)\right)
$$
uniquement déterminés par la propriété suivante : pour tout ouvert affine
$U \subset S$, en posant $A = \mathcal{A}(U)$, il existe un isomorphisme
$$
\theta_U :
i_U^*\left(
\bigoplus\nolimits_{n \in \mathbf{Z}} \mathcal{O}_{\text{Proj}(A)}(n)
\right)
\longrightarrow
\left(
\bigoplus\nolimits_{n \in \mathbf{Z}}
\mathcal{O}_{\underline{\text{Proj}}_S(\mathcal{A})}(n)
\right)|_{\pi^{-1}(U)}
$$
d'algèbres $\mathbf{Z}$-graduées sur
$\mathcal{O}_{\pi^{-1}(U)}$ tel que le
diagramme
$$
\xymatrix{
A_n
\ar[rr]_\psi
\ar[dr]_-{(\ref{constructions-equation-global-sections})}
& &
\Gamma(\pi^{-1}(U),
\mathcal{O}_{\underline{\text{Proj}}_S(\mathcal{A})}(n)) \\
&
\Gamma(\text{Proj}(A),
\mathcal{O}_{\text{Proj}(A)}(n))
\ar[ru]_-{\theta_U}
&
}
$$
soit commutatif.
\end{lemma}

\begin{proof}
Nous allons appliquer le Lemme \ref{constructions-lemma-relative-glueing-sheaves} pour
recoller les faisceaux d'algèbres $\mathbf{Z}$-graduées
$\bigoplus_{n \in \mathbf{Z}} \mathcal{O}_{\text{Proj}(A)}(n)$ pour
$A = \mathcal{A}(U)$, $U \subset S$ ouvert affine, sur le schéma
$\underline{\text{Proj}}_S(\mathcal{A})$.
Nous avons construit les données nécessaires dans le
Lemme \ref{constructions-lemma-proj-inclusion} et vérifié la condition (d) du
Lemme \ref{constructions-lemma-relative-glueing-sheaves} dans le
Lemme \ref{constructions-lemma-transitive-proj}. On obtient ainsi
le faisceau d'algèbres $\mathbf{Z}$-graduées sur
$\mathcal{O}_{\underline{\text{Proj}}_S(\mathcal{A})}$
$\bigoplus_{n \in \mathbf{Z}}
\mathcal{O}_{\underline{\text{Proj}}_S(\mathcal{A})}(n)$,
muni des isomorphismes $\theta_U$ pour tout ouvert affine $U \subset S$
et tout $n \in \mathbf{Z}$. Pour chaque ouvert affine $U \subset S$,
en posant $A = \mathcal{A}(U)$, on dispose d'un homomorphisme
$A \to \Gamma(\text{Proj}(A),
\bigoplus_{n \geq 0} \mathcal{O}_{\text{Proj}(A)}(n))$.
Le morphisme $\psi$ existe donc par fonctorialité du recollement relatif ;
voir la Remarque \ref{constructions-remark-relative-glueing-functorial}.
Le diagramme du lemme est commutatif par construction. Cette propriété
caractérise le faisceau d'algèbres $\mathbf{Z}$-graduées sur
$\mathcal{O}_{\underline{\text{Proj}}_S(\mathcal{A})}$
$\bigoplus \mathcal{O}_{\underline{\text{Proj}}_S(\mathcal{A})}(n)$ :
la démonstration du Lemme \ref{constructions-lemma-morphism-proj} montre en effet que la
commutativité de ces diagrammes détermine de façon unique les morphismes
$\theta_U$. Nous omettons certains détails.
\end{proof}















\section{Le Proj relatif comme foncteur}
\label{constructions-section-relative-proj}

\noindent
Plaçons-nous dans la Situation \ref{constructions-situation-relative-proj}.
Ainsi, $S$ est un schéma et
$\mathcal{A} = \bigoplus_{d \geq 0} \mathcal{A}_d$ est une
$\mathcal{O}_S$-algèbre graduée quasi-cohérente. Nous allons donner une
version relative de la construction de $\text{Proj}$, en définissant un
foncteur que représentera le spectre homogène relatif. Nous construirons
ainsi un morphisme de schémas
$$
\underline{\text{Proj}}_S(\mathcal{A}) \longrightarrow S
$$
qui, au-dessus des ouverts affines de $S$, sera le spectre homogène d'un
anneau gradué. Nous suivrons la démarche adoptée pour le spectre relatif
dans la section \ref{constructions-section-spec}. La méthode plus simple consistant à
recoller les schémas $\text{Proj}(A)$, où $A = \Gamma(U, \mathcal{A})$ et
$U \subset S$ est un ouvert affine, est exposée dans la
section \ref{constructions-section-relative-proj-via-glueing}. Nous montrerons que les
deux constructions donnent des schémas isomorphes.

\medskip\noindent
Fixons pour l'instant un entier $d \geq 1$. Nous notons
$\mathcal{A}^{(d)} = \bigoplus_{n \geq 0} \mathcal{A}_{nd}$, par analogie
avec la notation d'Algèbre, section \ref{algebra-section-graded}.
Soit $T$ un schéma. Considérons les {\it quadruplets
$(d, f : T \to S, \mathcal{L}, \psi)$ sur $T$} où
\begin{enumerate}
\item $d$ est l'entier fixé ci-dessus ;
\item $f : T \to S$ est un morphisme de schémas ;
\item $\mathcal{L}$ est un $\mathcal{O}_T$-module inversible ;
\item
$\psi : f^*\mathcal{A}^{(d)} \to \bigoplus_{n \geq 0}\mathcal{L}^{\otimes n}$
est un homomorphisme de $\mathcal{O}_T$-algèbres graduées tel que
$f^*\mathcal{A}_d \to \mathcal{L}$ soit surjectif.
\end{enumerate}
Étant donnés un morphisme $h : T' \to T$ et un quadruplet
$(d, f, \mathcal{L}, \psi)$ sur $T$, on peut en prendre l'image inverse,
qui est le quadruplet $(d, f \circ h, h^*\mathcal{L}, h^*\psi)$ sur $T'$.
Deux quadruplets $(d, f, \mathcal{L}, \psi)$ et
$(d, f', \mathcal{L}', \psi')$ sur $T$, associés au même entier $d$, sont
dits {\it strictement équivalents} si
$f = f'$ et s'il existe un isomorphisme
$\beta : \mathcal{L} \to \mathcal{L}'$ tel que $\beta \circ \psi = \psi'$
comme morphismes de $\mathcal{O}_T$-algèbres graduées
$f^*\mathcal{A}^{(d)} \to \bigoplus_{n \geq 0} (\mathcal{L}')^{\otimes n}$.

\medskip\noindent
Pour tout entier $d \geq 1$, définissons
\begin{eqnarray*}
F_d : \Sch^{opp} & \longrightarrow & \textit{Ensembles}, \\
T & \longmapsto &
\{\text{strict equivalence classes of }
(d, f : T \to S, \mathcal{L}, \psi)
\text{ as above}\}
\end{eqnarray*}
avec les images inverses définies ci-dessus.

\begin{lemma}
\label{constructions-lemma-proj-base-change}
Dans la Situation \ref{constructions-situation-relative-proj},
soit $d \geq 1$ et soit
$F_d$ le foncteur associé ci-dessus à $(S, \mathcal{A})$.
Soit $g : S' \to S$ un morphisme de schémas. Posons
$\mathcal{A}' = g^*\mathcal{A}$ et
soit $F_d'$ le foncteur associé à
$(S', \mathcal{A}')$. Il existe un
isomorphisme canonique de foncteurs
$$
F'_d \cong h_{S'} \times_{h_S} F_d
$$
\end{lemma}

\begin{proof}
Se donner un quadruplet
$(d, f' : T \to S', \mathcal{L}',
\psi' : (f')^*(\mathcal{A}')^{(d)} \to
\bigoplus_{n \geq 0} (\mathcal{L}')^{\otimes n})$
revient à se donner un quadruplet
$(d, f, \mathcal{L},
\psi : f^*\mathcal{A}^{(d)} \to
\bigoplus_{n \geq 0} \mathcal{L}^{\otimes n})$
et une factorisation de $f$ sous la forme $f = g \circ f'$. La correspondance est donnée par
$f = g \circ f'$, $\mathcal{L} = \mathcal{L}'$ et $\psi = \psi'$, via les
identifications
$(f')^*(\mathcal{A}')^{(d)} = (f')^*g^*(\mathcal{A}^{(d)}) =
f^*\mathcal{A}^{(d)}$. D'où le lemme.
\end{proof}

\begin{lemma}
\label{constructions-lemma-relative-proj-affine}
Dans la Situation \ref{constructions-situation-relative-proj}, soit $F_d$ le foncteur
associé ci-dessus à $(d, S, \mathcal{A})$.
Si $S$ est affine, alors $F_d$
est représentable par le sous-schéma ouvert $U_d$ de (\ref{constructions-equation-Ud})
du schéma $\text{Proj}(\Gamma(S, \mathcal{A}))$.
\end{lemma}

\begin{proof}
Écrivons $S = \Spec(R)$ et $A = \Gamma(S, \mathcal{A})$.
Alors $A$ est une $R$-algèbre graduée et $\mathcal{A} = \widetilde A$.
Il nous faut identifier $F_d$ au
foncteur $F_d^{triples}$ des triplets
défini dans la section \ref{constructions-section-morphisms-proj}.

\medskip\noindent
Soit $(d, f : T \to S, \mathcal{L}, \psi)$ un quadruplet. On peut voir
$\psi$ comme un morphisme de $\mathcal{O}_S$-modules
$\mathcal{A}^{(d)} \to \bigoplus_{n \geq 0} f_*\mathcal{L}^{\otimes n}$.
Puisque $\mathcal{A}^{(d)}$ est quasi-cohérente, cela revient à se donner
un homomorphisme $R$-linéaire d'anneaux gradués
$A^{(d)} \to \Gamma(S, \bigoplus_{n \geq 0} f_*\mathcal{L}^{\otimes n})$.
Or $\Gamma(S, \bigoplus_{n \geq 0} f_*\mathcal{L}^{\otimes n}) =
\Gamma_*(T, \mathcal{L})$.
On peut donc associer au quadruplet le triplet $(d, \mathcal{L}, \psi)$.

\medskip\noindent
Réciproquement, soit $(d, \mathcal{L}, \psi)$ un triplet.
La composée $R \to A_0 \to \Gamma(T, \mathcal{O}_T)$ détermine un morphisme
$f : T \to S = \Spec(R)$ ; voir Schémas,
Lemme \ref{schemes-lemma-morphism-into-affine}. Avec ce choix de $f$,
l'homomorphisme
$A^{(d)} \to \Gamma(S, \bigoplus_{n \geq 0} f_*\mathcal{L}^{\otimes n})$
est $R$-linéaire et correspond donc à un morphisme $\psi$ donnant un
quadruplet $(d, f : T \to S, \mathcal{L}, \psi)$.
Nous omettons la vérification que ces constructions définissent un
isomorphisme de foncteurs $F_d = F_d^{triples}$.
\end{proof}

\begin{lemma}
\label{constructions-lemma-relative-proj-d}
Dans la Situation \ref{constructions-situation-relative-proj}, le foncteur $F_d$ est
représentable par un schéma.
\end{lemma}

\begin{proof}
Nous allons appliquer Schémas, Lemme \ref{schemes-lemma-glue-functors}.

\medskip\noindent
Vérifions d'abord que $F_d$ est un faisceau pour la topologie de Zariski.
Soient $T$ un schéma, $T = \bigcup_{i \in I} U_i$
un recouvrement ouvert et
$(d, f_i, \mathcal{L}_i, \psi_i) \in F_d(U_i)$ tels que
$(d, f_i, \mathcal{L}_i, \psi_i)|_{U_i \cap U_j}$ et
$(d, f_j, \mathcal{L}_j, \psi_j)|_{U_i \cap U_j}$ soient strictement
équivalents. Les morphismes $f_i : U_i \to S$ se recollent alors en un
morphisme de schémas $f : T \to S$ tel que $f|_{I_i} = f_i$ ; voir
Schémas, section \ref{schemes-section-glueing-schemes}.
On a donc $f_i^*\mathcal{A}^{(d)} = f^*\mathcal{A}^{(d)}|_{U_i}$.
Il existe aussi des isomorphismes
$\beta_{ij} : \mathcal{L}_i|_{U_i \cap U_j} \to \mathcal{L}_j|_{U_i \cap U_j}$
tels que $\beta_{ij} \circ \psi_i = \psi_j$ sur $U_i \cap U_j$.
Cette condition détermine les $\beta_{ij}$ de façon unique puisque les
morphismes $f_i^*\mathcal{A}_d \to \mathcal{L}_i$ sont surjectifs.
En particulier, $\beta_{jk} \circ \beta_{ij} = \beta_{ik}$ sur
$U_i \cap U_j \cap U_k$. Par Faisceaux,
section \ref{sheaves-section-glueing-sheaves}, les faisceaux inversibles
$\mathcal{L}_i$ se recollent donc en un $\mathcal{O}_T$-module inversible
$\mathcal{L}$, et les morphismes $\psi_i$ se recollent en un morphisme de
$\mathcal{O}_T$-algèbres
$\psi : f^*\mathcal{A}^{(d)} \to \bigoplus_{n \geq 0} \mathcal{L}^{\otimes n}$.
Cela prouve que $F_d$ satisfait la
condition de faisceau pour la topologie
de Zariski.

\medskip\noindent
Soit $S = \bigcup_{i \in I} U_i$ un
recouvrement par des ouverts affines.
Soit $F_{d, i} \subset F_d$ le sous-foncteur constitué des couples
$(f : T \to S, \varphi)$ tels que $f(T) \subset U_i$.

\medskip\noindent
Montrons que chaque $F_{d, i}$ est représentable. D'après le
Lemme \ref{constructions-lemma-proj-base-change}, le
foncteur $F_{d, i}$ s'identifie au foncteur
associé à $U_i$ muni de la $\mathcal{O}_{U_i}$-algèbre graduée
quasi-cohérente $\mathcal{A}|_{U_i}$. La représentabilité résulte donc du
Lemme \ref{constructions-lemma-relative-proj-affine}.

\medskip\noindent
Montrons ensuite que $F_{d, i} \subset F_d$ est représentable par des
immersions ouvertes. Soit $(f : T \to S, \varphi) \in F_d(T)$ et posons
$V_i = f^{-1}(U_i)$. Par définition de $F_{d, i}$, pour un morphisme
$a : T' \to T$, on a $a^*(f, \varphi) \in F_{d, i}(T')$ si et seulement si
$a(T') \subset V_i$. C'est ce qu'il fallait montrer.

\medskip\noindent
Enfin, montrons que la famille $(F_{d, i})_{i \in I}$ recouvre $F_d$.
Soit $(f : T \to S, \varphi) \in F_d(T)$ et posons
$V_i = f^{-1}(U_i)$. Puisque $S = \bigcup_{i \in I} U_i$ est un
recouvrement ouvert de $S$, on obtient un recouvrement ouvert
$T = \bigcup_{i \in I} V_i$ de $T$. De plus, $(f, \varphi)|_{V_i} \in F_{d, i}(V_i)$.
Cela achève la démonstration du lemme.
\end{proof}

\noindent
Nous pouvons maintenant reprendre, dans le cadre relatif, la fin de la
section \ref{constructions-section-morphisms-proj}
et définir un foncteur représentable
par $\underline{\text{Proj}}_S(\mathcal{A})$.
Introduisons pour cela la notion d'équivalence entre deux quadruplets
$(d, f : T \to S, \mathcal{L}, \psi)$ et
$(d', f' : T \to S, \mathcal{L}', \psi')$, les entiers $d, d'$
pouvant être distincts. Ils sont dits {\it équivalents} si $f = f'$ et
s'il existe un isomorphisme
$\beta : \mathcal{L}^{\otimes d'} \to (\mathcal{L}')^{\otimes d}$
tel que
$\beta \circ \psi|_{f^*\mathcal{A}^{(dd')}} = \psi'|_{f^*\mathcal{A}^{(dd')}}$.
Le lemme suivant montre que l'on
obtient bien une relation d'équivalence
(ce qui n'est pas tout à fait immédiat).

\begin{lemma}
\label{constructions-lemma-equivalent-relative}
Dans la Situation \ref{constructions-situation-relative-proj}, soient $T$ un schéma et
$(d, f, \mathcal{L}, \psi)$, $(d', f', \mathcal{L}', \psi')$ deux
quadruplets sur $T$. Les assertions suivantes sont équivalentes :
\begin{enumerate}
\item Posons $m = \text{lcm}(d, d')$ et écrivons $m = ad = a'd'$.
On a $f = f'$ et il existe un isomorphisme
$\beta : \mathcal{L}^{\otimes a} \to (\mathcal{L}')^{\otimes a'}$
tel que $\beta \circ \psi|_{f^*\mathcal{A}^{(m)}}$ et
$\psi'|_{f^*\mathcal{A}^{(m)}}$
coïncident comme homomorphismes d'anneaux
gradués
$f^*\mathcal{A}^{(m)} \to \bigoplus_{n \geq 0} (\mathcal{L}')^{\otimes mn}$.
\item Les quadruplets $(d, f, \mathcal{L}, \psi)$ et
$(d', f', \mathcal{L}', \psi')$ sont équivalents.
\item On a $f = f'$ et, pour un certain entier strictement positif
$m = ad = a'd'$, il existe un isomorphisme
$\beta : \mathcal{L}^{\otimes a} \to (\mathcal{L}')^{\otimes a'}$
tel que $\beta \circ \psi|_{f^*\mathcal{A}^{(m)}}$ et
$\psi'|_{f^*\mathcal{A}^{(m)}}$
coïncident comme homomorphismes d'anneaux
gradués
$f^*\mathcal{A}^{(m)} \to \bigoplus_{n \geq 0} (\mathcal{L}')^{\otimes mn}$.
\end{enumerate}
\end{lemma}

\begin{proof}
Il est clair que (1) entraîne (2) et que
(2) entraîne (3), en se restreignant
à des degrés plus divisibles et en prenant des puissances des faisceaux
inversibles. Supposons (3) pour un entier $m = ad = a'd'$.
Posons $m_0 = \text{lcm}(d, d')$ et écrivons $m_0 = a_0d = a'_0d'$.
Nous disposons d'un isomorphisme
$\beta : \mathcal{L}^{\otimes a} \to (\mathcal{L}')^{\otimes a'}$
ayant la propriété (3), et cherchons un isomorphisme
$\beta_0 : \mathcal{L}^{\otimes a_0} \to (\mathcal{L}')^{\otimes a'_0}$
ayant la même propriété. Les morphismes
$\psi : f^*\mathcal{A}_d \to \mathcal{L}$ et
$\psi' : (f')^*\mathcal{A}_{d'} \to \mathcal{L}'$ étant surjectifs par
hypothèse, il en est de même des morphismes
$\psi : f^*\mathcal{A}_{m_0} \to \mathcal{L}^{\otimes a_0}$ et
$\psi' : (f')^*\mathcal{A}_{m_0} \to (\mathcal{L}')^{\otimes a_0}$.
Ainsi, s'il existe, $\beta_0$ est uniquement déterminé par la condition
$\beta_0 \circ \psi = \psi'$. On peut donc travailler localement sur $T$,
et supposer que $f = f' : T \to S$ est à valeurs dans un ouvert affine,
autrement dit que $S$ est affine. Le résultat découle alors de l'assertion
correspondante pour les triplets (Lemme \ref{constructions-lemma-equivalent}) et de la
correspondance entre triplets et quadruplets lorsque la base est affine
(voir la démonstration du Lemme \ref{constructions-lemma-relative-proj-affine}).
\end{proof}

\noindent
Supposons $d' = ad$. Considérons la transformation de foncteurs
$F_d \to F_{d'}$ qui associe au quadruplet $(d, f, \mathcal{L}, \psi)$ sur
$T$ le quadruplet
$(d', f, \mathcal{L}^{\otimes a}, \psi|_{f^*\mathcal{A}^{(d')}})$.
Une des implications du Lemme \ref{constructions-lemma-equivalent-relative} montre que
la transformation $F_d \to F_{d'}$ est
injective. Pour un schéma quasi-compact $T$,
définissons
$$
F(T) = \bigcup\nolimits_{d \in \mathbf{N}} F_d(T)
$$
avec les applications de transition ci-dessus. On obtient ainsi un foncteur
contravariant de la catégorie des schémas quasi-compacts vers celle des
ensembles. Pour un schéma quelconque $T$, définissons
$$
F(T)
=
\lim_{V \subset T\text{ quasi-compact open}} F(V).
$$
Autrement dit, un élément $\xi$ de $F(T)$ correspond à une famille
compatible d'éléments $\xi_V \in F(V)$, où $V$ parcourt les ouverts
quasi-compacts de $T$. Nous omettons
la définition de l'application d'image
inverse $F(T) \to F(T')$ associée à un morphisme de schémas $T' \to T$.
Nous avons ainsi défini le foncteur
\begin{equation}
\label{constructions-equation-proj}
F : \Sch^{opp} \longrightarrow \textit{Ensembles}
\end{equation}

\begin{lemma}
\label{constructions-lemma-relative-proj}
Dans la Situation \ref{constructions-situation-relative-proj}, le foncteur $F$ ci-dessus
est représentable par un schéma.
\end{lemma}

\begin{proof}
Soit $U_d \to S$ le schéma représentant le foncteur $F_d$ défini ci-dessus,
et soit $\mathcal{L}_d$,
$\psi^d : \pi_d^*\mathcal{A}^{(d)} \to
\bigoplus_{n \geq 0} \mathcal{L}_d^{\otimes n}$ l'objet universel.
Si $d | d'$, on peut considérer le quadruplet
$(d', \pi_d, \mathcal{L}_d^{\otimes d'/d}, \psi^d|_{\mathcal{A}^{(d')}})$,
qui détermine un morphisme canonique $U_d \to U_{d'}$ au-dessus de $S$.
Par construction, ce morphisme correspond à la transformation de foncteurs
$F_d \to F_{d'}$ définie ci-dessus.

\medskip\noindent
Pour tout ouvert affine $\Spec(R) = V \subset S$, en posant
$A = \Gamma(V, \mathcal{A})$, on dispose d'une identification canonique
du changement de base $U_{d, V}$ avec le sous-schéma ouvert correspondant
de $\text{Proj}(A)$ ; voir le Lemme \ref{constructions-lemma-relative-proj-affine}.
De plus, les morphismes $U_{d, V} \to U_{d', V}$ construits ci-dessus
correspondent aux inclusions d'ouverts dans $\text{Proj}(A)$.
On en déduit que $U_d \to U_{d'}$ est une immersion ouverte.

\medskip\noindent
On peut donc construire $X$ en recollant les schémas $U_d$ suivant les
immersions ouvertes $U_d \to U_{d'}$.
Il est commode de choisir une suite
$d_1 | d_2 | d_3 | \ldots$ telle que
tout entier strictement positif divise
l'un des $d_i$, puis de poser
$X = \bigcup U_{d_i}$ au moyen des immersions
ouvertes ci-dessus. Il est alors aisé de montrer que $X$ représente $F$.
\end{proof}

\begin{lemma}
\label{constructions-lemma-glueing-gives-functor-proj}
Dans la Situation \ref{constructions-situation-relative-proj}, le schéma
$\pi : \underline{\text{Proj}}_S(\mathcal{A}) \to S$ construit dans le
Lemme \ref{constructions-lemma-glue-relative-proj} est canoniquement
isomorphe au schéma représentant le foncteur $F$,
comme schéma sur $S$.
\end{lemma}

\begin{proof}
Soit $X$ le schéma représentant $F$.
$X$ est un schéma sur $S$, puisque $F$
est muni d'une transformation naturelle $F \to h_S$.
Posons $Y = \underline{\text{Proj}}_S(\mathcal{A})$.
Nous devons montrer
que $X \cong Y$ comme $S$-schémas. Donnons deux démonstrations.

\medskip\noindent
La première utilise la construction de $X$ comme réunion des schémas $U_d$
représentant $F_d$ dans la démonstration
du Lemme \ref{constructions-lemma-relative-proj}.
Au-dessus de chaque ouvert affine de $S$, on peut identifier $X$ au spectre
homogène des sections de $\mathcal{A}$
sur cet ouvert, puisque c'est le cas
pour les ouverts $U_d$. Ces identifications sont de plus compatibles avec
les restrictions à des ouverts affines plus petits. Or $Y$ a été construit
en recollant ces spectres homogènes. On
peut donc recoller ces isomorphismes
en un isomorphisme entre $X$ et $\underline{\text{Proj}}_S(\mathcal{A})$.
Nous omettons les détails.

\medskip\noindent
Voici la seconde démonstration. Le
Lemme \ref{constructions-lemma-glue-relative-proj-twists} fournit un morphisme d'algèbres
graduées
$$
\psi : \pi^*\mathcal{A}
\longrightarrow
\bigoplus\nolimits_{n \geq 0} \mathcal{O}_Y(n)
$$
sur $Y$ qui, sur les sections au-dessus des ouverts affines de $S$, coïncide
avec (\ref{constructions-equation-global-sections}). Pour tout $y \in Y$, il existe donc
un voisinage ouvert $V \subset Y$ de $y$ et un entier $d \geq 1$ tels que,
pour $d | n$, le faisceau $\mathcal{O}_Y(n)|_V$ soit inversible et que les
morphismes de multiplication
$\mathcal{O}_Y(n)|_V \otimes_{\mathcal{O}_V} \mathcal{O}_Y(m)|_V
\to \mathcal{O}_Y(n + m)|_V$ soient des isomorphismes.
La restriction de $\psi$ à $\pi^*\mathcal{A}^{(d)}|_V$ donne ainsi un
élément de $F_d(V)$. Les ouverts $V$ recouvrant $Y$, on voit que « $\psi$ »
définit un élément de $F(Y)$, donc un
morphisme canonique $Y \to X$ sur $S$.
La construction étant entièrement canonique, on peut vérifier localement
sur $S$ qu'il s'agit d'un isomorphisme. On se ramène donc au cas où $S$ est
affine, dans lequel le résultat est clair.
\end{proof}

\begin{definition}
\label{constructions-definition-relative-proj}
Soient $S$ un schéma et $\mathcal{A}$ un faisceau quasi-cohérent de
$\mathcal{O}_S$-algèbres graduées. Le {\it spectre homogène relatif de
$\mathcal{A}$ sur $S$}, ou {\it spectre
homogène de $\mathcal{A}$ sur $S$},
ou encore {\it Proj relatif de $\mathcal{A}$ sur $S$}, est le schéma
construit dans le Lemme \ref{constructions-lemma-glue-relative-proj} et représentant le
foncteur $F$ de (\ref{constructions-equation-proj}) ; voir le
Lemme \ref{constructions-lemma-glueing-gives-functor-proj}.
On le note $\pi : \underline{\text{Proj}}_S(\mathcal{A}) \to S$.
\end{definition}

\noindent
Le Proj relatif est muni d'un faisceau quasi-cohérent d'algèbres
$\mathbf{Z}$-graduées
$\bigoplus_{n \in \mathbf{Z}}
\mathcal{O}_{\underline{\text{Proj}}_S(\mathcal{A})}(n)$
(les faisceaux tordus du faisceau structural) et d'un homomorphisme
« universel » d'algèbres graduées
$$
\psi_{univ} :
\mathcal{A}
\longrightarrow
\pi_*\left(
\bigoplus\nolimits_{n \geq 0}
\mathcal{O}_{\underline{\text{Proj}}_S(\mathcal{A})}(n)
\right)
$$
voir le Lemme \ref{constructions-lemma-glue-relative-proj-twists}.
On peut également le
considérer comme un homomorphisme
$$
\psi_{univ} :
\pi^*\mathcal{A}
\longrightarrow
\bigoplus\nolimits_{n \geq 0}
\mathcal{O}_{\underline{\text{Proj}}_S(\mathcal{A})}(n)
$$
Le lemme suivant exprime l'universalité de cet objet.

\begin{lemma}
\label{constructions-lemma-tie-up-psi}
Dans la Situation \ref{constructions-situation-relative-proj}, soit
$(f : T \to S, d, \mathcal{L}, \psi)$
un quadruplet et soit
$r_{d, \mathcal{L}, \psi} : T \to \underline{\text{Proj}}_S(\mathcal{A})$
le $S$-morphisme associé. Il existe un isomorphisme
d'algèbres $\mathbf{Z}$-graduées sur $\mathcal{O}_T$
$$
\theta :
r_{d, \mathcal{L}, \psi}^*\left(
\bigoplus\nolimits_{n \in \mathbf{Z}}
\mathcal{O}_{\underline{\text{Proj}}_S(\mathcal{A})}(nd)
\right)
\longrightarrow
\bigoplus\nolimits_{n \in \mathbf{Z}} \mathcal{L}^{\otimes n}
$$
rendant commutatif le diagramme suivant
$$
\xymatrix{
\mathcal{A}^{(d)} \ar[rr]_-{\psi}
 \ar[rd]_-{\psi_{univ}} & &
f_*\left(
\bigoplus\nolimits_{n \in \mathbf{Z}}
\mathcal{L}^{\otimes n}
\right) \\
 &
\pi_*\left(
\bigoplus\nolimits_{n \geq 0}
\mathcal{O}_{\underline{\text{Proj}}_S(\mathcal{A})}(nd)
\right) \ar[ru]_\theta
}
$$
La commutativité de ce diagramme
détermine $\theta$ de façon unique.
\end{lemma}

\begin{proof}
Le quadruplet $(f : T \to S, d, \mathcal{L}, \psi)$ définit un élément de
$F_d(T)$. Soit
$U_d \subset \underline{\text{Proj}}_S(\mathcal{A})$ le lieu
où le faisceau
$\mathcal{O}_{\underline{\text{Proj}}_S(\mathcal{A})}(d)$ est
inversible et engendré par l'image de
$\psi_{univ} : \pi^*\mathcal{A}_d \to
\mathcal{O}_{\underline{\text{Proj}}_S(\mathcal{A})}(d)$.
Rappelons que $U_d$ représente $F_d$ ; voir la démonstration du
Lemme \ref{constructions-lemma-relative-proj}. Il
suffit donc de montrer que le quadruplet
$(U_d \to S, d, \mathcal{O}_{U_d}(d), \psi_{univ}|_{\mathcal{A}^{(d)}})$
est la famille universelle, c'est-à-dire l'objet représentant dans
$F_d(U_d)$. On peut le vérifier après restriction à un ouvert affine de
$S$, car (a) la formation des foncteurs
$F_d$ commute au changement de base
(Lemme \ref{constructions-lemma-proj-base-change}) et (b) le couple
$(\bigoplus_{n \in \mathbf{Z}}
\mathcal{O}_{\underline{\text{Proj}}_S(\mathcal{A})}(n),
\psi_{univ})$
est construit par recollement au-dessus des ouverts affines de $S$
(Lemme \ref{constructions-lemma-glue-relative-proj-twists}).
On peut donc supposer $S$ affine. Alors le foncteur $F_d$ des quadruplets
coïncide avec le foncteur $F_d$ des triplets (voir la démonstration du
Lemme \ref{constructions-lemma-relative-proj-affine}), et le
Lemme \ref{constructions-lemma-proj-functor-strict} montre que
$(d, \mathcal{O}_{U_d}(d), \psi^d)$ est le triplet universel sur $U_d$.
En remontant les identifications de la démonstration du
Lemme \ref{constructions-lemma-relative-proj-affine}, on voit que
$(U_d \to S, d, \mathcal{O}_{U_d}(d), \psi_{univ}|_{\mathcal{A}^{(d)}})$
est bien le quadruplet universel cherché.
\end{proof}

\begin{lemma}
\label{constructions-lemma-relative-proj-separated}
Soient $S$ un schéma et $\mathcal{A}$
un faisceau quasi-cohérent de
$\mathcal{O}_S$-algèbres graduées. Le morphisme
$\pi : \underline{\text{Proj}}_S(\mathcal{A}) \to S$ est séparé.
\end{lemma}

\begin{proof}
Pour vérifier qu'un morphisme est séparé, on peut travailler localement sur
la base ; voir Schémas, section \ref{schemes-section-separation-axioms}.
Par construction, $\underline{\text{Proj}}_S(\mathcal{A})$ est isomorphe,
au-dessus de chaque ouvert affine $U \subset S$, à $\text{Proj}(A)$, où
$A = \mathcal{A}(U)$. Par le Lemme \ref{constructions-lemma-proj-separated},
$\text{Proj}(A)$ est séparé. Le morphisme $\text{Proj}(A) \to U$ est donc
séparé (voir Schémas, Lemme \ref{schemes-lemma-compose-after-separated}),
comme voulu.
\end{proof}

\begin{lemma}
\label{constructions-lemma-relative-proj-base-change}
Soient $S$ un schéma et $\mathcal{A}$ un faisceau quasi-cohérent de
$\mathcal{O}_S$-algèbres graduées. Pour tout morphisme de schémas
$g : S' \to S$, il existe un isomorphisme canonique
$$
r :
\underline{\text{Proj}}_{S'}(g^*\mathcal{A})
\longrightarrow
S' \times_S \underline{\text{Proj}}_S(\mathcal{A})
$$
ainsi qu'un isomorphisme correspondant
$$
\theta :
r^*\text{pr}_2^*\left(\bigoplus\nolimits_{d \in \mathbf{Z}}
\mathcal{O}_{\underline{\text{Proj}}_S(\mathcal{A})}(d)\right)
\longrightarrow
\bigoplus\nolimits_{d \in \mathbf{Z}}
\mathcal{O}_{\underline{\text{Proj}}_{S'}(g^*\mathcal{A})}(d)
$$
d'algèbres $\mathbf{Z}$-graduées sur
$\mathcal{O}_{\underline{\text{Proj}}_{S'}(g^*\mathcal{A})}$.
\end{lemma}

\begin{proof}
Cela résulte du Lemme \ref{constructions-lemma-proj-base-change} et de la construction
de $\underline{\text{Proj}}_S(\mathcal{A})$, dans le
Lemme \ref{constructions-lemma-relative-proj}, comme
réunion des schémas $U_d$ représentant
les foncteurs $F_d$. Du point de vue de la construction par recollement,
cet isomorphisme est donné par les isomorphismes du
Lemme \ref{constructions-lemma-base-change-map-proj}, qui fournit aussi $\theta$.
Nous omettons certains détails.
\end{proof}

\begin{lemma}
\label{constructions-lemma-apply-relative}
Soient $S$ un schéma et $\mathcal{A}$ un faisceau quasi-cohérent de
$\mathcal{O}_S$-modules gradués, engendré comme $\mathcal{A}_0$-algèbre par
$\mathcal{A}_1$. Alors $X = \underline{\text{Proj}}_S(\mathcal{A})$
représente le foncteur $F_1$ qui associe à un schéma $f : T \to S$ sur $S$
l'ensemble des couples $(\mathcal{L}, \psi)$ où
\begin{enumerate}
\item $\mathcal{L}$ est un $\mathcal{O}_T$-module inversible ;
\item
$\psi : f^*\mathcal{A} \to \bigoplus_{n \geq 0} \mathcal{L}^{\otimes n}$
est un homomorphisme de $\mathcal{O}_T$-algèbres graduées tel que
$f^*\mathcal{A}_1 \to \mathcal{L}$ soit surjectif,
\end{enumerate}
à équivalence stricte près au sens
ci-dessus. De plus, dans ce cas, tous les
faisceaux quasi-cohérents
$\mathcal{O}_{\underline{\text{Proj}}(\mathcal{A})}(n)$ sont des
$\mathcal{O}_{\underline{\text{Proj}}(\mathcal{A})}$-modules inversibles,
et les morphismes de multiplication induisent des isomorphismes
$
\mathcal{O}_{\underline{\text{Proj}}(\mathcal{A})}(n)
\otimes_{\mathcal{O}_{\underline{\text{Proj}}(\mathcal{A})}}
\mathcal{O}_{\underline{\text{Proj}}(\mathcal{A})}(m) =
\mathcal{O}_{\underline{\text{Proj}}(\mathcal{A})}(n + m)$.
\end{lemma}

\begin{proof}
Sous les hypothèses du lemme, les faisceaux
$\mathcal{O}_{\underline{\text{Proj}}(\mathcal{A})}(n)$ sont inversibles
et les morphismes de multiplication sont des isomorphismes, par le
Lemme \ref{constructions-lemma-relative-proj} et le Lemme \ref{constructions-lemma-apply} appliqués
au-dessus des ouverts affines de $S$.
Ainsi, $X$ représente bien $F_1$ ;
voir la démonstration du Lemme \ref{constructions-lemma-relative-proj}.
\end{proof}












\section{Faisceaux quasi-cohérents sur le Proj relatif}
\label{constructions-section-quasi-coherent-relative-proj}

\noindent
Nous indiquons brièvement comment traiter les modules gradués dans le cadre
relatif.

\medskip\noindent
Plaçons-nous dans la Situation \ref{constructions-situation-relative-proj} : $S$ est un
schéma et $\mathcal{A}$ une
$\mathcal{O}_S$-algèbre graduée quasi-cohérente.
Soit $\mathcal{M} = \bigoplus_{n \in \mathbf{Z}} \mathcal{M}_n$ un
$\mathcal{A}$-module gradué, quasi-cohérent comme $\mathcal{O}_S$-module.
Nous allons décrire le faisceau quasi-cohérent de modules associé sur
$\underline{\text{Proj}}_S(\mathcal{A})$.
Nous décrivons d'abord son image
inverse sur les schémas $T$ munis d'un morphisme vers le Proj relatif.

\medskip\noindent
Soient $T$ un schéma et $(d, f : T \to S, \mathcal{L}, \psi)$ un quadruplet
sur $T$, au sens de la section \ref{constructions-section-relative-proj}.
Définissons un faisceau quasi-cohérent $\widetilde{\mathcal{M}}_T$ de
$\mathcal{O}_T$-modules par
\begin{equation}
\label{constructions-equation-widetilde-M}
\widetilde{\mathcal{M}}_T =
\left(
f^*\mathcal{M}^{(d)}
\otimes_{f^*\mathcal{A}^{(d)}}
\left(\bigoplus\nolimits_{n \in \mathbf{Z}} \mathcal{L}^{\otimes n}\right)
\right)_0
\end{equation}
Ainsi, $\widetilde{\mathcal{M}}_T$ est
la composante de degré $0$ du produit
tensoriel des $f^*\mathcal{A}^{(d)}$-modules gradués
$\mathcal{M}^{(d)}$ et
$\bigoplus\nolimits_{n \in \mathbf{Z}} \mathcal{L}^{\otimes n}$.
Le faisceau $\widetilde{\mathcal{M}}_T$ dépend du
quadruplet, bien que la notation ne l'indique pas.
Cette construction a la propriété agréable suivante : pour tout morphisme
$g : T' \to T$, on a
$\widetilde{\mathcal{M}}_{T'} = g^*\widetilde{\mathcal{M}}_T$,
où $\widetilde{\mathcal{M}}_{T'}$ désigne
le faisceau quasi-cohérent associé
au quadruplet image inverse $(d, f \circ g, g^*\mathcal{L}, g^*\psi)$.

\medskip\noindent
Tous les faisceaux de (\ref{constructions-equation-widetilde-M}) étant quasi-cohérents,
on peut expliciter la construction sur un ouvert affine
$\Spec(C) = V \subset T$ dont l'image est contenue dans un ouvert affine
$\Spec(R) = U \subset S$. Supposons que $\mathcal{A}|_U$ corresponde à la
$R$-algèbre graduée $A$, que $\mathcal{M}|_U$ corresponde au $A$-module
gradué $M$ et que $\mathcal{L}|_V$
corresponde au $C$-module inversible $L$.
Le morphisme $\psi$ induit un homomorphisme de $R$-algèbres graduées
$\gamma : A^{(d)} \to \bigoplus_{n \geq 0} L^{\otimes n}$
(les puissances tensorielles de $L$ sont prises sur $C$).
Alors $(\widetilde{\mathcal{M}}_T)|_V$ est le faisceau quasi-cohérent
associé au $C$-module
$$
N_{R, C, A, M, \gamma} =
\left(
M^{(d)} \otimes_{A^{(d)}, \gamma}
\left(\bigoplus\nolimits_{n \in \mathbf{Z}} L^{\otimes n}\right)
\right)_0
$$
Par hypothèse, on peut même recouvrir $T$ par des ouverts affines $V$ tels
qu'il existe $a \in A_d$ dont l'image $\gamma(a) \in L$ soit une base du
$C$-module $L$. Dans ce cas, tout élément de $N_{R, C, A, M, \gamma}$ est
une somme de tenseurs purs $\sum m_i \otimes \gamma(a)^{-n_i}$, avec
$m_i \in M_{n_id}$. En multipliant
chaque $m_i$ par une puissance positive
convenable de $a$, puis en regroupant
les termes, on voit que tout élément
de $N_{R, C, A, M, \gamma}$ s'écrit $m \otimes \gamma(a)^{-n}$ avec
$m \in M_{nd}$ et $n \gg 0$. Autrement dit, dans ce cas,
$$
N_{R, C, A, M, \gamma} = M_{(a)} \otimes_{A_{(a)}} C
$$
où $A_{(a)} \to C$ est donné par
$x/a^n \mapsto \gamma(x)/\gamma(a)^n$.
C'est donc la valeur de $\widetilde{M}$ sur
$D_{+}(a) \subset \text{Proj}(A)$, tirée en arrière sur $\Spec(C)$ par le
morphisme $\Spec(C) \to D_{+}(a)$ défini par $\gamma$.

\begin{lemma}
\label{constructions-lemma-relative-proj-modules}
Dans la Situation \ref{constructions-situation-relative-proj}, à tout faisceau
quasi-cohérent de $\mathcal{A}$-modules gradués $\mathcal{M}$ sur $S$ est
canoniquement associé un faisceau de
$\mathcal{O}_{\underline{\text{Proj}}_S(\mathcal{A})}$-modules
$\widetilde{\mathcal{M}}$ ayant les
propriétés suivantes :
\begin{enumerate}
\item Pour tout schéma $T$ et tout quadruplet
$(T \to S, d, \mathcal{L}, \psi)$ sur $T$ correspondant à un morphisme
$h : T \to \underline{\text{Proj}}_S(\mathcal{A})$, il existe un
isomorphisme canonique
$\widetilde{\mathcal{M}}_T = h^*\widetilde{\mathcal{M}}$,
où $\widetilde{\mathcal{M}}_T$ est défini par (\ref{constructions-equation-widetilde-M}).
\item Les isomorphismes de (1) sont compatibles avec les images inverses.
\item Il existe un morphisme canonique
$$
\pi^*\mathcal{M}_0 \longrightarrow \widetilde{\mathcal{M}}.
$$
\item La construction $\mathcal{M} \mapsto \widetilde{\mathcal{M}}$ est
fonctorielle en $\mathcal{M}$.
\item La construction $\mathcal{M} \mapsto \widetilde{\mathcal{M}}$ est
exacte.
\item Il existe des morphismes canoniques
$$
\widetilde{\mathcal{M}}
\otimes_{\mathcal{O}_{\underline{\text{Proj}}_S(\mathcal{A})}}
\widetilde{\mathcal{N}}
\longrightarrow
\widetilde{\mathcal{M} \otimes_\mathcal{A} \mathcal{N}}
$$
comme dans le Lemme \ref{constructions-lemma-widetilde-tensor}.
\item Il existe des morphismes canoniques
$$
\pi^*\mathcal{M}
\longrightarrow
\bigoplus\nolimits_{n \in \mathbf{Z}}
\widetilde{\mathcal{M}(n)}
$$
généralisant (\ref{constructions-equation-global-sections-more-generally}).
\item La formation de $\widetilde{\mathcal{M}}$
commute au changement de base.
\end{enumerate}
\end{lemma}

\begin{proof}
La démonstration est omise. Il conviendrait de scinder ce lemme en plusieurs énoncés et de les démontrer séparément.
\end{proof}












\section{Fonctorialité du Proj relatif}
\label{constructions-section-functoriality-relative-proj}

\noindent
Cette section est l'analogue de la section \ref{constructions-section-functoriality-proj}
pour le Proj relatif. Soit $S$ un schéma. Un morphisme de
$\mathcal{O}_S$-algèbres graduées $\psi : \mathcal{A} \to \mathcal{B}$
n'induit pas toujours un morphisme entre les Proj relatifs associés. L'énoncé correct est le suivant.

\begin{lemma}
\label{constructions-lemma-morphism-relative-proj}
Soient $S$ un schéma et $\mathcal{A}$, $\mathcal{B}$ deux
$\mathcal{O}_S$-algèbres graduées quasi-cohérentes. Posons
$p : X = \underline{\text{Proj}}_S(\mathcal{A}) \to S$ et
$q : Y = \underline{\text{Proj}}_S(\mathcal{B}) \to S$.
Soit $\psi : \mathcal{A} \to \mathcal{B}$ un homomorphisme de
$\mathcal{O}_S$-algèbres graduées. Il existe un ouvert canonique
$U(\psi) \subset Y$ et un morphisme canonique de schémas
$$
r_\psi :
U(\psi)
\longrightarrow
X
$$ sur $S$
et un morphisme d'algèbres $\mathbf{Z}$-graduées sur
$\mathcal{O}_{U(\psi)}$
$$
\theta = \theta_\psi :
r_\psi^*\left(
\bigoplus\nolimits_{d \in \mathbf{Z}} \mathcal{O}_X(d)
\right)
\longrightarrow
\bigoplus\nolimits_{d \in \mathbf{Z}} \mathcal{O}_{U(\psi)}(d).
$$
Le triplet $(U(\psi), r_\psi, \theta)$ est caractérisé par la propriété
suivante : pour tout ouvert affine $W \subset S$, le triplet
$$
(U(\psi) \cap p^{-1}W,\quad
r_\psi|_{U(\psi) \cap p^{-1}W} : U(\psi) \cap p^{-1}W \to q^{-1}W,\quad
\theta|_{U(\psi) \cap p^{-1}W})
$$ est égal au triplet associé à
$\psi : \mathcal{A}(W) \to \mathcal{B}(W)$ dans le
Lemme \ref{constructions-lemma-morphism-proj}, via les identifications
$p^{-1}W = \text{Proj}(\mathcal{A}(W))$ et
$q^{-1}W = \text{Proj}(\mathcal{B}(W))$ de la
section \ref{constructions-section-relative-proj-via-glueing}.
\end{lemma}

\begin{proof}
Le lemme s'obtient en recollant les triplets locaux.
\end{proof}

\begin{lemma}
\label{constructions-lemma-morphism-relative-proj-transitive}
Soient $S$ un schéma et $\mathcal{A}$, $\mathcal{B}$, $\mathcal{C}$ des
$\mathcal{O}_S$-algèbres graduées quasi-cohérentes. Posons
$X = \underline{\text{Proj}}_S(\mathcal{A})$,
$Y = \underline{\text{Proj}}_S(\mathcal{B})$ et
$Z = \underline{\text{Proj}}_S(\mathcal{C})$.
Soient $\varphi : \mathcal{A} \to \mathcal{B}$ et
$\psi : \mathcal{B} \to \mathcal{C}$ des morphismes de
$\mathcal{O}_S$-algèbres graduées. On a
$$
U(\psi \circ \varphi) = r_\varphi^{-1}(U(\psi))
\quad
\text{et}
\quad
r_{\psi \circ \varphi}
=
r_\varphi \circ r_\psi|_{U(\psi \circ \varphi)}.
$$
De plus,
$$
\theta_\psi \circ r_\psi^*\theta_\varphi
=
\theta_{\psi \circ \varphi}
$$
avec les notations évidentes.
\end{lemma}

\begin{proof}
La démonstration est omise.
\end{proof}

\begin{lemma}
\label{constructions-lemma-surjective-graded-rings-map-relative-proj}
Avec les hypothèses et notations du Lemme \ref{constructions-lemma-morphism-relative-proj},
supposons que $\mathcal{A}_d \to \mathcal{B}_d$ soit surjectif pour
$d \gg 0$. Alors
\begin{enumerate}
\item $U(\psi) = Y$ ;
\item $r_\psi : Y \to X$ est une immersion fermée ;
\item les morphismes $\theta : r_\psi^*\mathcal{O}_X(n) \to \mathcal{O}_Y(n)$
sont surjectifs, mais ne sont pas des isomorphismes en général, même si
$\mathcal{A} \to \mathcal{B}$ est surjectif.
\end{enumerate}
\end{lemma}

\begin{proof}
Cela résulte des Lemmes
\ref{constructions-lemma-morphism-relative-proj}
et
\ref{constructions-lemma-surjective-graded-rings-map-proj}.
\end{proof}

\begin{lemma}
\label{constructions-lemma-eventual-iso-graded-rings-map-relative-proj}
Avec les hypothèses et notations du Lemme \ref{constructions-lemma-morphism-relative-proj},
supposons que $\mathcal{A}_d \to \mathcal{B}_d$ soit un isomorphisme pour
tout $d \gg 0$. Alors
\begin{enumerate}
\item $U(\psi) = Y$ ;
\item $r_\psi : Y \to X$ est un isomorphisme ;
\item les morphismes $\theta : r_\psi^*\mathcal{O}_X(n) \to \mathcal{O}_Y(n)$
sont des isomorphismes.
\end{enumerate}
\end{lemma}

\begin{proof}
Cela résulte des Lemmes
\ref{constructions-lemma-morphism-relative-proj}
et
\ref{constructions-lemma-eventual-iso-graded-rings-map-proj}.
\end{proof}

\begin{lemma}
\label{constructions-lemma-surjective-generated-degree-1-map-relative-proj}
Avec les hypothèses et notations du Lemme \ref{constructions-lemma-morphism-relative-proj},
supposons que $\mathcal{A}_d \to \mathcal{B}_d$ soit surjectif pour
$d \gg 0$ et que $\mathcal{A}$ soit engendrée par $\mathcal{A}_1$ sur
$\mathcal{A}_0$. Alors
\begin{enumerate}
\item $U(\psi) = Y$ ;
\item $r_\psi : Y \to X$ est une immersion fermée ;
\item les morphismes $\theta : r_\psi^*\mathcal{O}_X(n) \to \mathcal{O}_Y(n)$
sont des isomorphismes.
\end{enumerate}
\end{lemma}

\begin{proof}
Cela résulte des Lemmes
\ref{constructions-lemma-morphism-relative-proj}
et
\ref{constructions-lemma-surjective-graded-rings-generated-degree-1-map-proj}.
\end{proof}

















\section{Faisceaux inversibles et morphismes vers le Proj relatif}
\label{constructions-section-invertible-relative-proj}

\noindent
Le lemme suivant pourra nous être utile. Voici les données considérées :
\begin{enumerate}
\item $S$ est un schéma.
\item $\mathcal{A}$ est une $\mathcal{O}_S$-algèbre graduée quasi-cohérente.
\item On note
$\pi : \underline{\text{Proj}}_S(\mathcal{A}) \to S$ le spectre
homogène relatif sur $S$.
\item $f : X \to S$ est un morphisme de schémas.
\item $\mathcal{L}$ est un $\mathcal{O}_X$-module inversible.
\item $\psi : f^*\mathcal{A} \to
\bigoplus_{d \geq 0} \mathcal{L}^{\otimes d}$ est un homomorphisme de
$\mathcal{O}_X$-algèbres graduées.
\end{enumerate}
Ces données étant fixées, posons
$$
U(\psi) = \bigcup\nolimits_{(U, V, a)} U_{\psi(a)}
$$
où $(U, V, a)$ vérifie les conditions suivantes :
\begin{enumerate}
\item $V \subset S$ est un ouvert affine ;
\item $U = f^{-1}(V)$ ;
\item $a \in \mathcal{A}(V)_{+}$ est homogène.
\end{enumerate}
On a alors $\psi(a) \in \Gamma(U, \mathcal{L}^{\otimes \deg(a)})$, et
$U_{\psi(a)}$ est l'ouvert correspondant (voir Modules,
Lemme \ref{modules-lemma-s-open}).

\begin{lemma}
\label{constructions-lemma-invertible-map-into-relative-proj}
Avec les hypothèses et notations ci-dessus, le morphisme $\psi$ induit un
morphisme canonique de schémas sur $S$
$$
r_{\mathcal{L}, \psi} :
U(\psi) \longrightarrow \underline{\text{Proj}}_S(\mathcal{A})
$$
muni d'un morphisme de $\mathcal{O}_{U(\psi)}$-algèbres graduées
$$
\theta :
r_{\mathcal{L}, \psi}^*\left(
\bigoplus\nolimits_{d \geq 0}
\mathcal{O}_{\underline{\text{Proj}}_S(\mathcal{A})}(d)
\right)
\longrightarrow
\bigoplus\nolimits_{d \geq 0} \mathcal{L}^{\otimes d}|_{U(\psi)}
$$
caractérisés par les propriétés suivantes :
\begin{enumerate}
\item Pour tout ouvert $V \subset S$ et tout $d \geq 0$, le diagramme
$$
\xymatrix{
\mathcal{A}_d(V) \ar[d]_{\psi} \ar[r]_{\psi} &
\Gamma(f^{-1}(V), \mathcal{L}^{\otimes d}) \ar[d]^{restrict} \\
\Gamma(\pi^{-1}(V),
\mathcal{O}_{\underline{\text{Proj}}_S(\mathcal{A})}(d)) \ar[r]^{\theta} &
\Gamma(f^{-1}(V) \cap U(\psi), \mathcal{L}^{\otimes d})
}
$$
est commutatif.
\item Pour tout $d \geq 1$ et tout sous-schéma ouvert $W \subset X$ tels que
$\psi|_W : f^*\mathcal{A}_d|_W \to \mathcal{L}^{\otimes d}|_W$ soit
surjectif, la restriction du morphisme $r_{\mathcal{L}, \psi}$ coïncide avec
le morphisme $W \to \underline{\text{Proj}}_S(\mathcal{A})$ fourni par la
construction du spectre homogène relatif ; voir la
Définition \ref{constructions-definition-relative-proj}.
\item Pour tout ouvert affine $V \subset S$, la restriction
$$
(U(\psi) \cap f^{-1}(V), r_{\mathcal{L}, \psi}|_{U(\psi) \cap f^{-1}(V)},
\theta|_{U(\psi) \cap f^{-1}(V)})
$$
coïncide, via $i_V$ (voir le Lemme \ref{constructions-lemma-glue-relative-proj}), avec le
triplet $(U(\psi'), r_{\mathcal{L}, \psi'}, \theta')$ du
Lemme \ref{constructions-lemma-invertible-map-into-proj} associé à l'homomorphisme
$\psi' : A = \mathcal{A}(V) \to \Gamma_*(f^{-1}(V), \mathcal{L}|_{f^{-1}(V)})$
induit par $\psi$.
\end{enumerate}
\end{lemma}

\begin{proof}
On utilise la caractérisation (3) pour construire $r_{\mathcal{L}, \psi}$
et $\theta$ localement sur $S$, puis l'unicité du
Lemme \ref{constructions-lemma-invertible-map-into-proj} pour montrer que ces constructions
se recollent. Nous omettons les détails.
\end{proof}







\section{Torsion par des faisceaux inversibles et Proj relatif}
\label{constructions-section-twisting-and-proj}

\noindent
Soient $S$ un schéma,
$\mathcal{A} = \bigoplus_{d \geq 0} \mathcal{A}_d$ une
$\mathcal{O}_S$-algèbre graduée
quasi-cohérente et $\mathcal{L}$ un faisceau
inversible sur $S$. On obtient une autre $\mathcal{O}_S$-algèbre graduée
quasi-cohérente en posant
$$
\mathcal{B}
=
\bigoplus\nolimits_{d \geq 0}
\mathcal{A}_d \otimes_{\mathcal{O}_S} \mathcal{L}^{\otimes d}
$$
Les spectres homogènes relatifs de $\mathcal{A}$ et de $\mathcal{B}$ sont
isomorphes.

\begin{lemma}
\label{constructions-lemma-twisting-and-proj}
Avec les notations $S$, $\mathcal{A}$,
$\mathcal{L}$ et $\mathcal{B}$
ci-dessus, il existe un isomorphisme canonique
$$
\xymatrix{
P = \underline{\text{Proj}}_S(\mathcal{A})
\ar[rr]_g \ar[rd]_\pi & &
\underline{\text{Proj}}_S(\mathcal{B}) = P'
\ar[ld]^{\pi'} \\
& S &
}
$$
ayant les propriétés suivantes :
\begin{enumerate}
\item Il existe des isomorphismes
$\theta_n : g^*\mathcal{O}_{P'}(n)
\to
\mathcal{O}_P(n) \otimes \pi^*\mathcal{L}^{\otimes n}$
qui forment un isomorphisme d'algèbres $\mathbf{Z}$-graduées
$$
\theta :
g^*\left(
\bigoplus\nolimits_{n \in \mathbf{Z}} \mathcal{O}_{P'}(n)
\right)
\longrightarrow
\bigoplus\nolimits_{n \in \mathbf{Z}} \mathcal{O}_P(n)
\otimes \pi^*\mathcal{L}^{\otimes n}
$$
\item Pour tout ouvert $V \subset S$, les diagrammes
$$
\xymatrix{
\mathcal{A}_n(V) \otimes \mathcal{L}^{\otimes n}(V)
\ar[r]_{multiply} \ar[d]^{\psi \otimes \pi^*}
&
\mathcal{B}_n(V) \ar[dd]^\psi \\
\Gamma(\pi^{-1}V, \mathcal{O}_P(n)) \otimes
\Gamma(\pi^{-1}V, \pi^*\mathcal{L}^{\otimes n})
\ar[d]^{multiply} \\
\Gamma(\pi^{-1}V, \mathcal{O}_P(n) \otimes \pi^*\mathcal{L}^{\otimes n})
&
\Gamma(\pi'^{-1}V, \mathcal{O}_{P'}(n)) \ar[l]_-{\theta_n}
}
$$
sont commutatifs.
\item Ajouter d'autres propriétés ici au besoin.
\end{enumerate}
\end{lemma}

\begin{proof}
Lorsque $\mathcal{L} \cong \mathcal{O}_S$, il s'agit de l'identité.
Dans le cas général, on choisit un recouvrement ouvert de $S$ sur les
éléments duquel $\mathcal{L}$ se trivialise, puis on recolle les morphismes
correspondants. Nous omettons les détails.
\end{proof}
















\section{Fibrés projectifs}
\label{constructions-section-projective-bundle}

\noindent
Soient $S$ un schéma et $\mathcal{E}$ un faisceau quasi-cohérent de
$\mathcal{O}_S$-modules. Par Modules,
Lemme \ref{modules-lemma-whole-tensor-algebra-permanence}, l'algèbre
symétrique $\text{Sym}(\mathcal{E})$
de $\mathcal{E}$ sur $\mathcal{O}_S$
est un faisceau quasi-cohérent de $\mathcal{O}_S$-algèbres. Elle est
engendrée en degré $1$ sur $\mathcal{O}_S$.
On peut donc lui appliquer la
construction de la section précédente, notamment les Lemmes
\ref{constructions-lemma-relative-proj} et \ref{constructions-lemma-apply-relative}.

\begin{definition}
\label{constructions-definition-projective-bundle}
Soient $S$ un schéma et $\mathcal{E}$ un $\mathcal{O}_S$-module
quasi-cohérent\footnote{Le lecteur pourrait s'attendre à la condition que
$\mathcal{E}$ soit localement libre de rang fini. Nous ne l'imposons pas,
afin de rester cohérents avec \cite[II, définition 4.1.1]{EGA}.}.
On note
$$
\pi :
\mathbf{P}(\mathcal{E}) = \underline{\text{Proj}}_S(\text{Sym}(\mathcal{E}))
\longrightarrow
S
$$
et on l'appelle le {\it fibré projectif associé à $\mathcal{E}$}.
Le symbole $\mathcal{O}_{\mathbf{P}(\mathcal{E})}(n)$ désigne le
$\mathcal{O}_{\mathbf{P}(\mathcal{E})}$-module inversible du
Lemme \ref{constructions-lemma-apply-relative}, appelé
le {\it faisceau structural tordu}
d'indice $n$.
\end{definition}

\noindent
D'après le Lemme \ref{constructions-lemma-glue-relative-proj-twists}, il existe des
homomorphismes canoniques de $\mathcal{O}_S$-modules
$$
\text{Sym}^n(\mathcal{E})
\longrightarrow
\pi_*\mathcal{O}_{\mathbf{P}(\mathcal{E})}(n)
\quad\text{equivalently}\quad
\pi^*\text{Sym}^n(\mathcal{E})
\longrightarrow
\mathcal{O}_{\mathbf{P}(\mathcal{E})}(n)
$$
pour tout $n \geq 0$. En particulier, pour $n = 1$, on a
$$
\mathcal{E}
\longrightarrow
\pi_*\mathcal{O}_{\mathbf{P}(\mathcal{E})}(1)
\quad\text{equivalently}\quad
\pi^*\mathcal{E}
\longrightarrow
\mathcal{O}_{\mathbf{P}(\mathcal{E})}(1)
$$
et $\pi^*\mathcal{E} \to \mathcal{O}_{\mathbf{P}(\mathcal{E})}(1)$ est
surjectif par le Lemme \ref{constructions-lemma-apply-relative}. C'est une bonne façon
de retenir la convention adoptée pour
construire $\mathbf{P}(\mathcal{E})$.

\medskip\noindent
Attention : certains ouvrages ne définissent $\mathbf{P}(\mathcal{E})$
que lorsque $\mathcal{E}$ est localement libre de rang fini sur $S$.
De plus, ils notent parfois $\mathbf{P}(\mathcal{E})$ ce que nous notons
$\mathbf{P}(\mathcal{E}^\vee)$, où $\mathcal{E}^\vee$ est le dual de
$\mathcal{E}$ (là encore, seulement lorsque $\mathcal{E}$ est localement libre de rang
fini).

\medskip\noindent
Soient $S$, $\mathcal{E}$ et $\mathbf{P}(\mathcal{E}) \to S$ comme dans la
Définition \ref{constructions-definition-projective-bundle}. Soit $f : T \to S$ un
schéma sur $S$ et soit $\psi : f^*\mathcal{E} \to \mathcal{L}$ une
surjection, où $\mathcal{L}$ est un $\mathcal{O}_T$-module inversible.
Le morphisme induit de $\mathcal{O}_T$-algèbres graduées
$$
f^*\text{Sym}(\mathcal{E}) = \text{Sym}(f^*\mathcal{E}) \to
\text{Sym}(\mathcal{L}) = \bigoplus\nolimits_{n \geq 0} \mathcal{L}^{\otimes n}
$$ correspond à un morphisme
$$
\varphi_{\mathcal{L}, \psi} : T \longrightarrow \mathbf{P}(\mathcal{E})
$$
sur $S$, d'après la construction du Proj relatif comme schéma représentant
le foncteur $F$ dans la section \ref{constructions-section-relative-proj}.
Réciproquement, pour un morphisme $\varphi : T \to \mathbf{P}(\mathcal{E})$
sur $S$, on pose
$\mathcal{L} = \varphi^*\mathcal{O}_{\mathbf{P}(\mathcal{E})}(1)$ et on
prend pour $\psi : f^*\mathcal{E} \to \mathcal{L}$ l'image inverse par
$\varphi$ de la surjection canonique
$\pi^*\mathcal{E} \to \mathcal{O}_{\mathbf{P}(\mathcal{E})}(1)$.
Par le Lemme \ref{constructions-lemma-apply-relative}, ces constructions sont des
bijections réciproques entre les classes d'isomorphisme de couples
$(\mathcal{L}, \psi)$ et les morphismes
$\varphi : T \to \mathbf{P}(\mathcal{E})$ sur $S$.
Ainsi, $\mathbf{P}(\mathcal{E})$ représente le foncteur qui associe à
$f : T \to S$ l'ensemble des quotients comme $\mathcal{O}_T$-modules de
$f^*\mathcal{E}$ qui sont localement libres de rang $1$.

\begin{example}[Espace projectif d'un espace vectoriel]
\label{constructions-example-projective-space}
Soient $k$ un corps et $V$ un $k$-espace
vectoriel. L'{\it espace projectif}
correspondant est le $k$-schéma
$$
\mathbf{P}(V) = \text{Proj}(\text{Sym}(V))
$$
où $\text{Sym}(V)$ est l'algèbre symétrique de $V$ sur $k$.
On a bien sûr $\mathbf{P}(V) \cong \mathbf{P}^n_k$ si $\dim(V) = n + 1$,
puisque l'algèbre symétrique de $V$ est alors isomorphe à un anneau de
polynômes en $n + 1$ variables. Si l'on considère $V$ comme un module
quasi-cohérent sur $\Spec(k)$, alors $\mathbf{P}(V)$ est le fibré projectif
correspondant sur $\Spec(k)$. D'après ce qui précède, un
$k$-point $p$ de $\mathbf{P}(V)$ correspond à une surjection de
$k$-espaces vectoriels $V \to L_p$, avec $\dim(L_p) = 1$.
Plus généralement, soient $X$ un schéma sur $k$, $\mathcal{L}$ un
$\mathcal{O}_X$-module inversible et
$\psi : V \to \Gamma(X, \mathcal{L})$ une application $k$-linéaire telle
que $\mathcal{L}$ soit engendré, comme $\mathcal{O}_X$-module, par les
sections appartenant à l'image de $\psi$.
On obtient un morphisme canonique
$$
\varphi_{\mathcal{L}, \psi} : X \longrightarrow \mathbf{P}(V)
$$
de schémas sur $k$, muni d'un isomorphisme
$\theta : \varphi_{\mathcal{L}, \psi}^*\mathcal{O}_{\mathbf{P}(V)}(1)
\to \mathcal{L}$ tel que $\psi$ coïncide avec la composée
$$
V \to
\Gamma(\mathbf{P}(V), \mathcal{O}_{\mathbf{P}(V)}(1))
\to
\Gamma(X, \varphi_{\mathcal{L}, \psi}^*\mathcal{O}_{\mathbf{P}(V)}(1))
\to
\Gamma(X, \mathcal{L})
$$
Voir le Lemme \ref{constructions-lemma-invertible-map-into-proj}.
Si $V \subset \Gamma(X, \mathcal{L})$ est un sous-espace, nous noterons
simplement $\varphi_{\mathcal{L}, V}$ le morphisme ainsi construit.
Si $\dim(V) = n + 1$ et si l'on choisit
une base $v_0, \ldots, v_n$ de $V$,
le diagramme
$$
\xymatrix{
X \ar@{=}[d] \ar[rr]_{\varphi_{\mathcal{L}, \psi}} & &
\mathbf{P}(V) \ar[d]^{\cong} \\
X \ar[rr]^{\varphi_{(\mathcal{L}, (s_0, \ldots, s_n))}} & &
\mathbf{P}^n_k
}
$$
est commutatif, où $s_i = \psi(v_i) \in \Gamma(X, \mathcal{L})$,
où $\varphi_{(\mathcal{L}, (s_0, \ldots, s_n))}$ est défini comme dans la
section \ref{constructions-section-projective-space}, et où la flèche verticale de
droite correspond à l'isomorphisme
$k[T_0, \ldots, T_n] \to \text{Sym}(V)$ envoyant $T_i$ sur $v_i$.
\end{example}

\begin{example}
\label{constructions-example-projective-bundle}
Le morphisme $\text{Sym}^n(\mathcal{E}) \to
\pi_*(\mathcal{O}_{\mathbf{P}(\mathcal{E})}(n))$ est un isomorphisme lorsque
$\mathcal{E}$ est localement libre, mais pas nécessairement dans le cas
général. Nous allons donner un exemple où il n'est pas injectif pour
$n = 1$. Posons $S = \Spec(A)$, avec
$$
A = k[u, v, s_1, s_2, t_1, t_2]/I
$$
où $k$ est un corps et
$$
I = (-us_1 + vt_1 + ut_2, vs_1 + us_2 - vt_2, vs_2, ut_1).
$$
Notons $\overline{u}$ la classe de $u$
dans $A$, et de même pour les autres
variables. Soit $M = (Ax \oplus Ay)/A(\overline{u}x + \overline{v}y)$,
de sorte que
$$
\text{Sym}(M) = A[x, y]/(\overline{u}x + \overline{v}y)
= k[x, y, u, v, s_1, s_2, t_1, t_2]/J
$$
où
$$
J = (-us_1 + vt_1 + ut_2, vs_1 + us_2 - vt_2, vs_2, ut_1, ux + vy).
$$
Le fibré projectif associé au faisceau quasi-cohérent
$\mathcal{E} = \widetilde{M}$ sur $S = \Spec(A)$ est alors le schéma
$$
P =
\text{Proj}(\text{Sym}(M)).
$$
Il admet le recouvrement ouvert affine $P = D_{+}(x) \cup D_{+}(y)$.
Considérons l'élément $m \in M$ qui est l'image de $us_1x + vt_2y$.
On a
$$
x(us_1x + vt_2y) = (s_1x + s_2y)(ux + vy) \bmod I
$$
et
$$
y(us_1x + vt_2y) = (t_1x + t_2y)(ux + vy) \bmod I.
$$
La première égalité montre que $m$ s'envoie sur la section nulle de
$\mathcal{O}_P(1)$ sur $D_{+}(x)$ ; la seconde montre
que son image dans $\mathcal{O}_P(1)$ est nulle
sur $D_{+}(y)$. Ainsi, l'image de $m$ dans
$\Gamma(P, \mathcal{O}_P(1))$ est nulle.
Montrons pourtant que $m \not = 0$ : $m$ donnera ainsi une section globale
non nulle de $\mathcal{E}$ dont l'image dans
$\Gamma(P, \mathcal{O}_P(1))$ est nulle. Supposons $m = 0$ pour obtenir
une contradiction. Il existe alors
$f \in k[u, v, s_1, s_2, t_1, t_2]$ tel que
$$
us_1x + vt_2y = f(ux + vy) \bmod I
$$
Puisque $I$ est engendré par des polynômes homogènes de degré $2$, on peut
décomposer $f$ en composantes homogènes et ne conserver que celle de degré
1. Autrement dit, on peut supposer
$$
f = au + bv + \alpha_1s_1 + \alpha_2s_2 + \beta_1t_1 + \beta_2t_2
$$
avec $a, b, \alpha_1, \alpha_2, \beta_1, \beta_2 \in k$.
On obtient les conditions
$$
\begin{matrix}
us_1 - u(au + bv + \alpha_1s_1 + \alpha_2s_2 + \beta_1t_1 + \beta_2t_2)
\in I \\
vt_2 - v(au + bv + \alpha_1s_1 + \alpha_2s_2 + \beta_1t_1 + \beta_2t_2)
\in I
\end{matrix}
$$
Il n'y a aucun terme $u^2, uv, v^2$ dans les générateurs de $I$ ;
on a donc $a = b = 0$. Il reste les relations
$$
\begin{matrix}
us_1 - u(\alpha_1s_1 + \alpha_2s_2 + \beta_1t_1 + \beta_2t_2)
\in I \\
vt_2 - v(\alpha_1s_1 + \alpha_2s_2 + \beta_1t_1 + \beta_2t_2)
\in I
\end{matrix}
$$
Le premier générateur de $I$ permet de
remplacer $us_1$ par $vt_1 + ut_2$,
le deuxième générateur de $I$ de remplacer $vs_1$
par $-us_2 + vt_2$, le troisième d'éliminer
$vs_2$, et le troisième d'éliminer $ut_1$. On obtient
$$
\begin{matrix}
(1 - \alpha_1)vt_1 + (1 - \alpha_1)ut_2 - \alpha_2us_2 - \beta_2ut_2 = 0 \\
(1 - \alpha_1)vt_2 + \alpha_1us_2 - \beta_1vt_1 - \beta_2vt_2 = 0
\end{matrix}
$$
Ces relations imposent à $\alpha_1$ d'être à la fois égal à $0$ et à $1$,
ce qui donne la contradiction cherchée.
\end{example}


\begin{lemma}
\label{constructions-lemma-projective-bundle-separated}
Soit $S$ un schéma. Le morphisme structural
$\mathbf{P}(\mathcal{E}) \to S$ d'un
fibré projectif sur $S$ est séparé.
\end{lemma}

\begin{proof}
Cela résulte immédiatement du Lemme \ref{constructions-lemma-relative-proj-separated}.
\end{proof}

\begin{lemma}
\label{constructions-lemma-projective-space-bundle}
Soient $S$ un schéma et $n \geq 0$. Alors $\mathbf{P}^n_S$ est un fibré
projectif sur $S$.
\end{lemma}

\begin{proof}
On a
$$
\mathbf{P}^n_{\mathbf{Z}} =
\text{Proj}(\mathbf{Z}[T_0, \ldots, T_n]) =
\underline{\text{Proj}}_{\Spec(\mathbf{Z})}
\left(\widetilde{\mathbf{Z}[T_0, \ldots, T_n]}\right)
$$
où l'anneau $\mathbf{Z}[T_0, \ldots, T_n]$ est gradué par
$\deg(T_i) = 1$, les éléments de $\mathbf{Z}$ étant de degré $0$.
Rappelons que $\mathbf{P}^n_S$ est défini comme
$\mathbf{P}^n_{\mathbf{Z}} \times_{\Spec(\mathbf{Z})} S$.
De plus, la formation du spectre homogène relatif commute au changement de
base ; voir le Lemme \ref{constructions-lemma-relative-proj-base-change}.
Pour tout schéma $g : S \to \Spec(\mathbf{Z})$, on a
$g^*\mathcal{O}_{\Spec(\mathbf{Z})}[T_0, \ldots, T_n]
= \mathcal{O}_S[T_0, \ldots, T_n]$.
En combinant ces observations, on obtient
$$
\mathbf{P}^n_S = \underline{\text{Proj}}_S(\mathcal{O}_S[T_0, \ldots, T_n]).
$$
Enfin,
$\mathcal{O}_S[T_0, \ldots, T_n] = \text{Sym}(\mathcal{O}_S^{\oplus n + 1})$.
Ainsi, $\mathbf{P}^n_S$ est un fibré projectif sur $S$.
\end{proof}



\section{Grassmanniennes}
\label{constructions-section-grassmannian}

\noindent
Dans cette section, nous introduisons les foncteurs grassmanniens usuels et
montrons qu'ils sont représentables par des schémas. Fixons des entiers
$k$, $n$ tels que $0 < k < n$. Nous allons construire un foncteur
\begin{equation}
\label{constructions-equation-gkn}
G(k, n) : \Sch \longrightarrow \textit{Ensembles}
\end{equation}
qui, en termes approximatifs, paramétrera les sous-espaces de dimension $k$
d'un espace de dimension $n$. Pour des raisons techniques, il est cependant
plus commode de paramétrer les quotients de dimension $(n - k)$ ; c'est la
convention que nous adopterons.

\medskip\noindent
Plus précisément, $G(k, n)$ associe à un schéma $S$ l'ensemble
$G(k, n)(S)$ des classes d'isomorphisme de surjections
$$
q : \mathcal{O}_S^{\oplus n} \longrightarrow \mathcal{Q}
$$
où $\mathcal{Q}$ est un $\mathcal{O}_S$-module localement libre de rang
$n - k$. Il s'agit bien d'un ensemble : cela résulte, par exemple,
de Modules,
Lemme \ref{modules-lemma-set-isomorphism-classes-finite-type-modules},
ou du fait que la classe d'isomorphisme de $q$ est déterminée par le noyau de $q$
(et que les sous-faisceaux d'un faisceau donné forment un ensemble).
À un morphisme de schémas $f : T \to S$, on associe l'application
$G(k, n)(f) : G(k, n)(S) \to G(k, n)(T)$ envoyant la classe de
$q : \mathcal{O}_S^{\oplus n} \longrightarrow \mathcal{Q}$ sur celle de
$f^*q : \mathcal{O}_T^{\oplus n} \longrightarrow f^*\mathcal{Q}$.
Cette définition a un sens, car (1) $f^*\mathcal{O}_S = \mathcal{O}_T$,
(2) $f^*$ est additif, (3) $f^*$ préserve les modules localement libres
(Modules, Lemme \ref{modules-lemma-pullback-locally-free}) et (4) $f^*$
est exact à droite (Modules,
Lemme \ref{modules-lemma-exactness-pushforward-pullback}).

\begin{lemma}
\label{constructions-lemma-gkn-representable}
Soit $0 < k < n$. Le foncteur $G(k, n)$ de (\ref{constructions-equation-gkn}) est
représentable par un schéma.
\end{lemma}

\begin{proof}
Posons $F = G(k, n)$ et appliquons le critère de Schémas,
Lemme \ref{schemes-lemma-glue-functors}.
Le foncteur $F$ est un faisceau pour la topologie de Zariski parce que
l'on peut recoller les faisceaux ; voir Faisceaux,
section \ref{sheaves-section-glueing-sheaves} (certains détails sont omis).

\medskip\noindent
La famille de sous-foncteurs $F_i$.
Soit $I$ l'ensemble des parties de
$\{1, \ldots, n\}$ de cardinal $n - k$.
Pour un schéma $S$ et $j \in \{1, \ldots, n\}$, notons $e_j$ la section
globale
$$
e_j = (0, \ldots, 0, 1, 0, \ldots, 0)\quad(1\text{ in }j\text{th spot})
$$
de $\mathcal{O}_S^{\oplus n}$. Ces sections forment une base de
$\mathcal{O}_S^{\oplus n}$. De même, pour $j \in \{1, \ldots, n - k\}$,
notons $f_j$ la section globale de $\mathcal{O}_S^{\oplus n - k}$ dont
toutes les composantes sont nulles, sauf la $j$-ième qui vaut $1$.
Pour $i \in I$, soit
$$
s_i : \mathcal{O}_S^{\oplus n - k} \longrightarrow \mathcal{O}_S^{\oplus n}
$$
la somme directe des injections canoniques
$\mathcal{O}_S \to \mathcal{O}_S^{\oplus n}$ correspondant aux éléments
de $I$. Plus précisément, si $i = \{i_1, \ldots, i_{n - k}\}$ avec
$i_1 < i_2 < \ldots < i_{n - k}$, le morphisme $s_i$ envoie $f_j$ sur
$e_{i_j}$ pour $j \in \{1, \ldots, n - k\}$. On pose
$$
F_i(S) = \{q : \mathcal{O}_S^{\oplus n} \to \mathcal{Q} \in F(S) \mid
q \circ s_i \text{ is surjective}\}
\subset F(S)
$$
Pour un morphisme de schémas $f : T \to S$, l'image inverse $f^*s_i$ est
le morphisme correspondant sur $T$. Puisque $f^*$ est exact à droite
(Modules, Lemme \ref{modules-lemma-exactness-pushforward-pullback}),
on en déduit que $F_i$ est un sous-foncteur de $F$.

\medskip\noindent
Représentabilité de $F_i$. Après renumérotation, on peut supposer
$i = \{1, \ldots, n - k\}$ ; ainsi, $s_i$ est l'inclusion des
$n - k$ premiers facteurs directs. Si $q \circ s_i$ est surjectif, alors $q \circ s_i$ est
un isomorphisme, puisqu'il s'agit d'un morphisme surjectif entre modules
localement libres de même rang fini (Modules,
Lemme \ref{modules-lemma-map-finite-locally-free}).
Pour $q : \mathcal{O}_S^{\oplus n} \to \mathcal{Q}$ dans $F_i(S)$,
on peut donc utiliser $q \circ s_i$ pour identifier $\mathcal{Q}$ à
$\mathcal{O}_S^{\oplus n - k}$. On obtient ainsi
$$
q : \mathcal{O}_S^{\oplus n} \longrightarrow \mathcal{O}_S^{\oplus n - k}
$$
qui envoie $e_j$ sur $f_j$ pour $j = 1, \ldots, n - k$, avec les notations
ci-dessus. Pour déterminer entièrement $q$, il reste à choisir les images
$q(e_{n - k + 1}), \ldots, q(e_n)$ dans
$\Gamma(S, \mathcal{O}_S^{\oplus n - k})$.
Ainsi, $F_i$ est isomorphe au foncteur
$$
S \longmapsto
\prod\nolimits_{j = n - k + 1, \ldots, n}
\Gamma(S,  \mathcal{O}_S^{\oplus n - k})
$$
qui est lui-même isomorphe au produit de $k(n - k)$ copies du foncteur
$S \mapsto \Gamma(S, \mathcal{O}_S)$. Ce dernier est représentable par
$\mathbf{A}^1_\mathbf{Z}$, d'après Schémas,
Exemple \ref{schemes-example-global-sections}. Par conséquent, $F_i$ est
représentable par $\mathbf{A}^{k(n - k)}_\mathbf{Z}$, puisque le produit
fibré sur $\Spec(\mathbf{Z})$ est le produit dans la catégorie des schémas.

\medskip\noindent
L'inclusion $F_i \subset F$ est représentable par des immersions ouvertes.
Soient $S$ un schéma et $q : \mathcal{O}_S^{\oplus n} \to \mathcal{Q}$
un élément de $F(S)$. Par Modules,
Lemme \ref{modules-lemma-finite-type-surjective-on-stalk}, l'ensemble
$U_i = \{s \in S \mid (q \circ s_i)_s\text{ surjective}\}$ est ouvert
dans $S$. Puisque $\mathcal{O}_{S, s}$ est un anneau local et
$\mathcal{Q}_s$ un $\mathcal{O}_{S, s}$-module de type fini, le lemme de
Nakayama (Algèbre, Lemme \ref{algebra-lemma-NAK}) donne
$$
s \in U_i \Leftrightarrow
\left(
\text{the map }
\kappa(s)^{\oplus n - k} \to \mathcal{Q}_s/\mathfrak m_s\mathcal{Q}_s
\text{ induced by }
(q \circ s_i)_s
\text{ is surjective}
\right)
$$
Soient $f : T \to S$ un morphisme de schémas et $t \in T$ un point
d'image $s \in S$. On a
$(f^*\mathcal{Q})_t =
\mathcal{Q}_s \otimes_{\mathcal{O}_{S, s}} \mathcal{O}_{T, t}$
(Faisceaux, Lemme \ref{sheaves-lemma-stalk-pullback-modules}), et de même
pour les morphismes correspondants. Ainsi, l'application
$$
\kappa(t)^{\oplus n - k} \to (f^*\mathcal{Q})_t/\mathfrak m_t(f^*\mathcal{Q})_t
$$
induite par $(f^*q \circ f^*s_i)_t$ s'obtient par changement de base de
l'application
$\kappa(s)^{\oplus n - k} \to \mathcal{Q}_s/\mathfrak m_s\mathcal{Q}_s$
ci-dessus suivant l'extension de corps $\kappa(t)/\kappa(s)$.
Il en résulte que $s \in U_i$ si et seulement si $t$ appartient à l'ouvert
correspondant pour $f^*q$. En particulier, $T \to S$ se factorise par
$U_i$ si et seulement si $f^*q \in F_i(T)$, comme voulu.

\medskip\noindent
La famille des $F_i$, $i \in I$, recouvre $F$.
Soit $q : \mathcal{O}_S^{\oplus n} \to \mathcal{Q}$ un élément de $F(S)$.
Il faut montrer que, pour tout point $s$ de $S$, il existe $i \in I$ tel
que $s_i$ soit surjectif dans un voisinage de $s$.
Il suffit donc de montrer que l'une des composées
$$
\kappa(s)^{\oplus n - k} \xrightarrow{s_i}
\kappa(s)^{\oplus n} \rightarrow
\mathcal{Q}_s/\mathfrak m_s\mathcal{Q}_s
$$
est surjective (voir le paragraphe précédent). Comme
$\mathcal{Q}_s/\mathfrak m_s\mathcal{Q}_s$ est un espace vectoriel de
dimension $n - k$, cela résulte de l'algèbre linéaire.
\end{proof}

\begin{definition}
\label{constructions-definition-grassmannian}
Soit $0 < k < n$. Le schéma $\mathbf{G}(k, n)$ représentant le foncteur
$G(k, n)$ est appelé la {\it grassmannienne sur $\mathbf{Z}$}.
Son changement de base $\mathbf{G}(k, n)_S$ à un schéma $S$ est appelé la
{\it grassmannienne sur $S$}. Si $R$ est un anneau, le changement de base
à $\Spec(R)$ est noté $\mathbf{G}(k, n)_R$ et appelé la
{\it grassmannienne sur $R$}.
\end{definition}

\noindent
Cette définition a un sens, puisque la
représentabilité de ces foncteurs a
été démontrée dans le Lemme \ref{constructions-lemma-gkn-representable}.

\begin{lemma}
\label{constructions-lemma-projective-space-grassmannian}
Soit $n \geq 1$. Il existe un isomorphisme canonique
$\mathbf{G}(n, n + 1) = \mathbf{P}^n_\mathbf{Z}$.
\end{lemma}

\begin{proof}
D'après le Lemme \ref{constructions-lemma-projective-space}, le schéma
$\mathbf{P}^n_\mathbf{Z}$ représente le foncteur qui associe à un schéma
$S$ l'ensemble des classes d'isomorphisme de couples
$(\mathcal{L}, (s_0, \ldots, s_n))$ formés d'un module inversible
$\mathcal{L}$ et de $(n + 1)$ sections
globales qui engendrent $\mathcal{L}$.
Un tel couple définit un quotient
$$
\mathcal{O}_S^{\oplus n + 1} \longrightarrow \mathcal{L},\quad
(h_0, \ldots, h_n) \longmapsto \sum h_i s_i.
$$
Réciproquement, un élément
$q : \mathcal{O}_S^{\oplus n + 1} \to \mathcal{Q}$ de $G(n, n + 1)(S)$
donne le couple $(\mathcal{Q}, (q(e_1), \ldots, q(e_{n + 1})))$.
Ici, les $e_i$, $i = 1, \ldots, n + 1$, sont les sections de la base
canonique du module libre $\mathcal{O}_S^{\oplus n + 1}$.
Nous omettons la vérification que ces constructions définissent des
transformations de foncteurs réciproques.
\end{proof}






% OK, start here.
%

\chapter{Propriétés des schémas}



\phantomsection
\label{properties-section-phantom}


\section{Introduction}
\label{properties-section-introduction}

\noindent
Dans ce chapitre, nous introduisons quelques propriétés absolues des schémas.
Une référence fondamentale est \cite{EGA}.




\section{Ensembles constructibles}
\label{properties-section-constructible}

\noindent
Les ensembles constructibles et localement constructibles sont introduits dans
Topologie, section \ref{topology-section-constructible}.
On peut caractériser les parties localement constructibles des schémas de la
manière suivante.

\begin{lemma}
\label{properties-lemma-locally-constructible}
Soit $X$ un schéma.
Une partie $E$ de $X$ est localement constructible dans $X$ si et seulement si
$E \cap U$ est constructible dans $U$ pour tout ouvert affine $U$ de $X$.
\end{lemma}

\begin{proof}
Supposons $E$ localement constructible. Il existe alors un recouvrement ouvert
$X = \bigcup U_i$ tel que $E \cap U_i$ soit constructible dans $U_i$
pour tout $i$. Soit $V \subset X$ un ouvert affine quelconque. On peut trouver un
recouvrement ouvert affine fini $V = V_1 \cup \ldots \cup V_m$ tel que, pour tout $j$,
on ait $V_j \subset U_i$ pour un certain $i = i(j)$. D'après le
Lemme \ref{topology-lemma-open-immersion-constructible-inverse-image} de Topologie,
chaque $E \cap V_j$ est constructible dans $V_j$. Comme les inclusions
$V_j \to V$ sont quasi-compactes (voir
Schémas, Lemme \ref{schemes-lemma-quasi-compact-affine}),
on conclut que $E \cap V$ est constructible dans $V$ d'après le
Lemme \ref{topology-lemma-collate-constructible} de Topologie.
L'implication réciproque est immédiate.
\end{proof}

\begin{lemma}
\label{properties-lemma-generic-point-in-constructible}
Soient $X$ un schéma et $E \subset X$ un sous-ensemble localement constructible.
Soit $\xi \in X$ un point générique d'une composante irréductible de $X$.
\begin{enumerate}
\item Si $\xi \in E$, alors un voisinage ouvert de
$\xi$ est contenu dans $E$.
\item Si $\xi \not \in E$, alors un voisinage ouvert
de $\xi$ est disjoint de $E$.
\end{enumerate}
\end{lemma}

\begin{proof}
Le complémentaire d'un sous-ensemble localement constructible étant
localement constructible, il suffit de montrer (2). On peut supposer $X$ affine,
et donc $E$ constructible (lemme \ref{properties-lemma-locally-constructible}).
Dans ce cas, $X$ est un espace
spectral (Algèbre, lemme \ref{algebra-lemma-spec-spectral}).
Alors $\xi \not \in E$ entraîne $\xi \not \in \overline{E}$ d'après
Topologie, lemme \ref{topology-lemma-constructible-stable-specialization-closed},
car aucun point de $X$ distinct de $\xi$
ne se spécialise en $\xi$.
\end{proof}

\begin{lemma}
\label{properties-lemma-quasi-separated-quasi-compact-open-retrocompact}
Soit $X$ un schéma quasi-séparé. L'intersection de deux
ouverts quasi-compacts de $X$ est un ouvert quasi-compact de $X$.
Tout ouvert quasi-compact de $X$ est rétrocompact dans $X$.
\end{lemma}

\begin{proof}
Si $U$ et $V$ sont des ouverts quasi-compacts, alors
$U \cap V = \Delta^{-1}(U \times V)$, où $\Delta : X \to X \times X$
est la diagonale. Comme $X$ est quasi-séparé, $\Delta$ est
quasi-compact. Ainsi $U \cap V$ est quasi-compact, puisque
$U \times V$ l'est (détails omis ; utiliser
Schémas, lemme \ref{schemes-lemma-affine-covering-fibre-product},
pour voir que $U \times V$ est une réunion finie d'ouverts
affines). Les autres assertions résultent de la première et
de Topologie, lemme \ref{topology-lemma-topology-quasi-separated-scheme}.
\end{proof}

\begin{lemma}
\label{properties-lemma-quasi-compact-quasi-separated-spectral}
Soit $X$ un schéma quasi-compact et quasi-séparé.
Alors l'espace topologique sous-jacent à $X$ est spectral.
\end{lemma}

\begin{proof}
D'après Topologie, définition \ref{topology-definition-spectral-space},
il faut vérifier que $X$ est sobre, quasi-compact, admet une
base d'ouverts quasi-compacts, et que l'intersection de deux
ouverts quasi-compacts est quasi-compacte. Cela résulte de
Schémas, lemmes \ref{schemes-lemma-scheme-sober} et
\ref{schemes-lemma-basis-affine-opens},
ainsi
que du lemme \ref{properties-lemma-quasi-separated-quasi-compact-open-retrocompact}
ci-dessus.
\end{proof}

\begin{lemma}
\label{properties-lemma-constructible-quasi-compact-quasi-separated}
Soit $X$ un schéma quasi-compact et quasi-séparé. Tout
sous-ensemble localement constructible de $X$ est constructible.
\end{lemma}

\begin{proof}
Comme $X$ est quasi-compact, on peut choisir un recouvrement ouvert affine fini
$X = V_1 \cup \ldots \cup V_m$. Comme $X$ est quasi-séparé, chaque $V_i$ est
rétrocompact dans $X$ d'après
le lemme \ref{properties-lemma-quasi-separated-quasi-compact-open-retrocompact}.
Ainsi,
par Topologie, lemme \ref{topology-lemma-collate-constructible},
un sous-ensemble $E \subset X$ est constructible dans $X$ si et seulement si
$E \cap V_j$ est constructible dans $V_j$. On conclut par le
lemme \ref{properties-lemma-locally-constructible}.
\end{proof}

\begin{lemma}
\label{properties-lemma-retrocompact}
Soit $X$ un schéma. Un sous-ensemble $E$ de $X$ est rétrocompact dans $X$ si et seulement si
$E \cap U$ est quasi-compact pour tout ouvert affine $U$ de $X$.
\end{lemma}

\begin{proof}
Cela résulte immédiatement du fait que tout ouvert quasi-compact de $X$ est une
réunion finie d'ouverts affines.
\end{proof}

\begin{lemma}
\label{properties-lemma-stratification-locally-finite-constructible}
Une partition $X = \coprod_{i \in I} X_i$ d'un schéma $X$ dont
les parties sont rétrocompactes est localement finie si et seulement si ses
parties sont localement constructibles.
\end{lemma}

\begin{proof}
Voir Topologie, définitions
\ref{topology-definition-quasi-compact},
\ref{topology-definition-paritition}et
\ref{topology-definition-locally-finite},
pour les définitions de rétrocompact, de partition et de localement fini.

\medskip\noindent
Si la partition est localement finie et si $U \subset X$ est
un ouvert affine, alors $U = \coprod_{i \in I} U \cap X_i$
est une partition finie (plus précisément, toutes ses parties sauf
un nombre fini sont vides). Ainsi $U \cap X_i$ est quasi-compact
et son complémentaire est rétrocompact dans $U$, étant une réunion
finie de parties rétrocompactes. Donc $U \cap X_i$ est constructible
d'après Topologie, lemme \ref{topology-lemma-locally-closed-constructible-image}.
Il s'ensuit que $X_i$ est localement constructible par le
lemme \ref{properties-lemma-locally-constructible}.

\medskip\noindent
Supposons les parties localement constructibles. Pour tout ouvert
affine $U \subset X$, on obtient alors un recouvrement $U = \coprod X_i \cap U$
par des sous-ensembles constructibles. La topologie constructible étant
quasi-compacte, voir
Topologie, lemme \ref{topology-lemma-constructible-hausdorff-quasi-compact},
ce recouvrement admet un raffinement fini ; autrement
dit, la partition est localement finie.
\end{proof}




\section{Schémas intègres, irréductibles et réduits}
\label{properties-section-integral}

\begin{definition}
\label{properties-definition-integral}
Soit $X$ un schéma. On dit que $X$ est {\it intègre} s'il est non vide et
si, pour tout ouvert affine non vide $\Spec(R) = U \subset X$, l'anneau $R$
est intègre.
\end{definition}

\begin{lemma}
\label{properties-lemma-characterize-reduced}
Soit $X$ un schéma. Les
assertions suivantes sont équivalentes.
\begin{enumerate}
\item Le schéma $X$ est réduit, voir
Schémas, définition \ref{schemes-definition-reduced}.
\item Il existe un recouvrement ouvert affine $X = \bigcup U_i$
tel que chaque $\Gamma(U_i, \mathcal{O}_X)$ soit réduit.
\item Pour tout ouvert affine $U \subset X$, l'anneau
$\mathcal{O}_X(U)$ est réduit.
\item Pour tout ouvert $U \subset X$, l'anneau $\mathcal{O}_X(U)$ est réduit.
\end{enumerate}
\end{lemma}

\begin{proof}
Voir Schémas, lemmes \ref{schemes-lemma-reduced} et
\ref{schemes-lemma-affine-reduced}.
\end{proof}

\begin{lemma}
\label{properties-lemma-characterize-irreducible}
Soit $X$ un schéma. Les
assertions suivantes sont équivalentes.
\begin{enumerate}
\item Le schéma $X$ est irréductible.
\item Il existe un recouvrement ouvert affine $X = \bigcup_{i \in I} U_i$
tel que $I$ soit non vide, que $U_i$ soit irréductible pour tout $i \in I$ et que
$U_i \cap U_j \not = \emptyset$ pour tous $i, j \in I$.
\item Le schéma $X$ est non vide et tout ouvert affine non vide
$U \subset X$ est irréductible.
\end{enumerate}
\end{lemma}

\begin{proof}
Supposons (1). D'après Schémas, lemme \ref{schemes-lemma-scheme-sober},
$X$ possède un unique point générique $\eta$. On a alors
$X = \overline{\{\eta\}}$. Par conséquent, $\eta$ appartient à
tout ouvert affine non vide $U \subset X$. Cela implique que le
point $\eta \in U$ est dense, donc que $U$ est irréductible.
Cela implique aussi que deux ouverts affines non vides se
rencontrent. Ainsi (1) entraîne à la fois (2) et (3).

\medskip\noindent
Supposons (2). Soit $X = Z_1 \cup Z_2$ une réunion de deux sous-ensembles
fermés. Pour tout $i$, on a soit $U_i \subset Z_1$, soit $U_i \subset Z_2$.
Choisissons un $i \in I$ et supposons $U_i \subset Z_1$ (quitte à permuter
$Z_1$, $Z_2$). Pour tout $j \in I$, l'ouvert $U_i \cap U_j$ est dense dans
$U_j$ et contenu dans le fermé $Z_1 \cap U_j$. On en déduit
que $U_j \subset Z_1$ également. Donc $X = Z_1$, comme voulu.

\medskip\noindent
Supposons (3). Choisissons un recouvrement ouvert affine $X = \bigcup_{i \in I} U_i$.
On peut supposer chaque $U_i$ non vide.
Puisque $X$ est non vide, $I$ est non vide.
Par hypothèse, chaque $U_i$ est irréductible.
Supposons $U_i \cap U_j = \emptyset$ pour un couple $i, j \in I$.
Alors l'ouvert $U_i \amalg  U_j = U_i \cup U_j$ est affine, voir
Schémas, lemme \ref{schemes-lemma-disjoint-union-affines}.
Il est donc irréductible par hypothèse, ce qui est absurde. On en conclut que (3)
entraîne (2). Le lemme est démontré.
\end{proof}

\begin{lemma}
\label{properties-lemma-characterize-integral}
Un schéma $X$ est intègre si et seulement s'il est réduit et irréductible.
\end{lemma}

\begin{proof}
Si $X$ est irréductible, tout ouvert affine $\Spec(R) = U \subset X$
est irréductible. Si $X$ est réduit, $R$ est réduit d'après
le lemme \ref{properties-lemma-characterize-reduced} ci-dessus. Ainsi $R$ est réduit
et $(0)$ est un idéal premier ; autrement dit, $R$ est un anneau intègre.

\medskip\noindent
Si $X$ est intègre, alors, pour tout ouvert affine non vide
$\Spec(R) = U \subset X$, l'anneau $R$ est
réduit, donc $X$ est réduit par le lemme \ref{properties-lemma-characterize-reduced}.
De plus, tout ouvert affine non vide est irréductible. Par
conséquent, $X$ est irréductible, voir le lemme \ref{properties-lemma-characterize-irreducible}.
\end{proof}

\noindent
Dans Exemples, section
\ref{examples-section-connected-locally-integral-not-integral},
on construit un schéma affine connexe dont tous les anneaux locaux sont intègres,
mais qui n'est pas intègre.











\section{Types de schémas définis par des propriétés d'anneaux}
\label{properties-section-properties-rings}

\noindent
Dans cette section, nous étudions quelles propriétés
d'anneaux permettent de définir des propriétés locales des schémas.

\begin{definition}
\label{properties-definition-property-local}
Soit $P$ une propriété des
anneaux. On dit que $P$ est {\it locale} si les conditions suivantes sont satisfaites :
\begin{enumerate}
\item Pour tout anneau $R$ et tout $f \in R$, on a
$P(R) \Rightarrow P(R_f)$.
\item Pour tout anneau $R$ et tous $f_i \in R$ tels que
$(f_1, \ldots, f_n) = R$, on a
$\forall i, P(R_{f_i}) \Rightarrow P(R)$.
\end{enumerate}
\end{definition}

\begin{definition}
\label{properties-definition-locally-P}
Soient $P$ une propriété des anneaux et $X$ un schéma. On
dit que $X$ est {\it localement $P$} si, pour tout $x \in X$,
il existe un voisinage ouvert affine $U$ de $x$
dans $X$ tel que $\mathcal{O}_X(U)$ possède la propriété $P$.
\end{definition}

\noindent
Cette notion n'est satisfaisante que si la propriété est locale.
Même lorsque $P$ est une propriété locale, nous
n'utiliserons pas automatiquement cette définition pour dire qu'un schéma est «
localement $P$», à moins d'énoncer aussi explicitement cette
définition ailleurs.

\begin{lemma}
\label{properties-lemma-locally-P}
Soient $X$ un schéma et $P$ une propriété locale des anneaux. Les
assertions suivantes sont équivalentes :
\begin{enumerate}
\item Le schéma $X$ est localement $P$.
\item Pour tout ouvert affine $U \subset X$, la propriété
$P(\mathcal{O}_X(U))$ est satisfaite.
\item Il existe un recouvrement ouvert affine $X = \bigcup U_i$ tel que
chaque $\mathcal{O}_X(U_i)$ satisfasse $P$.
\item Il existe un recouvrement ouvert $X = \bigcup X_j$
tel que chaque sous-schéma ouvert $X_j$ soit localement $P$.
\end{enumerate}
De plus, si $X$ est localement $P$, tout sous-schéma ouvert est
localement $P$.
\end{lemma}

\begin{proof}
On a évidemment (1) $\Leftrightarrow$ (3) et (2) $\Rightarrow$ (1).
Si (3) $\Rightarrow$ (2), alors la dernière assertion du lemme est
vraie, et il en résulte aisément que (4) est aussi équivalente à (1).
Montrons donc (3) $\Rightarrow$ (2).

\medskip\noindent
Soit $X = \bigcup U_i$ un recouvrement ouvert affine, avec
$U_i = \Spec(R_i)$. Supposons $P(R_i)$.
Soit $\Spec(R) = U \subset X$ un ouvert affine quelconque.
D'après Schémas, lemme \ref{schemes-lemma-good-subcover},
il existe un recouvrement standard de $U = \Spec(R)$ par
des ouverts principaux $D(f_j)$ tel que chaque anneau $R_{f_j}$ soit
une localisation principale d'un des anneaux $R_i$. Par
la définition \ref{properties-definition-property-local} (1), on obtient $P(R_{f_j})$.
D'où $P(R)$ par la définition \ref{properties-definition-property-local} (2).
\end{proof}

\noindent
Voici un exemple d'application.

\begin{lemma}
\label{properties-lemma-reduced-is-locally-reduced}
Soit $X$ un schéma. Alors $X$ est réduit si et seulement si $X$ est
« localement réduit » au sens de la définition \ref{properties-definition-locally-P}.
\end{lemma}

\begin{proof}
Cela résulte immédiatement du lemme \ref{properties-lemma-characterize-reduced}.
\end{proof}

\begin{lemma}
\label{properties-lemma-properties-local}
Les propriétés suivantes d'un anneau $R$ sont locales.
\begin{enumerate}
\item (Cohen-Macaulay.)
L'anneau $R$ est noethérien et de
Cohen--Macaulay, voir Algèbre, définition \ref{algebra-definition-ring-CM}.
\item (Régulier.)
L'anneau $R$ est noethérien et régulier, voir
Algèbre, définition \ref{algebra-definition-regular}.
\item (Absolument noethérien.)
L'anneau $R$ est de type fini sur $Z$.
\item Ajouter ici d'autres propriétés au besoin.\footnote{Nous ne répertorions ici que
les propriétés qui n'ont pas déjà été traitées séparément ailleurs.}
\end{enumerate}
\end{lemma}

\begin{proof}
Omis.
\end{proof}















\section{Schémas noethériens}
\label{properties-section-noetherian}

\noindent
Rappelons qu'un anneau $R$ est {\it noethérien} s'il satisfait à la condition de
chaîne ascendante sur les idéaux. De manière équivalente, tout idéal de $R$ est de
type fini.

\begin{definition}
\label{properties-definition-noetherian}
Soit $X$ un schéma.
\begin{enumerate}
\item On dit que $X$ est {\it localement noethérien} si tout
$x \in X$ possède un voisinage ouvert affine
$\Spec(R) = U \subset X$ tel que l'anneau $R$ soit noethérien.
\item On dit que $X$ est {\it noethérien} si $X$ est localement
noethérien et quasi-compact.
\end{enumerate}
\end{definition}

\noindent
Voici la caractérisation usuelle des schémas localement noethériens.

\begin{lemma}
\label{properties-lemma-locally-Noetherian}
Soit $X$ un schéma. Les assertions suivantes sont équivalentes :
\begin{enumerate}
\item Le schéma $X$ est localement noethérien.
\item Pour tout ouvert affine $U \subset X$, l'anneau $\mathcal{O}_X(U)$
est noethérien.
\item Il existe un recouvrement ouvert affine $X = \bigcup U_i$ tel que
chaque $\mathcal{O}_X(U_i)$ soit noethérien.
\item Il existe un recouvrement ouvert $X = \bigcup X_j$
tel que chaque sous-schéma ouvert $X_j$ soit localement noethérien.
\end{enumerate}
De plus, si $X$ est localement noethérien, tout sous-schéma ouvert
est localement noethérien.
\end{lemma}

\begin{proof}
Il suffit de montrer qu'être noethérien est une propriété locale
des anneaux, voir le lemme \ref{properties-lemma-locally-P}.
Toute localisation d'un anneau noethérien est noethérienne,
voir Algèbre, lemme \ref{algebra-lemma-Noetherian-permanence}.
Le lemme \ref{algebra-lemma-cover} du chapitre Algèbre donne
la seconde condition de la définition \ref{properties-definition-property-local}.
\end{proof}

\begin{lemma}
\label{properties-lemma-immersion-into-noetherian}
Toute immersion $Z \to X$, avec $X$ localement noethérien, est quasi-compacte.
\end{lemma}

\begin{proof}
Une immersion fermée est évidemment
quasi-compacte. Un composé de morphismes quasi-compacts est quasi-compact,
voir Topologie, lemme \ref{topology-lemma-composition-quasi-compact}.
Il suffit donc de montrer qu'une immersion ouverte dans un
schéma localement noethérien est quasi-compacte. À
l'aide de Schémas, lemme \ref{schemes-lemma-quasi-compact-affine},
on se ramène au cas où $X$ est affine.
Tout ouvert du spectre d'un anneau noethérien est
quasi-compact (on peut, par exemple,
combiner Algèbre, lemme \ref{algebra-lemma-Noetherian-topology},
et Topologie, lemmes \ref{topology-lemma-Noetherian} et
\ref{topology-lemma-Noetherian-quasi-compact}).
\end{proof}

\begin{lemma}
\label{properties-lemma-locally-Noetherian-quasi-separated}
Un schéma localement noethérien est quasi-séparé.
\end{lemma}

\begin{proof}
D'après Schémas, lemme \ref{schemes-lemma-characterize-quasi-separated},
il faut montrer que l'intersection $U \cap V$ de deux
ouverts affines de $X$ est quasi-compacte. Cela résulte du
lemme \ref{properties-lemma-immersion-into-noetherian} ci-dessus,
en considérant, par exemple, l'immersion ouverte $U \cap V \to U$. (Au fond,
cela vient simplement de ce que tout ouvert du spectre d'un
anneau noethérien est quasi-compact.)
\end{proof}

\begin{lemma}
\label{properties-lemma-Noetherian-topology}
L'espace topologique sous-jacent à un schéma
(localement) noethérien est (localement) noethérien,
voir Topologie, définition \ref{topology-definition-noetherian}.
\end{lemma}

\begin{proof}
Cela résulte du fait qu'un schéma noethérien est une réunion finie de spectres d'anneaux
noethériens, ainsi que des résultats
suivants : Algèbre, lemme \ref{algebra-lemma-Noetherian-topology}, et
Topologie, lemme \ref{topology-lemma-finite-union-Noetherian}.
\end{proof}

\begin{lemma}
\label{properties-lemma-locally-closed-in-Noetherian}
Tout sous-schéma localement fermé d'un schéma (localement)
noethérien est (localement) noethérien.
\end{lemma}

\begin{proof}
Omis. Indication : tout quotient et toute localisation d'un anneau
noethérien sont noethériens. Dans le cas noethérien, utiliser de
nouveau le fait que tout sous-ensemble d'un espace noethérien est
noethérien pour la topologie induite.
\end{proof}

\begin{lemma}
\label{properties-lemma-Noetherian-irreducible-components}
Un schéma noethérien possède un nombre fini de composantes irréductibles.
\end{lemma}

\begin{proof}
L'espace topologique sous-jacent à un schéma noethérien est noethérien
(lemme \ref{properties-lemma-Noetherian-topology}),
et l'on conclut parce qu'un espace topologique noethérien
n'a qu'un nombre fini de composantes irréductibles
(Topologie, lemme \ref{topology-lemma-Noetherian}).
\end{proof}

\begin{lemma}
\label{properties-lemma-morphism-Noetherian-schemes-quasi-compact}
Tout morphisme de schémas $f : X \to Y$, avec $X$ noethérien,
est quasi-compact.
\end{lemma}

\begin{proof}
Utiliser le lemme \ref{properties-lemma-Noetherian-topology}
et le fait que tout sous-ensemble d'un espace topologique
noethérien est quasi-compact (voir Topologie,
lemmes \ref{topology-lemma-Noetherian} et
\ref{topology-lemma-Noetherian-quasi-compact}).
\end{proof}

\noindent
Voici un lemme
intéressant. Il affirme que tout schéma localement noethérien possède
beaucoup de points fermés (au moins un dans tout fermé non vide).

\begin{lemma}
\label{properties-lemma-locally-Noetherian-closed-point}
Tout schéma localement noethérien non vide possède un point
fermé. Tout fermé non vide d'un schéma localement noethérien possède un point fermé. De
manière équivalente, tout point d'un schéma localement noethérien se
spécialise en un point fermé.
\end{lemma}

\begin{proof}
La seconde assertion découle de la première (à l'aide
de Schémas, lemme \ref{schemes-lemma-reduced-closed-subscheme},
et du lemme \ref{properties-lemma-locally-closed-in-Noetherian}).
Considérons un ouvert affine non vide $U \subset X$.
Soit $x \in U$ un point fermé. Si $x$ est fermé
dans $X$, la démonstration est terminée. Sinon, soit $X_0 \subset X$ le
sous-schéma fermé réduit induit sur $\overline{\{x\}}$.
Alors $U_0 = U \cap X_0$ est un ouvert affine de $X_0$ d'après
Schémas, lemme \ref{schemes-lemma-closed-subspace-scheme}, et
$U_0 = \{x\}$. Soit $y \in X_0$, $y \not = x$, une spécialisation de $x$.
Considérons l'anneau local $R = \mathcal{O}_{X_0, y}$.
C'est un anneau local noethérien, puisque $X_0$ est noethérien
d'après le lemme \ref{properties-lemma-locally-closed-in-Noetherian}. Notons $V \subset \Spec(R)$
l'image inverse de $U_0$ dans $\Spec(R)$ par le morphisme canonique
$\Spec(R) \to X_0$ (voir Schémas, section \ref{schemes-section-points}).
Par construction, $V$ est un singleton dont l'unique point correspond à $x$ (utiliser
Schémas, lemme \ref{schemes-lemma-specialize-points}).
D'après
Algèbre, lemme \ref{algebra-lemma-Noetherian-local-domain-dim-2-infinite-opens},
on a $\dim(R) = 1$.
Autrement dit, $y$ est une spécialisation immédiate
de $x$ (voir Topologie, définition \ref{topology-definition-dimension-function}).
Ainsi, tout
point $y \not = x$ tel que $x \leadsto y$ est une
spécialisation immédiate de $x$. Chacun de ces points est
évidemment fermé, comme voulu.
\end{proof}

\begin{lemma}
\label{properties-lemma-locally-Noetherian-specialization-dvr}
Soit $X$ un schéma localement noethérien.
Soit $x' \leadsto x$ une spécialisation de points de $X$. Alors
\begin{enumerate}
\item il existe un anneau de valuation discrète $R$ et un morphisme
$f : \Spec(R) \to X$ tels que le point générique $\eta$ de
$\Spec(R)$ ait pour image $x'$ et le point fermé ait pour image $x$;
\item si $x \not = x'$, étant donnée une extension de corps
de type fini $K/\kappa(x')$, on peut faire en sorte que l'extension
$\kappa(\eta)/\kappa(x')$
induite par $f$ soit isomorphe à l'extension donnée.
\end{enumerate}
\end{lemma}

\begin{proof}
Supposons d'abord que $x' \leadsto x$ soit une spécialisation dans $X$
avec $x \not = x'$, et soit $K/\kappa(x')$ une extension
de corps de type fini. D'après
Schémas, lemme \ref{schemes-lemma-specialize-points},
et la discussion qui suit
Schémas, lemme \ref{schemes-lemma-characterize-points},
on obtient des homomorphismes d'anneaux $\mathcal{O}_{X, x} \to \kappa(x') \to K$.
Puisque $x \not = x'$, l'image de $\mathcal{O}_{X, x}$ dans $K$
n'est pas un corps (détails omis).
Soit $R \subset K$ un anneau de valuation discrète dont le corps des fractions est
$K$ et qui domine l'image de $\mathcal{O}_{X, x} \to K$, voir
Algèbre, lemme \ref{algebra-lemma-exists-dvr}.
L'homomorphisme d'anneaux $\mathcal{O}_{X, x} \to R$ induit le morphisme
$f : \Spec(R) \to X$, voir
Schémas, lemme \ref{schemes-lemma-morphism-from-spec-local-ring}.
Par construction, ce morphisme possède toutes les propriétés voulues.
Si $x = x'$, on peut poser $R = \kappa(x)[t]_{(t)}$.
\end{proof}

\begin{lemma}
\label{properties-lemma-thin-infinite-sequence}
Soient $S$ un schéma noethérien et $T \subset S$ un sous-ensemble infini.
Il existe alors un sous-ensemble infini $T' \subset T$
tel qu'il n'y ait aucune spécialisation non triviale entre les points de $T'$.
\end{lemma}

\begin{proof}
Soit $T_0 \subset T$ l'ensemble des $t \in T$ qui ne se spécialisent en
aucun autre point de $T$. Si $T_0$ est infini, $T' = T_0$ convient.
On peut donc supposer $T_0$ fini.
Pour $i > 0$, considérons, par récurrence, l'ensemble $T_i \subset T$
des $t \in T$ tels que
\begin{enumerate}
\item $t \not \in T_{i - 1} \cup T_{i - 2} \cup \ldots \cup T_0$;
\item il existe une spécialisation non triviale $t \leadsto t'$ avec
$t' \in T_{i - 1}$;
\item pour toute spécialisation non triviale
$t \leadsto t'$ avec $t' \in T$, on a
$t' \in T_{i - 1} \cup T_{i - 2} \cup \ldots \cup T_0$.
\end{enumerate}
Là encore, si $T_i$ est infini, $T' = T_i$ convient.
Soit $d$ le maximum des dimensions des anneaux locaux
$\mathcal{O}_{S, t}$ pour $t \in T_0$; alors $d$ est un
entier, car $T_0$ est fini et les dimensions des anneaux locaux sont
finies d'après Algèbre, proposition \ref{algebra-proposition-dimension}.
On a $T_i = \emptyset$ pour $i > d$.
En effet, si $t \in T_i$, on peut trouver une suite
de spécialisations non triviales
$t = t_i \leadsto t_{i - 1} \leadsto \ldots \leadsto t_0$
avec $t_0 \in T_0$.
Comme les points $t = t_i, t_{i - 1}, \ldots, t_0$ appartiennent à
$\Spec(\mathcal{O}_{S, t_0})$
(Schémas, lemme \ref{schemes-lemma-specialize-points}),
on obtient $i \leq d$.
Ainsi $\bigcup T_i = T_d \cup \ldots \cup T_0$ est un sous-ensemble fini de $T$.

\medskip\noindent
Supposons qu'un $t \in T$ n'appartienne pas à $\bigcup T_i$. Il doit alors exister une spécialisation
$t \leadsto t'$ avec $t' \in T$ et $t' \not \in \bigcup T_i$. (En effet, si
toutes les spécialisations de $t$ appartiennent à l'ensemble fini $T_d \cup \ldots \cup T_0$,
il existe un plus grand $i$ tel qu'il y ait une spécialisation
$t \leadsto t'$ avec $t' \in T_i$, et alors $t \in T_{i + 1}$ par construction.)
On obtient donc une suite infinie
$$
t \leadsto t' \leadsto t'' \leadsto \ldots
$$
de spécialisations non triviales entre points de $T \setminus \bigcup T_i$.
C'est impossible, car l'espace topologique sous-jacent à $S$
est noethérien d'après le lemme \ref{properties-lemma-locally-Noetherian-quasi-separated}.
\end{proof}

\begin{lemma}
\label{properties-lemma-maximal-points}
Soient $S$ un schéma noethérien et $T \subset S$ un sous-ensemble. Soit
$T_0 \subset T$ l'ensemble des $t \in T$ tels qu'il n'existe aucune
spécialisation non triviale $t' \leadsto t$ avec $t' \in T'$. Alors (a) il n'y a
aucune spécialisation entre les points de $T_0$, (b) tout point de
$T$ est une spécialisation d'un point de $T_0$, et (c) les adhérences
de $T$ et de $T_0$ sont égales.
\end{lemma}

\begin{proof}
Rappelons que $\dim(\mathcal{O}_{S, s}) < \infty$ pour tout $s \in S$,
voir Algèbre, proposition \ref{algebra-proposition-dimension}. Soit $t \in T$.
Si $t' \leadsto t$, la théorie de la dimension donne
$\dim(\mathcal{O}_{S, t'}) \leq \dim(\mathcal{O}_{S, t})$,
avec égalité si et seulement si $t' = t$. Ainsi, si l'on choisit $t' \leadsto t$
avec $\dim(\mathcal{O}_{T, t'})$ minimal, alors $t' \in T_0$.
Autrement dit, 
tout $t \in T$ est une spécialisation d'un élément de $T_0$.
\end{proof}

\begin{lemma}
\label{properties-lemma-countable-dense-subset}
Soient $S$ un schéma noethérien et $T \subset S$ un sous-ensemble infini dense. Il
existe alors un sous-ensemble dénombrable $E \subset T$ dense dans $S$.
\end{lemma}

\begin{proof}
Soit $T'$ l'ensemble des points $s \in S$ tels que $\overline{\{s\}} \cap T$
contienne un sous-ensemble dénombrable d'adhérence $\overline{\{s\}}$.
Tout ensemble fini étant dénombrable, on a $T \subset T'$.
Pour $s \in T'$, choisissons un tel sous-ensemble dénombrable
$E_s \subset \overline{\{s\}} \cap T$.
Soit $E' = \{s_1, s_2, s_3, \ldots\} \subset T'$
un sous-ensemble dénombrable. Alors l'adhérence de $E'$ dans $S$ est
l'adhérence du sous-ensemble dénombrable $\bigcup_n E_{s_n}$ de $T$.
Il s'ensuit que, si $Z$
est une composante irréductible de l'adhérence de $E'$, le point
générique de $Z$ appartient à $T'$.

\medskip\noindent
Notons $T'_0 \subset T'$ le sous-ensemble des $t \in T'$ tels qu'il
n'existe aucune spécialisation non triviale $t' \leadsto t$ avec $t' \in T'$,
comme dans le lemme \ref{properties-lemma-maximal-points}, dont nous utiliserons les
résultats sans le mentionner à nouveau. Si $T'_0$ est infini, choisissons un
sous-ensemble dénombrable $E' \subset T'_0$. D'après l'argument du premier
paragraphe, les points génériques des composantes irréductibles de
l'adhérence de $E'$ appartiennent à $T'$. Mais l'un de ces points se spécialise en une
infinité d'éléments distincts de $E' \subset T'_0$,
ce qui est contradictoire. Ainsi $T'_0$ est fini ; écrivons
$T'_0 = \{s_1, \ldots, s_m\}$. Il en résulte que $S$, qui est
l'adhérence de $T$, est contenu dans l'adhérence de
$\{s_1, \ldots, s_m\}$, elle-même contenue dans l'adhérence du
sous-ensemble dénombrable $E_{s_1} \cup \ldots \cup E_{s_m} \subset T$,
comme voulu.
\end{proof}








\section{Schémas de Jacobson}
\label{properties-section-jacobson}

\noindent
Rappelons qu'un espace est dit {\it de Jacobson} si ses points fermés sont
denses dans tout sous-ensemble fermé,
voir Topologie, section \ref{topology-section-space-jacobson}.

\begin{definition}
\label{properties-definition-jacobson}
Un schéma $S$ est dit {\it de Jacobson} si son espace topologique
sous-jacent est de Jacobson.
\end{definition}

\noindent
Rappelons qu'un anneau $R$ est de Jacobson si tout idéal radical de $R$
est une intersection d'idéaux maximaux, voir
Algèbre, définition \ref{algebra-definition-ring-jacobson}.

\begin{lemma}
\label{properties-lemma-affine-jacobson}
Un schéma affine $\Spec(R)$ est de Jacobson si et seulement
si l'anneau $R$ est de Jacobson.
\end{lemma}

\begin{proof}
C'est le lemme \ref{algebra-lemma-jacobson}du chapitre Algèbre.
\end{proof}

\noindent
Voici la caractérisation usuelle des schémas de Jacobson.
Intuitivement, elle affirme : de Jacobson $\Leftrightarrow$ localement de Jacobson.

\begin{lemma}
\label{properties-lemma-locally-jacobson}
Soit $X$ un schéma. Les assertions suivantes sont équivalentes :
\begin{enumerate}
\item Le schéma $X$ est de Jacobson.
\item Le schéma $X$ est « localement de Jacobson » au sens de la
définition \ref{properties-definition-locally-P}.
\item Pour tout ouvert affine $U \subset X$, l'anneau $\mathcal{O}_X(U)$
est de Jacobson.
\item Il existe un recouvrement ouvert affine $X = \bigcup U_i$ tel que
chaque $\mathcal{O}_X(U_i)$ soit de Jacobson.
\item Il existe un recouvrement ouvert $X = \bigcup X_j$
tel que chaque sous-schéma ouvert $X_j$ soit de Jacobson.
\end{enumerate}
De plus, si $X$ est de Jacobson, tout sous-schéma ouvert
est de Jacobson.
\end{lemma}

\begin{proof}
La dernière assertion du lemme résulte de
Topologie, lemme \ref{topology-lemma-jacobson-inherited}.
L'équivalence de (5) et de (1) est donnée
par Topologie, lemme \ref{topology-lemma-jacobson-local}.
Ainsi, à l'aide du lemme \ref{properties-lemma-affine-jacobson},
on voit que (1) $\Leftrightarrow$ (2).
Pour achever la démonstration, il suffit de montrer qu'être de Jacobson est
une propriété locale des anneaux, voir le lemme \ref{properties-lemma-locally-P}.
Toute localisation d'un anneau de Jacobson en un élément est de Jacobson,
voir Algèbre, lemme \ref{algebra-lemma-Jacobson-invert-element}.
Soit $R$ un anneau ; supposons que $f_1, \ldots, f_n \in R$ engendrent l'idéal
unité et que chaque $R_{f_i}$ soit de Jacobson. Alors
$\Spec(R) = \bigcup D(f_i)$ est une réunion
d'ouverts tous de Jacobson ; donc $\Spec(R)$ est de Jacobson,
de nouveau par Topologie, lemme \ref{topology-lemma-jacobson-local}. Cela
démontre la seconde condition de la définition \ref{properties-definition-property-local}.
\end{proof}

\noindent
Beaucoup de schémas couramment utilisés en géométrie algébrique sont de Jacobson,
voir Morphismes, lemme \ref{morphisms-lemma-ubiquity-Jacobson-schemes}.
Mentionnons ici le cas intéressant suivant.

\begin{lemma}
\label{properties-lemma-complement-closed-point-Jacobson}
Exemples de schémas noethériens de Jacobson.
\begin{enumerate}
\item Si $(R, \mathfrak m)$ est un anneau local noethérien,
alors le spectre épointé $\Spec(R) \setminus \{\mathfrak m\}$
est un schéma de Jacobson.
\item Si $R$ est un anneau noethérien de radical de Jacobson $\text{rad}(R)$,
alors $\Spec(R) \setminus V(\text{rad}(R))$ est un schéma de Jacobson.
\item Si $(R, I)$ est un couple de Zariski (Compléments sur l'algèbre, définition
\ref{more-algebra-definition-zariski-pair}),
avec $R$ noethérien, alors $\Spec(R) \setminus V(I)$ est
un schéma de Jacobson.
\end{enumerate}
\end{lemma}

\begin{proof}
Démonstration de (3). Remarquons que $\Spec(R) - V(I)$ est recouvert par les ouverts
affines $\Spec(R_f)$ pour $f \in I$. Les anneaux $R_f$ sont de Jacobson
d'après Compléments sur l'algèbre, lemme
\ref{more-algebra-lemma-noetherian-zariski-jacobson-complement}.
Donc $\Spec(R) \setminus V(I)$ est de Jacobson par
le lemme \ref{properties-lemma-locally-jacobson}.
Les assertions (1) et (2) sont des cas particuliers de (3).

\medskip\noindent
Démonstration directe de (1).
Puisque $\Spec(R)$ est un schéma noethérien,
$S$ est un schéma noethérien (lemme \ref{properties-lemma-locally-closed-in-Noetherian}).
Ainsi $S$ est un espace topologique sobre et noethérien
(utiliser Schémas, lemme \ref{schemes-lemma-scheme-sober}).
Supposons que $S$ ne soit pas de Jacobson et
cherchons une contradiction.
D'après Topologie, lemme \ref{topology-lemma-non-jacobson-Noetherian-characterize},
il existe un point non fermé $\xi \in S$
tel que $\{\xi\}$ soit localement fermé. Il correspond à un
idéal premier $\mathfrak p \subset R$ tel que (1) il existe un idéal
premier $\mathfrak q$, $\mathfrak p \subset \mathfrak q \subset \mathfrak m$,
les deux inclusions étant strictes, et (2) $\{\mathfrak p\}$ soit ouvert dans
$\Spec(R/\mathfrak p)$. C'est impossible d'après Algèbre,
lemme \ref{algebra-lemma-Noetherian-local-domain-dim-2-infinite-opens}.
\end{proof}







\section{Schémas normaux}
\label{properties-section-normal}

\noindent
Rappelons qu'un anneau $R$ est dit normal si tous ses anneaux locaux sont des
anneaux intègres
normaux, voir Algèbre, définition \ref{algebra-definition-ring-normal}.
Un anneau intègre normal est un anneau intègre intégralement clos dans son
corps des fractions,
voir Algèbre, définition \ref{algebra-definition-domain-normal}.
Il est donc naturel de définir un schéma normal comme suit.

\begin{definition}
\label{properties-definition-normal}
Un schéma $X$ est {\it normal} si et seulement si, pour tout $x \in X$, l'anneau local
$\mathcal{O}_{X, x}$ est un anneau intègre normal.
\end{definition}

\noindent
Cela semble être la définition utilisée dans EGA, voir \cite[0, 4.1.4]{EGA}.
Supposons $X = \Spec(A)$, avec $A$ réduit. Dire que $X$ est
normal n'équivaut pas à dire que $A$ est intégralement clos dans son anneau
total des fractions. Cette équivalence est toutefois vraie si $A$ est noethérien (voir
Algèbre, lemme \ref{algebra-lemma-characterize-reduced-ring-normal}).

\begin{lemma}
\label{properties-lemma-locally-normal}
Soit $X$ un schéma. Les assertions suivantes sont équivalentes :
\begin{enumerate}
\item Le schéma $X$ est normal.
\item Pour tout ouvert affine $U \subset X$, l'anneau $\mathcal{O}_X(U)$
est normal.
\item Il existe un recouvrement ouvert affine $X = \bigcup U_i$ tel que
chaque $\mathcal{O}_X(U_i)$ soit normal.
\item Il existe un recouvrement ouvert $X = \bigcup X_j$
tel que chaque sous-schéma ouvert $X_j$ soit normal.
\end{enumerate}
De plus, si $X$ est normal, tout sous-schéma ouvert
est normal.
\end{lemma}

\begin{proof}
Cela résulte immédiatement des définitions.
\end{proof}

\begin{lemma}
\label{properties-lemma-normal-reduced}
Un schéma normal est réduit.
\end{lemma}

\begin{proof}
Cela résulte immédiatement des définitions.
\end{proof}

\begin{lemma}
\label{properties-lemma-integral-normal}
Soit $X$ un schéma intègre.
Alors $X$ est normal si et seulement si, pour tout ouvert affine non vide
$U \subset X$, l'anneau $\mathcal{O}_X(U)$ est un anneau intègre normal.
\end{lemma}

\begin{proof}
Cela résulte de
Algèbre, lemme \ref{algebra-lemma-normality-is-local}.
\end{proof}

\begin{lemma}
\label{properties-lemma-normal-locally-finite-nr-irreducibles}
Soit $X$ un schéma dont tout ouvert quasi-compact possède un nombre fini de
composantes irréductibles. Les assertions suivantes sont équivalentes :
\begin{enumerate}
\item $X$ est normal ;
\item $X$ est une réunion disjointe de schémas intègres normaux.
\end{enumerate}
\end{lemma}

\begin{proof}
Il résulte immédiatement des définitions que (2) entraîne (1).
Soit $X$ un schéma normal dont tout ouvert quasi-compact
possède un nombre fini de composantes irréductibles.
Si $X$ est affine, $X$ satisfait (2)
d'après Algèbre, lemme \ref{algebra-lemma-characterize-reduced-ring-normal}.
Pour $X$ quelconque, soit $X = \bigcup X_i$ un
recouvrement ouvert affine. Remarquons que chaque $X_i$ n'a
lui aussi qu'un nombre fini de composantes irréductibles et que le lemme
s'applique à chaque $X_i$. Soit $T \subset X$ une composante irréductible.
D'après le cas affine, chaque intersection $T \cap X_i$ est ouverte dans $X_i$
et est un schéma intègre normal.
Ainsi $T \subset X$ est ouvert et est un schéma intègre
normal. Cela prouve que $X$ est la réunion disjointe de ses composantes irréductibles,
qui sont des schémas intègres normaux.
\end{proof}

\begin{lemma}
\label{properties-lemma-normal-Noetherian}
Soit $X$ un schéma noethérien. Les
assertions suivantes sont équivalentes :
\begin{enumerate}
\item $X$ est normal ;
\item $X$ est une réunion disjointe finie de schémas intègres normaux.
\end{enumerate}
\end{lemma}

\begin{proof}
C'est un cas particulier du
lemme \ref{properties-lemma-normal-locally-finite-nr-irreducibles}, car l'espace
topologique sous-jacent à un schéma noethérien est noethérien
(lemme \ref{properties-lemma-Noetherian-topology}
et
Topologie, lemme \ref{topology-lemma-Noetherian}).
\end{proof}

\begin{lemma}
\label{properties-lemma-normal-locally-Noetherian}
Soit $X$ un schéma localement noethérien. Les
assertions suivantes sont équivalentes :
\begin{enumerate}
\item $X$ est normal ;
\item $X$ est une réunion disjointe de schémas intègres normaux.
\end{enumerate}
\end{lemma}

\begin{proof}
Omis. Indication : cela se déduit du
lemme \ref{properties-lemma-normal-Noetherian}par un raisonnement purement topologique.
\end{proof}

\begin{remark}
\label{properties-remark-normal-connected-irreducible}
Soit $X$ un schéma normal. Si $X$ est localement noethérien,
$X$ est intègre si et seulement si $X$ est connexe, voir le
lemme \ref{properties-lemma-normal-locally-Noetherian}.
Mais il existe un schéma affine connexe $X$ tel que
$\mathcal{O}_{X, x}$ soit intègre pour tout $x \in X$, sans que $X$ soit
irréductible, voir Exemples, section
\ref{examples-section-connected-locally-integral-not-integral}.
Cet exemple est même un schéma normal (démonstration omise) ; il faut donc y prendre garde !
\end{remark}

\begin{lemma}
\label{properties-lemma-normal-integral-sections}
\begin{slogan}
L'anneau des fonctions sur un schéma normal est normal.
\end{slogan}
Soit $X$ un schéma intègre normal.
Alors $\Gamma(X, \mathcal{O}_X)$ est un anneau intègre normal.
\end{lemma}

\begin{proof}
Posons $R = \Gamma(X, \mathcal{O}_X)$.
Il est clair que $R$ est un anneau intègre.
Supposons que $f = a/b$ soit un élément de son corps des
fractions entier sur $R$. Écrivons
$f^d + \sum_{i = 0, \ldots, d - 1} a_i f^i = 0$, avec
$a_i \in R$. Soit $U \subset X$ un ouvert affine non vide.
Puisque $b \in R$ est non nul et que $X$ est
intègre, $b|_U \in \mathcal{O}_X(U)$ est également non
nul. Ainsi $a/b$ est un élément du corps des fractions de
$\mathcal{O}_X(U)$ entier sur $\mathcal{O}_X(U)$
(car on peut utiliser la même équation polynomiale
$f^d + \sum_{i = 0, \ldots, d - 1} a_i|_U f^i = 0$ sur $U$).
Comme $\mathcal{O}_X(U)$ est un anneau intègre
normal (lemme \ref{properties-lemma-integral-normal}), on obtient
$f_U = (a|_U)/(b|_U) \in \mathcal{O}_X(U)$. Il est clair
que $f_U|_V = f_V$ dès que $V \subset U \subset X$ sont
des ouverts affines non vides. Les sections locales $f_U$ se recollent donc en un élément
$g \in R = \Gamma(X, \mathcal{O}_X)$. Alors $bg$ et $a$
se restreignent au même élément de $\mathcal{O}_X(U)$ pour
tout $U$ comme ci-dessus ; donc $bg = a$. Autrement dit, $g$ a pour image $f$
dans le corps des fractions de $R$.
\end{proof}










\section{Schémas de Cohen--Macaulay}
\label{properties-section-Cohen-Macaulay}

\noindent
Rappelons, voir Algèbre, définition \ref{algebra-definition-local-ring-CM},
qu'un anneau local noethérien $(R, \mathfrak m)$ est
dit de Cohen--Macaulay si $\text{depth}_{\mathfrak m}(R) = \dim(R)$.
Rappelons qu'un anneau noethérien $R$ est dit de Cohen--Macaulay si
tout anneau local $R_{\mathfrak p}$ de $R$ est de
Cohen--Macaulay, voir Algèbre, définition \ref{algebra-definition-ring-CM}.

\begin{definition}
\label{properties-definition-Cohen-Macaulay}
Soit $X$ un schéma. On dit que $X$ est {\it de Cohen--Macaulay} si,
pour tout $x \in X$, il existe un voisinage ouvert affine
$U \subset X$ de $x$ tel que l'anneau $\mathcal{O}_X(U)$ soit
noethérien et de Cohen--Macaulay.
\end{definition}

\begin{lemma}
\label{properties-lemma-characterize-Cohen-Macaulay}
Soit $X$ un schéma. Les assertions suivantes sont équivalentes :
\begin{enumerate}
\item $X$ est de Cohen--Macaulay ;
\item $X$ est localement noethérien et tous ses anneaux locaux sont de
Cohen--Macaulay ;
\item $X$ est localement noethérien et, pour tout point fermé $x \in X$,
l'anneau local $\mathcal{O}_{X, x}$ est de Cohen--Macaulay.
\end{enumerate}
\end{lemma}

\begin{proof}
Algèbre, lemme \ref{algebra-lemma-localize-CM}, affirme que toute localisation d'un
anneau local de Cohen--Macaulay est de Cohen--Macaulay. Le lemme s'en déduit en
combinant ce résultat avec le lemme \ref{properties-lemma-locally-Noetherian},
l'existence de points fermés
sur les schémas localement noethériens
(lemme \ref{properties-lemma-locally-Noetherian-closed-point}) et
les définitions.
\end{proof}

\begin{lemma}
\label{properties-lemma-locally-Cohen-Macaulay}
Soit $X$ un schéma. Les assertions suivantes sont équivalentes :
\begin{enumerate}
\item Le schéma $X$ est de Cohen--Macaulay.
\item Pour tout ouvert affine $U \subset X$, l'anneau $\mathcal{O}_X(U)$
est noethérien et de Cohen--Macaulay.
\item Il existe un recouvrement ouvert affine $X = \bigcup U_i$ tel que
chaque $\mathcal{O}_X(U_i)$ soit noethérien et de Cohen--Macaulay.
\item Il existe un recouvrement ouvert $X = \bigcup X_j$
tel que chaque sous-schéma ouvert $X_j$ soit de Cohen--Macaulay.
\end{enumerate}
De plus, si $X$ est de Cohen--Macaulay, tout sous-schéma ouvert
est de Cohen--Macaulay.
\end{lemma}

\begin{proof}
Combiner les lemmes \ref{properties-lemma-locally-Noetherian}
et \ref{properties-lemma-characterize-Cohen-Macaulay}.
\end{proof}

\noindent
On trouvera davantage de résultats sur les schémas de Cohen--Macaulay et la profondeur
dans Cohomologie des schémas, section \ref{coherent-section-depth}.








\section{Schémas réguliers}
\label{properties-section-regular}

\noindent
Rappelons, voir Algèbre, définition \ref{algebra-definition-regular-local},
qu'un anneau local noethérien $(R, \mathfrak m)$ est
dit {\it régulier} si $\mathfrak m$ peut être engendré
par $\dim(R)$ éléments.
Rappelons qu'un anneau noethérien $R$ est dit {\it régulier} si
tout anneau local $R_{\mathfrak p}$ de $R$ est
régulier, voir Algèbre, définition \ref{algebra-definition-regular}.

\begin{definition}
\label{properties-definition-regular}
Soit $X$ un schéma. On dit que $X$ est {\it régulier}, ou {\it non singulier},
si, pour tout $x \in X$, il existe un voisinage ouvert affine
$U \subset X$ de $x$ tel que l'anneau $\mathcal{O}_X(U)$ soit
noethérien et régulier.
\end{definition}

\begin{lemma}
\label{properties-lemma-characterize-regular}
Soit $X$ un schéma. Les assertions suivantes sont équivalentes :
\begin{enumerate}
\item $X$ est régulier ;
\item $X$ est localement noethérien et tous ses anneaux locaux sont
réguliers ;
\item $X$ est localement noethérien et, pour tout point fermé $x \in X$,
l'anneau local $\mathcal{O}_{X, x}$ est régulier.
\end{enumerate}
\end{lemma}

\begin{proof}
D'après la discussion précédant, dans le chapitre Algèbre, la définition
\ref{algebra-definition-regular}, toute localisation d'un anneau
local régulier est régulière. Le lemme s'en déduit en
combinant ce résultat avec le lemme \ref{properties-lemma-locally-Noetherian},
l'existence de points fermés
sur les schémas localement noethériens
(lemme \ref{properties-lemma-locally-Noetherian-closed-point}) et
les définitions.
\end{proof}

\begin{lemma}
\label{properties-lemma-locally-regular}
Soit $X$ un schéma. Les assertions suivantes sont équivalentes :
\begin{enumerate}
\item Le schéma $X$ est régulier.
\item Pour tout ouvert affine $U \subset X$, l'anneau $\mathcal{O}_X(U)$
est noethérien et régulier.
\item Il existe un recouvrement ouvert affine $X = \bigcup U_i$ tel que
chaque $\mathcal{O}_X(U_i)$ soit noethérien et régulier.
\item Il existe un recouvrement ouvert $X = \bigcup X_j$
tel que chaque sous-schéma ouvert $X_j$ soit régulier.
\end{enumerate}
De plus, si $X$ est régulier, tout sous-schéma ouvert est régulier.
\end{lemma}

\begin{proof}
Combiner les lemmes \ref{properties-lemma-locally-Noetherian}
et \ref{properties-lemma-characterize-regular}.
\end{proof}

\begin{lemma}
\label{properties-lemma-regular-normal}
Un schéma régulier est normal.
\end{lemma}

\begin{proof}
Voir
Algèbre, lemme \ref{algebra-lemma-regular-normal}.
\end{proof}




\section{Dimension}
\label{properties-section-dimension}

\noindent
La dimension d'un schéma est simplement la dimension de son espace
topologique sous-jacent.

\begin{definition}
\label{properties-definition-dimension}
Soit $X$ un schéma.
\begin{enumerate}
\item La {\it dimension} de $X$ est la dimension de $X$
en tant qu'espace topologique,
voir Topologie, définition \ref{topology-definition-Krull}.
\item Pour $x \in X$, on note $\dim_x(X)$ la dimension de l'espace
topologique sous-jacent à $X$ en $x$, au
sens de Topologie, définition \ref{topology-definition-Krull}.
On appelle $\dim_x(X)$ la {\it dimension de $X$ en $x$}.
\end{enumerate}
\end{definition}

\noindent
L'espace topologique sous-jacent à un schéma étant sobre
(Schémas, lemme \ref{schemes-lemma-scheme-sober}),
on peut calculer la dimension de $X$ comme la borne supérieure des longueurs $n$
des chaînes
$$
T_0 \subset T_1 \subset \ldots \subset T_n
$$
de fermés irréductibles de $X$, ou comme la borne supérieure des longueurs $n$
des chaînes de spécialisations
$$
\xi_n \leadsto \xi_{n - 1} \leadsto \ldots \leadsto \xi_0
$$
de points de $X$.

\begin{lemma}
\label{properties-lemma-dimension}
Soit $X$ un schéma. Les quantités suivantes sont égales :
\begin{enumerate}
\item la dimension de $X$;
\item la borne supérieure des dimensions des anneaux locaux de $X$;
\item la borne supérieure des $\dim_x(X)$ pour $x \in X$.
\end{enumerate}
\end{lemma}

\begin{proof}
Remarquons qu'étant donnée une chaîne de spécialisations
$$
\xi_n \leadsto \xi_{n - 1} \leadsto \ldots \leadsto \xi_0
$$
de points de $X$, tous les points $\xi_i$ correspondent à des idéaux
premiers de l'anneau local de $X$ en $\xi_0$,
d'après Schémas, lemme \ref{schemes-lemma-specialize-points}.
On voit donc que la dimension de $X$ est la borne supérieure des dimensions
de ses anneaux locaux. En particulier, $\dim_x(X) \geq \dim(\mathcal{O}_{X, x})$,
puisque $\dim_x(X)$ est le minimum des dimensions des voisinages ouverts de
$x$. Ainsi $\sup_{x \in X} \dim_x(X) \geq \dim(X)$. D'autre part, il
est clair que $\sup_{x \in X} \dim_x(X) \leq \dim(X)$,
car $\dim(U) \leq \dim(X)$ pour tout ouvert de $X$.
\end{proof}

\begin{lemma}
\label{properties-lemma-codimension-local-ring}
Soient $X$ un schéma et $Y \subset X$ un sous-ensemble fermé
irréductible. Soit $\xi \in Y$ son point générique. Alors
$$
\text{codim}(Y, X) = \dim(\mathcal{O}_{X, \xi})
$$
où la codimension est celle de
Topologie, définition \ref{topology-definition-codimension}.
\end{lemma}

\begin{proof}
D'après Topologie, lemme \ref{topology-lemma-codimension-at-generic-point},
on peut remplacer $X$ par un voisinage ouvert affine de $\xi$. Dans ce
cas, le résultat découle aisément de
Algèbre, lemme \ref{algebra-lemma-irreducible-components-containing-x}.
\end{proof}

\begin{lemma}
\label{properties-lemma-generic-point}
Soient $X$ un schéma et $x \in X$. Alors $x$ est un point générique
d'une composante irréductible de $X$ si et seulement si $\dim(\mathcal{O}_{X, x}) = 0$.
\end{lemma}

\begin{proof}
Cela résulte, par exemple, du lemme \ref{properties-lemma-codimension-local-ring}.
\end{proof}

\begin{lemma}
\label{properties-lemma-locally-Noetherian-dimension-0}
Un schéma localement noethérien de dimension $0$ est une réunion
disjointe de spectres d'anneaux locaux artiniens.
\end{lemma}

\begin{proof}
Un anneau noethérien de dimension $0$ est un produit fini d'anneaux locaux
artiniens,
voir Algèbre, proposition \ref{algebra-proposition-dimension-zero-ring}.
L'espace topologique sous-jacent à un ouvert affine d'un schéma localement noethérien $X$ de dimension
$0$ est donc discret. Il en résulte que la
topologie de $X$ est discrète. Le lemme se déduit aisément de ces
remarques.
\end{proof}

\begin{lemma}
\label{properties-lemma-dimension-zero}
\begin{reference}
Courriel d'Ofer Gabber du 4 juin 2016
\end{reference}
Soit $X$ un schéma de dimension zéro. Les assertions suivantes sont équivalentes :
\begin{enumerate}
\item $X$ est quasi-séparé ;
\item $X$ est séparé ;
\item $X$ est un espace topologique séparé ;
\item tout ouvert affine est fermé.
\end{enumerate}
Dans ce cas, les composantes connexes de $X$ sont des points et tout
ouvert quasi-compact de $X$ est affine. En particulier, si $X$
est quasi-compact, $X$ est affine.
\end{lemma}

\begin{proof}
Comme $X$ est de dimension zéro, pour tout ouvert affine
$U \subset X$, l'espace $U$ est profini et possède
diverses autres propriétés que nous utiliserons librement ci-dessous, voir
Algèbre, lemme \ref{algebra-lemma-ring-with-only-minimal-primes}.
Choisissons un recouvrement ouvert affine $X = \bigcup U_i$.

\medskip\noindent
Si (4) est satisfaite, $U_i \cap U_j$ est un fermé de
$U_i$, donc est quasi-compact. Par conséquent, $X$ est quasi-séparé,
d'après Schémas, lemme \ref{schemes-lemma-characterize-quasi-separated},
et (1) est satisfaite.

\medskip\noindent
Si (1) est satisfaite, $U_i \cap U_j$ est un ouvert quasi-compact
de $U_i$, donc est fermé dans $U_i$. Alors $U_i \cap U_j \to U_i$
est une immersion ouverte d'image fermée, et donc une
immersion fermée. En particulier, $U_i \cap U_j$ est affine
et $\mathcal{O}(U_i) \to \mathcal{O}_X(U_i \cap U_j)$ est surjectif.
Ainsi $X$ est séparé
d'après Schémas, lemme \ref{schemes-lemma-characterize-separated},
et (2) est satisfaite.

\medskip\noindent
Supposons (2) et soient $x, y \in X$. Supposons $x \in U_i$. Si $y \in U_i$
également, on peut trouver des voisinages ouverts disjoints de $x$ et de $y$,
car $U_i$ est un espace topologique séparé. Supposons $y \not \in U_i$ et $y \in U_j$.
Alors $y \not \in U_i \cap U_j$, qui est un ouvert affine de $U_j$,
donc fermé dans $U_j$. On peut ainsi trouver un voisinage ouvert
de $y$ ne rencontrant pas $U_i$, et l'on en déduit que $X$ est un espace topologique
séparé ; donc (3) est satisfaite.

\medskip\noindent
Supposons (3). Soit $U \subset X$ un ouvert affine.
Alors $U$ est fermé dans $X$ d'après Topologie, lemme
\ref{topology-lemma-quasi-compact-in-Hausdorff}.
Cela démontre (4).

\medskip\noindent
Supposons que $X$ satisfasse les conditions équivalentes (1) -- (4).
Démontrons les dernières assertions du lemme. Soient $x, y \in X$
avec $x \not = y$. Puisque $y$ ne se spécialise pas en $x$, on peut choisir
un ouvert affine $U \subset X$ avec $x \in U$ et $y \not \in U$.
Alors $X = U \amalg (X \setminus U)$ est une décomposition
en parties ouvertes et fermées, ce qui montre que $x$ et $y$ n'appartiennent
pas à la même composante connexe de $X$. Supposons maintenant que $U \subset X$
soit un ouvert quasi-compact. Écrivons $U = U_1 \cup \ldots \cup U_n$ comme
réunion d'ouverts affines. Montrons par récurrence sur $n$ que $U$ est
affine. On se ramène immédiatement au cas $n = 2$. Dans ce cas, on a
$U = (U_1 \setminus U_2) \amalg (U_1 \cap U_2) \amalg (U_2 \setminus U_1)$,
et les arguments précédents montrent que chacune de ces parties est affine.
\end{proof}

\begin{lemma}
\label{properties-lemma-isolated-dim-0-locally-noetherian}
Soit $x$ un point d'un schéma localement noethérien $X$.
Alors $\dim_x(X) = 0$ si et seulement si $x$ est un point isolé de $X$.
\end{lemma}

\begin{proof}
Si $x$ est un point isolé, $\{x\}$ est un ouvert de $X$
(Topologie, définition \ref{topology-definition-isolated-point}),
et donc $\dim_x(X) = \dim(\{x\}) = 0$ par définition. Réciproquement, si
$\dim_x(X) = 0$, il existe un voisinage ouvert $U \subset X$
de $x$ tel que $\dim(U) = 0$
(Topologie, définition \ref{topology-definition-Krull}).
D'après le lemme \ref{properties-lemma-locally-Noetherian-dimension-0},
la topologie de $U$ est discrète, et donc
$x$ est un point isolé.
\end{proof}





\section{Schémas caténaires}
\label{properties-section-catenary}

\noindent
Rappelons qu'un espace topologique $X$ est dit {\it caténaire} si,
pour tout couple de fermés irréductibles $T \subset T'$,
il existe une chaîne maximale de fermés irréductibles
$$
T = T_0 \subset T_1 \subset \ldots \subset T_e = T'
$$
et si toutes ces chaînes ont la même longueur. Voir
Topologie, définition \ref{topology-definition-catenary}.

\begin{definition}
\label{properties-definition-catenary}
Soit $S$ un schéma. On dit que $S$ est {\it caténaire} si
l'espace topologique sous-jacent à $S$ est caténaire.
\end{definition}

\noindent
Rappelons qu'un anneau $A$ est dit {\it caténaire} si,
pour tout couple d'idéaux premiers $\mathfrak p \subset \mathfrak q$,
il existe une chaîne maximale d'idéaux premiers
$$
\mathfrak p =
\mathfrak p_0 \subset \ldots \subset \mathfrak p_e
= \mathfrak q
$$
et si toutes ces chaînes ont la même longueur.
Voir Algèbre, définition \ref{algebra-definition-catenary}.

\begin{lemma}
\label{properties-lemma-catenary-local}
Soit $S$ un schéma. Les assertions suivantes sont équivalentes :
\begin{enumerate}
\item $S$ est caténaire ;
\item il existe un recouvrement ouvert de $S$ dont tous les membres sont des
schémas caténaires ;
\item pour tout ouvert affine $\Spec(R) = U \subset S$, l'anneau
$R$ est caténaire ;
\item il existe un recouvrement ouvert affine $S = \bigcup U_i$ tel
que chaque $U_i$ soit le spectre d'un anneau caténaire.
\end{enumerate}
De plus, dans ce cas, tout sous-schéma localement fermé de $S$ est lui aussi
caténaire.
\end{lemma}

\begin{proof}
Combiner Topologie, lemme \ref{topology-lemma-catenary}, et
Algèbre, lemme \ref{algebra-lemma-catenary}.
\end{proof}

\begin{lemma}
\label{properties-lemma-catenary-dimension-function}
Soit $S$ un schéma localement noethérien. Les
assertions suivantes sont équivalentes :
\begin{enumerate}
\item $S$ est caténaire ;
\item il existe, localement pour la topologie de Zariski, une fonction de dimension
sur $S$ (voir Topologie, définition \ref{topology-definition-dimension-function}).
\end{enumerate}
\end{lemma}

\begin{proof}
Cela résulte des résultats
suivants : Topologie, lemmes
\ref{topology-lemma-catenary},
\ref{topology-lemma-dimension-function-catenary}et
\ref{topology-lemma-locally-dimension-function},
Schémas, lemme \ref{schemes-lemma-scheme-sober},
et enfin le lemme \ref{properties-lemma-Noetherian-topology}.
\end{proof}

\noindent
Un schéma est caténaire si et seulement si ses anneaux
locaux sont caténaires.

\begin{lemma}
\label{properties-lemma-catenary-local-rings-catenary}
Soit $X$ un schéma. Les assertions suivantes sont équivalentes :
\begin{enumerate}
\item $X$ est caténaire ;
\item pour tout $x \in X$, l'anneau local $\mathcal{O}_{X, x}$ est
caténaire.
\end{enumerate}
\end{lemma}

\begin{proof}
Supposons $X$ caténaire. Soit $x \in X$. D'après le lemme \ref{properties-lemma-catenary-local},
on peut remplacer $X$ par un voisinage ouvert affine de $x$;
alors $\Gamma(X, \mathcal{O}_X)$ est un anneau caténaire.
D'après Algèbre, lemme \ref{algebra-lemma-localization-catenary},
toute localisation d'un anneau caténaire
est caténaire. Donc $\mathcal{O}_{X, x}$ est caténaire.

\medskip\noindent
Réciproquement, supposons tous les anneaux locaux de $X$ caténaires.
Soit $Y \subset Y'$ une inclusion de fermés
irréductibles de $X$. Soit $\xi \in Y$ le point générique.
Soit $\mathfrak p \subset \mathcal{O}_{X, \xi}$ l'idéal
premier correspondant au point générique de $Y'$, voir
Schémas, lemme \ref{schemes-lemma-specialize-points}. D'après ce
même lemme, les fermés irréductibles de $X$ compris entre $Y$ et $Y'$
correspondent aux idéaux premiers $\mathfrak q \subset \mathcal{O}_{X, \xi}$
tels que $\mathfrak p \subset \mathfrak q \subset \mathfrak m_{\xi}$.
Toutes leurs chaînes maximales sont donc finies et de même
longueur, puisque $\mathcal{O}_{X, \xi}$ est un anneau caténaire.
\end{proof}






\section{Conditions de Serre}
\label{properties-section-Rk}

\noindent
Voici deux notions techniques souvent utiles. Voir
aussi Cohomologie des schémas, section \ref{coherent-section-depth}.

\begin{definition}
\label{properties-definition-Rk}
Soit $X$ un schéma localement noethérien. Soit $k \geq 0$.
\begin{enumerate}
\item On dit que $X$ est {\it régulier en codimension $k$},
ou que $X$ possède la propriété {\it $(R_k)$}, si, pour tout $x \in X$,
on a
$$
\dim(\mathcal{O}_{X, x}) \leq k
\Rightarrow
\mathcal{O}_{X, x}\text{ est régulier}
$$
\item On dit que $X$ possède la propriété {\it $(S_k)$} si, pour tout $x \in X$, on a
$\text{depth}(\mathcal{O}_{X, x}) \geq \min(k, \dim(\mathcal{O}_{X, x}))$.
\end{enumerate}
\end{definition}

\noindent
L'expression « régulier en codimension $k$» a un sens, car nous avons vu, dans
la section \ref{properties-section-catenary}, que, si $Y \subset X$ est un fermé
irréductible de point générique $x$, alors
$\dim(\mathcal{O}_{X, x}) = \text{codim}(Y, X)$. Par exemple, la condition
$(R_0)$ signifie que, pour tout point générique $\eta \in X$ d'une composante
irréductible de $X$, l'anneau local $\mathcal{O}_{X, \eta}$ est un
corps. Mais, pour un schéma noethérien quelconque, le lieu régulier
de $X$ peut mal se comporter ; il faut donc être prudent.

\begin{lemma}
\label{properties-lemma-scheme-regular-iff-all-Rk}
Soit $X$ un schéma localement noethérien.
Alors $X$ est régulier si et seulement si $X$ satisfait $(R_k)$ pour tout $k \geq 0$.
\end{lemma}

\begin{proof}
Cela résulte du lemme \ref{properties-lemma-characterize-regular} et des définitions.
\end{proof}

\begin{lemma}
\label{properties-lemma-scheme-CM-iff-all-Sk}
Soit $X$ un schéma localement noethérien.
Alors $X$ est de Cohen--Macaulay si et seulement si $X$ satisfait $(S_k)$ pour tout $k \geq 0$.
\end{lemma}

\begin{proof}
Le lemme \ref{properties-lemma-characterize-Cohen-Macaulay}
permet de se ramener aux anneaux locaux.
Le résultat vient donc de ce qu'un anneau local noethérien est de Cohen--Macaulay
si et seulement si sa profondeur est égale à sa dimension.
\end{proof}

\begin{lemma}
\label{properties-lemma-criterion-reduced}
Soit $X$ un schéma localement noethérien.
Alors $X$ est réduit si et seulement si $X$ possède les propriétés $(S_1)$ et $(R_0)$.
\end{lemma}

\begin{proof}
C'est le lemme \ref{algebra-lemma-criterion-reduced}du chapitre Algèbre.
\end{proof}

\begin{lemma}
\label{properties-lemma-criterion-normal}
Soit $X$ un schéma localement noethérien.
Alors $X$ est normal si et seulement si $X$ possède les propriétés $(S_2)$ et $(R_1)$.
\end{lemma}

\begin{proof}
C'est le lemme \ref{algebra-lemma-criterion-normal}du chapitre Algèbre.
\end{proof}

\begin{lemma}
\label{properties-lemma-normal-dimension-1-regular}
Soit $X$ un schéma localement noethérien, normal et
de dimension $\leq 1$. Alors $X$ est régulier.
\end{lemma}

\begin{proof}
Cela résulte du lemme \ref{properties-lemma-criterion-normal} et des définitions.
\end{proof}

\begin{lemma}
\label{properties-lemma-normal-dimension-2-Cohen-Macaulay}
Soit $X$ un schéma localement noethérien, normal et
de dimension $\leq 2$. Alors $X$ est de Cohen--Macaulay.
\end{lemma}

\begin{proof}
Cela résulte du lemme \ref{properties-lemma-criterion-normal} et des définitions.
\end{proof}







\section{Schémas japonais et schémas de Nagata}
\label{properties-section-nagata}

\noindent
Les notions considérées dans cette section ne font pas l'objet de définitions explicites bien mises en
évidence dans EGA. Un « schéma universellement japonais » est mentionné et défini dans
\cite[IV Corollary 5.11.4]{EGA}. Un « schéma japonais » est mentionné dans
\cite[IV remarque 10.4.14 (ii)]{EGA}, mais sans définition. La notion
de schéma de Nagata (définie ci-dessous) apparaît à
quelques endroits dans la littérature (voir, par exemple, \cite[définition 8.2.30]{Liu} et
\cite[Page 142]{Greco}).

\medskip\noindent
Rappelons brièvement qu'un anneau intègre $R$ est dit {\it japonais} si la fermeture
intégrale de $R$ dans toute extension finie de son corps des fractions est finie sur
$R$. Un anneau $R$ est dit {\it universellement japonais} si, pour tout homomorphisme d'anneaux
de type fini $R \to S$ avec $S$ intègre, $S$ est japonais. Un anneau $R$ est dit
{\it de Nagata} s'il est noethérien et si $R/\mathfrak p$ est japonais pour tout idéal
premier $\mathfrak p$ de $R$.

\begin{definition}
\label{properties-definition-nagata}
Soit $X$ un schéma.
\begin{enumerate}
\item Supposons $X$ intègre. On dit que $X$ est {\it japonais}
si, pour tout $x \in X$, il existe un
voisinage ouvert affine $x \in U \subset X$ tel que l'anneau
$\mathcal{O}_X(U)$ soit japonais
(voir Algèbre, définition \ref{algebra-definition-N}).
\item On dit que $X$ est {\it universellement japonais} si, pour tout $x \in X$,
il existe un voisinage ouvert affine $x \in U \subset X$ tel que
l'anneau $\mathcal{O}_X(U)$ soit universellement japonais
(voir Algèbre, définition \ref{algebra-definition-nagata}).
\item On dit que $X$ est {\it de Nagata} si, pour tout $x \in X$, il existe un
voisinage ouvert affine $x \in U \subset X$ tel que l'anneau
$\mathcal{O}_X(U)$ soit de Nagata
(voir Algèbre, définition \ref{algebra-definition-nagata}).
\end{enumerate}
\end{definition}

\noindent
Être de Nagata équivaut à être localement noethérien et
universellement japonais, voir le
lemme \ref{properties-lemma-nagata-universally-Japanese}.

\begin{remark}
\label{properties-remark-non-integral-Japanese}
Dans \cite{Hoobler-finite}, un schéma (localement noethérien) $X$ est dit
japonais si, pour tout $x \in X$ et tout idéal premier associé $\mathfrak p$
de $\mathcal{O}_{X, x}$, l'anneau $\mathcal{O}_{X, x}/\mathfrak p$ est
japonais. Nous n'utilisons pas cette définition, car il existe un anneau
intègre noethérien de dimension un dont les anneaux locaux sont excellents
(en particulier japonais), mais dont la normalisation n'est pas finie. Voir
\cite[exemple 1]{Hochster-loci}, \cite{Heinzer-Levy} ou
\cite[Exposé XIX]{Traveaux}.
On pourrait toutefois contourner ce problème en disant qu'un schéma
$X$ est japonais si, pour tout ouvert affine $\Spec(A) \subset X$, l'anneau
$A/\mathfrak p$ est japonais pour tout idéal premier associé $\mathfrak p$ de $A$.
\end{remark}

\begin{lemma}
\label{properties-lemma-nagata-locally-Noetherian}
Un schéma de Nagata est localement noethérien.
\end{lemma}

\begin{proof}
En effet, un anneau de Nagata est noethérien par définition.
\end{proof}

\begin{lemma}
\label{properties-lemma-locally-Japanese}
Soit $X$ un schéma intègre. Les assertions suivantes sont équivalentes :
\begin{enumerate}
\item Le schéma $X$ est japonais.
\item Pour tout ouvert affine $U \subset X$, l'anneau intègre $\mathcal{O}_X(U)$
est japonais.
\item Il existe un recouvrement ouvert affine $X = \bigcup U_i$
tel que chaque $\mathcal{O}_X(U_i)$ soit japonais.
\item Il existe un recouvrement ouvert $X = \bigcup X_j$
tel que chaque sous-schéma ouvert $X_j$ soit japonais.
\end{enumerate}
De plus, si $X$ est japonais, tout sous-schéma ouvert
est japonais.
\end{lemma}

\begin{proof}
Cela résulte du lemme \ref{properties-lemma-locally-P} et des
lemmes suivants du chapitre Algèbre : \ref{algebra-lemma-localize-N} et
\ref{algebra-lemma-Japanese-local}.
\end{proof}

\begin{lemma}
\label{properties-lemma-locally-universally-Japanese}
Soit $X$ un schéma. Les assertions suivantes sont équivalentes :
\begin{enumerate}
\item Le schéma $X$ est universellement japonais.
\item Pour tout ouvert affine $U \subset X$, l'anneau $\mathcal{O}_X(U)$
est universellement japonais.
\item Il existe un recouvrement ouvert affine $X = \bigcup U_i$
tel que chaque $\mathcal{O}_X(U_i)$ soit universellement japonais.
\item Il existe un recouvrement ouvert $X = \bigcup X_j$
tel que chaque sous-schéma ouvert $X_j$ soit universellement japonais.
\end{enumerate}
De plus, si $X$ est universellement japonais, tout sous-schéma ouvert
est universellement japonais.
\end{lemma}

\begin{proof}
Cela résulte du lemme \ref{properties-lemma-locally-P} et des
lemmes suivants du chapitre Algèbre : \ref{algebra-lemma-universally-japanese} et
\ref{algebra-lemma-nagata-local}.
\end{proof}

\begin{lemma}
\label{properties-lemma-locally-nagata}
Soit $X$ un schéma. Les assertions suivantes sont équivalentes :
\begin{enumerate}
\item Le schéma $X$ est de Nagata.
\item Pour tout ouvert affine $U \subset X$, l'anneau $\mathcal{O}_X(U)$
est de Nagata.
\item Il existe un recouvrement ouvert affine $X = \bigcup U_i$
tel que chaque $\mathcal{O}_X(U_i)$ soit de Nagata.
\item Il existe un recouvrement ouvert $X = \bigcup X_j$
tel que chaque sous-schéma ouvert $X_j$ soit de Nagata.
\end{enumerate}
De plus, si $X$ est de Nagata, tout sous-schéma ouvert est de Nagata.
\end{lemma}

\begin{proof}
Cela résulte du lemme \ref{properties-lemma-locally-P} et des
lemmes suivants du chapitre Algèbre : \ref{algebra-lemma-nagata-localize} et
\ref{algebra-lemma-nagata-local}.
\end{proof}

\begin{lemma}
\label{properties-lemma-characterize-nagata}
Soit $X$ un schéma localement noethérien.
Alors $X$ est de Nagata si et seulement si tout sous-schéma fermé intègre
$Z \subset X$ est japonais.
\end{lemma}

\begin{proof}
Supposons $X$ de Nagata. Soit $Z \subset X$ un sous-schéma fermé intègre.
Soit $z \in Z$.
Soit $\Spec(A) = U \subset X$ un ouvert affine contenant $z$
tel que $A$ soit de Nagata. Alors
$Z \cap U \cong \Spec(A/\mathfrak p)$ pour un certain idéal premier $\mathfrak p$,
voir Schémas, lemme \ref{schemes-lemma-closed-subspace-scheme} (et
la définition \ref{properties-definition-integral}).
D'après Algèbre, définition \ref{algebra-definition-nagata},
 $A/\mathfrak p$ est japonais. Donc $Z$ est japonais par définition.

\medskip\noindent
Supposons tout sous-schéma fermé intègre de $X$ japonais.
Soit $\Spec(A) = U \subset X$ un ouvert affine quelconque.
Comme $X$ est localement noethérien, $A$ est noethérien
(lemme \ref{properties-lemma-locally-Noetherian}). Soit $\mathfrak p \subset A$
un idéal premier. Il faut montrer que $A/\mathfrak p$ est japonais.
Soit $T \subset U$ le fermé $V(\mathfrak p) \subset \Spec(A)$.
Soit $\overline{T} \subset X$ son adhérence. Alors $\overline{T}$ est
irréductible, étant l'adhérence d'un sous-ensemble irréductible. Le
sous-schéma fermé réduit défini par $\overline{T}$ est donc un sous-schéma
fermé intègre (encore noté $\overline{T}$ ), voir
Schémas, lemme \ref{schemes-lemma-reduced-closed-subscheme}.
Autrement dit, $\Spec(A/\mathfrak p)$ est un ouvert
affine d'un sous-schéma fermé intègre de $X$. Ce sous-schéma est
japonais par hypothèse ; le lemme \ref{properties-lemma-locally-Japanese} montre donc que
$A/\mathfrak p$ est japonais.
\end{proof}

\begin{lemma}
\label{properties-lemma-nagata-universally-Japanese}
Soit $X$ un schéma. Les
assertions suivantes sont équivalentes :
\begin{enumerate}
\item $X$ est de Nagata ;
\item $X$ est localement noethérien et universellement japonais.
\end{enumerate}
\end{lemma}

\begin{proof}
C'est
Algèbre, proposition \ref{algebra-proposition-nagata-universally-japanese}.
\end{proof}

\noindent
Cette discussion sera poursuivie dans
Morphismes, section \ref{morphisms-section-nagata}.



\section{G-schémas}
\label{properties-section-G-scheme}

\noindent
Rappelons qu'un anneau $R$ est un G-anneau si $R$ est noethérien et
si, pour tout idéal premier $\mathfrak p$ de $R$, l'homomorphisme d'anneaux
$R_\mathfrak p \to (R_\mathfrak p)^\wedge$ est régulier.

\begin{definition}
\label{properties-definition-G-scheme}
Soit $X$ un schéma. On dit que $X$ est un {\it G-schéma}\footnote{Cette
terminologie n'est peut-être pas usuelle. Si $G$ est un groupe fini ou un schéma
en groupes, un $G$-schéma désigne parfois un schéma muni d'une
action de $G$.} si, pour tout $x \in X$, il existe un
voisinage ouvert affine $x \in U \subset X$ tel que l'anneau
$\mathcal{O}_X(U)$ soit un G-anneau (voir
Compléments sur l'algèbre, définition \ref{more-algebra-definition-G-ring}).
\end{definition}

\begin{lemma}
\label{properties-lemma-characterize-G-scheme}
Soit $X$ un schéma. Les assertions suivantes sont équivalentes :
\begin{enumerate}
\item $X$ est un G-schéma ;
\item $X$ est localement noethérien et, pour tout $x \in X$, l'homomorphisme
d'anneaux $\mathcal{O}_{X, x} \to \mathcal{O}_{X, x}^\wedge$ est régulier ;
\item $X$ est localement noethérien et, pour tout point fermé $x \in X$,
l'homomorphisme d'anneaux $\mathcal{O}_{X, x} \to \mathcal{O}_{X, x}^\wedge$ est régulier.
\end{enumerate}
\end{lemma}

\begin{proof}
Il est clair que (1) entraîne (2) et que (2) entraîne (3).
Supposons (3). Soit $x \in X$ un point quelconque. D'après
le lemme \ref{properties-lemma-locally-Noetherian-closed-point},
il existe une spécialisation $x \leadsto x'$ avec $x'$ fermé
dans $X$. Alors $\mathcal{O}_{X, x'}$ est un G-anneau
d'après Compléments sur l'algèbre, lemme \ref{more-algebra-lemma-check-G-ring-maximal-ideals}.
Comme $\mathcal{O}_{X, x}$ est la localisation de $\mathcal{O}_{X, x'}$
en un idéal premier (Schémas, lemme \ref{schemes-lemma-specialize-points}),
l'homomorphisme d'anneaux $\mathcal{O}_{X, x} \to \mathcal{O}_{X, x}^\wedge$
est régulier. Le point $x \in X$ étant quelconque, on en déduit
que, pour tout ouvert affine $U \subset X$, l'anneau
noethérien (lemme \ref{properties-lemma-locally-Noetherian}) $\mathcal{O}_X(U)$ est un
G-anneau. Cela démontre (1).
\end{proof}

\begin{lemma}
\label{properties-lemma-locally-G-scheme}
Soit $X$ un schéma. Les assertions suivantes sont équivalentes :
\begin{enumerate}
\item Le schéma $X$ est un G-schéma.
\item Pour tout ouvert affine $U \subset X$, l'anneau $\mathcal{O}_X(U)$
est un G-anneau.
\item Il existe un recouvrement ouvert affine $X = \bigcup U_i$ tel que
chaque $\mathcal{O}_X(U_i)$ soit un G-anneau.
\item Il existe un recouvrement ouvert $X = \bigcup X_j$
tel que chaque sous-schéma ouvert $X_j$ soit un G-schéma.
\end{enumerate}
De plus, si $X$ est un G-schéma, tout sous-schéma ouvert est un
G-schéma.
\end{lemma}

\begin{proof}
Combiner les lemmes \ref{properties-lemma-locally-Noetherian}
et \ref{properties-lemma-characterize-G-scheme}.
\end{proof}




\section{Le lieu singulier}
\label{properties-section-singular-locus}

\noindent
Voici la définition.

\begin{definition}
\label{properties-definition-singular-locus}
Soit $X$ un schéma localement noethérien. Le {\it lieu régulier}
$\text{Reg}(X)$ de $X$ est l'ensemble des $x \in X$ tels que $\mathcal{O}_{X, x}$
soit un anneau local régulier. Le {\it lieu singulier} $\text{Sing}(X)$ est le
complémentaire $X \setminus \text{Reg}(X)$, c'est-à-dire l'ensemble des points $x \in X$
tels que $\mathcal{O}_{X, x}$ ne soit pas un anneau local régulier.
\end{definition}

\noindent
Le lieu régulier d'un schéma localement noethérien est stable
par générisation, voir la discussion précédant
Algèbre, définition \ref{algebra-definition-regular}.
Toutefois, pour un schéma localement noethérien quelconque, le lieu régulier n'est pas
nécessairement ouvert. Dans
Compléments sur l'algèbre, section \ref{more-algebra-section-singular-locus},
on trouvera quelques critères assurant qu'il l'est. Nous
reviendrons sur cette question dans
Morphismes, section \ref{morphisms-section-singular-locus}.






\section{Irréductibilité locale}
\label{properties-section-unibranch}

\noindent
Rappelons que, dans Compléments sur l'algèbre, section \ref{more-algebra-section-unibranch},
nous avons introduit la notion d'anneau local (géométriquement) unibranche.

\begin{definition}
\label{properties-definition-unibranch}
\begin{reference}
\cite[Chapter IV (6.15.1)]{EGA4}
\end{reference}
Soient $X$ un schéma et $x \in X$. On dit que $X$ est {\it unibranche en $x$}
si l'anneau local $\mathcal{O}_{X, x}$ est unibranche. On dit que $X$ est
{\it géométriquement unibranche en $x$}
si l'anneau local $\mathcal{O}_{X, x}$ est géométriquement unibranche. On
dit que $X$ est {\it unibranche} si $X$ est unibranche en chacun de ses points. On
dit que $X$ est {\it géométriquement unibranche} si $X$ est
géométriquement unibranche en chacun de ses points.
\end{definition}

\noindent
Il peut effectivement arriver qu'un anneau local $A$ soit géométriquement unibranche
(au sens de
Compléments sur l'algèbre, définition \ref{more-algebra-definition-unibranch}),
sans que le schéma $\Spec(A)$ soit géométriquement unibranche au sens de
la définition \ref{properties-definition-unibranch}. C'est le cas, par exemple,
si $A$ est l'anneau local au sommet du cône sur une courbe
plane irréductible ayant un point double ordinaire (un nœud).

\begin{lemma}
\label{properties-lemma-normal-geometrically-unibranch}
Un schéma normal est géométriquement unibranche.
\end{lemma}

\begin{proof}
Cela résulte des définitions. En effet, un schéma est normal
si ses anneaux locaux sont des anneaux intègres normaux. Il résulte
immédiatement de Compléments sur l'algèbre, définition \ref{more-algebra-definition-unibranch},
qu'un anneau local intègre normal est géométriquement unibranche.
\end{proof}

\begin{lemma}
\label{properties-lemma-geometrically-unibranch}
\begin{reference}
Comparer avec \cite[Proposition 2.3]{Etale-coverings}
\end{reference}
Soit $X$ un schéma noethérien. Les assertions suivantes sont équivalentes :
\begin{enumerate}
\item $X$ est géométriquement unibranche (définition \ref{properties-definition-unibranch}) ;
\item pour tout point $x \in X$ qui n'est pas le point générique
d'une composante irréductible de $X$, le spectre épointé de
l'hensélisé strict $\mathcal{O}_{X, x}^{sh}$ est connexe.
\end{enumerate}
\end{lemma}

\begin{proof}
Compléments sur l'algèbre, lemme \ref{more-algebra-lemma-geometrically-unibranch},
montre que (1) entraîne que les spectres épointés de
(2) sont irréductibles, donc connexes.

\medskip\noindent
Supposons (2). Soit $x \in X$. Il faut montrer que $\mathcal{O}_{X, x}$
est géométriquement unibranche. Par récurrence sur $\dim(\mathcal{O}_{X, x})$,
on peut supposer le résultat vrai pour toute générisation non triviale de $x$.
On peut remplacer $X$ par $\Spec(\mathcal{O}_{X, x})$. Autrement dit,
on peut supposer que $X = \Spec(A)$, avec $A$ local, et que
$A_\mathfrak p$ est géométriquement unibranche pour tout idéal premier non
maximal $\mathfrak p \subset A$.

\medskip\noindent
Soit $A^{sh}$ l'hensélisé strict de $A$. Si
$\mathfrak q \subset A^{sh}$ est un idéal premier au-dessus de $\mathfrak p \subset A$,
alors $A_\mathfrak p \to A^{sh}_\mathfrak q$ est
une colimite filtrante d'algèbres étales. Les hensélisés stricts de
$A_\mathfrak p$ et de $A^{sh}_\mathfrak q$ sont donc isomorphes.
Ainsi, d'après Compléments sur l'algèbre, lemme \ref{more-algebra-lemma-geometrically-unibranch},
$A^{sh}_\mathfrak q$
possède un unique idéal premier minimal pour tout idéal premier non maximal $\mathfrak q$ de
$A^{sh}$.

\medskip\noindent
Soient $\mathfrak q_1, \ldots, \mathfrak q_r$ les idéaux premiers minimaux
de $A^{sh}$. Il faut montrer que $r = 1$. D'après ce qui
précède, on a $V(\mathfrak q_1) \cap V(\mathfrak q_j) = \{\mathfrak m^{sh}\}$
pour $j = 2, \ldots, r$. Ainsi $V(\mathfrak q_1) \setminus \{\mathfrak m^{sh}\}$
est un ouvert fermé du spectre épointé de $A^{sh}$,
ce qui contredit l'hypothèse selon laquelle ce spectre épointé est
connexe, à moins que $r = 1$.
\end{proof}

\begin{definition}
\label{properties-definition-number-of-branches}
Soient $X$ un schéma et $x \in X$. Le {\it nombre de branches de $X$
en $x$} est le nombre de branches de l'anneau local $\mathcal{O}_{X, x}$
au sens de
Compléments sur l'algèbre, définition \ref{more-algebra-definition-number-of-branches}.
Le {\it nombre de branches géométriques de $X$ en $x$} est le nombre de
branches géométriques de l'anneau local $\mathcal{O}_{X, x}$ au sens de
Compléments sur l'algèbre, définition \ref{more-algebra-definition-number-of-branches}.
\end{definition}

\noindent
On souhaite souvent comparer ce nombre à celui des branches de l'anneau local
complété, mais cette comparaison n'est généralement pas immédiate ; on trouvera
quelques résultats à ce sujet dans Compléments sur l'algèbre, section
\ref{more-algebra-section-branches-completion}.

\begin{lemma}
\label{properties-lemma-number-of-branches-irreducible-components}
Soient $X$ un schéma et $x \in X$. Soient $X_i$, $i \in I$, les
composantes irréductibles de $X$ passant par $x$.
Alors le nombre de branches (géométriques) de $X$ en $x$
est la somme, sur $i \in I$, des nombres de branches
(géométriques) de $X_i$ en $x$.
\end{lemma}

\begin{proof}
On considère les $X_i$ comme des sous-schémas fermés intègres de $X$,
voir Schémas, définition \ref{schemes-definition-reduced-induced-scheme}, et
le lemme \ref{properties-lemma-characterize-integral}.
Remarquons que le nombre de branches (géométriques) de $X_i$ en $x$
est au moins $1$ pour tout $i$ (essentiellement par définition).
Rappelons que les $X_i$ correspondent $1$-à-$1$ aux idéaux
premiers minimaux $\mathfrak p_i \subset \mathcal{O}_{X, x}$, voir
Algèbre, lemme \ref{algebra-lemma-irreducible-components-containing-x}.
Ainsi, si $I$ est infini, $\mathcal{O}_{X, x}$ possède une
infinité d'idéaux premiers minimaux, et donc $\mathcal{O}_{X, x}^h$
et $\mathcal{O}_{X, x}^{sh}$ possèdent tous deux une infinité
d'idéaux premiers minimaux (combiner Algèbre, lemmes
\ref{algebra-lemma-injective-minimal-primes-in-image} et
\ref{algebra-lemma-minimal-prime-image-minimal-prime},
avec l'injectivité des homomorphismes
$\mathcal{O}_{X, x} \to  \mathcal{O}_{X, x}^h \to \mathcal{O}_{X, x}^{sh}$).
Dans ce cas, le nombre de branches (géométriques) de $X$ en $x$
est, par définition, $\infty$, tout comme la somme considérée.
On peut donc supposer $I$ fini.
Soit $A'$ la fermeture intégrale de $\mathcal{O}_{X, x}$
dans l'anneau total des fractions $Q$ de $(\mathcal{O}_{X, x})_{red}$.
Soit $A'_i$ la fermeture intégrale de $\mathcal{O}_{X, x}/\mathfrak p_i$
dans l'anneau total des fractions $Q_i$ de $\mathcal{O}_{X, x}/\mathfrak p_i$.
D'après Algèbre, lemme \ref{algebra-lemma-total-ring-fractions-no-embedded-points},
on a $Q = \prod_{i \in I} Q_i$. Ainsi $A' = \prod A'_i$.
L'égalité du lemme résulte alors de
Compléments sur l'algèbre, lemme \ref{more-algebra-lemma-number-of-branches-1},
qui exprime le nombre de branches (géométriques) en
fonction des idéaux maximaux de $A'$.
\end{proof}

\begin{lemma}
\label{properties-lemma-number-of-branches-1}
Soient $X$ un schéma et $x \in X$.
\begin{enumerate}
\item Le nombre de branches de $X$ en $x$ est $1$ si et seulement si
$X$ est unibranche en $x$.
\item Le nombre de branches géométriques de $X$ en $x$ est $1$ si et seulement si
$X$ est géométriquement unibranche en $x$.
\end{enumerate}
\end{lemma}

\begin{proof}
Ce lemme résulte immédiatement des définitions et du résultat correspondant pour
les anneaux, voir Compléments sur l'algèbre, lemme
\ref{more-algebra-lemma-number-of-branches-1}.
\end{proof}






\section{Caractérisation des modules de type fini et de présentation finie}
\label{properties-section-characterizing-finite-type-presentation}

\noindent
Soit $X$ un schéma.
Soit $\mathcal{F}$ un $\mathcal{O}_X$-module
quasi-cohérent. Le lemme suivant entraîne que $\mathcal{F}$ est de type fini
(voir Modules, définition \ref{modules-definition-finite-type})
si et seulement si $\mathcal{F}$ est,
sur tout ouvert affine $\Spec(A) = U \subset X$,
de la forme $\widetilde M$ pour un $A$-module de type fini $M$.
De même, $\mathcal{F}$ est de présentation finie
(voir Modules, définition \ref{modules-definition-finite-presentation})
si et seulement si $\mathcal{F}$ est,
sur tout ouvert affine $\Spec(A) = U \subset X$,
de la forme $\widetilde M$ pour un $A$-module de présentation finie $M$.

\begin{lemma}
\label{properties-lemma-finite-type-module}
Soit $X = \Spec(R)$ un schéma affine. Le
faisceau quasi-cohérent de $\mathcal{O}_X$-modules
$\widetilde M$ est un $\mathcal{O}_X$-module de
type fini si et seulement si $M$ est un $R$-module de type fini.
\end{lemma}

\begin{proof}
Supposons que $\widetilde M$ soit un $\mathcal{O}_X$-module de
type fini. Cela signifie qu'il existe un recouvrement ouvert de $X$ tel que la restriction de
$\widetilde M$ à chacun des membres de ce recouvrement soit
engendrée par un nombre fini de sections globales.
Il existe donc aussi un recouvrement ouvert standard
$X = \bigcup_{i = 1, \ldots, n} D(f_i)$ tel que $\widetilde M|_{D(f_i)}$
soit engendré par un nombre fini de sections. Ainsi $M_{f_i}$ est de
type fini pour tout $i$. On conclut par
Algèbre, lemme \ref{algebra-lemma-cover}.
\end{proof}

\begin{lemma}
\label{properties-lemma-finite-presentation-module}
Soit $X = \Spec(R)$ un schéma affine. Le faisceau quasi-cohérent
de $\mathcal{O}_X$-modules $\widetilde M$ est un $\mathcal{O}_X$-module de
présentation finie si et seulement si $M$ est un $R$-module de présentation finie.
\end{lemma}

\begin{proof}
Supposons que $\widetilde M$ soit un $\mathcal{O}_X$-module de présentation finie.
Le lemme \ref{properties-lemma-finite-type-module} montre que $M$ est un $R$-module de
type fini. Choisissons une surjection $R^n \to M$ de noyau $K$.
D'après Schémas, lemme \ref{schemes-lemma-spec-sheaves},
on a une suite exacte courte
$$
0 \to \widetilde{K} \to
\bigoplus \mathcal{O}_X^{\oplus n} \to
\widetilde{M} \to 0
$$
D'après
Modules, lemme
\ref{modules-lemma-kernel-surjection-finite-free-onto-finite-presentation},
$\widetilde{K}$ est un $\mathcal{O}_X$-module de type
fini. Donc, de nouveau par le lemme \ref{properties-lemma-finite-type-module},
$K$ est un $R$-module de type
fini. Ainsi $M$ est un $R$-module de présentation finie.
\end{proof}







\section{sections sur les ouverts principaux}
\label{properties-section-principal-opens}

\noindent
Voici un résultat typique de cette nature. Dans les lemmes suivants, nous utiliserons
une démonstration plus élémentaire, mais plus directe.

\begin{lemma}
\label{properties-lemma-invert-f-sections}
\begin{slogan}
Les sections des faisceaux quasi-cohérents n'ont que des singularités
méromorphes à l'infini.
\end{slogan}
Soient $X$ un schéma et $f \in \Gamma(X, \mathcal{O}_X)$.
Notons $X_f \subset X$ l'ouvert où $f$ est inversible, voir
Schémas, lemme \ref{schemes-lemma-f-open}.
Si $X$ est quasi-compact et quasi-séparé, l'homomorphisme canonique
$$
\Gamma(X, \mathcal{O}_X)_f \longrightarrow \Gamma(X_f, \mathcal{O}_X)
$$
est un isomorphisme. De plus, si $\mathcal{F}$ est un faisceau
quasi-cohérent de $\mathcal{O}_X$-modules, l'homomorphisme
$$
\Gamma(X, \mathcal{F})_f \longrightarrow \Gamma(X_f, \mathcal{F})
$$
est un isomorphisme.
\end{lemma}

\begin{proof}
Écrivons $R = \Gamma(X, \mathcal{O}_X)$. Considérons le morphisme canonique
$$
\varphi : X \longrightarrow \Spec(R)
$$
de schémas, voir
Schémas, lemme
\ref{schemes-lemma-morphism-into-affine}.
L'image inverse de l'ouvert principal $D(f)$ du membre
de droite est alors l'ouvert $X_f$ du membre de gauche.
De plus, comme $X$ est supposé quasi-compact et quasi-séparé,
le morphisme $\varphi$ est quasi-compact et quasi-séparé,
voir Schémas, lemmes \ref{schemes-lemma-quasi-compact-affine} et
\ref{schemes-lemma-compose-after-separated}. Ainsi,
d'après Schémas, lemme \ref{schemes-lemma-push-forward-quasi-coherent},
$\varphi_*\mathcal{F}$ est
quasi-cohérent. On a donc $\varphi_*\mathcal{F} = \widetilde M$,
où $M = \Gamma(X, \mathcal{F})$, considéré comme $R$-module.
Il en résulte que
$$
\Gamma(X_f, \mathcal{F}) =
\Gamma(D(f), \varphi_*\mathcal{F}) =
\Gamma(D(f), \widetilde M) = M_f
$$
ce qui est précisément l'assertion du lemme. Le premier isomorphisme
affiché s'obtient en prenant $\mathcal{F} = \mathcal{O}_X$.
\end{proof}

\noindent
Rappelons qu'étant donnés un schéma $X$, un faisceau inversible $\mathcal{L}$
sur $X$ et un faisceau de $\mathcal{O}_X$-modules $\mathcal{F}$,
on obtient un anneau gradué
$\Gamma_*(X, \mathcal{L}) =
\bigoplus\nolimits_{n \geq 0} \Gamma(X, \mathcal{L}^{\otimes n})$
et un $\Gamma_*(X, \mathcal{L})$-module gradué
$\Gamma_*(X, \mathcal{L}, \mathcal{F}) =
\bigoplus\nolimits_{n \in \mathbf{Z}}
\Gamma(X, \mathcal{F} \otimes_{\mathcal{O}_X} \mathcal{L}^{\otimes n})$,
voir Modules, définition \ref{modules-definition-gamma-star}.
Si l'on se donne de plus une section $s \in \Gamma(X, \mathcal{L})$, on
obtient une application
\begin{equation}
\label{properties-equation-module-invert-s}
\Gamma_*(X, \mathcal{L}, \mathcal{F})_{(s)}
\longrightarrow
\Gamma(X_s, \mathcal{F}|_{X_s})
\end{equation}
qui envoie $t/s^n$, où
$t \in \Gamma(X, \mathcal{F} \otimes_{\mathcal{O}_X} \mathcal{L}^{\otimes n})$,
sur $t|_{X_s} \otimes s|_{X_s}^{-n}$. Cela a un
sens, car $X_s \subset X$ est, par définition, l'ouvert sur lequel
$s$ possède un inverse, voir Modules, lemme \ref{modules-lemma-s-open}.

\begin{lemma}
\label{properties-lemma-invert-s-sections}
\begin{reference}
\cite[section 6.8]{EGA1-second}
\end{reference}
Soit $X$ un schéma. Soit $\mathcal{L}$ un faisceau inversible sur $X$.
Soit $s \in \Gamma(X, \mathcal{L})$. Soit $\mathcal{F}$ un
$\mathcal{O}_X$-module quasi-cohérent.
\begin{enumerate}
\item Si $X$ est quasi-compact, l'application (\ref{properties-equation-module-invert-s})
est injective ;
\item si $X$ est quasi-compact et quasi-séparé, l'application
(\ref{properties-equation-module-invert-s}) est un isomorphisme.
\end{enumerate}
En particulier, l'homomorphisme canonique
$$
\Gamma_*(X, \mathcal{L})_{(s)}
\longrightarrow
\Gamma(X_s, \mathcal{O}_X),\quad
a/s^n \longmapsto a \otimes s^{-n}
$$
est un isomorphisme si $X$ est quasi-compact et quasi-séparé.
\end{lemma}

\begin{proof}
Supposons $X$ quasi-compact. Choisissons un recouvrement ouvert affine fini
$X = U_1 \cup \ldots \cup U_m$, avec $U_j$ affine et
$\mathcal{L}|_{U_j} \cong \mathcal{O}_{U_j}$. Par cet isomorphisme,
la section $s|_{U_j}$ correspond à un élément
$f_j \in \Gamma(U_j, \mathcal{O}_{U_j})$. Alors
$X_s \cap U_j = D(f_j)$.

\medskip\noindent
Démonstration de (1). Soit $t/s^n$ un élément du noyau de
(\ref{properties-equation-module-invert-s}). Alors $t|_{X_s} = 0$.
Donc $(t|_{U_j})|_{D(f_j)} = 0$. Par
le lemme \ref{properties-lemma-invert-f-sections}, on obtient
$f_j^{e_j} t|_{U_j} = 0$ pour un certain
$e_j \geq 0$. Soit $e = \max(e_j)$. Alors $t \otimes s^e$
se restreint à zéro sur $U_j$ pour tout $j$, et est donc nul. Comme $t/s^n$
est égal à $t \otimes s^e/s^{n + e}$ dans
$\Gamma_*(X, \mathcal{L}, \mathcal{F})_{(s)}$, on en déduit que $t/s^n = 0$,
comme voulu.

\medskip\noindent
Démonstration de (2). Supposons $X$ quasi-compact et quasi-séparé.
Alors $U_j \cap U_{j'}$ est quasi-compact pour tout couple $j, j'$, voir
Schémas, lemme \ref{schemes-lemma-characterize-quasi-separated}.
D'après (1), on sait que (\ref{properties-equation-module-invert-s}) est injective.
Soit $t' \in \Gamma(X_s, \mathcal{F}|_{X_s})$. Pour tout $j$, il existe un
entier $e_j \geq 0$ et une section $t'_j \in \Gamma(U_j, \mathcal{F}|_{U_j})$ tels que
$t'|_{D(f_j)}$ corresponde à $t'_j/f_j^{e_j}$
par l'isomorphisme du lemme \ref{properties-lemma-invert-f-sections}.
Posons $e = \max(e_j)$ et
$$
t_j = f_j^{e - e_j} t'_j \otimes q_j^e \in
\Gamma(U_j,
(\mathcal{F} \otimes_{\mathcal{O}_X} \mathcal{L}^{\otimes e})|_{U_j})
$$
où $q_j \in \Gamma(U_j, \mathcal{L}|_{U_j})$ est la section
trivialisante provenant de l'isomorphisme
$\mathcal{L}|_{U_j} \cong \mathcal{O}_{U_j}$. En particulier, on a
$s|_{U_j} = f_j q_j$. Un calcul utilisant cette égalité montre que
$t_j|_{U_j \cap U_{j'}}$ et $t_{j'}|_{U_j \cap U_{j'}}$
ont pour image la même section de $\mathcal{F}$ sur $U_j \cap U_{j'} \cap X_s$.
Par quasi-compacité de $U_j \cap U_{j'}$ et d'après (1), il existe un
entier $e' \geq 0$ tel que
$$
t_j|_{U_j \cap U_{j'}} \otimes s^{e'}|_{U_j \cap U_{j'}} =
t_{j'}|_{U_j \cap U_{j'}} \otimes s^{e'}|_{U_j \cap U_{j'}}
$$
comme sections de $\mathcal{F} \otimes \mathcal{L}^{\otimes e + e'}$ sur
$U_j \cap U_{j'}$. On peut choisir un même $e'$ convenant à tous les couples
$j, j'$. La condition de faisceau fournit alors une section
$t \in \Gamma(X, \mathcal{F} \otimes \mathcal{L}^{\otimes e + e'})$
dont la restriction à $U_j$ est $t_j \otimes s^{e'}|_{U_j}$.
Un calcul immédiat montre que $t/s^{e + e'}$ a pour image $t'$,
comme voulu.
\end{proof}

\noindent
Soit $X$ un schéma. Soit $\mathcal{L}$ un $\mathcal{O}_X$-module inversible.
Soient $\mathcal{F}$ et $\mathcal{G}$ des $\mathcal{O}_X$-modules
quasi-cohérents. Considérons le $\Gamma_*(X, \mathcal{L})$-module gradué
$$
M = \bigoplus\nolimits_{n \in \mathbf{Z}} \Hom_{\mathcal{O}_X}(\mathcal{F},
\mathcal{G} \otimes_{\mathcal{O}_X} \mathcal{L}^{\otimes n})
$$
Soit maintenant $s \in \Gamma(X, \mathcal{L})$ une section. On dispose d'une
application canonique
\begin{equation}
\label{properties-equation-hom-invert-s}
M_{(s)} \longrightarrow
\Hom_{\mathcal{O}_{X_s}}(\mathcal{F}|_{X_s}, \mathcal{G}|_{X_s})
\end{equation}
qui envoie $\alpha/s^n$ sur le morphisme $\alpha|_{X_s} \otimes s|_{X_s}^{-n}$.
Le lemme suivant, combiné au
lemme \ref{properties-lemma-extend-finite-presentation},
affirme en substance que, si $X$ est quasi-compact et
quasi-séparé, la catégorie des $\mathcal{O}_{X_s}$-modules de
présentation finie s'obtient à partir de la catégorie des $\mathcal{O}_X$-modules de
présentation finie en inversant le système multiplicatif de morphismes
$s^n: \mathcal{F} \to
\mathcal{F} \otimes_{\mathcal{O}_X} \mathcal{L}^{\otimes n}$.

\begin{lemma}
\label{properties-lemma-section-maps-backwards}
Soit $X$ un schéma. Soit $\mathcal{L}$ un $\mathcal{O}_X$-module
inversible. Soit $s \in \Gamma(X, \mathcal{L})$ une section.
Soient $\mathcal{F}$, $\mathcal{G}$ des $\mathcal{O}_X$-modules quasi-cohérents.
\begin{enumerate}
\item Si $X$ est quasi-compact et si $\mathcal{F}$ est de type fini,
l'application (\ref{properties-equation-hom-invert-s}) est injective ;
\item si $X$ est quasi-compact et quasi-séparé et si $\mathcal{F}$
est de présentation finie, l'application
(\ref{properties-equation-hom-invert-s})
est bijective.
\end{enumerate}
\end{lemma}

\begin{proof}
Démontrons d'abord le lemme lorsque $X = \Spec(A)$ est affine
et $\mathcal{L} = \mathcal{O}_X$. Dans ce cas, $s$ correspond
à un élément $f \in A$. Écrivons
$\mathcal{F} = \widetilde{M}$ et $\mathcal{G} = \widetilde{N}$
pour des $A$-modules $M$ et $N$. Le lemme se traduit alors, à
l'aide des lemmes \ref{properties-lemma-finite-type-module} et
\ref{properties-lemma-finite-presentation-module}, par
les assertions algébriques suivantes :
\begin{enumerate}
\item Si $M$ est un $A$-module de type fini et si $\varphi : M \to N$ est un
morphisme de $A$-modules tel que le morphisme induit $M_f \to N_f$ soit nul,
alors $f^n\varphi = 0$ pour un certain $n$.
\item Si $M$ est un $A$-module de présentation finie, alors
$\Hom_A(M, N)_f = \Hom_{A_f}(M_f, N_f)$.
\end{enumerate}
La seconde assertion est
Algèbre, lemme \ref{algebra-lemma-hom-from-finitely-presented} ; nous
omettons la démonstration de la première.

\medskip\noindent
Démontrons maintenant (1) pour $X$ quelconque.
Supposons $X$ quasi-compact et choisissons un recouvrement ouvert affine fini
$X = U_1 \cup \ldots \cup U_m$, avec $U_j$ affine et
$\mathcal{L}|_{U_j} \cong \mathcal{O}_{U_j}$. Par cet isomorphisme,
la section $s|_{U_j}$ correspond à un élément
$f_j \in \Gamma(U_j, \mathcal{O}_{U_j})$. Alors
$X_s \cap U_j = D(f_j)$.
Soit $\alpha/s^n$ un élément du noyau de
(\ref{properties-equation-hom-invert-s}). Alors $\alpha|_{X_s} = 0$.
Donc $(\alpha|_{U_j})|_{D(f_j)} = 0$. D'après le cas affine traité
ci-dessus, on obtient $f_j^{e_j} \alpha|_{U_j} = 0$ pour un certain
$e_j \geq 0$. Soit $e = \max(e_j)$. Alors $\alpha \otimes s^e$
se restreint à zéro sur $U_j$ pour tout $j$, et est donc nul. Comme $\alpha/s^n$
est égal à $\alpha \otimes s^e/s^{n + e}$ dans $M_{(s)}$, on en déduit que
$\alpha/s^n = 0$, comme voulu.

\medskip\noindent
Démonstration de (2). Comme $\mathcal{F}$ est de présentation finie, le
faisceau $\SheafHom_{\mathcal{O}_X}(\mathcal{F}, \mathcal{G})$ est
quasi-cohérent, voir Schémas, section \ref{schemes-section-quasi-coherent}.
De plus, il est clair que
$$
\SheafHom_{\mathcal{O}_X}(\mathcal{F},
\mathcal{G} \otimes_{\mathcal{O}_X} \mathcal{L}^{\otimes n}) =
\SheafHom_{\mathcal{O}_X}(\mathcal{F}, \mathcal{G})
\otimes_{\mathcal{O}_X} \mathcal{L}^{\otimes n}
$$
pour tout $n$. Dans ce cas, l'assertion résulte donc du
lemme \ref{properties-lemma-invert-s-sections} appliqué à
$\SheafHom_{\mathcal{O}_X}(\mathcal{F}, \mathcal{G})$.
\end{proof}





\section{Schémas quasi-affines}
\label{properties-section-quasi-affine}

\begin{definition}
\label{properties-definition-quasi-affine}
Un schéma $X$ est dit {\it quasi-affine} s'il est quasi-compact
et isomorphe à un sous-schéma ouvert d'un schéma affine.
\end{definition}

\begin{lemma}
\label{properties-lemma-quasi-coherent-quasi-affine}
Soient $A$ un anneau et $U \subset \Spec(A)$ un sous-schéma
ouvert quasi-compact. Pour $\mathcal{F}$ quasi-cohérent sur $U$, le morphisme canonique
$$
\widetilde{H^0(U, \mathcal{F})}|_U \to \mathcal{F}
$$
est un isomorphisme.
\end{lemma}

\begin{proof}
Notons $j : U \to \Spec(A)$ le morphisme d'inclusion. Alors
$H^0(U, \mathcal{F}) = H^0(\Spec(A), j_*\mathcal{F})$, et
$j_*\mathcal{F}$ est quasi-cohérent
d'après Schémas, lemme \ref{schemes-lemma-push-forward-quasi-coherent}.
Ainsi $j_*\mathcal{F} = \widetilde{H^0(U, \mathcal{F})}$ par
Schémas, lemme \ref{schemes-lemma-equivalence-quasi-coherent}.
En restreignant à $U$, on obtient le lemme.
\end{proof}

\begin{lemma}
\label{properties-lemma-invert-f-affine}
Soient $X$ un schéma et $f \in \Gamma(X, \mathcal{O}_X)$.
Supposons $X$ quasi-compact et quasi-séparé, et supposons que
$X_f$ soit affine. Alors le morphisme canonique
$$
j : X \longrightarrow \Spec(\Gamma(X, \mathcal{O}_X))
$$
de Schémas, lemme \ref{schemes-lemma-morphism-into-affine},
induit un isomorphisme de $X_f = j^{-1}(D(f))$ sur l'ouvert affine
principal $D(f) \subset \Spec(\Gamma(X, \mathcal{O}_X))$.
\end{lemma}

\begin{proof}
C'est immédiat, puisque $j$ induit un isomorphisme d'anneaux
$\Gamma(X, \mathcal{O}_X)_f \to \mathcal{O}_X(X_f)$ d'après
le lemme \ref{properties-lemma-invert-f-sections} ci-dessus.
\end{proof}

\begin{lemma}
\label{properties-lemma-quasi-affine}
Soit $X$ un schéma. Alors $X$ est quasi-affine si et seulement
si le morphisme canonique
$$
X \longrightarrow \Spec(\Gamma(X, \mathcal{O}_X))
$$
de Schémas, lemme \ref{schemes-lemma-morphism-into-affine},
est une immersion ouverte quasi-compacte.
\end{lemma}

\begin{proof}
Si le morphisme affiché est une immersion ouverte quasi-compacte,
$X$ est isomorphe à un sous-schéma ouvert quasi-compact de
$\Spec(\Gamma(X, \mathcal{O}_X))$, et $X$ est évidemment quasi-affine.

\medskip\noindent
Supposons $X$ quasi-affine ; écrivons $X \subset \Spec(R)$ comme
un ouvert quasi-compact. Cela entraîne en particulier que $X$ est
séparé, voir
Schémas, lemme \ref{schemes-lemma-subscheme-of-separated-scheme}.
Soit $A = \Gamma(X, \mathcal{O}_X)$.
Considérons l'homomorphisme d'anneaux $R \to A$ provenant de
$R = \Gamma(\Spec(R), \mathcal{O}_{\Spec(R)})$
et de l'application de restriction du faisceau $\mathcal{O}_{\Spec(R)}$.
D'après Schémas, lemme \ref{schemes-lemma-morphism-into-affine},
on obtient une factorisation
$$
X \longrightarrow
\Spec(A) \longrightarrow
\Spec(R)
$$
du morphisme d'inclusion. Soit $x \in X$. Choisissons $r \in R$ tel que
$x \in D(r)$ et $D(r) \subset X$. Notons $f \in A$ l'image de $r$
dans $A$. L'ouvert $X_f$ du lemme \ref{properties-lemma-invert-f-sections}
ci-dessus est égal à $D(r) \subset X$ ; donc $A_f \cong R_r$ d'après
ce lemme.
Ainsi $D(r) \to \Spec(A)$ est un isomorphisme sur
l'ouvert affine principal $D(f)$ de $\Spec(A)$. Comme $X$
peut être recouvert par de tels ouverts affines $D(f)$, on conclut.
\end{proof}

\begin{lemma}
\label{properties-lemma-cartesian-diagram-quasi-affine}
Soit $U \to V$ une immersion ouverte de schémas quasi-affines. Alors
$$
\xymatrix{
U \ar[d] \ar[rr]_-j & & \Spec(\Gamma(U, \mathcal{O}_U)) \ar[d] \\
U \ar[r] & V \ar[r]^-{j'} & \Spec(\Gamma(V, \mathcal{O}_V))
}
$$
est cartésien.
\end{lemma}

\begin{proof}
Le diagramme est commutatif d'après Schémas, lemme
\ref{schemes-lemma-morphism-into-affine}.
Écrivons $A = \Gamma(U, \mathcal{O}_U)$ et $B = \Gamma(V, \mathcal{O}_V)$. Soit
$g \in B$ tel que $V_g$ soit affine et contenu dans $U$. Cela
signifie que, si $f$ est l'image de $g$ dans $A$, alors $U_f = V_g$. D'après le lemme
\ref{properties-lemma-invert-f-affine}, $j'$ induit un isomorphisme de
$V_g$ sur l'ouvert principal $D(g)$ de $\Spec(B)$. Ainsi
$V_g \times_{\Spec(B)} \Spec(A) \to \Spec(A)$ est
un isomorphisme sur $D(f) \subset \Spec(A)$. De nouveau par le lemme \ref{properties-lemma-invert-f-affine},
$j$ envoie $U_f$ isomorphiquement sur $D(f)$. On obtient donc
$U_f = U_f \times_{\Spec(B)} \Spec(A)$. Puisque le
lemme \ref{properties-lemma-quasi-affine} permet de recouvrir $U$ par des $V_g = U_f$ comme
ci-dessus, on voit que $U \to U \times_{\Spec(B)} \Spec(A)$ est un isomorphisme.
\end{proof}

\begin{lemma}
\label{properties-lemma-quasi-affine-presentation}
Soit $X$ un schéma quasi-affine. Il existe un entier $n \geq 0$,
un schéma affine $T$ et un morphisme $T \to X$ tels que, pour tout
morphisme $X' \to X$ avec $X'$ affine, le produit fibré $X' \times_X T$
soit isomorphe à $\mathbf{A}^n_{X'}$ au-dessus de $X'$.
\end{lemma}

\begin{proof}
Par définition, il existe un anneau $A$ tel que $X$ soit isomorphe à un
sous-schéma ouvert quasi-compact $U \subset \Spec(A)$. Rappelons que les ouverts
principaux $D(f) \subset \Spec(A)$ forment une base de la topologie,
voir Algèbre, section \ref{algebra-section-spectrum-ring}. Comme $U$ est
quasi-compact, on peut choisir $f_1, \ldots, f_n \in A$ tels que
$U = D(f_1) \cup \ldots \cup D(f_n)$. On peut donc supposer
$X = \Spec(A) \setminus V(I)$, où $I = (f_1, \ldots, f_n)$. Posons
$$
T = \Spec(A[t, x_1, \ldots, x_n]/(f_1 x_1 + \ldots + f_n x_n - 1))
$$
Le morphisme structural $T \to \Spec(A)$ se factorise par l'ouvert $X$
et fournit le morphisme $T \to X$. Si $X' = \Spec(A')$ et si le morphisme
$X' \to X$ correspond à l'homomorphisme d'anneaux $A \to A'$, les images
$f'_1, \ldots, f'_n \in A'$ de $f_1, \ldots, f_n$
engendrent l'idéal unité de $A'$.
Écrivons $1 = f'_1 a'_1 + \ldots + f'_n a'_n$.
Le changement de base $X' \times_X T$ est le spectre de
$A'[t, x_1, \ldots, x_n]/(f'_1 x_1 + \ldots + f'_n x_n - 1)$.
Nous affirmons que l'homomorphisme de $A'$-algèbres
$$
\varphi :
A'[y_1, \ldots, y_n]
\longrightarrow
A'[t, x_1, \ldots, x_n, x_{n + 1}]/(f'_1 x_1 + \ldots + f'_n x_n - 1)
$$
envoyant $y_i$ sur $a'_i t + x_i$ est un isomorphisme. Cette assertion  
achève la démonstration du lemme. L'inverse de $\varphi$ est l'homomorphisme de $A'$-algèbres donné
ci-dessous
$$
\psi :
A'[t, x_1, \ldots, x_n, x_{n + 1}]/(f'_1 x_1 + \ldots + f'_n x_n - 1)
\longrightarrow
A'[y_1, \ldots, y_n]
$$
envoyant $t$ sur $-1 + f'_1 y_1 + \ldots + f'_n y_n$ et $x_i$ sur
$y_i + a'_i - a'_i(f'_1 y_1 + \ldots + f'_n y_n)$ pour $i = 1, \ldots, n$.
Cela a un sens, puisque $\sum f'_ix_i$ est envoyé sur
$$
\begin{matrix}
\sum f'_i(y_i + a'_i - a'_i(\sum f'_j y_j)) =
(\sum f'_iy_i) + 1 - (\sum f'_j y_j) = 1
\end{matrix}
$$
Pour vérifier que ces homomorphismes sont inverses l'un de l'autre, on calcule :
$$
\begin{matrix}
\varphi(\psi(t) = \varphi(-1 +  \sum f'_i y_i) =
-1 + \sum f'_i (a'_i t + x_i) = t \\
\varphi(\psi(x_i)) = \varphi(y_i + a'_i - a'_i(\sum f'_j y_j)) =
a'_i t + x_i + a'_i - a'_i(\sum f'_ja'_jt + f'_jx_j) = x_i \\
\psi(\varphi(y_i)) = \psi(a'_i t + x_i) =
a'_i(-1 + \sum f'_j y_j) + y_i + a'_i - a'_i(\sum f'_j y_j) = y_i
\end{matrix}
$$
Cela achève la démonstration.
\end{proof}




\section{Modules plats}
\label{properties-section-flat}

\noindent
Sur tout espace annelé $(X, \mathcal{O}_X)$,
on sait ce que signifie, pour un $\mathcal{O}_X$-module,
être plat (en un point), voir
Modules, définition \ref{modules-definition-flat}
(définition \ref{modules-definition-flat-at-point}).
Pour les faisceaux quasi-cohérents sur un schéma affine, cette notion coïncide avec
celle définie dans le chapitre d'algèbre.

\begin{lemma}
\label{properties-lemma-flat-module}
\begin{slogan}
La platitude est la même pour les modules et les faisceaux.
\end{slogan}
Soit $X = \Spec(R)$ un schéma affine.
Soit $\mathcal{F} = \widetilde{M}$ pour un $R$-module $M$.
Le faisceau quasi-cohérent $\mathcal{F}$ est un
$\mathcal{O}_X$-module plat si et seulement si $M$ est un $R$-module plat.
\end{lemma}

\begin{proof}
La platitude de $\mathcal{F}$ se vérifie sur les germes, voir
Modules, lemme \ref{modules-lemma-flat-stalks-flat}.
Il en va de même pour les modules sur un anneau, voir
Algèbre, lemme \ref{algebra-lemma-flat-localization}.
Puisque $\mathcal{F}_x = M_{\mathfrak p}$ lorsque $x$ correspond
à $\mathfrak p$, le lemme en résulte.
\end{proof}





\section{Modules localement libres}
\label{properties-section-finite-locally-free}

\noindent
Sur tout espace annelé, on sait ce que signifie, pour un $\mathcal{O}_X$-module,
être localement libre (de rang fini). Sur un schéma affine, cette notion coïncide avec
celle définie dans le chapitre d'algèbre.

\begin{lemma}
\label{properties-lemma-locally-free-module}
Soit $X = \Spec(R)$ un schéma affine.
Soit $\mathcal{F} = \widetilde{M}$ pour un $R$-module $M$.
Le faisceau quasi-cohérent $\mathcal{F}$ est un
$\mathcal{O}_X$-module localement libre (de rang fini) si et seulement si $M$ est
un $R$-module localement libre (de rang fini).
\end{lemma}

\begin{proof}
Cela résulte des définitions :
voir Modules, définition \ref{modules-definition-locally-free},
ainsi que
Algèbre, définition \ref{algebra-definition-locally-free}.
\end{proof}

\noindent
On peut caractériser les modules localement libres de rang fini de plusieurs façons.

\begin{lemma}
\label{properties-lemma-finite-locally-free}
Soit $X$ un schéma.
Soit $\mathcal{F}$ un $\mathcal{O}_X$-module
quasi-cohérent. Les assertions suivantes sont équivalentes :
\begin{enumerate}
\item $\mathcal{F}$ est un $\mathcal{O}_X$-module plat de présentation finie ;
\item $\mathcal{F}$ est un $\mathcal{O}_X$-module de présentation finie et,
pour tout $x \in X$, le germe $\mathcal{F}_x$ est un
$\mathcal{O}_{X, x}$-module libre ;
\item $\mathcal{F}$ est un $\mathcal{O}_X$-module localement libre et de type fini ;
\item $\mathcal{F}$ est un $\mathcal{O}_X$-module localement libre de rang fini ;
\item $\mathcal{F}$ est un $\mathcal{O}_X$-module de type fini,
pour tout $x \in X$, le germe $\mathcal{F}_x$ est un
$\mathcal{O}_{X, x}$-module libre, et la fonction
$$
\rho_\mathcal{F} : X \to \mathbf{Z}, \quad
x \longmapsto
\dim_{\kappa(x)} \mathcal{F}_x \otimes_{\mathcal{O}_{X, x}} \kappa(x)
$$
est localement constante pour la topologie de Zariski de $X$.
\end{enumerate}
\end{lemma}

\begin{proof}
On se ramène immédiatement au cas affine. Dans
ce cas, le lemme est une reformulation de
Algèbre, lemme \ref{algebra-lemma-finite-projective}.
Le passage d'une formulation à l'autre
utilise les lemmes \ref{properties-lemma-finite-type-module},
\ref{properties-lemma-finite-presentation-module},
\ref{properties-lemma-flat-module}et
\ref{properties-lemma-locally-free-module}.
\end{proof}

\begin{lemma}
\label{properties-lemma-finite-locally-free-reduced}
Soit $X$ un schéma réduit. Soit $\mathcal{F}$ un
$\mathcal{O}_X$-module quasi-cohérent. Les conditions équivalentes du
lemme \ref{properties-lemma-finite-locally-free} sont alors aussi équivalentes à la suivante :
\begin{enumerate}
\item[(6)] $\mathcal{F}$ est un $\mathcal{O}_X$-module de type fini et
la fonction
$$
\rho_\mathcal{F} : X \to \mathbf{Z}, \quad
x \longmapsto
\dim_{\kappa(x)} \mathcal{F}_x \otimes_{\mathcal{O}_{X, x}} \kappa(x)
$$
est localement constante pour la topologie de Zariski de $X$.
\end{enumerate}
\end{lemma}

\begin{proof}
On se ramène immédiatement au cas affine. Dans
ce cas, le lemme est une reformulation de
Algèbre, lemme \ref{algebra-lemma-finite-projective-reduced}.
\end{proof}




\section{Modules localement projectifs}
\label{properties-section-locally-projective}

\noindent
Les résultats établis dans le chapitre d'algèbre permettent
de définir les modules localement projectifs comme suit.

\begin{definition}
\label{properties-definition-locally-projective}
Soit $X$ un schéma. Soit $\mathcal{F}$ un
$\mathcal{O}_X$-module quasi-cohérent. On dit que $\mathcal{F}$ est {\it localement projectif}
si, pour tout ouvert affine $U \subset X$, le $\mathcal{O}_X(U)$-module
$\mathcal{F}(U)$ est projectif.
\end{definition}

\begin{lemma}
\label{properties-lemma-locally-projective}
Soit $X$ un schéma.
Soit $\mathcal{F}$ un $\mathcal{O}_X$-module
quasi-cohérent. Les assertions suivantes sont équivalentes :
\begin{enumerate}
\item $\mathcal{F}$ est localement projectif ;
\item il existe un recouvrement ouvert affine $X = \bigcup U_i$
tel que le $\mathcal{O}_X(U_i)$-module
$\mathcal{F}(U_i)$ soit projectif pour tout $i$.
\end{enumerate}
En particulier, si $X = \Spec(A)$ et $\mathcal{F} = \widetilde{M}$,
alors $\mathcal{F}$ est localement projectif si et seulement si $M$ est un
$A$-module projectif.
\end{lemma}

\begin{proof}
Remarquons d'abord que, si $M$ est un $A$-module projectif et si $A \to B$ est un
homomorphisme d'anneaux, alors $M \otimes_A B$ est un $B$-module projectif,
voir Algèbre, lemme \ref{algebra-lemma-ascend-properties-modules}.
Ainsi, si $U$ est un ouvert affine tel que $\mathcal{F}(U)$ soit un
$\mathcal{O}_X(U)$-module projectif, l'ouvert principal $D(f)$ est
un ouvert affine tel que $\mathcal{F}(D(f))$ soit un
$\mathcal{O}_X(D(f))$-module projectif, pour tout $f \in \mathcal{O}_X(U)$.
Supposons (2). Soit $U \subset X$ un ouvert affine quelconque. On
peut trouver un recouvrement ouvert $U = \bigcup_{j = 1, \ldots, m} D(f_j)$
par un nombre fini d'ouverts principaux $D(f_j)$ tels que, pour tout
$j$, l'ouvert $D(f_j)$ soit un ouvert principal d'un certain $U_i$, voir
Schémas, lemme \ref{schemes-lemma-standard-open-two-affines}.
Ainsi, si l'on pose $A = \mathcal{O}_X(U)$ et si $M$ est un $A$-module
tel que $\mathcal{F}|_U$ corresponde à $M$, alors
$M_{f_j}$ est un $A_{f_j}$-module projectif. Il en résulte que
$A \to B = \prod A_{f_j}$ est un homomorphisme d'anneaux fidèlement
plat tel que $M \otimes_A B$ soit un $B$-module
projectif. Donc $M$ est projectif
d'après Algèbre, théorème \ref{algebra-theorem-ffdescent-projectivity}.
\end{proof}

\begin{lemma}
\label{properties-lemma-locally-projective-pullback}
Soit $f : X \to Y$ un morphisme de schémas.
Soit $\mathcal{G}$ un $\mathcal{O}_Y$-module
quasi-cohérent. Si $\mathcal{G}$ est localement projectif sur $Y$, alors $f^*\mathcal{G}$
est localement projectif sur $X$.
\end{lemma}

\begin{proof}
Cela résulte
de Algèbre, lemme \ref{algebra-lemma-ascend-properties-modules},
et du lemme \ref{properties-lemma-locally-projective}.
\end{proof}






\section{Prolongement des faisceaux quasi-cohérents}
\label{properties-section-extending-quasi-coherent-sheaves}

\noindent
Il est parfois utile de montrer qu'un faisceau quasi-cohérent donné sur
un sous-schéma ouvert se prolonge au schéma tout entier.

\begin{lemma}
\label{properties-lemma-extend-trivial}
Soit $j : U \to X$ une immersion ouverte quasi-compacte de schémas.
\begin{enumerate}
\item Tout faisceau quasi-cohérent sur $U$ se prolonge en un faisceau
quasi-cohérent sur $X$.
\item Soit $\mathcal{F}$ un faisceau quasi-cohérent sur $X$.
Soit $\mathcal{G} \subset \mathcal{F}|_U$ un sous-faisceau
quasi-cohérent. Il existe un sous-faisceau quasi-cohérent $\mathcal{H}$ de
$\mathcal{F}$ tel que $\mathcal{H}|_U = \mathcal{G}$
comme sous-faisceaux de $\mathcal{F}|_U$.
\item Soit $\mathcal{F}$ un faisceau quasi-cohérent sur $X$.
Soit $\mathcal{G}$ un faisceau quasi-cohérent sur $U$.
Soit $\varphi : \mathcal{G} \to \mathcal{F}|_U$ un morphisme
de $\mathcal{O}_U$-modules. Il existe un faisceau quasi-cohérent $\mathcal{H}$
de $\mathcal{O}_X$-modules et un morphisme $\psi : \mathcal{H} \to \mathcal{F}$
tels que $\mathcal{H}|_U = \mathcal{G}$ et que
$\psi|_U = \varphi$.
\end{enumerate}
\end{lemma}

\begin{proof}
Une immersion est séparée
(voir Schémas, lemme \ref{schemes-lemma-immersions-monomorphisms}),
et $j$ est quasi-compact par hypothèse.
Ainsi, pour tout faisceau quasi-cohérent $\mathcal{G}$ sur $U$, le faisceau
$j_*\mathcal{G}$ est un prolongement à $X$. Voir
Schémas, lemme \ref{schemes-lemma-push-forward-quasi-coherent}, et
Faisceaux, section \ref{sheaves-section-open-immersions}.

\medskip\noindent
Supposons $\mathcal{F}$ et $\mathcal{G}$ comme dans (2).
Alors $j_*\mathcal{G}$ est un faisceau quasi-cohérent sur $X$ (voir ci-dessus).
C'est un sous-faisceau de $j_*j^*\mathcal{F}$.
Le noyau
$$
\mathcal{H} =
\Ker(\mathcal{F} \oplus j_* \mathcal{G}
\longrightarrow j_*j^*\mathcal{F})
$$
est donc lui aussi quasi-cohérent,
voir Schémas, section \ref{schemes-section-quasi-coherent}.
On vérifie formellement que $\mathcal{H} \subset \mathcal{F}$ et que
$\mathcal{H}|_U = \mathcal{G}$ (en utilisant de nouveau les
résultats de Faisceaux, section \ref{sheaves-section-open-immersions} ).

\medskip\noindent
L'assertion (3) se démontre de la même manière que (2). Il suffit de prendre
$\mathcal{H} = \Ker(\mathcal{F} \oplus j_* \mathcal{G}
\to j_*j^*\mathcal{F})$, muni de son morphisme évident vers $\mathcal{F}$
et de son identification évidente avec $\mathcal{G}$ sur $U$.
\end{proof}

\begin{lemma}
\label{properties-lemma-extend}
Soit $X$ un schéma quasi-compact et quasi-séparé.
Soit $U \subset X$ un ouvert quasi-compact.
Soit $\mathcal{F}$ un $\mathcal{O}_X$-module
quasi-cohérent. Soit $\mathcal{G} \subset \mathcal{F}|_U$ un
$\mathcal{O}_U$-sous-module quasi-cohérent de type fini.
Alors il existe un sous-module quasi-cohérent $\mathcal{G}' \subset \mathcal{F}$
de type fini tel que $\mathcal{G}'|_U = \mathcal{G}$.
\end{lemma}

\begin{proof}
Soit $n$ le nombre minimal d'ouverts affines $U_i \subset X$,
$i = 1, \ldots, n$, tels que $X = U \cup \bigcup U_i$.
(On utilise ici la quasi-compacité de $X$.) Supposons que l'on
sache démontrer le lemme dans le cas $n = 1$. On peut alors prolonger
successivement $\mathcal{G}$
en un faisceau $\mathcal{G}_1$ sur $U \cup U_1$,
puis en un faisceau $\mathcal{G}_2$ sur $U \cup U_1 \cup U_2$,
puis en un faisceau $\mathcal{G}_3$ sur $U \cup U_1 \cup U_2 \cup U_3$,
et ainsi de
suite. On se ramène donc au cas $n = 1$.

\medskip\noindent
On peut ainsi supposer que $X = U \cup V$, avec $V$ affine.
Comme $X$ est quasi-séparé et que $U$, $V$ sont des ouverts
quasi-compacts, $U \cap V$ est un ouvert quasi-compact. Il suffit de démontrer le
lemme pour le système $(V, U \cap V, \mathcal{F}|_V, \mathcal{G}|_{U \cap V})$,
car on peut recoller le faisceau $\mathcal{G}'$ ainsi obtenu sur $V$
au faisceau donné $\mathcal{G}$ sur $U$, le long de leur restriction
commune à $U \cap V$.
On se ramène donc au cas où $X$ est affine.

\medskip\noindent
Supposons $X = \Spec(R)$. Écrivons $\mathcal{F} = \widetilde M$
pour un $R$-module $M$. Le lemme \ref{properties-lemma-extend-trivial} ci-dessus permet
de trouver un sous-faisceau quasi-cohérent $\mathcal{H} \subset \mathcal{F}$
qui se restreigne à $\mathcal{G}$ sur $U$.
Écrivons $\mathcal{H} = \widetilde N$ pour un $R$-module $N$.
Pour tout $u \in U$, il existe un $f \in R$ tel que
$u \in D(f) \subset U$ et que $N_f$ soit de type fini, voir
le lemme \ref{properties-lemma-finite-type-module}.
Comme $U$ est quasi-compact, on peut le recouvrir par un nombre fini
d'ouverts $D(f_i)$ tels que $N_{f_i}$ soit engendré par
un nombre fini d'éléments, notés $x_{i, 1}/f_i^N, \ldots, x_{i, r_i}/f_i^N$.
Soit $N' \subset N$ le sous-module engendré par les éléments
$x_{i, j}$. Le sous-faisceau
$\mathcal{G}' = \widetilde{N'} \subset \mathcal{H} \subset \mathcal{F}$
convient.
\end{proof}

\begin{lemma}
\label{properties-lemma-quasi-coherent-colimit-finite-type}
Soit $X$ un schéma quasi-compact et quasi-séparé. Tout
faisceau quasi-cohérent de $\mathcal{O}_X$-modules
est la colimite filtrante de ses
$\mathcal{O}_X$-sous-modules quasi-cohérents de type fini.
\end{lemma}

\begin{proof}
La colimite est filtrante, car, si $\mathcal{G}_1$ et $\mathcal{G}_2$
sont des sous-faisceaux quasi-cohérents de type fini, l'image de
$\mathcal{G}_1 \oplus \mathcal{G}_2 \to \mathcal{F}$ est
un sous-module quasi-cohérent de type fini.
Soit $U \subset X$ un ouvert affine quelconque, et soit
$s \in \Gamma(U, \mathcal{F})$ une section quelconque.
Soit $\mathcal{G} \subset \mathcal{F}|_U$ le
sous-faisceau engendré par $s$. Alors $\mathcal{G}$
est évidemment quasi-cohérent et de type fini comme $\mathcal{O}_U$-module.
Le lemme \ref{properties-lemma-extend} montre que $\mathcal{G}$ est la restriction
d'un sous-faisceau quasi-cohérent $\mathcal{G}' \subset \mathcal{F}$
de type fini. Comme $X$ possède une base de la topologie formée
d'ouverts affines, on en déduit que toute section locale de
$\mathcal{F}$ appartient localement à un sous-module quasi-cohérent de type
fini. Cela conclut la démonstration.
\end{proof}

\begin{lemma}
\label{properties-lemma-extend-finite-presentation}
Soit $X$ un schéma quasi-compact et quasi-séparé.
Soit $\mathcal{F}$ un $\mathcal{O}_X$-module
quasi-cohérent. Soit $U \subset X$ un ouvert quasi-compact.
Soit $\mathcal{G}$ un $\mathcal{O}_U$-module de présentation finie.
Soit $\varphi : \mathcal{G} \to \mathcal{F}|_U$ un morphisme de
$\mathcal{O}_U$-modules.
Il existe alors un $\mathcal{O}_X$-module
$\mathcal{G}'$ de présentation finie et un morphisme
de $\mathcal{O}_X$-modules $\varphi' : \mathcal{G}' \to \mathcal{F}$
tels que $\mathcal{G}'|_U = \mathcal{G}$ et que
$\varphi'|_U = \varphi$.
\end{lemma}

\begin{proof}
Le début de la démonstration reprend celui de la
démonstration du lemme \ref{properties-lemma-extend}. Nous le détaillons néanmoins.

\medskip\noindent
Soit $n$ le nombre minimal d'ouverts affines $U_i \subset X$,
$i = 1, \ldots, n$, tels que $X = U \cup \bigcup U_i$.
(On utilise ici la quasi-compacité de $X$.) Supposons que l'on
sache démontrer le lemme dans le cas $n = 1$. On peut alors prolonger
successivement le couple $(\mathcal{G}, \varphi)$
en un couple $(\mathcal{G}_1, \varphi_1)$ sur $U \cup U_1$,
puis en un couple $(\mathcal{G}_2, \varphi_2)$ sur $U \cup U_1 \cup U_2$,
puis en un couple $(\mathcal{G}_3, \varphi_3)$ sur $U \cup U_1 \cup U_2 \cup U_3$,
et ainsi de
suite. On se ramène donc au cas $n = 1$.

\medskip\noindent
On peut ainsi supposer que $X = U \cup V$, avec $V$ affine.
Comme $X$ est quasi-séparé et que $U$ est
quasi-compact, $U \cap V \subset V$ est quasi-compact.
Supposons que l'on démontre le lemme pour le système
$(V, U \cap V, \mathcal{F}|_V, \mathcal{G}|_{U \cap V}, \varphi|_{U \cap V})$,
en obtenant un couple $(\mathcal{G}', \varphi')$ sur $V$.
On peut alors recoller $\mathcal{G}'$ sur $V$ au faisceau donné $\mathcal{G}$
sur $U$, le long de leur restriction commune à $U \cap V$; de même, on peut recoller le
morphisme $\varphi'$ au morphisme $\varphi$ le long de leur restriction commune à
$U \cap V$. On se ramène donc au cas où $X$ est affine.

\medskip\noindent
Supposons $X = \Spec(R)$.
Le lemme \ref{properties-lemma-extend-trivial} ci-dessus permet
de trouver un faisceau quasi-cohérent $\mathcal{H}$ muni
d'un morphisme $\psi : \mathcal{H} \to \mathcal{F}$ sur $X$,
qui se restreignent à $\mathcal{G}$ et à $\varphi$ sur $U$.
Le lemme \ref{properties-lemma-extend} permet de
trouver un $\mathcal{O}_X$-sous-module quasi-cohérent de type fini
$\mathcal{H}' \subset \mathcal{H}$
tel que $\mathcal{H}'|_U = \mathcal{G}$. Ainsi, en
remplaçant $\mathcal{H}$ par $\mathcal{H}'$
et $\psi$ par la restriction de $\psi$ à $\mathcal{H}'$,
on peut supposer $\mathcal{H}$ de type fini. D'après
le lemme \ref{properties-lemma-finite-presentation-module},
on a $\mathcal{H} = \widetilde{N}$, avec
$N$ un $R$-module de type fini. Il existe donc une
surjection figurant dans la suite exacte
courte suivante de $\mathcal{O}_X$-modules quasi-cohérents :
$$
0 \to \mathcal{K} \to \mathcal{O}_X^{\oplus n} \to \mathcal{H} \to 0
$$
où $\mathcal{K}$ est défini comme le noyau.
Comme $\mathcal{G}$ est de présentation finie et que
$\mathcal{H}|_U = \mathcal{G}$,
d'après Modules, lemme
\ref{modules-lemma-kernel-surjection-finite-free-onto-finite-presentation},
la restriction $\mathcal{K}|_U$ est
un $\mathcal{O}_U$-module de type fini. Ainsi, de nouveau par le lemme \ref{properties-lemma-extend},
il existe un
$\mathcal{O}_X$-sous-module quasi-cohérent de type fini $\mathcal{K}' \subset \mathcal{K}$ tel
que $\mathcal{K}'|_U = \mathcal{K}|_U$. Pour résoudre le problème
posé dans le lemme, on prend
$$
\mathcal{G}' = \mathcal{O}_X^{\oplus n}/\mathcal{K}'
$$
qui est évidemment de présentation finie et dont la restriction est $\mathcal{G}$
sur $U$, avec $\varphi'$ égal au composé
$$
\mathcal{G}' = \mathcal{O}_X^{\oplus n}/\mathcal{K}'
\to \mathcal{O}_X^{\oplus n}/\mathcal{K} = \mathcal{H} \xrightarrow{\psi}
\mathcal{F}.
$$
Cela achève la démonstration du lemme.
\end{proof}

\begin{lemma}
\label{properties-lemma-lift-finite-presentation}
Soit $X$ un schéma quasi-compact et quasi-séparé. Soit $U \subset X$
un ouvert quasi-compact. Soit $\mathcal{G}$ un $\mathcal{O}_U$-module.
\begin{enumerate}
\item Si $\mathcal{G}$ est quasi-cohérent et de
type fini, il existe un $\mathcal{O}_X$-module quasi-cohérent $\mathcal{G}'$
de type fini tel que $\mathcal{G}'|_U = \mathcal{G}$.
\item Si $\mathcal{G}$ est de présentation finie,
il existe un $\mathcal{O}_X$-module $\mathcal{G}'$
de présentation finie tel que $\mathcal{G}'|_U = \mathcal{G}$.
\end{enumerate}
\end{lemma}

\begin{proof}
L'assertion (2) est le cas particulier du lemme \ref{properties-lemma-extend-finite-presentation}
où $\mathcal{F} = 0$. Pour (1), on écrit d'abord
$\mathcal{G} = \mathcal{F}|_U$ pour un $\mathcal{O}_X$-module
quasi-cohérent, d'après le lemme \ref{properties-lemma-extend-trivial},
puis on applique le lemme \ref{properties-lemma-extend} avec $\mathcal{G} = \mathcal{F}|_U$.
\end{proof}

\noindent
Le lemme suivant affirme que tout faisceau quasi-cohérent sur un schéma
quasi-compact et quasi-séparé est une colimite filtrante de $\mathcal{O}$-modules
de présentation finie. Dans le lemme qui suit, nous reformulerons ce résultat en termes de
colimites indexées par des ensembles ordonnés filtrants, formulation peut-être plus familière.

\begin{lemma}
\label{properties-lemma-directed-colimit-diagram-finite-presentation}
\begin{slogan}
Les modules quasi-cohérents sur les schémas quasi-compacts et quasi-séparés
sont des colimites filtrantes de modules de présentation finie.
\end{slogan}
Soit $X$ un schéma. Supposons $X$ quasi-compact et quasi-séparé.
Soit $\mathcal{F}$ un $\mathcal{O}_X$-module
quasi-cohérent. Il existe
\begin{enumerate}
\item une catégorie d'indices filtrante $\mathcal{I}$ (voir
Catégories, définition \ref{categories-definition-directed}) ;
\item un diagramme $\mathcal{I} \to \textit{Mod}(\mathcal{O}_X)$ (voir
Catégories, section \ref{categories-section-limits}),
$i \mapsto \mathcal{F}_i$;
\item des morphismes de $\mathcal{O}_X$-modules
$\varphi_i : \mathcal{F}_i \to \mathcal{F}$
\end{enumerate}
tels que chaque $\mathcal{F}_i$ soit de présentation
finie et que les morphismes $\varphi_i$ induisent un isomorphisme
$$
\colim_i \mathcal{F}_i
=
\mathcal{F}.
$$
\end{lemma}

\begin{proof}
Choisissons un ensemble $I$ et, pour tout $i \in I$, un $\mathcal{O}_X$-module
de présentation finie muni d'un homomorphisme de $\mathcal{O}_X$-modules
$\varphi_i : \mathcal{F}_i \to \mathcal{F}$, de façon à satisfaire la
propriété suivante : pour tout $\psi : \mathcal{G} \to \mathcal{F}$ avec $\mathcal{G}$
de présentation finie, il existe un $i \in I$ et un
isomorphisme $\alpha : \mathcal{F}_i \to \mathcal{G}$ tels que
$\varphi_i = \psi \circ \alpha$. Il résulte clairement
de Modules, lemme \ref{modules-lemma-set-isomorphism-classes-finite-type-modules},
qu'un tel ensemble existe (voir aussi sa
démonstration). Notons $\mathcal{I}$ la catégorie
dont $\Ob(\mathcal{I}) = I$ et dont, pour $i, i' \in I$,
les ensembles de morphismes sont définis par
$$
\Mor_\mathcal{I}(i, i') =
\{\alpha : \mathcal{F}_i \to \mathcal{F}_{i'} \mid
\alpha \circ \varphi_{i'} = \varphi_i
\}.
$$
Nous affirmons que $\mathcal{I}$ est une catégorie filtrante et que
$\mathcal{F} = \colim_i \mathcal{F}_i$.

\medskip\noindent
Soient $i, i' \in I$. Considérons le morphisme
$$
\mathcal{F}_i \oplus \mathcal{F}_{i'} \longrightarrow \mathcal{F}
$$
qui est la somme directe de $\varphi_i$ et de $\varphi_{i'}$.
Comme une somme directe de $\mathcal{O}_X$-modules
de présentation finie est de présentation finie, il existe un $i'' \in I$
tel que $\varphi_{i''} : \mathcal{F}_{i''} \to \mathcal{F}$
soit isomorphe à la flèche affichée ci-dessus vers $\mathcal{F}$.
Puisque l'on dispose des diagrammes commutatifs
$$
\xymatrix{
\mathcal{F}_i \ar[r] \ar[d] & \mathcal{F} \ar@{=}[d] \\
\mathcal{F}_i \oplus \mathcal{F}_{i'} \ar[r] & \mathcal{F}
}
\quad
\text{et}
\quad
\xymatrix{
\mathcal{F}_{i'} \ar[r] \ar[d] & \mathcal{F} \ar@{=}[d] \\
\mathcal{F}_i \oplus \mathcal{F}_{i'} \ar[r] & \mathcal{F}
}
$$
il existe des morphismes $i \to i''$ et $i' \to i''$
dans $\mathcal{I}$. Supposons maintenant donnés $i, i' \in I$ et
des morphismes $\alpha, \beta : i \to i'$ (correspondant à des morphismes de $\mathcal{O}_X$-modules
 $\alpha, \beta : \mathcal{F}_i \to \mathcal{F}_{i'}$).
Considérons leur coégalisateur
$$
\mathcal{G} =
\Coker(
\mathcal{F}_i \xrightarrow{\alpha - \beta} \mathcal{F}_{i'}
)
$$
Remarquons que $\mathcal{G}$ est un $\mathcal{O}_X$-module de présentation finie.
Puisque, par définition des morphismes de la catégorie $\mathcal{I}$,
on a $\varphi_{i'} \circ \alpha = \varphi_{i'} \circ \beta$,
on obtient un morphisme induit $\psi : \mathcal{G} \to \mathcal{F}$.
Le couple $(\mathcal{G}, \psi)$ est donc encore isomorphe
à un couple $(\mathcal{F}_{i''}, \varphi_{i''})$ pour un certain $i''$.
Il existe donc un morphisme $i' \to i''$ dans
$\mathcal{I}$ qui égalise $\alpha$ et $\beta$. Nous avons
ainsi montré que la catégorie $\mathcal{I}$ est filtrante.

\medskip\noindent
Il reste à montrer que la colimite du diagramme est $\mathcal{F}$.
Par définition de la colimite et par notre définition de la catégorie
$\mathcal{I}$, on dispose d'un morphisme canonique
$$
\varphi :
\colim_i \mathcal{F}_i
\longrightarrow
\mathcal{F}.
$$
Soit $x \in X$. Montrons que $\varphi_x$ est un isomorphisme.
Rappelons que
$$
(\colim_i \mathcal{F}_i)_x
=
\colim_i \mathcal{F}_{i, x},
$$
voir
Faisceaux, section \ref{sheaves-section-limits-sheaves}.
Montrons d'abord que $\varphi_x$ est
injectif. Soit $s \in \mathcal{F}_{i, x}$ un élément
tel que $s$ ait une image nulle dans $\mathcal{F}_x$. Il existe alors un
ouvert quasi-compact $U$ tel que $s$ provienne d'une section $s \in \mathcal{F}_i(U)$
et que $\varphi_i(s) = 0$ dans $\mathcal{F}(U)$.
D'après le lemme \ref{properties-lemma-extend},
on peut trouver un sous-faisceau quasi-cohérent de type fini
$\mathcal{K} \subset \Ker(\varphi_i)$ dont la
restriction soit le $\mathcal{O}_U$-sous-module quasi-cohérent de $\mathcal{F}_i$
engendré par $s$:
$\mathcal{K}|_U = \mathcal{O}_U\cdot s \subset \mathcal{F}_i|_U$.
Il est clair que $\mathcal{F}_i/\mathcal{K}$ est de présentation finie et que le
morphisme $\varphi_i$ se factorise par le morphisme quotient
$\mathcal{F}_i \to \mathcal{F}_i/\mathcal{K}$. On peut donc trouver
un $i' \in I$ et un morphisme $\alpha : \mathcal{F}_i \to \mathcal{F}_{i'}$
dans $\mathcal{I}$ qui s'identifie au morphisme quotient
$\mathcal{F}_i \to \mathcal{F}_i/\mathcal{K}$. Il en résulte
que la section $s$ a une image nulle dans $\mathcal{F}_{i'}(U)$ et,
en particulier, dans
$(\colim_i \mathcal{F}_i)_x =
\colim_i \mathcal{F}_{i, x}$.
L'injectivité en
découle. Montrons enfin que $\varphi_x$ est surjectif.
Soit $s \in \mathcal{F}_x$. Choisissons un voisinage ouvert quasi-compact
$U \subset X$ de $x$ tel que $s$ corresponde à une section
$s \in \mathcal{F}(U)$. Considérons le morphisme
$s : \mathcal{O}_U \to \mathcal{F}$ (multiplication par $s$).
D'après le lemme \ref{properties-lemma-extend-finite-presentation},
il existe un $\mathcal{O}_X$-module $\mathcal{G}$
de présentation finie et un morphisme de $\mathcal{O}_X$-modules
$\mathcal{G} \to \mathcal{F}$ tels que $\mathcal{G}|_U \to \mathcal{F}|_U$
s'identifie à
$s : \mathcal{O}_U \to \mathcal{F}$.
De nouveau, par définition de $\mathcal{I}$, il existe un $i \in I$
tel que $\mathcal{G} \to \mathcal{F}$ soit isomorphe à
$\varphi_i : \mathcal{F}_i \to \mathcal{F}$. Il existe évidemment
une section $s' \in \mathcal{F}_i(U)$ d'image $s \in \mathcal{F}(U)$.
Cela prouve la surjectivité et achève la démonstration du lemme.
\end{proof}

\begin{lemma}
\label{properties-lemma-directed-colimit-finite-presentation}
Soit $X$ un schéma. Supposons $X$ quasi-compact et quasi-séparé.
Soit $\mathcal{F}$ un $\mathcal{O}_X$-module
quasi-cohérent. Il existe
\begin{enumerate}
\item un ensemble ordonné filtrant $I$ (voir
Catégories, définition \ref{categories-definition-directed-set}) ;
\item un système $(\mathcal{F}_i, \varphi_{ii'})$
sur $I$ dans $\textit{Mod}(\mathcal{O}_X)$ (voir
Catégories, définition \ref{categories-definition-system-over-poset}) ;
\item des morphismes de $\mathcal{O}_X$-modules
$\varphi_i : \mathcal{F}_i \to \mathcal{F}$
\end{enumerate}
tels que chaque $\mathcal{F}_i$ soit de présentation
finie et que les morphismes $\varphi_i$ induisent un isomorphisme
$$
\colim_i \mathcal{F}_i
=
\mathcal{F}.
$$
\end{lemma}

\begin{proof}
C'est une conséquence directe du
lemme \ref{properties-lemma-directed-colimit-diagram-finite-presentation} et
de Catégories, lemme \ref{categories-lemma-directed-category-system},
combinés au fait que les
colimites existent dans la catégorie des faisceaux de $\mathcal{O}_X$-modules, voir
Faisceaux, section \ref{sheaves-section-limits-sheaves}.
\end{proof}

\begin{lemma}
\label{properties-lemma-finite-directed-colimit-surjective-maps}
Soit $X$ un schéma. Supposons $X$ quasi-compact et quasi-séparé.
Soit $\mathcal{F}$ un $\mathcal{O}_X$-module
quasi-cohérent de type fini. On peut alors écrire $\mathcal{F} = \colim \mathcal{F}_i$, avec $\mathcal{F}_i$
de présentation finie et tous les morphismes de transition
$\mathcal{F}_i \to \mathcal{F}_{i'}$ surjectifs.
\end{lemma}

\begin{proof}
Écrivons $\mathcal{F} = \colim \mathcal{G}_i$ comme une colimite
filtrante de $\mathcal{O}_X$-modules de
présentation finie (lemme \ref{properties-lemma-directed-colimit-finite-presentation}).
Nous affirmons que $\mathcal{G}_i \to \mathcal{F}$ est surjectif pour un certain $i$.
En effet, choisissons un recouvrement ouvert affine fini $X = U_1 \cup \ldots \cup U_m$.
Choisissons des sections $s_{jl} \in \mathcal{F}(U_j)$ engendrant
$\mathcal{F}|_{U_j}$, voir le lemme \ref{properties-lemma-finite-type-module}.
D'après Faisceaux, lemme \ref{sheaves-lemma-directed-colimits-sections},
$s_{jl}$ appartient à l'image de $\mathcal{G}_i \to \mathcal{F}$
pour $i$ assez grand. Ainsi $\mathcal{G}_i \to \mathcal{F}$ est surjectif
pour $i$ assez grand. Choisissons un tel $i$ et soit
$\mathcal{K} \subset \mathcal{G}_i$ le noyau du morphisme
$\mathcal{G}_i \to \mathcal{F}$. Écrivons $\mathcal{K} = \colim \mathcal{K}_a$
comme la colimite filtrante de ses sous-modules quasi-cohérents de type fini
(lemme \ref{properties-lemma-quasi-coherent-colimit-finite-type}). Alors
$\mathcal{F} = \colim \mathcal{G}_i/\mathcal{K}_a$ résout le
problème posé dans le lemme.
\end{proof}

\begin{lemma}
\label{properties-lemma-application-directed-colimit}
Soit $X$ un schéma quasi-compact et quasi-séparé.
Soit $\mathcal{F}$ un $\mathcal{O}_X$-module quasi-cohérent de type
fini. Soit $U \subset X$ un ouvert quasi-compact tel que $\mathcal{F}|_U$
soit de présentation finie. Il existe alors un morphisme de $\mathcal{O}_X$-modules
$\varphi : \mathcal{G} \to \mathcal{F}$ tel
que (a) $\mathcal{G}$ soit de présentation finie,
(b) $\varphi$ soit surjectif et
(c) $\varphi|_U$ soit un isomorphisme.
\end{lemma}

\begin{proof}
Écrivons $\mathcal{F} = \colim \mathcal{F}_i$ comme une colimite filtrante,
avec chaque $\mathcal{F}_i$ de présentation finie,
voir le lemme \ref{properties-lemma-directed-colimit-finite-presentation}.
Choisissons un recouvrement ouvert affine fini $X = \bigcup V_j$ et
choisissons un nombre fini de sections $s_{jl} \in \mathcal{F}(V_j)$ engendrant
$\mathcal{F}|_{V_j}$, voir le lemme \ref{properties-lemma-finite-type-module}.
D'après Faisceaux, lemme \ref{sheaves-lemma-directed-colimits-sections},
$s_{jl}$ appartient à l'image de $\mathcal{F}_i \to \mathcal{F}$
pour $i$ assez grand. Ainsi $\mathcal{F}_i \to \mathcal{F}$ est surjectif
pour $i$ assez grand. Choisissons un tel $i$ et soit
$\mathcal{K} \subset \mathcal{F}_i$ le noyau du morphisme
$\mathcal{F}_i \to \mathcal{F}$. Comme $\mathcal{F}_U$ est de
présentation finie, $\mathcal{K}|_U$ est de type fini, voir
Modules, lemme
\ref{modules-lemma-kernel-surjection-finite-free-onto-finite-presentation}.
On peut donc trouver un sous-module quasi-cohérent de type fini
$\mathcal{K}' \subset \mathcal{K}$ avec $\mathcal{K}'|_U = \mathcal{K}|_U$,
voir le lemme \ref{properties-lemma-extend}. Alors
$\mathcal{G} = \mathcal{F}_i/\mathcal{K}'$,
muni du morphisme donné $\mathcal{G} \to \mathcal{F}$, convient.
\end{proof}

\noindent
Soit $X$ un schéma. Dans le lemme suivant, nous utilisons la
notion de {\it  $\mathcal{O}_X$-algèbre quasi-cohérente $\mathcal{A}$
de présentation finie}. Cela signifie que, pour tout ouvert affine
$\Spec(R) \subset X$, on a $\mathcal{A} = \widetilde{A}$,
où $A$ est une $R$-algèbre (commutative) de présentation finie
comme $R$-algèbre.

\begin{lemma}
\label{properties-lemma-algebra-directed-colimit-finite-presentation}
Soit $X$ un schéma. Supposons $X$ quasi-compact et quasi-séparé.
Soit $\mathcal{A}$ une $\mathcal{O}_X$-algèbre
quasi-cohérente. Il existe
\begin{enumerate}
\item un ensemble ordonné filtrant $I$ (voir
Catégories, définition \ref{categories-definition-directed-set}) ;
\item un système $(\mathcal{A}_i, \varphi_{ii'})$
sur $I$ dans la catégorie des $\mathcal{O}_X$-algèbres ;
\item des morphismes de $\mathcal{O}_X$-algèbres
$\varphi_i : \mathcal{A}_i \to \mathcal{A}$
\end{enumerate}
tels que chaque $\mathcal{A}_i$ soit une $\mathcal{O}_X$-algèbre
quasi-cohérente de présentation finie et que les morphismes $\varphi_i$
induisent un isomorphisme
$$
\colim_i \mathcal{A}_i
=
\mathcal{A}.
$$
\end{lemma}

\begin{proof}
Écrivons d'abord $\mathcal{A} = \colim_i \mathcal{F}_i$ comme une
colimite filtrante de faisceaux quasi-cohérents de présentation finie, comme dans
le lemme \ref{properties-lemma-directed-colimit-finite-presentation}.
Pour tout $i$, soit $\mathcal{B}_i = \text{Sym}(\mathcal{F}_i)$ l'algèbre
symétrique de $\mathcal{F}_i$ sur $\mathcal{O}_X$. Écrivons
$\mathcal{I}_i = \Ker(\mathcal{B}_i \to \mathcal{A})$. Écrivons
$\mathcal{I}_i = \colim_j \mathcal{F}_{i, j}$, où
$\mathcal{F}_{i, j}$ est un sous-module quasi-cohérent de type fini de
$\mathcal{I}_i$, voir
le lemme \ref{properties-lemma-quasi-coherent-colimit-finite-type}.
Définissons $\mathcal{I}_{i, j} \subset \mathcal{I}_i$
comme le $\mathcal{B}_i$-idéal engendré par $\mathcal{F}_{i, j}$.
Posons $\mathcal{A}_{i, j} = \mathcal{B}_i/\mathcal{I}_{i, j}$.
Alors $\mathcal{A}_{i, j}$ est une
$\mathcal{O}_X$-algèbre quasi-cohérente de présentation finie. Définissons $(i, j) \leq (i', j')$ par les conditions
$i \leq i'$ et le fait que le morphisme $\mathcal{B}_i \to \mathcal{B}_{i'}$
envoie l'idéal $\mathcal{I}_{i, j}$ dans l'idéal $\mathcal{I}_{i', j'}$.
Il est alors clair que $\mathcal{A} = \colim_{i, j} \mathcal{A}_{i, j}$.
\end{proof}

\noindent
Soit $X$ un schéma. Dans le lemme suivant, nous utilisons la
notion de {\it  $\mathcal{O}_X$-algèbre quasi-cohérente $\mathcal{A}$
de type fini}. Cela signifie que, pour tout ouvert affine
$\Spec(R) \subset X$, on a $\mathcal{A} = \widetilde{A}$,
où $A$ est une $R$-algèbre (commutative) de type fini
comme $R$-algèbre.

\begin{lemma}
\label{properties-lemma-algebra-directed-colimit-finite-type}
Soit $X$ un schéma. Supposons $X$ quasi-compact et quasi-séparé.
Soit $\mathcal{A}$ une $\mathcal{O}_X$-algèbre
quasi-cohérente. Alors $\mathcal{A}$ est la colimite
filtrante de ses $\mathcal{O}_X$-sous-algèbres quasi-cohérentes de type fini.
\end{lemma}

\begin{proof}
Si $\mathcal{A}_1, \mathcal{A}_2 \subset \mathcal{A}$ sont
des $\mathcal{O}_X$-sous-algèbres quasi-cohérentes de type fini, l'image de
$\mathcal{A}_1 \otimes_{\mathcal{O}_X} \mathcal{A}_2 \to \mathcal{A}$
est encore une $\mathcal{O}_X$-sous-algèbre quasi-cohérente de
type fini (certains détails sont omis), qui contient à la fois $\mathcal{A}_1$ et $\mathcal{A}_2$.
Cela montre que le système est filtrant. Pour
montrer que $\mathcal{A}$ en est la colimite, écrivons
$\mathcal{A} = \colim_i \mathcal{A}_i$ comme
une colimite filtrante de $\mathcal{O}_X$-algèbres quasi-cohérentes de présentation finie,
comme dans le lemme \ref{properties-lemma-algebra-directed-colimit-finite-presentation}.
Les images $\mathcal{A}'_i = \Im(\mathcal{A}_i \to \mathcal{A})$ sont
alors des sous-algèbres quasi-cohérentes de type fini de $\mathcal{A}$. Puisque
$\mathcal{A}$ en est la colimite, le résultat en découle.
\end{proof}

\noindent
Soit $X$ un schéma. Dans le lemme suivant, nous utilisons la
notion d' {\it algèbre finie (resp.\ entière) quasi-cohérente sur
$\mathcal{O}_X$, notée $\mathcal{A}$}. Cela signifie que, pour tout
ouvert affine $\Spec(R) \subset X$, on a
$\mathcal{A} = \widetilde{A}$, où $A$ est une $R$-algèbre
(commutative) finie (resp.\ entière) comme $R$-algèbre.

\begin{lemma}
\label{properties-lemma-finite-algebra-directed-colimit-finite-finitely-presented}
Soit $X$ un schéma. Supposons $X$ quasi-compact et quasi-séparé.
Soit $\mathcal{A}$ une $\mathcal{O}_X$-algèbre quasi-cohérente
finie. Alors $\mathcal{A} = \colim \mathcal{A}_i$ est une
colimite filtrante de $\mathcal{O}_X$-algèbres
quasi-cohérentes finies et de présentation finie, telle que tous les morphismes de transition $\mathcal{A}_{i'} \to \mathcal{A}_i$
soient surjectifs.
\end{lemma}

\begin{proof}
D'après le lemme \ref{properties-lemma-finite-directed-colimit-surjective-maps},
il existe un $\mathcal{O}_X$-module de présentation finie
$\mathcal{F}$ et une surjection $\mathcal{F} \to \mathcal{A}$.
La structure d'algèbre fournit une surjection
$$
\text{Sym}^*_{\mathcal{O}_X}(\mathcal{F}) \longrightarrow \mathcal{A}
$$
Notons $\mathcal{J}$ son noyau. Écrivons $\mathcal{J} = \colim \mathcal{E}_i$
comme une colimite filtrante de $\mathcal{O}_X$-sous-modules de type fini
$\mathcal{E}_i$ (lemme \ref{properties-lemma-quasi-coherent-colimit-finite-type}). Posons
$$
\mathcal{A}_i = \text{Sym}^*_{\mathcal{O}_X}(\mathcal{F})/(\mathcal{E}_i)
$$
où $(\mathcal{E}_i)$ désigne le faisceau d'idéaux engendré
par l'image de $\mathcal{E}_i \to \text{Sym}^*_{\mathcal{O}_X}(\mathcal{F})$.
Chaque $\mathcal{A}_i$ est alors une $\mathcal{O}_X$-algèbre de
présentation finie, les morphismes de transition sont
surjectifs et $\mathcal{A} = \colim \mathcal{A}_i$. Pour achever la démonstration, il
reste à montrer que $\mathcal{A}_i$ est une $\mathcal{O}_X$-algèbre finie
pour $i$ assez grand. Pour cela, choisissons un recouvrement
ouvert affine $X = U_1 \cup \ldots \cup U_m$. Prenons des générateurs
$f_{j, 1}, \ldots, f_{j, N_j} \in \Gamma(U_i, \mathcal{F})$.
Comme $\mathcal{A}(U_j)$ est une $\mathcal{O}_X(U_j)$-algèbre
finie, pour tout $k$, il existe un polynôme unitaire
$P_{j, k} \in \mathcal{O}(U_j)[T]$ tel que $P_{j, k}(f_{j, k})$
soit nul dans $\mathcal{A}(U_j)$. Puisque
$\mathcal{A} = \colim \mathcal{A}_i$ par construction,
on a $P_{j, k}(f_{j, k}) = 0$ dans $\mathcal{A}_i(U_j)$
pour tout $i$ assez grand. Pour de tels $i$, les algèbres
$\mathcal{A}_i$ sont finies.
\end{proof}

\begin{lemma}
\label{properties-lemma-integral-algebra-directed-colimit-finite}
Soit $X$ un schéma. Supposons $X$ quasi-compact et quasi-séparé.
Soit $\mathcal{A}$ une $\mathcal{O}_X$-algèbre quasi-cohérente
entière. Alors
\begin{enumerate}
\item $\mathcal{A}$ est la colimite
filtrante de ses $\mathcal{O}_X$-sous-algèbres quasi-cohérentes finies ;
\item $\mathcal{A}$ est une colimite
filtrante de $\mathcal{O}_X$-algèbres quasi-cohérentes finies et de présentation finie.
\end{enumerate}
\end{lemma}

\begin{proof}
D'après le lemme \ref{properties-lemma-algebra-directed-colimit-finite-type}, on a
$\mathcal{A} = \colim \mathcal{A}_i$, où
$\mathcal{A}_i \subset \mathcal{A}$ parcourt
les $\mathcal{O}_X$-algèbres
quasi-cohérentes de type fini. Toute $\mathcal{O}_X$-sous-algèbre quasi-cohérente de type
fini de $\mathcal{A}$ est finie (appliquer Algèbre, lemme
\ref{algebra-lemma-characterize-finite-in-terms-of-integral},
à $\mathcal{A}_i(U) \subset \mathcal{A}(U)$ pour les ouverts affines $U$
de $X$). Cela démontre (1).

\medskip\noindent
Pour démontrer (2), écrivons $\mathcal{A} = \colim \mathcal{F}_i$
comme une colimite de $\mathcal{O}_X$-modules de présentation finie à
l'aide du lemme \ref{properties-lemma-directed-colimit-finite-presentation}.
Pour tout $i$, soit $\mathcal{J}_i$ le noyau du morphisme
$$
\text{Sym}^*_{\mathcal{O}_X}(\mathcal{F}_i) \longrightarrow \mathcal{A}
$$
Pour $i' \geq i$, on a un morphisme induit $\mathcal{J}_i \to \mathcal{J}_{i'}$,
et l'on a $\mathcal{A} =
\colim \text{Sym}^*_{\mathcal{O}_X}(\mathcal{F}_i)/\mathcal{J}_i$.
De plus, les $\mathcal{O}_X$-algèbres quasi-cohérentes
$\text{Sym}^*_{\mathcal{O}_X}(\mathcal{F}_i)/\mathcal{J}_i$
sont finies (voir ci-dessus). Écrivons $\mathcal{J}_i = \colim \mathcal{E}_{ik}$
comme une colimite de $\mathcal{O}_X$-modules de présentation
finie. Étant donnés $i' \geq i$ et $k$, il existe un $k'$ tel que l'on
dispose d'un morphisme $\mathcal{E}_{ik} \to \mathcal{E}_{i'k'}$
rendant
$$
\xymatrix{
\mathcal{J}_i \ar[r] & \mathcal{J}_{i'} \\
\mathcal{E}_{ik} \ar[u] \ar[r] & \mathcal{E}_{i'k'} \ar[u]
}
$$
commutatif. Cela résulte de
Modules, lemme \ref{modules-lemma-finite-presentation-quasi-compact-colimit}.
On en déduit un morphisme
$$
\mathcal{A}_{ik} =
\text{Sym}^*_{\mathcal{O}_X}(\mathcal{F}_i)/(\mathcal{E}_{ik})
\longrightarrow
\text{Sym}^*_{\mathcal{O}_X}(\mathcal{F}_{i'})/(\mathcal{E}_{i'k'}) =
\mathcal{A}_{i'k'}
$$
où $(\mathcal{E}_{ik})$ désigne l'idéal engendré par $\mathcal{E}_{ik}$.
Les $\mathcal{O}_X$-algèbres quasi-cohérentes $\mathcal{A}_{ki}$
sont de présentation finie et sont finies pour $k$ assez grand (voir
la démonstration
du lemme \ref{properties-lemma-finite-algebra-directed-colimit-finite-finitely-presented}).
Enfin, on a
$$
\colim \mathcal{A}_{ik} = \colim \mathcal{A}_i = \mathcal{A}
$$
En effet, la première égalité a été montrée dans la démonstration du
lemme \ref{properties-lemma-finite-algebra-directed-colimit-finite-finitely-presented},
et la seconde résulte de ce que $\mathcal{A}$ est la colimite
des modules $\mathcal{F}_i$.
\end{proof}





\section{Le résultat de Gabber}
\label{properties-section-gabber}

\noindent
Dans cette section, nous démontrons un résultat de Gabber assurant que, sur
tout schéma, il existe un cardinal $\kappa$ tel que tout module
quasi-cohérent $\mathcal{F}$ soit la réunion de ses sous-faisceaux quasi-cohérents
$\kappa$-engendrés. Il en résulte que la catégorie des faisceaux quasi-cohérents sur un
schéma est une catégorie abélienne de Grothendieck, qui possède des
limites et suffisamment d'injectifs\footnote{Ce résultat est bien expliqué dans un
\href{https://amathew.wordpress.com/2011/07/30/quasi-coherent-sheaves-presentable-categories-and-a-result-of-gabber/}{billet de blog}
d'Akhil Mathew.}.

\begin{definition}
\label{properties-definition-kappa-generated}
Soit $(X, \mathcal{O}_X)$ un espace annelé. Soit $\kappa$ un cardinal
infini. On dit qu'un faisceau de $\mathcal{O}_X$-modules $\mathcal{F}$ est
{\it $\kappa$-engendré} s'il existe un recouvrement ouvert
$X = \bigcup U_i$ tel que $\mathcal{F}|_{U_i}$ soit engendré par un
sous-ensemble $R_i \subset \mathcal{F}(U_i)$ de cardinal au
plus $\kappa$.
\end{definition}

\noindent
Remarquons qu'une somme directe d'au plus $\kappa$ modules $\kappa$-engendrés est
encore $\kappa$-engendrée, car $\kappa \otimes \kappa = \kappa$, voir
Ensembles, section \ref{sets-section-cardinals}.
C'est en particulier le cas de la somme directe de deux modules $\kappa$-engendrés. De plus,
un quotient d'un faisceau $\kappa$-engendré est $\kappa$-engendré. (Il
n'en va pas nécessairement de même pour les sous-modules.)

\begin{lemma}
\label{properties-lemma-set-of-iso-classes}
Soit $(X, \mathcal{O}_X)$ un espace annelé. Soit $\kappa$ un cardinal. Il
existe un ensemble $T$ et une famille $(\mathcal{F}_t)_{t \in T}$ de modules
$\kappa$-engendrés sur $\mathcal{O}_X$ tels que tout module $\kappa$-engendré sur
$\mathcal{O}_X$ soit isomorphe à l'un des $\mathcal{F}_t$.
\end{lemma}

\begin{proof}
Les recouvrements de $X$ forment un ensemble si l'on interdit les répétitions.
Supposons que $X = \bigcup U_i$ soit un recouvrement et que $\mathcal{F}_i$
soit un $\mathcal{O}_{U_i}$-module. Les classes
d'isomorphisme de $\mathcal{O}_X$-modules $\mathcal{F}$ tels
que $\mathcal{F}|_{U_i} \cong \mathcal{F}_i$ forment alors un ensemble, car
les morphismes de recollement forment un ensemble. On se ramène donc à montrer que les classes
d'isomorphisme de quotients
$\oplus_{k \in \kappa} \mathcal{O}_X \to \mathcal{F}$
forment un ensemble pour tout espace annelé $X$. C'est clair.
\end{proof}

\noindent
Voici le résultat auquel fait référence le titre de cette section.

\begin{lemma}
\label{properties-lemma-colimit-kappa}
Soit $X$ un schéma. Il existe un cardinal $\kappa$ tel que
tout module quasi-cohérent $\mathcal{F}$ soit la colimite filtrante de
ses sous-modules quasi-cohérents $\kappa$-engendrés.
\end{lemma}

\begin{proof}
Choisissons un recouvrement ouvert affine $X = \bigcup_{i \in I} U_i$. Pour tout couple
$i, j$, choisissons un recouvrement ouvert affine
$U_i \cap U_j = \bigcup_{k \in I_{ij}} U_{ijk}$.
Écrivons $U_i = \Spec(A_i)$ et $U_{ijk} = \Spec(A_{ijk})$.
Soit $\kappa$ un cardinal infini $\geq$ au cardinal de
chacun des ensembles $I$, $I_{ij}$.

\medskip\noindent
Soit $\mathcal{F}$ un faisceau quasi-cohérent. Posons $M_i = \mathcal{F}(U_i)$
et $M_{ijk} = \mathcal{F}(U_{ijk})$. Remarquons que
$$
M_i \otimes_{A_i} A_{ijk} = M_{ijk} = M_j \otimes_{A_j} A_{ijk}.
$$
Voir
Schémas, lemme \ref{schemes-lemma-widetilde-pullback}.
À l'aide de l'axiome du choix, choisissons une application
$$
(i, j, k, m) \mapsto S(i, j, k, m)
$$
qui associe à tous $i, j \in I$, $k \in I_{ij}$ et $m \in M_i$
un sous-ensemble fini $S(i, j, k, m) \subset M_j$ tel que l'on ait
$$
m \otimes 1 = \sum\nolimits_{m' \in S(i, j, k, m)} m' \otimes a_{m'}
$$
dans $M_{ijk}$ pour certains $a_{m'} \in A_{ijk}$. Convenons de plus
que $S(i, i, k, m) = \{m\}$ pour tous $i, j = i, k, m$ comme ci-dessus.
Fixons une telle application.

\medskip\noindent
Étant donnée une famille $\mathcal{S} = (S_i)_{i \in I}$ de sous-ensembles
$S_i \subset M_i$ de cardinal au plus $\kappa$, posons
$\mathcal{S}' = (S'_i)$, où
$$
S'_j = \bigcup\nolimits_{(i, k, m)\text{ tels que }m \in S_i}
S(i, j, k, m)
$$
On a $S_i \subset S'_i$. De plus, $S'_i$ a un cardinal au plus
$\kappa$, car c'est une réunion, indexée par un ensemble de cardinal au plus $\kappa$,
d'ensembles finis. Posons $\mathcal{S}^{(0)} = \mathcal{S}$,
$\mathcal{S}^{(1)} = \mathcal{S}'$ et, par récurrence,
$\mathcal{S}^{(n + 1)} = (\mathcal{S}^{(n)})'$. Posons ensuite
$\mathcal{S}^{(\infty)} = \bigcup_{n \geq 0} \mathcal{S}^{(n)}$.
En écrivant $\mathcal{S}^{(\infty)} = (S^{(\infty)}_i)$, on voit que,
pour tout élément $m \in S^{(\infty)}_i$, l'image de $m$ dans
$M_{ijk}$ peut s'écrire comme une somme finie $\sum m' \otimes a_{m'}$
avec $m' \in S_j^{(\infty)}$. Ainsi, en posant
$$
N_i = A_i\text{-sous-module de }M_i\text{ engendré par }S^{(\infty)}_i
$$
on a
$$
N_i \otimes_{A_i} A_{ijk} = N_j \otimes_{A_j} A_{ijk}.
$$
comme sous-modules de $M_{ijk}$. Il existe donc un sous-faisceau quasi-cohérent
$\mathcal{G} \subset \mathcal{F}$ tel que $\mathcal{G}(U_i) = N_i$.
De plus, par construction, le faisceau $\mathcal{G}$ est $\kappa$-engendré.

\medskip\noindent
Soit $\{\mathcal{G}_t\}_{t \in T}$ l'ensemble des sous-faisceaux quasi-cohérents $\kappa$-engendrés.
Si $t, t' \in T$, alors
$\mathcal{G}_t + \mathcal{G}_{t'}$ est encore un sous-faisceau quasi-cohérent $\kappa$-engendré,
étant l'image du morphisme
$\mathcal{G}_t \oplus \mathcal{G}_{t'} \to \mathcal{F}$.
Le système ordonné par inclusion est donc filtrant. Les
arguments précédents montrent que toute section de $\mathcal{F}$ sur $U_i$
appartient à l'un des $\mathcal{G}_t$ (car on peut partir d'une famille $\mathcal{S}$
telle que la section donnée soit un élément de $S_i$). Ainsi
$\colim_t \mathcal{G}_t \to \mathcal{F}$ est à la fois injectif et surjectif,
comme voulu.
\end{proof}

\begin{proposition}
\label{properties-proposition-coherator}
Soit $X$ un schéma.
\begin{enumerate}
\item La catégorie $\QCoh(\mathcal{O}_X)$ est une catégorie
abélienne de Grothendieck. Par conséquent, $\QCoh(\mathcal{O}_X)$
possède suffisamment d'injectifs et toutes les limites.
\item Le foncteur d'inclusion
$\QCoh(\mathcal{O}_X) \to \textit{Mod}(\mathcal{O}_X)$
possède un adjoint à droite\footnote{Ce foncteur est parfois appelé
le {\it cohérateur}.}
$$
Q : \textit{Mod}(\mathcal{O}_X) \longrightarrow \QCoh(\mathcal{O}_X)
$$
tel que, pour tout faisceau quasi-cohérent $\mathcal{F}$, le morphisme d'adjonction
$Q(\mathcal{F}) \to \mathcal{F}$ soit un isomorphisme.
\end{enumerate}
\end{proposition}

\begin{proof}
L'assertion (1) signifie que $\QCoh(\mathcal{O}_X)$ (a) possède toutes les
colimites, (b) a des colimites filtrantes exactes et (c) possède un générateur,
voir Injectifs, section \ref{injectives-section-grothendieck-conditions}.
D'après Schémas, section \ref{schemes-section-quasi-coherent},
les colimites dans $\QCoh(\mathcal{O}_X)$ existent et coïncident
avec les colimites dans $\textit{Mod}(\mathcal{O}_X)$.
D'après Modules, lemme \ref{modules-lemma-limits-colimits},
les colimites filtrantes sont exactes. Ainsi (a) et (b) sont
satisfaites. Pour construire un générateur $U$, choisissons un cardinal $\kappa$ comme dans
le lemme \ref{properties-lemma-colimit-kappa}. Choisissons une famille
$(\mathcal{F}_t)_{t \in T}$ de faisceaux quasi-cohérents $\kappa$-engendrés comme dans le
lemme \ref{properties-lemma-set-of-iso-classes}. Posons
$U = \bigoplus_{t \in T} \mathcal{F}_t$. Comme tout objet de
$\QCoh(\mathcal{O}_X)$ est une colimite filtrante de modules quasi-cohérents $\kappa$-engendrés,
c'est-à-dire d'objets isomorphes aux $\mathcal{F}_t$,
il est clair que $U$ est un générateur.
Les assertions relatives aux limites et aux injectifs sont vraies
dans toute catégorie abélienne de Grothendieck,
voir Injectifs, théorème
\ref{injectives-theorem-injective-embedding-grothendieck} et
lemme \ref{injectives-lemma-grothendieck-products}.

\medskip\noindent
Démonstration de (2). Pour construire $Q$, on utilise le procédé général suivant.
Étant donné un objet $\mathcal{F}$ de $\textit{Mod}(\mathcal{O}_X)$,
considérons le foncteur
$$
\QCoh(\mathcal{O}_X)^{opp} \longrightarrow \textit{Ensembles},\quad
\mathcal{G} \longmapsto \Hom_X(\mathcal{G}, \mathcal{F})
$$
Ce foncteur transforme les colimites en limites ;
il est donc représentable,
voir Injectifs, lemme \ref{injectives-lemma-grothendieck-brown}.
Il existe ainsi un faisceau quasi-cohérent $Q(\mathcal{F})$
et un isomorphisme fonctoriel
$\Hom_X(\mathcal{G}, \mathcal{F}) = \Hom_X(\mathcal{G}, Q(\mathcal{F}))$
pour $\mathcal{G}$ dans $\QCoh(\mathcal{O}_X)$. Par le lemme de
Yoneda (Catégories, lemme \ref{categories-lemma-yoneda}),
la construction $\mathcal{F} \leadsto Q(\mathcal{F})$ est
fonctorielle en $\mathcal{F}$. Par construction, $Q$ est
adjoint à droite du foncteur
d'inclusion. Le fait que $Q(\mathcal{F}) \to \mathcal{F}$ soit un isomorphisme
lorsque $\mathcal{F}$ est quasi-cohérent découle formellement du fait
que le foncteur d'inclusion
$\QCoh(\mathcal{O}_X) \to \textit{Mod}(\mathcal{O}_X)$
est pleinement fidèle.
\end{proof}







\section{sections à support dans un fermé}
\label{properties-section-sections-with-support-in-closed}

\noindent
Étant donnés un espace topologique $X$, un fermé $Z \subset X$ et un
faisceau abélien $\mathcal{F}$, on peut considérer le sous-faisceau des sections
dont le support est contenu dans $Z$. Si $X$ est un schéma, $Z$ un
sous-schéma fermé et $\mathcal{F}$ un module quasi-cohérent, il existe une variante
où l'on considère les sections dont le support est schématiquement contenu
dans $Z$. Il faut toutefois être prudent dans le cadre des
schémas, car le $\mathcal{O}_X$-module obtenu n'est pas nécessairement quasi-cohérent.

\begin{lemma}
\label{properties-lemma-quasi-coherent-finite-type-ideals}
Soit $X$ un schéma quasi-compact et quasi-séparé.
Soit $U \subset X$ un sous-schéma ouvert. Les assertions suivantes sont équivalentes :
\begin{enumerate}
\item $U$ est rétrocompact dans $X$;
\item $U$ est quasi-compact ;
\item $U$ est une réunion finie d'ouverts affines ;
\item il existe un faisceau d'idéaux quasi-cohérent de type fini
$\mathcal{I} \subset \mathcal{O}_X$ tel que $X \setminus U = V(\mathcal{I})$
ensemblistement.
\end{enumerate}
\end{lemma}

\begin{proof}
L'équivalence de (1), (2) et (3) résulte du
lemme \ref{properties-lemma-quasi-separated-quasi-compact-open-retrocompact}.
Supposons (1), (2), (3). Soit $T = X \setminus U$.
D'après Schémas, lemme \ref{schemes-lemma-reduced-closed-subscheme}, il existe
un unique faisceau d'idéaux quasi-cohérent $\mathcal{J}$ définissant
la structure de sous-schéma fermé réduit induite sur $T$.
Remarquons que $\mathcal{J}|_U = \mathcal{O}_U$, qui est un
$\mathcal{O}_U$-module de type fini. Par
le lemme \ref{properties-lemma-extend}, il existe un sous-faisceau quasi-cohérent
$\mathcal{I} \subset \mathcal{J}$ de type
fini tel que $\mathcal{I}|_U = \mathcal{J}|_U$.
Alors $X \setminus U = V(\mathcal{I})$, ce qui donne (4). Réciproquement,
si $\mathcal{I}$ est comme dans (4) et si $W = \Spec(R) \subset X$ est un ouvert
affine, alors $\mathcal{I}|_W = \widetilde{I}$ pour un idéal de type
fini $I \subset R$, voir le lemme \ref{properties-lemma-finite-type-module}.
Il en résulte que $U \cap W = \Spec(R) \setminus V(I)$ est quasi-compact,
voir Algèbre, lemme \ref{algebra-lemma-qc-open}. Ainsi $U \subset X$
est rétrocompact par le lemme \ref{properties-lemma-retrocompact}.
\end{proof}

\begin{lemma}
\label{properties-lemma-sections-annihilated-by-ideal}
Soit $X$ un schéma.
Soit $\mathcal{I} \subset \mathcal{O}_X$ un faisceau d'idéaux quasi-cohérent.
Soit $\mathcal{F}$ un $\mathcal{O}_X$-module
quasi-cohérent. Considérons le faisceau de $\mathcal{O}_X$-modules $\mathcal{F}'$
qui associe à tout ouvert $U \subset X$ l'ensemble
$$
\mathcal{F}'(U)
=
\{s \in \mathcal{F}(U) \mid
\mathcal{I}s = 0\}
$$
Supposons $\mathcal{I}$ de type fini. Alors
\begin{enumerate}
\item $\mathcal{F}'$ est un faisceau quasi-cohérent de $\mathcal{O}_X$-modules ;
\item sur tout ouvert affine $U \subset X$, on a
$\mathcal{F}'(U) = \{s \in \mathcal{F}(U) \mid \mathcal{I}(U)s = 0\}$;
\item $\mathcal{F}'_x = \{s \in \mathcal{F}_x \mid \mathcal{I}_x s = 0\}$.
\end{enumerate}
\end{lemma}

\begin{proof}
Il est clair que la règle définissant $\mathcal{F}'$ donne un sous-faisceau
de $\mathcal{F}$ (la condition de faisceau se vérifie aisément). On peut donc
travailler localement sur $X$ pour vérifier les autres assertions. Autrement
dit, on peut supposer $X = \Spec(A)$, $\mathcal{F} = \widetilde{M}$
et $\mathcal{I} = \widetilde{I}$. Il est clair que, dans ce cas,
$\mathcal{F}'(U) = \{x \in M \mid Ix = 0\} =: M'$, car $\widetilde{I}$
est engendré par ses sections globales $I$ ; cela démontre (2). Pour
montrer que $\mathcal{F}'$ est quasi-cohérent, il suffit de montrer que,
pour tout $f \in A$, on a
$\{x \in M_f \mid I_f x = 0\} = (M')_f$.
Écrivons $I = (g_1, \ldots, g_t)$, ce qui est possible puisque $\mathcal{I}$
est de type fini, voir le lemme \ref{properties-lemma-finite-type-module}.
Si $x = y/f^n$ et $I_fx = 0$, cela signifie que, pour tout $i$,
il existe un $m \geq 0$ tel que $f^mg_ix = 0$.
On peut choisir un même $m$ convenant à tous les $i$ ; c'est ici que l'on utilise
le fait que $I$ est de type fini. Alors $f^mx \in M'$
et $x/f^n = f^mx/f^{n + m}$ dans $(M')_f$, comme
voulu. La démonstration de (3) est analogue et omise.
\end{proof}

\begin{definition}
\label{properties-definition-subsheaf-sections-annihilated-by-ideal}
Soit $X$ un schéma.
Soit $\mathcal{I} \subset \mathcal{O}_X$ un faisceau d'idéaux quasi-cohérent
de type fini.
Soit $\mathcal{F}$ un $\mathcal{O}_X$-module
quasi-cohérent. Le sous-faisceau $\mathcal{F}' \subset \mathcal{F}$ défini dans
le lemme \ref{properties-lemma-sections-annihilated-by-ideal} ci-dessus est
appelé le {\it sous-faisceau des sections annulées par $\mathcal{I}$}.
\end{definition}

\begin{lemma}
\label{properties-lemma-push-sections-annihilated-by-ideal}
Soit $f : X \to Y$ un morphisme quasi-compact et quasi-séparé de
schémas. Soit $\mathcal{I} \subset \mathcal{O}_Y$ un faisceau
d'idéaux quasi-cohérent de type fini. Soit $\mathcal{F}$ un
$\mathcal{O}_X$-module quasi-cohérent. Soit $\mathcal{F}' \subset \mathcal{F}$
le sous-faisceau des sections annulées par $f^{-1}\mathcal{I}\mathcal{O}_X$.
Alors $f_*\mathcal{F}' \subset f_*\mathcal{F}$ est le
sous-faisceau des sections annulées par $\mathcal{I}$.
\end{lemma}

\begin{proof}
Démonstration omise. (Indication : l'hypothèse que $f$ est quasi-compact
et quasi-séparé entraîne que $f_*\mathcal{F}$ est quasi-cohérent, si
bien que le lemme \ref{properties-lemma-sections-annihilated-by-ideal} s'applique
à $\mathcal{I}$ et à $f_*\mathcal{F}$.)
\end{proof}

\noindent
Pour un faisceau abélien sur un espace topologique, nous avons étudié le
sous-faisceau des sections à support dans un fermé dans
Modules, remarque \ref{modules-remark-sections-support-in-closed}.
Pour les modules quasi-cohérents, ce sous-module n'est pas toujours
quasi-cohérent ; il l'est toutefois si le complémentaire du fermé est
rétrocompact.

\begin{lemma}
\label{properties-lemma-sections-supported-on-closed-subset}
Soient $X$ un schéma et $Z \subset X$ un fermé.
Soit $\mathcal{F}$ un $\mathcal{O}_X$-module
quasi-cohérent. Considérons le faisceau de $\mathcal{O}_X$-modules $\mathcal{F}'$
qui associe à tout ouvert $U \subset X$ l'ensemble
$$
\mathcal{F}'(U)
=
\{s \in \mathcal{F}(U) \mid
\text{le support de }s\text{ est contenu dans }Z \cap U\}
$$
Si $X \setminus Z$ est un ouvert rétrocompact de $X$, alors
\begin{enumerate}
\item pour tout ouvert affine $U \subset X$, il existe un idéal de type
fini $I \subset \mathcal{O}_X(U)$ tel que $Z \cap U = V(I)$;
\item pour $U$ et $I$ comme dans (1), on a
$\mathcal{F}'(U) = \{x \in \mathcal{F}(U) \mid
I^nx = 0 \text{ pour un certain } n\}$;
\item $\mathcal{F}'$ est un faisceau quasi-cohérent de $\mathcal{O}_X$-modules.
\end{enumerate}
\end{lemma}

\begin{proof}
L'assertion (1) est Algèbre, lemme \ref{algebra-lemma-qc-open}.
Soient $U = \Spec(A)$ et $I$ comme dans (1).
Alors $\mathcal{F}|_U$ est le faisceau quasi-cohérent associé
à un $A$-module $M$. On a
$$
\mathcal{F}'(U) = \{x \in M \mid x = 0\text{ dans }M_\mathfrak p
\text{ pour tout }\mathfrak p \not \in Z\}.
$$
d'après Modules, définition \ref{modules-definition-support}. Ainsi
$x \in \mathcal{F}'(U)$ si et seulement si $V(\text{Ann}(x)) \subset V(I)$, voir
Algèbre, lemme \ref{algebra-lemma-support-element}. Comme $I$ est
de type fini, cela équivaut à $I^n x = 0$ pour un certain $n$.
Cela démontre (2).

\medskip\noindent
Démonstration de (3). Remarquons que, pour tout ouvert $U \subset X$, on a une suite exacte
$$
0 \to \mathcal{F}'(U) \to \mathcal{F}(U) \to \mathcal{F}(U \setminus Z)
$$
Si l'on note $j : X \setminus Z \to X$ le morphisme
d'inclusion, $\mathcal{F}(U \setminus Z)$ est l'ensemble des sections
du module $j_*(\mathcal{F}|_{X \setminus Z})$ sur $U$. On a donc
une suite exacte
$$
0 \to \mathcal{F}' \to \mathcal{F} \to j_*(\mathcal{F}|_{X \setminus Z})
$$
La restriction $\mathcal{F}|_{X \setminus Z}$ est quasi-cohérente.
Donc $j_*(\mathcal{F}|_{X \setminus Z})$ est quasi-cohérent,
d'après Schémas, lemme \ref{schemes-lemma-push-forward-quasi-coherent},
et l'hypothèse que $j$ est quasi-compact (toute immersion
ouverte est séparée). Ainsi $\mathcal{F}'$ est quasi-cohérent,
étant le noyau d'un morphisme de modules quasi-cohérents,
voir Schémas, section \ref{schemes-section-quasi-coherent}.
\end{proof}

\begin{definition}
\label{properties-definition-subsheaf-sections-supported-on-closed}
Soit $X$ un schéma.
Soit $T \subset X$ un fermé dont le complémentaire est
rétrocompact dans $X$.
Soit $\mathcal{F}$ un $\mathcal{O}_X$-module
quasi-cohérent. Le sous-faisceau quasi-cohérent $\mathcal{F}' \subset \mathcal{F}$ défini dans
le lemme \ref{properties-lemma-sections-supported-on-closed-subset} est appelé
le {\it sous-faisceau des sections à support dans $T$}.
\end{definition}

\begin{lemma}
\label{properties-lemma-push-sections-supported-on-closed-subset}
Soit $f : X \to Y$ un morphisme quasi-compact et quasi-séparé de
schémas. Soit $Z \subset Y$ un fermé tel que
$Y \setminus Z$ soit rétrocompact dans $Y$. Soit $\mathcal{F}$ un
$\mathcal{O}_X$-module quasi-cohérent. Soit $\mathcal{F}' \subset \mathcal{F}$
le sous-faisceau des sections à support dans $f^{-1}Z$.
Alors $f_*\mathcal{F}' \subset f_*\mathcal{F}$ est le
sous-faisceau des sections à support dans $Z$.
\end{lemma}

\begin{proof}
Démonstration omise. (Indication : montrer d'abord que $X \setminus f^{-1}Z$ est rétrocompact
dans $X$, puisque $Y \setminus Z$ est rétrocompact dans $Y$. Ainsi,
le lemme \ref{properties-lemma-sections-supported-on-closed-subset}
s'applique à $f^{-1}Z$ et à $\mathcal{F}$. Comme $f$ est
quasi-compact et quasi-séparé, $f_*\mathcal{F}$ est
quasi-cohérent. Le lemme \ref{properties-lemma-sections-supported-on-closed-subset}
s'applique donc à $Z$ et à $f_*\mathcal{F}$. Enfin, identifier directement les
faisceaux.)
\end{proof}








\section{sections des faisceaux quasi-cohérents}
\label{properties-section-sections-quasi-coherent}

\noindent
Voici un calcul des sections d'un faisceau quasi-cohérent sur un ouvert
quasi-compact d'un spectre affine.

\begin{lemma}
\label{properties-lemma-sections-over-quasi-compact-open-in-affine}
Soit $A$ un anneau.
Soit $I \subset A$ un idéal de type fini.
Soit $M$ un $A$-module. Il
existe alors un homomorphisme canonique
$$
\colim_n \Hom_A(I^n, M)
\longrightarrow
\Gamma(\Spec(A) \setminus V(I), \widetilde{M}).
$$
Cet homomorphisme est toujours injectif.
Si, pour tout $x \in M$, on a $Ix = 0 \Rightarrow x = 0$,
il est un isomorphisme. Dans le cas général, posons
$M_n = \{x \in M \mid I^nx = 0\}$; on a alors un
isomorphisme
$$
\colim_n \Hom_A(I^n, M/M_n)
\longrightarrow
\Gamma(\Spec(A) \setminus V(I), \widetilde{M}).
$$
\end{lemma}

\begin{proof}
Comme $I^{n + 1} \subset I^n$ et $M_n \subset M_{n + 1}$, on peut
composer avec ces homomorphismes pour obtenir des homomorphismes canoniques de $A$-modules
$$
\Hom_A(I^n, M)
\longrightarrow
\Hom_A(I^{n + 1}, M)
$$
et
$$
\Hom_A(I^n, M/M_n)
\longrightarrow
\Hom_A(I^{n + 1}, M/M_{n + 1})
$$
que l'on utilisera comme morphismes de transition des systèmes. Étant donné un morphisme de
$A$-modules $\varphi : I^n \to M$, on obtient un morphisme de
faisceaux $\widetilde{\varphi} : \widetilde{I^n} \to \widetilde{M}$
que l'on peut restreindre à l'ouvert $\Spec(A) \setminus V(I)$.
Comme la restriction de $\widetilde{I^n}$ à cet ouvert est le faisceau structural,
on obtient un élément de
$\Gamma(\Spec(A) \setminus V(I), \widetilde{M})$.
Nous omettons de vérifier la compatibilité avec les morphismes de
transition du système $\Hom_A(I^n, M)$. Cela donne la première flèche.
Pour obtenir la seconde, remarquons que
$\widetilde{M}$ et $\widetilde{M/M_n}$ coïncident sur l'ouvert
$\Spec(A) \setminus V(I)$, puisque le faisceau $\widetilde{M_n}$
a évidemment son support dans $V(I)$. On peut donc procéder de la même
manière.

\medskip\noindent
Décrivons maintenant cette flèche en termes algébriques.
Écrivons $I = (f_1, \ldots, f_t)$. Alors
$\Spec(A) \setminus V(I) = \bigcup_{i = 1, \ldots, t} D(f_i)$.
Ainsi la suite
$$
0 \to
\Gamma(\Spec(A) \setminus V(I), \widetilde{M}) \to
\bigoplus\nolimits_i M_{f_i} \to
\bigoplus\nolimits_{i, j} M_{f_if_j}
$$
est exacte. Soit $\varphi : I^n \to M$ un morphisme de $A$-modules.
Considérons le vecteur dont les composantes sont $\varphi(f_i^n)/f_i^n \in M_{f_i}$.
On voit aisément que son image est nulle dans la seconde
somme directe de la suite exacte ci-dessus. Il définit donc un élément
de $\Gamma(\Spec(A) \setminus V(I), \widetilde{M})$.
Nous omettons de vérifier que cette description coïncide avec celle donnée
précédemment.

\medskip\noindent
Montrons à l'aide de cette description que la première flèche est
injective. Si $\varphi$ a une image nulle, alors, pour tout $i$, l'élément
$\varphi(f_i^n)/f_i^n$ est nul dans $M_{f_i}$. Autrement
dit, pour tout $i$, on a $f_i^m\varphi(f_i^n) = 0$ pour un certain $m \geq 0$.
On peut choisir un même $m$ convenant à tous les $i$. Alors
$\varphi(f_i^{n + m}) = 0$ pour tout $i$. On voit aisément que
cela entraîne $\varphi|_{I^{t(n + m - 1) + 1}} = 0$ ;
autrement dit, $\varphi$ a une image nulle dans le terme d'indice $t(n + m - 1) + 1$ de
la colimite. L'injectivité en découle.

\medskip\noindent
Remarquons que tous les $M_n = 0$ si l'on a
$Ix = 0 \Rightarrow x = 0$ pour $x \in M$. Pour
achever la démonstration du lemme, il suffit donc de montrer que
la seconde flèche est un isomorphisme.

\medskip\noindent
Cherchons à construire un inverse du second homomorphisme du lemme.
Soit $s \in \Gamma(\Spec(A) \setminus V(I), \widetilde{M})$.
Cette section correspond à un vecteur $x_i/f_i^n$, avec $x_i \in M$, de la
première somme directe de la suite exacte ci-dessus.
Ainsi, pour tous $i, j$, il existe un $m \geq 0$
tel que $f_i^m f_j^m (f_j^n x_i - f_i^n x_j) = 0$ dans $M$.
On peut choisir un même $m$ convenant à tous les couples $i, j$.
En remplaçant $x_i$ par $f_i^mx_i$ et $n$ par $n + m$, on
obtient $f_j^nx_i = f_i^nx_j$ dans $M$ pour tous $i, j$.
Introduisons
$$
K_n = \{x \in M \mid f_1^nx = \ldots = f_t^nx = 0\}
$$
Nous affirmons qu'il existe un morphisme de $A$-modules
$$
\varphi :
I^{t(n - 1) + 1}
\longrightarrow
M/K_n
$$
qui envoie le monôme
$f_1^{e_1} \ldots f_t^{e_t}$, avec $\sum e_i = t(n - 1) + 1$,
sur la classe modulo $K_n$ de l'expression
$f_1^{e_1} \ldots f_i^{e_i - n} \ldots f_t^{e_t}x_i$,
où $i$ est choisi de sorte que $e_i \geq n$ (il existe
au moins un tel $i$).
Pour le vérifier, supposons que
$$
\sum\nolimits_{E = (e_1, \ldots, e_t), |E| = t(n - 1) + 1}
a_E f_1^{e_1} \ldots f_t^{e_t} = 0
$$
soit une relation entre les monômes à coefficients $a_E$ dans $A$.
Cette relation serait envoyée sur
$$
z =
\sum\nolimits_{E = (e_1, \ldots, e_t), |E| = t(n - 1) + 1}
a_E f_1^{e_1} \ldots f_{i(E)}^{e_{i(E)} - n} \ldots f_t^{e_t}x_{i(E)}
$$
où, pour chaque multi-indice $E$, on a choisi un $i(E)$
tel que $e_{i(E)} \geq n$.
Remarquons qu'en multipliant cette expression par $f_j^n$, pour un $j$ quelconque,
on obtient zéro : en effet, les relations $f_j^nx_i = f_i^nx_j$ ci-dessus donnent
\begin{align*}
f_j^nz & = \sum\nolimits_{E = (e_1, \ldots, e_t), |E| = t(n - 1) + 1}
a_E f_1^{e_1} \ldots f_j^{e_j + n}
\ldots f_{i(E)}^{e_{i(E)} - n} \ldots f_t^{e_t}x_{i(E)} \\
& =
\sum\nolimits_{E = (e_1, \ldots, e_t), |E| = t(n - 1) + 1}
a_E f_1^{e_1} \ldots f_t^{e_t}x_j
= 0.
\end{align*}
Ainsi $z \in K_n$, et toute relation est envoyée sur zéro
dans $M/K_n$. L'assertion est démontrée.

\medskip\noindent
Remarquons que $K_n \subset M_{t(n - 1) + 1}$. Le
morphisme $\varphi$ fournit donc en particulier un morphisme de $A$-modules
$I^{t(n - 1) + 1} \to M/M_{t(n - 1) + 1}$.
Cela montre que la seconde flèche du lemme est surjective. Nous
omettons la démonstration de l'injectivité.
\end{proof}

\begin{example}
\label{properties-example-does-not-work-in-general}
Donnons deux exemples où le premier homomorphisme affiché du
lemme \ref{properties-lemma-sections-over-quasi-compact-open-in-affine}
n'est pas un isomorphisme.

\medskip\noindent
Soit $k$ un corps. Considérons l'anneau
$$
A = k[x, y, z_1, z_2, \ldots]/(x^nz_n).
$$
Posons $I = (x)$ et soit $M = A$. L'élément $y/x$ définit alors une
section du faisceau structural de $\Spec(A)$ sur $D(x) = \Spec(A) \setminus V(I)$.
Nous affirmons que $y/x$ n'appartient pas à l'image de l'homomorphisme canonique
$\colim \Hom_A(I^n, A) \to A_x = \mathcal{O}(D(x))$. Sinon, il
proviendrait d'un homomorphisme $\varphi : I^n \to A$ pour un certain $n$.
Posons $a = \varphi(x^n)$. On aurait alors $x^m(xa - x^ny) = 0$ pour un certain
$m > 0$, donc $x^{m + 1}a = x^{m + n}y$, puis
$\varphi(x^{n + m + 1}) = x^{m + n}y$. Cela conduit à une
contradiction, car on aurait
$$
0 = \varphi(0) = \varphi(z_{n + m + 1} x^{n + m + 1}) =
x^{m + n}y z_{n + m + 1}
$$
ce qui est faux dans l'anneau $A$.

\medskip\noindent
Soit $k$ un corps. Considérons l'anneau
$$
A = k[f, g, x, y, \{a_n, b_n\}_{n \geq 1}]/
(fy - gx, \{a_nf^n + b_ng^n\}_{n \geq 1}).
$$
Posons $I = (f, g)$ et soit $M = A$. Alors $x/f \in A_f$ et $y/g \in A_g$
ont la même image dans $A_{fg}$. Ils définissent donc une section $s$
du faisceau structural de $\Spec(A)$ sur
$D(f) \cup D(g) = \Spec(A) \setminus V(I)$. Toutefois, il n'existe aucun
$n \geq 0$ tel que $s$ provienne d'un morphisme de $A$-modules
$\varphi : I^n \to A$ appartenant à la source de la première flèche
affichée du lemme \ref{properties-lemma-sections-over-quasi-compact-open-in-affine}.
En effet, étant donné un tel morphisme de modules, posons
$x_n = \varphi(f^n)$ et $y_n = \varphi(g^n)$.
Alors $f^mx_n = f^{n + m - 1}x$ et
$g^my_n = g^{n + m - 1}y$ pour un certain $m \geq 0$ (voir la
démonstration du lemme). Mais on aurait alors
$0 = \varphi(0) =
\varphi(a_{n + m}f^{n + m} + b_{n + m}g^{n + m}) =
a_{n + m}f^{n + m - 1}x + b_{n + m}g^{n + m - 1}y$, ce qui est faux
dans l'anneau $A$.
\end{example}

\noindent
Nous préciserons le lemme suivant dans le cas noethérien, voir
Cohomologie des schémas, lemme \ref{coherent-lemma-homs-over-open}.

\begin{lemma}
\label{properties-lemma-sections-over-quasi-compact-open}
Soit $X$ un schéma quasi-compact.
Soit $\mathcal{I} \subset \mathcal{O}_X$ un
faisceau d'idéaux quasi-cohérent de type fini.
Soit $Z \subset X$ le sous-schéma fermé
défini par $\mathcal{I}$, et posons $U = X \setminus Z$.
Soit $\mathcal{F}$ un $\mathcal{O}_X$-module
quasi-cohérent. L'homomorphisme canonique
$$
\colim_n \Hom_{\mathcal{O}_X}(\mathcal{I}^n,
\mathcal{F})
\longrightarrow
\Gamma(U, \mathcal{F})
$$
est injectif. Supposons de plus $X$ quasi-séparé.
Soit $\mathcal{F}_n \subset \mathcal{F}$
le sous-faisceau des sections annulées par $\mathcal{I}^n$.
L'homomorphisme canonique
$$
\colim_n \Hom_{\mathcal{O}_X}(\mathcal{I}^n,
\mathcal{F}/\mathcal{F}_n)
\longrightarrow
\Gamma(U, \mathcal{F})
$$
est un isomorphisme.
\end{lemma}

\begin{proof}
Soit $\Spec(A) = W \subset X$ un ouvert affine. Écrivons
$\mathcal{F}|_W = \widetilde{M}$ pour un $A$-module $M$
et $\mathcal{I}|_W = \widetilde{I}$ pour un idéal de type
fini $I \subset A$. En restreignant le premier homomorphisme
affiché du lemme à $W$, on obtient le premier homomorphisme affiché
du lemme \ref{properties-lemma-sections-over-quasi-compact-open-in-affine}.
Comme on peut recouvrir $X$ par un nombre fini d'ouverts affines, cela
montre que le premier homomorphisme affiché du lemme est injectif.

\medskip\noindent
On a $\mathcal{F}_n|_W = \widetilde{M_n}$, où
$M_n \subset M$ est défini comme dans
le lemme \ref{properties-lemma-sections-over-quasi-compact-open-in-affine}
(détails omis). Ce lemme fournit une bijection
$$
\colim_n  \Hom_{\mathcal{O}_W}(
\mathcal{I}^n|_W, (\mathcal{F}/\mathcal{F}_n)|_W)
\longrightarrow
\Gamma(U \cap W, \mathcal{F})
$$
pour tout ouvert affine $W$ comme ci-dessus.

\medskip\noindent
Pour montrer que la seconde flèche affichée du lemme est bijective,
choisissons un recouvrement ouvert affine fini $X = \bigcup_{j = 1, \ldots, m} W_j$.
L'injectivité découle immédiatement de ce qui précède et de la finitude du
recouvrement. Si $X$ est quasi-séparé, pour tout couple
$j, j'$, choisissons un recouvrement ouvert affine fini
$$
W_j \cap W_{j'} = \bigcup\nolimits_{k = 1, \ldots, m_{jj'}} W_{jj'k}.
$$
Soit $s \in \Gamma(U, \mathcal{F})$. Comme on l'a vu, pour tout $j$, il existe
un $n_j$ et un morphisme
$\varphi_j : \mathcal{I}^{n_j}|_{W_j} \to
(\mathcal{F}/\mathcal{F}_{n_j})|_{W_j}$
correspondant à $s|_{U \cap W_j}$.
De même, pour tout triplet $(j, j', k)$, il existe un entier
$n_{jj'k}$ tel que les restrictions de $\varphi_j$ et de $\varphi_{j'}$,
comme morphismes $\mathcal{I}^{n_{jj'k}} \to \mathcal{F}/\mathcal{F}_{n_{jj'k}}$,
coïncident sur $W_{jj'k}$. Posons $n = \max\{n_j, n_{jj'k}\}$ ; alors
les $\varphi_j$ se recollent en un morphisme
$\mathcal{I}^n \to \mathcal{F}/\mathcal{F}_n$ sur $X$.
Cela démontre la surjectivité de l'homomorphisme.
\end{proof}













\section{Faisceaux inversibles amples}
\label{properties-section-ample}

\noindent
Rappelons d'après Modules, lemme \ref{modules-lemma-s-open},
qu'étant donnés un faisceau inversible $\mathcal{L}$ sur un espace localement
annelé $X$ et une section globale $s$ de $\mathcal{L}$,
l'ensemble $X_s = \{x \in X \mid s \not \in \mathfrak m_x\mathcal{L}_x\}$
est ouvert. On remarque en général que
$X_s \cap X_{s'} = X_{ss'}$, où $ss'$ désigne
la section $s \otimes s' \in \Gamma(X, \mathcal{L} \otimes \mathcal{L}')$.

\begin{definition}
\label{properties-definition-ample}
\begin{reference}
\cite[II Definition 4.5.3]{EGA}
\end{reference}
Soit $X$ un schéma.
Soit $\mathcal{L}$ un $\mathcal{O}_X$-module
inversible. On dit que $\mathcal{L}$ est {\it ample} si
\begin{enumerate}
\item $X$ est quasi-compact, et
\item pour tout $x \in X$, il existe un entier $n \geq 1$
et $s \in \Gamma(X, \mathcal{L}^{\otimes n})$ tels
que $x \in X_s$ et que $X_s$ soit affine.
\end{enumerate}
\end{definition}

\begin{lemma}
\label{properties-lemma-ample-power-ample}
\begin{reference}
\cite[II Proposition 4.5.6(i)]{EGA}
\end{reference}
Soit $X$ un schéma. Soit $\mathcal{L}$ un $\mathcal{O}_X$-module
inversible. Soit $n \geq 1$. Alors $\mathcal{L}$ est ample si et seulement si
$\mathcal{L}^{\otimes n}$ est ample.
\end{lemma}

\begin{proof}
Cela résulte de l'égalité $X_{s^n} = X_s$.
\end{proof}

\begin{lemma}
\label{properties-lemma-ample-on-closed}
Soit $X$ un schéma.
Soit $\mathcal{L}$ un $\mathcal{O}_X$-module
inversible ample. Pour tout sous-schéma fermé $Z \subset X$, la restriction de
$\mathcal{L}$ à $Z$ est ample.
\end{lemma}

\begin{proof}
C'est clair, car une partie fermée d'un espace quasi-compact est quasi-compacte
et un sous-schéma fermé d'un schéma affine est affine (voir
Schémas, lemme \ref{schemes-lemma-closed-immersion-affine-case}).
\end{proof}

\begin{lemma}
\label{properties-lemma-affine-cap-s-open}
Soit $X$ un schéma. Soit $\mathcal{L}$ un $\mathcal{O}_X$-module
inversible. Soit $s \in \Gamma(X, \mathcal{L})$. Pour tout ouvert affine $U \subset X$,
l'intersection $U \cap X_s$ est affine.
\end{lemma}

\begin{proof}
Cela se traduit par le problème d'algèbre suivant.
Soit $R$ un anneau. Soit $N$ un $R$-module
inversible (c'est-à-dire localement libre de rang 1). Soit $s \in N$.
Alors $V = \{\mathfrak p \mid s \not \in \mathfrak p N\}$ est
un ouvert affine de $\Spec(R)$.

\medskip\noindent
Soit $A = \bigoplus_{n \geq 0} A_n$ l'algèbre symétrique de $N$
(qui est commutative) et considérons $s$ comme un élément de $A_1$.
Posons $B = A/(s - 1)A$. C'est une $R$-algèbre dont la construction
commute à tout changement de base $R \to R'$. Ainsi $B' = B \otimes_R R'$
est l'anneau nul si $s$ a pour image zéro dans $N' = N \otimes_R R'$. Il
s'ensuit que, si $x \in \Spec(R) \setminus V$, alors $B \otimes_R \kappa(x) = 0$.
On en conclut que $\Spec(B) \to \Spec(R)$ se factorise par $V$, car les
fibres aux points $x \not \in V$ sont vides. D'autre part, si
$\Spec(R') \subset V$ est un ouvert affine, alors $s$ a pour image un
élément de base de $N'$ et l'on voit que $B' = R'[s]/(s - 1) \cong R'$.
Il en résulte que $\Spec(B) \to V$ est un isomorphisme et que $V$ est
bien affine.
\end{proof}

\begin{lemma}
\label{properties-lemma-ample-tensor-globally-generated}
\begin{reference}
\cite[II Proposition 4.5.6(ii)]{EGA}
\end{reference}
Soit $X$ un schéma. Soient $\mathcal{L}$ et $\mathcal{M}$
des $\mathcal{O}_X$-modules inversibles. Si
\begin{enumerate}
\item $\mathcal{L}$ est ample, et
\item les ouverts $X_t$, où $t \in \Gamma(X, \mathcal{M}^{\otimes m})$
pour $m > 0$, recouvrent $X$,
\end{enumerate}
alors $\mathcal{L} \otimes \mathcal{M}$ est ample.
\end{lemma}

\begin{proof}
Vérifions les conditions de la définition \ref{properties-definition-ample}.
Comme $\mathcal{L}$ est ample, $X$ est quasi-compact.
Soit $x \in X$. Choisissons $n \geq 1$, $m \geq 1$,
$s \in \Gamma(X, \mathcal{L}^{\otimes n})$ et
$t \in \Gamma(X, \mathcal{M}^{\otimes m})$
tels que $x \in X_s$, $x \in X_t$ et que $X_s$ soit affine.
Alors $s^mt^n \in \Gamma(X, (\mathcal{L} \otimes \mathcal{M})^{\otimes nm})$,
$x \in X_{s^mt^n}$, et $X_{s^mt^n}$ est affine d'après
le lemme \ref{properties-lemma-affine-cap-s-open}.
\end{proof}

\begin{lemma}
\label{properties-lemma-affine-s-opens}
Soit $X$ un schéma. Soit $\mathcal{L}$ un $\mathcal{O}_X$-module
inversible. Supposons que les ouverts $X_s$, où $s \in \Gamma(X, \mathcal{L}^{\otimes n})$
et $n \geq 1$, forment une base de la topologie de $X$.
Alors, parmi ces ouverts, ceux des $X_s$ qui sont affines
forment une base de la topologie de $X$.
\end{lemma}

\begin{proof}
Soit $x \in X$. Choisissons un voisinage ouvert affine
$\Spec(R) = U \subset X$ de $x$.
Par hypothèse, il existe un
entier $n \geq 1$ et une section $s \in \Gamma(X, \mathcal{L}^{\otimes n})$
tels que $X_s \subset U$. D'après le lemme \ref{properties-lemma-affine-cap-s-open} ci-dessus,
l'intersection $X_s = U \cap X_s$ est affine. Comme on
peut choisir $U$ arbitrairement petit, cela donne la conclusion.
\end{proof}

\begin{lemma}
\label{properties-lemma-affine-s-opens-cover-quasi-separated}
Soient $X$ un schéma et $\mathcal{L}$ un $\mathcal{O}_X$-module
inversible. Supposons que, pour tout point $x$ de $X$, il existe $n \geq 1$ et
$s \in \Gamma(X, \mathcal{L}^{\otimes n})$ tels que
$x \in X_s$ et que $X_s$ soit affine. Alors $X$ est séparé.
\end{lemma}

\begin{proof}
Montrons d'abord que $X$ est quasi-séparé. Par hypothèse, on peut
trouver un recouvrement de $X$ par des ouverts affines de la forme
$X_s$. D'après le lemme \ref{properties-lemma-affine-cap-s-open},
l'intersection de deux tels ouverts est affine ; Schémas, lemme
\ref{schemes-lemma-characterize-quasi-separated}, entraîne
donc que $X$ est quasi-séparé.

\medskip\noindent
Pour montrer que $X$ est séparé, on peut utiliser le critère
valuatif, Schémas, lemme
\ref{schemes-lemma-valuative-criterion-separatedness}.
Soit donc $A$ un anneau de valuation de corps des fractions $K$,
et considérons deux morphismes $f, g : \Spec(A) \to X$ tels que
les deux composés $\Spec(K) \to \Spec(A) \to X$ coïncident.
Comme $A$ est local, il existe $p, q \ge 1$, $s \in \Gamma(X,
\mathcal{L}^{\otimes p})$et $ t \in \Gamma(X,
\mathcal{L}^{\otimes q})$ tels que $X_s$ et $X_t$ soient
affines, $f(\Spec A) \subseteq X_s$et $g(\Spec A)
\subseteq X_t$. Remplaçons maintenant $s$ par $s^q$, $t$ par $t^p$et
 $\mathcal{L}$ par $\mathcal{L}^{\otimes pq}$. Cela ne
change rien, puisque $X_s = X_{s^q}$ et $X_t = X_{t^p}$, et désormais $s$
et $t$ sont des sections du même faisceau $\mathcal{L}$.

\medskip\noindent
Le module quasi-cohérent $f^*\mathcal{L}$ correspond à un $A$-module $M$ et
$g^*\mathcal{L}$ correspond à un $A$-module $N$,
d'après notre classification des modules quasi-cohérents sur les schémas
affines (Schémas, lemme \ref{schemes-lemma-quasi-coherent-affine}).
Les $A$-modules $M$ et $N$ sont localement libres de rang
$1$ (lemme \ref{properties-lemma-locally-free-module}) et, puisque $A$ est
local, ils sont libres (Algèbre, lemme
\ref{algebra-lemma-K0-local}). On peut donc identifier
$M$ et $N$ à des $A$-sous-modules de $M \otimes_A K$ et $N
\otimes_A K$. L'égalité $f|_{\Spec(K)} = g|_{\Spec(K)}$
détermine un isomorphisme $\phi \colon M \otimes_A K \to N
\otimes_A K$.

\medskip\noindent
Soient $x \in M$ et $y \in N$ les éléments correspondant aux
images réciproques de $s$ par $f$ et $g$, respectivement. Ils
vérifient $\phi(x \otimes 1) = y \otimes 1$. L'image de $f$
est contenue dans $X_s$, donc $x \not\in \mathfrak{m}_A M$,
autrement dit $x$ engendre $M$. Ainsi $\phi$ détermine un
isomorphisme de $M$ sur le sous-module de $N$ engendré par
$y$. En raisonnant symétriquement à l'aide de $t$, on voit que $\phi^{-1}$
détermine un isomorphisme de $N$ sur un sous-module de $M$.
Par conséquent, $\phi$ se restreint à un isomorphisme de $M$ sur
$N$. Puisque $x$ engendre $M$, son image $y$ engendre
$N$, d'où $y \not\in \mathfrak{m}_A N$. On a donc
$g(\Spec(A)) \subseteq X_s$. Comme $X_s$ est affine, il est
séparé d'après Schémas, lemme \ref{schemes-lemma-affine-separated},
et l'on en conclut que $f = g$.
\end{proof}

\begin{lemma}
\label{properties-lemma-ample-separated}
Soit $X$ un schéma. S'il existe un faisceau inversible ample sur $X$,
alors $X$ est séparé.
\end{lemma}

\begin{proof}
Cela résulte immédiatement du
lemme \ref{properties-lemma-affine-s-opens-cover-quasi-separated} et
de la définition \ref{properties-definition-ample}.
\end{proof}

\begin{lemma}
\label{properties-lemma-map-into-proj}
Soit $X$ un schéma.
Soit $\mathcal{L}$ un $\mathcal{O}_X$-module inversible.
Posons $S = \Gamma_*(X, \mathcal{L})$, considéré comme anneau
gradué. Si tout point de $X$ appartient à l'un
des ouverts $X_s$, pour un élément homogène $s \in S_{+}$, alors il
existe un morphisme canonique de schémas
$$
f : X \longrightarrow Y = \text{Proj}(S),
$$
à valeurs dans le spectre homogène de $S$ (voir
Constructions, section \ref{constructions-section-proj}).
Ce morphisme possède les propriétés suivantes :
\begin{enumerate}
\item $f^{-1}(D_{+}(s)) = X_s$ pour tout élément homogène $s \in S_{+}$,
\item il existe des homomorphismes de $\mathcal{O}_X$-modules
$f^*\mathcal{O}_Y(n) \to \mathcal{L}^{\otimes n}$
compatibles aux applications de multiplication,
voir Constructions, équation (\ref{constructions-equation-multiply}),
\item le composé
$S_n \to \Gamma(Y, \mathcal{O}_Y(n)) \to \Gamma(X, \mathcal{L}^{\otimes n})$
est l'application identique, et
\item pour tout $x \in X$, il existe un entier $d \geq 1$
et un voisinage ouvert $U \subset X$ de $x$
tels que $f^*\mathcal{O}_Y(dn)|_U \to \mathcal{L}^{\otimes dn}|_U$
soit un isomorphisme pour tout $n \in \mathbf{Z}$.
\end{enumerate}
\end{lemma}

\begin{proof}
Notons $\psi : S \to \Gamma_*(X, \mathcal{L})$ l'application
identique. Nous allons utiliser le triplet
$(U(\psi), r_{\mathcal{L}, \psi}, \theta)$ de
Constructions, lemme \ref{constructions-lemma-invertible-map-into-proj}.
Par hypothèse, le sous-schéma ouvert $U(\psi)$ est égal à $X$. Ainsi
$r_{\mathcal{L}, \psi} : U(\psi) \to Y$ est défini sur tout $X$.
Posons $f = r_{\mathcal{L}, \psi}$.
Les applications de (2) sont les composantes de $\theta$.
L'assertion (3) résulte de la condition (2) du lemme cité
ci-dessus. L'assertion (1) résulte de (3), combiné à la condition (1) du lemme
cité ci-dessus. L'assertion (4) résulte de la dernière assertion
de Constructions, lemme \ref{constructions-lemma-invertible-map-into-proj},
puisque l'application $\alpha$ qui y figure est un isomorphisme.
\end{proof}

\begin{lemma}
\label{properties-lemma-map-into-proj-quasi-compact}
Soit $X$ un schéma. Soit $\mathcal{L}$ un $\mathcal{O}_X$-module inversible.
Posons $S = \Gamma_*(X, \mathcal{L})$.
Supposons (a) que tout point de $X$ appartienne à l'un
des ouverts $X_s$, pour un élément homogène $s \in S_{+}$, et (b)
que $X$ soit quasi-compact. Alors le morphisme canonique de schémas
$f : X \longrightarrow \text{Proj}(S)$ du lemme \ref{properties-lemma-map-into-proj}
ci-dessus est quasi-compact et d'image dense.
\end{lemma}

\begin{proof}
Pour montrer que $f$ est quasi-compact, il suffit de montrer que $f^{-1}(D_{+}(s))$
est quasi-compact pour tout élément homogène $s \in S_{+}$. Écrivons
$X = \bigcup_{i = 1, \ldots, n} X_i$ comme réunion finie
d'ouverts affines. D'après le lemme \ref{properties-lemma-affine-cap-s-open}, chaque intersection
$X_s \cap X_i$ est affine. Ainsi $X_s = \bigcup_{i = 1, \ldots, n} X_s \cap X_i$
est quasi-compact. Supposons que l'image de $f$ ne soit pas dense et
cherchons une contradiction. Puisque les ouverts $D_+(s)$, avec $s \in S_+$ homogène,
forment une base de la topologie de $\text{Proj}(S)$, on peut trouver un
tel $s$ avec $D_+(s) \not = \emptyset$ et $f(X) \cap D_+(s) = \emptyset$.
D'après le lemme \ref{properties-lemma-map-into-proj},
cela signifie que $X_s = \emptyset$. D'après le lemme \ref{properties-lemma-invert-s-sections},
une puissance $s^n$ est alors la section nulle de
$\mathcal{L}^{\otimes n\deg(s)}$.
Cela entraîne à son tour que $D_+(s) = \emptyset$,
contradiction cherchée.
\end{proof}

\begin{lemma}
\label{properties-lemma-ample-immersion-into-proj}
Soit $X$ un schéma. Soit $\mathcal{L}$ un $\mathcal{O}_X$-module inversible.
Posons $S = \Gamma_*(X, \mathcal{L})$.
Supposons $\mathcal{L}$ ample. Alors le morphisme canonique de schémas
$f : X \longrightarrow \text{Proj}(S)$ du lemme \ref{properties-lemma-map-into-proj}
est une immersion ouverte d'image dense.
\end{lemma}

\begin{proof}
D'après le lemme \ref{properties-lemma-affine-s-opens-cover-quasi-separated},
 $X$ est quasi-séparé. Choisissons un nombre fini d'éléments homogènes
$s_1, \ldots, s_n \in S_{+}$ tels
que les $X_{s_i}$ soient affines et que $X = \bigcup X_{s_i}$.
Supposons $s_i$ de degré $d_i$. L'image réciproque de
$D_{+}(s_i)$ par $f$ est $X_{s_i}$, voir le lemme \ref{properties-lemma-map-into-proj}.
D'après le lemme \ref{properties-lemma-invert-s-sections}, l'homomorphisme d'anneaux
$$
(S^{(d_i)})_{(s_i)} = \Gamma(D_{+}(s_i), \mathcal{O}_{\text{Proj}(S)})
\longrightarrow
\Gamma(X_{s_i}, \mathcal{O}_X)
$$
est un isomorphisme. Ainsi $f$ induit un isomorphisme
$X_{s_i} \to D_{+}(s_i)$. Donc $f$ est un isomorphisme de $X$ sur le
sous-schéma ouvert $\bigcup_{i = 1, \ldots, n} D_{+}(s_i)$ de $\text{Proj}(S)$.
Son image est dense d'après le lemme \ref{properties-lemma-map-into-proj-quasi-compact}.
\end{proof}

\begin{lemma}
\label{properties-lemma-open-in-proj-ample}
Soit $X$ un schéma.
Soit $S$ un anneau gradué. Supposons $X$ quasi-compact et
supposons qu'il existe une immersion ouverte
$$
j : X \longrightarrow Y = \text{Proj}(S).
$$
Alors $j^*\mathcal{O}_Y(d)$ est un faisceau inversible ample pour
un entier $d > 0$.
\end{lemma}

\begin{proof}
C'est Constructions, lemme \ref{constructions-lemma-ample-on-proj}.
\end{proof}

\begin{proposition}
\label{properties-proposition-characterize-ample}
Soit $X$ un schéma quasi-compact.
Soit $\mathcal{L}$ un faisceau inversible sur $X$.
Posons $S = \Gamma_*(X, \mathcal{L})$.
Les conditions suivantes sont équivalentes :
\begin{enumerate}
\item
\label{properties-item-ample}
$\mathcal{L}$ est ample,
\item
\label{properties-item-immersion}
les ouverts $X_s$, avec $s \in S_{+}$ homogène,
recouvrent $X$ et le morphisme associé $X \to \text{Proj}(S)$
est une immersion ouverte,
\item
\label{properties-item-s-basis}
les ouverts $X_s$, avec $s \in S_{+}$ homogène,
forment une base de la topologie de $X$,
\item
\label{properties-item-s-affine-basis}
les ouverts $X_s$, avec $s \in S_{+}$ homogène,
qui sont affines forment une base de la topologie de $X$,
\item
\label{properties-item-qc-gg}
pour tout faisceau quasi-cohérent $\mathcal{F}$ sur $X$,
la somme des images des applications canoniques
$$
\Gamma(X, \mathcal{F} \otimes_{\mathcal{O}_X} \mathcal{L}^{\otimes n})
\otimes_{\mathbf{Z}} \mathcal{L}^{\otimes -n}
\longrightarrow
\mathcal{F}
$$
pour $n \geq 1$ est égale à $\mathcal{F}$,
\item
\label{properties-item-qc-i-gg}
même propriété qu'en (\ref{properties-item-qc-gg}), où $\mathcal{F}$
parcourt tous les faisceaux quasi-cohérents d'idéaux,
\item
\label{properties-item-c-gg}
$X$ est quasi-séparé et,
pour tout faisceau quasi-cohérent $\mathcal{F}$ de type fini sur $X$,
il existe un entier $n_0$ tel que
$\mathcal{F} \otimes_{\mathcal{O}_X} \mathcal{L}^{\otimes n}$
soit engendré par ses sections globales pour tout $n \geq n_0$,
\item
\label{properties-item-c-q}
$X$ est quasi-séparé et,
pour tout faisceau quasi-cohérent $\mathcal{F}$ de type fini sur $X$,
il existe des entiers $n > 0$, $k \geq 0$ tels que
$\mathcal{F}$ soit quotient d'une somme directe de $k$ exemplaires de
$\mathcal{L}^{\otimes - n}$, et
\item
\label{properties-item-c-i-q}
même propriété qu'en (\ref{properties-item-c-q}), où $\mathcal{F}$ parcourt tous les
faisceaux d'idéaux de type fini sur $X$.
\end{enumerate}
\end{proposition}

\begin{proof}
Le lemme \ref{properties-lemma-ample-immersion-into-proj} donne
(\ref{properties-item-ample}) $\Rightarrow$ (\ref{properties-item-immersion}).
Les lemmes \ref{properties-lemma-ample-power-ample} et \ref{properties-lemma-open-in-proj-ample}
donnent l'implication
(\ref{properties-item-ample}) $\Leftarrow$ (\ref{properties-item-immersion}).
Les implications (\ref{properties-item-immersion}) $\Rightarrow$
(\ref{properties-item-s-affine-basis}) $\Rightarrow$ (\ref{properties-item-s-basis})
résultent clairement de Constructions, section \ref{constructions-section-proj}.
Le lemme \ref{properties-lemma-affine-s-opens} donne
(\ref{properties-item-s-basis}) $\Rightarrow$ (\ref{properties-item-ample}).
Les quatre premières conditions sont donc équivalentes.

\medskip\noindent
Supposons vérifiées les conditions équivalentes (1) --
(4). Notons qu'en particulier $X$ est séparé (comme
sous-schéma ouvert du schéma séparé $\text{Proj}(S)$).
Soit $\mathcal{F}$ un faisceau quasi-cohérent sur $X$.
Choisissons $s \in S_{+}$ homogène tel que $X_s$ soit affine.
Nous affirmons que toute section $m \in \Gamma(X_s, \mathcal{F})$
appartient à l'image de l'une des applications figurant dans
(\ref{properties-item-qc-gg}) ci-dessus. Cela entraînera (\ref{properties-item-qc-gg}),
puisque ces ouverts affines $X_s$ recouvrent $X$.
En effet, d'après le lemme \ref{properties-lemma-invert-s-sections}, on peut écrire
$m$ comme image de $m' \otimes s^{-n}$ pour un certain
$n \geq 1$ et un certain
$m' \in \Gamma(X, \mathcal{F} \otimes \mathcal{L}^{\otimes n})$.
Cela prouve l'affirmation.

\medskip\noindent
Clairement, (\ref{properties-item-qc-gg}) $\Rightarrow$ (\ref{properties-item-qc-i-gg}).
Supposons (\ref{properties-item-qc-i-gg}) et montrons que $\mathcal{L}$ est
ample. Choisissons $x \in X$. Soit $U \subset X$ un ouvert
affine contenant $x$. Posons $Z = X \setminus U$. On peut considérer
$Z$ comme un sous-schéma fermé réduit,
voir Schémas, section \ref{schemes-section-reduced}.
Soit $\mathcal{I} \subset \mathcal{O}_X$ le faisceau
quasi-cohérent d'idéaux correspondant au sous-schéma fermé $Z$.
D'après l'hypothèse (\ref{properties-item-qc-i-gg}), il existe un entier $n \geq 1$ et une section
$s \in \Gamma(X, \mathcal{I} \otimes \mathcal{L}^{\otimes n})$
tels que $s$ ne s'annule pas en $x$ (plus précisément, tels que
$s \not \in \mathfrak m_x \mathcal{I}_x \otimes \mathcal{L}_x^{\otimes n}$).
On peut considérer $s$ comme une section de $\mathcal{L}^{\otimes n}$.
Puisqu'elle s'annule clairement le long de $Z$, on a
$X_s \subset U$. Ainsi $X_s$ est affine, voir
le lemme \ref{properties-lemma-affine-cap-s-open}.
Cela prouve que $\mathcal{L}$ est ample.
Nous avons donc montré à ce stade que (1) -- (6) sont équivalentes.

\medskip\noindent
Supposons vérifiées les conditions équivalentes (1) -- (6). Dans la suite,
nous utiliserons sans autre mention le fait que le produit tensoriel de deux
faisceaux de modules engendrés par leurs sections globales est lui-même engendré par
ses sections globales
(voir Modules, lemme \ref{modules-lemma-tensor-product-globally-generated}).
D'après (1), on peut trouver des éléments $s_i \in S_{d_i}$ avec $d_i \geq 1$
tels que $X = \bigcup_{i = 1, \ldots, n} X_{s_i}$.
Posons $d = d_1\ldots d_n$. Il en résulte que $\mathcal{L}^{\otimes d}$
est engendré par les sections globales
$$
s_1^{d/d_1}, \ldots, s_n^{d/d_n}.
$$
Ainsi, si $\mathcal{L}^{\otimes j}$ est engendré par ses sections globales,
il en est de même de $\mathcal{L}^{\otimes j + dn}$ pour tout $n \geq 0$.
Fixons $j \in \{0, \ldots, d - 1\}$. Pour tout point $x \in X$, il
existe un entier $n \geq 1$ et une section globale $s$ de $\mathcal{L}^{j + dn}$
qui ne s'annule pas en $x$, comme il résulte de (\ref{properties-item-qc-gg}) appliqué
à $\mathcal{F} = \mathcal{L}^{\otimes j}$ et au faisceau inversible
ample $\mathcal{L}^{\otimes d}$. Comme $X$ est quasi-compact, on
peut trouver une famille finie d'entiers $n_i$ et de sections globales
$s_i$ de $\mathcal{L}^{\otimes j + dn_i}$ qui n'ont aucun zéro commun
sur $X$. Puisque $\mathcal{L}^{\otimes d}$ est engendré par ses sections globales,
$\mathcal{L}^{\otimes j + dn}$ est engendré par ses sections globales pour $n = \max\{n_i\}$.
Cela étant établi pour chaque classe de congruence modulo $d$,
on en conclut qu'il existe un entier $n_0 = n_0(\mathcal{L})$ tel que
$\mathcal{L}^{\otimes n}$ soit engendré par ses sections globales pour tout $n \geq n_0$.
On voit dès lors que, si $\mathcal{F}$ est engendré par ses sections globales, il en est
de même de $\mathcal{F} \otimes \mathcal{L}^{\otimes n}$ pour tout
$n \geq n_0$.

\medskip\noindent
Supposons toujours vérifiées les conditions équivalentes (1) -- (6).
Soit $\mathcal{F}$ un faisceau
quasi-cohérent de $\mathcal{O}_X$-modules de type fini.
Notons $\mathcal{F}_n \subset \mathcal{F}$ l'image de l'application
canonique de (\ref{properties-item-qc-gg}). Par construction,
$\mathcal{F}_n \otimes \mathcal{L}^{\otimes n}$ est
engendré par ses sections globales. D'après (\ref{properties-item-qc-gg}),
$\mathcal{F}$ est la somme des sous-faisceaux $\mathcal{F}_n$,
$n \geq 1$.
D'après Modules, lemme \ref{modules-lemma-finite-type-quasi-compact-colimit},
on a $\mathcal{F} = \sum_{n = 1, \ldots, N} \mathcal{F}_n$
pour un entier $N \geq 1$. Il en résulte que
$\mathcal{F} \otimes \mathcal{L}^{\otimes n}$ est engendré par
ses sections globales dès que $n \geq N + n_0(\mathcal{L})$, où $n_0(\mathcal{L})$
est défini ci-dessus. Ainsi (1) -- (6) entraînent (\ref{properties-item-c-gg}).

\medskip\noindent
Supposons (\ref{properties-item-c-gg}). Soit $\mathcal{F}$ un faisceau
quasi-cohérent de $\mathcal{O}_X$-modules de type fini.
D'après (\ref{properties-item-c-gg}), il existe un entier $n \geq 1$ tel que
l'application canonique
$$
\Gamma(X, \mathcal{F} \otimes_{\mathcal{O}_X} \mathcal{L}^{\otimes n})
\otimes_{\mathbf{Z}} \mathcal{L}^{\otimes -n}
\longrightarrow
\mathcal{F}
$$
soit surjective. Soit $I$ l'ensemble des parties finies de
$\Gamma(X, \mathcal{F} \otimes_{\mathcal{O}_X} \mathcal{L}^{\otimes n})$,
ordonné par inclusion. Alors $I$ est un ensemble ordonné filtrant.
Pour $i = \{s_1, \ldots, s_{r(i)}\}$, soit $\mathcal{F}_i \subset \mathcal{F}$
l'image de l'application
$$
\bigoplus\nolimits_{j = 1, \ldots, r(i)} \mathcal{L}^{\otimes -n}
\longrightarrow
\mathcal{F}
$$
qui est la multiplication par $s_j$ sur le $j$-ième facteur. La surjectivité
ci-dessus entraîne que $\mathcal{F} = \colim_{i \in I} \mathcal{F}_i$.
On peut donc appliquer Modules, lemme \ref{modules-lemma-finite-type-quasi-compact-colimit},
et l'on en conclut que
$\mathcal{F} = \mathcal{F}_i$ pour un certain $i$.
Nous avons donc établi (\ref{properties-item-c-q}). Autrement dit,
(\ref{properties-item-c-gg}) $\Rightarrow$ (\ref{properties-item-c-q}).

\medskip\noindent
L'implication (\ref{properties-item-c-q}) $\Rightarrow$ (\ref{properties-item-c-i-q}) est immédiate.

\medskip\noindent
Supposons enfin (\ref{properties-item-c-i-q}).
Soit $\mathcal{I} \subset \mathcal{O}_X$ un faisceau quasi-cohérent
d'idéaux. D'après le lemme \ref{properties-lemma-quasi-coherent-colimit-finite-type}
(c'est ici qu'intervient l'hypothèse que $X$ est
quasi-séparé), on a $\mathcal{I} = \colim_\alpha I_\alpha$, où
chaque $I_\alpha$ est quasi-cohérent de type fini. Puisque, par hypothèse,
chaque $I_\alpha$ est quotient d'une somme de puissances tensorielles négatives de
$\mathcal{L}$, il en est de même de $\mathcal{I}$ (mais, bien
entendu, sans finitude ni borne sur les puissances). On en
conclut donc que (\ref{properties-item-c-i-q}) entraîne (\ref{properties-item-qc-i-gg}).
Cela achève la démonstration de la proposition.
\end{proof}

\begin{lemma}
\label{properties-lemma-ample-on-locally-closed}
Soit $X$ un schéma. Soit $\mathcal{L}$ un
$\mathcal{O}_X$-module inversible ample. Soit $i : X' \to X$ un morphisme de schémas.
Supposons vérifiée au moins l'une des conditions suivantes :
\begin{enumerate}
\item $i$ est une immersion quasi-compacte,
\item $X'$ est quasi-compact et $i$ est une immersion,
\item $i$ est quasi-compact et induit un
homéomorphisme entre $X'$ et $i(X')$,
\item $X'$ est quasi-compact et $i$ induit un
homéomorphisme entre $X'$ et $i(X')$.
\end{enumerate}
Alors $i^*\mathcal{L}$ est ample sur $X'$.
\end{lemma}

\begin{proof}
Observons que, dans les cas (1) et (3), le schéma $X'$ est quasi-compact,
puisque $X$ est quasi-compact d'après la définition \ref{properties-definition-ample}.
Il suffit donc de traiter les cas (2) et (4). Comme (2) est un cas
particulier de (4), il suffit de traiter (4).

\medskip\noindent
Supposons vérifiée la condition (4). Pour $s \in \Gamma(X, \mathcal{L}^{\otimes d})$,
notons $s' = i^*s$ l'image réciproque de $s$ sur $X'$. Alors
$s'$ est une section de $(i^*\mathcal{L})^{\otimes d}$.
D'après la proposition \ref{properties-proposition-characterize-ample},
les ouverts $X_s$, pour $s \in \Gamma(X, \mathcal{L}^{\otimes d})$,
forment une base de la topologie de $X$. Comme $X'_{s'} = i^{-1}(X_s)$ d'après
Modules, remarque \ref{modules-remark-pullback-s-open},
et comme $X' \to i(X')$ est un homéomorphisme, on en conclut que
les ouverts $X'_{s'}$ forment une base de la topologie de $X'$. Ainsi
$i^*\mathcal{L}$ est ample
d'après la proposition \ref{properties-proposition-characterize-ample}.
\end{proof}

\begin{lemma}
\label{properties-lemma-ample-on-product}
Soit $S$ un schéma quasi-séparé. Soient $X$, $Y$ des schémas sur $S$.
Soit $\mathcal{L}$ un $\mathcal{O}_X$-module
inversible ample et soit $\mathcal{N}$ un $\mathcal{O}_Y$-module inversible
ample. Alors $\mathcal{M} = \text{pr}_1^*\mathcal{L}
\otimes_{\mathcal{O}_{X \times_S Y}} \text{pr}_2^*\mathcal{N}$
est un faisceau inversible ample sur $X \times_S Y$.
\end{lemma}

\begin{proof}
Le morphisme $i : X \times_S Y \to X \times Y$ est une immersion
quasi-compacte, voir Schémas, lemme \ref{schemes-lemma-fibre-product-after-map}.
D'autre part, $\mathcal{M}$ est l'image réciproque par
$i$ du module inversible correspondant sur $X \times Y$.
D'après le lemme \ref{properties-lemma-ample-on-locally-closed}, il suffit de démontrer le
lemme pour $X \times Y$. Vérifions (1) et (2) de la
définition \ref{properties-definition-ample} pour $\mathcal{M}$ sur $X \times Y$.

\medskip\noindent
Puisque $X$ et $Y$ sont quasi-compacts, il en est de même de $X \times Y$.
Soit $z \in X \times Y$ un point. Soient $x \in X$ et $y \in Y$
ses projections. Choisissons $n > 0$ et
$s \in \Gamma(X, \mathcal{L}^{\otimes n})$
tels que $X_s$ soit un voisinage ouvert affine de $x$.
Choisissons $m > 0$ et
$t \in \Gamma(Y, \mathcal{N}^{\otimes m})$
tels que $Y_t$ soit un voisinage ouvert affine de $y$.
Alors $r = \text{pr}_1^*s \otimes \text{pr}_2^*t$ est une section
de $\mathcal{M}$ avec $(X \times Y)_r = X_s \times Y_t$.
C'est un voisinage ouvert affine de $z$, ce qui achève la démonstration.
\end{proof}












\section{Schémas affines et quasi-affines}
\label{properties-section-affine-quasi-affine}

\begin{lemma}
\label{properties-lemma-quasi-affine-O-ample}
Soit $X$ un schéma.
Alors $X$ est quasi-affine si et seulement si $\mathcal{O}_X$ est ample.
\end{lemma}

\begin{proof}
Supposons $X$ quasi-affine. Posons $A = \Gamma(X, \mathcal{O}_X)$.
Considérons l'immersion ouverte
$$
j : X \longrightarrow \Spec(A)
$$
du lemme \ref{properties-lemma-quasi-affine}. Notons que
$\Spec(A) = \text{Proj}(A[T])$,
voir Constructions, exemple \ref{constructions-example-trivial-proj}.
On peut donc appliquer le lemme \ref{properties-lemma-open-in-proj-ample}
pour en déduire que $\mathcal{O}_X$ est ample.

\medskip\noindent
Supposons $\mathcal{O}_X$ ample.
Notons que $\Gamma_*(X, \mathcal{O}_X) \cong A[T]$
comme anneaux gradués. Le résultat découle donc des lemmes
\ref{properties-lemma-ample-immersion-into-proj} et \ref{properties-lemma-quasi-affine},
compte tenu de l'égalité
$\Spec(A) = \text{Proj}(A[T])$ pour tout anneau $A$,
établie ci-dessus.
\end{proof}

\begin{lemma}
\label{properties-lemma-quasi-affine-locally-closed}
Soit $X$ un schéma quasi-affine. Pour toute immersion quasi-compacte
$i : X' \to X$, le schéma $X'$ est quasi-affine.
\end{lemma}

\begin{proof}
On peut le démontrer directement, sans faire appel aux résultats sur les
faisceaux inversibles amples ; nous invitons le lecteur à en esquisser une preuve.
Comme $X$ est quasi-affine, $\mathcal{O}_X$ est ample d'après
le lemme \ref{properties-lemma-quasi-affine-O-ample}.
Alors $\mathcal{O}_{X'}$ est ample d'après
le lemme \ref{properties-lemma-ample-on-locally-closed}. Donc $X'$ est quasi-affine d'après
le lemme \ref{properties-lemma-quasi-affine-O-ample}.
\end{proof}

\begin{lemma}
\label{properties-lemma-characterize-affine}
Soit $X$ un schéma. Supposons qu'il existe un nombre fini d'éléments
$f_1, \ldots, f_n \in \Gamma(X, \mathcal{O}_X)$ tels que
\begin{enumerate}
\item chaque $X_{f_i}$ soit un ouvert affine de $X$, et
\item l'idéal engendré par $f_1, \ldots, f_n$ dans
$\Gamma(X, \mathcal{O}_X)$ soit l'idéal unité.
\end{enumerate}
Alors $X$ est affine.
\end{lemma}

\begin{proof}
Supposons donnés $f_1, \ldots, f_n$ comme dans l'énoncé.
On peut écrire $1 = \sum g_i f_i$ avec $g_j \in \Gamma(X, \mathcal{O}_X)$,
d'où clairement $X = \bigcup X_{f_i}$. (Les $f_i$ ne peuvent
pas tous s'annuler en un même point.) Puisque chaque $X_{f_i}$
est quasi-compact (étant affine), $X$ est quasi-compact.
Il en résulte que $X$ est quasi-affine d'après
le lemme \ref{properties-lemma-quasi-affine-O-ample} ci-dessus.
Considérons l'immersion ouverte
$$
j : X \to \Spec(\Gamma(X, \mathcal{O}_X)),
$$
voir le lemme \ref{properties-lemma-quasi-affine}. L'image réciproque de l'ouvert standard
$D(f_i)$ du membre de droite est $X_{f_i}$ dans le
membre de gauche, et le morphisme $j$ induit un isomorphisme
$X_{f_i} \cong D(f_i)$, voir
le lemme \ref{properties-lemma-invert-f-affine}. Puisque les $f_i$ engendrent l'idéal
unité, on a $\Spec(\Gamma(X, \mathcal{O}_X))
= \bigcup_{i = 1, \ldots, n} D(f_i)$. Ainsi $j$ est un isomorphisme.
\end{proof}







\section{Faisceaux quasi-cohérents et faisceaux inversibles amples}
\label{properties-section-ample-quasi-coherent}

\noindent
Thème de cette section : en présence d'un faisceau inversible ample,
tout faisceau quasi-cohérent provient d'un module gradué.

\begin{situation}
\label{properties-situation-ample}
Soit $X$ un schéma.
Soit $\mathcal{L}$ un faisceau inversible ample sur $X$.
Posons $S = \Gamma_*(X, \mathcal{L})$, considéré comme anneau gradué.
Posons $Y = \text{Proj}(S)$.
Soit $f : X \to Y$ le morphisme canonique du lemme \ref{properties-lemma-map-into-proj}.
Il est muni d'un homomorphisme $\mathbf{Z}$-gradué de $\mathcal{O}_X$-algèbres
$\bigoplus f^*\mathcal{O}_Y(n) \to \bigoplus \mathcal{L}^{\otimes n}$.
\end{situation}

\noindent
Le lemme suivant est en fait un cas particulier de celui
qui le suit, mais il paraît utile de l'établir d'abord.

\begin{lemma}
\label{properties-lemma-ample-gcd-is-one}
Plaçons-nous dans la situation \ref{properties-situation-ample}.
Le morphisme canonique $f : X \to Y$
envoie $X$ dans le sous-schéma ouvert $W = W_1 \subset Y$
sur lequel $\mathcal{O}_Y(1)$ est inversible et toutes les
applications de multiplication
$\mathcal{O}_Y(n) \otimes_{\mathcal{O}_Y} \mathcal{O}_Y(m) \to
\mathcal{O}_Y(n + m)$
sont des isomorphismes
(voir Constructions, lemme \ref{constructions-lemma-where-invertible}).
De plus, les applications $f^*\mathcal{O}_Y(n) \to \mathcal{L}^{\otimes n}$
sont toutes des isomorphismes.
\end{lemma}

\begin{proof}
D'après la proposition \ref{properties-proposition-characterize-ample}, il existe un entier
$n_0$ tel que $\mathcal{L}^{\otimes n}$ soit engendré par ses sections globales pour tout
$n \geq n_0$. Soit $x \in X$. D'après ce qui précède, on peut trouver
$a \in S_{n_0}$ et $b \in S_{n_0 + 1}$ tels que
$a$ et $b$ ne s'annulent pas en $x$. Ainsi
$f(x) \in D_{+}(a) \cap D_{+}(b) = D_{+}(ab)$.
D'après Constructions, lemme \ref{constructions-lemma-where-invertible},
on a $f(x) \in W_1$, comme voulu.
D'après Constructions, lemme \ref{constructions-lemma-invertible-map-into-proj},
utilisé dans la construction de l'application $f$,
les applications
$f^*\mathcal{O}_Y(n_0) \to \mathcal{L}^{\otimes n_0}$ et
$f^*\mathcal{O}_Y(n_0 + 1) \to \mathcal{L}^{\otimes n_0 + 1}$
sont des isomorphismes dans un voisinage de $x$. Par compatibilité
avec la structure d'algèbre et du fait que $f$ prend ses valeurs dans $W$,
on en conclut que toutes les applications
$f^*\mathcal{O}_Y(n) \to \mathcal{L}^{\otimes n}$ sont des
isomorphismes dans un voisinage de $x$, d'où la conclusion.
\end{proof}

\noindent
Rappelons d'après Modules, définition \ref{modules-definition-gamma-star},
qu'étant donnés un espace localement annelé $X$, un faisceau inversible $\mathcal{L}$ et
un $\mathcal{O}_X$-module $\mathcal{F}$, on dispose du
$\Gamma_*(X, \mathcal{L})$-module gradué
$$
\Gamma_*(X, \mathcal{L}, \mathcal{F}) =
\bigoplus\nolimits_{n \in \mathbf{Z}}
\Gamma(X, \mathcal{F} \otimes_{\mathcal{O}_X} \mathcal{L}^{\otimes n}).
$$
Le lemme suivant affirme que, dans la situation \ref{properties-situation-ample},
on peut retrouver un $\mathcal{O}_X$-module quasi-cohérent $\mathcal{F}$
à partir de ce module gradué. Voir aussi
Constructions, lemme \ref{constructions-lemma-quasi-coherent-projective-space},
où ce lemme est démontré dans le cas particulier $X = \mathbf{P}^n_R$.

\begin{lemma}
\label{properties-lemma-ample-quasi-coherent}
Plaçons-nous dans la situation \ref{properties-situation-ample}.
Soit $\mathcal{F}$ un faisceau quasi-cohérent sur $X$.
Posons $M = \Gamma_*(X, \mathcal{L}, \mathcal{F})$, considéré comme $S$-module
gradué. Il existe des isomorphismes
$$
f^*\widetilde{M} \longrightarrow \mathcal{F}
$$
fonctoriels en $\mathcal{F}$, tels que
$M_0 \to \Gamma(\text{Proj}(S), \widetilde{M}) \to \Gamma(X, \mathcal{F})$
soit l'application identique.
\end{lemma}

\begin{proof}
Soit $s \in S_{+}$ homogène tel que $X_s$ soit un ouvert affine de $X$.
Rappelons que $\widetilde{M}|_{D_{+}(s)}$ correspond au
$S_{(s)}$-module $M_{(s)}$,
voir Constructions, lemme \ref{constructions-lemma-proj-sheaves}.
Rappelons que $f^{-1}(D_{+}(s)) = X_s$.
Puisque $X$ possède un faisceau inversible ample, il est quasi-compact et
quasi-séparé, voir la section \ref{properties-section-ample}.
D'après le lemme \ref{properties-lemma-invert-s-sections}, il existe un isomorphisme canonique
$M_{(s)} = \Gamma_*(X, \mathcal{L}, \mathcal{F})_{(s)} \to
\Gamma(X_s, \mathcal{F})$.
Puisque $\mathcal{F}$ est quasi-cohérent, cela donne
un isomorphisme canonique
$$
f^*\widetilde{M}|_{X_s} \to \mathcal{F}|_{X_s}
$$
Comme $\mathcal{L}$ est ample sur $X$, on sait que $X$ est
recouvert par les ouverts affines de la forme $X_s$. Il suffit donc de montrer que les
applications ci-dessus se recollent sur les intersections. Nous en omettons la
démonstration.
\end{proof}

\begin{remark}
\label{properties-remark-neurotic}
Sous les hypothèses et avec les notations du lemme \ref{properties-lemma-ample-quasi-coherent},
notons $\theta_\mathcal{F}$ l'application
qui y figure. Notons que l'isomorphisme $f^*\mathcal{O}_Y(n) \to \mathcal{L}^{\otimes n}$
du lemme \ref{properties-lemma-ample-gcd-is-one} n'est autre que
$\theta_{\mathcal{L}^{\otimes n}}$.
Considérons les applications de multiplication
$$
\widetilde{M} \otimes_{\mathcal{O}_Y} \mathcal{O}_Y(n)
\longrightarrow
\widetilde{M(n)}
$$
voir
Constructions, équation (\ref{constructions-equation-multiply-more-generally}).
Prenons-en l'image réciproque sur $X$ et considérons
$$
\xymatrix{
f^*\widetilde{M} \otimes_{\mathcal{O}_X} f^*\mathcal{O}_Y(n)
\ar[r]
\ar[d]_{\theta_\mathcal{F} \otimes \theta_{\mathcal{L}^{\otimes n}}}
&
f^*\widetilde{M(n)}
\ar[d]^{\theta_{\mathcal{F} \otimes \mathcal{L}^{\otimes n}}}
\\
\mathcal{F} \otimes \mathcal{L}^{\otimes n} \ar[r]^{\text{id}} &
\mathcal{F} \otimes \mathcal{L}^{\otimes n}
}
$$
On a utilisé ici l'identification évidente
$M(n) = \Gamma_*(X, \mathcal{L}, \mathcal{F} \otimes \mathcal{L}^{\otimes n})$.
Ce diagramme est commutatif. Démonstration omise.
\end{remark}

\noindent
On devrait pouvoir déduire le lemme suivant du
lemme \ref{properties-lemma-ample-quasi-coherent} (ou réciproquement), mais il paraît
plus simple d'en reprendre la démonstration.

\begin{lemma}
\label{properties-lemma-proj-quasi-coherent}
Soit $S$ un anneau gradué tel que $X = \text{Proj}(S)$ soit quasi-compact.
Soit $\mathcal{F}$ un $\mathcal{O}_X$-module quasi-cohérent. Posons
$M = \bigoplus_{n \in \mathbf{Z}} \Gamma(X, \mathcal{F}(n))$,
considéré comme $S$-module gradué,
voir Constructions, section \ref{constructions-section-invertible-on-proj}.
L'application
$$
\widetilde{M} \longrightarrow \mathcal{F}
$$
de Constructions, lemme
\ref{constructions-lemma-comparison-proj-quasi-coherent},
est un isomorphisme.
Si $X$ est recouvert par des ouverts standards $D_+(f)$, où $f$ est de degré $1$,
alors les applications induites
$M_n \to \Gamma(X, \mathcal{F}(n))$ sont les applications identiques.
\end{lemma}

\begin{proof}
Comme $X$ est quasi-compact, on peut trouver des éléments homogènes
$f_1, \ldots, f_n \in S$ de degrés strictement positifs tels que
$X = D_+(f_1) \cup \ldots \cup D_+(f_n)$. Soit $d$ le
plus petit commun multiple des degrés de $f_1, \ldots, f_n$.
En remplaçant chaque $f_i$ par une puissance, on peut supposer que chaque
$f_i$ est de degré $d$. Alors $\mathcal{L} = \mathcal{O}_X(d)$ est
inversible, les applications de multiplication
$\mathcal{O}_X(ad) \otimes \mathcal{O}_X(bd) \to \mathcal{O}_X((a + b)d)$
sont des isomorphismes, et chaque $f_i$ détermine une section globale $s_i$
de $\mathcal{L}$ telle que $X_{s_i} = D_+(f_i)$,
voir Constructions, lemmes \ref{constructions-lemma-where-invertible} et
\ref{constructions-lemma-principal-open}.
Ainsi $\Gamma(X, \mathcal{F}(ad)) =
\Gamma(X, \mathcal{F} \otimes \mathcal{L}^{\otimes a})$.
Rappelons que $\widetilde{M}|_{D_{+}(f_i)}$ correspond au
$S_{(f_i)}$-module $M_{(f_i)}$,
voir Constructions, lemme \ref{constructions-lemma-proj-sheaves}.
Comme le degré de $f_i$ est $d$, la classe d'isomorphisme de
$M_{(f_i)}$ ne dépend que des composantes homogènes de $M$ de
degré divisible par $d$. Plus précisément, la classe d'isomorphisme de
$M_{(f_i)}$ ne dépend que du $\Gamma_*(X, \mathcal{L})$-module gradué
$\Gamma_*(X, \mathcal{L}, \mathcal{F})$
et de l'image $s_i$ de $f_i$ dans $\Gamma_*(X, \mathcal{L})$.
Le schéma $X$ est quasi-compact par hypothèse et
séparé d'après Constructions, lemme \ref{constructions-lemma-proj-separated}.
D'après le lemme \ref{properties-lemma-invert-s-sections}, il existe un isomorphisme canonique
$$
M_{(f_i)} = \Gamma_*(X, \mathcal{L}, \mathcal{F})_{(s_i)} \to
\Gamma(X_{s_i}, \mathcal{F}).
$$
La construction de l'application de Constructions, lemme
\ref{constructions-lemma-comparison-proj-quasi-coherent},
montre alors qu'elle est un isomorphisme sur $D_+(f_i)$,
donc un isomorphisme puisque ces ouverts recouvrent $X$. Nous omettons la
démonstration de la dernière assertion.
\end{proof}




\section{Recherche d'ouverts affines convenables}
\label{properties-section-finding-affine-opens}

\noindent
Dans cette section, nous rassemblons quelques résultats sur l'existence
d'ouverts affines dans des situations plus ou moins générales.

\begin{lemma}
\label{properties-lemma-maximal-points-affine}
Soit $X$ un schéma quasi-séparé.
Soient $Z_1, \ldots, Z_n$ des composantes irréductibles deux à deux distinctes de $X$,
voir Topologie, section \ref{topology-section-irreducible-components}.
Soient $\eta_i \in Z_i$ leurs points génériques, voir
Schémas, lemme \ref{schemes-lemma-scheme-sober}.
Il existe des voisinages ouverts affines $\eta_i \in U_i$
tels que $U_i \cap U_j = \emptyset$ pour tous $i \not = j$.
En particulier, $U = U_1 \cup \ldots \cup U_n$ est un
ouvert affine contenant tous les points $\eta_1, \ldots, \eta_n$.
\end{lemma}

\begin{proof}
Soit $V_i$ un ouvert affine quelconque contenant $\eta_i$
et disjoint du fermé $Z_1 \cup \ldots \hat Z_i \ldots \cup Z_n$.
Comme $X$ est quasi-séparé, pour chaque $i$, la réunion
$W_i = \bigcup_{j, j \not = i} V_i \cap V_j$ est un ouvert
quasi-compact de $V_i$ ne contenant pas $\eta_i$.
On peut trouver des voisinages ouverts $U_i \subset V_i$
contenant $\eta_i$ et disjoints de $W_i$ d'après
Algèbre, lemme \ref{algebra-lemma-standard-open-containing-maximal-point}.
Enfin, $U$ est affine, car c'est le spectre de
l'anneau $R_1 \times \ldots \times R_n$, où $R_i = \mathcal{O}_X(U_i)$,
voir Schémas, lemme \ref{schemes-lemma-disjoint-union-affines}.
\end{proof}

\begin{remark}
\label{properties-remark-maximal-points-affine}
Le lemme \ref{properties-lemma-maximal-points-affine} ci-dessus est faux si $X$
n'est pas quasi-séparé. En voici un exemple. Prenons
$R = \mathbf{Q}[x, y_1, y_2, \ldots]/((x-i)y_i)$.
Considérons l'idéal premier minimal $\mathfrak p = (y_1, y_2, \ldots)$
de $R$. Recollons deux copies de $\Spec(R)$ le long de
l'ouvert (non quasi-compact) $\Spec(R) \setminus V(\mathfrak p)$
pour obtenir un schéma $X$ (par le recollement
décrit dans Schémas, exemple \ref{schemes-example-affine-space-zero-doubled}).
Alors les deux points maximaux de $X$ correspondant à $\mathfrak p$
ne sont contenus dans aucun ouvert affine commun. En
effet, tout ouvert de $\Spec(R)$ contenant $\mathfrak p$
contient une infinité des « droites » $x = i$, $y_j = 0$,
$j \not = i$, de paramètre $y_i$. Détails omis.
\end{remark}

\noindent
Malgré l'exemple ci-dessus, pour « la plupart » des familles finies de fermés
irréductibles, on peut appliquer le lemme \ref{properties-lemma-maximal-points-affine} ci-dessus,
du moins si $X$ est quasi-compact. Cela tient au fait que $X$ contient un
ouvert dense qui est séparé.

\begin{lemma}
\label{properties-lemma-quasi-compact-dense-open-separated}
Soit $X$ un schéma quasi-compact.
Il existe un ouvert dense $V \subset X$ qui est séparé.
\end{lemma}

\begin{proof}
Écrivons $X = \bigcup_{i = 1, \ldots, n} U_i$ comme réunion de $n$ sous-schémas
ouverts affines. Nous allons démontrer le lemme par récurrence sur $n$. Il est immédiat pour
$n = 1$. Soit $V' \subset \bigcup_{i = 1, \ldots, n - 1} U_i$ un sous-schéma
ouvert dense séparé, dont l'existence résulte de l'hypothèse de récurrence. Considérons
$$
V = V' \amalg (U_n \setminus \overline{V'}).
$$
Il est clair que $V$ est séparé et est un sous-schéma ouvert dense de $X$.
\end{proof}

\noindent
Même si $X$ est à la fois quasi-séparé et quasi-compact, il
n'existe pas nécessairement d'ouvert dense séparé et quasi-compact,
voir Exemples, lemme
\ref{examples-lemma-no-dense-separated-quasi-compact-open-in-qcqs}.
Voici une légère précision du lemme \ref{properties-lemma-maximal-points-affine} ci-dessus.

\begin{lemma}
\label{properties-lemma-point-and-maximal-points-affine}
Soit $X$ un schéma quasi-séparé. Soient $Z_1, \ldots, Z_n$ des
composantes irréductibles deux à deux distinctes de $X$. Soient $\eta_i \in Z_i$ leurs
points génériques. Soit $x \in X$ un point
quelconque. Il existe un ouvert affine $U \subset X$ contenant
$x$ et tous les $\eta_i$.
\end{lemma}

\begin{proof}
Supposons que $x \in Z_1 \cap \ldots \cap Z_r$ et que
$x \not \in Z_{r + 1}, \ldots, Z_n$. On peut alors choisir
un ouvert affine $W \subset X$ tel que
$x \in W$ et $W \cap Z_i = \emptyset$ pour
$i = r + 1, \ldots, n$. Il est clair que $\eta_i \in W$
pour $i = 1, \ldots, r$. D'après le lemme \ref{properties-lemma-maximal-points-affine},
on peut choisir des ouverts affines $U_i \subset X$ deux à deux
disjoints tels que $\eta_i \in U_i$ pour $i = r + 1, \ldots, n$.
Comme $X$ est quasi-séparé, les ouverts $W \cap U_i$
sont quasi-compacts et ne contiennent pas $\eta_i$ pour
$i = r + 1, \ldots, n$. Ainsi,
d'après Algèbre, lemme \ref{algebra-lemma-standard-open-containing-maximal-point},
on peut rétrécir $U_i$ de sorte que $W \cap U_i = \emptyset$
pour $i = r + 1, \ldots, n$. Alors la réunion
$U = W \cup \bigcup_{i = r + 1, \ldots, n} U_i$ est disjointe, donc
(d'après Schémas, lemme \ref{schemes-lemma-disjoint-union-affines})
c'est un ouvert affine convenable.
\end{proof}

\begin{lemma}
\label{properties-lemma-ample-finite-set-in-affine}
Soit $X$ un schéma. Supposons vérifiée l'une des conditions suivantes :
\begin{enumerate}
\item Le schéma $X$ est quasi-affine.
\item Le schéma $X$ est isomorphe à un sous-schéma localement
fermé d'un schéma affine.
\item Il existe un faisceau inversible ample sur $X$.
\item Le schéma $X$ est isomorphe à un sous-schéma localement fermé
de $\text{Proj}(S)$ pour un anneau gradué $S$.
\end{enumerate}
Alors, pour toute partie finie $E \subset X$, il existe un
ouvert affine $U \subset X$ tel que $E \subset U$.
\end{lemma}

\begin{proof}
D'après Propriétés, définition \ref{properties-definition-quasi-affine},
un schéma quasi-affine est un sous-schéma ouvert
quasi-compact d'un schéma affine. Tout schéma affine $\Spec(R)$ est isomorphe à
$\text{Proj}(R[X])$, où $R[X]$ est gradué en posant $\deg(X) = 1$.
D'après la proposition \ref{properties-proposition-characterize-ample},
si $X$ possède un faisceau inversible ample, alors $X$ est isomorphe à un
sous-schéma ouvert de $\text{Proj}(S)$ pour un anneau gradué $S$.
Il suffit donc de démontrer le lemme dans le cas
(4). (Nous invitons le lecteur à démontrer directement le cas (2).)

\medskip\noindent
Supposons donc que $X \subset \text{Proj}(S)$ soit un sous-schéma localement
fermé, où $S$ est un anneau gradué. Soit $T = \overline{X} \setminus X$.
Rappelons que les ouverts standards $D_{+}(f)$ forment une base de la
topologie de $\text{Proj}(S)$. Comme $E$ est fini, on peut choisir un nombre fini
d'éléments homogènes $f_i \in S_{+}$ tels que
$$
E \subset
D_{+}(f_1) \cup \ldots \cup D_{+}(f_n) \subset
\text{Proj}(S) \setminus T
$$
Supposons que $E = \{\mathfrak p_1, \ldots, \mathfrak p_m\}$,
considéré comme partie de $\text{Proj}(S)$.
Considérons l'idéal $I = (f_1, \ldots, f_n) \subset S$.
Puisque $I \not \subset \mathfrak p_j$ pour tous $j = 1, \ldots, m$,
Algèbre, lemme \ref{algebra-lemma-graded-silly},
montre qu'il existe un élément homogène $f \in I$, avec $f \not \in \mathfrak p_j$
pour tous $j = 1, \ldots, m$. Alors $E \subset D_{+}(f) \subset
D_{+}(f_1) \cup \ldots \cup D_{+}(f_n)$. Comme $D_{+}(f)$ ne
rencontre pas $T$, on voit que $X \cap D_{+}(f)$ est un sous-schéma fermé du
schéma affine $D_{+}(f)$, donc un ouvert affine de $X$, comme voulu.
\end{proof}

\begin{lemma}
\label{properties-lemma-ample-finite-set-in-principal-affine}
Soit $X$ un schéma. Soit $\mathcal{L}$ un faisceau inversible ample sur $X$.
Soient
$$
E \subset W \subset X
$$
avec $E$ fini et $W$ ouvert dans $X$. Alors il existe un entier $n > 0$
et une section $s \in \Gamma(X, \mathcal{L}^{\otimes n})$ tels que
$X_s$ soit affine et que $E \subset X_s \subset W$.
\end{lemma}

\begin{proof}
Le lecteur peut adapter la démonstration du lemme \ref{properties-lemma-ample-finite-set-in-affine}
pour établir celui-ci ; nous allons plutôt l'en déduire. D'après
le lemme \ref{properties-lemma-ample-finite-set-in-affine}, on peut choisir un ouvert
affine $U \subset W$ tel que $E \subset U$.
Considérons l'anneau gradué $S = \Gamma_*(X, \mathcal{L}) =
\bigoplus_{n \geq 0} \Gamma(X, \mathcal{L}^{\otimes n})$.
Pour chaque $x \in E$, soit $\mathfrak p_x \subset S$ l'idéal gradué des
sections qui s'annulent en $x$. Il est clair que $\mathfrak p_x$ est
un idéal premier et, puisqu'une puissance de $\mathcal{L}$ est engendrée
par ses sections globales, il est clair que $S_{+} \not \subset \mathfrak p_x$.
Soit $I \subset S$ l'idéal gradué des sections qui s'annulent en tous les
points de $X \setminus U$. Comme les ensembles $X_s$ forment une
base de la topologie, on a $I \not \subset \mathfrak p_x$ pour
tous $x \in E$.
Par le lemme d'évitement des idéaux premiers dans le cas gradué (Algèbre, lemme \ref{algebra-lemma-graded-silly}),
on peut trouver $s \in I$ homogène tel
que $s \not \in \mathfrak p_x$ pour tous $x \in E$.
Alors $E \subset X_s \subset U$ et $X_s$ est affine d'après
le lemme \ref{properties-lemma-affine-cap-s-open}.
\end{proof}

\begin{lemma}
\label{properties-lemma-quasi-affine-invertible-nonvanishing-section}
Soit $X$ un schéma quasi-affine. Soit $\mathcal{L}$ un
$\mathcal{O}_X$-module inversible. Soient $E \subset W \subset X$, avec $E$ fini
et $W$ ouvert. Alors il existe $s \in \Gamma(X, \mathcal{L})$
tel que $X_s$ soit affine et que $E \subset X_s \subset W$.
\end{lemma}

\begin{proof}
La démonstration de ce lemme présente de nombreux points communs avec celle donnée
dans Algèbre, lemme \ref{algebra-lemma-silly}.
Écrivons $E = \{x_1, \ldots, x_n\}$. Si $E = W = \emptyset$, alors $s = 0$
convient. Si $W \not = \emptyset$, on peut supposer $E \not = \emptyset$
en ajoutant un point au besoin. On peut donc supposer $n \geq 1$.
Nous allons démontrer le lemme par récurrence sur $n$.

\medskip\noindent
Cas initial : $n = 1$. En remplaçant $W$ par un voisinage ouvert affine
de $x_1$ dans $W$, on peut supposer $W$ affine. En combinant
le lemme \ref{properties-lemma-quasi-affine-O-ample} et
la proposition \ref{properties-proposition-characterize-ample},
on voit que tout
$\mathcal{O}_X$-module quasi-cohérent est engendré par ses
sections globales. Il existe
donc une section globale $s$ de $\mathcal{L}$ qui ne s'annule pas en $x_1$.
D'autre part, soit $Z \subset X$ le
sous-schéma fermé réduit induit sur $X \setminus W$.
En appliquant la propriété d'engendrement par les sections globales au faisceau quasi-cohérent d'idéaux
$\mathcal{I}$ de $Z$, on trouve une section globale $f$ de $\mathcal{I}$
qui ne s'annule pas en $x_1$. Alors $s' = fs$ est une section globale
de $\mathcal{L}$ qui ne s'annule pas en $x_1$ et telle que
$X_{s'} \subset W$. Alors $X_{s'}$ est affine d'après
le lemme \ref{properties-lemma-affine-cap-s-open}.

\medskip\noindent
Étape de récurrence pour $n > 1$. S'il existe une spécialisation
$x_i \leadsto x_j$ pour $i \not = j$, il suffit de
démontrer le lemme pour $\{x_1, \ldots, x_n\} \setminus \{x_i\}$,
et l'on conclut par récurrence. On peut donc supposer qu'il n'existe
aucune spécialisation entre les $x_i$.
D'après le lemme \ref{properties-lemma-ample-finite-set-in-affine} ou
le lemme \ref{properties-lemma-ample-finite-set-in-principal-affine},
on peut supposer $W$ affine. Par
récurrence, on peut trouver une section globale
$s$ de $\mathcal{L}$ telle que $X_s \subset W$ soit affine et contienne
$x_1, \ldots, x_{n - 1}$. Si $x_n \in X_s$, la démonstration est terminée.
Supposons que $s$ s'annule en $x_n$. D'après le cas $n = 1$, on peut
trouver une section globale $s'$ de $\mathcal{L}$ telle que
$\{x_n\} \subset X_{s'} \subset
W \setminus \overline{\{x_1, \ldots, x_{n - 1}\}}$.
On utilise ici le fait que $x_n$ n'est une spécialisation d'aucun des $x_1, \ldots, x_{n - 1}$.
Alors $s + s'$
est une section globale de $\mathcal{L}$ qui ne s'annule pas
en $x_1, \ldots, x_n$, avec $X_{s + s'} \subset W$, et
l'on conclut comme précédemment.
\end{proof}

\begin{lemma}
\label{properties-lemma-ring-affine-open-injective-local-ring}
Soient $X$ un schéma et $x \in X$ un point. Il existe un voisinage
ouvert affine $U \subset X$ de $x$ tel que l'application canonique
$\mathcal{O}_X(U) \to \mathcal{O}_{X, x}$ soit injective dans
chacun des cas suivants :
\begin{enumerate}
\item $X$ est intègre,
\item $X$ est localement noethérien,
\item $X$ est réduit et possède un nombre fini de composantes irréductibles.
\end{enumerate}
\end{lemma}

\begin{proof}
Après traduction en termes d'algèbre, on applique
Algèbre, lemme \ref{algebra-lemma-subring-of-local-ring}.
\end{proof}

\begin{lemma}
\label{properties-lemma-standard-open-two-affines}
Soient $U$, $V$ des schémas affines et soient $W \to U$ et $W \to V$
des immersions ouvertes. Pour tout $w \in W$, il existe un
voisinage ouvert affine $W' \subset W$ de $w$ tel que $W'$
ait pour image un ouvert standard de $U$ et de $V$.
\end{lemma}

\begin{proof}
Soit $X$ le schéma obtenu en recollant $U$ et $V$ le long de $W$ ;
appliquons Schémas, lemme \ref{schemes-lemma-standard-open-two-affines}.
\end{proof}










% OK, start here.
%

\chapter{Morphismes de schémas}



\phantomsection
\label{morphisms-section-phantom}


\section{Introduction}
\label{morphisms-section-introduction}

\noindent
Dans ce chapitre, nous introduisons certains types de morphismes de schémas.
Une référence fondamentale est \cite{EGA}.




















\section{Immersions fermées}
\label{morphisms-section-closed-immersions}

\noindent
Dans cette section, nous précisons certains des résultats obtenus précédemment sur les
immersions fermées de schémas. Rappelons qu'un morphisme de schémas $i : Z \to X$
est par définition une immersion fermée si (a) $i$ induit un homéomorphisme sur
une partie fermée de $X$, (b) $i^\sharp : \mathcal{O}_X \to i_*\mathcal{O}_Z$
est surjectif et (c) le noyau de $i^\sharp$ est engendré localement par des
sections, voir Schémas, Définitions \ref{schemes-definition-immersion}
et \ref{schemes-definition-closed-immersion-locally-ringed-spaces}. Il s'avère
que, puisque $Z$ et $X$ sont des schémas, il existe de nombreuses
façons de caractériser une immersion fermée.

\begin{lemma}
\label{morphisms-lemma-closed-immersion}
Soit $i : Z \to X$ un morphisme de schémas.
Les conditions suivantes sont équivalentes :
\begin{enumerate}
\item Le morphisme $i$ est une immersion fermée.
\item Pour tout ouvert affine $\Spec(R) = U \subset X$,
il existe un idéal $I \subset R$ tel que
$i^{-1}(U) = \Spec(R/I)$ en tant que schémas sur $U = \Spec(R)$.
\item Il existe un recouvrement ouvert affine $X = \bigcup_{j \in J} U_j$,
$U_j = \Spec(R_j)$ et, pour tout $j \in J$, il existe
un idéal $I_j \subset R_j$ tel que
$i^{-1}(U_j) = \Spec(R_j/I_j)$ en tant que schémas sur $U_j = \Spec(R_j)$.
\item Le morphisme $i$ induit un homéomorphisme de $Z$ sur une partie fermée
de $X$ et $i^\sharp : \mathcal{O}_X \to i_*\mathcal{O}_Z$ est surjectif.
\item Le morphisme $i$ induit un homéomorphisme de $Z$ sur une partie fermée
de $X$, l'application $i^\sharp : \mathcal{O}_X \to i_*\mathcal{O}_Z$ est surjective,
et le noyau $\Ker(i^\sharp)\subset \mathcal{O}_X$ est un faisceau quasi-cohérent
d'idéaux.
\item Le morphisme $i$ induit un homéomorphisme de $Z$ sur une partie fermée
de $X$, l'application $i^\sharp : \mathcal{O}_X \to i_*\mathcal{O}_Z$ est surjective,
et le noyau $\Ker(i^\sharp)\subset \mathcal{O}_X$ est un
faisceau d'idéaux engendré localement par des sections.
\end{enumerate}
\end{lemma}

\begin{proof}
La condition (6) est notre définition d'une immersion fermée, voir Schémas,
Définitions \ref{schemes-definition-closed-immersion-locally-ringed-spaces}
et \ref{schemes-definition-immersion}.
Ainsi (6) $\Leftrightarrow$ (1). Nous avons (1) $\Rightarrow$ (2) d'après
Schémas, Lemme \ref{schemes-lemma-closed-subspace-scheme}.
Il est immédiat que (2) $\Rightarrow$ (3).

\medskip\noindent
Supposons (3). Chacun des morphismes
$\Spec(R_j/I_j) \to \Spec(R_j)$ est
une immersion fermée, voir
Schémas, Exemple \ref{schemes-example-closed-immersion-affines}.
Ainsi $i^{-1}(U_j) \to U_j$ est un homéomorphisme sur son image
et $i^\sharp|_{U_j}$ est surjectif. Ainsi $i$ est un homéomorphisme
sur son image et $i^\sharp$ est surjectif, puisque cela peut se
vérifier localement. Nous concluons que (3) $\Rightarrow$ (4).

\medskip\noindent
L'implication (4) $\Rightarrow$ (1) est
Schémas, Lemme \ref{schemes-lemma-characterize-closed-immersions}.
L'implication (5) $\Rightarrow$ (6) est immédiate.
Et l'implication (6) $\Rightarrow$ (5) découle
de Schémas, Lemme \ref{schemes-lemma-closed-subspace-scheme}.
\end{proof}

\begin{lemma}
\label{morphisms-lemma-closed-immersion-ideals}
Soit $X$ un schéma. Soient $i : Z \to X$ et
$i' : Z' \to X$ des immersions fermées et considérons
les faisceaux d'idéaux $\mathcal{I} = \Ker(i^\sharp)$ et $\mathcal{I}' = \Ker((i')^\sharp)$
de $\mathcal{O}_X$.
\begin{enumerate}
\item Le morphisme $i : Z \to X$ se factorise en $Z \to Z' \to X$ pour
certains $a : Z \to Z'$ si et seulement si $\mathcal{I}' \subset \mathcal{I}$. Si cela se produit,
alors $a$ est une immersion fermée.
\item Nous avons $Z \cong Z'$ sur $X$ si et
seulement si $\mathcal{I} = \mathcal{I}'$.
\end{enumerate}
\end{lemma}

\begin{proof}
Cela résulte de notre discussion sur les sous-espaces
fermés dans Schémas, section \ref{schemes-section-closed-immersion}
en particulier
Schémas, Lemmes \ref{schemes-lemma-closed-immersion}
et \ref{schemes-lemma-characterize-closed-subspace}. Cela découle également
de manière simple de la caractérisation (3)
dans le lemme \ref{morphisms-lemma-closed-immersion} ci-dessus.
\end{proof}

\begin{lemma}
\label{morphisms-lemma-closed-immersion-bijection-ideals}
Soit $X$ un schéma.
Soit $\mathcal{I} \subset \mathcal{O}_X$ un faisceau d'idéaux. Les conditions suivantes
sont équivalentes :
\begin{enumerate}
\item $\mathcal{I}$ est engendré localement par des
sections comme un faisceau de $\mathcal{O}_X$-modules,
\item $\mathcal{I}$ est quasi-cohérent comme
un faisceau de $\mathcal{O}_X$-modules, et
\item existe une immersion fermée $i : Z \to X$ de schémas dont le
faisceau d'idéaux correspondant $\Ker(i^\sharp)$ est égal à $\mathcal{I}$.
\end{enumerate}
\end{lemma}

\begin{proof}
L'équivalence de (1) et (2) découle immédiatement
d'après Schémas, Lemme \ref{schemes-lemma-closed-subspace-scheme}.
Si (1) est vrai, alors il y a un sous-espace fermé
$i : Z \to X$ avec $\mathcal{I} = \Ker(i^\sharp)$
par Schémas, Définition \ref{schemes-definition-closed-subspace} et Exemple \ref{schemes-example-closed-subspace}.
Par Schémas, Lemme \ref{schemes-lemma-closed-subspace-scheme}
c'est une immersion fermée de schémas et (3) tient.
Inversement, si (3) est satisfaite, alors
(2) est vrai par Schémas, Lemme
\ref{schemes-lemma-closed-subspace-scheme} (ce qui s'applique car une immersion
fermée de schémas est a fortiori une immersion fermée
d'espaces localement annelés).
\end{proof}

\begin{lemma}
\label{morphisms-lemma-base-change-closed-immersion}
Le changement de base d'une immersion fermée est une immersion fermée.
\end{lemma}

\begin{proof}
Voir Schémas, Lemme \ref{schemes-lemma-base-change-immersion}.
\end{proof}

\begin{lemma}
\label{morphisms-lemma-composition-closed-immersion}
Une composition d'immersions fermées est une immersion fermée.
\end{lemma}

\begin{proof}
Nous avons vu cela
dans Schémas, Lemme \ref{schemes-lemma-composition-immersion}, mais voici une
autre démonstration.
En effet, cela découle de la caractérisation (3) des immersions
fermées dans le lemme \ref{morphisms-lemma-closed-immersion}. En effet, si
$I \subset R$ est un idéal, et $\overline{J} \subset R/I$ est un idéal,
alors $\overline{J} = J/I$ pour un idéal $J \subset R$ qui contient
$I$ et $(R/I)/\overline{J} = R/J$.
\end{proof}

\begin{lemma}
\label{morphisms-lemma-closed-immersion-quasi-compact}
Une immersion fermée est quasi-compacte.
\end{lemma}

\begin{proof}
Ce lemme fait double
emploi avec Schémas, Lemme \ref{schemes-lemma-closed-immersion-quasi-compact}.
\end{proof}

\begin{lemma}
\label{morphisms-lemma-closed-immersion-separated}
Une immersion fermée est séparée.
\end{lemma}

\begin{proof}
Ce lemme est un cas particulier
d'après Schémas, Lemme \ref{schemes-lemma-immersions-monomorphisms}.
\end{proof}





\section{Immersions}
\label{morphisms-section-immersions}

\noindent
Dans cette section, nous collectons quelques faits sur les immersions.

\begin{lemma}
\label{morphisms-lemma-immersion-permanence}
Soit $Z \to Y \to X$ des morphismes de schémas.
\begin{enumerate}
\item Si $Z \to X$ est une immersion, alors $Z \to Y$ est une immersion.
\item Si $Z \to X$ est une immersion quasi-compacte et $Y \to X$ est
quasi-séparé, alors $Z \to Y$ est une immersion quasi-compacte.
\item Si $Z \to X$ est une immersion fermée et que $Y \to X$ est séparé, alors
$Z \to Y$ est une immersion fermée.
\end{enumerate}
\end{lemma}

\begin{proof}
Dans chaque cas, la démonstration consiste à considérer le diagramme commutatif
$$
\xymatrix{
Z \ar[r] \ar[rd] & Y \times_X Z \ar[r] \ar[d] & Z \ar[d] \\
& Y \ar[r] & X
}
$$
où la composition des flèches horizontales supérieures est
l'identité. Démontrons (1). La première flèche horizontale
est une coupe de $Y \times_X Z \to Z$, d'où
une immersion d'après Schémas, Lemme \ref{schemes-lemma-section-immersion}.
La flèche $Y \times_X Z \to Y$ est un changement de base de $Z \to X$
donc une immersion (Schémas, Lemme \ref{schemes-lemma-base-change-immersion}). Enfin,
une composition d'immersions est une immersion
(Schémas, Lemme \ref{schemes-lemma-composition-immersion}). Cela prouve (1). Les deux autres résultats
se prouvent exactement de la même manière.
\end{proof}

\begin{lemma}
\label{morphisms-lemma-factor-quasi-compact-immersion}
Soit $h : Z \to X$ une immersion. Si
$h$ est quasi-compact, alors on
peut factoriser $h = i \circ j$ avec $j : Z \to \overline{Z}$
une immersion ouverte et $i : \overline{Z} \to X$ une immersion fermée.
\end{lemma}

\begin{proof}
Notons que $h$ est quasi-compact
et quasi-séparé (voir Schémas, Lemme \ref{schemes-lemma-immersions-monomorphisms}).
Ainsi $h_*\mathcal{O}_Z$ est un faisceau quasi-cohérent
de $\mathcal{O}_X$-modules d'après Schémas, Lemme
\ref{schemes-lemma-push-forward-quasi-coherent}.
Cela implique que $\mathcal{I} = \Ker(\mathcal{O}_X \to h_*\mathcal{O}_Z)$ est un
faisceau quasi-cohérent d'idéaux,
voir Schémas, section \ref{schemes-section-quasi-coherent}.
Soit $\overline{Z} \subset X$ le sous-schéma fermé correspondant
à $\mathcal{I}$, voir Lemme \ref{morphisms-lemma-closed-immersion-bijection-ideals}. Par Schémas,
Lemme \ref{schemes-lemma-characterize-closed-subspace} le morphisme $h$
se factorise en
$h = i \circ j$ où $i : \overline{Z} \to X$ est le morphisme d'inclusion. Pour voir que $j$
est une immersion ouverte, choisissez un sous-schéma ouvert
$U \subset X$ tel que $h$ induit une immersion fermée de
$Z$ dans $U$. Il est alors clair que $\mathcal{I}|_U$
est le faisceau d'idéaux correspondant à l'immersion fermée
$Z \to U$. Nous voyons donc que $Z = \overline{Z} \cap U$.
\end{proof}

\begin{lemma}
\label{morphisms-lemma-factor-reduced-immersion}
Soit $h : Z \to X$ une immersion. Si
$Z$ est réduit, alors on
peut factoriser $h = i \circ j$ avec $j : Z \to \overline{Z}$ une
immersion ouverte et $i : \overline{Z} \to X$ une immersion fermée.
\end{lemma}

\begin{proof}
Soit $\overline{Z} \subset X$ l'adhérence de $h(Z)$ munie de la structure de sous-schéma
fermé induit réduit, voir
Schémas, Définition \ref{schemes-definition-reduced-induced-scheme}. Par Schémas, Lemme \ref{schemes-lemma-map-into-reduction}
le morphisme $h$ se factorise en $h = i \circ j$
avec $i : \overline{Z} \to X$
le morphisme d'inclusion et $j : Z \to \overline{Z}$. De la définition d'une
immersion on voit qu'il existe un sous-schéma ouvert
$U \subset X$ tel que $h$ se factorise par une
immersion fermée dans $U$. Ainsi $\overline{Z} \cap U$
et $h(Z)$ sont des sous-schémas fermés réduits
de $U$ avec le même ensemble fermé sous-jacent.
Ainsi par l'unicité dans Schémas, Lemme \ref{schemes-lemma-reduced-closed-subscheme}
on voit que $h(Z) \cong \overline{Z} \cap U$. Ainsi $j$ induit
un isomorphisme de $Z$ avec $\overline{Z} \cap U$. En d'autres termes,
$j$ est une immersion ouverte.
\end{proof}



\begin{example}
\label{morphisms-example-thibaut}
Voici un exemple d'immersion qui n'est pas une composition d'une
immersion ouverte suivie d'une immersion
fermée. Soit
$k$ un corps. Soit $X = \Spec(k[x_1, x_2, x_3, \ldots])$.
Soit $U = \bigcup_{n = 1}^{\infty} D(x_n)$. Alors $U \to X$ est
une immersion ouverte. Considérons
les idéaux
$$
I_n =
(x_1^n, x_2^n, \ldots, x_{n - 1}^n, x_n - 1, x_{n + 1}, x_{n + 2}, \ldots)
\subset
k[x_1, x_2, x_3, \ldots][1/x_n].
$$
Notez que $I_n k[x_1, x_2, x_3, \ldots][1/x_nx_m] = (1)$ pour tout $m \not = n$. D'où
les idéaux quasi-cohérents $\widetilde I_n$
sur $D(x_n)$ s'accordent sur $D(x_nx_m)$, à
savoir $\widetilde I_n|_{D(x_nx_m)} = \mathcal{O}_{D(x_n x_m)}$ si $n \not = m$. Ces idéaux se
collent donc à un faisceau quasi-cohérent d'idéaux $\mathcal{I} \subset \mathcal{O}_U$.
Soit $Z \subset U$ le sous-schéma
fermé correspondant à $\mathcal{I}$. Ainsi
$Z \to X$ est une immersion.

\medskip\noindent
Nous affirmons que nous ne pouvons pas
factoriser $Z \to X$ en $Z \to \overline{Z} \to X$, où $\overline{Z} \to X$
est une immersion fermée et $Z \to \overline{Z}$ une immersion ouverte. En effet, $\overline{Z}$
devrait être défini par un idéal $I \subset k[x_1, x_2, x_3, \ldots]$ tel que
$I_n = I k[x_1, x_2, x_3, \ldots][1/x_n]$. Mais le seul élément $f \in k[x_1, x_2, x_3, \ldots]$
qui se retrouve dans tous les $I_n$ est $0$
! Par conséquent, $I$ n'existe pas.
\end{example}

\begin{lemma}
\label{morphisms-lemma-check-immersion}
Soit $f : Y \to X$ un morphisme de schémas. Si pour tout $y \in Y$
il existe un sous-schéma ouvert $f(y) \in U \subset X$ tel que
$f|_{f^{-1}(U)} : f^{-1}(U) \to U$ est une immersion, alors $f$ est une
immersion.
\end{lemma}

\begin{proof}
Cette affirmation découle facilement de la discussion sur les sous-schémas
fermés à la fin de Schémas, section \ref{schemes-section-immersions} mais nous
donnerons également une démonstration
détaillée. Soit $Z \subset X$ l'adhérence de $f(Y)$. Puisque la formation
de l'adhérence commute à la restriction aux ouverts, l'hypothèse montre que
$f(Y) \subset Z$ est ouvert. Par conséquent,
$Z' = Z \setminus f(Y)$ est fermé. Ainsi $X' = X \setminus Z'$ est un sous-schéma
ouvert de $X$ et $f$ se factorise en
$f : Y \to X'$ suivi de l'inclusion. Si $y \in Y$ et $U \subset X$
sont comme dans l'énoncé du lemme, alors $U' = X' \cap U$
est un voisinage ouvert de $f'(y)$ tel que $(f')^{-1}(U') \to U'$
est une immersion (Lemme \ref{morphisms-lemma-immersion-permanence}) à image fermée.
Il s'agit donc d'une immersion fermée,
voir Schémas, Lemme \ref{schemes-lemma-immersion-when-closed}. Puisqu'être une
immersion fermée est locale sur la cible (par exemple
par le lemme \ref{morphisms-lemma-closed-immersion}), nous concluons que
$f'$ est une immersion fermée comme voulu.
\end{proof}








\section{Immersions fermées et faisceaux quasi-cohérents}
\label{morphisms-section-closed-immersions-quasi-coherent}

\noindent
Le lemme suivant fait finalement pour les faisceaux quasi-cohérents
sur les schémas ce que Modules, Lemme \ref{modules-lemma-i-star-exact} fait pour les faisceaux
abéliens. Voir également
la discussion dans Modules, section \ref{modules-section-closed-immersion}.

\begin{lemma}
\label{morphisms-lemma-i-star-equivalence}
Soit $i : Z \to X$ une immersion fermée de schémas.
Soit $\mathcal{I} \subset \mathcal{O}_X$ le faisceau quasi-cohérent d'idéaux définissant
$Z$. Le foncteur
$$
i_* :
\QCoh(\mathcal{O}_Z)
\longrightarrow
\QCoh(\mathcal{O}_X)
$$
est exact, pleinement fidèle, avec image essentielle des $\mathcal{O}_X$-modules
quasi-cohérents $\mathcal{G}$ tels que $\mathcal{I}\mathcal{G} = 0$.
\end{lemma}

\begin{proof}
Une immersion fermée est quasi-compacte et
séparée, voir Lemmes \ref{morphisms-lemma-closed-immersion-quasi-compact} et \ref{morphisms-lemma-closed-immersion-separated}.
D'où Schémas, Lemme \ref{schemes-lemma-push-forward-quasi-coherent}
s'applique et l'image directe d'un faisceau quasi-cohérent
sur $Z$ est bien un
faisceau quasi-cohérent sur $X$.

\medskip\noindent
Par Modules, Lemme \ref{modules-lemma-i-star-equivalence} le foncteur $i_*$
est pleinement fidèle.

\medskip\noindent
Passons maintenant à la description de l'image essentielle
du foncteur $i_*$. On a $\mathcal{I}(i_*\mathcal{F}) = 0$
pour tout module $\mathcal{O}_Z$ quasi-cohérent, par exemple
par Modules, Lemme \ref{modules-lemma-i-star-equivalence}. Supposons
ensuite que $\mathcal{G}$ soit
un module $\mathcal{O}_X$ quasi-cohérent tel que
$\mathcal{I}\mathcal{G} = 0$. Il suffit de montrer que l'application canonique
$$
\mathcal{G} \longrightarrow i_* i^*\mathcal{G}
$$
est un isomorphisme\footnote{Cela a été prouvé dans une situation
plus générale dans la démonstration de Modules, Lemme \ref{modules-lemma-i-star-equivalence}.}. Dans
le cas de schémas et modules quasi-cohérents, en travaillant localement sur
les affines de $X$ et en utilisant Lemme \ref{morphisms-lemma-closed-immersion}
et Schémas, Lemme \ref{schemes-lemma-widetilde-pullback} il suffit de prouver l'énoncé
algébrique suivant : Étant donné un anneau $R$, un idéal $I$
et un $R$-module $N$ tel que $IN = 0$ l'application canonique
$$
N \longrightarrow N \otimes_R R/I,\quad
n \longmapsto n \otimes 1
$$
est un isomorphisme de $R$-modules. La démonstration de ce fait algébrique simple est omise.
\end{proof}

\noindent
Soit $i : Z \to X$ une immersion fermée. En raison du lemme ci-dessus, nous désignons
souvent, par abus de notation, $\mathcal{F}$ le faisceau $i_*\mathcal{F}$ sur $X$.

\begin{lemma}
\label{morphisms-lemma-largest-quasi-coherent-subsheaf}
Soit $X$ un schéma. Soit $\mathcal{F}$ un module
$\mathcal{O}_X$ quasi-cohérent. Soit $\mathcal{G} \subset \mathcal{F}$ un $\mathcal{O}_X$-sous-module.
Il existe un unique $\mathcal{O}_X$-sous-module
$\mathcal{G}' \subset \mathcal{G}$ quasi-cohérent avec la propriété
suivante : Pour tout $\mathcal{O}_X$-module $\mathcal{H}$ quasi-cohérent
l'application
$$
\Hom_{\mathcal{O}_X}(\mathcal{H}, \mathcal{G}')
\longrightarrow
\Hom_{\mathcal{O}_X}(\mathcal{H}, \mathcal{G})
$$
est bijective. En particulier, $\mathcal{G}'$ est le plus grand $\mathcal{O}_X$-sous-module
quasi-cohérent de $\mathcal{F}$ contenu dans $\mathcal{G}$.
\end{lemma}

\begin{proof}
Soit $\mathcal{G}_a$, $a \in A$ l'ensemble des $\mathcal{O}_X$-sous-modules
quasi-cohérents contenus dans $\mathcal{G}$.
Alors l'image $\mathcal{G}'$ de
$$
\bigoplus\nolimits_{a \in A} \mathcal{G}_a \longrightarrow \mathcal{F}
$$
est quasi-cohérente comme l'image d'une application de faisceaux quasi-cohérents
sur $X$ est quasi-cohérente et puisqu'une somme directe
de faisceaux quasi-cohérents
est quasi-cohérente, voir Schémas, section \ref{schemes-section-quasi-coherent}.
Le module $\mathcal{G}'$ est contenu dans $\mathcal{G}$. Il s'agit donc
du plus grand module $\mathcal{O}_X$ quasi-cohérent contenu dans $\mathcal{G}$.

\medskip\noindent
Pour prouver la formule, soit $\mathcal{H}$ un module
$\mathcal{O}_X$ quasi-cohérent et soit $\alpha : \mathcal{H} \to \mathcal{G}$ une application
de module $\mathcal{O}_X$. L'image de la composition $\mathcal{H} \to \mathcal{G} \to \mathcal{F}$
est quasi-cohérente comme l'image d'une application
de faisceaux quasi-cohérents. Il est donc contenu dans $\mathcal{G}'$.
Ainsi, $\alpha$ se factorise par $\mathcal{G}'$ comme
souhaité.
\end{proof}

\begin{lemma}
\label{morphisms-lemma-i-upper-shriek}
Soit $i : Z \to X$ une immersion fermée de schémas.
Il existe un foncteur\footnote{Il s'agit probablement
d'une notation non standard.}
$i^! : \QCoh(\mathcal{O}_X) \to \QCoh(\mathcal{O}_Z)$ qui est un adjoint droit de
$i_*$. (Comparer Modules, Lemme \ref{modules-lemma-i-star-right-adjoint}.)
\end{lemma}

\begin{proof}
Étant donné le $\mathcal{O}_X$-module $\mathcal{G}$ quasi-cohérent,
nous considérons le sous-faisceau $\mathcal{H}_Z(\mathcal{G})$ de $\mathcal{G}$ de sections
locales annihilées par
$\mathcal{I}$. Par le Lemme \ref{morphisms-lemma-largest-quasi-coherent-subsheaf} il existe
un $\mathcal{O}_X$-sous-module quasi-cohérent canonique le plus grand
$\mathcal{H}_Z(\mathcal{G})'$. Par construction, nous avons
$$
\Hom_{\mathcal{O}_X}(i_*\mathcal{F}, \mathcal{H}_Z(\mathcal{G})')
=
\Hom_{\mathcal{O}_X}(i_*\mathcal{F}, \mathcal{G})
$$
pour tout $\mathcal{O}_Z$-module quasi-cohérent $\mathcal{F}$.
Nous pouvons donc définir $i^!\mathcal{G} = i^*(\mathcal{H}_Z(\mathcal{G})')$. Détails
omis.
\end{proof}

\noindent
À l'aide des $1$-to-$1$ correspondant entre faisceaux
quasi-cohérents d'idéaux et sous-schémas
fermés (voir Lemme \ref{morphisms-lemma-closed-immersion-bijection-ideals}) on peut définir
des intersections schématiques et des unions de sous-schémas
fermés.

\begin{definition}
\label{morphisms-definition-scheme-theoretic-intersection-union}
Soit $X$ un schéma. Soient $Z, Y \subset X$ des sous-schémas fermés
correspondant à des faisceaux d'idéaux quasi-cohérents
$\mathcal{I}, \mathcal{J} \subset \mathcal{O}_X$. L'{\it intersection
schématique} de $Z$ et $Y$ est le
sous-schéma fermé de $X$ découpé par $\mathcal{I} + \mathcal{J}$. La {\it
réunion schématique} de $Z$
et $Y$ est le sous-schéma fermé de $X$
découpé par $\mathcal{I} \cap \mathcal{J}$.
\end{definition}

\begin{lemma}
\label{morphisms-lemma-scheme-theoretic-intersection}
Soit $X$ un schéma. Soient $Z, Y \subset X$ des sous-schémas fermés.
Soit $Z \cap Y$ l'intersection schématique de $Z$ et $Y$.
Alors $Z \cap Y \to Z$ et $Z \cap Y \to Y$ sont des immersions fermées
et
$$
\xymatrix{
Z \cap Y \ar[r] \ar[d] & Z \ar[d] \\
Y \ar[r] & X
}
$$
est un diagramme cartésien de schémas, c'est-à-dire $Z \cap Y = Z \times_X Y$.
\end{lemma}

\begin{proof}
Les morphismes $Z \cap Y \to Z$ et $Z \cap Y \to Y$ sont des immersions fermées
par le lemme \ref{morphisms-lemma-closed-immersion-ideals}. Soit $U = \Spec(A)$
un ouvert affine de $X$ et soit $Z \cap U$ et $Y \cap U$
correspondre aux idéaux $I \subset A$ et $J \subset A$.
Alors $Z \cap Y \cap U$ correspond à $I + J \subset A$. Puisque
$A/I \otimes_A A/J = A/(I + J)$ nous voyons que le schéma est cartésien
par notre description des produits fibrés
de schémas dans Schémas, section \ref{schemes-section-fibre-products}.
\end{proof}

\begin{lemma}
\label{morphisms-lemma-scheme-theoretic-union}
Soit $S$ un schéma. Soient $X, Y \subset S$ des sous-schémas fermés.
Soit $X \cup Y$ la réunion schématique de $X$ et $Y$.
Soit $X \cap Y$ l'intersection schématique de $X$ et $Y$.
Alors $X \to X \cup Y$ et $Y \to X \cup Y$ sont des immersions fermées, il existe une
suite exacte courte
$$
0 \to \mathcal{O}_{X \cup Y} \to \mathcal{O}_X \times \mathcal{O}_Y
\to \mathcal{O}_{X \cap Y} \to 0
$$
de $\mathcal{O}_S$-modules, et le diagramme
$$
\xymatrix{
X \cap Y \ar[r] \ar[d] & X \ar[d] \\
Y \ar[r] & X \cup Y
}
$$
est cocartésien dans la catégorie des schémas,
c'est-à-dire $X \cup Y = X \amalg_{X \cap Y} Y$.
\end{lemma}

\begin{proof}
Les morphismes $X \to X \cup Y$ et $Y \to X \cup Y$ sont des immersions fermées par
le lemme \ref{morphisms-lemma-closed-immersion-ideals}. Dans la suite exacte courte nous utilisons l'équivalence
du lemme \ref{morphisms-lemma-i-star-equivalence} pour penser les modules quasi-cohérents
sur sous-schémas fermés de $S$ comme des modules quasi-cohérents
sur $S$. Pour la première application de la séquence
on utilise les applications canoniques $\mathcal{O}_{X \cup Y} \to \mathcal{O}_X$
et $\mathcal{O}_{X \cup Y} \to \mathcal{O}_Y$ et pour la deuxième
application on utilise l'application canonique
$\mathcal{O}_X \to \mathcal{O}_{X \cap Y}$ et le négatif de l'application
canonique $\mathcal{O}_Y \to \mathcal{O}_{X \cap Y}$. Ensuite,
pour vérifier l'exactitude, nous pouvons travailler
localement sur les affines.
Soit $U = \Spec(A)$ un ouvert affine de $S$ et soit $X \cap U$ et
$Y \cap U$ correspondre aux idéaux $I \subset A$ et $J \subset A$.
Alors $(X \cup Y) \cap U$ correspond à $I \cap J \subset A$
et $X \cap Y \cap U$ correspond à $I + J \subset A$. Ainsi l'exactitude
découle de l'exactitude de
$$
0 \to A/I \cap J \to A/I \times A/J \to A/(I + J) \to 0
$$
Pour montrer que le diagramme est cocartésien, supposons qu'on nous donne
un schéma $T$ et des morphismes de schémas $f : X \to T$, $g : Y \to T$
qui s'accordent comme morphismes $X \cap Y \to T$. Il s'agit de montrer
qu'il existe un morphisme unique $h : X \cup Y \to T$ en accord
avec $f$ et $g$. Pour construire $h$ on peut
travailler localement sur les affines de $X \cup Y$, voir
Schémas, section \ref{schemes-section-glueing-schemes}. Si $s \in X$, $s \not \in Y$, alors
$X \to X \cup Y$ est un isomorphisme dans un voisinage de $s$
et il est clair comment construire $h$. De même pour
$s \in Y$, $s \not \in X$. Pour $s \in X \cap Y$ nous pouvons
choisir un $V = \Spec(B) \subset T$ ouvert affine contenant $f(s) = g(s)$.
On peut alors choisir un $U = \Spec(A) \subset S$ ouvert affine
contenant $s$ de telle sorte que $f(X \cap U)$ et $g(Y \cap U)$ soient
contenus dans $V$. Les morphismes $f|_{X \cap U}$ et $g|_{Y \cap V}$
dans $V$ correspondent à des homomorphismes d'anneaux
$$
B \to A/I
\quad\text{et}\quad
B \to A/J
$$
dont les composés vers $A/(I + J)$ coïncident. Par la suite exacte courte affichée
ci-dessus, il y a un relèvement unique de ces homomorphismes d'anneaux vers
un homomorphisme d'anneaux $B \to A/I \cap J$ comme souhaité.
\end{proof}









\section{Supports des modules}
\label{morphisms-section-support}

\noindent
Dans cette section nous collectons quelques résultats élémentaires
sur les supports de modules quasi-cohérents
sur des schémas. Rappelons que le support d'un faisceau
de modules a été défini dans Modules, section \ref{modules-section-support}.
En revanche, le support d'un module a été défini
dans Algèbre, section \ref{algebra-section-support}. Ceux-ci
correspondent.

\begin{lemma}
\label{morphisms-lemma-support-affine-open}
Soit $X$ un schéma. Soit $\mathcal{F}$ un faisceau quasi-cohérent sur
$X$. Soit $\Spec(A) = U \subset X$ un ouvert affine et définissons
$M = \Gamma(U, \mathcal{F})$. Soit $x \in U$,
et soit $\mathfrak p \subset A$ l'idéal premier correspondant. Les conditions suivantes
sont équivalentes
\begin{enumerate}
\item $\mathfrak p$ est dans le support de $M$, et
\item $x$ est dans le support de $\mathcal{F}$.
\end{enumerate}
\end{lemma}

\begin{proof}
Cela résulte de l'égalité $\mathcal{F}_x = M_{\mathfrak p}$, voir Schémas,
Lemme \ref{schemes-lemma-spec-sheaves} et les
définitions.
\end{proof}

\begin{lemma}
\label{morphisms-lemma-support-closed-specialization}
Soit $X$ un schéma.
Soit $\mathcal{F}$ un faisceau quasi-cohérent sur $X$.
Le support de $\mathcal{F}$ est clôturé sous spécialisation.
\end{lemma}

\begin{proof}
Si $x' \leadsto x$ est une spécialisation et $\mathcal{F}_x = 0$
alors $\mathcal{F}_{x'}$ vaut zéro, car $\mathcal{F}_{x'}$ est une localisation
du module $\mathcal{F}_x$. Le complément
de $\text{Supp}(\mathcal{F})$ est donc fermé sous génération.
\end{proof}

\noindent
Pour les modules de type fini quasi-cohérents le support est fermé,
vérifiable sur fibres, et commute avec changement de base.

\begin{lemma}
\label{morphisms-lemma-support-finite-type}
Soit $\mathcal{F}$ un module de type fini quasi-cohérent sur un schéma
$X$. Ensuite,
\begin{enumerate}
\item Le support de $\mathcal{F}$ est fermé.
\item Pour $x \in X$ nous avons
$$
x \in \text{Supp}(\mathcal{F})
\Leftrightarrow
\mathcal{F}_x \not = 0
\Leftrightarrow
\mathcal{F}_x \otimes_{\mathcal{O}_{X, x}} \kappa(x) \not = 0.
$$
\item Pour tout morphisme de schémas $f : Y \to X$
le pullback $f^*\mathcal{F}$ est également
de type fini et nous avons $\text{Supp}(f^*\mathcal{F}) = f^{-1}(\text{Supp}(\mathcal{F}))$.
\end{enumerate}
\end{lemma}

\begin{proof}
La partie (1) est une reformulation
des modules, Lemme \ref{modules-lemma-support-finite-type-closed}. Vous pouvez
également combiner
Lemme \ref{morphisms-lemma-support-affine-open}, Propriétés,
Lemme \ref{properties-lemma-finite-type-module} et Algèbre, Lemme \ref{algebra-lemma-support-closed}
pour
voir cela. La première équivalence
dans (2) est la définition du support, et la deuxième
équivalence découle du lemme de Nakayama,
voir Algèbre,
Lemme \ref{algebra-lemma-NAK}. Soit $f : Y \to X$
un morphisme de schémas. Notez que
$f^*\mathcal{F}$ est de type fini par
Modules, Lemme \ref{modules-lemma-pullback-finite-type}. Pour l'assertion finale,
soit $y \in Y$ avec l'image $x \in X$. Rappelons
que
$$
(f^*\mathcal{F})_y =
\mathcal{F}_x \otimes_{\mathcal{O}_{X, x}} \mathcal{O}_{Y, y},
$$
voir
Faisceaux, Lemme \ref{sheaves-lemma-stalk-pullback-modules}. Donc $(f^*\mathcal{F})_y \otimes \kappa(y)$
est différent de zéro si et seulement
si $\mathcal{F}_x \otimes \kappa(x)$ est différent de zéro. Par
(2) cela implique $x \in \text{Supp}(\mathcal{F})$ si et seulement
si $y \in \text{Supp}(f^*\mathcal{F})$, qui est le contenu de l'assertion
(3).
\end{proof}

\begin{lemma}
\label{morphisms-lemma-scheme-theoretic-support}
Soit $\mathcal{F}$ un module de type fini quasi-cohérent
sur un schéma $X$. Il existe un plus petit
sous-schéma fermé $i : Z \to X$ tel qu'il existe
un $\mathcal{O}_Z$-module $\mathcal{G}$ quasi-cohérent
avec $i_*\mathcal{G} \cong \mathcal{F}$. De plus :
\begin{enumerate}
\item Si $\Spec(A) \subset X$ est un ouvert
affine, et $\mathcal{F}|_{\Spec(A)} = \widetilde{M}$
alors $Z \cap \Spec(A) = \Spec(A/I)$ où $I = \text{Ann}_A(M)$.
\item Le faisceau quasi-cohérent $\mathcal{G}$ est unique à isomorphisme près
unique.
\item Le faisceau quasi-cohérent $\mathcal{G}$ est de type fini.
\item Le support de $\mathcal{G}$ et de $\mathcal{F}$ est $Z$.
\end{enumerate}
\end{lemma}

\begin{proof}
Supposons que $i' : Z' \to X$ soit un sous-schéma fermé qui satisfait
à la description sur ouverts affines
du lemme. Puis par le lemme \ref{morphisms-lemma-i-star-equivalence}
on voit que $\mathcal{F} \cong i'_*\mathcal{G}'$ pour un faisceau unique
quasi-cohérent $\mathcal{G}'$ sur $Z'$. De plus, il
est clair que $Z'$ est le plus petit sous-schéma fermé
avec cette propriété
(par le même lemme). Enfin, en utilisant Propriétés,
Lemme
\ref{properties-lemma-finite-type-module} et Algèbre, Lemme \ref{algebra-lemma-finite-over-subring}
Il s'ensuit que $\mathcal{G}'$ est de type fini.
Nous avons $\text{Supp}(\mathcal{G}') = Z$ par
Algèbre, Lemme \ref{algebra-lemma-support-closed}. Par conséquent, pour
prouver le lemme, il suffit de montrer que la caractérisation
en (1) définit effectivement un sous-schéma fermé. Et,
pour ce faire, il suffit de prouver que la règle donnée
produit un faisceau quasi-cohérent d'idéaux,
voir Lemme \ref{morphisms-lemma-closed-immersion-bijection-ideals}. Cela se résume au fait
algébrique suivant : si $A$ est un anneau, $f \in A$ et $M$
est un module $A$
fini, alors $\text{Ann}_A(M)_f = \text{Ann}_{A_f}(M_f)$. Nous omettons la
démonstration.
\end{proof}

\begin{definition}
\label{morphisms-definition-scheme-theoretic-support}
Soit $X$ un schéma. Soit $\mathcal{F}$ un $\mathcal{O}_X$-module
quasi-cohérent de type fini. Le {\it support
schématique de $\mathcal{F}$} est le sous-schéma fermé $Z \subset X$
construit dans le lemme \ref{morphisms-lemma-scheme-theoretic-support}.
\end{definition}

\noindent
Dans cette situation on pense souvent à $\mathcal{F}$ comme un faisceau
quasi-cohérent de type fini sur $Z$ (via l'équivalence
des catégories du lemme \ref{morphisms-lemma-i-star-equivalence}).













\section{Image schématique}
\label{morphisms-section-scheme-theoretic-image}

\noindent
Attention : Certains éléments de cette section sont ultra-généraux et se comportent
différemment de ce à quoi vous pourriez vous attendre.

\begin{lemma}
\label{morphisms-lemma-scheme-theoretic-image}
Soit $f : X \to Y$ un morphisme de schémas. Il existe un sous-schéma fermé
$Z \subset Y$ tel que $f$ se factorise par $Z$ et tel que
pour tout autre sous-schéma fermé $Z' \subset Y$ tel que $f$
se factorise par $Z'$ on a $Z \subset Z'$.
\end{lemma}

\begin{proof}
Soit $\mathcal{I} = \Ker(\mathcal{O}_Y \to f_*\mathcal{O}_X)$. Si $\mathcal{I}$ est quasi-cohérent
alors on prend simplement $Z$ comme étant le
sous-schéma fermé déterminé par $\mathcal{I}$,
voir Lemme \ref{morphisms-lemma-closed-immersion-bijection-ideals}. Cela fonctionne par le lemme
\ref{morphisms-lemma-closed-immersion-ideals}. De manière générale,
le même lemme nous oblige à montrer qu'il existe
un plus grand faisceau quasi-cohérent d'idéaux $\mathcal{I}'$
contenu
dans $\mathcal{I}$. Cela résulte du lemme \ref{morphisms-lemma-largest-quasi-coherent-subsheaf}.
\end{proof}

\begin{definition}
\label{morphisms-definition-scheme-theoretic-image}
Soit $f : X \to Y$ un morphisme de schémas. L'{\it image schématique}
de $f$ est le plus petit sous-schéma fermé $Z \subset Y$ à travers lequel
$f$ se factorise par, voir Lemme \ref{morphisms-lemma-scheme-theoretic-image} ci-dessus.
\end{definition}

\noindent
Pour un morphisme $f : X \to Y$ de schémas à image schématique $Z$ on note
souvent $f : X \to Z$ la factorisation de $f$ à travers
son image schématique. Si le morphisme $f$ n'est pas
quasi-compact, alors (en général)
\begin{enumerate}
\item l'inclusion théorique des ensembles $\overline{f(X)} \subset Z$ n'est pas une
égalité, c'est-à-dire que $f(X) \subset Z$ n'est pas un sous-ensemble dense, et
\item la construction de l'image schématique ne commute pas avec restriction
aux sous-schémas ouverts à $Y$.
\end{enumerate}
Dans Exemples, section \ref{examples-section-scheme-theoretic-image}, le lecteur trouve un exemple
pour les deux phénomènes. Ces phénomènes
peuvent survenir même en immersion, voir Exemples,
section \ref{examples-section-strange-immersion}. Cependant, le lemme suivant
montre que les deux désastres sont évités lorsque le
morphisme est quasi-compact.

\begin{lemma}
\label{morphisms-lemma-quasi-compact-scheme-theoretic-image}
Soit $f : X \to Y$ un morphisme de schémas. Soit
$Z \subset Y$ l'image schématique de $f$. Si $f$
est quasi-compact alors
\begin{enumerate}
\item le faisceau
d'idéaux $\mathcal{I} = \Ker(\mathcal{O}_Y \to f_*\mathcal{O}_X)$ est
quasi-cohérent,
\item l'image schématique $Z$ est le sous-schéma fermé
déterminé par $\mathcal{I}$,
\item pour tout ouvert $U \subset Y$ l'image schématique
de $f|_{f^{-1}(U)} : f^{-1}(U) \to U$ est égale à $Z \cap U$, et
\item l'image $f(X) \subset Z$ est un sous-ensemble dense de $Z$,
autrement dit le morphisme $X \to Z$ est
dominant (voir Définition \ref{morphisms-definition-dominant}).
\end{enumerate}
\end{lemma}

\begin{proof}
La partie (4) découle de la partie (3). Pour montrer (3) il suffit
de prouver (1) puisque la formation de $\mathcal{I}$ commute avec restriction aux
sous-schémas ouverts de $Y$. Et si (1) est vérifié alors dans
la démonstration du lemme \ref{morphisms-lemma-scheme-theoretic-image} nous
avons montré (2). Il suffit donc de prouver que $\mathcal{I}$ est quasi-cohérent. Puisque
la propriété d'être quasi-cohérent est locale, nous
pouvons supposer que $Y$ est affine. Comme $f$
est quasi-compact, on peut trouver un recouvrement
ouvert fini affine $X = \bigcup_{i = 1, \ldots, n} U_i$. Notons $f'$ la composition
$$
X' = \coprod U_i \longrightarrow X \longrightarrow Y.
$$
Alors $f_*\mathcal{O}_X$ est un sous-faisceau de $f'_*\mathcal{O}_{X'}$,
et donc $\mathcal{I} = \Ker(\mathcal{O}_Y \to f'_*\mathcal{O}_{X'})$. Par Schémas, Lemme
\ref{schemes-lemma-push-forward-quasi-coherent} le faisceau $f'_*\mathcal{O}_{X'}$ est quasi-cohérent
sur $Y$. On obtient donc le résultat.
\end{proof}

\begin{example}
\label{morphisms-example-scheme-theoretic-image}
Si $A \to B$ est un homomorphisme d'anneaux à noyau $I$,
alors l'image schématique de $\Spec(B) \to \Spec(A)$ est
le sous-schéma fermé $\Spec(A/I)$ de
$\Spec(A)$. Cela résulte du lemme \ref{morphisms-lemma-quasi-compact-scheme-theoretic-image}.
\end{example}

\noindent
Si le morphisme est quasi-compact, alors l'image schématique ajoute uniquement
des points qui sont des spécialisations de points dans l'image.

\begin{lemma}
\label{morphisms-lemma-reach-points-scheme-theoretic-image}
Soit $f : X \to Y$ un morphisme quasi-compact. Soit $Z$
l'image schématique de $f$. Soit $z \in Z$\footnote{Par
Lemme
\ref{morphisms-lemma-quasi-compact-scheme-theoretic-image}, au sens ensembliste $Z$ coïncide avec l'adhérence de
$f(X)$ dans $Y$.}. Il existe un anneau
de valorisation $A$ avec corps des
fractions $K$ et un diagramme commutatif
$$
\xymatrix{
\Spec(K) \ar[rr] \ar[d] & & X \ar[d] \ar[ld] \\
\Spec(A) \ar[r] & Z \ar[r] & Y
}
$$
tel que le point fermé de $\Spec(A)$ s'envoie sur $z$. En particulier tout
point de $Z$ est la spécialisation d'un point de $f(X)$.
\end{lemma}

\begin{proof}
Soit $z \in \Spec(R) = V \subset Y$ un voisinage ouvert affine de
$z$. Par le Lemme
\ref{morphisms-lemma-quasi-compact-scheme-theoretic-image} l'intersection $Z \cap V$ est l'image
schématique de $f^{-1}(V) \to V$. Nous pouvons donc
remplacer $Y$ par $V$ et supposer
que $Y = \Spec(R)$ est affine.
Dans ce cas, $X$ est quasi-compact comme
$f$ est quasi-compact.
Disons que $X = U_1 \cup \ldots \cup U_n$ est un recouvrement ouvert
affine fini. Écrivez $U_i = \Spec(A_i)$. Soit
$I = \Ker(R \to A_1 \times \ldots \times A_n)$. Par le Lemme \ref{morphisms-lemma-quasi-compact-scheme-theoretic-image} encore on voit
que $Z$ correspond au sous-schéma fermé
$\Spec(R/I)$ sur $Y$. Si $\mathfrak p \subset R$
est l'idéal premier correspondant à
$z$, alors on voit que $I_{\mathfrak p} \subset R_{\mathfrak p}$
n'est pas une égalité. Par conséquent
(la localisation étant exacte, voir Algèbre, Proposition
\ref{algebra-proposition-localization-exact}),
le morphisme $R_{\mathfrak p} \to
(A_1)_{\mathfrak p} \times \ldots \times (A_n)_{\mathfrak p}$ n'est pas nul. Donc
l'un des anneaux $(A_i)_{\mathfrak p}$ n'est pas nul. Il existe
donc un $i$ et un premier $\mathfrak q_i \subset A_i$ superposés
à un premier $\mathfrak p_i \subset \mathfrak p$. Par Algèbre,
Lemme \ref{algebra-lemma-dominate} on peut choisir un anneau de valorisation
$A \subset K = \kappa(\mathfrak q_i)$ dominant l'anneau
local
$R_{\mathfrak p}/\mathfrak p_iR_{\mathfrak p} \subset \kappa(\mathfrak q_i)$. Cela donne le diagramme souhaité.
Certains détails omis.
\end{proof}

\begin{lemma}
\label{morphisms-lemma-factor-factor}
Soit
$$
\xymatrix{
X_1 \ar[d] \ar[r]_{f_1} & Y_1 \ar[d] \\
X_2 \ar[r]^{f_2} & Y_2
}
$$
un diagramme commutatif de schémas. Soit $Z_i \subset Y_i$, $i = 1, 2$ l'image
schématique de $f_i$. Alors le morphisme $Y_1 \to Y_2$
induit un morphisme $Z_1 \to Z_2$ et un diagramme
commutatif
$$
\xymatrix{
X_1 \ar[r] \ar[d] & Z_1 \ar[d] \ar[r] & Y_1 \ar[d] \\
X_2 \ar[r] & Z_2 \ar[r] & Y_2
}
$$
\end{lemma}

\begin{proof}
L'image réciproque schématique de $Z_2$ dans $Y_1$ est
un sous-schéma fermé de $Y_1$ par lequel
$f_1$ se factorise. Par conséquent, $Z_1$ est contenu dans
ceci. Cela prouve le lemme.
\end{proof}

\begin{lemma}
\label{morphisms-lemma-scheme-theoretic-image-reduced}
Soit $f : X \to Y$ un morphisme de schémas. Si $X$
est réduit, alors l'image schématique de $f$ est
la structure schématique induite réduite sur $\overline{f(X)}$.
\end{lemma}

\begin{proof}
Ceci est vrai car la structure schématique induite réduite sur $\overline{f(X)}$ est clairement
le plus petit sous-schéma fermé de $Y$ à travers lequel $f$
se
factorise par, voir Schémas, Lemme \ref{schemes-lemma-map-into-reduction}.
\end{proof}

\begin{lemma}
\label{morphisms-lemma-scheme-theoretic-image-of-partial-section}
Soit $f : X \to Y$ un morphisme de schémas séparé. Soit
$V \subset Y$ un ouvert rétrocompact. Soit $s : V \to X$ un morphisme
tel que $f \circ s = \text{id}_V$. Soit $Y'$ l'image
schématique de $s$. Alors $Y' \to Y$
est un isomorphisme sur $V$.
\end{lemma}

\begin{proof}
L'hypothèse que $V$ est rétrocompact
dans $Y$ (Topologie, Définition \ref{topology-definition-quasi-compact})
signifie que $V \to Y$ est un morphisme
quasi-compact. Par Schémas, Lemme \ref{schemes-lemma-quasi-compact-permanence} le morphisme
$s : V \to X$ est quasi-compact. D'où
la construction de l'image schématique $Y'$
de $s$ navette avec restriction
aux ouverts par le lemme \ref{morphisms-lemma-quasi-compact-scheme-theoretic-image}. En particulier,
on voit que $Y' \cap f^{-1}(V)$ est l'image schématique
d'une section du morphisme séparé $f^{-1}(V) \to V$.
Puisqu'une section d'un morphisme séparé
est une immersion fermée
(Schémas, Lemme \ref{schemes-lemma-section-immersion}), Nous concluons
que $Y' \cap f^{-1}(V) \to V$
est un isomorphisme souhaité.
\end{proof}







\section{Adhérence schématique et densité}
\label{morphisms-section-scheme-theoretic-closure}

\noindent
On reprend la définition suivante de \cite[IV, Définition 11.10.2]{EGA}.

\begin{definition}
\label{morphisms-definition-scheme-theoretically-dense}
Soit $X$ un schéma. Soit $U \subset X$ un sous-schéma ouvert.
\begin{enumerate}
\item L'image schématique du morphisme $U \to X$ est appelée {\it
adhérence schématique de $U$ dans $X$}.
\item On dit que $U$ est {\it schématiquement dense dans
$X$} si pour tout $V \subset X$ ouvert l'adhérence schématique de $U \cap V$
dans $V$ est égale à $V$.
\end{enumerate}
\end{definition}

\noindent
Avec cette définition,
il n'est {\bf pas} que $U$ soit
schématiquement dense dans $X$ si et
seulement si l'adhérence schématique de $U$
est $X$, voir Exemple
\ref{morphisms-example-scheme-theretically-dense-not-dense}. C'est quelque peu inélégant ; mais voir Lemmes
\ref{morphisms-lemma-scheme-theoretically-dense-quasi-compact} et \ref{morphisms-lemma-reduced-scheme-theoretically-dense} ci-dessous.
En revanche, avec cette définition $U$ est schématiquement
dense en $X$ si et seulement si pour chaque $V \subset X$
ouvert l'homomorphisme d'anneaux $\mathcal{O}_X(V) \to \mathcal{O}_X(U \cap V)$ est
injectif, voir Lemme \ref{morphisms-lemma-characterize-scheme-theoretically-dense} ci-dessous. En
particulier on voit que schématiquement dense implique dense
ce qui est agréable.

\begin{example}
\label{morphisms-example-scheme-theretically-dense-not-dense}
Voici un exemple où le fait que l'adhérence schématique soit $X$ n'implique
pas la densité des espaces topologiques sous-jacents.
Soit $k$
un corps. Posons $A = k[x, z_1, z_2, \ldots]/(x^n z_n)$.
Posons $I = (z_1, z_2, \ldots) \subset A$. Considérons le schéma
affine $X = \Spec(A)$ et le sous-schéma ouvert $U = X \setminus V(I)$.
Comme $A \to \prod_n A_{z_n}$ est injectif,
l'adhérence schématique de $U$ est $X$. Considérons
le morphisme $X \to \Spec(k[x])$. Ce morphisme
est surjectif (poser tous les $z_n = 0$ permet
de le voir). Mais sa restriction à $U$
n'est pas surjective, car son image est réduite
au point $x = 0$. Par conséquent, $U$ ne peut pas être topologiquement dense dans
$X$.
\end{example}

\begin{lemma}
\label{morphisms-lemma-scheme-theoretically-dense-quasi-compact}
Soit $X$ un schéma.
Soit $U \subset X$ un sous-schéma ouvert. Si le morphisme
d'inclusion $U \to X$ est quasi-compact, alors $U$ est schématiquement
dense dans $X$ si et seulement si l'adhérence schématique de $U$
dans $X$ est $X$.
\end{lemma}

\begin{proof}
Cela résulte de la partie (3) du Lemme \ref{morphisms-lemma-quasi-compact-scheme-theoretic-image}.
\end{proof}

\begin{example}
\label{morphisms-example-scheme-theoretic-closure}
Soit $A$ un anneau et $X = \Spec(A)$.
Soient $f_1, \ldots, f_n \in A$ et posons $U = D(f_1) \cup \ldots \cup D(f_n)$.
Posons $I = \Ker(A \to \prod A_{f_i})$. Alors l'adhérence schématique de $U$
dans $X$ est le sous-schéma fermé
$\Spec(A/I)$ de $X$. Notons que
$U \to X$ est quasi-compact. Ainsi par
le lemme \ref{morphisms-lemma-scheme-theoretically-dense-quasi-compact} on voit que $U$ est schématiquement
dense en $X$ si et seulement si $I = 0$.
\end{example}

\begin{lemma}
\label{morphisms-lemma-characterize-scheme-theoretically-dense}
Soit $j : U \to X$ une immersion ouverte de schémas.
Alors $U$ est schématiquement dense dans $X$
si et seulement si $\mathcal{O}_X \to j_*\mathcal{O}_U$ est injectif.
\end{lemma}

\begin{proof}
Si $\mathcal{O}_X \to j_*\mathcal{O}_U$ est injectif, alors il en va de même lorsqu'il
est restreint à n'importe quel $V$ ouvert de $X$.
D'où l'adhérence schématique de $U \cap V$ dans $V$
est égale à $V$, voir démonstration du lemme \ref{morphisms-lemma-scheme-theoretic-image}.
A l'inverse, supposons que l'adhérence schématique
de $U \cap V$ soit égale à $V$ pour toutes les
ouvertures $V$. Supposons que $\mathcal{O}_X \to j_*\mathcal{O}_U$ ne soit pas injectif.
Ensuite, nous pouvons trouver un ouvert affine, disons $\Spec(A) = V \subset X$
et un élément non nul $f \in A$ tel que l'image de $f$ soit
nulle dans $\Gamma(V \cap U, \mathcal{O}_X)$. Dans ce cas l'adhérence schématique de $V \cap U$
dans $V$ est clairement contenue dans $\Spec(A/(f))$ une
contradiction.
\end{proof}

\begin{lemma}
\label{morphisms-lemma-intersection-scheme-theoretically-dense}
Soit $X$ un schéma. Si $U$, $V$ sont schématiquement denses
sous-schémas ouverts de $X$, alors $U \cap V$.
\end{lemma}

\begin{proof}
Soit $W \subset X$ un ouvert
quelconque. Considérons le morphisme
$\mathcal{O}_X(W) \to \mathcal{O}_X(W \cap V)
\to \mathcal{O}_X(W \cap V \cap U)$. D'après le Lemme \ref{morphisms-lemma-characterize-scheme-theoretically-dense},
les deux applications sont injectives.
Le composite est donc injectif. Donc
par le lemme \ref{morphisms-lemma-characterize-scheme-theoretically-dense} $U \cap V$ est schématiquement
dense en $X$.
\end{proof}

\begin{lemma}
\label{morphisms-lemma-quasi-compact-immersion}
Soit $h : Z \to X$ une immersion. Supposons que $h$ soit quasi-compact
ou que $Z$ soit réduit. Soit $\overline{Z} \subset X$ l'image schématique de
$h$. Alors le morphisme $Z \to \overline{Z}$ est une immersion ouverte qui identifie
$Z$ à un sous-schéma ouvert schématiquement
dense de $\overline{Z}$. De plus, $Z$ est topologiquement
dense dans $\overline{Z}$.
\end{lemma}

\begin{proof}
Par le Lemme \ref{morphisms-lemma-factor-quasi-compact-immersion} ou Lemme \ref{morphisms-lemma-factor-reduced-immersion}, nous
pouvons factoriser $Z \to X$ comme $Z \to \overline{Z}_1 \to X$
avec $Z \to \overline{Z}_1$ ouvert et $\overline{Z}_1 \to X$ fermé. Par contre, soit
$Z \to \overline{Z} \subset X$ l'image schématique de $Z \to X$. Nous
concluons que $\overline{Z} \subset \overline{Z}_1$. Puisque $Z$ est un sous-schéma
ouvert de $\overline{Z}_1$ Il s'ensuit que $Z$
est également un sous-schéma ouvert de $\overline{Z}$.
Dans le cas où $Z$ est réduit nous savons
que $Z \subset \overline{Z}_1$ est topologiquement dense par la construction
de $\overline{Z}_1$ dans la démonstration du lemme \ref{morphisms-lemma-factor-reduced-immersion}.
Par conséquent, $\overline{Z}_1$ et $\overline{Z}$ ont
les mêmes espaces topologiques sous-jacents
Ainsi $\overline{Z} \subset \overline{Z}_1$ est une immersion fermée dans un schéma réduit
qui induit une bijection sur les espaces topologiques sous-jacents,
et donc un isomorphisme. Dans le cas où $Z \to X$
est quasi-compact nous argumentons comme suit
: L'affirmation selon laquelle $Z$ est
schématiquement dense
dans $\overline{Z}$ découle de la partie (3) du lemme \ref{morphisms-lemma-quasi-compact-scheme-theoretic-image}.
La dernière assertion
découle de la partie (4) du lemme \ref{morphisms-lemma-quasi-compact-scheme-theoretic-image}.
\end{proof}

\begin{lemma}
\label{morphisms-lemma-reduced-scheme-theoretically-dense}
Soit $X$ un schéma réduit et soit $U \subset X$ un sous-schéma ouvert. Alors Les
conditions suivantes sont équivalentes
\begin{enumerate}
\item $U$ est topologiquement dense dans $X$,
\item l'adhérence schématique de $U$ dans $X$ est $X$, et
\item $U$ est schématiquement dense en $X$.
\end{enumerate}
\end{lemma}

\begin{proof}
Cela résulte du lemme
\ref{morphisms-lemma-quasi-compact-immersion} et le fait qu'un sous-schéma
fermé $Z$ de $X$ dont l'espace topologique
sous-jacent est égal à $X$ doit être égal à $X$ comme
schéma.
\end{proof}

\begin{lemma}
\label{morphisms-lemma-reduced-subscheme-closure}
Soit $X$ un schéma et soit $U \subset X$ un sous-schéma ouvert réduit. Alors Les
conditions suivantes sont équivalentes
\begin{enumerate}
\item l'adhérence schématique de $U$ dans $X$ est $X$, et
\item $U$ est schématiquement dense dans $X$.
\end{enumerate}
Si cela est vrai alors $X$ est un schéma réduit.
\end{lemma}

\begin{proof}
Cela résulte du
lemme \ref{morphisms-lemma-quasi-compact-immersion} et du fait que
l'adhérence schématique de $U$ dans $X$
est
réduite par le lemme \ref{morphisms-lemma-scheme-theoretic-image-reduced}.
\end{proof}


\begin{lemma}
\label{morphisms-lemma-equality-of-morphisms}
Soit $S$ un schéma. Soient $X$, $Y$ des schémas
sur $S$. Soit $f, g : X \to Y$ des morphismes de schémas sur
$S$. Soit $U \subset X$ un sous-schéma ouvert tel que $f|_U = g|_U$.
Si l'adhérence schématique de $U$ dans $X$
est $X$ et que $Y \to S$ est séparé, alors $f = g$.
\end{lemma}

\begin{proof}
Dérivé des définitions
et Schémas, Lemme \ref{schemes-lemma-where-are-they-equal}.
\end{proof}




\section{Morphismes dominants}
\label{morphisms-section-dominant}

\noindent
La définition d'un morphisme de schémas dominant est un peu différente
de ce à quoi on pourrait s'attendre si l'on est habitué à la notion
de morphisme dominant des variétés.

\begin{definition}
\label{morphisms-definition-dominant}
Un morphisme $f : X \to S$ de schémas est dit {\it dominant} si l'image de
$f$ est un sous-ensemble dense de $S$.
\end{definition}

\noindent
Ainsi, par exemple, si $k$ est un corps infini et $\lambda_1, \lambda_2, \ldots$ est une collection
dénombrable d'éléments distincts de $k$, alors le morphisme
$$
\coprod\nolimits_{i = 1, 2, \ldots } \Spec(k)
\longrightarrow
\Spec(k[x])
$$
avec le $i$ème facteur envoyé au point $x = \lambda_i$ est dominant.

\begin{lemma}
\label{morphisms-lemma-generic-points-in-image-dominant}
Soit $f : X \to S$ un morphisme de schémas. Si chaque
point générique de chaque composante irréductible de $S$ est à l'image
de $f$, alors $f$ est dominant.
\end{lemma}

\begin{proof}
Il s'agit d'un fait topologique qui découle directement du fait
que l'espace topologique sous-jacent à un schéma est
sobre, voir Schémas, Lemme \ref{schemes-lemma-scheme-sober}, et que tout point
de $S$ est contenu dans une composante irréductible
de $S$, voir Topologie, Lemme \ref{topology-lemma-irreducible}.
\end{proof}

\noindent
L'attente selon laquelle les morphismes sont dominants uniquement si des points génériques
de la cible sont dans l'image est valable si le morphisme est quasi-compact.

\begin{lemma}
\label{morphisms-lemma-quasi-compact-dominant}
\begin{slogan}
Les morphismes dont l'image contient les points génériques sont dominants
\end{slogan}
Soit $f : X \to S$ un morphisme de schémas quasi-compact. Alors $f$
est dominant si et seulement si pour chaque composante irréductible
$Z \subset S$ le point générique de $Z$ est à l'image de
$f$.
\end{lemma}

\begin{proof}
Soit $V \subset S$ un ouvert affine. Parce
que $f$ est quasi-compact, nous pouvons choisir un nombre
fini d'ouverts affines $U_i \subset f^{-1}(V)$, $i = 1, \ldots, n$ couvrant $f^{-1}(V)$.
Considérons le morphisme des affines
$$
f' :
\coprod\nolimits_{i = 1, \ldots, n} U_i
\longrightarrow
V.
$$
Une union disjointe d'affines est affine, voir
Schémas, Lemme \ref{schemes-lemma-disjoint-union-affines}. Les points génériques des composantes
irréductibles de $V$ sont exactement
les points génériques des composantes irréductibles de $S$ qui
répondent à $V$. De plus, $f$ est dominant si et seulement
si $f'$ est dominant quels que soient les choix de $V, n, U_i$
que nous faisons ci-dessus. Nous avons ainsi réduit le lemme au cas
d'un morphisme de
schémas affines. Le cas affine est l'Algèbre, Lemme \ref{algebra-lemma-image-dense-generic-points}.
\end{proof}

\begin{lemma}
\label{morphisms-lemma-base-change-dominant}
Soit $f : X \to S$ un morphisme de schémas dominant quasi-compact. Soit $g : S' \to S$
un morphisme de schémas, et soit $f' : X' \to S'$ le changement de base de $f$
par $g$. Si les générisations s'élèvent le long de $g$, alors
$f'$ est dominant.
\end{lemma}

\noindent
Nous verrons que les générisations s'élèvent le long de $g$
si $g$ est ouvert
(Lemme \ref{morphisms-lemma-open-generizing}) ou plat (Lemme \ref{morphisms-lemma-generalizations-lift-flat}).

\begin{proof}
Observez que $f'$ est quasi-compact par Schémas,
Lemme \ref{schemes-lemma-quasi-compact-preserved-base-change}. Soit $\eta' \in S'$ le point générique
d'une composante irréductible de $S'$.
Si les générisations s'élèvent le long de $g$,
alors $\eta = g(\eta')$ est le point générique d'une composante
irréductible de $S$. Par le Lemme \ref{morphisms-lemma-quasi-compact-dominant}
on voit que $\eta$ est à l'image de
$f$. Ainsi $\eta'$ est à
l'image de $f'$ d'après Schémas, Lemme \ref{schemes-lemma-points-fibre-product}.
Il s'ensuit que $f'$ est dominant par le lemme \ref{morphisms-lemma-quasi-compact-dominant}.
\end{proof}

\noindent
Cela vaut en partie pour les morphismes dominants arbitraires (pas nécessairement
quasi-compacts).

\begin{lemma}
\label{morphisms-lemma-open-base-change-dominant}
Soit $f : X \to S$ un morphisme de schémas dominant. Soit $g : S' \to S$
un morphisme de schémas ouvert, et soit $f' : X' \to S'$ le changement de base
de $f$ par $g$. Alors $f'$ est dominant.
\end{lemma}

\begin{proof}
Si $f'$ n'est pas dominant, alors $f'(X')$ est contenu dans un sous-ensemble
fermé $B$ différent de $S'$. Ensuite, $S'-B$ est ouvert dans $S'$
et non vide. Puisque $g$ est ouvert, $g(S'-B)$ n'est pas vide et ouvert
dans $S$. Puisque $f(X)$ est dense dans $S$, il existe un
point $x\in X$ avec $f(x) \in g(S'-B)$. Soit
$W = \Spec(\kappa (f(x)))$, et écrivez $X_W$, $S_W$, $S'_W$
pour les retraits de ces schémas de $S$ vers $W$. Alors $X_W$ et
$(S'-B)_W$ ne sont pas vides, et donc $X_W \times_W (S'-B)_W$ n'est pas vide, car $W$
est le spectre d'un corps. Cela contredit le fait que $f'(X \times_S S')$ soit contenu
dans $B \subset S'$. Donc en fait $f'$ est dominant.
\end{proof}

\begin{lemma}
\label{morphisms-lemma-quasi-compact-generic-point-not-in-image}
Soit $f : X \to S$ un morphisme de schémas quasi-compact.
Soit $\eta \in S$ un point générique d'une composante
irréductible de $S$. Si $\eta \not \in f(X)$ alors
il existe un voisinage ouvert $V \subset S$ de
$\eta$ tel que $f^{-1}(V) = \emptyset$.
\end{lemma}

\begin{proof}
Soit $Z \subset S$ l'image schématique de $f$.
Nous devons montrer que $\eta \not \in Z$.
Cela résulte du
lemme \ref{morphisms-lemma-reach-points-scheme-theoretic-image} mais peut aussi être vu comme suit.
Par le Lemme \ref{morphisms-lemma-quasi-compact-scheme-theoretic-image} le morphisme
$X \to Z$ est dominant, ce qui par le lemme
\ref{morphisms-lemma-quasi-compact-dominant} signifie que tous les points
génériques de toutes les composantes
irréductibles de $Z$ sont à l'image de $X \to Z$.
Par hypothèse on voit que $\eta \not \in Z$ puisque $\eta$
serait le point générique d'une composante
irréductible de $Z$ s'il était dans $Z$.
\end{proof}

\noindent
Il existe un autre cas où la dominante revient à avoir tous les points
génériques des composantes irréductibles dans l'image.

\begin{lemma}
\label{morphisms-lemma-dominant-finite-number-irreducible-components}
Soit $f : X \to S$ un morphisme de schémas. Supposons que
$X$ ait un nombre fini de composantes irréductibles. Alors $f$
est dominant (si et) seulement si pour toute composante irréductible
$Z \subset S$ le point générique de $Z$ est à l'image de $f$.
Si tel est le cas, alors $S$ a également un nombre fini de composantes
irréductibles.
\end{lemma}

\begin{proof}
Supposons que $f$
soit dominant. Disons que $X = Z_1 \cup Z_2 \cup \ldots \cup Z_n$ est la décomposition de $X$
en composantes irréductibles. Soit $\xi_i \in Z_i$ son
point générique, donc $Z_i = \overline{\{\xi_i\}}$. Notez que $f(Z_i)$
est un sous-ensemble irréductible de
$S$.
$$
S = \overline{f(X)} = \bigcup \overline{f(Z_i)} =
\bigcup \overline{\{f(\xi_i)\}}
$$
est donc une union finie de sous-ensembles irréductibles dont les points
génériques sont à l'image de $f$. Le lemme suit.
\end{proof}

\begin{lemma}
\label{morphisms-lemma-dominant-between-integral}
Soit $f : X \to Y$ un morphisme de schémas intègres. Les conditions suivantes sont
équivalentes
\begin{enumerate}
\item $f$ est dominant,
\item $f$ envoie le point générique de $X$ sur le point générique de $Y$,
\item il existe des ouverts affines non vides $U \subset X$ et
$V \subset Y$ tels que $f(U) \subset V$ et l'homomorphisme d'anneaux $\mathcal{O}_Y(V) \to \mathcal{O}_X(U)$ est
injectif,
\item pour tous les ouverts affines non vides $U \subset X$ et
$V \subset Y$ tels que $f(U) \subset V$, l'homomorphisme d'anneaux $\mathcal{O}_Y(V) \to \mathcal{O}_X(U)$ est
injectif,
\item il existe $x \in X$, d'image $y = f(x) \in Y$, tel que l'homomorphisme
d'anneaux locaux $\mathcal{O}_{Y, y} \to \mathcal{O}_{X, x}$ soit injectif, et
\item pour tout $x \in X$, d'image $y = f(x) \in Y$, l'homomorphisme
d'anneaux locaux $\mathcal{O}_{Y, y} \to \mathcal{O}_{X, x}$ est injectif.
\end{enumerate}
\end{lemma}

\begin{proof}
L'équivalence de (1) et (2) découle du lemme
\ref{morphisms-lemma-dominant-finite-number-irreducible-components}. Soient $U \subset X$ et $V \subset Y$ des ouverts affines
non vides avec $f(U) \subset V$. Rappelons que les anneaux $A = \mathcal{O}_X(U)$
et $B = \mathcal{O}_Y(V)$ sont des domaines intègres. Le morphisme
$f|_U : U \to V$ correspond à un homomorphisme
d'anneaux $\varphi : B \to A$. Les points génériques de $X$
et $Y$ correspondent aux idéaux premiers $(0) \subset A$
et $(0) \subset B$. Ainsi (2) équivaut à la condition $(0) = \varphi^{-1}((0))$,
c'est-à-dire à la condition que $\varphi$ soit injectif.
De cette façon, nous voyons que (2), (3) et
(4) sont équivalents. De même, étant donné $x$
et $y$ comme dans (5), les anneaux locaux
$\mathcal{O}_{X, x}$ et $\mathcal{O}_{Y, y}$ sont des domaines et les idéaux
premiers $(0) \subset \mathcal{O}_{X, x}$ et $(0) \subset \mathcal{O}_{Y, y}$ correspondent
aux points génériques de $X$ et $Y$
(via l'identification du spectre de l'anneau local à
$x$ avec l'ensemble des points
spécialisés à $x$, voir Schémas,
Lemme \ref{schemes-lemma-specialize-points}). Nous pouvons donc argumenter exactement
de la même manière que ci-dessus pour voir
que (2), (5) et (6) sont équivalents.
\end{proof}







\section{Morphismes surjectifs}
\label{morphisms-section-surjective}

\begin{definition}
\label{morphisms-definition-surjective}
Un morphisme de schémas est dit {\it surjectif}
s'il est surjectif sur les espaces topologiques
sous-jacents.
\end{definition}

\begin{lemma}
\label{morphisms-lemma-composition-surjective}
La composition des morphismes surjectifs est surjectif.
\end{lemma}

\begin{proof}
Omis.
\end{proof}

\begin{lemma}
\label{morphisms-lemma-when-point-maps-to-pair}
Soient $X$ et $Y$ des schémas sur un schéma de base $S$. Étant donné les points $x \in X$
et $y \in Y$, il y a un point de $X \times_S Y$ envoyé sur $x$ et $y$ sous les projections
si et seulement si $x$ et $y$ se trouvent au-dessus du même point de $S$.
\end{lemma}

\begin{proof}
La condition est évidemment nécessaire, et la réciproque découle de
la démonstration d'après Schémas, Lemme \ref{schemes-lemma-points-fibre-product}.
\end{proof}

\begin{lemma}
\label{morphisms-lemma-base-change-surjective}
Le changement de base d'un morphisme surjectif est surjectif.
\end{lemma}

\begin{proof}
Soit $f: X \to Y$ un morphisme de schémas sur un schéma de base $S$.
Si $S' \to S$ est un morphisme de schémas, soit $p: X_{S'} \to X$
et $q: Y_{S'} \to Y$ les projections canoniques. Le carré
commutatif
$$
\xymatrix{
X_{S'} \ar[d]_{f_{S'}} \ar[r]_p & X \ar[d]^{f} \\
Y_{S'} \ar[r]^{q} & Y.
}
$$
identifie $X_{S'}$ comme un produit fibré de $X \to Y$
et $Y_{S'} \to Y$. Soit $Z$ un sous-ensemble de l'espace
topologique sous-jacent de $X$. Puis $q^{-1}(f(Z)) = f_{S'}(p^{-1}(Z))$,
car $y' \in q^{-1}(f(Z))$ si et seulement si $q(y') = f(x)$ pour certains $x \in Z$,
si et seulement si, par le lemme \ref{morphisms-lemma-when-point-maps-to-pair}, il existe des $x' \in X_{S'}$
tels que $f_{S'}(x') = y'$ et $p(x') = x$. En particulier en prenant $Z = X$
on voit que si $f$ est surjectif le changement de base
$f_{S'}: X_{S'} \to Y_{S'}$ l'est aussi.
\end{proof}

\begin{example}
\label{morphisms-example-injective-not-preserved-base-change}
La bijectivité n'est pas stable par changement
de base, et l'injectivité ne l'est
pas non plus. Par exemple, considérons
la bijection
$\Spec(\mathbf{C}) \to \Spec(\mathbf{R})$. Le changement de base
$\Spec(\mathbf{C} \otimes_{\mathbf{R}} \mathbf{C}) \to
\Spec(\mathbf{C})$ n'est pas injectif, car il existe un isomorphisme
$\mathbf{C} \otimes_{\mathbf{R}} \mathbf{C} \cong \mathbf{C} \times \mathbf{C}$
(la décomposition
provient de l'idempotent $\frac{1 \otimes 1 + i \otimes i}{2}$) ;
ainsi, $\Spec(\mathbf{C} \otimes_{\mathbf{R}} \mathbf{C})$ a deux points.
\end{example}

\begin{lemma}
\label{morphisms-lemma-surjection-from-quasi-compact}
Soit
$$
\xymatrix{
X \ar[rr]_f \ar[rd]_p & &
Y \ar[dl]^q \\
& Z
}
$$
être un diagramme commutatif de morphismes de schémas. Si $f$
est surjectif et $p$ est quasi-compact, alors $q$ est quasi-compact.
\end{lemma}

\begin{proof}
Soit $W \subset Z$ un ouvert quasi-compact. Par hypothèse $p^{-1}(W)$
est quasi-compact.
Donc par Topologie, Lemme \ref{topology-lemma-image-quasi-compact} l'image
inverse $q^{-1}(W) = f(p^{-1}(W))$ est également quasi-compacte. Cela
prouve le lemme.
\end{proof}




\section{Morphismes radiciels et universellement injectifs}
\label{morphisms-section-radicial}

\noindent
Dans cette section, nous définissons ce que signifie pour un morphisme de schémas d'être {\it
radical} et ce que signifie pour un morphisme de schémas d'être {\it
universellement injectif}. Nous montrons ensuite que ces notions
concordent. La raison pour laquelle nous introduisons les deux est que dans le cas des espaces algébriques,
il existe des notions correspondantes qui ne concordent pas toujours.

\begin{definition}
\label{morphisms-definition-universally-injective}
Soit $f : X \to S$ un morphisme.
\begin{enumerate}
\item On dit que $f$ est {\it universellement injectif}
si et seulement si pour tout morphisme de schémas $S' \to S$ le changement
de base $f' : X_{S'} \to S'$ est injectif (sur les espaces topologiques sous-jacents).
\item On dit que $f$ est {\it radical} si $f$
est injectif comme application d'espaces topologiques, et pour tout
$x \in X$ l'extension du corps $\kappa(x)/\kappa(f(x))$ est purement inséparable.
\end{enumerate}
\end{definition}

\begin{lemma}
\label{morphisms-lemma-universally-injective}
\begin{reference}
\cite[Proposition 3.7.1]{EGA1-second}
\end{reference}
Soit $f : X \to S$ un morphisme de schémas. Les conditions
suivantes sont équivalentes :
\begin{enumerate}
\item Pour chaque corps $K$ l'application
induite $\Mor(\Spec(K), X) \to \Mor(\Spec(K), S)$ est
injectif.
\item Le morphisme $f$ est universellement injectif.
\item Le morphisme $f$ est radical.
\item Le morphisme diagonal $\Delta_{X/S} : X \longrightarrow X \times_S X$ est
surjectif.
\end{enumerate}
\end{lemma}

\begin{proof}
Soit $K$ un corps, et $s : \Spec(K) \to S$ un morphisme.
Donner un morphisme $x : \Spec(K) \to X$ tel que $f \circ x = s$ revient
à donner une section de la projection
$X_K = \Spec(K) \times_S X \to \Spec(K)$, ce qui revient à donner un point
$x \in X_K$ dont le corps résiduel est $K$. On voit
donc que (2) implique (1).

\medskip\noindent Inversement,
 supposons que (1) soit vérifié. Supposons que $x, x' \in X_{S'}$ s'envoient
sur le même point $s' \in S'$. Choisissons un diagramme commutatif
$$
\xymatrix{
K & \kappa(x) \ar[l] \\
\kappa(x') \ar[u] & \kappa(s') \ar[l] \ar[u]
}
$$
de corps. Par Schémas, Lemme \ref{schemes-lemma-characterize-points} on obtient deux morphismes
$a, a' : \Spec(K) \to X_{S'}$. L'un correspondant au point $x$ et à l'intégration
$\kappa(x) \subset K$ et l'autre correspondant au point $x'$
et à l'intégration $\kappa(x') \subset K$. Nous avons également $f' \circ a = f' \circ a'$.
La condition (1) implique désormais que les compositions
de $a$ et $a'$ avec $X_{S'} \to X$ sont égales. Puisque
$X_{S'}$ est le produit fibré de $S'$ et $X$
sur $S$ on voit que $a = a'$. D'où $x = x'$. Ainsi
(1) implique (2).

\medskip\noindent
S'il y a deux points différents $x, x' \in X$ envoyés au même point de $s$
alors (2) est violé.
Si pour certains $s = f(x)$, $x \in X$ l'extension
de corps $\kappa(x)/\kappa(s)$ n'est pas purement inséparable,
alors on peut trouver une extension de corps $K/\kappa(s)$
telle que $\kappa(x)$ possède deux $\kappa(s)$-homomorphismes
en $K$. Par Schémas, Lemme \ref{schemes-lemma-characterize-points} cela
implique que l'application
$\Mor(\Spec(K), X) \to \Mor(\Spec(K), S)$ n'est pas injectif,
et donc (1) est violé. Nous voyons donc
que les conditions équivalentes (1) et (2) impliquent que
$f$ est radical, c'est-à-dire qu'elles impliquent (3).

\medskip\noindent
Supposons
(3). Par Schémas, Lemme \ref{schemes-lemma-characterize-points} un morphisme
$\Spec(K) \to X$ est donné par un couple $(x, \kappa(x) \to K)$. La propriété
(3) dit exactement qu'en associant au couple $(x, \kappa(x) \to K)$
le couple $(s, \kappa(s) \to \kappa(x) \to K)$ est injectif. En d'autres termes,
(1) est valable. À ce stade, nous savons que (1), (2)
et (3) sont tous équivalents.

\medskip\noindent
Enfin, nous prouvons l'équivalence de (4) avec (1),
(2) et (3). Un point de $X \times_S X$ est
donné par un quadruple $(x_1, x_2, s, \mathfrak p)$, où $x_1, x_2 \in X$,
$f(x_1) = f(x_2) = s$ et $\mathfrak p \subset \kappa(x_1) \otimes_{\kappa(s)} \kappa(x_2)$
est un idéal premier, voir Schémas, Lemme \ref{schemes-lemma-points-fibre-product}.
Si $f$
est universellement injectif, alors en prenant
$S'=X$ dans la définition
de universellement injectif, $\Delta_{X/S}$ doit être
surjectif puisqu'il s'agit d'une section du morphisme
injectif
$X \times_S X  \longrightarrow X$. Inversement, si $\Delta_{X/S}$
est surjectif,
alors toujours $x_1 = x_2 = x$ et il existe exactement un tel idéal
premier $\mathfrak p$, ce qui signifie que $\kappa(s) \subset \kappa(x)$ est purement
inséparable. Donc $f$ est radical.
Alternativement, si $\Delta_{X/S}$
est surjectif, alors pour tout $S' \to S$
le changement de base $\Delta_{X_{S'}/S'}$
est surjectif ce qui implique que $f$ est universellement
injectif. Ceci termine la démonstration du lemme.
\end{proof}

\begin{lemma}
\label{morphisms-lemma-universally-injective-separated}
Un morphisme universellement injectif est séparé.
\end{lemma}

\begin{proof}
Combiner
Lemme \ref{morphisms-lemma-universally-injective} avec la remarque
que $X \to S$ est séparé si et seulement si l'image
de $\Delta_{X/S}$ est fermée dans $X \times_S X$,
voir Schémas, Définition \ref{schemes-definition-separated} et la
discussion qui en suit.
\end{proof}

\begin{lemma}
\label{morphisms-lemma-base-change-universally-injective}
Un changement de base d'un morphisme universellement injectif est universellement injectif.
\end{lemma}

\begin{proof}
Ceci est formel.
\end{proof}

\begin{lemma}
\label{morphisms-lemma-composition-universally-injective}
Une composition de morphismes radiciels est radicale, et il en va donc de même
pour la condition équivalente d'être universellement injectif.
\end{lemma}

\begin{proof}
Omis.
\end{proof}









\section{Morphismes affines}
\label{morphisms-section-affine}

\begin{definition}
\label{morphisms-definition-affine}
Un morphisme de schémas $f : X \to S$ est appelé {\it affine}
si l'image réciproque de chaque ouvert affine de $S$ est un ouvert affine
de $X$.
\end{definition}

\begin{lemma}
\label{morphisms-lemma-affine-separated}
Un morphisme affine est séparé et quasi-compact.
\end{lemma}

\begin{proof}
Soit $f : X \to S$ affine. La quasi-compacité est immédiate chez Schémas,
Lemme \ref{schemes-lemma-quasi-compact-affine}. Nous allons montrer que $f$
est séparé à l'aide d'après
Schémas, Lemme \ref{schemes-lemma-characterize-separated}. Soient $x_1, x_2 \in X$ des points de $X$
qui ont la même image $s \in S$. Choisissons un
ouvert affine quelconque $W \subset S$ contenant $s$. Par hypothèse,
$f^{-1}(W)$ est affine. Appliquer le lemme cité avec $U = V = f^{-1}(W)$.
\end{proof}

\begin{lemma}
\label{morphisms-lemma-characterize-affine}
\begin{reference}
\cite[II, Corollaire 1.3.2]{EGA}
\end{reference}
Soit $f : X \to S$ un morphisme de schémas. Les conditions
suivantes sont équivalentes
\begin{enumerate}
\item Le morphisme $f$ est affine.
\item Il existe un recouvrement ouvert affine $S = \bigcup W_j$ tel
que chaque $f^{-1}(W_j)$ soit affine.
\item Il existe un faisceau quasi-cohérent de $\mathcal{O}_S$-algèbres
$\mathcal{A}$ et un
isomorphisme $X \cong \underline{\Spec}_S(\mathcal{A})$ de schémas
sur $S$. Voir
Constructions, section \ref{constructions-section-spec} pour la notation.
\end{enumerate}
De plus, dans ce cas $X = \underline{\Spec}_S(f_*\mathcal{O}_X)$.
\end{lemma}

\begin{proof}
Il est évident que (1) implique (2).

\medskip\noindent
Supposons que $S = \bigcup_{j \in J} W_j$ soit un recouvrement ouvert
affine tel que chaque $f^{-1}(W_j)$
soit affine. Par Schémas, Lemme \ref{schemes-lemma-quasi-compact-affine}
on voit que $f$
est quasi-compact. Par Schémas, Lemme \ref{schemes-lemma-characterize-quasi-separated}
on voit que le morphisme $f$ est
quasi-séparé. Ainsi par Schémas, Lemme \ref{schemes-lemma-push-forward-quasi-coherent} le
faisceau $\mathcal{A} = f_*\mathcal{O}_X$ est un faisceau quasi-cohérent
de $\mathcal{O}_S$-algèbres. On a donc le schéma
$g : Y = \underline{\Spec}_S(\mathcal{A}) \to S$ sur $S$. L'application
d'identité
$\text{id} : \mathcal{A} = f_*\mathcal{O}_X \to f_*\mathcal{O}_X$ fournit, via la définition du spectre
relatif, un morphisme $can : X \to Y$ sur $S$,
voir Constructions, Lemme
\ref{constructions-lemma-canonical-morphism}. Par hypothèse et le lemme qui vient
d'être cité la restriction
$can|_{f^{-1}(W_j)} : f^{-1}(W_j) \to g^{-1}(W_j)$ est un isomorphisme. Ainsi $can$
est un isomorphisme. Nous avons montré
que (2) implique (3).

\medskip\noindent
Supposons (3). Par Constructions, Lemme \ref{constructions-lemma-spec-properties} nous voyons que l'image
inverse de tout ouvert affine est affine, et donc le morphisme
est affine par définition.
\end{proof}

\begin{remark}
\label{morphisms-remark-direct-argument}
Nous pouvons également affirmer directement
que (2) implique (1) dans le lemme \ref{morphisms-lemma-characterize-affine}
ci-dessus comme suit. Supposons que $S = \bigcup W_j$
soit un recouvrement ouvert
affine tel que chaque $f^{-1}(W_j)$ soit affine. Démontrons
d'abord que $\mathcal{A} = f_*\mathcal{O}_X$
est quasi-cohérent comme dans
la démonstration ci-dessus. Soit $\Spec(R) = V \subset S$
ouvert affine. Il faut montrer que $f^{-1}(V)$ est affine.
Posons $A = \mathcal{A}(V) = f_*\mathcal{O}_X(V) = \mathcal{O}_X(f^{-1}(V))$. Par Schémas, Lemme \ref{schemes-lemma-morphism-into-affine} il
existe un morphisme canonique $\psi : f^{-1}(V) \to \Spec(A)$ sur
$\Spec(R) = V$.
Par Schémas, Lemme \ref{schemes-lemma-good-subcover} il existe un entier $n \geq 0$,
un recouvrement standard ouvert
$V = \bigcup_{i = 1, \ldots, n} D(h_i)$, $h_i \in R$, et une application $a : \{1, \ldots, n\} \to J$
telle que chaque $D(h_i)$ est aussi un ouvert
principal du schéma affine $W_{a(i)}$. L'image réciproque
d'un principal ouvert sous un morphisme de schémas affines
est principal ouvert, voir Algèbre, Lemme \ref{algebra-lemma-spec-functorial}.
On voit donc que $f^{-1}(D(h_i))$ est un principe ouvert
de $f^{-1}(W_{a(i)})$, en particulier que $f^{-1}(D(h_i))$ est affine.
Parce que $\mathcal{A}$
est quasi-cohérent, nous avons $A_{h_i} = \mathcal{A}(D(h_i)) = \mathcal{O}_X(f^{-1}(D(h_i)))$, donc $f^{-1}(D(h_i))$
est le spectre de $A_{h_i}$.
Il s'ensuit que le morphisme $\psi$ induit un isomorphisme
de l'ouvert $f^{-1}(D(h_i))$ avec l'ouvert
$\Spec(A_{h_i})$ de $\Spec(A)$. Puisque $f^{-1}(V) = \bigcup f^{-1}(D(h_i))$
et $\Spec(A) = \bigcup \Spec(A_{h_i})$ on obtient le résultat.
\end{remark}

\begin{lemma}
\label{morphisms-lemma-affine-equivalence-algebras}
Soit $S$ un schéma. Il existe une anti-équivalence de catégories
$$
\begin{matrix}
\text{Schémas affines} \\
\text{sur }S
\end{matrix}
\longleftrightarrow
\begin{matrix}
\text{faisceaux quasi-cohérents} \\
\text{de }\mathcal{O}_S\text{-algèbres}
\end{matrix}
$$
qui associe à $f : X \to S$ le faisceau $f_*\mathcal{O}_X$. De plus, cette
équivalence est compatible avec un changement de base arbitraire.
\end{lemma}

\begin{proof}
Le foncteur de droite à gauche est donné par $\underline{\Spec}_S$.
Les deux foncteurs sont mutuellement
inverses par le lemme
\ref{morphisms-lemma-characterize-affine} et Constructions, Lemme \ref{constructions-lemma-spec-properties} partie (3).
L'instruction
finale est Constructions, Lemme \ref{constructions-lemma-spec-properties} partie (2).
\end{proof}

\begin{lemma}
\label{morphisms-lemma-quasi-coherence-scalar-restriction}
\begin{reference}
\cite[I, Proposition 9.6.1]{EGA}
\end{reference}
Soit $S$ un schéma et $\mathcal{A}$ une $\mathcal{O}_S$-algèbre
quasi-cohérente. Un $\mathcal{A}$-module est quasi-cohérent
comme un $\mathcal{O}_S$-module si et seulement si il
est quasi-cohérent comme un $\mathcal{A}$-module.
\end{lemma}

\begin{proof}
Soit $\mathcal{F}$ un module $\mathcal{A}$. Si $\mathcal{F}$
est quasi-cohérent en tant que module $\mathcal{A}$, alors pour chaque
$s \in S$ il existe un voisinage ouvert $U$ de $s$
et une suite exacte
$$
\bigoplus\nolimits_J\mathcal{A}|_U\to
\bigoplus\nolimits_I\mathcal{A}|_U\to
\mathcal{F}|_U
$$
de $\mathcal{A}|_U$-modules. Alors
c'est aussi une suite exacte de $\mathcal{O}_U$-modules. Ainsi
$\mathcal{F}|_U$ est quasi-cohérent comme le conoyau d'un morphisme de $\mathcal{O}_U$-modules
quasi-cohérents sur un schéma. Il s'ensuit que
$\mathcal{F}$ est quasi-cohérent en tant que $\mathcal{O}_X$-module.

\medskip\noindent
Inversement, supposons que $\mathcal{F}$ est quasi-cohérent en tant que module $\mathcal{O}_X$.
Choisissez un $\Spec(R) = U \subset S$ ouvert affine. Nous avons des isomorphismes
de $\mathcal{O}_U$-modules $\mathcal{A}|_U \cong \widetilde{A}$ et $\mathcal{F}|_U \cong \widetilde{M}$, pour certaines $R$-algèbre
$A$ et certains $R$-module $M$.
La structure de module $\mathcal{A}$ sur $\mathcal{F}$ se traduit par une structure
de module $A$ sur $M$ compatible avec la structure
de module $R$ donnée (détails omis). Choisissez une suite exacte
$$
\bigoplus\nolimits_J A \to
\bigoplus\nolimits_I A \to
M \to 0
$$
de $A$-modules. Puisque le foncteur $\widetilde{\ }$ est exact, cela
produit une suite exacte
$$
\bigoplus\nolimits_J \widetilde{A}\to
\bigoplus\nolimits_I \widetilde{A}\to
\widetilde{M} \to 0
$$
de $\widetilde{A}$-modules. Cela signifie que $\mathcal{F}$
est quasi-cohérent en tant que module $\mathcal{A}$.
\end{proof}

\begin{lemma}
\label{morphisms-lemma-affine-equivalence-modules}
Soit $f : X \to S$ un morphisme affine de
schémas. Soit $\mathcal{A} = f_*\mathcal{O}_X$.
Le foncteur $\mathcal{F} \mapsto f_*\mathcal{F}$ induit une équivalence
de catégories
$$
\left\{
\begin{matrix}
\text{catégorie des}\\
\mathcal{O}_X\text{-modules quasi-cohérents}
\end{matrix}
\right\}
\longrightarrow
\left\{
\begin{matrix}
\text{catégorie des}\\
\mathcal{A}\text{-modules quasi-cohérents}
\end{matrix}
\right\}
$$
De plus, un $\mathcal{A}$-module est quasi-cohérent
comme un $\mathcal{O}_S$-module si et seulement si il
est quasi-cohérent comme un $\mathcal{A}$-module.
\end{lemma}

\begin{proof}
L'instruction finale est Lemme \ref{morphisms-lemma-quasi-coherence-scalar-restriction}. Par le Lemme \ref{morphisms-lemma-affine-separated}
et Schémas, Lemme \ref{schemes-lemma-push-forward-quasi-coherent}
l'image directe $f_*\mathcal{F}$ d'un $\mathcal{O}_X$-module
quasi-cohérent $\mathcal{F}$ est un $\mathcal{O}_S$-module
quasi-cohérent.
D'où un foncteur comme dans
l'énoncé du lemme.

\medskip\noindent Nous
allons construire un $h$ quasi-inverse à ce foncteur. Soit $\mathcal{G}$ un
module $\mathcal{A}$ quasi-cohérent. Ensuite, nous définissons
$$
h(\mathcal{G}) = f^*\mathcal{G} \otimes_{f^*\mathcal{A}} \mathcal{O}_X =
f^*\mathcal{G} \otimes_{f^*f_*\mathcal{O}_X} \mathcal{O}_X
$$
Élucidation : le retrait $f^*\mathcal{A} = f^*f_*\mathcal{O}_X$ est une $\mathcal{O}_X$-algèbre,
l'application d'adjonction $f^*f_*\mathcal{O}_X \to \mathcal{O}_X$ est
un homomorphisme d'algèbre et le retrait $f^*\mathcal{G}$ est un $f^*\mathcal{A}$-module.
Observez que $h(\mathcal{G})$ est quasi-cohérent car la quasi-cohérence
est préservée par les retraits et le changement
d'anneaux. Observons qu'il existe un morphisme fonctoriel
$$
h(f_*\mathcal{F}) =
f^*f_*\mathcal{F} \otimes_{f^*f_*\mathcal{O}_X} \mathcal{O}_X \to \mathcal{F}
$$
provenant du morphisme d'adjonction $f^*f_*\mathcal{F} \to \mathcal{F}$ et un morphisme
fonctoriel
$$
\mathcal{G} \to f_*h(\mathcal{G}) \quad\text{adjoint du morphisme}\quad
f^*\mathcal{G} \to f^*\mathcal{G} \otimes_{f^*f_*\mathcal{O}_X} \mathcal{O}_X
$$
qui envoie une section locale $s$ de $f^*\mathcal{F}$ à $s \otimes 1$. Pour terminer la
démonstration il suffit de montrer que ces applications sont des isomorphismes pour
$\mathcal{F}$ et $\mathcal{G}$ comme ci-dessus. Ceci peut être vérifié sur les membres d'un revêtement
affine, c'est-à-dire lorsque $X$ et $S$ sont affines.

\medskip\noindent L'observation
algébrique clé qui fait ce travail est la suivante : Soit
$R \to A$ un homomorphisme d'anneaux. Soit $N$ un module $A$. Alors
$$
(N \otimes_R A) \otimes_{(A \otimes_R A)} A = N
$$
À savoir, le côté gauche de cette égalité est l'effet de l'application de
$h$ au module $\widetilde{A}$ quasi-cohérent $\widetilde{N}$ sur $\Spec(R)$. Nous omettons
les détails.
\end{proof}

\begin{lemma}
\label{morphisms-lemma-composition-affine}
La composition des morphismes affines est affine.
\end{lemma}

\begin{proof}
Soient $f : X \to Y$ et $g : Y \to Z$ des morphismes affines.
Soit $U \subset Z$ ouvert affine. Alors $g^{-1}(U)$ est affine
par hypothèse sur $g$. Sur quoi $f^{-1}(g^{-1}(U))$ est affine
par hypothèse sur $f$. Donc $(g \circ f)^{-1}(U)$ est affine.
\end{proof}

\begin{lemma}
\label{morphisms-lemma-base-change-affine}
Le changement de base d'un morphisme affine est affine.
\end{lemma}

\begin{proof}
Soit $f : X \to S$ un morphisme affine. Soit $S' \to S$ n'importe quel
morphisme. Notons $f' : X_{S'} = S' \times_S X \to S'$ le changement de base de $f$.
Pour chaque $s' \in S'$ il existe un voisinage ouvert affine
$s' \in V \subset S'$ qui s'envoie dans un certain ouvert affine $U \subset S$.
Par hypothèse, $f^{-1}(U)$ est affine. D'après le
matériel de Schémas, section \ref{schemes-section-fibre-products}, nous voyons que
$f^{-1}(U)_V = V \times_U f^{-1}(U)$ est affine et égal à $(f')^{-1}(V)$. Cela prouve
que $S'$ a un recouvrement ouvert par affines dont
l'image réciproque par $f'$ est affine. Nous concluons
par le lemme \ref{morphisms-lemma-characterize-affine} ci-dessus.
\end{proof}

\begin{lemma}
\label{morphisms-lemma-closed-immersion-affine}
Une immersion fermée est affine.
\end{lemma}

\begin{proof}
La première indication
en est Schémas, Lemme \ref{schemes-lemma-closed-immersion-affine-case}. Voir Schémas,
Lemme \ref{schemes-lemma-closed-subspace-scheme} pour un énoncé
complet.
\end{proof}

\begin{lemma}
\label{morphisms-lemma-affine-s-open}
Soit $X$ un
schéma. Soit $\mathcal{L}$ un module $\mathcal{O}_X$ inversible.
Soit $s \in \Gamma(X, \mathcal{L})$. Le morphisme
d'inclusion $j : X_s \to X$ est affine.
\end{lemma}

\begin{proof}
Cela résulte de Propriétés, Lemme \ref{properties-lemma-affine-cap-s-open} et la
définition.
\end{proof}

\begin{lemma}
\label{morphisms-lemma-affine-permanence}
Supposons que $g : X \to Y$ soit un morphisme de schémas sur $S$.
\begin{enumerate}
\item Si $X$ est affine sur $S$ et que $\Delta : Y \to Y \times_S Y$ est affine, alors $g$
est affine.
\item Si $X$ est affine sur $S$ et que $Y$ est séparé sur $S$, alors
$g$ est affine.
\item Un morphisme d'un schéma affine vers un schéma à diagonale
affine est affine.
\item Un morphisme d'un schéma affine vers un schéma séparé est affine.
\end{enumerate}
\end{lemma}

\begin{proof}
Démonstration de (1). Le changement de base $X \times_S Y \to Y$ est affine
par le lemme \ref{morphisms-lemma-base-change-affine}. Le morphisme
$(1, g) : X \to X \times_S Y$ est le changement de base de $Y \to Y \times_S Y$
par le morphisme $X \times_S Y \to Y \times_S Y$. Il est donc affine par le lemme
\ref{morphisms-lemma-base-change-affine}. La
composition des morphismes affines
est affine (voir Lemme \ref{morphisms-lemma-composition-affine})
et (1) s'ensuit. La partie (2) découle de (1) car une
immersion fermée est affine (voir Lemme \ref{morphisms-lemma-closed-immersion-affine})
et $Y/S$ séparé signifie $\Delta$ est une immersion
fermée. Les parties (3) et (4) sont des cas particuliers
de (1) et (2).
\end{proof}

\begin{lemma}
\label{morphisms-lemma-morphism-affines-affine}
Un morphisme entre schémas affines est affine.
\end{lemma}

\begin{proof}
Immédiat du lemme \ref{morphisms-lemma-affine-permanence} avec $S = \Spec(\mathbf{Z})$.
Cela découle également directement de l'équivalence
de (1) et (2) dans le lemme \ref{morphisms-lemma-characterize-affine}.
\end{proof}

\begin{lemma}
\label{morphisms-lemma-Artinian-affine}
Soit $S$ un schéma.
Soit $A$ un anneau artinien.
Tout morphisme $\Spec(A) \to S$ est affine.
\end{lemma}

\begin{proof}
Omis.
\end{proof}

\begin{lemma}
\label{morphisms-lemma-get-affine}
Soit $j : Y \to X$ une immersion de schémas. Supposons
qu'il existe un $U \subset X$ ouvert avec un complément
$Z = X \setminus U$ tel que
\begin{enumerate}
\item $U \to X$ est affine,
\item $j^{-1}(U) \to U$ est affine et
\item $j(Y) \cap Z$ est fermé.
\end{enumerate}
Alors $j$ est affine. En particulier, si $X$ est affine, $Y$ l'est également.
\end{lemma}

\begin{proof}
Par Schémas, Définition \ref{schemes-definition-immersion} il existe un sous-schéma ouvert $W \subset X$
tel que $j$ se factorise en une immersion fermée $i : Y \to W$
suivi du morphisme d'inclusion $W \to X$. Puisqu'une
immersion fermée est affine (Lemme
\ref{morphisms-lemma-closed-immersion-affine}), on voit que pour tout $x \in W$
il existe un voisinage ouvert affine de $x$
dans $X$ dont l'image réciproque par $j$ est affine.
Si $x \in U$, alors la même chose est vraie par l'hypothèse
(2). Enfin, supposons $x \in Z$ et $x \not \in W$. Puis $x \not \in j(Y) \cap Z$.
Par l'hypothèse (3) on peut trouver un voisinage affine
ouvert $V \subset X$ de $x$ qui ne satisfait pas
à $j(Y) \cap Z$. Puis $j^{-1}(V) = j^{-1}(V \cap U)$ qui est affine
par les hypothèses (1) et (2). Il s'ensuit que $j$
est affine par le lemme \ref{morphisms-lemma-characterize-affine}.
\end{proof}







\section{Familles de modules inversibles amples}
\label{morphisms-section-families-ample-invertible-modules}

\noindent
Une courte section sur la notion de famille de modules amples inversibles.

\begin{definition}
\label{morphisms-definition-family-ample-invertible-modules}
\begin{reference}
\cite[II Définition 2.2.4]{SGA6}
\end{reference}
Soit $X$ un schéma. Soit $\{\mathcal{L}_i\}_{i \in I}$ une famille de $\mathcal{O}_X$-modules
inversibles. Nous disons que $\{\mathcal{L}_i\}_{i \in I}$ est
une {\it grande famille
de modules inversibles sur $X$} si
\begin{enumerate}
\item $X$ est quasi-compact, et
\item pour chaque $x \in X$ il existe un $i \in I$, un $n \geq 1$
et $s \in \Gamma(X, \mathcal{L}_i^{\otimes n})$ de telle sorte que $x \in X_s$
et $X_s$ soient affines.
\end{enumerate}
\end{definition}

\noindent
Si $\{\mathcal{L}_i\}_{i \in I}$ est une ample famille de modules inversibles sur un
schéma $X$, alors il existe un sous-ensemble fini $I' \subset I$ tel que
$\{\mathcal{L}_i\}_{i \in I'}$ est une ample famille de modules inversibles sur $X$
(découle immédiatement de la quasi-compacité). Un schéma
ayant une ample famille de modules inversibles a une diagonale
affine par le lemme suivant et est donc a fortiori quasi-séparé.

\begin{lemma}
\label{morphisms-lemma-affine-diagonal}
Soit $X$ un schéma tel que pour tout point $x \in X$ il existe
un $\mathcal{O}_X$-module $\mathcal{L}$ inversible et une section
globale $s \in \Gamma(X, \mathcal{L})$ telle que $x \in X_s$ et $X_s$ sont affines.
Alors la diagonale de $X$ est un morphisme affine.
\end{lemma}

\begin{proof}
Étant donné les $\mathcal{O}_X$-modules $\mathcal{L}$, $\mathcal{M}$
et les sections globales $s \in \Gamma(X, \mathcal{L})$,
$t \in \Gamma(X, \mathcal{M})$ tels que $X_s$ et $X_t$ sont
affines, nous devons prouver que $X_s \cap X_t$
est affine. À savoir, alors
Lemme \ref{morphisms-lemma-characterize-affine} s'applique à $\Delta : X \to X \times X$
et le fait que $\Delta^{-1}(X_s \times X_t) = X_s \cap X_t$ montre que $\Delta$
est affine. Le fait que $X_s \cap X_t$ soit affine
découle de Propriétés, Lemme \ref{properties-lemma-affine-cap-s-open}.
\end{proof}

\begin{remark}
\label{morphisms-remark-affine-s-opens-cover-family}
Dans Propriétés, Lemme \ref{properties-lemma-affine-s-opens-cover-quasi-separated}, on voit qu'un schéma qui possède un module inversible ample
est séparé. Ceci est faux pour les schémas possédant
une famille ample de modules inversibles. En effet,
soit $X$ comme dans Schémas, Exemple \ref{schemes-example-affine-space-zero-doubled},
avec $n = 1$, c'est-à-dire la droite affine à zéro doublée.
Nous utilisons la notation de cet exemple sauf que nous écrivons $x$
pour $x_1$ et $y$ pour $y_1$. Il existe,
pour tout entier $n$, un faisceau inversible $\mathcal{L}_n$
sur $X$ qui est trivial sur $X_1$ et $X_2$
et dont la fonction de transition $U_{12} \to U_{21}$ est $f(x) \mapsto y^n f(y)$.
Les sections globales de $\mathcal{L}_n$ sont
des couples $(f(x), g(y)) \in k[x] \oplus
k[y]$ tels que $y^n f(y) = g(y)$. Les sections $s = (1,
y)$ de $\mathcal{L}_1$ et $t = (x, 1)$ de
$\mathcal{L}_{-1}$ déterminent un recouvrement par des ouverts affines, car
$X_s = X_1$ et $X_t = X_2$. Ainsi, $X$ possède
une famille ample de modules inversibles, mais il n'est pas séparé.
\end{remark}









\section{Morphismes quasi-affines}
\label{morphisms-section-quasi-affine}

\noindent
Rappelons qu'un schéma $X$ est dit {\it quasi-affine}
s'il est quasi-compact et isomorphe à un sous-schéma ouvert
d'un schéma affine, voir Propriétés, Définition \ref{properties-definition-quasi-affine}.

\begin{definition}
\label{morphisms-definition-quasi-affine}
Un morphisme de schémas $f : X \to S$ est appelé {\it quasi-affine} si
l'image réciproque de chaque ouvert affine de $S$ est un schéma quasi-affine.
\end{definition}

\begin{lemma}
\label{morphisms-lemma-quasi-affine-separated}
Un morphisme quasi-affine est séparé et quasi-compact.
\end{lemma}

\begin{proof}
Soit $f : X \to S$ quasi-affine. La
quasi-compacité est immédiate
chez Schémas, Lemme \ref{schemes-lemma-quasi-compact-affine}. Soit $U \subset S$ un
ouvert affine. Si l'on peut montrer que $f^{-1}(U)$
est un schéma séparé, alors $f$ est séparé
(Schémas, Lemme \ref{schemes-lemma-characterize-separated} montre que la séparation
est locale sur la base). Par hypothèse $f^{-1}(U)$
est isomorphe à un sous-schéma ouvert d'un schéma
affine. Un schéma affine est séparé et donc chaque sous-schéma
ouvert d'un schéma affine est séparé à volonté.
\end{proof}

\begin{lemma}
\label{morphisms-lemma-characterize-quasi-affine}
Soit $f : X \to S$ un morphisme de schémas. Les conditions
suivantes sont équivalentes
\begin{enumerate}
\item Le morphisme $f$ est quasi-affine.
\item Il existe un recouvrement ouvert affine $S = \bigcup W_j$ tel
que chaque $f^{-1}(W_j)$ soit quasi-affine.
\item Il existe un faisceau quasi-cohérent de $\mathcal{O}_S$-algèbres $\mathcal{A}$
et une immersion ouverte quasi-compacte
$$
\xymatrix{
X \ar[rr] \ar[rd] & & \underline{\Spec}_S(\mathcal{A}) \ar[dl] \\
& S &
}
$$
sur $S$.
\item Idem qu'en (3) mais avec $\mathcal{A} = f_*\mathcal{O}_X$ et
la flèche horizontale le morphisme
canonique des Constructions, Lemme \ref{constructions-lemma-canonical-morphism}.
\end{enumerate}
\end{lemma}

\begin{proof}
Il est évident que (1) implique (2) et que (4) implique (3).

\medskip\noindent
Supposons que $S = \bigcup_{j \in J} W_j$ soit un recouvrement ouvert
affine tel que chaque $f^{-1}(W_j)$
soit quasi-affine. Par Schémas, Lemme \ref{schemes-lemma-quasi-compact-affine}
on voit que $f$
est quasi-compact. Par Schémas, Lemme \ref{schemes-lemma-characterize-quasi-separated}
on voit que le morphisme $f$ est
quasi-séparé. Ainsi par Schémas, Lemme \ref{schemes-lemma-push-forward-quasi-coherent}
le faisceau $\mathcal{A} = f_*\mathcal{O}_X$ est un faisceau quasi-cohérent
de $\mathcal{O}_X$-algèbres. On a donc le schéma
$g : Y = \underline{\Spec}_S(\mathcal{A}) \to S$ sur $S$. L'application
d'identité
$\text{id} : \mathcal{A} = f_*\mathcal{O}_X \to f_*\mathcal{O}_X$ fournit, via la définition du spectre relatif,
un morphisme $can : X \to Y$ sur $S$,
voir Constructions, Lemme
\ref{constructions-lemma-canonical-morphism}. Par hypothèse, le lemme qui vient d'être
cité, et Propriétés, Lemme
\ref{properties-lemma-quasi-affine} la restriction $can|_{f^{-1}(W_j)} : f^{-1}(W_j) \to g^{-1}(W_j)$ est
une immersion ouverte quasi-compacte. Ainsi $can$
est une immersion ouverte quasi-compacte. Nous
avons montré que (2) implique (4).

\medskip\noindent
Supposons (3). Choisissez n'importe quel $U \subset S$
ouvert affine. Par Constructions, Lemme \ref{constructions-lemma-spec-properties} on voit que l'image
inverse de $U$ dans le spectre relatif est affine.
Nous concluons donc que $f^{-1}(U)$ est quasi-affine (notez
que la quasi-compacité est également codée dans (3)). Ainsi
(3) implique (1).
\end{proof}

\begin{lemma}
\label{morphisms-lemma-composition-quasi-affine}
La composition des morphismes quasi-affines est quasi-affine.
\end{lemma}

\begin{proof}
Soient $f : X \to Y$ et $g : Y \to Z$ des morphismes quasi-affines.
Soit $U \subset Z$ ouvert affine. Alors $g^{-1}(U)$
est quasi-affine par hypothèse sur $g$.
Soit $j : g^{-1}(U) \to V$ une immersion ouverte
quasi-compacte dans un schéma affine
$V$. Par le Lemme \ref{morphisms-lemma-characterize-quasi-affine}
ci-dessus on voit que $f^{-1}(g^{-1}(U))$ est un
sous-schéma ouvert quasi-compact du
spectre relatif $\underline{\Spec}_{g^{-1}(U)}(\mathcal{A})$ pour un faisceau quasi-cohérent
de $\mathcal{O}_{g^{-1}(U)}$-algèbres
$\mathcal{A}$. Par Schémas, Lemme \ref{schemes-lemma-push-forward-quasi-coherent}
le faisceau $\mathcal{A}' = j_*\mathcal{A}$
est un faisceau quasi-cohérent de $\mathcal{O}_V$-algèbres
avec la propriété que $j^*\mathcal{A}' = \mathcal{A}$. On obtient
donc un diagramme commutatif
$$
\xymatrix{
f^{-1}(g^{-1}(U)) \ar[r] &
\underline{\Spec}_{g^{-1}(U)}(\mathcal{A})
\ar[r] \ar[d] &
\underline{\Spec}_V(\mathcal{A}') \ar[d] \\
& g^{-1}(U) \ar[r]^j & V
}
$$
avec le carré étant un carré de
fibre, voir Constructions, Lemme \ref{constructions-lemma-spec-properties}. Notez que
le coin supérieur droit est un schéma affine.
Donc $(g \circ f)^{-1}(U)$ est quasi-affine.
\end{proof}

\begin{lemma}
\label{morphisms-lemma-base-change-quasi-affine}
Le changement de base d'un morphisme quasi-affine est quasi-affine.
\end{lemma}

\begin{proof}
Soit $f : X \to S$ un morphisme quasi-affine.
Par le Lemme \ref{morphisms-lemma-characterize-quasi-affine} ci-dessus on trouve un
faisceau quasi-cohérent de $\mathcal{O}_S$-algèbres
$\mathcal{A}$ et une immersion ouverte quasi-compacte
$X \to \underline{\Spec}_S(\mathcal{A})$ sur $S$. Soit $g : S' \to S$
n'importe
quel morphisme. Notons
$f' : X_{S'} = S' \times_S X \to S'$ le changement de base de $f$. Puisque le
changement de base d'une immersion ouverte quasi-compacte
est une immersion ouverte quasi-compacte
on voit que $X_{S'} \to \underline{\Spec}_{S'}(g^*\mathcal{A})$ est une immersion
ouverte quasi-compacte
(nous avons utilisé
Schémas, Lemmes \ref{schemes-lemma-quasi-compact-preserved-base-change} et \ref{schemes-lemma-base-change-immersion}
et Constructions,
Lemme \ref{constructions-lemma-spec-properties}). Par le Lemme \ref{morphisms-lemma-characterize-quasi-affine}
encore Nous concluons que $X_{S'} \to S'$
est quasi-affine.
\end{proof}

\begin{lemma}
\label{morphisms-lemma-quasi-compact-immersion-quasi-affine}
Une immersion quasi-compacte est quasi-affine.
\end{lemma}

\begin{proof}
Soit $X \to S$ une immersion quasi-compacte. Il faut montrer que l'image
l'image réciproque de tout ouvert affine est quasi-affine. Ainsi, en
supposant que $S$ est un schéma affine, il faut
montrer que $X$ est quasi-affine. Par le Lemme \ref{morphisms-lemma-quasi-compact-immersion} le morphisme
$X \to S$ se factorise en $X \to Z \to S$ où $Z$ est un sous-schéma fermé
de $S$ et $X \subset Z$ est un ouvert quasi-compact. Puisque
$S$ est affine, le Lemme \ref{morphisms-lemma-closed-immersion} implique que $Z$ est
affine, ce qui conclut.
\end{proof}

\begin{lemma}
\label{morphisms-lemma-affine-quasi-affine}
Soit $S$ un schéma. Soit $X$ un schéma affine.
Un morphisme $f : X \to S$ est quasi-affine si et seulement si il est quasi-compact.
En particulier, tout morphisme d'un schéma affine à un schéma quasi-séparé
est quasi-affine.
\end{lemma}

\begin{proof}
Soit $V \subset S$ un ouvert affine. Alors $f^{-1}(V)$ est un sous-schéma
ouvert du schéma affine $X$, donc quasi-affine si et
seulement si il est quasi-compact. Cela prouve la première affirmation.
La quasi-compacité de tout $f : X \to S$ où $X$ est affine
et $S$ quasi-séparé découle d'après Schémas, Lemme
\ref{schemes-lemma-quasi-compact-permanence} appliqué à $X \to S \to \Spec(\mathbf{Z})$.
\end{proof}

\begin{lemma}
\label{morphisms-lemma-quasi-affine-permanence}
Supposons que $g : X \to Y$ soit un morphisme de schémas sur $S$.
Si $X$ est quasi-affine sur $S$ et $Y$ est quasi-séparé sur
$S$, alors $g$ est quasi-affine. En particulier, tout morphisme
d'un schéma quasi-affine à un schéma quasi-séparé est quasi-affine.
\end{lemma}

\begin{proof}
Le changement de base $X \times_S Y \to Y$ est quasi-affine
par le lemme \ref{morphisms-lemma-base-change-quasi-affine}. Le morphisme
$X \to X \times_S Y$ est une immersion
quasi-compacte comme $Y \to S$ est quasi-séparé,
voir Schémas, Lemme \ref{schemes-lemma-section-immersion}. Une
immersion quasi-compacte est quasi-affine
par le lemme \ref{morphisms-lemma-quasi-compact-immersion-quasi-affine} et la composition
des morphismes quasi-affines est quasi-affine
(voir Lemme \ref{morphisms-lemma-composition-quasi-affine}). Ainsi on obtient le résultat.
\end{proof}











\section{Types de morphismes définis par des propriétés d'homomorphismes d'anneaux}
\label{morphisms-section-properties-ring-maps}

\noindent
Dans cette section nous étudions quelles propriétés des homomorphismes d'anneaux
permettent de définir des propriétés locales des morphismes de schémas.

\begin{definition}
\label{morphisms-definition-property-local}
Soit $P$ une propriété d'homomorphismes d'anneaux.
\begin{enumerate}
\item On dit que $P$ est {\it local} si ce qui suit est valable :
\begin{enumerate}
\item Pour tout homomorphisme d'anneaux $R \to A$, et
tout $f \in R$ on a $P(R \to A) \Rightarrow P(R_f \to A_f)$.
\item Pour tous les anneaux $R$, $A$, tout $f \in R$, $a\in A$
et tout homomorphisme d'anneaux $R_f \to A$ nous avons $P(R_f \to A) \Rightarrow P(R \to A_a)$.
\item Pour tout homomorphisme d'anneaux $R \to A$,
et $a_i \in A$
tel que $(a_1, \ldots, a_n) = A$ puis $\forall i, P(R \to A_{a_i}) \Rightarrow P(R \to A)$.
\end{enumerate}
\item On dit que $P$ est {\it stable sous changement
de base} si pour tout homomorphismes
d'anneaux $R \to A$, $R \to R'$ on a $P(R \to A) \Rightarrow P(R' \to R' \otimes_R A)$.
\item On dit que $P$ est {\it stable sous composition}
si pour tout homomorphismes
d'anneaux $A \to B$, $B \to C$ on a $P(A \to B) \wedge P(B \to C) \Rightarrow P(A \to C)$.
\end{enumerate}
\end{definition}

\begin{definition}
\label{morphisms-definition-locally-P}
Soit $P$ une propriété d'homomorphismes
d'anneaux. Soit $f : X \to S$ un morphisme de schémas. On dit
que $f$ est {\it localement de type $P$}
si pour tout $x \in X$ il existe un voisinage ouvert affine $U$
de $x$ dans $X$ qui s'envoie dans un voisinage ouvert affine
$V \subset S$ tel que l'homomorphisme d'anneaux induit $\mathcal{O}_S(V) \to \mathcal{O}_X(U)$ possède la
propriété $P$.
\end{definition}

\noindent
Il ne s'agit pas d'une ``bonne définition'' sauf si la propriété $P$
est une propriété locale. Même si $P$ est une propriété locale, nous n'utiliserons
pas automatiquement cette définition pour dire qu'un morphisme est ``localement
de type $P$'' à moins que nous n'indiquions également explicitement
la définition ailleurs.

\begin{lemma}
\label{morphisms-lemma-locally-P}
Soit $f : X \to S$ un morphisme de schémas. Soit $P$
une propriété des homomorphismes d'anneaux.
Soit $U$ un ouvert affine
de $X$, et $V$ un ouvert affine de
$S$ telle que
$f(U) \subset V$. Si $f$ est localement de type $P$
et que $P$ est local, alors $P(\mathcal{O}_S(V) \to \mathcal{O}_X(U))$ est valable.
\end{lemma}

\begin{proof}
Comme $f$ est localement de type $P$ pour
chaque $u \in U$, il existe un mappage $U_u \subset X$ ouvert affine
dans un $V_u \subset S$ ouvert affine tel que $P(\mathcal{O}_S(V_u) \to \mathcal{O}_X(U_u))$
est valable. Choisissez un voisinage ouvert $U'_u \subset U \cap U_u$
de $u$ qui est ouvert affine standard
dans $U$ et $U_u$, voir Schémas, Lemme
\ref{schemes-lemma-standard-open-two-affines}. Par Définition \ref{morphisms-definition-property-local} (1)(b)
nous voyons que $P(\mathcal{O}_S(V_u) \to \mathcal{O}_X(U'_u))$ est vrai. Nous pouvons donc
supposer que $U_u \subset U$ est un ouvert affine standard.
Choisissez un voisinage ouvert $V'_u \subset V \cap V_u$ de
$f(u)$ qui est ouvert affine standard dans $V$
et $V_u$, voir Schémas, Lemme \ref{schemes-lemma-standard-open-two-affines}.
Alors $U'_u = f^{-1}(V'_u) \cap U_u$ est un ouvert affine standard de $U_u$
(donc de $U$) et
on a $P(\mathcal{O}_S(V'_u) \to \mathcal{O}_X(U'_u))$ par définition \ref{morphisms-definition-property-local} (1)(a).
Par conséquent, nous pouvons supposer
que $U_u \subset U$ et $V_u \subset V$ sont tous deux affines
ouverts standard. En appliquant
la Définition \ref{morphisms-definition-property-local} (1)(b) une fois
de plus Nous concluons que $P(\mathcal{O}_S(V) \to \mathcal{O}_X(U_u))$ est valable. Parce que $U$
est quasi-compact, nous pouvons choisir un nombre
fini de points $u_1, \ldots, u_n \in U$ tel que
$$
U = U_{u_1} \cup \ldots \cup U_{u_n}.
$$
Par Définition \ref{morphisms-definition-property-local} (1)(c)
Nous concluons que $P(\mathcal{O}_S(V) \to \mathcal{O}_X(U))$ est vrai.
\end{proof}

\begin{lemma}
\label{morphisms-lemma-locally-P-characterize}
Soit $P$ une propriété locale des homomorphismes
d'anneaux. Soit $f : X \to S$ un morphisme de schémas. Les conditions
suivantes sont équivalentes
\begin{enumerate}
\item Le morphisme $f$ est localement de type $P$.
\item Pour chaque ouverture affine $U \subset X$,
$V \subset S$ avec $f(U) \subset V$ nous avons $P(\mathcal{O}_S(V) \to \mathcal{O}_X(U))$.
\item Il existe un recouvrement ouvert $S = \bigcup_{j \in J} V_j$ et des
recouvrements ouverts $f^{-1}(V_j) = \bigcup_{i \in I_j} U_i$ tels que chacun des
morphismes $U_i \to V_j$, $j\in J, i\in I_j$ est localement de
type $P$.
\item Il existe un recouvrement ouvert affine $S = \bigcup_{j \in J} V_j$
et un recouvrement ouvert affine $f^{-1}(V_j) = \bigcup_{i \in I_j} U_i$
tel que $P(\mathcal{O}_S(V_j) \to \mathcal{O}_X(U_i))$ est valable, pour tout
$j\in J, i\in I_j$.
\end{enumerate}
De plus, si $f$ est localement de type $P$
alors pour tout sous-schémas ouverts $U \subset X$, $V \subset S$ avec $f(U) \subset V$
la restriction $f|_U : U \to V$ est localement de type $P$.
\end{lemma}

\begin{proof}
Cela résulte du lemme \ref{morphisms-lemma-locally-P} ci-dessus.
\end{proof}

\begin{lemma}
\label{morphisms-lemma-composition-type-P}
Soit $P$ une propriété d'homomorphismes d'anneaux.
Supposons que $P$ soit local et stable en termes de composition.
La composition des morphismes localement de type $P$ est localement
de type $P$.
\end{lemma}

\begin{proof}
Soit $f : X \to Y$ et $g : Y \to Z$ des morphismes localement de type
$P$. Soit $x \in X$. Choisir un voisinage affine
ouvert $W \subset Z$ de $g(f(x))$. Choisir un voisinage affine
ouvert $V \subset g^{-1}(W)$ de $f(x)$. Choisir un voisinage affine
ouvert $U \subset f^{-1}(V)$ de $x$. Par le Lemme \ref{morphisms-lemma-locally-P-characterize} les
homomorphismes d'anneaux $\mathcal{O}_Z(W) \to \mathcal{O}_Y(V)$ et
$\mathcal{O}_Y(V) \to \mathcal{O}_X(U)$ satisfont $P$. Par conséquent,
$\mathcal{O}_Z(W) \to \mathcal{O}_X(U)$ satisfait à $P$ car $P$
est supposé stable sous composition.
\end{proof}

\begin{lemma}
\label{morphisms-lemma-base-change-type-P}
Soit $P$ une propriété d'homomorphismes d'anneaux.
Supposons que $P$ soit local et stable sous changement de base.
Le changement de base d'un morphisme localement de type $P$ est
localement de type $P$.
\end{lemma}

\begin{proof}
Soit $f : X \to S$ un morphisme localement de type $P$.
Soit $S' \to S$ n'importe quel
morphisme. Notons $f' : X_{S'} = S' \times_S X \to S'$ le changement de base de
$f$. Pour chaque $s' \in S'$ il existe un voisinage affine
ouvert $s' \in V' \subset S'$ qui s'envoie dans un ouvert affine $V \subset S$.
D'après le lemme \ref{morphisms-lemma-locally-P-characterize}, l'ouvert $f^{-1}(V)$ est une
réunion d'ouverts affines $U_i$ telle que les homomorphismes d'anneaux
$\mathcal{O}_S(V) \to \mathcal{O}_X(U_i)$ vérifient tous $P$.
D'après Schémas, section \ref{schemes-section-fibre-products},
$f^{-1}(U)_{V'} = V' \times_V f^{-1}(V)$ est la réunion
des ouverts affines $V' \times_V U_i$.
Comme $\mathcal{O}_{X_{S'}}(V' \times_V U_i) =
\mathcal{O}_{S'}(V') \otimes_{\mathcal{O}_S(V)} \mathcal{O}_X(U_i)$,
nous voyons que les homomorphismes d'anneaux
$\mathcal{O}_{S'}(V') \to \mathcal{O}_{X_{S'}}(V' \times_V U_i)$
vérifient $P$, puisque $P$ est supposée stable par changement de base.
\end{proof}

\begin{lemma}
\label{morphisms-lemma-properties-local}
Les propriétés suivantes d'un homomorphisme d'anneaux $R \to A$ sont locales.
\begin{enumerate}
\item (Isomorphisme sur anneaux locaux.)
Pour tout premier $\mathfrak q$ de $A$ couché sur $\mathfrak p \subset R$
l'homomorphisme d'anneaux $R \to A$
induit un isomorphisme $R_{\mathfrak p} \to A_{\mathfrak q}$.
\item (Immersion ouverte.)
Pour tout $\mathfrak q$ premier de $A$ il existe un
$f \in R$, $\varphi(f) \not \in \mathfrak q$ tel que l'homomorphisme d'anneaux $\varphi : R \to A$ induit
un isomorphisme $R_f \to A_f$.
\item (Fibres réduites.)
Pour chaque nombre premier $\mathfrak p$ de $R$,
l'anneau de fibres $A \otimes_R \kappa(\mathfrak p)$ est réduit.
\item (Fibres de dimension au plus $n$.)
Pour chaque nombre premier $\mathfrak p$ de $R$ l'anneau
de fibres $A \otimes_R \kappa(\mathfrak p)$ a une dimension de Krull au plus $n$.
\item (Localement noéthérien sur la cible.) L'homomorphisme
d'anneaux $R \to A$ a la propriété que $A$ est noéthérien.
\item Ajoutez-en davantage ici si nécessaire\footnote{Mais uniquement les propriétés
qui ne sont pas déjà traitées séparément ailleurs.}.
\end{enumerate}
\end{lemma}

\begin{proof}
Omis.
\end{proof}

\begin{lemma}
\label{morphisms-lemma-properties-base-change}
Les propriétés suivantes des homomorphismes d'anneaux sont stables sous changement de base.
\begin{enumerate}
\item (Isomorphisme sur anneaux locaux.)
Pour tout premier $\mathfrak q$ de $A$ couché sur $\mathfrak p \subset R$
l'homomorphisme d'anneaux $R \to A$
induit un isomorphisme $R_{\mathfrak p} \to A_{\mathfrak q}$.
\item (Immersion ouverte.)
Pour tout $\mathfrak q$ premier de $A$ il existe un
$f \in R$, $\varphi(f) \not \in \mathfrak q$ tel que l'homomorphisme d'anneaux $\varphi : R \to A$ induit
un isomorphisme $R_f \to A_f$.
\item Ajoutez-en davantage ici si nécessaire\footnote{Mais uniquement les propriétés
qui ne sont pas déjà traitées séparément ailleurs.}.
\end{enumerate}
\end{lemma}

\begin{proof}
Omis.
\end{proof}

\begin{lemma}
\label{morphisms-lemma-properties-composition}
Les propriétés suivantes des homomorphismes d'anneaux sont stables sous composition.
\begin{enumerate}
\item (Isomorphisme sur anneaux locaux.)
Pour tout premier $\mathfrak q$ de $A$ couché sur $\mathfrak p \subset R$
l'homomorphisme d'anneaux $R \to A$
induit un isomorphisme $R_{\mathfrak p} \to A_{\mathfrak q}$.
\item (Immersion ouverte.)
Pour tout $\mathfrak q$ premier de $A$ il existe un
$f \in R$, $\varphi(f) \not \in \mathfrak q$ tel que l'homomorphisme d'anneaux $\varphi : R \to A$ induit
un isomorphisme $R_f \to A_f$.
\item (Localement noethérien sur la cible.) L'homomorphisme
d'anneaux $R \to A$ a la propriété que $A$ est noéthérien.
\item Ajoutez-en davantage ici si nécessaire\footnote{Mais uniquement les propriétés
qui ne sont pas déjà traitées séparément ailleurs.}.
\end{enumerate}
\end{lemma}

\begin{proof}
Omis.
\end{proof}








\section{Morphismes de type fini}
\label{morphisms-section-finite-type}

\noindent
Rappelons qu'un homomorphisme d'anneaux $R \to A$ est dit de type
fini si $A$ est isomorphe à un quotient de $R[x_1, \ldots, x_n]$ comme une $R$-algèbre,
voir Algèbre, Définition \ref{algebra-definition-finite-type}.

\begin{definition}
\label{morphisms-definition-finite-type}
Soit $f : X \to S$ un morphisme de schémas.
\begin{enumerate}
\item On dit que $f$ est de type {\it fini
à $x \in X$} s'il existe un voisinage ouvert affine $\Spec(A) = U \subset X$
de $x$ et un ouvert affine $\Spec(R) = V \subset S$
à $f(U) \subset V$ tel que l'homomorphisme d'anneaux induit
$R \to A$ est de type fini.
\item On dit que $f$ est {\it localement de type fini}
s'il est de type fini en tout point de $X$.
\item On dit que $f$ est de type {\it fini} s'il est localement
de type fini et quasi-compact.
\end{enumerate}
\end{definition}

\begin{lemma}
\label{morphisms-lemma-locally-finite-type-characterize}
Soit $f : X \to S$ un morphisme de schémas. Les conditions
suivantes sont équivalentes
\begin{enumerate}
\item Le morphisme $f$ est localement de type fini.
\item Pour tous les ouverts affines $U \subset X$, $V \subset S$
avec $f(U) \subset V$ l'homomorphisme
d'anneaux $\mathcal{O}_S(V) \to \mathcal{O}_X(U)$ est de type fini.
\item Il existe un recouvrement ouvert $S = \bigcup_{j \in J} V_j$ et
des recouvrements ouverts $f^{-1}(V_j) = \bigcup_{i \in I_j} U_i$ tels que chacun
des morphismes $U_i \to V_j$, $j\in J, i\in I_j$ est localement
de type fini.
\item Il existe un recouvrement ouvert affine $S = \bigcup_{j \in J} V_j$
et des recouvrements ouverts affines $f^{-1}(V_j) = \bigcup_{i \in I_j} U_i$ tels
que l'homomorphisme d'anneaux $\mathcal{O}_S(V_j) \to \mathcal{O}_X(U_i)$ soit de
type fini, pour tout $j\in J, i\in I_j$.
\end{enumerate}
De plus, si $f$ est localement de type fini alors
pour tout sous-schémas ouverts $U \subset X$, $V \subset S$ avec $f(U) \subset V$
la restriction $f|_U : U \to V$ est localement de type fini.
\end{lemma}

\begin{proof}
Cela résulte du lemme \ref{morphisms-lemma-locally-P} si l'on montre que la propriété
``$R \to A$ est de type fini''
est locale. On vérifie les conditions (a), (b)
et (c) de la Définition \ref{morphisms-definition-property-local}.
Par Algèbre, Lemme \ref{algebra-lemma-base-change-finiteness}, la propriété d'être de type fini est
stable par changement de base ; on en déduit
la condition (a). Par Algèbre, Lemme
\ref{algebra-lemma-compose-finite-type}, elle est stable par composition,
et pour tout anneau $R$ l'homomorphisme
d'anneaux $R \to R_f$ est de type fini.
On en déduit la condition (b). Enfin, la
condition (c) résulte d'Algèbre, Lemme \ref{algebra-lemma-cover-upstairs}.
\end{proof}

\begin{lemma}
\label{morphisms-lemma-composition-finite-type}
La composition de deux morphismes qui sont localement de type fini est localement
de type fini. Il en va de même pour les morphismes de type fini.
\end{lemma}

\begin{proof}
Dans la démonstration du lemme \ref{morphisms-lemma-locally-finite-type-characterize} nous avons vu qu'être
de type fini est une propriété locale des homomorphismes
d'anneaux. D'où la première affirmation du lemme
découle du lemme \ref{morphisms-lemma-composition-type-P} combinée
au fait qu'être de type fini est une propriété des homomorphismes
d'anneaux qui est stable sous
composition, voir Algèbre, Lemme \ref{algebra-lemma-compose-finite-type}. Par
ce qui précède et le fait que les compositions
de morphismes quasi-compacts sont
quasi-compacts, voir Schémas, Lemme \ref{schemes-lemma-composition-quasi-compact} on voit
que la composition de morphismes de type fini est de
type fini.
\end{proof}

\begin{lemma}
\label{morphisms-lemma-base-change-finite-type}
Le changement de base d'un morphisme qui est localement de type fini
est localement de type fini. Il en va de même pour les morphismes de
type fini.
\end{lemma}

\begin{proof}
Dans la démonstration du lemme \ref{morphisms-lemma-locally-finite-type-characterize} nous avons vu qu'être
de type fini est une propriété locale des homomorphismes
d'anneaux. D'où le premier énoncé du lemme découle
du lemme \ref{morphisms-lemma-base-change-type-P} combiné au fait
qu'être de type fini est une propriété des homomorphismes d'anneaux
qui est stable par changement
de base, voir Algèbre, Lemme \ref{algebra-lemma-base-change-finiteness}. Par
ce qui précède et le fait qu'un changement de
base d'un morphisme quasi-compact est
quasi-compact, voir Schémas, Lemme \ref{schemes-lemma-quasi-compact-preserved-base-change} on voit que le changement
de base d'un morphisme de type fini est un morphisme
de type fini.
\end{proof}

\begin{lemma}
\label{morphisms-lemma-immersion-locally-finite-type}
Une immersion fermée est de type fini. Une
immersion est localement de type fini.
\end{lemma}

\begin{proof}
Ceci est vrai car une immersion ouverte est un isomorphisme local,
et une immersion fermée est évidemment de type fini.
\end{proof}

\begin{lemma}
\label{morphisms-lemma-finite-type-noetherian}
Soit $f : X \to S$ un morphisme. Si $S$
est (localement) noéthérien et $f$ (localement) de type fini alors
$X$ est (localement) noéthérien.
\end{lemma}

\begin{proof}
Cela découle immédiatement du fait qu'un anneau
de type fini sur un anneau noéthérien est noéthérien,
voir Algèbre, Lemme \ref{algebra-lemma-Noetherian-permanence}. (Aussi : utiliser
le fait que la source d'un morphisme quasi-compact à cible
quasi-compacte est quasi-compacte.)
\end{proof}

\begin{lemma}
\label{morphisms-lemma-finite-type-Noetherian-quasi-separated}
Soit $f : X \to S$ localement de type fini avec $S$ localement noethérien. Alors
$f$ est quasi-séparé.
\end{lemma}

\begin{proof}
En fait, il est vrai que $X$ est
quasi-séparé, voir Propriétés, Lemme \ref{properties-lemma-locally-Noetherian-quasi-separated} et Lemme
\ref{morphisms-lemma-finite-type-noetherian} ci-dessus. Appliquez ensuite
Schémas, Lemme \ref{schemes-lemma-compose-after-separated} pour conclure que
$f$ est quasi-séparé.
\end{proof}

\begin{lemma}
\label{morphisms-lemma-permanence-finite-type}
Soit $X \to Y$ un morphisme de schémas sur un schéma de base $S$. Si
$X$ est localement de type fini sur $S$, alors $X \to Y$
est localement de type fini.
\end{lemma}

\begin{proof}
Via Lemme \ref{morphisms-lemma-locally-finite-type-characterize} cela se traduit par le fait algébrique suivant
: Étant donné les
homomorphismes d'anneaux $A \to B \to C$ tels que $A \to C$
est de type fini, alors $B \to C$ est
de
type fini. (Voir Algèbre, Lemme \ref{algebra-lemma-compose-finite-type}).
\end{proof}







\section{Points de type fini et schémas de Jacobson}
\label{morphisms-section-points-finite-type}

\noindent
Soit $S$ un schéma. Un point de type fini $s$ de $S$
est un point tel que le morphisme $\Spec(\kappa(s)) \to S$ est de type fini. La raison
de cette étude est que les points de type fini peuvent remplacer les
points fermés dans un certain sens et dans certaines situations.
Il y en a toujours assez par exemple. De plus, un schéma est Jacobson
si et seulement si tous les points de type fini sont des points fermés.

\begin{lemma}
\label{morphisms-lemma-point-finite-type}
Soit $S$ un schéma. Soit $k$ un corps.
Soit $f : \Spec(k) \to S$ un morphisme. Les conditions
suivantes sont équivalentes :
\begin{enumerate}
\item Le morphisme $f$ est de type fini.
\item Le morphisme $f$ est localement de type fini.
\item Il existe un $U = \Spec(R)$ ouvert affine de $S$ tel que
$f$ correspond à un homomorphisme d'anneaux fini $R \to k$.
\item Il existe un $U = \Spec(R)$ ouvert affine de $S$ tel que
l'image de $f$ est constitué d'un point fermé $u$ dans $U$
et l'extension du corps $k/\kappa(u)$ est finie.
\end{enumerate}
\end{lemma}

\begin{proof}
L'équivalence de (1) et (2) est évidente car $\Spec(k)$ est un
singleton et donc tout morphisme qui en découle est quasi-compact.

\medskip\noindent
Supposons que $f$ soit localement de type fini.
Choisissez n'importe quel $\Spec(R) = U \subset S$ affine
ouvert tel que l'image de $f$ soit
contenue dans $U$, et l'homomorphisme d'anneaux
$R \to k$ soit de type fini. Soit $\mathfrak p \subset R$
le noyau. Alors $R/\mathfrak p \subset k$ est de type fini. En algèbre,
Lemme \ref{algebra-lemma-field-finite-type-over-domain} il existe un $\overline{f} \in R/\mathfrak p$ tel que
$(R/\mathfrak p)_{\overline{f}}$ est un corps et $(R/\mathfrak p)_{\overline{f}} \to k$ est
une extension de corps fini. Si $f \in R$ est un
lift de $\overline{f}$, alors nous voyons que $k$
est un $R_f$-module fini. Ainsi (2) $\Rightarrow$ (3).

\medskip\noindent
Supposons que $\Spec(R) = U \subset S$ soit un ouvert affine tel que
$f$ correspond à un homomorphisme d'anneaux fini
$R \to k$. Alors $f$ est
localement de type fini par le lemme \ref{morphisms-lemma-locally-finite-type-characterize}.
Ainsi (3) $\Rightarrow$ (2).

\medskip\noindent
Supposons que $R \to k$ soit fini. L'image de $R \to k$ est
un corps sur lequel $k$
est fini par l'Algèbre, Lemme \ref{algebra-lemma-integral-under-field}. Le
noyau de $R \to k$ est donc un idéal maximal. Ainsi
(3) $\Rightarrow$ (4).

\medskip\noindent
L'implication (4) $\Rightarrow$ (3) est immédiate.
\end{proof}

\begin{lemma}
\label{morphisms-lemma-artinian-finite-type}
Soit $S$ un schéma.
Soit $A$ un anneau local artinien à corps résiduel
$\kappa$. Soit $f : \Spec(A) \to S$ un morphisme de
schémas. Alors $f$ est de type fini si et seulement
si la composition $\Spec(\kappa) \to \Spec(A) \to S$ est de type fini.
\end{lemma}

\begin{proof}
Puisque le morphisme $\Spec(\kappa) \to \Spec(A)$ est de type fini, il
est clair que si $f$ est de type fini la composition
$\Spec(\kappa) \to S$ l'est également (voir Lemme \ref{morphisms-lemma-composition-finite-type}). A l'inverse,
notons que les applications $\Spec(A) \to S$ dans certains affines
ouverts $U = \Spec(B)$ de $S$ comme $\Spec(A)$ n'ont
qu'un seul point.
Pour finir, appliquez l'algèbre, Lemme \ref{algebra-lemma-essentially-of-finite-type-into-artinian-local} à
$B \to A$.
\end{proof}

\noindent
Rappelons que étant donné un point $s$ d'un schéma $S$
il existe un morphisme canonique $\Spec(\kappa(s)) \to S$,
voir Schémas, section \ref{schemes-section-points}.

\begin{definition}
\label{morphisms-definition-finite-type-point}
Soit $S$ un schéma.
Disons qu'un point $s$ de $S$ est un {\it
point de type fini} si le morphisme canonique $\Spec(\kappa(s)) \to S$ est de type
fini. On note $S_{\text{ft-pts}}$ l'ensemble des points de type fini $S$.
\end{definition}

\noindent
Nous pouvons décrire l'ensemble de points de type fini comme suit.

\begin{lemma}
\label{morphisms-lemma-identify-finite-type-points}
Soit $S$ un schéma. On a
$$
S_{\text{ft-pts}} = \bigcup\nolimits_{U \subset S\text{ ouvert}} U_0
$$
où $U_0$ est l'ensemble des points fermés de $U$. Ici, nous pouvons
faire parcourir à $U$ tous les ouverts, ou seulement tous les ouverts affines de $S$.
\end{lemma}

\begin{proof}
Immédiat du lemme \ref{morphisms-lemma-point-finite-type}.
\end{proof}

\begin{lemma}
\label{morphisms-lemma-finite-type-points-morphism}
Soit $f : T \to S$ un morphisme de schémas.
Si $f$ est localement
de type fini, alors $f(T_{\text{ft-pts}}) \subset S_{\text{ft-pts}}$.
\end{lemma}

\begin{proof}
Si $T$ est le spectre d'un corps, c'est Lemme \ref{morphisms-lemma-point-finite-type}. En général,
cela s'ensuit puisque la composition des morphismes localement
de type fini est localement de type fini (Lemme \ref{morphisms-lemma-composition-finite-type}).
\end{proof}

\begin{lemma}
\label{morphisms-lemma-finite-type-points-surjective-morphism}
Soit $f : T \to S$ un morphisme de schémas.
Si $f$ est localement de type fini
et surjectif, alors $f(T_{\text{ft-pts}}) = S_{\text{ft-pts}}$.
\end{lemma}

\begin{proof}
Nous avons $f(T_{\text{ft-pts}}) \subset S_{\text{ft-pts}}$ par le lemme \ref{morphisms-lemma-finite-type-points-morphism}. Soit
$s \in S$ un point de type fini.
Comme $f$ est surjectif le schéma $T_s = \Spec(\kappa(s)) \times_S T$ est
non vide, donc a un point de type fini $t \in T_s$
par le lemme \ref{morphisms-lemma-identify-finite-type-points}. Or
$T_s \to T$ est un morphisme de type fini
comme changement de base de $s \to S$ (Lemme
\ref{morphisms-lemma-base-change-finite-type}).
Ainsi l'image de $t$ dans
$T$ est un point de type fini par le
lemme \ref{morphisms-lemma-finite-type-points-morphism} qui s'applique à $s$
par construction.
\end{proof}

\begin{lemma}
\label{morphisms-lemma-enough-finite-type-points}
Soit $S$ un schéma.
Pour tout sous-ensemble localement fermé $T \subset S$ on a
$$
T \not = \emptyset
\Rightarrow
T \cap S_{\text{ft-pts}} \not = \emptyset.
$$
En particulier, pour tout sous-ensemble fermé $T \subset S$
on voit que $T \cap S_{\text{ft-pts}}$ est dense dans $T$.
\end{lemma}

\begin{proof}
Notons que $T$ porte une structure
schématique (voir Schémas, Lemme \ref{schemes-lemma-reduced-closed-subscheme}) telle
que $T \to S$ est une immersion localement
fermée. Toute immersion localement fermée
est localement de type fini, voir Lemme
\ref{morphisms-lemma-immersion-locally-finite-type}. Ainsi par le lemme \ref{morphisms-lemma-finite-type-points-morphism} nous
voyons $T_{\text{ft-pts}} \subset S_{\text{ft-pts}}$. Enfin, tout ouvert affine
non vide de $T$ possède au moins un point fermé
qui est un point de type fini de
$T$ par le lemme \ref{morphisms-lemma-identify-finite-type-points}.
\end{proof}

\noindent
Il s'ensuit que la plupart du matériel de
Topologie, section \ref{topology-section-space-jacobson} passe en revue avec l'ensemble des points
fermés remplacé par l'ensemble des points de type fini. En
fait, si $S$
est Jacobson alors on récupère les points fermés comme
points de type fini.

\begin{lemma}
\label{morphisms-lemma-jacobson-finite-type-points}
Soit $S$ un schéma. Les conditions suivantes sont équivalentes :
\begin{enumerate}
\item le schéma $S$ est Jacobson,
\item $S_{\text{ft-pts}}$ est l'ensemble des points fermés de $S$,
\item pour tout morphisme $T \to S$ localement de type fini, les points
fermés s'envoient sur des points fermés, et
\item pour tout morphisme $T \to S$ localement de type
fini, les points fermés $t \in T$ s'envoient sur des points fermés
$s \in S$ tels que l'extension $\kappa(s) \subset \kappa(t)$ soit finie.
\end{enumerate}
\end{lemma}

\begin{proof}
On a trivialement (4) $\Rightarrow$ (3) $\Rightarrow$ (2).
Lemme \ref{morphisms-lemma-enough-finite-type-points} montre que (2) implique (1). Il suffit donc de
montrer que (1) implique (4). Supposons
que $T \to S$ soit localement de type fini.
Choisissons un point fermé $t \in T$ et notons $s \in S$
son image. Choisissons des voisinages ouverts affines $\Spec(R) = U \subset S$
de $s$ et $\Spec(A) = V \subset T$ de $t$, tels que
$V$ s'envoie dans $U$. L'homomorphisme
d'anneaux induit $R \to A$ est de type fini (voir Lemme
\ref{morphisms-lemma-locally-finite-type-characterize}) et $R$ est Jacobson par Propriétés,
Lemme \ref{properties-lemma-locally-jacobson}. Ainsi
le résultat découle de l'Algèbre, Proposition \ref{algebra-proposition-Jacobson-permanence}.
\end{proof}

\begin{lemma}
\label{morphisms-lemma-Jacobson-universally-Jacobson}
Soit $S$ un Schéma de Jacobson.
Tout schéma localisé de type fini sur $S$ est Jacobson.
\end{lemma}

\begin{proof}
Cela ressort
clairement de l'algèbre, proposition
\ref{algebra-proposition-Jacobson-permanence} (et des propriétés, Lemme \ref{properties-lemma-locally-jacobson}
et Lemme \ref{morphisms-lemma-locally-finite-type-characterize}).
\end{proof}

\begin{lemma}
\label{morphisms-lemma-ubiquity-Jacobson-schemes}
Les types de schémas suivants sont de Jacobson.
\begin{enumerate}
\item Tout schéma local de type fini sur un corps.
\item Tout schéma localisé de type fini sur $\mathbf{Z}$.
\item Tout schéma local de type fini sur un domaine noéthérien $1$-dimensionnel
avec une infinité de nombres premiers.
\item Un schéma de la forme $\Spec(R) \setminus \{\mathfrak m\}$ où $(R, \mathfrak m)$
est un anneau noéthérien local. Également tout
schéma local de type fini dessus.
\end{enumerate}
\end{lemma}

\begin{proof}
Nous utiliserons Lemme \ref{morphisms-lemma-Jacobson-universally-Jacobson} sans mention. Le spectre
d'un corps est clairement
Jacobson. Le spectre de $\mathbf{Z}$
est Jacobson, voir Algèbre,
Lemme
\ref{algebra-lemma-pid-jacobson}. Pour (3) voir Algèbre, Lemme
\ref{algebra-lemma-noetherian-dim-1-Jacobson}.
Pour (4) voir Propriétés, Lemme \ref{properties-lemma-complement-closed-point-Jacobson}.
\end{proof}







\section{Schémas universellement caténaires}
\label{morphisms-section-universally-catenary}

\noindent
Rappelons qu'un espace topologique $X$ est appelé {\it caténaire}
si pour chaque paire de sous-ensembles fermés irréductibles $T \subset T'$
il existe une chaîne maximale de sous-ensembles fermés irréductibles
$$
T = T_0 \subset T_1 \subset \ldots \subset T_e = T'
$$
et que chaque chaîne a la même longueur.
Voir Topologie, Définition \ref{topology-definition-catenary}. Rappelons
qu'un schéma est caténaire si son espace topologique sous-jacent
est caténaire. Voir Propriétés, Définition \ref{properties-definition-catenary}.

\begin{definition}
\label{morphisms-definition-universally-catenary}
Soit $S$ un schéma. Supposons que $S$ soit localement noethérien.
On dit que $S$ est {\it universellement caténaire}
si pour tout morphisme $X \to S$ localement de type fini le schéma $X$ est caténaire.
\end{definition}

\noindent
C'est une ``meilleure'' notion que la caténaire car il existe
des schémas noéthériens qui sont caténaires mais pas
universellement caténaires. Voir Exemples, section \ref{examples-section-non-catenary-Noetherian-local}. De nombreux
schémas sont universellement caténaire,
voir Lemme \ref{morphisms-lemma-ubiquity-uc} ci-dessous.

\medskip\noindent Rappelons
qu'un anneau $A$ est appelé {\it
caténaire} si pour toute paire d'idéaux premiers $\mathfrak p \subset \mathfrak q$ il existe
une chaîne maximale de nombres premiers
$$
\mathfrak p =
\mathfrak p_0 \subset \ldots \subset \mathfrak p_e
= \mathfrak q
$$
et tous ont la même longueur. Voir Algèbre,
Définition \ref{algebra-definition-catenary}. Nous avons vu la relation entre les
schémas caténaires et les anneaux caténaires dans Propriétés,
section \ref{properties-section-catenary}. Rappelons qu'un anneau $A$ est appelé {\it
universellement caténaire} si $A$
est noéthérien et pour tout type fini homomorphisme d'anneaux
$A \to B$ l'anneau $B$
est caténaire. Voir Algèbre, Définition \ref{algebra-definition-universally-catenary}. De nombreux anneaux
intéressants rencontrés en géométrie
algébrique satisfont à cette propriété.

\begin{lemma}
\label{morphisms-lemma-universally-catenary-local}
Soit $S$ un schéma localement noethérien. Les conditions suivantes sont équivalentes
\begin{enumerate}
\item $S$ est universellement caténaire,
\item il existe un recouvrement ouvert de $S$ dont tous les membres sont des
schémas universellement caténaires,
\item pour chaque ouvert affine $\Spec(R) = U \subset S$ l'anneau $R$ est
universellement caténaire, et
\item il existe un recouvrement ouvert affine $S = \bigcup U_i$ tel que chaque
$U_i$ soit le spectre d'un anneau universellement caténaire.
\end{enumerate}
De plus, dans ce cas tout schéma localisé de type fini sur $S$ est également
universellement caténaire.
\end{lemma}

\begin{proof}
Par le Lemme \ref{morphisms-lemma-immersion-locally-finite-type} une immersion ouverte est localement de type
fini. Une composition de morphismes localement de type fini est
localement de type fini (Lemme \ref{morphisms-lemma-composition-finite-type}).
Il est donc clair que si $S$ est universellement caténaire
alors tout ouvert et tout schéma localisé de type fini sur $S$
est également universellement caténaire. Cela prouve l'énoncé
final du lemme et que (1) implique (2).

\medskip\noindent Si
$\Spec(R)$ est un schéma universellement caténaire, alors tout schéma
$\Spec(A)$ avec $A$ un type fini $R$-algèbre
est caténaire. Donc tous ces anneaux $A$ sont caténaires
par Algèbre, Lemme \ref{algebra-lemma-catenary}. Ainsi $R$
est universellement caténaire. Combiné avec les remarques ci-dessus Nous concluons
que (1) implique (3), et (2) implique (4). Bien sûr, (3) implique
(4) de manière triviale.

\medskip\noindent
Pour terminer la démonstration nous montrons
que (4) implique (1). Supposons (4) et soit $X \to S$ un morphisme
localement de type fini. On peut trouver un recouvrement
recouvrement par des ouverts affines $X = \bigcup V_j$ tel que chaque $V_j \to S$
ait son image dans l'un des $U_i$. Par le Lemme
\ref{morphisms-lemma-locally-finite-type-characterize} l'homomorphisme d'anneaux induit $\mathcal{O}(U_i) \to \mathcal{O}(V_j)$ est
de type fini. Donc $\mathcal{O}(V_j)$ est une caténaire. Par
conséquent, $X$ est caténaire par Propriétés, Lemme \ref{properties-lemma-catenary-local}.
\end{proof}

\begin{lemma}
\label{morphisms-lemma-universally-catenary-local-rings-universally-catenary}
Soit $S$ un schéma localement noethérien. Les conditions
suivantes sont équivalentes :
\begin{enumerate}
\item $S$ est universellement caténaire, et
\item tous les anneaux locaux $\mathcal{O}_{S, s}$ de $S$ sont universellement caténaire.
\end{enumerate}
\end{lemma}

\begin{proof}
Supposons que tous les anneaux locaux de $S$ soient universellement
caténaires. Soit $f : X \to S$ localement
de type fini. On sait que $X$ est caténaire si et
seulement si $\mathcal{O}_{X, x}$ est caténaire pour tous $x \in X$. Si $f(x) = s$,
alors $\mathcal{O}_{X, x}$ est essentiellement de type fini sur $\mathcal{O}_{S, s}$.
Donc $\mathcal{O}_{X, x}$ est caténaire en supposant que
$\mathcal{O}_{S, s}$ est universellement caténaire.

\medskip\noindent A
l'inverse, supposons que $S$ soit universellement caténaire. Soit
$s \in S$. On peut remplacer $S$ par un voisinage affine
ouvert de $s$ par le lemme \ref{morphisms-lemma-universally-catenary-local}. Disons que $S = \Spec(R)$
et $s$ correspond à l'idéal premier $\mathfrak p$. Tout type
fini $R_{\mathfrak p}$-algèbre $A'$ est de la forme
$A_{\mathfrak p}$ pour certains types fini $R$-algèbre $A$.
Par hypothèse (et Lemme \ref{morphisms-lemma-universally-catenary-local} si vous préférez) l'anneau $A$
est caténaire, et donc $A'$ (une localisation de $A$) est
caténaire. Ainsi $R_{\mathfrak p}$ est universellement caténaire.
\end{proof}

\begin{lemma}
\label{morphisms-lemma-catenary-check-irreducible}
Soit $S$ un schéma localement noethérien. Alors $S$ est universellement caténaire
si et seulement si les composantes irréductibles de $S$ sont universellement caténaire.
\end{lemma}

\begin{proof}
Omis. Pour le cas affine, veuillez
consulter Algèbre, Lemme \ref{algebra-lemma-catenary-check-irreducible}.
\end{proof}

\begin{lemma}
\label{morphisms-lemma-ubiquity-uc}
Les types de schémas suivants sont universellement caténaire.
\begin{enumerate}
\item Tout schéma local de type fini sur un corps.
\item Tout schéma localisé de type fini sur un schéma De Cohen--Macaulay.
\item Tout schéma localisé de type fini sur $\mathbf{Z}$.
\item Tout schéma localisé de type fini sur un domaine noéthérien
$1$-dimensionnel.
\item Et ainsi de suite.
\end{enumerate}
\end{lemma}

\begin{proof}
Tout cela découle du fait qu'un anneau
de De Cohen--Macaulay est universellement
caténaire, voir Algèbre, Lemme \ref{algebra-lemma-CM-ring-catenary}.
Utilisez également la dernière
assertion du lemme \ref{morphisms-lemma-universally-catenary-local}. Certains
détails omis.
\end{proof}























\section{Schémas de Nagata, reprise}
\label{morphisms-section-nagata}

\noindent
Voir Propriétés, section \ref{properties-section-nagata} pour les définitions et propriétés
de base des schémas Nagata et universellement japonais.

\begin{lemma}
\label{morphisms-lemma-finite-type-nagata}
Soit $f : X \to S$ un morphisme. Si $S$
est Nagata et $f$ localement de type fini alors $X$ est Nagata. Si $S$
est universellement japonais et
$f$ localement de type fini alors $X$ est universellement japonais.
\end{lemma}

\begin{proof}
Pour ``universellement japonais''
Cela résulte de Algèbre, Lemme \ref{algebra-lemma-universally-japanese}.
Pour ``Nagata''
Cela résulte de Algèbre, Proposition \ref{algebra-proposition-nagata-universally-japanese}.
\end{proof}

\begin{lemma}
\label{morphisms-lemma-ubiquity-nagata}
Les types de schémas suivants sont des Nagata.
\begin{enumerate}
\item Tout schéma local de type fini sur un corps.
\item Tout schéma local de type fini sur un anneau local complet noéthérien.
\item Tout schéma localisé de type fini sur $\mathbf{Z}$.
\item Tout schéma localement de type fini sur un anneau de Dedekind de
caractéristique zéro.
\item Et ainsi de suite.
\end{enumerate}
\end{lemma}

\begin{proof}
Par le Lemme \ref{morphisms-lemma-finite-type-nagata} il suffit de montrer que les anneaux
mentionnés ci-dessus sont des anneaux Nagata.
Pour cela, voir Algèbre, Proposition \ref{algebra-proposition-ubiquity-nagata}.
\end{proof}


\section{Le lieu singulier, reprise}
\label{morphisms-section-singular-locus}

\noindent
Nous recherchons un critère qui implique l'ouverture du lieu régulier
pour tout schéma local de type fini sur la base. Voici la définition.

\begin{definition}
\label{morphisms-definition-J}
Soit $X$ un schéma localement noethérien. On dit que $X$ est {\it
J-2} si pour tout morphisme $Y \to X$ qui est localement de type
fini le locus régulier $\text{Reg}(Y)$ est ouvert dans $Y$.
\end{definition}

\noindent
C'est l'analogue de la notion correspondante pour les anneaux noéthériens,
voir Compléments d'algèbre, Définition \ref{more-algebra-definition-J}.

\begin{lemma}
\label{morphisms-lemma-J}
Soit $X$ un schéma localement noethérien. Les conditions suivantes sont équivalentes
\begin{enumerate}
\item $X$ est J-2,
\item il existe un recouvrement ouvert de $X$ dont tous les membres sont des
schémas J-2,
\item pour chaque ouvert affine $\Spec(R) = U \subset X$ l'anneau $R$
est J-2, et
\item il existe un recouvrement ouvert affine $S = \bigcup U_i$ tel que chaque
$\mathcal{O}(U_i)$ soit J-2 pour tous les $i$.
\end{enumerate}
De plus, dans ce cas, tout schéma localisé de type fini sur $X$ est
également J-2.
\end{lemma}

\begin{proof}
Par le Lemme \ref{morphisms-lemma-immersion-locally-finite-type} une immersion ouverte est localement de type
fini. Une composition de morphismes localement de type fini
est localement de type fini (Lemme
\ref{morphisms-lemma-composition-finite-type}). Il est donc clair que si $X$ est J-2 alors tout
ouvert et tout schéma localisé de type fini sur
$X$ est également J-2. Cela prouve l'énoncé
final du lemme.

\medskip\noindent Si
$\Spec(R)$ est J-2, alors pour tout type fini $R$-algèbre $A$
le lieu régulier du schéma $\Spec(A)$ est ouvert. Donc $R$
est J-2, par définition
(voir Compléments d'algèbre, Définition \ref{more-algebra-definition-J}).
Combiné avec les remarques ci-dessus Nous concluons que (1) implique (3),
et (2) implique (4). Bien sûr (1) $\Rightarrow$ (2) et
(3) $\Rightarrow$ (4) trivialement.

\medskip\noindent
Pour terminer la démonstration nous montrons
que (4) implique (1). Supposons (4) et soit $Y \to X$
un morphisme localement de type fini. On peut trouver un
recouvrement par des ouverts affines $Y = \bigcup V_j$ tel
que chaque $V_j \to X$ ait son image dans l'un des $U_i$.
Par le Lemme \ref{morphisms-lemma-locally-finite-type-characterize}, l'homomorphisme d'anneaux
induit $\mathcal{O}(U_i) \to \mathcal{O}(V_j)$ est de type fini.
Le lieu régulier de $V_j = \Spec(\mathcal{O}(V_j))$ est donc
ouvert. Puisque $\text{Reg}(Y) \cap V_j = \text{Reg}(V_j)$ Nous concluons que $\text{Reg}(Y)$
est ouvert comme souhaité.
\end{proof}

\begin{lemma}
\label{morphisms-lemma-ubiquity-J-2}
Les types de schémas suivants sont J-2.
\begin{enumerate}
\item Tout schéma local de type fini sur un corps.
\item Tout schéma local de type fini sur un anneau local complet noéthérien.
\item Tout schéma localisé de type fini sur $\mathbf{Z}$.
\item Tout schéma local de type fini sur un anneau noéthérien local de dimension
$1$.
\item Tout schéma local de type fini sur un anneau Nagata de dimension $1$.
\item Tout schéma localisé de type fini sur un anneau de Dedekind de caractéristique
zéro.
\item Et ainsi de suite.
\end{enumerate}
\end{lemma}

\begin{proof}
Par le Lemme \ref{morphisms-lemma-J} il suffit de montrer que
les anneaux mentionnés ci-dessus sont J-2. Pour
cela, voir Compléments d'algèbre, Proposition \ref{more-algebra-proposition-ubiquity-J-2}.
\end{proof}




\section{Schémas excellents}
\label{morphisms-section-excellent}

\noindent
On rappelle qu'un anneau est quasi-excellent s'il est un G-anneau et J-2.
Un anneau est excellent s'il est quasi-excellent et universellement caténaire.

\begin{definition}
\label{morphisms-definition-excellent}
Soit $X$ un schéma.
\begin{enumerate}
\item On dit que $X$ est {\it quasi-excellent}
si pour tout $x \in X$ il existe un voisinage ouvert affine $x \in U \subset X$
tel que l'anneau $\mathcal{O}_X(U)$ est quasi-excellent
(voir Compléments d'algèbre, Définition \ref{more-algebra-definition-excellent}).
\item On dit que $X$ est {\it excellent} si
pour tout $x \in X$ il existe un voisinage ouvert affine $x \in U \subset X$
tel que l'anneau $\mathcal{O}_X(U)$ est
excellent (voir Compléments d'algèbre, Définition \ref{more-algebra-definition-excellent}).
\end{enumerate}
\end{definition}

\noindent
Les schémas quasi-excellents sont localement noethériens (les anneaux G sont noethériens).

\begin{lemma}
\label{morphisms-lemma-characterize-quasi-excellent}
Soit $X$ un schéma. Les conditions suivantes sont équivalentes
\begin{enumerate}
\item $X$ est quasi-excellent, et
\item $X$ est un G-schéma et J-2.
\end{enumerate}
\end{lemma}

\begin{proof}
L'implication (1) $\Rightarrow$ (2) résulte de la combinaison
de la définition des quasi-schémas
excellents (Définition \ref{morphisms-definition-excellent}),
de la définition des anneaux quasi-excellents
(Compléments d'algèbre, Définition \ref{more-algebra-definition-excellent}), de
la définition des G-schémas
(Propriétés, Définition \ref{properties-definition-G-scheme}) et du lemme \ref{morphisms-lemma-J}.
Inversement, si
$X$ est un G-schéma et J-2, alors pour tout $x \in X$
on peut trouver un voisinage ouvert affine $x \in U \subset X$
tel que $\mathcal{O}_X(U)$ soit un G-anneau. Par le lemme \ref{morphisms-lemma-J},
l'anneau $\mathcal{O}_X(U)$ est également J-2, donc quasi-excellent.
\end{proof}

\begin{lemma}
\label{morphisms-lemma-quasi-excellent-nagata}
Un quasi-schéma excellent est Nagata.
\end{lemma}

\begin{proof}
Compléments d'algèbre, Lemme \ref{more-algebra-lemma-quasi-excellent-nagata}.
\end{proof}

\begin{lemma}
\label{morphisms-lemma-characterize-excellent}
Soit $X$ un schéma. Les conditions suivantes sont équivalentes
\begin{enumerate}
\item $X$ est excellent, et
\item $X$ est quasi-excellent et universellement caténaire.
\end{enumerate}
\end{lemma}

\begin{proof}
Supposons (1). Puisque les anneaux excellents sont quasi-excellents,
il est clair que $X$ est quasi-excellent. Puisque les anneaux
excellents sont universellement caténaire, il s'ensuit également
que $X$ est universellement caténaire par le lemme
\ref{morphisms-lemma-universally-catenary-local}. Supposons (2). Soit $x \in X$. Puisque $X$
est quasi-excellent il existe un voisinage
ouvert affine $x \in U \subset X$ tel que $\mathcal{O}_X(U)$ est quasi-excellent.
Puisque $X$ est universellement caténaire, l'anneau
$\mathcal{O}_X(U)$ est aussi universellement caténaire d'après le lemme \ref{morphisms-lemma-universally-catenary-local},
donc excellent.
\end{proof}

\begin{lemma}
\label{morphisms-lemma-locally-excellent}
Soit $X$ un schéma. Les conditions suivantes sont équivalentes :
\begin{enumerate}
\item Le schéma $X$ est (quasi-)excellent.
\item Pour chaque $U \subset X$ ouvert affine, l'anneau $\mathcal{O}_X(U)$ est
(quasi-)excellent.
\item Il existe un recouvrement ouvert affine $X = \bigcup U_i$ tel que
chaque $\mathcal{O}_X(U_i)$ soit (quasi-)excellent.
\item Il existe un recouvrement ouvert $X = \bigcup X_j$ tel que
chaque sous-schéma ouvert $X_j$ est (quasi-)excellent.
\end{enumerate}
De plus, si $X$ est (quasi-)excellent alors tout sous-schéma ouvert
est (quasi-)excellent.
\end{lemma}

\begin{proof}
Par le Lemme \ref{morphisms-lemma-characterize-quasi-excellent} être quasi-excellent
équivaut à être un G-schéma et J-2. Ainsi
pour le cas quasi-excellent, le lemme
découle du lemme
\ref{morphisms-lemma-J} et Propriétés, Lemme \ref{properties-lemma-locally-G-scheme}.
Pour le cas excellent, utilisez Lemme \ref{morphisms-lemma-characterize-excellent},
le cas quasi-excellent, et Lemme \ref{morphisms-lemma-universally-catenary-local}.
\end{proof}

\begin{lemma}
\label{morphisms-lemma-finite-type-excellent}
Soit $f : X \to S$ un morphisme. Si $S$
est (quasi-)excellent et $f$ localement de type fini alors $X$
est (quasi-)excellent.
\end{lemma}

\begin{proof}
Compléments d'algèbre, Lemme \ref{more-algebra-lemma-finite-type-over-excellent}.
\end{proof}

\begin{lemma}
\label{morphisms-lemma-ubiquity-excellent}
Les types de schémas suivants sont excellents.
\begin{enumerate}
\item Tout schéma local de type fini sur un corps.
\item Tout schéma local de type fini sur un anneau local complet noéthérien.
\item Tout schéma localisé de type fini sur $\mathbf{Z}$.
\item Tout schéma localement de type fini sur un anneau de Dedekind de
caractéristique zéro.
\end{enumerate}
\end{lemma}

\begin{proof}
Par les Lemmes \ref{morphisms-lemma-locally-excellent} et \ref{morphisms-lemma-finite-type-excellent} il suffit de montrer que
les anneaux mentionnés ci-dessus sont excellents.
Pour cela,
voir Compléments d'algèbre, Proposition \ref{more-algebra-proposition-ubiquity-excellent}.
\end{proof}






\section{Morphismes quasi-finis}
\label{morphisms-section-quasi-finite}

\noindent
Un traitement solide des morphismes quasi-finis est à la base de nombreux développements
ultérieurs. Cela conduira à différentes versions du Théorème principal
de Zariski, comportement des dimensions des fibres, descente pour morphismes
étales, etc, etc. Avant de lire cette section, il peut être judicieux
de jeter un œil aux résultats d'algèbre dans Algèbre, section \ref{algebra-section-quasi-finite}.

\medskip\noindent
Rappelons qu'un type fini homomorphisme d'anneaux $R \to A$
est quasi-fini en un nombre premier $\mathfrak q$ si $\mathfrak q$
définit un point isolé de sa fibre, voir Algèbre, Définition \ref{algebra-definition-quasi-finite}.

\begin{definition}
\label{morphisms-definition-quasi-finite}
\begin{reference}
\cite[II Définition 6.2.3]{EGA}
\end{reference}
Soit $f : X \to S$ un morphisme de schémas.
\begin{enumerate}
\item On dit que $f$ est {\it quasi-fini en un point
$x \in X$} s'il existe un voisinage affine $\Spec(A) = U \subset X$
de $x$ et un ouvert affine $\Spec(R) = V \subset S$ tel que $f(U) \subset V$,
l'homomorphisme d'anneaux $R \to A$ est de type fini, et $R \to A$
est quasi-fini au premier de $A$ correspondant à $x$ (voir
ci-dessus).
\item On dit que $f$ est {\it localement quasi-fini} si $f$
est quasi-fini en tout point $x$ de $X$.
\item On dit que $f$ est {\it quasi-fini} si $f$ est
de type fini et chaque point $x$ est un point isolé de sa fibre.
\end{enumerate}
\end{definition}

\noindent
Trivialement, un morphisme quasi-fini localement est localement de type fini. Nous
verrons plus loin qu'un morphisme $f$ qui est localement de type fini
est quasi-fini à $x$ si et seulement si $x$ est isolé dans sa
fibre. De plus, l'ensemble des points auxquels un morphisme est quasi-fini est ouvert
; nous le verrons dans la section \ref{morphisms-section-Zariski} sur le Théorème principal de Zariski.

\begin{lemma}
\label{morphisms-lemma-algebraic-residue-field-extension-closed-point-fibre}
Soit $f : X \to S$ un morphisme de schémas. Soit
$x \in X$ un point. Posons $s = f(x)$.
Si $\kappa(x)/\kappa(s)$ est une
extension de corps algébrique, alors
\begin{enumerate}
\item $x$ est un point fermé de sa fibre, et
\item si en plus $s$ est un point fermé de $S$, alors $x$
est un point fermé de $X$.
\end{enumerate}
\end{lemma}

\begin{proof}
La deuxième instruction découle de la première par topologie élémentaire.
D'après Schémas, Lemme \ref{schemes-lemma-fibre-topological} pour prouver la première
affirmation on peut
remplacer $X$ par $X_s$ et $S$ par
$\Spec(\kappa(s))$. On peut donc supposer que $S = \Spec(k)$ est le spectre
d'un corps. Dans ce cas, soit $\Spec(A) = U \subset X$ un ouvert affine quelconque
contenant $x$. Le point $x$ correspond
à un idéal premier $\mathfrak q \subset A$ tel que $\kappa(\mathfrak q)/k$ est
une extension algébrique de corps.
D'après Algèbre, Lemme \ref{algebra-lemma-finite-residue-extension-closed}, $\mathfrak q$ est
un idéal maximal, c'est-à-dire que $x \in U$ est un point fermé.
Puisque les ouverts affines forment
une base de la topologie de $X$, nous concluons que $\{x\}$ est fermé.
\end{proof}

\noindent
Le lemme suivant est une version du Hilbert Nullstellensatz.

\begin{lemma}
\label{morphisms-lemma-closed-point-fibre-locally-finite-type}
Soit $f : X \to S$ un morphisme de schémas. Soit $x \in X$
un point. Posons $s = f(x)$. Supposons
que $f$ soit localement de type
fini. Alors $x$ est un point fermé de
sa fibre si et seulement si $\kappa(x)/\kappa(s)$ est
une extension finie de corps.
\end{lemma}

\begin{proof}
Si l'extension est finie, alors $x$ est
un point
fermé de la fibre par le lemme \ref{morphisms-lemma-algebraic-residue-field-extension-closed-point-fibre} ci-dessus. À l'inverse,
supposons que $x$ soit un point
fermé de sa fibre. Choisissons des ouverts affines
$\Spec(A) = U \subset X$ et $\Spec(R) = V \subset S$ tels que $f(U) \subset V$.
Par le Lemme \ref{morphisms-lemma-locally-finite-type-characterize} l'homomorphisme d'anneaux $R \to A$
est de type fini. Soit $\mathfrak q \subset A$, resp.\ $\mathfrak p \subset R$
l'idéal premier correspondant à
$x$, resp.\ $s$.
Considérez l'anneau de fibre $\overline{A} = A \otimes_R \kappa(\mathfrak p)$.
Soit $\overline{\mathfrak q}$ le premier de $\overline{A}$
correspondant à $\mathfrak q$. L'hypothèse selon
laquelle $x$ est un point fermé de sa fibre
implique que $\overline{\mathfrak q}$ est un idéal maximal
de $\overline{A}$. Puisque $\overline{A}$ est une algèbre
de type fini sur le corps
$\kappa(\mathfrak p)$ nous voyons par le Hilbert Nullstellensatz,
voir Algèbre, Théorème \ref{algebra-theorem-nullstellensatz},
que $\kappa(\overline{\mathfrak q})$
est une extension finie de $\kappa(\mathfrak p)$.
Puisque $\kappa(s) = \kappa(\mathfrak p)$ et $\kappa(x) = \kappa(\mathfrak q) = \kappa(\overline{\mathfrak q})$ nous
gagnons.
\end{proof}

\begin{lemma}
\label{morphisms-lemma-base-change-closed-point-fibre-locally-finite-type}
Soit $f : X \to S$ un morphisme de schémas qui est localement de type fini. Soit $g : S' \to S$
un morphisme quelconque. Notons $f' : X' \to S'$ le changement de base. Si $x' \in X'$
s'envoie sur un point $x \in X$ qui est fermé dans $X_{f(x)}$, alors
$x'$ est fermé dans $X'_{f'(x')}$.
\end{lemma}

\begin{proof}
Le corps résiduel $\kappa(x')$ est un quotient
de $\kappa(f'(x')) \otimes_{\kappa(f(x))} \kappa(x)$, voir Schémas, Lemme
\ref{schemes-lemma-points-fibre-product}. Il s'agit donc d'une extension finie
de $\kappa(f'(x'))$ comme $\kappa(x)$ est une
extension finie de $\kappa(f(x))$ par le lemme
\ref{morphisms-lemma-closed-point-fibre-locally-finite-type}. On voit donc que $x'$ est
fermé dans sa fibre en appliquant une nouvelle fois
ce lemme.
\end{proof}

\begin{lemma}
\label{morphisms-lemma-residue-field-quasi-finite}
Soit $f : X \to S$ un morphisme de schémas. Soit
$x \in X$ un point. Posons $s = f(x)$. Si
$f$ est quasi-fini en $x$, alors l'extension
du corps résiduel $\kappa(x)/\kappa(s)$ est finie.
\end{lemma}

\begin{proof}
Cela ressort clairement de l'algèbre, Définition \ref{algebra-definition-quasi-finite}.
\end{proof}

\begin{lemma}
\label{morphisms-lemma-quasi-finite-at-point-characterize}
Soit $f : X \to S$ un morphisme de schémas. Soit $x \in X$
un point. Posons $s = f(x)$. Soit $X_s$ la fibre
de $f$ en $s$. Supposons que $f$
soit localement de type fini. Les conditions
suivantes sont équivalentes :
\begin{enumerate}
\item Le morphisme $f$ est quasi-fini en $x$.
\item Le point $x$ est isolé dans $X_s$.
\item Le point $x$ est fermé dans $X_s$
et il n'y a aucun point $x' \in X_s$, $x' \not = x$ qui se
spécialise dans $x$.
\item Pour toute paire d'affines
ouvertes $\Spec(A) = U \subset X$, $\Spec(R) = V \subset S$ avec $f(U) \subset V$
et $x \in U$ correspondant à $\mathfrak q \subset A$ l'homomorphisme d'anneaux
$R \to A$ est quasi-fini en $\mathfrak q$.
\end{enumerate}
\end{lemma}

\begin{proof}
Supposons que $f$ soit quasi-fini à $x$. Par hypothèse il existe
des ouvertures $U \subset X$, $V \subset S$ telles que $f(U) \subset V$, $x \in U$
et $x$ un point isolé de $U_s$. Par conséquent, $\{x\} \subset U_s$ est un sous-ensemble
ouvert. Puisque $U_s = U \cap X_s \subset X_s$ est également ouvert, nous concluons que $\{x\} \subset X_s$
est également un sous-ensemble ouvert. Ainsi Nous concluons que $x$
est un point isolé de $X_s$.

\medskip\noindent Notez
que $X_s$ est un Schéma de Jacobson
du lemme \ref{morphisms-lemma-ubiquity-Jacobson-schemes} (et Lemme \ref{morphisms-lemma-base-change-finite-type}).
Si
$x$ est isolé dans $X_s$, c'est-à-dire
que $\{x\} \subset X_s$ est ouvert, alors $\{x\}$ contient un point
fermé (par la propriété Jacobson), donc $x$ est fermé dans $X_s$.
Il est clair qu'il n'y a aucun intérêt pour $x' \in X_s$, distinct de $x$,
de se spécialiser dans $x$.

\medskip\noindent
Supposons que $x$ est fermé dans $X_s$ et qu'il n'y
a aucun point $x' \in X_s$, distinct de $x$, spécialisé dans
$x$. Considérons une paire d'ouverts affines
$\Spec(A) = U \subset X$, $\Spec(R) = V \subset S$ avec $f(U) \subset V$ et $x \in U$. Soit
$\mathfrak q \subset A$ correspond à $x$ et $\mathfrak p \subset R$
correspond à $s$. Par le Lemme \ref{morphisms-lemma-locally-finite-type-characterize} l'homomorphisme
d'anneaux $R \to A$ est de type fini.
Considérez l'anneau de fibre $\overline{A} = A \otimes_R \kappa(\mathfrak p)$. Soit
$\overline{\mathfrak q}$ le premier de $\overline{A}$ correspondant à $\mathfrak q$.
Puisque $\Spec(\overline{A})$ est un sous-schéma ouvert de la fibre $X_s$
on voit que $\overline{q}$ est un idéal maximal
de $\overline{A}$ et qu'il ne sert à rien de spécialiser $\Spec(\overline{A})$
sur $\overline{\mathfrak q}$. Cela implique
que $\dim(\overline{A}_{\overline{q}}) = 0$. Ainsi par Algèbre, Définition
\ref{algebra-definition-quasi-finite}
nous voyons que $R \to A$ est quasi-fini
à $\mathfrak q$, c'est-à-dire que $X \to S$ est quasi-fini
à $x$ par définition.

\medskip\noindent À
ce stade, nous avons montré que les conditions (1) à (3) sont toutes équivalentes.
Il est clair que (4) implique (1). Et il est clair aussi que (2) implique
(4) puisque si $x$ est un point isolé de $X_s$ alors
c'est aussi un point isolé de $U_s$ pour tout $U$ ouvert
qui le contient.
\end{proof}

\begin{lemma}
\label{morphisms-lemma-finite-fibre}
Soit $f : X \to S$ un morphisme de schémas. Soit
$s \in S$. Supposons que
\begin{enumerate}
\item $f$ est localement de type fini et que
\item $f^{-1}(\{s\})$ est un ensemble fini.
\end{enumerate}
Alors $X_s$ est un espace topologique fini discret, et $f$
est quasi-fini en chaque point de $X$ recouvrant $s$.
\end{lemma}

\begin{proof}
Supposons que $T$ soit un schéma qui (a) est localement
de type fini sur un corps $k$, et (b) comporte
un nombre fini de points. Alors Lemme \ref{morphisms-lemma-ubiquity-Jacobson-schemes} montre
que $T$ est un Schéma de Jacobson. Un espace
de Jacobson fini est discret, voir Topologie, Lemme
\ref{topology-lemma-finite-jacobson}. Appliquez cette remarque à la fibre $X_s$ qui est localement
de type fini sur $\Spec(\kappa(s))$ pour voir la première instruction.
Enfin, appliquez Lemme \ref{morphisms-lemma-quasi-finite-at-point-characterize} pour voir la seconde.
\end{proof}

\begin{lemma}
\label{morphisms-lemma-locally-quasi-finite-fibres}
\begin{slogan}
Les morphismes de type fini avec des fibres discrètes sont quasi-finis.
\end{slogan}
Soit $f : X \to S$ un morphisme de schémas. Supposons que $f$
soit localement de type fini. Alors Les
conditions suivantes sont équivalentes
\begin{enumerate}
\item $f$ est localement quasi-fini,
\item pour tout $s \in S$ la fibre $X_s$ est un espace topologique discret, et
\item pour tout morphisme $\Spec(k) \to S$ où $k$ est un corps le changement
de base $X_k$ a un espace topologique discret sous-jacent.
\end{enumerate}
\end{lemma}

\begin{proof}
Il est immédiat que (3) implique
(2). Lemme \ref{morphisms-lemma-quasi-finite-at-point-characterize} montre que (2) est équivalent
à (1). Supposons (2)
et soit $\Spec(k) \to S$ comme dans (3).
Notons $s \in S$ l'image de $\Spec(k) \to S$.
Alors $X_k$ est le changement
de base de $X_s$ via $\Spec(k) \to \Spec(\kappa(s))$.
Ainsi chaque point
de $X_k$ est fermé par le lemme \ref{morphisms-lemma-base-change-closed-point-fibre-locally-finite-type}. Comme $X_k \to \Spec(k)$
est localement de type fini (par
le lemme \ref{morphisms-lemma-base-change-finite-type}), nous pouvons
appliquer
Lemme \ref{morphisms-lemma-quasi-finite-at-point-characterize} pour conclure que chaque point de
$X_k$ est isolé, c'est-à-dire que $X_k$ a un espace
topologique discret sous-jacent.
\end{proof}

\begin{lemma}
\label{morphisms-lemma-quasi-finite-locally-quasi-compact}
Soit $f : X \to S$ un morphisme de schémas. Alors $f$
est quasi-fini si et seulement si $f$
est localement quasi-fini et quasi-compact.
\end{lemma}

\begin{proof}
Supposons que $f$ soit quasi-fini. Il est quasi-compact
par Définition \ref{morphisms-definition-finite-type}. Soit $x \in X$.
On voit que $f$ est quasi-fini
en $x$ par le lemme \ref{morphisms-lemma-quasi-finite-at-point-characterize}. Donc $f$
est quasi-compact et localement quasi-fini.

\medskip\noindent
Supposons que $f$ soit quasi-compact et localement
quasi-fini. Alors $f$ est de type fini.
Soit $x \in X$ un point. Par le Lemme \ref{morphisms-lemma-quasi-finite-at-point-characterize} on voit
que $x$ est un point isolé de sa fibre. Le
lemme est prouvé.
\end{proof}

\begin{lemma}
\label{morphisms-lemma-quasi-finite}
Soit $f : X \to S$ un morphisme de schémas. Les conditions
suivantes sont équivalentes :
\begin{enumerate}
\item $f$ est quasi-fini, et
\item $f$ est localement de type fini, quasi-compact, et possède des fibres finies.
\end{enumerate}
\end{lemma}

\begin{proof}
Supposons que $f$ soit quasi-fini. En particulier $f$ est localement
de type fini et quasi-compact (puisqu'il est de type fini). Soit $s \in S$.
Puisque chaque $x \in X_s$ est isolé dans $X_s$ on
voit que $X_s = \bigcup_{x \in X_s} \{x\}$ est un recouvrement ouvert. Comme $f$ est
quasi-compact, la fibre $X_s$ est quasi-compacte. On voit donc que
$X_s$ est fini.

\medskip\noindent
À l'inverse, supposons que $f$ est localement de type
fini, quasi-compact et possède des fibres finies. Alors
il est localement quasi-fini par le lemme \ref{morphisms-lemma-finite-fibre}.
Il est donc quasi-fini par le lemme \ref{morphisms-lemma-quasi-finite-locally-quasi-compact}.
\end{proof}

\noindent
Rappelons qu'un homomorphisme d'anneaux $R \to A$ est quasi-fini
s'il est de type fini et quasi-fini à {\it tous les}
premiers de $A$, voir Algèbre, Définition \ref{algebra-definition-quasi-finite}.

\begin{lemma}
\label{morphisms-lemma-locally-quasi-finite-characterize}
Soit $f : X \to S$ un morphisme de schémas. Les conditions
suivantes sont équivalentes
\begin{enumerate}
\item Le morphisme $f$ est localement quasi-fini.
\item Pour toute paire d'ouverts affines $U \subset X$, $V \subset S$
avec $f(U) \subset V$ l'homomorphisme
d'anneaux $\mathcal{O}_S(V) \to \mathcal{O}_X(U)$ est quasi-fini.
\item Il existe un recouvrement ouvert $S = \bigcup_{j \in J} V_j$ et
des recouvrements ouverts $f^{-1}(V_j) = \bigcup_{i \in I_j} U_i$ tels que chacun
des morphismes $U_i \to V_j$, $j\in J, i\in I_j$ est localement
quasi-fini.
\item Il existe un recouvrement ouvert affine $S = \bigcup_{j \in J} V_j$
et des recouvrements ouverts affines $f^{-1}(V_j) = \bigcup_{i \in I_j} U_i$ tels
que l'homomorphisme d'anneaux $\mathcal{O}_S(V_j) \to \mathcal{O}_X(U_i)$ est quasi-fini,
pour tout $j\in J, i\in I_j$.
\end{enumerate}
De plus, si $f$ est localement quasi-fini alors
pour tout sous-schémas ouverts $U \subset X$, $V \subset S$ avec $f(U) \subset V$
la restriction $f|_U : U \to V$ est localement quasi-finie.
\end{lemma}

\begin{proof}
Pour un homomorphisme d'anneaux $R \to A$
définissons $P(R \to A)$ comme signifiant que ``$R \to A$
est quasi-fini''
(voir remarque ci-dessus lemme). Nous affirmons
que $P$ est une propriété locale des homomorphismes
d'anneaux. On vérifie les conditions (a), (b) et
(c) de la Définition \ref{morphisms-definition-property-local}. Dans la démonstration
du lemme \ref{morphisms-lemma-locally-finite-type-characterize} nous avons vu que (a), (b) et (c) sont
valables pour la propriété d'être ``de type fini''. A noter
que, pour un homomorphisme d'anneaux de type fini $R \to A$, la propriété
$R \to A$ est quasi-fini en $\mathfrak q$ ne dépend que
de l'anneau local $A_{\mathfrak q}$ comme une algèbre sur $R_{\mathfrak p}$ où $\mathfrak p = R \cap \mathfrak q$
(abus de notation habituel). À l'aide de ces remarques (a), (b) et
(c) de la Définition \ref{morphisms-definition-property-local} suivent immédiatement. Par exemple,
supposons que $R \to A$ soit un homomorphisme
d'anneaux tel que tous les homomorphismes d'anneaux $R \to A_{a_i}$
soient quasi-finis pour $a_1, \ldots, a_n \in A$ générant l'idéal
unitaire. Nous concluons que $R \to A$ est de type
fini. De plus, pour tout $\mathfrak q \subset A$ premier, l'anneau local $A_{\mathfrak q}$
est isomorphe en tant qu'algèbre $R$
à l'anneau local $(A_{a_i})_{\mathfrak q_i}$ pour certains $i$
et certains $\mathfrak q_i \subset A_{a_i}$. D'où Nous concluons que
$R \to A$ est quasi-fini en $\mathfrak q$.

\medskip\noindent Nous
concluons que Lemme \ref{morphisms-lemma-locally-P} s'applique à $P$ comme dans
le paragraphe précédent. Il
suffit donc de prouver que $f$ est localement quasi-fini équivaut
à $f$ est localement de type $P$. Puisque $P(R \to A)$
est ``$R \to A$ est quasi-fini'' ce qui
signifie que $R \to A$ est quasi-fini à chaque nombre
premier de $A$, cela résulte du lemme \ref{morphisms-lemma-quasi-finite-at-point-characterize}.
\end{proof}

\begin{lemma}
\label{morphisms-lemma-composition-quasi-finite}
La composition de deux morphismes localement quasi-finis est localement
quasi-finie. Il en va de même pour les morphismes quasi-finis.
\end{lemma}

\begin{proof}
Dans la démonstration du lemme \ref{morphisms-lemma-locally-quasi-finite-characterize} nous avons vu que $P = $``quasi-fini''
est une propriété locale des
homomorphismes d'anneaux, et qu'un morphisme de schémas est localement
quasi-fini si et seulement si il est localement de type
$P$ comme dans Définition \ref{morphisms-definition-locally-P}. D'où
le premier énoncé du lemme découle du lemme
\ref{morphisms-lemma-composition-type-P} combiné au fait qu'être quasi-fini est une propriété
des homomorphismes d'anneaux
qui est stable sous composition, voir Algèbre, Lemme \ref{algebra-lemma-quasi-finite-composition}.
Par ce qui précède, Lemme \ref{morphisms-lemma-quasi-finite-locally-quasi-compact} et le fait que
les compositions de morphismes
quasi-compacts sont quasi-compacts, voir
Schémas, Lemme \ref{schemes-lemma-composition-quasi-compact} on voit que la composition
de morphismes quasi-finis est
quasi-fini.
\end{proof}

\noindent
Nous verrons plus loin (Lemme \ref{morphisms-lemma-quasi-finite-points-open}) que l'ensemble
$U$ du lemme suivant est ouvert.

\begin{lemma}
\label{morphisms-lemma-base-change-quasi-finite}
\begin{slogan}
(Localement) les morphismes quasi-finis sont stables sous changement de base.
\end{slogan}
Soit $f : X \to S$ un morphisme de schémas. Soit $g : S' \to S$
un morphisme de schémas. Notons $f' : X' \to S'$
le changement de base de $f$ par $g$ et notons
$g' : X' \to X$ la projection. Supposons que $X$
soit localement de type fini sur $S$.
\begin{enumerate}
\item Soit $U \subset X$ (resp.\ $U' \subset X'$)
l'ensemble des points où $f$ (resp.\ $f'$)
est quasi-fini. Puis $U' = U \times_S S' = (g')^{-1}(U)$.
\item Le changement de base d'un morphisme localement quasi-fini est
localement quasi-fini.
\item Le changement de base d'un morphisme quasi-fini est
quasi-fini.
\end{enumerate}
\end{lemma}

\begin{proof}
La première et la deuxième assertion découlent du
résultat algébrique
correspondant, voir Algèbre, Lemme \ref{algebra-lemma-quasi-finite-base-change}
(combiné au fait que $f'$ est également localement
de type fini par le lemme \ref{morphisms-lemma-base-change-finite-type}).
Par ce qui précède, Lemme \ref{morphisms-lemma-quasi-finite-locally-quasi-compact} et le fait
qu'un changement de base d'un
morphisme quasi-compact est quasi-compact,
voir Schémas, Lemme \ref{schemes-lemma-quasi-compact-preserved-base-change} on voit que le changement
de base d'un morphisme quasi-fini est
quasi-fini.
\end{proof}

\begin{lemma}
\label{morphisms-lemma-quasi-finite-at-a-finite-number-of-points}
Soit $f : X \to S$ un morphisme de schémas de type fini. Soit $s \in S$.
Il y a au plus un nombre fini de points de $X$ situés
sur $s$ auxquels $f$ est quasi-fini.
\end{lemma}

\begin{proof}
La fibre $X_s$ est un schéma de type fini
sur un corps, donc noéthérien (Lemme \ref{morphisms-lemma-finite-type-noetherian}).
Ainsi la topologie sur $X_s$ est noéthérienne
(Propriétés, Lemme \ref{properties-lemma-Noetherian-topology})
et peut avoir au plus un nombre fini de points isolés (par
topologie élémentaire). Ainsi
notre lemme découle du lemme \ref{morphisms-lemma-quasi-finite-at-point-characterize}.
\end{proof}

\begin{lemma}
\label{morphisms-lemma-monomorphism-loc-finite-type-loc-quasi-finite}
Soit $f : X \to Y$ un morphisme de schémas. Si $f$
est localement de type fini et un monomorphisme, alors $f$
est séparé et localement quasi-fini.
\end{lemma}

\begin{proof}
Un monomorphisme est séparé par Schémas,
Lemme \ref{schemes-lemma-monomorphism-separated}. Un monomorphisme
est injectif, on obtient donc que
$f$ est quasi-fini à chaque $x \in X$
par exemple par le lemme \ref{morphisms-lemma-quasi-finite-at-point-characterize}.
\end{proof}

\begin{lemma}
\label{morphisms-lemma-immersion-locally-quasi-finite}
Toute immersion est localement quasi-finie.
\end{lemma}

\begin{proof}
Cela est vrai car une immersion ouverte est un isomorphisme local et
une immersion fermée est clairement quasi-finie.
\end{proof}

\begin{lemma}
\label{morphisms-lemma-permanence-quasi-finite}
Soit $X \to Y$ un morphisme de schémas sur un schéma de base $S$. Soit
$x \in X$. Si $X \to S$ est quasi-fini à $x$, alors $X \to Y$
est quasi-fini à $x$. Si $X$
est localement quasi-fini sur $S$, alors $X \to Y$
est localement quasi-fini.
\end{lemma}

\begin{proof}
Via Lemme \ref{morphisms-lemma-locally-quasi-finite-characterize} cela se traduit par le fait algébrique suivant
: Étant donné des homomorphismes
d'anneaux $A \to B \to C$ tels que $A \to C$ est
quasi-fini, alors $B \to C$ est quasi-fini.
Cela résulte
de Algèbre, Lemme \ref{algebra-lemma-four-rings} avec $R = A$,
$S = S' = C$ et $R' = B$.
\end{proof}

\begin{lemma}
\label{morphisms-lemma-quasi-finite-local-source}
Soient $f : X \to Y$ et $g : Y \to S$ des morphismes de schémas. Si $f$
est surjectif, $g \circ f$ localement quasi-fini, et $g$
localement de type fini, alors $g : Y \to S$ est localement quasi-fini.
\end{lemma}

\begin{proof}
Soit $x \in X$ avec les images $y \in Y$ et
$s \in S$. Puisque $g \circ f$ est localement
quasi-fini par le lemme \ref{morphisms-lemma-residue-field-quasi-finite}
l'extension $\kappa(x)/\kappa(s)$ est finie.
Donc $\kappa(y)/\kappa(s)$ est fini. Donc $y$
est un point fermé de $Y_s$
par le lemme \ref{morphisms-lemma-algebraic-residue-field-extension-closed-point-fibre}. Puisque $f$ est surjectif,
on voit que chaque point de $Y$ est fermé
dans sa fibre sur $S$.
Ainsi par le lemme \ref{morphisms-lemma-quasi-finite-at-point-characterize} Nous concluons
que $g$ est quasi-fini en tout point.
\end{proof}














\section{Morphismes de présentation finie}
\label{morphisms-section-finite-presentation}

\noindent
Rappelons qu'un homomorphisme d'anneaux $R \to A$ est de présentation
finie si $A$ est isomorphe à $R[x_1, \ldots, x_n]/(f_1, \ldots, f_m)$ comme une
$R$-algèbre pour certains $n, m$ et certains polynômes
$f_j$, voir Algèbre, Définition \ref{algebra-definition-finite-type}.

\begin{definition}
\label{morphisms-definition-finite-presentation}
Soit $f : X \to S$ un morphisme de schémas.
\begin{enumerate}
\item On dit que $f$ est de {\it présentation finie
à $x \in X$} s'il existe un voisinage ouvert affine $\Spec(A) = U \subset X$
de $x$ et ouvert affine $\Spec(R) = V \subset S$ avec
$f(U) \subset V$ tel que l'homomorphisme d'anneaux induit
$R \to A$ est de présentation finie.
\item On dit que $f$ est {\it localement de présentation finie}
s'il est de présentation finie en tout point de $X$.
\item On dit que $f$ est de {\it présentation finie} s'il est
localement de présentation finie, quasi-compact et quasi-séparé.
\end{enumerate}
\end{definition}

\noindent
Notez qu'un morphisme de présentation finie n'est {\bf pas} simplement
un morphisme quasi-compact qui est localement de présentation
finie. Plus tard, nous caractériserons les morphismes
qui sont localement de présentation finie comme les morphismes tels que
$$
\colim \Mor_S(T_i, X) = \Mor_S(\lim T_i, X)
$$
pour tout système orienté de schémas affines $T_i$ sur
$S$.
Voir Limites, Proposition \ref{limits-proposition-characterize-locally-finite-presentation}. Dans Limites, section \ref{limits-section-descending-relative}
nous montrons que, si $S = \lim_i S_i$ est une limite de schémas
affines, tout schéma $X$ de présentation finie
sur $S$ descend à un schéma $X_i$ sur $S_i$
pour certains $i$.

\begin{lemma}
\label{morphisms-lemma-locally-finite-presentation-characterize}
Soit $f : X \to S$ un morphisme de schémas. Les conditions
suivantes sont équivalentes
\begin{enumerate}
\item Le morphisme $f$ est localement de présentation finie.
\item Pour chaque ouvert affines $U \subset X$, $V \subset S$
avec $f(U) \subset V$ l'homomorphisme
d'anneaux $\mathcal{O}_S(V) \to \mathcal{O}_X(U)$ est de présentation finie.
\item Il existe un recouvrement ouvert $S = \bigcup_{j \in J} V_j$ et
des recouvrements ouverts $f^{-1}(V_j) = \bigcup_{i \in I_j} U_i$ tels que chacun
des morphismes $U_i \to V_j$, $j\in J, i\in I_j$ est localement
de présentation finie.
\item Il existe un recouvrement ouvert affine $S = \bigcup_{j \in J} V_j$
et des recouvrements ouverts affines $f^{-1}(V_j) = \bigcup_{i \in I_j} U_i$ tels
que l'homomorphisme d'anneaux $\mathcal{O}_S(V_j) \to \mathcal{O}_X(U_i)$ soit de présentation
finie, pour tout $j\in J, i\in I_j$.
\end{enumerate}
De plus, si $f$ est localement de présentation finie alors
pour tout sous-schémas ouverts $U \subset X$, $V \subset S$ avec $f(U) \subset V$
la restriction $f|_U : U \to V$ est localement de présentation finie.
\end{lemma}

\begin{proof}
Cela résulte du lemme \ref{morphisms-lemma-locally-P-characterize} si l'on montre que la propriété ``$R \to A$
est de présentation finie'' est locale.
On vérifie les conditions (a), (b) et (c) de
la Définition \ref{morphisms-definition-property-local}. Par
Algèbre, Lemme \ref{algebra-lemma-base-change-finiteness}, la propriété d'être de présentation finie est
stable par changement de base ; on en déduit la condition (a).
Par Algèbre, Lemme \ref{algebra-lemma-compose-finite-type},
elle est stable par composition, et
pour tout anneau $R$ l'homomorphisme
d'anneaux $R \to R_f$ est de présentation finie.
On en déduit la condition (b). Enfin, la
condition (c) résulte d'Algèbre, Lemme \ref{algebra-lemma-cover-upstairs}.
\end{proof}

\begin{lemma}
\label{morphisms-lemma-composition-finite-presentation}
La composition de deux morphismes qui sont localement de présentation finie est localement
de présentation finie. Il
en va de même pour les morphismes de présentation finie.
\end{lemma}

\begin{proof}
Dans la démonstration du lemme \ref{morphisms-lemma-locally-finite-presentation-characterize} nous avons vu qu'être de
présentation finie est une propriété locale des homomorphismes d'anneaux.
D'où la première affirmation du lemme
découle du lemme \ref{morphisms-lemma-composition-type-P} combinée au
fait qu'être de présentation finie est une
propriété des homomorphismes
d'anneaux qui est stable
sous composition, voir Algèbre, Lemme \ref{algebra-lemma-compose-finite-type}.
Par ce qui précède et le fait que les compositions
de quasi-compacts, quasi-morphismes
séparés sont quasi-compacts
et quasi-séparés, voir Schémas, Lemmes \ref{schemes-lemma-composition-quasi-compact}
et \ref{schemes-lemma-separated-permanence} on voit que la composition
de morphismes de présentation finie est de présentation
finie.
\end{proof}

\begin{lemma}
\label{morphisms-lemma-base-change-finite-presentation}
Le changement de base d'un morphisme qui est localement de présentation finie
est localement de présentation finie. Il en va de même pour les morphismes de
présentation finie.
\end{lemma}

\begin{proof}
Dans la démonstration du lemme \ref{morphisms-lemma-locally-finite-presentation-characterize} nous avons vu qu'être de
présentation finie est une propriété locale des homomorphismes
d'anneaux. D'où le premier énoncé du lemme
découle du lemme \ref{morphisms-lemma-composition-type-P} combiné au
fait qu'être de présentation finie est une
propriété des homomorphismes
d'anneaux qui est stable
sous changement de base, voir Algèbre, Lemme \ref{algebra-lemma-base-change-finiteness}.
Par ce qui précède et le fait qu'un
changement de base d'un quasi-compact, quasi-morphisme
séparé est quasi-compact
et quasi-séparé, voir Schémas, Lemmes \ref{schemes-lemma-quasi-compact-preserved-base-change}
et \ref{schemes-lemma-separated-permanence} on voit que le changement
de base d'un morphisme de présentation finie est un morphisme
de présentation finie.
\end{proof}

\begin{lemma}
\label{morphisms-lemma-open-immersion-locally-finite-presentation}
Toute immersion ouverte est localement de présentation finie.
\end{lemma}

\begin{proof}
Ceci est vrai car une immersion ouverte est un isomorphisme local.
\end{proof}

\begin{lemma}
\label{morphisms-lemma-quasi-compact-open-immersion-finite-presentation}
Toute immersion ouverte est de présentation finie si et seulement si elle
est quasi-compacte.
\end{lemma}

\begin{proof}
Nous avons vu (Lemme \ref{morphisms-lemma-open-immersion-locally-finite-presentation}) qu'une immersion ouverte
est localement de présentation finie.
Nous avons vu (Schémas, Lemme \ref{schemes-lemma-immersions-monomorphisms}) qu'une immersion
est séparée et donc quasi-séparée. De ceci
et de la Définition \ref{morphisms-definition-finite-presentation} découle le lemme.
\end{proof}

\begin{lemma}
\label{morphisms-lemma-closed-immersion-finite-presentation}
\begin{slogan}
Les immersions fermées de présentation finie correspondent
à des faisceaux quasi-cohérents d'idéaux de type fini.
\end{slogan}
Une immersion fermée $i : Z \to X$ est de présentation finie si
et seulement si le faisceau quasi-cohérent
d'idéaux associé $\mathcal{I} = \Ker(\mathcal{O}_X \to i_*\mathcal{O}_Z)$ est de type
fini (en tant que $\mathcal{O}_X$-module).
\end{lemma}

\begin{proof}
Sur tout $\Spec(R) \subset X$ ouvert affine, nous avons
$i^{-1}(\Spec(R)) = \Spec(R/I)$ et $\mathcal{I} = \widetilde{I}$. De plus,
$\mathcal{I}$ est de type fini si et seulement si $I$
est un $R$-module fini pour chaque ouvert affine
de ce type (voir Propriétés,
Lemme \ref{properties-lemma-finite-type-module}). Et $R/I$ est de présentation finie
sur $R$ si et seulement si
$I$ est un $R$-module fini. On obtient donc le résultat.
\end{proof}

\begin{lemma}
\label{morphisms-lemma-finite-presentation-finite-type}
Un morphisme qui est localement de présentation finie est localement de type fini.
Un morphisme de présentation finie est de type fini.
\end{lemma}

\begin{proof}
Omis.
\end{proof}

\begin{lemma}
\label{morphisms-lemma-noetherian-finite-type-finite-presentation}
\begin{slogan}
Sur une base localement noethérien, le type fini est présentation finie.
\end{slogan}
Soit $f : X \to S$ un morphisme.
\begin{enumerate}
\item Si $S$ est localement noethérien et $f$ localement de type fini
alors $f$ est localement de présentation finie.
\item Si $S$ est localement noethérien et $f$ de type fini
alors $f$ est de présentation finie.
\end{enumerate}
\end{lemma}

\begin{proof}
La première affirmation découle du fait qu'un
anneau de type fini sur un anneau noéthérien est de présentation
finie, voir Algèbre, Lemme \ref{algebra-lemma-Noetherian-finite-type-is-finite-presentation}. Supposons que $f$
soit de type fini et que $S$ soit localement
noethérien. Alors $f$ est quasi-compact et
localement de fini présenté par (1). Il suffit
donc de prouver que $f$ est quasi-séparé. Cela résulte
du lemme \ref{morphisms-lemma-finite-type-Noetherian-quasi-separated} (et Lemme \ref{morphisms-lemma-finite-presentation-finite-type}).
\end{proof}

\begin{lemma}
\label{morphisms-lemma-finite-presentation-quasi-compact-quasi-separated}
Soit $S$ un schéma quasi-compact et quasi-séparé. Si $X$
est de présentation finie sur $S$, alors $X$ est quasi-compact
et quasi-séparé.
\end{lemma}

\begin{proof}
Omis.
\end{proof}

\begin{lemma}
\label{morphisms-lemma-finite-presentation-permanence}
Soit $f : X \to Y$ un morphisme de schémas sur $S$.
\begin{enumerate}
\item Si $X$ est localement de présentation finie sur $S$ et $Y$
est localement de type fini sur $S$, alors $f$ est localement
de présentation finie.
\item Si $X$ est de présentation finie sur $S$ et $Y$ est quasi-séparé
et localement de type fini sur $S$, alors $f$ est de présentation finie.
\end{enumerate}
\end{lemma}

\begin{proof}
Démonstration de (1). Via Lemme \ref{morphisms-lemma-locally-finite-presentation-characterize} cela se traduit par le fait algébrique
suivant : Étant donné les homomorphismes
d'anneaux $A \to B \to C$ tels que $A \to C$
est de présentation finie et $A \to B$ est de
type fini, alors $B \to C$ est de présentation
finie. Voir Algèbre, Lemme \ref{algebra-lemma-compose-finite-type}.

\medskip\noindent
La partie (2)
découle de (1) et des Schémas, Lemmes
\ref{schemes-lemma-compose-after-separated} et \ref{schemes-lemma-quasi-compact-permanence}.
\end{proof}

\begin{lemma}
\label{morphisms-lemma-diagonal-morphism-finite-type}
Soit $f : X \to Y$ un morphisme de schémas de diagonale $\Delta : X \to X \times_Y X$.
Si $f$ est localement de type fini alors $\Delta$
est localement de présentation finie. Si $f$ est quasi-séparé et
localement de type fini, alors $\Delta$ est de présentation finie.
\end{lemma}

\begin{proof}
Notez que $\Delta$ est un morphisme de schémas sur $X$
(via la deuxième projection $X \times_Y X \to X$). Supposons que $f$
soit localement de type fini. Notez que $X$ est de présentation
finie sur $X$ et $X \times_Y X$ est localement de type fini sur $X$
(par le lemme \ref{morphisms-lemma-base-change-finite-type}).
Ainsi, la première affirmation du lemme
\ref{morphisms-lemma-finite-presentation-permanence} est valable. Le deuxième énoncé découle du premier,
des définitions, et du fait qu'un morphisme diagonal est un monomorphisme,
donc séparé (Schémas, Lemme \ref{schemes-lemma-monomorphism-separated}).
\end{proof}








\section{Ensembles constructibles}
\label{morphisms-section-constructible}

\noindent
Les ensembles constructibles et localement constructibles des schémas ont été discutés
dans Propriétés, section \ref{properties-section-constructible}. Dans cette section nous prouvons
quelques résultats concernant les images et les images inverses d'ensembles
(localement) constructibles. Le résultat principal est le théorème de Chevalley
qui affirme que l'image d'un ensemble localement constructible sous un morphisme
de présentation finie est localement constructible.

\begin{lemma}
\label{morphisms-lemma-inverse-image-constructible}
Soit $f : X \to Y$ un morphisme de schémas. Soit $E \subset Y$
un sous-ensemble. Si $E$ est
(localement) constructible dans $Y$, alors $f^{-1}(E)$ est (localement) constructible
dans $X$.
\end{lemma}

\begin{proof}
Pour montrer que l'image réciproque de chaque sous-ensemble constructible est constructible,
il suffit de montrer que l'image réciproque de chaque $V$ ouvert rétrocompact
de $Y$ est rétrocompacte
dans $X$, voir Topologie, Lemme \ref{topology-lemma-inverse-images-constructibles}. L'importance
du fait que $V$ soit rétrocompact dans
$Y$ est simplement que l'immersion ouverte $V \to Y$ est quasi-compacte.
Ainsi le changement de base $f^{-1}(V) = X \times_Y V \to X$ est également quasi-compact,
voir
Schémas, Lemme \ref{schemes-lemma-quasi-compact-preserved-base-change}. Nous voyons donc que $f^{-1}(V)$ est rétrocompact
dans $X$. Supposons que $E$ soit localement
constructible dans $Y$. Choisissez
$x \in X$. Choisir un voisinage affine $V$ de $f(x)$ et
un voisinage affine $U \subset X$ de $x$ tel que $f(U) \subset V$. Ainsi
nous pensons à $f|_U : U \to V$ comme un morphisme en $V$.
Par Propriétés, Lemme \ref{properties-lemma-locally-constructible} nous voyons
que $E \cap V$ est constructible dans $V$. Par le cas constructible
nous voyons que $(f|_U)^{-1}(E \cap V)$ est constructible
dans $U$. Puisque $(f|_U)^{-1}(E \cap V) = f^{-1}(E) \cap U$ on obtient le résultat.
\end{proof}

\begin{lemma}
\label{morphisms-lemma-chevalley}
Soit $f : X \to Y$ un morphisme de schémas. Supposons
que
\begin{enumerate}
\item $f$ est quasi-compact et localement de présentation finie, et que
\item $Y$ est quasi-compact et quasi-séparé.
\end{enumerate}
Alors l'image de chaque sous-ensemble constructible de $X$ est constructible dans $Y$.
\end{lemma}

\begin{proof}
Par
Propriétés,
Lemme \ref{properties-lemma-constructible-quasi-compact-quasi-separated} il suffit de prouver ce lemme dans le cas
où $Y$ est affine. Dans ce cas, $X$
est quasi-compact. On peut donc écrire
$X = U_1 \cup \ldots \cup U_n$ avec chaque $U_i$ ouvert affine dans $X$.
Si $E \subset X$ est constructible, alors chaque $E \cap U_i$
est
également
constructible, voir Topologie, Lemme \ref{topology-lemma-open-immersion-constructible-inverse-image}. Donc,
puisque $f(E) = \bigcup f(E \cap U_i)$ et que les unions finies d'ensembles constructibles
sont constructibles, cela nous ramène au cas où
$X$ est affine. Dans ce cas, le
résultat est Algèbre, Théorème \ref{algebra-theorem-chevalley}.
\end{proof}

\begin{theorem}[Théorème de Chevalley]
\label{morphisms-theorem-chevalley}
\begin{reference}
\cite[IV, Théorème 1.8.4]{EGA}
\end{reference}
Soit $f : X \to Y$ un morphisme de schémas. Supposons
que $f$ soit quasi-compact et localement de présentation finie. Alors
l'image de chaque sous-ensemble localement constructible est localement constructible.
\end{theorem}

\begin{proof}
Soit $E \subset X$ localement constructible. Il
faut montrer que $f(E)$ est également constructible
localement. Nous allons montrer que $f(E) \cap V$ est constructible
pour tout $V \subset Y$ ouvert affine. On se réduit donc au cas où
$Y$ est affine. Dans ce cas, $X$ est quasi-compact.
On peut donc écrire $X = U_1 \cup \ldots \cup U_n$ avec chaque $U_i$
ouvert affine dans $X$. Si $E \subset X$ est localement constructible,
alors chaque
$E \cap U_i$ est constructible, voir Propriétés, Lemme \ref{properties-lemma-locally-constructible}.
Donc, puisque $f(E) = \bigcup f(E \cap U_i)$ et que les unions finies d'ensembles constructibles
sont constructibles, cela nous ramène au cas où $X$ est
affine. Dans ce cas, le résultat est
Algèbre, Théorème \ref{algebra-theorem-chevalley}.
\end{proof}

\begin{lemma}
\label{morphisms-lemma-constructible-containing-open}
Soit $X$ un schéma. Soit $x \in X$. Soit $E \subset X$ un sous-ensemble
localement constructible. Si $\{x' \mid x' \leadsto x\} \subset E$, alors $E$
contient un voisinage ouvert de $x$.
\end{lemma}

\begin{proof}
Supposons $\{x' \mid x' \leadsto x\} \subset E$. Nous pouvons supposer
que $X$ est
affine. Dans ce cas, $E$
est constructible, voir Propriétés, Lemme \ref{properties-lemma-locally-constructible}.
En particulier, le complément $E^c$ est
également constructible. D'après Algèbre, Lemme \ref{algebra-lemma-constructible-is-image},
on peut trouver un morphisme de schémas affines
$f : Y \to X$ tel que $E^c = f(Y)$. Soit $Z \subset X$ l'image schématique
de $f$. Par le Lemme \ref{morphisms-lemma-reach-points-scheme-theoretic-image}
et l'hypothèse $\{x' \mid x' \leadsto x\} \subset E$ on voit que $x \not \in Z$.
Ainsi $X \setminus Z \subset E$ est un voisinage ouvert de $x$
contenu dans $E$.
\end{proof}




\section{Morphismes ouverts}
\label{morphisms-section-open}

\begin{definition}
\label{morphisms-definition-open}
Soit $f : X \to S$ un morphisme.
\begin{enumerate}
\item On dit que $f$ est {\it ouvert} si l'application sur les espaces
topologiques sous-jacents est ouverte.
\item On dit que $f$ est {\it universellement ouvert} si
pour tout morphisme de schémas $S' \to S$ le changement de base $f' : X_{S'} \to S'$ est ouvert.
\end{enumerate}
\end{definition}

\noindent
Selon la topologie,
les générisations du lemme \ref{topology-lemma-closed-open-map-specialization} s'élèvent le long
de certains types d'applications ouvertes d'espaces topologiques.
En fait, les générisations suivent tout morphisme de schémas ouvert
(voir
Lemme \ref{morphisms-lemma-open-generizing}). Aussi,
nous verrons que les générisations s'élèvent le long de morphismes
de schémas plats (Lemme \ref{morphisms-lemma-generalizations-lift-flat}). Cela implique
parfois à son tour que le morphisme est ouvert.

\begin{lemma}
\label{morphisms-lemma-locally-finite-presentation-universally-open}
Soit $f : X \to S$ un morphisme.
\begin{enumerate}
\item Si $f$ est localement de présentation finie et que les générisations s'élèvent le long de $f$,
alors $f$ est ouvert.
\item Si $f$ est localement de présentation finie et que les générisations s'élèvent à chaque
changement de base de $f$, alors $f$ est universellement ouvert.
\end{enumerate}
\end{lemma}

\begin{proof}
Il suffit de prouver la première assertion.
Cela se réduit au cas où $X$ et
$S$ sont affines. Dans
ce cas, le résultat découle de l'algèbre, du Lemme \ref{algebra-lemma-going-up-down-specialization}
et de la proposition \ref{algebra-proposition-fppf-open}.
\end{proof}

\noindent
Voir aussi Lemme \ref{morphisms-lemma-fppf-open} pour le cas d'un morphisme plat
de présentation finie.

\begin{lemma}
\label{morphisms-lemma-composition-open}
Une composition de morphismes (universellement) ouverts est (universellement) ouverte.
\end{lemma}

\begin{proof}
Omis.
\end{proof}

\begin{lemma}
\label{morphisms-lemma-scheme-over-field-universally-open}
Soit $k$ un corps. Soit $X$ un schéma sur $k$.
Le morphisme de structure $X \to \Spec(k)$ est universellement ouvert.
\end{lemma}

\begin{proof}
Soit $S \to \Spec(k)$ un morphisme. Il faut
montrer que le changement de base $X_S \to S$ est ouvert.
La question est locale sur $S$ et $X$, on peut
donc supposer que $S$ et $X$ sont affines.
Dans ce cas, le résultat est Algèbre, Lemme \ref{algebra-lemma-map-into-tensor-algebra-open}.
\end{proof}

\begin{lemma}
\label{morphisms-lemma-open-generizing}
\begin{reference}
découle de l'implication (a) $\Rightarrow$ (b) dans \cite[IV, Corollaire
1.10.4]{EGA}
\end{reference}
Soit $\varphi : X \to Y$ un morphisme de schémas. Si $\varphi$ est ouvert,
alors $\varphi$ est en train de générer (c'est-à-dire que
les générisations s'élèvent le long de $\varphi$). Si
$\varphi$ est universellement ouvert, alors $\varphi$ est universellement
génératrice.
\end{lemma}

\begin{proof}
Supposons que $\varphi$ soit
ouvert. Soit $y' \leadsto y$ une spécialisation de points de
$Y$. Soit $x \in X$ avec $\varphi(x) = y$.
Choisissons des ouverts affines $U \subset X$ et $V \subset Y$
tels que $\varphi(U) \subset V$ et $x \in U$. Alors $y' \in V$ aussi. Nous pouvons
donc remplacer $X$ par $U$ et $Y$ par $V$
et supposer
que $X$, $Y$ sont affines. Le cas affine
est l'Algèbre,
Lemme \ref{algebra-lemma-open-going-down} (combiné avec l'Algèbre, Lemme \ref{algebra-lemma-going-up-down-specialization}).
\end{proof}

\begin{lemma}
\label{morphisms-lemma-descent-quasi-compact}
Soit $f : X \to Y$ un morphisme de schémas. Soit
$g : Y' \to Y$ ouvert et surjectif tel que le changement de base $f' : X' \to Y'$
soit quasi-compact. Alors $f$ est quasi-compact.
\end{lemma}

\begin{proof}
Soit $V \subset Y$ un ouvert quasi-compact. Comme $g$ est ouvert
et surjectif on peut trouver un $W' \subset W$ ouvert quasi-compact tel
que $g(W') = V$. Par hypothèse $(f')^{-1}(W')$ est quasi-compact.
L'image de $(f')^{-1}(W')$ dans $X$ est
égale à $f^{-1}(V)$, voir Lemme \ref{morphisms-lemma-when-point-maps-to-pair}.
Ainsi $f^{-1}(V)$ est quasi-compact comme l'image d'un espace quasi-compact,
voir Topologie, Lemme \ref{topology-lemma-image-quasi-compact}. Ainsi $f$
est quasi-compact.
\end{proof}





\section{Morphismes submersifs}
\label{morphisms-section-submersive}

\begin{definition}
\label{morphisms-definition-submersive}
Soit $f : X \to Y$ un morphisme de schémas.
\begin{enumerate}
\item Nous disons que $f$ est {\it submersif}\footnote{Ceci
est très différent de la notion de submersion de
variétés différentielles.} si l'application continue des espaces topologiques
sous-jacents est submersive, voir Topologie, Définition \ref{topology-definition-submersive}.
\item On dit que $f$ est {\it universellement submersif}
si pour tout morphisme de schémas $Y' \to Y$ le
changement de base $Y' \times_Y X \to Y'$ est submersif.
\end{enumerate}
\end{definition}

\noindent
On remarque qu'un morphisme submersif est notamment surjectif.

\begin{lemma}
\label{morphisms-lemma-base-change-universally-submersive}
Le changement de base d'un morphisme de schémas universellement submersif
par tout morphisme de schémas est universellement submersif.
\end{lemma}

\begin{proof}
Cela découle directement de la définition.
\end{proof}

\begin{lemma}
\label{morphisms-lemma-composition-universally-submersive}
La composition d'une paire de morphismes de schémas (universellement) submersifs
est (universellement) submersive.
\end{lemma}

\begin{proof}
Omis.
\end{proof}










\section{Morphismes plats}
\label{morphisms-section-flat}

\noindent
La planéité est l'un des outils techniques les plus importants en
géométrie algébrique. Dans cette section, nous introduisons cette notion.
Nous limitons intentionnellement la discussion à des observations
simples, en dehors du lemme \ref{morphisms-lemma-fppf-open}. Une classe
de résultats
très importante, à savoir les critères de
planéité, est discutée dans Algèbre,
sections \ref{algebra-section-criteria-flatness}, \ref{algebra-section-flatness-artinian}, \ref{algebra-section-more-flatness-criteria},
et plus sur les
morphismes, section \ref{more-morphisms-section-criterion-flat-fibres}. Il existe un
chapitre consacré au matériel avancé sur les morphismes de
schémas plats, à savoir Compléments sur la platitude, section \ref{flat-section-introduction}.

\medskip\noindent Rappelons
qu'un module $M$ sur un anneau $R$ est {\it plat}
si le foncteur $-\otimes_R M : \text{Mod}_R \to \text{Mod}_R$ est exact. Un homomorphisme d'anneaux $R \to A$
est dit {\it plat} si $A$ est plat comme $R$-module.
Voir
Algèbre, Définition \ref{algebra-definition-flat}.

\begin{definition}
\label{morphisms-definition-flat}
Soit $f : X \to S$ un morphisme de schémas.
Soit $\mathcal{F}$ un faisceau quasi-cohérent de $\mathcal{O}_X$-modules.
\begin{enumerate}
\item On dit que $f$ est {\it plat en un point
$x \in X$} si l'anneau local $\mathcal{O}_{X, x}$ est plat sur
l'anneau local $\mathcal{O}_{S, f(x)}$.
\item On dit que $\mathcal{F}$ est {\it plat sur $S$ en un point
$x \in X$} si la tige $\mathcal{F}_x$ est un module $\mathcal{O}_{S, f(x)}$ plat.
\item Nous disons que $f$ est {\it plat} si $f$ est plat en chaque point de $X$.
\item On dit que $\mathcal{F}$ est {\it plat sur $S$}
si $\mathcal{F}$ est plat sur $S$ en chaque point $x$ de $X$.
\end{enumerate}
\end{definition}

\noindent
On voit donc que $f$ est plat si et seulement
si le faisceau structurel $\mathcal{O}_X$ est plat sur $S$.

\begin{lemma}
\label{morphisms-lemma-flat-module-characterize}
Soit $f : X \to S$ un morphisme de schémas. Soit
$\mathcal{F}$ un faisceau quasi-cohérent de $\mathcal{O}_X$-modules. Les conditions
suivantes sont équivalentes
\begin{enumerate}
\item Le faisceau $\mathcal{F}$ est plat sur $S$.
\item Pour chaque ouverture affine $U \subset X$, $V \subset S$
avec $f(U) \subset V$ le $\mathcal{O}_S(V)$-module $\mathcal{F}(U)$ est plat.
\item Il existe un recouvrement ouvert $S = \bigcup_{j \in J} V_j$ et
des revêtements ouverts $f^{-1}(V_j) = \bigcup_{i \in I_j} U_i$ tels que chacun des
modules $\mathcal{F}|_{U_i}$ soit à plat sur $V_j$,
pour tous les $j\in J, i\in I_j$.
\item Il existe un recouvrement ouvert affine $S = \bigcup_{j \in J} V_j$
et des recouvrements ouverts affines $f^{-1}(V_j) = \bigcup_{i \in I_j} U_i$ tels
que $\mathcal{F}(U_i)$ soit un $\mathcal{O}_S(V_j)$-module plat, pour tout
$j\in J, i\in I_j$.
\end{enumerate}
De plus, si $\mathcal{F}$ est plat sur $S$ alors
pour tout sous-schémas ouverts $U \subset X$, $V \subset S$ avec $f(U) \subset V$ la
restriction $\mathcal{F}|_U$ est plate sur $V$.
\end{lemma}

\begin{proof}
Soit $R \to A$ un homomorphisme d'anneaux. Soit $M$
un module $A$. Si $M$ est $R$-plat,
alors pour tous les nombres premiers $\mathfrak q$ le module $M_{\mathfrak q}$
est plat sur $R_{\mathfrak p}$ avec $\mathfrak p$ le premier de $R$ se trouvant sous $\mathfrak q$.
Inversement, si $M_{\mathfrak q}$ est plat sur $R_{\mathfrak p}$ pour tous les nombres premiers
$\mathfrak q$ de $A$, alors $M$
est plat sur $R$. Voir Algèbre, Lemme \ref{algebra-lemma-flat-localization}.
Cette équivalence implique facilement les énoncés du lemme.
\end{proof}

\begin{lemma}
\label{morphisms-lemma-flat-characterize}
Soit $f : X \to S$ un morphisme de schémas. Les conditions
suivantes sont équivalentes
\begin{enumerate}
\item Le morphisme $f$ est plat.
\item Pour chaque ouvert affines $U \subset X$, $V \subset S$
avec $f(U) \subset V$ l'homomorphisme
d'anneaux $\mathcal{O}_S(V) \to \mathcal{O}_X(U)$ est plat.
\item Il existe un recouvrement ouvert $S = \bigcup_{j \in J} V_j$ et
des recouvrements ouverts $f^{-1}(V_j) = \bigcup_{i \in I_j} U_i$ tels que chacun
des morphismes $U_i \to V_j$, $j\in J, i\in I_j$ soit
plat.
\item Il existe un recouvrement ouvert affine $S = \bigcup_{j \in J} V_j$
et un recouvrement ouvert affine $f^{-1}(V_j) = \bigcup_{i \in I_j} U_i$
tel que $\mathcal{O}_S(V_j) \to \mathcal{O}_X(U_i)$ soit plat, pour tout
$j\in J, i\in I_j$.
\end{enumerate}
De plus, si $f$ est plat alors
pour tout sous-schémas ouverts $U \subset X$, $V \subset S$ avec $f(U) \subset V$
la restriction $f|_U : U \to V$ est plate.
\end{lemma}

\begin{proof}
Ceci est un cas particulier du lemme \ref{morphisms-lemma-flat-module-characterize}
ci-dessus.
\end{proof}

\begin{lemma}
\label{morphisms-lemma-pushforward-flat-affine}
Soit $f : X \to Y$ un morphisme affine de schémas sur un schéma de base $S$.
Soit $\mathcal{F}$ un module $\mathcal{O}_X$ quasi-cohérent. Alors $\mathcal{F}$
est plat sur $S$ si et seulement si $f_*\mathcal{F}$
est plat sur $S$.
\end{lemma}

\begin{proof}
Par le Lemme \ref{morphisms-lemma-flat-module-characterize} et le fait que $f$ soit un morphisme
affine, cela nous ramène au cas affine. Supposons que $X \to Y \to S$
corresponde aux homomorphismes d'anneaux $C \leftarrow B \leftarrow A$. Soit $N$
le module $C$ correspondant à $\mathcal{F}$.
Rappelons que $f_*\mathcal{F}$ correspond à $N$ vu comme un $B$-module,
voir Schémas, Lemme \ref{schemes-lemma-widetilde-pullback}. Le résultat
est donc clair.
\end{proof}

\begin{lemma}
\label{morphisms-lemma-composition-module-flat}
Soit $X \to Y \to Z$ des morphismes de schémas. Soit $\mathcal{F}$ un module $\mathcal{O}_X$ quasi-cohérent.
Soit $x \in X$ avec l'image $y$ dans $Y$. Si $\mathcal{F}$ est plat
sur $Y$ à $x$ et que $Y$ est plat sur $Z$ à $y$,
alors $\mathcal{F}$ est plat sur $Z$ à $x$.
\end{lemma}

\begin{proof}
Voir Algèbre, Lemme \ref{algebra-lemma-composition-flat}.
\end{proof}

\begin{lemma}
\label{morphisms-lemma-composition-flat}
La composition des morphismes plats est plate.
\end{lemma}

\begin{proof}
Ceci est un cas particulier du lemme \ref{morphisms-lemma-composition-module-flat}.
\end{proof}

\begin{lemma}
\label{morphisms-lemma-base-change-module-flat}
Soit $f : X \to S$ un morphisme de schémas. Soit
$\mathcal{F}$ un faisceau quasi-cohérent de $\mathcal{O}_X$-modules. Soit
$g : S' \to S$ un morphisme de schémas. Notons
$g' : X' = X_{S'} \to X$ la projection. Soit $x' \in X'$ un
point d'image $x = g'(x') \in X$. Si $\mathcal{F}$ est plat sur
$S$ à $x$, alors $(g')^*\mathcal{F}$ est
plat sur $S'$ à $x'$. En particulier,
si $\mathcal{F}$ est plat sur $S$, alors
$(g')^*\mathcal{F}$ est plat sur $S'$.
\end{lemma}

\begin{proof}
Voir Algèbre, Lemme \ref{algebra-lemma-flat-base-change}.
\end{proof}

\begin{lemma}
\label{morphisms-lemma-base-change-flat}
Le changement de base d'un morphisme plat est plat.
\end{lemma}

\begin{proof}
Ceci est un cas particulier du lemme \ref{morphisms-lemma-base-change-module-flat}.
\end{proof}

\begin{lemma}
\label{morphisms-lemma-generalizations-lift-flat}
Soit $f : X \to S$ un morphisme de schémas plat.
Puis les générisations s'élèvent le
long de $f$, voir Topologie, Définition \ref{topology-definition-lift-specializations}.
\end{lemma}

\begin{proof}
Voir Algèbre, section \ref{algebra-section-going-up}.
\end{proof}

\begin{lemma}
\label{morphisms-lemma-fppf-open}
A morphisme plat localement de fini présentation est universellement ouvert.
\end{lemma}

\begin{proof}
Cela résulte de Lemmes \ref{morphisms-lemma-generalizations-lift-flat} et Lemme \ref{morphisms-lemma-locally-finite-presentation-universally-open}
ci-dessus. Nous pouvons également argumenter
directement comme suit.

\medskip\noindent
Soit $f : X \to S$ soit plat et localement de présentation finie.
Par les Lemmes \ref{morphisms-lemma-base-change-flat} et \ref{morphisms-lemma-base-change-finite-presentation}
tout changement de base de $f$ est plat et localement de présentation
finie. Il suffit donc de montrer que $f$ est ouvert. Pour
montrer que $f$ est ouvert, il suffit de montrer
que l'on peut couvrir $X$ par des ouverts affines
$X = \bigcup U_i$ telles que $U_i \to S$ soit ouvert. On peut recouvrir $X$
par des ouverts affines $U_i \subset X$ tels que chaque
$U_i$ ait son image dans un ouvert affine $V_i \subset S$ et tels
que l'homomorphisme d'anneaux induit $\mathcal{O}_S(V_i) \to \mathcal{O}_X(U_i)$ soit plat et de présentation
finie (Lemmes \ref{morphisms-lemma-flat-characterize} et \ref{morphisms-lemma-locally-finite-presentation-characterize}).
Alors $U_i \to V_i$ est ouvert
par Algèbre, Proposition \ref{algebra-proposition-fppf-open} et la démonstration
est complète.
\end{proof}

\begin{lemma}
\label{morphisms-lemma-pf-flat-module-open}
Soit $f : X \to Y$ un morphisme de schémas. Soit
$\mathcal{F}$ un module $\mathcal{O}_X$ quasi-cohérent. Supposons
que $f$ localement présentation finie, $\mathcal{F}$
de type fini, $X = \text{Supp}(\mathcal{F})$ et $\mathcal{F}$ plat sur
$Y$. Alors $f$ est universellement ouvert.
\end{lemma}

\begin{proof}
Par les Lemmes \ref{morphisms-lemma-base-change-module-flat}, \ref{morphisms-lemma-base-change-finite-presentation},
et \ref{morphisms-lemma-support-finite-type} les hypothèses
sont conservées sous
changement de base. Par le Lemme
\ref{morphisms-lemma-locally-finite-presentation-universally-open} il suffit de montrer que les générisations
s'élèvent le long de $f$.
Cela résulte de Algèbre, Lemme \ref{algebra-lemma-going-down-flat-module}.
\end{proof}

\begin{lemma}
\label{morphisms-lemma-fpqc-quotient-topology}
\begin{reference}
\cite[Expose VIII, Corollaire 4.3]{SGA1} et \cite[IV,
Corollaire 2.3.12]{EGA}
\end{reference}
Soit $f : X \to Y$ être un quasi-compact, surjectif, morphisme plat. Un sous-ensemble
$T \subset Y$ est ouvert (resp.\ fermé) si et seulement $f^{-1}(T)$
est ouvert (resp.\ fermé). Autrement dit, $f$ est
un morphisme submersif.
\end{lemma}

\begin{proof}
La question est locale sur $Y$, on peut donc supposer que $Y$
est affine. Dans ce cas, $X$ est quasi-compact
tout comme $f$ est quasi-compact. Écrivez $X = X_1 \cup \ldots \cup X_n$ comme
une union finie d'ouverts affines. Alors $f' : X' = X_1 \amalg \ldots \amalg X_n \to Y$ est un
morphisme surjectif plat de schémas affines. Notez que pour $T \subset Y$
nous avons $(f')^{-1}(T) = f^{-1}(T) \cap X_1 \amalg \ldots \amalg f^{-1}(T) \cap X_n$. Par conséquent, $f^{-1}(T)$ est ouvert si et seulement
si $(f')^{-1}(T)$ est ouvert. Ainsi, nous pouvons supposer
que $X$ et $Y$ sont affines.

\medskip\noindent
Soit $f : \Spec(B) \to \Spec(A)$ un morphisme surjectif de schémas affines
correspondant à un homomorphisme plat d'anneaux $A \to B$.
Supposons que $f^{-1}(T)$ soit fermé, disons $f^{-1}(T) = V(J)$
pour $J \subset B$ un idéal. Alors $T = f(f^{-1}(T)) = f(V(J))$ est l'image
de $\Spec(B/J) \to \Spec(A)$ (ici on utilise que $f$ est surjectif).
En revanche, les générisations s'élèvent le long
de $f$ (Lemme \ref{morphisms-lemma-generalizations-lift-flat}). Ainsi par Topologie, Lemme
\ref{topology-lemma-lift-specializations-images} on voit que $Y \setminus T = f(X \setminus f^{-1}(T))$ est stable par générisation.
Ainsi $T$ est stable par spécialisation
(Topologie, Lemme \ref{topology-lemma-open-closed-specialization}). Ainsi $T$
est fermé par
Algèbre, Lemme \ref{algebra-lemma-image-stable-specialization-closed}.
\end{proof}

\begin{lemma}
\label{morphisms-lemma-flat-permanence}
Soit $h : X \to Y$ un morphisme de schémas sur $S$. Soit
$\mathcal{G}$ un faisceau quasi-cohérent sur $Y$. Soit
$x \in X$ avec $y = h(x) \in Y$. Si $h$ est plat en $x$, alors
$$
\mathcal{G}\text{ est plat sur }S\text{ en }y
\Leftrightarrow
h^*\mathcal{G}\text{ est plat sur }S\text{ en }x.
$$
En particulier : Si $h$ est surjectif et plat, alors $\mathcal{G}$
est plat sur $S$, si et seulement si $h^*\mathcal{G}$
est plat sur $S$. Si $h$ est surjectif et plat, et que $X$
est plat sur $S$, alors $Y$ est plat sur $S$.
\end{lemma}

\begin{proof}
Vous pouvez le prouver
en appliquant l'algèbre, Lemme \ref{algebra-lemma-flatness-descends-more-general}.
Voici une démonstration directe. Soit $s \in S$
l'image de
$y$. Considérez les applications locales
anneau $\mathcal{O}_{S, s} \to \mathcal{O}_{Y, y} \to \mathcal{O}_{X, x}$. Par hypothèse l'homomorphisme
d'anneaux $\mathcal{O}_{Y, y} \to \mathcal{O}_{X, x}$
est fidèlement plat, voir
Algèbre, Lemme \ref{algebra-lemma-local-flat-ff}.
Soit $N = \mathcal{G}_y$. A noter que $h^*\mathcal{G}_x = N \otimes_{\mathcal{O}_{Y, y}} \mathcal{O}_{X, x}$, voir
Faisceaux, Lemme \ref{sheaves-lemma-stalk-pullback-modules}. Soit $M' \to M$
une injection de $\mathcal{O}_{S, s}$-modules.
Par la planéité fidèle mentionnée ci-dessus, nous avons
\begin{align*}
\Ker(
M' \otimes_{\mathcal{O}_{S, s}} N \to M \otimes_{\mathcal{O}_{S, s}} N)
\otimes_{\mathcal{O}_{Y, y}} \mathcal{O}_{X, x} \\
=
\Ker(
M' \otimes_{\mathcal{O}_{S, s}} N
\otimes_{\mathcal{O}_{Y, y}} \mathcal{O}_{X, x}
\to
M \otimes_{\mathcal{O}_{S, s}} N
\otimes_{\mathcal{O}_{Y, y}} \mathcal{O}_{X, x})
\end{align*}
Par conséquent, l'équivalence du lemme découle de la deuxième caractérisation
de la platitude donnée dans
Algèbre, Lemme \ref{algebra-lemma-flat}.
\end{proof}

\begin{lemma}
\label{morphisms-lemma-flat-pullback-support}
Soit $f : Y \to X$ un morphisme de schémas. Soit $\mathcal{F}$ un
$\mathcal{O}_X$-module quasi-cohérent de type fini avec support
schématique $Z \subset X$. Si $f$ est plat,
alors $f^{-1}(Z)$ est le support schématique de $f^*\mathcal{F}$.
\end{lemma}

\begin{proof}
En utilisant la caractérisation schématique des supports
sur affines telle que donnée dans le lemme \ref{morphisms-lemma-scheme-theoretic-support}
nous réduisons à l'Algèbre, Lemme \ref{algebra-lemma-annihilator-flat-base-change}.
\end{proof}

\begin{lemma}
\label{morphisms-lemma-flat-morphism-scheme-theoretically-dense-open}
Soit $f : X \to Y$ un morphisme de schémas plat. Soit $V \subset Y$ un ouvert
rétrocompact qui est schématiquement dense. Alors $f^{-1}V$ est schématiquement
dense en $X$.
\end{lemma}

\begin{proof}
Nous utiliserons la caractérisation
du lemme \ref{morphisms-lemma-characterize-scheme-theoretically-dense}. Il faut montrer
que pour tout $U \subset X$ ouvert l'application
$\mathcal{O}_X(U) \to \mathcal{O}_X(U \cap f^{-1}V)$ est injectif. Il suffit
de le prouver lorsque $U$ est un ouvert affine
qui s'envoie dans un ouvert affine $W \subset Y$.
Posons $W = \Spec(A)$ et $U = \Spec(B)$. Puis
$V \cap W = D(f_1) \cup \ldots \cup D(f_n)$
pour certains $f_i \in A$, voir
Algèbre, Lemme \ref{algebra-lemma-qc-open}.
Il faut donc montrer que $B \to B_{f_1} \times \ldots \times B_{f_n}$
est injectif. On sait que $A \to A_{f_1} \times \ldots \times A_{f_n}$ est injectif
et que $A \to B$ est plat. Puisque $B_{f_i} = A_{f_i} \otimes_A B$ on obtient le résultat.
\end{proof}

\begin{lemma}
\label{morphisms-lemma-flat-base-change-scheme-theoretic-image}
\begin{slogan}
La formation des images schématiques commute aux changements de base plats
dans le cas quasi-compact
\end{slogan}
Soit $f : X \to Y$ un morphisme plat de schémas. Soit $g : V \to Y$ un
morphisme de schémas quasi-compact. Soit $Z \subset Y$ l'image schématique
de $g$ et soit $Z' \subset X$ l'image schématique du changement
de base $V \times_Y X \to X$. Puis $Z' = f^{-1}Z$.
\end{lemma}

\begin{proof}
Rappelons que $Z$ est défini par
$\mathcal{I} = \Ker(\mathcal{O}_Y \to g_*\mathcal{O}_V)$, tandis que $Z'$ est défini par
$\mathcal{I}' = \Ker(\mathcal{O}_X \to
(V \times_Y X \to X)_*\mathcal{O}_{V \times_Y X})$,
voir le Lemme \ref{morphisms-lemma-quasi-compact-scheme-theoretic-image}. La question est locale
sur $X$ et $Y$, et l'on peut donc supposer
que $X$ et $Y$ sont affines. Nous pouvons remplacer $V$
par $\coprod V_i$ où $V = V_1 \cup \ldots \cup V_n$ est un recouvrement
ouvert affine fini. Nous pouvons donc supposer
que $g$ est affine. Dans ce cas,
$(V \times_Y X \to X)_*\mathcal{O}_{V \times_Y X}$ est le retrait de $g_*\mathcal{O}_V$ par $f$.
Puisque $f$ est plat, nous concluons que
$f^*\mathcal{I} = \mathcal{I}'$ et le lemme est vrai.
\end{proof}




\section{Immersions fermées plates}
\label{morphisms-section-flat-closed-immersions}

\noindent
Les composants connectés des schémas ne sont pas toujours ouverts. Mais ils ont toujours
une structure schématique canonique. Nous expliquons cela dans cette section.

\begin{lemma}
\label{morphisms-lemma-characterize-flat-closed-immersions}
Soit $X$ un schéma. La règle qui associe à un sous-schéma fermé de $X$
son sous-ensemble fermé sous-jacent définit une bijection
$$
\left\{
\begin{matrix}
\text{sous-schémas fermés }Z \subset X \\
\text{tels que }Z \to X\text{ soit plat}
\end{matrix}
\right\}
\leftrightarrow
\left\{
\begin{matrix}
\text{parties fermées }Z \subset X \\
\text{stables par générisation}
\end{matrix}
\right\}
$$
Si $Z \subset X$ est un tel sous-schéma fermé, tout morphisme de schémas $g : Y \to X$
tel que $g(Y) \subset Z$ ensemblistement se factorise
(schématiquement) par $Z$.
\end{lemma}

\begin{proof}
Le cas affine de la bijection est
l'Algèbre, Lemme \ref{algebra-lemma-pure-open-closed-specializations}. Pour les schémas généraux $X$
la bijection s'ensuit en recouvrant $X$ par affines et collage.
Détails omis. Pour l'assertion finale, notons que la projection
$Z \times_{X, g} Y \to Y$ est une immersion fermée plate
(Lemme \ref{morphisms-lemma-base-change-flat}) qui est bijective sur les espaces topologiques
sous-jacents et doit donc être un isomorphisme
par la bijection établie dans la première partie de la démonstration.
\end{proof}

\begin{lemma}
\label{morphisms-lemma-flat-closed-immersions-finite-presentation}
Une immersion fermée plate de présentation finie est
l'immersion ouverte d'un ouvert et sous-schéma fermé.
\end{lemma}

\begin{proof}
Le cas affine est
l'Algèbre, Lemme \ref{algebra-lemma-finitely-generated-pure-ideal}. En général, le lemme suit
en couvrant $X$ par des affines. Détails
omis.
\end{proof}

\noindent
Notons qu'un composant connexe $T$ d'un schéma $X$ est un
sous-ensemble fermé stable par génération. La définition suivante est donc
logique.

\begin{definition}
\label{morphisms-definition-scheme-structure-connected-component}
Soit $X$ un schéma. Soit $T \subset X$ un composant connexe. La structure
de schéma canonique {\it sur $T$}
est la structure de schéma unique sur $T$ telle
que l'immersion
fermée $T \to X$ est plate, voir Lemme \ref{morphisms-lemma-characterize-flat-closed-immersions}.
\end{definition}

\noindent
Il s'avère que nous pouvons déterminer quand chaque module $\mathcal{O}_X$
plat fini est fini localement libre en utilisant le lemme précédent.

\begin{lemma}
\label{morphisms-lemma-finite-flat-is-finite-locally-free}
Soit $X$ un schéma. Les conditions suivantes sont équivalentes
\begin{enumerate}
\item tout module $\mathcal{O}_X$ plat fini quasi-cohérent est fini
localement libre, et
\item tout sous-ensemble fermé $Z \subset X$ qui est fermé sous générisations est
ouvert.
\end{enumerate}
\end{lemma}

\begin{proof}
Dans le cas affine, il
s'agit de l'algèbre, Lemme \ref{algebra-lemma-finite-flat-module-finitely-presented}. Le cas schéma ne découle
pas directement du cas affine, nous répétons donc simplement
les arguments.

\medskip\noindent
Supposons (1). Considérons une immersion fermée $i : Z \to X$ telle que
$i$ soit plate. Alors $i_*\mathcal{O}_Z$ est quasi-cohérent et plat,
donc fini localement libre par (1). Ainsi $Z = \text{Supp}(i_*\mathcal{O}_Z)$ est également
ouvert et on voit que (2) est vrai. D'où l'implication (1) $\Rightarrow$
(2) découle de la caractérisation des immersions
plates fermées dans le lemme \ref{morphisms-lemma-characterize-flat-closed-immersions}.

\medskip\noindent
À l'inverse, supposons que $X$
satisfait (2). Soit $\mathcal{F}$ un module $\mathcal{O}_X$
plat fini, quasi-cohérent. Le support $Z = \text{Supp}(\mathcal{F})$ de $\mathcal{F}$
est fermé, voir Modules, Lemme \ref{modules-lemma-support-finite-type-closed}.
En revanche, si $x \leadsto x'$ est une spécialisation,
alors par Algèbre, Lemme \ref{algebra-lemma-finite-flat-local}
le module $\mathcal{F}_{x'}$ est libre sur $\mathcal{O}_{X, x'}$, et
$$
\mathcal{F}_x =
\mathcal{F}_{x'} \otimes_{\mathcal{O}_{X, x'}} \mathcal{O}_{X, x}.
$$
Donc
$x' \in \text{Supp}(\mathcal{F}) \Rightarrow x \in \text{Supp}(\mathcal{F})$, autrement dit, le support est fermé sous
générisation. Comme $X$ satisfait
(2), nous voyons que le support de $\mathcal{F}$
est ouvert et fermé. Les modules $\wedge^i(\mathcal{F})$, $i = 1, 2, 3, \ldots$
sont également des modules $\mathcal{O}_X$
plats
finis et quasi-cohérents, voir
Modules,
section \ref{modules-section-symmetric-exterior}. Notons que
$\text{Supp}(\wedge^{i + 1}(\mathcal{F})) \subset
\text{Supp}(\wedge^i(\mathcal{F}))$. On voit donc qu'il
existe une décomposition
$$
X = U_0 \amalg U_1 \amalg U_2 \amalg \ldots
$$
par sous-ensembles ouverts et fermés telle
que le support de $\wedge^i(\mathcal{F})$ est $U_i \cup U_{i + 1} \cup \ldots$ pour tout $i$.
Soit $x$ un point de $X$
et supposons $x \in U_r$. Notons que
$\wedge^i(\mathcal{F})_x \otimes \kappa(x) =
\wedge^i(\mathcal{F}_x \otimes \kappa(x))$. Par conséquent, $x \in U_r$
implique que $\mathcal{F}_x \otimes \kappa(x)$ est un espace vectoriel de dimension
$r$. Par le lemme de Nakayama, voir
Algèbre, Lemme \ref{algebra-lemma-NAK}, on peut choisir un voisinage ouvert affine $U \subset U_r \subset X$ de $x$
et des sections $s_1, \ldots, s_r \in \mathcal{F}(U)$ tels que l'application
induite
$$
\mathcal{O}_U^{\oplus r} \longrightarrow \mathcal{F}|_U, \quad
(f_1, \ldots, f_r) \longmapsto \sum f_i s_i
$$
soit surjectif. Cela signifie
que $\wedge^r(\mathcal{F}|_U)$ est un module $\mathcal{O}_U$
plat fini et quasi-cohérent dont le
support est entièrement $U$. Par ce
qui précède, il est généré
par un seul élément, à savoir $s_1 \wedge \ldots \wedge s_r$. D'où
$\wedge^r(\mathcal{F}|_U) \cong \mathcal{O}_U/\mathcal{I}$ pour un faisceau quasi-cohérent
d'idéaux $\mathcal{I}$ tel que $\mathcal{O}_U/\mathcal{I}$ est plat
sur $\mathcal{O}_U$ et tel
que $V(\mathcal{I}) = U$. Il s'ensuit que
$\mathcal{I} = 0$ en appliquant le Lemme \ref{morphisms-lemma-characterize-flat-closed-immersions}.
Ainsi $s_1 \wedge \ldots \wedge s_r$ est une base pour
$\wedge^r(\mathcal{F}|_U)$ et Il s'ensuit que l'application affichée est aussi bien injective
que surjective. Cela prouve que $\mathcal{F}$ est fini localement libre comme
souhaité.
\end{proof}






\section{Platitude générique}
\label{morphisms-section-generic-flatness}

\noindent
Un schéma de type fini sur une base intègre est plat sur un ouvert dense
ouverte de la base.
En Algèbre, section \ref{algebra-section-generic-flatness} nous avons prouvé une version
noéthérienne, une version pour les morphismes de présentation finie et une
version générale. Nous ne faisons ici qu'énoncer et prouver la version générale.
Cependant, il s'avère que cela sera remplacé
par la proposition \ref{morphisms-proposition-generic-flatness-reduced} qui montre que le résultat
est valable si nous supposons uniquement que la base est réduite.

\begin{proposition}[Platitude générique]
\label{morphisms-proposition-generic-flatness}
Soit $f : X \to S$ un morphisme de schémas. Soit
$\mathcal{F}$ un faisceau quasi-cohérent de $\mathcal{O}_X$-modules. Supposons
que
\begin{enumerate}
\item $S$ est intègre,
\item $f$ est de type fini et
\item $\mathcal{F}$ est un module de type fini $\mathcal{O}_X$.
\end{enumerate}
Alors il existe un sous-schéma dense ouvert $U \subset S$ tel que
$X_U \to U$ est plat et de présentation finie et tel
que $\mathcal{F}|_{X_U}$ est plat sur $U$ et de présentation
finie sur $\mathcal{O}_{X_U}$.
\end{proposition}

\begin{proof}
Comme $S$ est intègre, il est irréductible
(voir Propriétés, Lemme \ref{properties-lemma-characterize-integral}) et tout ouvert non vide
est dense. On peut donc remplacer $S$ par un
ouvert affine de $S$ et supposer que $S = \Spec(A)$ est affine.
Comme $S$ est intègre, nous voyons que $A$
est un domaine. Comme $f$ est de type fini, il est quasi-compact,
donc $X$ est quasi-compact. On peut
donc trouver un recouvrement ouvert affine fini $X = \bigcup_{i = 1, \ldots, n} X_i$.
Écrivons $X_i = \Spec(B_i)$. Alors $B_i$ est une algèbre
de type fini $A$, voir Lemme \ref{morphisms-lemma-locally-finite-type-characterize}. De
plus, il existe des $B_i$-modules
$M_i$ de type fini tels que $\mathcal{F}|_{X_i}$ est le
faisceau quasi-cohérent associé au $B_i$-module $M_i$,
voir Propriétés, Lemme \ref{properties-lemma-finite-type-module}. Ensuite,
pour chaque paire d'indices $i, j$ choisissez un $I_{ij} \subset B_i$
idéal tel que $X_i \setminus X_i \cap X_j = V(I_{ij})$ à l'intérieur de $X_i = \Spec(B_i)$.
Définissez $M_{ij} = B_i/I_{ij}$ et considérez-le comme
un module $B_i$. Alors $V(I_{ij}) = \text{Supp}(M_{ij})$ et $M_{ij}$
sont un module $B_i$ fini.

\medskip\noindent
À ce stade, nous
appliquons l'algèbre, Lemme \ref{algebra-lemma-generic-flatness} aux
paires $(A \to B_i, M_{ij})$ et aux paires $(A \to B_i, M_i)$.
Ainsi on obtient un $f_{ij}, f_i \in A$
non nul tel que (a) $A_{f_{ij}} \to B_{i, f_{ij}}$ est plat et de présentation
finie et $M_{ij, f_{ij}}$ est plat sur $A_{f_{ij}}$ et de
présentation finie sur $B_{i, f_{ij}}$, et (b) $B_{i, f_i}$ est plat
et de présentation finie sur $A_f$ et $M_{i, f_i}$
est plat et de présentation finie sur $B_{i, f_i}$. Posons $f = (\prod f_i) (\prod f_{ij})$.
Nous affirmons que
prendre $U = D(f)$ fonctionne.

\medskip\noindent
Pour prouver notre affirmation, nous pouvons remplacer
$A$ par $A_f$, c'est-à-dire effectuer
le changement de base par $U = \Spec(A_f) \to S$. Après ce changement
de base on voit que chacun des $A \to B_i$ est plat et de présentation
finie et que $M_i$, $M_{ij}$
sont plats sur $A$ et de présentation finie
sur $B_i$. Cela prouve déjà que $X \to S$ est quasi-compact,
localement de présentation finie, plate, et que $\mathcal{F}$
est plat sur $S$ et de présentation finie sur
$\mathcal{O}_X$,
voir Lemme \ref{morphisms-lemma-locally-finite-presentation-characterize} et Propriétés, Lemme \ref{properties-lemma-finite-presentation-module}.
Puisque $M_{ij}$ est de présentation finie
sur $B_i$ on voit que $X_i \cap X_j = X_i \setminus \text{Supp}(M_{ij})$ est un ouvert
quasi-compact
de $X_i$, voir Algèbre, Lemme \ref{algebra-lemma-support-finite-presentation-constructible}. On voit
donc que $X \to S$ est quasi-séparé
par Schémas, Lemme \ref{schemes-lemma-characterize-quasi-separated}. Cela prouve
la proposition.
\end{proof}

\noindent
Il s'avère effectivement qu'il existe également une version de platitude
générique sur une base réduite arbitraire. C'est ici.

\begin{proposition}[Platitude générique, cas réduit]
\label{morphisms-proposition-generic-flatness-reduced}
Soit $f : X \to S$ un morphisme de schémas. Soit
$\mathcal{F}$ un faisceau quasi-cohérent de $\mathcal{O}_X$-modules. Supposons
que
\begin{enumerate}
\item $S$ est réduit,
\item $f$ est de type fini et
\item $\mathcal{F}$ est un module de type fini $\mathcal{O}_X$-module.
\end{enumerate}
Il existe alors un sous-schéma dense ouvert $U \subset S$ tel que
$X_U \to U$ est plat et de présentation finie et tel
que $\mathcal{F}|_{X_U}$ est plat sur $U$ et de présentation
finie sur $\mathcal{O}_{X_U}$.
\end{proposition}

\begin{proof}
Pour le lecteur impatient : cette démonstration est une
répétition de la démonstration de la
proposition
\ref{morphisms-proposition-generic-flatness} utilisant l'algèbre, Lemme \ref{algebra-lemma-generic-flatness-reduced}
au
lieu de l'algèbre, Lemme \ref{algebra-lemma-generic-flatness}.

\medskip\noindent
Puisque le fait d'être plat et d'être de présentation finie est
local
sur la base, voir Lemmes \ref{morphisms-lemma-flat-module-characterize} et \ref{morphisms-lemma-locally-finite-presentation-characterize},
on peut travailler localement sur les affines
de $S$. Ainsi on peut supposer que $S = \Spec(A)$,
où $A$ est un anneau réduit (voir Propriétés,
Lemme \ref{properties-lemma-characterize-reduced}). Comme $f$ est de type fini,
il est quasi-compact, donc $X$ est quasi-compact. On peut
donc trouver un recouvrement ouvert affine
fini $X = \bigcup_{i = 1, \ldots, n} X_i$. Écrivons $X_i = \Spec(B_i)$. Alors $B_i$ est une
algèbre de type fini $A$, voir
Laisse-moi \ref{morphisms-lemma-locally-finite-type-characterize}. De plus il existe des
$B_i$-modules $M_i$
de type fini tels que $\mathcal{F}|_{X_i}$ est le faisceau
quasi-cohérent associé au $B_i$-module $M_i$,
voir Propriétés, Lemme \ref{properties-lemma-finite-type-module}. Ensuite, pour
chaque paire d'indices $i, j$ choisissez un idéal $I_{ij} \subset B_i$
tel que $X_i \setminus X_i \cap X_j = V(I_{ij})$ à l'intérieur de $X_i = \Spec(B_i)$.
Définissez $M_{ij} = B_i/I_{ij}$ et considérez-le comme
un module $B_i$. Alors $V(I_{ij}) = \text{Supp}(M_{ij})$ et $M_{ij}$
sont un module $B_i$ fini.

\medskip\noindent
À ce stade, nous
appliquons l'algèbre, Lemme \ref{algebra-lemma-generic-flatness-reduced} aux paires
$(A \to B_i, M_{ij})$ et aux paires $(A \to B_i, M_i)$. On obtient ainsi
des ouverts denses $U(A \to B_i, M_{ij}) \subset S$
et des ouverts denses $U(A \to B_i, M_i) \subset S$
avec une notation comme en Algèbre, Equation
(\ref{algebra-equation-good-locus}). Puisqu'une intersection
finie d'ouverts denses est ouverte dense, on voit que
$$
U =
\bigcap\nolimits_{i, j} U(A \to B_i, M_{ij})
\quad\cap\quad
\bigcap\nolimits_i U(A \to B_i, M_i)
$$
est ouverte et dense dans $S$. Nous affirmons que $U$ est l'ouverture souhaitée.

\medskip\noindent
Choisissons $u \in U$. Par définition des lieux $U(A \to B_i, M_{ij})$
et $U(A \to B, M_i)$, il existe $f_{ij}, f_i \in A$ tels que (a)
$u \in D(f_i)$ et $u \in D(f_{ij})$, (b) $A_{f_{ij}} \to B_{i, f_{ij}}$
est plat et de présentation finie et $M_{ij, f_{ij}}$ est plat
sur $A_{f_{ij}}$ et de présentation finie sur $B_{i, f_{ij}}$, et (c)
$B_{i, f_i}$ est plat
et de présentation finie sur $A_{f_i}$ et $M_{i, f_i}$ est de
présentation finie sur $B_{i, f_i}$ et plat sur $A_{f_i}$.
Définissez $f = (\prod f_i) (\prod f_{ij})$.
Il suffit maintenant
de prouver que $X \to S$ est plat et de présentation finie
sur $D(f)$ et que $\mathcal{F}$ restreint à $X_{D(f)}$ est plat
sur $D(f)$ et de présentation finie sur le faisceau structural de
$X_{D(f)}$.

\medskip\noindent
On peut donc remplacer $A$ par
$A_f$, c'est-à-dire effectuer le changement
de base par $\Spec(A_f) \to S$. Après ce changement de
base on voit que chacun des $A \to B_i$ est plat et de présentation finie
et que $M_i$, $M_{ij}$ sont
plats sur $A$ et de présentation finie sur
$B_i$. Cela prouve déjà que $X \to S$ est quasi-compact,
localement de présentation finie, plate, et que $\mathcal{F}$
est plat sur $S$ et de présentation finie sur
$\mathcal{O}_X$,
voir Lemme \ref{morphisms-lemma-locally-finite-presentation-characterize} et Propriétés, Lemme \ref{properties-lemma-finite-presentation-module}.
Puisque $M_{ij}$ est de présentation finie
sur $B_i$ on voit que $X_i \cap X_j = X_i \setminus \text{Supp}(M_{ij})$ est un ouvert
quasi-compact
de $X_i$, voir Algèbre, Lemme \ref{algebra-lemma-support-finite-presentation-constructible}. On voit
donc que $X \to S$ est quasi-séparé
par Schémas, Lemme \ref{schemes-lemma-characterize-quasi-separated}. Cela prouve
la proposition.
\end{proof}









\begin{remark}
\label{morphisms-remark-flattening}
Les résultats ci-dessus sont un premier pas vers des techniques d'aplatissement plus raffinées
des morphismes de schémas. L'article \cite{GruRay} de Raynaud et Gruson contient
de nombreux résultats merveilleux dans ce sens.
\end{remark}







\section{Morphismes et dimensions des fibres}
\label{morphisms-section-dimension-fibres}

\noindent
Soit $X$ un espace topologique, et $x \in X$.
Rappelons que nous avons défini $\dim_x(X)$ comme le minimum
des dimensions des voisinages ouverts de $x$
dans $X$. Voir Topologie, Définition \ref{topology-definition-Krull}.

\begin{lemma}
\label{morphisms-lemma-dimension-fibre-at-a-point}
Soit $f : X \to S$ un morphisme de schémas. Soit $x \in X$
et définissez $s = f(x)$. Supposons
que $f$ soit localement de type fini.
Puis
$$
\dim_x(X_s) =
\dim(\mathcal{O}_{X_s, x}) + \text{trdeg}_{\kappa(s)}(\kappa(x)).
$$
\end{lemma}

\begin{proof}
Cela se réduit immédiatement au cas $S = s$, et
$X$ affine. Dans ce
cas, le résultat découle de l'algèbre, Lemme \ref{algebra-lemma-dimension-at-a-point-finite-type-field}.
\end{proof}

\begin{lemma}
\label{morphisms-lemma-dimension-fibre-at-a-point-additive}
Soient $f : X \to Y$ et $g : Y \to S$ des morphismes de schémas. Soit $x \in X$
et définissez $y = f(x)$, $s = g(y)$. Supposons que $f$
et $g$ localement de type fini.
Alors
$$
\dim_x(X_s) \leq \dim_x(X_y) + \dim_y(Y_s).
$$
De plus, l'égalité est vraie si $\mathcal{O}_{X_s, x}$ est
plat sur $\mathcal{O}_{Y_s, y}$, ce qui est vrai par exemple si $\mathcal{O}_{X, x}$
est plat sur $\mathcal{O}_{Y, y}$.
\end{lemma}

\begin{proof}
Notons que $\text{trdeg}_{\kappa(s)}(\kappa(x)) =
\text{trdeg}_{\kappa(y)}(\kappa(x)) + \text{trdeg}_{\kappa(s)}(\kappa(y))$. Ainsi, par le Lemme \ref{morphisms-lemma-dimension-fibre-at-a-point},
l'instruction est équivalente
à
$$
\dim(\mathcal{O}_{X_s, x})
\leq
\dim(\mathcal{O}_{X_y, x}) + \dim(\mathcal{O}_{Y_s, y}).
$$
Pour cela voir Algèbre, Lemme \ref{algebra-lemma-dimension-base-fibre-total}.
Pour le cas
plat, voir Algèbre, Lemme \ref{algebra-lemma-dimension-base-fibre-equals-total}.
\end{proof}

\begin{lemma}
\label{morphisms-lemma-dimension-fibre-after-base-change}
Soit
$$
\xymatrix{
X' \ar[r]_{g'} \ar[d]_{f'} & X \ar[d]^f \\
S' \ar[r]^g & S
}
$$
un diagramme produit fibré de schémas. Supposons que $f$ localement de
type fini. Supposons que $x' \in X'$, $x = g'(x')$, $s' = f'(x')$
et $s = g(s') = f(x)$. Alors
\begin{enumerate}
\item $\dim_x(X_s) = \dim_{x'}(X'_{s'})$,
\item si $F$ est la fibre du morphisme $X'_{s'} \to X_s$ sur $x$,
alors
$$
\dim(\mathcal{O}_{F, x'}) =
\dim(\mathcal{O}_{X'_{s'}, x'}) - \dim(\mathcal{O}_{X_s, x}) =
\text{trdeg}_{\kappa(s)}(\kappa(x)) -
\text{trdeg}_{\kappa(s')}(\kappa(x'))
$$
notamment $\dim(\mathcal{O}_{X'_{s'}, x'}) \geq \dim(\mathcal{O}_{X_s, x})$
et $\text{trdeg}_{\kappa(s)}(\kappa(x)) \geq
\text{trdeg}_{\kappa(s')}(\kappa(x'))$.
\item étant donné $s', s, x$,
il existe un choix de $x'$
tel que $\dim(\mathcal{O}_{X'_{s'}, x'}) = \dim(\mathcal{O}_{X_s, x})$ et $\text{trdeg}_{\kappa(s)}(\kappa(x)) = \text{trdeg}_{\kappa(s')}(\kappa(x'))$.
\end{enumerate}
\end{lemma}

\begin{proof}
La partie (1) découle
immédiatement
de l'algèbre, Lemme \ref{algebra-lemma-dimension-at-a-point-preserved-field-extension}.
Parties (2)
et (3) d'Algèbre, Lemme \ref{algebra-lemma-inequalities-under-field-extension}.
\end{proof}

\noindent
Le lemme suivant découle d'un résultat algébrique non trivial. À savoir
la version algébrique du Théorème principal de Zariski.

\begin{lemma}
\label{morphisms-lemma-openness-bounded-dimension-fibres}
\begin{reference}
\cite[IV Théorème 13.1.3]{EGA}
\end{reference}
Soit $f : X \to S$ un morphisme de schémas. Soit $n \geq 0$.
Supposons que $f$ soit localement de type fini.
L'ensemble
$$
U_n = \{x \in X \mid \dim_x X_{f(x)} \leq n\}
$$
est ouvert dans $X$.
\end{lemma}

\begin{proof}
Ceci est immédiat
de
l'Algèbre, Lemme \ref{algebra-lemma-dimension-fibres-bounded-open-upstairs}
\end{proof}

\begin{lemma}
\label{morphisms-lemma-morphism-finite-type-bounded-dimension}
Soit $f : X \to Y$ un morphisme de type fini avec $Y$ quasi-compact. Alors
la dimension des fibres de $f$ est bornée.
\end{lemma}

\begin{proof}
Par le Lemme \ref{morphisms-lemma-openness-bounded-dimension-fibres} l'ensemble $U_n \subset X$ des points où
la dimension de la fibre est $\leq n$ est ouvert. Puisque $f$
est de type fini, chaque point est contenu dans un $U_n$ (car
la dimension d'une algèbre de type fini sur un corps est finie).
Puisque $Y$ est quasi-compact et $f$ est de type fini, on voit
que $X$ est quasi-compact. D'où $X = U_n$ pour certains
$n$.
\end{proof}

\begin{lemma}
\label{morphisms-lemma-openness-bounded-dimension-fibres-finite-presentation}
Soit $f : X \to S$ un morphisme de schémas. Soit $n \geq 0$.
Supposons que $f$ soit localement de présentation finie.
Le
$$
U_n = \{x \in X \mid \dim_x X_{f(x)} \leq n\}
$$
ouvert du lemme \ref{morphisms-lemma-openness-bounded-dimension-fibres} est rétrocompact dans $X$.
(Voir Topologie, Définition \ref{topology-definition-quasi-compact}.)
\end{lemma}

\begin{proof}
L'espace topologique $X$ a une base pour sa topologie
constituée d'ouverts affines $U \subset X$
telles que le morphisme induit $f|_U : U \to S$ se factorise
par un ouvert affine $V \subset S$. Il suffit donc de montrer
que $U \cap U_n$ est quasi-compact pour
un tel $U$. Notez que $U_n \cap U$ est identique
au $\{x \in U \mid \dim_x U_{f(x)} \leq n\}$ ouvert. Cela nous ramène au cas où $X$
et
$S$ sont affines. Dans ce cas, le lemme découle de
l'algèbre, Lemme \ref{algebra-lemma-dimension-fibres-bounded-quasi-compact-open-upstairs} (et Lemme \ref{morphisms-lemma-locally-finite-presentation-characterize}).
\end{proof}

\begin{lemma}
\label{morphisms-lemma-dimension-fibre-specialization}
Soit $f : X \to S$ un morphisme de schémas. Soit $x \leadsto x'$
une spécialisation non triviale de points dans $X$ situés sur
le même point $s \in S$. Supposons que $f$ soit localement de type fini.
Puis
\begin{enumerate}
\item $\dim_x(X_s) \leq \dim_{x'}(X_s)$,
\item $\dim(\mathcal{O}_{X_s, x}) < \dim(\mathcal{O}_{X_s, x'})$ et
\item $\text{trdeg}_{\kappa(s)}(\kappa(x)) >
\text{trdeg}_{\kappa(s)}(\kappa(x'))$.
\end{enumerate}
\end{lemma}

\begin{proof}
La partie (1) découle du fait que toute ouverture de $X_s$ contenant
$x'$ contient également $x$. La partie (2) suit puisque
$\mathcal{O}_{X_s, x}$ est une localisation de $\mathcal{O}_{X_s, x'}$ à un idéal premier,
donc toute chaîne d'idéaux premiers dans $\mathcal{O}_{X_s, x}$ fait partie d'une
chaîne strictement plus longue de nombres premiers dans $\mathcal{O}_{X_s, x'}$. La dernière
inégalité découle de l'algèbre, Lemme \ref{algebra-lemma-tr-deg-specialization}.
\end{proof}






\section{Morphismes de dimension relative donnée}
\label{morphisms-section-relative-dimension}

\noindent
Afin de pouvoir parler confortablement de morphismes d'une dimension
relative donnée nous introduisons la notion suivante.

\begin{definition}
\label{morphisms-definition-relative-dimension-d}
Soit $f : X \to S$ un morphisme de schémas. Supposons que
$f$ soit localement de type fini.
\begin{enumerate}
\item Nous disons que $f$ est de dimension {\it relative $\leq d$ à
$x$} si $\dim_x(X_{f(x)}) \leq d$.
\item Nous disons que $f$ est de dimension {\it relative
$\leq d$} si $\dim_x(X_{f(x)}) \leq d$ pour tout $x \in X$.
\item On dit que $f$ est de dimension {\it relative $d$}
si toutes les fibres non vides $X_s$ sont équidimensionnelles de dimension $d$.
\end{enumerate}
\end{definition}

\noindent
Ce n'est pas une notion particulièrement efficace, mais elle fonctionne bien dans un
certain nombre de situations.

\begin{lemma}
\label{morphisms-lemma-base-change-relative-dimension-d}
Soit $f : X \to S$ un morphisme de schémas qui est localement de type fini. Si $f$
a une dimension relative $d$, tout changement de base de $f$ aussi.
Idem pour la cote relative $\leq d$.
\end{lemma}

\begin{proof}
Ceci est immédiat
à partir du lemme \ref{morphisms-lemma-dimension-fibre-after-base-change}.
\end{proof}

\begin{lemma}
\label{morphisms-lemma-composition-relative-dimension-d}
Soit $f : X \to Y$, $g : Y \to Z$ être localement de type fini. Si $f$
a une dimension relative $\leq d$ et $g$ a une dimension relative $\leq e$
alors $g \circ f$ a une dimension relative $\leq d + e$.
Si
\begin{enumerate}
\item $f$ a une dimension relative $d$,
\item $g$ a une dimension relative $e$ et que
\item $f$ est plat,
\end{enumerate}
alors $g \circ f$ a une dimension relative $d + e$.
\end{lemma}

\begin{proof}
Ceci est immédiat à partir du lemme \ref{morphisms-lemma-dimension-fibre-at-a-point-additive}.
\end{proof}

\noindent
En général, il n'est pas possible de décomposer un morphisme
en morceaux où la dimension relative est donnée. Cependant,
c'est possible si le morphisme comporte des fibres de De Cohen--Macaulay
et est plat de présentation finie.

\begin{lemma}
\label{morphisms-lemma-flat-finite-presentation-CM-fibres-relative-dimension}
\begin{slogan}
Les morphismes de Cohen--Macaulay se décomposent en ouverts et fermés de dimension relative pure
\end{slogan}
Soit $f : X \to S$ un morphisme de schémas. Supposons
que
\begin{enumerate}
\item $f$ est plat,
\item $f$ est localement de présentation finie, et
\item pour tout $s \in S$ la fibre $X_s$ est
De Cohen--Macaulay (Propriétés, Définition \ref{properties-definition-Cohen-Macaulay})
\end{enumerate}
Alors il existe des ouverts et sous-schémas fermés $X_d \subset X$
tels que $X = \coprod_{d \geq 0} X_d$ et $f|_{X_d} : X_d \to S$ ont une dimension
relative $d$.
\end{lemma}

\begin{proof}
Ceci est immédiat
de l'algèbre,
Lemme \ref{algebra-lemma-relative-dimension-CM}.
\end{proof}

\begin{lemma}
\label{morphisms-lemma-locally-quasi-finite-rel-dimension-0}
Soit $f : X \to S$ un morphisme de schémas. Supposons que
$f$ soit localement de type fini. Soit
$x \in X$ avec $s = f(x)$. Alors $f$
est quasi-fini à $x$ si et seulement si $\dim_x(X_s) = 0$. En particulier,
$f$ est localement quasi-fini si et seulement si $f$ a une dimension relative
$0$.
\end{lemma}

\begin{proof}
Première démonstration.
Si $f$ est quasi-fini en $x$ alors
$\kappa(x)$ est
une extension finie de $\kappa(s)$ (par
le lemme \ref{morphisms-lemma-residue-field-quasi-finite}) et $x$
est isolé dans $X_s$ (par le lemme
\ref{morphisms-lemma-quasi-finite-at-point-characterize}), donc $\dim_x(X_s) = 0$
par le lemme \ref{morphisms-lemma-dimension-fibre-at-a-point}. A l'inverse,
si $\dim_x(X_s) = 0$ alors par le
lemme \ref{morphisms-lemma-dimension-fibre-at-a-point} on voit que $\kappa(s) \subset \kappa(x)$
est algébrique et il n'y a pas d'autres points
de $X_s$ spécialisés dans $x$.
Ainsi $x$ est fermé
dans sa fibre par le lemme \ref{morphisms-lemma-algebraic-residue-field-extension-closed-point-fibre} et par le lemme
\ref{morphisms-lemma-quasi-finite-at-point-characterize}
(3) Nous concluons que $f$
est quasi-fini en $x$.

\medskip\noindent
Deuxième démonstration. La fibre $X_s$ est un schéma
localement de type fini sur un corps, donc localement
noethérien (Lemme
\ref{morphisms-lemma-finite-type-noetherian}). Le résultat découle maintenant
du lemme \ref{morphisms-lemma-quasi-finite-at-point-characterize} et de Propriétés, Lemme \ref{properties-lemma-isolated-dim-0-locally-noetherian}.
\end{proof}

\begin{lemma}
\label{morphisms-lemma-rel-dimension-dimension}
Soit $f : X \to Y$ un morphisme de schémas localement noethériens qui est plat,
localement de type fini et de dimension relative $d$. Pour chaque
point $x$ dans $X$ avec l'image
$y$ dans $Y$ nous avons $\dim_x(X) = \dim_y(Y) + d$.
\end{lemma}

\begin{proof}
Après avoir réduit $X$ et $Y$ pour ouvrir les
voisinages de $x$ et $y$, on peut supposer que
$\dim(X) = \dim_x(X)$ et $\dim(Y) = \dim_y(Y)$, par définition de la dimension
d'un schéma en un point (Propriétés, Définition
\ref{properties-definition-dimension}). Le morphisme $f$
est ouvert par les Lemmes \ref{morphisms-lemma-noetherian-finite-type-finite-presentation} et \ref{morphisms-lemma-fppf-open}.
On peut donc réduire
$Y$ pour faire en sorte que $f$ soit
surjectif. Reste à montrer que $\dim(X) = \dim(Y) + d$.

\medskip\noindent
Soit $a$ un point dans $X$ avec l'image
$b$ dans $Y$. Par Algèbre, Lemme \ref{algebra-lemma-dimension-base-fibre-equals-total},
$$
\dim(\mathcal{O}_{X,a}) = \dim(\mathcal{O}_{Y,b}) + \dim(\mathcal{O}_{X_b, a}).
$$
Prenant la suprématie sur tous les points $a$ dans
$X$, Il s'ensuit que $\dim(X) = \dim(Y) + d$, comme
on veut, voir Propriétés, Lemme \ref{properties-lemma-dimension}.
\end{proof}











\section{Morphismes syntomiques}
\label{morphisms-section-syntomic}

\noindent
Une algèbre $A$ sur un corps $k$ est appelée
une {\it intersection globale
complète sur $k$} si $A \cong k[x_1, \ldots, x_n]/(f_1, \ldots, f_c)$ et $\dim(A) = n - c$.
Une algèbre $A$ sur un corps $k$ est appelée une {\it
intersection complète locale} si $\Spec(A)$
peut être couverte par des ouvertures standard dont chacune est
une intersection complète globale sur $k$. Voir Algèbre,
section \ref{algebra-section-lci}. Rappelons qu'un homomorphisme d'anneaux $R \to A$
est {\it syntomique} s'il est
de présentation finie, plat avec des anneaux d'intersection complets
locaux comme fibres, voir Algèbre, Définition \ref{algebra-definition-lci}.

\begin{definition}
\label{morphisms-definition-syntomic}
Soit $f : X \to S$ un morphisme de schémas.
\begin{enumerate}
\item On dit que $f$ est {\it syntomique
à $x \in X$} s'il existe un voisinage ouvert affine $\Spec(A) = U \subset X$
de $x$ et ouvert affine $\Spec(R) = V \subset S$ avec
$f(U) \subset V$ tel que l'homomorphisme d'anneaux induit $R \to A$
soit syntomique.
\item On dit que $f$ est {\it syntomique} s'il est syntomique
en tout point de $X$.
\item Si $S = \Spec(k)$ et $f$ sont syntomiques, alors on dit que $X$ est une intersection
locale complète {\it sur $k$}.
\item Un morphisme de schémas affines $f : X \to S$
est appelé {\it syntomique
standard} s'il
existe une intersection complète relative
globale
$R \to R[x_1, \ldots, x_n]/(f_1, \ldots, f_c)$ (voir Algèbre, Définition \ref{algebra-definition-relative-global-complete-intersection}) telle que
$X \to S$ est isomorphe à
$$
\Spec(R[x_1, \ldots, x_n]/(f_1, \ldots, f_c)) \to \Spec(R).
$$
\end{enumerate}
\end{definition}

\noindent
Dans la littérature, un morphisme syntomique est parfois appelé {\it
morphisme d'intersection complet local plat}. Il
s'avère qu'il s'agit d'une classe de morphismes pratique. Par exemple, on peut définir
une topologie syntomique en utilisant celles-ci, qui est plus fine que les topologies
lisses et étales, mais possède bon nombre des mêmes propriétés formelles.

\medskip\noindent Une
intersection complète relative globale (que nous avons utilisée pour définir les homomorphismes
syntomiques d'anneaux standards) est particulièrement
plate. Dans En savoir plus sur les morphismes, section \ref{more-morphisms-section-lci},
nous considérerons les morphismes $X \to S$ qui sont localement de la forme
$$
\Spec(R[x_1, \ldots, x_n]/(f_1, \ldots, f_c)) \to \Spec(R).
$$
pour certaines séquences régulières de Koszul $f_1, \ldots, f_r$ dans $R[x_1, \ldots, x_n]$. Un tel morphisme
sera appelé un {\it morphisme d'intersection complet local}.
Une fois que nous aurons cette définition en place, un morphisme sera syntomique
si et seulement si c'est un morphisme d'intersection complet, plat et local.

\medskip\noindent Notons
qu'il n'y a pas d'hypothèse de séparation ou de quasi-compacité dans la
définition d'un morphisme syntomique. La question du caractère syntomique est
donc de nature locale sur la source. Voici le résultat précis.

\begin{lemma}
\label{morphisms-lemma-syntomic-characterize}
Soit $f : X \to S$ un morphisme de schémas. Les conditions
suivantes sont équivalentes
\begin{enumerate}
\item Le morphisme $f$ est syntomique.
\item Pour toute ouverture affine $U \subset X$, $V \subset S$
avec $f(U) \subset V$ l'homomorphisme
d'anneaux $\mathcal{O}_S(V) \to \mathcal{O}_X(U)$ est syntomique.
\item Il existe un recouvrement ouvert $S = \bigcup_{j \in J} V_j$ et
des recouvrements ouverts $f^{-1}(V_j) = \bigcup_{i \in I_j} U_i$ tels que chacun
des morphismes $U_i \to V_j$, $j\in J, i\in I_j$ est
syntomique.
\item Il existe un recouvrement ouvert affine $S = \bigcup_{j \in J} V_j$
et ouvert affine revêtements $f^{-1}(V_j) = \bigcup_{i \in I_j} U_i$ tel que l'homomorphisme
d'anneaux $\mathcal{O}_S(V_j) \to \mathcal{O}_X(U_i)$ soit syntomique,
pour tout $j\in J, i\in I_j$.
\end{enumerate}
De plus, si $f$ est syntomique
alors pour tout sous-schémas ouverts $U \subset X$, $V \subset S$ avec $f(U) \subset V$
la restriction $f|_U : U \to V$ est syntomique.
\end{lemma}

\begin{proof}
Cela résulte du lemme \ref{morphisms-lemma-locally-P} si l'on montre que la propriété
``$R \to A$ est syntomique''
est locale. On vérifie les conditions (a), (b)
et (c) de la Définition \ref{morphisms-definition-property-local}.
Par Algèbre, Lemme \ref{algebra-lemma-base-change-syntomic}, la propriété d'être syntomique est stable
par changement de base ; on en déduit
la condition (a). Par
Algèbre, Lemme \ref{algebra-lemma-composition-syntomic}, elle est stable
par composition, et pour tout anneau $R$ l'homomorphisme
d'anneaux $R \to R_f$ est syntomique.
On en déduit la condition (b). Enfin, la
condition (c) résulte d'Algèbre, Lemme \ref{algebra-lemma-local-syntomic}.
\end{proof}

\begin{lemma}
\label{morphisms-lemma-composition-syntomic}
La composition de deux morphismes syntomiques est syntomique.
\end{lemma}

\begin{proof}
Dans la démonstration du lemme \ref{morphisms-lemma-syntomic-characterize} nous avons
vu qu'être syntomique est une propriété locale des homomorphismes
d'anneaux. D'où la première affirmation
du lemme découle du lemme \ref{morphisms-lemma-composition-type-P} combinée
au fait qu'être syntomique est une propriété des homomorphismes
d'anneaux qui est stable
sous composition, voir Algèbre, Lemme \ref{algebra-lemma-composition-syntomic}.
\end{proof}

\begin{lemma}
\label{morphisms-lemma-base-change-syntomic}
Le changement de base d'un morphisme qui est syntomique est syntomique.
\end{lemma}

\begin{proof}
Dans la démonstration du lemme \ref{morphisms-lemma-syntomic-characterize} nous avons
vu qu'être syntomique est une propriété locale des homomorphismes
d'anneaux. D'où
le lemme découle du lemme \ref{morphisms-lemma-composition-type-P} combiné
au fait qu'être syntomique est une propriété des homomorphismes
d'anneaux stable par
changement de base, voir Algèbre, Lemme \ref{algebra-lemma-base-change-syntomic}.
\end{proof}

\begin{lemma}
\label{morphisms-lemma-open-immersion-syntomic}
Toute immersion ouverte est syntomique.
\end{lemma}

\begin{proof}
Ceci est vrai car une immersion ouverte est un isomorphisme local.
\end{proof}

\begin{lemma}
\label{morphisms-lemma-syntomic-locally-finite-presentation}
Un morphisme syntomique est localement de présentation finie.
\end{lemma}

\begin{proof}
Vrai car un homomorphisme syntomique d'anneaux est de présentation finie par
définition.
\end{proof}

\begin{lemma}
\label{morphisms-lemma-syntomic-flat}
Un morphisme syntomique est plat.
\end{lemma}

\begin{proof}
Vrai car un homomorphisme syntomique d'anneaux est plat par définition.
\end{proof}

\begin{lemma}
\label{morphisms-lemma-syntomic-open}
Un morphisme syntomique est universellement ouvert.
\end{lemma}

\begin{proof}
Combine
Lemmes \ref{morphisms-lemma-syntomic-locally-finite-presentation},
\ref{morphisms-lemma-syntomic-flat} et
\ref{morphisms-lemma-fppf-open}.
\end{proof}

\noindent
Soit $k$ un corps. Soit $A$ une algèbre $k$ locale
essentiellement de type fini sur $k$. Rappelons que
$A$ est appelé une {\it intersection
complète sur $k$} si l'on peut écrire $A \cong R/(f_1, \ldots, f_c)$ où $R$
est un anneau local régulier essentiellement de type fini sur $k$,
et $f_1, \ldots, f_c$ est un séquence
régulière dans $R$, voir Algèbre, Définition \ref{algebra-definition-lci-local-ring}.

\begin{lemma}
\label{morphisms-lemma-local-complete-intersection}
Soit $k$ un corps.
Soit $X$ un schéma local de type fini sur $k$. Les conditions
suivantes sont équivalentes :
\begin{enumerate}
\item $X$ est une intersection complète locale sur $k$,
\item pour tout $x \in X$ il existe un voisinage
ouvert affine $U = \Spec(R) \subset X$ de $x$
tel que $R \cong k[x_1, \ldots, x_n]/(f_1, \ldots, f_c)$ est une intersection
complète globale sur $k$, et
\item pour chaque $x \in X$ l'anneau local $\mathcal{O}_{X, x}$ est une intersection
complète sur $k$.
\end{enumerate}
\end{lemma}

\begin{proof}
Les résultats d'algèbre correspondants peuvent
être trouvés dans Algèbre, Lemmes \ref{algebra-lemma-lci-at-prime}
et \ref{algebra-lemma-lci-global}.
\end{proof}

\noindent
Le lemme suivant dit localement que tout morphisme syntomique est standard syntomique. Nous pouvons
donc utiliser les morphismes syntomiques standards comme {\it modèle
local} pour un morphisme syntomique. De plus, il dit qu'un morphisme plat de présentation
finie est syntomique si et seulement si les fibres sont des schémas d'intersection
complets locaux.

\begin{lemma}
\label{morphisms-lemma-syntomic-locally-standard-syntomic}
Soit $f : X  \to S$ un morphisme de schémas. Soit $x \in X$ un point d'image
$s = f(x)$. Soit $V \subset S$ un voisinage ouvert affine de $s$. Les conditions
suivantes sont équivalentes
\begin{enumerate}
\item Le morphisme $f$ est syntomique en $x$.
\item Il existe un $U \subset X$ ouvert affine avec $x \in U$ et
$f(U) \subset V$ tel que $f|_U : U \to V$ soit syntomique standard.
\item Le morphisme $f$ est de présentation finie
à $x$, l'homomorphisme d'anneaux locaux
$\mathcal{O}_{S, s} \to \mathcal{O}_{X, x}$ est plat et $\mathcal{O}_{X, x}/\mathfrak m_s \mathcal{O}_{X, x}$ est
une intersection complète sur $\kappa(s)$
(voir Algèbre, Définition \ref{algebra-definition-lci-local-ring}).
\end{enumerate}
\end{lemma}

\begin{proof}
Dérivé des définitions et
de l'algèbre, Lemme \ref{algebra-lemma-syntomic}.
\end{proof}

\begin{lemma}
\label{morphisms-lemma-syntomic-flat-fibres}
Soit $f : X \to S$ un morphisme de schémas. Si $f$
est plat, localement de présentation finie, et que toutes les fibres $X_s$
sont des intersections locales complètes, alors $f$ est
syntomique.
\end{lemma}

\begin{proof}
Cela résulte des Lemmes
\ref{morphisms-lemma-local-complete-intersection} et
\ref{morphisms-lemma-syntomic-locally-standard-syntomic}, ainsi que des
isomorphismes d'anneaux locaux
$
\mathcal{O}_{X, x}/\mathfrak m_s \mathcal{O}_{X, x}
\cong
\mathcal{O}_{X_s, x}
$.
\end{proof}

\begin{lemma}
\label{morphisms-lemma-set-points-where-fibres-lci}
Soit $f : X \to S$ un morphisme de schémas. Supposons
que $f$ soit localement de type fini. La formation de l'ensemble
$$
T = \{x \in X \mid \mathcal{O}_{X_{f(x)}, x}
\text{ est d'intersection complète sur }\kappa(f(x))\}
$$
commute à tout changement de base :
pour tout morphisme $g : S' \to S$, considérons
le changement de base $f' : X' \to S'$ de $f$ et la
projection $g' : X' \to X$. Alors l'ensemble correspondant $T'$
pour le morphisme $f'$ est égal à $T' = (g')^{-1}(T)$. En particulier,
si $f$ est supposé plat, et localement de présentation finie
alors il en va de même pour l'ensemble ouvert de points où $f$
est syntomique.
\end{lemma}

\begin{proof}
Soit $s' \in S'$ un point et soit $s = g(s')$. On a alors
$$
X'_{s'} =
\Spec(\kappa(s')) \times_{\Spec(\kappa(s))} X_s
$$
Autrement dit les fibres du changement de base sont les changements
de base des fibres. La première partie
est donc équivalente à l'Algèbre, Lemme \ref{algebra-lemma-lci-field-change-local}.
La deuxième partie découle de la première car dans
ce cas $T$ est l'ensemble des points où $f$
est syntomique d'après le Lemme \ref{morphisms-lemma-syntomic-locally-standard-syntomic}.
\end{proof}

\begin{lemma}
\label{morphisms-lemma-standard-syntomic-relative-dimension}
Soit $R$ un
anneau. Soit $R \to A = R[x_1, \ldots, x_n]/(f_1, \ldots, f_c)$ une intersection complète relative
globale. Définissez $S = \Spec(R)$ et $X = \Spec(A)$.
Considérons le morphisme $f : X \to S$
associé à l'homomorphisme d'anneaux $R \to A$.
La fonction $x \mapsto \dim_x(X_{f(x)})$ est constante de valeur $n - c$.
\end{lemma}

\begin{proof}
En algèbre,
la Définition \ref{algebra-definition-relative-global-complete-intersection} $R \to A$ étant une intersection complète
globale relative signifie que tous les anneaux de
fibres non nuls ont une dimension $n - c$.
Ainsi, pour un $\mathfrak p$ premier de $R$
l'anneau de fibres $\kappa(\mathfrak p)[x_1, \ldots, x_n]/(\overline{f}_1, \ldots, \overline{f}_c)$ est soit nul, soit un anneau d'intersection
complet global de dimension $n - c$. D'après la discussion
qui suit l'Algèbre,
Définition \ref{algebra-definition-lci-field} cela implique qu'elle
est équidimensionnelle de dimension $n - c$. D'où
le lemme.
\end{proof}

\begin{lemma}
\label{morphisms-lemma-syntomic-relative-dimension}
Soit $f : X \to S$ un morphisme syntomique. La fonction
$x \mapsto \dim_x(X_{f(x)})$ est localement constante sur $X$.
\end{lemma}

\begin{proof}
Par le Lemme \ref{morphisms-lemma-syntomic-locally-standard-syntomic} le morphisme $f$
ressemble localement à un morphisme
syntomique d'affines standard. Le
résultat découle donc du lemme \ref{morphisms-lemma-standard-syntomic-relative-dimension}.
\end{proof}

\noindent
Lemme \ref{morphisms-lemma-syntomic-relative-dimension} dit que la définition
suivante a du sens.

\begin{definition}
\label{morphisms-definition-syntomic-relative-dimension}
Soit $d \geq 0$ un nombre entier. On dit qu'un morphisme de schémas $f : X \to S$ est {\it
syntomique de dimension relative $d$} si $f$
est syntomique et la fonction $\dim_x(X_{f(x)}) = d$ pour tout $x \in X$.
\end{definition}

\noindent
Autrement dit, $f$ est syntomique et les fibres non vides sont équidimensionnelles de dimension
$d$.

\begin{lemma}
\label{morphisms-lemma-syntomic-permanence}
Soit
$$
\xymatrix{
X \ar[rr]_f \ar[rd]_p & &
Y \ar[dl]^q \\
& S
}
$$
un diagramme commutatif de morphismes de schémas. Supposons que
\begin{enumerate}
\item $f$ est surjectif et syntomique,
\item $p$ est syntomique, et
\item $q$ est localement de présentation finie\footnote{En fait,
si $f$ est surjectif, plat et localement de présentation finie et $p$
est syntomique, alors $q$ et $f$
sont syntomiques, voir Descente, Lemme \ref{descent-lemma-syntomic-permanence}.}.
\end{enumerate}
Alors $q$ est syntomique.
\end{lemma}

\begin{proof}
Par le Lemme \ref{morphisms-lemma-flat-permanence} on voit que $q$ est plat. Il suffit
donc de montrer que les fibres de $Y \to S$ sont
des intersections locales complètes, voir Lemme \ref{morphisms-lemma-syntomic-flat-fibres}. Soit $s \in S$.
Considérons le morphisme $X_s \to Y_s$. Il
s'agit d'un changement de base du morphisme $X \to Y$
et donc surjectif, et syntomique (Lemme \ref{morphisms-lemma-base-change-syntomic}). Pour la
même raison, $X_s$ est syntomique par rapport à
$\kappa(s)$. De plus, $Y_s$ est localement de type
fini sur $\kappa(s)$ (Lemme \ref{morphisms-lemma-base-change-finite-type}). On se réduit
ainsi au cas où $S$ est le spectre d'un corps.

\medskip\noindent
Supposons que $S = \Spec(k)$. Soit $y \in Y$.
Choisissez un voisinage $\Spec(A) \subset Y$ ouvert affine
de $y$. Soit $\Spec(B) \subset X$ un ouvert
affine telle que $f(\Spec(B)) \subset \Spec(A)$, contenant
un point $x \in X$ tel que $f(x) = y$. Choisissez
une surjection $k[x_1, \ldots, x_n] \to A$ avec
noyau $I$. Choisissez une surjection $A[y_1, \ldots, y_m] \to B$,
qui donne lieu à son tour à une surjection $k[x_i, y_j] \to B$
à noyau $J$. Soit $\mathfrak q \subset k[x_i, y_j]$ le premier
correspondant à $y \in \Spec(B)$ et soit $\mathfrak p \subset k[x_i]$ le premier
correspondant à $x \in \Spec(A)$.
Puisque $x$ s'envoie sur $y$, nous avons $\mathfrak p = \mathfrak q \cap k[x_i]$. Considérons
le diagramme commutatif d'anneaux locaux suivant :
$$
\xymatrix{
\mathcal{O}_{X, x} \ar@{=}[r] &
B_{\mathfrak q} &
k[x_1, \ldots, x_n, y_1, \ldots, y_m]_{\mathfrak q} \ar[l] \\
\mathcal{O}_{Y, y} \ar@{=}[r] \ar[u] & A_{\mathfrak p} \ar[u] &
k[x_1, \ldots, x_n]_{\mathfrak p} \ar[l] \ar[u]
}
$$
Nous affirmons que les hypothèses
d'algèbre du lemme \ref{algebra-lemma-lci-permanence-initial} sont satisfaites. Les
conditions (1) et (2) sont triviales. La condition
(4) s'ensuit car $X \to Y$ est plat. La
condition (3) s'ensuit
car les anneaux $\mathcal{O}_{Y, y}$ et $\mathcal{O}_{X_y, x} = \mathcal{O}_{X, x}/\mathfrak m_y\mathcal{O}_{X, x}$ sont des anneaux
d'intersection complets d'après nos hypothèses
selon lesquelles $f$
et $p$ sont syntomiques, voir
Lemme \ref{morphisms-lemma-syntomic-locally-standard-syntomic}. Le résultat d'Algèbre, Lemme \ref{algebra-lemma-lci-permanence-initial}
est exactement que $\mathcal{O}_{Y, y}$ est un anneau d'intersection
complet ! Par conséquent, en Lemme \ref{morphisms-lemma-syntomic-locally-standard-syntomic}
, nous voyons à nouveau que $Y$ est syntomique sur $k$ à $y$ comme souhaité.
\end{proof}









\section{Faisceau conormal d'une immersion}
\label{morphisms-section-conormal-sheaf}

\noindent
Soit $i : Z \to X$ une immersion fermée.
Soit $\mathcal{I} \subset \mathcal{O}_X$ le faisceau quasi-cohérent d'idéaux correspondant.
Considérons la suite exacte courte
$$
0 \to \mathcal{I}^2 \to \mathcal{I} \to \mathcal{I}/\mathcal{I}^2 \to 0
$$
de faisceaux quasi-cohérents sur $X$. Le faisceau $\mathcal{I}/\mathcal{I}^2$ étant
annihilé par $\mathcal{I}$ cela correspond à un faisceau sur $Z$
par le lemme \ref{morphisms-lemma-i-star-equivalence}. Ce module $\mathcal{O}_Z$ quasi-cohérent
est appelé le {\it faisceau conormal de
$Z$ dans $X$}
et est souvent simplement noté $\mathcal{I}/\mathcal{I}^2$ par l'abus
de notation mentionné dans la section \ref{morphisms-section-closed-immersions-quasi-coherent}.

\medskip\noindent
Dans le cas où $i : Z \to X$ est une immersion (localement
fermée) nous définissons le faisceau conormal
de $i$ comme le faisceau conormal
de l'immersion fermée $i : Z \to X \setminus \partial Z$, où $\partial Z = \overline{Z} \setminus Z$.
Il est souvent noté $\mathcal{I}/\mathcal{I}^2$ où $\mathcal{I}$ est le faisceau
d'idéaux de l'immersion fermée $i : Z \to X \setminus \partial Z$.

\begin{definition}
\label{morphisms-definition-conormal-sheaf}
Soit $i : Z \to X$ une immersion. Le {\it faisceau conormal $\mathcal{C}_{Z/X}$
de $Z$ dans $X$} ou le {\it faisceau conormal
de $i$} est le faisceau conormal quasi-cohérent $\mathcal{O}_Z$-module $\mathcal{I}/\mathcal{I}^2$
décrit ci-dessus.
\end{definition}

\noindent
Dans \cite[IV Définition 16.1.2]{EGA} ce faisceau
est noté $\mathcal{N}_{Z/X}$. Nous ne suivrons pas cette convention puisque nous souhaiterions
réserver la notation $\mathcal{N}_{Z/X}$ au {\it
faisceau normal de l'immersion}. Il est défini comme
$$
\mathcal{N}_{Z/X} =
\SheafHom_{\mathcal{O}_Z}(\mathcal{C}_{Z/X}, \mathcal{O}_Z) =
\SheafHom_{\mathcal{O}_Z}(\mathcal{I}/\mathcal{I}^2, \mathcal{O}_Z)
$$
à condition que le faisceau conormal soit de présentation finie (sinon le faisceau
normal pourrait même ne pas être quasi-cohérent). Nous reviendrons plus tard sur
le faisceau normal (insérer la référence future ici).

\begin{lemma}
\label{morphisms-lemma-affine-conormal}
Soit $i : Z \to X$ une immersion. Le faisceau conormal de $i$
a les propriétés suivantes :
\begin{enumerate}
\item Soit $U \subset X$ tout sous-schéma ouvert tel
que $i$ se factorise en $Z \xrightarrow{i'} U \to X$ où $i'$
est une immersion fermée. Soit $\mathcal{I} = \Ker((i')^\sharp) \subset \mathcal{O}_U$.
Alors
$$
\mathcal{C}_{Z/X} = (i')^*\mathcal{I}\quad\text{et}\quad
i'_*\mathcal{C}_{Z/X} = \mathcal{I}/\mathcal{I}^2
$$
\item
Pour tout $\Spec(R) = U \subset X$ affine
ouvert tel que $Z \cap U = \Spec(R/I)$ il
existe un isomorphisme
canonique $\Gamma(Z \cap U, \mathcal{C}_{Z/X}) = I/I^2$.
\end{enumerate}
\end{lemma}

\begin{proof}
Généralement clair d'après les définitions. Notez que étant donné un anneau $R$
et un $I$ idéal de $R$ nous avons $I/I^2 = I \otimes_R R/I$. Détails omis.
\end{proof}

\begin{lemma}
\label{morphisms-lemma-conormal-functorial}
Soit
$$
\xymatrix{
Z \ar[r]_i \ar[d]_f & X \ar[d]^g \\
Z' \ar[r]^{i'} & X'
}
$$
un diagramme commutatif dans la catégorie des schémas. Supposons
des immersions $i$, $i'$. Il existe une application
canonique de $\mathcal{O}_Z$-modules
$$
f^*\mathcal{C}_{Z'/X'}
\longrightarrow
\mathcal{C}_{Z/X}
$$
caractérisée par la propriété suivante : Pour toute
paire d'ouverts affines $(\Spec(R) = U \subset X, \Spec(R') = U' \subset X')$ à
$g(U) \subset U'$ telles que
$Z \cap U = \Spec(R/I)$ et $Z' \cap U' = \Spec(R'/I')$ l'application
induite
$$
\Gamma(Z' \cap U', \mathcal{C}_{Z'/X'}) = I'/I'^2
\longrightarrow
I/I^2 = \Gamma(Z \cap U, \mathcal{C}_{Z/X})
$$
est celle induite par l'homomorphisme d'anneaux $g^\sharp : R' \to R$
qui possède la propriété $g^\sharp(I') \subset I$.
\end{lemma}

\begin{proof}
Soient $\partial Z' = \overline{Z'} \setminus Z'$ et $\partial Z = \overline{Z} \setminus Z$. Ce sont
des sous-ensembles fermés de $X'$ et de $X$. En remplaçant
$X'$ par $X' \setminus \partial Z'$ et $X$ par
$X \setminus \big(g^{-1}(\partial Z') \cup \partial Z\big)$ on voit que l'on peut supposer que $i$
et $i'$ sont des immersions fermées.

\medskip\noindent
Le fait que $g \circ i$ se factorise
par $i'$ implique que les
$g^*\mathcal{I}'$ applications en
$\mathcal{I}$
sous l'application canonique $g^*\mathcal{I}' \to \mathcal{O}_X$,
voir Schémas, Lemmes
\ref{schemes-lemma-characterize-closed-subspace} et \ref{schemes-lemma-restrict-map-to-closed}. On
obtient donc une application induite de
faisceaux quasi-cohérents
$g^*(\mathcal{I}'/(\mathcal{I}')^2) \to \mathcal{I}/\mathcal{I}^2$. Retirer de $i$ donne
$i^*g^*(\mathcal{I}'/(\mathcal{I}')^2) \to i^*(\mathcal{I}/\mathcal{I}^2)$. Notez que $i^*(\mathcal{I}/\mathcal{I}^2) = \mathcal{C}_{Z/X}$.
D'autre part,
$i^*g^*(\mathcal{I}'/(\mathcal{I}')^2) =
f^*(i')^*(\mathcal{I}'/(\mathcal{I}')^2) = f^*\mathcal{C}_{Z'/X'}$. Cela donne l'application
souhaitée.

\medskip\noindent
Vérifier que l'application est décrite localement comme l'application
donnée $I'/(I')^2 \to I/I^2$ consiste à dérouler les définitions et
est omise. Une autre observation est que étant donné tout
$x \in i(Z)$ il existe des voisinages ouverts affines $U$,
$U'$ avec $f(U) \subset U'$ et $Z \cap U$ ainsi que $U' \cap Z'$ fermés
tels que $x \in U$. Démonstration omise. L'exigence du lemme
caractérise donc bien l'application (et aurait pu être utilisée pour
la définir).
\end{proof}

\begin{lemma}
\label{morphisms-lemma-conormal-functorial-flat}
Soit
$$
\xymatrix{
Z \ar[r]_i \ar[d]_f & X \ar[d]^g \\
Z' \ar[r]^{i'} & X'
}
$$
un diagramme produit fibré dans la catégorie des schémas
à immersions $i$, $i'$.
Alors l'application canonique $f^*\mathcal{C}_{Z'/X'} \to \mathcal{C}_{Z/X}$ du
lemme \ref{morphisms-lemma-conormal-functorial} est surjectif. Si
$g$ est plat, alors c'est un isomorphisme.
\end{lemma}

\begin{proof}
Soit $R' \to R$ un homomorphisme d'anneaux, et $I' \subset R'$
un idéal. Posons $I = I'R$. Alors $I'/(I')^2 \otimes_{R'} R \to I/I^2$ est
surjectif. Si $R' \to R$ est plat, alors $I = I' \otimes_{R'} R$ et
$I^2 = (I')^2 \otimes_{R'} R$ et on voit que l'application est un
isomorphisme.
\end{proof}

\begin{lemma}
\label{morphisms-lemma-transitivity-conormal}
Soit $Z \to Y \to X$ des immersions de schémas. Il existe alors une suite exacte
canonique
$$
i^*\mathcal{C}_{Y/X} \to
\mathcal{C}_{Z/X} \to
\mathcal{C}_{Z/Y} \to 0
$$
dont les applications proviennent
du lemme \ref{morphisms-lemma-conormal-functorial} et
$i : Z \to Y$ est le premier morphisme.
\end{lemma}

\begin{proof}
Via
Lemme \ref{morphisms-lemma-conormal-functorial} cela se traduit
par le fait algébrique suivant. Supposons que $C \to B \to A$
soient des homomorphismes surjectifs d'anneaux. Soit $I = \Ker(B \to A)$,
$J = \Ker(C \to A)$ et $K = \Ker(C \to B)$. Il existe
alors une suite exacte
$$
K/K^2 \otimes_B A \to J/J^2 \to I/I^2 \to 0.
$$
Cela découle immédiatement de l'observation que $I = J/K$.
\end{proof}










\section{Faisceau des différentielles d'un morphisme}
\label{morphisms-section-sheaf-differentials}

\noindent
Nous suggérons au lecteur de consulter la section correspondante
dans le chapitre sur l'algèbre
commutative (Algèbre, section \ref{algebra-section-differentials})
et la section correspondante dans le chapitre sur les faisceaux
de modules (Modules, section \ref{modules-section-differentials}).

\begin{definition}
\label{morphisms-definition-sheaf-differentials}
Soit $f : X \to S$ un morphisme de schémas. Le {\it
faisceau des différentielles $\Omega_{X/S}$ de $X$ sur $S$}
est le faisceau des différentielles de $f$ vu comme un morphisme d'espaces
annelés (Modules, Définition \ref{modules-definition-differentials}) équipé de son {\it
dérivation $S$ universelle}
$$
\text{d}_{X/S} : \mathcal{O}_X \longrightarrow \Omega_{X/S}.
$$
\end{definition}

\noindent
Il s'avère que $\Omega_{X/S}$ est un $\mathcal{O}_X$-module quasi-cohérent par exemple
car il est isomorphe au faisceau conormal de la diagonale
du morphisme $\Delta : X \to X \times_S X$ (Lemme \ref{morphisms-lemma-differentials-diagonal}). Nous avons
défini le module des différentielles de $X$
sur $S$ en utilisant une propriété universelle, à savoir
comme réceptacle de la dérivation universelle. Si vous avez une autre
construction du faisceau de différentielles relatives qui satisfait cette propriété
universelle alors, par le lemme de Yoneda, elle sera canoniquement
isomorphe à celle définie ci-dessus. Pour plus de commodité, nous reformulons
ici la propriété universelle.

\begin{lemma}
\label{morphisms-lemma-universal-derivation-universal}
Soit $f : X \to S$ un morphisme de schémas. L'application
$$
\Hom_{\mathcal{O}_X}(\Omega_{X/S}, \mathcal{F})
\longrightarrow
\text{Der}_S(\mathcal{O}_X, \mathcal{F}),\quad
\alpha \longmapsto \alpha \circ \text{d}_{X/S}
$$
est un isomorphisme de foncteurs $\textit{Mod}(\mathcal{O}_X) \to \textit{Ensembles}$.
\end{lemma}

\begin{proof}
Il s'agit simplement d'une reformulation de la définition.
\end{proof}

\begin{lemma}
\label{morphisms-lemma-differentials-restrict-open}
Soit $f : X \to S$ un morphisme de schémas.
Soit $U \subset X$, $V \subset S$ être des sous-schémas
ouverts tels que $f(U) \subset V$. Il existe
alors un unique isomorphisme $\Omega_{X/S}|_U = \Omega_{U/V}$ des $\mathcal{O}_U$-modules
tel que $\text{d}_{X/S}|_U = \text{d}_{U/V}$.
\end{lemma}

\begin{proof}
Ceci est un cas particulier
de Modules, Lemme \ref{modules-lemma-localize-differentials}
si l'on utilise l'identification
canonique $f^{-1}\mathcal{O}_S|_U = (f|_U)^{-1}\mathcal{O}_V$.
\end{proof}

\noindent
Désormais nous utiliserons ces identifications canoniques et écrirons
simplement $\Omega_{U/S}$ ou $\Omega_{U/V}$ pour la restriction de $\Omega_{X/S}$
à $U$.

\begin{lemma}
\label{morphisms-lemma-affine-case-derivation}
Soit $R \to A$ un homomorphisme d'anneaux. Soit
$\mathcal{F}$ un faisceau de $\mathcal{O}_X$-modules
sur $X = \Spec(A)$. Posons $S = \Spec(R)$.
La règle qui associe à une $S$-dérivation sur $\mathcal{F}$ son
action sur les sections globales définit une bijection entre
l'ensemble des $S$-dérivations de $\mathcal{F}$ et l'ensemble
des $R$-dérivations sur $M = \Gamma(X, \mathcal{F})$.
\end{lemma}

\begin{proof}
Soit $D : A \to M$ une dérivation $R$. Nous devons montrer qu'il existe
une unique dérivation $S$ sur $\mathcal{F}$ qui donne naissance
à $D$ sur les sections globales. Soit $U = D(f) \subset X$ un ouvert affine standard.
Tout élément de $\Gamma(U, \mathcal{O}_X)$ est de la forme $a/f^n$ pour
certains $a \in A$ et $n \geq 0$. Par la règle de Leibniz, nous
avons
$$
D(a)|_U = a/f^n D(f^n)|_U + f^n D(a/f^n)
$$
dans $\Gamma(U, \mathcal{F})$. Puisque $f$ agit de manière
inversible sur $\Gamma(U, \mathcal{F})$, cela détermine complètement
la valeur de $D(a/f^n) \in \Gamma(U, \mathcal{F})$. Cela prouve l'unicité.
L'existence s'ensuit en définissant simplement
$$
D(a/f^n) := (1/f^n) D(a)|_U - a/f^{2n} D(f^n)|_U
$$
et en prouvant qu'il possède toutes les propriétés souhaitées (sur la base des
ouvertures standard de $X$). Détails omis.
\end{proof}

\begin{lemma}
\label{morphisms-lemma-differentials-affine}
Soit $f : X \to S$ un morphisme de schémas. Pour toute paire d'affines
ouvertes $\Spec(A) = U \subset X$, $\Spec(R) = V \subset S$ avec $f(U) \subset V$ il existe
un isomorphisme unique
$$
\Gamma(U, \Omega_{X/S}) = \Omega_{A/R}.
$$
compatible avec $\text{d}_{X/S}$ et $\text{d} : A \to \Omega_{A/R}$.
\end{lemma}

\begin{proof}
Par le Lemme \ref{morphisms-lemma-differentials-restrict-open} on peut remplacer $X$ et $S$
par $U$ et $V$. Ainsi nous pouvons supposer
$X = \Spec(A)$ et $S = \Spec(R)$ et nous devons montrer le
lemme avec $U = X$ et $V = S$. Considérons le module
$A$ $M = \Gamma(X, \Omega_{X/S})$ avec la dérivation $R$ $\text{d}_{X/S} : A \to M$.
Soit $N$ un autre $A$-module et désignons $\widetilde{N}$
le module $\mathcal{O}_X$ quasi-cohérent
associé à $N$, voir Schémas, section \ref{schemes-section-quasi-coherent-affine}.
En précomposant par $\text{d}_{X/S} : A \to M$ on obtient une flèche
$$
\alpha : \Hom_A(M, N) \longrightarrow \text{Der}_R(A, N)
$$
En utilisant les lemmes \ref{morphisms-lemma-universal-derivation-universal} et \ref{morphisms-lemma-affine-case-derivation}
on obtient les identifications
$$
\Hom_{\mathcal{O}_X}(\Omega_{X/S}, \widetilde{N}) =
\text{Der}_S(\mathcal{O}_X, \widetilde{N}) =
\text{Der}_R(A, N)
$$
La prise de sections globales détermine
une flèche $\Hom_{\mathcal{O}_X}(\Omega_{X/S}, \widetilde{N}) \to \Hom_R(M, N)$. En combinant cette flèche et les
identifications ci-dessus nous obtenons une flèche
$$
\beta : \text{Der}_R(A, N) \longrightarrow \Hom_R(M, N)
$$
En vérifiant ce qui se passe sur les sections globales, nous
constatons que $\alpha$ et $\beta$ sont l'un
l'autre inverses. Nous voyons donc
que $\text{d}_{X/S} : A \to M$ satisfait la même propriété universelle que
$\text{d} : A \to \Omega_{A/R}$, voir Algèbre, Lemme \ref{algebra-lemma-universal-omega}. Ainsi le lemme de Yoneda (Catégories,
Lemme \ref{categories-lemma-yoneda}) implique qu'il existe
un unique isomorphisme des $A$-modules
$M \cong \Omega_{A/R}$ compatible avec les dérivations.
\end{proof}

\begin{remark}
\label{morphisms-remark-differentials-glue}
Le lemme ci-dessus donne une deuxième façon de construire le module des
différentielles. En effet, soit $f : X \to S$ un morphisme de schémas. Considérons
la famille de tous les ouverts affines $U \subset X$ dont l'image est contenue
dans un ouvert affine de $S$. Ceux-ci forment une base
de la topologie sur $X$. Il suffit donc de définir $\Gamma(U, \Omega_{X/S})$
pour un tel $U$. Nous définissons simplement $\Gamma(U, \Omega_{X/S}) = \Omega_{A/R}$ si $A$,
$R$ sont comme dans le lemme \ref{morphisms-lemma-differentials-affine} ci-dessus. Cela fonctionne,
mais il faut des préliminaires un peu plus algébriques pour
construire les mappages de restriction et vérifier la condition
du faisceau avec cet ansatz.
\end{remark}

\noindent
Le lemme suivant donne encore une autre façon de définir le faisceau des différentielles
et il montre en particulier que $\Omega_{X/S}$ est quasi-cohérent
si $X$ et $S$ sont des schémas.

\begin{lemma}
\label{morphisms-lemma-differentials-diagonal}
Soit $f : X \to S$ un morphisme de schémas. Il existe
un isomorphisme canonique entre $\Omega_{X/S}$ et le
faisceau conormal du morphisme diagonal $\Delta_{X/S} : X \longrightarrow X \times_S X$.
\end{lemma}

\begin{proof}
Nous établissons d'abord l'existence de quelques ``globaux'' faisceaux
et applications globales de faisceaux, et plus loin nous décrivons
les constructions sur certaines ouverts affines.

\medskip\noindent
Rappelons que $\Delta = \Delta_{X/S} : X \to X \times_S X$ est une immersion,
voir Schémas, Lemme \ref{schemes-lemma-diagonal-immersion}. Soit $\mathcal{J}$ le faisceau d'idéaux
de l'immersion qui habite un sous-schéma
ouvert $W$ de $X \times_S X$ tel que $\Delta(X) \subset W$
soit fermé. Prenons celui que l'on retrouve dans la
démonstration d'après
Schémas, Lemme \ref{schemes-lemma-diagonal-immersion}. Notez que le faisceau
d'anneaux $\mathcal{O}_W/\mathcal{J}^2$ est supporté sur $\Delta(X)$.
De plus il se situe dans une suite exacte
courte de faisceaux
$$
0 \to \mathcal{J}/\mathcal{J}^2
\to \mathcal{O}_W/\mathcal{J}^2
\to \Delta_*\mathcal{O}_X
\to 0.
$$
En utilisant $\Delta^{-1}$ on peut penser à cela comme
une surjection de faisceaux de $f^{-1}\mathcal{O}_S$-algèbres
avec noyau le faisceau conormal de $\Delta$ (voir Définition
\ref{morphisms-definition-conormal-sheaf} et Lemme \ref{morphisms-lemma-affine-conormal}).
$$
0 \to \mathcal{C}_{X/X \times_S X}
\to \Delta^{-1}(\mathcal{O}_W/\mathcal{J}^2)
\to \mathcal{O}_X
\to 0
$$
Cela nous place dans la
situation des Modules, Lemme \ref{modules-lemma-double-structure-gives-derivation}. Les morphismes
de projection $p_i : X \times_S X \to X$, $i = 1, 2$ induisent
des applications
de faisceaux d'anneaux $(p_i)^\sharp : (p_i)^{-1}\mathcal{O}_X \to \mathcal{O}_{X \times_S X}$. Nous pouvons
nous limiter à $W$ et quotient par
$\mathcal{J}^2$ pour obtenir $(p_i)^{-1}\mathcal{O}_X \to \mathcal{O}_W/\mathcal{J}^2$. Puisque
$\Delta^{-1}p_i^{-1}\mathcal{O}_X = \mathcal{O}_X$, nous obtenons des
applications
$$
s_i : \mathcal{O}_X \to \Delta^{-1}(\mathcal{O}_W/\mathcal{J}^2).
$$
$s_1$ et $s_2$ sont des
sections de l'application $\Delta^{-1}(\mathcal{O}_W/\mathcal{J}^2) \to \mathcal{O}_X$, comme dans
Modules, Lemme \ref{modules-lemma-double-structure-gives-derivation}. Nous obtenons ainsi une dérivation
$S$ $\text{d} = s_2 - s_1 : \mathcal{O}_X \to \mathcal{C}_{X/X \times_S X}$.
Par la propriété universelle du module des
différentielles nous trouvons une application unique
$\mathcal{O}_X$-linéaire
$$
\Omega_{X/S} \longrightarrow \mathcal{C}_{X/X \times_S X},\quad
f\text{d}g \longmapsto fs_2(g) - fs_1(g)
$$
Pour voir que l'application est un isomorphisme, calculons cela sur
des ouverts affines appropriées. On peut couvrir $X$
par des ouverts affines $\Spec(A) = U \subset X$ dont l'image est contenue
dans un ouvert affine $\Spec(R) = V \subset S$. D'après la démonstration d'après
Schémas, Lemme \ref{schemes-lemma-diagonal-immersion} $U \times_V U \subset X \times_S X$ est un ouvert
affine contenue dans l'ouvert $W$
mentionné ci-dessus. Également
$U \times_V U = \Spec(A \otimes_R A)$. Le faisceau $\mathcal{J}$ correspond
à l'idéal $J = \Ker(A \otimes_R A \to A)$. La
suite exacte courte à la suite exacte courte
des $A \otimes_R A$-modules
$$
0 \to J/J^2 \to (A \otimes_R A)/J^2 \to A \to 0
$$
Les sections $s_i$ correspondent aux homomorphismes d'anneaux
$$
A \longrightarrow (A \otimes_R A)/J^2,\quad
s_1 : a \mapsto a \otimes 1,\quad
s_2 : a \mapsto 1 \otimes a.
$$
Par le Lemme \ref{morphisms-lemma-affine-conormal} nous avons
$\Gamma(U, \mathcal{C}_{X/X \times_S X}) = J/J^2$ et par le lemme
\ref{morphisms-lemma-differentials-affine} nous avons $\Gamma(U, \Omega_{X/S}) = \Omega_{A/R}$.
L'application ci-dessus
est l'application $a \text{d}b \mapsto a \otimes b - ab \otimes 1$, dont
Algèbre, Lemme \ref{algebra-lemma-differentials-diagonal}, montre qu'elle est
un isomorphisme.
\end{proof}

\begin{lemma}
\label{morphisms-lemma-functoriality-differentials}
Soit
$$
\xymatrix{
X' \ar[d] \ar[r]_f & X \ar[d] \\
S' \ar[r] & S
}
$$
être un diagramme commutatif de schémas. L'application
canonique $\mathcal{O}_X \to f_*\mathcal{O}_{X'}$ composée avec l'application
$f_*\text{d}_{X'/S'} : f_*\mathcal{O}_{X'} \to f_*\Omega_{X'/S'}$ est une $S$-dérivation. On obtient
donc une application canonique des $\mathcal{O}_X$-modules $\Omega_{X/S} \to f_*\Omega_{X'/S'}$,
et par adjointité de $f_*$
et $f^*$ un homomorphisme
canonique des $\mathcal{O}_{X'}$-modules
$$
c_f : f^*\Omega_{X/S} \longrightarrow \Omega_{X'/S'}.
$$
Il est uniquement caractérisé par la propriété que
les applications $f^*\text{d}_{X/S}(h)$ à $\text{d}_{X'/S'}(f^* h)$ pour toute
section locale $h$ de $\mathcal{O}_X$.
\end{lemma}

\begin{proof}
Ceci est un cas particulier
de Modules, Lemme \ref{modules-lemma-functoriality-differentials-ringed-spaces}. Dans le cas des schémas, on peut
aussi utiliser la fonctorialité des faisceaux conormaux
(voir le lemme \ref{morphisms-lemma-conormal-functorial}) et le lemme \ref{morphisms-lemma-differentials-diagonal} pour définir
$c_f$. Ou nous pouvons utiliser la caractérisation
de la dernière ligne du lemme pour coller des
applications définies sur des
ouverts affines (voir Algèbre, Equation (\ref{algebra-equation-functorial-omega})).
\end{proof}

\begin{lemma}
\label{morphisms-lemma-triangle-differentials}
Soient $f : X \to Y$, $g : Y \to S$ des morphismes de schémas. Il
existe alors une suite exacte canonique
$$
f^*\Omega_{Y/S} \to \Omega_{X/S} \to \Omega_{X/Y} \to 0
$$
où les applications proviennent
d'applications du lemme \ref{morphisms-lemma-functoriality-differentials}.
\end{lemma}

\begin{proof}
Il s'agit de la version simplifiée
de l'algèbre, Lemme \ref{algebra-lemma-exact-sequence-differentials}. Alternativement,
il existe une version générale pour les morphismes
des espaces annelés, voir Modules, Lemme \ref{modules-lemma-triangle-differentials}.
\end{proof}

\begin{lemma}
\label{morphisms-lemma-base-change-differentials}
Soit $X \to S$ un morphisme de schémas. Soit
$g : S' \to S$ un morphisme de schémas. Soit $X' = X_{S'}$
le changement de base de $X$.
Notons $g' : X' \to X$ la projection. Alors
l'application
$$
(g')^*\Omega_{X/S} \to \Omega_{X'/S'}
$$
du lemme \ref{morphisms-lemma-functoriality-differentials} est un isomorphisme.
\end{lemma}

\begin{proof}
Il s'agit de la version
simplifiée de l'algèbre, Lemme \ref{algebra-lemma-differentials-base-change}.
\end{proof}

\begin{lemma}
\label{morphisms-lemma-differential-product}
Soient $f : X \to S$ et $g : Y \to S$ des morphismes de schémas
de même cible. Soit $p : X \times_S Y \to X$ et $q : X \times_S Y \to Y$ les morphismes
de projection. Les
applications du lemme \ref{morphisms-lemma-functoriality-differentials}
$$
p^*\Omega_{X/S} \oplus q^*\Omega_{Y/S}
\longrightarrow
\Omega_{X \times_S Y/S}
$$
donnent un isomorphisme.
\end{lemma}

\begin{proof}
Par le Lemme \ref{morphisms-lemma-base-change-differentials} la composition
$p^*\Omega_{X/S} \to \Omega_{X \times_S Y/S} \to \Omega_{X \times_S Y/Y}$ est un isomorphisme, et de même
pour $q$. De plus, le conoyau
de $p^*\Omega_{X/S} \to \Omega_{X \times_S Y/S}$ est $\Omega_{X \times_S Y/X}$
par le lemme
\ref{morphisms-lemma-triangle-differentials}. Le résultat suit.
\end{proof}

\begin{lemma}
\label{morphisms-lemma-finite-type-differentials}
Soit $f : X \to S$ un morphisme de schémas. Si $f$
est localement de type fini, alors $\Omega_{X/S}$ est
un $\mathcal{O}_X$-module de type fini.
\end{lemma}

\begin{proof}
Immédiat
de l'algèbre, Lemme \ref{algebra-lemma-differentials-finitely-generated},
Lemme \ref{morphisms-lemma-differentials-affine},
Lemme \ref{morphisms-lemma-locally-finite-type-characterize} et Propriétés,
Lemme \ref{properties-lemma-finite-type-module}.
\end{proof}

\begin{lemma}
\label{morphisms-lemma-finite-presentation-differentials}
Soit $f : X \to S$ un morphisme de schémas. Si
$f$ est localement de présentation finie, alors $\Omega_{X/S}$
est un $\mathcal{O}_X$-module de présentation finie.
\end{lemma}

\begin{proof}
Immédiat
de l'algèbre, Lemme \ref{algebra-lemma-differentials-finitely-presented},
Lemme \ref{morphisms-lemma-differentials-affine},
Lemme \ref{morphisms-lemma-locally-finite-presentation-characterize} et Propriétés,
Lemme \ref{properties-lemma-finite-presentation-module}.
\end{proof}

\begin{lemma}
\label{morphisms-lemma-immersion-differentials}
Si $X \to S$ est une immersion, ou plus généralement un monomorphisme, alors
$\Omega_{X/S}$ est nul.
\end{lemma}

\begin{proof}
Ceci est vrai car $\Delta_{X/S}$ est un isomorphisme
dans ce cas et a donc un faisceau conormal trivial.
D'où $\Omega_{X/S} = 0$ par le lemme \ref{morphisms-lemma-differentials-diagonal}. La version
algébrique est Algèbre, Lemme \ref{algebra-lemma-trivial-differential-surjective}.
\end{proof}

\begin{lemma}
\label{morphisms-lemma-differentials-relative-immersion}
Soit $i : Z \to X$ une immersion de schémas sur $S$. Il
existe une suite exacte canonique
$$
\mathcal{C}_{Z/X} \to i^*\Omega_{X/S} \to \Omega_{Z/S} \to 0
$$
où la première flèche est induite par $\text{d}_{X/S}$
et la deuxième flèche vient du lemme \ref{morphisms-lemma-functoriality-differentials}.
\end{lemma}

\begin{proof}
Il s'agit de la version simplifiée
de l'algèbre, Lemme \ref{algebra-lemma-differential-seq}. Cependant, nous
devons nous assurer que nous pouvons définir la
première flèche de manière globale. Nous expliquons donc ici
la signification de `` induit par $\text{d}_{X/S}$''.
À savoir, on peut supposer
que $i$ est une immersion fermée par rétrécissement
$X$. Soit $\mathcal{I} \subset \mathcal{O}_X$ le
faisceau d'idéaux correspondant à $Z \subset X$. Puis
$\text{d}_{X/S} : \mathcal{I} \to \Omega_{X/S}$ applique le sous-faisceau $\mathcal{I}^2 \subset \mathcal{I}$
à $\mathcal{I}\Omega_{X/S}$. Il induit donc une application $\mathcal{I}/\mathcal{I}^2 \to \Omega_{X/S}/\mathcal{I}\Omega_{X/S}$
qui est $\mathcal{O}_X/\mathcal{I}$-linéaire.
Par le Lemme \ref{morphisms-lemma-i-star-equivalence} cela correspond à une application
$\mathcal{C}_{Z/X} \to i^*\Omega_{X/S}$ souhaitée.
\end{proof}

\begin{lemma}
\label{morphisms-lemma-differentials-relative-immersion-section}
Soit $i : Z \to X$ une immersion de schémas sur $S$, et supposons
que $i$ (localement) a un inverse gauche. Alors la séquence
canonique
$$
0 \to \mathcal{C}_{Z/X} \to i^*\Omega_{X/S} \to \Omega_{Z/S} \to 0
$$
du
lemme \ref{morphisms-lemma-differentials-relative-immersion} est (localement) divisée
exactement. En particulier, si $s : S \to X$ est une section
du morphisme de structure $X \to S$ alors
l'application $\mathcal{C}_{S/X} \to s^*\Omega_{X/S}$ induite par $\text{d}_{X/S}$
est un isomorphisme.
\end{lemma}

\begin{proof}
Il résulte de l'algèbre, Lemme \ref{algebra-lemma-differential-seq-split} que le résultat est valable
lorsque $X, Z, S$ sont affines. En général, nous concluons
que la séquence est exacte localement divisée. Si $g : X \to Z$
est un inverse gauche de $i$, alors
$i^*c_g$ est un inverse droit
de l'application $i^*\Omega_{X/S} \to \Omega_{Z/S}$ par Modules, Lemme \ref{modules-lemma-check-functoriality-differentials}.
Cela implique que la séquence est
divisée par homologie, Lemme \ref{homology-lemma-ses-split}.
Enfin, si $s$ est une section, alors c'est une immersion
$s : Z = S \to X$ sur $S$ (voir Schémas, Lemme \ref{schemes-lemma-section-immersion})
et dans ce cas $\Omega_{Z/S} = 0$.
\end{proof}

\begin{remark}
\label{morphisms-remark-differentials-diagonal}
Soit $X \to S$ un morphisme de schémas.
D'après le Lemme \ref{morphisms-lemma-differential-product}
on a
$$
\Omega_{X \times_S X/S} =
\text{pr}_1^*\Omega_{X/S} \oplus \text{pr}_2^*\Omega_{X/S}
$$
Par contre, le morphisme diagonal $\Delta : X \to X \times_S X$ est une immersion,
qui a localement un inverse gauche.
Ainsi par le lemme \ref{morphisms-lemma-differentials-relative-immersion-section} on obtient une
suite canonique exacte courte
$$
0 \to \mathcal{C}_{X/X \times_S X} \to \Omega_{X/S} \oplus \Omega_{X/S}
\to \Omega_{X/S} \to 0
$$
Notons que la flèche droite est $(1, 1)$ qui
est bien une surjection fractionnée. Par
contre, par le lemme \ref{morphisms-lemma-differentials-diagonal} nous avons une
identification $\Omega_{X/S} = \mathcal{C}_{X/X \times_S X}$. Comme nous avons choisi $\text{d}_{X/S}(f) = s_2(f) - s_1(f)$ dans cette
identification
il apparaît que la flèche de gauche est l'application
$(-1, 1)$\footnote{Plus précisément,
la section locale $\text{d}_{X/S}(f) = 1 \otimes f - f \otimes 1$ du faisceau d'idéaux de $\Delta$
s'envoie par $\text{d}_{X \times_S X/X}$ sur la section locale
$1 \otimes 1 \otimes 1 \otimes f - 1 \otimes f \otimes 1 \otimes 1
-1 \otimes 1 \otimes f \otimes 1 + f \otimes 1 \otimes 1 \otimes 1 =
\text{pr}_2^*\text{d}_{X/S}(f) - \text{pr}_1^*\text{d}_{X/S}(f)$.}.
\end{remark}

\begin{lemma}
\label{morphisms-lemma-two-immersions}
Soit
$$
\xymatrix{
Z \ar[r]_i \ar[rd]_j & X \ar[d] \\
& Y
}
$$
un diagramme commutatif de schémas où $i$ et $j$ sont des immersions.
Il existe alors une suite exacte canonique
$$
\mathcal{C}_{Z/Y} \to
\mathcal{C}_{Z/X} \to
i^*\Omega_{X/Y} \to 0
$$
où la première flèche
vient du lemme \ref{morphisms-lemma-conormal-functorial}
et la
seconde du lemme \ref{morphisms-lemma-differentials-relative-immersion}.
\end{lemma}

\begin{proof}
La version algébrique de
ceci est Algèbre, Lemme \ref{algebra-lemma-application-NL}.
\end{proof}










\section{Opérateurs différentiels d'ordre fini}
\label{morphisms-section-differential-operators}

\noindent
Nous suggérons au lecteur de consulter la section correspondante
dans le chapitre sur
l'algèbre commutative (Algèbre, section \ref{algebra-section-differential-operators})
et la section correspondante dans le chapitre sur les
faisceaux de modules (Modules, section \ref{modules-section-differential-operators}).

\begin{lemma}
\label{morphisms-lemma-affine-case-differential-operators}
Soit $R \to A$ un homomorphisme d'anneaux. Notons $f : X \to S$ le morphisme
correspondant des schémas affines. Soit $\mathcal{F}$ et $\mathcal{G}$ être
des $\mathcal{O}_X$-modules. Si $\mathcal{F}$ est quasi-cohérent alors
l'application
$$
\text{Diff}^k_{X/S}(\mathcal{F}, \mathcal{G}) \to
\text{Diff}^k_{A/R}(\Gamma(X, \mathcal{F}), \Gamma(X, \mathcal{G}))
$$
envoyant un opérateur différentiel à son action sur les sections globales
est bijective.
\end{lemma}

\begin{proof}
Écrivez $\mathcal{F} = \widetilde{M}$ pour certains $A$-module $M$.
Posons $N = \Gamma(X, \mathcal{G})$. Soit $D : M \to N$ un opérateur différentiel
d'ordre $k$. Nous devons montrer
qu'il existe un unique opérateur différentiel $\mathcal{F} \to \mathcal{G}$ d'ordre $k$
qui donne naissance à $D$ sur les sections globales.
Soit $U = D(f) \subset X$ un ouvert affine standard. Alors $\mathcal{F}(U) = M_f$
est la localisation. Par Algèbre, Lemme \ref{algebra-lemma-invert-system-differential-operators} l'opérateur différentiel
$D$ s'étend à un unique opérateur différentiel
$$
D_f : \mathcal{F}(U) = \widetilde{M}(U) = M_f \to N_f = \widetilde{N}(U)
$$
L'unicité montre que ces applications $D_f$ collent pour donner
une application de faisceaux $\widetilde{M} \to \widetilde{N}$ sur la base de toutes
les ouvertures standards de $X$.
Nous obtenons donc une application unique de faisceaux $\widetilde{D} : \widetilde{M} \to \widetilde{N}$
en accord avec ces applications par le matériel de Faisceaux,
section \ref{sheaves-section-bases}. Puisque $\widetilde{D}$ est donné par des
opérateurs différentiels d'ordre $k$ sur le standard ouvre,
on constate que $\widetilde{D}$ est un
opérateur différentiel d'ordre $k$ (petit détail omis). Enfin,
on peut post-composer avec
l'application canonique $\mathcal{O}_X$-module $c : \widetilde{N} \to \mathcal{G}$
(Schémas, Lemme \ref{schemes-lemma-compare-constructions}) pour obtenir $c \circ \widetilde{D} : \mathcal{F} \to \mathcal{G}$
qui est un opérateur différentiel d'ordre $k$
par Modules, Lemme \ref{modules-lemma-composition-differential-operators}. Cela prouve l'existence.
Nous omettons la démonstration de l'unicité.
\end{proof}

\begin{lemma}
\label{morphisms-lemma-base-change-differential-operators}
Soient $a : X \to S$ et $b : Y \to S$ des morphismes de schémas.
Soit $\mathcal{F}$ et $\mathcal{F}'$ des modules $\mathcal{O}_X$ quasi-cohérents. Soit
$D : \mathcal{F} \to \mathcal{F}'$ un opérateur différentiel d'ordre $k$
sur $X/S$. Soit $\mathcal{G}$ un module $\mathcal{O}_Y$ quasi-cohérent.
Il existe alors un unique opérateur
différentiel
$$
D' :
\text{pr}_1^*\mathcal{F} \otimes_{\mathcal{O}_{X \times_S Y}}
\text{pr}_2^*\mathcal{G}
\longrightarrow
\text{pr}_1^*\mathcal{F}' \otimes_{\mathcal{O}_{X \times_S Y}}
\text{pr}_2^*\mathcal{G}
$$
d'ordre $k$ sur $X \times_S Y / Y$ tel que
$
D'(s \otimes t) = D(s) \otimes t
$
pour les sections locales $s$ de $\mathcal{F}$ et $t$ de $\mathcal{G}$.
\end{lemma}

\begin{proof}
Dans le cas où $X$, $Y$ et $S$
sont affines, cela découle, via Lemme \ref{morphisms-lemma-affine-case-differential-operators},
du résultat algébrique correspondant,
voir Algèbre, Lemme \ref{algebra-lemma-base-change-differential-operators}. En général, on
utilise des revêtements par affines
(par exemple comme dans Schémas,
Lemme \ref{schemes-lemma-affine-covering-fibre-product}) pour construire $D'$
globalement. Détails omis.
\end{proof}

\begin{remark}
\label{morphisms-remark-base-change-differential-operators}
Soient $a : X \to S$ et $b : Y \to S$ des morphismes de schémas.
Notons $p : X \times_S Y \to X$ et $q : X \times_S Y \to Y$ les projections.
Dans cette remarque, étant donné un $\mathcal{O}_X$-module $\mathcal{F}$
et un $\mathcal{O}_Y$-module $\mathcal{G}$ fixons
$$
\mathcal{F} \boxtimes \mathcal{G} =
p^*\mathcal{F} \otimes_{\mathcal{O}_{X \times_S Y}} q^*\mathcal{G}
$$
Désignons $\mathcal{A}_{X/S}$ la catégorie additive
dont les objets sont des $\mathcal{O}_X$-modules quasi-cohérents
et dont les morphismes sont des opérateurs différentiels
d'ordre fini sur $X/S$. De même pour $\mathcal{A}_{Y/S}$
et $\mathcal{A}_{X \times_S Y/S}$. La construction du lemme \ref{morphisms-lemma-base-change-differential-operators} détermine
un foncteur
$$
\boxtimes : 
\mathcal{A}_{X/S} \times \mathcal{A}_{Y/S} \longrightarrow
\mathcal{A}_{X \times_S Y/S},
\quad
(\mathcal{F}, \mathcal{G}) \longmapsto \mathcal{F} \boxtimes \mathcal{G}
$$
qui est bilinéaire sur les morphismes. Si $X = \Spec(A)$, $Y = \Spec(B)$,
et $S = \Spec(R)$, alors via l'identification de
faisceaux quasi-cohérents avec des modules ce foncteur
est donné par $(M, N) \mapsto M \otimes_R N$ sur les objets et envoie
le morphisme $(D, D') : (M, N) \to (M', N')$ à $D \otimes D' : M \otimes_R N \to M' \otimes_R N'$.
\end{remark}








\section{Morphismes lisses}
\label{morphisms-section-smooth}

\noindent
Soit $f : X \to Y$ une application continue d'espaces topologiques.
Considérons la condition
suivante : Pour tout $x \in X$ il existe des voisinages
ouverts $x \in U \subset X$ et $f(x) \in V \subset Y$, et un entier $d$
tel que $f(U) \subset V$ et tel que l'on obtient un diagramme
commutatif
$$
\xymatrix{
X \ar[d] &
U \ar[l] \ar[d] \ar[r]_-\pi &
V \times \mathbf{R}^d \ar[ld] \\
Y & V \ar[l]
}
$$
où $\pi$ est un homéomorphisme sur un sous-ensemble
ouvert. Les morphismes de schémas lisses
sont l'analogue de ces applications
dans la catégorie des schémas.
Voir
Lemme \ref{morphisms-lemma-smooth-locally-standard-smooth} et Lemme \ref{morphisms-lemma-smooth-etale-over-affine-space}.

\medskip\noindent Contrairement
aux attentes (peut-être) la notion d'homomorphisme
d'anneaux lisse n'est pas définie uniquement en termes
de module des différentielles. À savoir, rappelons que $R \to A$
est un {\it homomorphisme d'anneaux lisse} si $A$ est de
présentation finie sur $R$ et si le complexe cotangent naïf de $A$ sur
$R$ est quasi-isomorphe à un module projectif placé
en degré $0$, voir Algèbre, Définition \ref{algebra-definition-smooth}.

\begin{definition}
\label{morphisms-definition-smooth}
Soit $f : X \to S$ un morphisme de schémas.
\begin{enumerate}
\item On dit que $f$ est {\it lisse
à $x \in X$} s'il existe un voisinage ouvert affine $\Spec(A) = U \subset X$
de $x$ et ouvert affine $\Spec(R) = V \subset S$ avec
$f(U) \subset V$ tel que l'homomorphisme d'anneaux induit $R \to A$
soit lisse.
\item On dit que $f$ est {\it lisse} s'il est lisse en chaque point de $X$.
\item Un morphisme de schémas affines $f : X \to S$
est appelé {\it standard lisse} s'il
existe un homomorphisme d'anneaux lisse standard
$R \to R[x_1, \ldots, x_n]/(f_1, \ldots, f_c)$ (voir Algèbre, Définition \ref{algebra-definition-standard-smooth}) tel
que $X \to S$ est isomorphe à
$$
\Spec(R[x_1, \ldots, x_n]/(f_1, \ldots, f_c)) \to \Spec(R).
$$
\end{enumerate}
\end{definition}

\noindent
Une caractéristique intéressante de cette définition est que l'ensemble des points
où un morphisme est lisse est automatiquement ouvert.

\medskip\noindent
A noter qu'il n'y a pas d'hypothèse de séparation ou de quasi-compacité dans
la définition. La question de la fluidité est donc de nature locale sur
la source. Voici le résultat précis.

\begin{lemma}
\label{morphisms-lemma-smooth-characterize}
Soit $f : X \to S$ un morphisme de schémas. Les conditions
suivantes sont équivalentes
\begin{enumerate}
\item Le morphisme $f$ est lisse.
\item Pour chaque ouverture affine $U \subset X$, $V \subset S$
avec $f(U) \subset V$ l'homomorphisme
d'anneaux $\mathcal{O}_S(V) \to \mathcal{O}_X(U)$ est lisse.
\item Il existe un recouvrement ouvert $S = \bigcup_{j \in J} V_j$ et
des recouvrements ouverts $f^{-1}(V_j) = \bigcup_{i \in I_j} U_i$ tels que chacun
des morphismes $U_i \to V_j$, $j\in J, i\in I_j$ soit
lisse.
\item Il existe un recouvrement ouvert affine $S = \bigcup_{j \in J} V_j$
et ouvert affine revêtements $f^{-1}(V_j) = \bigcup_{i \in I_j} U_i$ tel que
l'homomorphisme d'anneaux $\mathcal{O}_S(V_j) \to \mathcal{O}_X(U_i)$ soit lisse,
pour tout $j\in J, i\in I_j$.
\end{enumerate}
De plus, si $f$ est lisse alors
pour tout sous-schémas ouverts $U \subset X$, $V \subset S$ avec $f(U) \subset V$
la restriction $f|_U : U \to V$ est lisse.
\end{lemma}

\begin{proof}
Cela résulte du lemme \ref{morphisms-lemma-locally-P} si l'on montre que la propriété
``$R \to A$ est lisse''
est locale. On vérifie les conditions (a), (b)
et (c) de la Définition \ref{morphisms-definition-property-local}.
Par Algèbre, Lemme \ref{algebra-lemma-base-change-smooth}, la lissité est stable
par changement de base ; on en déduit
la condition (a). Par
Algèbre, Lemme \ref{algebra-lemma-compose-smooth}, la lissité est stable
par composition, et pour tout anneau $R$ l'homomorphisme
d'anneaux $R \to R_f$ est (standard) lisse.
On en déduit la condition (b). Enfin, la
condition (c) résulte d'Algèbre, Lemme \ref{algebra-lemma-locally-smooth}.
\end{proof}

\noindent
Le lemme suivant caractérise un morphisme lisse comme un
morphisme plat, de présentation finie, avec des fibres
lisses. A noter que les schémas lissant un corps sont abordés plus
en détail dans Variétés, section \ref{varieties-section-smooth}.

\begin{lemma}
\label{morphisms-lemma-smooth-flat-smooth-fibres}
Soit $f : X \to S$ un morphisme de schémas. Si $f$
est plat, localement de présentation finie, et que toutes les fibres
$X_s$ sont lisses, alors $f$
est lisse.
\end{lemma}

\begin{proof}
Cela résulte d'Algèbre, Lemme \ref{algebra-lemma-flat-fibre-smooth}.
\end{proof}

\begin{lemma}
\label{morphisms-lemma-composition-smooth}
La composition de deux morphismes lisses est lisse.
\end{lemma}

\begin{proof}
Dans la démonstration du lemme \ref{morphisms-lemma-smooth-characterize} nous avons
vu qu'être lisse est une propriété locale des homomorphismes
d'anneaux. D'où la première affirmation
du lemme découle du lemme \ref{morphisms-lemma-composition-type-P} combinée
au fait qu'être lisse est une propriété des homomorphismes d'anneaux
qui est stable sous
composition, voir Algèbre, Lemme \ref{algebra-lemma-compose-smooth}.
\end{proof}

\begin{lemma}
\label{morphisms-lemma-base-change-smooth}
Le changement de base d'un morphisme qui est lisse est lisse.
\end{lemma}

\begin{proof}
Dans la démonstration du lemme \ref{morphisms-lemma-smooth-characterize} nous avons
vu qu'être lisse est une propriété locale des homomorphismes
d'anneaux. D'où
le lemme découle du lemme \ref{morphisms-lemma-composition-type-P} combiné
au fait qu'être lisse est une propriété des homomorphismes
d'anneaux stables par changement
de base, voir Algèbre, Lemme \ref{algebra-lemma-base-change-smooth}.
\end{proof}

\begin{lemma}
\label{morphisms-lemma-open-immersion-smooth}
Toute immersion ouverte est fluide.
\end{lemma}

\begin{proof}
Ceci est vrai car une immersion ouverte est un isomorphisme local.
\end{proof}

\begin{lemma}
\label{morphisms-lemma-smooth-syntomic}
Un morphisme lisse est syntomique.
\end{lemma}

\begin{proof}
Voir Algèbre, Lemme \ref{algebra-lemma-smooth-syntomic}.
\end{proof}

\begin{lemma}
\label{morphisms-lemma-smooth-locally-finite-presentation}
Un morphisme lisse est localement de présentation finie.
\end{lemma}

\begin{proof}
Vrai car un homomorphisme d'anneaux lisse est de présentation finie par
définition.
\end{proof}

\begin{lemma}
\label{morphisms-lemma-smooth-flat}
Un morphisme lisse est plat.
\end{lemma}

\begin{proof}
Cela résulte des lemmes \ref{morphisms-lemma-syntomic-flat} et \ref{morphisms-lemma-smooth-syntomic}.
\end{proof}

\begin{lemma}
\label{morphisms-lemma-smooth-open}
Un morphisme lisse est universellement ouvert.
\end{lemma}

\begin{proof}
Combine
Lemmes \ref{morphisms-lemma-smooth-flat},
\ref{morphisms-lemma-smooth-locally-finite-presentation} et
\ref{morphisms-lemma-fppf-open}.
Ou bien, combinez
Lemmes \ref{morphisms-lemma-smooth-syntomic},
\ref{morphisms-lemma-syntomic-open}.
\end{proof}

\noindent
Le lemme suivant dit localement que tout morphisme lisse est standard lisse. Nous pouvons donc utiliser
les morphismes lisses standards comme {\it modèle local}
pour un morphisme lisse.

\begin{lemma}
\label{morphisms-lemma-smooth-locally-standard-smooth}
\begin{slogan}
Les morphismes lisses sont des lisses localement standards.
\end{slogan}
Soit $f : X  \to S$ un morphisme de schémas. Soit $x \in X$
un point. Soit $V \subset S$
un voisinage ouvert affine de $f(x)$. Les conditions
suivantes sont équivalentes
\begin{enumerate}
\item Le morphisme $f$ est lisse en $x$.
\item Il existe un $U \subset X$ ouvert affine,
avec $x \in U$ et $f(U) \subset V$ tel que le morphisme
induit $f|_U : U \to V$ est lisse standard.
\end{enumerate}
\end{lemma}

\begin{proof}
Dérivé des définitions
et de l'algèbre, Lemmes \ref{algebra-lemma-standard-smooth}
et \ref{algebra-lemma-smooth-syntomic}.
\end{proof}

\begin{lemma}
\label{morphisms-lemma-smooth-omega-finite-locally-free}
Soit $f : X \to S$ un morphisme de schémas. Supposons que
$f$ soit fluide.
Alors le module des différentielles $\Omega_{X/S}$ de $X$ sur $S$
est fini localement libre et
$$
\text{rang}_x(\Omega_{X/S}) = \dim_x(X_{f(x)})
$$
pour chaque $x \in X$.
\end{lemma}

\begin{proof}
L'instruction est locale sur $X$
et $S$. D'après le Lemme \ref{morphisms-lemma-smooth-locally-standard-smooth}
ci-dessus, nous pouvons supposer que $f$ est un morphisme
lisse standard d'affines.
Dans ce cas le résultat découle de l'Algèbre,
Lemme \ref{algebra-lemma-standard-smooth} (et de la définition d'une intersection
complète
relative globale, voir Algèbre, Définition \ref{algebra-definition-relative-global-complete-intersection}).
\end{proof}

\noindent
Lemme \ref{morphisms-lemma-smooth-omega-finite-locally-free} dit que la définition
suivante a du sens.

\begin{definition}
\label{morphisms-definition-smooth-relative-dimension}
Soit $d \geq 0$ un nombre entier. On dit qu'un morphisme de schémas $f : X \to S$ est {\it
lisse de dimension relative $d$} si $f$ est
lisse et $\Omega_{X/S}$ est fini localement libre de rang constant $d$.
\end{definition}

\noindent
Autrement dit, $f$ est lisse et les fibres non vides sont équidimensionnelles
de dimension $d$. Par le Lemme \ref{morphisms-lemma-smooth-at-point} ci-dessous cela revient
également à exiger : (a) $f$ est localement de présentation finie, (b) $f$
est plat, (c) toutes les fibres non vides équidimensionnelles de dimension
$d$, et (d) $\Omega_{X/S}$ fini localement libre de rang $d$.
Il ne suffit pas de simplement supposer que $f$ est plat, de présentation
finie, et que $\Omega_{X/S}$ est fini localement libre de rang
$d$. Un contre-exemple est donné par $\Spec(\mathbf{F}_p[t]) \to \Spec(\mathbf{F}_p[t^p])$.

\medskip\noindent Voici
un critère différentiel de régularité en un point. Il existe de nombreuses
variantes de ce résultat, qui peuvent
toutes être utiles à un moment donné. Nous les ajouterons simplement ici
si nécessaire.

\begin{lemma}
\label{morphisms-lemma-smooth-at-point}
Soit $f : X \to S$ un morphisme de schémas. Soit $x \in X$.
Posons $s = f(x)$.
Supposons que $f$
soit localement de présentation finie. Les conditions
suivantes sont équivalentes :
\begin{enumerate}
\item Le morphisme $f$ est lisse en $x$.
\item L'homomorphisme d'anneaux locaux $\mathcal{O}_{S, s} \to \mathcal{O}_{X, x}$ est plat
et $X_s \to \Spec(\kappa(s))$ est lisse à $x$.
\item L'homomorphisme local d'anneaux $\mathcal{O}_{S, s} \to \mathcal{O}_{X, x}$ est plat
et le $\mathcal{O}_{X, x}$-module $\Omega_{X/S, x}$ peut être
généré par au plus $\dim_x(X_{f(x)})$ éléments.
\item L'homomorphisme d'anneaux locaux $\mathcal{O}_{S, s} \to \mathcal{O}_{X, x}$ est plat et
l'espace vectoriel $\kappa(x)$-
$$
\Omega_{X_s/s, x} \otimes_{\mathcal{O}_{X_s, x}} \kappa(x) =
\Omega_{X/S, x} \otimes_{\mathcal{O}_{X, x}} \kappa(x)
$$
peut être généré par au plus $\dim_x(X_{f(x)})$ éléments.
\item Il existe des ouverts affines $U \subset X$,
et $V \subset S$ telles que $x \in U$, $f(U) \subset V$ et le
morphisme induit $f|_U : U \to V$ est lisse standard.
\item Il existe des ouverts affines $\Spec(A) = U \subset X$
et $\Spec(R) = V \subset S$ avec $x \in U$ correspondant
à $\mathfrak q \subset A$, et $f(U) \subset V$ telles qu'il
existe une présentation
$$
A = R[x_1, \ldots, x_n]/(f_1, \ldots, f_c)
$$
avec
$$
g =
\det
\left(
\begin{matrix}
\partial f_1/\partial x_1 &
\partial f_2/\partial x_1 &
\ldots &
\partial f_c/\partial x_1 \\
\partial f_1/\partial x_2 &
\partial f_2/\partial x_2 &
\ldots &
\partial f_c/\partial x_2 \\
\ldots & \ldots & \ldots & \ldots \\
\partial f_1/\partial x_c &
\partial f_2/\partial x_c &
\ldots &
\partial f_c/\partial x_c
\end{matrix}
\right)
$$
mappant à un élément de $A$ pas dans $\mathfrak q$.
\end{enumerate}
\end{lemma}

\begin{proof}
Notez que si $f$ est lisse à $x$, alors
nous voyons d'après le Lemme \ref{morphisms-lemma-smooth-locally-standard-smooth} que (5) est valable et (6)
est une version légèrement affaiblie de (5). De plus, si $f$
est lisse, l'homomorphisme d'anneaux
$\mathcal{O}_{S, s} \to \mathcal{O}_{X, x}$ est plat (voir Lemme \ref{morphisms-lemma-smooth-flat}) et que
$\Omega_{X/S}$ est fini localement
libre de rang égal à $\dim_x(X_s)$ (voir Lemme \ref{morphisms-lemma-smooth-omega-finite-locally-free}). Ainsi
(1) implique (3) et (4). Par le Lemme \ref{morphisms-lemma-base-change-smooth} on voit
aussi que (1) implique (2).

\medskip\noindent
Par le Lemme \ref{morphisms-lemma-base-change-differentials} le module des différentielles
$\Omega_{X_s/s}$ de la fibre $X_s$ sur $\kappa(s)$
est le retrait du module des différentielles $\Omega_{X/S}$
de $X$ sur $S$. D'où l'égalité affichée
dans la partie (4) du lemme. Par le Lemme \ref{morphisms-lemma-finite-type-differentials} ces modules
sont de type fini. Par conséquent, le nombre minimal de
générateurs des modules
$\Omega_{X/S, x}$ et $\Omega_{X_s/s, x}$ est le même et égal à la dimension
de cet espace vectoriel $\kappa(x)$ par le Lemme de Nakayama
(Algèbre, Lemme \ref{algebra-lemma-NAK}). Cela montre notamment que (3) et
(4) sont équivalents.

\medskip\noindent
Algèbre, Lemme \ref{algebra-lemma-flat-fibre-smooth} montre que (2)
implique
(1). Algèbre, Lemme \ref{algebra-lemma-characterize-smooth-over-field} montre que (3)
et (4) impliquent (2). Enfin, (6) implique
(5) voir par exemple Algèbre, Exemple \ref{algebra-example-make-standard-smooth}
et (5) implique (1) par Algèbre, Lemme \ref{algebra-lemma-standard-smooth}.
\end{proof}

\begin{lemma}
\label{morphisms-lemma-set-points-where-fibres-smooth}
Soit
$$
\xymatrix{
X' \ar[r]_{g'} \ar[d]_{f'} & X \ar[d]^f \\
S' \ar[r]^g & S
}
$$
un diagramme cartésien de schémas. Soit $W \subset X$, resp.\ $W' \subset X'$ le sous-schéma
ouvert des points où $f$, resp.\ $f'$ est
lisse. Alors $W' = (g')^{-1}(W)$ si
\begin{enumerate}
\item $f$ est plat et localement de présentation finie, ou
\item $f$ est localement de présentation finie et $g$ est plat.
\end{enumerate}
\end{lemma}

\begin{proof}
Supposons d'abord que $f$ soit localement de type fini. Considérons l'ensemble
$$
T = \{x \in X \mid X_{f(x)}\text{ est lisse sur }\kappa(f(x))\text{ en }x\}
$$
et l'ensemble correspondant $T' \subset X'$ pour $f'$. Nous affirmons alors que
$T' = (g')^{-1}(T)$. En effet, soit $s' \in S'$ un point, et posons
$s = g(s')$. On a alors
$$
X'_{s'} =
\Spec(\kappa(s')) \times_{\Spec(\kappa(s))} X_s
$$
Autrement dit les fibres du changement de base sont les changements
de base des fibres. L'affirmation
est donc équivalente à l'algèbre, Lemme \ref{algebra-lemma-smooth-field-change-local}.

\medskip\noindent
Ainsi le cas (1) en résulte, car $T$ est alors l'ensemble (ouvert)
des points où $f$ est lisse, d'après le lemme \ref{morphisms-lemma-smooth-at-point}.

\medskip\noindent
Dans le cas (2), soit $x' \in W'$. Alors $g'$
est plat en $x'$ (Lemme \ref{morphisms-lemma-base-change-module-flat})
et $g \circ f'$ est plat en $x'$ (Lemme \ref{morphisms-lemma-composition-module-flat}). Il s'ensuit
que $f$ est plat en $x = g'(x')$
par le lemme \ref{morphisms-lemma-flat-permanence}. D'autre part, puisque $x' \in T'$
(Lemme \ref{morphisms-lemma-base-change-smooth}) on voit que $x \in T$.
Par conséquent, $f$ est lisse en $x$
par le lemme \ref{morphisms-lemma-smooth-at-point}.
\end{proof}

\noindent
Voici un lemme qui utilise en fait la disparition de $H^{-1}$ du complexe
cotangent naïf pour un homomorphisme d'anneaux lisse.

\begin{lemma}
\label{morphisms-lemma-triangle-differentials-smooth}
Soient $f : X \to Y$, $g : Y \to S$ des morphismes de schémas. Supposons que
$f$ soit fluide. Alors
$$
0 \to f^*\Omega_{Y/S} \to \Omega_{X/S} \to \Omega_{X/Y} \to 0
$$
(voir Lemme \ref{morphisms-lemma-triangle-differentials}) est court et exact.
\end{lemma}

\begin{proof}
La version algébrique de ce lemme est la suivante : Étant donnés
les homomorphismes d'anneaux $A \to B \to C$ avec $B \to C$ lisses, alors la séquence
$$
0 \to C \otimes_B \Omega_{B/A} \to \Omega_{C/A} \to \Omega_{C/B} \to 0
$$
d'Algèbre,
Lemme \ref{algebra-lemma-exact-sequence-differentials} est exacte.
C'est
de l'algèbre, Lemme \ref{algebra-lemma-triangle-differentials-smooth}.
\end{proof}

\begin{lemma}
\label{morphisms-lemma-differentials-relative-immersion-smooth}
Soit $i : Z \to X$ une immersion de schémas sur $S$. Supposons
que $Z$ soit lisse sur $S$. Alors
la suite exacte canonique
$$
0 \to \mathcal{C}_{Z/X} \to i^*\Omega_{X/S} \to \Omega_{Z/S} \to 0
$$
du
lemme \ref{morphisms-lemma-differentials-relative-immersion} est exacte
courte.
\end{lemma}

\begin{proof}
La version algébrique de ce lemme est la suivante : Étant donné
les homomorphismes d'anneaux $A \to B \to C$ avec $A \to C$ lisse et $B \to C$
surjectif avec noyau $J$, alors la séquence
$$
0 \to J/J^2 \to C \otimes_B \Omega_{B/A} \to \Omega_{C/A} \to 0
$$
d'Algèbre,
Lemme \ref{algebra-lemma-differential-seq} est exacte.
C'est
de l'algèbre, Lemme \ref{algebra-lemma-differential-seq-smooth}.
\end{proof}

\begin{lemma}
\label{morphisms-lemma-two-immersions-smooth}
Soit
$$
\xymatrix{
Z \ar[r]_i \ar[rd]_j & X \ar[d] \\
& Y
}
$$
un diagramme commutatif de schémas où $i$ et $j$ sont des immersions
et $X \to Y$ est lisse.
Alors la séquence canonique exacte
$$
0 \to  \mathcal{C}_{Z/Y} \to \mathcal{C}_{Z/X} \to i^*\Omega_{X/Y} \to 0
$$
de
Lemme \ref{morphisms-lemma-two-immersions}
est exacte.
\end{lemma}

\begin{proof}
La version algébrique de ce lemme est la suivante : Étant donnés
les homomorphismes d'anneaux $A \to B \to C$ avec $A \to C$ surjectif et $A \to B$
lisse, alors la séquence
$$
0 \to I/I^2 \to J/J^2 \to C \otimes_B \Omega_{B/A} \to 0
$$
d'Algèbre,
Lemme \ref{algebra-lemma-application-NL} est
exacte. C'est
de l'algèbre, Lemme \ref{algebra-lemma-application-NL-smooth}.
\end{proof}

\begin{lemma}
\label{morphisms-lemma-smooth-permanence}
Soit
$$
\xymatrix{
X \ar[rr]_f \ar[rd]_p & &
Y \ar[dl]^q \\
& S
}
$$
un diagramme commutatif de morphismes de schémas. Supposons que
\begin{enumerate}
\item $f$ est surjectif, et lisse,
\item $p$ est lisse, et
\item $q$ est localement de présentation finie\footnote{En
fait, cela est
impliqué par (1) et (2), voir Descente, Lemme \ref{descent-lemma-flat-finitely-presented-permanence}. De plus,
il suffit de supposer que $f$ est surjectif, plat et
localement de fini présentation,
voir Descente, Lemme \ref{descent-lemma-smooth-permanence}.}.
\end{enumerate}
Alors $q$ est lisse.
\end{lemma}

\begin{proof}
Par le Lemme \ref{morphisms-lemma-flat-permanence} on voit que $q$ est plat. Choisissez
un point $y \in Y$. Choisissez un point $x \in X$ envoyé
sur $y$. Supposons que $f$ ait une dimension relative
$a$ à $x$ et que $p$ ait une dimension relative $b$
à $x$. Par le Lemme \ref{morphisms-lemma-smooth-omega-finite-locally-free} cela signifie que $\Omega_{X/S, x}$
est libre de rang $b$ et $\Omega_{X/Y, x}$ est
libre de rang $a$. Par la suite exacte courte
du lemme \ref{morphisms-lemma-triangle-differentials-smooth} cela signifie que $(f^*\Omega_{Y/S})_x$
est libre de rang $b - a$. D'après le Lemme de
Nakayama, cela implique que $\Omega_{Y/S, y}$ peut être engendré
par $b - a$ éléments. De plus, le lemme \ref{morphisms-lemma-dimension-fibre-at-a-point-additive} montre
que $\dim_y(Y_s) = b - a$. Nous concluons donc que
$Y \to S$ est lisse en $y$ d'après le Lemme \ref{morphisms-lemma-smooth-at-point}, partie (2).
\end{proof}

\noindent
Dans la situation du lemme suivant l'image de $\sigma$
est localement sur $X$ découpée par
une séquence régulière, voir Diviseurs, Lemme \ref{divisors-lemma-section-smooth-regular-immersion}.

\begin{lemma}
\label{morphisms-lemma-section-smooth-morphism}
Soit $f : X \to S$ un morphisme de schémas. Soit
$\sigma : S \to X$ une section de $f$.
Soit $s \in S$ un point tel que $f$ soit lisse en $x = \sigma(s)$.
Alors il existe des voisinages affines
ouverts $\Spec(A) = U \subset S$ de $s$ et $\Spec(B) = V \subset X$ de $x$
tels que
\begin{enumerate}
\item $f(V) \subset U$ et $\sigma(U) \subset V$,
\item avec $I = \Ker(\sigma^\# : B \to A)$ le module $I/I^2$ est un module
$A$ libre, et
\item $B^\wedge \cong A[[x_1, \ldots, x_d]]$ comme $A$-algèbres où $B^\wedge$ désigne
l'achèvement de $B$ par rapport à $I$.
\end{enumerate}
\end{lemma}

\begin{proof}
Choisissez un ouvert affine $U \subset S$ contenant
$s$ Choisissez un ouvert affine $V \subset f^{-1}(U)$ contenant
$x$. Choisissez un $U' \subset \sigma^{-1}(V)$ ouvert affine contenant
$s$. Notez que $V' = f^{-1}(U') \cap V$ est affine car il est égal au produit
fibré $V' = U' \times_U V$. Alors $U'$ et $V'$ satisfont (1).
Écrivez $U' = \Spec(A')$ et $V' = \Spec(B')$. Par algèbre,
Lemme \ref{algebra-lemma-section-smooth} le module $I'/(I')^2$ est fini
localement libre en tant que $A'$-module. Ainsi, après
avoir remplacé $U'$ par un ouvert affine plus petit $U'' \subset U'$
et $V'$ par $V'' = V' \cap f^{-1}(U'')$ nous obtenons la situation où $I''/(I'')^2$
est libre, c'est-à-dire que (2) est valable. Dans ce cas (3) est également
valable pour l'algèbre, Lemme \ref{algebra-lemma-section-smooth}.
\end{proof}

\noindent
La dimension d'un schéma $X$ en un
point $x$ (Propriétés, Définition \ref{properties-definition-dimension}) est juste
la dimension de $X$ en $x$ comme espace
topologique, voir Topologie, Définition \ref{topology-definition-Krull}.
Ce n'est pas la dimension de l'anneau local $\mathcal{O}_{X,x}$, en général.

\begin{lemma}
\label{morphisms-lemma-smoothness-dimension}
Soit $f : X \to Y$ un morphisme lisse de schémas localement noethériens. Pour tout point $x$
de $X$, d'image $y$ dans $Y$,
$$
\dim_x(X) = \dim_y(Y) + \dim_x(X_y),
$$
où $X_y$ désigne la fibre en $y$.
\end{lemma}

\begin{proof}
Après avoir remplacé $X$ par un voisinage
ouvert de $x$, il existe un nombre
naturel $d$ tel que toutes les fibres de
$X \to Y$ ont la dimension $d$ en tout
point, voir Lemme \ref{morphisms-lemma-smooth-omega-finite-locally-free}. Alors $f$
est plat (Lemme \ref{morphisms-lemma-smooth-flat}), localement de type fini (Lemme \ref{morphisms-lemma-smooth-locally-finite-presentation}),
et de dimension relative $d$. Le résultat
découle donc du lemme \ref{morphisms-lemma-rel-dimension-dimension}.
\end{proof}












\section{Morphismes non ramifiés}
\label{morphisms-section-unramified}

\noindent
Nous discutons brièvement des morphismes non ramifiés avant la classe (peut-être) plus intéressante
des morphismes étales. Rappelons qu'un homomorphisme d'anneaux $R \to A$
est {\it non ramifié} s'il est de type fini et $\Omega_{A/R} = 0$
(c'est la définition de \cite{Henselian}). Un homomorphisme d'anneaux $R \to A$ est appelé
{\it G-unramified} s'il est de présentation finie et $\Omega_{A/R} = 0$
(c'est la
définition de \cite{EGA}). Voir Algèbre, Définition \ref{algebra-definition-unramified}.

\begin{definition}
\label{morphisms-definition-unramified}
Soit $f : X \to S$ un morphisme de schémas.
\begin{enumerate}
\item On dit que $f$ est {\it non ramifié
en $x \in X$} s'il existe un voisinage ouvert affine $\Spec(A) = U \subset X$
de $x$ et ouvert affine $\Spec(R) = V \subset S$ avec $f(U) \subset V$
tel que l'homomorphisme d'anneaux induit $R \to A$
ne soit pas ramifié.
\item On dit que $f$ est {\it G-non ramifié
en $x \in X$} s'il existe un voisinage ouvert affine $\Spec(A) = U \subset X$
de $x$ et ouvert affine $\Spec(R) = V \subset S$ avec
$f(U) \subset V$ tel que l'homomorphisme d'anneaux induit $R \to A$
soit G-non ramifié.
\item On dit que $f$ est {\it non ramifié} s'il est non ramifié
en tout point de $X$.
\item On dit que $f$ est {\it G-non ramifié} s'il est G-non ramifié
en tout point de $X$.
\end{enumerate}
\end{definition}

\noindent
Notons qu'un G-morphisme non ramifié est non ramifié. Ainsi, tout résultat pour
les morphismes non ramifiés implique le résultat correspondant pour les G-morphismes
non ramifiés. De plus, si $S$ est localement noethérien alors il n'y a pas
de différence entre G-unramifié et morphismes non ramifiés,
voir Lemme \ref{morphisms-lemma-noetherian-unramified}. Une caractéristique
intéressante de cette définition est que l'ensemble des points où un morphisme
est non ramifié (resp.\ G-non ramifié) est automatiquement ouvert.

\begin{lemma}
\label{morphisms-lemma-unramified-omega-zero}
Soit $f : X \to S$ un morphisme de schémas. Ensuite,
\begin{enumerate}
\item $f$ est non ramifié si et seulement si $f$ est localement de
type fini et $\Omega_{X/S} = 0$, et
\item $f$ est G-unramifié si et seulement si $f$ est localement de présentation
finie et $\Omega_{X/S} = 0$.
\end{enumerate}
\end{lemma}

\begin{proof}
Par définition un homomorphisme d'anneaux $R \to A$ est unramifié (resp.\ G-unramifié)
si et seulement si il est de type fini (resp.\ présentation
finie) et $\Omega_{A/R} = 0$.
Le lemme découle donc directement des définitions et du Lemme \ref{morphisms-lemma-differentials-affine}.
\end{proof}

\noindent
Notez qu'il n'y a pas d'hypothèse de séparation ou de quasi-compacité dans
la définition. La question de la non-ramification est donc de nature locale
sur la source. Voici le résultat précis.

\begin{lemma}
\label{morphisms-lemma-unramified-characterize}
Soit $f : X \to S$ un morphisme de schémas. Les conditions
suivantes sont équivalentes
\begin{enumerate}
\item Le morphisme $f$ est non ramifié (resp.\ G-non ramifié).
\item Pour tout ouvert affine $U \subset X$, $V \subset S$
avec $f(U) \subset V$ l'homomorphisme
d'anneaux $\mathcal{O}_S(V) \to \mathcal{O}_X(U)$ est non ramifié (resp.\ G-unramifié).
\item Il existe un recouvrement ouvert $S = \bigcup_{j \in J} V_j$ et des
recouvrements ouverts $f^{-1}(V_j) = \bigcup_{i \in I_j} U_i$ tels que chacun des
morphismes $U_i \to V_j$, $j\in J, i\in I_j$ soit non ramifié (resp.\ G-non
ramifié).
\item Il existe un recouvrement ouvert affine $S = \bigcup_{j \in J} V_j$
et ouvert affine revêtements $f^{-1}(V_j) = \bigcup_{i \in I_j} U_i$ tel que l'homomorphisme
d'anneaux $\mathcal{O}_S(V_j) \to \mathcal{O}_X(U_i)$ soit non ramifié (resp.\ G-unramifié),
pour tout $j\in J, i\in I_j$.
\end{enumerate}
De plus, si $f$ est non ramifié (resp.\ G-unramifié)
alors pour tout sous-schémas ouverts $U \subset X$, $V \subset S$ avec $f(U) \subset V$
la restriction $f|_U : U \to V$ est non ramifiée (resp.\ G-non ramifié).
\end{lemma}

\begin{proof}
Cela résulte du lemme \ref{morphisms-lemma-locally-P} si l'on montre que la propriété
``$R \to A$ est non ramifié''
est locale. On vérifie les conditions (a),
(b) et (c) de la Définition
\ref{morphisms-definition-property-local}. Ces propriétés sont
prouvées dans Algèbre, Lemme \ref{algebra-lemma-unramified}.
\end{proof}

\begin{lemma}
\label{morphisms-lemma-composition-unramified}
La composition de deux morphismes non ramifiés est non ramifiée. Il en va
de même pour les G-morphismes non ramifiés.
\end{lemma}

\begin{proof}
La démonstration du lemme \ref{morphisms-lemma-unramified-characterize} montre qu'être
non ramifié (resp.\ G-non ramifié) est
une propriété locale des homomorphismes
d'anneaux. D'où le premier énoncé du lemme
découle du lemme \ref{morphisms-lemma-composition-type-P} combiné au fait
qu'être non ramifié (resp.\ G-non ramifié)
est une propriété des homomorphismes d'anneaux stable sous composition,
voir Algèbre, Lemme \ref{algebra-lemma-unramified}.
\end{proof}

\begin{lemma}
\label{morphisms-lemma-base-change-unramified}
Le changement de base d'un morphisme non ramifié est non ramifié. Il en
va de même pour les G-morphismes non ramifiés.
\end{lemma}

\begin{proof}
La démonstration du lemme \ref{morphisms-lemma-unramified-characterize} montre qu'être non
ramifié (resp.\ G-non ramifié) est une
propriété locale des homomorphismes d'anneaux. D'où le lemme
découle du lemme \ref{morphisms-lemma-base-change-type-P} combiné au fait
qu'être non ramifié (resp.\ G-non ramifié) est
une propriété des homomorphismes d'anneaux stable par changement
de base, voir Algèbre, Lemme \ref{algebra-lemma-unramified}.
\end{proof}

\begin{lemma}
\label{morphisms-lemma-noetherian-unramified}
Soit $f : X \to S$ un morphisme de schémas. Supposons que $S$ soit localement noethérien. Alors $f$
est non ramifié si et seulement si $f$ est G-non ramifié.
\end{lemma}

\begin{proof}
Découle des définitions
et du lemme \ref{morphisms-lemma-noetherian-finite-type-finite-presentation}.
\end{proof}


\begin{lemma}
\label{morphisms-lemma-open-immersion-unramified}
Toute immersion ouverte est G-non ramifiée.
\end{lemma}

\begin{proof}
Ceci est vrai car une immersion ouverte est un isomorphisme local.
\end{proof}

\begin{lemma}
\label{morphisms-lemma-closed-immersion-unramified}
Une immersion fermée $i : Z \to X$ n'est pas
ramifiée. Il est G-non ramifié si et seulement si le faisceau
quasi-cohérent d'idéaux associé $\mathcal{I} = \Ker(\mathcal{O}_X \to i_*\mathcal{O}_Z)$ est de type
fini (en tant que $\mathcal{O}_X$-module).
\end{lemma}

\begin{proof}
Cela résulte du Lemme \ref{morphisms-lemma-closed-immersion-finite-presentation} et d'Algèbre,
Lemme \ref{algebra-lemma-unramified}.
\end{proof}

\begin{lemma}
\label{morphisms-lemma-unramified-locally-finite-type}
An morphisme non ramifié est localement de type fini. Un
G-morphisme non ramifié est localement de présentation finie.
\end{lemma}

\begin{proof}
Un homomorphisme d'anneaux non ramifié est de type fini par définition.
Un homomorphisme d'anneaux G-non ramifié est de présentation finie par définition.
\end{proof}

\begin{lemma}
\label{morphisms-lemma-unramified-quasi-finite}
Soit $f : X \to S$ un morphisme de schémas. Si $f$
est non ramifié en $x$ alors $f$ est quasi-fini en $x$.
En particulier, un morphisme non ramifié est localement quasi-fini.
\end{lemma}

\begin{proof}
Voir Algèbre, Lemme \ref{algebra-lemma-unramified-quasi-finite}.
\end{proof}

\begin{lemma}
\label{morphisms-lemma-unramified-over-field}
Fibres de morphismes non ramifiés.
\begin{enumerate}
\item Soit $X$ un schéma sur un corps $k$.
Le morphisme de structure $X \to \Spec(k)$ est non ramifié si et seulement
si $X$ est une union disjointe de spectres d'extensions de corps finis
séparables de $k$.
\item Si $f : X \to S$ est un morphisme non ramifié alors pour tout $s \in S$ la fibre
$X_s$ est une union disjointe de spectres d'extensions de corps finis
séparables de $\kappa(s)$.
\end{enumerate}
\end{lemma}

\begin{proof}
La partie (2) découle de la partie
(1) et Lemme \ref{morphisms-lemma-base-change-unramified}. Démontrons
la partie (1). Nous utilisons
d'abord l'algèbre, Lemme \ref{algebra-lemma-characterize-unramified}. Ce lemme implique que
si $X$ est une union disjointe de spectres d'extensions de corps
finis séparables de $k$ alors $X \to \Spec(k)$ est non ramifié.
À l'inverse, supposons que $X \to \Spec(k)$ ne soit
pas ramifié. Par Algèbre, Lemme \ref{algebra-lemma-unramified-at-prime} pour tout $x \in X$ l'extension
du corps résiduel $\kappa(x)/k$ est fini
séparable. Puisque $X \to \Spec(k)$ est localement quasi-fini
(Lemme \ref{morphisms-lemma-unramified-quasi-finite}), nous voyons que tous les
points de $X$ sont des points fermés isolés, voir
Lemme \ref{morphisms-lemma-quasi-finite-at-point-characterize}. Ainsi $X$ est un espace
discret, en particulier l'union disjointe des spectres
de ses anneaux locaux. Par algèbre,
Lemme \ref{algebra-lemma-unramified-at-prime} encore une fois, ces anneaux locaux sont
des corps, et on obtient le résultat.
\end{proof}

\noindent
Le lemme suivant caractérise les morphismes non ramifiés comme des
morphismes localement de type fini à fibres non ramifiées.

\begin{lemma}
\label{morphisms-lemma-unramified-etale-fibres}
Soit $f : X \to S$ un morphisme de schémas.
\begin{enumerate}
\item Si $f$ est non ramifié, alors pour tout $x \in X$ l'extension
du corps $\kappa(x)/\kappa(f(x))$ est finie séparable.
\item Si $f$ est localement de type fini, et pour chaque
$s \in S$ la fibre $X_s$ est une union disjointe de spectres d'extensions de corps finis
séparables de $\kappa(s)$ alors $f$ est non ramifié.
\item Si $f$ est localement de présentation finie, et pour chaque
$s \in S$ la fibre $X_s$ est une union disjointe de spectres d'extensions de corps
finis séparables de $\kappa(s)$ alors $f$ est G-non ramifié.
\end{enumerate}
\end{lemma}

\begin{proof}
Cela résulte d'Algèbre,
les Lemmes \ref{algebra-lemma-unramified-at-prime}
et \ref{algebra-lemma-characterize-unramified}.
\end{proof}

\noindent
Voici une caractérisation des morphismes non ramifiés en termes de diagonale
du morphisme.

\begin{lemma}
\label{morphisms-lemma-diagonal-unramified-morphism}
Soit $f : X \to S$ un morphisme.
\begin{enumerate}
\item Si $f$ est non ramifié, alors le morphisme
diagonal $\Delta : X \to X \times_S X$ est une immersion ouverte.
\item Si $f$ est localement de type fini et
$\Delta$ est une immersion ouverte, alors $f$ est non ramifié.
\item Si $f$ est localement de présentation finie et $\Delta$ est une immersion
ouverte, alors $f$ est G-non ramifié.
\end{enumerate}
\end{lemma}

\begin{proof}
La première instruction découle
de l'algèbre, Lemme \ref{algebra-lemma-diagonal-unramified}. La deuxième
affirmation vient du fait que $\Omega_{X/S}$ est le
faisceau conormal du morphisme diagonal
(Lemme \ref{morphisms-lemma-differentials-diagonal}) et donc clairement
nul si $\Delta$ est une immersion ouverte.
\end{proof}

\begin{lemma}
\label{morphisms-lemma-unramified-at-point}
Soit $f : X \to S$ un morphisme de schémas. Soit $x \in X$.
Posons $s = f(x)$.
Supposons que $f$
est localement de type fini (resp.\ localement de présentation finie). Les conditions
suivantes sont équivalentes :
\begin{enumerate}
\item Le morphisme $f$ est non ramifié (resp.\ G-non ramifié) en $x$.
\item La fibre $X_s$ est non ramifiée sur $\kappa(s)$ à $x$.
\item Le module $\mathcal{O}_{X, x}$ $\Omega_{X/S, x}$ est nul.
\item Le module $\mathcal{O}_{X_s, x}$ $\Omega_{X_s/s, x}$ est nul.
\item L'espace vectoriel $\kappa(x)$
$$
\Omega_{X_s/s, x} \otimes_{\mathcal{O}_{X_s, x}} \kappa(x) =
\Omega_{X/S, x} \otimes_{\mathcal{O}_{X, x}} \kappa(x)
$$
est nul.
\item Nous avons $\mathfrak m_s\mathcal{O}_{X, x} = \mathfrak m_x$ et l'extension du corps
$\kappa(x)/\kappa(s)$ est séparable de manière
finie.
\end{enumerate}
\end{lemma}

\begin{proof}
Notons que si $f$ est non ramifié en
$x$, alors on voit que $\Omega_{X/S} = 0$ à proximité de
$x$ par les définitions et les résultats sur modules
de différentiels dans la section \ref{morphisms-section-sheaf-differentials}. Par conséquent (1)
implique (3) et la disparition de l'espace vectoriel de droite
dans (5). Cela implique également
(2) car par le lemme \ref{morphisms-lemma-base-change-differentials} le module
des différentielles $\Omega_{X_s/s}$ de la fibre $X_s$ sur
$\kappa(s)$ est le retrait du module des différentielles $\Omega_{X/S}$
de $X$ sur $S$. Ce fait sur les modules de différentielles
implique également l'égalité affichée des espaces vectoriels
dans la partie (4). Par le Lemme \ref{morphisms-lemma-finite-type-differentials}
les modules $\Omega_{X/S, x}$ et $\Omega_{X_s/s, x}$ sont de type fini. D'où les modules
$\Omega_{X/S, x}$ et $\Omega_{X_s/s, x}$ sont nuls si et seulement si l'espace vectoriel
$\kappa(x)$ correspondant dans (4) est nul par le Lemme
de Nakayama (Algèbre,
Lemme \ref{algebra-lemma-NAK}). Cela montre
notamment que (3), (4) et (5) sont équivalents. Le support
de $\Omega_{X/S}$ est terminé dans $X$, voir
Modules, Lemme \ref{modules-lemma-support-finite-type-closed}. L'hypothèse (3) implique que
$x$ n'est pas dans le support. Donc $\Omega_{X/S}$
est nul au voisinage de $x$, ce qui implique
(1). L'équivalence de (1) et (3) appliquée à $X_s \to s$ implique
l'équivalence de (2) et (4).
À ce stade, nous avons vu que (1) -- (5) sont équivalents.

\medskip\noindent
Vous pouvez également utiliser Algèbre, Lemme \ref{algebra-lemma-unramified} pour voir
l'équivalence de (1) -- (5) plus directement.

\medskip\noindent
L'équivalence de (1) et (6)
découle du lemme \ref{morphisms-lemma-unramified-etale-fibres}.
Cela découle aussi
plus directement de l'Algèbre, de Lemmes
\ref{algebra-lemma-unramified-at-prime} et \ref{algebra-lemma-characterize-unramified}.
\end{proof}

\begin{lemma}
\label{morphisms-lemma-set-points-where-fibres-unramified}
Soit $f : X \to S$ un morphisme de schémas. Supposons
que $f$ soit localement de type fini. La formation de l'ouvert
\begin{align*}
T
& =
\{x \in X \mid X_{f(x)}\text{ est non ramifié sur }\kappa(f(x))\text{ en }x\} \\
& =
\{x \in X \mid X\text{ est non ramifié sur }S\text{ en }x\}
\end{align*}
commute à tout changement de base :
pour tout morphisme $g : S' \to S$, considérons
le changement de base $f' : X' \to S'$ de $f$
et la projection $g' : X' \to X$. Alors l'ensemble correspondant
$T'$ pour le morphisme $f'$ est égal à $T' = (g')^{-1}(T)$. Si $f$
est supposé localement de présentation finie alors il en va de même pour
l'ensemble ouvert de points où $f$ est G-non ramifié.
\end{lemma}

\begin{proof}
Soit $s' \in S'$ un point et soit $s = g(s')$. On a alors
$$
X'_{s'} =
\Spec(\kappa(s')) \times_{\Spec(\kappa(s))} X_s
$$
Autrement dit les fibres du changement de base sont les changements de base
des fibres. En particulier
$$
\Omega_{X_s/s, x} \otimes_{\mathcal{O}_{X_s, x}} \kappa(x')
=
\Omega_{X'_{s'}/s', x'} \otimes_{\mathcal{O}_{X'_{s'}, x'}} \kappa(x')
$$
voir
Lemme \ref{morphisms-lemma-base-change-differentials}. D'où $x' \in T'$ si
et seulement si $x \in T$ par le lemme
\ref{morphisms-lemma-unramified-at-point}. La deuxième partie
découle de la première car dans ce cas $T$
est l'ensemble (ouvert) de points où $f$ est G-non
ramifié d'après le Lemme \ref{morphisms-lemma-unramified-at-point}.
\end{proof}

\begin{lemma}
\label{morphisms-lemma-unramified-permanence}
Soit $f : X \to Y$ un morphisme de schémas sur $S$.
\begin{enumerate}
\item Si $X$ est non ramifié sur $S$, alors $f$ est non ramifié.
\item Si $X$ est G-non ramifié sur $S$ et que $Y$ est localement de type fini sur
$S$, alors $f$ est G-non ramifié.
\end{enumerate}
\end{lemma}

\begin{proof}
Supposons que $X$ est non ramifié
sur $S$. Par le Lemme \ref{morphisms-lemma-permanence-finite-type} on voit que $f$
est localement de type fini.
Par hypothèse, nous avons $\Omega_{X/S} = 0$. D'où $\Omega_{X/Y} = 0$
du lemme \ref{morphisms-lemma-triangle-differentials}. Ainsi, $f$ est non ramifié. Si
$X$ est G-non ramifié sur $S$ et $Y$ est localement
de type fini sur $S$, alors
par le lemme \ref{morphisms-lemma-finite-presentation-permanence} on voit que $f$
est localement de présentation finie et Nous concluons que $f$
est G-non ramifié.
\end{proof}

\begin{lemma}
\label{morphisms-lemma-value-at-one-point}
Soit $S$ un schéma.
Soient $X$, $Y$ des schémas sur
$S$. Soit $f, g : X \to Y$ des morphismes sur $S$. Soit $x \in X$.
Supposons que
\begin{enumerate}
\item le morphisme de structure $Y \to S$ est non ramifié,
\item $f(x) = g(x)$ dans $Y$, disons $y = f(x) = g(x)$, et
\item les applications induites $f^\sharp, g^\sharp : \kappa(y) \to \kappa(x)$ sont
égales.
\end{enumerate}
Alors il existe un voisinage ouvert de $x$ dans $X$ sur lequel $f$
et $g$ sont égaux.
\end{lemma}

\begin{proof}
Considérons le morphisme $(f, g) : X \to Y \times_S Y$. Par l'hypothèse (1) et Lemme \ref{morphisms-lemma-diagonal-unramified-morphism}
l'image réciproque de $\Delta_{Y/S}(Y)$ est ouverte
dans $X$. Et les hypothèses (2) et (3) impliquent
que $x$ se trouve dans ce sous-ensemble ouvert.
\end{proof}










\section{Morphismes étales}
\label{morphisms-section-etale}

\noindent
La topologie Zariski d'un schéma est une topologie très grossière.
Cela est particulièrement clair lorsque l'on examine les variétés supérieures
à $\mathbf{C}$. Il s'avère que déclarer un morphisme étale comme l'analogue
d'un isomorphisme local en topologie introduit une topologie beaucoup
plus fine. Sur les variétés supérieures à $\mathbf{C}$, cette topologie donne
lieu aux ``corrects'' nombres de Betti lors du calcul de
cohomologie à coefficients finis. Un autre observable est que si $f : X \to Y$ est un
morphisme étale des variétés sur $\mathbf{C}$, et si
$x$ est un point fermé
de $X$, alors $f$ induit un isomorphisme $\mathcal{O}^{\wedge}_{Y, f(x)} \to \mathcal{O}^{\wedge}_{X, x}$ d'anneaux
locaux complets.

\medskip\noindent Dans
cette section, nous commençons notre étude de ces questions. En fait, nous limitons
délibérément notre discussion au minimum puisque nous discuterons ailleurs de résultats
plus intéressants. Rappelons qu'un homomorphisme d'anneaux $R \to A$ est dit {\it
étale} s'il est lisse
et $\Omega_{A/R} = 0$, voir Algèbre, Définition \ref{algebra-definition-etale}.

\begin{definition}
\label{morphisms-definition-etale}
Soit $f : X \to S$ un morphisme de schémas.
\begin{enumerate}
\item On dit que $f$ est {\it étale
à $x \in X$} s'il existe un voisinage ouvert affine $\Spec(A) = U \subset X$
de $x$ et ouvert affine $\Spec(R) = V \subset S$ avec $f(U) \subset V$
tel que l'homomorphisme d'anneaux induit $R \to A$
soit étale.
\item On dit que $f$ est {\it étale} s'il est étale en tout point de $X$.
\item Un morphisme de schémas affines $f : X \to S$ est appelé {\it
standard étale} si $X \to S$ est isomorphe à
$$
\Spec(R[x]_h/(g)) \to \Spec(R)
$$
où $R \to R[x]_h/(g)$ est un standard étale homomorphisme d'anneaux,
voir Algèbre, Définition \ref{algebra-definition-standard-etale}, c'est-à-dire que $g$
est monique et $g'$ inversible dans $R[x]_h/(g)$.
\end{enumerate}
\end{definition}

\noindent
Un morphisme est étale si et seulement si il est lisse de dimension relative
$0$ (voir Définition \ref{morphisms-definition-smooth-relative-dimension}). Une caractéristique intéressante
de la définition est que l'ensemble des points où un morphisme
est étale est automatiquement ouvert.

\medskip\noindent Notez
qu'il n'y a pas d'hypothèse de séparation ou de quasi-compacité dans la définition.
La question de l'existence de étale est donc de nature locale
sur la source. Voici le résultat précis.

\begin{lemma}
\label{morphisms-lemma-etale-characterize}
Soit $f : X \to S$ un morphisme de schémas. Les conditions
suivantes sont équivalentes
\begin{enumerate}
\item Le morphisme $f$ est étale.
\item Pour chaque ouvert affines $U \subset X$, $V \subset S$
avec $f(U) \subset V$ l'homomorphisme
d'anneaux $\mathcal{O}_S(V) \to \mathcal{O}_X(U)$ est étale.
\item Il existe un recouvrement ouvert $S = \bigcup_{j \in J} V_j$ et
des recouvrements ouverts $f^{-1}(V_j) = \bigcup_{i \in I_j} U_i$ tels que chacun
des morphismes $U_i \to V_j$, $j\in J, i\in I_j$ est
étale.
\item Il existe un recouvrement ouvert affine $S = \bigcup_{j \in J} V_j$
et ouvert affine revêtements $f^{-1}(V_j) = \bigcup_{i \in I_j} U_i$ tel que l'homomorphisme
d'anneaux $\mathcal{O}_S(V_j) \to \mathcal{O}_X(U_i)$ est étale,
pour tout $j\in J, i\in I_j$.
\end{enumerate}
De plus, si $f$ est étale
alors pour tout sous-schémas ouverts $U \subset X$, $V \subset S$ avec $f(U) \subset V$
la restriction $f|_U : U \to V$ est étale.
\end{lemma}

\begin{proof}
Cela résulte du lemme \ref{morphisms-lemma-locally-P} si l'on montre que la propriété
``$R \to A$ est étale''
est locale. On vérifie les conditions (a),
(b) et (c) de la Définition
\ref{morphisms-definition-property-local}. Tout cela découle de l'algèbre, Lemme \ref{algebra-lemma-etale}.
\end{proof}

\begin{lemma}
\label{morphisms-lemma-composition-etale}
La composition de deux morphismes qui sont étale est étale.
\end{lemma}

\begin{proof}
Dans la démonstration du lemme \ref{morphisms-lemma-etale-characterize} nous avons vu
qu'être étale est une propriété locale des homomorphismes
d'anneaux. D'où la première affirmation du
lemme découle du lemme \ref{morphisms-lemma-composition-type-P} combinée au fait
qu'être étale est une propriété des homomorphismes d'anneaux
qui est stable sous composition,
voir Algèbre, Lemme \ref{algebra-lemma-etale}.
\end{proof}

\begin{lemma}
\label{morphisms-lemma-base-change-etale}
Le changement de base d'un morphisme qui est étale est étale.
\end{lemma}

\begin{proof}
Dans la démonstration du lemme \ref{morphisms-lemma-etale-characterize} nous avons
vu qu'être étale est une propriété locale des homomorphismes
d'anneaux. D'où le
lemme découle du lemme \ref{morphisms-lemma-composition-type-P} combiné au
fait qu'être étale est une propriété des homomorphismes d'anneaux
stable par changement
de base, voir Algèbre, Lemme \ref{algebra-lemma-etale}.
\end{proof}

\begin{lemma}
\label{morphisms-lemma-etale-smooth-unramified}
Soit $f : X \to S$ un morphisme de schémas. Soit $x \in X$.
Alors $f$ est étale à $x$ si et seulement si $f$ est
lisse et non ramifié à $x$.
\end{lemma}

\begin{proof}
Cela découle immédiatement des définitions.
\end{proof}

\begin{lemma}
\label{morphisms-lemma-etale-locally-quasi-finite}
Un morphisme étale est localement quasi-fini.
\end{lemma}

\begin{proof}
Par
le lemme \ref{morphisms-lemma-etale-smooth-unramified}, un morphisme
étale est non ramifié.
Par le Lemme \ref{morphisms-lemma-unramified-quasi-finite} un morphisme
non ramifié est localement quasi-fini.
\end{proof}

\begin{lemma}
\label{morphisms-lemma-etale-over-field}
\begin{slogan}
Description des schémas étales sur corps et fibres de morphismes
étales.
\end{slogan}
Fibres de morphismes étales.
\begin{enumerate}
\item Soit $X$ un schéma sur un corps $k$.
Le morphisme de structure $X \to \Spec(k)$ est étale si et seulement
si $X$ est une union disjointe de spectres d'extensions de corps finis
séparables de $k$.
\item Si $f : X \to S$ est un morphisme étale, alors pour tout $s \in S$ la fibre
$X_s$ est une union disjointe de spectres d'extensions de corps finis
séparables de $\kappa(s)$.
\end{enumerate}
\end{lemma}

\begin{proof}
Vous pouvez le déduire du lemme \ref{morphisms-lemma-unramified-over-field} via Lemme
\ref{morphisms-lemma-etale-smooth-unramified} ci-dessus. Voici une démonstration
directe.

\medskip\noindent
Nous utiliserons l'algèbre, Lemme \ref{algebra-lemma-etale-over-field}. Il est donc clair que
si $X$ est une union disjointe de spectres d'extensions de corps finis
séparables de $k$ alors $X \to \Spec(k)$ est étale. À
l'inverse, supposons que $X \to \Spec(k)$ soit étale. Alors pour tout $U \subset X$
ouvert affine, nous voyons que $U$ est une union disjointe
finie de spectres d'extensions de corps finis séparables de $k$.
Ainsi tous les points de
$X$ sont des points fermés (voir Lemme \ref{morphisms-lemma-algebraic-residue-field-extension-closed-point-fibre} par exemple).
Ainsi $X$ est un espace discret et on obtient le résultat.
\end{proof}

\noindent
Le lemme suivant caractérise un morphisme étale comme un morphisme plat
de présentation finie avec des ``fibres étales''.

\begin{lemma}
\label{morphisms-lemma-etale-flat-etale-fibres}
Soit $f : X \to S$ un morphisme de schémas. Si $f$
est plat, localement de présentation finie, et pour tout $s \in S$ la fibre $X_s$
est une union disjointe de spectres d'extensions finies de corps séparables
de $\kappa(s)$, alors $f$ est étale.
\end{lemma}

\begin{proof}
Vous pouvez le déduire
de l'algèbre, Lemme \ref{algebra-lemma-characterize-etale}. Voici une autre
démonstration.

\medskip\noindent
Par le Lemme \ref{morphisms-lemma-etale-over-field} une fibre $X_s$ est étale
et donc lisse sur $s$. Par le Lemme \ref{morphisms-lemma-smooth-flat-smooth-fibres} on
voit que $X \to S$ est lisse.
Par le Lemme \ref{morphisms-lemma-unramified-etale-fibres} nous voyons
que $f$ est non ramifié. Nous
concluons par le lemme \ref{morphisms-lemma-etale-smooth-unramified}.
\end{proof}

\begin{lemma}
\label{morphisms-lemma-open-immersion-etale}
Toute immersion ouverte est étale.
\end{lemma}

\begin{proof}
Ceci est vrai car une immersion ouverte est un isomorphisme local.
\end{proof}

\begin{lemma}
\label{morphisms-lemma-etale-syntomic}
Un morphisme étale est syntomique.
\end{lemma}

\begin{proof}
Voir Algèbre, Lemme \ref{algebra-lemma-smooth-syntomic} et utiliser qu'un morphisme
étale est le même qu'un morphisme lisse de dimension relative $0$.
\end{proof}

\begin{lemma}
\label{morphisms-lemma-etale-locally-finite-presentation}
Un morphisme étale est localement de présentation finie.
\end{lemma}

\begin{proof}
Vrai car un étale homomorphisme d'anneaux est de présentation finie par
définition.
\end{proof}

\begin{lemma}
\label{morphisms-lemma-etale-flat}
Un morphisme étale est plat.
\end{lemma}

\begin{proof}
Cela résulte des lemmes \ref{morphisms-lemma-syntomic-flat} et \ref{morphisms-lemma-etale-syntomic}.
\end{proof}

\begin{lemma}
\label{morphisms-lemma-etale-open}
Un morphisme étale est ouvert.
\end{lemma}

\begin{proof}
Cela résulte des lemmes \ref{morphisms-lemma-etale-flat},
\ref{morphisms-lemma-etale-locally-finite-presentation} et
\ref{morphisms-lemma-fppf-open}.
\end{proof}

\noindent
Le lemme suivant dit localement que tout morphisme étale est un étale standard.
Il s'agit en fait d'un résultat difficile à prouver en toute généralité. Les parties délicates
sont cachées dans le chapitre sur l'algèbre commutative. Ainsi, un morphisme
étale standard est un {\it modèle local} pour un morphisme
étale général.

\begin{lemma}
\label{morphisms-lemma-etale-locally-standard-etale}
Soit $f : X  \to S$ un morphisme de schémas. Soit $x \in X$
un point. Soit $V \subset S$
un voisinage ouvert affine de $f(x)$. Les conditions
suivantes sont équivalentes
\begin{enumerate}
\item Le morphisme $f$ est étale à $x$.
\item Il existe un $U \subset X$ ouvert affine avec
$x \in U$ et $f(U) \subset V$ tel que le morphisme
induit $f|_U : U \to V$ est standard étale
(voir Définition \ref{morphisms-definition-etale}).
\end{enumerate}
\end{lemma}

\begin{proof}
Dérivé des définitions
et de l'algèbre, proposition \ref{algebra-proposition-etale-locally-standard}.
\end{proof}

\noindent
Voici un critère différentiel de étalité en un point. Il existe
de nombreuses variantes de ce résultat, qui peuvent
toutes être utiles à un moment donné. Nous les ajouterons simplement ici
si nécessaire.

\begin{lemma}
\label{morphisms-lemma-etale-at-point}
Soit $f : X \to S$ un morphisme de schémas. Soit $x \in X$.
Posons $s = f(x)$.
Supposons que $f$
soit localement de présentation finie. Les conditions
suivantes sont équivalentes :
\begin{enumerate}
\item Le morphisme $f$ est étale à $x$.
\item L'application anneau local $\mathcal{O}_{S, s} \to \mathcal{O}_{X, x}$ est plate et
$X_s \to \Spec(\kappa(s))$ est étale à $x$.
\item L'homomorphisme local d'anneaux $\mathcal{O}_{S, s} \to \mathcal{O}_{X, x}$ est plate
et $X_s \to \Spec(\kappa(s))$ est non ramifiée en $x$.
\item L'homomorphisme d'anneaux locaux $\mathcal{O}_{S, s} \to \mathcal{O}_{X, x}$ est
plat et le $\mathcal{O}_{X, x}$-module $\Omega_{X/S, x}$ est
nul.
\item L'homomorphisme d'anneaux locaux $\mathcal{O}_{S, s} \to \mathcal{O}_{X, x}$ est plat et
l'espace vectoriel $\kappa(x)$-
$$
\Omega_{X_s/s, x} \otimes_{\mathcal{O}_{X_s, x}} \kappa(x) =
\Omega_{X/S, x} \otimes_{\mathcal{O}_{X, x}} \kappa(x)
$$
est nul.
\item L'homomorphisme local d'anneaux $\mathcal{O}_{S, s} \to \mathcal{O}_{X, x}$
est plat, nous avons $\mathfrak m_s\mathcal{O}_{X, x} = \mathfrak m_x$ et l'extension
corps $\kappa(x)/\kappa(s)$ est fini
séparable.
\item Il existe des ouverts affines $U \subset X$,
et $V \subset S$ telles que $x \in U$, $f(U) \subset V$ et le morphisme
induit $f|_U : U \to V$ est standard lisse de dimension
relative $0$.
\item Il existe des ouverts affines $\Spec(A) = U \subset X$
et $\Spec(R) = V \subset S$ avec $x \in U$ correspondant
à $\mathfrak q \subset A$, et $f(U) \subset V$ telles qu'il
existe une présentation
$$
A = R[x_1, \ldots, x_n]/(f_1, \ldots, f_n)
$$
avec
$$
g =
\det
\left(
\begin{matrix}
\partial f_1/\partial x_1 &
\partial f_2/\partial x_1 &
\ldots &
\partial f_n/\partial x_1 \\
\partial f_1/\partial x_2 &
\partial f_2/\partial x_2 &
\ldots &
\partial f_n/\partial x_2 \\
\ldots & \ldots & \ldots & \ldots \\
\partial f_1/\partial x_n &
\partial f_2/\partial x_n &
\ldots &
\partial f_n/\partial x_n
\end{matrix}
\right)
$$
mappant à un élément de $A$ pas dans $\mathfrak q$.
\item Il existe des ouverts affines $U \subset X$,
et $V \subset S$ telles que $x \in U$, $f(U) \subset V$ et le morphisme
induit $f|_U : U \to V$ est standard étale.
\item Il existe des ouverts affines $\Spec(A) = U \subset X$
et $\Spec(R) = V \subset S$ avec $x \in U$ correspondant
à $\mathfrak q \subset A$, et $f(U) \subset V$ telles qu'il
existe une présentation
$$
A = R[x]_Q/(P) = R[x, 1/Q]/(P)
$$
avec $P, Q \in R[x]$, $P$ monique et $P' = \text{d}P/\text{d}x$ mappage à un élément
de $A$ pas dans $\mathfrak q$.
\end{enumerate}
\end{lemma}

\begin{proof}
Utilisez Lemme \ref{morphisms-lemma-etale-locally-standard-etale} et les définitions
pour voir que (1) implique toutes les autres
conditions. Pour chacune des conditions (2)
-- (10), combinons les Lemmes \ref{morphisms-lemma-smooth-at-point} et \ref{morphisms-lemma-unramified-at-point}
pour voir que (1) est valable en montrant
que $f$ est à la fois lisse et non ramifié
à $x$ et en appliquant le Lemme \ref{morphisms-lemma-etale-smooth-unramified}. Certains
détails omis.
\end{proof}

\begin{lemma}
\label{morphisms-lemma-flat-unramified-etale}
Un morphisme est étale en un point si et seulement s'il est plat et G-non ramifié
en ce point. Un morphisme
est étale si et seulement si il est plat et G-non ramifié.
\end{lemma}

\begin{proof}
Cela ressort clairement des Lemmes \ref{morphisms-lemma-etale-at-point}
et \ref{morphisms-lemma-unramified-at-point}.
\end{proof}

\begin{lemma}
\label{morphisms-lemma-set-points-where-fibres-etale}
Soit
$$
\xymatrix{
X' \ar[r]_{g'} \ar[d]_{f'} & X \ar[d]^f \\
S' \ar[r]^g & S
}
$$
un diagramme cartésien de schémas. Soit $W \subset X$, resp.\ $W' \subset X'$ le sous-schéma
ouvert des points où $f$, resp.\ $f'$ est étale.
Alors $W' = (g')^{-1}(W)$ si
\begin{enumerate}
\item $f$ est plat et localement de présentation finie, ou
\item $f$ est localement de présentation finie et $g$ est plat.
\end{enumerate}
\end{lemma}

\begin{proof}
Supposons d'abord que $f$ localement de type fini. Considérons l'ensemble
$$
T = \{x \in X \mid f\text{ est non ramifié en }x\}
$$
et l'ensemble correspondant $T' \subset X'$ pour
$f'$. Puis
$T' = (g')^{-1}(T)$ par le lemme \ref{morphisms-lemma-set-points-where-fibres-unramified}.

\medskip\noindent Ainsi
le cas (1) suit car dans le cas (1) $T$ est l'ensemble (ouvert) de
points où $f$ est étale par le lemme \ref{morphisms-lemma-flat-unramified-etale}.

\medskip\noindent
Dans le cas (2), soit $x' \in W'$. Alors $g'$
est plat en $x'$ (Lemme \ref{morphisms-lemma-base-change-module-flat})
et $g \circ f'$ est plat en $x'$ (Lemme \ref{morphisms-lemma-composition-module-flat}). Il s'ensuit
que $f$ est plat à $x = g'(x')$
par le lemme \ref{morphisms-lemma-flat-permanence}. D'autre part, puisque $x' \in T'$
(Lemme \ref{morphisms-lemma-base-change-smooth}) on voit que $x \in T$.
Ainsi $f$ est étale à $x$
par le lemme \ref{morphisms-lemma-etale-at-point}.
\end{proof}

\begin{lemma}
\label{morphisms-lemma-etale-permanence-better}
\begin{reference}
Voir par exemple \cite[Remarque 18.30 (2).]{GWII}.
\end{reference}
Soit $f : X \to Y$ un morphisme de schémas sur $S$. Si $X$ est étale
sur $S$ et que $Y$ est non ramifié sur $S$, alors $f$
est étale.
\end{lemma}

\begin{proof}
Considérons la factorisation $X \to X \times_S Y \to Y$, où la première flèche est donnée par
$\text{id}_X$ et $f$ et la deuxième flèche est la projection. Nous affirmons
que les deux flèches sont étale et donc $f$
est étale du lemme \ref{morphisms-lemma-composition-etale}. À savoir, la projection est
étale car c'est le changement de base de $X \to S$, voir Lemme \ref{morphisms-lemma-base-change-etale}.
La première flèche est le changement de base du morphisme
diagonal $Y \to Y \times_S Y$ car le carré
$$
\xymatrix{
X \ar[d] \ar[r] & X \times_S Y \ar[d] \\
Y \ar[r] & Y \times_S Y
}
$$
est cartésien. La diagonale $Y \to Y \times_S Y$ est une immersion ouverte
du lemme \ref{morphisms-lemma-diagonal-unramified-morphism}. Le changement de base d'une
immersion ouverte est une immersion
ouverte (Schémas, Lemme \ref{schemes-lemma-base-change-immersion}) et une
immersion ouverte est étale
(Lemme \ref{morphisms-lemma-open-immersion-etale}).
\end{proof}

\begin{lemma}
\label{morphisms-lemma-etale-permanence}
\begin{slogan}
Loi d'annulation pour les morphismes étales
\end{slogan}
Soit $f : X \to Y$ un morphisme de schémas sur $S$. Si $X$ et $Y$
sont étales sur $S$, alors $f$
est étale.
\end{lemma}

\begin{proof}
Puisque $Y \to S$ est non ramifié d'après le lemme \ref{morphisms-lemma-etale-smooth-unramified},
nous pouvons appliquer le Lemme \ref{morphisms-lemma-etale-permanence-better}.
Voir aussi Algèbre, Lemme \ref{algebra-lemma-map-between-etale}.
\end{proof}

\begin{lemma}
\label{morphisms-lemma-etale-permanence-two}
Soit
$$
\xymatrix{
X \ar[rr]_f \ar[rd]_p & &
Y \ar[dl]^q \\
& S
}
$$
un diagramme commutatif de morphismes de schémas. Supposons que
\begin{enumerate}
\item $f$ est surjectif et étale,
\item $p$ est étale, et
\item $q$ est localement de présentation finie\footnote{En
fait cela
est impliqué par (1) et (2), voir Descente, Lemme \ref{descent-lemma-flat-finitely-presented-permanence}. De plus,
il suffit de supposer que $f$ est surjectif, plat
et localement de présentation
finie, voir Descente, Lemme \ref{descent-lemma-smooth-permanence}.}.
\end{enumerate}
Alors $q$ est étale.
\end{lemma}

\begin{proof}
Par le Lemme \ref{morphisms-lemma-smooth-permanence} nous voyons que $q$ est lisse. Il suffit
donc de voir que $q$ a une dimension relative
$0$. Cela résulte du lemme \ref{morphisms-lemma-dimension-fibre-at-a-point-additive} et du fait que $f$
et $p$ ont une dimension relative $0$.
\end{proof}

\noindent
Une dernière caractérisation des morphismes lisses est qu'un morphisme lisse
$f : X \to S$ est localement la composition d'un morphisme étale par une
projection $\mathbf{A}_S^d \to S$.

\begin{lemma}
\label{morphisms-lemma-smooth-etale-over-affine-space}
\begin{slogan}
Les schémas lisses sont des espaces affines de type étalement-local.
\end{slogan}
Soit $\varphi : X \to Y$ un morphisme de schémas. Soit $x \in X$. Soit
$V \subset Y$ un voisinage ouvert affine de $\varphi(x)$. Si $\varphi$
est lisse à $x$, alors il existe un entier $d \geq 0$
et un ouvert affine $U \subset X$ avec $x \in U$
et $\varphi(U) \subset V$ tel qu'il existe un diagramme commutatif
$$
\xymatrix{
X \ar[d] & U \ar[l] \ar[d] \ar[r]_-\pi & \mathbf{A}^d_V \ar[ld] \\
Y & V \ar[l]
}
$$
où $\pi$ est étale.
\end{lemma}

\begin{proof}
Par
Lemme \ref{morphisms-lemma-smooth-locally-standard-smooth} on peut trouver un affine
ouvert $U$ comme dans le lemme tel que $\varphi|_U : U \to V$
est lisse standard. Écrivez $U = \Spec(A)$
et $V = \Spec(R)$ afin que nous puissions écrire
$$
A = R[x_1, \ldots, x_n]/(f_1, \ldots, f_c)
$$
avec le mappage
$$
g =
\det
\left(
\begin{matrix}
\partial f_1/\partial x_1 &
\partial f_2/\partial x_1 &
\ldots &
\partial f_c/\partial x_1 \\
\partial f_1/\partial x_2 &
\partial f_2/\partial x_2 &
\ldots &
\partial f_c/\partial x_2 \\
\ldots & \ldots & \ldots & \ldots \\
\partial f_1/\partial x_c &
\partial f_2/\partial x_c &
\ldots &
\partial f_c/\partial x_c
\end{matrix}
\right)
$$
sur un élément inversible de $A$. Il est alors
clair que $R[x_{c + 1}, \ldots, x_n] \to A$ est un lissage standard de dimension
relative $0$. Il est donc lisse de dimension
relative $0$. Autrement dit l'homomorphisme
d'anneaux $R[x_{c + 1}, \ldots, x_n] \to A$ est étale. Comme $\mathbf{A}^{n - c}_V = \Spec(R[x_{c + 1}, \ldots, x_n])$ le
lemme avec $d = n - c$.
\end{proof}
















\section{Faisceaux relativement amples}
\label{morphisms-section-relatively-ample}

\noindent
Soit $X$ un schéma et $\mathcal{L}$ un faisceau inversible sur
$X$. Alors $\mathcal{L}$ est suffisant sur $X$ si $X$
est quasi-compact et que chaque point de $X$ est
contenu dans un ouvert affine de la forme $X_s$, où $s \in \Gamma(X, \mathcal{L}^{\otimes n})$ et $n \geq 1$,
voir Propriétés, Définition \ref{properties-definition-ample}. Nous transformons
cela en une notion relative comme suit.

\begin{definition}
\label{morphisms-definition-relatively-ample}
\begin{reference}
\cite[II Définition 4.6.1]{EGA}
\end{reference}
Soit $f : X \to S$ un morphisme de schémas. Soit $\mathcal{L}$ un
module $\mathcal{O}_X$ inversible. On dit que $\mathcal{L}$ est {\it
relativement ample}, ou {\it $f$-relativement ample}, ou {\it
amplement sur $X/S$}, ou {\it $f$-ample}
si $f : X \to S$ est quasi-compact, et si pour tout ouvert
affine $V \subset S$ la restriction de $\mathcal{L}$ au sous-schéma ouvert $f^{-1}(V)$
de $X$ est suffisant.
\end{definition}

\noindent
On remarque que l'existence d'un faisceau relativement ample sur $X$ ne
force pas le morphisme $X \to S$ à être de type fini.

\begin{lemma}
\label{morphisms-lemma-ample-power-ample}
Soit $X \to S$ un morphisme de schémas.
Soit $\mathcal{L}$ un module $\mathcal{O}_X$ inversible. Soit $n \geq 1$.
Alors $\mathcal{L}$ est $f$-ample si et seulement
si $\mathcal{L}^{\otimes n}$ est $f$-ample.
\end{lemma}

\begin{proof}
Cela résulte de Propriétés, Lemme \ref{properties-lemma-ample-power-ample}.
\end{proof}

\begin{lemma}
\label{morphisms-lemma-relatively-ample-separated}
Soit $f : X \to S$ un morphisme de schémas. S'il existe
un $f$-faisceau ample inversible, alors $f$
est séparé.
\end{lemma}

\begin{proof}
Être séparé est local sur la base (voir
Schémas, Lemme \ref{schemes-lemma-characterize-separated} par exemple ; cela découle aussi facilement
de la définition). Nous pouvons donc
supposer que $S$ est affine et
que $X$ a un faisceau ample inversible.
Dans ce cas,
le résultat découle de Propriétés, Lemme \ref{properties-lemma-ample-separated}.
\end{proof}

\noindent
Il existe de nombreuses façons de caractériser des faisceaux inversibles
relativement amples, analogues aux conditions
équivalentes dans Propriétés, Proposition \ref{properties-proposition-characterize-ample}. Nous les
ajouterons ici si nécessaire.

\begin{lemma}
\label{morphisms-lemma-characterize-relatively-ample}
\begin{reference}
\cite[II, Proposition 4.6.3]{EGA}
\end{reference}
Soit $f : X \to S$ un morphisme de schémas quasi-compact. Soit $\mathcal{L}$
un faisceau inversible sur $X$. Les conditions
suivantes sont équivalentes :
\begin{enumerate}
\item Le faisceau inversible $\mathcal{L}$ est $f$-ample.
\item Il existe un recouvrement ouvert $S = \bigcup V_i$
tel que chaque $\mathcal{L}|_{f^{-1}(V_i)}$ soit ample
par rapport à $f^{-1}(V_i) \to V_i$.
\item Il existe un recouvrement ouvert affine $S = \bigcup V_i$
tel que chaque $\mathcal{L}|_{f^{-1}(V_i)}$ soit ample.
\item Il existe une $\mathcal{O}_S$-algèbre graduée
$\mathcal{A}$ et une application de $\mathcal{O}_X$-algèbres
graduées $\psi : f^*\mathcal{A} \to \bigoplus_{d \geq 0} \mathcal{L}^{\otimes d}$ telle que
$U(\psi) = X$ et
$$
r_{\mathcal{L}, \psi} :
X
\longrightarrow
\underline{\text{Proj}}_S(\mathcal{A})
$$
est une immersion ouverte (voir
Constructions, Lemme \ref{constructions-lemma-invertible-map-into-relative-proj} pour la notation).
\item Le morphisme $f$ est quasi-séparé
et la partie
(4) ci-dessus tient avec $\mathcal{A} = f_*(\bigoplus_{d \geq 0} \mathcal{L}^{\otimes d})$ et $\psi$
le mappage d'adjonction.
\item Identique à (4) mais exigeant simplement que $r_{\mathcal{L}, \psi}$ soit
une immersion.
\end{enumerate}
\end{lemma}

\begin{proof}
Il ressort immédiatement de la définition que (1) implique (2) et
(2) implique (3). Il est clair que (5) implique (4).

\medskip\noindent
Supposons que (3) soit valable pour le recouvrement
ouvert affine $S = \bigcup V_i$.
Nous allons montrer (5) prises. Puisque
chaque $f^{-1}(V_i)$ possède
un faisceau ample inversible on voit que $f^{-1}(V_i)$
est séparé
(Propriétés, Lemme \ref{properties-lemma-ample-separated}). Par conséquent,
$f$ est séparé. Par Schémas, Lemme \ref{schemes-lemma-push-forward-quasi-coherent} nous voyons
que $\mathcal{A} = f_*(\bigoplus_{d \geq 0} \mathcal{L}^{\otimes d})$ est une $\mathcal{O}_S$-algèbre
graduée quasi-cohérente. Désignons $\psi : f^*\mathcal{A} \to \bigoplus_{d \geq 0} \mathcal{L}^{\otimes d}$ le mappage
d'adjonction.
La description du $U(\psi)$
ouvert dans Constructions,
section \ref{constructions-section-invertible-relative-proj} et
la définition de l'ampleur
de $\mathcal{L}|_{f^{-1}(V_i)}$ montrent que $U(\psi) = X$.
De plus,
Constructions, Lemme \ref{constructions-lemma-invertible-map-into-relative-proj} partie (3) montre
que la restriction de $r_{\mathcal{L}, \psi}$ à
$f^{-1}(V_i)$ est la même que le morphisme
de Propriétés, Lemme \ref{properties-lemma-map-into-proj}
qui est une immersion
ouverte d'après Propriétés, Lemme \ref{properties-lemma-ample-immersion-into-proj}.
Donc (5) est valable.

\medskip\noindent
Montrons que (4) implique (1).
Supposons (4). Notons $\pi : \underline{\text{Proj}}_S(\mathcal{A}) \to S$ le morphisme
de structure. Choisissez $V \subset S$ ouvert
affine. Par Constructions, Définition \ref{constructions-definition-relative-proj}
on voit que $\pi^{-1}(V) \subset \underline{\text{Proj}}_S(\mathcal{A})$ est égal à $\text{Proj}(A)$ où
$A = \mathcal{A}(V)$ comme anneau gradué. Ainsi $r_{\mathcal{L}, \psi}$
applique $f^{-1}(V)$ de manière
isomorphe sur un ouvert
quasi-compact de $\text{Proj}(A)$.
De plus, $\mathcal{L}^{\otimes d}$ est isomorphe au retrait de
$\mathcal{O}_{\text{Proj}(A)}(d)$ pour certains $d \geq 1$.
(Voir la partie (3)
de Constructions, Lemme \ref{constructions-lemma-invertible-map-into-relative-proj} et
l'énoncé final de Constructions,
Lemme \ref{constructions-lemma-invertible-map-into-proj}.) Cela implique
que $\mathcal{L}|_{f^{-1}(V)}$ est amplement
rempli par les Propriétés, Lemmes
\ref{properties-lemma-open-in-proj-ample} et \ref{properties-lemma-ample-power-ample}.

\medskip\noindent
Supposons (6). Par l'équivalence de (1) - (5) ci-dessus on
voit que la propriété d'être relativement ample sur $X/S$
est locale sur $S$. On peut donc supposer que
$S$ est affine, et il faut montrer que $\mathcal{L}$
est ample sur $X$. Dans ce cas le morphisme $r_{\mathcal{L}, \psi}$
est identifié au morphisme, également noté $r_{\mathcal{L}, \psi} : X \to \text{Proj}(A)$ associé
à l'application $\psi : A = \mathcal{A}(V) \to \Gamma_*(X, \mathcal{L})$. (Voir les références
ci-dessus.) Comme ci-dessus, nous voyons également
que $\mathcal{L}^{\otimes d}$ est le retrait du faisceau $\mathcal{O}_{\text{Proj}(A)}(d)$
pour certains $d \geq 1$. De plus, puisque
$X$ est quasi-compact nous
voyons que $X$ s'identifie à un sous-schéma fermé
d'un sous-schéma ouvert quasi-compact $Y \subset \text{Proj}(A)$.
Par
Constructions,
Lemme \ref{constructions-lemma-ample-on-proj} (voir aussi Propriétés,
Lemme
\ref{properties-lemma-open-in-proj-ample})
on voit que $\mathcal{O}_Y(d')$ est un faisceau
ample inversible sur $Y$ pour certains $d' \geq 1$.
Puisque la restriction d'un faisceau suffisant à un
sous-schéma fermé est ample, voir Propriétés,
Lemme
\ref{properties-lemma-ample-on-closed} Nous concluons que le
retrait de $\mathcal{O}_Y(d')$ est
ample. En combinant ces résultats avec Propriétés,
Lemme
\ref{properties-lemma-ample-power-ample} Nous concluons que $\mathcal{L}$
est suffisant comme souhaité.
\end{proof}

\begin{lemma}
\label{morphisms-lemma-ample-over-affine}
\begin{reference}
\cite[II Corollaire 4.6.6]{EGA}
\end{reference}
Soit $f : X \to S$ un morphisme de schémas. Soit $\mathcal{L}$
un module $\mathcal{O}_X$ inversible. Supposons que $S$ soit
affine. Alors $\mathcal{L}$
est $f$-relativement amplement si et seulement si $\mathcal{L}$
est amplement sur $X$.
\end{lemma}

\begin{proof}
Immédiat du lemme \ref{morphisms-lemma-characterize-relatively-ample} et des
définitions.
\end{proof}

\begin{lemma}
\label{morphisms-lemma-quasi-affine-O-ample}
\begin{reference}
\cite[II Proposition 5.1.6]{EGA}
\end{reference}
Soit $f : X \to S$ un morphisme de schémas. Alors $f$ est quasi-affine si
et seulement si $\mathcal{O}_X$ est $f$-relativement ample.
\end{lemma}

\begin{proof}
Cela résulte de Propriétés, Lemme \ref{properties-lemma-quasi-affine-O-ample} et des
définitions.
\end{proof}

\begin{lemma}
\label{morphisms-lemma-pullback-ample-tensor-relatively-ample}
Soit $f : X \to Y$ un morphisme de schémas, $\mathcal{M}$
un $\mathcal{O}_Y$-module inversible, et $\mathcal{L}$ un
$\mathcal{O}_X$-module inversible.
\begin{enumerate}
\item Si $\mathcal{L}$ est suffisant pour $f$
et que $\mathcal{M}$ est suffisant, alors $\mathcal{L} \otimes f^*\mathcal{M}^{\otimes a}$ est suffisant pour
$a \gg 0$.
\item Si $\mathcal{M}$ est ample et
$f$ quasi-affine, alors $f^*\mathcal{M}$ est ample.
\end{enumerate}
\end{lemma}

\begin{proof}
Supposons que $\mathcal{L}$ est $f$-ample et $\mathcal{M}$
ample. Par hypothèse $Y$
et $f$ sont quasi-compacts (voir Définition
\ref{morphisms-definition-relatively-ample} et Propriétés, Définition \ref{properties-definition-ample}).
Donc $X$ est quasi-compact.
Par Propriétés, Lemme \ref{properties-lemma-ample-separated} le schéma $Y$
est séparé et par le lemme \ref{morphisms-lemma-relatively-ample-separated} le morphisme
$f$ est séparé. Ainsi $X$ est
séparé par Schémas, Lemme \ref{schemes-lemma-separated-permanence}. Choisissez
$x \in X$. On peut choisir $m \geq 1$
et $t \in \Gamma(Y, \mathcal{M}^{\otimes m})$ tels que $Y_t$ soit affine et
$f(x) \in Y_t$. Puisque $\mathcal{L}$ se limite à un faisceau
amplement inversible sur $f^{-1}(Y_t) = X_{f^*t}$, nous
pouvons choisir $n \geq 1$ et $s \in \Gamma(X_{f^*t}, \mathcal{L}^{\otimes n})$ avec $x \in (X_{f^*t})_s$ avec
$(X_{f^*t})_s$ affines. Par Propriétés, Lemme
\ref{properties-lemma-invert-s-sections} partie (2) dont les hypothèses sont satisfaites
par ce qui précède, il existe
un entier $e \geq 1$ et une section
$s' \in \Gamma(X, \mathcal{L}^{\otimes n} \otimes f^*\mathcal{M}^{\otimes em})$ qui se restreint à $s(f^*t)^e$ sur $X_{f^*t}$. Pour
tout $b > 0$, considérez la section $s'' = s'(f^*t)^b$
de $\mathcal{L}^{\otimes n} \otimes f^*\mathcal{M}^{\otimes (e + b)m}$. Alors $X_{s''} = (X_{f^*t})_s$
est un ouvert affine de $X$ contenant $x$.
En choisissant $b$ de telle sorte que $n$
divise $e + b$, nous voyons
$\mathcal{L}^{\otimes n} \otimes f^*\mathcal{M}^{\otimes (e + b)m}$ est la $n$ième puissance de $\mathcal{L} \otimes f^*\mathcal{M}^{\otimes a}$
pour certains $a$ et nous pouvons obtenir n'importe quel
$a$ divisible par $m$ et assez grand. Puisque
$X$ est quasi-compact, un nombre fini de
ces ouverts affines couvrent $X$. Nous concluons
que pour certains $a$
suffisamment divisibles et assez grands le faisceau inversible
$\mathcal{L} \otimes f^*\mathcal{M}^{\otimes a}$ est amplement suffisant sur $X$.
D'un autre côté, nous savons que $\mathcal{M}^{\otimes c}$ (et donc
son retrait vers $X$) est généré globalement pour tous
les $c \gg 0$ par Propriétés, Proposition \ref{properties-proposition-characterize-ample}.
Ainsi $\mathcal{L} \otimes f^*\mathcal{M}^{\otimes a + c}$ est suffisant (Propriétés, Lemme \ref{properties-lemma-ample-tensor-globally-generated})
pour $c \gg 0$ et (1) est prouvé.

\medskip\noindent
La partie (2) découle du lemme \ref{morphisms-lemma-quasi-affine-O-ample},
Propriétés, Lemme \ref{properties-lemma-ample-power-ample} et de la partie
(1).
\end{proof}

\begin{lemma}
\label{morphisms-lemma-ample-composition}
Soient $g : Y \to S$ et $f : X \to Y$ des morphismes de schémas.
Soit $\mathcal{M}$ un module $\mathcal{O}_Y$ inversible.
Soit $\mathcal{L}$ un module $\mathcal{O}_X$ inversible. Si
$S$ est quasi-compact, $\mathcal{M}$ est $g$-ample
et $\mathcal{L}$ est
$f$-ample, alors $\mathcal{L} \otimes f^*\mathcal{M}^{\otimes a}$ est $g \circ f$-ample
pour $a \gg 0$.
\end{lemma}

\begin{proof}
Soit $S = \bigcup_{i = 1, \ldots, n} V_i$ un recouvrement ouvert affine
fini. Par le Lemme \ref{morphisms-lemma-characterize-relatively-ample}
il suffit de
prouver que $\mathcal{L} \otimes f^*\mathcal{M}^{\otimes a}$ suffit
sur $(g \circ f)^{-1}(V_i)$ pour $i = 1, \ldots, n$.
Ainsi le lemme
découle du lemme \ref{morphisms-lemma-pullback-ample-tensor-relatively-ample}.
\end{proof}

\begin{lemma}
\label{morphisms-lemma-ample-base-change}
Soit $f : X \to S$ un morphisme de schémas. Soit
$\mathcal{L}$ un module $\mathcal{O}_X$ inversible. Soit $S' \to S$
un morphisme de schémas. Soit $f' : X' \to S'$
le changement de base de $f$ et notons $\mathcal{L}'$
le retrait de $\mathcal{L}$ vers $X'$. Si $\mathcal{L}$
est $f$-ample, alors $\mathcal{L}'$ est $f'$-ample.
\end{lemma}

\begin{proof}
Par le Lemme \ref{morphisms-lemma-characterize-relatively-ample} il suffit de trouver un recouvrement
ouvert affine $S' = \bigcup U'_i$ tel que $\mathcal{L}'$
se restreint à un faisceau ample inversible
sur $(f')^{-1}(U_i')$ pour tout $i$. Nous pouvons
choisir le mappage $U'_i$ dans un $U_i \subset S$
ouvert affine. Dans ce cas le morphisme $(f')^{-1}(U'_i) \to f^{-1}(U_i)$
est affine comme un changement de base
du morphisme affine $U'_i \to U_i$ (Lemme
\ref{morphisms-lemma-base-change-affine}). Ainsi $\mathcal{L}'|_{(f')^{-1}(U'_i)}$ est
amplement suffisant par le lemme \ref{morphisms-lemma-pullback-ample-tensor-relatively-ample}.
\end{proof}

\begin{lemma}
\label{morphisms-lemma-ample-permanence}
Soient $g : Y \to S$ et $f : X \to Y$ des morphismes de schémas. Soit
$\mathcal{L}$ un module $\mathcal{O}_X$ inversible. Si $\mathcal{L}$
est $g \circ f$-ample et $f$ est quasi-compact\footnote{Cela
s'ensuit si $g$ est quasi-séparé
par Schémas, Lemme \ref{schemes-lemma-quasi-compact-permanence}.} alors $\mathcal{L}$
est $f$-ample.
\end{lemma}

\begin{proof}
Supposons que $f$ soit quasi-compact et que $\mathcal{L}$
soit $g \circ f$-ample. Soit $U \subset S$
être un ouvert affine et soit
$V \subset Y$ être un ouvert affine avec
$g(V) \subset U$. Alors $\mathcal{L}|_{(g \circ f)^{-1}(U)}$
est suffisant sur $(g \circ f)^{-1}(U)$ par hypothèse.
Puisque $f^{-1}(V) \subset (g \circ f)^{-1}(U)$ on voit que $\mathcal{L}|_{f^{-1}(V)}$
est suffisant sur $f^{-1}(V)$ par Propriétés, Lemme
\ref{properties-lemma-ample-on-locally-closed}. À savoir, $f^{-1}(V) \to (g \circ f)^{-1}(U)$ est une immersion
ouverte quasi-compacte
par Schémas, Lemme \ref{schemes-lemma-quasi-compact-permanence} puisque
$(g \circ f)^{-1}(U)$ est séparé
(Propriétés, Lemme \ref{properties-lemma-ample-separated}) et $f^{-1}(V)$ est
quasi-compact (car $f$ est quasi-compact).
Ainsi Nous concluons que $\mathcal{L}$ est $f$-ample
par le lemme \ref{morphisms-lemma-characterize-relatively-ample}.
\end{proof}







\section{Faisceaux très amples}
\label{morphisms-section-very-ample}

\noindent
Rappelons que étant donné un faisceau quasi-cohérent
$\mathcal{E}$ sur un schéma $S$ le fibré projectif
{\it }
associé à $\mathcal{E}$ est le morphisme $\mathbf{P}(\mathcal{E}) \to S$, où
$\mathbf{P}(\mathcal{E}) = \underline{\text{Proj}}_S(\text{Sym}(\mathcal{E}))$,
voir Constructions, Définition \ref{constructions-definition-projective-bundle}.

\begin{definition}
\label{morphisms-definition-very-ample}
Soit $f : X \to S$ un morphisme de schémas. Soit $\mathcal{L}$
un module $\mathcal{O}_X$ inversible. On dit que $\mathcal{L}$
est {\it relativement très ample} ou plus
précisément {\it $f$-relativement
très ample}, ou encore {\it
très ample sur $X/S$}, ou {\it
$f$-très ample} s'il existe un $\mathcal{O}_S$-module
$\mathcal{E}$ quasi-cohérent
et une immersion $i : X \to \mathbf{P}(\mathcal{E})$ sur $S$ tel que $\mathcal{L} \cong i^*\mathcal{O}_{\mathbf{P}(\mathcal{E})}(1)$.
\end{definition}



\noindent
Puisqu'il n'y a pas d'hypothèse de quasi-compacité dans cette définition, il n'est pas vrai en général
qu'un ample inversible à faisceau relativement très large soit un ample inversible à faisceau
relativement important.

\begin{lemma}
\label{morphisms-lemma-ample-very-ample}
\begin{reference}
\cite[II, Proposition 4.6.2]{EGA}
\end{reference}
Soit $f : X \to S$ un morphisme de schémas. Soit
$\mathcal{L}$ un module $\mathcal{O}_X$ inversible. Si $f$
est quasi-compact et $\mathcal{L}$ est un ample inversible
relativement très gros, alors $\mathcal{L}$ est un ample inversible
relativement gros.
\end{lemma}

\begin{proof}
Par définition il existe un $\mathcal{O}_S$-module
$\mathcal{E}$ quasi-cohérent et une immersion
$i : X \to \mathbf{P}(\mathcal{E})$
sur $S$ telle que $\mathcal{L} \cong i^*\mathcal{O}_{\mathbf{P}(\mathcal{E})}(1)$.
Posons $\mathcal{A} = \text{Sym}(\mathcal{E})$, donc
$\mathbf{P}(\mathcal{E}) = \underline{\text{Proj}}_S(\mathcal{A})$ par définition. L'algèbre $\mathcal{O}_S$
graduée $\mathcal{A}$ est munie d'un
homomorphisme
$$
\psi :
\mathcal{A} \to
\bigoplus\nolimits_{n \geq 0}
\pi_*\mathcal{O}_{\mathbf{P}(\mathcal{E})}(n) \to
\bigoplus\nolimits_{n \geq 0}
f_*\mathcal{L}^{\otimes n}
$$
où la deuxième flèche utilise l'identification
$\mathcal{L} \cong i^*\mathcal{O}_{\mathbf{P}(\mathcal{E})}(1)$. Par adjointité de $f_*$
et $f^*$ on obtient un morphisme
$\psi : f^*\mathcal{A} \to \bigoplus_{n \geq 0}\mathcal{L}^{\otimes n}$. Nous omettons de vérifier que le
morphisme $r_{\mathcal{L}, \psi}$ associé à cette application
est exactement l'immersion $i$.
Le résultat découle
donc de la partie (6) du lemme \ref{morphisms-lemma-characterize-relatively-ample}.
\end{proof}

\noindent
Pour formuler une réciproque correcte de ce lemme, demandons-nous
si, étant donné un faisceau inversible relativement ample
$\mathcal{L}$, il existe un entier $n \geq 1$
tel que $\mathcal{L}^{\otimes n}$ soit relativement très ample ? En général, c'est
faux. Il y a plusieurs choses qui empêchent que cela soit vrai :
\begin{enumerate}
\item Même si $S$ est affine, il peut arriver qu'aucun entier
fini $n$ ne fonctionne car $X \to S$ n'est
pas de type fini, voir Exemple \ref{morphisms-example-not-finite-type-proj}.
\item Si la base n'est pas quasi-compacte, le résultat peut être faux
même lorsque $f$ est de type fini. En effet, étant
donné un corps $k$ il existe un schéma $X_d$ de type fini
sur $k$ avec un faisceau inversible ample $\mathcal{O}_{X_d}(1)$ de sorte que
la plus petite puissance tensorielle de $\mathcal{O}_{X_d}(1)$ qui est très ample soit la
$d$ième puissance. Voir Exemple \ref{morphisms-example-not-bounded}.
Prendre pour $f$ la réunion disjointe des schémas $X_d$, envoyée
sur la réunion disjointe de copies de $\Spec(k)$, donne un exemple.
\end{enumerate}
Pour voir notre version de l'inverse, jetez un œil
dans le lemme \ref{morphisms-lemma-finite-type-ample-very-ample} ci-dessous. Nous ferons quelques
travaux préliminaires avant de le prouver.

\begin{example}
\label{morphisms-example-very-ample}
Soit $S$ un schéma. Soit $\mathcal{A}$ une algèbre $\mathcal{O}_S$
graduée quasi-cohérente générée par $\mathcal{A}_1$ sur $\mathcal{A}_0$.
Posons $X = \underline{\text{Proj}}_S(\mathcal{A})$. Dans ce cas $\mathcal{O}_X(1)$ est inversible
et très ample sur $X/S$. À savoir, le morphisme
associé à l'homomorphisme d'algèbres graduées $\mathcal{O}_S$
$$
\text{Sym}_{\mathcal{O}_X}^*(\mathcal{A}_1)
\longrightarrow
\mathcal{A}
$$
est une immersion fermée $X \to \mathbf{P}(\mathcal{A}_1)$
qui ramène $\mathcal{O}_{\mathbf{P}(\mathcal{A}_1)}(1)$ à $\mathcal{O}_X(1)$,
voir
Constructions, Lemme \ref{constructions-lemma-surjective-generated-degree-1-map-relative-proj}.
\end{example}

\begin{example}
\label{morphisms-example-not-finite-type-proj}
Soit $k$ un corps.
Considérons l'algèbre $k$ graduée
$$
A = k[U, V, Z_1, Z_2, Z_3, \ldots]/I
\quad
\text{avec}
\quad
I = (U^2 - Z_1^2, U^4 - Z_2^2, U^6 - Z_3^2, \ldots)
$$
avec la graduation définie par $\deg(U) = \deg(V) = \deg(Z_1) = 1$ et $\deg(Z_d) = d$.
Notez que $X = \text{Proj}(A)$
est couvert par $D_{+}(U)$ et $D_{+}(V)$. Ainsi les faisceaux
$\mathcal{O}_X(n)$ sont tous inversibles et isomorphes
à $\mathcal{O}_X(1)^{\otimes n}$. En particulier $\mathcal{O}_X(1)$ est ample et
$f$-ample pour le morphisme $f : X \to \Spec(k)$. Nous
affirmons qu'aucune puissance de $\mathcal{O}_X(1)$
n'est $f$-relativement très ample. En effet, il est facile
de voir que $\Gamma(X, \mathcal{O}_X(n))$ est la composante homogène de degré $n$
de l'algèbre $A$. Donc si $\mathcal{O}_X(n)$ était très ample, alors
$X$ serait un sous-schéma fermé d'un espace projectif sur
$k$ et donc de type fini sur $k$. Par contre
$D_{+}(V)$ est le
spectre de $k[t, t_1, t_2, \ldots]/(t^2 - t_1^2, t^4 - t_2^2, t^6 - t_3^2, \ldots)$ qui n'est pas de type
fini sur $k$.
\end{example}

\begin{example}
\label{morphisms-example-not-bounded}
Soit $k$ un corps infini. Soit $\lambda_1, \lambda_2, \lambda_3, \ldots$ des éléments distincts par paires
de $k^*$. (Ce n'est pas strictement nécessaire, et en fait l'exemple
fonctionne parfaitement même si tous les $\lambda_i$ sont égaux à $1$.)
Considérons
l'algèbre $k$ graduée
$$
A_d = k[U, V, Z]/I_d
\quad
\text{avec}
\quad
I_d = (Z^2 - \prod\nolimits_{i = 1}^{2d} (U - \lambda_i V)).
$$
avec la graduation définie par $\deg(U) = \deg(V) = 1$ et $\deg(Z) = d$. Ensuite $X_d = \text{Proj}(A_d)$ a un faisceau inversible ample
$\mathcal{O}_{X_d}(1)$. Nous affirmons que si $\mathcal{O}_{X_d}(n)$ est
très ample, alors $n \geq d$. La raison en est que $Z$
a le diplôme $d$, et donc $\Gamma(X_d, \mathcal{O}_{X_d}(n)) =
k[U, V]_n$
pour $n < d$. Détails omis.
\end{example}

\begin{lemma}
\label{morphisms-lemma-relatively-very-ample-separated}
Soit $f : X \to S$ un morphisme de schémas. Soit $\mathcal{L}$
un faisceau inversible sur $X$. Si $\mathcal{L}$
est relativement très ample sur $X/S$ alors $f$
est séparé.
\end{lemma}

\begin{proof}
La séparation est locale sur la base
(voir Schémas, section \ref{schemes-section-separation-axioms}). Une immersion
est séparée
(voir Schémas, Lemme \ref{schemes-lemma-immersions-monomorphisms}). D'où le lemme s'ensuit
puisque localement $X$ a une immersion
dans le spectre homogène d'un anneau gradué qui est
séparé, voir Constructions, Lemme \ref{constructions-lemma-proj-separated}.
\end{proof}

\begin{lemma}
\label{morphisms-lemma-relatively-very-ample}
Soit $f : X \to S$ un morphisme de schémas. Soit $\mathcal{L}$
un faisceau inversible sur $X$. Supposons que $f$
soit quasi-compact. Les conditions suivantes sont
équivalentes
\begin{enumerate}
\item $\mathcal{L}$ est relativement très ample sur $X/S$,
\item il existe un recouvrement ouvert $S = \bigcup V_j$ tel
que $\mathcal{L}|_{f^{-1}(V_j)}$ est relativement très ample sur $f^{-1}(V_j)/V_j$
pour tout $j$,
\item il existe un faisceau quasi-cohérent
d'algèbres $\mathcal{O}_S$ graduées
$\mathcal{A}$ générées en degré $1$ sur $\mathcal{O}_S$
et une application d'algèbres $\mathcal{O}_X$
graduées $\psi : f^*\mathcal{A} \to \bigoplus_{n \geq 0} \mathcal{L}^{\otimes n}$ telle que $f^*\mathcal{A}_1 \to \mathcal{L}$
est surjectif
et le morphisme associé $r_{\mathcal{L}, \psi} : X \to \underline{\text{Proj}}_S(\mathcal{A})$ est une immersion,
et
\item $f$ est quasi-séparé,
l'application canonique $\psi : f^*f_*\mathcal{L} \to \mathcal{L}$ est surjectif,
et l'application associée $r_{\mathcal{L}, \psi} : X \to \mathbf{P}(f_*\mathcal{L})$ est une
immersion.
\end{enumerate}
\end{lemma}

\begin{proof}
Il est clair que (1) implique (2). Il est également
clair que (4) implique (1) ; l'hypothèse de
quasi-séparation dans (4) est utilisée pour garantir que $f_*\mathcal{L}$
est quasi-cohérent via Schémas, Lemme \ref{schemes-lemma-push-forward-quasi-coherent}.

\medskip\noindent
Supposons (2). Nous
allons démontrer (4). Soit $S = \bigcup V_j$ un
recouvrement ouvert comme dans (2). Définissez
$X_j = f^{-1}(V_j)$ et $f_j : X_j \to V_j$ la restriction
de $f$. On voit que $f$ est
séparé par le lemme \ref{morphisms-lemma-relatively-very-ample-separated} (car la séparation
est locale sur la base). Par hypothèse il existe un $\mathcal{O}_{V_j}$-module
$\mathcal{E}_j$ quasi-cohérent
et une immersion $i_j : X_j \to \mathbf{P}(\mathcal{E}_j)$ avec $\mathcal{L}|_{X_j} \cong i_j^*\mathcal{O}_{\mathbf{P}(\mathcal{E}_j)}(1)$. Le morphisme
$i_j$ correspond à une surjection
$f_j^*\mathcal{E}_j \to \mathcal{L}|_{X_j}$, voir Constructions
, section \ref{constructions-section-projective-bundle}. Cette application est adjointe
à une application $\mathcal{E}_j \to f_*\mathcal{L}|_{V_j}$ telle que
la composition
$$
f_j^*\mathcal{E}_j \to (f^*f_*\mathcal{L})|_{X_j} \to \mathcal{L}|_{X_j}
$$
est surjectif. Nous concluons
que $\psi : f^*f_*\mathcal{L} \to \mathcal{L}$ est surjectif. Soit $r_{\mathcal{L}, \psi} : X \to \mathbf{P}(f_*\mathcal{L})$
le morphisme associé. Il reste à montrer que
$r_{\mathcal{L}, \psi}$ est une immersion ; nous invitons le lecteur
à le prouver par lui-même. L'application $\mathcal{O}_{V_j}$-module
$\mathcal{E}_j \to f_*\mathcal{L}|_{V_j}$ détermine un homomorphisme sur les
algèbres symétriques, qui à son tour définit un
morphisme
$$
\mathbf{P}(f_*\mathcal{L}|_{V_j}) \supset U_j \longrightarrow
\mathbf{P}(\mathcal{E}_j)
$$
où $U_j$ est le sous-schéma
ouvert des Constructions, Lemme \ref{constructions-lemma-morphism-relative-proj}.
La compatibilité de $\psi$ avec $\mathcal{E}_j \to f_*\mathcal{L}|_{V_j}$
montre que $r_{\mathcal{L}, \psi}(X_j) \subset U_j$ et qu'il existe une
factorisation
$$
\xymatrix{
X_j \ar[r]^-{r_{\mathcal{L}, \psi}} & U_j \ar[r] &
\mathbf{P}(\mathcal{E}_j)
}
$$
Nous omettons la vérification.
Cela montre que $r_{\mathcal{L}, \psi}$ est une immersion.

\medskip\noindent
À ce stade, nous voyons que (1), (2) et (4) sont équivalents.
Clairement (4) implique (3). Supposons (3). Nous
allons démontrer (1). Soit $\mathcal{A}$ un faisceau quasi-cohérent de $\mathcal{O}_S$-algèbres
graduées générées en degré $1$ sur $\mathcal{O}_S$.
Considérons l'application des $\mathcal{O}_S$-algèbres graduées $\text{Sym}(\mathcal{A}_1) \to \mathcal{A}$. Ceci
est surjectif par hypothèse et induit donc une immersion fermée
$$
\underline{\text{Proj}}_S(\mathcal{A})
\longrightarrow
\mathbf{P}(\mathcal{A}_1)
$$
qui ramène $\mathcal{O}(1)$ à $\mathcal{O}(1)$,
voir Constructions,
Lemme \ref{constructions-lemma-surjective-generated-degree-1-map-relative-proj}. Il est donc clair
que (3) implique (1).
\end{proof}

\begin{lemma}
\label{morphisms-lemma-very-ample-base-change}
Soit $f : X \to S$ un morphisme de schémas. Soit
$\mathcal{L}$ un module $\mathcal{O}_X$ inversible. Soit $S' \to S$
un morphisme de schémas. Soit $f' : X' \to S'$
le changement de base de $f$ et notons $\mathcal{L}'$
le retrait de $\mathcal{L}$ vers $X'$. Si $\mathcal{L}$
est $f$-très ample, alors $\mathcal{L}'$ est $f'$-très ample.
\end{lemma}

\begin{proof}
Par Définition \ref{morphisms-definition-very-ample} il existe un $\mathcal{O}_S$-module
$\mathcal{E}$ quasi-cohérent et une
immersion $i : X \to \mathbf{P}(\mathcal{E})$ sur $S$
telle que $\mathcal{L} \cong i^*\mathcal{O}_{\mathbf{P}(\mathcal{E})}(1)$. Le changement de
base de $\mathbf{P}(\mathcal{E})$ vers $S'$
est le fibré projectif associé à l'image réciproque
$\mathcal{E}'$ de $\mathcal{E}$,
et l'image réciproque de $\mathcal{O}_{\mathbf{P}(\mathcal{E})}(1)$
est
$\mathcal{O}_{\mathbf{P}(\mathcal{E}')}(1)$, voir Constructions,
Lemme \ref{constructions-lemma-relative-proj-base-change}. Enfin, le changement
de base
d'une immersion est une
immersion (Schémas, Lemme \ref{schemes-lemma-base-change-immersion}).
\end{proof}







\section{Faisceaux amples et très amples relativement aux morphismes de type fini}
\label{morphisms-section-ample-finite-type}

\noindent
En fait, la majeure partie du contenu de cette section concerne la notion de
projet de (quasi-)morphisme que nous n'avons pas encore défini.

\begin{lemma}
\label{morphisms-lemma-very-ample-finite-type-over-affine}
Soit $f : X \to S$ un morphisme de schémas. Soit $\mathcal{L}$
un faisceau inversible sur $X$. Supposons
que
\begin{enumerate}
\item le faisceau inversible $\mathcal{L}$ est très ample sur $X/S$,
\item le morphisme $X \to S$ est de type fini, et
\item $S$ est affine.
\end{enumerate}
Alors il existe un $n \geq 0$
et une immersion $i : X \to \mathbf{P}^n_S$ sur
$S$ tel que $\mathcal{L} \cong i^*\mathcal{O}_{\mathbf{P}^n_S}(1)$.
\end{lemma}

\begin{proof}
Supposons (1), (2) et
(3). La condition (3) signifie $S = \Spec(R)$ pour
certains anneaux $R$.
La condition (1) signifie par définition
qu'il existe un $\mathcal{O}_S$-module $\mathcal{E}$ quasi-cohérent
et une immersion $\alpha : X \to \mathbf{P}(\mathcal{E})$ telle que $\mathcal{L} = \alpha^*\mathcal{O}_{\mathbf{P}(\mathcal{E})}(1)$. Écrivez $\mathcal{E} = \widetilde{M}$
pour certains $R$-module $M$.
On a donc
$$
\mathbf{P}(\mathcal{E}) = \text{Proj}(\text{Sym}_R(M)).
$$
Puisque $\alpha$ est une immersion, et puisque
la topologie de $\text{Proj}(\text{Sym}_R(M))$ est générée par le standard
ouvre $D_{+}(f)$, $f \in \text{Sym}_R^d(M)$, $d \geq 1$,
on peut trouver pour chaque
$x \in X$ un $f \in \text{Sym}_R^d(M)$, $d \geq 1$, avec $\alpha(x) \in D_{+}(f)$ tel
que
$$
\alpha|_{\alpha^{-1}(D_{+}(f))} : \alpha^{-1}(D_{+}(f)) \to D_{+}(f)
$$
soit une immersion fermée.
La condition (2) implique que $X$ est quasi-compact.
On peut donc trouver
une collection finie d'éléments $f_j \in \text{Sym}_R^{d_j}(M)$,
$d_j \geq 1$ telle que pour chaque $f = f_j$ l'application
affichée ci-dessus est une immersion fermée et telle que
$\alpha(X) \subset \bigcup D_{+}(f_j)$. Écrivez $U_j = \alpha^{-1}(D_{+}(f_j))$. Notez que $U_j$ est affine
comme sous-schéma fermé du schéma affine $D_{+}(f_j)$.
Écrivez $U_j = \Spec(A_j)$. La condition (2) implique également
que $A_j$ est de type
fini sur $R$, voir Lemme \ref{morphisms-lemma-locally-finite-type-characterize}. Choisissez
un nombre fini de $x_{j, k} \in A_j$ que
génère $A_j$ en tant qu'algèbre $R$. Puisque $\alpha|_{U_j}$ est une
immersion fermée on voit que $x_{j, k}$ est l'image d'un élément
$$
f_{j, k}/f_j^{e_{j, k}} \in \text{Sym}_R(M)_{(f_j)}
=
\Gamma(D_{+}(f_j), \mathcal{O}_{\text{Proj}(\text{Sym}_R(M))}).
$$
Enfin, choisir $n \geq 1$ et les éléments $y_0, \ldots, y_n \in M$ tels que chacun des polynômes
$f_j, f_{j, k} \in \text{Sym}_R(M)$ soit un polynôme dans les éléments $y_t$ à coefficients
en $R$. Considérons l'homomorphisme
gradué d'anneaux
$$
\psi : R[Y_0, \ldots, Y_n] \longrightarrow \text{Sym}_R(M),
\quad Y_i \longmapsto y_i.
$$
Désignons $F_j$, $F_{j, k}$ les éléments de $R[Y_0, \ldots, Y_n]$
tels que $\psi(F_j) = f_j$ et $\psi(F_{j, k}) = f_{j, k}$.
Par Constructions, Lemme \ref{constructions-lemma-morphism-proj} on obtient un
sous-schéma ouvert
$$
U(\psi) \subset \text{Proj}(\text{Sym}_R(M))
$$
et un morphisme
$r_\psi : U(\psi) \to \mathbf{P}^n_R$. Ce morphisme satisfait
à $r_\psi^{-1}(D_{+}(F_j)) = D_{+}(f_j)$, et nous voyons donc que
$\alpha(X) \subset U(\psi)$. De plus, il est clair
que
$$
i = r_\psi \circ \alpha : X \longrightarrow \mathbf{P}^n_R
$$
est encore une immersion
puisque $i^\sharp(F_{j, k}/F_j^{e_{j, k}}) = x_{j, k} \in
A_j = \Gamma(U_j, \mathcal{O}_X)$ par construction.
De plus, le morphisme $r_\psi$ est équipé d'une application
$\theta : r_\psi^*\mathcal{O}_{\mathbf{P}^n_R}(1)
\to \mathcal{O}_{\text{Proj}(\text{Sym}_R(M))}(1)|_{U(\psi)}$ qui est dans ce
cas un isomorphisme (pour la construction
$\theta$ voir le lemme cité ci-dessus
; certains détails omis). Puisque
l'application d'origine $\alpha$
était
supposée avoir la propriété que $\mathcal{L} = \alpha^*\mathcal{O}_{\text{Proj}(\text{Sym}_R(M))}(1)$
on obtient le résultat.
\end{proof}

\begin{lemma}
\label{morphisms-lemma-quasi-affine-finite-type-over-S}
Soit $\pi : X \to S$ un morphisme de schémas. Supposons
que $X$ soit quasi-affine et que $\pi$ soit localement de
type fini. Il existe alors $n \geq 0$ et une immersion $i : X \to \mathbf{A}^n_S$ sur
$S$.
\end{lemma}

\begin{proof}
Soit $A = \Gamma(X, \mathcal{O}_X)$. Par hypothèse $X$ est quasi-compact
et s'identifie à un sous-schéma ouvert de $\Spec(A)$,
voir Propriétés, Lemme \ref{properties-lemma-quasi-affine}. De plus,
l'ensemble des ouvertures $X_f$, pour celles $f \in A$
telles que $X_f$ est affine, constitue une base
pour la topologie de $X$, voir la démonstration
de Propriétés, Lemme \ref{properties-lemma-quasi-affine}. On peut donc trouver un
nombre fini de $f_j \in A$, $j = 1, \ldots, m$ tel que $X = \bigcup X_{f_j}$,
et tel que $\pi(X_{f_j}) \subset V_j$ pour certains
$V_j \subset S$ ouverts affines. Par le Lemme
\ref{morphisms-lemma-locally-finite-type-characterize} les homomorphismes d'anneaux $\mathcal{O}(V_j) \to \mathcal{O}(X_{f_j}) = A_{f_j}$ sont
de type fini. Ainsi, nous pouvons choisir $a_1, \ldots, a_N \in A$ tel que
les éléments $a_1, \ldots, a_N, 1/f_j$ génèrent $A_{f_j}$
sur $\mathcal{O}(V_j)$ pour chaque $j$. Prenons $n = m + N$ et
soit
$$
i : X \longrightarrow \mathbf{A}^n_S
$$
le morphisme donné par les sections
globales $f_1, \ldots, f_m, a_1, \ldots, a_N$ du faisceau structural de $X$.
Soit $D(x_j) \subset \mathbf{A}^n_S$ le sous-schéma ouvert où la $j$ème
fonction de coordonnées
est non nulle. Alors pour $1 \leq j \leq m$ nous avons
$i^{-1}(D(x_j)) = X_{f_j}$ et le morphisme induit $X_{f_j} \to D(x_j)$ se factorise
par l'ouvert affine $\Spec(\mathcal{O}(V_j)[x_1, \ldots, x_n, 1/x_j])$ de $D(x_j)$.
Puisque l'homomorphisme
d'anneaux $\mathcal{O}(V_j)[x_1, \ldots, x_n, 1/x_j] \to A_{f_j}$ est surjectif
par construction Nous concluons que $i^{-1}(D(x_j)) \to D(x_j)$ est une
immersion comme voulu.
\end{proof}

\begin{lemma}
\label{morphisms-lemma-quasi-projective-finite-type-over-S}
Soit $f : X \to S$ un morphisme de schémas. Soit $\mathcal{L}$
un faisceau inversible sur $X$. Supposons
que
\begin{enumerate}
\item le faisceau inversible $\mathcal{L}$ est suffisant sur $X$, et que
\item le morphisme $X \to S$ est localement de type fini.
\end{enumerate}
Alors il existe un $d_0 \geq 1$ tel que pour tout
$d \geq d_0$ il existe un
$n \geq 0$ et une immersion $i : X \to \mathbf{P}^n_S$
sur $S$ tel que $\mathcal{L}^{\otimes d} \cong i^*\mathcal{O}_{\mathbf{P}^n_S}(1)$.
\end{lemma}

\begin{proof}
Soit
$A = \Gamma_*(X, \mathcal{L}) =
\bigoplus_{d \geq 0} \Gamma(X, \mathcal{L}^{\otimes d})$. Par Propriétés,
la Proposition \ref{properties-proposition-characterize-ample} l'ensemble des ouvertes affines $X_a$ avec $a \in A_{+}$ homogène forme
une base pour la topologie de $X$. Par conséquent,
nous pouvons trouver un nombre fini
d'éléments de ce type $a_0, \ldots, a_n \in A_{+}$ tels que
\begin{enumerate}
\item nous avons $X = \bigcup_{i = 0, \ldots, n} X_{a_i}$,
\item chaque $X_{a_i}$ est affine, et
\item chaque application $X_{a_i}$ dans un $V_i \subset S$ ouvert affine.
\end{enumerate}
Par le Lemme \ref{morphisms-lemma-locally-finite-type-characterize} on voit que les homomorphismes
d'anneaux $\mathcal{O}_S(V_i) \to \mathcal{O}_X(X_{a_i})$
sont de type fini. Nous pouvons donc
trouver un nombre fini d'éléments
$f_{ij} \in \mathcal{O}_X(X_{a_i})$, $j = 1, \ldots, n_i$ qui génèrent $\mathcal{O}_X(X_{a_i})$ comme une
$\mathcal{O}_S(V_i)$-algèbre. Par Propriétés, Lemme \ref{properties-lemma-invert-s-sections}, nous
pouvons écrire chaque $f_{ij}$ comme $a_{ij}/a_i^{e_{ij}}$
pour certains
$a_{ij} \in A_{+}$ homogènes. Soit $N$
un entier positif qui est un multiple commun de tous
les degrés des éléments $a_i$, $a_{ij}$.
Considérez les éléments
$$
a_i^{N/\deg(a_i)}, \ a_{ij}a_i^{(N/\deg(a_i)) - e_{ij}} \in A_N.
$$
Par construction ceux-ci génèrent le faisceau inversible
$\mathcal{L}^{\otimes N}$ sur $X$. Ils donnent donc lieu à
un morphisme
$$
j : X \longrightarrow
\mathbf{P}_S^{m}
\quad
\text{avec } m = n + \sum n_i
$$
sur $S$, voir Constructions, Lemme \ref{constructions-lemma-projective-space} et
Définition \ref{constructions-definition-projective-space}. De plus, $j^*\mathcal{O}_{\mathbf{P}_S}(1) = \mathcal{L}^{\otimes N}$. On
nomme les coordonnées homogènes $T_0, \ldots, T_n, T_{ij}$ au lieu
de $T_0, \ldots, T_m$. Pour $i = 0, \ldots, n$, nous avons $j^{-1}(D_{+}(T_i)) = X_{a_i}$.
De plus, en
retirant l'élément $T_{ij}/T_i$ via $j^\sharp$
on obtient l'élément $f_{ij} \in \mathcal{O}_X(X_{a_i})$. D'où le morphisme
$j$ restreint à $X_{a_i}$
donne une immersion fermée de $X_{a_i}$
dans l'ouvert affine $D_{+}(T_i) \cap \mathbf{P}^m_{V_i}$ de
$\mathbf{P}^N_S$. Nous concluons donc que le
morphisme $j$ est une immersion.
Cela implique que le lemme est valable pour certains $d$ et $n$, ce qui est suffisant
dans pratiquement toutes les applications.

\medskip\noindent
Cela prouve que pour un $d_2 \geq 1$
(à savoir $d_2 = N$ ci-dessus), certains
$m \geq 0$ il existe une certaine immersion $j : X \to \mathbf{P}^m_S$
donnée par les sections globales $s'_0, \ldots, s'_m \in \Gamma(X, \mathcal{L}^{\otimes d_2})$. Par
Propriétés, Proposition \ref{properties-proposition-characterize-ample} nous savons qu'il existe
un entier $d_1$ tel
que $\mathcal{L}^{\otimes d}$ est généré globalement pour tous les
$d \geq d_1$. Posons $d_0 = d_1 + d_2$. Nous affirmons que
le lemme est valable avec cette valeur de $d_0$. À savoir, étant donné un entier $d \geq d_0$,
nous pouvons choisir $s''_1, \ldots, s''_t
\in \Gamma(X, \mathcal{L}^{\otimes d - d_2})$ qui génère
$\mathcal{L}^{\otimes d - d_2}$ sur $X$. Définissez $k = (m + 1)t$ et
désignez $s_0, \ldots, s_k$ l'ensemble des sections
$s'_\alpha s''_\beta$, $\alpha = 0, \ldots, m$, $\beta = 1, \ldots, t$.
Ceux-ci génèrent $\mathcal{L}^{\otimes d}$ sur $X$ et
définissent donc un morphisme
$$
i : X \longrightarrow \mathbf{P}^{k - 1}_S
$$
tel que $i^*\mathcal{O}_{\mathbf{P}^n_S}(1) \cong \mathcal{L}^{\otimes d}$. Pour voir que $i$ est une immersion,
observez que $i$ est la composition
$$
X \longrightarrow \mathbf{P}^m_S \times_S \mathbf{P}^{t - 1}_S
\longrightarrow \mathbf{P}^{k - 1}_S
$$
où le premier morphisme est $(j, j')$ avec
$j'$ donné par $s''_1, \ldots, s''_t$
et le deuxième morphisme
est le plongement de Segre (Constructions, Lemme
\ref{constructions-lemma-segre-embedding}). Puisque $j$
est une immersion, $(j, j')$ l'est
aussi (appliquez Lemme \ref{morphisms-lemma-immersion-permanence} à $X \to \mathbf{P}^m_S \times_S \mathbf{P}^{t - 1}_S
\to \mathbf{P}^m_S$).
Ainsi $i$ est une composition d'immersions et donc
une immersion (Schémas, Lemme \ref{schemes-lemma-composition-immersion}).
\end{proof}

\begin{lemma}
\label{morphisms-lemma-finite-type-over-affine-ample-very-ample}
Soit $f : X \to S$ un morphisme de schémas. Soit $\mathcal{L}$
un module $\mathcal{O}_X$ inversible. Supposons que $S$
affine et $f$ de type fini. Les conditions
suivantes sont équivalentes
\begin{enumerate}
\item $\mathcal{L}$ est suffisant sur $X$,
\item $\mathcal{L}$ est $f$-ample,
\item $\mathcal{L}^{\otimes d}$ est $f$-très suffisant pour certains $d \geq 1$,
\item $\mathcal{L}^{\otimes d}$ est $f$-très suffisant pour tous les $d \gg 1$,
\item pour certains $d \geq 1$ il existe
$n \geq 1$ et une immersion
$i : X \to \mathbf{P}^n_S$ telle que $\mathcal{L}^{\otimes d} \cong i^*\mathcal{O}_{\mathbf{P}^n_S}(1)$, et
\item pour tout $d \gg 1$ il existe $n \geq 1$
et une immersion
$i : X \to \mathbf{P}^n_S$ telle que $\mathcal{L}^{\otimes d} \cong i^*\mathcal{O}_{\mathbf{P}^n_S}(1)$.
\end{enumerate}
\end{lemma}

\begin{proof}
L'équivalence de (1) et (2) est Lemme \ref{morphisms-lemma-ample-over-affine}. L'implication
(2) $\Rightarrow$ (6) est Lemme
\ref{morphisms-lemma-quasi-projective-finite-type-over-S}. Trivialement (6) implique (5).
Comme $\mathbf{P}^n_S$
est un fibré projectif sur $S$ (voir Constructions,
Lemme \ref{constructions-lemma-projective-space-bundle}) on voit que (5) implique
(3) et (6)
implique (4) de la définition d'un faisceau relativement
très ample.
Trivialement (4) implique (3). Pour finir il
faut montrer que (3) implique (2) qui découle du lemme
\ref{morphisms-lemma-ample-very-ample} et Lemme \ref{morphisms-lemma-ample-power-ample}.
\end{proof}

\begin{lemma}
\label{morphisms-lemma-finite-type-ample-very-ample}
Soit $f : X \to S$ un morphisme de schémas. Soit $\mathcal{L}$
un module $\mathcal{O}_X$ inversible. Supposons $S$ quasi-compact
et $f$ de type fini. Les conditions
suivantes sont équivalentes
\begin{enumerate}
\item $\mathcal{L}$ est $f$-ample,
\item $\mathcal{L}^{\otimes d}$ est $f$-très suffisant pour certains $d \geq 1$,
\item $\mathcal{L}^{\otimes d}$ est $f$-très suffisant pour tous les $d \gg 1$.
\end{enumerate}
\end{lemma}

\begin{proof}
Trivialement (3) implique (2). Lemme \ref{morphisms-lemma-ample-very-ample} garantit que (2) implique
(1) puisqu'un morphisme de type fini est quasi-compact
par définition. Supposons que $\mathcal{L}$ soit $f$-ample.
Choisissons un recouvrement fini par des ouverts affines $S = V_1 \cup \ldots \cup V_m$. Écrivons
$X_i = f^{-1}(V_i)$. Par le Lemme \ref{morphisms-lemma-finite-type-over-affine-ample-very-ample} ci-dessus on voit qu'il existe
un $d_0$ tel que $\mathcal{L}^{\otimes d}$ est relativement
très ample sur $X_i/V_i$ pour tout $d \geq d_0$. Nous concluons
donc que (1) implique (3) par le lemme \ref{morphisms-lemma-relatively-very-ample}.
\end{proof}

\noindent
Les deux lemmes suivants fournissent les caractérisations les plus utilisées
et les plus utiles des faisceaux inversibles relativement très amples
et relativement amples lorsque le morphisme est de type fini.

\begin{lemma}
\label{morphisms-lemma-characterize-very-ample-on-finite-type}
Soit $f : X \to S$ un morphisme de schémas. Soit $\mathcal{L}$
un faisceau inversible sur $X$. Supposons que $f$
soit de type fini. Les conditions
suivantes sont équivalentes :
\begin{enumerate}
\item $\mathcal{L}$ est $f$-relativement très ample, et
\item il existe un recouvrement ouvert $S = \bigcup V_j$, pour chaque
$j$ un entier $n_j$, et des immersions
$$
i_j :
X_j = f^{-1}(V_j) = V_j \times_S X
\longrightarrow
\mathbf{P}^{n_j}_{V_j}
$$
sur $V_j$
tel que $\mathcal{L}|_{X_j} \cong i_j^*\mathcal{O}_{\mathbf{P}^{n_j}_{V_j}}(1)$.
\end{enumerate}
\end{lemma}

\begin{proof}
On voit que (1) implique (2) en prenant un recouvrement ouvert
affine de $S$ et en appliquant le Lemme \ref{morphisms-lemma-very-ample-finite-type-over-affine} à chacune
des restrictions de $f$
et $\mathcal{L}$. On voit que (2) implique
(1) par le lemme \ref{morphisms-lemma-relatively-very-ample}.
\end{proof}

\begin{lemma}
\label{morphisms-lemma-characterize-ample-on-finite-type}
Soit $f : X \to S$ un morphisme de schémas. Soit $\mathcal{L}$
un faisceau inversible sur $X$. Supposons que $f$
soit de type fini. Les conditions
suivantes sont équivalentes :
\begin{enumerate}
\item $\mathcal{L}$ est $f$-relativement ample, et
\item il existe un recouvrement ouvert $S = \bigcup V_j$, pour chaque
$j$ un entier $d_j \geq 1$, $n_j \geq 0$,
et des immersions
$$
i_j :
X_j = f^{-1}(V_j) = V_j \times_S X
\longrightarrow
\mathbf{P}^{n_j}_{V_j}
$$
sur $V_j$ tel que
$\mathcal{L}^{\otimes d_j}|_{X_j} \cong
i_j^*\mathcal{O}_{\mathbf{P}^{n_j}_{V_j}}(1)$.
\end{enumerate}
\end{lemma}

\begin{proof}
On voit que (1) implique (2) en prenant un recouvrement ouvert
affine de $S$ et en appliquant le Lemme \ref{morphisms-lemma-finite-type-over-affine-ample-very-ample} à
chacune des restrictions de $f$
et $\mathcal{L}$. On voit que (2)
implique (1) par le lemme \ref{morphisms-lemma-characterize-relatively-ample}.
\end{proof}

\begin{lemma}
\label{morphisms-lemma-invertible-add-enough-ample-very-ample}
Soit $f : X \to S$ un morphisme de schémas. Soit
$\mathcal{N}$, $\mathcal{L}$ des modules $\mathcal{O}_X$ inversibles. Supposons
que $S$ est quasi-compact, $f$ est de type fini
et $\mathcal{L}$ est
$f$-ample. Alors $\mathcal{N} \otimes_{\mathcal{O}_X} \mathcal{L}^{\otimes d}$ est $f$-très suffisant
pour tous les $d \gg 1$.
\end{lemma}

\begin{proof}
Par le Lemme \ref{morphisms-lemma-characterize-very-ample-on-finite-type} on se réduit au cas
$S$ est affine. En combinant
Lemme \ref{morphisms-lemma-finite-type-over-affine-ample-very-ample} et Propriétés, Proposition
\ref{properties-proposition-characterize-ample}, nous pouvons trouver un entier
$d_0$ tel que $\mathcal{N} \otimes \mathcal{L}^{\otimes d_0}$
est généré globalement. Choisissez
les sections globales $s_0, \ldots, s_n$
de $\mathcal{N} \otimes \mathcal{L}^{\otimes d_0}$ qui le génèrent. Ceci détermine
un morphisme $j : X \to \mathbf{P}^n_S$
sur $S$. Par le Lemme
\ref{morphisms-lemma-finite-type-over-affine-ample-very-ample} on peut aussi choisir un entier
$d_1$ tel que pour tout $d \geq d_1$
il existe des sections $t_{d, 0}, \ldots, t_{d, n(d)}$ de $\mathcal{L}^{\otimes d}$
qui le génèrent et définissent une
immersion
$$
j_d = \varphi_{\mathcal{L}^{\otimes d}, t_{d, 0}, \ldots, t_{d, n(d)}} :
X
\longrightarrow
\mathbf{P}^{n(d)}_S
$$
sur $S$. Alors pour $d \geq d_0 + d_1$ on peut considérer le
morphisme
$$
\varphi_{\mathcal{N} \otimes \mathcal{L}^{\otimes d},
s_j \otimes t_{d - d_0, i}} :
X
\longrightarrow
\mathbf{P}^{(n + 1)(n(d - d_0) + 1) - 1}_S
$$
Ce morphisme est une immersion car c'est la composition
$$
X \to \mathbf{P}^n_S \times_S \mathbf{P}^{n(d - d_0)}_S
\to \mathbf{P}^{(n + 1)(n(d - d_0) + 1) - 1}_S
$$
où le premier morphisme est $(j, j_{d - d_0})$ et
le second est le plongement
de Segre (Constructions, Lemme \ref{constructions-lemma-segre-embedding}). Puisque
$j$ est une immersion, $(j, j_{d - d_0})$
l'est aussi (appliquer le Lemme \ref{morphisms-lemma-immersion-permanence}). Nous avons une composition
d'immersions et
donc une immersion (Schémas, Lemme \ref{schemes-lemma-composition-immersion}).
\end{proof}










\section{Morphismes quasi-projectifs}
\label{morphisms-section-quasi-projective}

\noindent
La discussion dans la section précédente suggère les définitions suivantes. Nous reprenons
notre définition du quasi-projectif de \cite{EGA}. La version avec la lettre
``H'' est la définition dans \cite{H}.

\begin{definition}
\label{morphisms-definition-quasi-projective}
\begin{reference}
\cite[II, Définition 5.3.1]{EGA} et \cite[page 103]{H}
\end{reference}
Soit $f : X \to S$ un morphisme de schémas.
\begin{enumerate}
\item On dit que $f$ est {\it quasi-projectif} si $f$
est de type fini et qu'il existe un $f$-relativement ample module inversible $\mathcal{O}_X$.
\item On dit que $f$ est {\it H-quasi-projectif}
s'il existe une immersion quasi-compacte $X \to \mathbf{P}^n_S$ sur $S$
pour un certain $n$.\footnote{Ce n'est pas exactement la
même définition que celle dans Hartshorne. À savoir, la définition dans Hartshorne
(8e impression corrigée, 1997) est que $f$ devrait être la composition d'une
immersion ouverte suivie d'un H-morphisme projectif (voir Définition \ref{morphisms-definition-projective}),
ce qui n'implique pas que
$f$ soit quasi-compact. Voir le Lemme \ref{morphisms-lemma-H-quasi-projective-open-H-projective} pour
l'implication dans l'autre sens.}
\item On dit que $f$ est {\it localement quasi-projectif}
s'il existe un recouvrement ouvert $S = \bigcup V_j$ tel que chaque $f^{-1}(V_j) \to V_j$
soit quasi-projectif.
\end{enumerate}
\end{definition}

\noindent
Comme cette définition le suggère, la propriété d'être quasi-projectif
n'est pas locale sur $S$. À un stade ultérieur, nous
pourrons en dire plus sur la catégorie des schémas
quasi-projectifs, voir Compléments sur les morphismes, section \ref{more-morphisms-section-quasi-projective}.

\begin{lemma}
\label{morphisms-lemma-base-change-quasi-projective}
Un changement de base d'un morphisme quasi-projectif est quasi-projectif.
\end{lemma}

\begin{proof}
Cela résulte
des Lemme \ref{morphisms-lemma-base-change-finite-type}
et \ref{morphisms-lemma-ample-base-change}.
\end{proof}

\begin{lemma}
\label{morphisms-lemma-composition-quasi-projective}
Soient $f : X \to Y$ et $g : Y \to S$ des morphismes de schémas. Si $S$
est quasi-compact et $f$ et $g$ sont quasi-projectifs,
alors $g \circ f$ est quasi-projectif.
\end{lemma}

\begin{proof}
Cela résulte
des Lemme \ref{morphisms-lemma-composition-finite-type}
et \ref{morphisms-lemma-ample-composition}.
\end{proof}

\begin{lemma}
\label{morphisms-lemma-quasi-projective-properties}
Soit $f : X \to S$ un morphisme de schémas. Si $f$ est quasi-projectif, ou
H-quasi-projectif ou localement quasi-projectif, alors $f$
est séparé de type fini.
\end{lemma}

\begin{proof}
Omis.
\end{proof}

\begin{lemma}
\label{morphisms-lemma-H-quasi-projective-quasi-projective}
Un H-morphisme quasi-projectif est quasi-projectif.
\end{lemma}

\begin{proof}
Omis.
\end{proof}

\begin{lemma}
\label{morphisms-lemma-characterize-locally-quasi-projective}
Soit $f : X \to S$ un morphisme de schémas. Les conditions
suivantes sont équivalentes :
\begin{enumerate}
\item Le morphisme $f$ est localement quasi-projectif.
\item Il existe un recouvrement ouvert $S = \bigcup V_j$ tel
que chaque $f^{-1}(V_j) \to V_j$ soit H-quasi-projectif.
\end{enumerate}
\end{lemma}

\begin{proof}
D'après le Lemme \ref{morphisms-lemma-H-quasi-projective-quasi-projective} nous voyons que (2) implique (1).
Supposons (1). La question est locale sur
$S$ et par conséquent nous pouvons supposer que $S$ est affine,
$X$ de type fini sur
$S$ et $\mathcal{L}$ est un faisceau inversible relativement ample
sur $X/S$. D'après le Lemme \ref{morphisms-lemma-finite-type-over-affine-ample-very-ample} nous pouvons supposer
que $\mathcal{L}$ est ample sur $X$. D'après
le Lemme \ref{morphisms-lemma-quasi-projective-finite-type-over-S} nous voyons qu'il existe une immersion de $X$ dans
un espace projectif sur $S$,
c'est-à-dire que $X$ est H-quasi-projectif sur $S$ comme
souhaité.
\end{proof}

\begin{lemma}
\label{morphisms-lemma-quasi-affine-finite-type-quasi-projective}
\begin{reference}
\cite[II, Proposition 5.3.4 (i)]{EGA}
\end{reference}
Un morphisme quasi-affine de type fini est quasi-projectif.
\end{lemma}

\begin{proof}
Cela résulte du Lemme \ref{morphisms-lemma-quasi-affine-O-ample}.
\end{proof}

\begin{lemma}
\label{morphisms-lemma-quasi-projective-permanence}
Soient $g : Y \to S$ et $f : X \to Y$ des morphismes de schémas. Si $g \circ f$
est quasi-projectif et $f$ est
quasi-compact\footnote{Cela découle si $g$ est quasi-séparé
d'après Schémas, Lemme \ref{schemes-lemma-quasi-compact-permanence}.}, alors $f$
est quasi-projectif.
\end{lemma}

\begin{proof}
Observez que $f$ est de
type fini d'après le Lemme \ref{morphisms-lemma-permanence-finite-type}.
Ainsi, le lemme découle du Lemme \ref{morphisms-lemma-ample-permanence} et des
définitions.
\end{proof}





\section{Morphismes propres}
\label{morphisms-section-proper}

\noindent
La notion de morphisme propre joue un rôle important en géométrie algébrique.
Un exemple important de morphisme propre sera le morphisme
de structure $\mathbf{P}^n_S \to S$ de l'espace projectif $n$, et
ceci est en fait l'exemple motivant qui a conduit à la définition.

\begin{definition}
\label{morphisms-definition-proper}
Soit $f : X \to S$ un morphisme de schémas. Nous disons que $f$
est {\it propre} si $f$ est séparé, de type fini
et universellement fermé.
\end{definition}

\noindent
Le morphisme de la droite affine avec zéro doublé vers la droite affine est de type
fini et universellement fermé, donc la condition de séparation est nécessaire
dans la définition ci-dessus. Dans le
reste de cette section, nous prouvons certaines des propriétés de base des
morphismes propres et des morphismes universellement fermés.

\begin{lemma}
\label{morphisms-lemma-universally-closed-local-on-the-base}
Soit $f : X \to S$ un morphisme de schémas. Les conditions
suivantes sont équivalentes :
\begin{enumerate}
\item Le morphisme $f$ est universellement fermé.
\item Il existe un recouvrement ouvert $S = \bigcup V_j$ tel que $f^{-1}(V_j) \to V_j$
soit universellement fermé pour tous les indices $j$.
\end{enumerate}
\end{lemma}

\begin{proof}
Ceci est clair d'après la définition.
\end{proof}

\begin{lemma}
\label{morphisms-lemma-proper-local-on-the-base}
Soit $f : X \to S$ un morphisme de schémas. Les conditions
suivantes sont équivalentes :
\begin{enumerate}
\item Le morphisme $f$ est propre.
\item Il existe un recouvrement ouvert $S = \bigcup V_j$ tel que
$f^{-1}(V_j) \to V_j$ soit propre pour tous les indices $j$.
\end{enumerate}
\end{lemma}

\begin{proof}
Omis.
\end{proof}

\begin{lemma}
\label{morphisms-lemma-composition-proper}
La composition de morphismes propres est propre. Il en
va de même pour les morphismes universellement fermés.
\end{lemma}

\begin{proof}
Une composition de morphismes fermés est
fermée. Si $X \to Y \to Z$ sont des morphismes universellement
fermés et $Z' \to Z$ est un morphisme
quelconque, alors nous voyons que $Z' \times_Z X = (Z' \times_Z Y) \times_Y X  \to Z' \times_Z Y$ est fermé
et $Z' \times_Z Y \to Z'$ est fermé. D'où le résultat pour
les morphismes universellement fermés. Nous
avons vu que ``séparé'' et ``de
type fini'' sont
conservés sous les compositions (Schémas, Lemme \ref{schemes-lemma-separated-permanence}
et Lemme \ref{morphisms-lemma-composition-finite-type}). D'où le résultat pour
les morphismes propres.
\end{proof}

\begin{lemma}
\label{morphisms-lemma-base-change-proper}
Le changement de base d'un morphisme propre est propre. Il
en est de même pour les morphismes universellement fermés.
\end{lemma}

\begin{proof}
Ceci est vrai par définition pour les morphismes universellement
fermés. C'est vrai pour les
morphismes séparés (Schémas, Lemme \ref{schemes-lemma-separated-permanence}). C'est
vrai pour les morphismes de type fini
(Lemme \ref{morphisms-lemma-base-change-finite-type}). Par conséquent, c'est
vrai pour les morphismes propres.
\end{proof}

\begin{lemma}
\label{morphisms-lemma-closed-immersion-proper}
Une immersion fermée est propre, donc a fortiori universellement fermée.
\end{lemma}

\begin{proof}
Le changement de base d'une immersion fermée est une
immersion fermée (Schémas, Lemme \ref{schemes-lemma-base-change-immersion}). Par conséquent,
elle est universellement
fermée. Une immersion
fermée est séparée (Schémas, Lemme \ref{schemes-lemma-immersions-monomorphisms}). Une
immersion fermée est de type
fini (Lemme \ref{morphisms-lemma-immersion-locally-finite-type}). Par conséquent, une
immersion fermée est propre.
\end{proof}

\begin{lemma}
\label{morphisms-lemma-image-proper-scheme-closed}
Supposons donné un diagramme commutatif de schémas
$$
\xymatrix{
X \ar[rr] \ar[rd] & &
Y \ar[ld] \\
& S &
}
$$
avec $Y$ séparé sur $S$.
\begin{enumerate}
\item Si $X \to S$ est universellement fermé, alors le morphisme $X \to Y$
est universellement fermé.
\item Si $X$ est propre sur $S$, alors le morphisme $X \to Y$ est propre.
\end{enumerate}
En particulier, dans les deux cas, l'image de $X$ dans $Y$ est fermée.
\end{lemma}

\begin{proof}
Supposons que $X \to S$ soit universellement fermé (resp. \ propre).
Nous factorisons le morphisme comme
$X \to X \times_S Y \to Y$. Le premier morphisme est une immersion
fermée, voir Schémas, Lemme \ref{schemes-lemma-semi-diagonal}. Par
conséquent, le premier morphisme est propre (Lemme \ref{morphisms-lemma-closed-immersion-proper}). La projection $X \times_S Y \to Y$
est le changement de base d'un morphisme universellement
fermé (resp. \ propre) et donc universellement
fermé (resp. \ propre), voir Lemme \ref{morphisms-lemma-base-change-proper}. Ainsi $X \to Y$
est universellement fermé (resp. \ propre) comme la composition
de morphismes universellement fermés (resp.
\ propres) (Lemme \ref{morphisms-lemma-composition-proper}).
\end{proof}

\noindent
La démonstration du lemme suivant est due à Bjorn Poonen ; voir
\href{https://mathoverflow.net/questions/23337/is-a-universally-closed-morphism-of-schemes-quasi-compact/23528#23528}{cette page}.

\begin{lemma}
\label{morphisms-lemma-universally-closed-quasi-compact}
\begin{reference}
Due à Bjorn Poonen.
\end{reference}
Un morphisme universellement fermé de schémas est quasi-compact.
\end{lemma}

\begin{proof}
Soit $f : X \to S$ un morphisme. Supposons que $f$ n'est pas quasi-compact.
Notre objectif est de montrer que $f$ n'est pas universellement
fermé. D'après Schémas, Lemme \ref{schemes-lemma-quasi-compact-affine} il existe
un ouvert affine $V \subset S$ tel que $f^{-1}(V)$ n'est pas quasi-compact.
Pour atteindre notre objectif il suffit de montrer que $f^{-1}(V) \to V$
n'est pas universellement fermé, donc nous pouvons supposer que $S = \Spec(A)$
pour un certain anneau $A$.

\medskip\noindent
Écrivons $X = \bigcup_{i \in I} X_i$ où les $X_i$ sont des sous-schémas
ouverts affines de $X$. Soit
$T = \Spec(A[y_i ; i \in I])$. Soit $T_i = D(y_i) \subset T$. Soit $Z$
l'ensemble fermé $(X \times_S T) - \bigcup_{i \in I} (X_i \times_S T_i)$. Il suffit de prouver que
l'image $f_T(Z)$ de $Z$ sous $f_T : X \times_S T \to T$ n'est pas
fermée.

\medskip\noindent Il
existe un point $s \in S$ tel qu'il n'existe aucun voisinage
$U$ de $s$ dans $S$ tel que $X_U$ soit quasi-compact.
Sinon, nous pourrions recouvrir $S$ avec un nombre fini
de tels $U$ et Schémas, le Lemme \ref{schemes-lemma-quasi-compact-affine} impliquerait
que $f$ est quasi-compact. Fixons un tel $s \in S$.

\medskip\noindent
D'abord, nous vérifions que $f_T(Z_s) \ne T_s$. Soit $t \in T$ le point situé
au-dessus de $s$ avec $\kappa(t) = \kappa(s)$ tel que $y_i = 1$ dans
$\kappa(t)$ pour tous
$i$. Alors $t \in T_i$ pour tous $i$, et la fibre de
$Z_s \to T_s$ au-dessus de $t$ est isomorphe à $(X - \bigcup_{i \in I} X_i)_s$, qui est
vide. Ainsi $t \in T_s - f_T(Z_s)$.

\medskip\noindent Supposons
que $f_T(Z)$ soit fermé dans $T$. Alors il existe
un élément $g \in A[y_i; i \in I]$ avec $f_T(Z) \subset V(g)$ mais $t \not \in V(g)$. Par conséquent, l'image de
$g$ dans $\kappa(t)$ est non nulle. En particulier, un certain coefficient
de $g$ a une image non nulle dans $\kappa(s)$. Par conséquent, ce coefficient
est inversible sur un certain voisinage $U$ de $s$. Soit $J$
l'ensemble fini de $j \in I$ tel que $y_j$ apparaisse dans $g$.
Puisque $X_U$ n'est pas quasi-compact, nous pouvons choisir un point $x \in X - \bigcup_{j \in J} X_j$
situé au-dessus de certain $u \in U$. Puisque $g$ a un coefficient
qui est inversible sur $U$, nous pouvons trouver un point $t' \in T$
situé au-dessus de $u$ tel que $t' \not \in V(g)$ et $t' \in V(y_i)$
pour tous $i \notin J$. Ceci est vrai parce que $V(y_i; i \in I, i \not\in J) = \Spec(A[y_j; j\in J])$.
et l'ensemble des points de ce schéma situés au-dessus de $u$
est bijectif avec $\Spec(\kappa(u)[y_j; j \in J])$. En d'autres termes $t' \notin T_i$ pour
chaque $i \notin J$. D'après
Schémas, Lemme \ref{schemes-lemma-points-fibre-product}, nous pouvons trouver
un point $z$ de $X \times_S T$ se projetant sur $x \in X$ et sur
$t' \in T$. Puisque $x \not \in X_j$ pour $j \in J$ et $t' \not \in T_i$ pour $i \in I \setminus J$,
nous voyons que $z \in Z$. D'autre part, $f_T(z) = t' \not \in V(g)$,
ce qui contredit $f_T(Z) \subset V(g)$. Ainsi, l'hypothèse ``$f_T(Z)$
clôt'' est fausse et nous concluons en effet que $f_T$
n'est pas fermée, comme souhaité.
\end{proof}

\noindent
Le lemme suivant dit que l'image d'un schéma propre (dans un schéma séparé
de type fini sur la base) est propre.

\begin{lemma}
\label{morphisms-lemma-image-proper-is-proper}
Soit $S$ un schéma. Soit $f : X \to Y$ un morphisme de schémas au-dessus de $S$. Si $X$ est
universellement fermé au-dessus de $S$ et que $f$ est surjectif, alors $Y$
est universellement fermé au-dessus de $S$. En particulier, si également $Y$ est séparé
et localement de type fini sur $S$, alors $Y$ est propre sur $S$.
\end{lemma}

\begin{proof}
Supposons que $X$ soit universellement fermé
et que $f$ soit surjectif. Notons $p : X \to S$,
$q : Y \to S$ les morphismes de structure. Soit $S' \to S$
un morphisme de schémas. Le changement
de base $f' : X_{S'} \to Y_{S'}$ est surjectif (Lemme \ref{morphisms-lemma-base-change-surjective}),
et le changement de base $p' : X_{S'} \to S'$
est fermé. Si $T \subset Y_{S'}$ est fermé, alors $(f')^{-1}(T) \subset X_{S'}$
est fermé, donc $p'((f')^{-1}(T)) = q'(T)$ est fermé. Ainsi $q'$
est fermé. Cela prouve la première affirmation.
Donc $Y \to S$ est quasi-compact
par le Lemme \ref{morphisms-lemma-universally-closed-quasi-compact} et par conséquent
$Y \to S$ est propre par définition
si de plus $Y \to S$ est localement de type fini et séparé.
\end{proof}

\begin{lemma}
\label{morphisms-lemma-scheme-theoretic-image-is-proper}
Supposons donné un diagramme commutatif de schémas
$$
\xymatrix{
X \ar[rr]_h \ar[rd]_f & & Y \ar[ld]^g \\
& S
}
$$
Supposons
\begin{enumerate}
\item $X \to S$ est un morphisme universellement fermé (par exemple propre), et
\item $Y \to S$ est séparé et localement de type fini.
\end{enumerate}
Alors l'image schématique $Z \subset Y$ de $h$ est propre
sur $S$ et $X \to Z$ est surjectif.
\end{lemma}

\begin{proof}
L'image schématique de $h$ est construite
dans la section \ref{morphisms-section-scheme-theoretic-image}. Puisque $f$ est
quasi-compact (Lemme \ref{morphisms-lemma-universally-closed-quasi-compact}), nous
trouvons que $h$ est
quasi-compact (Schémas, Lemme \ref{schemes-lemma-quasi-compact-permanence}). Par
conséquent, $h(X) \subset Z$
est dense (Lemme \ref{morphisms-lemma-quasi-compact-scheme-theoretic-image}). D'autre part, $h(X)$
est fermé dans $Y$
(Lemme \ref{morphisms-lemma-image-proper-scheme-closed}), donc $X \to Z$
est surjectif. Ainsi,
$Z \to S$ est un morphisme propre (Lemme \ref{morphisms-lemma-image-proper-is-proper}).
\end{proof}

\noindent
Le but d'un schéma séparé sous un morphisme surjectif
universellement fermé est séparé.

\begin{lemma}
\label{morphisms-lemma-image-universally-closed-separated}
Soit $S$ un schéma. Soit $f : X \to Y$ un morphisme surjectif universellement fermé
de schémas sur $S$.
\begin{enumerate}
\item Si $X$ est quasi-séparé, alors $Y$ est quasi-séparé.
\item Si $X$ est séparé, alors $Y$ est séparé.
\item Si $X$ est quasi-séparé sur $S$, alors $Y$ est quasi-séparé sur $S$.
\item Si $X$ est séparé sur $S$, alors $Y$ est séparé sur $S$.
\end{enumerate}
\end{lemma}

\begin{proof}
Les parties (1) et (2) sont une conséquence de (3)
et (4) pour $S = \Spec(\mathbf{Z})$
(voir Schémas, Définition \ref{schemes-definition-separated}). Considérons
le diagramme commutatif
$$
\xymatrix{
X \ar[d] \ar[rr]_{\Delta_{X/S}} & & X \times_S X \ar[d] \\
Y \ar[rr]^{\Delta_{Y/S}} & & Y \times_S Y
}
$$
La flèche verticale de gauche est surjective (c'est-à-dire universellement
surjective). La flèche verticale de droite est universellement
fermée comme composition
des morphismes universellement fermés $X \times_S X \to X \times_S Y \to Y \times_S Y$. Par conséquent,
elle est également
quasi-compacte, voir Lemme \ref{morphisms-lemma-universally-closed-quasi-compact}.

\medskip\noindent
Supposons que $X$ soit quasi-séparé sur $S$, c'est-à-dire
que $\Delta_{X/S}$ soit quasi-compact. Si $V \subset Y \times_S Y$ est
un ouvert quasi-compact, alors $V \times_{Y \times_S Y} X \to \Delta_{Y/S}^{-1}(V)$ est surjectif et $V \times_{Y \times_S Y} X$
est quasi-compact d'après nos remarques ci-dessus. Nous concluons que
$\Delta_{Y/S}$ est quasi-compact, c'est-à-dire que $Y$ est quasi-séparé
sur $S$.

\medskip\noindent
Supposons que $X$ soit séparé sur $S$, c'est-à-dire que
$\Delta_{X/S}$ soit une immersion fermée. Alors $X \to Y \times_S Y$
est fermé comme composition de morphismes fermés.
Puisque $X \to Y$ est surjectif, il s'ensuit que $\Delta_{Y/S}(Y)$ est fermé dans
$Y \times_S Y$. Ainsi $Y$ est séparé sur $S$ d'après
la discussion suivant Schémas, Définition \ref{schemes-definition-separated}.
\end{proof}








\section{Critères valuatifs}
\label{morphisms-section-valuative-criteria}

\noindent
Nous avons déjà discuté du critère valuatif pour la fermeture universelle
et pour la séparabilité dans
les schémas, sections \ref{schemes-section-valuative-criterion-universal-closedness} et \ref{schemes-section-valuative-separatedness}. Dans cette
section, nous discuterons de quelques
conséquences et variantes. Dans Limites, section \ref{limits-section-Noetherian-valuative-criterion},
nous montrerons qu'il suffit de considérer les anneaux
de valuation discrète lorsqu'on travaille avec
des schémas localement noethériens et des morphismes
de type fini.

\begin{lemma}[Critère valuatif pour la propreté]
\label{morphisms-lemma-characterize-proper}
\begin{reference}
\cite[II Théorème 7.3.8]{EGA}
\end{reference}
Soit $S$ un schéma. Soit $f : X \to Y$ un morphisme de schémas sur $S$.
Supposons que $f$ soit de type fini et quasi-séparé. Alors les
conditions suivantes sont équivalentes
\begin{enumerate}
\item $f$ est propre,
\item $f$ satisfait le critère
valuatif (Schémas, Définition \ref{schemes-definition-valuative-criterion}),
\item donné tout diagramme solide commutatif
$$
\xymatrix{
\Spec(K) \ar[r] \ar[d] & X \ar[d] \\
\Spec(A) \ar[r] \ar@{-->}[ru] & Y
}
$$
où $A$ est un anneau de valuation avec corps de fractions $K$, il existe une
flèche pointillée unique rendant le diagramme commutatif.
\end{enumerate}
\end{lemma}

\begin{proof}
La partie (3) est une reformulation de
(2). Ainsi,
le lemme est une conséquence formelle de Schémas,
Proposition \ref{schemes-proposition-characterize-universally-closed} et Lemme
\ref{schemes-lemma-separated-implies-valuative}, Lemme \ref{schemes-lemma-valuative-criterion-separatedness}, et des
définitions.
\end{proof}

\noindent
On n'a généralement pas besoin de considérer tous les diagrammes possibles
lorsqu'on teste le critère valuatif. Nous appellerons un critère valuatif
comme dans le lemme suivant un ``critère valuatif raffiné''.

\begin{lemma}
\label{morphisms-lemma-refined-valuative-criterion-universally-closed}
Soient $f : X \to S$ et $h : U \to X$ des morphismes de schémas. Supposons que $f$
et $h$ soient quasi-compacts et que $h(U)$ soit dense dans $X$. Si,
pour tout diagramme commutatif solide donné
$$
\xymatrix{
\Spec(K) \ar[r] \ar[d] & U \ar[r]^h & X \ar[d]^f \\
\Spec(A) \ar[rr] \ar@{-->}[rru] & & S
}
$$
où $A$ est un anneau de valuation avec corps des fractions $K$, il
existe un unique trait pointillé rendant le diagramme commutatif, alors $f$
est universellement fermé. Si de plus $f$ est quasi-séparé, alors $f$
est séparé.
\end{lemma}

\begin{proof}
Pour prouver que $f$ est universellement fermé, nous allons vérifier
la partie existence du critère valuatif pour
$f$, ce qui suffit d'après Schémas, Proposition \ref{schemes-proposition-characterize-universally-closed}. Pour ce faire,
considérons un diagramme commutatif
$$
\xymatrix{
\Spec(K) \ar[r] \ar[d] & X \ar[d] \\
\Spec(A) \ar[r] & S
}
$$
où $A$ est un anneau de valuation et $K$ est le corps des fractions
de $A$. Notez que, puisque les anneaux de valuation et
les corps sont réduits, nous pouvons remplacer $U$, $X$
et $S$ par leurs réductions respectives d'après Schémas,
Lemme \ref{schemes-lemma-map-into-reduction}. Dans ce cas, l'hypothèse selon laquelle
$h(U)$ est dense signifie que l'image schématique de $h : U \to X$
est $X$, voir Lemme \ref{morphisms-lemma-scheme-theoretic-image-reduced}. Nous pouvons
aussi remplacer $S$ par un ouvert affine par lequel
le morphisme $\Spec(A) \to S$ se factorise. Ainsi, nous pouvons
supposer que $S = \Spec(R)$.

\medskip\noindent
Soit $\Spec(B) \subset X$ un ouvert affine par lequel le
morphisme $\Spec(K) \to X$ se factorise. Choisissons une algèbre
polynomiale $P$ sur $B$ et une surjection de
$B$-algèbres $P \to K$. Alors $\Spec(P) \to X$
est plat. Par conséquent, l'image schématique du morphisme
$U \times_X \Spec(P) \to \Spec(P)$ est $\Spec(P)$ par le Lemme \ref{morphisms-lemma-flat-base-change-scheme-theoretic-image}.
D'après le Lemme \ref{morphisms-lemma-reach-points-scheme-theoretic-image} nous pouvons trouver
un diagramme commutatif
$$
\xymatrix{
\Spec(K') \ar[r] \ar[d] & U \times_X \Spec(P) \ar[d] \\
\Spec(A') \ar[r] & \Spec(P)
}
$$
où $A'$ est un anneau de valuation et $K'$ est le corps des fractions de $A'$,
tel que le point fermé de $\Spec(A')$ s'envoie
dans $\Spec(K) \subset \Spec(P)$. En d'autres termes,
il existe un homomorphisme de $B$-algèbres $\varphi : K \to A'/\mathfrak m_{A'}$.
Choisissons un anneau de valuation $A'' \subset A'/\mathfrak m_{A'}$ dominant
$\varphi(A)$ et ayant pour corps des fractions
$K'' = A'/\mathfrak m_{A'}$ (Algèbre, Lemme \ref{algebra-lemma-dominate}). On pose
$$
C = \{\lambda \in A' \mid \lambda \bmod \mathfrak m_{A'} \in A''\}.
$$
qui est un anneau de valuation
par l'Algèbre, Lemme \ref{algebra-lemma-stack-valuation-rings}. Comme $C$ est une $R$-algèbre
avec corps des fractions $K'$, nous obtenons un diagramme
commutatif
$$
\xymatrix{
\Spec(K') \ar[r] \ar[d] & U \ar[r] & X \ar[d] \\
\Spec(C) \ar[rr] \ar@{-->}[rru] & & S
}
$$
comme dans l'énoncé du lemme. Ainsi, une flèche pointillée
s'insérant dans le diagramme comme indiqué. D'après
l'hypothèse d'unicité du lemme, la composition $\Spec(A') \to \Spec(C) \to X$
coïncide avec le morphisme donné $\Spec(A') \to \Spec(P) \to \Spec(B) \subset X$. Par
conséquent, la restriction du morphisme au spectre
de $C/\mathfrak m_{A'} = A''$ induit le morphisme donné
$\Spec(K'') = \Spec(A'/\mathfrak m_{A'}) \to \Spec(K) \to X$. Soit $x \in X$ l'image du point fermé
de $\Spec(A'') \to X$. L'image de l'homomorphisme d'anneaux induit
$\mathcal{O}_{X, x} \to A''$ est un sous-anneau local contenu dans
$K \subset K''$. Comme $A$ est maximal pour
la relation de domination dans $K$ et comme
$A \subset A''$, nous avons $A = K \cap A''$. Nous concluons
que $\mathcal{O}_{X, x} \to A''$ se factorise par $A \subset A''$. De cette manière,
nous obtenons la flèche désirée $\Spec(A) \to X$.

\medskip\noindent
Enfin, supposons que $f$ soit quasi-séparé. Alors $\Delta : X \to X \times_S X$ est quasi-compact.
Étant donné un diagramme solide
$$
\xymatrix{
\Spec(K) \ar[r] \ar[d] & U \ar[r]^h & X \ar[d]^\Delta \\
\Spec(A) \ar[rr] \ar@{-->}[rru] & & X \times_S X
}
$$
où $A$ est un anneau de valuation avec corps de fractions $K$,
il existe un unique flèche pointillée rendant le diagramme commutatif.
À savoir, la flèche horizontale inférieure est la même chose qu'une paire
de morphismes $\Spec(A) \to X$ qui peut servir de flèche pointillée dans le diagramme
du lemme. Ainsi, l'unicité requise montre que la flèche horizontale
inférieure se factorise par $\Delta$. Par conséquent,
nous pouvons appliquer le résultat que nous
venons de démontrer à $\Delta : X \to X \times_S X$ et $h : U \to X$ et conclure que $\Delta$
est universellement fermé. Clairement, cela signifie que $f$
est séparé.
\end{proof}

\begin{remark}
\label{morphisms-remark-check-val-on-open}
L'hypothèse sur l'unicité des flèches
pointillées dans le Lemme \ref{morphisms-lemma-refined-valuative-criterion-universally-closed} est nécessaire
(détails omis). Bien sûr, l'unicité est garantie
si $f$
est séparé (Schémas, Lemme \ref{schemes-lemma-separated-implies-valuative}).
\end{remark}

\begin{lemma}
\label{morphisms-lemma-morphism-defined-local-ring}
Soit $S$ un schéma. Soient $X$, $Y$ des schémas sur $S$.
Soient $s \in S$ et $x \in X$, $y \in Y$ des points au-dessus de $s$.
\begin{enumerate}
\item Soient $f, g : X \to Y$ des morphismes sur $S$
tels que
$f(x) = g(x) = y$ et $f^\sharp_x = g^\sharp_x : \mathcal{O}_{Y, y} \to \mathcal{O}_{X, x}$. Alors il existe un voisinage
ouvert $U \subset X$ contenant $f|_U = g|_U$
dans les cas suivants
\begin{enumerate}
\item $Y$ est localement de type fini sur $S$,
\item $X$ est intègre,
\item $X$ est localement noethérien, ou
\item $X$ est réduit avec un nombre fini de composantes irréductibles.
\end{enumerate}
\item Soit $\varphi : \mathcal{O}_{Y, y} \to \mathcal{O}_{X, x}$ un homomorphisme local
de $\mathcal{O}_{S, s}$-algèbres. Alors il existe un voisinage ouvert $U \subset X$
de $x$ et un morphisme $f : U \to Y$ envoyant $x$
vers $y$ avec $f^\sharp_x = \varphi$ dans les cas suivants
\begin{enumerate}
\item $Y$ est localement de présentation finie sur $S$,
\item $Y$ est localement de type fini et $X$ est intègre,
\item $Y$ est localement de type fini et $X$ est localement noethérien, ou
\item $Y$ est localement de type fini et $X$ est réduit avec un nombre
fini de composantes irréductibles.
\end{enumerate}
\end{enumerate}
\end{lemma}

\begin{proof}
Démonstration de (1). Nous pouvons remplacer $X$, $Y$, $S$ par des voisinages
ouverts affines appropriés de $x$, $y$, $s$ et réduire au problème algébrique
suivant : étant donné un anneau $R$, deux applications d'algèbres $R$ $\varphi, \psi : B \to A$
telles que
\begin{enumerate}
\item $R \to B$ est de type fini, ou $A$ est un domaine, ou $A$
est noethérien, ou $A$ est réduit et a un nombre fini d'idéaux premiers minimaux,
\item les deux applications $B \to A_\mathfrak p$ sont les mêmes pour
un certain premier $\mathfrak p \subset A$,
\end{enumerate}
montrez que $\varphi, \psi$ définit la même application $B \to A_g$
pour un $g \in A$, $g \not \in \mathfrak p$ approprié. Si $R \to B$ est de type
fini, soit $t_1, \ldots, t_m \in B$ des générateurs de $B$ comme $R$-algèbre.
Pour chaque $j$ nous
pouvons trouver $g_j \in A$, $g_j \not \in \mathfrak p$
tels que $\varphi(t_j)$ et $\psi(t_j)$ aient la même image dans
$A_{g_j}$. Alors nous définissons $g = \prod g_j$.
Dans les autres cas (si $A$ est un domaine, de
type de Noether ou réduit avec un nombre fini de premiers minimaux),
nous pouvons trouver un $g \in A$, $g \not \in \mathfrak p$ tels que
$A_g \subset A_\mathfrak p$. Voir Algèbre, Lemme \ref{algebra-lemma-subring-of-local-ring}. Ainsi les applications
$B \to A_g$ sont égales comme souhaité.

\medskip\noindent
Démonstration de (2). Pour ce faire, nous pouvons remplacer $X$, $Y$ et $S$
par des ouverts affines appropriés. Disons $X = \Spec(A)$,
$Y = \Spec(B)$ et $S = \Spec(R)$. Soit $\mathfrak p \subset A$ l'idéal premier correspondant
à $x$. Soit $\mathfrak q \subset B$ l'idéal premier correspondant à $y$.
Alors $\varphi$ est un homomorphisme local de $R$-algèbres
$\varphi : B_\mathfrak q \to A_\mathfrak p$. Si $R \to B$
est un homomorphisme d'anneaux de présentation finie, alors il existe
un $g \in A \setminus \mathfrak p$ et un homomorphisme de $R$-algèbres $B \to A_g$ tel que
$$
\xymatrix{
B_\mathfrak q \ar[r]_\varphi & A_\mathfrak p \\
B \ar[u] \ar[r] & A_g \ar[u]
}
$$
commute, voir
Algèbre, Lemmes \ref{algebra-lemma-characterize-finite-presentation} et \ref{algebra-lemma-localization-colimit}. Le morphisme induit
$\Spec(A_g) \to \Spec(B)$ fonctionne. Si $B$
est de type fini sur $R$, soit $t_1, \ldots, t_m \in B$
des générateurs de $B$ comme une $R$-algèbre.
Alors nous pouvons choisir $g_j \in A$, $g_j \not \in \mathfrak p$
de telle sorte que $\varphi(t_j) \in \Im(A_{g_j} \to A_\mathfrak p)$.
Ainsi, après avoir remplacé $A$ par $A[1/\prod g_j]$,
nous pouvons supposer que $B$ s'envoie
dans l'image de $A \to A_\mathfrak p$. Si nous pouvons
trouver un $g \in A$, $g \not \in \mathfrak p$ de telle
sorte que $A_g \to A_\mathfrak p$ soit injectif, alors
nous obtiendrons l'application de $R$-algèbre
désirée $B \to A_g$. Une nouvelle application du
Lemme \ref{algebra-lemma-subring-of-local-ring} d'Algèbre achève donc la démonstration.
\end{proof}

\begin{lemma}
\label{morphisms-lemma-extend-across}
Soit $S$ un schéma. Soient $X$, $Y$ des schémas sur $S$. Soit
$x \in X$. Soit $U \subset X$ un ouvert et soit $f : U \to Y$ un morphisme sur $S$.
Supposons
\begin{enumerate}
\item $x$ est dans l'adhérence de $U$,
\item $X$ est réduit avec un nombre fini de composantes irréductibles ou $X$
est noethérien,
\item $\mathcal{O}_{X, x}$ est un anneau de valuation,
\item $Y \to S$ est propre
\end{enumerate}
Alors il existe un ouvert $U \subset U' \subset X$ contenant $x$ et un
$S$-morphisme $f' : U' \to Y$ prolongeant $f$.
\end{lemma}

\begin{proof}
Il est sans danger de remplacer $X$ par un voisinage ouvert de
$x$ dans $X$ (petit détail
omis). Par les Propriétés, Lemme \ref{properties-lemma-ring-affine-open-injective-local-ring}, nous pouvons supposer
que $X$ est affine
avec $\Gamma(X, \mathcal{O}_X) \subset \mathcal{O}_{X, x}$. En particulier $X$ est intègre
avec un point générique unique $\xi$ dont le corps résiduel
est le corps des fractions $K$ de l'anneau
de valuation $\mathcal{O}_{X, x}$. Puisque
$x$ est dans l'adhérence de $U$, nous voyons
que $U$ n'est pas vide, donc $U$ contient $\xi$.
Ainsi, par le critère valuatif de propreté
(Lemme \ref{morphisms-lemma-characterize-proper}), il existe un morphisme $t : \Spec(\mathcal{O}_{X, x}) \to Y$
s'insérant dans un diagramme commutatif
$$
\xymatrix{
\Spec(K) \ar[d]_\xi \ar[r] & \Spec(\mathcal{O}_{X, x}) \ar[d]_t \\
U \ar[r]^f & Y
}
$$
de morphismes de schémas sur $S$. En
appliquant le Lemme \ref{morphisms-lemma-morphism-defined-local-ring} avec $y = t(x)$
et $\varphi = t^\sharp_x$, nous obtenons un voisinage ouvert $V \subset X$
de $x$ et un morphisme $g : V \to Y$ sur $S$ qui
envoie $x$ vers $y$ et tel que $g^\sharp_x = t^\sharp_x$. Comme $Y \to S$
est séparé, l'égaliseur $E$ de $f|_{U \cap V}$ et $g|_{U \cap V}$
est un sous-schéma fermé de $U \cap V$, voir Schémas, Lemme
\ref{schemes-lemma-where-are-they-equal}. Puisque $f$ et $g$ déterminent
le même morphisme $\Spec(K) \to Y$ par construction, nous voyons
que $E$ contient le point générique du schéma intègre
$U \cap V$. Ainsi $E = U \cap V$ et Nous concluons que $f$
et $g$ se recollent pour donner un morphisme $U' = U \cup V \to Y$
comme souhaité.
\end{proof}





\section{Morphismes projectifs}
\label{morphisms-section-projective}

\noindent
Nous utiliserons la définition d'un morphisme projectif de \cite{EGA}. La version
de la définition avec le ``H'' est celle de
\cite{H}. Les définitions résultantes sont différentes. Les deux sont utiles.

\begin{definition}
\label{morphisms-definition-projective}
Soit $f : X \to S$ un morphisme de schémas.
\begin{enumerate}
\item On dit que $f$ est {\it projectif}
si $X$ est isomorphe en tant que $S$-schéma à
un sous-schéma fermé d'un fibré projectif
$\mathbf{P}(\mathcal{E})$ pour un certain $\mathcal{O}_S$-module quasi-cohérent de type fini $\mathcal{E}$.
\item On dit que $f$ est {\it H-projectif} s'il existe un entier
$n$ et une immersion fermée $X \to \mathbf{P}^n_S$ sur $S$.
\item On dit que $f$ est {\it localement projectif}
s'il existe un recouvrement ouvert $S = \bigcup U_i$ tel que chaque $f^{-1}(U_i) \to U_i$ soit
projectif.
\end{enumerate}
\end{definition}

\noindent
Comme prévu, un morphisme projectif est quasi-projectif,
voir le Lemme \ref{morphisms-lemma-projective-quasi-projective}. Inversement,
les morphismes quasi-projectifs sont souvent des compositions
d'immersions ouvertes et de morphismes
projectifs, voir le Lemme \ref{morphisms-lemma-quasi-projective-open-projective}. Pour
un aperçu des propriétés des morphismes projectifs
sur une base quasi-projective,
voir Compléments sur les morphismes, section \ref{more-morphisms-section-projective}.

\begin{example}
\label{morphisms-example-projective}
Soit $S$ un schéma. Soit $\mathcal{A}$ une $\mathcal{O}_S$-algèbre
quasi-cohérente graduée générée par $\mathcal{A}_1$ sur $\mathcal{A}_0$.
Supposons de plus que $\mathcal{A}_1$ soit de type fini
sur $\mathcal{O}_S$. Posons $X = \underline{\text{Proj}}_S(\mathcal{A})$. Dans ce cas, $X \to S$
est projectif. À savoir, le morphisme
associé à l'application de $\mathcal{O}_S$-algèbres graduées
$$
\text{Sym}_{\mathcal{O}_X}^*(\mathcal{A}_1)
\longrightarrow
\mathcal{A}
$$
est une immersion
fermée,
voir Constructions, Lemme \ref{constructions-lemma-surjective-generated-degree-1-map-relative-proj}.
\end{example}
 
\begin{lemma}
\label{morphisms-lemma-H-projective}
Un H-morphisme projectif est H-quasi-projectif. Un
H-morphisme projectif est projectif.
\end{lemma}

\begin{proof}
La première affirmation est immédiate d'après les
définitions. La seconde tient car $\mathbf{P}^n_S$ est un fibré projectif
sur $S$, voir Constructions, Lemme \ref{constructions-lemma-projective-space-bundle}.
\end{proof}

\begin{lemma}
\label{morphisms-lemma-characterize-locally-projective}
Soit $f : X \to S$ un morphisme de schémas. Les conditions
suivantes sont équivalentes :
\begin{enumerate}
\item Le morphisme $f$ est localement projectif.
\item Il existe un recouvrement ouvert $S = \bigcup U_i$ tel
que chaque $f^{-1}(U_i) \to U_i$ soit H-projectif.
\end{enumerate}
\end{lemma}

\begin{proof}
D'après le Lemme \ref{morphisms-lemma-H-projective}, nous voyons que (2) implique (1). Supposons (1).
Pour chaque point $s \in S$ nous pouvons trouver $\Spec(R) = U \subset S$ un
voisinage ouvert affine de $s$ tel que $X_U$ soit isomorphe
à un sous-schéma fermé de $\mathbf{P}(\mathcal{E})$ pour un certain faisceau quasi-cohérent
de type fini de $\mathcal{O}_U$-modules $\mathcal{E}$. Écrivons $\mathcal{E} = \widetilde{M}$
pour un certain module $R$ de type fini $M$
(voir Propriétés,
Lemme \ref{properties-lemma-finite-type-module}). Choisissons des générateurs $x_0, \ldots, x_n \in M$
de $M$ en tant que $R$-module. Considérons l'application
surjective de $R$-algèbres graduées
$$
R[X_0, \ldots, X_n] \longrightarrow \text{Sym}_R(M).
$$
Selon Constructions,
Lemme \ref{constructions-lemma-surjective-graded-rings-map-proj} le morphisme
correspondant
$$
\mathbf{P}(\mathcal{E}) \to \mathbf{P}^n_R
$$
est une immersion fermée. Par conséquent, nous concluons que $f^{-1}(U)$ est isomorphe à
un sous-schéma fermé de $\mathbf{P}^n_U$ (en tant que schéma sur $U$). En d'autres
termes : (2) est vérifié.
\end{proof}

\begin{lemma}
\label{morphisms-lemma-locally-projective-proper}
Un morphisme localement projectif est propre.
\end{lemma}

\begin{proof}
Soit $f : X \to S$ localement projectif.
Pour montrer que $f$ est propre, nous
pouvons travailler localement sur la
base, voir le Lemme \ref{morphisms-lemma-proper-local-on-the-base}. Par conséquent,
d'après le Lemme \ref{morphisms-lemma-characterize-locally-projective} ci-dessus, nous pouvons supposer
qu'il existe une immersion
fermée $X \to \mathbf{P}^n_S$. D'après les Lemmes \ref{morphisms-lemma-composition-proper}
et \ref{morphisms-lemma-closed-immersion-proper}, il suffit de prouver
que $\mathbf{P}^n_S \to S$ est propre.
Puisque $\mathbf{P}^n_S \to S$ est le changement de base de
$\mathbf{P}^n_{\mathbf{Z}} \to \Spec(\mathbf{Z})$, il suffit de montrer que $\mathbf{P}^n_{\mathbf{Z}} \to \Spec(\mathbf{Z})$
est propre, voir le Lemme \ref{morphisms-lemma-base-change-proper}.
Par Constructions, Lemme \ref{constructions-lemma-proj-separated} le schéma $\mathbf{P}^n_{\mathbf{Z}}$
est séparé. Par Constructions,
Lemme \ref{constructions-lemma-proj-quasi-compact} le schéma $\mathbf{P}^n_{\mathbf{Z}}$ est quasi-compact.
Il est clair que $\mathbf{P}^n_{\mathbf{Z}} \to \Spec(\mathbf{Z})$
est localement de type fini puisque $\mathbf{P}^n_{\mathbf{Z}}$
est recouvert par les ouverts affines $D_{+}(X_i)$,
chacun d'eux étant le spectre de l'algèbre
de type fini $\mathbf{Z}$.
$$
\mathbf{Z}[X_0/X_i, \ldots, X_n/X_i].
$$
Enfin, nous devons
montrer que $\mathbf{P}^n_{\mathbf{Z}} \to \Spec(\mathbf{Z})$ est universellement
fermé. Cela
résulte de Constructions, Lemme \ref{constructions-lemma-proj-valuative-criterion}
et du critère valuatif
(voir Schémas, Proposition \ref{schemes-proposition-characterize-universally-closed}).
\end{proof}

\begin{lemma}
\label{morphisms-lemma-proper-ample-locally-projective}
Soit $f : X \to S$ un morphisme propre de schémas. S'il existe un $f$-faisceau
inversible ample sur $X$, alors $f$ est localement projectif.
\end{lemma}

\begin{proof}
S'il existe un $f$-faisceau inversible ample, alors
nous pouvons localement sur $S$ trouver une
immersion $i : X \to \mathbf{P}^n_S$, voir Lemme \ref{morphisms-lemma-finite-type-over-affine-ample-very-ample}. Comme $X \to S$
est propre, le morphisme $i$ est une immersion
fermée, voir Lemme \ref{morphisms-lemma-image-proper-scheme-closed}.
\end{proof}

\begin{lemma}
\label{morphisms-lemma-H-projective-composition}
Une composition de H-morphismes projectifs est H-projective.
\end{lemma}

\begin{proof}
Supposons que $X \to Y$ et $Y \to Z$ soient H-projectifs.
Alors il existe des immersions fermées $X \to \mathbf{P}^n_Y$ sur
$Y$, et $Y \to \mathbf{P}^m_Z$ sur $Z$. Considérons
le diagramme suivant
$$
\xymatrix{
X \ar[r] \ar[d] &
\mathbf{P}^n_Y \ar[r] \ar[dl] &
\mathbf{P}^n_{\mathbf{P}^m_Z} \ar[dl] \ar@{=}[r] &
\mathbf{P}^n_Z \times_Z \mathbf{P}^m_Z \ar[r] &
\mathbf{P}^{nm + n + m}_Z \ar[ddllll] \\
Y \ar[r] \ar[d] & \mathbf{P}^m_Z \ar[dl] & \\
Z & &
}
$$
Ici, la flèche horizontale supérieure la plus à droite est l'immersion
de Segre, voir Constructions, Lemme \ref{constructions-lemma-segre-embedding}. Le diagramme
identifie $X$
comme un sous-schéma fermé de $\mathbf{P}^{nm + n + m}_Z$ comme souhaité.
\end{proof}

\begin{lemma}
\label{morphisms-lemma-H-projective-base-change}
Un changement de base d'un H-morphisme projectif est H-projectif.
\end{lemma}

\begin{proof}
C'est vrai parce que le changement de base d'un espace
projectif sur un schéma est un espace projectif, et le fait
que le changement de base d'une immersion fermée est
une immersion fermée, voir Schémas, Lemme \ref{schemes-lemma-base-change-immersion}.
\end{proof}

\begin{lemma}
\label{morphisms-lemma-base-change-projective}
Un changement de base d'un morphisme projectif (localement) est projectif (localement).
\end{lemma}

\begin{proof}
Ceci est vrai parce que le changement de base d'un fibré
projectif sur un schéma est un fibré projectif,
le tiré en arrière d'un $\mathcal{O}$-module de type
fini est de type fini (Modules, Lemme \ref{modules-lemma-pullback-finite-type})
et le fait que le changement
de base d'une immersion fermée est une immersion
fermée, voir Schémas, Lemme \ref{schemes-lemma-base-change-immersion}. Certains
détails sont omis.
\end{proof}

\begin{lemma}
\label{morphisms-lemma-projective-quasi-projective}
Un morphisme projectif est quasi-projectif.
\end{lemma}

\begin{proof}
Soit $f : X \to S$ un morphisme projectif. Choisissez une immersion fermée
$i : X \to \mathbf{P}(\mathcal{E})$ où $\mathcal{E}$ est un $\mathcal{O}_S$-module quasi-cohérent et
de type fini. Alors $\mathcal{L} = i^*\mathcal{O}_{\mathbf{P}(\mathcal{E})}(1)$ est
$f$-très ample. Puisque $f$ est propre (Lemme \ref{morphisms-lemma-locally-projective-proper}), il
est quasi-compact. Par conséquent, le Lemme \ref{morphisms-lemma-ample-very-ample} implique
que $\mathcal{L}$ est $f$-ample. Puisque $f$ est propre,
il est de type fini. Ainsi, nous avons vérifié que toutes les propriétés
définissant un quasi-projectif sont respectées et nous
avons gagné.
\end{proof}

\begin{lemma}
\label{morphisms-lemma-H-quasi-projective-open-H-projective}
Soit $f : X \to S$ un H-morphisme quasi-projectif. Alors
$f$ se factorise en $X \to X' \to S$ où $X \to X'$ est une immersion
ouverte et $X' \to S$ est H-projectif.
\end{lemma}

\begin{proof}
Par définition, nous pouvons factoriser $f$ comme
une immersion quasi-compacte $i : X \to \mathbf{P}^n_S$ suivie de la projection $\mathbf{P}^n_S \to S$.
D'après le Lemme \ref{morphisms-lemma-quasi-compact-immersion}, il existe un sous-schéma fermé
$X' \subset \mathbf{P}^n_S$ tel que $i$ se factorise par une immersion
ouverte $X \to X'$. Le lemme s'ensuit.
\end{proof}

\begin{lemma}
\label{morphisms-lemma-quasi-projective-open-projective}
Soit $f : X \to S$ un morphisme quasi-projectif avec $S$ quasi-compact
et quasi-séparé. Alors $f$ se factorise en $X \to X' \to S$ où $X \to X'$ est
une immersion ouverte et $X' \to S$ est projectif.
\end{lemma}

\begin{proof}
Soit $\mathcal{L}$ $f$-ample. Comme $f$ est de
type fini et $S$ est quasi-compact, $\mathcal{L}^{\otimes n}$ est $f$-très
ample pour un certain $n > 0$,
voir Lemme \ref{morphisms-lemma-finite-type-ample-very-ample}. Remplacez $\mathcal{L}$
par $\mathcal{L}^{\otimes n}$. Écrivez $\mathcal{F} = f_*\mathcal{L}$. C'est un
$\mathcal{O}_S$-module quasi-cohérent
d'après Schémas, Lemme \ref{schemes-lemma-push-forward-quasi-coherent} (les
morphismes quasi-projectifs sont quasi-compacts
et séparés, voir Lemme \ref{morphisms-lemma-quasi-projective-properties}). D'après Propriétés,
Lemme \ref{properties-lemma-directed-colimit-finite-presentation}, nous pouvons trouver un ensemble dirigé
$I$ et un système de modules
quasi-cohérents de type fini $\mathcal{O}_S$ $\mathcal{E}_i$
sur $I$ tels que $\mathcal{F} = \colim \mathcal{E}_i$. Considérez
les compositions
$\psi_i : f^*\mathcal{E}_i \to f^*\mathcal{F} \to \mathcal{L}$. Choisissez un recouvrement ouvert
affine fini $S = \bigcup_{j = 1, \ldots, m} V_j$. Pour chaque $j$, nous pouvons
choisir des sections
$$
s_{j, 0}, \ldots, s_{j, n_j} \in
\Gamma(f^{-1}(V_j), \mathcal{L}) = f_*\mathcal{L}(V_j) = \mathcal{F}(V_j)
$$
qui génèrent $\mathcal{L}$ sur $f^{-1}V_j$ et définissent une immersion
$$
f^{-1}V_j \longrightarrow \mathbf{P}^{n_j}_{V_j},
$$
voir Lemme \ref{morphisms-lemma-very-ample-finite-type-over-affine}. Choisissons $i$ de sorte
qu'il existe des sections $e_{j, t} \in \mathcal{E}_i(V_j)$ dont les images soient $s_{j, t}$
dans $\mathcal{F}$ pour tous $j = 1, \ldots, m$ et $t = 1, \ldots, n_j$. Alors
l'application $\psi_i$ est surjective puisque les
sections $f^*e_{j, t}$ ont la même image que les sections $s_{j, t}$
qui génèrent $\mathcal{L}|_{f^{-1}V_j}$. D'où nous obtenons un morphisme
$$
r_{\mathcal{L}, \psi_i} : X \longrightarrow \mathbf{P}(\mathcal{E}_i)
$$
sur $S$ tel que sur $V_j$ nous avons une factorisation
$$
f^{-1}V_j \to \mathbf{P}(\mathcal{E}_i)|_{V_j} \to \mathbf{P}^{n_j}_{V_j}
$$
de l'immersion donnée ci-dessus. Il s'ensuit que $r_{\mathcal{L}, \psi_i}|_{V_j}$ est
une immersion, voir Lemme \ref{morphisms-lemma-immersion-permanence}. Puisque
$S = \bigcup V_j$ Nous concluons que $r_{\mathcal{L}, \psi_i}$ est une
immersion.
Notez que $r_{\mathcal{L}, \psi_i}$ est quasi-compact car
$X \to S$ est quasi-compact et $\mathbf{P}(\mathcal{E}_i) \to S$ est séparé (voir
Schémas, Lemme \ref{schemes-lemma-quasi-compact-permanence}). Par le Lemme \ref{morphisms-lemma-quasi-compact-immersion}
il existe un sous-schéma fermé $X' \subset \mathbf{P}(\mathcal{E}_i)$ tel que $i$
se factorise par une immersion ouverte $X \to X'$.
Alors $X' \to S$ est projectif par définition
et on obtient le résultat.
\end{proof}

\begin{lemma}
\label{morphisms-lemma-projective-is-quasi-projective-proper}
Soit $S$ un schéma quasi-compact et quasi-séparé.
Soit $f : X \to S$ un morphisme de schémas. Alors
\begin{enumerate}
\item $f$ est projectif si et seulement si $f$ est quasi-projectif et propre, et
\item $f$ est H-projectif si et seulement si $f$ est H-quasi-projectif et propre.
\end{enumerate}
\end{lemma}

\begin{proof}
Si $f$ est projectif, alors $f$ est quasi-projectif
d'après le Lemme \ref{morphisms-lemma-projective-quasi-projective} et propre d'après le
Lemme \ref{morphisms-lemma-locally-projective-proper}. Réciproquement, si $X \to S$ est quasi-projectif
et propre, alors nous pouvons choisir une immersion
ouverte $X \to X'$ avec $X' \to S$ projectif d'après le
Lemme \ref{morphisms-lemma-quasi-projective-open-projective}. Puisque $X \to S$ est propre, nous
voyons que $X$ est fermé dans $X'$ (Lemme \ref{morphisms-lemma-image-proper-scheme-closed}),
c'est-à-dire que $X \to X'$ est une immersion
fermée (et ouverte). Comme $X'$ est
isomorphe à un sous-schéma fermé d'un fibré projectif sur $S$
(Définition \ref{morphisms-definition-projective}), nous voyons que la même
chose est vraie pour $X$, c'est-à-dire
$X \to S$ est un morphisme projectif. Cela
prouve (1). La démonstration de (2) est
la même, sauf qu'elle utilise
les Lemmes \ref{morphisms-lemma-H-projective} et \ref{morphisms-lemma-H-quasi-projective-open-H-projective}.
\end{proof}

\begin{lemma}
\label{morphisms-lemma-composition-projective}
Soient $f : X \to Y$ et $g : Y \to S$ des morphismes de schémas. Si $S$
est quasi-compact et quasi-séparé et que $f$ et $g$ sont projectifs,
alors $g \circ f$ est projectif.
\end{lemma}

\begin{proof}
D'après les Lemmes \ref{morphisms-lemma-projective-quasi-projective} et \ref{morphisms-lemma-locally-projective-proper},
nous voyons que $f$
et $g$ sont quasi-projectifs et
propres. D'après les Lemmes \ref{morphisms-lemma-composition-proper}
et \ref{morphisms-lemma-composition-quasi-projective}, nous voyons
que $g \circ f$ est propre et quasi-projectif.
Ainsi, $g \circ f$
est projectif d'après le Lemme \ref{morphisms-lemma-projective-is-quasi-projective-proper}.
\end{proof}

\begin{lemma}
\label{morphisms-lemma-projective-permanence}
Soient $g : Y \to S$ et $f : X \to Y$ des morphismes de schémas. Si $g \circ f$
est projectif et $g$ est séparé, alors $f$
est projectif.
\end{lemma}

\begin{proof}
Choisissez une immersion fermée $X \to \mathbf{P}(\mathcal{E})$ où $\mathcal{E}$ est un $\mathcal{O}_S$-module
quasi-cohérent de type fini. Alors nous
obtenons un morphisme $X \to \mathbf{P}(\mathcal{E}) \times_S Y$. Ce morphisme est
une immersion fermée parce qu'il est la composition
$$
X \to X \times_S Y \to \mathbf{P}(\mathcal{E}) \times_S Y
$$
où le premier morphisme est une immersion fermée
d'après Schémas, Lemme \ref{schemes-lemma-semi-diagonal} (et le
fait que $g$ est séparé) et
le second comme le changement de base d'une immersion
fermée. Enfin, le produit fibré $\mathbf{P}(\mathcal{E}) \times_S Y$ est isomorphe
à $\mathbf{P}(g^*\mathcal{E})$ et le tiré en arrière préserve les
modules quasi-cohérents de type fini.
\end{proof}

\begin{lemma}
\label{morphisms-lemma-projective-over-quasi-projective-is-H-projective}
Soit $S$ un schéma qui admet un faisceau inversible ample. Alors
\begin{enumerate}
\item tout morphisme projectif $X \to S$ est H-projectif, et
\item tout morphisme quasi-projectif $X \to S$ est H-quasi-projectif.
\end{enumerate}
\end{lemma}

\begin{proof}
Les hypothèses sur $S$ impliquent que $S$ est
quasi-compact et séparé, voir Propriétés, Définition
\ref{properties-definition-ample} et Lemme \ref{properties-lemma-ample-immersion-into-proj} et Constructions,
Lemme \ref{constructions-lemma-proj-separated}. Ainsi, le Lemme \ref{morphisms-lemma-quasi-projective-open-projective}
s'applique et nous voyons que (1) implique
(2). Soit $\mathcal{E}$ un $\mathcal{O}_S$-module
quasi-cohérent de type fini. D'après notre définition
de morphismes projectifs, il suffit de montrer que
$\mathbf{P}(\mathcal{E}) \to S$ est H-projectif. Si $\mathcal{E}$
est engendré par un nombre fini de sections globales,
alors la surjection correspondante $\mathcal{O}_S^{\oplus n} \to \mathcal{E}$ induit
une immersion fermée
$$
\mathbf{P}(\mathcal{E}) \longrightarrow
\mathbf{P}(\mathcal{O}_S^{\oplus n}) = \mathbf{P}^{n - 1}_S
$$
comme souhaité. En général, soit $\mathcal{L}$ un faisceau
inversible ample sur $S$. D'après Propriétés,
Proposition \ref{properties-proposition-characterize-ample},
il existe un entier $n$ tel que
$\mathcal{E} \otimes_{\mathcal{O}_S} \mathcal{L}^{\otimes n}$ soit globalement engendré
par un nombre
fini de sections. Puisque $\mathbf{P}(\mathcal{E}) =
\mathbf{P}(\mathcal{E} \otimes_{\mathcal{O}_S} \mathcal{L}^{\otimes n})$ d'après Constructions,
Lemme \ref{constructions-lemma-twisting-and-proj}, cela termine la démonstration.
\end{proof}

\begin{lemma}
\label{morphisms-lemma-proper-ample-is-proj}
Soit $f : X \to S$ un morphisme universellement fermé. Soit
$\mathcal{L}$ un $f$-ample faisceau inversible $\mathcal{O}_X$. Alors
le morphisme canonique
$$
r : X
\longrightarrow
\underline{\text{Proj}}_S
\left(
\bigoplus\nolimits_{d \geq 0} f_*\mathcal{L}^{\otimes d}
\right)
$$
du Lemme \ref{morphisms-lemma-characterize-relatively-ample} est un isomorphisme.
\end{lemma}

\begin{proof}
Remarquez que $f$ est quasi-compact parce que l'existence
d'un $f$-module inversible ample oblige $f$
à être quasi-compact. D'après le lemme cité, le morphisme
$r$ est une immersion ouverte. D'autre
part, l'image de $r$ est fermée
par le Lemme \ref{morphisms-lemma-image-proper-scheme-closed} (la cible de $r$ est séparée
au-dessus de $S$ d'après Constructions, Lemme
\ref{constructions-lemma-relative-proj-separated}). Enfin, l'image de $r$
est dense par Propriétés, Lemme \ref{properties-lemma-ample-immersion-into-proj} (ici nous utilisons
également qu'il a été montré dans la démonstration
du Lemme \ref{morphisms-lemma-characterize-relatively-ample} que le morphisme $r$
sur des ouverts affines de $S$
est donné par le morphisme canonique de
Propriétés, Lemme \ref{properties-lemma-map-into-proj}). Ainsi nous concluons que
$r$ est une immersion ouverte surjective, c'est-à-dire un
isomorphisme.
\end{proof}

\begin{lemma}
\label{morphisms-lemma-proper-ample-delete-affine}
Soit $f : X \to S$ un morphisme universellement fermé. Soit $\mathcal{L}$
un $f$-ample module inversible $\mathcal{O}_X$.
Soit $s \in \Gamma(X, \mathcal{L})$. Alors $X_s \to S$ est un morphisme affine.
\end{lemma}

\begin{proof}
La question est locale sur $S$ (Lemme \ref{morphisms-lemma-characterize-affine}) donc
nous pouvons supposer que $S$ est affine. D'après le Lemme
\ref{morphisms-lemma-proper-ample-is-proj}, nous pouvons écrire $X = \text{Proj}(A)$ où $A$
est un anneau gradué et $s$
correspond à $f \in A_1$ et $X_s = D_+(f)$ (Propriétés, Lemme \ref{properties-lemma-map-into-proj})
ce qui prouve le lemme par construction de $\text{Proj}(A)$,
voir Constructions, section \ref{constructions-section-proj}.
\end{proof}











\section{Morphismes entiers et finis}
\label{morphisms-section-integral}

\noindent
Rappelons qu'un homomorphisme d'anneaux $R \to A$ est dit {\it
entier} si chaque élément de $A$ satisfait à une équation
monique avec des coefficients dans $R$. Rappelez-vous qu'un homomorphisme
d'anneaux $R \to A$ est dit {\it fini} si
$A$ est fini en tant que $R$-module. Voir Algèbre, Définition \ref{algebra-definition-integral-ring-map}.

\begin{definition}
\label{morphisms-definition-integral}
Soit $f : X \to S$ un morphisme de schémas.
\begin{enumerate}
\item On dit que $f$ est {\it entier}
si $f$ est affine et si pour tout ouvert affine
$\Spec(R) = V \subset S$ avec image réciproque $\Spec(A) = f^{-1}(V) \subset X$, l'homomorphisme
d'anneaux associé $R \to A$ est entier.
\item On dit que $f$ est {\it fini}
si $f$ est affine et si pour tout ouvert affine
$\Spec(R) = V \subset S$ avec image réciproque $\Spec(A) = f^{-1}(V) \subset X$ l'homomorphisme
d'anneaux associé $R \to A$ est fini.
\end{enumerate}
\end{definition}

\noindent
Il est clair que les morphismes entiers/finis sont séparés
et quasi-compacts. Il est également clair qu'un morphisme
fini est un morphisme de type fini. La plupart des lemmes de
cette section sont complètement
standards. Mais notez le lemme amusant \ref{morphisms-lemma-integral-universally-closed}
à la fin de la section.

\begin{lemma}
\label{morphisms-lemma-integral-local}
Soit $f : X \to S$ un morphisme de schémas. Les conditions
suivantes sont équivalentes :
\begin{enumerate}
\item Le morphisme $f$ est entier.
\item Il existe un recouvrement ouvert affine $S = \bigcup U_i$
tel que chaque $f^{-1}(U_i)$
soit affine et $\mathcal{O}_S(U_i) \to \mathcal{O}_X(f^{-1}(U_i))$ soit entier.
\item Il existe un recouvrement ouvert $S = \bigcup U_i$
tel que chaque $f^{-1}(U_i) \to U_i$ soit entier.
\end{enumerate}
De plus, si $f$ est entier alors pour tout sous-schéma
ouvert $U \subset S$ le morphisme $f : f^{-1}(U) \to U$ est entier.
\end{lemma}

\begin{proof}
Voir Algèbre, Lemme \ref{algebra-lemma-integral-local}. Certains détails
sont omis.
\end{proof}

\begin{lemma}
\label{morphisms-lemma-finite-local}
Soit $f : X \to S$ un morphisme de schémas. Les conditions
suivantes sont équivalentes :
\begin{enumerate}
\item Le morphisme $f$ est fini.
\item Il existe un recouvrement ouvert affine $S = \bigcup U_i$
tel que chaque $f^{-1}(U_i)$
soit affine et $\mathcal{O}_S(U_i) \to \mathcal{O}_X(f^{-1}(U_i))$ soit fini.
\item Il existe un recouvrement ouvert $S = \bigcup U_i$
tel que chaque $f^{-1}(U_i) \to U_i$ soit fini.
\end{enumerate}
De plus, si $f$ est fini alors pour chaque sous-schéma
ouvert $U \subset S$ le morphisme $f : f^{-1}(U) \to U$ est fini.
\end{lemma}

\begin{proof}
Voir Algèbre, Lemme \ref{algebra-lemma-integral-local}. Certains détails
sont omis.
\end{proof}

\begin{lemma}
\label{morphisms-lemma-finite-integral}
Un morphisme fini est entier.
Un morphisme entier qui est localement de type fini est fini.
\end{lemma}

\begin{proof}
Voir Algèbre, Lemme \ref{algebra-lemma-finite-is-integral}
et Lemme \ref{algebra-lemma-characterize-finite-in-terms-of-integral}.
\end{proof}

\begin{lemma}
\label{morphisms-lemma-composition-finite}
Une composition de morphismes finis est finie. Il en
est de même pour les morphismes entiers.
\end{lemma}

\begin{proof}
Voir Algèbre, Lemme \ref{algebra-lemma-finite-transitive}
et \ref{algebra-lemma-integral-transitive}.
\end{proof}

\begin{lemma}
\label{morphisms-lemma-base-change-finite}
Un changement de base d'un morphisme fini est fini. Il
en est de même pour les morphismes entiers.
\end{lemma}

\begin{proof}
Voir Algèbre, Lemme \ref{algebra-lemma-base-change-integral}.
\end{proof}

\begin{lemma}
\label{morphisms-lemma-integral-universally-closed}
\begin{slogan}
entier $=$ affine $+$ universellement fermé
\end{slogan}
Soit $f : X \to S$ un morphisme de schémas. Les conditions
suivantes sont équivalentes
\begin{enumerate}
\item $f$ est entier, et
\item $f$ est affine et universellement fermé.
\end{enumerate}
\end{lemma}

\begin{proof}
Supposons (1). Un morphisme entier est affine par
définition. Un changement de base d'un morphisme entier est entier,
donc pour prouver (2) il suffit de montrer
qu'un morphisme entier est fermé. Cela résulte de Algèbre,
Lemmes \ref{algebra-lemma-integral-going-up} et \ref{algebra-lemma-going-up-closed}.

\medskip\noindent Supposons
(2). Nous pouvons supposer $f$ est
le morphisme $f : \Spec(A) \to \Spec(R)$ provenant d'un homomorphisme
d'anneaux $R \to A$. Soit $a$ un élément de $A$.
Nous devons montrer que $a$ est entier sur
$R$, c'est-à-dire que dans le noyau $I$ de l'application
$R[x] \to A$ envoyant $x$
à $a$ il existe un polynôme unitaire. Considérons
l'anneau $B = A[x]/(ax -1)$ et soit $J$ le noyau de la composition
$R[x]\to A[x] \to B$. Si $f\in J$ il existe $q\in A[x]$
tel que $f = (ax-1)q$ dans $A[x]$ donc si $f = \sum_i f_ix^i$ et $q = \sum_iq_ix^i$,
pour tout $i \geq 0$ nous avons $f_i = aq_{i-1} - q_i$. Pour $n \geq \deg q + 1$ le polynôme
$$
\sum\nolimits_{i \geq 0} f_i x^{n - i} =
\sum\nolimits_{i \geq 0} (a q_{i - 1} - q_i) x^{n - i} =
(a - x) \sum\nolimits_{i \geq 0} q_i x^{n - i - 1}
$$
est clairement dans $I$ ; si $f_0 = 1$ ce polynôme est aussi unitaire, nous
sommes donc réduits à prouver que $J$ contient un polynôme avec un terme constant $1$.
Nous le faisons en prouvant que $\Spec(R[x]/(J + (x))$ est vide.


\medskip\noindent
Puisque $f$ est universellement fermé, le changement
de base $\Spec(A[x]) \to \Spec(R[x])$ est fermé. Par conséquent, l'image du sous-ensemble
fermé $\Spec(B) \subset \Spec(A[x])$ est le sous-ensemble
fermé $\Spec(R[x]/J) \subset \Spec(R[x])$, voir Exemple \ref{morphisms-example-scheme-theoretic-image}
et Lemme \ref{morphisms-lemma-quasi-compact-scheme-theoretic-image}. En particulier, $\Spec(B) \to \Spec(R[x]/J)$
est surjectif. Considérons le diagramme suivant
où chaque carré est un produit fibré :
$$
\xymatrix{
\Spec(B) \ar@{->>}[r]^g &
\Spec(R[x]/J)  \ar[r] &
\Spec(R[x])\\
\emptyset \ar[u] \ar[r] &
\Spec(R[x]/(J + (x)))\ar[u] \ar[r] &
\Spec(R) \ar[u]^0
}
$$
Le coin inférieur gauche est vide car il s'agit du spectre de $R\otimes_{R[x]} B$
où l'application $R[x]\to B$ envoie $x$ sur un élément inversible
et $R[x]\to R$ envoie $x$ sur $0$. Comme $g$
est surjectif, cela implique que $\Spec(R[x]/(J + (x)))$ est vide, comme nous voulions
le montrer.
\end{proof}

\begin{lemma}
\label{morphisms-lemma-integral-fibres}
Soit $f : X \to S$ un morphisme entier. Alors chaque
point de $X$ est fermé dans sa fibre.
\end{lemma}

\begin{proof}
Voir Algèbre, Lemme \ref{algebra-lemma-integral-no-inclusion}.
\end{proof}

\begin{lemma}
\label{morphisms-lemma-integral-dimension}
Soit $f : X \to Y$ un morphisme entier. Alors $\dim(X) \leq \dim(Y)$. Si $f$
est surjectif alors $\dim(X) = \dim(Y)$.
\end{lemma}

\begin{proof}
Puisque la dimension de $X$ et $Y$ est le supremum
des dimensions des éléments d'un recouvrement ouvert affine,
nous pouvons supposer que $Y$
et $X$ sont affines. L'inégalité
découle de l'Algèbre, Lemme
\ref{algebra-lemma-integral-dim-up}. L'égalité découle ensuite de l'Algèbre,
Lemmes \ref{algebra-lemma-dimension-going-up} et \ref{algebra-lemma-integral-going-up}.
\end{proof}

\begin{lemma}
\label{morphisms-lemma-finite-quasi-finite}
Un morphisme fini est quasi-fini.
\end{lemma}

\begin{proof}
Ceci est impliqué par l'Algèbre, Lemme \ref{algebra-lemma-quasi-finite} et Lemme
\ref{morphisms-lemma-quasi-finite-locally-quasi-compact}. Alternativement, tous les points
des fibres sont des points fermés par le Lemme
\ref{morphisms-lemma-integral-fibres} (et le fait qu'un morphisme fini est entier)
et utiliser le Lemme
\ref{morphisms-lemma-quasi-finite-at-point-characterize} (3) pour voir que $f$ est quasi-fini
en $x$ pour tous $x \in X$.
\end{proof}

\begin{lemma}
\label{morphisms-lemma-finite-proper}
Soit $f : X \to S$ un morphisme de schémas. Les conditions suivantes sont équivalentes
\begin{enumerate}
\item $f$ est fini, et
\item $f$ est affine et propre.
\end{enumerate}
\end{lemma}

\begin{proof}
Cela découle formellement
du Lemme \ref{morphisms-lemma-integral-universally-closed}, du fait qu'un morphisme
fini est entier et séparé, du fait qu'un
morphisme propre est la même chose qu'un morphisme
de type fini, séparé, universellement
fermé, et du fait qu'un morphisme entier de type
fini est fini (Lemme \ref{morphisms-lemma-finite-integral}).
\end{proof}

\begin{lemma}
\label{morphisms-lemma-closed-immersion-finite}
Une immersion fermée est finie (et a fortiori entière).
\end{lemma}

\begin{proof}
Vrai car une immersion fermée est affine
(Lemme \ref{morphisms-lemma-closed-immersion-affine}) et un homomorphisme
surjectif d'anneaux est fini et entier.
\end{proof}

\begin{lemma}
\label{morphisms-lemma-finite-union-finite}
Soient $X_i \to Y$, $i = 1, \ldots, n$ des morphismes finis de
schémas. Alors $X_1 \amalg \ldots \amalg X_n \to Y$ est fini aussi.
\end{lemma}

\begin{proof}
Cela découle du fait algébrique que si $R \to A_i$, $i = 1, \ldots, n$ sont
des homomorphismes finis d'anneaux, alors $R \to A_1 \times \ldots \times A_n$ est fini aussi.
\end{proof}

\begin{lemma}
\label{morphisms-lemma-finite-permanence}
Soient $f : X \to Y$ et $g : Y \to Z$ des morphismes.
\begin{enumerate}
\item Si $g \circ f$ est fini et $g$ séparé, alors $f$ est fini.
\item Si $g \circ f$ est entier et $g$ séparé, alors $f$ est entier.
\end{enumerate}
\end{lemma}

\begin{proof}
Supposons que $g \circ f$ soit fini (resp.\ entier) et que $g$ soit
séparé. Le changement de base $X \times_Z Y \to Y$ est fini (resp.\
entier) d'après le lemme
\ref{morphisms-lemma-base-change-finite}. Le morphisme $X \to X \times_Z Y$
est une immersion fermée car $Y \to Z$ est séparé,
voir Schémas, Lemme \ref{schemes-lemma-section-immersion}. Une immersion
fermée est finie (resp. \  entière),
voir Lemme \ref{morphisms-lemma-closed-immersion-finite}. La composition de morphismes
finis (resp. \  entiers) est finie (resp.
\  entière),
voir Lemme \ref{morphisms-lemma-composition-finite}. Ainsi nous avons gagné.
\end{proof}

\begin{lemma}
\label{morphisms-lemma-finite-monomorphism-closed}
Soit $f : X \to Y$ un morphisme de schémas. Si $f$
est fini et un monomorphisme, alors $f$ est une immersion fermée.
\end{lemma}

\begin{proof}
Cela se réduit
à Algèbre, Lemme \ref{algebra-lemma-finite-epimorphism-surjective}.
\end{proof}

\begin{lemma}
\label{morphisms-lemma-finite-projective}
Un morphisme fini est projectif.
\end{lemma}

\begin{proof}
Soit $f : X \to S$ un morphisme fini. Alors $f_*\mathcal{O}_X$ est
un $\mathcal{O}_S$-module quasi-cohérent
(Lemme \ref{morphisms-lemma-affine-equivalence-algebras}) de type fini (d'après
notre
définition des morphismes finis et
Propriétés, Lemme \ref{properties-lemma-finite-type-module}). Nous affirmons qu'il
existe une immersion fermée
$$
\sigma :
X
\longrightarrow
\mathbf{P}(f_*\mathcal{O}_X) =
\underline{\text{Proj}}_S(\text{Sym}^*_{\mathcal{O}_S}(f_*\mathcal{O}_X))
$$
sur $S$, ce qui termine
la démonstration. À savoir, nous prenons $\sigma$ comme étant
le morphisme qui correspond (via Constructions, Lemme \ref{constructions-lemma-apply-relative})
à la surjection
$$
f^*f_*\mathcal{O}_X \longrightarrow \mathcal{O}_X
$$
provenant de l'application d'adjonction $f^*f_* \to \text{id}$. Alors $\sigma$
est une immersion fermée. En effet, localement sur les
affines de $S$, nous pouvons écrire $X = \Spec(A)$ et $S = \Spec(R)$.
Puisque $X$ est fini sur $S$ nous pouvons choisir
$a_1, \ldots, a_n \in A$ générant $A$ en tant que $R$-module.
Alors l'homomorphisme de $R$-algèbres $R[T_0, T_1, \ldots, T_n] \to \text{Sym}_R^*(A)$ envoyant $T_0$
sur $1$ et $T_i$ sur $a_i$ pour $1 \leq i \leq n$
est surjective. D'où $\mathbf{P}(f_*\mathcal{O}_X)$ est
un sous-schéma fermé de $\mathbf{P}^n_S$ selon Constructions,
Lemme \ref{constructions-lemma-surjective-graded-rings-map-proj}. Ainsi, il suffit de prouver
que le morphisme induit $X \to \mathbf{P}^n_S$ est une
immersion fermée (voir par exemple, Lemme \ref{morphisms-lemma-closed-immersion-ideals}).
Le lecteur vérifie que $X \to \mathbf{P}^n_S$ a une image contenue
dans le $D_+(T_0) \cong \mathbf{A}^n_S$ ouvert et que $X \to \mathbf{A}^n_S$ correspond à l'application
d'algèbre $R$ surjective
$R[x_1, \ldots, x_n] \to \text{Sym}_R^*(A)$ envoyant $x_i$ vers $a_i$.
D'où $X \to \mathbf{P}^n_S$ est une immersion (comme une composition
d'une immersion fermée et d'une
immersion ouverte). Comme $X$ est fini sur
$S$, l'image de cette immersion
est fermée (cela utilise le Lemme
\ref{morphisms-lemma-finite-proper}, le Lemme \ref{morphisms-lemma-image-proper-scheme-closed} et Constructions, Lemme \ref{constructions-lemma-projective-space-separated}).
Nous concluons par Schémas, Lemme \ref{schemes-lemma-immersion-when-closed}.
\end{proof}







\section{Homéomorphismes universels}
\label{morphisms-section-universal-homeomorphisms}

\noindent
La définition suivante est vraiment superflue puisque un homéomorphisme
universel n'est vraiment qu'un morphisme entier,
universellement injectif
et surjectif, voir le Lemme \ref{morphisms-lemma-universal-homeomorphism}.

\begin{definition}
\label{morphisms-definition-universal-homeomorphism}
Un morphisme $f : X \to Y$ de schémas est appelé un {\it homéomorphisme universel}
si le changement de base $f' : Y' \times_Y X \to Y'$ est un homéomorphisme pour
tout morphisme $Y' \to Y$.
\end{definition}

\noindent
Tout d'abord, nous énonçons les lemmes obligatoires.

\begin{lemma}
\label{morphisms-lemma-base-change-universal-homeomorphism}
Le changement de base d'un homéomorphisme universel de schémas
par tout morphisme de schémas est un homéomorphisme universel.
\end{lemma}

\begin{proof}
Cela découle immédiatement de la définition.
\end{proof}

\begin{lemma}
\label{morphisms-lemma-composition-universal-homeomorphism}
La composition d'une paire d'homéomorphismes universels de schémas
est un homéomorphisme universel.
\end{lemma}

\begin{proof}
Omis.
\end{proof}

\noindent
Le lemme simple suivant est la clé pour caractériser les homéomorphismes
universels.

\begin{lemma}
\label{morphisms-lemma-homeomorphism-affine}
Soit $f : X \to Y$ un morphisme de schémas. Si $f$ est un homéomorphisme sur un sous-ensemble
fermé de $Y$, alors $f$ est affine.
\end{lemma}

\begin{proof}
Soit $y \in Y$ un point. Si $y \not \in f(X)$, alors il existe un voisinage
affine de $y$ qui est disjoint de $f(X)$.
Si $y \in f(X)$, soit $x \in X$ l'unique point de $X$ s'envoyant
sur $y$. Soit $y \in V$ un voisinage
ouvert affine. Soit $U \subset X$ un voisinage ouvert affine de $x$
qui s'envoie dans $V$. Puisque $f(U) \subset V \cap f(X)$ est ouvert dans la
topologie induite par notre hypothèse
sur $f$, nous pouvons choisir un $h \in \Gamma(V, \mathcal{O}_Y)$
tel que $y \in D(h)$ et $D(h) \cap f(X) \subset f(U)$. Notons $h' \in \Gamma(U, \mathcal{O}_X)$ la restriction
de $f^\sharp(h)$ à $U$. Alors nous voyons que $D(h') \subset U$
est égal à $f^{-1}(D(h))$. En d'autres termes, chaque point de $Y$
a un voisinage ouvert dont l'image réciproque est affine.
Ainsi $f$ est
affine, voir Lemme \ref{morphisms-lemma-characterize-affine}.
\end{proof}

\begin{lemma}
\label{morphisms-lemma-universal-homeomorphism}
Soit $f : X \to Y$ un morphisme de schémas. Les conditions suivantes sont équivalentes
:
\begin{enumerate}
\item $f$ est un homéomorphisme universel, et
\item $f$ est entier, universellement injectif et surjectif.
\end{enumerate}
\end{lemma}

\begin{proof}
Supposons que $f$ soit un homéomorphisme universel.
D'après le Lemme \ref{morphisms-lemma-homeomorphism-affine}, nous
voyons que $f$ est affine. Puisque $f$ est clairement universellement
fermé, nous voyons que $f$
est entier d'après le Lemme \ref{morphisms-lemma-integral-universally-closed}. Il
est aussi clair que $f$ est universellement injectif et surjectif.

\medskip\noindent
Supposons que $f$ soit entier, universellement injectif et
surjectif. D'après le Lemme \ref{morphisms-lemma-integral-universally-closed}, $f$
est universellement fermé. Puisqu'il est aussi universellement bijectif
(voir le Lemme \ref{morphisms-lemma-base-change-surjective}), nous voyons
que c'est un homéomorphisme universel.
\end{proof}

\begin{lemma}
\label{morphisms-lemma-reduction-universal-homeomorphism}
Soit $X$ un schéma. L'immersion fermée canonique $X_{red} \to X$
(voir Schémas, Définition \ref{schemes-definition-reduced-induced-scheme}) est un homéomorphisme
universel.
\end{lemma}

\begin{proof}
Omis.
\end{proof}

\begin{lemma}
\label{morphisms-lemma-check-closed-infinitesimally}
Soient $f : X \to S$ et $S' \to S$ des morphismes de schémas.
Supposons
\begin{enumerate}
\item $S' \to S$ est une immersion fermée,
\item $S' \to S$ est bijective sur les points,
\item $X \times_S S' \to S'$ est une immersion fermée, et
\item $X \to S$ est de type fini ou $S' \to S$ est de présentation finie.
\end{enumerate}
Alors $f : X \to S$ est une immersion fermée.
\end{lemma}

\begin{proof}
Les hypothèses (1) et (2) impliquent que $S' \to S$ est un homéomorphisme universel
(par exemple parce que $S_{red} = S'_{red}$ et en utilisant
le Lemme \ref{morphisms-lemma-reduction-universal-homeomorphism}). Par conséquent, (3) implique
que $X \to S$ est un homéomorphisme sur un sous-ensemble
fermé de $S$. Ensuite, $X \to S$
est affine d'après le Lemme \ref{morphisms-lemma-homeomorphism-affine}.
Soit $U \subset S$ un ouvert affine, disons $U = \Spec(A)$. Alors $S' = \Spec(A/I)$ par (1) pour
un idéal nilpotent localement $I$ par (2). Comme $f$ est affine,
on voit que $f^{-1}(U) = \Spec(B)$.
L'hypothèse (4) nous dit que $B$ est une $A$-algèbre
de type fini (Lemme \ref{morphisms-lemma-locally-finite-type-characterize}) ou que
$I$ est de type fini
(Lemme \ref{morphisms-lemma-closed-immersion-finite-presentation}). L'hypothèse (3) est que $A/I \to B/IB$
est surjectif. D'après l'Algèbre, Lemme
\ref{algebra-lemma-surjective-mod-locally-nilpotent} si $A \to B$ est de type fini ou Algèbre,
Lemme \ref{algebra-lemma-NAK} si $I$
est de type fini et donc nilpotent, nous déduisons
que $A \to B$ est surjectif. Cela signifie
que $f$ est une immersion fermée,
voir Lemme \ref{morphisms-lemma-closed-immersion}.
\end{proof}

\begin{lemma}
\label{morphisms-lemma-permanence-universal-homeomorphism}
Soit $f : X \to Z$ la composition de deux morphismes $g : X \to Y$
et $h : Y \to Z$. Si deux des
morphismes $\{f, g, h\}$ sont des homéomorphismes universels, le troisième
morphisme l'est aussi.
\end{lemma}

\begin{proof}
Si $g$ et $h$ sont tous deux des homéomorphismes universels,
$f$ l'est aussi d'après le Lemme \ref{morphisms-lemma-composition-universal-homeomorphism}.

\medskip\noindent Supposons
que $f$ et $g$ soient des homéomorphismes universels. Nous voulons montrer
que $h$ l'est aussi. Maintenant, faisons un changement de base du diagramme le long d'un
morphisme arbitraire $\alpha : Z' \to Z$ de schémas, nous obtenons le diagramme suivant avec tous les
carrés cartésiens :
$$
\xymatrix{
X' \ar[r]^{g'} \ar[d] & Y' \ar[r]^{h'} \ar[d] & Z' \ar[d] \\
X \ar[r]^{g} & Y \ar[r]^{h} & Z.
}
$$
Notre hypothèse implique que la composition $f'= h' \circ g' : X' \to Z'$ et $g' : X' \to Y'$ est un
homéomorphisme, donc $h'$ l'est également. Cela termine la
démonstration que $h$ est un homéomorphisme universel.

\medskip\noindent Enfin,
supposons que $f$ et $h$ soient
des homéomorphismes universels. Nous voulons montrer que $g$ est un homéomorphisme
universel. Soit $\beta : Y' \to Y$ un morphisme de schémas quelconque. Nous
obtenons le diagramme suivant avec tous les carrés cartésiens :
$$
\xymatrix{
X' \ar[r]^{g'} \ar[d] & Y' \ar[d]^{\gamma} & \\
X'' \ar[r]^{g''} \ar[d] & Y'' \ar[r]^{h''} \ar[d] & Y' \ar[d]^{h \circ \beta} \\
X \ar[r]^{g} & Y \ar[r]^{h} & Z.
}
$$
Ici, le morphisme $\gamma : Y' \to Y''$ est défini par la propriété universelle des produits
fibrés et les deux morphismes $id_{Y'} : Y' \to Y'$
et $\beta : Y' \to Y$. Nous allons prouver que
$g'$ est un homéomorphisme. Comme la
propriété d'être un homéomorphisme a la propriété 2-sur-3, nous voyons que
$g''$ est un homéomorphisme.
En regardant le carré du haut, il suffit de prouver que $\gamma$ est
un homéomorphisme universel. Comme $h''$ est un homéomorphisme,
nous voyons que c'est un morphisme affine d'après le
Lemme \ref{morphisms-lemma-homeomorphism-affine} et a fortiori séparé (Lemme \ref{morphisms-lemma-affine-separated}). Comme
$h'' \circ \gamma$ est l'identité, nous voyons que $\gamma$ est une
immersion fermée par
Schémas, Lemme \ref{schemes-lemma-section-immersion}. Puisque $h''$ est
bijectif, il s'ensuit que $\gamma$ est une immersion
fermée bijective et donc un homéomorphisme universel (par exemple
par la caractérisation dans le Lemme \ref{morphisms-lemma-universal-homeomorphism}) comme désiré.
\end{proof}




\section{Homéomorphismes universels de schémas affines}
\label{morphisms-section-universal-homeomorphisms-affine}

\noindent
Dans cette section, nous caractérisons les homéomorphismes universels des schémas affines.

\begin{lemma}
\label{morphisms-lemma-subalgebra-inherits}
Soit $A \to B$ un homomorphisme d'anneaux tel que le morphisme de schémas
induit $f : \Spec(B) \to \Spec(A)$ soit un homéomorphisme universel, resp. \ 
un homéomorphisme universel induisant des isomorphismes sur les corps résiduels, resp.
\  universellement
fermé, resp. \  universellement fermé et universellement
injectif. Alors, pour toute $A$-sous-algèbre
$B' \subset B$, la même chose est vraie pour $f' : \Spec(B') \to \Spec(A)$.
\end{lemma}

\begin{proof}
Si $f$ est universellement fermé, alors
$B$ est entier sur $A$ d'après le Lemme \ref{morphisms-lemma-integral-universally-closed}. Par conséquent, $B'$
est entier sur $A$ et $f'$
est universellement fermé (d'après le même lemme).
Cela prouve le cas où $f$ est universellement fermé.

\medskip\noindent
En continuant, nous voyons que $B$
est entier sur $B'$ (Algèbre, Lemme \ref{algebra-lemma-integral-permanence})
ce qui implique que $\Spec(B) \to \Spec(B')$ est surjectif
(Algèbre, Lemme \ref{algebra-lemma-integral-overring-surjective}). Ainsi, si $A \to B$ induit
des extensions purement inséparables des corps résiduels, alors il
en est de même pour $A \to B'$. Cela prouve le cas
où $f$ est universellement fermé et universellement
injectif, voir Lemme \ref{morphisms-lemma-universally-injective}.

\medskip\noindent
Le cas où $f$ est un homéomorphisme universel découle des
remarques ci-dessus, du Lemme \ref{morphisms-lemma-universal-homeomorphism}, et de l'observation
évidente que si $f$ est surjectif, alors $f'$ l'est aussi.

\medskip\noindent Si $A \to B$
induit des isomorphismes sur les corps résiduels, alors il en
est de même pour $A \to B'$ (voir l'argument au deuxième paragraphe).
De cette façon, nous voyons que le lemme est valable dans le cas restant.
\end{proof}

\begin{lemma}
\label{morphisms-lemma-colimit-inherits}
Soit $A$ un anneau. Soit $B = \colim B_\lambda$ une colimite filtrante d'$A$-algèbres.
Si chaque $f_\lambda : \Spec(B_\lambda) \to \Spec(A)$ est un homéomorphisme universel, respectivement
\  un homéomorphisme
universel induisant des isomorphismes sur les corps résiduels, respectivement
\  universellement
fermé, respectivement \  universellement fermé et
universellement injectif, alors il en va de même pour $f : \Spec(B) \to \Spec(A)$.
\end{lemma}

\begin{proof}
Si $f_\lambda$ est universellement fermé, alors $B_\lambda$
est entier sur $A$ d'après
le Lemme \ref{morphisms-lemma-integral-universally-closed}. Par conséquent, $B$
est entier sur $A$ et $f$
est universellement fermé (d'après le même lemme).
Cela prouve le cas où chaque $f_\lambda$ est universellement fermé.

\medskip\noindent
Pour un idéal premier $\mathfrak q \subset B$ au-dessus de $\mathfrak p \subset A$,
notons $\mathfrak q_\lambda \subset B_\lambda$ l'image réciproque. Alors $\kappa(\mathfrak q) = \colim \kappa(\mathfrak q_\lambda)$.
Ainsi, si $A \to B_\lambda$ induit des extensions
purement inséparables des corps résiduels, il
en est de même pour $A \to B$. Cela prouve le cas où $f_\lambda$
est universellement fermé et universellement injectif,
voir Lemme \ref{morphisms-lemma-universally-injective}.

\medskip\noindent
Le cas où $f$ est un homéomorphisme universel découle
des remarques ci-dessus et du Lemme \ref{morphisms-lemma-universal-homeomorphism} combiné au fait
que les idéaux premiers dans $B$ sont la même chose que les
suites compatibles d'idéaux premiers dans tous les $B_\lambda$.

\medskip\noindent Si $A \to B_\lambda$
induit des isomorphismes sur les corps résiduels, alors il en est de même
pour $A \to B$ (voir l'argument dans le deuxième paragraphe).
De cette manière, nous voyons que le lemme est valable dans le cas restant.
\end{proof}

\begin{lemma}
\label{morphisms-lemma-special-elements-and-localization}
Soit $A \subset B$ une extension d'anneaux. Soit $S \subset A$ un
sous-ensemble multiplicatif. Soient $n \geq 1$
et $b_i \in B$ pour $1 \leq i \leq n$. Tout $x \in S^{-1}B$ tel que
$$
x \not \in S^{-1}A\text{ et } b_i x^i \in S^{-1}A\text{ pour }i = 1, \ldots, n
$$
est égal à $s^{-1}y$ avec $s \in S$ et $y \in B$ tels que
$$
y \not \in A\text{ et } b_i y^i \in A\text{ pour }i = 1, \ldots, n
$$
\end{lemma}

\begin{proof}
Omis. Astuce : clarifier les dénominateurs.
\end{proof}

\begin{lemma}
\label{morphisms-lemma-nth-and-nplusone-implies-square-and-cube}
Soit $A \subset B$ une extension d'anneaux. S'il existe
$b \in B$, $b \not \in A$ et un entier $n \geq 2$
avec $b^n \in A$ et $b^{n + 1} \in A$, alors il existe
un $b' \in B$, $b' \not \in A$
avec $(b')^2 \in A$ et $(b')^3 \in A$.
\end{lemma}

\begin{proof}
Soient $b$ et $n$ comme dans
le lemme. Alors toutes les puissances suffisamment grandes
de $b$ sont dans $A$. À savoir, $(b^n)^k(b^{n + 1})^i = b^{(k + i)n + i}$
ce qui implique que toute puissance $b^m$ avec $m \geq n^2$
est dans $A$. Ainsi, si $i \geq 1$ est le plus
grand entier tel que $b^i \not \in A$, alors $(b^i)^2 \in A$ et $(b^i)^3 \in A$.
\end{proof}

\begin{lemma}
\label{morphisms-lemma-square-and-cube}
Soit $A \subset B$ une extension d'anneau telle que $\Spec(B) \to \Spec(A)$ soit un homéomorphisme
universel induisant des isomorphismes sur les corps résiduels.
Si $A \not = B$, alors il existe un $b \in B$, $b \not \in A$ avec
$b^2 \in A$ et $b^3 \in A$.
\end{lemma}

\begin{proof}
Rappelons que l'extension $A \subset B$ est entière
(Lemme \ref{morphisms-lemma-integral-universally-closed}). D'après le Lemme
\ref{morphisms-lemma-subalgebra-inherits}, nous pouvons supposer
que $B$ est généré par un seul élément
sur $A$. Ainsi $B$
est fini sur $A$ (Algèbre, Lemme \ref{algebra-lemma-characterize-finite-in-terms-of-integral}). Par
conséquent, le support de $B/A$
en tant que module $A$
est fermé et non vide (Algèbre, Lemmes \ref{algebra-lemma-support-closed} et \ref{algebra-lemma-support-zero}).
Soit $\mathfrak p \subset A$ un idéal premier
minimal du support. Après avoir remplacé
$A \subset B$ par $A_\mathfrak p \subset B_\mathfrak p$ (permis par
Lemme \ref{morphisms-lemma-special-elements-and-localization}) nous pouvons supposer que $(A, \mathfrak m)$
est un anneau local, que $B$ est fini
sur $A$, et que $B/A$ a pour support $\{\mathfrak m\}$ en
tant que $A$-module. Puisque $B/A$
est un module fini, nous voyons que $I = \text{Ann}_A(B/A)$ satisfait $\mathfrak m = \sqrt{I}$
(Algèbre, Lemme \ref{algebra-lemma-support-closed}). Soit $\mathfrak m' \subset B$
l'idéal premier unique au-dessus de $\mathfrak m$. Comme
$\Spec(B) \to \Spec(A)$ est un homéomorphisme, nous constatons que $\mathfrak m' = \sqrt{IB}$.
Pour $f \in \mathfrak m'$ choisissons
$n \geq 1$ de telle sorte que $f^n \in IB$. Alors
aussi $f^{n + 1} \in IB$. Puisque $IB \subset A$
par notre choix de $I$ Nous concluons que
$f^n, f^{n + 1} \in A$. En utilisant
le Lemme \ref{morphisms-lemma-nth-and-nplusone-implies-square-and-cube}, nous concluons que notre lemme
est vrai si $\mathfrak m' \not \subset A$. Cependant, si $\mathfrak m' \subset A$,
alors $\mathfrak m' = \mathfrak m$ et nous concluons que $A = B$ puisque les
corps résiduels sont également isomorphes
par hypothèse. Cette contradiction termine la
démonstration.
\end{proof}

\begin{lemma}
\label{morphisms-lemma-pth-power-and-multiple}
Soit $A \subset B$ une extension d'anneau telle que $\Spec(B) \to \Spec(A)$ soit un homéomorphisme
universel. Si
$A \not = B$, alors soit il existe un $b \in B$, $b \not \in A$ avec $b^2 \in A$
et $b^3 \in A$, soit il existe un nombre premier $p$
et un $b \in B$, $b \not \in A$ avec $pb \in A$ et $b^p \in A$.
\end{lemma}

\begin{proof}
L'argument est presque exactement le même que dans la démonstration
du Lemme \ref{morphisms-lemma-square-and-cube}, mais nous écrivons tout pour être
sûrs que cela fonctionne.

\medskip\noindent
Rappelons que l'extension $A \subset B$ est
entière (Lemme \ref{morphisms-lemma-integral-universally-closed}). D'après le Lemme
\ref{morphisms-lemma-subalgebra-inherits}, nous pouvons supposer
que $B$ est engendré par un seul élément
sur $A$. Par conséquent,
$B$ est fini sur $A$ (Algèbre, Lemme \ref{algebra-lemma-characterize-finite-in-terms-of-integral}). Ainsi,
le support de $B/A$ en tant que
$A$-module est fermé et
non vide (Algèbre, Lemmes \ref{algebra-lemma-support-closed} et \ref{algebra-lemma-support-zero}). Soit $\mathfrak p \subset A$
un idéal premier minimal du support.
Après avoir remplacé $A \subset B$ par
$A_\mathfrak p \subset B_\mathfrak p$ (permis par le Lemme \ref{morphisms-lemma-special-elements-and-localization}) nous
pouvons supposer que $(A, \mathfrak m)$ est un anneau
local, que $B$ est fini sur $A$,
et que $B/A$ a pour support $\{\mathfrak m\}$ en tant que
$A$-module. Puisque $B/A$ est
un module fini, nous voyons que $I = \text{Ann}_A(B/A)$ satisfait $\mathfrak m = \sqrt{I}$
(Algèbre, Lemme \ref{algebra-lemma-support-closed}). Soit $\mathfrak m' \subset B$
l'idéal premier unique au-dessus de $\mathfrak m$.
Comme $\Spec(B) \to \Spec(A)$ est un homéomorphisme, nous trouvons que
$\mathfrak m' = \sqrt{IB}$. Pour $f \in \mathfrak m'$ choisissez
$n \geq 1$ de telle sorte que $f^n \in IB$.
Alors également $f^{n + 1} \in IB$.
Puisque $IB \subset A$ par notre choix de $I$ Nous
concluons que $f^n, f^{n + 1} \in A$.
En utilisant le Lemme \ref{morphisms-lemma-nth-and-nplusone-implies-square-and-cube} nous concluons
que notre lemme est vrai si $\mathfrak m' \not \subset A$. Si $\mathfrak m' \subset A$,
alors $\mathfrak m' = \mathfrak m$. Puisque $A \not = B$ nous concluons
que l'application $\kappa = A/\mathfrak m \to B/\mathfrak m' = \kappa'$ du corps
résiduel ne peut pas être un isomorphisme.
Par le Lemme \ref{morphisms-lemma-universally-injective} Nous concluons
que la caractéristique de $\kappa$ est un
nombre premier $p$ et que l'extension $\kappa'/\kappa$
est purement inséparable. Choisissez $b \in B$
dont l'image dans $\kappa'$ est un élément
non contenu dans $\kappa$ mais dont la puissance $p$
est dans $\kappa$. Alors $b \not \in A$, $b^p \in A$,
et $pb \in A$ (parce que $pb \in \mathfrak m' = \mathfrak m \subset A$) comme
souhaité.
\end{proof}

\begin{proposition}
\label{morphisms-proposition-universal-homeomorphism-equal-residue-fields}
Soit $A \subset B$ une extension d'anneau. Les conditions suivantes sont équivalentes
\begin{enumerate}
\item $\Spec(B) \to \Spec(A)$ est un homéomorphisme universel induisant des isomorphismes
sur les corps résiduels, et
\item chaque sous-ensemble fini $E \subset B$ est contenu dans une extension
$$
A[b_1, \ldots, b_n] \subset B
$$
telle que $b_i^2, b_i^3 \in A[b_1, \ldots, b_{i - 1}]$ pour $i = 1, \ldots, n$.
\end{enumerate}
\end{proposition}

\begin{proof}
Supposons (1). En utilisant la récurrence transfinie, nous construisons
pour chaque ordinal $\alpha$ un $A$-sous-algèbre $B_\alpha \subset B$
comme suit. Posons $B_0 = A$. Si $\alpha$ est un ordinal limite,
alors nous posons $B_\alpha = \colim_{\beta < \alpha} B_\beta$.
Si $\alpha = \beta + 1$, alors soit $B_\beta = B$ auquel cas nous
posons $B_\alpha = B_\beta$ ou $B_\beta \not = B$, auquel cas nous appliquons le Lemme
\ref{morphisms-lemma-square-and-cube} pour choisir un $b_\alpha \in B$, $b_\alpha \not \in B_\beta$
avec $b_\alpha^2, b_\alpha^3 \in B_\beta$ et nous posons $B_\alpha = B_\beta[b_\alpha] \subset B$.
Il est clair que $B = \colim B_\alpha$ (en fait
$B = B_\alpha$ pour un certain ordinal $\alpha$ comme on
le voit en regardant les cardinalités). Nous prouverons,
par induction transfinie, que (2) tient pour
$A \to B_\alpha$ pour chaque ordinal $\alpha$. Cela est
clair pour $\alpha = 0$. Supposons que l'énoncé soit
vrai pour chaque $\beta < \alpha$ et soit $E \subset B_\alpha$
un sous-ensemble fini.
Si $\alpha$ est un ordinal limite, alors
$B_\alpha = \bigcup_{\beta < \alpha} B_\beta$ et nous voyons que $E \subset B_\beta$ pour
un certain $\beta < \alpha$ ce qui prouve le résultat
dans ce cas. Si $\alpha = \beta + 1$, alors $B_\alpha = B_\beta[b_\alpha]$.
Ainsi tout $e \in E$ peut être écrit comme
un polynôme $e = \sum d_{e, i}b_\alpha^i$ avec $d_{e, i} \in B_\beta$.
Soit $D \subset B_\beta$ l'ensemble $D = \{d_{e, i}\} \cup \{b_\alpha^2, b_\alpha^3\}$. Par hypothèse
d'induction
il existe une $A$-sous-algèbre $A[b_1, \ldots, b_n] \subset B_\beta$ comme dans
l'énoncé du lemme contenant $D$.
Alors $A[b_1, \ldots, b_n, b_\alpha] \subset B_\alpha$ est une $A$-sous-algèbre
de $B_\alpha$ comme dans l'énoncé du lemme
contenant $E$.

\medskip\noindent
Supposons (2). Écrivons $B = \colim B_\lambda$ comme le colimite de ses $A$-sous-algèbres
finies.
D'après le Lemme \ref{morphisms-lemma-colimit-inherits}, il suffit de montrer que $\Spec(B_\lambda) \to \Spec(A)$
est un homéomorphisme universel induisant
des isomorphismes sur les corps résiduels. Les compositions de
morphismes universellement fermés sont universellement fermées et il
en va de même pour les morphismes qui induisent des isomorphismes sur les corps résiduels.
Ainsi, il suffit de montrer que si $A \subset B$ et $B$ sont
générés par un seul élément $b$ avec $b^2, b^3 \in A$, alors
(1) tient. Une telle extension est intégrale et donc $\Spec(B) \to \Spec(A)$ est universellement
fermé et surjectif.
(Lemme \ref{morphisms-lemma-integral-universally-closed} et Algèbre, Lemme \ref{algebra-lemma-integral-overring-surjective}).
Notez que $(b^2)^3 = (b^3)^2$ dans $A$. Pour
tout homomorphisme d'anneaux
$\varphi : A \to K$ vers un corps $K$, nous voyons qu'il
existe un $\lambda \in K$ avec $\varphi(b^2) = \lambda^2$ et $\varphi(b^3) = \lambda^3$.
À savoir, $\lambda = 0$ si $\varphi(b^2) = 0$ et $\lambda = \varphi(b^3)/\varphi(b^2)$ sinon.
Ainsi, $B \otimes_A K$ est un quotient de $K[x]/(x^2 - \lambda^2, x^3 - \lambda^3)$.
Cet anneau a exactement un premier avec corps résiduel
$K$. Cela implique que $\Spec(B) \to \Spec(A)$
est bijectif et induit des isomorphismes sur les
corps résiduels. Combiné avec la clôture universelle,
cela montre que (1) est vrai, voir lemmes \ref{morphisms-lemma-universal-homeomorphism}
et \ref{morphisms-lemma-universally-injective}.
\end{proof}

\begin{proposition}
\label{morphisms-proposition-universal-homeomorphism}
Soit $A \subset B$ une extension d'anneau. Les conditions suivantes sont équivalentes
\begin{enumerate}
\item $\Spec(B) \to \Spec(A)$ est un homéomorphisme universel, et
\item chaque sous-ensemble fini $E \subset B$ est contenu dans une extension
$$
A[b_1, \ldots, b_n] \subset B
$$
telle que pour $i = 1, \ldots, n$ nous avons
\begin{enumerate}
\item $b_i^2, b_i^3 \in A[b_1, \ldots, b_{i - 1}]$, ou
\item il existe un nombre premier
$p$ avec $pb_i, b_i^p \in A[b_1, \ldots, b_{i - 1}]$.
\end{enumerate}
\end{enumerate}
\end{proposition}

\begin{proof}
La démonstration est exactement la même
que la démonstration de la Proposition \ref{morphisms-proposition-universal-homeomorphism-equal-residue-fields} sauf pour
les changements suivants :
\begin{enumerate}
\item Utilisez le Lemme \ref{morphisms-lemma-pth-power-and-multiple} au lieu du Lemme
\ref{morphisms-lemma-square-and-cube} ce qui signifie que pour chaque ordinal
successeur $\alpha = \beta + 1$ nous avons
soit $b_\alpha^2, b_\alpha^3 \in B_\beta$ ou nous avons un nombre premier
$p$ et $pb_\alpha, b_\alpha^p \in B_\beta$.
\item Si $\alpha$ est un ordinal successeur,
alors prenez $D = \{d_{e, i}\} \cup \{b_\alpha^2, b_\alpha^3\}$ ou
prenez $D = \{d_{e, i}\} \cup \{pb_\alpha, b_\alpha^p\}$ selon le cas dans lequel
$\alpha$ se trouve.
\item Dans la démonstration de (2) $\Rightarrow$ (1), nous devons également considérer
le cas où $B$ est généré sur $A$ par un seul élément
$b$ avec $pb, b^p \in B$ pour un certain nombre premier $p$.
Ici, $A \subset B$ induit un homéomorphisme universel sur les spectres
par exemple d'après Algèbre, Lemme \ref{algebra-lemma-p-ring-map}.
\end{enumerate}
Cela termine la démonstration.
\end{proof}

\begin{lemma}
\label{morphisms-lemma-universal-homeo-iso-if-invert-p}
Soit $p$ un nombre premier. Soit $A \to B$ un homomorphisme d'anneaux
qui induit un isomorphisme $A[1/p] \to B[1/p]$ (par exemple si
$p$ est nilpotent dans $A$). Les conditions
suivantes sont équivalentes
\begin{enumerate}
\item $\Spec(B) \to \Spec(A)$ est un homéomorphisme universel, et
\item le noyau de $A \to B$ est un idéal localement nilpotent et pour tout $b \in B$
il existe une puissance $p$-ième $q$ avec $qb$ et $b^q$
dans l'image de $A \to B$.
\end{enumerate}
\end{lemma}

\begin{proof}
Si (2) est vrai, alors (1) est vrai par Algèbre, Lemme \ref{algebra-lemma-p-ring-map}. Supposons
(1). Alors le noyau de $A \to B$ consiste en des
éléments nilpotents par Algèbre, Lemme \ref{algebra-lemma-image-dense-generic-points}. Ainsi, nous
pouvons remplacer $A$ par l'image de $A \to B$ et
supposer que $A \subset B$. Par Algèbre, Lemme \ref{algebra-lemma-help-with-powers}
l'ensemble
$$
B' = \{b \in B \mid p^nb, b^{p^n} \in A\text{ pour un certain }n \geq 0\}
$$
est une $A$-sous-algèbre de $B$ (le fait qu'elle soit fermée
sous les produits est trivial). Nous
devons montrer $B' = B$. Sinon, d'après le Lemme
\ref{morphisms-lemma-pth-power-and-multiple}, il existe un $b \in B$, $b \not \in B'$ avec soit
$b^2, b^3 \in B'$, soit il existe un nombre premier $\ell$
avec $\ell b, b^\ell \in B'$. Nous allons
montrer que les deux cas mènent à une contradiction, prouvant ainsi le lemme.

\medskip\noindent Puisque
$A[1/p] = B[1/p]$, nous pouvons choisir une $p$-puissance $q$ telle que $qb \in A$.

\medskip\noindent Si
$b^2, b^3 \in B'$, alors aussi $b^q \in B'$. Par définition de $B'$, nous
trouvons que $(b^q)^{q'} \in A$ pour une certaine $p$-puissance
$q'$. Alors $qq'b, b^{qq'} \in A$, d'où $b \in B'$ ce qui est une contradiction.

\medskip\noindent
Supposons maintenant qu'il existe un nombre premier $\ell$ avec $\ell b, b^\ell \in B'$.
Si $\ell \not = p$ alors $\ell b \in B'$ et $qb \in A \subset B'$ impliquent $b \in B'$, une
contradiction. Ainsi $\ell = p$ et $b^p \in B'$ et nous obtenons
une contradiction exactement comme avant.
\end{proof}

\begin{lemma}
\label{morphisms-lemma-make-universal-homeo}
Soit $A$ un anneau. Soit $x, y \in A$.
\begin{enumerate}
\item Si $x^3 = y^2$ est dans $A$, alors $A \to B = A[t]/(t^2 - x, t^3 - y)$ induit des bijections
sur les corps résiduels et un homéomorphisme
universel sur les spectres.
\item S'il existe un nombre premier $p$ tel que $p^px = y^p$ dans $A$,
alors $A \to B = A[t]/(t^p - x, pt - y)$ induit un homéomorphisme universel
sur les spectres.
\end{enumerate}
\end{lemma}

\begin{proof}
Nous utiliserons le critère du Lemme \ref{morphisms-lemma-universal-homeomorphism} pour vérifier cela. Dans les deux
cas, l'homomorphisme d'anneaux est entier. Il suffit donc de montrer que, étant
donné un corps $k$ et un homomorphisme d'anneaux $\varphi : A \to k$,
l'algèbre $k$-$B \otimes_A k$ possède un idéal premier
unique dont le corps résiduel est égal à $k$ dans
le cas (1) et purement inséparable sur $k$
dans le cas (2). Voir le Lemme \ref{morphisms-lemma-universally-injective}.

\medskip\noindent
Dans le cas (1), posons $\lambda = 0$ si
$\varphi(x) = 0$ et posons $\lambda = \varphi(y)/\varphi(x)$ sinon.
Alors $B = k[t]/(t^2 - \lambda^2, t^3 - \lambda^2)$. Ainsi le résultat
est clair.

\medskip\noindent Dans
le cas (2), si la caractéristique de $k$ est $p$,
alors nous obtenons $\varphi(y) = 0$ et $B = k[t]/(t^p - \varphi(x))$ qui est une algèbre
locale artinienne $k$ dont le corps résiduel est soit
$k$, soit une extension purement inséparable de degré $p$
de $k$. Si la caractéristique de $k$ n'est pas
$p$, alors en posant $\lambda = \varphi(y)/p$ nous voyons $B = k[t]/(t - \lambda) = k$ et
nous concluons également.
\end{proof}

\begin{lemma}
\label{morphisms-lemma-universal-homeo-limit}
Soit $A \to B$ un homomorphisme d'anneaux.
\begin{enumerate}
\item Si $A \to B$ induit un homéomorphisme universel sur les spectres, alors
$B = \colim B_i$ est une colimite filtrante d'algèbres $A$ de présentation finie $B_i$
telle que $A \to B_i$ induit un homéomorphisme universel sur les spectres.
\item Si $A \to B$ induit des isomorphismes sur les corps résiduels et un
homéomorphisme universel sur les spectres, alors $B = \colim B_i$ est une colimite
filtrante d'algèbres $A$ de présentation finie $B_i$ telles que $A \to B_i$
induit des isomorphismes sur les corps résiduels et un homéomorphisme
universel sur les spectres.
\end{enumerate}
\end{lemma}

\begin{proof}
Démonstration de (1). Nous utiliserons le critère
de Algèbre, Lemme \ref{algebra-lemma-when-colimit}. Soit
$A \to C$ de présentation finie et soit $\varphi : C \to B$ un homomorphisme
d'$A$-algèbres. Soit $B' = \varphi(C) \subset B$ l'image.
Alors $A \to B'$ induit un homéomorphisme universel sur
les spectres par le Lemme \ref{morphisms-lemma-subalgebra-inherits}.
Par Algèbre, Lemme \ref{algebra-lemma-ring-colimit-fp} nous pouvons
écrire $B' = \colim_{i \in I} B_i$ avec $A \to B_i$
de présentation finie et des homomorphismes
de transition
surjectives. Par Algèbre, Lemme \ref{algebra-lemma-characterize-finite-presentation} nous pouvons choisir
un indice $0 \in I$
et une factorisation $C \to B_0 \to B'$ de l'application $C \to B'$. Nous affirmons
que $\Spec(B_i) \to \Spec(A)$ est un homéomorphisme universel pour $i$ suffisamment
grand. Cette affirmation termine la démonstration de (1).

\medskip\noindent
Démonstration de l'affirmation. D'après le Lemme \ref{morphisms-lemma-reduction-universal-homeomorphism} l'homomorphisme
d'anneaux $A_{red} \to B'_{red}$ induit un homéomorphisme
universel sur les spectres. Ainsi $A_{red} \subset B'_{red}$
par l'Algèbre, Lemme \ref{algebra-lemma-image-dense-generic-points}. En posant
$A' = \Im(A \to B')$ nous avons des surjections $A \to A' \to A_{red}$ induisant
des bijections $\Spec(A_{red}) = \Spec(A') = \Spec(A)$. Ainsi $A' \subset B'$ induit
un homéomorphisme universel sur les spectres. D'après
la Proposition \ref{morphisms-proposition-universal-homeomorphism} et le fait que $B'$
est de type fini sur $A'$, nous
pouvons trouver $n$ et $b'_1, \ldots, b'_n \in B'$
tels que $B' = A'[b'_1, \ldots, b'_n]$
et tel que pour $j = 1, \ldots, n$ nous avons
\begin{enumerate}
\item $(b'_j)^2, (b'_j)^3 \in A'[b'_1, \ldots, b'_{j - 1}]$, ou
\item il existe un nombre premier
$p$ avec $pb'_j, (b'_j)^p \in A'[b'_1, \ldots, b'_{j - 1}]$.
\end{enumerate}
Choisissez $b_1, \ldots, b_n \in B_0$ en relevant $b'_1, \ldots, b'_n$. Pour $i \geq 0$ notez
$b_{j, i}$ l'image de $b_j$ dans $B_i$. Pour un $i$
suffisamment grand nous aurons pour $j = 1, \ldots, n$
\begin{enumerate}
\item $b_{j, i}^2, b_{j, i}^3 \in A_i[b_{1, i}, \ldots, b_{j - 1, i}]$, ou
\item il existe un nombre
premier $p$ avec $pb_{j, i}, b_{j, i}^p \in A_i[b_{1, i}, \ldots, b_{j - 1, i}]$.
\end{enumerate}
Voici $A_i \subset B_i$, l'image de $A \to B_i$. Observez que $A \to A_i$ est un
homomorphisme surjectif d'anneaux dont le noyau est un idéal
localement nilpotent. En augmentant encore $i$ si
nécessaire, nous pouvons supposer que $B_i$ est engendré
par $b_1, \ldots, b_n$ sur $A_i$,
en d'autres termes $B_i = A_i[b_1, \ldots, b_n]$. D'après
l'Algèbre, les Lemmes \ref{algebra-lemma-p-ring-map} et \ref{algebra-lemma-2-3-ring-map}
Nous concluons que $A \to A_i \to A_i[b_1] \to \ldots \to A_i[b_1, \ldots, b_n] = B_i$ induit des homéomorphismes universels
sur les spectres. Cela termine la démonstration de
l'assertion.

\medskip\noindent La
démonstration de (2) est exactement la même.
\end{proof}







\section{Normalisation faible absolue et semi-normalisation}
\label{morphisms-section-seminormalization}

\noindent
Motivés par les résultats prouvés dans la section précédente,
nous donnons la définition suivante.

\begin{definition}
\label{morphisms-definition-seminormal-ring}
Soit $A$ un anneau.
\begin{enumerate}
\item On dit que $A$ est {\it semi-normal} si pour
tout $x, y \in A$ avec $x^3 = y^2$ il existe un unique $a \in A$
avec $x = a^2$ et $y = a^3$.
\item On dit que $A$ est {\it absolument faiblement
normal} si (a) $A$
est semi-normal et (b) pour tout nombre premier $p$ et $x, y \in A$ avec
$p^px = y^p$ il existe un unique $a \in A$ avec $x = a^p$ et $y = pa$.
\end{enumerate}
\end{definition}

\noindent
Une observation amusante, voir \cite{Costa},
est que dans la définition des anneaux semi-normaux, il suffit\footnote{Soit
$A$ un anneau tel que pour tout $x, y \in A$ avec
$x^3 = y^2$, il existe un $a \in A$ avec $x = a^2$ et $y = a^3$.
Alors $A$ est réduit : si $x^2 = 0$, alors $x^2 = x^3$
et donc il existe un $a$ tel que $x = a^3$ et $x = a^2$. Puis
$x = a^3 = ax = a^4 = x^2 = 0$. Enfin, si $a_1^2 = a_2^2$ et $a_1^3 = a_2^3$ pour $a_1, a_2$ dans un
anneau réduit, alors $(a_1 - a_2)^3 = a_1^3 - 3a_1^2a_2 +3a_1a_2^2 - a_2^3 =
(1 - 3 + 3 - 1)a_1^3 = 0$ et donc $a_1 = a_2$.}
d'assumer l'existence de $a$. Les schémas absolument faiblement normaux ont été définis dans \cite[Appendix
B]{rydh_descent}.

\begin{lemma}
\label{morphisms-lemma-seminormal-local-property}
Être semi-normal ou être absolument faiblement
normal est une propriété locale
des anneaux, voir Propriétés, Définition \ref{properties-definition-property-local}.
\end{lemma}

\begin{proof}
Supposons que $A$ soit semi-normal et $f \in A$. Soit $x', y' \in A_f$ avec
$(x')^3 = (y')^2$. Écrivons $x' = x/f^{2n}$ et $y' = y/f^{3n}$ pour certains $n \geq 0$ et
$x, y \in A$. Après avoir remplacé $x, y$ par $f^{2m}x, f^{3m}y$ et $n$
par $n + m$, nous voyons que $x^3 = y^2$ dans $A$.
Puis nous trouvons un unique $a \in A$ avec $x = a^2$ et
$y = a^3$. En posant $a' = a/f^n$, nous obtenons $x' = (a')^2$
et $y' = (a')^3$ comme souhaité. L'unicité de $a'$ découle
de l'unicité de $a$. De la même manière
exactement, le lecteur montre que si $A$ est absolument
faiblement normal, alors $A_f$
est absolument faiblement normal.

\medskip\noindent Supposons
que $A$ soit un anneau et que $f_1, \ldots, f_n \in A$ génère l'idéal unité.
Supposons que $A_{f_i}$ soit semi-normal pour chaque
$i$. Soit $x, y \in A$ avec $x^3 = y^2$. Pour chaque $i$ nous
trouvons un unique $a_i \in A_{f_i}$ avec $x = a_i^2$ et $y = a_i^3$ dans $A_{f_i}$.
Par unicité et le résultat du premier paragraphe (qui nous dit que
$A_{f_if_j}$ est semi-normal) nous voyons que $a_i$ et
$a_j$ s'envoient sur le même élément de $A_{f_if_j}$. Par l'algèbre,
Lemme \ref{algebra-lemma-standard-covering} nous trouvons un unique $a \in A$ se
projetant sur $a_i$ dans
$A_{f_i}$ pour tous les $i$. Ensuite $x = a^2$
et $y = a^3$ par le même principe. Clairement, ce $a$ est unique.
Ainsi $A$ est semi-normal. Si nous supposons que $A_{f_i}$ est absolument faiblement normal,
alors le même argument exact montre que $A$ est absolument faiblement normal.
\end{proof}

\noindent
Ensuite, nous définissons les schémas semi-normaux et les schémas absolument faiblement normaux.

\begin{definition}
\label{morphisms-definition-seminormal}
Soit $X$ un schéma.
\begin{enumerate}
\item Nous disons que $X$ est {\it semi-normal}
si chaque $x \in X$ possède un voisinage ouvert affine $\Spec(R) = U \subset X$
tel que l'anneau $R$ soit semi-normal.
\item Nous disons que $X$ est {\it absolument faiblement normal}
si chaque $x \in X$ possède un voisinage ouvert affine $\Spec(R) = U \subset X$
tel que l'anneau $R$ soit absolument faiblement normal.
\end{enumerate}
\end{definition}

\noindent
Voici le lemme obligatoire.

\begin{lemma}
\label{morphisms-lemma-locally-seminormal}
Soit $X$ un schéma. Les conditions suivantes sont équivalentes :
\begin{enumerate}
\item Le schéma $X$ est semi-normal.
\item Pour tout ouvert affine $U \subset X$, l'anneau $\mathcal{O}_X(U)$ est
semi-normal.
\item Il existe un recouvrement ouvert affine $X = \bigcup U_i$ tel que
chaque $\mathcal{O}_X(U_i)$ soit semi-normal.
\item Il existe un recouvrement ouvert $X = \bigcup X_j$ tel que chaque
sous-schéma ouvert $X_j$ soit semi-normal.
\end{enumerate}
De plus, si $X$ est semi-normal, alors
tout sous-schéma ouvert est semi-normal. Les mêmes affirmations
sont vraies lorsque ``semi-normal'' est remplacé par ``absolument
faiblement normal''.
\end{lemma}

\begin{proof}
Combiner les Propriétés, Lemme \ref{properties-lemma-locally-P}
et Lemme \ref{morphisms-lemma-seminormal-local-property}.
\end{proof}

\begin{lemma}
\label{morphisms-lemma-seminormal-reduced}
Un schéma ou anneau semi-normal est réduit. A fortiori, il en est de même
pour les schémas ou anneaux absolument faiblement normaux.
\end{lemma}

\begin{proof}
Soit $A$ un anneau.
Si $a \in A$ est non nul mais $a^2 = 0$, alors $a^2 = 0^2$ et $a^3 = 0^3$
et donc $A$ ne sont pas semi-normaux.
\end{proof}

\begin{lemma}
\label{morphisms-lemma-seminormalization-ring}
Soit $A$ un anneau.
\begin{enumerate}
\item La catégorie des homomorphismes d'anneaux $A \to B$ induisant
un homéomorphisme universel sur les spectres a un objet final $A \to A^{awn}$.
\item Étant donné $A \to B$ dans la catégorie de (1) l'application résultante
$B \to A^{awn}$ est un isomorphisme si et seulement si $B$ est
absolument faiblement normal.
\item La catégorie des homomorphismes d'anneaux $A \to B$ induisant des isomorphismes
sur les corps résiduels et un homéomorphisme universel sur les spectres a un
objet final $A \to A^{sn}$.
\item Étant donné $A \to B$ dans la catégorie de (3) l'application résultante
$B \to A^{sn}$ est un isomorphisme si et seulement si $B$ est semi-normal.
\end{enumerate}
Pour tout homomorphisme d'anneaux $\varphi : A \to A'$ il
existe des applications uniques $\varphi^{awn} : A^{awn} \to (A')^{awn}$
et $\varphi^{sn} : A^{sn} \to (A')^{sn}$ compatibles avec $\varphi$.
\end{lemma}

\begin{proof}
Nous prouvons (1) et (2) et nous omettons la démonstration de (3) et (4) ainsi
que l'énoncé final. Considérons la catégorie des $A$-algèbres de la forme
$$
B = A[x_1, \ldots, x_n]/J
$$
où $J$ est un idéal de type fini tel que $A \to B$ définit un homéomorphisme
universel sur les spectres. Nous affirmons que cette catégorie
est dirigée (Catégories, Définition \ref{categories-definition-directed}). À savoir, étant donné
$$
B = A[x_1, \ldots, x_n]/J
\quad\text{et}\quad
B' = A[x_1, \ldots, x_{n'}]/J'
$$
alors nous pouvons considérer
$$
B'' = A[x_1, \ldots, x_{n + n'}]/J''
$$
où $J''$ est généré par les éléments de $J$ et
les éléments $f(x_{n + 1}, \ldots, x_{n + n'})$ où $f \in J'$. Ensuite,
nous avons des homomorphismes d'algèbre $A$ $B \to B''$
et $B' \to B''$ qui induisent un isomorphisme
$B \otimes_A B' \to B''$.
Il découle des Lemme \ref{morphisms-lemma-base-change-universal-homeomorphism} et \ref{morphisms-lemma-composition-universal-homeomorphism}
que $\Spec(B'') \to \Spec(A)$ est un homéomorphisme universel
et donc $A \to B''$ est dans notre catégorie. Enfin,
étant donné $\varphi, \varphi' : B \to B'$ dans notre
catégorie avec $B$ comme affiché
ci-dessus, nous considérons alors
le quotient $B''$ de $B'$
par l'idéal engendré par $\varphi(x_i) - \varphi'(x_i)$, $i = 1, \ldots, n$. Puisque $\Spec(B') = \Spec(B)$ nous
voyons que $\Spec(B'') \to \Spec(B')$ est une immersion bijective fermée, donc
un homéomorphisme universel. Ainsi $B''$ est dans notre
catégorie et $\varphi, \varphi'$ sont égalisés par $B' \to B''$. Ceci complète
la démonstration de notre affirmation. Nous posons
$$
A^{awn} = \colim B
$$
où la colimite est sur la catégorie précédemment décrite. Observez que
$A \to A^{awn}$ induit un homéomorphisme universel sur les spectres selon
le Lemme \ref{morphisms-lemma-colimit-inherits} (c'est ici que nous utilisons que la catégorie est
dirigée).

\medskip\noindent Étant
donné un homomorphisme d'anneaux $A \to B$ de présentation finie induisant
un homéomorphisme universel sur les spectres, nous obtenons une application canonique
$B \to A^{awn}$ par la construction même de $A^{awn}$. Puisque chaque
$A \to B$ comme dans (1) est une colimite filtrante de $A \to B$ comme
dans (1) de présentation finie (Lemme \ref{morphisms-lemma-universal-homeo-limit}), nous voyons
que $A \to A^{awn}$ est final dans la catégorie (1).

\medskip\noindent
Soient $x, y \in A^{awn}$ des éléments tels que $x^3 = y^2$.
Alors $A^{awn} \to A^{awn}[t]/(t^2 - x, t^3 - y)$ induit un homéomorphisme universel
sur les spectres par le Lemme \ref{morphisms-lemma-make-universal-homeo}. Ainsi $A \to A^{awn}[t]/(t^2 - x, t^3 - y)$ est dans
la catégorie (1) et nous obtenons une unique application
d'$A$-algèbre $A^{awn}[t]/(t^2 - x, t^3 - y) \to A^{awn}$.
L'image $a \in A^{awn}$ de $t$ est donc
l'élément unique tel que $a^2 = x$ et $a^3 = y$
dans $A^{awn}$. De la même manière, donné un
premier $p$ et $x, y \in A^{awn}$ avec $p^px = y^p$, nous trouvons
un unique $a \in A^{awn}$ avec $a^p = x$ et $pq = y$.
Ainsi, $A^{awn}$ est absolument faiblement normal
par définition.

\medskip\noindent
Enfin, soit $A \to B$ dans la catégorie (1) avec $B$ absolument
faiblement normal. Puisque $A^{awn} \to B^{awn}$ induit un homéomorphisme
universel sur les spectres et puisque $A^{awn}$ est réduit
(Lemme \ref{morphisms-lemma-seminormal-reduced}),
nous trouvons $A^{awn} \subset B^{awn}$ (voir Algèbre, Lemme \ref{algebra-lemma-image-dense-generic-points}).
Si cette inclusion n'est pas une égalité,
alors le Lemme \ref{morphisms-lemma-pth-power-and-multiple} implique
qu'il existe un élément $b \in B^{awn}$, $b \not \in A^{awn}$ tel
que soit $b^2, b^3 \in A^{awn}$ soit $pb, b^p \in A^{awn}$ pour un certain
nombre premier $p$. Cependant, par l'existence et l'unicité
dans la Définition \ref{morphisms-definition-seminormal-ring}, cela force $b \in A^{awn}$
et donc nous obtenons la contradiction qui termine la démonstration.
\end{proof}

\begin{lemma}
\label{morphisms-lemma-seminormalization}
Soit $X$ un schéma.
\begin{enumerate}
\item La catégorie des homéomorphismes universels $Y \to X$
a un objet initial $X^{awn} \to X$.
\item Étant donné $Y \to X$ dans la catégorie de (1) le morphisme résultant $X^{awn} \to Y$
est un isomorphisme si et seulement si $Y$ est absolument
faiblement normal.
\item La catégorie des homéomorphismes universels $Y \to X$ qui induisent
des isomorphismes sur les corps résiduels a un objet initial
$X^{sn} \to X$.
\item Étant donné $Y \to X$ dans la catégorie de (3) le morphisme résultant $X^{sn} \to Y$
est un isomorphisme si et seulement si $Y$ est semi-normal.
\end{enumerate}
Pour tout morphisme $h : X' \to X$ de schémas, il existe des
morphismes uniques $h^{awn} : (X')^{awn} \to X^{awn}$ et $h^{sn} : (X')^{sn} \to X^{sn}$ compatibles avec
$h$.
\end{lemma}

\begin{proof}
Nous allons prouver (1) et (2) et omettre la démonstration
de (3) et (4). Soit $h : X' \to X$ un morphisme de schémas. Si (1)
vaut pour $X$ et $X'$, alors $X' \times_X X^{awn} \to X'$ est un homéomorphisme
universel et nous obtenons donc un morphisme unique $(X')^{awn} \to X' \times_X X^{awn}$
au-dessus de $X'$ par la propriété universelle de $(X')^{awn} \to X'$.
Composé avec la projection $X' \times_X X^{awn} \to X^{awn}$, nous obtenons $h^{awn}$.
Si de plus (2) vaut pour $X$ et $X'$ et que $h$ est
une immersion ouverte, alors $X' \times_X X^{awn}$ est
absolument faiblement normal (Lemme \ref{morphisms-lemma-locally-seminormal}) et nous
en déduisons que $(X')^{awn} \to X' \times_X X^{awn}$ est un isomorphisme.

\medskip\noindent
Rappelez-vous que tout homéomorphisme universel
est affine, voir Lemme \ref{morphisms-lemma-homeomorphism-affine}. Ainsi, si $X$
est affine, alors (1) et (2) en découlent
immédiatement d'après le Lemme \ref{morphisms-lemma-seminormalization-ring}.
Soit $X$ un schéma et soit $\mathcal{B}$ l'ensemble
des ouverts affines de $X$. Pour chaque $U \in \mathcal{B}$
nous obtenons $U^{awn} \to U$ et pour $V \subset U$, $V, U \in \mathcal{B}$
nous obtenons un isomorphisme canonique $\rho_{V, U} : V^{awn} \to V \times_U U^{awn}$ selon la discussion
du paragraphe précédent. Ainsi, par recollement
relatif (Constructions, Lemme \ref{constructions-lemma-relative-glueing}), nous obtenons
un morphisme $X^{awn} \to X$ qui se restreint
à $U^{awn}$ sur $U$ de manière compatible avec le
$\rho_{V, U}$. Ensuite, soit $Y \to X$ un homéomorphisme
universel. Alors $U \times_X Y \to U$ est un homéomorphisme universel pour $U \in \mathcal{B}$
et nous obtenons un morphisme unique $g_U : U^{awn} \to U \times_X Y$ sur $U$. Ces $g_U$
sont compatibles avec les morphismes $\rho_{V, U}$ ; détails omis.
Par conséquent, il existe un morphisme unique $g : X^{awn} \to Y$ sur
$X$ concordant avec $g_U$ sur
$U$, voir Constructions, Remarque \ref{constructions-remark-relative-glueing-functorial}. Cela prouve (1) pour $X$.
La partie (2) s'ensuit car elle tient localement sur les affines.
\end{proof}

\begin{definition}
\label{morphisms-definition-seminormalization}
Soit $X$ un schéma.
\begin{enumerate}
\item Le morphisme $X^{sn} \to X$ construit dans le Lemme \ref{morphisms-lemma-seminormalization}
est la {\it semi-normalisation}
de $X$.
\item Le morphisme $X^{awn} \to X$ construit dans le Lemme \ref{morphisms-lemma-seminormalization}
est la {\it normalisation
faible absolue} de $X$.
\end{enumerate}
\end{definition}

\noindent
Pour être sûr, la semi-normalisation $X^{sn}$ de $X$ est un schéma
semi-normal et la normalisation faible absolue $X^{awn}$ est un
schéma absolument faiblement normal.
De plus, pour tout morphisme $h : Y \to X$ de schémas, nous obtenons un
diagramme commutatif canonique
$$
\xymatrix{
Y^{awn} \ar[d]^{h^{awn}} \ar[r] &
Y^{sn} \ar[d]^{h^{sn}} \ar[r] &
Y \ar[d]^h \\
X^{awn} \ar[r] &
X^{sn} \ar[r] &
X
}
$$
de schémas ; les flèches $h^{sn}$ et $h^{awn}$ sont les uniques compatibles
avec $h$.

\begin{lemma}
\label{morphisms-lemma-characterize-seminormal}
Soit $X$ un schéma. Les conditions suivantes sont équivalentes
\begin{enumerate}
\item $X$ est semi-normal,
\item $X$ est égal à sa propre semi-normalisation, c'est-à-dire que le morphisme
$X^{sn} \to X$ est un isomorphisme,
\item si $\pi : Y \to X$ est un homéomorphisme universel induisant des isomorphismes sur les corps
résiduels avec $Y$ réduit, alors $\pi$ est un isomorphisme.
\end{enumerate}
\end{lemma}

\begin{proof}
L'équivalence de (1) et (2) est claire d'après
le Lemme \ref{morphisms-lemma-seminormalization}. Si (3) est vrai, alors $X^{sn} \to X$ est un isomorphisme
et nous voyons que (2) est vrai.

\medskip\noindent
Supposons que (2) vaut et soit $\pi : Y \to X$ un homéomorphisme universel
induisant des isomorphismes sur les corps résiduels avec $Y$ réduit.
Alors il existe une factorisation $X \to Y \to X$ de $\text{id}_X$
par le Lemme \ref{morphisms-lemma-seminormalization}. Alors $X \to Y$ est une immersion fermée
(par Schémas, Lemme \ref{schemes-lemma-section-immersion} et le fait que $\pi$
est séparé, par exemple par le Lemme \ref{morphisms-lemma-universally-injective-separated}).
Puisque $X \to Y$ est également une bijection sur les points, la caractéristique
d'être réduit de $Y$ montre qu'il doit s'agir
d'un isomorphisme. Cela termine la démonstration.
\end{proof}

\begin{lemma}
\label{morphisms-lemma-characterize-absolutely-weakly-normal}
Soit $X$ un schéma. Les conditions suivantes sont équivalentes
\begin{enumerate}
\item $X$ est faiblement normal absolu,
\item $X$ est égal à sa propre normalisation faible absolue, c'est-à-dire que le morphisme
$X^{awn} \to X$ est un isomorphisme,
\item si $\pi : Y \to X$ est un homéomorphisme universel avec $Y$ réduit, alors $\pi$
est un isomorphisme.
\end{enumerate}
\end{lemma}

\begin{proof}
Cela se démontre exactement de la même
manière que le Lemme \ref{morphisms-lemma-characterize-seminormal}.
\end{proof}





\section{Morphismes finis localement libres}
\label{morphisms-section-finite-locally-free}

\noindent
Dans de nombreux articles, les auteurs utilisent des morphismes finis plats alors qu'ils
veulent réellement dire des morphismes finis localement libres. La raison en
est que si la base est localement noethérienne alors c'est la même chose. Mais
en général ce n'est pas le cas, voir Exercices, Exercice \ref{exercises-exercise-flat-not-projective}.

\begin{definition}
\label{morphisms-definition-finite-locally-free}
Soit $f : X \to S$ un morphisme de schémas. On dit que $f$
est {\it fini localement libre} si $f$
est affine et que $f_*\mathcal{O}_X$ est un $\mathcal{O}_S$-module fini
localement libre. Dans ce cas, on dit que $f$ a {\it
rang} ou {\it degré}
$d$ si le faisceau $f_*\mathcal{O}_X$ est fini localement libre de degré $d$.
\end{definition}

\noindent
Remarquez que si $f : X \to S$ est fini localement libre alors $S$ est l'union disjointe
de ouverts et de sous-schémas fermés $S_d$ tels que $f^{-1}(S_d) \to S_d$ soit fini
localement libre de degré $d$.

\begin{lemma}
\label{morphisms-lemma-finite-flat}
Soit $f : X \to S$ un morphisme de schémas. Les conditions
suivantes sont équivalentes :
\begin{enumerate}
\item $f$ est fini localement libre,
\item $f$ est fini, plat et localement de présentation finie.
\end{enumerate}
Si $S$ est localement noethérien, cela équivaut également à
\begin{enumerate}
\item[(3)] $f$ est fini et plat.
\end{enumerate}
\end{lemma}

\begin{proof}
Soit $V \subset S$ ouvert affine. Dans les trois cas, le morphisme
est affine, donc $f^{-1}(V)$ est affine. Ainsi, nous pouvons
écrire $V = \Spec(R)$ et $f^{-1}(V) = \Spec(A)$ pour une certaine $R$-algèbre $A$.
Supposons (1). Cela
signifie que nous pouvons couvrir $S$ par des ouverts affines
$V = \Spec(R)$ tels que $A$ est libre de type fini en tant que $R$-module.
Alors $R \to A$ est de présentation
finie par l'Algèbre, Lemme \ref{algebra-lemma-finite-finite-type}. Ainsi (2) est vérifié. Réciproquement,
supposons (2). Pour chaque ouvert
affine $V = \Spec(R)$ de $S$, le homomorphisme d'anneaux $R \to A$
est fini et de présentation finie et $A$ est plat en tant que
$R$-module. Par
Algèbre, Lemme \ref{algebra-lemma-finite-finitely-presented-extension} nous voyons que $A$ est de présentation
finie en tant que $R$-module.
Ainsi Algèbre, Lemme \ref{algebra-lemma-finite-projective} implique que $A$
est fini localement libre. Ainsi (1)
est vrai. Le cas noethérien découle du fait qu'un module
fini sur un anneau noethérien est un module
de présentation finie, voir Algèbre, Lemme \ref{algebra-lemma-Noetherian-finite-type-is-finite-presentation}.
\end{proof}

\begin{lemma}
\label{morphisms-lemma-composition-finite-locally-free}
Une composition de morphismes finis localement libres est finie localement libre.
\end{lemma}

\begin{proof}
Omis.
\end{proof}

\begin{lemma}
\label{morphisms-lemma-base-change-finite-locally-free}
Un changement de base d'un morphisme fini localement libre est fini localement libre.
\end{lemma}

\begin{proof}
Omis.
\end{proof}

\begin{lemma}
\label{morphisms-lemma-finite-locally-free}
Soit $f : X \to S$ un morphisme fini localement libre de schémas. Il
existe une décomposition en union disjointe
$S = \coprod_{d \geq 0} S_d$ par des ouverts et sous-schémas fermés telle qu'en
posant $X_d = f^{-1}(S_d)$, les restrictions $f|_{X_d}$ soient
des morphismes finis localement libres $X_d \to S_d$ de degré
$d$.
\end{lemma}

\begin{proof}
Cela est vrai car un faisceau fini localement libre a localement
un rang bien défini. Détails omis.
\end{proof}

\begin{lemma}
\label{morphisms-lemma-massage-finite}
Soit $f : Y \to X$ un morphisme fini avec $X$ affine.
Il existe un diagramme
$$
\xymatrix{
Z' \ar[rd] &
Y' \ar[l]^i \ar[d] \ar[r] &
Y \ar[d] \\
 & X' \ar[r] & X
}
$$
où
\begin{enumerate}
\item $Y' \to Y$ et $X' \to X$ sont surjectifs, finis et localement libres,
\item $Y' = X' \times_X Y$,
\item $i : Y' \to Z'$ est une immersion fermée,
\item $Z' \to X'$ est fini localement libre, et
\item $Z' = \bigcup_{j = 1, \ldots, m} Z'_j$ est une union finie (au sens des ensembles) de sous-schémas
fermés, chacun desquels se projette isomorphiquement sur
$X'$.
\end{enumerate}
\end{lemma}

\begin{proof}
Écrivez $X = \Spec(A)$ et $Y = \Spec(B)$. Voir aussi Plus
sur l'Algèbre, section \ref{more-algebra-section-descent-flatness-integral}. Soit $x_1, \ldots, x_n \in B$ des générateurs de $B$
sur $A$. Pour chaque $i$, nous pouvons choisir
un polynôme unitaire $P_i(T) \in A[T]$ tel que $P(x_i) = 0$ soit dans
$B$. D'après l'Algèbre,
Lemme \ref{algebra-lemma-adjoin-roots} (appliqué $n$ fois) il
existe une extension d'anneaux fini localement libre $A \subset A'$
telle que chaque $P_i$ se décompose complètement :
$$
P_i(T) = \prod\nolimits_{k = 1, \ldots, d_i} (T - \alpha_{ik})
$$
pour certains $\alpha_{ik} \in A'$. Posons
$$
C = A'[T_1, \ldots, T_n]/(P_1(T_1), \ldots, P_n(T_n))
$$
et $B' = A' \otimes_A B$. L'application $C \to B'$, $T_i \mapsto 1 \otimes x_i$ est une surjection d'$A'$-algèbres.
En posant $X' = \Spec(A')$, $Y' = \Spec(B')$ et $Z' = \Spec(C)$,
nous voyons que (1) -- (4) sont vérifiées. La partie
(5) est vérifiée parce que, du point de vue des ensembles,
$\Spec(C)$ est l'union des sous-schémas fermés
coupés par les idéaux
$$
(T_1 - \alpha_{1k_1}, T_2 - \alpha_{2k_2}, \ldots, T_n - \alpha_{nk_n})
$$
pour tout $1 \leq k_i \leq d_i$.
\end{proof}

\noindent
Le lemme suivant est énoncé dans la généralité
correcte dans le Lemme \ref{morphisms-lemma-image-nowhere-dense-quasi-finite} ci-dessous.

\begin{lemma}
\label{morphisms-lemma-image-nowhere-dense-finite}
Soit $f : Y \to X$ un morphisme fini de schémas. Soit $T \subset Y$
une partie fermée nulle part dense de $Y$. Alors $f(T) \subset X$
est une partie fermée nulle part dense de $X$.
\end{lemma}

\begin{proof}
D'après le Lemme \ref{morphisms-lemma-finite-proper} nous savons que $f(T) \subset X$ est fermé. Soit $X = \bigcup X_i$
un recouvrement affine. Comme $T$ est
nulle part dense dans $Y$, nous voyons que $T \cap f^{-1}(X_i)$ est également
nulle part dense dans $f^{-1}(X_i)$. Ainsi, si nous pouvons prouver le théorème dans
le cas affine, alors nous voyons que $f(T) \cap X_i$ est nulle part dense.
Cela implique alors que $T$ est nulle part
dense dans $X$ selon la Topologie, Lemme \ref{topology-lemma-nowhere-dense-local}.

\medskip\noindent Supposons
que $X$ soit affine. Choisissez un diagramme
$$
\xymatrix{
Z' \ar[rd] &
Y' \ar[l]^i \ar[d]^{f'} \ar[r]_a &
Y \ar[d]^f \\
 & X' \ar[r]^b & X
}
$$
comme dans le Lemme \ref{morphisms-lemma-massage-finite}. Les morphismes $a$,
$b$ sont ouverts puisqu'ils
sont finis localement libres (Lemmes \ref{morphisms-lemma-finite-flat} et \ref{morphisms-lemma-fppf-open}).
Par conséquent, $T' = a^{-1}(T)$
est nulle part dense, voir Topologie, Lemme \ref{topology-lemma-open-inverse-image-closed-nowhere-dense}. Le morphisme
$b$ est surjectif et ouvert.
Par conséquent, si nous pouvons prouver que
$f'(T') = b^{-1}(f(T))$ est nulle part dense, alors $f(T)$
est nulle part dense, voir Topologie, Lemme \ref{topology-lemma-open-inverse-image-closed-nowhere-dense}. Comme
$i$ est une immersion
fermée, d'après Topologie, Lemme \ref{topology-lemma-closed-image-nowhere-dense} nous voyons
que $i(T') \subset Z'$ est fermée et nulle part dense.
Ainsi, nous avons réduit le problème au cas discuté dans
le paragraphe suivant.

\medskip\noindent Supposons
que $Y = \bigcup_{i = 1, \ldots, n} Y_i$ soit une union finie de parties fermées, chacune s'envoyant
isomorphiquement sur $X$. Considérons $T_i = Y_i \cap T$. Si chacun
des $T_i$ est nulle part dense dans $Y_i$, alors chacun des $f(T_i)$
est nulle part dense dans $X$ car $Y_i \to X$ est un isomorphisme. Par conséquent,
$f(T) = f(T_i)$ est une union finie de sous-ensembles fermés nulle part
denses de $X$ et nous avons
gagné, voir Topologie, Lemme \ref{topology-lemma-nowhere-dense}. Supposons que ce ne
soit pas le cas, disons que $T_1$ contient un $V \subset Y_1$ ouvert non
vide. Nous allons montrer que cela conduit à une contradiction.
Considérons $Y_2 \cap V \subset V$. Ceci est soit un sous-ensemble
fermé propre, soit égal à $V$. Dans le premier cas, nous remplaçons
$V$ par $V \setminus V \cap Y_2$, donc $V \subset T_1$ est ouvert dans $Y_1$ et ne rencontre
pas $Y_2$. Dans le second cas, nous avons
$V \subset Y_1 \cap Y_2$ est ouvert à la fois dans $Y_1$ et $Y_2$.
Répétez séquentiellement avec $i = 3, \ldots, n$. Le résultat est une décomposition
en union disjointe
$$
\{1, \ldots, n\} = I_1 \amalg I_2, \quad 1 \in I_1
$$
et un $V$ ouvert de $Y_1$ contenu dans $T_1$ tel que
$V \subset Y_i$ pour $i \in I_1$ et $V \cap Y_i = \emptyset$ pour $i \in I_2$. Posons $U = f(V)$.
C'est un ouvert de $X$ puisque $f|_{Y_1} : Y_1 \to X$ est un
isomorphisme. Alors
$$
f^{-1}(U) = V\ \amalg\ \bigcup\nolimits_{i \in I_2} (Y_i \cap f^{-1}(U))
$$
Comme $\bigcup_{i \in I_2} Y_i$ est fermé, cela implique que $V \subset f^{-1}(U)$ est
ouvert, donc $V \subset Y$ est ouvert. Cela contredit l'hypothèse
que $T$ est quelque part dense dans $Y$, comme souhaité.
\end{proof}







\section{Applications rationnelles}
\label{morphisms-section-rational-maps}

\noindent
Soit $X$ un schéma. Remarquez que si $U$, $V$ sont des ouverts denses dans $X$,
alors il en est de même de $U \cap V$.

\begin{definition}
\label{morphisms-definition-rational-map}
Soit $X$, $Y$ des schémas.
\begin{enumerate}
\item Soit $f : U \to Y$, $g : V \to Y$ des morphismes de schémas définis sur des ouverts denses $U$,
$V$ de $X$. On dit que $f$ est {\it
équivalent} à $g$ si $f|_W = g|_W$ pour un certain $W \subset U \cap V$ ouvert
dense dans $X$.
\item Une {\it application rationnelle de $X$
vers $Y$} est une classe d'équivalence pour la relation d'équivalence définie en (1).
\item Si $X$, $Y$ sont des schémas sur un schéma de base $S$, nous disons qu'une
application rationnelle de $X$ vers $Y$ est une {\it $S$-application
rationnelle de $X$ vers $Y$} s'il existe un représentant $f : U \to Y$ de la classe
d'équivalence qui est un $S$-morphisme.
\end{enumerate}
\end{definition}

\noindent
Nous disons que deux morphismes $f$, $g$ comme dans (1) de la définition définissent
la même application rationnelle au lieu de dire qu'ils sont équivalents. Dans certains cas,
les applications rationnelles sont déterminées par des applications sur les anneaux locaux
aux points génériques.

\begin{lemma}
\label{morphisms-lemma-rational-map-finite-presentation}
Soit $S$ un schéma. Soient $X$ et $Y$ des schémas sur $S$. Supposons
que $X$ ait un nombre fini de composantes irréductibles avec des
points génériques $x_1, \ldots, x_n$. Soit $s_i \in S$ l'image de $x_i$. Considérons
l'application
$$
\left\{
\begin{matrix}
S\text{-applications rationnelles} \\
\text{de }X\text{ vers }Y
\end{matrix}
\right\}
\longrightarrow
\left\{
\begin{matrix}
(y_1, \varphi_1, \ldots, y_n, \varphi_n)\text{ où }
y_i \in Y\text{ est au-dessus de }s_i\text{ et}\\
\varphi_i : \mathcal{O}_{Y, y_i} \to \mathcal{O}_{X, x_i}
\text{ est un homomorphisme local de }\mathcal{O}_{S, s_i}\text{-algèbres}
\end{matrix}
\right\}
$$
qui envoie $f : U \to Y$ sur le $2n$-uplet
avec $y_i = f(x_i)$ et $\varphi_i = f^\sharp_{x_i}$. Alors
\begin{enumerate}
\item Si $Y \to S$ est localement de type fini, alors l'application est injective.
\item Si $Y \to S$ est localement de présentation finie, alors l'application est bijective.
\item Si $Y \to S$ est localement de type fini et que $X$ est réduit, alors
l'application est bijective.
\end{enumerate}
\end{lemma}

\begin{proof}
Remarquez que tout ouvert dense de $X$ contient les points $x_i$,
donc la construction a du sens. Pour prouver (1)
ou (2), nous pouvons remplacer $X$ par tout ouvert dense. Ainsi, si
$Z_1, \ldots, Z_n$ sont les composantes irréductibles de $X$, alors
nous pouvons remplacer $X$ par $X \setminus \bigcup_{i \not = j} Z_i \cap Z_j$. Après avoir
fait cela, $X$ est l'union disjointe de ses composantes irréductibles
(considérées comme ouvertes et sous-schémas fermés). Ensuite,
à la fois le côté droit et le côté gauche de la flèche sont des produits
sur les composantes irréductibles et nous réduisons au cas où $X$
est irréductible.

\medskip\noindent Supposons
que $X$ est irréductible avec le point générique $x$ situé
au-dessus de $s \in S$. La partie
(1) découle de la partie (1) du Lemme \ref{morphisms-lemma-morphism-defined-local-ring}. Les
parties (2) et (3) découlent de la partie (2 du même lemme.
\end{proof}

\begin{definition}
\label{morphisms-definition-rational-function}
Soit $X$ un schéma. Une fonction rationnelle {\it sur $X$} est une
application rationnelle de $X$ vers $\mathbf{A}^1_{\mathbf{Z}}$.
\end{definition}

\noindent
Voir Constructions, Définition \ref{constructions-definition-affine-n-space} pour la définition de la droite affine
$\mathbf{A}^1$. Soit $X$ un schéma sur $S$. Pour tout ouvert
$U \subset X$, un morphisme $U \to \mathbf{A}^1_{\mathbf{Z}}$ est le
même qu'un morphisme $U \to \mathbf{A}^1_S$ sur $S$. Par conséquent,
une fonction rationnelle est également la même qu'une
$S$-application rationnelle de $X$ vers $\mathbf{A}^1_S$.

\medskip\noindent
Rappelons que nous avons l'identification
canonique $\Mor(T, \mathbf{A}^1_{\mathbf{Z}}) = \Gamma(T, \mathcal{O}_T)$ pour tout schéma $T$,
voir Schémas, Exemple \ref{schemes-example-global-sections}. Ainsi $\mathbf{A}^1_{\mathbf{Z}}$ est
un objet anneau dans la catégorie des schémas.
Plus précisément, les morphismes
\begin{eqnarray*}
+ : \mathbf{A}^1_{\mathbf{Z}} \times \mathbf{A}^1_{\mathbf{Z}}
& \longrightarrow
& \mathbf{A}^1_{\mathbf{Z}}
\\ (f, g) & \longmapsto & f + g \\ *
: \mathbf{A}^1_{\mathbf{Z}} \times \mathbf{A}^1_{\mathbf{Z}}
& \longrightarrow
& \mathbf{A}^1_{\mathbf{Z}}
\\ (f, g) & \longmapsto & fg
\end{eqnarray*}
satisfaire tous les axiomes de l'addition et de la multiplication dans un anneau (commutatif
avec $1$ comme toujours). Par conséquent, l'ensemble des applications
rationnelles dans $\mathbf{A}^1_{\mathbf{Z}}$ possède également une structure d'anneau naturelle.

\begin{definition}
\label{morphisms-definition-ring-of-rational-functions}
Soit $X$ un schéma. L'anneau {\it des fonctions rationnelles sur $X$}
est l'anneau $R(X)$ dont les éléments sont des fonctions rationnelles avec
l'addition et la multiplication telles que décrites ci-dessus.
\end{definition}

\noindent
Pour les schémas avec un nombre fini de composantes irréductibles, nous pouvons le calculer.

\begin{lemma}
\label{morphisms-lemma-integral-scheme-rational-functions}
Soit $X$ un schéma avec un nombre fini de composantes irréductibles
$X_1, \ldots, X_n$. Si $\eta_i \in X_i$ est le point générique, alors
$$
R(X) = \mathcal{O}_{X, \eta_1} \times \ldots \times \mathcal{O}_{X, \eta_n}
$$
Si $X$ est réduit, cela équivaut à $\prod \kappa(\eta_i)$.
Si $X$ est intègre, alors $R(X) = \mathcal{O}_{X, \eta} = \kappa(\eta)$ est un
corps.
\end{lemma}

\begin{proof}
Soit $U \subset X$ un ouvert dense. Alors
$U_i = (U \cap X_i) \setminus (\bigcup_{j \not = i} X_j)$ est non vide et ouvert car il contient
$\eta_i$, contenu dans $X_i$, et $\bigcup U_i \subset U \subset X$
est dense. Ainsi, l'identification
dans le lemme provient de la chaîne d'égalités
\begin{align*}
R(X)
& =
\colim_{U \subset X\text{ ouvert dense}} \Mor(U, \mathbf{A}^1_\mathbf{Z}) \\
& =
\colim_{U \subset X\text{ ouvert dense}} \mathcal{O}_X(U) \\
& =
\colim_{\eta_i \in U_i \subset X\text{ ouvert}} \prod \mathcal{O}_X(U_i) \\
& =
\prod \colim_{\eta_i \in U_i \subset X\text{ ouvert}} \mathcal{O}_X(U_i) \\
& =
\prod \mathcal{O}_{X, \eta_i}
\end{align*}
où la deuxième égalité
est Schémas, Exemple \ref{schemes-example-global-sections}.
L'énoncé final
découle d'Algèbre, Lemme \ref{algebra-lemma-minimal-prime-reduced-ring}.
\end{proof}

\begin{definition}
\label{morphisms-definition-function-field}
Soit $X$ un schéma intègre. La fonction {\it
corps}, ou le {\it corps des fonctions rationnelles }
de $X$ est le corps $R(X)$.
\end{definition}

\noindent
Nous pouvons occasionnellement indiquer ce corps $k(X)$ au lieu de
$R(X)$. Nous pouvons utiliser la notion de corps de fonctions
pour élucider la condition de séparation sur
un schéma intègre. Notez que d'après le Lemme \ref{morphisms-lemma-integral-scheme-rational-functions} sur un schéma
intègre, chaque anneau local $\mathcal{O}_{X, x}$ peut être considéré comme un sous-anneau
local de $R(X)$.

\begin{lemma}
\label{morphisms-lemma-distinct-local-rings}
Soit $X$ un schéma intègre séparé.
Soient $Z_1$ et $Z_2$ des sous-ensembles fermés irréductibles
distincts de $X$. Soit
$\eta_i$ le point générique
de $Z_i$. Si $Z_1 \not\subset Z_2$, alors $\mathcal{O}_{X, \eta_1} \not \subset \mathcal{O}_{X, \eta_2}$ comme sous-anneaux
de $R(X)$. En
particulier, si $Z_1 = \{x\}$ consiste en un point fermé $x$,
il existe une fonction régulière dans un voisinage de
$x$ qui n'est pas dans $\mathcal{O}_{X, \eta_{2}}$.
\end{lemma}

\begin{proof}
Observez d'abord que sous l'hypothèse que $X$ est séparé,
il existe une application unique
de schémas $\Spec(\mathcal{O}_{X, \eta_2}) \to X$ sur $X$ telle
que la composition
$$
\Spec(R(X)) \longrightarrow
\Spec(\mathcal{O}_{X, \eta_2}) \longrightarrow X
$$
soit l'application canonique $\Spec(R(X)) \to X$.
À savoir, il existe l'application
canonique $can : \Spec(\mathcal{O}_{X, \eta_2}) \to X$, voir Schémas,
Équation (\ref{schemes-equation-canonical-morphism}). Étant donné un second morphisme $a$
vers $X$, nous avons que $a$ coïncide
avec $can$ sur le
point générique de $\Spec(\mathcal{O}_{X, \eta_2})$ par hypothèse.
Maintenant, le fait que $X$ soit séparé
garantit que le sous-ensemble
dans $\Spec(\mathcal{O}_{X, \eta_2})$ où ces deux applications
coïncident est fermé, voir Schémas, Lemme
\ref{schemes-lemma-where-are-they-equal}. Ainsi $a = can$ sur tout $\Spec(\mathcal{O}_{X, \eta_2})$.

\medskip\noindent
Supposons $Z_1 \not \subset Z_2$ et supposons au contraire
que $\mathcal{O}_{X, \eta_{1}} \subset \mathcal{O}_{X, \eta_{2}}$ comme sous-anneaux de $R(X)$.
Alors nous obtiendrions un second morphisme
$$
\Spec(\mathcal{O}_{X, \eta_{2}}) \longrightarrow
\Spec(\mathcal{O}_{X, \eta_{1}}) \longrightarrow
X.
$$
D'après ce qui précède, cette composition devrait être égale
à $can$. Cela implique que $\eta_2$ se spécialise
en $\eta_1$ (voir Schémas, Lemme \ref{schemes-lemma-specialize-points}).
Mais cela contredit notre hypothèse $Z_1 \not \subset Z_2$.
\end{proof}

\begin{definition}
\label{morphisms-definition-domain-of-definition}
Soit $\varphi$ une application rationnelle entre deux schémas $X$ et $Y$. Nous disons
que $\varphi$ est {\it défini en un point $x \in X$} s'il
existe un représentant $(U, f)$ de $\varphi$ avec $x \in U$. Le {\it
domaine de définition} de $\varphi$ est l'ensemble de tous les
points où $\varphi$ est défini.
\end{definition}

\noindent
Avec cette définition, il n'est pas vrai en général que $\varphi$
possède un représentant défini sur tout le domaine de définition.

\begin{lemma}
\label{morphisms-lemma-rational-map-from-reduced-to-separated}
Soient $X$ et $Y$ des schémas. Supposons $X$ réduit et $Y$ séparé.
Soit $\varphi$ une application rationnelle de $X$ vers $Y$ avec domaine de
définition $U \subset X$. Alors il existe un unique morphisme $f : U \to Y$ représentant $\varphi$.
Si $X$ et $Y$ sont des schémas sur un schéma séparé $S$ et
si $\varphi$ est une application rationnelle $S$, alors $f$ est un morphisme
sur $S$.
\end{lemma}

\begin{proof}
Soient $(V, g)$ et $(V', g')$ des représentants de $\varphi$.
Alors $g, g'$ coïncident sur un sous-schéma ouvert dense
$W \subset V \cap V'$. D'un autre côté, l'égaliseur $E$ de $g|_{V \cap V'}$ et $g'|_{V \cap V'}$
est un sous-schéma fermé de $V \cap V'$ (Schémas,
Lemme \ref{schemes-lemma-where-are-they-equal}). Maintenant $W \subset E$ implique que
$E = V \cap V'$ d'un point de vue des ensembles. Comme $V \cap V'$
est réduit, nous concluons $E = V \cap V'$ d'un point de vue schématique,
c'est-à-dire, $g|_{V \cap V'} = g'|_{V \cap V'}$. Il s'ensuit que nous pouvons recoller
les représentants $g : V \to Y$ de $\varphi$ en un morphisme $f : U \to Y$,
voir Schémas, Lemme \ref{schemes-lemma-glue}. Nous
omettons la démonstrations de l'énoncé final.
\end{proof}

\noindent
En général, il n'est pas logique de composer des applications rationnelles. La raison est
que l'image d'un représentant de la première application rationnelle peut avoir une intersection
vide avec le domaine de définition de la seconde. Cependant, si nous supposons
que nos schémas sont irréductibles et que nous considérons des applications rationnelles
dominantes, alors nous pouvons composer des applications rationnelles.

\begin{definition}
\label{morphisms-definition-dominant-rational}
Soient $X$ et $Y$ des schémas irréductibles. Une application rationnelle de $X$ vers
$Y$ est appelée {\it dominante} si tout représentant $f : U \to Y$ est un morphisme
de schémas dominant.
\end{definition}

\noindent
D'après le Lemme \ref{morphisms-lemma-dominant-finite-number-irreducible-components}, il est équivalent d'exiger que le point générique
$\eta \in X$ s'envoie sur le point générique $\xi$ de $Y$,
c'est-à-dire $f(\eta) = \xi$ pour tout représentant $f : U \to Y$.
Nous pouvons composer une application rationnelle dominante $\varphi$ entre
des schémas irréductibles $X$ et $Y$ avec une application rationnelle
arbitraire $\psi$ de $Y$ vers
$Z$. À savoir, choisir des représentants $f : U \to Y$
avec $U \subset X$ ouvert dense et $g : V \to Z$ avec $V \subset Y$
ouvert dense. Alors $W = f^{-1}(V) \subset X$ est ouvert non vide (car il contient
le point générique de $X$) et nous laissons $\psi \circ \varphi$
être la classe d'équivalence de $g \circ f|_W : W \to Z$. Nous omettons la vérification
que cela est bien défini.

\medskip\noindent De cette
manière, nous obtenons une catégorie dont les objets sont des schémas irréductibles et dont les morphismes
sont des applications rationnelles dominantes. Étant donné un schéma de base $S$, nous pouvons
de manière similaire définir une catégorie dont les objets sont des schémas irréductibles au-dessus de $S$
et dont les morphismes sont des $S$-applications rationnelles dominantes.

\begin{definition}
\label{morphisms-definition-birational-schemes}
Soient $X$ et $Y$ des schémas irréductibles.
\begin{enumerate}
\item Nous disons que $X$ et $Y$ sont {\it birationnels} si $X$ et $Y$
sont isomorphes dans la catégorie des schémas irréductibles et des applications rationnelles dominantes.
\item Supposons que $X$ et $Y$ sont des schémas sur un schéma de base $S$. On dit
que $X$ et $Y$ sont {\it $S$-birationnels} si $X$
et $Y$ sont isomorphes dans la catégorie des schémas irréductibles sur $S$ et des applications
rationnelles dominantes $S$.
\end{enumerate}
\end{definition}

\noindent
Si $X$ et $Y$ sont des schémas irréductibles birationnels, alors l'ensemble des applications
rationnelles de $X$ vers $Z$ est bijectif avec l'ensemble des applications rationnelles
de $Y$ vers $Z$ pour tous les schémas $Z$ (fonctoriellement en $Z$).
Pour des schémas irréductibles ``généraux'', ceci
n'est qu'une définition possible. Une autre serait de demander que $X$ et $Y$
aient des anneaux de fonctions rationnelles isomorphes.
Pour les variétés, ces conditions sont équivalentes, voir Lemme \ref{morphisms-lemma-criterion-birational-finite-presentation}.

\begin{lemma}
\label{morphisms-lemma-birational-integral}
Soient $X$ et $Y$ des schémas irréductibles.
\begin{enumerate}
\item Les schémas $X$ et $Y$ sont birationnels si et seulement si ils ont des
ouverts non vides isomorphes.
\item Supposons que $X$ et $Y$ soient des schémas sur un schéma de base $S$.
Alors $X$ et $Y$ sont $S$-birationnels si et seulement s'il
existe des ouverts non vides $U \subset X$ et $V \subset Y$ qui sont $S$-isomorphes.
\end{enumerate}
\end{lemma}

\begin{proof}
Supposons que $X$ et $Y$ soient birationnels. Soient $f : U \to Y$ et $g : V \to X$
définissant des applications rationnelles dominantes inverses de $X$ vers $Y$
et de $Y$ vers $X$. Nous pouvons supposer $V$ affine. Nous pouvons remplacer
$U$ par un ouvert affine de $f^{-1}(V)$. Comme $g \circ f$ est l'identité
en tant qu'application rationnelle dominante, nous voyons que la composition $U \to V \to X$
est l'identité sur un ouvert dense de $U$. Ainsi, après avoir remplacé $U$
par un ouvert affine plus petit, nous pouvons supposer que $U \to V \to X$ est
l'inclusion de $U$ dans
$X$. Il s'ensuit que $U \to V$ est une immersion (appliquer Schémas,
Lemme \ref{schemes-lemma-section-immersion} à $U \to g^{-1}(U) \to U$).
Cependant, en inversant les rôles de $U$ et $V$ et en refaisant
l'argument ci-dessus, nous voyons qu'il existe un ouvert affine non vide $V' \subset V$
tel que l'inclusion se factorise en $V' \to U \to V$. Alors $V' \to U$ est nécessairement
une immersion ouverte. En effet, $V' \to f^{-1}(V') \to V'$ sont des monomorphismes (Schémas,
Lemme
\ref{schemes-lemma-immersions-monomorphisms}) se composant pour donner l'identité, donc des isomorphismes.
Ainsi $V'$ est isomorphe à un ouvert à
la fois de $X$ et $Y$. Dans le cas des $S$-applications
rationnelles, le même argument exact fonctionne.
\end{proof}

\begin{remark}
\label{morphisms-remark-generalize-category}
Voici une généralisation de la catégorie des schémas irréductibles et des applications
rationnelles dominantes. Pour un schéma $X$, notons $X^0$ l'ensemble
des points $x \in X$ avec $\dim(\mathcal{O}_{X, x}) = 0$, en d'autres termes, $X^0$ est l'ensemble
des points génériques des composantes irréductibles de $X$. Ensuite,
nous pouvons considérer la catégorie avec
\begin{enumerate}
\item objets étant les schémas $X$ tels que chaque ouvert quasi-compact possède
un nombre fini de composantes irréductibles, et
\item morphismes de $X$ vers $Y$ étant des applications rationnelles $f : U \to Y$
de $X$ vers $Y$ telles que $f(U^0) = Y^0$.
\end{enumerate}
Si $U \subset X$ est un ouvert dense d'un schéma, alors $U^0 \subset X^0$
n'a pas besoin d'être une égalité, mais si $X$ est un objet
de notre catégorie, alors c'est le cas. Ainsi, étant
donnés deux morphismes dans notre catégorie, la composition est
bien définie et un morphisme dans notre catégorie.
\end{remark}

\begin{remark}
\label{morphisms-remark-pseudo-morphisms}
Il existe une variante de la Définition \ref{morphisms-definition-rational-map} où nous considérons
uniquement les morphismes $U \to Y$ définis sur des sous-schémas
ouverts schématiquement denses $U \subset X$. Nous utilisons le
Lemme \ref{morphisms-lemma-intersection-scheme-theoretically-dense} pour voir que nous obtenons une relation d'équivalence.
Une classe d'équivalence de ceux-ci est appelée
un {\it pseudo-morphisme de $X$
à $Y$}. Si $X$
est réduit, les deux notions coïncident.
\end{remark}






\section{Morphismes birationnels}
\label{morphisms-section-birational}

\noindent
Vous êtes peut-être habitué à la notion d'application birationnelle
de variétés ayant la propriété d'être un isomorphisme sur
un ouvert de la cible. Cependant, en général, un morphisme
birationnel peut ne pas être un isomorphisme sur
un ouvert non vide, voir Exemple \ref{morphisms-example-birational-not-iso-over-open}. Voici
la définition formelle.

\begin{definition}
\label{morphisms-definition-birational}
\begin{reference}
\cite[(2.2.9)]{EGA1}
\end{reference}
Soient $X$, $Y$ des schémas. Supposons que $X$ et $Y$ aient un nombre
fini de composantes irréductibles. On dit qu'un morphisme $f : X \to Y$ est {\it
birationnel} si
\begin{enumerate}
\item $f$ induit une bijection entre l'ensemble des points génériques des
composantes irréductibles de $X$ et l'ensemble des points génériques
des composantes irréductibles de $Y$, et
\item pour chaque point générique $\eta \in X$ d'une composante irréductible
de $X$,
l'homomorphisme d'anneaux locaux $\mathcal{O}_{Y, f(\eta)} \to \mathcal{O}_{X, \eta}$ est un
isomorphisme.
\end{enumerate}
\end{definition}

\noindent
Nous verrons ci-dessous que les fibres d'un morphisme birationnel au-dessus
des points génériques sont des singletons. De plus, nous verrons que dans la
plupart des cas rencontrés en pratique, l'existence d'un morphisme birationnel entre des schémas irréductibles
$X$ et $Y$ implique que $X$ et
$Y$ sont des schémas birationnels.

\begin{lemma}
\label{morphisms-lemma-birational-dominant}
Soit $f : X \to Y$ un morphisme de schémas ayant un nombre fini de composantes
irréductibles. Si $f$ est birationnel alors $f$
est dominant.
\end{lemma}

\begin{proof}
Cela découle du Lemme \ref{morphisms-lemma-generic-points-in-image-dominant} et de la
définition.
\end{proof}

\begin{lemma}
\label{morphisms-lemma-birational-generic-fibres}
Soit $f : X \to Y$ un morphisme birationnel de schémas ayant un
nombre fini de composantes irréductibles. Si $y \in Y$
est le point générique d'une composante
irréductible, alors le changement de base $X \times_Y \Spec(\mathcal{O}_{Y, y}) \to \Spec(\mathcal{O}_{Y, y})$ est un
isomorphisme.
\end{lemma}

\begin{proof}
Nous pouvons supposer que $Y = \Spec(B)$ est affine et irréductible.
Alors $X$ est irréductible aussi. Si nous prouvons
le résultat pour tout ouvert affine non vide $U \subset X$,
alors le résultat vaut pour $X$ (petit argument
omis). Par conséquent, nous
pouvons supposer que $X$ est affine aussi, disons $X = \Spec(A)$.
Soit $y \in Y$ correspondant au premier idéal minimal $\mathfrak q \subset B$.
D'après l'hypothèse, $A$ a un unique premier idéal minimal $\mathfrak p$
au-dessus de $\mathfrak q$ et $B_\mathfrak q \to A_\mathfrak p$ est un isomorphisme. Il
s'ensuit que $A_\mathfrak q \to \kappa(\mathfrak p)$ est surjectif, donc $\mathfrak p A_\mathfrak q$ est un idéal maximal.
D'autre part, $\mathfrak p A_\mathfrak q$ est le seul premier idéal minimal de $A_\mathfrak q$.
Nous concluons que $\mathfrak p A_\mathfrak q$ est le seul idéal premier de
$A_\mathfrak q$ et que $A_\mathfrak q = A_\mathfrak p$. Puisque
$A_\mathfrak q = A \otimes_B B_\mathfrak q$, le lemme s'ensuit.
\end{proof}

\begin{example}
\label{morphisms-example-birational-not-iso-over-open}
Voici deux exemples de morphismes birationnels qui ne sont pas des isomorphismes
sur aucun ouvert de la cible.

\medskip\noindent
Premier exemple. Soit $k$ un corps infini. Soit
$A = k[x]$. Soit $B = k[x, \{y_{\alpha}\}_{\alpha \in k}]/
((x-\alpha)y_\alpha, y_\alpha y_\beta)$. Il existe
une inclusion $A \subset B$ et une rétraction $B \to A$ définissant tous les $y_\alpha$
égaux à zéro. Le morphisme $\Spec(A) \to \Spec(B)$
ainsi que le morphisme $\Spec(B) \to \Spec(A)$
sont birationnels mais ne sont pas des
isomorphismes sur aucun ouvert.

\medskip\noindent Second
exemple. Soit $A$ un domaine. Soit $S \subset A$ un sous-ensemble multiplicatif
ne contenant pas
$0$. Avec $B = S^{-1}A$, le morphisme $f : \Spec(B) \to \Spec(A)$ est birationnel.
S'il existe un ouvert $U$ de $\Spec(A)$ tel
que $f^{-1}(U) \to U$ soit un isomorphisme, alors il existe un $a \in A$ tel
que chaque élément de $S$ devienne inversible dans la localisation
principale $A_a$. En prenant $A = \mathbf{Z}$ et $S$,
l'ensemble des entiers impairs donne un contre-exemple.
\end{example}

\begin{lemma}
\label{morphisms-lemma-birational-birational}
Soit $f : X \to Y$ un morphisme birationnel de schémas ayant un nombre fini de composantes
irréductibles sur un schéma de base $S$. Supposons que l'une des
conditions suivantes soit remplie
\begin{enumerate}
\item $f$ est localement de type fini et $Y$ réduit,
\item $f$ est localement de présentation finie.
\end{enumerate}
Alors il existe des ouverts denses $U \subset X$ et $V \subset Y$ tels que $f(U) \subset V$
et $f|_U : U \to V$ est un isomorphisme. En particulier, si $X$ et
$Y$ sont irréductibles, alors $X$ et $Y$
sont $S$-birationnels.
\end{lemma}

\begin{proof}
Il y a une réduction immédiate au cas où $X$ et $Y$ sont irréductibles
que nous omettons. De plus, après avoir encore restreint, nous pouvons
supposer que $X$ et $Y$ sont affines,
disons $X = \Spec(A)$ et $Y = \Spec(B)$. Par hypothèse, $A$, resp.\ $B$
a un unique idéal premier minimal $\mathfrak p$, resp.\ $\mathfrak q$,
le premier $\mathfrak p$
est au-dessus de $\mathfrak q$, et $B_\mathfrak q = A_\mathfrak p$.
Par le Lemme \ref{morphisms-lemma-birational-generic-fibres} nous avons $B_\mathfrak q = A_\mathfrak q = A_\mathfrak p$.

\medskip\noindent
Supposons que $B \to A$ soit de type fini, disons $A = B[x_1, \ldots, x_n]$.
Il existe un $b_i \in B$ et $g_i \in B \setminus \mathfrak q$ tels que $b_i/g_i$ s'envoie
à l'image de $x_i$ dans $A_\mathfrak q$. Par conséquent,
$b_i  - g_ix_i$ s'envoie sur zéro dans $A_{g_i'}$ pour
un certain $g_i' \in B \setminus \mathfrak q$. En posant $g = \prod g_i g'_i$ nous voyons que
$B_g \to A_g$ est surjectif. Si de plus $Y$ est réduit, alors
l'application $B_g \to B_\mathfrak q$ est injective et donc $B_g \to A_g$
est également injective. Cela prouve le cas (1).

\medskip\noindent
Démonstration de (2). Par l'argument donné dans le paragraphe précédent,
nous pouvons supposer que $B \to A$ est surjectif. Comme $f$
est localement de présentation finie, le noyau $J \subset B$
est un idéal de type fini. Disons $J = (b_1, \ldots, b_r)$. Puisque
$B_\mathfrak q = A_\mathfrak q$, il existe $g_i \in B \setminus \mathfrak q$ tels que $g_i b_i = 0$.
En posant $g = \prod g_i$, nous voyons que $B_g \to A_g$ est un
isomorphisme.
\end{proof}

\begin{lemma}
\label{morphisms-lemma-criterion-birational-finite-presentation}
Soit $S$ un schéma. Soient $X$ et $Y$ des schémas irréductibles
localement de présentation finie sur $S$. Soient $x \in X$ et $y \in Y$
les points génériques. Les conditions suivantes sont équivalentes
\begin{enumerate}
\item $X$ et $Y$ sont $S$-birationnels,
\item il existe des ouverts non vides de $X$ et $Y$
qui sont $S$-isomorphes, et
\item $x$ et $y$ s'envoient sur le même point
$s$ de $S$ et $\mathcal{O}_{X, x}$ et $\mathcal{O}_{Y, y}$ sont isomorphes
en tant qu'algèbres $\mathcal{O}_{S, s}$.
\end{enumerate}
\end{lemma}

\begin{proof}
Nous avons vu l'équivalence de (1) et (2) dans
le Lemme \ref{morphisms-lemma-birational-integral}. Il est immédiat
que (2) implique (3). Pour finir,
nous supposons que (3) est vrai et nous
prouvons (1). Par le Lemme \ref{morphisms-lemma-rational-map-finite-presentation}, il existe une
application rationnelle $f : U \to Y$
qui envoie $x \in U$ sur $y$
et induit l'isomorphisme donné $\mathcal{O}_{Y, y} \cong \mathcal{O}_{X, x}$. Ainsi $f$ est
un morphisme birationnel et induit donc un isomorphisme
sur des ouverts non
vides par le Lemme \ref{morphisms-lemma-birational-birational}. Cela termine
la démonstration.
\end{proof}

\begin{lemma}
\label{morphisms-lemma-common-open}
Soit $S$ un schéma. Soient $X$ et $Y$ des schémas intègres localement
de type fini sur $S$. Soient $x \in X$ et $y \in Y$ les points génériques. Les conditions
suivantes sont équivalentes
\begin{enumerate}
\item $X$ et $Y$ sont $S$-birationnels,
\item il existe des ouverts non vides de $X$ et $Y$ qui sont $S$-isomorphes, et
\item $x$ et $y$ ont la même image
$s \in S$ et $\kappa(x) \cong \kappa(y)$ en tant qu'extensions $\kappa(s)$.
\end{enumerate}
\end{lemma}

\begin{proof}
Nous avons vu l'équivalence de (1) et
(2) dans le Lemme \ref{morphisms-lemma-birational-integral}. Il
est immédiat que (2) implique
(3). Pour finir, nous supposons que
(3) est vrai et nous prouvons (1). Observez
que $\mathcal{O}_{X, x} = \kappa(x)$ et $\mathcal{O}_{Y, y} = \kappa(y)$
par l'Algèbre, Lemme \ref{algebra-lemma-minimal-prime-reduced-ring}. Par le Lemme
\ref{morphisms-lemma-rational-map-finite-presentation} il existe une application rationnelle
$f : U \to Y$ qui envoie $x \in U$
sur $y$ et induit
l'isomorphisme donné $\mathcal{O}_{Y, y} \cong \mathcal{O}_{X, x}$. Ainsi $f$ est
un morphisme birationnel et induit donc
un isomorphisme sur des ouverts
non vides par le Lemme \ref{morphisms-lemma-birational-birational}. Cela termine
la démonstration.
\end{proof}







\section{Morphismes génériquement finis}
\label{morphisms-section-generically-finite}

\noindent
Dans cette section, nous caractérisons les applications entre schémas
qui sont localement de type fini et qui sont ``génériquement fini'' d'une
certaine manière.

\begin{lemma}
\label{morphisms-lemma-generically-finite}
Soient $X$, $Y$
des schémas. Soit $f : X \to Y$ localement de type fini. Soit
$\eta \in Y$ un point générique d'une composante irréductible de $Y$.
Les conditions suivantes sont équivalentes :
\begin{enumerate}
\item l'ensemble $f^{-1}(\{\eta\})$ est fini,
\item il existe des ouverts affines $U_i \subset X$,
$i = 1, \ldots, n$ et $V \subset Y$ avec $f(U_i) \subset V$,
$\eta \in V$ et $f^{-1}(\{\eta\}) \subset \bigcup U_i$ tels
que chaque $f|_{U_i} : U_i \to V$ soit fini.
\end{enumerate}
Si $f$ est quasi-séparé, alors ceci est également équivalent à
\begin{enumerate}
\item[(3)] il existe des ouverts affines
$V \subset Y$ et $U \subset X$ avec $f(U) \subset V$,
$\eta \in V$ et $f^{-1}(\{\eta\}) \subset U$ tels
que $f|_U : U \to V$ soit fini.
\end{enumerate}
Si $f$ est quasi-compact et quasi-séparé, alors
ceux-ci sont également équivalents à
\begin{enumerate}
\item[(4)] il existe un ouvert affine $V \subset Y$, $\eta \in V$
tel que $f^{-1}(V) \to V$ soit fini.
\end{enumerate}
\end{lemma}

\begin{proof}
La question est locale sur la base. Par conséquent, nous pouvons remplacer $Y$ par
un voisinage affine de $\eta$, et nous pouvons et faisons l'hypothèse tout au long de
la démonstration ci-dessous que $Y$ est affine, disons $Y = \Spec(R)$.

\medskip\noindent
Il est clair que (2) implique (1).
Supposons que $f^{-1}(\{\eta\}) = \{\xi_1, \ldots, \xi_n\}$ soit fini. Choisissez des ouverts affines
$U_i \subset X$ avec $\xi_i \in U_i$. Par Algèbre, Lemme \ref{algebra-lemma-generically-finite}
nous voyons qu'après avoir remplacé $Y$ par
un ouvert principal dans $Y$, chacun
des morphismes $U_i \to Y$ est fini. En d'autres
termes, (2) est vrai.

\medskip\noindent Il
est clair que (3) implique (1). Supposons que $f$ soit quasi-séparé et
(1). Écrivons $f^{-1}(\{\eta\}) = \{\xi_1, \ldots, \xi_n\}$. Il n'y a pas de spécialisations parmi les $\xi_i$
d'après le Lemme \ref{morphisms-lemma-finite-fibre}. Puisque chaque $\xi_i$ s'envoie sur le
point générique $\eta$ d'une composante
irréductible de $Y$, il ne peut y avoir de spécialisation non triviale
$\xi \leadsto \xi_i$ dans $X$ (puisque $\xi$ s'enverrait
également sur $\eta$). Nous concluons
que chaque $\xi_i$ est un point générique d'une
composante irréductible de $X$.
Puisque $Y$ est affine et $f$ quasi-séparé, nous voyons que
$X$ est quasi-séparé. D'après les Propriétés, Lemme \ref{properties-lemma-maximal-points-affine}.
nous pouvons trouver un ouvert affine $U \subset X$ contenant chaque $\xi_i$.
Par l'Algèbre, Lemme \ref{algebra-lemma-generically-finite} nous voyons qu'après avoir remplacé
$Y$ par un ouvert principal dans $Y$
les morphismes $U \to Y$ sont finis. En
autres termes (3) est vrai.

\medskip\noindent Il
est clair que (4) implique toutes les conditions (1) -- (3) sans autre hypothèse
sur $f$. Supposons que $f$ soit quasi-compact et quasi-séparé. Nous
devons montrer que les conditions équivalentes (1) -- (3) impliquent
(4). Soient $U$, $V$ comme dans (3). Remplacez $Y$ par $V$.
Puisque $f$ est quasi-compact et $Y$ est quasi-compact (étant
affine), nous voyons que $X$ est quasi-compact. Par conséquent,
$Z = X \setminus U$ est quasi-compact, donc le morphisme $f|_Z : Z \to Y$ est quasi-compact.
Par la construction de $Z$,
nous voyons que $\eta \not \in f(Z)$. Par conséquent, d'après le Lemme
\ref{morphisms-lemma-quasi-compact-generic-point-not-in-image}, nous voyons qu'il existe un
ouvert affine voisinage $V'$ de $\eta$ dans $Y$ tel que $f^{-1}(V') \cap Z = \emptyset$.
Alors nous avons $f^{-1}(V') \subset U$ et cela signifie
que $f^{-1}(V') \to V'$ est fini.
\end{proof}

\begin{example}
\label{morphisms-example-counter-generically-finite}
Soit $A = \prod_{n \in \mathbf{N}} \mathbf{F}_2$. Chaque élément de $A$ est idempotent.
Ainsi, tout idéal premier est maximal avec corps résiduel $\mathbf{F}_2$.
Ainsi, la topologie sur $X = \Spec(A)$
est totalement déconnectée et quasi-compacte. Les applications
de projection $A \to \mathbf{F}_2$ définissent des points ouverts de $\Spec(A)$. Il
ne peut pas être le cas que tous les points de $X$ soient ouverts
puisque $X$ est quasi-compact. Soit $x \in X$ un point fermé
qui n'est pas ouvert. Alors nous pouvons former un schéma $Y$ constitué
de deux copies de $X$ collées le long de $X \setminus \{x\}$. En
d'autres termes, il s'agit de $X$
avec $x$ doublé, comparer Schémas, Exemple \ref{schemes-example-affine-space-zero-doubled}. Le
morphisme
$f : Y \to X$ est quasi-compact, de type fini et a des fibres finies mais n'est
pas quasi-séparé. Le point
$x \in X$ est un point générique d'une composante irréductible de $X$
(puisque $X$ est totalement déconnecté). Mais les propriétés (3)
et (4) du Lemme \ref{morphisms-lemma-generically-finite} ne sont pas vérifiées. La raison est que pour tout
voisinage ouvert $x \in U \subset X$, l'image réciproque $f^{-1}(U)$ n'est pas affine
parce que les fonctions sur $f^{-1}(U)$ ne peuvent pas séparer les deux
points situés au-dessus de $x$ (démonstration omise ; c'est un
bel exercice). Ainsi, la condition que $f$ soit quasi-séparé est nécessaire
dans les parties (3) et (4) du lemme.
\end{example}

\begin{remark}
\label{morphisms-remark-quasi-finite-finite-over-dense-open}
Une alternative
au Lemme \ref{morphisms-lemma-generically-finite} est
l'énoncé qu'un morphisme quasi-fini est fini
sur un ouvert dense de la cible. Cela
sera démontré
dans Compléments sur les morphismes, Lemme \ref{more-morphisms-lemma-quasi-finite-finite-over-dense-open}.
\end{remark}

\begin{lemma}
\label{morphisms-lemma-quasi-finiteness-over-generic-point}
Soient $X$, $Y$ des schémas. Soit $f : X \to Y$ localement de type fini. Soient $X^0$,
resp. \ $Y^0$ l'ensemble des points génériques des composantes irréductibles
de $X$, resp. \ $Y$. Soit $\eta \in Y^0$. Les conditions suivantes sont
équivalentes
\begin{enumerate}
\item $f^{-1}(\{\eta\}) \subset X^0$,
\item $f$ est quasi-fini en tous les points au-dessus de $\eta$,
\item $f$ est quasi-fini en tous $\xi \in X^0$ au-dessus de $\eta$.
\end{enumerate}
\end{lemma}

\begin{proof}
La condition (1) implique qu'il n'y a pas de spécialisations parmi les points
de la fibre $X_\eta$. Par conséquent, (2)
est vrai par le Lemme \ref{morphisms-lemma-quasi-finite-at-point-characterize}. L'implication (2) $\Rightarrow$
(3) est immédiate. Puisque $\eta$ est un point
générique de $Y$, les points génériques de $X_\eta$ sont des points
génériques de $X$. Par conséquent,
(3) et le Lemme \ref{morphisms-lemma-quasi-finite-at-point-characterize} impliquent que les points génériques
de $X_\eta$ sont également fermés. Ainsi, tous les
points de $X_\eta$ sont génériques et nous voyons que (1) est vrai.
\end{proof}

\begin{lemma}
\label{morphisms-lemma-finite-over-dense-open}
Soient $X$, $Y$ des schémas. Soit $f : X \to Y$ localement de type fini. Soient $X^0$,
resp. \ $Y^0$ l'ensemble des points génériques des composantes irréductibles de
$X$, resp. \ $Y$. Supposons
\begin{enumerate}
\item $X^0$ et $Y^0$ sont finis et $f^{-1}(Y^0) = X^0$,
\item soit $f$ est quasi-compact, soit $f$ est séparé.
\end{enumerate}
Alors il existe un ouvert dense $V \subset Y$
tel que $f^{-1}(V) \to V$ soit fini.
\end{lemma}

\begin{proof}
Puisque $Y$ a un nombre fini de composantes irréductibles, nous pouvons trouver un ouvert
dense qui est une union disjointe de ses composantes irréductibles. Ainsi, nous pouvons supposer
que $Y$ est un affine irréductible avec un point générique $\eta$. Alors la fibre au-dessus
de $\eta$ est finie car $X^0$ est fini.

\medskip\noindent
Supposons que $f$ soit séparé et $Y$ affine irréductible.
Choisissez $V \subset Y$ et $U \subset X$ comme dans le Lemme \ref{morphisms-lemma-generically-finite} partie
(3). Puisque $f|_U : U \to V$ est fini, nous voyons que $U \subset f^{-1}(V)$
est fermé ainsi qu'ouvert (Lemmes \ref{morphisms-lemma-image-proper-scheme-closed} et \ref{morphisms-lemma-finite-proper}). Ainsi $f^{-1}(V) = U \amalg W$
pour un certain sous-schéma ouvert $W$ de $X$. Cependant,
puisque $U$ contient tous les points génériques de $X$,
nous concluons que $W = \emptyset$ comme souhaité.

\medskip\noindent
Supposons que $f$ soit quasi-compact et $Y$ affine
irréductible. Alors $X$ est quasi-compact, donc il existe
un sous-schéma ouvert
dense $U \subset X$ qui est séparé (Propriétés, Lemme \ref{properties-lemma-quasi-compact-dense-open-separated}). Comme
l'ensemble des points génériques $X^0$ est fini, nous voyons que $X^0 \subset U$.
Ainsi $\eta \not \in f(X \setminus U)$. Comme $X \setminus U \to Y$ est quasi-compact, nous
concluons qu'il existe un ouvert non vide $V \subset Y$ tel que $f^{-1}(V) \subset U$,
voir Lemme \ref{morphisms-lemma-quasi-compact-dominant}. Après avoir remplacé $X$ par $f^{-1}(V)$
et $Y$ par $V$, nous nous réduisons au cas
séparé que nous avons traité dans le paragraphe précédent.
\end{proof}

\begin{lemma}
\label{morphisms-lemma-birational-isomorphism-over-dense-open}
Soient $X$, $Y$ des schémas. Soit $f : X \to Y$ un morphisme birationnel
entre des schémas qui ont un nombre fini de composantes irréductibles.
Supposons
\begin{enumerate}
\item soit $f$ est quasi-compact, soit $f$ est séparé, et
\item soit $f$ est localement de type fini et $Y$ est réduit, soit $f$
est localement de présentation finie.
\end{enumerate}
Alors, il existe un ouvert dense $V \subset Y$
tel que $f^{-1}(V) \to V$ soit un isomorphisme.
\end{lemma}

\begin{proof}
Par le Lemme \ref{morphisms-lemma-finite-over-dense-open} nous pouvons supposer que $f$ est fini. Comme $Y$
a un nombre fini de composantes irréductibles, nous pouvons trouver un ouvert dense
qui est une union disjointe de ses composantes irréductibles. Ainsi nous pouvons
supposer que $Y$ est irréductible. Par le Lemme \ref{morphisms-lemma-birational-birational} nous trouvons un ouvert
non vide $U \subset X$ tel que $f|_U : U \to Y$ est une immersion ouverte. Après avoir retiré
le sous-ensemble fermé (comme $f$ est fini)
$f(X \setminus U)$ de $Y$ nous voyons que $f$ est un isomorphisme.
\end{proof}

\begin{lemma}
\label{morphisms-lemma-finite-degree}
Soient $X$, $Y$ des schémas intègres.
Soit $f : X \to Y$ localement de type fini. Supposons que
$f$ est dominant. Les conditions
suivantes sont équivalentes :
\begin{enumerate}
\item l'extension $R(Y) \subset R(X)$ a un degré de transcendance
$0$,
\item l'extension $R(Y) \subset R(X)$ est finie,
\item il existe des ouverts affines non vides $U \subset X$
et $V \subset Y$ tels que $f(U) \subset V$
et $f|_U : U \to V$ soit fini, et
\item le point générique de $X$ est le seul point de $X$ qui s'envoie
sur le point générique de $Y$.
\end{enumerate}
Si $f$ est séparé ou si $f$ est quasi-compact, alors cela équivaut
également à
\begin{enumerate}
\item[(5)] il existe un ouvert affine non vide $V \subset Y$ tel
que $f^{-1}(V) \to V$ soit fini.
\end{enumerate}
\end{lemma}

\begin{proof}
Choisissez des ouverts affines $\Spec(A) = U \subset X$ et $\Spec(R) = V \subset Y$
tels que $f(U) \subset V$. Alors $R$ et $A$
sont des anneaux intègres par définition. L'homomorphisme
d'anneaux $R \to A$ est
de type fini (Lemme \ref{morphisms-lemma-locally-finite-type-characterize}). D'après le Lemme \ref{morphisms-lemma-dominant-finite-number-irreducible-components},
le point générique de $X$ s'envoie sur le point générique
de $Y$, donc $R \to A$ est injectif. Soit
$K = R(Y)$ le corps des fractions de $R$ et $L = R(X)$ le
corps des fractions de $A$. Alors $L/K$ est une
extension de corps de type fini. Ainsi, nous voyons que (1)
est équivalent à (2).

\medskip\noindent
Supposons que (2) soit vrai. Soient $x_1, \ldots, x_n \in A$ des générateurs de
$A$ sur $R$. D'après l'hypothèse, il existe des polynômes
non nuls $P_i(X) \in R[X]$ tels que $P_i(x_i) = 0$. Soit $f_i \in R$ le coefficient
principal de $P_i$. Nous concluons alors
que $R_{f_1 \ldots f_n} \to A_{f_1 \ldots f_n}$ est fini, c'est-à-dire que (3) est vrai. Notez que (3) implique
(2). Nous voyons donc maintenant que (1), (2) et (3) sont tous
équivalents.

\medskip\noindent
Soit $\eta$ le point générique de $X$, et soit $\eta' \in Y$
le point générique de $Y$.
Supposons (4). Alors $\dim_\eta(X_{\eta'}) = 0$ et nous voyons que $R(X) = \kappa(\eta)$
a un degré de transcendance $0$ sur $R(Y) = \kappa(\eta')$
par le Lemme \ref{morphisms-lemma-dimension-fibre-at-a-point}. En d'autres
termes, (1) est vrai. Supposons les conditions équivalentes (1),
(2) et (3). Supposons que $x \in X$ soit un point se
mappant sur $\eta'$. Comme $x$
est une spécialisation de $\eta$, cela donne des inclusions
$R(Y) \subset \mathcal{O}_{X, x} \subset R(X)$, ce qui implique que $\mathcal{O}_{X, x}$
est un corps, voir Algèbre, Lemme \ref{algebra-lemma-integral-over-field}.
Ainsi $x = \eta$. Nous voyons donc que (1) -- (4) sont
tous équivalents.

\medskip\noindent
Il est clair que (5) implique (3) sans hypothèses supplémentaires
sur $f$. Il reste à prouver que si $f$ est soit
séparé soit quasi-compact, alors les conditions équivalentes (1)
-- (4) impliquent (5). Cela résulte du lemme \ref{morphisms-lemma-finite-over-dense-open}.
\end{proof}

\begin{definition}
\label{morphisms-definition-degree}
Soient $X$ et $Y$ des
schémas intègres. Soit $f : X \to Y$ localement de type
fini et dominant. Supposons $[R(X) : R(Y)] < \infty$, ou toute autre des conditions
équivalentes (1) -- (4) du Lemme \ref{morphisms-lemma-finite-degree}.
Alors l'entier positif
$$
\deg(X/Y) = [R(X) : R(Y)]
$$
est appelé le degré {\it de $X$ sur $Y$}.
\end{definition}

\noindent
Il est possible d'étendre cette notion à un morphisme $f : X \to Y$
si (a) $Y$ est intègre avec point générique $\eta$,
(b) $f$ est localement de type fini, et (c) $f^{-1}(\{\eta\})$ est fini. Dans ce
cas, nous pouvons définir
$$
\deg(X/Y)
=
\sum\nolimits_{\xi \in X, \ f(\xi) = \eta}
\dim_{R(Y)} (\mathcal{O}_{X, \xi}).
$$
À savoir, étant donné que $R(Y) = \kappa(\eta) = \mathcal{O}_{Y, \eta}$ (Lemme \ref{morphisms-lemma-integral-scheme-rational-functions})
les dimensions ci-dessus sont finies par
le Lemme \ref{morphisms-lemma-generically-finite} ci-dessus.
Cependant, pour la plupart des
applications, la définition donnée ci-dessus est
la bonne.

\begin{lemma}
\label{morphisms-lemma-degree-composition}
Soient $X$, $Y$, $Z$
des schémas intègres. Soient $f : X \to Y$ et $g : Y \to Z$ des morphismes
dominants localement de type fini. Supposons que $[R(X) : R(Y)] < \infty$
et $[R(Y) : R(Z)] < \infty$. Alors
$$
\deg(X/Z) = \deg(X/Y) \deg(Y/Z).
$$
\end{lemma}

\begin{proof}
Cela provient de la multiplicativité des degrés
dans les tours d'extensions
finies de corps, voir Corps, Lemme \ref{fields-lemma-multiplicativity-degrees}.
\end{proof}

\begin{remark}
\label{morphisms-remark-definition-generically-finite}
Soit $f : X \to Y$ un morphisme de schémas qui est localement de type fini. Il y a (au moins) deux
propriétés que nous pourrions utiliser pour définir les morphismes {\it
génériquement finis}. Celles-ci correspondent à la question de savoir
si vous voulez que la propriété soit locale sur la source ou locale sur le but :
\begin{enumerate}
\item (Locale sur le but ; suggéré par Ravi Vakil.) Supposons
que tout ouvert quasi-compact de $Y$ ait un
nombre fini de composantes irréductibles (par exemple si $Y$ est localement
noethérien). L'exigence est que l'image réciproque de chaque point générique
soit finie, voir Lemme \ref{morphisms-lemma-generically-finite}.
\item (Local sur la source.) L'exigence est qu'il existe un ouvert
dense $U \subset X$ tel que $U \to Y$ soit localement quasi-fini.
\end{enumerate}
Dans le cas (1), l'exigence peut être formulée sans la condition auxiliaire sur $Y$,
mais ne donne probablement pas la bonne notion pour les schémas en général.
La propriété (2) telle que formulée n'implique pas que les fibres au-dessus des points génériques
soient finies ; cependant, si $f$ est quasi-compact
et que $Y$ est comme dans (1), alors c'est le cas.
\end{remark}

\begin{definition}
\label{morphisms-definition-modification}
Soit $X$ un schéma intègre. Une modification {\it de $X$}
est un morphisme birationnel propre $f : X' \to X$ avec $X'$
intègre.
\end{definition}

\noindent
Soit $f : X' \to X$ une modification comme dans la définition. D'après
le Lemme \ref{morphisms-lemma-finite-degree} il existe un $U \subset X$ non vide tel que $f^{-1}(U) \to U$ soit
fini. Par platitude générique (Proposition \ref{morphisms-proposition-generic-flatness})
nous pouvons supposer que $f^{-1}(U) \to U$ est plat et
de présentation finie. Ainsi $f^{-1}(U) \to U$ est fini localement libre (Lemme
\ref{morphisms-lemma-finite-flat}). Puisque $f$ est birationnel, le degré de $X'$ sur
$X$ est $1$. Par conséquent, $f^{-1}(U) \to U$ est fini
localement libre de degré $1$, autrement dit c'est un isomorphisme.
Ainsi nous pouvons {\it redéfinir}
une modification pour être un morphisme propre $f : X' \to X$ de schémas
intègres tel que $f^{-1}(U) \to U$ soit un isomorphisme pour un ouvert non vide
$U \subset X$.

\begin{definition}
\label{morphisms-definition-alteration}
\begin{reference}
\cite[Définition 2.20]{alterations}
\end{reference}
Soit $X$ un schéma intègre. Une {\it altération
de $X$} est un morphisme propre dominant $f : Y \to X$ avec $Y$
intègre tel que $f^{-1}(U) \to U$ soit fini pour un certain ouvert non vide $U \subset X$.
\end{definition}

\noindent
Ceci est la définition donnée dans \cite{alterations}, sauf que ici nous ne supposons
pas que $X$ et $Y$ soient noethériens. En argumentant comme
ci-dessus, nous voyons qu'une altération est un morphisme propre dominant $f : Y \to X$
de schémas intègres qui induit une extension finie de
corps de fonctions, c'est-à-dire, telle que les conditions équivalentes
du Lemme \ref{morphisms-lemma-finite-degree} soient vérifiées.



\section{La formule de dimension}
\label{morphisms-section-dimension-formula}

\noindent
Pour les morphismes entre schémas noethériens, nous pouvons dire un peu plus
sur les dimensions des anneaux locaux. Voici un résultat important (et pas si
difficile à prouver). Rappelons que $R(X)$ désigne le corps des fonctions d'un
schéma intègre $X$.

\begin{lemma}
\label{morphisms-lemma-dimension-formula}
Soit $S$ un schéma.
Soit $f : X \to S$ un morphisme de schémas. Soit
$x \in X$, et posons $s = f(x)$.
Supposons
\begin{enumerate}
\item $S$ localement noethérien,
\item $f$ localement de type fini,
\item $X$ et $S$ intègres, et
\item $f$ dominant.
\end{enumerate}
Nous avons
\begin{equation}
\label{morphisms-equation-dimension-formula}
\dim(\mathcal{O}_{X, x})
\leq
\dim(\mathcal{O}_{S, s}) + \text{trdeg}_{R(S)}R(X)
- \text{trdeg}_{\kappa(s)} \kappa(x).
\end{equation}
De plus, l'égalité est vérifiée si $S$ est universellement caténaire.
\end{lemma}

\begin{proof}
L'énoncé algébrique correspondant
est Algèbre, Lemme \ref{algebra-lemma-dimension-formula}.
\end{proof}

\begin{lemma}
\label{morphisms-lemma-dimension-formula-general}
Soit $S$ un schéma. Soit $f : X \to S$ un morphisme de schémas. Soit $x \in X$,
et posons $s = f(x)$. Supposons que $S$ soit localement noethérien et que $f$
soit localement de type fini. Nous
avons
\begin{equation}
\label{morphisms-equation-dimension-formula-general}
\dim(\mathcal{O}_{X, x})
\leq
\dim(\mathcal{O}_{S, s}) + E - \text{trdeg}_{\kappa(s)} \kappa(x).
\end{equation}
où $E$ est le maximum de $\text{trdeg}_{\kappa(f(\xi))}(\kappa(\xi))$ lorsque $\xi$ parcourt les
points génériques des composantes irréductibles de $X$ contenant
$x$.
\end{lemma}

\begin{proof}
Soient $X_1, \ldots, X_n$ les composantes irréductibles de $X$
contenant $x$ munies de leur
structure de schéma induite réduite. Celles-ci
correspondent aux idéaux premiers minimaux
$\mathfrak q_i$ de $\mathcal{O}_{X, x}$ et il y en
a donc un nombre fini (Schémas, Lemme \ref{schemes-lemma-specialize-points} et
Algèbre,
Lemme \ref{algebra-lemma-Noetherian-irreducible-components}). Alors $\dim(\mathcal{O}_{X, x}) = \max \dim(\mathcal{O}_{X, x}/\mathfrak q_i)
= \max \dim(\mathcal{O}_{X_i, x})$. Les $\xi$
apparaissant dans la définition de $E$ sont exactement les
points génériques $\xi_i \in X_i$. Soit $Z_i = \overline{\{f(\xi_i)\}} \subset S$ muni de la
structure de schéma induite réduite.
La composition $X_i \to X \to S$ se factorise
par $Z_i$ (Schémas, Lemme \ref{schemes-lemma-map-into-reduction}).
Ainsi nous pouvons appliquer
la formule de dimension (Lemme \ref{morphisms-lemma-dimension-formula})
pour voir que $\dim(\mathcal{O}_{X_i, x}) \leq \dim(\mathcal{O}_{Z_i, x}) +
\text{trdeg}_{\kappa(f(\xi))}(\kappa(\xi)) -
\text{trdeg}_{\kappa(s)} \kappa(x)$.
En combinant le tout, nous obtenons le lemme.
\end{proof}

\noindent
Une application est la construction d'une fonction de dimension
sur tout schéma de type fini sur un schéma universellement
caténaire doté d'une fonction de dimension.
Pour la définition des fonctions
de dimension, voir Topologie, Définition \ref{topology-definition-dimension-function}.

\begin{lemma}
\label{morphisms-lemma-dimension-function-propagates}
Soit $S$ un schéma localement noethérien et universellement caténaire.
Soit $\delta : S \to \mathbf{Z}$ une fonction de dimension. Soit $f : X \to S$
un morphisme de schémas. Supposons $f$
localement de type fini. Alors
l'application
\begin{align*}
\delta = \delta_{X/S} : X & \longrightarrow \mathbf{Z} \\
x & \longmapsto \delta(f(x)) + \text{trdeg}_{\kappa(f(x))} \kappa(x)
\end{align*}
est une fonction de dimension sur $X$.
\end{lemma}

\begin{proof}
Soit $f : X \to S$ localement de type fini.
Soient $x \leadsto y$, $x \not = y$ une spécialisation dans
$X$. Nous devons montrer que $\delta_{X/S}(x) > \delta_{X/S}(y)$ et
que $\delta_{X/S}(x) = \delta_{X/S}(y) + 1$ si $y$ est une spécialisation
immédiate de $x$.

\medskip\noindent Choisissez
un ouvert affine $V \subset S$ contenant l'image de $y$
et choisissons un ouvert affine $U \subset X$ dont l'image est contenue dans $V$
et contenant $y$. Nous pouvons clairement remplacer
$X$ par $U$ et $S$ par $V$.
Ainsi nous pouvons supposer que $X = \Spec(A)$ et $S = \Spec(R)$
et que $f$ est donné par un homomorphisme d'anneaux
$R \to A$. L'anneau $R$ est universellement
caténaire (Lemme \ref{morphisms-lemma-universally-catenary-local}) et l'application
$R \to A$ est de type fini (Lemme \ref{morphisms-lemma-locally-finite-type-characterize}).

\medskip\noindent
Soit $\mathfrak q \subset A$ l'idéal premier correspondant au point $x$
et soit $\mathfrak p \subset R$ l'idéal premier correspondant à $f(x)$.
La restriction $\delta'$ de $\delta$ à $S' = \Spec(R/\mathfrak p) \subset S$ est une fonction
de dimension. L'anneau $R/\mathfrak p$ est universellement
caténaire. La restriction de $\delta_{X/S}$ à $X' = \Spec(A/\mathfrak q)$ est
clairement égale à la fonction $\delta_{X'/S'}$ construite à l'aide de
la fonction de dimension $\delta'$. Par conséquent, on peut
supposer, en plus de ce qui précède, que $R \subset A$ sont des
domaines, en d'autres termes que $X$ et $S$ sont
des schémas intègres, et que $x$ est le point générique
de $X$ et $f(x)$ est le point générique de $S$.

\medskip\noindent
Notez que $\mathcal{O}_{X, x} = R(X)$ et que puisque $x \leadsto y$,
$x \not = y$, le spectre de $\mathcal{O}_{X, y}$ a au
moins
deux points (Schémas, Lemme \ref{schemes-lemma-specialize-points})
donc $\dim(\mathcal{O}_{X, y}) > 0$. Si $y$
est une spécialisation immédiate
de $x$, alors $\Spec(\mathcal{O}_{X, y}) = \{x, y\}$
et $\dim(\mathcal{O}_{X, y}) = 1$.

\medskip\noindent Écrivez
$s = f(x)$ et $t = f(y)$. Nous calculons
\begin{align*}
\delta_{X/S}(x) - \delta_{X/S}(y)
& =
\delta(s) + \text{trdeg}_{\kappa(s)} \kappa(x)
- \delta(t) - \text{trdeg}_{\kappa(t)} \kappa(y) \\
& =
\delta(s) - \delta(t) +
\text{trdeg}_{R(S)} R(X) - \text{trdeg}_{\kappa(t)} \kappa(y) \\
& =
\delta(s) - \delta(t) + \dim(\mathcal{O}_{X, y})
- \dim(\mathcal{O}_{S, t})
\end{align*}
où nous utilisons l'égalité dans (\ref{morphisms-equation-dimension-formula}) à
la dernière étape. Puisque $\delta$ est une fonction
de dimension sur le schéma $S$
et $s \in S$ est le point générique, la différence $\delta(s) - \delta(t)$
est égale à $\text{codim}(\overline{\{t\}}, S)$ par Topologie, Lemme
\ref{topology-lemma-dimension-function-catenary}. C'est égal à $\dim(\mathcal{O}_{S, t})$
par Propriétés,
Lemme \ref{properties-lemma-codimension-local-ring}. Ainsi Nous
concluons que
$$
\delta_{X/S}(x) - \delta_{X/S}(y) = \dim(\mathcal{O}_{X, y})
$$
et le lemme découle de ce que nous avons dit ci-dessus à propos de $\dim(\mathcal{O}_{X, y})$.
\end{proof}

\noindent
Une autre application de la formule de la dimension est que la dimension ne change pas sous
``altérations'' (à définir plus tard).

\begin{lemma}
\label{morphisms-lemma-alteration-dimension}
Soit $f : X \to Y$ un morphisme de schémas. Supposons que
\begin{enumerate}
\item $Y$ soit localement noethérien,
\item $X$ et $Y$ soient des schémas intègres,
\item $f$ soit dominant, et
\item $f$ soit localement de type fini.
\end{enumerate}
Alors nous avons
$$
\dim(X) \leq \dim(Y) + \text{trdeg}_{R(Y)} R(X).
$$
Si $f$ est fermé\footnote{Par exemple si $f$ est propre,
voir Définition \ref{morphisms-definition-proper}.} alors l'égalité est vérifiée.
\end{lemma}

\begin{proof}
Soit $f : X \to Y$ comme dans le
lemme. Soit $\xi_0 \leadsto \xi_1 \leadsto \ldots \leadsto \xi_e$ une suite de spécialisations
dans $X$. Posons $x = \xi_e$ et
$y = f(x)$. Remarquons que $e \leq \dim(\mathcal{O}_{X, x})$ tel que les
spécialisations données apparaissent dans
le spectre de $\mathcal{O}_{X, x}$, voir Schémas, Lemme \ref{schemes-lemma-specialize-points}. Par la
formule de dimension, Lemme \ref{morphisms-lemma-dimension-formula}, nous voyons
que
\begin{align*}
e & \leq \dim(\mathcal{O}_{X, x}) \\
& \leq \dim(\mathcal{O}_{Y, y}) +
\text{trdeg}_{R(Y)} R(X) - \text{trdeg}_{\kappa(y)} \kappa(x) \\
& \leq \dim(\mathcal{O}_{Y, y}) + \text{trdeg}_{R(Y)} R(X)
\end{align*}
Par conséquent, nous concluons que $e \leq \dim(Y) + \text{trdeg}_{R(Y)} R(X)$ comme souhaité.

\medskip\noindent
Ensuite, supposons que $f$
soit également fermé. Supposons que $\overline{\xi}_0 \leadsto \overline{\xi}_1 \leadsto \ldots
\leadsto \overline{\xi}_d$ soit
une suite de spécialisations dans $Y$. Nous voulons montrer que $\dim(X) \geq d + r$. Nous pouvons
supposer que $\overline{\xi}_0 = \eta$ est le point générique de $Y$.
La fibre générique $X_\eta$ est un schéma localement de
type fini sur $\kappa(\eta) = R(Y)$. Elle est non vide puisque $f$
est dominant. Ainsi, d'après le Lemme \ref{morphisms-lemma-ubiquity-Jacobson-schemes}, c'est
un Schéma de Jacobson. Donc, d'après le Lemme
\ref{morphisms-lemma-jacobson-finite-type-points}, nous pouvons trouver un point fermé $\xi_0 \in X_\eta$
et l'extension $\kappa(\eta) \subset \kappa(\xi_0)$ est
une extension finie. Notez
que $\mathcal{O}_{X, \xi_0} = \mathcal{O}_{X_\eta, \xi_0}$ parce que $\eta$ est le point générique
de $Y$. Ainsi, nous voyons que
$\dim(\mathcal{O}_{X, \xi_0}) = r$ par le Lemme \ref{morphisms-lemma-dimension-formula} appliqué au schéma $X_\eta$
sur le schéma universellement caténaire $\Spec(\kappa(\eta))$
(voir Lemme \ref{morphisms-lemma-ubiquity-uc}) et le point $\xi_0$. Cela
signifie que nous pouvons trouver $\xi_{-r} \leadsto \ldots \leadsto \xi_{-1} \leadsto \xi_0$
dans $X$. D'autre part, comme $f$
est des spécialisations fermées, on peut relever
le long de $f$,
voir Topologie, Lemme \ref{topology-lemma-closed-open-map-specialization}. Ainsi, comme $\xi_0$
se situe
au-dessus de $\eta = \overline{\xi}_0$, nous pouvons trouver des spécialisations
$\xi_0 \leadsto \xi_1 \leadsto \ldots \leadsto \xi_d$ au-dessus de $\overline{\xi}_0 \leadsto \overline{\xi}_1 \leadsto \ldots
\leadsto \overline{\xi}_d$. En
d'autres termes, nous avons
$$
\xi_{-r} \leadsto \ldots \leadsto \xi_{-1} \leadsto \xi_0
\leadsto \xi_1 \leadsto \ldots \leadsto \xi_d
$$
ce qui signifie que $\dim(X) \geq d + r$ comme souhaité.
\end{proof}

\begin{lemma}
\label{morphisms-lemma-alteration-dimension-general}
Soit $f : X \to Y$ un morphisme de schémas. Supposons que $Y$ soit
localement noethérien et que $f$ soit localement de type fini.
Alors
$$
\dim(X) \leq \dim(Y) + E
$$
où $E$ est le supremum de $\text{trdeg}_{\kappa(f(\xi))}(\kappa(\xi))$ lorsque $\xi$ parcourt
les points génériques des composantes irréductibles de
$X$.
\end{lemma}

\begin{proof}
Conséquence immédiate du Lemme \ref{morphisms-lemma-dimension-formula-general} et
des Propriétés, Lemme \ref{properties-lemma-dimension}.
\end{proof}








\section{Normalisation relative}
\label{morphisms-section-normalization-X-in-Y}

\noindent
Dans cette section, nous construisons la normalisation {\it d'un schéma dans un autre }.

\begin{lemma}
\label{morphisms-lemma-integral-closure}
Soit $X$ un schéma. Soit $\mathcal{A}$ un faisceau quasi-cohérent
de $\mathcal{O}_X$-algèbres. Le sous-faisceau $\mathcal{A}' \subset \mathcal{A}$ défini
par la règle
$$
U \longmapsto \{f \in \mathcal{A}(U) \mid
f_x \in \mathcal{A}_x \text{ entier sur } \mathcal{O}_{X, x}
\text{ pour tout }x \in U\}
$$
est une $\mathcal{O}_X$-algèbre quasi-cohérente, le germe $\mathcal{A}'_x$
est la clôture intégrale de $\mathcal{O}_{X, x}$ dans $\mathcal{A}_x$,
et pour tout ouvert affine $U \subset X$
l'anneau $\mathcal{A}'(U) \subset \mathcal{A}(U)$ est la clôture
intégrale de $\mathcal{O}_X(U)$ dans $\mathcal{A}(U)$.
\end{lemma}

\begin{proof}
Il s'agit d'un sous-faisceau par la nature locale des conditions.
C'est une $\mathcal{O}_X$-algèbre
d'après Algèbre, Lemme \ref{algebra-lemma-integral-closure-is-ring}. Soit $U \subset X$ un
ouvert affine. Soit $U = \Spec(R)$ et supposons que $\mathcal{A}$
soit le faisceau quasi-cohérent associé à la $R$-algèbre
$A$. Alors, d'après Algèbre,
Lemme \ref{algebra-lemma-integral-closure-stalks}, la valeur de $\mathcal{A}'$ sur $U$
est donnée par la clôture intégrale $A'$ de
$R$ dans $A$. Cela prouve la dernière assertion
du lemme. Pour prouver que $\mathcal{A}'$ est quasi-cohérent,
il suffit de montrer que $\mathcal{A}'(D(f)) = A'_f$. Cela résulte du fait
que la clôture intégrale et la localisation commutent, voir
Algèbre, Lemme \ref{algebra-lemma-integral-closure-localize}. Le même fait montre
que les fibres sont comme annoncé.
\end{proof}

\begin{definition}
\label{morphisms-definition-integral-closure}
Soit $X$ un schéma. Soit $\mathcal{A}$ un faisceau quasi-cohérent
d'$\mathcal{O}_X$-algèbres. La clôture {\it intégrale
de $\mathcal{O}_X$ dans $\mathcal{A}$} est la $\mathcal{O}_X$-sous-algèbre
quasi-cohérente $\mathcal{A}' \subset \mathcal{A}$ construite
dans le Lemme \ref{morphisms-lemma-integral-closure} ci-dessus.
\end{definition}

\noindent
Dans le cadre de la définition ci-dessus, nous pouvons considérer le morphisme
des spectres relatifs
$$
\xymatrix{
Y = \underline{\Spec}_X(\mathcal{A}) \ar[rr] \ar[rd] & &
X' = \underline{\Spec}_X(\mathcal{A}') \ar[ld] \\
& X &
}
$$
voir Lemme \ref{morphisms-lemma-affine-equivalence-algebras}. Le schéma $X' \to X$ sera
la normalisation de $X$ dans le schéma $Y$. Voici un
contexte légèrement plus général. Supposons que nous
ayons un quasi-compact et un quasi-morphisme séparé
$f : Y \to X$ de schémas. Dans ce
cas, le faisceau de $\mathcal{O}_X$-algèbres $f_*\mathcal{O}_Y$ est quasi-cohérent,
voir Schémas, Lemme \ref{schemes-lemma-push-forward-quasi-coherent}. En prenant la clôture
intégrale $\mathcal{O}' \subset f_*\mathcal{O}_Y$ nous obtenons un faisceau quasi-cohérent
de $\mathcal{O}_X$-algèbres dont le spectre relatif est
la normalisation de $X$ dans $Y$. Voici la définition
formelle.

\begin{definition}
\label{morphisms-definition-normalization-X-in-Y}
Soit $f : Y \to X$ un morphisme de schémas quasi-compact et quasi-séparé. Soit $\mathcal{O}'$ la clôture
intégrale de $\mathcal{O}_X$ dans $f_*\mathcal{O}_Y$. La normalisation {\it
de $X$ dans $Y$} est le schéma\footnote{
Le schéma $X'$ n'a pas besoin d'être normal, par exemple si $Y = X$
et $f = \text{id}_X$, alors $X' = X$.}
$$
\nu : X' = \underline{\Spec}_X(\mathcal{O}') \to X
$$
au-dessus de $X$. Il est muni d'une factorisation naturelle
$$
Y \xrightarrow{f'} X' \xrightarrow{\nu} X
$$
du morphisme initial $f$.
\end{definition}

\noindent
La factorisation est la composition du morphisme
canonique $Y \to \underline{\Spec}_X(f_*\mathcal{O}_Y)$ (voir Constructions,
Lemme \ref{constructions-lemma-canonical-morphism})
et du morphisme des spectres
relatifs provenant de l'application d'inclusion $\mathcal{O}' \to f_*\mathcal{O}_Y$.
Nous pouvons caractériser la normalisation
comme suit.

\begin{lemma}
\label{morphisms-lemma-characterize-normalization}
Soit $f : Y \to X$ un morphisme de schémas quasi-compact et quasi-séparé. La factorisation
$f = \nu \circ f'$, où $\nu : X' \to X$ est la normalisation de $X$
dans $Y$, est caractérisée par les deux propriétés
suivantes :
\begin{enumerate}
\item le morphisme $\nu$ est entier, et
\item pour toute factorisation $f = \pi \circ g$, avec $\pi : Z \to X$ entier,
il existe un diagramme commutatif
$$
\xymatrix{
Y \ar[d]_{f'} \ar[r]_g & Z \ar[d]^\pi \\
X' \ar[ru]^h \ar[r]^\nu & X
}
$$
pour un morphisme unique $h : X' \to Z$.
\end{enumerate}
De plus, le morphisme $f' : Y \to X'$ est dominant et dans (2) le morphisme
$h : X' \to Z$ est la normalisation de $Z$ dans $Y$.
\end{lemma}

\begin{proof}
Soit $\mathcal{O}' \subset f_*\mathcal{O}_Y$ la clôture intégrale de $\mathcal{O}_X$ comme dans
la définition \ref{morphisms-definition-normalization-X-in-Y}. Le morphisme $\nu$ est entier
par construction, ce qui prouve (1). Supposons donnée
une factorisation $f = \pi \circ g$ avec $\pi : Z \to X$ entier comme
dans (2). Par la Définition \ref{morphisms-definition-integral} $\pi$
est affine, et donc $Z$ est le
spectre relatif d'un faisceau quasi-cohérent d'algèbres $\mathcal{O}_X$
$\mathcal{B}$. Le morphisme $g : Y \to Z$ correspond à un homomorphisme
d'$\mathcal{O}_X$-algèbres $\chi : \mathcal{B} \to f_*\mathcal{O}_Y$. Comme l'algèbre $\mathcal{B}(U)$ est entière
sur $\mathcal{O}_X(U)$ pour chaque ouvert affine $U \subset X$
(par la Définition \ref{morphisms-definition-integral}), nous
voyons à partir du Lemme \ref{morphisms-lemma-integral-closure}
cela $\chi(\mathcal{B}) \subset \mathcal{O}'$. Par la fonctorialité
du spectre relatif Lemme \ref{morphisms-lemma-affine-equivalence-algebras},
cela nous fournit un morphisme
unique $h : X' \to Z$. Nous omettons
la vérification que le diagramme commute.

\medskip\noindent
Il est clair que (1) et (2) caractérisent la factorisation $f = \nu \circ f'$ puisqu'elle
la caractérise comme un objet initial dans une catégorie.

\medskip\noindent D'après
la propriété universelle en (2), nous voyons que $f'$ ne se factorise pas par
un sous-schéma fermé propre de $X'$. Ainsi, l'image schématique de $f'$
est $X'$. Puisque $f'$ est quasi-compact (d'après
Schémas, Lemme \ref{schemes-lemma-quasi-compact-permanence} et le fait que $\nu$ est séparé en tant que morphisme
affine), nous voyons que $f'(Y)$ est dense dans $X'$. Par conséquent,
$f'$ est dominant.

\medskip\noindent
Observez que $g$ est quasi-compact et quasi-séparé
d'après Schémas, Lemme \ref{schemes-lemma-compose-after-separated} et \ref{schemes-lemma-quasi-compact-permanence}.
Ainsi, la dernière affirmation du lemme a du sens. Le morphisme
$h$ dans
(2) est entier d'après le Lemme \ref{morphisms-lemma-finite-permanence}. Étant donné une factorisation
$g = \pi' \circ g'$ avec $\pi' : Z' \to Z$ entier, nous obtenons
une factorisation $f = (\pi \circ \pi') \circ g'$ et nous obtenons un morphisme
$h' : X' \to Z'$. L'unicité implique que $\pi' \circ h' = h$. Par
conséquent, la caractérisation (1), (2) s'applique au
morphisme $h : X' \to Z$ ce qui donne l'affirmation finale du lemme.
\end{proof}

\begin{lemma}
\label{morphisms-lemma-functoriality-normalization}
Soit
$$
\xymatrix{
Y_2 \ar[d]_{f_2} \ar[r] & Y_1 \ar[d]^{f_1} \\
X_2 \ar[r] & X_1
}
$$
soit un diagramme commutatif de morphismes de schémas.
Supposons $f_1$, $f_2$ quasi-compacts et
quasi-séparés. Soient $f_i = \nu_i \circ f_i'$, $i = 1, 2$
les factorisations canoniques, où $\nu_i : X_i' \to X_i$ est la normalisation
de $X_i$ dans $Y_i$. Alors il existe un unique flèche
$X'_2 \to X'_1$ s'insérant dans un diagramme
commutatif
$$
\xymatrix{
Y_2 \ar[d]_{f_2'} \ar[r] & Y_1 \ar[d]^{f_1'} \\
X_2' \ar[d]_{\nu_2} \ar[r] & X_1' \ar[d]^{\nu_1} \\
X_2 \ar[r] & X_1
}
$$
\end{lemma}

\begin{proof}
D'après les Lemme \ref{morphisms-lemma-characterize-normalization} (1) et \ref{morphisms-lemma-base-change-finite},
le changement de base
$X_2 \times_{X_1} X'_1 \to X_2$ est entier. Notons que $f_2$
se factorise par ce morphisme. Ainsi on
obtient un morphisme unique
$X'_2 \to X_2 \times_{X_1} X'_1$ d'après le Lemme
\ref{morphisms-lemma-characterize-normalization} (2). Cela donne la flèche
$X'_2 \to X'_1$ s'insérant dans le diagramme
commutatif et l'unicité s'ensuit également.
\end{proof}

\begin{lemma}
\label{morphisms-lemma-normalization-localization}
Soit $f : Y \to X$ un morphisme quasi-compact et quasi-séparé de schémas. Soit $U \subset X$ un
sous-schéma ouvert et posons $V = f^{-1}(U)$. Alors la normalisation de $U$
dans $V$ est l'image réciproque de $U$ dans la normalisation
de $X$ dans $Y$.
\end{lemma}

\begin{proof}
Clair d'après la construction.
\end{proof}

\begin{lemma}
\label{morphisms-lemma-normalization-is-normalization}
Soit $f : Y \to X$ un morphisme quasi-compact et quasi-séparé de schémas. Soit $X'$ la normalisation
de $X$ dans $Y$. Alors la normalisation de $X'$ dans $Y$
est $X'$.
\end{lemma}

\begin{proof}
Si $Y \to X'' \to X'$ est la normalisation de $X'$ dans $Y$, alors
nous pouvons appliquer le Lemme \ref{morphisms-lemma-characterize-normalization} à la composition
$X'' \to X$ pour obtenir un morphisme canonique $h : X' \to X''$ sur
$X$. Nous omettons la vérification que les morphismes $h$
et $X'' \to X'$ sont mutuellement inverses (en utilisant l'unicité
de la factorisation dans le lemme).
\end{proof}

\begin{lemma}
\label{morphisms-lemma-normalization-in-reduced}
Soit $f : Y \to X$ un morphisme de schémas quasi-compact et quasi-séparé. Soit $X' \to X$ la normalisation
de $X$ dans $Y$. Si $Y$ est réduit, alors $X'$ l'est
également.
\end{lemma}

\begin{proof}
Cela résulte du fait qu'un sous-anneau d'un anneau réduit est réduit. Quelques
détails sont omis.
\end{proof}

\begin{lemma}
\label{morphisms-lemma-normalization-generic}
Soit $f : Y \to X$ un morphisme de schémas quasi-compact et quasi-séparé. Soit $X' \to X$
la normalisation de $X$ dans $Y$. Chaque point générique d'une composante
irréductible de $X'$ est l'image d'un point générique d'une composante
irréductible de $Y$.
\end{lemma}

\begin{proof}
D'après le Lemme \ref{morphisms-lemma-normalization-localization}, on peut supposer que $X = \Spec(A)$ est
affine. Choisissons un recouvrement fini par des ouverts affines $Y = \bigcup \Spec(B_i)$.
Ensuite $X' = \Spec(A')$ et les morphismes $\Spec(B_i) \to Y \to X'$ définissent
conjointement une application $A$-algèbre
injective $A' \to \prod B_i$.
Ainsi, le lemme découle de l'Algèbre, Lemme \ref{algebra-lemma-injective-minimal-primes-in-image}.
\end{proof}

\begin{lemma}
\label{morphisms-lemma-normalization-in-disjoint-union}
Soit $f : Y \to X$ un morphisme de schémas quasi-compact et quasi-séparé. Supposons que
$Y = Y_1 \amalg Y_2$ soit une réunion disjointe de deux schémas. Écrivons $f_i = f|_{Y_i}$.
Soit $X_i'$ la normalisation de $X$ dans $Y_i$. Alors $X_1' \amalg X_2'$
est la normalisation de $X$ dans $Y$.
\end{lemma}

\begin{proof}
En termes de clôtures intégrales, cela correspond au fait suivant : Soit $A \to B$ un
homomorphisme d'anneaux. Supposons que $B = B_1 \times B_2$. Soit $A_i'$ la clôture
intégrale de $A$ dans $B_i$. Alors $A_1' \times A_2'$ est la clôture
intégrale de $A$ dans $B$. La raison pour laquelle
cela fonctionne est que les éléments $(1, 0)$ et $(0, 1)$ de $B$ sont idempotents
et donc entiers sur $A$. Ainsi, la clôture intégrale $A'$ de $A$
dans $B$ est un produit et il n'est pas difficile de voir que les se factorise
par sont les clôtures intégrales $A'_i$ comme décrit ci-dessus (certains détails
omis).
\end{proof}

\begin{lemma}
\label{morphisms-lemma-normalization-in-universally-closed}
Soit $f : X \to S$ un morphisme de schémas quasi-compact, quasi-séparé
et universellement fermé. Alors
$f_*\mathcal{O}_X$ est entier sur $\mathcal{O}_S$. En d'autres termes, la normalisation
de $S$ dans $X$ est égale à la factorisation
$$
X \longrightarrow \underline{\Spec}_S(f_*\mathcal{O}_X)
\longrightarrow S
$$
de Constructions, Lemme \ref{constructions-lemma-canonical-morphism}.
\end{lemma}

\begin{proof}
La question est locale sur $S$, donc nous pouvons supposer que $S = \Spec(R)$
est affine. Soit $h \in \Gamma(X, \mathcal{O}_X)$. Nous devons montrer que $h$ satisfait
une équation unitaire sur $R$. Considérons $h$ comme un morphisme
comme dans le diagramme commutatif suivant
$$
\xymatrix{
X \ar[rr]_h \ar[rd]_f & & \mathbf{A}^1_S \ar[ld] \\
& S &
}
$$
Soit $Z \subset \mathbf{A}^1_S$ l'image schématique de $h$, voir
Définition \ref{morphisms-definition-scheme-theoretic-image}. Le morphisme $h$
est quasi-compact car $f$ est quasi-compact
et $\mathbf{A}^1_S \to S$ est séparé, voir
Schémas, Lemme \ref{schemes-lemma-quasi-compact-permanence}. D'après le Lemme \ref{morphisms-lemma-quasi-compact-scheme-theoretic-image},
le morphisme $X \to Z$ est dominant. D'après
le Lemme \ref{morphisms-lemma-image-proper-scheme-closed}, le
morphisme $X \to Z$ est fermé. Donc $h(X) = Z$
(au sens des ensembles). Ainsi nous pouvons utiliser
le Lemme
\ref{morphisms-lemma-image-proper-is-proper} pour conclure que $Z \to S$
est universellement fermé (et même propre).
Puisque $Z \subset \mathbf{A}^1_S$, nous voyons que $Z \to S$ est affine
et propre, donc entier d'après le Lemme \ref{morphisms-lemma-integral-universally-closed}. En écrivant $\mathbf{A}^1_S = \Spec(R[T])$
Nous concluons que l'idéal $I \subset R[T]$ de
$Z$ contient un polynôme unitaire $P(T) \in R[T]$.
Donc $P(h) = 0$ et nous avons terminé.
\end{proof}

\begin{lemma}
\label{morphisms-lemma-normalization-in-integral}
Soit $f : Y \to X$ un morphisme entier. Alors la normalisation
de $X$ dans $Y$ est égale à $Y$.
\end{lemma}

\begin{proof}
D'après le Lemme \ref{morphisms-lemma-integral-universally-closed} ceci est un cas particulier
du Lemme \ref{morphisms-lemma-normalization-in-universally-closed}.
\end{proof}

\begin{lemma}
\label{morphisms-lemma-normal-normalization}
Soit $f : Y \to X$ un morphisme de schémas quasi-compact et quasi-séparé. Soit
$X'$ la normalisation de $X$ dans $Y$. Supposons
\begin{enumerate}
\item $Y$ est un schéma normal,
\item les ouverts quasi-compacts de $Y$ ont un nombre fini de composantes irréductibles.
\end{enumerate}
Alors $X'$ est une réunion disjointe de schémas normaux intègres. De plus,
le morphisme $Y \to X'$ est dominant et induit une bijection entre les
composantes irréductibles.
\end{lemma}

\begin{proof}
Soit $U \subset X$ un ouvert affine. Considérons
l'image réciproque $U'$ de
$U$ dans $X'$. Posons $V = f^{-1}(U)$. D'après le Lemme
\ref{morphisms-lemma-normalization-localization}, le diagramme $V \to U' \to U$ est la normalisation de
$U$ dans $V$. Disons $U = \Spec(A)$. Alors $V$ est
quasi-compact, et possède donc un nombre
fini de composantes irréductibles par hypothèse. Ainsi $V = \coprod_{i = 1, \ldots n} V_i$
est une union
finie disjointe de schémas intègres normaux d'après les Propriétés, Lemme
\ref{properties-lemma-normal-locally-finite-nr-irreducibles}. D'après le Lemme \ref{morphisms-lemma-normalization-in-disjoint-union}, nous
voyons que $U' = \coprod_{i = 1, \ldots, n} U_i'$, où $U'_i$ est la
normalisation de $U$ dans $V_i$.
Par les propriétés, Lemme \ref{properties-lemma-normal-integral-sections} nous voyons que
$B_i = \Gamma(V_i, \mathcal{O}_{V_i})$ est un domaine normal. Notez que $U_i' = \Spec(A_i')$,
où $A_i' \subset B_i$ est la clôture intégrale de
$A$ dans $B_i$, voir
Lemme \ref{morphisms-lemma-integral-closure}. Par l'algèbre,
Lemme \ref{algebra-lemma-integral-closure-in-normal} nous voyons que $A_i' \subset B_i$
est un domaine normal. Ainsi $U' = \coprod U_i'$
est une union finie de schémas intègres normaux et est
donc normal.

\medskip\noindent
Comme $X'$ a un recouvrement ouvert par les schémas $U'$,
nous concluons d'après Propriétés, Lemme \ref{properties-lemma-locally-normal} que $X'$ est normal. D'autre
part, chaque $U'$ est une réunion finie disjointe de schémas
irréductibles, donc chaque ouvert quasi-compact de $X'$ a un nombre fini de
composantes irréductibles (par un argument topologique que nous omettons).
Ainsi $X'$ est une réunion disjointe
de schémas intègres normaux d'après Propriétés, Lemme \ref{properties-lemma-normal-locally-finite-nr-irreducibles}. Il est clair d'après
la description de $X'$ ci-dessus que $Y \to X'$ est dominant
et induit une bijection sur les composantes irréductibles $V \to U'$
pour chaque ouvert affine $U \subset X$. La bijection des composantes
irréductibles pour le morphisme $Y \to X'$
résulte de cela par un argument topologique (omis).
\end{proof}

\begin{lemma}
\label{morphisms-lemma-nagata-normalization-finite-general}
Soit $f : X \to S$ un morphisme. Supposons que
\begin{enumerate}
\item $S$ soit un schéma de Nagata,
\item $f$ soit quasi-compact et quasi-séparé,
\item les ouverts quasi-compacts de $X$ ont un nombre fini de composantes irréductibles,
\item si $x \in X$ est un point générique d'une composante irréductible,
alors l'extension de corps $\kappa(x)/\kappa(f(x))$ est de type
fini, et
\item $X$ est réduit.
\end{enumerate}
Alors la normalisation $\nu : S' \to S$ de $S$ dans $X$ est finie.
\end{lemma}

\begin{proof}
Il y a une réduction immédiate au cas $S = \Spec(R)$ où $R$
est un anneau de Nagata par l'hypothèse (1). Nous devons montrer
que la clôture intégrale $A$ de $R$ dans $\Gamma(X, \mathcal{O}_X)$
est finie sur $R$. Étant donné que $f$ est quasi-compact
par l'hypothèse (2), nous pouvons écrire $X = \bigcup_{i = 1, \ldots, n} U_i$ avec
chaque $U_i$ affine. Disons $U_i = \Spec(B_i)$. Chaque $B_i$
est réduit par l'hypothèse (5) et
possède un nombre fini de premiers minimaux
$\mathfrak q_{i1}, \ldots, \mathfrak q_{im_i}$ par l'hypothèse
(3) et l'Algèbre, Lemme \ref{algebra-lemma-irreducible}. Nous
avons
$$
\Gamma(X, \mathcal{O}_X) \subset B_1 \times \ldots \times B_n
\subset
\prod\nolimits_{i = 1, \ldots, n}
\prod\nolimits_{j = 1, \ldots, m_i} (B_i)_{\mathfrak q_{ij}}
$$
la deuxième inclusion
par l'algèbre, Lemme \ref{algebra-lemma-reduced-ring-sub-product-fields}. Nous avons $\kappa(\mathfrak q_{ij}) = (B_i)_{\mathfrak q_{ij}}$
par l'algèbre, Lemme \ref{algebra-lemma-minimal-prime-reduced-ring}. Par conséquent,
la clôture intégrale $A$ de $R$ dans
$\Gamma(X, \mathcal{O}_X)$ est contenue dans le produit des clôtures
intégrales $A_{ij}$ de $R$ dans $\kappa(\mathfrak q_{ij})$.
Comme $R$ est noethérien, il suffit
de montrer que $A_{ij}$ est un $R$-module fini
pour chaque $i, j$. Soit $\mathfrak p_{ij} \subset R$ l'image de $\mathfrak q_{ij}$.
Comme $\kappa(\mathfrak q_{ij})/\kappa(\mathfrak p_{ij})$ est une extension
de corps de type fini par l'hypothèse (4), nous
voyons que $R \to \kappa(\mathfrak q_{ij})$ est essentiellement de type fini.
Ainsi $R \to A_{ij}$ est fini
par l'algèbre, Lemme \ref{algebra-lemma-nagata-in-reduced-finite-type-finite-integral-closure}.
\end{proof}

\begin{lemma}
\label{morphisms-lemma-nagata-normalization-finite}
Soit $f : X \to S$ un morphisme. Supposons que
\begin{enumerate}
\item $S$ est un schéma de Nagata,
\item $f$ est de type fini,
\item $X$ est réduit.
\end{enumerate}
Alors la normalisation $\nu : S' \to S$ de $S$ dans $X$ est finie.
\end{lemma}

\begin{proof}
Ceci est un cas particulier
du Lemme \ref{morphisms-lemma-nagata-normalization-finite-general}. À savoir, (2) est vérifié
car le morphisme de type fini $f$ est quasi-compact par
définition et quasi-séparé par
le Lemme \ref{morphisms-lemma-finite-type-Noetherian-quasi-separated}. La condition (3) est vérifiée parce
que $X$ est localement noethérien par le
Lemme \ref{morphisms-lemma-finite-type-noetherian}. Enfin, la condition (4) est vérifiée car
un morphisme de type fini induit des extensions de corps résiduels
de type fini.
\end{proof}

\begin{lemma}
\label{morphisms-lemma-relative-normalization-normal-codim-1}
Soit $f : Y \to X$ un morphisme de schémas de type fini avec $Y$ réduit et $X$
Nagata. Soit $X'$ la normalisation de $X$ dans $Y$.
Soit $x' \in X'$ un point tel que
\begin{enumerate}
\item $\dim(\mathcal{O}_{X', x'}) = 1$, et
\item la fibre de $Y \to X'$ au-dessus de $x'$ est vide.
\end{enumerate}
Alors $\mathcal{O}_{X', x'}$ est un anneau de valuation discret.
\end{lemma}

\begin{proof}
Nous pouvons remplacer $X$ par un voisinage affine de l'image de
$x'$. Ainsi, nous pouvons supposer $X = \Spec(A)$
avec $A$ Nagata. Par le Lemme \ref{morphisms-lemma-nagata-normalization-finite} le
morphisme $X' \to X$ est fini. Ainsi, nous pouvons écrire
$X' = \Spec(A')$ pour une $A$-algèbre finie $A'$.
Par le Lemme \ref{morphisms-lemma-normalization-is-normalization} après avoir remplacé
$X$ par $X'$
nous revenons au cas décrit dans le paragraphe suivant.

\medskip\noindent
Le cas $X = X' = \Spec(A)$ avec $A$ noethérien. Soit $\mathfrak p \subset A$
l'idéal premier correspondant à notre point
$x'$. Choisissez $g \in \mathfrak p$ qui n'est contenu dans
aucun idéal premier minimal de $A$ (utilisez l'évitement des
idéaux premiers et le fait que $A$ possède un nombre
fini d'idéaux premiers minimaux, voir Algèbre, Lemmes
\ref{algebra-lemma-silly} et \ref{algebra-lemma-Noetherian-irreducible-components}). Posons $Z = f^{-1}V(g) \subset Y$ ; c'est
un sous-schéma fermé de $Y$. Alors $f(Z)$ ne contient aucun
point générique par choix de $g$ et ne contient pas $x'$
car $x'$ n'est pas dans l'image de $f$. L'adhérence
de $f(Z)$ est l'ensemble des spécialisations des points
de $f(Z)$ par le lemme \ref{morphisms-lemma-reach-points-scheme-theoretic-image}. Ainsi, l'adhérence de $f(Z)$
ne contient pas $x'$ car la condition $\dim(\mathcal{O}_{X', x'}) = 1$ implique
seulement que les points génériques de $X = X'$ se spécialisent
en $x'$. En d'autres termes, après avoir remplacé
$X$ par un voisinage ouvert affine de $x'$,
nous pouvons supposer que $f^{-1}V(g) = \emptyset$. Ainsi $g$ s'envoie
sur une fonction globale inversible sur $Y$ et nous obtenons
une factorisation
$$
A \to A_g \to \Gamma(Y, \mathcal{O}_Y)
$$
Puisque $X = X'$, cela implique que $A$ est égal
à la clôture intégrale de
$A$ dans $A_g$. Par l'algèbre, Lemme \ref{algebra-lemma-integral-closure-localize}
Nous concluons que $A_\mathfrak p$ est la clôture intégrale
de $A_\mathfrak p$ dans $A_\mathfrak p[1/g]$. Par
notre choix de $g$, puisque $\dim(A_\mathfrak p) = 1$ et
puisque $A$ est réduit, nous voyons que $A_\mathfrak p[1/g]$
est un produit fini de corps (le produit des corps résiduels
des idéaux premiers minimaux contenus dans $\mathfrak p$). Par conséquent
$A_\mathfrak p$ est normal (Algèbre, Lemme \ref{algebra-lemma-characterize-reduced-ring-normal}) et la démonstration
est terminée. Certains détails sont omis.
\end{proof}






\section{Normalisation}
\label{morphisms-section-normalization}

\noindent
Ensuite, nous arrivons à la normalisation d'un schéma $X$.
Nous ne le définissons/construisons que lorsque $X$ a localement un nombre fini de
composantes irréductibles. Soit $X$ un schéma tel que chaque ouvert quasi-compact
possède un nombre fini de composantes irréductibles.
Soit $X^{(0)} \subset X$ l'ensemble des points génériques des composantes irréductibles de $X$.
Soit
\begin{equation}
\label{morphisms-equation-generic-points}
f :
Y = \coprod\nolimits_{\eta \in X^{(0)}} \Spec(\kappa(\eta))
\longrightarrow
X
\end{equation}
l'inclusion des points génériques dans $X$ à
l'aide des applications canoniques de Schémas, section \ref{schemes-section-points}.
Notez que ce morphisme est quasi-compact par
hypothèse et quasi-séparé car $Y$
est séparé (voir Schémas, section \ref{schemes-section-separation-axioms}).

\begin{definition}
\label{morphisms-definition-normalization}
Soit $X$ un schéma tel que tout ouvert quasi-compact ait un nombre fini de
composantes irréductibles. Nous définissons la normalisation {\it
} de $X$ comme le morphisme
$$
\nu : X^\nu \longrightarrow X
$$
qui est la normalisation de $X$ dans le morphisme
$f : Y \to X$ (\ref{morphisms-equation-generic-points}) construit ci-dessus.
\end{definition}

\noindent
Tout schéma localement noethérien possède un ensemble localement fini de composantes
irréductibles et la définition s'y applique. Habituellement,
la normalisation est définie uniquement pour les schémas réduits.
Avec la définition ci-dessus, la normalisation de $X$ est la même
que la normalisation de la réduction $X_{red}$ de $X$.

\begin{lemma}
\label{morphisms-lemma-normalization-reduced}
Soit $X$ un schéma tel que tout ouvert quasi-compact ait
un nombre fini de composantes irréductibles. Le morphisme de normalisation
$\nu$ se factorise par la réduction $X_{red}$ et $X^\nu \to X_{red}$
est la normalisation de $X_{red}$.
\end{lemma}

\begin{proof}
Soit $f : Y \to X$ le morphisme (\ref{morphisms-equation-generic-points}). Nous obtenons
une factorisation $Y \to X_{red} \to X$ de $f$ provenant
des Schémas, Lemme \ref{schemes-lemma-map-into-reduction}. D'après le
Lemme \ref{morphisms-lemma-characterize-normalization}, nous obtenons un morphisme canonique $X^\nu \to X_{red}$
et que $X^\nu$
est la normalisation de $X_{red}$ dans $Y$.
Le lemme s'ensuit puisque $Y \to X_{red}$ est identique au
morphisme (\ref{morphisms-equation-generic-points}) construit pour $X_{red}$.
\end{proof}

\noindent
Si $X$ est réduit, alors la normalisation de
$X$ est la même que le spectre relatif de l'adhérence
intégrale de $\mathcal{O}_X$ dans le faisceau
de fonctions méromorphes $\mathcal{K}_X$ (voir Diviseurs, section
\ref{divisors-section-meromorphic-functions}). À savoir, $\mathcal{K}_X = f_*\mathcal{O}_Y$ dans ce cas,
voir Diviseurs, Lemme \ref{divisors-lemma-reduced-finite-irreducible} et sa démonstration.
Nous décrivons cela ici explicitement.

\begin{lemma}
\label{morphisms-lemma-description-normalization}
Soit $X$ un schéma réduit tel que chaque ouvert quasi-compact
possède un nombre fini de composantes irréductibles. Soit $\Spec(A) = U \subset X$
un ouvert affine. Alors
\begin{enumerate}
\item $A$ possède un nombre fini
de premiers minimaux $\mathfrak q_1, \ldots, \mathfrak q_t$,
\item l'anneau total des fractions
$Q(A)$ de $A$ est $Q(A/\mathfrak q_1) \times \ldots \times Q(A/\mathfrak q_t)$,
\item la clôture intégrale $A'$ de $A$ dans $Q(A)$
est le produit des clôtures intégrales des domaines $A/\mathfrak q_i$
dans le corps $Q(A/\mathfrak q_i)$, et
\item $\nu^{-1}(U)$ est identifié au spectre de $A'$ où $\nu : X^\nu \to X$
est le morphisme de normalisation.
\end{enumerate}
\end{lemma}

\begin{proof}
Les idéaux premiers minimaux correspondent
aux composantes irréductibles
(Algèbre, Lemme
\ref{algebra-lemma-irreducible}), donc nous avons (1) par hypothèse.
Ensuite $(0) = \mathfrak q_1 \cap \ldots \cap \mathfrak q_t$ parce que
$A$
est réduit (Algèbre, Lemme \ref{algebra-lemma-Zariski-topology}).
Ensuite nous avons $Q(A) = \prod A_{\mathfrak q_i} = \prod \kappa(\mathfrak q_i)$ par Algèbre,
Lemmes \ref{algebra-lemma-total-ring-fractions-no-embedded-points} et \ref{algebra-lemma-minimal-prime-reduced-ring}.
Cela prouve (2).
La partie (3) découle de Algèbre, Lemme \ref{algebra-lemma-characterize-reduced-ring-normal},
ou Lemme \ref{morphisms-lemma-normalization-in-disjoint-union}.
La partie (4) est vraie parce qu'il est clair que $f^{-1}(U) \to U$ est le morphisme
$$
\Spec\left(\prod \kappa(\mathfrak q_i)\right)
\longrightarrow
\Spec(A)
$$
où $f : Y \to X$ est le morphisme (\ref{morphisms-equation-generic-points}).
\end{proof}

\begin{lemma}
\label{morphisms-lemma-stalk-normalization}
Soit $X$ un schéma tel que tout ouvert quasi-compact ait un nombre
fini de composantes irréductibles. Soit $\nu : X^\nu \to X$ la normalisation
de $X$. Soit $x \in X$. Alors les suivants sont
canoniquement isomorphes en tant que $\mathcal{O}_{X, x}$-algèbres
\begin{enumerate}
\item le germe $(\nu_*\mathcal{O}_{X^\nu})_x$,
\item la clôture intégrale de $\mathcal{O}_{X, x}$ dans
l'anneau total des fractions de $(\mathcal{O}_{X, x})_{red}$,
\item la clôture intégrale de $\mathcal{O}_{X, x}$ dans le produit
des corps résiduels des idéaux premiers minimaux de $\mathcal{O}_{X, x}$
(et il y en a un nombre fini).
\end{enumerate}
\end{lemma}

\begin{proof}
Après avoir remplacé $X$ par un ouvert affine
voisinage de $x$, nous pouvons supposer que $X$
a un nombre fini de composantes irréductibles
et que $x$ est contenu dans chacune d'elles.
Alors le germe $(\nu_*\mathcal{O}_{X^\nu})_x$ est l'adhérence entière de $A = \mathcal{O}_{X, x}$
dans le produit $L$ des corps résiduels
des idéaux premiers minimaux de $A$. Cela résulte
de la construction de la normalisation
et du Lemme \ref{morphisms-lemma-integral-closure}.
Alternativement, vous pouvez utiliser
le Lemme \ref{morphisms-lemma-description-normalization} et le fait que la normalisation
commute avec la localisation (Algèbre, Lemme \ref{algebra-lemma-integral-closure-localize}).
Puisque $A_{red}$ a un nombre fini de premiers minimaux
(parce qu'ils correspondent exactement aux points génériques
des composantes irréductibles de $X$ passant
par $x$), nous voyons que $L$ est l'anneau total
des fractions de $A_{red}$ (Algèbre, Lemme \ref{algebra-lemma-total-ring-fractions-no-embedded-points}). Ainsi, notre anneau est
également la clôture intégrale de $A$ dans l'anneau total
des fractions de $A_{red}$.
\end{proof}

\begin{lemma}
\label{morphisms-lemma-normalization-normal}
Soit $X$ un schéma tel que tout ouvert quasi-compact ait un
nombre fini de composantes irréductibles.
\begin{enumerate}
\item La normalisation $X^\nu$ est une union disjointe de schémas intègres normaux.
\item Le morphisme $\nu : X^\nu \to X$ est entier, surjectif, et induit une
bijection sur les composantes irréductibles.
\item Pour tout morphisme entier $\alpha : X' \to X$ tel que pour $U \subset X$
ouvert quasi-compact l'image réciproque $\alpha^{-1}(U)$ ait un nombre fini
de composantes irréductibles et $\alpha|_{\alpha^{-1}(U)} : \alpha^{-1}(U) \to U$
soit birationnel\footnote{Cette formulation maladroite est
nécessaire car nous n'avons défini ce que signifie qu'un
morphisme soit birationnel que si la source et la cible ont
un nombre fini de composantes irréductibles. Il suffit
que $X'_{red} \to X_{red}$ satisfasse la condition.}
il existe une factorisation
$X^\nu \to X' \to X$ et $X^\nu \to X'$ est la normalisation de $X'$.
\item Pour tout morphisme $Z \to X$ avec $Z$ un schéma normal
tel que chaque composante irréductible de $Z$ domine une composante
irréductible de $X$ il existe une unique factorisation $Z \to X^\nu \to X$.
\end{enumerate}
\end{lemma}

\begin{proof}
Soit $f : Y \to X$ comme dans (\ref{morphisms-equation-generic-points}). Le schéma $X^\nu$
est une réunion disjointe de schémas intègres normaux
car $Y$ est normal et tout ouvert affine de
$Y$ possède un nombre
fini de composantes irréductibles, voir le Lemme
\ref{morphisms-lemma-normal-normalization}. Cela prouve (1). Alternativement,
on peut déduire (1) des
Lemmes \ref{morphisms-lemma-normalization-reduced} et \ref{morphisms-lemma-description-normalization}.

\medskip\noindent
Le morphisme $\nu$ est entier d'après le Lemme \ref{morphisms-lemma-characterize-normalization}. D'après le Lemme
\ref{morphisms-lemma-normal-normalization}, le morphisme $Y \to X^\nu$ induit
une bijection sur les composantes irréductibles, et par construction
de $Y$ cela implique que $X^\nu \to X$ induit une bijection
sur les composantes irréductibles. Par construction $f : Y \to X$
est dominant, donc également $\nu$ est dominant. Puisqu'un morphisme
entier est fermé (Lemme \ref{morphisms-lemma-integral-universally-closed}), cela implique que $\nu$
est surjectif. Cela prouve (2).

\medskip\noindent
Supposons que $\alpha : X' \to X$ soit comme dans (3). Il est clair
que $X'$ satisfait les hypothèses sous lesquelles
la normalisation est définie. Soit $f' : Y' \to X'$
le morphisme (\ref{morphisms-equation-generic-points}) construit à partir de $X'$.
Comme $\alpha$ est localement birationnel, il
est clair que $Y' = Y$ et $f = \alpha \circ f'$.
Par conséquent, la factorisation $X^\nu \to X' \to X$
existe et $X^\nu \to X'$ est la normalisation
de $X'$ d'après le Lemme \ref{morphisms-lemma-characterize-normalization}. Cela prouve (3).

\medskip\noindent
Soit $g : Z \to X$ un morphisme dont le domaine est un schéma normal
et tel que chaque composante irréductible domine une composante irréductible
de $X$. D'après le Lemme \ref{morphisms-lemma-normalization-reduced} nous avons
$X^\nu = X_{red}^\nu$ et d'après Schémas, Lemme
\ref{schemes-lemma-map-into-reduction} $Z \to X$ se factorise par $X_{red}$.
Par conséquent, nous pouvons remplacer $X$ par $X_{red}$
et supposer que $X$ est réduit. De plus, comme la factorisation
est unique, il suffit de la construire localement
sur $Z$. Soient $W \subset Z$ et $U \subset X$ des
ouverts affines tels que $g(W) \subset U$. Écrivons $U = \Spec(A)$
et $W = \Spec(B)$, avec $g|_W$ donné par $\varphi : A \to B$.
Nous utiliserons librement les résultats du Lemme \ref{morphisms-lemma-description-normalization}. Soient $\mathfrak p_1, \ldots, \mathfrak p_t$ les
premiers minimaux de $A$. Comme $Z$ est normal, nous voyons
que $B$ est un anneau normal, en particulier
réduit. De plus, d'après l'hypothèse, pour tout premier
minimal $\mathfrak q \subset B$, nous avons que $\varphi^{-1}(\mathfrak q)$ est un premier minimal de
$A$. Ainsi, si $x \in A$ est un non-diviseur de zéro, c'est-à-dire $x \not \in \bigcup \mathfrak p_i$,
alors $\varphi(x)$ est un non-diviseur de zéro dans $B$. Ainsi, nous
obtenons un homomorphisme canonique d'anneaux $Q(A) \to Q(B)$. Comme $B$
est normal, il est égal à sa clôture intégrale dans
$Q(B)$ (voir Algèbre, Lemme \ref{algebra-lemma-normal-ring-integrally-closed}). Ainsi, nous voyons
que la clôture intégrale $A' \subset Q(A)$ de $A$ s'envoie dans
$B$ par l'homomorphisme canonique $Q(A) \to Q(B)$.
Puisque $\nu^{-1}(U) = \Spec(A')$, cela donne la factorisation canonique
$W \to \nu^{-1}(U) \to U$ de $\nu|_W$. Nous omettons
la vérification qu'elle est unique.
\end{proof}

\begin{lemma}
\label{morphisms-lemma-normalization-in-terms-of-components}
Soit $X$ un schéma tel que chaque ouvert quasi-compact
ait un nombre fini de composantes irréductibles. Soient $Z_i \subset X$,
$i \in I$ les composantes irréductibles de $X$
munies de la structure réduite induite. Soit $Z_i^\nu \to Z_i$ la normalisation.
Alors $\coprod_{i \in I} Z_i^\nu \to X$ est la normalisation de $X$.
\end{lemma}

\begin{proof}
Nous pouvons supposer que $X$ est réduit, voir Lemme \ref{morphisms-lemma-normalization-reduced}.
Ensuite, le lemme découle soit de la description locale
dans le lemme \ref{morphisms-lemma-description-normalization} soit du Lemme
\ref{morphisms-lemma-normalization-normal} partie (3) car $\coprod Z_i \to X$ est entier et
localement birationnel (puisque $X$ est réduit et a localement
un nombre fini de composantes irréductibles).
\end{proof}

\begin{lemma}
\label{morphisms-lemma-normalization-birational}
Soit $X$ un schéma réduit avec un nombre fini de composantes irréductibles.
Alors le morphisme de normalisation $X^\nu \to X$ est birationnel.
\end{lemma}

\begin{proof}
La normalisation induit une bijection des composantes irréductibles
par le Lemme \ref{morphisms-lemma-normalization-normal}. Soit $\eta \in X$ un point générique
d'une composante irréductible de $X$ et soit
$\eta^\nu \in X^\nu$ le point générique de la composante irréductible correspondante
de $X^\nu$. Alors $\eta^\nu \mapsto \eta$ et pour terminer la démonstration,
nous devons montrer que $\mathcal{O}_{X, \eta} \to \mathcal{O}_{X^\nu, \eta^\nu}$ est
un isomorphisme, voir la Définition \ref{morphisms-definition-birational}. Comme
$X$ et $X^\nu$ sont réduits, nous voyons que les
deux anneaux locaux sont égaux
à leur corps résiduel (Algèbre, Lemme \ref{algebra-lemma-minimal-prime-reduced-ring}).
D'autre part, par la construction de la normalisation
comme la normalisation de $X$ dans $Y = \coprod \Spec(\kappa(\eta))$
nous voyons que
nous avons $\kappa(\eta) \subset \kappa(\eta^\nu) \subset \kappa(\eta)$ et la démonstration
est complète.
\end{proof}

\begin{lemma}
\label{morphisms-lemma-finite-birational-over-normal}
Un morphisme fini (ou même entier) birationnel $f : X \to Y$ de schémas
intègres avec $Y$ normal est un isomorphisme.
\end{lemma}

\begin{proof}
Soit $V \subset Y$ un ouvert affine avec
image réciproque $U \subset X$ qui est également un ouvert affine.
Comme $f$ est un morphisme birationnel de schémas intègres, l'homomorphisme
$\mathcal{O}_Y(V) \to \mathcal{O}_X(U)$ est un homomorphisme injectif de domaines intègres qui
induit un isomorphisme de corps de fractions. Comme $Y$
est normal, l'anneau $\mathcal{O}_Y(V)$ est intégralement clos dans le corps des
fractions. Comme $f$ est fini (ou entier), chaque élément
de $\mathcal{O}_X(U)$ est entier sur $\mathcal{O}_Y(V)$. Nous concluons
que $\mathcal{O}_Y(V) = \mathcal{O}_X(U)$. Cela prouve que $f$ est un isomorphisme
comme souhaité.
\end{proof}

\begin{lemma}
\label{morphisms-lemma-normalization-trivial}
Soit $X$ un schéma avec localement un nombre fini de composantes irréductibles.
Le morphisme de normalisation $\nu : X^\nu \to X$ est un isomorphisme si et seulement
si $X$ est normal.
\end{lemma}

\begin{proof}
Si $\nu$ est un isomorphisme, alors $X$
est normal d'après le Lemme
\ref{morphisms-lemma-normalization-normal}. Réciproquement,
supposons que $X$ soit normal. D'après
le Lemme \ref{morphisms-lemma-normalization-localization} et Propriétés, Lemme \ref{properties-lemma-normal-locally-finite-nr-irreducibles}, nous pouvons
supposer que $X$
est intègre. D'après le Lemme \ref{morphisms-lemma-normalization-normal}, le morphisme
$\nu$ est intègre et $X^\nu$ a une unique
composante irréductible
(donc c'est un schéma intègre). Nous concluons d'après
les Lemme \ref{morphisms-lemma-normalization-birational} et \ref{morphisms-lemma-finite-birational-over-normal}.
\end{proof}

\begin{lemma}
\label{morphisms-lemma-Japanese-normalization}
Soit $X$ un schéma intègre, japonais.
La normalisation $\nu : X^\nu \to X$ est un morphisme fini.
\end{lemma}

\begin{proof}
Cela découle de la définition
(Propriétés, Définition \ref{properties-definition-nagata}) et du Lemme
\ref{morphisms-lemma-description-normalization}. À savoir, dans ce cas, le lemme dit que
$\nu^{-1}(\Spec(A))$ est le spectre de la clôture intégrale
de $A$ dans son corps de fractions.
\end{proof}

\begin{lemma}
\label{morphisms-lemma-nagata-normalization}
Soit $X$ un Schéma de
Nagata. La normalisation $\nu : X^\nu \to X$ est un morphisme fini.
\end{lemma}

\begin{proof}
Notez qu'un schéma de Nagata est localement noethérien,
ainsi la Définition \ref{morphisms-definition-normalization}
s'applique. Le lemme est maintenant un
cas particulier du Lemme \ref{morphisms-lemma-nagata-normalization-finite-general} mais nous pouvons
également le prouver directement
comme suit. Écrivons $X^\nu \to X$ comme
la composition $X^\nu \to X_{red} \to X$. Comme $X_{red} \to X$ est une immersion
fermée, elle est finie. Par conséquent, il suffit de
prouver le lemme pour un schéma de Nagata réduit
(d'après le Lemme \ref{morphisms-lemma-composition-finite}). Soit $\Spec(A) = U \subset X$
un ouvert affine. D'après le Lemme \ref{morphisms-lemma-description-normalization}
nous avons $\nu^{-1}(U) = \Spec(\prod A_i')$ où $A_i'$ est la clôture
intégrale de $A/\mathfrak q_i$ dans son corps des fractions. Comme $A$
est un anneau de Nagata (voir Propriétés, Lemme \ref{properties-lemma-locally-nagata}) chacune
des extensions d'anneaux
$A/\mathfrak q_i \subset A'_i$ sont finies. Donc $A \to \prod A'_i$ est un homomorphisme d'anneaux
fini et on obtient le résultat.
\end{proof}

\begin{lemma}
\label{morphisms-lemma-normalization-geom-unibranch-univ-homeo}
Soit $X$ un schéma irréductible, géométriquement unibranche.
Le morphisme de normalisation $\nu : X^\nu \to X$ est un homéomorphisme
universel.
\end{lemma}

\begin{proof}
Nous devons montrer que $\nu$ est intègre, universellement
injectif et surjectif, voir Lemme \ref{morphisms-lemma-universal-homeomorphism}. D'après
le Lemme \ref{morphisms-lemma-normalization-normal}, le morphisme $\nu$ est
intègre. Soit $x \in X$ et posons $A = \mathcal{O}_{X, x}$. Puisque
$X$ est irréductible, nous voyons que
$A$ possède un seul idéal premier minimal $\mathfrak p$
et $A_{red} = A/\mathfrak p$. Par le Lemme \ref{morphisms-lemma-stalk-normalization},
le germe $A' = (\nu_*\mathcal{O}_{X^\nu})_x$ est la clôture intégrale de $A$
dans le corps des fractions de
$A_{red}$. D'après Compléments d'algèbre, Définition \ref{more-algebra-definition-unibranch},
nous voyons que $A'$ possède un seul idéal
premier $\mathfrak m'$ au-dessus de $\mathfrak m_x \subset A$ et $\kappa(\mathfrak m')/\kappa(x)$.
est purement inséparable. Ainsi $\nu$ est bijectif (donc
surjectif) et universellement injectif par le Lemme \ref{morphisms-lemma-universally-injective}.
\end{proof}







\section{Normalisation faible}
\label{morphisms-section-weak-normalization}

\noindent
Nous définirons seulement la normalisation faible d'un schéma lorsque celui-ci
a localement un nombre fini de composantes irréductibles ; similaire
au cas de la normalisation.

\begin{lemma}
\label{morphisms-lemma-relative-weak-normalization-algebra}
Soit $A \to B$ un homomorphisme d'anneaux induisant un morphisme
dominant $\Spec(B) \to \Spec(A)$ de spectres. Il existe une $A$-sous-algèbre
$B' \subset B$ telle que
\begin{enumerate}
\item $\Spec(B') \to \Spec(A)$ est un homéomorphisme universel,
\item étant donnée une factorisation $A \to C \to B$ telle que
$\Spec(C) \to \Spec(A)$ est un homéomorphisme universel, l'image de
$C \to B$ est contenue dans $B'$.
\end{enumerate}
\end{lemma}

\begin{proof}
Nous utiliserons Lemme \ref{morphisms-lemma-reduction-universal-homeomorphism} sans autre mention.
Considérons le diagramme commutatif
$$
\xymatrix{
B \ar[r] & B_{red} \\
A \ar[u] \ar[r] & A_{red} \ar[u]
}
$$
Pour toute factorisation $A \to C \to B$ de $A \to B$ comme dans (2),
nous voyons que $A_{red} \to C_{red} \to B_{red}$ est une factorisation de
$A_{red} \to B_{red}$ comme dans (2). Il s'ensuit que si le lemme
est vrai pour $A_{red} \to B_{red}$ et produit la $A_{red}$-sous-algèbre $B'_{red} \subset B_{red}$,
alors en posant $B' \subset B$ égal à l'image
réciproque de $B'_{red}$ résout le lemme pour $A \to B$.
Cela nous ramène au cas discuté dans le paragraphe suivant.

\medskip\noindent
Supposons que $A$ et $B$ sont réduits.
Dans ce cas $A \subset B$ par Algèbre, Lemme \ref{algebra-lemma-image-dense-generic-points}.
Soit $A \to C \to B$ une factorisation comme dans
(2). Alors nous pouvons appliquer la Proposition \ref{morphisms-proposition-universal-homeomorphism} à $A \subset C$
pour voir que chaque élément de $C$ est
contenu dans une extension $A[c_1, \ldots, c_n] \subset C$ telle que pour
$i = 1, \ldots, n$ nous avons
\begin{enumerate}
\item $c_i^2, c_i^3 \in A[c_1, \ldots, c_{i - 1}]$, ou
\item il existe un nombre premier
$p$ avec $pc_i, c_i^p \in A[c_1, \ldots, c_{i - 1}]$.
\end{enumerate}
Ainsi, la propriété (2) est vérifiée si nous définissons $B' \subset B$
comme le sous-ensemble des éléments $b \in B$ qui sont contenus dans
une extension $
A[b_1, \ldots, b_n] \subset B
$ telle
que (*) est vérifiée : pour $i = 1, \ldots, n$ nous avons
\begin{enumerate}
\item $b_i^2, b_i^3 \in A[b_1, \ldots, b_{i - 1}]$, ou
\item il existe un nombre premier
$p$ avec $pb_i, b_i^p \in A[b_1, \ldots, b_{i - 1}]$.
\end{enumerate}
Il n'y a que deux choses à vérifier : (a) $B'$ est une $A$-sous-algèbre,
et (b) $\Spec(B') \to \Spec(A)$ est un homéomorphisme
universel. La partie (a) découle du fait que étant donné
$n \geq 0$ et $b_1, \ldots, b_n \in B$ satisfaisant (*), et $m \geq 0$
et $b'_1, \ldots, b'_m \in B$ satisfaisant
(*), l'entier $n + m$ et $b_1, \ldots, b_n, b'_1, \ldots, b'_m \in B$ satisfait
également (*). Enfin, la partie (b)
tient par la Proposition \ref{morphisms-proposition-universal-homeomorphism} et notre construction de
$B'$.
\end{proof}

\begin{lemma}
\label{morphisms-lemma-relative-weak-normalization-localization}
Soit $A \to B$ un homomorphisme d'anneaux induisant un morphisme
dominant $\Spec(B) \to \Spec(A)$ des spectres. La formation de la $A$-sous-algèbre
$B' \subset B$ dans le Lemme \ref{morphisms-lemma-relative-weak-normalization-algebra} commute avec la
localisation (voir la démonstration pour explication).
\end{lemma}

\begin{proof}
Soit $S \subset A$ un sous-ensemble multiplicatif. Alors $S^{-1}A \to S^{-1}B$
est un homomorphisme d'anneaux qui induit également un morphisme
dominant $\Spec(S^{-1}B) \to \Spec(S^{-1}A)$ (voir Lemmes \ref{morphisms-lemma-base-change-dominant}
et \ref{morphisms-lemma-generalizations-lift-flat}). Par conséquent,
le Lemme \ref{morphisms-lemma-relative-weak-normalization-algebra} produit une $S^{-1}A$-sous-algèbre
$(S^{-1}B)' \subset S^{-1}B$. L'énoncé signifie
que $S^{-1}B' = (S^{-1}B)'$ en tant que $S^{-1}A$-sous-algèbres de
$S^{-1}B$.

\medskip\noindent
Pour voir que cela est vrai, nous utiliserons la construction
de $B'$ et $(S^{-1}B)'$
dans la démonstration du Lemme \ref{morphisms-lemma-relative-weak-normalization-algebra}. Dans
la première étape, nous voyons que $B'$ est l'image
réciproque de la $A_{red}$-sous-algèbre $B'_{red} \subset B_{red}$ construite
pour l'homomorphisme d'anneaux $A_{red} \to B_{red}$ et de même pour
$(S^{-1}B)'$. En notant que $S^{-1}B_{red} = (S^{-1}B)_{red}$, cela nous ramène
au cas discuté dans le paragraphe suivant.

\medskip\noindent Si
$A$ et $B$ sont réduits, nous avons construit $B'$
comme l'union des sous-algèbres $A[b_1, \ldots, b_n]$ telles que pour $i = 1, \ldots, n$ nous
avons
\begin{enumerate}
\item $b_i^2, b_i^3 \in A[b_1, \ldots, b_{i - 1}]$, ou
\item il existe un nombre premier
$p$ avec $pb_i, b_i^p \in A[b_1, \ldots, b_{i - 1}]$.
\end{enumerate}
De même pour $(S^{-1}B)' \subset S^{-1}B$. Ainsi il est clair que l'image de $B' \to B \to S^{-1}B$
est contenue dans $(S^{-1}B)'$. Pour montrer que
l'application correspondante $S^{-1}B' \to (S^{-1}B)'$ est surjective,
on utilise le Lemme \ref{morphisms-lemma-special-elements-and-localization} pour éliminer les dénominateurs
successivement ; nous omettons les détails.
\end{proof}

\begin{lemma}
\label{morphisms-lemma-relative-seminormalization-algebra}
Soit $A \to B$ un homomorphisme d'anneaux induisant un morphisme
dominant $\Spec(B) \to \Spec(A)$ des spectres. Il existe un $A$-sous-algèbre
$B' \subset B$ tel que
\begin{enumerate}
\item $\Spec(B') \to \Spec(A)$ est un homéomorphisme universel induisant des isomorphismes
sur les corps résiduels,
\item donnée une factorisation $A \to C \to B$ telle que $\Spec(C) \to \Spec(A)$ est un homéomorphisme
universel induisant des isomorphismes sur les corps résiduels, l'image
de $C \to B$ est contenue dans $B'$.
\end{enumerate}
\end{lemma}

\begin{proof}
Cette démonstration est exactement
la même que celle du Lemme \ref{morphisms-lemma-relative-weak-normalization-algebra}
sauf que nous utilisons la Proposition \ref{morphisms-proposition-universal-homeomorphism-equal-residue-fields}
au lieu de la Proposition \ref{morphisms-proposition-universal-homeomorphism}
\end{proof}

\begin{lemma}
\label{morphisms-lemma-relative-seminormalization-localization}
Soit $A \to B$ un homomorphisme d'anneaux induisant un morphisme
dominant $\Spec(B) \to \Spec(A)$ de spectres. La formation de la $A$-sous-algèbre
$B' \subset B$ dans le Lemme \ref{morphisms-lemma-relative-seminormalization-algebra} commute avec la
localisation (voir la démonstration pour explication).
\end{lemma}

\begin{proof}
La démonstration est la même
que la démonstration du Lemme \ref{morphisms-lemma-relative-weak-normalization-localization}.
\end{proof}

\begin{lemma}
\label{morphisms-lemma-relative-seminormalization}
Soit $f : Y \to X$ un morphisme de schémas quasi-compact, quasi-séparé
et dominant.
\begin{enumerate}
\item La catégorie des factorisations $Y \to X' \to X$ où $X' \to X$
est un homéomorphisme universel a un objet initial $Y \to X^{Y/wn} \to X$.
\item La catégorie des factorisations $Y \to X' \to X$ où $X' \to X$ est un
homéomorphisme universel induisant des isomorphismes sur les corps
résiduels a un objet initial $Y \to X^{Y/sn} \to X$.
\end{enumerate}
De plus, la formation de la factorisation $Y \to X^{Y/wn} \to X$ et $Y \to X^{Y/sn} \to X$ commute
avec le changement de base vers des sous-schémas ouverts de $X$.
\end{lemma}

\begin{proof}
Nous allons prouver (1) et omettre la démonstration de (2) ; de plus,
l'assertion finale découlera de la construction de la factorisation.
Nous utiliserons le Lemme \ref{morphisms-lemma-universal-homeomorphism}
sans mention supplémentaire. Tout d'abord, soit $Y \to X^{Y/n} \to X$ la normalisation
de $X$ dans $Y$, voir Définition \ref{morphisms-definition-normalization-X-in-Y}.
Pour $Y \to X' \to X$ comme en (1), nous obtenons un morphisme
unique $X^{Y/n} \to X'$ compatible avec les morphismes donnés,
voir Lemme \ref{morphisms-lemma-characterize-normalization}. Ainsi, il suffit de prouver le
lemme avec $f$ remplacé par $X^{Y/n} \to X$. Cela nous ramène
au cas étudié dans le paragraphe suivant.

\medskip\noindent
Supposons que $f$ soit entier (le reste de la démonstration fonctionne
plus généralement si $f$ est affine). Soit $U = \Spec(A)$
un ouvert affine de $X$ et soit $V = f^{-1}(U) = \Spec(B)$ l'image
réciproque dans $Y$. Alors $A \to B$ est un homomorphisme
d'anneaux qui induit un morphisme dominant sur les spectres. Par le
Lemme \ref{morphisms-lemma-relative-weak-normalization-algebra}, nous obtenons une $A$-sous-algèbre $B' \subset B$
telle qu'en posant $U^{V/wn} = \Spec(B')$, la factorisation $V \to U^{V/wn} \to U$
est initiale dans la catégorie des factorisations $V \to U' \to U$ où $U' \to U$
est un homéomorphisme universel.

\medskip\noindent
Si $U_1 \subset U_2 \subset X$ sont des ouverts affines, alors en posant $V_i = f^{-1}(U_i)$
nous obtenons un morphisme canonique
$$
\rho_{U_1}^{U_2} : U_1^{V_1/wn} \to U_1 \times_{U_2} U_2^{V_2/wn}
$$
sur $U_1$ par la propriété universelle de $U_1^{V_1/wn}$. Ces morphismes
satisfont une fonctorialité naturelle que nous laissons au lecteur
de formuler et de démontrer. De plus,
le morphisme $\rho_{U_1}^{U_2}$ est
un isomorphisme ; cela résulte du Lemme \ref{morphisms-lemma-relative-weak-normalization-localization}
à condition que $U_1 \subset U_2$ soit un ouvert
principal et, dans le cas général, peut être
réduit à ce cas par la nature fonctorielle
de ces applications et Schémas, Lemme \ref{schemes-lemma-standard-open-two-affines}
(détails omis). Ainsi, par recollement
relatif (Constructions, Lemme \ref{constructions-lemma-relative-glueing}), nous obtenons
un morphisme $X^{Y/wn} \to X$ qui se restreint
à $U^{V/wn} \to U$ sur $U$ de manière compatible avec le $\rho_{U_1}^{U_2}$.
Bien sûr, les morphismes $V \to U^{V/wn}$ se recollent pour former un
morphisme $Y \to X^{Y/wn}$ (voir
Constructions, Remarque \ref{constructions-remark-relative-glueing-functorial}) et nous obtenons
notre factorisation $Y \to X^{Y/wn} \to X$ où le deuxième morphisme est un
homéomorphisme universel.

\medskip\noindent
Enfin, soit $Y \to X' \to X$ une factorisation comme en (1).
Avec $V \to U^{V/wn} \to U \subset X$ comme ci-dessus, nous obtenons
une factorisation $V \to U \times_X X' \to U$ où la deuxième flèche
est un homéomorphisme universel et nous obtenons
un morphisme unique $g_U : U^{V/wn} \to U \times_X X'$ au-dessus de $U$ par la propriété
universelle de $U^{V/wn}$. Ces
$g_U$ sont compatibles avec les morphismes $\rho_{U_1}^{U_2}$ ; détails
omis. Ainsi il existe un morphisme unique $g : X^{Y/wn} \to X'$ au-dessus
de $X$ concordant avec $g_U$
au-dessus de $U$, voir Constructions, Remarque \ref{constructions-remark-relative-glueing-functorial}.
Cela prouve que $Y \to X^{Y/wn} \to X$ est initial dans notre catégorie
et la démonstration est terminée.
\end{proof}

\begin{definition}
\label{morphisms-definition-relative-seminormalization}
Soit $f : Y \to X$ un morphisme de schémas quasi-compact, quasi-séparé
et dominant.
\begin{enumerate}
\item La factorisation $Y \to X^{Y/sn} \to X$ construite dans le Lemme \ref{morphisms-lemma-relative-seminormalization}
partie (2) est la {\it semi-normalisation
de $X$ dans $Y$}.
\item La factorisation $Y \to X^{Y/wn} \to X$ construite dans le Lemme \ref{morphisms-lemma-relative-seminormalization}
partie (1) est la {\it normalisation
faible de $X$ dans $Y$}.
\end{enumerate}
\end{definition}

\noindent
Voici une façon de réinterpréter la semi-normalisation d'un schéma qui
localement a un nombre fini de composantes irréductibles.

\begin{lemma}
\label{morphisms-lemma-seminormal-and-normalization}
Soit $X$ un schéma tel que chaque ouvert quasi-compact possède
un nombre fini de composantes irréductibles. Soit $\nu : X^\nu \to X$
la normalisation de $X$. Alors la semi-normalisation de $X$
dans $X^\nu$ est la semi-normalisation de $X$.
En formule : $X^{sn} = X^{X^\nu/sn}$.
\end{lemma}

\begin{proof}
Soit $f : Y \to X$ comme dans (\ref{morphisms-equation-generic-points}) de sorte que $X^\nu$
soit la normalisation de $X$ dans $Y$.
La semi-normalisation $X^{sn} \to X$ de $X$ est l'objet initial dans la catégorie
des homéomorphismes universels $X' \to X$ induisant des isomorphismes
sur les corps résiduels. Puisque $Y$ est l'union disjointe des spectres
des corps résiduels aux points génériques des composantes irréductibles
de $X$, nous voyons que pour tout $X' \to X$ dans cette catégorie,
nous obtenons un relèvement canonique
$f' : Y \to X'$ de $f$. Alors, par le Lemme
\ref{morphisms-lemma-characterize-normalization}, nous obtenons un morphisme canonique
$X^\nu \to X'$. D'où, à son tour, un morphisme canonique $X^{X^\nu/sn} \to X'$
par la propriété universelle de $X^{X^\nu/sn}$.
De cette façon, nous voyons que $X^{X^\nu/sn}$ satisfait la même
propriété universelle que $X^{sn}$ possède et nous concluons.
\end{proof}

\noindent
Le lemme \ref{morphisms-lemma-seminormal-and-normalization} motive la définition suivante. Puisque
nous n'avons construit la normalisation que dans le cas
où $X$ a localement un nombre fini de composantes irréductibles,
nous nous limiterons également à ce cas pour la normalisation faible.

\begin{definition}
\label{morphisms-definition-weak-normalization}
Soit $X$ un schéma tel que tout ouvert quasi-compact ait un nombre fini
de composantes irréductibles. Nous définissons la {\it
normalisation faible} de $X$ comme la normalisation faible
$$
X^\nu \longrightarrow X^{wn} \longrightarrow X
$$
de $X$ dans la normalisation $X^\nu$
de $X$ (Définition \ref{morphisms-definition-normalization}).
En formule : $X^{wn} = X^{X^\nu/wn}$.
\end{definition}

\noindent
Combiné avec Lemme \ref{morphisms-lemma-seminormal-and-normalization}, nous voyons que pour un schéma $X$
qui a localement un nombre fini de composantes irréductibles, il existe
des morphismes canoniques
$$
X^\nu \to X^{wn} \to X^{sn} \to X
$$
Ayant fait cette définition, nous pouvons dire ce que cela signifie pour un schéma d'être
faiblement normal (à condition qu'il ait localement un nombre fini de composantes
irréductibles).

\begin{definition}
\label{morphisms-definition-weakly-normal}
Soit $X$ un schéma tel que chaque ouvert quasi-compact ait un
nombre fini de composantes irréductibles. Nous disons que $X$ est {\it
faiblement normal} si la normalisation faible
$X^{wn} \to X$ est un isomorphisme (Définition \ref{morphisms-definition-weak-normalization}).
\end{definition}

\noindent
Il découle immédiatement des définitions que pour un schéma $X$ tel que chaque
ouvert quasi-compact possède un nombre fini de composantes irréductibles, nous
avons
$$
X\text{ normal} \Rightarrow
X\text{ faiblement normal} \Rightarrow
X\text{ semi-normal}
$$
Nous pouvons déterminer le sens de la normalité faible dans le cas affine comme suit.

\begin{lemma}
\label{morphisms-lemma-affine-weakly-normal}
Soit $X = \Spec(A)$ un schéma affine qui a un nombre fini de composantes irréductibles.
Alors $X$ est faiblement normal si et seulement si
\begin{enumerate}
\item $A$ est semi-normal (Définition \ref{morphisms-definition-seminormal-ring}),
\item pour un nombre premier $p$ et $z, w \in A$ tels que (a) $z$
est un non-diviseur de zéro, (b) $w^p$ est divisible par $z^p$ et (c) $pw$
est divisible par $z$, alors $w$ est divisible par $z$.
\end{enumerate}
\end{lemma}

\begin{proof}
Supposons que $X$ soit faiblement normal. Puisqu'un schéma faiblement normal
est semi-normal, on voit que (1) est vérifiée (par définition des schémas
faiblement normaux). En particulier, $A$ est réduit.
Soit $p, z, w$ comme dans (2). Choisissez $x, y \in A$ de sorte
que $z^p x = w^p$ et $zy = pw$. Alors
$p^p x = y^p$. L'homomorphisme d'anneaux $A \to C = A[t]/(t^p - x, pt - y)$ induit un homéomorphisme universel
sur les spectres. La normalisation $X^\nu$
de $X$ est le spectre de la clôture
intégrale $A'$ de $A$ dans l'anneau total
des fractions de $A$, voir Lemme \ref{morphisms-lemma-description-normalization}.
Notons que $a = w/z \in A'$ puisque $a^p = x$.
Ainsi, nous avons un homomorphisme d'$A$-algèbres $A \to C \to A'$
envoyant $t$ sur $a$.
À ce stade, la propriété définissante $X = X^{wn} = X^{X^\nu/wn}$ d'être faiblement
normal nous dit que l'image de $C \to A'$ est contenue dans $A$. Ainsi,
nous trouvons $a \in A$ comme désiré.

\medskip\noindent
Inversement, supposons (1) et (2). Soit $A'$ comme dans le
paragraphe précédent. Nous devons montrer que $X^{X^\nu/wn} = X$. Par
construction dans la démonstration du Lemme \ref{morphisms-lemma-relative-weak-normalization-algebra}, le
schéma $X^{X^\nu/wn}$ est le spectre du sous-anneau
de $A'$ qui est l'union des sous-anneaux $A[a_1, \ldots, a_n] \subset A'$ tels
que pour $i = 1, \ldots, n$ nous avons
\begin{enumerate}
\item[(a)] $a_i^2, a_i^3 \in A[a_1, \ldots, a_{i - 1}]$, ou
\item[(b)] il existe un nombre premier
$p$ avec $pa_i, a_i^p \in A[a_1, \ldots, a_{i - 1}]$.
\end{enumerate}
Nous pouvons alors utiliser (1) et (2) pour voir de manière inductive que $a_1, \ldots, a_n \in A$ ; nous
omettons les détails. Par conséquent, nous avons $X = X^{X^\nu/wn}$ et donc $X$
est faiblement normal.
\end{proof}

\noindent
Voici le lemme obligatoire.

\begin{lemma}
\label{morphisms-lemma-locally-weakly-normal}
Soit $X$ un schéma tel que chaque ouvert quasi-compact possède un nombre
fini de composantes irréductibles. Les conditions suivantes sont équivalentes :
\begin{enumerate}
\item Le schéma $X$ est faiblement normal.
\item Pour chaque ouvert affine $U \subset X$, l'anneau $\mathcal{O}_X(U)$
satisfait les conditions (1) et (2 du Lemme \ref{morphisms-lemma-affine-weakly-normal}.
\item Il existe un recouvrement ouvert affine $X = \bigcup U_i$ tel
que chaque anneau $\mathcal{O}_X(U_i)$
satisfasse aux conditions (1) et (2) du Lemme \ref{morphisms-lemma-affine-weakly-normal}.
\item Il existe un recouvrement ouvert $X = \bigcup X_j$ tel que chaque
sous-schéma ouvert $X_j$ soit faiblement normal.
\end{enumerate}
De plus, si $X$ est faiblement normal, alors tout sous-schéma ouvert est faiblement normal.
\end{lemma}

\begin{proof}
La condition pour que $X$ soit faiblement normal est que
le morphisme $X^{wn} = X^{X^\nu/wn} \to X$ soit un isomorphisme. Puisque la construction de
$X^\nu \to X$ commute avec le changement de base vers les sous-schémas
ouverts et que la construction de $X^{X^\nu/wn}$ commute avec le
changement de base vers les sous-schémas ouverts de $X$ (Lemme \ref{morphisms-lemma-relative-seminormalization}), le
lemme est clair.
\end{proof}










\section{Théorème principal de Zariski (version algébrique)}
\label{morphisms-section-Zariski}

\noindent
Ceci est la version que vous pouvez prouver en utilisant uniquement
des méthodes algébriques. Avant de pouvoir prouver des versions
plus puissantes (pour les morphismes non affines), nous
devons développer plus d'outils. Voir Cohomologie
des
schémas, section \ref{coherent-section-applications-formal-functions}
et Compléments sur les morphismes, section \ref{more-morphisms-section-application-etale-neighbourhoods}.

\begin{theorem}[Version algébrique du Théorème principal de Zariski]
\label{morphisms-theorem-main-theorem}
Soit $f : Y \to X$ un morphisme affine de schémas. Supposons que $f$
soit de type fini. Soit $X'$
la normalisation de $X$ dans $Y$. Illustration :
$$
\xymatrix{
Y \ar[rd]_f \ar[rr]_{f'} & & X' \ar[ld]^\nu \\
& X &
}
$$
Alors il existe un sous-schéma ouvert $U' \subset X'$ tel que
\begin{enumerate}
\item $(f')^{-1}(U') \to U'$ soit un isomorphisme, et
\item $(f')^{-1}(U') \subset Y$ est l'ensemble des points en lesquels $f$
est quasi-fini.
\end{enumerate}
\end{theorem}

\begin{proof}
Il y a une réduction immédiate au cas où $X$ et donc $Y$
sont affines. Disons $X = \Spec(R)$ et $Y = \Spec(A)$.
Alors $X' = \Spec(A')$, où $A'$ est la clôture
intégrale de $R$ dans $A$, voir Définitions
\ref{morphisms-definition-integral-closure} et \ref{morphisms-definition-normalization-X-in-Y}. Par Algèbre,
Théorème \ref{algebra-theorem-main-theorem} pour tout $y \in Y$ en
lequel $f$ est quasi-fini, il existe un ouvert $U'_y \subset X'$
tel que $(f')^{-1}(U'_y) \to U'_y$ soit un isomorphisme. Posons
$U' = \bigcup U'_y$ où $y \in Y$ parcourt tous les points où $f$
est quasi-fini. Il reste à montrer que $f$ est
quasi-fini en tous les points de $(f')^{-1}(U')$. Si
$y \in (f')^{-1}(U')$ avec image $x \in X$, alors nous voyons que
$Y_x \to X'_x$ est un isomorphisme dans un voisinage de $y$. Par conséquent,
il n'existe aucun point de $Y_x$ qui se spécialise en $y$,
puisque cela est vrai pour $f'(y)$ dans $X'_x$, voir Lemme
\ref{morphisms-lemma-integral-fibres}. D'après le Lemme \ref{morphisms-lemma-quasi-finite-at-point-characterize} partie (3), cela implique que
$f$ est quasi-fini en $y$.
\end{proof}

\noindent
Nous pouvons utiliser la version algébrique du Théorème principal de Zariski pour montrer
que l'ensemble des points où un morphisme est quasi-fini est ouvert.

\begin{lemma}
\label{morphisms-lemma-quasi-finite-points-open}
\begin{slogan}
Le lieu localement quasi-fini d'un morphisme est ouvert
\end{slogan}
Soit $f : X \to S$ un morphisme de schémas. L'ensemble
des points de $X$ où $f$ est quasi-fini est un $U \subset X$
ouvert. Le morphisme induit $U \to S$ est localement quasi-fini.
\end{lemma}

\begin{proof}
Supposons que $f$ soit quasi-fini
en $x$. Soient $x \in U = \Spec(A) \subset X$, $V = \Spec(R) \subset S$ des ouverts
affines comme dans la Définition \ref{morphisms-definition-quasi-finite}. D'après le Théorème
\ref{morphisms-theorem-main-theorem} ci-dessus ou Algèbre, Lemme \ref{algebra-lemma-quasi-finite-open},
l'ensemble des idéaux premiers $\mathfrak q$ en lesquels
$R \to A$ est quasi-fini est ouvert dans $\Spec(A)$. Puisque
tous correspondent à des points de $X$ où $f$
est quasi-fini, on obtient la première affirmation. La seconde
affirmation est évidente.
\end{proof}

\noindent
Nous améliorerons le lemme suivant aux morphismes
quasi-finis séparés
plus tard,
voir Compléments sur les morphismes, Lemme \ref{more-morphisms-lemma-quasi-finite-separated-pass-through-finite}.

\begin{lemma}
\label{morphisms-lemma-quasi-finite-affine}
Soit $f : Y \to X$ un morphisme de schémas.
Supposons
\begin{enumerate}
\item $X$ et $Y$ soient affines, et
\item $f$ est quasi-fini.
\end{enumerate}
Alors il existe un diagramme
$$
\xymatrix{
Y \ar[rd]_f \ar[rr]_j & & Z \ar[ld]^\pi \\
& X &
}
$$
avec $Z$ affine, $\pi$ fini et $j$ une immersion ouverte.
\end{lemma}

\begin{proof}
Ceci
est l'Algèbre, Lemme \ref{algebra-lemma-quasi-finite-open-integral-closure} reformulé dans
le langage des schémas.
\end{proof}

\begin{lemma}
\label{morphisms-lemma-image-nowhere-dense-quasi-finite}
Soit $f : Y \to X$ un morphisme quasi-fini de schémas. Soit $T \subset Y$
une partie fermée nulle part dense de $Y$. Alors $f(T) \subset X$
est une partie nulle part dense de $X$.
\end{lemma}

\begin{proof}
Comme dans la démonstration du Lemme \ref{morphisms-lemma-image-nowhere-dense-finite}, cela se réduit immédiatement
au cas où la base $X$ est affine. Dans ce cas,
$Y = \bigcup_{i = 1, \ldots, n} Y_i$ est une union finie d'ouverts affines (puisque $f$
est quasi-compact). Comme chaque $T \cap Y_i$ est nulle part dense,
et qu'une union finie de parties nulle part denses est elle-même
nulle part dense (voir
Topologie, Lemme \ref{topology-lemma-nowhere-dense}), il suffit de prouver
que l'image $f(T \cap Y_i)$ est nulle part dense dans $X$. Cela nous ramène
au cas où $X$ et $Y$ sont tous deux affines. À ce stade, nous
appliquons le Lemme \ref{morphisms-lemma-quasi-finite-affine} ci-dessus pour obtenir un diagramme
$$
\xymatrix{
Y \ar[rd]_f \ar[rr]_j & & Z \ar[ld]^\pi \\
& X &
}
$$
avec $Z$ affine, $\pi$ fini et $j$
une immersion ouverte. Posons $\overline{T} = \overline{j(T)} \subset Z$.
Par la Topologie, Lemme \ref{topology-lemma-image-nowhere-dense-open} nous voyons
que $\overline{T}$ est nulle part dense
dans $Z$. Puisque $f(T) \subset \pi(\overline{T})$
le lemme découle du résultat correspondant dans
le cas fini, voir Lemme \ref{morphisms-lemma-image-nowhere-dense-finite}.
\end{proof}




\section{Fibres universellement bornées}
\label{morphisms-section-universally-bounded}

\noindent
Soit $X$ un schéma sur un corps $k$. Si $X$ est fini sur
$k$, alors $X = \Spec(A)$ où $A$ est une $k$-algèbre finie. Une
autre façon de dire cela est que $X$ est fini localement libre sur $\Spec(k)$,
voir Définition \ref{morphisms-definition-finite-locally-free}. Par conséquent $X \to \Spec(k)$ a un
{\it degré} qui est un entier $d \geq 0$, à savoir $d = \dim_k(A)$.
Nous appelons parfois cela le {\it degré} du schéma (fini) $X$
sur $k$.

\begin{definition}
\label{morphisms-definition-universally-bounded}
Soit $f : X \to Y$ un morphisme de schémas.
\begin{enumerate}
\item On dit que l'entier $n$ {\it bornes les degrés des fibres
de $f$} si pour
tout $y \in Y$ la fibre $X_y$ est un schéma fini sur $\kappa(y)$
dont le degré sur $\kappa(y)$ est $\leq n$.
\item On dit que les fibres {\it de $f$ sont universellement borné}\footnote{Ceci
est probablement une notation
non standard.} s'il existe un entier $n$ qui bornes les degrés des fibres de
$f$.
\end{enumerate}
\end{definition}

\noindent
Notez qu'en particulier le nombre de points dans une fibre est lui aussi borné par $n$.
(La réciproque n'est pas vraie, même si toutes les fibres sont des schémas réduits
finis.)

\begin{lemma}
\label{morphisms-lemma-characterize-universally-bounded}
Soit $f : X \to Y$ un morphisme de schémas. Soit $n \geq 0$. Les conditions
suivantes sont équivalentes :
\begin{enumerate}
\item l'entier $n$ borne les degrés des fibres de $f$, et
\item pour tout morphisme $\Spec(k) \to Y$, où $k$ est un corps,
le produit fibré $X_k = \Spec(k) \times_Y X$ est fini sur $k$ de degré
$\leq n$.
\end{enumerate}
Dans ce cas, les fibres de $f$ sont universellement bornées et les schémas
$X_k$ ont au plus $n$ points. Plus
précisément, si $X_k = \{x_1, \ldots, x_t\}$, alors nous avons
$$
n \geq \sum\nolimits_{i = 1, \ldots, t} [\kappa(x_i) : k]
$$
\end{lemma}

\begin{proof}
L'implication (2) $\Rightarrow$ (1) est triviale. L'autre implication
tient parce que si l'image de $\Spec(k) \to Y$ est $y$,
alors $X_k = \Spec(k) \times_{\Spec(\kappa(y))} X_y$. Par définition, les fibres de $f$
étant universellement bornées signifie qu'il existe
un(e) $n$. Enfin, supposons que $X_k = \Spec(A)$.
Alors $\dim_k A = n$. Par conséquent, $A$ est artinien, tous les
idéaux premiers sont des idéaux maximaux $\mathfrak m_i$, et $A$
est le produit des localisations en ces idéaux maximaux.
Voir Algèbre, Lemme \ref{algebra-lemma-finite-dimensional-algebra} et \ref{algebra-lemma-artinian-finite-length}. Alors
$\mathfrak m_i$ correspond à $x_i$,
nous avons $A_{\mathfrak m_i} = \mathcal{O}_{X_k, x_i}$ et donc
il existe une surjection
$A \to  \bigoplus \kappa(\mathfrak m_i) = \bigoplus \kappa(x_i)$ ce qui implique l'inégalité
dans l'énoncé du lemme par algèbre
linéaire.
\end{proof}

\begin{lemma}
\label{morphisms-lemma-finite-locally-free-universally-bounded}
Si $f$ est un morphisme fini localement libre de degré $d$, alors
$d$ borne le degré des fibres de $f$.
\end{lemma}

\begin{proof}
C'est vrai car tout changement de base de $f$
est fini localement libre de
degré $d$ (Lemme \ref{morphisms-lemma-base-change-finite-locally-free}) et donc les fibres
de $f$ ont toutes le degré $d$.
\end{proof}

\begin{lemma}
\label{morphisms-lemma-composition-universally-bounded}
Une composition de morphismes avec des fibres universellement bornées
est un morphisme avec des fibres universellement bornées. Plus
précisément, supposons que $n$ borne les degrés des fibres de $f : X \to Y$
et que $m$ borne les degrés de
$g : Y \to Z$. Alors $nm$ borne les degrés des fibres de $g \circ f : X \to Z$.
\end{lemma}

\begin{proof}
Soient $f : X \to Y$ et $g : Y \to Z$ des fibres universellement bornées. Disons
que $\deg(X_y/\kappa(y)) \leq n$ pour tout $y \in Y$, et que $\deg(Y_z/\kappa(z)) \leq m$ pour tout $z \in Z$.
Soit $z \in Z$ un point. Par hypothèse,
le schéma $Y_z$ est fini sur $\Spec(\kappa(z))$. En particulier,
l'espace topologique sous-jacent
de $Y_z$ est un ensemble fini discret. Les fibres du
morphisme $f_z : X_z \to Y_z$ sont les fibres de $f$ aux
points correspondants de $Y$, qui sont des ensembles finis
discrets d'après le raisonnement ci-dessus. Par conséquent, nous concluons
que l'espace topologique sous-jacent de $X_z$ est
également un ensemble fini discret. Ainsi, $X_z$ est un schéma affine
(c'est un bel exercice ; cela découle également par exemple de
Propriétés, Lemme \ref{properties-lemma-maximal-points-affine} appliqué à l'ensemble
de tous les points de $X_z$). Écrire $X_z = \Spec(A)$, $Y_z = \Spec(B)$,
et $k = \kappa(z)$. Ensuite $k \to B \to A$ et
nous savons que (a) $\dim_k(B) \leq m$, et (b) pour chaque idéal
maximal $\mathfrak m \subset B$ nous avons
$\dim_{\kappa(\mathfrak m)}(A/\mathfrak mA) \leq n$. Nous affirmons que cela implique
que $\dim_k(A) \leq nm$. Notez que $B$
est le produit de ses localisations $B_{\mathfrak m}$, par exemple
parce que $Y_z$ est une union disjointe de schémas
à points $1$, ou par l'Algèbre, Lemmes \ref{algebra-lemma-finite-dimensional-algebra}
et \ref{algebra-lemma-artinian-finite-length}. Ainsi, nous
voyons que
$\dim_k(B) = \sum_{\mathfrak m} \dim_k(B_{\mathfrak m})$ et $\dim_k(A) = \sum_{\mathfrak m} \dim_k(A_{\mathfrak m})$, où dans les deux cas
$\mathfrak m$ parcourt les idéaux maximaux de $B$
(et non de $A$). D'après ce
qui précède, et le lemme de Nakayama (Algèbre,
Lemme \ref{algebra-lemma-NAK}), nous
voyons que chaque $A_{\mathfrak m}$ est un quotient
de $B_{\mathfrak m}^{\oplus n}$ en tant que $B_{\mathfrak m}$-module.
Par conséquent $\dim_k(A_{\mathfrak m}) \leq n \dim_k(B_{\mathfrak m})$. En rassemblant le
tout, nous voyons que
$$
\dim_k(A) = \sum\nolimits_{\mathfrak m} \dim_ta	(A_{\mathfrak m})
\leq \sum\nolimits_{\mathfrak m} n \dim_k(B_{\mathfrak m})
= n \dim_k(B) \leq nm
$$
comme souhaité.
\end{proof}

\begin{lemma}
\label{morphisms-lemma-base-change-universally-bounded}
Un changement de base d'un morphisme à fibres universellement bornées
est un morphisme à fibres universellement bornées. Plus précisément,
si $n$ borne les degrés des fibres de $f : X \to Y$ et que $Y' \to Y$
est un morphisme quelconque, alors les degrés des fibres du changement de
base $f' : Y' \times_Y X \to Y'$ sont également bornés par $n$.
\end{lemma}

\begin{proof}
Ceci est clair d'après
le résultat du Lemme \ref{morphisms-lemma-characterize-universally-bounded}.
\end{proof}

\begin{lemma}
\label{morphisms-lemma-descent-universally-bounded}
Soit $f : X \to Y$ un morphisme de schémas. Soit $Y' \to Y$
un morphisme de schémas, et soit $f' : X' = X_{Y'} \to Y'$ le
changement de base de $f$. Si $Y' \to Y$ est surjectif
et $f'$ a des fibres universellement bornées, alors $f$ a des
fibres universellement bornées. Plus précisément, si $n$ borne le degré
des fibres de $f'$, alors $n$ borne également les degrés des
fibres de $f$.
\end{lemma}

\begin{proof}
Soit $n \geq 0$ un entier bornant les degrés des fibres de $f'$. Nous affirmons que
$n$ convient également pour $f$. C'est-à-dire, si $y \in Y$ est un point,
choisissez un point $y' \in Y'$ situé au-dessus de $y$ et observez que
$$
X'_{y'} = \Spec(\kappa(y')) \times_{\Spec(\kappa(y))} X_y.
$$
Puisque $X'_{y'}$ est supposé fini de degré $\leq n$ sur $\kappa(y')$, il s'ensuit
que également $X_y$ est fini de degré $\leq n$ sur $\kappa(y)$. (Certains
détails sont omis.)
\end{proof}

\begin{lemma}
\label{morphisms-lemma-immersion-universally-bounded}
Une immersion a des fibres universellement bornées.
\end{lemma}

\begin{proof}
L'entier $n = 1$ fonctionne dans la définition.
\end{proof}

\begin{lemma}
\label{morphisms-lemma-etale-universally-bounded}
Soit $f : X \to Y$ un morphisme étale de schémas. Soit $n \geq 0$.
Les conditions suivantes sont équivalentes
\begin{enumerate}
\item l'entier $n$ borne les degrés des fibres,
\item pour tout corps $k$ et morphisme $\Spec(k) \to Y$ le
changement de base $X_k = \Spec(k) \times_Y X$ a au plus $n$ points, et
\item pour chaque $y \in Y$ et chaque clôture algébrique
séparable $\kappa(y) \subset \kappa(y)^{sep}$, le schéma $X_{\kappa(y)^{sep}}$
a au plus $n$ points.
\end{enumerate}
\end{lemma}

\begin{proof}
Cela résulte du
Lemme \ref{morphisms-lemma-characterize-universally-bounded} et du fait que les fibres
$X_y$ sont des unions disjointes de spectres d'extensions
de corps finies et séparables de $\kappa(y)$,
voir Lemme \ref{morphisms-lemma-etale-over-field}.
\end{proof}

\noindent
Avoir des fibres universellement bornées est une notion absolue et non une notion relative.
C'est pourquoi la condition dans le lemme suivant est que $X$ est quasi-compact,
et non que $f$ est quasi-compact.

\begin{lemma}
\label{morphisms-lemma-locally-quasi-finite-qc-source-universally-bounded}
Soit $f : X \to Y$ un morphisme de schémas. Supposons
que
\begin{enumerate}
\item $f$ est localement quasi-fini, et
\item $X$ est quasi-compact.
\end{enumerate}
Alors $f$ a des fibres universellement bornées.
\end{lemma}

\begin{proof}
Puisque $X$ est quasi-compact, il existe un recouvrement
ouvert affine fini $X = \bigcup_{i = 1, \ldots, n} U_i$ et des ouverts affines $V_i \subset Y$,
$i = 1, \ldots, n$ tels que $f(U_i) \subset V_i$. En raison
de la nature locale de ``quasi-finitude
locale'' (voir le Lemme \ref{morphisms-lemma-quasi-finite-at-point-characterize} partie (4)), nous
voyons que les morphismes $f|_{U_i} : U_i \to V_i$ sont localement
des morphismes quasi-finis d'affines, donc quasi-finis,
voir le Lemme \ref{morphisms-lemma-quasi-finite-locally-quasi-compact}. Pour $y \in Y$,
il est clair que $X_y = \bigcup_{y \in V_i} (U_i)_y$ est un recouvrement ouvert. Il
suffit donc de prouver le lemme pour un morphisme
quasi-fini d'affines (à savoir, si $n_i$ fonctionne
pour le morphisme $f|_{U_i} : U_i \to V_i$, alors $\sum n_i$)
fonctionne pour $f$).

\medskip\noindent Supposons
que $f : X \to Y$ soit un morphisme quasi-fini de schémas affines.
D'après le Lemme \ref{morphisms-lemma-quasi-finite-affine}, nous pouvons
trouver un diagramme
$$
\xymatrix{
X \ar[rd]_f \ar[rr]_j & & Z \ar[ld]^\pi \\
& Y &
}
$$
avec $Z$ affine, $\pi$ fini et $j$ une immersion
ouverte. Puisque $j$
a des fibres universellement bornées (Lemme
\ref{morphisms-lemma-immersion-universally-bounded}), cela nous ramène à montrer que $\pi$ a
des fibres universellement bornées (Lemme \ref{morphisms-lemma-composition-universally-bounded}).

\medskip\noindent Cela
nous ramène à un morphisme de la forme $\Spec(B) \to \Spec(A)$ où $A \to B$ est fini.
Supposons que $B$ soit engendré par $x_1, \ldots, x_n$ comme $A$-module,
de sorte que
$$
A^{\oplus n} \longrightarrow B, \quad
(a_1, \ldots, a_n) \longmapsto \sum a_i x_i
$$
est une application de module $A$ surjective. Pour tout premier $\mathfrak p \subset A$, cela
induit une application surjective entre les espaces vectoriels $\kappa(\mathfrak p)$
$$
\kappa(\mathfrak p)^{\oplus n} \longrightarrow B \otimes_A \kappa(\mathfrak p)
$$
En d'autres termes, l'entier $n$ intervient dans la définition
d'un morphisme à fibres universellement bornées.
\end{proof}

\begin{lemma}
\label{morphisms-lemma-universally-bounded-permanence}
Considérons un diagramme commutatif de morphismes de schémas
$$
\xymatrix{
X \ar[rd]_g \ar[rr]_f & & Y \ar[ld]^h \\
& Z &
}
$$
Si $g$ a des fibres universellement bornées, et que $f$ est surjectif et plat,
alors $h$ a également des fibres universellement bornées. Plus précisément, si $n$
borne le degré des fibres de $g$, alors $n$ borne également
le degré des fibres de $h$.
\end{lemma}

\begin{proof}
Supposons que $g$ a des fibres universellement bornées, et que $f$
est surjectif et plat. Disons que le degré des fibres de $g$
est borné par $n \in \mathbf{N}$. Nous affirmons
que $n$ fonctionne également pour $h$. Soit
$z \in Z$. Considérons le morphisme de schémas $X_z \to Y_z$. Il est
plat et surjectif. Par hypothèse $X_z$ est un schéma
fini sur $\kappa(z)$,
en particulier il est le spectre d'un anneau artinien (par
Algèbre, Lemme \ref{algebra-lemma-finite-dimensional-algebra}). Par le Lemme \ref{morphisms-lemma-Artinian-affine} le morphisme $X_z \to Y_z$
est affine, en particulier quasi-compact. Il
en découle du Lemme \ref{morphisms-lemma-fpqc-quotient-topology} que $Y_z$
est un discret fini car cela vaut pour $X_z$.
Par conséquent, $Y_z$ est un schéma affine (c'est un bel exercice ; cela
découle également,
par exemple, de Propriétés, Lemme \ref{properties-lemma-maximal-points-affine} appliqué à
l'ensemble de tous les points de $Y_z$).
Écrivez $Y_z = \Spec(B)$ et $X_z = \Spec(A)$. Alors $A$
est fidèlement plat sur $B$, donc $B \subset A$.
Par conséquent, $\dim_k(B) \leq \dim_k(A) \leq n$ comme souhaité.
\end{proof}





\section{Divers}
\label{morphisms-section-misc}

\noindent
Résultats qui ne trouvent pas leur place ailleurs.

\begin{lemma}
\label{morphisms-lemma-factor-through-point}
Soit $f : Y \to X$ un morphisme de schémas. Soit $x \in X$ un point. Supposons
que $Y$ soit réduit et que $f(Y)$ soit contenu au sens d'ensemble
dans $\{x\}$. Alors $f$ se factorise par le morphisme
canonique $x = \Spec(\kappa(x)) \to X$.
\end{lemma}

\begin{proof}
Omis. Indices : en travaillant localement sur les schémas affines, on se ramène à
un lemme d'algèbre commutative. Étant donné un homomorphisme d'anneaux $A \to B$
avec $B$ réduit tel qu'il existe un unique idéal premier
$\mathfrak p \subset A$ dans l'image de $\Spec(B) \to \Spec(A)$, alors $A \to B$ se factorise par $\kappa(\mathfrak p)$.
C'est un bel exercice.
\end{proof}

\begin{lemma}
\label{morphisms-lemma-factor-through-set-of-points}
Soit $f : Y \to X$ un morphisme de schémas. Soit $E \subset X$. On suppose que
$X$ est localement noethérien, qu'il n'y a pas de spécialisations non triviales
parmi les éléments de $E$, que $Y$ est réduit, et
$f(Y) \subset E$. Alors $f$ se factorise par $\coprod_{x \in E} x \to X$.
\end{lemma}

\begin{proof}
Lorsque $E$ est un singleton, cela
résulte du Lemme \ref{morphisms-lemma-factor-through-point}. Si $E$ est fini, alors $E$
(avec la topologie induite de $X$) est un espace discret
fini selon notre hypothèse sur les spécialisations. Ainsi,
ce cas se réduit au cas du singleton. En général, il
y a une réduction au cas où $X$ et $Y$ sont des schémas affines.
Disons que $f : Y \to X$ correspond à l'homomorphisme d'anneaux $\varphi : A \to B$.
On note $A' \subset B$ l'image de $\varphi$. Soit $E' \subset \Spec(A') \subset \Spec(A)$ l'ensemble des
idéaux premiers minimaux de $A'$. Par l'Algèbre, Lemme \ref{algebra-lemma-injective-minimal-primes-in-image}, l'ensemble
$E'$ est contenu dans l'image
de $\Spec(B) \to \Spec(A') \subset \Spec(A)$. Nous concluons que $E' \subset E$. Comme $A'$ est noethérien,
nous avons que $E'$ est fini par l'Algèbre, Lemme
\ref{algebra-lemma-Noetherian-irreducible-components}. Puisque tout autre point de l'image de $\Spec(B) \to \Spec(A)$
est une spécialisation d'un élément de $E'$ et dans $E$,
nous concluons que l'image est contenue dans $E'$ (par
notre hypothèse sur les spécialisations entre les points de $E$).
Ainsi, nous réduisons au cas où $E$ est fini,
ce que nous avons traité ci-dessus.
\end{proof}








% OK, start here.
%

\chapter{Cohomologie des schémas}



\phantomsection
\label{coherent-section-phantom}


\section{Introduction}
\label{coherent-section-introduction}

\noindent
Dans ce chapitre, nous démontrons d'abord plusieurs résultats sur la cohomologie des
faisceaux quasi-cohérents. Une référence fondamentale est \cite{EGA}.
Cela fait, nous développerons la cohomologie des
faisceaux cohérents dans le cadre noethérien. Voir \cite{FAC}.







\section{Cohomologie de {\v C}ech des faisceaux quasi-cohérents}
\label{coherent-section-cech-quasi-coherent}

\noindent
Soit $X$ un schéma.
Soit $U \subset X$ un ouvert affine.
Rappelons qu'un {\it recouvrement standard} de $U$ est un recouvrement
de la forme $\mathcal{U} : U = \bigcup_{i = 1}^n D(f_i)$,
où $f_1, \ldots, f_n \in \Gamma(U, \mathcal{O}_X)$ engendrent
l'idéal unité, voir
Schémas, définition \ref{schemes-definition-standard-covering}.

\begin{lemma}
\label{coherent-lemma-cech-cohomology-quasi-coherent-trivial}
Soit $X$ un schéma.
Soit $\mathcal{F}$ un $\mathcal{O}_X$-module quasi-cohérent.
Soit $\mathcal{U} : U = \bigcup_{i = 1}^n D(f_i)$ un recouvrement
standard d'un ouvert affine de $X$.
Alors $\check{H}^p(\mathcal{U}, \mathcal{F}) = 0$ pour
tout $p > 0$.
\end{lemma}

\begin{proof}
Écrivons $U = \Spec(A)$ pour un anneau $A$.
Autrement dit, $f_1, \ldots, f_n$ sont des éléments de $A$
qui engendrent l'idéal unité de $A$.
Écrivons $\mathcal{F}|_U = \widetilde{M}$ pour un $A$-module $M$.
Clairement, le complexe de {\v C}ech
$\check{\mathcal{C}}^\bullet(\mathcal{U}, \mathcal{F})$
s'identifie au complexe
$$
\prod\nolimits_{i_0} M_{f_{i_0}} \to
\prod\nolimits_{i_0i_1} M_{f_{i_0}f_{i_1}} \to
\prod\nolimits_{i_0i_1i_2} M_{f_{i_0}f_{i_1}f_{i_2}} \to
\ldots
$$
Il s'agit de montrer que le complexe augmenté
\begin{equation}
\label{coherent-equation-extended}
0 \to
M \to
\prod\nolimits_{i_0} M_{f_{i_0}} \to
\prod\nolimits_{i_0i_1} M_{f_{i_0}f_{i_1}} \to
\prod\nolimits_{i_0i_1i_2} M_{f_{i_0}f_{i_1}f_{i_2}} \to
\ldots
\end{equation}
(dont nous avons étudié la troncature dans
Algèbre, lemme \ref{algebra-lemma-cover-module}) est exact.
Il suffit de montrer que (\ref{coherent-equation-extended})
est exact après localisation en un idéal premier $\mathfrak p$, voir
Algèbre, lemme \ref{algebra-lemma-characterize-zero-local}.
En fait, nous montrerons que le complexe augmenté localisé
en $\mathfrak p$ est homotope à zéro.

\medskip\noindent
Il existe un indice $i$ tel que $f_i \not \in \mathfrak p$.
Choisissons et fixons un tel indice $i_{\text{fix}}$. Remarquons que
$M_{f_{i_{\text{fix}}}, \mathfrak p} = M_{\mathfrak p}$. De même,
dans une localisation en un produit $f_{i_0} \ldots f_{i_p}$ et en $\mathfrak p$,
nous pouvons omettre tout $f_{i_j}$ tel que $i_j = i_{\text{fix}}$.
Définissons une homotopie
$$
h :
\prod\nolimits_{i_0 \ldots i_{p + 1}}
M_{f_{i_0} \ldots f_{i_{p + 1}}, \mathfrak p}
\longrightarrow
\prod\nolimits_{i_0 \ldots i_p}
M_{f_{i_0} \ldots f_{i_p}, \mathfrak p}
$$
par la formule
$$
h(s)_{i_0 \ldots i_p} = s_{i_{\text{fix}} i_0 \ldots i_p}
$$
(Ceci est ``dual'' de l'homotopie dans la démonstration de
Cohomologie, lemme \ref{cohomology-lemma-homology-complex}.)
Autrement dit, $h : \prod_{i_0} M_{f_{i_0}, \mathfrak p} \to M_\mathfrak p$
est la projection sur le facteur
$M_{f_{i_{\text{fix}}}, \mathfrak p} = M_{\mathfrak p}$ et, en général,
l'application $h$ est la projection sur les facteurs
$M_{f_{i_{\text{fix}}} f_{i_1} \ldots f_{i_{p + 1}}, \mathfrak p}
= M_{f_{i_1} \ldots f_{i_{p + 1}}, \mathfrak p}$. Calculons
\begin{align*}
(dh + hd)(s)_{i_0 \ldots i_p}
& =
\sum\nolimits_{j = 0}^p
(-1)^j
h(s)_{i_0 \ldots \hat i_j \ldots i_p}
+
d(s)_{i_{\text{fix}} i_0 \ldots i_p}\\
& =
\sum\nolimits_{j = 0}^p
(-1)^j
s_{i_{\text{fix}} i_0 \ldots \hat i_j \ldots i_p}
+
s_{i_0 \ldots i_p}
+
\sum\nolimits_{j = 0}^p
(-1)^{j + 1}
s_{i_{\text{fix}} i_0 \ldots \hat i_j \ldots i_p} \\
& =
s_{i_0 \ldots i_p}
\end{align*}
Cela montre que l'application identique est homotope à zéro, comme voulu.
\end{proof}

\noindent
Le lemme suivant dit notamment que pour tout schéma affine
$X$ et tout faisceau quasi-cohérent $\mathcal{F}$ sur $X$, on a
$$
H^p(X, \mathcal{F}) = 0
$$
pour tout $p > 0$.

\begin{lemma}
\label{coherent-lemma-quasi-coherent-affine-cohomology-zero}
\begin{slogan}
Théorème d'annulation de Serre : la cohomologie supérieure s'annule sur les schémas
affines pour les modules quasi-cohérents.
\end{slogan}
Soit $X$ un schéma.
Soit $\mathcal{F}$ un $\mathcal{O}_X$-module
quasi-cohérent. Pour tout ouvert affine $U \subset X$ nous avons
$H^p(U, \mathcal{F}) = 0$ pour tout $p > 0$.
\end{lemma}

\begin{proof}
Nous allons appliquer Cohomologie,
lemme \ref{cohomology-lemma-cech-vanish-basis}.
Comme base $\mathcal{B}$ pour la topologie de $X$, nous utiliserons les
ouverts affines de $X$.
Comme ensemble $\text{Cov}$ de recouvrements ouverts, nous allons utiliser les recouvrements ouverts
standards des ouverts affines de $X$.
Nous vérifions que les conditions (1), (2) et (3) de
Cohomologie, lemme \ref{cohomology-lemma-cech-vanish-basis}
sont satisfaites. Notez que l'intersection d'ouverts standards d'un ouvert affine est encore un
ouvert standard. Par conséquent, la propriété (1) est
vérifiée. Ces recouvrements forment un système cofinal de recouvrements ouverts de tout élément
de $\mathcal{B}$, voir
Schémas, lemme \ref{schemes-lemma-standard-open}.
Donc (2) est
vérifiée. Enfin, la condition (3) du lemme découle du
lemme \ref{coherent-lemma-cech-cohomology-quasi-coherent-trivial}.
\end{proof}

\noindent
Voici une version relative de l'annulation de la cohomologie des faisceaux quasi-cohérents sur
les schémas affines.

\begin{lemma}
\label{coherent-lemma-relative-affine-vanishing}
Soit $f : X \to S$ un morphisme de schémas.
Soit $\mathcal{F}$ un $\mathcal{O}_X$-module quasi-cohérent.
Si $f$ est affine, alors $R^if_*\mathcal{F} = 0$ pour tout $i > 0$.
\end{lemma}

\begin{proof}
D'après Cohomologie,
lemme \ref{cohomology-lemma-describe-higher-direct-images}
le faisceau
$R^if_*\mathcal{F}$ est le faisceau associé au préfaisceau
$V \mapsto H^i(f^{-1}(V), \mathcal{F}|_{f^{-1}(V)})$. Par
hypothèse, chaque fois que $V$ est affine, $f^{-1}(V)$ est
affine ; voir Morphismes, définition \ref{morphisms-definition-affine}.
Par le lemme \ref{coherent-lemma-quasi-coherent-affine-cohomology-zero} nous concluons que
$H^i(f^{-1}(V), \mathcal{F}|_{f^{-1}(V)}) = 0$
chaque fois que $V$ est affine. Puisque $S$ possède une base constituée d'ouverts affines, le
faisceau est nul.
\end{proof}

\begin{lemma}
\label{coherent-lemma-relative-affine-cohomology}
Soit $f : X \to S$ un morphisme affine de schémas.
Soit $\mathcal{F}$ un $\mathcal{O}_X$-module quasi-cohérent.
Alors $H^i(X, \mathcal{F}) = H^i(S, f_*\mathcal{F})$ pour tout $i \geq 0$.
\end{lemma}

\begin{proof}
Cela résulte du lemme \ref{coherent-lemma-relative-affine-vanishing}
et de la suite spectrale de Leray. Voir
Cohomologie, lemme \ref{cohomology-lemma-apply-Leray}.
\end{proof}

\noindent
Les deux lemmes suivants précisent quand la cohomologie de {\v C}ech peut
être utilisée pour calculer la cohomologie de modules quasi-cohérents.

\begin{lemma}
\label{coherent-lemma-affine-diagonal}
Soit $X$ un schéma. Les éléments suivants sont équivalents
\begin{enumerate}
\item $X$ a une diagonale affine $\Delta : X \to X \times X$,
\item pour tous ouverts affines $U, V \subset X$, l'intersection
$U \cap V$ est affine, et
\item il existe un recouvrement ouvert $\mathcal{U} : X = \bigcup_{i \in I} U_i$
tel que $U_{i_0 \ldots i_p}$ est ouvert affine pour tous les $p \ge 0$ et tous les
$i_0, \ldots, i_p \in I$.
\end{enumerate}
En particulier, ces conditions sont satisfaites si $X$ est séparé.
\end{lemma}

\begin{proof}
Supposons que $X$ ait une diagonale affine. Soient $U, V \subset X$ des ouverts affines.
Alors $U \cap V = \Delta^{-1}(U \times V)$ est affine. Ainsi (2) est
vraie. Il est immédiat que (2) implique (3). Inversement, s'il existe un
recouvrement de $X$ comme dans (3), alors $X \times X = \bigcup U_i \times U_{i'}$
est un recouvrement ouvert affine, et on voit que
$\Delta^{-1}(U_i \times U_{i'}) = U_i \cap U_{i'}$
est affine. Alors $\Delta$ est un morphisme affine par
Morphismes, lemme \ref{morphisms-lemma-characterize-affine}.
L'assertion finale découle de Schémas, Lemme
\ref{schemes-lemma-characterize-separated}.
\end{proof}

\begin{lemma}
\label{coherent-lemma-cech-cohomology-quasi-coherent}
Soit $X$ un schéma.
Soit $\mathcal{U} : X = \bigcup_{i \in I} U_i$ un recouvrement ouvert tel que
$U_{i_0 \ldots i_p}$ soit ouvert affine pour tout $p \ge 0$ et tout
$i_0, \ldots, i_p \in I$.
Dans ce cas, pour tout faisceau quasi-cohérent $\mathcal{F}$, on a
$$
\check{H}^p(\mathcal{U}, \mathcal{F}) = H^p(X, \mathcal{F})
$$
comme $\Gamma(X, \mathcal{O}_X)$-modules, pour tout $p$.
\end{lemma}

\begin{proof}
D'après le
lemme \ref{coherent-lemma-quasi-coherent-affine-cohomology-zero},
c'est un cas particulier de
Cohomologie, Lemme
\ref{cohomology-lemma-cech-spectral-sequence-application}.
\end{proof}







\section{Annulation de la cohomologie}
\label{coherent-section-vanishing}

\noindent
Nous avons vu que, sur un schéma affine, les groupes de cohomologie supérieurs
de tout faisceau quasi-cohérent s'annulent
(lemme \ref{coherent-lemma-quasi-coherent-affine-cohomology-zero}).
Il s'avère que cela caractérise aussi
les schémas affines. Nous donnons deux versions.

\begin{lemma}
\label{coherent-lemma-quasi-compact-h1-zero-covering}
\begin{reference}
\cite{Serre-criterion}, \cite[II, Théorème 5.2.1 (d') et IV (1.7.17)]{EGA}
\end{reference}
\begin{slogan}
Le critère d'affinité de Serre.
\end{slogan}
Soit $X$ un schéma.
Supposons que
\begin{enumerate}
\item $X$ soit quasi-compact,
\item pour tout faisceau quasi-cohérent d'idéaux
$\mathcal{I} \subset \mathcal{O}_X$ nous avons $H^1(X, \mathcal{I}) = 0$.
\end{enumerate}
Alors $X$ est affine.
\end{lemma}

\begin{proof}
Soit $x \in X$ un point fermé. Soit $U \subset X$ un voisinage ouvert
affine de $x$. Écrivons $U = \Spec(A)$ et soit
$\mathfrak m \subset A$ l'idéal maximal correspondant à $x$.
Posons $Z = X \setminus U$ et $Z' = Z \cup \{x\}$.
Par Schémas, lemme \ref{schemes-lemma-reduced-closed-subscheme},
il existe des faisceaux quasi-cohérents d'idéaux
$\mathcal{I}$, resp.\ $\mathcal{I}'$, définissant
les sous-schémas fermés réduits $Z$, resp.\ $Z'$.
Considérons la suite exacte courte
$$
0 \to \mathcal{I}' \to \mathcal{I} \to \mathcal{I}/\mathcal{I}' \to 0.
$$
Puisque $x$ est un point fermé de $X$ et $x \not \in Z$ nous voyons que
$\mathcal{I}/\mathcal{I}'$ est supporté en $x$. En fait, la restriction
de $\mathcal{I}/\mathcal{I'}$ à $U$ correspond au $A$-module
$A/\mathfrak m$. On voit donc que $\Gamma(X, \mathcal{I}/\mathcal{I'})
= A/\mathfrak m$. Puisque par hypothèse $H^1(X, \mathcal{I}') = 0$
nous voyons qu'il existe une section globale $f \in \Gamma(X, \mathcal{I})$
qui s'envoie sur l'élément $1 \in A/\mathfrak m$ en tant que section de
$\mathcal{I}/\mathcal{I'}$. Manifestement, nous avons
$x \in X_f \subset U$. Cela implique que $X_f = D(f_A)$ où
$f_A$ est l'image de $f$ dans $A = \Gamma(U, \mathcal{O}_X)$.
En particulier, $X_f$ est affine.

\medskip\noindent
Considérons la réunion $W = \bigcup X_f$ portant sur les $f \in \Gamma(X, \mathcal{O}_X)$
tels que $X_f$ soit affine. Il est clair que $W$ est ouvert dans $X$.
D'après les arguments ci-dessus, chaque point fermé de
$X$ est contenu dans $W$. Le sous-ensemble fermé $X \setminus W$ de $X$
est également quasi-compact
(voir Topologie, lemme \ref{topology-lemma-closed-in-quasi-compact}).
Il a donc un point fermé s'il n'est pas vide (voir
Topologie, lemme \ref{topology-lemma-quasi-compact-closed-point}).
Cela contredirait le fait que tous les points fermés appartiennent à
$W$. On conclut donc $X = W$.

\medskip\noindent
Choisissons un nombre fini de $f_1, \ldots, f_n \in \Gamma(X, \mathcal{O}_X)$
tels que $X = X_{f_1} \cup \ldots \cup X_{f_n}$ et tels que chaque
$X_{f_i}$ soit affine. C'est possible comme nous l'avons vu ci-dessus.
Par Propriétés, lemme \ref{properties-lemma-characterize-affine}
pour terminer la preuve il suffit
de montrer que $f_1, \ldots, f_n$ engendrent l'idéal unité dans
$\Gamma(X, \mathcal{O}_X)$. Considérons la suite exacte courte
$$
\xymatrix{
0 \ar[r] &
\mathcal{F} \ar[r] &
\mathcal{O}_X^{\oplus n} \ar[rr]^{f_1, \ldots, f_n} & &
\mathcal{O}_X \ar[r] &
0
}
$$
La flèche définie par $f_1, \ldots, f_n$ est surjective puisque
l'ouverture $X_{f_i}$ recouvrent $X$. Notons $\mathcal{F}$ le noyau de
cette application surjective.
Remarquons que $\mathcal{F}$ possède une filtration
$$
0 = \mathcal{F}_0 \subset \mathcal{F}_1 \subset
\ldots \subset \mathcal{F}_n = \mathcal{F}
$$
dont chaque sous-quotient $\mathcal{F}_i/\mathcal{F}_{i - 1}$ est
isomorphe à un faisceau quasi-cohérent d'idéaux. Plus
précisément, on peut prendre $\mathcal{F}_i$ égal à l'intersection de
$\mathcal{F}$ avec les $i$ premiers facteurs directs de
$\mathcal{O}_X^{\oplus n}$.
L'hypothèse
du lemme implique que $H^1(X, \mathcal{F}_i/\mathcal{F}_{i - 1}) = 0$
pour tout $i$. Cela implique que
$H^1(X, \mathcal{F}_2) = 0$, car ce groupe se place entre
$H^1(X, \mathcal{F}_1)$ et $H^1(X, \mathcal{F}_2/\mathcal{F}_1)$.
En continuant ainsi on en déduit que $H^1(X, \mathcal{F}) = 0$.
Nous concluons donc que l'application
$$
\xymatrix{
\bigoplus\nolimits_{i = 1, \ldots, n} \Gamma(X, \mathcal{O}_X)
\ar[rr]^{f_1, \ldots, f_n} & &
\Gamma(X, \mathcal{O}_X)
}
$$
est surjective comme souhaité.
\end{proof}

\noindent
Notons que, si $X$ est un schéma noethérien, tout faisceau quasi-cohérent
d'idéaux est automatiquement un faisceau cohérent d'idéaux, donc un faisceau quasi-cohérent
d'idéaux de type fini. Par conséquent, le
lemme précédent et le lemme suivant s'appliquent tous deux dans ce cas.

\begin{lemma}
\label{coherent-lemma-quasi-separated-h1-zero-covering}
\begin{reference}
\cite{Serre-criterion}, \cite[II, Théorème 5.2.1]{EGA}
\end{reference}
\begin{slogan}
Le critère d'affinité de Serre.
\end{slogan}
Soit $X$ un schéma. Supposons que
\begin{enumerate}
\item $X$ soit quasi-compact,
\item $X$ soit quasi-séparé, et
\item $H^1(X, \mathcal{I}) = 0$ pour tout faisceau
quasi-cohérent d'idéaux $\mathcal{I}$ de type fini.
\end{enumerate}
Alors $X$ est affine.
\end{lemma}

\begin{proof}
Par
Propriétés, lemme \ref{properties-lemma-quasi-coherent-colimit-finite-type},
tout faisceau quasi-cohérent d'idéaux est une limite inductive filtrante
de faisceaux quasi-cohérents d'idéaux de type fini.
Par Cohomologie, lemme \ref{cohomology-lemma-quasi-separated-cohomology-colimit}
prenant la cohomologie sur $X$ commute aux limites inductives
filtrantes. On voit donc que $H^1(X, \mathcal{I}) = 0$
pour tout faisceau quasi-cohérent d'idéaux sur $X$. Autrement
dit, le lemme \ref{coherent-lemma-quasi-compact-h1-zero-covering} s'applique.
\end{proof}

\noindent
Les arguments précédents fournissent une condition suffisante pour qu'un faisceau
inversible soit ample. Nous avertissons toutefois le lecteur que cette
condition n'est pas nécessaire.

\begin{lemma}
\label{coherent-lemma-quasi-compact-h1-zero-invertible}
Soit $X$ un schéma. Soit $\mathcal{L}$ un $\mathcal{O}_X$-module
inversible. Supposons que
\begin{enumerate}
\item $X$ soit quasi-compact,
\item pour tout faisceau quasi-cohérent d'idéaux
$\mathcal{I} \subset \mathcal{O}_X$
il existe un $n \geq 1$ tel que
$H^1(X, \mathcal{I} \otimes_{\mathcal{O}_X} \mathcal{L}^{\otimes n}) = 0$.
\end{enumerate}
Alors $\mathcal{L}$ est ample.
\end{lemma}

\begin{proof}
Cela se démontre exactement de la même manière que le
lemme \ref{coherent-lemma-quasi-compact-h1-zero-covering}.
Soit $x \in X$ un point fermé. Soit $U \subset X$ un voisinage ouvert
affine de $x$ tel que $\mathcal{L}|_U \cong \mathcal{O}_U$.
Écrivons $U = \Spec(A)$ et soit
$\mathfrak m \subset A$ l'idéal maximal correspondant à $x$.
Posons $Z = X \setminus U$ et $Z' = Z \cup \{x\}$.
Par Schémas, lemme \ref{schemes-lemma-reduced-closed-subscheme},
il existe des faisceaux quasi-cohérents d'idéaux
$\mathcal{I}$, resp.\ $\mathcal{I}'$, définissant
les sous-schémas fermés réduits $Z$, resp.\ $Z'$.
Considérons la suite exacte courte
$$
0 \to \mathcal{I}' \to \mathcal{I} \to \mathcal{I}/\mathcal{I}' \to 0.
$$
Pour chaque $n \geq 1$ nous obtenons une suite exacte courte
$$
0 \to \mathcal{I}' \otimes_{\mathcal{O}_X} \mathcal{L}^{\otimes n}
\to \mathcal{I} \otimes_{\mathcal{O}_X} \mathcal{L}^{\otimes n} \to
\mathcal{I}/\mathcal{I}' \otimes_{\mathcal{O}_X} \mathcal{L}^{\otimes n} \to 0.
$$
D'après notre hypothèse, nous pouvons choisir $n$ tel que
$H^1(X, \mathcal{I}' \otimes_{\mathcal{O}_X} \mathcal{L}^{\otimes n}) = 0$.
Puisque $x$ est un point fermé de $X$ et $x \not \in Z$, nous voyons que
$\mathcal{I}/\mathcal{I}'$ est supporté en $x$. En fait, la restriction
de $\mathcal{I}/\mathcal{I'}$ à $U$ correspond au $A$-module
$A/\mathfrak m$. Puisque $\mathcal{L}$ est trivial sur $U$,
nous voyons que la restriction de
$\mathcal{I}/\mathcal{I}' \otimes_{\mathcal{O}_X} \mathcal{L}^{\otimes n}$
à $U$ correspond également au $A$-module $A/\mathfrak m$.
Ainsi, nous voyons que
$\Gamma(X, \mathcal{I}/\mathcal{I'} \otimes_{\mathcal{O}_X}
\mathcal{L}^{\otimes n}) = A/\mathfrak m$.
Par notre choix de $n$, nous voyons qu'il existe une section globale
$s \in \Gamma(X, \mathcal{I} \otimes_{\mathcal{O}_X} \mathcal{L}^{\otimes n})$
qui s'envoie sur l'élément $1 \in A/\mathfrak m$. Manifestement, nous avons
$x \in X_s \subset U$ parce que $s$ s'annule aux points de $Z$.
Cela implique que $X_s = D(f)$ où
$f \in A$ est l'image de $s$ dans $A \cong \Gamma(U, \mathcal{L}^{\otimes n})$.
En particulier, $X_s$ est affine.

\medskip\noindent
Considérons la réunion $W = \bigcup X_s$ portant sur les
$s \in \Gamma(X, \mathcal{L}^{\otimes n})$ pour $n \geq 1$
tels que $X_s$ soit affine. Il est clair que $W$ est ouvert dans $X$.
D'après les arguments ci-dessus, chaque point fermé de
$X$ est contenu dans $W$. Le sous-ensemble fermé $X \setminus W$ de $X$
est également quasi-compact
(voir Topologie, lemme \ref{topology-lemma-closed-in-quasi-compact}).
Il possède donc un point fermé s'il n'est pas vide (voir
Topologie, lemme \ref{topology-lemma-quasi-compact-closed-point}).
Cela contredirait le fait que tous les points fermés appartiennent à
$W$. Nous concluons donc $X = W$. Cela signifie que $\mathcal{L}$
est ample par Propriétés, définition \ref{properties-definition-ample}.
\end{proof}

\noindent
Il existe une variante du lemme \ref{coherent-lemma-quasi-compact-h1-zero-invertible}
pour les faisceaux d'idéaux de type fini ; nous la formulerons et la démontrerons ici
si nous en avons besoin.

\begin{lemma}
\label{coherent-lemma-criterion-affine-morphism}
Soit $f : X \to Y$ un morphisme quasi-compact, avec $X$ et $Y$ quasi-séparés.
Si $R^1f_*\mathcal{I} = 0$ pour tout faisceau quasi-cohérent d'idéaux
$\mathcal{I}$ sur $X$, alors $f$ est affine.
\end{lemma}

\begin{proof}
Soit $V \subset Y$ un sous-schéma ouvert affine. Il suffit de montrer que
$U = f^{-1}(V)$ est affine. Le morphisme d'inclusion $V \to Y$ est quasi-compact
d'après Schémas, lemme \ref{schemes-lemma-quasi-compact-permanence}.
Par conséquent, le changement de base $U \to X$ est quasi-compact, voir
Schémas, lemme \ref{schemes-lemma-quasi-compact-preserved-base-change}.
Ainsi, tout faisceau quasi-cohérent d'idéaux $\mathcal{I}$ sur $U$
s'étend à un faisceau quasi-cohérent d'idéaux sur $X$, voir
Propriétés, lemme \ref{properties-lemma-extend-trivial}.
Puisque la formation de $R^1f_*$ est locale sur $Y$
(Cohomologie, section \ref{cohomology-section-locality}),
nous concluons que $R^1(U \to V)_*\mathcal{I} = 0$ par l'hypothèse du
lemme. Ainsi, par la suite spectrale de Leray
(Cohomologie, lemme \ref{cohomology-lemma-Leray}),
nous concluons que $H^1(U, \mathcal{I}) = H^1(V, (U \to V)_*\mathcal{I})$.
Puisque $(U \to V)_*\mathcal{I}$ est quasi-cohérent d'après
Schémas, lemme \ref{schemes-lemma-push-forward-quasi-coherent}, nous avons
$H^1(V, (U \to V)_*\mathcal{I}) = 0$ d'après le
lemme \ref{coherent-lemma-quasi-coherent-affine-cohomology-zero}.
Ainsi, on conclut que $U$ est affine d'après le
lemme \ref{coherent-lemma-quasi-compact-h1-zero-covering}.
\end{proof}















\section{Quasi-cohérence des images directes supérieures}
\label{coherent-section-quasi-coherence}

\noindent
Nous avons vu que les groupes de cohomologie supérieure d'un module quasi-cohérent sur un schéma
affine sont nuls. Pour les schémas quasi-compacts et (quasi-)séparés $X$, il
en résulte l'annulation des groupes de cohomologie des
faisceaux quasi-cohérents au-delà d'un certain degré. Toutefois, $X$ n'est pas nécessairement
de dimension cohomologique finie, car celle-ci est définie par l'annulation de la cohomologie
de {\it tous} les $\mathcal{O}_X$-modules.

\begin{lemma}[Principe de récurrence]
\label{coherent-lemma-induction-principle}
\begin{reference}
\cite[Proposition 3.3.1]{BvdB}
\end{reference}
Soit $X$ un schéma quasi-compact et quasi-séparé. Soit $P$ une propriété
des ouverts quasi-compacts de $X$. Supposons que
\begin{enumerate}
\item $P$ est vérifiée pour tout ouvert affine de $X$,
\item si $U$ est un ouvert quasi-compact, $V$ un ouvert affine, et si
$P$ est vérifiée pour $U$, $V$ et $U \cap V$, alors
$P$ est vérifiée pour $U \cup V$.
\end{enumerate}
Alors $P$ est vérifiée pour tout ouvert quasi-compact de $X$
et en particulier pour $X$.
\end{lemma}

\begin{proof}
Montrons d'abord par récurrence que $P$ est vérifiée pour les ouverts quasi-compacts {\it séparés}
$W \subset X$. En effet, un tel ouvert s'écrit
$W = U_1 \cup \ldots \cup U_n$ et l'on raisonne par récurrence sur $n$ en utilisant
la propriété (2) avec $U = U_1 \cup \ldots \cup U_{n - 1}$ et $V = U_n$.
Cela est possible car
$U \cap V = (U_1 \cap U_n) \cup \ldots \cup (U_{n - 1} \cap U_n)$
est aussi une réunion de $n - 1$ sous-schémas ouverts affines, d'après
Schémas, lemme \ref{schemes-lemma-characterize-separated},
appliqué aux ouverts affines $U_i$ et $U_n$ de $W$.
Cela étant, pour tout ouvert quasi-compact $W \subset X$, on
peut raisonner par récurrence sur le nombre d'ouverts affines nécessaires pour recouvrir $W$,
en employant le même procédé et en observant que l'ouvert quasi-compact
$U_i \cap U_n$ est séparé, comme sous-schéma ouvert du schéma affine $U_n$.
\end{proof}

\begin{lemma}
\label{coherent-lemma-vanishing-nr-affines}
\begin{slogan}
Pour les schémas à diagonale affine, la cohomologie des
modules quasi-cohérents s'annule en degrés supérieurs au nombre d'ouverts affines
nécessaires pour former un recouvrement.
\end{slogan}
Soit $X$ un schéma quasi-compact à diagonale affine (par exemple
si $X$ est séparé).
Soit $t = t(X)$ le nombre minimal d'ouverts affines nécessaires pour
recouvrir $X$. Alors $H^n(X, \mathcal{F}) = 0$ pour tout $n \geq t$ et tout
faisceau quasi-cohérent $\mathcal{F}$.
\end{lemma}

\begin{proof}
Première démonstration.
Raisonnons par récurrence sur $t$.
Si $t = 1$, le résultat découle du
lemme \ref{coherent-lemma-quasi-coherent-affine-cohomology-zero}.
Si $t > 1$, écrivons $X = U \cup V$, où $V$ est un ouvert affine et
$U = U_1 \cup \ldots \cup U_{t - 1}$ une réunion de $t - 1$ ouverts affines.
Notons que, dans ce cas,
$U \cap V =  (U_1 \cap V) \cup \ldots (U_{t - 1} \cap V)$
est aussi une réunion de $t - 1$ sous-schémas ouverts affines.
En effet, la diagonale étant affine, l'intersection de deux
ouverts affines est affine, voir lemme \ref{coherent-lemma-affine-diagonal}.
Appliquons la suite exacte longue de Mayer-Vietoris
$$
0 \to
H^0(X, \mathcal{F}) \to
H^0(U, \mathcal{F}) \oplus H^0(V, \mathcal{F}) \to
H^0(U \cap V, \mathcal{F}) \to
H^1(X, \mathcal{F}) \to \ldots
$$
voir Cohomologie, lemme \ref{cohomology-lemma-mayer-vietoris}.
Par récurrence, les groupes $H^i(U, \mathcal{F})$,
$H^i(V, \mathcal{F})$, $H^i(U \cap V, \mathcal{F})$ sont nuls pour
$i \geq t - 1$. Il en résulte immédiatement que $H^i(X, \mathcal{F})$
est nul pour $i \geq t$.

\medskip\noindent
Deuxième démonstration.
Soit $\mathcal{U} : X = \bigcup_{i = 1}^t U_i$ un recouvrement ouvert affine
fini. Puisque $X$ est à diagonale affine, les intersections multiples
$U_{i_0 \ldots i_p}$ sont toutes affines, voir
lemme \ref{coherent-lemma-affine-diagonal}.
D'après le lemme \ref{coherent-lemma-cech-cohomology-quasi-coherent}, les groupes de cohomologie de {\v C}ech
$\check{H}^p(\mathcal{U}, \mathcal{F})$
coïncident avec les groupes de cohomologie. D'après
Cohomologie, lemme \ref{cohomology-lemma-alternating-usual},
les groupes de cohomologie de {\v C}ech se calculent au moyen du complexe alterné de
{\v C}ech $\check{\mathcal{C}}_{alt}^\bullet(\mathcal{U}, \mathcal{F})$.
Le recouvrement ayant $t$ éléments, on voit immédiatement
que $\check{\mathcal{C}}_{alt}^p(\mathcal{U}, \mathcal{F}) = 0$
pour tout $p \geq t$. Le résultat s'ensuit.
\end{proof}

\begin{lemma}
\label{coherent-lemma-affine-diagonal-universal-delta-functor}
Soit $X$ un schéma quasi-compact à diagonale affine
(par exemple si $X$ est séparé). Alors
\begin{enumerate}
\item pour tout $\mathcal{O}_X$-module quasi-cohérent $\mathcal{F}$,
il existe une injection $\mathcal{F} \to \mathcal{F}'$
de $\mathcal{O}_X$-modules
quasi-cohérents telle que $H^p(X, \mathcal{F}') = 0$ pour tout $p \geq 1$, et
\item $\{H^n(X, -)\}_{n \geq 0}$
est un $\delta$-foncteur universel de $\QCoh(\mathcal{O}_X)$ dans
$\textit{Ab}$.
\end{enumerate}
\end{lemma}

\begin{proof}
Soit $X = \bigcup U_i$ un recouvrement ouvert affine fini.
Posons $U = \coprod U_i$ et notons $j : U \to X$
le morphisme induisant les immersions ouvertes données $U_i \to X$.
Puisque $U$ est un schéma affine et que $X$ est à diagonale affine,
le morphisme $j$ est affine, voir
Morphismes, lemme \ref{morphisms-lemma-affine-permanence}.
Pour tout $\mathcal{O}_X$-module $\mathcal{F}$, il
existe un morphisme canonique $\mathcal{F} \to j_*j^*\mathcal{F}$.
Ce morphisme est injectif, comme on le vérifie sur les fibres :
si $x \in U_i$, on a une factorisation
$$
\mathcal{F}_x \to (j_*j^*\mathcal{F})_x
\to (j^*\mathcal{F})_{x'} = \mathcal{F}_x
$$
où $x' \in U$ est le point $x$ considéré comme point de $U_i \subset U$.
Si maintenant $\mathcal{F}$ est quasi-cohérent, alors $j^*\mathcal{F}$
est quasi-cohérent sur le schéma affine $U$, donc sa cohomologie
supérieure est nulle d'après le
lemme \ref{coherent-lemma-quasi-coherent-affine-cohomology-zero}. On
a alors $H^p(X, j_*j^*\mathcal{F}) = 0$ pour
$p > 0$ d'après le lemme \ref{coherent-lemma-relative-affine-cohomology},
puisque $j$ est affine. Cela prouve (1).
Enfin, on voit que le morphisme
$H^p(X, \mathcal{F}) \to H^p(X, j_*j^*\mathcal{F})$
est nul, et l'assertion (2) découle de
Homologie, lemme \ref{homology-lemma-efface-implies-universal}.
\end{proof}

\begin{lemma}
\label{coherent-lemma-vanishing-nr-affines-quasi-separated}
Soit $X$ un schéma quasi-compact et quasi-séparé.
Soit $X = U_1 \cup \ldots \cup U_n$ un recouvrement ouvert
dont chaque $U_i$ est quasi-compact et séparé (par exemple affine).
Posons
$$
d = \max\nolimits_{I \subset \{1, \ldots, n\}}
\left(|I| + t(\bigcap\nolimits_{i \in I} U_i) - 1\right)
$$
où $t(U)$ est le nombre minimal d'ouverts affines nécessaires pour recouvrir
le schéma $U$. Alors $H^p(X, \mathcal{F}) = 0$ pour tout $p \geq d$ et tout
faisceau quasi-cohérent $\mathcal{F}$.
\end{lemma}

\begin{proof}
Notons que, puisque $X$ est quasi-séparé et les $U_i$ quasi-compacts, les nombres
$t(\bigcap_{i \in I} U_i)$ sont finis. Raisonnons par récurrence sur $n$.
Si $n = 1$, le résultat découle du lemme \ref{coherent-lemma-vanishing-nr-affines}.
Si $n > 1$, écrivons $X = U \cup V$, avec $U = U_1 \cup \ldots \cup U_{n - 1}$
et $V = U_n$. Appliquons la suite exacte longue de Mayer-Vietoris
$$
0 \to
H^0(X, \mathcal{F}) \to
H^0(U, \mathcal{F}) \oplus H^0(V, \mathcal{F}) \to
H^0(U \cap V, \mathcal{F}) \to
H^1(X, \mathcal{F}) \to \ldots
$$
voir Cohomologie, lemme \ref{cohomology-lemma-mayer-vietoris}.
Pour achever la démonstration lorsque $q \geq d$, montrons que
$H^q(V, \mathcal{F})$, $H^q(U, \mathcal{F})$ et
$H^{q - 1}(U \cap V, \mathcal{F})$ sont
nuls. Le cas $n = 1$ donne $H^q(V, \mathcal{F}) = 0$ pour
$q \geq t(V) = t(U_n)$. Comme $t(V) \leq d$, on obtient l'annulation voulue. Par
l'hypothèse de récurrence, on a $H^q(U, \mathcal{F}) = 0$ pour
$$
q \geq \max\nolimits_{I \subset \{1, \ldots, n - 1\}}
\left(|I| + t(\bigcap\nolimits_{i \in I} U_i) - 1\right)
$$
L'entier de droite étant inférieur ou égal à $d$, on
obtient l'annulation voulue. Enfin, appliquons l'hypothèse de récurrence à l'ouvert
$U \cap V = (U_1 \cap U_n) \cup \ldots \cup (U_{n - 1} \cap U_n)$ pour
obtenir l'annulation $H^q(U \cap V, \mathcal{F}) = 0$ pour
$$
q \geq \max\nolimits_{I \subset \{1, \ldots, n - 1\}}
\left(|I| + t(U_n \cap \bigcap\nolimits_{i \in I} U_i) - 1\right)
$$
L'entier de droite étant strictement inférieur à $d$,
le lemme s'ensuit.
\end{proof}

\begin{lemma}
\label{coherent-lemma-quasi-coherence-higher-direct-images}
Soit $f : X \to S$ un morphisme de schémas.
Supposons que $f$ soit quasi-séparé et quasi-compact.
\begin{enumerate}
\item Pour tout $\mathcal{O}_X$-module quasi-cohérent $\mathcal{F}$,
les images directes supérieures $R^pf_*\mathcal{F}$ sont quasi-cohérentes sur $S$.
\item Si $S$ est quasi-compact, il existe un entier $n = n(X, S, f)$
tel que $R^pf_*\mathcal{F} = 0$ pour tout $p \geq n$ et tout
faisceau quasi-cohérent $\mathcal{F}$ sur $X$.
\item En fait, si $S$ est quasi-compact, on peut trouver $n = n(X, S, f)$
tel que, pour tout
morphisme de schémas $S' \to S$, on ait $R^p(f')_*\mathcal{F}' = 0$
pour $p \geq n$ et tout faisceau quasi-cohérent $\mathcal{F}'$
sur $X'$. Ici $f' : X' = S' \times_S X \to S'$ est le changement de base de $f$.
\end{enumerate}
\end{lemma}

\begin{proof}
Prouvons d'abord (1). Notons que, sous les hypothèses du lemme, le faisceau
$R^0f_*\mathcal{F} = f_*\mathcal{F}$ est quasi-cohérent d'après
Schémas, lemme \ref{schemes-lemma-push-forward-quasi-coherent}.
D'après
Cohomologie, lemme \ref{cohomology-lemma-localize-higher-direct-images},
la formation des images directes supérieures commute à la
restriction aux sous-schémas ouverts. La quasi-cohérence étant locale sur $S$,
on se ramène au cas traité au paragraphe suivant.

\medskip\noindent
Démonstration de (1) lorsque $S$ est affine. Utilisons le principe de récurrence.
Puisque $f$ est quasi-compact et quasi-séparé, $X$
est quasi-compact et quasi-séparé. Pour un ouvert quasi-compact $U \subset X$
et $a = f|_U$, notons $P(U)$ la propriété suivante :
$R^pa_*\mathcal{F}$ est quasi-cohérent sur $S$ pour tout module quasi-cohérent
$\mathcal{F}$ sur $U$ et tout $p \geq 0$. Comme $P(X)$ est l'assertion (1), il suffit
de vérifier les conditions (1) et (2) du lemme \ref{coherent-lemma-induction-principle}.
Si $U$ est affine, alors $P(U)$ est vérifiée car $R^pa_*\mathcal{F} = 0$ pour
$p \geq 1$ (d'après le lemme \ref{coherent-lemma-relative-affine-vanishing} et
Morphismes, lemme \ref{morphisms-lemma-morphism-affines-affine}),
et nous avons déjà établi le résultat pour $p = 0$ au premier
paragraphe. Soient ensuite $U \subset X$ un ouvert quasi-compact et $V \subset X$
un ouvert affine; supposons $P(U)$, $P(V)$ et $P(U \cap V)$ vérifiées.
Posons $a = f|_U$, $b = f|_V$, $c = f|_{U \cap V}$ et $g = f|_{U \cup V}$.
Pour tout $\mathcal{O}_{U \cup V}$-module quasi-cohérent $\mathcal{F}$,
on dispose alors de la suite de Mayer-Vietoris relative
$$
0 \to
g_*\mathcal{F} \to
a_*(\mathcal{F}|_U) \oplus b_*(\mathcal{F}|_V) \to
c_*(\mathcal{F}|_{U \cap V}) \to
R^1g_*\mathcal{F} \to \ldots
$$
voir Cohomologie, lemme \ref{cohomology-lemma-relative-mayer-vietoris}.
D'après $P(U)$, $P(V)$ et $P(U \cap V)$, les faisceaux
$R^pa_*(\mathcal{F}|_U)$, $R^pb_*(\mathcal{F}|_V)$ et
$R^pc_*(\mathcal{F}|_{U \cap V})$ sont tous quasi-cohérents.
Les résultats sur les faisceaux quasi-cohérents de
Schémas, section \ref{schemes-section-quasi-coherent},
impliquent alors que chacun des faisceaux
$R^pg_*\mathcal{F}$ est quasi-cohérent : il est le terme médian d'une suite exacte
courte dont le terme de gauche est le conoyau d'un morphisme
de faisceaux quasi-cohérents et le terme de droite le noyau d'un morphisme de faisceaux quasi-cohérents.
D'où $P(U \cup V)$, ce qui achève la démonstration de (1).

\medskip\noindent
Prouvons ensuite (3), et a fortiori (2). Choisissons un recouvrement ouvert affine
fini $S = \bigcup_{j = 1, \ldots m} S_j$. Pour chaque $j$,
choisissons un recouvrement ouvert affine fini
$f^{-1}(S_j) = \bigcup_{i = 1, \ldots t_j} U_{ji} $.
Soit
$$
d_j = \max\nolimits_{I \subset \{1, \ldots, t_j\}}
\left(|I| + t(\bigcap\nolimits_{i \in I} U_{ji})\right)
$$
l'entier obtenu dans le
lemme \ref{coherent-lemma-vanishing-nr-affines-quasi-separated}.
Nous affirmons que $n(X, S, f) = \max d_j$ convient.

\medskip\noindent
En effet, soient $S' \to S$ un morphisme de schémas et
$\mathcal{F}'$ un faisceau quasi-cohérent sur $X' = S' \times_S X$.
Montrons que $R^pf'_*\mathcal{F}' = 0$ pour $p \geq n(X, S, f)$.
Cette question étant locale sur $S'$, on peut supposer que $S'$ est affine
et s'envoie dans $S_j$ pour un certain $j$. Alors $X' = S' \times_{S_j} f^{-1}(S_j)$
est recouvert par les ouverts affines $S' \times_{S_j} U_{ji}$, $i = 1, \ldots t_j$,
et les intersections
$$
\bigcap\nolimits_{i \in I} S' \times_{S_j} U_{ji} =
S' \times_{S_j} \bigcap\nolimits_{i \in I} U_{ji}
$$
sont recouvertes par le même nombre d'ouverts affines qu'avant le changement de base. En appliquant
le
lemme \ref{coherent-lemma-vanishing-nr-affines-quasi-separated},
on obtient $H^p(X', \mathcal{F}') = 0$. La première partie de la démonstration
montre déjà que chaque $R^qf'_*\mathcal{F}'$ est quasi-cohérent,
donc a ses groupes de cohomologie supérieure nuls sur le schéma affine $S'$;
ainsi $H^0(S', R^pf'_*\mathcal{F}') = H^p(X', \mathcal{F}') = 0$
d'après Cohomologie, lemme \ref{cohomology-lemma-apply-Leray}.
Puisque $R^pf'_*\mathcal{F}'$ est quasi-cohérent,
on en déduit que $R^pf'_*\mathcal{F}' = 0$.
\end{proof}

\begin{lemma}
\label{coherent-lemma-quasi-coherence-higher-direct-images-application}
Soit $f : X \to S$ un morphisme de schémas.
Supposons que $f$ soit quasi-séparé et quasi-compact.
Supposons $S$ affine.
Pour tout $\mathcal{O}_X$-module quasi-cohérent $\mathcal{F}$,
on a
$$
H^q(X, \mathcal{F}) = H^0(S, R^qf_*\mathcal{F})
$$
pour tout $q \in \mathbf{Z}$.
\end{lemma}

\begin{proof}
Considérons la suite spectrale de Leray $E_2^{p, q} = H^p(S, R^qf_*\mathcal{F})$
convergeant vers $H^{p + q}(X, \mathcal{F})$, voir
Cohomologie, lemme \ref{cohomology-lemma-Leray}.
D'après le lemme \ref{coherent-lemma-quasi-coherence-higher-direct-images},
les faisceaux $R^qf_*\mathcal{F}$ sont quasi-cohérents.
D'après le lemme \ref{coherent-lemma-quasi-coherent-affine-cohomology-zero},
on a $E_2^{p, q} = 0$ lorsque $p > 0$.
La suite spectrale dégénère donc en $E_2$, ce qui
conclut. Voir aussi
Cohomologie, lemme \ref{cohomology-lemma-apply-Leray} (2),
pour le principe général.
\end{proof}








\section{Cohomologie et changement de base, I}
\label{coherent-section-cohomology-and-base-change}

\noindent
Soit $f : X \to S$ un morphisme de schémas.
Soit $\mathcal{F}$ un faisceau quasi-cohérent sur $X$.
Soit de plus $g : S' \to S$ un morphisme quelconque de schémas. Notons
$X' = X_{S'} = S' \times_S X$ le changement de base de $X$ et
$f' : X' \to S'$ le changement de base de $f$.
Notons aussi $g' : X' \to X$ la projection,
et posons $\mathcal{F}' = (g')^*\mathcal{F}$.
La situation est représentée par le diagramme suivant :
\begin{equation}
\label{coherent-equation-base-change-diagram}
\vcenter{
\xymatrix{
\mathcal{F}' = (g')^*\mathcal{F} &
X' \ar[r]_{g'} \ar[d]_{f'} &
X \ar[d]^f &
\mathcal{F} \\
Rf'_*\mathcal{F}' &
S' \ar[r]^g &
S &
Rf_*\mathcal{F}
}
}
\end{equation}
Voici le cas le plus simple de la propriété de changement de base envisagée.

\begin{lemma}
\label{coherent-lemma-affine-base-change}
Soit $f : X \to S$ un morphisme de schémas.
Soit $\mathcal{F}$ un $\mathcal{O}_X$-module quasi-cohérent.
Supposons $f$ affine. Dans
ce cas, $f_*\mathcal{F} \cong Rf_*\mathcal{F}$ est
un faisceau quasi-cohérent et, pour tout diagramme de changement de base
(\ref{coherent-equation-base-change-diagram})
on a
$$
g^*f_*\mathcal{F} = f'_*(g')^*\mathcal{F}.
$$
\end{lemma}

\begin{proof}
L'annulation des images directes supérieures est donnée par le
lemme \ref{coherent-lemma-relative-affine-vanishing}.
L'énoncé est local sur $S$ et $S'$. On peut donc
supposer $X = \Spec(A)$, $S = \Spec(R)$,
$S' = \Spec(R')$ et $\mathcal{F} = \widetilde{M}$,
pour un certain $A$-module $M$.
Utilisons Schémas, lemme \ref{schemes-lemma-widetilde-pullback},
pour décrire les images inverses et directes de $\mathcal{F}$. On
a $X' = \Spec(R' \otimes_R A)$ et
$\mathcal{F}'$ est le faisceau quasi-cohérent associé
à $(R' \otimes_R A) \otimes_A M$.
Le lemme se ramène donc à
l'égalité
$$
(R' \otimes_R A) \otimes_A M = R' \otimes_R M
$$
de $R'$-modules.
\end{proof}

\noindent
Dans de nombreuses situations, il suffit de connaître le cas particulier suivant
de la cohomologie et du changement de base. Il découle immédiatement
des résultats plus forts de la
section \ref{coherent-section-cohomology-and-base-change-derived},
mais son importance justifie une démonstration distincte.

\begin{lemma}[Changement de base plat]
\label{coherent-lemma-flat-base-change-cohomology}
Considérons un diagramme cartésien de schémas
$$
\xymatrix{
X' \ar[d]_{f'} \ar[r]_{g'} & X \ar[d]^f \\
S' \ar[r]^g & S
}
$$
Soit $\mathcal{F}$ un $\mathcal{O}_X$-module
quasi-cohérent, d'image inverse $\mathcal{F}' = (g')^*\mathcal{F}$.
Supposons que $g$ soit plat et que $f$ soit quasi-compact et quasi-séparé.
Pour tout $i \geq 0$,
\begin{enumerate}
\item le morphisme de changement de base
de Cohomologie, lemme \ref{cohomology-lemma-base-change-map-flat-case},
est un isomorphisme
$$
g^*R^if_*\mathcal{F} \longrightarrow R^if'_*\mathcal{F}',
$$
\item si $S = \Spec(A)$ et $S' = \Spec(B)$, alors
$H^i(X, \mathcal{F}) \otimes_A B = H^i(X', \mathcal{F}')$.
\end{enumerate}
\end{lemma}

\begin{proof}
On peut appliquer Cohomologie, lemme \ref{cohomology-lemma-base-change-map-flat-case}, dans
(1) car $g'$ est plat d'après
Morphismes, lemme \ref{morphisms-lemma-base-change-flat}.
Cela étant, l'assertion (1) découle de (2). En effet,
l'assertion (1) est locale sur $S'$, et l'on peut donc supposer $S$
et $S'$ affines. Autrement dit, on a $S = \Spec(A)$
et $S' = \Spec(B)$ comme dans
(2). Puisque $R^if_*\mathcal{F}$ est quasi-cohérent
(lemme \ref{coherent-lemma-quasi-coherence-higher-direct-images}),
c'est le $\mathcal{O}_S$-module quasi-cohérent associé au
$A$-module $H^0(S, R^if_*\mathcal{F}) = H^i(X, \mathcal{F})$
(égalité donnée par
le lemme \ref{coherent-lemma-quasi-coherence-higher-direct-images-application}). De
même, $R^if'_*\mathcal{F}'$ est le
$\mathcal{O}_{S'}$-module quasi-cohérent associé au $B$-module
$H^i(X', \mathcal{F}')$. L'image inverse par $g$ correspondant au
foncteur $- \otimes_A B$ sur les
modules (Schémas, lemme \ref{schemes-lemma-widetilde-pullback}),
il suffit de prouver (2).

\medskip\noindent
Soit $A \to B$ un homomorphisme plat d'anneaux.
Soit $X$ un schéma quasi-compact et quasi-séparé sur $A$.
Soit $\mathcal{F}$ un $\mathcal{O}_X$-module quasi-cohérent.
Posons $X_B = X \times_{\Spec(A)} \Spec(B)$ et notons
$\mathcal{F}_B$ l'image inverse de $\mathcal{F}$.
Il s'agit de montrer que le morphisme
$$
H^i(X, \mathcal{F}) \otimes_A B \longrightarrow H^i(X_B, \mathcal{F}_B)
$$
(donné par la référence citée dans l'énoncé du lemme)
est un isomorphisme.

\medskip\noindent
Lorsque $X$ est séparé, choisissons un recouvrement ouvert affine
$\mathcal{U} : X = U_1 \cup \ldots \cup U_t$ et rappelons que
$$
\check{H}^p(\mathcal{U}, \mathcal{F}) = H^p(X, \mathcal{F}),
$$
voir
lemme \ref{coherent-lemma-cech-cohomology-quasi-coherent}. Pour
le recouvrement $\mathcal{U}_B : X_B = (U_1)_B \cup \ldots \cup (U_t)_B$ obtenu
par changement de base, chaque $(U_i)_B$ est encore
affine et $X_B$ est séparé. On obtient donc
$$
\check{H}^p(\mathcal{U}_B, \mathcal{F}_B) = H^p(X_B, \mathcal{F}_B).
$$
On a la relation suivante entre les complexes de {\v C}ech
$$
\check{\mathcal{C}}^\bullet(\mathcal{U}_B, \mathcal{F}_B) =
\check{\mathcal{C}}^\bullet(\mathcal{U}, \mathcal{F}) \otimes_A B
$$
comme il résulte du
lemme \ref{coherent-lemma-affine-base-change}.
Puisque $A \to B$ est plat, cette égalité subsiste en passant à la cohomologie.

\medskip\noindent
Lorsque $X$ est quasi-séparé, choisissons un recouvrement ouvert affine
$\mathcal{U} : X = U_1 \cup \ldots \cup U_t$. Utilisons la suite spectrale reliant la cohomologie de
{\v C}ech à la cohomologie, donnée
dans Cohomologie, lemme \ref{cohomology-lemma-cech-spectral-sequence}.
Le lecteur qui souhaite éviter cette suite spectrale peut utiliser
Mayer-Vietoris et raisonner par récurrence sur $t$, comme dans la démonstration du
lemme \ref{coherent-lemma-quasi-coherence-higher-direct-images}.
La suite spectrale a pour terme $E_2$
$E_2^{p, q} = \check{H}^p(\mathcal{U}, \underline{H}^q(\mathcal{F}))$
et converge vers $H^{p + q}(X, \mathcal{F})$.
De même, on dispose d'une suite spectrale de terme $E_2$
$E_2^{p, q} = \check{H}^p(\mathcal{U}_B, \underline{H}^q(\mathcal{F}_B))$
qui converge vers $H^{p + q}(X_B, \mathcal{F}_B)$.
Puisque les intersections $U_{i_0 \ldots i_p}$ sont quasi-compactes
et séparées, le résultat du deuxième paragraphe de la démonstration donne
$\check{H}^p(\mathcal{U}_B, \underline{H}^q(\mathcal{F}_B)) =
\check{H}^p(\mathcal{U}, \underline{H}^q(\mathcal{F})) \otimes_A B$.
La platitude de $A \to B$ permet de conclure que
$H^i(X, \mathcal{F}) \otimes_A B \to H^i(X_B, \mathcal{F}_B)$
est un isomorphisme pour tout $i$, ce qui conclut.
\end{proof}

\begin{lemma}[Changement de base fini localement libre]
\label{coherent-lemma-finite-locally-free-base-change-cohomology}
Considérons un diagramme cartésien de schémas
$$
\xymatrix{
Y \ar[d]_{g} \ar[r]_h & X \ar[d]^f \\
\Spec(B) \ar[r] & \Spec(A)
}
$$
Soit $\mathcal{F}$ un $\mathcal{O}_X$-module
quasi-cohérent, d'image inverse $\mathcal{G} = h^*\mathcal{F}$.
Si $B$ est un $A$-module localement libre de type fini, alors
$H^i(X, \mathcal{F}) \otimes_A B = H^i(Y, \mathcal{G})$.
\end{lemma}

\noindent
{\bf Attention} : N'utilisez ce lemme que si vous comprenez la différence entre
celui-ci et le lemme \ref{coherent-lemma-flat-base-change-cohomology}.

\begin{proof}
Lorsque $X$ est séparé, choisissons un recouvrement ouvert affine
$\mathcal{U} : X = \bigcup_{i \in I} U_i$ et rappelons que
$$
\check{H}^p(\mathcal{U}, \mathcal{F}) = H^p(X, \mathcal{F}),
$$
voir
lemme \ref{coherent-lemma-cech-cohomology-quasi-coherent}.
Soit $\mathcal{V} : Y = \bigcup_{i \in I} g^{-1}(U_i)$
le recouvrement ouvert affine correspondant de $Y$.
Les ouverts $V_i = g^{-1}(U_i) = U_i \times_{\Spec(A)} \Spec(B)$
sont affines et $Y$ est séparé. On obtient donc
$$
\check{H}^p(\mathcal{V}, \mathcal{G}) = H^p(Y, \mathcal{G}).
$$
Nous affirmons que le morphisme de complexes de {\v C}ech
$$
\check{\mathcal{C}}^\bullet(\mathcal{U}, \mathcal{F}) \otimes_A B
\longrightarrow
\check{\mathcal{C}}^\bullet(\mathcal{V}, \mathcal{G})
$$
est un isomorphisme. En effet, puisque $B$ est de présentation finie comme $A$-module,
le produit tensoriel avec $B$ sur $A$ commute aux produits, voir
Algèbre, proposition \ref{algebra-proposition-fp-tensor}.
Il suffit donc de montrer que les morphismes
$\Gamma(U_{i_0 \ldots i_p}, \mathcal{F}) \otimes_A B \to
\Gamma(V_{i_0 \ldots i_p}, \mathcal{G})$
sont des isomorphismes, ce qui résulte du
lemme \ref{coherent-lemma-affine-base-change}.
Puisque $A \to B$ est plat, cela reste vrai en passant à la cohomologie.

\medskip\noindent
Dans le cas général, on raisonne exactement de même en utilisant un recouvrement
ouvert affine $\mathcal{U} : X = \bigcup_{i \in I} U_i$ et le
recouvrement correspondant $\mathcal{V} : Y = \bigcup_{i \in I} V_i$,
avec $V_i = g^{-1}(U_i)$ comme ci-dessus. Utilisons la suite spectrale reliant la cohomologie de
{\v C}ech à la cohomologie, donnée
dans Cohomologie, lemme \ref{cohomology-lemma-cech-spectral-sequence}.
La suite spectrale a pour terme $E_2$
$E_2^{p, q} = \check{H}^p(\mathcal{U}, \underline{H}^q(\mathcal{F}))$
et converge vers $H^{p + q}(X, \mathcal{F})$.
De même, on dispose d'une suite spectrale de terme $E_2$
$E_2^{p, q} = \check{H}^p(\mathcal{V}, \underline{H}^q(\mathcal{G}))$
qui converge vers $H^{p + q}(Y, \mathcal{G})$.
Les intersections $U_{i_0 \ldots i_p}$ étant séparées, le résultat
du paragraphe précédent donne des isomorphismes
$\Gamma(U_{i_0 \ldots i_p}, \underline{H}^q(\mathcal{F})) \otimes_A B
\to \Gamma(V_{i_0 \ldots i_p}, \underline{H}^q(\mathcal{G}))$.
Le foncteur $- \otimes_A B$ commutant aux produits et étant exact, on en déduit
que
$\check{H}^p(\mathcal{U}, \underline{H}^q(\mathcal{F})) \otimes_A B
\to \check{H}^p(\mathcal{V}, \underline{H}^q(\mathcal{G}))$
est un isomorphisme. La platitude de $A \to B$ permet de conclure que
$H^i(X, \mathcal{F}) \otimes_A B \to H^i(Y, \mathcal{G})$
est un isomorphisme pour tout $i$, ce qui conclut.
\end{proof}









\section{Limites inductives et images directes supérieures}
\label{coherent-section-colimits}

\noindent
On trouvera des résultats généraux de cette nature dans
Cohomologie, section \ref{cohomology-section-limits},
Faisceaux, lemme \ref{sheaves-lemma-directed-colimits-sections}, et
Modules, lemme \ref{modules-lemma-finite-presentation-quasi-compact-colimit}.

\begin{lemma}
\label{coherent-lemma-colimit-cohomology}
Soit $f : X \to S$ un morphisme quasi-compact et quasi-séparé de schémas.
Soit $\mathcal{F} = \colim \mathcal{F}_i$ une limite inductive filtrante
de faisceaux abéliens sur $X$. Alors, pour tout $p \geq 0$, on a
$$
R^pf_*\mathcal{F} = \colim R^pf_*\mathcal{F}_i.
$$
\end{lemma}

\begin{proof}
Nous appliquons Cohomologie, Lemme
\ref{cohomology-lemma-higher-direct-image-colimit}.
Comme les ouverts affines forment une base de la topologie de $S$,
il suffit de montrer que, pour tout ouvert affine $U \subset S$, on a
$H^p(f^{-1}U, \mathcal{F}) = \colim H^p(f^{-1}U, \mathcal{F}_i)$.
Puisque $f^{-1}U$ est quasi-compact et quasi-séparé, cela résulte
de
Cohomologie, lemme \ref{cohomology-lemma-quasi-separated-cohomology-colimit} (la
base des ouverts affines de $f^{-1}U$
vérifie les hypothèses de ce lemme).
\end{proof}










\section{Cohomologie et changement de base, II}
\label{coherent-section-cohomology-and-base-change-derived}

\noindent
Soit $f : X \to S$ un morphisme de schémas et soit $\mathcal{F}$
un $\mathcal{O}_X$-module quasi-cohérent. Si $f$ est quasi-compact
et quasi-séparé, nous souhaitons représenter $Rf_*\mathcal{F}$
par un complexe de faisceaux quasi-cohérents sur $S$. Cela résulte
du fait que les faisceaux $R^if_*\mathcal{F}$ sont quasi-cohérents
lorsque $S$ est quasi-compact et à diagonale affine, en utilisant
l'équivalence de $D_\QCoh(S)$ et de
$D(\QCoh(\mathcal{O}_S))$ ; voir
Catégories dérivées des schémas, Proposition
\ref{perfect-proposition-quasi-compact-affine-diagonal}.

\medskip\noindent
Dans cette section, nous suivrons une autre approche, qui produit un complexe
explicite doté d'une bonne propriété de changement de base. La construction
est particulièrement simple si $f$ et $S$ sont séparés ou, plus généralement,
à diagonale affine. Comme il s'agit de loin
du cas le plus fréquemment utilisé, nous le traitons séparément.

\begin{lemma}
\label{coherent-lemma-separated-case-relative-cech}
Soit $f : X \to S$ un morphisme de schémas.
Soit $\mathcal{F}$ un $\mathcal{O}_X$-module quasi-cohérent.
Supposons que $X$ soit quasi-compact et que $X$ et $S$ aient une diagonale affine
(par exemple, si $X$ et $S$ sont séparés).
On peut alors calculer $Rf_*\mathcal{F}$ comme suit :
\begin{enumerate}
\item Choisir un recouvrement ouvert affine fini
$\mathcal{U} : X = \bigcup_{i = 1, \ldots, n} U_i$.
\item Pour $i_0, \ldots, i_p \in \{1, \ldots, n\}$, noter
$f_{i_0 \ldots i_p} : U_{i_0 \ldots i_p} \to S$ la restriction de $f$
à l'intersection $U_{i_0 \ldots i_p} = U_{i_0} \cap \ldots \cap U_{i_p}$.
\item Noter $\mathcal{F}_{i_0 \ldots i_p}$ la restriction
de $\mathcal{F}$ à $U_{i_0 \ldots i_p}$.
\item Poser
$$
\check{\mathcal{C}}^p(\mathcal{U}, f, \mathcal{F}) =
\bigoplus\nolimits_{i_0 \ldots i_p}
f_{i_0 \ldots i_p *} \mathcal{F}_{i_0 \ldots i_p}
$$
et définir les différentielles
$d : \check{\mathcal{C}}^p(\mathcal{U}, f, \mathcal{F})
\to \check{\mathcal{C}}^{p + 1}(\mathcal{U}, f, \mathcal{F})$
comme dans Cohomologie, Équation (\ref{cohomology-equation-d-cech}).
\end{enumerate}
Le complexe $\check{\mathcal{C}}^\bullet(\mathcal{U}, f, \mathcal{F})$
est alors un complexe de faisceaux quasi-cohérents sur $S$, muni d'un
isomorphisme
$$
\check{\mathcal{C}}^\bullet(\mathcal{U}, f, \mathcal{F})
\longrightarrow
Rf_*\mathcal{F}
$$
dans $D^{+}(S)$. Cet isomorphisme est fonctoriel en le faisceau
quasi-cohérent $\mathcal{F}$.
\end{lemma}

\begin{proof}
Considérons la résolution
$\mathcal{F} \to {\mathfrak C}^\bullet(\mathcal{U}, \mathcal{F})$
de Cohomologie, lemme \ref{cohomology-lemma-covering-resolution}.
On a une égalité de complexes
$\check{\mathcal{C}}^\bullet(\mathcal{U}, f, \mathcal{F}) =
f_*{\mathfrak C}^\bullet(\mathcal{U}, \mathcal{F})$
de $\mathcal{O}_S$-modules
quasi-cohérents. Les morphismes $j_{i_0 \ldots i_p} : U_{i_0 \ldots i_p} \to X$
et $f_{i_0 \ldots i_p} : U_{i_0 \ldots i_p} \to S$
sont affines d'après Morphismes, lemme \ref{morphisms-lemma-affine-permanence},
et le lemme \ref{coherent-lemma-affine-diagonal}. Par
conséquent, $R^qj_{i_0 \ldots i_p *}\mathcal{F}_{i_0 \ldots i_p}$
ainsi que $R^qf_{i_0 \ldots i_p *}\mathcal{F}_{i_0 \ldots i_p}$
sont nuls pour $q > 0$ (lemme \ref{coherent-lemma-relative-affine-vanishing}). En
utilisant $f \circ j_{i_0 \ldots i_p} = f_{i_0 \ldots i_p}$ et
la suite spectrale de
Cohomologie, lemme \ref{cohomology-lemma-relative-Leray},
on en déduit que
$R^qf_*(j_{i_0 \ldots i_p *}\mathcal{F}_{i_0 \ldots i_p}) = 0$
pour $q > 0$.
Comme les termes du complexe
${\mathfrak C}^\bullet(\mathcal{U}, \mathcal{F})$ sont des sommes
directes finies des faisceaux $j_{i_0 \ldots i_p *}\mathcal{F}_{i_0 \ldots i_p}$,
le lemme d'acyclicité de Leray
(Catégories dérivées, lemme \ref{derived-lemma-leray-acyclicity})
donne
$$
Rf_* \mathcal{F} = f_*{\mathfrak C}^\bullet(\mathcal{U}, \mathcal{F}) =
\check{\mathcal{C}}^\bullet(\mathcal{U}, f, \mathcal{F})
$$
comme souhaité.
\end{proof}

\noindent
Considérons maintenant ce qui se passe après un changement de base.

\begin{lemma}
\label{coherent-lemma-base-change-complex}
Avec les notations du diagramme (\ref{coherent-equation-base-change-diagram}), supposons
que $f : X \to S$ et $\mathcal{F}$ vérifient les hypothèses du
lemme \ref{coherent-lemma-separated-case-relative-cech}. Choisissons un recouvrement
ouvert affine fini $\mathcal{U} : X = \bigcup U_i$ de $X$.
Il existe un isomorphisme canonique
$$
g^*\check{\mathcal{C}}^\bullet(\mathcal{U}, f, \mathcal{F})
\longrightarrow
Rf'_*\mathcal{F}'
$$
dans $D^{+}(S')$. De plus, si $S' \to S$ est affine, on a en fait
$$
g^*\check{\mathcal{C}}^\bullet(\mathcal{U}, f, \mathcal{F})
=
\check{\mathcal{C}}^\bullet(\mathcal{U}', f', \mathcal{F}')
$$
avec $\mathcal{U}' : X' = \bigcup U_i'$, où
$U_i' = (g')^{-1}(U_i) = U_{i, S'}$ est également affine.
\end{lemma}

\begin{proof}
En fait, on peut définir $U_i' = (g')^{-1}(U_i) = U_{i, S'}$,
que $S'$ soit ou non affine sur $S$.
Soit $\mathcal{U}' : X' = \bigcup U_i'$ le recouvrement induit de $X'$.
Nous affirmons alors que
$$
g^*\check{\mathcal{C}}^\bullet(\mathcal{U}, f, \mathcal{F})
=
\check{\mathcal{C}}^\bullet(\mathcal{U}', f', \mathcal{F}')
$$
où $\check{\mathcal{C}}^\bullet(\mathcal{U}', f', \mathcal{F}')$
est défini exactement comme dans le
lemme \ref{coherent-lemma-separated-case-relative-cech}.
Cela résulte du cas des morphismes affines
(lemme \ref{coherent-lemma-affine-base-change}) en raisonnant localement sur $S'$.
De plus, exactement comme dans la démonstration du
lemme \ref{coherent-lemma-separated-case-relative-cech},
on voit qu'il existe un isomorphisme
$$
\check{\mathcal{C}}^\bullet(\mathcal{U}', f', \mathcal{F}')
\longrightarrow
Rf'_*\mathcal{F}'
$$
dans $D^{+}(S')$, car les morphismes $U_i' \to X'$ et $U_i' \to S'$
restent affines (ce sont des changements de base de morphismes affines). Nous
omettons les détails.
\end{proof}

\noindent
Le lemme précédent affirme que le complexe
$$
\mathcal{K}^\bullet = \check{\mathcal{C}}^\bullet(\mathcal{U}, f, \mathcal{F})
$$
est un complexe borné inférieurement de faisceaux quasi-cohérents sur $S$ qui calcule
{\it universellement} les images directes supérieures de $f : X \to S$.
Il s'agit d'une propriété de ce complexe particulier,
qui n'est pas préservée lorsqu'on remplace
$\check{\mathcal{C}}^\bullet(\mathcal{U}, f, \mathcal{F})$ par
un complexe quasi-isomorphe en général ! Autrement dit, cet énoncé
n'a pas de sens dans la catégorie dérivée. La
raison en est que l'image inverse $g^*\mathcal{K}^\bullet$ n'est
{\it pas} égale, en général, à l'image inverse dérivée $Lg^*\mathcal{K}^\bullet$
de $\mathcal{K}^\bullet$ !

\medskip\noindent
Voici un cas plus général où cet énoncé peut être démontré.
Remarquons que l'hypothèse de séparation de $S$ est inoffensive dans la
plupart des applications, car le résultat sert habituellement à établir
une propriété locale de l'image dérivée
totale. La démonstration est nettement plus délicate et
utilise les hyperrecouvrements ; c'est un bon exemple de leur emploi.

\begin{lemma}
\label{coherent-lemma-hypercoverings}
Soit $f : X \to S$ un morphisme de schémas.
Soit $\mathcal{F}$ un faisceau quasi-cohérent sur $X$.
Supposons que $f$ soit quasi-compact et quasi-séparé, et
que $S$ soit quasi-compact et
séparé. Il existe un complexe $\mathcal{K}^\bullet$,
borné inférieurement, de $\mathcal{O}_S$-modules quasi-cohérents, possédant la
propriété suivante : pour tout morphisme
$g : S' \to S$, le complexe $g^*\mathcal{K}^\bullet$
représente $Rf'_*\mathcal{F}'$, avec les notations du
diagramme (\ref{coherent-equation-base-change-diagram}).
\end{lemma}

\begin{proof}
(Si $f$ est lui aussi séparé, voir le
lemme \ref{coherent-lemma-base-change-complex}.)
Les hypothèses impliquent notamment que $X$
est quasi-compact et quasi-séparé en tant que schéma.
Soit $\mathcal{B}$ l'ensemble des ouverts affines de $X$. D'après
Hyperrecouvrements,
lemme \ref{hypercovering-lemma-quasi-separated-quasi-compact-hypercovering},
il existe un hyperrecouvrement $K = (I, \{U_i\})$ tel que chaque
$I_n$ soit fini et chaque $U_i$ un ouvert affine de $X$. D'après
Hyperrecouvrements, lemme \ref{hypercovering-lemma-cech-spectral-sequence},
il existe une suite spectrale dont la page $E_2$ est
$$
E_2^{p, q} = \check{H}^p(K, \underline{H}^q(\mathcal{F}))
$$
et qui converge vers $H^{p + q}(X, \mathcal{F})$. Remarquons que
$\check{H}^p(K, \underline{H}^q(\mathcal{F}))$ est le $p$-ième groupe
de cohomologie du complexe
$$
\prod\nolimits_{i \in I_0} H^q(U_i, \mathcal{F})
\to
\prod\nolimits_{i \in I_1} H^q(U_i, \mathcal{F})
\to
\prod\nolimits_{i \in I_2} H^q(U_i, \mathcal{F})
\to \ldots
$$
Chaque $U_i$ étant affine, ce groupe est nul sauf si $q = 0$,
auquel cas on obtient
$$
\prod\nolimits_{i \in I_0} \mathcal{F}(U_i)
\to
\prod\nolimits_{i \in I_1} \mathcal{F}(U_i)
\to
\prod\nolimits_{i \in I_2} \mathcal{F}(U_i)
\to \ldots
$$
On en conclut que $R\Gamma(X, \mathcal{F})$ est calculé par
ce complexe.

\medskip\noindent
Pour tout $n$ et tout $i \in I_n$, notons $f_i : U_i \to S$ la restriction de
$f$ à $U_i$. Puisque $S$ est séparé et $U_i$ affine, ce morphisme est
affine. Considérons le complexe de faisceaux quasi-cohérents
$$
\mathcal{K}^\bullet = (
\prod\nolimits_{i \in I_0} f_{i, *}\mathcal{F}|_{U_i}
\to
\prod\nolimits_{i \in I_1} f_{i, *}\mathcal{F}|_{U_i}
\to
\prod\nolimits_{i \in I_2} f_{i, *}\mathcal{F}|_{U_i}
\to \ldots )
$$
sur $S$. Comme dans
Hyperrecouvrements, lemme \ref{hypercovering-lemma-cech-spectral-sequence},
on obtient un morphisme $\mathcal{K}^\bullet \to Rf_*\mathcal{F}$ dans
$D(\mathcal{O}_S)$ en choisissant une résolution injective de $\mathcal{F}$
(nous omettons les détails). Considérons un schéma affine quelconque $V$ et un morphisme
$g : V \to S$. Le changement de base $X_V$ possède alors l'hyperrecouvrement
$K_V = (I, \{U_{i, V}\})$ obtenu par changement de base. De plus,
$g^*f_{i, *}\mathcal{F} = f_{i, V, *}(g')^*\mathcal{F}|_{U_{i, V}}$.
Les arguments précédents montrent donc que $\Gamma(V, g^*\mathcal{K}^\bullet)$
calcule $R\Gamma(X_V, (g')^*\mathcal{F})$.
Cela achève la démonstration du lemme, puisqu'il suffit d'établir l'égalité des complexes
localement pour la topologie de Zariski sur $S'$.
\end{proof}

\noindent
Le lemme suivant est une variante du changement de base plat.

\begin{lemma}
\label{coherent-lemma-base-change-tensor-with-flat}
Considérons un diagramme cartésien de schémas
$$
\xymatrix{
X' \ar[d]_{f'} \ar[r]_{g'} & X \ar[d]^f \\
S' \ar[r]^g & S
}
$$
Soit $\mathcal{F}$ un $\mathcal{O}_X$-module quasi-cohérent.
Soit $\mathcal{G}$ un $\mathcal{O}_{S'}$-module quasi-cohérent
et plat sur $S$. Supposons que $f$ soit quasi-compact et quasi-séparé.
Pour tout $i \geq 0$, il existe une identification
$$
\mathcal{G} \otimes_{\mathcal{O}_{S'}} g^*R^if_*\mathcal{F} =
R^if'_*\left((f')^*\mathcal{G}
\otimes_{\mathcal{O}_{X'}} (g')^*\mathcal{F}\right)
$$
\end{lemma}

\begin{proof}
Construisons un morphisme du membre de gauche vers celui de droite. On dispose d'abord du morphisme de changement de base
$Lg^*Rf_*\mathcal{F} \to Rf'_*L(g')^*\mathcal{F}$. On dispose aussi du morphisme
d'adjonction $\mathcal{G} \to Rf'_*L(f')^*\mathcal{G}$. À l'aide du cup-produit relatif,
on obtient
\begin{align*}
\mathcal{G}
\otimes_{\mathcal{O}_{S'}}^\mathbf{L}
Lg^*Rf_*\mathcal{F}
& \to
Rf'_*L(f')^*\mathcal{G}
\otimes_{\mathcal{O}_{S'}}^\mathbf{L}
Rf'_*L(g')^*\mathcal{F} \\
& \to
Rf'_*\left(L(f')^*\mathcal{G}
\otimes_{\mathcal{O}_{X'}}^\mathbf{L}
L(g')^*\mathcal{F}\right) \\
& \to
Rf'_*\left((f')^*\mathcal{G}
\otimes_{\mathcal{O}_{X'}} (g')^*\mathcal{F}\right)
\end{align*}
où la flèche centrale est donnée par le cup-produit relatif ; voir
Cohomologie, remarque \ref{cohomology-remark-cup-product}.
La source du composé est
$$
\mathcal{G} \otimes_{\mathcal{O}_{S'}}^\mathbf{L} Lg^*Rf_*\mathcal{F} =
\mathcal{G} \otimes_{g^{-1}\mathcal{O}_S}^\mathbf{L} g^{-1}Rf_*\mathcal{F}
$$
d'après Cohomologie, lemme \ref{cohomology-lemma-variant-derived-pullback}.
Comme $\mathcal{G}$ est plat en tant que faisceau de $g^{-1}\mathcal{O}_S$-modules
et que $g^{-1}$ est un foncteur exact, il s'agit d'un complexe dont le
$i$-ième faisceau de cohomologie est
$\mathcal{G} \otimes_{g^{-1}\mathcal{O}_S} g^{-1}R^if_*\mathcal{F} =
\mathcal{G} \otimes_{\mathcal{O}_{S'}} g^*R^if_*\mathcal{F}$.
On obtient ainsi, globalement, les morphismes du membre de gauche vers celui de droite
dans l'égalité du lemme. Pour montrer que ce morphisme est un isomorphisme, on peut
raisonner localement sur $S'$. Nous pouvons donc supposer, et supposons, que $S$ et $S'$
sont des schémas affines.

\medskip\noindent
Démonstration dans le cas où $S$ et $S'$ sont affines. Écrivons $S = \Spec(A)$ et $S' = \Spec(B)$,
et supposons que $\mathcal{G}$ corresponde au $B$-module $N$, supposé plat sur
$A$. Puisque $S$ est affine, $X$ est quasi-compact et quasi-séparé.
Nous achevons la démonstration à l'aide d'un hyperrecouvrement ; si $X$ est séparé
ou à diagonale affine, on peut utiliser un recouvrement de {\v C}ech.
Soit $\mathcal{B}$ l'ensemble des ouverts affines de $X$. D'après
Hyperrecouvrements,
lemme \ref{hypercovering-lemma-quasi-separated-quasi-compact-hypercovering},
il existe un hyperrecouvrement $K = (I, \{U_i\})$ de $X$ tel que chaque
$I_n$ soit fini et chaque $U_i$ un ouvert affine de $X$. D'après
Hyperrecouvrements, lemme \ref{hypercovering-lemma-cech-spectral-sequence},
il existe une suite spectrale dont la page $E_2$ est
$$
E_2^{p, q} = \check{H}^p(K, \underline{H}^q(\mathcal{F}))
$$
et qui converge vers $H^{p + q}(X, \mathcal{F})$. Comme chaque $U_i$ est affine et
que $\mathcal{F}$ est quasi-cohérent, la valeur de
$\underline{H}^q(\mathcal{F})$ sur $U_i$ est nulle pour $q > 0$. La suite
spectrale dégénère donc, et on en conclut que les modules de
cohomologie $H^q(X, \mathcal{F})$ sont calculés par
$$
\prod\nolimits_{i \in I_0} \mathcal{F}(U_i)
\to
\prod\nolimits_{i \in I_1} \mathcal{F}(U_i)
\to
\prod\nolimits_{i \in I_2} \mathcal{F}(U_i)
\to \ldots
$$
Remarquons ensuite que le changement de base de notre hyperrecouvrement à
$S'$ est un hyperrecouvrement de $X' = S' \times_S X$.
Les schémas $S' \times_S U_i$ sont eux aussi affines,
et l'on a
$$
\left((f')^*\mathcal{G} \otimes_{\mathcal{O}_{S'}} (g')^*\mathcal{F}\right)
(S' \times_S U_i) =
N \otimes_A \mathcal{F}(U_i)
$$
On en conclut que les modules de cohomologie
$H^q(X', (f')^*\mathcal{G} \otimes_{\mathcal{O}_{S'}} (g')^*\mathcal{F})$
sont calculés par
$$
N \otimes_A \left(
\prod\nolimits_{i \in I_0} \mathcal{F}(U_i)
\to
\prod\nolimits_{i \in I_1} \mathcal{F}(U_i)
\to
\prod\nolimits_{i \in I_2} \mathcal{F}(U_i)
\to \ldots
\right)
$$
Puisque $N$ est plat sur $A$, on en déduit que
$$
H^q(X', (f')^*\mathcal{G} \otimes_{\mathcal{O}_{S'}} (g')^*\mathcal{F})
=
N \otimes_A H^q(X, \mathcal{F})
$$
Il s'agit précisément de la traduction algébrique de l'énoncé
à démontrer, ce qui achève la démonstration.
\end{proof}







\section{Cohomologie de l'espace projectif}
\label{coherent-section-cohomology-projective-space}

\noindent
Dans cette section, nous calculons la cohomologie des faisceaux structuraux
tordus sur $\mathbf{P}^n_S$, où $S$ est un
schéma. Rappelons que $\mathbf{P}^n_S$ a été défini comme le produit fibré
$
\mathbf{P}^n_S = S \times_{\Spec(\mathbf{Z})} \mathbf{P}^n_{\mathbf{Z}}
$
dans Constructions, définition \ref{constructions-definition-projective-space}.
On a montré l'égalité
$$
\mathbf{P}^n_S = \underline{\text{Proj}}_S(\mathcal{O}_S[T_0, \ldots, T_n])
$$
dans Constructions, lemme \ref{constructions-lemma-projective-space-bundle}.
En particulier, l'espace projectif est un cas particulier de fibré projectif.
Si $S = \Spec(R)$ est affine, on a
$$
\mathbf{P}^n_S = \mathbf{P}^n_R = \text{Proj}(R[T_0, \ldots, T_n]).
$$
Toutes ces identifications sont compatibles entre elles ainsi qu'avec les
constructions des faisceaux structuraux tordus $\mathcal{O}_{\mathbf{P}^n_S}(d)$.

\medskip\noindent
Avant d'énoncer le résultat, introduisons quelques notations.
Soit $R$ un anneau.
Rappelons que $R[T_0, \ldots, T_n]$ est une
$R$-algèbre graduée dans laquelle chaque $T_i$ est homogène de degré $1$.
Notons $(R[T_0, \ldots, T_n])_d$ la composante homogène de degré $d$.
C'est un $R$-module libre de type fini de rang $\binom{n + d}{d}$
si $d \geq 0$, et nul
sinon. Il possède une base formée des monômes $T_0^{e_0} \ldots T_n^{e_n}$
tels que $\sum e_i = d$. Nous utiliserons aussi la notation suivante :
$R[\frac{1}{T_0}, \ldots, \frac{1}{T_n}]$ désigne l'anneau $\mathbf{Z}$-gradué
dans lequel $\frac{1}{T_i}$ est de degré $-1$. En particulier, le module
$\mathbf{Z}$-gradué sur $R[\frac{1}{T_0}, \ldots, \frac{1}{T_n}]$
$$
\frac{1}{T_0 \ldots T_n} R[\frac{1}{T_0}, \ldots, \frac{1}{T_n}]
$$
qui figure dans l'énoncé ci-dessous est nul en degrés
$\geq -n$, est libre de générateur $\frac{1}{T_0 \ldots T_n}$
en degré $-n - 1$, et est libre de rang $(-1)^n\binom{n + d}{d}$ pour
$d \leq -n - 1$.

\begin{lemma}
\label{coherent-lemma-cohomology-projective-space-over-ring}
\begin{reference}
\cite[III Proposition 2.1.12]{EGA}
\end{reference}
Soit $R$ un anneau.
Soit $n \geq 0$ un entier.
On a
$$
H^q(\mathbf{P}^n, \mathcal{O}_{\mathbf{P}^n_R}(d)) =
\left\{
\begin{matrix}
(R[T_0, \ldots, T_n])_d & \text{si} & q = 0,\ d \geq 0 \\
\left(\frac{1}{T_0 \ldots T_n} R[\frac{1}{T_0}, \ldots, \frac{1}{T_n}]\right)_d
& \text{si} & q = n,\ d < 0 \\
0 & \text{si} & q, n, d\text{ hors des cas ci-dessus}
\end{matrix}
\right.
$$
comme $R$-modules.
\end{lemma}

\begin{proof}
Utilisons le recouvrement ouvert affine standard
$$
\mathcal{U} : \mathbf{P}^n_R = \bigcup\nolimits_{i = 0}^n D_{+}(T_i)
$$
pour calculer la cohomologie au moyen du complexe de {\v C}ech.
Le lemme \ref{coherent-lemma-cech-cohomology-quasi-coherent}
le permet, puisque toute intersection d'un nombre fini d'ouverts affines $D_{+}(T_i)$ est encore un
ouvert affine standard (voir
Constructions, section \ref{constructions-section-proj}).
On peut même utiliser le complexe de {\v C}ech alterné ou ordonné, d'après
Cohomologie, lemmes \ref{cohomology-lemma-ordered-alternating} et
\ref{cohomology-lemma-alternating-usual}.

\medskip\noindent
Nous munissons $\{0, \ldots, n\}$ de son ordre
usuel. Le complexe considéré a donc pour termes
$$
\check{\mathcal{C}}_{ord}^p(\mathcal{U}, \mathcal{O}_{\mathbf{P}_R}(d))
=
\bigoplus\nolimits_{i_0 < \ldots < i_p}
(R[T_0, \ldots, T_n, \frac{1}{T_{i_0} \ldots T_{i_p}}])_d
$$
De plus, les différentielles sont données par la formule usuelle
$$
d(s)_{i_0 \ldots i_{p + 1}} =
\sum\nolimits_{j = 0}^{p + 1} (-1)^j s_{i_0 \ldots \hat i_j \ldots i_{p + 1}}
$$
voir Cohomologie, section \ref{cohomology-section-alternating-cech}.
Notons que chaque terme de ce complexe possède une
$\mathbf{Z}^{n + 1}$-graduation naturelle : on déclare que le monôme
$T_0^{e_0} \ldots T_n^{e_n}$ est homogène de poids
$(e_0, \ldots, e_n) \in \mathbf{Z}^{n + 1}$. Il est clair que la
différentielle ci-dessus respecte cette graduation. On a donc
$$
\check{\mathcal{C}}_{ord}^\bullet(\mathcal{U}, \mathcal{O}_{\mathbf{P}_R}(d))
=
\bigoplus\nolimits_{\vec{e} \in \mathbf{Z}^{n + 1}}
\check{\mathcal{C}}^\bullet(\vec{e})
$$
où tous les facteurs directs de droite n'interviennent pas nécessairement (voir
ci-dessous). Pour calculer les modules de
cohomologie du complexe, il suffit donc de calculer la cohomologie de ses
composantes homogènes, puis d'en prendre la somme directe.

\medskip\noindent
Fixons $\vec{e} = (e_0, \ldots, e_n) \in \mathbf{Z}^{n + 1}$. Pour que ce
poids intervienne dans le complexe ci-dessus, il faut supposer
$e_0 + \ldots + e_n = d$ (sinon, il intervient bien entendu pour un autre tordu
du faisceau structural). Sous cette hypothèse, posons
$$
NEG(\vec{e}) = \{i \in \{0, \ldots, n\} \mid e_i < 0\}.
$$
Avec cette notation, le facteur direct de poids $\vec{e}$
$\check{\mathcal{C}}^\bullet(\vec{e})$ du complexe de {\v C}ech ci-dessus a
pour termes
$$
\check{\mathcal{C}}^p(\vec{e})
=
\bigoplus\nolimits_{i_0 < \ldots < i_p,
\ NEG(\vec{e}) \subset \{i_0, \ldots, i_p\}}
R \cdot T_0^{e_0} \ldots T_n^{e_n}
$$
Autrement dit, les termes correspondant aux $i_0 < \ldots < i_p$ tels
que $NEG(\vec{e})$ ne soit pas contenu dans $\{i_0 \ldots i_p\}$ sont nuls.
La différentielle du complexe $\check{\mathcal{C}}^\bullet(\vec{e})$
est encore donnée par la même formule que ci-dessus.

\medskip\noindent
Supposons que $NEG(\vec{e}) = \{0, \ldots, n\}$, c'est-à-dire que tous les
exposants $e_i$ soient négatifs.
Dans ce cas, le complexe $\check{\mathcal{C}}^\bullet(\vec{e})$ ne
possède qu'un terme, à savoir $\check{\mathcal{C}}^n(\vec{e}) =
R \cdot \frac{1}{T_0^{-e_0} \ldots T_n^{-e_n}}$. On a donc, dans ce
cas,
$$
H^q(\check{\mathcal{C}}^\bullet(\vec{e})) =
\left\{
\begin{matrix}
R \cdot \frac{1}{T_0^{-e_0} \ldots T_n^{-e_n}} & \text{si} & q = n \\
0 & \text{si} & \text{sinon}
\end{matrix}
\right.
$$
La somme directe de tous ces termes donne manifestement
$$
\left(\frac{1}{T_0 \ldots T_n} R[\frac{1}{T_0}, \ldots, \frac{1}{T_n}]\right)_d
$$
en degré $n$, comme l'indique le lemme. De plus, ces termes ne contribuent
pas à la cohomologie en d'autres degrés (conformément, là encore, à l'énoncé du
lemme).

\medskip\noindent
Supposons $NEG(\vec{e}) = \emptyset$. Dans ce cas, le complexe
$\check{\mathcal{C}}^\bullet(\vec{e})$ possède un facteur direct $R$ correspondant
à chaque $i_0 < \ldots < i_p$.
Comparons le complexe $\check{\mathcal{C}}^\bullet(\vec{e})$
à un autre complexe. Considérons le recouvrement ouvert affine
$$
\mathcal{V} : \Spec(R) = \bigcup\nolimits_{i \in \{0, \ldots, n\}} V_i
$$
où $V_i = \Spec(R)$ pour tout $i$. Considérons le complexe de
{\v C}ech alterné
$$
\check{\mathcal{C}}_{ord}^\bullet(\mathcal{V}, \mathcal{O}_{\Spec(R)})
$$
Le même raisonnement que ci-dessus montre qu'il calcule la cohomologie du
faisceau structural sur $\Spec(R)$. Ainsi,
$H^p(
\check{\mathcal{C}}_{ord}^\bullet(\mathcal{V}, \mathcal{O}_{\Spec(R)})
) = R$ si $p = 0$, et vaut $0$ lorsque $p > 0$.
Pour ces faits, voir le
lemme \ref{coherent-lemma-cech-cohomology-quasi-coherent-trivial} et sa
démonstration. Notons que
$\check{\mathcal{C}}_{ord}^\bullet(\mathcal{V}, \mathcal{O}_{\Spec(R)})$
possède lui aussi un facteur direct $R$ pour chaque $i_0 < \ldots < i_p$, et a exactement la même
différentielle que $\check{\mathcal{C}}^\bullet(\vec{e})$. Autrement dit, ces
complexes sont isomorphes et ont donc la même cohomologie. On
en déduit que
$$
H^q(\check{\mathcal{C}}^\bullet(\vec{e})) =
\left\{
\begin{matrix}
R \cdot T_0^{e_0} \ldots T_n^{e_n} & \text{si} & q = 0 \\
0 & \text{si} & \text{sinon}
\end{matrix}
\right.
$$
lorsque $NEG(\vec{e}) = \emptyset$.
La somme directe de tous ces termes donne manifestement
$$
(R[T_0, \ldots, T_n])_d
$$
en degré $0$, comme l'indique le lemme. De plus, ces termes ne contribuent
pas à la cohomologie en d'autres degrés (conformément, là encore, à l'énoncé du
lemme).

\medskip\noindent
Pour achever la démonstration du lemme, il reste à montrer que les complexes
$\check{\mathcal{C}}^\bullet(\vec{e})$ sont acycliques lorsque
$NEG(\vec{e})$ n'est ni vide ni égal à $\{0, \ldots, n\}$.
Choisissons un indice $i_{\text{fix}} \not \in NEG(\vec{e})$ (il en existe un).
Considérons le morphisme
$$
h :
\check{\mathcal{C}}^{p + 1}(\vec{e})
\to
\check{\mathcal{C}}^p(\vec{e})
$$
défini, pour $i_0 < \ldots < i_p$, par
$$
h(s)_{i_0 \ldots i_p} =
\left\{
\begin{matrix}
0 & \text{si} & p \not \in \{0, \ldots, n - 1\} \\
0 & \text{si} & i_{\text{fix}} \in \{i_0, \ldots, i_p\} \\
s_{i_{\text{fix}} i_0 \ldots i_p} & \text{si} & i_{\text{fix}} < i_0 \\
(-1)^a s_{i_0 \ldots i_{a - 1} i_{\text{fix}} i_a \ldots i_p} &
\text{si} & i_{a - 1} < i_{\text{fix}} < i_a \\
(-1)^{p + 1} s_{i_0 \ldots i_p i_{\text{fix}}} &
\text{si} & i_p < i_{\text{fix}}
\end{matrix}
\right.
$$
On comparera avec la démonstration du
lemme \ref{coherent-lemma-cech-cohomology-quasi-coherent-trivial}.
Cette définition a un sens car
$$
NEG(\vec{e}) \subset \{i_0, \ldots, i_p\}
\Leftrightarrow
NEG(\vec{e}) \subset \{i_{\text{fix}}, i_0, \ldots, i_p\}
$$
Le même calcul combinatoire\footnote{
Par exemple, supposons que $i_0 < \ldots < i_p$ soit tel que
$i_{\text{fix}} \not \in \{i_0, \ldots, i_p\}$ et que
$i_{a - 1} < i_{\text{fix}} < i_a$ pour un certain $1 \leq a \leq p$. Alors
\begin{align*}
&
(dh + hd)(s)_{i_0 \ldots i_p} \\
& =
\sum\nolimits_{j = 0}^p
(-1)^j
h(s)_{i_0 \ldots \hat i_j \ldots i_p}
+
(-1)^a d(s)_{i_0 \ldots i_{a - 1} i_{\text{fix}} i_a \ldots i_p}\\
& =
\sum\nolimits_{j = 0}^{a - 1}
(-1)^{j + a - 1}
s_{i_0 \ldots \hat i_j \ldots i_{a - 1} i_{\text{fix}} i_a \ldots i_p} +
\sum\nolimits_{j = a}^p
(-1)^{j + a}
s_{i_0 \ldots i_{a - 1} i_{\text{fix}} i_a \ldots \hat i_j \ldots i_p}
+ \\
&
\sum\nolimits_{j = 0}^{a - 1}
(-1)^{a + j}
s_{i_0 \ldots \hat i_j \ldots i_{a - 1} i_{\text{fix}} i_a \ldots i_p} +
(-1)^{2a} s_{i_0 \ldots i_p} +
\sum\nolimits_{j = a}^p
(-1)^{a + j + 1}
s_{i_0 \ldots i_{a - 1} i_{\text{fix}} i_a \ldots \hat i_j \ldots i_p}
\\
& =
s_{i_0 \ldots i_p}
\end{align*}
comme voulu. Les autres cas sont analogues.}
que dans la démonstration du lemme \ref{coherent-lemma-cech-cohomology-quasi-coherent-trivial}
montre que
$$
(hd + dh)(s)_{i_0 \ldots i_p}
=
s_{i_0 \ldots i_p}
$$
On voit donc que le morphisme identité du complexe
$\check{\mathcal{C}}^\bullet(\vec{e})$ est homotope à zéro,
ce qui entraîne son acyclicité.
\end{proof}

\noindent
Dans le lemme suivant, nous utiliserons l'accouplement de
$R$-modules libres
$$
R[T_0, \ldots, T_n]
\times
\frac{1}{T_0 \ldots T_n} R[\frac{1}{T_0}, \ldots, \frac{1}{T_n}]
\longrightarrow
R
$$
défini par la règle
$$
(f, g)
\longmapsto
\text{coefficient de }
\frac{1}{T_0 \ldots T_n}
\text{ dans }fg.
$$
Autrement dit, l'accouplement de l'élément de base $T_0^{e_0} \ldots T_n^{e_n}$ avec
l'élément de base $T_0^{d_0} \ldots T_n^{d_n}$ vaut $1$ si et seulement
si $e_i + d_i = -1$ pour tout $i$, et vaut zéro dans tous les autres
cas. Cet accouplement donne une identification
$$
\left(\frac{1}{T_0 \ldots T_n} R[\frac{1}{T_0}, \ldots, \frac{1}{T_n}]\right)_d
=
\Hom_R((R[T_0, \ldots, T_n])_{-n - 1 - d}, R)
$$
On peut donc reformuler le résultat du
lemme \ref{coherent-lemma-cohomology-projective-space-over-ring} sous la forme
\begin{equation}
\label{coherent-equation-identify}
H^q(\mathbf{P}^n, \mathcal{O}_{\mathbf{P}^n_R}(d)) =
\left\{
\begin{matrix}
(R[T_0, \ldots, T_n])_d & \text{si} & q = 0 \\
0 & \text{si} & q \not = 0, n \\
\Hom_R((R[T_0, \ldots, T_n])_{-n - 1 - d}, R)
& \text{si} & q = n
\end{matrix}
\right.
\end{equation}

\begin{lemma}
\label{coherent-lemma-identify-functorially}
Les identifications de l'Équation (\ref{coherent-equation-identify}) sont compatibles
au changement de base par les homomorphismes\ d'anneaux $R \to R'$.
De plus, pour tout $f \in R[T_0, \ldots, T_n]$ homogène
de degré $m$, le morphisme de multiplication par $f$
$$
\mathcal{O}_{\mathbf{P}^n_R}(d)
\longrightarrow
\mathcal{O}_{\mathbf{P}^n_R}(d + m)
$$
induit, via les identifications de
l'Équation (\ref{coherent-equation-identify}), le morphisme sur la cohomologie donné par la multiplication par
$f$ sur $H^0$ et par la transposée de la multiplication par $f$
$$
(R[T_0, \ldots, T_n])_{-n - 1 - (d + m)}
\longrightarrow
(R[T_0, \ldots, T_n])_{-n - 1 - d}
$$
sur $H^n$.
\end{lemma}

\begin{proof}
Soit $R \to R'$ un homomorphisme d'anneaux.
Soit $\mathcal{U}$ le recouvrement ouvert affine standard de $\mathbf{P}^n_R$,
et soit $\mathcal{U}'$ le recouvrement ouvert affine standard de
$\mathbf{P}^n_{R'}$. Notons que $\mathcal{U}'$ est l'image inverse du recouvrement
$\mathcal{U}$ par le morphisme canonique
$\mathbf{P}^n_{R'} \to \mathbf{P}^n_R$. Il existe donc
un morphisme de complexes de {\v C}ech
$$
\gamma :
\check{\mathcal{C}}_{ord}^\bullet(\mathcal{U},
\mathcal{O}_{\mathbf{P}_R}(d))
\longrightarrow
\check{\mathcal{C}}_{ord}^\bullet(\mathcal{U}',
\mathcal{O}_{\mathbf{P}_{R'}}(d))
$$
compatible avec le morphisme induit sur la cohomologie, d'après
Cohomologie, lemme \ref{cohomology-lemma-functoriality-cech}.
Les calculs de la démonstration du
lemme \ref{coherent-lemma-cohomology-projective-space-over-ring}
montrent que ce morphisme de complexes de {\v C}ech est compatible avec les identifications des groupes
de cohomologie en question. (En effet, les éléments de base du complexe
de {\v C}ech sur $R$ sont simplement envoyés sur les éléments de base correspondants
du complexe de {\v C}ech sur $R'$.) D'où la première assertion du lemme.

\medskip\noindent
Fixons maintenant l'anneau $R$ et considérons deux polynômes homogènes
$f, g \in R[T_0, \ldots, T_n]$ de même degré $m$.
La cohomologie étant un foncteur additif, il est clair que le morphisme
induit par la multiplication par $f + g$ est la somme des morphismes
induits par la multiplication par $f$ et
par $g$. De plus, la cohomologie étant un foncteur,
un résultat analogue vaut pour la multiplication par un produit $fg$, où
$f, g$ sont homogènes (mais pas nécessairement de même degré). Pour
vérifier la deuxième assertion du lemme, il suffit donc de
la prouver lorsque $f = x \in R$ ou lorsque $f = T_i$.
Dans le cas de la multiplication par un élément $x \in R$, le
résultat découle du fait que tous les groupes ou complexes de cohomologie considérés
possèdent une structure de $R$-module ou de complexe de $R$-modules.
Enfin, considérons la multiplication par $T_i$
comme morphisme $\mathcal{O}_{\mathbf{P}^n_R}$-linéaire
$$
\mathcal{O}_{\mathbf{P}^n_R}(d)
\longrightarrow
\mathcal{O}_{\mathbf{P}^n_R}(d + 1)
$$
L'assertion concernant $H^0$ est claire. Pour celle concernant $H^n$,
considérons la multiplication par $T_i$ comme morphisme de complexes de {\v C}ech
$$
\check{\mathcal{C}}_{ord}^\bullet(\mathcal{U},
\mathcal{O}_{\mathbf{P}_R}(d))
\longrightarrow
\check{\mathcal{C}}_{ord}^\bullet(\mathcal{U},
\mathcal{O}_{\mathbf{P}_R}(d + 1))
$$
Utilisons les notations introduites dans la démonstration du
lemme \ref{coherent-lemma-cohomology-projective-space-over-ring}.
Considérons l'effet de la multiplication par $T_i$
sur les décompositions
$$
\check{\mathcal{C}}_{ord}^\bullet(\mathcal{U}, \mathcal{O}_{\mathbf{P}_R}(d))
=
\bigoplus\nolimits_{\vec{e} \in \mathbf{Z}^{n + 1}, \ \sum e_i = d}
\check{\mathcal{C}}^\bullet(\vec{e})
$$
et
$$
\check{\mathcal{C}}_{ord}^\bullet(\mathcal{U},
\mathcal{O}_{\mathbf{P}_R}(d + 1))
=
\bigoplus\nolimits_{\vec{e} \in \mathbf{Z}^{n + 1}, \ \sum e_i = d + 1}
\check{\mathcal{C}}^\bullet(\vec{e})
$$
Il est clair qu'elle envoie le sous-complexe
$\check{\mathcal{C}}^\bullet(\vec{e})$ dans le sous-complexe
$\check{\mathcal{C}}^\bullet(\vec{e} + \vec{b}_i)$, où
$\vec{b}_i = (0, \ldots, 0, 1, 0, \ldots, 0))$ est le $i$-ième vecteur de
base. Autrement dit, elle envoie le facteur direct de $H^n$ correspondant à
$\vec{e}$, avec $e_i < 0$ et $\sum e_i = d$,
dans le facteur direct de $H^n$ correspondant à
$\vec{e} + \vec{b}_i$ (qui est nul si $e_i + b_i \geq 0$).
On voit aisément que cela correspond exactement à l'action de la
transposée de la multiplication par $T_i$, considérée comme morphisme
$$
(R[T_0, \ldots, T_n])_{-n - 1 - (d + 1)}
\longrightarrow
(R[T_0, \ldots, T_n])_{-n - 1 - d}
$$
Cela prouve le lemme.
\end{proof}

\noindent
Avant d'énoncer la version relative, introduisons quelques notations.
Rappelons que $\mathcal{O}_S[T_0, \ldots, T_n]$ est un
$\mathcal{O}_S$-module gradué dans lequel chaque $T_i$ est homogène de degré $1$.
Notons $(\mathcal{O}_S[T_0, \ldots, T_n])_d$ la composante homogène de degré $d$.
C'est un faisceau localement libre de type fini de rang $\binom{n + d}{d}$ sur $S$.

\begin{lemma}
\label{coherent-lemma-cohomology-projective-space-over-base}
Soit $S$ un schéma.
Soit $n \geq 0$ un entier. Considérons
le morphisme structural
$$
f : \mathbf{P}^n_S \longrightarrow S.
$$
On a
$$
R^qf_*(\mathcal{O}_{\mathbf{P}^n_S}(d)) =
\left\{
\begin{matrix}
(\mathcal{O}_S[T_0, \ldots, T_n])_d & \text{si} & q = 0 \\
0 & \text{si} & q \not = 0, n \\
\SheafHom_{\mathcal{O}_S}(
(\mathcal{O}_S[T_0, \ldots, T_n])_{- n - 1 - d}, \mathcal{O}_S)
& \text{si} & q = n
\end{matrix}
\right.
$$
\end{lemma}

\begin{proof}
Démonstration omise. Indication : cela résulte de la compatibilité des identifications
(\ref{coherent-equation-identify}) au changement de base affine, donnée par
le lemme \ref{coherent-lemma-identify-functorially}.
\end{proof}

\noindent
Énonçons maintenant la version pour les fibrés projectifs associés aux faisceaux localement libres de
type fini. Soit $S$ un schéma. Soit $\mathcal{E}$
un $\mathcal{O}_S$-module localement libre de type fini de rang constant $n + 1$, voir
Modules, section \ref{modules-section-locally-free}.
On considère alors $\text{Sym}(\mathcal{E})$ comme un
$\mathcal{O}_S$-module gradué dont $\mathcal{E}$ est la composante homogène de degré $1$.
Et $\text{Sym}^d(\mathcal{E})$ est la composante homogène de degré $d$.
C'est un faisceau localement libre de type fini de rang $\binom{n + d}{d}$ sur $S$.
Rappelons notre convention :
$$
\pi :
\mathbf{P}(\mathcal{E})
=
\underline{\text{Proj}}_S(\text{Sym}(\mathcal{E}))
\longrightarrow
S
$$
et l'existence de morphismes naturels
$\text{Sym}^d(\mathcal{E}) \to \pi_*\mathcal{O}_{\mathbf{P}(\mathcal{E})}(d)$.

\begin{lemma}
\label{coherent-lemma-cohomology-projective-bundle}
Soit $S$ un schéma. Soit $n \geq 1$.
Soit $\mathcal{E}$
un $\mathcal{O}_S$-module localement libre de type fini de rang constant $n + 1$.
Considérons le morphisme structural
$$
\pi : \mathbf{P}(\mathcal{E}) \longrightarrow S.
$$
On a
$$
R^q\pi_*(\mathcal{O}_{\mathbf{P}(\mathcal{E})}(d)) =
\left\{
\begin{matrix}
\text{Sym}^d(\mathcal{E}) & \text{si} & q = 0 \\
0 & \text{si} & q \not = 0, n \\
\SheafHom_{\mathcal{O}_S}(
\text{Sym}^{- n - 1 - d}(\mathcal{E})
\otimes_{\mathcal{O}_S}
\wedge^{n + 1}\mathcal{E},
\mathcal{O}_S)
& \text{si} & q = n
\end{matrix}
\right.
$$
Ces identifications sont compatibles au changement de base et
aux isomorphismes de faisceaux localement libres.
\end{lemma}

\begin{proof}
Considérons le morphisme canonique
$$
\pi^*\mathcal{E} \longrightarrow \mathcal{O}_{\mathbf{P}(\mathcal{E})}(1)
$$
et prenons son tordu de degré opposé à $1$ pour obtenir
$$
\pi^*(\mathcal{E})(-1) \longrightarrow \mathcal{O}_{\mathbf{P}(\mathcal{E})}
$$
C'est un morphisme surjectif d'un faisceau localement libre de rang $n + 1$ sur le
faisceau structural. Le complexe de Koszul correspondant est
donc exact (Compléments d'algèbre, Lemme
\ref{more-algebra-lemma-homotopy-koszul-abstract}).
Autrement dit, on dispose d'un complexe exact
$$
0 \to
\pi^*(\wedge^{n + 1}\mathcal{E})(-n - 1) \to
\ldots \to
\pi^*(\wedge^i\mathcal{E})(-i) \to
\ldots \to
\pi^*\mathcal{E}(-1) \to
\mathcal{O}_{\mathbf{P}(\mathcal{E})} \to 0
$$
Nous plaçons le terme $\pi^*(\wedge^i\mathcal{E})(-i)$ en
degré $-i$.
Calculons les images directes supérieures de ce complexe
acyclique à l'aide de la première suite spectrale de
Catégories dérivées, lemme \ref{derived-lemma-two-ss-complex-functor}.
On dispose ainsi d'une suite spectrale de termes
$$
E_1^{p, q} = R^q\pi_*\left(\pi^*(\wedge^{-p}\mathcal{E})(p)\right)
$$
convergeant vers zéro ! La formule de projection
(Cohomologie, lemme \ref{cohomology-lemma-projection-formula})
donne
$$
E_1^{p, q} = \wedge^{-p} \mathcal{E} \otimes_{\mathcal{O}_S}
R^q\pi_*\left(\mathcal{O}_{\mathbf{P}(\mathcal{E})}(p)\right).
$$
Notons que, localement sur $S$, le faisceau $\mathcal{E}$ est trivial,
c'est-à-dire isomorphe à $\mathcal{O}_S^{\oplus n + 1}$; localement sur
$S$, le morphisme $\mathbf{P}(\mathcal{E}) \to S$ s'identifie donc
à $\mathbf{P}^n_S \to S$. Ainsi,
localement sur $S$, on peut utiliser les résultats des Lemmes
\ref{coherent-lemma-cohomology-projective-space-over-ring},
\ref{coherent-lemma-identify-functorially} ou
\ref{coherent-lemma-cohomology-projective-space-over-base}.
Il en résulte que $E_1^{p, q} = 0$, sauf si $(p, q)$ vaut $(0, 0)$
ou $(-n - 1, n)$. Les termes non nuls sont
\begin{align*}
E_1^{0, 0} & = \pi_*\mathcal{O}_{\mathbf{P}(\mathcal{E})} = \mathcal{O}_S \\
E_1^{-n - 1, n} & =
R^n\pi_*\left(\pi^*(\wedge^{n + 1}\mathcal{E})(-n - 1)\right) =
\wedge^{n + 1}\mathcal{E} \otimes_{\mathcal{O}_S}
R^n\pi_*\left(\mathcal{O}_{\mathbf{P}(\mathcal{E})}(-n - 1)\right)
\end{align*}
Il ne peut donc y avoir qu'une seule différentielle
non nulle dans la suite spectrale, à savoir le morphisme
$d_{n + 1}^{-n - 1, n} : E_{n + 1}^{-n - 1, n} \to E_{n + 1}^{0, 0}$,
qui est nécessairement un isomorphisme (puisque la suite spectrale converge vers
le faisceau $0$). Ainsi $E_1^{p, q} = E_{n + 1}^{p, q}$,
et l'on obtient un isomorphisme canonique
$$
\wedge^{n + 1}\mathcal{E} \otimes_{\mathcal{O}_S}
R^n\pi_*\left(\mathcal{O}_{\mathbf{P}(\mathcal{E})}(-n - 1)\right) =
R^n\pi_*\left(\pi^*(\wedge^{n + 1}\mathcal{E})(-n - 1)\right)
\xrightarrow{d_{n + 1}^{-n - 1, n}}
\mathcal{O}_S
$$
Puisque $\wedge^{n + 1}\mathcal{E}$ est un faisceau
inversible, il en résulte que
$R^n\pi_*\mathcal{O}_{\mathbf{P}(\mathcal{E})}(-n - 1)$ est lui
aussi inversible et canoniquement isomorphe à l'inverse de
$\wedge^{n + 1}\mathcal{E}$. Nous avons donc prouvé le cas
$d = - n - 1$ du lemme.

\medskip\noindent
En travaillant localement sur $S$, on déduit immédiatement des calculs de
cohomologie des lemmes \ref{coherent-lemma-cohomology-projective-space-over-ring},
\ref{coherent-lemma-identify-functorially} ou
\ref{coherent-lemma-cohomology-projective-space-over-base} les assertions d'annulation
du lemme. De plus, le résultat concernant $R^0\pi_*$ est lui
aussi clair, puisqu'il existe des morphismes canoniques
$\text{Sym}^d(\mathcal{E}) \to \pi_* \mathcal{O}_{\mathbf{P}(\mathcal{E})}(d)$
pour tout $d$. Il reste à vérifier la description de
$R^n\pi_*\mathcal{O}_{\mathbf{P}(\mathcal{E})}(d)$
lorsque $d < -n - 1$. Pour cela, considérons le morphisme
$$
\pi^*(\text{Sym}^{-d - n - 1}(\mathcal{E}))
\otimes_{\mathcal{O}_{\mathbf{P}(\mathcal{E})}}
\mathcal{O}_{\mathbf{P}(\mathcal{E})}(d)
\longrightarrow
\mathcal{O}_{\mathbf{P}(\mathcal{E})}(-n - 1)
$$
En appliquant $R^n\pi_*$ et la formule de projection (voir ci-dessus), on obtient un morphisme
$$
\text{Sym}^{-d - n - 1}(\mathcal{E})
\otimes_{\mathcal{O}_S}
R^n\pi_*(\mathcal{O}_{\mathbf{P}(\mathcal{E})}(d))
\longrightarrow
R^n\pi_*\mathcal{O}_{\mathbf{P}(\mathcal{E})}(-n - 1) =
(\wedge^{n + 1}\mathcal{E})^{\otimes -1}
$$
(la dernière égalité a été établie
ci-dessus). Les calculs locaux des Lemmes
\ref{coherent-lemma-cohomology-projective-space-over-ring},
\ref{coherent-lemma-identify-functorially} ou
\ref{coherent-lemma-cohomology-projective-space-over-base}
montrent à nouveau que ce morphisme induit un accouplement parfait entre
$R^n\pi_*(\mathcal{O}_{\mathbf{P}(\mathcal{E})}(d))$ et
$\text{Sym}^{-d - n - 1}(\mathcal{E}) \otimes \wedge^{n + 1}(\mathcal{E})$,
comme voulu.
\end{proof}

















\section{Faisceaux cohérents sur les schémas localement noethériens}
\label{coherent-section-coherent-sheaves}

\noindent
Nous avons défini la notion de module cohérent sur un espace annelé quelconque dans
Modules, section \ref{modules-section-coherent}.
Bien que l'on puisse considérer des faisceaux cohérents sur des schémas
non noethériens, nous supposerons toujours le schéma de base localement noethérien lorsque
nous étudierons des faisceaux cohérents. Voici une caractérisation de ces faisceaux
sur les schémas localement noethériens.

\begin{lemma}
\label{coherent-lemma-coherent-Noetherian}
Soit $X$ un schéma localement noethérien.
Soit $\mathcal{F}$ un $\mathcal{O}_X$-module.
Les conditions suivantes sont équivalentes :
\begin{enumerate}
\item $\mathcal{F}$ est cohérent,
\item $\mathcal{F}$ est un $\mathcal{O}_X$-module quasi-cohérent de type fini,
\item $\mathcal{F}$ est un $\mathcal{O}_X$-module de présentation finie,
\item pour tout ouvert affine $\Spec(A) = U \subset X$, on a
$\mathcal{F}|_U = \widetilde M$, où $M$ est un $A$-module de type fini, et
\item il existe un recouvrement ouvert affine $X = \bigcup U_i$,
$U_i = \Spec(A_i)$, tel que, pour chaque indice,
$\mathcal{F}|_{U_i} = \widetilde M_i$, où $M_i$ est un $A_i$-module de type fini.
\end{enumerate}
En particulier, $\mathcal{O}_X$ est cohérent, tout
$\mathcal{O}_X$-module inversible est cohérent et, plus
généralement, tout $\mathcal{O}_X$-module localement libre de type fini est cohérent.
\end{lemma}

\begin{proof}
Les implications (1) $\Rightarrow$ (2) et (1) $\Rightarrow$ (3) sont
vraies en général, voir
Modules, lemme \ref{modules-lemma-coherent-finite-presentation}.
Si $\mathcal{F}$ est de présentation finie, alors $\mathcal{F}$ est
quasi-cohérent, voir
Modules, lemme \ref{modules-lemma-finite-presentation-quasi-coherent}. On
a donc aussi (3) $\Rightarrow$ (2).

\medskip\noindent
Supposons que $\mathcal{F}$ soit un $\mathcal{O}_X$-module quasi-cohérent
de type
fini. D'après Propriétés, lemme \ref{properties-lemma-finite-type-module},
sur tout ouvert affine
$\Spec(A) = U \subset X$, on a $\mathcal{F}|_U = \widetilde M$,
où $M$ est un $A$-module de type fini. Puisque $A$ est noethérien,
$M$ admet une présentation finie
$$
A^{\oplus m} \to A^{\oplus n} \to M \to 0.
$$
Ainsi $\mathcal{F}$ est de présentation finie d'après
Propriétés, lemme \ref{properties-lemma-finite-presentation-module}.
Autrement dit, (2) $\Rightarrow$ (3).

\medskip\noindent
D'après Modules, lemme \ref{modules-lemma-coherent-structure-sheaf}, il
suffit de montrer que $\mathcal{O}_X$ est cohérent pour prouver que (3) implique
(1). Il faut donc montrer que, pour tout ouvert $U \subset X$ et
toute famille finie de sections $f_i \in \mathcal{O}_X(U)$,
$i = 1, \ldots, n$, le noyau du morphisme
$\bigoplus_{i = 1, \ldots, n} \mathcal{O}_U \to \mathcal{O}_U$
est de type fini. Cette propriété étant locale, il suffit
de la vérifier au voisinage de tout $x \in U$.
On peut donc supposer $U = \Spec(A)$ affine. Dans ce cas,
$f_1, \ldots, f_n \in A$ sont des éléments de $A$. Puisque $A$ est
noethérien, voir
Propriétés, lemme \ref{properties-lemma-locally-Noetherian},
le noyau $K$ du morphisme $\bigoplus_{i = 1, \ldots, n} A \to A$
est un $A$-module de type fini. Voir par
exemple Algèbre, lemme \ref{algebra-lemma-Noetherian-basic}.
Le foncteur\ $\widetilde{ }$\ étant exact, voir
Schémas, lemme \ref{schemes-lemma-spec-sheaves},
on obtient une suite exacte
$$
\widetilde K \to
\bigoplus\nolimits_{i = 1, \ldots, n} \mathcal{O}_U \to
\mathcal{O}_U
$$
et, en appliquant à
nouveau Propriétés, lemme \ref{properties-lemma-finite-type-module},
on voit que $\widetilde K$ est de type fini. On en
déduit l'équivalence de (1), (2) et (3).

\medskip\noindent
Il résulte de
Propriétés, lemme \ref{properties-lemma-finite-type-module},
que (2) implique (4). Il est immédiat que (4) implique (5).
La discussion de
Schémas, section \ref{schemes-section-quasi-coherent},
montre que (5) implique la
quasi-cohérence de $\mathcal{F}$, et il est clair que (5)
implique que $\mathcal{F}$ soit de type fini. Ainsi (5) implique
(2), ce qui conclut.
\end{proof}

\begin{lemma}
\label{coherent-lemma-coherent-abelian-Noetherian}
Soit $X$ un schéma localement noethérien.
La catégorie des $\mathcal{O}_X$-modules cohérents est
abélienne. Plus précisément, le noyau et le conoyau d'un morphisme de
$\mathcal{O}_X$-modules cohérents sont cohérents. Toute extension
de faisceaux cohérents est cohérente.
\end{lemma}

\begin{proof}
C'est une reformulation de
Modules, lemme \ref{modules-lemma-coherent-abelian},
dans un cas particulier.
\end{proof}

\noindent
Le lemme suivant n'est pas toujours vrai pour la catégorie des
$\mathcal{O}_X$-modules cohérents sur un espace annelé quelconque $X$.

\begin{lemma}
\label{coherent-lemma-coherent-Noetherian-quasi-coherent-sub-quotient}
Soit $X$ un schéma localement noethérien.
Soit $\mathcal{F}$ un $\mathcal{O}_X$-module
cohérent. Tout sous-module quasi-cohérent de $\mathcal{F}$ est cohérent.
Tout module quotient quasi-cohérent de $\mathcal{F}$ est cohérent.
\end{lemma}

\begin{proof}
On peut supposer $X$ affine, disons $X = \Spec(A)$.
Propriétés, lemme \ref{properties-lemma-locally-Noetherian},
implique que $A$ est noethérien. Le lemme \ref{coherent-lemma-coherent-Noetherian}
ramène alors l'énoncé à l'algèbre. Le fait
algébrique correspondant est qu'un quotient ou un sous-module d'un $A$-module de
type fini est un $A$-module de type fini, voir par
exemple Algèbre, lemme \ref{algebra-lemma-Noetherian-basic}.
\end{proof}

\begin{lemma}
\label{coherent-lemma-tensor-hom-coherent}
Soit $X$ un schéma localement noethérien.
Soient $\mathcal{F}$, $\mathcal{G}$ des $\mathcal{O}_X$-modules cohérents.
Les $\mathcal{O}_X$-modules $\mathcal{F} \otimes_{\mathcal{O}_X} \mathcal{G}$
et $\SheafHom_{\mathcal{O}_X}(\mathcal{F}, \mathcal{G})$ sont
cohérents.
\end{lemma}

\begin{proof}
On montre dans
Modules, lemme \ref{modules-lemma-internal-hom-locally-kernel-direct-sum}, que
$\SheafHom_{\mathcal{O}_X}(\mathcal{F}, \mathcal{G})$ est cohérent. Le
résultat pour les produits tensoriels est donné
par Modules, lemme \ref{modules-lemma-tensor-product-permanence}.
\end{proof}

\begin{lemma}
\label{coherent-lemma-local-isomorphism}
Soit $X$ un schéma localement noethérien.
Soient $\mathcal{F}$, $\mathcal{G}$ des $\mathcal{O}_X$-modules cohérents.
Soit $\varphi : \mathcal{G} \to \mathcal{F}$ un homomorphisme
de $\mathcal{O}_X$-modules. Soit $x \in X$.
\begin{enumerate}
\item Si $\mathcal{F}_x = 0$, il existe un voisinage ouvert
$U \subset X$ de $x$ tel que $\mathcal{F}|_U = 0$.
\item Si $\varphi_x : \mathcal{G}_x \to \mathcal{F}_x$ est injectif,
il existe un voisinage ouvert $U \subset X$ de $x$ tel que
$\varphi|_U$ soit injectif.
\item Si $\varphi_x : \mathcal{G}_x \to \mathcal{F}_x$ est surjectif,
il existe un voisinage ouvert $U \subset X$ de $x$ tel que
$\varphi|_U$ soit surjectif.
\item Si $\varphi_x : \mathcal{G}_x \to \mathcal{F}_x$ est bijectif,
il existe un voisinage ouvert $U \subset X$ de $x$ tel que
$\varphi|_U$ soit un isomorphisme.
\end{enumerate}
\end{lemma}

\begin{proof}
Voir Modules, Lemmes
\ref{modules-lemma-finite-type-surjective-on-stalk},
\ref{modules-lemma-finite-type-stalk-zero} et
\ref{modules-lemma-finite-type-to-coherent-injective-on-stalk}.
\end{proof}

\begin{lemma}
\label{coherent-lemma-map-stalks-local-map}
Soit $X$ un schéma localement noethérien.
Soient $\mathcal{F}$, $\mathcal{G}$ des $\mathcal{O}_X$-modules cohérents.
Soit $x \in X$. Supposons
que $\psi : \mathcal{G}_x \to \mathcal{F}_x$ soit un morphisme de
$\mathcal{O}_{X, x}$-modules.
Il existe alors un voisinage ouvert $U \subset X$ de $x$ et un morphisme
$\varphi : \mathcal{G}|_U \to \mathcal{F}|_U$ tels que
$\varphi_x = \psi$.
\end{lemma}

\begin{proof}
Compte tenu du lemme \ref{coherent-lemma-coherent-Noetherian},
c'est une reformulation de
Modules, lemme \ref{modules-lemma-stalk-internal-hom}.
\end{proof}

\begin{lemma}
\label{coherent-lemma-coherent-support-closed}
Soit $X$ un schéma localement noethérien. Soit $\mathcal{F}$ un
$\mathcal{O}_X$-module cohérent. Alors $\text{Supp}(\mathcal{F})$ est fermé, et
$\mathcal{F}$ provient d'un faisceau cohérent sur le support schématique
de $\mathcal{F}$, voir
Morphismes, définition \ref{morphisms-definition-scheme-theoretic-support}.
\end{lemma}

\begin{proof}
Soit $i : Z \to X$ le support schématique de $\mathcal{F}$, et
soit $\mathcal{G}$ le faisceau quasi-cohérent de type fini sur $Z$
tel que $i_*\mathcal{G} \cong \mathcal{F}$.
Puisque $Z = \text{Supp}(\mathcal{F})$, le support est fermé.
Le schéma $Z$ est localement noethérien d'après
Morphismes, lemmes \ref{morphisms-lemma-immersion-locally-finite-type}
et \ref{morphisms-lemma-finite-type-noetherian}.
Enfin, $\mathcal{G}$ est un $\mathcal{O}_Z$-module cohérent d'après
le lemme \ref{coherent-lemma-coherent-Noetherian}.
\end{proof}

\begin{lemma}
\label{coherent-lemma-i-star-equivalence}
Soit $i : Z \to X$ une immersion fermée de schémas localement noethériens.
Soit $\mathcal{I} \subset \mathcal{O}_X$ le faisceau quasi-cohérent d'idéaux
définissant $Z$. Le foncteur $i_*$ induit une équivalence entre
la catégorie des $\mathcal{O}_X$-modules cohérents annulés par $\mathcal{I}$
et la catégorie des $\mathcal{O}_Z$-modules cohérents.
\end{lemma}

\begin{proof}
Le foncteur est pleinement fidèle d'après
Morphismes, lemme \ref{morphisms-lemma-i-star-equivalence}.
Soit $\mathcal{F}$ un $\mathcal{O}_X$-module
cohérent annulé par $\mathcal{I}$. D'après
Morphismes, lemme \ref{morphisms-lemma-i-star-equivalence},
on peut écrire $\mathcal{F} = i_*\mathcal{G}$,
où $\mathcal{G}$ est un faisceau quasi-cohérent sur $Z$. D'après
Modules, lemme \ref{modules-lemma-i-star-reflects-finite-type},
le faisceau $\mathcal{G}$ est de type fini.
Ainsi $\mathcal{G}$ est cohérent d'après le
lemme \ref{coherent-lemma-coherent-Noetherian}.
Le foncteur est donc aussi essentiellement surjectif, comme voulu.
\end{proof}

\begin{lemma}
\label{coherent-lemma-finite-pushforward-coherent}
Soit $f : X \to Y$ un morphisme de schémas.
Soit $\mathcal{F}$ un $\mathcal{O}_X$-module quasi-cohérent.
Supposons $f$ fini et $Y$ localement noethérien.
Alors $R^pf_*\mathcal{F} = 0$ pour $p > 0$, et
$f_*\mathcal{F}$ est cohérent si $\mathcal{F}$ est cohérent.
\end{lemma}

\begin{proof}
Les images directes supérieures sont nulles d'après le
lemme \ref{coherent-lemma-relative-affine-vanishing}, et
parce qu'un morphisme fini est affine (par
définition). Notons que les hypothèses impliquent que $X$ est lui aussi localement
noethérien (voir Morphismes, lemme \ref{morphisms-lemma-finite-type-noetherian}),
de sorte que l'énoncé a un sens.
Soit $\Spec(A) = V \subset Y$ un ouvert affine. D'après
Morphismes, définition \ref{morphisms-definition-integral},
on a $f^{-1}(V) = \Spec(B)$, avec $A \to B$ fini. Le
lemme \ref{coherent-lemma-coherent-Noetherian}
ramène l'énoncé au fait algébrique suivant :
si $M$ est un $B$-module de type fini, alors $M$ est aussi de
type fini comme $A$-module, voir
Algèbre, lemme \ref{algebra-lemma-finite-module-over-finite-extension}.
\end{proof}

\noindent
Dans la situation du lemme, les images directes supérieures sont elles
aussi cohérentes puisqu'elles sont
nulles. Nous montrerons qu'il en est toujours ainsi pour un morphisme propre
entre schémas localement noethériens
(proposition \ref{coherent-proposition-proper-pushforward-coherent}).

\begin{lemma}
\label{coherent-lemma-coherent-support-dimension-0}
Soit $X$ un schéma localement noethérien. Soit $\mathcal{F}$
un faisceau cohérent tel que $\dim(\text{Supp}(\mathcal{F})) \leq 0$.
Alors $\mathcal{F}$ est engendré par ses sections globales, et
$H^i(X, \mathcal{F}) = 0$ pour $i > 0$.
\end{lemma}

\begin{proof}
D'après le lemme \ref{coherent-lemma-coherent-support-closed}, on a
$\mathcal{F} = i_*\mathcal{G}$, où $i : Z \to X$ est l'inclusion
du support schématique de $\mathcal{F}$ et où $\mathcal{G}$
est un $\mathcal{O}_Z$-module cohérent. Puisque la dimension de $Z$ est
$0$, le schéma $Z$ est une réunion disjointe de schémas affines (Propriétés, Lemme
\ref{properties-lemma-locally-Noetherian-dimension-0}).
Ainsi $\mathcal{G}$ est engendré par ses sections globales et les groupes
de cohomologie supérieure de $\mathcal{G}$ sont nuls
(lemme \ref{coherent-lemma-quasi-coherent-affine-cohomology-zero}). Par
conséquent, $\mathcal{F} = i_*\mathcal{G}$ est engendré par ses
sections globales. Puisque les cohomologies de $\mathcal{F}$ et de $\mathcal{G}$ coïncident (le
lemme \ref{coherent-lemma-relative-affine-cohomology} s'applique, car une
immersion fermée est
affine), les groupes de cohomologie supérieure de $\mathcal{F}$ sont nuls.
\end{proof}

\begin{lemma}
\label{coherent-lemma-pushforward-coherent-on-open}
Soit $X$ un schéma. Soit $j : U \to X$ l'inclusion d'un ouvert.
Soit $T \subset X$ un sous-ensemble fermé contenu dans $U$.
Si $\mathcal{F}$ est un $\mathcal{O}_U$-module cohérent
tel que $\text{Supp}(\mathcal{F}) \subset T$, alors
$j_*\mathcal{F}$ est un $\mathcal{O}_X$-module cohérent.
\end{lemma}

\begin{proof}
Considérons le recouvrement ouvert $X = U \cup (X \setminus T)$.
Alors $j_*\mathcal{F}|_U = \mathcal{F}$ est cohérent, et
$j_*\mathcal{F}|_{X \setminus T} = 0$ est également cohérent.
Ainsi $j_*\mathcal{F}$ est cohérent.
\end{proof}














\section{Faisceaux cohérents sur les schémas noethériens}
\label{coherent-section-coherent-quasi-compact}

\noindent
Dans cette section, nous donnons quelques propriétés des faisceaux cohérents sur les
schémas noethériens.

\begin{lemma}
\label{coherent-lemma-acc-coherent}
Soit $X$ un schéma noethérien.
Soit $\mathcal{F}$ un $\mathcal{O}_X$-module
cohérent. Les sous-modules quasi-cohérents
de $\mathcal{F}$ vérifient la condition de chaîne ascendante. Autrement dit, pour toute suite
$$
\mathcal{F}_1 \subset \mathcal{F}_2 \subset \ldots \subset \mathcal{F}
$$
de sous-modules quasi-cohérents, on a
$\mathcal{F}_n = \mathcal{F}_{n + 1} = \ldots $ pour un certain $n \geq 0$.
\end{lemma}

\begin{proof}
Choisissons un recouvrement ouvert affine fini.
Sur chaque ouvert du recouvrement, la suite stationne d'après
Algèbre, lemme \ref{algebra-lemma-Noetherian-basic}.
Le lemme s'ensuit.
\end{proof}

\begin{lemma}
\label{coherent-lemma-power-ideal-kills-sheaf}
Soit $X$ un schéma noethérien.
Soit $\mathcal{F}$ un faisceau cohérent sur $X$.
Soit $\mathcal{I} \subset \mathcal{O}_X$ un faisceau
quasi-cohérent d'idéaux correspondant à un sous-schéma fermé $Z \subset X$.
Il existe un entier $n \geq 0$ tel que $\mathcal{I}^n\mathcal{F} = 0$
si et seulement si $\text{Supp}(\mathcal{F}) \subset Z$ ensemblistement.
\end{lemma}

\begin{proof}
Cela résulte immédiatement
d'Algèbre, lemme \ref{algebra-lemma-Noetherian-power-ideal-kills-module},
puisque $X$ admet un recouvrement fini par des spectres d'anneaux noethériens.
\end{proof}

\begin{lemma}[Artin-Rees]
\label{coherent-lemma-Artin-Rees}
Soit $X$ un schéma noethérien.
Soit $\mathcal{F}$ un faisceau cohérent sur $X$.
Soit $\mathcal{G} \subset \mathcal{F}$ un sous-faisceau quasi-cohérent.
Soit $\mathcal{I} \subset \mathcal{O}_X$ un faisceau quasi-cohérent d'idéaux.
Il
existe alors un entier $c \geq 0$ tel que, pour tout $n \geq c$, on
ait
$$
\mathcal{I}^{n - c}(\mathcal{I}^c\mathcal{F} \cap \mathcal{G})
=
\mathcal{I}^n\mathcal{F} \cap \mathcal{G}
$$
\end{lemma}

\begin{proof}
Cela résulte immédiatement
d'Algèbre, lemme \ref{algebra-lemma-Artin-Rees},
puisque $X$ admet un recouvrement fini par des spectres d'anneaux noethériens.
\end{proof}

\begin{lemma}
\label{coherent-lemma-directed-colimit-coherent}
Soit $X$ un schéma noethérien. Tout $\mathcal{O}_X$-module
quasi-cohérent est la limite inductive filtrante de ses sous-modules cohérents.
\end{lemma}

\begin{proof}
C'est une reformulation de Propriétés, Lemme
\ref{properties-lemma-quasi-coherent-colimit-finite-type},
compte tenu du fait qu'un
$\mathcal{O}_X$-module quasi-cohérent de type fini est cohérent d'après
le lemme \ref{coherent-lemma-coherent-Noetherian}.
\end{proof}

\begin{lemma}
\label{coherent-lemma-homs-over-open}
Soit $X$ un schéma noethérien.
Soit $\mathcal{F}$ un $\mathcal{O}_X$-module quasi-cohérent.
Soit $\mathcal{G}$ un $\mathcal{O}_X$-module cohérent.
Soit $\mathcal{I} \subset \mathcal{O}_X$ un faisceau quasi-cohérent d'idéaux.
Notons $Z \subset X$ le sous-schéma fermé correspondant et
posons $U = X \setminus Z$.
Il existe un isomorphisme canonique
$$
\colim_n \Hom_{\mathcal{O}_X}(\mathcal{I}^n\mathcal{G}, \mathcal{F})
\longrightarrow
\Hom_{\mathcal{O}_U}(\mathcal{G}|_U, \mathcal{F}|_U).
$$
En particulier, on a un isomorphisme
$$
\colim_n \Hom_{\mathcal{O}_X}(
\mathcal{I}^n, \mathcal{F})
\longrightarrow
\Gamma(U, \mathcal{F}).
$$
\end{lemma}

\begin{proof}
Montrons d'abord que le deuxième morphisme est un isomorphisme. Il est injectif d'après
Propriétés, lemme \ref{properties-lemma-sections-over-quasi-compact-open}.
Puisque $\mathcal{F}$ est la réunion de ses sous-modules cohérents, voir
Propriétés, lemme \ref{properties-lemma-quasi-coherent-colimit-finite-type}
(et lemme \ref{coherent-lemma-coherent-Noetherian}),
on peut supposer $\mathcal{F}$ cohérent pour prouver la surjectivité.
Notons $\mathcal{F}_n$ le sous-faisceau quasi-cohérent de $\mathcal{F}$
formé des sections annulées par $\mathcal{I}^n$,
voir Propriétés, lemme \ref{properties-lemma-sections-over-quasi-compact-open}.
Puisque $\mathcal{F}_1 \subset \mathcal{F}_2 \subset \ldots$, on a
$\mathcal{F}_n = \mathcal{F}_{n + 1} = \ldots $ pour un certain $n \geq 0$
d'après le lemme \ref{coherent-lemma-acc-coherent}. Posons $\mathcal{H} = \mathcal{F}_n$
pour cet $n$. D'après Artin-Rees (lemme \ref{coherent-lemma-Artin-Rees}),
il existe un entier $c \geq 0$ tel que
$\mathcal{I}^m\mathcal{F} \cap \mathcal{H}
\subset \mathcal{I}^{m - c}\mathcal{H}$. En prenant $m = n + c$, on obtient
$\mathcal{I}^m\mathcal{F} \cap \mathcal{H} \subset \mathcal{I}^n\mathcal{H}
= 0$. Si l'on pose $\mathcal{F}' = \mathcal{I}^m\mathcal{F}$, on
voit donc que $\mathcal{F}' \cap \mathcal{F}_n = 0$ et
$\mathcal{F}'|_U = \mathcal{F}|_U$. En particulier, le sous-faisceau
$(\mathcal{F}')_N$ des sections annulées par $\mathcal{I}^N$ est nul
pour tout $N \geq 0$. D'après
Propriétés, lemme \ref{properties-lemma-sections-over-quasi-compact-open},
on en
déduit que la flèche horizontale supérieure du diagramme commutatif
suivant est une bijection :
$$
\xymatrix{
\colim_n \Hom_{\mathcal{O}_X}(
\mathcal{I}^n, \mathcal{F}')
\ar[r] \ar[d] &
\Gamma(U, \mathcal{F}') \ar[d] \\
\colim_n \Hom_{\mathcal{O}_X}(
\mathcal{I}^n, \mathcal{F})
\ar[r] &
\Gamma(U, \mathcal{F})
}
$$
La flèche verticale de droite étant elle aussi bijective, on en déduit que
la flèche horizontale inférieure est surjective, comme voulu.

\medskip\noindent
Montrons ensuite que la première flèche du lemme est une bijection. D'après
le lemme \ref{coherent-lemma-coherent-Noetherian}, le faisceau $\mathcal{G}$
est de présentation finie, donc le faisceau
$\mathcal{H} = \SheafHom_{\mathcal{O}_X}(\mathcal{G}, \mathcal{F})$
est quasi-cohérent, voir
Schémas, section \ref{schemes-section-quasi-coherent}.
Par définition, on a
$$
\mathcal{H}(U)
=
\Hom_{\mathcal{O}_U}(\mathcal{G}|_U, \mathcal{F}|_U)
$$
Choisissons un élément $\psi$ dans le but de la première flèche du
lemme, c'est-à-dire $\psi \in \mathcal{H}(U)$. Le résultat que nous venons de prouver s'applique
à $\mathcal{H}$; il existe donc un entier $n \geq 0$ et un morphisme
$\varphi : \mathcal{I}^n \to \mathcal{H}$ qui redonne
$\psi$ par restriction à $U$. D'après
Modules, lemme \ref{modules-lemma-internal-hom},
$\varphi$ correspond à un morphisme
$$
\varphi :
\mathcal{I}^n \otimes_{\mathcal{O}_X} \mathcal{G}
\longrightarrow
\mathcal{F}.
$$
C'est presque le morphisme cherché, à ceci près que sa source est
le produit tensoriel de $\mathcal{I}^n$ et de $\mathcal{G}$
et non leur produit. Nous allons montrer que, quitte à augmenter $n$,
cette différence disparaît. Considérons la suite exacte courte
$$
0 \to \mathcal{K} \to
\mathcal{I}^n \otimes_{\mathcal{O}_X} \mathcal{G} \to
\mathcal{I}^n\mathcal{G} \to 0
$$
où $\mathcal{K}$ est le noyau. Notons que
$\mathcal{I}^n\mathcal{K} = 0$ (démonstration omise). En appliquant
à nouveau Artin-Rees, on obtient
$$
\mathcal{K}
\cap
\mathcal{I}^m(\mathcal{I}^n \otimes_{\mathcal{O}_X} \mathcal{G})
=
0
$$
pour un certain $m$ suffisamment grand. Autrement dit, le morphisme
$$
\mathcal{I}^m(\mathcal{I}^n \otimes_{\mathcal{O}_X} \mathcal{G})
\longrightarrow
\mathcal{I}^{n + m}\mathcal{G}
$$
est un isomorphisme. Soit $\varphi'$ la restriction de
$\varphi$ à ce sous-module, considérée comme morphisme
$\mathcal{I}^{m + n}\mathcal{G} \to \mathcal{F}$.
Alors $\varphi'$ définit un élément
de la source de la première flèche du lemme, dont l'image par cette
flèche est $\psi$. Nous avons donc prouvé la surjectivité de la
flèche. Nous omettons la démonstration de son injectivité.
\end{proof}

\begin{lemma}
\label{coherent-lemma-extend-coherent}
Soit $X$ un schéma localement noethérien. Soient $\mathcal{F}$, $\mathcal{G}$
des $\mathcal{O}_X$-modules cohérents. Soit $U \subset X$ un ouvert, et
soit $\varphi : \mathcal{F}|_U \to \mathcal{G}|_U$ un morphisme de
$\mathcal{O}_U$-modules. Il existe alors un sous-module
cohérent $\mathcal{F}' \subset \mathcal{F}$ coïncidant avec
$\mathcal{F}$ sur $U$, tel que $\varphi$ se prolonge en
$\varphi' : \mathcal{F}' \to \mathcal{G}$.
\end{lemma}

\begin{proof}
Soit $\mathcal{I} \subset \mathcal{O}_X$ le faisceau cohérent d'idéaux
définissant la structure réduite induite sur $X \setminus U$.
Si $X$ est noethérien, le lemme \ref{coherent-lemma-homs-over-open} montre que
l'on peut prendre $\mathcal{F}' = \mathcal{I}^n\mathcal{F}$
pour un certain $n$. Le cas général s'en déduit au moyen du lemme de Zorn.

\medskip\noindent
Considérons l'ensemble des triplets $(U', \mathcal{F}', \varphi')$ tels que
$U \subset U' \subset X$ soit ouvert, que $\mathcal{F}' \subset \mathcal{F}|_{U'}$
soit un sous-faisceau cohérent coïncidant avec $\mathcal{F}$ sur $U$, et que
$\varphi' : \mathcal{F}' \to \mathcal{G}|_{U'}$ se restreigne à
$\varphi$ sur $U$. Posons
$(U'', \mathcal{F}'', \varphi'') \geq (U', \mathcal{F}', \varphi')$
si et seulement si $U'' \supset U'$, $\mathcal{F}''|_{U'} = \mathcal{F}'$
et $\varphi''|_{U'} = \varphi'$.
Il est clair que, pour une famille totalement ordonnée
de triplets $(U_i, \mathcal{F}_i, \varphi_i)$, on
peut recoller les $\mathcal{F}_i$ en un sous-faisceau $\mathcal{F}'$ de $\mathcal{F}$
sur $U' = \bigcup U_i$, et prolonger $\varphi$ en un morphisme
$\varphi' : \mathcal{F}' \to \mathcal{G}|_{U'}$.
Tout sous-ensemble totalement ordonné de triplets admet donc un
majorant. Enfin, soit $(U', \mathcal{F}', \varphi')$
un triplet tel que $U' \not = X$. On peut alors choisir un
ouvert affine $W \subset X$ qui n'est pas contenu dans $U'$.
D'après le résultat du premier paragraphe, on peut prolonger
le sous-faisceau $\mathcal{F}'|_{W \cap U'}$ et la restriction
$\varphi'|_{W \cap U'}$ en un sous-faisceau $\mathcal{F}'' \subset \mathcal{F}|_W$
et en un morphisme $\varphi'' : \mathcal{F}'' \to \mathcal{G}|_W$.
La coïncidence de $(\mathcal{F}', \varphi')$ et de
$(\mathcal{F}'', \varphi'')$ sur $W \cap U'$ signifie précisément que
l'on peut prolonger le triplet en $(U' \cup W, \mathcal{F}''', \varphi''')$.
Ainsi, tout triplet maximal $(U', \mathcal{F}', \varphi')$
(dont l'existence résulte du lemme de Zorn) vérifie $U' = X$, ce
qui achève la démonstration.
\end{proof}









\section{Profondeur}
\label{coherent-section-depth}

\noindent
Dans cette section, nous étudions brièvement la profondeur et la propriété
$(S_k)$ pour les modules cohérents sur les schémas localement
noethériens. Rappelons que cette notion a déjà été abordée pour les
schémas localement noethériens dans Propriétés, section \ref{properties-section-Rk}.

\begin{definition}
\label{coherent-definition-depth}
Soit $X$ un schéma localement noethérien.
Soit $\mathcal{F}$ un $\mathcal{O}_X$-module cohérent.
Soit $k \geq 0$ un entier.
\begin{enumerate}
\item On dit que $\mathcal{F}$ est de {\it profondeur $k$ en un point}
$x$ de $X$ si $\text{depth}_{\mathcal{O}_{X, x}}(\mathcal{F}_x) = k$.
\item On dit que $X$ est de {\it profondeur $k$ en un point} $x$ de $X$ si
$\text{depth}(\mathcal{O}_{X, x}) = k$.
\item On dit que $\mathcal{F}$ vérifie la propriété {\it $(S_k)$} si
$$
\text{depth}_{\mathcal{O}_{X, x}}(\mathcal{F}_x)
\geq \min(k, \dim(\text{Supp}(\mathcal{F}_x)))
$$
pour tout $x \in X$.
\item On dit que $X$ vérifie la propriété {\it $(S_k)$} si $\mathcal{O}_X$ vérifie la
propriété $(S_k)$.
\end{enumerate}
\end{definition}

\noindent
Tout faisceau cohérent vérifie la condition $(S_0)$. La
condition $(S_1)$ équivaut à l'absence de points associés
immergés, voir Diviseurs, lemme \ref{divisors-lemma-S1-no-embedded}.

\begin{lemma}
\label{coherent-lemma-hom-into-depth}
Soit $X$ un schéma localement noethérien. Soient $\mathcal{F}$, $\mathcal{G}$
des $\mathcal{O}_X$-modules cohérents et $x \in X$.
\begin{enumerate}
\item Si $\mathcal{G}_x$ est de profondeur $\geq 1$, alors
$\SheafHom_{\mathcal{O}_X}(\mathcal{F}, \mathcal{G})_x$
est de profondeur $\geq 1$.
\item Si $\mathcal{G}_x$ est de profondeur $\geq 2$, alors
$\Hom_{\mathcal{O}_X}(\mathcal{F}, \mathcal{G})_x$ est de profondeur $\geq 2$.
\end{enumerate}
\end{lemma}

\begin{proof}
Remarquons que $\SheafHom_{\mathcal{O}_X}(\mathcal{F}, \mathcal{G})$ est
un $\mathcal{O}_X$-module cohérent d'après le lemme \ref{coherent-lemma-tensor-hom-coherent}.
Les modules cohérents sont de présentation finie
(lemme \ref{coherent-lemma-coherent-Noetherian}); le passage aux fibres commute donc à
la formation de $\SheafHom$ et de $\Hom$, voir
Modules, lemme \ref{modules-lemma-stalk-internal-hom}.
On se ramène ainsi au cas des modules de type fini sur
les anneaux locaux, traité dans Compléments d'algèbre, lemme \ref{more-algebra-lemma-hom-into-depth}.
\end{proof}

\begin{lemma}
\label{coherent-lemma-hom-into-S2}
Soit $X$ un schéma localement noethérien. Soient $\mathcal{F}$, $\mathcal{G}$
des $\mathcal{O}_X$-modules cohérents.
\begin{enumerate}
\item Si $\mathcal{G}$ vérifie la propriété $(S_1)$, alors
$\SheafHom_{\mathcal{O}_X}(\mathcal{F}, \mathcal{G})$ vérifie la propriété $(S_1)$.
\item Si $\mathcal{G}$ vérifie la propriété $(S_2)$, alors
$\SheafHom_{\mathcal{O}_X}(\mathcal{F}, \mathcal{G})$ vérifie la propriété $(S_2)$.
\end{enumerate}
\end{lemma}

\begin{proof}
Cela résulte immédiatement du lemme \ref{coherent-lemma-hom-into-depth}
et des définitions.
\end{proof}

\noindent
Nous avons vu dans Propriétés, lemme \ref{properties-lemma-scheme-CM-iff-all-Sk},
qu'un schéma localement noethérien
est de Cohen-Macaulay si et seulement s'il vérifie $(S_k)$ pour tout $k$.
Il est donc naturel d'introduire la définition suivante, qui équivaut
à demander que toutes les fibres soient des modules de Cohen-Macaulay.

\begin{definition}
\label{coherent-definition-Cohen-Macaulay}
Soit $X$ un schéma localement noethérien.
Soit $\mathcal{F}$ un $\mathcal{O}_X$-module cohérent.
On dit que $\mathcal{F}$ est {\it de Cohen-Macaulay} si et seulement s'il
vérifie $(S_k)$ pour tout $k \geq 0$.
\end{definition}

\begin{lemma}
\label{coherent-lemma-Cohen-Macaulay-over-regular}
Soit $X$ un schéma régulier. Soit $\mathcal{F}$ un
$\mathcal{O}_X$-module cohérent. Les conditions suivantes sont équivalentes :
\begin{enumerate}
\item $\mathcal{F}$ est de Cohen-Macaulay et $\text{Supp}(\mathcal{F}) = X$,
\item $\mathcal{F}$ est localement libre de type fini de rang $> 0$.
\end{enumerate}
\end{lemma}

\begin{proof}
Soit $x \in X$. Si (2) est vérifiée, alors $\mathcal{F}_x$ est un
$\mathcal{O}_{X, x}$-module libre de rang $> 0$. Ainsi
$\text{depth}(\mathcal{F}_x) = \dim(\mathcal{O}_{X, x})$,
car un anneau local régulier est de Cohen-Macaulay
(Algèbre, lemme \ref{algebra-lemma-regular-ring-CM}).
Réciproquement, si (1) est vérifiée, alors $\mathcal{F}_x$ est un
module de Cohen-Macaulay maximal sur $\mathcal{O}_{X, x}$
(Algèbre, définition \ref{algebra-definition-maximal-CM}).
Ainsi $\mathcal{F}_x$ est libre d'après
Algèbre, lemme \ref{algebra-lemma-regular-mcm-free}.
\end{proof}









\section{Dévissage des faisceaux cohérents}
\label{coherent-section-devissage}

\noindent
Soit $X$ un schéma noethérien. Considérons un sous-schéma fermé intègre
$i : Z \to X$. Il est souvent commode de considérer les faisceaux cohérents de
la forme $i_*\mathcal{G}$, où $\mathcal{G}$ est un faisceau cohérent sur
$Z$. Nous nous intéressons en particulier au cas où $\mathcal{G}$
est sans torsion et de rang $1$. Par exemple, $\mathcal{G}$ peut être
un faisceau d'idéaux non nul sur $Z$, ou, plus particulièrement,
$\mathcal{G} = \mathcal{O}_Z$.

\medskip\noindent
Tout au long de cette section, nous utiliserons le fait qu'un faisceau
cohérent est la même chose qu'un faisceau quasi-cohérent de type fini, et
qu'un sous-quotient quasi-cohérent d'un faisceau cohérent est cohérent, voir
section \ref{coherent-section-coherent-sheaves}.
Le support d'un faisceau cohérent est fermé, voir
Modules, lemme \ref{modules-lemma-support-finite-type-closed}.

\begin{lemma}
\label{coherent-lemma-prepare-filter-support}
Soit $X$ un schéma noethérien.
Soit $\mathcal{F}$ un faisceau cohérent sur $X$.
Supposons $\text{Supp}(\mathcal{F}) = Z \cup Z'$, avec $Z$, $Z'$ fermés.
Il existe alors une suite exacte courte de faisceaux cohérents
$$
0 \to \mathcal{G}' \to \mathcal{F} \to \mathcal{G} \to 0
$$
avec $\text{Supp}(\mathcal{G}') \subset Z'$ et
$\text{Supp}(\mathcal{G}) \subset Z$.
\end{lemma}

\begin{proof}
Soit $\mathcal{I} \subset \mathcal{O}_X$ le faisceau d'idéaux définissant la
structure réduite induite de sous-schéma fermé sur $Z$, voir
Schémas, lemme \ref{schemes-lemma-reduced-closed-subscheme}.
Considérons les sous-faisceaux
$\mathcal{G}'_n = \mathcal{I}^n\mathcal{F}$ et les
quotients $\mathcal{G}_n = \mathcal{F}/\mathcal{I}^n\mathcal{F}$.
Pour tout $n$, on a une suite exacte courte
$$
0 \to \mathcal{G}'_n \to \mathcal{F} \to \mathcal{G}_n \to 0
$$
Pour tout point $x$ de $Z' \setminus Z$, on a
$\mathcal{I}_x = \mathcal{O}_{X, x}$
et donc $\mathcal{G}_{n, x} = 0$. Ainsi
$\text{Supp}(\mathcal{G}_n) \subset Z$. Notons que $X \setminus Z'$
est un schéma noethérien. D'après le lemme \ref{coherent-lemma-power-ideal-kills-sheaf},
il existe donc un entier $n$ tel que
$\mathcal{G}'_n|_{X \setminus Z'} =
\mathcal{I}^n\mathcal{F}|_{X \setminus Z'} = 0$.
Pour un tel $n$, on a $\text{Supp}(\mathcal{G}'_n) \subset Z'$. Il
suffit donc de poser
$\mathcal{G}' = \mathcal{G}'_n$ et $\mathcal{G} = \mathcal{G}_n$
pour conclure.
\end{proof}

\begin{lemma}
\label{coherent-lemma-prepare-filter-irreducible}
Soit $X$ un schéma noethérien.
Soit $i : Z \to X$ un sous-schéma fermé intègre.
Soit $\xi \in Z$ son point générique.
Soit $\mathcal{F}$ un faisceau cohérent sur $X$.
Supposons que $\mathcal{F}_\xi$ soit annulé par
$\mathfrak m_\xi$. Il existe alors un entier
$r \geq 0$, un faisceau cohérent d'idéaux $\mathcal{I} \subset \mathcal{O}_Z$
et un morphisme injectif de faisceaux cohérents
$$
i_*\left(\mathcal{I}^{\oplus r}\right) \to \mathcal{F}
$$
qui est un isomorphisme au voisinage de $\xi$.
\end{lemma}

\begin{proof}
Soit $\mathcal{J} \subset \mathcal{O}_X$ le faisceau d'idéaux de $Z$.
Soit $\mathcal{F}' \subset \mathcal{F}$ le sous-faisceau des
sections locales de $\mathcal{F}$ annulées par
$\mathcal{J}$. C'est un faisceau quasi-cohérent d'après
Propriétés, lemme \ref{properties-lemma-sections-annihilated-by-ideal}. De
plus, $\mathcal{F}'_\xi = \mathcal{F}_\xi$, car
$\mathcal{J}_\xi = \mathfrak m_\xi$, et d'après l'assertion (3)
de Propriétés, lemme \ref{properties-lemma-sections-annihilated-by-ideal}.
Le lemme \ref{coherent-lemma-local-isomorphism} montre que
$\mathcal{F}' \to \mathcal{F}$
induit un isomorphisme au voisinage de $\xi$.
On peut donc remplacer $\mathcal{F}$ par $\mathcal{F}'$ et supposer
que $\mathcal{F}$ est annulé par $\mathcal{J}$.

\medskip\noindent
Supposons $\mathcal{J}\mathcal{F} = 0$. D'après le
lemme \ref{coherent-lemma-i-star-equivalence}, on peut écrire
$\mathcal{F} = i_*\mathcal{G}$,
où $\mathcal{G}$ est un faisceau cohérent sur $Z$. Supposons que l'on puisse trouver un morphisme
$\mathcal{I}^{\oplus r} \to \mathcal{G}$ qui soit un
isomorphisme au voisinage du point générique $\xi$ de $Z$.
En appliquant $i_*$ (qui est exact à gauche), on obtient alors le résultat du
lemme. Nous sommes donc ramenés au cas $X = Z$.

\medskip\noindent
Supposons que $Z = X$ soit un schéma noethérien intègre de point générique $\xi$.
Dans ce cas, $\mathcal{O}_{X, \xi} = \kappa(\xi)$ est, pour $X$, le corps
des fonctions.
Puisque $\mathcal{F}_\xi$ est un $\mathcal{O}_\xi$-module de type fini,
l'entier $r = \dim_{\kappa(\xi)} \mathcal{F}_\xi$ est fini.
Les faisceaux $\mathcal{O}_X^{\oplus r}$ et $\mathcal{F}$
ont donc des fibres isomorphes en $\xi$.
D'après le lemme \ref{coherent-lemma-map-stalks-local-map}, il existe un ouvert non
vide $U \subset X$ et un morphisme
$\psi : \mathcal{O}_X^{\oplus r}|_U \to \mathcal{F}|_U$
qui est un isomorphisme
en $\xi$, et donc au voisinage de $\xi$, d'après
le lemme \ref{coherent-lemma-local-isomorphism}.
D'après Schémas, lemme \ref{schemes-lemma-reduced-closed-subscheme},
il existe un faisceau quasi-cohérent d'idéaux
$\mathcal{I} \subset \mathcal{O}_X$
dont le sous-schéma fermé associé $Y \subset X$ est le complémentaire
de $U$.
D'après le lemme \ref{coherent-lemma-homs-over-open}, il existe un entier $n \geq 0$ et un morphisme
$\mathcal{I}^n(\mathcal{O}_X^{\oplus r}) \to \mathcal{F}$
qui redonne $\psi$ sur $U$. Puisque
$\mathcal{I}^n(\mathcal{O}_X^{\oplus r}) = (\mathcal{I}^n)^{\oplus r}$,
on obtient un morphisme comme dans le lemme. Il est injectif, car $X$ est
intègre et il est injectif au point générique de $X$
(démonstration facile omise).
\end{proof}

\begin{lemma}
\label{coherent-lemma-coherent-filter}
Soit $X$ un schéma noethérien.
Soit $\mathcal{F}$ un faisceau cohérent sur $X$.
Il existe une filtration
$$
0 = \mathcal{F}_0 \subset \mathcal{F}_1 \subset
\ldots \subset \mathcal{F}_m = \mathcal{F}
$$
par des sous-faisceaux cohérents telle que, pour chaque $j = 1, \ldots, m$,
il existe un sous-schéma fermé intègre $Z_j \subset X$ et
un faisceau cohérent d'idéaux non nul $\mathcal{I}_j \subset \mathcal{O}_{Z_j}$
tels que
$$
\mathcal{F}_j/\mathcal{F}_{j - 1}
\cong (Z_j \to X)_* \mathcal{I}_j
$$
\end{lemma}

\begin{proof}
Considérons l'ensemble
$$
\mathcal{T} =
\left\{
\begin{matrix}
Z \subset X
\text{ fermé tel qu'il existe un faisceau cohérent }
\mathcal{F} \\
\text{ avec }
\text{Supp}(\mathcal{F}) = Z
\text{ pour lequel le lemme est faux}
\end{matrix}
\right\}
$$
Nous voulons montrer que $\mathcal{T}$ est vide. Sinon,
puisque $X$ est noethérien, on peut choisir un élément minimal
$Z \in \mathcal{T}$. Il existe alors un faisceau
cohérent $\mathcal{F}$ sur $X$, de support $Z$, pour lequel le lemme
est faux. Il est clair que $Z \not = \emptyset$, car le
seul faisceau de support vide est le faisceau nul, pour lequel le
lemme est vrai (avec $m = 0$).

\medskip\noindent
Si $Z$ n'est pas irréductible, on peut écrire $Z = Z_1 \cup Z_2$,
où $Z_1, Z_2$ sont fermés et strictement contenus dans $Z$.
On peut alors appliquer le lemme \ref{coherent-lemma-prepare-filter-support}
pour obtenir une suite exacte courte de faisceaux cohérents
$$
0 \to
\mathcal{G}_1 \to
\mathcal{F} \to
\mathcal{G}_2 \to 0
$$
avec $\text{Supp}(\mathcal{G}_i) \subset Z_i$. Par minimalité de
$Z$, chacun des $\mathcal{G}_i$ admet une filtration comme dans l'énoncé
du lemme. En considérant la filtration induite sur $\mathcal{F}$,
on aboutit à une contradiction. On en déduit
que $Z$ est irréductible.

\medskip\noindent
Supposons $Z$ irréductible. Soit $\mathcal{J}$ le faisceau d'idéaux définissant
la structure réduite induite de sous-schéma fermé sur $Z$,
voir Schémas, lemme \ref{schemes-lemma-reduced-closed-subscheme}.
D'après le lemme \ref{coherent-lemma-power-ideal-kills-sheaf}, il existe un
entier $n \geq 0$ tel que $\mathcal{J}^n\mathcal{F} = 0$. On obtient donc
une filtration
$$
0 = \mathcal{J}^n\mathcal{F} \subset \mathcal{J}^{n - 1}\mathcal{F}
\subset \ldots \subset \mathcal{J}\mathcal{F} \subset \mathcal{F}
$$
dont chaque sous-quotient successif est annulé par $\mathcal{J}$.
Si chacun de ces sous-quotients admet une filtration comme dans l'énoncé du lemme, il
en va de même pour $\mathcal{F}$. Autrement dit, on peut
supposer que $\mathcal{J}$ annule $\mathcal{F}$.

\medskip\noindent
Lorsque $Z$ est irréductible et $\mathcal{J}\mathcal{F} = 0$,
on peut appliquer le lemme \ref{coherent-lemma-prepare-filter-irreducible}.
On obtient ainsi une suite exacte courte
$$
0 \to
i_*(\mathcal{I}^{\oplus r}) \to
\mathcal{F} \to
\mathcal{Q} \to 0
$$
où $\mathcal{Q}$ est le quotient.
Puisque $\mathcal{Q}$ est nul au voisinage de $\xi$,
d'après le lemme qui vient d'être cité, le support de $\mathcal{Q}$
est strictement contenu dans $Z$. Ainsi $\mathcal{Q}$
admet une filtration du type voulu par minimalité de $Z$. Mais
il en va alors de même pour $\mathcal{F}$, d'où la contradiction finale.
\end{proof}

\begin{lemma}
\label{coherent-lemma-property-initial}
Soit $X$ un schéma noethérien.
Soit $\mathcal{P}$ une propriété des faisceaux cohérents sur $X$. Supposons que
\begin{enumerate}
\item pour toute suite exacte courte de faisceaux cohérents
$$
0 \to \mathcal{F}_1 \to \mathcal{F} \to \mathcal{F}_2 \to 0
$$
si les $\mathcal{F}_i$, $i = 1, 2$, vérifient la propriété $\mathcal{P}$,
il en va de même pour $\mathcal{F}$;
\item pour tout sous-schéma fermé intègre $Z \subset X$
et tout faisceau quasi-cohérent d'idéaux
$\mathcal{I} \subset \mathcal{O}_Z$, on ait
$\mathcal{P}$ pour $i_*\mathcal{I}$.
\end{enumerate}
Alors la propriété $\mathcal{P}$ est vérifiée par tout faisceau cohérent
sur $X$.
\end{lemma}

\begin{proof}
Notons d'abord que, si $\mathcal{F}$ est un faisceau cohérent muni d'une filtration
$$
0 = \mathcal{F}_0 \subset \mathcal{F}_1 \subset
\ldots \subset \mathcal{F}_m = \mathcal{F}
$$
par des sous-faisceaux cohérents telle que chaque $\mathcal{F}_i/\mathcal{F}_{i - 1}$
vérifie la propriété $\mathcal{P}$, il en va de même pour $\mathcal{F}$.
Cela résulte de la condition (1) sur $\mathcal{P}$.
D'autre part, le lemme \ref{coherent-lemma-coherent-filter}
permet de filtrer tout $\mathcal{F}$
avec des sous-quotients successifs comme dans
(2). Le lemme s'ensuit.
\end{proof}

\begin{lemma}
\label{coherent-lemma-property-irreducible}
Soit $X$ un schéma noethérien. Soit $Z_0 \subset X$ un sous-ensemble fermé irréductible
de point générique $\xi$. Soit $\mathcal{P}$ une propriété des faisceaux
cohérents sur $X$ dont le support est contenu dans $Z_0$, telle que
\begin{enumerate}
\item pour toute suite exacte courte de faisceaux cohérents, si
deux d'entre eux vérifient la propriété $\mathcal{P}$, il en va de même pour
le troisième;
\item pour tout sous-schéma fermé intègre $Z \subset Z_0 \subset X$,
$Z \not = Z_0$, et tout faisceau quasi-cohérent d'idéaux
$\mathcal{I} \subset \mathcal{O}_Z$, on ait
$\mathcal{P}$ pour $(Z \to X)_*\mathcal{I}$;
\item il existe un faisceau cohérent $\mathcal{G}$ sur $X$ tel que
\begin{enumerate}
\item $\text{Supp}(\mathcal{G}) = Z_0$,
\item $\mathcal{G}_\xi$ soit annulé par $\mathfrak m_\xi$,
\item $\dim_{\kappa(\xi)} \mathcal{G}_\xi = 1$, et
\item la propriété $\mathcal{P}$ soit vérifiée par $\mathcal{G}$.
\end{enumerate}
\end{enumerate}
Alors la propriété $\mathcal{P}$ est vérifiée par tout faisceau cohérent
$\mathcal{F}$ sur $X$ dont le support est contenu dans $Z_0$.
\end{lemma}

\begin{proof}
Notons d'abord que, si $\mathcal{F}$ est un faisceau cohérent de support
contenu dans $Z_0$, muni d'une filtration
$$
0 = \mathcal{F}_0 \subset \mathcal{F}_1 \subset
\ldots \subset \mathcal{F}_m = \mathcal{F}
$$
par des sous-faisceaux cohérents telle que chaque $\mathcal{F}_i/\mathcal{F}_{i - 1}$
vérifie la propriété $\mathcal{P}$, il en va de même pour $\mathcal{F}$. De même, si $\mathcal{F}$
vérifie la propriété $\mathcal{P}$ et que tous les
$\mathcal{F}_i/\mathcal{F}_{i - 1}$ sauf un vérifient $\mathcal{P}$,
le dernier la vérifie aussi. Cela résulte de l'hypothèse (1).

\medskip\noindent
On en déduit d'abord que tout faisceau cohérent dont le support
est strictement contenu dans $Z_0$ vérifie la propriété $\mathcal{P}$. En effet, un
tel faisceau admet une filtration (voir lemme \ref{coherent-lemma-coherent-filter})
dont les sous-quotients vérifient $\mathcal{P}$ d'après (2).

\medskip\noindent
Soit $\mathcal{G}$ comme dans (3). D'après le lemme \ref{coherent-lemma-prepare-filter-irreducible},
il existe un faisceau d'idéaux $\mathcal{I}$ sur $Z_0$, un
entier $r \geq 1$ et une suite exacte courte
$$
0 \to
\left((Z_0 \to X)_*\mathcal{I}\right)^{\oplus r} \to
\mathcal{G} \to
\mathcal{Q} \to 0
$$
où le support de $\mathcal{Q}$ est strictement contenu dans $Z_0$.
D'après (3)(c), on a $r = 1$. Puisque $\mathcal{Q}$ vérifie lui aussi la propriété $\mathcal{P}$,
on en déduit que $(Z_0 \to X)_*\mathcal{I}$ vérifie la propriété
$\mathcal{P}$.

\medskip\noindent
Supposons ensuite que $\mathcal{I}' \not = 0$ soit un autre
faisceau quasi-cohérent d'idéaux sur $Z_0$. On peut considérer l'intersection
$\mathcal{I}'' = \mathcal{I}' \cap \mathcal{I}$, ce qui
donne deux suites exactes courtes
$$
0 \to
(Z_0 \to X)_*\mathcal{I}'' \to
(Z_0 \to X)_*\mathcal{I} \to
\mathcal{Q} \to 0
$$
et
$$
0 \to
(Z_0 \to X)_*\mathcal{I}'' \to
(Z_0 \to X)_*\mathcal{I}' \to
\mathcal{Q}' \to 0.
$$
Notons que les supports des faisceaux cohérents $\mathcal{Q}$ et
$\mathcal{Q}'$ sont strictement contenus dans $Z_0$.
Ainsi $\mathcal{Q}$ et $\mathcal{Q}'$ vérifient la propriété $\mathcal{P}$
(voir ci-dessus). La condition (1) entraîne alors
que $(Z_0 \to X)_*\mathcal{I}''$ et $(Z_0 \to X)_*\mathcal{I}'$
vérifient eux aussi $\mathcal{P}$.

\medskip\noindent
Pour achever la démonstration, il suffit de noter que tout faisceau cohérent
$\mathcal{F}$ sur $X$ dont le support est contenu dans $Z_0$ admet une filtration (voir
à nouveau lemme \ref{coherent-lemma-coherent-filter}) dont les sous-quotients
vérifient tous la propriété $\mathcal{P}$ d'après ce qui précède.
\end{proof}

\begin{lemma}
\label{coherent-lemma-property}
Soit $X$ un schéma noethérien.
Soit $\mathcal{P}$ une propriété des faisceaux cohérents sur $X$ telle que
\begin{enumerate}
\item pour toute suite exacte courte de faisceaux cohérents, si
deux d'entre eux vérifient la propriété $\mathcal{P}$, il en va de même pour
le troisième;
\item pour tout sous-schéma fermé intègre $Z \subset X$
de point générique $\xi$, il
existe un faisceau cohérent $\mathcal{G}$ tel que
\begin{enumerate}
\item $\text{Supp}(\mathcal{G}) = Z$,
\item $\mathcal{G}_\xi$ soit annulé par $\mathfrak m_\xi$,
\item $\dim_{\kappa(\xi)} \mathcal{G}_\xi = 1$, et
\item la propriété $\mathcal{P}$ soit vérifiée par $\mathcal{G}$.
\end{enumerate}
\end{enumerate}
Alors la propriété $\mathcal{P}$ est vérifiée par tout faisceau cohérent
sur $X$.
\end{lemma}

\begin{proof}
D'après le lemme \ref{coherent-lemma-property-initial}, il suffit de montrer que,
pour tout sous-schéma fermé intègre $Z \subset X$ et tout faisceau
quasi-cohérent d'idéaux $\mathcal{I} \subset \mathcal{O}_Z$, la propriété $\mathcal{P}$
est vérifiée par $(Z \to X)_*\mathcal{I}$. Sinon, puisque $X$ est noethérien,
il existe un sous-schéma fermé intègre minimal $Z_0 \subset X$ tel que
$\mathcal{P}$ ne soit pas vérifiée par $(Z_0 \to X)_*\mathcal{I}_0$ pour un
certain faisceau quasi-cohérent d'idéaux $\mathcal{I}_0 \subset \mathcal{O}_{Z_0}$, mais
que $\mathcal{P}$ soit vérifiée par $(Z \to X)_*\mathcal{I}$ pour tout sous-schéma
fermé intègre $Z \subset Z_0$, $Z \not = Z_0$, et tout
faisceau quasi-cohérent d'idéaux $\mathcal{I} \subset \mathcal{O}_Z$. Puisque l'existence
de $\mathcal{G}$ pour $Z_0$ est assurée par (2), le
lemme \ref{coherent-lemma-property-irreducible} exclut cette éventualité.
\end{proof}

\begin{lemma}
\label{coherent-lemma-property-irreducible-higher-rank-cohomological}
Soit $X$ un schéma noethérien. Soit $Z_0 \subset X$ un sous-ensemble fermé
irréductible de point générique $\xi$. Soit $\mathcal{P}$ une propriété des
faisceaux cohérents sur $X$ telle que
\begin{enumerate}
\item pour toute suite exacte courte de faisceaux cohérents
$$
0 \to \mathcal{F}_1 \to \mathcal{F} \to \mathcal{F}_2 \to 0
$$
si les $\mathcal{F}_i$, $i = 1, 2$, vérifient la propriété $\mathcal{P}$,
il en va de même pour $\mathcal{F}$;
\item si $\mathcal{P}$ est vérifiée par $\mathcal{F}^{\oplus r}$ pour un
certain $r \geq 1$, elle est vérifiée par $\mathcal{F}$;
\item pour tout sous-schéma fermé intègre $Z \subset Z_0 \subset X$,
$Z \not = Z_0$, et tout faisceau quasi-cohérent d'idéaux
$\mathcal{I} \subset \mathcal{O}_Z$, on ait
$\mathcal{P}$ pour $(Z \to X)_*\mathcal{I}$;
\item il existe un faisceau cohérent $\mathcal{G}$ tel que
\begin{enumerate}
\item $\text{Supp}(\mathcal{G}) = Z_0$,
\item $\mathcal{G}_\xi$ soit annulé par $\mathfrak m_\xi$, et
\item pour tout faisceau quasi-cohérent d'idéaux
$\mathcal{J} \subset \mathcal{O}_X$ tel que
$\mathcal{J}_\xi = \mathcal{O}_{X, \xi}$, il existe un sous-faisceau
quasi-cohérent $\mathcal{G}' \subset \mathcal{J}\mathcal{G}$ avec
$\mathcal{G}'_\xi = \mathcal{G}_\xi$ et tel que
$\mathcal{P}$ soit vérifiée par $\mathcal{G}'$.
\end{enumerate}
\end{enumerate}
Alors la propriété $\mathcal{P}$ est vérifiée par tout faisceau cohérent
$\mathcal{F}$ sur $X$ dont le support est contenu dans $Z_0$.
\end{lemma}

\begin{proof}
Notons que, si $\mathcal{F}$ est un faisceau cohérent muni d'une filtration
$$
0 = \mathcal{F}_0 \subset \mathcal{F}_1 \subset
\ldots \subset \mathcal{F}_m = \mathcal{F}
$$
par des sous-faisceaux cohérents telle que chaque $\mathcal{F}_i/\mathcal{F}_{i - 1}$
vérifie la propriété $\mathcal{P}$, il en va de même pour $\mathcal{F}$.
Cela résulte de l'hypothèse (1).

\medskip\noindent
On en déduit d'abord que tout faisceau cohérent dont le support
est strictement contenu dans $Z_0$ vérifie la propriété $\mathcal{P}$. En effet, un
tel faisceau admet une filtration (voir lemme \ref{coherent-lemma-coherent-filter})
dont les sous-quotients vérifient $\mathcal{P}$ d'après (3).

\medskip\noindent
Notons $i : Z_0 \to X$ l'immersion fermée. Considérons
un faisceau cohérent $\mathcal{G}$ comme dans (4). D'après
le lemme \ref{coherent-lemma-prepare-filter-irreducible},
il existe un faisceau d'idéaux $\mathcal{I}$ sur $Z_0$ et
une suite exacte courte
$$
0 \to
i_*\mathcal{I}^{\oplus r} \to
\mathcal{G} \to
\mathcal{Q} \to 0
$$
où le support de $\mathcal{Q}$ est strictement contenu dans $Z_0$.
En particulier, $r > 0$ et $\mathcal{I}$ est non
nul, car le support de $\mathcal{G}$ est égal à $Z_0$.
Soit $\mathcal{I}' \subset \mathcal{I}$ un faisceau quasi-cohérent d'idéaux
non nul sur $Z_0$ contenu dans $\mathcal{I}$.
On obtient alors aussi une suite exacte courte
$$
0 \to
i_*(\mathcal{I}')^{\oplus r} \to
\mathcal{G} \to
\mathcal{Q}' \to 0
$$
où $\mathcal{Q}'$ a son support strictement contenu dans $Z_0$.
Soit $\mathcal{J} \subset \mathcal{O}_X$ un faisceau quasi-cohérent
d'idéaux définissant le support de $\mathcal{Q}'$ (par exemple
l'idéal correspondant à la structure réduite induite de sous-schéma
fermé sur le support de $\mathcal{Q}'$). Alors
$\mathcal{J}_\xi = \mathcal{O}_{X, \xi}$. D'après le
lemme \ref{coherent-lemma-power-ideal-kills-sheaf},
on a $\mathcal{J}^n\mathcal{Q}' = 0$ pour un certain $n$.
Ainsi $\mathcal{J}^n\mathcal{G} \subset i_*(\mathcal{I}')^{\oplus r}$.
L'hypothèse (4)(c) du lemme assure l'existence d'un
sous-faisceau quasi-cohérent $\mathcal{G}' \subset \mathcal{J}^n\mathcal{G}$
tel que $\mathcal{G}'_\xi = \mathcal{G}_\xi$
et qui vérifie la propriété $\mathcal{P}$.
On obtient donc une suite exacte courte
$$
0 \to \mathcal{G}' \to
i_*(\mathcal{I}')^{\oplus r} \to
\mathcal{Q}'' \to 0
$$
où $\mathcal{Q}''$ a son support strictement contenu dans $Z_0$.
D'après nos remarques initiales et la condition (1) du lemme, on
en déduit que $i_*(\mathcal{I}')^{\oplus r}$ vérifie
$\mathcal{P}$. Ainsi $i_*\mathcal{I}'$ vérifie
$\mathcal{P}$ d'après (2). Enfin, pour un faisceau
quasi-cohérent d'idéaux quelconque $\mathcal{I}'' \subset \mathcal{O}_{Z_0}$, on peut poser
$\mathcal{I}' = \mathcal{I}'' \cap \mathcal{I}$, ce qui
donne une suite exacte courte
$$
0 \to
i_*(\mathcal{I}') \to
i_*(\mathcal{I}'') \to
\mathcal{Q}''' \to 0
$$
où $\mathcal{Q}'''$ a son support strictement contenu dans $Z_0$.
On en déduit que la propriété $\mathcal{P}$ est également vérifiée par
$i_*\mathcal{I}''$.

\medskip\noindent
Pour achever la démonstration, il suffit de noter que tout faisceau cohérent
$\mathcal{F}$ sur $X$ dont le support est contenu dans $Z_0$ admet une filtration (voir
à nouveau lemme \ref{coherent-lemma-coherent-filter}) dont les sous-quotients
vérifient tous la propriété $\mathcal{P}$ d'après ce qui précède.
\end{proof}

\begin{lemma}
\label{coherent-lemma-property-higher-rank-cohomological}
Soit $X$ un schéma noethérien.
Soit $\mathcal{P}$ une propriété des faisceaux cohérents sur $X$ telle que
\begin{enumerate}
\item pour toute suite exacte courte de faisceaux cohérents
$$
0 \to \mathcal{F}_1 \to \mathcal{F} \to \mathcal{F}_2 \to 0
$$
si les $\mathcal{F}_i$, $i = 1, 2$, vérifient la propriété $\mathcal{P}$,
il en va de même pour $\mathcal{F}$;
\item si $\mathcal{P}$ est vérifiée par $\mathcal{F}^{\oplus r}$ pour un
certain $r \geq 1$, elle est vérifiée par $\mathcal{F}$;
\item pour tout sous-schéma fermé intègre $Z \subset X$
de point générique $\xi$, il
existe un faisceau cohérent $\mathcal{G}$ tel que
\begin{enumerate}
\item $\text{Supp}(\mathcal{G}) = Z$,
\item $\mathcal{G}_\xi$ soit annulé par $\mathfrak m_\xi$, et
\item pour tout faisceau quasi-cohérent d'idéaux
$\mathcal{J} \subset \mathcal{O}_X$ tel que
$\mathcal{J}_\xi = \mathcal{O}_{X, \xi}$, il existe un sous-faisceau
quasi-cohérent $\mathcal{G}' \subset \mathcal{J}\mathcal{G}$ avec
$\mathcal{G}'_\xi = \mathcal{G}_\xi$ et tel que
$\mathcal{P}$ soit vérifiée par $\mathcal{G}'$.
\end{enumerate}
\end{enumerate}
Alors la propriété $\mathcal{P}$ est vérifiée par tout faisceau cohérent
sur $X$.
\end{lemma}

\begin{proof}
Cela résulte du lemme \ref{coherent-lemma-property-irreducible-higher-rank-cohomological}
exactement comme le lemme \ref{coherent-lemma-property} résulte du
lemme \ref{coherent-lemma-property-irreducible}.
\end{proof}








\section{Morphismes finis et schémas affines}
\label{coherent-section-finite-affine}

\noindent
Dans cette section, nous utilisons les résultats des sections précédentes
pour montrer que l'image d'un schéma affine noethérien par un morphisme fini
est affine. Nous verrons plus loin que ce résultat vaut plus
généralement (voir Limites, lemme \ref{limits-lemma-affine} et
proposition \ref{limits-proposition-affine}).

\begin{lemma}
\label{coherent-lemma-finite-morphism-Noetherian}
Soit $f : Y \to X$ un morphisme de schémas.
Supposons $f$ fini et surjectif, et $X$ localement noethérien.
Soit $Z \subset X$ un sous-schéma fermé intègre de
point générique $\xi$. Il
existe alors un faisceau cohérent $\mathcal{F}$ sur $Y$
tel que le support de $f_*\mathcal{F}$ soit égal à $Z$
et que $(f_*\mathcal{F})_\xi$ soit annulé par $\mathfrak m_\xi$.
\end{lemma}

\begin{proof}
Notons que $Y$ est localement noethérien d'après
Morphismes, lemme \ref{morphisms-lemma-finite-type-noetherian}.
Puisque $f$ est surjectif, la fibre $Y_\xi$ est non vide.
Choisissons $\xi' \in Y$ d'image $\xi$. Posons $Z' = \overline{\{\xi'\}}$.
On peut considérer $Z' \subset Y$ comme un sous-schéma fermé réduit,
voir Schémas, lemme \ref{schemes-lemma-reduced-closed-subscheme}.
Ainsi le faisceau $\mathcal{F} = (Z' \to Y)_*\mathcal{O}_{Z'}$
est un faisceau cohérent sur $Y$ (voir
lemme \ref{coherent-lemma-finite-pushforward-coherent}).
Considérons le diagramme commutatif
$$
\xymatrix{
Z' \ar[r]_{i'} \ar[d]_{f'} &
Y \ar[d]^f \\
Z \ar[r]^i &
X
}
$$
On a $f_*\mathcal{F} = i_*f'_*\mathcal{O}_{Z'}$.
La fibre de $f_*\mathcal{F}$ en $\xi$ est donc celle
de $f'_*\mathcal{O}_{Z'}$ en $\xi$. Puisque $Z'$ est
intègre de point générique $\xi'$, on voit que
$\xi'$ est le seul point de $Z'$ au-dessus de $\xi$, voir
Algèbre, lemmes \ref{algebra-lemma-finite-is-integral} et
\ref{algebra-lemma-integral-no-inclusion}.
La fibre de $f'_*\mathcal{O}_{Z'}$ en $\xi$
est donc égale à $\mathcal{O}_{Z', \xi'} = \kappa(\xi')$. En particulier,
la fibre de $f_*\mathcal{F}$ en $\xi$ est non nulle.
Cela, joint au fait que $f_*\mathcal{F}$ est
de la forme $i_*f'_*(\text{quelque chose})$, entraîne le lemme.
\end{proof}

\begin{lemma}
\label{coherent-lemma-affine-morphism-projection-ideal}
Soit $f : Y \to X$ un morphisme de schémas.
Soit $\mathcal{F}$ un faisceau quasi-cohérent sur $Y$.
Soit $\mathcal{I}$ un faisceau quasi-cohérent d'idéaux sur $X$.
Si le morphisme $f$ est affine, alors
$\mathcal{I}f_*\mathcal{F} = f_*(f^{-1}\mathcal{I}\mathcal{F})$.
\end{lemma}

\begin{proof}
Précisons cette notation. Puisque $f^{-1}$ est un
foncteur exact, $f^{-1}\mathcal{I}$ est un
faisceau d'idéaux de $f^{-1}\mathcal{O}_X$. Par le morphisme
$f^\sharp : f^{-1}\mathcal{O}_X \to \mathcal{O}_Y$, il opère sur
$\mathcal{F}$. Alors $f^{-1}\mathcal{I}\mathcal{F}$ est le sous-faisceau
engendré par les sommes de sections locales de la forme $as$, où $a$
est une section locale de $f^{-1}\mathcal{I}$ et $s$ une section locale
de $\mathcal{F}$. C'est un sous-$\mathcal{O}_Y$-module quasi-cohérent
de $\mathcal{F}$, car c'est aussi l'image d'un morphisme naturel
$f^*\mathcal{I} \otimes_{\mathcal{O}_Y} \mathcal{F} \to \mathcal{F}$.

\medskip\noindent
Cela étant, la démonstration est immédiate. La question est locale, et
on peut donc supposer $X$ affine. Puisque $f$ est affine,
$Y$ est lui aussi affine. On peut donc écrire
$Y = \Spec(B)$, $X = \Spec(A)$, $\mathcal{F} = \widetilde{M}$
et $\mathcal{I} = \widetilde{I}$. Dans ce cas, l'assertion du
lemme se ramène à l'égalité
$$
I(M_A) = ((IB)M)_A
$$
où $M_A$ désigne le $A$-module sous-jacent au $B$-module $M$.
\end{proof}

\begin{lemma}
\label{coherent-lemma-image-affine-finite-morphism-affine-Noetherian}
Soit $f : Y \to X$ un morphisme de schémas. Supposons
que
\begin{enumerate}
\item $f$ soit fini,
\item $f$ soit surjectif,
\item $Y$ soit affine, et
\item $X$ soit noethérien.
\end{enumerate}
Alors $X$ est affine.
\end{lemma}

\begin{proof}
Nous allons montrer que, sous les hypothèses du lemme, pour tout
$\mathcal{O}_X$-module cohérent $\mathcal{F}$, on a $H^1(X, \mathcal{F}) = 0$.
Cela entraînera en particulier que $H^1(X, \mathcal{I}) = 0$
pour tout faisceau quasi-cohérent d'idéaux de $\mathcal{O}_X$. On en
déduira que $X$ est affine, soit par le
lemme \ref{coherent-lemma-quasi-compact-h1-zero-covering}, soit
par le lemme \ref{coherent-lemma-quasi-separated-h1-zero-covering}.

\medskip\noindent
Soit $\mathcal{P}$ la propriété des faisceaux cohérents
$\mathcal{F}$ sur $X$ définie par
$$
\mathcal{P}(\mathcal{F}) \Leftrightarrow H^1(X, \mathcal{F}) = 0.
$$
Nous allons appliquer le lemme \ref{coherent-lemma-property-higher-rank-cohomological}.
Il faut donc vérifier les conditions (1), (2) et (3) de ce lemme pour $\mathcal{P}$.
La propriété (1) résulte de la suite exacte longue de cohomologie associée à
une suite exacte courte de faisceaux. La propriété (2) résulte de ce que
$H^1(X, -)$ est un foncteur additif. Pour vérifier (3), soit $Z \subset X$ un
sous-schéma fermé intègre de point générique $\xi$.
Soit $\mathcal{F}$ un faisceau cohérent sur $Y$ tel que
le support de $f_*\mathcal{F}$ soit égal à $Z$
et que $(f_*\mathcal{F})_\xi$ soit annulé par $\mathfrak m_\xi$,
voir lemme \ref{coherent-lemma-finite-morphism-Noetherian}. Nous affirmons
que $\mathcal{G} = f_*\mathcal{F}$ convient. Il suffit de vérifier la
condition (3)(c) du lemme \ref{coherent-lemma-property-higher-rank-cohomological}.
Supposons donc que $\mathcal{J} \subset \mathcal{O}_X$ soit un
faisceau quasi-cohérent d'idéaux tel que
$\mathcal{J}_\xi = \mathcal{O}_{X, \xi}$.
Un morphisme fini est affine; d'après le
lemme \ref{coherent-lemma-affine-morphism-projection-ideal}, on a donc
$\mathcal{J}\mathcal{G} = f_*(f^{-1}\mathcal{J}\mathcal{F})$.
De plus, comme on l'a remarqué dans la démonstration du
lemme \ref{coherent-lemma-affine-morphism-projection-ideal}, le faisceau
$f^{-1}\mathcal{J}\mathcal{F}$ est un $\mathcal{O}_Y$-module quasi-cohérent.
Puisque $Y$ est affine, on a $H^1(Y, f^{-1}\mathcal{J}\mathcal{F}) = 0$,
voir lemme \ref{coherent-lemma-quasi-coherent-affine-cohomology-zero}.
Puisque $f$ est fini, donc affine, on a
$$
H^1(X, \mathcal{J}\mathcal{G}) =
H^1(X, f_*(f^{-1}\mathcal{J}\mathcal{F})) =
H^1(Y, f^{-1}\mathcal{J}\mathcal{F}) = 0
$$
d'après le lemme \ref{coherent-lemma-relative-affine-cohomology}.
Le sous-faisceau quasi-cohérent $\mathcal{G}' = \mathcal{J}\mathcal{G}$
vérifie donc $\mathcal{P}$. Cela établit la propriété (3)(c) du
lemme \ref{coherent-lemma-property-higher-rank-cohomological}, comme voulu.
\end{proof}












\section{Faisceaux cohérents sur Proj, I}
\label{coherent-section-coherent-proj}

\noindent
Dans cette section, nous étudions les faisceaux cohérents sur $\text{Proj}(A)$,
où $A$ est un anneau gradué noethérien engendré par $A_1$ sur $A_0$.
Dans la section suivante, nous examinerons le cas où $A$ n'est pas engendré par ses
éléments de degré $1$. Commençons par un résultat d'ensemble pour
l'espace projectif sur un anneau noethérien.

\begin{lemma}
\label{coherent-lemma-coherent-projective}
Soit $R$ un anneau noethérien.
Soit $n \geq 0$ un entier. Pour
tout faisceau cohérent $\mathcal{F}$ sur $\mathbf{P}^n_R$,
on a les propriétés suivantes :
\begin{enumerate}
\item Il existe $r \geq 0$,
$d_1, \ldots, d_r \in \mathbf{Z}$ et une surjection
$$
\bigoplus\nolimits_{j = 1, \ldots, r}
\mathcal{O}_{\mathbf{P}^n_R}(d_j)
\longrightarrow
\mathcal{F}.
$$
\item On a $H^i(\mathbf{P}^n_R, \mathcal{F}) = 0$ sauf si
$0 \leq i \leq n$.
\item Pour tout $i$, le groupe de cohomologie $H^i(\mathbf{P}^n_R, \mathcal{F})$
est un $R$-module de type fini.
\item Si $i > 0$, alors
$H^i(\mathbf{P}^n_R, \mathcal{F}(d)) = 0$ pour tout $d$ assez grand.
\item Pour tout $k \in \mathbf{Z}$, le $R[T_0, \ldots, T_n]$-module gradué
$$
\bigoplus\nolimits_{d \geq k} H^0(\mathbf{P}^n_R, \mathcal{F}(d))
$$
est un $R[T_0, \ldots, T_n]$-module de type fini.
\end{enumerate}
\end{lemma}

\begin{proof}
Nous utiliserons le fait que $\mathcal{O}_{\mathbf{P}^n_R}(1)$ est un faisceau inversible
ample sur
le schéma $\mathbf{P}^n_R$. Cela résulte directement de la définition,
puisque $\mathbf{P}^n_R$ est recouvert par les ouverts affines standard $D_{+}(T_i)$.
Ainsi, d'après
Propriétés, proposition \ref{properties-proposition-characterize-ample},
tout $\mathcal{O}_{\mathbf{P}^n_R}$-module
quasi-cohérent de type fini est quotient d'une somme directe finie de puissances tensorielles de
$\mathcal{O}_{\mathbf{P}^n_R}(1)$. D'autre part, les faisceaux cohérents
coïncident avec les faisceaux quasi-cohérents de type fini sur l'espace
projectif sur $R$, d'après le lemme \ref{coherent-lemma-coherent-Noetherian}. D'où (1).

\medskip\noindent
L'espace projectif de dimension $n$, $\mathbf{P}^n_R$, est recouvert par $n + 1$ ouverts
affines, à savoir les ouverts standard $D_{+}(T_i)$, $i = 0, \ldots, n$, voir Constructions,
lemme \ref{constructions-lemma-standard-covering-projective-space}.
Ainsi, pour tout faisceau
quasi-cohérent $\mathcal{F}$ sur $\mathbf{P}^n_R$,
on a $H^i(\mathbf{P}^n_R, \mathcal{F}) = 0$ pour $i \geq n + 1$,
voir lemme \ref{coherent-lemma-vanishing-nr-affines}. D'où (2).

\medskip\noindent
Démontrons (3) et (4) simultanément pour tous les faisceaux cohérents
sur $\mathbf{P}^n_R$, par récurrence descendante sur $i$. Le résultat est
vrai pour $i \geq n + 1$ d'après (2). Supposons-le connu pour
$i + 1$ et montrons-le pour $i$. (Si $i = 0$, l'assertion
(4) est sans objet.) Soit $\mathcal{F}$ un faisceau cohérent sur
$\mathbf{P}^n_R$. Choisissons une surjection comme dans (1) et notons
$\mathcal{G}$ son noyau, de sorte que l'on ait une suite exacte courte
$$
0 \to \mathcal{G} \to
\bigoplus\nolimits_{j = 1, \ldots, r}
\mathcal{O}_{\mathbf{P}^n_R}(d_j)
\to
\mathcal{F} \to 0
$$
D'après le lemme \ref{coherent-lemma-coherent-abelian-Noetherian},
$\mathcal{G}$ est cohérent. La suite exacte
longue de cohomologie donne une suite exacte
$$
H^i(\mathbf{P}^n_R, \bigoplus\nolimits_{j = 1, \ldots, r}
\mathcal{O}_{\mathbf{P}^n_R}(d_j))
\to
H^i(\mathbf{P}^n_R, \mathcal{F})
\to
H^{i + 1}(\mathbf{P}^n_R, \mathcal{G}).
$$
Par hypothèse de récurrence, le $R$-module de droite est de type fini et, d'après
le lemme \ref{coherent-lemma-cohomology-projective-space-over-ring}, le
$R$-module de gauche est de type fini. Puisque $R$ est noethérien, il en résulte aussitôt
que $H^i(\mathbf{P}^n_R, \mathcal{F})$ est un $R$-module
de type fini. Cela établit l'étape de récurrence pour l'assertion (3).
Puisque $\mathcal{O}_{\mathbf{P}^n_R}(d)$ est
inversible, tordre un faisceau sur $\mathbf{P}^n_R$ définit un foncteur exact (car
ce foncteur est donné par le produit tensoriel avec un faisceau inversible, voir
Constructions, définition \ref{constructions-definition-twist}).
Autrement dit, pour tout $d \in \mathbf{Z}$, la suite
$$
0 \to \mathcal{G}(d) \to
\bigoplus\nolimits_{j = 1, \ldots, r}
\mathcal{O}_{\mathbf{P}^n_R}(d_j + d)
\to
\mathcal{F}(d) \to 0
$$
est exacte courte. La suite de cohomologie qui en résulte est
$$
H^i(\mathbf{P}^n_R, \bigoplus\nolimits_{j = 1, \ldots, r}
\mathcal{O}_{\mathbf{P}^n_R}(d_j + d))
\to
H^i(\mathbf{P}^n_R, \mathcal{F}(d))
\to
H^{i + 1}(\mathbf{P}^n_R, \mathcal{G}(d)).
$$
Par hypothèse de récurrence, le module de droite est nul
pour $d \gg 0$ et, d'après le calcul du
lemme \ref{coherent-lemma-cohomology-projective-space-over-ring},
le module de gauche est nul dès que $d \geq -\min\{d_j\}$
et $i \geq 1$. D'où l'étape de récurrence pour l'assertion (4).
Cela achève la démonstration de (3) et (4).

\medskip\noindent
Pour démontrer (5), notons que, pour tout $d$
assez grand, l'application
$$
H^0(\mathbf{P}^n_R, \bigoplus\nolimits_{j = 1, \ldots, r}
\mathcal{O}_{\mathbf{P}^n_R}(d_j + d))
\to
H^0(\mathbf{P}^n_R, \mathcal{F}(d))
$$
est surjective, d'après l'annulation de $H^1(\mathbf{P}^n_R, \mathcal{G}(d))$
que nous venons de démontrer. Autrement dit, le module
$$
M_k
=
\bigoplus\nolimits_{d \geq k} H^0(\mathbf{P}^n_R, \mathcal{F}(d))
$$
est, pour $k$ assez grand, un quotient du module correspondant
$$
N_k
=
\bigoplus\nolimits_{d \geq k} H^0(\mathbf{P}^n_R,
\bigoplus\nolimits_{j = 1, \ldots, r}
\mathcal{O}_{\mathbf{P}^n_R}(d_j + d)
)
$$
Lorsque $k$ est assez petit (par exemple\ $k < -d_j$ pour tout $j$), on a
$$
N_k = \bigoplus\nolimits_{j = 1, \ldots, r}
R[T_0, \ldots, T_n](d_j)
$$
d'après les calculs de la section \ref{coherent-section-cohomology-projective-space}.
Il est donc en particulier de type fini.
Soit $k \in \mathbf{Z}$ quelconque.
Choisissons $k_{-} \ll k \ll k_{+}$.
Considérons le diagramme
$$
\xymatrix{
N_{k_{-}} & N_{k_{+}} \ar[d] \ar[l] \\
M_k & M_{k_{+}} \ar[l]
}
$$
où la flèche verticale est l'application surjective ci-dessus
et où les flèches horizontales sont les
inclusions évidentes. D'après ce qui précède, $N_{k_{-}}$ est
un $R[T_0, \ldots, T_n]$-module de type fini. Ainsi $N_{k_{+}}$ est
un $R[T_0, \ldots, T_n]$-module de type
fini, car c'est un sous-module d'un module de type fini et l'anneau
$R[T_0, \ldots, T_n]$ est noethérien. La flèche
verticale étant surjective, $M_{k_{+}}$ est
un $R[T_0, \ldots, T_n]$-module de type fini. Le quotient
$M_k/M_{k_{+}}$ est de type fini comme $R$-module, car c'est une
somme directe finie des $R$-modules de type fini
$H^0(\mathbf{P}^n_R, \mathcal{F}(d))$ pour $k \leq d < k_{+}$.
On utilise ici l'assertion (3) pour $i = 0$. Ainsi
$M_k/M_{k_{+}}$ est a fortiori un $R[T_0, \ldots, T_n]$-module
de type fini. Autrement dit, $M_k$ est extension de deux
$R[T_0, \ldots, T_n]$-modules de type fini, ce qui conclut.
\end{proof}

\begin{lemma}
\label{coherent-lemma-coherent-on-proj}
Soit $A$ un anneau gradué tel que $A_0$ soit noethérien et que
$A$ soit engendré par un nombre fini d'éléments de $A_1$ sur $A_0$.
Posons $X = \text{Proj}(A)$. Alors $X$ est un schéma noethérien.
Soit $\mathcal{F}$ un $\mathcal{O}_X$-module cohérent.
\begin{enumerate}
\item Il existe $r \geq 0$,
$d_1, \ldots, d_r \in \mathbf{Z}$ et une surjection
$$
\bigoplus\nolimits_{j = 1, \ldots, r} \mathcal{O}_X(d_j)
\longrightarrow \mathcal{F}.
$$
\item Pour tout $p$, le groupe de cohomologie $H^p(X, \mathcal{F})$ est un
$A_0$-module de type fini.
\item Si $p > 0$, alors $H^p(X, \mathcal{F}(d)) = 0$ pour tout $d$ assez grand.
\item Pour tout $k \in \mathbf{Z}$, le $A$-module gradué
$$
\bigoplus\nolimits_{d \geq k} H^0(X, \mathcal{F}(d))
$$
est un $A$-module de type fini.
\end{enumerate}
\end{lemma}

\begin{proof}
Par hypothèse, il existe une surjection de $A_0$-algèbres graduées
$$
A_0[T_0, \ldots, T_n] \longrightarrow A
$$
où $\deg(T_j) = 1$ pour $j = 0, \ldots, n$. D'après Constructions, Lemme
\ref{constructions-lemma-surjective-graded-rings-generated-degree-1-map-proj},
elle définit une immersion fermée $i : X \to \mathbf{P}^n_{A_0}$
telle que $i^*\mathcal{O}_{\mathbf{P}^n_{A_0}}(1) = \mathcal{O}_X(1)$.
En particulier, $X$ est noethérien, étant un sous-schéma fermé du schéma
noethérien $\mathbf{P}^n_{A_0}$. Nous affirmons que les résultats du lemme pour
$\mathcal{F}$ se déduisent des résultats
correspondants du lemme \ref{coherent-lemma-coherent-projective} pour le faisceau cohérent
$i_*\mathcal{F}$ (lemme \ref{coherent-lemma-i-star-equivalence}) sur
$\mathbf{P}^n_{A_0}$. Par exemple, ce lemme assure
l'existence d'une surjection
$$
\bigoplus\nolimits_{j = 1, \ldots, r} \mathcal{O}_{\mathbf{P}^n_{A_0}}(d_j)
\longrightarrow i_*\mathcal{F}.
$$
Par adjonction, celle-ci correspond à une application
$\bigoplus_{j = 1, \ldots, r} \mathcal{O}_X(d_j) \longrightarrow \mathcal{F}$,
qui est elle aussi surjective. Les assertions sur la cohomologie résultent
de l'égalité
$H^p(X, \mathcal{F}(d)) = H^p(\mathbf{P}^n_{A_0}, i_*\mathcal{F}(d))$,
d'après le lemme \ref{coherent-lemma-relative-affine-cohomology}.
\end{proof}

\begin{lemma}
\label{coherent-lemma-recover-tail-graded-module}
Soit $A$ un anneau gradué tel que $A_0$ soit noethérien et que $A$ soit engendré
par un nombre fini d'éléments de $A_1$ sur $A_0$. Soit $M$
un $A$-module gradué de type fini. Posons $X = \text{Proj}(A)$ et soit $\widetilde{M}$
le $\mathcal{O}_X$-module quasi-cohérent sur $X$ associé à $M$.
Les applications
$$
M_n \longrightarrow \Gamma(X, \widetilde{M}(n))
$$
de Constructions, lemme \ref{constructions-lemma-apply-modules},
sont des isomorphismes pour tout $n$ assez grand.
\end{lemma}

\begin{proof}
Puisque $M$ est un $A$-module de type fini,
$\widetilde{M}$ est un $\mathcal{O}_X$-module de
type fini, c'est-à-dire un $\mathcal{O}_X$-module cohérent.
Posons $N = \bigoplus_{n \in \mathbf{Z}} \Gamma(X, \widetilde{M}(n))$.
Il faut montrer que l'application $M \to N$ de $A$-modules
gradués est un isomorphisme en tout degré assez grand.
D'après Propriétés, lemme \ref{properties-lemma-proj-quasi-coherent},
on a un isomorphisme canonique $\widetilde{N} \to \widetilde{M}$
tel que les applications induites $N_n \to N_n = \Gamma(X, \widetilde{M}(n))$
soient les identités. On a donc des applications
$\widetilde{M} \to \widetilde{N} \to \widetilde{M}$
telles que, pour tout $n$, le diagramme
$$
\xymatrix{
M_n \ar[d] \ar[r] & N_n \ar[d] \ar@{=}[rd] \\
\Gamma(X, \widetilde{M}(n)) \ar[r] &
\Gamma(X, \widetilde{N}(n)) \ar[r]^{\cong} &
\Gamma(X, \widetilde{M}(n))
}
$$
soit commutatif. Cela signifie que le composé
$$
M_n \to \Gamma(X, \widetilde{M}(n))
\to \Gamma(X, \widetilde{N}(n)) \to \Gamma(X, \widetilde{M}(n))
$$
est l'application canonique $M_n \to \Gamma(X, \widetilde{M}(n))$.
Cela entraîne clairement que le composé
$\widetilde{M} \to \widetilde{N} \to \widetilde{M}$
est l'identité. Ainsi $\widetilde{M} \to \widetilde{N}$ est un isomorphisme.
Posons $K = \Ker(M \to N)$ et $Q = \Coker(M \to N)$.
Rappelons que le foncteur $M \mapsto \widetilde{M}$ est exact, voir
Constructions, lemme \ref{constructions-lemma-proj-sheaves}.
On a donc $\widetilde{K} = 0$ et $\widetilde{Q} = 0$.
Rappelons que $A$ est un anneau noethérien, que $M$ est un
$A$-module de type fini et que $N$ est un $A$-module gradué tel que
$N' = \bigoplus_{n \geq 0} N_n$ soit de type fini,
d'après la dernière assertion du lemme \ref{coherent-lemma-coherent-on-proj}.
Ainsi $K' = \bigoplus_{n \geq 0} K_n$ et
$Q' = \bigoplus_{n \geq 0} Q_n$ sont des $A$-modules de
type fini. Notons que $\widetilde{Q} = \widetilde{Q'}$ et
$\widetilde{K} = \widetilde{K'}$.
Pour achever la démonstration, il suffit donc
de montrer qu'un $A$-module de type fini $K$ tel que $\widetilde{K} = 0$
n'a qu'un nombre fini de composantes homogènes non nulles $K_d$ avec $d \geq 0$.
Pour cela, soient $x_1, \ldots, x_r \in K$ des générateurs homogènes,
de degrés respectifs $d_1, \ldots, d_r$.
Soient $f_1, \ldots, f_n \in A_1$ des éléments engendrant $A$ sur $A_0$.
Pour tout $i$ et tout $j$, il existe $n_{ij} \geq 0$ tel que
$f_i^{n_{ij}} x_j = 0$ dans $K_{d_j + n_{ij}}$ : sinon
$x_i/f_i^{d_i} \in K_{(f_i)}$ serait non nul, donc $\widetilde{K}$
serait non
nul. Alors $K_d$ est nul pour $d > \max_j(d_j + \sum_i n_{ij})$,
car tout élément de $K_d$ est une somme de termes, chacun étant un
monôme en les $f_i$ multiplié par l'un des $x_j$, de degré total $d$.
\end{proof}

\noindent
Soit $A$ un anneau gradué tel que $A_0$ soit noethérien et que $A$ soit engendré
par un nombre fini d'éléments de $A_1$ sur $A_0$. Rappelons que
$A_+ = \bigoplus_{n > 0} A_n$ est l'idéal des éléments de degré strictement
positif. Soit $M$ un $A$-module gradué. Rappelons que
$M$ est un module de torsion pour les puissances de $A_+$ si, pour tout $x \in M$,
il existe $n \geq 1$ tel que $(A_+)^n x = 0$, voir
Compléments d'algèbre, définition \ref{more-algebra-definition-f-power-torsion}.
Si $M$ est de type fini,
cela équivaut à $M_n = 0$ pour $n \gg 0$. Les modules de
torsion pour les puissances de $A_+$ sont parfois appelés modules de torsion. On
dit parfois qu'un $A$-module gradué $M$ est sans torsion si
$x \in M$ avec $(A_+)^n x = 0$, $n > 0$, entraîne $x = 0$.
Notons $\text{Mod}_A$ la catégorie des $A$-modules gradués,
$\text{Mod}^{fg}_A$ la sous-catégorie pleine des modules de type fini,
et $\text{Mod}^{fg}_{A, torsion}$ la sous-catégorie pleine des
modules $M$ tels que $M_n = 0$ pour $n \gg 0$.

\begin{proposition}
\label{coherent-proposition-coherent-modules-on-proj}
Soit $A$ un anneau gradué tel que $A_0$ soit noethérien et que $A$ soit engendré
par un nombre fini d'éléments de $A_1$ sur $A_0$.
Posons $X = \text{Proj}(A)$. Le foncteur $M \mapsto \widetilde M$
induit une équivalence
$$
\text{Mod}^{fg}_A/\text{Mod}^{fg}_{A, torsion}
\longrightarrow
\textit{Coh}(\mathcal{O}_X)
$$
dont un quasi-inverse est donné par
$\mathcal{F} \longmapsto \bigoplus_{n \geq 0} \Gamma(X, \mathcal{F}(n))$.
\end{proposition}

\begin{proof}
La sous-catégorie $\text{Mod}^{fg}_{A, torsion}$ est une sous-catégorie de Serre
de $\text{Mod}^{fg}_A$, voir
Homologie, définition \ref{homology-definition-serre-subcategory}.
Cela est clair d'après la description des objets donnée ci-dessus, mais résulte aussi
de Compléments d'algèbre, lemme \ref{more-algebra-lemma-I-power-torsion}.
La catégorie quotient à gauche de la flèche est donc définie
dans Homologie, lemme \ref{homology-lemma-serre-subcategory-is-kernel}.
Pour définir le foncteur de la proposition, il suffit de montrer que
le foncteur $M \mapsto \widetilde M$ envoie les modules de torsion sur $0$.
Cela est clair, car pour tout élément homogène $f \in A_+$, le
module $M_f$ est nul; la valeur $M_{(f)}$ de $\widetilde M$
sur $D_+(f)$ est donc elle aussi nulle.

\medskip\noindent
D'après le lemme \ref{coherent-lemma-coherent-on-proj}, le quasi-inverse proposé
est bien défini. En effet, le lemme montre que
$\mathcal{F} \longmapsto \bigoplus_{n \geq 0} \Gamma(X, \mathcal{F}(n))$
est un foncteur $\textit{Coh}(\mathcal{O}_X) \to \text{Mod}^{fg}_A$,
que l'on peut composer avec le foncteur quotient
$\text{Mod}^{fg}_A \to \text{Mod}^{fg}_A/\text{Mod}^{fg}_{A, torsion}$.

\medskip\noindent
D'après le lemme \ref{coherent-lemma-recover-tail-graded-module},
le composé de gauche à droite puis de droite à gauche est isomorphe au foncteur identité.

\medskip\noindent
Enfin, soit $\mathcal{F}$ un $\mathcal{O}_X$-module cohérent.
Posons $M = \bigoplus_{n \in \mathbf{Z}} \Gamma(X, \mathcal{F}(n))$,
considéré comme un $A$-module gradué, de sorte que notre foncteur envoie $\mathcal{F}$ sur
$M_{\geq 0} = \bigoplus_{n \geq 0} M_n$.
D'après Propriétés, lemme \ref{properties-lemma-proj-quasi-coherent},
l'application canonique $\widetilde M \to \mathcal{F}$
est un isomorphisme. Puisque l'inclusion
$M_{\geq 0} \to M$ définit un isomorphisme
$\widetilde{M_{\geq 0}} \to \widetilde M$, on en déduit
que le composé de droite à gauche puis de gauche à droite est lui aussi
isomorphe au foncteur identité.
\end{proof}





\section{Faisceaux cohérents sur Proj, II}
\label{coherent-section-coherent-proj-general}

\noindent
Dans cette section, nous étudions les faisceaux cohérents sur $\text{Proj}(A)$,
où $A$ est un anneau gradué noethérien. La plupart des résultats
se déduisent astucieusement du résultat correspondant de la section précédente,
où nous avons traité le cas où $A$ est engendré par des
éléments de degré $1$.

\begin{lemma}
\label{coherent-lemma-coherent-on-proj-general}
Soit $A$ un anneau gradué noethérien. Posons $X = \text{Proj}(A)$. Alors $X$
est un schéma noethérien. Soit $\mathcal{F}$ un $\mathcal{O}_X$-module cohérent.
\begin{enumerate}
\item Il existe $r \geq 0$, des entiers
$d_1, \ldots, d_r \in \mathbf{Z}$ et une surjection
$$
\bigoplus\nolimits_{j = 1, \ldots, r} \mathcal{O}_X(d_j)
\longrightarrow \mathcal{F}.
$$
\item Pour tout $p$, le groupe de cohomologie $H^p(X, \mathcal{F})$ est un
$A_0$-module de type fini.
\item Si $p > 0$, alors $H^p(X, \mathcal{F}(d)) = 0$ pour tout $d$ suffisamment grand.
\item Pour tout $k \in \mathbf{Z}$, le $A$-module gradué
$$
\bigoplus\nolimits_{d \geq k} H^0(X, \mathcal{F}(d))
$$
est un $A$-module de type fini.
\end{enumerate}
\end{lemma}

\begin{proof}
Nous ramenons l'énoncé au
lemme \ref{coherent-lemma-coherent-on-proj}.
D'après Algèbre, lemmes \ref{algebra-lemma-graded-Noetherian} et
\ref{algebra-lemma-S-plus-generated}, l'anneau $A_0$ est noethérien
et $A$ est engendré sur $A_0$ par un nombre fini d'éléments
$f_1, \ldots, f_r$ homogènes de degré strictement positif.
Soit $d$ un entier suffisamment divisible. Posons $A' = A^{(d)}$, avec
les notations d'Algèbre, section \ref{algebra-section-graded}.
Alors $A'$ est engendré sur $A'_0 = A_0$ par des éléments de
degré $1$ ; voir Algèbre, lemme \ref{algebra-lemma-uple-generated-degree-1}.
Le lemme \ref{coherent-lemma-coherent-on-proj} s'applique donc à $X' = \text{Proj}(A')$.

\medskip\noindent
Par Constructions, lemme \ref{constructions-lemma-d-uple}, il
existe un isomorphisme de schémas $i : X \to X'$ et des
isomorphismes $\mathcal{O}_X(nd) \to i^*\mathcal{O}_{X'}(n)$
compatibles avec le morphisme $A' \to A$ et les morphismes
$A_n \to H^0(X, \mathcal{O}_X(n))$ et $A'_n \to H^0(X', \mathcal{O}_{X'}(n))$.
Le lemme \ref{coherent-lemma-coherent-on-proj} implique donc que $X$ est noethérien et que
(1) et (2) sont satisfaits. Pour établir (3)
et (4), on utilise que, pour $k$, $p$ et $q$ fixés, on a
$$
\bigoplus\nolimits_{dn + q \geq k} H^p(X, \mathcal{F}(dn + q)) =
\bigoplus\nolimits_{dn + q \geq k} H^p(X', (i_*\mathcal{F}(q))(n)
$$
par les compatibilités précédentes. Si $p > 0$, le membre de droite
s'annule pour $k$ suffisamment grand en fonction de $q$, d'après le
lemme \ref{coherent-lemma-coherent-on-proj}. Comme il n'existe qu'un nombre fini
de classes de congruence d'entiers modulo $d$, on voit que (3) est satisfait pour
$\mathcal{F}$ sur $X$. Si $p = 0$, le membre de droite
est un $A'$-module de type fini d'après le lemme \ref{coherent-lemma-coherent-on-proj}. En
utilisant à nouveau la finitude des classes de congruence, on obtient que
$\bigoplus_{n \geq k} H^0(X, \mathcal{F}(n))$ est lui aussi un $A'$-module de type fini.
Comme la structure de $A'$-module provient de la structure de $A$-module (par
les compatibilités mentionnées plus haut), on conclut qu'il est également de type
fini comme $A$-module.
\end{proof}

\begin{lemma}
\label{coherent-lemma-recover-tail-graded-module-general}
Soit $A$ un anneau gradué noethérien et soit $d$ le ppcm des degrés d'un système de générateurs
de $A$ sur $A_0$. Soit $M$ un $A$-module gradué de
type fini. Posons $X = \text{Proj}(A)$ et soit $\widetilde{M}$
le $\mathcal{O}_X$-module quasi-cohérent sur $X$ associé à $M$.
Soit $k \in \mathbf{Z}$.
\begin{enumerate}
\item $N' = \bigoplus_{n \geq k} H^0(X, \widetilde{M(n)})$
est un $A$-module de type fini,
\item $N = \bigoplus_{n \geq k} H^0(X, \widetilde{M}(n))$
est un $A$-module de type fini,
\item il existe un morphisme canonique $N \to N'$,
\item si $k$ est suffisamment petit, il existe un morphisme canonique $M \to N'$,
\item le morphisme $M_n \to N'_n$ est un isomorphisme pour $n \gg 0$,
\item $N_n \to N'_n$ est un isomorphisme si $d | n$.
\end{enumerate}
\end{lemma}

\begin{proof}
Le morphisme $N \to N'$ de (3) provient de
Constructions, Équation (\ref{constructions-equation-multiply-more-generally}),
en prenant les sections globales.

\medskip\noindent
Par
Constructions, Équation
(\ref{constructions-equation-global-sections-more-generally}),
il existe un morphisme de $A$-modules gradués
$M \to \bigoplus_{n \in \mathbf{Z}} H^0(X, \widetilde{M(n)})$.
Si les générateurs de $M$ sont de degrés $\geq k$, alors l'image est
contenue dans le sous-module
$N' \subset \bigoplus_{n \in \mathbf{Z}} H^0(X, \widetilde{M(n)})$,
ce qui donne le morphisme de (4).

\medskip\noindent
Par Algèbre, lemmes \ref{algebra-lemma-graded-Noetherian} et
\ref{algebra-lemma-S-plus-generated}, l'anneau $A_0$ est noethérien
et $A$ est engendré sur $A_0$ par un nombre fini d'éléments
$f_1, \ldots, f_r$ homogènes de degré strictement positif.
Posons $d = \text{lcm}(\deg(f_i))$. Alors (6) résulte, par
exemple, de Constructions,
lemme \ref{constructions-lemma-where-invertible}.

\medskip\noindent
Comme $M$ est un $A$-module de type fini,
$\widetilde{M}$ est un $\mathcal{O}_X$-module de
type fini, c'est-à-dire un $\mathcal{O}_X$-module
cohérent. L'assertion (2) résulte donc du lemme \ref{coherent-lemma-coherent-on-proj-general}.

\medskip\noindent
Nous déduisons (1) de (2) par l'astuce suivante. Pour $q \in \{0, \ldots, d - 1\}$,
posons
$$
{}^qN = \bigoplus\nolimits_{n + q \geq k} H^0(X, \widetilde{M(q)}(n))
$$
Ce sont des $A$-modules de type fini par (2). L'anneau noethérien $A$
est fini sur $A^{(d)} = \bigoplus_{n \geq 0} A_{dn}$, car
il est engendré par les $f_i$ sur $A^{(d)}$ et $f_i^d \in A^{(d)}$. Par
conséquent, ${}^qN$ est un $A^{(d)}$-module de type
fini. De plus, $A^{(d)}$ est noethérien (cela résulte
d'Algèbre, lemme \ref{algebra-lemma-dehomogenize-finite-type}).
Il s'ensuit que le sous-$A^{(d)}$-module
${}^qN^{(d)} = \bigoplus_{n \in \mathbf{Z}} {}^qN_{dn}$
est un module de type fini sur $A^{(d)}$. À l'aide des isomorphismes
$\widetilde{M(dn + q)} = \widetilde{M(q)}(dn)$, on peut écrire
$$
N' =
\bigoplus\nolimits_{q \in \{0, \ldots, d - 1\}}
\bigoplus\nolimits_{dn + q \geq k}
H^0(X, \widetilde{M(q)}(dn)) =
\bigoplus\nolimits_{q \in \{0, \ldots, d - 1\}} {}^qN^{(d)}
$$
Ainsi, $N'$ est de type fini sur $A^{(d)}$ et, a fortiori, de type fini sur $A$.
Cela démontre (1).

\medskip\noindent
Soit $K$ un $A$-module de type fini tel que $\widetilde{K} = 0$. Nous affirmons
que $K_n = 0$ si $d|n$ et $n \gg 0$. En raisonnant comme ci-dessus, on voit que
$K^{(d)}$ est un $A^{(d)}$-module de type fini. Soient
$x_1, \ldots, x_m \in K$ des générateurs homogènes de $K^{(d)}$
sur $A^{(d)}$, de degrés respectifs $d_1, \ldots, d_m$, avec $d | d_j$.
Pour tout $i$ et tout $j$, il existe $n_{ij} \geq 0$ tel que
$f_i^{n_{ij}} x_j = 0$ dans $K_{d_j + n_{ij}}$ : sinon,
$x_j/f_i^{d_i/\deg(f_i)} \in K_{(f_i)}$ serait non nul, c'est-à-dire que
$\widetilde{K}$ serait non nul. Nous utilisons ici $\deg(f_i) | d | d_j$
pour tous $i, j$. On conclut que $K_n$ est nul pour tout
$n$ tel que $d | n$ et $n > \max_j (d_j + \sum_i n_{ij} \deg(f_i))$,
car tout élément de $K_n$ est une somme de termes, chacun étant le produit d'un
monôme en les $f_i$ par l'un des $x_j$, de degré total $n$.

\medskip\noindent
Pour achever la démonstration, il reste à montrer que $M \to N'$ est un isomorphisme
en tout degré suffisamment grand. Le morphisme $N \to N'$ induit un
isomorphisme $\widetilde{N} \to \widetilde{N'}$, car les modules correspondants sont isomorphes sur les
ouverts affines $D_+(f_i) = D_+(f_i^d)$ :
$N_{(f_i)} \cong N_{(f_i^d)} \cong N'_{(f_i^d)} \cong N'_{(f_i)}$
par la propriété (6).
D'après Propriétés, lemme \ref{properties-lemma-proj-quasi-coherent},
on a un isomorphisme canonique $\widetilde{N} \to \widetilde{M}$.
Le composé $\widetilde{N} \to \widetilde{M} \to \widetilde{N'}$
est l'isomorphisme précédent (démonstration omise ; indication : vérifier
sur les ouverts affines standards). Ainsi, le morphisme $M \to N'$ induit un isomorphisme
$\widetilde{M} \to \widetilde{N'}$.
Posons $K = \Ker(M \to N')$ et $Q = \Coker(M \to N')$.
Rappelons que le foncteur
$M \mapsto \widetilde{M}$ est exact ; voir
Constructions, lemme \ref{constructions-lemma-proj-sheaves}.
On a donc $\widetilde{K} = 0$ et $\widetilde{Q} = 0$.
D'après le paragraphe précédent, $K_n = 0$ et
$Q_n = 0$ si $d | n$ et $n \gg 0$. On voit ici l'avantage
de l'emploi de $N'$ plutôt que de $N$ : le foncteur $M \leadsto N'$
est compatible aux décalages (cela résulte immédiatement de la construction).
En répétant tout l'argument après avoir remplacé $M$ par $M(q)$,
on obtient $K_n = 0$ et $Q_n = 0$ si $n \equiv q \bmod d$
et $n \gg 0$. Comme il n'existe qu'un
nombre fini de classes de congruence modulo $n$, la démonstration est achevée.
\end{proof}

\noindent
Soit $A$ un anneau gradué noethérien. Rappelons que
$A_+ = \bigoplus_{n > 0} A_n$ est l'idéal non pertinent.
D'après Algèbre, lemmes \ref{algebra-lemma-graded-Noetherian} et
\ref{algebra-lemma-S-plus-generated}, l'anneau $A_0$ est noethérien
et $A$ est engendré sur $A_0$ par un nombre fini d'éléments
$f_1, \ldots, f_r$ homogènes de degré strictement positif.
Posons $d = \text{lcm}(\deg(f_i))$. Soit $M$ un $A$-module
gradué. Dans cette situation, nous dirons qu'un élément homogène $x \in M$
est {\it non pertinent}\footnote{Cette terminologie n'est pas standard.} si
$$
(A_+ x)_{nd} = 0\text{ pour tout }n \gg 0
$$
Si $x \in M$ est homogène et non pertinent et si $f \in A$
est homogène, alors $fx$ est encore non pertinent.
Par conséquent, les éléments non pertinents engendrent un sous-module gradué
$M_{irrelevant} \subset M$. Nous dirons que $M$ est {\it non pertinent}
si tout élément homogène de $M$ est non pertinent, autrement
dit si $M_{irrelevant} = M$.
Lorsque $M$ est de type fini,
cela équivaut à $M_{nd} = 0$ pour $n \gg 0$.
Notons $\text{Mod}_A$ la catégorie des $A$-modules gradués,
$\text{Mod}^{fg}_A$ sa sous-catégorie pleine formée de ceux qui sont de type
fini et $\text{Mod}^{fg}_{A, irrelevant}$ la sous-catégorie pleine des modules
non pertinents.

\begin{proposition}
\label{coherent-proposition-coherent-modules-on-proj-general}
Soit $A$ un anneau gradué noethérien. Posons $X = \text{Proj}(A)$. Le foncteur
$M \mapsto \widetilde M$ induit une équivalence
$$
\text{Mod}^{fg}_A/\text{Mod}^{fg}_{A, irrelevant}
\longrightarrow
\textit{Coh}(\mathcal{O}_X)
$$
dont un quasi-inverse est donné par
$\mathcal{F} \longmapsto \bigoplus_{n \geq 0} \Gamma(X, \mathcal{F}(n))$.
\end{proposition}

\begin{proof}
Nous recommandons au lecteur de commencer par lire la démonstration dans le cas où $A$
est engendré en degré $1$ ; voir la
proposition \ref{coherent-proposition-coherent-modules-on-proj}.
Soient $f_1, \ldots, f_r \in A$ des éléments homogènes de degré strictement positif
qui engendrent $A$ sur $A_0$. Soit $d$ le ppcm des degrés
$d_i$ des $f_i$. Soit $M$ un $A$-module
de type fini. Montrons que $\widetilde{M}$ est nul si et
seulement si $M$ est un $A$-module gradué non pertinent (au
sens défini avant l'énoncé de la proposition). Soit donc $x \in M$ un élément
homogène.
Choisissons $k \in \mathbf{Z}$ assez petit, et soient
$N \to N'$ et $M \to N'$ comme dans le
lemme \ref{coherent-lemma-recover-tail-graded-module-general}.
On peut aussi choisir $l$ assez grand pour que
$M_n \to N_n$ soit un isomorphisme pour $n \geq l$.
Si $\widetilde{M}$ est nul, alors $N = 0$.
Ainsi, pour tout $f \in A_+$ homogène tel que
$\deg(f) + \deg(x) = nd$ et $nd > l$, l'élément $fx$ est nul,
car $N_{nd} \to N'_{nd}$ et $M_{nd} \to N'_{nd}$ sont des isomorphismes.
Donc $x$ est non
pertinent. Réciproquement, supposons $M$ non pertinent. Alors $M_{nd}$ est nul
pour $n \gg 0$ (voir la discussion précédant la
proposition). Cela entraîne manifestement que $M_{(f_i)} = M_{(f_i^{d/\deg(f_i)})} = 0$,
d'où $\widetilde{M} = 0$ par construction.

\medskip\noindent
Il en résulte que la sous-catégorie $\text{Mod}^{fg}_{A, irrelevant}$
est une sous-catégorie de Serre de $\text{Mod}^{fg}_A$, puisqu'elle est
le noyau du foncteur exact $M \mapsto \widetilde M$ ; voir
Homologie, lemme \ref{homology-lemma-kernel-exact-functor},
et Constructions, lemme \ref{constructions-lemma-proj-sheaves}.
La catégorie quotient figurant à gauche de la flèche est donc celle définie
dans Homologie, lemme \ref{homology-lemma-serre-subcategory-is-kernel}.
Pour définir le foncteur de la proposition, il suffit de montrer que
le foncteur $M \mapsto \widetilde M$ envoie les modules non pertinents sur $0$,
ce que nous avons établi ci-dessus.

\medskip\noindent
Le quasi-inverse proposé est bien défini d'après le lemme \ref{coherent-lemma-coherent-on-proj-general}. En
effet, ce lemme montre que
$\mathcal{F} \longmapsto \bigoplus_{n \geq 0} \Gamma(X, \mathcal{F}(n))$
est un foncteur $\textit{Coh}(\mathcal{O}_X) \to \text{Mod}^{fg}_A$
que l'on peut composer avec le foncteur quotient
$\text{Mod}^{fg}_A \to \text{Mod}^{fg}_A/\text{Mod}^{fg}_{A, irrelevant}$.

\medskip\noindent
D'après le lemme \ref{coherent-lemma-recover-tail-graded-module-general},
le composé de gauche à droite puis de droite à gauche est isomorphe au foncteur identité.
En effet, soit $M$ un $A$-module gradué de type fini. Choisissons
$k \in \mathbf{Z}$ assez petit, et soient
$N \to N'$ et $M \to N'$ comme dans le
lemme \ref{coherent-lemma-recover-tail-graded-module-general}.
Le noyau et le conoyau de $M \to N'$ ne sont alors non
nuls que dans un nombre fini de degrés, et sont donc non pertinents.
De plus, le noyau et le conoyau de $N \to N'$ sont nuls en tout degré
assez grand divisible par $d$ ; ce sont donc eux aussi des modules non pertinents.
Ainsi $M \to N'$ et $N \to N'$ sont tous deux des isomorphismes dans la
catégorie quotient, comme voulu.

\medskip\noindent
Enfin, soit $\mathcal{F}$ un $\mathcal{O}_X$-module cohérent.
Posons $M = \bigoplus_{n \in \mathbf{Z}} \Gamma(X, \mathcal{F}(n))$
et considérons-le comme un $A$-module gradué, de sorte que notre foncteur envoie $\mathcal{F}$ sur
$M_{\geq 0} = \bigoplus_{n \geq 0} M_n$.
D'après Propriétés, lemme \ref{properties-lemma-proj-quasi-coherent},
le morphisme canonique $\widetilde M \to \mathcal{F}$
est un isomorphisme. Comme l'inclusion
$M_{\geq 0} \to M$ définit un isomorphisme
$\widetilde{M_{\geq 0}} \to \widetilde M$, on en conclut
que le composé de droite à gauche puis de gauche à droite est lui aussi
isomorphe au foncteur identité.
\end{proof}










\section{Images directes supérieures par les morphismes projectifs}
\label{coherent-section-projective-pushforward}

\noindent
Nous énonçons et démontrons d'abord un résultat lorsque la base est affine, puis
nous en déduisons quelques résultats sur les morphismes projectifs.

\begin{lemma}
\label{coherent-lemma-coherent-proper-ample}
Soit $R$ un anneau noethérien. Soit $X \to \Spec(R)$ un morphisme propre.
Soit $\mathcal{L}$ un faisceau inversible ample sur $X$. Soit $\mathcal{F}$
un $\mathcal{O}_X$-module cohérent.
\begin{enumerate}
\item L'anneau gradué
$A = \bigoplus_{d \geq 0} H^0(X, \mathcal{L}^{\otimes d})$
est une $R$-algèbre de type fini.
\item Il existe $r \geq 0$, des entiers
$d_1, \ldots, d_r \in \mathbf{Z}$ et une surjection
$$
\bigoplus\nolimits_{j = 1, \ldots, r} \mathcal{L}^{\otimes d_j}
\longrightarrow \mathcal{F}.
$$
\item Pour tout $p$, le groupe de cohomologie $H^p(X, \mathcal{F})$ est un
$R$-module de type fini.
\item Si $p > 0$, alors
$H^p(X, \mathcal{F} \otimes_{\mathcal{O}_X} \mathcal{L}^{\otimes d}) = 0$
pour tout $d$ assez grand.
\item Pour tout $k \in \mathbf{Z}$, le $A$-module gradué
$$
\bigoplus\nolimits_{d \geq k}
H^0(X, \mathcal{F} \otimes_{\mathcal{O}_X} \mathcal{L}^{\otimes d})
$$
est un $A$-module de type fini.
\end{enumerate}
\end{lemma}

\begin{proof}
D'après
Morphismes, lemme \ref{morphisms-lemma-finite-type-over-affine-ample-very-ample},
il existe $d > 0$ et une immersion $i : X \to \mathbf{P}^n_R$
telles que $\mathcal{L}^{\otimes d} \cong i^*\mathcal{O}_{\mathbf{P}^n_R}(1)$.
Puisque $X$ est propre sur $R$, le morphisme $i$ est une immersion fermée
(Morphismes, lemme \ref{morphisms-lemma-image-proper-scheme-closed}).
On a donc $H^i(X, \mathcal{G}) = H^i(\mathbf{P}^n_R, i_*\mathcal{G})$
pour tout faisceau quasi-cohérent $\mathcal{G}$ sur $X$
(d'après le lemme \ref{coherent-lemma-relative-affine-cohomology} et le fait que les
immersions fermées sont affines ; voir
Morphismes, lemme \ref{morphisms-lemma-closed-immersion-affine}).
De plus, si $\mathcal{G}$ est cohérent, alors $i_*\mathcal{G}$
l'est aussi (lemme \ref{coherent-lemma-i-star-equivalence}).
Nous emploierons désormais ces faits sans autre mention.

\medskip\noindent
Démonstration de (1). Posons $S = R[T_0, \ldots, T_n]$, de sorte que
$\mathbf{P}^n_R = \text{Proj}(S)$.
Remarquons que $A$ est une $S$-algèbre (mais le morphisme d'anneaux $S \to A$ n'est pas
un morphisme d'anneaux gradués, puisque $S_n$ s'envoie dans $A_{dn}$).
La formule de projection
(Cohomologie, lemme \ref{cohomology-lemma-projection-formula})
donne
$$
i_*(\mathcal{L}^{\otimes nd + q}) =
i_*(\mathcal{L}^{\otimes q})
\otimes_{\mathcal{O}_{\mathbf{P}^n_R}}
\mathcal{O}_{\mathbf{P}^n_R}(n)
$$
pour tout $n \in \mathbf{Z}$. On en conclut que $\bigoplus_{n \geq 0} A_{nd + q}$
est un $S$-module gradué de type fini d'après le lemme \ref{coherent-lemma-coherent-projective}.
Comme
$A = \bigoplus_{q \in \{0, \ldots, d - 1} \bigoplus_{n \geq 0} A_{nd + q}$
on voit que $A$ est une $S$-algèbre finie, ce qui prouve (1).

\medskip\noindent
Démonstration de (2). Cela résulte de
Propriétés, proposition \ref{properties-proposition-characterize-ample}.

\medskip\noindent
Démonstration de (3). On applique le lemme \ref{coherent-lemma-coherent-projective}
et l'égalité $H^p(X, \mathcal{F}) = H^p(\mathbf{P}^n_R, i_*\mathcal{F})$.

\medskip\noindent
Démonstration de (4). Fixons $p > 0$. La formule de projection donne
$$
i_*(\mathcal{F} \otimes_{\mathcal{O}_X} \mathcal{L}^{\otimes nd + q}) =
i_*(\mathcal{F} \otimes_{\mathcal{O}_X} \mathcal{L}^{\otimes q})
\otimes_{\mathcal{O}_{\mathbf{P}^n_R}}
\mathcal{O}_{\mathbf{P}^n_R}(n)
$$
pour tout $n \in \mathbf{Z}$. D'après le lemme \ref{coherent-lemma-coherent-projective},
on a donc $H^p(X, \mathcal{F} \otimes \mathcal{L}^{nd + q}) = 0$
pour $n \gg 0$. Comme il n'existe qu'un nombre fini de classes de
congruence d'entiers modulo $d$, ceci prouve (4).

\medskip\noindent
Démonstration de (5). Fixons un entier $k$. Posons
$M = \bigoplus_{n \geq k} H^0(X, \mathcal{F} \otimes \mathcal{L}^{\otimes n})$.
En raisonnant comme ci-dessus, on voit que $\bigoplus_{nd + q \geq k} M_{nd + q}$
est un $S$-module gradué de type fini. Comme
$M = \bigoplus_{q \in \{0, \ldots, d - 1\}}
\bigoplus_{nd + q \geq k} M_{nd + q}$
on voit que $M$ est un $S$-module de type fini. Puisque sa structure de $S$-module se
factorise par le morphisme d'anneaux $S \to A$, on conclut que $M$ est
un $A$-module de type fini.
\end{proof}

\begin{lemma}
\label{coherent-lemma-kill-by-twisting}
Soit $f : X \to S$ un morphisme de schémas.
Soit $\mathcal{F}$ un $\mathcal{O}_X$-module quasi-cohérent.
Soit $\mathcal{L}$ un faisceau inversible sur $X$.
Supposons que
\begin{enumerate}
\item $S$ soit noethérien,
\item $f$ soit propre,
\item $\mathcal{F}$ soit cohérent, et
\item $\mathcal{L}$ soit relativement ample sur $X/S$.
\end{enumerate}
Il existe alors $n_0$ tel que, pour tout $n \geq n_0$,
on ait
$$
R^pf_*\left(\mathcal{F} \otimes_{\mathcal{O}_X} \mathcal{L}^{\otimes n}\right)
=
0
$$
pour tout $p > 0$.
\end{lemma}

\begin{proof}
Choisissons un recouvrement ouvert affine fini $S = \bigcup V_j$ et
posons $X_j = f^{-1}(V_j)$.
Il suffit manifestement de résoudre la question pour chacun des systèmes, en nombre
fini, $(X_j \to V_j, \mathcal{L}|_{X_j}, \mathcal{F}|_{V_j})$.
On peut donc supposer $S$ affine. Dans
ce cas, l'annulation de
$R^pf_*(\mathcal{F} \otimes \mathcal{L}^{\otimes n})$
équivaut à celle de
$H^p(X, \mathcal{F} \otimes \mathcal{L}^{\otimes n})$ ; voir le
lemme \ref{coherent-lemma-quasi-coherence-higher-direct-images-application}.
L'annulation voulue résulte donc
du lemme \ref{coherent-lemma-coherent-proper-ample} (qui s'applique parce que
$\mathcal{L}$ est ample sur $X$ d'après Morphismes, Lemme
\ref{morphisms-lemma-finite-type-over-affine-ample-very-ample}).
\end{proof}

\begin{lemma}
\label{coherent-lemma-locally-projective-pushforward}
Soit $S$ un schéma localement noethérien.
Soit $f : X \to S$ un morphisme localement projectif.
Soit $\mathcal{F}$ un $\mathcal{O}_X$-module cohérent.
Alors $R^if_*\mathcal{F}$ est un $\mathcal{O}_S$-module
cohérent pour tout $i \geq 0$.
\end{lemma}

\begin{proof}
Remarquons d'abord qu'un morphisme localement projectif est propre
(Morphismes, lemme \ref{morphisms-lemma-locally-projective-proper})
et donc de type fini.
En particulier, $X$ est localement noethérien
(Morphismes, lemme \ref{morphisms-lemma-finite-type-noetherian}),
de sorte que l'énoncé a un sens. De
plus, d'après le lemme \ref{coherent-lemma-quasi-coherence-higher-direct-images},
les faisceaux $R^pf_*\mathcal{F}$ sont quasi-cohérents.

\medskip\noindent
Cela étant, l'énoncé est local sur $S$ (par exemple d'après
Cohomologie, lemme \ref{cohomology-lemma-localize-higher-direct-images}).
On peut donc supposer que $S = \Spec(R)$ est le spectre d'un
anneau noethérien et que $X$ est un sous-schéma fermé de
$\mathbf{P}^n_R$ pour un certain $n$ ; voir
Morphismes, lemme \ref{morphisms-lemma-characterize-locally-projective}.
Dans ce cas, les faisceaux $R^pf_*\mathcal{F}$ sont les faisceaux
quasi-cohérents associés aux $R$-modules $H^p(X, \mathcal{F})$ ; voir le
lemme \ref{coherent-lemma-quasi-coherence-higher-direct-images-application}.
Il suffit donc de montrer que les $R$-modules $H^p(X, \mathcal{F})$
sont des $R$-modules de type fini (lemme \ref{coherent-lemma-coherent-Noetherian}).
Cela résulte du lemme \ref{coherent-lemma-coherent-proper-ample},
car la restriction de $\mathcal{O}_{\mathbf{P}^n_R}(1)$
à $X$ est ample sur $X$.
\end{proof}








\section{Faisceaux inversibles amples et cohomologie}
\label{coherent-section-ample-cohomology}

\noindent
Voici un critère d'amplitude sur les schémas propres sur une base affine,
exprimé par l'annulation de la cohomologie des faisceaux tordus.

\begin{lemma}
\label{coherent-lemma-vanshing-gives-ample}
\begin{reference}
\cite[III Proposition 2.6.1]{EGA}
\end{reference}
Soit $R$ un anneau noethérien. Soit $f : X \to \Spec(R)$ un morphisme propre.
Soit $\mathcal{L}$ un $\mathcal{O}_X$-module inversible.
Les conditions suivantes sont équivalentes :
\begin{enumerate}
\item $\mathcal{L}$ est ample sur $X$ (cette condition admet bien d'autres formes
équivalentes, voir
Propriétés, proposition \ref{properties-proposition-characterize-ample} et
Morphismes, Lemme
\ref{morphisms-lemma-finite-type-over-affine-ample-very-ample}),
\item pour tout $\mathcal{O}_X$-module cohérent $\mathcal{F}$, il
existe $n_0 \geq 0$ tel que
$H^p(X, \mathcal{F} \otimes \mathcal{L}^{\otimes n}) = 0$ pour tout $n \geq n_0$
et tout $p > 0$, et
\item pour tout faisceau quasi-cohérent d'idéaux
$\mathcal{I} \subset \mathcal{O}_X$, il existe $n \geq 1$
tel que $H^1(X, \mathcal{I} \otimes \mathcal{L}^{\otimes n}) = 0$.
\end{enumerate}
\end{lemma}

\begin{proof}
L'implication (1) $\Rightarrow$ (2) résulte du
lemme \ref{coherent-lemma-coherent-proper-ample}.
L'implication (2) $\Rightarrow$ (3) est
triviale. L'implication (3) $\Rightarrow$ (1) est le
lemme \ref{coherent-lemma-quasi-compact-h1-zero-invertible}.
\end{proof}

\begin{lemma}
\label{coherent-lemma-surjective-finite-morphism-ample}
Soit $R$ un anneau noethérien. Soit $f : Y \to X$ un morphisme de
schémas propres sur $R$. Soit $\mathcal{L}$
un $\mathcal{O}_X$-module inversible. Supposons $f$ fini et surjectif.
Alors $\mathcal{L}$ est ample si et seulement si $f^*\mathcal{L}$ est ample.
\end{lemma}

\begin{proof}
L'image réciproque d'un faisceau inversible ample par un morphisme quasi-affine
est ample, voir Morphismes, Lemme
\ref{morphisms-lemma-pullback-ample-tensor-relatively-ample}.
Cela démontre l'une des implications, car un morphisme fini est affine
par définition.

\medskip\noindent
Supposons $f^*\mathcal{L}$ ample. Soit $P$ la propriété suivante
des $\mathcal{O}_X$-modules cohérents $\mathcal{F}$ : il existe $n_0$
tel que $H^p(X, \mathcal{F} \otimes \mathcal{L}^{\otimes n}) = 0$
pour tout $n \geq n_0$ et tout $p > 0$. Nous allons montrer que $P$ est
vérifiée par tout $\mathcal{O}_X$-module cohérent $\mathcal{F}$, ce qui entraînera que
$\mathcal{L}$ est ample d'après le lemme \ref{coherent-lemma-vanshing-gives-ample}.
Nous allons appliquer le lemme \ref{coherent-lemma-property-higher-rank-cohomological}.
Il faut donc vérifier les conditions (1), (2) et (3) de ce lemme pour $P$.
La propriété (1) résulte de la suite exacte longue de cohomologie associée
à une suite exacte courte de faisceaux et du fait que le produit tensoriel avec un
faisceau inversible est un foncteur exact. La propriété (2) résulte de ce que
$H^p(X, -)$ est un foncteur additif. Pour vérifier (3), soit $Z \subset X$ un
sous-schéma fermé intègre de point générique $\xi$.
Soit $\mathcal{F}$ un faisceau cohérent sur $Y$ tel que
le support de $f_*\mathcal{F}$ soit égal à $Z$
et que $(f_*\mathcal{F})_\xi$ soit annulé par $\mathfrak m_\xi$,
voir lemme \ref{coherent-lemma-finite-morphism-Noetherian}. Nous affirmons
que $\mathcal{G} = f_*\mathcal{F}$ convient. Il suffit de vérifier la
condition (3)(c) du lemme \ref{coherent-lemma-property-higher-rank-cohomological}.
Supposons donc que $\mathcal{J} \subset \mathcal{O}_X$ soit un
faisceau quasi-cohérent d'idéaux tel que
$\mathcal{J}_\xi = \mathcal{O}_{X, \xi}$.
Un morphisme fini est affine; d'après le
lemme \ref{coherent-lemma-affine-morphism-projection-ideal}, on a donc
$\mathcal{J}\mathcal{G} = f_*(f^{-1}\mathcal{J}\mathcal{F})$.
De plus, comme on l'a remarqué dans la démonstration du
lemme \ref{coherent-lemma-affine-morphism-projection-ideal}, le faisceau
$f^{-1}\mathcal{J}\mathcal{F}$ est un $\mathcal{O}_Y$-module cohérent.
Puisque $\mathcal{L}$ est ample, le lemme \ref{coherent-lemma-vanshing-gives-ample}
assure l'existence de $n_0$ tel que
$$
H^p(Y, f^{-1}\mathcal{J}\mathcal{F}
\otimes_{\mathcal{O}_Y} f^*\mathcal{L}^{\otimes n}) = 0,
$$
pour $n \geq n_0$ et $p > 0$. Puisque $f$ est fini, donc affine, on a
\begin{align*}
H^p(X, \mathcal{J}\mathcal{G} \otimes_{\mathcal{O}_X}
\mathcal{L}^{\otimes n})
& =
H^p(X, f_*(f^{-1}\mathcal{J}\mathcal{F}) \otimes_{\mathcal{O}_X}
\mathcal{L}^{\otimes n}) \\
& =
H^p(X, f_*(f^{-1}\mathcal{J}\mathcal{F} \otimes_{\mathcal{O}_Y}
f^*\mathcal{L}^{\otimes n})) \\
& =
H^p(Y, f^{-1}\mathcal{J}\mathcal{F} \otimes_{\mathcal{O}_Y}
f^*\mathcal{L}^{\otimes n}) = 0
\end{align*}
On a utilisé ici la formule de projection
(Cohomologie, lemme \ref{cohomology-lemma-projection-formula}) et le
lemme \ref{coherent-lemma-relative-affine-cohomology}.
Le sous-faisceau quasi-cohérent $\mathcal{G}' = \mathcal{J}\mathcal{G}$
vérifie donc $P$. Cela établit la propriété (3)(c) du
lemme \ref{coherent-lemma-property-higher-rank-cohomological}, comme voulu.
\end{proof}

\noindent
La cohomologie est fonctorielle. En particulier, étant donnés un espace annelé $X$,
un $\mathcal{O}_X$-module inversible $\mathcal{L}$ et une section
$s \in \Gamma(X, \mathcal{L})$, on obtient des applications
$$
H^p(X, \mathcal{F})
\longrightarrow
H^p(X, \mathcal{F} \otimes_{\mathcal{O}_X} \mathcal{L}), \quad
\xi \longmapsto s\xi
$$
induites par l'application
$\mathcal{F} \to \mathcal{F} \otimes_{\mathcal{O}_X} \mathcal{L}$
de multiplication par $s$. On pose
$\Gamma_*(X, \mathcal{L}) =
\bigoplus_{n \geq 0} \Gamma(X, \mathcal{L}^{\otimes n})$,
considéré comme un anneau gradué,
voir Modules, définition \ref{modules-definition-gamma-star}.
Étant donnés un faisceau de $\mathcal{O}_X$-modules $\mathcal{F}$ et un entier
$p \geq 0$, on pose
$$
H^p_*(X, \mathcal{L}, \mathcal{F}) =
\bigoplus\nolimits_{n \in \mathbf{Z}} H^p(X,
\mathcal{F} \otimes_{\mathcal{O}_X} \mathcal{L}^{\otimes n})
$$
La multiplication définie ci-dessus en fait un $\Gamma_*(X, \mathcal{L})$-module gradué. Attention:
la notation $H^p_*(X, \mathcal{L}, \mathcal{F})$
n'est pas usuelle.

\begin{lemma}
\label{coherent-lemma-invert-s-cohomology}
Soit $X$ un schéma. Soit $\mathcal{L}$ un faisceau inversible sur $X$.
Soit $s \in \Gamma(X, \mathcal{L})$.
Soit $\mathcal{F}$ un $\mathcal{O}_X$-module quasi-cohérent.
Si $X$ est quasi-compact et quasi-séparé, l'application canonique
$$
H^p_*(X, \mathcal{L}, \mathcal{F})_{(s)}
\longrightarrow
H^p(X_s, \mathcal{F})
$$
qui envoie $\xi/s^n$ sur $s^{-n}\xi$ est un isomorphisme.
\end{lemma}

\begin{proof}
Pour $p = 0$, il s'agit
de Propriétés, lemme \ref{properties-lemma-invert-s-sections}.
Nous allons démontrer l'assertion au moyen du principe de
récurrence (lemme \ref{coherent-lemma-induction-principle}), en notant, pour tout ouvert quasi-compact
$U \subset X$, $P(U)$ la propriété :
pour tout $p \geq 0$, l'application
$$
H^p_*(U, \mathcal{L}, \mathcal{F})_{(s)}
\longrightarrow
H^p(U_s, \mathcal{F})
$$
est un isomorphisme.

\medskip\noindent
Si $U$ est affine, les deux membres de la flèche ci-dessus
sont nuls pour $p > 0$, d'après
le lemme \ref{coherent-lemma-quasi-coherent-affine-cohomology-zero}
et
Propriétés, lemme \ref{properties-lemma-affine-cap-s-open},
et l'assertion est vraie. Si $P$ est vraie pour $U$, $V$ et $U \cap V$,
on peut utiliser les suites de Mayer-Vietoris
(Cohomologie, lemme \ref{cohomology-lemma-mayer-vietoris}) pour obtenir
un morphisme de suites exactes longues
$$
\xymatrix{
H^{p - 1}_*(U \cap V, \mathcal{L}, \mathcal{F})_{(s)} \ar[r] \ar[d] &
H^p_*(U \cup V, \mathcal{L}, \mathcal{F})_{(s)} \ar[r] \ar[d] &
H^p_*(U, \mathcal{L}, \mathcal{F})_{(s)}
\oplus
H^p_*(V, \mathcal{L}, \mathcal{F})_{(s)} \ar[d] \\
H^{p - 1}(U_s \cap V_s, \mathcal{F}) \ar[r]&
H^p(U_s \cup V_s, \mathcal{F}) \ar[r] &
H^p(U_s, \mathcal{F})
\oplus
H^p(V_s, \mathcal{F})
}
$$
(dont on ne représente qu'un fragment). Notons que $U_s \cap V_s = (U \cap V)_s$ et
que $U_s \cup V_s = (U \cup V)_s$. Les applications verticales de gauche et
de droite sont donc des isomorphismes (ainsi qu'une application supplémentaire à droite et une
à gauche, qui ne figurent pas dans le diagramme). On
en déduit que $P(U \cup V)$ est vérifiée, d'après
le lemme des cinq (Homologie, lemme \ref{homology-lemma-five-lemma}).
Cela achève la démonstration.
\end{proof}

\begin{lemma}
\label{coherent-lemma-section-affine-open-kills-classes}
Soit $X$ un schéma.
Soit $\mathcal{L}$ un $\mathcal{O}_X$-module inversible.
Soit $s \in \Gamma(X, \mathcal{L})$ une section.
Supposons que
\begin{enumerate}
\item $X$ soit quasi-compact et quasi-séparé, et
\item $X_s$ soit affine.
\end{enumerate}
Alors, pour tout $\mathcal{O}_X$-module quasi-cohérent $\mathcal{F}$,
tout $p > 0$ et tout $\xi \in H^p(X, \mathcal{F})$, il
existe $n \geq 0$ tel que $s^n\xi = 0$ dans
$H^p(X, \mathcal{F} \otimes_{\mathcal{O}_X} \mathcal{L}^{\otimes n})$.
\end{lemma}

\begin{proof}
Rappelons que $H^p(X_s, \mathcal{G})$ est nul pour tout module
quasi-cohérent $\mathcal{G}$, d'après
le lemme \ref{coherent-lemma-quasi-coherent-affine-cohomology-zero}.
Le lemme résulte donc du
lemme \ref{coherent-lemma-invert-s-cohomology}.
\end{proof}

\noindent
Pour une version plus générale du lemme suivant, voir
Limites, lemme \ref{limits-lemma-ample-on-reduction}.

\begin{lemma}
\label{coherent-lemma-ample-on-reduction}
Soit $i : Z \to X$ une immersion fermée de schémas noethériens induisant
un homéomorphisme des espaces topologiques sous-jacents.
Soit $\mathcal{L}$ un faisceau inversible sur $X$.
Alors $i^*\mathcal{L}$ est ample sur $Z$ si et seulement si
$\mathcal{L}$ est ample sur $X$.
\end{lemma}

\begin{proof}
Si $\mathcal{L}$ est ample, alors $i^*\mathcal{L}$ est ample, par
exemple d'après Morphismes, Lemme
\ref{morphisms-lemma-pullback-ample-tensor-relatively-ample}.
Supposons $i^*\mathcal{L}$ ample. Il faut montrer que $\mathcal{L}$
est ample sur $X$.
Soit $\mathcal{I} \subset \mathcal{O}_X$ le faisceau cohérent d'idéaux
définissant le sous-schéma fermé $Z$. Puisque $i(Z) = X$ ensemblistement,
on a $\mathcal{I}^n = 0$ pour un certain $n$, d'après
le lemme \ref{coherent-lemma-power-ideal-kills-sheaf}.
Considérons la suite
$$
X = Z_n \supset Z_{n - 1} \supset Z_{n - 2} \supset \ldots \supset Z_1 = Z
$$
de sous-schémas fermés définis par
$0 = \mathcal{I}^n \subset \mathcal{I}^{n - 1} \subset \ldots \subset
\mathcal{I}$. Chaque immersion fermée $Z_{i - 1} \to Z_{i}$
est alors définie par un faisceau cohérent d'idéaux de carré nul. On se
ramène ainsi au cas où $\mathcal{I}^2 = 0$.

\medskip\noindent
Considérons la suite exacte courte
$$
0 \to \mathcal{I} \to \mathcal{O}_X \to i_*\mathcal{O}_Z \to 0
$$
de $\mathcal{O}_X$-modules quasi-cohérents. En tensorisant par
$\mathcal{L}^{\otimes n}$, on obtient des suites exactes courtes
\begin{equation}
\label{coherent-equation-ses}
0 \to \mathcal{I} \otimes_{\mathcal{O}_X} \mathcal{L}^{\otimes n}
\to \mathcal{L}^{\otimes n} \to i_*i^*\mathcal{L}^{\otimes n} \to 0
\end{equation}
Puisque $\mathcal{I}^2 = 0$, on peut utiliser
Morphismes, lemme \ref{morphisms-lemma-i-star-equivalence},
pour considérer $\mathcal{I}$ comme un $\mathcal{O}_Z$-module quasi-cohérent;
on a alors $\mathcal{I} \otimes_{\mathcal{O}_X} \mathcal{L}^{\otimes n} =
\mathcal{I} \otimes_{\mathcal{O}_Z} i^*\mathcal{L}^{\otimes n}$,
avec un abus de notation
évident. De plus, la cohomologie de ce faisceau sur $Z$ s'identifie
canoniquement à sa cohomologie sur $X$ (puisque $i$ est un
homéomorphisme).

\medskip\noindent
Soit $x \in X$ un point, et notons $z \in Z$ le point correspondant.
Puisque $i^*\mathcal{L}$ est ample, il existe $n$ et une section
$s \in \Gamma(Z, i^*\mathcal{L}^{\otimes n})$ tels que $z \in Z_s$
et que $Z_s$ soit affine. L'obstruction au relèvement de $s$ en une section
de $\mathcal{L}^{\otimes n}$ sur $X$ est le cobord
$$
\xi = \partial s \in
H^1(X, \mathcal{I} \otimes_{\mathcal{O}_X} \mathcal{L}^{\otimes n}) =
H^1(Z, \mathcal{I} \otimes_{\mathcal{O}_Z} i^*\mathcal{L}^{\otimes n})
$$
issu de la suite exacte courte de faisceaux (\ref{coherent-equation-ses}).
Si l'on remplace $s$ par $s^{e + 1}$, alors $\xi$ est remplacé par
$\partial(s^{e + 1}) = (e + 1) s^e \xi$ dans
$H^1(Z, \mathcal{I} \otimes_{\mathcal{O}_Z} i^*\mathcal{L}^{\otimes (e + 1)n})$,
car l'application de cobord de
$$
0 \to
\bigoplus\nolimits_{m \geq 0}
\mathcal{I} \otimes_{\mathcal{O}_X} \mathcal{L}^{\otimes m} \to
\bigoplus\nolimits_{m \geq 0}
\mathcal{L}^{\otimes m} \to
\bigoplus\nolimits_{m \geq 0}
i_*i^*\mathcal{L}^{\otimes m} \to 0
$$
est une dérivation d'après Cohomologie, Lemme
\ref{cohomology-lemma-boundary-derivation-over-cup-product}. D'après le
lemme \ref{coherent-lemma-section-affine-open-kills-classes},
$s^e \xi$ est nul pour $e$ assez grand.
Ainsi, quitte à remplacer $s$ par une puissance, on peut supposer que $s$ est l'image
d'une section $s' \in \Gamma(X, \mathcal{L}^{\otimes n})$.
Alors $X_{s'}$ est un sous-schéma ouvert et $Z_s \to X_{s'}$ est une immersion
fermée surjective de schémas noethériens, avec $Z_s$ affine. Ainsi
$X_s$ est affine d'après le
lemme \ref{coherent-lemma-image-affine-finite-morphism-affine-Noetherian}, et
on en déduit que $\mathcal{L}$ est ample.
\end{proof}

\noindent
Pour une version plus générale du lemme suivant, voir
Limites, lemme \ref{limits-lemma-thickening-quasi-affine}.

\begin{lemma}
\label{coherent-lemma-thickening-quasi-affine}
Soit $i : Z \to X$ une immersion fermée de schémas noethériens induisant
un homéomorphisme des espaces topologiques sous-jacents.
Alors $X$ est quasi-affine si et seulement si $Z$ est quasi-affine.
\end{lemma}

\begin{proof}
Rappelons qu'un schéma est quasi-affine si
et seulement si son faisceau structural est ample, voir
Propriétés, lemme \ref{properties-lemma-quasi-affine-O-ample}.
Ainsi, si $Z$ est quasi-affine, alors $\mathcal{O}_Z$ est ample,
donc $\mathcal{O}_X$ est ample d'après le
lemme \ref{coherent-lemma-ample-on-reduction}, et par suite
$X$ est quasi-affine. La réciproque, que l'on peut
également établir de manière élémentaire, se démontre en lisant le
raisonnement précédent en sens inverse.
\end{proof}

\begin{lemma}
\label{coherent-lemma-affine-in-presence-ample}
Soit $X$ un schéma. Soit $\mathcal{L}$ un
$\mathcal{O}_X$-module inversible ample. Soit $n_0$ un entier.
Si $H^p(X, \mathcal{L}^{\otimes -n}) = 0$ pour $n \geq n_0$ et $p > 0$,
alors $X$ est affine.
\end{lemma}

\begin{proof}
Nous affirmons que $H^p(X, \mathcal{F}) = 0$ pour tout
$\mathcal{O}_X$-module quasi-cohérent et tout $p > 0$. Puisque $X$ est quasi-compact
d'après Propriétés, définition \ref{properties-definition-ample},
cette assertion achève la démonstration, en
vertu du lemme \ref{coherent-lemma-quasi-compact-h1-zero-covering}.
Le schéma $X$ est séparé d'après
Propriétés, lemme \ref{properties-lemma-ample-separated}. Supposons
que $X$ soit recouvert par $e + 1$ ouverts affines. Alors
$H^p(X, \mathcal{F}) = 0$ pour $p > e$, voir
lemme \ref{coherent-lemma-vanishing-nr-affines}. On peut donc procéder par récurrence
descendante sur $p$. En écrivant $\mathcal{F}$
comme limite inductive filtrante de modules quasi-cohérents de type
fini (Propriétés, Lemme
\ref{properties-lemma-quasi-coherent-colimit-finite-type})
et en utilisant Cohomologie, Lemme
\ref{cohomology-lemma-quasi-separated-cohomology-colimit},
on peut supposer $\mathcal{F}$ de type fini. On
peut alors choisir $n > n_0$ tel que
$\mathcal{F} \otimes_{\mathcal{O}_X} \mathcal{L}^{\otimes n}$
soit engendré par ses sections globales, voir Propriétés, Proposition
\ref{properties-proposition-characterize-ample}.
Cela signifie qu'il existe une suite exacte courte
$$
0 \to \mathcal{F}' \to
\bigoplus\nolimits_{i \in I} \mathcal{L}^{\otimes -n}
\to \mathcal{F} \to 0
$$
pour un certain ensemble $I$ (on peut en fait choisir $I$ fini). Par
hypothèse de récurrence, on a $H^{p + 1}(X, \mathcal{F}') = 0$
et, par hypothèse (en utilisant de nouveau la commutation
de la cohomologie aux limites inductives),
on a $H^p(X, \bigoplus_{i \in I} \mathcal{L}^{\otimes -n}) = 0$.
La suite exacte longue de cohomologie donne alors
$H^p(X, \mathcal{F}) = 0$, comme voulu.
\end{proof}

\begin{lemma}
\label{coherent-lemma-affine-if-quasi-affine}
Soit $X$ un schéma quasi-affine.
Si $H^p(X, \mathcal{O}_X) = 0$ pour $p > 0$,
alors $X$ est affine.
\end{lemma}

\begin{proof}
Puisque $\mathcal{O}_X$ est ample d'après
Propriétés, lemme \ref{properties-lemma-quasi-affine-O-ample},
cela résulte du lemme \ref{coherent-lemma-affine-in-presence-ample}.
\end{proof}







\section{Lemme de Chow}
\label{coherent-section-chows-lemma}

\noindent
Dans cette section, nous démontrons le lemme de Chow dans le cas
noethérien (lemme \ref{coherent-lemma-chow-Noetherian}).
Dans
Limites, section \ref{limits-section-chows-lemma},
nous en démontrons des variantes dans le cas non noethérien.

\begin{lemma}
\label{coherent-lemma-chow-Noetherian}
\begin{reference}
\cite[II théorème 5.6.1(a)]{EGA}
\end{reference}
Soit $S$ un schéma noethérien.
Soit $f : X \to S$ un morphisme séparé de type fini.
Il existe alors $n \geq 0$ et un diagramme
$$
\xymatrix{
X \ar[rd] & X' \ar[d] \ar[l]^\pi \ar[r] & \mathbf{P}^n_S \ar[dl] \\
& S &
}
$$
où $X' \to \mathbf{P}^n_S$ est une immersion et
$\pi : X' \to X$ est propre et surjectif. De plus, on
peut faire en sorte qu'il existe un sous-schéma ouvert dense
$U \subset X$ tel que $\pi^{-1}(U) \to U$ soit un isomorphisme.
\end{lemma}

\begin{proof}
Tous les schémas qui interviendront dans la suite de la
démonstration seront de type fini sur le schéma noethérien $S$,
donc noethériens
(voir Morphismes, lemme \ref{morphisms-lemma-finite-type-noetherian}).
Tous les morphismes entre eux seront automatiquement quasi-compacts, localement de
type fini et quasi-séparés, voir
Morphismes, lemme \ref{morphisms-lemma-permanence-finite-type} et
Propriétés,
lemmes \ref{properties-lemma-locally-Noetherian-quasi-separated} et
\ref{properties-lemma-morphism-Noetherian-schemes-quasi-compact}.

\medskip\noindent
Le schéma $X$ n'a qu'un nombre fini de composantes irréductibles
(Propriétés, lemme \ref{properties-lemma-Noetherian-irreducible-components}).
Soit $X = X_1 \cup \ldots \cup X_r$ la décomposition
de $X$ en composantes irréductibles.
Soit $\eta_i \in X_i$ le point générique.
Pour tout point $x \in X$, il existe un ouvert affine
$U_x \subset X$ contenant $x$ et chacun des points
génériques $\eta_i$. Voir
Propriétés, lemme \ref{properties-lemma-point-and-maximal-points-affine}.
Puisque $X$ est quasi-compact, il admet un recouvrement fini par des ouverts
affines $X = U_1 \cup \ldots \cup U_m$ tel que
chaque $U_i$ contienne $\eta_1, \ldots, \eta_r$.
On en déduit en particulier que l'ouvert
$U = U_1 \cap \ldots \cap U_m \subset X$ est
dense. C'est tout ce que nous utiliserons ci-dessous, avec le fait que les $U_i$ sont des ouverts affines
recouvrant $X$.

\medskip\noindent
Soit $X^* \subset X$ l'adhérence schématique de $U \to X$, voir
Morphismes, définition \ref{morphisms-definition-scheme-theoretic-image}.
Posons $U_i^* = X^* \cap U_i$. Notons que $U_i^*$ est un sous-schéma fermé
de $U_i$. Ainsi $U_i^*$ est affine. Puisque $U$ est dense dans $X$, le
morphisme $X^* \to X$ est une immersion fermée surjective. C'est un isomorphisme
au-dessus de $U$. On peut donc remplacer $X$ par $X^*$ et
$U_i$ par $U_i^*$, et supposer que $U$ est schématiquement dense
dans $X$, voir
Morphismes, définition \ref{morphisms-definition-scheme-theoretically-dense}.

\medskip\noindent
D'après Morphismes, lemme \ref{morphisms-lemma-quasi-projective-finite-type-over-S},
on peut trouver une immersion $j_i : U_i \to \mathbf{P}_S^{n_i}$
pour chaque $i$. D'après
Morphismes, lemme \ref{morphisms-lemma-quasi-compact-immersion}, on peut trouver
des sous-schémas fermés $Z_i \subset \mathbf{P}_S^{n_i}$
tels que $j_i : U_i \to Z_i$ soit une immersion ouverte
schématiquement dense. Notons que $Z_i \to S$ est propre, voir
Morphismes, lemme \ref{morphisms-lemma-locally-projective-proper}.
Considérons le morphisme
$$
j = (j_1|_U, \ldots, j_m|_U) : U \longrightarrow
\mathbf{P}_S^{n_1} \times_S \ldots \times_S \mathbf{P}_S^{n_m}.
$$
D'après le lemme cité ci-dessus, il existe un sous-schéma fermé
$Z$ de $\mathbf{P}_S^{n_1} \times_S \ldots \times_S \mathbf{P}_S^{n_m}$
tel que $j : U \to Z$ soit une immersion ouverte et que $U$
soit schématiquement dense dans $Z$. Le morphisme $Z \to S$
est propre. Considérons la $i$-ième projection
$$
\text{pr}_i|_Z : Z \longrightarrow \mathbf{P}^{n_i}_S.
$$
Ce morphisme se factorise par $Z_i$ (voir Morphismes,
lemme \ref{morphisms-lemma-factor-factor}). Notons $p_i : Z \to Z_i$
le morphisme induit. Ce morphisme est propre, voir par exemple
Morphismes, lemme \ref{morphisms-lemma-image-proper-scheme-closed}.
Nous avons donc des immersions ouvertes
$U \subset U_i \subset Z_i$ schématiquement denses.
De plus, on peut considérer $Z$ comme l'image
schématique du morphisme « diagonal »
$U \to Z_1 \times_S \ldots \times_S Z_m$.

\medskip\noindent
Posons $V_i = p_i^{-1}(U_i)$. Notons que $p_i|_{V_i} : V_i \to U_i$ est propre.
Posons $X' = V_1 \cup \ldots \cup V_m$. Par construction, $X'$ admet une immersion
dans le schéma
$\mathbf{P}^{n_1}_S \times_S \ldots \times_S \mathbf{P}^{n_m}_S$.
Ainsi, grâce au plongement de Segre (voir
Constructions, lemme \ref{constructions-lemma-segre-embedding}
), $X'$ admet
une immersion dans un espace projectif sur $S$.

\medskip\noindent
Nous affirmons que les morphismes $p_i|_{V_i}: V_i \to U_i$ se recollent en un morphisme
$X' \to X$. En effet, $p_i|_U$ est clairement l'application identité
de $U$ dans $U$. Puisque $U \subset X'$ est schématiquement dense
par construction, il est aussi schématiquement dense
dans le sous-schéma ouvert $V_i \cap V_j$. On a donc
$p_i|_{V_i \cap V_j} = p_j|_{V_i \cap V_j}$ comme morphismes dans
le $S$-schéma séparé $X$, voir
Morphismes, lemme \ref{morphisms-lemma-equality-of-morphisms}.
Notons $\pi : X' \to X$ le morphisme ainsi obtenu.

\medskip\noindent
Nous affirmons que $\pi^{-1}(U_i) = V_i$.
Puisque $\pi|_{V_i} = p_i|_{V_i}$, on a
$V_i \subset \pi^{-1}(U_i)$. Considérons le diagramme
$$
\xymatrix{
V_i \ar[r] \ar[rd]_{p_i|_{V_i}} & \pi^{-1}(U_i) \ar[d] \\
& U_i
}
$$
Puisque $V_i \to U_i$ est propre, l'image de la flèche
horizontale est fermée, voir
Morphismes, lemme \ref{morphisms-lemma-image-proper-scheme-closed}.
Puisque $V_i \subset \pi^{-1}(U_i)$ est schématiquement
dense (car il contient $U$),
on en déduit que $V_i = \pi^{-1}(U_i)$, comme annoncé.

\medskip\noindent
Cela montre que $\pi^{-1}(U_i) \to U_i$ s'identifie au morphisme
propre $p_i|_{V_i} : V_i \to U_i$. Ainsi $X$ admet un
recouvrement affine fini $X = \bigcup U_i$ tel que la restriction
de $\pi$ soit propre au-dessus de chaque ouvert du
recouvrement. D'après Morphismes, lemme \ref{morphisms-lemma-proper-local-on-the-base},
$\pi$ est donc propre.

\medskip\noindent
Enfin, il faut montrer que $\pi^{-1}(U) = U$. On raisonne pour cela comme
ci-dessus, à l'aide du diagramme
$$
\xymatrix{
U \ar[r] \ar[rd] & \pi^{-1}(U) \ar[d] \\
& U
}
$$
en utilisant le fait que $\text{id}_U : U \to U$ est propre et que
$U$ est schématiquement dense dans $\pi^{-1}(U)$.
\end{proof}

\begin{remark}
\label{coherent-remark-chow-Noetherian}
Dans la situation du Lemme de
Chow \ref{coherent-lemma-chow-Noetherian} :
\begin{enumerate}
\item Le morphisme $\pi$ est en fait H-projectif (donc projectif, voir
Morphismes, lemme \ref{morphisms-lemma-H-projective}),
car le morphisme $X' \to \mathbf{P}^n_S \times_S X = \mathbf{P}^n_X$
est une immersion fermée (on utilise le fait que $\pi$ est propre, voir
Morphismes, lemme \ref{morphisms-lemma-image-proper-scheme-closed}).
\item On peut supposer que $\pi^{-1}(U)$ est schématiquement dense
dans $X'$. En effet, il suffit de remplacer $X'$ par l'adhérence schématique
de $\pi^{-1}(U)$. On peut alors considérer $U$ comme un
sous-schéma ouvert schématiquement dense de $X'$.
Voir Morphismes, section \ref{morphisms-section-scheme-theoretic-image}.
\item Si $X$ est réduit, on peut choisir $X'$ réduit. Cela résulte clairement
de (2).
\end{enumerate}
\end{remark}






\section{Images directes supérieures des faisceaux cohérents}
\label{coherent-section-proper-pushforward}

\noindent
Dans cette section, nous démontrons le fait fondamental que les images
directes supérieures d'un faisceau cohérent par un morphisme propre
sont cohérentes.

\begin{proposition}
\label{coherent-proposition-proper-pushforward-coherent}
\begin{reference}
\cite[III théorème 3.2.1]{EGA}
\end{reference}
Soit $S$ un schéma localement noethérien.
Soit $f : X \to S$ un morphisme propre.
Soit $\mathcal{F}$ un $\mathcal{O}_X$-module cohérent.
Alors $R^if_*\mathcal{F}$ est un $\mathcal{O}_S$-module
cohérent pour tout $i \geq 0$.
\end{proposition}

\begin{proof}
La question étant locale sur $S$, on peut supposer que $S$ est
un schéma noethérien. Un morphisme propre
étant de type fini, $X$ est alors lui aussi un schéma
noethérien. Considérons la propriété $\mathcal{P}$ des faisceaux cohérents
sur $X$ définie par
$$
\mathcal{P}(\mathcal{F}) \Leftrightarrow
R^pf_*\mathcal{F}\text{ est cohérent pour tout }p \geq 0
$$
Nous allons utiliser le
lemme \ref{coherent-lemma-property} pour montrer que
$\mathcal{P}$ est vérifiée par tout faisceau cohérent sur $X$.

\medskip\noindent
Soit
$$
0 \to \mathcal{F}_1 \to \mathcal{F}_2 \to \mathcal{F}_3 \to 0
$$
une suite exacte courte de faisceaux cohérents sur $X$.
Considérons la suite exacte longue des images directes supérieures
$$
R^{p - 1}f_*\mathcal{F}_3 \to
R^pf_*\mathcal{F}_1 \to
R^pf_*\mathcal{F}_2 \to
R^pf_*\mathcal{F}_3 \to
R^{p + 1}f_*\mathcal{F}_1
$$
Il est alors clair que, si deux des trois faisceaux $\mathcal{F}_i$
vérifient la propriété $\mathcal{P}$, les images directes supérieures du troisième figurent,
dans ce complexe exact, entre deux faisceaux cohérents. Ces
images directes supérieures sont donc elles aussi cohérentes, d'après les
lemmes \ref{coherent-lemma-coherent-abelian-Noetherian} et
\ref{coherent-lemma-coherent-Noetherian-quasi-coherent-sub-quotient}.
La propriété $\mathcal{P}$ est donc aussi vérifiée par le troisième.

\medskip\noindent
Soit $Z \subset X$ un sous-schéma fermé intègre.
Il faut trouver un faisceau cohérent $\mathcal{F}$ sur $X$ dont le support soit
contenu dans $Z$ et dont la fibre au point générique $\xi$ de $Z$ soit un espace vectoriel de dimension
$1$ sur $\kappa(\xi)$, tel que $\mathcal{P}$
soit vérifiée par $\mathcal{F}$. Notons $g = f|_Z : Z \to S$ la restriction de $f$.
Supposons que l'on puisse trouver un faisceau cohérent $\mathcal{G}$ sur $Z$ tel
que
(a) $\mathcal{G}_\xi$ soit un espace vectoriel de dimension $1$ sur $\kappa(\xi)$,
(b) $R^pg_*\mathcal{G} = 0$ pour $p > 0$, et
(c) $g_*\mathcal{G}$ soit cohérent. On peut alors considérer
$\mathcal{F} = (Z \to X)_*\mathcal{G}$. Puisque $Z \to X$ est une
immersion fermée, $(Z \to X)_*\mathcal{G}$ est cohérent sur $X$
et $R^p(Z \to X)_*\mathcal{G} = 0$ pour $p > 0$
(lemme \ref{coherent-lemma-finite-pushforward-coherent}).
D'après la suite spectrale de Leray relative
(Cohomologie, lemme \ref{cohomology-lemma-relative-Leray}),
on aura donc $R^pf_*\mathcal{F} = R^pg_*\mathcal{G} = 0$ pour $p > 0$
et $f_*\mathcal{F} = g_*\mathcal{G}$ sera cohérent.
Enfin, $\mathcal{F}_\xi = ((Z \to X)_*\mathcal{G})_\xi = \mathcal{G}_\xi$,
ce qui vérifie la condition sur la fibre en $\xi$.
Tout se ramène donc à trouver un faisceau cohérent $\mathcal{G}$
sur $Z$ vérifiant les propriétés (a), (b) et (c).

\medskip\noindent
On peut appliquer le Lemme de Chow \ref{coherent-lemma-chow-Noetherian}
au morphisme $Z \to S$. On obtient ainsi un diagramme
$$
\xymatrix{
Z \ar[rd]_g & Z' \ar[d]^-{g'} \ar[l]^\pi \ar[r]_i & \mathbf{P}^m_S \ar[dl] \\
& S &
}
$$
comme dans l'énoncé du lemme de Chow. Soit aussi $U \subset Z$ le
sous-schéma ouvert dense tel que $\pi^{-1}(U) \to U$ soit un
isomorphisme. D'après la remarque \ref{coherent-remark-chow-Noetherian},
$i' = (i, \pi) : Z' \to \mathbf{P}^m_Z$ est
une immersion fermée. Ainsi
$$
\mathcal{L} = i^*\mathcal{O}_{\mathbf{P}^m_S}(1) \cong
(i')^*\mathcal{O}_{\mathbf{P}^m_Z}(1)
$$
est relativement ample pour $g'$ et pour $\pi$ (par exemple d'après Morphismes,
lemme \ref{morphisms-lemma-characterize-ample-on-finite-type}).
D'après le lemme \ref{coherent-lemma-kill-by-twisting},
il existe donc $n \geq 0$ tel que l'on ait à
la fois $R^p\pi_*\mathcal{L}^{\otimes n} = 0$ pour tout $p > 0$ et
$R^p(g')_*\mathcal{L}^{\otimes n} = 0$ pour tout $p > 0$.
Posons $\mathcal{G} = \pi_*\mathcal{L}^{\otimes n}$.
La propriété (a) est vérifiée, car $\pi_*\mathcal{L}^{\otimes n}|_U$ est
un faisceau inversible (puisque $\pi^{-1}(U) \to U$ est un isomorphisme).
Les propriétés (b) et (c) sont vérifiées, car la suite spectrale de
Leray relative
(Cohomologie, lemme \ref{cohomology-lemma-relative-Leray})
donne
$$
E_2^{p, q} = R^pg_* R^q\pi_*\mathcal{L}^{\otimes n}
\Rightarrow
R^{p + q}(g')_*\mathcal{L}^{\otimes n}
$$
et, par le choix de $n$, les seuls termes non nuls de $E_2^{p, q}$ sont
ceux pour lesquels $q = 0$, tandis que les seuls termes non nuls de
$R^{p + q}(g')_*\mathcal{L}^{\otimes n}$ sont ceux pour lesquels $p = q = 0$.
Cela entraîne que $R^pg_*\mathcal{G} = 0$ pour $p > 0$ et
que $g_*\mathcal{G} = (g')_*\mathcal{L}^{\otimes n}$.
Enfin, en appliquant le
lemme \ref{coherent-lemma-locally-projective-pushforward},
on voit que $g_*\mathcal{G} = (g')_*\mathcal{L}^{\otimes n}$ est
cohérent, comme voulu.
\end{proof}

\begin{lemma}
\label{coherent-lemma-proper-over-affine-cohomology-finite}
Soit $S = \Spec(A)$, où $A$ est un anneau noethérien.
Soit $f : X \to S$ un morphisme propre.
Soit $\mathcal{F}$ un $\mathcal{O}_X$-module cohérent.
Alors $H^i(X, \mathcal{F})$ est un $A$-module de type fini pour tout $i \geq 0$.
\end{lemma}

\begin{proof}
C'est simplement le cas affine de la
proposition \ref{coherent-proposition-proper-pushforward-coherent}. En
effet, d'après les lemmes \ref{coherent-lemma-quasi-coherence-higher-direct-images} et
\ref{coherent-lemma-quasi-coherence-higher-direct-images-application},
$R^if_*\mathcal{F}$ est le faisceau quasi-cohérent
associé au $A$-module $H^i(X, \mathcal{F})$;
d'après le lemme \ref{coherent-lemma-coherent-Noetherian}, ce
faisceau est cohérent si et seulement si $H^i(X, \mathcal{F})$
est un $A$-module de type fini.
\end{proof}

\begin{lemma}
\label{coherent-lemma-graded-finiteness}
Soit $A$ un anneau noethérien.
Soit $B$ une $A$-algèbre graduée de
type fini. Soit $f : X \to \Spec(A)$ un morphisme propre.
Posons $\mathcal{B} = f^*\widetilde B$.
Soit $\mathcal{F}$
un $\mathcal{B}$-module gradué quasi-cohérent de type fini.
\begin{enumerate}
\item Pour tout $p \geq 0$, le $B$-module gradué $H^p(X, \mathcal{F})$
est un $B$-module de type fini.
\item Si $\mathcal{L}$ est un $\mathcal{O}_X$-module
inversible ample, il existe un entier $d_0$ tel que
$H^p(X, \mathcal{F} \otimes \mathcal{L}^{\otimes d}) = 0$
pour tout $p > 0$ et tout $d \geq d_0$.
\end{enumerate}
\end{lemma}

\begin{proof}
Pour le démontrer, considérons le carré cartésien
$$
\xymatrix{
X' = \Spec(B) \times_{\Spec(A)} X
\ar[r]_-\pi \ar[d]_{f'} &
X \ar[d]^f \\
\Spec(B) \ar[r] &
\Spec(A)
}
$$
Notons que $f'$ est un morphisme propre, voir
Morphismes, lemme \ref{morphisms-lemma-base-change-proper}. De
plus, $B$ est de type fini comme $A$-algèbre, et est donc
noethérien (Algèbre, lemme \ref{algebra-lemma-Noetherian-permanence}).
Par suite, $X'$ est un schéma noethérien
(Morphismes, lemme \ref{morphisms-lemma-finite-type-noetherian}).
Notons que $X'$ est le spectre relatif de la
$\mathcal{O}_X$-algèbre quasi-cohérente $\mathcal{B}$,
d'après Constructions, lemme \ref{constructions-lemma-spec-properties}.
Puisque $\mathcal{F}$ est un $\mathcal{B}$-module
quasi-cohérent, il existe un unique
$\mathcal{O}_{X'}$-module quasi-cohérent $\mathcal{F}'$ tel que
$\pi_*\mathcal{F}' = \mathcal{F}$, voir
Morphismes, lemme \ref{morphisms-lemma-affine-equivalence-modules}.
Puisque $\mathcal{F}$ est de type fini comme $\mathcal{B}$
-module, $\mathcal{F}'$ est un
$\mathcal{O}_{X'}$-module de type fini (détails omis). Autrement dit,
$\mathcal{F}'$ est un $\mathcal{O}_{X'}$-module cohérent
(lemme \ref{coherent-lemma-coherent-Noetherian}).
Puisque le morphisme $\pi : X' \to X$ est affine, on a
$$
H^p(X, \mathcal{F}) = H^p(X', \mathcal{F}')
$$
d'après le lemme \ref{coherent-lemma-relative-affine-cohomology}.
Ainsi (1) résulte du
lemme \ref{coherent-lemma-proper-over-affine-cohomology-finite}. Étant
donné $\mathcal{L}$ comme dans (2), posons
$\mathcal{L}' = \pi^*\mathcal{L}$. Notons que $\mathcal{L}'$ est
ample sur $X'$,
d'après Morphismes, lemme \ref{morphisms-lemma-pullback-ample-tensor-relatively-ample}.
La formule de projection
(Cohomologie, lemme \ref{cohomology-lemma-projection-formula}) donne
$\pi_*(\mathcal{F}' \otimes \mathcal{L}') = \mathcal{F} \otimes \mathcal{L}$.
L'assertion (2) se déduit donc, par le même raisonnement que ci-dessus, du
lemme \ref{coherent-lemma-kill-by-twisting}.
\end{proof}












\section{Le théorème des fonctions formelles}
\label{coherent-section-theorem-formal-functions}

\noindent
Dans cette section, nous étudions le comportement de la cohomologie
des suites de faisceaux de la forme $\{I^n\mathcal{F}\}_{n \geq 0}$
ou de la forme $\{\mathcal{F}/I^n\mathcal{F}\}_{n \geq 0}$ lorsque
$n$ varie.

\medskip\noindent
Nous utiliserons ici et dans la suite la notation suivante.
Étant donnés un morphisme de schémas $f : X \to Y$, un faisceau quasi-cohérent
$\mathcal{F}$ sur $X$ et un faisceau quasi-cohérent d'idéaux
$\mathcal{I} \subset \mathcal{O}_Y$, on note
$\mathcal{I}^n\mathcal{F}$ le sous-faisceau quasi-cohérent engendré
par les produits de sections locales de $f^{-1}(\mathcal{I}^n)$ et de
$\mathcal{F}$. Autrement dit,
$$
\mathcal{I}^n\mathcal{F}
=
\Im\left(
f^*(\mathcal{I}^n) \otimes_{\mathcal{O}_X} \mathcal{F}
\longrightarrow
\mathcal{F}
\right).
$$
Notons qu'il existe des applications naturelles
$$
f^{-1}(\mathcal{I}^n) \otimes_{f^{-1}\mathcal{O}_Y} \mathcal{I}^m\mathcal{F}
\longrightarrow
f^*(\mathcal{I}^n) \otimes_{\mathcal{O}_X} \mathcal{I}^m\mathcal{F}
\longrightarrow
\mathcal{I}^{n + m}\mathcal{F}
$$
Ainsi, une section de $\mathcal{I}^n$ définit une
application $R^pf_*(\mathcal{I}^m\mathcal{F}) \to
R^pf_*(\mathcal{I}^{n + m}\mathcal{F})$ par fonctorialité
des images directes supérieures. Par localisation puis faisceautisation, on
obtient des applications de $\mathcal{O}_Y$-modules
$$
\mathcal{I}^n \otimes_{\mathcal{O}_Y} R^pf_*(\mathcal{I}^m\mathcal{F})
\longrightarrow
R^pf_*(\mathcal{I}^{n + m}\mathcal{F}).
$$
Autrement dit,
$\bigoplus_{n \geq 0} R^pf_*(\mathcal{I}^n\mathcal{F})$
est un $\bigoplus_{n \geq 0} \mathcal{I}^n$-module gradué.

\medskip\noindent
Si $Y = \Spec(A)$ et $\mathcal{I} = \widetilde{I}$, on note
$\mathcal{I}^n\mathcal{F}$ simplement $I^n\mathcal{F}$. Les applications
introduites ci-dessus munissent $M = \bigoplus H^p(X, I^n\mathcal{F})$ d'une
structure de module gradué sur $S = \bigoplus I^n$. Si $f$ est propre,
$A$ noethérien et $\mathcal{F}$ cohérent, ce module
est de type fini.

\begin{lemma}
\label{coherent-lemma-cohomology-powers-ideal-times-F}
\begin{reference}
\cite[III Cor 3.3.2]{EGA}
\end{reference}
Soit $A$ un anneau noethérien.
Soit $I \subset A$ un idéal.
Posons $B = \bigoplus_{n \geq 0} I^n$.
Soit $f : X \to \Spec(A)$ un morphisme propre.
Soit $\mathcal{F}$ un faisceau cohérent sur $X$.
Alors, pour tout $p \geq 0$, le $B$-module gradué
$\bigoplus_{n \geq 0} H^p(X, I^n\mathcal{F})$ est
un $B$-module de type fini.
\end{lemma}

\begin{proof}
Posons $\mathcal{B} = f^*\widetilde{B}$. Puisque $\mathcal{B}$ admet une
surjection sur $\bigoplus I^n\mathcal{O}_X$, il est clair que
$\bigoplus I^n\mathcal{F}$ est un $\mathcal{B}$-module gradué
de type fini. Le résultat découle donc du lemme \ref{coherent-lemma-graded-finiteness} (1).
\end{proof}

\begin{lemma}
\label{coherent-lemma-cohomology-powers-ideal-times-sheaf}
Soient un morphisme de schémas $f : X \to Y$, un faisceau quasi-cohérent
$\mathcal{F}$ sur $X$ et un faisceau quasi-cohérent d'idéaux
$\mathcal{I} \subset \mathcal{O}_Y$. Supposons $Y$ localement
noethérien, $f$ propre et $\mathcal{F}$ cohérent.
Alors
$$
\mathcal{M} =
\bigoplus\nolimits_{n \geq 0} R^pf_*(\mathcal{I}^n\mathcal{F})
$$
est un $\mathcal{A} = \bigoplus_{n \geq 0} \mathcal{I}^n$-module
gradué quasi-cohérent de type fini.
\end{lemma}

\begin{proof}
L'assertion est locale sur $Y$; on se ramène donc au
cas où $Y$ est affine. Dans ce cas, le résultat découle
du lemme \ref{coherent-lemma-cohomology-powers-ideal-times-F}.
Détails omis.
\end{proof}

\begin{lemma}
\label{coherent-lemma-cohomology-powers-ideal-application}
Soit $A$ un anneau noethérien.
Soit $I \subset A$ un idéal.
Soit $f : X \to \Spec(A)$ un morphisme propre.
Soit $\mathcal{F}$ un faisceau cohérent sur $X$.
Alors, pour tout $p \geq 0$, il existe un entier $c \geq 0$
tel que
\begin{enumerate}
\item l'application de multiplication
$I^{n - c} \otimes H^p(X, I^c\mathcal{F}) \to H^p(X, I^n\mathcal{F})$
soit surjective pour tout $n \geq c$;
\item l'image de $H^p(X, I^{n + m}\mathcal{F}) \to H^p(X, I^n\mathcal{F})$
soit contenue dans le sous-module $I^{m - e} H^p(X, I^n\mathcal{F})$,
où $e = \max(0, c - n)$, pour $n + m \geq c$, $n, m \geq 0$;
\item on ait
$$
\Ker(H^p(X, I^n\mathcal{F}) \to H^p(X, \mathcal{F})) =
\Ker(H^p(X, I^n\mathcal{F}) \to H^p(X, I^{n - c}\mathcal{F}))
$$
pour $n \geq c$;
\item il existe des applications $I^nH^p(X, \mathcal{F}) \to H^p(X, I^{n - c}\mathcal{F})$
pour $n \geq c$, telles que les composés
$$
H^p(X, I^n\mathcal{F}) \to
I^{n - c}H^p(X, \mathcal{F}) \to
H^p(X, I^{n - 2c}\mathcal{F})
$$
et
$$
I^nH^p(X, \mathcal{F}) \to
H^p(X, I^{n - c}\mathcal{F}) \to
I^{n - 2c}H^p(X, \mathcal{F})
$$
pour $n \geq 2c$ soient les applications canoniques; et
\item les systèmes projectifs $(H^p(X, I^n\mathcal{F}))$ et
$(I^nH^p(X, \mathcal{F}))$ soient pro-isomorphes.
\end{enumerate}
\end{lemma}

\begin{proof}
Posons $M_n = H^p(X, I^n\mathcal{F})$ pour $n \geq 1$ et
$M_0 = H^p(X, \mathcal{F})$, de sorte que l'on ait des applications
$\ldots \to M_3 \to M_2 \to M_1 \to M_0$. En posant
$B = \bigoplus_{n \geq 0} I^n$, on voit que $M = \bigoplus_{n \geq 0} M_n$
est un $B$-module gradué de type fini,
voir lemme \ref{coherent-lemma-cohomology-powers-ideal-times-F}.
Notons que les produits
$B_n \otimes M_m \to M_{m + n}$, $a \otimes m \mapsto a \cdot m$,
sont compatibles avec les applications de notre système projectif, au sens
où les diagrammes
$$
\xymatrix{
B_n \otimes_A M_m \ar[r] \ar[d] & M_{n + m} \ar[d] \\
B_n \otimes_A M_{m'} \ar[r] & M_{n + m'}
}
$$
commutent pour $n, m' \geq 0$ et $m \geq m'$.

\medskip\noindent
Démontrons (1). Choisissons $d_1, \ldots, d_t \geq 0$ et $x_i \in M_{d_i}$
tels que $M$ soit engendré par $x_1, \ldots, x_t$ sur $B$.
Pour tout $c \geq \max\{d_i\}$, on a alors
$B_{n - c} \cdot M_c = M_n$ pour $n \geq c$, ce qui établit (1).

\medskip\noindent
Démontrons (2). Soit $c$ comme dans la démonstration de (1). Soit $n + m \geq c$.
On a $M_{n + m} = B_{n + m - c} \cdot M_c$.
Si $c > n$, on utilise $M_c \to M_n$ et la compatibilité des produits
avec les applications de transition signalée ci-dessus pour en déduire
que l'image de $M_{n + m} \to M_n$ est contenue dans $I^{n + m - c}M_n$.
Si $c \leq n$, on écrit $M_{n + m} = B_m \cdot B_{n - c} \cdot M_c =
B_m \cdot M_n$, ce qui montre que l'image est contenue dans $I^m M_n$.
Cela démontre (2).

\medskip\noindent
Soit $K_n \subset M_n$ le noyau de l'application $M_n \to M_0$. La
compatibilité des produits avec les applications de transition signalée ci-dessus
montre que $K = \bigoplus K_n \subset M$ est un sous-$B$-module gradué.
Puisque $B$ est noethérien et que $M$ est un $B$-module
gradué de type fini, $K$ est un $B$-module gradué de
type fini. Choisissons $d'_1, \ldots, d'_{t'} \geq 0$ et $y_i \in K_{d'_i}$
tels que $K$ soit engendré par $y_1, \ldots, y_{t'}$ sur $B$.
Posons $c = \max(d'_i, d'_j)$. Puisque $y_i \in \Ker(M_{d'_i} \to M_0)$,
on a $B_n \cdot y_i \subset \Ker(M_{n + d'_i} \to M_n)$.
On en déduit que $K_n = \Ker(M_n \to M_{n - c})$ pour $n \geq c$.
Cela démontre (3).

\medskip\noindent
Considérons le diagramme commutatif suivant, dont les flèches pleines sont données :
$$
\xymatrix{
I^n \otimes_A M_0 \ar[r] \ar[d] &
I^nM_0 \ar[r] \ar@{..>}[d] &
M_0 \ar[d] \\
M_n \ar[r] &
M_{n - c} \ar[r] &
M_0
}
$$
Puisque le noyau de la flèche surjective $I^n \otimes_A M_0 \to I^nM_0$
a son image dans $K_n$, ce qui précède fournit la flèche pointillée, et le
composé $I^nM_0 \to M_{n - c} \to M_0$ est l'application canonique. Il est
alors clair que le composé $I^nM_0 \to M_{n - c} \to I^{n - 2c}M_0$
est l'application canonique pour $n \geq 2c$. Considérons le composé
$M_n \to I^{n - c}M_0 \to M_{n - 2c}$. La première application
envoie un élément de la forme $a \cdot m$,
avec $a \in I^{n - c}$ et $m \in M_c$,
sur $a m'$, où $m'$ est l'image de $m$ dans $M_0$.
La seconde application envoie cet élément sur $a \cdot m'$ dans $M_{n - 2c}$,
ce qui établit (4).

\medskip\noindent
L'assertion (5) résulte immédiatement de (4) et de la définition des
morphismes de pro-objets.
\end{proof}

\noindent
Dans la situation des lemmes \ref{coherent-lemma-cohomology-powers-ideal-times-F} et
\ref{coherent-lemma-cohomology-powers-ideal-application}, considérons le système
projectif
$$
\mathcal{F}/I\mathcal{F} \leftarrow
\mathcal{F}/I^2\mathcal{F} \leftarrow
\mathcal{F}/I^3\mathcal{F} \leftarrow \ldots
$$
Nous voudrions connaître le comportement des groupes de cohomologie.
Voici un premier résultat.

\begin{lemma}
\label{coherent-lemma-ML-cohomology-powers-ideal}
Soit $A$ un anneau noethérien.
Soit $I \subset A$ un idéal.
Soit $f : X \to \Spec(A)$ un morphisme propre.
Soit $\mathcal{F}$ un faisceau cohérent sur $X$.
Fixons $p \geq 0$. Il existe $c \geq 0$ tel que
\begin{enumerate}
\item pour tout $n \geq c$, on ait
$$
\Ker(H^p(X, \mathcal{F}) \to H^p(X, \mathcal{F}/I^n\mathcal{F})) \subset
I^{n - c}H^p(X, \mathcal{F}).
$$
\item le système projectif
$$
\left(H^p(X, \mathcal{F}/I^n\mathcal{F})\right)_{n \in \mathbf{N}}
$$
vérifie la condition de Mittag-Leffler (voir
Homologie, définition \ref{homology-definition-Mittag-Leffler}); et
\item on ait
$$
\Im(H^p(X, \mathcal{F}/I^k\mathcal{F})
\to H^p(X, \mathcal{F}/I^n\mathcal{F}))
=
\Im(H^p(X, \mathcal{F})
\to H^p(X, \mathcal{F}/I^n\mathcal{F}))
$$
pour tout $k \geq n + c$.
\end{enumerate}
\end{lemma}

\begin{proof}
Posons $c = \max\{c_p, c_{p + 1}\}$, où $c_p, c_{p + 1}$ sont les entiers obtenus
dans le lemme \ref{coherent-lemma-cohomology-powers-ideal-application} pour
$H^p$ et $H^{p + 1}$.

\medskip\noindent
Démontrons (1). Considérons la suite exacte courte
$$
0 \to I^n\mathcal{F} \to \mathcal{F} \to \mathcal{F}/I^n\mathcal{F} \to 0
$$
La suite exacte longue de cohomologie donne
$$
\Ker(
H^p(X, \mathcal{F}) \to H^p(X, \mathcal{F}/I^n\mathcal{F})
)
=
\Im(
H^p(X, I^n\mathcal{F}) \to H^p(X, \mathcal{F})
)
$$
D'après le lemme \ref{coherent-lemma-cohomology-powers-ideal-application} (2), ce
sous-module est donc contenu dans $I^{n - c}H^p(X, \mathcal{F})$ pour $n \geq c$.

\medskip\noindent
Notons que (3) entraîne (2), par définition des systèmes de
Mittag-Leffler.

\medskip\noindent
Démontrons (3). Fixons $n$. Considérons le diagramme commutatif
$$
\xymatrix{
0 \ar[r] &
I^n\mathcal{F} \ar[r] &
\mathcal{F} \ar[r] &
\mathcal{F}/I^n\mathcal{F} \ar[r] &
0 \\
0 \ar[r] &
I^{n + m}\mathcal{F} \ar[r] \ar[u] &
\mathcal{F} \ar[r] \ar[u] &
\mathcal{F}/I^{n + m}\mathcal{F} \ar[r] \ar[u] &
0
}
$$
Il induit le diagramme commutatif suivant :
$$
\xymatrix{
H^p(X, \mathcal{F}) \ar[r] &
H^p(X, \mathcal{F}/I^n\mathcal{F}) \ar[r]_\delta &
H^{p + 1}(X, I^n\mathcal{F}) \ar[r] &
H^{p + 1}(X, \mathcal{F}) \\
H^p(X, \mathcal{F}) \ar[r] \ar[u]^1 &
H^p(X, \mathcal{F}/I^{n + m}\mathcal{F}) \ar[r] \ar[u]^\gamma &
H^{p + 1}(X, I^{n + m}\mathcal{F}) \ar[u]^\alpha \ar[r]^-\beta &
H^{p + 1}(X, \mathcal{F}) \ar[u]_1
}
$$
dont les lignes sont exactes. D'après
le lemme \ref{coherent-lemma-cohomology-powers-ideal-application} (4), le noyau
de $\beta$ est égal à celui de $\alpha$ pour $m \geq c$.
Une chasse au diagramme montre alors que l'image de $\gamma$ est contenue
dans le noyau de $\delta$, ce qui établit (3) (en
posant $k = n + m$).
\end{proof}

\begin{theorem}[Théorème des fonctions formelles]
\label{coherent-theorem-formal-functions}
Soit $A$ un anneau noethérien.
Soit $I \subset A$ un idéal.
Soit $f : X \to \Spec(A)$ un morphisme propre.
Soit $\mathcal{F}$ un faisceau cohérent sur $X$.
Fixons $p \geq 0$.
Le système d'applications
$$
H^p(X, \mathcal{F})/I^nH^p(X, \mathcal{F})
\longrightarrow
H^p(X, \mathcal{F}/I^n\mathcal{F})
$$
définit un isomorphisme de limites
$$
H^p(X, \mathcal{F})^\wedge
\longrightarrow
\lim_n H^p(X, \mathcal{F}/I^n\mathcal{F})
$$
où le membre de gauche est le complété du $A$-module
$H^p(X, \mathcal{F})$ relativement à l'idéal $I$, voir
Algèbre, section \ref{algebra-section-completion}.
De plus, cet isomorphisme est un homéomorphisme pour les topologies de limite projective.
\end{theorem}

\begin{proof}
Cela se déduit du lemme \ref{coherent-lemma-ML-cohomology-powers-ideal} comme suit.
Posons $M = H^p(X, \mathcal{F})$, $M_n = H^p(X, \mathcal{F}/I^n\mathcal{F})$,
et notons $N_n = \Im(M \to M_n)$. D'après le
lemme \ref{coherent-lemma-ML-cohomology-powers-ideal} (2) et (3),
$(M_n)$ est un système de Mittag-Leffler dans lequel
$N_n \subset M_n$ est l'image de $M_k$ pour tout $k \gg n$.
Il en résulte que $\lim M_n = \lim N_n$ comme modules topologiques (munis des
topologies de limite projective). D'autre part, les $N_n$ forment un système projectif de
quotients du module $M$; par suite, $\lim N_n$ est le complété de $M$
pour la topologie définie par les noyaux $K_n = \Ker(M \to N_n)$.
D'après le lemme \ref{coherent-lemma-ML-cohomology-powers-ideal} (1), on
a $K_n \subset I^{n - c}M$ et, puisque $N_n \subset M_n$
est annulé par $I^n$, on a $I^n M \subset K_n$.
La topologie définie par les sous-modules $K_n$
comme système fondamental de voisinages ouverts de $0$
est donc la topologie $I$-adique, et l'application
induite $M^\wedge = \lim M/I^nM \to \lim N_n = \lim M_n$
est un isomorphisme de modules topologiques\footnote{
Précisons que la topologie de limite projective sur $M^\wedge$ coïncide avec sa topologie
$I$-adique, d'après Algèbre,
lemme \ref{algebra-lemma-hathat-finitely-generated}.
Voir Compléments d'algèbre, section \ref{more-algebra-section-topological-ring}.}.
\end{proof}

\begin{lemma}
\label{coherent-lemma-spell-out-theorem-formal-functions}
Soit $A$ un anneau. Soit $I \subset A$ un idéal. Supposons $A$ noethérien
et complet relativement à $I$.
Soit $f : X \to \Spec(A)$ un morphisme propre.
Soit $\mathcal{F}$ un faisceau cohérent sur $X$.
Alors
$$
H^p(X, \mathcal{F}) = \lim_n H^p(X, \mathcal{F}/I^n\mathcal{F})
$$
pour tout $p \geq 0$.
\end{lemma}

\begin{proof}
C'est une reformulation du théorème des fonctions formelles
(Théorème \ref{coherent-theorem-formal-functions}) dans le
cas d'un anneau de base noethérien complet. En effet, dans ce cas, le
$A$-module $H^p(X, \mathcal{F})$ est de type
fini (lemme \ref{coherent-lemma-proper-over-affine-cohomology-finite}), donc complet pour la topologie
$I$-adique (Algèbre, lemme \ref{algebra-lemma-completion-tensor});
il est donc inutile de compléter le membre de gauche.
\end{proof}

\begin{lemma}
\label{coherent-lemma-formal-functions-stalk}
Soient un morphisme de schémas $f : X \to Y$ et un faisceau quasi-cohérent
$\mathcal{F}$ sur $X$. Supposons que
\begin{enumerate}
\item $Y$ soit localement noethérien,
\item $f$ soit propre, et
\item $\mathcal{F}$ soit cohérent.
\end{enumerate}
Soit $y \in Y$ un point. Considérons les voisinages infinitésimaux
$$
\xymatrix{
X_n =
\Spec(\mathcal{O}_{Y, y}/\mathfrak m_y^n) \times_Y X
\ar[r]_-{i_n} \ar[d]_{f_n} &
X \ar[d]^f \\
\Spec(\mathcal{O}_{Y, y}/\mathfrak m_y^n) \ar[r]^-{c_n} & Y
}
$$
de la fibre $X_1 = X_y$ et posons $\mathcal{F}_n = i_n^*\mathcal{F}$.
On a alors
$$
\left(R^pf_*\mathcal{F}\right)_y^\wedge
\cong
\lim_n H^p(X_n, \mathcal{F}_n)
$$
comme $\mathcal{O}_{Y, y}^\wedge$-modules.
\end{lemma}

\begin{proof}
C'est simplement une reformulation d'un cas particulier du théorème des
fonctions formelles, Théorème \ref{coherent-theorem-formal-functions}.
Détaillons-la. Notons que $\mathcal{O}_{Y, y}$ est un anneau
local noethérien. Considérons le morphisme canonique
$c : \Spec(\mathcal{O}_{Y, y}) \to Y$, voir
Schémas, Équation (\ref{schemes-equation-canonical-morphism}).
Ce morphisme est plat, car il identifie les anneaux locaux.
Notons provisoirement $f' : X' \to \Spec(\mathcal{O}_{Y, y})$
le changement de base de $f$ à cet anneau local. On a
$c^*R^pf_*\mathcal{F} = R^pf'_*\mathcal{F}'$, d'après
le lemme \ref{coherent-lemma-flat-base-change-cohomology}.
De plus, les voisinages infinitésimaux
des fibres $X_y$ et $X'_y$ s'identifient (vérification omise; indication :
les morphismes $c_n$ se factorisent par $c$).

\medskip\noindent
On peut donc supposer que $Y = \Spec(A)$ est le spectre d'un
anneau local noethérien $A$ d'idéal maximal $\mathfrak m$,
et que $y \in Y$ est le point fermé (c'est-à-dire celui correspondant à $\mathfrak m$).
Il en résulte en particulier que
$$
\left(R^pf_*\mathcal{F}\right)_y =
\Gamma(Y, R^pf_*\mathcal{F}) =
H^p(X, \mathcal{F}).
$$

\medskip\noindent
Dans ce cas, les morphismes $c_n$ sont aussi des immersions
fermées. Leurs changements de base $i_n$ sont donc eux aussi des immersions
fermées. Notons que $i_{n, *}\mathcal{F}_n = i_{n, *}i_n^*\mathcal{F}
= \mathcal{F}/\mathfrak m^n\mathcal{F}$. La suite spectrale de Leray
pour $i_n$ et le lemme \ref{coherent-lemma-finite-pushforward-coherent} donnent
$$
H^p(X_n, \mathcal{F}_n) =
H^p(X, i_{n, *}\mathcal{F}_n) =
H^p(X, \mathcal{F}/\mathfrak m^n\mathcal{F})
$$
On peut donc appliquer le théorème des fonctions formelles
pour calculer la limite de l'énoncé, ce qui conclut.
\end{proof}

\noindent
Voici un lemme que nous généraliserons aux fibres de
dimension $ > 0$ dans le lemme suivant.

\begin{lemma}
\label{coherent-lemma-higher-direct-images-zero-finite-fibre}
Soit $f : X \to Y$ un morphisme de schémas.
Soit $y \in Y$. Supposons
que
\begin{enumerate}
\item $Y$ soit localement noethérien,
\item $f$ soit propre, et
\item $f^{-1}(\{y\})$ soit fini.
\end{enumerate}
Alors, pour tout faisceau cohérent $\mathcal{F}$ sur $X$, on a
$(R^pf_*\mathcal{F})_y = 0$ pour tout $p > 0$.
\end{lemma}

\begin{proof}
La fibre $X_y$ est finie et, d'après
Morphismes, lemme \ref{morphisms-lemma-finite-fibre},
c'est un espace discret fini. De
plus, chaque voisinage infinitésimal $X_n$ a le même
espace topologique sous-jacent. Chacun des schémas $X_n$ est donc affine, d'après
Schémas, lemme \ref{schemes-lemma-scheme-finite-discrete-affine}.
Il en résulte que $H^p(X_n, \mathcal{F}_n) = 0$ pour tout
$p > 0$. Ainsi $(R^pf_*\mathcal{F})_y^\wedge = 0$,
d'après le lemme \ref{coherent-lemma-formal-functions-stalk}.
Notons que $R^pf_*\mathcal{F}$ est cohérent, d'après la
proposition \ref{coherent-proposition-proper-pushforward-coherent}; par
suite, $R^pf_*\mathcal{F}_y$ est un
$\mathcal{O}_{Y, y}$-module de type fini. D'après le lemme
de Nakayama (Algèbre, lemme \ref{algebra-lemma-NAK}),
si le complété d'un module de type fini sur un anneau
local est nul, le module est nul. D'où
$(R^pf_*\mathcal{F})_y = 0$.
\end{proof}

\begin{lemma}
\label{coherent-lemma-higher-direct-images-zero-above-dimension-fibre}
Soit $f : X \to Y$ un morphisme de schémas.
Soit $y \in Y$. Supposons
que
\begin{enumerate}
\item $Y$ soit localement noethérien,
\item $f$ soit propre, et
\item $\dim(X_y) = d$.
\end{enumerate}
Alors, pour tout faisceau cohérent $\mathcal{F}$ sur $X$, on a
$(R^pf_*\mathcal{F})_y = 0$ pour tout $p > d$.
\end{lemma}

\begin{proof}
La fibre $X_y$ est de type fini sur $\Spec(\kappa(y))$.
Ainsi $X_y$ est un schéma noethérien, d'après
Morphismes, lemme \ref{morphisms-lemma-finite-type-noetherian}.
L'espace topologique sous-jacent de $X_y$ est donc noethérien, voir
Propriétés, lemme \ref{properties-lemma-Noetherian-topology}.
De plus, chaque voisinage
infinitésimal $X_n$ a le même espace topologique sous-jacent que $X_y$.
Ainsi $H^p(X_n, \mathcal{F}_n) = 0$ pour tout $p > d$,
d'après Cohomologie, proposition \ref{cohomology-proposition-vanishing-Noetherian}.
On en déduit que $(R^pf_*\mathcal{F})_y^\wedge = 0$,
d'après le lemme \ref{coherent-lemma-formal-functions-stalk}, pour $p > d$.
Notons que $R^pf_*\mathcal{F}$ est cohérent, d'après la
proposition \ref{coherent-proposition-proper-pushforward-coherent}; par
suite, $R^pf_*\mathcal{F}_y$ est un
$\mathcal{O}_{Y, y}$-module de type fini. D'après le lemme
de Nakayama (Algèbre, lemme \ref{algebra-lemma-NAK}),
si le complété d'un module de type fini sur un anneau
local est nul, le module est nul. D'où
$(R^pf_*\mathcal{F})_y = 0$.
\end{proof}








\section{Applications du théorème sur les fonctions formelles}
\label{coherent-section-applications-formal-functions}


\noindent
Nous compléterons cette section au besoin. Pour le moment, nous avons besoin de la
caractérisation suivante des morphismes finis dans le cas noethérien.

\begin{lemma}
\label{coherent-lemma-characterize-finite}
(Pour une version plus générale, voir Compléments
sur les morphismes, lemme \ref{more-morphisms-lemma-characterize-finite}.)
Soit $f : X \to S$ un morphisme de schémas. Supposons
que $S$ soit localement noethérien. Les
assertions suivantes sont équivalentes :
\begin{enumerate}
\item $f$ est fini, et
\item $f$ est propre à fibres finies.
\end{enumerate}
\end{lemma}

\begin{proof}
Un morphisme fini est propre d'après
Morphismes, lemme \ref{morphisms-lemma-finite-proper}.
Un morphisme fini est quasi-fini d'après
Morphismes, lemme \ref{morphisms-lemma-finite-quasi-finite}.
Un morphisme quasi-fini a des fibres finies, voir
Morphismes, lemme \ref{morphisms-lemma-quasi-finite}.
Ainsi, un morphisme fini est propre et a des fibres finies.

\medskip\noindent
Supposons que $f$ soit propre à fibres finies.
Nous voulons montrer que $f$ est fini.
En fait, il suffit de montrer que $f$ est affine.
En effet, si $f$ est affine, alors
$f$ est intégral d'après
Morphismes, lemme \ref{morphisms-lemma-integral-universally-closed},
et il résulte alors
de Morphismes, lemme \ref{morphisms-lemma-finite-integral},
que $f$ est fini.

\medskip\noindent
Pour montrer que $f$ est affine, nous pouvons supposer que $S$ est affine, et notre
but est alors de montrer que $X$ est également affine.
Puisque $f$ est propre, on voit que $X$ est séparé et quasi-compact.
Nous pouvons donc utiliser le critère du
lemme \ref{coherent-lemma-quasi-separated-h1-zero-covering} pour montrer que $X$
est affine. Pour cela, soit $\mathcal{I} \subset \mathcal{O}_X$
un faisceau d'idéaux de type fini. En particulier, $\mathcal{I}$ est
un faisceau cohérent sur $X$. D'après le
lemme \ref{coherent-lemma-higher-direct-images-zero-finite-fibre}, nous concluons que
$R^1f_*\mathcal{I}_s = 0$ pour tout $s \in S$.
Autrement dit, $R^1f_*\mathcal{I} = 0$. La suite spectrale
de Leray pour $f$ donne donc
$H^1(X , \mathcal{I}) = H^1(S, f_*\mathcal{I})$.
Comme $S$ est affine et que $f_*\mathcal{I}$ est quasi-cohérent
(Schémas, lemme \ref{schemes-lemma-push-forward-quasi-coherent}),
nous déduisons $H^1(S, f_*\mathcal{I}) = 0$
du lemme \ref{coherent-lemma-quasi-coherent-affine-cohomology-zero},
comme voulu. Ainsi, $H^1(X, \mathcal{I}) = 0$, comme voulu.
\end{proof}

\noindent
On en déduit le résultat utile suivant.

\begin{lemma}
\label{coherent-lemma-proper-finite-fibre-finite-in-neighbourhood}
\begin{slogan}
Un morphisme propre est fini au voisinage d'une fibre finie.
\end{slogan}
(Pour une version plus générale, voir Compléments
sur les morphismes,
lemme \ref{more-morphisms-lemma-proper-finite-fibre-finite-in-neighbourhood}.)
Soit $f : X \to S$ un morphisme de schémas.
Soit $s \in S$. Supposons
que
\begin{enumerate}
\item $S$ est localement noethérien,
\item $f$ est propre, et
\item $f^{-1}(\{s\})$ est un ensemble fini.
\end{enumerate}
Il existe alors un voisinage ouvert $V \subset S$ de $s$
tel que $f|_{f^{-1}(V)} : f^{-1}(V) \to V$ soit fini.
\end{lemma}

\begin{proof}
Le morphisme $f$ est quasi-fini en tout point de $f^{-1}(\{s\})$
d'après Morphismes, lemme \ref{morphisms-lemma-finite-fibre}.
D'après Morphismes, lemme \ref{morphisms-lemma-quasi-finite-points-open},
l'ensemble des points où $f$ est quasi-fini est un ouvert $U \subset X$.
Posons $Z = X \setminus U$. Alors $s \not \in f(Z)$. Puisque $f$ est propre,
l'ensemble $f(Z) \subset S$ est fermé. Choisissons un voisinage ouvert
$V \subset S$ de $s$ tel que $Z \cap V = \emptyset$. Alors
$f^{-1}(V) \to V$ est localement quasi-fini et propre.
Il est donc quasi-fini
(Morphismes, lemme \ref{morphisms-lemma-quasi-finite-locally-quasi-compact}),
donc à fibres finies
(Morphismes, lemme \ref{morphisms-lemma-quasi-finite}), et par
conséquent fini d'après le lemme \ref{coherent-lemma-characterize-finite}.
\end{proof}

\begin{lemma}
\label{coherent-lemma-ample-on-fibre}
Soit $f : X \to Y$ un morphisme propre de schémas, avec $Y$ noethérien.
Soit $\mathcal{L}$ un $\mathcal{O}_X$-module inversible.
Soit $\mathcal{F}$ un $\mathcal{O}_X$-module cohérent.
Soit $y \in Y$ un point tel que $\mathcal{L}_y$ soit ample sur $X_y$.
Il existe alors un entier $d_0$ tel que, pour tout $d \geq d_0$, on ait
$$
R^pf_*(\mathcal{F} \otimes_{\mathcal{O}_X} \mathcal{L}^{\otimes d})_y = 0
\text{ pour }p > 0
$$
et l'application
$$
f_*(\mathcal{F} \otimes_{\mathcal{O}_X} \mathcal{L}^{\otimes d})_y
\longrightarrow
H^0(X_y, \mathcal{F}_y \otimes_{\mathcal{O}_{X_y}} \mathcal{L}_y^{\otimes d})
$$
est surjective.
\end{lemma}

\begin{proof}
Remarquons que $\mathcal{O}_{Y, y}$ est un anneau local noethérien.
Considérons le morphisme canonique
$c : \Spec(\mathcal{O}_{Y, y}) \to Y$, voir
Schémas, Équation (\ref{schemes-equation-canonical-morphism}).
Ce morphisme est plat, puisqu'il identifie les anneaux locaux.
Notons provisoirement $f' : X' \to \Spec(\mathcal{O}_{Y, y})$
le changement de base de $f$ à cet anneau local. On a
$c^*R^pf_*\mathcal{F} = R^pf'_*\mathcal{F}'$ d'après le
lemme \ref{coherent-lemma-flat-base-change-cohomology}.
De plus, les fibres $X_y$ et $X'_y$ sont identifiées.
Nous pouvons donc supposer que $Y = \Spec(A)$ est le spectre d'un
anneau local noethérien $(A, \mathfrak m, \kappa)$ et que $y \in Y$
correspond à $\mathfrak m$. Dans ce cas,
$R^pf_*(\mathcal{F} \otimes_{\mathcal{O}_X} \mathcal{L}^{\otimes d})_y =
H^p(X, \mathcal{F} \otimes_{\mathcal{O}_X} \mathcal{L}^{\otimes d})$
pour tout $p \geq 0$. Notons $f_y : X_y \to \Spec(\kappa)$ la projection.

\medskip\noindent
Posons $B = \text{Gr}_\mathfrak m(A) =
\bigoplus_{n \geq 0} \mathfrak m^n/\mathfrak m^{n + 1}$.
Considérons le faisceau $\mathcal{B} = f_y^*\widetilde{B}$
de $\mathcal{O}_{X_y}$-algèbres
graduées quasi-cohérentes. Nous utiliserons les notations de la section \ref{coherent-section-theorem-formal-functions},
où $I$ est remplacé par $\mathfrak m$.
Comme $X_y$ est le sous-schéma fermé de $X$ défini par
$\mathfrak m\mathcal{O}_X$, nous pouvons considérer
$\mathfrak m^n\mathcal{F}/\mathfrak m^{n + 1}\mathcal{F}$
comme un $\mathcal{O}_{X_y}$-module cohérent, voir
le lemme \ref{coherent-lemma-i-star-equivalence}. Alors
$\bigoplus_{n \geq 0} \mathfrak m^n\mathcal{F}/\mathfrak m^{n + 1}\mathcal{F}$
est un $\mathcal{B}$-module gradué quasi-cohérent de type
fini, car il est engendré en degré zéro sur $\mathcal{B}$
et car sa partie de degré zéro est
$\mathcal{F}_y = \mathcal{F}/\mathfrak m \mathcal{F}$
qui est un $\mathcal{O}_{X_y}$-module cohérent.
Ainsi, d'après le lemme \ref{coherent-lemma-graded-finiteness}, partie
(2), on voit que
$$
H^p(X_y, \mathfrak m^n \mathcal{F}/ \mathfrak m^{n + 1}\mathcal{F}
\otimes_{\mathcal{O}_{X_y}} \mathcal{L}_y^{\otimes d}) = 0
$$
pour tous $p > 0$, $d \geq d_0$ et $n \geq 0$. D'après le
lemme \ref{coherent-lemma-relative-affine-cohomology},
cela revient à dire que
$
H^p(X, \mathfrak m^n \mathcal{F}/ \mathfrak m^{n + 1}\mathcal{F}
\otimes_{\mathcal{O}_X} \mathcal{L}^{\otimes d}) = 0
$
pour tous $p > 0$, $d \geq d_0$ et $n \geq 0$.

\medskip\noindent
Considérons les suites exactes courtes
$$
0 \to \mathfrak m^n\mathcal{F}/\mathfrak m^{n + 1} \mathcal{F}
\to \mathcal{F}/\mathfrak m^{n + 1} \mathcal{F}
\to \mathcal{F}/\mathfrak m^n \mathcal{F} \to 0
$$
de $\mathcal{O}_X$-modules cohérents. La tensorisation par $\mathcal{L}^{\otimes d}$
est un foncteur exact, et nous obtenons les suites exactes courtes
$$
0 \to
\mathfrak m^n\mathcal{F}/\mathfrak m^{n + 1} \mathcal{F}
\otimes_{\mathcal{O}_X} \mathcal{L}^{\otimes d}
\to \mathcal{F}/\mathfrak m^{n + 1} \mathcal{F}
\otimes_{\mathcal{O}_X} \mathcal{L}^{\otimes d}
\to \mathcal{F}/\mathfrak m^n \mathcal{F}
\otimes_{\mathcal{O}_X} \mathcal{L}^{\otimes d} \to 0
$$
À l'aide de la suite exacte longue de cohomologie et de l'annulation
ci-dessus, nous concluons (par récurrence) que
\begin{enumerate}
\item $H^p(X, \mathcal{F}/\mathfrak m^n \mathcal{F}
\otimes_{\mathcal{O}_X} \mathcal{L}^{\otimes d}) = 0$
pour tous $p > 0$, $d \geq d_0$ et $n \geq 0$, et
\item $H^0(X, \mathcal{F}/\mathfrak m^n \mathcal{F}
\otimes_{\mathcal{O}_X} \mathcal{L}^{\otimes d}) \to
H^0(X_y, \mathcal{F}_y \otimes_{\mathcal{O}_{X_y}} \mathcal{L}_y^{\otimes d})$
est surjective pour tous $d \geq d_0$ et $n \geq 1$.
\end{enumerate}
D'après le théorème sur les fonctions formelles (Théorème \ref{coherent-theorem-formal-functions}),
on voit que la complétion $\mathfrak m$-adique de
$H^p(X, \mathcal{F} \otimes_{\mathcal{O}_X} \mathcal{L}^{\otimes d})$
est nulle pour tous $d \geq d_0$ et $p > 0$.
Comme $H^p(X, \mathcal{F} \otimes_{\mathcal{O}_X} \mathcal{L}^{\otimes d})$
est un $A$-module de type fini d'après
le lemme \ref{coherent-lemma-proper-over-affine-cohomology-finite},
le lemme de Nakayama (Algèbre, lemme \ref{algebra-lemma-NAK}) entraîne
que $H^p(X, \mathcal{F} \otimes_{\mathcal{O}_X} \mathcal{L}^{\otimes d})$
est nul pour tous $d \geq d_0$ et $p > 0$.
Pour $p = 0$, nous déduisons du
lemme \ref{coherent-lemma-ML-cohomology-powers-ideal}, partie (3),
que $H^0(X, \mathcal{F} \otimes_{\mathcal{O}_X} \mathcal{L}^{\otimes d}) \to
H^0(X_y, \mathcal{F}_y \otimes_{\mathcal{O}_{X_y}} \mathcal{L}_y^{\otimes d})$
est surjective, ce qui donne la dernière assertion du lemme.
\end{proof}

\begin{lemma}
\label{coherent-lemma-ample-in-neighbourhood}
(Pour une version plus générale, voir Compléments
sur les morphismes,
lemme \ref{more-morphisms-lemma-ample-in-neighbourhood}.)
Soit $f : X \to Y$ un morphisme propre de schémas, avec $Y$ noethérien.
Soit $\mathcal{L}$ un $\mathcal{O}_X$-module inversible.
Soit $y \in Y$ un point tel que $\mathcal{L}_y$ soit ample
sur $X_y$. Il existe alors un voisinage ouvert $V \subset Y$
de $y$ tel que $\mathcal{L}|_{f^{-1}(V)}$ soit ample sur $f^{-1}(V)/V$.
\end{lemma}

\begin{proof}
Choisissons $d_0$ comme dans le lemme \ref{coherent-lemma-ample-on-fibre} pour
$\mathcal{F} = \mathcal{O}_X$. Choisissons $d \geq d_0$
de sorte que l'on puisse trouver $r \geq 0$ et des sections
$s_{y, 0}, \ldots, s_{y, r} \in H^0(X_y, \mathcal{L}_y^{\otimes d})$
qui définissent une immersion fermée
$$
\varphi_y =
\varphi_{\mathcal{L}_y^{\otimes d}, (s_{y, 0}, \ldots, s_{y, r})} :
X_y \to \mathbf{P}^r_{\kappa(y)}.
$$
Cela est possible d'après Morphismes, Lemme
\ref{morphisms-lemma-finite-type-over-affine-ample-very-ample},
mais nous utilisons également
Morphismes, lemme \ref{morphisms-lemma-image-proper-scheme-closed},
pour voir que $\varphi_y$ est une immersion fermée, et
Constructions, section \ref{constructions-section-projective-space},
pour la description des morphismes vers l'espace
projectif en termes de faisceaux inversibles et de sections.
D'après notre choix de $d_0$, après avoir remplacé $Y$ par un voisinage ouvert
de $y$, nous pouvons choisir
$s_0, \ldots, s_r \in H^0(X, \mathcal{L}^{\otimes d})$
dont les images sont $s_{y, 0}, \ldots, s_{y, r}$.
Soit $X_{s_i} \subset X$ l'ouvert où $s_i$
est un générateur de $\mathcal{L}^{\otimes d}$. Puisque
les $s_{y, i}$ engendrent $\mathcal{L}_y^{\otimes d}$, on voit que
$X_y \subset U = \bigcup X_{s_i}$. Puisque
le morphisme $X \to Y$ est fermé, il existe
un voisinage ouvert $y \in V \subset Y$
tel que $f^{-1}(V) \subset U$.
Après avoir remplacé $Y$ par $V$, nous pouvons supposer que
les $s_i$ engendrent $\mathcal{L}^{\otimes d}$. Nous obtenons
ainsi un morphisme
$$
\varphi = \varphi_{\mathcal{L}^{\otimes d}, (s_0, \ldots, s_r)} :
X \longrightarrow \mathbf{P}^r_Y
$$
tel que $\mathcal{L}^{\otimes d} \cong \varphi^*\mathcal{O}_{\mathbf{P}^r_Y}(1)$
et dont le changement de base à $y$ donne $\varphi_y$.

\medskip\noindent
Nous terminerons la démonstration par une astuce ; la démonstration « correcte »
consiste à montrer directement que $\varphi$ est une immersion
fermée après changement de base à un voisinage ouvert de $y$. En
effet, d'après le lemme \ref{coherent-lemma-proper-finite-fibre-finite-in-neighbourhood},
on voit que $\varphi$ est fini au-dessus d'un voisinage ouvert
de la fibre $\mathbf{P}^r_{\kappa(y)}$ de $\mathbf{P}^r_Y \to Y$
au-dessus de $y$. Puisque $\mathbf{P}^r_Y \to Y$ est fermé, quitte à
rétrécir $Y$, nous pouvons supposer que $\varphi$ est fini.
Alors $\mathcal{L}^{\otimes d} \cong \varphi^*\mathcal{O}_{\mathbf{P}^r_Y}(1)$
est ample d'après le résultat très général suivant:
Morphismes, lemme \ref{morphisms-lemma-pullback-ample-tensor-relatively-ample}.
\end{proof}





\section{Cohomologie et changement de base, III}
\label{coherent-section-cohomology-and-base-change-perfect}

\noindent
Dans cette section, nous démontrons le cas le plus simple d'un phénomène très
général qui sera étudié dans
Catégories dérivées des schémas, section
\ref{perfect-section-cohomology-and-base-change-perfect}.
Voir la remarque \ref{coherent-remark-explain-perfect-direct-image} pour une
traduction algébrique du lemme suivant.

\begin{lemma}
\label{coherent-lemma-perfect-direct-image}
Soit $A$ un anneau noethérien et posons $S = \Spec(A)$. Soit $f : X \to S$ un morphisme
propre de schémas. Soit $\mathcal{F}$ un
$\mathcal{O}_X$-module cohérent et plat sur $S$. Alors
\begin{enumerate}
\item $R\Gamma(X, \mathcal{F})$ est un objet parfait de $D(A)$, et
\item pour tout homomorphisme d'anneaux $A \to A'$, le morphisme de changement de base
$$
R\Gamma(X, \mathcal{F}) \otimes_A^{\mathbf{L}} A'
\longrightarrow
R\Gamma(X_{A'}, \mathcal{F}_{A'})
$$
est un isomorphisme.
\end{enumerate}
\end{lemma}

\begin{proof}
Choisissons un recouvrement ouvert affine fini $X = \bigcup_{i = 1, \ldots, n} U_i$.
D'après les lemmes \ref{coherent-lemma-separated-case-relative-cech} et
\ref{coherent-lemma-base-change-complex}, le complexe de {\v C}ech
$K^\bullet = {\check C}^\bullet(\mathcal{U}, \mathcal{F})$ vérifie
$$
K^\bullet \otimes_A A' = R\Gamma(X_{A'}, \mathcal{F}_{A'})
$$
pour tout homomorphisme d'anneaux $A \to A'$. Soit
$K_{alt}^\bullet = {\check C}_{alt}^\bullet(\mathcal{U}, \mathcal{F})$
le complexe de {\v C}ech alterné. D'après
Cohomologie, lemme \ref{cohomology-lemma-alternating-usual},
il existe une équivalence d'homotopie $K_{alt}^\bullet \to K^\bullet$
de $A$-modules. En particulier, on a
$$
K_{alt}^\bullet \otimes_A A' = R\Gamma(X_{A'}, \mathcal{F}_{A'})
$$
également. Comme $\mathcal{F}$ est plat sur $A$, on voit que chaque $K_{alt}^n$
est plat sur $A$ (voir
Morphismes, lemme \ref{morphisms-lemma-flat-module-characterize}).
Comme, de plus, $K_{alt}^\bullet$ est borné supérieurement (c'est pourquoi nous sommes passés
au complexe de {\v C}ech alterné),
$K_{alt}^\bullet \otimes_A A' = K_{alt}^\bullet \otimes_A^{\mathbf{L}} A'$
d'après la définition des produits tensoriels dérivés (voir
Compléments sur l'algèbre, section \ref{more-algebra-section-derived-tensor-product}).
D'après le
lemme \ref{coherent-lemma-proper-over-affine-cohomology-finite},
les groupes de cohomologie $H^i(K_{alt}^\bullet)$ sont des $A$-modules de
type fini. Comme $K_{alt}^\bullet$ est borné, nous en déduisons que $K_{alt}^\bullet$
est pseudo-cohérent, voir
Compléments sur l'algèbre, lemme \ref{more-algebra-lemma-Noetherian-pseudo-coherent}.
Pour tout $A$-module $M$, posons $A' = A \oplus M$, où $M$ est un idéal de carré
nul, c'est-à-dire $(a, m) \cdot (a', m') = (aa', am' + a'm)$. D'après ce qui précède,
on voit que la cohomologie de $K_{alt}^\bullet \otimes_A^\mathbf{L} A'$ est concentrée
dans les degrés $0, \ldots, n$. De même, celle de $K_{alt}^\bullet \otimes_A^\mathbf{L} M$
est concentrée dans les degrés $0, \ldots, n$. Par conséquent, $K_{alt}^\bullet$ a
une dimension de Tor finie,
voir Compléments sur l'algèbre, définition \ref{more-algebra-definition-tor-amplitude}.
La conclusion résulte de Compléments sur l'algèbre, lemme \ref{more-algebra-lemma-perfect}.
\end{proof}

\begin{remark}
\label{coherent-remark-explain-perfect-direct-image}
Une conséquence du lemme \ref{coherent-lemma-perfect-direct-image} est qu'il
existe un complexe fini de $A$-modules projectifs de type fini $M^\bullet$ tel
que l'on ait
$$
H^i(X_{A'}, \mathcal{F}_{A'}) = H^i(M^\bullet \otimes_A A')
$$
fonctoriellement en $A'$. La condition que $\mathcal{F}$ soit
plat sur $A$ est essentielle, voir \cite{Hartshorne}.
\end{remark}







\section{Modules formels cohérents}
\label{coherent-section-coherent-formal}

\noindent
Comme nous ne disposons pas encore de la théorie des schémas formels, nous développons
un peu de terminologie qui remplace la notion de ``module cohérent sur un
schéma formel adique noethérien''.

\medskip\noindent
Soit $X$ un schéma noethérien et soit $\mathcal{I} \subset \mathcal{O}_X$
un faisceau quasi-cohérent d'idéaux. Nous considérerons des systèmes
projectifs $(\mathcal{F}_n)$ de $\mathcal{O}_X$-modules cohérents tels que
\begin{enumerate}
\item $\mathcal{F}_n$ soit annulé par $\mathcal{I}^n$, et
\item les morphismes de transition induisent des isomorphismes
$\mathcal{F}_{n + 1}/\mathcal{I}^n\mathcal{F}_{n + 1} \to \mathcal{F}_n$.
\end{enumerate}
Un morphisme de tels systèmes projectifs est défini de la
manière habituelle. Notons la catégorie de ces systèmes projectifs
$\textit{Coh}(X, \mathcal{I})$. Nous allons établir une série de lemmes
sur les objets de cette catégorie. En fait, la plupart des
lemmes qui suivent sont des conséquences immédiates de la description
suivante de la catégorie dans le cas affine.

\begin{lemma}
\label{coherent-lemma-inverse-systems-affine}
Si $X = \Spec(A)$ est le spectre d'un anneau noethérien et si
$\mathcal{I}$ est le faisceau quasi-cohérent d'idéaux associé à l'idéal
$I \subset A$, alors $\textit{Coh}(X, \mathcal{I})$ est équivalente à
la catégorie des $A^\wedge$-modules de type fini, où $A^\wedge$ est le complété
de $A$ par rapport à $I$.
\end{lemma}

\begin{proof}
Soit $\text{Mod}^{fg}_{A, I}$ la catégorie des systèmes projectifs $(M_n)$
de $A$-modules de type fini satisfaisant : (1) $M_n$ est annulé par $I^n$ et (2)
$M_{n + 1}/I^nM_{n + 1} = M_n$. D'après la correspondance entre les faisceaux
cohérents sur $X$ et les $A$-modules de type fini (lemme \ref{coherent-lemma-coherent-Noetherian}),
il suffit de montrer que $\text{Mod}^{fg}_{A, I}$ est équivalente à la catégorie
des $A^\wedge$-modules de type fini. Pour le voir, il suffit de prouver qu'étant
donné un objet $(M_n)$ de $\text{Mod}^{fg}_{A, I}$, le module
$$
M = \lim M_n
$$
est un $A^\wedge$-module de type fini et que $M/I^nM = M_n$. Les morphismes de
transition étant surjectifs, nous voyons que $M \to M_1$ est surjectif.
Choisissons $x_1, \ldots, x_t \in M$ dont les images engendrent $M_1$.
Cela induit un morphisme de systèmes $(A/I^n)^{\oplus t} \to M_n$.
D'après le lemme de Nakayama (Algèbre, lemme \ref{algebra-lemma-NAK}), ces morphismes sont
surjectifs. Soit $K_n \subset (A/I^n)^{\oplus t}$ le noyau. La
propriété (2) implique que $K_{n + 1} \to K_n$ est surjectif ; en particulier,
le système $(K_n)$ satisfait la condition de Mittag-Leffler.
D'après Homologie, lemme \ref{homology-lemma-Mittag-Leffler},
nous obtenons une suite exacte
$0 \to K \to (A^\wedge)^{\oplus t} \to M \to 0$
où $K = \lim K_n$.
Ainsi, $M$ est un $A^\wedge$-module de
type fini. Comme $K \to K_n$ est surjectif, il s'ensuit que
$$
M/I^nM = \Coker(K \to (A/I^n)^{\oplus t}) =
(A/I^n)^{\oplus t}/K_n = M_n
$$
comme souhaité.
\end{proof}

\begin{lemma}
\label{coherent-lemma-inverse-systems-abelian}
Soit $X$ un schéma noethérien et soit $\mathcal{I} \subset \mathcal{O}_X$
un faisceau quasi-cohérent d'idéaux.
\begin{enumerate}
\item La catégorie $\textit{Coh}(X, \mathcal{I})$ est abélienne.
\item Pour tout ouvert $U \subset X$, le foncteur de restriction
$\textit{Coh}(X, \mathcal{I}) \to \textit{Coh}(U, \mathcal{I}|_U)$
est exact.
\item L'exactitude dans $\textit{Coh}(X, \mathcal{I})$ peut se vérifier
après restriction aux ouverts d'un recouvrement ouvert de $X$.
\end{enumerate}
\end{lemma}

\begin{proof}
Soit $\alpha =(\alpha_n) : (\mathcal{F}_n) \to (\mathcal{G}_n)$ un morphisme dans
$\textit{Coh}(X, \mathcal{I})$. Le conoyau de $\alpha$ est le système projectif
$(\Coker(\alpha_n))$ (détails omis). Pour décrire le noyau, posons
$$
\mathcal{K}'_{l, m} = \Im(\Ker(\alpha_l) \to \mathcal{F}_m)
$$
pour $l \geq m$. Nous
affirmons que :
\begin{enumerate}
\item[(a)] le système projectif $(\mathcal{K}'_{l, m})_{l \geq m}$ est constant
à partir d'un certain rang, disons de valeur $\mathcal{K}'_m$,
\item[(b)] le système $(\mathcal{K}'_m/\mathcal{I}^n\mathcal{K}'_m)_{m \geq n}$
est constant à partir d'un certain rang, disons de valeur $\mathcal{K}_n$,
\item[(c)] le système $(\mathcal{K}_n)$ constitue un objet de
$\textit{Coh}(X, \mathcal{I})$, et
\item[(d)] cet objet est le noyau de $\alpha$.
\end{enumerate}
Pour établir (a), (b) et (c), nous pouvons raisonner localement sur un ouvert affine, disons $X = \Spec(A)$,
et $\mathcal{I}$ correspond à l'idéal $I \subset A$. D'après le
lemme \ref{coherent-lemma-inverse-systems-affine},
$\alpha$ correspond à un morphisme $f : M \to N$ de $A^\wedge$-modules de
type fini. Notons $K = \Ker(f)$. Remarquons que $A^\wedge$ est un anneau
noethérien (Algèbre, lemme \ref{algebra-lemma-completion-Noetherian-Noetherian}).
Choisissons un entier $c \geq 0$ tel que
$K \cap I^n M \subset I^{n - c}K$ pour $n \geq c$
(Algèbre, lemme \ref{algebra-lemma-Artin-Rees})
et choisissons-le de sorte que le lemme \ref{algebra-lemma-map-AR}
d'Algèbre s'applique pour le morphisme $f$ et l'idéal $I^\wedge = IA^\wedge$. Alors
$\mathcal{K}'_{l, m}$ correspond au $A$-module
$$
K'_{l, m} = \frac{a^{-1}(I^lN) + I^mM}{I^mM} =
\frac{K + I^{l - c}f^{-1}(I^cN) + I^mM}{I^mM} =
\frac{K + I^mM}{I^mM}
$$
où la dernière égalité a lieu si $l \geq m + c$. Ainsi, $\mathcal{K}'_m$
correspond au $A$-module $K/K \cap I^mM$ et
$\mathcal{K}'_m/\mathcal{I}^n\mathcal{K}'_m$ correspond à
$$
\frac{K}{K \cap I^mM + I^nK} = \frac{K}{I^nK}
$$
pour $m \geq n + c$, d'après le choix de $c$ ci-dessus. Par conséquent, $\mathcal{K}_n$
correspond à $K/I^nK$.

\medskip\noindent
Démontrons (d). Il résulte clairement de la description affine ci-dessus que
la composée $(\mathcal{K}_n) \to (\mathcal{F}_n) \to (\mathcal{G}_n)$
est nulle. Soit $\beta : (\mathcal{H}_n) \to (\mathcal{F}_n)$
un morphisme tel que $\alpha \circ \beta = 0$. Alors l'image de
$\mathcal{H}_l \to \mathcal{F}_l$ est contenue dans $\Ker(\alpha_l)$.
Puisque $\mathcal{H}_m = \mathcal{H}_l/\mathcal{I}^m\mathcal{H}_l$
pour $l \geq m$, nous obtenons un système de morphismes
$\mathcal{H}_m \to \mathcal{K}'_{l, m}$. D'où un morphisme
$\mathcal{H}_m \to \mathcal{K}_m'$. Puisque
$\mathcal{H}_n = \mathcal{H}_m/\mathcal{I}^n\mathcal{H}_m$, nous
obtenons un système de morphismes $\mathcal{H}_n \to \mathcal{K}'_m/\mathcal{I}^n\mathcal{K}'_m$
et donc un morphisme $\mathcal{H}_n \to \mathcal{K}_n$, comme souhaité.

\medskip\noindent
Pour achever la preuve de (1), il reste à montrer que $\Coim = \Im$
dans $\textit{Coh}(X, \mathcal{I})$. Nous avons vu ci-dessus que la formation des
noyaux et des conoyaux est compatible, sur les ouverts affines, avec la description
de $\textit{Coh}(X, \mathcal{I})$ comme catégorie de modules. Puisque l'égalité
$\Im = \Coim$ est vérifiée dans la catégorie des modules,
nous obtenons $\Coim = \Im$ dans $\textit{Coh}(X, \mathcal{I})$.
Les assertions (2) et (3) du lemme découlent immédiatement de notre construction des
noyaux et des conoyaux.
\end{proof}

\begin{lemma}
\label{coherent-lemma-inverse-systems-surjective}
Soit $X$ un schéma noethérien et soit $\mathcal{I} \subset \mathcal{O}_X$
un faisceau quasi-cohérent d'idéaux. Un morphisme
$(\mathcal{F}_n) \to (\mathcal{G}_n)$ est surjectif dans
$\textit{Coh}(X, \mathcal{I})$
si et seulement si $\mathcal{F}_1 \to \mathcal{G}_1$ est surjectif.
\end{lemma}

\begin{proof}
Omis. Indication : se placer sur les ouverts affines, utiliser
le lemme \ref{coherent-lemma-inverse-systems-affine}, ainsi que
Algèbre, lemme \ref{algebra-lemma-NAK}.
\end{proof}

\noindent
Soit $X$ un schéma noethérien et soit $\mathcal{I} \subset \mathcal{O}_X$
un faisceau quasi-cohérent d'idéaux. Il existe un foncteur
\begin{equation}
\label{coherent-equation-completion-functor}
\textit{Coh}(\mathcal{O}_X) \longrightarrow \textit{Coh}(X, \mathcal{I}), \quad
\mathcal{F} \longmapsto \mathcal{F}^\wedge
\end{equation}
qui associe au $\mathcal{O}_X$-module cohérent $\mathcal{F}$
l'objet $\mathcal{F}^\wedge = (\mathcal{F}/\mathcal{I}^n\mathcal{F})$
de $\textit{Coh}(X, \mathcal{I})$.

\begin{lemma}
\label{coherent-lemma-exact}
Le foncteur (\ref{coherent-equation-completion-functor}) est exact.
\end{lemma}

\begin{proof}
Il suffit de le vérifier localement sur $X$. Nous pouvons donc supposer que $X$ est
affine, c'est-à-dire que nous sommes dans la situation du
lemme \ref{coherent-lemma-inverse-systems-affine}.
Le foncteur est le foncteur $\text{Mod}^{fg}_A \to \text{Mod}^{fg}_{A^\wedge}$
qui associe à un $A$-module de type fini $M$ sa complétion $M^\wedge$.
Le résultat découle donc de
Algèbre, lemme \ref{algebra-lemma-completion-flat}.
\end{proof}

\begin{lemma}
\label{coherent-lemma-completion-internal-hom}
Soit $X$ un schéma noethérien et soit $\mathcal{I} \subset \mathcal{O}_X$
un faisceau quasi-cohérent d'idéaux. Soient $\mathcal{F}$, $\mathcal{G}$
des $\mathcal{O}_X$-modules cohérents. Posons
$\mathcal{H} = \SheafHom_{\mathcal{O}_X}(\mathcal{G}, \mathcal{F})$.
Alors
$$
\lim H^0(X, \mathcal{H}/\mathcal{I}^n\mathcal{H}) =
\Mor_{\textit{Coh}(X, \mathcal{I})}
(\mathcal{G}^\wedge, \mathcal{F}^\wedge).
$$
\end{lemma}

\begin{proof}
Pour le démontrer, nous pouvons travailler localement sur des ouverts affines de $X$.
Nous pouvons donc supposer $X = \Spec(A)$ et $\mathcal{F}$, $\mathcal{G}$
donnés par les $A$-modules de type fini $M$ et $N$. Alors $\mathcal{H}$
correspond au $A$-module de type fini $H = \Hom_A(M, N)$.
L'énoncé du lemme devient l'égalité
$$
H^\wedge = \Hom_{A^\wedge}(M^\wedge, N^\wedge)
$$
par l'équivalence du lemme \ref{coherent-lemma-inverse-systems-affine}.
D'après Algèbre, lemme \ref{algebra-lemma-completion-flat}
(utilisé 3 fois), nous avons
$$
H^\wedge = \Hom_A(M, N) \otimes_A A^\wedge =
\Hom_{A^\wedge}(M \otimes_A A^\wedge, N \otimes_A A^\wedge) =
\Hom_{A^\wedge}(M^\wedge, N^\wedge)
$$
où la deuxième égalité utilise le fait que $A^\wedge$ est plat sur $A$
(voir Compléments d'algèbre, Lemme
\ref{more-algebra-lemma-pseudo-coherence-and-base-change-ext}).
Le lemme en découle.
\end{proof}

\noindent
Soit $X$ un schéma noethérien. Soit $\mathcal{I} \subset \mathcal{O}_X$
un faisceau quasi-cohérent d'idéaux. Nous disons qu'un objet $(\mathcal{F}_n)$ de
$\textit{Coh}(X, \mathcal{I})$ est {\it de torsion annulée par une puissance de $\mathcal{I}$}
ou est {\it annulé par une puissance de $\mathcal{I}$} s'il existe
un $c \geq 1$ tel que $\mathcal{F}_n = \mathcal{F}_c$ pour tout $n \geq c$.
Dans ce cas, nous dirons que $(\mathcal{F}_n)$ est {\it annulé par
$\mathcal{I}^c$}. Si $X = \Spec(A)$ est affine, alors, par l'équivalence du
lemme \ref{coherent-lemma-inverse-systems-affine},
ces objets correspondent exactement aux $A$-modules de
type fini annulés par une puissance de $I$ ou par $I^c$.

\begin{lemma}
\label{coherent-lemma-existence-easy}
Soit $X$ un schéma noethérien et soit $\mathcal{I} \subset \mathcal{O}_X$
un faisceau quasi-cohérent d'idéaux. Soit $\mathcal{G}$ un
$\mathcal{O}_X$-module cohérent. Soit $(\mathcal{F}_n)$ un objet de
$\textit{Coh}(X, \mathcal{I})$.
\begin{enumerate}
\item Si $\alpha : (\mathcal{F}_n) \to \mathcal{G}^\wedge$ est
un morphisme dont le noyau et le conoyau sont annulés par une puissance de $\mathcal{I}$,
alors il existe un triplet unique (à isomorphisme unique près)
$(\mathcal{F}, a, \beta)$ où
\begin{enumerate}
\item $\mathcal{F}$ est un $\mathcal{O}_X$-module cohérent,
\item $a : \mathcal{F} \to \mathcal{G}$ est un morphisme de $\mathcal{O}_X$-modules dont
le noyau et le conoyau sont annulés par une puissance de $\mathcal{I}$,
\item $\beta : (\mathcal{F}_n) \to \mathcal{F}^\wedge$ est un isomorphisme, et
\item $\alpha = a^\wedge \circ \beta$.
\end{enumerate}
\item Si $\alpha : \mathcal{G}^\wedge \to (\mathcal{F}_n)$ est
un morphisme dont le noyau et le conoyau sont annulés par une puissance de $\mathcal{I}$,
alors il existe un triplet unique (à isomorphisme unique près)
$(\mathcal{F}, a, \beta)$ où
\begin{enumerate}
\item $\mathcal{F}$ est un $\mathcal{O}_X$-module cohérent,
\item $a : \mathcal{G} \to \mathcal{F}$ est un morphisme de $\mathcal{O}_X$-modules dont
le noyau et le conoyau sont annulés par une puissance de $\mathcal{I}$,
\item $\beta : \mathcal{F}^\wedge \to (\mathcal{F}_n)$ est un isomorphisme, et
\item $\alpha = \beta \circ a^\wedge$.
\end{enumerate}
\end{enumerate}
\end{lemma}

\begin{proof}
Démonstration de (1). L'unicité montre qu'il suffit de construire
$(\mathcal{F}, a, \beta)$
localement pour la topologie de Zariski sur $X$. Nous pouvons donc supposer $X = \Spec(A)$ et que $\mathcal{I}$
correspond à l'idéal $I \subset A$. Dans cette situation, le
lemme \ref{coherent-lemma-inverse-systems-affine} s'applique.
Soit $M'$ le $A^\wedge$-module de type fini
correspondant à $(\mathcal{F}_n)$. Soit $N$ le $A$-module de type fini correspondant à
$\mathcal{G}$. Alors $\alpha$ correspond à un morphisme
$$
\varphi : M' \longrightarrow N^\wedge
$$
dont le noyau et le conoyau sont annulés par $I^t$ pour un certain $t$. Rappelons que
$N^\wedge = N \otimes_A A^\wedge$
(Algèbre, lemme \ref{algebra-lemma-completion-tensor}).
D'après Compléments d'algèbre, lemme \ref{more-algebra-lemma-application-formal-glueing},
il existe un morphisme de $A$-modules $\psi : M \to N$ dont le noyau et le conoyau sont de torsion annulée par une puissance de
$I$, et un isomorphisme
$M \otimes_A A^\wedge = M'$ compatible avec $\varphi$.
Comme $N$ et $M'$ sont des modules de type fini, nous concluons que $M$
est un $A$-module de type
fini, voir Compléments d'algèbre, remarque \ref{more-algebra-remark-formal-glueing-algebras}.
Ainsi $M \otimes_A A^\wedge = M^\wedge$. Nous omettons la vérification
que le triplet $(M, N \to M, M^\wedge \to M')$ ainsi obtenu
est unique à isomorphisme unique près.

\medskip\noindent
La démonstration de (2) est exactement la même et nous l'omettons.
\end{proof}

\begin{lemma}
\label{coherent-lemma-torsion-hom-ext}
Soit $X$ un schéma noethérien et soit $\mathcal{I} \subset \mathcal{O}_X$
un faisceau quasi-cohérent d'idéaux. Tout objet de
$\textit{Coh}(X, \mathcal{I})$ qui est annulé
par une puissance de $\mathcal{I}$ appartient à l'image essentielle de
(\ref{coherent-equation-completion-functor}).
De plus, si $\mathcal{F}$, $\mathcal{G}$ appartiennent à $\textit{Coh}(\mathcal{O}_X)$
et si $\mathcal{F}$ ou $\mathcal{G}$ est annulé par une puissance de
$\mathcal{I}$, alors les applications
$$
\xymatrix{
\Hom_X(\mathcal{F}, \mathcal{G}) \ar[d] &
\Ext_X(\mathcal{F}, \mathcal{G}) \ar[d] \\
\Hom_{\textit{Coh}(X, \mathcal{I})}(\mathcal{F}^\wedge, \mathcal{G}^\wedge) &
\Ext_{\textit{Coh}(X, \mathcal{I})}(\mathcal{F}^\wedge, \mathcal{G}^\wedge)
}
$$
sont des isomorphismes.
\end{lemma}

\begin{proof}
Supposons que $(\mathcal{F}_n)$ soit un objet de $\textit{Coh}(X, \mathcal{I})$
annulé par $\mathcal{I}^c$ pour un certain $c \geq 1$. Alors
$\mathcal{F}_n \to \mathcal{F}_c$ est un isomorphisme pour $n \geq c$.
Ainsi, si nous posons $\mathcal{F} = \mathcal{F}_c$, nous voyons que
$\mathcal{F}^\wedge \cong (\mathcal{F}_n)$. Cela prouve la première assertion.

\medskip\noindent
Soient $\mathcal{F}$, $\mathcal{G}$ des objets de $\textit{Coh}(\mathcal{O}_X)$
tels que $\mathcal{F}$ ou $\mathcal{G}$ soit annulé par
$\mathcal{I}^c$ pour un certain $c \geq 1$. Alors
$\mathcal{H} = \SheafHom_{\mathcal{O}_X}(\mathcal{G}, \mathcal{F})$
est un $\mathcal{O}_X$-module cohérent annulé par $\mathcal{I}^c$.
Nous voyons donc que
$$
\Hom_X(\mathcal{G}, \mathcal{F}) =
H^0(X, \mathcal{H}) =
\lim H^0(X, \mathcal{H}/\mathcal{I}^n\mathcal{H}) =
\Mor_{\textit{Coh}(X, \mathcal{I})}
(\mathcal{G}^\wedge, \mathcal{F}^\wedge).
$$
Cela résulte du lemme \ref{coherent-lemma-completion-internal-hom}.
Cela prouve l'énoncé concernant les homomorphismes.

\medskip\noindent
La notation $\Ext$ désigne les extensions définies dans
Homologie, section \ref{homology-section-extensions}.
L'injectivité de l'application sur les $\Ext$ découle immédiatement de
la bijectivité de l'application sur les $\Hom$. Pour
la surjectivité, supposons que $\mathcal{F}$ soit annulé
par une puissance de $I$. Alors la partie (1) du lemme \ref{coherent-lemma-existence-easy}
montre que, étant donnée une extension
$$
0 \to \mathcal{G}^\wedge \to (\mathcal{E}_n) \to \mathcal{F}^\wedge \to 0
$$
dans $\textit{Coh}(U, I\mathcal{O}_U)$,
le morphisme $\mathcal{G}^\wedge \to (\mathcal{E}_n)$ est
isomorphe à $\mathcal{G} \to \mathcal{E}^\wedge$
pour un certain $\mathcal{G} \to \mathcal{E}$ dans $\textit{Coh}(\mathcal{O}_U)$.
On procède de même dans l'autre cas.
\end{proof}

\begin{lemma}
\label{coherent-lemma-finite-over-rees-algebra}
Soit $X$ un schéma noethérien et soit $\mathcal{I} \subset \mathcal{O}_X$
un faisceau quasi-cohérent d'idéaux. Si $(\mathcal{F}_n)$ est un objet de
$\textit{Coh}(X, \mathcal{I})$, alors
$\bigoplus \Ker(\mathcal{F}_{n + 1} \to \mathcal{F}_n)$ est
un
$\bigoplus \mathcal{I}^n/\mathcal{I}^{n + 1}$-module gradué quasi-cohérent de type fini.
\end{lemma}

\begin{proof}
La question est locale sur $X$ ; nous pouvons donc supposer $X$ affine, c'est-à-dire que
nous sommes dans la situation du lemme \ref{coherent-lemma-inverse-systems-affine}.
Dans ce cas, si $(\mathcal{F}_n)$ correspond au $A^\wedge$
-module de type fini $M$, alors $\bigoplus \Ker(\mathcal{F}_{n + 1} \to \mathcal{F}_n)$
correspond à $\bigoplus I^nM/I^{n + 1}M$, qui est manifestement un module de
type fini sur $\bigoplus I^n/I^{n + 1}$.
\end{proof}

\begin{lemma}
\label{coherent-lemma-inverse-systems-pullback}
Soit $f : X \to Y$ un morphisme de schémas noethériens.
Soit $\mathcal{J} \subset \mathcal{O}_Y$ un faisceau quasi-cohérent
d'idéaux et posons $\mathcal{I} = f^{-1}\mathcal{J} \mathcal{O}_X$.
Il existe alors un foncteur exact à droite
$$
f^* : \textit{Coh}(Y, \mathcal{J}) \longrightarrow \textit{Coh}(X, \mathcal{I})
$$
qui envoie $(\mathcal{G}_n)$ sur $(f^*\mathcal{G}_n)$. Si $f$ est plat,
alors $f^*$ est un foncteur exact.
\end{lemma}

\begin{proof}
Puisque $f^* : \textit{Coh}(\mathcal{O}_Y) \to \textit{Coh}(\mathcal{O}_X)$
est exact à droite, nous avons
$$
f^*\mathcal{G}_n =
f^*(\mathcal{G}_{n + 1}/\mathcal{I}^n\mathcal{G}_{n + 1}) =
f^*\mathcal{G}_{n + 1}/f^{-1}\mathcal{I}^nf^*\mathcal{G}_{n + 1} =
f^*\mathcal{G}_{n + 1}/\mathcal{J}^nf^*\mathcal{G}_{n + 1}
$$
donc l'image inverse d'un système est un système. La construction des
conoyaux dans la démonstration du lemme \ref{coherent-lemma-inverse-systems-abelian}
montre que
$f^* : \textit{Coh}(Y, \mathcal{J}) \to \textit{Coh}(X, \mathcal{I})$
est toujours exact à droite. Si $f$ est plat, alors
$f^* : \textit{Coh}(\mathcal{O}_Y) \to \textit{Coh}(\mathcal{O}_X)$
est un foncteur exact. Il résulte de la construction des noyaux
dans la démonstration du lemme \ref{coherent-lemma-inverse-systems-abelian}
que, dans ce cas,
$f^* : \textit{Coh}(Y, \mathcal{J}) \to \textit{Coh}(X, \mathcal{I})$
envoie également les noyaux sur les noyaux.
\end{proof}

\begin{lemma}
\label{coherent-lemma-inverse-systems-pullback-equivalence}
Soit $f : X' \to X$ un morphisme de schémas noethériens. Soit $Z \subset X$
un sous-schéma fermé et notons $Z' = f^{-1}Z$ l'image inverse
schématique. Soient $\mathcal{I} \subset \mathcal{O}_X$,
$\mathcal{I}' \subset \mathcal{O}_{X'}$ les faisceaux
quasi-cohérents d'idéaux correspondants.
Si $f$ est plat et que le morphisme induit $Z' \to Z$
est un isomorphisme, alors le foncteur d'image inverse
$f^* : \textit{Coh}(X, \mathcal{I}) \to \textit{Coh}(X', \mathcal{I}')$
(lemme \ref{coherent-lemma-inverse-systems-pullback})
est une équivalence.
\end{lemma}

\begin{proof}
Si $X$ et $X'$ sont affines, cela découle immédiatement de
Compléments d'algèbre, lemme \ref{more-algebra-lemma-neighbourhood-equivalence}.
Pour le démontrer en général, soient $Z_n \subset X$, $Z'_n \subset X'$
les $n$-ièmes voisinages infinitésimaux de $Z$, $Z'$.
Le morphisme $Z_n \to Z'_n$ induit un homéomorphisme entre
les espaces topologiques sous-jacents. D'autre part, si $z' \in Z'$
s'envoie sur $z \in Z$, alors le morphisme d'anneaux
$\mathcal{O}_{X, z} \to \mathcal{O}_{X', z'}$ est plat
et induit un isomorphisme
$\mathcal{O}_{X, z}/\mathcal{I}_z \to \mathcal{O}_{X', z'}/\mathcal{I}'_{z'}$.
Il induit donc un isomorphisme
$\mathcal{O}_{X, z}/\mathcal{I}_z^n \to
\mathcal{O}_{X', z'}/(\mathcal{I}'_{z'})^n$
pour tout $n \geq 1$, par exemple
d'après Compléments d'algèbre, lemme \ref{more-algebra-lemma-neighbourhood-isomorphism}.
Ainsi $Z'_n \to Z_n$ est un isomorphisme de schémas.
Ainsi $f^*$ induit une équivalence entre
la catégorie des $\mathcal{O}_X$-modules cohérents annulés par $\mathcal{I}^n$
et la
catégorie des $\mathcal{O}_{X'}$-modules cohérents annulés par
$(\mathcal{I}')^n$, voir le
lemme \ref{coherent-lemma-i-star-equivalence}.
Cela entraîne clairement le lemme.
\end{proof}

\begin{lemma}
\label{coherent-lemma-inverse-systems-ideals-equivalence}
Soit $X$ un schéma noethérien. Soient
$\mathcal{I}, \mathcal{J} \subset \mathcal{O}_X$
des faisceaux quasi-cohérents d'idéaux.
Si $V(\mathcal{I}) = V(\mathcal{J})$ comme sous-ensembles fermés
de $X$, alors $\textit{Coh}(X, \mathcal{I})$ et $\textit{Coh}(X, \mathcal{J})$
sont équivalentes.
\end{lemma}

\begin{proof}
Supposons d'abord que $X = \Spec(A)$ soit affine. Soient $I, J \subset A$ les idéaux
correspondant à $\mathcal{I}, \mathcal{J}$. Alors $V(I) = V(J)$
implique que $I^c \subset J$ et $J^d \subset I$ pour certains $c, d \geq 1$
par les propriétés élémentaires de la topologie de
Zariski (voir Algèbre, section \ref{algebra-section-spectrum-ring} et
lemme \ref{algebra-lemma-Noetherian-power}).
Les complétions $I$-adique et $J$-adique de $A$ coïncident donc,
voir Algèbre, lemme \ref{algebra-lemma-change-ideal-completion}.
L'équivalence découle ainsi du lemme \ref{coherent-lemma-inverse-systems-affine}
dans ce cas.

\medskip\noindent
En général, d'après ce qui précède et puisque
$X$ est quasi-compact, on peut choisir $c, d \geq 1$ tels que
$\mathcal{I}^c \subset \mathcal{J}$ et $\mathcal{J}^d \subset \mathcal{I}$.
Alors, étant donné un objet $(\mathcal{F}_n)$ de
$\textit{Coh}(X, \mathcal{I})$, nous affirmons que le
système projectif
$$
(\mathcal{F}_{cn}/\mathcal{J}^n\mathcal{F}_{cn})
$$
appartient à $\textit{Coh}(X, \mathcal{J})$. Cela peut se vérifier sur
les éléments d'un recouvrement affine ; nous omettons les détails.
De la même manière, nous pouvons construire un objet de
$\textit{Coh}(X, \mathcal{I})$ à partir d'un objet de
$\textit{Coh}(X, \mathcal{J})$. Nous omettons la vérification que
ces constructions définissent des foncteurs quasi-inverses l'un de l'autre.
\end{proof}





\section{Théorème d'existence de Grothendieck, I}
\label{coherent-section-existence}

\noindent
Dans cette section, nous étudions le théorème d'existence de Grothendieck dans
le cas projectif. Nous utiliserons la notion de modules formels cohérents développée
à la section \ref{coherent-section-coherent-formal}. Le lecteur familier des
schémas formels est invité à lire l'énoncé et la démonstration
du théorème dans \cite{EGA}.

\begin{lemma}
\label{coherent-lemma-fully-faithful}
Soit $A$ un anneau noethérien complet pour la topologie $I$-adique.
Soit $f : X \to \Spec(A)$ un morphisme propre. Posons
$\mathcal{I} = I\mathcal{O}_X$.
Alors le foncteur (\ref{coherent-equation-completion-functor}) est pleinement fidèle.
\end{lemma}

\begin{proof}
Soient $\mathcal{F}$ et $\mathcal{G}$ des $\mathcal{O}_X$-modules cohérents.
Alors $\mathcal{H} = \SheafHom_{\mathcal{O}_X}(\mathcal{G}, \mathcal{F})$
est un $\mathcal{O}_X$-module cohérent,
voir Modules, lemme \ref{modules-lemma-internal-hom-locally-kernel-direct-sum}.
D'après le lemme \ref{coherent-lemma-completion-internal-hom}, l'application
$$
\lim_n H^0(X, \mathcal{H}/\mathcal{I}^n\mathcal{H})
\to
\Mor_{\textit{Coh}(X, \mathcal{I})}
(\mathcal{G}^\wedge, \mathcal{F}^\wedge)
$$
est bijective. Par conséquent, la pleine fidélité de
(\ref{coherent-equation-completion-functor}) résulte du théorème des fonctions
formelles (lemme \ref{coherent-lemma-spell-out-theorem-formal-functions})
appliqué au faisceau cohérent $\mathcal{H}$.
\end{proof}

\begin{lemma}
\label{coherent-lemma-vanishing-projective}
Soit $A$ un anneau noethérien et soit $I \subset A$ un idéal.
Soit $f : X \to \Spec(A)$ un morphisme propre et soit
$\mathcal{L}$ un faisceau inversible $f$-ample. Posons
$\mathcal{I} = I\mathcal{O}_X$. Soit $(\mathcal{F}_n)$ un objet
de $\textit{Coh}(X, \mathcal{I})$. Il existe alors un
entier $d_0$ tel que
$$
H^1(X, \Ker(\mathcal{F}_{n + 1} \to \mathcal{F}_n)
\otimes \mathcal{L}^{\otimes d} )
= 0
$$
pour tout $n \geq 0$ et tout $d \geq d_0$.
\end{lemma}

\begin{proof}
Posons $B = \bigoplus I^n/I^{n + 1}$ et
$\mathcal{B} = \bigoplus \mathcal{I}^n/\mathcal{I}^{n + 1} = f^*\widetilde{B}$.
D'après le lemme \ref{coherent-lemma-finite-over-rees-algebra}, le
$\mathcal{B}$-module gradué quasi-cohérent
$\mathcal{G} = \bigoplus \Ker(\mathcal{F}_{n + 1} \to \mathcal{F}_n)$
est de type fini. Le lemme découle donc de la partie (2) du
lemme \ref{coherent-lemma-graded-finiteness}.
\end{proof}

\begin{lemma}
\label{coherent-lemma-existence-projective}
Soit $A$ un anneau noethérien complet pour la topologie $I$-adique.
Soit $f : X \to \Spec(A)$ un morphisme projectif. Posons
$\mathcal{I} = I\mathcal{O}_X$.
Alors le foncteur (\ref{coherent-equation-completion-functor}) est une équivalence.
\end{lemma}

\begin{proof}
Nous avons déjà vu que (\ref{coherent-equation-completion-functor}) est
pleinement fidèle (lemme \ref{coherent-lemma-fully-faithful}). Il suffit donc
de montrer que ce foncteur est essentiellement surjectif.

\medskip\noindent
Montrons d'abord que tout objet $(\mathcal{F}_n)$ de
$\textit{Coh}(X, \mathcal{I})$ est le quotient d'un objet
appartenant à l'image de (\ref{coherent-equation-completion-functor}).
Soit $\mathcal{L}$ un faisceau inversible $f$-ample sur $X$.
Choisissons $d_0$ comme dans le lemme \ref{coherent-lemma-vanishing-projective}.
Choisissons $d \geq d_0$ tel que
$\mathcal{F}_1 \otimes \mathcal{L}^{\otimes d}$
soit globalement engendré par des sections $s_{1, 1}, \ldots, s_{t, 1}$.
Puisque les applications de transition du système
$$
H^0(X, \mathcal{F}_{n + 1} \otimes \mathcal{L}^{\otimes d})
\longrightarrow
H^0(X, \mathcal{F}_n \otimes \mathcal{L}^{\otimes d})
$$
sont surjectives par l'annulation de $H^1$, nous pouvons relever
$s_{1, 1}, \ldots, s_{t, 1}$ en un système compatible de sections globales
$s_{1, n}, \ldots, s_{t, n}$ de
$\mathcal{F}_n \otimes \mathcal{L}^{\otimes d}$.
Celles-ci déterminent un système compatible d'applications
$$
(s_{1, n}, \ldots, s_{t, n}) :
(\mathcal{L}^{\otimes -d})^{\oplus t} \longrightarrow \mathcal{F}_n
$$
En utilisant le lemme \ref{coherent-lemma-inverse-systems-surjective},
nous en déduisons une application surjective
$$
\left((\mathcal{L}^{\otimes -d})^{\oplus t}\right)^\wedge
\longrightarrow
(\mathcal{F}_n)
$$
comme voulu.

\medskip\noindent
Le résultat du paragraphe précédent et le fait que
$\textit{Coh}(X, \mathcal{I})$ est abélienne
(lemme \ref{coherent-lemma-inverse-systems-abelian})
impliquent que
tout objet de $\textit{Coh}(X, \mathcal{I})$ est le conoyau
d'une application entre des objets provenant de $\textit{Coh}(\mathcal{O}_X)$.
Comme (\ref{coherent-equation-completion-functor}) est pleinement fidèle et exact par les
lemmes \ref{coherent-lemma-fully-faithful} et \ref{coherent-lemma-exact},
nous concluons.
\end{proof}








\section{Théorème d'existence de Grothendieck, II}
\label{coherent-section-existence-proper}

\noindent
Dans cette section, nous étudions le théorème d'existence de Grothendieck dans le
cas propre. Avant d'en donner l'énoncé et la démonstration, nous devons développer quelque
peu la théorie des catégories $\textit{Coh}(X, \mathcal{I})$
de modules formels cohérents introduites
à la section \ref{coherent-section-coherent-formal}.

\begin{remark}
\label{coherent-remark-inverse-systems-kernel-cokernel-annihilated-by}
Soit $X$ un schéma noethérien et soient
$\mathcal{I}, \mathcal{K} \subset \mathcal{O}_X$
des faisceaux quasi-cohérents d'idéaux. Soit
$\alpha : (\mathcal{F}_n) \to (\mathcal{G}_n)$ un morphisme de
$\textit{Coh}(X, \mathcal{I})$.
Étant donné un ouvert affine $\Spec(A) = U \subset X$ tel que
$\mathcal{I}|_U, \mathcal{K}|_U$ correspondent aux idéaux $I, K \subset A$,
notons $\alpha_U : M \to N$ le morphisme de $A^\wedge$-modules finis
qui correspond à $\alpha|_U$ via le lemme \ref{coherent-lemma-inverse-systems-affine}.
Nous affirmons que les conditions suivantes sont équivalentes :
\begin{enumerate}
\item il existe un entier $t \geq 1$ tel que
$\Ker(\alpha_n)$ et $\Coker(\alpha_n)$
soient annulés par $\mathcal{K}^t$ pour tout $n \geq 1$,
\item pour tout ouvert affine $\Spec(A) = U \subset X$ comme ci-dessus,
les modules $\Ker(\alpha_U)$ et $\Coker(\alpha_U)$
sont annulés par $K^t$ pour un certain entier $t \geq 1$, et
\item il existe un recouvrement affine fini $X = \bigcup U_i$
tel que la conclusion de (2) soit vérifiée pour $\alpha_{U_i}$.
\end{enumerate}
Si ces conditions équivalentes sont vérifiées, nous dirons que
$\alpha$ est une
{\it application dont le noyau et le conoyau sont annulés par une puissance de
$\mathcal{K}$}.
Pour voir l'équivalence, utilisons le fait d'algèbre commutative suivant :
supposons donnée une suite exacte
$$
0 \to T \to M \to N \to Q \to 0
$$
de $A$-modules, où $T$ et $Q$ sont annulés par $K^t$ pour un certain
idéal $K \subset A$. Alors, pour tous $f, g \in K^t$, il existe une
application canonique $"fg": N \to M$ telle que $M \to N \to M$ soit la
multiplication par $fg$. En effet, pour $y \in N$, nous pouvons choisir $x \in M$
ayant pour image $fy$ dans $N$, puis poser $"fg"(y) = gx$. Il est alors
clair que $\Ker(M/JM \to N/JN)$ et $\Coker(M/JM \to N/JN)$
sont annulés par $K^{2t}$ pour tout idéal $J \subset A$.

\medskip\noindent
En appliquant ce fait d'algèbre commutative à $\alpha_{U_i}$ et $J = I^n$,
nous voyons que (3) implique (1). Réciproquement,
supposons (1) vérifiée et $M \to N$ égal à $\alpha_U$. Il existe
alors $t \geq 1$ tel que
$\Ker(M/I^nM \to N/I^nN)$ et $\Coker(M/I^nM \to N/I^nN)$
soient annulés par $K^t$ pour tout $n$. Nous obtenons des applications
$"fg" : N/I^nN \to M/I^nM$ qui induisent, par passage à la limite, une application $N \to M$,
car $N$ et $M$ sont complets pour la topologie $I$-adique. Comme le composé
$N \to M \to N$ est la multiplication par $fg$, nous concluons que $fg$
annule $T$ et $Q$. Autrement dit, $T$ et $Q$ sont annulés par
$K^{2t}$, comme voulu.
\end{remark}

\begin{lemma}
\label{coherent-lemma-existence-tricky}
Soit $X$ un schéma noethérien. Soient
$\mathcal{I}, \mathcal{K} \subset \mathcal{O}_X$
des faisceaux quasi-cohérents d'idéaux.
Soit $X_e \subset X$ le sous-schéma fermé défini par $\mathcal{K}^e$.
Posons $\mathcal{I}_e = \mathcal{I}\mathcal{O}_{X_e}$.
Soit $(\mathcal{F}_n)$ un objet de $\textit{Coh}(X, \mathcal{I})$. Supposons
que
\begin{enumerate}
\item le foncteur
$\textit{Coh}(\mathcal{O}_{X_e}) \to \textit{Coh}(X_e, \mathcal{I}_e)$
est une équivalence pour tout $e \geq 1$, et
\item il existe un faisceau cohérent $\mathcal{H}$ sur $X$ et une application
$\alpha : (\mathcal{F}_n) \to \mathcal{H}^\wedge$ dont
le noyau et le conoyau sont annulés par une puissance de $\mathcal{K}$.
\end{enumerate}
Alors $(\mathcal{F}_n)$ appartient à l'image essentielle de
(\ref{coherent-equation-completion-functor}).
\end{lemma}

\begin{proof}
Dans toute cette démonstration, nous utiliserons sans autre mention le fait que, pour une immersion
fermée $i : Z \to X$, le foncteur $i_*$ donne une équivalence entre la
catégorie des modules cohérents sur $Z$ et celle des modules cohérents sur $X$ annulés
par le faisceau d'idéaux de $Z$, voir le lemme \ref{coherent-lemma-i-star-equivalence}.
En particulier, nous pouvons identifier $\textit{Coh}(\mathcal{O}_{X_e})$
à la catégorie des $\mathcal{O}_X$-modules cohérents annulés par
$\mathcal{K}^e$, et $\textit{Coh}(X_e, \mathcal{I}_e)$ à la sous-catégorie pleine
de $\textit{Coh}(X, \mathcal{I})$ formée des objets annulés par $\mathcal{K}^e$.
De plus, (1) nous dit que le foncteur de
complétion (\ref{coherent-equation-completion-functor}) induit une équivalence entre ces catégories.

\medskip\noindent
En appliquant cette équivalence, nous obtenons un $\mathcal{O}_X$-module cohérent
$\mathcal{G}_e$, annulé par $\mathcal{K}^e$, qui correspond au système
$(\mathcal{F}_n/\mathcal{K}^e\mathcal{F}_n)$ de
$\textit{Coh}(X, \mathcal{I})$. Les applications
$\mathcal{F}_n/\mathcal{K}^{e + 1}\mathcal{F}_n \to
\mathcal{F}_n/\mathcal{K}^e\mathcal{F}_n$ correspondent à des applications canoniques
$\mathcal{G}_{e + 1} \to \mathcal{G}_e$ qui induisent des isomorphismes
$\mathcal{G}_{e + 1}/\mathcal{K}^e\mathcal{G}_{e + 1} \to \mathcal{G}_e$.
Ainsi, $(\mathcal{G}_e)$ est un objet de $\textit{Coh}(X, \mathcal{K})$.
L'application $\alpha$ induit un système d'applications
$$
\mathcal{F}_n/\mathcal{K}^e\mathcal{F}_n
\longrightarrow
\mathcal{H}/(\mathcal{I}^n + \mathcal{K}^e)\mathcal{H}
$$
et donc des applications $\mathcal{G}_e \to \mathcal{H}/\mathcal{K}^e\mathcal{H}$
(en utilisant à nouveau l'équivalence de catégories).
Soit $t \geq 1$ un entier, qui existe par l'hypothèse (2),
tel que $\mathcal{K}^t$ annule le noyau et le conoyau de toutes les applications
$\mathcal{F}_n \to \mathcal{H}/\mathcal{I}^n\mathcal{H}$.
Alors $\mathcal{K}^{2t}$ annule le noyau et le conoyau des applications
$\mathcal{F}_n/\mathcal{K}^e\mathcal{F}_n \to
\mathcal{H}/(\mathcal{I}^n + \mathcal{K}^e)\mathcal{H}$, voir la
remarque \ref{coherent-remark-inverse-systems-kernel-cokernel-annihilated-by}.
Nous en concluons que $\mathcal{K}^{4t}$ annule le noyau et
le conoyau des applications
$$
\mathcal{G}_e
\longrightarrow
\mathcal{H}/\mathcal{K}^e\mathcal{H},
$$
voir la remarque \ref{coherent-remark-inverse-systems-kernel-cokernel-annihilated-by}.
Nous appliquons le lemme \ref{coherent-lemma-existence-easy} pour obtenir un
$\mathcal{O}_X$-module cohérent $\mathcal{F}$, une application
$a : \mathcal{F} \to \mathcal{H}$ et un isomorphisme
$\beta : (\mathcal{G}_e) \to (\mathcal{F}/\mathcal{K}^e\mathcal{F})$
dans $\textit{Coh}(X, \mathcal{K})$. En remontant les constructions, pour $n$
donné, le triplet
$(\mathcal{F}/\mathcal{I}^n\mathcal{F}, a \bmod \mathcal{I}^n, \beta
\bmod \mathcal{I}^n)$ est un triplet comme dans le lemme pour le morphisme
$\alpha_n \bmod \mathcal{K}^e :
(\mathcal{F}_n/\mathcal{K}^e\mathcal{F}_n) \to
(\mathcal{H}/(\mathcal{I}^n + \mathcal{K}^e)\mathcal{H})$
de $\textit{Coh}(X, \mathcal{K})$. L'unicité dans le
lemme \ref{coherent-lemma-existence-easy}
donne donc un isomorphisme canonique
$\mathcal{F}/\mathcal{I}^n\mathcal{F} \to \mathcal{F}_n$
compatible avec tous les morphismes considérés. Ceci achève la
démonstration du lemme.
\end{proof}

\begin{lemma}
\label{coherent-lemma-inverse-systems-push-pull}
Soit $Y$ un schéma noethérien. Soient
$\mathcal{J}, \mathcal{K} \subset \mathcal{O}_Y$
des faisceaux quasi-cohérents d'idéaux.
Soit $f : X \to Y$ un morphisme propre qui est un isomorphisme au-dessus
de $V = Y \setminus V(\mathcal{K})$.
Posons $\mathcal{I} = f^{-1}\mathcal{J} \mathcal{O}_X$.
Soit $(\mathcal{G}_n)$ un objet de $\textit{Coh}(Y, \mathcal{J})$,
soit $\mathcal{F}$ un $\mathcal{O}_X$-module cohérent, et soit
$\beta : (f^*\mathcal{G}_n)  \to \mathcal{F}^\wedge$ un isomorphisme dans
$\textit{Coh}(X, \mathcal{I})$. Il existe alors une application
$$
\alpha :
(\mathcal{G}_n)
\longrightarrow
(f_*\mathcal{F})^\wedge
$$
dans $\textit{Coh}(Y, \mathcal{J})$ dont le noyau et le
conoyau sont annulés par une puissance de $\mathcal{K}$.
\end{lemma}

\begin{proof}
Puisque $f$ est un morphisme propre, $f_*\mathcal{F}$
est un $\mathcal{O}_Y$-module cohérent
(proposition \ref{coherent-proposition-proper-pushforward-coherent}).
L'énoncé du lemme a donc bien un sens.
Considérons les composés
$$
\gamma_n : \mathcal{G}_n \to
f_*f^*\mathcal{G}_n \to
f_*(\mathcal{F}/\mathcal{I}^n\mathcal{F}).
$$
Ici, la première application est l'application d'adjonction, et la seconde est $f_*\beta_n$.
Nous affirmons qu'il existe une unique $\alpha$ comme dans le lemme
telle que les composés
$$
\mathcal{G}_n \xrightarrow{\alpha_n}
f_*\mathcal{F}/\mathcal{J}^nf_*\mathcal{F} \to
f_*(\mathcal{F}/\mathcal{I}^n\mathcal{F})
$$
soient égaux à $\gamma_n$ pour tout $n$. En raison de l'unicité, nous pouvons supposer
que $Y = \Spec(B)$ est affine. Soit $J \subset B$ l'idéal correspondant
à $\mathcal{J}$. Posons
$$
M_n = H^0(X, \mathcal{F}/\mathcal{I}^n\mathcal{F})
\quad\text{et}\quad
M = H^0(X, \mathcal{F})
$$
D'après le lemme \ref{coherent-lemma-ML-cohomology-powers-ideal} et le
Théorème \ref{coherent-theorem-formal-functions},
la limite inverse des modules
$M_n$ est égale à la complétion $M^\wedge = \lim M/J^nM$.
Posons $N_n = H^0(Y, \mathcal{G}_n)$ et $N = \lim N_n$.
Par l'équivalence de catégories du
lemme \ref{coherent-lemma-inverse-systems-affine},
les $B^\wedge$-modules finis $N$ et $M^\wedge$ correspondent
à $(\mathcal{G}_n)$ et $f_*\mathcal{F}^\wedge$.
Il s'ensuit que $\alpha$ doit être le morphisme de
$\textit{Coh}(Y, \mathcal{J})$ correspondant à l'homomorphisme
$$
\lim \gamma_n : N = \lim_n N_n \longrightarrow \lim M_n = M^\wedge
$$
de $B^\wedge$-modules finis.

\medskip\noindent
Il reste à montrer que le noyau et le conoyau de $\alpha$ sont
annulés par une puissance de $\mathcal{K}$. Posons $Y' = \Spec(B^\wedge)$
et $X' = Y' \times_Y X$. Soient $\mathcal{K}'$, $\mathcal{J}'$, $\mathcal{G}'_n$
et $\mathcal{I}'$, $\mathcal{F}'$ les images inverses de
$\mathcal{K}$, $\mathcal{J}$, $\mathcal{G}_n$ et
$\mathcal{I}$, $\mathcal{F}$, sur $Y'$ et $X'$.
Le morphisme de projection $f' : X' \to Y'$ est le changement de base de
$f$ par $Y' \to Y$. Notons que $Y' \to Y$ est un morphisme plat de schémas,
car $B \to B^\wedge$ est plat d'après
Algèbre, lemme \ref{algebra-lemma-completion-flat}.
Ainsi, $f'_*\mathcal{F}'$, resp.\ $f'_*(f')^*\mathcal{G}_n'$,
est l'image inverse de $f_*\mathcal{F}$, resp.\ $f_*f^*\mathcal{G}_n$,
sur $Y'$ par le lemme \ref{coherent-lemma-flat-base-change-cohomology}.
L'unicité de notre construction montre que l'image inverse de $\alpha$ sur $Y'$
est l'application correspondante $\alpha'$ construite dans la situation sur $Y'$.
De plus, pour vérifier que le noyau et le conoyau de $\alpha$ sont
annulés par $\mathcal{K}^t$, il suffit de vérifier que le
noyau et le conoyau de $\alpha'$ sont annulés par
$(\mathcal{K}')^t$. En effet, il suffit pour cela de le vérifier pour
les noyaux et conoyaux des applications $\alpha_n$ et $\alpha'_n$
(voir la remarque \ref{coherent-remark-inverse-systems-kernel-cokernel-annihilated-by}),
et l'homomorphisme d'anneaux $B \to B^\wedge$ induit
une équivalence de catégories entre les modules annulés par
$J^n$ et ceux annulés par $(J')^n$, voir
Compléments sur l'algèbre, lemme \ref{more-algebra-lemma-neighbourhood-equivalence}.
Nous pouvons donc supposer que $B$ est complet pour la topologie $J$-adique.

\medskip\noindent
Supposons que $Y = \Spec(B)$ soit affine, que $\mathcal{J}$ corresponde à l'idéal
$J \subset B$, et que $B$ soit complet pour la topologie $J$-adique.
Dans ce cas, $(\mathcal{G}_n)$ appartient à l'image essentielle du foncteur
$\textit{Coh}(\mathcal{O}_Y) \to \textit{Coh}(Y, \mathcal{J})$.
Soit $\mathcal{G}$ un $\mathcal{O}_Y$-module cohérent tel que
$(\mathcal{G}_n) = \mathcal{G}^\wedge$. Notons que
$f^*(\mathcal{G}^\wedge) = (f^*\mathcal{G})^\wedge$. Le
lemme \ref{coherent-lemma-fully-faithful}
nous dit donc que $\beta$ provient d'un isomorphisme
$b : f^*\mathcal{G} \to \mathcal{F}$,
et que $\alpha$ s'obtient en appliquant le foncteur de complétion à
$$
\mathcal{G} \to f_*f^*\mathcal{G} \cong f_*\mathcal{F}
$$
Il s'agit donc de vérifier que le noyau et le conoyau de
l'application d'adjonction $c : \mathcal{G} \to f_*f^*\mathcal{G}$ sont annulés par
une puissance de $\mathcal{K}$. Or, puisque la restriction
$f|_{f^{-1}(V)} : f^{-1}(V) \to V$ est un isomorphisme,
nous voyons que $c|_V$ est un isomorphisme. Ainsi, les faisceaux cohérents
$\Ker(c)$ et $\Coker(c)$ sont supportés sur $V(\mathcal{K})$,
et sont donc annulés par une puissance de $\mathcal{K}$
(lemme \ref{coherent-lemma-power-ideal-kills-sheaf}), comme voulu.
\end{proof}

\noindent
La proposition suivante est la forme du théorème d'existence de
Grothendieck qui est la plus souvent utilisée en pratique.

\begin{proposition}
\label{coherent-proposition-existence-proper}
Soit $A$ un anneau noethérien complet pour la topologie $I$-adique.
Soit $f : X \to \Spec(A)$ un morphisme propre de schémas.
Posons $\mathcal{I} = I\mathcal{O}_X$.
Alors le foncteur (\ref{coherent-equation-completion-functor}) est une équivalence.
\end{proposition}

\begin{proof}
Nous avons déjà vu que (\ref{coherent-equation-completion-functor}) est
pleinement fidèle (lemme \ref{coherent-lemma-fully-faithful}). Il suffit donc
de montrer que ce foncteur est essentiellement surjectif.

\medskip\noindent
Considérons l'ensemble $\Xi$ des faisceaux quasi-cohérents d'idéaux
$\mathcal{K} \subset \mathcal{O}_X$ tels que tout objet
$(\mathcal{F}_n)$ annulé par $\mathcal{K}$ appartienne à l'image essentielle.
Nous voulons montrer que $(0)$ appartient à $\Xi$. Sinon, puisque $X$ est noethérien,
il existe un faisceau quasi-cohérent d'idéaux maximal $\mathcal{K}$
n'appartenant pas à $\Xi$, voir le
lemme \ref{coherent-lemma-acc-coherent}.
Après avoir remplacé $X$ par le sous-schéma fermé de $X$
correspondant à $\mathcal{K}$, nous pouvons supposer que tout faisceau quasi-cohérent d'idéaux
$\mathcal{K}$ non nul appartient à $\Xi$. (On utilise ici la correspondance entre
les modules cohérents annulés par $\mathcal{K}$ et les modules
cohérents sur le sous-schéma fermé correspondant à $\mathcal{K}$, voir le
lemme \ref{coherent-lemma-i-star-equivalence}.)
Soit $(\mathcal{F}_n)$ un objet de
$\textit{Coh}(X, \mathcal{I})$.
Nous allons montrer que cet objet appartient à l'image essentielle du
foncteur (\ref{coherent-equation-completion-functor}), achevant ainsi la
démonstration de la proposition.

\medskip\noindent
Appliquons le lemme de Chow (lemme \ref{coherent-lemma-chow-Noetherian}) pour obtenir un
morphisme propre surjectif $f : X' \to X$ qui soit un isomorphisme
au-dessus d'un ouvert dense $U \subset X$ et tel que $X'$ soit projectif sur $A$.
Soit $\mathcal{K}$ le faisceau quasi-cohérent d'idéaux définissant
le complémentaire réduit $X \setminus U$. Par le cas
projectif du théorème d'existence de Grothendieck
(lemme \ref{coherent-lemma-existence-projective}),
il existe un module cohérent $\mathcal{F}'$ sur $X'$ tel
que $(\mathcal{F}')^\wedge \cong (f^*\mathcal{F}_n)$. D'après la
proposition \ref{coherent-proposition-proper-pushforward-coherent},
le $\mathcal{O}_X$-module $\mathcal{H} = f_*\mathcal{F}'$ est cohérent
et, d'après le lemme \ref{coherent-lemma-inverse-systems-push-pull},
il existe un morphisme $(\mathcal{F}_n) \to \mathcal{H}^\wedge$
de $\textit{Coh}(X, \mathcal{I})$ dont le noyau et le conoyau
sont annulés par une puissance de $\mathcal{K}$. Toutes les puissances $\mathcal{K}^e$
appartiennent à $\Xi$, si bien que (\ref{coherent-equation-completion-functor})
est une équivalence pour les sous-schémas fermés $X_e = V(\mathcal{K}^e)$.
Nous concluons par le lemme \ref{coherent-lemma-existence-tricky}.
\end{proof}





\section{Propreté sur une base}
\label{coherent-section-proper-over-base}

\noindent
Il s'agit simplement d'une courte section destinée à signaler quelques propriétés utiles des sous-ensembles
fermés propres sur une base et des modules quasi-cohérents de type fini
dont le support est propre sur une base.

\begin{lemma}
\label{coherent-lemma-closed-proper-over-base}
Soit $f : X \to S$ un morphisme de schémas localement de type fini.
Soit $Z \subset X$ un sous-ensemble fermé. Les conditions suivantes sont équivalentes :
\begin{enumerate}
\item le morphisme $Z \to S$ est propre lorsque $Z$ est muni de la structure
de sous-schéma fermé réduit induit
(Schémas, définition \ref{schemes-definition-reduced-induced-scheme}),
\item pour une certaine structure de sous-schéma fermé sur $Z$, le morphisme $Z \to S$
est propre,
\item pour toute structure de sous-schéma fermé sur $Z$, le morphisme
$Z \to S$ est propre.
\end{enumerate}
\end{lemma}

\begin{proof}
Les implications (3) $\Rightarrow$ (1) et (1) $\Rightarrow$ (2)
sont immédiates. Il suffit donc de prouver que (2) implique (3).
Nous invitons le lecteur à trouver sa propre démonstration de ce fait.
Soient $Z'$ et $Z''$ des structures de sous-schémas fermés sur $Z$
telles que $Z' \to S$ soit propre. Il faut montrer que $Z'' \to S$ est propre.
Soit $Z''' = Z' \cup Z''$ la réunion schématique, voir
Morphismes, Définition
\ref{morphisms-definition-scheme-theoretic-intersection-union}.
Alors $Z'''$ est une autre structure de sous-schéma fermé sur $Z$.
Cela résulte par exemple de la description des réunions schématiques dans
Morphismes, lemme \ref{morphisms-lemma-scheme-theoretic-union}.
Puisque $Z'' \to Z'''$ est une immersion fermée, il suffit de prouver
que $Z''' \to S$ est propre (voir
Morphismes, lemmes \ref{morphisms-lemma-closed-immersion-proper} et
\ref{morphisms-lemma-composition-proper}).
Le morphisme $Z' \to Z'''$ est une immersion fermée bijective,
donc en particulier surjective et universellement
fermée. Le fait que $Z' \to S$ soit séparé implique alors que
$Z''' \to S$ est séparé, voir
Morphismes, lemme \ref{morphisms-lemma-image-universally-closed-separated}. De
plus, $Z''' \to S$ est localement de type fini
puisque $X \to S$ est localement de type fini
(Morphismes, lemmes \ref{morphisms-lemma-immersion-locally-finite-type} et
\ref{morphisms-lemma-composition-finite-type}).
Puisque $Z' \to S$ est quasi-compact et $Z' \to Z'''$ est un homéomorphisme,
nous voyons que $Z''' \to S$ est quasi-compact.
Enfin, puisque $Z' \to S$ est universellement fermé, nous voyons qu'il
en va de même pour $Z''' \to S$ d'après
Morphismes, lemme \ref{morphisms-lemma-image-proper-is-proper}.
Cela achève la démonstration.
\end{proof}

\begin{definition}
\label{coherent-definition-proper-over-base}
Soit $f : X \to S$ un morphisme de schémas localement de type fini.
Soit $Z \subset X$ un sous-ensemble fermé. Nous
disons que {\it $Z$ est propre sur $S$}
si les conditions équivalentes du lemme \ref{coherent-lemma-closed-proper-over-base}
sont satisfaites.
\end{definition}

\noindent
Le lemme utilisé dans la définition ci-dessus est faux si le morphisme
$f : X \to S$ n'est pas localement de type fini. Nous invitons donc le
lecteur à ne pas employer cette terminologie si $f$ n'est pas localement de
type fini.

\begin{lemma}
\label{coherent-lemma-closed-closed-proper-over-base}
Soit $f : X \to S$ un morphisme de schémas localement de type fini.
Soient $Y \subset Z \subset X$ des sous-ensembles fermés.
Si $Z$ est propre sur $S$, alors il en va de même pour $Y$.
\end{lemma}

\begin{proof}
Démonstration omise.
\end{proof}

\begin{lemma}
\label{coherent-lemma-base-change-closed-proper-over-base}
Considérons un diagramme cartésien de schémas
$$
\xymatrix{
X' \ar[d]_{f'} \ar[r]_{g'} & X \ar[d]^f \\
S' \ar[r]^g & S
}
$$
où $f$ est localement de type fini.
Si $Z$ est un sous-ensemble fermé de $X$ propre sur $S$, alors
$(g')^{-1}(Z)$ est un sous-ensemble fermé de $X'$ propre sur $S'$.
\end{lemma}

\begin{proof}
Remarquons que l'énoncé a un sens puisque $f'$ est localement de
type fini d'après Morphismes, lemme \ref{morphisms-lemma-base-change-finite-type}.
Munissons $Z$ de la structure de sous-schéma fermé réduit induit.
Notons $Z' = (g')^{-1}(Z)$ l'image inverse schématique (Schémas,
définition \ref{schemes-definition-inverse-image-closed-subscheme}).
Alors $Z' = X' \times_X Z = (S' \times_S X) \times_X Z = S' \times_S Z$
est propre sur $S'$ en tant que changement de base de $Z$ sur $S$
(Morphismes, lemme \ref{morphisms-lemma-base-change-proper}).
\end{proof}

\begin{lemma}
\label{coherent-lemma-functoriality-closed-proper-over-base}
Soit $S$ un schéma. Soit $f : X \to Y$ un morphisme entre schémas
localement de type fini sur $S$.
\begin{enumerate}
\item Si $Y$ est séparé sur $S$ et si $Z \subset X$ est un sous-ensemble fermé
propre sur $S$, alors $f(Z)$ est un sous-ensemble fermé de $Y$ propre sur $S$.
\item Si $f$ est universellement fermé et si $Z \subset X$ est un
sous-ensemble fermé propre sur $S$, alors $f(Z)$ est un sous-ensemble fermé
de $Y$ propre sur $S$.
\item Si $f$ est propre et si $Z \subset Y$ est un sous-ensemble fermé
propre sur $S$, alors $f^{-1}(Z)$ est un sous-ensemble fermé de $X$ propre sur $S$.
\end{enumerate}
\end{lemma}

\begin{proof}
Démonstration de (1). Supposons $Y$ séparé sur $S$ et $Z \subset X$
fermé et propre sur $S$. Munissons $Z$ de sa structure réduite
induite de sous-schéma fermé et appliquons
Morphismes, lemme \ref{morphisms-lemma-scheme-theoretic-image-is-proper},
à $Z \to Y$ sur $S$ pour conclure.

\medskip\noindent
Démonstration de (2). Supposons $f$ universellement fermé et $Z \subset X$ fermé
et propre sur $S$. Munissons $Z$ et $Z' = f(Z)$ de leurs structures
réduites induites de sous-schémas fermés. On obtient un morphisme
induit $Z \to Z'$.
Notons $Z'' = f^{-1}(Z')$ l'image inverse schématique (Schémas,
définition \ref{schemes-definition-inverse-image-closed-subscheme}).
Alors $Z'' \to Z'$ est universellement fermé, comme changement de base de $f$
(Morphismes, lemme \ref{morphisms-lemma-base-change-proper}).
Ainsi, $Z \to Z'$ est universellement fermé, comme composé de
l'immersion fermée $Z \to Z''$ et de $Z'' \to Z'$
(Morphismes, Lemmes
\ref{morphisms-lemma-closed-immersion-proper} et
\ref{morphisms-lemma-composition-proper}).
On en déduit que $Z' \to S$ est séparé d'après
Morphismes, lemme \ref{morphisms-lemma-image-universally-closed-separated}.
Puisque $Z \to S$ est quasi-compact et que $Z \to Z'$ est surjectif,
le morphisme $Z' \to S$ est quasi-compact.
Puisque $Z' \to S$ est le composé de $Z' \to Y$ et de $Y \to S$,
le morphisme $Z' \to S$ est localement de type fini
(Morphismes, lemmes \ref{morphisms-lemma-immersion-locally-finite-type} et
\ref{morphisms-lemma-composition-finite-type}).
Enfin, puisque $Z \to S$ est universellement fermé, il
en va de même pour $Z' \to S$ d'après
Morphismes, lemme \ref{morphisms-lemma-image-proper-is-proper}.
Cela achève la démonstration.

\medskip\noindent
Démonstration de (3). Supposons $f$ propre et $Z \subset Y$ fermé et propre
sur $S$. Munissons $Z$ de sa structure réduite induite de sous-schéma
fermé. Notons $Z' = f^{-1}(Z)$ l'image inverse schématique (Schémas,
définition \ref{schemes-definition-inverse-image-closed-subscheme}).
Alors $Z' \to Z$ est propre, comme changement de base de $f$
(Morphismes, lemme \ref{morphisms-lemma-base-change-proper}).
Ainsi $Z' \to S$ est propre, comme composé de $Z' \to Z$
et de $Z \to S$
(Morphismes, lemme \ref{morphisms-lemma-composition-proper}).
Cela achève la démonstration.
\end{proof}

\begin{lemma}
\label{coherent-lemma-union-closed-proper-over-base}
Soit $f : X \to S$ un morphisme de schémas localement de type fini.
Soient $Z_i \subset X$, $i = 1, \ldots, n$, des sous-ensembles fermés.
Si les $Z_i$, $i = 1, \ldots, n$, sont propres sur $S$, il en va de
même pour $Z_1 \cup \ldots \cup Z_n$.
\end{lemma}

\begin{proof}
Munissons les $Z_i$ de leurs structures réduites induites de sous-schémas fermés.
Le morphisme
$$
Z_1 \amalg \ldots \amalg Z_n \longrightarrow X
$$
est fini d'après Morphismes, Lemmes
\ref{morphisms-lemma-closed-immersion-finite} et
\ref{morphisms-lemma-finite-union-finite}.
Les morphismes finis étant universellement fermés
(Morphismes, lemme \ref{morphisms-lemma-finite-proper}),
et $Z_1 \amalg \ldots \amalg Z_n$ étant propre sur $S$,
on déduit du
lemme \ref{coherent-lemma-functoriality-closed-proper-over-base} (2)
que l'image $Z_1 \cup \ldots \cup Z_n$ est propre sur $S$.
\end{proof}

\noindent
Soit $f : X \to S$ un morphisme de schémas localement de
type fini. Soit $\mathcal{F}$ un
$\mathcal{O}_X$-module quasi-cohérent de type fini. Alors le support $\text{Supp}(\mathcal{F})$
de $\mathcal{F}$ est un sous-ensemble fermé de $X$, voir
Morphismes, lemme \ref{morphisms-lemma-support-finite-type}.
On peut donc dire que «
le support de $\mathcal{F}$ est propre sur $S$ ».

\begin{lemma}
\label{coherent-lemma-module-support-proper-over-base}
Soit $f : X \to S$ un morphisme de schémas localement de
type fini. Soit $\mathcal{F}$ un
$\mathcal{O}_X$-module quasi-cohérent de type fini. Les conditions suivantes sont équivalentes :
\begin{enumerate}
\item le support de $\mathcal{F}$ est propre sur $S$,
\item le support schématique de $\mathcal{F}$
(Morphismes, définition \ref{morphisms-definition-scheme-theoretic-support})
est propre sur $S$, et
\item il existe un sous-schéma fermé $Z \subset X$ et
un $\mathcal{O}_Z$-module quasi-cohérent de type fini
$\mathcal{G}$ tels que (a) $Z \to S$ soit propre et (b)
$(Z \to X)_*\mathcal{G} = \mathcal{F}$.
\end{enumerate}
\end{lemma}

\begin{proof}
Le support $\text{Supp}(\mathcal{F})$ de $\mathcal{F}$ est un sous-ensemble fermé
de $X$, voir Morphismes, lemme \ref{morphisms-lemma-support-finite-type}.
On peut donc appliquer la définition \ref{coherent-definition-proper-over-base}.
Puisque le support schématique de $\mathcal{F}$ est un sous-schéma
fermé dont le sous-ensemble fermé sous-jacent est $\text{Supp}(\mathcal{F})$,
les conditions (1) et (2) sont équivalentes d'après la
définition \ref{coherent-definition-proper-over-base}.
Il est clair que (2) implique (3).
Réciproquement, si (3) est vérifiée, alors
$\text{Supp}(\mathcal{F}) \subset Z$
(inclusion de sous-ensembles fermés de $X$),
et donc $\text{Supp}(\mathcal{F})$
est propre sur $S$, par exemple d'après
le lemme \ref{coherent-lemma-closed-closed-proper-over-base}.
\end{proof}

\begin{lemma}
\label{coherent-lemma-base-change-module-support-proper-over-base}
Considérons un diagramme cartésien de schémas
$$
\xymatrix{
X' \ar[d]_{f'} \ar[r]_{g'} & X \ar[d]^f \\
S' \ar[r]^g & S
}
$$
où $f$ est localement de type fini. Soit $\mathcal{F}$
un $\mathcal{O}_X$-module quasi-cohérent
de type fini. Si le support de $\mathcal{F}$ est propre sur $S$, alors
le support de $(g')^*\mathcal{F}$ est propre sur $S'$.
\end{lemma}

\begin{proof}
Remarquons que l'énoncé a un sens, car
$(g')*\mathcal{F}$ est de type fini d'après
Modules, lemme \ref{modules-lemma-pullback-finite-type}.
On a $\text{Supp}((g')^*\mathcal{F}) = (g')^{-1}(\text{Supp}(\mathcal{F}))$
d'après Morphismes, lemme \ref{morphisms-lemma-support-finite-type}.
Le lemme résulte donc du
lemme \ref{coherent-lemma-base-change-closed-proper-over-base}.
\end{proof}

\begin{lemma}
\label{coherent-lemma-cat-module-support-proper-over-base}
Soit $f : X \to S$ un morphisme de schémas localement de
type fini. Soient $\mathcal{F}$, $\mathcal{G}$
des $\mathcal{O}_X$-modules quasi-cohérents de type fini.
\begin{enumerate}
\item Si les supports de $\mathcal{F}$, $\mathcal{G}$
sont propres sur $S$, il en va de même
pour $\mathcal{F} \oplus \mathcal{G}$, pour toute extension
de $\mathcal{G}$ par $\mathcal{F}$, pour $\Im(u)$ et $\Coker(u)$,
quel que soit le morphisme de $\mathcal{O}_X$-modules $u : \mathcal{F} \to \mathcal{G}$,
et pour tout quotient quasi-cohérent de $\mathcal{F}$ ou de $\mathcal{G}$.
\item Si $S$ est localement noethérien, la catégorie
des $\mathcal{O}_X$-modules cohérents à support propre sur
$S$ est une sous-catégorie de Serre (Homologie, Définition
\ref{homology-definition-serre-subcategory})
de la catégorie abélienne
des $\mathcal{O}_X$-modules cohérents.
\end{enumerate}
\end{lemma}

\begin{proof}
Démonstration de (1). Soient $Z$, $Z'$ les supports de $\mathcal{F}$
et de $\mathcal{G}$. Tous les faisceaux mentionnés dans (1) ont
leur support contenu dans $Z \cup Z'$. L'assertion résulte
donc des lemmes \ref{coherent-lemma-closed-closed-proper-over-base} et
\ref{coherent-lemma-union-closed-proper-over-base}
dès que l'on a vérifié que ces faisceaux sont de type
fini et quasi-cohérents. Pour la quasi-cohérence, nous renvoyons le lecteur à
Schémas, section \ref{schemes-section-quasi-coherent}.
Pour la propriété « de type fini », nous invitons le lecteur à consulter
Modules, section \ref{modules-section-finite-type}.

\medskip\noindent
Démonstration de (2). On raisonne comme pour (1). Notons que les assertions
ont un sens, puisque $X$ est localement noethérien d'après
Morphismes, lemme \ref{morphisms-lemma-finite-type-noetherian},
et compte tenu de la description de la catégorie des modules
cohérents dans la section \ref{coherent-section-coherent-sheaves}.
\end{proof}

\begin{lemma}
\label{coherent-lemma-support-proper-over-base-pushforward}
Soit $S$ un schéma localement noethérien.
Soit $f : X \to S$ un morphisme de schémas localement de type fini.
Soit $\mathcal{F}$ un $\mathcal{O}_X$-module
cohérent à support propre sur $S$. Alors $R^pf_*\mathcal{F}$
est un $\mathcal{O}_S$-module cohérent pour tout $p \geq 0$.
\end{lemma}

\begin{proof}
D'après le lemme \ref{coherent-lemma-module-support-proper-over-base},
il existe une immersion fermée $i : Z \to X$ et
un $\mathcal{O}_Z$-module quasi-cohérent de type fini
$\mathcal{G}$ tels que (a) $g = f \circ i : Z \to S$ soit propre et (b)
$i_*\mathcal{G} = \mathcal{F}$.
Le faisceau $R^pg_*\mathcal{G}$ est cohérent sur $S$ d'après la
proposition \ref{coherent-proposition-proper-pushforward-coherent}.
D'autre part, $R^qi_*\mathcal{G} = 0$ pour $q > 0$
(lemme \ref{coherent-lemma-finite-pushforward-coherent}).
D'après Cohomologie, lemme \ref{cohomology-lemma-relative-Leray},
on obtient $R^pf_*\mathcal{F} = R^pg_*\mathcal{G}$, ce qui achève la démonstration.
\end{proof}

\begin{lemma}
\label{coherent-lemma-systems-with-proper-support}
Soit $S$ un schéma noethérien. Soit $f : X \to S$ un morphisme de type fini.
Soit $\mathcal{I} \subset \mathcal{O}_X$ un
faisceau quasi-cohérent d'idéaux. Les catégories suivantes sont des sous-catégories de Serre
de $\textit{Coh}(X, \mathcal{I})$ :
\begin{enumerate}
\item la sous-catégorie pleine de $\textit{Coh}(X, \mathcal{I})$
formée des objets $(\mathcal{F}_n)$ tels que
le support de $\mathcal{F}_1$ soit propre sur $S$,
\item la sous-catégorie pleine de $\textit{Coh}(X, \mathcal{I})$
formée des objets $(\mathcal{F}_n)$ tels qu'il
existe un sous-schéma fermé $Z \subset X$ propre sur $S$
avec $\mathcal{I}_Z \mathcal{F}_n = 0$ pour tout $n \geq 1$.
\end{enumerate}
\end{lemma}

\begin{proof}
Utilisons le critère de
Homologie, lemme \ref{homology-lemma-characterize-serre-subcategory}.
Nous utiliserons aussi le fait que, si
$0 \to (\mathcal{G}_n) \to (\mathcal{F}_n) \to (\mathcal{H}_n) \to 0$
est une suite exacte courte de $\textit{Coh}(X, \mathcal{I})$, alors
(a) $\mathcal{G}_n \to \mathcal{F}_n \to \mathcal{H}_n \to 0$
est exacte pour tout $n \geq 1$ et
(b) $\mathcal{G}_n$ est un quotient de $\Ker(\mathcal{F}_m \to \mathcal{H}_m)$
pour un certain $m \geq n$. Voir la démonstration du lemme \ref{coherent-lemma-inverse-systems-abelian}.

\medskip\noindent
Démonstration de (1). Soit $(\mathcal{F}_n)$ un objet de
$\textit{Coh}(X, \mathcal{I})$. Alors
$\text{Supp}(\mathcal{F}_n) = \text{Supp}(\mathcal{F}_1)$ pour tout $n \geq 1$.
Les remarques (a) et (b) ci-dessus montrent donc que,
pour toute suite exacte courte
$0 \to (\mathcal{G}_n) \to (\mathcal{F}_n) \to (\mathcal{H}_n) \to 0$
de $\textit{Coh}(X, \mathcal{I})$, on a
$\text{Supp}(\mathcal{G}_1) \cup \text{Supp}(\mathcal{H}_1) =
\text{Supp}(\mathcal{F}_1)$.
Cela prouve que la catégorie définie dans (1) est
une sous-catégorie de Serre de $\textit{Coh}(X, \mathcal{I})$.

\medskip\noindent
Démonstration de (2). On raisonne de même. Soit
$0 \to (\mathcal{G}_n) \to (\mathcal{F}_n) \to (\mathcal{H}_n) \to 0$
une suite exacte courte de $\textit{Coh}(X, \mathcal{I})$.
Si $Z \subset X$ est un sous-schéma fermé et si $\mathcal{I}_Z$
annule $\mathcal{F}_n$ pour tout $n$, alors
$\mathcal{I}_Z$ annule $\mathcal{G}_n$ et $\mathcal{H}_n$
pour tout $n$ d'après (a) et (b) ci-dessus.
Si donc $Z \to S$ est propre, on en déduit que la catégorie
définie dans (2) est stable par passage aux sous-objets et aux objets
quotients dans $\textit{Coh}(X, \mathcal{I})$.
Enfin, supposons que $Z \subset X$ et $Y \subset X$ soient
des sous-schémas fermés propres sur $S$ tels que
$\mathcal{I}_Z \mathcal{G}_n = 0$ et
$\mathcal{I}_Y \mathcal{H}_n = 0$ pour tout $n \geq 1$.
Il résulte alors de (a) ci-dessus que
$\mathcal{I}_{Z \cup Y} = \mathcal{I}_Z \cdot \mathcal{I}_Y$
annule $\mathcal{F}_n$ pour tout $n$.
D'après le lemme \ref{coherent-lemma-union-closed-proper-over-base}
(et la définition \ref{coherent-definition-proper-over-base},
qui permet de choisir une structure de schéma quelconque sur la
réunion), le morphisme $Z \cup Y \to S$ est propre, ce qui achève la démonstration.
\end{proof}








\section{Théorème d'existence de Grothendieck, III}
\label{coherent-section-existence-proper-support}

\noindent
Pour énoncer la version générale du théorème d'existence
de Grothendieck, introduisons quelques notations supplémentaires. Soit $A$ un anneau noethérien complet
pour la topologie définie par un idéal $I$. Soit $f : X \to \Spec(A)$
un morphisme de schémas séparé et de type fini. Posons
$\mathcal{I} = I\mathcal{O}_X$. Dans cette situation, notons
$$
\textit{Coh}_{\text{support propre sur } A}(\mathcal{O}_X)
$$
la sous-catégorie pleine de $\textit{Coh}(\mathcal{O}_X)$
formée des $\mathcal{O}_X$-modules cohérents dont
le support est propre sur $\Spec(A)$. C'est une sous-catégorie de Serre de
$\textit{Coh}(\mathcal{O}_X)$, voir
lemme \ref{coherent-lemma-cat-module-support-proper-over-base}.
De même, notons
$$
\textit{Coh}_{\text{support propre sur } A}(X, \mathcal{I})
$$
la sous-catégorie pleine de $\textit{Coh}(X, \mathcal{I})$
formée des objets $(\mathcal{F}_n)$ tels que
le support de $\mathcal{F}_1$ soit propre sur $\Spec(A)$.
C'est une sous-catégorie de Serre de $\textit{Coh}(X, \mathcal{I})$
d'après le lemme \ref{coherent-lemma-systems-with-proper-support} (1). Puisque
le support d'un module quotient est contenu dans celui
du module, le foncteur (\ref{coherent-equation-completion-functor})
induit un foncteur
\begin{equation}
\label{coherent-equation-completion-functor-proper-over-A}
\textit{Coh}_{\text{support propre sur }A}(\mathcal{O}_X)
\longrightarrow
\textit{Coh}_{\text{support propre sur }A}(X, \mathcal{I})
\end{equation}
Nous pouvons maintenant énoncer le théorème principal de cette section.

\begin{theorem}[Théorème d'existence de Grothendieck]
\label{coherent-theorem-grothendieck-existence}
\begin{reference}
\cite[III théorème 5.1.5]{EGA}
\end{reference}
Soit $A$ un anneau noethérien complet pour la topologie définie par un idéal $I$.
Soit $X$ un schéma séparé et de type fini sur $A$. Alors
le foncteur
(\ref{coherent-equation-completion-functor-proper-over-A})
$$
\textit{Coh}_{\text{support propre sur }A}(\mathcal{O}_X)
\longrightarrow
\textit{Coh}_{\text{support propre sur }A}(X, \mathcal{I})
$$
est une équivalence.
\end{theorem}

\begin{proof}
Nous utiliserons l'équivalence de catégories du
lemme \ref{coherent-lemma-i-star-equivalence}
sans autre mention.
Pour un sous-schéma fermé $Z \subset X$ propre sur $A$,
nous dirons, dans cette démonstration, qu'un module cohérent sur $X$ est
« porté par $Z$ » s'il est annulé par le faisceau
d'idéaux de $Z$, ou, de manière équivalente, s'il est l'image
directe d'un module cohérent sur $Z$.
D'après la proposition \ref{coherent-proposition-existence-proper}, le
résultat est vrai pour le
foncteur entre modules cohérents et systèmes de modules
cohérents portés par $Z$. Il suffit donc de montrer que
tout objet de
$\textit{Coh}_{\text{support propre sur }A}(\mathcal{O}_X)$
et tout objet de
$\textit{Coh}_{\text{support propre sur }A}(X, \mathcal{I})$ est
porté par un sous-schéma fermé $Z \subset X$ propre sur $A$.
C'est vrai par définition pour les objets de
$\textit{Coh}_{\text{support propre sur }A}(\mathcal{O}_X)$.
Nous allons le démontrer pour les objets de
$\textit{Coh}_{\text{support propre sur }A}(X, \mathcal{I})$
en utilisant la méthode de démonstration de la proposition \ref{coherent-proposition-existence-proper}.
Nous invitons le lecteur à lire d'abord cette démonstration.

\medskip\noindent
Considérons l'ensemble $\Xi$ des faisceaux quasi-cohérents d'idéaux
$\mathcal{K} \subset \mathcal{O}_X$ tels que l'assertion soit vraie
pour tout objet $(\mathcal{F}_n)$ de
$\textit{Coh}_{\text{support propre sur }A}(X, \mathcal{I})$
annulé par $\mathcal{K}$. Montrons que $(0)$ appartient à $\Xi$.
Sinon, puisque $X$ est noethérien, il existe un
faisceau quasi-cohérent d'idéaux $\mathcal{K}$ maximal parmi ceux qui n'appartiennent pas à $\Xi$, voir
lemme \ref{coherent-lemma-acc-coherent}.
En remplaçant $X$ par le sous-schéma fermé de $X$
correspondant à $\mathcal{K}$, on peut supposer que tout faisceau non nul
$\mathcal{K}$ appartient à $\Xi$. Soit $(\mathcal{F}_n)$ un objet de
$\textit{Coh}_{\text{support propre sur }A}(X, \mathcal{I})$.
Nous allons montrer que cet objet est porté par un sous-schéma fermé
$Z \subset X$ propre sur $A$, ce qui achèvera
la démonstration du théorème.

\medskip\noindent
Appliquons le lemme de Chow (lemme \ref{coherent-lemma-chow-Noetherian}) pour obtenir un
morphisme propre surjectif $f : Y \to X$ qui est un isomorphisme
au-dessus d'un ouvert dense $U \subset X$ et tel que $Y$ soit H-quasi-projectif
sur $A$. Choisissons une immersion ouverte $j : Y \to Y'$, où
$Y'$ est projectif sur $A$, voir
Morphismes, lemme \ref{morphisms-lemma-H-quasi-projective-open-H-projective}.
Remarquons que
$$
\text{Supp}(f^*\mathcal{F}_n) = f^{-1}\text{Supp}(\mathcal{F}_n) =
f^{-1}\text{Supp}(\mathcal{F}_1)
$$
La première égalité résulte de
Morphismes, lemme \ref{morphisms-lemma-support-finite-type}.
Par hypothèse et d'après le
lemme \ref{coherent-lemma-functoriality-closed-proper-over-base} (3),
le sous-ensemble $f^{-1}\text{Supp}(\mathcal{F}_1)$ est propre sur $A$.
L'image de $f^{-1}\text{Supp}(\mathcal{F}_1)$
par $j$ est donc fermée dans $Y'$ d'après le
lemme \ref{coherent-lemma-functoriality-closed-proper-over-base} (1).
Ainsi $\mathcal{F}'_n = j_*f^*\mathcal{F}_n$ est cohérent sur
$Y'$ d'après le lemme \ref{coherent-lemma-pushforward-coherent-on-open}.
Il en résulte que $(\mathcal{F}_n')$
est un objet de $\textit{Coh}(Y', I\mathcal{O}_{Y'})$.
D'après le cas projectif du théorème d'existence de Grothendieck
(lemme \ref{coherent-lemma-existence-projective}),
il existe un $\mathcal{O}_{Y'}$-module cohérent
$\mathcal{F}'$ et un isomorphisme
$(\mathcal{F}')^\wedge \cong (\mathcal{F}'_n)$ dans
$\textit{Coh}(Y', I\mathcal{O}_{Y'})$.
Puisque $\mathcal{F}'/I\mathcal{F}' = \mathcal{F}'_1$, on a
$$
\text{Supp}(\mathcal{F}') \cap V(I\mathcal{O}_{Y'}) =
\text{Supp}(\mathcal{F}'_1) = j(f^{-1}\text{Supp}(\mathcal{F}_1))
$$
Le morphisme structural $p' : Y' \to \Spec(A)$ est propre, donc
$p'(\text{Supp}(\mathcal{F}') \setminus j(Y))$
est fermé dans $\Spec(A)$. Un sous-ensemble fermé non vide de $\Spec(A)$
contient un point de $V(I)$, car $I$ est contenu dans le radical de Jacobson
de $A$ d'après Algèbre, lemme \ref{algebra-lemma-radical-completion}.
L'égalité ci-dessus montre que
$\text{Supp}(\mathcal{F}') \cap (p')^{-1}V(I) \subset j(Y)$
et l'on en déduit que $\text{Supp}(\mathcal{F}') \subset j(Y)$.
Ainsi $\mathcal{F}'|_Y = j^*\mathcal{F}'$
est porté par un sous-schéma fermé $Z'$ de $Y$ propre sur $A$,
et $(\mathcal{F}'|_Y)^\wedge = (f^*\mathcal{F}_n)$.

\medskip\noindent
Soit $\mathcal{K}$ le faisceau quasi-cohérent d'idéaux définissant
le complémentaire réduit $X \setminus U$. D'après la
proposition \ref{coherent-proposition-proper-pushforward-coherent},
le $\mathcal{O}_X$-module $\mathcal{H} = f_*(\mathcal{F}'|_Y)$ est cohérent,
et d'après le lemme \ref{coherent-lemma-inverse-systems-push-pull},
il existe un morphisme $\alpha : (\mathcal{F}_n) \to \mathcal{H}^\wedge$
de $\textit{Coh}(X, \mathcal{I})$ dont le noyau et le conoyau
sont annulés par une puissance $\mathcal{K}^t$ de $\mathcal{K}$.
On obtient une suite exacte
$$
0 \to \Ker(\alpha) \to (\mathcal{F}_n) \to
\mathcal{H}^\wedge \to \Coker(\alpha) \to 0
$$
dans $\textit{Coh}(X, \mathcal{I})$. Si $Z_0 \subset X$ est le support
schématique de $\mathcal{H}$, il est clair que $Z_0 \subset f(Z')$
ensemblistement. Ainsi $Z_0$ est propre sur $A$ d'après le
lemme \ref{coherent-lemma-closed-closed-proper-over-base} et le
lemme \ref{coherent-lemma-functoriality-closed-proper-over-base} (2). Par
conséquent, $\mathcal{H}^\wedge$ appartient à la sous-catégorie définie dans le
lemme \ref{coherent-lemma-systems-with-proper-support} (2), et
a fortiori à
$\textit{Coh}_{\text{support propre sur }A}(X, \mathcal{I})$.
On en déduit que $\Ker(\alpha)$ et $\Coker(\alpha)$
appartiennent à $\textit{Coh}_{\text{support propre sur }A}(X, \mathcal{I})$
d'après le lemme \ref{coherent-lemma-systems-with-proper-support} (1). Par
l'hypothèse de récurrence, plus précisément parce que $\mathcal{K}^t$ appartient à $\Xi$,
les objets $\Ker(\alpha)$ et $\Coker(\alpha)$ appartiennent à
la sous-catégorie définie dans le
lemme \ref{coherent-lemma-systems-with-proper-support} (2). Celle-ci
étant une sous-catégorie de Serre d'après le lemme, il en va
de même pour $(\mathcal{F}_n)$, comme voulu.
\end{proof}

\begin{remark}[Explicitation du théorème d'existence de Grothendieck]
\label{coherent-remark-reformulate-existence-theorem}
Soit $A$ un anneau noethérien complet pour la topologie définie par un idéal $I$.
Écrivons $S = \Spec(A)$ et $S_n = \Spec(A/I^n)$.
Soit $X \to S$ un morphisme séparé de type fini.
Pour $n \geq 1$, posons $X_n = X \times_S S_n$. On a le
diagramme suivant :
$$
\xymatrix{
X_1 \ar[r]_{i_1} \ar[d] & X_2 \ar[r]_{i_2} \ar[d] & X_3 \ar[r] \ar[d] &
\ldots & X \ar[d] \\
S_1 \ar[r] & S_2 \ar[r] & S_3 \ar[r] & \ldots & S
}
$$
Dans cette situation, considérons les systèmes $(\mathcal{F}_n, \varphi_n)$
tels que
\begin{enumerate}
\item $\mathcal{F}_n$ soit un $\mathcal{O}_{X_n}$-module cohérent,
\item $\varphi_n : i_n^*\mathcal{F}_{n + 1} \to \mathcal{F}_n$
soit un isomorphisme, et
\item $\text{Supp}(\mathcal{F}_1)$ soit propre sur $S_1$.
\end{enumerate}
Le Théorème \ref{coherent-theorem-grothendieck-existence} affirme que le foncteur
de complétion
$$
\begin{matrix}
\text{les }\mathcal{O}_X\text{-modules cohérents }\mathcal{F} \\
\text{à support propre sur }A
\end{matrix}
\quad
\longrightarrow
\quad
\begin{matrix}
\text{systèmes }(\mathcal{F}_n) \\
\text{comme ci-dessus}
\end{matrix}
$$
est une équivalence de catégories. Dans le cas particulier où $X$ est
propre sur $A$, on peut omettre les conditions portant sur les supports.
\end{remark}



\section{Théorème d'algébrisation de Grothendieck}
\label{coherent-section-algebraization}

\noindent
Notre premier résultat est une reformulation du théorème d'existence de
Grothendieck en termes de sous-schémas fermés et de morphismes finis.

\begin{lemma}
\label{coherent-lemma-algebraize-formal-closed-subscheme}
Soit $A$ un anneau noethérien complet pour la topologie $I$-adique.
Posons $S = \Spec(A)$ et $S_n = \Spec(A/I^n)$.
Soit $X \to S$ un morphisme séparé de type fini.
Pour $n \geq 1$, posons $X_n = X \times_S S_n$.
Supposons donné un diagramme commutatif
$$
\xymatrix{
Z_1 \ar[r] \ar[d] & Z_2 \ar[r] \ar[d] & Z_3 \ar[r] \ar[d] & \ldots \\
X_1 \ar[r]^{i_1} & X_2 \ar[r]^{i_2} & X_3 \ar[r] & \ldots
}
$$
de schémas dont les carrés sont cartésiens. Supposons que
\begin{enumerate}
\item $Z_1 \to X_1$ est une immersion fermée, et
\item $Z_1 \to S_1$ est propre.
\end{enumerate}
Il existe alors une immersion fermée de schémas $Z \to X$ telle que
$Z_n = Z \times_S S_n$. De plus, $Z$ est propre sur $S$.
\end{lemma}

\begin{proof}
Notons $j_n : Z_n \to X_n$ les morphismes verticaux.
Les carrés de l'énoncé étant cartésiens,
le changement de base de $j_n$ à $X_1$ est $j_1$.
Ainsi, Morphismes, lemme \ref{morphisms-lemma-check-closed-infinitesimally},
montre que $j_n$ est une immersion fermée.
Posons $\mathcal{F}_n = j_{n, *}\mathcal{O}_{Z_n}$, de sorte que
$j_n^\sharp$ soit une surjection $\mathcal{O}_{X_n} \to \mathcal{F}_n$.
En utilisant à nouveau le caractère cartésien des carrés, on voit que
l'image inverse de $\mathcal{F}_{n + 1}$ sur $X_n$ est $\mathcal{F}_n$.
Le théorème d'existence de Grothendieck, sous la forme de la
remarque \ref{coherent-remark-reformulate-existence-theorem},
assure donc l'existence d'un morphisme
$\mathcal{O}_X \to \mathcal{F}$
de $\mathcal{O}_X$-modules cohérents dont la restriction à
$X_n$ redonne $\mathcal{O}_{X_n} \to \mathcal{F}_n$.
De plus, le support de $\mathcal{F}$ est propre sur $S$.
Comme le foncteur de complétion est exact (lemme \ref{coherent-lemma-exact}),
le conoyau $\mathcal{Q}$ de $\mathcal{O}_X \to \mathcal{F}$
a une complétion nulle. Comme le support de $\mathcal{F}$ est propre
sur $S$, et qu'il en est de même de $\mathcal{Q}$, il en résulte que
$\mathcal{Q} = 0$ ; cela découle par exemple du fait que le
foncteur (\ref{coherent-equation-completion-functor-proper-over-A}) est une équivalence par
le théorème d'existence de Grothendieck.
Ainsi $\mathcal{F} = \mathcal{O}_X/\mathcal{J}$
pour un certain faisceau quasi-cohérent d'idéaux $\mathcal{J}$. Il
suffit de poser $Z = V(\mathcal{J})$ pour achever la démonstration.
\end{proof}

\noindent
Dans le lemme suivant, il suffit en fait de supposer que $Y_1 \to X_1$
est fini, car cela entraîne que $Y_n \to X_n$ est également fini
(voir Compléments sur les morphismes, Lemme
\ref{more-morphisms-lemma-thicken-property-morphisms-cartesian}).

\begin{lemma}
\label{coherent-lemma-algebraize-formal-scheme-finite-over-proper}
Soit $A$ un anneau noethérien complet pour la topologie $I$-adique.
Posons $S = \Spec(A)$ et $S_n = \Spec(A/I^n)$.
Soit $X \to S$ un morphisme séparé de type fini.
Pour $n \geq 1$, posons $X_n = X \times_S S_n$.
Supposons donné un diagramme commutatif
$$
\xymatrix{
Y_1 \ar[r] \ar[d] & Y_2 \ar[r] \ar[d] & Y_3 \ar[r] \ar[d] & \ldots \\
X_1 \ar[r]^{i_1} & X_2 \ar[r]^{i_2} & X_3 \ar[r] & \ldots
}
$$
de schémas dont les carrés sont cartésiens. Supposons que
\begin{enumerate}
\item $Y_n \to X_n$ est un morphisme fini, et
\item $Y_1 \to S_1$ est propre.
\end{enumerate}
Il existe alors un morphisme fini de schémas $Y \to X$ tel que
$Y_n = Y \times_S S_n$. De plus, $Y$ est propre sur $S$.
\end{lemma}

\begin{proof}
Notons $f_n : Y_n \to X_n$ les morphismes verticaux.
Posons $\mathcal{F}_n = f_{n, *}\mathcal{O}_{Y_n}$. C'est
un $\mathcal{O}_{X_n}$-module cohérent puisque $f_n$ est fini
(lemme \ref{coherent-lemma-finite-pushforward-coherent}).
Le caractère cartésien des carrés montre que
l'image inverse de $\mathcal{F}_{n + 1}$ sur $X_n$ est $\mathcal{F}_n$.
Le théorème d'existence de Grothendieck, sous la forme de la
remarque \ref{coherent-remark-reformulate-existence-theorem},
assure donc l'existence d'un $\mathcal{O}_X$-module cohérent
$\mathcal{F}$ dont la restriction à $X_n$ redonne $\mathcal{F}_n$.
De plus, le support de $\mathcal{F}$ est propre sur $S$.
Comme le foncteur de complétion est pleinement fidèle
(Théorème \ref{coherent-theorem-grothendieck-existence}),
les applications de multiplication
$\mathcal{F}_n \otimes_{\mathcal{O}_{X_n}} \mathcal{F}_n \to
\mathcal{F}_n$ se recollent en une structure d'algèbre sur $\mathcal{F}$. Il
suffit de poser $Y = \underline{\Spec}_X(\mathcal{F})$ pour achever la démonstration.
\end{proof}

\begin{lemma}
\label{coherent-lemma-algebraize-morphism}
Soit $A$ un anneau noethérien complet pour la topologie $I$-adique.
Posons $S = \Spec(A)$ et $S_n = \Spec(A/I^n)$. Soient $X$ et $Y$ des schémas
sur $S$. Pour $n \geq 1$, posons $X_n = X \times_S S_n$ et
$Y_n = Y \times_S S_n$. Supposons donné un système compatible de
diagrammes commutatifs
$$
\xymatrix{
& & X_{n + 1} \ar[rd] \ar[rr]_{g_{n + 1}} & & Y_{n + 1} \ar[ld] \\
X_n \ar[rru] \ar[rd] \ar[rr]_{g_n} & & Y_n \ar[rru] \ar[ld] & S_{n + 1} \\
& S_n \ar[rru]
}
$$
Supposons que
\begin{enumerate}
\item $X \to S$ est propre, et
\item $Y \to S$ est séparé de type fini.
\end{enumerate}
Il existe alors un unique morphisme de schémas $g : X \to Y$
au-dessus de $S$ tel que $g_n$ soit le changement de base de $g$ à $S_n$.
\end{lemma}

\begin{proof}
Les morphismes $(1, g_n) : X_n \to X_n \times_S Y_n$ sont des immersions fermées,
car $Y_n \to S_n$ est séparé
(Schémas, lemme \ref{schemes-lemma-section-immersion}). Le
lemme \ref{coherent-lemma-algebraize-formal-closed-subscheme}
fournit donc un sous-schéma fermé $Z \subset X \times_S Y$,
propre sur $S$, dont le changement de base à $S_n$ redonne
$X_n \subset X_n \times_S Y_n$. La première projection $p : Z \to X$
est un morphisme propre (car $Z$ est propre sur $S$, voir
Morphismes, lemme \ref{morphisms-lemma-image-proper-scheme-closed})
dont le changement de base à $S_n$ est un isomorphisme
pour tout $n$. En particulier, $p : Z \to X$ est fini au-dessus d'un voisinage
ouvert de $X_0$, d'après
le lemme \ref{coherent-lemma-proper-finite-fibre-finite-in-neighbourhood}.
Comme $X$ est propre sur $S$, ce voisinage ouvert est $X$
tout entier, et nous concluons que $p : Z \to X$ est fini.
En appliquant l'équivalence de la proposition \ref{coherent-proposition-existence-proper},
on obtient $p_*\mathcal{O}_Z = \mathcal{O}_X$, puisque cette égalité est vraie
modulo $I^n$ pour tout $n$. Par conséquent, $p$ est un isomorphisme, et
le morphisme $g$ est le composé $X \cong Z \to Y$.
Nous omettons la démonstration de l'unicité.
\end{proof}

\noindent
Pour démontrer un théorème d'algébrisation « abstrait », nous
devons supposer l'existence d'un faisceau inversible ample : sans cette hypothèse,
le résultat est faux.

\begin{theorem}[Théorème d'algébrisation de Grothendieck]
\label{coherent-theorem-algebraization}
Soit $A$ un anneau noethérien complet pour la topologie $I$-adique.
Posons $S = \Spec(A)$ et $S_n = \Spec(A/I^n)$. Considérons un diagramme
commutatif
$$
\xymatrix{
X_1 \ar[r]_{i_1} \ar[d] & X_2 \ar[r]_{i_2} \ar[d] & X_3 \ar[r] \ar[d] &
\ldots \\
S_1 \ar[r] & S_2 \ar[r] & S_3 \ar[r] & \ldots
}
$$
de schémas dont les carrés sont cartésiens. Supposons donnés des couples $(\mathcal{L}_n, \varphi_n)$,
où chaque $\mathcal{L}_n$ est un faisceau inversible sur $X_n$ et
$\varphi_n : i_n^*\mathcal{L}_{n + 1} \to \mathcal{L}_n$
est un isomorphisme. Si
\begin{enumerate}
\item $X_1 \to S_1$ est propre, et
\item $\mathcal{L}_1$ est ample sur $X_1$,
\end{enumerate}
alors il existe un morphisme propre de schémas $X \to S$,
un $\mathcal{O}_X$-module inversible ample $\mathcal{L}$,
ainsi que des isomorphismes $X_n \cong X \times_S S_n$ et
$\mathcal{L}_n \cong \mathcal{L}|_{X_n}$ compatibles avec
les morphismes $i_n$ et $\varphi_n$.
\end{theorem}

\begin{proof}
Les carrés du diagramme étant cartésiens et les morphismes
$S_n \to S_{n + 1}$ étant des immersions fermées, les morphismes
$i_n$ sont eux aussi des immersions fermées. Nous pouvons donc considérer
$X_m$ comme un sous-schéma fermé de $X_n$ pour $m < n$. Plus précisément, $X_m$ est
le sous-schéma fermé défini par le faisceau quasi-cohérent d'idéaux
$I^m\mathcal{O}_{X_n}$. De plus, les espaces topologiques
sous-jacents aux schémas $X_1, X_2, X_3, \ldots$ s'identifient tous ; nous pouvons
donc considérer, et considérerons, les faisceaux $\mathcal{O}_{X_n}$ comme définis sur un même
espace topologique sous-jacent, et de même pour les
$\mathcal{O}_{X_n}$-modules cohérents. Posons
$$
\mathcal{F}_n =
\Ker(\mathcal{O}_{X_{n + 1}} \to \mathcal{O}_{X_n})
$$
de sorte que l'on obtienne des suites exactes courtes
$$
0 \to \mathcal{F}_n \to \mathcal{O}_{X_{n + 1}} \to \mathcal{O}_{X_n} \to 0
$$
On a, d'après ce qui précède, $\mathcal{F}_n = I^n\mathcal{O}_{X_{n + 1}}$.
Il s'ensuit que $\mathcal{F}_n$ est un faisceau cohérent sur $X_{n + 1}$
annulé par $I$ ; nous pouvons donc le considérer, et le considérerons, comme
un $\mathcal{O}_{X_1}$-module cohérent. Remarquons que, pour $m > n$, le faisceau
$$
I^n\mathcal{O}_{X_m}/I^{n + 1}\mathcal{O}_{X_m}
$$
s'envoie isomorphiquement sur $\mathcal{F}_n$ par le morphisme
$\mathcal{O}_{X_m} \to \mathcal{O}_{X_{n + 1}}$. Étant donnés
$n_1, n_2 \geq 0$, on peut donc choisir $m > n_1 + n_2$ et considérer le morphisme
de multiplication
$$
I^{n_1}\mathcal{O}_{X_m} \times I^{n_2}\mathcal{O}_{X_m}
\longrightarrow
I^{n_1 + n_2}\mathcal{O}_{X_m} \to \mathcal{F}_{n_1 + n_2}
$$
Celui-ci induit une application $\mathcal{O}_{X_1}$-bilinéaire
$$
\mathcal{F}_{n_1} \times \mathcal{F}_{n_2} \longrightarrow
\mathcal{F}_{n_1 + n_2}
$$
qui définit à son tour une structure de $\mathcal{O}_{X_1}$-algèbre graduée
sur $\mathcal{F} = \bigoplus_{n \geq 0} \mathcal{F}_n$.

\medskip\noindent
Posons $B = \bigoplus I^n/I^{n + 1}$ ; c'est
une $A/I$-algèbre graduée de type fini. Posons $\mathcal{B} = (X_1 \to S_1)^*\widetilde{B}$.
La discussion précédente fournit une surjection canonique
$$
\mathcal{B} \longrightarrow \mathcal{F}
$$
de $\mathcal{O}_{X_1}$-algèbres graduées. En particulier,
$\mathcal{F}$ est un $\mathcal{B}$-module gradué quasi-cohérent
de type fini. Le lemme \ref{coherent-lemma-graded-finiteness} permet de choisir un entier $d_0$
tel que $H^1(X_1, \mathcal{F} \otimes \mathcal{L}^{\otimes d}) = 0$
pour tout $d \geq d_0$. Choisissons $d \geq d_0$ de sorte qu'il existe des sections
$s_{0, 1}, \ldots, s_{N, 1} \in \Gamma(X_1, \mathcal{L}_1^{\otimes d})$
induisant une immersion
$$
\psi_1 : X_1 \to \mathbf{P}^N_{S_1}
$$
au-dessus de $S_1$, voir
Morphismes, lemme \ref{morphisms-lemma-finite-type-over-affine-ample-very-ample}.
Comme $X_1$ est propre sur $S_1$, $\psi_1$
est une immersion fermée, voir
Morphismes, lemme \ref{morphisms-lemma-image-proper-scheme-closed},
ainsi que
Schémas, lemme \ref{schemes-lemma-immersion-when-closed}.
Nous allons « relever » $\psi_1$ en un système compatible
d'immersions fermées des $X_n$ dans $\mathbf{P}^N$.

\medskip\noindent
En tensorisant les suites exactes courtes du premier paragraphe
de la démonstration
par $\mathcal{L}_{n + 1}^{\otimes d}$, on obtient les suites exactes courtes
$$
0 \to \mathcal{F}_n \otimes \mathcal{L}_{n + 1}^{\otimes d} \to
\mathcal{L}_{n + 1}^{\otimes d} \to \mathcal{L}_{n + 1}^{\otimes d} \to 0
$$
À l'aide des isomorphismes $\varphi_n$, on obtient des isomorphismes
$\mathcal{L}_{n + 1} \otimes \mathcal{O}_{X_l} = \mathcal{L}_l$
pour $l \leq n$. La suite précédente devient donc
$$
0 \to \mathcal{F}_n \otimes \mathcal{L}_1^{\otimes d} \to
\mathcal{L}_{n + 1}^{\otimes d} \to \mathcal{L}_n^{\otimes d} \to 0
$$
L'annulation de $H^1(X, \mathcal{F}_n \otimes \mathcal{L}_1^{\otimes d})$
permet de relever par récurrence
$s_{0, 1}, \ldots, s_{N, 1} \in \Gamma(X_1, \mathcal{L}_1^{\otimes d})$
en des sections
$s_{0, n}, \ldots, s_{N, n} \in \Gamma(X_n, \mathcal{L}_n^{\otimes d})$.
On obtient ainsi un diagramme commutatif
$$
\xymatrix{
X_1 \ar[r]_{i_1} \ar[d]_{\psi_1} &
X_2 \ar[r]_{i_2} \ar[d]_{\psi_2} &
X_3 \ar[r] \ar[d]_{\psi_3} &
\ldots \\
\mathbf{P}^N_{S_1} \ar[r] &
\mathbf{P}^N_{S_2} \ar[r] &
\mathbf{P}^N_{S_3} \ar[r] & \ldots
}
$$
où
$\psi_n = \varphi_{(\mathcal{L}_n, (s_{0, n}, \ldots, s_{N, n}))}$
avec les notations de Constructions, section
\ref{constructions-section-projective-space}.
Les carrés de l'énoncé du théorème étant
cartésiens, ceux du diagramme ci-dessus le sont
aussi. Le lemme \ref{coherent-lemma-algebraize-formal-closed-subscheme} permet alors de conclure.
\end{proof}





















% OK, start here.
%

\chapter{Diviseurs}



\phantomsection
\label{divisors-section-phantom}


\section{Introduction}
\label{divisors-section-introduction}

\noindent
Dans ce chapitre, nous étudions quelques questions très élémentaires liées à la
définition des diviseurs, etc. Une référence de base est  \cite{EGA}.



\section{Points associés}
\label{divisors-section-associated}

\noindent
Soit  $R$  un anneau et soit  $M$  un  $R$-module. Rappelons
qu'un idéal premier  $\mathfrak p \subset R$  est  {\it  associé}  à
 $M$  s'il existe un élément de  $M$  dont l'annulateur est
 $\mathfrak p$. Voir Algèbre, Définition  \ref{algebra-definition-associated}. Voici la définition des
points associés des faisceaux quasi-cohérents sur les schémas telle qu'elle est donnée
dans  \cite[IV définition 3.1.1]{EGA}.

\begin{definition}
\label{divisors-definition-associated}
Soit  $X$  un schéma. Soit  $\mathcal{F}$  un faisceau quasi-cohérent sur
 $X$.
\begin{enumerate}
\item On dit que  $x \in X$  est  {\it  associé}  à
 $\mathcal{F}$  si l'idéal maximal  $\mathfrak m_x$  est associé au  $\mathcal{O}_{X, x}$-module
 $\mathcal{F}_x$.
\item Nous notons  $\text{Ass}(\mathcal{F})$  ou  $\text{Ass}_X(\mathcal{F})$
l'ensemble des points associés de
 $\mathcal{F}$.
\item Les points associés  {\it  de  $X$}  sont les
points associés de  $\mathcal{O}_X$.
\end{enumerate}
\end{definition}

\noindent
Ces définitions sont les plus utiles lorsque  $X$  est localement
noethérien et  $\mathcal{F}$  de type fini. Par exemple, il peut arriver qu'un point
générique d'une composante irréductible de  $X$  ne soit pas associé à
 $X$, voir l'exemple  \ref{divisors-example-no-associated-prime}. Dans le cas non noethérien, il peut
être plus pratique d'utiliser les points faiblement associés, voir la section
 \ref{divisors-section-weakly-associated}. Relions la notion schématique à la notion algébrique sur les
ouverts affines; notons que cette correspondance ne fonctionne parfaitement que
pour les schémas localement noethériens.

\begin{lemma}
\label{divisors-lemma-associated-affine-open}
Soit  $X$  un schéma. Soit  $\mathcal{F}$  un faisceau quasi-cohérent sur
 $X$. Soit  $\Spec(A) = U \subset X$  un ouvert affine, et posons  $M = \Gamma(U, \mathcal{F})$. Soit
 $x \in U$, et soit  $\mathfrak p \subset A$  l'idéal premier correspondant.
\begin{enumerate}
\item Si  $\mathfrak p$  est associé à  $M$, alors  $x$  est associé à
 $\mathcal{F}$.
\item Si  $\mathfrak p$  est de type fini, alors la réciproque est également vraie.
\end{enumerate}
En particulier, si  $X$  est localement noethérien, alors l'équivalence
$$
\mathfrak p \in \text{Ass}(M) \Leftrightarrow x \in \text{Ass}(\mathcal{F})
$$
vaut pour toutes les paires  $(\mathfrak p, x)$  comme ci-dessus.
\end{lemma}

\begin{proof}
Cela découle de Algèbre, Lemme  \ref{algebra-lemma-associated-primes-localize}. Mais nous pouvons aussi raisonner
directement comme suit. Supposons que  $\mathfrak p$  soit associé à  $M$.
Alors il existe un  $m \in M$  dont l'annulateur est  $\mathfrak p$. Comme la
localisation est exacte, nous voyons que  $\mathfrak pA_{\mathfrak p}$  est l'annulateur de
 $m/1 \in M_{\mathfrak p}$. Comme  $M_{\mathfrak p} = \mathcal{F}_x$  (Schémas, Lemme  \ref{schemes-lemma-spec-sheaves}), nous concluons que
 $x$  est associé à  $\mathcal{F}$.

\medskip\noindent
Inversement, supposons que  $x$  soit associé à  $\mathcal{F}$, et que
 $\mathfrak p$  soit de type fini. Comme  $x$  est associé à  $\mathcal{F}$ , il
existe un élément  $m' \in M_{\mathfrak p}$  dont l'annulateur est  $\mathfrak pA_{\mathfrak p}$. Écrivons
 $m' = m/f$  pour quelque  $f \in A$,  $f \not \in \mathfrak p$. L'annulateur  $I$  de
 $m$  est un idéal de  $A$  tel que  $IA_{\mathfrak p} = \mathfrak pA_{\mathfrak p}$. Par conséquent
 $I \subset \mathfrak p$, et  $(\mathfrak p/I)_{\mathfrak p} = 0$. Puisque  $\mathfrak p$  est de type fini, il existe un
 $g \in A$,  $g \not \in \mathfrak p$  tel que  $g(\mathfrak p/I) = 0$. Ainsi, l'annulateur de
 $gm$  est  $\mathfrak p$, d'où le résultat.

\medskip\noindent
Si  $X$  est localement noethérien, alors  $A$  est noethérien
(Propriétés, Lemme  \ref{properties-lemma-locally-Noetherian}) et  $\mathfrak p$  est toujours de type fini.
\end{proof}

\begin{lemma}
\label{divisors-lemma-ass-support}
Soit  $X$  un schéma. Soit  $\mathcal{F}$  un  $\mathcal{O}_X$-module
quasi-cohérent. Alors  $\text{Ass}(\mathcal{F}) \subset \text{Supp}(\mathcal{F})$.
\end{lemma}

\begin{proof}
Cela découle immédiatement des définitions.
\end{proof}

\begin{lemma}
\label{divisors-lemma-ses-ass}
Soit  $X$  un schéma. Soit  $0 \to \mathcal{F}_1 \to \mathcal{F}_2 \to \mathcal{F}_3 \to 0$  une suite exacte de
faisceaux quasi-cohérents sur  $X$. Alors  $\text{Ass}(\mathcal{F}_2) \subset
\text{Ass}(\mathcal{F}_1) \cup \text{Ass}(\mathcal{F}_3)$  et  $\text{Ass}(\mathcal{F}_1) \subset \text{Ass}(\mathcal{F}_2)$.
\end{lemma}

\begin{proof}
Pour chaque point  $x \in X$, la suite des germes  $0 \to \mathcal{F}_{1, x} \to \mathcal{F}_{2, x} \to \mathcal{F}_{3, x} \to 0$  est une suite
exacte de  $\mathcal{O}_{X, x}$-modules. Par conséquent, le lemme découle de
Algèbre, Lemme  \ref{algebra-lemma-ass}.
\end{proof}

\begin{lemma}
\label{divisors-lemma-finite-ass}
Soit  $X$  un schéma localement noethérien. Soit  $\mathcal{F}$  un
 $\mathcal{O}_X$-module cohérent. Alors  $\text{Ass}(\mathcal{F}) \cap U$  est fini pour tout ouvert
quasi-compact  $U \subset X$.
\end{lemma}

\begin{proof}
C'est vrai car l'ensemble des idéaux premiers associés d'un module fini sur un
anneau noethérien est fini, voir Algèbre, Lemme  \ref{algebra-lemma-finite-ass}. Pour passer des
schémas à l'algèbre, utilisons le fait que  $U$  est une réunion finie d'ouverts
affines, chacun de ces ouverts est le spectre d'un anneau noethérien (Propriétés,
Lemme  \ref{properties-lemma-locally-Noetherian}),  $\mathcal{F}$  correspond à un module fini sur cet anneau
(Cohomologie des schémas, Lemme  \ref{coherent-lemma-coherent-Noetherian}), et enfin le Lemme
 \ref{divisors-lemma-associated-affine-open}.
\end{proof}

\begin{lemma}
\label{divisors-lemma-ass-zero}
Soit  $X$  un schéma localement noethérien. Soit  $\mathcal{F}$  un
 $\mathcal{O}_X$-module quasi-cohérent. Alors
$$
\mathcal{F} = 0 \Leftrightarrow \text{Ass}(\mathcal{F}) = \emptyset.
$$
\end{lemma}

\begin{proof}
Si  $\mathcal{F} = 0$, alors  $\text{Ass}(\mathcal{F}) = \emptyset$  par définition. Réciproquement, si
 $\text{Ass}(\mathcal{F}) = \emptyset$, alors  $\mathcal{F} = 0$  par Algèbre, Lemme  \ref{algebra-lemma-ass-zero}. Pour passer des
schémas à l'algèbre, restreignons-nous à chaque ouvert affine et appliquons le Lemme
 \ref{divisors-lemma-associated-affine-open}.
\end{proof}

\begin{example}
\label{divisors-example-no-associated-prime}
Soit  $k$  un corps. L'anneau  $R = k[x_1, x_2, x_3, \ldots]/(x_i^2)$  est local avec un idéal maximal
localement nilpotent  $\mathfrak m$. Il n'existe aucun élément de  $R$  dont
l'annulateur soit  $\mathfrak m$. Donc  $\text{Ass}(R) = \emptyset$, et  $X = \Spec(R)$  est un exemple
d'un schéma qui n'a pas de points associés.
\end{example}

\begin{lemma}
\label{divisors-lemma-restriction-injective-open-contains-ass}
Soit  $X$  un schéma localement noethérien. Soit  $\mathcal{F}$  un
 $\mathcal{O}_X$-module quasi-cohérent. Si  $U \subset X$  est ouvert et  $\text{Ass}(\mathcal{F}) \subset U$,
alors  $\Gamma(X, \mathcal{F}) \to \Gamma(U, \mathcal{F})$  est injectif.
\end{lemma}

\begin{proof}
Soit  $s \in \Gamma(X, \mathcal{F})$  une section qui se restreint à zéro sur  $U$. Soit
 $\mathcal{F}' \subset \mathcal{F}$  l'image de l'application  $\mathcal{O}_X \to \mathcal{F}$  définie par  $s$. Alors
 $\text{Supp}(\mathcal{F}') \cap U = \emptyset$. D'autre part,  $\text{Ass}(\mathcal{F}') \subset \text{Ass}(\mathcal{F})$  d'après le Lemme  \ref{divisors-lemma-ses-ass}. Puisque
également  $\text{Ass}(\mathcal{F}') \subset \text{Supp}(\mathcal{F}')$  (Lemme  \ref{divisors-lemma-ass-support}), nous concluons  $\text{Ass}(\mathcal{F}') = \emptyset$. Par
conséquent,  $\mathcal{F}' = 0$  d'après le Lemme  \ref{divisors-lemma-ass-zero}.
\end{proof}

\begin{lemma}
\label{divisors-lemma-minimal-support-in-ass}
Soit  $X$  un schéma localement noethérien. Soit  $\mathcal{F}$  un
 $\mathcal{O}_X$-module quasi-cohérent. Soit  $x \in \text{Supp}(\mathcal{F})$  un point dans le support de
 $\mathcal{F}$  qui n'est pas une spécialisation d'un autre point de  $\text{Supp}(\mathcal{F})$.
Alors  $x \in \text{Ass}(\mathcal{F})$. En particulier, tout point générique d'une composante
irréductible de  $X$  est un point associé de  $X$.
\end{lemma}

\begin{proof}
Comme  $x \in \text{Supp}(\mathcal{F})$  le module  $\mathcal{F}_x$  n'est pas nul. Par conséquent
 $\text{Ass}(\mathcal{F}_x) \subset \Spec(\mathcal{O}_{X, x})$  n'est pas vide selon Algèbre, Lemme  \ref{algebra-lemma-ass-zero}. D'autre part, par
hypothèse  $\text{Supp}(\mathcal{F}_x) = \{\mathfrak m_x\}$. Comme  $\text{Ass}(\mathcal{F}_x) \subset \text{Supp}(\mathcal{F}_x)$  (Algèbre, Lemme  \ref{algebra-lemma-ass-support}) nous
voyons que  $\mathfrak m_x$  est associé à  $\mathcal{F}_x$, d'où le résultat.
\end{proof}

\noindent
Le lemme suivant est analogue à Compléments d'algèbre, Lemme  \ref{more-algebra-lemma-check-injective-on-ass}.

\begin{lemma}
\label{divisors-lemma-check-injective-on-ass}
Soit $X$ un schéma localement noethérien. Soit
$\varphi : \mathcal{F} \to \mathcal{G}$ un morphisme de
$\mathcal{O}_X$-modules quasi-cohérents. Supposons que, pour
chaque $x \in X$, au moins l'une des conditions suivantes soit vérifiée
\begin{enumerate}
\item $\mathcal{F}_x \to \mathcal{G}_x$  est injectif, ou
\item $x \not \in \text{Ass}(\mathcal{F})$.
\end{enumerate}
Alors  $\varphi$  est injectif.
\end{lemma}

\begin{proof}
Les hypothèses impliquent que  $\text{Ass}(\Ker(\varphi)) = \emptyset$  et donc  $\Ker(\varphi) = 0$  par le Lemme
 \ref{divisors-lemma-ass-zero}.
\end{proof}

\begin{lemma}
\label{divisors-lemma-check-isomorphism-via-depth-and-ass}
Soit $X$ un schéma localement noethérien. Soit $\varphi : \mathcal{F} \to \mathcal{G}$ un
morphisme de $\mathcal{O}_X$-modules quasi-cohérents. Supposons $\mathcal{F}$ cohérent
et que, pour tout $x \in X$, l'une des situations suivantes se produise
\begin{enumerate}
\item $\mathcal{F}_x \to \mathcal{G}_x$  est un isomorphisme, ou
\item $\text{depth}(\mathcal{F}_x) \geq 2$  et  $x \not \in \text{Ass}(\mathcal{G})$.
\end{enumerate}
Alors  $\varphi$  est un isomorphisme.
\end{lemma}

\begin{proof}
Ceci est la traduction de Compléments d'algèbre, Lemme  \ref{more-algebra-lemma-check-isomorphism-via-depth-and-ass}, dans le langage des
schémas.
\end{proof}




\section{Morphismes et points associés}
\label{divisors-section-morphisms-associated}

\noindent
Soit  $f : X \to S$  un morphisme de schémas. Soit  $\mathcal{F}$  un faisceau de
 $\mathcal{O}_X$-modules. Si  $s \in S$  est un point, il est souvent commode de
noter  $\mathcal{F}_s$  le  $\mathcal{O}_{X_s}$-module obtenu en prenant l'image réciproque de
 $\mathcal{F}$  par le morphisme  $i_s : X_s \to X$. Ici,  $X_s$  est la fibre au sens
des schémas de  $f$  au-dessus de  $s$. Dans une formule
$$
\mathcal{F}_s = i_s^*\mathcal{F}
$$
Bien sûr, cette notation entre en conflit avec la notation déjà existante pour la
fibre de  $\mathcal{F}$  en un point  $x \in X$  si  $f = \text{id}_X$. Cependant, la
notation est souvent pratique, comme dans la formulation du lemme suivant.

\begin{lemma}
\label{divisors-lemma-bourbaki}
Soit  $f : X \to S$  un morphisme de schémas. Soit  $\mathcal{F}$  un faisceau
quasi-cohérent sur  $X$  qui est plat sur  $S$. Soit  $\mathcal{G}$ 
un faisceau quasi-cohérent sur  $S$. Alors nous avons
$$
\text{Ass}_X(\mathcal{F} \otimes_{\mathcal{O}_X} f^*\mathcal{G})
\supset
\bigcup\nolimits_{s \in \text{Ass}_S(\mathcal{G})}
\text{Ass}_{X_s}(\mathcal{F}_s)
$$
et l'égalité tient si  $S$  est localement noethérien (pour la notation
 $\mathcal{F}_s$  voir ci-dessus).
\end{lemma}

\begin{proof}
Soit  $x \in X$  et soit  $s = f(x) \in S$. Posons  $B = \mathcal{O}_{X, x}$,  $A = \mathcal{O}_{S, s}$,
 $N = \mathcal{F}_x$ et  $M = \mathcal{G}_s$. Remarquez que la fibre de  $\mathcal{F} \otimes_{\mathcal{O}_X} f^*\mathcal{G}$  en
 $x$  est égale au  $B$-module  $M \otimes_A N$. Ainsi  $x \in \text{Ass}_X(\mathcal{F} \otimes_{\mathcal{O}_X} f^*\mathcal{G})$  si
et seulement si  $\mathfrak m_B$  est dans  $\text{Ass}_B(M \otimes_A N)$. De même  $s \in \text{Ass}_S(\mathcal{G})$  et
 $x \in \text{Ass}_{X_s}(\mathcal{F}_s)$  si et seulement si  $\mathfrak m_A \in \text{Ass}_A(M)$  et  $\mathfrak m_B/\mathfrak m_A B \in
\text{Ass}_{B \otimes \kappa(\mathfrak m_A)}(N \otimes \kappa(\mathfrak m_A))$. Ainsi, le lemme
découle de Algèbre, Lemme  \ref{algebra-lemma-bourbaki-fibres}.
\end{proof}





\section{Points immergés}
\label{divisors-section-embedded}

\noindent
Soit  $R$  un anneau et soit  $M$  un  $R$-module. Rappelons
qu'un idéal premier  $\mathfrak p \subset R$  est un  {\it  idéal premier associé
immergé}  de  $M$  s'il est un idéal premier associé à
 $M$  qui n'est pas minimal parmi les idéaux premiers associés de
 $M$. Voir Algèbre, Définition  \ref{algebra-definition-embedded-primes}. Voici la définition des
points associés immergés pour les faisceaux quasi-cohérents sur les schémas telle
qu'elle est donnée dans  \cite[IV définition 3.1.1]{EGA}.

\begin{definition}
\label{divisors-definition-embedded}
Soit  $X$  un schéma. Soit  $\mathcal{F}$  un faisceau quasi-cohérent sur
 $X$.
\begin{enumerate}
\item Un  {\it  point associé immergé}  de  $\mathcal{F}$  est
un point associé qui n'est pas maximal parmi les points associés de  $\mathcal{F}$,
c'est-à-dire qu'il est la spécialisation d'un autre point associé de  $\mathcal{F}$.
\item Un point  $x$  de  $X$  est appelé un  {\it  point
immergé}  si  $x$  est un point associé immergé de
 $\mathcal{O}_X$.
\item Une  {\it  composante immergée}  de  $X$  est un
sous-ensemble fermé irréductible  $Z = \overline{\{x\}}$  où  $x$  est un point
immergé de  $X$.
\end{enumerate}
\end{definition}

\noindent
Dans le cas noethérien lorsque  $\mathcal{F}$  est cohérent, nous avons ce qui suit.

\begin{lemma}
\label{divisors-lemma-embedded}
Soit  $X$  un schéma localement noethérien. Soit  $\mathcal{F}$  un
 $\mathcal{O}_X$-module cohérent. Alors
\begin{enumerate}
\item les points génériques des composantes irréductibles de  $\text{Supp}(\mathcal{F})$  sont des
points associés de  $\mathcal{F}$, et
\item un point associé de  $\mathcal{F}$  est immergé si et seulement s'il n'est pas un
point générique d'une composante irréductible de  $\text{Supp}(\mathcal{F})$.
\end{enumerate}
En particulier un point immergé de $X$  est un point associé de
 $X$ qui n'est pas un point générique d'une composante irréductible
de $X$.
\end{lemma}

\begin{proof}
Rappelons que dans ce cas $Z = \text{Supp}(\mathcal{F})$ est fermé, voir Morphismes,
Lemme \ref{morphisms-lemma-support-finite-type}  et que les points génériques des composantes irréductibles de
 $Z$ sont des points associés de $\mathcal{F}$, voir le Lemme  \ref{divisors-lemma-minimal-support-in-ass}.
Enfin, nous avons  $\text{Ass}(\mathcal{F}) \subset Z$, par le Lemme  \ref{divisors-lemma-ass-support}. Ces résultats, combinés au
fait que  $Z$ est un espace topologique sobre et donc chaque point
de $Z$ est une spécialisation d'un point générique de $Z$, implique
(1) et (2).
\end{proof}

\begin{lemma}
\label{divisors-lemma-S1-no-embedded}
Soit $X$ un schéma localement noethérien. Soit $\mathcal{F}$ un
faisceau cohérent sur $X$. Alors les assertions suivantes sont équivalentes:
\begin{enumerate}
\item $\mathcal{F}$  n'a pas de points associés immergés, et
\item $\mathcal{F}$  a la propriété  $(S_1)$.
\end{enumerate}
\end{lemma}

\begin{proof}
Ceci résulte du Lemme  \ref{algebra-lemma-criterion-no-embedded-primes} d'Algèbre, combiné au Lemme  \ref{divisors-lemma-associated-affine-open}
ci-dessus.
\end{proof}

\begin{lemma}
\label{divisors-lemma-noetherian-dim-1-CM-no-embedded-points}
Soit  $X$  un schéma localement noethérien de dimension  $\leq 1$. Les
énoncés suivants sont équivalents
\begin{enumerate}
\item $X$  est de Cohen-Macaulay, et
\item $X$  n'a pas de points immergés.
\end{enumerate}
\end{lemma}

\begin{proof}
Cela découle du lemme  \ref{divisors-lemma-S1-no-embedded}  et des définitions.
\end{proof}

\begin{lemma}
\label{divisors-lemma-scheme-theoretically-dense-contain-embedded-points}
Soit  $X$  un schéma localement noethérien. Soit  $U \subset X$  un
sous-schéma ouvert. Les énoncés suivants sont équivalents
\begin{enumerate}
\item $U$  est schématiquement dense dans  $X$  (Morphismes, Définition
 \ref{morphisms-definition-scheme-theoretically-dense}),
\item $U$  est dense dans  $X$  et  $U$  contient tous les points
immergés de  $X$, et
\item $U$  contient tous les points associés de  $X$.
\end{enumerate}
\end{lemma}

\begin{proof}
La question est locale sur  $X$, par conséquent nous pouvons supposer que
 $X = \Spec(A)$  où  $A$  est un anneau noethérien. Alors  $U$  est
quasi-compact (Propriétés, Lemme  \ref{properties-lemma-immersion-into-noetherian}) donc  $U = D(f_1) \cup \ldots \cup D(f_n)$  (Algèbre, Lemme
 \ref{algebra-lemma-qc-open}). Dans cette situation,  $U$  est schématiquement dense dans
 $X$  si et seulement si  $A \to A_{f_1} \times \ldots \times A_{f_n}$  est injectif, voir Morphismes,
Exemple  \ref{morphisms-example-scheme-theoretic-closure}. La condition (3) traduite en algèbre signifie que pour
chaque idéal premier associé  $\mathfrak p$  à  $A$  il existe un
 $i$  avec  $f_i \not \in \mathfrak p$.

\medskip\noindent
Supposons (1), c'est-à-dire que  $A \to A_{f_1} \times \ldots \times A_{f_n}$  est injectif. Soit  $a \in A$
un élément dont l'annulateur est l'idéal premier  $\mathfrak p$. Puisque  $a$  s'envoie sur un
élément non nul de  $A_{f_i}$  pour un certain  $i$  nous voyons que
 $f_i \not \in \mathfrak p$. Ainsi, (3) est vérifiée.

\medskip\noindent
Notons que  $U$  est dense dans  $\Spec(A)$  si et seulement s'il contient
tous les idéaux premiers minimaux. Comme l'ensemble des idéaux premiers associés à
 $A$  est égal à l'union de l'ensemble des idéaux premiers minimaux et de
l'ensemble des idéaux premiers associés immergés (voir Algèbre, Proposition
 \ref{algebra-proposition-minimal-primes-associated-primes}), nous voyons que (2) et (3) sont équivalents.

\medskip\noindent
Supposons (3). Cela signifie que chaque idéal premier associé  $\mathfrak p$  à
 $A$  correspond à un idéal premier de  $A_{f_i}$  pour un certain  $i$.
Alors  $A \to A_{f_1} \times \ldots \times A_{f_n}$  est injectif parce que  $A \to \prod_{\mathfrak p \in \text{Ass}(A)} A_\mathfrak p$  est injectif selon Algèbre,
Lemme  \ref{algebra-lemma-zero-at-ass-zero}. Ainsi nous voyons que (3) implique (1).
\end{proof}

\begin{lemma}
\label{divisors-lemma-remove-embedded-points}
Soit  $X$  un schéma localement noethérien. Soit  $\mathcal{F}$  un faisceau
cohérent sur  $X$. L'ensemble des sous-faisceaux cohérents
$$
\{
\mathcal{K} \subset \mathcal{F}
\mid
\text{Supp}(\mathcal{K})\text{ is nowhere dense in }\text{Supp}(\mathcal{F})
\}
$$
contient un élément maximal  $\mathcal{K}$. En posant  $\mathcal{F}' = \mathcal{F}/\mathcal{K}$ , nous avons ce
qui suit
\begin{enumerate}
\item $\text{Supp}(\mathcal{F}') = \text{Supp}(\mathcal{F})$,
\item $\mathcal{F}'$  n'a pas de points associés immergés, et
\item il existe un ouvert dense  $U \subset X$  tel que  $U \cap \text{Supp}(\mathcal{F})$  est dense dans
 $\text{Supp}(\mathcal{F})$  et  $\mathcal{F}'|_U \cong \mathcal{F}|_U$.
\end{enumerate}
\end{lemma}

\begin{proof}
Cela découle de l'Algèbre, Lemme  \ref{algebra-lemma-remove-embedded-primes}  et  \ref{algebra-lemma-remove-embedded-primes-localize}. Notons que
 $U$  peut être pris comme le complément de l'adhérence de l'ensemble des
points associés immergés de  $\mathcal{F}$.
\end{proof}

\begin{lemma}
\label{divisors-lemma-no-embedded-points-endos}
Soit  $X$  un schéma localement noethérien. Soit  $\mathcal{F}$  un
 $\mathcal{O}_X$-module cohérent sans points associés immergés. Posons
$$
\mathcal{I}
=
\Ker(\mathcal{O}_X
\longrightarrow
\SheafHom_{\mathcal{O}_X}(\mathcal{F}, \mathcal{F})).
$$
Il s'agit d'un faisceau cohérent d'idéaux qui définit un sous-schéma fermé
 $Z \subset X$  sans points immergés. De plus, il existe un faisceau cohérent
 $\mathcal{G}$  sur  $Z$  tel que (a)  $\mathcal{F} = (Z \to X)_*\mathcal{G}$, (b)  $\mathcal{G}$  n'a pas
de points associés immergés, et (c)  $\text{Supp}(\mathcal{G}) = Z$  (comme ensembles).
\end{lemma}

\begin{proof}
Certains de ces énoncés ont été vus dans la preuve de Cohomologie des
schémas, Lemme  \ref{coherent-lemma-coherent-support-closed}. Les autres découlent d'Algèbre, Lemme  \ref{algebra-lemma-no-embedded-primes-endos}.
\end{proof}



\section{Points faiblement associés}
\label{divisors-section-weakly-associated}

\noindent
Soit  $R$  un anneau et soit  $M$  un  $R$-module. Rappelons
qu'un idéal premier  $\mathfrak p \subset R$  est  {\it  faiblement
associé}  à  $M$  s'il existe un élément  $m$  de
 $M$  tel que  $\mathfrak p$  soit minimal parmi les idéaux premiers contenant
l'annulateur de  $m$. Voir Algèbre, Définition  \ref{algebra-definition-weakly-associated}. Si
 $R$  est un anneau local avec un idéal maximal  $\mathfrak m$, alors
 $\mathfrak m$  est faiblement associé à  $M$  si et seulement s'il existe un
élément  $m \in M$  dont l'annulateur a pour radical  $\mathfrak m$, voir Algèbre,
Lemme  \ref{algebra-lemma-weakly-ass-local}.

\begin{definition}
\label{divisors-definition-weakly-associated}
Soit  $X$  un schéma. Soit  $\mathcal{F}$  un faisceau quasi-cohérent sur
 $X$.
\begin{enumerate}
\item Nous disons que  $x \in X$  est  {\it faiblement
associé} à $\mathcal{F}$ si l'idéal maximal $\mathfrak m_x$ est faiblement
associé au $\mathcal{O}_{X, x}$-module $\mathcal{F}_x$.
\item Nous notons  $\text{WeakAss}(\mathcal{F})$  l'ensemble des points faiblement associés de
 $\mathcal{F}$.
\item Les points faiblement associés  {\it  de
 $X$}  sont les points faiblement associés de  $\mathcal{O}_X$.
\end{enumerate}
\end{definition}

\noindent
Dans ce cas, sur tout ouvert affine, cela correspond exactement aux idéaux
premiers faiblement associés tels que définis ci-dessus. Voici l'énoncé précis.

\begin{lemma}
\label{divisors-lemma-weakly-associated-affine-open}
Soit  $X$  un schéma. Soit  $\mathcal{F}$  un faisceau quasi-cohérent sur
 $X$. Soit  $\Spec(A) = U \subset X$  un ouvert affine, et posons  $M = \Gamma(U, \mathcal{F})$. Soit
 $x \in U$, et soit  $\mathfrak p \subset A$  l'idéal premier correspondant. Les énoncés suivants
sont équivalents
\begin{enumerate}
\item $\mathfrak p$  est faiblement associé à  $M$, et
\item $x$  est faiblement associé à  $\mathcal{F}$.
\end{enumerate}
\end{lemma}

\begin{proof}
Cela découle de Algèbre, Lemme  \ref{algebra-lemma-weakly-ass-local}.
\end{proof}

\begin{lemma}
\label{divisors-lemma-weakly-ass-support}
Soit  $X$  un schéma. Soit  $\mathcal{F}$  un  $\mathcal{O}_X$-module
quasi-cohérent. Alors
$$
\text{Ass}(\mathcal{F}) \subset \text{WeakAss}(\mathcal{F}) \subset
\text{Supp}(\mathcal{F}).
$$
\end{lemma}

\begin{proof}
Cela découle immédiatement des définitions.
\end{proof}

\begin{lemma}
\label{divisors-lemma-ses-weakly-ass}
Soit  $X$  un schéma. Soit  $0 \to \mathcal{F}_1 \to \mathcal{F}_2 \to \mathcal{F}_3 \to 0$  une suite exacte de
faisceaux quasi-cohérents sur  $X$. Alors  $\text{WeakAss}(\mathcal{F}_2) \subset
\text{WeakAss}(\mathcal{F}_1) \cup \text{WeakAss}(\mathcal{F}_3)$  et  $\text{WeakAss}(\mathcal{F}_1) \subset \text{WeakAss}(\mathcal{F}_2)$.
\end{lemma}

\begin{proof}
Pour chaque point  $x \in X$  la suite des germes  $0 \to \mathcal{F}_{1, x} \to \mathcal{F}_{2, x} \to \mathcal{F}_{3, x} \to 0$  est une suite
exacte de  $\mathcal{O}_{X, x}$-modules. Par conséquent, le lemme découle de
Algèbre, Lemme  \ref{algebra-lemma-weakly-ass}.
\end{proof}

\begin{lemma}
\label{divisors-lemma-weakly-ass-zero}
Soit  $X$  un schéma. Soit  $\mathcal{F}$  un  $\mathcal{O}_X$-module
quasi-cohérent. Alors
$$
\mathcal{F} = (0) \Leftrightarrow \text{WeakAss}(\mathcal{F}) = \emptyset
$$
\end{lemma}

\begin{proof}
Découle du Lemme  \ref{divisors-lemma-weakly-associated-affine-open}  et de Algèbre, Lemme  \ref{algebra-lemma-weakly-ass-zero}
\end{proof}

\begin{lemma}
\label{divisors-lemma-restriction-injective-open-contains-weakly-ass}
Soit  $X$  un schéma. Soit  $\mathcal{F}$  un  $\mathcal{O}_X$-module
quasi-cohérent. Si  $U \subset X$  est ouvert et  $\text{WeakAss}(\mathcal{F}) \subset U$, alors  $\Gamma(X, \mathcal{F}) \to \Gamma(U, \mathcal{F})$  est
injectif.
\end{lemma}

\begin{proof}
Soit  $s \in \Gamma(X, \mathcal{F})$  une section qui se restreint à zéro sur  $U$. Soit
 $\mathcal{F}' \subset \mathcal{F}$  l'image de l'application  $\mathcal{O}_X \to \mathcal{F}$  définie par  $s$. Alors
 $\text{Supp}(\mathcal{F}') \cap U = \emptyset$. D'autre part,  $\text{WeakAss}(\mathcal{F}') \subset \text{WeakAss}(\mathcal{F})$  d'après le Lemme  \ref{divisors-lemma-ses-weakly-ass}. Puisque
également  $\text{WeakAss}(\mathcal{F}') \subset \text{Supp}(\mathcal{F}')$  (Lemme  \ref{divisors-lemma-weakly-ass-support}), nous concluons  $\text{WeakAss}(\mathcal{F}') = \emptyset$. Par
conséquent,  $\mathcal{F}' = 0$  d'après le Lemme  \ref{divisors-lemma-weakly-ass-zero}.
\end{proof}

\begin{lemma}
\label{divisors-lemma-minimal-support-in-weakly-ass}
Soit  $X$  un schéma. Soit  $\mathcal{F}$  un  $\mathcal{O}_X$-module
quasi-cohérent. Soit  $x \in \text{Supp}(\mathcal{F})$  un point dans le support de  $\mathcal{F}$  qui
n'est pas une spécialisation d'un autre point de  $\text{Supp}(\mathcal{F})$. Alors  $x \in \text{WeakAss}(\mathcal{F})$.
En particulier, tout point générique d'une composante irréductible de  $X$ 
est faiblement associé à  $\mathcal{O}_X$.
\end{lemma}

\begin{proof}
Puisque  $x \in \text{Supp}(\mathcal{F})$  le module  $\mathcal{F}_x$  n'est pas nul. Par conséquent,
 $\text{WeakAss}(\mathcal{F}_x) \subset \Spec(\mathcal{O}_{X, x})$  est non vide d'après Algèbre, Lemme  \ref{algebra-lemma-weakly-ass-zero}. D'autre part, par
hypothèse  $\text{Supp}(\mathcal{F}_x) = \{\mathfrak m_x\}$. Comme  $\text{WeakAss}(\mathcal{F}_x) \subset \text{Supp}(\mathcal{F}_x)$  (Algèbre, Lemme  \ref{algebra-lemma-weakly-ass-support}), nous
voyons que  $\mathfrak m_x$  est faiblement associé à  $\mathcal{F}_x$, d'où le résultat.
\end{proof}

\begin{lemma}
\label{divisors-lemma-ass-weakly-ass}
Soit  $X$  un schéma. Soit  $\mathcal{F}$  un  $\mathcal{O}_X$-module
quasi-cohérent. Si  $\mathfrak m_x$  est un idéal de type fini de  $\mathcal{O}_{X, x}$, alors
$$
x \in \text{Ass}(\mathcal{F}) \Leftrightarrow
x \in \text{WeakAss}(\mathcal{F}).
$$
En particulier, si  $X$  est localement noethérien, alors  $\text{Ass}(\mathcal{F}) = \text{WeakAss}(\mathcal{F})$.
\end{lemma}

\begin{proof}
Voir Algèbre, Lemme  \ref{algebra-lemma-ass-weakly-ass}.
\end{proof}

\begin{lemma}
\label{divisors-lemma-weakass-pushforward}
Soit  $f : X \to S$  un morphisme quasi-compact et quasi-séparé de schémas. Soit
 $\mathcal{F}$  un  $\mathcal{O}_X$-module quasi-cohérent. Soit  $s \in S$  un point qui
n'est pas dans l'image de  $f$. Alors  $s$  n'est pas faiblement
associé à  $f_*\mathcal{F}$.
\end{lemma}

\begin{proof}
Considérons le changement de base  $f' : X' \to \Spec(\mathcal{O}_{S, s})$  de  $f$  par le morphisme
 $g : \Spec(\mathcal{O}_{S, s}) \to S$  et notons  $g' : X' \to X$  l'autre projection. Alors
$$
(f_*\mathcal{F})_s = (g^*f_*\mathcal{F})_s = (f'_*(g')^*\mathcal{F})_s
$$
La première égalité parce que  $g$  induit un isomorphisme sur les anneaux
locaux en  $s$  et la seconde par changement de base plat (Cohomologie des
schémas, Lemme  \ref{coherent-lemma-flat-base-change-cohomology}). Bien sûr,  $s \in \Spec(\mathcal{O}_{S, s})$  n'est pas dans l'image de
 $f'$. Ainsi, nous pouvons supposer que  $S$  est le spectre d'un
anneau local  $(A, \mathfrak m)$  et que  $s$  correspond à  $\mathfrak m$. D'après
Schémas, Lemme  \ref{schemes-lemma-push-forward-quasi-coherent} , le faisceau  $f_*\mathcal{F}$  est quasi-cohérent, disons
correspondant au  $A$-module  $M$. Comme  $s$  n'est pas
dans l'image de  $f$ , nous voyons que  $X = \bigcup_{a \in \mathfrak m} f^{-1}D(a)$  est un recouvrement
ouvert. Comme  $X$  est quasi-compact, nous pouvons trouver  $a_1, \ldots, a_n \in \mathfrak m$ 
tels que  $X = f^{-1}D(a_1) \cup \ldots \cup f^{-1}D(a_n)$. Il s'ensuit que
$$
M \to M_{a_1} \oplus \ldots \oplus M_{a_r}
$$
est injectif. Par conséquent, pour tout élément non nul  $m$  du localisé
 $M_\mathfrak p$ , il existe un  $i$  tel que  $a_i^n m$  est non nul pour tout
 $n \geq 0$. Ainsi,  $\mathfrak m$  n'est pas faiblement associé à  $M$.
\end{proof}

\begin{lemma}
\label{divisors-lemma-check-injective-on-weakass}
Soit  $X$  un schéma. Soit  $\varphi : \mathcal{F} \to \mathcal{G}$  une application de
 $\mathcal{O}_X$-modules quasi-cohérents. Supposons que pour chaque  $x \in X$  au
moins l'un des cas suivants se produise
\begin{enumerate}
\item $\mathcal{F}_x \to \mathcal{G}_x$  est injectif, ou
\item $x \not \in \text{WeakAss}(\mathcal{F})$.
\end{enumerate}
Alors  $\varphi$  est injectif.
\end{lemma}

\begin{proof}
Les hypothèses impliquent que  $\text{WeakAss}(\Ker(\varphi)) = \emptyset$  et donc  $\Ker(\varphi) = 0$  par le Lemme
 \ref{divisors-lemma-weakly-ass-zero}.
\end{proof}

\begin{lemma}
\label{divisors-lemma-depth-2-hartog}
Soit  $X$  un schéma localement noethérien. Soit  $\mathcal{F}$  un
 $\mathcal{O}_X$-module cohérent. Soit  $j : U \to X$  un sous-schéma ouvert tel que pour
 $x \in X \setminus U$  nous avons  $\text{depth}(\mathcal{F}_x) \geq 2$. Alors
$$
\mathcal{F} \longrightarrow j_*(\mathcal{F}|_U)
$$
est un isomorphisme et par conséquent  $\Gamma(X, \mathcal{F}) \to \Gamma(U, \mathcal{F})$  est également un isomorphisme.
\end{lemma}

\begin{proof}
Nous affirmons que le Lemme  \ref{divisors-lemma-check-isomorphism-via-depth-and-ass}  s'applique à l'application affichée dans le
lemme. Soit  $x \in X$. Si  $x \in U$, alors l'application est un isomorphisme sur
les germes puisque  $j_*(\mathcal{F}|_U)|_U = \mathcal{F}|_U$. Si  $x \in X \setminus U$, alors  $x \not \in \text{Ass}(j_*(\mathcal{F}|_U))$  (Lemmes
 \ref{divisors-lemma-weakass-pushforward}  et  \ref{divisors-lemma-weakly-ass-support}). Puisque nous avons supposé  $\text{depth}(\mathcal{F}_x) \geq 2$ , cela
termine la preuve.
\end{proof}

\begin{lemma}
\label{divisors-lemma-weakass-reduced}
Soit  $X$  un schéma réduit. Alors les points faiblement associés de
 $X$  sont exactement les points génériques des composantes irréductibles
de  $X$.
\end{lemma}

\begin{proof}
Cela découle de Algèbre, Lemme  \ref{algebra-lemma-reduced-weakly-ass-minimal}.
\end{proof}



\section{Morphismes et points faiblement associés}
\label{divisors-section-morphisms-weakly-associated}

\begin{lemma}
\label{divisors-lemma-weakly-ass-reverse-functorial}
Soit  $f : X \to S$  un morphisme affine de schémas. Soit  $\mathcal{F}$  un
 $\mathcal{O}_X$-module quasi-cohérent. Alors nous avons
$$
\text{WeakAss}_S(f_*\mathcal{F}) \subset f(\text{WeakAss}_X(\mathcal{F}))
$$
\end{lemma}

\begin{proof}
On peut supposer que  $X$  et  $S$  sont affines, donc  $X \to S$ 
provient d'une application d'anneaux  $A \to B$. Alors  $\mathcal{F} = \widetilde M$  pour un
certain  $B$-module  $M$. D'après le Lemme  \ref{divisors-lemma-weakly-associated-affine-open} , les
points faiblement associés de  $\mathcal{F}$  correspondent exactement aux idéaux
idéaux premiers faiblement associés de  $M$. De même, les points
faiblement associés de  $f_*\mathcal{F}$  correspondent exactement aux idéaux
premiers faiblement associés de  $M$  en tant que  $A$-module. Par
conséquent, le lemme découle de Algèbre, Lemme  \ref{algebra-lemma-weakly-ass-reverse-functorial}.
\end{proof}

\begin{lemma}
\label{divisors-lemma-ass-functorial-equal}
Soit  $f : X \to S$  un morphisme affine de schémas. Soit  $\mathcal{F}$  un
 $\mathcal{O}_X$-module quasi-cohérent. Si  $X$  est localement noethérien,
alors nous avons
$$
f(\text{Ass}_X(\mathcal{F})) =
\text{Ass}_S(f_*\mathcal{F}) =
\text{WeakAss}_S(f_*\mathcal{F}) =
f(\text{WeakAss}_X(\mathcal{F}))
$$
\end{lemma}

\begin{proof}
Nous pouvons supposer que  $X$  et  $S$  sont affines, donc
 $X \to S$  provient d'un morphisme d'anneaux  $A \to B$. Comme  $X$ 
est localement noethérien, l'anneau  $B$  est noethérien, voir Propriétés,
Lemme  \ref{properties-lemma-locally-Noetherian}. Écrivons  $\mathcal{F} = \widetilde M$  pour un certain  $B$-module
 $M$. D'après le Lemme  \ref{divisors-lemma-associated-affine-open} , les points associés de  $\mathcal{F}$ 
correspondent exactement aux idéaux premiers associés de  $M$, et tout
idéal premier associé de  $M$  en tant que  $A$-module est un point
associé de  $f_*\mathcal{F}$. Ainsi, l'inclusion
$$
f(\text{Ass}_X(\mathcal{F})) \subset \text{Ass}_S(f_*\mathcal{F})
$$
découle du Lemme  \ref{algebra-lemma-ass-functorial-Noetherian} d'Algèbre. Nous avons l'inclusion
$$
\text{Ass}_S(f_*\mathcal{F}) \subset \text{WeakAss}_S(f_*\mathcal{F})
$$
d'après le Lemme  \ref{divisors-lemma-weakly-ass-support}. Nous avons l'inclusion
$$
\text{WeakAss}_S(f_*\mathcal{F}) \subset f(\text{WeakAss}_X(\mathcal{F}))
$$
d'après le Lemme  \ref{divisors-lemma-weakly-ass-reverse-functorial}. Les ensembles extrêmes sont égaux d'après le Lemme
 \ref{divisors-lemma-ass-weakly-ass} , donc nous avons l'égalité partout.
\end{proof}

\begin{lemma}
\label{divisors-lemma-weakly-associated-finite}
Soit  $f : X \to S$  un morphisme fini de schémas. Soit  $\mathcal{F}$  un
 $\mathcal{O}_X$-module quasi-cohérent. Alors  $\text{WeakAss}(f_*\mathcal{F}) = f(\text{WeakAss}(\mathcal{F}))$.
\end{lemma}

\begin{proof}
Nous pouvons supposer que  $X$  et  $S$  sont affines, donc
 $X \to S$  provient d'un morphisme fini d'anneaux  $A \to B$. Notons
 $\mathcal{F} = \widetilde M$  pour un certain  $B$-module  $M$. D'après le Lemme
 \ref{divisors-lemma-weakly-associated-affine-open} , les points faiblement associés de  $\mathcal{F}$  correspondent
exactement aux idéaux premiers faiblement associés de  $M$. De
même, les points faiblement associés de  $f_*\mathcal{F}$  correspondent exactement aux
idéaux premiers faiblement associés de  $M$  en tant que
 $A$-module. Ainsi, le lemme découle du Lemme d'Algèbre  \ref{algebra-lemma-weakly-ass-finite-ring-map}.
\end{proof}

\begin{lemma}
\label{divisors-lemma-weakly-ass-pullback}
Soit  $f : X \to S$  un morphisme de schémas. Soit  $\mathcal{G}$  un
 $\mathcal{O}_S$-module quasi-cohérent. Soit  $x \in X$  avec  $s = f(x)$. Si
 $f$  est plat en  $x$, le point  $x$  est un point
générique de la fibre  $X_s$, et  $s \in \text{WeakAss}_S(\mathcal{G})$, alors  $x \in \text{WeakAss}(f^*\mathcal{G})$.
\end{lemma}

\begin{proof}
Soient  $A = \mathcal{O}_{S, s}$,  $B = \mathcal{O}_{X, x}$, et  $M = \mathcal{G}_s$. Soit  $m \in M$  un élément
dont l'annulateur  $I = \{a \in A \mid am = 0\}$  a un radical  $\mathfrak m_A$. Alors  $m \otimes 1$  a
pour annulateur  $I B$  car  $A \to B$  est fidèlement plat. Ainsi, il
suffit de voir que  $\sqrt{I B} = \mathfrak m_B$. Cela découle du fait que l'idéal maximal de
 $B/\mathfrak m_AB$  est localement nilpotent (voir Algèbre, Lemme  \ref{algebra-lemma-minimal-prime-reduced-ring}) et de
l'hypothèse que  $\sqrt{I} = \mathfrak m_A$. Certains détails sont omis.
\end{proof}

\begin{lemma}
\label{divisors-lemma-weakly-ass-change-fields}
Soit  $K/k$  une extension de corps. Soit  $X$  un schéma sur
 $k$. Soit  $\mathcal{F}$  un  $\mathcal{O}_X$-module quasi-cohérent. Soit
 $y \in X_K$  avec image  $x \in X$. Si  $y$  est un point faiblement
associé de l'image réciproque  $\mathcal{F}_K$, alors  $x$  est un point
faiblement associé de  $\mathcal{F}$.
\end{lemma}

\begin{proof}
Ceci est la traduction de Algèbre, Lemme  \ref{algebra-lemma-weakly-ass-change-fields}  dans le langage des schémas.
\end{proof}

\noindent
Voici un lemme simple où nous constatons que les images directes ont souvent une
profondeur d'au moins 2.

\begin{lemma}
\label{divisors-lemma-depth-pushforward}
Soit  $f : X \to S$  un morphisme de schémas quasi-compact et quasi-séparé. Soit
 $\mathcal{F}$  un  $\mathcal{O}_X$-module quasi-cohérent. Soit  $s \in S$.
\begin{enumerate}
\item Si  $s \not \in f(X)$, alors  $s$  n'est pas faiblement associé à  $f_*\mathcal{F}$.
\item Si  $s \not \in f(X)$  et  $\mathcal{O}_{S, s}$  est noethérien, alors  $s$  n'est pas
associé à  $f_*\mathcal{F}$.
\item Si  $s \not \in f(X)$,  $(f_*\mathcal{F})_s$  est un  $\mathcal{O}_{S, s}$-module fini, et  $\mathcal{O}_{S, s}$  est
noethérien, alors  $\text{depth}((f_*\mathcal{F})_s) \geq 2$.
\item Si  $\mathcal{F}$  est plat sur  $S$  et  $a \in \mathfrak m_s$  est un non-diviseur de
zéro, alors  $a$  est un non-diviseur de zéro dans  $(f_*\mathcal{F})_s$.
\item Si  $\mathcal{F}$  est plat sur  $S$  et si  $a, b \in \mathfrak m_s$  forment une suite régulière,
alors  $a$  est un non-diviseur de zéro dans  $(f_*\mathcal{F})_s$  et  $b$ 
est un non-diviseur de zéro dans  $(f_*\mathcal{F})_s/a(f_*\mathcal{F})_s$.
\item Si  $\mathcal{F}$  est plat sur  $S$  et  $(f_*\mathcal{F})_s$  est un
 $\mathcal{O}_{S, s}$-module fini, alors  $\text{depth}((f_*\mathcal{F})_s) \geq
\min(2, \text{depth}(\mathcal{O}_{S, s}))$.
\end{enumerate}
\end{lemma}

\begin{proof}
La partie (1) est le Lemme  \ref{divisors-lemma-weakass-pushforward}. La partie (2) découle de (1) et du Lemme
 \ref{divisors-lemma-ass-weakly-ass}.

\medskip\noindent
Preuve de la partie (3). Pour montrer que la profondeur est  $\geq 2$ , il
suffit de montrer que  $\Hom_{\mathcal{O}_{S, s}}(\kappa(s), (f_*\mathcal{F})_s) = 0$  et  $\Ext^1_{\mathcal{O}_{S, s}}(\kappa(s), (f_*\mathcal{F})_s) = 0$, voir Algèbre, Lemme
 \ref{algebra-lemma-depth-ext}. En utilisant la suite exacte  $0 \to \mathfrak m_s \to \mathcal{O}_{S, s} \to \kappa(s) \to 0$ , il suffit de prouver que
l'application
$$
\Hom_{\mathcal{O}_{S, s}}(\mathcal{O}_{S, s}, (f_*\mathcal{F})_s)
\to
\Hom_{\mathcal{O}_{S, s}}(\mathfrak m_s, (f_*\mathcal{F})_s)
$$
est un isomorphisme. Par changement de base plat (Cohomologie des schémas, Lemme
 \ref{coherent-lemma-flat-base-change-cohomology}), nous pouvons remplacer  $S$  par  $\Spec(\mathcal{O}_{S, s})$  et
 $X$  par  $\Spec(\mathcal{O}_{S, s}) \times_S X$. Notons  $\mathfrak m \subset \mathcal{O}_S$  le faisceau d'idéaux de
 $s$. Alors nous constatons que
$$
\Hom_{\mathcal{O}_{S, s}}(\mathfrak m_s, (f_*\mathcal{F})_s) =
\Hom_{\mathcal{O}_S}(\mathfrak m, f_*\mathcal{F}) =
\Hom_{\mathcal{O}_X}(f^*\mathfrak m, \mathcal{F})
$$
la première égalité car  $S$  est local avec le point fermé  $s$  et
la deuxième égalité par adjonction pour  $f^*, f_*$  sur des modules
quasi-cohérents. Cependant, puisque  $s \not \in f(X)$  nous voyons que  $f^*\mathfrak m = \mathcal{O}_X$. En
revenant sur les arguments, nous obtenons l'égalité désirée.

\medskip\noindent
Pour la preuve de (4), (5) et (6), nous utilisons le changement de base plat
(Cohomologie des schémas, Lemme  \ref{coherent-lemma-flat-base-change-cohomology}) pour réduire au cas où  $S$ 
est le spectre de  $\mathcal{O}_{S, s}$. Alors un non-diviseur de zéro  $a \in \mathcal{O}_{S, s}$  détermine une
suite exacte
$$
0 \to \mathcal{O}_S \xrightarrow{a} \mathcal{O}_S \to
\mathcal{O}_S/a \mathcal{O}_S \to 0
$$
Puisque  $\mathcal{F}$  est plat sur  $S$, nous obtenons une suite exacte
$$
0 \to \mathcal{F} \xrightarrow{a} \mathcal{F} \to
\mathcal{F}/a\mathcal{F} \to 0
$$
En prenant l'image directe, nous obtenons une suite exacte
$$
0 \to f_*\mathcal{F} \xrightarrow{a} f_*\mathcal{F} \to
f_*(\mathcal{F}/a\mathcal{F})
$$
Cela prouve (4) et montre que  $f_*\mathcal{F}/ af_*\mathcal{F} \subset f_*(\mathcal{F}/a\mathcal{F})$. Si  $b$  est un diviseur de
zéro non nul sur  $\mathcal{O}_{S, s}/a\mathcal{O}_{S, s}$, alors le même argument exact montre que
 $b : \mathcal{F}/a\mathcal{F} \to \mathcal{F}/a\mathcal{F}$  est injectif. En prenant l'image directe, nous concluons que
$$
b : f_*(\mathcal{F}/a\mathcal{F}) \to f_*(\mathcal{F}/a\mathcal{F})
$$
est injectif et donc aussi  $b : f_*\mathcal{F}/ af_*\mathcal{F} \to f_*\mathcal{F}/ af_*\mathcal{F}$  est injectif. Cela prouve (5). La partie
(6) découle de (4) et (5) et des définitions.
\end{proof}











\section{Assassin relatif}
\label{divisors-section-relative-assassin}

\noindent
Soit  $A \to B$  un morphisme d'anneaux. Soit  $N$  un
 $B$-module. Rappelons qu'un idéal premier  $\mathfrak q \subset B$  est dit appartenir
à l'assassin relatif de  $N$  sur  $B/A$  si  $\mathfrak q$  est un idéal
premier associé à  $N \otimes_A \kappa(\mathfrak p)$. Ici  $\mathfrak p = A \cap \mathfrak q$. Voir Algèbre, Définition
 \ref{algebra-definition-relative-assassin}. Voici la définition de l'assassin relatif pour les faisceaux
quasi-cohérents sur un morphisme de schémas.

\begin{definition}
\label{divisors-definition-relative-assassin}
Soit $f : X \to S$ un morphisme de schémas. Soit $\mathcal{F}$ un
$\mathcal{O}_X$-module quasi-cohérent. L'{\it assassin relatif
de $\mathcal{F}$ dans $X$ relativement à $S$} est l'ensemble
$$
\text{Ass}_{X/S}(\mathcal{F}) =
\bigcup\nolimits_{s \in S} \text{Ass}_{X_s}(\mathcal{F}_s)
$$
où  $\mathcal{F}_s = (X_s \to X)^*\mathcal{F}$  est la restriction de  $\mathcal{F}$  à la fibre de  $f$  en
 $s$.
\end{definition}

\noindent
Cette notion s'emploie surtout lorsque les fibres de  $f$  sont localement
noethériennes et que  $\mathcal{F}$  est de type fini. Dans le cas général, il
convient plutôt d'utiliser l'assassin faible relatif (défini dans la section
suivante). Relions la notion schématique à la notion algébrique sur les ouverts
affines; notons que cette correspondance
ne fonctionne parfaitement que pour les morphismes de schémas dont les fibres
sont localement noethériennes.

\begin{lemma}
\label{divisors-lemma-relative-assassin-affine-open}
Soit  $f : X \to S$  un morphisme de schémas. Soit  $\mathcal{F}$  un faisceau
quasi-cohérent sur  $X$. Soit  $U \subset X$  et  $V \subset S$  des ouverts
affines avec  $f(U) \subset V$. Écrivons  $U = \Spec(A)$,  $V = \Spec(R)$, et posons
 $M = \Gamma(U, \mathcal{F})$. Soit  $x \in U$, et soit  $\mathfrak p \subset A$  l'idéal premier
correspondant. Alors
$$
\mathfrak p \in \text{Ass}_{A/R}(M) \Rightarrow
x \in \text{Ass}_{X/S}(\mathcal{F})
$$
Si toutes les fibres  $X_s$  de  $f$  sont localement noethériennes,
alors  $\mathfrak p \in \text{Ass}_{A/R}(M) \Leftrightarrow
x \in \text{Ass}_{X/S}(\mathcal{F})$  pour toutes les paires  $(\mathfrak p, x)$  comme ci-dessus.
\end{lemma}

\begin{proof}
L'ensemble  $\text{Ass}_{A/R}(M)$  est défini en Algèbre, Définition  \ref{algebra-definition-relative-assassin}. Fixons
une paire  $(\mathfrak p, x)$. Soit  $s = f(x)$. Soit  $\mathfrak r \subset R$  l'idéal premier situé sous
 $\mathfrak p$, c'est-à-dire l'idéal premier correspondant à  $s$. Soit
 $\mathfrak p' \subset A \otimes_R \kappa(\mathfrak r)$  l'idéal premier dont l'image réciproque est  $\mathfrak p$, c'est-à-dire l'idéal
premier correspondant à  $x$  considéré comme un point de sa fibre
 $X_s$. Alors  $\mathfrak p \in \text{Ass}_{A/R}(M)$  si et seulement si  $\mathfrak p'$  est un idéal premier
associé à  $M \otimes_R \kappa(\mathfrak r)$, voir Algèbre, Lemme  \ref{algebra-lemma-compare-relative-assassins}. Notons que l'anneau
 $A \otimes_R \kappa(\mathfrak r)$  correspond à  $U_s$  et que le module  $M \otimes_R \kappa(\mathfrak r)$  correspond au
faisceau quasi-cohérent  $\mathcal{F}_s|_{U_s}$. Ainsi  $x$  est un point associé de
 $\mathcal{F}_s$  par le Lemme  \ref{divisors-lemma-associated-affine-open}. L'implication inverse est valable si
 $\mathfrak p'$  est de type fini, ce qui permet de voir que la dernière phrase est
vraie.
\end{proof}

\begin{lemma}
\label{divisors-lemma-base-change-relative-assassin}
Soit  $f : X \to S$  un morphisme de schémas. Soit  $\mathcal{F}$  un
 $\mathcal{O}_X$-module quasi-cohérent. Soit  $g : S' \to S$  un morphisme de schémas.
Considérons le diagramme de changement de base
$$
\xymatrix{
X' \ar[d] \ar[r]_{g'} & X \ar[d] \\
S' \ar[r]^g & S
}
$$
et posons  $\mathcal{F}' = (g')^*\mathcal{F}$. Soit  $x' \in X'$  un point avec les images  $x \in X$,
 $s' \in S'$  et  $s \in S$. Supposons  $f$  localement de type fini.
Alors  $x' \in \text{Ass}_{X'/S'}(\mathcal{F}')$  si et seulement si  $x \in \text{Ass}_{X/S}(\mathcal{F})$  et  $x'$  correspond à un
point générique d'une composante irréductible de  $\Spec(\kappa(s') \otimes_{\kappa(s)} \kappa(x))$.
\end{lemma}

\begin{proof}
Considérons le morphisme  $X'_{s'} \to X_s$  des fibres. Comme  $X_{s'} = X_s \times_{\Spec(\kappa(s))} \Spec(\kappa(s'))$ , il s'agit
d'un morphisme plat. De plus,  $\mathcal{F}'_{s'}$  est l'image réciproque de  $\mathcal{F}_s$ 
via ce morphisme. Comme  $X_s$  est localement de type fini sur le schéma
noethérien  $\Spec(\kappa(s))$ , nous avons que  $X_s$  est localement noethérien,
voir Morphismes, Lemme  \ref{morphisms-lemma-finite-type-noetherian}. Ainsi, nous pouvons appliquer le Lemme
 \ref{divisors-lemma-bourbaki}  et nous voyons que
$$
\text{Ass}_{X'_{s'}}(\mathcal{F}'_{s'}) =
\bigcup\nolimits_{x \in \text{Ass}(\mathcal{F}_s)} \text{Ass}((X'_{s'})_x).
$$
Ainsi, pour prouver le lemme, il suffit de montrer que les points associés de la
fibre  $(X'_{s'})_x$  du morphisme  $X'_{s'} \to X_s$  au-dessus de  $x$  sont ses
points génériques. Notons que  $(X'_{s'})_x = \Spec(\kappa(s') \otimes_{\kappa(s)} \kappa(x))$  en tant que schémas. D'après Algèbre,
Lemme  \ref{algebra-lemma-tensor-fields-CM} , l'anneau  $\kappa(s') \otimes_{\kappa(s)} \kappa(x)$  est un anneau de Cohen-Macaulay
noethérien. Par conséquent, ses idéaux premiers associés sont ses idéaux premiers
minimaux, voir Algèbre, Proposition  \ref{algebra-proposition-minimal-primes-associated-primes}  (les idéaux premiers minimaux
sont associés) et Algèbre, Lemme  \ref{algebra-lemma-criterion-no-embedded-primes}  (pas d'idéaux premiers immergés).
\end{proof}

\begin{remark}
\label{divisors-remark-base-change-relative-assassin}
Avec les notations et hypothèses comme dans le Lemme  \ref{divisors-lemma-base-change-relative-assassin} , nous voyons
qu'on a toujours  $(g')^{-1}(\text{Ass}_{X/S}(\mathcal{F})) \supset
\text{Ass}_{X'/S'}(\mathcal{F}')$. Si le morphisme  $S' \to S$  est
quasi-fini localement, alors nous avons en fait
$$
(g')^{-1}(\text{Ass}_{X/S}(\mathcal{F}))
=
\text{Ass}_{X'/S'}(\mathcal{F}')
$$
car dans ce cas les extensions de corps  $\kappa(s')/\kappa(s)$  sont toujours finies. En
fait, cela est vrai de manière plus générale pour tout morphisme  $g : S' \to S$  tel
que toutes les extensions de corps  $\kappa(s')/\kappa(s)$  soient algébriques, car dans ce
cas tous les idéaux premiers de  $\kappa(s') \otimes_{\kappa(s)} \kappa(x)$  sont des idéaux maximaux (et
minimaux), voir Algèbre, Lemme  \ref{algebra-lemma-integral-over-field}.
\end{remark}




\section{Assassin faible relatif}
\label{divisors-section-relative-weak-assassin}

\begin{definition}
\label{divisors-definition-relative-weak-assassin}
Soit  $f : X \to S$  un morphisme de schémas. Soit  $\mathcal{F}$  un
 $\mathcal{O}_X$-module quasi-cohérent. L'assassin faible relatif
 {\it  de  $\mathcal{F}$  dans  $X$  sur
 $S$}  est l'ensemble
$$
\text{WeakAss}_{X/S}(\mathcal{F}) =
\bigcup\nolimits_{s \in S} \text{WeakAss}(\mathcal{F}_s)
$$
où  $\mathcal{F}_s = (X_s \to X)^*\mathcal{F}$  est la restriction de  $\mathcal{F}$  à la fibre de  $f$  en
 $s$.
\end{definition}

\begin{lemma}
\label{divisors-lemma-relative-weak-assassin-assassin-finite-type}
Soit  $f : X \to S$  un morphisme de schémas qui est localement de type fini. Soit
 $\mathcal{F}$  un  $\mathcal{O}_X$-module quasi-cohérent. Alors  $\text{WeakAss}_{X/S}(\mathcal{F}) = \text{Ass}_{X/S}(\mathcal{F})$.
\end{lemma}

\begin{proof}
C'est vrai parce que les fibres de  $f$  sont des schémas localement
noethériens, et les points associés et faiblement associés coïncident sur les
schémas localement noethériens, voir le Lemme  \ref{divisors-lemma-ass-weakly-ass}.
\end{proof}

\begin{lemma}
\label{divisors-lemma-relative-weak-assassin-finite}
Soit  $f : X \to S$  un morphisme de schémas. Soit  $i : Z \to X$  un morphisme fini.
Soit  $\mathcal{F}$  un  $\mathcal{O}_Z$-module quasi-cohérent. Alors  $\text{WeakAss}_{X/S}(i_*\mathcal{F}) =
i(\text{WeakAss}_{Z/S}(\mathcal{F}))$.
\end{lemma}

\begin{proof}
Soit  $i_s : Z_s \to X_s$  le morphisme induit entre les fibres. Alors  $(i_*\mathcal{F})_s = i_{s, *}(\mathcal{F}_s)$  d'après
Cohomologie des schémas, Lemme  \ref{coherent-lemma-affine-base-change}  et le fait que  $i$  est
affine. Par conséquent, nous pouvons appliquer le Lemme  \ref{divisors-lemma-weakly-associated-finite}  pour
conclure.
\end{proof}





\section{Idéaux de Fitting}
\label{divisors-section-fitting-ideals}

\noindent
Cette section poursuit la discussion de Compléments d'algèbre,
section  \ref{more-algebra-section-fitting-ideals}. Soit  $S$  un schéma. Soit  $\mathcal{F}$  un
 $\mathcal{O}_S$-module quasi-cohérent de type fini. Dans cette situation, nous
pouvons construire les idéaux de Fitting
$$
0 = \text{Fit}_{-1}(\mathcal{F}) \subset \text{Fit}_0(\mathcal{F}) \subset
\text{Fit}_1(\mathcal{F}) \subset \ldots \subset \mathcal{O}_S
$$
comme la suite des idéaux quasi-cohérents caractérisés par la propriété suivante:
pour tout ouvert affine  $U = \Spec(A)$  de  $S$ , si  $\mathcal{F}|_U$  correspond
au  $A$-module  $M$, alors  $\text{Fit}_i(\mathcal{F})|_U$  correspond à l'idéal
 $\text{Fit}_i(M) \subset A$. Cela est bien défini et constitue un faisceau quasi-cohérent
d'idéaux parce que si  $f \in A$, alors le  $i$-ième idéal de Fitting de
 $M_f$  sur  $A_f$  est égal à  $\text{Fit}_i(M) A_f$  selon Compléments d'algèbre, Lemme
 \ref{more-algebra-lemma-fitting-ideal-basics}.

\medskip\noindent
Alternativement, nous pouvons construire les idéaux de Fitting en termes de
présentations locales de  $\mathcal{F}$. À savoir, si  $U \subset X$  est ouvert, et
$$
\bigoplus\nolimits_{i \in I} \mathcal{O}_U \to
\mathcal{O}_U^{\oplus n} \to \mathcal{F}|_U \to 0
$$
est une présentation de  $\mathcal{F}$  sur  $U$, alors  $\text{Fit}_r(\mathcal{F})|_U$  est
généré par les  $(n - r) \times (n - r)$ mineurs de la matrice définissant la première flèche
de la présentation. Cela est compatible avec la construction ci-dessus car c'est
ainsi que l'idéal de Fitting d'un module sur un anneau est réellement défini.
Quelques détails omis.

\begin{lemma}
\label{divisors-lemma-base-change-fitting-ideal}
Soit  $f : T \to S$  un morphisme de schémas. Soit  $\mathcal{F}$  un
 $\mathcal{O}_S$-module quasi-cohérent de type fini. Alors  $f^{-1}\text{Fit}_i(\mathcal{F}) \cdot \mathcal{O}_T =
\text{Fit}_i(f^*\mathcal{F})$.
\end{lemma}

\begin{proof}
Découle immédiatement de Compléments d'algèbre, Lemme  \ref{more-algebra-lemma-fitting-ideal-basics}  partie (3).
\end{proof}

\begin{lemma}
\label{divisors-lemma-fitting-ideal-of-finitely-presented}
Soit  $S$  un schéma. Soit  $\mathcal{F}$  un  $\mathcal{O}_S$-module de
présentation finie. Alors  $\text{Fit}_r(\mathcal{F})$  est un idéal quasi-cohérent de type fini.
\end{lemma}

\begin{proof}
Découle immédiatement de Compléments d'algèbre, Lemme  \ref{more-algebra-lemma-fitting-ideal-basics}  partie (4).
\end{proof}

\begin{lemma}
\label{divisors-lemma-on-subscheme-cut-out-by-Fit-0}
Soit  $S$  un schéma. Soit  $\mathcal{F}$  un  $\mathcal{O}_S$-module
quasi-cohérent de type fini. Soit  $Z_0 \subset S$  le sous-schéma fermé défini par
 $\text{Fit}_0(\mathcal{F})$. Soit  $Z \subset S$  le support schématique de  $\mathcal{F}$. Alors
\begin{enumerate}
\item $Z \subset Z_0 \subset S$  en tant que sous-schémas fermés,
\item $Z = Z_0 = \text{Supp}(\mathcal{F})$  en tant que sous-ensembles fermés,
\item il existe un  $\mathcal{O}_{Z_0}$-module quasi-cohérent de type fini  $\mathcal{G}_0$  avec
$$
(Z_0 \to X)_*\mathcal{G}_0 = \mathcal{F}.
$$
\end{enumerate}
\end{lemma}

\begin{proof}
Rappelons que  $Z$  est localement défini par l'annulateur de  $\mathcal{F}$,
voir Morphismes, Définition  \ref{morphisms-definition-scheme-theoretic-support}  (qui utilise Morphismes, Lemme
 \ref{morphisms-lemma-scheme-theoretic-support}  pour définir  $Z$). Par conséquent, nous voyons que
 $Z \subset Z_0$  schématiquement par Compléments d'algèbre, Lemme  \ref{more-algebra-lemma-fitting-ideal-basics}  partie (6).
D'autre part, nous avons  $Z = \text{Supp}(\mathcal{F})$  ensemblistement par Morphismes, Lemme
 \ref{morphisms-lemma-scheme-theoretic-support}  et nous avons  $Z_0 = Z$  ensemblistement par Compléments d'algèbre, Lemme
 \ref{more-algebra-lemma-fitting-ideal-basics}  partie (7). Enfin, pour obtenir  $\mathcal{G}_0$  comme dans la partie
(3), nous pouvons soit utiliser le fait que nous avons  $\mathcal{G}$  sur
 $Z$  comme dans Morphismes, Lemme  \ref{morphisms-lemma-scheme-theoretic-support}  et définir  $\mathcal{G}_0 = (Z \to Z_0)_*\mathcal{G}$ ,
soit utiliser Morphismes, Lemme  \ref{morphisms-lemma-i-star-equivalence}  et le fait que  $\text{Fit}_0(\mathcal{F})$  annule
 $\mathcal{F}$  selon Compléments d'algèbre, Lemme  \ref{more-algebra-lemma-fitting-ideal-basics}  partie (6).
\end{proof}

\begin{lemma}
\label{divisors-lemma-fitting-ideal-generate-locally}
Soit  $S$  un schéma. Soit  $\mathcal{F}$  un  $\mathcal{O}_S$-module
quasi-cohérent de type fini. Soit  $s \in S$. Alors  $\mathcal{F}$  peut être
généré par  $r$  éléments dans un voisinage de  $s$  si et seulement
si  $\text{Fit}_r(\mathcal{F})_s = \mathcal{O}_{S, s}$.
\end{lemma}

\begin{proof}
Découle immédiatement de Compléments d'algèbre, Lemme  \ref{more-algebra-lemma-fitting-ideal-generate-locally}.
\end{proof}

\begin{lemma}
\label{divisors-lemma-fitting-ideal-finite-locally-free}
Soit  $S$  un schéma. Soit  $\mathcal{F}$  un  $\mathcal{O}_S$-module
quasi-cohérent de type fini. Soit  $r \geq 0$. Les éléments suivants sont
équivalents
\begin{enumerate}
\item $\mathcal{F}$  est localement libre de type fini et de rang  $r$
\item $\text{Fit}_{r - 1}(\mathcal{F}) = 0$  et  $\text{Fit}_r(\mathcal{F}) = \mathcal{O}_S$, et
\item $\text{Fit}_k(\mathcal{F}) = 0$  pour  $k < r$  et  $\text{Fit}_k(\mathcal{F}) = \mathcal{O}_S$  pour  $k \geq r$.
\end{enumerate}
\end{lemma}

\begin{proof}
Découle immédiatement de Compléments d'algèbre, Lemme  \ref{more-algebra-lemma-fitting-ideal-finite-locally-free}.
\end{proof}

\begin{lemma}
\label{divisors-lemma-locally-free-rank-r-pullback}
Soit  $S$  un schéma. Soit  $\mathcal{F}$  un  $\mathcal{O}_S$-module
quasi-cohérent de type fini. Les sous-schémas fermés
$$
S = Z_{-1} \supset Z_0 \supset Z_1 \supset Z_2 \ldots
$$
définis par les idéaux de Fitting de  $\mathcal{F}$  ont les propriétés suivantes
\begin{enumerate}
\item L'intersection  $\bigcap Z_r$  est vide.
\item Le foncteur  $(\Sch/S)^{opp} \to \textit{Ensembles}$  défini par la règle
$$
T \longmapsto
\left\{
\begin{matrix}
\{*\} & \text{if }\mathcal{F}_T\text{ is locally generated by }
\leq r\text{ sections} \\
\emptyset & \text{otherwise}
\end{matrix}
\right.
$$
est représentable par le sous-schéma ouvert  $S \setminus Z_r$.
\item Le foncteur  $F_r : (\Sch/S)^{opp} \to \textit{Ensembles}$  défini par la règle
$$
T \longmapsto
\left\{
\begin{matrix}
\{*\} & \text{if }\mathcal{F}_T\text{ locally free rank }r\\
\emptyset & \text{otherwise}
\end{matrix}
\right.
$$
est représentable par le sous-schéma localement fermé  $Z_{r - 1} \setminus Z_r$  de
 $S$.
\end{enumerate}
Si  $\mathcal{F}$  est de présentation finie, alors  $Z_r \to S$,  $S \setminus Z_r \to S$ et
 $Z_{r - 1} \setminus Z_r \to S$  sont de présentation finie.
\end{lemma}

\begin{proof}
La partie (1) est vraie car sur chaque ouvert affine  $U$ , il existe un
entier  $n$  tel que  $\text{Fit}_n(\mathcal{F})|_U = \mathcal{O}_U$. À savoir, nous pouvons prendre
 $n$  comme étant le nombre de générateurs de  $\mathcal{F}$  sur
 $U$, voir Compléments d'algèbre, section  \ref{more-algebra-section-fitting-ideals}.

\medskip\noindent
Pour tout morphisme  $g : T \to S$ , nous voyons d'après les lemmes  \ref{divisors-lemma-base-change-fitting-ideal}  et
 \ref{divisors-lemma-fitting-ideal-generate-locally}  que  $\mathcal{F}_T$  est localement engendré par un nombre  $\leq r$  de sections si
et seulement si  $\text{Fit}_r(\mathcal{F}) \cdot \mathcal{O}_T = \mathcal{O}_T$. Cela prouve (2).

\medskip\noindent
Pour tout morphisme  $g : T \to S$ , nous voyons d'après les lemmes  \ref{divisors-lemma-base-change-fitting-ideal}  et
 \ref{divisors-lemma-fitting-ideal-finite-locally-free}  que  $\mathcal{F}_T$  est libre de rang  $r$  si et seulement si
 $\text{Fit}_r(\mathcal{F}) \cdot \mathcal{O}_T = \mathcal{O}_T$  et  $\text{Fit}_{r - 1}(\mathcal{F}) \cdot \mathcal{O}_T = 0$. Cela prouve (3).

\medskip\noindent
Supposons que  $\mathcal{F}$  soit de présentation finie. Alors chacun des morphismes
 $Z_r \to S$  est de présentation finie puisque  $\text{Fit}_r(\mathcal{F})$  est de type fini
(Lemme  \ref{divisors-lemma-fitting-ideal-of-finitely-presented}  et Morphismes, Lemme  \ref{morphisms-lemma-closed-immersion-finite-presentation}). Cela implique que
 $Z_{r - 1} \setminus Z_r$  est un ouvert rétrocompact dans  $Z_r$  (Propriétés, Lemme
 \ref{properties-lemma-quasi-coherent-finite-type-ideals}) et donc le morphisme  $Z_{r - 1} \setminus Z_r \to Z_r$  est également de présentation
finie.
\end{proof}

\noindent
Malgré le Lemme  \ref{divisors-lemma-locally-free-rank-r-pullback}, le lemme suivant n'est plus vrai si
 $\mathcal{F}$  est un module quasi-cohérent de type fini. À savoir, la
stratification existe toujours mais il n'est pas vrai qu'elle représente le
foncteur  $F_{flat}$  en général.

\begin{lemma}
\label{divisors-lemma-finite-presentation-module}
Soit  $S$  un schéma. Soit  $\mathcal{F}$  un  $\mathcal{O}_S$-module de
présentation finie. Soit  $S = Z_{-1} \supset Z_0 \supset Z_1 \supset \ldots$  comme dans le Lemme  \ref{divisors-lemma-locally-free-rank-r-pullback}. Posons
 $S_r = Z_{r - 1} \setminus Z_r$. Alors  $S' = \coprod_{r \geq 0} S_r$  représente le foncteur
$$
F_{flat} : \Sch/S \longrightarrow \textit{Ensembles},\quad\quad
T \longmapsto
\left\{
\begin{matrix}
\{*\} & \text{if }\mathcal{F}_T\text{ flat over }T\\
\emptyset & \text{otherwise}
\end{matrix}
\right.
$$
De plus,  $\mathcal{F}|_{S_r}$  est localement libre de rang  $r$  et les morphismes
 $S_r \to S$  et  $S' \to S$  sont de présentation finie.
\end{lemma}

\begin{proof}
Supposons que  $g : T \to S$  soit un morphisme de schémas tel que l'image réciproque
 $\mathcal{F}_T = g^*\mathcal{F}$  soit plat. Alors  $\mathcal{F}_T$  est un  $\mathcal{O}_T$-module plat de
présentation finie. Par conséquent,  $\mathcal{F}_T$  est localement libre de type fini,
voir Propriétés, Lemme  \ref{properties-lemma-finite-locally-free}. Ainsi  $T = \coprod_{r \geq 0} T_r$, où  $\mathcal{F}_T|_{T_r}$  est
localement libre de rang  $r$. Cela implique que
$$
F_{flat} = \coprod\nolimits_{r \geq 0} F_r
$$
dans la catégorie des faisceaux de Zariski sur  $\Sch/S$  où  $F_r$  est
comme dans le Lemme  \ref{divisors-lemma-locally-free-rank-r-pullback}. Il s'ensuit que  $F_{flat}$  est représenté par
 $\coprod_{r \geq 0} (Z_{r - 1} \setminus Z_r)$  où  $Z_r$  est comme dans le Lemme  \ref{divisors-lemma-locally-free-rank-r-pullback}. Les autres
énoncés découlent également du lemme.
\end{proof}

\begin{example}
\label{divisors-example-not-fp-MB}
Soit  $R = \prod_{n \in \mathbf{N}} \mathbf{F}_2$. Soit  $I \subset R$  l'idéal des éléments  $a = (a_n)_{n \in \mathbf{N}}$  dont
presque toutes les composantes sont nulles. Soit  $\mathfrak m$  un idéal maximal
contenant  $I$. Alors  $M = R/\mathfrak m$  est un  $R$-module fini et
plat, car  $R$  est absolument plat (Compléments d'algèbre, Lemme  \ref{more-algebra-lemma-product-fields-absolutely-flat}).
Posons  $S = \Spec(R)$  et  $\mathcal{F} = \widetilde{M}$. Les sous-schémas fermés du Lemme  \ref{divisors-lemma-locally-free-rank-r-pullback} 
sont  $S = Z_{-1}$,  $Z_0 = \Spec(R/\mathfrak m)$ et  $Z_i = \emptyset$  pour  $i > 0$. Mais
 $\text{id} : S \to S$  ne se factorise pas par  $(S \setminus Z_0) \amalg Z_0$  car  $\mathfrak m$  est un point
non isolé de  $S$. Ainsi, le Lemme  \ref{divisors-lemma-finite-presentation-module}  ne tient pas pour les
modules de type fini.
\end{example}







\section{Le lieu singulier d'un morphisme}
\label{divisors-section-singular-locus-morphism}

\noindent
Soit  $f : X \to S$  un morphisme de type fini de schémas. L'ensemble  $U$ 
des points où  $f$  est lisse est un ouvert de  $X$  (d'après
Morphismes, Définition  \ref{morphisms-definition-smooth}). Dans de nombreuses situations, il est utile
d'avoir un sous-schéma fermé canonique  $\text{Sing}(f) \subset X$  dont le complémentaire est
 $U$  et dont la formation commute avec tout changement de base.

\medskip\noindent
Si  $f$  est de présentation finie, un choix possible consiste à considérer le
sous-schéma fermé  $Z$  défini par les fonctions qui sont « strictement standard »
localement sur des ouverts affines, au sens de Lissage des morphismes d'anneaux, Définition
 \ref{smoothing-definition-strictly-standard}. Il découle de Lissage des morphismes d'anneaux, Lemme  \ref{smoothing-lemma-strictly-standard-base-change}  que si
 $f' : X' \to S'$  est le changement de base de  $f$  par un morphisme
 $S' \to S$, alors  $Z' \subset S' \times_S Z$  où  $Z'$  est le sous-schéma fermé de
 $X'$  défini par les fonctions qui sont strictement standard localement sur des
ouverts affines. Cependant, l'égalité n'est pas claire. La notion d'élément strictement
standard a été utile dans le chapitre sur le théorème de Popescu. Le sous-schéma
fermé défini par ces éléments n'est (autant que nous le sachions) pas utilisé dans
la littérature\footnote{Si  $f$  est un morphisme
localement d'intersection complète (Compléments sur les morphismes, Définition  \ref{more-morphisms-definition-lci}),
alors le sous-schéma fermé défini par les éléments localement strictement standard
est bien celui qu'il faut considérer.}.

\medskip\noindent
Si  $f$  est plat, de présentation finie, et que les fibres de  $f$ 
sont toutes équidimensionnelles de dimension  $d$, alors l'idéal de
Fitting d'indice  $d$ de  $\Omega_{X/S}$  fournit un bon sous-schéma
fermé. Pour tout morphisme de type fini, les sous-schémas fermés de  $X$ 
définis par les idéaux de Fitting de  $\Omega_{X/S}$  définissent une stratification
de  $X$  en termes de rang de  $\Omega_{X/S}$  dont la formation commute avec
le changement de base. Cela peut être utile; c'est lié aux dimensions d'immersion
des fibres, voir Variétés, section  \ref{varieties-section-embedding-dimension}.

\begin{lemma}
\label{divisors-lemma-base-change-and-fitting-ideal-omega}
Soit  $f : X \to S$  un morphisme de schémas qui est localement de type fini. Soit
 $X = Z_{-1} \supset Z_0 \supset Z_1 \supset \ldots$  les sous-schémas fermés définis par les idéaux de Fitting de
 $\Omega_{X/S}$. Alors la formation de  $Z_i$  commute avec tout changement de
base.
\end{lemma}

\begin{proof}
Remarquez que  $\Omega_{X/S}$  est un  $\mathcal{O}_X$-module quasi-cohérent de type fini
(Morphismes, Lemme  \ref{morphisms-lemma-finite-type-differentials}) donc les idéaux de Fitting sont définis. Si
 $f' : X' \to S'$  est le changement de base de  $f$  par  $g : S' \to S$, alors
 $\Omega_{X'/S'} = (g')^*\Omega_{X/S}$  où  $g' : X' \to X$  est la projection (Morphismes, Lemme  \ref{morphisms-lemma-base-change-differentials}).
D'où  $(g')^{-1}\text{Fit}_i(\Omega_{X/S}) \cdot \mathcal{O}_{X'} =
\text{Fit}_i(\Omega_{X'/S'})$. Cela signifie que
$$
Z'_i = (g')^{-1}(Z_i) = Z_i \times_X X'
$$
schématiquement et c'est le sens de l'énoncé du lemme.
\end{proof}

\noindent
Le  $0$-ième idéal de Fitting de  $\Omega$  définit le « lieu de
ramification » du morphisme.

\begin{lemma}
\label{divisors-lemma-zero-fitting-ideal-omega-unramified}
Soit  $f : X \to S$  un morphisme de schémas qui est localement de type fini. Le
sous-schéma fermé  $Z \subset X$  défini par le  $0$-ième idéal de Fitting de
 $\Omega_{X/S}$  est exactement l'ensemble des points où  $f$  n'est pas non ramifié.
\end{lemma}

\begin{proof}
D'après le Lemme  \ref{divisors-lemma-on-subscheme-cut-out-by-Fit-0} , le complément de  $Z$  est exactement le
lieu où  $\Omega_{X/S}$  est nul. Il s'agit exactement de l'ensemble des points où
 $f$  est non ramifié par les Morphismes, Lemme  \ref{morphisms-lemma-unramified-omega-zero}.
\end{proof}

\begin{lemma}
\label{divisors-lemma-d-fitting-ideal-omega-smooth}
Soit  $f : X \to S$  un morphisme de schémas. Soit  $d \geq 0$  un entier. Supposons
\begin{enumerate}
\item que $f$  est plat,
\item que $f$  est localement de présentation finie, et
\item que toute fibre non vide de  $f$  est équidimensionnelle de dimension
 $d$.
\end{enumerate}
Soit  $Z \subset X$  le sous-schéma fermé défini par le $d$-ième idéal de Fitting
de  $\Omega_{X/S}$. Alors  $Z$  est exactement l'ensemble des
points où  $f$  n'est pas lisse.
\end{lemma}

\begin{proof}
D'après le Lemme  \ref{divisors-lemma-locally-free-rank-r-pullback} , le complément de  $Z$  est exactement le
lieu où  $\Omega_{X/S}$  peut être engendré par au plus  $d$  éléments. Par
conséquent, le lemme découle de Morphismes, Lemme  \ref{morphisms-lemma-smooth-at-point}.
\end{proof}







\section{Modules sans torsion}
\label{divisors-section-torsion-free}

\noindent
Cette section est l'analogue de Compléments d'algèbre, section  \ref{more-algebra-section-torsion-flat}  pour les
modules quasi-cohérents.

\begin{lemma}
\label{divisors-lemma-torsion-sections}
Soit  $X$  un schéma intègre avec point générique  $\eta$. Soit
 $\mathcal{F}$  un  $\mathcal{O}_X$-module quasi-cohérent. Soit  $U \subset X$  ouvert non
vide et  $s \in \mathcal{F}(U)$. Les énoncés suivants sont équivalents
\begin{enumerate}
\item pour certains  $x \in U$  l'image de  $s$  dans  $\mathcal{F}_x$  est de
torsion,
\item pour tous  $x \in U$  l'image de  $s$  dans  $\mathcal{F}_x$  est de torsion,
\item l'image de  $s$  dans  $\mathcal{F}_\eta$  est nulle,
\item l'image de  $s$  dans  $j_*\mathcal{F}_\eta$  est nulle, où  $j : \eta \to X$  est le
morphisme d'inclusion.
\end{enumerate}
\end{lemma}

\begin{proof}
Démonstration omise.
\end{proof}

\begin{definition}
\label{divisors-definition-torsion}
Soit  $X$  un schéma intègre. Soit  $\mathcal{F}$  un  $\mathcal{O}_X$-module
quasi-cohérent.
\begin{enumerate}
\item On dit qu'une section locale de  $\mathcal{F}$  est  {\it 
de torsion}  si elle satisfait les conditions équivalentes du Lemme
 \ref{divisors-lemma-torsion-sections}.
\item On dit que  $\mathcal{F}$  est  {\it  sans torsion}  si
chaque section de torsion de  $\mathcal{F}$  est  $0$.
\end{enumerate}
\end{definition}

\noindent
Voici le lemme obligatoire comparant cela à la notion algébrique habituelle.

\begin{lemma}
\label{divisors-lemma-check-torsion-on-affines}
Soit  $X$  un schéma intègre. Soit  $\mathcal{F}$  un  $\mathcal{O}_X$-module
quasi-cohérent. Les conditions suivantes sont équivalentes
\begin{enumerate}
\item $\mathcal{F}$  est sans torsion,
\item pour  $U \subset X$  ouvert affine,  $\mathcal{F}(U)$  est un  $\mathcal{O}(U)$-module sans
torsion.
\end{enumerate}
\end{lemma}

\begin{proof}
Démonstration omise.
\end{proof}

\begin{lemma}
\label{divisors-lemma-torsion}
Soit  $X$  un schéma intègre. Soit  $\mathcal{F}$  un  $\mathcal{O}_X$-module
quasi-cohérent. Les sections de torsion de  $\mathcal{F}$  forment un
 $\mathcal{O}_X$-sous-module quasi-cohérent  $\mathcal{F}_{tors} \subset \mathcal{F}$. Le module quotient
 $\mathcal{F}/\mathcal{F}_{tors}$  est sans torsion.
\end{lemma}

\begin{proof}
Démonstration omise. Voir Compléments d'algèbre, Lemme  \ref{more-algebra-lemma-torsion}  pour l'analogue algébrique.
\end{proof}

\begin{lemma}
\label{divisors-lemma-flat-torsion-free}
Soit  $X$  un schéma intègre. Tout  $\mathcal{O}_X$-module quasi-cohérent plat
est sans torsion.
\end{lemma}

\begin{proof}
Démonstration omise. Voir Compléments d'algèbre, Lemme  \ref{more-algebra-lemma-flat-torsion-free}.
\end{proof}

\begin{lemma}
\label{divisors-lemma-flat-pullback-torsion}
Soit  $f : X \to Y$  un morphisme plat de schémas intègres. Soit  $\mathcal{G}$  un
 $\mathcal{O}_Y$-module quasi-cohérent sans torsion. Alors  $f^*\mathcal{G}$  est un
 $\mathcal{O}_X$-module sans torsion.
\end{lemma}

\begin{proof}
Démonstration omise. Voir Compléments d'algèbre, Lemme  \ref{more-algebra-lemma-flat-pullback-torsion}  pour l'analogue algébrique.
\end{proof}

\begin{lemma}
\label{divisors-lemma-flat-over-integral-integral-fibre}
Soit  $f : X \to Y$  un morphisme plat de schémas. Si  $Y$  est intègre et la
fibre générique de  $f$  est intègre, alors  $X$  est intègre.
\end{lemma}

\begin{proof}
L'analogue algébrique est le suivant: soit  $A$  un domaine avec le corps
des fractions  $K$  et soit  $B$  une  $A$-algèbre plate
telle que  $B \otimes_A K$  soit un domaine. Alors  $B$  est un domaine. Ceci
est vrai parce que  $B$  est sans torsion selon Compléments d'algèbre, Lemme
 \ref{more-algebra-lemma-flat-torsion-free}  et donc  $B \subset B \otimes_A K$.
\end{proof}

\begin{lemma}
\label{divisors-lemma-check-torsion}
Soit  $X$  un schéma intègre. Soit  $\mathcal{F}$  un  $\mathcal{O}_X$-module
quasi-cohérent. Alors  $\mathcal{F}$  est sans torsion si et seulement si
 $\mathcal{F}_x$  est un  $\mathcal{O}_{X, x}$-module sans torsion pour tout  $x \in X$.
\end{lemma}

\begin{proof}
Démonstration omise. Voir Compléments d'algèbre, Lemme  \ref{more-algebra-lemma-check-torsion}.
\end{proof}

\begin{lemma}
\label{divisors-lemma-extension-torsion-free}
Soit  $X$  un schéma intègre. Soit  $0 \to \mathcal{F} \to \mathcal{F}' \to \mathcal{F}'' \to 0$  une suite exacte de
 $\mathcal{O}_X$-modules quasi-cohérents. Si  $\mathcal{F}$  et  $\mathcal{F}''$  sont sans
torsion, alors  $\mathcal{F}'$  est sans torsion.
\end{lemma}

\begin{proof}
Démonstration omise. Voir Compléments d'algèbre, Lemme  \ref{more-algebra-lemma-extension-torsion-free}  pour l'analogue algébrique.
\end{proof}

\begin{lemma}
\label{divisors-lemma-torsion-free-finite-noetherian-domain}
Soit  $X$  un schéma intègre localement noethérien avec le point générique
 $\eta$. Soit  $\mathcal{F}$  un  $\mathcal{O}_X$-module cohérent non nul. Les
énoncés suivants sont équivalents
\begin{enumerate}
\item $\mathcal{F}$  est sans torsion,
\item $\eta$  est le seul idéal premier associé de  $\mathcal{F}$,
\item $\eta$  est dans le support de  $\mathcal{F}$  et  $\mathcal{F}$  possède la
propriété  $(S_1)$, et
\item $\eta$  est dans le support de  $\mathcal{F}$  et  $\mathcal{F}$  n'a aucun idéal
idéal premier associé immergé.
\end{enumerate}
\end{lemma}

\begin{proof}
Ceci est une traduction de Compléments d'algèbre, Lemme  \ref{more-algebra-lemma-torsion-free-finite-noetherian-domain}  dans le langage
des schémas. Nous omettons la traduction.
\end{proof}

\begin{lemma}
\label{divisors-lemma-torsion-free-over-regular-dim-1}
Soit  $X$  un schéma régulier intègre de dimension  $\leq 1$. Soit
 $\mathcal{F}$  un  $\mathcal{O}_X$-module cohérent. Les énoncés suivants sont
équivalents
\begin{enumerate}
\item $\mathcal{F}$  est sans torsion,
\item $\mathcal{F}$  est localement libre de type fini.
\end{enumerate}
\end{lemma}

\begin{proof}
Il est clair qu'un module localement libre de type fini est sans torsion. Pour la
réciproque, nous montrerons que si  $\mathcal{F}$  est sans torsion, alors
 $\mathcal{F}_x$  est un  $\mathcal{O}_{X, x}$-module libre pour tout  $x \in X$. Cela suffit
d'après Algèbre, Lemme  \ref{algebra-lemma-finite-projective}  et le fait que  $\mathcal{F}$  est cohérent. Si
 $\dim(\mathcal{O}_{X, x}) = 0$, alors  $\mathcal{O}_{X, x}$  est un corps et l'énoncé est clair. Si
 $\dim(\mathcal{O}_{X, x}) = 1$, alors  $\mathcal{O}_{X, x}$  est un anneau de valuation discrète (Algèbre,
Lemme  \ref{algebra-lemma-characterize-dvr}) et  $\mathcal{F}_x$  est sans torsion. Par conséquent,  $\mathcal{F}_x$ 
est libre d'après Compléments d'algèbre, Lemme  \ref{more-algebra-lemma-dedekind-torsion-free-flat}.
\end{proof}

\begin{lemma}
\label{divisors-lemma-hom-into-torsion-free}
Soit  $X$  un schéma intègre. Soient  $\mathcal{F}$,  $\mathcal{G}$  des
 $\mathcal{O}_X$-modules quasi-cohérents. Si  $\mathcal{G}$  est sans torsion et
 $\mathcal{F}$  est de présentation finie, alors  $\SheafHom_{\mathcal{O}_X}(\mathcal{F}, \mathcal{G})$  est sans torsion.
\end{lemma}

\begin{proof}
L'énoncé a du sens car  $\SheafHom_{\mathcal{O}_X}(\mathcal{F}, \mathcal{G})$  est quasi-cohérent selon Schémas, section
 \ref{schemes-section-quasi-coherent}. Pour voir que l'énoncé est vrai, voir Compléments d'algèbre, Lemme
 \ref{more-algebra-lemma-hom-into-torsion-free}. Certains détails sont omis.
\end{proof}

\begin{lemma}
\label{divisors-lemma-isom-depth-2-torsion-free}
Soit  $X$  un schéma intègre localement noethérien. Soit  $\varphi : \mathcal{F} \to \mathcal{G}$  une
application de  $\mathcal{O}_X$-modules quasi-cohérents. Supposons que  $\mathcal{F}$ 
soit cohérent, que  $\mathcal{G}$  soit sans torsion, et que pour chaque
 $x \in X$ , l'une des situations suivantes se produise
\begin{enumerate}
\item $\mathcal{F}_x \to \mathcal{G}_x$  est un isomorphisme, ou
\item $\text{depth}(\mathcal{F}_x) \geq 2$.
\end{enumerate}
Alors  $\varphi$  est un isomorphisme.
\end{lemma}

\begin{proof}
Ceci est une traduction de Compléments d'algèbre, Lemme  \ref{more-algebra-lemma-isom-depth-2-torsion-free}  dans le langage des
schémas.
\end{proof}








\section{Rangs des modules}
\label{divisors-section-rank}

\noindent
Cette section est l'analogue de Compléments d'algèbre, section  \ref{more-algebra-section-rank}  pour les
modules quasi-cohérents. Rappelons que si  $X$  est un schéma intègre avec
le point générique  $\eta$, alors l'anneau local  $\mathcal{O}_{X, \eta}$  est égal au
corps résiduel  $\kappa(\eta)$, voir Morphismes, Lemme  \ref{morphisms-lemma-integral-scheme-rational-functions}.

\begin{definition}
\label{divisors-definition-rank}
Soit $X$ un schéma intègre de point générique $\eta$.
Soit  $\mathcal{F}$ un $\mathcal{O}_X$-module quasi-cohérent. Le
 {\it rang}  de  $\mathcal{F}$  est  $\dim_{\kappa(\eta)} \mathcal{F}_\eta$.
\end{definition}

\noindent
Cela ne contredit pas la définition du rang d'un module localement libre (dans le
cas où les deux définitions s'appliquent). Il ne semble pas judicieux de définir
le rang d'un module quasi-cohérent si le schéma  $X$  n'est pas
irréductible ou si l'anneau local  $\mathcal{O}_{X, \eta}$  au point générique n'est pas un
corps\footnote{S'il  $X$  est irréductible et  $\mathcal{F}_\eta$ 
est libre sur  $\mathcal{O}_{X, \eta}$ , il semblerait naturel d'utiliser le rang de ce module
libre. Attention, certaines références définissent le rang en utilisant la
longueur de  $\mathcal{F}_\eta$.}

\begin{lemma}
\label{divisors-lemma-check-rank-on-affines}
Soit  $X$  un schéma intègre. Soit  $\mathcal{F}$  un  $\mathcal{O}_X$-module
quasi-cohérent. Soit  $r$  un cardinal. Les énoncés suivants sont
équivalents
\begin{enumerate}
\item $\mathcal{F}$  a pour rang  $r$,
\item pour tout ouvert affine non vide  $U \subset X$ ,  $\mathcal{F}(U)$  est un
 $\mathcal{O}_X(U)$-module de rang  $r$, et
\item pour un certain  $U \subset X$  ouvert affine non vide  $\mathcal{F}(U)$  est un
 $\mathcal{O}_X(U)$-module de rang  $r$.
\end{enumerate}
\end{lemma}

\begin{proof}
Démonstration omise.
\end{proof}

\begin{lemma}
\label{divisors-lemma-dominant-pullback-rank}
Soit  $f : X \to Y$  un morphisme dominant de schémas intègres. Soit  $\mathcal{G}$  un
 $\mathcal{O}_Y$-module quasi-cohérent. Alors le rang de  $\mathcal{G}$  sur  $Y$ 
est égal au rang de  $f^*\mathcal{G}$  sur  $X$.
\end{lemma}

\begin{proof}
Utilisons le Lemme  \ref{divisors-lemma-check-rank-on-affines}, Morphismes, Lemme  \ref{morphisms-lemma-dominant-between-integral}, et Compléments
d'algèbre, Lemme  \ref{more-algebra-lemma-dominant-pullback-rank}.
\end{proof}

\begin{lemma}
\label{divisors-lemma-rank-additive}
Soit  $X$  un schéma intègre. Soit  $0 \to \mathcal{F} \to \mathcal{F}' \to \mathcal{F}'' \to 0$  une suite exacte de
 $\mathcal{O}_X$-modules quasi-cohérents. Alors le rang de  $\mathcal{F}'$  est la somme
des rangs de  $\mathcal{F}$  et  $\mathcal{F}''$.
\end{lemma}

\begin{proof}
Démonstration omise. Voir Compléments d'algèbre, Lemme  \ref{more-algebra-lemma-rank-additive}.
\end{proof}

\begin{lemma}
\label{divisors-lemma-rank-tensor}
Soit  $X$  un schéma intègre. Soient  $\mathcal{F}$  et  $\mathcal{G}$  des
 $\mathcal{O}_X$-modules quasi-cohérents.
\begin{enumerate}
\item Le rang de  $\mathcal{F} \otimes_{\mathcal{O}_X} \mathcal{G}$  est le produit des rangs de  $\mathcal{F}$  et  $\mathcal{G}$.
\item Si  $\mathcal{F}$  est de présentation finie, alors le rang de  $\SheafHom_{\mathcal{O}_X}(\mathcal{F}, \mathcal{G})$  est le
produit des rangs de  $\mathcal{F}$  et  $\mathcal{G}$.
\end{enumerate}
\end{lemma}

\begin{proof}
Démonstration omise. Notons que la partie (2) a du sens car  $\SheafHom_{\mathcal{O}_X}(\mathcal{F}, \mathcal{G})$  est quasi-cohérent selon
l'article (10) dans Schémas, section  \ref{schemes-section-quasi-coherent}. La version algébrique de ce
lemme est Compléments d'algèbre, Lemme  \ref{more-algebra-lemma-rank-tensor}.
\end{proof}

\begin{lemma}
\label{divisors-lemma-finite-module-rank}
Soit  $X$  un schéma intègre. Soit  $\mathcal{F}$  un  $\mathcal{O}_X$-module
quasi-cohérent de type fini. Alors
\begin{enumerate}
\item $\mathcal{F}$  a un rang  $r \geq 0$, et
\item il existe un ouvert non vide  $U \subset X$  tel que  $\mathcal{F}|_U$  soit libre de rang
 $r$.
\end{enumerate}
\end{lemma}

\begin{proof}
Voir Compléments d'algèbre, Lemme  \ref{more-algebra-lemma-finite-module-rank}  pour la version algébrique.
\end{proof}











\section{Modules réflexifs}
\label{divisors-section-reflexive}

\noindent
Cette section est l'analogue de Compléments d'algèbre, section  \ref{more-algebra-section-reflexive}  pour les
modules cohérents sur les schémas localement noethériens. La raison de travailler
avec des modules cohérents est que  $\SheafHom_{\mathcal{O}_X}(\mathcal{F}, \mathcal{G})$  est cohérent pour chaque paire de
 $\mathcal{O}_X$-modules cohérents  $\mathcal{F}, \mathcal{G}$, voir Modules, Lemme  \ref{modules-lemma-internal-hom-locally-kernel-direct-sum}.

\begin{definition}
\label{divisors-definition-reflexive}
Soit  $X$  un schéma intègre localement noethérien. Soit  $\mathcal{F}$  un
 $\mathcal{O}_X$-module cohérent. L'{\it  enveloppe
réflexive}  de  $\mathcal{F}$  est le  $\mathcal{O}_X$-module
$$
\mathcal{F}^{**} = \SheafHom_{\mathcal{O}_X}(
\SheafHom_{\mathcal{O}_X}(\mathcal{F}, \mathcal{O}_X), \mathcal{O}_X)
$$
Nous disons que  $\mathcal{F}$  est  {\it  réflexif}  si
l'application naturelle  $j : \mathcal{F} \longrightarrow \mathcal{F}^{**}$  est un isomorphisme.
\end{definition}

\noindent
Il découle du Lemme  \ref{divisors-lemma-dual-reflexive}  que l'enveloppe réflexive est un
 $\mathcal{O}_X$-module réflexif. Vous pouvez utiliser la même définition pour définir
des modules réflexifs dans des situations plus générales, mais cela ne semble pas
très utile. Voici le lemme obligatoire comparant cela à la notion algébrique
habituelle.

\begin{lemma}
\label{divisors-lemma-check-reflexive-on-affines}
Soit  $X$  un schéma intègre localement noethérien. Soit  $\mathcal{F}$  un
 $\mathcal{O}_X$-module cohérent. Les énoncés suivants sont équivalents
\begin{enumerate}
\item $\mathcal{F}$  est réflexif,
\item pour  $U \subset X$  ouvert affine,  $\mathcal{F}(U)$  est un  $\mathcal{O}(U)$-module réflexif.
\end{enumerate}
\end{lemma}

\begin{proof}
Démonstration omise.
\end{proof}

\begin{remark}
\label{divisors-remark-different-reflexive}
Si  $X$  est un schéma de type fini sur un corps, alors parfois une notion
différente de modules réflexifs est utilisée (voir par exemple \cite[bas de la page 5 et définition 1.1.9]{HL}).
Cette autre notion utilise  $R\SheafHom$  dans un complexe dualisant  $\omega_X^\bullet$  au
 lieu de  $\mathcal{O}_X$  et devrait probablement avoir un nom différent car elle
peut être différente lorsque  $X$  n'est pas Gorenstein. Par exemple, si
 $X = \Spec(k[t^3, t^4, t^5])$, alors un calcul montre que le faisceau dualisant  $\omega_X$  n'est pas
réflexif selon notre définition, mais il est réflexif selon l'autre définition car
 $\omega_X \to \SheafHom(\SheafHom(\omega_X, \omega_X), \omega_X)$  est un isomorphisme.
\end{remark}

\begin{lemma}
\label{divisors-lemma-reflexive-torsion-free}
Soit  $X$  un schéma intègre localement noethérien. Soit  $\mathcal{F}$  un
 $\mathcal{O}_X$-module cohérent.
\begin{enumerate}
\item Si  $\mathcal{F}$  est réflexif, alors  $\mathcal{F}$  est sans torsion.
\item L'application  $j : \mathcal{F} \longrightarrow \mathcal{F}^{**}$  est injective si et seulement si  $\mathcal{F}$  est sans
torsion.
\end{enumerate}
\end{lemma}

\begin{proof}
Démonstration omise. Voir Compléments d'algèbre, Lemme  \ref{more-algebra-lemma-reflexive-torsion-free}.
\end{proof}

\begin{lemma}
\label{divisors-lemma-check-reflexive}
Soit  $X$  un schéma intègre localement noethérien. Soit  $\mathcal{F}$  un
 $\mathcal{O}_X$-module cohérent. Les énoncés suivants sont équivalents
\begin{enumerate}
\item $\mathcal{F}$  est réflexif,
\item $\mathcal{F}_x$  est un  $\mathcal{O}_{X, x}$-module réflexif pour tout  $x \in X$,
\item $\mathcal{F}_x$  est un  $\mathcal{O}_{X, x}$-module réflexif pour tous les points fermés
 $x \in X$.
\end{enumerate}
\end{lemma}

\begin{proof}
Par modules, le lemme  \ref{modules-lemma-stalk-internal-hom}  on voit que (1) et (2) sont équivalents.
Puisque chaque point de  $X$  se spécialise vers un point fermé
(Propriétés, lemme  \ref{properties-lemma-locally-Noetherian-closed-point}), on voit que (2) et (3) sont équivalents.
\end{proof}

\begin{lemma}
\label{divisors-lemma-flat-pullback-reflexive}
Soit  $f : X \to Y$  un morphisme plat de schémas localement noethériens intégraux.
Soit  $\mathcal{G}$  un  $\mathcal{O}_Y$-module réflexif cohérent. Alors  $f^*\mathcal{G}$  est
un  $\mathcal{O}_X$-module réflexif cohérent.
\end{lemma}

\begin{proof}
Démonstration omise. Voir Compléments d'algèbre, Lemme  \ref{more-algebra-lemma-flat-pullback-reflexive}  pour l'analogue algébrique.
\end{proof}

\begin{lemma}
\label{divisors-lemma-sequence-reflexive}
Soit  $X$  un schéma intègre localement noethérien. Soit  $0 \to \mathcal{F} \to \mathcal{F}' \to \mathcal{F}''$  une
suite exacte de  $\mathcal{O}_X$-modules cohérents. Si  $\mathcal{F}'$  est réflexif et
 $\mathcal{F}''$  est sans torsion, alors  $\mathcal{F}$  est réflexif.
\end{lemma}

\begin{proof}
Démonstration omise. Voir Compléments d'algèbre, Lemme  \ref{more-algebra-lemma-sequence-reflexive}.
\end{proof}

\begin{lemma}
\label{divisors-lemma-dual-reflexive}
Soit  $X$  un schéma intègre localement noethérien. Soit  $\mathcal{F}$,
 $\mathcal{G}$  des  $\mathcal{O}_X$-modules cohérents. Si  $\mathcal{G}$  est réflexif,
alors  $\SheafHom_{\mathcal{O}_X}(\mathcal{F}, \mathcal{G})$  est réflexif.
\end{lemma}

\begin{proof}
L'énoncé a un sens car  $\SheafHom_{\mathcal{O}_X}(\mathcal{F}, \mathcal{G})$  est cohérent selon Cohomologie des schémas,
Lemme  \ref{coherent-lemma-tensor-hom-coherent}. Pour voir que l'énoncé est vrai, voir Compléments d'algèbre, Lemme
 \ref{more-algebra-lemma-dual-reflexive}. Certains détails sont omis.
\end{proof}

\begin{remark}
\label{divisors-remark-tensor}
Soit  $X$  un schéma intègre localement noethérien. Grâce au Lemme
 \ref{divisors-lemma-dual-reflexive} , nous savons que l'enveloppe réflexive  $\mathcal{F}^{**}$  d'un
 $\mathcal{O}_X$-module cohérent est cohérent réflexif. Considérons la catégorie
 $\mathcal{C}$  des  $\mathcal{O}_X$-modules cohérents réflexifs. La prise des enveloppes
réflexives donne un adjoint à gauche du foncteur d'inclusion  $\mathcal{C} \to \textit{Coh}(\mathcal{O}_X)$.
Remarquons que  $\mathcal{C}$  est une catégorie additive avec des noyaux et des
conoyaux. En effet, étant donné  $\varphi : \mathcal{F} \to \mathcal{G}$  dans  $\mathcal{C}$, le noyau usuel
 $\Ker(\varphi)$  est réflexif (Lemme  \ref{divisors-lemma-sequence-reflexive}) et l'enveloppe réflexive
 $\Coker(\varphi)^{**}$  du conoyau usuel est le conoyau dans  $\mathcal{C}$. De plus,
 $\mathcal{C}$  hérite d'un produit tensoriel
$$
\mathcal{F} \otimes_\mathcal{C} \mathcal{G} =
(\mathcal{F} \otimes_{\mathcal{O}_X} \mathcal{G})^{**}
$$
qui est associatif et symétrique. Il existe un Hom interne au sens où pour trois
objets  $\mathcal{F}, \mathcal{G}, \mathcal{H}$  de  $\mathcal{C}$ , nous avons l'identité
$$
\Hom_\mathcal{C}(\mathcal{F} \otimes_\mathcal{C} \mathcal{G}, \mathcal{H}) =
\Hom_\mathcal{C}(\mathcal{F},
\SheafHom_{\mathcal{O}_X}(\mathcal{G}, \mathcal{H}))
$$
voir Modules, Lemme  \ref{modules-lemma-internal-hom}. Dans  $\mathcal{C}$ , chaque objet  $\mathcal{F}$  a
un objet  {\it  dual }   $\SheafHom_{\mathcal{O}_X}(\mathcal{F}, \mathcal{O}_X)$. Sans conditions
supplémentaires sur  $X$ , il peut arriver que
$$
\SheafHom_{\mathcal{O}_X}(\mathcal{F}, \mathcal{G}) \not \cong
\SheafHom_{\mathcal{O}_X}(\mathcal{F}, \mathcal{O}_X)
\otimes_\mathcal{C} \mathcal{G}
\quad\text{et}\quad
\mathcal{F} \otimes_\mathcal{C}
\SheafHom_{\mathcal{O}_X}(\mathcal{F}, \mathcal{O}_X)
\not \cong \mathcal{O}_X
$$
pour  $\mathcal{F}, \mathcal{G}$  de rang  $1$  dans  $\mathcal{C}$. Pour donner un exemple,
soient  $X = \Spec(R)$  où  $R$  est comme dans Compléments d'algèbre, Exemple
 \ref{more-algebra-example-ring-not-S2}  et soient  $\mathcal{F}, \mathcal{G}$  les modules correspondant à  $M$.
Calcul omis.
\end{remark}

\begin{lemma}
\label{divisors-lemma-reflexive-depth-2}
Soit  $X$  un schéma intègre localement noethérien. Soit  $\mathcal{F}$  un
 $\mathcal{O}_X$-module cohérent. Les énoncés suivants sont équivalents
\begin{enumerate}
\item $\mathcal{F}$  est réflexif,
\item pour chaque  $x \in X$  l'un des cas suivants se produit
\begin{enumerate}
\item $\mathcal{F}_x$  est un  $\mathcal{O}_{X, x}$-module réflexif, ou
\item $\text{depth}(\mathcal{F}_x) \geq 2$.
\end{enumerate}
\end{enumerate}
\end{lemma}

\begin{proof}
Démonstration omise. Voir Compléments d'algèbre, Lemme  \ref{more-algebra-lemma-reflexive-depth-2}.
\end{proof}

\begin{lemma}
\label{divisors-lemma-reflexive-S2}
Soit  $X$  un schéma intègre localement noethérien. Soit  $\mathcal{F}$  un
 $\mathcal{O}_X$-module cohérent réflexif. Soit  $x \in X$.
\begin{enumerate}
\item Si  $\text{depth}(\mathcal{O}_{X, x}) \geq 2$, alors  $\text{depth}(\mathcal{F}_x) \geq 2$.
\item Si  $X$  est  $(S_2)$, alors  $\mathcal{F}$  est  $(S_2)$.
\end{enumerate}
\end{lemma}

\begin{proof}
Démonstration omise. Voir Compléments d'algèbre, Lemme  \ref{more-algebra-lemma-reflexive-S2}.
\end{proof}

\begin{lemma}
\label{divisors-lemma-reflexive-S2-extend}
Soit  $X$  un schéma intègre localement noethérien. Soit  $j : U \to X$  un
sous-schéma ouvert avec complémentaire  $Z$. Supposons que  $\mathcal{O}_{X, z}$ 
ait une profondeur  $\geq 2$  pour tous  $z \in Z$. Alors  $j^*$  et
 $j_*$  définissent une équivalence de catégories entre la catégorie des
 $\mathcal{O}_X$-modules cohérents réflexifs et la catégorie des  $\mathcal{O}_U$-modules
cohérents réflexifs.
\end{lemma}

\begin{proof}
Soit  $\mathcal{F}$  un  $\mathcal{O}_X$-module cohérent réflexif. Pour  $z \in Z$ , la
germe  $\mathcal{F}_z$  a une profondeur  $\geq 2$  d'après le Lemme  \ref{divisors-lemma-reflexive-S2}.
Ainsi,  $\mathcal{F} \to j_*j^*\mathcal{F}$  est un isomorphisme d'après le Lemme  \ref{divisors-lemma-depth-2-hartog}.
Inversement, soit  $\mathcal{G}$  un  $\mathcal{O}_U$-module cohérent réflexif. Il suffit
de montrer que  $j_*\mathcal{G}$  est un  $\mathcal{O}_X$-module cohérent réflexif. Pour le
prouver, nous pouvons supposer que  $X$  est affine. D'après les
Propriétés, Lemme  \ref{properties-lemma-lift-finite-presentation}  il existe un  $\mathcal{O}_X$-module cohérent
 $\mathcal{F}$  avec  $\mathcal{G} = j^*\mathcal{F}$. Après avoir remplacé  $\mathcal{F}$  par son
enveloppe réflexive, nous pouvons supposer que  $\mathcal{F}$  est réflexif (voir la
discussion ci-dessus et en particulier le Lemme  \ref{divisors-lemma-dual-reflexive}). D'après ce qui
précède,  $j_*\mathcal{G} = j_*j^*\mathcal{F} = \mathcal{F}$  comme désiré.
\end{proof}

\noindent
Si le schéma est normal, alors réflexif est la même chose que sans torsion et
 $(S_2)$.

\begin{lemma}
\label{divisors-lemma-reflexive-over-normal}
Soit  $X$  un schéma normal intègre localement noethérien. Soit
 $\mathcal{F}$  un  $\mathcal{O}_X$-modèle cohérent. Les énoncés suivants sont
équivalents:
\begin{enumerate}
\item $\mathcal{F}$  est réflexif,
\item $\mathcal{F}$  est sans torsion et possède la propriété  $(S_2)$, et
\item il existe un sous-schéma ouvert  $j : U \to X$  tel que
\begin{enumerate}
\item chaque composante irréductible de  $X \setminus U$  a une codimension  $\geq 2$  dans
 $X$.
\item $j^*\mathcal{F}$  est localement libre de type fini, et
\item $\mathcal{F} = j_*j^*\mathcal{F}$.
\end{enumerate}
\end{enumerate}
\end{lemma}

\begin{proof}
En utilisant le Lemme  \ref{divisors-lemma-check-reflexive-on-affines} , l'équivalence de (1) et (2) découle de Compléments
d'algèbre, Lemme  \ref{more-algebra-lemma-reflexive-over-normal}. Soit  $U \subset X$  comme en (3). Par Propriétés, Lemme
 \ref{properties-lemma-criterion-normal}  nous voyons que  $\text{depth}(\mathcal{O}_{X, x}) \geq 2$  pour  $x \not \in U$. Comme un module
localement libre de type fini est réflexif, nous en concluons que (3) implique (1)
par le Lemme  \ref{divisors-lemma-reflexive-S2-extend}.

\medskip\noindent
Supposons (1). Soit  $U \subset X$  le sous-schéma ouvert maximal tel que
 $j^*\mathcal{F} = \mathcal{F}|_U$  soit localement libre de type fini. Donc (3)(b) est satisfait. Soit
 $x \in X$  un point. Si  $\mathcal{F}_x$  est un  $\mathcal{O}_{X, x}$-module libre, alors
 $x \in U$, voir Modules, Lemme  \ref{modules-lemma-finite-presentation-stalk-free}. Si  $\dim(\mathcal{O}_{X, x}) \leq 1$, alors
 $\mathcal{O}_{X, x}$  est soit un corps, soit un anneau de valuation discrète (Propriétés,
Lemme  \ref{properties-lemma-criterion-normal}) et donc  $\mathcal{F}_x$  est libre (Compléments d'algèbre, Lemme
 \ref{more-algebra-lemma-dedekind-torsion-free-flat}). Ainsi  $x \not \in U \Rightarrow \dim(\mathcal{O}_{X, x}) \geq 2$. Ensuite, Propriétés, Lemme  \ref{properties-lemma-codimension-local-ring}  montre
que (3)(a) est satisfait. Par le Lemme déjà utilisé des Propriétés,  \ref{properties-lemma-criterion-normal} ,
nous voyons également que  $\text{depth}(\mathcal{O}_{X, x}) \geq 2$  pour  $x \not \in U$  et donc (3)(c) découle du
Lemme  \ref{divisors-lemma-reflexive-S2-extend}.
\end{proof}

\begin{lemma}
\label{divisors-lemma-describe-reflexive-hull}
Soit  $X$  un schéma normal localement noethérien intègre de point
générique  $\eta$. Soient  $\mathcal{F}$,  $\mathcal{G}$  des  $\mathcal{O}_X$-modules
cohérents. Soit  $T : \mathcal{G}_\eta \to \mathcal{F}_\eta$  une application linéaire. Alors  $T$  se
prolonge en une application  $\mathcal{G} \to \mathcal{F}^{**}$  de  $\mathcal{O}_X$-modules si et seulement
si
\begin{itemize}
\item[(*)] pour tout  $x \in X$  avec  $\dim(\mathcal{O}_{X, x}) = 1$  nous avons
$$
T\left(\Im(\mathcal{G}_x \to \mathcal{G}_\eta)\right) \subset
\Im(\mathcal{F}_x \to \mathcal{F}_\eta).
$$
\end{itemize}
\end{lemma}

\begin{proof}
Comme  $\mathcal{F}^{**}$  est sans torsion et que  $\mathcal{F}_\eta = \mathcal{F}^{**}_\eta$, un prolongement, s'il
existe, est unique. Ainsi, il suffit de prouver le lemme sur les éléments d'un
recouvrement ouvert de  $X$, c'est-à-dire que nous pouvons supposer que
 $X$  est affine. Dans ce cas, nous posons la question algébrique suivante:
Soit  $R$  un anneau normal noethérien de corps des fractions  $K$,
soient  $M$,  $N$  des  $R$-modules finis, soit
 $T : M \otimes_R K \to N \otimes_R K$  une application  $K$-linéaire. Quand  $T$  se
prolonge-t-elle en une application  $N \to M^{**}$? Selon Compléments d'algèbre, Lemme
\ref{more-algebra-lemma-describe-reflexive-hull}  cela se produit si et seulement si  $N_\mathfrak p$  s'envoie dans
 $(M/M_{tors})_\mathfrak p$  pour chaque hauteur  $1$  premier  $\mathfrak p$  de
 $R$. C'est exactement la condition  $(*)$  du lemme.
\end{proof}

\begin{lemma}
\label{divisors-lemma-reflexive-over-regular-dim-2}
Soit  $X$  un schéma régulier de dimension  $\leq 2$. Soit  $\mathcal{F}$ 
un  $\mathcal{O}_X$-module cohérent. Les affirmations suivantes sont équivalentes
\begin{enumerate}
\item $\mathcal{F}$  est réflexif,
\item $\mathcal{F}$  est localement libre de type fini.
\end{enumerate}
\end{lemma}

\begin{proof}
Il est clair qu'un module localement libre de type fini est réflexif. Pour la réciproque,
nous montrerons que si  $\mathcal{F}$  est réflexif, alors  $\mathcal{F}_x$  est un
 $\mathcal{O}_{X, x}$-module libre pour tout  $x \in X$. Cela suffit d'après Algèbre,
Lemme  \ref{algebra-lemma-finite-projective}  et le fait que  $\mathcal{F}$  est cohérent. Si  $\dim(\mathcal{O}_{X, x}) = 0$,
alors  $\mathcal{O}_{X, x}$  est un corps et l'énoncé est clair. Si  $\dim(\mathcal{O}_{X, x}) = 1$, alors
 $\mathcal{O}_{X, x}$  est un anneau de valuation discrète (Algèbre, Lemme  \ref{algebra-lemma-characterize-dvr}) et
 $\mathcal{F}_x$  est sans torsion. Par conséquent,  $\mathcal{F}_x$  est libre d'après
 Compléments d'algèbre, Lemme  \ref{more-algebra-lemma-dedekind-torsion-free-flat}. Si  $\dim(\mathcal{O}_{X, x}) = 2$, alors  $\mathcal{O}_{X, x}$  est un
 anneau local régulier de dimension  $2$. D'après Compléments d'algèbre, Lemme
 \ref{more-algebra-lemma-reflexive-over-normal}  nous voyons que  $\mathcal{F}_x$  a une profondeur  $\geq 2$. Par
conséquent,  $\mathcal{F}_x$  est libre d'après Algèbre, Lemme  \ref{algebra-lemma-regular-mcm-free}.
\end{proof}







\section{Diviseurs de Cartier effectifs}
\label{divisors-section-effective-Cartier-divisors}

\noindent
Nous définissons la notion de diviseur de Cartier effectif avant tout autre type
de diviseur.

\begin{definition}
\label{divisors-definition-effective-Cartier-divisor}
Soit  $S$  un schéma.
\begin{enumerate}
\item Un {\it sous-schéma fermé localement principal}  de
 $S$ est un sous-schéma fermé dont le faisceau d'idéaux est localement
engendré par un seul élément.
\item Un {\it diviseur de Cartier effectif} sur $S$
est un sous-schéma fermé  $D \subset S$ dont le faisceau d'idéaux $\mathcal{I}_D \subset \mathcal{O}_S$  est
un $\mathcal{O}_S$-module inversible.
\end{enumerate}
\end{definition}

\noindent
Ainsi, un diviseur de Cartier effectif est un sous-schéma fermé localement
principal, mais la réciproque n'est pas toujours vraie. Les diviseurs de Cartier
effectifs sont des sous-schémas fermés de codimension pure  $1$  dans le
sens le plus fort possible. En effet, ils sont localement définis par un seul
élément qui est un non-diviseur de zéro. En particulier, ils sont nulle part
denses.

\begin{lemma}
\label{divisors-lemma-characterize-effective-Cartier-divisor}
Soit  $S$  un schéma. Soit  $D \subset S$  un sous-schéma fermé. Les
assertions suivantes sont équivalentes:
\begin{enumerate}
\item Le sous-schéma  $D$  est un diviseur de Cartier effectif sur  $S$.
\item Pour chaque  $x \in D$ , il existe un voisinage ouvert affine  $\Spec(A) = U \subset S$  de
 $x$  tel que  $U \cap D = \Spec(A/(f))$  avec  $f \in A$  un non-diviseur de zéro.
\end{enumerate}
\end{lemma}

\begin{proof}
Supposons (1). Pour chaque  $x \in D$ , il existe un voisinage ouvert affine
 $\Spec(A) = U \subset S$  de  $x$  tel que  $\mathcal{I}_D|_U \cong \mathcal{O}_U$. En d'autres termes, il existe
une section  $f \in \Gamma(U, \mathcal{I}_D)$  qui engendre librement la restriction  $\mathcal{I}_D|_U$. Ainsi
 $f \in A$, et l'application de multiplication  $f : A \to A$  est injective. De
plus, comme  $\mathcal{I}_D$  est quasi-cohérent, nous voyons que  $D \cap U = \Spec(A/(f))$.

\medskip\noindent
Supposons (2). Soit  $x \in D$. D'après l'hypothèse, il existe un voisinage
ouvert affine  $\Spec(A) = U \subset S$  de  $x$  tel que  $U \cap D = \Spec(A/(f))$  avec  $f \in A$ 
un non-diviseur de zéro. Alors  $\mathcal{I}_D|_U \cong \mathcal{O}_U$  puisque c'est égal à  $\widetilde{(f)} \cong \widetilde{A} \cong \mathcal{O}_U$.
Bien sûr,  $\mathcal{I}_D$  restreint au sous-schéma ouvert  $S \setminus D$  est isomorphe
à  $\mathcal{O}_{S \setminus D}$. Ainsi  $\mathcal{I}_D$  est un  $\mathcal{O}_S$-module inversible.
\end{proof}

\begin{lemma}
\label{divisors-lemma-complement-locally-principal-closed-subscheme}
Soit  $S$  un schéma. Soit  $Z \subset S$  un sous-schéma fermé localement
principal. Soit  $U = S \setminus Z$. Alors  $U \to S$  est un morphisme affine.
\end{lemma}

\begin{proof}
La question est locale sur  $S$, voir Morphismes, Lemme  \ref{morphisms-lemma-characterize-affine}.
Ainsi nous pouvons supposer  $S = \Spec(A)$  et  $Z = V(f)$  pour un certain
 $f \in A$. Dans ce cas,  $U = D(f) = \Spec(A_f)$  est affine donc  $U \to S$  est affine.
\end{proof}

\begin{lemma}
\label{divisors-lemma-complement-effective-Cartier-divisor}
Soit  $S$  un schéma. Soit  $D \subset S$  un diviseur de Cartier effectif.
Soit  $U = S \setminus D$. Alors  $U \to S$  est un morphisme affine et  $U$  est
schématiquement dense dans  $S$.
\end{lemma}

\begin{proof}
L'affinité est le Lemme  \ref{divisors-lemma-complement-locally-principal-closed-subscheme}. La question de la densité est locale sur
 $S$, voir Morphismes, Lemme  \ref{morphisms-lemma-characterize-scheme-theoretically-dense}. Ainsi nous pouvons supposer
 $S = \Spec(A)$  et  $D$  correspondant au non-diviseur de zéro  $f \in A$,
voir Lemme  \ref{divisors-lemma-characterize-effective-Cartier-divisor}. Ainsi  $A \subset A_f$  ce qui implique que  $U \subset S$  est
schématiquement dense, voir Morphismes, Exemple  \ref{morphisms-example-scheme-theoretic-closure}.
\end{proof}

\begin{lemma}
\label{divisors-lemma-effective-Cartier-makes-dimension-drop}
Soit  $S$  un schéma. Soit  $D \subset S$  un diviseur de Cartier effectif.
Soit  $s \in D$. Si  $\dim_s(S) < \infty$, alors  $\dim_s(D) < \dim_s(S)$.
\end{lemma}

\begin{proof}
Supposons  $\dim_s(S) < \infty$. Soit  $U = \Spec(A) \subset S$  un voisinage ouvert affine de
 $s$  tel que  $\dim(U) = \dim_s(S)$  et tel que  $D = V(f)$  pour un certain
non-diviseur de zéro  $f \in A$  (voir le Lemme  \ref{divisors-lemma-characterize-effective-Cartier-divisor}). Rappelons que  $\dim(U)$  est
la dimension de Krull de l'anneau  $A$  et que  $\dim(U \cap D)$  est la
dimension de Krull de l'anneau  $A/(f)$. Alors  $f$  n'est contenu dans
aucun idéal premier minimal de  $A$. Ainsi, toute chaîne maximale d'idéaux
premiers dans  $A/(f)$, considérée comme une chaîne d'idéaux premiers dans
 $A$, peut être étendue en ajoutant un idéal premier minimal.
\end{proof}

\begin{definition}
\label{divisors-definition-sum-effective-Cartier-divisors}
Soit  $S$  un schéma. Étant donné des diviseurs de Cartier effectifs
 $D_1$,  $D_2$  sur  $S$ , nous posons  $D = D_1 + D_2$  égal au
sous-schéma fermé de  $S$  correspondant au faisceau quasi-cohérent
d'idéaux  $\mathcal{I}_{D_1}\mathcal{I}_{D_2} \subset \mathcal{O}_S$. Nous appelons ce sous-schéma la  {\it  somme
des diviseurs de Cartier effectifs  $D_1$  et  $D_2$}.
\end{definition}

\noindent
Il est clair que nous pouvons définir la somme  $\sum n_iD_i$  donnée par un nombre
fini de diviseurs de Cartier effectifs  $D_i$  sur  $X$  et des
entiers non négatifs  $n_i$.

\begin{lemma}
\label{divisors-lemma-sum-effective-Cartier-divisors}
La somme de deux diviseurs de Cartier effectifs est un diviseur de Cartier
effectif.
\end{lemma}

\begin{proof}
 Preuve omise. Localement,  $f_1, f_2 \in A$  sont des non-diviseurs de zéro ; alors
 $f_1f_2 \in A$  est aussi un non-diviseur de zéro.
\end{proof}

\begin{lemma}
\label{divisors-lemma-difference-effective-Cartier-divisors}
Soit  $X$  un schéma. Soient  $D, D'$  deux diviseurs de Cartier
effectifs sur  $X$. Si  $D \subset D'$  (en tant que sous-schémas fermés de
 $X$), alors il existe un diviseur de Cartier effectif  $D''$  tel
que  $D' = D + D''$.
\end{lemma}

\begin{proof}
Démonstration omise.
\end{proof}

\begin{lemma}
\label{divisors-lemma-sum-closed-subschemes-effective-Cartier}
Soit  $X$  un schéma. Soient  $Z, Y$  deux sous-schémas fermés de
 $X$  avec des faisceaux d'idéaux  $\mathcal{I}$  et  $\mathcal{J}$. Si
 $\mathcal{I}\mathcal{J}$  définit un diviseur de Cartier effectif  $D \subset X$, alors
 $Z$  et  $Y$  sont des diviseurs de Cartier effectifs et
 $D = Z + Y$.
\end{lemma}

\begin{proof}
En appliquant le Lemme  \ref{divisors-lemma-characterize-effective-Cartier-divisor} , nous obtenons la situation algébrique
suivante:  $A$  est un anneau,  $I, J \subset A$  des idéaux et  $f \in A$  un
élément non-diviseur de zéro tel que  $IJ = (f)$. Ainsi, le résultat découle de
Algèbre, Lemme  \ref{algebra-lemma-product-ideals-principal}.
\end{proof}

\begin{lemma}
\label{divisors-lemma-sum-effective-Cartier-divisors-union}
Soit  $X$  un schéma. Soient  $D, D' \subset X$  des diviseurs de Cartier
effectifs tels que l'intersection schématique  $D \cap D'$  soit un diviseur de
Cartier effectif sur  $D'$. Alors  $D + D'$  est l'union schématique de
 $D$  et  $D'$.
\end{lemma}

\begin{proof}
Voir Morphismes, Définition  \ref{morphisms-definition-scheme-theoretic-intersection-union}  pour la définition de l'intersection et
de l'union schématiques. Pour prouver le lemme en travaillant localement (en
utilisant le Lemme  \ref{divisors-lemma-characterize-effective-Cartier-divisor}) nous obtenons le problème algébrique suivant:
Étant donné un anneau  $A$  et des non-diviseurs de zéro  $f_1, f_2 \in A$  tels
que  $f_1$  s'envoie sur un non-diviseur de zéro dans  $A/f_2A$, montrer
que  $f_1A \cap f_2A = f_1f_2A$. Nous omettons l'argument direct.
\end{proof}

\noindent
Rappelons que nous avons défini l'image réciproque d'un sous-schéma fermé par
tout morphisme de schémas dans Schémas, Définition  \ref{schemes-definition-inverse-image-closed-subscheme}.

\begin{lemma}
\label{divisors-lemma-pullback-locally-principal}
Soit  $f : S' \to S$  un morphisme de schémas. Soit  $Z \subset S$  un sous-schéma fermé
localement principal. Alors l'image réciproque  $f^{-1}(Z)$  est un sous-schéma fermé
localement principal de  $S'$.
\end{lemma}

\begin{proof}
Démonstration omise.
\end{proof}

\begin{definition}
\label{divisors-definition-pullback-effective-Cartier-divisor}
Soit  $f : S' \to S$  un morphisme de schémas. Soit  $D \subset S$  un diviseur de
Cartier effectif. Nous disons que  {\it  l'image réciproque de
 $D$  par  $f$  est définie}  si le sous-schéma fermé
 $f^{-1}(D) \subset S'$  est un diviseur de Cartier effectif. Dans ce cas, nous le notons soit
 $f^*D$  soit  $f^{-1}(D)$  et nous l'appelons l'{\it image réciproque
du diviseur de Cartier effectif}.
\end{definition}

\noindent
La condition selon laquelle  $f^{-1}(D)$  est un diviseur de Cartier effectif est
souvent satisfaite en pratique. Voici un lemme exemple.

\begin{lemma}
\label{divisors-lemma-pullback-effective-Cartier-defined}
Soit  $f : X \to Y$  un morphisme de schémas. Soit  $D \subset Y$  un diviseur de
Cartier effectif. L'image réciproque de  $D$  par  $f$  est définie
dans chacun des cas suivants:
\begin{enumerate}
\item $f(x) \not \in D$  pour tout point faiblement associé  $x$  de  $X$,
\item $X$,  $Y$  intègre et  $f$  dominant,
\item $X$  réduit et  $f(\xi) \not \in D$  pour tout point générique  $\xi$  de
toute composante irréductible de  $X$,
\item $X$  est localement noethérien et  $f(x) \not \in D$  pour tout point associé
 $x$  de  $X$,
\item $X$  est localement noethérien, n'a pas de points immergés, et
 $f(\xi) \not \in D$  pour tout point générique  $\xi$  d'une composante irréductible
de  $X$,
\item $f$  est plat, et
\item ajouter plus ici si nécessaire.
\end{enumerate}
\end{lemma}

\begin{proof}
La question est locale sur  $X$, et nous réduisons donc au cas où
 $X = \Spec(A)$,  $Y = \Spec(R)$,  $f$  est donné par  $\varphi : R \to A$  et
 $D = \Spec(R/(t))$  où  $t \in R$  est un non-diviseur de zéro. L'objectif dans
chaque cas est de montrer que  $\varphi(t) \in A$  est un non-diviseur de zéro.

\medskip\noindent
Dans le cas (1), cela découle du Lemme d'Algèbre  \ref{algebra-lemma-weakly-ass-zero-divisors}. Le cas (4) est un
cas particulier du (1) par le Lemme  \ref{divisors-lemma-ass-weakly-ass}. Le cas (5) découle du (4) et des
définitions. Le cas (3) est un cas particulier du (1) par le Lemme  \ref{divisors-lemma-weakass-reduced}.
Le cas (2) est un cas particulier du (3). Si  $R \to A$  est plat, alors
 $t : R \to R$  étant injectif montre que  $t : A \to A$  est injectif. Cela prouve (6).
\end{proof}

\begin{lemma}
\label{divisors-lemma-pullback-effective-Cartier-divisors-additive}
Soit  $f : S' \to S$  un morphisme de schémas. Soient  $D_1$,  $D_2$  des
diviseurs de Cartier effectifs sur  $S$. Si les images réciproques de  $D_1$ 
et  $D_2$  sont définies, alors l'image réciproque de  $D = D_1 + D_2$  est définie et
 $f^*D = f^*D_1 + f^*D_2$.
\end{lemma}

\begin{proof}
Démonstration omise.
\end{proof}





\section{Diviseurs de Cartier effectifs et faisceaux inversibles}
\label{divisors-section-effective-Cartier-invertible}

\noindent
Puisqu'un diviseur de Cartier effectif possède un faisceau d'idéaux inversible
(Définition  \ref{divisors-definition-effective-Cartier-divisor}), la définition suivante a un sens.

\begin{definition}
\label{divisors-definition-invertible-sheaf-effective-Cartier-divisor}
Soit  $S$  un schéma. Soit  $D \subset S$  un diviseur de Cartier effectif
avec le faisceau d'idéaux  $\mathcal{I}_D$.
\begin{enumerate}
\item Le  {\it faisceau inversible $\mathcal{O}_S(D)$ associé
à $D$} est défini par
$$
\mathcal{O}_S(D) =
\SheafHom_{\mathcal{O}_S}(\mathcal{I}_D, \mathcal{O}_S) =
\mathcal{I}_D^{\otimes -1}.
$$
\item La  {\it section canonique}, généralement notée
 $1$  ou  $1_D$, est la section globale de  $\mathcal{O}_S(D)$ correspondant
à l'application d'inclusion $\mathcal{I}_D \to \mathcal{O}_S$.
\item Nous écrivons  $\mathcal{O}_S(-D) = \mathcal{O}_S(D)^{\otimes -1} = \mathcal{I}_D$.
\item Étant donné un second diviseur de Cartier effectif  $D' \subset S$ , nous définissons
 $\mathcal{O}_S(D - D') =
\mathcal{O}_S(D) \otimes_{\mathcal{O}_S} \mathcal{O}_S(-D')$.
\end{enumerate}
\end{definition}

\noindent
Quelques commentaires. Nous verrons ci-dessous que l'assignation  $D \mapsto \mathcal{O}_S(D)$ 
transforme l'addition des diviseurs de Cartier effectifs (Définition  \ref{divisors-definition-sum-effective-Cartier-divisors})
en addition dans le groupe de Picard de  $S$  (Lemme  \ref{divisors-lemma-invertible-sheaf-sum-effective-Cartier-divisors}).
Cependant, l'expression  $D - D'$  dans la définition ci-dessus n'a aucun sens
géométrique. Plus précisément, nous pouvons considérer l'ensemble des diviseurs de
Cartier effectifs sur  $S$  comme un monoïde commutatif  $\text{EffCart}(S)$  dont
l'élément zéro est le diviseur de Cartier effectif vide. Alors l'assignation
 $(D, D') \mapsto \mathcal{O}_S(D - D')$  définit un homomorphisme de groupes
$$
\text{EffCart}(S)^{gp} \longrightarrow \Pic(S)
$$
où le côté gauche est l'achèvement en groupe de  $\text{EffCart}(S)$. En d'autres termes,
lorsque nous écrivons  $\mathcal{O}_S(D - D')$ , nous pouvons penser à  $D - D'$  comme un
élément de  $\text{EffCart}(S)^{gp}$.

\begin{lemma}
\label{divisors-lemma-conormal-effective-Cartier-divisor}
Soit  $S$  un schéma et soit  $D \subset S$  un diviseur de Cartier effectif.
Alors le faisceau conormal est  $\mathcal{C}_{D/S} = \mathcal{I}_D|_D =
\mathcal{O}_S(-D)|_D$  et le faisceau normal est
 $\mathcal{N}_{D/S} = \mathcal{O}_S(D)|_D$.
\end{lemma}

\begin{proof}
Cela découle de Morphismes, Lemme  \ref{morphisms-lemma-affine-conormal}.
\end{proof}

\begin{lemma}
\label{divisors-lemma-ses-add-divisor}
Soit  $X$  un schéma. Soient  $D, D' \subset X$  des diviseurs de Cartier
effectifs avec  $D \subset D'$  et soit  $C \subset X$  un diviseur de Cartier effectif
tel que  $D' = D + C$  (qui existe selon le Lemme  \ref{divisors-lemma-difference-effective-Cartier-divisors}). Alors il existe une
suite exacte
$$
0 \to \mathcal{O}_X(-D)|_C \to \mathcal{O}_{D'} \to \mathcal{O}_D \to 0
$$
de  $\mathcal{O}_X$-modules.
\end{lemma}

\begin{proof}
Dans l'énoncé du lemme et dans la démonstration, nous utilisons l'équivalence des
morphismes, Lemme  \ref{morphisms-lemma-i-star-equivalence} , pour penser aux modules quasi-cohérents sur les
sous-schémas fermés de  $X$  comme des modules quasi-cohérents sur
 $X$. Soit  $\mathcal{I}$  le faisceau d'idéaux de  $D$  dans
 $D'$  et soit  $\mathcal{I}_D = \mathcal{O}_X(-D)$  le faisceau d'idéaux de  $D$  dans
 $X$. Alors il existe une surjection canonique  $\mathcal{I}_D \to \mathcal{I}$. Ainsi, il
suffit de montrer que le noyau de cette application est égal à  $\mathcal{I}_C\mathcal{I}_D$  (avec
la notation évidente). Cela peut être vérifié localement, donc nous pouvons
supposer  $X = \Spec(A)$,  $D = \Spec(A/fA)$  et  $C = \Spec(A/gA)$  où  $f, g \in A$  sont des
non-diviseurs de zéro (Lemme  \ref{divisors-lemma-characterize-effective-Cartier-divisor}). Alors  $\mathcal{I}_D$  correspond à
 $fA$, le non-diviseur de zéro  $fg \in A$  génère l'idéal de  $D'$,
et  $\mathcal{I}$  correspond à  $I = fA/fgA$. Ainsi, le résultat est clair.
\end{proof}

\begin{lemma}
\label{divisors-lemma-invertible-sheaf-sum-effective-Cartier-divisors}
Soit  $S$  un schéma. Soient  $D_1$,  $D_2$  des diviseurs de
Cartier effectifs sur  $S$. Soit  $D = D_1 + D_2$. Alors il existe un
isomorphisme unique
$$
\mathcal{O}_S(D_1) \otimes_{\mathcal{O}_S} \mathcal{O}_S(D_2)
\longrightarrow
\mathcal{O}_S(D)
$$
qui envoie  $1_{D_1} \otimes 1_{D_2}$  sur  $1_D$.
\end{lemma}

\begin{proof}
Démonstration omise.
\end{proof}

\begin{lemma}
\label{divisors-lemma-pullback-effective-Cartier-divisors}
Soit  $f : S' \to S$  un morphisme de schémas. Soit  $D$  un diviseur de
Cartier effectif sur  $S$. Si l'image réciproque de  $D$  est
définie, alors  $f^*\mathcal{O}_S(D) = \mathcal{O}_{S'}(f^*D)$  et l'image réciproque de la section canonique  $1_D$  est
la section canonique  $1_{f^*D}$.
\end{lemma}

\begin{proof}
Démonstration omise.
\end{proof}

\begin{definition}
\label{divisors-definition-regular-section}
Soit  $(X, \mathcal{O}_X)$  un espace localement annelé. Soit  $\mathcal{L}$  un faisceau
inversible sur  $X$. Une section globale  $s \in \Gamma(X, \mathcal{L})$  est appelée
 {\it  section régulière}  si l'application
 $\mathcal{O}_X \to \mathcal{L}$,  $f \mapsto fs$  est injective.
\end{definition}

\begin{lemma}
\label{divisors-lemma-regular-section-structure-sheaf}
Soit  $X$  un espace localement annelé. Soit  $f \in \Gamma(X, \mathcal{O}_X)$. Les énoncés
suivants sont équivalents:
\begin{enumerate}
\item $f$  est une section régulière, et
\item pour tout  $x \in X$  l'image  $f \in \mathcal{O}_{X, x}$  est un non-diviseur de zéro.
\end{enumerate}
Si  $X$  est un schéma, cela équivaut également à
\begin{enumerate}
\item[(3)] pour tout ouvert affine  $\Spec(A) = U \subset X$  l'image  $f \in A$  est un non-diviseur de
zéro,
\item[(4)] il existe un recouvrement par ouverts affines  $X = \bigcup \Spec(A_i)$  tel que l'image de
 $f$  dans  $A_i$  soit un non-diviseur de zéro pour tout
 $i$.
\end{enumerate}
\end{lemma}

\begin{proof}
Démonstration omise.
\end{proof}

\noindent
Notez qu'une section globale  $s$  d'un  $\mathcal{O}_X$-module inversible
 $\mathcal{L}$  peut être vue comme une application de  $\mathcal{O}_X$-module
 $s : \mathcal{O}_X \to \mathcal{L}$. Son dual est donc une application  $s : \mathcal{L}^{\otimes -1} \to \mathcal{O}_X$. (Voir Modules,
Définition  \ref{modules-definition-powers}  pour la définition du faisceau inversible dual.)

\begin{definition}
\label{divisors-definition-zero-scheme-s}
Soit  $X$  un schéma. Soit  $\mathcal{L}$  un faisceau inversible. Soit
 $s \in \Gamma(X, \mathcal{L})$  une section globale. Le  {\it  schéma des
zéros}  de  $s$  est le sous-schéma fermé  $Z(s) \subset X$  défini par
le faisceau quasi-cohérent d'idéaux  $\mathcal{I} \subset \mathcal{O}_X$  qui est l'image de l'application
 $s : \mathcal{L}^{\otimes -1} \to \mathcal{O}_X$.
\end{definition}

\begin{lemma}
\label{divisors-lemma-zero-scheme}
Soit  $X$  un schéma. Soit  $\mathcal{L}$  un faisceau inversible. Soit
 $s \in \Gamma(X, \mathcal{L})$.
\begin{enumerate}
\item Considérons des immersions fermées  $i : Z \to X$  telles que  $i^*s \in \Gamma(Z, i^*\mathcal{L})$  soit
ordonné zéro par inclusion. Le schéma des zéros  $Z(s)$  est l'élément maximal
de cet ensemble ordonné.
\item Pour tout morphisme de schémas  $f : Y \to X$ , nous avons  $f^*s = 0$  dans
 $\Gamma(Y, f^*\mathcal{L})$  si et seulement si  $f$  se factorise par  $Z(s)$.
\item Le schéma des zéros  $Z(s)$  est un sous-schéma fermé localement principal.
\item Le schéma des zéros  $Z(s)$  est un diviseur de Cartier effectif si et seulement si
 $s$  est une section régulière de  $\mathcal{L}$.
\end{enumerate}
\end{lemma}

\begin{proof}
Démonstration omise.
\end{proof}

\begin{lemma}
\label{divisors-lemma-characterize-OD}
\begin{slogan}
Les diviseurs de Cartier effectifs sur un schéma s'identifient aux faisceaux
inversibles munis d'une section globale régulière.
\end{slogan}
Soit  $X$  un schéma.
\begin{enumerate}
\item Si  $D \subset X$  est un diviseur de Cartier effectif, alors la section canonique
 $1_D$  de  $\mathcal{O}_X(D)$  est régulière.
\item Inversement, si  $s$  est une section régulière du faisceau inversible
 $\mathcal{L}$, alors il existe un unique diviseur de Cartier effectif  $D = Z(s) \subset X$ 
et un unique isomorphisme  $\mathcal{O}_X(D) \to \mathcal{L}$  qui envoie  $1_D$  sur  $s$.
\end{enumerate}
Les constructions  $D \mapsto (\mathcal{O}_X(D), 1_D)$  et  $(\mathcal{L}, s) \mapsto Z(s)$  fournissent des applications
mutuellement inverses.
$$
\left\{
\begin{matrix}
\text{effective Cartier divisors on }X
\end{matrix}
\right\}
\leftrightarrow
\left\{
\begin{matrix}
\text{isomorphism classes of pairs }(\mathcal{L}, s)\\
\text{consisting of an invertible }
\mathcal{O}_X\text{-module}\\
\mathcal{L}\text{ et une section globale régulière }s
\end{matrix}
\right\}
$$
\end{lemma}

\begin{proof}
Démonstration omise.
\end{proof}

\begin{remark}
\label{divisors-remark-ses-regular-section}
Soit  $X$  un schéma,  $\mathcal{L}$  un  $\mathcal{O}_X$-module inversible, et
 $s$  une section régulière de  $\mathcal{L}$. Alors le schéma des zéros
 $D = Z(s)$  est un diviseur de Cartier effectif sur  $X$  et il existe des
suites exactes courtes
$$
0 \to \mathcal{O}_X \to \mathcal{L} \to i_*(\mathcal{L}|_D) \to 0
\quad\text{et}\quad
0 \to \mathcal{L}^{\otimes -1} \to \mathcal{O}_X \to i_*\mathcal{O}_D \to 0.
$$
Étant donné un diviseur de Cartier effectif  $D \subset X$  en utilisant les lemmes
 \ref{divisors-lemma-characterize-OD}  et  \ref{divisors-lemma-conormal-effective-Cartier-divisor} , nous obtenons
$$
0 \to \mathcal{O}_X \to \mathcal{O}_X(D) \to i_*(\mathcal{N}_{D/X}) \to 0
\quad\text{et}\quad
0 \to \mathcal{O}_X(-D) \to \mathcal{O}_X \to i_*(\mathcal{O}_D) \to 0
$$
\end{remark}






\section{Diviseurs de Cartier effectifs sur les schémas noethériens}
\label{divisors-section-Noetherian-effective-Cartier}

\noindent
Dans le cadre localement noethérien, la plupart des discussions sur les diviseurs
de Cartier effectifs et les sections régulières se simplifient quelque peu.

\begin{lemma}
\label{divisors-lemma-regular-section-associated-points}
Soit  $X$  un schéma localement noethérien. Soit  $\mathcal{L}$  un
 $\mathcal{O}_X$-module inversible. Soit  $s \in \Gamma(X, \mathcal{L})$. Alors  $s$  est une
section régulière si et seulement si  $s$  ne s'annule pas aux points
associés de  $X$.
\end{lemma}

\begin{proof}
 Preuve omise. Indication : se ramener au cas affine où  $\mathcal{L}$  est trivial, puis utiliser le
lemme  \ref{divisors-lemma-regular-section-structure-sheaf}  et Algèbre, lemme  \ref{algebra-lemma-ass-zero-divisors}.
\end{proof}

\begin{lemma}
\label{divisors-lemma-effective-Cartier-in-points}
Soit  $X$  un schéma localement noethérien. Soit  $D \subset X$  un
sous-schéma fermé correspondant au faisceau d'idéaux quasi-cohérent  $\mathcal{I} \subset \mathcal{O}_X$.
\begin{enumerate}
\item Si pour chaque  $x \in D$  l'idéal  $\mathcal{I}_x \subset \mathcal{O}_{X, x}$  peut être engendré par un seul
élément, alors  $D$  est localement principal.
\item Si pour chaque  $x \in D$  l'idéal  $\mathcal{I}_x \subset \mathcal{O}_{X, x}$  peut être engendré par un seul
élément non diviseur de zéro, alors  $D$  est un diviseur de Cartier
effectif.
\end{enumerate}
\end{lemma}

\begin{proof}
Soit  $\Spec(A)$  un voisinage affine d'un point  $x \in D$. Soit  $\mathfrak p \subset A$ 
l'idéal premier correspondant à  $x$. Soit  $I \subset A$  l'idéal
définissant la trace de  $D$  sur  $\Spec(A)$. Puisque  $A$  est
noethérien (car  $X$  est localement noethérien), l'idéal  $I$  est
engendré par un nombre fini d'éléments, disons  $I = (f_1, \ldots, f_r)$. Sous l'hypothèse de
(1), nous avons  $I_\mathfrak p = (f)$  pour un certain  $f \in A_\mathfrak p$. Alors  $f_i = g_i f$  pour
un certain  $g_i \in A_\mathfrak p$. Écrivons  $g_i = a_i/h_i$  et  $f = f'/h$  pour certains
 $a_i, h_i, f', h \in A$,  $h_i, h \not \in \mathfrak p$. Alors  $I_{h_1 \ldots h_r h} \subset A_{h_1 \ldots h_r h}$  est principal, car il est engendré
par  $f'$. Cela prouve (1). Pour (2) nous pouvons supposer  $I = (f)$.
L'hypothèse implique que l'image de  $f$  dans  $A_\mathfrak p$  est un élément
non diviseur de zéro. Alors  $f$  est un élément non diviseur de zéro dans
un voisinage de  $x$  par Algèbre, Lemme  \ref{algebra-lemma-regular-sequence-in-neighbourhood}. Cela prouve (2).
\end{proof}

\begin{lemma}
\label{divisors-lemma-effective-Cartier-codimension-1}
Soit  $X$  un schéma localement noethérien.
\begin{enumerate}
\item Soit  $D \subset X$  un sous-schéma fermé localement principal. Soit  $\xi \in D$  un
point générique d'une composante irréductible de  $D$. Alors  $\dim(\mathcal{O}_{X, \xi}) \leq 1$.
\item Soit  $D \subset X$  un diviseur de Cartier effectif. Soit  $\xi \in D$  un point
générique d'une composante irréductible de  $D$. Alors  $\dim(\mathcal{O}_{X, \xi}) = 1$.
\end{enumerate}
\end{lemma}

\begin{proof}
Preuve de (1). Par hypothèse, nous pouvons supposer  $X = \Spec(A)$  et  $D = \Spec(A/(f))$ 
où  $A$  est un anneau noethérien et  $f \in A$. Soit  $\xi$ 
correspondant à l'idéal premier  $\mathfrak p \subset A$. L'hypothèse que  $\xi$  est un
point générique d'une composante irréductible de  $D$  signifie que
 $\mathfrak p$  est minimal sur  $(f)$. Ainsi  $\dim(A_\mathfrak p) \leq 1$  par Algèbre, Lemme
 \ref{algebra-lemma-minimal-over-1}.

\medskip\noindent
Preuve de (2). D'après la partie (1), nous voyons que  $\dim(\mathcal{O}_{X, \xi}) \leq 1$. D'autre part,
l'équation locale  $f$  est un non-diviseur de zéro dans  $A_\mathfrak p$  par le Lemme
 \ref{divisors-lemma-characterize-effective-Cartier-divisor}  ce qui implique que la dimension est au moins  $1$  (parce
qu'il doit y avoir un idéal premier dans  $A_\mathfrak p$  ne contenant pas  $f$
selon le Lemme élémentaire de l'Algèbre  \ref{algebra-lemma-Zariski-topology}).
\end{proof}

\begin{lemma}
\label{divisors-lemma-integral-effective-Cartier-divisor-dvr}
Soit  $X$  un schéma noethérien. Soit  $D \subset X$  un sous-schéma fermé
intégral qui est aussi un diviseur de Cartier effectif. Alors l'anneau local de
 $X$  au point générique de  $D$  est un anneau de valuation
discrète.
\end{lemma}

\begin{proof}
D'après le Lemme  \ref{divisors-lemma-characterize-effective-Cartier-divisor} , nous pouvons supposer que  $X = \Spec(A)$  et
 $D = \Spec(A/(f))$  où  $A$  est un anneau noethérien et  $f \in A$  est un
non-diviseur de zéro. L'hypothèse que  $D$  est intégral signifie que
 $(f)$  est premier. Par conséquent, l'anneau local de  $X$  au point
générique est  $A_{(f)}$ , qui est un anneau local noethérien dont l'idéal
maximal est engendré par un non-diviseur de zéro. Ainsi, c'est un anneau de
valuation discrète d'après Algèbre, Lemme  \ref{algebra-lemma-characterize-dvr}.
\end{proof}

\begin{lemma}
\label{divisors-lemma-effective-Cartier-divisor-Sk}
Soit  $X$  un schéma localement noethérien. Soit  $D \subset X$  un diviseur
de Cartier effectif. Si  $X$  est  $(S_k)$, alors  $D$  est
 $(S_{k - 1})$.
\end{lemma}

\begin{proof}
Soit  $x \in D$. Alors  $\mathcal{O}_{D, x} = \mathcal{O}_{X, x}/(f)$  où  $f \in \mathcal{O}_{X, x}$  est un non-diviseur de zéro.
Par hypothèse, nous avons  $\text{depth}(\mathcal{O}_{X, x}) \geq \min(\dim(\mathcal{O}_{X, x}), k)$. Par Algèbre, Lemme  \ref{algebra-lemma-depth-drops-by-one}  nous avons
 $\text{depth}(\mathcal{O}_{D, x}) = \text{depth}(\mathcal{O}_{X, x}) - 1$  et par Algèbre, Lemme  \ref{algebra-lemma-one-equation}   $\dim(\mathcal{O}_{D, x}) = \dim(\mathcal{O}_{X, x}) - 1$. Il s'ensuit que
 $\text{depth}(\mathcal{O}_{D, x}) \geq \min(\dim(\mathcal{O}_{D, x}), k - 1)$  comme désiré.
\end{proof}

\begin{lemma}
\label{divisors-lemma-normal-effective-Cartier-divisor-S1}
Soit  $X$  un schéma normal localement noethérien. Soit  $D \subset X$  un
diviseur de Cartier effectif. Alors  $D$  est  $(S_1)$.
\end{lemma}

\begin{proof}
Par Propriétés, Lemme  \ref{properties-lemma-criterion-normal}  nous voyons que  $X$  est  $(S_2)$.
Ainsi, nous concluons par le Lemme  \ref{divisors-lemma-effective-Cartier-divisor-Sk}.
\end{proof}

\begin{lemma}
\label{divisors-lemma-weil-divisor-is-cartier-UFD}
Soit  $X$  un schéma localement noethérien. Soit  $D \subset X$  un
sous-schéma fermé intègre. Supposons que
\begin{enumerate}
\item $D$  ait une codimension  $1$  dans  $X$, et
\item $\mathcal{O}_{X, x}$  est un UFD pour tout  $x \in D$.
\end{enumerate}
Alors  $D$  est un diviseur de Cartier effectif.
\end{lemma}

\begin{proof}
Soit  $x \in D$  et soit  $A = \mathcal{O}_{X, x}$. Soit  $\mathfrak p \subset A$  correspondant au point
générique de  $D$. Alors  $A_\mathfrak p$  a une dimension  $1$  par
l'assumption (1). Ainsi  $\mathfrak p$  est un idéal premier de hauteur  $1$.
Puisque  $A$  est un UFD, cela implique que  $\mathfrak p = (f)$  pour un certain
 $f \in A$. Bien sûr,  $f$  est un non-diviseur de zéro et nous concluons
par le Lemme  \ref{divisors-lemma-effective-Cartier-in-points}.
\end{proof}

\begin{lemma}
\label{divisors-lemma-codim-1-part}
Soit  $X$  un schéma noethérien. Soit  $Z \subset X$  un sous-schéma fermé.
Supposons qu'il existe
\begin{enumerate}
\item une collection de diviseurs de Cartier effectifs intègres  $D_i \subset X$,
 $i \in I$
\item un sous-ensemble fermé  $Z' \subset X$  dont toutes les composantes irréductibles ont
une codimension  $\geq 2$  dans  $X$  (Topologie, Définition
 \ref{topology-definition-codimension})
\end{enumerate}
 de sorte que  $Z \subset Z' \cup \bigcup_{i \in I} D_i$  ensemblistement. Alors il
existe des entiers  $a_i \geq 0$  avec  $a_i = 0$  pour presque tous les
 $i \in I$  de sorte que le diviseur de Cartier effectif
$$
D = \sum a_i D_i
$$
est contenu dans  $Z$  et que le morphisme d'inclusion  $D \to Z$  est un
isomorphisme hors de la codimension  $2$  dans  $X$  (dans le sens
qu'il existe un ouvert  $U \subset Z$  tel que  $D \cap U \to Z \cap U$  est un isomorphisme et
que chaque composante irréductible de  $Z \setminus U$  a une codimension  $\geq 2$ 
dans  $X$). Lorsque  $Z$  est nulle part dense dans  $X$ , l'existence
de  $D_i$,  $i \in I$  et  $Z'$  est garantie si  $\mathcal{O}_{X, x}$  est un
UFD pour tous les  $x \in Z$  ou si  $X$  est régulier.
\end{lemma}

\begin{proof}
Soit  $\xi_i \in D_i$  le point générique et soit  $\mathcal{O}_i = \mathcal{O}_{X, \xi_i}$  l'anneau local qui est
un anneau de valuation discrète d'après le Lemme  \ref{divisors-lemma-integral-effective-Cartier-divisor-dvr}. Soit  $a_i \geq 0$ 
la valuation minimale d'un élément de  $\mathcal{I}_{Z, \xi_i} \subset \mathcal{O}_i$. Bien sûr  $a_i > 0$ 
seulement si  $D_i$  est une composante irréductible de  $Z$  et donc
 $a_i > 0$  seulement pour un nombre fini de  $i \in I$. Nous affirmons que le
diviseur de Cartier effectif  $D = \sum a_i D_i$  convient.

\medskip\noindent
À savoir, supposons que  $x \in X$. Soit  $A = \mathcal{O}_{X, x}$. Soient  $D_1, \ldots, D_n$  les
diviseurs mutuellement distincts  $D_i$  tels que  $x \in D_i$  et
 $a_i > 0$. Pour  $1 \leq i \leq n$ , soit  $f_i \in A$  une équation locale pour
 $D_i$. Alors  $f_i$  est un élément premier de  $A$  et
 $\mathcal{O}_i = A_{(f_i)}$. Soit  $I = \mathcal{I}_{Z, x} \subset A$  le germe du faisceau d'idéaux de  $Z$. Par
notre choix de  $a_i$ , nous avons  $I A_{(f_i)} = f_i^{a_i}A_{(f_i)}$. Nous affirmons que
 $I \subset (\prod_{i = 1, \ldots, n} f_i^{a_i})$.

\medskip\noindent
Preuve de l'assertion. L'application de localisation  $\varphi : A/(f_i) \to A_{(f_i)}/f_iA_{(f_i)}$  est injective puisque
l'idéal premier  $(f_i)$  est l'image réciproque de l'idéal maximal
 $f_iA_{(f_i)}$. Par induction sur  $n$ , nous déduisons que  $\varphi_n : A/(f_i^n)\to A_{(f_i)}/f_i^nA_{(f_i)}$  est
également injective. Puisque  $\varphi_{a_i}(I) = 0$, nous avons  $I \subset (f_i^{a_i})$. Ainsi, pour
tout  $x \in I$, nous pouvons écrire  $x = f_1^{a_1}x_1$  pour un certain  $x_1 \in A$.
Puisque  $D_1, \ldots, D_n$  sont deux à deux distincts,  $f_i$  est une unité dans
 $A_{(f_j)}$  pour  $i \not = j$. En comparant  $x$  et  $x_1$  en
 $A_{(f_i)}$  pour  $n \geq i > 1$, nous avons encore  $x_1 \in (f_i^{a_i})$. En répétant le
processus précédent, nous écrivons par induction  $x_i = f_{i + 1}^{a_{i + 1}}x_{i + 1}$  pour tout
 $n > i \geq 1$. En conclusion,  $x \in (\prod_{i = 1, \ldots n} f_i^{a_i})$  pour tout  $x \in I$  comme souhaité.

\medskip\noindent
L'affirmation montre que  $\mathcal{I}_Z \subset \mathcal{I}_D$, c'est-à-dire que  $D \subset Z$. De plus,
nous voyons également que  $D$  et  $Z$  coïncident aux points  $\xi_i$,
ce qui prouve que  $D \to Z$  est un isomorphisme hors d'un fermé dont les
composantes irréductibles ont codimension au moins  $2$  dans  $X$.

\medskip\noindent
Pour voir les dernières assertions, nous raisonnons comme suit. Un anneau local
régulier est un anneau factoriel (Compléments d'algèbre, Lemme  \ref{more-algebra-lemma-regular-local-UFD}), donc il suffit de
raisonner dans le cas factoriel. Dans ce cas, soient  $D_i$  les composantes
irréductibles de  $Z$  qui ont une codimension  $1$  dans
 $X$. Par le Lemme  \ref{divisors-lemma-weil-divisor-is-cartier-UFD} , chaque  $D_i$  est un diviseur de
Cartier effectif.
\end{proof}

\begin{lemma}
\label{divisors-lemma-codimension-1-is-effective-Cartier}
Soit  $Z \subset X$  un sous-schéma fermé d'un schéma noethérien. Supposons
\begin{enumerate}
\item $Z$  n'a pas de points immergés,
\item chaque composante irréductible de  $Z$  a une codimension  $1$  dans
 $X$.
\item chaque anneau local  $\mathcal{O}_{X, x}$,  $x \in Z$  est un UFD ou  $X$  est
régulier.
\end{enumerate}
Alors  $Z$  est un diviseur de Cartier effectif.
\end{lemma}

\begin{proof}
Soit  $D = \sum a_i D_i$  comme dans le Lemme  \ref{divisors-lemma-codim-1-part}  où  $D_i \subset Z$  sont les
composantes irréductibles de  $Z$. Si  $D \to Z$  n'est pas un
isomorphisme, alors  $\mathcal{O}_Z \to \mathcal{O}_D$  a un noyau non nul situé en codimension
 $\geq 2$. Cela signifierait que  $Z$  a des points immergés, ce qui
est interdit par l'hypothèse (1). Ainsi  $D \cong Z$  comme désiré.
\end{proof}

\begin{lemma}
\label{divisors-lemma-UFD-one-equation-CM}
Soit  $R$  un UFD noethérien. Soit  $I \subset R$  un idéal tel que
 $R/I$  n'a pas d'idéaux premiers immergés et tel que chaque idéal premier minimal
au-dessus de  $I$  a une hauteur  $1$. Alors  $I = (f)$  pour un
certain  $f \in R$.
\end{lemma}

\begin{proof}
D'après le Lemme  \ref{divisors-lemma-codimension-1-is-effective-Cartier} , le faisceau d'idéaux  $\tilde I$  est inversible
sur  $\Spec(R)$. D'après Compléments d'algèbre, Lemme  \ref{more-algebra-lemma-UFD-Pic-trivial} , il est engendré
par un seul élément.
\end{proof}

\begin{lemma}
\label{divisors-lemma-effective-Cartier-divisor-is-a-sum}
Soit  $X$  un schéma noethérien. Soit  $D \subset X$  un diviseur de Cartier
effectif. Supposons qu'il existe des diviseurs de Cartier effectifs intègres
 $D_i \subset X$  tels que  $D \subset \bigcup D_i$  ensemblistement. Alors
 $D = \sum a_i D_i$  pour un certain  $a_i \geq 0$. L'existence de  $D_i$  est
garantie si  $\mathcal{O}_{X, x}$  est un UFD pour tout  $x \in D$  ou si  $X$  est
régulier.
\end{lemma}

\begin{proof}
Choisissons  $a_i$  comme dans le Lemme  \ref{divisors-lemma-codim-1-part}  et définissez
 $D' = \sum a_i D_i$. Alors  $D' \to D$  est une inclusion de diviseurs de Cartier
effectifs qui est un isomorphisme en dehors de la codimension  $2$  sur
 $X$. Choisissons  $x \in X$. Posons  $A = \mathcal{O}_{X, x}$  et soient
 $f, f' \in A$  les non-diviseurs de zéro engendrant respectivement les idéaux de  $D, D'$  dans
 $A$. Alors  $f = gf'$  pour quelque  $g \in A$. De plus, pour chaque
idéal premier  $\mathfrak p$  de hauteur  $\leq 1$  de  $A$ , nous voyons
que  $g$  s'envoie sur une unité de  $A_\mathfrak p$. Cela implique que
 $g$  est une unité parce que les idéaux premiers minimaux au-dessus de
 $(g)$  ont une hauteur  $1$  (Algèbre, Lemme  \ref{algebra-lemma-minimal-over-1}).
\end{proof}

\begin{lemma}
\label{divisors-lemma-quasi-projective-Noetherian-pic-effective-Cartier}
\begin{slogan}
Sur un schéma projectif, tout faisceau inversible admet une section méromorphe
régulière.
\end{slogan}
Soit  $X$  un schéma noethérien qui possède un faisceau inversible ample.
Alors chaque  $\mathcal{O}_X$-module inversible est isomorphe à
$$
\mathcal{O}_X(D - D') =
\mathcal{O}_X(D) \otimes_{\mathcal{O}_X} \mathcal{O}_X(D')^{\otimes -1}
$$
pour certains diviseurs de Cartier effectifs  $D, D'$  dans  $X$. De
plus, étant donné un sous-ensemble fini  $E \subset X$ , nous pouvons choisir
 $D, D'$  tel que  $E \cap D = \emptyset$  et  $E \cap D' = \emptyset$. Si  $X$  est
quasi-affine, alors nous pouvons choisir  $D' = \emptyset$.
\end{lemma}

\begin{proof}
Soit  $x_1, \ldots, x_n$  les points associés de  $X$  (Lemme  \ref{divisors-lemma-finite-ass}).

\medskip\noindent
Si  $X$  est quasi-affine et  $\mathcal{N}$  est un module  $\mathcal{O}_X$
inversible quelconque, alors nous pouvons choisir une section  $t$  de
 $\mathcal{N}$  qui ne s'annule en aucun des points de  $E \cup \{x_1, \ldots, x_n\}$, voir Propriétés,
Lemme  \ref{properties-lemma-quasi-affine-invertible-nonvanishing-section}. Ensuite,  $t$  est une section régulière de
 $\mathcal{N}$  selon le Lemme  \ref{divisors-lemma-regular-section-associated-points}. Par conséquent,  $\mathcal{N} \cong \mathcal{O}_X(D)$  où
 $D = Z(t)$  est le diviseur de Cartier effectif correspondant à  $t$,
voir Lemme  \ref{divisors-lemma-characterize-OD}. Puisque  $E \cap D = \emptyset$  par construction, nous avons terminé
dans ce cas.

\medskip\noindent
En revenant au cas général, soit  $\mathcal{L}$  un faisceau inversible ample sur
 $X$. Il existe un  $m > 0$  et une section  $s \in \Gamma(X, \mathcal{L}^{\otimes m})$  tels que
 $X_s$  soit affine et tels que  $E \cup \{x_1, \ldots, x_n\} \subset X_s$  (Propriétés, Lemme
 \ref{properties-lemma-ample-finite-set-in-principal-affine}).

\medskip\noindent
Soit  $\mathcal{N}$  un  $\mathcal{O}_X$-module inversible arbitraire. Dans le cas
quasi-affine, nous pouvons trouver une section  $t \in \mathcal{N}(X_s)$  qui ne s'annule en
aucun point de  $E \cup \{x_1, \ldots, x_n\}$. D'après les Propriétés, Lemme  \ref{properties-lemma-invert-s-sections} , nous
voyons que pour un certain  $e \geq 0$ , la section  $s^e|_{X_s} t$  s'étend à une
section globale  $\tau$  de  $\mathcal{L}^{\otimes me} \otimes \mathcal{N}$. Ainsi,  $\mathcal{L}^{\otimes me} \otimes \mathcal{N}$  et  $\mathcal{L}^{\otimes me}$ 
sont toutes deux des faisceaux inversibles qui possèdent des sections globales ne
s'annulant en aucun point de  $E \cup \{x_1, \ldots, x_n\}$. Par conséquent, ce sont des sections
régulières d'après le Lemme  \ref{divisors-lemma-regular-section-associated-points}. Donc  $\mathcal{L}^{\otimes me} \otimes \mathcal{N} \cong \mathcal{O}_X(D)$  et  $\mathcal{L}^{\otimes me} \cong \mathcal{O}_X(D')$  pour
certains diviseurs de Cartier effectifs  $D$  et  $D'$, voir le
Lemme  \ref{divisors-lemma-characterize-OD}. Par construction  $E \cap D = \emptyset$  et  $E \cap D' = \emptyset$  et la
démonstration est terminée.
\end{proof}

\begin{lemma}
\label{divisors-lemma-wedge-product-ses}
Soit  $X$  un schéma régulier intègre de dimension  $2$. Soit
 $i : D \to X$  l'immersion d'un diviseur de Cartier effectif. Soit  $\mathcal{F} \to \mathcal{F}' \to i_*\mathcal{G} \to 0$  une
suite exacte de  $\mathcal{O}_X$-modules cohérents. Supposons que
\begin{enumerate}
\item $\mathcal{F}, \mathcal{F}'$  soient localement libres de rang  $r$  sur un ouvert non vide
de  $X$,
\item $D$  est un schéma intègre,
\item $\mathcal{G}$  est un  $\mathcal{O}_D$-module localement libre de rang
 $s$.
\end{enumerate}
Alors  $\mathcal{L} = (\wedge^r\mathcal{F})^{**}$  et  $\mathcal{L}' = (\wedge^r \mathcal{F}')^{**}$  sont des  $\mathcal{O}_X$-modules inversibles et
 $\mathcal{L}' \cong \mathcal{L}(k D)$  pour un certain  $k \in \{0, \ldots, \min(s, r)\}$.
\end{lemma}

\begin{proof}
La première assertion découle du Lemme  \ref{divisors-lemma-reflexive-over-regular-dim-2}  car l'hypothèse (1) implique
que  $\mathcal{L}$  et  $\mathcal{L}'$  ont un rang  $1$. Comme prendre
 $\wedge^r$  et les doubles duaux sont des foncteurs, nous obtenons donc une
application canonique  $\sigma : \mathcal{L} \to \mathcal{L}'$  qui est un isomorphisme sur l'ouvert non vide
de (1), et donc non nulle. Pour terminer la démonstration, il suffit de voir que
 $\sigma$  considéré comme une section globale de  $\mathcal{L}' \otimes \mathcal{L}^{\otimes -1}$  ne s'annule en
aucun point de codimension de  $X$, sauf au point générique de
 $D$  et là avec un ordre d'annulation au plus égal à  $\min(s, r)$.

\medskip\noindent
Traduit en algèbre, nous arrivons au problème suivant: Soit  $(A, \mathfrak m, \kappa)$  un anneau
de valuation discrète avec corps des fractions  $K$. Soit  $M \to M' \to N \to 0$  une
suite exacte de  $A$-modules finis avec  $\dim_K(M \otimes K) = \dim_K(M' \otimes K) = r$  et avec  $N \cong \kappa^{\oplus s}$.
Montrer que l'application induite  $L = \wedge^r(M)^{**} \to L' = \wedge^r(M')^{**}$  s'annule d'ordre au plus
 $\min(s, r)$. Nous utiliserons le théorème de structure pour les modules sur
 $A$, voir Compléments d'algèbre, Lemme  \ref{more-algebra-lemma-generalized-valuation-ring-modules}  ou  \ref{more-algebra-lemma-modules-PID}. Diviser
un  $A$-module fini par un sous-module de torsion ne change pas le double
dual. Ainsi, nous pouvons remplacer  $M$  par  $M/M_{tors}$  et  $M'$ 
par  $M'/\Im(M_{tors} \to M')$  et supposer que  $M$  est sans torsion. Alors  $M \to M'$ 
est injectif et  $M'_{tors} \to N$  est injectif. Par conséquent, nous pouvons remplacer
 $M'$  par  $M'/M'_{tors}$  et  $N$  par  $N/M'_{tors}$. Ainsi, nous
réduisons au cas où  $M$  et  $M'$  sont libres de rang  $r$ 
et  $N \cong \kappa^{\oplus s}$. Dans ce cas,  $\sigma$  est le déterminant de  $M \to M'$  et
s'annule d'ordre  $s$ , par exemple d'après Algèbre, Lemme  \ref{algebra-lemma-order-vanishing-determinant}.
\end{proof}









\section{Compléments d'ouverts affines}
\label{divisors-section-complement-affine-open}

\noindent
Dans cette section, nous discutons du résultat selon lequel le complémentaire d'un
ouvert affine dans une variété a une codimension pure  $1$.

\begin{lemma}
\label{divisors-lemma-affine-punctured-spec}
Soit  $(A, \mathfrak m)$  un anneau local noethérien. Le spectre ponctué  $U = \Spec(A) \setminus \{\mathfrak m\}$  de
 $A$  est affine si et seulement si  $\dim(A) \leq 1$.
\end{lemma}

\begin{proof}
Si  $\dim(A) = 0$, alors  $U$  est vide et donc affine (égale au spectre de
l'anneau  $0$ ). Si  $\dim(A) = 1$, alors nous pouvons choisir un élément
 $f \in \mathfrak m$  qui n'est contenu dans aucun des idéaux premiers minimaux, en nombre fini,
de  $A$  (Algèbre, Lemme  \ref{algebra-lemma-Noetherian-irreducible-components}  et  \ref{algebra-lemma-silly}). Alors  $U = \Spec(A_f)$ 
est affine.

\medskip\noindent
Le cas inverse est plus intéressant. Nous allons donner une preuve quelque peu non
standard et discuter de l'argument standard dans une remarque ci-dessous.
Supposons que  $U = \Spec(B)$  soit affine. Puisque l'affinité et la dimension ne sont
pas affectées en passant à la réduction, nous pouvons remplacer  $A$  par
le quotient par son idéal d'éléments nilpotents et supposer que  $A$  est
réduit. Posons  $Q = B/A$  considéré comme un  $A$-module. Le support de
 $Q$  est  $\{\mathfrak m\}$  car  $A_\mathfrak p = B_\mathfrak p$  pour tous les idéaux premiers non
maximaux  $\mathfrak p$  de  $A$. Nous pouvons supposer  $\dim(A) \geq 1$, donc
comme ci-dessus nous pouvons choisir  $f \in \mathfrak m$  qui n'est contenu dans aucun des
idéaux minimaux de  $A$. Puisque  $A$  est réduit, cela implique
que  $f$  est un non-diviseur de zéro dans  $A$. Le même
argument montre que  $f : B \to B$  est injectif car  $\Spec(B)$  est ouvert dans
 $\Spec(A)$. Notons que  $\dim(A/fA) = \dim(A) - 1$  d'après Algèbre, Lemme  \ref{algebra-lemma-one-equation}. En
appliquant le lemme du serpent à la multiplication par  $f$  sur la suite
exacte  $0 \to A \to B \to Q \to 0$ , nous obtenons
$$
0 \to Q[f] \to A/fA \to B/fB \to Q/fQ \to 0
$$
où  $Q[f] = \Ker(f : Q \to Q)$. Cela implique que  $Q[f]$  est un  $A$-module fini.
Puisque le support de  $Q[f]$  est  $\{\mathfrak m\}$ , nous voyons  $l = \text{length}_A(Q[f]) < \infty$ 
(Algèbre, Lemme  \ref{algebra-lemma-support-point}). Posons  $l_n = \text{length}_A(Q[f^n])$. La suite exacte
$$
0 \to Q[f^n] \to Q[f^{n + 1}] \xrightarrow{f^n} Q[f]
$$
montre par induction que  $l_n < \infty$  et que  $l_n \leq l_{n + 1}$. En considérant la suite
exacte
$$
0 \to Q[f] \to Q[f^{n + 1}] \xrightarrow{f} Q[f^n] \to Q/fQ
$$
et nous voyons que l'image de  $Q[f^n]$  dans  $Q/fQ$  a une longueur de
 $l_n - l_{n + 1} + l \leq l$. Puisque  $Q = \bigcup Q[f^n]$  nous trouvons que la longueur de  $Q/fQ$ 
est au plus  $l$, c'est-à-dire bornée. Ainsi  $Q/fQ$  est un
 $A$-module fini. Par conséquent,  $A/fA \to B/fB$  est un morphisme d'anneaux
fini, en particulier il induit une application fermée sur les spectres (Algèbre,
lemmes  \ref{algebra-lemma-integral-going-up}  et  \ref{algebra-lemma-going-up-closed}). D'autre part,  $\Spec(B/fB)$  est le spectre
ponctué de  $\Spec(A/fA)$. Ceci est une contradiction à moins de  $\Spec(B/fB) = \emptyset$ , ce qui
signifie que  $\dim(A/fA) = 0$  comme souhaité.
\end{proof}

\begin{remark}
\label{divisors-remark-affine-punctured-spectrum-standard-proof}
Si  $(A, \mathfrak m)$  est un anneau intègre local noethérien normal de dimension  $\geq 2$ 
et  $U$  est le spectre ponctué de  $A$, alors  $\Gamma(U, \mathcal{O}_U) = A$. Cette
version algébrique du théorème de Hartogs découle du fait que  $A = \bigcap_{\text{height}(\mathfrak p) = 1} A_\mathfrak p$  que
nous avons vu en Algèbre, Lemme  \ref{algebra-lemma-normal-domain-intersection-localizations-height-1}. Ainsi, dans ce cas,  $U$  ne
peut pas être affine (car cela obligerait  $\mathfrak m$  à être un point de
 $U$). Cela est souvent utilisé comme point de départ de la preuve du
Lemme  \ref{divisors-lemma-affine-punctured-spec}. Pour réduire le cas d'un anneau local noethérien général à ce
cas, nous complétons d'abord (pour obtenir un anneau local de Nagata), puis
remplaçons  $A$  par  $A/\mathfrak q$  pour un idéal premier minimal approprié,
puis normalisons. Chacune de ces étapes ne change pas la dimension et nous
obtenons une contradiction. Vous pouvez sauter l'étape de la complétion, mais
alors la normalisation n'est généralement pas un anneau intègre noethérien. Cependant,
c'est toujours un anneau de Krull de la même dimension (cela se prouve en
utilisant Krull-Akizuki) et on peut appliquer le même argument.
\end{remark}

\begin{remark}
\label{divisors-remark-affine-puctured-spectrum-general}
Il n'est pas clair comment caractériser les anneaux locaux non noethériens
 $(A, \mathfrak m)$  dont le spectre ponctué est affine. Un tel anneau possède un idéal
de type fini  $I$  avec  $\mathfrak m = \sqrt{I}$. Bien sûr, si nous pouvons prendre
 $I$  généré par  $1$  élément, alors  $A$  a un spectre
ponctué affine; cela donne de nombreux exemples non noethériens. Inversement, il
découle de l'argument dans la démonstration du Lemme  \ref{divisors-lemma-affine-punctured-spec}  qu'un tel anneau
ne peut pas posséder un non-diviseur de zéro  $f \in \mathfrak m$  avec  $H^0_I(A/fA) = 0$  (donc
 $A$  ne peut pas avoir une suite régulière de longueur  $2$). De
plus, il en va de même pour tout anneau  $A'$  qui est la cible d'un
homomorphisme local d'anneaux locaux  $A \to A'$  tel que  $\mathfrak m_{A'} = \sqrt{\mathfrak mA'}$.
\end{remark}

\begin{lemma}
\label{divisors-lemma-complement-affine-open-immersion}
\begin{reference}
\cite[EGA IV, Corollaire 21.12.7]{EGA4}
\end{reference}
Soit  $X$  un schéma localement noethérien. Soit  $U \subset X$  un
sous-schéma ouvert tel que le morphisme d'inclusion  $U \to X$  soit affine. Pour
tout point générique  $\xi$  d'une composante irréductible de  $X \setminus U$ ,
l'anneau local  $\mathcal{O}_{X, \xi}$  a une dimension  $\leq 1$. Si  $U$  est
dense ou si  $\xi$  se trouve dans l'adhérence de  $U$, alors
 $\dim(\mathcal{O}_{X, \xi}) = 1$.
\end{lemma}

\begin{proof}
Puisque  $\xi$  est un point générique de  $X \setminus U$, nous voyons que
$$
U_\xi = U \times_X \Spec(\mathcal{O}_{X, \xi}) \subset
\Spec(\mathcal{O}_{X, \xi})
$$
est le spectre ponctué de  $\mathcal{O}_{X, \xi}$  (indication: utiliser Schémas, Lemme
 \ref{schemes-lemma-specialize-points}). Comme  $U \to X$  est affine, nous voyons que  $U_\xi \to \Spec(\mathcal{O}_{X, \xi})$  est
affine (Morphismes, Lemme  \ref{morphisms-lemma-base-change-affine}) et nous concluons que  $U_\xi$  est
affine. Donc  $\dim(\mathcal{O}_{X, \xi}) \leq 1$  d'après le Lemme  \ref{divisors-lemma-affine-punctured-spec}. Si  $\xi \in \overline{U}$, alors il
existe une spécialisation  $\eta \to \xi$  où  $\eta \in U$  (il suffit de prendre
 $\eta$  un point générique d'une composante irréductible de  $\overline{U}$  qui
contient  $\xi$; puisque  $\overline{U}$  est localement noethérien, donc
localement a un nombre fini de composantes irréductibles, nous voyons que
 $\eta \in U$). Alors  $\eta \in \Spec(\mathcal{O}_{X, \xi})$  et nous voyons que la dimension ne peut pas être
 $0$.
\end{proof}

\begin{lemma}
\label{divisors-lemma-complement-affine-open}
Soit  $X$  un schéma séparé et localement noethérien. Soit  $U \subset X$  un
ouvert affine. Pour chaque point générique  $\xi$  d'une composante
irréductible de  $X \setminus U$ , l'anneau local  $\mathcal{O}_{X, \xi}$  a pour dimension
 $\leq 1$. Si  $U$  est dense ou si  $\xi$  est dans l'adhérence de
 $U$, alors  $\dim(\mathcal{O}_{X, \xi}) = 1$.
\end{lemma}

\begin{proof}
Ceci découle du Lemme  \ref{divisors-lemma-complement-affine-open-immersion}  car le morphisme  $U \to X$  est affine par
Morphismes, Lemme  \ref{morphisms-lemma-affine-permanence}.
\end{proof}

\noindent
Le lemme suivant peut parfois être utilisé pour produire des diviseurs de Cartier
effectifs.

\begin{lemma}
\label{divisors-lemma-complement-open-affine-effective-cartier-divisor}
Soit  $X$  un schéma séparé noethérien. Soit  $U \subset X$  un ouvert affine
dense. Si  $\mathcal{O}_{X, x}$  est un UFD pour tout  $x \in X \setminus U$, alors il existe un
diviseur de Cartier effectif  $D \subset X$  avec  $U = X \setminus D$.
\end{lemma}

\begin{proof}
Puisque  $X$  est noethérien, le complément  $X \setminus U$  a un nombre fini
de composantes irréductibles  $D_1, \ldots, D_r$  (Propriétés, Lemme  \ref{properties-lemma-Noetherian-irreducible-components}  appliqué
à la structure de sous-schéma réduit induite sur  $X \setminus U$). Chaque
 $D_i \subset X$  a une codimension  $1$  d'après le Lemme  \ref{divisors-lemma-complement-affine-open}  (et
Propriétés, Lemme  \ref{properties-lemma-codimension-local-ring}). Ainsi,  $D_i$  est un diviseur de Cartier
effectif d'après le Lemme  \ref{divisors-lemma-weil-divisor-is-cartier-UFD}. Par conséquent, nous pouvons prendre
 $D = D_1 + \ldots + D_r$.
\end{proof}

\begin{lemma}
\label{divisors-lemma-complement-open-affine-effective-cartier-divisor-bis}
Soit  $X$  un schéma noethérien avec diagonale affine. Soit  $U \subset X$  un
ouvert affine dense. Si  $\mathcal{O}_{X, x}$  est un UFD pour tout  $x \in X \setminus U$, alors il
existe un diviseur de Cartier effectif  $D \subset X$  avec  $U = X \setminus D$.
\end{lemma}

\begin{proof}
Puisque  $X$  est noethérien, le complément  $X \setminus U$  a un nombre fini
de composantes irréductibles  $D_1, \ldots, D_r$  (Propriétés, Lemme  \ref{properties-lemma-Noetherian-irreducible-components}  appliqué
à la structure de sous-schéma réduit induite sur  $X \setminus U$). Nous considérons
 $D_i$  comme un sous-schéma fermé réduit de  $X$. Soit  $X = \bigcup_{j \in J} X_j$ 
un recouvrement ouvert affine de  $X$. Pour tout  $j$  dans
 $J$, posons  $U_j = U \cap X_j$. Puisque  $X$  a une diagonale affine, le
schéma
$$
U_j = X \times_{(X \times X)} (U \times X_j)
$$
est affine. Par conséquent, comme  $X_j$  est séparé, il s'ensuit du Lemme
 \ref{divisors-lemma-complement-open-affine-effective-cartier-divisor}  et de sa démonstration que pour tous  $j \in J$  et  $1 \leq i \leq r$ ,
l'intersection  $D_i \cap X_j$  est soit vide, soit un diviseur de Cartier effectif
dans  $X_j$. Ainsi  $D_i \subset X$  est un diviseur de Cartier effectif (car il
s'agit d'une propriété locale). Par conséquent, nous pouvons prendre  $D = D_1 + \ldots + D_r$.
\end{proof}

\begin{lemma}
\label{divisors-lemma-regular-ample-family}
Soit  $X$  un schéma quasi-compact, régulier avec diagonale affine. Alors
 $X$  possède une famille abondante de modules inversibles (Morphismes,
Définition  \ref{morphisms-definition-family-ample-invertible-modules}).
\end{lemma}

\begin{proof}
Remarquons que  $X$  est une réunion finie disjointe de schémas intègres
(Propriétés, Lemmes  \ref{properties-lemma-regular-normal}  et  \ref{properties-lemma-normal-Noetherian}). Ainsi, on peut supposer que
 $X$  est intégral ainsi que Noethérien, régulier et ayant une diagonale
affine. Soit  $x \in X$. Choisissons un voisinage ouvert affine  $U \subset X$  de
 $x$. Puisque  $X$  est intégral,  $U$  est dense dans
 $X$. D'après Compléments d'algèbre, Lemme  \ref{more-algebra-lemma-regular-local-UFD} , les anneaux locaux de
 $X$  sont des UFD. Par conséquent, d'après le Lemme  \ref{divisors-lemma-complement-open-affine-effective-cartier-divisor-bis} , nous
pouvons trouver un diviseur de Cartier effectif  $D \subset X$  dont le complément
est  $U$. Alors la section canonique  $s = 1_D$  de  $\mathcal{L} = \mathcal{O}_X(D)$, voir la
Définition  \ref{divisors-definition-invertible-sheaf-effective-Cartier-divisor}, s'annule exactement le long de  $D$  donc
 $U = X_s$. Ainsi, les deux conditions dans Morphismes, Définition  \ref{morphisms-definition-family-ample-invertible-modules} 
sont satisfaites et nous avons terminé.
\end{proof}




\section{Normes}
\label{divisors-section-norms}

\noindent
Soit  $\pi : X \to Y$  un morphisme fini de schémas et soit  $d \geq 1$  un entier.
Disons qu'il existe une  {\it  norme de degré  $d$  pour
 $\pi$}\footnote{Ceci est une notation non
standard.}  s'il existe une application multiplicative
$$
\text{Norm}_\pi : \pi_*\mathcal{O}_X \to \mathcal{O}_Y
$$
de faisceaux telle que
\begin{enumerate}
\item la composition  $\mathcal{O}_Y \xrightarrow{\pi^\sharp} \pi_*\mathcal{O}_X
\xrightarrow{\text{Norm}_\pi} \mathcal{O}_Y$  égale  $g \mapsto g^d$, et
\item pour  $V \subset Y$  ouvert, si  $f \in \mathcal{O}_X(\pi^{-1}V)$  s'annule en  $x \in \pi^{-1}(V)$, alors
 $\text{Norm}_\pi(f)$  s'annule en  $\pi(x)$.
\end{enumerate}
Nous observons que la condition (1) oblige  $\pi$  à être surjective. Puisque
 $\text{Norm}_\pi$  est multiplicatif, il envoie les unités sur des unités; par
conséquent, étant donné  $y \in Y$, si  $f$  est une fonction régulière sur
 $X$  définie et non nulle en tout  $x \in X$  tel que  $\pi(x) = y$,
alors  $\text{Norm}_\pi(f)$  est définie et ne s'annule pas en  $y$. Cela reste vrai
sans supposer (2); en fait, les constructions dans cette section ne nécessiteront que
la condition (1) et seules certaines propriétés d'annulation (qui sont utilisées
en particulier dans la démonstration du Lemme  \ref{divisors-lemma-norm-ample}) nécessiteront la
propriété (2).

\begin{lemma}
\label{divisors-lemma-finite-trivialize-invertible-upstairs}
Soit  $\pi : X \to Y$  un morphisme fini de schémas. Soit  $\mathcal{L}$  un
 $\mathcal{O}_X$-module inversible. Soit  $y \in Y$. Il existe un voisinage ouvert
 $V \subset Y$  de  $y$  tel que  $\mathcal{L}|_{\pi^{-1}(V)}$  soit trivial.
\end{lemma}

\begin{proof}
De toute évidence, nous pouvons supposer que  $Y$  et donc  $X$ 
sont affines. Puisque  $\pi$  est fini, la fibre  $\pi^{-1}(\{y\})$  au-dessus de
 $y$  est finie. Puisque  $X$  est affine, nous pouvons choisir
 $s \in \Gamma(X, \mathcal{L})$  ne s'annulant en aucun point de  $\pi^{-1}(\{y\})$. Cela découle des
Propriétés, Lemme  \ref{properties-lemma-quasi-affine-invertible-nonvanishing-section} , mais nous donnons également un argument direct. À
savoir, nous pouvons choisir un ensemble fini  $E \subset X$  de points fermés tel
que chaque  $x \in \pi^{-1}(\{y\})$  se spécialise en un certain point de  $E$. Pour
 $x \in E$ , notons  $i_x : x \to X$  l'immersion fermée. Alors  $\mathcal{L} \to \bigoplus_{x \in E} i_{x, *}i_x^*\mathcal{L}$  est une
application surjective de  $\mathcal{O}_X$-modules quasi-cohérents, et donc
l'application
$$
\Gamma(X, \mathcal{L}) \to
\bigoplus\nolimits_{x \in E} \mathcal{L}_x/\mathfrak m_x\mathcal{L}_x
$$
est surjective (puisque la prise de sections globales est un foncteur exact sur la
catégorie des  $\mathcal{O}_X$-modules quasi-cohérents, voir Schémas, Lemme
 \ref{schemes-lemma-equivalence-quasi-coherent}). Ainsi, nous pouvons trouver un  $s \in \Gamma(X, \mathcal{L})$  ne s'annulant en aucun
point se spécialisant en un point de  $E$. Alors  $X_s \subset X$  est un
voisinage ouvert de  $\pi^{-1}(\{y\})$. Puisque  $\pi$  est fini, donc fermé, nous
en concluons qu'il existe un voisinage ouvert  $V \subset Y$  de  $y$  dont
l'image réciproque est contenue dans  $X_s$  comme souhaité.
\end{proof}

\begin{lemma}
\label{divisors-lemma-norm-invertible}
Soit  $\pi : X \to Y$  un morphisme fini de schémas. S'il existe une norme de degré
 $d$  pour  $\pi$, alors il existe un homomorphisme de groupes
abéliens
$$
\text{Norm}_\pi : \Pic(X) \to \Pic(Y)
$$
tel que  $\text{Norm}_\pi(\pi^*\mathcal{N}) \cong \mathcal{N}^{\otimes d}$  pour tous les  $\mathcal{O}_Y$-modules inversibles  $\mathcal{N}$.
\end{lemma}

\begin{proof}
Nous utiliserons la correspondance entre les classes d'isomorphisme de modules
 $\mathcal{O}_X$ inversibles et les éléments de  $H^1(X, \mathcal{O}_X^*)$  donnée dans Cohomologie,
Lemme  \ref{cohomology-lemma-h1-invertible}  sans autre mention. Nous expliquons comment prendre la norme
d'un module  $\mathcal{O}_X$ inversible  $\mathcal{L}$. À savoir, d'après le Lemme
 \ref{divisors-lemma-finite-trivialize-invertible-upstairs} , il existe un recouvrement ouvert  $Y = \bigcup V_j$  tel que  $\mathcal{L}|_{\pi^{-1}V_j}$ 
soit trivial. Choisissons une section génératrice  $s_j \in \mathcal{L}(\pi^{-1}V_j)$  pour chaque
 $j$. Sur les intersections  $\pi^{-1}V_j \cap \pi^{-1}V_{j'}$ , nous pouvons écrire
$$
s_j = u_{jj'} s_{j'}
$$
pour un  $u_{jj'} \in \mathcal{O}^*_X(\pi^{-1}V_j \cap \pi^{-1}V_{j'})$ unique. Ainsi, nous pouvons considérer les éléments
$$
v_{jj'} = \text{Norm}_\pi(u_{jj'}) \in \mathcal{O}_Y^*(V_j \cap V_{j'})
$$
Ces éléments satisfont la condition de cocycle (car  $u_{jj'}$  le fait et
 $\text{Norm}_\pi$  est multiplicatif) et définissent donc un  $\mathcal{O}_Y$-module
inversible. Nous omettons la vérification que: cela est bien défini, additif sur
les groupes de Picard, et satisfait la propriété  $\text{Norm}_\pi(\pi^*\mathcal{N}) \cong \mathcal{N}^{\otimes d}$  pour tous les
 $\mathcal{O}_Y$-modules inversibles  $\mathcal{N}$.
\end{proof}

\begin{lemma}
\label{divisors-lemma-norm-map-invertible}
Soit  $\pi : X \to Y$  un morphisme fini de schémas. Supposons qu'il existe une norme
de degré  $d$  pour  $\pi$. Pour tout morphisme  $\mathcal{O}_X$-linéaire
 $\varphi : \mathcal{L} \to \mathcal{L}'$  de  $\mathcal{O}_X$-modules inversibles, il existe un
morphisme  $\mathcal{O}_Y$-linéaire
$$
\text{Norm}_\pi(\varphi) :
\text{Norm}_\pi(\mathcal{L})
\longrightarrow
\text{Norm}_\pi(\mathcal{L}')
$$
avec  $\text{Norm}_\pi(\mathcal{L})$,  $\text{Norm}_\pi(\mathcal{L}')$  comme dans le Lemme  \ref{divisors-lemma-norm-invertible}. De plus, pour
 $y \in Y$ , les éléments suivants sont équivalents
\begin{enumerate}
\item $\varphi$  est nul en un point de  $x \in X$  avec  $\pi(x) = y$, et
\item $\text{Norm}_\pi(\varphi)$  est nul à  $y$.
\end{enumerate}
\end{lemma}

\begin{proof}
Nous choisissons un recouvrement ouvert  $Y = \bigcup V_j$  tel que  $\mathcal{L}$  et
 $\mathcal{L}'$  soient triviaux sur les ouverts  $\pi^{-1}V_j$. Cela est possible
d'après le Lemme  \ref{divisors-lemma-finite-trivialize-invertible-upstairs}. Choisissons des sections génératrices  $s_j$ 
et  $s'_j$  de  $\mathcal{L}$  et  $\mathcal{L}'$  sur les ouverts  $\pi^{-1}V_j$.
Alors  $\varphi(s_j) = f_js'_j$  pour certains  $f_j \in \mathcal{O}_X(\pi^{-1}V_j)$. Définissons  $\text{Norm}_\pi(\varphi)$  comme étant
la multiplication par  $\text{Norm}_\pi(f_j)$  sur  $V_j$. Un calcul simple impliquant
les cocycles utilisés pour construire  $\text{Norm}_\pi(\mathcal{L})$,  $\text{Norm}_\pi(\mathcal{L}')$  dans la
démonstration du Lemme  \ref{divisors-lemma-norm-invertible}  montre que cela définit une application comme
indiqué dans le lemme. L'énoncé final découle de la condition (2) dans la
définition d'une application norme de degré  $d$. Certains détails sont
omis.
\end{proof}

\begin{lemma}
\label{divisors-lemma-norm-ample}
Soit  $\pi : X \to Y$  un morphisme fini de schémas. Supposons que  $X$  possède
un faisceau inversible ample et qu'il existe une norme de degré  $d$  pour
 $\pi$. Alors  $Y$  possède un faisceau inversible ample.
\end{lemma}

\begin{proof}
Soit  $\mathcal{L}$  le faisceau inversible ample sur  $X$  qui nous est donné
par hypothèse. Nous prouverons que  $\mathcal{N} = \text{Norm}_\pi(\mathcal{L})$  est ample sur  $Y$.

\medskip\noindent
Puisque  $X$  est quasi-compact (Propriétés, Définition  \ref{properties-definition-ample}) et
que  $X \to Y$  est surjectif (par l'existence de  $\text{Norm}_\pi$), nous voyons que
 $Y$  est quasi-compact. Soit  $y \in Y$  un point. Pour terminer la
preuve, nous montrerons qu'il existe une section  $t$  d'une certaine
puissance tensorielle positive de  $\mathcal{N}$  qui ne s'annule pas en  $y$ 
de sorte que  $Y_t$  soit affine. Pour ce faire, choisissons un voisinage
ouvert affine  $V \subset Y$  de  $y$. Choisissons  $n \gg 0$  et une section
 $s \in \Gamma(X, \mathcal{L}^{\otimes n})$  telle que
$$
\pi^{-1}(\{y\}) \subset X_s \subset \pi^{-1}V
$$
par Propriétés, Lemme  \ref{properties-lemma-ample-finite-set-in-principal-affine}. Ensuite,  $t = \text{Norm}_\pi(s)$  est une section de
 $\mathcal{N}^{\otimes n}$  qui ne s'annule pas en  $x$  et avec  $Y_t \subset V$, voir Lemme
 \ref{divisors-lemma-norm-map-invertible}. Alors  $Y_t$  est affine par Propriétés, Lemme  \ref{properties-lemma-affine-cap-s-open}.
\end{proof}

\begin{lemma}
\label{divisors-lemma-norm-quasi-affine}
Soit  $\pi : X \to Y$  un morphisme fini de schémas. Supposons que  $X$  soit
quasi-affine et qu'il existe une norme de degré  $d$  pour  $\pi$.
Alors  $Y$  est quasi-affine.
\end{lemma}

\begin{proof}
D'après les Propriétés, Lemme  \ref{properties-lemma-quasi-affine-O-ample} , nous voyons que  $\mathcal{O}_X$  est un
faisceau inversible ample sur  $X$. La démonstration du Lemme  \ref{divisors-lemma-norm-ample} 
montre que  $\text{Norm}_\pi(\mathcal{O}_X) = \mathcal{O}_Y$  est un  $\mathcal{O}_Y$-module inversible ample. Par
conséquent, les Propriétés, Lemme  \ref{properties-lemma-quasi-affine-O-ample}  montrent que  $Y$  est
quasi-affine.
\end{proof}

\begin{lemma}
\label{divisors-lemma-finite-locally-free-has-norm}
Soit  $\pi : X \to Y$  un morphisme fini localement libre de degré  $d \geq 1$. Alors
il existe une norme canonique de degré  $d$  dont la formation commute avec
tout changement de base.
\end{lemma}

\begin{proof}
Soit  $V \subset Y$  un ouvert affine tel que  $(\pi_*\mathcal{O}_X)|_V$  soit libre de rang
 $d$. En choisissant une base, nous obtenons un isomorphisme
$$
\mathcal{O}_V^{\oplus d} \cong (\pi_*\mathcal{O}_X)|_V
$$
Pour tout  $f \in \pi_*\mathcal{O}_X(V) = \mathcal{O}_X(\pi^{-1}(V))$, la multiplication par  $f$  définit un endomorphisme
 $\mathcal{O}_V$-linéaire  $m_f$  du module libre affiché. Ainsi, nous
obtenons une matrice  $d \times d$   $M_f \in \text{Mat}(d \times d, \mathcal{O}_Y(V))$  et nous pouvons définir
$$
\text{Norm}_\pi(f) = \det(M_f)
$$
Comme le déterminant d'une matrice est indépendant du choix d'une base,
nous voyons que cela est bien défini, ce qui signifie également que cette
construction se recollera pour donner une application globale comme souhaité. La
compatibilité avec le changement de base découle directement de la construction.

\medskip\noindent
La propriété (1) découle du fait que le déterminant d'une matrice diagonale
 $d \times d$  avec des entrées  $g, g, \ldots, g$  est  $g^d$. Pour voir la
propriété (2), nous pouvons changer de base et supposer que  $Y$  est le
spectre d'un corps  $k$. Puis  $X = \Spec(A)$  avec  $A$  une
 $k$-algèbre avec  $\dim_k(A) = d$. S'il existe un  $x \in X$  tel que
 $f \in A$  s'annule en  $x$, alors il existe une application
 $A \to \kappa$  dans un corps telle que  $f$  s'envoie à zéro dans
 $\kappa$. Alors  $f : A \to A$  ne peut pas être surjectif, donc  $\det(f : A \to A) = 0$ 
comme souhaité.
\end{proof}

\begin{lemma}
\label{divisors-lemma-norm-in-normal-case}
Soit  $\pi : X \to Y$  un morphisme fini et surjectif avec  $X$  et  $Y$ 
intégraux et  $Y$  normal. Alors il existe une norme de degré  $[R(X) : R(Y)]$ 
pour  $\pi$.
\end{lemma}

\begin{proof}
Soit  $\Spec(B) \subset Y$  un ouvert affine et soit  $\Spec(A) \subset X$  son image réciproque.
Alors  $A$  et  $B$  sont des domaines. Soit  $K$  le corps
des fractions de  $A$  et  $L$  le corps des fractions de
 $B$. Considérons le diagramme
$$
\xymatrix{
L \ar[r] & K \\
B \ar[u] \ar[r] & A \ar[u]
}
$$
Étant donné que  $K/L$  est une extension finie, il existe une application
norme  $\text{Norm}_{K/L} : K^* \to L^*$  de degré  $d = [K : L]$; elle envoie  $f \in K$ 
sur  $\det_L(f : K \to K)$  comme dans la démonstration du Lemme  \ref{divisors-lemma-finite-locally-free-has-norm}. Observons que
le polynôme caractéristique de  $f : K \to K$  est une puissance du polynôme minimal
de  $f$  sur  $L$; en particulier  $\text{Norm}_{K/L}(f)$  est une puissance du
coefficient constant du polynôme minimal de  $f$  sur  $L$. Par
conséquent, d'après l'Algèbre, Lemme  \ref{algebra-lemma-minimal-polynomial-normal-domain} ,  $\text{Norm}_{K/L}$  envoie  $A$ 
dans  $B$. Cela détermine un système compatible de fonctions sur les
sections sur des affines et donc une application norme globale  $\text{Norm}_\pi$  de
degré  $d$.

\medskip\noindent
La propriété (1) découle immédiatement de la construction. Pour voir la propriété
(2), soit  $f \in A$  contenu dans l'idéal premier  $\mathfrak p \subset A$. Soit
 $f^m + b_1 f^{m - 1} + \ldots + b_m$  le polynôme minimal de  $f$  sur  $L$. Par Algèbre,
Lemme  \ref{algebra-lemma-minimal-polynomial-normal-domain} , nous avons  $b_i \in B$. Donc  $b_0 \in B \cap \mathfrak p$. Puisque
 $\text{Norm}_{K/L}(f) = b_0^{d/m}$  (voir ci-dessus), nous concluons que la norme s'annule au point image
de  $\mathfrak p$.
\end{proof}

\begin{lemma}
\label{divisors-lemma-Frobenius-gives-norm-for-reduction}
Soit  $X$  un schéma noethérien. Soit  $p$  un nombre premier tel
que  $p\mathcal{O}_X = 0$. Alors pour un certain  $e > 0$ , il existe une norme de degré
 $p^e$  pour  $X_{red} \to X$  où  $X_{red}$  est la réduction de  $X$.
\end{lemma}

\begin{proof}
Soit  $A$  un anneau noethérien avec  $pA = 0$. Soit  $I \subset A$ 
l'idéal des éléments nilpotents. Alors  $I^n = 0$  pour un certain  $n$ 
(Algèbre, Lemme  \ref{algebra-lemma-Noetherian-power}). Choisissons  $e$  tel que  $p^e \geq n$. Alors
$$
A/I \longrightarrow A,\quad
f \bmod I \longmapsto f^{p^e}
$$
est bien défini. Cela produit une norme de degré  $p^e$  pour  $\Spec(A/I) \to \Spec(A)$.
Maintenant, si  $X$  est obtenu en recollant certains schémas affines
 $\Spec(A_i)$  alors pour un certain  $e \gg 0$  ces applications se recollent en
une application de norme pour  $X_{red} \to X$. Détails omis.
\end{proof}

\begin{proposition}
\label{divisors-proposition-push-down-ample}
Soit  $\pi : X \to Y$  un morphisme fini surjectif de schémas. Supposons que
 $X$  ait un  $\mathcal{O}_X$-module inversible ample. Si
\begin{enumerate}
\item $\pi$  est fini et localement libre, soit
\item $Y$  est un schéma normal intègre, ou
\item $Y$  est noethérien,  $p\mathcal{O}_Y = 0$, et  $X = Y_{red}$,
\end{enumerate}
alors  $Y$  a un  $\mathcal{O}_Y$-module inversible ample.
\end{proposition}

\begin{proof}
Le cas (1) découle d'une combinaison des lemmes  \ref{divisors-lemma-finite-locally-free-has-norm}  et  \ref{divisors-lemma-norm-ample}. Le
cas (3) découle d'une combinaison des lemmes  \ref{divisors-lemma-Frobenius-gives-norm-for-reduction}  et  \ref{divisors-lemma-norm-ample}. Dans le
cas (2), nous remplaçons d'abord  $X$  par une composante irréductible de
 $X$  qui domine  $Y$  (considéré comme un sous-schéma fermé réduit
de  $X$). Ensuite, nous pouvons appliquer le lemme  \ref{divisors-lemma-norm-in-normal-case}.
\end{proof}

\begin{lemma}
\label{divisors-lemma-push-down-quasi-affine}
Soit  $\pi : X \to Y$  un morphisme fini surjectif de schémas. Supposons que
 $X$  soit quasi-affine. Si l'une des conditions suivantes est satisfaite :
\begin{enumerate}
\item $\pi$  est fini et localement libre, soit
\item $Y$  est un schéma intègre normal
\end{enumerate}
alors  $Y$  est quasi-affine.
\end{lemma}

\begin{proof}
Le cas (1) découle d'une combinaison des lemmes  \ref{divisors-lemma-finite-locally-free-has-norm}  et  \ref{divisors-lemma-norm-quasi-affine}. Dans
le cas (2), nous remplaçons d'abord  $X$  par une composante irréductible
de  $X$  qui domine  $Y$  (considéré comme un sous-schéma fermé
réduit de  $X$). Ensuite, nous pouvons appliquer le lemme  \ref{divisors-lemma-norm-in-normal-case}.
\end{proof}





\section{Diviseurs de Cartier effectifs relatifs}
\label{divisors-section-effective-Cartier-morphisms}

\noindent
Le lemme suivant montre qu'un diviseur de Cartier effectif plat sur la
base est bien une « famille de diviseurs de Cartier effectifs » sur cette base.
Par exemple, la restriction à toute fibre est un diviseur de Cartier effectif.

\begin{lemma}
\label{divisors-lemma-relative-Cartier}
Soit  $f : X \to S$  un morphisme de schémas. Soit  $D \subset X$  un sous-schéma fermé.
Supposons
\begin{enumerate}
\item $D$  est un diviseur de Cartier effectif, et
\item $D \to S$  est un morphisme plat.
\end{enumerate}
Puis, pour tout morphisme de schémas  $g : S' \to S$ , l'image réciproque  $(g')^{-1}D$ 
est un diviseur de Cartier effectif sur  $X' = S' \times_S X$  où  $g' : X' \to X$  est la
projection.
\end{lemma}

\begin{proof}
En utilisant le Lemme  \ref{divisors-lemma-characterize-effective-Cartier-divisor} , nous traduisons ceci comme suit en algèbre.
Soit  $A \to B$  une application d'anneaux et  $h \in B$. Supposons que
 $h$  soit un non-diviseur de zéro et que  $B/hB$  soit plat sur
 $A$. Alors
$$
0 \to B \xrightarrow{h} B \to B/hB \to 0
$$
est une suite exacte de  $A$-modules avec  $B/hB$  plat sur
 $A$. D'après le Lemme d'Algèbre  \ref{algebra-lemma-flat-tor-zero} , cette suite reste exacte
lors du produit tensoriel sur  $A$  avec n'importe quel module, en
particulier avec n'importe quelle  $A$-algèbre  $A'$.
\end{proof}

\noindent
Ce lemme est la motivation pour la définition suivante.

\begin{definition}
\label{divisors-definition-relative-effective-Cartier-divisor}
Soit  $f : X \to S$  un morphisme de schémas. Un  {\it  diviseur de
Cartier effectif relatif}  sur  $X/S$  est un diviseur de Cartier
effectif  $D \subset X$  tel que  $D \to S$  soit un morphisme plat de schémas.
\end{definition}

\noindent
Nous avertissons le lecteur que cela peut être une notation non standard. En
particulier, dans  \cite[IV, section 21.15]{EGA}  la notion de diviseur relatif est discutée
uniquement lorsque  $X \to S$  est plat et localement de présentation finie.
Notre définition est un peu plus générale. Cependant, il s'avère que si
 $x \in D$  alors  $X \to S$  est plat à  $x$  dans de nombreux cas (mais
pas toujours).

\begin{lemma}
\label{divisors-lemma-sum-relative-effective-Cartier-divisor}
Soit  $f : X \to S$  un morphisme de schémas. Si  $D_1, D_2 \subset X$  sont des diviseurs de
Cartier effectifs relatifs sur  $X/S$  alors il en est de même de
 $D_1 + D_2$  (Définition  \ref{divisors-definition-sum-effective-Cartier-divisors}).
\end{lemma}

\begin{proof}
Cela se traduit par le fait algébrique suivant: Soit  $A \to B$  une application
d'anneaux et  $h_1, h_2 \in B$. Supposons que les  $h_i$  sont des diviseurs de
zéro et que  $B/h_iB$  est plat sur  $A$. Alors  $h_1h_2$  est un
non-diviseur de zéro et  $B/h_1h_2B$  est plat sur  $A$. La raison est que
nous avons une suite exacte
$$
0 \to B/h_1B \to B/h_1h_2B \to B/h_2B \to 0
$$
où la première flèche est donnée par la multiplication par  $h_2$. Comme les
deux modules extérieurs sont plats sur  $A$, il en est de même pour le
module du milieu, voir Algèbre, Lemme  \ref{algebra-lemma-flat-ses}.
\end{proof}

\begin{lemma}
\label{divisors-lemma-difference-relative-effective-Cartier-divisor}
Soit  $f : X \to S$  un morphisme de schémas. Si  $D_1, D_2 \subset X$  sont des diviseurs de
Cartier effectifs relatifs sur  $X/S$  et  $D_1 \subset D_2$  comme sous-schémas
fermés, alors le diviseur de Cartier effectif  $D$  tel que  $D_2 = D_1 + D$ 
(Lemme  \ref{divisors-lemma-difference-effective-Cartier-divisors}) est un diviseur de Cartier effectif relatif sur  $X/S$.
\end{lemma}

\begin{proof}
Cela se traduit par le fait algébrique suivant: Soit  $A \to B$  un homomorphisme
d'anneaux et  $h_1, h_2 \in B$. Supposons que les  $h_i$  soient des diviseurs de
zéro non nuls, que  $B/h_iB$  soit plat sur  $A$, et que  $(h_2) \subset (h_1)$.
Alors nous pouvons écrire  $h_2 = h h_1$  où  $h \in B$  est un diviseur de zéro non
nul. Nous obtenons une suite exacte
$$
0 \to B/hB \to B/h_2B \to B/h_1B \to 0
$$
où la première flèche est donnée par la multiplication par  $h_1$. Puisque
les deux modules de droite sont plats sur  $A$, il en est de même pour
celui du milieu, voir Algèbre, Lemme  \ref{algebra-lemma-flat-ses}.
\end{proof}

\begin{lemma}
\label{divisors-lemma-flat-at-x}
Soit  $f : X \to S$  un morphisme de schémas. Soit  $D \subset X$  un diviseur de
Cartier effectif relatif sur  $X/S$. Si  $x \in D$  et si  $\mathcal{O}_{X, x}$  est
noethérien, alors  $f$  est plat en  $x$.
\end{lemma}

\begin{proof}
Posons  $A = \mathcal{O}_{S, f(x)}$  et  $B = \mathcal{O}_{X, x}$. Soit  $h \in B$  un élément qui génère
l'idéal de  $D$. Alors  $h$  est un non-diviseur de zéro dans
 $B$  tel que  $B/hB$  soit une  $A$-algèbre locale plate. Soit
 $I \subset A$  un idéal de type fini. Considérons le diagramme commutatif
$$
\xymatrix{
0 \ar[r] &
B \ar[r]_h &
B \ar[r] &
B/hB \ar[r] & 0 \\
0 \ar[r] &
B \otimes_A I \ar[r]^h \ar[u] &
B \otimes_A I \ar[r] \ar[u] &
B/hB \otimes_A I \ar[r] \ar[u] & 0
}
$$
La suite inférieure est exacte car  $B/hB$  est plat sur
 $A$, voir Algèbre, Lemme  \ref{algebra-lemma-flat-tor-zero}. La flèche verticale de droite est
injective car  $B/hB$  est plat sur  $A$, voir Algèbre, Lemme
 \ref{algebra-lemma-flat}. Par conséquent, la multiplication par  $h$  est surjective
sur le noyau  $K$  de la flèche verticale du milieu. Par le lemme de
Nakayama, voir Algèbre, Lemme  \ref{algebra-lemma-NAK} , nous concluons que  $K= 0$. Par
conséquent,  $B$  est plat sur  $A$, voir Algèbre, Lemme
 \ref{algebra-lemma-flat}.
\end{proof}

\noindent
Le lemme suivant repose sur la version algébrique de l'ouverture du lieu plat. La
version schématique peut être trouvée dans Compléments sur les morphismes, section
 \ref{more-morphisms-section-open-flat}.

\begin{lemma}
\label{divisors-lemma-flat-relative-Cartier-divisor}
Soit  $f : X \to S$  un morphisme de schémas. Soit  $D \subset X$  un diviseur de
Cartier effectif relatif. Si  $f$  est localement de présentation finie,
alors il existe un sous-schéma ouvert  $U \subset X$  tel que  $D \subset U$  et tel que
 $f|_U : U \to S$  soit plat.
\end{lemma}

\begin{proof}
Choisissons  $x \in D$. Il suffit de trouver un voisinage ouvert  $U \subset X$  de
 $x$  tel que  $f|_U$  soit plat. Ainsi, le lemme se réduit au cas où
 $X = \Spec(B)$  et  $S = \Spec(A)$  sont affines et où  $D$  est donné par un
élément non diviseur de zéro  $h \in B$. Par hypothèse,  $B$  est une
 $A$-algèbre de présentation finie et  $B/hB$  est une
 $A$-algèbre plate. Nous allons utiliser l'approximation absolue de
Noether.

\medskip\noindent
Écrivons  $B = A[x_1, \ldots, x_n]/(g_1, \ldots, g_m)$. Supposons que  $h$  est l'image de  $h' \in A[x_1, \ldots, x_n]$.
Choisissons une  $\mathbf{Z}$-sous-algèbre de type fini  $A_0 \subset A$  telle que tous les
coefficients des polynômes  $h', g_1, \ldots, g_m$  soient dans  $A_0$. Alors nous
pouvons définir  $B_0 = A_0[x_1, \ldots, x_n]/(g_1, \ldots, g_m)$  et  $h_0$  l'image de  $h'$  dans
 $B_0$. Alors  $B = B_0 \otimes_{A_0} A$  et  $B/hB = B_0/h_0B_0 \otimes_{A_0} A$. D'après le Lemme  \ref{algebra-lemma-flat-finite-presentation-limit-flat} 
d'Algèbre, nous pouvons, après avoir élargi  $A_0$, supposer que
 $B_0/h_0B_0$  est plat sur  $A_0$. Soit  $K_0 = \Ker(h_0 : B_0 \to B_0)$. Comme  $B_0$  est
de type fini sur  $\mathbf{Z}$ , nous voyons que  $K_0$  est un idéal de type
fini. Soit  $A_1 \subset A$  une  $\mathbf{Z}$-sous-algèbre de type fini contenant
 $A_0$  et notons  $B_1$,  $h_1$,  $K_1$  les objets
correspondants sur  $A_1$. D'après Compléments d'algèbre, Lemme  \ref{more-algebra-lemma-base-change-H1-regular} ,
l'application  $K_0 \otimes_{A_0} A_1 \to K_1$  est surjective. D'autre part, le noyau de  $h : B \to B$ 
est nul par hypothèse. Ainsi, chaque élément de  $K_0$  s'envoie sur zéro
dans  $K_1$  pour des sous-anneaux  $A_1 \subset A$ suffisamment grands. Comme
 $K_0$  est de type fini, nous en concluons que  $K_1 = 0$  pour un choix
approprié de  $A_1$.

\medskip\noindent
Posons  $f_1 : X_1 \to S_1$  égal à  $\Spec$  de l'application d'anneaux  $A_1 \to B_1$. Posons
 $D_1 = \Spec(B_1/h_1B_1)$. Puisque  $B = B_1 \otimes_{A_1} A$, c'est-à-dire  $X = X_1 \times_{S_1} S$, il suffit
maintenant de prouver le lemme pour  $X_1 \to S_1$  et le diviseur de Cartier
effectif relatif  $D_1$, voir Morphismes, Lemme  \ref{morphisms-lemma-base-change-module-flat}. Ainsi, nous
avons réduit au cas où  $A$  est un anneau noethérien. Dans ce cas, nous
savons que le morphisme d'anneaux  $A \to B$  est plat en chaque idéal premier
 $\mathfrak q$  de  $V(h)$  d'après le Lemme  \ref{divisors-lemma-flat-at-x}. Combiné avec le fait
que le lieu de platitude est ouvert dans ce cas, voir Algèbre, Théorème
 \ref{algebra-theorem-openness-flatness} , nous concluons.
\end{proof}

\noindent
Il y a aussi le lemme suivant (dont l'idée est apparemment due à Michael Artin,
voir  \cite{Nobile}) qui n'a besoin d'aucune hypothèse de finitude.

\begin{lemma}
\label{divisors-lemma-michael-artin}
Soit  $f : X \to S$  un morphisme de schémas. Soit  $D \subset X$  un diviseur de
Cartier effectif relatif sur  $X/S$. Si  $f$  est plat en tous les
points de  $X \setminus D$, alors  $f$  est plat.
\end{lemma}

\begin{proof}
Cela se traduit par le fait algébrique suivant: Soit  $A \to B$  un morphisme
d'anneaux et  $h \in B$. Supposons que  $h$  soit un non-diviseur de
zéro, que  $B/hB$  soit plat sur  $A$, et que la localisation
 $B_h$  soit plate sur  $A$. Alors  $B$  est plat sur
 $A$. La raison est que nous avons une suite exacte
$$
0 \to B \to B_h \to \colim_n (1/h^n)B/B \to 0
$$
et que le deuxième et le troisième termes sont plats sur  $A$, ce qui
implique que  $B$  est plat sur  $A$  (voir Algèbre, Lemme
 \ref{algebra-lemma-flat-ses}). Notez qu'une colimite filtrée de modules plats est plate (voir
Algèbre, Lemme  \ref{algebra-lemma-colimit-flat}) et que par induction sur  $n$  chaque
 $(1/h^n)B/B \cong B/h^nB$  est plat sur  $A$  puisqu'il s'insère dans la suite exacte
courte
$$
0 \to B/h^{n - 1}B \xrightarrow{h} B/h^nB \to B/hB \to 0
$$
Quelques détails omis.
\end{proof}

\begin{example}
\label{divisors-example-relative-cartier-ambient-space-not-flat}
Voici un exemple de diviseur de Cartier effectif relatif  $D$  où le schéma
ambiant n'est pas plat dans un voisinage de  $D$. À savoir, soit
 $A = k[t]$  et
$$
B = k[t, x, y, x^{-1}y, x^{-2}y, \ldots]/(ty, tx^{-1}y, tx^{-2}y, \ldots)
$$
Alors  $B$  n'est pas plat sur  $A$  mais  $B/xB \cong A$  est plat sur
 $A$. De plus,  $x$  est un non-diviseur de zéro et définit donc un
diviseur de Cartier effectif relatif dans  $\Spec(B)$  sur  $\Spec(A)$.
\end{example}

\noindent
Si le schéma ambiant est plat et localement de présentation finie sur la base,
alors nous pouvons caractériser un diviseur de Cartier effectif relatif en termes
de ses fibres. Voir aussi Compléments sur les morphismes, Lemme  \ref{more-morphisms-lemma-slice-given-element}  pour une approche
légèrement différente de ce lemme.

\begin{lemma}
\label{divisors-lemma-fibre-Cartier}
Soit  $\varphi : X \to S$  un morphisme plat qui est localement de présentation finie. Soit
 $Z \subset X$  un sous-schéma fermé. Soit  $x \in Z$  avec une image  $s \in S$.
\begin{enumerate}
\item Si  $Z_s \subset X_s$  est un diviseur de Cartier dans un voisinage de  $x$,
alors il existe un ouvert  $U \subset X$  et un diviseur de Cartier effectif relatif
 $D \subset U$  tels que  $Z \cap U \subset D$  et  $Z_s \cap U = D_s$.
\item Si  $Z_s \subset X_s$  est un diviseur de Cartier dans un voisinage de  $x$, que
le morphisme  $Z \to X$  est de présentation finie, et que  $Z \to S$  est plat
en  $x$, alors nous pouvons choisir  $U$  et  $D$  tels que
 $Z \cap U = D$.
\item Si  $Z_s \subset X_s$  est un diviseur de Cartier dans un voisinage de  $x$  et
 $Z$  est un sous-schéma fermé localement principal de  $X$  dans un
voisinage de  $x$, alors nous pouvons choisir  $U$  et  $D$ 
tels que  $Z \cap U = D$.
\end{enumerate}
En particulier, si  $Z \to S$  est localement de présentation finie et plat et
que toutes les fibres  $Z_s \subset X_s$  sont des diviseurs de Cartier effectifs, alors
 $Z$  est un diviseur de Cartier effectif relatif. De même, si  $Z$ 
est un sous-schéma fermé localement principal de  $X$  tel que toutes les
fibres  $Z_s \subset X_s$  sont des diviseurs de Cartier effectifs, alors  $Z$ 
est un diviseur de Cartier effectif relatif.
\end{lemma}

\begin{proof}
Choisissons des voisinages ouverts affines  $\Spec(A)$  de  $s$  et
 $\Spec(B)$  de  $x$  tels que  $\varphi(\Spec(B)) \subset \Spec(A)$. Soit  $\mathfrak p \subset A$  l'idéal
premier correspondant à  $s$. Soit  $\mathfrak q \subset B$  l'idéal premier
correspondant à  $x$. Soit  $I \subset B$  l'idéal correspondant à
 $Z$. D'après l'hypothèse initiale du lemme, nous savons que  $A \to B$ 
est plat et de présentation finie. L'hypothèse en (1) signifie que, après avoir
réduit  $\Spec(B)$, nous pouvons supposer que  $I(B \otimes_A \kappa(\mathfrak p))$  est engendré par un
seul élément qui est un diviseur de zéro dans  $B \otimes_A \kappa(\mathfrak p)$. Disons que
 $f \in I$  s'envoie sur ce générateur. Nous affirmons qu'après avoir inversé un
élément  $g \in B$,  $g \not \in \mathfrak q$  le sous-schéma fermé  $D = V(f) \subset \Spec(B_g)$  est un
diviseur de Cartier effectif relatif.

\medskip\noindent
D'après Algèbre, Lemme  \ref{algebra-lemma-flat-finite-presentation-limit-flat} , nous pouvons trouver un morphisme d'anneaux
de type fini plat  $A_0 \to B_0$  d'anneaux noethériens, un élément  $f_0 \in B_0$, un
morphisme d'anneaux  $A_0 \to A$  et un isomorphisme  $A \otimes_{A_0} B_0 \cong B$. Si  $\mathfrak p_0 = A_0 \cap \mathfrak p$ 
alors nous voyons que
$$
B \otimes_A \kappa(\mathfrak p) =
\left(B_0 \otimes_{A_0} \kappa(\mathfrak p_0)\right)
\otimes_{\kappa(\mathfrak p_0))} \kappa(\mathfrak p)
$$
, d'où  $f_0$  est un non-diviseur de zéro dans  $B_0 \otimes_{A_0} \kappa(\mathfrak p_0)$. D'après
Algèbre, Lemme  \ref{algebra-lemma-grothendieck} , nous voyons que  $f_0$  est un non-diviseur de
zéro dans  $(B_0)_{\mathfrak q_0}$  où  $\mathfrak q_0 = B_0 \cap \mathfrak q$  et que  $(B_0/f_0B_0)_{\mathfrak q_0}$  est plat sur
 $A_0$. Donc d'après Algèbre, Lemme  \ref{algebra-lemma-regular-sequence-in-neighbourhood}  et Théorème d'Algèbre
 \ref{algebra-theorem-openness-flatness} , il existe un  $g_0 \in B_0$,  $g_0 \not \in \mathfrak q_0$  tel que  $f_0$  est un
non-diviseur de zéro dans  $(B_0)_{g_0}$  et tel que  $(B_0/f_0B_0)_{g_0}$  est plat sur
 $A_0$. Par conséquent, nous voyons que  $D_0 = V(f_0) \subset \Spec((B_0)_{g_0})$  est un diviseur de
Cartier effectif relatif. Puisque nous savons que cette propriété est préservée
par changement de base, voir le Lemme  \ref{divisors-lemma-relative-Cartier}, nous obtenons l'affirmation
mentionnée ci-dessus avec  $g$  égal à l'image de  $g_0$  dans
 $B$.

\medskip\noindent
À ce stade, nous avons démontré (1). Pour voir (2), considérons l'immersion fermée
 $Z \to D$. L'application d'anneaux surjective  $u : \mathcal{O}_{D, x} \to \mathcal{O}_{Z, x}$  est une application
d'algèbres locales plates sur  $\mathcal{O}_{S, s}$, essentiellement de présentation finie,
et qui devient un isomorphisme après division par  $\mathfrak m_s$. Par conséquent,
c'est un isomorphisme, voir Algèbre, Lemme  \ref{algebra-lemma-mod-injective-general}. Il s'ensuit que
 $Z \to D$  est un isomorphisme dans un voisinage de  $x$, voir Algèbre,
Lemme  \ref{algebra-lemma-local-isomorphism}. Pour voir (3), après avoir éventuellement rétréci
 $U$ , nous pouvons supposer que l'idéal de  $D$  est engendré par
un seul non-diviseur de zéro  $f$  et que l'idéal de  $Z$  est
engendré par un élément  $g$. Alors  $f = gh$. Mais  $g|_{U_s}$  et
 $f|_{U_s}$  définissent le même diviseur de Cartier effectif dans un voisinage de
 $x$. Par conséquent,  $h|_{X_s}$  est une unité dans  $\mathcal{O}_{X_s, x}$, donc
 $h$  est une unité dans  $\mathcal{O}_{X, x}$  et donc  $h$  est une unité
dans un voisinage ouvert de  $x$. C.-à-d.,  $Z \cap U = D$  après avoir
rétréci  $U$.

\medskip\noindent
Les affirmations finales du lemme suivent immédiatement des parties (2) et (3),
combinées avec le fait que  $Z \to S$  est localement de présentation finie si et
seulement si  $Z \to X$  est de présentation finie, voir Morphismes, Lemme
 \ref{morphisms-lemma-composition-finite-presentation}  et  \ref{morphisms-lemma-finite-presentation-permanence}.
\end{proof}



\section{Le cône normal d'une immersion}
\label{divisors-section-normal-cone}

\noindent
Soit $i : Z \to X$ une immersion fermée. Soit $\mathcal{I} \subset \mathcal{O}_X$ le faisceau
quasi-cohérent d'idéaux correspondant. Considérons le faisceau quasi-cohérent
d'algèbres graduées sur $\mathcal{O}_X$ donné par $\bigoplus_{n \geq 0} \mathcal{I}^n/\mathcal{I}^{n + 1}$. Comme les faisceaux $\mathcal{I}^n/\mathcal{I}^{n + 1}$ sont tous
annulés par $\mathcal{I}$, cette algèbre graduée correspond à un faisceau quasi-cohérent
d'algèbres graduées sur $\mathcal{O}_Z$ par Morphismes, Lemme \ref{morphisms-lemma-i-star-equivalence}. Ce
faisceau d'algèbres graduées sur $\mathcal{O}_Z$ est appelé l'{\it algèbre conormale
de $Z$ dans $X$}  et est souvent simplement noté
 $\bigoplus_{n \geq 0} \mathcal{I}^n/\mathcal{I}^{n + 1}$ par l'abus de notation mentionné dans Morphismes, section \ref{morphisms-section-closed-immersions-quasi-coherent}.

\medskip\noindent
Soit  $f : Z \to X$  une immersion. Nous définissons l'algèbre conormale de
 $f$  comme l'algèbre conormale de l'immersion fermée  $i : Z \to X \setminus \partial Z$, où
 $\partial Z = \overline{Z} \setminus Z$. Elle est souvent notée  $\bigoplus_{n \geq 0} \mathcal{I}^n/\mathcal{I}^{n + 1}$  où  $\mathcal{I}$  est le faisceau
d'idéaux de l'immersion fermée  $i : Z \to X \setminus \partial Z$.

\begin{definition}
\label{divisors-definition-conormal-sheaf}
Soit $f : Z \to X$ une immersion. L'{\it algèbre
conormale $\mathcal{C}_{Z/X, *}$ de $Z$ dans $X$}, ou
l'{\it algèbre conormale de $f$}, est le faisceau quasi-cohérent d'algèbres
graduées sur $\mathcal{O}_Z$ donné par $\bigoplus_{n \geq 0} \mathcal{I}^n/\mathcal{I}^{n + 1}$ et décrit ci-dessus.
\end{definition}

\noindent
Ainsi,  $\mathcal{C}_{Z/X, 1} = \mathcal{C}_{Z/X}$  est le faisceau conormal de l'immersion. De plus,  $\mathcal{C}_{Z/X, 0} = \mathcal{O}_Z$ 
et  $\mathcal{C}_{Z/X, n}$  est un  $\mathcal{O}_Z$-module quasi-cohérent caractérisé par la
propriété
\begin{equation}
\label{divisors-equation-conormal-in-degree-n}
i_*\mathcal{C}_{Z/X, n} = \mathcal{I}^n/\mathcal{I}^{n + 1}
\end{equation}
où  $i : Z \to X \setminus \partial Z$  et  $\mathcal{I}$  est le faisceau d'idéaux de  $i$  comme
ci-dessus. Enfin, noter qu'il existe un morphisme surjectif canonique
\begin{equation}
\label{divisors-equation-conormal-algebra-quotient}
\text{Sym}^*(\mathcal{C}_{Z/X}) \longrightarrow \mathcal{C}_{Z/X, *}
\end{equation}
d'algèbres  $\mathcal{O}_Z$ graduées quasi-cohérentes qui est un isomorphisme dans les
degrés  $0$  et  $1$.

\begin{lemma}
\label{divisors-lemma-affine-conormal-sheaf}
Soit  $i : Z \to X$  une immersion. L'algèbre conormale de  $i$  possède les
propriétés suivantes:
\begin{enumerate}
\item Soit  $U \subset X$  un ouvert quelconque tel que  $i(Z)$  est un sous-ensemble
fermé de  $U$. Soit  $\mathcal{I} \subset \mathcal{O}_U$  le faisceau d'idéaux correspondant au
sous-schéma fermé  $i(Z) \subset U$. Alors
$$
\mathcal{C}_{Z/X, *} =
i^*\left(\bigoplus\nolimits_{n \geq 0} \mathcal{I}^n\right) =
i^{-1}\left(
\bigoplus\nolimits_{n \geq 0} \mathcal{I}^n/\mathcal{I}^{n + 1}
\right)
$$
\item Pour tout ouvert affine  $\Spec(R) = U \subset X$  tel que  $Z \cap U = \Spec(R/I)$  il existe un
isomorphisme canonique  $\Gamma(Z \cap U, \mathcal{C}_{Z/X, *}) = \bigoplus_{n \geq 0} I^n/I^{n + 1}$.
\end{enumerate}
\end{lemma}

\begin{proof}
 Cela découle essentiellement des définitions. Notons que, pour un anneau
 $R$  et un idéal  $I$  de  $R$, nous avons  $I^n/I^{n + 1} = I^n \otimes_R R/I$.
Détails omis.
\end{proof}

\begin{lemma}
\label{divisors-lemma-conormal-algebra-functorial}
Soit
$$
\xymatrix{
Z \ar[r]_i \ar[d]_f & X \ar[d]^g \\
Z' \ar[r]^{i'} & X'
}
$$
soit un diagramme commutatif dans la catégorie des schémas. Supposons
 $i$,  $i'$  immersions. Il existe un morphisme canonique d'algèbres
 $\mathcal{O}_Z$ graduées
$$
f^*\mathcal{C}_{Z'/X', *}
\longrightarrow
\mathcal{C}_{Z/X, *}
$$
caractérisé par la propriété suivante: Pour chaque paire d'ouverts affines
 $(\Spec(R) = U \subset X, \Spec(R') = U' \subset X')$  avec  $g(U) \subset U'$  tels que  $Z \cap U = \Spec(R/I)$  et  $Z' \cap U' = \Spec(R'/I')$  la flèche
induite
$$
\Gamma(Z' \cap U', \mathcal{C}_{Z'/X', *}) =
\bigoplus\nolimits (I')^n/(I')^{n + 1}
\longrightarrow
\bigoplus\nolimits_{n \geq 0} I^n/I^{n + 1} =
\Gamma(Z \cap U, \mathcal{C}_{Z/X, *})
$$
est celle induite par le morphisme d'anneaux  $f^\sharp : R' \to R$  qui a la propriété
 $f^\sharp(I') \subset I$.
\end{lemma}

\begin{proof}
Soient  $\partial Z' = \overline{Z'} \setminus Z'$  et  $\partial Z = \overline{Z} \setminus Z$. Ce sont des sous-ensembles fermés de
 $X'$  et de  $X$. En remplaçant  $X'$  par  $X' \setminus \partial Z'$  et
 $X$  par  $X \setminus \big(g^{-1}(\partial Z') \cup \partial Z\big)$ , nous voyons que nous pouvons supposer que
 $i$  et  $i'$  sont des immersions fermées.

\medskip\noindent
Le fait que  $g \circ i$  se factorise par  $i'$  implique que
 $g^*\mathcal{I}'$  s'envoie dans  $\mathcal{I}$  par l'application canonique
 $g^*\mathcal{I}' \to \mathcal{O}_X$, voir Schémas, Lemmes  \ref{schemes-lemma-characterize-closed-subspace}  et  \ref{schemes-lemma-restrict-map-to-closed}. Par conséquent,
nous obtenons une application induite de faisceaux quasi-cohérents  $g^*((\mathcal{I}')^n/(\mathcal{I}')^{n + 1}) \to
\mathcal{I}^n/\mathcal{I}^{n + 1}$.
L'image réciproque par  $i$  donne  $i^*g^*((\mathcal{I}')^n/(\mathcal{I}')^{n + 1}) \to
i^*(\mathcal{I}^n/\mathcal{I}^{n + 1})$. Notons que  $i^*(\mathcal{I}^n/\mathcal{I}^{n + 1}) = \mathcal{C}_{Z/X, n}$.
D'autre part,  $i^*g^*((\mathcal{I}')^n/(\mathcal{I}')^{n + 1}) =
f^*(i')^*((\mathcal{I}')^n/(\mathcal{I}')^{n + 1}) =
f^*\mathcal{C}_{Z'/X', n}$. Cela donne l'application désirée.

\medskip\noindent
Vérifier que l'application est localement décrite comme l'application donnée
 $(I')^n/(I')^{n + 1} \to I^n/I^{n + 1}$  est une question de déroulement des définitions et est omis. Une
autre observation est que pour tout  $x \in i(Z)$  donné, il existe effectivement
des voisinages ouverts affines  $U$,  $U'$  avec  $f(U) \subset U'$  et
 $Z \cap U$  ainsi que  $U' \cap Z'$  fermés tels que  $x \in U$. Preuve omise.
Par conséquent, l'exigence du lemme caractérise effectivement l'application (et
aurait pu être utilisée pour la définir).
\end{proof}

\begin{lemma}
\label{divisors-lemma-conormal-algebra-functorial-flat}
Soit
$$
\xymatrix{
Z \ar[r]_i \ar[d]_f & X \ar[d]^g \\
Z' \ar[r]^{i'} & X'
}
$$
soit un diagramme de produits fibrés dans la catégorie des schémas avec
 $i$,  $i'$  immersions. Alors l'application canonique  $f^*\mathcal{C}_{Z'/X', *} \to \mathcal{C}_{Z/X, *}$ 
du Lemme  \ref{divisors-lemma-conormal-algebra-functorial}  est surjective. Si  $g$  est plat, alors c'est un
isomorphisme.
\end{lemma}

\begin{proof}
Soit  $R' \to R$  un morphisme d'anneaux, et  $I' \subset R'$  un idéal. Posons
 $I = I'R$. Alors  $(I')^n/(I')^{n + 1} \otimes_{R'} R \to I^n/I^{n + 1}$  est surjective. Si  $R' \to R$  est plat, alors
 $I^n = (I')^n \otimes_{R'} R$  et nous voyons que l'application est un isomorphisme.
\end{proof}

\begin{definition}
\label{divisors-definition-normal-cone}
Soit  $i : Z \to X$  une immersion de schémas. Le {\it cône normal
 $C_ZX$}  de  $Z$  dans
 $X$  est
$$
C_ZX = \underline{\Spec}_Z(\mathcal{C}_{Z/X, *})
$$
voir Constructions, Définitions  \ref{constructions-definition-cone}  et  \ref{constructions-definition-abstract-cone}. Le {\it fibré normal}
de  $Z$  dans  $X$  est le
fibré vectoriel
$$
N_ZX = \underline{\Spec}_Z(\text{Sym}(\mathcal{C}_{Z/X}))
$$
voir Constructions, Définitions  \ref{constructions-definition-vector-bundle}  et  \ref{constructions-definition-abstract-vector-bundle}.
\end{definition}

\noindent
Ainsi,  $C_ZX \to Z$  est un cône sur  $Z$  et  $N_ZX \to Z$  est un fibré
vectoriel sur  $Z$  (rappelons que dans notre terminologie, cela n'implique
pas que le faisceau conormal soit un faisceau localement libre de type fini). De
plus, la surjection canonique (\ref{divisors-equation-conormal-algebra-quotient}) des algèbres graduées définit une
immersion fermée canonique
\begin{equation}
\label{divisors-equation-normal-cone-in-normal-bundle}
C_ZX \longrightarrow N_ZX
\end{equation}
des cônes sur  $Z$.





\section{Faisceaux d'idéaux réguliers}
\label{divisors-section-regular-ideal-sheaves}

\noindent
Dans cette section, nous généralisons la notion de diviseur de Cartier effectif à
des codimensions supérieures. Rappelons qu'une suite d'éléments  $f_1, \ldots, f_r$  d'un
anneau  $R$  est une  {\it  suite régulière} 
si pour chaque  $i = 1, \ldots, r$  l'élément  $f_i$  est un non-diviseur de zéro dans
 $R/(f_1, \ldots, f_{i - 1})$  et  $R/(f_1, \ldots, f_r) \not = 0$, voir Algèbre, Définition  \ref{algebra-definition-regular-sequence}. Il existe
trois conditions plus faibles et étroitement liées que nous pouvons imposer. La
première consiste à supposer que  $f_1, \ldots, f_r$  est une  {\it 
suite Koszul-régulière}, c'est-à-dire que  $H_i(K_\bullet(f_1, \ldots, f_r)) = 0$  pour
 $i > 0$, voir Compléments d'algèbre, Définition  \ref{more-algebra-definition-koszul-regular-sequence}. La suite est
appelée une  {\it  suite  $H_1$-régulière}  si
 $H_1(K_\bullet(f_1, \ldots, f_r)) = 0$. Une autre condition que nous pouvons imposer est qu'avec
 $J = (f_1, \ldots, f_r)$, l'application
$$
R/J[T_1, \ldots, T_r]
\longrightarrow
\bigoplus\nolimits_{n \geq 0}
J^n/J^{n + 1}
$$
qui associe  $T_i$  à  $f_i \bmod J^2$  est un isomorphisme. Dans ce cas, nous
disons que  $f_1, \ldots, f_r$  est une  {\it  suite
quasi-régulière}, voir Algèbre, Définition  \ref{algebra-definition-quasi-regular-sequence}. Étant donné un
 $R$-module  $M$ , il existe également une notion de
 $M$-séquence régulière et de  $M$-séquence quasi-régulière.

\medskip\noindent
Nous pouvons généraliser cela au cas des espaces annelés comme suit. Soit
 $X$  un espace annelé et soit  $f_1, \ldots, f_r \in \Gamma(X, \mathcal{O}_X)$. Nous disons que  $f_1, \ldots, f_r$ 
est une  {\it  séquence régulière}  si pour chaque
 $i = 1, \ldots, r$ , l'application
\begin{equation}
\label{divisors-equation-map-regular}
f_i :
\mathcal{O}_X/(f_1, \ldots, f_{i - 1})
\longrightarrow
\mathcal{O}_X/(f_1, \ldots, f_{i - 1})
\end{equation}
est une application injective de faisceaux. Nous disons que  $f_1, \ldots, f_r$  est une
 {\it  séquence Koszul-régulière}  si le complexe de
Koszul
\begin{equation}
\label{divisors-equation-koszul}
K_\bullet(\mathcal{O}_X, f_\bullet),
\end{equation}
voir Modules, Définition  \ref{modules-definition-koszul-complex}, est acyclique dans les degrés  $> 0$.
Nous disons que  $f_1, \ldots, f_r$  est une séquence  {\it 
 $H_1$-régulière}  si le complexe de Koszul  $K_\bullet(\mathcal{O}_X, f_\bullet)$  est exact
dans le degré  $1$. Enfin, nous disons que  $f_1, \ldots, f_r$  est une séquence
 {\it  quasi-régulière}  si l'application
\begin{equation}
\label{divisors-equation-map-quasi-regular}
\mathcal{O}_X/\mathcal{J}[T_1, \ldots, T_r]
\longrightarrow
\bigoplus\nolimits_{d \geq 0} \mathcal{J}^d/\mathcal{J}^{d + 1}
\end{equation}
est un isomorphisme de faisceaux où  $\mathcal{J} \subset \mathcal{O}_X$  est le faisceau d'idéaux engendré
par  $f_1, \ldots, f_r$. (Il existe également une notion de séquence
 $\mathcal{F}$-régulière et  $\mathcal{F}$-quasi-régulière pour un  $\mathcal{O}_X$-module
 $\mathcal{F}$  donné que nous introduirons ici si nous en avons besoin.)

\begin{lemma}
\label{divisors-lemma-types-regular-sequences-implications}
Soit  $X$  un espace annelé. Soit  $f_1, \ldots, f_r \in \Gamma(X, \mathcal{O}_X)$. Nous avons les implications
suivantes  $f_1, \ldots, f_r$  est une suite régulière  $\Rightarrow$   $f_1, \ldots, f_r$  est une
suite Koszul-régulière  $\Rightarrow$   $f_1, \ldots, f_r$  est une suite
 $H_1$-régulière  $\Rightarrow$   $f_1, \ldots, f_r$  est une suite quasi-régulière.
\end{lemma}

\begin{proof}
Puisque nous pouvons vérifier l'exactitude sur les fibres, une suite  $f_1, \ldots, f_r$ 
est une suite régulière si et seulement si les applications
$$
f_i :
\mathcal{O}_{X, x}/(f_1, \ldots, f_{i - 1})
\longrightarrow
\mathcal{O}_{X, x}/(f_1, \ldots, f_{i - 1})
$$
sont injectives pour tout  $x \in X$. Les autres types de régularité peuvent
également être vérifiés fibre par fibre (détails omis). Par conséquent, les
implications découlent de Compléments d'algèbre, Lemme  \ref{more-algebra-lemma-regular-koszul-regular},  \ref{more-algebra-lemma-koszul-regular-H1-regular}, et
 \ref{more-algebra-lemma-H1-regular-quasi-regular}.
\end{proof}

\begin{definition}
\label{divisors-definition-regular-ideal-sheaf}
\begin{reference}
Le concept de faisceau d'idéaux Koszul-régulier a été introduit dans  \cite[Expose VII, Definition 1.4]{SGA6} 
où il était appelé un faisceau d'idéaux régulier.
\end{reference}
Soit  $X$  un espace annelé. Soit  $\mathcal{J} \subset \mathcal{O}_X$  un faisceau d'idéaux.
\begin{enumerate}
\item On dit que  $\mathcal{J}$  est  {\it  régulier}  si pour
chaque  $x \in \text{Supp}(\mathcal{O}_X/\mathcal{J})$  il existe un voisinage ouvert  $x \in U \subset X$  et une suite
régulière  $f_1, \ldots, f_r \in \mathcal{O}_X(U)$  telle que  $\mathcal{J}|_U$  est engendré par  $f_1, \ldots, f_r$.
\item On dit que  $\mathcal{J}$  est  {\it  Koszul-régulier} 
si pour chaque  $x \in \text{Supp}(\mathcal{O}_X/\mathcal{J})$  il existe un voisinage ouvert  $x \in U \subset X$  et une
suite Koszul-régulière  $f_1, \ldots, f_r \in \mathcal{O}_X(U)$  telle que  $\mathcal{J}|_U$  est engendré par
 $f_1, \ldots, f_r$.
\item On dit que  $\mathcal{J}$  est  {\it 
 $H_1$-régulier}  si pour chaque  $x \in \text{Supp}(\mathcal{O}_X/\mathcal{J})$  il existe un
voisinage ouvert  $x \in U \subset X$  et une suite  $H_1$-régulière  $f_1, \ldots, f_r \in \mathcal{O}_X(U)$ 
telle que  $\mathcal{J}|_U$  est engendré par  $f_1, \ldots, f_r$.
\item Nous disons que  $\mathcal{J}$  est  {\it 
quasi-régulier}  si pour chaque  $x \in \text{Supp}(\mathcal{O}_X/\mathcal{J})$  il existe un voisinage
ouvert  $x \in U \subset X$  et une suite quasi-régulière  $f_1, \ldots, f_r \in \mathcal{O}_X(U)$  telle que
 $\mathcal{J}|_U$  est engendré par  $f_1, \ldots, f_r$.
\end{enumerate}
\end{definition}

\noindent
De nombreuses propriétés de cette notion découlent immédiatement des notions
correspondantes pour les suites régulières et quasi-régulières dans les anneaux.

\begin{lemma}
\label{divisors-lemma-regular-quasi-regular-scheme}
Soit  $X$  un espace annelé. Soit  $\mathcal{J}$  un faisceau d'idéaux. Nous
avons les implications suivantes:  $\mathcal{J}$  est régulier  $\Rightarrow$ 
 $\mathcal{J}$  est Koszul-régulier  $\Rightarrow$   $\mathcal{J}$  est
 $H_1$-régulier  $\Rightarrow$   $\mathcal{J}$  est quasi-régulier.
\end{lemma}

\begin{proof}
Le lemme se réduit immédiatement au Lemme  \ref{divisors-lemma-types-regular-sequences-implications}.
\end{proof}

\begin{lemma}
\label{divisors-lemma-quasi-regular-ideal}
Soit  $X$  un espace localement annelé. Soit  $\mathcal{J} \subset \mathcal{O}_X$  un faisceau
d'idéaux. Alors  $\mathcal{J}$  est quasi-régulier si et seulement si les conditions
suivantes sont satisfaites:
\begin{enumerate}
\item $\mathcal{J}$  est un  $\mathcal{O}_X$-module de type fini,
\item $\mathcal{J}/\mathcal{J}^2$  est un  $\mathcal{O}_X/\mathcal{J}$-module localement libre de type fini, et
\item les applications canoniques
$$
\text{Sym}^n_{\mathcal{O}_X/\mathcal{J}}(\mathcal{J}/\mathcal{J}^2)
\longrightarrow
\mathcal{J}^n/\mathcal{J}^{n + 1}
$$
sont des isomorphismes pour tout  $n \geq 0$.
\end{enumerate}
\end{lemma}

\begin{proof}
Il est clair que si  $U \subset X$  est un ouvert tel que  $\mathcal{J}|_U$  soit engendré
par une suite quasi-régulière  $f_1, \ldots, f_r \in \mathcal{O}_X(U)$  alors  $\mathcal{J}|_U$  est de type fini,
 $\mathcal{J}|_U/\mathcal{J}^2|_U$  est libre avec pour base  $f_1, \ldots, f_r$, et les applications de (3)
sont des isomorphismes car elles sont une formulation indépendante des coordonnées
de la partie de degré  $n$  de (\ref{divisors-equation-map-quasi-regular}). Il en découle qu'il est
clair qu'être quasi-régulier implique les conditions (1), (2) et (3).

\medskip\noindent
Inversement, supposons que (1), (2) et (3) soient vérifiés. Choisissons un point
 $x \in \text{Supp}(\mathcal{O}_X/\mathcal{J})$. Il existe alors un voisinage  $U \subset X$  de  $x$  tel que
 $\mathcal{J}|_U/\mathcal{J}^2|_U$  soit libre de rang  $r$  sur  $\mathcal{O}_U/\mathcal{J}|_U$. Après avoir
éventuellement restreint  $U$ , nous pouvons supposer qu'il existe
 $f_1, \ldots, f_r \in \mathcal{J}(U)$  qui s'envoient sur une base de  $\mathcal{J}|_U/\mathcal{J}^2|_U$  comme
 $\mathcal{O}_U/\mathcal{J}|_U$-module. En particulier, nous voyons que les images de  $f_1, \ldots, f_r$ 
dans  $\mathcal{J}_x/\mathcal{J}^2_x$  engendrent. Par conséquent, d'après le lemme de Nakayama
(Algèbre, Lemme  \ref{algebra-lemma-NAK}), nous voyons que  $f_1, \ldots, f_r$  engendrent la fibre
 $\mathcal{J}_x$. Ainsi, puisque  $\mathcal{J}$  est de type fini, d'après Modules, Lemme
 \ref{modules-lemma-finite-type-surjective-on-stalk} , après avoir restreint  $U$ , nous pouvons supposer que
 $f_1, \ldots, f_r$  engendrent  $\mathcal{J}$. Enfin, à partir de (3) et de l'isomorphisme
 $\mathcal{J}|_U/\mathcal{J}^2|_U = \bigoplus \mathcal{O}_U/\mathcal{J}|_U f_i$ , il est clair que  $f_1, \ldots, f_r \in \mathcal{O}_X(U)$  est une suite quasi-régulière.
\end{proof}

\begin{lemma}
\label{divisors-lemma-generate-regular-ideal}
Soit  $(X, \mathcal{O}_X)$  un espace localement annelé. Soit  $\mathcal{J} \subset \mathcal{O}_X$  un faisceau
d'idéaux. Soient  $x \in X$  et  $f_1, \ldots, f_r \in \mathcal{J}_x$  dont les images fournissent une base
de l'espace vectoriel  $\kappa(x)$  $\mathcal{J}_x/\mathfrak m_x\mathcal{J}_x$.
\begin{enumerate}
\item Si  $\mathcal{J}$  est quasi-régulier, alors il existe un voisinage ouvert tel que
 $f_1, \ldots, f_r \in \mathcal{O}_X(U)$  forme une suite quasi-régulière générant  $\mathcal{J}|_U$.
\item Si  $\mathcal{J}$  est  $H_1$-régulier, alors il existe un voisinage ouvert tel
que  $f_1, \ldots, f_r \in \mathcal{O}_X(U)$  forme une suite  $H_1$-régulière générant  $\mathcal{J}|_U$.
\item Si  $\mathcal{J}$  est Koszul-régulier, alors il existe un voisinage ouvert tel que
 $f_1, \ldots, f_r \in \mathcal{O}_X(U)$  forme une suite Koszul-régulière générant  $\mathcal{J}|_U$.
\end{enumerate}
\end{lemma}

\begin{proof}
Tout d'abord, supposons que  $\mathcal{J}$  est quasi-régulier. Nous pouvons choisir
un voisinage ouvert  $U \subset X$  de  $x$  et une suite quasi-régulière
 $g_1, \ldots, g_s \in \mathcal{O}_X(U)$  qui engendre  $\mathcal{J}|_U$. Notons que cela implique que  $\mathcal{J}/\mathcal{J}^2$ 
est libre de rang  $s$  sur  $\mathcal{O}_U/\mathcal{J}|_U$  (voir le Lemme  \ref{divisors-lemma-quasi-regular-ideal}  et sa
démonstration) et donc  $r = s$. Nous pouvons réduire  $U$  et supposer
 $f_1, \ldots, f_r \in \mathcal{J}(U)$. Ainsi nous pouvons écrire
$$
f_i = \sum a_{ij} g_j
$$
pour certains  $a_{ij} \in \mathcal{O}_X(U)$. Par hypothèse, la matrice  $A = (a_{ij})$  s'envoie sur
une matrice inversible sur  $\kappa(x)$. Par conséquent, après avoir encore réduit
 $U$ , nous pouvons supposer que  $(a_{ij})$  est inversible. Ainsi nous
voyons que  $f_1, \ldots, f_r$  donnent une base pour  $(\mathcal{J}/\mathcal{J}^2)|_U$  ce qui prouve que
 $f_1, \ldots, f_r$  est une suite quasi-régulière sur  $U$.

\medskip\noindent
Notons que pour prouver (2) et (3), nous pouvons, puisque les hypothèses de (2) et
(3) sont plus fortes que l'hypothèse de (1), déjà supposer que  $f_1, \ldots, f_r \in \mathcal{J}(U)$  et
 $f_i = \sum a_{ij}g_j$  avec  $(a_{ij})$  inversible comme ci-dessus, où maintenant
 $g_1, \ldots, g_r$  est une suite  $H_1$-régulière ou Koszul-régulière. Puisque le
complexe de Koszul sur  $f_1, \ldots, f_r$  est isomorphe au complexe de Koszul sur
 $g_1, \ldots, g_r$  via la matrice  $(a_{ij})$  (voir Compléments d'algèbre, Lemme
 \ref{more-algebra-lemma-change-basis}), nous concluons que  $f_1, \ldots, f_r$  est  $H_1$-régulier ou
Koszul-régulier comme souhaité.
\end{proof}

\begin{lemma}
\label{divisors-lemma-regular-ideal-sheaf-quasi-coherent}
Tout faisceau d'idéaux régulier, Koszul-régulier,  $H_1$-régulier ou
quasi-régulier sur un schéma est un faisceau d'idéaux quasi-cohérent de type fini.
\end{lemma}

\begin{proof}
Cela découle du fait qu'un tel faisceau d'idéaux est localement engendré par un
nombre fini de sections. Et tout faisceau d'idéaux localement engendré par des
sections sur un schéma est quasi-cohérent, voir Schémas, Lemme  \ref{schemes-lemma-closed-subspace-scheme}.
\end{proof}

\begin{lemma}
\label{divisors-lemma-regular-ideal-sheaf-scheme}
Soit  $X$  un schéma. Soit  $\mathcal{J}$  un faisceau d'idéaux. Alors
 $\mathcal{J}$  est régulier (resp. \ Koszul-régulier,
 $H_1$-régulier, quasi-régulier) si et seulement si pour chaque  $x \in \text{Supp}(\mathcal{O}_X/\mathcal{J})$ 
il existe un voisinage affine ouvert  $x \in U \subset X$,  $U = \Spec(A)$  tel que
 $\mathcal{J}|_U = \widetilde{I}$  et tel que  $I$  soit engendré par une suite régulière (resp.
\ Koszul-régulière,  $H_1$-régulière, quasi-régulière)
 $f_1, \ldots, f_r \in A$.
\end{lemma}

\begin{proof}
Par hypothèse, nous pouvons trouver un voisinage ouvert  $U$  de
 $x$  au-dessus duquel  $\mathcal{J}$  est engendré par une suite régulière
(resp. \ Koszul-régulière,  $H_1$-régulière, quasi-régulière)
 $f_1, \ldots, f_r \in \mathcal{O}_X(U)$. Après avoir restreint  $U$ , nous pouvons supposer que
 $U$  est affine, disons  $U = \Spec(A)$. Puisque  $\mathcal{J}$  est
quasi-cohérent d'après le Lemme  \ref{divisors-lemma-regular-ideal-sheaf-quasi-coherent} , nous voyons que  $\mathcal{J}|_U = \widetilde{I}$  pour un
certain idéal  $I \subset A$. Nous pouvons maintenant utiliser le fait que
$$
\widetilde{\ } : \text{Mod}_A \longrightarrow \QCoh(\mathcal{O}_U)
$$
est une équivalence de catégories qui préserve l'exactitude. Par exemple, le fait
que les fonctions  $f_i$  génèrent  $\mathcal{J}$  signifie que les
 $f_i$, vus comme des éléments de  $A$ , génèrent  $I$. Le
fait que (\ref{divisors-equation-map-regular}) soit injectif (resp. \ (\ref{divisors-equation-koszul}) soit
exact, (\ref{divisors-equation-koszul}) soit exact en degré  $1$, (\ref{divisors-equation-map-quasi-regular}) soit un
isomorphisme) implique la propriété correspondante de l'application  $A/(f_1, \ldots, f_{i - 1}) \to A/(f_1, \ldots, f_{i - 1})$ 
(resp. \  le complexe  $K_\bullet(A, f_1, \ldots, f_r)$, l'application  $A/I[T_1, \ldots, T_r] \to \bigoplus I^n/I^{n + 1}$). Ainsi,
 $f_1, \ldots, f_r \in A$  est une suite régulière (resp. \ Koszul-régulière,
 $H_1$-régulière, quasi-régulière) de l'anneau  $A$.
\end{proof}

\begin{lemma}
\label{divisors-lemma-Noetherian-scheme-regular-ideal}
Soit  $X$  un schéma localement noethérien. Soit  $\mathcal{J} \subset \mathcal{O}_X$  un faisceau
quasi-cohérent d'idéaux. Soit  $x$  un point du support de  $\mathcal{O}_X/\mathcal{J}$.
Les assertions suivantes sont équivalentes
\begin{enumerate}
\item $\mathcal{J}_x$  est généré par une suite régulière dans  $\mathcal{O}_{X, x}$,
\item $\mathcal{J}_x$  est généré par une suite Koszul-régulière dans  $\mathcal{O}_{X, x}$,
\item $\mathcal{J}_x$  est généré par une suite  $H_1$-régulière dans  $\mathcal{O}_{X, x}$,
\item $\mathcal{J}_x$  est généré par une suite quasi-régulière dans  $\mathcal{O}_{X, x}$,
\item il existe un voisinage affine  $U = \Spec(A)$  de  $x$  tel que  $\mathcal{J}|_U = \widetilde{I}$  et
 $I$  sont générés par une suite régulière dans  $A$, et
\item il existe un voisinage affine  $U = \Spec(A)$  de  $x$  tel que  $\mathcal{J}|_U = \widetilde{I}$  et
 $I$  sont générés par une suite Koszul-régulière dans  $A$, et
\item il existe un voisinage affine  $U = \Spec(A)$  de  $x$  tel que  $\mathcal{J}|_U = \widetilde{I}$  et
 $I$  sont générés par une suite  $H_1$-régulière dans  $A$,
et
\item il existe un voisinage affine  $U = \Spec(A)$  de  $x$  tel que  $\mathcal{J}|_U = \widetilde{I}$  et
 $I$  soit généré par une suite quasi-régulière dans  $A$,
\item il existe un voisinage  $U$  de  $x$  tel que  $\mathcal{J}|_U$  soit
régulier, et
\item il existe un voisinage  $U$  de  $x$  tel que  $\mathcal{J}|_U$  soit
Koszul-régulier, et
\item il existe un voisinage  $U$  de  $x$  tel que  $\mathcal{J}|_U$  soit
 $H_1$-régulier, et
\item il existe un voisinage  $U$  de  $x$  tel que  $\mathcal{J}|_U$  soit
quasi-régulier.
\end{enumerate}
En particulier, sur un schéma localement noethérien, les notions de faisceau
d'idéaux régulier, Koszul-régulier,  $H_1$-régulier, ou quasi-régulier
coïncident toutes.
\end{lemma}

\begin{proof}
Il résulte du Lemme  \ref{divisors-lemma-regular-ideal-sheaf-scheme}  que (5)  $\Leftrightarrow$  (9), (6)  $\Leftrightarrow$  (10),
(7)  $\Leftrightarrow$  (11), et (8)  $\Leftrightarrow$  (12). Il est clair que (5)  $\Rightarrow$ 
(1), (6)  $\Rightarrow$  (2), (7)  $\Rightarrow$  (3), et (8)  $\Rightarrow$  (4). Nous
avons (1)  $\Rightarrow$  (5) par Algèbre, Lemme  \ref{algebra-lemma-regular-sequence-in-neighbourhood}. Nous avons (9)
 $\Rightarrow$  (10)  $\Rightarrow$  (11)  $\Rightarrow$  (12) par le Lemme  \ref{divisors-lemma-regular-quasi-regular-scheme}.
Enfin, (4)  $\Rightarrow$  (1) par Algèbre, Lemme  \ref{algebra-lemma-quasi-regular-regular}. Maintenant, tous les
12 énoncés sont équivalents.
\end{proof}













\section{Immersions régulières}
\label{divisors-section-regular-immersions}

\noindent
Soit  $i : Z \to X$  une immersion de schémas. Par définition, cela signifie qu'il
existe un sous-schéma ouvert  $U \subset X$  tel que  $Z$  soit identifié à un
sous-schéma fermé de  $U$. Soit  $\mathcal{I} \subset \mathcal{O}_U$  le faisceau quasi-cohérent
d'idéaux correspondant. Supposons que  $U' \subset X$  soit un second sous-schéma
ouvert de ce type, et notons  $\mathcal{I}' \subset \mathcal{O}_{U'}$  le faisceau quasi-cohérent d'idéaux
correspondant. Alors  $\mathcal{I}|_{U \cap U'} = \mathcal{I}'|_{U \cap U'}$. De plus, le support de  $\mathcal{O}_U/\mathcal{I}$  est
 $Z$  qui est contenu dans  $U \cap U'$  et est également le support de
 $\mathcal{O}_{U'}/\mathcal{I}'$. Il s'ensuit donc d'après la Définition  \ref{divisors-definition-regular-ideal-sheaf}  que  $\mathcal{I}$ 
est un idéal régulier si et seulement si  $\mathcal{I}'$  est un idéal régulier. De
même pour être Koszul-régulier,  $H_1$-régulier, ou quasi-régulier.

\begin{definition}
\label{divisors-definition-regular-immersion}
\begin{reference}
Le concept d'une immersion Koszul-régulière a été introduit dans  \cite[Expose VII, Definition 1.4]{SGA6}  où
il était appelé une immersion régulière.
\end{reference}
Soit  $i : Z \to X$  une immersion de schémas. Choisissons un sous-schéma ouvert
 $U \subset X$  tel que  $i$  identifie  $Z$  avec un sous-schéma fermé
de  $U$  et désignons par  $\mathcal{I} \subset \mathcal{O}_U$  le faisceau quasi-cohérent d'idéaux
correspondant.
\begin{enumerate}
\item Nous disons que  $i$  est une  {\it  immersion
régulière}  si  $\mathcal{I}$  est régulière.
\item Nous disons que  $i$  est une  {\it  immersion
Koszul-régulière}  si  $\mathcal{I}$  est Koszul-régulière.
\item Nous disons que  $i$  est une  {\it  immersion
 $H_1$-régulière}  si  $\mathcal{I}$  est  $H_1$-régulière.
\item Nous disons que  $i$  est une  {\it  immersion
quasi-régulière}  si  $\mathcal{I}$  est quasi-régulière.
\end{enumerate}
\end{definition}

\noindent
La discussion ci-dessus montre que cela est indépendant du choix de  $U$.
Les conditions sont énumérées dans l'ordre décroissant de force, voir le Lemme
 \ref{divisors-lemma-regular-quasi-regular-immersion}. Une immersion fermée de Koszul-régulière est lisse localement une
immersion régulière, voir le Lemme  \ref{divisors-lemma-koszul-regular-smooth-locally-regular}. Dans le cas localement
noethérien, les quatre notions coïncident, voir le Lemme  \ref{divisors-lemma-Noetherian-scheme-regular-ideal}.

\begin{lemma}
\label{divisors-lemma-regular-quasi-regular-immersion}
Soit  $i : Z \to X$  une immersion de schémas. Nous avons les implications suivantes:
 $i$  est régulière  $\Rightarrow$   $i$  est Koszul-régulière
 $\Rightarrow$   $i$  est  $H_1$-régulière  $\Rightarrow$   $i$  est
quasi-régulière.
\end{lemma}

\begin{proof}
Le lemme se réduit immédiatement au Lemme  \ref{divisors-lemma-regular-quasi-regular-scheme}.
\end{proof}

\begin{lemma}
\label{divisors-lemma-regular-immersion-noetherian}
Soit  $i : Z \to X$  une immersion de schémas. Supposons que  $X$  soit
localement noethérien. Alors  $i$  est régulier  $\Leftrightarrow$   $i$ 
est de Koszul-régulier  $\Leftrightarrow$   $i$  est  $H_1$-régulier
 $\Leftrightarrow$   $i$  est quasi-régulier.
\end{lemma}

\begin{proof}
Suit immédiatement du Lemme  \ref{divisors-lemma-regular-quasi-regular-immersion}  et du Lemme  \ref{divisors-lemma-Noetherian-scheme-regular-ideal}.
\end{proof}

\begin{lemma}
\label{divisors-lemma-flat-base-change-regular-immersion}
Soit  $i : Z \to X$  une immersion régulière (resp. \  de
Koszul-régularité,  $H_1$-régulière, quasi-régulière). Soit  $X' \to X$  un
morphisme plat. Alors le changement de base  $i' : Z \times_X X' \to X'$  est une immersion
régulière (resp. \  de Koszul-régularité,  $H_1$-régulière,
quasi-régulière).
\end{lemma}

\begin{proof}
Via le Lemme  \ref{divisors-lemma-regular-ideal-sheaf-scheme}  cela se traduit en les énoncés algébriques dans Algèbre,
Lemmes  \ref{algebra-lemma-flat-increases-depth}  et  \ref{algebra-lemma-flat-base-change-quasi-regular}  et Compléments d'algèbre, Lemme  \ref{more-algebra-lemma-koszul-regular-flat-base-change}.
\end{proof}

\begin{lemma}
\label{divisors-lemma-quasi-regular-immersion}
Soit  $i : Z \to X$  une immersion de schémas. Alors  $i$  est une immersion
quasi-régulière si et seulement si les conditions suivantes sont satisfaites
\begin{enumerate}
\item $i$  est localement de présentation finie,
\item le faisceau conormal  $\mathcal{C}_{Z/X}$  est localement libre de type fini, et
\item l'application (\ref{divisors-equation-conormal-algebra-quotient}) est un isomorphisme.
\end{enumerate}
\end{lemma}

\begin{proof}
Une immersion ouverte est localement de présentation finie. Par conséquent, nous
pouvons remplacer  $X$  par un sous-schéma ouvert  $U \subset X$  tel que
 $i$  identifie  $Z$  avec un sous-schéma fermé de  $U$,
c'est-à-dire que nous pouvons supposer que  $i$  est une immersion fermée.
Soit  $\mathcal{I} \subset \mathcal{O}_X$  le faisceau d'idéaux quasi-cohérents correspondant. Rappelons,
voir Morphismes, Lemme  \ref{morphisms-lemma-closed-immersion-finite-presentation} , que  $\mathcal{I}$  est de type fini si et
seulement si  $i$  est localement de présentation finie. Par conséquent,
l'équivalence découle du Lemme  \ref{divisors-lemma-quasi-regular-ideal}  et du dépliement des définitions.
\end{proof}

\begin{lemma}
\label{divisors-lemma-transitivity-conormal-quasi-regular}
Soient  $Z \to Y \to X$  des immersions de schémas. Supposons que  $Z \to Y$  est
 $H_1$-régulier. Alors la suite canonique de Morphismes, Lemme  \ref{morphisms-lemma-transitivity-conormal}
$$
0 \to i^*\mathcal{C}_{Y/X} \to
\mathcal{C}_{Z/X} \to
\mathcal{C}_{Z/Y} \to 0
$$
est exacte et localement scindée.
\end{lemma}

\begin{proof}
Puisque  $\mathcal{C}_{Z/Y}$  est localement libre de type fini (voir le Lemme  \ref{divisors-lemma-quasi-regular-immersion} 
et le Lemme  \ref{divisors-lemma-regular-quasi-regular-scheme}), il suffit de prouver que la suite est exacte. D'après
ce qui a été prouvé dans Morphismes, Lemme  \ref{morphisms-lemma-transitivity-conormal} , il suffit de montrer que
la première application est injective. En travaillant localement de manière
affine, cela se réduit à la question suivante: Supposons que nous ayons un anneau
 $A$  et des idéaux  $I \subset J \subset A$. Supposons que  $J/I \subset A/I$  soit engendré
par une suite  $H_1$-régulière. Cela implique-t-il que  $I/I^2 \otimes_A A/J \to J/J^2$  est
injective? Notons que  $I/I^2 \otimes_A A/J = I/IJ$. Par conséquent, nous essayons de prouver que
 $I \cap J^2 = IJ$. Ceci est le résultat de Compléments d'algèbre, Lemme  \ref{more-algebra-lemma-conormal-sequence-H1-regular}.
\end{proof}

\noindent
Une composition d'immersions presque régulières peut ne pas être presque
régulière, voir Algèbre, Remarque  \ref{algebra-remark-join-quasi-regular-sequences}. Les autres types d'immersions
régulières sont préservés par composition.

\begin{lemma}
\label{divisors-lemma-composition-regular-immersion}
Soient  $i : Z \to Y$  et  $j : Y \to X$  des immersions de schémas.
\begin{enumerate}
\item Si  $i$  et  $j$  sont des immersions régulières, il en est de même
pour  $j \circ i$.
\item Si  $i$  et  $j$  sont des immersions Koszul-régulières, il en est
de même pour  $j \circ i$.
\item Si  $i$  et  $j$  sont des immersions  $H_1$-régulières, il en
est de même pour  $j \circ i$.
\item Si  $i$  est une immersion  $H_1$-régulière et  $j$  est une
immersion quasi-régulière, alors  $j \circ i$  est une immersion quasi-régulière.
\end{enumerate}
\end{lemma}

\begin{proof}
La version algébrique de (1) est Algèbre, Lemme  \ref{algebra-lemma-join-regular-sequences}. La version
algébrique de (2) est Compléments d'algèbre, Lemme  \ref{more-algebra-lemma-join-koszul-regular-sequences}. La version algébrique
de (3) est Compléments d'algèbre, Lemme  \ref{more-algebra-lemma-join-H1-regular-sequences}. La version algébrique de (4) est
Compléments d'algèbre, Lemme  \ref{more-algebra-lemma-join-quasi-regular-H1-regular}.
\end{proof}

\begin{lemma}
\label{divisors-lemma-permanence-regular-immersion}
Soient  $i : Z \to Y$  et  $j : Y \to X$  des immersions de schémas. Supposons que
 $j$  soit localement de présentation finie et que la séquence
$$
0 \to i^*\mathcal{C}_{Y/X} \to
\mathcal{C}_{Z/X} \to
\mathcal{C}_{Z/Y} \to 0
$$
des morphismes, le Lemme  \ref{morphisms-lemma-transitivity-conormal}  est exact et localement scindé.
\begin{enumerate}
\item Si  $j \circ i$  est une immersion quasi-régulière, il en est de même pour
 $i$.
\item Si  $j \circ i$  est une immersion  $H_1$-régulière, il en est de même pour
 $i$.
\item Si  $j$  et  $j \circ i$  sont toutes deux des immersions de
Koszul-régulières, il en est de même pour  $i$.
\end{enumerate}
\end{lemma}

\begin{proof}
Après avoir restreint  $Y$  et  $X$ , nous pouvons supposer que
 $i$  et  $j$  sont des immersions fermées. Notons  $\mathcal{I} \subset \mathcal{O}_X$  le
faisceau d'idéaux de  $Y$  et  $\mathcal{J} \subset \mathcal{O}_X$  le faisceau d'idéaux de
 $Z$. La suite conormale est  $0 \to \mathcal{I}/\mathcal{I}\mathcal{J}
\to \mathcal{J}/\mathcal{J}^2 \to
\mathcal{J}/(\mathcal{I} + \mathcal{J}^2) \to 0$. Soient  $z \in Z$  et posons
 $y = i(z)$,  $x = j(y) = j(i(z))$. Choisissons  $f_1, \ldots, f_n \in \mathcal{I}_x$  qui s'envoient sur une
base de  $\mathcal{I}_x/\mathfrak m_z\mathcal{I}_x$. Étendons cette famille à  $f_1, \ldots, f_n, g_1, \ldots, g_m \in \mathcal{J}_x$  qui s'envoient sur une base
de  $\mathcal{J}_x/\mathfrak m_z\mathcal{J}_x$. Cela est possible car nous avons supposé que la suite des
faisceaux conormaux est scindée dans un voisinage de  $z$, donc
 $\mathcal{I}_x/\mathfrak m_x\mathcal{I}_x \to
\mathcal{J}_x/\mathfrak m_x\mathcal{J}_x$  est injective.

\medskip\noindent
Preuve de (1). D'après le Lemme  \ref{divisors-lemma-generate-regular-ideal} , nous pouvons trouver un voisinage
affine ouvert  $U$  de  $x$  tel que  $f_1, \ldots, f_n, g_1, \ldots, g_m$  forme une suite
quasi-régulière générant  $\mathcal{J}$. Puisque  $i$  est localement de
présentation finie, nous pouvons en outre supposer que  $\mathcal{I}$  est engendré
par  $f_1, \ldots, f_n$  (utiliser Morphismes, Lemme  \ref{morphisms-lemma-closed-immersion-finite-presentation}, le lemme de Nakayama,
et Modules, Lemme  \ref{modules-lemma-finite-type-surjective-on-stalk}). Ainsi, d'après Algèbre, Lemme  \ref{algebra-lemma-truncate-quasi-regular} , nous
voyons que  $g_1, \ldots, g_m$  induit une suite quasi-régulière sur  $Y \cap U$ 
définissant  $Z$.

\medskip\noindent
Preuve de (2). Exactement la même que la preuve de (1), en utilisant toutefois Compléments
d'algèbre, Lemme  \ref{more-algebra-lemma-truncate-H1-regular}.

\medskip\noindent
Preuve de (3). D'après le Lemme  \ref{divisors-lemma-generate-regular-ideal}  (appliqué deux fois), nous pouvons
trouver un voisinage affine ouvert  $U$  de  $x$  tel que
 $f_1, \ldots, f_n$  forme une suite Koszul-régulière générant  $\mathcal{I}$  et
 $f_1, \ldots, f_n, g_1, \ldots, g_m$  forme une suite Koszul-régulière générant  $\mathcal{J}$. Par
conséquent, d'après Compléments d'algèbre, Lemme  \ref{more-algebra-lemma-truncate-koszul-regular} , nous voyons que
 $g_1, \ldots, g_m$  induit une suite Koszul-régulière sur  $Y \cap U$  coupant
 $Z$.
\end{proof}

\begin{lemma}
\label{divisors-lemma-extra-permanence-regular-immersion-noetherian}
Soient  $i : Z \to Y$  et  $j : Y \to X$  des immersions de schémas. Choisissons
 $z \in Z$  et désignons par  $y \in Y$,  $x \in X$  les points correspondants.
Supposons que  $X$  soit localement noethérien. Les conditions suivantes
sont équivalentes
\begin{enumerate}
\item $i$  est une immersion régulière dans un voisinage de  $z$  et
 $j$  est une immersion régulière dans un voisinage de  $y$.
\item $i$  et  $j \circ i$  sont des immersions régulières dans un voisinage de
 $z$,
\item $j \circ i$  est une immersion régulière dans un voisinage de  $z$  et la
suite conormale
$$
0 \to i^*\mathcal{C}_{Y/X} \to
\mathcal{C}_{Z/X} \to
\mathcal{C}_{Z/Y} \to 0
$$
est exacte scindée dans un voisinage de  $z$.
\end{enumerate}
\end{lemma}

\begin{proof}
Puisque  $X$  (et donc  $Y$) est localement noethérien, les 4
types d'immersions régulières coïncident, et de plus, nous pouvons vérifier si un
morphisme est une immersion régulière au niveau des anneaux locaux, voir le Lemme
 \ref{divisors-lemma-Noetherian-scheme-regular-ideal}. L'implication (1)  $\Rightarrow$  (2) est le Lemme  \ref{divisors-lemma-composition-regular-immersion}.
L'implication (2)  $\Rightarrow$  (3) est le Lemme  \ref{divisors-lemma-transitivity-conormal-quasi-regular}. Ainsi, il suffit de
prouver que (3) implique (1).

\medskip\noindent
Supposons (3). Posons  $A = \mathcal{O}_{X, x}$. Désignons par  $I \subset A$  le noyau de
l'application surjective  $\mathcal{O}_{X, x} \to \mathcal{O}_{Y, y}$  et par  $J \subset A$  le noyau de
l'application surjective  $\mathcal{O}_{X, x} \to \mathcal{O}_{Z, z}$. Notons que toute suite minimale d'éléments
générant  $J$  dans  $A$  est une suite quasi-régulière donc
régulière, voir le Lemme  \ref{divisors-lemma-generate-regular-ideal}. Par hypothèse, la suite conormale
$$
0 \to I/IJ \to J/J^2 \to J/(I + J^2) \to 0
$$
est exactement scindé en une suite de  $A/J$-modules. Par conséquent, nous
pouvons choisir un système minimal de générateurs  $f_1, \ldots, f_n, g_1, \ldots, g_m$  de  $J$ 
avec  $f_1, \ldots, f_n \in I$  un système minimal de générateurs de  $I$. Comme indiqué
ci-dessus,  $f_1, \ldots, f_n, g_1, \ldots, g_m$  est une suite régulière dans  $A$. Il ressort
directement de la définition d'une suite régulière que  $f_1, \ldots, f_n$  est une suite
régulière dans  $A$  et  $\overline{g}_1, \ldots, \overline{g}_m$  est une suite régulière dans
 $A/I$. Ainsi,  $j$  est une immersion régulière en  $y$  et
 $i$  est une immersion régulière en  $z$.
\end{proof}

\begin{remark}
\label{divisors-remark-not-always-extra-permanence}
Dans la situation du Lemme  \ref{divisors-lemma-extra-permanence-regular-immersion-noetherian} , les parties (1), (2), (3) sont
 {\bf  non }  équivalentes à « $j \circ i$  et
 $j$  sont des immersions régulières en  $z$  et  $y$ ». Un
exemple est  $X = \mathbf{A}^1_k = \Spec(k[x])$,  $Y = \Spec(k[x]/(x^2))$  et  $Z = \Spec(k[x]/(x))$.
\end{remark}

\begin{lemma}
\label{divisors-lemma-koszul-regular-smooth-locally-regular}
Soit  $i : Z \to X$  une immersion fermée koszul-régulière. Alors il existe un
morphisme lisse surjectif  $X' \to X$  tel que le changement de base  $i' : Z \times_X X' \to X'$ 
de  $i$  soit une immersion régulière.
\end{lemma}

\begin{proof}
Nous pouvons supposer que  $X$  est affine et que l'idéal de  $Z$ 
est engendré par une suite koszul-régulière en remplaçant  $X$  par les
éléments d'un recouvrement ouvert affine approprié (ouverts affines comme dans le
Lemme  \ref{divisors-lemma-regular-ideal-sheaf-scheme}). Le cas affine est Développement sur Algèbre, Lemme
 \ref{more-algebra-lemma-Koszul-regular-flat-locally-regular}.
\end{proof}

\begin{lemma}
\label{divisors-lemma-immersion-regular-regular-immersion}
Soit  $i : Z \to X$  une immersion. Si  $Z$  et  $X$  sont des schémas
réguliers, alors  $i$  est une immersion régulière.
\end{lemma}

\begin{proof}
Soit  $z \in Z$. D'après le Lemme  \ref{divisors-lemma-Noetherian-scheme-regular-ideal} , il suffit de montrer que le
noyau de  $\mathcal{O}_{X, z} \to \mathcal{O}_{Z, z}$  est engendré par une suite régulière. Cela découle des lemmes
d'Algèbre  \ref{algebra-lemma-regular-quotient-regular}  et  \ref{algebra-lemma-regular-ring-CM}.
\end{proof}





\section{Immersions régulières relatives}
\label{divisors-section-relative-regular-immersion}

\noindent
Dans cette section, nous considérons la propriété de changement de base pour les
immersions régulières. Le lemme suivant ne vaut pas pour les immersions
régulières ni pour les immersions Koszul-régulières, voir Exemples, Lemme  \ref{examples-lemma-base-change-regular-sequence}.

\begin{lemma}
\label{divisors-lemma-relative-regular-immersion}
Soit  $f : X \to S$  un morphisme de schémas. Soit  $i : Z \subset X$  une immersion.
Supposons
\begin{enumerate}
\item $i$  est un  $H_1$-régulier (resp.\ immersion
quasi-régulière, et
\item $Z \to S$  est un morphisme plat.
\end{enumerate}
Alors, pour tout morphisme de schémas  $g : S' \to S$ , le changement de base
 $Z' = S' \times_S Z \to X' = S' \times_S X$  est une immersion  $H_1$-régulière (resp. quasi-régulière)
\ .
\end{lemma}

\begin{proof}
En déroulant les définitions et en utilisant le Lemme  \ref{divisors-lemma-regular-ideal-sheaf-scheme} , cela se
traduit par Compléments d'algèbre, Lemme  \ref{more-algebra-lemma-relative-regular-immersion-algebra}.
\end{proof}

\noindent
Ce lemme est la motivation pour la définition suivante.

\begin{definition}
\label{divisors-definition-relative-H1-regular-immersion}
Soit  $f : X \to S$  un morphisme de schémas. Soit  $i : Z \to X$  une immersion.
\begin{enumerate}
\item On dit que  $i$  est une  {\it  immersion quasi-régulière
relative }  si  $Z \to S$  est plate et  $i$  est une immersion
quasi-régulière.
\item On dit que  $i$  est une  {\it  immersion
 $H_1$-régulière relative }  si  $Z \to S$  est plate et
 $i$  est une immersion  $H_1$-régulière.
\end{enumerate}
\end{definition}

\noindent
Nous avertissons le lecteur que cela peut être une notation non standard. Le lemme
 \ref{divisors-lemma-relative-regular-immersion}  garantit que les immersions quasi-régulières relatives (resp.
\ $H_1$-régulières) sont préservées sous tout changement de base.
Une immersion relative  $H_1$-régulière est une immersion quasi-régulière
relative, voir le lemme  \ref{divisors-lemma-regular-quasi-regular-immersion}. Voir le lemme  \ref{divisors-lemma-flat-relative-H1-regular}  (ou
le lemme  \ref{divisors-lemma-relative-regular-immersion-flat-in-neighbourhood}) qui montre que si  $Z \to X$  est une immersion relative
 $H_1$-régulière (ou quasi-régulière) et que le schéma ambiant est (plat et)
localement de présentation finie sur  $S$, alors  $Z \to X$  est en fait
une immersion régulière et il en reste de même après tout changement de base.

\begin{lemma}
\label{divisors-lemma-quasi-regular-immersion-flat-at-x}
Soit  $f : X \to S$  un morphisme de schémas. Soit  $Z \to X$  une immersion
quasi-régulière relative. Si  $x \in Z$  et  $\mathcal{O}_{X, x}$  sont noethériens, alors
 $f$  est plat en  $x$.
\end{lemma}

\begin{proof}
Soit  $f_1, \ldots, f_r \in \mathcal{O}_{X, x}$  une suite quasi-régulière définissant l'idéal de  $Z$  en
 $x$. D'après le Lemme d'Algèbre  \ref{algebra-lemma-quasi-regular-regular} , nous savons que
 $f_1, \ldots, f_r$  est une suite régulière. Par conséquent,  $f_r$  est un élément
non-diviseur de zéro dans  $\mathcal{O}_{X, x}/(f_1, \ldots, f_{r - 1})$  tel que le quotient est un
 $\mathcal{O}_{S, f(x)}$-module plat. D'après le Lemme  \ref{divisors-lemma-flat-at-x} , nous en concluons que
 $\mathcal{O}_{X, x}/(f_1, \ldots, f_{r - 1})$  est un  $\mathcal{O}_{S, f(x)}$-module plat. En continuant par induction, nous
trouvons que  $\mathcal{O}_{X, x}$  est un  $\mathcal{O}_{S, s}$-module plat.
\end{proof}

\begin{lemma}
\label{divisors-lemma-relative-regular-immersion-flat-in-neighbourhood}
Soit  $X \to S$  un morphisme de schémas. Soit  $Z \to X$  une immersion.
Supposons
\begin{enumerate}
\item que $X \to S$  est plat et localement de présentation finie,
\item que $Z \to X$  est une immersion quasi-régulière relative.
\end{enumerate}
Alors  $Z \to X$  est une immersion régulière et la même chose reste vraie après
tout changement de base.
\end{lemma}

\begin{proof}
Choisissons  $x \in Z$  avec l'image  $s \in S$. Pour le prouver, il suffit de
trouver un voisinage affine de  $x$  contenu dans  $U$  tel que le
résultat soit valide sur cet ouvert affine. Ainsi, nous pouvons supposer que
 $X$  est affine et qu'il existe une suite quasi-régulière  $f_1, \ldots, f_r \in \Gamma(X, \mathcal{O}_X)$ 
telle que  $Z = V(f_1, \ldots, f_r)$. D'après Compléments d'algèbre, Lemme  \ref{more-algebra-lemma-relative-regular-immersion-algebra} , la suite
 $f_1|_{X_s}, \ldots, f_r|_{X_s}$  est une suite quasi-régulière dans  $\Gamma(X_s, \mathcal{O}_{X_s})$. Comme  $X_s$ 
est noethérien, cela implique, éventuellement après avoir légèrement réduit
 $X$ , que  $f_1|_{X_s}, \ldots, f_r|_{X_s}$  est une suite régulière, voir Algèbre, Lemmes
 \ref{algebra-lemma-quasi-regular-regular}  et  \ref{algebra-lemma-regular-sequence-in-neighbourhood}. D'après le Lemme  \ref{divisors-lemma-fibre-Cartier} , il découle que
 $Z_1 = V(f_1) \subset X$  est un diviseur de Cartier effectif relatif, encore après avoir
éventuellement un peu rétréci  $X$ . En appliquant à nouveau le même lemme,
mais cette fois à  $Z_2 = V(f_1, f_2) \subset Z_1$ , nous voyons que  $Z_2 \subset Z_1$  est un diviseur de
Cartier effectif relatif. Et ainsi de suite jusqu'à ce que l'on atteigne
 $Z = Z_n = V(f_1, \ldots, f_n)$. Comme être un diviseur de Cartier effectif relatif se conserve par
tout changement de base arbitraire, voir le Lemme  \ref{divisors-lemma-relative-Cartier}, nous voyons
également que l'énoncé final du lemme est vrai.
\end{proof}

\begin{remark}
\label{divisors-remark-relative-regular-immersion-elements}
La codimension d'une immersion quasi-régulière relative, si elle est constante, ne
change pas après un changement de base. En fait, si nous avons un morphisme
d'anneaux  $A \to B$  et une suite quasi-régulière  $f_1, \ldots, f_r \in B$  telle que
 $B/(f_1, \ldots, f_r)$  est plate sur  $A$, alors pour tout morphisme d'anneaux
 $A \to A'$  nous avons une suite quasi-régulière  $f_1 \otimes 1, \ldots, f_r \otimes 1$  dans  $B' = B \otimes_A A'$ 
selon Compléments d'algèbre, Lemme  \ref{more-algebra-lemma-relative-regular-immersion-algebra}  (qui a été utilisé dans la démonstration
du Lemme  \ref{divisors-lemma-relative-regular-immersion}  ci-dessus). Maintenant, la démonstration du Lemme
 \ref{divisors-lemma-relative-regular-immersion-flat-in-neighbourhood}  montre que si  $A \to B$  est plate et localement de présentation
finie, alors pour tout idéal premier  $\mathfrak q' \subset B'$  la suite  $f_1 \otimes 1, \ldots, f_r \otimes 1$  est même
une suite régulière dans l'anneau local  $B'_{\mathfrak q'}$.
\end{remark}

\begin{lemma}
\label{divisors-lemma-flat-relative-H1-regular}
Soit  $X \to S$  un morphisme de schémas. Soit  $Z \to X$  une immersion
 $H_1$-régulière relative. Supposons que  $X \to S$  soit localement de
présentation finie. Alors
\begin{enumerate}
\item il existe un sous-schéma ouvert  $U \subset X$  tel que  $Z \subset U$  et tel que
 $U \to S$  soit plat, et
\item $Z \to X$  est une immersion régulière et il en reste de même après tout
changement de base.
\end{enumerate}
\end{lemma}

\begin{proof}
Choisissons  $x \in Z$. Pour prouver (1), il suffit de trouver un voisinage ouvert
 $U \subset X$  de  $x$  tel que  $U \to S$  soit plat. Ainsi, le lemme se
réduit au cas où  $X = \Spec(B)$  et  $S = \Spec(A)$  sont affines et où  $Z$  est
donné par une suite  $H_1$-régulière  $f_1, \ldots, f_r \in B$. Par hypothèse,
 $B$  est une  $A$-algèbre de présentation finie et  $B/(f_1, \ldots, f_r)B$ 
est une  $A$-algèbre plate. Nous allons utiliser l'approximation
noethérienne absolue.

\medskip\noindent
Écrivons  $B = A[x_1, \ldots, x_n]/(g_1, \ldots, g_m)$. Supposons que  $f_i$  soit l'image de  $f_i' \in A[x_1, \ldots, x_n]$.
Choisissons une  $\mathbf{Z}$-sous-algèbre de type fini  $A_0 \subset A$  telle que tous
les coefficients des polynômes  $f_1', \ldots, f_r', g_1, \ldots, g_m$  soient dans  $A_0$. Nous
définissons  $B_0 = A_0[x_1, \ldots, x_n]/(g_1, \ldots, g_m)$  et nous notons  $f_{i, 0}$  l'image de  $f_i'$  dans
 $B_0$. Ensuite  $B = B_0 \otimes_{A_0} A$  et
$$
B/(f_1, \ldots, f_r) =
B_0/(f_{0, 1}, \ldots, f_{0, r}) \otimes_{A_0} A.
$$
Par Algèbre, Lemme  \ref{algebra-lemma-flat-finite-presentation-limit-flat} , nous pouvons, après avoir agrandi  $A_0$,
supposer que  $B_0/(f_{0, 1}, \ldots, f_{0, r})$  est plat sur  $A_0$. Il se peut qu'à ce moment le
groupe de cohomologie de Koszul  $H_1(K_\bullet(B_0, f_{0, 1}, \ldots, f_{0, r}))$  ne soit pas nul. D'un autre côté,
comme  $B_0$  est Noethérien, c'est un  $B_0$-module de type fini.
Soient  $\xi_1, \ldots, \xi_n \in H_1(K_\bullet(B_0, f_{0, 1}, \ldots, f_{0, r}))$  des générateurs. Soit  $A_0 \subset A_1 \subset A$  une
 $\mathbf{Z}$-sous-algèbre de type fini plus grande de  $A$. Notons
 $f_{1, i}$  l'image de  $f_{0, i}$  dans  $B_1 = B_0 \otimes_{A_0} A_1$. Par Compléments d'algèbre,
Lemme  \ref{more-algebra-lemma-base-change-H1-regular} , l'application
$$
H_1(K_\bullet(B_0, f_{0, 1}, \ldots, f_{0, r})) \otimes_{A_0} A_1
\longrightarrow
H_1(K_\bullet(B_1, f_{1, 1}, \ldots, f_{1, r}))
$$
est surjective. De plus, il est clair que le colimite (sur toutes les choix de
 $A_1$  comme ci-dessus) des complexes  $K_\bullet(B_1, f_{1, 1}, \ldots, f_{1, r})$  est le complexe
 $K_\bullet(B, f_1, \ldots, f_r)$  qui est acyclique en degré  $1$. Par conséquent,
$$
\colim_{A_0 \subset A_1 \subset A}
H_1(K_\bullet(B_1, f_{1, 1}, \ldots, f_{1, r}))
= 0
$$
d'après Algèbre, Lemme  \ref{algebra-lemma-directed-colimit-exact}. Ainsi, nous pouvons trouver un choix de
 $A_1$  tel que  $\xi_1, \ldots, \xi_n$  s'envoient tous sur zéro dans  $H_1(K_\bullet(B_1, f_{1, 1}, \ldots, f_{1, r}))$. En
d'autres termes, le groupe de cohomologie de Koszul  $H_1(K_\bullet(B_1, f_{1, 1}, \ldots, f_{1, r}))$  est nul.

\medskip\noindent
Considérons le morphisme de schémas affines  $X_1 \to S_1$  égal à  $\Spec$  du
morphisme d'anneaux  $A_1 \to B_1$  et  $Z_1 = \Spec(B_1/(f_{1, 1}, \ldots, f_{1, r}))$. Puisque  $B = B_1 \otimes_{A_1} A$,
c'est-à-dire  $X = X_1 \times_{S_1} S$, et de même  $Z = Z_1 \times_S S_1$, il suffit maintenant de
prouver (1) pour  $X_1 \to S_1$  et l'immersion  $H_1$-relative
 $Z_1 \to X_1$-régulière, voir Morphismes, Lemme  \ref{morphisms-lemma-base-change-module-flat}. Nous avons donc
réduit au cas où  $X \to S$  est un morphisme de type fini de schémas de Noether.
Dans ce cas, nous savons que  $X \to S$  est plat en tout point de  $Z$ 
d'après le Lemme  \ref{divisors-lemma-quasi-regular-immersion-flat-at-x}. Combiné avec le fait que le lieu plat est ouvert
dans ce cas, voir Algèbre, Théorème  \ref{algebra-theorem-openness-flatness} , nous voyons que (1) est vrai. La
partie (2) découle alors d'une application du Lemme  \ref{divisors-lemma-relative-regular-immersion-flat-in-neighbourhood}.
\end{proof}

\noindent
Si le schéma ambiant est plat et localement de présentation finie sur la base,
alors nous pouvons caractériser une immersion quasi-régulière relative en termes
de ses fibres.

\begin{lemma}
\label{divisors-lemma-fibre-quasi-regular}
Soit  $\varphi : X \to S$  un morphisme plat qui est localement de présentation finie. Soit
 $T \subset X$  un sous-schéma fermé. Soit  $x \in T$  avec une image  $s \in S$.
\begin{enumerate}
\item Si  $T_s \subset X_s$  est une immersion quasi-régulière dans un voisinage de
 $x$, alors il existe un ouvert  $U \subset X$  et une immersion
quasi-régulière relative  $Z \subset U$  tels que  $Z_s = T_s \cap U_s$  et  $T \cap U \subset Z$.
\item Si  $T_s \subset X_s$  est une immersion quasi-régulière dans un voisinage de
 $x$, le morphisme  $T \to X$  est de présentation finie, et  $T \to S$ 
est plat en  $x$, alors nous pouvons choisir  $U$  et  $Z$ 
comme dans (1) tels que  $T \cap U = Z$.
\item Si  $T_s \subset X_s$  est une immersion quasi-régulière dans un voisinage de
 $x$, et  $T$  est définie par  $c$  équations dans un
voisinage de  $x$, où  $c = \dim_x(X_s) - \dim_x(T_s)$, alors nous pouvons choisir
 $U$  et  $Z$  comme dans (1) de telle sorte que  $T \cap U = Z$.
\end{enumerate}
Dans chaque cas,  $Z \to U$  est une immersion régulière selon le Lemme
 \ref{divisors-lemma-relative-regular-immersion-flat-in-neighbourhood}. En particulier, si  $T \to S$  est localement de présentation finie
et plate et que toutes les fibres  $T_s \subset X_s$  sont des immersions
quasi-régulières, alors  $T \to X$  est une immersion quasi-régulière relative.
\end{lemma}

\begin{proof}
Choisissons des voisinages ouverts affines  $\Spec(A)$  de  $s$  et
 $\Spec(B)$  de  $x$  tels que  $\varphi(\Spec(B)) \subset \Spec(A)$. Soit  $\mathfrak p \subset A$  l'idéal
premier correspondant à  $s$. Soit  $\mathfrak q \subset B$  l'idéal premier
correspondant à  $x$. Soit  $I \subset B$  l'idéal correspondant à
 $T$. D'après l'hypothèse initiale du lemme, nous savons que  $A \to B$ 
est plat et de présentation finie. L'hypothèse en (1) signifie que, après avoir
rétréci  $\Spec(B)$, nous pouvons supposer que  $I(B \otimes_A \kappa(\mathfrak p))$  est engendré par une
suite quasi-régulière d'éléments. Après avoir éventuellement localisé  $B$ 
en certains  $g \in B$,  $g \not \in \mathfrak q$ , nous pouvons supposer qu'il existe
 $f_1, \ldots, f_r \in I$  qui s'envoient sur une suite quasi-régulière dans  $B \otimes_A \kappa(\mathfrak p)$  qui
génère  $I(B \otimes_A \kappa(\mathfrak p))$. Par l'Algèbre, les lemmes  \ref{algebra-lemma-quasi-regular-regular}  et  \ref{algebra-lemma-regular-sequence-in-neighbourhood} , nous
pouvons supposer qu'après une autre localisation,  $f_1, \ldots, f_r \in I$  forment une suite
régulière dans  $B \otimes_A \kappa(\mathfrak p)$. Selon le lemme  \ref{divisors-lemma-fibre-Cartier} , il s'ensuit que
 $Z_1 = V(f_1) \subset \Spec(B)$  est un diviseur de Cartier effectif relatif, encore une fois après
avoir éventuellement localisé  $B$. En appliquant à nouveau le même lemme,
mais maintenant à  $Z_2 = V(f_1, f_2) \subset Z_1$ , nous voyons que  $Z_2 \subset Z_1$  est un diviseur de
Cartier effectif relatif. Et ainsi de suite jusqu'à atteindre  $Z = Z_n = V(f_1, \ldots, f_n)$.
Ensuite,  $Z \to \Spec(B)$  est une immersion régulière et  $Z$  est plat sur
 $S$, en particulier  $Z \to \Spec(B)$  est une immersion quasi-régulière
relative sur  $\Spec(A)$. Cela prouve (1).

\medskip\noindent
Pour voir (2), considérons l'immersion fermée  $Z \to D$. L'application
surjective d'anneaux  $u : \mathcal{O}_{D, x} \to \mathcal{O}_{Z, x}$  est une application d'algèbres locales plates
 $\mathcal{O}_{S, s}$ qui sont essentiellement de présentation finie, et qui devient un
isomorphisme après division par  $\mathfrak m_s$. Par conséquent, c'est un
isomorphisme, voir Algèbre, Lemme  \ref{algebra-lemma-mod-injective-general}. Il s'ensuit que  $Z \to D$  est
un isomorphisme dans un voisinage de  $x$, voir Algèbre, Lemme
 \ref{algebra-lemma-local-isomorphism}.

\medskip\noindent
Pour voir (3), après avoir éventuellement restreint  $U$ , nous pouvons
supposer que l'idéal de  $Z$  est engendré par une suite régulière
 $f_1, \ldots, f_r$  (voir notre construction de  $Z$  ci-dessus) et que l'idéal de
 $T$  est engendré par  $g_1, \ldots, g_c$. Nous affirmons que  $c = r$. À
savoir,
\begin{align*}
\dim_x(X_s) & = \dim(\mathcal{O}_{X_s, x}) +
\text{trdeg}_{\kappa(s)}(\kappa(x)), \\
\dim_x(T_s) & = \dim(\mathcal{O}_{T_s, x}) +
\text{trdeg}_{\kappa(s)}(\kappa(x)), \\
\dim(\mathcal{O}_{X_s, x}) & = \dim(\mathcal{O}_{T_s, x}) + r
\end{align*}
les deux premières égalités par l'Algèbre, Lemme  \ref{algebra-lemma-dimension-at-a-point-finite-type-field}  et la seconde par
 $r$  fois en appliquant l'Algèbre, Lemme  \ref{algebra-lemma-one-equation}. Comme  $T \subset Z$ 
nous voyons que  $f_i = \sum b_{ij} g_j$. Mais les idéaux de  $Z$  et  $T$ 
définissent le même sous-schéma fermé quasi-régulier de  $X_s$  dans un
voisinage de  $x$. Par conséquent, la matrice  $(b_{ij}) \bmod \mathfrak m_x$  est inversible
(certains détails omis). Ainsi  $(b_{ij})$  est inversible dans un voisinage
ouvert de  $x$. En d'autres termes,  $T \cap U = Z$  après avoir restreint
 $U$.

\medskip\noindent
Les assertions finales du lemme suivent immédiatement de la partie (2),
combinées avec le fait que  $Z \to S$  est localement de présentation finie si et
seulement si  $Z \to X$  est de présentation finie, voir Morphismes, Lemme
 \ref{morphisms-lemma-composition-finite-presentation}  et  \ref{morphisms-lemma-finite-presentation-permanence}.
\end{proof}

\noindent
Le lemme suivant est une amélioration de Morphismes, Lemme  \ref{morphisms-lemma-section-smooth-morphism}.

\begin{lemma}
\label{divisors-lemma-section-smooth-regular-immersion}
Soit  $f : X \to S$  un morphisme lisse de schémas. Soit  $\sigma : S \to X$  une section de
 $f$. Alors  $\sigma$  est une immersion régulière.
\end{lemma}

\begin{proof}
D'après Schémas, Lemme  \ref{schemes-lemma-semi-diagonal}  le morphisme  $\sigma$  est une immersion.
Après avoir remplacé  $X$  par un voisinage ouvert de  $\sigma(S)$ , nous
pouvons supposer que  $\sigma$  est une immersion fermée. Soit  $T = \sigma(S)$  le
sous-schéma fermé correspondant de  $X$. Puisque  $T \to S$  est un
isomorphisme, il est plat et de présentation finie. De plus, un morphisme lisse
est plat et localement de présentation finie, voir Morphismes, Lemme  \ref{morphisms-lemma-smooth-flat} 
et  \ref{morphisms-lemma-smooth-locally-finite-presentation}. Ainsi, selon le Lemme  \ref{divisors-lemma-fibre-quasi-regular}, il suffit de montrer que
 $T_s \subset X_s$  est un sous-schéma fermé quasi-régulier. Cela découle immédiatement
des morphismes, Lemme  \ref{morphisms-lemma-section-smooth-morphism} , mais nous pouvons aussi le voir directement
comme suit. Soit  $k$  un corps et soit  $A$  une
 $k$-algèbre lisse. Soit  $\mathfrak m \subset A$  un idéal maximal dont le corps
résiduel est  $k$. Alors  $\mathfrak m$  est engendré par une suite
quasi-régulière, éventuellement après avoir remplacé  $A$  par  $A_g$ 
pour certains  $g \in A$,  $g \not \in \mathfrak m$. Dans Algèbre, Lemme  \ref{algebra-lemma-characterize-smooth-over-field}  nous
avons prouvé que  $A_{\mathfrak m}$  est un anneau local régulier, donc  $\mathfrak mA_{\mathfrak m}$  est
engendré par une suite régulière. Cela implique en effet que  $\mathfrak m$  est
engendré par une suite régulière (après avoir remplacé  $A$  par
 $A_g$  pour certains  $g \in A$,  $g \not \in \mathfrak m$), voir Algèbre, Lemme
 \ref{algebra-lemma-regular-sequence-in-neighbourhood}.
\end{proof}

\noindent
Le lemme suivant a une sorte de réciproque, voir Lemme  \ref{divisors-lemma-push-regular-immersion-thru-smooth}.

\begin{lemma}
\label{divisors-lemma-lift-regular-immersion-to-smooth}
Soit
$$
\xymatrix{
Y \ar[rd]_j \ar[rr]_i & & X \ar[ld] \\
& S
}
$$
soit un diagramme commutatif de morphismes de schémas. Supposons que  $X \to S$ 
soit lisse, et que  $i$,  $j$  soient des immersions. Si
 $j$  est une immersion régulière (resp. \ Koszul-régulière,
 $H_1$-régulière, quasi-régulière), alors il en est de même pour
 $i$.
\end{lemma}

\begin{proof}
Nous pouvons écrire  $i$  comme la composition
$$
Y \to Y \times_S X \to X
$$
D'après le Lemme  \ref{divisors-lemma-section-smooth-regular-immersion} , la première flèche est une immersion régulière.
La deuxième flèche est un changement de base plat de  $Y \to S$, par conséquent
c'est une immersion régulière (resp. \ Koszul-régulière,
 $H_1$-régulière, quasi-régulière), voir Lemme  \ref{divisors-lemma-flat-base-change-regular-immersion}. Nous concluons
par une application du Lemme  \ref{divisors-lemma-composition-regular-immersion}.
\end{proof}

\begin{lemma}
\label{divisors-lemma-immersion-lci-into-smooth-regular-immersion}
Soit
$$
\xymatrix{
Y \ar[rd] \ar[rr]_i & & X \ar[ld] \\
& S
}
$$
un diagramme commutatif de morphismes de schémas. Supposons que  $Y \to S$  soit
syntomique,  $X \to S$  lisse, et  $i$  une immersion. Alors  $i$ 
est une immersion régulière.
\end{lemma}

\begin{proof}
Après avoir remplacé  $X$  par un voisinage ouvert de  $i(Y)$ , nous
pouvons supposer que  $i$  est une immersion fermée. Soit  $T = i(Y)$  le
sous-schéma fermé correspondant de  $X$. Comme  $T \cong Y$ , le morphisme
 $T \to S$  est plat et de présentation finie (Morphismes, Lemmes  \ref{morphisms-lemma-syntomic-locally-finite-presentation} 
et  \ref{morphisms-lemma-syntomic-flat}). De plus, un morphisme lisse est plat et localement de
présentation finie (Morphismes, Lemmes  \ref{morphisms-lemma-smooth-flat}  et  \ref{morphisms-lemma-smooth-locally-finite-presentation}). Ainsi,
selon le Lemme  \ref{divisors-lemma-fibre-quasi-regular}, il suffit de montrer que  $T_s \subset X_s$  est un
sous-schéma fermé quasi-régulier. Comme  $X_s$  est localement de type fini
sur un corps, il est noethérien (Morphismes, Lemme  \ref{morphisms-lemma-finite-type-noetherian}). Ainsi, nous
pouvons vérifier que  $T_s \subset X_s$  est une immersion quasi-régulière en les points,
voir le Lemme  \ref{divisors-lemma-Noetherian-scheme-regular-ideal}. Prenons  $t \in T_s$. Par Morphismes, Lemme  \ref{morphisms-lemma-local-complete-intersection} ,
l'anneau local  $\mathcal{O}_{T_s, t}$  est une intersection complète locale sur  $\kappa(s)$.
L'anneau local  $\mathcal{O}_{X_s, t}$  est régulier, voir Algèbre, Lemme  \ref{algebra-lemma-characterize-smooth-over-field}. Par
Algèbre, Lemme  \ref{algebra-lemma-lci-local} , nous voyons que le noyau de la surjection
 $\mathcal{O}_{X_s, t} \to \mathcal{O}_{T_s, t}$  est engendré par une suite régulière, ce que nous devions montrer.
\end{proof}

\begin{lemma}
\label{divisors-lemma-immersion-smooth-into-smooth-regular-immersion}
Soit
$$
\xymatrix{
Y \ar[rd] \ar[rr]_i & & X \ar[ld] \\
& S
}
$$
un diagramme commutatif de morphismes de schémas. Supposons que  $Y \to S$  est
lisse,  $X \to S$  lisse, et  $i$  une immersion. Alors  $i$  est
une immersion régulière.
\end{lemma}

\begin{proof}
Il s'agit d'un cas particulier du Lemme  \ref{divisors-lemma-immersion-lci-into-smooth-regular-immersion}  car un morphisme lisse est
syntomique, voir Morphismes, Lemme  \ref{morphisms-lemma-smooth-syntomic}.
\end{proof}

\begin{lemma}
\label{divisors-lemma-push-regular-immersion-thru-smooth}
Soit
$$
\xymatrix{
Y \ar[rd]_j \ar[rr]_i & & X \ar[ld] \\
& S
}
$$
un diagramme commutatif de morphismes de schémas. Supposons  $X \to S$  lisse et
 $i$  et  $j$  des immersions. Si  $i$  est une immersion de
Koszul-régulière (resp. \ $H_1$-régulière, quasi-régulière),
alors il en est de même pour  $j$.
\end{lemma}

\begin{proof}
Nous utiliserons le Lemme  \ref{divisors-lemma-regular-quasi-regular-immersion}  sans autre mention. Soit  $y \in Y$  un
point quelconque. Posons  $x = i(y)$  et posons  $s = j(y)$. Il suffit de prouver
le résultat après avoir remplacé  $X$  et  $S$  par des voisinages
ouverts  $U$  et  $V$  de  $x$  et  $s$  et
 $Y$  par un voisinage ouvert de  $y$  dans  $i^{-1}(U) \cap j^{-1}(V)$.

\medskip\noindent
Nous prouvons d'abord le résultat pour  $X = \mathbf{A}^n_S$. Après avoir remplacé
 $S$  par un ouvert affine  $V$  et remplacé  $Y$  par
 $j^{-1}(V)$ , nous pouvons supposer que  $j$  est une immersion fermée et
que  $S$  est affine. Écrivons  $S = \Spec(A)$. Alors  $j : Y \to S$  définit un
isomorphisme de  $Y$  vers le sous-schéma fermé  $\Spec(A/I)$  pour un
certain idéal  $I \subset A$. L'application  $i : Y = \Spec(A/I) \to
\mathbf{A}^n_S = \Spec(A[x_1, \ldots, x_n])$  correspond à un homomorphisme
d'$A$-algèbres  $i^\sharp : A[x_1, \ldots, x_n] \to A/I$. Choisissons  $a_i \in A$  dont l'image est
 $i^\sharp(x_i)$  dans  $A/I$. Observons que l'idéal de l'immersion fermée
 $i$  est
$$
J = (x_1 - a_1, \ldots, x_n - a_n) + IA[x_1, \ldots, x_n].
$$
Posons  $K = (x_1 - a_1, \ldots, x_n - a_n)$. Nous affirmons que la suite
$$
0 \to K/KJ \to J/J^2 \to J/(K + J^2) \to 0
$$
est exacte scindée. Pour le voir, notons que  $K/K^2$  est libre avec pour
base  $x_i - a_i$  sur l'anneau  $A[x_1, \ldots, x_n]/K \cong A$. Par conséquent,  $K/KJ$  est
libre avec la même base sur l'anneau  $A[x_1, \ldots, x_n]/J \cong A/I$. D'autre part, l'application différentielle donne une application
$$
\text{d}_{A[x_1, \ldots, x_n]/A} :
J/J^2
\longrightarrow
\Omega_{A[x_1, \ldots, x_n]/A} \otimes_{A[x_1, \ldots, x_n]}
A[x_1, \ldots, x_n]/J
$$
qui associe les générateurs  $x_i - a_i$  aux éléments de base  $\text{d}x_i$  du
module libre de droite. L'affirmation s'ensuit. De plus, notons que
 $x_1 - a_1, \ldots, x_n - a_n$  est une suite régulière dans  $A[x_1, \ldots, x_n]$  avec un anneau quotient
 $A[x_1, \ldots, x_n]/(x_1 - a_1, \ldots, x_n - a_n) \cong A$. Ainsi, nous avons une factorisation
$$
Y \to V(x_1 - a_1, \ldots, x_n - a_n) \to \mathbf{A}^n_S
$$
de notre immersion fermée  $i$  où la composition est Koszul-régulière
(resp. \ $H_1$-régulière, quasi-régulière), la deuxième flèche
est une immersion régulière, et la suite conormale associée est scindée.
Maintenant, le résultat s'ensuit du Lemme  \ref{divisors-lemma-permanence-regular-immersion}.

\medskip\noindent
Ensuite, nous prouvons que le résultat vaut si  $i$  est
 $H_1$-régulier. À savoir, en rétrécissant comme dans le premier paragraphe
de la preuve, nous pouvons supposer que  $Y$,  $X$ et  $S$ 
sont affines. Dans ce cas, nous pouvons choisir une immersion fermée  $h : X \to \mathbf{A}^n_S$ 
au-dessus de  $S$  pour un certain  $n$. Notons que  $h$  est
une immersion régulière d'après le Lemme  \ref{divisors-lemma-immersion-smooth-into-smooth-regular-immersion}. Par conséquent,
 $h \circ i$  est une immersion  $H_1$-régulière, voir le Lemme  \ref{divisors-lemma-composition-regular-immersion} 
 (notons que cette étape ne fonctionne pas dans le « cas quasi-régulier »). Ainsi,
nous sommes ramenés au cas où  $X = \mathbf{A}^n_S$  et  $S$  est affine, traité
ci-dessus.

\medskip\noindent
Enfin, supposons que  $i$  soit quasi-régulier. Après rétrécissement comme
dans le premier paragraphe de la démonstration, nous pouvons utiliser le
Lemme de Morphismes  \ref{morphisms-lemma-smooth-etale-over-affine-space}  pour factoriser  $f$  en  $X \to \mathbf{A}^n_S \to S$  où
le premier morphisme  $X \to \mathbf{A}^n_S$  est  étale. Cela réduit le problème
aux deux cas (a)  $X = \mathbf{A}^n_S$  et (b)  $f$  est  étale. Le cas
(a) a été traité dans le deuxième paragraphe de la démonstration. Le cas (b) est
traité dans le paragraphe suivant.

\medskip\noindent
Supposons que  $f$  soit  étale. Après réduction, nous pouvons
supposer  $X$,  $Y$ et  $S$  affines, et  $i$  et  $j$  des immersions
fermées (petit détail omis). Écrivons
 $S = \Spec(A)$,  $X = \Spec(B)$  et  $Y = \Spec(B/J) = \Spec(A/I)$. En rétrécissant encore, nous pouvons
supposer que  $J$  est engendré par une suite quasi-régulière.
L'application d'anneaux  $A \to B$  est  étale, donc formellement
 étale (Algèbre, Lemme  \ref{algebra-lemma-formally-etale-etale}). Ainsi  $\bigoplus I^n/I^{n + 1} \cong \bigoplus J^n/J^{n + 1}$  par Algèbre,
Lemme  \ref{algebra-lemma-formally-etale-lift-infinitesimal}. Comme  $J$  est engendré par une suite quasi-régulière,
il en est de même pour  $I$. Cela termine la démonstration.
\end{proof}











\section{Fonctions et sections méromorphes}
\label{divisors-section-meromorphic-functions}

\noindent
Cette section contient uniquement les définitions générales et quelques résultats
élémentaires. Voir  \cite{misconceptions}  pour certains pièges possibles \footnote
{Danger, Will Robinson! }.

\medskip\noindent
Soit  $(X, \mathcal{O}_X)$  un espace localement annelé. Pour tout ouvert  $U \subset X$ , nous
avons défini l'ensemble  $\mathcal{S}(U) \subset \mathcal{O}_X(U)$  des sections régulières de  $\mathcal{O}_X$  sur
 $U$, voir la Définition  \ref{divisors-definition-regular-section}. La restriction d'une section
régulière à un ouvert plus petit est régulière. Ainsi,  $\mathcal{S} : U \mapsto \mathcal{S}(U)$  est un
sous-faisceau (d'ensembles) de  $\mathcal{O}_X$. Nous notons parfois  $\mathcal{S} = \mathcal{S}_X$  si
nous voulons indiquer la dépendance à  $X$. De plus,  $\mathcal{S}(U)$  est un
sous-ensemble multiplicatif de l'anneau  $\mathcal{O}_X(U)$  pour chaque  $U$. Par
conséquent, nous pouvons considérer le préfaisceau d'anneaux
$$
U \longmapsto \mathcal{S}(U)^{-1} \mathcal{O}_X(U),
$$
voir Modules, Lemme  \ref{modules-lemma-simple-invert}.

\begin{definition}
\label{divisors-definition-sheaf-meromorphic-functions}
Soit  $(X, \mathcal{O}_X)$  un espace localement annelé. Le faisceau
 {\it  des fonctions méromorphes sur  $X$} 
est le faisceau  {\it   $\mathcal{K}_X$}  associé au
préfaisceau affiché ci-dessus. Une  {\it  fonction méromorphe
}  sur  $X$  est une section globale de  $\mathcal{K}_X$.
\end{definition}

\noindent
Comme chaque élément de chaque  $\mathcal{S}(U)$  est un non-diviseur de zéro dans
 $\mathcal{O}_X(U)$ , nous voyons que l'application naturelle de faisceaux d'anneaux
 $\mathcal{O}_X \to \mathcal{K}_X$  est injective.

\begin{example}
\label{divisors-example-no-change}
Soit  $A = \mathbf{C}[x, \{y_\alpha\}_{\alpha \in \mathbf{C}}]/
((x - \alpha)y_\alpha, y_\alpha y_\beta)$. Tout élément de  $A$  peut être écrit de manière unique
comme  $f(x) + \sum \lambda_\alpha y_\alpha$  avec  $f(x) \in \mathbf{C}[x]$  et  $\lambda_\alpha \in \mathbf{C}$. Soit  $X = \Spec(A)$. Dans ce
cas,  $\mathcal{O}_X = \mathcal{K}_X$, puisque sur tout ouvert affine  $D(f)$ , l'anneau
anneau  $A_f$, tout non-diviseur de zéro est une unité (preuve omise).
\end{example}

\noindent
Soit  $(X, \mathcal{O}_X)$  un espace localement annelé. Soit  $\mathcal{F}$  un faisceau de
 $\mathcal{O}_X$-modules. Considérons le préfaisceau  $U \mapsto \mathcal{S}(U)^{-1}\mathcal{F}(U)$. Sa faisceautisation
est le faisceau  $\mathcal{F} \otimes_{\mathcal{O}_X} \mathcal{K}_X$, voir Modules, Lemme  \ref{modules-lemma-simple-invert-module}.

\begin{definition}
\label{divisors-definition-meromorphic-section}
Soit  $X$  un espace localement annelé. Soit  $\mathcal{F}$  un faisceau de
 $\mathcal{O}_X$-modules.
\begin{enumerate}
\item Nous notons  $\mathcal{K}_X(\mathcal{F})$  le faisceau de  $\mathcal{K}_X$-modules qui est la
faisceautisation du préfaisceau  $U \mapsto \mathcal{S}(U)^{-1}\mathcal{F}(U)$. Équivalemment  $\mathcal{K}_X(\mathcal{F}) =
\mathcal{F} \otimes_{\mathcal{O}_X} \mathcal{K}_X$  (voir
ci-dessus).
\item Une section méromorphe  {\it  de  $\mathcal{F}$}  est
une section globale de  $\mathcal{K}_X(\mathcal{F})$.
\end{enumerate}
\end{definition}

\noindent
En particulier, nous avons
$$
\mathcal{K}_X(\mathcal{F})_x
=
\mathcal{F}_x \otimes_{\mathcal{O}_{X, x}} \mathcal{K}_{X, x}
=
\mathcal{S}_x^{-1}\mathcal{F}_x
$$
pour tout point  $x \in X$. Cependant, il faut faire attention car il se peut
 que  $\mathcal{S}_x$  ne soit pas l'ensemble des non-diviseurs de zéro dans l'anneau local
 $\mathcal{O}_{X, x}$. À savoir, il existe toujours une application injective
$$
\mathcal{K}_{X, x} \longrightarrow Q(\mathcal{O}_{X, x})
$$
vers l'anneau total des fractions. Elle est également surjective si et seulement si
 $\mathcal{S}_x$  est l'ensemble des non-diviseurs de zéro dans  $\mathcal{O}_{X, x}$. Les
faisceaux de sections méromorphes ne sont pas en général des modules
quasi-cohérents, mais ils possèdent certaines propriétés en commun avec les
modules quasi-cohérents.

\begin{definition}
\label{divisors-definition-pullback-meromorphic-sections}
Soit  $f : (X, \mathcal{O}_X) \to (Y, \mathcal{O}_Y)$  un morphisme d'espaces localement annelés. Nous disons que
 {\it  les images réciproques de fonctions méromorphes sont
définies pour  $f$}  si, pour chaque paire d'ouverts
 $U \subset X$,  $V \subset Y$  telle que  $f(U) \subset V$, et pour toute section
 $s \in \Gamma(V, \mathcal{S}_Y)$, l'image  $f^\sharp(s) \in \Gamma(U, \mathcal{O}_X)$  est un élément de  $\Gamma(U, \mathcal{S}_X)$.
\end{definition}

\noindent
Dans ce cas, il existe une application induite  $f^\sharp : f^{-1}\mathcal{K}_Y \to \mathcal{K}_X$, en d'autres termes,
nous obtenons un diagramme commutatif de morphismes d'espaces annelés
$$
\xymatrix{
(X, \mathcal{K}_X) \ar[r] \ar[d]^f &
(X, \mathcal{O}_X) \ar[d]^f \\
(Y, \mathcal{K}_Y) \ar[r] &
(Y, \mathcal{O}_Y)
}
$$
Nous notons parfois  $f^*(s) = f^\sharp(s)$  pour une section  $s \in \Gamma(Y, \mathcal{K}_Y)$.

\begin{lemma}
\label{divisors-lemma-pullback-meromorphic-sections-defined}
Soit  $f : X \to Y$  un morphisme de schémas. Dans chacun des cas suivants, les tirés
images réciproques des fonctions méromorphes sont définies.
\begin{enumerate}
\item chaque point faiblement associé de  $X$  s'envoie sur un point générique
d'une composante irréductible de  $Y$,
\item $X$,  $Y$  sont intègres et  $f$  est dominant,
\item $X$  est intègre et le point générique de  $X$  s'envoie sur le
point générique d'une composante irréductible de  $Y$,
\item $X$  est réduit et chaque point générique de chaque composante
irréductible de  $X$  s'envoie sur le point générique d'une composante
irréductible de  $Y$,
\item $X$  est localement noethérien et tout point associé de  $X$  se
projette sur un point générique d'une composante irréductible de  $Y$.
\item $X$  est localement noethérien, n'a pas de points immergés et tout point
générique d'une composante irréductible de  $X$  s'envoie sur le point
générique d'une composante irréductible de  $Y$, et
\item $f$  est plat.
\end{enumerate}
\end{lemma}

\begin{proof}
La question est locale sur  $X$  et  $Y$. Nous nous réduisons donc
au cas où  $X = \Spec(A)$,  $Y = \Spec(R)$  et où  $f$  est donné par un morphisme
d'anneaux  $\varphi : R \to A$. D'après la caractérisation des sections régulières du
faisceau structural dans le Lemme  \ref{divisors-lemma-regular-section-structure-sheaf} , nous devons montrer que
 $R \to A$  envoie les non-diviseurs de zéro sur des non-diviseurs de zéro. Soit
 $t \in R$  un non-diviseur de zéro.

\medskip\noindent
Si  $R \to A$  est plat, alors l'injectivité de  $t : R \to R$  montre que
 $t : A \to A$  est injectif. Cela prouve (7).

\medskip\noindent
Dans les autres cas, nous notons que  $t$  n'est contenu dans aucun des
idéaux premiers minimaux de  $R$  (car chaque élément d'un idéal premier
minimal dans un anneau est un diviseur de zéro). Ainsi, dans le cas (1), nous
voyons que  $\varphi(t)$  n'est contenu dans aucun idéal premier faiblement
associé de  $A$. Par conséquent, ce cas découle de Algèbre, Lemme
 \ref{algebra-lemma-weakly-ass-zero-divisors}. Le cas (5) est un cas particulier de (1) selon le Lemme
 \ref{divisors-lemma-ass-weakly-ass}. Le cas (6) découle de (5) et des définitions. Le cas (4) est un cas
particulier de (1) selon le Lemme  \ref{divisors-lemma-weakass-reduced}. Les cas (2) et (3) sont des cas
particuliers du (4).
\end{proof}


\begin{lemma}
\label{divisors-lemma-meromorphic-weakass-finite}
Soit  $X$  un schéma tel que
\begin{enumerate}
\item[(a)] chaque point faiblement associé de  $X$  est un point générique d'une
composante irréductible de  $X$, et
\item[(b)] tout ouvert quasi-compact a un nombre fini de composantes irréductibles.
\end{enumerate}
Soit  $X^0$  l'ensemble des points génériques des composantes irréductibles
de  $X$. Alors nous avons
$$
\mathcal{K}_X =
\bigoplus\nolimits_{\eta \in X^0} j_{\eta, *}\mathcal{O}_{X, \eta} =
\prod\nolimits_{\eta \in X^0} j_{\eta, *}\mathcal{O}_{X, \eta}
$$
où  $j_\eta : \Spec(\mathcal{O}_{X, \eta}) \to X$  est l'application canonique des Schémas, section  \ref{schemes-section-points}. De
plus,
\begin{enumerate}
\item $\mathcal{K}_X$  est un faisceau quasi-cohérent d'algèbres  $\mathcal{O}_X$,
\item pour tout  $\mathcal{O}_X$-module quasi-cohérent  $\mathcal{F}$  le faisceau
$$
\mathcal{K}_X(\mathcal{F}) =
\bigoplus\nolimits_{\eta \in X^0} j_{\eta, *}\mathcal{F}_\eta =
\prod\nolimits_{\eta \in X^0} j_{\eta, *}\mathcal{F}_\eta
$$
de sections méromorphes de  $\mathcal{F}$  est quasi-cohérent,
\item $\mathcal{S}_x \subset \mathcal{O}_{X, x}$  est l'ensemble des non-diviseurs de zéro pour tout  $x \in X$,
\item $\mathcal{K}_{X, x}$  est l'anneau total des fractions de  $\mathcal{O}_{X, x}$  pour tout
 $x \in X$,
\item $\mathcal{K}_X(U)$  est égal à l'anneau total des fractions de  $\mathcal{O}_X(U)$  pour tout
ouvert affine  $U \subset X$,
\item l'anneau des fonctions rationnelles de  $X$  (Morphismes, Définition
 \ref{morphisms-definition-rational-function}) est l'anneau des fonctions méromorphes sur  $X$, en formule:
 $R(X) = \Gamma(X, \mathcal{K}_X)$.
\end{enumerate}
\end{lemma}

\begin{proof}
Remarquons qu'une somme directe localement finie de faisceaux de modules est égale
au produit, ce que l'on peut par exemple vérifier sur les fibres. Ensuite,
puisque  $\mathcal{K}_X(\mathcal{F}) =
\mathcal{F} \otimes_{\mathcal{O}_X} \mathcal{K}_X$ , nous voyons que (2) découle des autres énoncés. De plus,
remarquez que la partie (6) découle de l'énoncé initial du lemme et du Morphismes,
Lemme  \ref{morphisms-lemma-integral-scheme-rational-functions}  lorsque  $X^0$  est fini; le cas général de (6) découle de
cela par recollement (argument omis).

\medskip\noindent
Soit  $j : Y = \coprod\nolimits_{\eta \in X^0} \Spec(\mathcal{O}_{X, \eta}) \to X$  le coproduit des morphismes  $j_\eta$. Nous devons montrer
que  $\mathcal{K}_X = j_*\mathcal{O}_Y$. Notons d'abord que  $\mathcal{K}_Y = \mathcal{O}_Y$  puisque  $Y$  est une
réunion disjointe de spectres d'anneaux locaux de dimension  $0$: dans un
anneau local de dimension zéro, tout élément qui n'est pas un diviseur de zéro est
une unité. Notons ensuite que les images réciproques de fonctions méromorphes sont
définies pour  $j$  d'après le Lemme  \ref{divisors-lemma-pullback-meromorphic-sections-defined}. Cela donne une
application
$$
\mathcal{K}_X \longrightarrow j_*\mathcal{O}_Y.
$$
Soit  $\Spec(A) = U \subset X$  un ouvert affine. Alors  $A$  est un anneau avec un
nombre fini de premiers minimaux  $\mathfrak q_1, \ldots, \mathfrak q_t$  et chaque idéal premier faiblement
associé de  $A$  est l'un des  $\mathfrak q_i$. Nous obtenons  $Q(A) = \prod A_{\mathfrak q_i}$  par
Algèbre, Lemme  \ref{algebra-lemma-total-ring-fractions-no-embedded-points}  et  \ref{algebra-lemma-weakly-ass-zero-divisors}. En d'autres termes, déjà la valeur du
préfaisceau  $U \mapsto \mathcal{S}(U)^{-1}\mathcal{O}_X(U)$  coïncide avec  $j_*\mathcal{O}_Y(U)$  sur notre ouvert affine
 $U$. Par conséquent, l'application affichée est un isomorphisme, ce qui
prouve la première égalité affichée dans l'énoncé du lemme.

\medskip\noindent
Enfin, nous prouvons (1), (3), (4) et (5). Pour la partie (5), nous avons vu au
cours de la démonstration que  $\mathcal{K}_X = j_*\mathcal{O}_Y$. Le morphisme  $j$  est
quasi-compact par notre hypothèse que l'ensemble des composantes irréductibles de
 $X$  est localement fini. Ainsi,  $j$  est quasi-compact et
quasi-séparé (car  $Y$  est séparé). D'après Schémas, Lemme  \ref{schemes-lemma-push-forward-quasi-coherent} ,
 $j_*\mathcal{O}_Y$  est quasi-cohérent. Cela prouve (1). Soit  $x \in X$. Nous pouvons
choisir un voisinage affine ouvert  $U = \Spec(A)$  de  $x$  dont toutes les
composantes irréductibles passent par  $x$. Alors  $A \subset A_\mathfrak p$  car chaque
idéal premier faiblement associé de  $A$  est contenu dans
 $\mathfrak p$ , donc les éléments de  $A \setminus \mathfrak p$  sont des non-diviseurs de zéro
d'après Algèbre, Lemme  \ref{algebra-lemma-weakly-ass-zero-divisors}. Il s'ensuit facilement que tout non-diviseur
de zéro de  $A_\mathfrak p$  est l'image d'un non-diviseur de zéro dans un voisinage
ouvert affine (possiblement plus petit) de  $x$. Cela prouve (3). La
partie (4) découle de la partie (3) par calcul des fibres.
\end{proof}

\begin{definition}
\label{divisors-definition-regular-meromorphic-section}
Soit  $X$  un espace localement annelé. Soit  $\mathcal{L}$  un
 $\mathcal{O}_X$-module inversible. Une section méromorphe  $s$  de
 $\mathcal{L}$  est dite  {\it  régulière}  si
l'application induite  $\mathcal{K}_X \to \mathcal{K}_X(\mathcal{L})$  est injective. En d'autres termes,  $s$ 
est une section régulière du  $\mathcal{K}_X$-module inversible  $\mathcal{K}_X(\mathcal{L})$, voir
Définition  \ref{divisors-definition-regular-section}.
\end{definition}

\noindent
Détaillons quand les sections méromorphes (régulières) peuvent être tirées en
arrière.

\begin{lemma}
\label{divisors-lemma-meromorphic-sections-pullback}
Soit  $f : X \to Y$  un morphisme d'espaces localement annelés. Supposons que les
images réciproques de fonctions méromorphes soient définies pour  $f$ 
(voir Définition  \ref{divisors-definition-pullback-meromorphic-sections}).
\begin{enumerate}
\item Soit  $\mathcal{F}$  un faisceau de  $\mathcal{O}_Y$-modules. Il existe une application
canonique d'image réciproque  $f^* : \Gamma(Y, \mathcal{K}_Y(\mathcal{F})) \to
\Gamma(X, \mathcal{K}_X(f^*\mathcal{F}))$  pour les sections méromorphes de
 $\mathcal{F}$.
\item Soit  $\mathcal{L}$  un  $\mathcal{O}_Y$-module inversible. L'image réciproque d'une section méromorphe
régulière  $s$  de  $\mathcal{L}$  est une section
méromorphe régulière  $f^*s$  de  $f^*\mathcal{L}$.
\end{enumerate}
\end{lemma}

\begin{proof}
Démonstration omise.
\end{proof}

\begin{lemma}
\label{divisors-lemma-regular-meromorphic-ideal-denominators}
Soit  $X$  un schéma. Soit  $\mathcal{L}$  un  $\mathcal{O}_X$-module inversible.
Soit  $s$  une section méromorphe régulière de  $\mathcal{L}$. Notons
 $\mathcal{I} \subset \mathcal{O}_X$  le faisceau d'idéaux défini par la règle
$$
\mathcal{I}(V)
=
\{f \in \mathcal{O}_X(V) \mid fs \in \mathcal{L}(V)\}.
$$
La formule a du sens puisque  $\mathcal{L}(V) \subset \mathcal{K}_X(\mathcal{L})(V)$. Ensuite,  $\mathcal{I}$  est un faisceau
quasi-cohérent d'idéaux et nous avons des applications injectives
$$
1 : \mathcal{I} \longrightarrow \mathcal{O}_X,
\quad
s : \mathcal{I} \longrightarrow \mathcal{L}
$$
dont les conoyaux ont leur support dans des sous-ensembles fermés nulle part denses
de  $X$.
\end{lemma}

\begin{proof}
La question est locale sur  $X$. Par conséquent, nous pouvons supposer que
 $X = \Spec(A)$ et  $\mathcal{L} = \mathcal{O}_X$. Après un rétrécissement supplémentaire, nous pouvons
supposer que  $s = a/b$  avec  $a, b \in A$   {\it  tous deux
}  non-diviseurs de zéro dans  $A$. Posons  $I = \{x \in A \mid x(a/b) \in A\}$.

\medskip\noindent
Pour montrer que  $\mathcal{I}$  est quasi-cohérent, nous devons montrer que
 $I_f = \{x \in A_f \mid x(a/b) \in A_f\}$  pour chaque  $f \in A$. Si  $c/f^n \in A_f$,  $(c/f^n)(a/b) \in A_f$, alors nous
voyons que  $f^mc(a/b) \in A$  pour certains  $m$, donc  $c/f^n \in I_f$.
Inversement, il est facile de voir que  $I_f$  est contenu dans  $\{x \in A_f \mid x(a/b) \in A_f\}$.
Cela prouve la quasi-cohérence.

\medskip\noindent
Prouvons l'assertion finale. Il est clair que  $(b) \subset I$. Par conséquent,
 $V(I) \subset V(b)$  est un sous-ensemble nulle part dense puisque  $b$  est un
non-diviseur de zéro. Ainsi, le conoyau de  $1$  a son support dans un
ensemble fermé nulle part dense. Le même argument fonctionne pour le conoyau de
 $s$  puisque  $s(b) = (a) \subset sI \subset A$.
\end{proof}

\begin{definition}
\label{divisors-definition-regular-meromorphic-ideal-denominators}
Soit  $X$  un schéma. Soit  $\mathcal{L}$  un  $\mathcal{O}_X$-module inversible.
Soit  $s$  une section méromorphe régulière de  $\mathcal{L}$. Le faisceau
d'idéaux  $\mathcal{I}$  construit dans le Lemme  \ref{divisors-lemma-regular-meromorphic-ideal-denominators}  est appelé le faisceau
d'idéaux de dénominateurs  {\it  de  $s$}.
\end{definition}




\section{Fonctions et sections méromorphes ; cas noethérien}
\label{divisors-section-meromorphic-noetherian}

\noindent
Pour les schémas localement noethériens, nous pouvons prouver certains résultats
concernant le faisceau de fonctions méromorphes. Cependant, il existe un exemple
dans  \cite{misconceptions}  montrant que  $\mathcal{K}_X$  n'est pas nécessairement quasi-cohérent
pour un schéma noethérien  $X$.

\begin{lemma}
\label{divisors-lemma-meromorphic-section-restricts-to-zero}
Soit  $X$  un schéma quasi-compact. Soient  $h \in \Gamma(X, \mathcal{O}_X)$  et  $f \in \Gamma(X, \mathcal{K}_X)$ 
tels que  $f$  se restreigne à zéro sur  $X_h$. Alors  $h^n f = 0$ 
pour un certain  $n \gg 0$.
\end{lemma}

\begin{proof}
Nous pouvons trouver un recouvrement de  $X$  par des ouverts affines
 $U$  tels que  $f|_U = s^{-1}a$  avec  $a \in \mathcal{O}_X(U)$  et  $s \in \mathcal{S}(U)$. Comme
 $X$  est quasi-compact, nous pouvons le recouvrir par un nombre fini
d'ouverts affines de cette forme. Ainsi, il suffit de prouver le lemme lorsque
 $X = \Spec(A)$  et  $f = s^{-1}a$. Notons que  $s \in A$  est un non-diviseur de zéro,
il suffit donc de prouver le résultat lorsque  $f = a$. La condition
 $f|_{X_h} = 0$  implique que  $a$  s'envoie à zéro dans  $A_h = \mathcal{O}_X(X_h)$  comme
 $\mathcal{O}_X \subset \mathcal{K}_X$. Ainsi  $h^na = 0$  pour un certain  $n > 0$  comme désiré.
\end{proof}

\begin{lemma}
\label{divisors-lemma-locally-Noetherian-K}
Soit  $X$  un schéma localement noethérien.
\begin{enumerate}
 \item Pour tout  $x \in X$,  $\mathcal{S}_x \subset \mathcal{O}_{X, x}$  est l'ensemble des
 non-diviseurs de zéro, et donc  $\mathcal{K}_{X, x}$  est l'anneau total des fractions de  $\mathcal{O}_{X, x}$.
 \item Pour tout ouvert affine  $U \subset X$ , l'anneau  $\mathcal{K}_X(U)$  est égal à l'anneau total
 des fractions de  $\mathcal{O}_X(U)$.
\end{enumerate}
\end{lemma}

\begin{proof}
Pour démontrer ce lemme, nous pouvons supposer que  $X$  est le spectre
d'un anneau noethérien  $A$. Disons que  $x \in X$  correspond à
 $\mathfrak p \subset A$.

\medskip\noindent
Preuve de (1). Il est clair que  $\mathcal{S}_x$  est contenu dans l'ensemble des
non-diviseurs de zéro de  $\mathcal{O}_{X, x} = A_\mathfrak p$. Pour la réciproque, soit  $f, g \in A$,
 $g \not \in \mathfrak p$  et supposons que  $f/g$  soit un non-diviseur de zéro dans
 $A_{\mathfrak p}$. Soit  $I = \{a \in A \mid af = 0\}$. Alors nous voyons que  $I_{\mathfrak p} = 0$  par
l'exactitude de la localisation. Puisque  $A$  est noethérien, nous voyons
que  $I$  est de type fini et donc que  $g'I = 0$  pour certains
 $g' \in A$,  $g' \not \in \mathfrak p$. Ainsi  $f$  est un non-diviseur de zéro dans
 $A_{g'}$, c'est-à-dire dans un voisinage ouvert de Zariski de  $\mathfrak p$.
Ainsi  $f/g$  est un élément de  $\mathcal{S}_x$.

\medskip\noindent
Preuve de (2). Soit  $f \in \Gamma(X, \mathcal{K}_X)$  une fonction méromorphe. Posons  $I = \{a \in A \mid af \in A\}$.
Fixons un idéal premier  $\mathfrak p \subset A$  correspondant au point  $x \in X$. D'après (1),
nous pouvons écrire l'image de  $f$  dans la fibre en  $\mathfrak p$  comme
 $a/b$,  $a, b \in A_{\mathfrak p}$  avec  $b \in A_{\mathfrak p}$  n'étant pas un diviseur de zéro.
Écrivons  $b = c/d$  avec  $c, d \in A$,  $d \not \in \mathfrak p$. Alors  $ad - cf$  est une
section de  $\mathcal{K}_X$  qui s'annule dans un voisinage ouvert de  $x$.
Disons qu'elle s'annule sur  $D(e)$  avec  $e \in A$,  $e \not \in \mathfrak p$. Alors
 $e^n(ad - cf) = 0$  pour un certain  $n \gg 0$  d'après le Lemme  \ref{divisors-lemma-meromorphic-section-restricts-to-zero}. Ainsi,
 $e^nc \in I$  et  $e^nc$  s'envoie sur un non-diviseur de zéro dans
 $A_{\mathfrak p}$. Écrivons  $\text{Ass}(A) = \{\mathfrak q_1, \ldots, \mathfrak q_t\}$  pour les idéaux premiers associés à  $A$. En
regardant  $IA_{\mathfrak q_i}$  et en utilisant Algèbre, Lemme  \ref{algebra-lemma-associated-primes-localize} , ce qui précède
dit que  $I \not \subset \mathfrak q_i$  pour chaque  $i$. Par Algèbre, Lemme  \ref{algebra-lemma-silly} , il
existe un élément  $z \in I$,  $z \not \in \bigcup \mathfrak q_i$. Par Algèbre, Lemme  \ref{algebra-lemma-ass-zero-divisors} ,
nous voyons que  $z$  n'est pas un diviseur de zéro dans  $A$. Par
conséquent,  $f = (zf)/z$  est un élément de l'anneau total des fractions de
 $A$. Cela prouve (2).
\end{proof}

\begin{lemma}
\label{divisors-lemma-quasi-coherent-K}
Soit  $X$  un schéma localement noethérien n'ayant pas de points immergés.
Soit  $X^0$  l'ensemble des points génériques des composantes irréductibles
de  $X$. Alors nous avons
$$
\mathcal{K}_X =
\bigoplus\nolimits_{\eta \in X^0} j_{\eta, *}\mathcal{O}_{X, \eta} =
\prod\nolimits_{\eta \in X^0} j_{\eta, *}\mathcal{O}_{X, \eta}
$$
où  $j_\eta : \Spec(\mathcal{O}_{X, \eta}) \to X$  est l'application canonique des Schémas, section  \ref{schemes-section-points}. De
plus,
\begin{enumerate}
\item $\mathcal{K}_X$  est un faisceau quasi-cohérent d'algèbres  $\mathcal{O}_X$,
\item pour tout  $\mathcal{O}_X$-module quasi-cohérent  $\mathcal{F}$  le faisceau
$$
\mathcal{K}_X(\mathcal{F}) =
\bigoplus\nolimits_{\eta \in X^0} j_{\eta, *}\mathcal{F}_\eta =
\prod\nolimits_{\eta \in X^0} j_{\eta, *}\mathcal{F}_\eta
$$
des sections méromorphes de  $\mathcal{F}$  est quasi-cohérente, et
\item l'anneau des fonctions rationnelles de  $X$  est l'anneau des fonctions
méromorphes sur  $X$, sous une formule:  $R(X) = \Gamma(X, \mathcal{K}_X)$.
\end{enumerate}
\end{lemma}

\begin{proof}
Ce lemme est un cas particulier du Lemme  \ref{divisors-lemma-meromorphic-weakass-finite}  car, dans le cas localement
noethérien, les points faiblement associés sont la même chose que les points
associés par le Lemme  \ref{divisors-lemma-ass-weakly-ass}.
\end{proof}

\begin{lemma}
\label{divisors-lemma-regular-meromorphic-section-exists-noetherian}
Soit  $X$  un schéma localement noethérien n'ayant pas de points immergés.
Soit  $\mathcal{L}$  un  $\mathcal{O}_X$-module inversible. Alors  $\mathcal{L}$  possède
une section méromorphe régulière.
\end{lemma}

\begin{proof}
Pour chaque point générique  $\eta$  de  $X$ , choisissons un générateur
 $s_\eta$  du module libre de rang  $1$   $\mathcal{L}_\eta$  sur l'anneau local
artinien  $\mathcal{O}_{X, \eta}$. Il découle immédiatement de la description de  $\mathcal{K}_X$ 
et  $\mathcal{K}_X(\mathcal{L})$  dans le Lemme  \ref{divisors-lemma-quasi-coherent-K}  que  $s = \prod s_\eta$  est une section
méromorphe régulière de  $\mathcal{L}$.
\end{proof}

\begin{lemma}
\label{divisors-lemma-make-maps-regular-section}
Supposons données
\begin{enumerate}
\item $X$  un schéma localement noethérien,
\item $\mathcal{L}$  un  $\mathcal{O}_X$-module inversible,
\item $s$  une section méromorphe régulière de  $\mathcal{L}$, et
\item $\mathcal{F}$  cohérent sur  $X$  sans points associés immergés et
 $\text{Supp}(\mathcal{F}) = X$.
\end{enumerate}
Soit  $\mathcal{I} \subset \mathcal{O}_X$  l'idéal des dénominateurs de  $s$. Soit  $T \subset X$ 
l'union des supports de  $\mathcal{O}_X/\mathcal{I}$  et  $\mathcal{L}/s(\mathcal{I})$ , qui est un sous-ensemble
fermé et nulle part dense  $T \subset X$  selon le Lemme  \ref{divisors-lemma-regular-meromorphic-ideal-denominators}. Alors, il existe des
applications injectives canoniques
$$
1 : \mathcal{I}\mathcal{F} \to \mathcal{F}, \quad
s : \mathcal{I}\mathcal{F} \to \mathcal{F} \otimes_{\mathcal{O}_X}\mathcal{L}
$$
dont les conoyaux ont leur support dans  $T$.
\end{lemma}

\begin{proof}
Réduisons-nous au cas affine avec  $\mathcal{L} \cong \mathcal{O}_X$, et  $s = a/b$  avec  $a, b \in A$  tous
deux non diviseurs de zéro. La preuve de l'étape de réduction est omise. Écrivons
 $\mathcal{F} = \widetilde{M}$. Soit  $I = \{x \in A \mid x(a/b) \in A\}$  de sorte que  $\mathcal{I} = \widetilde{I}$  (voir la preuve du Lemme
 \ref{divisors-lemma-regular-meromorphic-ideal-denominators}). Notons que  $T = V(I) \cup V((a/b)I)$. Pour tout  $A$-module
 $M$ , considérons l'application  $1 : IM \to M$; c'est l'application qui donne
naissance à l'application  $1$  du lemme. Considérons d'autre part l'application
 $\sigma : IM \to M_b, x \mapsto ax/b$. Puisque  $b$  n'est pas un diviseur de zéro dans
 $A$, et puisque  $M$  a pour support  $\Spec(A)$  et aucun idéal
premier associé, nous voyons que  $b$  est également un non-diviseur de
zéro sur  $M$ . Par conséquent  $M \subset M_b$. Par définition de
 $I$ , nous avons  $\sigma(IM) \subset M$  comme sous-modules de  $M_b$. Ainsi,
nous obtenons une application de  $A$-modules  $s : IM \to M$  (à savoir
l'application unique telle que  $s(z)/1 = \sigma(z)$  dans  $M_b$  pour tout
 $z \in IM$). Elle est injective parce que  $a$  est aussi un non-diviseur
de zéro (à la fois sur  $A$  et  $M$). Il est clair que
 $M/IM$  est annihilé par  $I$  et que  $M/s(IM)$  est annihilé par
 $(a/b)I$. Ainsi, le lemme s'ensuit.
\end{proof}





\section{Fonctions et sections méromorphes ; cas réduit}
\label{divisors-section-meromorphic-reduced}

\noindent
Pour un schéma qui est réduit et qui a localement un nombre fini de composantes
irréductibles, le faisceau des fonctions méromorphes est quasi-cohérent.

\begin{lemma}
\label{divisors-lemma-reduced-finite-irreducible}
Soit  $X$  un schéma réduit tel que tout ouvert quasi-compact possède un
nombre fini de composantes irréductibles. Soit  $X^0$  l'ensemble des points
génériques des composantes irréductibles de  $X$. Alors nous avons
$$
\mathcal{K}_X =
\bigoplus\nolimits_{\eta \in X^0} j_{\eta, *}\kappa(\eta) =
\prod\nolimits_{\eta \in X^0} j_{\eta, *}\kappa(\eta)
$$
où  $j_\eta : \Spec(\kappa(\eta)) \to X$  est l'application canonique des Schémas, section  \ref{schemes-section-points}. De
plus,
\begin{enumerate}
\item $\mathcal{K}_X$  est un faisceau quasi-cohérent d'algèbres  $\mathcal{O}_X$,
\item pour tout  $\mathcal{O}_X$-module quasi-cohérent  $\mathcal{F}$  le faisceau
$$
\mathcal{K}_X(\mathcal{F}) =
\bigoplus\nolimits_{\eta \in X^0} j_{\eta, *}\mathcal{F}_\eta =
\prod\nolimits_{\eta \in X^0} j_{\eta, *}\mathcal{F}_\eta
$$
de sections méromorphes de  $\mathcal{F}$  est quasi-cohérent,
\item $\mathcal{S}_x \subset \mathcal{O}_{X, x}$  est l'ensemble des non-diviseurs de zéro pour tout  $x \in X$,
\item $\mathcal{K}_{X, x}$  est l'anneau total des fractions de  $\mathcal{O}_{X, x}$  pour tout
 $x \in X$,
\item $\mathcal{K}_X(U)$  est égal à l'anneau total des fractions de  $\mathcal{O}_X(U)$  pour tout
ouvert affine  $U \subset X$,
\item l'anneau des fonctions rationnelles de  $X$  est l'anneau des fonctions
méromorphes sur  $X$, sous une formule:  $R(X) = \Gamma(X, \mathcal{K}_X)$.
\end{enumerate}
\end{lemma}

\begin{proof}
Ce lemme est un cas particulier du Lemme  \ref{divisors-lemma-meromorphic-weakass-finite}  car sur un schéma réduit,
les points faiblement associés sont les points génériques selon le Lemme
 \ref{divisors-lemma-weakass-reduced}.
\end{proof}

\begin{lemma}
\label{divisors-lemma-reduced-normalization}
Soit  $X$  un schéma. Supposons que  $X$  est réduit et que tout
ouvert quasi-compact  $U \subset X$  ait un nombre fini de composantes irréductibles.
Alors le morphisme de normalisation  $\nu : X^\nu \to X$  est le morphisme
$$
\underline{\Spec}_X(\mathcal{O}') \longrightarrow X
$$
où  $\mathcal{O}' \subset \mathcal{K}_X$  est la fermeture intégrale de  $\mathcal{O}_X$  dans le faisceau des
fonctions méromorphes.
\end{lemma}

\begin{proof}
Comparez la définition du morphisme de normalisation  $\nu : X^\nu \to X$  (voir
Morphismes, Définition  \ref{morphisms-definition-normalization}) avec la description de  $\mathcal{K}_X$  dans le
Lemme  \ref{divisors-lemma-reduced-finite-irreducible}  ci-dessus.
\end{proof}

\begin{lemma}
\label{divisors-lemma-meromorphic-functions-integral-scheme}
Soit  $X$  un schéma intègre avec point générique  $\eta$. Nous avons
\begin{enumerate}
\item le faisceau des fonctions méromorphes est isomorphe au faisceau constant de valeur
le corps des fonctions (voir Morphismes, Définition  \ref{morphisms-definition-function-field}) de  $X$.
\item Pour tout faisceau quasi-cohérent  $\mathcal{F}$  sur  $X$ , le faisceau
 $\mathcal{K}_X(\mathcal{F})$  est isomorphe au faisceau constant de valeur  $\mathcal{F}_\eta$.
\end{enumerate}
\end{lemma}

\begin{proof}
Démonstration omise.
\end{proof}

\noindent
Dans certains cas, nous pouvons montrer que des sections méromorphes régulières
existent.

\begin{lemma}
\label{divisors-lemma-regular-meromorphic-section-exists}
Soit  $X$  un schéma. Soit  $\mathcal{L}$  un  $\mathcal{O}_X$-module inversible.
Dans chacun des cas suivants,  $\mathcal{L}$  possède une section méromorphe
régulière:
\begin{enumerate}
\item $X$  est intègre,
\item $X$  est réduit et tout ouvert quasi-compact a un nombre fini de
composantes irréductibles,
\item $X$  est localement noethérien et n'a pas de points immergés.
\end{enumerate}
\end{lemma}

\begin{proof}
Dans le cas (1), soit  $\eta \in X$  le point générique. Nous avons vu dans le Lemme
 \ref{divisors-lemma-meromorphic-functions-integral-scheme}  que  $\mathcal{K}_X$, respectivement \ $\mathcal{K}_X(\mathcal{L})$ , est le
faisceau constant de valeur  $\kappa(\eta)$, respectivement \ $\mathcal{L}_\eta$.
Puisque  $\dim_{\kappa(\eta)} \mathcal{L}_\eta = 1$ , nous pouvons choisir un élément non nul  $s \in \mathcal{L}_\eta$. Il est
clair que  $s$  est une section méromorphe régulière de  $\mathcal{L}$. Dans
le cas (2), choisissons  $s_\eta \in \mathcal{L}_\eta$  non nul pour tous les points génériques
 $\eta$  de  $X$; cela est possible car  $\mathcal{L}_\eta$  est un espace
vectoriel de dimension  $1$ sur  $\kappa(\eta)$. Il découle immédiatement de
la description de  $\mathcal{K}_X$  et  $\mathcal{K}_X(\mathcal{L})$  dans le Lemme  \ref{divisors-lemma-reduced-finite-irreducible}  que
 $s = \prod s_\eta$  est une section méromorphe régulière de  $\mathcal{L}$. Le cas (3) est
le Lemme  \ref{divisors-lemma-regular-meromorphic-section-exists-noetherian}.
\end{proof}






\section{Diviseurs de Weil}
\label{divisors-section-Weil-divisors}

\noindent
Nous allons introduire les diviseurs de Weil et l'équivalence rationnelle des
diviseurs de Weil pour les schémas intègres localement noethériens. Comme nous ne
supposons pas que nos schémas sont quasi-compacts, nous devons être un peu
prudents lors de la définition des diviseurs de Weil. Nous devons permettre des
sommes infinies de diviseurs premiers car une fonction rationnelle peut avoir, par
exemple, un nombre infini de pôles. Pour les schémas quasi-compacts, nos diviseurs
de Weil sont des sommes finies comme d'habitude. Voici un lemme de base que nous
utiliserons souvent pour prouver que des collections de sous-schémas fermés sont
localement finies.

\begin{lemma}
\label{divisors-lemma-components-locally-finite}
Soit  $X$  un schéma localement noethérien. Soit  $Z \subset X$  un
sous-schéma fermé. La collection des composantes irréductibles de  $Z$  est
localement finie dans  $X$.
\end{lemma}

\begin{proof}
Soit  $U \subset X$  un sous-schéma ouvert quasi-compact. Alors  $U$  est un
schéma noethérien, et a donc un espace topologique sous-jacent noethérien
(Propriétés, Lemme  \ref{properties-lemma-Noetherian-topology}). Par conséquent, chaque sous-espace est
noethérien et possède un nombre fini de composantes irréductibles (voir Topologie,
Lemme  \ref{topology-lemma-Noetherian}).
\end{proof}

\noindent
Rappelons que si  $Z$  est un sous-ensemble fermé irréductible d'un schéma
 $X$, alors la codimension de  $Z$  dans  $X$  est égale à
la dimension de l'anneau local  $\mathcal{O}_{X, \xi}$, où  $\xi \in Z$  est le point
générique. Voir Propriétés, Lemme  \ref{properties-lemma-codimension-local-ring}.

\begin{definition}
\label{divisors-definition-Weil-divisor}
Soit  $X$  un schéma intègre localement noethérien.
\begin{enumerate}
\item Un  {\it  diviseur premier}  est un sous-schéma fermé
intègre  $Z \subset X$  de codimension  $1$.
\item Un {\it diviseur de Weil}  est une somme formelle
 $D = \sum n_Z Z$ où la somme porte sur les diviseurs premiers de $X$  et la
collection  $\{Z \mid n_Z \not = 0\}$ est localement finie (Topologie, Définition \ref{topology-definition-locally-finite}).
\end{enumerate}
Le groupe de tous les diviseurs de Weil sur  $X$  est noté  $\text{Div}(X)$.
\end{definition}

\noindent
Notre prochaine tâche est de définir le diviseur de Weil associé à une fonction
rationnelle. Pour ce faire, nous utilisons l'ordre d'annulation d'une fonction
rationnelle le long d'un diviseur premier, qui est défini comme suit.

\begin{definition}
\label{divisors-definition-order-vanishing}
Soit $X$ un schéma intègre localement noethérien. Soit $f \in R(X)^*$.
Pour chaque diviseur premier $Z \subset X$ nous définissons
l'{\it ordre d'annulation de $f$ le
long de $Z$} comme l'entier
$$
\text{ord}_Z(f) = \text{ord}_{\mathcal{O}_{X, \xi}}(f)
$$
où le membre de droite est la notion d'Algèbre, Définition  \ref{algebra-definition-ord}, et  $\xi$
est le point générique de  $Z$.
\end{definition}

\noindent
Notons que pour  $f, g \in R(X)^*$  nous avons
$$
\text{ord}_Z(fg) = \text{ord}_Z(f) + \text{ord}_Z(g).
$$
Bien sûr, il peut arriver que  $\text{ord}_Z(f) < 0$. Dans ce cas, nous disons que
 $f$  a un  {\it  pôle}  le long de
 $Z$  et que  $-\text{ord}_Z(f) > 0$  est le  {\it  ordre du pôle de
 $f$  le long de  $Z$}. Il est important de noter que la
condition  $\text{ord}_Z(f) \geq 0$  est  {\bf  pas}  équivalente à
la condition  $f \in \mathcal{O}_{X, \xi}$  sauf si l'anneau local  $\mathcal{O}_{X, \xi}$  est un anneau de
valuation discrète.

\begin{lemma}
\label{divisors-lemma-divisor-locally-finite}
Soit  $X$  un schéma intègre localement noethérien. Soit  $f \in R(X)^*$.
Alors les collections
$$
\{Z \subset X \mid Z\text{ est un diviseur premier de point générique }\xi
\text{ et }f\text{ n'appartient pas à }\mathcal{O}_{X, \xi}\}
$$
et
$$
\{Z \subset X \mid Z \text{ est un diviseur premier et }\text{ord}_Z(f) \not = 0\}
$$
sont localement finis dans  $X$.
\end{lemma}

\begin{proof}
Il existe un sous-schéma ouvert non vide  $U \subset X$  tel que  $f$ 
corresponde à une section de  $\Gamma(U, \mathcal{O}_X^*)$. Les diviseurs premiers qui
peuvent apparaître dans les ensembles du lemme sont tous des composantes
irréductibles de  $X \setminus U$. Le Lemme  \ref{divisors-lemma-components-locally-finite}  donne donc le résultat
souhaité.
\end{proof}

\noindent
Ce lemme nous permet de faire la définition suivante.

\begin{definition}
\label{divisors-definition-principal-divisor}
Soit  $X$  un schéma intègre localement noethérien. Soit  $f \in R(X)^*$. Le
diviseur de Weil principal  {\it  associé à
 $f$}  est le diviseur de Weil
$$
\text{div}(f) = \text{div}_X(f) = \sum \text{ord}_Z(f) [Z]
$$
où la somme porte sur les diviseurs premiers et  $\text{ord}_Z(f)$  est comme dans la
Définition  \ref{divisors-definition-order-vanishing}. Cette somme est bien définie grâce au Lemme  \ref{divisors-lemma-divisor-locally-finite}.
\end{definition}

\begin{lemma}
\label{divisors-lemma-div-additive}
Soit  $X$  un schéma intègre localement noethérien. Soit  $f, g \in R(X)^*$.
Alors
$$
\text{div}_X(fg) = \text{div}_X(f) + \text{div}_X(g)
$$
comme diviseurs de Weil sur  $X$.
\end{lemma}

\begin{proof}
Cela est clair à partir de l'additivité des fonctions  $\text{ord}$ .
\end{proof}

\noindent
Nous voyons d'après le lemme ci-dessus que l'ensemble des diviseurs de Weil
principaux forme un sous-groupe du groupe de tous les diviseurs de Weil. Cela
conduit à la définition suivante.

\begin{definition}
\label{divisors-definition-class-group}
Soit  $X$  un schéma intègre localement noethérien. Le
 {\it  groupe de classes de diviseurs de Weil }  de
 $X$  est le quotient du groupe des diviseurs de Weil par le sous-groupe
des diviseurs de Weil principaux. On le note  $\text{Cl}(X)$.
\end{definition}

\noindent
Par construction, nous obtenons un complexe exact
\begin{equation}
\label{divisors-equation-Weil-divisor-class}
R(X)^* \xrightarrow{\text{div}} \text{Div}(X) \to \text{Cl}(X) \to 0
\end{equation}
que nous pouvons considérer comme une présentation de  $\text{Cl}(X)$. Notre
prochaine tâche est de relier le groupe de classes de diviseurs de Weil au groupe
de Picard.





\section{La classe de diviseur de Weil associée à un module inversible}
\sectionmark{Classe de diviseur d'un module inversible}
\label{divisors-section-c1}

\noindent
Dans cette section, nous suivons exactement la même progression que dans la
section  \ref{divisors-section-Weil-divisors}  pour définir une application canonique  $\Pic(X) \to \text{Cl}(X)$  sur un
schéma intègre localement noethérien.

\medskip\noindent
Soit  $X$  un schéma. Soit  $\mathcal{L}$  un  $\mathcal{O}_X$-module inversible.
Soit  $\xi \in X$  un point. Si  $s_\xi, s'_\xi \in \mathcal{L}_\xi$  engendrent  $\mathcal{L}_\xi$  comme
 $\mathcal{O}_{X, \xi}$-module, alors il existe une unité  $u \in \mathcal{O}_{X, \xi}^*$  telle que
 $s_\xi = u s'_\xi$. La fibre du faisceau des sections méromorphes  $\mathcal{K}_X(\mathcal{L})$  de
 $\mathcal{L}$  en  $x$  est  $\mathcal{K}_{X, x} \otimes_{\mathcal{O}_{X, x}} \mathcal{L}_x$. Ainsi, l'image de toute
section méromorphe  $s$  de  $\mathcal{L}$  dans la fibre en  $x$  peut
s'écrire comme  $s = fs_\xi$  avec  $f \in \mathcal{K}_{X, x}$. Ci-dessous, nous abrégerons cela en
disant  $f = s/s_\xi$. De plus, si  $X$  est intègre, nous avons que
 $\mathcal{K}_{X, x} = R(X)$, le corps des fonctions de  $X$, donc  $s/s_\xi \in R(X)$.
Si  $s$  est une section méromorphe régulière, alors en fait  $s/s_\xi \in R(X)^*$.
Sur un schéma intègre, une section méromorphe régulière est la même chose qu'une
section méromorphe non nulle. Enfin, nous voyons que  $s/s_\xi$  est indépendant
du choix de  $s_\xi$  à multiplication près par une unité de
l'anneau local  $\mathcal{O}_{X, x}$. Tout ceci montre que la définition
suivante est bien définie.

\begin{definition}
\label{divisors-definition-order-vanishing-meromorphic}
Soit  $X$  un schéma intègre localement noethérien. Soit  $\mathcal{L}$  un
 $\mathcal{O}_X$-module inversible. Soit  $s \in \Gamma(X, \mathcal{K}_X(\mathcal{L}))$  une section méromorphe régulière
de  $\mathcal{L}$. Pour chaque diviseur premier  $Z \subset X$ , nous définissons
l'{\it ordre d'annulation de  $s$  le long de
 $Z$}  comme l'entier
$$
\text{ord}_{Z, \mathcal{L}}(s)
= \text{ord}_{\mathcal{O}_{X, \xi}}(s/s_\xi)
$$
où le membre de droite est la notion d'Algèbre, Définition  \ref{algebra-definition-ord},  $\xi \in Z$
est le point générique, et  $s_\xi \in \mathcal{L}_\xi$  est un générateur.
\end{definition}

\noindent
Comme dans le cas des diviseurs principaux, nous avons le lemme suivant.

\begin{lemma}
\label{divisors-lemma-divisor-meromorphic-locally-finite}
Soit  $X$  un schéma intègre localement noethérien. Soit  $\mathcal{L}$  un
 $\mathcal{O}_X$-module inversible. Soit  $s \in \mathcal{K}_X(\mathcal{L})$  une section méromorphe régulière
(c'est-à-dire non nulle) de  $\mathcal{L}$. Alors les ensembles
$$
\{Z \subset X \mid Z \text{ est un diviseur premier de point générique }\xi
\text{ et }s\text{ n'appartient pas à }\mathcal{L}_\xi\}
$$
et
$$
\{Z \subset X \mid Z \text{ est un diviseur premier et }
\text{ord}_{Z, \mathcal{L}}(s) \not = 0\}
$$
sont localement finis dans  $X$.
\end{lemma}

\begin{proof}
Il existe un sous-schéma ouvert non vide  $U \subset X$  tel que  $s$ 
corresponde à une section de  $\Gamma(U, \mathcal{L})$  qui génère  $\mathcal{L}$  sur
 $U$. Ainsi, les diviseurs premiers qui peuvent apparaître dans les
ensembles du lemme sont tous des composantes irréductibles de  $X \setminus U$. Ainsi,
le Lemme  \ref{divisors-lemma-components-locally-finite} donne le résultat souhaité.
\end{proof}

\begin{lemma}
\label{divisors-lemma-divisor-meromorphic-well-defined}
Soit  $X$  un schéma intègre localement noethérien. Soit  $\mathcal{L}$  un
 $\mathcal{O}_X$-module inversible. Soient  $s, s' \in \mathcal{K}_X(\mathcal{L})$  deux sections méromorphes non nulles
de  $\mathcal{L}$. Alors  $f = s/s'$  est un élément de  $R(X)^*$  et
$$
\sum \text{ord}_{Z, \mathcal{L}}(s)[Z]
=
\sum \text{ord}_{Z, \mathcal{L}}(s')[Z]
+
\text{div}(f)
$$
comme diviseurs de Weil.
\end{lemma}

\begin{proof}
Cela est clair d'après les définitions. Il faut noter que le Lemme  \ref{divisors-lemma-divisor-meromorphic-locally-finite} 
garantit que les sommes sont bien des diviseurs de Weil.
\end{proof}

\begin{definition}
\label{divisors-definition-divisor-invertible-sheaf}
Soit  $X$  un schéma intègre localement noethérien. Soit  $\mathcal{L}$  un
 $\mathcal{O}_X$-module inversible.
\begin{enumerate}
\item Pour toute section méromorphe non nulle  $s$  de  $\mathcal{L}$ , nous
définissons le diviseur de Weil  {\it  associé à
 $s$}  comme
$$
\text{div}_\mathcal{L}(s) =
\sum \text{ord}_{Z, \mathcal{L}}(s) [Z] \in \text{Div}(X)
$$
où la somme est effectuée sur les diviseurs premiers.
\item Nous définissons la  {\it classe de diviseur de Weil associée
à $\mathcal{L}$} comme l'image
de $\text{div}_\mathcal{L}(s)$ dans $\text{Cl}(X)$ où $s$ est toute section méromorphe non nulle
de $\mathcal{L}$ sur $X$. Cela est bien défini par le
Lemme \ref{divisors-lemma-divisor-meromorphic-well-defined}.
\end{enumerate}
\end{definition}

\noindent
Comme prévu, cette construction est additive dans le module inversible.

\begin{lemma}
\label{divisors-lemma-c1-additive}
Soit  $X$  un schéma intègre localement noethérien. Soient  $\mathcal{L}$,
 $\mathcal{N}$  des  $\mathcal{O}_X$-modules inversibles. Soient  $s$, resp.
\ $t$  une section méromorphe non nulle de  $\mathcal{L}$, resp.
\ $\mathcal{N}$. Alors  $st$  est une section méromorphe non nulle
de  $\mathcal{L} \otimes \mathcal{N}$, et
$$
\text{div}_{\mathcal{L} \otimes \mathcal{N}}(st)
=
\text{div}_\mathcal{L}(s) + \text{div}_\mathcal{N}(t)
$$
dans  $\text{Div}(X)$. En particulier, la classe de diviseur de Weil de  $\mathcal{L} \otimes_{\mathcal{O}_X} \mathcal{N}$ 
est la somme des classes de diviseurs de Weil de  $\mathcal{L}$  et  $\mathcal{N}$.
\end{lemma}

\begin{proof}
Soit  $s$, resp. \ $t$ , une section méromorphe non
nulle de  $\mathcal{L}$, resp. \ $\mathcal{N}$. Alors  $st$  est une
section méromorphe non nulle de  $\mathcal{L} \otimes \mathcal{N}$. Soit  $Z \subset X$  un diviseur
premier. Soit  $\xi \in Z$  son point générique. Choisissons les générateurs
 $s_\xi \in \mathcal{L}_\xi$ et  $t_\xi \in \mathcal{N}_\xi$. Alors  $s_\xi t_\xi$  est un générateur de
 $(\mathcal{L} \otimes \mathcal{N})_\xi$. Donc  $st/(s_\xi t_\xi) = (s/s_\xi)(t/t_\xi)$. Ainsi, nous voyons que
$$
\text{ord}_{Z, \mathcal{L} \otimes \mathcal{N}}(st)
=
\text{ord}_{Z, \mathcal{L}}(s) + \text{ord}_{Z, \mathcal{N}}(t)
$$
par l'additivité de la fonction  $\text{ord}_Z$ .
\end{proof}

\noindent
De cette manière, nous obtenons un homomorphisme de groupes abéliens
\begin{equation}
\label{divisors-equation-c1}
\Pic(X) \longrightarrow \text{Cl}(X)
\end{equation}
qui assigne à un module inversible sa classe de diviseur de Weil.

\begin{lemma}
\label{divisors-lemma-normal-c1-injective}
Soit  $X$  un schéma intègre localement noethérien. Si  $X$  est
normal, alors l'application (\ref{divisors-equation-c1})  $\Pic(X) \to \text{Cl}(X)$  est injective.
\end{lemma}

\begin{proof}
Soit  $\mathcal{L}$  un  $\mathcal{O}_X$-module inversible dont la classe de diviseur de
Weil associée est triviale. Soit  $s$  une section méromorphe régulière de
 $\mathcal{L}$. L'hypothèse signifie que  $\text{div}_\mathcal{L}(s) = \text{div}(f)$  pour un certain  $f \in R(X)^*$.
Alors nous voyons que  $t = f^{-1}s$  est une section méromorphe régulière de
 $\mathcal{L}$  avec  $\text{div}_\mathcal{L}(t) = 0$, voir le Lemme  \ref{divisors-lemma-divisor-meromorphic-well-defined}. Nous montrerons que
 $t$  définit une trivialisation de  $\mathcal{L}$ , ce qui termine la preuve
du lemme. Pour prouver cela, nous pouvons travailler localement sur  $X$.
Par conséquent, nous pouvons supposer que  $X = \Spec(A)$  est affine et que
 $\mathcal{L}$  est triviale. Alors  $A$  est un anneau noethérien intègre et normal, et
 $t$  est un élément de son corps de fractions tel que  $\text{ord}_{A_\mathfrak p}(t) = 0$  pour
tout idéal premier de hauteur  $1$,  $\mathfrak p$, de  $A$. Notre
objectif est de montrer que  $t$  est une unité de  $A$. Puisque
 $A_\mathfrak p$  est un anneau de valuation discrète pour les idéaux premiers de hauteur un
de  $A$  (Algèbre, Lemme  \ref{algebra-lemma-criterion-normal}), la condition signifie que
 $t \in A_\mathfrak p^*$  pour tous les idéaux premiers  $\mathfrak p$  de hauteur  $1$. Cela
implique  $t \in A$  et  $t^{-1} \in A$  d'après Algèbre, Lemme  \ref{algebra-lemma-normal-domain-intersection-localizations-height-1}, et la
preuve est complète.
\end{proof}

\begin{lemma}
\label{divisors-lemma-local-rings-UFD-c1-bijective}
Soit  $X$  un schéma intègre localement noethérien. Considérons l'application
(\ref{divisors-equation-c1})  $\Pic(X) \to \text{Cl}(X)$. Les conditions suivantes sont équivalentes
\begin{enumerate}
\item les anneaux locaux de  $X$  sont factoriels, et
\item $X$  est normal et  $\Pic(X) \to \text{Cl}(X)$  est surjective.
\end{enumerate}
Dans ce cas,  $\Pic(X) \to \text{Cl}(X)$  est un isomorphisme.
\end{lemma}

\begin{proof}
Si (1) est vrai, alors  $X$  est normal d'après Algèbre, Lemme  \ref{algebra-lemma-UFD-normal}.
Par conséquent, l'application (\ref{divisors-equation-c1}) est injective d'après le Lemme
 \ref{divisors-lemma-normal-c1-injective}. De plus, tout diviseur premier  $D \subset X$  est un diviseur de
Cartier effectif d'après le Lemme  \ref{divisors-lemma-weil-divisor-is-cartier-UFD}. Dans ce cas, la section canonique
 $1_D$  de  $\mathcal{O}_X(D)$  (Définition  \ref{divisors-definition-invertible-sheaf-effective-Cartier-divisor}) s'annule exactement le long
de  $D$  et nous voyons que la classe de  $D$  est l'image de
 $\mathcal{O}_X(D)$  sous l'application (\ref{divisors-equation-c1}). Ainsi, l'application est également
surjective.

\medskip\noindent
Supposons que (2) soit vrai. Choisissons un diviseur premier  $D \subset X$. Puisque
(\ref{divisors-equation-c1}) est surjective, il existe un faisceau inversible  $\mathcal{L}$, une
section méromorphe régulière  $s$, et  $f \in R(X)^*$  tels que  $\text{div}_\mathcal{L}(s) + \text{div}(f) = [D]$.
En d'autres termes,  $\text{div}_\mathcal{L}(fs) = [D]$. Soit  $x \in X$  et soit  $A = \mathcal{O}_{X, x}$. Ainsi
 $A$  est un anneau local noethérien intègre et normal, de corps des fractions
 $K = R(X)$. Chaque idéal premier de hauteur  $1$  de  $A$ 
correspond à un diviseur premier sur  $X$  et chaque  $\mathcal{O}_X$-module
inversible se restreint au module inversible trivial sur  $\Spec(A)$. Il s'ensuit
que pour chaque idéal premier de hauteur  $1$,  $\mathfrak p \subset A$, il existe un
élément  $f \in K$  tel que  $\text{ord}_{A_\mathfrak p}(f) = 1$  et  $\text{ord}_{A_{\mathfrak p'}}(f) = 0$  pour chaque autre idéal
premier  $\mathfrak p'$  de hauteur un. Alors  $f \in A$  par Algèbre, Lemme
 \ref{algebra-lemma-normal-domain-intersection-localizations-height-1}. En procédant de la même manière, nous voyons que chaque élément
 $g \in \mathfrak p$  est de la forme  $g = af$  pour un certain  $a \in A$. Ainsi,
nous voyons que chaque idéal premier de hauteur un de  $A$  est principal
et  $A$  est un anneau factoriel par Algèbre, Lemme  \ref{algebra-lemma-characterize-UFD-height-1}.
\end{proof}







\section{Compléments sur les modules inversibles}
\label{divisors-section-invertible-modules}

\noindent
Dans cette section, nous discutons de certaines propriétés des modules
inversibles.

\begin{lemma}
\label{divisors-lemma-in-image-pullback}
Soit  $\varphi : X \to Y$  un morphisme de schémas. Soit  $\mathcal{L}$  un
 $\mathcal{O}_X$-module inversible. Supposons que
\begin{enumerate}
\item $X$  soit localement noethérien,
\item $Y$  soit localement noethérien, intègre et normal,
\item $\varphi$  soit plat à fibres intègres (donc non vides),
\item $\varphi$  soit soit quasi-compact, soit localement de type fini,
\item $\mathcal{L}$  soit trivial lorsqu'on le restreint à la fibre générique de
 $\varphi$.
\end{enumerate}
Alors  $\mathcal{L} \cong \varphi^*\mathcal{N}$  pour un certain  $\mathcal{O}_Y$-module inversible  $\mathcal{N}$.
\end{lemma}

\begin{proof}
Soit  $\xi \in Y$  le point générique. Soit  $X_\xi$  la fibre schématique de
 $\varphi$  au-dessus de  $\xi$. On note  $\mathcal{L}_\xi$  l'image réciproque de
 $\mathcal{L}$  sur  $X_\xi$. L'hypothèse (5) signifie que  $\mathcal{L}_\xi$  est
trivial. Choisissons une section trivialisante  $s \in \Gamma(X_\xi, \mathcal{L}_\xi)$. On observe que
 $X$  est intègre d'après le Lemme  \ref{divisors-lemma-flat-over-integral-integral-fibre}. Ainsi, nous pouvons
considérer  $s$  comme une section méromorphe régulière de  $\mathcal{L}$.
Les images réciproques des fonctions méromorphes sont définies pour  $\varphi$ 
d'après le Lemme  \ref{divisors-lemma-pullback-meromorphic-sections-defined}. Soit  $\mathcal{N} \subset \mathcal{K}_Y$  le  $\mathcal{O}_Y$-module dont les
sections sur un ouvert  $V \subset Y$  sont ces fonctions méromorphes  $g \in \mathcal{K}_Y(V)$ 
telles que  $\varphi^*(g)s \in \mathcal{L}(\varphi^{-1}V)$. A priori  $\varphi^*(g)s$  est une section de  $\mathcal{K}_X(\mathcal{L})$  sur
 $\varphi^{-1}V$. Nous affirmons que  $\mathcal{N}$  est un  $\mathcal{O}_Y$-module
inversible et que l'application
$$
\varphi^*\mathcal{N} \longrightarrow \mathcal{L},\quad
g \longmapsto gs
$$
est un isomorphisme.

\medskip\noindent
Nous prouvons d'abord l'affirmation dans la situation suivante:  $X$  et
 $Y$  sont affines et  $\mathcal{L}$  est trivial. Disons que  $Y = \Spec(R)$,
 $X = \Spec(A)$, et  $s$  est donné par l'élément  $s \in A \otimes_R K$  où
 $K$  est le corps des fractions de  $R$. Nous pouvons écrire
 $s = a/r$  pour un certain  $r \in R$  non nul et  $a \in A$. Puisque
 $s$  engendre  $\mathcal{L}$  sur la fibre générique, nous voyons qu'il
existe un  $s' \in A \otimes_R K$  tel que  $ss' = 1$. Ainsi, nous voyons que  $s = r'/a'$ 
pour un certain  $r' \in R$  et  $a' \in A$ non nuls. Soit  $\mathfrak p_1, \ldots, \mathfrak p_n \subset R$  les
idéaux premiers minimaux au-dessus de  $rr'$. Chaque  $R_{\mathfrak p_i}$  est un
anneau de valuation discrète (Algèbre, Lemmes  \ref{algebra-lemma-minimal-over-1}  et  \ref{algebra-lemma-criterion-normal}). Par
hypothèse,  $\mathfrak q_i = \mathfrak p_i A$  est un idéal premier. Donc  $\mathfrak q_i A_{\mathfrak q_i}$  est engendré par
un seul élément et nous trouvons que  $A_{\mathfrak q_i}$  est aussi un anneau de valuation
discrète (Algèbre, Lemme  \ref{algebra-lemma-characterize-dvr}). Bien sûr,  $R_{\mathfrak p_i} \to A_{\mathfrak q_i}$  a un indice de
ramification  $1$. Soit  $e_i, e'_i \geq 0$  la valuation de  $a, a'$  dans
 $A_{\mathfrak q_i}$. Alors  $e_i + e'_i$  est la valuation de  $rr'$  dans
 $R_{\mathfrak p_i}$. Notons que
$$
\mathfrak p_1^{(e_1 + e'_1)} \cap \ldots \cap \mathfrak p_i^{(e_n + e'_n)} =
(rr')
$$
dans  $R$  d'après Algèbre, Lemme  \ref{algebra-lemma-normal-domain-intersection-localizations-height-1}. Posons
$$
I = \mathfrak p_1^{(e_1)} \cap \ldots \cap \mathfrak p_i^{(e_n)}
\quad\text{et}\quad
I' = \mathfrak p_1^{(e'_1)} \cap \ldots \cap \mathfrak p_i^{(e'_n)}
$$
de sorte que  $II' \subset (rr')$. Observons que
$$
IA =
(\mathfrak p_1^{(e_1)} \cap \ldots \cap \mathfrak p_i^{(e_n)})A =
(\mathfrak p_1A)^{(e_1)} \cap \ldots \cap (\mathfrak p_i A)^{(e_n)}
$$
d'après Algèbre, Lemmes  \ref{algebra-lemma-symbolic-power-flat-extension}  et  \ref{algebra-lemma-flat-intersect-ideals}. De même pour  $I'A$.
D'où  $a \in IA$  et  $a' \in I'A$. Nous concluons que  $IA \otimes_A I'A \to rr'A$  est surjectif.
Par la platitude fidèle de  $R \to A$ , nous trouvons que  $I \otimes_R I' \to (rr')$  est
également surjectif. Il s'ensuit que  $II' = (rr')$ ,  $I$  et  $I'$ 
sont finis localement libres de rang  $1$, voir Algèbre, Lemme
 \ref{algebra-lemma-product-ideals-principal}. Ainsi, localement pour la topologie de Zariski sur  $R$ ,
nous pouvons écrire  $I = (g)$  et  $I' = (g')$  avec  $gg' = rr'$. Ensuite,
 $a = ug$  et  $a' = u'g'$  pour certains  $u, u' \in A$. Nous concluons que
 $u, u'$  sont des unités. Ainsi, localement pour la topologie de Zariski sur
 $R$ , nous avons  $s = ug/r$  et l'affirmation s'ensuit dans ce cas.

\medskip\noindent
Soit  $y \in Y$  un point. Choisissons  $x \in X$  dont l'image est  $y$.
Nous pouvons appliquer le résultat du paragraphe précédent à  $\Spec(\mathcal{O}_{X, x}) \to \Spec(\mathcal{O}_{Y, y})$. Nous en
concluons qu'il existe un élément  $g \in R(Y)^*$  bien défini à un facteur près d'un
élément de  $\mathcal{O}_{Y, y}^*$  tel que  $\varphi^*(g)s$  génère  $\mathcal{L}_x$. Ainsi,
 $\varphi^*(g)s$  génère  $\mathcal{L}$  dans un voisinage  $U$  de  $x$.
Supposons que  $x'$  soit un second point au-dessus de  $y$  et que
 $g' \in R(Y)^*$  soit tel que  $\varphi^*(g')s$  génère  $\mathcal{L}$  dans un voisinage
ouvert  $U'$  de  $x'$. Alors nous pouvons choisir un point
 $x''$  dans  $U \cap U' \cap \varphi^{-1}(\{y\})$  car la fibre est irréductible. Par l'unicité pour
l'application d'anneaux  $\mathcal{O}_{Y, y} \to \mathcal{O}_{X, x''}$ , nous trouvons que  $g$  et
 $g'$  diffèrent (multiplicativement) par un élément de  $\mathcal{O}_{Y, y}^*$. Ainsi,
nous voyons que  $\varphi^*(g)s$  est un générateur pour  $\mathcal{L}$  sur un voisinage
ouvert de  $\varphi^{-1}(y)$. Soit  $Z \subset X$  l'ensemble des points  $z \in X$  tels
que  $\varphi^*(g)s$  ne génère pas  $\mathcal{L}_z$. Les arguments ci-dessus montrent que
 $Z$  est fermé et que  $Z = \varphi^{-1}(T)$  pour un certain sous-ensemble
 $T \subset Y$  avec  $y \not \in T$. Si nous pouvons montrer que  $T$  est
fermé, alors  $g$  sera un générateur pour  $\mathcal{N}$  en tant que
 $\mathcal{O}_Y$-module dans le voisinage ouvert  $Y \setminus T$  de  $y$ ,
ce qui achève la preuve (certains détails sont omis).

\medskip\noindent
Si  $\varphi$  est quasi-compact, alors  $T$  est fermé par les
morphismes, Lemme  \ref{morphisms-lemma-fpqc-quotient-topology}. Si  $\varphi$  est localement de type fini, alors
 $\varphi$  est ouvert par les morphismes, Lemme  \ref{morphisms-lemma-fppf-open}. Alors  $Y \setminus T$ 
est ouvert en tant qu'image de l'ouvert  $X \setminus Z$.
\end{proof}

\begin{lemma}
\label{divisors-lemma-closure-effective-cartier-divisor}
Soit  $X$  un schéma localement noethérien. Soit  $U \subset X$  un ouvert et
soit  $D \subset U$  un diviseur de Cartier effectif. Si  $\mathcal{O}_{X, x}$  est un anneau factoriel pour
tout  $x \in X \setminus U$, alors il existe un diviseur de Cartier effectif  $D' \subset X$ 
avec  $D = U \cap D'$.
\end{lemma}

\begin{proof}
Soit  $D' \subset X$  l'image schématique du morphisme  $D \to X$. Puisque
 $X$  est localement noethérien, le morphisme  $D \to X$  est
quasi-compact, voir Propriétés, Lemme  \ref{properties-lemma-immersion-into-noetherian}. Par conséquent, la formation
de  $D'$  commute avec le passage aux ouverts dans  $X$  d'après
Morphismes, Lemme  \ref{morphisms-lemma-quasi-compact-scheme-theoretic-image}. Ainsi, nous pouvons supposer que  $X = \Spec(A)$  est
affine. Soit  $I \subset A$  l'idéal correspondant à  $D'$. Soit  $\mathfrak p \subset A$ 
un idéal premier correspondant à un point de  $X \setminus U$. Pour terminer la
démonstration, il suffit de montrer que  $I_\mathfrak p$  est engendré par un élément,
voir Lemme  \ref{divisors-lemma-effective-Cartier-in-points}. Ainsi, nous pouvons remplacer  $X$  par
 $\Spec(A_\mathfrak p)$, voir Morphismes, Lemme  \ref{morphisms-lemma-flat-base-change-scheme-theoretic-image}. En d'autres termes, nous
pouvons supposer que  $X$  est le spectre d'un anneau local factoriel  $A$.
Alors tous les anneaux locaux de  $A$  sont factoriels. Il s'ensuit que
 $D = \sum a_i D_i$  avec  $D_i \subset U$  un diviseur de Cartier effectif et intégral, voir
Lemme  \ref{divisors-lemma-effective-Cartier-divisor-is-a-sum}. Les points génériques  $\xi_i$  de  $D_i$ 
correspondent aux idéaux premiers  $\mathfrak p_i \subset A$  de hauteur  $1$, voir Lemme
 \ref{divisors-lemma-effective-Cartier-codimension-1}. Alors  $\mathfrak p_i = (f_i)$  pour un certain élément premier  $f_i \in A$  et
nous concluons que  $D'$  est défini par  $\prod f_i^{a_i}$  comme souhaité.
\end{proof}

\begin{lemma}
\label{divisors-lemma-extend-invertible-module}
Soit  $X$  un schéma localement noethérien. Soit  $U \subset X$  un ouvert et
soit  $\mathcal{L}$  un  $\mathcal{O}_U$-module inversible. Si  $\mathcal{O}_{X, x}$  est un anneau factoriel
pour tout  $x \in X \setminus U$, alors il existe un  $\mathcal{O}_X$-module inversible
 $\mathcal{L}'$  avec  $\mathcal{L} \cong \mathcal{L}'|_U$.
\end{lemma}

\begin{proof}
Choisissons  $x \in X$,  $x \not \in U$. Nous montrerons qu'il existe un voisinage
ouvert affine  $W \subset X$, tel que  $\mathcal{L}|_{W \cap U}$  s'étend à un faisceau inversible
sur  $W$. Cela implique, par collage de faisceaux (Faisceaux, section
 \ref{sheaves-section-glueing-sheaves}), que nous pouvons étendre  $\mathcal{L}$  à l'ouvert strictement plus
grand  $U \cup W$. Soit  $W = \Spec(A)$  un voisinage ouvert affine. Puisque
 $U \cap W$  est quasi-affine, nous voyons que nous pouvons écrire  $\mathcal{L}|_{W \cap U}$ 
comme  $\mathcal{O}(D_1) \otimes \mathcal{O}(D_2)^{\otimes -1}$  pour certains diviseurs de Cartier effectifs  $D_1, D_2 \subset W \cap U$, voir
le Lemme  \ref{divisors-lemma-quasi-projective-Noetherian-pic-effective-Cartier}. Alors  $D_1$  et  $D_2$  s'étendent à des
diviseurs de Cartier effectifs de  $W$  par le Lemme  \ref{divisors-lemma-closure-effective-cartier-divisor} , ce qui
nous donne l'extension du faisceau inversible.

\medskip\noindent
Si  $X$  est noethérien (ce qui est le cas le plus utilisé en pratique), ce
qui précède, combiné à l'induction noethérienne, termine la preuve. Dans le cas
général, nous procédons comme suit. Tout d'abord, parce que chaque anneau local
d'un point à l'extérieur de  $U$  est un domaine et que  $X$  est
localement noethérien, nous voyons que l'adhérence de  $U$  dans
 $X$  est ouverte. Ainsi, nous pouvons supposer que  $U \subset X$  est dense
et schématiquement dense. Maintenant, nous considérons l'ensemble  $T$  des
triplets  $(U', \mathcal{L}', \alpha)$  où  $U \subset U' \subset X$  est un sous-schéma ouvert,  $\mathcal{L}'$  est
un  $\mathcal{O}_{U'}$-module inversible, et  $\alpha : \mathcal{L}'|_U \to \mathcal{L}$  est un isomorphisme. Nous
dotons  $T$  d'un ordre partiel  $\leq$  défini par la règle
 $(U', \mathcal{L}', \alpha) \leq (U'', \mathcal{L}'', \alpha')$  si et seulement si  $U' \subset U''$  et qu'il existe un isomorphisme
 $\beta : \mathcal{L}''|_{U'} \to \mathcal{L}'$  compatible avec  $\alpha$  et  $\alpha'$. Observons que
 $\beta$  est unique (si elle existe) car  $U \subset X$  est dense. La première
partie de la preuve montre que pour tout élément  $t = (U', \mathcal{L}', \alpha)$  de  $T$  avec
 $U' \not = X$ , il existe un  $t' \in T$  avec  $t' > t$. Ainsi, pour terminer la
preuve, il suffit de montrer que le lemme de Zorn s'applique. Considérons donc un
sous-ensemble totalement ordonné  $I \subset T$. Si  $i \in I$  correspond au
triple  $(U_i, \mathcal{L}_i, \alpha_i)$, alors nous pouvons construire un module inversible
 $\mathcal{L}'$  sur  $U' = \bigcup U_i$  de la manière suivante. Pour  $W \subset U'$  ouvert et
quasi-compact, nous voyons que  $W \subset U_i$  pour un certain  $i$  et nous
posons
$$
\mathcal{L}'(W) = \mathcal{L}_i(W)
$$
Comme applications de transition, nous utilisons les  $\beta$  (qui sont uniques
et se composent donc correctement). Cela définit un
 $\mathcal{O}$-module inversible $\mathcal{L}'$ sur la base des ouverts quasi-compacts
de $U'$, ce qui suffit à définir un module inversible (Faisceaux,
section  \ref{sheaves-section-bases}). Nous omettons les détails.
\end{proof}

\begin{lemma}
\label{divisors-lemma-open-subscheme-UFD}
Soit  $R$  un anneau factoriel. Les groupes de Picard des schémas suivants sont triviaux.
\begin{enumerate}
\item $\Spec(R)$  et tout sous-schéma ouvert de celui-ci.
\item $\mathbf{A}^n_R = \Spec(R[x_1, \ldots, x_n])$  et tout sous-schéma ouvert de celui-ci.
\end{enumerate}
En particulier, le groupe de Picard de tout sous-schéma ouvert de l'espace affine
 $n$  $\mathbf{A}^n_k$  sur un corps  $k$  est trivial.
\end{lemma}

\begin{proof}
Tout anneau de polynômes sur  $R$  est factoriel d'après Algèbre, Lemme
 \ref{algebra-lemma-polynomial-ring-UFD}. Ainsi, la question est de montrer que  $\Pic(U)$  est trivial si
 $U \subset \Spec(R)$  est un ouvert (non vide).

\medskip\noindent
Dans le cas noethérien, nous pouvons utiliser le Lemme  \ref{divisors-lemma-extend-invertible-module}  pour nous ramener à
montrer que  $\Pic(\Spec(R))$  est trivial. Ce dernier résultat est établi dans Compléments
d'algèbre, Lemme  \ref{more-algebra-lemma-UFD-Pic-trivial}. Nous conseillons au lecteur de sauter le reste de la
preuve.

\medskip\noindent
Soit  $U \subset \Spec(R)$  un ouvert non vide. Choisissons un  $f \in R$  non nul avec
 $D(f) \subset U$. Puisque  $f$  est un produit d'éléments premiers de
 $R$ , nous concluons qu'il existe un nombre fini d'éléments premiers
 $a_1, \ldots, a_n \in R$  tels que  $(a_i) \not \in U$  (ou, de manière équivalente,  $U \cap V(a_i) = \emptyset$). Après avoir
remplacé  $R$  par  $R_{a_1 \ldots a_n}$ , nous pouvons supposer que  $U$ 
contient tous les idéaux premiers principaux. (Notons qu'une localisation d'un anneau factoriel
est factorielle, par exemple d'après Algèbre, Lemme  \ref{algebra-lemma-invert-prime-elements}.)

\medskip\noindent
Supposons que  $U$  contienne tous les idéaux premiers principaux. Écrivons
 $U = \Spec(R) \setminus V(I)$  pour un certain idéal  $I \subset R$. Prenons  $f \in I$  non nul et
écrivons  $f = a_1 \ldots a_n$  comme un produit d'éléments premiers. Comme  $(a_i) \not \in V(I)$ ,
nous pouvons choisir  $g \in I$  avec  $g \not \in (a_i)$  pour tous les  $i$ 
(évitement des idéaux premiers). Considérons
$$
U' = D(f) \cup D(g) \subset U \subset \Spec(R)
$$
D'après Compléments d'algèbre, Lemme  \ref{more-algebra-lemma-UFD-Pic-trivial} , les groupes de Picard de
 $D(f)$  et  $D(g)$  sont triviaux et, par un argument élémentaire sur les
anneaux factoriels, l'application  $R_f^* \times R_g^* \to R_{fg}^*$,  $(u, v) \mapsto uv^{-1}$  est surjective. Nous concluons que
 $\Pic(U')$  est trivial.

\medskip\noindent
Ainsi, il suffit de montrer qu'un module inversible  $\mathcal{L}$  sur  $U$ 
dont l'image réciproque par  $j : U' \to U'$  est triviale est lui-même trivial. Pour ce faire, il
suffit de montrer que l'application  $\mathcal{L} \to j_*j^*\mathcal{L}$  est un isomorphisme. Pour cela à
son tour, soit  $u \in U \setminus U'$. Alors  $f, g \in \mathfrak m_u \subset \mathcal{O}_{U, u} = A$  sont des éléments d'un anneau local factoriel
qui n'ont pas de facteur premier en commun. Le lecteur montre facilement que cela
signifie que
$$
0 \to A \to A_f \times A_g \to A_{fg}
$$
est exacte en examinant les facteurs premiers. Puisque  $A = \mathcal{L}_u$  parce que
 $\mathcal{L}$  est inversible et puisque  $(j_*j^*\mathcal{L})_u$  est le noyau de  $A_f \times A_g \to A_{fg}$ ,
nous concluons.
\end{proof}

\begin{lemma}
\label{divisors-lemma-Pic-projective-space-UFD}
Soit  $R$  un anneau factoriel. Le groupe de Picard de  $\mathbf{P}^n_R$  est  $\mathbf{Z}$.
Plus précisément, il existe un isomorphisme
$$
\mathbf{Z} \longrightarrow \Pic(\mathbf{P}^n_R),\quad
m \longmapsto \mathcal{O}_{\mathbf{P}^n_R}(m)
$$
En particulier, le groupe de Picard de l'espace projectif  $\mathbf{P}^n_k$  sur un
corps  $k$  est  $\mathbf{Z}$.
\end{lemma}

\begin{proof}
D'après le Lemme  \ref{divisors-lemma-open-subscheme-UFD} , les groupes de Picard des ouverts  $D_+(T_i) \cong \mathbf{A}^n_R$  sont
triviaux. Ainsi, si  $\mathcal{L}$  est un module inversible sur  $\mathbf{P}^n_R$, alors
il est donné par un  $1$-cocycle à valeurs dans le faisceau des fonctions
inversibles  $\mathcal{O}^*$  pour le recouvrement ouvert
$$
\mathbf{P}^n_R = \bigcup\nolimits_{i = 0, \ldots, n} D_+(T_i)
$$
Remarquons que pour  $i \not = j$  nous avons
$$
\mathcal{O}^*(D_+(T_i) \cap D_+(T_j)) =
\mathcal{O}^*(D_+(T_iT_j)) = \left(R[T_0, \ldots, T_n]_{(T_iT_j)}\right)^* =
R^* \times (T_i/T_j)^\mathbf{Z}
$$
Ainsi, un tel cocycle  $(g_{ij})$  est donné par des unités
$$
g_{ij} = u_{ij} (T_i/T_j)^{e_{ij}}
$$
avec  $u_{ij} \in R^*$  et  $e_{ij} \in \mathbf{Z}$  satisfaisant à la condition de cocycle. La
condition de cocycle sur  $D_+(T_iT_jT_k)$  pour  $\#\{i, j, k\} = 3$  nous indique que
$$
u_{ik} = u_{ij} u_{jk} \quad\text{et}\quad
e_{ik} = e_{ij} = e_{jk}
$$
D'où  $u_{ij} = u_{i1} u_{j1}^{-1}$  est une frontière. Ainsi, toutes les classes d'isomorphisme de
modules inversibles sont obtenues en prenant le cocycle avec
$$
g_{ij} = (T_i/T_j)^e
$$
pour un certain  $e \in \mathbf{Z}$. Puisque  $\mathcal{O}(n)$  possède la section
trivialisante  $T_i^n$  sur  $D_+(T_i)$, nous voyons que le cocycle
correspondant à  $\mathcal{O}(n)$  est  $(T_i/T_j)^n$  et la preuve est complète.
\end{proof}






\section{Diviseurs de Weil sur les schémas normaux}
\label{divisors-section-weil-divisors-normal}

\noindent
Tout d'abord, nous discutons des propriétés des modules réflexifs.

\begin{lemma}
\label{divisors-lemma-reflexive-normal}
Soit  $X$  un schéma normal intègre et localement noethérien. Pour
 $\mathcal{F}$  et  $\mathcal{G}$  des  $\mathcal{O}_X$-modules réflexifs cohérents,
l'application
$$
(\SheafHom_{\mathcal{O}_X}(\mathcal{F}, \mathcal{O}_X)
\otimes_{\mathcal{O}_X} \mathcal{G})^{**} \to
\SheafHom_{\mathcal{O}_X}(\mathcal{F}, \mathcal{G})
$$
est un isomorphisme. La règle  $\mathcal{F}, \mathcal{G} \mapsto
(\mathcal{F} \otimes_{\mathcal{O}_X} \mathcal{G})^{**}$  définit une loi de groupe abélien sur
l'ensemble des classes d'isomorphisme des modules cohérents réflexifs de rang  $1$
sur  $\mathcal{O}_X$.
\end{lemma}

\begin{proof}
Bien que cela ne soit pas strictement nécessaire, nous recommandons de lire la
Remarque  \ref{divisors-remark-tensor}  avant de procéder à la démonstration. Choisissons un
sous-schéma ouvert  $j : U \to X$  de sorte que chaque composante irréductible de
 $X \setminus U$  ait codimension  $\geq 2$  dans  $X$  et de sorte que
 $j^*\mathcal{F}$  et  $j^*\mathcal{G}$  soient localement libres de type fini, voir le Lemme
 \ref{divisors-lemma-reflexive-over-normal}. L'application
$$
\SheafHom_{\mathcal{O}_U}(j^*\mathcal{F}, \mathcal{O}_U)
\otimes_{\mathcal{O}_U} j^*\mathcal{G} \to
\SheafHom_{\mathcal{O}_U}(j^*\mathcal{F}, j^*\mathcal{G})
$$
est un isomorphisme, car nous pouvons le vérifier localement et cela est clair
lorsque les modules sont libres de type fini. Observons que  $j^*$  appliqué à
la flèche affichée du lemme donne la flèche que nous venons de montrer comme étant
un isomorphisme (petit détail omis). Puisque  $j^*$  définit une équivalence
entre les modules réflexifs cohérents sur  $U$  et les modules réflexifs
cohérents sur  $X$  (d'après le Lemme  \ref{divisors-lemma-reflexive-S2-extend}  et le critère de Serre,
Lemme  \ref{properties-lemma-criterion-normal}), nous concluons que la flèche du lemme est aussi un
isomorphisme. Si  $\mathcal{F}$  a un rang  $1$, alors  $j^*\mathcal{F}$  est un
 $\mathcal{O}_U$-module inversible et le module réflexif  $\mathcal{F}^\vee = \SheafHom(\mathcal{F}, \mathcal{O}_X)$  se restreint à
son inverse. Il s'ensuit de la même manière qu'auparavant que  $(\mathcal{F} \otimes_{\mathcal{O}_X} \mathcal{F}^\vee)^{**} = \mathcal{O}_X$. De
cette façon, nous voyons que nous avons des inverses pour la loi de groupe donnée
dans l'énoncé du lemme.
\end{proof}

\begin{lemma}
\label{divisors-lemma-normal-class-group}
Soit $X$ un schéma normal, intègre et localement noethérien. Le
groupe des modules cohérents réflexifs de rang $1$ sur $\mathcal{O}_X$ est isomorphe au
groupe des classes de diviseurs de Weil $\text{Cl}(X)$  de  $X$.
\end{lemma}

\begin{proof}
Soit $\mathcal{F}$ un module cohérent réflexif de rang $1$ sur $\mathcal{O}_X$.
Choisissons un ouvert $U \subset X$  tel que chaque composante irréductible
de  $X \setminus U$ ait codimension $\geq 2$ dans $X$  et que
 $\mathcal{F}|_U$  soit inversible, voir le Lemme  \ref{divisors-lemma-reflexive-over-normal}. Remarquons que
 $\text{Cl}(U) = \text{Cl}(X)$, car le groupe des classes de diviseurs de Weil de $X$ ne
dépend que de son corps des fonctions rationnelles, de ses points de
codimension $1$ et de leurs anneaux locaux. Nous pouvons donc définir la
classe de diviseur de Weil de  $\mathcal{F}$ comme la classe de diviseur de Weil
de $\mathcal{F}|_U$ dans $\text{Cl}(U)$. Nous omettons la vérification que ceci est
indépendant du choix de  $U$.

\medskip\noindent
Désignons par  $\text{Cl}'(X)$  l'ensemble des classes d'isomorphisme des modules
cohérents réflexifs de rang  $1$  sur  $\mathcal{O}_X$. La construction ci-dessus donne
un homomorphisme de groupes
$$
\text{Cl}'(X) \longrightarrow \text{Cl}(X)
$$
car pour toute paire  $\mathcal{F}, \mathcal{G}$  d'éléments de  $\text{Cl}'(X)$ , nous pouvons choisir
un  $U$  qui convient pour les deux et l'application (\ref{divisors-equation-c1})
envoyant un module inversible à sa classe de diviseur de Weil est un
homomorphisme. Si  $\mathcal{F}$  est dans le noyau de cette application, alors nous
trouvons que  $\mathcal{F}|_U$  est trivial (Lemme  \ref{divisors-lemma-normal-c1-injective}) et donc  $\mathcal{F}$ 
est aussi trivial d'après le Lemme  \ref{divisors-lemma-reflexive-S2-extend}  et le critère de Serre Propriétés,
Lemme  \ref{properties-lemma-criterion-normal}. Pour terminer la preuve, il suffit de vérifier que
l'application est surjective.

\medskip\noindent
Soit  $D = \sum n_Z Z$  un diviseur de Weil sur  $X$. Nous affirmons qu'il
existe un ouvert  $U \subset X$  tel que chaque composante irréductible de
 $X \setminus U$  ait codimension  $\geq 2$  dans  $X$  et tel que
 $Z|_U$  soit un diviseur de Cartier effectif pour  $n_Z \not = 0$. Pour prouver
cette affirmation, nous pouvons supposer que  $X$  est affine. Ensuite,
nous pouvons supposer que  $D = n_1 Z_1 + \ldots + n_r Z_r$  est une somme finie avec  $Z_1, \ldots, Z_r$  deux
à deux distincts. Après avoir éliminé  $Z_i \cap Z_j$  pour  $i \not = j$ , nous pouvons
supposer que  $Z_1, \ldots, Z_r$  sont deux à deux disjoints. Cela nous ramène au cas d'un
seul diviseur premier  $Z$  sur  $X$. Comme  $X$  est
 vérifie la propriété  $(R_1)$  d'après Propriétés, Lemme  \ref{properties-lemma-criterion-normal}; l'anneau local
 $\mathcal{O}_{X, \xi}$  au point générique  $\xi$  de  $Z$  est un anneau de
valuation discrète. Soit  $f \in \mathcal{O}_{X, \xi}$  une uniformisante. Soit  $V \subset X$  un
voisinage ouvert de  $\xi$  tel que  $f$  soit l'image d'un élément
 $f \in \mathcal{O}_X(V)$. Après avoir réduit  $V$ , nous pouvons supposer que
 $Z \cap V = V(f)$  schématiquement, puisque c'est vrai dans l'anneau local en
 $\xi$. Dans ce cas, prendre
$$
U = X \setminus (Z \setminus V) = (X \setminus Z) \cup V
$$
donne l'ouvert souhaité, prouvant ainsi l'affirmation.

\medskip\noindent
Afin de montrer que la classe de diviseur de  $D$  est dans l'image, nous
pouvons écrire  $D = \sum_{n_Z < 0} n_Z Z - \sum_{n_Z > 0} (-n_Z) Z$. Par additivité de l'application construite ci-dessus,
nous pouvons et supposons  $n_Z \leq 0$  pour tous les diviseurs premiers
 $Z$  (cette étape peut être évitée si le lecteur le souhaite). Soit
 $U \subset X$  comme dans l'assertion ci-dessus. Si  $U$  est quasi-compact,
alors nous écrivons  $D|_U = -n_1 Z_1 - \ldots - n_r Z_r$  pour des diviseurs premiers deux à deux distincts
 $Z_i$  et  $n_i > 0$  et nous considérons le  $\mathcal{O}_U$-module inversible
$$
\mathcal{L} =
\mathcal{I}_1^{n_1} \ldots \mathcal{I}_r^{n_r} \subset \mathcal{O}_U
$$
où  $\mathcal{I}_i$  est le faisceau d'idéaux de  $Z_i$. C'est inversible par
notre choix de  $U$  et le Lemme  \ref{divisors-lemma-sum-effective-Cartier-divisors}. De plus,  $\text{div}_\mathcal{L}(1) = D|_U$.
Puisque  $\mathcal{L} = \mathcal{F}|_U$  pour un certain module cohérent réflexif de rang  $1$  sur
 $\mathcal{O}_X$, que nous notons  $\mathcal{F}$, le Lemme  \ref{divisors-lemma-reflexive-S2-extend}  montre que
 $D$  est dans l'image de notre application.

\medskip\noindent
Si  $U$  n'est pas quasi-compact, alors nous définissons  $\mathcal{L} \subset \mathcal{O}_U$ 
localement par la formule affichée ci-dessus. Le lecteur montre que la
construction se recolle et termine la preuve exactement comme auparavant. Détails
omis.
\end{proof}

\begin{lemma}
\label{divisors-lemma-structure-sheaf-Xs}
Soit  $X$  un schéma intègre, normal et localement noethérien. Soit
 $\mathcal{F}$  un  $\mathcal{O}_X$-module cohérent réflexif de rang 1. Soit
 $s \in \Gamma(X, \mathcal{F})$. Posons
$$
U = \{x \in X \mid s : \mathcal{O}_{X, x} \to \mathcal{F}_x
\text{ est un isomorphisme}\}
$$
Alors  $j : U \to X$  est un sous-schéma ouvert de  $X$  et
$$
j_*\mathcal{O}_U =
\colim (\mathcal{O}_X \xrightarrow{s} \mathcal{F}
\xrightarrow{s} \mathcal{F}^{[2]}
\xrightarrow{s} \mathcal{F}^{[3]}
\xrightarrow{s} \ldots)
$$
où  $\mathcal{F}^{[1]} = \mathcal{F}$  et par induction  $\mathcal{F}^{[n + 1]} =
(\mathcal{F} \otimes_{\mathcal{O}_X} \mathcal{F}^{[n]})^{**}$.
\end{lemma}

\begin{proof}
L'ensemble  $U$  est ouvert d'après Modules, Lemme  \ref{modules-lemma-finite-type-surjective-on-stalk}  et
 \ref{modules-lemma-finite-type-to-coherent-injective-on-stalk}. On observe que  $j$  est quasi-compact d'après Propriétés,
Lemme  \ref{properties-lemma-immersion-into-noetherian}. Pour prouver l'affirmation finale, il suffit de montrer que
pour tout ouvert quasi-compact  $W \subset X$ , il existe un isomorphisme
$$
\colim \Gamma(W, \mathcal{F}^{[n]})
\longrightarrow
\Gamma(U \cap W, \mathcal{O}_U)
$$
de  $\mathcal{O}_X(W)$-modules compatible avec les applications de restriction. Nous
omettrons la vérification des compatibilités. Après avoir remplacé  $X$ 
par  $W$  et réécrit ce qui précède en termes de morphismes, nous voyons qu'il
suffit de construire un isomorphisme
$$
\colim \Hom_{\mathcal{O}_X}(\mathcal{O}_X, \mathcal{F}^{[n]})
\longrightarrow
\Hom_{\mathcal{O}_U}(\mathcal{O}_U, \mathcal{O}_U)
$$
Choisissons un ouvert  $V \subset X$  tel que chaque composante irréductible de
 $X \setminus V$  ait une codimension  $\geq 2$  dans  $X$  et tel que
 $\mathcal{F}|_V$  soit inversible, voir le Lemme  \ref{divisors-lemma-reflexive-over-normal}. Alors la restriction
définit une équivalence de catégories entre les modules cohérents
réflexifs de rang  $1$  sur  $X$  et  $V$  et entre les
modules cohérents réflexifs de rang  $1$  sur  $U$  et  $V \cap U$.
Voir le Lemme  \ref{divisors-lemma-reflexive-S2-extend}  et le critère de Serre Propriétés, Lemme  \ref{properties-lemma-criterion-normal}.
Ainsi il suffit de construire un isomorphisme
$$
\colim \Gamma(V, (\mathcal{F}|_V)^{\otimes n}) \longrightarrow
\Gamma(V \cap U, \mathcal{O}_U)
$$
Puisque  $\mathcal{F}|_V$  est inversible et puisque  $U \cap V$  est l'ensemble
des points où  $s|_V$  génère ce module inversible, il s'agit d'un cas
particulier des Propriétés, Lemme  \ref{properties-lemma-invert-s-sections}  (il existe également une formule
explicite pour l'application).
\end{proof}

\begin{lemma}
\label{divisors-lemma-Xs-codim-complement}
Sous les hypothèses et avec les notations du Lemme  \ref{divisors-lemma-structure-sheaf-Xs}, si  $s$  est non
nul, alors chaque composante irréductible de  $X \setminus U$  est de codimension
 $1$  dans  $X$.
\end{lemma}

\begin{proof}
Soit  $\xi \in X$  un point générique d'une composante irréductible  $Z$  de
 $X \setminus U$. Après avoir remplacé  $X$  par un voisinage ouvert de
 $\xi$ , on peut supposer que  $Z = X \setminus U$  est irréductible. Puisque
 $s : \mathcal{O}_U \to \mathcal{F}|_U$  est un isomorphisme, si la codimension de  $Z$  dans
 $X$  est  $\geq 2$, alors  $s : \mathcal{O}_X \to \mathcal{F}$  est un isomorphisme d'après le
Lemme  \ref{divisors-lemma-reflexive-S2-extend}  et le critère de Serre Propriétés, Lemme  \ref{properties-lemma-criterion-normal}. Cela
signifierait que  $Z = \emptyset$, ce qui constitue une contradiction.
\end{proof}

\begin{remark}
\label{divisors-remark-structure-sheaf-Xs}
Soit  $A$  un anneau noethérien intègre et normal. Soit  $M$  un
module réflexif fini de rang  $1$  sur  $A$. Soit  $s \in M$  non nul. Soit
 $\mathfrak p_1, \ldots, \mathfrak p_r$  les idéaux premiers de hauteur  $1$  de  $A$  dans le
support de  $M/As$. Alors l'ouvert  $U$  du Lemme  \ref{divisors-lemma-structure-sheaf-Xs}  est
$$
U = \Spec(A) \setminus
\left(V(\mathfrak p_1) \cup \ldots \cup V(\mathfrak p_r)\right)
$$
d'après le Lemme  \ref{divisors-lemma-Xs-codim-complement}. De plus, si  $M^{[n]}$  désigne l'enveloppe
réflexive de  $M \otimes_A \ldots \otimes_A M$  ($n$ facteurs), alors
$$
\Gamma(U, \mathcal{O}_U) = \colim M^{[n]}
$$
d'après le Lemme  \ref{divisors-lemma-structure-sheaf-Xs}.
\end{remark}

\begin{lemma}
\label{divisors-lemma-affine-Xs}
Sous les hypothèses et avec les notations du Lemme  \ref{divisors-lemma-structure-sheaf-Xs}, les conditions suivantes
sont équivalentes
\begin{enumerate}
\item le morphisme d'inclusion  $j : U \to X$  est affine, et
\item pour chaque  $x \in X \setminus U$  il existe un  $n > 0$  tel que  $s^n \in \mathfrak m_x \mathcal{F}^{[n]}_x$.
\end{enumerate}
\end{lemma}

\begin{proof}
Supposons (1). Alors pour  $x \in X \setminus U$  l'image réciproque  $U_x$  de
 $U$  sous le morphisme canonique  $f_x : \Spec(\mathcal{O}_{X, x}) \to X$  est affine et ne contient
pas  $x$. Ainsi,  $\mathfrak m_x \Gamma(U_x, \mathcal{O}_{U_x})$  est l'idéal unité. En particulier, nous
voyons que nous pouvons écrire
$$
1 = \sum f_i g_i
$$
avec  $f_i \in \mathfrak m_x$  et  $g_i \in \Gamma(U_x, \mathcal{O}_{U_x})$. Par le Lemme  \ref{divisors-lemma-structure-sheaf-Xs}  nous avons
 $\Gamma(U_x, \mathcal{O}_{U_x}) = \colim \mathcal{F}^{[n]}_x$  avec des applications de transition données par la multiplication par
 $s$. Donc pour certains  $n > 0$  nous avons
$$
s^n = \sum f_i t_i
$$
pour certains  $t_i = s^ng_i \in \mathcal{F}^{[n]}_x$. Ainsi (2) est vrai.

\medskip\noindent
Inversement, supposons (2). Le caractère affine de  $j$  est
local sur  $X$, voir Morphismes, Lemme  \ref{morphisms-lemma-characterize-affine}. Nous pouvons donc
supposer que  $X$  est affine. Notre objectif est de montrer que
 $U$  est affine. D'après Cohomologie des schémas, Lemme  \ref{coherent-lemma-affine-if-quasi-affine}, il suffit
de montrer que  $H^p(U, \mathcal{O}_U) = 0$  pour  $p > 0$. Puisque  $H^p(U, \mathcal{O}_U) = H^0(X, R^pj_*\mathcal{O}_U)$  (Cohomologie
des schémas, Lemme  \ref{coherent-lemma-quasi-coherence-higher-direct-images-application}) et puisque  $R^pj_*\mathcal{O}_U$  est quasi-cohérente
(Cohomologie des schémas, Lemme  \ref{coherent-lemma-quasi-coherence-higher-direct-images}), il suffit de montrer que la fibre
 $(R^pj_*\mathcal{O}_U)_x$  en tout point  $x \in X$  est nul. Considérons le diagramme de
changement de base
$$
\xymatrix{
U_x \ar[d]_{j_x} \ar[r] & U \ar[d]^j \\
\Spec(\mathcal{O}_{X, x}) \ar[r] & X
}
$$
Par la Cohomologie des schémas, Lemme  \ref{coherent-lemma-flat-base-change-cohomology}  nous avons  $(R^pj_*\mathcal{O}_U)_x = R^pj_{x, *}\mathcal{O}_{U_x}$. Par
conséquent, nous pouvons supposer que  $X$  est local avec un point fermé
 $x$  et nous devons montrer que  $U$  est affine (car cela équivaut
à l'annulation souhaitée d'après la référence donnée ci-dessus). En particulier,
 $d = \dim(X)$  est fini (Algèbre, Proposition  \ref{algebra-proposition-dimension}). Si  $x \in U$, alors
 $U = X$  et le résultat est clair. Si  $d = 0$  et  $x \not \in U$, alors
 $U = \emptyset$  et le résultat est clair. Supposons maintenant  $d > 0$  et
 $x \not \in U$. Comme  $j_*\mathcal{O}_U = \colim \mathcal{F}^{[n]}$ , notre hypothèse signifie que nous pouvons écrire
$$
1 = \sum f_i g_i
$$
pour certains  $n > 0$,  $f_i \in \mathfrak m_x$, et  $g_i \in \mathcal{O}(U)$. Par induction sur
 $d$ , nous savons que  $D(f_i) \cap U$  est affine pour tout  $i$: en
parcourant tout l'argument donné avec  $X$  remplacé par  $D(f_i)$ , nous
obtenons des anneaux locaux noethériens dont la dimension est strictement
inférieure à  $d$. Par conséquent,  $U$  est affine d'après
Propriétés, Lemme  \ref{properties-lemma-characterize-affine}, comme voulu.
\end{proof}





\section{Proj relatif}
\label{divisors-section-relative-proj}

\noindent
Quelques résultats sur le Proj relatif. Commençons par des résultats très
élémentaires. Rappelons qu'un Proj relatif est toujours séparé sur la base, voir
Constructions, Lemme  \ref{constructions-lemma-relative-proj-separated}.

\begin{lemma}
\label{divisors-lemma-relative-proj-quasi-compact}
Soit  $S$  un schéma. Soit  $\mathcal{A}$  une  $\mathcal{O}_S$-algèbre graduée
quasi-cohérente. Soit  $p : X = \underline{\text{Proj}}_S(\mathcal{A}) \to S$  le Proj relatif de  $\mathcal{A}$. Si l'une des
conditions suivantes est satisfaite
\begin{enumerate}
\item $\mathcal{A}$  est de type fini comme faisceau de  $\mathcal{A}_0$-algèbres,
\item $\mathcal{A}$  est engendrée par  $\mathcal{A}_1$  comme  $\mathcal{A}_0$-algèbre et
 $\mathcal{A}_1$  est un  $\mathcal{A}_0$-module de type fini,
\item il existe un  $\mathcal{A}_0$-sous-module quasi-cohérent de type fini  $\mathcal{F} \subset \mathcal{A}_{+}$  tel
que  $\mathcal{A}_{+}/\mathcal{F}\mathcal{A}$  soit un faisceau d'idéaux localement nilpotent de  $\mathcal{A}/\mathcal{F}\mathcal{A}$,
\end{enumerate}
alors  $p$  est quasi-compact.
\end{lemma}

\begin{proof}
La question est locale sur la base, voir Schémas, Lemme  \ref{schemes-lemma-quasi-compact-affine}. Ainsi, nous
pouvons supposer que  $S$  est affine. Disons que  $S = \Spec(R)$  et
 $\mathcal{A}$  correspond à la  $R$-algèbre graduée  $A$. Alors
 $X = \text{Proj}(A)$, voir Constructions, section  \ref{constructions-section-relative-proj-via-glueing}. Dans le cas (1), nous
pouvons, quitte à localiser encore, supposer que  $A$ 
est généré par des éléments homogènes  $f_1, \ldots, f_n \in A_{+}$  sur  $A_0$. Alors
 $A_{+} = (f_1, \ldots, f_n)$  par Algèbre, Lemme  \ref{algebra-lemma-S-plus-generated}. Dans le cas (3), nous voyons que
 $\mathcal{F} = \widetilde{M}$  pour un certain  $A_0$-module de type fini  $M \subset A_{+}$. Disons
 $M = \sum A_0f_i$. Disons que  $f_i = \sum f_{i, j}$  est la décomposition en parties homogènes.
La condition du cas (3) signifie que  $A_{+} \subset \sqrt{(f_{i, j})}$. Ainsi, dans les deux cas, nous
concluons que  $\text{Proj}(A)$  est quasi-compact par Constructions, Lemme
 \ref{constructions-lemma-proj-quasi-compact}. Enfin, (2) découle de (1).
\end{proof}

\begin{lemma}
\label{divisors-lemma-relative-proj-finite-type}
Soit  $S$  un schéma. Soit  $\mathcal{A}$  une  $\mathcal{O}_S$-algèbre graduée
quasi-cohérente. Soit  $p : X = \underline{\text{Proj}}_S(\mathcal{A}) \to S$  le Proj relatif de  $\mathcal{A}$. Si
 $\mathcal{A}$  est de type fini en tant que faisceau de  $\mathcal{O}_S$-algèbres, alors
 $p$  est de type fini et  $\mathcal{O}_X(d)$  est un  $\mathcal{O}_X$-module de type
fini.
\end{lemma}

\begin{proof}
L'hypothèse implique que  $p$  est quasi-compact, voir le Lemme
 \ref{divisors-lemma-relative-proj-quasi-compact}. Il suffit donc de montrer que  $p$  est localement de type
fini. Ainsi, la question est locale sur la base et la cible, voir Morphismes,
Lemme  \ref{morphisms-lemma-locally-finite-type-characterize}. Supposons que  $S = \Spec(R)$  et que  $\mathcal{A}$  corresponde à
la  $R$-algèbre graduée  $A$. Après une nouvelle localisation sur
 $S$ , nous pouvons supposer que  $A$  est une  $R$-algèbre
de type fini. Le schéma  $X$  est construit en collant les spectres des
anneaux  $A_{(f)}$  pour  $f \in A_{+}$  homogène. Chacun d'eux est de type fini sur
 $R$  d'après Algèbre, Lemme  \ref{algebra-lemma-dehomogenize-finite-type}  partie (1). Ainsi  $\text{Proj}(A)$ 
est de type fini sur  $R$. Pour voir l'énoncé sur  $\mathcal{O}_X(d)$ , utilisons
la partie (2) de Algèbre, Lemme  \ref{algebra-lemma-dehomogenize-finite-type}.
\end{proof}

\begin{lemma}
\label{divisors-lemma-relative-proj-universally-closed}
Soit  $S$  un schéma. Soit  $\mathcal{A}$  une  $\mathcal{O}_S$-algèbre graduée
quasi-cohérente. Soit  $p : X = \underline{\text{Proj}}_S(\mathcal{A}) \to S$  le Proj relatif de  $\mathcal{A}$. Si
 $\mathcal{O}_S \to \mathcal{A}_0$  est un morphisme d'algèbres entier \footnote{ En
d'autres termes, la clôture intégrale de  $\mathcal{O}_S$  dans  $\mathcal{A}_0$, voir
Morphismes, Définition  \ref{morphisms-definition-integral-closure}, est égale à  $\mathcal{A}_0$. }  et
 $\mathcal{A}$  est de type fini comme  $\mathcal{A}_0$-algèbre, alors  $p$  est
universellement fermé.
\end{lemma}

\begin{proof}
La question est locale sur la base. Ainsi, nous pouvons supposer que  $X = \Spec(R)$ 
est affine. Soit  $\mathcal{A}$  la  $\mathcal{O}_X$-algèbre quasi-cohérente associée à la
 $R$-algèbre graduée  $A$. L'hypothèse est que  $R \to A_0$  est
entière et que  $A$  est de type fini sur  $A_0$. Écrivons
 $X \to \Spec(R)$  comme la composition  $X \to \Spec(A_0) \to \Spec(R)$. Comme  $R \to A_0$  est un
morphisme d'anneaux entier, nous voyons que  $\Spec(A_0) \to \Spec(R)$  est universellement
fermé, voir Morphismes, Lemme  \ref{morphisms-lemma-integral-universally-closed}. Le morphisme quasi-compact (voir
Constructions, Lemme  \ref{constructions-lemma-proj-quasi-compact})
$$
X = \text{Proj}(A) \to \Spec(A_0)
$$
satisfait la partie existence du critère valuatif par Constructions, Lemme
 \ref{constructions-lemma-proj-valuative-criterion}  et donc il est universellement fermé par Schémas, Proposition
 \ref{schemes-proposition-characterize-universally-closed}. Ainsi,  $X \to \Spec(R)$  est universellement fermé comme composition de
morphismes universellement fermés.
\end{proof}

\begin{lemma}
\label{divisors-lemma-relative-proj-proper}
Soit  $S$  un schéma. Soit  $\mathcal{A}$  une  $\mathcal{O}_S$-algèbre graduée
quasi-cohérente. Soit  $p : X = \underline{\text{Proj}}_S(\mathcal{A}) \to S$  le Proj relatif de  $\mathcal{A}$. Les conditions
suivantes sont équivalentes
\begin{enumerate}
\item $\mathcal{A}_0$  est un  $\mathcal{O}_S$-module de type fini et  $\mathcal{A}$  est de type
fini comme  $\mathcal{A}_0$-algèbre,
\item $\mathcal{A}_0$  est un  $\mathcal{O}_S$-module de type fini et  $\mathcal{A}$  est de type
fini comme  $\mathcal{O}_S$-algèbre
\end{enumerate}
Si ces conditions sont vérifiées, alors  $p$  est localement projectif et
en particulier propre.
\end{lemma}

\begin{proof}
Supposons que  $\mathcal{A}_0$  soit un  $\mathcal{O}_S$-module de type fini. Choisissons un
ouvert affine  $U = \Spec(R) \subset X$  tel que  $\mathcal{A}$  corresponde à une
 $R$-algèbre graduée  $A$  avec  $A_0$  un  $R$-module
fini. La condition (1) signifie que (après éventuellement une localisation
supplémentaire sur  $S$)  $A$  est une  $A_0$-algèbre de type
fini et la condition (2) signifie que (après éventuellement une localisation
supplémentaire sur  $S$)  $A$  est une  $R$-algèbre de type
fini. Ainsi, ces conditions s'impliquent mutuellement d'après le Lemme d'Algèbre
 \ref{algebra-lemma-compose-finite-type}.

\medskip\noindent
Un morphisme localement projectif est propre, voir Morphismes, Lemme  \ref{morphisms-lemma-locally-projective-proper}.
Ainsi, nous pouvons maintenant supposer que  $S = \Spec(R)$  et  $X = \text{Proj}(A)$  et que
 $A_0$  est fini sur  $R$  et  $A$  de type fini sur
 $R$. Nous allons montrer que  $X = \text{Proj}(A) \to \Spec(R)$  est projectif. Nous
encourageons le lecteur à le prouver lui-même, en construisant directement une
immersion fermée de  $X$  dans un espace projectif sur  $R$, plutôt
qu'en lisant l'argument que nous donnons ci-dessous.

\medskip\noindent
D'après le Lemme  \ref{divisors-lemma-relative-proj-finite-type} , nous voyons que  $X$  est de type fini sur
 $\Spec(R)$. Constructions, le Lemme  \ref{constructions-lemma-ample-on-proj}  nous dit que  $\mathcal{O}_X(d)$  est
ample sur  $X$  pour un certain  $d \geq 1$  (voir Propriétés, section
 \ref{properties-section-ample}). Par conséquent,  $X \to \Spec(R)$  est quasi-projectif (d'après
Morphismes, Définition  \ref{morphisms-definition-quasi-projective}). D'après Morphismes, le Lemme  \ref{morphisms-lemma-quasi-projective-open-projective} ,
nous concluons que  $X$  est isomorphe à un sous-schéma ouvert d'un schéma
projectif sur  $\Spec(R)$. Par conséquent, pour terminer la démonstration, il
suffit de montrer que  $X \to \Spec(R)$  est universellement fermé (utiliser Morphismes,
Lemme  \ref{morphisms-lemma-image-proper-scheme-closed}). Cela découle du Lemme  \ref{divisors-lemma-relative-proj-universally-closed}.
\end{proof}

\begin{lemma}
\label{divisors-lemma-relative-proj-projective}
Soit  $S$  un schéma. Soit  $\mathcal{A}$  une  $\mathcal{O}_S$-algèbre graduée
quasi-cohérente. Soit  $p : X = \underline{\text{Proj}}_S(\mathcal{A}) \to S$  le Proj relatif de  $\mathcal{A}$. Si
 $\mathcal{A}$  est engendrée par  $\mathcal{A}_1$  sur  $\mathcal{A}_0$  et  $\mathcal{A}_1$  est un
 $\mathcal{O}_S$-module de type fini, alors  $p$  est projectif.
\end{lemma}

\begin{proof}
À savoir, le morphisme associé au morphisme de  $\mathcal{O}_S$-algèbres graduées
$$
\text{Sym}_{\mathcal{O}_X}^*(\mathcal{A}_1)
\longrightarrow
\mathcal{A}
$$
est une immersion fermée  $X \to \mathbf{P}(\mathcal{A}_1)$, voir Constructions, Lemme  \ref{constructions-lemma-surjective-generated-degree-1-map-relative-proj}.
\end{proof}

\begin{lemma}
\label{divisors-lemma-relative-proj-flat}
Soit  $S$  un schéma. Soit  $\mathcal{A}$  une  $\mathcal{O}_S$-algèbre graduée
quasi-cohérente. Soit  $p : X = \underline{\text{Proj}}_S(\mathcal{A}) \to S$  le Proj relatif de  $\mathcal{A}$. Si
 $\mathcal{A}_d$  est un  $\mathcal{O}_S$-module plat pour  $d \gg 0$, alors  $p$ 
est plat et  $\mathcal{O}_X(d)$  est plat sur  $S$.
\end{lemma}

\begin{proof}
La vérification de la platitude de  $X$  sur  $S$  étant locale sur des ouverts affines, elle se ramène à
l'énoncé suivant : soit  $R$  un anneau, soit  $A$  une
 $R$-algèbre graduée avec  $A_d$  plat sur  $R$  pour
 $d \gg 0$, soit  $f \in A_d$  pour un certain  $d > 0$, alors  $A_{(f)}$ 
est plat sur  $R$. Comme  $A_{(f)} = \colim A_{nd}$  où les applications de transition sont
données par multiplication par  $f$, ceci découle du Lemme  \ref{algebra-lemma-colimit-flat}
d'Algèbre. Le même argument donne la platitude de
 $\mathcal{O}_X(d)$  sur  $S$.
\end{proof}

\begin{lemma}
\label{divisors-lemma-relative-proj-finite-presentation}
Soit  $S$  un schéma. Soit  $\mathcal{A}$  une  $\mathcal{O}_S$-algèbre graduée
quasi-cohérente. Soit  $p : X = \underline{\text{Proj}}_S(\mathcal{A}) \to S$  le Proj relatif de  $\mathcal{A}$. Si
 $\mathcal{A}$  est une  $\mathcal{O}_S$-algèbre de présentation finie, alors
 $p$  est de présentation finie et  $\mathcal{O}_X(d)$  est un  $\mathcal{O}_X$-module
de présentation finie.
\end{lemma}

\begin{proof}
Localement affine, cela se réduit à l'énoncé suivant: Soit  $R$  un anneau
et soit  $A$  une  $R$-algèbre graduée de présentation finie. Alors
 $\text{Proj}(A) \to \Spec(R)$  est de présentation finie et  $\mathcal{O}_{\text{Proj}(A)}(d)$  est un  $\mathcal{O}_{\text{Proj}(A)}$-module
de présentation finie. La condition de présentation finie implique que nous
pouvons choisir une présentation
$$
A = R[X_1, \ldots, X_n]/(F_1, \ldots, F_m)
$$
où  $R[X_1, \ldots, X_n]$  est un anneau de polynômes gradué en donnant des poids  $d_i$ 
à  $X_i$  et  $F_1, \ldots, F_m$  sont des polynômes homogènes de degré  $e_j$.
Soit  $R_0 \subset R$  le sous-anneau généré par les coefficients des polynômes
 $F_1, \ldots, F_m$. Ensuite, nous posons  $A_0 = R_0[X_1, \ldots, X_n]/(F_1, \ldots, F_m)$. Par construction,  $A = A_0 \otimes_{R_0} R$.
Ainsi, d'après Constructions, Lemme  \ref{constructions-lemma-base-change-map-proj} , il suffit de prouver le résultat
pour  $X_0 = \text{Proj}(A_0)$  sur  $R_0$. D'après le Lemme  \ref{divisors-lemma-relative-proj-finite-type} , nous savons que
 $X_0$  est de type fini sur  $R_0$  et que  $\mathcal{O}_{X_0}(d)$  est un
 $\mathcal{O}_{X_0}$-module quasi-cohérent de type fini. Puisque  $R_0$  est
noethérien (en tant qu'algèbre  $\mathbf{Z}$ de type fini), nous voyons que
 $X_0$  est de présentation finie sur  $R_0$  (Morphismes, Lemme
 \ref{morphisms-lemma-noetherian-finite-type-finite-presentation}) et que  $\mathcal{O}_{X_0}(d)$  est de présentation finie d'après Cohomologie des
schémas, Lemme  \ref{coherent-lemma-coherent-Noetherian}. Cela termine la démonstration.
\end{proof}










\section{Sous-schémas fermés du Proj relatif}
\label{divisors-section-closed-in-relative-proj}

\noindent
Quelques lemmes auxiliaires sur les sous-schémas fermés du Proj relatif.

\begin{lemma}
\label{divisors-lemma-closed-subscheme-proj}
Soit  $S$  un schéma. Soit  $\mathcal{A}$  une  $\mathcal{O}_S$-algèbre graduée
quasi-cohérente. Soit  $p : X = \underline{\text{Proj}}_S(\mathcal{A}) \to S$  le Proj relatif de  $\mathcal{A}$. Soit
 $i : Z \to X$  un sous-schéma fermé. On note  $\mathcal{I} \subset \mathcal{A}$  le noyau de l'application
canonique
$$
\mathcal{A}
\longrightarrow
\bigoplus\nolimits_{d \geq 0} p_*\left((i_*\mathcal{O}_Z)(d)\right).
$$
Si  $p$  est quasi-compact, alors il existe un isomorphisme  $Z = \underline{\text{Proj}}_S(\mathcal{A}/\mathcal{I})$.
\end{lemma}

\begin{proof}
Le morphisme  $p$  est séparé par Constructions, Lemme  \ref{constructions-lemma-relative-proj-separated}. Comme
 $p$  est quasi-compact,  $p_*$  transforme les modules
quasi-cohérents en modules quasi-cohérents, voir Schémas, Lemme  \ref{schemes-lemma-push-forward-quasi-coherent}. Par
conséquent,  $\mathcal{I}$  est un  $\mathcal{O}_S$-module quasi-cohérent. En
particulier,  $\mathcal{B} = \mathcal{A}/\mathcal{I}$  est une  $\mathcal{O}_S$-algèbre graduée quasi-cohérente. Le
morphisme de fonctorialité  $Z' = \underline{\text{Proj}}_S(\mathcal{B}) \to
\underline{\text{Proj}}_S(\mathcal{A})$  est défini partout et est une immersion
fermée, voir Constructions, Lemme  \ref{constructions-lemma-surjective-graded-rings-map-relative-proj}. Il suffit donc de prouver
 $Z = Z'$  en tant que sous-schémas fermés de  $X$.

\medskip\noindent
Ceci étant dit, la question est locale sur la base et nous pouvons supposer que
 $S = \Spec(R)$  et que  $X = \text{Proj}(A)$  pour une certaine  $R$-algèbre graduée
 $A$. Supposons  $\mathcal{I} = \widetilde{I}$  pour  $I \subset A$  un idéal gradué. D'après
Constructions, Lemme  \ref{constructions-lemma-proj-quasi-compact} , il existe  $f_0, \ldots, f_n \in A_{+}$  tels que  $A_{+} \subset \sqrt{(f_0, \ldots, f_n)}$ , en
d'autres termes  $X = \bigcup D_{+}(f_i)$. Par conséquent, il suffit de vérifier que
 $Z \cap D_{+}(f_i) = Z' \cap D_{+}(f_i)$  pour chaque  $i$. En renumérotant, nous pouvons supposer
 $i = 0$. Disons que  $Z \cap D_{+}(f_0)$, resp. \ $Z' \cap D_{+}(f_0)$  est défini
par l'idéal  $J$, resp. \ $J'$  de  $A_{(f_0)}$.

\medskip\noindent
Montrons l'inclusion  $J' \subset J$. Soit  $d$  le plus petit commun multiple de
 $\deg(f_0), \ldots, \deg(f_n)$. Notons que chacun des faisceaux tordus  $\mathcal{O}_X(nd)$  est inversible,
trivialisé par  $f_i^{nd/\deg(f_i)}$  sur  $D_{+}(f_i)$, et que pour tout module
quasi-cohérent  $\mathcal{F}$  sur  $X$ , les applications de multiplication
 $\mathcal{O}_X(nd) \otimes_{\mathcal{O}_X} \mathcal{F}(m)
\to \mathcal{F}(nd + m)$  sont des isomorphismes, voir Constructions, Lemme  \ref{constructions-lemma-when-invertible}.
Observons que  $J'$  est l'idéal engendré par les éléments  $g/f_0^e$  où
 $g \in I$  est homogène de degré  $e\deg(f_0)$  (voir la preuve de Constructions,
Lemme  \ref{constructions-lemma-surjective-graded-rings-map-proj}). Bien sûr, en remplaçant  $g$  par  $f_0^lg$  pour un
 $l$  convenable, nous pouvons toujours supposer que  $d | e$. Ensuite,
puisque  $g$  s'annule en tant que section de  $\mathcal{O}_X(e\deg(f_0))$  restreinte à
 $Z$ , nous voyons que  $g/f_0^d$  est un élément de  $J$. Ainsi
 $J' \subset J$.

\medskip\noindent
Inversement, supposons que  $g/f_0^e \in J$. Encore une fois, nous pouvons supposer
 $d | e$. Choisissons  $i \in \{1, \ldots, n\}$. Alors  $Z \cap D_{+}(f_i)$  est coupé par un certain
idéal  $J_i \subset A_{(f_i)}$. De plus,
$$
J \cdot A_{(f_0f_i)} = J_i \cdot A_{(f_0f_i)}.
$$
Le côté droit est la localisation de  $J_i$  par rapport à  $f_0^{\deg(f_i)}/f_i^{\deg(f_0)}$. Il
s'ensuit que
$$
f_0^{e_i}g/f_i^{(e_i + e)\deg(f_0)/\deg(f_i)} \in J_i
$$
pour un certain  $e_i \gg 0$  suffisamment divisible. Cela prouve que  $f_0^{\max(e_i)}g$ 
est un élément de  $I$, car sa restriction à chaque ouvert affine
 $D_{+}(f_i)$  s'annule sur le sous-schéma fermé  $Z \cap D_{+}(f_i)$. Ainsi  $g/f_0^e \in J'$  et
nous concluons  $J \subset J'$  comme souhaité.
\end{proof}

\begin{example}
\label{divisors-example-closed-subscheme-of-proj}
Soit  $A$  un anneau gradué. Soient  $X = \text{Proj}(A)$  et  $S = \Spec(A_0)$. Étant
donné un idéal gradué  $I \subset A$ , nous obtenons un sous-schéma fermé
 $V_+(I) = \text{Proj}(A/I) \to X$  par le Lemme de Constructions  \ref{constructions-lemma-surjective-graded-rings-map-proj}. En traduisant le résultat
du Lemme  \ref{divisors-lemma-closed-subscheme-proj} , nous voyons que si  $X$  est quasi-compact, alors
tout sous-schéma fermé  $Z$  est de la forme  $V_+(I(Z))$  où l'idéal gradué
 $I(Z) \subset A$  est donné par la règle
$$
I(Z) = \Ker(A \longrightarrow
\bigoplus\nolimits_{n \geq 0} \Gamma(Z, \mathcal{O}_Z(n)))
$$
Nous pouvons alors poser les deux questions naturelles suivantes:
\begin{enumerate}
\item Quels idéaux  $I$  sont de la forme  $I(Z)$?
\item Pouvons-nous décrire l'opération  $I \mapsto I(V_+(I))$?
\end{enumerate}
Nous y répondrons lorsque  $A$  est noethérien.

\medskip\noindent
Tout d'abord, supposons que  $A$  est généré par  $A_1$  sur
 $A_0$. Dans ce cas, pour tout idéal  $I \subset A$ , le noyau de l'application
 $A/I \to \bigoplus \Gamma(\text{Proj}(A/I), \mathcal{O})$  est constitué des éléments de torsion de  $A/I$, voir
Cohomologie des schémas, Proposition  \ref{coherent-proposition-coherent-modules-on-proj}. Par conséquent, nous concluons
que
$$
I(V_+(I)) = \{x \in A \mid A_n x \subset I\text{ pour un certain }n \geq 0\}
$$
L'idéal à droite est parfois appelé la saturation de  $I$. Cela répond à
(2) et la réponse à (1) est qu'un idéal est de la forme  $I(Z)$  si et
seulement s'il est saturé, c'est-à-dire égal à sa propre saturation.

\medskip\noindent
Si  $A$  est un anneau gradué noethérien général, alors nous utilisons
Cohomologie des schémas, Proposition  \ref{coherent-proposition-coherent-modules-on-proj-general}. Ainsi, nous voyons que pour
 $d$  égal au ppcm des degrés des générateurs de  $A$  sur
 $A_0$  nous obtenons
$$
I(V_+(I)) = \{x \in A \mid (Ax)_{nd} \subset I\text{ pour tout }n \gg 0\}
$$
Cela peut être différent de la saturation de  $I$  si  $d \not = 1$. Par
exemple, supposons que  $A = \mathbf{Q}[x, y]$  avec  $\deg(x) = 2$  et  $\deg(y) = 3$. Alors
 $d = 6$. Soit  $I = (y^2)$. Alors nous voyons  $y \in I(V_+(I))$  parce que pour
tout  $f \in A$  homogène tel que  $6 | \deg(fy)$  nous avons  $y | f$, donc
 $fy \in I$. Il s'ensuit que  $I(V_+(I)) = (y)$  mais  $x^n y \not \in I$  pour tout
 $n$  donc  $I(V_+(I))$  n'est pas égal à la saturation.
\end{example}

\begin{lemma}
\label{divisors-lemma-equation-codim-1-in-projective-space}
Soit  $R$  un anneau factoriel noethérien. Soit  $Z \subset \mathbf{P}^n_R$  un sous-schéma fermé qui
n'a aucun point associé immergé et dont chaque composante irréductible
 $Z$  ait une codimension  $1$  dans  $\mathbf{P}^n_R$. Alors l'idéal
 $I(Z) \subset R[T_0, \ldots, T_n]$  correspondant à  $Z$  est principal.
\end{lemma}

\begin{proof}
Remarquez que les anneaux locaux de  $X = \mathbf{P}^n_R$  sont factoriels car  $X$  est
recouvert par des ouverts affines isomorphes à  $\mathbf{A}^n_R$  et  $R[x_1, \ldots, x_n]$  est un
un anneau factoriel (Algèbre, Lemme  \ref{algebra-lemma-polynomial-ring-UFD}). Ainsi  $Z$  est un diviseur de Cartier
effectif d'après le Lemme  \ref{divisors-lemma-codimension-1-is-effective-Cartier}. Soit  $\mathcal{I} \subset \mathcal{O}_X$  le faisceau
quasi-cohérent d'idéaux correspondant à  $Z$. Choisissons un isomorphisme
 $\mathcal{O}(m) \to \mathcal{I}$  pour un certain  $m \in \mathbf{Z}$, voir le Lemme  \ref{divisors-lemma-Pic-projective-space-UFD}. Alors la
composition
$$
\mathcal{O}_X(m) \to \mathcal{I} \to \mathcal{O}_X
$$
est non nulle. Nous concluons que  $m \leq 0$  et que la section correspondante de
 $\mathcal{O}_X(m)^{\otimes -1} = \mathcal{O}_X(-m)$  est donnée par un certain  $F \in R[T_0, \ldots, T_n]$  de degré  $-m$, voir
Cohomologie des schémas, Lemme  \ref{coherent-lemma-cohomology-projective-space-over-ring}. Ainsi, sur l'ouvert standard
 $i$-ième  $U_i = D_+(T_i)$ , le sous-schéma fermé  $Z \cap U_i$  est découpé par
l'idéal
$$
(F(T_0/T_i, \ldots, T_n/T_i)) \subset R[T_0/T_i, \ldots, T_n/T_i]
$$
Ainsi, les éléments homogènes de l'idéal gradué  $I(Z) = \Ker(R[T_0, \ldots, T_n] \to \bigoplus \Gamma(\mathcal{O}_Z(m)))$  sont les
polynômes homogènes  $G$  tels que
$$
G(T_0/T_i, \ldots, T_n/T_i) \in (F(T_0/T_i, \ldots, T_n/T_i))
$$
pour  $i = 0, \ldots, n$. En éliminant les dénominateurs, nous voyons qu'il existe
 $e_i \geq 0$  tel que
$$
T_i^{e_i}G \in (F)
$$
pour  $i = 0, \ldots, n$. Comme  $R$  est factoriel, il en est de même pour
 $R[T_0, \ldots, T_n]$. Alors  $F | T_0^{e_0}G$  et  $F | T_1^{e_1}G$  impliquent  $F | G$  car
 $T_0^{e_0}$  et  $T_1^{e_1}$  n'ont pas de facteur commun. Ainsi  $I(Z) = (F)$.
\end{proof}

\noindent
Dans le cas où le sous-schéma fermé est localement défini par un nombre fini
d'équations, nous pouvons le définir par un faisceau d'idéaux de type fini de
 $\mathcal{A}$.

\begin{lemma}
\label{divisors-lemma-closed-subscheme-proj-finite}
Soit  $S$  un schéma quasi-compact et quasi-séparé. Soit  $\mathcal{A}$  une
 $\mathcal{O}_S$-algèbre graduée quasi-cohérente. Soit  $p : X = \underline{\text{Proj}}_S(\mathcal{A}) \to S$  le Proj relatif de
 $\mathcal{A}$. Soit  $i : Z \to X$  un sous-schéma fermé. Si  $p$  est
quasi-compact et  $i$  est de présentation finie, alors il existe un
 $d > 0$  et un  $\mathcal{O}_S$-submodule quasi-cohérent de type fini  $\mathcal{F} \subset \mathcal{A}_d$ 
tel que  $Z = \underline{\text{Proj}}_S(\mathcal{A}/\mathcal{F}\mathcal{A})$.
\end{lemma}

\begin{proof}
D'après le Lemme \ref{divisors-lemma-closed-subscheme-proj}, il existe un faisceau quasi-cohérent
d'idéaux gradués $\mathcal{I} \subset \mathcal{A}$ tel que $Z = \underline{\text{Proj}}(\mathcal{A}/\mathcal{I})$. Puisque  $S$ est
quasi-compact, nous pouvons choisir un recouvrement ouvert affine
fini $S = U_1 \cup \ldots \cup U_n$. Écrivons  $U_i = \Spec(R_i)$. Faisons correspondre  $\mathcal{A}|_{U_i}$ à
la $R_i$-algèbre graduée $A_i$ et  $\mathcal{I}|_{U_i}$ à l'idéal
gradué $I_i \subset A_i$. Notons que  $p^{-1}(U_i) = \text{Proj}(A_i)$ comme schémas sur $R_i$. Puisque
 $p$ est quasi-compact, nous pouvons choisir un nombre fini d'éléments
homogènes $f_{i, j} \in A_{i, +}$ tels que $p^{-1}(U_i) = \bigcup_j D_{+}(f_{i, j})$. La condition sur  $Z \to X$  signifie
que le faisceau d'idéaux de  $Z$  dans  $\mathcal{O}_X$  est de type fini, voir
Morphismes, Lemme  \ref{morphisms-lemma-closed-immersion-finite-presentation}. Par conséquent, nous pouvons trouver un nombre
fini d'éléments homogènes  $h_{i, j, k} \in I_i \cap A_{i, +}$  tels que l'idéal de  $Z \cap D_{+}(f_{i, j})$  soit
engendré par les éléments  $h_{i, j, k}/f_{i, j}^{e_{i, j, k}}$. Choisissons  $d > 0$  comme un multiple
commun de tous les entiers  $\deg(f_{i, j})$  et  $\deg(h_{i, j, k})$. D'après Propriétés, Lemme
 \ref{properties-lemma-quasi-coherent-colimit-finite-type} , il existe un faisceau quasi-cohérent de type fini  $\mathcal{F} \subset \mathcal{I}_d$  tel
que toutes les sections locales
$$
h_{i, j, k}f_{i, j}^{(d - \deg(h_{i, j, k}))/\deg(f_{i, j})}
$$
soient des sections de  $\mathcal{F}$. Par construction,  $\mathcal{F}$  convient.
\end{proof}

\noindent
La version suivante du Lemme  \ref{divisors-lemma-closed-subscheme-proj-finite}  sera utilisée dans la démonstration du
Lemme  \ref{divisors-lemma-composition-admissible-blowups}.

\begin{lemma}
\label{divisors-lemma-closed-subscheme-proj-finite-type}
Soit  $S$  un schéma quasi-compact et quasi-séparé. Soit  $\mathcal{A}$  une
 $\mathcal{O}_S$-algèbre graduée quasi-cohérente. Soit  $p : X = \underline{\text{Proj}}_S(\mathcal{A}) \to S$  le Proj relatif de
 $\mathcal{A}$. Soit  $i : Z \to X$  un sous-schéma fermé. Soit  $U \subset S$  un ouvert.
Supposons que
\begin{enumerate}
\item $p$  soit quasi-compact,
\item $i$  soit de présentation finie,
\item $U \cap p(i(Z)) = \emptyset$,
\item $U$  soit quasi-compact,
\item $\mathcal{A}_n$  soit un  $\mathcal{O}_S$-module de type fini pour tout  $n$.
\end{enumerate}
Alors il existe un  $d > 0$  et un  $\mathcal{O}_S$-sous-module quasi-cohérent de
type fini  $\mathcal{F} \subset \mathcal{A}_d$  tel que (a)  $Z = \underline{\text{Proj}}_S(\mathcal{A}/\mathcal{F}\mathcal{A})$  et (b) le support de  $\mathcal{A}_d/\mathcal{F}$  soit
disjoint de  $U$.
\end{lemma}

\begin{proof}
Soit  $\mathcal{I} \subset \mathcal{A}$  le faisceau d'idéaux gradués quasi-cohérents construit dans le
Lemme  \ref{divisors-lemma-closed-subscheme-proj}. Soient  $U_i$,  $R_i$,  $A_i$,  $I_i$,
 $f_{i, j}$,  $h_{i, j, k}$ et  $d$  tels que construits dans la
démonstration du Lemme  \ref{divisors-lemma-closed-subscheme-proj-finite}. Puisque  $U \cap p(i(Z)) = \emptyset$ , nous voyons que
 $\mathcal{I}_d|_U = \mathcal{A}_d|_U$  (par notre construction de  $\mathcal{I}$  en tant que noyau). Comme
 $U$  est quasi-compact, nous pouvons choisir un recouvrement ouvert affine
fini  $U = W_1 \cup \ldots \cup W_m$. Puisque  $\mathcal{A}_d$  est de type fini, nous pouvons trouver un
nombre fini de sections  $g_{t, s} \in \mathcal{A}_d(W_t)$  qui engendrent  $\mathcal{A}_d|_{W_t} = \mathcal{I}_d|_{W_t}$  en tant que
 $\mathcal{O}_{W_t}$-module. Pour terminer la preuve, remarquons que, d'après Propriétés,
Lemme  \ref{properties-lemma-quasi-coherent-colimit-finite-type} , il existe un  $\mathcal{F} \subset \mathcal{I}_d$  de type fini tel que toutes les
sections locales
$$
h_{i, j, k}f_{i, j}^{(d - \deg(h_{i, j, k}))/\deg(f_{i, j})}
\quad\text{et}\quad
g_{t, s}
$$
soient des sections de  $\mathcal{F}$. Par construction,  $\mathcal{F}$  convient.
\end{proof}

\begin{lemma}
\label{divisors-lemma-conormal-sheaf-section-projective-bundle}
Soit  $X$  un schéma. Soit  $\mathcal{E}$  un  $\mathcal{O}_X$-module
quasi-cohérent. Il existe une bijection
$$
\left\{
\begin{matrix}
\text{sections }\sigma\text{ du } \\
\text{morphisme } \mathbf{P}(\mathcal{E}) \to X
\end{matrix}
\right\}
\leftrightarrow
\left\{
\begin{matrix}
\text{surjections }\mathcal{E} \to \mathcal{L}\text{ où} \\
\mathcal{L}\text{ est un }\mathcal{O}_X\text{-module inversible}
\end{matrix}
\right\}
$$
Dans ce cas,  $\sigma$  est une immersion fermée et il existe un isomorphisme
canonique
$$
\Ker(\mathcal{E} \to \mathcal{L})
\otimes_{\mathcal{O}_X} \mathcal{L}^{\otimes -1}
\longrightarrow
\mathcal{C}_{\sigma(X)/\mathbf{P}(\mathcal{E})}
$$
La bijection et l'isomorphisme sont compatibles avec le changement de base.
\end{lemma}

\begin{proof}
Rappelons que  $\pi : \mathbf{P}(\mathcal{E}) \to X$  est le Proj relatif de l'algèbre symétrique sur
 $\mathcal{E}$, voir Constructions, Définition  \ref{constructions-definition-projective-bundle}. Par conséquent, les
descriptions des sections  $\sigma$  découlent immédiatement de la description du
foncteur des points de  $\mathbf{P}(\mathcal{E})$  dans Constructions, Lemme  \ref{constructions-lemma-apply-relative}.
Puisque  $\pi$  est séparé, toute section est une immersion fermée
(Constructions, Lemme  \ref{constructions-lemma-relative-proj-separated}  et Schémas, Lemme  \ref{schemes-lemma-section-immersion}). Soit
 $U \subset X$  un ouvert affine et  $k \in \mathcal{E}(U)$  et  $s \in \mathcal{E}(U)$  des sections locales
telles que  $k$  s'envoie sur zéro dans  $\mathcal{L}$  et que  $s$ 
s'envoie sur un générateur  $\overline{s}$  de  $\mathcal{L}$. Alors  $f = k/s$  est
une section de  $\mathcal{O}_{\mathbf{P}(\mathcal{E})}$  définie dans un voisinage ouvert  $D_+(s)$  de
 $s(U)$  dans  $\pi^{-1}(U)$. De plus, puisque  $k$  s'envoie sur zéro
dans  $\mathcal{L}$ , nous voyons que  $f$  est une section du faisceau
d'idéaux de  $s(U)$  dans  $\pi^{-1}(U)$. Ainsi, nous pouvons prendre l'image
 $\overline{f}$  de  $f$  dans  $\mathcal{C}_{\sigma(X)/\mathbf{P}(\mathcal{E})}(U)$. Nous affirmons (1) que l'image
 $\overline{f}$  dépend seulement des sections  $k$  et  $\overline{s}$  et non du
choix de  $s$  et (2) que nous obtenons un isomorphisme sur  $U$  de
cette manière (voir ci-dessous). Cependant, une fois (1) et (2) établis, nous
voyons que la construction est compatible avec le changement de base par
 $U' \to U$  où  $U'$  est affine, ce qui prouve que ces applications
locales se recollent et sont compatibles avec un changement de base arbitraire.

\medskip\noindent
Pour prouver (1) et (2), explicitons la situation. Plus précisément, supposons
que  $U = \Spec(A)$  et que  $\mathcal{E} \to \mathcal{L}$  correspond au morphisme de
 $A$-modules  $M \to N$. Alors  $k \in K = \Ker(M \to N)$  et  $s \in M$  s'envoie
sur un générateur  $\overline{s}$  de  $N$. D'où  $M = K \oplus A s$. Ainsi
$$
\text{Sym}(M) = \text{Sym}(K)[s]
$$
Considérons l'identification  $\text{Sym}(K) \to \text{Sym}(M)_{(s)}$  via la règle  $g \mapsto g/s^n$  pour
 $g \in \text{Sym}^n(K)$. Cela donne un isomorphisme  $D_+(s) = \Spec(\text{Sym}(K))$  tel que  $\sigma$ 
correspond à l'homomorphisme d'anneaux  $\text{Sym}(K) \to A$  qui annule  $K$.
Via cet isomorphisme, nous voyons que le module quasi-cohérent correspondant à
 $K$  est identifié avec  $\mathcal{C}_{\sigma(U)/D_+(s)}$ , prouvant (2). Enfin, supposons que
 $s' = k' + s$  pour un certain  $k' \in K$. Alors
$$
k/s' = (k/s) (s/s') = (k/s) (s'/s)^{-1} = (k/s) (1 + k'/s)^{-1}
$$
dans un voisinage ouvert de  $\sigma(U)$  dans  $D_+(s)$. Ainsi, nous voyons que
 $s'/s$  se restreint à  $1$  sur  $\sigma(U)$  et nous voyons que
 $k/s'$  a la même image que  $k/s$  dans le faisceau conormal,
prouvant ainsi (1).
\end{proof}





\section{Éclatements}
\label{divisors-section-blowing-up}

\noindent
L'éclatement est un outil important en géométrie algébrique.

\begin{definition}
\label{divisors-definition-blow-up}
Soit  $X$  un schéma. Soit  $\mathcal{I} \subset \mathcal{O}_X$  un faisceau quasi-cohérent
d'idéaux, et soit  $Z \subset X$  le sous-schéma fermé correspondant à  $\mathcal{I}$,
voir Schémas, Définition  \ref{schemes-definition-immersion}. L'éclatement  {\it  de
 $X$  de centre  $Z$}, ou l'éclatement
 {\it  de  $X$  relativement au faisceau d'idéaux
 $\mathcal{I}$}  est le morphisme
$$
b :
\underline{\text{Proj}}_X
\left(\bigoplus\nolimits_{n \geq 0} \mathcal{I}^n\right)
\longrightarrow
X
$$
Le  {\it diviseur exceptionnel} de l'éclatement est
l'image réciproque $b^{-1}(Z)$. Parfois  $Z$  est appelé le
 {\it centre}  de l'éclatement.
\end{definition}

\noindent
Nous verrons plus tard que le diviseur exceptionnel est un diviseur de Cartier
effectif. De plus, l'éclatement se caractérise comme le schéma « le plus petit »
au-dessus de $X$ tel que l'image réciproque de $Z$ est un diviseur
de Cartier effectif.

\medskip\noindent
Si  $b : X' \to X$  est l'éclatement de  $X$  de centre  $Z$, alors nous
notons souvent  $\mathcal{O}_{X'}(n)$  les faisceaux tordus du faisceau structural. Notons que ce sont
des  $\mathcal{O}_{X'}$-modules inversibles et que  $\mathcal{O}_{X'}(n) = \mathcal{O}_{X'}(1)^{\otimes n}$  car  $X'$  est le
Proj relatif d'une  $\mathcal{O}_X$-algèbre quasi-cohérente graduée qui est engendrée
en degré  $1$, voir Constructions, Lemme  \ref{constructions-lemma-apply-relative}. Notons que
 $\mathcal{O}_{X'}(1)$  est très ample relativement à  $b$, même si  $b$ 
n'est pas nécessairement de type fini, ni même quasi-compact, car  $X'$  est
muni d'une immersion fermée dans  $\mathbf{P}(\mathcal{I})$, voir Morphismes, Exemple
 \ref{morphisms-example-very-ample}.

\begin{lemma}
\label{divisors-lemma-blowing-up-affine}
Soit  $X$  un schéma. Soit  $\mathcal{I} \subset \mathcal{O}_X$  un faisceau quasi-cohérent
d'idéaux. Soit  $U = \Spec(A)$  un sous-schéma ouvert affine de  $X$  et soit
 $I \subset A$  l'idéal correspondant à  $\mathcal{I}|_U$. Si  $b : X' \to X$  est
l'éclatement de  $X$  relativement à  $\mathcal{I}$, alors il existe un isomorphisme
canonique
$$
b^{-1}(U) = \text{Proj}(\bigoplus\nolimits_{d \geq 0} I^d)
$$
de  $b^{-1}(U)$  avec le spectre homogène de l'algèbre de Rees de  $I$  dans
 $A$. De plus,  $b^{-1}(U)$  possède un recouvrement ouvert affine par des
spectres des algèbres d'éclatement affines  $A[\frac{I}{a}]$.
\end{lemma}

\begin{proof}
La première affirmation découle de la construction du Proj relatif par collage,
voir Constructions, section  \ref{constructions-section-relative-proj-via-glueing}. Pour  $a \in I$ , notons  $a^{(1)}$ 
l'élément  $a$  considéré comme un élément de degré  $1$  dans
l'algèbre de Rees  $\bigoplus_{n \geq 0} I^n$. Puisque ces éléments génèrent l'algèbre de Rees
sur  $A$ , nous voyons que  $\text{Proj}(\bigoplus_{d \geq 0} I^d)$  est couvert par les ouverts affines
 $D_{+}(a^{(1)})$. Le schéma affine  $D_{+}(a^{(1)})$  est le spectre de l'algèbre d'éclatement
affine  $A' = A[\frac{I}{a}]$, voir Algèbre, Définition  \ref{algebra-definition-blow-up}. Cela termine la
preuve.
\end{proof}

\begin{lemma}
\label{divisors-lemma-flat-base-change-blowing-up}
\begin{slogan}
L'éclatement commute avec le changement de base plat.
\end{slogan}
Soit  $X_1 \to X_2$  un morphisme plat de schémas. Soit  $Z_2 \subset X_2$  un sous-schéma
fermé. Soit  $Z_1$  l'image réciproque de  $Z_2$  dans  $X_1$.
Soit  $X'_i$  l'éclatement de centre  $Z_i$  du schéma  $X_i$. Alors il existe un
diagramme cartésien
$$
\xymatrix{
X_1' \ar[r] \ar[d] & X_2' \ar[d] \\
X_1 \ar[r] & X_2
}
$$
de schémas.
\end{lemma}

\begin{proof}
Soit  $\mathcal{I}_2$  le faisceau d'idéaux de  $Z_2$  dans  $X_2$. On note
 $g : X_1 \to X_2$  le morphisme donné. Alors le faisceau d'idéaux  $\mathcal{I}_1$  de
 $Z_1$  est l'image de  $g^*\mathcal{I}_2 \to \mathcal{O}_{X_1}$  (par définition de l'image réciproque,
voir Schémas, Définition  \ref{schemes-definition-inverse-image-closed-subscheme}). D'après Constructions, Lemme  \ref{constructions-lemma-relative-proj-base-change} ,
nous voyons que  $X_1 \times_{X_2} X_2'$  est le Proj relatif de  $\bigoplus_{n \geq 0} g^*\mathcal{I}_2^n$. Comme
 $g$  est plat, l'application  $g^*\mathcal{I}_2^n \to \mathcal{O}_{X_1}$  est injective avec image
 $\mathcal{I}_1^n$. Ainsi nous voyons que  $X_1 \times_{X_2} X_2' = X_1'$.
\end{proof}

\begin{lemma}
\label{divisors-lemma-blowing-up-gives-effective-Cartier-divisor}
Soit  $X$  un schéma. Soit  $Z \subset X$  un sous-schéma fermé. L'éclatement
 $b : X' \to X$  de centre  $Z$  dans  $X$  possède les propriétés suivantes:
\begin{enumerate}
\item $b|_{b^{-1}(X \setminus Z)} : b^{-1}(X \setminus Z) \to X \setminus Z$  est un isomorphisme,
\item le diviseur exceptionnel  $E = b^{-1}(Z)$  est un diviseur de Cartier effectif sur
 $X'$,
\item il existe un isomorphisme canonique  $\mathcal{O}_{X'}(-1) = \mathcal{O}_{X'}(E)$
\end{enumerate}
\end{lemma}

\begin{proof}
Comme l'éclatement commute avec les restrictions aux sous-schémas ouverts (Lemme
 \ref{divisors-lemma-flat-base-change-blowing-up}), la première assertion signifie simplement que  $X' = X$  si
 $Z = \emptyset$. Dans ce cas, nous éclatons dans le faisceau d'idéaux  $\mathcal{I} = \mathcal{O}_X$  et
le résultat découle de Constructions, Exemple  \ref{constructions-example-trivial-proj}.

\medskip\noindent
La deuxième assertion est locale sur  $X$, donc nous pouvons supposer
 $X$  affine. Écrivons  $X = \Spec(A)$  et  $Z = \Spec(A/I)$. Par le Lemme
 \ref{divisors-lemma-blowing-up-affine} , nous voyons que  $X'$  est couvert par les spectres des
algèbres d'éclatement affines  $A' = A[\frac{I}{a}]$. Alors  $IA' = aA'$  et  $a$  s'envoie
sur un élément non-diviseur de zéro dans  $A'$  d'après Algèbre,
Lemme  \ref{algebra-lemma-affine-blowup}. Cela prouve le lemme car l'image réciproque de  $Z$ 
dans  $\Spec(A')$  correspond à  $\Spec(A'/IA') \subset \Spec(A')$.

\medskip\noindent
Considérons l'application canonique  $\psi_{univ, 1} : b^*\mathcal{I} \to \mathcal{O}_{X'}(1)$, voir la discussion suivant
Constructions, Définition  \ref{constructions-definition-relative-proj}. Nous affirmons que ce morphisme se factorise en passant par
un isomorphisme  $\mathcal{I}_E \to \mathcal{O}_{X'}(1)$  (ce qui prouve l'affirmation finale). À
savoir, sur l'ouvert affine correspondant à l'algèbre de l'éclatement  $A' = A[\frac{I}{a}]$ 
mentionnée ci-dessus,  $\psi_{univ, 1}$  correspond au morphisme de
 $A'$-modules
$$
I \otimes_A A'
\longrightarrow
\left(\Big(\bigoplus\nolimits_{d \geq 0} I^d\Big)_{a^{(1)}}\right)_1
$$
où  $a^{(1)}$  est comme dans Algèbre, Définition  \ref{algebra-definition-blow-up}. Nous omettons la
vérification que ceci est l'application  $I \otimes_A A' \to IA' = aA'$.
\end{proof}

\begin{lemma}[Propriété universelle de l'éclatement]
\label{divisors-lemma-universal-property-blowing-up}
Soit  $X$  un schéma. Soit  $Z \subset X$  un sous-schéma fermé. Soit
 $\mathcal{C}$  la sous-catégorie pleine de  $(\Sch/X)$  consistant en  $Y \to X$ 
tels que l'image réciproque de  $Z$  est un diviseur de Cartier effectif sur
 $Y$. Alors l'éclatement  $b : X' \to X$  de centre  $Z$  du schéma  $X$ 
est un objet final de  $\mathcal{C}$.
\end{lemma}

\begin{proof}
On voit que  $b : X' \to X$  est un objet de  $\mathcal{C}$  d'après le Lemme  \ref{divisors-lemma-blowing-up-gives-effective-Cartier-divisor}.
Soit  $f : Y \to X$  un objet de  $\mathcal{C}$. Nous devons montrer qu'il existe un
unique morphisme  $Y \to X'$  sur  $X$. Soit  $D = f^{-1}(Z)$. Soit
 $\mathcal{I} \subset \mathcal{O}_X$  le faisceau d'idéaux de  $Z$  et soit  $\mathcal{I}_D$  le faisceau
d'idéaux de  $D$. Alors  $f^*\mathcal{I} \to \mathcal{I}_D$  est une surjection sur un
 $\mathcal{O}_Y$-module inversible. Cela s'étend à un morphisme  $\psi : \bigoplus f^*\mathcal{I}^d \to \bigoplus \mathcal{I}_D^d$  de
 $\mathcal{O}_Y$-algèbres graduées. (On observe que  $\mathcal{I}_D^d = \mathcal{I}_D^{\otimes d}$  puisque  $D$  est un
diviseur de Cartier effectif.) D'après Constructions, section
 \ref{constructions-section-relative-proj} , le triple  $(1, f : Y \to X, \psi)$  définit un morphisme  $Y \to X'$  sur
 $X$. La restriction
$$
Y \setminus D \longrightarrow X' \setminus b^{-1}(Z) = X \setminus Z
$$
est unique. L'ouvert  $Y \setminus D$  est schématiquement dense dans  $Y$ 
d'après le Lemme  \ref{divisors-lemma-complement-effective-Cartier-divisor}. Ainsi, le morphisme  $Y \to X'$  est unique d'après
Morphismes, Lemme  \ref{morphisms-lemma-equality-of-morphisms}  (de même  $b$  est séparé par Constructions,
Lemme  \ref{constructions-lemma-relative-proj-separated}).
\end{proof}

\begin{lemma}
\label{divisors-lemma-characterize-affine-blowup}
Soit  $b : X' \to X$  l'éclatement du schéma  $X$  de centre un sous-schéma fermé
 $Z$. Soit  $U = \Spec(A)$  un ouvert affine de  $X$  et soit
 $I \subset A$  l'idéal correspondant à  $Z \cap U$. Soit  $a \in I$  et soit
 $x' \in X'$  un point se projetant sur un point de  $U$. Alors
 $x'$  est un point de l'ouvert affine  $U' = \Spec(A[\frac{I}{a}])$  si et seulement si
l'image de  $a$  dans  $\mathcal{O}_{X', x'}$  définit le diviseur exceptionnel.
\end{lemma}

\begin{proof}
Puisque le diviseur exceptionnel au-dessus de  $U'$  est défini par l'image
de  $a$  dans  $A' = A[\frac{I}{a}]$ , une implication est claire. Inversement,
supposons que l'image de  $a$  dans  $\mathcal{O}_{X', x'}$  définisse  $E$.
Puisque chaque élément de  $I$  s'envoie à un élément de l'idéal
définissant  $E$  sur  $b^{-1}(U)$ , nous voyons que les éléments de
 $I$  deviennent divisibles par  $a$  dans  $\mathcal{O}_{X', x'}$. Ainsi, pour
 $f \in I^n$ , nous pouvons écrire  $f = \psi(f) a^n$  pour un certain  $\psi(f) \in \mathcal{O}_{X', x'}$.
Comme  $a$  s'envoie sur un non-diviseur de zéro de
 $\mathcal{O}_{X', x'}$ , cette relation caractérise l'élément  $\psi(f)$  de manière unique.
On définit alors
$$
A' \longrightarrow \mathcal{O}_{X', x'},\quad
f/a^n \longmapsto \psi(f)
$$
Nous utilisons ici la description des algèbres d'éclatement donnée après
Algèbre, Définition  \ref{divisors-definition-blow-up}. L'unicité mentionnée ci-dessus montre qu'il
s'agit d'un homomorphisme de  $A$-algèbres. Cela donne un morphisme
 $\Spec(\mathcal{O}_{X', x"}) \to \Spec(A') = U'$. Par la propriété universelle de l'éclatement (Lemme  \ref{divisors-lemma-universal-property-blowing-up}), il
s'agit d'un morphisme au-dessus de  $X'$, ce qui implique bien sûr que
 $x' \in U'$.
\end{proof}

\begin{lemma}
\label{divisors-lemma-blow-up-effective-Cartier-divisor}
Soit  $X$  un schéma. Soit  $Z \subset X$  un diviseur de Cartier effectif. L'
éclatement de  $X$  de centre  $Z$  est le morphisme identité de
 $X$.
\end{lemma}

\begin{proof}
Immédiat d'après la propriété universelle des éclatements (Lemme  \ref{divisors-lemma-universal-property-blowing-up}).
\end{proof}

\begin{lemma}
\label{divisors-lemma-blow-up-reduced-scheme}
Soit  $X$  un schéma. Soit  $\mathcal{I} \subset \mathcal{O}_X$  un faisceau quasi-cohérent
d'idéaux. Si  $X$  est réduit, alors l'éclatement  $X'$  de
 $X$  dans  $\mathcal{I}$  est réduit.
\end{lemma}

\begin{proof}
Il suffit de combiner le Lemme  \ref{divisors-lemma-blowing-up-affine}  avec Algèbre, Lemme  \ref{algebra-lemma-blowup-reduced}.
\end{proof}

\begin{lemma}
\label{divisors-lemma-blow-up-integral-scheme}
Soit  $X$  un schéma. Soit  $\mathcal{I} \subset \mathcal{O}_X$  un faisceau quasi-cohérent d'idéaux
non nul. Si  $X$  est intègre, alors l'éclatement  $X'$  de  $X$ 
relativement à  $\mathcal{I}$  est intègre.
\end{lemma}

\begin{proof}
Il suffit de combiner le Lemme  \ref{divisors-lemma-blowing-up-affine}  avec Algèbre, Lemme  \ref{algebra-lemma-blowup-domain}.
\end{proof}

\begin{lemma}
\label{divisors-lemma-blow-up-and-irreducible-components}
Soit  $X$  un schéma. Soit  $Z \subset X$  un sous-schéma fermé. Soit
 $b : X' \to X$  l'éclatement de  $X$  de centre  $Z$. Alors  $b$ 
induit une application bijective de l'ensemble des points génériques des
composantes irréductibles de  $X'$  vers l'ensemble des points génériques
des composantes irréductibles de  $X$  qui ne sont pas dans  $Z$.
\end{lemma}

\begin{proof}
Le diviseur exceptionnel  $E \subset X'$  est un diviseur de Cartier effectif et
 $X' \setminus E \to X \setminus Z$  est un isomorphisme, voir le Lemme  \ref{divisors-lemma-blowing-up-gives-effective-Cartier-divisor}. Il suffit donc de
montrer ce qui suit: étant donné un diviseur de Cartier effectif  $D \subset S$  d'un
schéma  $S$ , aucun des points génériques des composantes irréductibles de
 $S$  n'est contenu dans  $D$. Pour voir ceci, nous pouvons
remplacer  $S$  par les membres d'un recouvrement ouvert affine. Ainsi,
d'après le Lemme  \ref{divisors-lemma-characterize-effective-Cartier-divisor} , nous pouvons supposer  $S = \Spec(A)$  et  $D = V(f)$ 
où  $f \in A$  est un non-diviseur de zéro. Ensuite, nous devons montrer que
 $f$  n'est contenu dans aucun idéal premier minimal  $\mathfrak p \subset A$. Si
c'était le cas, alors  $f$  s'enverrait sur un non-diviseur de zéro
contenu dans l'idéal maximal de  $R_\mathfrak p$ , ce qui est en contradiction avec
Algèbre, Lemme  \ref{algebra-lemma-minimal-prime-reduced-ring}.
\end{proof}

\begin{lemma}
\label{divisors-lemma-blow-up-pullback-effective-Cartier}
Soit  $X$  un schéma. Soit  $b : X' \to X$  un éclatement de  $X$  dont le centre est
un sous-schéma fermé. L'image réciproque  $b^{-1}D$  est définie pour tous les
diviseurs de Cartier effectifs  $D \subset X$  et les images réciproques de fonctions
méromorphes sont définies pour  $b$  (Définitions  \ref{divisors-definition-pullback-effective-Cartier-divisor}  et
 \ref{divisors-definition-pullback-meromorphic-sections}).
\end{lemma}

\begin{proof}
D'après les lemmes  \ref{divisors-lemma-blowing-up-affine}  et  \ref{divisors-lemma-characterize-effective-Cartier-divisor} , cela se réduit au fait algébrique
suivant : soit  $A$  un anneau,  $I \subset A$  un idéal,  $a \in I$, et
 $x \in A$  un non-diviseur de zéro. Alors l'image de  $x$  dans
 $A[\frac{I}{a}]$  est un non-diviseur de zéro. Supposons en effet que  $x (y/a^n) = 0$ 
dans  $A[\frac{I}{a}]$. Alors  $a^mxy = 0$  dans  $A$  pour un certain
 $m$. Par conséquent  $a^my = 0$  car  $x$  est un non-diviseur de
zéro. Par conséquent,  $y/a^n$  est nul dans  $A[\frac{I}{a}]$, comme voulu.
\end{proof}

\begin{lemma}
\label{divisors-lemma-blowing-up-two-ideals}
Soit  $X$  un schéma. Soient  $\mathcal{I}, \mathcal{J} \subset \mathcal{O}_X$  des faisceaux quasi-cohérents
d'idéaux. Soit  $b : X' \to X$  l'éclatement de  $X$  relativement à  $\mathcal{I}$. Soit
 $b' : X'' \to X'$  l'éclatement de  $X'$  relativement à  $b^{-1}\mathcal{J} \mathcal{O}_{X'}$. Alors  $X'' \to X$ 
est canoniquement isomorphe à l'éclatement de  $X$  relativement à  $\mathcal{I}\mathcal{J}$.
\end{lemma}

\begin{proof}
Soit  $E \subset X'$  le diviseur exceptionnel de  $b$  qui est un diviseur de
Cartier effectif d'après le Lemme  \ref{divisors-lemma-blowing-up-gives-effective-Cartier-divisor}. Alors  $(b')^{-1}E$  est un diviseur
de Cartier effectif sur  $X''$  d'après le Lemme  \ref{divisors-lemma-blow-up-pullback-effective-Cartier}. Soit
 $E' \subset X''$  le diviseur exceptionnel de  $b'$  (également un diviseur de
Cartier effectif). Considérons le diviseur de Cartier effectif  $E'' = E' + (b')^{-1}E$. Par
construction, l'idéal de  $E''$  est  $(b \circ b')^{-1}\mathcal{I} (b \circ b')^{-1}\mathcal{J} \mathcal{O}_{X''}$. Par conséquent, d'après le
Lemme  \ref{divisors-lemma-universal-property-blowing-up} , il existe un morphisme canonique de  $X''$  vers
l'éclatement  $c : Y \to X$  de  $X$  relativement à  $\mathcal{I}\mathcal{J}$. Inversement, comme l'image réciproque de
 $\mathcal{I}\mathcal{J}$  est un idéal inversible, nous voyons que  $c^{-1}\mathcal{I}\mathcal{O}_Y$ 
définit un diviseur de Cartier effectif, voir le Lemme  \ref{divisors-lemma-sum-closed-subschemes-effective-Cartier}. Il existe donc un
morphisme  $c' : Y \to X'$  sur  $X$  par le Lemme  \ref{divisors-lemma-universal-property-blowing-up}. Puis
 $(c')^{-1}b^{-1}\mathcal{J}\mathcal{O}_Y = c^{-1}\mathcal{J}\mathcal{O}_Y$, qui définit également un diviseur de Cartier effectif. Il existe donc un
morphisme  $c'' : Y \to X''$  sur  $X'$. Nous omettons la vérification que ce
morphisme est l'inverse du morphisme  $X'' \to Y$  construit précédemment.
\end{proof}

\begin{lemma}
\label{divisors-lemma-blowing-up-projective}
Soit  $X$  un schéma. Soit  $\mathcal{I} \subset \mathcal{O}_X$  un faisceau quasi-cohérent
d'idéaux. Soit  $b : X' \to X$  l'éclatement de  $X$  relativement au faisceau d'idéaux
 $\mathcal{I}$. Si  $\mathcal{I}$  est de type fini, alors
\begin{enumerate}
\item $b : X' \to X$  est un morphisme projectif, et
\item $\mathcal{O}_{X'}(1)$  est un faisceau inversible ample relativement à  $b$.
\end{enumerate}
\end{lemma}

\begin{proof}
La surjection des  $\mathcal{O}_X$-algèbres graduées
$$
\text{Sym}_{\mathcal{O}_X}^*(\mathcal{I})
\longrightarrow
\bigoplus\nolimits_{d \geq 0} \mathcal{I}^d
$$
définit via Constructions, Lemme  \ref{constructions-lemma-surjective-generated-degree-1-map-relative-proj}  une immersion fermée
$$
X' = \underline{\text{Proj}}_X (\bigoplus\nolimits_{d \geq 0} \mathcal{I}^d)
\longrightarrow
\mathbf{P}(\mathcal{I}).
$$
Par conséquent,  $b$  est projectif, voir Morphismes, Définition
 \ref{morphisms-definition-projective}. La deuxième affirmation découle par exemple de la caractérisation
des faisceaux inversibles relativement amples dans Morphismes, Lemme  \ref{morphisms-lemma-characterize-relatively-ample}.
Certains détails sont omis.
\end{proof}

\begin{lemma}
\label{divisors-lemma-composition-finite-type-blowups}
\begin{slogan}
La composition d'éclatements est un éclatement.
\end{slogan}
Soit  $X$  un schéma quasi-compact et quasi-séparé. Soit  $Z \subset X$  un
sous-schéma fermé de présentation finie. Soit  $b : X' \to X$  l'éclatement de centre
 $Z$. Soit  $Z' \subset X'$  un sous-schéma fermé de présentation finie. Soit
 $X'' \to X'$  l'éclatement de centre  $Z'$. Il existe un sous-schéma fermé
 $Y \subset X$  de présentation finie, tel que
\begin{enumerate}
\item $Y = Z \cup b(Z')$  ensemblistement, et
\item la composition  $X'' \to X$  est isomorphe à l'éclatement de  $X$  de centre
 $Y$.
\end{enumerate}
\end{lemma}

\begin{proof}
La condition que  $Z \to X$  soit de présentation finie signifie que  $Z$ 
est découpé par un faisceau quasi-cohérent d'idéaux de type fini  $\mathcal{I} \subset \mathcal{O}_X$,
voir Morphismes, Lemme  \ref{morphisms-lemma-closed-immersion-finite-presentation}. Écrivons  $\mathcal{A} = \bigoplus_{n \geq 0} \mathcal{I}^n$  de sorte que
 $X' = \underline{\text{Proj}}(\mathcal{A})$. Remarquons que  $X \setminus Z$  est un ouvert quasi-compact de
 $X$  d'après Propriétés, Lemme  \ref{properties-lemma-quasi-coherent-finite-type-ideals}. Puisque  $b^{-1}(X \setminus Z) \to X \setminus Z$  est un
isomorphisme (Lemme  \ref{divisors-lemma-blowing-up-gives-effective-Cartier-divisor}), le même résultat montre que  $b^{-1}(X \setminus Z) \setminus Z'$  est un
ouvert quasi-compact de  $X'$. Ainsi,  $U = X \setminus (Z \cup b(Z'))$  est un ouvert
quasi-compact de  $X$. D'après le Lemme  \ref{divisors-lemma-closed-subscheme-proj-finite-type} , il existe un
 $d > 0$  et un sous-$\mathcal{O}_X$-module de type fini  $\mathcal{F} \subset \mathcal{I}^d$  tel que
 $Z' = \underline{\text{Proj}}(\mathcal{A}/\mathcal{F}\mathcal{A})$  et tel que le support de  $\mathcal{I}^d/\mathcal{F}$  soit contenu dans  $X \setminus U$.

\medskip\noindent
Puisque  $\mathcal{F} \subset \mathcal{I}^d$  est un sous-$\mathcal{O}_X$-module, nous pouvons considérer
 $\mathcal{F} \subset \mathcal{I}^d \subset \mathcal{O}_X$  comme un faisceau quasi-cohérent d'idéaux de type fini sur
 $X$. Notons-le  $\mathcal{J} \subset \mathcal{O}_X$  pour éviter toute confusion. Puisque
 $\mathcal{I}^d / \mathcal{J}$  et  $\mathcal{O}/\mathcal{I}^d$  ont leur support dans  $X \setminus U$ , nous voyons que
 $V(\mathcal{J})$  est contenu dans  $X \setminus U$. Inversement, comme  $\mathcal{J} \subset \mathcal{I}^d$ ,
nous voyons que  $Z \subset V(\mathcal{J})$. Sur  $X \setminus Z \cong X' \setminus b^{-1}(Z)$ , le faisceau d'idéaux  $\mathcal{J}$ 
définit  $Z'$  (voir la formule affichée ci-dessous). Par conséquent,
 $V(\mathcal{J})$  est égal à  $Z \cup b(Z')$. Il s'ensuit que  $V(\mathcal{I}\mathcal{J}) = Z \cup b(Z')$  également,
ensemblistement. De plus,  $\mathcal{I}\mathcal{J}$  est un idéal de type fini puisqu'il
est le produit de deux idéaux de type fini. Nous affirmons que  $X'' \to X$  est isomorphe à
l'éclatement de  $X$  relativement à  $\mathcal{I}\mathcal{J}$ , ce qui termine la démonstration du
lemme en posant  $Y = V(\mathcal{I}\mathcal{J})$.

\medskip\noindent
Tout d'abord, rappelons que l'éclatement de  $X$  relativement à  $\mathcal{I}\mathcal{J}$  est le
même que l'éclatement de  $X'$  relativement à  $b^{-1}\mathcal{J} \mathcal{O}_{X'}$, voir le Lemme
 \ref{divisors-lemma-blowing-up-two-ideals}. Il suffit donc de montrer que l'éclatement de  $X'$  relativement à
 $b^{-1}\mathcal{J} \mathcal{O}_{X'}$  coïncide avec l'éclatement de  $X'$  de centre  $Z'$. Nous
allons montrer que
$$
b^{-1}\mathcal{J} \mathcal{O}_{X'} = \mathcal{I}_E^d \mathcal{I}_{Z'}
$$
en tant que faisceaux d'idéaux sur  $X''$. Cela prouvera ce que nous voulons
car  $\mathcal{I}_E^d$  découpe le diviseur de Cartier effectif  $dE$  et nous
pouvons utiliser les Lemmes  \ref{divisors-lemma-blow-up-effective-Cartier-divisor}  et  \ref{divisors-lemma-blowing-up-two-ideals}.

\medskip\noindent
Pour voir l'égalité affichée des idéaux, nous pouvons travailler localement. Avec
la notation  $A$,  $I$,  $a \in I$  comme dans le Lemme
 \ref{divisors-lemma-blowing-up-affine} , nous voyons que  $\mathcal{F}$  correspond à un
 $R$-sous-module  $M \subset I^d$  s'envoyant isomorphiquement sur un idéal
 $J \subset R$. La condition  $Z' = \underline{\text{Proj}}(\mathcal{A}/\mathcal{F}\mathcal{A})$  signifie que  $Z' \cap \Spec(A[\frac{I}{a}])$  est découpé par
l'idéal généré par les éléments  $m/a^d$,  $m \in M$. Écrivons que l'élément
 $m \in M$  correspond à la fonction  $f \in J$. Alors, dans l'algèbre affine
de l'éclatement  $A' = A[\frac{I}{a}]$ , nous voyons que  $f = (a^dm)/a^d = a^d (m/a^d)$. Ainsi, l'égalité
est vérifiée.
\end{proof}






\section{Transformée stricte}
\label{divisors-section-strict-transform}

\noindent
Dans cette section, nous discutons brièvement de la transformée stricte lors d'un
éclatement. Soit  $S$  un schéma et soit  $Z \subset S$  un sous-schéma fermé.
Soit  $b : S' \to S$  l'éclatement de  $S$  de centre  $Z$  et notons
 $E \subset S'$  le diviseur exceptionnel  $E = b^{-1}Z$. Dans ce qui suit, nous
considérerons souvent un schéma  $X$  sur  $S$  et formerons le
diagramme cartésien
$$
\xymatrix{
\text{pr}_{S'}^{-1}E \ar[r] \ar[d] &
X \times_S S' \ar[r]_-{\text{pr}_X} \ar[d]_{\text{pr}_{S'}} &
X \ar[d]^f \\
E \ar[r] & S' \ar[r] & S
}
$$
Puisque  $E$  est un diviseur de Cartier effectif (Lemme  \ref{divisors-lemma-blowing-up-gives-effective-Cartier-divisor}),
nous voyons que  $\text{pr}_{S'}^{-1}E \subset X \times_S S'$  est localement principal (Lemme  \ref{divisors-lemma-pullback-locally-principal}). Ainsi,
le complément de  $\text{pr}_{S'}^{-1}E$  dans  $X \times_S S'$  est rétrocompact (Lemme
 \ref{divisors-lemma-complement-locally-principal-closed-subscheme}). Par conséquent, pour un  $\mathcal{O}_{X \times_S S'}$-module quasi-cohérent
 $\mathcal{G}$ , le sous-faisceau des sections à support dans  $\text{pr}_{S'}^{-1}E$  est un
sous-module quasi-cohérent, voir Propriétés, Lemme  \ref{properties-lemma-sections-supported-on-closed-subset}. Si  $\mathcal{G}$ 
est un faisceau quasi-cohérent d'algèbres, par exemple  $\mathcal{G} = \mathcal{O}_{X \times_S S'}$, alors ce
sous-faisceau est un idéal de  $\mathcal{G}$.

\begin{definition}
\label{divisors-definition-strict-transform}
Avec  $Z \subset S$  et  $f : X \to S$  comme ci-dessus.
\begin{enumerate}
\item Étant donné un  $\mathcal{O}_X$-module quasi-cohérent  $\mathcal{F}$ , la
 {\it  transformée stricte}  de  $\mathcal{F}$  par
rapport à l'éclatement de  $S$  de centre  $Z$  est le quotient  $\mathcal{F}'$
de  $\text{pr}_X^*\mathcal{F}$  par le sous-module des sections à support dans  $\text{pr}_{S'}^{-1}E$.
\item La  {\it transformée stricte}  de  $X$ est
le sous-schéma fermé $X' \subset X \times_S S'$  découpé par l'idéal quasi-cohérent de sections de
 $\mathcal{O}_{X \times_S S'}$ à support dans $\text{pr}_{S'}^{-1}E$.
\end{enumerate}
\end{definition}

\noindent
Notons que prendre la transformée stricte par un éclatement dépend du
sous-schéma fermé utilisé pour l'éclatement (et pas seulement du
morphisme $S' \to S$). Cette notion est souvent utilisée pour les sous-schémas
fermés de  $S$. Il s'avère que la transformée stricte de  $X$  est
un éclatement de  $X$.

\begin{lemma}
\label{divisors-lemma-strict-transform}
Dans la situation de la Définition  \ref{divisors-definition-strict-transform}.
\begin{enumerate}
\item La transformée stricte $X'$  de  $X$  est l'éclatement de
 $X$ dont le centre est le sous-schéma fermé $f^{-1}Z$  de  $X$.
\item Pour un  $\mathcal{O}_X$-module quasi-cohérent  $\mathcal{F}$ , la transformée stricte
 $\mathcal{F}'$  est canoniquement isomorphe à l'image directe par  $X' \to X \times_S S'$ 
de la transformée stricte de  $\mathcal{F}$  relative à l'éclatement  $X' \to X$.
\end{enumerate}
\end{lemma}

\begin{proof}
Soit  $X'' \to X$  l'éclatement de  $X$  de centre  $f^{-1}Z$. Par la propriété
universelle de l'éclatement (Lemme  \ref{divisors-lemma-universal-property-blowing-up}) il existe un diagramme commutatif
$$
\xymatrix{
X'' \ar[r] \ar[d] & X \ar[d] \\
S' \ar[r] & S
}
$$
d'où un morphisme  $X'' \to X \times_S S'$. Ainsi, la première affirmation est que ce
morphisme est une immersion fermée avec image  $X'$. La question est locale
sur  $X$. Ainsi, nous pouvons supposer que  $X$  et  $S$ 
sont affines. Écrivons  $S = \Spec(A)$,  $X = \Spec(B)$ et supposons que  $Z$  soit défini
par l'idéal  $I \subset A$. Posons  $J = IB$. L'application  $B \otimes_A \bigoplus_{n \geq 0} I^n \to \bigoplus_{n \geq 0} J^n$  définit
une immersion fermée  $X'' \to X \times_S S'$, voir Constructions, lemmes  \ref{constructions-lemma-base-change-map-proj}  et
 \ref{constructions-lemma-surjective-graded-rings-generated-degree-1-map-proj}. Nous omettons la vérification que ce morphisme est le même que celui
construit ci-dessus à partir de la propriété universelle. Choisissons  $a \in I$ 
correspondant à l'ouvert affine  $\Spec(A[\frac{I}{a}]) \subset S'$, voir le Lemme  \ref{divisors-lemma-blowing-up-affine}. L'image
réciproque de  $\Spec(A[\frac{I}{a}])$  dans la transformée stricte  $X'$  de
 $X$  est le spectre de
$$
B' = (B \otimes_A A[\textstyle{\frac{I}{a}}])/a\text{-torsion par puissances}
$$
voir Propriétés, Lemme  \ref{properties-lemma-sections-supported-on-closed-subset}. D'autre part, si  $b \in J$  désigne
l'image de  $a$ , nous voyons que  $\Spec(B[\frac{J}{b}])$  est l'image réciproque de
 $\Spec(A[\frac{I}{a}])$  dans  $X''$. Par Algèbre, Lemme  \ref{algebra-lemma-blowup-base-change} , l'ouvert
 $\Spec(B[\frac{J}{b}])$  s'envoie isomorphiquement sur le sous-schéma ouvert  $\text{pr}_{S'}^{-1}(\Spec(A[\frac{I}{a}]))$  de
 $X'$. Ainsi  $X'' \to X'$  est un isomorphisme.

\medskip\noindent
Dans la notation ci-dessus, supposons que  $\mathcal{F}$  corresponde au  $B$-module
 $N$. La transformée stricte de  $\mathcal{F}$  correspond au
 $B \otimes_A A[\frac{I}{a}]$-module
$$
N' = (N \otimes_A A[\textstyle{\frac{I}{a}}])/a\text{-torsion par puissances}
$$
voir Propriétés, Lemme  \ref{properties-lemma-sections-supported-on-closed-subset}. La transformée stricte de  $\mathcal{F}$ 
relativement à l'éclatement de  $X$  de centre  $f^{-1}Z$  correspond au
 $B[\frac{J}{b}]$-module  $N \otimes_B B[\frac{J}{b}]/b\text{-torsion par puissances}$. De la même manière que ci-dessus, on
démontre que ces deux modules sont isomorphes. Les détails sont omis.
\end{proof}

\begin{lemma}
\label{divisors-lemma-strict-transform-flat}
Dans la situation de la Définition  \ref{divisors-definition-strict-transform}.
\begin{enumerate}
\item Si  $X$  est plat sur  $S$  en tous les points au-dessus de
 $Z$, alors la transformée stricte de  $X$  est égale au changement
de base  $X \times_S S'$.
\item Soit  $\mathcal{F}$  un  $\mathcal{O}_X$-module quasi-cohérent. Si  $\mathcal{F}$  est plat
sur  $S$  en tous les points au-dessus de  $Z$, alors la
transformée stricte  $\mathcal{F}'$  de  $\mathcal{F}$  est égale à l'image réciproque
 $\text{pr}_X^*\mathcal{F}$.
\end{enumerate}
\end{lemma}

\begin{proof}
Nous prouverons la partie (2) car elle implique la partie (1) par la définition de
la transformée stricte d'un schéma sur  $S$. La question est locale sur
 $X$. Ainsi, nous pouvons supposer que  $S = \Spec(A)$,  $X = \Spec(B)$, et que
 $\mathcal{F}$  correspond au  $B$-module  $N$. Alors  $\mathcal{F}'$ 
sur l'ouvert  $\Spec(B \otimes_A A[\frac{I}{a}])$  de  $X \times_S S'$  correspond au module
$$
N' = (N \otimes_A A[\textstyle{\frac{I}{a}}])/a\text{-torsion par puissances}
$$
voir Propriétés, Lemme  \ref{properties-lemma-sections-supported-on-closed-subset}. Ainsi, nous devons montrer que la
torsion annulée par une puissance de  $a$  dans  $N \otimes_A A[\frac{I}{a}]$  est nulle. Soit  $y \in N \otimes_A A[\frac{I}{a}]$  avec
 $a^n y = 0$. Si  $\mathfrak q \subset B$  est un idéal premier et  $a \not \in \mathfrak q$, alors
 $y$  s'envoie à zéro dans  $(N \otimes_A A[\frac{I}{a}])_\mathfrak q$. D'autre part, si  $a \in \mathfrak q$,
alors  $N_\mathfrak q$  est un  $A$-module plat et nous voyons que  $N_\mathfrak q \otimes_A A[\frac{I}{a}]
=(N \otimes_A A[\frac{I}{a}])_\mathfrak q$ 
n'a pas de torsion annulée par une puissance de  $a$ (car  $A[\frac{I}{a}]$  n'en a pas). Par
conséquent,  $y$  s'envoie également à zéro dans cette localisation. Nous
concluons que  $y$  est nul d'après Algèbre, Lemme  \ref{algebra-lemma-characterize-zero-local}.
\end{proof}

\begin{lemma}
\label{divisors-lemma-strict-transform-affine}
Soit  $S$  un schéma. Soit  $Z \subset S$  un sous-schéma fermé. Soit
 $b : S' \to S$  l'éclatement de centre  $Z$  du schéma  $S$. Soit  $g : X \to Y$  un
morphisme affine de schémas sur  $S$. Soit  $\mathcal{F}$  un faisceau
quasi-cohérent sur  $X$. Soit  $g' : X \times_S S' \to Y \times_S S'$  le changement de base de
 $g$. Soit  $\mathcal{F}'$  la transformée stricte de  $\mathcal{F}$  par rapport
à  $b$. Alors  $g'_*\mathcal{F}'$  est la transformée stricte de  $g_*\mathcal{F}$.
\end{lemma}

\begin{proof}
Observons que  $g'_*\text{pr}_X^*\mathcal{F} = \text{pr}_Y^*g_*\mathcal{F}$  par Cohomologie des schémas, Lemme  \ref{coherent-lemma-affine-base-change}. Soit
 $\mathcal{K} \subset \text{pr}_X^*\mathcal{F}$  le sous-faisceau des sections à support dans l'image réciproque de
 $Z$  dans  $X \times_S S'$. D'après Propriétés, Lemme  \ref{properties-lemma-push-sections-supported-on-closed-subset} , l'
image directe  $g'_*\mathcal{K}$  est le sous-faisceau des sections de  $\text{pr}_Y^*g_*\mathcal{F}$ 
à support dans l'image réciproque de  $Z$  dans  $Y \times_S S'$. Comme
 $g'$  est affine (Morphismes, Lemme  \ref{morphisms-lemma-base-change-affine}), nous voyons que
 $g'_*$  est exact, donc nous concluons.
\end{proof}

\begin{lemma}
\label{divisors-lemma-strict-transform-different-centers}
Soit  $S$  un schéma. Soit  $Z \subset S$  un sous-schéma fermé. Soit
 $D \subset S$  un diviseur de Cartier effectif. Soit  $Z' \subset S$  le sous-schéma
fermé défini par le produit des faisceaux d'idéaux de  $Z$  et de
 $D$. Soit  $S' \to S$  l'éclatement de  $S$  de centre  $Z$.
\begin{enumerate}
\item L'éclatement de  $S$  de centre  $Z'$  est isomorphe à  $S' \to S$.
\item Soit  $f : X \to S$  un morphisme de schémas et soit  $\mathcal{F}$  un
 $\mathcal{O}_X$-module quasi-cohérent. Si  $\mathcal{F}$  n'a pas de sections locales
non nulles à support dans  $f^{-1}D$, alors la transformée stricte de
 $\mathcal{F}$  par rapport à l'éclatement de centre  $Z$  coïncide avec la
transformée stricte de  $\mathcal{F}$  par rapport à l'éclatement de  $S$  de centre
 $Z'$.
\end{enumerate}
\end{lemma}

\begin{proof}
La première affirmation découle de la combinaison des lemmes  \ref{divisors-lemma-blowing-up-two-ideals}  et
 \ref{divisors-lemma-blow-up-effective-Cartier-divisor}. En utilisant le lemme  \ref{divisors-lemma-blowing-up-affine} , la deuxième affirmation se
traduit par le problème algébrique suivant. Soit  $A$  un anneau,
 $I \subset A$  un idéal,  $x \in A$  un élément non diviseur de zéro, et
 $a \in I$. Soit  $M$  un  $A$-module sans  $x$-torsion. Il faut montrer que la torsion
annulée par une puissance de  $a$ dans  $M \otimes_A A[\frac{I}{a}]$  est égale à la torsion
annulée par une puissance de  $xa$. La raison en est que le noyau et le conoyau de
l'application  $A \to A[\frac{I}{a}]$  sont annulés par une puissance de  $a$, donc cette application devient un isomorphisme après
inversion de  $a$. Ainsi, le noyau et le conoyau de  $M \to M \otimes_A A[\frac{I}{a}]$  sont
également annulés par une puissance de  $a$. Cela implique le résultat.
\end{proof}

\begin{lemma}
\label{divisors-lemma-strict-transform-composition-blowups}
Soit  $S$  un schéma. Soit  $Z \subset S$  un sous-schéma fermé. Soit
 $b : S' \to S$  l'éclatement de centre  $Z$. Soit  $Z' \subset S'$  un
sous-schéma fermé. Soit  $S'' \to S'$  l'éclatement de centre  $Z'$.
Soit  $Y \subset S$  un sous-schéma fermé tel que  $Y = Z \cup b(Z')$
ensemblistement et que la composition  $S'' \to S$  soit isomorphe à l'éclatement
de  $S$  de centre  $Y$. Dans cette situation, pour tout schéma
 $X$  sur  $S$  et  $\mathcal{F} \in \QCoh(\mathcal{O}_X)$ , nous avons
\begin{enumerate}
\item la transformée stricte de  $\mathcal{F}$  par rapport à l'éclatement de
 $S$  de centre  $Y$  est égale à la transformée stricte, par rapport à
l'éclatement  $S'' \to S'$  de centre  $Z'$, de la transformée stricte de
 $\mathcal{F}$  par rapport à l'éclatement  $S' \to S$  de  $S$  de centre
 $Z$, et
\item la transformée stricte de  $X$  par rapport à l'éclatement de  $S$  de centre
 $Y$  est égale à la transformée stricte par rapport à l'éclatement
 $S'' \to S'$  de centre  $Z'$  de la transformée stricte de  $X$  par
rapport à l'éclatement  $S' \to S$  de  $S$  de centre  $Z$.
\end{enumerate}
\end{lemma}

\begin{proof}
Soit  $\mathcal{F}'$  la transformée stricte de  $\mathcal{F}$  par rapport à l'éclatement
 $S' \to S$  de  $S$  de centre  $Z$. Soit  $\mathcal{F}''$  la transformée
stricte de  $\mathcal{F}'$  par rapport à l'éclatement  $S'' \to S'$  de  $S'$  de centre
 $Z'$. Soit  $\mathcal{G}$  la transformée stricte de  $\mathcal{F}$  par rapport
à l'éclatement  $S'' \to S$  de  $S$  de centre  $Y$. Notons
également les morphismes
$$
\xymatrix{
X \times_S S'' \ar[r]_q \ar[d]^{f''} &
X \times_S S' \ar[r]_p \ar[d]^{f'} &
X \ar[d]^f \\
S'' \ar[r] & S' \ar[r] & S
}
$$
Par définition, il existe une surjection  $p^*\mathcal{F} \to \mathcal{F}'$  et une surjection
 $q^*\mathcal{F}' \to \mathcal{F}''$  qui se combinent par exactitude à droite de  $q^*$  en une
surjection  $(p \circ q)^*\mathcal{F} \to \mathcal{F}''$. De plus, nous avons la surjection  $(p \circ q)^*\mathcal{F} \to \mathcal{G}$. Il suffit
donc de prouver que ces deux surjections ont le même noyau.

\medskip\noindent
Le noyau de la surjection  $p^*\mathcal{F} \to \mathcal{F}'$  a son support dans  $(f \circ p)^{-1}Z$, donc cette
application est un isomorphisme aux points dans le complément. Par conséquent, le
noyau de  $q^*p^*\mathcal{F} \to q^*\mathcal{F}'$  a son support dans  $(f \circ p \circ q)^{-1}Z$. Le noyau de  $q^*\mathcal{F}' \to \mathcal{F}''$  a
son support dans  $(f' \circ q)^{-1}Z'$. Ensemble, ces faits montrent que le noyau de  $(p \circ q)^*\mathcal{F} \to \mathcal{F}''$  a
son support dans  $(f \circ p \circ q)^{-1}Z \cup (f' \circ q)^{-1}Z' =
(f \circ p \circ q)^{-1}Y$. Par construction de  $\mathcal{G}$ , nous voyons que nous
obtenons une factorisation  $(p \circ q)^*\mathcal{F} \to \mathcal{F}'' \to \mathcal{G}$. Pour terminer la preuve, il suffit de
montrer que  $\mathcal{F}''$  n'a pas de sections (locales) non nulles à support dans
 $(f \circ p \circ q)^{-1}(Y) =
(f \circ p \circ q)^{-1}Z \cup (f' \circ q)^{-1}Z'$. Cela découle du Lemme  \ref{divisors-lemma-strict-transform-different-centers}  appliqué à  $\mathcal{F}'$  sur
 $X \times_S S'$  au-dessus de  $S'$, le sous-schéma fermé  $Z'$  et le
diviseur de Cartier effectif  $b^{-1}Z$.
\end{proof}

\begin{lemma}
\label{divisors-lemma-strict-transform-universally-injective}
Dans la situation de la Définition  \ref{divisors-definition-strict-transform}. Supposons que
$$
0 \to \mathcal{F}_1 \to \mathcal{F}_2 \to \mathcal{F}_3 \to 0
$$
soit une suite exacte de faisceaux quasi-cohérents sur  $X$  qui reste
exacte après tout changement de base  $T \to S$. Alors les transformées strictes
$\mathcal{F}_i'$  relativement à tout éclatement  $S' \to S$  forment également une
suite exacte courte  $0 \to \mathcal{F}'_1 \to \mathcal{F}'_2 \to \mathcal{F}'_3 \to 0$ .
\end{lemma}

\begin{proof}
Nous pouvons nous localiser sur  $S$  et  $X$  et supposer que les
deux sont affines. Ensuite, nous pouvons pousser  $\mathcal{F}_i$  vers  $S$,
voir le Lemme  \ref{divisors-lemma-strict-transform-affine}. Nous pouvons supposer que notre éclatement est le
morphisme  $1 : S \to S$  associé à un diviseur de Cartier effectif  $D \subset S$.
La traduction en algèbre est alors la suivante. Supposons que  $A$  soit un
anneau et que  $0 \to M_1 \to M_2 \to M_3 \to 0$  soit une suite universellement exacte de
 $A$-modules. Soit  $a\in A$. Alors la suite
$$
0 \to
M_1/a\text{-torsion par puissances} \to
M_2/a\text{-torsion par puissances} \to
M_3/a\text{-torsion par puissances} \to 0
$$
est également exacte. En effet, la surjectivité de la dernière application et
l'injectivité de la première application sont immédiates. Le problème est
l'exactitude au milieu. Supposons que  $x \in M_2$  s'envoie à zéro dans
 $M_3/a\text{-torsion par puissances}$. Alors  $y = a^n x \in M_1$  pour un certain  $n$. Ensuite,  $y$ 
s'envoie à zéro dans  $M_2/a^nM_2$. Comme  $M_1 \to M_2$  est universellement
injectif, nous voyons que  $y$  s'envoie à zéro dans  $M_1/a^nM_1$. Ainsi,
 $y = a^n z$  pour un certain  $z \in M_1$. Ainsi,  $a^n(x - y) = 0$. Par conséquent,
 $y$  s'envoie à la classe de  $x$  dans  $M_2/a\text{-torsion par puissances}$  comme
souhaité.
\end{proof}







\section{Éclatements admissibles}
\label{divisors-section-admissible-blowups}

\noindent
Pour avoir un peu plus de contrôle sur nos éclatements, nous introduisons la
terminologie standard suivante.

\begin{definition}
\label{divisors-definition-admissible-blowup}
Soit  $X$  un schéma. Soit  $U \subset X$  un sous-schéma ouvert. Un morphisme
 $X' \to X$  est appelé un  {\it éclatement
$U$-admissible}  s'il existe une immersion fermée  $Z \to X$  de
présentation finie avec  $Z$  disjointe de  $U$  telle que
 $X'$  est isomorphe à l'éclatement de  $X$  de centre  $Z$.
\end{definition}

\noindent
Nous rappelons que  $Z \to X$  est de présentation finie si et seulement si le
faisceau d'idéaux  $\mathcal{I}_Z \subset \mathcal{O}_X$  est de type fini, voir Morphismes, Lemme
 \ref{morphisms-lemma-closed-immersion-finite-presentation}. En particulier, un éclatement  $U$-admissible est un
morphisme projectif, voir Lemme  \ref{divisors-lemma-blowing-up-projective}. Notez qu'il peut y avoir plusieurs
centres donnant lieu au même morphisme. Par conséquent, l'exigence est simplement
l'existence d'un certain centre disjoint de  $U$  qui produit  $X'$.
Enfin, comme le morphisme  $b : X' \to X$  est un isomorphisme au-dessus de
 $U$  (voir Lemme  \ref{divisors-lemma-blowing-up-gives-effective-Cartier-divisor}), nous abuserons souvent de la notation et
considérerons  $U$  comme un sous-schéma ouvert de  $X'$  également.

\begin{lemma}
\label{divisors-lemma-composition-admissible-blowups}
\begin{slogan}
Les éclatements admissibles sont stables par composition.
\end{slogan}
Soit  $X$  un schéma quasi-compact et quasi-séparé. Soit  $U \subset X$  un
sous-schéma ouvert quasi-compact. Soit  $b : X' \to X$  un éclatement
 $U$-admissible. Soit  $X'' \to X'$  un éclatement  $U$-admissible.
Alors la composition  $X'' \to X$  est un éclatement  $U$-admissible.
\end{lemma}

\begin{proof}
Immédiat à partir du lemme plus précis  \ref{divisors-lemma-composition-finite-type-blowups}.
\end{proof}

\begin{lemma}
\label{divisors-lemma-extend-admissible-blowups}
Soit  $X$  un schéma quasi-compact et quasi-séparé. Soit  $U, V \subset X$  des
sous-schémas ouverts quasi-compacts. Soit  $b : V' \to V$  un éclatement
 $U \cap V$-admissible. Alors il existe un éclatement  $U$-admissible
 $X' \to X$  dont la restriction à  $V$  est  $V'$.
\end{lemma}

\begin{proof}
Soit  $\mathcal{I} \subset \mathcal{O}_V$  le faisceau quasi-cohérent d'idéaux de type fini tel que
 $V(\mathcal{I})$  soit disjoint de  $U \cap V$  et tel que  $V'$  soit isomorphe
à l'éclatement de  $V$  relativement à  $\mathcal{I}$. Soit  $\mathcal{I}' \subset \mathcal{O}_{U \cup V}$  le faisceau
quasi-cohérent d'idéaux dont la restriction à  $U$  est  $\mathcal{O}_U$  et
dont la restriction à  $V$  est  $\mathcal{I}$  (voir Faisceaux, section
 \ref{sheaves-section-glueing-sheaves}). D'après Propriétés, Lemme  \ref{properties-lemma-extend} , il existe un faisceau
quasi-cohérent d'idéaux de type fini  $\mathcal{J} \subset \mathcal{O}_X$  dont la restriction à
 $U \cup V$  est  $\mathcal{I}'$. Le lemme s'ensuit.
\end{proof}

\begin{lemma}
\label{divisors-lemma-dominate-admissible-blowups}
Soit  $X$  un schéma quasi-compact et quasi-séparé. Soit  $U \subset X$  un
sous-schéma ouvert quasi-compact. Soient  $b_i : X_i \to X$,  $i = 1, \ldots, n$  des
éclatements  $U$-admissibles. Il existe un éclatement  $U$-admissible
 $b : X' \to X$  tel que (a)  $b$  se factorise sous la forme
 $X' \to X_i \to X$  pour  $i = 1, \ldots, n$  et (b) chacun des morphismes  $X' \to X_i$  est un
éclatement admissible  $U$.
\end{lemma}

\begin{proof}
Soit  $\mathcal{I}_i \subset \mathcal{O}_X$  le faisceau quasi-cohérent d'idéaux de type fini tel que
 $V(\mathcal{I}_i)$  soit disjoint de  $U$  et tel que  $X_i$  soit isomorphe
à l'éclatement de  $X$  relativement à  $\mathcal{I}_i$. Posons  $\mathcal{I} = \mathcal{I}_1 \cdot \ldots \cdot \mathcal{I}_n$  et soit
 $X'$  l'éclatement de  $X$  relativement à  $\mathcal{I}$. Alors  $X' \to X$  se
factorise en passant par  $b_i$  d'après le Lemme  \ref{divisors-lemma-blowing-up-two-ideals}.
\end{proof}

\begin{lemma}
\label{divisors-lemma-separate-disjoint-opens-by-blowing-up}
\begin{slogan}
Séparer les composantes irréductibles en effectuant un éclatement.
\end{slogan}
Soit  $X$  un schéma quasi-compact et quasi-séparé. Soient  $U, V$  des
sous-schémas ouverts disjoints quasi-compacts de  $X$. Alors il existe un
éclatement  $U \cup V$-admissible  $b : X' \to X$  tel que  $X'$  soit une
union disjointe de sous-schémas ouverts  $X' = X'_1 \amalg X'_2$  avec  $b^{-1}(U) \subset X'_1$  et
 $b^{-1}(V) \subset X'_2$.
\end{lemma}

\begin{proof}
Choisissons un faisceau d'idéaux quasi-cohérent de type fini  $\mathcal{I}$,
respectivement \ $\mathcal{J}$  tel que  $X \setminus U = V(\mathcal{I})$, respectivement
\ $X \setminus V = V(\mathcal{J})$, voir Propriétés, Lemme  \ref{properties-lemma-quasi-coherent-finite-type-ideals}. Alors  $V(\mathcal{I}\mathcal{J}) = X$ 
au sens des ensembles, donc  $\mathcal{I}\mathcal{J}$  est un faisceau d'idéaux localement
nilpotent. Puisque  $\mathcal{I}$  et  $\mathcal{J}$  sont de type fini et que
 $X$  est quasi-compact, il existe un  $n > 0$  tel que  $\mathcal{I}^n \mathcal{J}^n = 0$.
 Nous pouvons donc remplacer  $\mathcal{I}$  par  $\mathcal{I}^n$  et  $\mathcal{J}$  par
 $\mathcal{J}^n$. Dès lors,  $\mathcal{I} \mathcal{J} = 0$. Soit  $b : X' \to X$  l'éclatement relativement à  $\mathcal{I} + \mathcal{J}$.
Ceci est  $U \cup V$-admissible car  $V(\mathcal{I} + \mathcal{J}) = X \setminus U \cup V$. Nous montrerons que
 $X'$  est une réunion disjointe de sous-schémas ouverts  $X' = X'_1 \amalg X'_2$  tels
que  $b^{-1}\mathcal{I}|_{X'_2} = 0$  et  $b^{-1}\mathcal{J}|_{X'_1} = 0$ , ce qui prouvera le lemme.

\medskip\noindent
Nous utiliserons la description de l'éclatement donnée dans le Lemme  \ref{divisors-lemma-blowing-up-affine}.
Supposons que  $U = \Spec(A) \subset X$  soit un ouvert affine tel que  $\mathcal{I}|_U$, resp.
\ $\mathcal{J}|_U$  correspond à l'idéal de type fini  $I \subset A$, resp.
\ $J \subset A$. Alors
$$
b^{-1}(U) = \text{Proj}(A \oplus (I + J) \oplus (I + J)^2 \oplus \ldots)
$$
Celui-ci est recouvert par les ouverts affines correspondant aux algèbres  $A[\frac{I + J}{x}]$  et
 $A[\frac{I + J}{y}]$  avec  $x \in I$  et  $y \in J$. Puisque  $x \in I$  est un
non-diviseur de zéro dans  $A[\frac{I + J}{x}]$  et  $IJ = 0$ , nous voyons que
 $J A[\frac{I + J}{x}] = 0$. Puisque  $y \in J$  est un non-diviseur de zéro dans  $A[\frac{I + J}{y}]$ 
et  $IJ = 0$ , nous voyons que  $I A[\frac{I + J}{y}] = 0$. De plus,
$$
\Spec(A[\textstyle{\frac{I + J}{x}}]) \cap
\Spec(A[\textstyle{\frac{I + J}{y}}]) =
\Spec(A[\textstyle{\frac{I + J}{xy}}]) = \emptyset
$$
parce que  $xy$  est à la fois un non-diviseur de zéro et nul. Ainsi,
 $b^{-1}(U)$  est l'union disjointe du sous-schéma ouvert  $U_1$  défini comme
l'union des ouverts standards  $\Spec(A[\frac{I + J}{x}])$  pour  $x \in I$  et du sous-schéma
ouvert  $U_2$  qui est l'union des ouverts affines  $\Spec(A[\frac{I + J}{y}])$  pour
 $y \in J$. Nous avons vu que  $b^{-1}\mathcal{I}\mathcal{O}_{X'}$  se restreint à zéro sur  $U_2$ 
et  $b^{-1}\mathcal{I}\mathcal{O}_{X'}$  se restreint à zéro sur  $U_1$. Nous omettons la
vérification que ces sous-schémas ouverts se recollent en sous-schémas ouverts
globaux  $X'_1$  et  $X'_2$.
\end{proof}

\begin{lemma}
\label{divisors-lemma-blowing-up-denominators}
Soit  $X$  un schéma localement noethérien. Soit  $\mathcal{L}$  un
 $\mathcal{O}_X$-module inversible. Soit  $s$  une section méromorphe régulière
de  $\mathcal{L}$. Soit  $U \subset X$  le plus grand sous-schéma ouvert tel que
 $s$  corresponde à une section de  $\mathcal{L}$  sur  $U$.
L'éclatement  $b : X' \to X$  relativement à l'idéal des dénominateurs de  $s$  est
 $U$-admissible. Il existe un diviseur de Cartier effectif  $D \subset X'$  et
un isomorphisme
$$
b^*\mathcal{L} = \mathcal{O}_{X'}(D - E),
$$
où  $E \subset X'$  est le diviseur exceptionnel et où la section méromorphe
 $b^*s$  correspond, via l'isomorphisme, à la section méromorphe
 $1_D \otimes (1_E)^{-1}$.
\end{lemma}

\begin{proof}
D'après la définition de l'idéal des dénominateurs dans la Définition
 \ref{divisors-definition-regular-meromorphic-ideal-denominators} , nous voyons immédiatement que  $b$  est un éclatement
 $U$-admissible. Pour la notation  $1_D$,  $1_E$ et
 $\mathcal{O}_{X'}(D - E)$ , voir la Définition  \ref{divisors-definition-invertible-sheaf-effective-Cartier-divisor}. L'image réciproque
 $b^*s$  est définie par les lemmes  \ref{divisors-lemma-blow-up-pullback-effective-Cartier}  et  \ref{divisors-lemma-meromorphic-sections-pullback}.
Ainsi, l'énoncé du lemme a un sens. Nous pouvons réinterpréter la dernière
assertion en disant que  $b^*s$  est une section régulière globale de
 $b^*\mathcal{L}(E)$  dont le schéma des zéros est  $D$. Cela définit de manière
unique  $D$ , donc pour prouver le lemme, nous pouvons travailler
localement de manière affine sur  $X$  et  $X'$. Supposons que
 $X = \Spec(A)$  soit affine et  $\mathcal{L} = \mathcal{O}_X$. Alors  $s$  est une fonction
méromorphe régulière et en rétrécissant davantage nous pouvons supposer
 $s = a'/a$  avec  $a', a \in A$  des non-diviseurs de zéro. Alors l'idéal des
dénominateurs de  $s$  correspond à l'idéal  $I = \{x \in A \mid xa' \in aA\}$. Rappelons que
 $X'$  est recouvert par les spectres des algèbres d'éclatement affines
 $A' = A[\frac{I}{x}]$  avec  $x \in I$  (Lemme  \ref{divisors-lemma-blowing-up-affine}). Fixons  $x \in I$  et
écrivons  $xa' = a a''$  pour un certain  $a'' \in A$. Le diviseur  $E \subset X'$  est
défini par  $x \in A'$  sur le spectre de  $A'$  et donc  $1/x$  est un
générateur de  $\mathcal{O}_{X'}(E)$  sur  $\Spec(A')$. Enfin, dans l'anneau total de
fractions de  $A'$  nous avons  $a'/a = a''/x$. Par conséquent,  $b^*s = a'/a$  se
restreint à une section régulière de  $\mathcal{O}_{X'}(E)$  qui, sur
 $\Spec(A')$, est donnée par  $a''/x$. Cela termine la preuve. (Le diviseur
 $D \cap \Spec(A')$  est découpé par l'image de  $a''$  dans  $A'$.)
\end{proof}





\section{Éclatements et platitude}
\label{divisors-section-blowup-flat}

\noindent
Nous continuons la discussion commencée dans Compléments d'algèbre, section
 \ref{more-algebra-section-blowup-flat}. Nous prouverons des résultats supplémentaires dans Compléments sur la
platitude, section  \ref{flat-section-blowup-flat}.


\begin{lemma}
\label{divisors-lemma-strict-transform-blowup-fitting-ideal}
Soit  $S$  un schéma. Soit  $\mathcal{F}$  un  $\mathcal{O}_S$-module
quasi-cohérent de type fini. Soit  $Z_k \subset S$  le sous-schéma fermé découpé par
 $\text{Fit}_k(\mathcal{F})$, voir section  \ref{divisors-section-fitting-ideals}. Soit  $S' \to S$  l'éclatement de
 $S$  de centre  $Z_k$  et soit  $\mathcal{F}'$  la transformée stricte de
 $\mathcal{F}$. Alors  $\mathcal{F}'$  peut être engendré localement par  $\leq k$ 
sections.
\end{lemma}

\begin{proof}
Rappelons que  $\mathcal{F}'$  peut être engendré localement par  $\leq k$  sections
si et seulement si  $\text{Fit}_k(\mathcal{F}') = \mathcal{O}_{S'}$, voir le Lemme  \ref{divisors-lemma-fitting-ideal-generate-locally}. Ainsi ce
lemme est une traduction de Compléments d'algèbre, Lemme  \ref{more-algebra-lemma-blowup-fitting-ideal}.
\end{proof}

\begin{lemma}
\label{divisors-lemma-strict-transform-blowup-fitting-ideal-locally-free}
Soit  $S$  un schéma. Soit  $\mathcal{F}$  un  $\mathcal{O}_S$-module
quasi-cohérent de type fini. Soit  $Z_k \subset S$  le sous-schéma fermé découpé par
 $\text{Fit}_k(\mathcal{F})$, voir la section  \ref{divisors-section-fitting-ideals}. Supposons que  $\mathcal{F}$  soit
localement libre de rang  $k$  sur  $S \setminus Z_k$. Soit  $S' \to S$ 
l'éclatement de  $S$  de centre  $Z_k$  et soit  $\mathcal{F}'$  la transformée
stricte de  $\mathcal{F}$. Alors  $\mathcal{F}'$  est localement libre de rang
 $k$.
\end{lemma}

\begin{proof}
Traduction de Compléments d'algèbre, Lemme  \ref{more-algebra-lemma-blowup-fitting-ideal-locally-free}.
\end{proof}

\begin{lemma}
\label{divisors-lemma-blowup-fitting-ideal}
Soit  $X$  un schéma. Soit  $\mathcal{F}$  un  $\mathcal{O}_X$-module de
présentation finie. Soit  $U \subset X$  un ouvert schématiquement dense tel que
 $\mathcal{F}|_U$  soit localement libre de rang constant  $r$. Alors
\begin{enumerate}
\item l'éclatement  $b : X' \to X$  de  $X$  relativement au  $r$-ième idéal de
Fitting de  $\mathcal{F}$  est  $U$-admissible,
\item la transformée stricte  $\mathcal{F}'$  de  $\mathcal{F}$  par rapport à  $b$  est
localement libre de rang  $r$,
\item le noyau  $\mathcal{K}$  de la surjection  $b^*\mathcal{F} \to \mathcal{F}'$  est de présentation finie et
 $\mathcal{K}|_U = 0$,
\item $b^*\mathcal{F}$  et  $\mathcal{K}$  sont des  $\mathcal{O}_{X'}$-modules parfaits de dimension
de Tor  $\leq 1$.
\end{enumerate}
\end{lemma}

\begin{proof}
L'idéal  $\text{Fit}_r(\mathcal{F})$  est de type fini d'après le Lemme  \ref{divisors-lemma-fitting-ideal-of-finitely-presented}  et sa
restriction à  $U$  est égale à  $\mathcal{O}_U$  d'après le Lemme  \ref{divisors-lemma-fitting-ideal-finite-locally-free}.
Par conséquent,  $b : X' \to X$  est  $U$-admissible, voir la Définition
 \ref{divisors-definition-admissible-blowup}.

\medskip\noindent
D'après le Lemme  \ref{divisors-lemma-fitting-ideal-finite-locally-free} , la restriction de  $\text{Fit}_{r - 1}(\mathcal{F})$  à  $U$  est
nulle, et puisque  $U$  est schématiquement dense, nous concluons que
 $\text{Fit}_{r - 1}(\mathcal{F}) = 0$  sur tout  $X$. Ainsi il découle du Lemme  \ref{divisors-lemma-fitting-ideal-finite-locally-free}  que
 $\mathcal{F}$  est localement libre de rang  $r$  sur le complément du
sous-schéma défini par le  $r$-ième idéal de Fitting de  $\mathcal{F}$  (ce
complément peut être plus grand que  $U$ , ce qui explique pourquoi nous
avons dû effectuer cette étape dans l'argument). Par conséquent, d'après le Lemme
 \ref{divisors-lemma-strict-transform-blowup-fitting-ideal-locally-free} , la transformée stricte
$$
b^*\mathcal{F} \longrightarrow \mathcal{F}'
$$
est localement libre de rang  $r$. Le noyau  $\mathcal{K}$  de cette
application a son support dans le diviseur exceptionnel de l'éclatement  $b$ 
et donc  $\mathcal{K}|_U = 0$. Enfin, puisque  $\mathcal{F}'$  est localement libre de type
fini et que la flèche affichée est surjective, nous pouvons localement sur
 $X'$  écrire  $b^*\mathcal{F}$  comme la somme directe de  $\mathcal{K}$  et
 $\mathcal{F}'$. Puisque  $b^*\mathcal{F}'$  est de présentation finie (Modules, Lemme
 \ref{modules-lemma-pullback-finite-presentation}), il en est de même pour  $\mathcal{K}$.

\medskip\noindent
L'énoncé sur la dimension de Tor découle de Compléments d'algèbre, Lemme  \ref{more-algebra-lemma-fitting-ideals-and-pd1}.
\end{proof}





\section{Modifications}
\label{divisors-section-modifications}

\noindent
Dans cette section, nous rassemblerons des résultats affirmant qu'après une
modification, certaines propriétés deviennent vraies. Nous verrons plus tard qu'une
modification peut être dominée par un éclatement (Compléments sur la platitude, Lemme
 \ref{flat-lemma-dominate-modification-by-blowup}).

\begin{lemma}
\label{divisors-lemma-filter-after-modification}
Soit  $X$  un schéma intègre. Soit  $\mathcal{E}$  un  $\mathcal{O}_X$-module
localement libre de type fini. Il existe une modification  $f : X' \to X$  telle que
 $f^*\mathcal{E}$  possède une filtration dont les quotients successifs sont des
 $\mathcal{O}_{X'}$-modules inversibles.
\end{lemma}

\begin{proof}
Nous prouvons cela par induction sur le rang  $r$  de  $\mathcal{E}$. Si
 $r = 1$  ou  $r = 0$ , le lemme est évident. Supposons  $r > 1$. Soit
 $P = \mathbf{P}(\mathcal{E})$  avec un morphisme de structure  $\pi : P \to X$, voir Constructions,
section  \ref{constructions-section-projective-bundle}. Alors  $\pi$  est propre (Lemme  \ref{divisors-lemma-relative-proj-proper}). Il
existe une surjection canonique
$$
\pi^*\mathcal{E} \to \mathcal{O}_P(1)
$$
dont le noyau est localement libre de rang  $r - 1$. Choisissons un sous-schéma
ouvert non vide  $U \subset X$  tel que  $\mathcal{E}|_U \cong \mathcal{O}_U^{\oplus r}$. Alors  $P_U = \pi^{-1}(U)$  est
isomorphe à  $\mathbf{P}^{r - 1}_U$. En particulier, il existe une section  $s : U \to P_U$  de
 $\pi$. Soit  $X' \subset P$  l'image schématique du morphisme  $U \to P_U \to P$.
Alors  $X'$  est intègre (Morphismes, Lemme  \ref{morphisms-lemma-scheme-theoretic-image-reduced}), le morphisme
 $f = \pi|_{X'} : X' \to X$  est propre (Morphismes, Lemme  \ref{morphisms-lemma-closed-immersion-proper}  et  \ref{morphisms-lemma-composition-proper}), et
 $f^{-1}(U) \to U$  est un isomorphisme. Par conséquent,  $f$  est une
modification (Morphismes, Définition  \ref{morphisms-definition-modification}). Par construction, l'image réciproque
 $f^*\mathcal{E}$  possède une filtration en deux étapes dont le quotient est
inversible car il est égal à  $\mathcal{O}_P(1)|_{X'}$  et dont le sous-module  $\mathcal{E}'$  est
localement libre de rang  $r - 1$. Par récurrence, nous pouvons trouver une
modification  $g : X'' \to X'$  telle que  $g^*\mathcal{E}'$  possède une filtration comme dans
l'énoncé du lemme. Ainsi,  $f \circ g : X'' \to X$  est la modification requise.
\end{proof}

\begin{lemma}
\label{divisors-lemma-extend-rational-map-after-modification}
Soit  $S$  un schéma. Soient  $X$,  $Y$  des schémas sur
 $S$. Supposons que  $X$  soit noethérien et que  $Y$  soit
propre sur  $S$. Étant donnée une application rationnelle sur  $S$,
 $f : U \to Y$, de  $X$  vers  $Y$ , il existe un morphisme
 $p : X' \to X$  et un  $S$-morphisme  $f' : X' \to Y$  tels que
\begin{enumerate}
\item $p$  soit propre et  $p^{-1}(U) \to U$  soit un isomorphisme,
\item $f'|_{p^{-1}(U)}$  soit égal à  $f \circ p|_{p^{-1}(U)}$.
\end{enumerate}
\end{lemma}

\begin{proof}
On note  $j : U \to X$  le morphisme d'inclusion. Soit  $X' \subset Y \times_S X$  l'image
schématique de  $(f, j) : U \to Y \times_S X$  (Morphismes, Définition  \ref{morphisms-definition-scheme-theoretic-image}). La projection
 $g : Y \times_S X \to X$  est propre (Morphismes, Lemme  \ref{morphisms-lemma-base-change-proper}). La composition
 $p : X' \to X$  de  $X' \to Y \times_S X$  et  $g$  est propre (Morphismes, Lemmes
 \ref{morphisms-lemma-closed-immersion-proper}  et  \ref{morphisms-lemma-composition-proper}). Comme  $g$  est séparé et  $U \subset X$  est
rétrocompact (puisque  $X$  est noethérien), nous concluons que
 $p^{-1}(U) \to U$  est un isomorphisme d'après Morphismes, Lemme  \ref{morphisms-lemma-scheme-theoretic-image-of-partial-section}. D'autre
part, la composition  $f' : X' \to Y$  de  $X' \to Y \times_S X$  et de la projection  $Y \times_S X \to Y$ 
coïncide avec  $f$  sur  $p^{-1}(U)$.
\end{proof}











% OK, start here.
%

\chapter{Limites de schémas}



\phantomsection
\label{limits-section-phantom}


\section{Introduction}
\label{limits-section-introduction}

\noindent
Dans ce chapitre, nous rassemblons des résultats relatifs aux limites de
schémas. Nous étudions principalement les limites de systèmes projectifs
indexés par des ensembles ordonnés filtrants
(Catégories, Définition \ref{categories-definition-directed-set})
dont les morphismes de transition sont affines. Nous traitons de
l'approximation noethérienne absolue. Nous caractérisons les schémas
localement de présentation finie sur une base comme ceux dont le foncteur
des points associé commute aux limites. Comme application de
l'approximation noethérienne absolue, nous démontrons que l'image d'un
schéma affine par un morphisme entier est affine. Nous établissons en outre
quelques variantes très générales du lemme de Chow. Une référence de base
est \cite{EGA}.




\section{Limites projectives de schémas à morphismes de transition affines}
\label{limits-section-limits}

\noindent
Dans cette section, nous construisons la limite.

\begin{lemma}
\label{limits-lemma-directed-inverse-system-affine-schemes-has-limit}
Soit $I$ un ensemble ordonné filtrant. Soit $(S_i, f_{ii'})$ un système
projectif de schémas indexé par $I$. Si tous les schémas $S_i$ sont
affines, alors la limite $S = \lim_i S_i$ existe dans la catégorie des
schémas. En fait, $S$ est affine et $S = \Spec(\colim_i R_i)$, où
$R_i = \Gamma(S_i, \mathcal{O})$.
\end{lemma}

\begin{proof}
Il suffit de poser $S = \Spec(\colim_i R_i)$.
Il résulte de Schémas, Lemme \ref{schemes-lemma-morphism-into-affine}
que $S$ est même la limite dans la catégorie des espaces localement
annelés.
\end{proof}

\begin{lemma}
\label{limits-lemma-directed-inverse-system-has-limit}
Soit $I$ un ensemble ordonné filtrant. Soit $(S_i, f_{ii'})$ un système
projectif de schémas indexé par $I$. Si tous les morphismes
$f_{ii'} : S_i \to S_{i'}$ sont affines, alors la limite $S = \lim_i S_i$
existe dans la catégorie des schémas. De plus,
\begin{enumerate}
\item chacun des morphismes $f_i : S \to S_i$ est affine,
\item pour un élément $0 \in I$ et tout sous-schéma ouvert
$U_0 \subset S_0$, on a
$$
f_0^{-1}(U_0) = \lim_{i \geq 0} f_{i0}^{-1}(U_0)
$$
dans la catégorie des schémas.
\end{enumerate}
\end{lemma}

\begin{proof}
Choisissons un élément $0 \in I$. Remarquons que $I$ est non vide puisqu'il
est filtrant. Pour tout $i \geq 0$, considérons le faisceau quasi-cohérent
d'$\mathcal{O}_{S_0}$-algèbres
$\mathcal{A}_i = f_{i0, *}\mathcal{O}_{S_i}$. Rappelons que
$S_i = \underline{\Spec}_{S_0}(\mathcal{A}_i)$ ; voir Morphismes,
Lemme \ref{morphisms-lemma-characterize-affine}. Posons
$\mathcal{A} = \colim_{i \geq 0} \mathcal{A}_i$. C'est un faisceau
quasi-cohérent d'$\mathcal{O}_{S_0}$-algèbres ; voir Schémas,
section \ref{schemes-section-quasi-coherent}. Posons
$S = \underline{\Spec}_{S_0}(\mathcal{A})$. Par Morphismes,
Lemme \ref{morphisms-lemma-affine-equivalence-algebras}, nous obtenons,
pour $i \geq 0$, des morphismes $f_i : S \to S_i$ compatibles avec les
morphismes de transition. Les morphismes $f_i$ sont affines, par exemple
par Morphismes, Lemme \ref{morphisms-lemma-affine-permanence}. D'après le
Lemme \ref{limits-lemma-directed-inverse-system-affine-schemes-has-limit} ci-dessus,
pour tout ouvert affine $U_0 \subset S_0$, l'image réciproque
$U = f_0^{-1}(U_0) \subset S$ est la limite, dans la catégorie des schémas,
du système des ouverts $U_i = f_{i0}^{-1}(U_0)$, $i \geq 0$.

\medskip\noindent
Soit $T$ un schéma. Soit $g_i : T \to S_i$ un système compatible de
morphismes. Pour montrer que $S = \lim_i S_i$, il faut établir qu'il existe
un unique morphisme $g : T \to S$ tel que $g_i = f_i \circ g$ pour tout
$i \in I$. Pour tout $t \in T$, il existe un ouvert affine
$U_0 \subset S_0$ contenant $g_0(t)$. Soit
$V \subset g_0^{-1}(U_0)$ un voisinage ouvert affine de $t$.
D'après les remarques précédentes, nous obtenons un unique morphisme
$g_V : V \to U = f_0^{-1}(U_0)$ tel que $f_i \circ g_V = g_i|_{U_i}$
pour tout $i$. Les ouverts $V \subset T$ ainsi construits forment une base
de la topologie de $T$. Les morphismes $g_V$ se recollent en un morphisme
$g : T \to S$ grâce à la propriété d'unicité. On obtient ainsi le morphisme
cherché $g : T \to S$.

\medskip\noindent
La dernière assertion résulte clairement de la construction précédente de
la limite.
\end{proof}

\begin{lemma}
\label{limits-lemma-scheme-over-limit}
Soit $I$ un ensemble ordonné filtrant.
Soit $(S_i, f_{ii'})$ un système projectif de schémas indexé par $I$.
Supposons tous les morphismes $f_{ii'} : S_i \to S_{i'}$ affines.
Soit $S = \lim_i S_i$. Soit $0 \in I$. Supposons que $T$ soit un schéma
sur $S_0$. Alors
$$
T \times_{S_0} S = \lim_{i \geq 0} T \times_{S_0} S_i
$$
\end{lemma}

\begin{proof}
Le membre de droite est un schéma d'après le
Lemme \ref{limits-lemma-directed-inverse-system-has-limit}. L'égalité est formelle ;
voir Catégories, Lemme \ref{categories-lemma-colimits-commute}.
\end{proof}



\section{Produits infinis}
\label{limits-section-inifinite-products}

\noindent
En général, les produits infinis de schémas n'existent pas. Par exemple,
il est montré dans Exemples, section \ref{examples-section-not-algebraic},
qu'un produit infini de copies de $\mathbf{P}^1$ n'est même pas un espace
algébrique.

\medskip\noindent
En revanche, les produits infinis de schémas affines existent et sont
affines. Par Schémas, Lemme \ref{schemes-lemma-morphism-into-affine}, cela
correspond au fait que la catégorie des anneaux admet des coproduits
infinis : si $I$ est un ensemble et si $R_i$ est un anneau pour tout $i$,
on peut considérer l'anneau
$$
R = \otimes R_i =
\colim_{\{i_1, \ldots, i_n\} \subset I}
R_{i_1} \otimes_\mathbf{Z} \ldots \otimes_\mathbf{Z} R_{i_n}
$$
Étant donné un autre anneau $A$, une application $R \to A$ équivaut à une
famille d'homomorphismes d'anneaux $R_i \to A$, pour tout $i \in I$, comme
il résulte de la propriété correspondante des produits tensoriels finis.

\begin{lemma}
\label{limits-lemma-infinite-product}
\begin{slogan}
Les produits infinis de schémas affines existent et sont affines.
\end{slogan}
Soit $S$ un schéma. Soit $I$ un ensemble et, pour tout $i \in I$, soit
$f_i : T_i \to S$ un morphisme affine. Alors le produit $T = \prod T_i$
existe dans la catégorie des schémas sur $S$. En fait, on a
$$
T = \lim_{\{i_1, \ldots, i_n\} \subset I}
T_{i_1} \times_S \ldots \times_S T_{i_n}
$$
et les morphismes de projection
$T \to T_{i_1} \times_S \ldots \times_S T_{i_n}$ sont affines.
\end{lemma}

\begin{proof}
Omis. Indication : raisonner comme dans la discussion qui précède le lemme
et utiliser le Lemme \ref{limits-lemma-directed-inverse-system-has-limit} pour
l'existence de la limite.
\end{proof}

\begin{lemma}
\label{limits-lemma-infinite-product-surjective}
Soit $S$ un schéma. Soit $I$ un ensemble et, pour tout $i \in I$, soit
$f_i : T_i \to S$ un morphisme affine surjectif. Alors le produit
$T = \prod T_i$ dans la catégorie des schémas sur $S$
(Lemme \ref{limits-lemma-infinite-product}) se projette surjectivement sur $S$.
\end{lemma}

\begin{proof}
Soit $s \in S$. Choisissons $t_i \in T_i$ d'image $s$.
Choisissons une extension de corps suffisamment grande $K/\kappa(s)$ telle
que $\kappa(s_i)$ se plonge dans $K$ pour tout $i$. Nous obtenons alors des
morphismes $\Spec(K) \to T_i$ d'image $s_i$ qui coïncident comme morphismes
vers $S$. Il en résulte un morphisme $\Spec(K) \to T$, ce qui prouve qu'il
existe un point de $T$ d'image $s$.
\end{proof}

\begin{lemma}
\label{limits-lemma-infinite-product-integral}
Soit $S$ un schéma. Soit $I$ un ensemble et, pour tout $i \in I$, soit
$f_i : T_i \to S$ un morphisme entier. Alors le produit $T = \prod T_i$
dans la catégorie des schémas sur $S$ (Lemme \ref{limits-lemma-infinite-product})
est entier sur $S$.
\end{lemma}

\begin{proof}
Omis. Indication : sur des morceaux affines, cela se ramène au fait
algébrique suivant : si $A \to B_i$ est entier pour tout $i$, alors
$A \to \otimes_A B_i$ est entier.
\end{proof}





\section{Descente des propriétés}
\label{limits-section-descent}

\noindent
D'abord quelques lemmes élémentaires décrivant la topologie d'une limite.

\begin{lemma}
\label{limits-lemma-inverse-limit-sets}
Soit $S = \lim S_i$ la limite d'un système projectif filtrant de schémas
dont les morphismes de transition sont affines
(Lemme \ref{limits-lemma-directed-inverse-system-has-limit}). Alors
$S_{set} = \lim_i S_{i, set}$, où $S_{set}$ désigne l'ensemble sous-jacent
au schéma $S$.
\end{lemma}

\begin{proof}
Fixons $i \in I$ et prenons un ouvert affine $U_i \subset S_i$.
Notons $U_{i'} = f_{i'i}^{-1}(U_i)$ et $U = f_i^{-1}(U_i)$. Ici,
$f_{i'i} : S_{i'} \to S_i$ est le morphisme de transition et
$f_i : S \to S_i$ la projection. D'après le
Lemme \ref{limits-lemma-directed-inverse-system-has-limit}, on a
$U = \lim_{i' \geq i} U_i$. Supposons que l'on puisse montrer
$U_{set} = \lim_{i' \geq i} U_{i', set}$. Le lemme en résulte alors par un
argument simple utilisant un recouvrement affine de $S_i$. Nous pouvons
donc supposer que tous les $S_i$ et $S$ sont affines. Cela nous ramène à la
question algébrique examinée au paragraphe suivant.

\medskip\noindent
Soit donné un système d'anneaux $(A_i, \varphi_{ii'})$ indexé par $I$.
Posons $A = \colim_i A_i$, avec les homomorphismes canoniques
$\varphi_i : A_i \to A$. Alors
$$
\Spec(A) = \lim_i \Spec(A_i)
$$
En effet, supposons donnés des idéaux premiers
$\mathfrak p_i \subset A_i$ tels que
$\mathfrak p_i = \varphi_{ii'}^{-1}(\mathfrak p_{i'})$ pour tout
$i' \geq i$. Il suffit alors de poser
$$
\mathfrak p =
\{x \in A
\mid
\exists i, x_i \in \mathfrak p_i \text{ avec }\varphi_i(x_i) = x\}
$$
Il est clair que c'est un idéal et qu'il vérifie
$\varphi_i^{-1}(\mathfrak p) = \mathfrak p_i$. On en déduit aisément que
c'est également un idéal premier.
\end{proof}

\begin{lemma}
\label{limits-lemma-inverse-limit-top}
\begin{reference}
\cite[IV, Proposition 8.2.9]{EGA}
\end{reference}
Soit $S = \lim S_i$ la limite d'un système projectif filtrant de schémas
dont les morphismes de transition sont affines
(Lemme \ref{limits-lemma-directed-inverse-system-has-limit}). Alors
$S_{top} = \lim_i S_{i, top}$, où $S_{top}$ désigne l'espace topologique
sous-jacent au schéma $S$.
\end{lemma}

\begin{proof}
Nous utiliserons le critère de Topologie,
Lemme \ref{topology-lemma-characterize-limit}. Nous avons vu que
$S_{set} = \lim_i S_{i, set}$ dans le
Lemme \ref{limits-lemma-inverse-limit-sets}. Les applications $f_i : S \to S_i$ sont des
morphismes de schémas, donc sont continues. Ainsi, $f_i^{-1}(U_i)$ est
ouvert pour tout ouvert $U_i \subset S_i$. Enfin, soient $s \in S$ et
$s \in V \subset S$ un voisinage ouvert. Choisissons $0 \in I$ et un
voisinage ouvert affine $U_0 \subset S_0$ de l'image de $s$. Alors
$f_0^{-1}(U_0) = \lim_{i \geq 0} f_{i0}^{-1}(U_0)$ ; voir le
Lemme \ref{limits-lemma-directed-inverse-system-has-limit}. Les schémas
$f_0^{-1}(U_0)$ et $f_{i0}^{-1}(U_0)$ sont affines et
$$
\mathcal{O}_S(f_0^{-1}(U_0)) =
\colim_{i \geq 0} \mathcal{O}_{S_i}(f_{i0}^{-1}(U_0))
$$
soit d'après la démonstration du
Lemme \ref{limits-lemma-directed-inverse-system-has-limit}, soit d'après le
Lemme \ref{limits-lemma-directed-inverse-system-affine-schemes-has-limit}.
Choisissons $a \in \mathcal{O}_S(f_0^{-1}(U_0))$ tel que
$s \in D(a) \subset V$. C'est possible puisque les ouverts principaux
forment une base de la topologie du schéma affine $f_0^{-1}(U_0)$.
On peut alors choisir $i \geq 0$ et
$a_i \in \mathcal{O}_{S_i}(f_{i0}^{-1}(U_0))$ d'image $a$. Il en résulte
que $D(a_i) \subset f_{i0}^{-1}(U_0) \subset S_i$ est un ouvert dont
l'image réciproque dans $S$ est $D(a)$. Cela achève la démonstration.
\end{proof}

\begin{lemma}
\label{limits-lemma-limit-nonempty}
Soit $S = \lim S_i$ la limite d'un système projectif filtrant de schémas
dont les morphismes de transition sont affines
(Lemme \ref{limits-lemma-directed-inverse-system-has-limit}). Si tous les schémas
$S_i$ sont non vides et quasi-compacts, alors la limite
$S = \lim_i S_i$ est non vide.
\end{lemma}

\begin{proof}
Choisissons $0 \in I$. Remarquons que $I$ est non vide puisqu'il est
filtrant. Choisissons un recouvrement ouvert affine
$S_0 = \bigcup_{j = 1, \ldots, m} U_j$. Comme $I$ est filtrant, il existe
$j \in \{1, \ldots, m\}$ tel que
$f_{i0}^{-1}(U_j) \not = \emptyset$ pour tout $i \geq 0$. Par suite,
$\lim_{i \geq 0} f_{i0}^{-1}(U_j)$ n'est pas vide, car une limite
inductive filtrante d'anneaux non nuls est non nulle (puisque
$1 \not = 0$). Comme $\lim_{i \geq 0} f_{i0}^{-1}(U_j)$ est un sous-schéma
ouvert de la limite, on conclut.
\end{proof}

\begin{lemma}
\label{limits-lemma-inverse-limit-irreducibles}
Soit $S = \lim S_i$ la limite d'un système projectif filtrant de schémas
dont les morphismes de transition sont affines
(Lemme \ref{limits-lemma-directed-inverse-system-has-limit}). Soit $s \in S$
d'images $s_i \in S_i$. Alors
\begin{enumerate}
\item $s = \lim s_i$ comme schémas, c'est-à-dire
$\kappa(s) = \colim \kappa(s_i)$,
\item $\overline{\{s\}} = \lim \overline{\{s_i\}}$ comme ensembles, et
\item $\overline{\{s\}} = \lim \overline{\{s_i\}}$ comme schémas, où
$\overline{\{s\}}$ et $\overline{\{s_i\}}$ sont munis de leur structure
induite de schéma réduit.
\end{enumerate}
\end{lemma}

\begin{proof}
Choisissons $0 \in I$ et un recouvrement ouvert affine
$S_0 = \bigcup_{j \in J} U_{0, j}$. Pour $i \geq 0$, posons
$U_{i, j} = f_{i, 0}^{-1}(U_{0, j})$ et
$U_j = f_0^{-1}(U_{0, j})$. Ici, $f_{i'i} : S_{i'} \to S_i$ est le
morphisme de transition et $f_i : S \to S_i$ la projection. Pour
$j \in J$, les conditions suivantes sont équivalentes :
(a) $s \in U_j$, (b) $s_0 \in U_{0, j}$,
(c) $s_i \in U_{i, j}$ pour tout $i \geq 0$. Soit $J' \subset J$
l'ensemble des indices pour lesquels (a), (b) et (c) sont vérifiées.
Alors $\overline{\{s\}} = \bigcup_{j \in J'} (\overline{\{s\}} \cap U_j)$,
et de même pour $\overline{\{s_i\}}$ lorsque $i \geq 0$. Remarquons que
$\overline{\{s\}} \cap U_j$ est l'adhérence de $\{s\}$ dans l'espace
topologique $U_j$. Il en va de même de
$\overline{\{s_i\}} \cap U_{i, j}$ pour $i \geq 0$. Il suffit donc de
démontrer le lemme lorsque $S$ et $S_i$ sont affines pour tout $i$. Cela nous
ramène à la question algébrique du paragraphe suivant.

\medskip\noindent
Soit donné un système d'anneaux $(A_i, \varphi_{ii'})$ indexé par $I$.
Posons $A = \colim_i A_i$, avec les homomorphismes canoniques
$\varphi_i : A_i \to A$. Soit $\mathfrak p \subset A$ un idéal premier et
posons $\mathfrak p_i = \varphi_i^{-1}(\mathfrak p)$. Alors
$$
V(\mathfrak p) = \lim_i V(\mathfrak p_i)
$$
Cela résulte du Lemme \ref{limits-lemma-inverse-limit-sets}, puisque
$A/\mathfrak p = \colim A_i/\mathfrak p_i$. Cette égalité d'anneaux montre
aussi la dernière assertion concernant les structures induites de schéma
réduit. L'égalité $\kappa(\mathfrak p) = \colim \kappa(\mathfrak p_i)$
résulte également de cette assertion.
\end{proof}

\noindent
Dans la suite de cette section, nous nous plaçons dans la situation suivante.

\begin{situation}
\label{limits-situation-descent}
Soit $S = \lim_{i \in I} S_i$ la limite d'un système projectif filtrant de schémas
dont les morphismes de transition $f_{i'i} : S_{i'} \to S_i$ sont affines
(Lemme \ref{limits-lemma-directed-inverse-system-has-limit}). Nous supposons que
$S_i$ est quasi-compact et quasi-séparé pour tout $i \in I$. Nous notons
$f_i : S \to S_i$ la projection. Nous choisissons aussi un élément
$0 \in I$.
\end{situation}

\noindent
Dans cette situation, le morphisme $S \to S_0$ est affine. Il s'ensuit que
$S$ est quasi-compact et quasi-séparé\footnote{Cela résulte de Morphismes,
Lemme \ref{morphisms-lemma-affine-separated}, Topologie,
Définition \ref{topology-definition-quasi-compact}, et Schémas,
Lemme \ref{schemes-lemma-separated-permanence}.}. Le type de résultat que
nous recherchons est le suivant : si l'on dispose d'un objet sur $S$, alors,
pour un certain $i$, il existe un objet analogue sur $S_i$.

\begin{lemma}
\label{limits-lemma-topology-limit}
Dans la Situation \ref{limits-situation-descent}.
\begin{enumerate}
\item On a $S_{set} = \lim_i S_{i, set}$, où $S_{set}$ désigne l'ensemble
sous-jacent au schéma $S$.
\item On a $S_{top} = \lim_i S_{i, top}$, où $S_{top}$ désigne l'espace
topologique sous-jacent au schéma $S$.
\item Si $s, s' \in S$ et si $s'$ n'est pas une spécialisation de $s$,
alors, pour un certain $i \in I$, l'image $s'_i \in S_i$ de $s'$ n'est pas
une spécialisation de l'image $s_i \in S_i$ de $s$.
\item Ajouter ici d'autres faits élémentaires sur la topologie de $S$.
(Exigence : tout ajout doit être facile dans le cas affine.)
\end{enumerate}
\end{lemma}

\begin{proof}
L'assertion (1) est un cas particulier du Lemme \ref{limits-lemma-inverse-limit-sets}.

\medskip\noindent
L'assertion (2) est un cas particulier du Lemme \ref{limits-lemma-inverse-limit-top}.

\medskip\noindent
L'assertion (3) est un cas particulier du
Lemme \ref{limits-lemma-inverse-limit-irreducibles}.
\end{proof}

\begin{lemma}
\label{limits-lemma-descend-section}
Dans la Situation \ref{limits-situation-descent}. Supposons que $\mathcal{F}_0$
soit un faisceau quasi-cohérent sur $S_0$. Posons
$\mathcal{F}_i = f_{i0}^*\mathcal{F}_0$ pour $i \geq 0$ et
$\mathcal{F} = f_0^*\mathcal{F}_0$. Alors
$$
\Gamma(S, \mathcal{F}) = \colim_{i \geq 0} \Gamma(S_i, \mathcal{F}_i)
$$
\end{lemma}

\begin{proof}
Écrivons $\mathcal{A}_j = f_{i0, *} \mathcal{O}_{S_i}$. C'est un faisceau
quasi-cohérent d'$\mathcal{O}_{S_0}$-algèbres (voir Morphismes,
Lemme \ref{morphisms-lemma-affine-equivalence-algebras}) et $S_i$ est le
spectre relatif de $\mathcal{A}_i$ sur $S_0$. Dans la démonstration du
Lemme \ref{limits-lemma-directed-inverse-system-has-limit}, nous avons construit
$S$ comme le spectre relatif de
$\mathcal{A} = \colim_{i \geq 0} \mathcal{A}_i$ sur $S_0$. Posons
$$
\mathcal{M}_i = \mathcal{F}_0 \otimes_{\mathcal{O}_{S_0}} \mathcal{A}_i
$$
et
$$
\mathcal{M} = \mathcal{F}_0 \otimes_{\mathcal{O}_{S_0}} \mathcal{A}.
$$
On a alors $f_{i0, *} \mathcal{F}_i = \mathcal{M}_i$ et
$f_{0, *}\mathcal{F} = \mathcal{M}$. Comme $\mathcal{A}$ est la limite
inductive des faisceaux $\mathcal{A}_i$ et que le produit tensoriel commute
aux limites inductives filtrantes, on en déduit
$\mathcal{M} = \colim_{i \geq 0} \mathcal{M}_i$. Puisque $S_0$ est
quasi-compact et quasi-séparé, on obtient
\begin{eqnarray*}
\Gamma(S, \mathcal{F})
& = &
\Gamma(S_0, \mathcal{M}) \\
& = &
\Gamma(S_0, \colim_{i \geq 0} \mathcal{M}_i) \\
& = &
\colim_{i \geq 0} \Gamma(S_0, \mathcal{M}_i) \\
& = &
\colim_{i \geq 0} \Gamma(S_i, \mathcal{F}_i)
\end{eqnarray*}
voir Faisceaux, Lemme \ref{sheaves-lemma-directed-colimits-sections}, et
Topologie, Lemme \ref{topology-lemma-topology-quasi-separated-scheme}, pour
l'égalité intermédiaire.
\end{proof}

\begin{lemma}
\label{limits-lemma-limit-closed-nonempty}
Dans la Situation \ref{limits-situation-descent}. Supposons donné, pour chaque $i$,
un fermé non vide $Z_i \subset S_i$ tel que
$f_{i'i}(Z_{i'}) \subset Z_i$ pour tout $i' \geq i$. Alors il existe un
point $s \in S$ tel que $f_i(s) \in Z_i$ pour tout $i$.
\end{lemma}

\begin{proof}
Notons encore $Z_i \subset S_i$ le sous-schéma fermé réduit associé à
$Z_i$ ; voir Schémas,
Définition \ref{schemes-definition-reduced-induced-scheme}. Une immersion
fermée est affine et une composée de morphismes affines est affine (voir
Morphismes, Lemmes \ref{morphisms-lemma-closed-immersion-affine} et
\ref{morphisms-lemma-composition-affine}) ; par conséquent,
$Z_{i'} \to S_i$ est affine lorsque $i' \geq i$. On en déduit que le morphisme
$f_{i'i} : Z_{i'} \to Z_i$ est affine par
Morphismes, Lemme \ref{morphisms-lemma-affine-permanence}. Chacun des schémas $Z_i$ est
quasi-compact, car c'est un sous-schéma fermé d'un schéma quasi-compact.
Nous pouvons donc appliquer le Lemme \ref{limits-lemma-limit-nonempty} pour voir
que $Z = \lim_i Z_i$ est non vide. Comme il existe un morphisme canonique
$Z \to S$, on conclut.
\end{proof}

\begin{lemma}
\label{limits-lemma-limit-fibre-product-empty}
Dans la Situation \ref{limits-situation-descent}. Supposons donnés un $i$ et un
morphisme $T \to S_i$ tels que
\begin{enumerate}
\item $T \times_{S_i} S = \emptyset$, et
\item $T$ est quasi-compact.
\end{enumerate}
Alors $T \times_{S_i} S_{i'} = \emptyset$ pour tout $i'$ assez grand.
\end{lemma}

\begin{proof}
D'après le Lemme \ref{limits-lemma-scheme-over-limit}, on a
$T \times_{S_i} S = \lim_{i' \geq i} T \times_{S_i} S_{i'}$. Le résultat
découle donc du Lemme \ref{limits-lemma-limit-nonempty}.
\end{proof}

\begin{lemma}
\label{limits-lemma-limit-contained-in-constructible}
Dans la Situation \ref{limits-situation-descent}. Supposons donnés un $i$ et une
partie localement constructible $E \subset S_i$ telle que
$f_i(S) \subset E$. Alors $f_{i'i}(S_{i'}) \subset E$ pour tout $i'$ assez
grand.
\end{lemma}

\begin{proof}
En écrivant $S_i$ comme réunion finie de sous-schémas ouverts affines, on
se ramène au cas où $S_i$ est affine et $E$ constructible ; voir le
Lemme \ref{limits-lemma-directed-inverse-system-has-limit} et Propriétés,
Lemme \ref{properties-lemma-locally-constructible}. Dans ce cas, le
complémentaire $S_i \setminus E$ est lui aussi constructible. Il existe donc
un schéma affine $T$ et un morphisme $T \to S_i$ dont l'image est
$S_i \setminus E$ ; voir Algèbre,
Lemme \ref{algebra-lemma-constructible-is-image}. D'après le
Lemme \ref{limits-lemma-limit-fibre-product-empty}, le schéma
$T \times_{S_i} S_{i'}$ est vide pour tout $i'$ assez grand, et par
conséquent $f_{i'i}(S_{i'}) \subset E$ pour tout $i'$ assez grand.
\end{proof}

\begin{lemma}
\label{limits-lemma-descend-opens}
Dans la Situation \ref{limits-situation-descent}, on a les assertions suivantes :
\begin{enumerate}
\item Pour tout ouvert quasi-compact $V \subset S = \lim_i S_i$, il existe
$i \in I$ et un ouvert quasi-compact $V_i \subset S_i$ tels que
$f_i^{-1}(V_i) = V$.
\item Étant donnés des ouverts quasi-compacts $V_i \subset S_i$ et
$V_{i'} \subset S_{i'}$ tels que
$f_i^{-1}(V_i) = f_{i'}^{-1}(V_{i'})$, il existe un indice
$i'' \geq i, i'$ tel que
$f_{i''i}^{-1}(V_i) = f_{i''i'}^{-1}(V_{i'})$.
\item Si $V_{1, i}, \ldots, V_{n, i} \subset S_i$ sont des ouverts
quasi-compacts et si
$S = f_i^{-1}(V_{1, i}) \cup \ldots \cup f_i^{-1}(V_{n, i})$, alors
$S_{i'} = f_{i'i}^{-1}(V_{1, i}) \cup \ldots \cup f_{i'i}^{-1}(V_{n, i})$
pour un certain $i' \geq i$.
\end{enumerate}
\end{lemma}

\begin{proof}
Choisissons $i_0 \in I$. Remarquons que $I$ est non vide puisqu'il est
filtrant. Pour simplifier les notations, écrivons $S_0 = S_{i_0}$ et
$i_0 = 0$. Choisissons un recouvrement ouvert affine
$S_0 = U_{1, 0} \cup \ldots \cup U_{m, 0}$. Notons
$U_{j, i} \subset S_i$ l'image réciproque de $U_{j, 0}$ par le morphisme de
transition pour $i \geq 0$. Notons $U_j$ l'image réciproque de $U_{j, 0}$ dans $S$.
Remarquons que $U_j = \lim_i U_{j, i}$ est une limite de schémas affines.

\medskip\noindent
Nous démontrons d'abord l'assertion d'unicité. Soient des ouverts
quasi-compacts $V_i \subset S_i$ et $V_{i'} \subset S_{i'}$ tels que
$f_i^{-1}(V_i) = f_{i'}^{-1}(V_{i'})$. Il suffit de montrer que
$f_{i''i}^{-1}(V_i \cap U_{j, i''})$ et
$f_{i''i'}^{-1}(V_{i'} \cap U_{j, i''})$ deviennent égaux pour $i''$ assez
grand. On se ramène ainsi au cas d'une limite de schémas affines. Dans ce
cas, écrivons $S = \Spec(R)$ et $S_i = \Spec(R_i)$ pour tout $i \in I$.
On peut écrire $V_i = S_i \setminus V(h_1, \ldots, h_m)$ et
$V_{i'} = S_{i'} \setminus V(g_1, \ldots, g_n)$. L'hypothèse signifie que
les idéaux $\sum g_jR$ et $\sum h_jR$ ont le même radical dans $R$.
Cela signifie que $g_j^N = \sum a_{jj'}h_{j'}$ et
$h_j^N = \sum b_{jj'} g_{j'}$ pour un certain $N \gg 0$ et certains
$a_{jj'}$, $b_{jj'}$ dans $R$. Puisque $R = \colim_i R_i$, nous pouvons
choisir un indice $i'' \geq i$ tel que les égalités
$g_j^N = \sum a_{jj'}h_{j'}$ et
$h_j^N = \sum b_{jj'} g_{j'}$ soient vérifiées dans $R_{i''}$ pour certains
$a_{jj'}$ et $b_{jj'}$ dans $R_{i''}$. Il en résulte que les idéaux
$\sum g_jR_{i''}$ et $\sum h_jR_{i''}$ ont le même radical dans $R_{i''}$,
comme voulu.

\medskip\noindent
Montrons l'existence. Si $S_0$ est affine, alors
$S_i = \Spec(R_i)$ pour tout $i \geq 0$ et $S = \Spec(R)$, où
$R = \colim R_i$. On a alors
$V = S \setminus V(g_1, \ldots, g_n)$ pour certains
$g_1, \ldots, g_n \in R$. Choisissons $i$ assez grand pour que chacun des
$g_j$ provienne d'un élément $g_{j, i} \in R_i$, et prenons
$V_i = S_i \setminus V(g_{1, i}, \ldots, g_{n, i})$. Pour $S_0$ quelconque,
les ouverts $V \cap U_j$ sont quasi-compacts puisque $S$ est quasi-séparé.
Le cas affine montre donc que, pour chaque $j = 1, \ldots, m$, il existe
$i_j \in I$ et un ouvert quasi-compact $V_{i_j} \subset U_{j, i_j}$ dont
l'image réciproque dans $U_j$ est $V \cap U_j$. Posons
$i = \max(i_1, \ldots, i_m)$ et
$V_i = \bigcup f_{ii_j}^{-1}(V_{i_j})$.

\medskip\noindent
L'assertion sur les recouvrements découle de l'assertion d'unicité appliquée
aux ouverts $V_{1, i} \cup \ldots \cup V_{n, i}$ et $S_i$ de $S_i$.
\end{proof}

\begin{lemma}
\label{limits-lemma-limit-quasi-affine}
Dans la Situation \ref{limits-situation-descent}, si $S$ est quasi-affine, alors,
pour un certain $i_0 \in I$, les schémas $S_i$ sont quasi-affines pour tout
$i \geq i_0$.
\end{lemma}

\begin{proof}
Choisissons $i_0 \in I$. Remarquons que $I$ est non vide puisqu'il est
filtrant. Pour simplifier les notations, écrivons $S_0 = S_{i_0}$ et
$i_0 = 0$. Soit $s \in S$. Nous pouvons choisir un ouvert affine
$U_0 \subset S_0$ contenant $f_0(s)$. Puisque $S$ est quasi-affine, nous
pouvons choisir $a \in \Gamma(S, \mathcal{O}_S)$ tel que
$s \in D(a) \subset f_0^{-1}(U_0)$ et tel que $D(a)$ soit affine. D'après le
Lemme \ref{limits-lemma-descend-section}, il existe $i \geq 0$ tel que $a$ provienne
d'un élément $a_i \in \Gamma(S_i, \mathcal{O}_{S_i})$. Pour tout indice
$j \geq i$, notons $a_j$ l'image de $a_i$ dans les sections globales du
faisceau structural de $S_j$. Considérons les ouverts $D(a_j) \subset S_j$
et $U_j = f_{j0}^{-1}(U_0)$. Remarquons que $U_j$ est affine et que
$D(a_j)$ est un ouvert quasi-compact de $S_j$ ; voir par exemple Propriétés,
Lemme \ref{properties-lemma-affine-cap-s-open}. Nous pouvons donc appliquer
le Lemme \ref{limits-lemma-descend-opens} aux ouverts $U_j$ et
$U_j \cup D(a_j)$ pour conclure que $D(a_j) \subset U_j$ pour un certain
$j \geq i$. Pour un tel indice $j$, $D(a_j) \subset S_j$ est un ouvert
affine (car $D(a_j)$ est un ouvert affine standard de l'ouvert affine $U_j$)
contenant l'image $f_j(s)$.

\medskip\noindent
Nous en déduisons que, pour tout $s \in S$, il existe un indice $i \in I$
et une section globale $a \in \Gamma(S_i, \mathcal{O}_{S_i})$ tels que
$D(a) \subset S_i$ soit un ouvert affine contenant $f_i(s)$. Comme $S$ est
quasi-compact, nous pouvons choisir un seul indice $i \in I$ et des sections
globales $a_1, \ldots, a_m \in \Gamma(S_i, \mathcal{O}_{S_i})$ tels que
chaque $D(a_j) \subset S_i$ soit un ouvert affine et que l'image de
$f_i : S \to S_i$ soit contenue dans la réunion
$W_i = \bigcup_{j = 1, \ldots, m} D(a_j)$. Pour $i' \geq i$, posons
$W_{i'} = f_{i'i}^{-1}(W_i)$. Puisque $f_i^{-1}(W_i)$ est égal à $S$, une nouvelle
application du Lemme \ref{limits-lemma-descend-opens} montre que
pour un $i' \geq i$ convenable, on a $S_{i'} = W_{i'}$. Nous pouvons donc remplacer
$i$ par $i'$ et supposer
$S_i = \bigcup_{j = 1, \ldots, m} D(a_j)$. Il en résulte que
$\mathcal{O}_{S_i}$ est un faisceau inversible ample sur $S_i$ (voir
Propriétés, Définition \ref{properties-definition-ample}) et donc que $S_i$
est quasi-affine ; voir Propriétés,
Lemme \ref{properties-lemma-quasi-affine-O-ample}. On conclut.
\end{proof}

\begin{lemma}
\label{limits-lemma-limit-affine}
Dans la Situation \ref{limits-situation-descent}, si $S$ est affine, alors, pour
un certain $i_0 \in I$, les schémas $S_i$ sont affines pour tout
$i \geq i_0$.
\end{lemma}

\begin{proof}
D'après le Lemme \ref{limits-lemma-limit-quasi-affine}, nous pouvons supposer que
$S_0$ est quasi-affine pour un certain $0 \in I$. Posons
$R_0 = \Gamma(S_0, \mathcal{O}_{S_0})$. Alors $S_0$ est un ouvert
quasi-compact de $T_0 = \Spec(R_0)$. Notons $j_0 : S_0 \to T_0$
l'immersion ouverte quasi-compacte correspondante. Pour $i \geq 0$, posons
$\mathcal{A}_i = f_{i0, *}\mathcal{O}_{S_i}$. Comme $f_{i0}$ est affine,
on a $S_i = \underline{\Spec}_{S_0}(\mathcal{A}_i)$. Posons
$T_i = \underline{\Spec}_{T_0}(j_{0, *}\mathcal{A}_i)$. Alors
$T_i \to T_0$ est affine, donc $T_i$ est affine. Ainsi, $T_i$ est le spectre
de
$$
R_i = \Gamma(T_0, j_{0, *}\mathcal{A}_i) = \Gamma(S_0, \mathcal{A}_i) =
\Gamma(S_i, \mathcal{O}_{S_i}).
$$
Écrivons $S = \Spec(R)$. On a $R = \colim_i R_i$
d'après le Lemme \ref{limits-lemma-descend-section}. On a donc aussi $S = \lim_i T_i$. Comme la formation du
spectre relatif commute au changement de base, l'image réciproque de
l'ouvert $S_0 \subset T_0$ dans $T_i$ est $S_i$. Soit
$Z_0 = T_0 \setminus S_0$, et soit $Z_i \subset T_i$ l'image réciproque de
$Z_0$. Puisque $S_i = T_i \setminus Z_i$, il suffit de montrer que $Z_i$ est
vide pour un certain $i$. Supposons, en vue d'une contradiction, que $Z_i$
soit non vide pour tout $i$. D'après le
Lemme \ref{limits-lemma-limit-closed-nonempty}, il existe un point $s$ de
$S = \lim T_i$ dont l'image appartient à $Z_i$ pour tout $i$. Mais
$S = \lim_i S_i$, ce qui contredit le Lemme \ref{limits-lemma-topology-limit}.
\end{proof}

\begin{lemma}
\label{limits-lemma-limit-separated}
Dans la Situation \ref{limits-situation-descent}, si $S$ est séparé, alors, pour un
certain $i_0 \in I$, les schémas $S_i$ sont séparés pour tout $i \geq i_0$.
\end{lemma}

\begin{proof}
Choisissons un recouvrement ouvert affine fini
$S_0 = U_{0, 1} \cup \ldots \cup U_{0, m}$. Notons
$U_{i, j} \subset S_i$ et $U_j \subset S$ les images réciproques de
$U_{0, j}$. Remarquons que $U_{i, j}$ et $U_j$ sont affines. Comme $S$ est
séparé, les intersections $U_{j_1} \cap U_{j_2}$ sont affines. Puisque
$U_{j_1} \cap U_{j_2} = \lim_{i \geq 0} U_{i, j_1} \cap U_{i, j_2}$,
on voit que $U_{i, j_1} \cap U_{i, j_2}$ est affine pour $i$ assez grand,
d'après le Lemme \ref{limits-lemma-limit-affine}. Pour montrer
que $S_i$ est séparé lorsque $i$ est assez grand, il suffit maintenant de
montrer que
$$
\mathcal{O}_{S_i}(U_{i, j_1})
\otimes_{\mathcal{O}_S(S)}
\mathcal{O}_{S_i}(U_{i, j_2})
\longrightarrow
\mathcal{O}_{S_i}(U_{i, j_1} \cap U_{i, j_2})
$$
est surjectif pour $i$ assez grand (Schémas,
Lemme \ref{schemes-lemma-characterize-separated}).

\medskip\noindent
Pour alléger les indices, supposons donnés des ouverts affines
$U, V \subset S_0$ tels que $U \cap V$ soit également affine. Notons
$U_i, V_i \subset S_i$, resp.\ $U, V \subset S$, leurs images
réciproques. Il faut montrer que
$\mathcal{O}(U_i) \otimes \mathcal{O}(V_i) \to
\mathcal{O}(U_i \cap V_i)$
est surjectif pour $i$ assez grand, sachant que
$\mathcal{O}(U) \otimes \mathcal{O}(V) \to \mathcal{O}(U \cap V)$
est surjectif. Remarquons que
$\mathcal{O}(U_0) \otimes \mathcal{O}(V_0) \to
\mathcal{O}(U_0 \cap V_0)$
est de type fini, car le morphisme diagonal $S_i \to S_i \times S_i$ est une
immersion (Schémas, Lemme \ref{schemes-lemma-diagonal-immersion}), donc est
localement de type fini (Morphismes,
Lemmes \ref{morphisms-lemma-locally-finite-type-characterize} et
\ref{morphisms-lemma-immersion-locally-finite-type}). Nous pouvons donc
choisir des éléments
$f_{0, 1}, \ldots, f_{0, n} \in \mathcal{O}(U_0 \cap V_0)$
qui engendrent $\mathcal{O}(U_0 \cap V_0)$ sur
$\mathcal{O}(U_0) \otimes \mathcal{O}(V_0)$. Observons que, pour $i \geq 0$,
le diagramme de schémas
$$
\xymatrix{
U_i \cap V_i \ar[r] \ar[d] & U_i \ar[d] \\
U_0 \cap V_0 \ar[r] & U_0
}
$$
est cartésien. Nous voyons donc que les images
$f_{i, 1}, \ldots, f_{i, n} \in \mathcal{O}(U_i \cap V_i)$
engendrent $\mathcal{O}(U_i \cap V_i)$ sur
$\mathcal{O}(U_i) \otimes \mathcal{O}(V_0)$ et, a fortiori, sur
$\mathcal{O}(U_i) \otimes \mathcal{O}(V_i)$. Par hypothèse, les images
$f_1, \ldots, f_n \in \mathcal{O}(U \otimes V)$ appartiennent à l'image de
l'application
$\mathcal{O}(U) \otimes \mathcal{O}(V) \to \mathcal{O}(U \cap V)$.
Puisque
$\mathcal{O}(U) \otimes \mathcal{O}(V) =
\colim \mathcal{O}(U_i) \otimes \mathcal{O}(V_i)$
nous voyons qu'elles appartiennent à l'image de l'application à un certain
niveau fini, ce qui démontre le lemme.
\end{proof}

\begin{lemma}
\label{limits-lemma-limit-ample}
Dans la Situation \ref{limits-situation-descent}, soit $\mathcal{L}_0$ un faisceau
inversible de modules sur $S_0$. Si son image réciproque $\mathcal{L}$ sur
$S$ est ample, alors, pour un certain $i \in I$, son image réciproque
$\mathcal{L}_i$ sur $S_i$ est ample.
\end{lemma}

\begin{proof}
L'hypothèse signifie qu'il existe un nombre fini de sections
$s_1, \ldots, s_m \in \Gamma(S, \mathcal{L})$ telles que $S_{s_j}$ soit
affine et que $S = \bigcup S_{s_j}$ ; voir Propriétés,
Définition \ref{properties-definition-ample}. D'après le
Lemme \ref{limits-lemma-descend-section}, nous pouvons trouver $i \in I$ et des
sections $s_{i, j} \in \Gamma(S_i, \mathcal{L}_i)$ d'images $s_j$.
Après avoir augmenté $i$, le Lemme \ref{limits-lemma-limit-affine} permet de
supposer que $(S_i)_{s_{i, j}}$ est affine pour $j = 1, \ldots, m$.
Après avoir augmenté $i$ une dernière fois, le
Lemme \ref{limits-lemma-descend-opens} permet de supposer que
$S_i = \bigcup (S_i)_{s_{i, j}}$. Alors $\mathcal{L}_i$ est ample par
définition.
\end{proof}

\begin{lemma}
\label{limits-lemma-finite-type-eventually-closed}
Soit $S$ un schéma. Soit $X = \lim X_i$ une limite projective de schémas sur
$S$ dont les morphismes de transition sont affines. Soit $Y \to X$ un
morphisme de schémas sur $S$.
\begin{enumerate}
\item Si $Y \to X$ est une immersion fermée, si $X_i$ est quasi-compact et
si $Y$ est localement de type fini sur $S$, alors $Y \to X_i$ est une
immersion fermée pour $i$ assez grand.
\item Si $Y \to X$ est une immersion, si $X_i$ est quasi-séparé, si
$Y \to S$ est localement de type fini et si $Y$ est quasi-compact, alors
$Y \to X_i$ est une immersion pour $i$ assez grand.
\item Si $Y \to X$ est un isomorphisme, si $X_i$ est quasi-compact, si
$X_i \to S$ est localement de type fini, si les morphismes de transition
$X_{i'} \to X_i$ sont des immersions fermées et si $Y \to S$ est localement
de présentation finie, alors $Y \to X_i$ est un isomorphisme pour $i$ assez
grand.
\end{enumerate}
\end{lemma}

\begin{proof}
Démonstration de (1). Choisissons $0 \in I$ et un recouvrement ouvert affine
fini $X_0 = U_{0, 1} \cup \ldots \cup U_{0, m}$ tel que $U_{0, j}$ s'envoie
dans un ouvert affine $W_j \subset S$. Soient $V_j \subset Y$,
resp.\ $U_{i, j} \subset X_i$, $i \geq 0$, respectivement
$U_j \subset X$, les images réciproques de $U_{0, j}$. Il suffit de montrer
que $V_j \to U_{i, j}$ est une immersion fermée pour $i$ assez grand, et
nous savons que $V_j \to U_j$ est une immersion fermée. On se ramène ainsi
au fait algébrique suivant : si $A = \colim A_i$ est une limite inductive
filtrante de $R$-algèbres, si $A \to B$ est une surjection de $R$-algèbres
et si $B$ est une $R$-algèbre de type fini, alors $A_i \to B$ est surjectif
pour $i$ assez grand.

\medskip\noindent
Démonstration de (2). Choisissons $0 \in I$. Choisissons un ouvert
quasi-compact $X'_0 \subset X_0$ tel que $Y \to X_0$ se factorise par $X'_0$.
Après avoir remplacé $X_i$ par l'image réciproque de
$X'_0$ pour $i \geq 0$, nous pouvons supposer tous les $X_i'$ quasi-compacts et
quasi-séparés. Soit $U \subset X$ un ouvert quasi-compact tel que
$Y \to X$ se factorise par une immersion fermée $Y \to U$ (un tel $U$
existe puisque $Y$ est quasi-compact). D'après le
Lemme \ref{limits-lemma-descend-opens}, nous pouvons supposer
$U = \lim U_i$, avec $U_i \subset X_i$ ouvert quasi-compact. L'assertion (1)
montre que $Y \to U_i$ est une immersion fermée pour un certain $i$. Ainsi,
(2) est vérifiée.

\medskip\noindent
Démonstration de (3). En travaillant localement sur des ouverts affines de
$X_0$, pour un certain $0 \in I$, comme dans la démonstration de (1), on se
ramène au fait algébrique suivant : si $A = \lim A_i$ est une limite
inductive filtrante de $R$-algèbres dont les homomorphismes de transition
sont surjectifs, et si $A$ est de présentation finie sur $A_0$, alors
$A = A_i$ pour un certain $i$. En effet, écrivons
$A = A_0/(f_1, \ldots, f_n)$. Choisissons $i$ tel que
$f_1, \ldots, f_n$ aient une image nulle par l'homomorphisme surjectif
$A_0 \to A_i$.
\end{proof}

\begin{lemma}
\label{limits-lemma-eventually-separated}
Soit $S$ un schéma. Soit $X = \lim X_i$ une limite projective de schémas sur
$S$ dont les morphismes de transition sont affines. Supposons que
\begin{enumerate}
\item $S$ soit quasi-séparé,
\item $X_i$ soit quasi-compact et quasi-séparé,
\item $X \to S$ soit séparé.
\end{enumerate}
Alors $X_i \to S$ est séparé pour tout $i$ assez grand.
\end{lemma}

\begin{proof}
Soit $0 \in I$. Remarquons que $I$ est non vide puisqu'il est filtrant.
Comme $X_0$ est quasi-compact, nous pouvons trouver un nombre fini d'ouverts
affines $U_1, \ldots, U_n \subset S$ tels que l'image de $X_0 \to S$ soit
contenue dans $U_1 \cup \ldots \cup U_n$. Notons $h_i : X_i \to S$ le
morphisme structural. Il suffit de vérifier que, pour un certain $i \geq 0$,
les morphismes $h_i^{-1}(U_j) \to U_j$ sont séparés pour
$j = 1, \ldots,  n$. Puisque $S$ est quasi-séparé, les morphismes
$U_j \to S$ sont quasi-compacts. Par conséquent, $h_i^{-1}(U_j)$ est
quasi-compact et quasi-séparé. On se ramène ainsi au cas où $S$ est affine.
Dans ce cas, il faut montrer que $X_i$ est séparé, sachant que $X$ l'est.
Le lemme résulte donc du Lemme \ref{limits-lemma-limit-separated}.
\end{proof}

\begin{lemma}
\label{limits-lemma-eventually-affine}
Soit $S$ un schéma. Soit $X = \lim X_i$ une limite projective de schémas sur
$S$ dont les morphismes de transition sont affines. Supposons que
\begin{enumerate}
\item $S$ soit quasi-compact et quasi-séparé,
\item $X_i$ soit quasi-compact et quasi-séparé,
\item $X \to S$ soit affine.
\end{enumerate}
Alors $X_i \to S$ est affine pour $i$ assez grand.
\end{lemma}

\begin{proof}
Choisissons un recouvrement ouvert affine fini
$S = \bigcup_{j = 1, \ldots, n} V_j$. Notons $f : X \to S$ et
$f_i : X_i \to S$ les morphismes structuraux. Pour chaque $j$, le schéma
$f^{-1}(V_j) = \lim_i f_i^{-1}(V_j)$ est affine (car un morphisme fini est
affine par définition). D'après le Lemme \ref{limits-lemma-limit-affine}, il existe
donc $i \in I$ tel que chaque $f_i^{-1}(V_j)$ soit affine. Autrement dit,
$f_i : X_i \to S$ est affine pour $i$ assez grand ; voir Morphismes,
Lemme \ref{morphisms-lemma-characterize-affine}.
\end{proof}

\begin{lemma}
\label{limits-lemma-eventually-finite}
Soit $S$ un schéma. Soit $X = \lim X_i$ une limite projective de schémas sur
$S$ dont les morphismes de transition sont affines. Supposons que
\begin{enumerate}
\item $S$ soit quasi-compact et quasi-séparé,
\item $X_i$ soit quasi-compact et quasi-séparé,
\item les morphismes de transition $X_{i'} \to X_i$ soient finis,
\item $X_i \to S$ soit localement de type fini,
\item $X \to S$ soit entier.
\end{enumerate}
Alors $X_i \to S$ est fini pour $i$ assez grand.
\end{lemma}

\begin{proof}
D'après le Lemme \ref{limits-lemma-eventually-affine}, nous pouvons supposer que
$X_i \to S$ est affine pour tout $i$. Choisissons un recouvrement ouvert
affine fini $S = \bigcup_{j = 1, \ldots, n} V_j$. Notons $f : X \to S$ et
$f_i : X_i \to S$ les morphismes structuraux. Il suffit de montrer qu'il
existe $i$ tel que $f_i^{-1}(V_j)$ soit fini sur $V_j$ pour
$j = 1, \ldots, m$ (Morphismes,
Lemme \ref{morphisms-lemma-finite-local}). En effet, pour $i' \geq i$, la
composée $X_{i'} \to X_i \to S$ sera finie en tant que composée de morphismes
finis (Morphismes, Lemme \ref{morphisms-lemma-composition-finite}). On se
ramène ainsi au cas affine. Soit $R$ un anneau et soit $A = \colim A_i$, avec
$R \to A$ entier et $A_i \to A_{i'}$ fini pour tout $i \leq i'$. De plus,
$R \to A_i$ est de type fini pour tout $i$. Il faut montrer que $A_i$ est
fini sur $R$ pour un certain $i$. Pour cela, choisissons $i \in I$ et des
générateurs $x_1, \ldots, x_m \in A_i$ de $A_i$ comme $R$-algèbre. Puisque
$A$ est entier sur $R$, nous pouvons trouver des polynômes unitaires
$P_j \in R[T]$ tels que $P_j(x_j) = 0$ dans $A$. Il existe donc
$i' \geq i$ tel que $P_j(x_j) = 0$ dans $A_{i'}$ pour
$j = 1, \ldots, m$. Alors l'image $A'_i$ de $A_i$ dans $A_{i'}$ est finie
sur $R$ d'après Algèbre,
Lemme \ref{algebra-lemma-characterize-finite-in-terms-of-integral}. Comme
$A'_i \subset A_{i'}$ est également fini, on en déduit qu'$A_{i'}$ est fini
sur $R$ par Algèbre, Lemme \ref{algebra-lemma-finite-transitive}.
\end{proof}

\begin{lemma}
\label{limits-lemma-eventually-closed-immersion}
Soit $S$ un schéma. Soit $X = \lim X_i$ une limite projective de schémas sur
$S$ dont les morphismes de transition sont affines. Supposons que
\begin{enumerate}
\item $S$ soit quasi-compact et quasi-séparé,
\item $X_i$ soit quasi-compact et quasi-séparé,
\item les morphismes de transition $X_{i'} \to X_i$ soient des immersions
fermées,
\item $X_i \to S$ soit localement de type fini,
\item $X \to S$ soit une immersion fermée.
\end{enumerate}
Alors $X_i \to S$ est une immersion fermée pour $i$ assez grand.
\end{lemma}

\begin{proof}
D'après le Lemme \ref{limits-lemma-eventually-affine}, nous pouvons supposer que
$X_i \to S$ est affine pour tout $i$. Choisissons un recouvrement ouvert
affine fini $S = \bigcup_{j = 1, \ldots, n} V_j$. Notons $f : X \to S$ et
$f_i : X_i \to S$ les morphismes structuraux. Il suffit de montrer qu'il
existe $i$ tel que $f_i^{-1}(V_j)$ soit un sous-schéma fermé de $V_j$ pour
$j = 1, \ldots, m$ (Morphismes,
Lemme \ref{morphisms-lemma-closed-immersion}). On se ramène ainsi au cas
affine. Soit $R$ un anneau et soit $A = \colim A_i$, avec $R \to A$
surjectif et $A_i \to A_{i'}$ surjectif pour tout $i \leq i'$. De plus,
$R \to A_i$ est de type fini pour tout $i$. Il faut montrer que
$R \to A_i$ est surjectif pour un certain $i$. Pour cela, choisissons
$i \in I$ et des générateurs $x_1, \ldots, x_m \in A_i$ de $A_i$ comme
$R$-algèbre. Puisque $R \to A$ est surjectif, nous pouvons trouver
$r_j \in R$ tels que $r_j$ ait pour image $x_j$ dans $A$. Il existe donc $i' \geq i$ tel
que $r_j$ ait pour image celle de $x_j$ dans $A_{i'}$ pour
$j = 1, \ldots, m$. Comme $A_i \to A_{i'}$ est surjectif, il en résulte que
$R \to A_{i'}$ est surjectif.
\end{proof}

\begin{lemma}
\label{limits-lemma-eventually-immersion}
Soit $S$ un schéma. Soit $X = \lim X_i$ une limite projective de schémas sur
$S$ dont les morphismes de transition sont affines. Supposons que
\begin{enumerate}
\item $S$ soit quasi-séparé,
\item $X_i$ soit quasi-compact et quasi-séparé,
\item les morphismes de transition $X_{i'} \to X_i$ soient des immersions
fermées,
\item $X_i \to S$ soit localement de type fini, et
\item $X \to S$ soit une immersion.
\end{enumerate}
Alors $X_i \to S$ est une immersion pour $i$ assez grand.
\end{lemma}

\begin{proof}
Choisissons un sous-schéma ouvert $U \subset S$ tel que $X \to S$ soit la
composée d'une immersion fermée $X \to U$ et du morphisme d'inclusion
$U \to S$. Puisque $X$ est quasi-compact, nous pouvons rétrécir $U$ et
supposer $U$ quasi-compact. Notons $V_i \subset X_i$ l'image réciproque de $U$.
Comme l'image réciproque de $V_i$ est $X$, le
Lemme \ref{limits-lemma-descend-opens} montre que $V_i = X_i$ pour tout $i$ assez
grand. Nous pouvons donc supposer $X = \lim X_i$ dans la catégorie des
schémas sur $U$. Alors $X_i \to U$ est une immersion fermée pour $i$ assez grand,
d'après le Lemme \ref{limits-lemma-eventually-closed-immersion}. Cela
démontre le lemme.
\end{proof}










\section{Approximation noethérienne absolue}
\label{limits-section-approximation}

\noindent
Une bonne référence pour cette section est l'appendice C de l'article de
Thomason et Trobaugh \cite{TT}. Pour nos conventions sur les systèmes
filtrants, voir Catégories, section \ref{categories-section-posets-limits}.
Nous utiliserons sans autre mention le résultat d'existence et les propriétés
de la limite établis à la section \ref{limits-section-limits}.

\begin{lemma}
\label{limits-lemma-quasi-affine-finite-type-over-Z}
Soit $W$ un schéma quasi-affine de type fini sur $\mathbf{Z}$. Supposons que
$W \to \Spec(R)$ soit une immersion ouverte dans un schéma affine. Il existe
une $\mathbf{Z}$-algèbre de type fini $A \subset R$ qui induit une immersion
ouverte $W \to \Spec(A)$. De plus, $R$ est la limite inductive filtrante de
telles sous-algèbres.
\end{lemma}

\begin{proof}
Choisissons un recouvrement ouvert affine
$W = \bigcup_{i = 1, \ldots, n} W_i$ tel que chaque $W_i$ soit un ouvert
affine standard de $\Spec(R)$. Autrement dit, si l'on écrit
$W_i = \Spec(R_i)$, alors $R_i = R_{f_i}$ pour un certain $f_i \in R$.
Choisissons un nombre fini d'éléments $x_{ij} \in R_i$ qui engendrent $R_i$
sur $\mathbf{Z}$. Prenons $N \gg 0$ tel que chaque $f_i^Nx_{ij}$ provienne
d'un élément de $R$, disons $y_{ij} \in R$. Soit $A$ la
$\mathbf{Z}$-algèbre engendrée par les $f_i$ et les $y_{ij}$, ainsi que,
éventuellement, par un nombre fini d'éléments supplémentaires de $R$.
Alors $A$ convient. Les détails sont omis.
\end{proof}

\begin{lemma}
\label{limits-lemma-diagram}
Supposons donné un diagramme cartésien d'anneaux
$$
\xymatrix{
B \ar[r]_s & R \\
B'\ar[u] \ar[r] & R' \ar[u]_t
}
$$
Soit $W' \subset \Spec(R')$ un ouvert de la forme
$W' = D(f_1) \cup \ldots \cup D(f_n)$ tel que $t(f_i) = s(g_i)$ pour un
certain $g_i \in B$ et que $B_{g_i} \cong R_{s(g_i)}$. Alors $B' \to R'$
induit une immersion ouverte de $W'$ dans $\Spec(B')$.
\end{lemma}

\begin{proof}
Posons $h_i = (g_i, f_i) \in B'$. Compléments d'algèbre,
Lemme \ref{more-algebra-lemma-diagram-localize}, montre que
$(B')_{h_i} \cong (R')_{f_i}$, comme voulu.
\end{proof}

\noindent
Le lemme suivant donne un énoncé précis de l'approximation noethérienne.

\begin{lemma}
\label{limits-lemma-approximate}
Soit $S$ un schéma quasi-compact et quasi-séparé. Soit $V \subset S$ un
ouvert quasi-compact. Soit $I$ un ensemble ordonné filtrant et soit
$(V_i, f_{ii'})$ un système projectif de schémas indexé par $I$, à
morphismes de transition affines, chaque $V_i$ étant de type fini sur
$\mathbf{Z}$, et tel que $V = \lim V_i$. Il existe alors
\begin{enumerate}
\item un ensemble ordonné filtrant $J$,
\item un système projectif de schémas $(S_j, g_{jj'})$ indexé par $J$,
\item une application croissante $\alpha : J \to I$,
\item des sous-schémas ouverts $V'_j \subset S_j$, et
\item des isomorphismes $V'_j \to V_{\alpha(j)}$
\end{enumerate}
tels que
\begin{enumerate}
\item les morphismes de transition $g_{jj'} : S_j \to S_{j'}$ sont affines,
\item chaque $S_j$ est de type fini sur $\mathbf{Z}$,
\item $g_{jj'}^{-1}(V'_{j'}) = V'_j$,
\item $S = \lim S_j$ et $V = \lim V'_j$, et
\item les diagrammes
$$
\vcenter{
\xymatrix{
V \ar[d] \ar[rd] \\
V'_j \ar[r] & V_{\alpha(j)}
}
}
\quad\text{et}\quad
\vcenter{
\xymatrix{
V'_j \ar[r] \ar[d] & V_{\alpha(j)} \ar[d] \\
V'_{j'} \ar[r] & V_{\alpha(j')}
}
}
$$
sont commutatifs.
\end{enumerate}
\end{lemma}

\begin{proof}
Posons $Z = S \setminus V$. Choisissons des ouverts affines
$U_1, \ldots, U_m \subset S$ tels que
$Z \subset \bigcup_{l = 1, \ldots, m} U_l$. Considérons les ouverts
$$
V \subset V \cup U_1 \subset V \cup U_1 \cup U_2 \subset
\ldots \subset V \cup \bigcup\nolimits_{l = 1, \ldots, m} U_l = S
$$
Si nous pouvons démontrer le lemme successivement dans chacun des cas
$$
V \cup U_1 \cup \ldots \cup U_l
\subset
V \cup U_1 \cup \ldots \cup U_{l + 1}
$$
alors le lemme en résultera pour $V \subset S$. Dans chaque cas, nous
ajoutons un seul ouvert affine. Nous pouvons donc supposer que
\begin{enumerate}
\item $S = U \cup V$,
\item $U$ est un ouvert affine de $S$,
\item $V$ est un ouvert quasi-compact de $S$, et
\item $V = \lim_i V_i$, où $(V_i, f_{ii'})$ est un système projectif indexé
par un ensemble ordonné filtrant $I$, chaque $f_{ii'}$ étant affine et
chaque $V_i$ de type fini sur $\mathbf{Z}$.
\end{enumerate}
Notons $f_i : V \to V_i$ les projections. Posons $W = U \cap V$. Comme $S$
est quasi-séparé, c'est un ouvert quasi-compact de $V$. D'après le
Lemme \ref{limits-lemma-descend-opens} (et après avoir restreint $I$), nous pouvons
supposer qu'il existe des ouverts $W_i \subset V_i$ tels que
$f_{ii'}^{-1}(W_{i'}) = W_i$ et $f_i^{-1}(W_i) = W$. Puisque $W$ est un
ouvert quasi-compact de $U$, il est quasi-affine. Nous pouvons donc supposer,
après avoir restreint $I$ une nouvelle fois, que $W_i$ est quasi-affine pour
tout $i$ ; voir le Lemme \ref{limits-lemma-limit-quasi-affine}.

\medskip\noindent
Écrivons $U = \Spec(B)$. Posons $R = \Gamma(W, \mathcal{O}_W)$ et
$R_i = \Gamma(W_i, \mathcal{O}_{W_i})$. D'après le
Lemme \ref{limits-lemma-descend-section}, on a $R = \colim_i R_i$. Nous disposons
maintenant des homomorphismes d'anneaux
$$
\xymatrix{
B \ar[r]_s & R \\
& R_i \ar[u]_{t_i}
}
$$
Posons $B_i = \{(b, r) \in B \times R_i \mid s(b) = t_i(r)\}$, de sorte que
nous ayons un diagramme cartésien
$$
\xymatrix{
B \ar[r]_s & R \\
B_i \ar[u] \ar[r] & R_i \ar[u]_{t_i}
}
$$
pour chaque $i$. Les homomorphismes de transition $R_i \to R_{i'}$
induisent des homomorphismes $B_i \to B_{i'}$. Il est clair que
$B = \colim_i B_i$. Au paragraphe suivant, nous montrons que, pour tout $i$
assez grand, la composée $W_i \to \Spec(R_i) \to \Spec(B_i)$ est une
immersion ouverte.

\medskip\noindent
Comme $W$ est un ouvert quasi-compact de $U = \Spec(B)$, nous pouvons trouver
un nombre fini d'éléments $g_l \in B$, $l = 1, \ldots, m$, tels que
$D(g_l) \subset W$ et $W = \bigcup_{l = 1, \ldots, m} D(g_l)$. Remarquons
que cela implique $D(g_l) = W_{s(g_l)}$ comme ouverts de $U$, où
$W_{s(g_l)}$ désigne le plus grand ouvert de $W$ sur lequel $s(g_l)$ est
inversible. Par conséquent,
$$
B_{g_l} =
\Gamma(D(g_l), \mathcal{O}_U) =
\Gamma(W_{s(g_l)}, \mathcal{O}_W) = R_{s(g_l)},
$$
où la dernière égalité découle de Propriétés,
Lemme \ref{properties-lemma-invert-f-sections}. Comme $W_{s(g_l)}$ est
affine, cela implique aussi que $D(s(g_l)) = W_{s(g_l)}$ comme ouverts de
$\Spec(R)$. Puisque $R = \colim_i R_i$, nous pouvons, après avoir restreint
$I$, supposer qu'il existe des $g_{l, i} \in R_i$ pour tout $i \in I$ tels
que $s(g_l) = t_i(g_{l, i})$. Bien entendu, nous choisissons les $g_{l, i}$
de sorte que $g_{l, i}$ s'envoie sur $g_{l, i'}$ par les homomorphismes de
transition $R_i \to R_{i'}$. D'après le Lemme \ref{limits-lemma-descend-opens}, nous
pouvons alors, après avoir restreint $I$ une nouvelle fois, supposer que les
ouverts correspondants $D(g_{l, i}) \subset \Spec(R_i)$ sont contenus dans
$W_i$, pour $l = 1, \ldots, m$, et recouvrent $W_i$. Nous en déduisons que
le morphisme $W_i \to \Spec(R_i) \to \Spec(B_i)$ est une immersion ouverte ;
voir le Lemme \ref{limits-lemma-diagram}.

\medskip\noindent
D'après le Lemme \ref{limits-lemma-quasi-affine-finite-type-over-Z}, nous pouvons
écrire $B_i$ comme une limite inductive filtrante de sous-algèbres
$A_{i, p} \subset B_i$, $p \in P_i$, chacune de type fini sur $\mathbf{Z}$,
et telles que $W_i$ s'identifie à un sous-schéma ouvert de
$\Spec(A_{i, p})$. Soit $S_{i, p}$ le schéma obtenu en recollant $V_i$ et
$\Spec(A_{i, p})$ le long de l'ouvert $W_i$ ; voir Schémas,
section \ref{schemes-section-glueing-schemes}. Voici le diagramme commutatif
de schémas qui en résulte :
$$
\xymatrix{
& & V \ar[lld] \ar[d] & W \ar[l] \ar[lld] \ar[d] \\
V_i \ar[d] & W_i \ar[l] \ar[d] & S \ar[lld] & U \ar[lld] \ar[l] \\
S_{i, p} & \Spec(A_{i, p}) \ar[l]
}
$$
Le morphisme $S \to S_{i, p}$ existe parce que le carré supérieur droit est
un coproduit amalgamé dans la catégorie des schémas. Remarquons que
$S_{i, p}$ est de type fini sur $\mathbf{Z}$ puisqu'il admet un recouvrement
ouvert affine fini dont les membres sont les spectres de
$\mathbf{Z}$-algèbres de type fini. Nous définissons un préordre sur
$J = \coprod_{i \in I} P_i$ par la règle suivante :
$(i', p') \geq (i, p)$ si et seulement si $i' \geq i$ et si
l'homomorphisme $B_i \to B_{i'}$ envoie $A_{i, p}$ dans $A_{i', p'}$.
C'est exactement la condition requise pour définir un morphisme
$S_{i', p'} \to S_{i, p}$ : on forme en effet un diagramme commutatif comme
ci-dessus à l'aide des morphismes de transition $V_{i'} \to V_i$ et
$W_{i'} \to W_i$, ainsi que du morphisme
$\Spec(A_{i', p'}) \to \Spec(A_{i, p})$ induit par l'homomorphisme d'anneaux
$A_{i, p} \to A_{i', p'}$. Les commutativités nécessaires ont été intégrées
aux constructions. Nous affirmons que $S$ est la limite projective des
schémas $S_{i, p}$. Puisque, par construction, les schémas $V_i$ ont pour
limite $V$, cela se ramène au fait que $B$ est la limite des anneaux
$A_{i, p}$, ce qui est vrai par construction. L'application
$\alpha : J \to I$ est donnée par la règle $j = (i, p) \mapsto i$. Le
sous-schéma ouvert $V'_j$ est simplement l'image de $V_i \to S_{i, p}$
ci-dessus. La commutativité des diagrammes de (5) résulte clairement de la
construction. Cela achève la démonstration du lemme.
\end{proof}

\begin{proposition}
\label{limits-proposition-approximate}
Soit $S$ un schéma quasi-compact et quasi-séparé. Il existe un ensemble
ordonné filtrant $I$ et un système projectif de schémas $(S_i, f_{ii'})$
indexé par $I$ tels que
\begin{enumerate}
\item les morphismes de transition $f_{ii'}$ sont affines,
\item chaque $S_i$ est de type fini sur $\mathbf{Z}$, et
\item $S = \lim_i S_i$.
\end{enumerate}
\end{proposition}

\begin{proof}
C'est un cas particulier du Lemme \ref{limits-lemma-approximate}, avec
$V = \emptyset$.
\end{proof}






\section{Limites et morphismes de présentation finie}
\label{limits-section-finite-presentation}

\noindent
Le résultat suivant généralise Algèbre,
Lemme \ref{algebra-lemma-characterize-finite-presentation}.

\begin{proposition}
\label{limits-proposition-characterize-locally-finite-presentation}
\begin{reference}
\cite[IV, Proposition 8.14.2]{EGA}
\end{reference}
Soit $f : X \to S$ un morphisme de schémas. Les conditions suivantes sont
équivalentes :
\begin{enumerate}
\item Le morphisme $f$ est localement de présentation finie.
\item Pour tout ensemble ordonné filtrant $I$ et tout système projectif
$(T_i, f_{ii'})$ de $S$-schémas indexé par $I$, chaque $T_i$ étant affine,
on a
$$
\Mor_S(\lim_i T_i, X) =
\colim_i \Mor_S(T_i, X)
$$
\item Pour tout ensemble ordonné filtrant $I$ et tout système projectif
$(T_i, f_{ii'})$ de $S$-schémas indexé par $I$, chaque $f_{ii'}$ étant affine
et chaque $T_i$ étant quasi-compact et quasi-séparé comme schéma, on a
$$
\Mor_S(\lim_i T_i, X) =
\colim_i \Mor_S(T_i, X)
$$
\end{enumerate}
\end{proposition}

\begin{proof}
Il est clair que (3) implique (2).

\medskip\noindent
Montrons que (2) implique (1). Supposons (2). Choisissons des ouverts affines
$U \subset X$ et $V \subset S$ tels que $f(U) \subset V$. Il faut montrer
que $\mathcal{O}_S(V) \to \mathcal{O}_X(U)$ est de présentation finie.
Soit $(A_i, \varphi_{ii'})$ un système inductif filtrant
d'$\mathcal{O}_S(V)$-algèbres. Posons $A = \colim_i A_i$. D'après Algèbre,
Lemme \ref{algebra-lemma-characterize-finite-presentation}, il faut montrer
que
$$
\Hom_{\mathcal{O}_S(V)}(\mathcal{O}_X(U), A) =
\colim_i \Hom_{\mathcal{O}_S(V)}(\mathcal{O}_X(U), A_i)
$$
Considérons les schémas $T_i = \Spec(A_i)$. Ils forment un système projectif
de $V$-schémas indexé par $I$, dont les morphismes de transition
$f_{ii'} : T_i \to T_{i'}$ sont induits par les homomorphismes
d'$\mathcal{O}_S(V)$-algèbres $\varphi_{i'i}$. Posons
$T := \Spec(A) = \lim_i T_i$. En termes d'ensembles de morphismes de
schémas, la formule ci-dessus devient
$$
\Mor_V(\lim_i T_i, U) =
\colim_i \Mor_V(T_i, U).
$$
Observons d'abord que
$\Mor_V(T_i, U) = \Mor_S(T_i, U)$
et
$\Mor_V(T, U) = \Mor_S(T, U)$.
Il faut donc montrer que
$$
\Mor_S(\lim_i T_i, U) =
\colim_i \Mor_S(T_i, U)
$$
et l'on sait que
$$
\Mor_S(\lim_i T_i, X) =
\colim_i \Mor_S(T_i, X).
$$
Il suffit donc de prouver ceci : étant donné un morphisme
$g_i : T_i \to X$ sur $S$ tel que la composée $T \to T_i \to X$ ait son
image dans $U$, il existe $i' \geq i$ tel que la composée
$g_{i'} : T_{i'} \to T_i \to X$ ait son image dans $U$. Notons
$Z_{i'} = g_{i'}^{-1}(X \setminus U)$. Supposons tous les $Z_{i'}$ non vides,
en vue d'une contradiction. D'après le
Lemme \ref{limits-lemma-limit-closed-nonempty}, il existe un point $t$ de $T$ qui
s'envoie dans $Z_{i'}$ pour tout $i' \geq i$. Un tel point ne s'envoie pas
dans $U$, contradiction.

\medskip\noindent
Enfin, montrons que (1) implique (3). Supposons (1). Soit donné un système
projectif filtrant $(T_i, f_{ii'})$ de $S$-schémas. Supposons les morphismes
$f_{ii'}$ affines et chaque $T_i$ quasi-compact et quasi-séparé comme schéma.
Soit $T = \lim_i T_i$. Notons $f_i : T \to T_i$ les morphismes de
projection. Il faut montrer :
\begin{enumerate}
\item[(a)] étant donnés des morphismes $g_i, g'_i : T_i \to X$ sur $S$ tels
que $g_i \circ f_i = g'_i \circ f_i$, il existe $i' \geq i$ tel que
$g_i \circ f_{i'i} = g'_i \circ f_{i'i}$ ;
\item[(b)] étant donné tout morphisme $g : T \to X$ sur $S$, il existe
$i \in I$ et un morphisme $g_i : T_i \to X$ tels que
$g = f_i \circ g_i$.
\end{enumerate}

\noindent
Montrons d'abord l'unicité (a). Soient $g_i, g'_i : T_i \to X$ des
morphismes tels que $g_i \circ f_i = g'_i \circ f_i$. Pour tout
$i' \geq i$, posons $g_{i'} = g_i \circ f_{i'i}$ et
$g'_{i'} = g'_i \circ f_{i'i}$. Posons également
$g = g_i \circ f_i = g'_i \circ f_i$. Considérons le morphisme
$(g_i, g'_i) : T_i \to X \times_S X$. Posons
$$
W =
\bigcup\nolimits_{U \subset X\text{ ouvert affine},
V \subset S\text{ ouvert affine}, f(U) \subset V}
U \times_V U.
$$
C'est un ouvert de $X \times_S X$ tel que le morphisme $\Delta_{X/S}$ se
factorise par une immersion fermée dans $W$ ; voir la démonstration de
Schémas, Lemme \ref{schemes-lemma-diagonal-immersion}. Remarquons que la
composée $(g_i, g'_i) \circ f_i : T \to X \times_S X$ est un morphisme dans
$W$, puisqu'elle se factorise par la diagonale par hypothèse. Posons
$Z_{i'} = (g_{i'}, g'_{i'})^{-1}(X \times_S X \setminus W)$. Si tous les
$Z_{i'}$ sont non vides, le Lemme \ref{limits-lemma-limit-closed-nonempty} fournit
un point $t \in T$ qui s'envoie dans $Z_{i'}$ pour tout $i' \geq i$. Cela
contredit le fait que $T$ s'envoie dans $W$. Nous pouvons donc augmenter $i$
et supposer que $(g_i, g'_i) : T_i \to X \times_S X$ est un morphisme dans
$W$. Par construction de $W$, et puisque $T_i$ est quasi-compact, nous
pouvons trouver un recouvrement ouvert affine fini
$T_i = T_{1, i} \cup \ldots \cup T_{n, i}$ tel que
$(g_i, g'_i)|_{T_{j, i}}$ soit un morphisme dans $U \times_V U$ pour une
certaine paire $(U, V)$ intervenant dans la définition de $W$ ci-dessus.
Comme il suffit de prouver que $g_{i'}$ et $g'_{i'}$ coïncident sur chacun
des $f_{i'i}^{-1}(T_{j, i})$, on se ramène au cas affine. Celui-ci résulte
d'Algèbre, Lemme \ref{algebra-lemma-characterize-finite-presentation}, et du
fait que l'homomorphisme d'anneaux
$\mathcal{O}_S(V) \to \mathcal{O}_X(U)$ est de présentation finie (voir
Morphismes,
Lemme \ref{morphisms-lemma-locally-finite-presentation-characterize}).

\medskip\noindent
Enfin, montrons l'existence (b). Soit $g : T \to X$ un morphisme de schémas
sur $S$. Nous pouvons trouver un recouvrement ouvert affine fini
$T = W_1 \cup \ldots \cup W_n$ tel que, pour tout
$j \in \{1, \ldots, n\}$, il existe des ouverts affines $U_j \subset X$ et
$V_j \subset S$ avec $f(U_j) \subset V_j$ et $g(W_j) \subset U_j$.
D'après les Lemmes \ref{limits-lemma-descend-opens} et \ref{limits-lemma-limit-affine}, et
après avoir éventuellement restreint $I$, nous pouvons supposer qu'il existe
des recouvrements ouverts affines
$T_i = W_{1, i} \cup \ldots \cup W_{n, i}$ compatibles aux morphismes de
transition et tels que $W_j = \lim_i W_{j, i}$. Appliquons Algèbre,
Lemme \ref{algebra-lemma-characterize-finite-presentation}, aux anneaux
correspondant aux schémas affines $U_j$, $V_j$, $W_{j, i}$ et $W_j$, en
utilisant que $\mathcal{O}_S(V_j) \to \mathcal{O}_X(U_j)$ est de présentation
finie (voir Morphismes,
Lemme \ref{morphisms-lemma-locally-finite-presentation-characterize}). Pour
chaque $j$, nous pouvons ainsi trouver un indice $i_j \in I$ et un morphisme
$g_{j, i_j} : W_{j, i_j} \to X$ tels que
$g_{j, i_j} \circ f_i|_{W_j} : W_j \to W_{j, i} \to X$ soit égal à
$g|_{W_j}$. Par l'assertion (a) démontrée ci-dessus, et en utilisant la
quasi-compacité de $W_{j_1, i} \cap W_{j_2, i}$, qui résulte de ce que $T_i$
est quasi-séparé, nous pouvons trouver un indice $i' \in I$ supérieur à tous
les $i_j$ tel que
$$
g_{j_1, i_{j_1}} \circ f_{i'i_{j_1}}|_{W_{j_1, i'} \cap W_{j_2, i'}} =
g_{j_2, i_{j_2}} \circ f_{i'i_{j_2}}|_{W_{j_1, i'} \cap W_{j_2, i'}}
$$
pour tous $j_1, j_2 \in \{1, \ldots, n\}$. Les morphismes
$g_{j, i_j} \circ f_{i'i_j}|_{W_{j, i'}}$ se recollent donc pour donner le
morphisme cherché $T_{i'} \to X$.
\end{proof}

\begin{remark}
\label{limits-remark-limit-preserving}
Soit $S$ un schéma. Disons qu'un foncteur
$F : (\Sch/S)^{opp} \to \textit{Ensembles}$ est
{\it compatible aux limites} si, pour tout système projectif filtrant
$\{T_i\}_{i \in I}$ de schémas affines de limite $T$, on a
$F(T) = \colim_i F(T_i)$. Soit $X$ un schéma sur $S$ et soit
$h_X : (\Sch/S)^{opp} \to \textit{Ensembles}$ son foncteur des points ; voir
Schémas, section \ref{schemes-section-representable}. Dans cette
terminologie, la
Proposition \ref{limits-proposition-characterize-locally-finite-presentation}
affirme qu'un schéma $X$ est localement de présentation finie sur $S$ si et
seulement si $h_X$ est compatible aux limites.
\end{remark}

\begin{lemma}
\label{limits-lemma-surjection-is-enough}
Soit $f : X \to S$ un morphisme de schémas. Si, pour toute limite projective
$T = \lim_{i \in I} T_i$ de schémas affines sur $S$, l'application
$$
\colim \Mor_S(T_i, X) \longrightarrow \Mor_S(T, X)
$$
est surjective, alors $f$ est localement de présentation finie. Autrement
dit, dans les assertions (2) et (3) de la
Proposition \ref{limits-proposition-characterize-locally-finite-presentation}, il
suffit de vérifier la surjectivité de l'application.
\end{lemma}

\begin{proof}
La démonstration est exactement la même que celle de l'implication
« (2) implique (1) » dans la
Proposition \ref{limits-proposition-characterize-locally-finite-presentation}.
Choisissons des ouverts affines $U \subset X$ et $V \subset S$ tels que
$f(U) \subset V$. Il faut montrer que
$\mathcal{O}_S(V) \to \mathcal{O}_X(U)$ est de présentation finie. Soit
$(A_i, \varphi_{ii'})$ un système inductif filtrant d'$\mathcal{O}_S(V)$-algèbres.
Posons $A = \colim_i A_i$. D'après Algèbre,
Lemme \ref{algebra-lemma-characterize-finite-presentation}, il suffit de
montrer que
$$
\colim_i \Hom_{\mathcal{O}_S(V)}(\mathcal{O}_X(U), A_i) \to
\Hom_{\mathcal{O}_S(V)}(\mathcal{O}_X(U), A)
$$
est surjective. Considérons les schémas $T_i = \Spec(A_i)$. Ils forment un
système projectif de $V$-schémas indexé par $I$, dont les morphismes de
transition $f_{ii'} : T_i \to T_{i'}$ sont induits par les homomorphismes
d'$\mathcal{O}_S(V)$-algèbres $\varphi_{i'i}$. Posons
$T := \Spec(A) = \lim_i T_i$. En termes d'ensembles de morphismes de
schémas, la formule ci-dessus devient
$$
\colim_i \Mor_V(T_i, U) \to \Mor_V(\lim_i T_i, U)
$$
Observons d'abord que
$\Mor_V(T_i, U) = \Mor_S(T_i, U)$
et
$\Mor_V(T, U) = \Mor_S(T, U)$.
Il faut donc montrer que
$$
\colim_i \Mor_S(T_i, U) \to
\Mor_S(\lim_i T_i, U)
$$
est surjective, et l'on sait que
$$
\colim_i \Mor_S(T_i, X) \to
\Mor_S(\lim_i T_i, X)
$$
est surjective. Il suffit donc de prouver ceci : étant donné un morphisme
$g_i : T_i \to X$ sur $S$ tel que la composée $T \to T_i \to X$ ait son
image dans $U$, il existe $i' \geq i$ tel que la composée
$g_{i'} : T_{i'} \to T_i \to X$ ait son image dans $U$. Notons
$Z_{i'} = g_{i'}^{-1}(X \setminus U)$. Supposons tous les $Z_{i'}$ non vides,
en vue d'une contradiction. D'après le
Lemme \ref{limits-lemma-limit-closed-nonempty}, il existe un point $t$ de $T$ qui
s'envoie dans $Z_{i'}$ pour tout $i' \geq i$. Un tel point ne s'envoie pas
dans $U$, contradiction.
\end{proof}

\noindent
Voici un exemple d'application de la
Proposition \ref{limits-proposition-characterize-locally-finite-presentation}.

\begin{lemma}
\label{limits-lemma-morphism-glueing-near-closed-point}
Soit $S$ un schéma. Soient $X$ et $Y$ des schémas sur $S$. Supposons que
$Y$ soit localement de présentation finie sur $S$. Soit $x \in X$ un point
fermé tel que $U = X \setminus \{x\} \to X$ soit quasi-compact. En posant
$V = \Spec(\mathcal{O}_{X, x}) \setminus \{x\}$, on a une bijection
$$
\left\{
\begin{matrix}
\text{morphismes }X \to Y\text{ sur }S
\end{matrix}
\right\}
\longrightarrow
\left\{
\begin{matrix}
(a, b)\text{ où }
a : U \to Y\text{ et }b : \Spec(\mathcal{O}_{X, x}) \to Y\\
\text{ sont des morphismes sur }S
\text{ qui coïncident sur }V
\end{matrix}
\right\}
$$
\end{lemma}

\begin{proof}
Soit $W \subset X$ un voisinage ouvert de $x$. Par recollement des schémas,
voir Schémas, section \ref{schemes-section-glueing-schemes}, le résultat est
valable si l'on considère des paires de morphismes $a : U \to Y$ et
$c : W \to Y$ qui coïncident sur $U \cap W$. On a
$\mathcal{O}_{X, x} = \colim \mathcal{O}_W(W)$, où $W$ parcourt les
voisinages ouverts affines de $x$ dans $X$. Ainsi,
$\Spec(\mathcal{O}_{X, x}) = \lim W$, où $W$ parcourt les voisinages ouverts
affines de $s$. D'après la
Proposition \ref{limits-proposition-characterize-locally-finite-presentation}, tout
morphisme $b : \Spec(\mathcal{O}_{X, x}) \to Y$ sur $S$ provient donc d'un
morphisme $c : W \to Y$ pour un certain $W$ comme ci-dessus (et $c$ est
unique après un éventuel rétrécissement supplémentaire de $W$). Pour tout
ouvert affine $x \in W$, l'ouvert $U \cap W$ est quasi-compact puisque
$U \to X$ est quasi-compact. Par conséquent,
$V = \lim W \cap U = \lim W \setminus \{x\}$ est une limite de schémas
quasi-compacts et quasi-séparés (voir le
Lemme \ref{limits-lemma-directed-inverse-system-has-limit}). Ainsi, si $a$ et $b$
coïncident sur $V$, alors, après avoir rétréci $W$, $a$ et $c$ coïncident sur
$U \cap W$ (par la même proposition). Le lemme en résulte.
\end{proof}







\section{Approximation relative}
\label{limits-section-relative-approximation}

\noindent
Nous examinons des variantes sur une base de la
Proposition \ref{limits-proposition-approximate}.

\begin{lemma}
\label{limits-lemma-approximate-morphism}
Soit $f : X \to S$ un morphisme de schémas quasi-compacts et quasi-séparés.
Il existe alors un ensemble ordonné filtrant $I$ et un système projectif
$(f_i : X_i \to S_i)$ de morphismes de schémas indexé par $I$, tels que les
morphismes de transition $X_i \to X_{i'}$ et $S_i \to S_{i'}$ soient
affines, que $X_i$ et $S_i$ soient de type fini sur $\mathbf{Z}$, et que
$(X \to S) = \lim (X_i \to S_i)$.
\end{lemma}

\begin{proof}
Écrivons $X = \lim_{a \in A} X_a$ et $S = \lim_{b \in B} S_b$ comme dans la
Proposition \ref{limits-proposition-approximate}, c'est-à-dire avec $X_a$ et $S_b$
de type fini sur $\mathbf{Z}$ et des morphismes de transition affines.

\medskip\noindent
Fixons $b \in B$. En appliquant la
Proposition \ref{limits-proposition-characterize-locally-finite-presentation} à
$S_b$ et à $X = \lim X_a$ sur $\mathbf{Z}$, nous trouvons $a \in A$ et un
morphisme $f_{a, b} : X_a \to S_b$ rendant commutatif le diagramme
$$
\xymatrix{
X \ar[d] \ar[r] & S \ar[d] \\
X_a \ar[r] & S_b
}
$$
Soit $I$ l'ensemble des triplets $(a, b, f_{a, b})$ ainsi obtenus.

\medskip\noindent
Soient $(a, b, f_{a, b})$ et $(a', b', f_{a', b'})$ dans $I$. Soit
$b'' \leq \min(b, b')$. Une nouvelle application de la
Proposition \ref{limits-proposition-characterize-locally-finite-presentation}
fournit $a'' \geq \max(a, a')$ tel que les composées
$X_{a''} \to X_a \to S_b \to S_{b''}$ et
$X_{a''} \to X_{a'} \to S_{b'} \to S_{b''}$ soient égales. Munissons $I$
du préordre
$$
(a, b, f_{a, b}) \geq (a', b', f_{a', b'})
\Leftrightarrow
a \geq a',\ b \geq b',\text{ et }
g_{b, b'} \circ f_{a, b} = f_{a', b'} \circ h_{a, a'}
$$
où $h_{a, a'} : X_a \to X_{a'}$ et $g_{b, b'} : S_b \to S_{b'}$ sont les
morphismes de transition. Les remarques précédentes montrent que $I$ est
filtrant et que les applications $I \to A$,
$(a, b, f_{a, b}) \mapsto a$, et $I \to B$, $(a, b, f_{a, b})$, sont
cofinales. Si, pour $i = (a, b, f_{a, b})$, nous posons $X_i = X_a$,
$S_i = S_b$ et $f_i = f_{a, b}$, nous obtenons un système projectif de
morphismes indexé par $I$, et nous avons
$$
\lim_{i \in I} X_i = \lim_{a \in A} X_a = X
\quad\text{et}\quad
\lim_{i \in I} S_i = \lim_{b \in B} S_b = S
$$
d'après Catégories, Lemme \ref{categories-lemma-initial} (rappelons que les
limites indexées par $I$ sont en réalité des limites sur la catégorie opposée
associée à $I$, de sorte que « cofinal » devient « initial »). Cela achève la
démonstration.
\end{proof}

\begin{lemma}
\label{limits-lemma-relative-approximation}
Soit $f : X \to S$ un morphisme de schémas. Supposons que
\begin{enumerate}
\item $X$ soit quasi-compact et quasi-séparé, et
\item $S$ soit quasi-séparé.
\end{enumerate}
Alors $X = \lim X_i$ est la limite d'un système projectif filtrant de schémas $X_i$
de présentation finie sur $S$, dont les morphismes de transition sur $S$
sont affines.
\end{lemma}

\begin{proof}
Comme $f(X)$ est quasi-compact, nous pouvons remplacer $S$ par un ouvert
quasi-compact contenant $f(X)$. Nous pouvons donc supposer $S$
quasi-compact. D'après le Lemme \ref{limits-lemma-approximate-morphism}, on peut
écrire $(X \to S) = \lim (X_i \to S_i)$ pour un certain système projectif
filtrant de morphismes entre schémas de type fini sur $\mathbf{Z}$, à
morphismes de transition affines. Comme les limites commutent aux limites
(Catégories, Lemme \ref{categories-lemma-colimits-commute}), on a
$X = \lim X_i \times_{S_i} S$. Soient $i \geq i'$ dans $I$. Le morphisme
$X_i \times_{S_i} S \to X_{i'} \times_{S_{i'}} S$ est affine, car c'est la
composée
$$
X_i \times_{S_i} S \to X_i \times_{S_{i'}} S \to X_{i'} \times_{S_{i'}} S
$$
où le premier morphisme est une immersion fermée (d'après Schémas,
Lemme \ref{schemes-lemma-fibre-product-after-map}) et le second est le
changement de base d'un morphisme affine (Morphismes,
Lemme \ref{morphisms-lemma-base-change-affine}) ; or la composée de morphismes
affines est affine (Morphismes,
Lemme \ref{morphisms-lemma-composition-affine}). Les morphismes $f_i$ sont de
présentation finie (Morphismes, Lemmes
\ref{morphisms-lemma-noetherian-finite-type-finite-presentation} et
\ref{morphisms-lemma-finite-presentation-permanence}) ; par conséquent,
leurs changements de base $X_i \times_{f_i, S_i} S \to S$ sont de
présentation finie (Morphismes,
Lemme \ref{morphisms-lemma-base-change-finite-presentation}).
\end{proof}

\begin{lemma}
\label{limits-lemma-integral-limit-finite-and-finite-presentation}
Soit $X \to S$ un morphisme entier, où $S$ est quasi-compact et
quasi-séparé. Alors $X = \lim X_i$, avec $X_i \to S$ fini et de présentation
finie.
\end{lemma}

\begin{proof}
Considérons le faisceau $\mathcal{A} = f_*\mathcal{O}_X$. C'est un faisceau
quasi-cohérent d'$\mathcal{O}_S$-algèbres ; voir Schémas,
Lemme \ref{schemes-lemma-push-forward-quasi-coherent}. D'après
Propriétés,
Lemme \ref{properties-lemma-integral-algebra-directed-colimit-finite}, nous
pouvons écrire $\mathcal{A} = \colim_i \mathcal{A}_i$ comme limite inductive
filtrante d'$\mathcal{O}_S$-algèbres finies et de présentation finie. Alors
$$
X_i = \underline{\Spec}_S(\mathcal{A}_i)
\longrightarrow
S
$$
est un morphisme de schémas fini et de présentation finie. Par construction,
$X = \lim_i X_i$, ce qui démontre le lemme.
\end{proof}








\section{Descente des propriétés des morphismes}
\label{limits-section-descent-of-properties}

\noindent
Cette section est l'analogue de la section \ref{limits-section-descent} pour les
propriétés des morphismes sur $S$. Nous nous placerons dans la situation
suivante.

\begin{situation}
\label{limits-situation-descent-property}
Soit $S = \lim S_i$ la limite d'un système projectif filtrant de schémas dont les
morphismes de transition sont affines
(Lemme \ref{limits-lemma-directed-inverse-system-has-limit}). Soit $0 \in I$ et soit
$f_0 : X_0 \to Y_0$ un morphisme de schémas sur $S_0$. Supposons $S_0$,
$X_0$ et $Y_0$ quasi-compacts et quasi-séparés. Soit
$f_i : X_i \to Y_i$ le changement de base de $f_0$ à $S_i$, et soit
$f : X \to Y$ le changement de base de $f_0$ à $S$.
\end{situation}

\begin{lemma}
\label{limits-lemma-descend-affine-finite-presentation}
Notations et hypothèses comme dans la
Situation \ref{limits-situation-descent-property}. Si $f$ est affine, il existe un
indice $i \geq 0$ tel que $f_i$ soit affine.
\end{lemma}

\begin{proof}
Soit $Y_0 = \bigcup_{j = 1, \ldots, m} V_{j, 0}$ un recouvrement ouvert
affine fini. Posons $U_{j, 0} = f_0^{-1}(V_{j, 0})$. Pour $i \geq 0$, notons
$V_{j, i}$ l'image réciproque de $V_{j, 0}$ dans $Y_i$ et
$U_{j, i} = f_i^{-1}(V_{j, i})$. De même, on a $U_j = f^{-1}(V_j)$. Alors
$U_j = \lim_{i \geq 0} U_{j, i}$ (voir le
Lemme \ref{limits-lemma-directed-inverse-system-has-limit}). Comme $U_j$ est affine
par hypothèse, le Lemme \ref{limits-lemma-limit-affine} montre que chaque $U_{j, i}$
est affine pour $i$ assez grand. Puisqu'il n'y a qu'un nombre fini de $j$,
nous pouvons choisir un $i$ qui convienne à tous les $j$. Ainsi, $f_i$ est affine
pour $i$ assez grand ; voir Morphismes,
Lemme \ref{morphisms-lemma-characterize-affine}.
\end{proof}

\begin{lemma}
\label{limits-lemma-descend-finite-finite-presentation}
Notations et hypothèses comme dans la
Situation \ref{limits-situation-descent-property}. Si
\begin{enumerate}
\item $f$ est un morphisme fini, et
\item $f_0$ est localement de type fini,
\end{enumerate}
alors il existe $i \geq 0$ tel que $f_i$ soit fini.
\end{lemma}

\begin{proof}
Un morphisme fini est affine ; voir Morphismes,
Définition \ref{morphisms-definition-integral}. D'après le
Lemme \ref{limits-lemma-descend-affine-finite-presentation} ci-dessus, nous pouvons,
après avoir augmenté $0$, supposer que $f_0$ est affine. En écrivant $Y_0$
comme réunion finie d'ouverts affines, on se ramène à démontrer le résultat
lorsque $X_0$ et $Y_0$ sont affines et s'envoient dans un même ouvert affine
$W \subset S_0$. L'énoncé algébrique correspondant résulte d'Algèbre,
Lemme \ref{algebra-lemma-colimit-finite}.
\end{proof}

\begin{lemma}
\label{limits-lemma-descend-unramified}
Notations et hypothèses comme dans la
Situation \ref{limits-situation-descent-property}. Si
\begin{enumerate}
\item $f$ est non ramifié, et
\item $f_0$ est localement de type fini,
\end{enumerate}
alors il existe $i \geq 0$ tel que $f_i$ soit non ramifié.
\end{lemma}

\begin{proof}
Choisissons un recouvrement ouvert affine fini
$Y_0 = \bigcup_{j = 1, \ldots, m} Y_{j, 0}$ tel que chaque $Y_{j, 0}$
s'envoie dans un ouvert affine $S_{j, 0} \subset S_0$. Pour chaque $j$, soit
$f_0^{-1}Y_{j, 0} = \bigcup_{k = 1, \ldots, n_j} X_{k, 0}$ un recouvrement
ouvert affine fini. Comme la propriété d'être non ramifié est locale, il
suffit de démontrer le lemme pour les morphismes entre schémas affines
$X_{k, i} \to Y_{j, i} \to S_{j, i}$ obtenus par changement de base
de $X_{k, 0} \to Y_{j, 0} \to S_{j, 0}$ à $S_i$. On se ramène ainsi au cas où
$X_0, Y_0, S_0$ sont affines.

\medskip\noindent
Dans le cas affine, on se ramène au résultat algébrique suivant. Supposons
$R = \colim_{i \in I} R_i$. Pour un certain $0 \in I$, supposons donné un
homomorphisme de type fini de $R_0$-algèbres $A_i \to B_i$. Si
$R \otimes_{R_0} A_0 \to R \otimes_{R_0} B_0$ est non ramifié, alors, pour
un certain $i \geq 0$, l'homomorphisme
$R_i \otimes_{R_0} A_0 \to R_i \otimes_{R_0} B_0$ est non ramifié. Cela
résulte d'Algèbre, Lemme \ref{algebra-lemma-colimit-unramified}.
\end{proof}

\begin{lemma}
\label{limits-lemma-descend-closed-immersion-finite-presentation}
Notations et hypothèses comme dans la
Situation \ref{limits-situation-descent-property}. Si
\begin{enumerate}
\item $f$ est une immersion fermée, et
\item $f_0$ est localement de type fini,
\end{enumerate}
alors il existe $i \geq 0$ tel que $f_i$ soit une immersion fermée.
\end{lemma}

\begin{proof}
Une immersion fermée est affine ; voir Morphismes,
Lemme \ref{morphisms-lemma-closed-immersion-affine}. D'après le
Lemme \ref{limits-lemma-descend-affine-finite-presentation} ci-dessus, nous pouvons,
après avoir augmenté $0$, supposer que $f_0$ est affine. En écrivant $Y_0$
comme réunion finie d'ouverts affines, on se ramène à démontrer le résultat
lorsque $X_0$ et $Y_0$ sont affines et s'envoient dans un même ouvert affine
$W \subset S_0$. L'énoncé algébrique correspondant est une conséquence
d'Algèbre, Lemme \ref{algebra-lemma-colimit-surjective}.
\end{proof}

\begin{lemma}
\label{limits-lemma-descend-separated-finite-presentation}
Notations et hypothèses comme dans la
Situation \ref{limits-situation-descent-property}. Si $f$ est séparé, alors $f_i$
est séparé pour un certain $i \geq 0$.
\end{lemma}

\begin{proof}
Appliquons le Lemme \ref{limits-lemma-descend-closed-immersion-finite-presentation}
au morphisme diagonal
$\Delta_{X_0/S_0} : X_0 \to X_0 \times_{S_0} X_0$. (C'est permis puisque
les morphismes diagonaux sont localement de type fini et que le produit
fibré $X_0 \times_{S_0} X_0$ est quasi-compact et quasi-séparé ; voir
Schémas, Lemme \ref{schemes-lemma-diagonal-immersion}, Morphismes,
Lemme \ref{morphisms-lemma-immersion-locally-finite-type}, et Schémas,
Remarque \ref{schemes-remark-quasi-compact-and-quasi-separated}.
\end{proof}

\begin{lemma}
\label{limits-lemma-descend-flat-finite-presentation}
Notations et hypothèses comme dans la
Situation \ref{limits-situation-descent-property}. Si
\begin{enumerate}
\item $f$ est plat,
\item $f_0$ est localement de présentation finie,
\end{enumerate}
alors $f_i$ est plat pour un certain $i \geq 0$.
\end{lemma}

\begin{proof}
Choisissons un recouvrement ouvert affine fini
$Y_0 = \bigcup_{j = 1, \ldots, m} Y_{j, 0}$ tel que chaque $Y_{j, 0}$
s'envoie dans un ouvert affine $S_{j, 0} \subset S_0$. Pour chaque $j$, soit
$f_0^{-1}Y_{j, 0} = \bigcup_{k = 1, \ldots, n_j} X_{k, 0}$ un recouvrement
ouvert affine fini. Comme la propriété d'être plat est locale, il suffit de
démontrer le lemme pour les morphismes entre schémas affines
$X_{k, i} \to Y_{j, i} \to S_{j, i}$ obtenus par changement de base
de $X_{k, 0} \to Y_{j, 0} \to S_{j, 0}$ à $S_i$. On se ramène ainsi au cas où
$X_0, Y_0, S_0$ sont affines.

\medskip\noindent
Dans le cas affine, on se ramène au résultat algébrique suivant. Supposons
$R = \colim_{i \in I} R_i$. Pour un certain $0 \in I$, supposons donné un
homomorphisme de présentation finie de $R_0$-algèbres $A_i \to B_i$. Si
$R \otimes_{R_0} A_0 \to R \otimes_{R_0} B_0$ est plat, alors, pour un
certain $i \geq 0$, l'homomorphisme
$R_i \otimes_{R_0} A_0 \to R_i \otimes_{R_0} B_0$ est plat. Cela résulte
d'Algèbre, Lemme \ref{algebra-lemma-flat-finite-presentation-limit-flat},
assertion (3).
\end{proof}

\begin{lemma}
\label{limits-lemma-descend-finite-locally-free}
Notations et hypothèses comme dans la
Situation \ref{limits-situation-descent-property}. Si
\begin{enumerate}
\item $f$ est fini et localement libre (de degré $d$),
\item $f_0$ est localement de présentation finie,
\end{enumerate}
alors $f_i$ est fini et localement libre (de degré $d$) pour un certain
$i \geq 0$.
\end{lemma}

\begin{proof}
D'après les Lemmes \ref{limits-lemma-descend-flat-finite-presentation} et
\ref{limits-lemma-descend-finite-finite-presentation}, il existe $i$ tel que $f_i$
soit plat et fini. D'autre part, $f_i$ est localement de présentation finie.
Ainsi, $f_i$ est fini et localement libre d'après Morphismes,
Lemme \ref{morphisms-lemma-finite-flat}. Si, de plus, $f$ est fini et
localement libre de degré $d$, alors l'image de $Y \to Y_i$ est contenue dans
le lieu ouvert et fermé $W_d \subset Y_i$ au-dessus duquel $f_i$ est de
degré $d$. D'après le Lemme \ref{limits-lemma-limit-contained-in-constructible},
pour un certain $i' \geq i$, l'image de $Y_{i'} \to Y_i$ est contenue dans $W_d$.
Alors $f_{i'}$ est fini et localement libre de degré $d$.
\end{proof}

\begin{lemma}
\label{limits-lemma-descend-smooth}
Notations et hypothèses comme dans la
Situation \ref{limits-situation-descent-property}. Si
\begin{enumerate}
\item $f$ est lisse,
\item $f_0$ est localement de présentation finie,
\end{enumerate}
alors $f_i$ est lisse pour un certain $i \geq 0$.
\end{lemma}

\begin{proof}
La lissité est locale sur la source et sur le but (Morphismes,
Lemme \ref{morphisms-lemma-smooth-characterize}) ; nous pouvons donc supposer
$S_0, X_0, Y_0$ affines (détails omis). Le fait algébrique correspondant
est Algèbre, Lemme \ref{algebra-lemma-colimit-smooth}.
\end{proof}

\begin{lemma}
\label{limits-lemma-descend-etale}
Notations et hypothèses comme dans la
Situation \ref{limits-situation-descent-property}. Si
\begin{enumerate}
\item $f$ est étale,
\item $f_0$ est localement de présentation finie,
\end{enumerate}
alors $f_i$ est étale pour un certain $i \geq 0$.
\end{lemma}

\begin{proof}
La propriété d'être étale est locale sur la source et sur le but (Morphismes,
Lemme \ref{morphisms-lemma-etale-characterize}) ; nous pouvons donc supposer
$S_0, X_0, Y_0$ affines (détails omis). Le fait algébrique correspondant
est Algèbre, Lemme \ref{algebra-lemma-colimit-etale}.
\end{proof}

\begin{lemma}
\label{limits-lemma-descend-isomorphism}
Notations et hypothèses comme dans la
Situation \ref{limits-situation-descent-property}. Si
\begin{enumerate}
\item $f$ est un isomorphisme, et
\item $f_0$ est localement de présentation finie,
\end{enumerate}
alors $f_i$ est un isomorphisme pour un certain $i \geq 0$.
\end{lemma}

\begin{proof}
D'après les Lemmes \ref{limits-lemma-descend-etale} et
\ref{limits-lemma-descend-closed-immersion-finite-presentation}, nous pouvons
trouver $i$ tel que $f_i$ soit plat et soit une immersion fermée. Alors $f_i$
identifie $X_i$ à un sous-schéma ouvert et fermé de $Y_i$ ; voir Morphismes,
Lemme \ref{morphisms-lemma-flat-closed-immersions-finite-presentation}. Par
hypothèse, l'image de $Y \to Y_i$ s'envoie dans $f_i(X_i)$. Le
Lemme \ref{limits-lemma-limit-contained-in-constructible} montre donc que $Y_{i'}$
s'envoie dans $f_i(X_i)$ pour un certain $i' \geq i$. Il en résulte que
$X_{i'} \to Y_{i'}$ est surjectif, et on conclut.
\end{proof}

\begin{lemma}
\label{limits-lemma-descend-open-immersion}
Notations et hypothèses comme dans la
Situation \ref{limits-situation-descent-property}. Si
\begin{enumerate}
\item $f$ est une immersion ouverte, et
\item $f_0$ est localement de présentation finie,
\end{enumerate}
alors $f_i$ est une immersion ouverte pour un certain $i \geq 0$.
\end{lemma}

\begin{proof}
D'après le Lemme \ref{limits-lemma-descend-etale}, nous pouvons trouver $i$ tel que
$f_i$ soit étale. Alors $V_i = f_i(X_i)$ est un sous-schéma ouvert
quasi-compact de $Y_i$ (Morphismes,
Lemme \ref{morphisms-lemma-etale-open}). Soient $V$ et $V_{i'}$, pour
$i' \geq i$, les images réciproques de $V_i$ dans $Y$ et $Y_{i'}$. Alors
$f : X \to V$ est un isomorphisme (c'est en effet une immersion ouverte
surjective). D'après le Lemme \ref{limits-lemma-descend-isomorphism}, le morphisme
$X_{i'} \to V_{i'}$ est donc un isomorphisme pour un certain $i' \geq i$,
comme voulu.
\end{proof}

\begin{lemma}
\label{limits-lemma-descend-immersion}
Notations et hypothèses comme dans la
Situation \ref{limits-situation-descent-property}. Si
\begin{enumerate}
\item $f$ est une immersion, et
\item $f_0$ est localement de type fini,
\end{enumerate}
alors $f_i$ est une immersion pour un certain $i \geq 0$.
\end{lemma}

\begin{proof}
Il existe un ouvert $V \subset Y$ tel que le morphisme $f$ se factorise en
$X \to V \to Y$ et que $X \to V$ soit une immersion fermée ; voir la
discussion de Schémas, section \ref{schemes-section-immersions}. Puisque $X$
est quasi-compact, nous pouvons et allons supposer que $V$ est un ouvert
quasi-compact de $Y$. Après avoir augmenté $0$, le
Lemme \ref{limits-lemma-descend-opens} fournit un ouvert quasi-compact
$V_0 \subset Y_0$ tel que $V$ soit l'image réciproque de $V_0$. L'image réciproque de
$V_0$ dans $X_0$ est alors un ouvert quasi-compact dont l'image réciproque
dans $X$ est $X$. En appliquant le même lemme à $X = \lim X_i$, et après
avoir augmenté $0$, nous pouvons donc supposer donnée la factorisation
$X_0 \to V_0 \to Y_0$. Pour $i \geq 0$ assez grand, le morphisme
$X_i \to V_i$, où $V_i = Y_i \times_{Y_0} V_0$, est alors une immersion
fermée d'après le
Lemme \ref{limits-lemma-descend-closed-immersion-finite-presentation}. La
démonstration est achevée.
\end{proof}

\begin{lemma}
\label{limits-lemma-descend-monomorphism}
Notations et hypothèses comme dans la
Situation \ref{limits-situation-descent-property}. Si
\begin{enumerate}
\item $f$ est un monomorphisme, et
\item $f_0$ est localement de type fini,
\end{enumerate}
alors $f_i$ est un monomorphisme pour un certain $i \geq 0$.
\end{lemma}

\begin{proof}
Rappelons qu'un morphisme de schémas $V \to W$ est un monomorphisme si et
seulement si la diagonale $V \to V \times_W V$ est un isomorphisme (Schémas,
Lemme \ref{schemes-lemma-monomorphism}). Le morphisme
$X_0 \to X_0 \times_{Y_0} X_0$ est localement de présentation finie d'après
Morphismes, Lemme \ref{morphisms-lemma-diagonal-morphism-finite-type}.
Puisque $X_0 \times_{Y_0} X_0$ est quasi-compact et quasi-séparé (Schémas,
Remarque \ref{schemes-remark-quasi-compact-and-quasi-separated}), le
Lemme \ref{limits-lemma-descend-isomorphism} montre que
$\Delta_i : X_i \to X_i \times_{Y_i} X_i$ est un isomorphisme pour un
certain $i \geq 0$. Pour cet $i$, le morphisme $f_i$ est un monomorphisme.
\end{proof}

\begin{lemma}
\label{limits-lemma-descend-surjective}
Notations et hypothèses comme dans la
Situation \ref{limits-situation-descent-property}. Si
\begin{enumerate}
\item $f$ est surjectif, et
\item $f_0$ est localement de présentation finie,
\end{enumerate}
alors il existe $i \geq 0$ tel que $f_i$ soit surjectif.
\end{lemma}

\begin{proof}
Le morphisme $f_0$ est de présentation finie. Ainsi,
$E = f_0(X_0)$ est une partie constructible de $Y_0$ ; voir Morphismes,
Lemme \ref{morphisms-lemma-chevalley}. Comme $f_i$ est le changement de base
de $f_0$ par $Y_i \to Y_0$, l'image de $f_i$ est l'image réciproque de $E$
dans $Y_i$. De plus, nous savons que $Y \to Y_0$ s'envoie dans $E$. On
conclut donc par le Lemme \ref{limits-lemma-limit-contained-in-constructible}.
\end{proof}

\begin{lemma}
\label{limits-lemma-descend-syntomic}
Notations et hypothèses comme dans la
Situation \ref{limits-situation-descent-property}. Si
\begin{enumerate}
\item $f$ est syntomique, et
\item $f_0$ est localement de présentation finie,
\end{enumerate}
alors il existe $i \geq 0$ tel que $f_i$ soit syntomique.
\end{lemma}

\begin{proof}
Choisissons un recouvrement ouvert affine fini
$Y_0 = \bigcup_{j = 1, \ldots, m} Y_{j, 0}$ tel que chaque $Y_{j, 0}$
s'envoie dans un ouvert affine $S_{j, 0} \subset S_0$. Pour chaque $j$, soit
$f_0^{-1}Y_{j, 0} = \bigcup_{k = 1, \ldots, n_j} X_{k, 0}$ un recouvrement
ouvert affine fini. Comme la propriété d'être syntomique est locale, il
suffit de démontrer le lemme pour les morphismes entre schémas affines
$X_{k, i} \to Y_{j, i} \to S_{j, i}$ obtenus par changement de base
de $X_{k, 0} \to Y_{j, 0} \to S_{j, 0}$ à $S_i$. On se ramène ainsi au cas où
$X_0, Y_0, S_0$ sont affines.

\medskip\noindent
Dans le cas affine, on se ramène au résultat algébrique suivant. Supposons
$R = \colim_{i \in I} R_i$. Pour un certain $0 \in I$, supposons donné un
homomorphisme de présentation finie de $R_0$-algèbres $A_i \to B_i$. Si
$R \otimes_{R_0} A_0 \to R \otimes_{R_0} B_0$ est syntomique, alors, pour
un certain $i \geq 0$, l'homomorphisme
$R_i \otimes_{R_0} A_0 \to R_i \otimes_{R_0} B_0$ est syntomique. Cela
résulte d'Algèbre, Lemme \ref{algebra-lemma-colimit-lci}.
\end{proof}





\section{Immersion fermée d'un morphisme de type fini dans un morphisme de présentation finie}
\label{limits-section-finite-type-closed-in-finite-presentation}

\noindent
Un résultat de ce type est \cite[Satz 2.10]{Kiehl}. Une autre référence est
\cite{Conrad-Nagata}.

\begin{lemma}
\label{limits-lemma-locally-finite-type-in-finite-presentation}
Soit $f : X \to S$ un morphisme de schémas. Faisons les hypothèses suivantes :
\begin{enumerate}
\item le morphisme $f$ est localement de type fini ;
\item le schéma $X$ est quasi-compact et quasi-séparé.
\end{enumerate}
Alors il existe un morphisme de présentation finie $f' : X' \to S$ et une
immersion $X \to X'$ de schémas sur $S$.
\end{lemma}

\begin{proof}
D'après la Proposition \ref{limits-proposition-approximate}, on peut écrire
$X = \lim_i X_i$, chaque $X_i$ étant de type fini sur $\mathbf{Z}$ et les
morphismes de transition $f_{ii'} : X_i \to X_{i'}$ étant affines.
Considérons le diagramme commutatif
$$
\xymatrix{
X \ar[r] \ar[rd] & X_{i, S} \ar[r] \ar[d] & X_i \ar[d] \\
& S \ar[r] & \Spec(\mathbf{Z})
}
$$
Remarquons que $X_i$ est de présentation finie sur $\Spec(\mathbf{Z})$ ;
voir Morphismes,
Lemme \ref{morphisms-lemma-noetherian-finite-type-finite-presentation}.
Ainsi, le changement de base $X_{i, S} \to S$ est de présentation finie par
Morphismes, Lemme \ref{morphisms-lemma-base-change-finite-presentation}. Il
suffit donc de montrer que la flèche $X \to X_{i, S}$ est une immersion pour
$i$ assez grand.

\medskip\noindent
Pour cela, choisissons un recouvrement ouvert affine fini
$X = V_1 \cup \ldots \cup V_n$ tel que $f$ envoie chaque $V_j$ dans un ouvert
affine $U_j \subset S$. Soit $h_{j, a} \in \mathcal{O}_X(V_j)$ une famille
finie d'éléments qui engendrent $\mathcal{O}_X(V_j)$ comme
$\mathcal{O}_S(U_j)$-algèbre ; voir Morphismes,
Lemme \ref{morphisms-lemma-locally-finite-type-characterize}. D'après les
Lemmes \ref{limits-lemma-descend-opens} et \ref{limits-lemma-limit-affine}, et après avoir
éventuellement restreint $I$, nous pouvons supposer qu'il existe des
recouvrements ouverts affines $X_i = V_{1, i} \cup \ldots \cup V_{n, i}$,
compatibles aux morphismes de transition et tels que
$V_j = \lim_i V_{j, i}$. D'après le Lemme \ref{limits-lemma-descend-section}, nous
pouvons choisir $i$ assez grand pour que chaque $h_{j, a}$ provienne d'un
élément $h_{j, a, i} \in \mathcal{O}_{X_i}(V_{j, i})$. Ainsi, la flèche dans
$$
V_j \longrightarrow U_j \times_{\Spec(\mathbf{Z})} V_{j, i} =
(V_{j, i})_{U_j} \subset (V_{j, i})_S \subset X_{i, S}
$$
est une immersion fermée. Puisque $\bigcup (V_{j, i})_{U_j}$ est un ouvert
de $X_{i, S}$ et que l'image réciproque de $(V_{j, i})_{U_j}$ dans $X$ est
$V_j$, il en résulte que $X \to X_{i, S}$ est une immersion.
\end{proof}

\begin{remark}
\label{limits-remark-cannot-do-better}
On ne peut améliorer cet énoncé sans hypothèses supplémentaires sur $S$ et
sur le morphisme $f : X \to S$. Par exemple, il ne sera généralement pas
possible de trouver comme dans le lemme une immersion {\it fermée}
$X \to X'$. En effet, cela impliquerait que $f$ soit quasi-compact, ce qui
peut ne pas être le cas. On peut prendre pour exemple $S$ égal à l'espace
affine de dimension infinie dont le point $0$ est dédoublé, et $X$ égal à
l'un des deux espaces affines de dimension infinie.
\end{remark}

\begin{lemma}
\label{limits-lemma-finite-type-closed-in-finite-presentation}
Soit $f : X \to S$ un morphisme de schémas. Faisons les hypothèses suivantes :
\begin{enumerate}
\item le morphisme $f$ est localement de type fini ;
\item le schéma $X$ est quasi-compact et quasi-séparé ;
\item le schéma $S$ est quasi-séparé.
\end{enumerate}
Alors il existe un morphisme de présentation finie $f' : X' \to S$ et une
immersion fermée $X \to X'$ de schémas sur $S$.
\end{lemma}

\begin{proof}
D'après le Lemme \ref{limits-lemma-locally-finite-type-in-finite-presentation}
ci-dessus, il existe un morphisme de présentation finie $Y \to S$ et une
immersion $i : X \to Y$ de schémas sur $S$. Pour tout point $x \in X$, il
existe un ouvert affine $V_x \subset Y$ tel que
$i^{-1}(V_x) \to V_x$ soit une immersion fermée. Puisque $X$ est
quasi-compact, nous pouvons trouver un nombre fini d'ouverts affines
$V_1, \ldots, V_n \subset Y$ tels que
$i(X) \subset V_1 \cup \ldots \cup V_n$ et que
$i^{-1}(V_j) \to V_j$ soit une immersion fermée. Autrement dit,
$i : X \to X' = V_1 \cup \ldots \cup V_n$ est une immersion fermée de
schémas sur $S$. Comme $S$ est quasi-séparé et $Y$ quasi-séparé sur $S$, on
en déduit que $Y$ est quasi-séparé ; voir Schémas,
Lemme \ref{schemes-lemma-separated-permanence}. Ainsi, l'immersion ouverte
$X' = V_1 \cup \ldots \cup V_n \to Y$ est quasi-compacte. Il en résulte que
$X' \to Y$ est de présentation finie ; voir Morphismes,
Lemme \ref{morphisms-lemma-quasi-compact-open-immersion-finite-presentation}.
On conclut, car $X' \to Y \to S$ est alors une composée de morphismes de
présentation finie, donc est de présentation finie (voir Morphismes,
Lemme \ref{morphisms-lemma-composition-finite-presentation}).
\end{proof}

\begin{lemma}
\label{limits-lemma-closed-is-limit-closed-and-finite-presentation}
Soit $X \to Y$ une immersion fermée de schémas. Supposons $Y$ quasi-compact
et quasi-séparé. Alors $X$ peut s'écrire comme une limite projective
$X = \lim X_i$ de schémas sur $Y$, où $X_i \to Y$ est une immersion fermée
de présentation finie.
\end{lemma}

\begin{proof}
Soit $\mathcal{I} \subset \mathcal{O}_Y$ le faisceau quasi-cohérent d'idéaux
qui définit $X$ comme sous-schéma fermé de $Y$. D'après Propriétés,
Lemme \ref{properties-lemma-quasi-coherent-colimit-finite-type}, nous pouvons
écrire $\mathcal{I}$ comme limite inductive filtrante
$\mathcal{I} = \colim_{i \in I} \mathcal{I}_i$ de ses faisceaux
quasi-cohérents d'idéaux de type fini. Soit $X_i \subset Y$ le sous-schéma
fermé défini par $\mathcal{I}_i$. Ils forment un système projectif de schémas
indexé par $I$. Les morphismes de transition $X_i \to X_{i'}$ sont affines
parce que ce sont des immersions fermées. Chaque $X_i$ est quasi-compact et
quasi-séparé, puisqu'il est un sous-schéma fermé de $Y$ et que $Y$ est
quasi-compact et quasi-séparé par hypothèse. On a $X = \lim_i X_i$, comme il
résulte directement du fait que
$\mathcal{I} = \colim_{i \in I} \mathcal{I}_a$. Chacun des morphismes
$X_i \to Y$ est de présentation finie ; voir Morphismes,
Lemme \ref{morphisms-lemma-closed-immersion-finite-presentation}.
\end{proof}

\begin{lemma}
\label{limits-lemma-finite-type-is-limit-finite-presentation}
Soit $f : X \to S$ un morphisme de schémas. Faisons les hypothèses suivantes :
\begin{enumerate}
\item le morphisme $f$ est localement de type fini ;
\item le schéma $X$ est quasi-compact et quasi-séparé ;
\item le schéma $S$ est quasi-séparé.
\end{enumerate}
Alors $X = \lim X_i$, où les $X_i \to S$ sont de présentation finie, les
$X_i$ sont quasi-compacts et quasi-séparés, et les morphismes de transition
$X_{i'} \to X_i$ sont des immersions fermées (ce qui implique que
$X \to X_i$ est une immersion fermée pour tout $i$).
\end{lemma}

\begin{proof}
D'après le Lemme \ref{limits-lemma-finite-type-closed-in-finite-presentation}, il
existe une immersion fermée $X \to Y$, où $Y \to S$ est de présentation
finie. Le schéma $Y$ est alors quasi-séparé par Schémas,
Lemme \ref{schemes-lemma-separated-permanence}. Comme $X$ est quasi-compact,
nous pouvons supposer $Y$ quasi-compact en remplaçant $Y$ par un ouvert
quasi-compact contenant $X$. On a $X = \lim X_i$, avec
$X_i \to Y$ immersion fermée de présentation finie, d'après le
Lemme \ref{limits-lemma-closed-is-limit-closed-and-finite-presentation}.
Les morphismes $X_i \to S$ sont de présentation finie par Morphismes,
Lemme \ref{morphisms-lemma-composition-finite-presentation}.
\end{proof}

\begin{proposition}
\label{limits-proposition-separated-closed-in-finite-presentation}
Soit $f : X \to S$ un morphisme de schémas. Faisons les hypothèses suivantes :
\begin{enumerate}
\item $f$ est de type fini et séparé, et
\item $S$ est quasi-compact et quasi-séparé.
\end{enumerate}
Alors il existe un morphisme séparé de présentation finie $f' : X' \to S$
et une immersion fermée $X \to X'$ de schémas sur $S$.
\end{proposition}

\begin{proof}
Appliquons le Lemme \ref{limits-lemma-finite-type-is-limit-finite-presentation} et
remarquons que $X_i \to S$ est séparé pour $i$ assez grand d'après le
Lemme \ref{limits-lemma-eventually-separated}, puisque nous avons supposé
$X \to S$ séparé.
\end{proof}

\begin{lemma}
\label{limits-lemma-finite-closed-in-finite-finite-presentation}
Soit $f : X \to S$ un morphisme de schémas. Faisons les hypothèses suivantes :
\begin{enumerate}
\item $f$ est fini, et
\item $S$ est quasi-compact et quasi-séparé.
\end{enumerate}
Alors il existe un morphisme $f' : X' \to S$ fini et de présentation finie,
ainsi qu'une immersion fermée $X \to X'$ de schémas sur $S$.
\end{lemma}

\begin{proof}
Nous pouvons écrire $X = \lim X_i$ comme dans le
Lemme \ref{limits-lemma-finite-type-is-limit-finite-presentation}. En appliquant le
Lemme \ref{limits-lemma-eventually-finite}, nous voyons que $X_i \to S$ est fini
pour $i$ assez grand.
\end{proof}

\begin{lemma}
\label{limits-lemma-finite-in-finite-and-finite-presentation}
Soit $f : X \to S$ un morphisme de schémas. Faisons les hypothèses suivantes :
\begin{enumerate}
\item $f$ est fini, et
\item $S$ est quasi-compact et quasi-séparé.
\end{enumerate}
Alors $X$ est une limite projective $X = \lim X_i$, dont les morphismes de
transition sont des immersions fermées et dont les objets $X_i$ sont finis
et de présentation finie sur $S$.
\end{lemma}

\begin{proof}
Nous pouvons écrire $X = \lim X_i$ comme dans le
Lemme \ref{limits-lemma-finite-type-is-limit-finite-presentation}. En appliquant le
Lemme \ref{limits-lemma-eventually-finite}, nous voyons que $X_i \to S$ est fini
pour $i$ assez grand.
\end{proof}








\section{Descente des objets relatifs}
\label{limits-section-descending-relative}

\noindent
Le lemme suivant est représentatif des résultats de cette section.
Nous en donnons intégralement la démonstration « standard ». Il sera
peut-être plus rapide de se convaincre du résultat que de lire cette démonstration.

\begin{lemma}
\label{limits-lemma-descend-finite-presentation}
Soit $I$ un ensemble ordonné filtrant.
Soit $(S_i, f_{ii'})$ un système projectif de schémas indexé par $I$.
Supposons que
\begin{enumerate}
\item les morphismes $f_{ii'} : S_i \to S_{i'}$ soient affines,
\item les schémas $S_i$ soient quasi-compacts et quasi-séparés.
\end{enumerate}
Soit $S = \lim_i S_i$. On a alors les propriétés suivantes :
\begin{enumerate}
\item Pour tout morphisme de présentation finie $X \to S$,
il existe un indice $i \in I$ et un morphisme de présentation finie
$X_i \to S_i$ tels que $X \cong X_{i, S}$ comme schémas sur $S$.
\item Étant donnés un indice $i \in I$, des schémas
$X_i$, $Y_i$ de présentation finie sur $S_i$ et un morphisme
$\varphi : X_{i, S} \to Y_{i, S}$ sur $S$, il existe un indice
$i' \geq i$ et un morphisme
$\varphi_{i'} : X_{i, S_{i'}} \to Y_{i, S_{i'}}$
dont le changement de base à $S$ est $\varphi$.
\item Étant donnés un indice $i \in I$, des schémas $X_i$, $Y_i$
de présentation finie sur $S_i$ et deux morphismes
$\varphi_i, \psi_i : X_i \to Y_i$
dont les changements de base $\varphi_{i, S} = \psi_{i, S}$ sont égaux,
il existe un indice $i' \geq i$ tel que
$\varphi_{i, S_{i'}} = \psi_{i, S_{i'}}$.
\end{enumerate}
Autrement dit, la catégorie des schémas de présentation finie sur
$S$ est la colimite, indexée par $I$, des catégories de schémas de
présentation finie sur $S_i$.
\end{lemma}

\begin{proof}
Si tous les schémas $S_i$ sont affines et si l'on ne considère
que les schémas affines de présentation finie sur $S_i$, resp.\ $S$,
ce lemme équivaut à
Algèbre, Lemme \ref{algebra-lemma-colimit-category-fp-algebras}.
Nous affirmons que le cas affine entraîne le cas général.

\medskip\noindent
Démontrons (3). Donnons-nous un indice $i \in I$, des schémas
$X_i$, $Y_i$ de présentation finie sur $S_i$ et deux morphismes
$\varphi_i, \psi_i : X_i \to Y_i$. Supposons leurs changements de base
égaux : $\varphi_{i, S} = \psi_{i, S}$. Nous noterons
$X_{i'} = X_{i, S_{i'}}$ et $Y_{i'} = Y_{i, S_{i'}}$ pour
$i' \geq i$. Posons aussi $X = X_{i, S}$ et $Y = Y_{i, S}$.
D'après le Lemme \ref{limits-lemma-scheme-over-limit}, on a
$X = \lim_{i' \geq i} X_{i'}$, et de même pour $Y$.
Notons en outre $\varphi_{i'}$ et $\psi_{i'}$
(resp.\ $\varphi$ et $\psi$)
les changements de base de $\varphi_i$ et $\psi_i$ à $S_{i'}$
(resp.\ $S$). Notre hypothèse signifie donc que $\varphi = \psi$.
Puisque $Y_i$ et $X_i$ sont de présentation finie
sur $S_i$ et que $S_i$ est quasi-compact et quasi-séparé,
$X_i$ et $Y_i$ sont eux aussi quasi-compacts et quasi-séparés
(voir Morphismes,
Lemme \ref{morphisms-lemma-finite-presentation-quasi-compact-quasi-separated}).
On peut donc choisir un recouvrement ouvert affine fini
$Y_i = \bigcup V_{j, i}$ tel que chaque $V_{j, i}$ ait son image dans
un ouvert affine de $S_i$. Comme ci-dessus, notons $V_{j, i'}$ l'image
réciproque de $V_{j, i}$ dans $Y_{i'}$ et $V_j$ son image réciproque dans $Y$.
Les immersions $V_{j, i'} \to Y_{i'}$ sont quasi-compactes, et les images réciproques
$U_{j, i'} = \varphi_i^{-1}(V_{j, i'})$ et
$U_{j, i'}' = \psi_i^{-1}(V_{j, i'})$
sont des ouverts quasi-compacts de $X_{i'}$. Par hypothèse, les images réciproques
de $V_j$ par $\varphi$ et $\psi$ dans $X$ sont égales.
D'après le Lemme \ref{limits-lemma-descend-opens},
il existe donc un indice $i' \geq i$ tel que
$U_{j, i'} = U_{j, i'}'$ dans $X_{i'}$.
Choisissons un recouvrement ouvert affine fini
$U_{j, i'} = U_{j, i'}' = \bigcup W_{j, k, i'}$
qui induit des recouvrements $U_{j, i''} = U_{j, i''}' = \bigcup W_{j, k, i''}$
pour tout $i'' \geq i'$.
D'après le cas affine, il existe
un indice $i''$ tel que
$\varphi_{i''}|_{W_{j, k, i''}} = \psi_{i''}|_{W_{j, k, i''}}$
pour tous $j, k$. Alors $i''$ vérifie
$\varphi_{i''} = \psi_{i''}$, ce qui démontre (3).

\medskip\noindent
Démontrons (2). Donnons-nous un indice $i \in I$, des schémas
$X_i$, $Y_i$ de présentation finie sur $S_i$ et un morphisme
$\varphi : X_{i, S} \to Y_{i, S}$. Nous noterons
$X_{i'} = X_{i, S_{i'}}$ et $Y_{i'} = Y_{i, S_{i'}}$ pour
$i' \geq i$. Posons aussi $X = X_{i, S}$ et $Y = Y_{i, S}$.
D'après le Lemme \ref{limits-lemma-scheme-over-limit}, on a
$X = \lim_{i' \geq i} X_{i'}$, et de même pour $Y$.
Puisque $Y_i$ et $X_i$ sont de présentation finie
sur $S_i$ et que $S_i$ est quasi-compact et quasi-séparé,
$X_i$ et $Y_i$ sont eux aussi quasi-compacts et quasi-séparés
(voir Morphismes,
Lemme \ref{morphisms-lemma-finite-presentation-quasi-compact-quasi-separated}).
On peut donc choisir un recouvrement ouvert affine fini
$Y_i = \bigcup V_{j, i}$ tel que chaque $V_{j, i}$ ait son image dans
un ouvert affine de $S_i$. Comme ci-dessus, notons $V_{j, i'}$ l'image
réciproque de $V_{j, i}$ dans $Y_{i'}$ et $V_j$ son image réciproque dans $Y$.
Les immersions $V_j \to Y$ sont quasi-compactes, et les images réciproques
$U_j = \varphi^{-1}(V_j)$ sont des ouverts quasi-compacts de $X$.
D'après le Lemme \ref{limits-lemma-descend-opens}, il existe donc un indice
$i' \geq i$ et des ouverts quasi-compacts $U_{j, i'}$ de $X_{i'}$
dont l'image réciproque dans $X$ est $U_j$. Choisissons un recouvrement ouvert affine fini
$U_{j, i'} = \bigcup W_{j, k, i'}$ qui induit des recouvrements ouverts affines
$U_{j, i''} = \bigcup W_{j, k, i''}$
pour tout $i'' \geq i'$ et un recouvrement ouvert affine
$U_j = \bigcup W_{j, k}$. D'après le cas affine, il existe
un indice $i''$ et des morphismes
$\varphi_{j, k, i''} : W_{j, k, i''} \to V_{j, i''}$
tels que
$\varphi|_{W_{j, k}} = \varphi_{j, k, i'', S}$ pour tous $j, k$.
D'après (3), démontré ci-dessus, il existe un indice $i''' \geq i''$
tel que
$$
\varphi_{j_1, k_1, i'', S_{i'''}}|_{W_{j_1, k_1, i'''} \cap W_{j_2, k_2, i'''}}
=
\varphi_{j_2, k_2, i'', S_{i'''}}|_{W_{j_1, k_1, i'''} \cap W_{j_2, k_2, i'''}}
$$
pour tous $j_1, j_2, k_1, k_2$. Alors $i'''$ est un indice pour lequel
il existe un morphisme $\varphi_{i'''} : X_{i'''} \to Y_{i'''}$
dont le changement de base à $S$ est $\varphi$. Ainsi (2) est établi.

\medskip\noindent
Démontrons (1). Donnons-nous un schéma $X$ de présentation finie
sur $S$. Puisque $X$ est de présentation finie
sur $S$ et que $S$ est quasi-compact et quasi-séparé,
$X$ est lui aussi quasi-compact et quasi-séparé
(voir Morphismes,
Lemme \ref{morphisms-lemma-finite-presentation-quasi-compact-quasi-separated}).
Choisissons un recouvrement ouvert affine fini $X = \bigcup U_j$
tel que chaque $U_j$ ait son image dans un ouvert affine $V_j \subset S$.
Notons $U_{j_1j_2} = U_{j_1} \cap U_{j_2}$ et
$U_{j_1j_2j_3} = U_{j_1} \cap U_{j_2} \cap U_{j_3}$.
D'après les Lemmes \ref{limits-lemma-descend-opens} et \ref{limits-lemma-limit-affine},
on peut trouver un indice $i_1$ et des ouverts affines $V_{j, i_1} \subset S_{i_1}$
tels que chaque $V_j$ soit leur image réciproque dans $S$.
Soit $V_{j, i}$ l'image réciproque de $V_{j, i_1}$ dans $S_i$ pour
$i \geq i_1$. D'après le cas affine, on peut trouver un indice $i_2 \geq i_1$ et
des schémas affines $U_{j, i_2} \to V_{j, i_2}$ tels
que $U_j = S \times_{S_{i_2}} U_{j, i_2}$ soit le changement de base.
Notons $U_{j, i} = S_i \times_{S_{i_2}} U_{j, i_2}$ pour $i \geq i_2$.
D'après le Lemme \ref{limits-lemma-descend-opens}, il existe un indice
$i_3 \geq i_2$ et des sous-schémas ouverts
$W_{j_1, j_2, i_3} \subset U_{j_1, i_3}$
dont le changement de base à $S$ est égal à $U_{j_1j_2}$.
Notons $W_{j_1, j_2, i} = S_i \times_{S_{i_3}} W_{j_1, j_2, i_3}$
pour $i \geq i_3$. D'après (2), démontré ci-dessus, il existe un indice
$i_4 \geq i_3$ et des morphismes
$\varphi_{j_1, j_2, i_4} : W_{j_1, j_2, i_4} \to W_{j_2, j_1, i_4}$
dont le changement de base à $S$ est le morphisme identité
$U_{j_1j_2} = U_{j_2j_1}$ pour tous $j_1, j_2$.
Pour tout $i \geq i_4$, notons
$\varphi_{j_1, j_2, i} = \text{id}_S \times \varphi_{j_1, j_2, i_4}$
le changement de base. Nous affirmons que, pour un certain $i_5 \geq i_4$, le système
$((U_{j, i_5})_j, (W_{j_1, j_2, i_5})_{j_1, j_2},
(\varphi_{j_1, j_2, i_5})_{j_1, j_2})$ est une donnée de recollement
au sens de Schémas, section \ref{schemes-section-glueing-schemes}.
Pour le voir, il faut vérifier que, pour $i$ assez grand,
on a
$$
\varphi_{j_1, j_2, i}^{-1}(W_{j_2, j_1, i} \cap W_{j_2, j_3, i})
=
W_{j_1, j_2, i} \cap W_{j_1, j_3, i}
$$
et que, pour $i$ assez grand, la condition de cocycle est satisfaite.
La première condition résulte du Lemme \ref{limits-lemma-descend-opens}
et de l'égalité $U_{j_2j_1j_3} = U_{j_1j_2j_3}$.
La seconde résulte de (3), démontré ci-dessus, et du fait
que la condition de cocycle est satisfaite par les applications
$\text{id} : U_{j_1j_2} \to U_{j_2j_1}$.
On peut donc appliquer Schémas, Lemme \ref{schemes-lemma-glue-schemes}
pour recoller le système
$((U_{j, i_5})_j, (W_{j_1, j_2, i_5})_{j_1, j_2},
(\varphi_{j_1, j_2, i_5})_{j_1, j_2})$ et obtenir un schéma
$X_{i_5} \to S_{i_5}$. Par construction, le changement de base de
$X_{i_5}$ à $S$ est obtenu en recollant les ouverts affines
$U_j$ le long des ouverts $U_{j_1} \leftarrow U_{j_1j_2} \rightarrow U_{j_2}$.
Ainsi $S \times_{S_{i_5}} X_{i_5} \cong X$, comme voulu.
\end{proof}

\begin{lemma}
\label{limits-lemma-descend-modules-finite-presentation}
Soit $I$ un ensemble ordonné filtrant.
Soit $(S_i, f_{ii'})$ un système projectif de schémas indexé par $I$.
Supposons que
\begin{enumerate}
\item tous les morphismes $f_{ii'} : S_i \to S_{i'}$ soient affines,
\item tous les schémas $S_i$ soient quasi-compacts et quasi-séparés.
\end{enumerate}
Soit $S = \lim_i S_i$. On a alors les propriétés suivantes :
\begin{enumerate}
\item Pour tout faisceau de $\mathcal{O}_S$-modules
$\mathcal{F}$ de présentation finie, il existe un indice
$i \in I$ et un faisceau de $\mathcal{O}_{S_i}$-modules de présentation finie
$\mathcal{F}_i$ tels que
$\mathcal{F} \cong f_i^*\mathcal{F}_i$.
\item Donnons-nous un indice $i \in I$, des faisceaux
de $\mathcal{O}_{S_i}$-modules $\mathcal{F}_i$, $\mathcal{G}_i$
de présentation finie et un morphisme
$\varphi : f_i^*\mathcal{F}_i \to f_i^*\mathcal{G}_i$ sur $S$.
Il existe alors un indice $i' \geq i$ et un morphisme
$\varphi_{i'} : f_{i'i}^*\mathcal{F}_i \to f_{i'i}^*\mathcal{G}_i$
dont le changement de base à $S$ est $\varphi$.
\item Donnons-nous un indice $i \in I$, des faisceaux de $\mathcal{O}_{S_i}$-modules
$\mathcal{F}_i$, $\mathcal{G}_i$ de présentation finie
et deux morphismes $\varphi_i, \psi_i : \mathcal{F}_i \to \mathcal{G}_i$.
Supposons leurs changements de base égaux : $f_i^*\varphi_i = f_i^*\psi_i$.
Il existe alors un indice $i' \geq i$ tel que
$f_{i'i}^*\varphi_i = f_{i'i}^*\psi_i$.
\end{enumerate}
Autrement dit, la catégorie des modules
de présentation finie sur $S$ est la colimite, indexée par $I$,
des catégories de modules de présentation finie sur $S_i$.
\end{lemma}

\begin{proof}
Nous esquissons deux démonstrations, en omettant les détails.

\medskip\noindent
Première démonstration. Si $S$ et les $S_i$ sont affines, ce lemme
équivaut à
Algèbre, Lemme \ref{algebra-lemma-colimit-category-fp-modules}.
Dans le cas général, on se ramène au cas affine par recollement de Zariski.

\medskip\noindent
Seconde démonstration. On utilise les faits suivants :
\begin{enumerate}
\item il existe une équivalence de catégories entre les
$\mathcal{O}_S$-modules quasi-cohérents et les fibrés vectoriels sur $S$ ; voir
Constructions, section \ref{constructions-section-vector-bundle},
\item un fibré vectoriel $\mathbf{V}(\mathcal{F}) \to S$ est
de présentation finie sur $S$ si et seulement si $\mathcal{F}$
est un $\mathcal{O}_S$-module de présentation finie.
\end{enumerate}
Ceci étant, on peut utiliser le Lemme \ref{limits-lemma-descend-finite-presentation}
pour montrer que la catégorie des fibrés vectoriels de présentation finie
sur $S$ est la colimite, indexée par $I$,
des catégories de fibrés vectoriels sur $S_i$.
\end{proof}

\begin{lemma}
\label{limits-lemma-descend-invertible-modules}
Soit $S = \lim S_i$ la limite d'un système projectif filtrant de schémas
quasi-compacts et quasi-séparés $S_i$, à morphismes de transition affines. Alors
\begin{enumerate}
\item tout $\mathcal{O}_S$-module localement libre de rang fini est l'image inverse
d'un $\mathcal{O}_{S_i}$-module localement libre de rang fini pour un certain $i$,
\item tout $\mathcal{O}_S$-module inversible est l'image inverse d'un
$\mathcal{O}_{S_i}$-module inversible pour un certain $i$,
\item tout idéal quasi-cohérent de type fini $\mathcal{I} \subset \mathcal{O}_S$
est de la forme $\mathcal{I}_i \cdot \mathcal{O}_S$ pour un certain $i$ et un
idéal quasi-cohérent de type fini $\mathcal{I}_i \subset \mathcal{O}_{S_i}$.
\end{enumerate}
\end{lemma}

\begin{proof}
Soit $\mathcal{E}$ un $\mathcal{O}_S$-module localement libre de rang fini. Puisque
les modules localement libres de rang fini sont de présentation finie, on peut trouver un $i$
et un $\mathcal{O}_{S_i}$-module $\mathcal{E}_i$ de présentation finie
tels que $f_i^*\mathcal{E}_i \cong \mathcal{E}$ ; voir
Lemme \ref{limits-lemma-descend-modules-finite-presentation}.
En augmentant $i$, on peut supposer que $\mathcal{E}_i$ est un
$\mathcal{O}_{S_i}$-module plat ; voir
Algèbre, Lemme \ref{algebra-lemma-flat-finite-presentation-limit-flat}.
(L'emploi de ce lemme n'est pas nécessaire, mais commode.)
Alors $\mathcal{E}_i$ est localement libre de rang fini d'après
Algèbre, Lemme \ref{algebra-lemma-finite-projective}.

\medskip\noindent
Si $\mathcal{L}$ est un $\mathcal{O}_S$-module inversible,
ce qui précède fournit un $i$ et des
$\mathcal{O}_{S_i}$-modules localement libres de rang fini $\mathcal{L}_i$
et $\mathcal{N}_i$ dont les images inverses sont $\mathcal{L}$ et
$\mathcal{L}^{\otimes -1}$. Après avoir éventuellement augmenté
$i$, l'application
$\mathcal{L} \otimes_{\mathcal{O}_X} \mathcal{L}^{\otimes -1}
\to \mathcal{O}_X$ descend en une application
$\mathcal{L}_i \otimes_{\mathcal{O}_{S_i}} \mathcal{N}_i \to
\mathcal{O}_{S_i}$. En augmentant encore $i$, on
peut supposer que c'est un isomorphisme. Il s'ensuit que
$\mathcal{L}_i$ est un module inversible
(Modules, Lemme \ref{modules-lemma-invertible}), ce qui
achève la démonstration de (2).

\medskip\noindent
Étant donné $\mathcal{I}$ comme dans (3), on voit que
$\mathcal{O}_S \to \mathcal{O}_S/\mathcal{I}$
est une application de $\mathcal{O}_S$-modules de présentation finie.
D'après le Lemme \ref{limits-lemma-descend-modules-finite-presentation},
c'est donc l'image inverse d'une
application $\mathcal{O}_{S_i} \to \mathcal{F}_i$ de
$\mathcal{O}_{S_i}$-modules de présentation finie. En augmentant $i$,
on peut supposer cette application surjective (détails omis ; indication : appliquer
Algèbre, Lemme \ref{algebra-lemma-module-map-property-in-colimit}
sur un recouvrement ouvert affine). Le noyau de $\mathcal{O}_{S_i} \to \mathcal{F}_i$
est alors un idéal quasi-cohérent de type fini de $\mathcal{O}_{S_i}$
dont l'image inverse donne $\mathcal{I}$.
\end{proof}

\begin{lemma}
\label{limits-lemma-descend-module-flat-finite-presentation}
Reprenons les notations et les hypothèses du
Lemme \ref{limits-lemma-descend-finite-presentation}.
Soit $i \in I$.
Supposons que $\varphi_i : X_i \to Y_i$ soit un morphisme de schémas
de présentation finie sur $S_i$ et que $\mathcal{F}_i$ soit un
$\mathcal{O}_{X_i}$-module quasi-cohérent de présentation finie.
Si l'image inverse de $\mathcal{F}_i$ sur $X_i \times_{S_i} S$ est plate
sur $Y_i \times_{S_i} S$, il existe un indice $i' \geq i$
tel que l'image inverse de $\mathcal{F}_i$ sur $X_i \times_{S_i} S_{i'}$
soit plate sur $Y_i \times_{S_i} S_{i'}$.
\end{lemma}

\begin{proof}
(Ce lemme est l'analogue du
Lemme \ref{limits-lemma-descend-flat-finite-presentation}
pour les modules.)
Pour $i' \geq i$, notons $X_{i'} = S_{i'} \times_{S_i} X_i$,
$\mathcal{F}_{i'} = (X_{i'} \to X_i)^*\mathcal{F}_i$, et de même
pour $Y_{i'}$. Notons $\varphi_{i'}$ le changement de base
de $\varphi_i$ à $S_{i'}$. Posons aussi $X = S \times_{S_i} X_i$,
$Y =S \times_{S_i} X_i$, $\mathcal{F} = (X \to X_i)^*\mathcal{F}_i$
et notons $\varphi$ le changement de base de $\varphi_i$ à $S$.
Soit $Y_i = \bigcup_{j = 1, \ldots, m} V_{j, i}$ un recouvrement ouvert affine
fini tel que chaque $V_{j, i}$ ait son image dans un ouvert affine de $S_i$.
Pour chaque $j = 1, \ldots m$, soit
$\varphi_i^{-1}(V_{j, i}) = \bigcup_{k = 1, \ldots, m(j)} U_{k, j, i}$
un recouvrement ouvert affine fini. Pour $i' \geq i$, notons
$V_{j, i'}$ l'image réciproque de $V_{j, i}$ dans $Y_{i'}$ et
$U_{k, j, i'}$ celle de $U_{k, j, i}$ dans $X_{i'}$.
On a de même $U_{k, j} \subset X$ et $V_j \subset Y$.
Alors $U_{k, j} = \lim_{i' \geq i} U_{k, j, i'}$
et $V_j = \lim_{i' \geq i} V_j$
(voir Lemme \ref{limits-lemma-directed-inverse-system-has-limit}).
Puisque $X_{i'} = \bigcup_{k, j} U_{k, j, i'}$ est un recouvrement ouvert fini,
il suffit de démontrer le lemme pour chacun des morphismes
$U_{k, j, i} \to V_{j, i}$ et le faisceau $\mathcal{F}_i|_{U_{k, j, i}}$.
On se ramène donc au cas où $X_i$ et
$Y_i$ sont affines et ont leur image dans un ouvert affine de $S_i$, c'est-à-dire
que l'on peut aussi supposer $S$ affine.

\medskip\noindent
Dans le cas affine, on se ramène au résultat d'algèbre suivant.
Supposons que $R = \colim_{i \in I} R_i$. Pour un certain $i \in I$,
donnons-nous une application $A_i \to B_i$ de $R_i$-algèbres de présentation finie.
Soit $N_i$ un $B_i$-module de présentation finie.
Si $R \otimes_{R_i} N_i$ est plat sur $R \otimes_{R_i} A_i$,
alors, pour un certain $i' \geq i$, le module
$R_{i'} \otimes_{R_i} N_i$ est plat sur $R_{i'} \otimes_{R_i} A$.
C'est exactement le résultat démontré dans
Algèbre,
Lemme \ref{algebra-lemma-flat-finite-presentation-limit-flat} (3).
\end{proof}

\begin{lemma}
\label{limits-lemma-descend-finite-presentation-variant}
Pour un schéma $T$, notons $\mathcal{C}_T$ la sous-catégorie pleine des
schémas $W$ sur $T$ tels que $W$ soit quasi-compact et quasi-séparé
et que le morphisme structural $W \to T$ soit
localement de présentation finie.
Soit $S = \lim S_i$ une limite projective de schémas à morphismes
de transition affines. Il existe alors une équivalence
de catégories
$$
\colim \mathcal{C}_{S_i} \longrightarrow \mathcal{C}_S
$$
donnée par les foncteurs de changement de base.
\end{lemma}

\noindent
Attention : n'utilisez pas ce lemme sans comprendre la différence
entre celui-ci et le Lemme \ref{limits-lemma-descend-finite-presentation}.

\begin{proof}
Pleine fidélité. Donnons-nous $i \in I$ et des objets
$X_i$, $Y_i$ de $\mathcal{C}_{S_i}$. Notons
$X = X_i \times_{S_i} S$ et $Y = Y_i \times_{S_i} S$.
Donnons-nous un morphisme $f : X \to Y$ sur $S$.
On peut choisir un recouvrement ouvert affine fini
$Y_i = V_{i, 1} \cup \ldots \cup V_{i, m}$
tel que $V_{i, j} \to Y_i \to S_i$ ait son image dans un
ouvert affine $W_{i, j}$ de $S_i$. Notons $Y = V_1 \cup \ldots \cup V_m$ le
recouvrement ouvert affine induit sur $Y$.
Puisque $f : X \to Y$ est quasi-compact
(Schémas, Lemme \ref{schemes-lemma-quasi-compact-permanence}),
on peut, en augmentant $i$, supposer qu'il existe un
recouvrement ouvert fini $X_i = U_{i, 1} \cup \ldots \cup U_{i, m}$
par des ouverts quasi-compacts tels que l'image réciproque de
$U_{i, j}$ dans $Y$ soit $f^{-1}(V_j)$ ; voir
Lemme \ref{limits-lemma-descend-opens}.
D'après le Lemme \ref{limits-lemma-descend-finite-presentation}
appliqué à $f|_{f^{-1}(V_j)}$ sur $W_j$, on peut supposer, en
augmentant $i$, qu'il existe un morphisme
$f_{i, j} : V_{i, j} \to U_{i, j}$ sur $S$ dont le changement de base
à $S$ est $f|_{f^{-1}(V_j)}$.
En augmentant encore $i$, on peut supposer que $f_{i, j}$ et $f_{i, j'}$
coïncident sur l'ouvert quasi-compact $U_{i, j} \cap U_{i, j'}$.
On peut alors recoller ces morphismes pour obtenir
le morphisme cherché $f_i : X_i \to Y_i$.
Ce morphisme est unique (quitte à augmenter $i$),
car c'est le cas des morphismes $f_{i, j}$.

\medskip\noindent
Pour montrer que le foncteur est essentiellement surjectif, on raisonne
exactement de la même manière. Supposons en effet que
$X$ soit un objet de $\mathcal{C}_S$. Choisissons $i \in I$.
On peut choisir un recouvrement ouvert affine fini
$X = U_1 \cup \ldots \cup U_m$ tel que $U_j \to X \to S \to S_i$
se factorise par un ouvert affine $W_{i, j} \subset S_i$.
Posons $W_j = W_{i, j} \times_{S_i} S$. C'est un ouvert affine de $S$.
D'après le Lemme \ref{limits-lemma-descend-finite-presentation},
on peut, en augmentant $i$, supposer qu'il existe des morphismes
$U_{i, j} \to W_{i, j}$ de présentation finie
dont le changement de base à $W_j$ est $U_j$.
En augmentant $i$, on peut supposer qu'il existe
des ouverts quasi-compacts $U_{i, j, j'} \subset U_{i, j}$
dont les changements de base à $S$ sont égaux à $U_j \cap U_{j'}$.
Nous affirmons qu'en augmentant $i$, on peut supposer que l'image du morphisme
$U_{i, j, j'} \to U_{i, j} \to W_{i, j}$
est contenue dans $W_{i, j} \cap W_{i, j'}$.
En effet, le complémentaire de $W_{i, j} \cap W_{i, j'}$
est fermé dans le schéma affine $W_{i, j}$, donc affine.
Puisque $U_j \cap U_{j'} = \lim U_{i, j, j'}$ a bien son image dans
$W_{i, j} \cap W_{i, j'}$,
on peut appliquer le Lemme \ref{limits-lemma-limit-fibre-product-empty}
pour obtenir l'assertion. Ainsi, on peut considérer
$$
U_{i, j, j'} \quad\text{et}\quad U_{i, j', j}
$$
comme deux schémas sur $W_{i, j'}$ dont les changements de base à $W_{j'}$
redonnent $U_j \cap U_{j'}$. En augmentant $i$ et en utilisant le
Lemme \ref{limits-lemma-descend-finite-presentation},
on peut donc supposer qu'il existe des isomorphismes
$U_{i, j, j'} \to U_{i, j', j}$ sur $W_{i, j'}$, donc
sur $S_i$. En augmentant encore $i$ (détails omis),
on peut supposer que ces isomorphismes
satisfont la condition de cocycle mentionnée dans
Schémas, section \ref{schemes-section-glueing-schemes}.
En appliquant Schémas, Lemme \ref{schemes-lemma-glue},
on obtient un objet $X_i$ de $\mathcal{C}_{S_i}$ dont le
changement de base à $S$ est isomorphe à $X$ ; nous omettons
certaines vérifications.
\end{proof}












\section{Caractérisation des schémas affines}
\label{limits-section-affine}

\noindent
Si $f : X \to S$ est un morphisme entier surjectif de schémas
tel que $X$ soit affine, alors $S$ est lui aussi affine.
Voir \cite[A.2]{Conrad-Nagata}. Notre démonstration repose
sur le cas noethérien, énoncé et démontré dans Cohomologie des schémas,
Lemme \ref{coherent-lemma-image-affine-finite-morphism-affine-Noetherian}.
Voir aussi \cite[II 6.7.1]{EGA}.

\begin{lemma}
\label{limits-lemma-affine}
\begin{slogan}
Un schéma qui admet un morphisme fini surjectif provenant d'un schéma affine est affine.
\end{slogan}
Soit $f : X \to S$ un morphisme de schémas.
Supposons que $f$ soit surjectif et fini, et que $X$ soit affine.
Alors $S$ est affine.
\end{lemma}

\begin{proof}
Puisque $f$ est surjectif et $X$ quasi-compact, $S$ est
quasi-compact. Puisque $X$ est séparé et $f$ surjectif et
universellement fermé (Morphismes, Lemme
\ref{morphisms-lemma-integral-universally-closed}), $S$
est séparé (Morphismes, Lemme
\ref{morphisms-lemma-image-universally-closed-separated}).

\medskip\noindent
D'après le Lemme \ref{limits-lemma-finite-in-finite-and-finite-presentation},
on peut écrire $X = \lim_a X_a$, où $X_a \to S$ est fini et de présentation
finie. D'après le Lemme \ref{limits-lemma-limit-affine}, $X_a$
est affine pour un certain $a \in A$. En remplaçant $X$ par $X_a$, on peut supposer
que $X \to S$ est surjectif, fini et de présentation finie, et
que $X$ est affine.

\medskip\noindent
D'après la Proposition \ref{limits-proposition-approximate}, on peut écrire
$S = \lim_{i \in I} S_i$ comme
une limite projective de schémas de type fini sur $\mathbf{Z}$.
D'après le Lemme \ref{limits-lemma-descend-finite-presentation}, on peut,
en restreignant $I$, supposer qu'il existe des schémas $X_i \to S_i$
de présentation finie tels que $X_{i'} = X_i \times_S S_{i'}$
pour $i' \geq i$ et que $X = \lim_i X_i$. D'après le
Lemme \ref{limits-lemma-descend-finite-finite-presentation}, on peut
également supposer que $X_i \to S_i$ est fini pour tout $i \in I$.
En appliquant de nouveau le Lemme \ref{limits-lemma-limit-affine}, on peut supposer $X_i$
affine pour tout $i \in I$. Le résultat découle donc du
cas noethérien ; voir Cohomologie des schémas,
Lemme \ref{coherent-lemma-image-affine-finite-morphism-affine-Noetherian}.
\end{proof}

\begin{proposition}
\label{limits-proposition-affine}
\begin{slogan}
Un schéma qui admet un morphisme entier surjectif provenant d'un schéma affine est affine.
\end{slogan}
Soit $f : X \to S$ un morphisme de schémas. Supposons $X$ affine
et $f$ surjectif et universellement fermé\footnote{Un morphisme entier
est universellement fermé ; voir
Morphismes, Lemme \ref{morphisms-lemma-integral-universally-closed}.}.
Alors $S$ est affine.
\end{proposition}

\begin{proof}
D'après Morphismes, Lemme \ref{morphisms-lemma-image-universally-closed-separated},
le schéma $S$ est séparé. D'après
Morphismes, Lemme \ref{morphisms-lemma-affine-permanence},
$f$ est alors affine. Dès lors, d'après
Morphismes, Lemme \ref{morphisms-lemma-integral-universally-closed},
$f$ est entier.

\medskip\noindent
D'après le paragraphe précédent, on peut supposer $f : X \to S$
surjectif et entier, $X$ affine et $S$ séparé.
Puisque $f$ est surjectif et $X$ quasi-compact, $S$ est aussi
quasi-compact.

\medskip\noindent
D'après le Lemme \ref{limits-lemma-integral-limit-finite-and-finite-presentation},
on peut écrire $X = \lim_i X_i$, où $X_i \to S$ est fini. D'après le
Lemme \ref{limits-lemma-limit-affine},
pour $i$ assez grand, le schéma $X_i$ est affine.
De plus, puisque $X \to S$ se factorise par chaque $X_i$,
$X_i \to S$ est surjectif. On conclut donc que $S$ est affine par le
Lemme \ref{limits-lemma-affine}.
\end{proof}

\begin{lemma}
\label{limits-lemma-affines-glued-in-closed-affine}
Soit $X$ un schéma qui est, ensemblistement, réunion d'un
nombre fini de sous-schémas fermés affines. Alors $X$ est affine.
\end{lemma}

\begin{proof}
Soient $Z_i \subset X$, $i = 1, \ldots, n$, des sous-schémas fermés affines tels que
$X = \bigcup Z_i$ ensemblistement. Alors $\coprod Z_i \to X$ est surjectif
et entier, de source affine. Ainsi $X$ est affine par la
Proposition \ref{limits-proposition-affine}.
\end{proof}

\begin{lemma}
\label{limits-lemma-ample-on-reduction}
Soit $i : Z \to X$ une immersion fermée de schémas
induisant un homéomorphisme des espaces topologiques sous-jacents.
Soit $\mathcal{L}$ un faisceau inversible sur $X$.
Alors $i^*\mathcal{L}$ est ample sur $Z$ si et seulement si $\mathcal{L}$
est ample sur $X$.
\end{lemma}

\begin{proof}
Si $\mathcal{L}$ est ample, alors $i^*\mathcal{L}$ est ample, par
exemple d'après Morphismes, Lemme
\ref{morphisms-lemma-pullback-ample-tensor-relatively-ample}.
Supposons $i^*\mathcal{L}$ ample. Alors $Z$ est quasi-compact
(Propriétés, Définition \ref{properties-definition-ample})
et séparé
(Propriétés, Lemme \ref{properties-lemma-ample-separated}).
Puisque $i$ est surjectif, $X$ est quasi-compact.
Puisque $i$ est universellement fermé et surjectif,
$X$ est séparé (Morphismes, Lemme
\ref{morphisms-lemma-image-universally-closed-separated}).

\medskip\noindent
D'après la Proposition \ref{limits-proposition-approximate}, on peut écrire
$X = \lim X_i$ comme une limite projective de schémas de type fini sur $\mathbf{Z}$
à morphismes de transition affines. On peut trouver un $i$ et un faisceau
inversible $\mathcal{L}_i$ sur $X_i$ dont l'image inverse sur $X$ est isomorphe à
$\mathcal{L}$ ; voir Lemme \ref{limits-lemma-descend-modules-finite-presentation}.

\medskip\noindent
Pour chaque $i$, soit $Z_i \subset X_i$ l'image schématique
du morphisme $Z \to X_i$. Si $\Spec(A_i) \subset X_i$ est un sous-schéma ouvert
affine dont l'image réciproque est $\Spec(A)$ dans $X$ et si
$Z \cap \Spec(A)$ est défini par l'idéal $I \subset A$,
alors $Z_i \cap \Spec(A_i)$ est défini par l'idéal $I_i \subset A_i$,
image réciproque de $I$ dans $A_i$ par l'homomorphisme
$A_i \to A$ ; voir
Morphismes, Exemple \ref{morphisms-example-scheme-theoretic-image}.
Puisque $\colim A_i/I_i = A/I$, on en déduit $\lim Z_i = Z$.
D'après le Lemme \ref{limits-lemma-limit-ample}, $\mathcal{L}_i|_{Z_i}$
est ample pour un certain $i$. Puisque $Z$, et donc $X$, a son image dans $Z_i$
ensemblistement, $X_{i'} \to X_i$ a son image dans $Z_i$
ensemblistement pour un certain $i' \geq i$ ; voir
Lemme \ref{limits-lemma-limit-contained-in-constructible}.
(Comme $X_i$ est noethérien, toute partie fermée
de $X_i$ est constructible.) Soit $T \subset X_{i'}$
l'image réciproque schématique de $Z_i$ dans $X_{i'}$.
Remarquons que $\mathcal{L}_{i'}|_T$ est l'image inverse
de $\mathcal{L}_i|_{Z_i}$, donc est ample d'après
Morphismes, Lemme \ref{morphisms-lemma-pullback-ample-tensor-relatively-ample},
car $T \to Z_i$ est un morphisme affine.
Ainsi $\mathcal{L}_{i'}$ est ample sur $X_{i'}$
d'après Cohomologie des schémas, Lemme \ref{coherent-lemma-ample-on-reduction}.
En prenant l'image inverse sur $X$ (à l'aide du même lemme que ci-dessus),
on trouve que $\mathcal{L}$ est ample.
\end{proof}

\begin{lemma}
\label{limits-lemma-thickening-quasi-affine}
Soit $i : Z \to X$ une immersion fermée de schémas
induisant un homéomorphisme des espaces topologiques sous-jacents.
Alors $X$ est quasi-affine si et seulement si $Z$ est quasi-affine.
\end{lemma}

\begin{proof}
Rappelons qu'un schéma est quasi-affine
si et seulement si son faisceau structural est ample ; voir
Propriétés, Lemme \ref{properties-lemma-quasi-affine-O-ample}.
Ainsi, si $Z$ est quasi-affine, $\mathcal{O}_Z$ est ample,
donc $\mathcal{O}_X$ est ample par le
Lemme \ref{limits-lemma-ample-on-reduction}, et
$X$ est donc quasi-affine. La réciproque, qui
se démontre aussi de manière élémentaire, s'obtient
en lisant le raisonnement précédent en sens inverse.
\end{proof}

\noindent
Le lemme suivant n'a pas vraiment sa place dans cette section.

\begin{lemma}
\label{limits-lemma-ample-profinite-set-in-principal-affine}
Soit $X$ un schéma. Soit $\mathcal{L}$ un faisceau inversible ample sur $X$.
Supposons donnés des morphismes de schémas
$$
\Spec(k) \leftarrow \Spec(A) \to W \subset X
$$
où $k$ est un corps, $A$ une $k$-algèbre entière et $W$ un ouvert de $X$.
Il existe alors un $n > 0$ et une section
$s \in \Gamma(X, \mathcal{L}^{\otimes n})$ tels que
$X_s$ soit affine, $X_s \subset W$, et que $\Spec(A) \to W$ se factorise par $X_s$.
\end{lemma}

\begin{proof}
Puisque $\Spec(A)$ est quasi-compact, on peut remplacer $W$ par un ouvert
quasi-compact contenant encore l'image de $\Spec(A) \to X$.
Rappelons que $X$ est quasi-séparé et quasi-compact du fait
qu'il possède un faisceau inversible ample ; voir
Propriétés, Définition \ref{properties-definition-ample} et
Lemme \ref{properties-lemma-affine-s-opens-cover-quasi-separated}.
D'après la Proposition \ref{limits-proposition-approximate}, on peut
écrire $X = \lim X_i$ comme la limite d'un système projectif filtrant
de schémas de type fini sur $\mathbf{Z}$
à morphismes de transition affines.
Pour un certain $i$, le faisceau inversible ample $\mathcal{L}$ sur $X$
descend en un faisceau inversible ample $\mathcal{L}_i$ sur $X_i$,
et l'ouvert $W$ est l'image réciproque d'un ouvert
quasi-compact $W_i \subset X_i$ ; voir
Lemmes \ref{limits-lemma-limit-ample}, \ref{limits-lemma-descend-invertible-modules} et
\ref{limits-lemma-descend-opens}.
On peut remplacer $X, W, \mathcal{L}$ par $X_i, W_i, \mathcal{L}_i$
et supposer $X$ de présentation finie sur $\mathbf{Z}$.
Écrivons $A = \colim A_j$ comme la limite inductive de ses $k$-sous-algèbres finies.
Pour un certain $j$, le morphisme $\Spec(A) \to X$ se factorise alors par
un morphisme $\Spec(A_j) \to X$ ; voir
Proposition \ref{limits-proposition-characterize-locally-finite-presentation}.
Puisque $\Spec(A_j)$ est fini, on est ramené à
Propriétés, Lemme \ref{properties-lemma-ample-finite-set-in-principal-affine}.
\end{proof}



















\section{Variantes du lemme de Chow}
\label{limits-section-chows-lemma}

\noindent
Dans cette section, nous démontrons plusieurs variantes du lemme de Chow.
La version la plus intéressante est probablement le cas noethérien,
énoncé et démontré dans
Cohomologie des schémas, section \ref{coherent-section-chows-lemma}.

\begin{lemma}
\label{limits-lemma-chow-finite-type}
Soit $S$ un schéma quasi-compact et quasi-séparé.
Soit $f : X \to S$ un morphisme séparé de type fini.
Il existe alors un $n \geq 0$ et un diagramme
$$
\xymatrix{
X \ar[rd] & X' \ar[d] \ar[l]^\pi \ar[r] & \mathbf{P}^n_S \ar[dl] \\
& S &
}
$$
où $X' \to \mathbf{P}^n_S$ est une immersion et
$\pi : X' \to X$ est propre et surjectif.
\end{lemma}

\begin{proof}
D'après la Proposition \ref{limits-proposition-separated-closed-in-finite-presentation},
on peut trouver une immersion fermée $X \to Y$, où $Y$ est séparé
et de présentation finie sur $S$. Si l'on démontre l'assertion
pour $Y$, le résultat pour $X$ s'en déduit immédiatement. On peut donc supposer
$X$ de présentation finie sur $S$.

\medskip\noindent
Écrivons $S = \lim_i S_i$ comme une limite projective de schémas noethériens ; voir
Proposition \ref{limits-proposition-approximate}. D'après le
Lemme \ref{limits-lemma-descend-finite-presentation}, on peut
trouver un indice $i \in I$ et un schéma $X_i \to S_i$ de présentation finie
tels que $X = S \times_{S_i} X_i$.
D'après le Lemme \ref{limits-lemma-descend-separated-finite-presentation},
on peut supposer $X_i \to S_i$ séparé.
Si l'on démontre l'assertion pour
$X_i$ sur $S_i$, elle vaut évidemment pour $X$. Le cas
$X_i \to S_i$ est traité dans
Cohomologie des schémas, Lemme \ref{coherent-lemma-chow-Noetherian}.
\end{proof}

\begin{remark}
\label{limits-remark-chow-finite-type}
Dans la situation du Lemme de Chow \ref{limits-lemma-chow-finite-type} :
\begin{enumerate}
\item Le morphisme $\pi$ est en fait H-projectif (donc projectif ; voir
Morphismes, Lemme \ref{morphisms-lemma-H-projective}),
puisque le morphisme $X' \to \mathbf{P}^n_S \times_S X = \mathbf{P}^n_X$
est une immersion fermée (utiliser le fait que $\pi$ est propre ; voir
Morphismes, Lemme \ref{morphisms-lemma-image-proper-scheme-closed}).
\item On peut supposer $X'$ réduit, car on peut remplacer $X'$ par
son réduit sans changer les autres assertions du lemme.
\item On peut supposer $X' \to X$ de présentation finie
sans changer les autres assertions du lemme. Cela se déduit
de la démonstration du Lemme \ref{limits-lemma-chow-finite-type}, mais on peut aussi le démontrer
directement comme suit. D'après (1), on a une immersion fermée
$X' \to \mathbf{P}^n_X$. D'après le
Lemme \ref{limits-lemma-closed-is-limit-closed-and-finite-presentation}, on peut écrire
$X' = \lim X'_i$, où $X'_i \to \mathbf{P}^n_X$ est une immersion
fermée de présentation finie. En particulier, $X'_i \to X$ est
de présentation finie, propre et surjectif. Pour
$i$ assez grand, le morphisme $X'_i \to \mathbf{P}^n_S$ est une immersion par le
Lemme \ref{limits-lemma-finite-type-eventually-closed}. En remplaçant
$X'$ par $X'_i$, on obtient le résultat voulu.
\end{enumerate}
Bien entendu, on ne peut pas, en général, obtenir simultanément (2) et (3).
\end{remark}

\noindent
Voici une variante du lemme de Chow dans laquelle on suppose que le schéma
source possède un nombre fini de composantes irréductibles.

\begin{lemma}
\label{limits-lemma-chow-EGA}
Soit $S$ un schéma quasi-compact et quasi-séparé.
Soit $f : X \to S$ un morphisme séparé de type fini.
Supposons que $X$ possède un nombre fini de composantes irréductibles.
Il existe alors un $n \geq 0$ et un diagramme
$$
\xymatrix{
X \ar[rd] & X' \ar[d] \ar[l]^\pi \ar[r] & \mathbf{P}^n_S \ar[dl] \\
& S &
}
$$
où $X' \to \mathbf{P}^n_S$ est une immersion et
$\pi : X' \to X$ est propre et surjectif. De plus, il existe
un sous-schéma ouvert dense $U \subset X$ tel que $\pi^{-1}(U) \to U$
soit un isomorphisme de schémas.
\end{lemma}

\begin{proof}
Soit $X = Z_1 \cup \ldots \cup Z_n$ la décomposition de $X$
en composantes irréductibles. Soit $\eta_j \in Z_j$ le point générique.

\medskip\noindent
Il y a au moins deux façons de poursuivre la démonstration.
La première consiste à reprendre la démonstration de
Cohomologie des schémas, Lemme \ref{coherent-lemma-chow-Noetherian},
en utilisant le résultat général
Propriétés, Lemme \ref{properties-lemma-point-and-maximal-points-affine},
pour trouver des ouverts affines convenables de $X$. (C'est la démonstration « standard ».)
La seconde consiste à utiliser l'approximation noethérienne absolue, comme dans
la démonstration du Lemme \ref{limits-lemma-chow-finite-type} ci-dessus.
C'est ce que nous allons faire ici.

\medskip\noindent
D'après la Proposition \ref{limits-proposition-separated-closed-in-finite-presentation},
on peut trouver une immersion fermée $X \to Y$, où $Y$ est séparé
et de présentation finie sur $S$.
Écrivons $S = \lim_i S_i$ comme une limite projective de schémas noethériens ; voir
Proposition \ref{limits-proposition-approximate}. D'après le
Lemme \ref{limits-lemma-descend-finite-presentation}, on peut
trouver un indice $i \in I$ et un schéma $Y_i \to S_i$ de présentation finie
tels que $Y = S \times_{S_i} Y_i$.
D'après le Lemme \ref{limits-lemma-descend-separated-finite-presentation},
on peut supposer $Y_i \to S_i$ séparé.
On a le diagramme suivant :
$$
\xymatrix{
\eta_j \in Z_j \ar[r] & X \ar[r] \ar[rd] & Y \ar[r] \ar[d] & Y_i \ar[d] \\
& & S \ar[r] & S_i
}
$$
Notons $h : X \to Y_i$ la composée.

\medskip\noindent
Pour $i' \geq i$, posons $Y_{i'} = S_{i'} \times_{S_i} Y_i$.
Alors $Y = \lim_{i' \geq i} Y_{i'}$ ; voir
Lemme \ref{limits-lemma-scheme-over-limit}.
Choisissons $j, j' \in \{1, \ldots, n\}$, $j \not = j'$.
Remarquons que $\eta_j$ n'est pas une spécialisation de $\eta_{j'}$.
D'après le Lemme \ref{limits-lemma-topology-limit},
on peut remplacer $i$ par un indice plus grand et supposer
que $h(\eta_j)$ n'est pas une spécialisation de $h(\eta_{j'})$
pour tous les couples $(j, j')$ ci-dessus.
Pour un tel indice, soit
$Y' \subset Y_i$ l'image schématique de
$h : X \to Y_i$ ; voir
Morphismes, Définition \ref{morphisms-definition-scheme-theoretic-image}.
Le morphisme $h$ est quasi-compact, étant la composée des morphismes quasi-compacts
$X \to Y$ et $Y \to Y_i$ (ce dernier est affine).
D'après
Morphismes, Lemme \ref{morphisms-lemma-quasi-compact-scheme-theoretic-image},
le morphisme $X \to Y'$ est donc dominant. Ainsi, les points génériques
de $Y'$ appartiennent tous à l'ensemble
$\{h(\eta_1), \ldots, h(\eta_n)\}$ ; voir
Morphismes, Lemme \ref{morphisms-lemma-quasi-compact-dominant}.
Puisqu'aucun des $h(\eta_j)$ n'est une spécialisation d'un autre,
les points $h(\eta_1), \ldots, h(\eta_n)$ sont deux à deux
distincts et chacun est un point générique de $Y'$.

\medskip\noindent
Appliquons Cohomologie des schémas, Lemme
\ref{coherent-lemma-chow-Noetherian}, au morphisme
$Y' \to S_i$. On obtient un diagramme
$$
\xymatrix{
Y' \ar[rd] & Y^* \ar[d] \ar[l]^\pi \ar[r] & \mathbf{P}^n_{S_i} \ar[dl] \\
& S_i &
}
$$
tel que $\pi$ soit propre et surjectif, et un isomorphisme au-dessus
d'un sous-schéma ouvert dense $V \subset Y'$. Par le choix de $i$ ci-dessus,
on sait que $h(\eta_1), \ldots, h(\eta_n) \in V$. Considérons
le diagramme commutatif
$$
\xymatrix{
X' \ar@{=}[r] &
X \times_{Y'} Y^* \ar[r] \ar[d] &
Y^* \ar[r] \ar[d] &
\mathbf{P}^n_{S_i} \ar[ddl] \\
& X \ar[r] \ar[d] & Y' \ar[d] & \\
& S \ar[r] & S_i &
}
$$
Remarquons que $X' \to X$ est un isomorphisme au-dessus du sous-schéma ouvert
$U = h^{-1}(V)$, qui contient chacun des $\eta_j$, donc est
dense dans $X$. On conclut que $X \leftarrow X' \rightarrow \mathbf{P}^n_S$
résout le problème posé dans le lemme.
\end{proof}













\section{Applications du lemme de Chow}
\label{limits-section-apply-chow}

\noindent
Voici une première application du lemme de Chow.

\begin{lemma}
\label{limits-lemma-eventually-proper}
\begin{slogan}
Si le changement de base d'un schéma à une limite est propre, alors
le changement de base est déjà propre à un étage fini.
\end{slogan}
Reprenons les hypothèses et les notations de la Situation \ref{limits-situation-descent-property}.
Si
\begin{enumerate}
\item $f$ est propre,
\item $f_0$ est localement de type fini,
\end{enumerate}
alors il existe un $i$ tel que $f_i$ soit propre.
\end{lemma}

\begin{proof}
D'après le Lemme \ref{limits-lemma-descend-separated-finite-presentation},
$f_i$ est séparé pour un certain $i \geq 0$. En remplaçant
$0$ par $i$, on peut supposer $f_0$ séparé.
Remarquons que $f_0$ est quasi-compact ; voir
Schémas, Lemme \ref{schemes-lemma-quasi-compact-permanence}.
D'après le Lemme \ref{limits-lemma-chow-finite-type}, on peut choisir un diagramme
$$
\xymatrix{
X_0 \ar[rd] & X_0' \ar[d] \ar[l]^\pi \ar[r] & \mathbf{P}^n_{Y_0} \ar[dl] \\
& Y_0 &
}
$$
où $X_0' \to \mathbf{P}^n_{Y_0}$ est une immersion et
$\pi : X_0' \to X_0$ est propre et surjectif. Posons
$X' = X_0' \times_{Y_0} Y$ et $X_i' = X_0' \times_{Y_0} Y_i$.
D'après Morphismes, Lemmes \ref{morphisms-lemma-composition-proper} et
\ref{morphisms-lemma-base-change-proper},
$X' \to Y$ est propre. Ainsi $X' \to \mathbf{P}^n_Y$ est
une immersion fermée (Morphismes, Lemme
\ref{morphisms-lemma-image-proper-scheme-closed}). D'après
Morphismes, Lemme \ref{morphisms-lemma-image-proper-is-proper},
il suffit de montrer que $X'_i \to Y_i$ est propre pour un certain $i$.
D'après le Lemme \ref{limits-lemma-descend-closed-immersion-finite-presentation},
$X'_i \to \mathbf{P}^n_{Y_i}$ est
une immersion fermée pour $i$ assez grand. Alors $X'_i \to Y_i$
est propre, ce qui conclut.
\end{proof}

\begin{lemma}
\label{limits-lemma-proper-limit-of-proper-finite-presentation}
Soit $f : X \to S$ un morphisme propre, où $S$ est quasi-compact et
quasi-séparé. Alors $X = \lim X_i$ est une limite projective de schémas
$X_i$ propres et de présentation finie sur $S$, telle que
tous les morphismes de transition et les morphismes $X \to X_i$ soient des immersions
fermées.
\end{lemma}

\begin{proof}
D'après la Proposition \ref{limits-proposition-separated-closed-in-finite-presentation},
on peut trouver une immersion fermée $X \to Y$, où $Y$ est séparé et de
présentation finie sur $S$. D'après le Lemme \ref{limits-lemma-chow-finite-type},
on peut trouver un diagramme
$$
\xymatrix{
Y \ar[rd] & Y' \ar[d] \ar[l]^\pi \ar[r] & \mathbf{P}^n_S \ar[dl] \\
& S &
}
$$
où $Y' \to \mathbf{P}^n_S$ est une immersion et
$\pi : Y' \to Y$ est propre et surjectif. D'après le
Lemme \ref{limits-lemma-closed-is-limit-closed-and-finite-presentation},
on peut écrire $X = \lim X_i$, où $X_i \to Y$ est une immersion fermée de
présentation finie. Notons $X'_i \subset Y'$, resp.\ $X' \subset Y'$,
l'image réciproque schématique de $X_i \subset Y$, resp.\ $X \subset Y$.
Alors $\lim X'_i = X'$. Puisque $X' \to S$ est propre
(Morphismes, Lemmes \ref{morphisms-lemma-composition-proper}),
$X' \to \mathbf{P}^n_S$ est une immersion fermée (Morphismes, Lemme
\ref{morphisms-lemma-image-proper-scheme-closed}). Pour $i$ assez grand,
$X'_i \to \mathbf{P}^n_S$ est donc une immersion fermée d'après le
Lemme \ref{limits-lemma-eventually-closed-immersion}.
Ainsi $X'_i$ est propre sur $S$.
Pour un tel $i$, le morphisme $X_i \to S$ est propre d'après
Morphismes, Lemme \ref{morphisms-lemma-image-proper-is-proper}.
\end{proof}

\begin{lemma}
\label{limits-lemma-proper-limit-of-proper-finite-presentation-noetherian}
Soit $f : X \to S$ un morphisme propre, où $S$ est quasi-compact et
quasi-séparé. Il existe alors un ensemble ordonné filtrant $I$ et un
système projectif $(f_i : X_i \to S_i)$ de morphismes de schémas indexé par $I$,
tels que les morphismes de transition $X_i \to X_{i'}$ et $S_i \to S_{i'}$
soient affines, que $f_i$ soit propre, que $S_i$ soit de type
fini sur $\mathbf{Z}$ et que
$(X \to S) = \lim (X_i \to S_i)$.
\end{lemma}

\begin{proof}
D'après le Lemme \ref{limits-lemma-proper-limit-of-proper-finite-presentation},
on peut écrire $X = \lim_{k \in K} X_k$, où $X_k \to S$ est propre et
de présentation finie. Ensuite, par approximation noethérienne absolue
(Proposition \ref{limits-proposition-approximate}), on peut
écrire $S = \lim_{j \in J} S_j$, où $S_j$ est de type fini sur $\mathbf{Z}$.
Pour chaque $k$, il existe un $j$ et un morphisme $X_{k, j} \to S_j$
de présentation finie tels que $X_k \cong S \times_{S_j} X_{k, j}$
comme schémas sur $S$ ; voir
Lemme \ref{limits-lemma-descend-finite-presentation}.
En augmentant $j$, on peut supposer $X_{k, j} \to S_j$ propre ; voir
Lemme \ref{limits-lemma-eventually-proper}. L'ensemble $I$ sera constitué
de ces couples $(k, j)$, et le morphisme correspondant sera $X_{k, j} \to S_j$.
Pour tout $k' \geq k$, on peut trouver un $j' \geq j$ et un morphisme
$X_{j', k'} \to X_{j, k}$ au-dessus de $S_{j'} \to S_j$ dont le changement de base à $S$
donne le morphisme $X_{k'} \to X_k$ (cela découle encore du
Lemme \ref{limits-lemma-descend-finite-presentation}).
Ces morphismes constituent les morphismes de transition du système. Certains détails
sont omis.
\end{proof}

\begin{lemma}
\label{limits-lemma-finite-type-eventually-proper}
Soit $S$ un schéma. Soit $X = \lim X_i$ une limite projective de
schémas sur $S$ à morphismes de transition affines. Soit $Y \to X$
un morphisme de schémas sur $S$.
Si $Y \to X$ est propre, si les $X_i$ sont quasi-compacts et quasi-séparés et si
$Y$ est localement de type fini sur $S$, alors $Y \to X_i$ est propre
pour $i$ assez grand.
\end{lemma}

\begin{proof}
Choisissons une immersion fermée $Y \to Y'$, où $Y'$ est propre et de présentation
finie sur $X$ ; voir
Lemme \ref{limits-lemma-proper-limit-of-proper-finite-presentation}.
Choisissons ensuite un $i$ et un morphisme propre $Y'_i \to X_i$
tels que $Y' = X \times_{X_i} Y'_i$. C'est possible d'après les
Lemmes \ref{limits-lemma-descend-finite-presentation} et
\ref{limits-lemma-eventually-proper}. En remplaçant $i$
par un indice plus grand, $Y \to Y'_i$ est alors une immersion
fermée ; voir Lemme \ref{limits-lemma-finite-type-eventually-closed}.
\end{proof}

\noindent
Rappelons la notion de support schématique
d'un module quasi-cohérent de type fini ; voir
Morphismes, Définition \ref{morphisms-definition-scheme-theoretic-support}.

\begin{lemma}
\label{limits-lemma-eventually-proper-support}
Reprenons les hypothèses et les notations de la Situation \ref{limits-situation-descent-property}.
Soit $\mathcal{F}_0$ un $\mathcal{O}_{X_0}$-module quasi-cohérent.
Notons $\mathcal{F}$ et $\mathcal{F}_i$ les images inverses de
$\mathcal{F}_0$ sur $X$ et $X_i$. Faisons les hypothèses suivantes :
\begin{enumerate}
\item $f_0$ est localement de type fini,
\item $\mathcal{F}_0$ est de type fini,
\item le support schématique de $\mathcal{F}$ est propre sur $Y$.
\end{enumerate}
Alors le support schématique de $\mathcal{F}_i$ est propre sur $Y_i$
pour un certain $i$.
\end{lemma}

\begin{proof}
On peut remplacer $X_0$ par le support schématique de $\mathcal{F}_0$.
D'après Morphismes, Lemme \ref{morphisms-lemma-support-finite-type}, cela
garantit que $X_i$ est le support de $\mathcal{F}_i$ et $X$ le
support de $\mathcal{F}$. Si $Z \subset X$ désigne le support
schématique de $\mathcal{F}$, alors $Z \to X$ est un homéomorphisme
universel. On conclut que $X \to Y$ est propre, puisque
$Z \to Y$ l'est par hypothèse ; voir
Morphismes, Lemme \ref{morphisms-lemma-image-proper-is-proper}.
D'après le Lemme \ref{limits-lemma-eventually-proper}, $X_i \to Y$ est propre
pour un certain $i$. Le support schématique $Z_i$ de
$\mathcal{F}_i$ est donc propre sur $Y$ d'après
Morphismes, Lemmes \ref{morphisms-lemma-closed-immersion-proper} et
\ref{morphisms-lemma-composition-proper}.
\end{proof}











\section{Morphismes universellement fermés}
\label{limits-section-universally-closed}

\noindent
Dans cette section, nous étudions quand un morphisme quasi-compact (mais
pas nécessairement séparé) est universellement fermé. Nous démontrons d'abord
un lemme permettant de tester cette propriété après un changement de base
localement de présentation finie.

\begin{lemma}
\label{limits-lemma-separate}
Soit $f : X \to S$ un morphisme quasi-compact de schémas.
Soit $g : T \to S$ un morphisme de schémas.
Soit $t \in T$ un point et $Z \subset X_T$ un sous-schéma
fermé tel que $Z \cap X_t = \emptyset$.
Il existe alors un voisinage ouvert
$V \subset T$ de $t$, un diagramme commutatif
$$
\xymatrix{
V \ar[d] \ar[r]_a & T' \ar[d]^b \\
T \ar[r]^g & S,
}
$$
et un sous-schéma fermé $Z' \subset X_{T'}$ tels que
\begin{enumerate}
\item le morphisme $b : T' \to S$ soit localement de présentation finie,
\item en posant $t' = a(t)$, on ait $Z' \cap X_{t'} = \emptyset$,
\item $Z \cap X_V$ ait son image dans $Z'$ par le morphisme $X_V \to X_{T'}$.
\end{enumerate}
De plus, on peut supposer $V$ et $T'$ affines.
\end{lemma}

\begin{proof}
Posons $s = g(t)$. Au cours de la démonstration, on peut toujours remplacer $T$ par un
voisinage ouvert de $t$. On peut donc aussi remplacer $S$ par un voisinage ouvert
de $s$. On peut ainsi supposer, et l'on suppose, $T$ et $S$ affines.
Écrivons $S = \Spec(A)$, $T = \Spec(B)$ ; soit $g$ donné par
l'homomorphisme $A \to B$, et soit $t$ associé à l'idéal premier
$\mathfrak q \subset B$.

\medskip\noindent
Puisque $X \to S$ est quasi-compact et $S$ affine, on peut écrire
$X = \bigcup_{i = 1, \ldots, n} U_i$ comme réunion finie d'ouverts affines.
Écrivons $U_i = \Spec(C_i)$. On a en particulier
$X_T = \bigcup_{i = 1, \ldots, n} U_{i, T} =
\bigcup_{i = 1, \ldots n} \Spec(C_i \otimes_A B)$.
Soit $I_i \subset C_i \otimes_A B$ l'idéal correspondant au
sous-schéma fermé $Z \cap U_{i, T}$. La condition
$Z \cap X_t = \emptyset$ signifie que $I_i$ engendre
l'idéal unité de l'anneau
$$
C_i \otimes_A \kappa(\mathfrak q) =
(B \setminus \mathfrak q)^{-1}\left(
C_i \otimes_A B/\mathfrak q C_i \otimes_A B \right)
$$
Puisque $I_i (B \setminus \mathfrak q)^{-1}(C_i \otimes_A B) =
(B \setminus \mathfrak q)^{-1} I_i$, cela signifie que $1 = x_i/g_i$
pour certains $x_i \in I_i$ et $g_i \in B$, $g_i \not \in \mathfrak q$.
En chassant les dénominateurs, on trouve donc une relation de la forme
$$
x_i + \sum\nolimits_j f_{i, j}c_{i, j} = g_i
$$
avec $x_i \in I_i$, $f_{i, j} \in \mathfrak q$, $c_{i, j} \in C_i \otimes_A B$
et $g_i \in B$, $g_i \not \in \mathfrak q$. En remplaçant $B$ par
$B_{g_1 \ldots g_n}$, c'est-à-dire $T$ par un voisinage affine plus petit
de $t$, on peut supposer les équations de la forme
$$
x_i + \sum\nolimits_j f_{i, j}c_{i, j} = 1
$$
avec $x_i \in I_i$, $f_{i, j} \in \mathfrak q$, $c_{i, j} \in C_i \otimes_A B$.

\medskip\noindent
Pour terminer, écrivons $B$ comme limite inductive de
$A$-algèbres de présentation finie $B_\lambda$ indexée par un ensemble ordonné filtrant $\Lambda$.
Pour chaque $\lambda$, posons
$\mathfrak q_\lambda = (B_\lambda \to B)^{-1}(\mathfrak q)$.
Pour $\lambda \in \Lambda$ assez grand, on peut trouver
\begin{enumerate}
\item un élément
$x_{i, \lambda} \in C_i \otimes_A B_\lambda$ d'image $x_i$,
\item des éléments $f_{i, j, \lambda} \in \mathfrak q_{i, \lambda}$
d'images $f_{i, j}$,
\item des éléments $c_{i, j, \lambda} \in C_i \otimes_A B_\lambda$
d'images $c_{i, j}$.
\end{enumerate}
En augmentant encore $\lambda$, l'équation
$$
x_{i, \lambda} + \sum\nolimits_j f_{i, j, \lambda}c_{i, j, \lambda} = 1
$$
sera satisfaite. Fixons un tel $\lambda$ et posons $T' = \Spec(B_\lambda)$.
Alors $t' \in T'$ est le point associé à l'idéal premier $\mathfrak q_\lambda$.
Enfin, soit $Z' \subset X_{T'}$ l'image schématique de
$Z \to X_T \to X_{T'}$. Puisque $X_T \to X_{T'}$ est affine, on peut calculer $Z'$
sur les ouverts affines $U_{i, T'}$ comme le sous-schéma fermé associé
à $\Ker(C_i \otimes_A B_\lambda \to C_i \otimes_A B/I_i)$ ; voir
Morphismes, Exemple \ref{morphisms-example-scheme-theoretic-image}.
Ainsi $x_{i, \lambda}$ appartient à l'idéal définissant $Z'$. La dernière
équation affichée montre donc que $Z' \cap X_{t'}$ est vide.
\end{proof}

\begin{lemma}
\label{limits-lemma-test-universally-closed}
Soit $f : X \to S$ un morphisme quasi-compact de schémas.
Les conditions suivantes sont équivalentes :
\begin{enumerate}
\item $f$ est universellement fermé,
\item pour tout morphisme $S' \to S$ localement de présentation finie,
le changement de base $X_{S'} \to S'$ est fermé,
\item pour tout $n$, le morphisme
$\mathbf{A}^n \times X \to \mathbf{A}^n \times S$
est fermé.
\end{enumerate}
\end{lemma}

\begin{proof}
Il est clair que (1) entraîne (2). Montrons que (2) entraîne (1).
Supposons que le changement de base $X_T \to T$ ne soit pas fermé pour un certain
schéma $T$ sur $S$. D'après
Schémas, Lemme \ref{schemes-lemma-quasi-compact-closed},
il existe alors une spécialisation $t_1 \leadsto t$ dans
$T$ et un point $\xi \in X_T$ d'image $t_1$ tel que $\xi$ ne
se spécialise en aucun point de la fibre au-dessus de $t$. Posons
$Z = \overline{\{\xi\}} \subset X_T$. Alors $Z \cap X_t = \emptyset$. Appliquons le
Lemme \ref{limits-lemma-separate}.
On trouve un voisinage ouvert $V \subset T$ de $t$, un diagramme commutatif
$$
\xymatrix{
V \ar[d] \ar[r]_a & T' \ar[d]^b \\
T \ar[r]^g & S,
}
$$
et un sous-schéma fermé $Z' \subset X_{T'}$ tels que
\begin{enumerate}
\item le morphisme $b : T' \to S$ soit localement de présentation finie,
\item en posant $t' = a(t)$, on ait $Z' \cap X_{t'} = \emptyset$,
\item $Z \cap X_V$ ait son image dans $Z'$ par le morphisme $X_V \to X_{T'}$.
\end{enumerate}
Cela signifie que $X_{T'} \to T'$ envoie le fermé $Z'$
sur une partie de $T'$ contenant $a(t_1)$ mais non $t' = a(t)$.
Puisque $a(t_1) \leadsto a(t) = t'$, on conclut que $X_{T'} \to T'$
n'est pas fermé. Ainsi, si $X \to S$ n'est pas universellement fermé,
$X_{T'} \to T'$ n'est pas fermé pour un certain $T' \to S$
localement de présentation finie. Autrement dit, (2)
entraîne (1).

\medskip\noindent
Supposons que $\mathbf{A}^n \times X \to \mathbf{A}^n \times S$ soit
fermé pour tout entier $n$. Nous voulons montrer que $X_T \to T$ est
fermé pour tout schéma $T$ localement de présentation finie
sur $S$. On peut bien sûr supposer $T$ affine et son image contenue dans
un ouvert affine $V$ de $S$ (car le fait que $X_T \to T$ soit fermé est local sur $T$).
Dans ce cas, il existe une immersion fermée $T \to \mathbf{A}^n \times V$,
car $\mathcal{O}_T(T)$ est une
$\mathcal{O}_S(V)$-algèbre de présentation finie ; voir
Morphismes,
Lemme \ref{morphisms-lemma-locally-finite-presentation-characterize}.
Alors $T \to \mathbf{A}^n \times S$ est une immersion localement fermée.
On obtient donc un diagramme cartésien
$$
\xymatrix{
X_T \ar[d]_{f_T} \ar[r] & \mathbf{A}^n \times X \ar[d]^{f_n} \\
T \ar[r] & \mathbf{A}^n \times S
}
$$
de schémas dont les flèches horizontales sont des immersions localement fermées.
Toute partie fermée $Z \subset X_T$ peut donc s'écrire
$X_T \cap Z'$ pour une partie fermée $Z' \subset \mathbf{A}^n \times X$.
Alors $f_T(Z) = T \cap f_n(Z')$, et si $f_n$ est fermé,
$f_T$ l'est aussi.
\end{proof}

\begin{lemma}
\label{limits-lemma-limited-base-change}
Soit $S$ un schéma.
Soit $f : X \to S$ un morphisme séparé de type fini.
Les conditions suivantes sont équivalentes :
\begin{enumerate}
\item Le morphisme $f$ est propre.
\item Pour tout morphisme $S' \to S$ localement de type fini,
le changement de base $X_{S'} \to S'$ est fermé.
\item Pour tout $n \geq 0$, le morphisme
$\mathbf{A}^n \times X \to \mathbf{A}^n \times S$ est fermé.
\end{enumerate}
\end{lemma}

\begin{proof}[Première démonstration]
Puisqu'un morphisme propre est exactement un morphisme
séparé, de type fini et universellement fermé, ce
lemme est un cas particulier du Lemme \ref{limits-lemma-test-universally-closed}.
\end{proof}

\begin{proof}[Seconde démonstration]
Il est clair que (1) entraîne (2) et que (2) entraîne (3) ; il suffit donc de montrer que (3)
entraîne (1).
Ramenons-nous d'abord au cas où $S$ est affine. Supposons que (3) entraîne (1)
lorsque la base est affine. Soit maintenant $f: X \to S$ un morphisme séparé de
type fini. La propreté est locale sur la base
(voir Morphismes, Lemme \ref{morphisms-lemma-proper-local-on-the-base}) ; si
$S = \bigcup_\alpha S_\alpha$ est un recouvrement ouvert affine et si
l'on note $X_\alpha := f^{-1}(S_\alpha)$, il
suffit donc de montrer que $f|_{X_\alpha}: X_\alpha \to S_\alpha$ est propre
pour tout $\alpha$. Puisque $S_\alpha$ est affine, si
l'application $f|_{X_\alpha}$ satisfait (3), elle satisfera (1)
par hypothèse, donc sera propre. Pour achever la réduction au
cas où $S$ est affine, il faut montrer que, si $f: X \to S$ est séparé de
type fini et satisfait (3), alors $f|_{X_\alpha} : X_\alpha \to S_\alpha$
est séparé de type fini et satisfait (3). Le caractère séparé et le type
fini sont clairs. Pour (3), remarquons que
$\mathbf{A}^n \times X_\alpha$ est l'image réciproque ouverte de
$\mathbf{A}^n \times S_\alpha$ par l'application $1 \times f$. Fixons un
fermé $Z \subset \mathbf A^n \times X_\alpha$. Notons $\bar Z$
l'adhérence de $Z$ dans $\mathbf{A}^n \times X$. Pour des raisons
topologiques, on a
$$
1 \times f(\bar Z) \cap \mathbf{A}^n \times S_\alpha  = 1 \times f(Z).
$$
Ainsi $1 \times f(Z)$ est fermé, et la démonstration de
(3) $\Rightarrow$ (1) est ramenée au cas affine.

\medskip\noindent
Supposons $S$ affine et $f : X \to S$ séparé de type fini.
On peut appliquer le Lemme de Chow \ref{limits-lemma-chow-finite-type}
pour obtenir $\pi : X' \to X$ propre surjectif et $X' \to \mathbf{P}^n_S$
une immersion. Si $X$ est propre sur $S$, alors $X' \to S$ est propre
(Morphismes, Lemme \ref{morphisms-lemma-composition-proper}). Puisque
$\mathbf{P}^n_S \to S$ est séparé, on conclut que $X' \to
\mathbf{P}^n_S$ est propre
(Morphismes, Lemme \ref{morphisms-lemma-image-proper-scheme-closed}),
donc est une immersion fermée
(Schémas, Lemme \ref{schemes-lemma-immersion-when-closed}).
Réciproquement, supposons que $X' \to \mathbf{P}^n_S$ soit une immersion fermée.
Considérons le diagramme :
\begin{equation}
\label{limits-equation-check-proper}
\xymatrix{
X' \ar[r] \ar@{->>}[d]_{\pi} &
\mathbf{P}^n_S \ar[d] \\
X \ar[r]^f & S
}
\end{equation}
Tous les morphismes sont a priori propres, sauf $X \to S$.
On en conclut que $X \to S$ est propre par
Morphismes, Lemme \ref{morphisms-lemma-image-proper-is-proper}.
Nous avons donc montré que $X \to S$ est propre si et seulement si
$X' \to \mathbf{P}^n_S$ est une immersion fermée.

\medskip\noindent
Supposons $S$ affine et (3) satisfaite, et soient $n, X', \pi$ comme ci-dessus.
Puisque le caractère fermé d'un morphisme est local sur la base, l'application
$X \times \mathbf{P}^n \to S \times \mathbf{P}^n$ est fermée : d'après (3),
$X \times \mathbf{A}^n \to S \times \mathbf{A}^n$ est fermée, et
l'espace projectif est recouvert par des copies de l'espace affine de dimension $n$ ; voir
Constructions,
Lemme \ref{constructions-lemma-standard-covering-projective-space}.
D'après Morphismes, Lemme \ref{morphisms-lemma-base-change-proper},
le morphisme
$$
X' \times_S \mathbf{P}^n_S
\to
X \times_S \mathbf{P}^n_S =
X \times \mathbf{P}^n
$$
est propre. Puisque $\mathbf{P}^n$ est séparé, la projection
$$
X' \times_S \mathbf{P}^n_S = \mathbf{P}^n_{X'} \to X'
$$
est séparée, étant simplement un changement de base d'un morphisme
séparé. Par conséquent, l'application $X' \to X' \times_S \mathbf{P}^n_S$ est propre,
car elle est une section d'un morphisme séparé (voir
Schémas, Lemme \ref{schemes-lemma-section-immersion}).
En composant ces morphismes
$$
X' \to X' \times_S \mathbf{P}^n_S \to X \times_S \mathbf{P}^n_S
= X \times \mathbf{P}^n \to S \times \mathbf{P}^n = \mathbf{P}^n_S
$$
on trouve que l'immersion $X' \to \mathbf{P}^n_S$ est fermée, donc
est une immersion fermée.
\end{proof}





\section{Critère valuatif noethérien}
\label{limits-section-Noetherian-valuative-criterion}

\noindent
Si la base est noethérienne, on peut établir le critère valuatif
en ne considérant que des anneaux de valuation discrète.

\medskip\noindent
Beaucoup des résultats de cette section peuvent (et devraient peut-être)
se démontrer à l'aide du lemme suivant, bien que nous ne l'ayons
pas toujours fait.

\begin{lemma}
\label{limits-lemma-reach-point-closure-Noetherian}
Soit $f : X \to Y$ un morphisme de schémas.
Supposons $f$ de type fini et $Y$ localement noethérien.
Soit $y \in Y$ un point de l'adhérence de l'image de $f$.
Il existe alors un diagramme commutatif
$$
\xymatrix{
\Spec(K) \ar[r] \ar[d] & X \ar[d]^f \\
\Spec(A) \ar[r] & Y
}
$$
où $A$ est un anneau de valuation discrète et $K$ son corps des fractions,
envoyant le point fermé de $\Spec(A)$ sur $y$. De plus, on peut supposer
que le point image de $\Spec(K) \to X$ est un point générique $\eta$
d'une composante irréductible de $X$ et que $K = \kappa(\eta)$.
\end{lemma}

\begin{proof}
D'après la version non noethérienne de ce lemme
(Morphismes, Lemme \ref{morphisms-lemma-reach-points-scheme-theoretic-image}),
il existe un point $x \in X$ tel que $f(x)$ se spécialise en $y$.
On peut remplacer $x$ par tout point se spécialisant en $x$ ; on peut donc
supposer que $x$ est un point générique d'une composante irréductible de $X$.
On obtient ainsi un homomorphisme $\mathcal{O}_{Y, y} \to \kappa(x)$
(voir Schémas, section \ref{schemes-section-points}).
Soit $R \subset \kappa(x)$ son image. Alors $R$ est noethérien comme quotient
de l'anneau local noethérien $\mathcal{O}_{Y, y}$.
D'autre part, $\kappa(x)$ est
une extension de type fini du corps des fractions de $R$,
puisque $f$ est de type fini.
Il existe donc un anneau de valuation discrète $A \subset \kappa(x)$
de corps des fractions $\kappa(x)$ dominant $R$, d'après
Algèbre, Lemme \ref{algebra-lemma-exists-dvr}. Alors
$$
\xymatrix{
\Spec(\kappa(x)) \ar[d] \ar[rrr] & & & X \ar[d] \\
\Spec(A) \ar[r] & \Spec(R) \ar[r] & \Spec(\mathcal{O}_{Y, y}) \ar[r] & Y
}
$$
est le diagramme cherché.
\end{proof}

\noindent
Énonçons d'abord le résultat concernant
le caractère séparé. Nous utiliserons souvent des diagrammes commutatifs de morphismes
de schémas, dont les flèches pleines sont données, de la forme suivante :
\begin{equation}
\label{limits-equation-valuative}
\vcenter{
\xymatrix{
\Spec(K) \ar[r] \ar[d] & X \ar[d] \\
\Spec(A) \ar[r] \ar@{-->}[ru] & S
}
}
\end{equation}
où $A$ est un anneau de valuation et $K$ son corps des fractions.

\begin{lemma}
\label{limits-lemma-Noetherian-dvr-valuative-separation}
Soit $S$ un schéma localement noethérien.
Soit $f : X \to S$ un morphisme de schémas.
Supposons $f$ localement de type fini.
Les conditions suivantes sont équivalentes :
\begin{enumerate}
\item Le morphisme $f$ est séparé.
\item Pour tout diagramme (\ref{limits-equation-valuative}), il existe au plus
une flèche pointillée.
\item Pour tout diagramme (\ref{limits-equation-valuative}) où $A$ est un anneau
de valuation discrète, il existe au plus une flèche pointillée.
\item Pour toute composante irréductible $X_0$ de $X$ de
point générique $\eta \in X_0$, tout anneau de valuation discrète
$A \subset K = \kappa(\eta)$ de corps des fractions $K$ et tout
diagramme (\ref{limits-equation-valuative}) tel que
le morphisme $\Spec(K) \to X$ soit le morphisme canonique
(voir Schémas, section \ref{schemes-section-points}),
il existe au plus une flèche pointillée.
\end{enumerate}
\end{lemma}

\begin{proof}
Il est clair que (1) entraîne (2), que (2) entraîne (3) et que (3) entraîne (4).
Il reste à montrer que (4) entraîne (1). Supposons (4).
Ramenons-nous d'abord au cas où $S$ est affine. Le caractère séparé est local
sur la base (voir
Schémas, Lemme \ref{schemes-lemma-characterize-separated}).
Il suffit donc de montrer que, lorsque
$X \to S$ satisfait (4), la restriction $X_\alpha \to S_\alpha$ satisfait (4),
où $S_\alpha \subset S$ est un ouvert (affine) et $X_\alpha :=
f^{-1}(S_\alpha)$. Les
points génériques des composantes irréductibles de $X_\alpha$ sont des
points génériques de composantes irréductibles de $X$, puisque $X_\alpha$ est
ouvert dans $X$. Deux flèches pointillées distinctes dans le diagramme
\begin{equation}
\label{limits-equation-valuative-alpha}
\xymatrix{
\Spec(K) \ar[r] \ar[d] & X_\alpha \ar[d] \\
\Spec(A) \ar[r] \ar@{-->}[ru] & S_\alpha
}
\end{equation}
donneraient donc deux flèches distinctes dans le diagramme
(\ref{limits-equation-valuative}), par les morphismes $X_\alpha \to X$ et
$S_\alpha \to S$, ce qui est contradictoire. On est ainsi ramené
au cas où $S$ est affine. Remarquons qu'au cours de cette
réduction, nous montrons que, si $X \to S$ satisfait (4), la restriction $U
\to V$ satisfait (4) pour des ouverts $U \subset X$ et $V \subset S$ tels que
$f(U) \subset V$.

\medskip\noindent
Ramenons-nous ensuite au cas où $X \to S$ est de type fini. Supposons
que (4) entraîne (1) lorsque $X$ est de type fini. Puisque
$S$ est noethérien et $X$ localement de type fini sur $S$,
$X$ est lui aussi localement noethérien (voir Morphismes,
Lemme \ref{morphisms-lemma-finite-type-noetherian}).
Ainsi $X \to S$ est quasi-séparé (voir
Propriétés, Lemme \ref{properties-lemma-locally-Noetherian-quasi-separated}),
et l'on peut donc appliquer le critère valuatif pour déterminer si $X$
est séparé (voir
Schémas, Lemme \ref{schemes-lemma-valuative-criterion-separatedness}).
Soit $X = \bigcup_\alpha X_\alpha$ un recouvrement ouvert
affine de $X$. Étant données deux flèches pointillées dans un diagramme
(\ref{limits-equation-valuative}), les images du point fermé de
$\Spec A$
appartiennent à deux ouverts $X_\alpha$ et $X_\beta$. Puisque $X_\alpha \cup
X_\beta$ est ouvert, il contient, pour des raisons topologiques, l'image de
$\Spec(A)$ par chacune des deux applications. Les deux flèches pointillées se factorisent donc
par $X_\alpha \cup X_\beta \to X$, dont la source est un schéma de type fini sur
$S$. Puisque $X_\alpha \cup X_\beta$ est un ouvert de $X$, notre
remarque précédente montre que $X_\alpha \cup X_\beta$ satisfait (4), donc
est séparé par hypothèse. Les deux flèches pointillées
sont donc égales. On est ainsi ramené au cas où $X \to S$ est de type fini.

\medskip\noindent
Supposons $X \to S$ de type fini et (4) satisfaite.
Puisque $X \to S$ est de type fini et $S$ un schéma affine
noethérien, $X$ est lui aussi noethérien (voir
Morphismes, Lemme \ref{morphisms-lemma-finite-type-noetherian}).
Ainsi $X \to X \times_S X$ est
une immersion quasi-compacte de schémas noethériens. Raisonnons par
l'absurde. Supposons que $X \to X \times_S X$ ne soit pas fermé.
Il existe alors un $y \in X \times_S X$ dans l'adhérence de l'image, mais
hors de l'image. Comme $X$ est noethérien, il possède un nombre fini de composantes
irréductibles. Le point $y$ appartient donc à l'adhérence de l'image d'une
composante irréductible $X_0 \subset X$. Munissons $X_0$ de la structure
réduite induite. La composée $X_0 \to X \to X \times_S X$
se factorise par le sous-schéma fermé $X_0 \times_S X_0 \subset X \times_S X$.
Notons l'adhérence de $\Delta(X_0)$ dans $X_0 \times_S X_0$
par $\bar X_0$ (encore munie de la structure de sous-schéma fermé réduit). Ainsi $y \in \bar X_0$.
Puisque $X_0 \to X_0 \times_S X_0$ est une immersion, l'image de $X_0$
est ouverte dans $\bar X_0$. Ainsi $X_0$ et $\bar X_0$ sont
birationnels. Puisque $\bar{X}_0$ est un sous-schéma fermé d'un
schéma noethérien, il est noethérien. L'anneau local
$\mathcal O_{{\bar X_0, y}}$ est donc un anneau local noethérien intègre de corps des
fractions $K$ égal au corps des fonctions de $X_0$. D'après le théorème de Krull--Akizuki
(voir Algèbre, Lemme \ref{algebra-lemma-exists-dvr}), il existe un
anneau de valuation discrète $A$ dominant $\mathcal O_{{\bar X_0, y}}$
et de corps des fractions $K$. Cela permet de construire un diagramme :
\begin{equation}
\label{limits-equation-valuative-generic}
\xymatrix{
\Spec(K) \ar[r] \ar[d] & X_0 \ar[d]^{\Delta} \\
\Spec(A) \ar[r] \ar@{-->}[ur]& X_0 \times_S X_0 \\
}
\end{equation}
qui envoie $\Spec K$ sur le point générique de $\Delta(X_0)$ et
le point fermé de $A$ sur $y \in X_0 \times_S X_0$ (utiliser
Schémas, section \ref{schemes-section-points}, pour construire les flèches).
Il ne peut même pas exister de
flèche pointillée ensembliste, puisque $y$ n'est pas dans l'image de
$\Delta$, par le choix de $y$. Par un argument catégorique, l'existence
de la flèche pointillée dans le diagramme ci-dessus équivaut à l'unicité
de la flèche pointillée dans le diagramme suivant :
\begin{equation}
\label{limits-equation-valuative-nonexistent}
\xymatrix{
\Spec(K) \ar[r] \ar[d] & X_0 \ar[d]\\
\Spec(A) \ar[r] \ar@{-->}[ur] & S \\
}
\end{equation}
La non-existence dans le premier diagramme donne donc la non-unicité dans
ce dernier. Ainsi $X_0$ ne satisfait pas
l'unicité pour les anneaux de valuation discrète ; puisque $X_0$ est une
composante irréductible de $X$, le morphisme $X \to S$ ne satisfait pas
(4). Nous avons donc montré que (4) entraîne (1).
\end{proof}

\begin{lemma}
\label{limits-lemma-Noetherian-dvr-valuative-proper}
Soit $S$ un schéma localement noethérien.
Soit $f : X \to S$ un morphisme de type fini.
Les conditions suivantes sont équivalentes :
\begin{enumerate}
\item Le morphisme $f$ est propre.
\item Pour tout diagramme (\ref{limits-equation-valuative}), il existe exactement
une flèche pointillée.
\item Pour tout diagramme (\ref{limits-equation-valuative}) où $A$ est un anneau
de valuation discrète, il existe exactement une flèche pointillée.
\item Pour toute composante irréductible $X_0$ de $X$ de
point générique $\eta \in X_0$, tout anneau de valuation discrète
$A \subset K = \kappa(\eta)$ de corps des fractions $K$ et tout
diagramme (\ref{limits-equation-valuative}) tel que
le morphisme $\Spec(K) \to X$ soit le morphisme canonique
(voir
Schémas, section \ref{schemes-section-points}),
il existe exactement une flèche pointillée.
\end{enumerate}
\end{lemma}

\begin{proof}
(1) entraîne (2), qui entraîne (3), qui entraîne (4). Montrons maintenant que (4) entraîne
(1). Comme dans la démonstration du Lemme \ref{limits-lemma-Noetherian-dvr-valuative-separation},
on se ramène au
cas où $S$ est affine : la propreté est locale sur la base, et si $X
\to S$ satisfait (4), alors $X_\alpha \to S_\alpha$ la satisfait aussi pour
un ouvert $S_\alpha \subset S$ et $X_\alpha = f^{-1}(S_\alpha)$.

\medskip\noindent
Le schéma $S$ est maintenant noethérien ; $X$ l'est donc aussi, puisque $X \to
S$ est de type fini. On peut alors utiliser le lemme de Chow
(Cohomologie des schémas, Lemme \ref{coherent-lemma-chow-Noetherian})
pour obtenir un morphisme surjectif, propre et birationnel
$X' \to X$ et une immersion $X' \to \mathbf{P}^n_S$. Nous voulons
montrer que $X \to S$ est universellement fermé. Comme dans la démonstration du Lemme
\ref{limits-lemma-limited-base-change}, il suffit de vérifier que
$X' \to \mathbf{P}^n_S$ est une immersion fermée.
Raisonnons par l'absurde et supposons que $X' \to
\mathbf{P}^n_S$ ne soit pas une immersion fermée. Il existe alors un $y
\in \mathbf{P}^n_S$ dans l'adhérence de l'image de $X'$,
mais hors de l'image. Ainsi $y$ est dans l'adhérence de l'image d'une
composante irréductible $X_0'$ de $X'$, mais hors de l'image.
Soit $\bar X_0' \subset \mathbf{P}^n_S$ l'adhérence de
l'image de $X_0'$. Puisque $X' \to \mathbf{P}^n_S$ est une immersion
de schémas noethériens, le morphisme $X'_0 \to \bar X_0'$ est
ouvert et dense. D'après
Algèbre, Lemme \ref{algebra-lemma-exists-dvr},
ou
Propriétés, Lemme \ref{properties-lemma-locally-Noetherian-specialization-dvr},
on peut trouver un anneau de valuation discrète $A$ dominant
$\mathcal{O}_{\bar X_0', y}$ et de même corps
des fractions $K$. Il est clair que
$K$ est le corps résiduel du point générique de $X_0'$.
On obtient donc le diagramme commutatif de flèches pleines
\begin{equation}
\label{limits-equation-solid}
\xymatrix{
\Spec K \ar[r] \ar[d] & X' \ar [r] \ar[d] &
\mathbf{P}^n_S \ar[d] \\
\Spec A \ar@{-->}[r] \ar@{-->}[ru] \ar[urr] & X \ar[r] & S\\
}
\end{equation}
Remarquons que le point fermé de $A$ a pour image $y \in \mathbf{P}^n_S$. Par
construction, il n'existe pas de relèvement ensembliste à $X'$.
Puisque $X' \to X$ est birationnel, l'image de $X'_0$ dans $X$ est une
composante irréductible $X_0$ de $X$, et $K$ s'identifie aussi
au corps des fonctions de $X_0$. Puisque $X \to S$ satisfait (4) par hypothèse,
la flèche pointillée $\Spec(A) \to X$ existe.
Puisque $X' \to X$ est propre, cette flèche pointillée
se relève en une flèche pointillée $\Spec(A) \to X'$ (utiliser Schémas,
Proposition \ref{schemes-proposition-characterize-universally-closed}).
En la composant avec l'immersion $X' \to \mathbf{P}^n_S$, on obtient
un autre morphisme (non représenté dans le diagramme)
$\Spec(A) \to \mathbf{P}^n_S$. Puisque $\mathbf{P}^n_S$
est propre sur $S$, il satisfait (2), et ces deux morphismes
coïncident. C'est une contradiction, car nous avons construit
le relèvement interdit de l'application initiale $\Spec(A) \to \mathbf{P}^n_S$
à $X'$.
\end{proof}

\begin{lemma}
\label{limits-lemma-check-universally-closed-Noetherian}
Soit $f : X \to S$ un morphisme de type fini de schémas.
Supposons $S$ localement noethérien. Les conditions suivantes sont alors équivalentes :
\begin{enumerate}
\item $f$ est universellement fermé,
\item pour tout $n$, le morphisme
$\mathbf{A}^n \times X \to \mathbf{A}^n \times S$ est fermé,
\item pour tout diagramme (\ref{limits-equation-valuative}), il existe une
flèche pointillée,
\item pour tout diagramme (\ref{limits-equation-valuative}) où $A$ est un anneau
de valuation discrète, il existe une flèche pointillée.
\end{enumerate}
\end{lemma}

\begin{proof}
L'équivalence de (1) et (2) est un cas particulier du
Lemme \ref{limits-lemma-test-universally-closed}.
L'équivalence de (1) et (3) est un cas particulier de
Schémas, Proposition \ref{schemes-proposition-characterize-universally-closed}.
Il est évident que (3) entraîne (4).
Il suffit donc de montrer que (4) entraîne (2).
Nous montrerons que $\mathbf{A}^n \times X \to \mathbf{A}^n \times S$
est fermé à l'aide du critère de
Schémas, Lemme \ref{schemes-lemma-quasi-compact-closed}.
Choisissons $n$, une spécialisation non triviale $z \leadsto z'$ de points
de $\mathbf{A}^n \times S$ et un point $y \in \mathbf{A}^n \times X$
au-dessus de $z$. Remarquons que $\kappa(y)$ est une extension de corps de type fini
de $\kappa(z)$, puisque $\mathbf{A}^n \times X \to \mathbf{A}^n \times S$
est de type fini. D'après
Propriétés, Lemme \ref{properties-lemma-locally-Noetherian-specialization-dvr},
ou
Algèbre, Lemme \ref{algebra-lemma-exists-dvr},
il existe donc un anneau de valuation discrète $A \subset \kappa(y)$
de corps des fractions $\kappa(z)$ dominant l'image de
$\mathcal{O}_{\mathbf{A}^n \times S, z'}$ dans $\kappa(z)$.
Cela donne un diagramme commutatif
$$
\xymatrix{
\Spec(\kappa(y)) \ar[r] \ar[d] &
\mathbf{A}^n \times X \ar[d] \ar[r] & X \ar[d] \\
\Spec(A) \ar[r] & \mathbf{A}^n \times S \ar[r] & S
}
$$
La propriété (4) entraîne alors l'existence d'un morphisme
$\Spec(A) \to X$ s'insérant dans ce diagramme.
Puisque la flèche horizontale inférieure gauche fournit déjà
le morphisme $\Spec(A) \to \mathbf{A}^n$,
on obtient aussi un morphisme
$\Spec(A) \to \mathbf{A}^n \times X$ s'insérant dans le
carré de gauche. L'image $y' \in \mathbf{A}^n \times X$
du point fermé est donc une spécialisation de $y$ au-dessus de $z'$.
Cela montre que les spécialisations se relèvent le long de
$\mathbf{A}^n \times X \to \mathbf{A}^n \times S$,
ce qui conclut.
\end{proof}









\section{Critères valuatifs noethériens affinés}
\label{limits-section-refined-valuative-criteria}

\noindent
Il n'est généralement pas nécessaire de considérer tous les diagrammes possibles
avec des anneaux de valuation pour vérifier les critères valuatifs. Un exemple
est donné par Morphismes, Lemme
\ref{morphisms-lemma-refined-valuative-criterion-universally-closed}.
Dans le cadre noethérien, nous l'avons également vu dans les
Lemmes \ref{limits-lemma-Noetherian-dvr-valuative-separation} et
\ref{limits-lemma-Noetherian-dvr-valuative-proper}. Voici une autre variante.

\begin{lemma}
\label{limits-lemma-refined-valuative-criterion-proper}
Soient $f : X \to S$ et $h : U \to X$ des morphismes de schémas.
Supposons $S$ localement noethérien, $f$ et $h$ de type fini,
$f$ séparé et $h(U)$ dense dans $X$.
Si, pour tout diagramme commutatif de flèches pleines
$$
\xymatrix{
\Spec(K) \ar[r] \ar[d] & U \ar[r]^h & X \ar[d]^f \\
\Spec(A) \ar[rr] \ar@{-->}[rru] & & S
}
$$
où $A$ est un anneau de valuation discrète de corps des fractions $K$,
il existe une flèche pointillée rendant le diagramme commutatif, alors $f$ est propre.
\end{lemma}

\begin{proof}
On se ramène immédiatement au cas où $S$ est affine.
Alors $U$ est quasi-compact.
Soit $U = U_1 \cup \ldots \cup U_n$ un recouvrement ouvert affine.
On peut remplacer $U$ par $U_1 \amalg \ldots \amalg U_n$ sans
changer les hypothèses ; on peut donc supposer $U$ affine.
On peut alors trouver une immersion ouverte $U \to Y$ sur $X$,
où $Y$ est propre sur $X$. (Plonger d'abord $U$ dans $\mathbf{A}^n_X$
à l'aide de Morphismes, Lemme \ref{morphisms-lemma-quasi-affine-finite-type-over-S},
puis prendre l'adhérence dans $\mathbf{P}^n_X$, ou utiliser directement
Morphismes, Lemme \ref{morphisms-lemma-quasi-projective-open-projective}.)
On peut supposer $U$ dense dans $Y$ (remplacer $Y$ par l'adhérence schématique
de $U$ si nécessaire ; voir Morphismes, section
\ref{morphisms-section-scheme-theoretic-closure}).
Remarquons que $g : Y \to X$ est surjectif, car son image est fermée
et contient la partie dense $h(U)$.
Nous allons montrer que $Y \to S$ est propre. Cela entraînera que
$X \to S$ est propre par
Morphismes, Lemme \ref{morphisms-lemma-image-proper-is-proper},
et achèvera la démonstration.
Pour montrer que $Y \to S$ est propre, nous utiliserons
le point (4) du Lemme \ref{limits-lemma-Noetherian-dvr-valuative-proper}.
Considérons à cet effet un diagramme
$$
\xymatrix{
\Spec(K) \ar[r]_y \ar[d] & Y \ar[d]^{f \circ g} \\
\Spec(A) \ar[r] \ar@{..>}[ru] & S
}
$$
où $A$ est un anneau de valuation discrète de corps des fractions $K$
et où $y : \Spec(K) \to Y$ est l'inclusion d'un point générique.
Il faut montrer qu'il existe une unique flèche pointillée.
L'unicité résulte de la réciproque du critère valuatif
de séparation
(Schémas, Lemme \ref{schemes-lemma-separated-implies-valuative}),
puisque $Y \to S$ est séparé comme
composée des morphismes séparés $Y \to X$ et $X \to S$
(Schémas, Lemme \ref{schemes-lemma-separated-permanence}).
L'existence se voit comme suit.
Puisque $y$ est un point générique de $Y$, il appartient à $U$.
L'hypothèse du lemme
fournit un morphisme $a : \Spec(A) \to X$ tel que
$$
\xymatrix{
\Spec(K) \ar[r]_y \ar[d] & U \ar[r] & X \ar[d]^f \\
\Spec(A) \ar[rr] \ar[rru]^a & & S
}
$$
soit commutatif. Puisque $Y \to X$ est propre, on peut ensuite
appliquer le critère valuatif de propreté
(Morphismes, Lemme \ref{morphisms-lemma-characterize-proper})
pour trouver un morphisme $b : \Spec(A) \to Y$ tel que
$$
\xymatrix{
\Spec(K) \ar[r]_y \ar[d] & Y \ar[d]^g \\
\Spec(A) \ar[r]^a \ar[ru]^b & X
}
$$
soit commutatif. Cela achève la démonstration, puisque
$b$ peut servir de flèche pointillée ci-dessus.
\end{proof}

\begin{lemma}
\label{limits-lemma-refined-valuative-criterion-separated}
Soient $f : X \to S$ et $h : U \to X$ des morphismes de schémas.
Supposons $S$ localement noethérien, $f$ localement de type fini,
$h$ de type fini et $h(U)$ dense dans $X$.
Si, pour tout diagramme commutatif de flèches pleines
$$
\xymatrix{
\Spec(K) \ar[r] \ar[d] & U \ar[r]^h & X \ar[d]^f \\
\Spec(A) \ar[rr] \ar@{-->}[rru] & & S
}
$$
où $A$ est un anneau de valuation discrète de corps des fractions $K$,
il existe au plus une flèche pointillée rendant le diagramme commutatif, alors $f$ est
séparé.
\end{lemma}

\begin{proof}
Nous allons appliquer le Lemme \ref{limits-lemma-refined-valuative-criterion-proper}
aux morphismes $U \to X$ et $\Delta : X \to X \times_S X$.
Vérifions les conditions. Remarquons que $\Delta$ est quasi-compact d'après
Propriétés, Lemme \ref{properties-lemma-locally-Noetherian-quasi-separated}
(et Schémas, Lemme \ref{schemes-lemma-compose-after-separated}).
Bien entendu, $\Delta$ est localement de type fini et séparé (ce qui est vrai
de tout morphisme diagonal).
Enfin, donnons-nous un diagramme commutatif de flèches pleines
$$
\xymatrix{
\Spec(K) \ar[r] \ar[d] & U \ar[r]^h & X \ar[d]^\Delta \\
\Spec(A) \ar[rr]^{(a, b)} \ar@{-->}[rru] & & X \times_S X
}
$$
où $A$ est un anneau de valuation discrète de corps des fractions $K$.
Alors $a$ et $b$ donnent deux flèches pointillées dans le diagramme du lemme
et doivent être égales. On peut donc prendre $a = b$ comme flèche pointillée,
ce qui donne l'existence et achève la démonstration.
\end{proof}

\begin{lemma}
\label{limits-lemma-refined-valuative-criterion-universally-closed}
Soient $f : X \to S$ et $h : U \to X$ des morphismes de schémas.
Supposons $S$ localement noethérien, $f$ et $h$ de type fini et
$h(U)$ dense dans $X$. Si, pour tout diagramme commutatif de flèches pleines
$$
\xymatrix{
\Spec(K) \ar[r] \ar[d] & U \ar[r]^h & X \ar[d]^f \\
\Spec(A) \ar[rr] \ar@{-->}[rru] & & S
}
$$
où $A$ est un anneau de valuation discrète de corps des fractions $K$,
il existe une unique flèche pointillée rendant le diagramme commutatif, alors $f$ est propre.
\end{lemma}

\begin{proof}
Combiner les Lemmes \ref{limits-lemma-refined-valuative-criterion-separated} et
\ref{limits-lemma-refined-valuative-criterion-proper}.
\end{proof}









\section{Critères valuatifs sur une base de Nagata}
\label{limits-section-nagata-valuative}

\noindent
Pour les schémas localement de type fini sur une
base de Nagata, on peut se limiter aux anneaux de valuation discrète
essentiellement de type fini sur la base.
Les résultats suivants ne sont que quelques exemples de ce que l'on peut obtenir.

\begin{lemma}
\label{limits-lemma-essentially-finite-type-criterion-universally-closed}
Soit $S$ un schéma de Nagata (donc, en particulier, localement noethérien).
Soit $f : X \to Y$ un morphisme quasi-compact de schémas localement
de type fini sur $S$. Les conditions suivantes sont équivalentes :
\begin{enumerate}
\item $f$ est universellement fermé,
\item pour tout $n$, le morphisme
$\mathbf{A}^n \times X \to \mathbf{A}^n \times Y$ est fermé,
\item pour tout diagramme commutatif
$$
\xymatrix{
U \ar[r] \ar[d] & X \ar[d]^f \\
C \ar[r] \ar@{..>}[ru] & Y
}
$$
de schémas sur $S$ tel que
\begin{enumerate}
\item $C$ soit un schéma intègre normal de type fini sur $S$,
\item $U = C \setminus \{c\}$ pour un point fermé $c \in C$,
\item $A = \mathcal{O}_{C, c}$ soit de dimension $1$\footnote{Il s'ensuit
que $A$ est un anneau de valuation discrète ; voir
Algèbre, Lemme \ref{algebra-lemma-characterize-dvr}.
De plus, $c$ a pour image un point de type fini $s \in S$, et
$A$ est essentiellement de type fini sur $\mathcal{O}_{S, s}$.}
\end{enumerate}
alors, dans le diagramme commutatif
$$
\xymatrix{
\Spec(K) \ar[r] \ar[d] & X \ar[d]^f \\
\Spec(A) \ar[r] \ar@{-->}[ru] & Y
}
$$
où $K = \text{Frac}(A)$, il existe une flèche pointillée\footnote{D'après le
Lemme \ref{limits-lemma-morphism-glueing-near-closed-point}, cela
équivaut à demander l'existence d'une flèche pointillée
rendant commutatif le premier diagramme.}
rendant le diagramme commutatif.
\end{enumerate}
\end{lemma}

\begin{proof}
Nous avons vu l'équivalence de (1) et (2), ainsi que le
fait que ces conditions entraînent (3), dans le
Lemme \ref{limits-lemma-check-universally-closed-Noetherian}.
Il suffit donc de montrer que (3) entraîne (2).
Remarquons que, si la condition (3) est satisfaite par $f : X \to Y$,
elle l'est aussi par
$1 \times f : \mathbf{A}^n \times X \to \mathbf{A}^n \times Y$
(voir l'argument de la démonstration du
Lemme \ref{limits-lemma-check-universally-closed-Noetherian}).
Il suffit donc de montrer que (3) entraîne que $f$ est fermé.

\medskip\noindent
Réduction au cas où $Y$ et $S$ sont affines ; on peut
sauter ce paragraphe.
Soit $S' \subset S$ un ouvert affine et $Y' \subset Y$ un
ouvert affine dont l'image est contenue dans $S'$. Posons $X' = f^{-1}(Y')$. Nous
affirmons que la restriction $f' : X' \to Y'$ de $f$, considérée comme morphisme
de schémas sur $S'$, possède aussi la propriété (3). Nous omettons les détails.
Si l'on montre que $f'$ est fermé pour tous les choix de
$S'$ et $Y'$, il s'ensuit que $f$ est fermé. Cela nous ramène
au cas traité au paragraphe suivant.

\medskip\noindent
Supposons $S$ et $Y$ affines. Soit $Z \subset X$ une partie fermée.
Munissons $Z$ de la structure de sous-schéma fermé réduit de $X$.
Il faut montrer que $E = f(Z)$ est fermé. Choisissons un point fermé $y \in Y$
dans l'adhérence de $f(Z)$. Il suffit de montrer $y \in E$.
Supposons $y \not \in E$ pour obtenir une contradiction.
L'image $s \in S$ de $y$ est un point de type fini de $S$ ; voir
Morphismes, Lemme \ref{morphisms-lemma-finite-type-points-morphism}.
Rappelons que $E$ est constructible
(Morphismes, Lemme \ref{morphisms-lemma-chevalley}).
Considérons l'intersection $\Spec(\mathcal{O}_{Y, y}) \cap E$.
C'est une partie constructible du spectre
(Morphismes, Lemme \ref{morphisms-lemma-inverse-image-constructible})
qui ne contient pas le point fermé.
Puisque le spectre épointé
$\Spec(\mathcal{O}_{Y, y}) \setminus \{y\}$ est de Jacobson
(Morphismes, Lemme \ref{morphisms-lemma-ubiquity-Jacobson-schemes}), on trouve un
point fermé $t \in \Spec(\mathcal{O}_{Y, y}) \setminus \{y\}$
tel que $t \in E$ (voir Topologie, Lemme \ref{topology-lemma-jacobson-inherited}).
Autrement dit, $t \in E$ est un point de $Y$
admettant une spécialisation immédiate $t \leadsto y$.
Puisque $t \in E$, la fibre schématique $Z_t$ est non vide.
Choisissons un point fermé $x \in Z_t$. On a en particulier
$[\kappa(x) : \kappa(t)] < \infty$ par le Nullstellensatz de Hilbert
(Morphismes, Lemme
\ref{morphisms-lemma-closed-point-fibre-locally-finite-type}).

\medskip\noindent
Notons $T = \overline{\{t\}} \subset Y$ le sous-schéma fermé intègre
dont l'espace topologique sous-jacent est celui indiqué
(Schémas, Définition \ref{schemes-definition-reduced-induced-scheme}).
Alors $t \in T$ est le point générique.
Notons $C \to T$ la normalisation de $T$ dans $\kappa(x)$ ; voir
Morphismes, section \ref{morphisms-section-normalization-X-in-Y}
(plus précisément, $C \to T$ est la normalisation de $T$ dans $x$,
où l'on considère $x = \Spec(\kappa(x)) \to T$ comme schéma sur $T$).
Puisque $S$ est un schéma de Nagata, $T$ l'est aussi
(Morphismes, Lemme \ref{morphisms-lemma-finite-type-nagata}).
Ainsi $C \to T$ est fini
(Morphismes, Lemme \ref{morphisms-lemma-nagata-normalization-finite-general}).
Puisque $t$ appartient à l'image, $C \to T$ est surjectif (car
l'image est fermée et $T$ est l'adhérence de $t$ dans $Y$).
Choisissons un point $c \in C$ d'image $y \in T$.
Puisque $y$ est un point fermé de $T$, $c$ est
un point fermé de $C$. Puisque $\dim(\mathcal{O}_{T, y}) = 1$,
on a $\dim(\mathcal{O}_{C, c}) = 1$ (la dimension est
au moins $1$, car $c$ n'est pas le point générique de $C$, et
au plus $1$, car $C \to T$ est fini).
Puisque le corps des fonctions de $C$ est $\kappa(x)$ et que $x$
est un point de $X$, on a une application rationnelle sur $Y$ de $C$ vers $X$
(voir par exemple Morphismes, Lemme
\ref{morphisms-lemma-rational-map-finite-presentation}).
Soit $C \supset U \to X$ un représentant (en particulier, $U$ est non vide).
On peut supposer $c \not \in U$ (remplacer $U$ par $U \setminus \{c\}$).
Puisque $c$ est un point fermé de codimension $1$ du schéma intègre $C$,
on a $C = U \amalg \{c\} \amalg \Sigma$
pour une partie fermée propre $\Sigma \subset C$.
En remplaçant $C$ par $C \setminus \Sigma$,
on obtient un diagramme commutatif comme dans (3).
D'après la seconde note de bas de page de l'énoncé,
l'existence de la flèche pointillée prolonge
l'application rationnelle à tout $C$ ; on obtient une contradiction,
car l'image de $c$ est alors un point de $Z$ d'image $y$.
\end{proof}

\begin{lemma}
\label{limits-lemma-essentially-finite-type-criterion-separation}
Soit $S$ un schéma de Nagata (donc, en particulier, localement noethérien).
Soit $f : X \to Y$ un morphisme de schémas localement de type fini sur $S$.
Les conditions suivantes sont équivalentes :
\begin{enumerate}
\item $f$ est séparé,
\item pour tout diagramme commutatif
$$
\xymatrix{
U \ar[r] \ar[d] & X \ar[d]^f \\
C \ar[r] \ar@{..>}[ru] & Y
}
$$
de schémas sur $S$ tel que
\begin{enumerate}
\item $C$ soit un schéma intègre normal de type fini sur $S$,
\item $U = C \setminus \{c\}$ pour un point fermé $c \in C$,
\item $A = \mathcal{O}_{C, c}$ soit de dimension $1$\footnote{Il s'ensuit
que $A$ est un anneau de valuation discrète ; voir
Algèbre, Lemme \ref{algebra-lemma-characterize-dvr}.
De plus, $c$ a pour image un point de type fini $s \in S$, et
$A$ est essentiellement de type fini sur $\mathcal{O}_{S, s}$.}
\end{enumerate}
alors, dans le diagramme commutatif
$$
\xymatrix{
\Spec(K) \ar[r] \ar[d] & X \ar[d]^f \\
\Spec(A) \ar[r] \ar@{-->}[ru] & Y
}
$$
où $K = \text{Frac}(A)$, il existe au plus une flèche pointillée\footnote{D'après le
Lemme \ref{limits-lemma-morphism-glueing-near-closed-point}, cela
équivaut à demander qu'il existe au plus une flèche pointillée
rendant commutatif le premier diagramme.}
rendant le diagramme commutatif.
\end{enumerate}
\end{lemma}

\begin{proof}
D'après le Lemme \ref{limits-lemma-Noetherian-dvr-valuative-separation},
(1) entraîne (2). Supposons (2). Pour montrer que $f$ est séparé,
il faut montrer que $\Delta : X \to X \times_Y X$ est fermé. D'après
Morphismes, Lemme \ref{morphisms-lemma-finite-type-Noetherian-quasi-separated},
le morphisme $\Delta$ est quasi-compact.
D'après le Lemme \ref{limits-lemma-essentially-finite-type-criterion-universally-closed},
il suffit de montrer que, pour tout diagramme commutatif
$$
\xymatrix{
U \ar[rr] \ar[d] & & X \ar[d]^\Delta \\
C \ar[rr]^{(a_1, a_2)} \ar@{..>}[rru] & & X \times_Y X
}
$$
de schémas sur $S$ tel que
\begin{enumerate}
\item $C$ soit un schéma intègre normal de type fini sur $S$,
\item $U = C \setminus \{c\}$ pour un point fermé $c \in C$,
\item $A = \mathcal{O}_{C, c}$ soit de dimension $1$,
\end{enumerate}
alors, dans le diagramme commutatif
$$
\xymatrix{
\Spec(K) \ar[r] \ar[d] & X \ar[d]^\Delta \\
\Spec(A) \ar[r] \ar@{-->}[ru] & X \times_Y X
}
$$
où $K = \text{Frac}(A)$, il existe une flèche pointillée
rendant le diagramme commutatif. D'après le
Lemme \ref{limits-lemma-morphism-glueing-near-closed-point},
l'existence de la flèche pointillée dans le second diagramme
équivaut à celle de la flèche pointillée
dans le premier. De plus, cette existence
équivaut à demander $a_1 = a_2$. Or $a_1|_U = a_2|_U$ ;
l'hypothèse d'unicité (2) donne donc cette égalité,
ce qui achève la démonstration.
\end{proof}

\begin{lemma}
\label{limits-lemma-essentially-finite-type-criterion-proper}
Soit $S$ un schéma de Nagata (donc, en particulier, localement noethérien).
Soit $f : X \to Y$ un morphisme quasi-compact de schémas localement
de type fini sur $S$. Les conditions suivantes sont équivalentes :
\begin{enumerate}
\item $f$ est propre,
\item pour tout diagramme commutatif
$$
\xymatrix{
U \ar[r] \ar[d] & X \ar[d]^f \\
C \ar[r] \ar@{..>}[ru] & Y
}
$$
de schémas sur $S$ tel que
\begin{enumerate}
\item $C$ soit un schéma intègre normal de type fini sur $S$,
\item $U = C \setminus \{c\}$ pour un point fermé $c \in C$,
\item $A = \mathcal{O}_{C, c}$ soit de dimension $1$\footnote{Il s'ensuit
que $A$ est un anneau de valuation discrète ; voir
Algèbre, Lemme \ref{algebra-lemma-characterize-dvr}.
De plus, $c$ a pour image un point de type fini $s \in S$, et
$A$ est essentiellement de type fini sur $\mathcal{O}_{S, s}$.}
\end{enumerate}
alors, dans le diagramme commutatif
$$
\xymatrix{
\Spec(K) \ar[r] \ar[d] & X \ar[d]^f \\
\Spec(A) \ar[r] \ar@{-->}[ru] & Y
}
$$
où $K = \text{Frac}(A)$, il existe exactement une flèche pointillée\footnote{D'après le
Lemme \ref{limits-lemma-morphism-glueing-near-closed-point}, cela
équivaut à demander l'existence et l'unicité de la flèche pointillée
rendant commutatif le premier diagramme.}
rendant le diagramme commutatif.
\end{enumerate}
\end{lemma}

\begin{proof}
Cela découle formellement des Lemmes
\ref{limits-lemma-essentially-finite-type-criterion-universally-closed} et
\ref{limits-lemma-essentially-finite-type-criterion-separation}
et de la définition d'un morphisme propre : un morphisme de type fini, séparé
et universellement fermé.
\end{proof}






\section{Limites et dimensions des fibres}
\label{limits-section-limits-dimension}

\noindent
Le lemme suivant est le plus souvent utilisé dans la situation du
Lemme \ref{limits-lemma-descend-finite-presentation}
pour garantir que, si les fibres de la limite sont de dimension $\leq d$,
les fibres à un certain étage fini sont aussi de dimension $\leq d$.

\begin{lemma}
\label{limits-lemma-limit-dimension}
Soit $I$ un ensemble ordonné filtrant.
Soit $(f_i : X_i \to S_i)$ un système projectif de morphismes de schémas
indexé par $I$. Faisons les hypothèses suivantes :
\begin{enumerate}
\item tous les morphismes $S_{i'} \to S_i$ sont affines,
\item tous les schémas $S_i$ sont quasi-compacts et quasi-séparés,
\item les morphismes $f_i$ sont de type fini,
\item les morphismes $X_{i'} \to X_i \times_{S_i} S_{i'}$ sont des immersions
fermées.
\end{enumerate}
Soit $f : X = \lim_i X_i \to S = \lim_i S_i$ la limite.
Soit $d \geq 0$.
Si toute fibre de $f$ est de dimension $\leq d$, alors, pour un certain $i$,
toute fibre de $f_i$ est de dimension $\leq d$.
\end{lemma}

\begin{proof}
Pour chaque $i$, soit $U_i = \{x \in X_i \mid \dim_x((X_i)_{f_i(x)}) \leq d\}$.
C'est un ouvert de $X_i$ ; voir
Morphismes, Lemme \ref{morphisms-lemma-openness-bounded-dimension-fibres}.
Posons $Z_i = X_i \setminus U_i$ (muni de la structure réduite induite).
Il faut montrer que $Z_i = \emptyset$ pour un certain $i$.
Sinon, $Z = \lim Z_i \not = \emptyset$ ; voir
Lemme \ref{limits-lemma-limit-nonempty}.
Soit $z \in Z$ un point. Remarquons que $Z \subset X$ est un sous-schéma fermé.
Posons $s = f(z)$. Pour chaque $i$, soit $s_i \in S_i$ l'image
de $s$. Remarquons que $Z_s$ est la limite des schémas $(Z_i)_{s_i}$,
et que $Z_s$ est aussi la limite des schémas $(Z_i)_{s_i}$ après changement
de base à $\kappa(s)$. De plus, tous les morphismes
$$
Z_s
\longrightarrow
(Z_{i'})_{s_{i'}} \times_{\Spec(\kappa(s_{i'}))} \Spec(\kappa(s))
\longrightarrow
(Z_i)_{s_i} \times_{\Spec(\kappa(s_i))} \Spec(\kappa(s))
\longrightarrow
X_s
$$
sont des immersions fermées par l'hypothèse (4). Ainsi $Z_s$ est l'intersection
schématique des sous-schémas fermés
$(Z_i)_{s_i} \times_{\Spec(\kappa(s_i))} \Spec(\kappa(s))$
dans $X_s$. Puisque toutes les composantes irréductibles des schémas
$(Z_i)_{s_i} \times_{\Spec(\kappa(s_i))} \Spec(\kappa(s))$
sont de dimension $> d$ et contiennent $z$, on conclut que
$Z_s$ contient une composante irréductible de dimension $> d$ passant
par $z$, ce qui contredit $Z_s \subset X_s$ et
$\dim(X_s) \leq d$.
\end{proof}

\begin{lemma}
\label{limits-lemma-descend-quasi-finite}
Reprenons les notations et les hypothèses de la Situation \ref{limits-situation-descent-property}.
Si
\begin{enumerate}
\item $f$ est un morphisme quasi-fini,
\item $f_0$ est localement de type fini,
\end{enumerate}
alors il existe un $i \geq 0$ tel que $f_i$ soit quasi-fini.
\end{lemma}

\begin{proof}
Cela découle immédiatement du Lemme \ref{limits-lemma-limit-dimension}.
\end{proof}

\begin{lemma}
\label{limits-lemma-eventually-relative-dimension}
Reprenons les hypothèses et les notations de la Situation \ref{limits-situation-descent-property}.
Soit $d \geq 0$. Si
\begin{enumerate}
\item $f$ est de dimension relative $\leq d$
(Morphismes, Définition \ref{morphisms-definition-relative-dimension-d}),
\item $f_0$ est localement de type fini,
\end{enumerate}
alors il existe un $i$ tel que $f_i$ soit de dimension relative $\leq d$.
\end{lemma}

\begin{proof}
Cela découle immédiatement du Lemme \ref{limits-lemma-limit-dimension}.
\end{proof}

\begin{lemma}
\label{limits-lemma-descend-dimension-d}
Reprenons les notations et les hypothèses de la Situation \ref{limits-situation-descent-property}.
Si
\begin{enumerate}
\item $f$ est de dimension relative $d$,
\item $f_0$ est localement de présentation finie,
\end{enumerate}
alors il existe un $i \geq 0$ tel que $f_i$
soit de dimension relative $d$.
\end{lemma}

\begin{proof}
D'après le Lemme \ref{limits-lemma-limit-dimension}, on peut supposer que toutes les fibres
de $f_0$ sont de dimension $\leq d$. D'après Morphismes, Lemme
\ref{morphisms-lemma-openness-bounded-dimension-fibres-finite-presentation},
l'ensemble $U_0 \subset X_0$ des points $x \in X_0$ tels que
la dimension de la fibre de $X_0 \to Y_0$ en $x$ soit $\leq d - 1$
est ouvert et rétrocompact dans $X_0$. Son complémentaire
$E = X_0 \setminus U_0$ est donc constructible.
De plus, l'image de $X \to X_0$ est contenue dans $E$
d'après Morphismes, Lemme \ref{morphisms-lemma-dimension-fibre-after-base-change}.
Ainsi, pour $i \gg 0$, l'image
de $X_i \to X_0$ est contenue dans $E$
(Lemme \ref{limits-lemma-limit-contained-in-constructible}). Toutes les fibres
de $X_i \to Y_i$ sont alors de dimension $d$ d'après le résultat déjà cité
Morphismes, Lemme \ref{morphisms-lemma-dimension-fibre-after-base-change}.
\end{proof}

\begin{lemma}
\label{limits-lemma-approximate-given-relative-dimension}
Soit $S$ un schéma quasi-compact et quasi-séparé.
Soit $f : X \to S$ un morphisme de présentation finie.
Soit $d \geq 0$ un entier.
Si $Z \subset X$ est un sous-schéma fermé tel que
$\dim(Z_s) \leq d$ pour tout $s \in S$, il existe un
sous-schéma fermé $Z' \subset X$ tel que
\begin{enumerate}
\item $Z \subset Z'$,
\item $Z' \to X$ soit de présentation finie,
\item $\dim(Z'_s) \leq d$ pour tout $s \in S$.
\end{enumerate}
\end{lemma}

\begin{proof}
D'après la
Proposition \ref{limits-proposition-approximate},
on peut écrire $S = \lim S_i$ comme la limite d'un système projectif
filtrant de schémas noethériens à morphismes de transition affines. D'après le
Lemme \ref{limits-lemma-descend-finite-presentation},
on peut supposer qu'il existe un système de morphismes
$f_i : X_i \to S_i$ de présentation finie tels que
$X_{i'} = X_i \times_{S_i} S_{i'}$
pour tout $i' \geq i$ et que $X = X_i \times_{S_i} S$.
Soit $Z_i \subset X_i$ l'image schématique de
$Z \to X \to X_i$. Pour $i' \geq i$, le morphisme $X_{i'} \to X_i$
envoie alors $Z_{i'}$ dans $Z_i$, et le morphisme induit
$Z_{i'} \to Z_i \times_{S_i} S_{i'}$ est une immersion fermée. D'après le
Lemme \ref{limits-lemma-limit-dimension},
les fibres de $Z_i \to S_i$
sont toutes de dimension $\leq d$ pour un indice convenable $i \in I$.
Fixons un tel $i$ et posons $Z' = Z_i \times_{S_i} S \subset X$.
Puisque $S_i$ est noethérien, $X_i$ l'est aussi, et
le morphisme $Z_i \to X_i$ est donc de présentation finie.
Son changement de base $Z' \to X$ est donc lui aussi de présentation finie.
De plus, les fibres de $Z' \to S$ sont des changements de base des fibres
de $Z_i \to S_i$, donc sont de dimension $\leq d$.
\end{proof}






\section{Changement de base en degré maximal}
\label{limits-section-top-degree}

\noindent
Pour un morphisme propre et un module quasi-cohérent de type fini,
l'application de changement de base est un isomorphisme en degré maximal.

\begin{lemma}
\label{limits-lemma-top-cohomology-functor}
Soit $f : X \to Y$ un morphisme de schémas. Soit $d \geq 0$. Faisons les hypothèses suivantes :
\begin{enumerate}
\item $X$ et $Y$ sont quasi-compacts et quasi-séparés,
\item $R^if_*\mathcal{F} = 0$ pour $i > d$ et
tout $\mathcal{O}_X$-module quasi-cohérent $\mathcal{F}$.
\end{enumerate}
On a alors les propriétés suivantes :
\begin{enumerate}
\item[(a)] pour tout diagramme de changement de base
$$
\xymatrix{
X' \ar[d]_{f'} \ar[r]_{g'} & X \ar[d]^f \\
Y' \ar[r]^g & Y
}
$$
on a $R^if'_*\mathcal{F}' = 0$ pour $i > d$ et tout
$\mathcal{O}_{X'}$-module quasi-cohérent $\mathcal{F}'$,
\item[(b)]
$R^df'_*(\mathcal{F}' \otimes_{\mathcal{O}_{X'}} (f')^*\mathcal{G}') =
R^df'_*\mathcal{F}' \otimes_{\mathcal{O}_{Y'}} \mathcal{G}'$
pour tout $\mathcal{O}_{Y'}$-module quasi-cohérent $\mathcal{G}'$,
\item[(c)] la formation de $R^df'_*\mathcal{F}'$ commute à tout
changement de base ultérieur (voir l'explication dans la démonstration).
\end{enumerate}
\end{lemma}

\begin{proof}
Avant de démontrer ces assertions, expliquons le sens de (c). Supposons
qu'un carré cartésien supplémentaire
$$
\xymatrix{
X'' \ar[d]_{f''} \ar[r]_{h'} &
X' \ar[d]_{f'} \ar[r]_{g'} &
X \ar[d]^f \\
Y'' \ar[r]^h &
Y' \ar[r]^g &
Y
}
$$
soit accolé au diagramme donné. Si (a) est satisfaite, il existe une application
canonique $\gamma : h^*R^df'_*\mathcal{F}' \to R^df''_*(h')^*\mathcal{F}'$.
En effet, $\gamma$ est l'application sur les faisceaux de cohomologie de degré $d$
induite par la composée
$$
Lh^*Rf'_*\mathcal{F}' \longrightarrow
Rf''_*L(h')^*\mathcal{F}' \longrightarrow Rf''_*(h')^*\mathcal{F}'
$$
Ici, la première flèche est l'application de changement de base
(Cohomologie, Remarque \ref{cohomology-remark-base-change}), et
la seconde provient de l'application canonique
$L(g')^*\mathcal{F} \to (g')^*\mathcal{F}$.
De même, puisque $Rf'_*\mathcal{F}$ n'a pas de faisceau de cohomologie non nul en
degré $> d$ d'après (a), on a
$H^d(Lh^*Rf_*\mathcal{F}') = h^*R^df_*\mathcal{F}$.
L'assertion (c) signifie que $\gamma$ est un isomorphisme.

\medskip\noindent
Ceci étant, on peut vérifier (a), (b) et (c) localement sur $Y'$ et $Y''$.
Soit $V \subset Y$ un sous-schéma ouvert quasi-compact. Nous
affirmons que (1) et (2) sont satisfaites par $f|_{f^{-1}(V)} : f^{-1}(V) \to V$.
En effet, (1) est immédiate, et (2) résulte du fait que tout module quasi-cohérent
sur $f^{-1}(V)$ est la restriction d'un module quasi-cohérent
sur $X$ (Propriétés, Lemme \ref{properties-lemma-extend-trivial}) et que la formation
des images directes supérieures commute à la restriction aux ouverts.
On peut donc aussi travailler localement sur $Y$. Autrement dit, on peut
supposer $Y''$, $Y'$ et $Y$ affines.

\medskip\noindent
Démonstration de (a) lorsque $Y'$ et $Y$ sont affines. Dans ce cas, les morphismes
$g$ et $g'$ sont affines. Ainsi $g_* = Rg_*$ et $g'_* = Rg'_*$
(Cohomologie des schémas, Lemme \ref{coherent-lemma-relative-affine-vanishing}),
et $g_*$ s'identifie au foncteur de restriction des scalaires sur les modules
(Schémas, Lemme \ref{schemes-lemma-widetilde-pullback}).
Alors
$$
g_*(R^if'_*\mathcal{F}') = H^i(Rg_*Rf'_*\mathcal{F}') =
H^i(Rf_*Rg'_*\mathcal{F}') =
H^i(Rf_*g'_*\mathcal{F}') =
Rf^i_*g'_*\mathcal{F}'
$$
est nul par l'hypothèse (2). D'où (a), compte tenu de notre description de $g_*$.

\medskip\noindent
Démonstration de (b) lorsque $Y'$ est affine, disons $Y' = \Spec(R')$.
D'après (a), on a $H^{d + 1}(X', \mathcal{F}') = 0$
pour tout $\mathcal{O}_{X'}$-module quasi-cohérent $\mathcal{F}'$ ; voir
Cohomologie des schémas, Lemme
\ref{coherent-lemma-quasi-coherence-higher-direct-images-application}.
Considérons le foncteur $F$ sur les $R'$-modules défini par
$$
F(M) = H^d(X', \mathcal{F}' \otimes_{\mathcal{O}_{X'}} (f')^*\widetilde{M})
$$
D'après Cohomologie, Lemme \ref{cohomology-lemma-quasi-separated-cohomology-colimit},
ce foncteur commute aux sommes directes (c'est ici qu'intervient le fait que
$X$, et donc $X'$, est quasi-compact et quasi-séparé).
D'autre part, si $M_1 \to M_2 \to M_3 \to 0$ est une suite exacte,
alors
$$
\mathcal{F}' \otimes_{\mathcal{O}_{X'}} (f')^*\widetilde{M}_1 \to
\mathcal{F}' \otimes_{\mathcal{O}_{X'}} (f')^*\widetilde{M}_2 \to
\mathcal{F}' \otimes_{\mathcal{O}_{X'}} (f')^*\widetilde{M}_3 \to 0
$$
est une suite exacte de modules quasi-cohérents sur $X'$,
et l'annulation de la cohomologie supérieure établie ci-dessus donne une suite exacte
$$
F(M_1) \to F(M_2) \to F(M_3) \to 0
$$
Autrement dit, $F$ est exact à droite. Tout foncteur $R'$-linéaire exact à droite
$F : \text{Mod}_{R'} \to \text{Mod}_{R'}$
qui commute aux sommes directes est donné par le produit tensoriel avec un $R'$-module
(démonstration omise et laissée en exercice).
On obtient ainsi $F(M) = H^d(X', \mathcal{F}') \otimes_{R'} M$.
Puisque $R^d(f')_*\mathcal{F}'$ et
$R^d(f')_*(\mathcal{F}' \otimes_{\mathcal{O}_{X'}} (f')^*\widetilde{M})$
sont quasi-cohérents (Cohomologie des schémas, Lemme
\ref{coherent-lemma-quasi-coherence-higher-direct-images}),
l'égalité $F(M) = H^d(X', \mathcal{F}') \otimes_{R'} M$
se traduit par l'énoncé de (b).

\medskip\noindent
Démonstration de (c) lorsque $Y'' \to Y' \to Y$ sont des morphismes de schémas affines.
Écrivons $Y'' = \Spec(R'')$ et $Y' = \Spec(R')$.
Alors $R^df''_*(h')^*\mathcal{F}'$
est le module quasi-cohérent sur $Y'$ associé au $R''$-module
$H^d(X'', (h')^*\mathcal{F}')$. Or $h' : X'' \to X'$ est affine ;
ainsi $H^d(X'', (h')^*\mathcal{F}') = H^d(X, h'_*(h')^*\mathcal{F}')$
d'après le résultat déjà utilisé
Cohomologie des schémas, Lemme \ref{coherent-lemma-relative-affine-cohomology}.
On a
$$
h'_*(h')^*\mathcal{F}' =
\mathcal{F}' \otimes_{\mathcal{O}_{X'}} (f')^*\widetilde{R''}
$$
comme on le vérifie sur un recouvrement ouvert affine. Ainsi
$H^d(X'', (h')^*\mathcal{F}') = H^d(X', \mathcal{F}') \otimes_{R'} R''$
d'après (b), appliqué à $f'$, ce qui achève la démonstration.
\end{proof}

\begin{lemma}
\label{limits-lemma-higher-direct-images-zero-above-dimension-fibre}
Soit $f : X \to Y$ un morphisme de schémas. Soit $y \in Y$.
Supposons $f$ propre et $\dim(X_y) = d$.
Alors
\begin{enumerate}
\item pour $\mathcal{F} \in \QCoh(\mathcal{O}_X)$,
on a $(R^if_*\mathcal{F})_y = 0$ pour tout $i > d$,
\item il existe un voisinage ouvert affine $V \subset Y$
de $y$ tel que $f^{-1}(V) \to V$ et $d$ satisfassent les
hypothèses et les conclusions du Lemme \ref{limits-lemma-top-cohomology-functor}.
\end{enumerate}
\end{lemma}

\begin{proof}
D'après Morphismes, Lemme
\ref{morphisms-lemma-openness-bounded-dimension-fibres},
et puisque $f$ est fermé, on peut trouver un voisinage ouvert affine
$V$ de $y$ tel que les fibres au-dessus des points de $V$ soient toutes de dimension
$\leq d$. On peut donc supposer que $X \to Y$ est un morphisme propre dont
toutes les fibres sont de dimension $\leq d$, avec $Y$ affine.
Nous montrerons (2), ce qui entraînera immédiatement (1)
pour tout $y \in Y$.

\medskip\noindent
D'après le Lemme \ref{limits-lemma-proper-limit-of-proper-finite-presentation},
on peut écrire $X = \lim X_i$ comme une limite projective, où $X_i \to Y$ est propre
et de présentation finie, et où les morphismes $X \to X_i$
ainsi que les morphismes de transition sont des immersions fermées.
Pour un certain $i$, les fibres de $X_i \to Y$ sont de dimension $\leq d$ ;
voir Lemme \ref{limits-lemma-limit-dimension}.
Pour un $\mathcal{O}_X$-module quasi-cohérent $\mathcal{F}$, on a
$R^pf_*\mathcal{F} = R^pf_{i, *}(X \to X_i)_*\mathcal{F}$ d'après
Cohomologie des schémas, Lemme \ref{coherent-lemma-relative-affine-vanishing},
et Leray (Cohomologie, Lemme \ref{cohomology-lemma-relative-Leray}).
On peut donc remplacer $X$ par $X_i$ et
se ramener au cas traité au paragraphe suivant.

\medskip\noindent
Supposons $Y$ affine et $f : X \to Y$ propre et de présentation finie,
toutes les fibres étant de dimension $\leq d$. Il suffit de montrer que
$H^p(X, \mathcal{F}) = 0$ pour $p > d$. En effet, d'après
Cohomologie des schémas, Lemme
\ref{coherent-lemma-quasi-coherence-higher-direct-images-application},
on a $H^p(X, \mathcal{F}) = H^0(Y, R^pf_*\mathcal{F})$.
D'autre part, $R^pf_*\mathcal{F}$ est quasi-cohérent sur $Y$
d'après Cohomologie des schémas, Lemme
\ref{coherent-lemma-quasi-coherence-higher-direct-images} ;
l'annulation de ses sections globales entraîne donc son annulation.
Écrivons $Y = \lim_{i \in I} Y_i$ comme une limite projective de schémas affines,
où $Y_i$ est le spectre d'un anneau noethérien
(par exemple d'une $\mathbf{Z}$-algèbre de type fini).
On peut choisir un élément $0 \in I$ et un morphisme de type fini
$X_0 \to Y_0$ tels que $X \cong Y \times_{Y_0} X_0$ ; voir
Lemme \ref{limits-lemma-descend-finite-presentation}.
En augmentant $0$, on peut supposer $X_0 \to Y_0$ propre
(Lemme \ref{limits-lemma-eventually-proper})
et les fibres de $X_0 \to Y_0$ de dimension $\leq d$
(Lemme \ref{limits-lemma-limit-dimension}).
Puisque $X \to X_0$ est affine, on a
$H^p(X, \mathcal{F}) = H^p(X_0, (X \to X_0)_*\mathcal{F})$ d'après
Cohomologie des schémas, Lemme \ref{coherent-lemma-relative-affine-cohomology}.
Cela nous ramène au cas traité au paragraphe suivant.

\medskip\noindent
Supposons $Y$ affine noethérien et $f : X \to Y$ propre,
toutes les fibres étant de dimension $\leq d$.
Dans ce cas, on peut écrire $\mathcal{F} = \colim \mathcal{F}_i$
comme une limite inductive filtrante de $\mathcal{O}_X$-modules cohérents ; voir
Propriétés, Lemme
\ref{properties-lemma-directed-colimit-finite-presentation}.
Alors $H^p(X, \mathcal{F}) = \colim H^p(X, \mathcal{F}_i)$ d'après
Cohomologie, Lemme \ref{cohomology-lemma-quasi-separated-cohomology-colimit}.
On peut donc supposer $\mathcal{F}$ cohérent.
Dans ce cas, on a $(R^pf_*\mathcal{F})_y = 0$ pour
tout $y \in Y$ d'après Cohomologie des schémas, Lemme
\ref{coherent-lemma-higher-direct-images-zero-above-dimension-fibre}.
Ainsi $R^pf_*\mathcal{F} = 0$, donc
$H^p(X, \mathcal{F}) = 0$ (voir ci-dessus), ce qui conclut.
\end{proof}

\begin{lemma}
\label{limits-lemma-proper-top-cohomology-finite-type}
Soit $f : X \to Y$ un morphisme de schémas. Soit $d \geq 0$. Soit $\mathcal{F}$
un $\mathcal{O}_X$-module. Faisons les hypothèses suivantes :
\begin{enumerate}
\item $f$ est un morphisme propre dont toutes les fibres sont de dimension $\leq d$,
\item $\mathcal{F}$ est un $\mathcal{O}_X$-module quasi-cohérent de type fini.
\end{enumerate}
Alors $R^df_*\mathcal{F}$ est un $\mathcal{O}_X$-module quasi-cohérent
de type fini.
\end{lemma}

\begin{proof}
Le module $R^df_*\mathcal{F}$ est quasi-cohérent d'après
Cohomologie des schémas, Lemme
\ref{coherent-lemma-quasi-coherence-higher-direct-images}.
La question est locale sur $Y$ ; on peut donc supposer $Y$ affine.
Écrivons $Y = \Spec(R)$. Il suffit alors de montrer que $H^d(X, \mathcal{F})$
est un $R$-module de type fini.

\medskip\noindent
D'après le Lemme \ref{limits-lemma-proper-limit-of-proper-finite-presentation},
on peut écrire $X = \lim X_i$ comme une limite projective, où $X_i \to Y$ est propre
et de présentation finie, et où les morphismes $X \to X_i$
ainsi que les morphismes de transition sont des immersions fermées.
Pour un certain $i$, les fibres de $X_i \to Y$ sont de dimension $\leq d$ ;
voir Lemme \ref{limits-lemma-limit-dimension}. On a
$R^pf_*\mathcal{F} = R^pf_{i, *}(X \to X_i)_*\mathcal{F}$ d'après
Cohomologie des schémas, Lemme \ref{coherent-lemma-relative-affine-vanishing},
et Leray (Cohomologie, Lemme \ref{cohomology-lemma-relative-Leray}).
On peut donc remplacer $X$ par $X_i$ et
se ramener au cas traité au paragraphe suivant.

\medskip\noindent
Supposons $Y$ affine et $f : X \to Y$ propre et de présentation finie,
toutes les fibres étant de dimension $\leq d$.
On peut écrire $\mathcal{F}$ comme quotient d'un
$\mathcal{O}_X$-module de présentation finie $\mathcal{F}'$ ; voir
Propriétés, Lemme
\ref{properties-lemma-finite-directed-colimit-surjective-maps}.
L'application $H^d(X, \mathcal{F}') \to H^d(X, \mathcal{F})$ est
surjective, car $H^{d + 1}(X, \Ker(\mathcal{F}' \to \mathcal{F})) = 0$
d'après l'annulation de la cohomologie supérieure établie dans le
Lemme \ref{limits-lemma-higher-direct-images-zero-above-dimension-fibre}
(ou sa démonstration). On est ainsi ramené au cas traité au paragraphe suivant.

\medskip\noindent
Supposons $Y = \Spec(R)$ affine et $f : X \to Y$
propre et de présentation finie,
toutes les fibres étant de dimension $\leq d$, et $\mathcal{F}$ un
$\mathcal{O}_X$-module de présentation finie.
Écrivons $Y = \lim_{i \in I} Y_i$ comme une limite projective de schémas affines,
où $Y_i = \Spec(R_i)$ est le spectre d'un anneau noethérien
(par exemple d'une $\mathbf{Z}$-algèbre de type fini).
On peut choisir un élément $0 \in I$ et un morphisme de type fini
$X_0 \to Y_0$ tels que $X \cong Y \times_{Y_0} X_0$ ; voir
Lemme \ref{limits-lemma-descend-finite-presentation}.
En augmentant $0$, on peut supposer $X_0 \to Y_0$ propre
(Lemme \ref{limits-lemma-eventually-proper})
et les fibres de $X_0 \to Y_0$ de dimension $\leq d$
(Lemme \ref{limits-lemma-limit-dimension}).
En augmentant $0$, on peut supposer qu'il existe un
$\mathcal{O}_{X_0}$-module cohérent $\mathcal{F}_0$ dont l'image
inverse est $\mathcal{F}$ ; voir
Lemme \ref{limits-lemma-descend-modules-finite-presentation}.
D'après le Lemme \ref{limits-lemma-top-cohomology-functor},
on a
$$
H^d(X, \mathcal{F}) = H^d(X_0, \mathcal{F}_0) \otimes_{R_0} R
$$
Cela achève la démonstration, puisque le module de cohomologie
$H^d(X_0, \mathcal{F}_0)$ est de type fini d'après
Cohomologie des schémas, Lemme
\ref{coherent-lemma-proper-over-affine-cohomology-finite}.
\end{proof}

\begin{lemma}
\label{limits-lemma-proper-top-cohomology-finite-presentation}
Soit $f : X \to Y$ un morphisme de schémas. Soit $d \geq 0$.
Soit $\mathcal{F}$ un $\mathcal{O}_X$-module. Faisons les hypothèses suivantes :
\begin{enumerate}
\item $f$ est un morphisme propre de présentation finie
dont toutes les fibres sont de dimension $\leq d$,
\item $\mathcal{F}$ est un $\mathcal{O}_X$-module de présentation finie.
\end{enumerate}
Alors $R^df_*\mathcal{F}$ est un $\mathcal{O}_X$-module
de présentation finie.
\end{lemma}

\begin{proof}
La démonstration est exactement la même que celle du
Lemme \ref{limits-lemma-proper-top-cohomology-finite-type},
à ceci près que l'on peut omettre le troisième paragraphe.
Nous omettons les détails.
\end{proof}











\section{Recollement au voisinage des fibres fermées}
\label{limits-section-change-over-closed-points}

\noindent
En appliquant la théorie précédente au spectre d'un anneau local, on obtient
le résultat de recollement suivant pour les schémas relatifs.

\begin{lemma}
\label{limits-lemma-glueing-near-closed-point}
Soit $S$ un schéma. Soit $s \in S$ un point fermé tel que
$U = S \setminus \{s\} \to S$ soit quasi-compact. En posant
$V = \Spec(\mathcal{O}_{S, s}) \setminus \{s\}$, on a
une équivalence de catégories
$$
\left\{
\begin{matrix}
X \to S\text{ de présentation finie}
\end{matrix}
\right\}
\longrightarrow
\left\{
\vcenter{
\xymatrix{
X' \ar[d] & Y' \ar[d] \ar[l] \ar[r] & Y \ar[d] \\
U & V \ar[l] \ar[r] & \Spec(\mathcal{O}_{S, s})
}
}
\right\}
$$
où, à droite, on considère les diagrammes commutatifs
dont les carrés sont cartésiens et les flèches verticales
de présentation finie.
\end{lemma}

\begin{proof}
Soit $W \subset S$ un voisinage ouvert de $s$. Par
recollement des schémas relatifs, voir
Constructions, section \ref{constructions-section-relative-glueing},
le foncteur
$$
\left\{
\begin{matrix}
X \to S\text{ de présentation finie}
\end{matrix}
\right\}
\longrightarrow
\left\{
\vcenter{
\xymatrix{
X' \ar[d] & Y' \ar[d] \ar[l] \ar[r] & Y \ar[d] \\
U & W \setminus \{s\} \ar[l] \ar[r] & W
}
}
\right\}
$$
est une équivalence de catégories. On a
$\mathcal{O}_{S, s} = \colim \mathcal{O}_W(W)$, où
$W$ parcourt les voisinages ouverts affines de $s$.
Ainsi $\Spec(\mathcal{O}_{S, s}) = \lim W$, où $W$
parcourt les voisinages ouverts affines de $s$.
La catégorie des schémas de présentation finie
sur $\Spec(\mathcal{O}_{S, s})$ est donc la limite des
catégories de schémas de présentation finie sur
$W$, où $W$ parcourt les voisinages ouverts affines
de $s$ ; voir
Lemme \ref{limits-lemma-descend-finite-presentation}.
Pour tout ouvert affine $s \in W$, l'intersection $U \cap W$
est quasi-compacte, puisque $U \to S$ est quasi-compact.
Ainsi $V = \lim W \cap U = \lim W \setminus \{s\}$ est une limite de
schémas quasi-compacts et quasi-séparés (voir
Lemme \ref{limits-lemma-directed-inverse-system-has-limit}).
La catégorie des schémas de présentation finie
sur $V$ est donc elle aussi la limite des
catégories de schémas de présentation finie sur
$W \cap U$, où $W$ parcourt les voisinages ouverts affines
de $s$. Le lemme découle formellement de la combinaison
de ces résultats.
\end{proof}

\begin{lemma}
\label{limits-lemma-glueing-near-closed-point-modules}
Soit $S$ un schéma. Soit $s \in S$ un point fermé tel que
$U = S \setminus \{s\} \to S$ soit quasi-compact. En posant
$V = \Spec(\mathcal{O}_{S, s}) \setminus \{s\}$, on a
une équivalence de catégories
$$
\left\{
\mathcal{O}_S\text{-modules }\mathcal{F}\text{ de présentation finie}
\right\}
\longrightarrow
\left\{
(\mathcal{G}, \mathcal{H}, \alpha)
\right\}
$$
où, à droite, on considère les triplets
constitués d'un $\mathcal{O}_U$-module $\mathcal{G}$ de
présentation finie, d'un $\mathcal{O}_{\Spec(\mathcal{O}_{S, s})}$-module
$\mathcal{H}$ de présentation finie et d'un isomorphisme
$\alpha : \mathcal{G}|_V \to \mathcal{H}|_V$ de
$\mathcal{O}_V$-modules.
\end{lemma}

\begin{proof}
On peut démontrer ce résultat
en reprenant la démonstration du Lemme \ref{limits-lemma-glueing-near-closed-point}
à l'aide du Lemme \ref{limits-lemma-descend-modules-finite-presentation},
ou le déduire du Lemme \ref{limits-lemma-glueing-near-closed-point}
à l'aide de l'équivalence entre les modules quasi-cohérents et
les « fibrés vectoriels » de
Constructions, section \ref{constructions-section-vector-bundle}.
Nous omettons les détails.
\end{proof}

\begin{lemma}
\label{limits-lemma-glueing-near-point}
Soit $S$ un schéma. Soit $U \subset S$ un ouvert rétrocompact.
Soit $s \in S$ un point du complémentaire de $U$. En posant
$V = \Spec(\mathcal{O}_{S, s}) \cap U$, on a
une équivalence de catégories
$$
\colim_{s \in U' \supset U\text{ ouvert}}
\left\{
\vcenter{
\xymatrix{
X \ar[d] \\
U'
}
}
\right\}
\longrightarrow
\left\{
\vcenter{
\xymatrix{
X' \ar[d] & Y' \ar[d] \ar[l] \ar[r] & Y \ar[d] \\
U & V \ar[l] \ar[r] & \Spec(\mathcal{O}_{S, s})
}
}
\right\}
$$
où, à gauche, la flèche verticale est de présentation
finie, et où, à droite, on considère les diagrammes commutatifs
dont les carrés sont cartésiens et les flèches verticales
de présentation finie.
\end{lemma}

\begin{proof}
Soit $W \subset S$ un voisinage ouvert de $s$. Par
recollement des schémas relatifs, voir
Constructions, section \ref{constructions-section-relative-glueing},
le foncteur
$$
\left\{
\begin{matrix}
X \to U' = U \cup W \text{ de présentation finie}
\end{matrix}
\right\}
\longrightarrow
\left\{
\vcenter{
\xymatrix{
X' \ar[d] & Y' \ar[d] \ar[l] \ar[r] & Y \ar[d] \\
U & W \cap U \ar[l] \ar[r] & W
}
}
\right\}
$$
est une équivalence de catégories. On a
$\mathcal{O}_{S, s} = \colim \mathcal{O}_W(W)$, où
$W$ parcourt les voisinages ouverts affines de $s$.
Ainsi $\Spec(\mathcal{O}_{S, s}) = \lim W$, où $W$
parcourt les voisinages ouverts affines de $s$.
La catégorie des schémas de présentation finie
sur $\Spec(\mathcal{O}_{S, s})$ est donc la limite des
catégories de schémas de présentation finie sur
$W$, où $W$ parcourt les voisinages ouverts affines
de $s$ ; voir
Lemme \ref{limits-lemma-descend-finite-presentation}.
Pour tout ouvert affine $s \in W$, l'intersection $U \cap W$
est quasi-compacte, puisque $U \to S$ est quasi-compact.
Ainsi $V = \lim W \cap U$ est une limite de
schémas quasi-compacts et quasi-séparés (voir
Lemme \ref{limits-lemma-directed-inverse-system-has-limit}).
La catégorie des schémas de présentation finie
sur $V$ est donc elle aussi la limite des
catégories de schémas de présentation finie sur
$W \cap U$, où $W$ parcourt les voisinages ouverts affines
de $s$. Le lemme découle formellement de la combinaison
de ces résultats.
\end{proof}

\begin{lemma}
\label{limits-lemma-glueing-near-point-properties}
Reprenons les notations et les hypothèses du Lemme \ref{limits-lemma-glueing-near-point}.
Soit $U \subset U' \subset X$ un ouvert contenant $s$.
\begin{enumerate}
\item Soit $f' : X \to U'$ correspondant à $f : X' \to U$ et
$g : Y \to \Spec(\mathcal{O}_{S, s})$
par l'équivalence. Si $f$ et $g$ sont séparés, propres, finis ou étales,
alors, quitte à rétrécir $U'$, le morphisme $f'$ possède la même propriété.
\item Soit $a : X_1 \to X_2$
un morphisme de schémas de présentation finie sur $U'$,
dont les changements de base sont $a' : X'_1 \to X'_2$ sur $U$ et
$b : Y_1 \to Y_2$ sur $\Spec(\mathcal{O}_{S, s})$.
Si $a'$ et $b$ sont séparés, propres, finis ou étales,
alors, quitte à rétrécir $U'$, le morphisme $a$ possède la même propriété.
\end{enumerate}
\end{lemma}

\begin{proof}
Démonstration de (1). Rappelons que $\Spec(\mathcal{O}_{S, s})$ est la limite des
voisinages ouverts affines de $s$ dans $S$. Puisque $g$ possède la propriété
considérée, la restriction de $f'$ à l'un de ces
voisinages ouverts affines la possède aussi ; voir
Lemmes \ref{limits-lemma-descend-separated-finite-presentation},
\ref{limits-lemma-eventually-proper},
\ref{limits-lemma-descend-finite-finite-presentation} et
\ref{limits-lemma-descend-etale}.
Puisque $f'$ possède la propriété donnée sur $U$, car $f$ la possède,
on conclut en vérifiant cette propriété localement sur la base.

\medskip\noindent
Démonstration de (2). Si l'on écrit $\Spec(\mathcal{O}_{S, s}) = \lim W$,
où $W$ parcourt les voisinages ouverts affines de $s$ dans $S$,
on a $Y_i = \lim W \times_S X_i$. On peut donc utiliser
exactement les mêmes arguments que dans la démonstration de (1).
\end{proof}

\begin{lemma}
\label{limits-lemma-glueing-near-multiple-closed-points}
Soit $S$ un schéma. Soient $s_1, \ldots, s_n \in S$ des points fermés
deux à deux distincts tels que
$U = S \setminus \{s_1, \ldots, s_n\} \to S$ soit quasi-compact. En posant
$S_i = \Spec(\mathcal{O}_{S, s_i})$ et $U_i = S_i \setminus \{s_i\}$,
on a une équivalence de catégories
$$
FP_S \longrightarrow
FP_U \times_{(FP_{U_1} \times \ldots \times FP_{U_n})}
(FP_{S_1} \times \ldots \times FP_{S_n})
$$
où $FP_T$ est la catégorie des schémas de présentation finie sur
le schéma $T$.
\end{lemma}

\begin{proof}
Pour $n = 1$, c'est le Lemme \ref{limits-lemma-glueing-near-closed-point}.
Pour $n > 1$, le lemme peut se démontrer exactement de la même manière ou
s'en déduire. Supposons par exemple que les $f_i : X_i \to S_i$
soient des objets de $FP_{S_i}$, que $f : X \to U$ soit un objet
de $FP_U$ et que l'on se donne des isomorphismes $X_i \times_{S_i} U_i = X \times_U U_i$.
D'après le Lemme \ref{limits-lemma-glueing-near-closed-point}, on peut trouver
un morphisme $f' : X' \to U' = S \setminus \{s_1, \ldots, s_{n - 1}\}$
de présentation finie, isomorphe à
$X_i$ sur $S_i$ et à $X$ sur $U$, ces
isomorphismes étant compatibles avec l'isomorphisme donné
$X_i \times_{S_n} U_n = X \times_U U_n$.
On peut alors appliquer l'hypothèse de récurrence à
$f_i : X_i \to S_i$, $i \leq n - 1$,
$f' : X' \to U'$ et aux isomorphismes
induits $X_i \times_{S_i} U_i = X' \times_{U'} U_i$, $i \leq n - 1$.
Cela montre la surjectivité essentielle. Nous omettons la démonstration de
la pleine fidélité.
\end{proof}






\section{Application aux modifications}
\label{limits-section-modifications-at-a-point}

\noindent
À l'aide des résultats de la section \ref{limits-section-change-over-closed-points},
on peut décrire la catégorie des modifications d'un
schéma au-dessus d'un point fermé en termes d'anneau local.

\begin{lemma}
\label{limits-lemma-modifications}
Soit $S$ un schéma. Soit $s \in S$ un point fermé tel que
$U = S \setminus \{s\} \to S$ soit quasi-compact. En posant
$V = \Spec(\mathcal{O}_{S, s}) \setminus \{s\}$, le foncteur de changement de base
$$
\left\{
\begin{matrix}
f : X \to S\text{ de présentation finie} \\
f^{-1}(U) \to U\text{ est un isomorphisme}
\end{matrix}
\right\}
\longrightarrow
\left\{
\begin{matrix}
g : Y \to \Spec(\mathcal{O}_{S, s})\text{ de présentation finie} \\
g^{-1}(V) \to V\text{ est un isomorphisme}
\end{matrix}
\right\}
$$
est une équivalence de catégories.
\end{lemma}

\begin{proof}
C'est un cas particulier du Lemme \ref{limits-lemma-glueing-near-closed-point}.
\end{proof}

\begin{lemma}
\label{limits-lemma-modifications-properties}
Reprenons les notations et les hypothèses du Lemme \ref{limits-lemma-modifications}.
Soit $f : X \to S$ correspondant à $g : Y \to \Spec(\mathcal{O}_{S, s})$
par l'équivalence. Alors $f$ est séparé, propre, fini, étale,
et d'autres propriétés à ajouter ici, si et seulement si $g$ l'est.
\end{lemma}

\begin{proof}
Le caractère séparé, propre, entier, fini, etc.,
est stable par changement de base. Voir
Schémas, Lemme \ref{schemes-lemma-separated-permanence},
et
Morphismes, Lemmes \ref{morphisms-lemma-base-change-proper} et
\ref{morphisms-lemma-base-change-finite}.
Ainsi, si $f$ possède la propriété, $g$ la possède aussi.
La réciproque découle du Lemme \ref{limits-lemma-glueing-near-point-properties},
mais nous en donnons aussi une démonstration directe.
En effet, si $g$ possède la propriété, $f$ la possède dans un voisinage de $s$ d'après les
Lemmes \ref{limits-lemma-descend-separated-finite-presentation},
\ref{limits-lemma-eventually-proper},
\ref{limits-lemma-descend-finite-finite-presentation} et
\ref{limits-lemma-descend-etale}.
Puisque $f$ possède évidemment la propriété considérée sur $S \setminus \{s\}$,
on conclut en vérifiant cette propriété localement sur la base.
\end{proof}

\begin{remark}
\label{limits-remark-more-general-modification}
Le lemme précédent se généralise comme suit. Soit $S$ un schéma et
soit $T \subset S$ une partie fermée. Supposons qu'il existe un système cofinal
de voisinages ouverts $T \subset W_i$ tel que
(1) $W_i \setminus T$ soit quasi-compact et
(2) $W_i \subset W_j$ soit un morphisme affine.
Alors $W = \lim W_i$ est un schéma contenant $T$
comme sous-schéma fermé. Posons $U = X \setminus T$ et $V = W \setminus T$.
Le foncteur de changement de base
$$
\left\{
\begin{matrix}
f : X \to S\text{ de présentation finie} \\
f^{-1}(U) \to U\text{ est un isomorphisme}
\end{matrix}
\right\}
\longrightarrow
\left\{
\begin{matrix}
g : Y \to W\text{ de présentation finie} \\
g^{-1}(V) \to V\text{ est un isomorphisme}
\end{matrix}
\right\}
$$
est alors une équivalence de catégories. Si nous en avons besoin un jour,
nous transformerons cette remarque en lemme et en donnerons une démonstration détaillée.
\end{remark}



\section{Descente des schémas de type fini}
\label{limits-section-finite-type-quasi-separated}

\noindent
Cette section poursuit le thème de la
section \ref{limits-section-finite-type-closed-in-finite-presentation}
dans l'esprit des résultats de la
section \ref{limits-section-descending-relative}.

\begin{situation}
\label{limits-situation-limit-noetherian}
Soit $S = \lim_{i \in I} S_i$ la limite d'un système projectif filtrant de
schémas noethériens à morphismes de transition affines $S_{i'} \to S_i$
pour $i' \geq i$.
\end{situation}

\begin{lemma}
\label{limits-lemma-good-diagram}
Plaçons-nous dans la Situation \ref{limits-situation-limit-noetherian}.
Soit $X \to S$ quasi-séparé et de type fini.
Il existe alors un $i \in I$ et un diagramme
\begin{equation}
\label{limits-equation-good-diagram}
\vcenter{
\xymatrix{
X \ar[r] \ar[d] & W \ar[d] \\
S \ar[r] & S_i
}
}
\end{equation}
tels que $W \to S_i$ soit de type fini et que
le morphisme induit $X \to S \times_{S_i} W$ soit une immersion
fermée.
\end{lemma}

\begin{proof}
D'après le Lemme \ref{limits-lemma-finite-type-closed-in-finite-presentation},
on peut trouver une immersion fermée $X \to X'$
sur $S$, où $X'$ est un schéma de présentation finie sur $S$.
D'après le Lemme \ref{limits-lemma-descend-finite-presentation},
on peut trouver un $i$ et un morphisme de présentation finie
$X'_i \to S_i$ dont l'image inverse est $X'$. Posons $W = X'_i$.
\end{proof}

\begin{lemma}
\label{limits-lemma-limit-from-good-diagram}
Plaçons-nous dans la Situation \ref{limits-situation-limit-noetherian}.
Soit $X \to S$ quasi-séparé et de type fini.
Étant donnés $i \in I$ et un diagramme
$$
\vcenter{
\xymatrix{
X \ar[r] \ar[d] & W \ar[d] \\
S \ar[r] & S_i
}
}
$$
comme dans (\ref{limits-equation-good-diagram}), pour $i' \geq i$, soit
$X_{i'}$ l'image schématique de $X \to S_{i'} \times_{S_i} W$.
Alors $X = \lim_{i' \geq i} X_{i'}$.
\end{lemma}

\begin{proof}
Puisque $X$ est quasi-compact et quasi-séparé, la formation de
l'image schématique de $X \to S_{i'} \times_{S_i} W$
commute à la restriction aux sous-schémas ouverts
(Morphismes, Lemme \ref{morphisms-lemma-quasi-compact-scheme-theoretic-image}).
On peut donc supposer, et l'on suppose, $W$ affine et son image contenue dans un ouvert affine
$U_i$ de $S_i$. Soient $U \subset S$, $U_{i'} \subset S_{i'}$
les images réciproques de $U_i$. Alors $U$, $U_{i'}$,
$S_{i'} \times_{S_i} W = U_{i'} \times_{U_i} W$ et
$S \times_{S_i} W = U \times_{U_i} W$ sont tous affines.
Il s'ensuit que $X$ est affine, puisque $X \to S \times_{S_i} W$ est
une immersion fermée. Cela montre aussi que l'homomorphisme
$$
\mathcal{O}(U) \otimes_{\mathcal{O}(U_i)} \mathcal{O}(W) \to
\mathcal{O}(X)
$$
est surjectif. Soit $I$ son noyau. Alors $X_{i'}$
est le spectre de l'anneau
$$
\mathcal{O}(X_{i'}) =
\mathcal{O}(U_{i'}) \otimes_{\mathcal{O}(U_i)} \mathcal{O}(W)/I_{i'}
$$
où $I_{i'}$ est l'image réciproque de l'idéal $I$ (voir
Morphismes, Exemple \ref{morphisms-example-scheme-theoretic-image}).
Puisque $\mathcal{O}(U) = \colim \mathcal{O}(U_{i'})$,
on a $I = \colim I_{i'}$, et l'on conclut
que $\colim \mathcal{O}(X_{i'}) = \mathcal{O}(X)$.
\end{proof}

\begin{lemma}
\label{limits-lemma-morphism-good-diagram}
Plaçons-nous dans la Situation \ref{limits-situation-limit-noetherian}.
Soit $f : X \to Y$ un morphisme de schémas quasi-séparés
et de type fini sur $S$. Soient
$$
\vcenter{
\xymatrix{
X \ar[r] \ar[d] & W \ar[d] \\
S \ar[r] & S_{i_1}
}
}
\quad\text{et}\quad
\vcenter{
\xymatrix{
Y \ar[r] \ar[d] & V \ar[d] \\
S \ar[r] & S_{i_2}
}
}
$$
des diagrammes comme dans (\ref{limits-equation-good-diagram}). Soient
$X = \lim_{i \geq i_1} X_i$ et
$Y = \lim_{i \geq i_2} Y_i$ les descriptions
correspondantes comme limites, données par le Lemme \ref{limits-lemma-limit-from-good-diagram}.
Il existe alors un $i_0 \geq \max(i_1, i_2)$ et un morphisme
$$
(f_i)_{i \geq i_0} : (X_i)_{i \geq i_0} \to (Y_i)_{i \geq i_0}
$$
de systèmes projectifs sur $(S_i)_{i \geq i_0}$ tel que
$f = \lim_{i \geq i_0} f_i$.
Si $(g_i)_{i \geq i_0} : (X_i)_{i \geq i_0} \to (Y_i)_{i \geq i_0}$
est un second morphisme de systèmes projectifs sur $(S_i)_{i \geq i_0}$ tel que
$f = \lim_{i \geq i_0} g_i$,
alors $f_i = g_i$ pour tout $i \gg i_0$.
\end{lemma}

\begin{proof}
Puisque $V \to S_{i_2}$ est de présentation finie et
$X = \lim_{i \geq i_1} X_i$, on peut appliquer la Proposition
\ref{limits-proposition-characterize-locally-finite-presentation}
pour trouver un $i_0 \geq \max(i_1, i_2)$ et un morphisme $h : X_{i_0} \to V$
sur $S_{i_2}$ tels que $X \to X_{i_0} \to V$ soit égal à $X \to Y \to V$.
Pour $i \geq i_0$, on obtient un diagramme commutatif de flèches pleines
$$
\xymatrix{
X \ar[d] \ar[r] &
X_i \ar[r] \ar@{..>}[d] \ar@/_2pc/[dd] |!{[d];[ld]}\hole &
X_{i_0} \ar[d]^h \\
Y \ar[r] \ar[d] & Y_i \ar[r] \ar[d] & V \ar[d] \\
S \ar[r] & S_i \ar[r] & S_{i_0}
}
$$
Puisque $X \to X_i$ a une image schématiquement dense
et que $Y_i$ est l'image schématique de
$Y \to S_i \times_{S_{i_2}} V$,
le morphisme $X_i \to S_i \times_{S_{i_2}} V$
induit par le diagramme
se factorise par $Y_i$ (Morphismes, Lemme \ref{morphisms-lemma-factor-factor}).
Cela démontre l'existence.

\medskip\noindent
Unicité. Soit $E_i \subset X_i$ l'égalisateur de $f_i$ et $g_i$
pour $i \geq i_0$. D'après
Schémas, Lemme \ref{schemes-lemma-where-are-they-equal},
$E_i$ est un sous-schéma localement fermé de $X_i$.
Puisque $X_i$ est un sous-schéma fermé de $S_i \times_{S_{i_0}} X_{i_0}$,
et de même pour $Y_i$, on a
$$
E_i = X_i \times_{(S_i \times_{S_{i_0}} X_{i_0})} (S_i \times_{S_{i_0}} E_{i_0})
$$
Pour achever la démonstration, il suffit donc de montrer que $X_i \to X_{i_0}$
se factorise par $E_{i_0}$ pour un certain $i \geq i_0$.
Pour cela, on utilise le fait que $X \to X_{i_0}$ se factorise par $E_{i_0}$,
puisque $f_{i_0}$ et $g_{i_0}$ sont tous deux compatibles avec $f$.
Puisque $X_i$ est noethérien, l'espace topologique
sous-jacent $|E_{i_0}|$ est une partie constructible de $|X_{i_0}|$
(Topologie, Lemme \ref{topology-lemma-constructible-Noetherian-space}).
Ainsi $X_i \to X_{i_0}$ se factorise ensemblistement par $E_{i_0}$
pour $i$ assez grand, d'après le Lemme \ref{limits-lemma-limit-contained-in-constructible}.
Pour un tel $i$, l'image réciproque schématique
$(X_i \to X_{i_0})^{-1}(E_{i_0})$ est un sous-schéma fermé de $X_i$
par lequel $X$ se factorise ; elle est donc égale à $X_i$, puisque
$X \to X_i$ a une image schématiquement dense par construction.
Cela conclut la démonstration.
\end{proof}

\begin{remark}
\label{limits-remark-finite-type-gives-well-defined-system}
Dans la Situation \ref{limits-situation-limit-noetherian},
les Lemmes \ref{limits-lemma-good-diagram}, \ref{limits-lemma-limit-from-good-diagram} et
\ref{limits-lemma-morphism-good-diagram}
montrent que la catégorie des schémas quasi-séparés et
de type fini sur $S$ est équivalente à certains types de
systèmes projectifs de schémas sur $(S_i)_{i \in I}$, à savoir
ceux obtenus en appliquant le Lemme \ref{limits-lemma-limit-from-good-diagram}
à un diagramme de la forme (\ref{limits-equation-good-diagram}).
Par exemple, étant donné $X \to S$ de type fini et quasi-séparé,
si l'on choisit deux diagrammes différents $X \to V_1 \to S_{i_1}$
et $X \to V_2 \to S_{i_2}$ comme dans (\ref{limits-equation-good-diagram}),
le Lemme \ref{limits-lemma-morphism-good-diagram} appliqué à $\text{id}_X$
(dans les deux sens)
montre que les descriptions correspondantes de
$X$ comme limite sont canoniquement isomorphes (quitte à restreindre
l'ensemble ordonné filtrant $I$). Et ainsi de suite.
\end{remark}

\begin{lemma}
\label{limits-lemma-morphism-good-diagram-flat}
Reprenons les notations et les hypothèses du Lemme \ref{limits-lemma-morphism-good-diagram}.
Si $f$ est plat et de présentation finie,
il existe un $i_3 \geq i_0$ tel que, pour $i \geq i_3$,
$f_i$ soit plat, $X_i = Y_i \times_{Y_{i_3}} X_{i_3}$ et
$X = Y \times_{Y_{i_3}} X_{i_3}$.
\end{lemma}

\begin{proof}
D'après le Lemme \ref{limits-lemma-descend-finite-presentation},
on peut choisir un $i \geq i_2$ et un morphisme
$U \to Y_i$ de présentation finie tels que $X = Y \times_{Y_i} U$
(c'est ici qu'intervient la présentation finie de $f$).
En augmentant $i$, on peut supposer $U \to Y_i$ plat ; voir
Lemme \ref{limits-lemma-descend-flat-finite-presentation}.
Comme expliqué dans la Remarque \ref{limits-remark-finite-type-gives-well-defined-system},
on peut remplacer, et l'on remplace, le diagramme initial définissant le système
$(X_i)_{i \geq i_1}$ par le système correspondant à
$X \to U \to S_i$. Ainsi $X_{i'}$, pour $i' \geq i$, est défini comme
l'image schématique de $X \to S_{i'} \times_{S_i} U$.

\medskip\noindent
Puisque $U \to Y_i$ est plat (c'est ici qu'intervient la platitude de $f$),
que $X = Y \times_{Y_i} U$ et
que l'image schématique de $Y \to Y_i$ est $Y_i$,
l'image schématique de $X \to U$ est $U$
(Morphismes, Lemme
\ref{morphisms-lemma-flat-base-change-scheme-theoretic-image}).
Remarquons que $Y_{i'} \to S_{i'} \times_{S_i} Y_i$ est une immersion
fermée pour $i' \geq i$, par construction du système des $Y_j$.
Le même argument que ci-dessus montre alors que l'image schématique
de $X \to S_{i'} \times_{S_i} U$
est égale au sous-schéma fermé $Y_{i'} \times_{Y_i} U$.
Ainsi $X_{i'} = Y_{i'} \times_{Y_i} U$ pour tout $i' \geq i$,
et le lemme est satisfait avec $i_3 = i$.
\end{proof}

\begin{lemma}
\label{limits-lemma-morphism-good-diagram-smooth}
Reprenons les notations et les hypothèses du Lemme \ref{limits-lemma-morphism-good-diagram}.
Si $f$ est lisse, il existe un $i_3 \geq i_0$ tel que, pour
$i \geq i_3$, le morphisme $f_i$ soit lisse.
\end{lemma}

\begin{proof}
Combiner les Lemmes \ref{limits-lemma-morphism-good-diagram-flat} et
\ref{limits-lemma-descend-smooth}.
\end{proof}

\begin{lemma}
\label{limits-lemma-morphism-good-diagram-proper}
Reprenons les notations et les hypothèses du Lemme \ref{limits-lemma-morphism-good-diagram}.
Si $f$ est propre, il existe un $i_3 \geq i_0$ tel que, pour
$i \geq i_3$, le morphisme $f_i$ soit propre.
\end{lemma}

\begin{proof}
D'après la discussion de la
Remarque \ref{limits-remark-finite-type-gives-well-defined-system},
le choix de $i_1$ et de $W$ s'insérant dans un diagramme comme dans
(\ref{limits-equation-good-diagram}) n'influe pas sur la validité
du lemme. Choisissons donc $W$ comme suit.
Choisissons d'abord une immersion fermée $X \to X'$
avec $X' \to S$ propre et de présentation finie ; voir
Lemme \ref{limits-lemma-proper-limit-of-proper-finite-presentation}.
Choisissons ensuite un $i_3 \geq i_2$ et un morphisme propre $W \to Y_{i_3}$
tels que $X' = Y \times_{Y_{i_3}} W$. C'est possible puisque
$Y = \lim_{i \geq i_2} Y_i$, d'après les
Lemmes \ref{limits-lemma-descend-finite-presentation} et
\ref{limits-lemma-eventually-proper}.
Avec ce choix de $W$, la construction montre immédiatement que,
pour $i \geq i_3$, le schéma $X_i$ est un sous-schéma fermé de
$Y_i \times_{Y_{i_3}} W \subset S_i \times_{S_{i_3}} W$,
donc est propre sur $Y_i$.
\end{proof}

\begin{lemma}
\label{limits-lemma-good-diagram-fibre-product}
Dans la Situation \ref{limits-situation-limit-noetherian}, supposons donné un
diagramme cartésien
$$
\xymatrix{
X^1 \ar[r]_p \ar[d]_q & X^3 \ar[d]^a \\
X^2 \ar[r]^b & X^4
}
$$
de schémas quasi-séparés et de type fini sur $S$.
Pour chaque $j = 1, 2, 3, 4$, choisissons $i_j \in I$ et un diagramme
$$
\xymatrix{
X^j \ar[r] \ar[d] & W^j \ar[d] \\
S \ar[r] & S_{i_j}
}
$$
comme dans (\ref{limits-equation-good-diagram}). Soit
$X^j = \lim_{i \geq i_j} X^j_i$ la description correspondante comme limite,
donnée par le Lemme \ref{limits-lemma-morphism-good-diagram}.
Soient $(a_i)_{i \geq i_5}$, $(b_i)_{i \geq i_6}$, $(p_i)_{i \geq i_7}$ et
$(q_i)_{i \geq i_8}$ les morphismes de systèmes correspondants, construits
dans le Lemme \ref{limits-lemma-morphism-good-diagram}. Il existe alors un
$i_9 \geq \max(i_5, i_6, i_7, i_8)$ tel que, pour $i \geq i_9$,
$a_i \circ p_i = b_i \circ q_i$ et que
$$
(q_i, p_i) : X^1_i \longrightarrow X^2_i \times_{b_i, X^4_i, a_i} X^3_i
$$
soit une immersion fermée.
Si $a$ et $b$ sont plats et de présentation finie, il existe un
$i_{10} \geq \max(i_5, i_6, i_7, i_8, i_9)$ tel que, pour $i \geq i_{10}$,
le dernier morphisme affiché soit un isomorphisme.
\end{lemma}

\begin{proof}
D'après la discussion de la
Remarque \ref{limits-remark-finite-type-gives-well-defined-system},
le choix de $W^1$ s'insérant dans un diagramme comme dans
(\ref{limits-equation-good-diagram}) n'influe pas sur la validité
du lemme. On peut donc choisir $W^1 = W^2 \times_{W^4} W^3$.
La construction de $X^1_i$ montre alors immédiatement que
$a_i \circ p_i = b_i \circ q_i$ et que
$$
(q_i, p_i) : X^1_i \longrightarrow X^2_i \times_{b_i, X^4_i, a_i} X^3_i
$$
est une immersion fermée.

\medskip\noindent
Si $a$ et $b$ sont plats et de présentation finie, il en va de même pour
$p$ et $q$, qui sont des changements de base de $a$ et $b$. On peut donc appliquer le
Lemme \ref{limits-lemma-morphism-good-diagram-flat}
à chacun des morphismes $a$, $b$, $p$, $q$ et $a \circ p = b \circ q$.
Il s'ensuit qu'il existe un $i_9 \in I$ tel que
$$
(q_i, p_i) : X^1_i \to X^2_i \times_{X^4_i} X^3_i
$$
soit le changement de base de $(q_{i_9}, p_{i_9})$ par le morphisme
$X^4_i \to X^4_{i_9}$ pour tout $i \geq i_9$.
On conclut que $(q_i, p_i)$ est un isomorphisme pour tout
$i$ assez grand, d'après le Lemme \ref{limits-lemma-descend-isomorphism}.
\end{proof}











% OK, start here.
%

\chapter{Variétés}



\phantomsection
\label{varieties-section-phantom}


\section{Introduction}
\label{varieties-section-introduction}

\noindent
Dans ce chapitre, nous commençons à étudier les variétés et, plus
généralement, les schémas sur un corps. Une référence fondamentale est \cite{EGA}.








\section{Notation}
\label{varieties-section-notation}

\noindent
Tout au long de ce chapitre, nous utilisons la lettre $k$ pour désigner le corps de base.










\section{Variétés}
\label{varieties-section-varieties}

\noindent
Dans le projet Champs, nous adoptons la définition suivante
d'une variété.

\begin{definition}
\label{varieties-definition-variety}
Soit $k$ un corps. Une {\it variété} est un schéma $X$ sur $k$
qui est intègre et dont le morphisme structural
$X \to \Spec(k)$ est séparé et de type fini.
\end{definition}

\noindent
Cette définition présente l'inconvénient suivant. Soit $k'/k$
une extension de corps et soit $X$ une variété sur $k$. Le changement
de base $X_{k'} = X \times_{\Spec(k)} \Spec(k')$ n'est alors
pas nécessairement une variété sur $k'$. Ce phénomène sera étudié
en détail, dans un caunere plus général, dans les sections suivantes.
Le produit de deux variétés n'est pas nécessairement d variété non
plus (il s'agit en fait du même phénomène).
En voici un exemple.


\begin{example}
\label{varieties-example-product-not-a-variety}
Soit $k = \mathbf{Q}$. Posons $X = \Spec(\mathbf{Q}(i))$
et $Y = \Spec(\mathbf{Q}(i))$. Alors le produit
$X \times_{\Spec(k)} Y$ des variétés $X$ et $Y$
n'est pas d variété, car il est réductible. (Il est
isomorphe à la réunion disjointe de deux copies de $X$.)
\end{example}

\noindent
Toutefois, si le corps de base est algébriquement clos, le produit de
variétés est une variété. Cela résulte du chapitre sur l'algèbre, où
l'on traite des situations bien plus générales. Nous en donnons également
ici une démonstration directe et simple.

\begin{lemma}
\label{varieties-lemma-product-varieties}
\begin{slogan}
Les produits de variétés sont des variétés sur des corps algébriquement clos.
\end{slogan}
Soit $k$ un corps algébriquement clos.
Soient $X$ et $Y$ des variétés sur $k$.
Alors $X \times_{\Spec(k)} Y$ est d variété sur $k$.
\end{lemma}

\begin{proof}
Le morphisme $X \times_{\Spec(k)} Y \to \Spec(k)$ est de type
fini et séparé : c'est le composé des morphismes
$X \times_{\Spec(k)} Y \to Y \to \Spec(k)$, qui sont séparés
et de type fini ; voir Morphismes, lemmes
\ref{morphisms-lemma-base-change-finite-type} et
\ref{morphisms-lemma-composition-finite-type}, ainsi que
Schémas, lemme \ref{schemes-lemma-separated-permanence}.
Il reste à montrer que $X \times_{\Spec(k)} Y$ est intègre.
Soient $X = \bigcup_{i = 1, \ldots, n} U_i$ et
$Y = \bigcup_{j = 1, \ldots, m} V_j$ des recouvrements
ouverts affines finis. Il suffit que chacun des schémas
$U_i \times_{\Spec(k)} V_j$ soit intègre, d'après
Propriétés, lemmes \ref{properties-lemma-characterize-reduced},
\ref{properties-lemma-characterize-irreducible} et
\ref{properties-lemma-characterize-integral}.
On est ainsi ramené au cas affine.

\medskip\noindent
Le cas affine équivaut à l'énoncé algébrique suivant : si $A$ et
$B$ sont des anneaux intègres et des $k$-algèbres de type fini, alors $A
\otimes_k B$ est un anneau intègre. Raisonnons par l'absurde et supposons que
$$
(\sum\nolimits_{i = 1, \ldots, n} a_i \otimes b_i)
(\sum\nolimits_{j = 1, \ldots, m} c_j \otimes une_j) = 0
$$
dans $A \otimes_k B$, les deux facteurs étant non nuls dans $A
\otimes_k B$. On peut supposer que $b_1, \ldots, b_n$ sont
linéairement indépendants sur $k$ dans $B$, et de même pour
$d_1, \ldots, d_m$. On peut aussi supposer que $a_1$ et $c_1$
sont non nuls dans $A$. Ainsi, $D(a_1c_1) \subset \Spec(A)$ est
non vide. D'après le théorème des zéros de Hqubert
(Algèbre, théorème \ref{algebra-theorem-nullstellensatz}), il existe
un idéal maximal $\mathfrak m \in D(a_1c_1)$ tel que
$A/\mathfrak m = k$, puisque $k$ est algébriquement clos. Notons
$\overline{a}_i$ et $\overline{c}_j$ les classes de $a_i$ et $c_j$
dans $A/\mathfrak m = k$. L'égalité précédente devient
$$
(\sum\nolimits_{i = 1, \ldots, n} \overline{a}_i b_i)
(\sum\nolimits_{j = 1, \ldots, m} \overline{c}_j une_j) = 0
$$
ce qui contredit $\mathfrak m \in D(a_1c_1)$, l'indépendance linéaire
des famqules $b_1, \ldots, b_n$ et $une_1, \ldots, une_m$, ainsi que
le fait que $B$ est un anneau intègre.
\end{proof}






\section{Variétés et applications rationnelles}
\label{varieties-section-varieties-rational-maps}

\noindent
Soit $k$ un corps et soient $X$ et $Y$ des variétés sur $k$.
Par {\it application rationnelle de variétés de $X$ vers $Y$}, nous entendrons
une application rationnelle sur $\Spec(k)$ du schéma $X$ vers le schéma $Y$,
au sens de Morphismes, définition \ref{morphisms-definition-rational-map}.
Selon l'usage, l'expression «~application rationnelle de variétés~» ne
mentionne pas le corps de base (commun) des variétés, bien que, pour des
schémas généraux, nous distinguions les applications rationnelles des
applications rationnelles sur une base donnée.

\medskip\noindent
Le titre de cette section fait référence au théorème fondamental suivant.

\begin{theorem}
\label{varieties-theorem-varieties-rational-maps}
Soit $k$ un corps. La catégorie des variétés
et des applications rationnelles dominantes est équivalente
à la catégorie des extensions de corps de type fini $K/k$.
\end{theorem}

\begin{proof}
Soient $X$ et $Y$ des variétés de points génériques respectifs
$x \in X$ et $y \in Y$. Rappelons que les applications rationnelles
dominantes de $X$ vers $Y$ sont précisément les applications
rationnelles qui envoient $x$ sur $y$ (Morphismes,
définition \ref{morphisms-definition-dominant-rational} et discussion
qui suit). Toute application rationnelle dominante $X
\supset U \to Y$ induit donc un homomorphisme de corps de fonctions
$$
k(Y) = \kappa(y) = \mathcal{O}_{Y, y}
\longrightarrow
\mathcal{O}_{X, x} = \kappa(x) = k(X)
$$
Inversement, un tel homomorphisme de $k$-algèbres (automatiquement
local, puisque sa source et son but sont des corps) détermine de
manière unique d application rationnelle dominante, d'après Morphismes,
lemme \ref{morphisms-lemma-rational-map-finite-presentation}. On obtient ainsi
un foncteur pleinement fidèle. Il reste à voir que toute extension de
corps de type fini $K/k$ appartient à l'image essentielle. Comme $K/k$
est de type fini, qu existe d $k$-algèbre de type fini $A \subset K$
dont $K$ est le corps des fractions. Alors $X = \Spec(A)$ est d variété
de corps de fonctions $K$.
\end{proof}

\noindent
Soit $k$ un corps et soient $X$ et $Y$ des variétés sur $k$.
Nous dirons que {\it $X$ et $Y$ sont des variétés birationnelles}
si $X$ et $Y$ sont birationnels sur $\Spec(k)$, au sens de
Morphismes, définition \ref{morphisms-definition-birational}.
Selon l'usage, l'expression «~variétés birationnelles~» ne mentionne pas
le corps de base (commun) des variétés, bien que, pour des schémas
irréductibles généraux, nous distinguions le fait d'être birationnels
du fait de l'être sur une base donnée.

\begin{lemma}
\label{varieties-lemma-birational-varieties}
Soient $X$ et $Y$ des variétés sur un corps $k$.
Les conditions suivantes sont équivalentes :
\begin{enumerate}
\item $X$ et $Y$ sont des variétés birationnelles,
\item les corps de fonctions $k(X)$ et $k(Y)$ sont isomorphes,
\item il existe des ouverts non vides de $X$ et de $Y$ qui sont
isomorphes comme variétés,
\item il existe un ouvert $U \subset X$ et un morphisme
birationnel de variétés $U \to Y$.
\end{enumerate}
\end{lemma}

\begin{proof}
C'est un cas particulier de Morphismes, lemme
\ref{morphisms-lemma-criterion-birational-finite-presentation}.
\end{proof}







\section{Changement de corps et anneaux locaux}
\label{varieties-section-local-rings}

\noindent
Voici quelques résultats préliminaires sur le comportement des
anneaux locaux par extension du corps de base.

\begin{lemma}
\label{varieties-lemma-change-fields-flat}
Soit $K/k$ d extension de corps. Soit $X$ un schéma sur $k$
et posons $Y = X_K$. Si $y \in Y$ a pour image $x \in X$, alors
\begin{enumerate}
\item $\mathcal{O}_{X, x} \to \mathcal{O}_{Y, y}$ est un
homomorphisme local fidèlement plat,
\item si $\mathfrak p_0 = \Ker(\kappa(x) \otimes_k K \to \kappa(y))$,
alors $\kappa(y) = \kappa(\mathfrak p_0)$,
\item $\mathcal{O}_{Y, y} =
(\mathcal{O}_{X, x} \otimes_k K)_\mathfrak p$, où
$\mathfrak p \subset \mathcal{O}_{X, x} \otimes_k K$ est l'image
réciproque de $\mathfrak p_0$,
\item on a
$\mathcal{O}_{Y, y}/\mathfrak m_x\mathcal{O}_{Y, y} =
(\kappa(x) \otimes_k K)_{\mathfrak p_0}$.
\end{enumerate}
\end{lemma}

\begin{proof}
On peut supposer $X = \Spec(A)$ affine. Alors
$Y = \Spec(A \otimes_k K)$. Puisque $K$ est plat sur $k$,
$A \to A \otimes_k K$ est plat. Il en va donc de même de $Y \to X$,
et la première assertion résulte aussi d'Algèbre, lemme
\ref{algebra-lemma-local-flat-ff}. La deuxième résulte de
Schémas, lemme \ref{schemes-lemma-points-fibre-product}. Le point
$y$ correspond à un idéal premier $\mathfrak q \subset A \otimes_k K$
et $x$ à $\mathfrak r = A \cap \mathfrak q$. Alors $\mathfrak p_0$
est le noyau de l'homomorphisme induit
$\kappa(\mathfrak r) \otimes_k K \to \kappa(\mathfrak q)$.
L'homomorphisme d'anneaux locaux est
$$
A_\mathfrak r \longrightarrow (A \otimes_k K)_\mathfrak q
$$
Il se factorise par
$A_\mathfrak r \otimes_k K =
(A \otimes_k K)_{\mathfrak r}$ sous la forme
$$
A_\mathfrak r \longrightarrow A_\mathfrak r \otimes_k K
\longrightarrow (A \otimes_k K)_\mathfrak q
$$
où la deuxième flèche est la localisation en un certain idéal
premier. Celui-ci est l'image réciproque de $\mathfrak p_0$
(nous omettons les détails), ce qui prouve (3). Pour (4), on utilise
(3) et l'exactitude de la localisation et du foncteur $- \otimes_k K$.
\end{proof}

\begin{lemma}
\label{varieties-lemma-change-fields-algebraic-dim}
Avec les notations du lemme \ref{varieties-lemma-change-fields-flat},
si $X$ est localement de type fini sur $k$, alors
$$
\dim(\mathcal{O}_{Y, y}/\mathfrak m_x\mathcal{O}_{Y, y}) =
\text{trdeg}_k(\kappa(x)) - \text{trdeg}_K(\kappa(y)) =
\dim(\mathcal{O}_{Y, y}) - \dim(\mathcal{O}_{X, x})
$$
\end{lemma}

\begin{proof}
Il s'agit d'une reformulation d'Algèbre, lemme
\ref{algebra-lemma-inequalities-under-field-extension}.
\end{proof}

\begin{lemma}
\label{varieties-lemma-change-fields-algebraic-unramified}
Avec les notations du lemme
\ref{varieties-lemma-change-fields-flat}, supposons $X$ localement de type fini
sur $k$, $\dim(\mathcal{O}_{X, x}) = \dim(\mathcal{O}_{Y, y})$
et $\kappa(x) \otimes_k K$ réduit (par exemple si
$\kappa(x)/k$ est séparable ou si $K/k$ est séparable). Alors
$\mathfrak m_x \mathcal{O}_{Y, y} = \mathfrak m_y$.
\end{lemma}

\begin{proof}
(L'assertion entre parenthèses résulte d'Algèbre, lemme
\ref{algebra-lemma-separable-extension-preserves-reducedness}.)
En combinant les lemmes \ref{varieties-lemma-change-fields-flat} et
\ref{varieties-lemma-change-fields-algebraic-dim}, on voit que
$\mathcal{O}_{Y, y}/\mathfrak m_x \mathcal{O}_{Y, y}$ a une
dimension $0$ et est réduit. C'est donc un corps.
\end{proof}





\section{Schémas géométriquement réduits}
\label{varieties-section-geometrically-reduced}

\noindent
Si $X$ est un schéma réduit sur un corps, il peut arriver que
$X$ devienne non réduit après l'extension du corps de base.
Cela ne se produit pas pour les schémas géométriquement réduits.

\begin{definition}
\label{varieties-definition-geometrically-reduced}
Soit $k$ un corps.
Soit $X$ un schéma sur $k$.
\begin{enumerate}
\item Soit $x \in X$. On dit que $X$ est
{\it géométriquement réduit en $x$} si, pour toute
extension de corps $k'/k$ et tout point $x' \in X_{k'}$
au-dessus de $x$, l'anneau local $\mathcal{O}_{X_{k'}, x'}$
est réduit.
\item On dit que $X$ est {\it géométriquement réduit} sur $k$
si $X$ est géométriquement réduit en tout point de $X$.
\end{enumerate}
\end{definition}

\noindent
Cette définition peut paraître quelque peu mystérieuse, mais elle
coïncide en réalité avec la notion étudiée dans le chapitre sur
l'algèbre. Les résultats élémentaires qui suivent précisent ce lien.
Voir aussi le lemme
\ref{varieties-lemma-geometrically-reduced-dense-smooth-open}, d'après lequel
«~géométriquement réduit~» équivaut à «~génériquement lisse~» pour les variétés sur un corps.

\begin{lemma}
\label{varieties-lemma-geometrically-reduced-at-point}
\begin{slogan}
La propriété d'être géométriquement réduit se vérifie sur les anneaux locaux.
\end{slogan}
Soit $k$ un corps.
Supposons que $X$ soit un schéma
sur $k$. Soit $x \in X$. Les
énoncés suivants sont équivalents
\begin{enumerate}
\item $X$ est géométriquement réduit en $x$, et
\item l'anneau $\mathcal{O}_{X, x}$ est
géométriquement réduit sur $k$ (voir Algèbre,
définition \ref{algebra-definition-geometrically-reduced}).
\end{enumerate}
\end{lemma}

\begin{proof}
Supposons (1). Cela implique en particulier que
$\mathcal{O}_{X, x}$ est réduit. Supposons que $k'/k$ soit une
extension de corps finie purement inséparable. Considérons
l'anneau $\mathcal{O}_{X, x} \otimes_k k'$. D'après l'algèbre,
le lemme \ref{algebra-lemma-p-ring-map}, son spectre est le
même que le spectre de $\mathcal{O}_{X, x}$. Par conséquent,
c'est également un anneau local (Algèbre, lemme
\ref{algebra-lemma-characterize-local-ring}). Donc, il existe un
point unique $x' \in X_{k'}$ au-dessus de $x$ et de
$\mathcal{O}_{X_{k'}, x'} \cong \mathcal{O}_{X, x} \otimes_k k'$, voir
le lemme \ref{varieties-lemma-change-fields-flat}. Par hypothèse, il s'agit d'un
anneau réduit. Ainsi, nous déduisons (2) d'après Algèbre, lemme
\ref{algebra-lemma-geometrically-reduced-finite-purely-inseparable-extension}.

\medskip\noindent
Supposons (2). Supposons que $k'/k$ soit une extension
de corps. Puisque $\Spec(k') \to \Spec(k)$ est surjectif,
$X_{k'} \to X$ est aussi surjectif (Morphismes, lemme
\ref{morphisms-lemma-base-change-surjective}). Supposons
que $x' \in X_{k'}$ soit un point quelconque situé
au-dessus de $x$. L'anneau local $\mathcal{O}_{X_{k'}, x'}$
est une localisation de l'anneau $\mathcal{O}_{X, x}
\otimes_k k'$ d'après le lemme \ref{varieties-lemma-change-fields-flat}.
Par conséquent, il est réduit par hypothèse et (1) est prouvé.
\end{proof}

\noindent
La notion n'est pas intéressante en caractéristique zéro.

\begin{lemma}
\label{varieties-lemma-perfect-reduced}
Supposons que $X$ soit un schéma sur un corps parfait $k$ (par
exemple \ $k$ a une caractéristique nulle). Soit $x \in X$. Si
$\mathcal{O}_{X, x}$ est réduit, alors $X$ est géométriquement réduit
en $x$. Si $X$ est réduit, alors $X$ est géométriquement réduit sur $k$.
\end{lemma}

\begin{proof}
La première affirmation découle du lemme
\ref{varieties-lemma-geometrically-reduced-at-point} et de
l'Algèbre, du lemme
\ref{algebra-lemma-separable-extension-preserves-reducedness} et de
la définition d'un corps parfait (Algèbre, définition
\ref{algebra-definition-perfect}). La seconde affirmation découle de la première.
\end{proof}

\begin{lemma}
\label{varieties-lemma-geometrically-reduced}
Supposons que $k$ soit un corps de caractéristique $p > 0$. Supposons
que $X$ soit un schéma sur $k$. Les énoncés suivants sont équivalents
\begin{enumerate}
\item $X$ est géométriquement réduit, \item
$X_{k'}$ est réduit pour toute extension de
corps $k'/k$, \item $X_{k'}$ est réduit pour
toute extension de corps finie purement
inséparable $k'/k$, \item $X_{k^{1/p}}$ est
réduit, \item $X_{k^{perf}}$ est réduit, \item
$X_{\bar k}$ est réduit, \item pour tout ouvert
affine $U \subset X$ l'anneau $\mathcal{O}_X(U)$
est géométriquement réduit (voir Algèbre,
définition \ref{algebra-definition-geometrically-reduced}).
\end{enumerate}
\end{lemma}

\begin{proof}
Supposons (1). Alors pour toute extension de corps $k'/k$ et
chaque point $x' \in X_{k'}$, l'anneau local de $X_{k'}$ en $x'$
est réduit. En d'autres termes, $X_{k'}$ est réduit. D'où (2).

\medskip\noindent
Supposons (2). Supposons que $U \subset X$ soit un ouvert affine.
Alors, pour toute extension de corps $k'/k$, le schéma $X_{k'}$
est réduit, donc $U_{k'} = \Spec(\mathcal{O}(U)\otimes_k k')$ est
réduit, donc $\mathcal{O}(U)\otimes_k k'$ est réduit (voir
Propriétés, section \ref{properties-section-integral}). En d'autres
termes, $\mathcal{O}(U)$ est géométriquement réduit, donc (7) est vérifié.

\medskip\noindent
Supposons (7). Pour toute extension de corps $k'/k$, le
changement de base $X_{k'}$ s'obtient par recollement des spectres
des anneaux $\mathcal{O}_X(U) \otimes_k k'$ où $U$ est un
ouvert affine dans $X$ (voir Schémas, section
\ref{schemes-section-fibre-products}). Donc $X_{k'}$ est réduit. Ainsi (1) est vrai.

\medskip\noindent
Cela prouve que (1), (2) et (7) sont équivalents. Ceux-ci sont
équivalents à (3), (4), (5) et (6) car nous pouvons appliquer
Algèbre, lemme
\ref{algebra-lemma-geometrically-reduced-finite-purely-inseparable-extension} à $\mathcal{O}_X(U)$ pour $U \subset X$ ouvert affine.
\end{proof}

\begin{lemma}
\label{varieties-lemma-check-only-finite-inseparable-extensions}
Supposons que $k$ soit un corps de caractéristique $p > 0$. Supposons que $X$
soit un schéma sur $k$. Soit $x \in X$. Les énoncés suivants sont équivalents
\begin{enumerate}
\item $X$ est géométriquement réduit en $x$, \item
$\mathcal{O}_{X_{k'}, x'}$ est réduit pour chaque
extension de corps finie purement inséparable $k'$ de
$k$ et $x' \in X_{k'}$ le point unique au-dessus de $x$,
\item $\mathcal{O}_{X_{k^{1/p}}, x'}$ est réduit pour
$x' \in X_{k^{1/p}}$ le point unique au-dessus de $x$,
et \item $\mathcal{O}_{X_{k^{perf}}, x'}$ est réduit pour
$x' \in X_{k^{perf}}$ le point unique au-dessus de $x$.
\end{enumerate}
\end{lemma}

\begin{proof}
Notons que si $k'/k$ est purement inséparable, alors $X_{k'}
\to X$ induit un homéomorphisme sur les espaces
topologiques sous-jacents, voir Algèbre, lemme
\ref{algebra-lemma-p-ring-map}. D'où l'unicité de $x'$
au-dessus de $x$ mentionnée dans l'énoncé. De plus, dans ce cas
$\mathcal{O}_{X_{k'}, x'} = \mathcal{O}_{X, x} \otimes_k k'$.
Par conséquent, le lemme découle du lemme
\ref{varieties-lemma-geometrically-reduced-at-point} ci-dessus et du lemme
\ref{algebra-lemma-geometrically-reduced-finite-purely-inseparable-extension} d'Algèbre.
\end{proof}

\begin{lemma}
\label{varieties-lemma-geometrically-reduced-upstairs}
Soit $k$ un corps. Supposons que $X$
soit un schéma sur $k$. Supposons que $k'/k$ soit
une extension de corps. Supposons que $x \in X$ soit
un point, et soit $x' \in X_{k'}$ un point situé
au-dessus de $x$. Les énoncés suivants sont équivalents
\begin{enumerate}
\item $X$ est géométriquement réduit en $x$,
\item $X_{k'}$ est géométriquement réduit en $x'$.
\end{enumerate}
En particulier, $X$ est géométriquement réduit sur $k$ si et
seulement si $X_{k'}$ est géométriquement réduit sur $k'$.
\end{lemma}

\begin{proof}
Il est clair que (1) implique (2). Supposons (2). Supposons que
$k''/k$ soit une extension de corps finie purement inséparable et
soit $x'' \in X_{k''}$ un point au-dessus de $x$ (en fait, il
est unique). Choisissons une extension de corps commune $k'''/k$ de
$k''/k$ et $k'/k$, voir corps, lemme
\ref{fields-lemma-common-extension-field}. Choisissons un point $x''' \in X_{k'''}$
au-dessus à la fois de $x'$ et $x''$. Considérons le morphisme d'anneaux locaux
$$
\mathcal{O}_{X_{k''}, x''} \longrightarrow \mathcal{O}_{X_{k'''}, x'''}.
$$
Ceci est un homomorphisme local plat d'anneaux et donc fidèlement
plat. Par (2), on voit que l'anneau local à droite est réduit.
Ainsi, d'après L'Algèbre, lemme
\ref{algebra-lemma-descent-reduced}, nous concluons que
$\mathcal{O}_{X_{k''}, x''}$ est réduit. Ainsi, d'après le lemme
\ref{varieties-lemma-check-only-finite-inseparable-extensions}, nous concluons que $X$ est géométriquement réduit en $x$.
\end{proof}

\begin{lemma}
\label{varieties-lemma-geometrically-reduced-any-base-change}
Soit $k$ un corps. Supposons que $X$, $Y$ soient des schémas sur $k$.
\begin{enumerate}
\item Si $X$ est géométriquement réduit en $x$, et
$Y$ est réduit, alors $X \times_k Y$ est réduit
en tout point au-dessus de $x$. \item Si $X$ est
géométriquement réduit sur $k$ et $Y$ est réduit,
alors $X \times_k Y$ est réduit. \item Si $X$ et
$Y$ sont géométriquement réduits sur $k$, alors $X
\times_k Y$ est géométriquement réduit. \item Si
$k$ est parfait et $X$ et $Y$ sont réduits, alors
$X \times_k Y$ est réduit. \item Ajouter plus ici.
\end{enumerate}
\end{lemma}

\begin{proof}
Pour démontrer (1), on combine le lemme
\ref{varieties-lemma-geometrically-reduced-at-point} avec le lemme
d'Algèbre
\ref{algebra-lemma-geometrically-reduced-any-reduced-base-change}.
Pour démontrer (2), on combine le lemme
\ref{varieties-lemma-geometrically-reduced} avec le lemme d'Algèbre
\ref{algebra-lemma-geometrically-reduced-any-reduced-base-change}.
Pour démontrer (3), notons que $(X \times_k Y)_{\overline{k}} =
X_{\overline{k}} \times_{\overline{k}} Y_{\overline{k}}$ et utilisez
(2) ainsi que le lemme \ref{varieties-lemma-geometrically-reduced}. Pour démontrer
(4), utilisez (3) combiné avec le lemme \ref{varieties-lemma-perfect-reduced}.
\end{proof}

\begin{lemma}
\label{varieties-lemma-generic-points-geometrically-reduced}
Soit $k$ un corps.
Soit $X$ un schéma sur $k$.
\begin{enumerate}
\item Si $x' \leadsto x$ est une spécialisation et $X$
est géométriquement réduit en $x$, alors $X$ est
géométriquement réduit en $x'$. \item Si $x \in X$ tel
que (a) $\mathcal{O}_{X, x}$ est réduit, et (b) pour
chaque spécialisation $x' \leadsto x$ où $x'$ est un
point générique d'une composante irréductible de $X$, le
schéma $X$ est géométriquement réduit en $x'$, alors $X$
est géométriquement réduit en $x$. \item Si $X$ est
réduit et géométriquement réduit en tous les points
génériques des composantes irréductibles de $X$, alors $X$
est géométriquement réduit. \item Si $X$ est une variété sur
$k$, alors $X$ est géométriquement réduit si et seulement si
son corps des fonctions est une extension séparable de $k$.
\end{enumerate}
\end{lemma}

\begin{proof}
La partie (1) découle du lemme
\ref{varieties-lemma-geometrically-reduced-at-point} et du fait que
si $A$ est une $k$-algèbre géométriquement réduite, alors
$S^{-1}A$ est une $k$-algèbre géométriquement réduite pour
tout sous-ensemble multiplicatif $S$ de $A$, voir Algèbre,
lemme \ref{algebra-lemma-geometrically-reduced-permanence}.

\medskip\noindent
Soit $A = \mathcal{O}_{X, x}$. Les hypothèses (a) et (b) de (2)
impliquent que $A$ est réduit, et que $A_{\mathfrak q}$ est
géométriquement réduit sur $k$ pour chaque idéal premier minimal
$\mathfrak q$ de $A$. Par conséquent, $A$ est géométriquement
réduit sur $k$, voir Algèbre, lemme
\ref{algebra-lemma-generic-points-geometrically-reduced}. Ainsi, $X$ est
géométriquement réduit en $x$, voir lemme \ref{varieties-lemma-geometrically-reduced-at-point}.

\medskip\noindent
La partie (3) découle trivialement de la partie (2).
Enfin, (4) découle de (3) par Algèbre, lemme
\ref{algebra-lemma-characterize-separable-field-extensions}.
\end{proof}

\begin{lemma}
\label{varieties-lemma-Noetherian-geometrically-reduced-at-point}
Soit $k$ un corps. Supposons
que $X$ soit un schéma sur $k$. Soit $x \in
X$. Supposons que $X$ soit localement
noethérien et géométriquement réduit en $x$.
Alors il existe un voisinage ouvert $U \subset
X$ de $x$ qui est géométriquement réduit sur $k$.
\end{lemma}

\begin{proof}
Supposons que $X$ soit localement noethérien et
géométriquement réduit en $x$. D'après le lemme des
propriétés
\ref{properties-lemma-ring-affine-open-injective-local-ring}, nous
pouvons trouver un voisinage affine ouvert $U \subset X$ de $x$
tel que $R = \mathcal{O}_X(U) \to \mathcal{O}_{X, x}$ soit
injectif. D'après le lemme \ref{varieties-lemma-geometrically-reduced-at-point},
l'hypothèse signifie que $\mathcal{O}_{X, x}$ est géométriquement
réduit sur $k$. D'après Algèbre, lemme
\ref{algebra-lemma-subalgebra-separable}, cela implique que $R$ est
géométriquement réduit sur $k$, ce qui implique à son tour que $U$ est géométriquement réduit.
\end{proof}

\begin{example}
\label{varieties-example-not-geometrically-reduced}
Soit $k = \mathbf{F}_p(s, t)$, c'est-à-dire une
extension purement transcendantale du corps premier.
Considérons la variété $X = \Spec(k[x, y]/(1 + sx^p +
ty^p))$. Soit $k'/k$ une extension quelconque telle que $s$
et $t$ possèdent toutes deux une racine $p$-ième dans $k'$.
Alors le changement de base $X_{k'}$ n'est pas réduit. En
effet, l'anneau $k'[x, y]/(1 + s x^p + ty^p)$ contient
l'élément $1 + s^{1/p}x + t^{1/p}y$, dont la puissance $p$-ième est
nulle alors que lui-même ne l'est pas (l'idéal $(1 + sx^p + ty^p)$
ne contient certainement aucun élément non nul de degré $< p$).
\end{example}

\begin{lemma}
\label{varieties-lemma-finite-extension-geometrically-reduced}
Soit $k$ un corps. Soit $X \to \Spec(k)$ un morphisme
localement de type fini. Supposons que $X$ n'ait il'un nombre
fini de composantes irréductibles. Il existe alors une
extension finie purement inséparable $k'/k$ telle que
$(X_{k'})_{red}$ soit géométriquement réduit sur $k'$.
\end{lemma}

\begin{proof}
On peut remplacer $X$ par sa réduction $X_{red}$, et donc
supposer $X$ réduit et localement de type fini sur $k$. Soient
$x_1, \ldots, x_n \in X$ les points génériques de ses composantes
irréductibles. Pour toute extension algébrique purement
inséparable $k'/k$, le morphisme
$(X_{k'})_{red} \to X$ est un
homéomorphisme ; voir Algèbre, lemme
\ref{algebra-lemma-p-ring-map}. Les points $x'_1, \ldots, x'_n$
au-dessus de $x_1, \ldots, x_n$ sont donc les points génériques
des composantes irréductibles de $(X_{k'})_{red}$. Comme $X$ est
réduit, les anneaux locaux $K_i = \mathcal{O}_{X, x_i}$ sont des
corps ; voir Algèbre, lemme
\ref{algebra-lemma-minimal-prime-reduced-ring}. Puisque $X$ est
localement de type fini sur $k$, les extensions de corps $K_i/k$
sont de type fini. Enfin, les anneaux locaux
$\mathcal{O}_{(X_{k'})_{red}, x'_i}$ sont les corps
$(K_i \otimes_k k')_{red}$. D'après Algèbre, lemme
\ref{algebra-lemma-make-separable}, il existe une extension finie
purement inséparable $k'/k$ telle que les corps
$(K_i \otimes_k k')_{red}$ soient séparables sur $k'$. Chacun est
donc géométriquement réduit sur $k'$, d'après
Algèbre, lemme
\ref{algebra-lemma-characterize-separable-field-extensions}.
Le lemme \ref{varieties-lemma-generic-points-geometrically-reduced}, partie (3),
montre alors que $(X_{k'})_{red}$ est géométriquement réduit.
\end{proof}







\section{Schémas géométriquement connexes}
\label{varieties-section-geometrically-connected}

\noindent
Si $X$ est un schéma connexe sur un corps, il peut arriver
que $X$ devienne non connexe après l'extension du corps de
base. Cela n'arrive pas pour les schémas géométriquement connexes.

\begin{definition}
\label{varieties-definition-geometrically-connected}
Soit $X$ un schéma sur le corps $k$. On dit que
$X$ est {\it géométriquement connexe} sur $k$ si le schéma
$X_{k'}$ est connexe pour toute extension de corps $k'$ de $k$.
\end{definition}

\noindent
Par convention, un espace topologique connexe est non vide ; par
conséquent a fortiori, les schémas géométriquement connexes sont non
vides. Voici un exemple d'une variété qui n'est pas géométriquement connexe.

\begin{example}
\label{varieties-example-not-geometrically-irreducible}
Soit $k = \mathbf{Q}$. Le schéma $X =
\Spec(\mathbf{Q}(i))$ est une variété sur
$\Spec(\mathbf{Q})$. Mais le changement de base
$X_{\mathbf{C}}$ est le spectre de $\mathbf{C}
\otimes_{\mathbf{Q}} \mathbf{Q}(i) \cong \mathbf{C} \times \mathbf{C}$
qui est l'union disjointe de deux copies de $\Spec(\mathbf{C})$. Donc
en fait, c'est un exemple d'une variété non géométriquement connexe.
\end{example}

\begin{lemma}
\label{varieties-lemma-geometrically-connected-check-after-extension}
Supposons que $X$ soit un schéma sur le corps $k$.
Supposons que $k'/k$ soit une extension de corps.
Alors $X$ est géométriquement connexe sur $k$ si et
seulement si $X_{k'}$ est géométriquement connexe sur $k'$.
\end{lemma}

\begin{proof}
Si $X$ est géométriquement connexe sur $k$, alors il est clair
que $X_{k'}$ est géométriquement connexe sur $k'$. Pour la
réciproque, pour toute extension de corps $k''/k$, il existe
une extension de corps commune $k'''/k'$ et $k'''/k''$, voir
corps, lemme \ref{fields-lemma-common-extension-field}. Comme
le morphisme $X_{k'''} \to X_{k''}$ est surjectif (en tant que
changement de base d'un morphisme surjectif entre spectres de
corps), on voit que la connexité de $X_{k'''}$ implique la
connexité de $X_{k''}$. Ainsi, si $X_{k'}$ est géométriquement
connexe sur $k'$, alors $X$ est géométriquement connexe sur $k$.
\end{proof}

\begin{lemma}
\label{varieties-lemma-bijection-connected-components}
Soit $k$ un corps. Supposons
que $X$, $Y$ soient des schémas sur $k$.
Supposons que $X$ soit géométriquement
connexe sur $k$. Alors le morphisme de projection
$$
p : X \times_k Y \longrightarrow Y
$$
induit une bijection entre les composantes connexes.
\end{lemma}

\begin{proof}
Les fibres schématiques de $p$ sont connexes,
puisque ce sont des changements de base du
schéma géométriquement connexe $X$ par
extensions de corps. De plus, les fibres
schématiques sont homéomorphes aux fibres
ensemblistes, voir Schémas, lemme
\ref{schemes-lemma-fibre-topological}. Par morphismes,
lemme
\ref{morphisms-lemma-scheme-over-field-universally-open},
l'application $p$ est ouverte. Ainsi, nous pouvons appliquer Topologie,
lemme \ref{topology-lemma-connected-fibres-connected-components} pour conclure.
\end{proof}

\begin{lemma}
\label{varieties-lemma-affine-geometrically-connected}
Soit $k$ un corps. Supposons que $A$
soit une $k$-algèbre. Alors $X = \Spec(A)$ est
géométriquement connexe sur $k$ si et seulement si $A$
est géométriquement connexe sur $k$ (voir Algèbre,
définition \ref{algebra-definition-geometrically-connected} ).
\end{lemma}

\begin{proof}
Immédiat d'après les définitions.
\end{proof}

\begin{lemma}
\label{varieties-lemma-separably-closed-field-connected-components}
Supposons que $k'/k$ soit une extension de corps.
Soit $X$ un schéma sur $k$. Supposons
que $k$ soit séparablement algébriquement clos. Alors
le morphisme $X_{k'} \to X$ induit une bijection des
composantes connexes. En particulier, $X$ est
géométriquement connexe sur $k$ si et seulement si $X$ est connexe.
\end{lemma}

\begin{proof}
Puisque $k$ est séparablement algébriquement clos, nous
voyons que $k'$ est géométriquement connexe sur $k$, voir
Algèbre, lemme
\ref{algebra-lemma-separably-closed-connected-implies-geometric}.
Par conséquent, $Z = \Spec(k')$ est géométriquement connexe sur
$k$ d'après le lemme \ref{varieties-lemma-affine-geometrically-connected}
ci-dessus. Puisque $X_{k'} = Z \times_k X$, le résultat est un cas
particulier du lemme \ref{varieties-lemma-bijection-connected-components}.
\end{proof}

\begin{lemma}
\label{varieties-lemma-characterize-geometrically-connected}
Soit $k$ un corps. Supposons que $X$
soit un schéma sur $k$. Supposons que $\overline{k}$
soit une clôture algébrique séparable de $k$.
Alors $X$ est géométriquement connexe si et seulement
si le changement de base $X_{\overline{k}}$ est connexe.
\end{lemma}

\begin{proof}
Supposons que $X_{\overline{k}}$ soit connexe. Supposons
que $k'/k$ soit une extension de corps. Il existe une
extension de corps $\overline{k}'/\overline{k}$ telle que
$k'$ se plonge dans $\overline{k}'$ comme extension de
$k$. D'après le lemme
\ref{varieties-lemma-separably-closed-field-connected-components}, on voit que
$X_{\overline{k}'}$ est connexe. Puisque $X_{\overline{k}'} \to X_{k'}$
est surjectif, nous concluons que $X_{k'}$ est connexe comme voulu.
\end{proof}

\begin{lemma}
\label{varieties-lemma-descend-open}
Soit $k$ un corps. Supposons
que $X$ soit un schéma sur $k$. Supposons
que $A$ soit une $k$-algèbre. Supposons que
$V \subset X_A$ soit un ouvert
quasi-compact. Alors il existe une
$k$-subalgèbre de type fini $A' \subset A$ et un ouvert
quasi-compact $V' \subset X_{A'}$ tels que $V = V'_A$.
\end{lemma}

\begin{proof}
Nous remarquons que si $X$ est également quasi-séparé, cela
découle de Limites, lemme \ref{limits-lemma-descend-opens}.
Choisissons un nombre fini d'ouverts affines $U_1, \ldots, U_n$
de $X$ tels que $V \subset \bigcup U_{i,
A}$. Disons $U_i = \Spec(R_i)$. Puisque $V$ est quasi-compact,
nous pouvons trouver un nombre fini d'éléments $f_{ij} \in R_i
\otimes_k A$, $j = 1, \ldots, n_i$ tels que $V = \bigcup_i
\bigcup_{j = 1, \ldots, n_i} D(f_{ij})$ où $D(f_{ij}) \subset U_{i,
A}$ est l'ouvert principal correspondant. (Nous ne prétendons pas
que $V \cap U_{i, A}$ est l'union des $D(f_{ij})$, $j = 1, \ldots,
n_i$.) Il est clair que nous pouvons trouver une $k$-sous-algèbre
de type fini $A' \subset A$ telle que $f_{ij}$ soit l'image de
certains $f_{ij}' \in R_i \otimes_k A'$. Posons $V' = \bigcup
D(f_{ij}')$ qui est un ouvert quasi-compact de $X_{A'}$. Notons $\pi
: X_A \to X_{A'}$ le morphisme canonique. Nous avons $\pi(V) \subset
V'$ comme $\pi(D(f_{ij})) \subset D(f_{ij}')$. Si $x \in X_A$ avec
$\pi(x) \in V'$, alors $\pi(x) \in D(f_{ij}')$ pour certains $i, j$
et on voit que $x \in D(f_{ij})$ en tant que $f_{ij}'$ s'envoie
vers $f_{ij}$. Ainsi, on voit que $V = \pi^{-1}(V')$ comme voulu.
\end{proof}

\noindent
Soit $k$ un corps.
Supposons que $\overline{k}/k$ soit
une extension galoisienne
(éventuellement infinie). Par
exemple, $\overline{k}$ pourrait
être la clôture algébrique séparable
de $k$. Pour tout $\sigma \in
\text{Gal}(\overline{k}/k)$, nous
obtenons un automorphisme
correspondant $ \Spec(\sigma) :
\Spec(\overline{k}) \longrightarrow
\Spec(\overline{k}) $. Notons que
$\Spec(\sigma) \circ \Spec(\tau) = \Spec(\tau \circ
\sigma)$. Par conséquent, nous obtenons une action
$$
\text{Gal}(\overline{k}/k)^{opp} \times \Spec(\overline{k})
\longrightarrow
\Spec(\overline{k})
$$
du groupe opposé sur le schéma
$\Spec(\overline{k})$. Supposons que $X$ soit un
schéma sur $k$. Puisque $X_{\overline{k}} =
\Spec(\overline{k}) \times_{\Spec(k)} X$ par définition,
on voit que l'action ci-dessus induit une action canonique
\begin{equation}
\label{varieties-equation-galois-action-base-change-kbar}
\text{Gal}(\overline{k}/k)^{opp} \times X_{\overline{k}}
\longrightarrow
X_{\overline{k}}.
\end{equation}

\begin{lemma}
\label{varieties-lemma-Galois-action-quasi-compact-open}
Soit $k$ un corps. Supposons que $X$ soit un
schéma sur $k$. Supposons que $\overline{k}$ soit une
extension de Galois (éventuellement infinie) de $k$. Supposons que
$V \subset X_{\overline{k}}$ soit un ouvert quasi-compact. Alors
\begin{enumerate}
\item il existe une sous-extension finie
$\overline{k}/k'/k$ et un ouvert quasi-compact $V'
\subset X_{k'}$ tels que $V = (V')_{\overline{k}}$, \item il
existe un sous-groupe ouvert $H \subset
\text{Gal}(\overline{k}/k)$ tels que $\sigma(V) = V$ pour tout $\sigma \in H$.
\end{enumerate}
\end{lemma}

\begin{proof}
D'après le lemme \ref{varieties-lemma-descend-open}, il existe une sous-extension
finie $k'/k \subset \overline{k}$ et un ouvert $V' \subset X_{k'}$ dont
l'image réciproque est $V$. Cela prouve (1). Puisque $\text{Gal}(\overline{k}/k')$ est
ouvert dans $\text{Gal}(\overline{k}/k)$, la partie (2) est également claire.
\end{proof}

\begin{lemma}
\label{varieties-lemma-closed-fixed-by-Galois}
Soit $k$ un corps. Supposons que $\overline{k}/k$ soit une
extension de Galois (éventuellement infinie). Supposons que $X$ soit un schéma
sur $k$. Soit $\overline{T} \subset X_{\overline{k}}$ vérifiant les propriétés suivantes
\begin{enumerate}
\item $\overline{T}$ est un sous-ensemble fermé de
$X_{\overline{k}}$, \item pour chaque $\sigma \in
\text{Gal}(\overline{k}/k)$ nous avons $\sigma(\overline{T}) = \overline{T}$.
\end{enumerate}
Alors il existe un sous-ensemble fermé $T \subset X$ dont
l'image réciproque dans $X_{\overline{k}}$ est $\overline{T}$.
\end{lemma}

\begin{proof}
Ce lemme se réduit immédiatement au cas où $X = \Spec(A)$
est affine. Dans ce cas, soit $\overline{I} \subset A
\otimes_k \overline{k}$ l'idéal radical correspondant à
$\overline{T}$. L'hypothèse (2) implique que
$\sigma(\overline{I}) = \overline{I}$ pour tout $\sigma \in
\text{Gal}(\overline{k}/k)$. Choisissons $x \in \overline{I}$. Il
existe une extension galoisienne finie $k'/k$ contenue dans
$\overline{k}$ telle que $x \in A \otimes_k k'$. Posons $G = \text{Gal}(k'/k)$. Posons
$$
P(T) = \prod\nolimits_{\sigma \in G} (T - \sigma(x)) \in (A \otimes_k k')[T]
$$
Il est clair que $P(T)$ est unitaire et est en fait un élément
de $(A \otimes_k k')^G[T] = A[T]$ (selon la théorie de Galois de
base). De plus, si nous écrivons $P(T) = T^d + a_1T^{d - 1} +
\ldots + a_d$ alors on voit que $a_i \in I := A \cap
\overline{I}$. En combinant $P(x) = 0$ et $a_i \in I$ nous
trouvons $x^d = - a_1 x^{d - 1} - \ldots - a_d \in I(A \otimes_k
\overline{k})$. Ainsi, $x$ est contenu dans le radical de $I(A
\otimes_k \overline{k})$. Par conséquent, $\overline{I}$ est le radical
de $I(A \otimes_k \overline{k})$; le choix $T = V(I)$ convient.
\end{proof}

\begin{lemma}
\label{varieties-lemma-characterize-geometrically-disconnected}
Soit $k$ un corps.
Supposons que $X$ soit un schéma sur
$k$. Les conditions suivantes sont équivalentes
\begin{enumerate}
\item $X$ est géométriquement connexe,
\item pour toute extension de corps finie
séparable $k'/k$ le schéma $X_{k'}$ est connexe.
\end{enumerate}
\end{lemma}

\begin{proof}
Il découle immédiatement de la définition que (1)
implique (2). Supposons que $X$ ne soit pas
géométriquement connexe. Supposons que $k \subset
\overline{k}$ soit une clôture algébrique séparable de $k$.
D'après le lemme
\ref{varieties-lemma-characterize-geometrically-connected}, il s'ensuit que
$X_{\overline{k}}$ est non connexe. Disons $X_{\overline{k}} = \overline{U} \amalg
\overline{V}$ avec $\overline{U}$ et $\overline{V}$ ouverts, fermés et non vides.

\medskip\noindent
Supposons que $W \subset X$ soit un ouvert quasi-compact
quelconque. Alors $W_{\overline{k}} \cap \overline{U}$ et
$W_{\overline{k}} \cap \overline{V}$ sont ouverts et
fermés dans $W_{\overline{k}}$. En particulier,
$W_{\overline{k}} \cap \overline{U}$ et $W_{\overline{k}}
\cap \overline{V}$ sont quasi-compacts, et d'après le lemme
\ref{varieties-lemma-Galois-action-quasi-compact-open},
$W_{\overline{k}} \cap \overline{U}$ et $W_{\overline{k}} \cap
\overline{V}$ sont définis sur une sous-extension finie et
invariants sous un sous-groupe ouvert de
$\text{Gal}(\overline{k}/k)$. Nous utiliserons cela sans autre mention dans ce qui suit.

\medskip\noindent
Choisissons $W_0 \subset X$ quasi-compact ouvert de
telle sorte que $W_{0, \overline{k}} \cap
\overline{U}$ et $W_{0, \overline{k}} \cap
\overline{V}$ soient tous deux non vides. Choisissons
une sous-extension finie $\overline{k}/k'/k$ et une
décomposition $W_{0, k'} = U_0' \amalg V_0'$ en
parties ouvertes et fermées de sorte que $W_{0,
\overline{k}} \cap \overline{U} =
(U'_0)_{\overline{k}}$ et $W_{0, \overline{k}} \cap
\overline{V} = (V'_0)_{\overline{k}}$. Soit $H =
\text{Gal}(\overline{k}/k') \subset \text{Gal}(\overline{k}/k)$.
En particulier $\sigma(W_{0, \overline{k}} \cap \overline{U}) =
W_{0, \overline{k}} \cap \overline{U}$ et de même pour $\overline{V}$.

\medskip\noindent
Ayant choisi $W_0$, $k'$ comme ci-dessus, pour
chaque ouvert quasi-compact $W \subset X$ nous posons
$$
U_W =
\bigcap\nolimits_{\sigma \in H} \sigma(W_{\overline{k}} \cap \overline{U}),
\quad
V_W =
\bigcup\nolimits_{\sigma \in H} \sigma(W_{\overline{k}} \cap \overline{V}).
$$
Maintenant, puisque $W_{\overline{k}} \cap \overline{U}$ et
$W_{\overline{k}} \cap \overline{V}$ sont fixés par un sous-groupe ouvert
de $\text{Gal}(\overline{k}/k)$, on voit que l'union et
l'intersection ci-dessus sont finies. Par conséquent, $U_W$ et $V_W$ sont à la
fois ouverts et fermés. De plus, par construction $W_{\bar k} = U_W \amalg V_W$.

\medskip\noindent
Nous affirmons que si $W \subset W' \subset X$ sont
quasi-compacts ouverts, alors $W_{\overline{k}} \cap U_{W'}
= U_W$ et $W_{\overline{k}} \cap V_{W'} = V_W$. Vérification
omise. Ainsi, nous voyons qu'en définissant $U = \bigcup_{W
\subset X} U_W$ et $V = \bigcup_{W \subset X} V_W$ nous
obtenons $X_{\overline{k}} = U \amalg V$ est une union
disjointe de ouverts et fermés. Il est évident que $V$ est
non vide car il est construit en prenant des unions
(localement). D'autre part, $U$ est non vide puisqu'il contient
$W_0 \cap \overline{U}$ par construction. Enfin, $U, V \subset
X_{\bar k}$ sont fermés et $H$-invariants par construction.
Ainsi, par le lemme \ref{varieties-lemma-closed-fixed-by-Galois} nous
avons $U = (U')_{\bar k}$, et $V = (V')_{\bar k}$ pour un certain
$U', V' \subset X_{k'}$ fermé. Il est clair que $X_{k'} = U'
\amalg V'$ et on voit que $X_{k'}$ est non connexe comme voulu.
\end{proof}

\begin{lemma}
\label{varieties-lemma-tricky}
Soit $k$ un corps. Supposons que $\overline{k}/k$ soit une
extension de Galois (éventuellement infinie). Supposons que $f : T \to X$
soit un morphisme de schémas sur $k$. Supposons que $T_{\overline{k}}$
soit connexe et $X_{\overline{k}}$ soit non connexe. Alors $X$ est non connexe.
\end{lemma}

\begin{proof}
Écrivons $X_{\overline{k}} = \overline{U} \amalg \overline{V}$
avec $\overline{U}$ et $\overline{V}$ ouverts et fermés. Notons
$\overline{f} : T_{\overline{k}} \to X_{\overline{k}}$ le
changement de base de $f$. Puisque $T_{\overline{k}}$ est connexe,
on voit que $T_{\overline{k}}$ est contenu soit dans
$\overline{f}^{-1}(\overline{U})$, soit dans
$\overline{f}^{-1}(\overline{V})$. Disons que $T_{\overline{k}} \subset \overline{f}^{-1}(\overline{U})$.

\medskip\noindent
Fixons un ouvert quasi-compact $W \subset X$. Il
existe une sous-extension galoisienne finie
$\overline{k}/k'/k$ telle que $\overline{U} \cap
W_{\overline{k}}$ et $\overline{V} \cap
W_{\overline{k}}$ proviennent d'ouverts quasi-compacts $U', V'
\subset W_{k'}$. Alors aussi $W_{k'} = U' \amalg V'$. Considérons
$$
U'' = \bigcap\nolimits_{\sigma \in \text{Gal}(k'/k)} \sigma(U'),
\quad
V'' = \bigcup\nolimits_{\sigma \in \text{Gal}(k'/k)} \sigma(V').
$$
Ces parties sont invariantes par Galois, ouvertes et
fermés, et $W_{k'} = U'' \amalg V''$.
D'après le lemme
\ref{varieties-lemma-closed-fixed-by-Galois}, nous obtenons des
ouverts et fermés $U_W, V_W \subset W$ tels que $U'' =
(U_W)_{k'}$, $V'' = (V_W)_{k'}$ et $W = U_W \amalg V_W$.

\medskip\noindent
Nous affirmons que si $W \subset W' \subset X$ sont
quasi-compacts ouverts, alors $W \cap U_{W'} = U_W$ et $W
\cap V_{W'} = V_W$. Nous omettons la vérification. Par conséquent,
nous voyons qu'en définissant $U = \bigcup_{W \subset X}
U_W$ et $V = \bigcup_{W \subset X} V_W$, nous obtenons $X =
U \amalg V$. Il est clair que $V$ est non vide car il est
construit en prenant des unions (localement). D'autre part,
$U$ est non vide puisqu'il contient $f(T)$ par construction.
\end{proof}

\begin{lemma}
\label{varieties-lemma-geometrically-connected-criterion}
\begin{reference}
\cite[IV Corollary 4.5.13.1(i)]{EGA}
\end{reference}
Soit $k$ un corps. Soit $T \to X$ un
morphisme de schémas sur $k$. Supposons $T$
géométriquement connexe et $X$ connexe. Alors $X$ est géométriquement connexe.
\end{lemma}

\begin{proof}
Ceci est une reformulation
du lemme \ref{varieties-lemma-tricky}.
\end{proof}

\begin{lemma}
\label{varieties-lemma-geometrically-connected-if-connected-and-point}
Soit $k$ un corps. Supposons que $X$ soit
un schéma sur $k$. Supposons que $X$ soit connexe et
possède un point $x$ tel que $k$ soit algébriquement clos
dans $\kappa(x)$. Alors $X$ est géométriquement connexe.
En particulier, si $X$ possède un point $k$-rationnel et
que $X$ est connexe, alors $X$ est géométriquement connexe.
\end{lemma}

\begin{proof}
Posons $T = \Spec(\kappa(x))$. Supposons que
$\overline{k}$ soit une clôture algébrique séparable
de $k$. L'hypothèse sur $\kappa(x)/k$ implique que
$T_{\overline{k}}$ est irréductible, voir Algèbre, lemme
\ref{algebra-lemma-field-extension-geometrically-irreducible}.
Donc, d'après le lemme
\ref{varieties-lemma-geometrically-connected-criterion}, on voit que
$X_{\overline{k}}$ est connexe. D'après le lemme
\ref{varieties-lemma-characterize-geometrically-connected}, nous concluons que $X$ est géométriquement connexe.
\end{proof}

\begin{lemma}
\label{varieties-lemma-inverse-image-connected-component}
Soit $K/k$ une extension de corps.
Soit $X$ un schéma sur $k$. Pour chaque
composante connexe $T$ de $X$, l'image réciproque $T_K
\subset X_K$ est une union de composantes connexes de $X_K$.
\end{lemma}

\begin{proof}
C'est un énoncé purement topologique. Notons $p
: X_K \to X$ le morphisme de projection. Soit $T \subset
X$ une composante connexe de $X$. Soit $t \in T_K =
p^{-1}(T)$. Soit $C \subset X_K$ une composante connexe
contenant $t$. Alors $p(C)$ est un sous-ensemble connexe de $X$ qui
rencontre $T$, donc $p(C) \subset T$. Par conséquent $C \subset T_K$.
\end{proof}

\noindent
Le lemme suivant sera supplanté par le lemme plus
fort \ref{varieties-lemma-image-connected-component} ci-dessous.

\begin{lemma}
\label{varieties-lemma-image-connected-component-finite-extension}
Supposons que $K/k$ soit une extension finie de corps et
soit $X$ un schéma sur $k$. Notons $p : X_K \to X$ le
morphisme de projection. Pour chaque composante connexe $T$
de $X_K$, l'image $p(T)$ est une composante connexe de $X$.
\end{lemma}

\begin{proof}
L'image $p(T)$ est contenue dans une certaine composante
connexe $X'$ de $X$. Considérons $X'$ comme un
sous-schéma fermé de $X$ de quelque manière que ce soit.
Alors $T$ est également une composante connexe de $X'_K =
p^{-1}(X')$ et nous pouvons donc supposer que $X$ est
connexe. Le morphisme $p$ est ouvert (Morphismes, lemme
\ref{morphisms-lemma-scheme-over-field-universally-open}),
fermé (Morphismes, lemme
\ref{morphisms-lemma-integral-universally-closed}) et les fibres de
$p$ sont des ensembles finis (Morphismes, lemme
\ref{morphisms-lemma-finite-quasi-finite}). Ainsi, nous pouvons appliquer
Topologie, lemme \ref{topology-lemma-finite-fibre-connected-components} pour conclure.
\end{proof}

\begin{lemma}[Gabber]
\label{varieties-lemma-image-connected-component}
\begin{reference}
Email de Ofer Gabber daté du 4 juin 2016
\end{reference}
Soit $K/k$ une extension de corps. Supposons que $X$
soit un schéma sur $k$. Notons $p : X_K \to X$ le morphisme de
projection. Supposons que $\overline{T} \subset X_K$ soit une
composante connexe. Alors $p(\overline{T})$ est une composante connexe de $X$.
\end{lemma}

\begin{proof}
Lorsque $K/k$ est fini, c'est le lemme
\ref{varieties-lemma-image-connected-component-finite-extension}.
En général, la preuve est plus difficile.

\medskip\noindent
Soit $T \subset X$ soit la composante connexe de $X$
contenant l'image de $\overline{T}$. Nous pouvons remplacer $X$ par
$T$ (avec la structure de sous-schéma réduit induite). Ainsi, nous
pouvons supposer que $X$ est connexe. Supposons que $A = H^0(X,
\mathcal{O}_X)$. Soit $L \subset A$ la $k$-algèbre maximale faiblement
étale, voir Compléments d'algèbre, lemme
\ref{more-algebra-lemma-max-weakly-etale-subalgebra}. Comme $A$ n'a pas
d'idempotents non triviaux, on voit que $L$ est un corps et une
extension algébrique séparable de $k$ d'après Compléments d'algèbre, lemme
\ref{more-algebra-lemma-class-weakly-etale-over-field}. Observons que $L$
est également l'algèbre $L$ maximale faiblement étale de $A$ (car toute
algèbre faiblement étale $L$ est faiblement étale sur $k$ d'après
Compléments d'algèbre, lemme \ref{more-algebra-lemma-composition-weakly-etale}).
D'après Schémas, lemme \ref{schemes-lemma-morphism-into-affine}, nous
obtenons une factorisation $X \to \Spec(L) \to \Spec(k)$ du morphisme de structure.

\medskip\noindent
Supposons que $L'/L$ soit une extension
finie séparable. D'après Cohomologie des
Schémas, lemme
\ref{coherent-lemma-finite-locally-free-base-change-cohomology} nous avons
$$
A \otimes_L L' =
H^0(X \times_{\Spec(L)} \Spec(L'), \mathcal{O}_{X \times_{\Spec(L)} \Spec(L')})
$$
La $L'$-sous-algèbre faiblement étale maximale de $A \otimes_L
L'$ est $L \otimes_L L' = L'$, d'après Compléments d'algèbre, lemme
\ref{more-algebra-lemma-change-fields-max-weakly-etale-subalgebra}.
En particulier, $A \otimes_L L'$ n'a pas d'idempotents non
triviaux (un tel idempotent engendrerait une sous-algèbre faiblement
étale) et nous concluons que $X \times_{\Spec(L)} \Spec(L')$ est
connexe. Par le lemme
\ref{varieties-lemma-characterize-geometrically-disconnected}, nous concluons que $X$ est géométriquement connexe sur $L$.

\medskip\noindent
Donnons à $\overline{T}$ la structure de
schéma induit réduite et considérons la composition
$$
\overline{T} \xrightarrow{i} X_K = X \times_{\Spec(k)} \Spec(K)
\xrightarrow{\pi}
\Spec(L \otimes_k K)
$$
L'image est contenue dans une composante connexe de $\Spec(L
\otimes_k K)$. Puisque $K \to L \otimes_k K$ est intégral,
on voit que les composantes connexes de $\Spec(L \otimes_k
K)$ sont des points et tous les points sont fermés, voir
Algèbre, lemme \ref{algebra-lemma-integral-over-field}. Ainsi,
nous obtenons un corps quotient $L \otimes_k K \to E$ tel que
$\overline{T}$ s'envoie dans $\Spec(E) \subset \Spec(L \otimes_k
K)$. Par conséquent $i(\overline{T}) \subset \pi^{-1}(\Spec(E))$. Mais
$$
\pi^{-1}(\Spec(E)) =
(X \times_{\Spec(k)} \Spec(K)) \times_{\Spec(L \otimes_k K)} \Spec(E) =
X \times_{\Spec(L)} \Spec(E)
$$
qui est connexe parce que $X$ est géométriquement connexe sur
$L$. Ensuite, nous obtenons l'égalité $\overline{T} = X
\times_{\Spec(L)} \Spec(E)$ (au sens des ensembles) et nous
concluons que $\overline{T} \to X$ est surjective comme voulu.
\end{proof}

\noindent
Soit $X$ un schéma. On note
$\pi_0(X)$ l'ensemble des composantes connexes de $X$.

\begin{lemma}
\label{varieties-lemma-galois-action-connected-components}
Soit $k$ un corps, muni d'une clôture
algébrique séparable $\overline{k}$. Supposons
que $X$ soit un schéma sur $k$. Il existe une action
$$
\text{Gal}(\overline{k}/k)^{opp} \times \pi_0(X_{\overline{k}})
\longrightarrow
\pi_0(X_{\overline{k}})
$$
avec les propriétés suivantes :
\begin{enumerate}
\item Un élément $\overline{T} \in
\pi_0(X_{\overline{k}})$ est fixé par l'action si et
seulement s'il existe une composante connexe $T \subset X$, qui
est géométriquement connexe sur $k$, telle que
$T_{\overline{k}} = \overline{T}$. \item Pour toute extension de corps
$k'/k$ avec clôture algébrique séparable $\overline{k}'$ le diagramme
$$
\xymatrix{
\text{Gal}(\overline{k}'/k') \times \pi_0(X_{\overline{k}'})
\ar[r] \ar[d] &
\pi_0(X_{\overline{k}'}) \ar[d] \\
\text{Gal}(\overline{k}/k) \times \pi_0(X_{\overline{k}})
\ar[r] &
\pi_0(X_{\overline{k}})
}
$$
est commutatif (où la flèche verticale droite est une bijection
d'après le lemme \ref{varieties-lemma-separably-closed-field-connected-components}).
\end{enumerate}
\end{lemma}

\begin{proof}
L'action (\ref{varieties-equation-galois-action-base-change-kbar})
de $\text{Gal}(\overline{k}/k)$ sur $X_{\overline{k}}$
induit une action sur ses composantes connexes. Les
composantes connexes sont toujours fermées (Topologie,
lemme \ref{topology-lemma-connected-components}). Donc si
$\overline{T}$ est comme dans (1), alors par le lemme
\ref{varieties-lemma-closed-fixed-by-Galois} il existe un sous-ensemble
fermé $T \subset X$ tel que $\overline{T} =
T_{\overline{k}}$. Notons que $T$ est géométriquement connexe
sur $k$, voir le lemme
\ref{varieties-lemma-characterize-geometrically-connected}. Pour voir que $T$ est
une composante connexe de $X$, supposons que $T \subset T'$, $T \not =
T'$ où $T'$ est une composante connexe de $X$. Dans ce cas, $T'_{k'}$
contient strictement $\overline{T}$ et est donc non connexe. Par le lemme
\ref{varieties-lemma-tricky} cela signifie que $T'$ est non connexe ! Contradiction.

\medskip\noindent
Nous omettons la preuve de la fonctorialité en (2).
\end{proof}

\begin{lemma}
\label{varieties-lemma-galois-action-connected-components-continuous}
Soit $k$ un corps, muni d'une clôture
algébrique séparable $\overline{k}$.
Soit $X$ un schéma sur $k$. Supposons
\begin{enumerate}
\item $X$ est quasi-compact, et \item les
composantes connexes de $X_{\overline{k}}$ sont ouvertes.
\end{enumerate}
Alors
\begin{enumerate}
\item [(a)] $\pi_0(X_{\overline{k}})$ est fini,
et \item [(b)] l'action de
$\text{Gal}(\overline{k}/k)$ sur $\pi_0(X_{\overline{k}})$ est continue.
\end{enumerate}
De plus, les hypothèses (1) et (2) sont
satisfaites lorsque $X$ est de type fini sur $k$.
\end{lemma}

\begin{proof}
Comme les composantes connexes, qui sont ouvertes, recouvrent
$X_{\overline{k}}$ (Topologie, lemme
\ref{topology-lemma-connected-components}) et
$X_{\overline{k}}$ est quasi-compact, nous concluons qu'il
n'y en a qu'un nombre fini. Ainsi, (a) est vérifié. D'après
le lemme \ref{varieties-lemma-descend-open}, ces composantes connexes
sont chacune définies sur une sous-extension finie de
$\overline{k}/k$ et nous obtenons (b). Si $X$ est de type
fini sur $k$, alors $X_{\overline{k}}$ est de type fini sur
$\overline{k}$ (Morphismes, lemme
\ref{morphisms-lemma-base-change-finite-type}). Par conséquent,
$X_{\overline{k}}$ est un schéma noethérien (Morphismes, lemme
\ref{morphisms-lemma-finite-type-noetherian}). Ainsi,
$X_{\overline{k}}$ a un nombre fini de composantes irréductibles
(Propriétés, lemme \ref{properties-lemma-Noetherian-irreducible-components})
et a fortiori un nombre fini de composantes connexes (qui sont donc ouvertes).
\end{proof}









\section{Schémas géométriquement irréductibles}
\label{varieties-section-geometrically-irreducible}

\noindent
Si $X$ est un schéma irréductible sur un corps, il peut arriver
que $X$ devienne réductible après l'extension du corps de base.
Cela n'arrive pas pour les schémas géométriquement irréductibles.

\begin{definition}
\label{varieties-definition-geometrically-irreducible}
Soit $X$ un schéma sur le corps $k$. On dit que
$X$ est {\it géométriquement irréductible} sur $k$ si le
schéma $X_{k'}$ est irréductible \footnote { Un espace
irréductible est non vide. } pour toute extension de corps $k'$ de $k$.
\end{definition}

\begin{lemma}
\label{varieties-lemma-geometrically-irreducible-check-after-extension}
Supposons que $X$ soit un schéma sur le corps $k$.
Supposons que $k'/k$ soit une extension de corps. Alors
$X$ est géométriquement irréductible sur $k$ si et
seulement si $X_{k'}$ est géométriquement irréductible sur $k'$.
\end{lemma}

\begin{proof}
Si $X$ est géométriquement irréductible sur $k$, alors il est
clair que $X_{k'}$ est géométriquement irréductible sur $k'$.
Pour la réciproque, pour toute extension de corps $k''/k$, il
existe une extension de corps commune $k'''/k'$ et $k'''/k''$,
voir corps, lemme \ref{fields-lemma-common-extension-field}. Comme
le morphisme $X_{k'''} \to X_{k''}$ est surjectif (en tant que
changement de base d'un morphisme surjectif entre spectres de
corps), on voit que l'irréductibilité de $X_{k'''}$ implique
l'irréductibilité de $X_{k''}$. Ainsi, si $X_{k'}$ est géométriquement
irréductible sur $k'$, alors $X$ est géométriquement irréductible sur $k$.
\end{proof}

\begin{lemma}
\label{varieties-lemma-separably-closed-irreducible}
Supposons que $X$ soit un schéma sur un corps
séparablement clos $k$. Si $X$ est irréductible, alors
$X_K$ est irréductible pour toute extension de corps
$K/k$. C.-à-d., $X$ est géométriquement irréductible sur $k$.
\end{lemma}

\begin{proof}
Cela résulte de Propriétés, lemme
\ref{properties-lemma-characterize-irreducible} et Algèbre, lemme \ref{algebra-lemma-separably-closed-irreducible}.
\end{proof}

\begin{lemma}
\label{varieties-lemma-bijection-irreducible-components}
Soit $k$ un corps. Supposons
que $X$, $Y$ soient des schémas sur $k$.
Supposons que $X$ soit géométriquement
irréductible sur $k$. Ensuite, le morphisme de projection
$$
p : X \times_k Y \longrightarrow Y
$$
induit une bijection entre les composantes irréductibles.
\end{lemma}

\begin{proof}
Tout d'abord, notons que les fibres schématiques de $p$
sont irréductibles, puisqu'elles sont des changements de
base du schéma géométriquement irréductible $X$ par des
extensions de corps. De plus, les fibres schématiques
sont homéomorphes aux fibres ensemblistes, voir Schémas,
lemme \ref{schemes-lemma-fibre-topological}. Selon
Morphismes, lemme
\ref{morphisms-lemma-scheme-over-field-universally-open}, l'application
$p$ est ouverte. Ainsi, nous pouvons appliquer Topologie, lemme
\ref{topology-lemma-irreducible-fibres-irreducible-components} pour conclure.
\end{proof}

\begin{lemma}
\label{varieties-lemma-geometrically-irreducible-local}
\begin{slogan}
L'irréductibilité géométrique est Zariski locale modulo la connexité.
\end{slogan}
Soit $k$ un corps. Supposons que $X$ soit
un schéma sur $k$. Les conditions suivantes sont équivalentes
\begin{enumerate}
\item $X$ est géométriquement irréductible sur $k$, \item
pour chaque ouvert affine non vide $U$ la $k$-algèbre
$\mathcal{O}_X(U)$ est géométriquement irréductible sur $k$
(voir Algèbre, définition
\ref{algebra-definition-geometrically-irreducible}), \item $X$
est irréductible et il existe un recouvrement ouvert affine $X =
\bigcup U_i$ tel que chaque $k$-algèbre $\mathcal{O}_X(U_i)$ est
géométriquement irréductible, et \item il existe un recouvrement
ouvert $X = \bigcup_{i \in I} X_i$ avec $I \not = \emptyset$
tel que $X_i$ est géométriquement irréductible pour chaque $i$ et
tel que $X_i \cap X_j \not = \emptyset$ pour tout $i, j \in I$.
\end{enumerate}
De plus, si $X$ est géométriquement irréductible, il en
est de même pour tout sous-schéma ouvert non vide de $X$.
\end{lemma}

\begin{proof}
Un schéma affine $\Spec(A)$ sur $k$ est géométriquement
irréductible si et seulement si $A$ est géométriquement
irréductible sur $k$ ; cela découle immédiatement des
définitions. Rappelons que si un schéma est irréductible, alors
tout sous-schéma ouvert non vide de $X$ l'est aussi, et que deux
ouverts non vides ont une intersection non vide. De plus, si tout
ouvert affine est irréductible, alors le schéma est irréductible,
voir Propriétés, lemme
\ref{properties-lemma-characterize-irreducible}. Ainsi, l'énoncé final du
lemme est clair, ainsi que les implications (1) $\Rightarrow$ (2), (2)
$\Rightarrow$ (3), et (3) $\Rightarrow$ (4). Si (4) est vrai, alors pour
toute extension de corps $k'/k$, le schéma $X_{k'}$ possède un recouvrement
par des ouverts irréductibles dont les intersections deux à deux sont non
vides. Par conséquent, $X_{k'}$ est irréductible. Ainsi, (4) implique (1).
\end{proof}

\begin{lemma}
\label{varieties-lemma-geometrically-irreducible-function-field}
Supposons que $X$ soit un schéma irréductible sur le corps $k$. Supposons que
$\xi \in X$ soit son point générique. Les énoncés suivants sont équivalents :
\begin{enumerate}
\item $X$ est géométriquement irréductible sur $k$ ; \item
$\kappa(\xi)$ est géométriquement irréductible sur $k$ ; \item la
clôture algébrique séparable de $k$ dans $\kappa(\xi)$ est égale à $k$.
\end{enumerate}
\end{lemma}

\begin{proof}
D'après Algèbre, lemme
\ref{algebra-lemma-geometrically-irreducible-separable-elements},
(2) et (3) sont équivalents.

\medskip\noindent
Supposons (1). Rappelons que $\mathcal{O}_{X, \xi}$ est
la colimite filtrante de $\mathcal{O}_X(U)$ où $U$
parcourt les sous-schémas affines ouverts non vides de
$X$. En combinant le lemme
\ref{varieties-lemma-geometrically-irreducible-local} et le lemme
d'Algèbre
\ref{algebra-lemma-subalgebra-geometrically-irreducible}, nous voyons
que $\mathcal{O}_{X, \xi}$ est géométriquement irréductible sur $k$.
Puisque $\mathcal{O}_{X, \xi} \to \kappa(\xi)$ est une surjection avec
noyau localement nilpotent (voir le lemme d'Algèbre
\ref{algebra-lemma-minimal-prime-reduced-ring}), il s'ensuit que $\kappa(\xi)$ est
géométriquement irréductible, voir le lemme d'Algèbre \ref{algebra-lemma-p-ring-map}.

\medskip\noindent
Supposons (2). On peut supposer que $X$ est réduit. Supposons que
$U \subset X$ est un ouvert affine non vide. Alors $U = \Spec(A)$ où $A$
est un anneau intègre avec corps des fractions $\kappa(\xi)$. Ainsi $A$ est
une $k$-sous-algèbre d'une $k$-algèbre géométriquement irréductible. Par
conséquent, d'après le lemme d'Algèbre
\ref{algebra-lemma-subalgebra-geometrically-irreducible}, on voit que $A$ est
géométriquement irréductible sur $k$. D'après le lemme
\ref{varieties-lemma-geometrically-irreducible-local}, nous en concluons que $X$ est géométriquement irréductible sur $k$.
\end{proof}

\begin{lemma}
\label{varieties-lemma-geometrically-irreducible-self-product}
Un schéma $X$ sur un corps $k$ est géométriquement
irréductible sur $k$ si et seulement si $X\times_k X$ est irréductible.
\end{lemma}

\begin{proof}
Si $X$ est géométriquement irréductible sur $k$,
alors $X\times_k X$ est irréductible selon le
lemme \ref{varieties-lemma-bijection-irreducible-components}.

\medskip\noindent
Inversement, supposons que $X \times_k X$ soit
irréductible. En remplaçant $X$ par le schéma réduit
sous-jacent, on peut supposer que $X$ est réduit sans
changer le problème. Puisque la première projection $X
\times_k X \to X$ est surjective, $X$ est également
irréductible, donc intégral. Supposons qu'il existe un élément
$\alpha$ dans le corps des fonctions $k(X)$ qui soit
algébrique séparable sur $k$ mais pas dans $k$. Soit $E = k(\alpha)$.

\medskip\noindent
Puis il existe un sous-schéma affine ouvert non vide $U$ de $X$
tel que $k \subset E \subset O(U)$, et donc $E \otimes_k E
\subset O(U)\otimes_k O(U) = O(U\times_k U)$. Ainsi nous avons
un morphisme dominant $U \times_k U \to \Spec(E\otimes_k E)$.
Puisque $X\times_k X$ est irréductible, $U \times_k U$ l'est
aussi, et donc $\Spec(E \otimes_k E)$ est connexe. Si $E$ est
strictement plus grand que $k$, alors $E \otimes_k E$ est un
produit de corps incluant $E$ comme un facteur. Il a une
dimension $(\dim_k(E))^2>\dim_k(E)$ comme espace vectoriel sur
$k$, et donc $E\otimes_k E$ est un produit d'au moins deux corps,
ce qui contredit que $\Spec(E \otimes_k E)$ est connexe. Donc en
fait $E$ est égal à $k$. D'après le lemme
\ref{varieties-lemma-geometrically-irreducible-function-field}, $X$ est géométriquement irréductible sur $k$.
\end{proof}

\begin{lemma}
\label{varieties-lemma-separably-closed-field-irreducible-components}
Supposons que $k'/k$ soit une extension de corps. Supposons que
$X$ est un schéma sur $k$. Posons $X' = X_{k'}$. Supposons que
$k$ est séparablement algébriquement clos. Alors le morphisme
$X' \to X$ induit une bijection des composantes irréductibles.
\end{lemma}

\begin{proof}
Puisque $k$ est séparablement algébriquement clos, on voit que
$k'$ est géométriquement irréductible sur $k$, voir Algèbre, lemme
\ref{algebra-lemma-separably-closed-irreducible-implies-geometric}.
Par conséquent, $Z = \Spec(k')$ est géométriquement irréductible
sur $k$, d'après le lemme
\ref{varieties-lemma-geometrically-irreducible-local} ci-dessus. Puisque $X' = Z \times_k X$, le
résultat est un cas particulier du lemme \ref{varieties-lemma-bijection-irreducible-components}.
\end{proof}

\begin{lemma}
\label{varieties-lemma-characterize-geometrically-irreducible}
\begin{slogan}
L'irréductibilité géométrique peut être testée sur
une clôture algébrique séparable du corps de base.
\end{slogan}
Soit $k$ un corps. Supposons que $X$
soit un schéma sur $k$. Les conditions suivantes sont équivalentes :
\begin{enumerate}
\item $X$ est géométriquement irréductible sur $k$,
\item pour toute extension de corps finie séparable
$k'/k$ le schéma $X_{k'}$ est irréductible, et \item
$X_{\overline{k}}$ est irréductible, où $k \subset
\overline{k}$ est une clôture algébrique séparable de $k$.
\end{enumerate}
\end{lemma}

\begin{proof}
Supposons que $X_{\overline{k}}$ soit irréductible, c'est-à-dire,
supposons (3). Supposons que $k'/k$ soit une extension de corps.
Il existe une extension de corps $\overline{k}'/\overline{k}$
telle que $k'$ se plonge dans $\overline{k}'$ comme une extension
de $k$. D'après le lemme
\ref{varieties-lemma-separably-closed-field-irreducible-components}, on voit que
$X_{\overline{k}'}$ est irréductible. Puisque $X_{\overline{k}'} \to X_{k'}$ est
surjective, nous concluons que $X_{k'}$ est irréductible. Par conséquent, (1) est vrai.

\medskip\noindent
Supposons que $k \subset \overline{k}$ soit une clôture
algébrique séparable de $k$. Supposons que non (3),
c'est-à-dire supposons que $X_{\overline{k}}$ soit
réductible. Notre but est de montrer que $X_{k'}$ est
également réductible pour une certaine sous-extension
finie $\overline{k}/k'/k$. Supposons que $X = \bigcup_{i
\in I} U_i$ soit un recouvrement affine ouvert avec $U_i$
non vide. Si pour un certain $i$ le schéma $U_i$ est
réductible, ou si pour une certaine paire $i \not = j$
l'intersection $U_i \cap U_j$ est vide, alors $X$ est
réductible (Propriétés, lemme
\ref{properties-lemma-characterize-irreducible}) et nous
avons terminé. En particulier, on peut supposer que
$U_{i, \overline{k}} \cap U_{j, \overline{k}}$ pour tout $i, j
\in I$ est non vide et nous en concluons que $U_{i,
\overline{k}}$ doit être réductible pour un certain $i$. Selon
Algèbre, lemme \ref{algebra-lemma-geometrically-irreducible},
cela signifie que $U_{i, k'}$ est réductible pour une certaine
extension de corps finie séparable $k'/k$. Par conséquent, $X_{k'}$
est également réductible. Ainsi, on voit que (2) implique (3).

\medskip\noindent
L'implication (1) $\Rightarrow$ (2)
est immédiate. Cela prouve le lemme.
\end{proof}

\begin{lemma}
\label{varieties-lemma-inverse-image-irreducible}
Soit $K/k$ une extension de corps.
Soit $X$ un schéma sur $k$. Pour chaque
composante irréductible $T$ de $X$, l'image réciproque $T_K
\subset X_K$ est une union de composantes irréductibles de $X_K$.
\end{lemma}

\begin{proof}
Soit $T \subset X$ une composante irréductible de
$X$. Le morphisme $T_K \to T$ est plat, donc les généralisations
se relèvent le long de $T_K \to T$. Ainsi, chaque $\xi \in T_K$
qui est un point générique d'une composante irréductible de
$T_K$ s'envoie sur le point générique $\eta$ de $T$. Si $\xi'
\leadsto \xi$ est une spécialisation dans $X_K$, alors $\xi'$ se
mappe à $\eta$ puisqu'il n'y a pas de points se spécialisant en
$\eta$ dans $X$. Ainsi $\xi' \in T_K$ et nous concluons que $\xi =
\xi'$. En d'autres termes, $\xi$ est le point générique d'une
composante irréductible de $X_K$. Cela signifie que les composantes
irréductibles de $T_K$ sont toutes des composantes irréductibles de $X_K$.
\end{proof}

\noindent
Pour un schéma $X$, nous notons
$\text{IrredComp}(X)$ l'ensemble des composantes irréductibles de $X$.

\begin{lemma}
\label{varieties-lemma-image-irreducible}
Soit $K/k$ une extension de corps.
Soit $X$ un schéma sur $k$. Pour chaque
composante irréductible $\overline{T} \subset X_K$,
l'image de $\overline{T}$ dans $X$ est une composante
irréductible dans $X$. Cela définit une application canonique
$$
\text{IrredComp}(X_K)
\longrightarrow
\text{IrredComp}(X)
$$
qui est surjective.
\end{lemma}

\begin{proof}
Considérons le diagramme
$$
\xymatrix{
X_K \ar[d] & X_{\overline{K}} \ar[d] \ar[l] \\
X & X_{\overline{k}} \ar[l]
}
$$
où $\overline{K}$ est la clôture algébrique séparable de $K$,
et où $\overline{k}$ est la clôture algébrique séparable de
$k$. D'après le lemme
\ref{varieties-lemma-separably-closed-field-irreducible-components}, le
morphisme $X_{\overline{K}} \to X_{\overline{k}}$ induit une bijection
entre les composantes irréductibles. Il suffit donc de démontrer le
lemme pour les morphismes $X_{\overline{k}} \to X$ et $X_{\overline{K}}
\to X_K$. En d'autres termes, on peut supposer que $K = \overline{k}$.

\medskip\noindent
Le morphisme $p : X_{\overline{k}} \to X$ est intégral, plat et
surjectif. La platitude implique que les généralisations se relèvent
le long de $p$, voir Morphismes, lemme
\ref{morphisms-lemma-generalizations-lift-flat}. Par conséquent, les
points génériques des composantes irréductibles de $X_{\overline{k}}$ se
projettent sur des points génériques des composantes irréductibles de
$X$. L'intégralité implique que $p$ est universellement fermé, voir
Morphismes, lemme \ref{morphisms-lemma-integral-universally-closed}.
Ainsi, nous concluons que l'image $p(\overline{T})$ d'une composante
irréductible est un sous-ensemble fermé irréductible qui contient un point
générique d'une composante irréductible de $X$, donc $p(\overline{T})$ est
une composante irréductible de $X$. Cela prouve la première assertion. Si
$T \subset X$ est une composante irréductible, alors $p^{-1}(T) =T_K$ est
une union non vide de composantes irréductibles, voir lemme
\ref{varieties-lemma-inverse-image-irreducible}. Chacune d'elles se projette nécessairement
sur $T$ selon la première partie. Par conséquent, l'application est surjective.
\end{proof}

\begin{lemma}
\label{varieties-lemma-irreducible-dense-rational-points-geometrically-irreducible}
Soit $k$ un corps. Supposons que $X$ soit un schéma
sur $k$. Si $X$ est irréductible et possède un ensemble dense de
points $k$-rationnels, alors $X$ est géométriquement irréductible.
\end{lemma}

\begin{proof}
Supposons que $k'/k$ soit une extension finie de corps et soit
$Z, Z' \subset X_{k'}$ des composantes irréductibles. Il
suffit de montrer $Z = Z'$, voir le lemme
\ref{varieties-lemma-characterize-geometrically-irreducible}. Par le lemme
\ref{varieties-lemma-image-irreducible}, nous avons $p(Z) = p(Z') = X$ où
$p : X_{k'} \to X$ est la projection. Si $Z \not = Z'$ alors $Z
\cap Z'$ est nulle part dense dans $X_{k'}$ et donc $p(Z \cap
Z')$ n'est pas dense d'après Morphismes, lemme
\ref{morphisms-lemma-image-nowhere-dense-finite} ; ici nous
utilisons aussi que $p$ est un morphisme fini comme le changement
de base du morphisme fini $\Spec(k') \to \Spec(k)$, voir
Morphismes, lemme \ref{morphisms-lemma-base-change-finite}. Ainsi
nous pouvons choisir un point $k$-rationnel $x \in X$ avec $x \not \in
p(Z \cap Z')$. Puisque le corps résiduel de $x$ est $k$, nous voyons
que $p^{-1}(\{x\}) = \{x'\}$ où $x' \in X_{k'}$ est un point dont le
corps résiduel est $k'$. Puisque $x \in p(Z) = p(Z')$, nous concluons
que $x' \in Z \cap Z'$, ce qui est la contradiction que nous cherchions.
\end{proof}

\begin{lemma}
\label{varieties-lemma-galois-action-irreducible-components}
Soit $k$ un corps, muni d'une clôture
algébrique séparable $\overline{k}$. Supposons
que $X$ soit un schéma sur $k$. Il existe une action
$$
\text{Gal}(\overline{k}/k)^{opp} \times \text{IrredComp}(X_{\overline{k}})
\longrightarrow
\text{IrredComp}(X_{\overline{k}})
$$
avec les propriétés suivantes :
\begin{enumerate}
\item Un élément $\overline{T} \in
\text{IrredComp}(X_{\overline{k}})$ est fixé par l'action si et
seulement s'il existe une composante irréductible $T \subset X$,
qui est géométriquement irréductible sur $k$, telle que
$T_{\overline{k}} = \overline{T}$. \item Pour toute extension de corps
$k'/k$ avec clôture algébrique séparable $\overline{k}'$, le diagramme
$$
\xymatrix{
\text{Gal}(\overline{k}'/k') \times \text{IrredComp}(X_{\overline{k}'})
\ar[r] \ar[d] &
\text{IrredComp}(X_{\overline{k}'}) \ar[d] \\
\text{Gal}(\overline{k}/k) \times \text{IrredComp}(X_{\overline{k}})
\ar[r] &
\text{IrredComp}(X_{\overline{k}})
}
$$
est commutatif (où la flèche verticale droite est une bijection selon
le lemme \ref{varieties-lemma-separably-closed-field-irreducible-components}).
\end{enumerate}
\end{lemma}

\begin{proof}
L'action ( \ref{varieties-equation-galois-action-base-change-kbar} ) de
$\text{Gal}(\overline{k}/k)$ sur $X_{\overline{k}}$ induit
une action sur ses composantes irréductibles. Les composantes
irréductibles sont toujours fermées (Topologie, lemme
\ref{topology-lemma-connected-components}). Par conséquent, si
$\overline{T}$ est comme dans (1), alors par le lemme
\ref{varieties-lemma-closed-fixed-by-Galois} il existe un sous-ensemble
fermé $T \subset X$ tel que $\overline{T} = T_{\overline{k}}$.
Notons que $T$ est géométriquement irréductible sur $k$, voir le
lemme \ref{varieties-lemma-characterize-geometrically-irreducible}. Pour
voir que $T$ est une composante irréductible de $X$, supposons
que $T \subset T'$, $T \not = T'$ où $T'$ est une composante
irréductible de $X$. Supposons que $\overline{\eta}$ soit le
point générique de $\overline{T}$. Il s'envoie sur le point
générique $\eta$ de $T$. Ensuite, le point générique $\xi \in
T'$ se spécialise en $\eta$. Comme $X_{\overline{k}} \to X$ est
plat, il existe un point $\overline{\xi} \in X_{\overline{k}}$
qui s'envoie sur $\xi$ et se spécialise en $\overline{\eta}$. Il
s'ensuit que l'adhérence du singleton $\{\overline{\xi}\}$ est un
sous-ensemble fermé irréductible de $X_{\overline{\xi}}$ qui
contient strictement $\overline{T}$. C'est la contradiction souhaitée.

\medskip\noindent
Nous omettons la preuve de la fonctorialité en (2).
\end{proof}

\begin{lemma}
\label{varieties-lemma-orbit-irreducible-components}
Soit $k$ un corps, muni d'une clôture
algébrique séparable $\overline{k}$. Supposons que
$X$ soit un schéma sur $k$. Les fibres de l'application
$$
\text{IrredComp}(X_{\overline{k}})
\longrightarrow
\text{IrredComp}(X)
$$
du lemme \ref{varieties-lemma-image-irreducible}
sont exactement les orbites de
$\text{Gal}(\overline{k}/k)$ sous l'action du lemme
\ref{varieties-lemma-galois-action-irreducible-components}.
\end{lemma}

\begin{proof}
Soit $T \subset X$ une composante
irréductible de $X$. Supposons que $\eta \in T$ soit
son point générique. D'après les lemmes
\ref{varieties-lemma-inverse-image-irreducible} et
\ref{varieties-lemma-image-irreducible}, les points génériques des
composantes irréductibles de $\overline{T}$ qui se
projettent sur $T$ se projettent sur $\eta$. D'après le lemme
d'algèbre \ref{algebra-lemma-Galois-orbit}, le groupe de
Galois agit transitivement sur tous les points de
$X_{\overline{k}}$ se projetant sur $\eta$. D'où le lemme en découle.
\end{proof}

\begin{lemma}
\label{varieties-lemma-galois-action-irreducible-components-locally-finite-type}
Soit $k$ un corps.
Supposons que $X \to \Spec(k)$ soit
localement de type fini. Dans ce cas
\begin{enumerate}
\item l'action
$$
\text{Gal}(\overline{k}/k)^{opp} \times \text{IrredComp}(X_{\overline{k}})
\longrightarrow
\text{IrredComp}(X_{\overline{k}})
$$
est continue si nous donnons à
$\text{IrredComp}(X_{\overline{k}})$ la topologie
discrète, \item chaque composante irréductible de
$X_{\overline{k}}$ peut être définie sur une extension finie
de $k$, et \item étant donnée une composante irréductible
$T \subset X$, le schéma $T_{\overline{k}}$ est une union
finie de composantes irréductibles de $X_{\overline{k}}$ qui
sont toutes dans le même $\text{Gal}(\overline{k}/k)$-orbite.
\end{enumerate}
\end{lemma}

\begin{proof}
Supposons que $\overline{T}$ soit une composante irréductible
de $X_{\overline{k}}$. Nous pouvons choisir un ouvert affine $U
\subset X$ tel que $\overline{T} \cap U_{\overline{k}}$ ne soit
pas vide. Écrivons $U = \Spec(A)$, ainsi $A$ est une
$k$-algèbre de type fini, voir Morphismes, lemme
\ref{morphisms-lemma-locally-finite-type-characterize}. Par
conséquent, $A_{\overline{k}}$ est une $\overline{k}$-algèbre de
type fini, et en particulier noethérienne. Supposons que $\mathfrak
p = (f_1, \ldots, f_n)$ soit l'idéal premier correspondant à
$\overline{T} \cap U_{\overline{k}}$. Comme $A_{\overline{k}} = A
\otimes_k \overline{k}$, nous voyons qu'il existe une sous-extension
finie $\overline{k}/k'/k$ telle que chaque $f_i \in A_{k'}$. Il est clair
que $\text{Gal}(\overline{k}/k')$ fixe $\overline{T}$, ce qui prouve (1).

\medskip\noindent
La partie (2) suit en appliquant le lemme
\ref{varieties-lemma-galois-action-irreducible-components} (1) à
la situation sur $k'$, ce qui implique que la composante
irréductible $\overline{T}$ est de la forme
$T'_{\overline{k}}$ pour un certain irréductible $T' \subset X_{k'}$.

\medskip\noindent
Pour démontrer (3), soit $T \subset X$ une composante
irréductible. Choisissons une composante irréductible
$\overline{T} \subset X_{\overline{k}}$ qui s'envoie sur
$T$, voir le lemme \ref{varieties-lemma-image-irreducible}. D'après ce
qui précède, l'orbite de $\overline{T}$ est finie, disons
qu'elle est $\overline{T}_1, \ldots, \overline{T}_n$. Alors
$\overline{T}_1 \cup \ldots \cup \overline{T}_n$ est un
sous-ensemble fermé invariant par $\text{Gal}(\overline{k}/k)$
de $X_{\overline{k}}$ et donc de la forme $W_{\overline{k}}$
pour un certain $W \subset X$ fermé d'après le lemme
\ref{varieties-lemma-closed-fixed-by-Galois}. Clairement $W = T$ et cela conclut.
\end{proof}

\begin{lemma}
\label{varieties-lemma-finite-extension-geometrically-irreducible-components}
Soit $k$ un corps. Supposons que $X
\to \Spec(k)$ soit localement de type fini.
Supposons que $X$ ait un nombre fini de composantes
irréductibles. Alors il existe une extension finie
séparable $k'/k$ telle que chaque composante irréductible
de $X_{k'}$ soit géométriquement irréductible sur $k'$.
\end{lemma}

\begin{proof}
Supposons que $\overline{k}$ soit une clôture algébrique
séparable de $k$. L'hypothèse que $X$ a un nombre fini de
composantes irréductibles combinée avec le lemme
\ref{varieties-lemma-galois-action-irreducible-components-locally-finite-type}
(3) montre que $X_{\overline{k}}$ a un nombre fini de
composantes irréductibles $\overline{T}_1, \ldots,
\overline{T}_n$. Par le lemme
\ref{varieties-lemma-galois-action-irreducible-components-locally-finite-type} (2), il existe
une extension finie $\overline{k}/k'/k$ et des composantes irréductibles $T_i
\subset X_{k'}$ telles que $\overline{T}_i = T_{i, \overline{k}}$ et nous avons terminé.
\end{proof}

\begin{lemma}
\label{varieties-lemma-irreducible-components-geometrically-irreducible}
Supposons que $X$ soit un schéma sur le corps
$k$. Supposons que $X$ ait un nombre fini de
composantes irréductibles, toutes géométriquement
irréductibles. Alors $X$ a un nombre fini de
composantes connexes, chacune étant géométriquement connexe.
\end{lemma}

\begin{proof}
Ceci est clair parce qu'une composante connexe est
une union de composantes irréductibles. Nous omettons les détails.
\end{proof}







\section{Schémas géométriquement intègres}
\label{varieties-section-geometrically-integral}

\noindent
Si $X$ est un schéma intègre sur un corps, il peut arriver que $X$
devienne soit non réduit soit réductible après extension du corps de
base. Cela n'arrive pas pour les schémas géométriquement intègres.

\begin{definition}
\label{varieties-definition-geometrically-integral}
Soit $X$ un schéma sur le corps $k$.
\begin{enumerate}
\item Soit $x \in X$. Nous disons que $X$ est {\it
géométriquement ponctuellement intègre en $x$ } si pour toute
extension de corps $k'/k$ et tout $x' \in X_{k'}$ au-dessus de $x$,
l'anneau local $\mathcal{O}_{X_{k'}, x'}$ est intègre. \item Nous
disons que $X$ est {\it géométriquement ponctuellement intègre} si
$X$ est géométriquement ponctuellement intègre en tout point. \item
Nous disons que $X$ est {\it géométriquement intègre} sur $k$ si le
schéma $X_{k'}$ est intègre pour toute extension de corps $k'$ de $k$.
\end{enumerate}
\end{definition}

\noindent
La distinction entre les notions (2) et (3) est nécessaire.
Par exemple, si $k = \mathbf{R}$ et $X =
\Spec(\mathbf{C}[x])$, alors $X$ est géométriquement ponctuellement
intégral sur $\mathbf{R}$ mais bien sûr pas géométriquement intègre.

\begin{lemma}
\label{varieties-lemma-geometrically-integral}
Soit $k$ un corps. Supposons
que $X$ soit un schéma sur $k$. Alors $X$ est
géométriquement intègre sur $k$ si et
seulement si $X$ est à la fois géométriquement
réduit et géométriquement irréductible sur $k$.
\end{lemma}

\begin{proof}
Voir Propriétés, lemme \ref{properties-lemma-characterize-integral}.
\end{proof}

\begin{lemma}
\label{varieties-lemma-proper-geometrically-reduced-global-sections}
Soit $k$ un corps. Supposons que $X$ soit un schéma propre sur $k$.
\begin{enumerate}
\item $A = H^0(X, \mathcal{O}_X)$ est une $k$-algèbre
de dimension finie, \item $A = \prod_{i = 1, \ldots,
n} A_i$ est un produit d'$k$-algèbres locales
artiniennes, un facteur pour chaque composante connexe
de $X$, \item si $X$ est réduit, alors $A = \prod_{i =
1, \ldots, n} k_i$ est un produit de corps, chacun
étant une extension finie de $k$, \item si $X$ est
géométriquement réduit, alors $k_i$ est finie et séparable
sur $k$, \item si $X$ est géométriquement connexe, alors
$A$ est géométriquement irréductible sur $k$, \item si
$X$ est géométriquement irréductible, alors $A$ est
géométriquement irréductible sur $k$, \item si $X$ est
géométriquement réduit et géométriquement connexe, alors $A =
k$, et \item si $X$ est géométriquement intègre, alors $A = k$.
\end{enumerate}
\end{lemma}

\begin{proof}
D'après Cohomologie des Schémas, lemme
\ref{coherent-lemma-proper-over-affine-cohomology-finite},
on voit que $A = H^0(X, \mathcal{O}_X)$ est
une $k$-algèbre de dimension finie. Cela prouve (1).

\medskip\noindent
Alors $A$ est un produit de $k$-algèbres artiniennes locales par
Algèbre, lemme
\ref{algebra-lemma-finite-dimensional-algebra} et la proposition
\ref{algebra-proposition-dimension-zero-ring}. Si $X = Y \amalg Z$
avec $Y$ et $Z$ ouverts dans $X$, alors nous obtenons un idempotent $e
\in A$ en prenant la section de $\mathcal{O}_X$ qui est $1$ sur $Y$ et
$0$ sur $Z$. Réciproquement, si $e \in A$ est un idempotent, alors nous
obtenons une décomposition correspondante de $X$. Enfin, comme $X$ a un
espace topologique sous-jacent noethérien, ses composantes connexes
sont ouvertes. Par conséquent, les composantes connexes de $X$
correspondent de $1$ à $1$ avec des idempotents primitifs de $A$. Cela prouve (2).

\medskip\noindent
Si $X$ est réduit, alors $A$ est réduit. Par
conséquent, les anneaux locaux $A_i = k_i$ sont réduits
et donc corps (par exemple d'après Algèbre, lemme
\ref{algebra-lemma-minimal-prime-reduced-ring}). Cela prouve (3).

\medskip\noindent
Si $X$ est géométriquement réduit, alors $A
\otimes_k \overline{k} = H^0(X_{\overline{k}},
\mathcal{O}_{X_{\overline{k}}})$ (égalité par la
cohomologie des Schémas, lemme
\ref{coherent-lemma-flat-base-change-cohomology})
est réduit. Cela implique que $k_i \otimes_k
\overline{k}$ est un produit de corps et donc $k_i/k$
est séparable, par exemple d'après Algèbre, lemmes
\ref{algebra-lemma-characterize-separable-field-extensions}
et
\ref{algebra-lemma-geometrically-reduced-finite-purely-inseparable-extension}. Cela prouve (4).

\medskip\noindent
Si $X$ est géométriquement connexe, alors $A \otimes_k
\overline{k} = H^0(X_{\overline{k}},
\mathcal{O}_{X_{\overline{k}}})$ est un anneau local de
dimension zéro d'après la partie (2) et donc son spectre a un
point, en particulier il est irréductible. Ainsi $A$ est
géométriquement irréductible. Cela prouve (5). Bien sûr, (5) implique (6).

\medskip\noindent
Si $X$ est géométriquement réduit et géométriquement connexe, alors
$A = k_1$ est un corps et l'extension $k_1/k$ est finie séparable et
géométriquement irréductible. Cependant, alors $k_1 \otimes_k
\overline{k}$ est un produit de $[k_1 : k]$ copies de $\overline{k}$ et
nous concluons que $k_1 = k$. Cela prouve (7). Bien sûr, (7) implique (8).
\end{proof}

\noindent
Voici une version simplifiée de la factorisation de Stein ; la
factorisation de Stein proprement dite sera abordée dans Plus sur les
Morphismes, section \ref{more-morphisms-section-stein-factorization}.

\begin{lemma}
\label{varieties-lemma-baby-stein}
Supposons que $X$ soit un schéma propre sur un corps $k$.
Posons $A = H^0(X, \mathcal{O}_X)$. Les fibres du morphisme
canonique $X \to \Spec(A)$ sont géométriquement connexes.
\end{lemma}

\begin{proof}
Posons $S = \Spec(A)$. Le morphisme canonique $X \to S$
est le morphisme correspondant à $\Gamma(S, \mathcal{O}_S)
= A = \Gamma(X, \mathcal{O}_X)$ d'après Schémas, lemme
\ref{schemes-lemma-morphism-into-affine}. La $k$-algèbre
$A$ est un produit fini $A = \prod A_i$ de $k$-algèbres
locales artiniennes finies sur $k$, voir le lemme
\ref{varieties-lemma-proper-geometrically-reduced-global-sections}.
Désignons par $s_i \in S$ le point correspondant à l'idéal
maximal de $A_i$. Choisissons une clôture algébrique
$\overline{k}$ de $k$ et posons $\overline{A} = A \otimes_k
\overline{k}$. Choisissons un plongement $\kappa(s_i) \to
\overline{k}$ sur $k$ ; cela détermine un morphisme d'algèbres $\overline{k}$.
$$
\sigma_i : \overline{A} = A \otimes_k \overline{k} \to
\kappa(s_i) \otimes_k \overline{k} \to \overline{k}
$$
Considérons le changement de base
$$
\xymatrix{
\overline{X} \ar[r] \ar[d] & X \ar[d] \\
\overline{S} \ar[r] & S
}
$$
de $X$ à $\overline{S} = \Spec(\overline{A})$. Par la
cohomologie des Schémas, lemme
\ref{coherent-lemma-flat-base-change-cohomology} nous avons
$\Gamma(\overline{X}, \mathcal{O}_{\overline{X}}) =
\overline{A}$. Si $\overline{s}_i \in \Spec(\overline{A})$
désigne le point $\overline{k}$-rationnel correspondant à
$\sigma_i$, alors on voit que $\overline{s}_i$ s'envoie sur $s_i
\in S$ et $\overline{X}_{\overline{s}_i}$ est le changement de
base de $X_{s_i}$ par $\Spec(\sigma_i)$. Ainsi, nous voyons qu'il
suffit de prouver le lemme dans le cas où $k$ est algébriquement clos.

\medskip\noindent
Supposons que $k$ soit algébriquement clos. Dans ce cas,
$\kappa(s_i)$ est algébriquement clos et nous devons montrer
que $X_{s_i}$ est connexe. La décomposition en produit $A =
\prod A_i$ correspond à une décomposition en union disjointe
$\Spec(A) = \coprod \Spec(A_i)$, voir Algèbre, lemme
\ref{algebra-lemma-spec-product}. Notons $X_i$ l'image
inverse de $\Spec(A_i)$. Il découle du lemme
\ref{varieties-lemma-proper-geometrically-reduced-global-sections},
partie (2), que $A_i = \Gamma(X_i, \mathcal{O}_{X_i})$. Il faut
observer que $X_{s_i} \to X_i$ est une immersion fermée
induisant un isomorphisme sur les espaces topologiques
sous-jacents (parce que $\Spec(A_i)$ est un singleton). Ainsi, si
$X_{s_i}$ n'est pas connexe, alors $X_i$ ne l'est pas non plus.
Donc soit $X_i$ est vide et $A_i = 0$ ou $X_i$ peut être écrit
comme $U \amalg V$ avec $U$ et $V$ ouverts et non vides, ce qui
impliquerait que $A_i$ possède un idempotent non trivial. Puisque
$A_i$ est local, c'est une contradiction et la démonstration est complète.
\end{proof}

\begin{lemma}
\label{varieties-lemma-geometrically-reduced-stein}
Soit $k$ un corps. Supposons que $X$ soit un schéma propre
géométriquement réduit sur $k$. Les conditions suivantes sont équivalentes
\begin{enumerate}
\item $H^0(X, \mathcal{O}_X) = k$, et
\item $X$ est géométriquement connexe.
\end{enumerate}
\end{lemma}

\begin{proof}
D'après le lemme \ref{varieties-lemma-baby-stein}, nous avons (1)
$\Rightarrow$ (2). D'après le lemme
\ref{varieties-lemma-proper-geometrically-reduced-global-sections}, nous avons (2) $\Rightarrow$ (1).
\end{proof}




\section{Schémas géométriquement normaux}
\label{varieties-section-geometrically-normal}

\noindent
Dans Propriétés, définition
\ref{properties-definition-normal} nous avons défini la
notion de schéma normal. Cette notion est définie même pour
des schémas non noethériens. Ainsi, contrairement à notre
discussion des schémas «~géométriquement réguliers~», nous
considérons toutes les extensions de corps du corps de base.

\begin{definition}
\label{varieties-definition-geometrically-normal}
Soit $X$ un schéma sur le corps $k$.
\begin{enumerate}
\item Soit $x \in X$. On dit que $X$ est {\it
géométriquement normal en $x$ } si pour toute extension
de corps $k'/k$ et tout $x' \in X_{k'}$ au-dessus de $x$
l'anneau local $\mathcal{O}_{X_{k'}, x'}$ est normal.
\item On dit que $X$ est {\it géométriquement normal} sur
$k$ si $X$ est géométriquement normal en tout $x \in X$.
\end{enumerate}
\end{definition}

\begin{lemma}
\label{varieties-lemma-geometrically-normal-at-point}
Soit $k$ un corps.
Supposons que $X$ soit un schéma
sur $k$. Soit $x \in X$. Les
énoncés suivants sont équivalents
\begin{enumerate}
\item $X$ est géométriquement normal en $x$, \item
pour toute extension de corps finie purement
inséparable $k'$ de $k$ et $x' \in X_{k'}$ au-dessus
de $x$, l'anneau local $\mathcal{O}_{X_{k'}, x'}$
est normal, et \item l'anneau $\mathcal{O}_{X, x}$
est géométriquement normal sur $k$ (voir Algèbre,
définition \ref{algebra-definition-geometrically-normal}).
\end{enumerate}
\end{lemma}

\begin{proof}
Il est clair que (1) implique (2). Supposons (2).
Supposons que $k'/k$ soit une extension de corps finie
purement inséparable (par exemple $k = k'$). Considérons
l'anneau $\mathcal{O}_{X, x} \otimes_k k'$. D'après
Algèbre, lemme \ref{algebra-lemma-p-ring-map}, son spectre
est le même que le spectre de $\mathcal{O}_{X, x}$. Par
conséquent, c'est aussi un anneau local (Algèbre, lemme
\ref{algebra-lemma-characterize-local-ring}). Il existe donc
un point unique $x' \in X_{k'}$ au-dessus de $x$ et
$\mathcal{O}_{X_{k'}, x'} \cong \mathcal{O}_{X, x} \otimes_k k'$.
Par hypothèse, c'est un anneau normal. Nous en déduisons donc (3)
d'après Algèbre, lemme \ref{algebra-lemma-geometrically-normal}.

\medskip\noindent
Supposons (3). Supposons que $k'/k$ soit une extension de
corps. Comme $\Spec(k') \to \Spec(k)$ est surjectif,
$X_{k'} \to X$ est aussi surjectif (Morphismes, lemme
\ref{morphisms-lemma-base-change-surjective}). Soit
$x' \in X_{k'}$ un point quelconque au-dessus
de $x$. L'anneau local $\mathcal{O}_{X_{k'}, x'}$ est une
localisation de l'anneau $\mathcal{O}_{X, x} \otimes_k k'$.
Par conséquent, il est normal par hypothèse et (1) est prouvé.
\end{proof}

\begin{lemma}
\label{varieties-lemma-geometrically-normal}
Soit $k$ un corps.
Supposons que $X$ soit un schéma sur
$k$. Les conditions suivantes sont équivalentes
\begin{enumerate}
\item $X$ est géométriquement normal, \item $X_{k'}$
est un schéma normal pour toute extension de corps
$k'/k$, \item $X_{k'}$ est un schéma normal pour
chaque extension de corps de type fini $k'/k$, \item
$X_{k'}$ est un schéma normal pour chaque extension
de corps finie purement inséparable $k'/k$, \item
pour chaque ouvert affine $U \subset X$ l'anneau
$\mathcal{O}_X(U)$ est géométriquement normal (voir
Algèbre, définition
\ref{algebra-definition-geometrically-normal}), et \item $X_{k^{perf}}$ est un schéma normal.
\end{enumerate}
\end{lemma}

\begin{proof}
Supposons (1). Alors pour toute extension de corps
$k'/k$ et tout point $x' \in X_{k'}$, l'anneau
local de $X_{k'}$ en $x'$ est normal. Par
définition, cela signifie que $X_{k'}$ est normal. D'où (2).

\medskip\noindent
Il est clair que (2) implique (3) implique (4).

\medskip\noindent
Supposons (4) et soit $U \subset X$ un sous-schéma affine
ouvert. Alors $U_{k'}$ est un schéma normal pour toute
extension finie purement inséparable $k'/k$ (y compris $k =
k'$). Cela signifie que $k' \otimes_k \mathcal{O}(U)$ est un
anneau normal pour toutes les extensions finies purement
inséparables $k'/k$. Par conséquent, $\mathcal{O}(U)$ est une algèbre
géométriquement normale $k$ par définition. Ainsi, (4) implique (5).

\medskip\noindent
Supposons (5). Pour toute extension de corps $k'/k$, le
changement de base $X_{k'}$ s'obtient par recollement des spectres
des anneaux $\mathcal{O}_X(U) \otimes_k k'$ où $U$ est un
ouvert affine dans $X$ (voir Schémas, section
\ref{schemes-section-fibre-products}). Donc $X_{k'}$ est normal. Ainsi (1) est vrai.

\medskip\noindent
L'équivalence de (5) et (6) découle de la
définition des algèbres géométriquement normales
et de l'équivalence que l'on vient de démontrer de (3) et (4).
\end{proof}

\begin{lemma}
\label{varieties-lemma-geometrically-normal-upstairs}
Soit $k$ un corps. Supposons que $X$
soit un schéma sur $k$. Supposons que $k'/k$ soit
une extension de corps. Supposons que $x \in X$ soit
un point, et soit $x' \in X_{k'}$ un point situé
au-dessus de $x$. Les énoncés suivants sont équivalents
\begin{enumerate}
\item $X$ est géométriquement normal en $x$,
\item $X_{k'}$ est géométriquement normal en $x'$.
\end{enumerate}
En particulier, $X$ est géométriquement normal sur $k$ si et
seulement si $X_{k'}$ est géométriquement normal sur $k'$.
\end{lemma}

\begin{proof}
Il est clair que (1) implique (2). Supposons (2). Supposons que
$k''/k$ soit une extension de corps finie purement inséparable et
soit $x'' \in X_{k''}$ un point au-dessus de $x$ (en fait, il
est unique). Choisissons une extension de corps commune $k'''/k$ de
$k'/k$ et $k''/k$, voir corps, lemme
\ref{fields-lemma-common-extension-field}. Choisissons un point $x''' \in X_{k'''}$
au-dessus à la fois de $x'$ et $x''$. Considérons le morphisme d'anneaux locaux
$$
\mathcal{O}_{X_{k''}, x''} \longrightarrow \mathcal{O}_{X_{k'''}, x''''}.
$$
Ceci est un homomorphisme local plat d'anneaux et donc
fidèlement plat. D'après (2), on voit que l'anneau local à
droite est normal. Ainsi, d'après L'Algèbre, lemme
\ref{algebra-lemma-descent-normal}, nous concluons que
$\mathcal{O}_{X_{k''}, x''}$ est normal. D'après le lemme
\ref{varieties-lemma-geometrically-normal-at-point}, on voit que $X$ est géométriquement normal en $x$.
\end{proof}

\begin{lemma}
\label{varieties-lemma-fibre-product-normal}
Soit $k$ un corps. Supposons que $X$ soit un
schéma géométriquement normal sur $k$ et soit $Y$ un schéma
normal sur $k$. Alors $X \times_k Y$ est un schéma normal.
\end{lemma}

\begin{proof}
Cela se ramène à l'algèbre, lemme
\ref{algebra-lemma-geometrically-normal-tensor-normal}
par le lemme
\ref{varieties-lemma-geometrically-normal}.
\end{proof}

\begin{lemma}
\label{varieties-lemma-base-change-normal-by-separable}
Soit $k$ un corps. Supposons que $X$ soit un schéma normal sur $k$.
Supposons que $K/k$ soit une extension de corps séparable. Alors $X_K$ est un schéma normal.
\end{lemma}

\begin{proof}
Cela résulte du lemme
\ref{varieties-lemma-fibre-product-normal} et du lemme d'Algèbre
\ref{algebra-lemma-separable-field-extension-geometrically-normal}.
\end{proof}

\begin{lemma}
\label{varieties-lemma-geometrically-normal-stein}
Soit $k$ un corps. Supposons que $X$ soit un schéma propre
géométriquement normal sur $k$. Les conditions suivantes sont équivalentes
\begin{enumerate}
\item $H^0(X, \mathcal{O}_X) = k$, \item
$X$ est géométriquement connexe, \item
$X$ est géométriquement irréductible, et
\item $X$ est géométriquement intègre.
\end{enumerate}
\end{lemma}

\begin{proof}
D'après le lemme
\ref{varieties-lemma-geometrically-reduced-stein}, nous avons
l'équivalence de (1) et (2). Un schéma normal localement
noethérien (tel que $X_{\overline{k}}$) est une réunion
disjointe de ses composantes irréductibles (Propriétés, lemme
\ref{properties-lemma-normal-Noetherian}). Ainsi, on voit que
(2) et (3) sont équivalents. Puisque $X_{\overline{k}}$ est
supposé réduit, on voit que (3) et (4) sont également équivalents.
\end{proof}









\section{Changement de corps et schémas localement noethériens}
\label{varieties-section-locally-Noetherian}

\noindent
Soit $X$ un schéma localement noethérien sur un corps $k$. Il n'est pas
toujours vrai que $X_{k'}$ soit aussi localement noethérien. Par
exemple, si $X = \Spec(\overline{\mathbf{Q}})$ et $k = \mathbf{Q}$,
alors $X_{\overline{\mathbf{Q}}}$ est le spectre de
$\overline{\mathbf{Q}} \otimes_{\mathbf{Q}} \overline{\mathbf{Q}}$ qui n'est
pas noethérien. (Indice : Il a trop d'idempotents). Mais si nous ne faisons
que des changements de base utilisant des extensions de corps de type fini,
alors la propriété noethérienne est préservée. (Ou si $X$ est localement de
type fini sur $k$, puisque cette propriété est préservée par changement de base.)

\begin{lemma}
\label{varieties-lemma-locally-Noetherian-base-change}
Soit $k$ un corps. Supposons
que $X$ soit un schéma sur $k$. Supposons que
$k'/k$ soit une extension de corps de type
fini. Alors $X$ est localement noethérien si et
seulement si $X_{k'}$ est localement noethérien.
\end{lemma}

\begin{proof}
D'après Propriétés, lemme
\ref{properties-lemma-locally-Noetherian} nous réduisons au cas où
$X$ est affine, disons $X = \Spec(A)$. Dans ce cas, nous devons
prouver que $A$ est noethérien si et seulement si $A_{k'}$ est
noethérien. Puisque $A \to A_{k'} = k' \otimes_k A$ est fidèlement
plat, on voit que si $A_{k'}$ est noethérien, alors $A$ l'est
aussi, d'après Algèbre, lemme
\ref{algebra-lemma-descent-Noetherian}. Réciproquement, si $A$ est noethérien, alors
$A_{k'}$ est noethérien d'après Algèbre, lemme \ref{algebra-lemma-Noetherian-field-extension}.
\end{proof}







\section{Schémas géométriquement réguliers}
\label{varieties-section-geometrically-regular}

\noindent
Un schéma géométriquement régulier sur un corps $k$ est un schéma
localement noethérien sur $k$ qui reste régulier après des changements
appropriés de corps de base. Un schéma de type fini sur $k$ est
géométriquement régulier si et seulement s'il est lisse sur $k$ (voir le lemme
\ref{varieties-lemma-geometrically-regular-smooth}). La notion de régularité géométrique
est particulièrement intéressante dans les situations où la lissité ne peut pas
être utilisée, comme pour les fibres formelles (insérer référence future ici).

\medskip\noindent
Dans la définition suivante, nous nous limitons aux schémas
localement noethériens, puisque la propriété d'être un anneau
local régulier n'est définie que pour les anneaux locaux
noethériens. D'après le lemme
\ref{varieties-lemma-locally-Noetherian-base-change} ci-dessus, si nous nous limitons
aux extensions de corps de type fini, alors cette propriété est préservée par
changement de corps de base. Ce commentaire sera utilisé sans référence
supplémentaire dans cette section. En particulier, la définition suivante a un sens.

\begin{definition}
\label{varieties-definition-geometrically-regular}
Soit $k$ un corps. Supposons que $X$ soit un schéma localement noethérien sur $k$.
\begin{enumerate}
\item Soit $x \in X$. On dit que $X$ est {\it
géométriquement régulier en $x$} sur $k$ si pour toute
extension de corps de type fini $k'/k$ et tout $x' \in X_{k'}$
au-dessus de $x$, l'anneau local $\mathcal{O}_{X_{k'}, x'}$ est
régulier. \item On dit que $X$ est {\it géométriquement régulier
sur $k$} si $X$ est géométriquement régulier en tout point.
\end{enumerate}
\end{definition}

\noindent
Une définition analogue permet de définir les schémas géométriquement
de Cohen--Macaulay, $(R_k)$ et $(S_k)$ sur un corps. Nous ajouterons
des sections séparées lorsque cela sera nécessaire.

\begin{lemma}
\label{varieties-lemma-geometrically-regular-at-point}
Soit $k$ un corps.
Supposons que $X$ soit un schéma
localement noethérien sur $k$. Soit $x \in
X$. Les assertions suivantes sont équivalentes
\begin{enumerate}
\item $X$ est géométriquement régulier en $x$, \item
pour toute extension de corps finie purement
inséparable $k'$ de $k$ et $x' \in X_{k'}$ située
au-dessus de $x$ l'anneau local $\mathcal{O}_{X_{k'},
x'}$ est régulier, et \item l'anneau $\mathcal{O}_{X,
x}$ est géométriquement régulier sur $k$ (voir Algèbre,
définition \ref{algebra-definition-geometrically-regular}).
\end{enumerate}
\end{lemma}

\begin{proof}
Il est clair que (1) implique (2). Supposons (2). Cela implique
en particulier que $\mathcal{O}_{X, x}$ est un anneau local
régulier. Supposons que $k'/k$ soit une extension de corps finie
purement inséparable. Considérons l'anneau $\mathcal{O}_{X, x}
\otimes_k k'$. D'après Algèbre, lemme
\ref{algebra-lemma-p-ring-map}, son spectre est le même que le
spectre de $\mathcal{O}_{X, x}$. Par conséquent, c'est aussi un
anneau local (Algèbre, lemme
\ref{algebra-lemma-characterize-local-ring}). Donc, il y a un point unique $x'
\in X_{k'}$ au-dessus de $x$ et $\mathcal{O}_{X_{k'}, x'} \cong
\mathcal{O}_{X, x} \otimes_k k'$. Par hypothèse, il s'agit d'un anneau régulier.
Ainsi, nous déduisons (3) de la définition d'un anneau géométriquement régulier.

\medskip\noindent
Supposons (3). Supposons que $k'/k$ soit une extension de
corps. Comme $\Spec(k') \to \Spec(k)$ est surjectif,
$X_{k'} \to X$ est aussi surjectif (Morphismes, lemme
\ref{morphisms-lemma-base-change-surjective}). Supposons que
$x' \in X_{k'}$ est un point quelconque situé au-dessus de
$x$. L'anneau local $\mathcal{O}_{X_{k'}, x'}$ est une
localisation de l'anneau $\mathcal{O}_{X, x} \otimes_k k'$.
Par conséquent, il est régulier par hypothèse et (1) est prouvé.
\end{proof}

\begin{lemma}
\label{varieties-lemma-geometrically-regular}
Soit $k$ un corps. Supposons que
$X$ soit un schéma localement noethérien sur
$k$. Les affirmations suivantes sont équivalentes
\begin{enumerate}
\item $X$ est géométriquement régulier, \item $X_{k'}$
est un schéma régulier pour toute extension de corps
de type fini $k'/k$, \item $X_{k'}$ est un schéma
régulier pour toute extension de corps finie purement
inséparable $k'/k$, \item pour tout ouvert affine $U
\subset X$ l'anneau $\mathcal{O}_X(U)$ est
géométriquement régulier (voir Algèbre, définition
\ref{algebra-definition-geometrically-regular}), et \item il
existe un recouvrement ouvert affine $X = \bigcup U_i$ tel que
chaque $\mathcal{O}_X(U_i)$ est géométriquement régulier sur $k$.
\end{enumerate}
\end{lemma}

\begin{proof}
Supposons (1). Alors pour toute extension de corps de
type fini $k'/k$ et pour tout point $x' \in X_{k'}$,
l'anneau local de $X_{k'}$ en $x'$ est régulier. D'après
les Propriétés, lemme
\ref{properties-lemma-characterize-regular}, cela signifie que $X_{k'}$ est régulier. D'où (2).

\medskip\noindent
Il est clair que (2) implique (3).

\medskip\noindent
Supposons (3) et soit $U \subset X$ un sous-schéma affine
ouvert. Alors $U_{k'}$ est un schéma régulier pour toute
extension finie purement inséparable $k'/k$ (y compris $k =
k'$). Cela signifie que $k' \otimes_k \mathcal{O}(U)$ est un
anneau régulier pour toutes les extensions finies purement
inséparables $k'/k$. Par conséquent, $\mathcal{O}(U)$ est une
algèbre $k$ géométriquement régulière et on voit que (4) est vrai.

\medskip\noindent
Il est clair que (4) implique (5). Supposons que $X = \bigcup
U_i$ soit un recouvrement affine ouvert comme dans (5). Pour toute
extension de corps $k'/k$, le changement de base $X_{k'}$ est
obtenu en recollant les spectres des anneaux $\mathcal{O}_X(U_i)
\otimes_k k'$ (voir Schémas, section
\ref{schemes-section-fibre-products}). Par conséquent, $X_{k'}$ est régulier. Donc (1) est vérifié.
\end{proof}

\begin{lemma}
\label{varieties-lemma-geometrically-regular-upstairs}
Soit $k$ un corps. Supposons que $X$ soit
un schéma sur $k$. Supposons que $k'/k$ est une
extension de corps de type fini. Supposons que $x \in
X$ est un point, et soit $x' \in X_{k'}$ un point situé
au-dessus de $x$. Les énoncés suivants sont équivalents
\begin{enumerate}
\item $X$ est géométriquement régulier en $x$,
\item $X_{k'}$ est géométriquement régulier en $x'$.
\end{enumerate}
En particulier, $X$ est géométriquement régulier sur $k$ si et
seulement si $X_{k'}$ est géométriquement régulier sur $k'$.
\end{lemma}

\begin{proof}
Il est clair que (1) implique (2). Supposons (2). Supposons que
$k''/k$ soit une extension de corps finie purement inséparable
et soit $x'' \in X_{k''}$ un point au-dessus de $x$ (en fait il
est unique). Nous pouvons trouver une extension de corps
commune, de type fini, $k'''/k$ de $k'/k$ et $k''/k$, voir
corps, lemme \ref{fields-lemma-common-extension-field} et sa
démonstration, et un point $x''' \in X_{k'''}$ se trouvant au-dessus
à la fois de $x'$ et $x''$. Considérons le morphisme d'anneaux locaux
$$
\mathcal{O}_{X_{k''}, x''} \longrightarrow \mathcal{O}_{X_{k'''}, x''''}.
$$
Ceci est un homomorphisme plat local d'anneaux locaux
noethériens et donc fidèlement plat. D'après (2), nous
voyons que l'anneau local à droite est régulier. Ainsi,
d'après Algèbre, lemme
\ref{algebra-lemma-flat-under-regular}, nous concluons que
$\mathcal{O}_{X_{k''}, x''}$ est régulier. D'après le lemme
\ref{varieties-lemma-geometrically-regular-at-point}, on voit que $X$ est géométriquement régulier en $x$.
\end{proof}

\noindent
Le lemme suivant est une variante géométrique de
l'Algèbre, lemme \ref{algebra-lemma-geometrically-regular-descent}.

\begin{lemma}
\label{varieties-lemma-flat-under-geometrically-regular}
Soit $k$ un corps. Supposons que $f : X \to Y$
soit un morphisme de schémas localement noethériens sur $k$.
Supposons que $x \in X$ soit un point et posons $y = f(x)$. Si
$X$ est géométriquement régulier en $x$ et que $f$ est plat en
$x$ alors $Y$ est géométriquement régulier en $y$. En
particulier, si $X$ est géométriquement régulier sur $k$ et que $f$
est plat et surjectif, alors $Y$ est géométriquement régulier sur $k$.
\end{lemma}

\begin{proof}
Supposons que $k'$ soit une extension finie purement
inséparable de $k$. Supposons que $f' : X_{k'} \to
Y_{k'}$ soit le changement de base de $f$. Supposons que
$x' \in X_{k'}$ soit le point unique situé au-dessus de
$x$. Si nous montrons que $Y_{k'}$ est régulier en $y' =
f'(x')$, alors $Y$ est géométriquement régulier sur $k$
en $y'$, voir le lemme
\ref{varieties-lemma-geometrically-regular}. D'après Morphismes, lemme
\ref{morphisms-lemma-base-change-module-flat} le morphisme $X_{k'}
\to Y_{k'}$ est plat en $x'$. Par conséquent, l'homomorphisme d'anneaux
$$
\mathcal{O}_{Y_{k'}, y'}
\longrightarrow
\mathcal{O}_{X_{k'}, x'}
$$
est un homomorphisme local plat d'anneaux locaux
noethériens avec le côté droit régulier par hypothèse. Par
conséquent, le côté gauche est un anneau local régulier
d'après Algèbre, lemme \ref{algebra-lemma-flat-under-regular}.
\end{proof}

\begin{lemma}
\label{varieties-lemma-geometrically-regular-smooth}
Soit $k$ un corps. Supposons que $X$
soit un schéma localement de type fini sur $k$. Soit $x
\in X$. Alors $X$ est géométriquement régulier en $x$
si et seulement si $X \to \Spec(k)$ est lisse en $x$
(Morphismes, définition \ref{morphisms-definition-smooth}).
\end{lemma}

\begin{proof}
La question est locale autour de $x$, donc
on peut supposer que $X = \Spec(A)$
pour une certaine $k$-algèbre de type fini.
Soit $\mathfrak p$ l'idéal premier correspondant à $x$.

\medskip\noindent
Si $A$ est lisse sur $k$ à $\mathfrak p$, alors nous pouvons
localiser $A$ et supposer que $A$ est lisse sur $k$. Dans ce cas,
$k' \otimes_k A$ est lisse sur $k'$ pour tout corps d'extension
$k'/k$, et chacun de ces anneaux noethériens est régulier selon
l'Algèbre, lemme \ref{algebra-lemma-characterize-smooth-over-field}.

\medskip\noindent
Supposons que $X$ soit géométriquement régulier en $x$.
Considérons le corps résiduel $K := \kappa(x) =
\kappa(\mathfrak p)$ de $x$. C'est une extension de type
fini de $k$. D'après Algèbre, lemme
\ref{algebra-lemma-make-separable}, il existe une
extension finie purement inséparable $k'/k$ telle que le
compositum $k'K$ soit une extension de corps séparable de
$k'$. Supposons que $\mathfrak p' \subset A' = k'
\otimes_k A$ soit un idéal premier au-dessus de $\mathfrak
p$. C'est le seul premier au-dessus de $\mathfrak p$, voir
Algèbre, lemme \ref{algebra-lemma-p-ring-map}. Ainsi, le
corps résiduel $K' := \kappa(\mathfrak p')$ est le
compositum $k'K$. Par hypothèse, l'anneau local
$(A')_{\mathfrak p'}$ est régulier. Ainsi, d'après l'Algèbre,
lemme \ref{algebra-lemma-separable-smooth}, on voit que $k'
\to A'$ est lisse en $\mathfrak p'$. Cela implique à son tour
que $k \to A$ est lisse en $\mathfrak p$ d'après Algèbre, lemme
\ref{algebra-lemma-smooth-field-change-local}. Le lemme est prouvé.
\end{proof}


\begin{example}
\label{varieties-example-geometrically-reduced-not-normal}
Soit $k =\mathbf{F}_p(t)$. Il est assez facile de donner un
exemple d'une variété régulière $V$ sur $k$ qui n'est pas
géométriquement réduite. Par exemple, nous pouvons prendre
$\Spec(k[x]/(x^p - t))$. En fait, il existe un exemple d'une
variété régulière $V$ qui est géométriquement réduite, mais pas même
géométriquement normale. À savoir, prenons pour $p > 2$ le schéma $V
= \Spec(k[x, y]/(y^2 - x^p + t))$. C'est une variété puisque le
polynôme $y^2 - x^p + t \in k[x, y]$ est irréductible. Le morphisme
$V \to \Spec(k)$ est lisse en tous les points sauf au point $v_0 \in
V$ correspondant à l'idéal maximal $(y, x^p - t)$ (car $2y$ est
inversible). En particulier, on voit que $V$ est (géométriquement)
régulière en tous les points, sauf éventuellement $v_0$. L'anneau local
$$
\mathcal{O}_{V, v_0} = \left(k[x, y]/(y^2 - x^p + t)\right)_{(y, x^p - t)}
$$
est un anneau intègre de dimension $1$. Son idéal maximal est engendré
par un seul élément, à savoir $y$. Par conséquent, c'est un anneau
à valuation discrète et régulier. Soit $k' = k[t^{1/p}]$. On note
$t' = t^{1/p} \in k'$, $V' = V_{k'}$, $v'_0 \in V'$ le point
unique au-dessus de $v_0$. Au-dessus de $k'$, nous pouvons écrire
$x^p - t = (x - t')^p$, mais le polynôme $y^2 - (x - t')^p$ est
toujours irréductible et $V'$ est toujours une variété. Mais l'élément
$$
\frac{y}{x - t'} \in (\text{fraction field of }\mathcal{O}_{V', v'_0})
$$
est intégral sur $\mathcal{O}_{V', v'_0}$ (il suffit
de calculer son carré) et n'y est pas contenu, donc
$V'$ n'est pas normal en $v'_0$. Cela conclut l'exemple.
\end{example}







\section{Changement de corps et propriété de Cohen--Macaulay}
\label{varieties-section-CM}

\noindent
Le lemme suivant dit qu'il n'a pas de sens de définir
des schémas géométriquement de Cohen--Macaulay, puisque
ceux-ci seraient les mêmes que les schémas de Cohen-Macaulay.

\begin{lemma}
\label{varieties-lemma-CM-base-change}
Supposons que $X$ soit un schéma localement noethérien sur le
corps $k$. Supposons que $k'/k$ soit une extension de corps de
type fini. Supposons que $x \in X$ soit un point, et soit $x'
\in X_{k'}$ un point au-dessus de $x$. Alors nous avons
$$
\mathcal{O}_{X, x}\text{ is Cohen-Macaulay}
\Leftrightarrow
\mathcal{O}_{X_{k'}, x'}\text{ is Cohen-Macaulay}
$$
Si $X$ est localement de type fini sur $k$, il en
va de même pour toute extension de corps $k'/k$.
\end{lemma}

\begin{proof}
Le premier cas du lemme découle de l'Algèbre,
lemme \ref{algebra-lemma-CM-geometrically-CM}. Le
deuxième cas du lemme est équivalent à l'Algèbre,
lemme \ref{algebra-lemma-extend-field-CM-locus}.
\end{proof}







\section{Changement de corps et propriété de Jacobson}
\label{varieties-section-overfield}

\noindent
Un schéma localement de type fini sur un corps a beaucoup
de points fermés, à savoir qu'il est de Jacobson. De plus, les
corps résiduels sont des extensions finies du corps de base.

\begin{lemma}
\label{varieties-lemma-locally-finite-type-Jacobson}
Supposons que $X$ soit un schéma qui est
localement de type fini sur $k$. Alors
\begin{enumerate}
\item pour tout point fermé $x \in X$
l'extension $\kappa(x)/k$ est algébrique, et
\item $X$ est un schéma de Jacobson (Propriétés,
définition \ref{properties-definition-jacobson} ).
\end{enumerate}
\end{lemma}

\begin{proof}
Un schéma est de Jacobson si et seulement s'il possède un
recouvrement par des ouverts affines de schémas de Jacobson,
voir Propriétés, lemme
\ref{properties-lemma-locally-jacobson}. La propriété sur les corps
résiduels aux points fermés est également locale sur $X$. Par
conséquent, on peut supposer que $X$ est affine. Dans ce cas,
le résultat est une conséquence du Nullstellensatz de Hilbert, voir
Algèbre, théorème \ref{algebra-theorem-nullstellensatz}. Il découle
également d'une combinaison de Morphismes, lemmes
\ref{morphisms-lemma-jacobson-finite-type-points},
\ref{morphisms-lemma-Jacobson-universally-Jacobson} et \ref{morphisms-lemma-ubiquity-Jacobson-schemes}.
\end{proof}

\noindent
Il s'avère que si $X$ n'est pas localement de type fini, alors nous pouvons obtenir
le même résultat après avoir fait une extension du corps de base suffisamment grande.

\begin{lemma}
\label{varieties-lemma-make-Jacobson}
Supposons que $X$ soit un schéma sur un corps
$k$. Pour toute extension de corps $K/k$ dont la
cardinalité est suffisamment grande, nous avons
\begin{enumerate}
\item pour tout point fermé $x \in X_K$
l'extension $\kappa(x)/K$ est algébrique, et
\item $X_K$ est un schéma de Jacobson (Propriétés,
définition \ref{properties-definition-jacobson} ).
\end{enumerate}
\end{lemma}

\begin{proof}
Choisissons un recouvrement ouvert affine $X =
\bigcup U_i$. D'après Algèbre, lemme
\ref{algebra-lemma-base-change-Jacobson} et les
Propriétés, lemme
\ref{properties-lemma-affine-jacobson}, il existe des
cardinaux $\kappa_i$ tels que $U_{i, K}$ possède les
propriétés souhaitées sur $K$ si $\#(K) \geq \kappa_i$.
Posons $\kappa = \max\{\kappa_i\}$. Alors, si la cardinalité
de $K$ est supérieure à $\kappa$, on voit que chaque
$U_{i, K}$ satisfait les conclusions du lemme. Par
conséquent, $X_K$ est Jacobson d'après Propriétés, lemme
\ref{properties-lemma-locally-jacobson}. L'énoncé sur les corps
résiduels aux points fermés de $X_K$ découle des énoncés
correspondants pour les corps résiduels des points fermés de $U_{i, K}$.
\end{proof}





\section{Changement de corps et faisceaux inversibles amples}
\label{varieties-section-change-fields-ample}

\noindent
Le résultat suivant est typique des résultats de cette section.

\begin{lemma}
\label{varieties-lemma-ample-after-field-extension}
Soit $k$ un corps. Supposons que $X$ soit
un schéma sur $k$. S'il existe un faisceau inversible
ample sur $X_K$ pour une certaine extension de corps
$K/k$, alors $X$ possède un faisceau inversible ample.
\end{lemma}

\begin{proof}
Supposons que $K/k$ soit une extension de corps telle que
$X_K$ possède un faisceau inversible ample $\mathcal{L}$. Le
morphisme $X_K \to X$ est surjectif. Par conséquent, $X$ est
quasi-compact en tant qu'image d'un schéma quasi-compact
(Propriétés, définition \ref{properties-definition-ample}).
Puisque $X_K$ est quasi-séparé (d'après Propriétés, lemme
\ref{properties-lemma-affine-s-opens-cover-quasi-separated}), nous
voyons que $X$ est quasi-séparé : si $U, V \subset X$ sont ouverts
affines, alors $(U \cap V)_K = U_K \cap V_K$ est quasi-compact et
$(U \cap V)_K \to U \cap V$ est surjectif. Ainsi, le lemme
\ref{schemes-lemma-characterize-quasi-separated} sur les schémas s'applique.

\medskip\noindent
Écrivons $K = \colim A_i$ comme colimite des sous-algèbres de
$K$ qui sont de type fini sur $k$. Notons $X_i = X
\times_{\Spec(k)} \Spec(A_i)$. Puisque $X_K = \lim X_i$, nous
trouvons un $i$ et un faisceau inversible $\mathcal{L}_i$ sur
$X_i$ dont le image réciproque sur $X_K$ est $\mathcal{L}$
(Limites, lemme \ref{limits-lemma-descend-invertible-modules} ;
ici et ci-dessous, nous utilisons que $X$ est quasi-compact et
quasi-séparé comme montré précédemment). D'après Limites, lemme
\ref{limits-lemma-limit-ample}, on peut supposer que
$\mathcal{L}_i$ est ample après avoir éventuellement augmenté
$i$. Fixons un tel $i$ et soit $\mathfrak m \subset A_i$ un idéal
maximal. Par le Nullstellensatz de Hilbert (Algèbre, théorème
\ref{algebra-theorem-nullstellensatz}), le corps résiduel $k' =
A_i/\mathfrak m$ est une extension finie de $k$. Ainsi $X_{k'}
\subset X_i$ est un sous-schéma fermé et possède donc un faisceau
inversible ample (Propriétés, lemme
\ref{properties-lemma-ample-on-closed}). Puisque $X_{k'} \to X$ est localement
libre de rang fini, nous en concluons que $X$ possède un faisceau inversible
ample d'après Diviseurs, proposition \ref{divisors-proposition-push-down-ample}.
\end{proof}

\begin{lemma}
\label{varieties-lemma-quasi-affine-after-field-extension}
Soit $k$ un corps. Soit $X$ un schéma sur $k$. Si $X_K$ est
quasi-affine pour une certaine extension de corps $K/k$, alors $X$ est quasi-affine.
\end{lemma}

\begin{proof}
Supposons que $K/k$ soit une extension de corps telle que $X_K$ soit
quasi-affine. Le morphisme $X_K \to X$ est surjectif. Ainsi $X$ est
quasi-compact comme l'image d'un schéma quasi-compact (Propriétés,
définition \ref{properties-definition-quasi-affine}). Puisque $X_K$
est quasi-séparé (en tant que sous-schéma ouvert d'un schéma affine),
on voit que $X$ est quasi-séparé : si $U, V \subset X$ sont
ouverts affines, alors $(U \cap V)_K = U_K \cap V_K$ est quasi-compact
et $(U \cap V)_K \to U \cap V$ est surjectif. Ainsi, le lemme sur les
schémas, \ref{schemes-lemma-characterize-quasi-separated}, s'applique.

\medskip\noindent
Écrivons $K = \colim A_i$ comme colimite des
sous-algèbres de $K$ qui sont de type fini sur $k$. Noter
$X_i = X \times_{\Spec(k)} \Spec(A_i)$. Puisque $X_K = \lim
X_i$, nous trouvons un $i$ tel que $X_i$ soit quasi-affine
(Limites, lemme \ref{limits-lemma-limit-quasi-affine} ; ici
nous utilisons que $X$ est quasi-compact et quasi-séparé
comme montré précédemment). Par le Nullstellensatz de
Hilbert (Algèbre, théorème
\ref{algebra-theorem-nullstellensatz}), le corps résiduel $k' =
A_i/\mathfrak m$ est une extension finie de $k$. Ainsi $X_{k'}
\subset X_i$ est un sous-schéma fermé donc est quasi-affine
(Propriétés, lemme
\ref{properties-lemma-quasi-affine-locally-closed}). Puisque $X_{k'} \to X$ est localement libre
de type fini, nous concluons par Diviseurs, lemme \ref{divisors-lemma-push-down-quasi-affine}.
\end{proof}

\begin{lemma}
\label{varieties-lemma-quasi-projective-after-field-extension}
Soit $k$ un corps. Supposons que $X$ soit un schéma
sur $k$. Si $X_K$ est quasi-projectif sur $K$ pour une certaine
extension de corps $K/k$, alors $X$ est quasi-projectif sur $k$.
\end{lemma}

\begin{proof}
Par définition, un morphisme de schémas $g : Y \to T$ est
quasi-projectif s'il est localement de type fini, quasi-compact,
et s'il existe un faisceau inversible $g$-ample sur $Y$.
Supposons que $K/k$ soit une extension de corps telle que $X_K$
soit quasi-projectif sur $K$. Supposons que $\Spec(A) \subset X$
soit un ouvert affine. Alors $U_K$ est un sous-schéma ouvert
affine de $X_K$, donc $A_K$ est une $K$-algèbre de type fini.
Alors $A$ est une $k$-algèbre de type fini d'après le lemme
d'algèbre \ref{algebra-lemma-finite-type-descends}. Par conséquent,
$X \to \Spec(k)$ est localement de type fini. Puisque $X_K \to
\Spec(K)$ est quasi-compact, on voit que $X_K$ est
quasi-compact, donc $X$ est quasi-compact, donc $X \to \Spec(k)$ est
de type fini. D'après le lemme sur les morphismes
\ref{morphisms-lemma-finite-type-over-affine-ample-very-ample}, nous
voyons que $X_K$ possède un faisceau inversible ample. Ensuite, $X$ possède
un faisceau inversible ample d'après le lemme
\ref{varieties-lemma-ample-after-field-extension}. Par conséquent, $X \to \Spec(k)$ est
quasi-projectif d'après Morphismes, lemme \ref{morphisms-lemma-finite-type-over-affine-ample-very-ample}.
\end{proof}

\noindent
Le lemme suivant est un cas particulier de la Descente,
lemme \ref{descent-lemma-descending-property-proper}.

\begin{lemma}
\label{varieties-lemma-proper-after-field-extension}
Soit $k$ un corps. Supposons que $X$ soit un
schéma sur $k$. Si $X_K$ est propre sur $K$ pour une
certaine extension de corps $K/k$, alors $X$ est propre sur $k$.
\end{lemma}

\begin{proof}
Supposons que $K/k$ soit une extension de corps telle que $X_K$ soit
propre sur $K$. Rappelons que cela implique que $X_K$ est séparé et
quasi-compact (Morphismes, définition
\ref{morphisms-definition-proper}). Le morphisme $X_K \to X$ est
surjectif. Ainsi, $X$ est quasi-compact comme l'image d'un schéma
quasi-compact (Propriétés, définition \ref{properties-definition-ample}).
Puisque $X_K$ est séparé, on voit que $X$ est quasi-séparé : si $U, V
\subset X$ sont ouverts affines, alors $(U \cap V)_K = U_K \cap V_K$ est
quasi-compact et $(U \cap V)_K \to U \cap V$ est surjectif. Ainsi, le
lemme des schémas \ref{schemes-lemma-characterize-quasi-separated} s'applique.

\medskip\noindent
Écrivons $K = \colim A_i$ comme colimite des
sous-algèbres de $K$ qui sont de type fini sur $k$.
Notons $X_i = X \times_{\Spec(k)} \Spec(A_i)$. D'après
Limites, lemme \ref{limits-lemma-eventually-proper}, il
existe un $i$ tel que $X_i \to \Spec(A_i)$ soit propre.
Ici, nous utilisons que $X$ est quasi-compact et
quasi-séparé comme montré précédemment. Choisissons un
idéal maximal $\mathfrak m \subset A_i$. D'après le
Nullstellensatz de Hilbert (Algèbre, théorème
\ref{algebra-theorem-nullstellensatz}), le corps résiduel
$k' = A_i/\mathfrak m$ est une extension finie de $k$. Le
changement de base $X_{k'} \to \Spec(k')$ est propre
(Morphismes, lemme
\ref{morphisms-lemma-base-change-proper}). Puisque $k'/k$ est
fini, à la fois $X_{k'} \to X$ et la composition $X_{k'} \to
\Spec(k)$ sont également propres (Morphismes, lemmes
\ref{morphisms-lemma-finite-proper},
\ref{morphisms-lemma-base-change-proper} et
\ref{morphisms-lemma-composition-proper}). Le premier implique que $X$ est
séparé sur $k$ puisque $X_{k'}$ est séparé (Morphismes, lemme
\ref{morphisms-lemma-image-universally-closed-separated}). Le second implique que $X \to
\Spec(k)$ est propre d'après Morphismes, lemme \ref{morphisms-lemma-image-proper-is-proper}.
\end{proof}

\begin{lemma}
\label{varieties-lemma-projective-after-field-extension}
Soit $k$ un corps. Supposons que $X$ soit un
schéma sur $k$. Si $X_K$ est projectif sur $K$ pour une
certaine extension de corps $K/k$, alors $X$ est projectif sur $k$.
\end{lemma}

\begin{proof}
Un schéma sur $k$ est projectif sur $k$ si et
seulement s'il est quasi-projectif et propre sur $k$.
Voir Morphismes, lemme
\ref{morphisms-lemma-projective-is-quasi-projective-proper}.
Ainsi, le lemme découle des lemmes
\ref{varieties-lemma-quasi-projective-after-field-extension} et \ref{varieties-lemma-proper-after-field-extension}.
\end{proof}




\section{Espaces tangents}
\label{varieties-section-tangent-spaces}

\noindent
Dans cette section, nous définissons l'espace tangent d'un morphisme de schémas
en un point de la source à l'aide de points à valeurs dans les nombres duaux.

\begin{definition}
\label{varieties-definition-dual-numbers}
Pour tout anneau $R$, l'{\it algèbre des nombres duaux} sur $R$
est la $R$-algèbre notée $R[\epsilon]$. Comme $R$-module,
elle est libre de base $1$, $\epsilon$; la structure de
$R$-algèbre est définie par la relation $\epsilon^2 = 0$.
\end{definition}

\noindent
Soit $f : X \to S$ un morphisme de schémas.
Soit $x \in X$ un point d'image $s =
f(x)$ dans $S$. Considérons le diagramme commutatif formé de flèches pleines
\begin{equation}
\label{varieties-equation-tangent-space}
\vcenter{
\xymatrix{
\Spec(\kappa(x)) \ar[r] \ar[dr] \ar@/^1pc/[rr] &
\Spec(\kappa(x)[\epsilon]) \ar@{.>}[r] \ar[d]&
X \ar[d] \\
&
\Spec(\kappa(s)) \ar[r] &
S
}
}
\end{equation}
où la flèche courbe est le morphisme
canonique de $\Spec(\kappa(x))$ vers $X$.

\begin{lemma}
\label{varieties-lemma-tangent-space}
L'ensemble des flèches pointillées rendant commutatif le diagramme
(\ref{varieties-equation-tangent-space}) possède une structure canonique de $\kappa(x)$-espace vectoriel.
\end{lemma}

\begin{proof}
Posons $\kappa = \kappa(x)$. Observons que la somme
amalgamée suivante existe dans la catégorie des schémas :
$$
\Spec(\kappa[\epsilon]) \amalg_{\Spec(\kappa)} \Spec(\kappa[\epsilon])
= \Spec(\kappa[\epsilon_1, \epsilon_2])
$$
où $\kappa[\epsilon_1, \epsilon_2]$ est la $\kappa$-algèbre
libre comme module, de base $1, \epsilon_1, \epsilon_2$, avec
$\epsilon_1^2 = \epsilon_1\epsilon_2 =
\epsilon_2^2 = 0$. Cela découle immédiatement du
résultat correspondant pour les anneaux et de la
description des morphismes du spectre d'un anneau
local vers un schéma; voir Schémas, lemme
\ref{schemes-lemma-morphism-from-spec-local-ring}.
Étant donné deux flèches $\theta_1, \theta_2 :
\Spec(\kappa[\epsilon]) \to X$, nous pouvons considérer le morphisme
$$
\theta_1 + \theta_2 :
\Spec(\kappa[\epsilon]) \to
\Spec(\kappa[\epsilon_1, \epsilon_2])
\xrightarrow{\theta_1, \theta_2} X
$$
où la première flèche est donnée par $\epsilon_i \mapsto \epsilon$.
D'autre part, pour $\lambda \in \kappa$, il existe un endomorphisme de
$\Spec(\kappa[\epsilon])$ correspondant à l'endomorphisme
de $\kappa$-algèbres de $\kappa[\epsilon]$ qui envoie $\epsilon$ sur
$\lambda \epsilon$. La précomposition de $\theta : \Spec(\kappa[\epsilon]) \to X$ par
cet endomorphisme donne $\lambda \theta$. On vérifie
les axiomes d'un espace vectoriel au moyen de diagrammes
commutatifs convenables de schémas. Nous omettons les détails. (Une
autre preuve consisterait à tout exprimer en termes d'anneaux locaux, puis à
vérifier les axiomes de l'espace vectoriel au niveau des homomorphismes d'anneaux.)
\end{proof}

\begin{definition}
\label{varieties-definition-tangent-space}
Soit $f : X \to S$ un morphisme de schémas. Soit $x \in X$.
L'ensemble des flèches pointillées rendant commutatif le diagramme \ref{varieties-equation-tangent-space},
muni de sa structure canonique d'espace vectoriel sur $\kappa(x)$, est
appelé le {\it espace tangent de $X$ sur $S$ en $x$} et nous le notons $T_{X/S,
x}$. Un élément de cet espace est appelé un {\it vecteur tangent} de $X/S$ en $x$.
\end{definition}

\noindent
Puisque les vecteurs tangents en $x \in X$ sont contenus dans la fibre schématique $X_s$ de
$f : X \to S$ au-dessus de $s = f(x)$, nous obtenons une identification canonique
\begin{equation}
\label{varieties-equation-tangent-space-fibre}
T_{X/S, x} = T_{X_s/s, x}
\end{equation}
Cette agréable définition faisant intervenir le
foncteur des points admet la description algébrique
suivante, qui suggère de définir {\it l'espace cotangent de
$X$ sur $S$ en $x$} comme l'espace vectoriel sur $\kappa(x)$
$$
T^*_{X/S, x} = \Omega_{X/S, x} \otimes_{\mathcal{O}_{X, x}} \kappa(x)
$$
puisqu'il est canoniquement le dual, sur
$\kappa(x)$, de l'espace tangent de $X$ sur $S$ en $x$.

\begin{lemma}
\label{varieties-lemma-tangent-space-cotangent-space}
Soit $f : X \to S$ un morphisme de
schémas. Soit $x \in X$. Il existe un isomorphisme canonique
$$
T_{X/S, x} = \Hom_{\mathcal{O}_{X, x}}(\Omega_{X/S, x}, \kappa(x))
$$
d'espaces vectoriels sur $\kappa(x)$.
\end{lemma}

\begin{proof}
Posons $\kappa = \kappa(x)$. Étant donné
$\theta \in T_{X/S, x}$, nous obtenons un homomorphisme
$$
\theta^*\Omega_{X/S} \to
\Omega_{\Spec(\kappa[\epsilon])/\Spec(\kappa(s))} \to
\Omega_{\Spec(\kappa[\epsilon])/\Spec(\kappa)}
$$
En prenant les sections, nous obtenons un morphisme
$\mathcal{O}_{X, x}$-linéaire $\xi_\theta : \Omega_{X/S, x}
\to \kappa \text{d}\epsilon$, c'est-à-dire un élément du
membre de droite de la formule du lemme. Pour montrer que
$\theta \mapsto \xi_\theta$ est un isomorphisme, nous
pouvons remplacer $S$ par $s$ et $X$ par la fibre schématique
$X_s$. En effet, les deux membres de la formule ne dépendent
que de la fibre schématique ; cela est clair pour $T_{X/S,
x}$ et, pour le membre de droite, voir Morphismes, lemme
\ref{morphisms-lemma-base-change-differentials}. Nous pouvons
aussi remplacer $X$ par le spectre de $\mathcal{O}_{X, x}$,
car cela ne change ni $T_{X/S, x}$ (Schémas, lemme
\ref{schemes-lemma-morphism-from-spec-local-ring}) ni $\Omega_{X/S,
x}$ (Modules, lemme \ref{modules-lemma-stalk-module-differentials}).

\medskip\noindent
Soit $(A, \mathfrak m, \kappa)$ un anneau local sur un
corps $k$. Il faut montrer que tout
homomorphisme $A$-linéaire $\xi : \Omega_{A/k} \to \kappa$ provient
d'un unique morphisme de $k$-algèbres $\varphi : A \to
\kappa[\epsilon]$ coïncidant avec le morphisme canonique $c : A \to
\kappa$ modulo $\epsilon$. Écrivons $\varphi(a) = c(a) + D(a) \epsilon$
: on constate que $a \mapsto D(a)$ est une $k$-dérivation. En
utilisant la propriété universelle de $\Omega_{A/k}$, on voit que
chaque $D$ correspond à un $\xi$ unique et vice versa. Cela termine la preuve.
\end{proof}

\begin{lemma}
\label{varieties-lemma-tangent-space-rational-point}
Soit $f : X \to S$ un morphisme de schémas. Soit
$x \in X$ un point et posons $s = f(x) \in S$. Supposons que
$\kappa(x) = \kappa(s)$. Alors il existe des isomorphismes canoniques
$$
\mathfrak m_x/(\mathfrak m_x^2 + \mathfrak m_s\mathcal{O}_{X, x})
=
\Omega_{X/S, x} \otimes_{\mathcal{O}_{X, x}} \kappa(x)
$$
et
$$
T_{X/S, x} =
\Hom_{\kappa(x)}(
\mathfrak m_x/(\mathfrak m_x^2 + \mathfrak m_s\mathcal{O}_{X, x}),
\kappa(x))
$$
Cela fonctionne de manière plus générale si
$\kappa(x)/\kappa(s)$ est une extension algébrique séparable.
\end{lemma}

\begin{proof}
Le deuxième isomorphisme découle du premier par le lemme
\ref{varieties-lemma-tangent-space-cotangent-space}. Pour le premier, nous
pouvons remplacer $S$ par $s$ et $X$ par $X_s$, voir Morphismes,
lemme \ref{morphisms-lemma-base-change-differentials}. Nous
pouvons également remplacer $X$ par le spectre de $\mathcal{O}_{X,
x}$, voir Modules, lemme
\ref{modules-lemma-stalk-module-differentials}. Il suffit donc de démontrer le fait
algébrique suivant : soit $(A, \mathfrak m, \kappa)$ un anneau local sur un corps
$k$ tel que $\kappa/k$ soit algébrique séparable. Alors l'application canonique
$$
\mathfrak m/\mathfrak m^2
\longrightarrow
\Omega_{A/k} \otimes \kappa
$$
est un isomorphisme. Observons que $\mathfrak
m/\mathfrak m^2 = H_1(\NL_{\kappa/A})$. D'après
Algèbre, lemme \ref{algebra-lemma-exact-sequence-NL},
il suffit de montrer que $\Omega_{\kappa/k} = 0$ et
$H_1(\NL_{\kappa/k}) = 0$. Comme $\kappa$ est la réunion des
extensions finies séparables de $k$ qu'il contient, il suffit de
prouver cela lorsque $\kappa$ est une extension séparable
finie de $k$ (Algèbre, lemme
\ref{algebra-lemma-colimits-NL}). Dans ce cas, l'homomorphisme
d'anneaux $k \to \kappa$ est étale et donc $\NL_{\kappa/k} = 0$ (à peu
près par définition, voir Algèbre, section \ref{algebra-section-etale}).
\end{proof}

\begin{lemma}
\label{varieties-lemma-map-tangent-spaces}
Soit $f : X \to Y$ un morphisme de schémas sur un schéma de
base $S$. Soit $x \in X$ un point. Posons $y = f(x)$. Si
$\kappa(y) = \kappa(x)$, alors $f$ induit une application linéaire naturelle
$$
\text{d}f : T_{X/S, x} \longrightarrow T_{Y/S, y}
$$
qui est duale de l'application linéaire $\Omega_{Y/S, y} \otimes
\kappa(y) \to \Omega_{X/S, x} \otimes \kappa(x)$ via les
identifications du lemme \ref{varieties-lemma-tangent-space-cotangent-space}.
\end{lemma}

\begin{proof}
Démonstration omise.
\end{proof}

\begin{lemma}
\label{varieties-lemma-tangent-space-product}
Soient $X$, $Y$ des schémas sur une base $S$. Soient $x \in
X$ et $y \in Y$ avec le même point image $s \in S$ tel que $\kappa(s) =
\kappa(x)$ et $\kappa(s) = \kappa(y)$. Il existe un isomorphisme canonique
$$
T_{X \times_S Y/S, (x, y)} = T_{X/S, x} \oplus T_{Y/S, y}
$$
L'application de gauche à droite est induite par les
applications sur les espaces tangents provenant des
projections $X \times_S Y \to X$ et $X \times_S Y \to Y$. La
application de droite à gauche est induite par les applications $1
\times y : X_s \to X_s \times_s Y_s$ et $x \times 1 : Y_s \to X_s
\times_s Y_s$ via l'identification
(\ref{varieties-equation-tangent-space-fibre}) des espaces tangents avec les espaces tangents des fibres.
\end{lemma}

\begin{proof}
La décomposition en somme directe découle de Morphismes,
lemme \ref{morphisms-lemma-differential-product} via le lemme
\ref{varieties-lemma-tangent-space-rational-point}. La compatibilité avec
les applications provient du lemme \ref{varieties-lemma-map-tangent-spaces}.
\end{proof}

\begin{lemma}
\label{varieties-lemma-injective-tangent-spaces-unramified}
Supposons que $f : X \to Y$ soit un morphisme de schémas localement de type fini sur
un schéma de base $S$. Supposons que $x \in X$ soit un point. Posons $y = f(x)$ et
supposons que $\kappa(y) = \kappa(x)$. Alors les assertions suivantes sont équivalentes
\begin{enumerate}
\item $\text{d}f : T_{X/S, x} \longrightarrow T_{Y/S, y}$ est injective, et
\item $f$ est non ramifié en $x$.
\end{enumerate}
\end{lemma}

\begin{proof}
Le morphisme $f$ est localement de type fini d'après
Morphismes, lemme
\ref{morphisms-lemma-permanence-finite-type}. L'application
$\text{d}f$ est injective si et seulement si $\Omega_{Y/S, y}
\otimes \kappa(y) \to \Omega_{X/S, x} \otimes \kappa(x)$ est
surjective (lemme \ref{varieties-lemma-map-tangent-spaces}). La suite
exacte $f^*\Omega_{Y/S} \to \Omega_{X/S} \to \Omega_{X/Y} \to
0$ (Morphismes, lemme
\ref{morphisms-lemma-triangle-differentials}) montre alors que cela se
produit si et seulement si $\Omega_{X/Y, x} \otimes \kappa(x) = 0$. Le
résultat découle donc de Morphismes, lemme \ref{morphisms-lemma-unramified-at-point}.
\end{proof}







\section{Morphismes génériquement finis}
\label{varieties-section-generically-finite}

\noindent
Dans cette section, nous revenons sur la notion de
morphisme génériquement fini de schémas telle qu'elle est étudiée dans
Morphismes, section \ref{morphisms-section-generically-finite}.

\begin{lemma}
\label{varieties-lemma-quasi-finite-in-codim-1}
Soit $f : X \to Y$ un morphisme localement de type fini. Soit $y
\in Y$ un point tel que $\mathcal{O}_{Y, y}$ soit noethérien de dimension
$\leq 1$. Supposons en outre que l'une des conditions suivantes soit satisfaite :
\begin{enumerate}
\item pour tout point générique $\eta$ d'une composante irréductible
de $X$, l'extension de corps $\kappa(\eta)/\kappa(f(\eta))$
est finie (ou algébrique),
\item pour tout point générique $\eta$ d'une composante irréductible
de $X$ tel que $f(\eta) \leadsto y$, l'extension de corps
$\kappa(\eta)/\kappa(f(\eta))$ est finie (ou algébrique),
\item $f$ est quasi-fini en tout point générique d'une
composante irréductible de $X$,
\item $Y$ est localement noethérien et $f$
est quasi-fini sur une partie dense de $X$,
\item ajouter d'autres conditions ici.
\end{enumerate}
Alors $f$ est quasi-fini en tout point de $X$ au-dessus de $y$.
\end{lemma}

\begin{proof}
La condition (4) implique que $X$ est localement noethérien
(Morphismes, lemme \ref{morphisms-lemma-finite-type-noetherian}).
L'ensemble des points en lesquels le morphisme est quasi-fini est
ouvert (Morphismes, lemme
\ref{morphisms-lemma-quasi-finite-points-open}). Un ouvert dense d'un
schéma localement noethérien contient tous les points génériques des
composantes irréductibles, donc (4) implique (3). La condition (3)
implique la condition (1) d'après Morphismes, lemme
\ref{morphisms-lemma-residue-field-quasi-finite}. La condition (1) implique la condition
(2). Ainsi, il suffit de prouver le lemme dans le cas où la condition (2) est vérifiée.

\medskip\noindent
Supposons (2). Rappelons que
$\Spec(\mathcal{O}_{Y, y})$ est l'ensemble des points de $Y$
se spécialisant en $y$, voir Schémas, lemme
\ref{schemes-lemma-specialize-points}. En combinant avec Morphismes,
lemme \ref{morphisms-lemma-base-change-quasi-finite}, on voit
que nous pouvons remplacer $Y$ par $\Spec(\mathcal{O}_{Y, y})$.
Ainsi nous pouvons supposer $Y = \Spec(B)$ où $B$ est un anneau
local noethérien de dimension $\leq 1$ et $y$ est le point fermé.

\medskip\noindent
Soit $X = \bigcup X_i$ la décomposition en composantes irréductibles
de $X$, munies de leur structure réduite. Si nous pouvons
montrer que chaque fibre $X_{i, y}$ est un espace discret, alors
$X_y = \bigcup X_{i, y}$ est également discret et nous en concluons
que $X \to Y$ est quasi-fini en tout point de $X_y$ d'après
Morphismes, lemme
\ref{morphisms-lemma-quasi-finite-at-point-characterize}. Ainsi on peut supposer que $X$ est un schéma intègre.

\medskip\noindent
Si $X \to Y$ envoie le point générique $\eta$ de $X$ à
$y$, alors $X$ est le spectre d'une extension finie de
$\kappa(y)$ et le résultat est vrai. Supposons que $X$
envoie $\eta$ à un point correspondant à un idéal premier
minimal $\mathfrak q$ de $B$ différent de $\mathfrak m_B$.
Nous obtenons une factorisation $X \to \Spec(B/\mathfrak
q) \to \Spec(B)$. Supposons que $x \in X$ soit un point
au-dessus de $y$. Par la formule de la dimension
(Morphismes, lemme \ref{morphisms-lemma-dimension-formula}) nous avons
$$
\dim(\mathcal{O}_{X, x}) \leq \dim(B/\mathfrak q) +
\text{trdeg}_{\kappa(\mathfrak q)}(R(X)) - \text{trdeg}_{\kappa(y)} \kappa(x)
$$
Nous savons que $\dim(B/\mathfrak q) = 1$, que le point
générique de $X$ n'est pas égal à $x$ et se spécialise en $x$ et que
$R(X)$ est algébrique sur $\kappa(\mathfrak q)$. Ainsi nous obtenons
$$
1 \leq 1 - \text{trdeg}_{\kappa(y)} \kappa(x)
$$
Ainsi, chaque point $x$ de $X_y$ est fermé dans $X_y$ par
Morphismes, lemme
\ref{morphisms-lemma-algebraic-residue-field-extension-closed-point-fibre}
et donc $X \to Y$ est quasi-fini en tout point $x$ de $X_y$ par
Morphismes, lemme
\ref{morphisms-lemma-quasi-finite-at-point-characterize} (ce qui implique également que $X_y$ est un espace topologique discret).
\end{proof}

\begin{lemma}
\label{varieties-lemma-finite-in-codim-1}
Soit $f : X \to Y$ un morphisme propre. Soit $y \in Y$
un point tel que $\mathcal{O}_{Y, y}$ soit noethérien de dimension
$\leq 1$. Supposons en outre que l'une des conditions suivantes soit satisfaite :
\begin{enumerate}
\item pour tout point générique $\eta$ d'une composante irréductible
de $X$, l'extension de corps $\kappa(\eta)/\kappa(f(\eta))$
est finie (ou algébrique),
\item pour tout point générique $\eta$ d'une composante irréductible
de $X$ tel que $f(\eta) \leadsto y$, l'extension de corps
$\kappa(\eta)/\kappa(f(\eta))$ est finie (ou algébrique),
\item $f$ est quasi-fini en tout point générique de $X$,
\item $Y$ est localement noethérien et $f$ est quasi-fini
sur une partie dense de $X$,
\item ajouter d'autres conditions ici.
\end{enumerate}
Alors il existe un voisinage ouvert $V \subset
Y$ de $y$ tel que $f^{-1}(V) \to V$ soit fini.
\end{lemma}

\begin{proof}
D'après le lemme \ref{varieties-lemma-quasi-finite-in-codim-1}, le
morphisme $f$ est quasi-fini en tout point de la fibre
$X_y$. Par conséquent, $X_y$ est un espace topologique
discret (Morphismes, lemme
\ref{morphisms-lemma-quasi-finite-at-point-characterize}). Comme $f$
est propre, la fibre $X_y$ est quasi-compacte, c'est-à-dire finie.
Ainsi, nous pouvons appliquer Cohomologie des schémas, lemme
\ref{coherent-lemma-proper-finite-fibre-finite-in-neighbourhood} pour conclure.
\end{proof}

\begin{lemma}
\label{varieties-lemma-modification-normal-iso-over-codimension-1}
Soit $X$ un schéma noethérien. Soit $f : Y \to X$ un
morphisme propre birationnel, avec $Y$ réduit. Soit $U \subset
X$ le plus grand ouvert sur lequel $f$ est un isomorphisme. Alors $U$ contient
\begin{enumerate}
\item tout point de codimension $0$ dans $X$,
\item tout $x \in X$ de codimension $1$ dans $X$ tel que
$\mathcal{O}_{X, x}$ soit un anneau de valuation discrète,
\item tout $x \in X$ tel que la fibre de $Y \to X$ au-dessus
de $x$ soit finie et que $\mathcal{O}_{X, x}$ soit normal, et
\item tout $x \in X$ tel que $f$ soit quasi-fini en un
$y \in Y$ au-dessus de $x$ et que $\mathcal{O}_{X, x}$ soit normal.
\end{enumerate}
\end{lemma}

\begin{proof}
La partie (1) découle de Morphismes, lemme
\ref{morphisms-lemma-birational-isomorphism-over-dense-open}. La
partie (2) découle de la partie (3) et du lemme
\ref{varieties-lemma-finite-in-codim-1} (et du fait que les morphismes finis ont des fibres finies).

\medskip\noindent
La partie (3) découle de la partie (4) et de Morphismes,
lemme \ref{morphisms-lemma-finite-fibre}, mais nous allons
également donner une preuve directe. Soit $x \in
X$ comme dans (3). D'après Cohomologie des schémas,
lemme
\ref{coherent-lemma-proper-finite-fibre-finite-in-neighbourhood},
on peut supposer que $f$ est fini. Nous pouvons supposer
que $X$ est affine. Cela nous ramène au cas d'un morphisme
birationnel fini de schémas affines noethériens $Y \to X$ et $x
\in X$ tels que $\mathcal{O}_{X, x}$ est un anneau intègre normal.
Comme $\mathcal{O}_{X, x}$ est un anneau intègre et $X$ est
noethérien, nous pouvons remplacer $X$ par un voisinage ouvert affine de
$x$ qui est intègre. Ensuite, puisque $Y \to X$ est birationnel
et $Y$ est réduit, on voit que $Y$ est intègre. Écrivons
$X = \Spec(A)$ et $Y = \Spec(B)$, on voit que $A \subset B$
est une inclusion finie d'anneaux intègres ayant le même corps de
fractions. Si $\mathfrak p \subset A$ est l'idéal premier correspondant
à $x$, alors le fait que $A_\mathfrak p$ soit normal implique que
$A_\mathfrak p \subset B_\mathfrak p$ est une égalité. Puisque $B$
est un $A$-module fini, nous voyons qu'il existe un $a \in A$, $a
\not \in \mathfrak p$ tel que $A_a \to B_a$ est un isomorphisme.

\medskip\noindent
Soient $x \in X$ et $y \in Y$ comme dans (4). Après
avoir remplacé $X$ par un voisinage affine ouvert, nous pouvons
supposer $X = \Spec(A)$ et $A \subset \mathcal{O}_{X, x}$, voir
Propriétés, lemme
\ref{properties-lemma-ring-affine-open-injective-local-ring}.
Ensuite, $A$ est un anneau intègre et donc $X$ est intègre. Puisque $f$
est birationnel et $Y$ est réduit, il s'ensuit que $Y$ est
également intègre. Considérons l'homomorphisme d'anneaux
$\mathcal{O}_{X, x} \to \mathcal{O}_{Y, y}$. Il s'agit d'un
homomorphisme d'anneaux qui est essentiellement de type fini,
l'extension des corps résiduels est finie, et $\dim(\mathcal{O}_{Y,
y}/\mathfrak m_x\mathcal{O}_{Y, y}) = 0$ (pour le voir, remontons
aux définitions des morphismes quasi-finis dans
Morphismes, définition \ref{morphisms-definition-quasi-finite} et
Algèbre, définition \ref{algebra-definition-quasi-finite}). D'après
Algèbre, lemme
\ref{algebra-lemma-essentially-finite-type-fibre-dim-zero},
$\mathcal{O}_{Y, y}$ est la localisation d'une
$\mathcal{O}_{X, x}$-algèbre finie $B$. Bien sûr, nous pouvons remplacer $B$ par
l'image de $B$ dans $\mathcal{O}_{Y, y}$ et supposer que $B$ est un
anneau intègre ayant le même corps des fractions que $\mathcal{O}_{Y, y}$. Comme
Les anneaux de l'inclusion $\mathcal{O}_{X, x} \subset B$ ont le même corps de fractions puisque $f$ est birationnel.
Comme $\mathcal{O}_{X, x}$ est normal, nous concluons que
$\mathcal{O}_{X, x} = B$ (car fini implique intégral), en particulier, nous
voyons que $\mathcal{O}_{X, x} = \mathcal{O}_{Y, y}$. D'après Morphismes,
lemme \ref{morphisms-lemma-morphism-defined-local-ring}, après avoir rétréci
$X$, nous pouvons supposer qu'il existe une section $X \to Y$ de $f$
envoyant $x$ vers $y$ et induisant l'isomorphisme donné sur les anneaux
locaux. Comme $X \to Y$ est une immersion fermée (d'après Schémas, lemme
\ref{schemes-lemma-section-immersion}) et envoie nécessairement le point
générique de $X$ sur le point générique de $Y$, il s'ensuit que l'image
de $X \to Y$ est $Y$. Ainsi $Y = X$, ce qu'il fallait démontrer.
\end{proof}






\section{Variantes de la normalisation de Noether}
\label{varieties-section-noether-normalization}

\noindent
La normalisation de Noether affirme que si $k$ est un corps
et $A$ une $k$-algèbre de type fini de dimension $d$, alors
il existe un homomorphisme injectif et fini de $k$-algèbres
$k[x_1,
\ldots, x_d] \to A$. Voir Algèbre, lemme
\ref{algebra-lemma-Noether-normalization}.
Géométriquement, cela signifie qu'il existe un morphisme
surjectif fini $\Spec(A) \to \mathbf{A}^d_k$ au-dessus de $\Spec(k)$.

\begin{lemma}
\label{varieties-lemma-noether-normalization}
Soit $f : X \to S$ un morphisme de schémas.
Soit $x \in X$, d'image $s \in S$. Soit $V \subset S$
un voisinage ouvert affine de $s$. Si $f$ est localement de type
fini et si $\dim_x(X_s) = d$, alors il existe un ouvert affine $U
\subset X$, avec $x \in U$ et $f(U) \subset V$, et une factorisation
$$
U \xrightarrow{\pi} \mathbf{A}^d_V \to V
$$
de $f|_U : U \to V$ telle que $\pi$ soit quasi-fini.
\end{lemma}

\begin{proof}
Cela résulte d'Algèbre, lemme
\ref{algebra-lemma-quasi-finite-over-polynomial-algebra}.
\end{proof}

\begin{lemma}
\label{varieties-lemma-noether-normalization-affine}
Soit $f : X \to S$ un morphisme de type fini de schémas
affines. Soit $s \in S$. Si $\dim(X_s) = d$, alors il existe une factorisation
$$
X \xrightarrow{\pi} \mathbf{A}^d_S \to S
$$
de $f$ de telle sorte que le morphisme $\pi_s : X_s \to
\mathbf{A}^d_{\kappa(s)}$ des fibres au-dessus de $s$ soit fini.
\end{lemma}

\begin{proof}
Écrivons $S = \Spec(A)$ et $X = \Spec(B)$, et notons $A \to B$
l'homomorphisme d'anneaux correspondant à $f$. Soit
$\mathfrak p \subset A$ l'idéal premier correspondant à $s$.
Nous pouvons choisir une surjection $A[x_1, \ldots, x_r] \to
B$. Par Algèbre, lemme
\ref{algebra-lemma-Noether-normalization}, il existe des éléments
$y_1, \ldots, y_d \in A$ appartenant à la $\mathbf{Z}$-sous-algèbre de $A$
engendrée par $x_1, \ldots, x_r$ tels que l'homomorphisme de $A$-algèbres
$A[t_1, \ldots, t_d] \to B$ envoyant $t_i$ sur $y_i$ induise un
homomorphisme fini de $\kappa(\mathfrak p)$-algèbres $\kappa(\mathfrak
p)[t_1, \ldots, t_d] \to B \otimes_A \kappa(\mathfrak p)$. Ceci démontre le lemme.
\end{proof}

\begin{lemma}
\label{varieties-lemma-geometric-structure-unramified}
Soit $f : X \to S$ un morphisme de schémas.
Soit $x \in X$. Soit $V = \Spec(A)$ un voisinage
ouvert affine de $f(x)$ dans $S$. Si $f$ est non ramifié en
$x$, alors il existe un ouvert affine $U \subset X$ avec $x \in
U$ et $f(U) \subset V$ tel que nous ayons un diagramme commutatif
$$
\xymatrix{
X \ar[d] & U \ar[l] \ar[rd] \ar[r]^-j &
\Spec(A[t]_{g'}/(g)) \ar[d] \ar[r] &
\Spec(A[t]) = \mathbf{A}^1_V \ar[ld] \\
Y & & V \ar[ll]
}
$$
où $j$ est une immersion, $g \in A[t]$ est un polynôme unitaire, et $g'$ est
la dérivée de $g$ par rapport à $t$. Si $f$ est étale en $x$, alors nous
pouvons choisir le diagramme de manière à ce que $j$ soit une immersion ouverte.
\end{lemma}

\begin{proof}
Le cas non ramifié est la traduction géométrique d'Algèbre,
proposition
\ref{algebra-proposition-unramified-locally-standard}. Dans le cas
étale, on traduit de même Algèbre, proposition
\ref{algebra-proposition-etale-locally-standard}; de façon
équivalente, le résultat découle de Morphismes, lemme
\ref{morphisms-lemma-etale-locally-standard-etale} bien que les énoncés diffèrent légèrement.
\end{proof}

\begin{lemma}
\label{varieties-lemma-unramfied-over-affine}
Soit $f : X \to S$ un morphisme de type fini
de schémas affines. Soit $x \in X$, d'image $s \in S$. Posons
$$
r =
\dim_{\kappa(x)} \Omega_{X/S, x} \otimes_{\mathcal{O}_{X, x}} \kappa(x) =
\dim_{\kappa(x)} \Omega_{X_s/s, x} \otimes_{\mathcal{O}_{X_s, x}} \kappa(x) =
\dim_{\kappa(x)} T_{X/S, x}
$$
Alors il existe une factorisation
$$
X \xrightarrow{\pi} \mathbf{A}^r_S \to S
$$
de $f$ de telle sorte que $\pi$ soit non ramifié en $x$.
\end{lemma}

\begin{proof}
D'après Morphismes, lemme
\ref{morphisms-lemma-finite-type-differentials} la
première dimension est finie. La première égalité résulte
du fait que la restriction de $\Omega_{X/S}$ à la fibre
est le module des différentielles, d'après Morphismes, lemme
\ref{morphisms-lemma-base-change-differentials}. La
dernière égalité découle du lemme
\ref{varieties-lemma-tangent-space-cotangent-space}. Ainsi, on voit que l'énoncé a du sens.

\medskip\noindent
Pour démontrer le lemme, écrivons $S = \Spec(A)$ et $X = \Spec(B)$ et
notons $A \to B$ l'homomorphisme d'anneaux correspondant à $f$.
Soit $\mathfrak q \subset B$ l'idéal premier
correspondant à $x$. Choisissons une surjection de $A$-algèbres
$A[x_1, \ldots, x_t] \to B$. Puisque $\Omega_{B/A}$ est engendré
par $\text{d}x_1, \ldots, \text{d}x_t$, on voit que leurs images
dans $\Omega_{X/S, x} \otimes_{\mathcal{O}_{X, x}} \kappa(x)$
engendrent cet espace comme espace vectoriel sur $\kappa(x)$. Après
renumérotation, on peut supposer que $\text{d}x_1, \ldots,
\text{d}x_r$ s'envoient sur une base de $\Omega_{X/S, x}
\otimes_{\mathcal{O}_{X, x}} \kappa(x)$. Nous affirmons que $P = A[x_1,
\ldots, x_r] \to B$ est non ramifié en $\mathfrak q$. Pour le voir, il
suffit de montrer que $\Omega_{B/P, \mathfrak q} = 0$ (Algèbre, lemme
\ref{algebra-lemma-unramified}). Notons que $\Omega_{B/P}$ est le
quotient de $\Omega_{B/A}$ par le sous-module généré par $\text{d}x_1,
\ldots, \text{d}x_r$. Par conséquent, $\Omega_{B/P, \mathfrak q}
\otimes_{B_\mathfrak q} \kappa(\mathfrak q) = 0$ par notre choix de $x_1,
\ldots, x_r$. Par le lemme de Nakayama, plus précisément Algèbre, lemme
\ref{algebra-lemma-NAK} partie (2) qui s'applique car $\Omega_{B/P}$ est fini
(voir référence ci-dessus), nous concluons que $\Omega_{B/P, \mathfrak q} = 0$.
\end{proof}

\begin{lemma}
\label{varieties-lemma-immersion-into-affine}
Soit $f : X \to S$ un morphisme de
schémas. Soit $x \in X$, d'image $s
\in S$. Soit $V \subset S$ un voisinage ouvert affine
ouvert de $s$. Si $f$ est localement de type fini et
$$
r =
\dim_{\kappa(x)} \Omega_{X/S, x} \otimes_{\mathcal{O}_{X, x}} \kappa(x) =
\dim_{\kappa(x)} \Omega_{X_s/s, x} \otimes_{\mathcal{O}_{X_s, x}} \kappa(x) =
\dim_{\kappa(x)} T_{X/S, x}
$$
alors il existe
\begin{enumerate}
\item un ouvert affine $U \subset X$ tel que $x \in U$ et $f(U) \subset V$, et une
factorisation
$$
U \xrightarrow{j} \mathbf{A}^{r + 1}_V \to V
$$
de $f|_U$ tel que $j$ soit une immersion, ou
\item un ouvert affine $U \subset X$ tel que $x \in U$ et $f(U) \subset V$, et une
factorisation
$$
U \xrightarrow{j} D \to V
$$
de $f|_U$ telle que $j$ soit une immersion fermée et $D \to V$
soit lisse de dimension relative $r$.
\end{enumerate}
\end{lemma}

\begin{proof}
Choisissons un ouvert affine quelconque $U \subset X$ avec $x \in U$ et
$f(U) \subset V$. Appliquons le lemme
\ref{varieties-lemma-unramfied-over-affine} à $U \to V$ pour obtenir $U \to
\mathbf{A}^r_V \to V$ comme dans l'énoncé de ce lemme. D'après le lemme
\ref{varieties-lemma-geometric-structure-unramified}, nous obtenons une factorisation
$$
U \xrightarrow{j} D \xrightarrow{j'} \mathbf{A}^{r + 1}_V
\xrightarrow{p} \mathbf{A}^r_V \to V
$$
où $j$ et $j'$ sont des immersions, $p$ est la
projection, et $p \circ j'$ est étale standard.
Ainsi, les assertions (1) et (2) sont vérifiées.
\end{proof}





\section{Dimension des fibres}
\label{varieties-section-dimension-fibres}

\noindent
Nous avons déjà vu que la dimension des fibres des morphismes de type fini peut
augmenter. Dans cette section, nous observons qu'en codimension $1$,
ce phénomène ne se produit pas. Plus généralement, nous étudions l'amplitude
possible de ce saut de dimension. Voici une liste de résultats connexes :
\begin{enumerate}
\item Pour un morphisme de type fini $X \to S$, la partie formée des
$x \in X$ tels que $\dim_x(X_{f(x)}) \leq d$ est ouverte; voir
Algèbre, lemme \ref{algebra-lemma-dimension-fibres-bounded-open-upstairs}
et
Morphismes, lemme \ref{morphisms-lemma-openness-bounded-dimension-fibres}.
\item On dispose de la formule de dimension; voir
Algèbre, lemme \ref{algebra-lemma-dimension-formula} et
Morphismes, lemme \ref{morphisms-lemma-dimension-formula}.
\item Les fibres d'un schéma intègre de type fini dominant
le spectre d'un anneau de valuation sont de dimension constante; voir Algèbre, lemme
\ref{algebra-lemma-finite-type-domain-over-valuation-ring-dim-fibres}.
\item Si $X \to S$ est de type fini et quasi-fini en tout
point générique de $X$, alors $X \to S$ est quasi-fini en codimension $1$; voir
Algèbre, lemme \ref{algebra-lemma-finite-in-codim-1} et
lemme \ref{varieties-lemma-quasi-finite-in-codim-1}.
\end{enumerate}
Le dernier résultat mentionné ci-dessus se généralise comme suit.

\begin{lemma}
\label{varieties-lemma-dimension-fibre-in-codim-1}
Soit $f : X \to Y$ un morphisme localement de type fini. Soit $x \in
X$ un point d'image $y \in Y$ tel que $\mathcal{O}_{Y, y}$ soit
noethérien de dimension $\leq 1$. Soit $d \geq 0$ un entier tel
que pour tout point générique $\eta$ d'une composante irréductible de $X$
contenant $x$, nous ayons $\dim_\eta(X_{f(\eta)}) = d$. Alors $\dim_x(X_y) = d$.
\end{lemma}

\begin{proof}
Rappelons que $\Spec(\mathcal{O}_{Y, y})$ est l'ensemble
des points de $Y$ se spécialisant en $y$, voir Schémas,
lemme \ref{schemes-lemma-specialize-points}. Ainsi, nous
pouvons remplacer $Y$ par $\Spec(\mathcal{O}_{Y, y})$ et
supposer $Y = \Spec(B)$ où $B$ est un anneau local
noethérien de dimension $\leq 1$ et $y$ est le point fermé. Nous
pouvons également remplacer $X$ par un voisinage affine de $x$.

\medskip\noindent
Soit $X = \bigcup X_i$ la décomposition de $X$ en composantes irréductibles,
munies de leur structure réduite. S'il est possible de montrer que
chaque fibre $X_{i, y}$ a dimension $d$, alors $X_y = \bigcup X_{i, y}$ a
également dimension $d$. Ainsi, on peut supposer que $X$ est un schéma intègre.

\medskip\noindent
Si $X \to Y$ envoie le point générique $\eta$ de $X$ sur
$y$, alors $X$ est un schéma sur $\kappa(y)$ et le
résultat est vrai par hypothèse. Supposons que $X$ envoie
$\eta$ sur un point $\xi \in Y$ correspondant à un idéal
premier minimal $\mathfrak q$ de $B$ différent de $\mathfrak
m_B$. Nous obtenons une factorisation $X \to
\Spec(B/\mathfrak q) \to \Spec(B)$. Par la formule de dimension
(Morphismes, lemme \ref{morphisms-lemma-dimension-formula}) nous avons
$$
\dim(\mathcal{O}_{X, x}) + \text{trdeg}_{\kappa(y)} \kappa(x) \leq
\dim(B/\mathfrak q) + \text{trdeg}_{\kappa(\mathfrak q)}(R(X))
$$
Nous avons $\dim(B/\mathfrak q) = 1$. Nous avons
$\text{trdeg}_{\kappa(\mathfrak q)}(R(X)) = d$ par
hypothèse, puisque $\dim_\eta(X_\xi) = d$; voir
Morphismes, lemme
\ref{morphisms-lemma-dimension-fibre-at-a-point}. Puisque
$\mathcal{O}_{X, x} \to \mathcal{O}_{X_s, x}$ a un noyau
(comme $\eta \mapsto \xi \not = y$) et puisque
$\mathcal{O}_{X, x}$ est un anneau intègre noethérien, on voit que
$\dim(\mathcal{O}_{X, x}) > \dim(\mathcal{O}_{X_y, x})$. Nous concluons que
$$
\dim_x(X_s) =
\dim(\mathcal{O}_{X_s, x}) + \text{trdeg}_{\kappa(y)} \kappa(x) \leq d
$$
(Morphismes, lemme
\ref{morphisms-lemma-dimension-fibre-at-a-point}). D'autre part,
nous avons $\dim_x(X_s) \geq \dim_\eta(X_{f(\eta)}) = d$ d'après
Morphismes, lemme \ref{morphisms-lemma-openness-bounded-dimension-fibres}.
\end{proof}

\begin{lemma}
\label{varieties-lemma-dominate-valuation-ring-dimension-fibres}
Soit $f : X \to \Spec(R)$ un morphisme d'un schéma
irréductible vers le spectre d'un anneau de valuation. Si $f$ est
localement de type fini et surjectif, alors la fibre spéciale est
équidimensionnelle de dimension égale à la dimension de la fibre générique.
\end{lemma}

\begin{proof}
Nous pouvons remplacer $X$ par sa réduction car cela ne change
pas la dimension de $X$ ni celle de la fibre spéciale. Alors
$X$ est intègre et le lemme résulte d'Algèbre, lemme
\ref{algebra-lemma-finite-type-domain-over-valuation-ring-dim-fibres}.
\end{proof}

\noindent
Le lemme suivant généralise le lemme \ref{varieties-lemma-dimension-fibre-in-codim-1}.

\begin{lemma}
\label{varieties-lemma-dimension-fibre-in-higher-codimension}
Soit $f : X \to Y$ un morphisme localement de type fini. Soit
$x \in X$ un point d'image $y \in Y$ tel que $\mathcal{O}_{Y,
y}$ soit noethérien. Soit $d \geq 0$ un entier tel que
pour tout point générique $\eta$ d'une composante irréductible de $X$
qui contient $x$, nous ayons $f(\eta) \not = y$ et
$\dim_\eta(X_{f(\eta)}) = d$. Alors $\dim_x(X_y) \leq d + \dim(\mathcal{O}_{Y, y}) - 1$.
\end{lemma}

\begin{proof}
Comme dans la preuve du lemme
\ref{varieties-lemma-dimension-fibre-in-codim-1}, nous réduisons au cas $X =
\Spec(A)$ avec $A$ un anneau intègre et $Y = \Spec(B)$ où $B$ est un anneau
local noethérien dont l'idéal maximal correspond à $y$. Après avoir
remplacé $B$ par $B/\Ker(B \to A)$, on peut supposer que $B$ est un anneau
intègre et que $B \subset A$. Ensuite, nous utilisons la formule de
dimension (Morphismes, lemme \ref{morphisms-lemma-dimension-formula}) pour obtenir
$$
\dim(\mathcal{O}_{X, x}) + \text{trdeg}_{\kappa(y)} \kappa(x) \leq
\dim(B) + \text{trdeg}_B(A)
$$
Nous avons $\text{trdeg}_B(A) = d$ par hypothèse,
puisque $\dim_\eta(X_\xi) = d$; voir
Morphismes, lemme
\ref{morphisms-lemma-dimension-fibre-at-a-point}. Puisque
$\mathcal{O}_{X, x} \to \mathcal{O}_{X_y, x}$ a un noyau
(comme $f(\eta) \not = y$) et puisque $\mathcal{O}_{X, x}$
est un anneau intègre noethérien, on voit que
$\dim(\mathcal{O}_{X, x}) > \dim(\mathcal{O}_{X_y, x})$. Nous concluons que
$$
\dim_x(X_y) =
\dim(\mathcal{O}_{X_y, x}) + \text{trdeg}_{\kappa(y)} \kappa(x)
< \dim(B) + d
$$
(l'égalité vient de Morphismes, lemme
\ref{morphisms-lemma-dimension-fibre-at-a-point}), ce qui donne le résultat.
\end{proof}




\section{Schémas algébriques}
\label{varieties-section-algebraic-schemes}

\noindent
La définition suivante est tirée de
\cite[I, définition 6.4.1]{EGA}.

\begin{definition}
\label{varieties-definition-algebraic-scheme}
Soit $k$ un corps. Un {\it $k$-schéma
algébrique} est un schéma $X$ sur $k$ tel que le morphisme de
structure $X \to \Spec(k)$ soit de type fini. Un {\it $k$-schéma
localement algébrique} est un schéma $X$ sur $k$ tel que le
morphisme de structure $X \to \Spec(k)$ soit localement de type fini.
\end{definition}

\noindent
Notons que tout $k$-schéma (localement) algébrique est (localement)
noethérien, voir Morphismes, lemme
\ref{morphisms-lemma-finite-type-noetherian}. La catégorie des $k$-schémas
algébriques possède tous les produits et produits fibrés (contrairement à la catégorie
des variétés sur $k$). De même pour la catégorie des $k$-schémas localement algébriques.

\begin{lemma}
\label{varieties-lemma-algebraic-scheme-dim-0}
Soit $k$ un corps. Soit $X$ un $k$-schéma
localement algébrique de dimension $0$. Alors $X$ est une réunion
disjointe de spectres de $k$-algèbres locales artiniennes $A$ avec
$\dim_k(A) < \infty$. Si $X$ est un $k$-schéma algébrique de dimension
$0$, alors en plus $X$ est affine et le morphisme $X \to \Spec(k)$ est fini.
\end{lemma}

\begin{proof}
Soit $X$ un $k$-schéma localement algébrique
de dimension $0$. Soit $U = \Spec(A) \subset X$
un sous-schéma ouvert affine. Comme $\dim(X) = 0$, on
a $\dim(A) = 0$. Par normalisation de Noether,
voir Algèbre, lemme
\ref{algebra-lemma-Noether-normalization}, nous voyons qu'il
existe une injection finie $k \to A$, c'est-à-dire $\dim_k(A) <
\infty$. Par conséquent, $A$ est artinien, voir Algèbre, lemme
\ref{algebra-lemma-finite-dimensional-algebra}. Cela implique que
$A = A_1 \times \ldots \times A_r$ est un produit d'un nombre
fini d'anneaux locaux artiniens, voir Algèbre, lemme
\ref{algebra-lemma-artinian-finite-length}. Bien sûr, $\dim_k(A_i) <
\infty$ pour chaque $i$ puisque la somme de ces dimensions égale $\dim_k(A)$.

\medskip\noindent
Les arguments ci-dessus montrent que $X$ possède un recouvrement ouvert
dont les membres sont des espaces topologiques discrets finis. Par
conséquent, $X$ est un espace topologique discret. Il s'ensuit que $X$ est
isomorphe à la réunion disjointe de ses composantes connexes, chacune étant un
espace à un point. Puisqu'un tel schéma est affine, nous concluons (d'après
les résultats du paragraphe précédent) que chacune de ces composantes est le
spectre d'une $k$-algèbre locale artinienne $A$ telle que $\dim_k(A) < \infty$.

\medskip\noindent
Enfin, si $X$ est un $k$-schéma algébrique de
dimension $0$, alors $X$ est quasi-compact donc
est une réunion disjointe finie $X = \Spec(A_1)
\amalg \ldots \amalg \Spec(A_r)$ donc affine (voir
Schémas, lemme
\ref{schemes-lemma-disjoint-union-affines}); nous avons vu la
finitude de $X \to \Spec(k)$ dans le premier paragraphe de la démonstration.
\end{proof}

\noindent
Le lemme suivant rassemble quelques énoncés sur la théorie
de la dimension pour les schémas localement algébriques.

\begin{lemma}
\label{varieties-lemma-dimension-locally-algebraic}
Soit $k$ un corps. Soit $X$ un $k$-schéma localement algébrique.
\begin{enumerate}
\item
\label{varieties-item-catenary}
L'espace topologique sous-jacent à $X$ est caténaire
(Topologie, définition \ref{topology-definition-catenary}).
\item
\label{varieties-item-dimension-at-closed-point}
Pour tout point fermé $x \in X$, on a $\dim_x(X) = \dim(\mathcal{O}_{X, x})$.
\item
\label{varieties-item-dimension-irreducible}
Si $X$ est irréductible, on a $\dim(X) = \dim(U)$ pour tout
ouvert non vide $U \subset X$, et $\dim(X) = \dim_x(X)$
pour tout $x \in X$.
\item
\label{varieties-item-irreducible-maximal-chains}
Si $X$ est irréductible, toute chaîne de parties fermées irréductibles peut être
prolongée en une chaîne maximale, et toutes les chaînes maximales de parties
fermées irréductibles sont de longueur $\dim(X)$.
\item
\label{varieties-item-dimension-irreducibles-passing-through}
Pour $x \in X$, on a
$\dim_x(X) = \max \dim(Z) = \min \dim(\mathcal{O}_{X, x'})$,
où le maximum porte sur les composantes irréductibles
$Z \subset X$ contenant $x$, et le minimum sur les
spécialisations $x \leadsto x'$ telles que $x'$ soit fermé dans $X$.
\item
\label{varieties-item-dimension-irreducible-trdeg}
Si $X$ est irréductible de point générique $x$, alors
$\dim(X) = \text{trdeg}_k(\kappa(x))$.
\item
\label{varieties-item-immediate-specialization}
Si $x \leadsto x'$ est une spécialisation immédiate
de points de $X$, alors
$\text{trdeg}_k(\kappa(x)) = \text{trdeg}_k(\kappa(x')) + 1$.
\item
\label{varieties-item-dimension-sup-trdeg}
La dimension de $X$ est le suprémum des nombres
$\text{trdeg}_k(\kappa(x))$, où $x$ parcourt les
points génériques des composantes irréductibles de $X$.
\item
\label{varieties-item-specialization}
Si $x \leadsto x'$ est une spécialisation non triviale de points de $X$, alors
\begin{enumerate}
\item $\dim_x(X) \leq \dim_{x'}(X)$,
\item $\dim(\mathcal{O}_{X, x}) < \dim(\mathcal{O}_{X, x'})$,
\item $\text{trdeg}_k(\kappa(x)) > \text{trdeg}_k(\kappa(x'))$, et
\item toute chaîne maximale de spécialisations non triviales
$x = x_0 \leadsto x_1 \leadsto \ldots \leadsto x_n = x'$ est de
longueur $n = \text{trdeg}_k(\kappa(x)) - \text{trdeg}_k(\kappa(x'))$.
\end{enumerate}
\item
\label{varieties-item-dimension-formula}
Pour $x \in X$, on a
$\dim_x(X) = \text{trdeg}_k(\kappa(x)) + \dim(\mathcal{O}_{X, x})$.
\item
\label{varieties-item-immediate-specialization-local-ring}
Si $x \leadsto x'$ est une spécialisation immédiate
de points de $X$ et si $X$ est irréductible ou équidimensionnel, alors
$\dim(\mathcal{O}_{X, x'}) = \dim(\mathcal{O}_{X, x}) + 1$.
\end{enumerate}
\end{lemma}

\begin{proof}
Au lieu de nous appuyer sur les résultats plus généraux
démontrés précédemment, nous allons ramener ces énoncés aux
énoncés correspondants sur les $k$-algèbres de type fini et
citer des résultats du chapitre d'algèbre commutative.

\medskip\noindent
Preuve de (\ref{varieties-item-catenary}). Cette propriété est locale sur $X$ d'après
Topologie, lemme \ref{topology-lemma-catenary}. Ainsi, nous
pouvons supposer $X = \Spec(A)$ où $A$ est une $k$-algèbre de
type fini. Nous devons montrer que $A$ est caténaire (Algèbre,
lemme \ref{algebra-lemma-catenary}). Nous pouvons réduire à
$k[x_1, \ldots, x_n]$ en utilisant Algèbre, lemme
\ref{algebra-lemma-quotient-catenary} et ensuite appliquer
Algèbre, lemme
\ref{algebra-lemma-dimension-height-polynomial-ring}. On peut aussi conclure
du fait que $k$ est de Cohen--Macaulay (trivialement) et que les anneaux de Cohen--Macaulay
sont universellement caténaires (Algèbre, lemme \ref{algebra-lemma-CM-ring-catenary}).

\medskip\noindent
Preuve de (\ref{varieties-item-dimension-at-closed-point}). Choisissons un
voisinage affine $U = \Spec(A)$ de $x$. Alors $\dim_x(X) =
\dim_x(U)$. Par conséquent, nous nous réduisons au cas affine, qui est
Algèbre, lemme \ref{algebra-lemma-dimension-closed-point-finite-type-field}.

\medskip\noindent
Preuve de (\ref{varieties-item-dimension-irreducible}). Il suffit de
montrer que deux ouverts affines non vides $U, U' \subset X$
quelconques ont la même dimension (toute chaîne finie de
sous-ensembles irréductibles rencontre un ouvert affine).
Choisissons un point fermé $x$ de $X$ avec $x \in U \cap U'$. Cela
est possible car $X$ est irréductible, donc $U \cap U'$ est non
vide, donc il existe un tel point fermé car $X$ est Jacobson
d'après le lemme \ref{varieties-lemma-locally-finite-type-Jacobson}. Ensuite
$\dim(U) = \dim(\mathcal{O}_{X, x}) = \dim(U')$ d'après Algèbre,
lemme \ref{algebra-lemma-dimension-spell-it-out} (strictement
parlant, il faut remplacer $X$ par sa réduction avant d'appliquer le lemme).

\medskip\noindent
Preuve de (\ref{varieties-item-irreducible-maximal-chains}). Étant
donné une chaîne de sous-ensembles fermés irréductibles, nous
pouvons trouver un ouvert affine $U \subset X$ qui rencontre
le plus petit. Ainsi, l'énoncé résulte d'Algèbre, lemme
\ref{algebra-lemma-dimension-spell-it-out} et $\dim(U) =
\dim(X)$, comme on l'a vu dans (\ref{varieties-item-dimension-irreducible}).

\medskip\noindent
Preuve de
(\ref{varieties-item-dimension-irreducibles-passing-through}).
Choisissons un voisinage affine $U = \Spec(A)$ de $x$. Alors
$\dim_x(X) = \dim_x(U)$. La règle $Z \mapsto Z \cap U$ est une
bijection entre les composantes irréductibles de $X$ passant par
$x$ et les composantes irréductibles de $U$ passant par $x$. De
plus, $\dim(Z \cap U) = \dim(Z)$ pour un tel $Z$ d'après
(\ref{varieties-item-dimension-irreducible}). Ainsi, l'énoncé découle d'Algèbre,
lemme \ref{algebra-lemma-dimension-at-a-point-finite-type-over-field}.

\medskip\noindent
Preuve de (\ref{varieties-item-dimension-irreducible-trdeg}).
D'après (\ref{varieties-item-dimension-irreducible}), cela se
réduit au cas où $X = \Spec(A)$ est affine. Dans ce
cas, cela résulte d'Algèbre, lemme
\ref{algebra-lemma-dimension-prime-polynomial-ring} appliqué à $A_{red}$.

\medskip\noindent
Preuve de (\ref{varieties-item-immediate-specialization}).
Soit $Z = \overline{\{x\}} \supset Z' =
\overline{\{x'\}}$. Il résulte de
(\ref{varieties-item-irreducible-maximal-chains}) que $Z \supset Z'$
est le début d'une chaîne maximale de parties fermées
irréductibles dans $Z$ et par conséquent $\dim(Z) = \dim(Z') +
1$. On conclut par (\ref{varieties-item-dimension-irreducible-trdeg}).

\medskip\noindent
Preuve de (\ref{varieties-item-dimension-sup-trdeg}). Un argument
topologique simple montre que $\dim(X) = \sup \dim(Z)$ où le
suprémum porte sur les composantes irréductibles de $X$ (indication :
utiliser Topologie, lemme \ref{topology-lemma-irreducible}).
Le résultat découle donc de (\ref{varieties-item-dimension-irreducible-trdeg}).

\medskip\noindent
Preuve de (\ref{varieties-item-specialization}). La partie (a) découle du fait
que tout ouvert $U \subset X$ contenant $x'$ contient également $x$. La
partie (b) découle du fait que $\mathcal{O}_{X, x}$ est une localisation
de $\mathcal{O}_{X, x'}$; toute chaîne d'idéaux premiers de
$\mathcal{O}_{X, x}$ correspond donc à une chaîne d'idéaux premiers de $\mathcal{O}_{X,
x'}$, qui peut être prolongée en ajoutant $\mathfrak m_{x'}$ à la fin. Les
parties (c) et (d) découlent formellement de (\ref{varieties-item-immediate-specialization}).

\medskip\noindent
Preuve de (\ref{varieties-item-dimension-formula}). Choisissons un
voisinage affine $U = \Spec(A)$ de $x$. Alors $\dim_x(X) =
\dim_x(U)$. Par conséquent, nous nous réduisons au cas affine, qui est
Algèbre, lemme \ref{algebra-lemma-dimension-at-a-point-finite-type-field}.

\medskip\noindent
Preuve de
(\ref{varieties-item-immediate-specialization-local-ring}). Si $X$ est
équidimensionnel (Topologie, définition
\ref{topology-definition-equidimensional}) alors $\dim(X)$ est égal à la
dimension de chaque composante irréductible de $X$, d'où $\dim_x(X) =
\dim(X) = \dim_{x'}(X)$ d'après
(\ref{varieties-item-dimension-irreducibles-passing-through}). Cela découle donc de (\ref{varieties-item-immediate-specialization}).
\end{proof}

\begin{lemma}
\label{varieties-lemma-dimension-fibres-locally-algebraic}
Soit $k$ un corps. Soit $f : X \to Y$ un morphisme de
$k$-schémas localement algébriques.
\begin{enumerate}
\item Pour $y \in Y$, la fibre $X_y$ est un schéma localement
algébrique sur $\kappa(y)$; tous les résultats du
lemme \ref{varieties-lemma-dimension-locally-algebraic} s'appliquent donc.
\item Supposons $X$ irréductible. Posons $Z = \overline{f(X)}$ et
$d = \dim(X) - \dim(Z)$. Alors
\begin{enumerate}
\item $\dim_x(X_{f(x)}) \geq d$ pour tout $x \in X$,
\item la partie des $x \in X$ tels que $\dim_x(X_{f(x)}) = d$ est ouverte dense,
\item si $\dim(\mathcal{O}_{Z, f(x)}) \geq 1$, alors
$\dim_x(X_{f(x)}) \leq d + \dim(\mathcal{O}_{Z, f(x)}) - 1$,
\item si $\dim(\mathcal{O}_{Z, f(x)}) = 1$, alors $\dim_x(X_{f(x)}) = d$,
\end{enumerate}
\item Pour $x \in X$ avec $y = f(x)$ nous
avons $\dim_x(X_y) \geq \dim_x(X) - \dim_y(Y)$.
\end{enumerate}
\end{lemma}

\begin{proof}
Le morphisme $f$ est localement de type fini d'après
Morphismes, lemme
\ref{morphisms-lemma-permanence-finite-type}. Par conséquent,
le changement de base $X_y \to \Spec(\kappa(y))$ est localement
de type fini. Cela prouve (1). Dans le reste de la
démonstration, nous utiliserons librement les résultats du lemme
\ref{varieties-lemma-dimension-locally-algebraic} pour $X$, $Y$, et les fibres de $f$.

\medskip\noindent
Preuve de (2). Soit $\eta \in X$ le point générique et
posons $\xi = f(\eta)$. Alors $Z = \overline{\{\xi\}}$. Par conséquent
$$
d = \dim(X) - \dim(Z) =
\text{trdeg}_k \kappa(\eta) - \text{trdeg}_k \kappa(\xi) =
\text{trdeg}_{\kappa(\xi)} \kappa(\eta) = \dim_\eta(X_\xi)
$$
Ainsi, les parties (2)(a) et (2)(b) découlent des
Morphismes, lemme
\ref{morphisms-lemma-openness-bounded-dimension-fibres}. Les
parties (2)(c) et (2)(d) découlent des lemmes
\ref{varieties-lemma-dimension-fibre-in-higher-codimension} et \ref{varieties-lemma-dimension-fibre-in-codim-1}.

\medskip\noindent
Preuve de (3). Soit $x \in X$. Soit $X' \subset X$ une
composante irréductible de $X$ passant par $x$ de dimension
$\dim_x(X)$. Alors (2) implique que $\dim_x(X_y) \geq \dim(X') -
\dim(Z')$ où $Z' \subset Y$ est l'adhérence de l'image de $X'$. Cela prouve (3).
\end{proof}

\begin{lemma}
\label{varieties-lemma-dimension-product-locally-algebraic}
\begin{slogan}
La dimension du produit est la somme des dimensions.
\end{slogan}
Soit $k$ un corps. Soient $X$, $Y$ des $k$-schémas localement algébriques.
\begin{enumerate}
\item Pour $z \in X \times Y$ au-dessus de $(x, y)$,
nous avons $\dim_z(X \times Y) = \dim_x(X) + \dim_y(Y)$.
\item Nous avons $\dim(X \times Y) = \dim(X) + \dim(Y)$.
\end{enumerate}
\end{lemma}

\begin{proof}
Preuve de (1). Considérons la factorisation
$$
X \times Y \longrightarrow Y \longrightarrow \Spec(k)
$$
du morphisme de structure. Le premier morphisme $p : X \times Y \to
Y$ est plat par changement de base du morphisme plat $X \to
\Spec(k)$ d'après Morphismes, lemme
\ref{morphisms-lemma-base-change-flat}. De plus, nous avons
$\dim_z(p^{-1}(y)) = \dim_x(X)$ d'après Morphismes, lemme
\ref{morphisms-lemma-dimension-fibre-after-base-change}. Par conséquent $\dim_z(X
\times Y) = \dim_x(X) + \dim_y(Y)$ d'après Morphismes, lemme
\ref{morphisms-lemma-dimension-fibre-at-a-point-additive}. La partie (2) est une conséquence directe de (1).
\end{proof}






\section{Anneaux locaux complets}
\label{varieties-section-complete-local-rings}

\noindent
Quelques résultats sur les anneaux locaux complets des schémas sur un corps.

\begin{lemma}
\label{varieties-lemma-complete-local-ring-structure-as-algebra}
Soit $k$ un corps. Soit $X$ un
schéma localement noethérien sur $k$. Soit $x \in X$
un point de corps résiduel $\kappa$. Il existe un isomorphisme
\begin{equation}
\label{varieties-equation-complete-local-ring}
\kappa[[x_1, \ldots, x_n]]/I \longrightarrow \mathcal{O}_{X, x}^\wedge
\end{equation}
induisant l'identité sur les corps résiduels. En général, on ne
peut choisir (\ref{varieties-equation-complete-local-ring}) comme
isomorphisme de $k$-algèbres. Toutefois, si l'extension
$\kappa/k$ est séparable, on peut choisir
(\ref{varieties-equation-complete-local-ring}) comme isomorphisme de $k$-algèbres.
\end{lemma}

\begin{proof}
L'existence de l'isomorphisme est une conséquence immédiate du
théorème de structure de Cohen\footnote{Notons que si $\kappa$ a
pour caractéristique $p$, alors le théorème dit simplement que nous
obtenons une surjection $\Lambda[[x_1, \ldots, x_n]] \to
\mathcal{O}_{X, x}^\wedge$ où $\Lambda$ est un anneau de Cohen pour
$\kappa$. Mais bien sûr, dans ce cas, l'application se factorise par
$\Lambda/p\Lambda[[x_1, \ldots, x_n]]$ et $\Lambda/p\Lambda = \kappa$.
} (Algèbre, théorème \ref{algebra-theorem-cohen-structure-theorem}).

\medskip\noindent
Soit $p$ un nombre premier impair, posons $k =
\mathbf{F}_p(t)$ et $A = k[x, y]/(y^2 + x^p - t)$. Alors le
complété $A^\wedge$ de $A$ pour l'idéal maximal $\mathfrak m
= (y)$ est isomorphe à $k(t^{1/p})[[z]]$ comme anneau,
mais non comme $k$-algèbre. En effet, $A^\wedge$
ne contient pas d'élément dont la $p$-ième puissance est $t$
(comme le lecteur peut le voir en calculant modulo $y^2$). Cela
montre également qu'aucun isomorphisme
(\ref{varieties-equation-complete-local-ring}) ne peut être un isomorphisme de $k$-algèbres.

\medskip\noindent
Si $\kappa/k$ est séparable, alors il existe un
homomorphisme de $k$-algèbres $\kappa \to \mathcal{O}_{X,
x}^\wedge$ induisant l'identité sur les corps résiduels
d'après Compléments d'algèbre, lemme
\ref{more-algebra-lemma-lift-residue-field}. Soient $f_1,
\ldots, f_n \in \mathfrak m_x$ des générateurs. Considérons l'homomorphisme
$$
\kappa[[x_1, \ldots, x_n]] \longrightarrow \mathcal{O}_{X, x}^\wedge,\quad
x_i \longmapsto f_i
$$
Puisque les deux membres sont complets pour la topologie
$(x_1, \ldots, x_n)$-adique (le membre de droite d'après
Algèbre, lemme
\ref{algebra-lemma-hathat-finitely-generated}), cette
application est surjective d'après Algèbre, lemme
\ref{algebra-lemma-completion-generalities} car elle est
surjective modulo $(x_1, \ldots, x_n)$ par construction.
\end{proof}

\begin{lemma}
\label{varieties-lemma-base-change-complete-local-ring}
Soit $K/k$ une extension de corps. Soit
$X$ un $k$-schéma localement algébrique. Posons $Y = X_K$.
Soit $y \in Y$ un point d'image $x \in X$.
Supposons que $\dim(\mathcal{O}_{X, x}) = \dim(\mathcal{O}_{Y,
y})$ et que $\kappa(x)/k$ soit séparable. Choisissons un isomorphisme
$$
\kappa(x)[[x_1, \ldots, x_n]]/(g_1, \ldots, g_m) \longrightarrow
\mathcal{O}_{X, x}^\wedge
$$
de $k$-algèbres comme dans
(\ref{varieties-equation-complete-local-ring}). On a alors un isomorphisme
$$
\kappa(y)[[x_1, \ldots, x_n]]/(g_1, \ldots, g_m) \longrightarrow
\mathcal{O}_{Y, y}^\wedge
$$
de $K$-algèbres comme dans
(\ref{varieties-equation-complete-local-ring}). On utilise ici $\kappa(x)
\to \kappa(y)$ pour considérer $g_j$ comme une série formelle sur $\kappa(y)$.
\end{lemma}

\begin{proof}
L'homomorphisme local $\mathcal{O}_{X, x}
\to \mathcal{O}_{Y, y}$ induit un homomorphisme
local $\mathcal{O}_{X, x}^\wedge \to
\mathcal{O}_{Y, y}^\wedge$. L'homomorphisme induit
$$
\kappa(x) \to \kappa(x)[[x_1, \ldots, x_n]]/(g_1, \ldots, g_m)
\to \mathcal{O}_{X, x}^\wedge \to \mathcal{O}_{Y, y}^\wedge
$$
composée avec la projection vers $\kappa(y)$ est
l'homomorphisme canonique $\kappa(x) \to \kappa(y)$. D'après
le lemme \ref{varieties-lemma-change-fields-flat}, le corps résiduel
$\kappa(y)$ est une localisation de $\kappa(x) \otimes_k K$
au noyau $\mathfrak p_0$ de $\kappa(x) \otimes_k K \to
\kappa(y)$. D'autre part, d'après le lemme
\ref{varieties-lemma-change-fields-algebraic-unramified}, l'anneau local $(\kappa(x)
\otimes_k K)_{\mathfrak p_0}$ est égal à $\kappa(y)$. Ainsi, l'application
$$
\kappa(x) \otimes_k K \to \mathcal{O}_{Y, y}^\wedge
$$
se factorise canoniquement par
$\kappa(y)$. Nous obtenons un diagramme commutatif
$$
\xymatrix{
\kappa(y) \ar[rr] & & \mathcal{O}_{Y, y}^\wedge \\
\kappa(x) \ar[r] \ar[u] &
\kappa(x)[[x_1, \ldots, x_n]]/(g_1, \ldots, g_m) \ar[r] &
\mathcal{O}_{X, x}^\wedge \ar[u]
}
$$
Soit $f_i \in \mathfrak m_x^\wedge \subset
\mathcal{O}_{X, x}^\wedge$ soit l'image de $x_i$.
Observons que $\mathfrak m_x^\wedge = (f_1, \ldots, f_n)$ puisque
l'homomorphisme est surjectif. Considérons l'homomorphisme
$$
\kappa(y)[[x_1, \ldots, x_n]] \longrightarrow \mathcal{O}_{Y, y}^\wedge,\quad
x_i \longmapsto f_i
$$
où ici $f_i$ signifie vraiment l'image de $f_i$ dans
$\mathfrak m_y^\wedge$. Puisque $\mathfrak m_x
\mathcal{O}_{Y, y} = \mathfrak m_y$ d'après le lemme
\ref{varieties-lemma-change-fields-algebraic-unramified}, nous voyons
que le membre de droite est complet pour la topologie $(x_1, \ldots,
x_n)$-adique (on utilise Algèbre, lemme
\ref{algebra-lemma-hathat-finitely-generated} pour voir que le membre
de droite est un anneau local complet). Puisque les deux membres sont complets pour
la topologie $(x_1, \ldots, x_n)$-adique, notre application est
surjective par Algèbre, lemme
\ref{algebra-lemma-completion-generalities} car elle est surjective modulo
$(x_1, \ldots, x_n)$. Bien sûr, les séries formelles $g_1, \ldots, g_m$ sont
envoyées sur zéro par cet homomorphisme, puisqu'elles le sont déjà
dans $\mathcal{O}_{X, x}^\wedge$. On obtient donc le diagramme commutatif
$$
\xymatrix{
\kappa(y)[[x_1, \ldots, x_n]]/(g_1, \ldots, g_m) \ar[r] &
\mathcal{O}_{Y, y}^\wedge \\
\kappa(x)[[x_1, \ldots, x_n]]/(g_1, \ldots, g_m) \ar[r] \ar[u] &
\mathcal{O}_{X, x}^\wedge \ar[u]
}
$$
Nous devons encore montrer que la flèche horizontale supérieure
est un isomorphisme. Nous savons déjà qu'elle est surjective.
Nous savons que $\mathcal{O}_{X, x} \to \mathcal{O}_{Y, y}$ est
plat (lemme \ref{varieties-lemma-change-fields-flat}), ce qui implique
que $\mathcal{O}_{X, x}^\wedge \to \mathcal{O}_{Y, y}^\wedge$
est plat (Compléments d'algèbre, lemme
\ref{more-algebra-lemma-flat-completion}). Ainsi, nous pouvons
appliquer Algèbre, lemme \ref{algebra-lemma-mod-injective} avec $R =
\kappa(x)[[x_1, \ldots, x_n]]/(g_1, \ldots, g_m)$, avec $S =
\kappa(y)[[x_1, \ldots, x_n]]/(g_1, \ldots, g_m)$, avec $M =
\mathcal{O}_{Y, y}^\wedge$ et avec $N = S$ pour conclure que l'application est injective.
\end{proof}


\section{Engendrement par les sections globales}
\label{varieties-section-global-generation}

\noindent
Quelques lemmes sur l'engendrement des modules quasi-cohérents par leurs sections globales.

\begin{lemma}
\label{varieties-lemma-globally-generated-base-change}
Soit $X \to \Spec(A)$ un morphisme de schémas. Soit $A \subset
A'$ une inclusion d'anneaux fidèlement plate. Soit $\mathcal{F}$
un $\mathcal{O}_X$-module quasi-cohérent. Alors $\mathcal{F}$ est engendré par ses sections globales
si et seulement si $\mathcal{F}_{A'}$ est engendré par ses sections globales.
\end{lemma}

\begin{proof}
Plus précisément, posons $X_{A'} = X \times_{\Spec(A)}
\Spec(A')$. Soit $\mathcal{F}_{A'} = p^*\mathcal{F}$ où $p :
X_{A'} \to X$ est la projection. D'après Cohomologie des
schémas, lemme
\ref{coherent-lemma-flat-base-change-cohomology} nous avons
$H^0(X_{k'}, \mathcal{F}_{A'}) = H^0(X, \mathcal{F}) \otimes_A
A'$. Ainsi, si $s_i$, $i \in I$ sont des générateurs de $H^0(X,
\mathcal{F})$ comme module $A$, alors leurs images dans
$H^0(X_{A'}, \mathcal{F}_{A'})$ sont des générateurs de
$H^0(X_{A'}, \mathcal{F}_{A'})$ comme module $A'$. Ainsi, nous
devons montrer que l'application $\alpha : \bigoplus_{i \in I}
\mathcal{O}_X \to \mathcal{F}$, $(f_i) \mapsto \sum f_i s_i$ est
surjective si et seulement si $p^*\alpha$ est surjective. Cela,
nous pouvons le vérifier sur un ouvert affine $U = \Spec(B)$ de
$X$. Alors $\mathcal{F}|_U$ correspond à un $B$-module $M$ et
$s_i|_U$ aux éléments $x_i \in M$. Ainsi, nous devons montrer que
$\bigoplus_{i \in I} B \to M$ est surjective si et seulement si le
changement de base $\bigoplus_{i \in I} B \otimes_A A' \to M \otimes_A
A'$ est surjective. C'est vrai parce que $A \to A'$ est fidèlement plat.
\end{proof}

\begin{lemma}
\label{varieties-lemma-very-ample-vanish-at-point}
Soit $k$ un corps infini. Soit $X$
un schéma de type fini sur $k$. Soit $\mathcal{L}$
un faisceau inversible très ample sur $X$. Soient
$n \geq 0$ et $x, x_1, \ldots, x_n \in X$, avec
$x$ un point $k$-rationnel, c'est-à-dire $\kappa(x) = k$,
et $x \not = x_i$ pour $i = 1, \ldots, n$. Alors il existe un $s
\in H^0(X, \mathcal{L})$ qui s'annule en $x$ mais en aucun des $x_i$.
\end{lemma}

\begin{proof}
Si $n = 0$ le résultat est trivial, nous supposons donc $n >
0$. Par définition d'un faisceau inversible très ample, le
lemme se réduit immédiatement au cas où $X = \mathbf{P}^r_k$
pour un certain $r > 0$ et $\mathcal{L} = \mathcal{O}_X(1)$.
Écrivons $\mathbf{P}^r_k = \text{Proj}(k[T_0, \ldots, T_r])$.
Posons $V = H^0(X, \mathcal{L}) = kT_0 \oplus \ldots \oplus
kT_r$. Puisque $x$ est un point $k$-rationnel, on voit que
l'ensemble des $s \in V$ qui s'annulent en $x$ est un sous-espace
de codimension $1$, à savoir $W \subset V$, et que $W$ engendre l'idéal premier
homogène correspondant à $x$. Puisque $x_i \not = x$, l'idéal premier
homogène correspondant $\mathfrak p_i \subset k[T_0, \ldots, T_r]$
ne contient pas $W$. Comme $k$ est infini, nous voyons alors que $W
\not = \bigcup W \cap \mathfrak q_i$ et la démonstration est complète.
\end{proof}

\begin{lemma}
\label{varieties-lemma-generated-by-dim-plus-1-sections}
Soit $k$ un corps infini. Soit $X$ un $k$-schéma
algébrique. Soit $\mathcal{L}$ un $\mathcal{O}_X$-module
inversible. Soit $V \to \Gamma(X, \mathcal{L})$ une application
linéaire de $k$-espaces vectoriels dont l'image génère $\mathcal{L}$. Alors il existe
un sous-espace $W \subset V$ avec $\dim_k(W) \leq \dim(X) + 1$ qui génère $\mathcal{L}$.
\end{lemma}

\begin{proof}
Tout au long de la démonstration, nous utiliserons
que pour tout $x \in X$, l'application linéaire
$$
\psi_x : V \to \Gamma(X, \mathcal{L}) \to \mathcal{L}_x \to
\mathcal{L}_x \otimes_{\mathcal{O}_{X, x}} \kappa(x)
$$
n'est pas nulle. La démonstration se fait par récurrence sur $\dim(X)$.

\medskip\noindent
Le cas de base est $\dim(X) = 0$. Dans ce cas, $X$ a un
nombre fini de points $X = \{x_1, \ldots, x_n\}$ (voir par
exemple le lemme \ref{varieties-lemma-algebraic-scheme-dim-0}). Puisque
$k$ est infini, il existe un vecteur $v \in V$ tel que
$\psi_{x_i}(v) \not = 0$ pour tout $i$. Alors $W = k\cdot v$ convient.

\medskip\noindent
Supposons $\dim(X) > 0$. Soient $X_i \subset X$ les composantes
irréductibles de dimension égale à $\dim(X)$. Comme $X$
est noethérien, elles sont en nombre fini. Pour tout $i$,
choisissons un point $x_i \in X_i$. Comme ci-dessus, choisissons $v
\in V$ tel que $\psi_{x_i}(v) \not = 0$ pour tout $i$. Soit
$Z \subset X$ le schéma des zéros de l'image de $v$ dans $\Gamma(X,
\mathcal{L})$, voir Diviseurs, définition
\ref{divisors-definition-zero-scheme-s}. Par construction $\dim(Z) < \dim(X)$.
Par induction, nous pouvons trouver $W \subset V$ avec $\dim(W) \leq \dim(X)$
tel que $W$ génère $\mathcal{L}|_Z$. Alors $W + k\cdot v$ génère $\mathcal{L}$.
\end{proof}






\section{Séparation des points et des vecteurs tangents}
\label{varieties-section-separating-points-tangent-vectors}

\noindent
C'est simplement le résultat suivant.

\begin{lemma}
\label{varieties-lemma-separate-points-tangent-vectors}
Soit $k$ un corps algébriquement clos. Soit
$X$ un $k$-schéma propre. Soit $\mathcal{L}$
un $\mathcal{O}_X$-module inversible. Soit $V
\subset H^0(X, \mathcal{L})$ soit un $k$-sous-espace vectoriel. Si
\begin{enumerate}
\item pour toute paire de points fermés distincts $x, y \in X$, il
existe une section $s \in V$ qui s'annule en $x$ mais non en $y$, et
\item pour tout point fermé $x \in X$ et tout vecteur tangent non
nul $\theta \in T_{X/k, x}$, il existe une section $s \in V$ qui
s'annule en $x$ mais dont l'image réciproque par $\theta$ est non nulle,
\end{enumerate}
alors $\mathcal{L}$ est très ample et le
morphisme canonique $\varphi_{\mathcal{L}, V} :
X \to \mathbf{P}(V)$ est une immersion fermée.
\end{lemma}

\begin{proof}
La condition (1) implique en particulier que les éléments de $V$ engendrent
$\mathcal{L}$ sur $X$. On obtient donc un morphisme canonique
$$
\varphi = \varphi_{\mathcal{L}, V} :
X
\longrightarrow
\mathbf{P}(V)
$$
par Constructions, exemple
\ref{constructions-example-projective-space}. Le morphisme
$\varphi$ est propre d'après Morphismes, lemme
\ref{morphisms-lemma-image-proper-scheme-closed}. D'après (1),
le morphisme $\varphi$ est injectif sur les points fermés (calcul
omis). En particulier, la fibre au-dessus de tout point fermé de
$\mathbf{P}(V)$ est réduite à un point (petit détail omis). Ainsi, nous voyons
que $\varphi$ est fini, par exemple en utilisant Cohomologie des
schémas, lemme
\ref{coherent-lemma-proper-finite-fibre-finite-in-neighbourhood}. Pour terminer la démonstration, il suffit de montrer que l'application
$$
\varphi^\sharp :
\mathcal{O}_{\mathbf{P}(V)}
\longrightarrow
\varphi_*\mathcal{O}_X
$$
est surjective. Nous pouvons le vérifier sur les germes aux
points fermés. Soit $x \in X$ un point fermé d'image le point
fermé $p = \varphi(x) \in \mathbf{P}(V)$. Puisque
$\varphi^{-1}(\{p\}) = \{x\}$ par (1) et que $\varphi$ est
propre (donc fermé), on voit que $\varphi^{-1}(U)$ décrit un
système fondamental de voisinages ouverts de $x$ lorsque $U$
décrit un système fondamental de voisinages ouverts de $p$. Nous
concluons qu'en passant aux germes en $p$, nous obtenons l'homomorphisme
$$
\varphi^\sharp_x :
\mathcal{O}_{\mathbf{P}(V), p}
\longrightarrow
\mathcal{O}_{X, x}
$$
En particulier, $\mathcal{O}_{X, x}$ est un
$\mathcal{O}_{\mathbf{P}(V), p}$-module fini. De plus, les
corps résiduels de $x$ et $p$ sont égaux à $k$ (comme $k$
est algébriquement clos --- on utilise le Nullstellensatz de
Hilbert). Enfin, la condition (2) implique que l'application
$$
T_{X/k, x} \longrightarrow T_{\mathbf{P}(V)/k, p}
$$
est injective puisque tout $\theta$ non nul dans le noyau de cette
application ne saurait satisfaire la conclusion de (2). En
termes du morphisme d'anneaux locaux ci-dessus, cela signifie que
$$
\mathfrak m_p/\mathfrak m_p^2 \longrightarrow
\mathfrak m_x/\mathfrak m_x^2
$$
est surjective, voir le lemme
\ref{varieties-lemma-tangent-space-rational-point}. La preuve s'achève
alors par Algèbre, lemme \ref{algebra-lemma-when-surjective-local}.
\end{proof}

\begin{lemma}
\label{varieties-lemma-variant-separate-points-tangent-vectors}
Soit $k$ un corps algébriquement clos.
Soit $X$ un $k$-schéma propre. Soit
$\mathcal{L}$ un $\mathcal{O}_X$-module
inversible. Supposons que, pour tout sous-schéma fermé $Z
\subset X$ de dimension $0$ et de degré $2$ sur $k$, l'application
$$
H^0(X, \mathcal{L}) \longrightarrow H^0(Z, \mathcal{L}|_Z)
$$
est surjective. Alors $\mathcal{L}$ est très ample sur $X$ relativement à $k$.
\end{lemma}

\begin{proof}
C'est une reformulation du lemme
\ref{varieties-lemma-separate-points-tangent-vectors}. En effet,
étant donnés deux points fermés distincts $x, y
\in X$, le choix $Z = x \cup y$ (considéré comme
sous-schéma fermé) donne la condition (1) du
lemme. Étant donné un vecteur tangent non nul
$\theta \in T_{X/k, x}$ le morphisme $\theta :
\Spec(k[\epsilon]) \to X$ est une immersion fermée. En
posant $Z = \Im(\theta)$ nous obtenons la condition (2) du lemme.
\end{proof}









\section{Adhérences de produits}
\label{varieties-section-closure-of-products}

\noindent
Quelques résultats sur les rapports entre l'adhérence et les produits.

\begin{lemma}
\label{varieties-lemma-closure-of-product}
Soit $k$ un corps.
Soient $X$, $Y$ des schémas sur $k$, et soient
$A \subset X$, $B \subset Y$ des
parties. Posons
$$
AB =
\{z \in X \times_k Y \mid \text{pr}_X(z) \in A, \ \text{pr}_Y(z) \in B\}
\subset X \times_k Y
$$
Alors, ensemblistement, nous avons
$$
\overline{A} \times_k \overline{B} = \overline{AB}
$$
\end{lemma}

\begin{proof}
L'inclusion $\overline{AB} \subset \overline{A}
\times_k \overline{B}$ est immédiate. Nous pouvons
remplacer $X$ et $Y$ par les sous-schémas fermés réduits
$\overline{A}$ et $\overline{B}$. Soit $W
\subset X \times_k Y$ soit un ouvert non vide. Par
Morphismes, lemme
\ref{morphisms-lemma-scheme-over-field-universally-open}, le
sous-ensemble $U = \text{pr}_X(W)$ est ouvert non vide dans
$X$. Ainsi $A \cap U$ n'est pas vide. Choisissons $a \in A \cap
U$. Notons $Y_a = \{a\} \times_k Y \cong
Y_{\kappa(a)}$ la fibre de $\text{pr}_X : X \times_k Y \to X$
au-dessus de $a$. D'après Morphismes, lemme
\ref{morphisms-lemma-scheme-over-field-universally-open}, le morphisme
$Y_a \to Y$ est encore ouvert puisque
$\Spec(\kappa(a)) \to \Spec(k)$ est universellement ouvert. Ainsi
l'ouvert non vide $W_a = W \times_{X \times_k Y} Y_a$ se projette
sur un ouvert non vide de $Y$. Nous en concluons qu'il existe un $b
\in B$ dans l'image. Ainsi $AB \cap W \not = \emptyset$ comme voulu.
\end{proof}

\begin{lemma}
\label{varieties-lemma-closure-image-product-map}
Soit $k$ un corps. Soient $f : A
\to X$, $g : B \to Y$ des morphismes de schémas
sur $k$. Alors, ensemblistement, on a
$$
\overline{f(A)} \times_k \overline{g(B)} =
\overline{(f \times g)(A \times_k B)}
$$
\end{lemma}

\begin{proof}
Cela découle du lemme
\ref{varieties-lemma-closure-of-product}
puisque l'image de $f \times g$ est
$f(A)g(B)$ dans la notation de ce lemme.
\end{proof}

\begin{lemma}
\label{varieties-lemma-scheme-theoretic-image-product-map}
Soit $k$ un corps. Soient $f : A \to X$, $g
: B \to Y$ des morphismes quasi-compacts de schémas sur
$k$. Soit $Z \subset X$ l'image schématique de $f$;
voir Morphismes, définition
\ref{morphisms-definition-scheme-theoretic-image}. De même, soit $Z' \subset Y$
l'image schématique de $g$. Alors $Z \times_k Z'$ est l'image schématique de $f \times g$.
\end{lemma}

\begin{proof}
Rappelons que $Z$ est le plus petit sous-schéma fermé de $X$ par
lequel $f$ se factorise. De même pour $Z'$. Soit $W \subset X
\times_k Y$ l'image schématique de $f \times g$. Comme $f \times g$ se
factorise par $Z \times_k Z'$, on voit que $W \subset Z \times_k Z'$.

\medskip\noindent
Pour démontrer l'autre inclusion, soient $U \subset X$ et $V
\subset Y$ des ouverts affines. D'après Morphismes, lemme
\ref{morphisms-lemma-quasi-compact-scheme-theoretic-image},
le schéma $Z \cap U$ est l'image schématique de
$f|_{f^{-1}(U)} : f^{-1}(U) \to U$, et de même pour $Z'
\cap V$ et $W \cap U \times_k V$. On peut donc
supposer que $X$ et $Y$ sont affines. Comme $f$ et
$g$ sont quasi-compacts, cela implique que $A = \bigcup
U_i$ est une union finie d'affines et $B = \bigcup V_j$ est
une union finie d'affines. On peut alors remplacer
$A$ par $\coprod U_i$ et $B$ par $\coprod V_j$,
c'est-à-dire qu'on peut supposer que $A$ et $B$ sont
également affines. Dans ce cas, $Z$ est défini par
$\Ker(\Gamma(X, \mathcal{O}_X) \to \Gamma(A, \mathcal{O}_A))$ et de
même pour $Z'$ et $W$. Par conséquent, le résultat découle de l'égalité
$$
\Gamma(A \times_k B, \mathcal{O}_{A \times_k B})
=
\Gamma(A, \mathcal{O}_A) \otimes_k \Gamma(B, \mathcal{O}_B)
$$
qui est vraie puisque $A$ et $B$ sont affines. Nous omettons les détails.
\end{proof}





\section{Schémas lisses sur des corps} 
\label{varieties-section-smooth}

\noindent
 Voici deux lemmes caractérisant les schémas lisses sur les corps.

\begin{lemma}
\label{varieties-lemma-char-zero-differentials-free-smooth}
 Soit $k$ un corps. Soit $X$ un schéma sur $k$.
Supposons 
\begin{enumerate}
\item $X$ est localement de type fini sur $k$, 
\item $\Omega_{X/k}$ est localement libre, et 
\item $k$ a caractéristique zéro. 
\end{enumerate}
 Alors le morphisme de structure $X \to \Spec(k)$ est lisse. 
\end{lemma}

\begin{proof}
 Cela résulte
d'Algèbre, lemme \ref{algebra-lemma-characteristic-zero-local-smooth}. 
\end{proof}

\noindent
 En caractéristique positive, il existe des schémas de type fini
non réduits dont le faisceau des différentielles est libre, par exemple 
$\Spec(\mathbf{F}_p[t]/(t^p))$ sur $\Spec(\mathbf{F}_p)$.
Si le corps de base $k$ est non parfait de caractéristique $p$,
il existe des schémas réduits $X/k$ avec $\Omega_{X/k}$ libre qui
ne sont pas lisses, par exemple $\Spec(k[t]/(t^p-a)$ où $a \in k$
 n'est pas une puissance $p$.

\begin{lemma}
\label{varieties-lemma-char-p-differentials-free-smooth}
 Soit $k$ un corps. Soit $X$ un schéma sur $k$.
Supposons 
\begin{enumerate}
\item $X$ est localement de type fini sur $k$, 
\item $\Omega_{X/k}$ est localement libre, 
\item $X$ est réduit, et 
\item $k$ est parfait. 
\end{enumerate}
 Alors le morphisme de structure $X \to \Spec(k)$ est lisse. 
\end{lemma}

\begin{proof}
 Soit $x \in X$ un point. Comme $X$ est localement noethérien (voir
Morphismes, lemme \ref{morphisms-lemma-finite-type-noetherian}),
il y a un nombre fini de composantes irréductibles 
$X_1, \ldots, X_n$ passant par $x$ (voir
Propriétés, lemme \ref{properties-lemma-Noetherian-topology} et
Topologie, lemme \ref{topology-lemma-Noetherian}).
Soit $\eta_i \in X_i$ le point générique. Comme $X$ est réduit, on a 
$\mathcal{O}_{X, \eta_i} = \kappa(\eta_i)$,
voir Algèbre, lemme \ref{algebra-lemma-minimal-prime-reduced-ring}.
De plus, $\kappa(\eta_i)$ est une extension de corps de type
fini du corps parfait $k$, donc séparablement engendrée sur $k$ (voir
Algèbre, section \ref{algebra-section-separability}).
Il s'ensuit que $\Omega_{X/k, \eta_i} = \Omega_{\kappa(\eta_i)/k}$
 est libre de rang égal au degré de transcendance de $\kappa(\eta_i)$ sur $k$.
D'après Morphismes, lemme \ref{morphisms-lemma-dimension-fibre-at-a-point},
nous concluons que $\dim_{\eta_i}(X_i) = \text{rank}_{\eta_i}(\Omega_{X/k})$.
Puisque $x \in X_1 \cap \ldots \cap X_n$, on voit que 
$$
\text{rank}_x(\Omega_{X/k}) = \text{rank}_{\eta_i}(\Omega_{X/k}) = \dim(X_i).
$$
 Par conséquent $\dim_x(X) = \text{rank}_x(\Omega_{X/k})$,
voir Algèbre, lemme \ref{algebra-lemma-dimension-at-a-point-finite-type-over-field}.
Il s'ensuit que $X \to \Spec(k)$ est lisse en $x$ par exemple
d'après Algèbre, lemme \ref{algebra-lemma-characterize-smooth-over-field}. 
\end{proof}

\begin{lemma}
\label{varieties-lemma-smooth-regular}
\begin{slogan}
 Être lisse sur un corps implique être régulier 
\end{slogan}
 Soit $X \to \Spec(k)$ un morphisme lisse, où $k$ est un corps.
Alors $X$ est un schéma régulier. 
\end{lemma}

\begin{proof}
 (Voir aussi
le lemme \ref{varieties-lemma-geometrically-regular-smooth}.)

D'après Algèbre, lemme \ref{algebra-lemma-characterize-smooth-over-field},
tout anneau local $\mathcal{O}_{X, x}$ est régulier.
Comme $X$ est localement de type fini sur $k$, il est
localement noethérien. Par conséquent, $X$ est régulier
d'après Propriétés, lemme \ref{properties-lemma-characterize-regular}. 
\end{proof}

\begin{lemma}
\label{varieties-lemma-smooth-geometrically-normal}
 Soit $X \to \Spec(k)$ un morphisme lisse, où $k$ est un corps.
Alors $X$ est géométriquement régulier, géométriquement normal et
géométriquement réduit sur $k$. 
\end{lemma}

\begin{proof}
 (Voir aussi
le lemme \ref{varieties-lemma-geometrically-regular-smooth}.)
Soit $k'$ une extension finie purement inséparable de $k$.
Il suffit de montrer que $X_{k'}$ est régulier, normal et réduit ; voir les
lemmes \ref{varieties-lemma-geometrically-regular}, 
\ref{varieties-lemma-geometrically-normal}, et 
\ref{varieties-lemma-check-only-finite-inseparable-extensions}.

D'après Morphismes, lemme \ref{morphisms-lemma-base-change-smooth},
le morphisme $X_{k'} \to \Spec(k')$ est également lisse.
Par conséquent, il suffit de montrer qu'un schéma $X$ lisse sur un corps est régulier,
normal et réduit. On voit que $X$ est régulier d'après
le lemme \ref{varieties-lemma-smooth-regular}.

Par conséquent, Propriétés, lemme \ref{properties-lemma-regular-normal}
 garantit que $X$ est normal. 
\end{proof}

\begin{lemma}
\label{varieties-lemma-affine-space-over-field}
 Soit $k$ un corps, soit $d \geq 0$ et soit $W \subset \mathbf{A}^d_k$
 ouvert non vide. Alors il existe un point fermé $w \in W$ tel que 
$k \subset \kappa(w)$ soit finie et séparable. 
\end{lemma}

\begin{proof}
 Après avoir éventuellement rétréci $W$, on peut supposer que 
$W = \mathbf{A}^d_k \setminus V(f)$ pour certains $f \in k[x_1, \ldots, x_d]$.
Si le lemme est faux, alors $f(a_1, \ldots, a_d) = 0$ pour tous les 
$(a_1, \ldots, a_d) \in (k^{sep})^d$. C'est absurde car $k^{sep}$
 est un corps infini. 
\end{proof}

\begin{lemma}
\label{varieties-lemma-smooth-separable-closed-points-dense}
 Soit $k$ un corps. Si $X$ est lisse sur $\Spec(k)$ alors
l'ensemble 
$$
\{x \in X\text{ fermé tel que }k \subset \kappa(x)
\text{ est finie séparable}\}
$$
 est dense en $X$. 
\end{lemma}

\begin{proof}
 Il suffit de montrer que, pour tout schéma lisse non vide $X$ sur $k$,
il existe au moins un point fermé dont le corps résiduel
est fini et séparable sur $k$. Pour le voir, choisissons un diagramme 
$$
\xymatrix{
X & U \ar[l] \ar[r]^-\pi & \mathbf{A}^d_k
}
$$
 où $\pi$ est étale, voir
Morphismes, lemme \ref{morphisms-lemma-smooth-etale-over-affine-space}.
Le morphisme $\pi : U \to \mathbf{A}^d_k$ est ouvert,
voir Morphismes, lemme \ref{morphisms-lemma-etale-open}.
D'après
le lemme \ref{varieties-lemma-affine-space-over-field},
nous pouvons choisir un point fermé $w \in \pi(U)$ dont le corps résiduel
est fini et séparable sur $k$. Choisissons un $x \in U$ avec $\pi(x) = w$.
D'après Morphismes, lemme \ref{morphisms-lemma-etale-over-field},
l'extension de corps $\kappa(x)/\kappa(w)$ est finie et séparable. Par
conséquent, $\kappa(x)/k$ est finie et séparable. Le point $x$ est
un point fermé de $X$
 d'après Morphismes, lemme 
\ref{morphisms-lemma-algebraic-residue-field-extension-closed-point-fibre}. 
\end{proof}

\begin{lemma}
\label{varieties-lemma-geometrically-reduced-dense-smooth-open}
 Soit $X$ un schéma réduit localement de type fini sur un corps $k$.
Alors $X$ est géométriquement réduit sur $k$ si et seulement si 
$X$ contient un ouvert dense qui est lisse sur $k$. 
\end{lemma}

\begin{proof}
 La question étant locale sur $X$, on peut supposer que $X$ est quasi-compact. Soit 
$X = X_1 \cup \ldots \cup X_n$ la réunion des composantes irréductibles de $X$.
Supposons d'abord que $X$ contient un ouvert dense qui est lisse sur $k$.
Alors $X$ est géométriquement réduit au point générique de chaque $X_i$
 (par exemple d'après le lemme \ref{varieties-lemma-smooth-geometrically-normal}).
D'après le lemme \ref{varieties-lemma-generic-points-geometrically-reduced},
il s'ensuit que $X$ est géométriquement réduit.

\medskip\noindent
 Réciproquement, supposons que $X$ soit géométriquement réduit.
Comme $Z = \bigcup_{i \not = j} X_i \cap X_j$ est nulle part dense dans $X$,
nous pouvons remplacer $X$ par $X \setminus Z$. Comme $X \setminus Z$ est
une réunion disjointe de schémas irréductibles, cela nous ramène au cas
où $X$ est irréductible. Comme $X$ est irréductible et réduit, il
est intègre,
voir Propriétés, lemme \ref{properties-lemma-characterize-integral}.
Soit $\eta \in X$ son point générique.
Alors le corps des fonctions $K = k(X) = \kappa(\eta) = \mathcal{O}_{X, \eta}$
 est géométriquement réduit sur $k$, donc séparable sur $k$,
voir Algèbre, lemme \ref{algebra-lemma-characterize-separable-field-extensions}.
Soit $U = \Spec(A) \subset X$ un ouvert affine non vide quelconque, de
sorte que $K = A_{(0)}$ soit le corps des fractions de $A$.
Appliquons Algèbre, lemme \ref{algebra-lemma-separable-smooth}
 pour en déduire que $A$ est lisse en $(0)$ sur $k$.
Par définition, cela signifie qu'une certaine localisation principale de 
$A$ est lisse sur $k$, ce qui achève la démonstration. 
\end{proof}

\begin{lemma}
\label{varieties-lemma-dense-smooth-open-variety-over-perfect-field}
 Soit $k$ un corps parfait. Soit $X$
un $k$-schéma réduit localement algébrique, par exemple une variété sur $k$. Alors nous avons 
$$
\{x \in X \mid X \to \Spec(k)\text{ is smooth at }x\} =
\{x \in X \mid \mathcal{O}_{X, x}\text{ is regular}\}
$$
 et c'est un ouvert dense de $X$. 
\end{lemma}

\begin{proof}
 L'égalité des deux ensembles résulte immédiatement
d'Algèbre, lemme \ref{algebra-lemma-separable-smooth} et des définitions
(voir Algèbre, définition \ref{algebra-definition-perfect} pour la définition
d'un corps parfait). L'ensemble est ouvert car l'ensemble des points où
un morphisme de schémas est lisse est
ouvert, voir Morphismes, définition \ref{morphisms-definition-smooth}.
Enfin, nous donnons deux arguments pour voir qu'il est dense :
(1) Les points génériques de $X$ sont dans l'ensemble car les anneaux locaux
aux points génériques sont des corps (Algèbre, lemme 
\ref{algebra-lemma-minimal-prime-reduced-ring}), donc réguliers. (2)
Nous utilisons que $X$ est géométriquement réduit d'après le
lemme \ref{varieties-lemma-perfect-reduced} et donc
le lemme \ref{varieties-lemma-geometrically-reduced-dense-smooth-open} s'applique. 
\end{proof}

\begin{lemma}
\label{varieties-lemma-flat-under-smooth}
 Soit $k$ un corps. Soit $f : X \to Y$ un morphisme de schémas localement
de type fini sur $k$. Soit $x \in X$ un point et posons $y = f(x)$.
Si $X \to \Spec(k)$ est lisse en $x$ et que $f$ est plat en $x$
 alors $Y \to \Spec(k)$ est lisse en $y$. En particulier, si $X$
 est lisse sur $k$ et que $f$ est plat et surjectif, alors $Y$ est lisse sur $k$. 
\end{lemma}

\begin{proof}
 Il suffit de montrer que $Y$ est géométriquement régulier en $y$, voir
le lemme \ref{varieties-lemma-geometrically-regular-smooth}.
Cela découle du
lemme \ref{varieties-lemma-flat-under-geometrically-regular}
 (et
du lemme \ref{varieties-lemma-geometrically-regular-smooth}
 appliqué à $(X, x)$). 
\end{proof}

\begin{lemma}
\label{varieties-lemma-variety-with-smooth-rational-point}
 Soit $k$ un corps. Soit $X$ une variété sur $k$ possédant un point 
$k$-rationnel $x$ tel que $X$ soit lisse en $x$.
Alors $X$ est géométriquement intègre sur $k$. 
\end{lemma}

\begin{proof}
 Soit $U \subset X$ le lieu lisse de $X$. Par hypothèse, $U$ n'est pas
vide et est donc dense et schématiquement dense. Alors 
$U_{\overline{k}} \subset X_{\overline{k}}$ est également
dense et schématiquement dense (certains détails omis). Par conséquent,
il suffit de montrer que $U$ est géométriquement intègre.
Comme $U$ possède un point $k$-rationnel, il est géométriquement connexe d'après le
lemme \ref{varieties-lemma-geometrically-connected-if-connected-and-point}.
D'autre part, $U_{\overline{k}}$ est réduit et normal
(lemme \ref{varieties-lemma-smooth-geometrically-normal}).
Puisqu'un schéma noethérien connexe et
normal est intègre (Propriétés, lemme \ref{properties-lemma-normal-Noetherian}),
la démonstration est complète. 
\end{proof}

\begin{lemma}
\label{varieties-lemma-smooth-stratification}
 Soit $X$ un schéma de type fini sur un corps $k$.
Il existe une extension finie purement inséparable $k'/k$,
un entier $t \geq 0$, et des sous-schémas fermés 
$$
X_{k'} \supset Z_0 \supset Z_1 \supset \ldots \supset Z_t = \emptyset
$$
 de sorte que $Z_0 = (X_{k'})_{red}$ et $Z_i \setminus Z_{i + 1}$
 soient lisses sur $k'$ pour tout $i$. 
\end{lemma}

\begin{proof}
 Procédons par récurrence sur $\dim(X)$. D'après
le lemme \ref{varieties-lemma-finite-extension-geometrically-reduced},
nous pouvons trouver une extension finie purement inséparable $k'/k$
 telle que $(X_{k'})_{red}$ soit géométriquement réduit sur $k'$.
D'après le lemme \ref{varieties-lemma-geometrically-reduced-dense-smooth-open},
il existe un sous-schéma fermé nulle part dense $X' \subset (X_{k'})_{red}$
 tel que $(X_{k'})_{red} \setminus X'$ soit lisse sur $k'$.
Ensuite $\dim(X') < \dim(X)$. Par hypothèse d'induction, il existe
une extension finie purement inséparable $k''/k'$,
un entier $t' \geq 0$, et des sous-schémas fermés 
$$
X'_{k''} \supset Y_0 \supset Y_1 \supset \ldots \supset Y_{t'} = \emptyset
$$
 tels que $Y_0 = (X'_{k''})_{red}$ et $Y_i \setminus Y_{i + 1}$
 soit lisse sur $k''$ pour tout $i$. Posons alors $t = t' + 1$
 et nous considérons 
$$
X_{k''} \supset Z_0 \supset Z_1 \supset \ldots \supset Z_t = \emptyset
$$
 où $Z_0 = (X_{k''})_{red}$ et $Z_i = Y_{i - 1}$ pour $i > 0$;
cela a du sens car $X'_{k''}$ est un sous-schéma fermé de $X_{k''}$.
Nous omettons de vérifier toutes les propriétés énoncées. 
\end{proof}




\section{Types de variétés} 
\label{varieties-section-types}

\noindent
 Cette courte section traite de quelques propriétés globales élémentaires des variétés.

\begin{definition}
\label{varieties-definition-variety-type}
 Soit $k$ un corps. Soit $X$ une variété sur $k$. 
\begin{enumerate}
\item On dit que $X$ est une {\it variété affine} si $X$ est un schéma
affine. Cela équivaut à exiger que $X$ soit isomorphe à un sous-schéma
fermé de $\mathbf{A}^n_k$ pour un certain $n$. 
\item On dit que $X$ est une {\it variété projective} si le
morphisme de structure $X \to \Spec(k)$ est projectif.
D'après Morphismes, lemme \ref{morphisms-lemma-characterize-locally-projective},
cela est vrai si et seulement si $X$ est isomorphe à un sous-schéma
fermé de $\mathbf{P}^n_k$ pour un certain $n$. 
\item On dit que $X$ est une {\it variété quasi-projective} si le
morphisme de structure $X \to \Spec(k)$ est quasi-projectif. D'après Morphismes, lemme
\ref{morphisms-lemma-characterize-locally-quasi-projective},
cela est vrai si et seulement si $X$ est isomorphe à
un sous-schéma localement fermé de $\mathbf{P}^n_k$ pour un certain $n$. 
\item Une {\it variété propre} est une variété telle que
le morphisme $X \to \Spec(k)$ est propre. 
\item Une {\it variété lisse} est une variété telle que
le morphisme $X \to \Spec(k)$ est lisse. 
\end{enumerate}
\end{definition}

\noindent
 Notons qu'une variété projective est une variété propre,
voir Morphismes, lemme \ref{morphisms-lemma-locally-projective-proper}.
De plus, une variété affine est quasi-projective car $\mathbf{A}^n_k$
 est isomorphe à un sous-schéma ouvert de $\mathbf{P}^n_k$,
voir Constructions,
lemme \ref{constructions-lemma-standard-covering-projective-space}.

\begin{lemma}
\label{varieties-lemma-regular-functions-proper-variety}
 Soit $X$ une variété propre sur $k$. Alors 
\begin{enumerate}
\item $K = H^0(X, \mathcal{O}_X)$ est un corps qui
est une extension finie du corps $k$, 
\item si $X$ est géométriquement réduit, alors $K/k$ est séparable, 
\item si $X$ est géométriquement irréductible, alors $K/k$
 est purement inséparable, 
\item si $X$ est géométriquement intègre, alors $K = k$. 
\end{enumerate}
\end{lemma}

\begin{proof}
 Ceci est un cas particulier du
lemme \ref{varieties-lemma-proper-geometrically-reduced-global-sections}. 
\end{proof}




\section{Normalisation} 
\label{varieties-section-normalization}

\noindent
 Quelques questions relatives à la normalisation.

\begin{lemma}
\label{varieties-lemma-normalization-locally-algebraic}
 Soit $k$ un corps. Soit $X$ un schéma localement algébrique sur $k$.
Soit $\nu : X^\nu \to X$ le morphisme de normalisation,
voir Morphismes, définition \ref{morphisms-definition-normalization}.
Alors 
\begin{enumerate}
\item $\nu$ est fini, dominant, et $X^\nu$ est une
réunion disjointe de schémas irréductibles normaux localement algébriques sur $k$, 
\item $\nu$ se factorise en $X^\nu \to X_{red} \to X$ et le
premier morphisme est le morphisme de normalisation de $X_{red}$, 
\item si $X$ est un schéma algébrique réduit, alors $\nu$
 est birationnel, 
\item si $X$ est une variété, alors $X^\nu$ est une variété et 
$\nu$ est un morphisme birationnel fini de variétés. 
\end{enumerate}
\end{lemma}

\begin{proof}
 Puisque $X$ est localement de type fini sur un corps, on voit que 
$X$ est localement noethérien
(Morphismes, lemme \ref{morphisms-lemma-finite-type-noetherian}),
donc chaque ouvert quasi-compact a un nombre fini
de composantes irréductibles (Propriétés, lemme 
\ref{properties-lemma-Noetherian-irreducible-components}).
Ainsi, Morphismes, définition \ref{morphisms-definition-normalization}, s'applique.
La normalisation $X^\nu$ est toujours une réunion disjointe de schémas
intègres normaux et le morphisme de normalisation $\nu$ est toujours dominant,
voir Morphismes, lemme \ref{morphisms-lemma-normalization-normal}.
Comme $X$ est universellement Nagata
(Morphismes, lemme \ref{morphisms-lemma-ubiquity-nagata}),
on voit que $\nu$ est
fini (Morphismes, lemme \ref{morphisms-lemma-nagata-normalization}).
Ainsi $X^\nu$ est également localement algébrique. À
ce stade, nous avons prouvé (1).

\medskip\noindent
 La partie (2) est donnée par Morphismes, lemme \ref{morphisms-lemma-normalization-reduced}.

\medskip\noindent
 La partie (3) est donnée par Morphismes, lemme \ref{morphisms-lemma-normalization-birational}.

\medskip\noindent
 La partie (4) découle de (1), (2), (3), et du fait que $X^\nu$
 est séparé en tant que schéma fini sur un schéma séparé. 
\end{proof}

\begin{lemma}
\label{varieties-lemma-normalize-projective}
 Soit $k$ un corps. Soit $X$ un schéma propre sur $k$.
Soit $\nu : X^\nu \to X$ le morphisme de normalisation,
voir Morphismes, définition \ref{morphisms-definition-normalization}.
Alors $X^\nu$ est propre sur $k$. Si $X$ est projectif sur $k$,
alors $X^\nu$ est projectif sur $k$. 
\end{lemma}

\begin{proof}
 D'après le lemme \ref{varieties-lemma-normalization-locally-algebraic} le morphisme $\nu$
 est fini. Par conséquent, $X^\nu$ est propre sur $k$
 d'après Morphismes, lemmes \ref{morphisms-lemma-finite-proper} et 
\ref{morphisms-lemma-composition-proper}.
Le morphisme $\nu$ est projectif d'après
Morphismes, lemme \ref{morphisms-lemma-finite-projective}
 et donc si $X$ est projectif sur $k$, alors $X^\nu$ est projectif sur $k$
 d'après Morphismes, le lemme \ref{morphisms-lemma-composition-projective}. 
\end{proof}

\begin{lemma}
\label{varieties-lemma-relative-normalization-finite}
 Soit $k$ un corps. Soit $f : Y \to X$ un morphisme
quasi-compact de schémas localement algébriques sur $k$. Soit $X'$
 la normalisation de $X$ dans $Y$. Si $Y$ est réduit, alors 
$X' \to X$ est fini. 
\end{lemma}

\begin{proof}
 Puisque $Y$ est quasi-séparé (par
Propriétés, lemme \ref{properties-lemma-locally-Noetherian-quasi-separated}
 et Morphismes, lemme \ref{morphisms-lemma-finite-type-noetherian}),
le morphisme $f$ est quasi-séparé
(Schémas, lemme \ref{schemes-lemma-compose-after-separated}).
Ainsi, Morphismes, définition \ref{morphisms-definition-normalization-X-in-Y}
 s'applique. Le résultat découle de Morphismes, lemme 
\ref{morphisms-lemma-nagata-normalization-finite-general}.
Cela utilise que les schémas localement algébriques sont localement noethériens
(et ont donc localement un nombre fini de composantes irréductibles)
et que les schémas localement algébriques sont de
Nagata (Morphismes, lemme \ref{morphisms-lemma-ubiquity-nagata}).
Quelques petits détails sont omis. 
\end{proof}

\begin{lemma}
\label{varieties-lemma-finite-extension-geometrically-normal}
 Soit $k$ un corps. Soit $X$ un $k$-schéma
algébrique. Alors il existe une extension finie purement inséparable $k'/k$
 telle que la normalisation $Y$ de $X_{k'}$ soit géométriquement normale sur $k'$. 
\end{lemma}

\begin{proof}
 Soit $K = k^{perf}$ la clôture parfaite. Soit $Y_K$
la normalisation de $X_K$, voir le lemme \ref{varieties-lemma-normalization-locally-algebraic}.
D'après Limites, lemme \ref{limits-lemma-descend-finite-presentation},
il existe une sous-extension finie $K/k'/k$ et un morphisme 
$\nu : Y \to X_{k'}$ de présentation finie dont le changement de base vers $K$
 est le morphisme de normalisation $\nu_K : Y_K \to X_K$.
On observe que $Y$ est géométriquement normal sur $k'$
 (lemme \ref{varieties-lemma-geometrically-normal}).
Après avoir augmenté $k'$, on peut supposer que $Y \to X_{k'}$ est fini
(Limites, lemme \ref{limits-lemma-descend-finite-finite-presentation}).
Comme $\nu_K : Y_K \to X_K$ est le morphisme de normalisation,
il induit un morphisme birationnel $Y_K \to (X_K)_{red}$.
Par conséquent, il existe un ouvert dense $V_K \subset X_K$ tel que 
$\nu_K^{-1}(V_K) \to V_K$ soit une immersion
fermée (induisant un isomorphisme de $\nu_K^{-1}(V_K)$ avec $V_{K, red}$,
voir par exemple Morphismes, lemme 
\ref{morphisms-lemma-birational-isomorphism-over-dense-open}).
Après avoir augmenté $k'$, nous trouvons que $V_K$ est le changement de base d'un ouvert dense 
$V \subset Y$ et que le morphisme $\nu^{-1}(V) \to V$ est une immersion
fermée (Limites, lemmes \ref{limits-lemma-descend-opens} et 
\ref{limits-lemma-descend-closed-immersion-finite-presentation}).
Il s'ensuit facilement de cela que $\nu$ est le morphisme de
normalisation et la preuve est complète. 
\end{proof}

\begin{lemma}
\label{varieties-lemma-normalization-and-change-of-fields}
 Soit $k$ un corps. Soit $X$ un $k$-schéma localement algébrique.
Soit $K/k$ une extension de corps. Soit $\nu : X^\nu \to X$
 la normalisation de $X$, et soit $Y^\nu \to X_K$ la
normalisation du changement de base. Alors le morphisme canonique 
$$
Y^\nu \longrightarrow X^\nu \times_{\Spec(k)} \Spec(K)
$$
 est un isomorphisme si $K/k$ est séparable et un homéomorphisme universel
en général. 
\end{lemma}

\begin{proof}
 Posons $Y = X_K$. Soient $X^{(0)}$, resp.\ $Y^{(0)}$, les ensembles des points
génériques des composantes irréductibles de $X$, resp. \ $Y$. Alors le morphisme de projection 
$\pi : Y \to X$ satisfait $\pi(Y^{(0)}) = X^{(0)}$. Cela est vrai car 
$\pi$ est surjectif, ouvert et générisant,
voir Morphismes, lemme \ref{morphisms-lemma-scheme-over-field-universally-open} et 
\ref{morphisms-lemma-open-generizing}.
Si nous considérons $X^{(0)}$, resp. \ $Y^{(0)}$ comme des schémas (réduits), alors 
$X^\nu$, resp. \ $Y^\nu$ est la normalisation de $X$, resp. \ $Y$ dans 
$X^{(0)}$, resp. \ $Y^{(0})$.
Ainsi, Morphismes, lemme \ref{morphisms-lemma-functoriality-normalization}
 donne un morphisme canonique $Y^\nu \to X^\nu$ sur $Y \to X$ qui
à son tour donne le morphisme canonique du lemme par la
propriété universelle du produit fibré.

\medskip\noindent
 Pour montrer que ce morphisme possède les propriétés énoncées dans le lemme, on peut
supposer que $X = \Spec(A)$ est affine. Soit $Q(A_{red})$
l'anneau total des fractions de $A_{red}$. Alors $X^\nu$ est le spectre
de la clôture intégrale $A'$ de $A$ dans $Q(A_{red})$,
voir Morphismes, lemmes \ref{morphisms-lemma-normalization-reduced} et 
\ref{morphisms-lemma-description-normalization}.
De même, $Y^\nu$ est le spectre de la clôture intégrale $B'$ de 
$A \otimes_k K$ dans $Q((A \otimes_k K)_{red})$. Il existe un homomorphisme
canonique $Q(A_{red}) \to Q((A \otimes_k K)_{red})$, un homomorphisme canonique 
$A' \to B'$, et le morphisme du lemme correspond à l'homomorphisme induit. 
$$
A' \otimes_k K \longrightarrow B'
$$
 de $K$-algèbres. Le noyau est constitué d'éléments nilpotents
car le noyau de $Q(A_{red}) \otimes_k K \to Q((A \otimes_k K)_{red})$
 est l'ensemble des éléments nilpotents.

\medskip\noindent
 Si $K/k$ est séparable, alors $A' \otimes_k K$ est normal d'après
le lemme \ref{varieties-lemma-base-change-normal-by-separable}. En particulier,
il est réduit, donc $Q((A \otimes_k K)_{red}) = Q(A' \otimes_k K)$
 et $B' = A' \otimes_k K$
 d'après Algèbre, lemme \ref{algebra-lemma-characterize-reduced-ring-normal}.

\medskip\noindent
 Supposons que $K/k$ ne soit pas séparable. Alors la caractéristique de $k$ est $p > 0$.
Nous allons montrer que pour chaque $b \in B'$ il existe une puissance $q$ de $p$
 telle que $b^q$ appartienne à l'image de $A' \otimes_k K$. Cela montrera
que l'application affichée est un homéomorphisme universel d'après le
lemme d'Algèbre \ref{algebra-lemma-p-ring-map}.
Pour $b$ fixé, il existe un sous-corps $F \subset K$
 avec $F/k$ de type fini tel que $b$ soit contenu dans 
$Q((A \otimes_k F)_{red})$ et soit entier sur $A \otimes_k F$.
Choisissons un polynôme unitaire $P = T^d + \alpha_1 T^{d - 1} + \ldots + \alpha _d$
 avec $P(b) = 0$ et $\alpha_i \in A \otimes_k F$.
Choisissons une base de transcendance $t_1, \ldots, t_r$ de $F$ sur $k$.
Soit $F/F'/k(t_1, \ldots, t_r)$ la sous-extension
séparable maximale (Corps, lemme \ref{fields-lemma-separable-first}).
Comme $F/F'$ est finie purement inséparable, il existe un $q$ tel
que $\lambda^q \in F'$ pour tout $\lambda \in F$. Alors $b^q$ est
dans $Q((A \otimes_k F')_{red})$ et satisfait le polynôme 
$T^d + \alpha_1^q T^{d - 1} + \ldots + \alpha _d^q$ avec 
$\alpha_i^q \in A \otimes_k F'$. Par le cas séparable,
nous voyons que $b^q \in A' \otimes_k F'$ et la démonstration est complète. 
\end{proof}

\begin{lemma}
\label{varieties-lemma-geometrically-normal-in-codim-1}
 Soit $k$ un corps. Soit $X$ un $k$-schéma localement algébrique. Soit
$\nu : X^\nu \to X$ la normalisation de $X$.
Soit $x \in X$ un point tel que (a) $\mathcal{O}_{X, x}$
 soit réduit, (b) $\dim(\mathcal{O}_{X, x}) = 1$, et
(c) pour tout $x' \in X^\nu$ avec $\nu(x') = x$ l'extension 
$\kappa(x')/k$ soit séparable. Alors $X$ est géométriquement réduit en $x$
 et $X^\nu$ est géométriquement régulier en $x'$ avec $\nu(x') = x$. 
\end{lemma}

\begin{proof}
 Nous utiliserons les résultats du lemme \ref{varieties-lemma-normalization-locally-algebraic}
 sans autre mention. Soit $x' \in X^\nu$ un point au-dessus de $x$.
Par la théorie de la dimension (section \ref{varieties-section-algebraic-schemes}), nous avons 
$\dim(\mathcal{O}_{X^\nu, x'}) = 1$. Puisque $X^\nu$ est normal, on
voit que $\mathcal{O}_{X^\nu, x'}$ est un anneau de valuation
discrète (Propriétés, lemme \ref{properties-lemma-criterion-normal}).
Ainsi, $\mathcal{O}_{X^\nu, x'}$ est une $k$-algèbre locale
régulière dont le corps résiduel est séparable sur $k$. Par conséquent, 
$k \to \mathcal{O}_{X^\nu, x'}$ est formellement lisse dans la
topologie $\mathfrak m_{x'}$-adique, voir
Compléments d'algèbre, lemme \ref{more-algebra-lemma-regular-implies-fs}.
Alors $\mathcal{O}_{X^\nu, x'}$ est géométriquement régulier
sur $k$
 par Compléments d'algèbre, théorème \ref{more-algebra-theorem-regular-fs}.
Ainsi $X^\nu$ est géométriquement régulier en $x'$ d'après
le lemme \ref{varieties-lemma-geometrically-regular-at-point}.

\medskip\noindent
 Puisque $\mathcal{O}_{X, x}$ est réduit, la famille d'homomorphismes 
$\mathcal{O}_{X, x} \to \mathcal{O}_{X^\nu, x'}$ est injective.
Puisque $\mathcal{O}_{X^\nu, x'}$ est une 
$k$-algèbre géométriquement réduite, il s'ensuit immédiatement que $\mathcal{O}_{X, x}$
 est une $k$-algèbre géométriquement réduite. Par conséquent, $X$ est géométriquement
réduit en $x$ d'après le lemme \ref{varieties-lemma-geometrically-reduced-at-point}. 
\end{proof}









\section{Groupes de fonctions inversibles} 
\label{varieties-section-units}

\noindent
 Il arrive souvent (mais pas toujours) que $\mathcal{O}^*(X)/k^*$
 soit un groupe abélien de type fini si $X$ est une variété sur $k$.
Nous le montrons par une série
de lemmes. Tout repose sur le cas spécial suivant.

\begin{lemma}
\label{varieties-lemma-open-in-normal-proper}
 Soit $k$ un corps algébriquement clos. Soit
$\overline{X}$ une variété propre sur $k$.
Soit $X \subset \overline{X}$ un sous-schéma ouvert.
Supposons que $X$ soit normal.
Alors $\mathcal{O}^*(X)/k^*$ est un groupe abélien de type fini. 
\end{lemma}

\begin{proof}
 Puisque l'énoncé ne concerne que $X$, nous pouvons remplacer $\overline{X}$
 par une autre variété propre sur $k$. Soit $\nu : \overline{X}^\nu \to
\overline{X}$ le morphisme de normalisation. D'après le
lemme \ref{varieties-lemma-normalization-locally-algebraic},
on sait que $\nu$ est fini et que $\overline{X}^\nu$ est
une variété. Comme $X$ est normale, on voit que $\nu^{-1}(X) \to X$ est
un isomorphisme (petit détail omis). Enfin, on voit que 
$\overline{X}^\nu$ est propre sur $k$, car un morphisme fini
est propre (Morphismes, lemme \ref{morphisms-lemma-finite-proper})
et les compositions de morphismes
propres sont propres (Morphismes, lemme \ref{morphisms-lemma-composition-proper}).
Ainsi, nous pouvons et faisons l'hypothèse que $\overline{X}$ est normal.

\medskip\noindent
 Nous utiliserons sans plus le rappeler que, pour tout ouvert affine $U$ de 
$\overline{X}$, l'anneau $\mathcal{O}(U)$ est une 
$k$-algèbre de type fini, qui est noethérienne, intègre et normale,
voir Algèbre, lemme \ref{algebra-lemma-Noetherian-permanence},
Propriétés, définition \ref{properties-definition-integral},
Propriétés, lemmes \ref{properties-lemma-locally-Noetherian} et 
\ref{properties-lemma-locally-normal},
Morphismes, lemme \ref{morphisms-lemma-locally-finite-type-characterize}.

\medskip\noindent
 Soient $\xi_1, \ldots, \xi_r$ les points génériques du complément de $X$
 dans $\overline{X}$. Ils sont en nombre fini puisque $\overline{X}$ possède
un espace topologique sous-jacent noethérien (voir
Morphismes, lemme \ref{morphisms-lemma-finite-type-noetherian},
Propriétés, lemme \ref{properties-lemma-Noetherian-topology},
et Topologie, lemme \ref{topology-lemma-Noetherian}).
Pour chaque $i$, l'anneau local $\mathcal{O}_i = \mathcal{O}_{X, \xi_i}$
 est un anneau intègre local normal noethérien (comme une localisation d'un
anneau intègre normal noethérien). Soit $J \subset \{1, \ldots, r\}$ l'ensemble
des indices $i$ tels que $\dim(\mathcal{O}_i) = 1$. Pour $j \in J$,
l'anneau local $\mathcal{O}_j$ est un anneau de valuation discrète,
voir Algèbre, lemme \ref{algebra-lemma-characterize-dvr}.
Par conséquent, nous obtenons une valuation. 
$$
v_j : k(\overline{X})^* \longrightarrow \mathbf{Z}
$$
 avec la propriété que $v_j(f) \geq 0 \Leftrightarrow f \in \mathcal{O}_j$.

\medskip\noindent
 Identifions $\mathcal{O}(X)$ à une sous-$k$-algèbre de $k(X) = k(\overline{X})$.
Nous affirmons que le noyau de l'application 
$$
\mathcal{O}(X)^* \longrightarrow
\prod\nolimits_{j \in J} \mathbf{Z},
\quad
f \longmapsto \prod v_j(f)
$$
 est $k^*$. Il est clair que cette affirmation prouve le lemme.
 En effet, soit $f \in \mathcal{O}(X)$ un élément du noyau. Soit
$U = \Spec(B) \subset \overline{X}$ un ouvert affine quelconque.
Alors $B$ est un anneau intègre normal noethérien. Pour
chaque idéal premier de hauteur un $\mathfrak q \subset B$ avec le point
correspondant $\xi \in X$, on a ou bien $\xi = \xi_j$ pour un certain $j \in J$,
ou bien $\xi \in X$. En effet, 
$\text{codim}(\overline{\{\xi\}}, \overline{X}) = 1$ d'après
Propriétés, lemme \ref{properties-lemma-codimension-local-ring},
 et donc si $\xi \in \overline{X} \setminus X$, cela doit
être un point générique de $\overline{X} \setminus X$, donc égal à un certain 
$\xi_j$, $j \in J$.
Nous concluons que $f \in \mathcal{O}_{X, \xi} = B_{\mathfrak q}$
 dans les deux cas puisque $f$ est dans le noyau de l'application. Ainsi 
$f \in \bigcap_{\text{ht}(\mathfrak q) = 1} B_{\mathfrak q} = B$,
voir Algèbre, lemme 
\ref{algebra-lemma-normal-domain-intersection-localizations-height-1}.
En d'autres termes, on voit que 
$f \in \Gamma(\overline{X}, \mathcal{O}_{\overline{X}})$.
Mais puisque $k$ est algébriquement clos, nous concluons que 
$f \in k$ d'après
le lemme \ref{varieties-lemma-regular-functions-proper-variety}. 
\end{proof}

\noindent
 Nous généralisons maintenant le cas précédent par quelques arguments élémentaires, en
supposant toujours le corps algébriquement clos.

\begin{lemma}
\label{varieties-lemma-units-integral-finite-type-algebraically-closed}
 Soit $k$ un corps algébriquement clos. Soit
$X$ un schéma intègre localement de type fini sur $k$.
Alors $\mathcal{O}^*(X)/k^*$ est un groupe abélien de type fini. 
\end{lemma}

\begin{proof}
 Comme $X$ est intègre, l'homomorphisme de restriction 
$\mathcal{O}(X) \to \mathcal{O}(U)$ est injectif pour tout
sous-schéma ouvert non vide $U \subset X$. Par conséquent, on peut supposer
que $X$ est affine. Choisissons une immersion fermée 
$X \to \mathbf{A}^n_k$
 et notons $\overline{X}$ l'adhérence de $X$ dans $\mathbf{P}^n_k$
 via l'immersion usuelle $\mathbf{A}^n_k \to \mathbf{P}^n_k$.
Ainsi, on peut supposer que $X$ est un ouvert affine d'une variété
projective $\overline{X}$.

\medskip\noindent
 Soit $\nu : \overline{X}^\nu \to \overline{X}$ le morphisme
de normalisation,
voir Morphismes, définition \ref{morphisms-definition-normalization}.
Nous savons que $\nu$ est fini, dominant, et que $\overline{X}^\nu$
 est un schéma irréductible normal,
voir Morphismes, lemmes \ref{morphisms-lemma-normalization-normal}, 
\ref{morphisms-lemma-Japanese-normalization}, et 
\ref{morphisms-lemma-ubiquity-nagata}.
Il s'ensuit que $\overline{X}^\nu$ est une variété propre,
parce que $\overline{X} \to \Spec(k)$ est propre comme composition d'un
morphisme fini et d'un morphisme propre (voir les
résultats dans Morphismes, sections \ref{morphisms-section-proper} et 
\ref{morphisms-section-integral}).
Il s'ensuit également que $\nu$ est un morphisme surjectif, parce que
l'image de $\nu$ est fermée et contient le point générique de $\overline{X}$.
Ainsi, en posant $X^\nu = \nu^{-1}(X)$, nous voyons qu'il suffit de prouver le
résultat pour $X^\nu$. En d'autres termes, on peut supposer que $X$ est un ouvert
non vide d'une variété propre normale $\overline{X}$. Ce cas est traité par
le lemme \ref{varieties-lemma-open-in-normal-proper}. 
\end{proof}

\noindent
 Le lemme précédent implique la légère généralisation suivante.

\begin{lemma}
\label{varieties-lemma-units-general-algebraically-closed}
 Soit $k$ un corps algébriquement clos. Soit
$X$ un schéma réduit connexe localement de type fini
sur $k$ avec un nombre fini de composantes irréductibles.
Alors $\mathcal{O}^*(X)/k^*$ est un groupe abélien de type fini. 
\end{lemma}

\begin{proof}
 Écrivons $X = \bigcup X_i$ pour la décomposition en composantes irréductibles. Par le
lemme \ref{varieties-lemma-units-integral-finite-type-algebraically-closed},
nous voyons que $\mathcal{O}(X_i)^*/k^*$ est un groupe abélien de
type fini. Soit $f \in \mathcal{O}(X)^*$ dans le noyau
de l'application 
$$
\mathcal{O}(X)^* \longrightarrow \prod \mathcal{O}(X_i)^*/k^*.
$$
 Alors pour chaque $i$ il existe un élément $\lambda_i \in k$
 tel que $f|_{X_i} = \lambda_i$.
En se restreignant à $X_i \cap X_j$, nous concluons que 
$\lambda_i = \lambda_j$ si $X_i \cap X_j \not = \emptyset$.
Puisque $X$ est connexe, nous concluons que tous les $\lambda_i$ sont
égaux et donc que $f \in k^*$. Cela prouve que 
$$
\mathcal{O}(X)^*/k^* \subset \prod \mathcal{O}(X_i)^*/k^*
$$
 et le lemme en découle, car le membre de droite est un produit d'un
nombre fini de groupes abéliens de type fini. 
\end{proof}

\begin{lemma}
\label{varieties-lemma-integral-closure-ground-field}
 Soit $k$ un corps. Soit
$X$ un schéma sur $k$, connexe et réduit.
Alors la clôture intégrale de $k$ dans $\Gamma(X, \mathcal{O}_X)$
 est un corps. 
\end{lemma}

\begin{proof}
 Soit $k' \subset \Gamma(X, \mathcal{O}_X)$ la clôture intégrale de 
$k$. Alors $X \to \Spec(k)$ se factorise à travers $\Spec(k')$,
voir Schémas, lemme \ref{schemes-lemma-morphism-into-affine}.
Comme $X$ est réduit, on voit que $k'$ n'a pas d'éléments nilpotents non nuls. Comme 
$k \to k'$ est intègre, on voit que tout idéal premier de $k'$ est
à la fois un idéal maximal et un premier minimal,
et $\Spec(k')$ est totalement discontinu,
voir Algèbre, lemmes \ref{algebra-lemma-integral-no-inclusion} et 
\ref{algebra-lemma-ring-with-only-minimal-primes}.
Comme $X$ est connexe, le morphisme $X \to \Spec(k')$
 est constant, d'image le point correspondant à 
$\mathfrak p \subset k'$. Alors tout $f \in k'$, $f \not \in \mathfrak p$
 s'envoie sur un élément inversible de $\mathcal{O}_X$. Par définition
de $k'$, il s'ensuit que $f$ est une unité de $k'$. Ainsi, on voit
que $k'$ est local avec idéal maximal $\mathfrak p$,
voir Algèbre, lemme \ref{algebra-lemma-characterize-local-ring}.
Puisque nous avons déjà vu que $k'$ est réduit, cela implique que 
$k'$ est un corps,
voir Algèbre, lemme \ref{algebra-lemma-minimal-prime-reduced-ring}. 
\end{proof}

\begin{proposition}
\label{varieties-proposition-units-general}
 Soit $k$ un corps. Soit $X$ un schéma sur $k$. Supposons que $X$
 soit localement de type fini sur $k$, connexe et réduit, et possède un nombre
fini de composantes irréductibles. Alors $\mathcal{O}(X)^*/k^*$ est un groupe abélien
de type fini si, en plus des conditions ci-dessus, au
moins l'une des conditions suivantes est satisfaite : 
\begin{enumerate}
\item la clôture intégrale de $k$ dans $\Gamma(X, \mathcal{O}_X)$ est $k$, 
\item $X$ a un point $k$-rationnel, ou 
\item $X$ est géométriquement intègre. 
\end{enumerate}
\end{proposition}

\begin{proof}
 Soit $\overline{k}$ une clôture algébrique de $k$.
Soit $Y$ une composante connexe de $(X_{\overline{k}})_{red}$.
Notons que le morphisme canonique $p : Y \to X$ est ouvert
(d'après Morphismes, lemme \ref{morphisms-lemma-scheme-over-field-universally-open})
et fermé
(d'après Morphismes, lemme \ref{morphisms-lemma-integral-universally-closed}).
Ainsi, $p(Y) = X$, puisque $X$ est connexe. En particulier, comme 
$X$ est réduit, cela implique $\mathcal{O}(X) \subset \mathcal{O}(Y)$. D'après
le lemme \ref{varieties-lemma-galois-action-irreducible-components-locally-finite-type},
nous voyons que $Y$ possède un nombre fini de composantes irréductibles (le corps 
$\overline{k}$ dans le lemme est la clôture algébrique séparable, mais cela
n'a pas d'importance car les deux changements de base $X_{\overline{k}}$
 ont les mêmes composantes irréductibles d'après le
lemme \ref{varieties-lemma-separably-closed-field-irreducible-components}). Ainsi,
le lemme \ref{varieties-lemma-units-general-algebraically-closed}
 s'applique à $Y$. Cela implique que si $\mathcal{O}(X)^*/k^*$
 n'est pas un groupe abélien de type fini, alors il existe un élément 
$f \in \mathcal{O}(X)$, $f \not \in k$, qui s'envoie sur un élément de 
$\overline{k}$ via l'application $\mathcal{O}(X) \to \mathcal{O}(Y)$.
Dans ce cas, $f$ est algébrique sur $k$,
donc intégral sur $k$. Ainsi, si la condition (1) est vérifiée, cela ne peut pas se
produire. Pour terminer la démonstration, nous montrons que les conditions (2) et (3) impliquent (1).

\medskip\noindent
 Soit $k \subset k' \subset \Gamma(X, \mathcal{O}_X)$ la clôture
intégrale de $k$ dans $\Gamma(X, \mathcal{O}_X)$. D'après
le lemme \ref{varieties-lemma-integral-closure-ground-field},
on voit que $k'$ est un corps.
Si $e : \Spec(k) \to X$ est un point $k$-rationnel, alors 
$e^\sharp : \Gamma(X, \mathcal{O}_X) \to k$ est une section
de l'application d'inclusion $k \to \Gamma(X, \mathcal{O}_X)$. En particulier,
la restriction de $e^\sharp$ à $k'$ est un homomorphisme de corps $k' \to k$ sur $k$,
ce qui montre clairement que (2) implique (1).

\medskip\noindent
 Si la clôture intégrale $k'$ de $k$ dans $\Gamma(X, \mathcal{O}_X)$
 n'est pas triviale, alors on voit que $X$ n'est pas géométriquement connexe 
(si $k'/k$ n'est pas purement inséparable), ou bien que $X$ n'est pas
géométriquement réduit (si $k'/k$ est purement inséparable non triviale). Nous omettons
les détails. Par conséquent, (3) implique (1). 
\end{proof}

\begin{lemma}
\label{varieties-lemma-units-variety}
 Soit $k$ un corps. Soit
$X$ une variété sur $k$.
Le groupe $\mathcal{O}(X)^*/k^*$ est un groupe abélien de type fini
à condition qu'au moins une des conditions suivantes soit remplie : 
\begin{enumerate}
\item $k$ est intégralement clos dans $\Gamma(X, \mathcal{O}_X)$, 
\item $k$ est algébriquement clos dans $k(X)$, 
\item $X$ est géométriquement intègre sur $k$, ou 
\item $k$ est l'intersection `` '' des extensions de corps 
$\kappa(x)/k$ où $x$ parcourt les points fermés de $x$. 
\end{enumerate}
\end{lemma}

\begin{proof}
 On voit que (1) suffit
d'après la proposition \ref{varieties-proposition-units-general}.
Nous omettons la vérification que chacun de (2), (3), (4) implique (1). 
\end{proof}







\section{Formule de K\"unneth, I} 
\label{varieties-section-kunneth}

\noindent
 Dans cette section, nous démontrons la formule de K\"unneth lorsque
la base est un corps et que nous considérons la cohomologie
des modules quasi-cohérents. Pour une version plus générale, voir
Catégories dérivées de schémas, section \ref{perfect-section-kunneth}.

\begin{lemma}
\label{varieties-lemma-kunneth}
 Soit $k$ un corps. Soient $X$ et $Y$ des schémas sur $k$, et
soient $\mathcal{F}$, resp.\ $\mathcal{G}$, un module quasi-cohérent sur 
$\mathcal{O}_X$, resp.\ sur $\mathcal{O}_Y$.
Alors nous avons un isomorphisme canonique 
$$
H^n(X \times_{\Spec(k)} Y, \text{pr}_1^*\mathcal{F}
\otimes_{\mathcal{O}_{X \times_{\Spec(k)} Y}} \text{pr}_2^*\mathcal{G}) =
\bigoplus\nolimits_{p + q = n}
H^p(X, \mathcal{F}) \otimes_k H^q(Y, \mathcal{G})
$$
 si $X$ et $Y$ sont quasi-compacts et ont une
diagonale affine\footnote{Le cas où $X$ et $Y$ sont quasi-séparés sera
traité au lemme \ref{varieties-lemma-kunneth-general} ci-dessous.}
(par exemple si $X$ et $Y$ sont séparés). 
\end{lemma}

\begin{proof}
 Dans cette démonstration, les produits non décorés et les produits tensoriels sont pris sur $k$. Pour construire les applications 
$$
H^p(X, \mathcal{F}) \otimes H^q(Y, \mathcal{G})
\longrightarrow
H^n(X \times Y, \text{pr}_1^*\mathcal{F}
\otimes_{\mathcal{O}_{X \times Y}} \text{pr}_2^*\mathcal{G})
$$
 on utilise la fonctorialité de la cohomologie pour obtenir les homomorphismes 
$H^p(X, \mathcal{F}) \to H^p(X \times Y, \text{pr}_1^*\mathcal{F})$ et 
$H^p(Y, \mathcal{G}) \to H^p(X \times Y, \text{pr}_2^*\mathcal{G})$,
puis le cup-produit 
$$
\cup :
H^p(X \times Y, \text{pr}_1^*\mathcal{F})
\otimes H^q(X \times Y, \text{pr}_2^*\mathcal{G})
\longrightarrow
H^n(X \times Y, \text{pr}_1^*\mathcal{F}
\otimes_{\mathcal{O}_{X \times Y}} \text{pr}_2^*\mathcal{G})
$$
 Le résultat est vrai lorsque $X$ et $Y$ sont affines en raison de la
nullité des groupes de cohomologie supérieure sur les affines (Cohomologie des Schémas, lemme 
\ref{coherent-lemma-quasi-coherent-affine-cohomology-zero})
et des définitions (des images réciproques de modules quasi-cohérents et
des produits tensoriels de modules quasi-cohérents).

\medskip\noindent
 Choisissons des recouvrements affines ouverts finis 
$\mathcal{U} : X = \bigcup_{i \in I} U_i$ et 
$\mathcal{V} : Y = \bigcup_{j \in J} V_j$.
Cela détermine un recouvrement affine ouvert 
$\mathcal{W} : X \times Y = \bigcup_{(i, j) \in I \times J} U_i \times V_j$.
Notons que $\mathcal{W}$ est un raffinement de 
$\text{pr}_1^{-1}\mathcal{U}$ et de $\text{pr}_2^{-1}\mathcal{V}$.
 Ainsi, d'après Cohomologie, lemme \ref{cohomology-lemma-functoriality-cech},
on obtient des homomorphismes 
$$
\check{\mathcal{C}}^\bullet(\mathcal{U}, \mathcal{F}) \to
\check{\mathcal{C}}^\bullet(\mathcal{W}, \text{pr}_1^*\mathcal{F})
\quad\text{et}\quad
\check{\mathcal{C}}^\bullet(\mathcal{V}, \mathcal{G}) \to
\check{\mathcal{C}}^\bullet(\mathcal{W}, \text{pr}_2^*\mathcal{G})
$$
 compatibles avec les homomorphismes d'image réciproque en cohomologie. Dans Cohomologie, équation ( 
\ref{cohomology-equation-needs-signs})
nous avons construit une application de complexes 
$$
\text{Tot}(
\check{\mathcal{C}}^\bullet(\mathcal{W}, \text{pr}_1^*\mathcal{F})
\otimes
\check{\mathcal{C}}^\bullet(\mathcal{W}, \text{pr}_2^*\mathcal{G}))
\longrightarrow
\check{\mathcal{C}}^\bullet(\mathcal{W},
\text{pr}_1^*\mathcal{F} \otimes_{\mathcal{O}_{X \times Y}}
\text{pr}_2^*\mathcal{G})
$$
 qui définit le cup-produit en cohomologie. En combinant ce qui précède, nous
obtenons une application de complexes 
\begin{equation}
\label{varieties-equation-kunneth-on-cech}
\text{Tot}(
\check{\mathcal{C}}^\bullet(\mathcal{U}, \mathcal{F})
\otimes
\check{\mathcal{C}}^\bullet(\mathcal{V}, \mathcal{G}))
\longrightarrow
\check{\mathcal{C}}^\bullet(\mathcal{W},
\text{pr}_1^*\mathcal{F} \otimes_{\mathcal{O}_{X \times Y}}
\text{pr}_2^*\mathcal{G})
\end{equation}
 Cette application n'est pas
un isomorphisme de complexes. Rappelons que nous pouvons calculer les
cohomologies de nos faisceaux quasi-cohérents en utilisant nos
recouvrements (Cohomologie des Schémas, lemmes 
\ref{coherent-lemma-affine-diagonal} et 
\ref{coherent-lemma-cech-cohomology-quasi-coherent}).
Ainsi, en cohomologie, (\ref{varieties-equation-kunneth-on-cech}) reproduit l'application
du lemme.

\medskip\noindent
 Considérons une suite exacte courte
$0 \to \mathcal{F} \to \mathcal{F}' \to \mathcal{F}'' \to 0$
 de modules quasi-cohérents. Puisque la construction de ( 
\ref{varieties-equation-kunneth-on-cech} ) est fonctorielle en $\mathcal{F}$ et puisque
la formation des complexes de {\v C}ech concernés est exacte
dans la variable $\mathcal{F}$ (car nous prenons des sections sur des
ouverts affines), on obtient un morphisme entre deux suites exactes courtes
de complexes. 
$$
\xymatrix{
\text{Tot}(
\check{\mathcal{C}}^\bullet(\mathcal{U}, \mathcal{F})
\otimes
\check{\mathcal{C}}^\bullet(\mathcal{V}, \mathcal{G})) \ar[r] \ar[d] &
\text{Tot}(
\check{\mathcal{C}}^\bullet(\mathcal{U}, \mathcal{F}')
\otimes
\check{\mathcal{C}}^\bullet(\mathcal{V}, \mathcal{G})) \ar[r] \ar[d] &
\text{Tot}(
\check{\mathcal{C}}^\bullet(\mathcal{U}, \mathcal{F}'')
\otimes
\check{\mathcal{C}}^\bullet(\mathcal{V}, \mathcal{G})) \ar[d] \\
\check{\mathcal{C}}^\bullet(\mathcal{W},
\text{pr}_1^*\mathcal{F} \otimes_{\mathcal{O}_{X \times Y}}
\text{pr}_2^*\mathcal{G}) \ar[r] &
\check{\mathcal{C}}^\bullet(\mathcal{W},
\text{pr}_1^*\mathcal{F}' \otimes_{\mathcal{O}_{X \times Y}}
\text{pr}_2^*\mathcal{G}) \ar[r] &
\check{\mathcal{C}}^\bullet(\mathcal{W},
\text{pr}_1^*\mathcal{F}'' \otimes_{\mathcal{O}_{X \times Y}}
\text{pr}_2^*\mathcal{G})
}
$$
 (les zéros extérieurs ont été omis). En
considérant les suites exactes longues de cohomologie, on constate que, si le résultat
du lemme vaut pour $2$ des $3$ faisceaux 
$\mathcal{F}, \mathcal{F}', \mathcal{F}''$, alors il vaut aussi
pour le troisième.

\medskip\noindent
 Notons que $X$ est de dimension cohomologique finie pour
les modules quasi-cohérents, voir Cohomologie des Schémas, lemme 
\ref{coherent-lemma-vanishing-nr-affines}.
En utilisant l'induction sur 
$d(\mathcal{F}) = \max \{d \mid H^d(X, \mathcal{F}) \not = 0\}$,
nous allons réduire au cas $d(\mathcal{F}) = 0$.
Supposons $d(\mathcal{F}) > 0$.
 D'après Cohomologie des Schémas, lemme 
\ref{coherent-lemma-affine-diagonal-universal-delta-functor},
nous avons
vu qu'il existe un plongement $\mathcal{F} \to \mathcal{F}'$
 tel que $H^p(X, \mathcal{F}') = 0$ pour tout $p \geq 1$.
En posant $\mathcal{F}'' = \Coker(\mathcal{F} \to \mathcal{F}')$,
on voit que $d(\mathcal{F}'') < d(\mathcal{F})$.
Ensuite, nous pouvons appliquer le résultat du paragraphe précédent pour
voir qu'il suffit de prouver le lemme pour $\mathcal{F}'$
 et $\mathcal{F}''$, ce qui établit l'étape d'induction.

\medskip\noindent
 En raisonnant de même pour $\mathcal{G}$, on constate
qu'on peut supposer que $\mathcal{F}$ et $\mathcal{G}$
 n'ont une cohomologie non nulle qu'au degré $0$.
 Soit $V \subset Y$ un ouvert affine. Considérons le recouvrement affine ouvert 
$\mathcal{U}_V : X \times V = \bigcup_{i \in I} U_i \times V$.
Il est immédiat que 
$$
\check{\mathcal{C}}^\bullet(\mathcal{U}, \mathcal{F}) \otimes \mathcal{G}(V) =
\check{\mathcal{C}}^\bullet(\mathcal{U}_V,
\text{pr}_1^*\mathcal{F} \otimes_{\mathcal{O}_{X \times Y}}
\text{pr}_2^*\mathcal{G})
$$
 (égalité des complexes). Nous concluons que 
$$
R\text{pr}_{2, *}(\text{pr}_1^*\mathcal{F} \otimes_{\mathcal{O}_{X \times Y}}
\text{pr}_2^*\mathcal{G})
\cong
\Gamma(X, \mathcal{F}) \otimes_k \mathcal{G} \cong
\bigoplus\nolimits_{\alpha \in A} \mathcal{G}
$$
 sur $Y$. Ici $A$ est une base du $k$-espace vectoriel $\Gamma(X, \mathcal{F})$.
La cohomologie sur $Y$ commute avec les sommes
directes (Cohomologie, lemme \ref{cohomology-lemma-quasi-separated-cohomology-colimit}).
En utilisant la suite spectrale de Leray pour $\text{pr}_2$
 (via Cohomologie, lemme \ref{cohomology-lemma-apply-Leray}),
nous concluons que 
$H^n(X \times Y, \text{pr}_1^*\mathcal{F} \otimes_{\mathcal{O}_{X \times Y}}
\text{pr}_2^*\mathcal{G})$
 est nul pour $n > 0$ et isomorphe à 
$H^0(X, \mathcal{F}) \otimes H^0(Y, \mathcal{G})$ pour $n = 0$.
 Cela achève la démonstration (sauf qu'il faudrait vérifier que
l'isomorphisme est bien donné par le cup-produit en degré $0$ ;
nous omettons la vérification). 
\end{proof}

\begin{lemma}
\label{varieties-lemma-kunneth-general}
 Soit $k$ un corps. Soient $X$ et $Y$ des schémas sur $k$, et
soient $\mathcal{F}$, resp.\ $\mathcal{G}$, un module quasi-cohérent sur 
$\mathcal{O}_X$, resp.\ sur $\mathcal{O}_Y$.
Alors nous avons un isomorphisme canonique 
$$
H^n(X \times_{\Spec(k)} Y, \text{pr}_1^*\mathcal{F}
\otimes_{\mathcal{O}_{X \times_{\Spec(k)} Y}} \text{pr}_2^*\mathcal{G}) =
\bigoplus\nolimits_{p + q = n}
H^p(X, \mathcal{F}) \otimes_k H^q(Y, \mathcal{G})
$$
 à condition que $X$ et $Y$ soient quasi-compacts et quasi-séparés. 
\end{lemma}

\begin{proof}
 Si $X$ et $Y$ sont séparés ou, plus généralement, ont une diagonale affine, voir
le lemme \ref{varieties-lemma-kunneth} pour une « meilleure » démonstration
(dont l'avantage sur celle-ci est d'identifier les
applications comme des images réciproques suivies de cup-produits). Soient
$X'$, resp.\ $Y'$, les épaississements infinitésimaux de $X$, resp.\ $Y$,
de faisceaux de structure 
$\mathcal{O}_{X'} = \mathcal{O}_X \oplus \mathcal{F}$,
respectivement \ $\mathcal{O}_{Y'} = \mathcal{O}_Y \oplus \mathcal{G}$, où 
$\mathcal{F}$, resp.\ $\mathcal{G}$, est un idéal de carré nul.
Alors 
$$
\mathcal{O}_{X' \times Y'} =
\mathcal{O}_{X \times Y}
\oplus \text{pr}_1^*\mathcal{F}
\oplus \text{pr}_2^*\mathcal{G}
\oplus \text{pr}_1^*\mathcal{F}
\otimes_{\mathcal{O}_{X \times Y}} \text{pr}_2^*\mathcal{G}
$$
 en tant que faisceaux sur $X \times Y$. On voit ainsi qu'il suffit de
prouver que 
$$
H^n(X \times Y, \mathcal{O}_{X \times Y}) =
\bigoplus\nolimits_{p + q = n}
H^p(X, \mathcal{O}_X) \otimes_k H^q(Y, \mathcal{O}_Y)
$$
 pour toute paire de schémas quasi-compacts et quasi-séparés sur $k$.
Certains détails sont omis.

\medskip\noindent
 Pour démontrer cette assertion, nous utilisons la cohomologie et le changement de base sous
la forme de Cohomologie des Schémas, lemme \ref{coherent-lemma-hypercoverings}.
Ce lemme fournit un complexe de $k$-espaces vectoriels borné inférieurement, c'est-à-dire
un complexe $\mathcal{K}^\bullet$ de modules quasi-cohérents sur $\Spec(k)$,
qui calcule universellement la cohomologie de $Y$ sur $\Spec(k)$.
En particulier, on voit que 
$$
R\text{pr}_{1, *}(\mathcal{O}_{X \times Y}) \cong
(X \to \Spec(k))^*\mathcal{K}^\bullet
$$
 dans $D(\mathcal{O}_X)$. À homotopie près, le complexe $\mathcal{K}^\bullet$
 est isomorphe à $\bigoplus_{q \geq 0} H^q(Y, \mathcal{O}_Y)[-q]$ car
cela est vrai pour tout complexe d'espaces vectoriels sur un corps.
Nous concluons que 
$$
R\text{pr}_{1, *}(\mathcal{O}_{X \times Y}) \cong
\bigoplus\nolimits_{q \geq 0}
H^q(Y, \mathcal{O}_Y)[-q] \otimes_k \mathcal{O}_X
$$
 dans $D(\mathcal{O}_X)$. Puis nous avons 
\begin{align*}
R\Gamma(X \times Y, \mathcal{O}_{X \times Y})
& =
R\Gamma(X, R\text{pr}_{1, *}(\mathcal{O}_{X \times Y})) \\
& =
R\Gamma(X, \bigoplus\nolimits_{q \geq 0}
H^q(Y, \mathcal{O}_Y)[-q] \otimes_k \mathcal{O}_X) \\
& =
\bigoplus\nolimits_{q \geq 0}
R\Gamma(X,  H^q(Y, \mathcal{O}_Y) \otimes \mathcal{O}_X)[-q] \\
& =
\bigoplus\nolimits_{q \geq 0}
R\Gamma(X, \mathcal{O}_X) \otimes_k H^q(Y, \mathcal{O}_Y)[-q] \\
& =
\bigoplus\nolimits_{p, q \geq 0}
H^p(X, \mathcal{O}_X)[-p] \otimes_k H^q(Y, \mathcal{O}_Y)[-q]
\end{align*}
 comme voulu. La première égalité résulte de Leray pour $\text{pr}_1$
 (Cohomologie, lemme \ref{cohomology-lemma-before-Leray}).
La deuxième résulte de la décomposition de l'image directe totale donnée ci-dessus. La
troisième vient de ce que la cohomologie commute toujours avec les sommes directes finies (et
la cohomologie de $Y$ s'annule en degré suffisamment grand d'après
Cohomologie des Schémas, lemme 
\ref{coherent-lemma-vanishing-nr-affines-quasi-separated} ).
La quatrième vient de ce que la cohomologie sur $X$ commute avec les sommes
directes infinies d'après Cohomologie, lemme 
\ref{cohomology-lemma-quasi-separated-cohomology-colimit}.
Enfin, la dernière égalité résulte de la remarque faite plus haut sur la catégorie dérivée
d'un corps. 
\end{proof}








\section{Groupes de Picard des variétés} 
\label{varieties-section-picard-groups}

\noindent
 Dans cette section, nous recueillons quelques résultats élémentaires sur les groupes de
Picard des variétés algébriques.

\begin{lemma}
\label{varieties-lemma-change-rings-pic-pre}
 Soit $A \to B$ un homomorphisme d'anneaux fidèlement plat. Soit $X$ un schéma quasi-compact
et quasi-séparé sur $A$. Soit $\mathcal{L}$ un 
$\mathcal{O}_X$-module inversible dont l'image réciproque sur $X_B$ est triviale. Alors 
$H^0(X, \mathcal{L})$ et $H^0(X, \mathcal{L}^{\otimes -1})$ sont des 
$H^0(X, \mathcal{O}_X)$-modules inversibles,
et l'application de multiplication induit un isomorphisme 
$$
H^0(X, \mathcal{L}) \otimes_{H^0(X, \mathcal{O}_X)}
H^0(X, \mathcal{L}^{\otimes -1}) \longrightarrow
H^0(X, \mathcal{O}_X)
$$
\end{lemma}

\begin{proof}
 Notons $\mathcal{L}_B$ l'image réciproque de $\mathcal{L}$ sur $X_B$.
Choisissons un isomorphisme $\mathcal{L}_B \to \mathcal{O}_{X_B}$.
Posons $R = H^0(X, \mathcal{O}_X)$, $M = H^0(X, \mathcal{L})$ et considérons 
$M$ comme un $R$-module. Pour tout $\mathcal{O}_X$-module quasi-cohérent 
$\mathcal{F}$, dont $\mathcal{F}_B$ désigne le tiré en arrière sur $X_B$, il
existe un isomorphisme canonique 
$H^0(X_B, \mathcal{F}_B) = H^0(X, \mathcal{F}) \otimes_A B$,
voir Cohomologie des Schémas, lemme \ref{coherent-lemma-flat-base-change-cohomology}.
Ainsi, nous avons 
$$
M \otimes_R (R \otimes_A B) =
M \otimes_A B = H^0(X_B, \mathcal{L}_B) \cong
H^0(X_B, \mathcal{O}_{X_B}) = R \otimes_A B
$$
 Puisque $R \to R \otimes_A B$ est fidèlement plat (en tant que changement de
base de l'application fidèlement plate $A \to B$), nous concluons
que $M$ est un $R$-module inversible
d'après Algèbre, proposition \ref{algebra-proposition-ffdescent-finite-projectivity}.
De même, $N = H^0(X, \mathcal{L}^{\otimes -1})$ est un $R$-module inversible.
 Pour l'assertion sur les produits tensoriels, on utilise qu'elle devient vraie après
passage à $X_B$, ainsi que la fidèle platitude de $R \to R \otimes_A B$.
Quelques détails sont omis. 
\end{proof}

\begin{lemma}
\label{varieties-lemma-change-rings-pic}
 Soit $A \to B$ un homomorphisme d'anneaux fidèlement plat. Soit
$X$ un schéma sur $A$ tel que 
\begin{enumerate}
\item $X$ est quasi-compact et quasi-séparé, et 
\item $R = H^0(X, \mathcal{O}_X)$ est un anneau semi-local. 
\end{enumerate}
 Alors l'homomorphisme induit par image réciproque $\Pic(X) \to \Pic(X_B)$ est injectif. 
\end{lemma}

\begin{proof}
 Soit $\mathcal{L}$ un $\mathcal{O}_X$-module inversible
dont l'image réciproque $\mathcal{L}'$ sur $X_B$ est triviale.
Posons $M = H^0(X, \mathcal{L})$ et $N = H^0(X, \mathcal{L}^{\otimes - 1})$.
D'après le lemme \ref{varieties-lemma-change-rings-pic-pre}, les $R$-modules $M$ et $N$
 sont inversibles. Puisque $R$ est semi-local, $M \cong R$ et $N \cong R$,
voir Algèbre, lemme \ref{algebra-lemma-locally-free-semi-local-free}.
Choisissons des générateurs $s \in M$ et $t \in N$. Alors 
$st \in R = H^0(X, \mathcal{O}_X)$ est une unité d'après la dernière assertion
du lemme \ref{varieties-lemma-change-rings-pic-pre}.
Nous concluons que $s$ et $t$ définissent des trivialisations de $\mathcal{L}$ et 
$\mathcal{L}^{\otimes -1}$ sur $X$. 
\end{proof}

\begin{lemma}
\label{varieties-lemma-change-fields-pic}
 Soit $k'/k$ une extension de corps. Soit
$X$ un schéma sur $k$ tel que 
\begin{enumerate}
\item $X$ est quasi-compact et quasi-séparé, et 
\item $R = H^0(X, \mathcal{O}_X)$ est semi-local, par exemple, si $\dim_k R < \infty$. 
\end{enumerate}
 Alors l'homomorphisme induit par image réciproque $\Pic(X) \to \Pic(X_{k'})$ est injectif. 
\end{lemma}

\begin{proof}
 Cas particulier du lemme \ref{varieties-lemma-change-rings-pic}.
Si $\dim_k R < \infty$, alors 
$R$ est artinien et donc semi-local (Algèbre, lemmes 
\ref{algebra-lemma-finite-dimensional-algebra} et 
\ref{algebra-lemma-artinian-finite-nr-max}). 
\end{proof}

\begin{example}
\label{varieties-example-need-condition}
 Le lemme \ref{varieties-lemma-change-fields-pic} devient faux sans hypothèse supplémentaire
sur le schéma $X$ sur le corps $k$. Voici un exemple. Soit
$k$ un corps. Soit $t \in \mathbf{P}^1_k$ un point fermé.
Posons $X = \mathbf{P}^1 \setminus \{t\}$. Alors nous avons une surjection 
$$
\mathbf{Z} = \Pic(\mathbf{P}^1_k) \longrightarrow \Pic(X)
$$
 La première égalité résulte du lemme
des Diviseurs \ref{divisors-lemma-Pic-projective-space-UFD},
 et la surjectivité résulte du lemme
l'homomorphisme résulte de Diviseurs, lemme \ref{divisors-lemma-extend-invertible-module}
 (puisque $\mathbf{P}^1_k$ est lisse de dimension $1$ sur $k$ et
donc tous ses anneaux locaux sont des anneaux de valuation discrète). Si 
$\mathcal{L}$ est dans le noyau de l'application affichée,
alors $\mathcal{L} \cong \mathcal{O}_{\mathbf{P}^1_k}(nt)$
 pour un certain $n \in \mathbf{Z}$. Nous laissons au lecteur le soin
de montrer que 
$\mathcal{O}_{\mathbf{P}^1_k}(t) \cong \mathcal{O}_{\mathbf{P}^1_k}(d)$
 où $d = [\kappa(t) : k]$. Par conséquent 
$$
\Pic(X) = \mathbf{Z}/d\mathbf{Z}
$$
 Ainsi, si $t$ n'est pas un point $k$-rationnel, alors $d > 1$ et ce
groupe de Picard est non
nul. En revanche, si l'on étend le corps de base $k$ à une extension
de corps $k'$ telle qu'il existe un $k$-plongement 
$\kappa(t) \to k'$, alors $\mathbf{P}^1_{k'} \setminus X_{k'}$
 possède un point $k'$-rationnel $t'$. Par conséquent, 
$\mathcal{O}_{\mathbf{P}^1_{k'}}(1) = \mathcal{O}_{\mathbf{P}^1_{k'}}(t')$
 sera dans le noyau de l'application $\mathbf{Z} \to \Pic(X_{k'})$
 et il s'ensuivra de la même manière que ci-dessus que 
$\Pic(X_{k'}) = 0$. 
\end{example}

\noindent
 Le lemme suivant affirme que les «
modules inversibles rationnellement équivalents » sont isomorphes sur les variétés normales.

\begin{lemma}
\label{varieties-lemma-rational-equivalence-for-Pic}
 Soit $k$ un corps. Soit $X$ une variété normale sur $k$.
Soit $U \subset \mathbf{A}^n_k$ un sous-schéma ouvert possédant des points 
$k$-rationnels $p, q \in U(k)$. Pour chaque module
inversible $\mathcal{L}$ sur $X \times_{\Spec(k)} U$, les restrictions 
$\mathcal{L}|_{X \times p}$ et $\mathcal{L}|_{X \times q}$
 sont isomorphes. 
\end{lemma}

\begin{proof}
 Les fibres de $X \times_{\Spec(k)} U \to X$ sont des sous-schémas ouverts de
l'espace affine de dimension $n$ sur un corps. Par conséquent, ces fibres ont
des groupes de Picard triviaux
d'après Diviseurs, lemme \ref{divisors-lemma-open-subscheme-UFD}.
En appliquant Diviseurs, lemme \ref{divisors-lemma-in-image-pullback},
on voit que $\mathcal{L}$ est l'image réciproque d'un module
inversible $\mathcal{N}$ sur $X$. 
\end{proof}










\section{Unicité du corps de base} 
\label{varieties-section-base-field}

\noindent
 La phrase « soit $X$ un schéma sur $k$ » signifie que $X$ est un
schéma muni d'un morphisme $X \to \Spec(k)$. On peut alors
se demander si le corps $k$ est déterminé de manière unique par le schéma $X$.
Bien sûr, ce n'est pas le cas, car par exemple 
$\mathbf{A}^1_{\mathbf{C}}$ que nous considérons habituellement comme un schéma
sur le corps $\mathbf{C}$ des nombres complexes, pourrait également être considéré comme
un schéma sur $\mathbf{Q}$. Mais que se passe-t-il si nous demandons que le morphisme 
$X \to \Spec(k)$ ne se factorise pas sous la forme 
$X \to \Spec(k') \to \Spec(k)$ pour une extension de corps
non triviale $k'/k$ ? En d'autres termes, nous demandons que $k$ soit
d'une certaine manière maximal pour que $X$ soit défini sur $k$.

\medskip\noindent
 Un exemple montrant que cela ne garantit toujours pas l'unicité de $k$
 est le schéma 
$$
X =
\Spec\left(
\mathbf{Q}(x)[y]\left[\frac{1}{P(y)}, P \in \mathbf{Q}[y], P \not = 0\right]
\right)
$$
 À première vue, cela semble être un schéma sur $\mathbf{Q}(x)$, mais
à y regarder de plus près, il est clair que c'est aussi un schéma sur $\mathbf{Q}(y)$.
De plus, les corps $\mathbf{Q}(x)$ et $\mathbf{Q}(y)$ sont des sous-corps
de $R = \Gamma(X, \mathcal{O}_X)$ qui sont maximaux parmi les sous-corps de 
$R$ (détails omis). En particulier, $\mathbf{Q}(x)$ et 
$\mathbf{Q}(y)$ sont tous deux maximaux au sens ci-dessus. Remarquons que les deux morphismes 
$X \to \Spec(\mathbf{Q}(x))$
 et $X \to \Spec(\mathbf{Q}(y))$ sont « essentiellement de type fini
» (c'est-à-dire que l'homomorphisme d'anneaux correspondant est essentiellement de type fini).
Ainsi $X$ est un schéma noethérien de dimension finie, c'est-à-dire
qu'il n'est pas complètement pathologique.

\medskip\noindent
 Un autre problème qui peut empêcher l'unicité est que le schéma $X$ peut
ne pas être réduit. Dans ce cas, il peut y avoir de nombreux morphismes différents
de $X$ vers le spectre d'un corps donné. À titre d'exemple explicite, considérons
l'algèbre des nombres duaux 
$D = \mathbf{C}[y]/(y^2) = \mathbf{C} \oplus \epsilon \mathbf{C}$.
Étant donnée toute dérivation $\theta : \mathbf{C} \to \mathbf{C}$ sur $\mathbf{Q}$,
nous obtenons un homomorphisme d'anneaux 
$$
\mathbf{C} \longrightarrow D, \quad
c \longmapsto c + \epsilon \theta(c).
$$
 Le sous-corps de $\mathbf{C}$ sur lequel tous ces homomorphismes sont les
mêmes est la clôture algébrique de $\mathbf{Q}$. Cela signifie que prendre
l'intersection de tous les corps sur lesquels $X$ peut être défini peut
être un très petit corps si $X$ n'est pas réduit.

\medskip\noindent
Une observation à ce sujet est la suivante : étant donné un corps
$k$ et deux sous-corps $k_1, k_2$ de $k$ tels que $k$ soit fini
sur $k_1$ et sur $k_2$, il n'est alors en général {\it pas} vrai
que $k$ soit fini sur $k_1 \cap k_2$. Un exemple est le corps
$k = \mathbf{Q}(t)$ et ses sous-corps $k_1 = \mathbf{Q}(t^2)$ et
$\mathbf{Q}((t + 1)^2)$. En effet, nous avons $k_1 \cap k_2 = \mathbf{Q}$
dans ce cas. Ainsi, dans la suite, il faut être prudent lorsqu'on prend
des intersections de corps.

\medskip\noindent
 Cela dit, nous montrons maintenant que si $X$ est localement de type
fini sur un corps, alors une certaine unicité s'applique. Voici le
résultat précis.

\begin{proposition}
\label{varieties-proposition-unique-base-field}
 Soit $X$ un schéma. Soient $a : X \to \Spec(k_1)$ et 
$b : X \to \Spec(k_2)$ des morphismes de $X$ vers des spectres de corps.
Supposons que $a, b$ soient localement de type fini et que 
$X$ soit réduit et connexe. Alors nous avons 
$k_1' = k_2'$, où $k_i' \subset \Gamma(X, \mathcal{O}_X)$
 est la clôture intégrale de $k_i$ dans $\Gamma(X, \mathcal{O}_X)$. 
\end{proposition}

\begin{proof}
 Admettons d'abord le résultat lorsque $X$ est quasi-compact (nous établirons
ce cas ci-dessous). Comme 
$X$ est localement de type fini sur un corps, il est localement noethérien,
voir Morphismes, lemme \ref{morphisms-lemma-finite-type-noetherian}.
En particulier, cela signifie qu'il est localement
connexe, les composantes connexes des ouverts sont
ouvertes, et les intersections d'ouverts quasi-compacts sont quasi-compactes,
voir Propriétés, lemme \ref{properties-lemma-Noetherian-topology},
Topologie, lemme \ref{topology-lemma-locally-connected},
Topologie, section \ref{topology-section-noetherian},
et Topologie, lemme \ref{topology-lemma-constructible-Noetherian-space}.
Choisissons un recouvrement ouvert $X = \bigcup_{i \in I} U_i$
 tel que chaque $U_i$ soit quasi-compact et connexe. Pour
chaque $i$, soit $K_i \subset \mathcal{O}_X(U_i)$ la clôture
intégrale de $k_1$ et de $k_2$.
Pour chaque paire $i, j \in I$ nous décomposons 
$$
U_i \cap U_j = \coprod U_{i, j, l}
$$
 en ses composantes connexes, qui sont en nombre fini. Écrivons 
$K_{i, j, l} \subset \mathcal{O}(U_{i, j, l})$
 pour la clôture intégrale de $k_1$ et de $k_2$. D'après
le lemme \ref{varieties-lemma-integral-closure-ground-field},
les anneaux $K_i$ et $K_{i, j, l}$ sont des corps.
Nous affirmons maintenant que $k_1'$ et $k_2'$ sont tous deux égaux au noyau de l'application 
$$
\prod K_i \longrightarrow \prod K_{i, j, l}, \quad
(x_i)_i \longmapsto x_i|_{U_{i, j, l}} - x_j|_{U_{i, j, l}}
$$
ce qui donne le résultat. En effet, il est clair que
$k_1'$ est contenu dans ce noyau. D'autre part, supposons que
$(x_i)_i$ soit dans le noyau. Par la condition de faisceau,
$(x_i)_i$ correspond à $f \in \mathcal{O}(X)$.
Choisissons $i_0 \in I$ et un polynôme unitaire $P(T) \in k_1[T]$
tel que $P(x_{i_0}) = 0$. Nous affirmons alors que $P(f) =
0$, ce qui prouve que $f \in k_1$. Il suffit de
montrer que $P(x_i) = 0$ pour tout $i$. Choisissons $i \in I$.
Comme $X$ est connexe, il existe une suite $i_0, i_1, \ldots, i_n =
i \in I$ telle que $U_{i_t} \cap U_{i_{t + 1}} \not = \emptyset$.
Cela signifie que, pour chaque $t$, il existe un $l_t$
tel que $x_{i_t}$ et $x_{i_{t + 1}}$ se projettent sur le même
élément du corps $K_{i_t, i_{t + 1}, l_t}$. Par conséquent, si
$P(x_{i_t}) = 0$, alors $P(x_{i_{t + 1}}) = 0$. Par induction, en
partant de $P(x_{i_0}) = 0$, nous obtenons $P(x_i) = 0$, comme annoncé.

\medskip\noindent
Pour terminer la démonstration du lemme, nous prouvons
le lemme sous l'hypothèse supplémentaire que $X$ est
quasi-compact. D'après le lemme
\ref{varieties-lemma-integral-closure-ground-field}, après avoir
remplacé $k_i$ par $k_i'$, on peut supposer que $k_i$
est intégralement clos dans $\Gamma(X, \mathcal{O}_X)$.
Cela implique que $\mathcal{O}(X)^*/k_i^*$ est un groupe
abélien de type fini, voir la proposition
\ref{varieties-proposition-units-general}. Soit $k_{12} = k_1 \cap k_2$ un
sous-anneau de $\mathcal{O}(X)$. Notons que $k_{12}$ est un corps. Puisque
$$
k_1^*/k_{12}^* \longrightarrow \mathcal{O}(X)^*/k_2^*
$$
on voit que $k_1^*/k_{12}^*$ est également un groupe abélien
de type fini. Par conséquent, il existe $\alpha_1, \ldots,
\alpha_n \in k_1^*$ tels que chaque élément $\lambda \in k_1$ ait la forme
$$
\lambda = c \alpha_1^{e_1} \ldots \alpha_n^{e_n}
$$
pour certains $e_i \in \mathbf{Z}$ et $c \in
k_{12}$. En particulier, l'homomorphisme d'anneaux
$$
k_{12}[x_1, \ldots, x_n, \frac{1}{x_1 \ldots x_n}] \longrightarrow k_1, \quad
x_i \longmapsto \alpha_i
$$
est surjectif. D'après le Nullstellensatz de Hilbert, Algèbre,
théorème \ref{algebra-theorem-nullstellensatz}, nous concluons
que $k_1$ est une extension finie de $k_{12}$. De la même
manière, nous concluons que $k_2$ est une extension finie de
$k_{12}$. En particulier, à la fois $k_1$ et $k_2$ sont contenus
dans la clôture intégrale $k_{12}'$ de $k_{12}$ dans $\Gamma(X,
\mathcal{O}_X)$. Mais puisque $k_{12}'$ est un corps selon le
lemme \ref{varieties-lemma-integral-closure-ground-field} et que nous avons
choisi $k_i$ pour être intégralement clos dans $\Gamma(X,
\mathcal{O}_X)$, nous concluons que $k_1 = k_{12} = k_2$, comme annoncé.
\end{proof}





\section{Automorphismes}
\label{varieties-section-automorphisms}

\noindent
Une section sur les automorphismes des schémas sur les
corps. Pour certaines informations sur les automorphismes
(infinitésimaux) des courbes, voir Courbes Algébriques,
section \ref{curves-section-vector-fields} et Modules de
Courbes, section \ref{moduli-curves-section-finite-aut}.

\begin{lemma}
\label{varieties-lemma-automorphism}
Supposons que $X$ soit un schéma de type fini réduit sur un corps $k$.
Supposons que $f : X \to X$ soit un automorphisme sur $k$ qui induit
l'identité sur l'espace topologique sous-jacent de $X$. Alors
\begin{enumerate}
\item $f^*\mathcal{F} \cong \mathcal{F}$ pour
chaque $\mathcal{O}_X$-module cohérent, et
\item si $\dim(Z) > 0$ pour chaque composante
irréductible $Z \subset X$, alors $f$ est l'identité.
\end{enumerate}
\end{lemma}

\begin{proof}
La partie (1) découle de la partie (2) et du fait que les
composantes connexes de $X$ de dimension $0$ sont des spectres de corps.

\medskip\noindent
Supposons que $Z \subset X$ soit une composante irréductible considérée
comme un sous-schéma fermé intègre. Évidemment, $f(Z) \subset Z$
et $f|_Z : Z \to Z$ sont un automorphisme sur $k$ qui induit
l'identité sur l'espace topologique sous-jacent de $Z$.
Puisque $X$ est réduit, il suffit de montrer que les flèches $f|_Z : Z \to Z$
sont l'identité. Cela nous ramène au cas discuté dans le paragraphe suivant.

\medskip\noindent
Supposons que $X$ soit irréductible de dimension $> 0$.
Choisissons un ouvert affine non vide $U \subset X$. Puisque
$f(U) \subset U$ et puisque $U \subset X$ est schématiquement
dense, il suffit de prouver que $f|_U : U \to U$ est l'identité.

\medskip\noindent
Supposons que $X = \Spec(A)$ soit affine, irréductible, de dimension $> 0$ et que
$k$ soit un corps infini. Soit $g \in A$ non constant. L'ensemble
$$
S = \bigcup\nolimits_{\lambda \in k} V(g - \lambda)
$$
est dense dans $X$ parce que c'est l'image réciproque du
sous-ensemble dense $\mathbf{A}^1_k(k)$ par le morphisme
non constant $g : X \to \mathbf{A}^1_k$. Si $x \in S$,
alors l'image $g(x)$ de $g$ dans $\kappa(x)$ est dans
l'image de $k \to \kappa(x)$. Par conséquent, $f^\sharp :
\kappa(x) \to \kappa(x)$ fixe $g(x)$. Ainsi, l'image de
$f^\sharp(g)$ dans $\kappa(x)$ est égale à $g(x)$. Nous en concluons que
$$
S \subset V(g - f^\sharp(g))
$$
et puisque $X$ est réduit et $S$ est dense, nous concluons
$g=f^\sharp(g)$. Cela prouve $f^\sharp = \text{id}_A$ car $A$ est
engendré comme une $k$-algèbre par les éléments $g$ comme ci-dessus (détails
omis ; indication : l'ensemble des fonctions constantes est un $k$-sous-espace
vectoriel de dimension finie de $A$). Nous concluons que $f = \text{id}_X$.

\medskip\noindent
Supposons que $X = \Spec(A)$ soit affine, irréductible, de
dimension $> 0$ et que $k$ soit un corps fini. Si pour tout
sous-schéma fermé intègre de dimension $1$, $C \subset X$, la
restriction $f|_C : C \to C$ est l'identité, alors $f$ est
l'identité. Cela nous ramène au cas où $X$ est une courbe. Une courbe
sur un corps fini possède un groupe d'automorphismes fini (détails
omis). Ainsi, $f$ a un ordre fini, disons $n$. Ensuite, nous choisissons
$g : X \to \mathbf{A}^1_k$ non constant comme ci-dessus et considérons
$$
S = \{x \in X\text{ fermé tel que }[\kappa(g(x)) : k]
\text{ est premier à }n\}
$$
En argumentant comme précédemment, nous trouvons que
$S$ est dense dans $X$. Puisque pour $x \in X$ fermé,
l'application $f^\sharp : \kappa(x) \to \kappa(x)$
est un automorphisme d'ordre divisant $n$, nous
voyons que pour $x \in S$, cet automorphisme agit
trivialement sur le sous-corps engendré par l'image de
$g$ dans $\kappa(x)$. Ainsi nous concluons que $S \subset
V(g - f^\sharp(g))$ et on conclut comme précédemment.
\end{proof}






\section{Caractéristiques d'Euler--Poincaré}
\label{varieties-section-euler}

\noindent
Dans cette section, nous prouvons quelques propriétés élémentaires des
caractéristiques d'Euler--Poincaré des faisceaux cohérents sur des schémas propres sur un corps.

\begin{definition}
\label{varieties-definition-euler-characteristic}
Soit $k$ un corps. Soit $X$ un schéma propre sur $k$. Soit
$\mathcal{F}$ un $\mathcal{O}_X$-module cohérent. Dans cette
situation, la {\it caractéristique d'Euler--Poincaré de $\mathcal{F}$ } est l'entier
$$
\chi(X, \mathcal{F}) = \sum\nolimits_i (-1)^i \dim_k H^i(X, \mathcal{F}).
$$
Pour la justification de la formule, voir ci-dessous.
\end{definition}

\noindent
Dans le cas de la définition, seul un nombre fini des
espaces vectoriels $H^i(X, \mathcal{F})$ sont non nuls
(Cohomologie des Schémas, lemme
\ref{coherent-lemma-quasi-coherence-higher-direct-images})
et chacun de ces espaces est de dimension finie (Cohomologie
des Schémas, lemme
\ref{coherent-lemma-proper-over-affine-cohomology-finite}). Ainsi
$\chi(X, \mathcal{F}) \in \mathbf{Z}$ est bien défini. Observons que cette
définition dépend du corps $k$ et pas seulement de la paire $(X, \mathcal{F})$.

\begin{lemma}
\label{varieties-lemma-euler-characteristic-additive}
Soit $k$ un corps. Supposons que $X$ soit un schéma propre
sur $k$. Supposons que $0 \to \mathcal{F}_1 \to \mathcal{F}_2 \to
\mathcal{F}_3 \to 0$ soit une suite exacte courte de modules cohérents sur $X$. Alors
$$
\chi(X, \mathcal{F}_2) = \chi(X, \mathcal{F}_1) + \chi(X, \mathcal{F}_3)
$$
\end{lemma}

\begin{proof}
Considérons la longue suite exacte de cohomologie
$$
0 \to H^0(X, \mathcal{F}_1) \to H^0(X, \mathcal{F}_2) \to
H^0(X, \mathcal{F}_3) \to H^1(X, \mathcal{F}_1) \to \ldots
$$
associée à la suite exacte courte du lemme. Le théorème
du rang en algèbre linéaire montre que
$$
0 = \dim H^0(X, \mathcal{F}_1) - \dim H^0(X, \mathcal{F}_2)
+ \dim H^0(X, \mathcal{F}_3) - \dim H^1(X, \mathcal{F}_1) + \ldots
$$
Cela implique immédiatement le lemme.
\end{proof}

\begin{lemma}
\label{varieties-lemma-chi-tensor-finite}
Soit $k$ un corps. Supposons que $X$ soit un
schéma propre sur $k$. Supposons que $\mathcal{F}$ soit un
faisceau cohérent avec $\dim(\text{Supp}(\mathcal{F})) \leq 0$. Alors
\begin{enumerate}
\item $\mathcal{F}$ est engendré par ses sections
globales, \item $H^0(X, \mathcal{F}) =
\bigoplus_{x \in \text{Supp}(\mathcal{F})}
\mathcal{F}_x$, \item $H^i(X, \mathcal{F}) = 0$ pour
$i > 0$, \item $\chi(X, \mathcal{F}) = \dim_k
H^0(X, \mathcal{F})$, et \item $\chi(X, \mathcal{F}
\otimes \mathcal{E}) = n\chi(X, \mathcal{F})$ pour
chaque module localement libre $\mathcal{E}$ de rang $n$.
\end{enumerate}
\end{lemma}

\begin{proof}
Par Cohomologie des Schémas, lemme
\ref{coherent-lemma-coherent-support-closed}, on voit que
$\mathcal{F} = i_*\mathcal{G}$ où $i : Z \to X$ est l'inclusion
du support schématique de $\mathcal{F}$ et où $\mathcal{G}$ est
un $\mathcal{O}_Z$-module cohérent. Par définition du support
schématique, l'espace topologique sous-jacent de $Z$ est
$\text{Supp}(\mathcal{F})$. Puisque la dimension de $Z$ est
$0$, on voit que $Z$ est affine (Propriétés, lemme
\ref{properties-lemma-locally-Noetherian-dimension-0}). Ainsi
$\mathcal{G}$ est engendré par ses sections globales et les groupes de
cohomologie supérieurs de $\mathcal{G}$ sont nuls (Cohomologie
des Schémas, lemme
\ref{coherent-lemma-quasi-coherent-affine-cohomology-zero}). En
fait, par le lemme \ref{varieties-lemma-algebraic-scheme-dim-0}, le schéma $Z$
est une réunion finie disjointe de spectres d'anneaux locaux
artiniens. Par conséquent, $H^0(Z, \mathcal{G}) =
\bigoplus_{z \in Z} \mathcal{G}_z$. Les cohomologies de $\mathcal{F}$
et $\mathcal{G}$ coïncident d'après Cohomologie des Schémas, lemme
\ref{coherent-lemma-relative-affine-cohomology}. Ainsi $H^i(X,
\mathcal{F}) = 0$ pour $i > 0$ et $H^0(X, \mathcal{F}) = H^0(Z,
\mathcal{G})$. En particulier, (3) est vérifiée. Pour $z \in Z$
correspondant à $x \in \text{Supp}(\mathcal{F})$ nous avons
$\mathcal{G}_z = (i_*\mathcal{G})_x = \mathcal{F}_x$. Nous concluons que (2)
est vérifiée. Évidemment, (2) implique (1). Nous obtenons (4) par définition de la
caractéristique d'Euler--Poincaré $\chi(X, \mathcal{F})$ et (3). Par la formule de
projection (Cohomologie, lemme \ref{cohomology-lemma-projection-formula}) nous avons
$$
i_*(\mathcal{G} \otimes i^*\mathcal{E}) = \mathcal{F} \otimes \mathcal{E}.
$$
Puisque $Z$ a pour dimension $0$, le faisceau localement
libre $i^*\mathcal{E}$ est isomorphe à $\mathcal{O}_Z^{\oplus
n}$ et, en raisonnant comme ci-dessus, on voit que (5) est vérifiée.
\end{proof}

\begin{lemma}
\label{varieties-lemma-euler-characteristic-extend-base-field}
Supposons que $k'/k$ soit une extension de corps. Supposons que $X$ soit un
schéma propre sur $k$. Supposons que $\mathcal{F}$ soit un faisceau
cohérent sur $X$. Supposons que $\mathcal{F}'$ soit l'image réciproque de
$\mathcal{F}$ sur $X_{k'}$. Alors $\chi(X, \mathcal{F}) = \chi(X_{k'}, \mathcal{F}')$.
\end{lemma}

\begin{proof}
En effet,
$$
H^i(X_{k'}, \mathcal{F}') = H^i(X, \mathcal{F}) \otimes_k k'
$$
par changement de base plat, voir Cohomologie des
Schémas, lemme \ref{coherent-lemma-flat-base-change-cohomology}.
\end{proof}

\begin{lemma}
\label{varieties-lemma-euler-characteristic-morphism}
Soit $k$ un corps. Supposons que $f : Y \to X$ soit un morphisme de schémas
propres sur $k$. Supposons que $\mathcal{G}$ soit un $\mathcal{O}_Y$-module cohérent. Alors
$$
\chi(Y, \mathcal{G}) = \sum (-1)^i \chi(X, R^if_*\mathcal{G})
$$
\end{lemma}

\begin{proof}
La formule a du sens : les faisceaux
$R^if_*\mathcal{G}$ sont cohérents et seul un nombre
fini d'entre eux est non nul, voir Cohomologie des
schémas, proposition
\ref{coherent-proposition-proper-pushforward-coherent} et lemme
\ref{coherent-lemma-quasi-coherence-higher-direct-images}. D'après
Cohomologie, lemme \ref{cohomology-lemma-Leray}, il existe une suite spectrale avec
$$
E_2^{p, q} = H^p(X, R^qf_*\mathcal{G})
$$
convergeant vers $H^{p + q}(Y, \mathcal{G})$. Par la finitude de la
cohomologie sur $X$, les termes $E_2^{p, q}$ non nuls sont en nombre
fini et chaque $E_2^{p, q}$ est un espace vectoriel de dimension finie.
Il en va de même des termes $E_r^{p, q}$ pour $r \geq 2$, et
$$
\sum (-1)^{p + q} \dim_k E_r^{p, q}
$$
est indépendant de $r$. Puisque pour $r$ suffisamment grand
nous avons $E_r^{p, q} = E_\infty^{p, q}$ et puisque la
convergence signifie qu'il existe une filtration sur $H^n(Y,
\mathcal{G})$ dont les morceaux gradués sont $E_\infty^{p, q}$
avec $p + q = n$ (c'est le sens de la convergence de la suite
spectrale), nous concluons. Voir aussi le lemme plus général du chapitre
Homologie générale, lemme \ref{homology-lemma-biregular-ss-relation-in-K0}.
\end{proof}








\section{Espace projectif}
\label{varieties-section-projective-space}

\noindent
Quelques résultats sur l'espace projectif sur un corps.

\begin{lemma}
\label{varieties-lemma-projective-space-smooth}
\begin{slogan}
L'espace projectif est lisse.
\end{slogan}
Supposons que $k$ soit un corps et $n \geq 0$. Alors
$\mathbf{P}^n_k$ est une variété projective lisse de dimension $n$ sur $k$.
\end{lemma}

\begin{proof}
Démonstration omise.
\end{proof}

\begin{lemma}
\label{varieties-lemma-intersection-in-affine-space}
Supposons que $k$ soit un corps et $n \geq 0$. Supposons que $X, Y
\subset \mathbf{A}^n_k$ soit des sous-ensembles fermés. Supposons que $X$
et $Y$ soient équidimensionnels, $\dim(X) = r$ et $\dim(Y) = s$. Alors
chaque composante irréductible de $X \cap Y$ a une dimension $\geq r + s - n$.
\end{lemma}

\begin{proof}
Considérons le sous-schéma fermé $X \times Y \subset
\mathbf{A}^{2n}_k$ où nous utilisons les coordonnées
$x_1, \ldots, x_n, y_1, \ldots, y_n$. Alors $X \cap Y = X
\times Y \cap V(x_1 - y_1, \ldots, x_n - y_n)$. Supposons
que $t \in X \cap Y \subset X \times Y$ soit un point
fermé. Par le lemme
\ref{varieties-lemma-dimension-product-locally-algebraic}, nous avons $\dim_t(X
\times Y) = \dim(X) + \dim(Y)$. Ainsi $\dim(\mathcal{O}_{X \times Y,
t}) = r + s$ par le lemme \ref{varieties-lemma-dimension-locally-algebraic}.
Par Algèbre, lemme \ref{algebra-lemma-one-equation}, nous concluons que
$$
\dim(\mathcal{O}_{X \cap Y, t}) =
\dim(\mathcal{O}_{X \times Y, t}/(x_1 - y_1, \ldots, x_n - y_n)) \geq
r + s - n
$$
Cela implique le résultat par le lemme \ref{varieties-lemma-dimension-locally-algebraic}.
\end{proof}

\begin{lemma}
\label{varieties-lemma-intersection-in-projective-space}
Supposons que $k$ soit un corps et $n \geq 0$. Supposons que
$X, Y \subset \mathbf{P}^n_k$ soit des sous-ensembles fermés
non vides. Si $\dim(X) = r$ et $\dim(Y) = s$ et $r + s \geq n$,
alors $X \cap Y$ est non vide et $\dim(X \cap Y) \geq r + s - n$.
\end{lemma}

\begin{proof}
Écrivons $\mathbf{A}^{n + 1} = \Spec(k[x_0, \ldots, x_n])$ et
$\mathbf{P}^n = \text{Proj}(k[T_0, \ldots, T_n])$.
Considérons le morphisme $\pi : \mathbf{A}^{n + 1}
\setminus \{0\} \to \mathbf{P}^n$ qui envoie $(x_0,
\ldots, x_n)$ vers le point $[x_0 : \ldots : x_n]$. Plus
précisément, c'est le morphisme associé à la paire
$(\mathcal{O}_{\mathbf{A}^{n + 1} \setminus \{0\}}, (x_0,
\ldots, x_n))$, voir Constructions, lemme
\ref{constructions-lemma-projective-space}. Sur l'ouvert affine standard
$D_+(T_i)$, nous obtenons le morphisme associé à l'homomorphisme d'anneaux
$$
k\left[\frac{T_0}{T_i}, \ldots, \frac{T_n}{T_i}\right]
\longrightarrow
k\left[T_0, \ldots, T_n, \frac{1}{T_i}\right] \cong
k\left[\frac{T_0}{T_i}, \ldots, \frac{T_n}{T_i}\right]
\left[T_i, \frac{1}{T_i}\right]
$$
qui est surjectif et lisse de dimension relative $1$, à fibres
irréductibles (détails omis). Par conséquent, $\pi^{-1}(X)$
et $\pi^{-1}(Y)$ sont des sous-ensembles fermés non vides de
dimension $r + 1$ et $s + 1$. Choisissons une composante
irréductible $V \subset \pi^{-1}(X)$ de dimension $r + 1$ et une
composante irréductible $W \subset \pi^{-1}(Y)$ de dimension $s + 1$.
Observons que cela implique que $V$ et $W$ contiennent chaque fibre
de $\pi$ qu'elles rencontrent (puisque $\pi$ a des fibres
irréductibles de dimension $1$ et puisque le lemme
\ref{varieties-lemma-dimension-fibres-locally-algebraic} dit que les fibres de $V
\to \pi(V)$ et $W \to \pi(W)$ ont pour dimension $\geq 1$). Supposons
que $\overline{V}$ et $\overline{W}$ soient l'adhérence de $V$ et $W$
dans $\mathbf{A}^{n + 1}$. Puisque $0 \in \mathbf{A}^{n + 1}$ est dans
l'adhérence de chaque fibre de $\pi$, on voit que $0 \in
\overline{V} \cap \overline{W}$. Par le lemme
\ref{varieties-lemma-intersection-in-affine-space}, nous avons $\dim(\overline{V} \cap
\overline{W}) \geq r + s - n + 1$. En raisonnant comme ci-dessus et en appliquant à
nouveau le lemme \ref{varieties-lemma-dimension-fibres-locally-algebraic}, nous concluons que
$\pi(V \cap W) \subset X \cap Y$ est de dimension au moins $r + s - n$, comme annoncé.
\end{proof}

\begin{lemma}
\label{varieties-lemma-equation-codim-1-in-projective-space}
Soit $k$ un corps. Supposons que $Z \subset
\mathbf{P}^n_k$ soit un sous-schéma fermé sans points immergés
et dont chaque composante irréductible de $Z$ est de dimension $n - 1$. Alors
l'idéal $I(Z) \subset k[T_0, \ldots, T_n]$ correspondant à $Z$ est principal.
\end{lemma}

\begin{proof}
Ceci est un cas particulier des Diviseurs, lemme
\ref{divisors-lemma-equation-codim-1-in-projective-space}.
\end{proof}




\section{Faisceaux cohérents sur l'espace projectif}
\label{varieties-section-coherent-Pn}

\noindent
Dans cette section, nous prouvons certains résultats sur la
cohomologie des faisceaux cohérents sur $\mathbf{P}^n$ sur un
corps, qui peuvent être trouvés dans \cite{Mum}. Ceux-ci seront
utiles plus tard lors de la discussion des schémas de Quot et de Hilbert.

\subsection{Préliminaires}
\label{varieties-subsection-preliminaries}

\noindent
Soit $k$ un corps, $n \geq 1$, $d \geq 1$, et soit $s \in
\Gamma(\mathbf{P}_k^n, \mathcal{O}(d))$ une section non
nulle. Dans cette section, nous noterons $\mathcal{O}(d)$ le
$d$-ième tordu du faisceau structural sur l'espace projectif
(Constructions, définitions
\ref{constructions-definition-twist} et
\ref{constructions-definition-projective-space}). Comme $\mathbf{P}^n_k$
est une variété, cette section est régulière ; ainsi $s$ est une section
régulière de $\mathcal{O}(d)$ et définit un diviseur de Cartier effectif
$H = Z(s) \subset \mathbf{P}^n_k$, voir Diviseurs, section
\ref{divisors-section-effective-Cartier-divisors}. Un tel diviseur $H$ est
appelé une {\it hypersurface} et, si $d = 1$, il est appelé un {\it hyperplan}.

\begin{lemma}
\label{varieties-lemma-hyperplane}
Soit $k$ un corps. Supposons
que $n \geq 1$. Soit $i : H \to \mathbf{P}^n_k$
un hyperplan. Alors il existe un isomorphisme
$$
\varphi : \mathbf{P}^{n - 1}_k \longrightarrow H
$$
tel que l'image réciproque de $i^*\mathcal{O}(1)$ soit isomorphe à $\mathcal{O}(1)$.
\end{lemma}

\begin{proof}
Nous avons $\mathbf{P}^n_k = \text{Proj}(k[T_0,
\ldots, T_n])$. La section $s$ correspond à une
forme homogène dans $T_0, \ldots, T_n$ de degré
$1$, voir Cohomologie des schémas, section
\ref{coherent-section-cohomology-projective-space}.
Écrivons $s = \sum a_i T_i$. Constructions, le lemme
\ref{constructions-lemma-closed-in-projective-space}
montre que $H = \text{Proj}(k[T_0, \ldots, T_n]/I)$
pour l'idéal gradué $I$ défini en posant $I_d$ égal
au noyau de l'application $\Gamma(\mathbf{P}^n_k,
\mathcal{O}(d)) \to \Gamma(H, i^*\mathcal{O}(d))$.
D'après notre construction de $Z(s)$ dans Diviseurs,
définition \ref{divisors-definition-zero-scheme-s},
on voit que sur $D_{+}(T_j)$ l'idéal de $H$ est
engendré par $\sum a_i T_i/T_j$ dans l'anneau de polynômes
$k[T_0/T_j, \ldots, T_n/T_j]$. Ainsi, il est clair que
$I$ est l'idéal engendré par $\sum a_i T_i$. Remarquez que
$$
k[T_0, \ldots, T_n]/I = k[T_0, \ldots, T_n]/(\sum a_i T_i) \cong
k[S_0, \ldots, S_{n - 1}]
$$
en tant qu'anneaux gradués. Par exemple, si $a_n \not = 0$,
alors envoyer $S_i$ sur la classe de $T_i$ convient.
Nous obtenons l'isomorphisme annoncé par fonctorialité de
$\text{Proj}$. L'égalité des tordus du faisceau structural
résulte par exemple du chapitre Constructions, lemme
\ref{constructions-lemma-surjective-graded-rings-generated-degree-1-map-proj}.
\end{proof}

\begin{lemma}
\label{varieties-lemma-exact-sequence-induction}
Supposons que $k$ soit un corps infini.
Supposons que $n \geq 1$. Soit $\mathcal{F}$ un
module cohérent sur $\mathbf{P}^n_k$. Alors il
existe une section non nulle $s \in
\Gamma(\mathbf{P}^n_k, \mathcal{O}(1))$ et une suite exacte courte
$$
0 \to \mathcal{F}(-1) \to \mathcal{F} \to i_*\mathcal{G} \to 0
$$
où $i : H \to \mathbf{P}^n_k$ est l'hyperplan $H$
associé à $s$ et $\mathcal{G} = i^*\mathcal{F}$.
\end{lemma}

\begin{proof}
L'application $\mathcal{F}(-1) \to \mathcal{F}$ provient
de Constructions, Équation (
\ref{constructions-equation-multiply-on-sheaf} ) avec $n =
1$, $m = -1$ et la section $s$ de $\mathcal{O}(1)$. Voyons à
quoi ressemble cette application si nous la restreignons à
$D_{+}(T_0)$. Écrivons $D_{+}(T_0) = \Spec(k[x_1, \ldots,
x_n])$ avec $x_i = T_i/T_0$. Identifions
$\mathcal{O}(1)|_{D_{+}(T_0)}$ avec $\mathcal{O}$ à l'aide de la
section $T_0$. Par conséquent, si $s = \sum a_iT_i$ alors
$s|_{D_{+}(T_0)} = a_0 + \sum a_ix_i$ avec l'identification
choisie ci-dessus. De plus, supposons que
$\mathcal{F}|_{D_{+}(T_0)}$ corresponde au $k[x_1, \ldots, x_n]$-module
fini $M$. Via l'identification $\mathcal{F}(-1) = \mathcal{F} \otimes
\mathcal{O}(-1)$ et notre trivialisation choisie de $\mathcal{O}(1)$,
on voit que $\mathcal{F}(-1)$ correspond également à $M$. Ainsi,
restreindre $\mathcal{F}(-1) \to \mathcal{F}$ à $D_{+}(T_0)$ donne l'application
$$
M \xrightarrow{a_0 + \sum a_ix_i} M
$$
Pour que la flèche soit injective, il suffit de choisir $a_0
+ \sum a_ix_i$ n'appartenant à aucun des idéaux premiers
associés de $M$, voir Algèbre, lemme
\ref{algebra-lemma-ass-zero-divisors}. D'après Algèbre, lemme
\ref{algebra-lemma-finite-ass}, l'ensemble $\text{Ass}(M)$ des
idéaux premiers associés de $M$ est fini. Notons que pour
$\mathfrak p \in \text{Ass}(M)$, l'intersection $\mathfrak p \cap
\{a_0 + \sum a_i x_i\}$ est un sous-espace vectoriel propre de $k$.
Nous en concluons qu'il existe une famille finie de sous-espaces
vectoriels propres $V_1, \ldots, V_m \subset \Gamma(\mathbf{P}^n_k,
\mathcal{O}(1))$ tels que si nous prenons $s$ en dehors de $\bigcup
V_i$, alors la multiplication par $s$ est injective sur
$D_{+}(T_0)$. De même pour la restriction à $D_{+}(T_j)$ pour $j =
1, \ldots, n$. Comme $k$ est infini, une union finie de sous-espaces
vectoriels propres n'est jamais égale à l'espace tout entier ; nous
pouvons donc choisir $s$ de sorte que l'application soit
injective. Le conoyau de $\mathcal{F}(-1) \to \mathcal{F}$ est annulé
par $\Im(s : \mathcal{O}(-1) \to \mathcal{O})$ qui est l'idéal
faisceau de $H$ selon Diviseurs, définition
\ref{divisors-definition-zero-scheme-s}. Ainsi, nous obtenons $\mathcal{G}$ sur $H$ en
utilisant Cohomologie des Schémas, lemme \ref{coherent-lemma-i-star-equivalence}.
\end{proof}

\begin{remark}
\label{varieties-remark-exact-sequence-induction}
Supposons que $k$ soit un corps infini. Soit $n \geq 1$. Étant
donné un nombre fini de modules cohérents $\mathcal{F}_i$ sur
$\mathbf{P}^n_k$, nous pouvons choisir un seul $s \in
\Gamma(\mathbf{P}^n_k, \mathcal{O}(1))$ de sorte que l'énoncé du lemme
\ref{varieties-lemma-exact-sequence-induction} fonctionne pour chacun d'eux. Pour
le prouver, il suffit d'appliquer le lemme à $\bigoplus \mathcal{F}_i$.
\end{remark}

\begin{remark}
\label{varieties-remark-exact-sequence-induction-cohomology}
Dans la situation des lemmes \ref{varieties-lemma-hyperplane} et
\ref{varieties-lemma-exact-sequence-induction}, nous avons $H \cong
\mathbf{P}^{n - 1}_k$ avec les tordus de Serre
$\mathcal{O}_H(d) = i^*\mathcal{O}_{\mathbf{P}^n_k}(d)$. Pour
chaque $d \in \mathbf{Z}$, nous avons une suite exacte courte
$$
0 \to \mathcal{F}(d - 1) \to \mathcal{F}(d) \to i_*(\mathcal{G}(d)) \to 0
$$
En effet, le produit tensoriel par
$\mathcal{O}_{\mathbf{P}^n_k}(d)$ est un foncteur exact et,
d'après la formule de projection (Cohomologie, lemme
\ref{cohomology-lemma-projection-formula}), nous avons
$i_*(\mathcal{G}(d)) = i_*\mathcal{G} \otimes
\mathcal{O}_{\mathbf{P}^n_k}(d)$. Nous obtenons des suites exactes longues correspondantes
$$
H^i(\mathbf{P}^n_k, \mathcal{F}(d - 1)) \to
H^i(\mathbf{P}^n_k, \mathcal{F}(d)) \to
H^i(H, \mathcal{G}(d)) \to
H^{i + 1}(\mathbf{P}^n_k, \mathcal{F}(d - 1))
$$
Cela résulte de ce qui précède et de l'égalité
$H^i(\mathbf{P}^n_k, i_*\mathcal{G}(d)) = H^i(H,
\mathcal{G}(d))$ d'après Cohomologie des Schémas, lemme
\ref{coherent-lemma-relative-affine-cohomology} (les immersions fermées sont affines).
\end{remark}





\subsection{Régularité}
\label{varieties-subsection-regularity}

\noindent
Voici la définition.

\begin{definition}
\label{varieties-definition-regularity}
Soit $k$ un corps. Supposons que $n \geq 0$.
Soit $\mathcal{F}$ un faisceau cohérent sur
$\mathbf{P}^n_k$. Nous disons que $\mathcal{F}$ est {\it $m$-régulier} si
$$
H^i(\mathbf{P}^n_k, \mathcal{F}(m - i)) = 0
$$
pour $i = 1, \ldots, n$.
\end{definition}

\noindent
Notons que $\mathcal{F} = \mathcal{O}(d)$ est $m$-régulier si
et seulement si $d \geq -m$. Cela découle du calcul des groupes
de cohomologie dans Cohomologie des schémas, Équation (
\ref{coherent-equation-identify} ). En effet, on voit que
$H^n(\mathbf{P}^n_k, \mathcal{O}(d)) = 0$ si et seulement si $d \geq -n$.

\begin{lemma}
\label{varieties-lemma-m-regular-extend-base-field}
Supposons que $k'/k$ soit une extension de corps. Supposons que
$n \geq 0$. Soit $\mathcal{F}$ un faisceau cohérent sur
$\mathbf{P}^n_k$. Supposons que $\mathcal{F}'$ soit l'image réciproque de
$\mathcal{F}$ sur $\mathbf{P}^n_{k'}$. Alors $\mathcal{F}$
est $m$-régulier si et seulement si $\mathcal{F}'$ est $m$-régulier.
\end{lemma}

\begin{proof}
Cela résulte de l'isomorphisme
$$
H^i(\mathbf{P}^n_{k'}, \mathcal{F}') =
H^i(\mathbf{P}^n_k, \mathcal{F}) \otimes_k k'
$$
par changement de base plat, voir Cohomologie des
Schémas, lemme \ref{coherent-lemma-flat-base-change-cohomology}.
\end{proof}

\begin{lemma}
\label{varieties-lemma-m-regular}
Dans la situation du lemme
\ref{varieties-lemma-exact-sequence-induction}, si $\mathcal{F}$ est $m$-régulier,
alors $\mathcal{G}$ est $m$-régulier sur $H \cong \mathbf{P}^{n - 1}_k$.
\end{lemma}

\begin{proof}
Rappelons que $H^i(\mathbf{P}^n_k, i_*\mathcal{G}) = H^i(H,
\mathcal{G})$ d'après Cohomologie des Schémas, lemme
\ref{coherent-lemma-relative-affine-cohomology}. Ainsi, pour $i \geq 1$, nous obtenons
$$
H^i(\mathbf{P}^n_k, \mathcal{F}(m - i)) \to
H^i(H, \mathcal{G}(m - i)) \to
H^{i + 1}(\mathbf{P}^n_k, \mathcal{F}(m - 1 - i))
$$
par la Remarque
\ref{varieties-remark-exact-sequence-induction-cohomology}. Le lemme suit.
\end{proof}

\begin{lemma}
\label{varieties-lemma-m-regular-up}
Soit $k$ un corps. Supposons que $n
\geq 0$. Soit $\mathcal{F}$ un faisceau cohérent
sur $\mathbf{P}^n_k$. Si $\mathcal{F}$ est
$m$-régulier, alors $\mathcal{F}$ est $(m + 1)$-régulier.
\end{lemma}

\begin{proof}
Nous procédons par récurrence sur $n$. Si $n = 0$, chaque
faisceau est $m$-régulier pour tout $m$, il n'y a rien à
démontrer. D'après le lemme
\ref{varieties-lemma-m-regular-extend-base-field}, nous pouvons remplacer $k$ par
une extension infinie et supposer que $k$ est infini. Nous pouvons alors
appliquer le lemme \ref{varieties-lemma-exact-sequence-induction}. D'après le lemme
\ref{varieties-lemma-m-regular}, nous savons que $\mathcal{G}$ est $m$-régulier.
Par induction sur $n$, on voit que $\mathcal{G}$ est $(m +
1)$-régulier. En considérant la suite exacte longue de cohomologie associée à la suite
$$
0 \to \mathcal{F}(m - i) \to \mathcal{F}(m + 1 - i)
\to i_*\mathcal{G}(m + 1 - i) \to 0
$$
comme dans la remarque
\ref{varieties-remark-exact-sequence-induction-cohomology}, on déduit
pour $i \geq 1$ l'annulation de $H^i(\mathbf{P}^n_k,
\mathcal{F}(m + 1 - i))$ à partir de l'annulation de
$H^i(\mathbf{P}^n_k, \mathcal{F}(m - i))$ et $H^i(\mathbf{P}^n_k, \mathcal{G}(m + 1 - i))$.
\end{proof}

\begin{lemma}
\label{varieties-lemma-m-regular-multiply}
Soit $k$ un corps. Supposons que $n \geq 0$. Soit
$\mathcal{F}$ un faisceau cohérent sur $\mathbf{P}^n_k$. Si
$\mathcal{F}$ est $m$-régulier, alors l'application de multiplication
$$
H^0(\mathbf{P}^n_k, \mathcal{F}(m)) \otimes_k
H^0(\mathbf{P}^n_k, \mathcal{O}(1))
\longrightarrow
H^0(\mathbf{P}^n_k, \mathcal{F}(m + 1))
$$
est surjective.
\end{lemma}

\begin{proof}
Supposons que $k'/k$ soit une extension de corps. Supposons que
$\mathcal{F}'$ soit comme dans le lemme
\ref{varieties-lemma-m-regular-extend-base-field}. D'après le chapitre Cohomologie des schémas,
lemme \ref{coherent-lemma-flat-base-change-cohomology}, l'application linéaire obtenue
par changement de base vers $k'$ coïncide avec l'application linéaire associée au
faisceau $\mathcal{F}'$. Puisque $k \to k'$ est fidèlement plat, il suffit de
prouver le lemme sur $k'$, c'est-à-dire que l'on peut supposer $k$ infini.

\medskip\noindent
Supposons que $k$ soit infini. Nous prouvons le lemme par
induction sur $n$. Le cas $n = 0$ est trivial car
$\mathcal{O}(1) \cong \mathcal{O}$ est généré par $T_0$. Pour
$n > 0$, appliquons le lemme \ref{varieties-lemma-exact-sequence-induction}
et tensorisons la suite par $\mathcal{O}(m + 1)$ ; on obtient
$$
0 \to \mathcal{F}(m) \xrightarrow{s} \mathcal{F}(m + 1) \to
i_*\mathcal{G}(m + 1) \to 0
$$
voir Remarque
\ref{varieties-remark-exact-sequence-induction-cohomology}. Soit $t \in
H^0(\mathbf{P}^n_k, \mathcal{F}(m + 1))$. Par induction, l'image
$\overline{t} \in H^0(H, \mathcal{G}(m + 1))$ est l'image de $\sum g_i
\otimes \overline{s}_i$ avec $\overline{s}_i \in \Gamma(H,
\mathcal{O}(1))$ et $g_i \in H^0(H, \mathcal{G}(m))$. Puisque $\mathcal{F}$
est $m$-régulier, nous avons $H^1(\mathbf{P}^n_k, \mathcal{F}(m - 1)) = 0$,
d'où la suite exacte longue de cohomologie associée à la suite exacte courte
$$
0 \to \mathcal{F}(m - 1) \xrightarrow{s} \mathcal{F}(m) \to
i_*\mathcal{G}(m) \to 0
$$
montre que nous pouvons relever $g_i$ en $f_i \in
H^0(\mathbf{P}^n_k, \mathcal{F}(m))$. Nous pouvons également
relever $\overline{s}_i$ en $s_i \in H^0(\mathbf{P}^n_k,
\mathcal{O}(1))$ (voir la preuve du lemme \ref{varieties-lemma-hyperplane} par
exemple). Après avoir soustrait l'image de $\sum f_i \otimes s_i$ de
$t$, on voit que nous pouvons supposer $\overline{t} = 0$. Mais
cela signifie exactement que $t$ est l'image de $f \otimes s$ pour un
certain $f \in H^0(\mathbf{P}^n_k, \mathcal{F}(m))$, comme annoncé.
\end{proof}

\begin{lemma}
\label{varieties-lemma-m-regular-globally-generated}
Soit $k$ un corps. Supposons que $n
\geq 0$. Soit $\mathcal{F}$ un faisceau cohérent
sur $\mathbf{P}^n_k$. Si $\mathcal{F}$ est
$m$-régulier, alors $\mathcal{F}(m)$ est engendré par ses sections globales.
\end{lemma}

\begin{proof}
Pour tout $d \gg 0$, le faisceau $\mathcal{F}(d)$ est engendré par ses
sections globales. Cela découle par exemple de la première partie de
Cohomologie des Schémas, lemme
\ref{coherent-lemma-coherent-projective}. Choisissons $d \geq m$ de telle sorte
que $\mathcal{F}(d)$ soit engendré par ses sections globales. Choisissons une base $f_1,
\ldots, f_r \in H^0(\mathbf{P}^n_k, \mathcal{F})$. D'après le lemme
\ref{varieties-lemma-m-regular-multiply}, chaque élément $f \in H^0(\mathbf{P}^n_k,
\mathcal{F}(d))$ peut être écrit comme $f = \sum P_if_i$ pour un certain $P_i \in
k[T_0, \ldots, T_n]$ homogène de degré $d - m$. Puisque les sections $f$ génèrent
$\mathcal{F}(d)$, il en découle que les sections $f_i$ génèrent $\mathcal{F}(m)$.
\end{proof}

\subsection{Polynômes de Hilbert}
\label{varieties-subsection-hilbert}

\noindent
Le lemme suivant sera englobé par le lemme plus général
\ref{varieties-lemma-numerical-polynomial-from-euler}.

\begin{lemma}
\label{varieties-lemma-hilbert-polynomial}
Soit $k$ un corps. Supposons que $n \geq 0$. Soit
$\mathcal{F}$ un faisceau cohérent sur $\mathbf{P}^n_k$. La fonction
$$
d \longmapsto \chi(\mathbf{P}^n_k, \mathcal{F}(d))
$$
est un polynôme.
\end{lemma}

\begin{proof}
Nous procédons par récurrence sur $n$. Si $n =
0$, alors $\mathbf{P}^n_k = \Spec(k)$ et
$\mathcal{F}(d) = \mathcal{F}$. Par conséquent, dans
ce cas, la fonction est constante, c'est-à-dire un
polynôme de degré $0$. Supposons $n > 0$. Par le lemme
\ref{varieties-lemma-euler-characteristic-extend-base-field},
on peut supposer que $k$ est infini. Appliquons
le lemme \ref{varieties-lemma-exact-sequence-induction}. En
appliquant le lemme
\ref{varieties-lemma-euler-characteristic-additive} aux suites tordues $0 \to
\mathcal{F}(d - 1) \to \mathcal{F}(d) \to i_*\mathcal{G}(d) \to 0$, nous obtenons
$$
\chi(\mathbf{P}^n_k, \mathcal{F}(d)) -
\chi(\mathbf{P}^n_k, \mathcal{F}(d - 1)) =
\chi(H, \mathcal{G}(d))
$$
Voir la Remarque
\ref{varieties-remark-exact-sequence-induction-cohomology}. Comme $H
\cong \mathbf{P}^{n - 1}_k$ par induction, le côté droit est un
polynôme. On conclut en observant que toute fonction $f :
\mathbf{Z} \to \mathbf{Z}$ ayant la propriété que l'application
$d \mapsto f(d) - f(d - 1)$ est un polynôme, est elle-même un
polynôme. Nous omettons la preuve de ce fait (indication : comparer
avec Algèbre, lemme \ref{algebra-lemma-numerical-polynomial}).
\end{proof}

\begin{definition}
\label{varieties-definition-hilbert-polynomial}
Soit $k$ un corps. Soit $n \geq 0$. Soit $\mathcal{F}$
un faisceau cohérent sur $\mathbf{P}^n_k$. La fonction
$d \mapsto \chi(\mathbf{P}^n_k, \mathcal{F}(d))$ est
appelée le {\it polynôme de Hilbert} de $\mathcal{F}$.
\end{definition}

\noindent
Le polynôme de Hilbert a des coefficients dans
$\mathbf{Q}$ et pas en général dans $\mathbf{Z}$. Par exemple,
le polynôme de Hilbert de $\mathcal{O}_{\mathbf{P}^n_k}$ est
$$
d \longmapsto {d + n \choose n} = \frac{d^n}{n!} + \ldots
$$
Cela découle du lemme suivant et du fait que
$$
H^0(\mathbf{P}^n_k, \mathcal{O}_{\mathbf{P}^n_k}(d)) = k[T_0, \ldots, T_n]_d
$$
(sa partie de degré $d$), dont la dimension sur $k$ est ${d + n \choose n}$.

\begin{lemma}
\label{varieties-lemma-hilbert-polynomial-H0}
Soit $k$ un corps. Supposons que $n \geq 0$.
Soit $\mathcal{F}$ un faisceau cohérent sur
$\mathbf{P}^n_k$ avec le polynôme de Hilbert $P \in \mathbf{Q}[t]$. Alors
$$
P(d) = \dim_k H^0(\mathbf{P}^n_k, \mathcal{F}(d))
$$
pour tous les $d \gg 0$.
\end{lemma}

\begin{proof}
Cela résulte de l'annulation de la cohomologie
des tordus suffisamment élevés de $\mathcal{F}$.
Voir Cohomologie des schémas,
lemme \ref{coherent-lemma-coherent-projective}.
\end{proof}



\subsection{Bornitude des quotients}
\label{varieties-subsection-boundedness}

\noindent
Dans cette sous-section, nous bornons la
régularité des quotients d'un faisceau cohérent donné
sur $\mathbf{P}^n$ en termes du polynôme de Hilbert.

\begin{lemma}
\label{varieties-lemma-bound-quotients-free}
Soit $k$ un corps. Soit $n \geq 0$. Soit $r
\geq 1$. Soit $P \in \mathbf{Q}[t]$. Il existe un entier $m$
dépendant de $n$, $r$ et $P$ ayant la propriété suivante : si
$$
0 \to \mathcal{K} \to \mathcal{O}^{\oplus r} \to \mathcal{F} \to 0
$$
est une suite exacte courte de faisceaux cohérents
sur $\mathbf{P}^n_k$ et $\mathcal{F}$ a pour polynôme
de Hilbert $P$, alors $\mathcal{K}$ est $m$-régulier.
\end{lemma}

\begin{proof}
Nous procédons par récurrence sur $n$. Si $n = 0$, alors
$\mathbf{P}^n_k = \Spec(k)$ et tout module cohérent est
$0$-régulier et tout épimorphisme induit une surjection
sur les sections globales. Supposons $n > 0$. Considérons une
suite exacte comme dans le lemme. Supposons que $P' \in
\mathbf{Q}[t]$ est le polynôme $P'(t) = P(t) - P(t - 1)$.
Supposons que $m'$ soit l'entier qui fonctionne pour $n - 1$,
$r$ et $P'$. D'après les lemmes
\ref{varieties-lemma-m-regular-extend-base-field} et
\ref{varieties-lemma-euler-characteristic-extend-base-field}, nous pouvons
remplacer $k$ par une extension de corps, donc nous pouvons
supposer que $k$ est infini. Appliquons le lemme
\ref{varieties-lemma-exact-sequence-induction} au faisceau cohérent
$\mathcal{F}$. Le polynôme de Hilbert de $\mathcal{F}' =
i^*\mathcal{F}$ est $P'$ (voir la preuve du lemme
\ref{varieties-lemma-hilbert-polynomial}). Puisque $i^*$ est exact à droite, nous
voyons que $\mathcal{F}'$ est un quotient de $\mathcal{O}_H^{\oplus r}
= i^*\mathcal{O}^{\oplus r}$. Ainsi, l'hypothèse d'induction s'applique
à $\mathcal{F}'$ sur $H \cong \mathbf{P}^{n - 1}_k$ (lemme
\ref{varieties-lemma-hyperplane}). Notons que l'application $\mathcal{K}(-1) \to
\mathcal{K}$ est injective comme $\mathcal{K} \subset \mathcal{O}^{\oplus
r}$ et a un conoyau $i_*\mathcal{H}$ où $\mathcal{H} = i^*\mathcal{K}$.
Par le lemme du serpent (Homologie, lemme \ref{homology-lemma-snake}) nous
obtenons un diagramme commutatif avec des colonnes et des lignes exactes.
$$
\xymatrix{
& 0 \ar[d] & 0 \ar[d] & 0 \ar[d] \\
0 \ar[r] &
\mathcal{K}(-1) \ar[r] \ar[d] &
\mathcal{O}^{\oplus r}(-1) \ar[r] \ar[d] &
\mathcal{F}(-1) \ar[d] \ar[r] & 0 \\
0 \ar[r] &
\mathcal{K} \ar[r] \ar[d] &
\mathcal{O}^{\oplus r} \ar[r] \ar[d] &
\mathcal{F} \ar[d] \ar[r] & 0\\
0 \ar[r] &
i_*\mathcal{H} \ar[r] \ar[d] &
i_*\mathcal{O}_H^{\oplus r} \ar[r] \ar[d] &
i_*\mathcal{F}' \ar[r] \ar[d] & 0 \\
& 0 & 0 & 0
}
$$
Ainsi, l'hypothèse d'induction s'applique à la suite exacte
$0 \to \mathcal{H} \to \mathcal{O}_H^{\oplus r} \to
\mathcal{F}' \to 0$ sur $H \cong \mathbf{P}^{n - 1}_k$ (lemme
\ref{varieties-lemma-hyperplane}) et $\mathcal{H}$ est $m'$-régulier.
Rappelons que cela implique que $\mathcal{H}$ est
$d$-régulier pour tous les $d \geq m'$ (lemme \ref{varieties-lemma-m-regular-up}).

\medskip\noindent
Soient $i \geq 2$ et $d \geq m'$. Il découle de la
suite exacte longue de cohomologie associée à la
colonne de gauche du diagramme ci-dessus et de la
annulation de $H^{i - 1}(H, \mathcal{H}(d))$ que l'application
$$
H^i(\mathbf{P}^n_k, \mathcal{K}(d - 1))
\longrightarrow
H^i(\mathbf{P}^n_k, \mathcal{K}(d))
$$
est injective. Comme ces groupes sont nuls pour
$d \gg 0$ (Cohomologie des Schémas, lemme
\ref{coherent-lemma-coherent-projective}), nous en
concluons que $H^i(\mathbf{P}^n_k, \mathcal{K}(d))$
sont nuls pour tous les $d \geq m'$ et $i \geq 2$.

\medskip\noindent
Nous devons encore contrôler $H^1$. Tout d'abord, nous observons que toutes les applications
$$
H^1(\mathbf{P}^n_k, \mathcal{K}(m' - 1)) \to
H^1(\mathbf{P}^n_k, \mathcal{K}(m')) \to
H^1(\mathbf{P}^n_k, \mathcal{K}(m' + 1)) \to \ldots
$$
sont surjectives en raison de l'annulation de $H^1(H,
\mathcal{H}(d))$ pour $d \geq m'$. Supposons que $d > m'$ soit tel que
$$
H^1(\mathbf{P}^n_k, \mathcal{K}(d - 1))
\longrightarrow
H^1(\mathbf{P}^n_k, \mathcal{K}(d))
$$
est injective. Ensuite, $H^0(\mathbf{P}^n_k,
\mathcal{K}(d)) \to H^0(H, \mathcal{H}(d))$ est
surjective. Considérons le diagramme commutatif
$$
\xymatrix{
H^0(\mathbf{P}^n_k, \mathcal{K}(d)) \otimes_k
H^0(\mathbf{P}^n_k, \mathcal{O}(1))
\ar[r] \ar[d] &
H^0(\mathbf{P}^n_k, \mathcal{K}(d + 1)) \ar[d] \\
H^0(H, \mathcal{H}(d)) \otimes_k
H^0(H, \mathcal{O}_H(1))
\ar[r] &
H^0(H, \mathcal{H}(d + 1))
}
$$
D'après le lemme \ref{varieties-lemma-m-regular-multiply}, nous voyons
que la flèche horizontale du bas est surjective. Par conséquent,
la flèche verticale de droite est surjective. Nous concluons que
$$
H^1(\mathbf{P}^n_k, \mathcal{K}(d))
\longrightarrow
H^1(\mathbf{P}^n_k, \mathcal{K}(d + 1))
$$
est injective. Par induction, on voit que
$$
H^1(\mathbf{P}^n_k, \mathcal{K}(d - 1)) \to
H^1(\mathbf{P}^n_k, \mathcal{K}(d)) \to
H^1(\mathbf{P}^n_k, \mathcal{K}(d + 1)) \to \ldots
$$
sont toutes injectives et nous concluons que
$H^1(\mathbf{P}^n_k, \mathcal{K}(d - 1)) = 0$ en
raison du fait que ces groupes finissent par s'annuler.
Ainsi, les dimensions des groupes $H^1(\mathbf{P}^n_k,
\mathcal{K}(d))$ pour $d \geq m'$ diminuent strictement
jusqu'à devenir nulles. Il s'ensuit que la régularité
de $\mathcal{K}$ est bornée par $m' + \dim_k
H^1(\mathbf{P}^n_k, \mathcal{K}(m'))$. D'autre part, par la
annulation des groupes de cohomologie supérieurs, nous avons
$$
\dim_k H^1(\mathbf{P}^n_k, \mathcal{K}(m')) = 
- \chi(\mathbf{P}^n_k, \mathcal{K}(m')) +
\dim_k H^0(\mathbf{P}^n_k, \mathcal{K}(m'))
$$
Notons que le $H^0$ a une dimension bornée par la dimension
de $H^0(\mathbf{P}^n_k, \mathcal{O}^{\oplus r}(m'))$ qui
est au maximum $r{n + m' \choose n}$ si $m' > 0$ et zéro
sinon. Enfin, le terme $\chi(\mathbf{P}^n_k,
\mathcal{K}(m'))$ est égal à $r{n + m' \choose n} - P(m')$. Cela
donne une borne du type voulu et achève la démonstration du lemme.
\end{proof}







\section{Frobenius}
\label{varieties-section-frobenius}

\noindent
Supposons que $p$ soit un nombre premier. Si $X$ est un schéma, alors nous disons
que `` $X$ a pour caractéristique $p$ '', ou que `` $X$ est de caractéristique $p$
'', ou que `` $X$ est en caractéristique $p$ '' si $p$ est nul dans $\mathcal{O}_X$.

\begin{definition}
\label{varieties-definition-absolute-frobenius}
Soit $p$ un nombre premier. Soit $X$ un schéma de caractéristique $p$.
Le {\it Frobenius absolu de $X$} est le morphisme $F_X : X \to X$
donné par l'identité sur l'espace topologique sous-jacent et dont
$F_X^\sharp : \mathcal{O}_X \to \mathcal{O}_X$ est donné par $g \mapsto g^p$.
\end{definition}

\noindent
Cela a du sens car pour tout anneau $A$ de caractéristique $p$,
l'application $F_A : A \to A$, $a \mapsto a^p$ est un
endomorphisme d'anneau qui induit l'identité sur $\Spec(A)$. De
plus, si $A$ est local, alors $F_A$ est un homomorphisme local.
On voit ainsi que le Frobenius absolu de $X$ est
un endomorphisme de $X$ dans la catégorie des schémas. Il s'avère
que le Frobenius absolu définit un endomorphisme du
foncteur identité sur la catégorie des schémas en caractéristique $p$.

\begin{lemma}
\label{varieties-lemma-frobenius-endomorphism-identity}
Supposons que $p > 0$ soit un nombre premier.
Supposons que $f : X \to Y$ soit un morphisme de
schémas en caractéristique $p$. Alors le diagramme
$$
\xymatrix{
X \ar[d]_f \ar[r]_{F_X} & X \ar[d]^f \\
Y \ar[r]^{F_Y} & Y
}
$$
est commutatif.
\end{lemma}

\begin{proof}
Cela découle du fait algébrique trivial suivant : si
$\varphi : A \to B$ est un homomorphisme d'anneaux de
caractéristique $p$, alors $\varphi(a^p) = \varphi(a)^p$.
\end{proof}

\begin{lemma}
\label{varieties-lemma-frobenius}
Supposons que $p > 0$ soit un nombre premier. Supposons que $X$
soit un schéma en caractéristique $p$. Alors le Frobenius absolu
$F_X : X \to X$ est un homéomorphisme universel, est intégral, et
induit des extensions purement inséparables des corps résiduels.
\end{lemma}

\begin{proof}
Cela découle des résultats correspondants pour l'endomorphisme de
Frobenius $F_A : A \to A$ d'un anneau $A$ de caractéristique $p > 0$.
Voir la discussion dans Algèbre, section
\ref{algebra-section-universal-homeomorphism}, par exemple le lemme \ref{algebra-lemma-p-ring-map}.
\end{proof}

\noindent
Si nous travaillons avec des schémas sur une base fixe,
alors il existe une version relative du morphisme de Frobenius.

\begin{definition}
\label{varieties-definition-relative-frobenius}
Supposons que $p > 0$ soit un nombre premier. Supposons que $S$ soit un schéma en
caractéristique $p$. Supposons que $X$ soit un schéma sur $S$. Nous définissons
$$
X^{(p)} = X^{(p/S)} = X \times_{S, F_S} S
$$
vu comme un schéma sur $S$. En appliquant le lemme
\ref{varieties-lemma-frobenius-endomorphism-identity}, nous voyons
qu'il existe un unique morphisme $F_{X/S} : X
\longrightarrow X^{(p)}$ sur $S$ s'insérant dans le diagramme commutatif
$$
\xymatrix{
X \ar[rr]_{F_{X/S}} \ar[rrd] \ar@/^1em/[rrrr]^{F_X}
& & X^{(p)} \ar[rr] \ar[d] & & X \ar[d] \\
& & S \ar[rr]^{F_S} & & S
}
$$
où le carré de droite est cartésien. Le morphisme $F_{X/S}$
est appelé le {\it morphisme de Frobenius relatif de $X/S$}.
\end{definition}

\noindent
Remarquez que $X \mapsto X^{(p)}$ est un foncteur ; c'est le foncteur de
changement de base pour le morphisme de Frobenius absolu $F_S : S \to S$. Nous
obtenons les analogues des lemmes précédents pour le morphisme de Frobenius relatif.

\begin{lemma}
\label{varieties-lemma-relative-frobenius-endomorphism-identity}
Supposons que $p > 0$ soit un nombre premier. Supposons que
$S$ soit un schéma en caractéristique $p$. Supposons que $f : X
\to Y$ soit un morphisme de schémas sur $S$. Alors le diagramme
$$
\xymatrix{
X \ar[d]_f \ar[r]_{F_{X/S}} & X^{(p)} \ar[d]^{f^{(p)}} \\
Y \ar[r]^{F_{Y/S}} & Y^{(p)}
}
$$
est commutatif.
\end{lemma}

\begin{proof}
Cela découle du lemme
\ref{varieties-lemma-frobenius-endomorphism-identity} et des définitions.
\end{proof}

\begin{lemma}
\label{varieties-lemma-relative-frobenius}
Supposons que $p > 0$ soit un nombre premier. Supposons que
$S$ soit un schéma en caractéristique $p$. Supposons que $X$
soit un schéma sur $S$. Alors le Frobenius relatif $F_{X/S} : X
\to X^{(p)}$ est un homéomorphisme universel, est intégral, et
induit des extensions purement inséparables des corps résiduels.
\end{lemma}

\begin{proof}
D'après le lemme \ref{varieties-lemma-frobenius}, les morphismes $F_X : X
\to X$ et le changement de base $h : X^{(p)} \to X$ de $F_S$ sont
des homéomorphismes universels. Puisque $h \circ F_{X/S} = F_X$,
nous concluons que $F_{X/S}$ est un homéomorphisme universel
(Morphismes, lemme
\ref{morphisms-lemma-permanence-universal-homeomorphism}). D'après les lemmes
\ref{morphisms-lemma-universal-homeomorphism} et
\ref{morphisms-lemma-universally-injective} de Morphismes, nous concluons que $F_{X/S}$ possède également les autres propriétés.
\end{proof}

\begin{lemma}
\label{varieties-lemma-relative-frobenius-omega}
Supposons que $p > 0$ soit un nombre premier. Supposons que $S$ soit un schéma en
caractéristique $p$. Supposons que $X$ soit un schéma sur $S$. Alors $\Omega_{X/S} = \Omega_{X/X^{(p)}}$.
\end{lemma}

\begin{proof}
Cela se traduit par le fait algébrique suivant. Supposons
que $A \to B$ soit un homomorphisme d'anneaux de
caractéristique $p$. Posons $B' = B \otimes_{A, F_A} A$ et
considérons l'homomorphisme d'anneaux $F_{B/A} : B' \to B$,
$b \otimes a \mapsto b^pa$. Alors notre assertion est que
$\Omega_{B/A} = \Omega_{B/B'}$. En effet,
$\text{d}(b^pa) = 0$ si $\text{d} : B \to \Omega_{B/A}$ est la
dérivation universelle et donc $\text{d}$ est une dérivation $B'$.
\end{proof}

\begin{lemma}
\label{varieties-lemma-relative-frobenius-finite}
Supposons que $p > 0$ soit un nombre premier. Supposons que $S$ soit un
schéma en caractéristique $p$. Supposons que $X$ soit un schéma sur
$S$. Si $X \to S$ est localement de type fini, alors $F_{X/S}$ est fini.
\end{lemma}

\begin{proof}
Cela se traduit par le fait algébrique suivant. Supposons que $A
\to B$ soit un homomorphisme de type fini d'anneaux de
caractéristique $p$. Posons $B' = B \otimes_{A, F_A} A$ et
considérons l'homomorphisme d'anneaux $F_{B/A} : B' \to B$, $b
\otimes a \mapsto b^pa$. Alors notre assertion est que $F_{B/A}$ est
fini. En effet, si $x_1, \ldots, x_n \in B$ sont des générateurs sur
$A$, alors $x_i$ est intégral sur $B'$ parce que $x_i^p = F_{B/A}(x_i
\otimes 1)$. Par conséquent, $F_{B/A} : B' \to B$ est fini d'après le
lemme d'Algèbre \ref{algebra-lemma-characterize-finite-in-terms-of-integral}.
\end{proof}

\begin{lemma}
\label{varieties-lemma-geometrically-reduced-p}
Supposons que $k$ soit un corps de caractéristique $p > 0$. Supposons que $X$ soit un
schéma sur $k$. Alors $X$ est géométriquement réduit si et seulement si $X^{(p)}$ est réduit.
\end{lemma}

\begin{proof}
Considérons le Frobenius absolu $F_k : k \to k$. Alors
$F_k(k) = k^p$, en d'autres termes, $F_k : k \to k$ est
s'identifie à l'inclusion de $k$ dans $k^{1/p}$. Ainsi, le
lemme découle du lemme \ref{varieties-lemma-geometrically-reduced}.
\end{proof}

\begin{lemma}
\label{varieties-lemma-inseparable-deg-p-smooth}
Supposons que $k$ soit un corps de caractéristique $p > 0$. Supposons
que $X$ soit une variété sur $k$. Les énoncés suivants sont équivalents
\begin{enumerate}
\item $X^{(p)}$ est réduit, \item $X$ est
géométriquement réduit, \item il y a un
ouvert non vide $U \subset X$ lisse sur $k$.
\end{enumerate}
Dans ce cas, $X^{(p)}$ est une variété sur $k$ et $F_{X/k} : X
\to X^{(p)}$ est un morphisme dominant fini de degré $p^{\dim(X)}$.
\end{lemma}

\begin{proof}
L'équivalence de (1) et (2) résulte du
lemme \ref{varieties-lemma-geometrically-reduced-p}. Nous avons
vu dans le lemme
\ref{varieties-lemma-geometrically-reduced-dense-smooth-open} que (2) implique (3). Si
(3) est vrai, alors $U$ est géométriquement réduit (voir
par exemple le lemme
\ref{varieties-lemma-geometrically-regular-smooth}) et donc $X$ est
géométriquement réduit par le lemme
\ref{varieties-lemma-generic-points-geometrically-reduced}. On voit ainsi que (1), (2) et (3) sont équivalents.

\medskip\noindent
Supposons que (1), (2) et (3) soient vérifiés. Puisque
$F_{X/k}$ est un homéomorphisme (lemme
\ref{varieties-lemma-relative-frobenius}), on voit que $X^{(p)}$ est
une variété. Ensuite, $F_{X/k}$ est fini par le lemme
\ref{varieties-lemma-relative-frobenius-finite}. Il est dominant car il
est surjectif. Pour calculer le degré (Morphismes, définition
\ref{morphisms-definition-degree}), il suffit de calculer le
degré de $F_{U/k} : U \to U^{(p)}$ (comme $F_{U/k} = F_{X/k}|_U$
par le lemme
\ref{varieties-lemma-relative-frobenius-endomorphism-identity}). Après avoir
réduit un peu $U$, nous pouvons supposer qu'il existe un morphisme
étale $h : U \to \mathbf{A}^n_k$, voir Morphismes, lemme
\ref{morphisms-lemma-smooth-etale-over-affine-space}. On a nécessairement $n =
\dim(U)$ parce que $\mathbf{A}^n_k \to \Spec(k)$ est lisse de
dimension relative $n$, le morphisme étale $h$ est lisse de
dimension relative $0$, et $U \to \Spec(k)$ est lisse de dimension
relative $\dim(U)$ et les dimensions relatives s'additionnent
correctement (Morphismes, lemme
\ref{morphisms-lemma-composition-relative-dimension-d}). Remarquons que $h$ est un
morphisme génériquement fini et dominant de variétés, et donc $\deg(h)$ est défini. D'après le
lemme \ref{varieties-lemma-relative-frobenius-endomorphism-identity}, nous avons un diagramme commutatif
$$
\xymatrix{
U \ar[rr]_{F_{X/k}} \ar[d]_h & & U^{(p)} \ar[d]^{h^{(p)}} \\
\mathbf{A}^n_k \ar[rr]^{F_{\mathbf{A}^n_k/k}} & &
(\mathbf{A}^n_k)^{(p)}
}
$$
Puisque $h^{(p)}$ est un changement de base de $h$, il
est également étale et il s'ensuit que $h^{(p)}$ est
également un morphisme dominant génériquement fini de
variétés. Le degré de $h^{(p)}$ est le degré de
l'extension $k(X^{(p)})/k((\mathbf{A}^n_k)^{(p)})$, qui est
le même que le degré de l'extension
$k(X)/k(\mathbf{A}^n_k)$ parce que $h^{(p)}$ est le changement
de base de $h$ (un détail technique est omis). Par la multiplicativité des
degrés (Morphismes, lemme
\ref{morphisms-lemma-degree-composition}), il suffit de montrer que le
degré de $F_{\mathbf{A}^n_k/k}$ est $p^n$. Pour le voir, observons que
$(\mathbf{A}^n_k)^{(p)} = \mathbf{A}^n_k$ et que $F_{\mathbf{A}^n_k/k}$
est donné par l'application qui envoie les coordonnées sur leurs puissances $p$.
\end{proof}

\begin{remark}
\label{varieties-remark-n-fold-relative-frobenius}
Supposons que $p > 0$ soit un nombre premier. Supposons que $S$ soit un schéma
en caractéristique $p$. Supposons que $X$ soit un schéma sur $S$. Pour $n \geq 1$
$$
X^{(p^n)} = X^{(p^n/S)} = X \times_{S, F_S^n} S
$$
considéré comme un schéma sur $S$. Observons que $X \mapsto
X^{(p^n)}$ est un foncteur. En appliquant le lemme
\ref{varieties-lemma-frobenius-endomorphism-identity}, on voit que
$F_{X/S, n} = (F_X^n, \text{id}_S) : X \longrightarrow X^{(p^n)}$
est un morphisme sur $S$ s'insérant dans le diagramme commutatif
$$
\xymatrix{
X \ar[rr]_{F_{X/S, n}} \ar[rrd] \ar@/^1em/[rrrr]^{F_X^n}
& & X^{(p^n)} \ar[rr] \ar[d] & & X \ar[d] \\
& & S \ar[rr]^{F_S^n} & & S
}
$$
où le carré de droite est cartésien. Le
morphisme $F_{X/S, n}$ est parfois appelé le {\it
morphisme de Frobenius relatif itéré $n$ fois de
$X/S$ }. Cela a du sens car nous avons la formule
$$
F_{X/S, n} =
F_{X^{(p^{n - 1})}/S} \circ \ldots \circ F_{X^{(p)}/S} \circ F_{X/S}
$$
ce qui montre que $F_{X/S, n}$ est la composition
de $n$ morphismes de Frobenius relatifs. Puisque nous avons
$$
F_{X^{(p^m)}/S} = F_{X^{(p^{m - 1})}/S}^{(p)} = \ldots = F_{X/S}^{(p^m)}
$$
(détails omis) nous obtenons également que
$$
F_{X/S, n} =
F_{X/S}^{(p^{n - 1})} \circ \ldots \circ F_{X/S}^{(p)} \circ F_{X/S}
$$
\end{remark}




\section{Recollement d'anneaux de dimension un}
\label{varieties-section-glueing-dim-1}

\noindent
Cette section contient quelques préliminaires algébriques
pour démontrer qu'un ensemble fini de points de codimension
$1$ d'un schéma séparé est contenu dans un ouvert affine.

\begin{situation}
\label{varieties-situation-glue}
On se donne un diagramme commutatif d'anneaux
$$
\xymatrix{
A \ar[r] & K \\
R \ar[u] \ar[r] & B \ar[u]
}
$$
où $K$ est un corps et $A$, $B$ sont des sous-anneaux de $K$ avec
le corps des fractions $K$. Enfin, $R = A \times_K B = A \cap B$.
\end{situation}

\begin{lemma}
\label{varieties-lemma-glue-valuation-ring}
Dans la situation \ref{varieties-situation-glue}, supposons que $B$ soit un anneau de valuation.
Alors pour chaque élément inversible $u$ de $A$, soit $u \in R$ soit $u^{-1} \in R$.
\end{lemma}

\begin{proof}
En effet, si l'image $c$ de $u$ dans $K$ est
dans $B$, alors $u \in R$. Sinon, $c^{-1} \in
B$ (Algèbre, lemme
\ref{algebra-lemma-valuation-ring-x-or-x-inverse}) et $u^{-1} \in R$.
\end{proof}

\noindent
Le lemme suivant explique la signification de la
condition `` $A \otimes B \to K$ est surjective
'' qui revient assez souvent dans ce qui suit.

\begin{lemma}
\label{varieties-lemma-glue-separated}
Dans la situation \ref{varieties-situation-glue}, supposons que $A$ soit un anneau
noethérien de dimension $1$. Les assertions suivantes sont équivalentes
\begin{enumerate}
\item $A \otimes B \to K$ n'est pas surjective,
\item il existe un anneau de valuation discrète
$\mathcal{O} \subset K$ contenant à la fois $A$ et $B$.
\end{enumerate}
\end{lemma}

\begin{proof}
Il est clair que (2) implique (1). D'autre part, si $A \otimes B \to
K$ n'est pas surjective, alors l'image $C \subset K$ n'est pas un
corps, donc $C$ possède un idéal maximal non nul $\mathfrak m$.
Choisissons un anneau de valuation $\mathcal{O} \subset K$ dominant
$C_\mathfrak m$. D'après le lemme d'Algèbre
\ref{algebra-lemma-krull-akizuki} appliqué à $A \subset \mathcal{O}$, l'anneau
$\mathcal{O}$ est noethérien. Par conséquent, $\mathcal{O}$ est un anneau de valuation
discrète d'après le lemme d'Algèbre \ref{algebra-lemma-valuation-ring-Noetherian-discrete}.
\end{proof}

\begin{lemma}
\label{varieties-lemma-semi-local}
Dans la situation \ref{varieties-situation-glue}, supposons
\begin{enumerate}
\item $A$ est un anneau intègre semi-local noethérien de
dimension $1$, \item $B$ est un anneau de valuation discrète,
\end{enumerate}
Alors nous avons les deux possibilités suivantes
\begin{enumerate}
\item [(a)] Si $A^*$ n'est pas contenu dans $R$, alors
$\Spec(A) \to \Spec(R)$ et $\Spec(B) \to \Spec(R)$ sont
des immersions ouvertes recouvrant $\Spec(R)$ et $K = A
\otimes_R B$. \item [(b)] Si $A^*$ est contenu dans $R$,
alors $B$ domine l'un des anneaux locaux de $A$ en un
idéal maximal et $A \otimes B \to K$ n'est pas surjectif.
\end{enumerate}
\end{lemma}

\begin{proof}
L'hypothèse (a) implique qu'il existe une unité $u$ de $A$ dont
l'image dans $K$ appartienne à l'idéal maximal de $B$. Alors $u$
est un non-diviseur de zéro de $R$ et pour chaque $a \in A$ il
existe un $n$ tel que $u^n a \in R$. Il s'ensuit que $A = R_u$.

\medskip\noindent
Supposons que $\mathfrak m_A$ soit le radical de Jacobson de $A$.
Supposons que $x \in \mathfrak m_A$ soit un élément non nul. Puisque
$\dim(A) = 1$, on voit que $K = A_x$. Après avoir remplacé $x$ par $x^n
u^m$ pour certains $n \geq 1$ et $m \in \mathbf{Z}$, on peut supposer que
l'image de $x$ soit une unité de $B$. Pour chaque $b \in B$, le produit $x^nb$
appartient à l'image de $R$ pour un entier $n$ convenable. Ainsi $B = R_x$.

\medskip\noindent
Soit $z \in R$. Si $z \not \in \mathfrak m_A$ et si l'image de $z$ n'appartient pas
à $\mathfrak m_B$, alors $z$ est inversible. Ainsi, $x
+ u$ est inversible dans $R$. Par conséquent, $\Spec(R) = D(x) \cup
D(u)$. Nous avons vu ci-dessus que $D(u) = \Spec(A)$ et $D(x) = \Spec(B)$.

\medskip\noindent
Cas (b). Si $x \in \mathfrak m_A$, alors $1 + x$ est une
unité et donc $1 + x \in R$, c'est-à-dire $x \in R$. Ainsi,
on voit que $\mathfrak m_A \subset R \subset A$. En
fait, dans ce cas, $A$ est intégral sur $R$. En effet,
écrivons $A/\mathfrak m_A = \kappa_1 \times \ldots \times
\kappa_n$ comme un produit de corps. Disons que $x = (c_1,
\ldots, c_r, 0, \ldots, 0)$ est un élément avec $c_i \not = 0$. Alors
$$
x^2 - x(c_1, \ldots, c_r, 1, \ldots, 1)  = 0
$$
Puisque $R$ contient toutes les unités, on voit que
$A/\mathfrak m_A$ est intégral sur l'image de $R$ en lui,
et donc $A$ est intégral sur $R$. Il s'ensuit que $R \subset
A \subset B$ comme $B$ est intégralement clos. De plus, si
$x \in \mathfrak m_A$ est non nul, alors $K = A_x = \bigcup
x^{-n}A = \bigcup x^{-n}R$. Par conséquent $x^{-1} \not \in
B$, c'est-à-dire $x \in \mathfrak m_B$. Nous concluons
$\mathfrak m_A \subset \mathfrak m_B$. Ainsi, $A \cap
\mathfrak m_B$ est un idéal maximal de $A$, ce qui achève la démonstration.
\end{proof}

\begin{lemma}
\label{varieties-lemma-semi-local-dimension-one-conductor}
Supposons que $B$ soit un anneau intègre noethérien semi-local de
dimension $1$. Soit $B'$ la clôture
intégrale de $B$ dans son corps des fractions. Alors $B'$
est un anneau de Dedekind semi-local. Supposons que $x$
soit un élément non nul du radical de Jacobson de $B'$. Alors
pour chaque $y \in B'$ il existe un $n$ tel que $x^n y \in B$.
\end{lemma}

\begin{proof}
Supposons que $\mathfrak m_B$ soit le radical de Jacobson
de $B$. La structure de $B'$ résulte de l'Algèbre, lemme
\ref{algebra-lemma-integral-closure-Dedekind}. Étant
donné $x, y \in B'$ comme dans l'énoncé du lemme,
considérons le sous-anneau $B \subset A \subset B'$
engendré par $x$ et $y$. Alors $A$ est fini sur $B$
(Algèbre, lemme
\ref{algebra-lemma-characterize-finite-in-terms-of-integral}).
Puisque le corps des fractions de $B$ et $A$ est le même, nous
voyons que le module fini $A/B$ a son support dans l'ensemble des
points fermés de $B$. Ainsi $\mathfrak m_B^n A \subset B$ pour
un $n$ approprié. De plus, $\Spec(B') \to \Spec(A)$ est
surjectif (Algèbre, lemme
\ref{algebra-lemma-integral-overring-surjective}), donc $A$ est également
semi-local. Il en découle également que $x$ est dans le radical de Jacobson
$\mathfrak m_A$ de $A$. Notons que $\mathfrak m_A = \sqrt{\mathfrak m_B A}$.
Ainsi $x^m y \in \mathfrak m_B A$ pour un certain $m$. Puis $x^{nm} y \in B$.
\end{proof}

\begin{lemma}
\label{varieties-lemma-semi-local-both-side}
Dans la situation \ref{varieties-situation-glue}, supposons
\begin{enumerate}
\item $A$ est un anneau intègre semi-local noethérien de
dimension $1$, \item $B$ est un anneau intègre semi-local
noethérien de dimension $1$, \item $A \otimes B \to K$ est surjectif.
\end{enumerate}
Alors $\Spec(A) \to \Spec(R)$ et $\Spec(B) \to \Spec(R)$ sont des
immersions ouvertes couvrant $\Spec(R)$ et $K = A \otimes_R B$.
\end{lemma}

\begin{proof}
Cas particulier : $B$ est intégralement clos dans $K$.
Cela signifie que $B$ est un anneau de Dedekind (Algèbre,
lemme \ref{algebra-lemma-characterize-Dedekind}) d'où toutes
ses localisations en des idéaux maximaux sont des anneaux
de valuation discrète. Soient $\mathfrak m_1, \ldots,
\mathfrak m_r$ les idéaux maximaux de $B$. Posons
$$
R_1 = A \times_K B_{\mathfrak m_1}
$$
En observant que $A \otimes_{R_1} B_{\mathfrak m_1} \to K$ est surjectif,
nous concluons d'après le lemme \ref{varieties-lemma-semi-local} que $A$ et
$B_{\mathfrak m_1}$ définissent des sous-schémas ouverts couvrant $\Spec(R_1)$
et que $K = A \otimes_{R_1} B_{\mathfrak m_1}$. En particulier, $R_1$ est un
anneau noethérien semi-local de dimension $1$. Par induction, nous définissons
$$
R_{i + 1} = R_i \times_K B_{\mathfrak m_{i + 1}}
$$
pour $i = 1, \ldots, r - 1$. Observons que $R = R_r$ parce
que $B = B_{\mathfrak m_1} \cap \ldots \cap B_{\mathfrak
m_r}$ (voir Algèbre, lemme
\ref{algebra-lemma-normal-domain-intersection-localizations-height-1}).
Il découle de la procédure inductive que $R \to A$ définit
une immersion ouverte $\Spec(A) \to \Spec(R)$. D'autre part,
les idéaux maximaux $\mathfrak n_i$ de $R$ ne se trouvant pas
dans cet ouvert correspondent aux idéaux maximaux $\mathfrak
m_i$ de $B$ et en fait le homomorphisme d'anneaux $R \to B$
définit un isomorphisme $R_{\mathfrak n_i} \to B_{\mathfrak
m_i}$ (détails omis ; indication : à chaque étape nous avons ajouté
exactement un idéal maximal à $\Spec(R_i)$). Il s'ensuit que
$\Spec(B) \to \Spec(R)$ est une immersion ouverte, comme annoncé.

\medskip\noindent
Cas général. Soit $B' \subset K$ la clôture intégrale de $B$.
Voir le lemme \ref{varieties-lemma-semi-local-dimension-one-conductor}.
Alors le cas particulier s'applique à $R' = A \times_K B'$.
Choisissons $x \in R'$ qui n'appartient à aucun des idéaux maximaux
de $A$ et appartient à tous les idéaux maximaux de $B'$
(voir Algèbre, lemme \ref{algebra-lemma-chinese-remainder}).
D'après le lemme \ref{varieties-lemma-semi-local-dimension-one-conductor},
il existe un entier $n$ tel que $x^n \in R = A \times_K B$.
Remplaçons $x$ par $x^n$ afin que $x \in R$. Pour tout $y \in R'$, il existe
un entier $n$ tel que $x^n y \in R$. D'autre part,
il est clair que $R'_x = A$. Ainsi $R_x = A$.
En échangeant les rôles de $A$ et de $B$, on trouve aussi un $y \in R$
tel que $B = R_y$. Remarquons qu'en inversant à la fois $x$ et $y$, il
ne reste aucun idéal premier hormis $(0)$. Ainsi $K = R_{xy} = R_x \otimes_R R_y$.
Ceci achève la démonstration.
\end{proof}

\begin{lemma}
\label{varieties-lemma-glue-a-bunch-of-local-rings}
Soit $K$ un corps. Soient $A_1, \ldots, A_r \subset K$ des anneaux
semi-locaux noethériens de dimension $1$ et de corps des fractions $K$. Si
$A_i \otimes A_j \to K$ est surjectif pour tous $i \not = j$, alors
il existe un anneau intègre semi-local noethérien $A \subset K$
de dimension $1$, contenu dans $A_1, \ldots, A_r$, tel que
\begin{enumerate}
\item $A \to A_i$ induit une immersion ouverte $j_i : \Spec(A_i) \to \Spec(A)$,
\item $\Spec(A)$ est la réunion des ouverts $j_i(\Spec(A_i))$,
\item chaque point fermé de $\Spec(A)$ appartient à exactement un de ces
ouverts.
\end{enumerate}
\end{lemma}

\begin{proof}
On peut en effet prendre $A = A_1 \cap \ldots \cap A_r$. Remarquons d'abord que (3),
une fois (1) et (2) démontrés, résulte de l'hypothèse suivant laquelle
$A_i \otimes A_j \to K$ est surjectif : en effet, si
$\mathfrak m \in j_i(\Spec(A_i)) \cap j_j(\Spec(A_j))$, alors
$A_i \otimes A_j \to K$ se factorise par $A_\mathfrak m$.
Pour démontrer (1) et (2), raisonnons par récurrence sur $r$.
Si $r > 1$, l'hypothèse de récurrence donne (1) et (2) pour
$B = A_2 \cap \ldots \cap A_r$. Le
lemme \ref{varieties-lemma-semi-local-both-side} montre alors qu'ils valent pour
$A = A_1 \cap B$.
\end{proof}

\begin{lemma}
\label{varieties-lemma-create-globally-generated}
Soit $A$ un anneau intègre de corps des fractions $K$. Soient $B_1, \ldots, B_r \subset K$
des anneaux intègres semi-locaux noethériens de dimension $1$ dont les corps des
fractions sont égaux à $K$. Si $A \otimes B_i \to K$ est surjectif pour $i = 1, \ldots, r$,
alors il existe un $x \in A$ tel que $x^{-1}$ appartienne au radical de Jacobson de
$B_i$ pour $i = 1, \ldots, r$.
\end{lemma}

\begin{proof}
Notons $B_i'$ la clôture intégrale de $B_i$ dans $K$. Supposons qu'il existe un
$x \in A$ non nul tel que $x^{-1}$ appartienne au radical de Jacobson de $B'_i$ pour
$i = 1, \ldots, r$. Alors, d'après le
lemme \ref{varieties-lemma-semi-local-dimension-one-conductor},
quitte à remplacer $x$ par une puissance, on obtient $x^{-1} \in B_i$.
Comme $\Spec(B'_i) \to \Spec(B_i)$ est surjectif, il en résulte que
$x^{-1}$ appartient aussi au radical de Jacobson de $B_i$.
On peut donc supposer que chaque $B_i$ est un anneau de Dedekind semi-local.

\medskip\noindent
Si $B_i$ n'est pas local, supprimons $B_i$ de la liste et
remplaçons-le par la famille finie d'anneaux locaux $(B_i)_\mathfrak m$.
On peut ainsi supposer que $B_i$ est un anneau de valuation discrète pour
$i = 1, \ldots, r$.

\medskip\noindent
Soient $v_i : K \to \mathbf{Z}$, $i = 1, \ldots, r$,
les valuations discrètes correspondantes (voir
Algèbre, lemme \ref{algebra-lemma-characterize-Dedekind}).
Nous cherchons un $x \in A$ non nul tel que $v_i(x) < 0$ pour
$i = 1, \ldots, r$. Nous allons procéder par récurrence sur $r$.

\medskip\noindent
Si $r = 1$ et si l'assertion est fausse, alors $A \subset B$ et le morphisme
$A \otimes B \to K$ n'est pas surjectif, ce qui est contradictoire.

\medskip\noindent
Si $r > 1$, l'hypothèse de récurrence fournit un $x \in A$ non nul tel que
$v_i(x) < 0$ pour $i = 1, \ldots, r - 1$. Si $v_r(x) < 0$, le résultat est
acquis; on peut donc supposer que $v_r(x) \geq 0$. Le cas initial fournit un
$y \in A$ non nul tel que $v_r(y) < 0$. Quitte à remplacer $x$ par une
puissance, on peut supposer que $v_i(x) < v_i(y)$ pour $i = 1, \ldots, r - 1$.
Alors $x + y$ est l'élément recherché.
\end{proof}

\begin{lemma}
\label{varieties-lemma-power-equal}
Soit $A$ un anneau local noethérien de dimension $1$.
Soit $L = \prod A_\mathfrak p$, le produit portant sur
les idéaux premiers minimaux de $A$. Soient $a_1, a_2 \in \mathfrak m_A$
ayant même image dans $L$. Alors $a_1^n = a_2^n$ pour
un certain $n > 0$.
\end{lemma}

\begin{proof}
Écrivons $a_1 = a_2 + x$. Alors $x$ a une image nulle dans $L$. Par conséquent,
$x$ est un élément nilpotent de $A$, car $\bigcap \mathfrak p$
est le radical de $(0)$, et l'annulateur $I$ de $x$ contient
une puissance de l'idéal maximal, puisque $\mathfrak p \not \in V(I)$
pour tout idéal premier minimal. Écrivons $x^k = 0$ et $\mathfrak m^n \subset I$.
Alors
$$
a_1^{k + n} = a_2^{k + n} + {n + k \choose 1} a_2^{n + k - 1} x +
{n + k \choose 2} a_2^{n + k - 2} x^2 + \ldots +
{n + k \choose k - 1} a_2^{n + 1} x^{k - 1} = a_2^{n + k}
$$
car $a_2 \in \mathfrak m_A$.
\end{proof}

\begin{lemma}
\label{varieties-lemma-power-works}
Soit $A$ un anneau local noethérien de dimension $1$.
Soient $L = \prod A_\mathfrak p$ et $I = \bigcap \mathfrak p$,
le produit et l'intersection portant sur
les idéaux premiers minimaux de $A$. Soit $f \in L$ un élément
de la forme $f = i + a$, où $a \in \mathfrak m_A$ et
$i \in IL$. Alors une puissance de $f$ appartient à l'image de $A \to L$.
\end{lemma}

\begin{proof}
Comme $A$ est noethérien, on a $I^t = 0$ pour un certain $t > 0$.
Supposons que $f = a + i$ avec $i \in I^kL$.
Alors $f^n = a^n + na^{n - 1}i \bmod I^{k + 1}L$.
Il suffit donc de montrer que $na^{n - 1}i$ appartient à
l'image de $I^k \to I^kL$ pour un certain $n \gg 0$.
Pour cela, choisissons un $g \in A$ tel que $\mathfrak m_A = \sqrt{(g)}$
(Algèbre, lemme \ref{algebra-lemma-height-1}). Alors $L = A_g$, par exemple d'après
Algèbre, proposition \ref{algebra-proposition-dimension-zero-ring}.
D'autre part, il existe un $n$ tel que $a^n \in (g)$.
On peut donc chasser les dénominateurs des éléments de $L$
en multipliant par une puissance suffisamment grande de $a$.
\end{proof}

\begin{lemma}
\label{varieties-lemma-good-intersection}
Soit $A$ un anneau local noethérien de dimension $1$.
Soit $L = \prod A_\mathfrak p$, le produit portant sur
les idéaux premiers minimaux de $A$. Soit $K \to L$ un homomorphisme d'anneaux entier.
Alors il existe $a \in \mathfrak m_A$ et $x \in K$
ayant même image dans $L$ et tels que $\mathfrak m_A = \sqrt{(a)}$.
\end{lemma}

\begin{proof}
D'après le lemme \ref{varieties-lemma-power-works}, on peut remplacer $A$ par
$A/(\bigcap \mathfrak p)$ et supposer que $A$ est un anneau réduit
(certains détails sont omis).
On peut aussi remplacer $K$ par l'image de $K \to L$.
Alors $K$ est un anneau réduit. Le morphisme $\Spec(L) \to \Spec(K)$ est
surjectif et fermé (détails omis). Ainsi, $\Spec(K)$ est un espace fini
discret. Il en résulte que $K$ est un produit fini de corps.

\medskip\noindent
Soient $\mathfrak p_j$, $j = 1, \ldots, m$, les idéaux premiers minimaux de $A$.
Notons $L_j$ le corps des fractions de $A_j$, de sorte que
$L = \prod_{j = 1, \ldots, m} L_j$.
Soit $A_j$ la normalisation de $A/\mathfrak p_j$. Alors
$A_j$ est un anneau de Dedekind semi-local possédant au moins
un idéal maximal; voir
Algèbre, lemme \ref{algebra-lemma-integral-closure-Dedekind}.
Soit $n$ la somme des nombres d'idéaux maximaux
de $A_1, \ldots, A_m$. Pour chaque idéal maximal $\mathfrak m \subset A_j$ de ce type,
considérons l'application
$$
v_{\mathfrak m} : L \to L_j \to \mathbf{Z} \cup \{\infty\}
$$
où la seconde flèche est la valuation discrète associée
à l'anneau de valuation discrète $(A_j)_{\mathfrak m}$, prolongée
en envoyant $0$ sur $\infty$. On obtient ainsi $n$ applications
$v_1, \ldots, v_n : L \to \mathbf{Z} \cup \{\infty\}$.
Nous allons trouver un élément $x \in K$ tel que $v_i(x) < 0$
pour tout $i = 1, \ldots, n$.

\medskip\noindent
Affirmons d'abord que, pour chaque $i$, il existe un élément $x \in K$
tel que $v_i(x) < 0$. Supposons en effet que $v_i$ corresponde à
$\mathfrak m \subset A_j$. Si $v_i(x) \geq 0$ pour tout $x \in K$,
alors $K$ s'envoie dans $(A_j)_{\mathfrak m}$ à l'intérieur du corps des fractions
$L_j$ de $A_j$.
L'image de $K$ dans $L_j$ est un corps sur lequel $L_j$ est
algébrique, d'après Algèbre, lemme \ref{algebra-lemma-integral-under-field}.
Avec ce qui précède, cela contredit Algèbre, lemme
\ref{algebra-lemma-valuation-ring-cap-field-finite}.

\medskip\noindent
Supposons avoir trouvé un élément $x \in K$ tel que
$v_1(x) < 0, \ldots, v_r(x) < 0$ pour un certain $r < n$. Si $v_{r + 1}(x) < 0$,
alors $x$ convient pour $r + 1$. Sinon, choisissons un $y \in K$ tel que
$v_{r + 1}(y) < 0$, ce qui est possible par le paragraphe précédent.
Après avoir remplacé $x$ par $x^n$ pour un certain $n > 0$,
on peut supposer que $v_i(x) < v_i(y)$ pour $i = 1, \ldots, r$. Alors
$v_j(x + y) = v_j(x) < 0$ pour $j = 1, \ldots, r$, par les propriétés
des valuations, et de même $v_{r + 1}(x + y) = v_{r + 1}(y) < 0$.
En raisonnant par récurrence, on trouve
$x \in K$ tel que $v_i(x) < 0$ pour $i = 1, \ldots, n$.

\medskip\noindent
En particulier, l'élément $x \in K$ a une projection non nulle
dans chaque facteur de $K$ (rappelons que $K$ est un produit fini de
corps et que, si une composante de $x$ était nulle, l'une
des valeurs $v_i(x)$ serait égale à $\infty$). Ainsi, $x$ est
inversible, et $x^{-1} \in K$ est un élément tel que
$\infty > v_i(x^{-1}) > 0$ pour tout $i$. Il résulte du
lemme \ref{varieties-lemma-semi-local-dimension-one-conductor} que,
pour un certain $e < 0$, l'élément $x^e \in K$ s'envoie sur un élément de
$\mathfrak m_A/\mathfrak p_j \subset A/\mathfrak p_j$ pour tout
$j = 1, \ldots, m$. Remarquons que le conoyau de l'application
$\mathfrak m_A \to \prod \mathfrak m_A/\mathfrak p_j$ est annulé
par une puissance de $\mathfrak m_A$. Ainsi, quitte à remplacer
$e$ par un $e$ plus négatif, on trouve un élément $a \in \mathfrak m_A$
dont l'image dans $\mathfrak m_A/\mathfrak p_j$ est égale à
l'image de $x^e$. La paire $(a, x^e)$ satisfait les
conclusions du lemme.
\end{proof}

\begin{lemma}
\label{varieties-lemma-localization-semi-local}
Soit $A$ un anneau. Soient $\mathfrak p_1, \ldots, \mathfrak p_r$
un ensemble fini d'idéaux premiers de $A$. Posons
$S = A \setminus \bigcup \mathfrak p_i$. Alors $S$ est un système multiplicatif
et $S^{-1}A$ est un anneau semi-local dont les idéaux maximaux
correspondent aux éléments maximaux de l'ensemble $\{\mathfrak p_i\}$.
\end{lemma}

\begin{proof}
Si $a, b \in A$ et $a, b \in S$, alors $a, b \not \in \mathfrak p_i$,
donc $ab \not \in \mathfrak p_i$, et par suite $ab \in S$. De plus, $1 \in S$.
Ainsi, $S$ est une partie multiplicative de $A$. D'après la description de
$\Spec(S^{-1}A)$ donnée dans
Algèbre, lemme \ref{algebra-lemma-spec-localization},
et d'après
Algèbre, lemme \ref{algebra-lemma-silly},
les idéaux premiers de $S^{-1}A$ correspondent aux idéaux premiers de
$A$ contenus dans l'un des $\mathfrak p_i$.
Les idéaux maximaux de $S^{-1}A$ correspondent donc bijectivement aux
éléments maximaux (pour l'inclusion)\ parmi les éléments de l'ensemble
$\{\mathfrak p_1, \ldots, \mathfrak p_r\}$.
\end{proof}




\section{Schémas noethériens de dimension un}
\label{varieties-section-dimension-one}

\noindent
Le résultat principal de cette section affirme qu'un
schéma noethérien séparé de dimension $1$ possède un
faisceau inversible ample. Voir la
proposition \ref{varieties-proposition-dim-1-noetherian-separated-has-ample}.

\begin{lemma}
\label{varieties-lemma-affine}
Soit $X$ un schéma dont tous les anneaux locaux sont noethériens de dimension
$\leq 1$. Soit $U \subset X$ un ouvert rétrocompact. Notons
$j : U \to X$ le morphisme d'inclusion. Alors $R^pj_*\mathcal{F} = 0$, $p > 0$
pour tout $\mathcal{O}_U$-module quasi-cohérent $\mathcal{F}$.
\end{lemma}

\begin{proof}
Il suffit de vérifier l'annulation de $R^pj_*\mathcal{F}$ sur les germes.
La formation de $R^qj_*$ commute au changement de base plat; voir
Cohomologie des schémas, lemme
\ref{coherent-lemma-flat-base-change-cohomology}.
On peut donc supposer que $X$ est le spectre d'un anneau local noethérien
de dimension $\leq 1$. Dans ce cas, $X$ possède un point fermé
$x$ et un nombre fini d'autres points $x_1, \ldots, x_n$ qui se spécialisent
en $x$, mais non les uns dans les autres (voir
Algèbre, lemme \ref{algebra-lemma-Noetherian-irreducible-components}).
Si $x \in U$, alors $U = X$ et le résultat est clair. Sinon,
$U = \{x_1, \ldots, x_r\}$ pour un certain $r$, après avoir éventuellement renuméroté
les points. Alors $U$ est affine
(Schémas, lemme \ref{schemes-lemma-scheme-finite-discrete-affine}).
Le résultat découle donc de Cohomologie des schémas, lemme
\ref{coherent-lemma-relative-affine-vanishing}.
\end{proof}

\begin{lemma}
\label{varieties-lemma-open-in-affine-curve-affine}
Soit $X$ un schéma affine dont tous les anneaux locaux sont noethériens
de dimension $\leq 1$. Alors tout ouvert quasi-compact $U \subset X$ est affine.
\end{lemma}

\begin{proof}
Notons $j : U \to X$ le morphisme d'inclusion. Soit $\mathcal{F}$
un $\mathcal{O}_U$-module quasi-cohérent.
D'après le lemme \ref{varieties-lemma-affine}, les images directes supérieures
$R^pj_*\mathcal{F}$ sont nulles. Le $\mathcal{O}_X$-module $j_*\mathcal{F}$
est quasi-cohérent
(Schémas, lemme \ref{schemes-lemma-push-forward-quasi-coherent}).
Ses groupes de cohomologie supérieurs sont donc nuls d'après
Cohomologie des schémas, lemme
\ref{coherent-lemma-quasi-coherent-affine-cohomology-zero}.
La suite spectrale de Leray
Cohomologie, lemme \ref{cohomology-lemma-apply-Leray},
donne $H^p(U, \mathcal{F}) = 0$ pour tout $p > 0$.
Ainsi, $U$ est affine, par exemple d'après
Cohomologie des schémas, lemme
\ref{coherent-lemma-quasi-compact-h1-zero-covering}.
\end{proof}

\begin{lemma}
\label{varieties-lemma-complement-codim-1-closed-points}
Soit $X$ un schéma. Soit $U \subset X$ un ouvert. Supposons que
\begin{enumerate}
\item $U$ soit un ouvert rétrocompact de $X$,
\item $X \setminus U$ soit discret, et
\item pour $x \in X \setminus U$, l'anneau local
$\mathcal{O}_{X, x}$ soit noethérien de dimension $\leq 1$.
\end{enumerate}
Alors (1) il existe un $\mathcal{O}_X$-module inversible $\mathcal{L}$
et une section $s$ tels que $U = X_s$, et (2) le morphisme
$\Pic(X) \to \Pic(U)$ est surjectif.
\end{lemma}

\begin{proof}
Écrivons $X \setminus U = \{x_i; i \in I\}$.
Choisissons des ouverts affines $U_i \subset X$ tels que $x_i \in U_i$ et
$x_j \not \in U_i$ pour $j \not = i$. C'est possible par la condition (2).
Écrivons $U_i = \Spec(A_i)$. Soit $\mathfrak m_i \subset A_i$ l'idéal maximal
correspondant à $x_i$. Par notre hypothèse sur les anneaux locaux,
il n'y a qu'un nombre fini d'idéaux premiers
$\mathfrak q \subset \mathfrak m_i$,
$\mathfrak q \not = \mathfrak m_i$ (voir
Algèbre, lemme \ref{algebra-lemma-Noetherian-irreducible-components}).
Le lemme d'évitement des idéaux premiers (Algèbre, lemme
\ref{algebra-lemma-silly}) fournit donc un $f_i \in \mathfrak m_i$
qui n'appartient à aucun de ces idéaux premiers. Alors
$V(f_i) = \{\mathfrak m_i\} \amalg Z_i$ pour un certain sous-ensemble fermé
$Z_i \subset U_i$, car $Z_i$ est un ouvert rétrocompact de
$V(f_i)$ stable par spécialisation; voir
Algèbre, lemme \ref{algebra-lemma-constructible-stable-specialization-closed}.
Quitte à rétrécir $U_i$, on peut supposer que $V(f_i) = \{x_i\}$. Alors
$$
\mathcal{U} : X = U \cup \bigcup U_i
$$
est un recouvrement ouvert de $X$. Considérons le $2$-cocycle à valeurs
dans $\mathcal{O}_X^*$ donné par $f_i$ sur $U \cap U_i$ et par
$f_i/f_j$ sur $U_i \cap U_j$. Il définit un faisceau inversible
$\mathcal{L}$ tel que la section $s$, définie par $1$ sur $U$
et par $f_i$ sur $U_i$, possède les propriétés énoncées dans le lemme.

\medskip\noindent
Soit $\mathcal{N}$ un $\mathcal{O}_U$-module inversible.
Soit $N_i$ le $(A_i)_{f_i}$-module inversible tel que
$\mathcal{N}|_{U \cap U_i}$ soit égal à $\tilde N_i$.
Remarquons que $(A_{\mathfrak m_i})_{f_i}$ est un anneau artinien
(car il s'agit d'un anneau noethérien de dimension zéro; voir
Algèbre, lemme \ref{algebra-lemma-Noetherian-dimension-0}).
C'est donc un produit d'anneaux locaux
(Algèbre, lemme \ref{algebra-lemma-artinian-finite-length}) et
son groupe de Picard est par conséquent trivial. En rétrécissant $U_i$
(c'est-à-dire en remplaçant $A_i$ par $(A_i)_g$ pour un $g \in A_i$,
$g \not \in \mathfrak m_i$),
on peut donc supposer que $N_i = (A_i)_{f_i}$, c'est-à-dire que
$\mathcal{N}|_{U \cap U_i}$ est trivial. Dans ce cas, on voit clairement
comment prolonger $\mathcal{N}$ en un faisceau inversible sur $X$
(en le prolongeant par un module inversible trivial sur chaque $U_i$).
\end{proof}

\begin{lemma}
\label{varieties-lemma-find-globally-generated}
Soit $X$ un schéma intègre séparé. Soit $U \subset X$ un ouvert affine non vide
tel que $X \setminus U$ soit un ensemble fini de points
$x_1, \ldots, x_r$ et que $\mathcal{O}_{X, x_i}$ soit noethérien de dimension $1$.
Alors il existe un $\mathcal{O}_X$-module inversible engendré par ses sections globales
$\mathcal{L}$ et une section $s$ tels que $U = X_s$.
\end{lemma}

\begin{proof}
Écrivons $U = \Spec(A)$ et notons $K$ le corps des fonctions rationnelles de $X$.
Posons $B_i = \mathcal{O}_{X, x_i}$ et $\mathfrak m_i = \mathfrak m_{x_i}$.
Comme $x_i \not \in U$, l'ouvert
$U \times_X \Spec(B_i)$ de $\Spec(B_i)$ ne possède qu'un point, c'est-à-dire
$U \times_X \Spec(B_i) = \Spec(K)$.
Puisque $X$ est séparé, $\Spec(K)$ est un sous-schéma fermé
de $U \times \Spec(B_i)$, c'est-à-dire que le morphisme $A \otimes B_i \to K$ est surjectif.
D'après le lemme \ref{varieties-lemma-create-globally-generated}, il existe un
$f \in A$ non nul tel que $f^{-1} \in \mathfrak m_i$ pour $i = 1, \ldots, r$.
Choisissons des ouverts $x_i \in U_i \subset X$ tels que $f^{-1} \in \mathcal{O}(U_i)$.
Alors
$$
\mathcal{U} : X = U \cup \bigcup U_i
$$
est un recouvrement ouvert de $X$. Considérons le $2$-cocycle à valeurs
dans $\mathcal{O}_X^*$ donné par $f$ sur $U \cap U_i$ et par
$1$ sur $U_i \cap U_j$. Il définit un faisceau inversible $\mathcal{L}$
muni de deux sections :
\begin{enumerate}
\item une section $s$, définie par $1$ sur $U$
et par $f^{-1}$ sur $U_i$, qui possède les propriétés énoncées dans le lemme, et
\item une section $t$, définie par $f$ sur $U$ et
par $1$ sur $U_i$.
\end{enumerate}
Remarquons que $X_t \supset U_1 \cup \ldots \cup U_r$.
Ainsi, $s, t$ engendrent $\mathcal{L}$, ce qui démontre le lemme.
\end{proof}

\begin{lemma}
\label{varieties-lemma-enough-globally-generated-ample}
Soit $X$ un schéma quasi-compact. Si, pour tout $x \in X$,
il existe une paire $(\mathcal{L}, s)$ composée d'un faisceau inversible
$\mathcal{L}$ engendré par ses sections globales et d'une section globale $s$ telle que
$x \in X_s$ et que $X_s$ soit affine, alors $X$ possède un faisceau inversible
ample.
\end{lemma}

\begin{proof}
Comme $X$ est quasi-compact, on peut trouver une famille finie
$(\mathcal{L}_i, s_i)$, $i = 1, \ldots, n$,
de paires telles que $\mathcal{L}_i$ soit engendré par ses sections globales, que $X_{s_i}$
soit affine et que $X = \bigcup X_{s_i}$.
Toujours parce que $X$ est quasi-compact, on peut trouver, pour chaque $i$, une famille
finie de sections $t_{i, j}$ de $\mathcal{L}_i$, $j = 1, \ldots, m_i$,
telle que $X = \bigcup X_{t_{i, j}}$. Posons $t_{i, 0} = s_i$.
Considérons le faisceau inversible
$$
\mathcal{L} = \mathcal{L}_1 
\otimes_{\mathcal{O}_X} \ldots
\otimes_{\mathcal{O}_X} \mathcal{L}_n
$$
et les sections globales
$$
\tau_J = t_{1, j_1} \otimes \ldots \otimes t_{n, j_n}
$$
D'après Propriétés, lemme \ref{properties-lemma-affine-cap-s-open},
l'ouvert $X_{\tau_J}$ est affine dès que $j_i = 0$ pour un certain $i$.
Il est immédiat que ces ouverts recouvrent $X$.
Ainsi, $\mathcal{L}$ est ample par définition.
\end{proof}

\begin{lemma}
\label{varieties-lemma-dim-1-noetherian-integral-separated-has-ample}
Soit $X$ un schéma noethérien intègre séparé de dimension $1$.
Alors $X$ possède un faisceau inversible ample.
\end{lemma}

\begin{proof}
Choisissons un recouvrement ouvert affine $X = U_1 \cup \ldots \cup U_n$.
Comme $X$ est noethérien, chacun des ensembles $X \setminus U_i$ est fini.
Le lemme \ref{varieties-lemma-find-globally-generated}
fournit donc une paire $(\mathcal{L}_i, s_i)$
composée d'un faisceau inversible $\mathcal{L}_i$ engendré par ses sections globales
et d'une section globale $s_i$ telle que $U_i = X_{s_i}$.
On conclut que $X$ possède un faisceau inversible ample grâce au
lemme \ref{varieties-lemma-enough-globally-generated-ample}.
\end{proof}

\begin{lemma}
\label{varieties-lemma-surjective-pic-birational-finite}
Soit $f : X \to Y$ un morphisme fini de schémas. Supposons qu'il
existe un ouvert $V \subset Y$ tel que $f^{-1}(V) \to V$ soit un
isomorphisme et que $Y \setminus V$ soit un espace discret. Alors tout
$\mathcal{O}_X$-module inversible est l'image réciproque d'un
$\mathcal{O}_Y$-module inversible.
\end{lemma}

\begin{proof}
Nous utiliserons l'égalité $\Pic(X) = H^1(X, \mathcal{O}_X^*)$; voir
Cohomologie, lemme \ref{cohomology-lemma-h1-invertible}.
Considérons la suite spectrale de Leray associée au faisceau abélien $\mathcal{O}_X^*$
et à $f$; voir Cohomologie, lemme \ref{cohomology-lemma-Leray}.
Considérons le morphisme induit
$$
H^1(X, \mathcal{O}_X^*) \longrightarrow H^0(Y, R^1f_*\mathcal{O}_X^*)
$$
Diviseurs, lemme \ref{divisors-lemma-finite-trivialize-invertible-upstairs},
affirme précisément que ce morphisme est nul. La suite spectrale de Leray donne donc
$H^1(X, \mathcal{O}_X^*) = H^1(Y, f_*\mathcal{O}_X^*)$.
Considérons ensuite le morphisme
$$
f^\sharp : \mathcal{O}_Y^* \longrightarrow f_*\mathcal{O}_X^*
$$
Par hypothèse, le noyau et le conoyau de ce morphisme sont à support
dans le sous-ensemble fermé $T = Y \setminus V$ de $Y$.
Comme $T$ est un espace topologique discret par hypothèse,
les groupes de cohomologie supérieurs de tout faisceau abélien sur $Y$ à support
dans $T$ sont nuls (cela résulte de
Cohomologie, lemme \ref{cohomology-lemma-cohomology-and-closed-immersions},
Modules, lemme \ref{modules-lemma-i-star-exact}, et
du fait que $H^i(T, \mathcal{F}) = 0$ pour tout $i > 0$
et tout faisceau abélien $\mathcal{F}$ sur $T$).
En décomposant le morphisme affiché en suites exactes courtes
$$
0 \to \Ker(f^\sharp) \to \mathcal{O}_Y^* \to \Im(f^\sharp) \to 0,\quad
0 \to \Im(f^\sharp) \to f_*\mathcal{O}_X^* \to \Coker(f^\sharp) \to 0
$$
on conclut d'abord que $H^1(Y, \mathcal{O}_Y^*) \to H^1(Y, \Im(f^\sharp))$
est surjectif, puis que
$H^1(Y, \Im(f^\sharp)) \to H^1(Y, f_*\mathcal{O}_X^*)$ est surjectif.
En combinant tout ce qui précède, on voit que $H^1(Y, \mathcal{O}_Y^*) \to
H^1(X, \mathcal{O}_X^*)$ est surjectif, comme voulu.
\end{proof}

\begin{lemma}
\label{varieties-lemma-glue-invertible-sheaves}
Soit $X$ un schéma. Soient $Z_1, \ldots, Z_n \subset X$ des sous-schémas
fermés. Soit $\mathcal{L}_i$ un faisceau inversible sur $Z_i$.
Supposons que
\begin{enumerate}
\item $X$ soit réduit,
\item $X = \bigcup Z_i$ ensemblistement, et
\item $Z_i \cap Z_j$ soit un espace topologique discret pour $i \not = j$.
\end{enumerate}
Alors il existe un faisceau inversible $\mathcal{L}$ sur $X$ dont la restriction
à $Z_i$ est $\mathcal{L}_i$. De plus, si l'on se donne des sections
$s_i \in \Gamma(Z_i, \mathcal{L}_i)$ qui ne s'annulent pas aux
points de $Z_i \cap Z_j$, alors on peut choisir $\mathcal{L}$ de telle sorte
qu'il existe $s \in \Gamma(X, \mathcal{L})$ vérifiant
$s|_{Z_i} = s_i$ pour tout $i$.
\end{lemma}

\begin{proof}
L'existence de $\mathcal{L}$ peut se déduire du
lemme \ref{varieties-lemma-surjective-pic-birational-finite},
mais nous donnerons aussi une preuve directe et nous utiliserons
celle-ci pour établir l'assertion concernant les sections.
Posons $T = \bigcup_{i \not = j} Z_i \cap Z_j$. Comme $X$ est réduit, on a
$$
X \setminus T = \bigcup (Z_i \setminus T)
$$
en tant que schémas. L'hypothèse (3) implique que $T$ est une partie discrète de $X$.
Ainsi, pour chaque $t \in T$, il existe un ouvert $U_t \subset X$
tel que $t \in U_t$, mais $t' \not \in U_t$ si $t' \in T$, $t' \not = t$.
En restreignant $U_t$ au besoin, on peut supposer qu'il existe des isomorphismes
$\varphi_{t, i} : \mathcal{L}_i|_{U_t \cap Z_i} \to
\mathcal{O}_{U_t \cap Z_i}$. De plus, pour chaque $i$, choisissons un recouvrement ouvert
$$
Z_i \setminus T = \bigcup\nolimits_j U_{ij}
$$
tel qu'il existe des isomorphismes
$\varphi_{i, j} : \mathcal{L}_i|_{U_{ij}} \cong \mathcal{O}_{U_{ij}}$.
Observons que
$$
\mathcal{U} : X = \bigcup U_t \cup \bigcup U_{ij}
$$
est un recouvrement ouvert de $X$. Nous affirmons que les isomorphismes
$\varphi_{t, i}$ et $\varphi_{i, j}$ permettent de définir un $2$-cocycle à valeurs
dans $\mathcal{O}_X^*$ relativement à ce recouvrement, qui définit $\mathcal{L}$ comme
dans l'énoncé du lemme.

\medskip\noindent
Plus précisément, si $i \not = i'$, alors $U_{i, j} \cap U_{i', j'} = \emptyset$
et il n'y a rien à faire. Sur $U_{i, j} \cap U_{i, j'}$, on a
$\mathcal{O}_X(U_{i, j} \cap U_{i, j'}) =
\mathcal{O}_{Z_i}(U_{i, j} \cap U_{i, j'})$ d'après la première remarque de la preuve.
Ainsi, la fonction de transition de $\mathcal{L}_i$ (plus précisément
$\varphi_{i, j} \circ \varphi_{i, j'}^{-1}$) fournit la valeur de notre
cocycle sur cette intersection.
Sur $U_t \cap U_{i, j}$, on peut procéder de même.
Enfin, si $t \not = t'$, on a
$$
U_t \cap U_{t'} = \coprod (U_t \cap U_{t'}) \cap Z_i
$$
et, de plus, l'intersection $U_t \cap U_{t'} \cap Z_i$ est contenue
dans $Z_i \setminus T$. Par le même raisonnement que précédemment, on voit donc que
$$
\mathcal{O}_X(U_t \cap U_{t'}) =
\prod \mathcal{O}_{Z_i}(U_t \cap U_{t'} \cap Z_i)
$$
et l'on peut utiliser les fonctions de transition de $\mathcal{L}_i$ (plus précisément
$\varphi_{t, i} \circ \varphi_{t', i}^{-1}$) pour définir la valeur de
notre cocycle sur $U_t \cap U_{t'}$. Ceci achève la preuve de l'existence
de $\mathcal{L}$.

\medskip\noindent
Étant données des sections $s_i$ comme dans la dernière assertion du lemme, choisissons,
dans l'argument ci-dessus, $U_t$ de sorte que $s_i|_{U_t \cap Z_i}$ ne s'annule pas et
choisissons $\varphi_{t, i}$ de sorte que $\varphi_{t, i}(s_i|_{U_t \cap Z_i}) = 1$.
Alors l'emploi de $1$ sur $U_t$ et de $\varphi_{i, j}(s_i|_{U_{i, j}})$ sur
$U_{i, j}$ définit une section de $\mathcal{L}$ dont la restriction
est égale à $s_i$ sur $Z_i$.
\end{proof}

\begin{remark}
\label{varieties-remark-conductor}
Soit $A$ un anneau réduit. Soient $I, J$ des idéaux de $A$
tels que $V(I) \cup V(J) = \Spec(A)$. Posons $B = A/J$.
Alors $I \to IB$ est un isomorphisme de $A$-modules. En effet, on
a $IB = I + J/J = I/(I \cap J)$, et $I \cap J$ est nul car
$A$ est réduit et $\Spec(A) = V(I) \cup V(J) = V(I \cap J)$.
Ainsi, pour tout $A$-module projectif $P$, on a aussi $IP = I(P/JP)$.
\end{remark}

\begin{lemma}
\label{varieties-lemma-dim-1-noetherian-reduced-separated-has-ample}
Soit $X$ un schéma noethérien, réduit, séparé et de dimension $1$.
Alors $X$ possède un faisceau inversible ample.
\end{lemma}

\begin{proof}
Soient $Z_i$, $i = 1, \ldots, n$, les composantes irréductibles de $X$.
Nous les considérons comme des sous-schémas fermés réduits de $X$.
D'après le lemme \ref{varieties-lemma-dim-1-noetherian-integral-separated-has-ample},
il existe des faisceaux inversibles amples $\mathcal{L}_i$ sur $Z_i$.
Posons $T = \bigcup_{i \not = j} Z_i \cap Z_j$. Comme $X$ est noethérien
de dimension $1$, l'ensemble $T$ est fini et constitué de points fermés
de $X$. Pour chaque $i$, quitte à remplacer
$\mathcal{L}_i$ par une puissance, on peut choisir $s_i \in \Gamma(Z_i, \mathcal{L}_i)$
tel que $(Z_i)_{s_i}$ soit affine et contienne $T \cap Z_i$; voir
Propriétés, lemme \ref{properties-lemma-ample-finite-set-in-principal-affine}.

\medskip\noindent
D'après le lemme \ref{varieties-lemma-glue-invertible-sheaves}, il existe un faisceau inversible
$\mathcal{L}$ sur $X$ et $s \in \Gamma(X, \mathcal{L})$
tels que $(\mathcal{L}, s)|_{Z_i} = (\mathcal{L}_i, s_i)$.
Observons que $X_s$ contient $T$ et coïncide ensemblistement
avec la réunion des sous-schémas fermés affines $(Z_i)_{s_i}$. Il est donc affine d'après
Limites, lemme \ref{limits-lemma-affines-glued-in-closed-affine}.
Pour achever la preuve, il suffit de trouver, pour tout $x \in X$, $x \not \in T$,
un entier $m > 0$ et une section $t \in \Gamma(X, \mathcal{L}^{\otimes m})$
tels que $X_t$ soit affine et que $x \in X_t$. Puisque $x \not \in T$,
on voit que $x \in Z_i$ pour un unique $i$, disons $i = 1$.
Soit $Z \subset X$ le sous-schéma fermé réduit dont l'espace
topologique sous-jacent est $Z_2 \cup \ldots \cup Z_n$.
Soit $\mathcal{I} \subset \mathcal{O}_X$ le faisceau
d'idéaux de $Z$. Notons $\mathcal{I}_1 \subset \mathcal{O}_{Z_1}$
l'image inverse de ce faisceau d'idéaux par le morphisme
d'inclusion $Z_1 \to X$. Observons que
$$
\Gamma(X, \mathcal{I}\mathcal{L}^{\otimes m}) =
\Gamma(Z_1, \mathcal{I}_1 \mathcal{L}_1^{\otimes m})
$$
voir la remarque \ref{varieties-remark-conductor}. Il suffit donc de trouver $m > 0$
et $t \in \Gamma(Z_1, \mathcal{I}_1 \mathcal{L}_1^{\otimes m})$
tels que $x \in (Z_1)_t$, cet ouvert étant affine. Puisque $\mathcal{L}_1$ est ample
et que $x$ n'appartient pas à $Z_1 \cap T = V(\mathcal{I}_1)$,
on peut trouver une section
$t_1 \in \Gamma(Z_1, \mathcal{I}_1 \mathcal{L}_1^{\otimes m_1})$
telle que $x \in (Z_1)_{t_1}$; voir Propriétés, proposition
\ref{properties-proposition-characterize-ample}.
Puisque $\mathcal{L}_1$ est ample, on peut trouver une section
$t_2 \in \Gamma(Z_1, \mathcal{L}_1^{\otimes m_2})$
telle que $x \in (Z_1)_{t_2}$ et que $(Z_1)_{t_2}$ soit affine; voir
Propriétés, définition \ref{properties-definition-ample}.
Posons $m = m_1 + m_2$ et $t = t_1 t_2$. Alors
$t \in \Gamma(Z_1, \mathcal{I}_1 \mathcal{L}_1^{\otimes m})$
et $x \in (Z_1)_t$ par construction, tandis que
$(Z_1)_t$ est affine d'après Propriétés, lemme
\ref{properties-lemma-affine-cap-s-open}.
\end{proof}

\begin{lemma}
\label{varieties-lemma-lift-line-bundle-from-reduction-dimension-1}
Soit $i : Z \to X$ une immersion fermée de schémas.
Si l'espace topologique sous-jacent à $X$ est noethérien et si
$\dim(X) \leq 1$, alors $\Pic(X) \to \Pic(Z)$ est surjective.
\end{lemma}

\begin{proof}
Considérons la suite exacte courte
$$
0 \to (1 + \mathcal{I}) \cap \mathcal{O}_X^* \to
\mathcal{O}^*_X \to i_*\mathcal{O}^*_Z \to 0
$$
de faisceaux de groupes abéliens sur $X$, où $\mathcal{I}$
est le faisceau quasi-cohérent d'idéaux correspondant à $Z$.
Puisque $\dim(X) \leq 1$, on a $H^2(X, \mathcal{F}) = 0$
pour tout faisceau abélien $\mathcal{F}$; voir
Cohomologie, proposition \ref{cohomology-proposition-vanishing-Noetherian}.
Par conséquent, l'application $H^1(X, \mathcal{O}^*_X) \to H^1(X, i_*\mathcal{O}_Z^*)$
est surjective. D'après
Cohomologie, lemme \ref{cohomology-lemma-cohomology-and-closed-immersions},
on a $H^1(X, i_*\mathcal{O}_Z^*) = H^1(Z, \mathcal{O}_Z^*)$.
Cela démontre le lemme d'après
Cohomologie, lemme \ref{cohomology-lemma-h1-invertible}.
\end{proof}

\begin{proposition}
\label{varieties-proposition-dim-1-noetherian-separated-has-ample}
Soit $X$ un schéma noethérien séparé de dimension $1$.
Alors $X$ possède un faisceau inversible ample.
\end{proposition}

\begin{proof}
Soit $Z \subset X$ la réduction de $X$. D'après le
lemme \ref{varieties-lemma-dim-1-noetherian-reduced-separated-has-ample},
le schéma $Z$ possède un faisceau inversible ample.
Ainsi, d'après le lemme \ref{varieties-lemma-lift-line-bundle-from-reduction-dimension-1},
il existe un $\mathcal{O}_X$-module inversible $\mathcal{L}$
sur $X$ dont la restriction à $Z$ est ample.
Alors $\mathcal{L}$ est ample d'après
Cohomologie des schémas, lemme \ref{coherent-lemma-ample-on-reduction}.
\end{proof}

\begin{remark}
\label{varieties-remark-useless-generalization}
En fait, si $X$ est un schéma dont la réduction est un schéma noethérien
séparé de dimension $1$, alors $X$ possède un faisceau inversible
ample. La démonstration est la même que celle de la
proposition \ref{varieties-proposition-dim-1-noetherian-separated-has-ample},
à ceci près que l'on utilise
Limites, lemme \ref{limits-lemma-ample-on-reduction},
au lieu de
Cohomologie des schémas, lemme \ref{coherent-lemma-ample-on-reduction}.
\end{remark}

\noindent
Le lemme suivant vaut en fait pour les morphismes séparés quasi-finis,
comme le lecteur peut le constater en utilisant le théorème principal de Zariski
(Compléments sur les morphismes, lemme
\ref{more-morphisms-lemma-quasi-finite-separated-pass-through-finite})
et le lemme \ref{varieties-lemma-complement-codim-1-closed-points}.

\begin{lemma}
\label{varieties-lemma-surjection-on-pic-quasi-finite}
Soit $f : X \to Y$ un morphisme de schémas. Supposons que
$Y$ soit noethérien de dimension $\leq 1$, que $f$ soit fini et
qu'il existe un ouvert dense $V \subset Y$ tel que
$f^{-1}(V) \to V$ soit une immersion fermée. Alors tout
$\mathcal{O}_X$-module inversible est l'image réciproque d'un $\mathcal{O}_Y$-module inversible.
\end{lemma}

\begin{proof}
Factorisons $f$ en $X \to Z \to Y$, où $Z$ est l'image schématique
de $f$. Alors $X \to Z$ est un isomorphisme au-dessus de $V \cap Z$
et le lemme \ref{varieties-lemma-surjective-pic-birational-finite} s'applique.
D'autre part, le
lemme \ref{varieties-lemma-lift-line-bundle-from-reduction-dimension-1}
s'applique à $Z \to Y$. Nous omettons certains détails.
\end{proof}








\section{L'invariant delta}
\label{varieties-section-delta-invariant}

\noindent
Dans cette section, nous définissons l'invariant $\delta$ d'un point singulier
d'un schéma de Nagata réduit de dimension $1$.

\begin{lemma}
\label{varieties-lemma-pre-pre-delta-invariant}
Soit $(A, \mathfrak m)$ un anneau local noethérien de dimension $1$.
Soit $f \in \mathfrak m$. Les conditions suivantes sont équivalentes :
\begin{enumerate}
\item $\mathfrak m = \sqrt{(f)}$,
\item $f$ n'appartient à aucun idéal premier minimal de $A$, et
\item $A_f = \prod_{\mathfrak p\text{ minimal}} A_\mathfrak p$ comme $A$-algèbres.
\end{enumerate}
Un tel $f \in \mathfrak m$ existe. Si $\text{depth}(A) = 1$ (par exemple si
$A$ est réduit), alors (1) -- (3) sont aussi équivalentes à
\begin{enumerate}
\item[(4)] $f$ est un non-diviseur de zéro,
\item[(5)] $A_f$ est l'anneau total des fractions de $A$.
\end{enumerate}
Si $A$ est réduit, alors (1) -- (5) sont aussi équivalentes à
\begin{enumerate}
\item[(6)] $A_f$ est le produit des corps résiduels aux idéaux
premiers minimaux de $A$.
\end{enumerate}
\end{lemma}

\begin{proof}
Le spectre de $A$ ne possède qu'un nombre fini d'idéaux premiers
$\mathfrak p_1, \ldots, \mathfrak p_n$ autres que $\mathfrak m$,
et ceux-ci sont tous minimaux; voir
Algèbre, lemme \ref{algebra-lemma-Noetherian-irreducible-components}.
L'équivalence de (1) et (2) résulte alors de
Algèbre, lemme \ref{algebra-lemma-Zariski-topology}.
Il est clair que (3) implique (2). Réciproquement, si (2) est vraie,
alors le spectre de $A_f$ est la partie
$\{\mathfrak p_1, \ldots, \mathfrak p_n\}$ de $\Spec(A)$
munie de la topologie induite; voir
Algèbre, lemme \ref{algebra-lemma-spec-localization}.
C'est un espace topologique discret fini.
Par conséquent, $A_f = \prod_{\mathfrak p\text{ minimal}} A_\mathfrak p$ d'après
Algèbre, proposition \ref{algebra-proposition-dimension-zero-ring}.
L'existence d'un tel $f$ est affirmée dans
Algèbre, lemme \ref{algebra-lemma-height-1}.

\medskip\noindent
Supposons que $A$ soit de profondeur $1$. (C'est la valeur maximale d'après
Algèbre, lemme \ref{algebra-lemma-bound-depth}, et elle est atteinte si $A$ est réduit d'après
Algèbre, lemme \ref{algebra-lemma-criterion-reduced}.)
Alors $\mathfrak m$ n'est pas un idéal premier associé de $A$.
Tout idéal premier minimal de $A$ est un idéal premier associé
(Algèbre, proposition
\ref{algebra-proposition-minimal-primes-associated-primes}).
Par conséquent, les non-diviseurs de zéro de $A$ sont exactement les éléments
qui n'appartiennent à aucun des idéaux premiers minimaux, d'après
Algèbre, lemme \ref{algebra-lemma-ass-zero-divisors}.
Ainsi, (4) est équivalente à (2).
L'assertion (5) est équivalente à (3) d'après
Algèbre, lemme \ref{algebra-lemma-total-ring-fractions-no-embedded-points}.

\medskip\noindent
Alors $A_\mathfrak p$ est un corps pour tout idéal premier
$\mathfrak p \subset A$ minimal; voir
Algèbre, lemme \ref{algebra-lemma-minimal-prime-reduced-ring}.
Ainsi, (3) est équivalente à (6).
\end{proof}

\begin{lemma}
\label{varieties-lemma-pre-delta-invariant}
Soit $(A, \mathfrak m)$ un anneau local de Nagata réduit de dimension $1$.
Soit $A'$ la clôture intégrale de $A$ dans l'anneau total des fractions
de $A$. Alors $A'$ est un anneau de Nagata normal, $A \to A'$ est fini, et
$A'/A$ est un $A$-module de longueur finie.
\end{lemma}

\begin{proof}
L'anneau total des fractions est essentiellement de type fini sur $A$;
ainsi, $A \to A'$ est fini puisque $A$ est de Nagata; voir Algèbre, lemme
\ref{algebra-lemma-nagata-in-reduced-finite-type-finite-integral-closure}.
L'anneau $A'$ est normal, par exemple d'après
Algèbre, lemme \ref{algebra-lemma-characterize-reduced-ring-normal} et
\ref{algebra-lemma-Noetherian-irreducible-components}.
L'anneau $A'$ est de Nagata, par exemple d'après
Algèbre, lemme \ref{algebra-lemma-quasi-finite-over-nagata}.
Choisissons $f \in \mathfrak m$ comme dans le lemme \ref{varieties-lemma-pre-pre-delta-invariant}.
Comme $A' \subset A_f$, il est clair que $A_f = A'_f$. Le support du
$A$-module fini $A'/A$ est donc contenu dans $\{\mathfrak m\}$.
Il en résulte qu'il est de longueur finie d'après
Algèbre, lemme \ref{algebra-lemma-support-point}.
\end{proof}

\begin{definition}
\label{varieties-definition-delta-invariant-algebra}
Soit $A$ un anneau local de Nagata réduit de dimension $1$.
L'{\it invariant $\delta$ de $A$} est $\text{length}_A(A'/A)$,
où $A'$ est défini comme dans le lemme \ref{varieties-lemma-pre-delta-invariant}.
\end{definition}

\noindent
Nous démontrons quelques lemmes sur le comportement de cet invariant.

\begin{lemma}
\label{varieties-lemma-delta-invariant-is-zero}
Soit $A$ un anneau local de Nagata réduit de dimension $1$.
L'invariant $\delta$ de $A$ est $0$ si et seulement si
$A$ est un anneau de valuation discrète.
\end{lemma}

\begin{proof}
Si $A$ est un anneau de valuation discrète, alors $A$ est normal et
l'anneau $A'$ est égal à $A$. Réciproquement, si
l'invariant $\delta$ de $A$ est $0$, alors $A$ est intégralement
clos dans son anneau total des fractions, ce qui implique que
$A$ est normal
(Algèbre, lemme \ref{algebra-lemma-characterize-reduced-ring-normal}),
et cela force $A$ à être un anneau de valuation discrète d'après
Algèbre, lemme \ref{algebra-lemma-characterize-dvr}.
\end{proof}

\begin{lemma}
\label{varieties-lemma-normalization-same-after-completion}
Soit $A$ un anneau local de Nagata réduit de dimension $1$.
Soit $A \to A'$ comme dans le lemme \ref{varieties-lemma-pre-delta-invariant}.
Soient $A^h$, $A^{sh}$, resp.\ $A^\wedge$
l'hensélisé, l'hensélisé strict, resp.\ le complété de $A$.
Alors $A^h$, $A^{sh}$, resp. $A^\wedge$ est un anneau local de Nagata
réduit de dimension $1$, et
$A' \otimes_A A^h$, $A' \otimes_A A^{sh}$, resp. $A' \otimes_A A^\wedge$
est la clôture intégrale de $A^h$, $A^{sh}$, resp.\ $A^\wedge$
dans son anneau total des fractions.
\end{lemma}

\begin{proof}
Observons que $A^\wedge$ est réduit; voir
Compléments d'algèbre, lemme \ref{more-algebra-lemma-completion-reduced}.
Les anneaux $A^h$ et $A^{sh}$ sont réduits d'après
Compléments d'algèbre, lemme \ref{more-algebra-lemma-henselization-reduced}.
Les dimensions de $A$, $A^h$, $A^{sh}$ et $A^\wedge$ sont les mêmes
d'après Compléments d'algèbre, lemmes
\ref{more-algebra-lemma-completion-dimension} et
\ref{more-algebra-lemma-henselization-dimension}.

\medskip\noindent
Rappelons qu'un anneau local noethérien est de Nagata si et seulement si les
fibres formelles de $A$ sont géométriquement réduites; voir
Compléments d'algèbre, lemme \ref{more-algebra-lemma-Nagata-local-ring}.
Cette propriété se transmet à $A^h$ et à $A^{sh}$; voir
les résultats de Compléments d'algèbre, section
\ref{more-algebra-section-properties-formal-fibres},
et en particulier le lemme \ref{more-algebra-lemma-henselization-P-ring}.
Le complété est de Nagata d'après
Algèbre, lemme \ref{algebra-lemma-Noetherian-complete-local-Nagata}.

\medskip\noindent
Passons maintenant à l'assertion sur les clôtures intégrales. Avant de poursuivre,
choisissons $f \in \mathfrak m$ comme dans le
lemme \ref{varieties-lemma-pre-pre-delta-invariant}.
Alors l'image de $f$ dans $A^h$, $A^{sh}$ et $A^\wedge$
est manifestement un élément qui vérifie les propriétés (1) -- (6)
dans l'anneau considéré.

\medskip\noindent
Puisque $A \to A'$ est fini, on voit que
$A' \otimes_A A^h$ et $A' \otimes_A A^{sh}$
sont des produits d'anneaux locaux henséliens finis sur $A^h$ et
$A^{sh}$; voir Algèbre, lemme \ref{algebra-lemma-finite-over-henselian}.
Chacun de ces anneaux locaux est l'hensélisé de $A'$ en un
idéal maximal $\mathfrak m' \subset A'$ situé au-dessus de $\mathfrak m$; voir
Algèbre, lemme \ref{algebra-lemma-quasi-finite-henselization} ou
\ref{algebra-lemma-quasi-finite-strict-henselization}.
Ces anneaux locaux sont donc des anneaux intègres normaux d'après
Compléments d'algèbre, lemme \ref{more-algebra-lemma-henselization-normal}.
Il en résulte que $A' \otimes_A A^h$ et $A' \otimes_A A^{sh}$
sont des anneaux normaux. Puisque $A^h \to A' \otimes_A A^h$ et
$A^{sh} \to A' \otimes_A A^{sh}$ sont finis (donc entiers) et puisque
$A' \otimes_A A^h \subset (A^h)_f = Q(A^h)$ et
$A' \otimes_A A^{sh} \subset (A^{sh})_f = Q(A^{sh})$,
on conclut que $A' \otimes_A A^h$ et $A' \otimes_A A^{sh}$ sont
les clôtures intégrales cherchées.

\medskip\noindent
Pour le complété, raisonnons exactement de la même manière. Tout d'abord,
d'après Algèbre, lemme \ref{algebra-lemma-completion-finite-extension}, on a
$$
A' \otimes_A A^\wedge = (A')^\wedge =
\prod\nolimits (A'_{\mathfrak m'})^\wedge
$$
Les anneaux locaux $A'_{\mathfrak m'}$ sont normaux et de dimension $1$
(d'après Algèbre, lemme \ref{algebra-lemma-finite-in-codim-1}, par exemple, ou
la discussion de la
section Algèbre \ref{algebra-section-homomorphism-dimension}).
Ainsi, $A'_{\mathfrak m'}$ est un anneau de valuation discrète; voir
Algèbre, lemme \ref{algebra-lemma-characterize-dvr}.
Par conséquent, $(A'_{\mathfrak m'})^\wedge$ est un anneau de valuation discrète
d'après Compléments d'algèbre, lemme \ref{more-algebra-lemma-completion-dvr}.
Il en résulte que $A' \otimes_A A^\wedge$ est un anneau normal, et
on peut conclure exactement de la même manière que précédemment.
\end{proof}

\begin{lemma}
\label{varieties-lemma-delta-same-after-completion}
Soit $A$ un anneau local de Nagata réduit de dimension $1$.
L'invariant $\delta$ de $A$ est égal à
l'invariant $\delta$ de l'hensélisé, de l'hensélisé strict
ou du complété de $A$.
\end{lemma}

\begin{proof}
Traitons le cas du complété $B = A^\wedge$;
les autres cas se démontrent exactement de la même manière.
Soient $A'$, resp.\ $B'$, la clôture intégrale de $A$, resp.\ de $B$,
dans son anneau total des fractions.
Alors $B' = A' \otimes_A B$ d'après le
lemme \ref{varieties-lemma-normalization-same-after-completion}.
Par conséquent, $B'/B = (A'/A) \otimes_A B$.
L'égalité résulte maintenant du
Algèbre, lemme \ref{algebra-lemma-pullback-module},
et du fait que $B \otimes_A \kappa_A = \kappa_B$.
\end{proof}

\begin{definition}
\label{varieties-definition-delta-invariant}
Soit $k$ un corps. Soit $X$ un $k$-schéma localement algébrique.
Soit $x \in X$ un point tel que $\mathcal{O}_{X, x}$
soit réduit et que $\dim(\mathcal{O}_{X, x}) = 1$.
L'{\it invariant $\delta$ de $X$ en $x$} est
l'invariant $\delta$ de $\mathcal{O}_{X, x}$ défini dans la
définition \ref{varieties-definition-delta-invariant-algebra}.
\end{definition}

\noindent
Cette définition a un sens parce que l'anneau local d'un schéma localement
algébrique est de Nagata d'après
Algèbre, proposition \ref{algebra-proposition-ubiquity-nagata}.
Plus généralement, on peut bien entendu poser cette définition
chaque fois que $x \in X$ est un point d'un schéma tel que
l'anneau local $\mathcal{O}_{X, x}$ soit un anneau de Nagata réduit de dimension $1$.
Il résulte du lemme \ref{varieties-lemma-delta-same-after-completion}
que l'invariant $\delta$ de $X$ en $x$ est
$$
\delta\text{-invariant of }X\text{ at }x =
\delta\text{-invariant of }\mathcal{O}_{X, x}^h =
\delta\text{-invariant of }\mathcal{O}_{X, x}^\wedge
$$
On en conclut que l'invariant $\delta$ est un invariant
de l'anneau local complet du point.

\begin{lemma}
\label{varieties-lemma-delta-invariant-and-change-of-fields}
Soit $k$ un corps. Soit $X$ un $k$-schéma localement algébrique.
Soit $K/k$ une extension de corps et posons $Y = X_K$.
Soit $y \in Y$ d'image $x \in X$.
Supposons que $X$ soit géométriquement réduit en $x$ et que
$\dim(\mathcal{O}_{X, x}) = \dim(\mathcal{O}_{Y, y}) = 1$.
Alors
$$
\delta\text{-invariant of }X\text{ at }x \leq
\delta\text{-invariant of }Y\text{ at }y
$$
\end{lemma}

\begin{proof}
Posons $A = \mathcal{O}_{X, x}$ et $B = \mathcal{O}_{Y, y}$.
D'après le lemme \ref{varieties-lemma-geometrically-reduced-at-point},
on voit que $A$ est géométriquement réduit.
Ainsi, $B$ est une localisation de $A \otimes_k K$.
Soit $A \to A'$ comme dans le lemme \ref{varieties-lemma-pre-delta-invariant}.
Alors
$$
B' = B \otimes_{(A \otimes_k K)} (A' \otimes_k K)
$$
est fini sur $B$, et $B \to B'$ induit
un isomorphisme sur les anneaux totaux des fractions. En effet, choisissons $f \in \mathfrak m_A$
vérifiant (1) -- (6) du lemme \ref{varieties-lemma-pre-pre-delta-invariant};
puisque $\dim(B) = 1$, on voit que $f \in \mathfrak m_B$.
joue le même rôle pour $B$ et l'on voit que $B_f = B'_f$ puisque $A_f = A'_f$.
Soit $B''$ la clôture intégrale de $B$ dans son
anneau total des fractions, comme dans le lemme \ref{varieties-lemma-pre-delta-invariant}.
Alors $B' \subset B''$. Ainsi, l'invariant $\delta$ de $Y$ en $y$ est
$\text{length}_B(B''/B)$ et
\begin{align*}
\text{length}_B(B''/B)
& \geq
\text{length}_B(B'/B) \\
& =
\text{length}_B((A'/A) \otimes_A B) \\
& =
\text{length}_B(B/\mathfrak m_A B) \text{length}_A(A'/A)
\end{align*}
d'après Algèbre, lemme \ref{algebra-lemma-pullback-module},
puisque $A \to B$ est plat (étant une localisation de $A \to A \otimes_k K$).
Comme $\text{length}_A(A'/A)$ est l'invariant $\delta$ de $X$ en $x$
et que
$\text{length}_B(B/\mathfrak m_A B) \geq 1$, le lemme est démontré.
\end{proof}

\begin{lemma}
\label{varieties-lemma-delta-invariant-and-change-of-fields-better}
Soit $k$ un corps. Soit $X$ un $k$-schéma localement algébrique.
Soit $K/k$ une extension de corps, et posons $Y = X_K$.
Soit $y \in Y$, d'image $x \in X$.
Supposons que les hypothèses (a), (b), (c) du
lemme \ref{varieties-lemma-geometrically-normal-in-codim-1}
soient satisfaites pour $x \in X$ et que $\dim(\mathcal{O}_{Y, y}) = 1$.
Alors l'invariant $\delta$
de $X$ en $x$ est l'invariant $\delta$ de $Y$ en $y$.
\end{lemma}

\begin{proof}
Posons $A = \mathcal{O}_{X, x}$ et $B = \mathcal{O}_{Y, y}$.
D'après le lemme \ref{varieties-lemma-geometrically-normal-in-codim-1},
$A$ est géométriquement réduit.
Ainsi, $B$ est une localisation de $A \otimes_k K$.
Soit $A \to A'$ comme dans le lemme \ref{varieties-lemma-pre-delta-invariant}.
D'après le lemme \ref{varieties-lemma-geometrically-normal-in-codim-1},
$A' \otimes_k K$ est normal.
Par conséquent,
$$
B' = B \otimes_{(A \otimes_k K)} (A' \otimes_k K)
$$
est normal, fini sur $B$, et $B \to B'$ induit
un isomorphisme sur les anneaux totaux des fractions. En effet, choisissons $f \in \mathfrak m_A$
satisfaisant (1) -- (6) du lemme \ref{varieties-lemma-pre-pre-delta-invariant} ;
comme $\dim(B) = 1$, on a $f \in \mathfrak m_B$, qui
joue le même rôle pour $B$, et l'on a $B_f = B'_f$ puisque $A_f = A'_f$.
Il s'ensuit que $B \to B'$ est comme dans le lemme \ref{varieties-lemma-pre-delta-invariant}
pour $B$. Il faut donc montrer que
$\text{length}_A(A'/A) =
\text{length}_B(B'/B) = \text{length}_B((A'/A) \otimes_A B)$.
Comme $A \to B$ est plat (étant une localisation de $A \to A \otimes_k K$)
et que $\mathfrak m_B = \mathfrak m_A B$ (car
$B/\mathfrak m_A B$ est de dimension zéro d'après ce qui précède, et est
une localisation de $K \otimes_k \kappa(x)$, lequel est réduit puisque
$\kappa(x)$ est séparable sur $k$), on conclut par
Algèbre, lemme \ref{algebra-lemma-pullback-module}.
\end{proof}





\section{Le nombre de branches}
\label{varieties-section-number-of-branches}

\noindent
Nous avons défini le nombre de branches d'un schéma en un point
dans Propriétés, section \ref{properties-section-unibranch}.

\begin{lemma}
\label{varieties-lemma-number-of-branches}
Soit $X$ un schéma. Supposons que tout ouvert quasi-compact de $X$ possède
un nombre fini de composantes irréductibles. Soit $\nu : X^\nu \to X$
le morphisme de normalisation de $X$. Soit $x \in X$.
\begin{enumerate}
\item Le nombre de branches de $X$ en $x$ est le nombre des
antécédents de $x$ dans $X^\nu$.
\item Le nombre de branches géométriques de $X$ en $x$ est
$\sum_{\nu(x^\nu) = x} [\kappa(x^\nu) : \kappa(x)]_s$.
\end{enumerate}
\end{lemma}

\begin{proof}
Remarquons d'abord que l'hypothèse sur $X$ signifie exactement que la
normalisation est définie, voir
Morphismes, définition \ref{morphisms-definition-normalization}.
Alors le germe $A' = (\nu_*\mathcal{O}_{X^\nu})_x$ est la
clôture intégrale de $A = \mathcal{O}_{X, x}$ dans l'anneau total
des fractions de $A_{red}$, voir
Morphismes, lemme \ref{morphisms-lemma-stalk-normalization}.
Puisque $\nu$ est un morphisme intégral,
les points de $X^\nu$ situés au-dessus de $x$ correspondent
aux idéaux premiers de $A'$ situés au-dessus de l'idéal maximal $\mathfrak m$
de $A$. Comme $A \to A'$ est intégral, cela revient aux
idéaux maximaux de $A'$
(Algèbre, lemmes \ref{algebra-lemma-integral-no-inclusion} et
\ref{algebra-lemma-integral-going-up}).
Le lemme découle donc
de son analogue algébrique :
Compléments d'algèbre, lemme \ref{more-algebra-lemma-number-of-branches-1}.
\end{proof}

\begin{lemma}
\label{varieties-lemma-geometric-branches-and-change-of-fields}
Soit $k$ un corps. Soit $X$ un $k$-schéma localement algébrique.
Soit $K/k$ une extension de corps. Soit $y \in X_K$ un
point d'image $x$ dans $X$. Alors le nombre de
branches géométriques de $X$ en $x$ est le nombre de branches géométriques
de $X_K$ en $y$.
\end{lemma}

\begin{proof}
Écrivons $Y = X_K$ et soient $X^\nu$, resp.\ $Y^\nu$, les normalisés
de $X$, resp.\ $Y$. Considérons le diagramme commutatif
$$
\xymatrix{
Y^\nu \ar[r] \ar[d] & X^\nu_K \ar[r] \ar[d]_{\nu_K} & X^\nu \ar[d]_\nu \\
Y \ar@{=}[r] & Y \ar[r] & X
}
$$
D'après le lemme \ref{varieties-lemma-normalization-and-change-of-fields},
la flèche horizontale supérieure gauche est un homéomorphisme universel.
Elle induit donc des extensions purement inséparables des corps résiduels, voir
Morphismes, lemmes \ref{morphisms-lemma-universal-homeomorphism} et
\ref{morphisms-lemma-universally-injective}.
Ainsi, le nombre de branches géométriques de $Y$ en $y$ est
$\sum_{\nu_K(y') = y} [\kappa(y') : \kappa(y)]_s$
d'après le lemme \ref{varieties-lemma-number-of-branches}. De même,
$\sum_{\nu(x') = x} [\kappa(x') : \kappa(x)]_s$ est le nombre de
branches géométriques de $X$ en $x$. À l'aide de Schémas, lemme
\ref{schemes-lemma-points-fibre-product},
notre assertion découle du fait algébrique suivant :
étant données une extension de corps $l/\kappa$ et
une extension algébrique de corps $m/\kappa$, on a
$$
\sum\nolimits_{m \otimes_\kappa l \to m'} [m' : l']_s = [m : \kappa]_s
$$
où la somme porte sur les corps quotients de $m \otimes_\kappa l$.
On peut le démontrer de manière élémentaire, ou appliquer
le lemme \ref{varieties-lemma-separably-closed-field-connected-components}
à
$$
\Spec(m \otimes_\kappa l) \times_{\Spec(l)} \Spec(\overline{l}) =
\Spec(m) \otimes_{\Spec(\kappa)} \Spec(\overline{l})
\longrightarrow
\Spec(m) \times_{\Spec(\kappa)} \Spec(\overline{\kappa})
$$
car on peut interpréter $[m : \kappa]_s$ comme le nombre de
composantes connexes du membre de droite, et la somme
$\sum_{m \otimes_\kappa l \to m'} [m' : l']_s$
comme le nombre de composantes connexes du membre de gauche.
\end{proof}

\begin{lemma}
\label{varieties-lemma-geometrically-unibranch-and-change-of-fields}
Soit $k$ un corps. Soit $X$ un $k$-schéma localement algébrique.
Soit $K/k$ une extension de corps. Soit $y \in X_K$ un
point d'image $x$ dans $X$. Alors $X$ est géométriquement unibranche
en $x$ si et seulement si $X_K$ est géométriquement unibranche en $y$.
\end{lemma}

\begin{proof}
Cela découle immédiatement du
lemme \ref{varieties-lemma-geometric-branches-and-change-of-fields}
et de Compléments d'algèbre, lemme \ref{more-algebra-lemma-number-of-branches-1}.
\end{proof}

\begin{definition}
\label{varieties-definition-wedge}
Soient $A$ et $A_i$, $1 \leq i \leq n$, des anneaux locaux. On dit
que {\it $A$ est un bouquet de $A_1, \ldots, A_n$}
s'il existe des isomorphismes
$$
\kappa_{A_1} \to \kappa_{A_2} \to \ldots \to \kappa_{A_n}
$$
et si $A$ est isomorphe à l'anneau constitué des $n$-uplets
$(a_1, \ldots, a_n) \in A_1 \times \ldots \times A_n$ qui ont la
même image dans $\kappa_{A_n}$.
\end{definition}

\noindent
Si l'on se donne un anneau de base $\Lambda$ et que $A$ et $A_i$ sont des $\Lambda$-algèbres,
on exige alors que $\kappa_{A_i} \to \kappa_{A_{i + 1}}$
soient des isomorphismes de $\Lambda$-algèbres,
et que $A$ soit isomorphe, comme $\Lambda$-algèbre, à la $\Lambda$-algèbre
constituée des $n$-uplets
$(a_1, \ldots, a_n) \in A_1 \times \ldots \times A_n$ qui ont la
même image dans $\kappa_{A_n}$. En particulier, si $\Lambda = k$ est un corps
et que les applications $k \to \kappa_{A_i}$ sont des isomorphismes, il existe
un choix unique des isomorphismes
$\kappa_{A_i} \to \kappa_{A_{i + 1}}$, et
on parle souvent du {\it bouquet de $A_1, \ldots, A_n$}.

\begin{lemma}
\label{varieties-lemma-delta-number-branches-inequality-sh}
Soit $(A, \mathfrak m)$ un anneau local de Nagata réduit,
strictement hensélien et de dimension $1$. Alors
$$
\delta\text{-invariant of }A \geq \text{number of geometric branches of }A - 1
$$
Si l'égalité a lieu, alors $A$ est un bouquet de $n \geq 1$ anneaux locaux
de valuation discrète strictement henséliens.
\end{lemma}

\begin{proof}
Le nombre de branches géométriques est égal au nombre de branches de $A$
(cela découle immédiatement de
Compléments d'algèbre, définition \ref{more-algebra-definition-number-of-branches}).
Soit $A \to A'$ comme dans le lemme \ref{varieties-lemma-pre-delta-invariant}.
Remarquons que le nombre de branches de $A$ est le nombre
d'idéaux maximaux de $A'$, voir
Compléments d'algèbre, lemme \ref{more-algebra-lemma-number-of-branches-1}.
Il existe une surjection
$$
A'/A \longrightarrow
\left(\prod\nolimits_{\mathfrak m'} \kappa(\mathfrak m')\right)/
\kappa(\mathfrak m)
$$
Comme $\dim_{\kappa(\mathfrak m)} \prod \kappa(\mathfrak m')$
est $\geq$ au nombre de branches, l'inégalité est évidente.

\medskip\noindent
Si l'égalité a lieu, alors $\kappa(\mathfrak m') = \kappa(\mathfrak m)$
pour tout $\mathfrak m' \subset A'$, et la flèche affichée ci-dessus
est un isomorphisme. Comme $A$ est hensélien et que
$A \to A'$ est fini, $A'$ est un produit d'anneaux locaux
henséliens, voir Algèbre, lemme \ref{algebra-lemma-finite-over-henselian}.
Les facteurs sont les anneaux locaux $A'_{\mathfrak m'}$ et, comme
$A'$ est normal, ces facteurs sont des anneaux de valuation discrète
(Algèbre, lemme \ref{algebra-lemma-characterize-dvr}).
La flèche affichée étant un isomorphisme,
$A$ est bien le bouquet de ces anneaux locaux.
\end{proof}

\begin{lemma}
\label{varieties-lemma-delta-number-branches-inequality}
Soit $(A, \mathfrak m)$ un anneau local de Nagata réduit de dimension $1$. Alors
$$
\delta\text{-invariant of }A \geq \text{number of geometric branches of }A - 1
$$
\end{lemma}

\begin{proof}
On peut remplacer $A$ par l'hensélisé strict de $A$ sans
changer l'invariant $\delta$
(lemme \ref{varieties-lemma-delta-same-after-completion}) ni
le nombre de branches géométriques de $A$
(cela découle immédiatement de la définition, voir
Compléments d'algèbre, définition \ref{more-algebra-definition-number-of-branches}).
On peut donc supposer que $A$ est strictement hensélien, puis
appliquer le lemme \ref{varieties-lemma-delta-number-branches-inequality-sh}.
\end{proof}






\section{Normalisation des schémas de dimension un}
\label{varieties-section-normalization-one-dimensional}

\noindent
Le morphisme de normalisation d'un schéma noethérien de dimension
$1$ possède des propriétés remarquablement bonnes grâce au théorème de Krull--Akizuki.

\begin{lemma}
\label{varieties-lemma-normalize-noetherian-dim-1}
Soit $X$ un schéma localement noethérien de dimension $1$.
Soit $\nu : X^\nu \to X$ le morphisme de normalisation. Alors
\begin{enumerate}
\item $\nu$ est intégral, surjectif et induit une bijection
sur les composantes irréductibles,
\item il existe une factorisation $X^\nu \to X_{red} \to X$
et le morphisme $X^\nu \to X_{red}$ est la normalisation
de $X_{red}$,
\item $X^\nu \to X_{red}$ est birationnel,
\item pour tout point fermé $x \in X$, le germe
$(\nu_*\mathcal{O}_{X^\nu})_x$ est la clôture intégrale
de $\mathcal{O}_{X, x}$ dans l'anneau total des fractions
de $(\mathcal{O}_{X, x})_{red} = \mathcal{O}_{X_{red}, x}$,
\item les fibres de $\nu$ sont finies et les extensions de
corps résiduels sont finies,
\item $X^\nu$ est une union disjointe de schémas intègres normaux
localement noethériens, et chaque ouvert affine est le spectre d'un produit
fini d'anneaux de Dedekind.
\end{enumerate}
\end{lemma}

\begin{proof}
Beaucoup de ces résultats sont en fait des propriétés générales du morphisme
de normalisation, voir
Morphismes, lemmes \ref{morphisms-lemma-normalization-reduced},
\ref{morphisms-lemma-stalk-normalization},
\ref{morphisms-lemma-normalization-normal} et
\ref{morphisms-lemma-normalization-birational}.
Ce qui n'est pas clair, c'est que les fibres sont finies,
que les extensions induites des corps résiduels sont finies, et
que $X^\nu$ est localement le spectre d'un
anneau de Dedekind (et est donc noethérien).
Pour le voir, on peut supposer que $X = \Spec(A)$ est
affine, noethérien, de dimension $1$, et que $A$ est réduit.
La description donnée dans
Morphismes, lemme \ref{morphisms-lemma-description-normalization},
permet alors de se ramener au cas où $A$ est un anneau intègre noethérien
de dimension $1$. Dans ce cas, les propriétés voulues
découlent du théorème de Krull--Akizuki sous la forme énoncée dans
Algèbre, lemme \ref{algebra-lemma-integral-closure-Dedekind}.
\end{proof}

\noindent
Bien entendu, le lemme suivant admet une variante lorsque
$X$ n'est pas réduit.

\begin{lemma}
\label{varieties-lemma-prepare-delta-invariant}
Soit $X$ un schéma de Nagata réduit de dimension $1$. Soit $\nu : X^\nu \to X$
le morphisme de normalisation. Soit $x \in X$ un point fermé. Alors
\begin{enumerate}
\item $\nu : X^\nu \to X$ est fini, surjectif et birationnel,
\item $\mathcal{O}_X \subset \nu_*\mathcal{O}_{X^\nu}$ et
$\nu_*\mathcal{O}_{X^\nu}/\mathcal{O}_X$ est une somme directe de
faisceaux gratte-ciel, de valeur $\mathcal{Q}_x$ en $x$, laquelle est non nulle
si et seulement si $x$ est un point singulier de $X$,
\item $A' = (\nu_*\mathcal{O}_{X^\nu})_x$ est la clôture intégrale
de $A = \mathcal{O}_{X, x}$ dans son anneau total des fractions,
\item $\mathcal{Q}_x = A'/A$ est de longueur finie, égale à
l'invariant $\delta$ de $X$ en $x$,
\item $A'$ est un anneau semi-local qui est un produit fini
d'anneaux de Dedekind,
\item $A^\wedge$ est un anneau local complet noethérien
réduit, de dimension $1$,
\item $(A')^\wedge$ est la clôture intégrale de $A^\wedge$
dans son anneau total des fractions,
\item $(A')^\wedge$ est un produit fini d'anneaux
de valuation discrète complets, et
\item $A'/A \cong (A')^\wedge/A^\wedge$.
\end{enumerate}
\end{lemma}

\begin{proof}
Nous pouvons utiliser, et utiliserons, tous les résultats du
lemme \ref{varieties-lemma-normalize-noetherian-dim-1}.
La finitude de $\nu$ découle de
Morphismes, lemme \ref{morphisms-lemma-nagata-normalization}.
Comme $X$ est réduit, de Nagata, et de dimension $1$, le
l'ensemble des points réguliers est un ouvert dense $U \subset X$ d'après
Compléments d'algèbre, proposition \ref{more-algebra-proposition-ubiquity-J-2}.
Comme un schéma régulier est normal, cela montre que $\nu$
est un isomorphisme au-dessus de $U$.
Puisque $\dim(X) \leq 1$, l'ensemble des points au-dessus desquels $\nu$ n'est pas
un isomorphisme est un ensemble discret de points fermés $x \in X$.
En particulier, on obtient une suite exacte courte
$$
0 \to \mathcal{O}_X \to \nu_*\mathcal{O}_{X^\nu} \to
\bigoplus\nolimits_{x \in X \setminus U} \mathcal{Q}_x \to 0
$$
La description des germes de $\nu_*\mathcal{O}_{X^\nu}$
donnée par le lemme \ref{varieties-lemma-normalize-noetherian-dim-1} montre que
$Q_x = A'/A$ est bien de longueur égale à l'invariant $\delta$ de $X$ en $x$.
On a $Q_x \not = 0$ exactement lorsque $x$ est un point singulier, par
exemple d'après le lemme \ref{varieties-lemma-delta-invariant-is-zero}.
La description de $A'$ comme produit de domaines de Dedekind semi-locaux découle
elle aussi du lemme \ref{varieties-lemma-normalize-noetherian-dim-1}.
La relation entre $A$, $A'$ et $(A')^\wedge$
a été vue dans le lemme \ref{varieties-lemma-normalization-same-after-completion}
(et dans sa démonstration).
\end{proof}




\section{Recherche d'ouverts affines}
\label{varieties-section-finding-affine-opens}

\noindent
Nous poursuivons la discussion commencée dans
Propriétés, section \ref{properties-section-finding-affine-opens}.
Il s'avère que l'on peut trouver des ouverts affines contenant un ensemble fini donné
de points de codimension $1$ sur un schéma séparé. Voir
la proposition \ref{varieties-proposition-finite-set-of-points-of-codim-1-in-affine}.

\medskip\noindent
Nous améliorerons le lemme suivant dans
Descente, lemme \ref{descent-lemma-characterize-open-immersion}.

\begin{lemma}
\label{varieties-lemma-characterize-open-immersion}
Soit $f : X \to Y$ un morphisme de schémas. Soit $X^0$ l'ensemble
des points génériques des composantes irréductibles de $X$. Si
\begin{enumerate}
\item $f$ est séparé,
\item il existe un recouvrement ouvert $X = \bigcup U_i$ tel que
$f|_{U_i} : U_i \to Y$ soit une immersion ouverte, et
\item si $\xi, \xi' \in X^0$, $\xi \not = \xi'$, alors $f(\xi) \not = f(\xi')$,
\end{enumerate}
alors $f$ est une immersion ouverte.
\end{lemma}

\begin{proof}
Supposons que $y = f(x) = f(x')$. Choisissons une spécialisation $y_0 \leadsto y$
où $y_0$ est un point générique d'une composante irréductible de $Y$.
Comme $f$ est localement, sur la source, un isomorphisme, on peut choisir des spécialisations
$x_0 \leadsto x$ et $x'_0 \leadsto x'$ qui s'envoient sur $y_0 \leadsto y$.
Remarquons que $x_0, x'_0 \in X^0$. L'hypothèse (3) donne donc $x_0 = x'_0$.
Comme $f$ est séparé, on en conclut que $x = x'$. Ainsi, $f$ est une immersion ouverte.
\end{proof}

\begin{lemma}
\label{varieties-lemma-local-isomorphism}
Soit $X \to S$ un morphisme de schémas. Soit $x \in X$ un point
d'image $s \in S$. Si
\begin{enumerate}
\item $\mathcal{O}_{X, x} = \mathcal{O}_{S, s}$,
\item $S$ est réduit,
\item $X \to S$ est de type fini, et
\item $S$ possède un nombre fini de composantes irréductibles,
\end{enumerate}
alors il existe un voisinage ouvert $U$
de $x$ tel que $f|_U$ soit une immersion ouverte.
\end{lemma}

\begin{proof}
On peut supprimer les composantes irréductibles, en nombre fini, de $S$
qui ne contiennent pas $s$. On peut remplacer $S$ par un voisinage ouvert affine
de $s$. On peut remplacer $X$ par un voisinage ouvert affine
de $x$. Écrivons $S = \Spec(A)$ et $X = \Spec(B)$. Soient $\mathfrak q \subset B$,
resp.\ $\mathfrak p \subset A$, les idéaux premiers correspondant à $x$,
resp.\ $s$. Comme $A$ est réduit et que tous les idéaux premiers minimaux de
$A$ sont contenus dans $\mathfrak p$, on a $A \subset A_\mathfrak p$.
Comme $X \to S$ est de type fini, $B$ est de type fini sur $A$.
Soient $b_1, \ldots, b_n \in B$ des éléments qui engendrent $B$ sur $A$.
Puisque $A_\mathfrak p = B_\mathfrak q$, il existe
$f \in A$, $f \not \in \mathfrak p$,
et $a_i \in A$ tels que $b_i$ et $a_i/f$ aient la même image
dans $B_\mathfrak q$. Il existe donc $g \in B$, $g \not \in \mathfrak q$,
tel que $g(fb_i - a_i) = 0$ dans $B$. Il s'ensuit que l'image de
$A_f \to B_{fg}$ contient les images de $b_1, \ldots, b_n$, et
en particulier l'image de $g$.
Choisissons $n \geq 0$ et $f' \in A$ tels que $f'/f^n$ ait
pour image l'image de $g$ dans $B_{fg}$. Puisque $A_\mathfrak p = B_\mathfrak q$,
on a $f' \not \in \mathfrak p$.
On en conclut que $A_{ff'} \to B_{fg}$ est surjective.
Enfin, comme $A_{ff'} \subset A_\mathfrak p = B_\mathfrak q$ (voir ci-dessus),
l'application $A_{ff'} \to B_{fg}$ est injective, donc bijective.
\end{proof}

\begin{lemma}
\label{varieties-lemma-points-in-affine}
Soit $f : T \to X$ un morphisme de schémas. Soient $X^0$, resp.\ $T^0$,
les ensembles des points génériques des composantes irréductibles.
Soit $t_1, \ldots, t_m \in T$ un ensemble fini de points
d'images $x_j = f(t_j)$. Si
\begin{enumerate}
\item $T$ est affine,
\item $X$ est quasi-séparé,
\item $X^0$ est fini,
\item $f(T^0) \subset X^0$ et $f : T^0 \to X^0$ est injective, et
\item $\mathcal{O}_{X, x_j} = \mathcal{O}_{T, t_j}$,
\end{enumerate}
alors il existe un ouvert affine de $X$ contenant $x_1, \ldots, x_r$.
\end{lemma}

\begin{proof}
À l'aide de Limites, proposition \ref{limits-proposition-affine},
on se ramène immédiatement au cas où $X$ et $T$ sont réduits.
Nous omettons les détails.

\medskip\noindent
Supposons $X$ et $T$ réduits. On peut écrire $T = \lim_{i \in I} T_i$
comme limite projective de schémas de présentation finie sur $X$
à morphismes de transition affines, voir
Limites, lemme \ref{limits-lemma-relative-approximation}.
Choisissons $i \in I$ tel que $T_i$ soit affine, voir
Limites, lemme \ref{limits-lemma-limit-affine}.
Écrivons $T_i = \Spec(R_i)$ et $T = \Spec(R)$.
Soit $R' \subset R$ l'image de $R_i \to R$.
Alors $T' = \Spec(R')$ est affine, réduit, de type fini sur $X$,
et $T \to T'$ est dominant. Pour $j = 1, \ldots, r$, soit $t'_j \in T'$
l'image de $t_j$. Considérons les homomorphismes d'anneaux locaux
$$
\mathcal{O}_{X, x_j} \to
\mathcal{O}_{T', t'_j} \to
\mathcal{O}_{T, t_j}
$$
Notons $(T')^0$ l'ensemble des points génériques des composantes irréductibles
de $T'$. Soit $\xi \leadsto t'_j$ une spécialisation telle que
$\xi \in (T')^0$. Comme $T \to T'$ est dominant, on peut choisir $\eta \in T^0$ 
d'image $\xi$ (attention : a priori, on ne sait pas si $\eta$ se spécialise
en $t_j$). L'hypothèse (3), appliquée à $\eta$, dit que l'image $\theta$
de $\xi$ dans $X$ correspond à un idéal premier minimal de $\mathcal{O}_{X, x_j}$.
En relevant $\xi$ par l'isomorphisme de (5), on obtient une spécialisation
$\eta' \leadsto t_j$ avec $\eta' \in T^0$ d'image $\theta \leadsto x_j$.
L'injectivité de (4) montre que $\eta = \eta'$. Ainsi,
tout idéal premier minimal de $\mathcal{O}_{T', t'_j}$ est contenu dans
un idéal premier minimal de $\mathcal{O}_{T, t_j}$. On en conclut que
$\mathcal{O}_{T', t'_j} \to \mathcal{O}_{T, t_j}$ est injective,
donc les deux homomorphismes ci-dessus sont des isomorphismes.

\medskip\noindent
D'après le lemme \ref{varieties-lemma-local-isomorphism}, il existe un ouvert
$U \subset T'$ contenant tous les points $t'_j$ tel que
$U \to X$ soit un isomorphisme local comme dans le
lemme \ref{varieties-lemma-characterize-open-immersion}.
D'après ce lemme, $U \to X$ est une immersion ouverte.
Enfin, d'après
Propriétés, lemme \ref{properties-lemma-ample-finite-set-in-affine},
il existe un ouvert $W \subset U \subset T'$
contenant tous les $t'_j$. L'image de $W$ dans $X$ est
l'ouvert affine recherché.
\end{proof}

\begin{lemma}
\label{varieties-lemma-finite-set-codim-1-points-in-affine}
Soit $X$ un schéma intègre séparé. Soit $x_1, \ldots, x_r \in X$
un ensemble fini de points tel que l'anneau local $\mathcal{O}_{X, x_i}$
soit noethérien de dimension $\leq 1$. Alors il existe un sous-schéma
ouvert affine de $X$ contenant tous les $x_1, \ldots, x_r$.
\end{lemma}

\begin{proof}
Soit $K$ le corps des fonctions rationnelles de $X$.
Posons $A_i = \mathcal{O}_{X, x_i}$. Alors $A_i \subset K$ et $K$
est le corps des fractions de $A_i$. Puisque $X$ est séparé et que
$x_i \not = x_j$, il ne peut exister d'anneau de valuation $\mathcal{O} \subset K$
dominant à la fois $A_i$ et $A_j$. En effet, si l'on considère le
diagramme
$$
\xymatrix{
\Spec(\mathcal{O}) \ar[r] \ar[d] & \Spec(A_1) \ar[d] \\
\Spec(A_2) \ar[r] & X
}
$$
et que l'on applique le critère valuatif de séparation
(Schémas, lemme \ref{schemes-lemma-separated-implies-valuative}),
on obtiendrait $x_i = x_j$. Ainsi,
le lemme \ref{varieties-lemma-glue-separated} montre
que $A_i \otimes A_j \to K$ est surjectif pour tout $i \not = j$.
D'après le lemme \ref{varieties-lemma-glue-a-bunch-of-local-rings},
$A = A_1 \cap \ldots \cap A_r$ est un anneau noethérien
semi-local possédant exactement $r$ idéaux maximaux
$\mathfrak m_1, \ldots, \mathfrak m_r$ tels que $A_i = A_{\mathfrak m_i}$.
De plus,
$$
\Spec(A) = \Spec(A_1) \cup \ldots \cup \Spec(A_r)
$$
est un recouvrement ouvert, et l'intersection de deux membres distincts de ce
recouvrement est $\Spec(K)$. Par conséquent, les morphismes $\Spec(A_i) \to X$
donnés se recollent en un morphisme de schémas
$$
\Spec(A) \longrightarrow X
$$
qui envoie $\mathfrak m_i$ sur $x_i$ et induit des isomorphismes d'anneaux locaux.
Le résultat découle donc du lemme \ref{varieties-lemma-points-in-affine}.
\end{proof}

\begin{lemma}
\label{varieties-lemma-extra-silly}
Soient $A$ un anneau, $I \subset A$ un idéal,
$\mathfrak p_1, \ldots, \mathfrak p_r$ des idéaux premiers de $A$, et
$\overline{f} \in A/I$ un élément. Si $I \not \subset \mathfrak p_i$
pour tout $i$, alors il existe $f \in A$, $f \not \in \mathfrak p_i$,
qui s'envoie sur $\overline{f}$ dans $A/I$.
\end{lemma}

\begin{proof}
On peut supposer qu'il n'y a aucune relation d'inclusion entre les $\mathfrak p_i$
(quitte à supprimer les idéaux premiers les plus petits). Choisissons d'abord un $f \in A$ relevant
$\overline{f}$. Soit $S$ l'ensemble des $s \in \{1, \ldots, r\}$ tels
que $f \in \mathfrak p_s$. Si $S$ est vide, la démonstration est terminée. Sinon,
considérons l'idéal $J = I \prod_{i \not \in S} \mathfrak p_i$.
Notons que $J$ n'est contenu dans aucun $\mathfrak p_s$ pour $s \in S$,
car il n'y a aucune relation d'inclusion entre les $\mathfrak p_i$ et
que $I$ n'est contenu dans aucun $\mathfrak p_i$.
On peut donc choisir $g \in J$, $g \not \in \mathfrak p_s$ pour
$s \in S$, d'après Algèbre, lemme \ref{algebra-lemma-silly}.
Alors $f + g$ répond aux conditions du lemme.
\end{proof}

\begin{lemma}
\label{varieties-lemma-finite-set-codim-1-points-in-affine-per-component}
Soit $X$ un schéma. Soit $T \subset X$ un ensemble fini de points. Supposons que
\begin{enumerate}
\item $X$ possède un nombre fini de composantes irréductibles $Z_1, \ldots, Z_t$, et
\item $Z_i \cap T$ soit contenu dans un ouvert affine du sous-schéma
réduit induit correspondant à $Z_i$.
\end{enumerate}
Alors il existe un sous-schéma ouvert affine de $X$ contenant $T$.
\end{lemma}

\begin{proof}
La proposition \ref{limits-proposition-affine} de Limites
permet de se ramener immédiatement au cas où $X$
est réduit. Nous omettons les détails. Dans la suite de la démonstration, nous
munissons tout sous-ensemble fermé de $X$ de la structure induite de sous-schéma
fermé réduit.

\medskip\noindent
Nous allons montrer par récurrence que l'on peut trouver un ouvert affine
$U \subset Z_1 \cup \ldots \cup Z_r$ contenant
$T \cap (Z_1 \cup \ldots \cup Z_r)$. Pour $r = 1$, cela résulte de l'hypothèse.
Soit $r > 1$, et soit $U \subset Z_1 \cup \ldots \cup Z_{r -  1}$
un ouvert affine contenant $T \cap (Z_1 \cup \ldots \cup Z_{r - 1})$.
Soit $V \subset X_r$ un ouvert affine contenant $T \cap Z_r$ (il existe par
hypothèse). Alors $U \cap V$ contient
$T \cap ( Z_1 \cup \ldots \cup Z_{r - 1} ) \cap Z_r$.
Par conséquent,
$$
\Delta = (U \cap Z_r) \setminus (U \cap V)
$$
ne contient aucun élément de $T$. Notons que $\Delta$ est un sous-ensemble fermé
de $U$. Par évitement des idéaux premiers (Algèbre, lemme \ref{algebra-lemma-silly}),
on peut trouver un ouvert principal $U'$ de $U$ contenant $T \cap U$ et évitant
$\Delta$, c'est-à-dire tel que $U' \cap Z_r \subset U \cap V$.
Après avoir remplacé $U$ par $U'$, on peut supposer que $U \cap V$ est
fermé dans $U$.

\medskip\noindent
Le même raisonnement montre en outre que l'ensemble
$\Delta' = (U \cap (Z_1 \cup \ldots \cup Z_{r - 1})) \setminus (U \cap V)$
ne contient aucun élément de $T$; il existe donc $h \in \mathcal{O}(V)$
tel que $D(h) \subset V$ contienne $T \cap V$ et que
$U \cap D(h) \subset U \cap V$. Comme $U \cap V$ est fermé dans $U$,
on peut appliquer le lemme \ref{varieties-lemma-extra-silly}
pour trouver un élément $g \in \mathcal{O}(U)$ dont la restriction à $U \cap V$
est égale à la restriction de $h$ à $U \cap V$ et tel que
$T \cap U \subset D(g)$. On peut alors remplacer $U$ par $D(g)$
et $V$ par $D(h)$, de manière que $U \cap V$ soit
fermé à la fois dans $U$ et dans $V$. Dans ce cas, le schéma $U \cup V$ est
affine d'après Limites, lemme \ref{limits-lemma-affines-glued-in-closed-affine}.
Cela établit l'étape de récurrence, et donc le lemme.
\end{proof}

\noindent
Voici une conséquence des résultats qui précèdent.

\begin{proposition}
\label{varieties-proposition-finite-set-of-points-of-codim-1-in-affine}
Soit $X$ un schéma séparé tel que tout ouvert quasi-compact
possède un nombre fini de composantes irréductibles. Soient $x_1, \ldots, x_r \in X$
des points tels que l'anneau local $\mathcal{O}_{X, x_i}$ soit noethérien
de dimension $\leq 1$. Alors il existe un sous-schéma ouvert affine
de $X$ contenant tous les $x_1, \ldots, x_r$.
\end{proposition}

\begin{proof}
On peut remplacer $X$ par un ouvert quasi-compact contenant $x_1, \ldots, x_r$;
on peut donc supposer que $X$ possède un nombre fini de composantes irréductibles.
Le lemme \ref{varieties-lemma-finite-set-codim-1-points-in-affine-per-component}
ramène au cas où $X$ est intègre. Ce cas est traité par le
lemme \ref{varieties-lemma-finite-set-codim-1-points-in-affine}.
\end{proof}





\section{Courbes}
\label{varieties-section-curves}

\noindent
Dans le projet Champs, nous adopterons la définition
suivante d'une courbe.

\begin{definition}
\label{varieties-definition-curve}
Soit $k$ un corps. Une {\it courbe} est une variété de dimension $1$ sur $k$.
\end{definition}

\noindent
Deux exemples classiques de courbes sur $k$ sont la droite affine $\mathbf{A}^1_k$
et la droite projective $\mathbf{P}^1_k$. Le schéma $X = \Spec(k[x, y]/(f))$
est une courbe si et seulement si $f \in k[x, y]$ est irréductible.

\medskip\noindent
Notre définition d'une courbe présente les mêmes difficultés que celle d'une
variété; voir la discussion qui suit la définition \ref{varieties-definition-variety}.
Elle signifie en outre que toute courbe est munie d'un corps de définition spécifié.
Par exemple, $X = \Spec(\mathbf{C}[x])$ est une courbe sur $\mathbf{C}$,
mais on peut aussi le considérer comme une courbe sur $\mathbf{R}$. Le schéma
$\Spec(\mathbf{Z})$ n'est pas une courbe, bien que les schémas $\Spec(\mathbf{Z})$
et $\mathbf{A}^1_{\mathbf{F}_p}$ se comportent de manière analogue à bien des égards.

\begin{lemma}
\label{varieties-lemma-proper-minus-point}
Soit $X$ un schéma séparé irréductible de dimension $> 0$ sur un corps $k$.
Soit $x \in X$ un point fermé. Le sous-schéma ouvert $X \setminus \{x\}$
n'est pas propre sur $k$.
\end{lemma}

\begin{proof}
Puisque $X$ est irréductible, $U = X \setminus \{x\}$ n'est pas fermé dans $X$.
En particulier, l'immersion $U \to X$ n'est pas propre.
D'après Morphismes, lemme \ref{morphisms-lemma-image-proper-scheme-closed}
(on utilise ici que $X$ est séparé), $U \to \Spec(k)$ n'est pas propre non plus.
\end{proof}

\begin{lemma}
\label{varieties-lemma-dim-1-quasi-projective}
Soit $X$ un schéma séparé de type fini sur un corps $k$.
Si $\dim(X) \leq 1$, alors $X$ est H-quasi-projectif sur $k$.
\end{lemma}

\begin{proof}
D'après la proposition \ref{varieties-proposition-dim-1-noetherian-separated-has-ample},
le schéma $X$ possède un faisceau inversible ample $\mathcal{L}$.
D'après Morphismes, lemme \ref{morphisms-lemma-quasi-projective-finite-type-over-S},
on voit que $X$ est isomorphe à un sous-schéma localement
fermé de $\mathbf{P}^n_k$ au-dessus de $\Spec(k)$. C'est
la définition de H-quasi-projectif sur $k$; voir
Morphismes, définition \ref{morphisms-definition-quasi-projective}.
\end{proof}

\begin{lemma}
\label{varieties-lemma-dim-1-proper-projective}
Soit $X$ un schéma propre sur un corps $k$.
Si $\dim(X) \leq 1$, alors $X$ est H-projectif sur $k$.
\end{lemma}

\begin{proof}
D'après le lemme \ref{varieties-lemma-dim-1-quasi-projective}, $X$ est un
sous-schéma localement fermé de $\mathbf{P}^n_k$ pour un certain corps $k$.
Puisque $X$ est propre sur $k$, il s'ensuit que $X$ est un sous-schéma fermé
de $\mathbf{P}^n_k$
(Morphismes, lemme \ref{morphisms-lemma-image-proper-scheme-closed}).
\end{proof}

\begin{lemma}
\label{varieties-lemma-dim-1-projective-completion}
Soit $X$ un schéma séparé de type fini sur $k$.
Si $\dim(X) \leq 1$, alors il existe une immersion ouverte
$j : X \to \overline{X}$ possédant les propriétés suivantes :
\begin{enumerate}
\item $\overline{X}$ est H-projectif sur $k$, c'est-à-dire que $\overline{X}$
est un sous-schéma fermé de $\mathbf{P}^d_k$ pour un certain $d$,
\item $j(X) \subset \overline{X}$ est dense et schématiquement
dense,
\item $\overline{X} \setminus X = \{x_1, \ldots, x_n\}$
pour certains points fermés $x_i \in \overline{X}$.
\end{enumerate}
\end{lemma}

\begin{proof}
D'après le lemme \ref{varieties-lemma-dim-1-quasi-projective}, on peut supposer que $X$ est un
sous-schéma localement fermé de $\mathbf{P}^d_k$ pour un certain $d$. Soit
$\overline{X} \subset \mathbf{P}^d_k$ l'image schématique
de $X \to \mathbf{P}^d_k$; voir Morphismes, définition
\ref{morphisms-definition-scheme-theoretic-image}.
La description donnée dans
Morphismes, lemme \ref{morphisms-lemma-quasi-compact-immersion},
entraîne les propriétés (1) et (2).
On a alors $\dim(X) = 1 \Rightarrow \dim(\overline{X}) = 1$, par exemple en
considérant les points génériques; voir
le lemme \ref{varieties-lemma-dimension-locally-algebraic}.
Comme $\overline{X}$ est noethérien, on en déduit que
$\overline{X} \setminus X = \{x_1, \ldots, x_n\}$
est un ensemble fini de points fermés.
\end{proof}

\begin{lemma}
\label{varieties-lemma-reduced-dim-1-projective-completion}
Soit $X$ un schéma séparé de type fini sur $k$.
Si $X$ est réduit et $\dim(X) \leq 1$, alors il existe
une immersion ouverte $j : X \to \overline{X}$ telle que
\begin{enumerate}
\item $\overline{X}$ est H-projectif sur $k$, c'est-à-dire que $\overline{X}$
est un sous-schéma fermé de $\mathbf{P}^d_k$ pour un certain $d$,
\item $j(X) \subset \overline{X}$ est dense et schématiquement
dense,
\item $\overline{X} \setminus X = \{x_1, \ldots, x_n\}$
pour certains points fermés $x_i \in \overline{X}$,
\item les anneaux locaux $\mathcal{O}_{\overline{X}, x_i}$
sont des anneaux de valuation discrète pour $i = 1, \ldots, n$.
\end{enumerate}
\end{lemma}

\begin{proof}
Soit $j : X \to \overline{X}$ comme dans le
lemme \ref{varieties-lemma-dim-1-projective-completion}.
Considérons la normalisation $X'$ de $\overline{X}$ dans
$X$. D'après le lemme \ref{varieties-lemma-relative-normalization-finite},
le morphisme $X' \to \overline{X}$ est fini.
D'après Morphismes, lemme \ref{morphisms-lemma-finite-projective},
$X' \to \overline{X}$ est projectif. D'après Morphismes, lemme
\ref{morphisms-lemma-projective-over-quasi-projective-is-H-projective},
$X' \to \overline{X}$ est H-projectif.
D'après Morphismes, lemme \ref{morphisms-lemma-H-projective-composition},
$X' \to \Spec(k)$ est H-projectif.
Soit $\{x'_1, \ldots, x'_m\} \subset X'$ l'image réciproque
de $\{x_1, \ldots, x_n\} = \overline{X} \setminus X$.
Alors $\dim(\mathcal{O}_{X', x'_i}) = 1$ pour tout $1 \leq i \leq m$.
Par conséquent, les anneaux locaux $\mathcal{O}_{X', x'}$
sont des anneaux de valuation discrète d'après
Morphismes, lemme
\ref{morphisms-lemma-relative-normalization-normal-codim-1}.
Alors $X \to X'$ et $\{x'_1, \ldots, x'_m\}$ répondent aux conditions voulues.
\end{proof}

\begin{lemma}
\label{varieties-lemma-dim-1-projective-completion-after-insep}
Soit $X$ un schéma séparé de type fini sur $k$ tel que $\dim(X) \leq 1$.
Alors il existe un diagramme commutatif
$$
\xymatrix{
\overline{Y}_1 \amalg \ldots \amalg \overline{Y}_n \ar[rd] &
Y_1 \amalg \ldots \amalg Y_n \ar[r]_-\nu \ar[d] \ar[l]^j &
X_{k'} \ar[r] \ar[d] &
X \ar[d]^f \\
& \Spec(k'_1) \amalg \ldots \amalg \Spec(k'_n) \ar[r] &
\Spec(k') \ar[r] &
\Spec(k)
}
$$
de schémas possédant les propriétés suivantes :
\begin{enumerate}
\item $k'/k$ est une extension finie purement inséparable de corps,
\item $\nu$ est la normalisation de $X_{k'}$,
\item $j$ est une immersion ouverte d'image dense,
\item $k'_i/k'$ est une extension finie séparable pour $i = 1, \ldots, n$,
\item $\overline{Y}_i$ est lisse, projectif, géométriquement irréductible
et de dimension $\leq 1$ sur $k'_i$.
\end{enumerate}
\end{lemma}

\begin{proof}
Puisque l'on peut remplacer $X$ par sa réduction, on peut et on va supposer $X$ réduit.
Choisissons $X \to \overline{X}$ comme dans le
lemme \ref{varieties-lemma-reduced-dim-1-projective-completion}.
Si le lemme est établi pour $\overline{X}$, il en découle
pour $X$ (détails omis).
On peut donc et on va supposer $X$ projectif.

\medskip\noindent
Choisissons une extension finie purement inséparable $k'/k$ telle que la normalisation de
$X_{k'}$ soit géométriquement normale sur $k'$; voir
le lemme \ref{varieties-lemma-finite-extension-geometrically-normal}.
Notons $Y = (X_{k'})^\nu$ la normalisation; pour les propriétés
de la normalisation, voir la section \ref{varieties-section-normalization}.
Alors $Y$ est géométriquement régulier, car normalité et régularité coïncident
en dimension $\leq 1$; voir
Propriétés, lemme \ref{properties-lemma-normal-dimension-1-regular}.
Ainsi, $Y$ est lisse sur $k'$ d'après
le lemme \ref{varieties-lemma-geometrically-regular-smooth}.
Soit $Y = Y_1 \amalg \ldots \amalg Y_n$ la décomposition
de $Y$ en composantes irréductibles.
Posons $k'_i = \Gamma(Y_i, \mathcal{O}_{Y_i})$.
Ce sont des extensions finies séparables de $k'$ d'après
le lemme \ref{varieties-lemma-proper-geometrically-reduced-global-sections}.
Le lemme \ref{varieties-lemma-baby-stein} achève la démonstration.
\end{proof}

\begin{lemma}
\label{varieties-lemma-regular-point-on-curve}
Soit $k$ un corps. Soit $X$ une courbe sur $k$. Soit $x \in X$ un point
fermé. Considérons $x$ comme un sous-schéma fermé (réduit) de $X$, de faisceau
d'idéaux $\mathcal{I}$. Les conditions suivantes sont équivalentes :
\begin{enumerate}
\item $\mathcal{O}_{X, x}$ est régulier,
\item $\mathcal{O}_{X, x}$ est normal,
\item $\mathcal{O}_{X, x}$ est un anneau de valuation discrète,
\item $\mathcal{I}$ est un $\mathcal{O}_X$-module inversible,
\item $x$ est un diviseur de Cartier effectif sur $X$.
\end{enumerate}
Si $k$ est parfait ou si $\kappa(x)$ est séparable sur $k$,
ces conditions sont aussi équivalentes à
\begin{enumerate}
\item[(6)] $X \to \Spec(k)$ est lisse en $x$.
\end{enumerate}
\end{lemma}

\begin{proof}
Puisque $X$ est une courbe, l'anneau local $\mathcal{O}_{X, x}$ est un anneau local
noethérien intègre de dimension $1$ (lemme \ref{varieties-lemma-dimension-locally-algebraic}).
Les conditions (4) et (5) sont équivalentes par définition et équivalent au fait que
$\mathcal{I}_x = \mathfrak m_x \subset \mathcal{O}_{X, x}$ soit engendré par un élément
(Diviseurs, lemme \ref{divisors-lemma-effective-Cartier-in-points}).
L'équivalence de (1), (2), (3), (4) et (5) découle donc du
lemme \ref{algebra-lemma-characterize-dvr} d'Algèbre. La dernière assertion
découle du lemme \ref{varieties-lemma-dense-smooth-open-variety-over-perfect-field}
lorsque $k$ est parfait. Si $\kappa(x)/k$ est séparable, l'équivalence
découle du lemme \ref{algebra-lemma-separable-smooth} d'Algèbre.
\end{proof}

\begin{remark}
\label{varieties-remark-smooth-projective-compactification}
Soit $k$ un corps. Soit $X$ une courbe régulière sur $k$. D'après les
lemmes \ref{varieties-lemma-regular-point-on-curve} et
\ref{varieties-lemma-reduced-dim-1-projective-completion},
il existe une courbe projective non singulière $\overline{X}$
qui est une compactification de $X$, c'est-à-dire qu'il existe une
immersion ouverte $j : X \to \overline{X}$ dont le complément
est formé d'un nombre fini de points fermés. Si $k$ est parfait,
alors $X$ et $\overline{X}$ sont lisses sur $k$ et
$\overline{X}$ est une compactification projective lisse de $X$.
\end{remark}

\noindent
Observons que si un schéma affine $X$ sur $k$ est propre sur $k$,
alors $X$ est fini sur $k$ (Morphismes, lemme
\ref{morphisms-lemma-finite-proper}) et est donc de
dimension $0$
(Algèbre, lemme \ref{algebra-lemma-finite-dimensional-algebra} et
proposition \ref{algebra-proposition-dimension-zero-ring}).
Ainsi, un schéma de dimension $> 0$ sur $k$ ne peut être à la fois affine et
propre sur $k$. Les deux possibilités du lemme suivant
s'excluent donc mutuellement.

\begin{lemma}
\label{varieties-lemma-curve-affine-projective}
Soit $X$ une courbe sur $k$. Alors soit $X$ est un schéma affine, soit $X$
est H-projectif sur $k$.
\end{lemma}

\begin{proof}
Choisissons $X \to \overline{X}$ tel que
$\overline{X} \setminus X = \{x_1, \ldots, x_r\}$ comme dans le
lemme \ref{varieties-lemma-reduced-dim-1-projective-completion}.
Alors $\overline{X}$ est également une courbe.
Si $r = 0$, alors $X = \overline{X}$ est H-projectif sur $k$.
On peut donc supposer $r \geq 1$; il s'agit de montrer que
$X$ est affine. D'après le lemme \ref{varieties-lemma-open-in-affine-curve-affine},
il suffit de montrer que $\overline{X} \setminus \{x_1\}$ est affine.
Cela nous ramène à l'assertion énoncée au paragraphe suivant.

\medskip\noindent
Soit $X$ une courbe H-projective sur $k$. Soit $x \in X$ un point fermé
tel que $\mathcal{O}_{X, x}$ soit un anneau de valuation discrète. Affirmation :
$U = X \setminus \{x\}$ est affine. D'après le lemme \ref{varieties-lemma-regular-point-on-curve},
le point $x$ définit un diviseur de Cartier effectif de $X$.
Pour $n \geq 1$, notons $nx = x + \ldots + x$ la somme de $n$ termes; voir
Diviseurs, définition \ref{divisors-definition-sum-effective-Cartier-divisors}.
Notons $\mathcal{O}_{nx}$ le faisceau structural de $nx$, considéré comme un
module cohérent sur $X$.
Tout module inversible sur le schéma local $nx$ étant trivial,
la première suite exacte courte de
Diviseurs, remarque \ref{divisors-remark-ses-regular-section},
s'écrit
$$
0 \to \mathcal{O}_X
\xrightarrow{1} \mathcal{O}_X(nx) \to \mathcal{O}_{nx} \to 0
$$
dans notre cas. Notons que $\dim_k H^0(X, \mathcal{O}_{nx}) \geq n$.
En effet, d'après le lemme \ref{varieties-lemma-chi-tensor-finite}, on a
$H^0(X, \mathcal{O}_{nx}) = \mathcal{O}_{X, x}/(\pi^n)$, où $\pi$ est une uniformisante de
$\mathcal{O}_{X, x}$, et les puissances $\pi^i$ ont des images
$k$-linéairement indépendantes dans $\mathcal{O}_{X, x}/(\pi^n)$
pour $i = 0, 1, \ldots, n - 1$. On a $\dim_k H^1(X, \mathcal{O}_X) < \infty$
d'après Cohomologie des schémas, lemme
\ref{coherent-lemma-proper-over-affine-cohomology-finite}.
Si $n > \dim_k H^1(X, \mathcal{O}_X)$,
la suite exacte longue de cohomologie montre
qu'il existe $s \in \Gamma(X, \mathcal{O}_X(nx))$
qui n'est pas une section de $\mathcal{O}_X$.
Si l'on choisit $n$ minimal avec cette propriété, alors $s$
a pour image un générateur du germe
$\left(\mathcal{O}_X(nx)\right)_x$, faute de quoi il
définirait une section de $\mathcal{O}_X((n - 1)x) \subset \mathcal{O}_X(nx)$.
Pour ce $n$, on en déduit que $s_0 = 1$ et $s_1 = s$
engendrent le module inversible $\mathcal{L} = \mathcal{O}_X(nx)$.

\medskip\noindent
Considérons le morphisme correspondant
$f = \varphi_{\mathcal{L}, (s_0, s_1)} : X \to \mathbf{P}^1_k$
de Constructions, section \ref{constructions-section-projective-space}.
Notons que l'image réciproque de $D_{+}(T_0)$ est $U = X \setminus \{x\}$,
car la section $s_0$ de $\mathcal{L}$ ne s'annule qu'en $x$.
En particulier, $f$ n'est pas constant, c'est-à-dire que $\Im(f)$ contient
plus d'un point. Par conséquent, $f$ envoie nécessairement le point générique
$\eta$ de $X$ sur le point générique de $\mathbf{P}^1_k$.
Ainsi, si $y \in \mathbf{P}^1_k$ est un point fermé, $f^{-1}(\{y\})$
est un fermé de $X$ ne contenant pas $\eta$, et est donc fini.
Enfin, $f$ est propre\footnote{En effet, une variété H-projective
est une variété propre d'après Morphismes, lemme
\ref{morphisms-lemma-projective-is-quasi-projective-proper}.
Un morphisme de variétés dont la source est une variété propre est un morphisme propre
d'après Morphismes, lemme \ref{morphisms-lemma-image-proper-scheme-closed}.}.
D'après Cohomologie des schémas, lemme
\ref{coherent-lemma-proper-finite-fibre-finite-in-neighbourhood}\footnote{On
peut éviter d'utiliser ce lemme, qui repose sur le théorème des fonctions
formelles. En effet, $X$ est projectif, et il suffit donc de montrer
qu'un morphisme propre $f : X \to Y$ à fibres finies entre des schémas quasi-projectifs
sur $k$ est fini. Pour cela, on choisit un ouvert affine de $X$
contenant la fibre de $f$ au-dessus d'un point $y$, en utilisant le fait que tout ensemble fini de
points d'un schéma quasi-projectif sur $k$ est contenu dans un ouvert affine.
Quitte à restreindre $Y$ à un petit voisinage affine de $y$, on se ramène au
cas d'un morphisme propre entre schémas affines. Un tel morphisme est fini d'après
Morphismes, lemme \ref{morphisms-lemma-integral-universally-closed}.},
on conclut que $f$ est fini. Ainsi, $U = f^{-1}(D_{+}(T_0))$
est affine.
\end{proof}

\noindent
Le lemme suivant, combiné au lemme \ref{varieties-lemma-proper-minus-point},
montre que, si $X$ est un schéma séparé de dimension $1$ et
de type fini sur $k$, alors $X \setminus Z$ est affine dès que le
sous-ensemble fermé $Z$ rencontre toute composante irréductible de $X$.

\begin{lemma}
\label{varieties-lemma-dim-1-nonproper-affine}
Soit $X$ un schéma séparé de type fini sur $k$.
Si $\dim(X) \leq 1$ et si aucune composante irréductible de $X$
n'est propre de dimension $1$, alors $X$ est affine.
\end{lemma}

\begin{proof}
Soit $X = \bigcup X_i$ la décomposition de $X$ en composantes irréductibles.
Nous considérons $X_i$ comme un schéma intègre (muni de la structure
de schéma réduit induite, voir Schémas, définition
\ref{schemes-definition-reduced-induced-scheme}).
En particulier, $X_i$ est un singleton (donc affine) ou une courbe,
donc est affine par le lemme \ref{varieties-lemma-curve-affine-projective}.
Alors $\coprod X_i \to X$ est fini surjectif et $\coprod X_i$ est affine.
Nous voyons donc que $X$ est affine par
Cohomologie des schémas, lemme
\ref{coherent-lemma-image-affine-finite-morphism-affine-Noetherian}.
\end{proof}






\section{Degrés sur les courbes}
\label{varieties-section-divisors-curves}

\noindent
Nous commençons par définir le degré d'un faisceau inversible et, plus généralement,
celui d'un faisceau localement libre sur un schéma propre
de dimension $1$ sur un corps. Dans la section \ref{varieties-section-euler},
nous avons défini la caractéristique d'Euler-Poincaré
d'un faisceau cohérent $\mathcal{F}$ sur un schéma propre $X$ sur un corps
$k$ par la formule
$$
\chi(X, \mathcal{F}) = \sum (-1)^i \dim_k H^i(X, \mathcal{F}).
$$

\begin{definition}
\label{varieties-definition-degree-invertible-sheaf}
Soit $k$ un corps, soit $X$ un schéma propre de dimension $\leq 1$
sur $k$, et soit $\mathcal{L}$ un $\mathcal{O}_X$-module inversible.
Le {\it degré} de $\mathcal{L}$ est défini par
$$
\deg(\mathcal{L}) = \chi(X, \mathcal{L}) - \chi(X, \mathcal{O}_X)
$$
Plus généralement, si $\mathcal{E}$ est un faisceau localement libre de rang $n$,
nous définissons le {\it degré} de $\mathcal{E}$ par
$$
\deg(\mathcal{E}) = \chi(X, \mathcal{E}) - n\chi(X, \mathcal{O}_X)
$$
\end{definition}

\noindent
Remarquons que celui-ci dépend du triplet $\mathcal{E}/X/k$.
Si $X$ n'est pas connexe et si $\mathcal{E}$ est localement libre de type fini (mais
pas nécessairement de rang constant), on peut modifier la définition en sommant les degrés
des restrictions de $\mathcal{E}$ aux composantes connexes de $X$.
Si $\mathcal{E}$ est seulement un faisceau cohérent, il existe plusieurs
manières d'étendre la définition\footnote{Si $X$ est une courbe propre
et si $\mathcal{F}$ est un faisceau cohérent sur $X$, on définit souvent
le degré comme $\chi(X, \mathcal{F}) - r\chi(X, \mathcal{O}_X)$,
où $r = \dim_{\kappa(\xi)} \mathcal{F}_\xi$ est le rang de $\mathcal{F}$
au point générique $\xi$ de $X$.}.
Nous montrons dans une série de lemmes que cette définition possède toutes les propriétés
attendues du degré.

\begin{lemma}
\label{varieties-lemma-degree-base-change}
Soit $k'/k$ une extension de corps. Soit $X$ un schéma propre de
dimension $\leq 1$ sur $k$. Soit $\mathcal{E}$ un
$\mathcal{O}_X$-module localement libre de rang constant $n$. Alors le degré de
$\mathcal{E}/X/k$ est égal au degré de
$\mathcal{E}_{k'}/X_{k'}/k'$.
\end{lemma}

\begin{proof}
Plus précisément, posons $X_{k'} = X \times_{\Spec(k)} \Spec(k')$.
Soit $\mathcal{E}_{k'} = p^*\mathcal{E}$, où $p : X_{k'} \to X$
est la projection. Par
Cohomologie des schémas, lemme \ref{coherent-lemma-flat-base-change-cohomology},
nous avons
$H^i(X_{k'}, \mathcal{E}_{k'}) = H^i(X, \mathcal{E}) \otimes_k k'$
et
$H^i(X_{k'}, \mathcal{O}_{X_{k'}}) = H^i(X, \mathcal{O}_X) \otimes_k k'$.
Nous voyons donc que les caractéristiques d'Euler-Poincaré sont inchangées, et par conséquent
que le degré est inchangé.
\end{proof}

\begin{lemma}
\label{varieties-lemma-degree-additive}
Soit $k$ un corps. Soit $X$ un schéma propre de dimension $\leq 1$
sur $k$. Soit $0 \to \mathcal{E}_1 \to \mathcal{E}_2 \to \mathcal{E}_3 \to 0$
une suite exacte courte de $\mathcal{O}_X$-modules localement libres,
tous de type fini et de rang constant. Alors
$$
\deg(\mathcal{E}_2) = \deg(\mathcal{E}_1) + \deg(\mathcal{E}_3)
$$
\end{lemma}

\begin{proof}
Cela résulte immédiatement de l'additivité des caractéristiques d'Euler-Poincaré
(lemme \ref{varieties-lemma-euler-characteristic-additive})
et de l'additivité des rangs.
\end{proof}

\begin{lemma}
\label{varieties-lemma-degree-birational-pullback}
Soit $k$ un corps. Soit $f : X' \to X$ un morphisme birationnel de
schémas propres de dimension $\leq 1$ sur $k$. Alors
$$
\deg(f^*\mathcal{E}) = \deg(\mathcal{E})
$$
pour tout faisceau localement libre de type fini et de rang constant. Plus généralement,
il suffit que $f$ induise une bijection entre les composantes irréductibles
de dimension $1$ et des isomorphismes entre les anneaux locaux aux
points génériques correspondants.
\end{lemma}

\begin{proof}
Le morphisme $f$ est propre
(Morphismes, lemme \ref{morphisms-lemma-image-proper-scheme-closed})
et ses fibres sont de dimension $\leq 0$. Ainsi, $f$ est fini
(Cohomologie des schémas, lemme
\ref{coherent-lemma-proper-finite-fibre-finite-in-neighbourhood}).
Par conséquent,
$$
Rf_*f^*\mathcal{E} = f_*f^*\mathcal{E} =
\mathcal{E} \otimes_{\mathcal{O}_X} f_*\mathcal{O}_{X'}
$$
Comme $f$ induit un isomorphisme sur les anneaux locaux aux points génériques de
toutes les composantes irréductibles de dimension $1$, nous voyons que le noyau
et le conoyau dans
$$
0 \to \mathcal{K} \to \mathcal{O}_X \to f_*\mathcal{O}_{X'}
\to \mathcal{Q} \to 0
$$
ont des supports de dimension $\leq 0$. Remarquons que tensoriser cette suite par
$\mathcal{E}$ donne encore une suite exacte puisque $\mathcal{E}$ est localement libre.
Nous obtenons
\begin{align*}
\chi(X, \mathcal{E}) - \chi(X', f^*\mathcal{E})
& =
\chi(X, \mathcal{E}) - \chi(X, f_*f^*\mathcal{E}) \\
& =
\chi(X, \mathcal{E}) - \chi(X, \mathcal{E} \otimes f_*\mathcal{O}_{X'}) \\
& =
\chi(X, \mathcal{K} \otimes \mathcal{E}) -
\chi(X, \mathcal{Q} \otimes \mathcal{E}) \\
& =
n\chi(X, \mathcal{K}) -
n\chi(X, \mathcal{Q}) \\
& =
n\chi(X, \mathcal{O}_X) - n\chi(X, f_*\mathcal{O}_{X'}) \\
& =
n\chi(X, \mathcal{O}_X) - n\chi(X', \mathcal{O}_{X'})
\end{align*}
ce qui prouve l'assertion. La première égalité vaut parce que $f$ est fini, voir
Cohomologie des schémas, lemme \ref{coherent-lemma-relative-affine-cohomology}.
La deuxième égalité résulte de la formule de projection, voir
Cohomologie, lemme \ref{cohomology-lemma-projection-formula}.
La troisième résulte de l'additivité des caractéristiques d'Euler-Poincaré, voir
le lemme \ref{varieties-lemma-euler-characteristic-additive}.
La quatrième résulte du lemme \ref{varieties-lemma-chi-tensor-finite}.
\end{proof}

\begin{lemma}
\label{varieties-lemma-degree-on-proper-curve}
Soit $k$ un corps. Soit $X$ une courbe propre sur $k$ de point générique
$\xi$. Soit $\mathcal{E}$ un $\mathcal{O}_X$-module localement libre de rang $n$
et soit $\mathcal{F}$ un $\mathcal{O}_X$-module cohérent. Alors
$$
\chi(X, \mathcal{E} \otimes \mathcal{F}) =
r \deg(\mathcal{E}) + n \chi(X, \mathcal{F})
$$
où $r = \dim_{\kappa(\xi)} \mathcal{F}_\xi$ est le rang de $\mathcal{F}$.
\end{lemma}

\begin{proof}
Soit $\mathcal{P}$ la propriété d'un faisceau cohérent $\mathcal{F}$
sur $X$ selon laquelle la formule du lemme est satisfaite. Nous affirmons que
les hypothèses (1) et (2) de
Cohomologie des schémas, lemme \ref{coherent-lemma-property},
sont satisfaites pour $\mathcal{P}$. En effet, (1) l'est parce que la caractéristique d'Euler-Poincaré
et le rang $r$ sont additifs dans les suites exactes courtes de faisceaux cohérents.
La condition (2) l'est également : si $Z = X$, nous pouvons prendre
$\mathcal{G} = \mathcal{O}_X$, et $\mathcal{P}(\mathcal{O}_X)$
est vraie par définition du degré. Si $i : Z \to X$ est l'inclusion
d'un point fermé, nous pouvons prendre $\mathcal{G} = i_*\mathcal{O}_Z$,
et $\mathcal{P}$ est satisfaite par le lemme \ref{varieties-lemma-chi-tensor-finite}
et par le fait que $r = 0$ dans ce cas.
\end{proof}

\noindent
Soit $k$ un corps. Soit $X$ un schéma de type fini sur $k$ de
dimension $\leq 1$. Soient $C_i \subset X$, $i = 1, \ldots, t$, les
composantes irréductibles de dimension $1$. Nous munissons $C_i$ de la
structure de schéma réduit induite. Soit $\xi_i \in C_i$ le
point générique. La {\it multiplicité de $C_i$ dans $X$} est définie
comme la longueur
$$
m_i = \text{longueur}_{\mathcal{O}_{X, \xi_i}} \mathcal{O}_{X, \xi_i}
$$
Cette définition a un sens parce que $\mathcal{O}_{X, \xi_i}$ est un anneau local
noethérien de dimension zéro, donc est de longueur finie sur lui-même
(Algèbre, proposition \ref{algebra-proposition-dimension-zero-ring}).
Voir Homologie de Chow, section \ref{chow-section-cycle-of-closed-subscheme},
pour de plus amples informations. En fait, le degré d'un faisceau localement libre
ne dépend que de ses restrictions aux composantes irréductibles.

\begin{lemma}
\label{varieties-lemma-degree-in-terms-of-components}
Soit $k$ un corps. Soit $X$ un schéma propre de dimension $\leq 1$ sur $k$.
Soit $\mathcal{E}$ un $\mathcal{O}_X$-module localement libre de rang $n$.
Alors
$$
\deg(\mathcal{E}) = \sum m_i \deg(\mathcal{E}|_{C_i})
$$
où $C_i \subset X$, $i = 1, \ldots, t$, sont les composantes irréductibles
de dimension $1$, munies de la structure de schéma réduit induite, et $m_i$ est la
multiplicité de $C_i$ dans $X$.
\end{lemma}

\begin{proof}
Remarquons que l'énoncé a un sens parce que $C_i \to \Spec(k)$
est propre de dimension $1$ ; nous utiliserons ci-dessous le fait que tout sous-schéma fermé
de $X$ est propre sur $k$ par Morphismes, lemmes
\ref{morphisms-lemma-closed-immersion-proper} et
\ref{morphisms-lemma-composition-proper}. Considérons le sous-schéma ouvert
$U_i = X \setminus (\bigcup_{j \not = i} C_j)$ et soit $X_i \subset X$
l'adhérence schématique de $U_i$. Remarquons que $X_i \cap U_i = U_i$
(schématiquement) et que $X_i \cap U_j = \emptyset$
(ensemblistement) pour $i \not = j$ ; cela résulte de la description
de l'adhérence schématique dans
Morphismes, lemme \ref{morphisms-lemma-quasi-compact-immersion}.
Nous pouvons donc appliquer le lemme \ref{varieties-lemma-degree-birational-pullback} au morphisme
$X' = \coprod X_i \to X$. Comme il est clair que $C_i \subset X_i$
(schématiquement) et que la multiplicité de $C_i$ dans $X_i$ est égale
à la multiplicité de $C_i$ dans $X$, nous voyons que nous sommes ramenés au cas
étudié dans le paragraphe suivant.

\medskip\noindent
Supposons $X$ irréductible, de point générique $\xi$. Soit
$C = X_{red}$, de multiplicité $m$.
Nous devons montrer que $\deg(\mathcal{E}) = m \deg(\mathcal{E}|_C)$.
Soit $\mathcal{I} \subset \mathcal{O}_X$ l'idéal définissant le sous-schéma fermé
$C$. Soit $e \geq 0$ minimal tel que $\mathcal{I}^{e + 1} = 0$
(Cohomologie des schémas, lemme \ref{coherent-lemma-power-ideal-kills-sheaf}).
Nous raisonnons par récurrence sur $e$. Si $e = 0$, alors $X = C$ et le résultat
est immédiat. Sinon, posons $\mathcal{F} = \mathcal{I}^e$, considéré comme
un $\mathcal{O}_C$-module cohérent (Cohomologie des schémas, lemme
\ref{coherent-lemma-i-star-equivalence}).
Soit $X' \subset X$ le sous-schéma fermé défini par
l'idéal cohérent $\mathcal{I}^e$, et soit $m'$ la multiplicité
de $C$ dans $X'$. En prenant les germes en $\xi$ de la suite exacte courte
$$
0 \to \mathcal{F} \to \mathcal{O}_X \to \mathcal{O}_{X'} \to 0
$$
nous trouvons (en utilisant Algèbre, lemmes \ref{algebra-lemma-length-additive},
\ref{algebra-lemma-dimension-is-length} et
\ref{algebra-lemma-length-independent}) que
$$
m = \text{longueur}_{\mathcal{O}_{X, \xi}} \mathcal{O}_{X, \xi}
= \dim_{\kappa(\xi)} \mathcal{F}_\xi +
\text{longueur}_{\mathcal{O}_{X', \xi}} \mathcal{O}_{X', \xi}
= r + m'
$$
où $r$ est le rang de $\mathcal{F}$ en tant que faisceau cohérent sur $C$.
En tensorisant par $\mathcal{E}$, nous obtenons une suite exacte courte
$$
0 \to \mathcal{E}|_C \otimes \mathcal{F} \to \mathcal{E} \to
\mathcal{E} \otimes \mathcal{O}_{X'} \to 0
$$
Par récurrence, nous avons
$\deg(\mathcal{E}|_{X'}) = m' \deg(\mathcal{E}|_C)$.
Par le lemme \ref{varieties-lemma-degree-on-proper-curve}, nous avons
$\chi(\mathcal{E}|_C \otimes \mathcal{F}) =
r \deg(\mathcal{E}|_C) + n \chi(\mathcal{F})$.
En réunissant ces égalités, nous obtenons le résultat.
\end{proof}

\begin{lemma}
\label{varieties-lemma-degree-tensor-product}
Soit $k$ un corps, soit $X$ un schéma propre de dimension $\leq 1$
sur $k$, et soient $\mathcal{E}$, $\mathcal{V}$ des
$\mathcal{O}_X$-modules localement libres de type fini et de rang constant. Alors
$$
\deg(\mathcal{E} \otimes \mathcal{V}) =
\text{rang}(\mathcal{E}) \deg(\mathcal{V}) +
\text{rang}(\mathcal{V}) \deg(\mathcal{E})
$$
\end{lemma}

\begin{proof}
Par le lemme \ref{varieties-lemma-degree-in-terms-of-components} et un calcul
élémentaire, nous nous ramenons au cas d'une courbe propre.
Ce cas résulte du lemme \ref{varieties-lemma-degree-on-proper-curve}.
\end{proof}

\begin{lemma}
\label{varieties-lemma-degree-and-det}
Soit $k$ un corps, soit $X$ un schéma propre de dimension $\leq 1$
sur $k$, et soit $\mathcal{E}$ un
$\mathcal{O}_X$-module localement libre de rang $n$. Alors
$$
\deg(\mathcal{E}) = \deg(\wedge^n(\mathcal{E})) = \deg(\det(\mathcal{E}))
$$
\end{lemma}

\begin{proof}
Par le lemme \ref{varieties-lemma-degree-in-terms-of-components} et un calcul
élémentaire, nous nous ramenons au cas d'une courbe propre.
Il existe alors une modification $f : X' \to X$ telle que
$f^*\mathcal{E}$ admette une filtration dont les quotients
successifs sont des modules inversibles, voir
Diviseurs, lemme \ref{divisors-lemma-filter-after-modification}.
Par le lemme \ref{varieties-lemma-degree-birational-pullback}, nous pouvons travailler sur $X'$.
Nous pouvons donc supposer que nous avons une filtration
$$
0 = \mathcal{E}_0 \subset \mathcal{E}_1 \subset \mathcal{E}_2 \subset
\ldots \subset \mathcal{E}_n = \mathcal{E}
$$
par des $\mathcal{O}_X$-modules localement libres
telle que $\mathcal{L}_i = \mathcal{E}_i/\mathcal{E}_{i - 1}$ soit inversible.
Par Modules, lemme \ref{modules-lemma-det-ses}, et par récurrence, nous trouvons
$\det(\mathcal{E}) = \mathcal{L}_1 \otimes \ldots \otimes \mathcal{L}_n$.
L'égalité résulte donc du lemme \ref{varieties-lemma-degree-tensor-product}
et de l'additivité (lemme \ref{varieties-lemma-degree-additive}).
\end{proof}

\begin{lemma}
\label{varieties-lemma-degree-effective-Cartier-divisor}
Soit $k$ un corps, soit $X$ un schéma propre de dimension $\leq 1$
sur $k$. Soit $D$ un diviseur de Cartier effectif sur $X$.
Alors $D$ est fini sur $\Spec(k)$, de degré
$\deg(D) = \dim_k \Gamma(D, \mathcal{O}_D)$. Pour un faisceau localement libre
$\mathcal{E}$ de rang $n$, nous avons
$$
\deg(\mathcal{E}(D)) = n\deg(D) + \deg(\mathcal{E})
$$
où $\mathcal{E}(D) = \mathcal{E} \otimes_{\mathcal{O}_X} \mathcal{O}_X(D)$.
\end{lemma}

\begin{proof}
Comme $D$ est nulle part dense dans $X$ (Diviseurs, lemme
\ref{divisors-lemma-complement-effective-Cartier-divisor}),
nous voyons que $\dim(D) \leq 0$. Par conséquent, $D$ est fini sur $k$
par le lemme \ref{varieties-lemma-algebraic-scheme-dim-0}. Comme $k$ est un corps,
le morphisme $D \to \Spec(k)$ est fini localement libre et possède donc
un degré
(Morphismes, définition \ref{morphisms-definition-finite-locally-free}),
manifestement égal à $\dim_k \Gamma(D, \mathcal{O}_D)$,
comme l'affirme le lemme. Par
Diviseurs, définition
\ref{divisors-definition-invertible-sheaf-effective-Cartier-divisor},
il existe une suite exacte courte
$$
0 \to \mathcal{O}_X \to \mathcal{O}_X(D) \to i_*i^*\mathcal{O}_X(D) \to 0
$$
où $i : D \to X$ est l'immersion fermée. En tensorisant par $\mathcal{E}$,
nous obtenons une suite exacte courte
$$
0 \to \mathcal{E} \to \mathcal{E}(D) \to i_*i^*\mathcal{E}(D) \to 0
$$
L'égalité du lemme résulte de l'additivité de
la caractéristique d'Euler-Poincaré (lemme \ref{varieties-lemma-euler-characteristic-additive})
et du lemme \ref{varieties-lemma-chi-tensor-finite}.
\end{proof}

\begin{lemma}
\label{varieties-lemma-divisible}
Soit $k$ un corps. Soit $X$ un schéma propre sur $k$,
réduit et connexe.
Soit $\kappa = H^0(X, \mathcal{O}_X)$. Alors $\kappa/k$ est une
extension finie de corps et $w = [\kappa : k]$ divise
\begin{enumerate}
\item $\deg(\mathcal{E})$ pour tout $\mathcal{O}_X$-module localement libre
$\mathcal{E}$,
\item $[\kappa(x) : k]$ pour tout point fermé $x \in X$, et
\item $\deg(D)$ pour tout sous-schéma fermé $D \subset X$
de dimension zéro.
\end{enumerate}
\end{lemma}

\begin{proof}
Voir le lemme \ref{varieties-lemma-proper-geometrically-reduced-global-sections}
pour les assertions concernant $\kappa$.
Pour tout $\mathcal{O}_X$-module quasi-cohérent, les
$k$-espaces vectoriels $H^i(X, \mathcal{F})$ sont des $\kappa$-espaces vectoriels.
Les divisibilités résultent immédiatement de cette assertion et des
définitions.
\end{proof}

\begin{lemma}
\label{varieties-lemma-degree-pullback-map-proper-curves}
Soit $k$ un corps. Soit $f : X \to Y$ un morphisme non constant de
courbes propres sur $k$. Soit $\mathcal{E}$ un
$\mathcal{O}_Y$-module localement libre. Alors
$$
\deg(f^*\mathcal{E}) = \deg(X/Y) \deg(\mathcal{E})
$$
\end{lemma}

\begin{proof}
Le degré de $X$ sur $Y$ est défini dans
Morphismes, définition \ref{morphisms-definition-degree}.
Ainsi, $f_*\mathcal{O}_X$ est un $\mathcal{O}_Y$-module cohérent
de rang $\deg(X/Y)$, c'est-à-dire
$\deg(X/Y) = \dim_{\kappa(\xi)} (f_*\mathcal{O}_X)_\xi$, où $\xi$
est le point générique de $Y$. Nous obtenons donc
\begin{align*}
\chi(X, f^*\mathcal{E})
& =
\chi(Y, f_*f^*\mathcal{E}) \\
& =
\chi(Y, \mathcal{E} \otimes f_*\mathcal{O}_X) \\
& =
\deg(X/Y) \deg(\mathcal{E}) + n \chi(Y, f_*\mathcal{O}_X) \\
& =
\deg(X/Y) \deg(\mathcal{E}) + n \chi(X, \mathcal{O}_X)
\end{align*}
comme voulu. La première égalité vaut parce que $f$ est fini, voir
Cohomologie des schémas, lemme \ref{coherent-lemma-relative-affine-cohomology}.
La deuxième égalité résulte de la formule de projection, voir
Cohomologie, lemme \ref{cohomology-lemma-projection-formula}.
La troisième égalité résulte du lemme \ref{varieties-lemma-degree-on-proper-curve}.
\end{proof}

\noindent
Ce qui suit est une conséquence évidente, mais importante, des
résultats ci-dessus sur les degrés.

\begin{lemma}
\label{varieties-lemma-check-invertible-sheaf-trivial}
Soit $k$ un corps. Soit $X$ une courbe propre sur $k$.
Soit $\mathcal{L}$ un $\mathcal{O}_X$-module inversible.
\begin{enumerate}
\item Si $\mathcal{L}$ possède une section non nulle, alors
$\deg(\mathcal{L}) \geq 0$.
\item Si $\mathcal{L}$ possède une section non nulle $s$ qui s'annule
en un point, alors $\deg(\mathcal{L}) > 0$.
\item Si $\mathcal{L}$ et $\mathcal{L}^{-1}$ possèdent des sections non nulles, alors
$\mathcal{L} \cong \mathcal{O}_X$.
\item Si $\deg(\mathcal{L}) \leq 0$ et si $\mathcal{L}$ possède une section non nulle,
alors $\mathcal{L} \cong \mathcal{O}_X$.
\item Si $\mathcal{N} \to \mathcal{L}$ est un morphisme non nul de
$\mathcal{O}_X$-modules inversibles, alors $\deg(\mathcal{L}) \geq \deg(\mathcal{N})$,
et, s'il y a égalité, ce morphisme est un isomorphisme.
\end{enumerate}
\end{lemma}

\begin{proof}
Soit $s$ une section non nulle de $\mathcal{L}$. Comme $X$ est une courbe, nous
voyons que $s$ est une section régulière. Il existe donc un diviseur de
Cartier effectif $D \subset X$ et un isomorphisme
$\mathcal{L} \to \mathcal{O}_X(D)$ envoyant $s$ sur la section canonique
$1$ de $\mathcal{O}_X(D)$, voir
Diviseurs, lemme \ref{divisors-lemma-characterize-OD}.
Alors $\deg(\mathcal{L}) = \deg(D)$ par
le lemme \ref{varieties-lemma-degree-effective-Cartier-divisor}.
Comme $\deg(D) \geq 0$ et que l'égalité $= 0$ a lieu
si et seulement si $D = \emptyset$, cela prouve (1) et (2).
Dans le cas (3), nous voyons que $\deg(\mathcal{L}) = 0$ et que
$D = \emptyset$. De même pour (4). Pour prouver (5), appliquons
(1) et (4) au faisceau inversible
$$
\mathcal{L} \otimes_{\mathcal{O}_X} \mathcal{N}^{\otimes -1} =
\SheafHom_{\mathcal{O}_X}(\mathcal{N}, \mathcal{L})
$$
dont le degré est $\deg(\mathcal{L}) - \deg(\mathcal{N})$
par le lemme \ref{varieties-lemma-degree-tensor-product}.
\end{proof}

\begin{lemma}
\label{varieties-lemma-no-sections-dual-nef}
Soit $k$ un corps. Soit $X$ un schéma propre sur $k$,
réduit, connexe et équidimensionnel de dimension $1$.
Soit $\mathcal{L}$ un $\mathcal{O}_X$-module inversible.
Si $\deg(\mathcal{L}|_C) \leq 0$ pour toute composante irréductible $C$
de $X$, alors ou bien $H^0(X, \mathcal{L}) = 0$, ou bien
$\mathcal{L} \cong \mathcal{O}_X$.
\end{lemma}

\begin{proof}
Soit $s \in H^0(X, \mathcal{L})$ non nul. Comme $X$ est réduit, il existe
une composante irréductible $C$ de $X$ telle que $s|_C \not = 0$.
Mais si $s|_C$ est non nul, alors
$s$ ne s'annule nulle part sur $C$ par
le lemme \ref{varieties-lemma-check-invertible-sheaf-trivial}.
Cela implique à son tour que $s$ ne s'annule nulle part sur
toute composante irréductible de $X$ rencontrant $C$.
Comme $X$ est connexe, nous en concluons que $s$
ne s'annule nulle part, et le lemme en résulte.
\end{proof}

\begin{lemma}
\label{varieties-lemma-ample-curve}
Soit $k$ un corps. Soit $X$ une courbe propre sur $k$.
Soit $\mathcal{L}$ un $\mathcal{O}_X$-module inversible.
Alors $\mathcal{L}$ est ample si et seulement si $\deg(\mathcal{L}) > 0$.
\end{lemma}

\begin{proof}
Si $\mathcal{L}$ est ample, il existe un $n > 0$ et une section
$s \in H^0(X, \mathcal{L}^{\otimes n})$ tels que $X_s$ soit affine. Comme
$X$ n'est pas affine (sinon, d'après
Morphismes, lemme \ref{morphisms-lemma-finite-proper},
$X$ serait fini), on voit que $s$ s'annule en un point.
Ainsi, $\deg(\mathcal{L}^{\otimes n}) > 0$ d'après le
lemme \ref{varieties-lemma-check-invertible-sheaf-trivial}.
D'après le lemme \ref{varieties-lemma-degree-tensor-product},
on conclut que $\deg(\mathcal{L}) = 1/n\deg(\mathcal{L}^{\otimes n}) > 0$.

\medskip\noindent
Supposons que $\deg(\mathcal{L}) > 0$. Alors
$$
\dim_k H^0(X, \mathcal{L}^{\otimes n}) \geq \chi(X, \mathcal{L}^n)
= n\deg(\mathcal{L}) + \chi(X, \mathcal{O}_X)
$$
croît linéairement en $n$. Pour toute famille finie de points
fermés $x_1, \ldots, x_t$ de $X$, on peut donc trouver un $n$ tel que
$\dim_k H^0(X, \mathcal{L}^{\otimes n}) > \sum \dim_k \kappa(x_i)$.
(Rappelons que, d'après le théorème des zéros de Hilbert, les extensions
$\kappa(x_i)/k$ sont finies, voir par exemple
Morphismes, lemme \ref{morphisms-lemma-closed-point-fibre-locally-finite-type}.)
On peut donc trouver une section non nulle $s \in H^0(X, \mathcal{L}^{\otimes n})$
qui s'annule en $x_1, \ldots, x_t$. En particulier, si l'on choisit
$x_1, \ldots, x_t$ de telle sorte que $X \setminus \{x_1, \ldots, x_t\}$
soit affine, alors $X_s$ est également affine
(par exemple d'après Propriétés, lemme \ref{properties-lemma-affine-cap-s-open},
bien que, si l'on choisit l'ensemble fini de telle sorte que
$\mathcal{L}|_{X \setminus \{x_1, \ldots, x_t\}}$ soit trivial,
ce soit immédiat). On peut donc trouver un $n > 0$
et une section non nulle $s \in H^0(X, \mathcal{L}^{\otimes n})$ telle
que $X_s$ soit affine.

\medskip\noindent
Nous allons montrer que, pour tout faisceau quasi-cohérent d'idéaux $\mathcal{I}$,
il existe un $m > 0$ tel que
$H^1(X, \mathcal{I} \otimes \mathcal{L}^{\otimes m})$ soit nul.
Cela achèvera la démonstration d'après
Cohomologie des schémas, lemme \ref{coherent-lemma-vanshing-gives-ample}.
Pour cela, considérons les applications
$$
\mathcal{I} \xrightarrow{s}
\mathcal{I} \otimes \mathcal{L}^{\otimes n} \xrightarrow{s}
\mathcal{I} \otimes \mathcal{L}^{\otimes 2n} \xrightarrow{s} \ldots
$$
Puisque $\mathcal{I}$ est sans torsion, ces applications sont injectives et
sont des isomorphismes sur $X_s$ ; leurs conoyaux ont donc un $H^1$ nul
(par exemple d'après Cohomologie des schémas, lemme
\ref{coherent-lemma-coherent-support-dimension-0}).
On en conclut que les applications d'espaces vectoriels
$$
H^1(X, \mathcal{I}) \to
H^1(X, \mathcal{I} \otimes \mathcal{L}^{\otimes n}) \to
H^1(X, \mathcal{I} \otimes \mathcal{L}^{\otimes 2n}) \to \ldots
$$
sont surjectives. D'autre part, $H^1(X, \mathcal{I})$ est de dimension
finie, et tout élément a une image nulle à partir d'un certain rang d'après
Cohomologie des schémas, lemme
\ref{coherent-lemma-section-affine-open-kills-classes}.
Il existe donc un $e > 0$ tel que
$H^1(X, \mathcal{I} \otimes \mathcal{L}^{\otimes en})$ soit nul.
Cela achève la démonstration.
\end{proof}

\begin{lemma}
\label{varieties-lemma-ampleness-in-terms-of-degrees-components}
Soit $k$ un corps. Soit $X$ un schéma propre de dimension $\leq 1$
sur $k$. Soit $\mathcal{L}$ un $\mathcal{O}_X$-module inversible.
Soient $C_i \subset X$, $i = 1, \ldots, t$, les composantes irréductibles
de dimension $1$. Les conditions suivantes sont équivalentes :
\begin{enumerate}
\item $\mathcal{L}$ est ample, et
\item $\deg(\mathcal{L}|_{C_i}) > 0$ pour $i = 1, \ldots, t$.
\end{enumerate}
\end{lemma}

\begin{proof}
Soient $x_1, \ldots, x_r \in X$ les points fermés isolés.
Considérons $x_i = \Spec(\kappa(x_i))$ comme un schéma.
Considérons le morphisme de schémas
$$
f : C_1 \amalg \ldots \amalg C_t \amalg x_1 \amalg \ldots \amalg x_r
\longrightarrow X
$$
C'est un morphisme fini surjectif entre des schémas propres sur $k$
(détails omis). Ainsi, $\mathcal{L}$ est ample si et seulement si
$f^*\mathcal{L}$ est ample (Cohomologie des schémas, lemme
\ref{coherent-lemma-surjective-finite-morphism-ample}).
On conclut donc par le lemme \ref{varieties-lemma-ample-curve}.
\end{proof}

\begin{lemma}
\label{varieties-lemma-general-degree-g-line-bundle}
Soit $k$ un corps algébriquement clos. Soit $X$ une courbe propre sur $k$.
Il existe alors
\begin{enumerate}
\item un $\mathcal{O}_X$-module inversible $\mathcal{L}$ tel que
$\dim_k H^0(X, \mathcal{L}) = 1$ et $H^1(X, \mathcal{L}) = 0$, et
\item un $\mathcal{O}_X$-module inversible $\mathcal{N}$ tel que
$\dim_k H^0(X, \mathcal{N}) = 0$ et $H^1(X, \mathcal{N}) = 0$.
\end{enumerate}
\end{lemma}

\begin{proof}
Choisissons une immersion fermée $i : X \to \mathbf{P}^n_k$
(lemme \ref{varieties-lemma-dim-1-proper-projective}).
En posant $\mathcal{L} = i^*\mathcal{O}_{\mathbf{P}^n}(d)$
pour $d \gg 0$, on voit qu'il existe un faisceau inversible
$\mathcal{L}$ tel que $H^0(X, \mathcal{L}) \not = 0$ et
$H^1(X, \mathcal{L}) = 0$ (voir
Cohomologie des schémas, lemme \ref{coherent-lemma-vanshing-gives-ample},
pour l'annulation, et les références qui y sont données pour la non-annulation).
Nous allons achever la démonstration de (1) par récurrence descendante sur
$t = \dim_k H^0(X, \mathcal{L})$. Le cas initial $t = 1$ est immédiat.
Supposons que $t > 1$.

\medskip\noindent
Soit $U \subset X$ l'ouvert non vide des points non singuliers
étudié dans le lemme \ref{varieties-lemma-dense-smooth-open-variety-over-perfect-field}.
Soit $s \in H^0(X, \mathcal{L})$ non nul. Il existe un point
fermé $x \in U$ tel que $s$ ne s'annule pas en $x$. Soit $\mathcal{I}$
le faisceau d'idéaux de $i : x \to X$ comme dans le
lemme \ref{varieties-lemma-regular-point-on-curve}. Considérons la
suite exacte courte
$$
0 \to \mathcal{I} \otimes_{\mathcal{O}_X} \mathcal{L} \to
\mathcal{L} \to i_*i^*\mathcal{L} \to 0
$$
Remarquons que $H^0(X, i_*i^*\mathcal{L}) = H^0(x, i^*\mathcal{L})$
est de dimension $1$, puisque $x$ est un point $k$-rationnel ($k$ est
algébriquement clos). Comme $s$ ne s'annule pas en $x$, on en conclut que
$$
H^0(X, \mathcal{L}) \longrightarrow H^0(X, i_*i^*\mathcal{L})
$$
est surjective. Par conséquent,
$\dim_k H^0(X, \mathcal{I} \otimes_{\mathcal{O}_X} \mathcal{L}) = t - 1$.
Enfin, la suite exacte longue de cohomologie montre également que
$H^1(X, \mathcal{I} \otimes_{\mathcal{O}_X} \mathcal{L}) = 0$,
ce qui achève la démonstration de l'étape de récurrence.

\medskip\noindent
Pour obtenir un faisceau inversible comme en (2), prenons un faisceau inversible
$\mathcal{L}$ comme en (1) et appliquons l'argument du paragraphe précédent
une fois encore.
\end{proof}

\begin{lemma}
\label{varieties-lemma-vanishing-degree-2g-and-1-line-bundle}
Soit $k$ un corps algébriquement clos. Soit $X$ une courbe propre sur $k$.
Posons $g = \dim_k H^1(X, \mathcal{O}_X)$. Pour tout
$\mathcal{O}_X$-module inversible $\mathcal{L}$ tel que $\deg(\mathcal{L}) \geq 2g - 1$,
on a $H^1(X, \mathcal{L}) = 0$.
\end{lemma}

\begin{proof}
Soit $\mathcal{N}$ le module inversible obtenu dans
le lemme \ref{varieties-lemma-general-degree-g-line-bundle}, partie (2).
Le degré de $\mathcal{N}$ est
$\chi(X, \mathcal{N}) - \chi(X, \mathcal{O}_X) = 0 - (1 - g) = g - 1$.
Le degré de $\mathcal{L} \otimes \mathcal{N}^{\otimes - 1}$
est donc $\deg(\mathcal{L}) - (g - 1) \geq g$.
Par conséquent,
$\chi(X, \mathcal{L} \otimes \mathcal{N}^{\otimes -1}) \geq g + 1 - g = 1$.
Il existe ainsi une section globale non nulle $s$ dont le schéma des zéros est un
diviseur de Cartier effectif $D$ de degré $\deg(\mathcal{L}) - (g - 1)$.
On obtient une suite exacte courte
$$
0 \to \mathcal{N} \xrightarrow{s} \mathcal{L} \to i_*(\mathcal{L}|_D) \to 0
$$
où $i : D \to X$ est le morphisme d'inclusion. On en conclut que
$H^0(X, \mathcal{L})$ s'applique isomorphiquement sur $H^0(D, \mathcal{L}|_D)$,
qui est de dimension $\deg(\mathcal{L}) - (g - 1)$. Le résultat découle
de la définition du degré.
\end{proof}






\section{Intersections numériques}
\label{varieties-section-num}

\noindent
Dans cette section, nous allons exploiter la caractéristique d'Euler-Poincaré des
faisceaux cohérents sur les schémas propres afin d'obtenir des nombres d'intersection
numériques pour les modules inversibles. Notre principal outil sera le
lemme suivant.

\begin{lemma}
\label{varieties-lemma-numerical-polynomial-from-euler}
Soit $k$ un corps. Soit $X$ un schéma propre sur $k$. Soit $\mathcal{F}$
un $\mathcal{O}_X$-module cohérent. Soient
$\mathcal{L}_1, \ldots, \mathcal{L}_r$ des $\mathcal{O}_X$-modules inversibles.
L'application
$$
(n_1, \ldots, n_r) \longmapsto
\chi(X, \mathcal{F} \otimes
\mathcal{L}_1^{\otimes n_1} \otimes \ldots \otimes
\mathcal{L}_r^{\otimes n_r})
$$
est un polynôme numérique en $n_1, \ldots, n_r$, dont le degré total est au
plus la dimension du support de $\mathcal{F}$.
\end{lemma}

\begin{proof}
Nous procédons par récurrence sur $\dim(\text{Supp}(\mathcal{F}))$.
Si ce nombre est nul, la fonction est constante de valeur
$\dim_k \Gamma(X, \mathcal{F})$ d'après le lemme \ref{varieties-lemma-chi-tensor-finite}.
Supposons que $\dim(\text{Supp}(\mathcal{F})) > 0$.

\medskip\noindent
Si $\mathcal{F}$ a des points associés immergés, on peut considérer
la suite exacte courte
$0 \to \mathcal{K} \to \mathcal{F} \to \mathcal{F}' \to 0$
construite dans Diviseurs, lemme \ref{divisors-lemma-remove-embedded-points}.
Comme la dimension du support de $\mathcal{K}$ est strictement plus petite,
le résultat vaut pour $\mathcal{K}$ par l'hypothèse de récurrence, avec
un degré total strictement plus petit.
Par additivité de la caractéristique d'Euler-Poincaré
(lemme \ref{varieties-lemma-euler-characteristic-additive}),
il suffit d'établir le résultat pour $\mathcal{F}'$. On peut donc supposer que
$\mathcal{F}$ n'a pas de points associés immergés.

\medskip\noindent
Si $i : Z \to X$ est une immersion fermée et $\mathcal{F} = i_*\mathcal{G}$,
le résultat pour $X$, $\mathcal{F}$,
$\mathcal{L}_1, \ldots, \mathcal{L}_r$ équivaut au résultat
pour $Z$, $\mathcal{G}$, $i^*\mathcal{L}_1, \ldots, i^*\mathcal{L}_r$
(car les groupes de cohomologie coïncident, voir
Cohomologie des schémas, lemme \ref{coherent-lemma-relative-affine-cohomology}).
En appliquant Diviseurs, lemme \ref{divisors-lemma-no-embedded-points-endos},
on peut supposer que $X$ n'a pas de composantes immergées et que
$X = \text{Supp}(\mathcal{F})$.

\medskip\noindent
Choisissons une section méromorphe régulière $s$ de $\mathcal{L}_1$, voir
Diviseurs, lemme \ref{divisors-lemma-regular-meromorphic-section-exists}.
Soit $\mathcal{I} \subset \mathcal{O}_X$ l'idéal des
dénominateurs de $s$, et considérons les applications
$$
\mathcal{I}\mathcal{F} \to \mathcal{F},\quad
\mathcal{I}\mathcal{F} \to \mathcal{F} \otimes \mathcal{L}_1
$$
de Diviseurs, lemme \ref{divisors-lemma-make-maps-regular-section}.
Elles sont injectives et ont pour conoyaux $\mathcal{Q}$, $\mathcal{Q}'$,
à support dans des sous-schémas fermés rares de $X = \text{Supp}(\mathcal{F})$.
Le produit tensoriel par le module inversible
$\mathcal{L}_1^{\otimes n_1} \otimes \ldots \otimes \mathcal{L}_r^{\otimes n_r}$
est exact ; par conséquent, en utilisant de nouveau l'additivité,
on voit que
\begin{align*}
&\chi(X, \mathcal{F} \otimes \mathcal{L}_1^{\otimes n_1} \otimes \ldots \otimes
\mathcal{L}_r^{\otimes n_r}) -
\chi(X, \mathcal{F} \otimes \mathcal{L}_1^{\otimes n_1 + 1}
\otimes \ldots \otimes \mathcal{L}_r^{\otimes n_r}) \\
& =
\chi(\mathcal{Q} \otimes \mathcal{L}_1^{\otimes n_1} \otimes \ldots \otimes
\mathcal{L}_r^{\otimes n_r}) -
\chi(\mathcal{Q}' \otimes \mathcal{L}_1^{\otimes n_1} \otimes \ldots \otimes
\mathcal{L}_r^{\otimes n_r})
\end{align*}
Ainsi, la fonction $P(n_1, \ldots, n_r)$ du lemme
a la propriété que
$$
P(n_1 + 1, n_2, \ldots, n_r) - P(n_1, \ldots, n_r)
$$
est un polynôme numérique dont le degré total est $<$ la dimension
du support de $\mathcal{F}$. Bien entendu, par symétrie, il en va
de même de
$$
P(n_1, \ldots, n_{i - 1}, n_i + 1, n_{i + 1}, \ldots, n_r)
- P(n_1, \ldots, n_r)
$$
pour tout $i \in \{1, \ldots, r\}$. Un argument arithmétique simple montre
que $P$ est un polynôme numérique de degré total au plus
$\dim(\text{Supp}(\mathcal{F}))$.
\end{proof}

\noindent
Le lemme suivant montre, en gros, que le coefficient dominant ne dépend que
de la longueur du module cohérent aux points génériques de son
support.

\begin{lemma}
\label{varieties-lemma-numerical-polynomial-leading-term}
Soit $k$ un corps. Soit $X$ un schéma propre sur $k$. Soit
$\mathcal{F}$ un $\mathcal{O}_X$-module cohérent. Soient
$\mathcal{L}_1, \ldots, \mathcal{L}_r$ des $\mathcal{O}_X$-modules inversibles.
Posons $d = \dim(\text{Supp}(\mathcal{F}))$.
Soient $Z_i \subset X$ les composantes irréductibles
de $\text{Supp}(\mathcal{F})$ de dimension $d$. Soit $\xi_i \in Z_i$
le point générique, et posons
$m_i = \text{longueur}_{\mathcal{O}_{X, \xi_i}}(\mathcal{F}_{\xi_i})$.
Alors
$$
\chi(X, \mathcal{F} \otimes \mathcal{L}_1^{\otimes n_1} \otimes \ldots \otimes
\mathcal{L}_r^{\otimes n_r}) -
\sum\nolimits_i
m_i\ \chi(Z_i, \mathcal{L}_1^{\otimes n_1} \otimes \ldots \otimes
\mathcal{L}_r^{\otimes n_r}|_{Z_i})
$$
est un polynôme numérique en $n_1, \ldots, n_r$ de degré total $< d$.
\end{lemma}

\begin{proof}
Considérons les couples $(\xi , Z)$ où $Z \subset X$ est un
sous-schéma fermé intègre de dimension $d$ et où $\xi$ est son point générique.
Alors le $\mathcal{O}_{X, \xi}$-module de type fini $\mathcal{F}_\xi$
a son support contenu dans $\{\xi\}$ ; sa longueur
$m_Z = \text{longueur}_{\mathcal{O}_{X, \xi}}(\mathcal{F}_\xi)$
est donc finie (Algèbre, lemme \ref{algebra-lemma-support-point})
et est nulle sauf si $Z = Z_i$ pour un certain $i$. L'expression
du lemme peut donc s'écrire
$$
E(\mathcal{F}) =
\chi(X, \mathcal{F} \otimes \mathcal{L}_1^{\otimes n_1} \otimes \ldots \otimes
\mathcal{L}_r^{\otimes n_r}) -
\sum\nolimits
m_Z\ \chi(Z, \mathcal{L}_1^{\otimes n_1} \otimes \ldots \otimes
\mathcal{L}_r^{\otimes n_r}|_Z)
$$
où la somme porte sur les sous-schémas fermés intègres $Z \subset X$
de dimension $d$. L'application $\mathcal{F} \mapsto E(\mathcal{F})$
est additive pour les suites exactes courtes
$0 \to \mathcal{F} \to \mathcal{F}' \to \mathcal{F}'' \to 0$
de $\mathcal{O}_X$-modules cohérents dont le support est de dimension
$\leq d$. Cela résulte de l'additivité des caractéristiques d'Euler-Poincaré
(lemme \ref{varieties-lemma-euler-characteristic-additive})
et de l'additivité des longueurs
(Algèbre, lemme \ref{algebra-lemma-length-additive}).
Appliquons Cohomologie des schémas, lemme \ref{coherent-lemma-coherent-filter},
pour obtenir une filtration
$$
0 = \mathcal{F}_0 \subset \mathcal{F}_1 \subset
\ldots \subset \mathcal{F}_m = \mathcal{F}
$$
par des sous-faisceaux cohérents tels que, pour chaque $j = 1, \ldots, m$,
il existe un sous-schéma fermé intègre $V_j \subset X$
et un faisceau d'idéaux non nul $\mathcal{I}_j \subset \mathcal{O}_{V_j}$
tels que
$$
\mathcal{F}_j/\mathcal{F}_{j - 1}
\cong (V_j \to X)_* \mathcal{I}_j
$$
Il en résulte que $V_j \subset \text{Supp}(\mathcal{F})$ et
donc que $\dim(V_j) \leq d$.
Par l'additivité relevée ci-dessus, il suffit
d'établir le résultat pour chacun des sous-quotients
$\mathcal{F}_j/\mathcal{F}_{j - 1}$. Il suffit donc d'établir
le résultat lorsque $\mathcal{F} = (V \to X)_*\mathcal{I}$, où
$V \subset X$ est un sous-schéma fermé intègre de dimension $\leq d$
et $\mathcal{I} \subset \mathcal{O}_V$ est un faisceau cohérent
d'idéaux non nul. Si $\dim(V) < d$, et plus généralement pour $\mathcal{F}$
dont le support est de dimension $< d$, le premier terme
de $E(\mathcal{F})$ est de degré total $< d$ d'après le
lemme \ref{varieties-lemma-numerical-polynomial-from-euler},
et le second terme est nul. Si $\dim(V) = d$, on peut utiliser la
suite exacte courte
$$
0 \to (V \to X)_*\mathcal{I} \to (V \to X)_*\mathcal{O}_V
\to (V \to X)_*(\mathcal{O}_V/\mathcal{I}) \to 0
$$
Le résultat vaut pour le faisceau du milieu, car
le seul $Z$ qui figure dans la somme est $Z = V$,
avec $m_Z = 1$, et parce que
$$
H^i(X, ((V \to X)_*\mathcal{O}_V) \otimes 
 \mathcal{L}_1^{\otimes n_1} \otimes \ldots \otimes
\mathcal{L}_r^{\otimes n_r}) =
H^i(V,  \mathcal{L}_1^{\otimes n_1} \otimes \ldots \otimes
\mathcal{L}_r^{\otimes n_r}|_V)
$$
d'après la formule de projection
(Cohomologie, section \ref{cohomology-section-projection-formula}) et
Cohomologie des schémas, lemme
\ref{coherent-lemma-relative-affine-cohomology} ;
dans ce cas, on a donc en fait $E(\mathcal{F}) = 0$.
Le résultat vaut pour le faisceau de droite, car son support
est de dimension $< d$. Il vaut donc pour le faisceau
de gauche, ce qui démontre le lemme.
\end{proof}

\begin{definition}
\label{varieties-definition-intersection-number}
Soit $k$ un corps. Soit $X$ un schéma propre sur $k$. Soit
$i : Z \to X$ un sous-schéma fermé de dimension $d$. Soient
$\mathcal{L}_1, \ldots, \mathcal{L}_d$ des
$\mathcal{O}_X$-modules inversibles. On définit le {\it nombre d'intersection}
$(\mathcal{L}_1 \cdots \mathcal{L}_d \cdot Z)$
comme le coefficient du monôme $n_1 \ldots n_d$ dans le polynôme numérique
$$
\chi(X, i_*\mathcal{O}_Z \otimes \mathcal{L}_1^{\otimes n_1} \otimes
\ldots \otimes \mathcal{L}_d^{\otimes n_d}) =
\chi(Z, \mathcal{L}_1^{\otimes n_1} \otimes
\ldots \otimes \mathcal{L}_d^{\otimes n_d}|_Z)
$$
Dans le cas particulier
où $\mathcal{L}_1 = \ldots = \mathcal{L}_d = \mathcal{L}$,
on écrit $(\mathcal{L}^d \cdot Z)$.
\end{definition}

\noindent
L'égalité affichée dans la définition résulte de
la formule de projection
(Cohomologie, section \ref{cohomology-section-projection-formula}) et de
Cohomologie des schémas, lemme
\ref{coherent-lemma-relative-affine-cohomology}.
Nous allons établir quelques lemmes sur ces nombres d'intersection.

\begin{lemma}
\label{varieties-lemma-intersection-number-integer}
Dans la situation de la définition \ref{varieties-definition-intersection-number},
le nombre d'intersection
$(\mathcal{L}_1 \cdots \mathcal{L}_d \cdot Z)$
est un entier.
\end{lemma}

\begin{proof}
Tout polynôme numérique de degré $e$ en $n_1, \ldots, n_d$ s'écrit
de manière unique comme combinaison $\mathbf{Z}$-linéaire des fonctions
${n_1 \choose k_1}{n_2 \choose k_2} \ldots {n_d \choose k_d}$ avec
$k_1 + \ldots + k_d \leq e$. Appliquons cela avec $e = d$.
Nous laissons les détails en exercice.
\end{proof}

\begin{lemma}
\label{varieties-lemma-intersection-number-additive}
Dans la situation de la définition \ref{varieties-definition-intersection-number},
le nombre d'intersection
$(\mathcal{L}_1 \cdots \mathcal{L}_d \cdot Z)$
est additif : si $\mathcal{L}_i = \mathcal{L}_i' \otimes \mathcal{L}_i''$,
alors
$$
(\mathcal{L}_1 \cdots \mathcal{L}_i \cdots \mathcal{L}_d \cdot Z) =
(\mathcal{L}_1 \cdots \mathcal{L}_i' \cdots \mathcal{L}_d \cdot Z) +
(\mathcal{L}_1 \cdots \mathcal{L}_i'' \cdots \mathcal{L}_d \cdot Z)
$$
\end{lemma}

\begin{proof}
Cela résulte de ce que, d'après le lemme \ref{varieties-lemma-numerical-polynomial-from-euler},
la fonction
$$
(n_1, \ldots, n_{i - 1}, n_i', n_i'', n_{i + 1}, \ldots, n_d)
\mapsto
\chi(Z, \mathcal{L}_1^{\otimes n_1} \otimes
\ldots \otimes (\mathcal{L}_i')^{\otimes n_i'} \otimes
(\mathcal{L}_i'')^{\otimes n_i''} \otimes \ldots \otimes
\mathcal{L}_d^{\otimes n_d}|_Z)
$$
est un polynôme numérique de degré total au plus $d$ en $d + 1$ variables.
\end{proof}

\begin{lemma}
\label{varieties-lemma-intersection-number-in-terms-of-components}
Dans la situation de la définition \ref{varieties-definition-intersection-number},
soient $Z_i \subset Z$ les composantes irréductibles de dimension $d$. Posons
$m_i = \text{longueur}_{\mathcal{O}_{X, \xi_i}}(\mathcal{O}_{Z, \xi_i})$,
où $\xi_i \in Z_i$ est le point générique. Alors
$$
(\mathcal{L}_1 \cdots \mathcal{L}_d \cdot Z) =
\sum m_i(\mathcal{L}_1 \cdots \mathcal{L}_d \cdot Z_i)
$$
\end{lemma}

\begin{proof}
Cela résulte immédiatement du lemme \ref{varieties-lemma-numerical-polynomial-leading-term}
et des définitions.
\end{proof}

\begin{lemma}
\label{varieties-lemma-intersection-number-and-pullback}
Soit $k$ un corps. Soit $f : Y \to X$ un morphisme de schémas propres sur $k$.
Soit $Z \subset Y$ un sous-schéma fermé intègre de dimension $d$, et soient
$\mathcal{L}_1, \ldots, \mathcal{L}_d$ des $\mathcal{O}_X$-modules inversibles.
Alors
$$
(f^*\mathcal{L}_1 \cdots f^*\mathcal{L}_d \cdot Z) =
\deg(f|_Z : Z \to f(Z)) (\mathcal{L}_1 \cdots \mathcal{L}_d \cdot f(Z))
$$
où $\deg(Z \to f(Z))$ est défini comme dans
Morphismes, définition \ref{morphisms-definition-degree},
ou vaut $0$ si $\dim(f(Z)) < d$.
\end{lemma}

\begin{proof}
Le membre de gauche est calculé à l'aide du coefficient de $n_1 \ldots n_d$
dans la fonction
$$
\chi(Y, \mathcal{O}_Z \otimes f^*\mathcal{L}_1^{\otimes n_1} \otimes
\ldots \otimes f^*\mathcal{L}_d^{\otimes n_d}) =
\sum (-1)^i
\chi(X, R^if_*\mathcal{O}_Z \otimes
\mathcal{L}_1^{\otimes n_1} \otimes \ldots \otimes
\mathcal{L}_d^{\otimes n_d})
$$
L'égalité résulte du lemme \ref{varieties-lemma-euler-characteristic-morphism}
et de la formule de projection
(Cohomologie, lemme \ref{cohomology-lemma-projection-formula}).
Si $f(Z)$ est de dimension $< d$, alors le membre de droite
est un polynôme de degré total $< d$ d'après le
lemme \ref{varieties-lemma-numerical-polynomial-from-euler},
et le résultat est établi. Supposons que $\dim(f(Z)) = d$. Soit
$\xi \in f(Z)$ le point générique. D'après la
théorie de la dimension (voir les lemmes \ref{varieties-lemma-dimension-locally-algebraic} et
\ref{varieties-lemma-dimension-fibres-locally-algebraic}),
le point générique de $Z$ est l'unique point de $Z$ envoyé sur $\xi$.
Alors $f : Z \to f(Z)$ est fini au-dessus d'un ouvert non vide de $f(Z)$, voir
Morphismes, lemme \ref{morphisms-lemma-generically-finite}.
Ainsi $\deg(f : Z \to f(Z))$ est défini et est en fait égal
à la longueur de la fibre de $f_*\mathcal{O}_Z$ en $\xi$, considérée
comme module sur $\mathcal{O}_{X, \xi}$. De plus, la fibre de
$R^if_*\mathcal{O}_X$ en $\xi$ est nulle pour $i > 0$, car
nous venons de voir que $f|_Z$ est fini dans un voisinage de $\xi$
(de sorte que Cohomologie des Schémas, lemme
\ref{coherent-lemma-finite-pushforward-coherent} donne cette annulation).
Ainsi, les termes $\chi(X, R^if_*\mathcal{O}_Z \otimes
\mathcal{L}_1^{\otimes n_1} \otimes \ldots \otimes
\mathcal{L}_d^{\otimes n_d})$ avec $i > 0$ sont de degré total
$< d$, et
$$
\chi(X, f_*\mathcal{O}_Z \otimes
\mathcal{L}_1^{\otimes n_1} \otimes \ldots \otimes
\mathcal{L}_d^{\otimes n_d})
=
\deg(f : Z \to f(Z)) \chi(f(Z),
\mathcal{L}_1^{\otimes n_1} \otimes \ldots \otimes
\mathcal{L}_d^{\otimes n_d}|_{f(Z)})
$$
modulo un polynôme de degré total $< d$, d'après le
lemme \ref{varieties-lemma-numerical-polynomial-leading-term}.
Le résultat voulu en résulte.
\end{proof}

\begin{lemma}
\label{varieties-lemma-numerical-intersection-effective-Cartier-divisor}
Soit $k$ un corps. Soit $X$ propre sur $k$. Soit $Z \subset X$ un
sous-schéma fermé de dimension $d$. Soient $\mathcal{L}_1, \ldots, \mathcal{L}_d$
des $\mathcal{O}_X$-modules inversibles. Supposons qu'il existe un
diviseur de Cartier effectif $D \subset Z$ tel que
$\mathcal{L}_1|_Z \cong \mathcal{O}_Z(D)$. Alors
$$
(\mathcal{L}_1 \cdots \mathcal{L}_d \cdot Z) =
(\mathcal{L}_2 \cdots \mathcal{L}_d \cdot D)
$$
\end{lemma}

\begin{proof}
On peut remplacer $X$ par $Z$ et $\mathcal{L}_i$ par $\mathcal{L}_i|_Z$.
On peut donc supposer que $X = Z$ et $\mathcal{L}_1 = \mathcal{O}_X(D)$.
Alors $\mathcal{L}_1^{-1}$ est le faisceau d'idéaux de $D$ et l'on peut
considérer la suite exacte courte
$$
0 \to \mathcal{L}_1^{\otimes -1} \to \mathcal{O}_X \to \mathcal{O}_D \to 0
$$
Posons
$P(n_1, \ldots, n_d) =
\chi(X, \mathcal{L}_1^{\otimes n_1} \otimes \ldots \otimes
\mathcal{L}_d^{\otimes n_d})$
et
$Q(n_1, \ldots, n_d) =
\chi(D, \mathcal{L}_1^{\otimes n_1} \otimes \ldots \otimes
\mathcal{L}_d^{\otimes n_d}|_D)$.
L'additivité donne alors
$$
P(n_1, \ldots, n_d) - P(n_1 - 1, n_2, \ldots, n_d) =
Q(n_1, \ldots, n_d)
$$
Comme le degré total de $P$ est au plus $d$, on voit que
le coefficient de $n_1 \ldots n_d$ dans $P$ est égal au coefficient
de $n_2 \ldots n_d$ dans $Q$.
\end{proof}

\begin{lemma}
\label{varieties-lemma-ample-positive}
Soit $k$ un corps. Soit $X$ propre sur $k$. Soit $Z \subset X$ un
sous-schéma fermé de dimension $d$. Si $\mathcal{L}_1, \ldots, \mathcal{L}_d$
sont amples, alors $(\mathcal{L}_1 \cdots \mathcal{L}_d \cdot Z)$ est strictement positif.
\end{lemma}

\begin{proof}
Raisonnons par récurrence sur $d$. Le cas $d = 0$
résulte du lemme \ref{varieties-lemma-chi-tensor-finite}. Supposons $d > 0$.
D'après le lemme \ref{varieties-lemma-intersection-number-in-terms-of-components},
on peut supposer que $Z$ est un sous-schéma fermé intègre.
En fait, on peut remplacer $X$ par $Z$ et $\mathcal{L}_i$
par $\mathcal{L}_i|_Z$ pour se ramener au cas où $Z = X$ est une
variété propre de dimension $d$.
D'après le lemme \ref{varieties-lemma-intersection-number-additive},
on peut remplacer $\mathcal{L}_1$ par une puissance tensorielle positive.
On peut donc supposer qu'il existe une section non nulle
$s \in \Gamma(X, \mathcal{L}_1)$
telle que $X_s$ soit affine (on utilise ici la définition d'un
faisceau inversible ample, voir
Propriétés, définition \ref{properties-definition-ample}).
Remarquons que $X$ n'est pas affine, car être propre et affine
implique être fini (Morphismes, lemme \ref{morphisms-lemma-finite-proper}),
ce qui contredit $d > 0$. Il s'ensuit que $s$ possède un schéma des zéros non vide
$Z(s) \subset X$. Comme $X$ est une variété, $s$ est une section régulière
de $\mathcal{L}_1$ ; ainsi $Z(s)$ est un diviseur de Cartier effectif,
donc $Z(s)$ est de codimension $1$ dans $X$, et
par conséquent $Z(s)$ est de dimension $d - 1$ (on utilise ici les résultats de
Diviseurs, sections \ref{divisors-section-effective-Cartier-divisors},
\ref{divisors-section-effective-Cartier-invertible} et
\ref{divisors-section-Noetherian-effective-Cartier},
ainsi que la théorie de la dimension comme dans le lemme \ref{varieties-lemma-dimension-locally-algebraic}).
D'après le lemme \ref{varieties-lemma-numerical-intersection-effective-Cartier-divisor},
on a
$$
(\mathcal{L}_1 \cdots \mathcal{L}_d \cdot X) =
(\mathcal{L}_2 \cdots \mathcal{L}_d \cdot Z(s))
$$
Le membre de droite est strictement positif par récurrence, ce qui achève la preuve.
\end{proof}

\begin{definition}
\label{varieties-definition-degree}
Soit $k$ un corps. Soit $X$ un schéma propre sur $k$. Soit
$\mathcal{L}$ un $\mathcal{O}_X$-module inversible ample.
Pour tout sous-schéma fermé $Z$, son {\it degré relativement à
$\mathcal{L}$}, noté $\deg_\mathcal{L}(Z)$, est
le nombre d'intersection $(\mathcal{L}^d \cdot Z)$,
où $d = \dim(Z)$.
\end{definition}

\noindent
D'après le lemme \ref{varieties-lemma-ample-positive}, le degré d'un sous-schéma est toujours un
entier strictement positif. Remarquons que $\deg_\mathcal{L}(Z) = d$ si et seulement si
$$
\chi(Z, \mathcal{L}^{\otimes n}|_Z) = \frac{d}{\dim(Z)!} n^{\dim(Z)} + l.o.t
$$
ce que l'on voit en utilisant l'identité
$$
(n_1 + \ldots + n_{\dim(Z)})^{\dim(Z)} =
\dim(Z)!\ n_1 \ldots n_{\dim(Z)} + \text{autres termes}
$$

\begin{lemma}
\label{varieties-lemma-degree-finite-morphism-in-terms-degrees}
Soit $k$ un corps. Soit $f : Y \to X$ un morphisme fini
dominant de variétés propres sur $k$. Soit $\mathcal{L}$
un $\mathcal{O}_X$-module inversible ample.
Alors
$$
\deg_{f^*\mathcal{L}}(Y) = \deg(f) \deg_\mathcal{L}(X)
$$
où $\deg(f)$ est défini dans
Morphismes, définition \ref{morphisms-definition-degree}.
\end{lemma}

\begin{proof}
L'énoncé a un sens, car $f^*\mathcal{L}$ est ample d'après
Morphismes, lemme \ref{morphisms-lemma-pullback-ample-tensor-relatively-ample}.
Cela étant, le résultat est un cas particulier du
lemme \ref{varieties-lemma-intersection-number-and-pullback}.
\end{proof}

\noindent
Enfin, nous mettons en relation le nombre d'intersection avec une courbe et la notion
de degré des modules inversibles sur les courbes introduite dans la
section \ref{varieties-section-divisors-curves}.

\begin{lemma}
\label{varieties-lemma-intersection-numbers-and-degrees-on-curves}
Soit $k$ un corps. Soit $X$ un schéma propre sur $k$.
Soit $Z \subset X$ un sous-schéma fermé de dimension $\leq 1$.
Soit $\mathcal{L}$ un $\mathcal{O}_X$-module inversible.
Alors
$$
(\mathcal{L} \cdot Z) = \deg(\mathcal{L}|_Z)
$$
où $\deg(\mathcal{L}|_Z)$ est défini dans la
définition \ref{varieties-definition-degree-invertible-sheaf}.
Si $\mathcal{L}$ est ample, alors
$\deg_\mathcal{L}(Z) = \deg(\mathcal{L}|_Z)$.
\end{lemma}

\begin{proof}
Cela résulte du fait que la fonction
$n \mapsto \chi(Z, \mathcal{L}|_Z^{\otimes n})$ est de degré $1$
et que son coefficient dominant est donc la différence de valeurs consécutives.
\end{proof}

\begin{proposition}[Riemann-Roch asymptotique]
\label{varieties-proposition-asymptotic-riemann-roch}
Soit $k$ un corps. Soit $X$ un schéma propre sur $k$ de dimension $d$.
Soit $\mathcal{L}$ un $\mathcal{O}_X$-module inversible ample.
Alors
$$
\dim_k \Gamma(X, \mathcal{L}^{\otimes n}) \sim c n^d + l.o.t.
$$
où $c = \deg_\mathcal{L}(X)/d!$ est une constante strictement positive.
\end{proposition}

\begin{proof}
Cela résulte des définitions,
du lemme \ref{varieties-lemma-ample-positive} et de l'annulation
de la cohomologie en degrés supérieurs donnée dans
Cohomologie des Schémas, lemme \ref{coherent-lemma-vanshing-gives-ample}.
\end{proof}






\section{Dimension de plongement}
\label{varieties-section-embedding-dimension}

\noindent
Il existe plusieurs façons de définir la dimension de plongement, mais, pour
les points fermés des schémas algébriques sur des corps algébriquement clos,
toutes ces définitions sont équivalentes à la suivante.

\begin{definition}
\label{varieties-definition-embed-dim}
Soit $k$ un corps algébriquement clos. Soit $X$ un
$k$-schéma localement algébrique et soit $x \in X$ un point fermé. La
{\it dimension de plongement de $X$ en $x$} est
$\dim_k \mathfrak m_x/\mathfrak m_x^2$.
\end{definition}

\noindent
Quelques faits sur la dimension de plongement.
Soient $k, X, x$ comme dans la définition \ref{varieties-definition-embed-dim}.
\begin{enumerate}
\item La dimension de plongement de $X$ en $x$ est
la dimension de
l'espace tangent $T_{X/k, x}$ (définition \ref{varieties-definition-tangent-space})
comme $k$-espace vectoriel.
\item La dimension de plongement de $X$ en $x$ est le plus petit entier
$d \geq 0$ tel qu'il existe une surjection
$$
k[[x_1, \ldots, x_d]] \longrightarrow \mathcal{O}_{X, x}^\wedge
$$
de $k$-algèbres.
\item La dimension de plongement de $X$ en $x$ est le plus petit entier
$d \geq 0$ tel qu'il existe un voisinage ouvert $U \subset X$
de $x$ et une immersion fermée $U \to Y$, où $Y$ est une variété lisse
de dimension $d$ sur $k$.
\item La dimension de plongement de $X$ en $x$ est le plus petit entier
$d \geq 0$ tel qu'il existe un voisinage ouvert $U \subset X$
de $x$ et un morphisme non ramifié $U \to \mathbf{A}^d_k$.
\item Si l'on se donne une immersion fermée $X \to Y$ avec $Y$ lisse
sur $k$, alors la dimension de plongement de $X$ en $x$ est le plus petit
entier $d \geq 0$ tel qu'il existe un sous-schéma fermé $Z \subset Y$
tel que $X \subset Z$, que $Z \to \Spec(k)$ soit lisse en $x$ et que
$\dim_x(Z) = d$.
\end{enumerate}
Si ces faits nous sont nécessaires, nous formulerons un énoncé précis et en donnerons
une démonstration.

\medskip\noindent
Corps de base non algébriquement clos ou points non fermés.
Soit $k$ un corps et soit $X$ un $k$-schéma localement algébrique.
Si $x \in X$ est un point, plusieurs choix sont possibles pour la
dimension de plongement de $X$ en $x$. On peut en effet prendre
\begin{enumerate}
\item $\dim_{\kappa(x)}(\mathfrak m_x/\mathfrak m_x^2)$,
\item $\dim_{\kappa(x)}(T_{X/k, x}) =
\dim_{\kappa(x)}(\Omega_{X/k, x} \otimes_{\mathcal{O}_{X, x}} \kappa(x))$
(lemme \ref{varieties-lemma-tangent-space-cotangent-space}),
\item le plus petit entier $d \geq 0$ tel qu'il existe un
voisinage ouvert $U \subset X$ de $x$ et une immersion fermée
$U \to Y$, où $Y$ est une variété lisse de dimension $d$ sur $k$.
\end{enumerate}
En caractéristique zéro, (1) $=$ (2) si $x$ est un point fermé ;
plus généralement, cette égalité vaut si $\kappa(x)$ est algébrique séparable
sur $k$, voir le
lemme \ref{varieties-lemma-tangent-space-rational-point}.
La définition géométrique (3) semble correspondre le mieux à
l'intuition géométrique qu'évoque l'expression
``dimension de plongement''.
Puisque l'on peut montrer que (3) et (2) définissent le même nombre
(cela résulte du lemme \ref{varieties-lemma-immersion-into-affine})
c'est ce que nous utiliserons.
Dans notre terminologie, nous préciserons que nous
prenons la dimension de plongement relativement au corps de base.

\begin{definition}
\label{varieties-definition-embedding-dimension}
Soit $k$ un corps. Soit $X$ un $k$-schéma localement algébrique.
Soit $x \in X$ un point. La {\it dimension de plongement de $X/k$ en $x$}
est $\dim_{\kappa(x)}(T_{X/k, x})$.
\end{definition}

\noindent
Si $(A, \mathfrak m, \kappa)$ est un anneau local noethérien,
la {\it dimension de plongement de $A$} est parfois définie comme la dimension de
$\mathfrak m/\mathfrak m^2$ sur $\kappa$. Nous avons vu ci-dessus que,
si $A$ est donné comme algèbre sur un corps $k$, il peut être préférable
d'utiliser $\dim_\kappa(\Omega_{A/k} \otimes_A \kappa)$. Appelons
cette quantité la {\it dimension de plongement de $A/k$}.
Avec cette terminologie, on a
$$
\begin{aligned}
\text{dimension de plongement de }X/k\text{ en }x
&= \text{dimension de plongement de }\mathcal{O}_{X, x}/k \\
&= \text{dimension de plongement de }\mathcal{O}_{X, x}^\wedge/k
\end{aligned}
$$
si $k, X, x$ sont comme dans la définition \ref{varieties-definition-embedding-dimension}.






\section{Théorèmes de Bertini}
\label{varieties-section-bertini}

\noindent
Dans cette section, nous démontrons des résultats de la forme suivante : étant donnée une variété
projective lisse $X$ sur un corps $k$, il existe un diviseur ample $H \subset X$
qui est lisse.

\begin{lemma}
\label{varieties-lemma-generate-over-complement}
Soit $k$ un corps. Soit $X$ un schéma propre sur $k$. Soit $\mathcal{L}$
un $\mathcal{O}_X$-module inversible ample. Soit $Z \subset X$ un
sous-schéma fermé. Alors il existe un entier $n_0$ tel que, pour tout
$n \geq n_0$, le noyau $V_n$ de
$\Gamma(X, \mathcal{L}^{\otimes n}) \to \Gamma(Z, \mathcal{L}^{\otimes n}|_Z)$
engendre $\mathcal{L}^{\otimes n}|_{X \setminus Z}$ et
le morphisme canonique
$$
X \setminus Z \longrightarrow \mathbf{P}(V_n)
$$
est une immersion de schémas sur $k$.
\end{lemma}

\begin{proof}
Soit $\mathcal{I} \subset \mathcal{O}_X$ le faisceau quasi-cohérent d'idéaux de
$Z$. Remarquons que, via l'inclusion
$\mathcal{I} \otimes_{\mathcal{O}_X} \mathcal{L}^{\otimes n} \subset
\mathcal{L}^{\otimes n}$, on a
$V_n = \Gamma(X, \mathcal{I} \otimes_{\mathcal{O}_X} \mathcal{L}^{\otimes n})$.
Choisissons $n_1$ tel que, pour $n \geq n_1$, le faisceau
$\mathcal{I} \otimes \mathcal{L}^{\otimes n}$ soit engendré par ses sections globales, voir
Propriétés, proposition \ref{properties-proposition-characterize-ample}.
Il en résulte que $V_n$ engendre $\mathcal{L}^{\otimes n}|_{X \setminus Z}$
pour $n \geq n_1$.

\medskip\noindent
Pour $n \geq n_1$, notons
$\psi_n : V_n \to
\Gamma(X \setminus Z, \mathcal{L}^{\otimes n}|_{X \setminus Z})$
l'application de restriction. On obtient un morphisme canonique
$$
\varphi = \varphi_{\mathcal{L}^{\otimes n}|_{X \setminus Z}, \psi_n} :
X \setminus Z
\longrightarrow
\mathbf{P}(V_n)
$$
d'après Constructions, exemple \ref{constructions-example-projective-space}.
Choisissons $n_2$ tel que, pour tout $n \geq n_2$, le faisceau inversible
$\mathcal{L}^{\otimes n}$ soit très ample sur $X$.
Nous affirmons que $n_0 = n_1 + n_2$ convient.

\medskip\noindent
Preuve de l'affirmation. Supposons $n \geq n_0$ et écrivons $n = n_1 + n'$.
Pour $x \in X \setminus Z$, on peut choisir $s_1 \in V_1$ ne
s'annulant pas en $x$. Posons $V' = \Gamma(X, \mathcal{L}^{\otimes n'})$.
Par le choix de $n$ et de $n'$, le morphisme correspondant
$\varphi' : X \to \mathbf{P}(V')$ est une immersion fermée. Ainsi, si l'on choisit
$s' \in \Gamma(X, \mathcal{L}^{\otimes n'})$ ne s'annulant pas en $x$,
alors $X_{s'} = (\varphi')^{-1}(D_+(s'))$ (voir
Constructions, lemme \ref{constructions-lemma-invertible-map-into-proj})
est affine et $X_{s'} \to D_+(s')$ est une immersion fermée.
Alors $s = s_1 \otimes s' \in V_n$ ne s'annule pas en $x$.
Si $D_+(s) \subset \mathbf{P}(V_n)$ désigne
l'ouvert affine correspondant de notre espace projectif, alors
$\varphi^{-1}(D_+(s)) = X_s \subset X \setminus Z$ (voir la référence ci-dessus).
L'ouvert $X_s = X_{s'} \cap X_{s_1}$ est affine, voir
Propriétés, lemme \ref{properties-lemma-affine-cap-s-open}.
Considérons l'homomorphisme d'anneaux
$$
\text{Sym}(V)_{(s)} \longrightarrow \mathcal{O}_X(X_s)
$$
qui définit le morphisme $X_s \to D_+(s)$. Comme $X_{s'} \to D_+(s')$
est une immersion fermée, les images des éléments
$$
\frac{s_1 \otimes t'}{s_1 \otimes s'}
$$
où $t' \in V'$, engendrent l'image de
$\mathcal{O}_X(X_{s'}) \to \mathcal{O}_X(X_s)$.
Puisque $X_s \to X_{s'}$ est une immersion ouverte,
il s'ensuit que $X_s \to D_+(s)$ est une immersion de schémas affines
(voir ci-dessous). Ainsi, $\varphi_n$ est une immersion d'après
Morphismes, lemme \ref{morphisms-lemma-check-immersion}.

\medskip\noindent
Soient $a : A' \to A$ et $c : B \to A$ des homomorphismes d'anneaux tels que
$\Spec(a)$ soit une immersion et $\Im(a) \subset \Im(c)$.
Posons $B' = A' \times_A B$, avec les projections $b : B' \to B$
et $c' : B' \to A'$. Par hypothèse, $c'$
est surjectif, de sorte que $\Spec(c')$ est une immersion fermée.
Par conséquent, $\Spec(c') \circ \Spec(a)$ est une immersion
(Schémas, lemme \ref{schemes-lemma-composition-immersion}).
Alors $\Spec(c)$ est nécessairement une immersion, car il factorise
l'immersion $\Spec(c') \circ \Spec(a) = \Spec(b) \circ \Spec(c)$, voir
Morphismes, lemme \ref{morphisms-lemma-immersion-permanence}.
\end{proof}

\begin{situation}
\label{varieties-situation-family-divisors}
Soit $k$ un corps,
soit $X$ un schéma sur $k$,
soit $\mathcal{L}$ un $\mathcal{O}_X$-module inversible,
soit $V$ un $k$-espace vectoriel de dimension finie, et
soit $\psi : V \to \Gamma(X, \mathcal{L})$ une application $k$-linéaire.
Supposons $\dim(V) = r$ et fixons une base $v_1, \ldots, v_r$ de $V$.
On obtient alors un ``diviseur universel''
$$
H_{univ} = Z(s_{univ}) \subset  \mathbf{A}^r \times_k X
$$
qui est le schéma des zéros
(Diviseurs, définition \ref{divisors-definition-zero-scheme-s}) de la section
$$
s_{univ} = \sum\nolimits_{i = 1, \ldots, r} x_i \psi(v_i) \in
\Gamma(\mathbf{A}^r \times_k X, \text{pr}_2^*\mathcal{L})
$$
Pour une extension de corps $k'/k$, les $k'$-points $v \in \mathbf{A}^r_k(k')$
correspondent aux vecteurs $(a_1, \ldots, a_r)$ d'éléments de $k'$.
Ainsi, on peut d'une part considérer $v$ comme l'élément
$v = \sum_{i = 1, \ldots, r} a_i v_i \in V \otimes_k k'$
et, d'autre part, associer à $v$ la section
$$
\psi(v) = \sum\nolimits_{i = 1, \ldots, r} a_i \psi(v_i) \in
\Gamma(X_{k'}, \mathcal{L}|_{X_{k'}})
$$
Avec ces notations, il est clair que la fibre de $H_{univ}$ au-dessus de
$v \in V \otimes k'$ est le schéma des zéros de $\psi(v)$.
En formule :
$$
H_v = H_{univ, v} = Z(\psi(v))
$$
Nous désignerons cette valeur commune par $H_v$, comme indiqué. Enfin, dans
cette situation, soit $P$ une propriété des vecteurs $v \in V \otimes_k k'$
pour une extension de corps arbitraire $k'/k$\footnote{Par exemple, nous pourrions
considérer la condition que $H_v$ soit lisse sur $k'$, ou géométriquement
irréductible sur $k'$.}. Nous disons que $P$ {\it est vérifiée pour un vecteur général}
$v \in V \otimes_k k'$ s'il existe un ouvert de Zariski non vide
$U \subset \mathbf{A}^r_k$ tel que, si $v$ correspond à un $k'$-point
de $U$, quelle que soit l'extension $k'/k$, alors $P(v)$ est vérifiée.
\end{situation}

\begin{lemma}
\label{varieties-lemma-bertini}
Dans la situation \ref{varieties-situation-family-divisors}, supposons que
\begin{enumerate}
\item $X$ soit lisse sur $k$,
\item l'image de $\psi : V \to \Gamma(X, \mathcal{L})$
engendre $\mathcal{L}$,
\item le morphisme correspondant
$\varphi_{\mathcal{L}, \psi} : X \to \mathbf{P}(V)$ soit
une immersion.
\end{enumerate}
Alors, pour un vecteur général $v \in V \otimes_k k'$,
le schéma $H_v$ est lisse sur $k'$.
\end{lemma}

\begin{proof}
(Nous observons que $X$ est séparé et de type fini comme sous-schéma
localement fermé d'un espace projectif.)
Utilisons les notations introduites avant l'énoncé du lemme.
Considérons les projections
$$
\xymatrix{
\mathbf{A}^r_k \times_k X \ar[d] &
H_{univ} \ar[l] \ar[ld]^p \ar[r] \ar[rd]_q &
\mathbf{A}^r_k \times_k X \ar[d] \\
X & &
\mathbf{A}^r_k
}
$$
Soit $\Sigma \subset H_{univ}$ le lieu singulier du morphisme
$q : H_{univ} \to \mathbf{A}^r_k$, c'est-à-dire l'ensemble des points où
$q$ n'est pas lisse. Alors $\Sigma$ est fermé, car le lieu lisse
d'un morphisme est ouvert par définition. Comme toute fibre d'un
morphisme lisse est lisse, il suffit de montrer que $q(\Sigma)$
est contenu dans un sous-ensemble fermé strict de $\mathbf{A}^r_k$.
Puisque $\Sigma$ (muni de la structure de schéma induite réduite) est un
schéma de type fini sur $k$, il suffit de montrer que $\dim(\Sigma) < r$.
Cela résulte du lemme \ref{varieties-lemma-dimension-fibres-locally-algebraic}.
Comme les dimensions ne changent pas lorsqu'on remplace $k$ par un corps plus grand
(Morphismes, lemme \ref{morphisms-lemma-dimension-fibre-after-base-change}),
nous pouvons supposer, et nous supposons désormais, que $k$ est algébriquement clos.
D'après la théorie de la dimension (lemme \ref{varieties-lemma-dimension-fibres-locally-algebraic}),
il suffit de montrer que, pour tout point fermé $x \in X \setminus Z$,
$p^{-1}(\{x\}) \cap \Sigma$ est de dimension $< r - \dim(X')$,
où $X'$ est l'unique composante irréductible de $X$ contenant $x$.
Comme $X$ est lisse sur $k$ et que $x$ est un point fermé, on a
$\dim(X') = \dim \mathfrak m_x/\mathfrak m_x^2$
(Morphismes, lemme \ref{morphisms-lemma-smooth-omega-finite-locally-free}
et Algèbre, lemme \ref{algebra-lemma-rank-omega}).
Il suffit donc que
$$
\dim p^{-1}(x) \cap \Sigma < r - \dim \mathfrak m_x/\mathfrak m_x^2
$$
pour tout point fermé $x \in X$.

\medskip\noindent
Puisque les sections globales provenant de $V$ engendrent $\mathcal{L}$, pour toute composante irréductible
$X'$ de $X$, il existe un ouvert de Zariski non vide de $\mathbf{A}^r$ tel que les
fibres de $q$ au-dessus de cet ouvert ne contiennent pas $X'$. (Par exemple, si $x' \in X'$
est un point fermé, on peut prendre l'ouvert correspondant aux
vecteurs $v \in V$ tels que $\psi(v)$ ne s'annule pas en $x'$. Cet
ouvert est le complémentaire d'un hyperplan dans $\mathbf{A}^r_k$.)
Soit $U \subset \mathbf{A}^r$ l'intersection (non vide) de ces ouverts.
Alors les fibres de $q^{-1}(U) \to U$ sont des diviseurs de Cartier effectifs
sur les fibres de $U \times_k X \to U$ (car une section non nulle
d'un module inversible sur un schéma intègre est une section régulière). Par suite, le
morphisme $q^{-1}(U) \to U$ est plat d'après
Diviseurs, lemme \ref{divisors-lemma-fibre-Cartier}.
Ainsi, pour tout point fermé $x \in X$ et tout $v \in V = \mathbf{A}^r_k(k)$,
si $(x, v) \in H_{univ}$, c'est-à-dire si $x \in H_v$,
alors $q$ est lisse en $(x, v)$ si et seulement si la fibre
$H_v$ est lisse en $x$, voir
Morphismes, lemme \ref{morphisms-lemma-smooth-at-point}.

\medskip\noindent
Considérons l'image $\psi(v)_x$ dans la fibre $\mathcal{L}_x$
de la section correspondant à $v \in V$. On a
$$
x \in H_v \Leftrightarrow \psi(v)_x \in \mathfrak m_x\mathcal{L}_x
$$
Si cela est vrai, alors on a
$$
H_v\text{ est singulier en }x \Leftrightarrow
\psi(v)_x \in \mathfrak m_x^2\mathcal{L}_x
$$
En effet, $\psi(v)_x$ n'appartient pas à $\mathfrak m_x^2\mathcal{L}_x$
$\Leftrightarrow$
l'équation locale de $H_v \subset X$ en $x$ n'appartient pas
à $\mathfrak m_x^2$
$\Leftrightarrow$
$\mathcal{O}_{H_v, x}$ est régulier
(Algèbre, lemme \ref{algebra-lemma-regular-ring-CM})
$\Leftrightarrow$
$H_v$ est lisse en $x$ sur $k$
(Algèbre, lemme \ref{algebra-lemma-separable-smooth}).
On conclut que les points fermés de $p^{-1}(x) \cap \Sigma$ correspondent
aux $v \in V$ tels que $\psi(v)_x \in \mathfrak m_x^2\mathcal{L}_x$.
Cependant, comme $\varphi_{\mathcal{L}, \psi}$ est une immersion,
l'application
$$
V \longrightarrow \mathcal{L}_x/\mathfrak m_x^2\mathcal{L}_x
$$
est surjective (un petit détail est omis). D'après ce qui précède, les points fermés
du lieu $p^{-1}(x) \cap \Sigma$, considéré comme un sous-espace de $V$,
forment le noyau de cette application, lequel a donc pour dimension
$r - \dim \mathfrak m_x/\mathfrak m_x^2 - 1$, comme voulu.
\end{proof}






\section{Enriques-Severi-Zariski}
\label{varieties-section-vanishing-negative}

\noindent
Dans cette section, nous démontrons des résultats de la forme suivante : tensoriser par
un module inversible « très négatif » annule la cohomologie en
petits degrés. Nous en déduisons aussi la connexité
d'une section par une hypersurface d'un schéma propre normal de
dimension $\geq 2$.

\begin{lemma}
\label{varieties-lemma-vanishin-h0-negative}
Soit $k$ un corps. Soit $X$ un schéma propre sur $k$. Soit $\mathcal{L}$
un $\mathcal{O}_X$-module inversible ample. Soit $\mathcal{F}$ un
$\mathcal{O}_X$-module cohérent. Si $\text{Ass}(\mathcal{F})$ ne
contient aucun point fermé, alors
$\Gamma(X, \mathcal{F} \otimes_{\mathcal{O}_X} \mathcal{L}^{\otimes n}) = 0$
pour $n \ll 0$.
\end{lemma}

\begin{proof}
Pour un $\mathcal{O}_X$-module cohérent $\mathcal{F}$,
notons $\mathcal{P}(\mathcal{F})$ la propriété suivante : il
existe un $n_0 \in \mathbf{Z}$ tel que, pour $n \leq n_0$,
toute section $s$ de
$\mathcal{F} \otimes_{\mathcal{O}_X} \mathcal{L}^{\otimes n}$
ait un support constitué uniquement de points fermés.
Comme $\text{Ass}(\mathcal{F}) =
\text{Ass}(\mathcal{F} \otimes_{\mathcal{O}_X} \mathcal{L}^{\otimes n})$,
on voit qu'il suffit de démontrer que $\mathcal{P}$
est satisfaite par tout module cohérent sur $X$.
Pour ce faire, nous allons démontrer que les conditions (1), (2) et (3) de
Cohomologie des schémas, lemme
\ref{coherent-lemma-property-higher-rank-cohomological},
sont satisfaites.

\medskip\noindent
Pour vérifier la condition (1), supposons que
$$
0 \to \mathcal{F}_1 \to \mathcal{F} \to \mathcal{F}_2 \to 0
$$
soit une suite exacte courte de $\mathcal{O}_X$-modules cohérents
telle que $\mathcal{P}$ soit satisfaite par $\mathcal{F}_i$, $i = 1, 2$.
Soient $n_1, n_2$ les bornes ainsi obtenues.
Soit $\mathcal{F}'_2 \subset \mathcal{F}_2$ le plus grand
sous-module cohérent dont le support est un ensemble fini de points
fermés. Soit $\mathcal{I} \subset \mathcal{O}_X$ l'annulateur
de $\mathcal{F}'_2$. Comme $\mathcal{L}$ est ample, il existe un $e > 0$
tel que $\mathcal{I} \otimes_{\mathcal{O}_X} \mathcal{L}^{\otimes e}$
soit engendré par ses sections globales. Posons $n_0 = \min(n_2, n_1 - e)$. Soit $n \leq n_0$
et soit $t$ une section globale de
$\mathcal{F} \otimes \mathcal{L}^{\otimes n}$.
L'image de $t$ dans $\mathcal{F}_2 \otimes \mathcal{L}^{\otimes n}$
appartient à $\mathcal{F}'_2 \otimes \mathcal{L}^{\otimes n}$ puisque
$n \leq n_2$. Par conséquent, pour tout
$s \in \Gamma(X, \mathcal{I} \otimes_{\mathcal{O}_X} \mathcal{L}^{\otimes e})$,
le produit $t \otimes s$ appartient à
$\mathcal{F}_1 \otimes \mathcal{L}^{\otimes n + e}$.
Ainsi, le support de $t \otimes s$ est contenu dans
l'ensemble fini des points fermés de $\text{Ass}(\mathcal{F}_1)$
puisque $n + e \leq n_1$. Par notre choix de $e$,
pour tout point on peut choisir $s$ qui y soit inversible, sauf s'il
appartient au support de $\mathcal{F}'_2$ ;
on conclut que le support de $t$ est contenu dans la réunion de
l'ensemble fini des points fermés de $\text{Ass}(\mathcal{F}_1)$ et de
l'ensemble fini des points fermés de $\text{Ass}(\mathcal{F}_2)$.
Ceci achève la démonstration de la condition (1).

\medskip\noindent
La condition (2) est immédiate.

\medskip\noindent
Pour la condition (3), choisissons $\mathcal{G} = \mathcal{O}_Z$.
Dans ce cas, si $Z$ est un point fermé de $X$, il n'y a
rien à démontrer. Si $\dim(Z) > 0$, nous allons montrer que
$\Gamma(Z, \mathcal{L}^{\otimes n}|_Z) = 0$
pour $n < 0$. En effet, soit $s$ une section non nulle d'une puissance négative
de $\mathcal{L}|_Z$. Choisissons une section non nulle $t$ d'une puissance positive
de $\mathcal{L}|_Z$ (ce choix est possible puisque $\mathcal{L}$ est ample, voir
Propriétés, proposition \ref{properties-proposition-characterize-ample}).
Alors $s^{\deg(t)} \otimes t^{\deg(s)}$
est une section globale non nulle de $\mathcal{O}_Z$
(puisque $Z$ est intègre), et donc
une unité (lemme \ref{varieties-lemma-proper-geometrically-reduced-global-sections}).
Cela implique que $t$ est une section trivialisante
d'une puissance positive de $\mathcal{L}$.
Ainsi, la fonction $n \mapsto \dim_k \Gamma(X, \mathcal{L}^{\otimes n})$
est bornée sur un ensemble infini d'entiers positifs, ce qui contredit
le théorème de Riemann-Roch asymptotique
(proposition \ref{varieties-proposition-asymptotic-riemann-roch}),
puisque $\dim(Z) > 0$.
\end{proof}

\begin{lemma}[Enriques-Severi-Zariski]
\label{varieties-lemma-vanishin-h1-negative}
Soit $k$ un corps. Soit $X$ un schéma propre sur $k$. Soit $\mathcal{L}$
un $\mathcal{O}_X$-module inversible ample. Soit $\mathcal{F}$ un
$\mathcal{O}_X$-module cohérent. Supposons que,
pour tout point fermé $x \in X$, on ait $\text{profondeur}(\mathcal{F}_x) \geq 2$.
Alors
$H^1(X, \mathcal{F} \otimes_{\mathcal{O}_X} \mathcal{L}^{\otimes m}) = 0$
pour $m \ll 0$.
\end{lemma}

\begin{proof}
Choisissons une immersion fermée $i : X \to \mathbf{P}^n_k$ telle que
$i^*\mathcal{O}(1) \cong \mathcal{L}^{\otimes e}$ pour un certain $e > 0$
(voir Morphismes, lemme
\ref{morphisms-lemma-finite-type-over-affine-ample-very-ample}).
Il suffit alors de démontrer le lemme pour
$$
\mathcal{G} =
i_*(\mathcal{F} \oplus \mathcal{F} \otimes \mathcal{L} \oplus \ldots
\oplus \mathcal{F} \otimes \mathcal{L}^{\otimes e - 1})
\quad\text{et}\quad \mathcal{O}(1)
$$
sur $\mathbf{P}^n_k$. En effet, on a
$$
H^1(\mathbf{P}^n_k, \mathcal{G}(m)) =
\bigoplus\nolimits_{j = 0, \ldots, e - 1}
H^1(X, \mathcal{F} \otimes \mathcal{L}^{\otimes j + me})
$$
d'après
Cohomologie des schémas, lemme \ref{coherent-lemma-relative-affine-cohomology}.
De plus, si $y \in \mathbf{P}^n_k$ est un point fermé, alors
$\text{profondeur}(\mathcal{G}_y) = \infty$ si $y \not \in i(X)$
et $\text{profondeur}(\mathcal{G}_y) = \text{profondeur}(\mathcal{F}_x)$
si $y = i(x)$ puisque, dans ce cas,
$\mathcal{G}_y \cong \mathcal{F}_x^{\oplus e}$ comme module sur
$\mathcal{O}_{\mathbf{P}^n_k, x}$ et que l'on peut, par exemple, utiliser
Algèbre, lemme \ref{algebra-lemma-depth-goes-down-finite},
pour obtenir l'égalité.

\medskip\noindent
Supposons que $X = \mathbf{P}^n_k$, que $\mathcal{L} = \mathcal{O}(1)$ et que
$k$ soit infini. Choisissons $s \in H^0(\mathbf{P}^1_k, \mathcal{O}(1))$
qui détermine une suite exacte
$$
0 \to \mathcal{F}(-1) \xrightarrow{s} \mathcal{F} \to \mathcal{G} \to 0
$$
comme au lemme \ref{varieties-lemma-exact-sequence-induction}. Puisque le morphisme
$\mathcal{F}(-1) \to \mathcal{F}$ est, localement sur des ouverts affines, donné par
la multiplication par un non-diviseur de zéro dans $\mathcal{F}$,
on voit que, pour tout point fermé $x \in \mathbf{P}^n_k$, on a
$\text{profondeur}(\mathcal{G}_x) \geq 1$, voir
Algèbre, lemme \ref{algebra-lemma-depth-drops-by-one}.
Ainsi, d'après le lemme \ref{varieties-lemma-vanishin-h0-negative},
on a $H^0(\mathcal{G}(m)) = 0$ pour $m \ll 0$.
En considérant la suite exacte longue de cohomologie après tensorisation
(voir remarque \ref{varieties-remark-exact-sequence-induction-cohomology}),
on constate que la suite de nombres
$$
\dim H^1(\mathbf{P}^n_k, \mathcal{F}(m))
$$
se stabilise pour $m \leq m_0$, pour un certain entier $m_0$.
Soit $N$ la dimension commune de ces
espaces pour $m \leq m_0$. Il faut montrer que $N = 0$.

\medskip\noindent
Pour $d > 0$ et $m \leq m_0$, considérons l'application bilinéaire
$$
H^0(\mathbf{P}^n_k, \mathcal{O}(d)) \times
H^1(\mathbf{P}^n_k, \mathcal{F}(m - d))
\longrightarrow
H^1(\mathbf{P}^n_k, \mathcal{F}(m))
$$
D'après l'algèbre linéaire, il existe un sous-espace de codimension $\leq N^2$ que l'on note
$V_m \subset H^0(\mathbf{P}^n_k, \mathcal{O}(d))$ et tel que
la multiplication par tout $s' \in V_m$ annule
$H^1(\mathbf{P}^n_k, \mathcal{F}(m - d))$.
Notons que, pour $m' < m \leq m_0$, le diagramme
$$
\xymatrix{
H^0(\mathbf{P}^n_k, \mathcal{O}(d)) \times
H^1(\mathbf{P}^n_k, \mathcal{F}(m' - d)) \ar[r] \ar[d]^{1 \times s^{m' - m}} &
H^1(\mathbf{P}^n_k, \mathcal{F}(m')) \ar[d]^{s^{m' - m}}\\
H^0(\mathbf{P}^n_k, \mathcal{O}(d)) \times
H^1(\mathbf{P}^n_k, \mathcal{F}(m - d)) \ar[r] &
H^1(\mathbf{P}^n_k, \mathcal{F}(m))
}
$$
est commutatif, et les flèches verticales sont des isomorphismes. Ainsi, $V_m = V$ est
indépendant de $m \leq m_0$. Pour $x \in \text{Ass}(\mathcal{F})$,
posons $Z = \overline{\{x\}}$. Pour $d$ suffisamment grand, l'application linéaire
$$
H^0(\mathbf{P}^n_k, \mathcal{O}(d)) \to H^0(Z, \mathcal{O}(d)|_Z)
$$
est de rang $> N^2$ puisque $\dim(Z) \geq 1$ (cela résulte, par exemple,
du théorème de Riemann-Roch asymptotique et de l'amplitude de $\mathcal{O}(1)$ ; détails
omis). On peut donc trouver $s' \in V$ tel que $s'$ ne s'annule
en aucun point associé de $\mathcal{F}$ (on utilise le fait que l'ensemble
des points associés est fini). On obtient alors
$$
0 \to \mathcal{F}(-d) \xrightarrow{s'} \mathcal{F} \to \mathcal{G}' \to 0
$$
et, comme précédemment, on conclut que la multiplication par $s'$
sur $H^1(\mathbf{P}^n_k, \mathcal{F}(m - d))$ est injective
pour $m \ll 0$. Cela contredit le choix de $s'$, à moins que
$N = 0$, comme voulu.

\medskip\noindent
Il reste à traiter le cas où $k$ est fini.
Dans ce cas, soit $K/k$ une extension algébrique infinie.
Notons $\mathcal{F}_K$ et $\mathcal{L}_K$ les images réciproques
de $\mathcal{F}$ et $\mathcal{L}$ sur $X_K = \Spec(K) \times_{\Spec(k)} X$.
On a
$$
H^1(X_K, \mathcal{F}_K \otimes \mathcal{L}_K^{\otimes m}) =
H^1(X, \mathcal{F} \otimes \mathcal{L}^{\otimes m}) \otimes_k K
$$
d'après Cohomologie des schémas, lemme
\ref{coherent-lemma-flat-base-change-cohomology}.
D'autre part, un point fermé $x_K$ de $X_K$ a pour image un point fermé
$x$ de $X$ puisque $K/k$ est une extension algébrique.
L'homomorphisme d'anneaux $\mathcal{O}_{X, x} \to \mathcal{O}_{X_K, x_K}$
est plat (lemme \ref{varieties-lemma-change-fields-flat}). Par conséquent, on a
$$
\text{profondeur}(\mathcal{F}_{x_K}) =
\text{profondeur}(\mathcal{F}_x \otimes_{\mathcal{O}_{X, x}}
\mathcal{O}_{X_K, x_K}) \geq
\text{profondeur}(\mathcal{F}_x)
$$
d'après Algèbre, lemme \ref{algebra-lemma-apply-grothendieck-module}
(en fait, on a égalité ici, mais nous n'en avons pas besoin).
Le résultat sur $k$ découle donc
du résultat sur le corps infini $K$, ce qui achève la démonstration.
\end{proof}

\begin{lemma}
\label{varieties-lemma-connectedness-ample-divisor}
Soit $k$ un corps. Soit $X$ un schéma propre sur $k$. Soit $\mathcal{L}$
un $\mathcal{O}_X$-module inversible ample. Soit
$s \in \Gamma(X, \mathcal{L})$. Supposons que
\begin{enumerate}
\item $s$ soit une section régulière
(Diviseurs, définition \ref{divisors-definition-regular-section}),
\item pour tout point fermé $x \in X$, on ait
$\text{profondeur}(\mathcal{O}_{X, x}) \geq 2$, et
\item $X$ soit connexe.
\end{enumerate}
Alors le schéma des zéros $Z(s)$ de $s$ est connexe.
\end{lemma}

\begin{proof}
Puisque $s$ est une section régulière, il en est de même de
$s^n \in \Gamma(X, \mathcal{L}^{\otimes n})$ pour tout $n > 1$.
De plus, le morphisme d'inclusion $Z(s) \to Z(s^n)$ induit une bijection
sur les espaces topologiques sous-jacents. Par conséquent, si $Z(s)$ est non connexe,
$Z(s^n)$ l'est aussi. Considérons maintenant la suite exacte courte canonique
$$
0 \to \mathcal{L}^{\otimes -n} \xrightarrow{s^n}
\mathcal{O}_X \to \mathcal{O}_{Z(s^n)} \to 0
$$
Considérons la $k$-algèbre $R_n = \Gamma(X, \mathcal{O}_{Z(s^n)})$.
Si $Z(s)$ est non connexe, c'est-à-dire si $Z(s^n)$ est non connexe,
alors soit $R_n$ est nul lorsque $Z(s^n) = \emptyset$, soit
$R_n$ contient un idempotent non trivial lorsque $Z(s^n) = U \amalg V$, avec
$U, V \subset Z(s^n)$ ouverts et non vides (le lecteur pourra consulter le
lemme \ref{varieties-lemma-proper-geometrically-reduced-global-sections}).
Ainsi, l'homomorphisme $\Gamma(X, \mathcal{O}_X) \to R_n$ ne peut être un isomorphisme.
Il en résulte que l'un des deux groupes $H^0(X, \mathcal{L}^{\otimes -n})$ et
$H^1(X, \mathcal{L}^{\otimes -n})$ est non nul pour une infinité
d'entiers positifs $n$. Cela contredit le lemme \ref{varieties-lemma-vanishin-h0-negative} ou
\ref{varieties-lemma-vanishin-h1-negative},
ce qui achève la démonstration.
\end{proof}











\bibliography{my}
\bibliographystyle{amsalpha}
\end{document}
